\documentclass[12pt,a4paper,oneside]{book}

% Metadades pròpies del llibre: una única font de veritat.
% Dades pròpies del llibre. Es carrega abans del preàmbul compartit.
% Una sola font de veritat per a la portada, la pàgina de crèdits i les
% metadades del PDF. Actualitzeu versió i data a cada versió publicada.
\newcommand{\llibretitol}{Una introducció a la teoria de categories}
% El subtítol només surt a la portada; a la resta de llocs n'hi ha prou amb el títol.
\newcommand{\llibresubtitol}{Formació prèvia per aprendre lògica categòrica}
\newcommand{\llibreautor}{Miquel Àngel Perelló}
\newcommand{\llibreversio}{1.7}
% Els enllaços ja inclouen \url (no s'han d'embolcallar de nou).
% Data editorial fixa (reproduïble), en lloc de \today.
\newcommand{\llibredata}{Octubre de 2026}
\newcommand{\llibreany}{2026}
\newcommand{\llibreweb}{\url{https://aprendes.com/logcat/}}
\newcommand{\llibrerepositori}{\url{https://github.com/maperello/aprendes}}
\newcommand{\llibreerrates}{\url{https://github.com/maperello/aprendes/issues}}

\title{\llibretitol\\[0.8em]\large\llibresubtitol}
\author{\llibreautor}
\date{\llibredata}

% S'ha de cridar després de carregar hyperref (és a preamble.tex).
\newcommand{\hypermetadata}{%
  \hypersetup{%
    pdftitle={\llibretitol},
    pdfauthor={\llibreautor},
    pdfsubject={Introducció a la teoria de categories orientada cap a la lògica categòrica. Versió \llibreversio\ (\llibredata). Llicència CC BY-SA 4.0},
    pdfkeywords={teoria de categories, lògica categòrica, topos, Yoneda, adjuncions}%
  }%
}


% ============================================================
% INTERRUPTOR DE COMPILACIÓ
% ------------------------------------------------------------
% Per defecte es genera el LLIBRE COMPLET (capítols 1-10). El mode fascicle
% (només els capítols 1-3 en les quatre parts, amb apèndixs, glossari,
% bibliografia, taula de notació i índex) s'activa de qualsevol d'aquestes
% dues maneres:
%
%   (a) definint \fascicle a la línia d'ordres --- ho fa el Makefile
%       (make fascicle):
%           xelatex "\def\fascicle{}\documentclass[12pt,a4paper,oneside]{book}

% Metadades pròpies del llibre: una única font de veritat.
% Dades pròpies del llibre. Es carrega abans del preàmbul compartit.
% Una sola font de veritat per a la portada, la pàgina de crèdits i les
% metadades del PDF. Actualitzeu versió i data a cada versió publicada.
\newcommand{\llibretitol}{Una introducció a la teoria de categories}
% El subtítol només surt a la portada; a la resta de llocs n'hi ha prou amb el títol.
\newcommand{\llibresubtitol}{Formació prèvia per aprendre lògica categòrica}
\newcommand{\llibreautor}{Miquel Àngel Perelló}
\newcommand{\llibreversio}{1.7}
% Els enllaços ja inclouen \url (no s'han d'embolcallar de nou).
% Data editorial fixa (reproduïble), en lloc de \today.
\newcommand{\llibredata}{Octubre de 2026}
\newcommand{\llibreany}{2026}
\newcommand{\llibreweb}{\url{https://aprendes.com/logcat/}}
\newcommand{\llibrerepositori}{\url{https://github.com/maperello/aprendes}}
\newcommand{\llibreerrates}{\url{https://github.com/maperello/aprendes/issues}}

\title{\llibretitol\\[0.8em]\large\llibresubtitol}
\author{\llibreautor}
\date{\llibredata}

% S'ha de cridar després de carregar hyperref (és a preamble.tex).
\newcommand{\hypermetadata}{%
  \hypersetup{%
    pdftitle={\llibretitol},
    pdfauthor={\llibreautor},
    pdfsubject={Introducció a la teoria de categories orientada cap a la lògica categòrica. Versió \llibreversio\ (\llibredata). Llicència CC BY-SA 4.0},
    pdfkeywords={teoria de categories, lògica categòrica, topos, Yoneda, adjuncions}%
  }%
}


% ============================================================
% INTERRUPTOR DE COMPILACIÓ
% ------------------------------------------------------------
% Per defecte es genera el LLIBRE COMPLET (capítols 1-10). El mode fascicle
% (només els capítols 1-3 en les quatre parts, amb apèndixs, glossari,
% bibliografia, taula de notació i índex) s'activa de qualsevol d'aquestes
% dues maneres:
%
%   (a) definint \fascicle a la línia d'ordres --- ho fa el Makefile
%       (make fascicle):
%           xelatex "\def\fascicle{}\documentclass[12pt,a4paper,oneside]{book}

% Metadades pròpies del llibre: una única font de veritat.
% Dades pròpies del llibre. Es carrega abans del preàmbul compartit.
% Una sola font de veritat per a la portada, la pàgina de crèdits i les
% metadades del PDF. Actualitzeu versió i data a cada versió publicada.
\newcommand{\llibretitol}{Una introducció a la teoria de categories}
% El subtítol només surt a la portada; a la resta de llocs n'hi ha prou amb el títol.
\newcommand{\llibresubtitol}{Formació prèvia per aprendre lògica categòrica}
\newcommand{\llibreautor}{Miquel Àngel Perelló}
\newcommand{\llibreversio}{1.7}
% Els enllaços ja inclouen \url (no s'han d'embolcallar de nou).
% Data editorial fixa (reproduïble), en lloc de \today.
\newcommand{\llibredata}{Octubre de 2026}
\newcommand{\llibreany}{2026}
\newcommand{\llibreweb}{\url{https://aprendes.com/logcat/}}
\newcommand{\llibrerepositori}{\url{https://github.com/maperello/aprendes}}
\newcommand{\llibreerrates}{\url{https://github.com/maperello/aprendes/issues}}

\title{\llibretitol\\[0.8em]\large\llibresubtitol}
\author{\llibreautor}
\date{\llibredata}

% S'ha de cridar després de carregar hyperref (és a preamble.tex).
\newcommand{\hypermetadata}{%
  \hypersetup{%
    pdftitle={\llibretitol},
    pdfauthor={\llibreautor},
    pdfsubject={Introducció a la teoria de categories orientada cap a la lògica categòrica. Versió \llibreversio\ (\llibredata). Llicència CC BY-SA 4.0},
    pdfkeywords={teoria de categories, lògica categòrica, topos, Yoneda, adjuncions}%
  }%
}


% ============================================================
% INTERRUPTOR DE COMPILACIÓ
% ------------------------------------------------------------
% Per defecte es genera el LLIBRE COMPLET (capítols 1-10). El mode fascicle
% (només els capítols 1-3 en les quatre parts, amb apèndixs, glossari,
% bibliografia, taula de notació i índex) s'activa de qualsevol d'aquestes
% dues maneres:
%
%   (a) definint \fascicle a la línia d'ordres --- ho fa el Makefile
%       (make fascicle):
%           xelatex "\def\fascicle{}\documentclass[12pt,a4paper,oneside]{book}

% Metadades pròpies del llibre: una única font de veritat.
\input{metadata.tex}

% ============================================================
% INTERRUPTOR DE COMPILACIÓ
% ------------------------------------------------------------
% Per defecte es genera el LLIBRE COMPLET (capítols 1-10). El mode fascicle
% (només els capítols 1-3 en les quatre parts, amb apèndixs, glossari,
% bibliografia, taula de notació i índex) s'activa de qualsevol d'aquestes
% dues maneres:
%
%   (a) definint \fascicle a la línia d'ordres --- ho fa el Makefile
%       (make fascicle):
%           xelatex "\def\fascicle{}\input{main.tex}"
%
%   (b) descomentant la línia \publicacioparcialtrue de més avall.
%
% Recordeu recompilar amb XeLaTeX + makeindex + XeLaTeX x2.
% ============================================================
\newif\ifpublicacioparcial
\ifdefined\fascicle\publicacioparcialtrue\else\publicacioparcialfalse\fi
% \publicacioparcialtrue % <- descomenteu per forçar el fascicle sense \fascicle

% Mode de sortida. Per defecte: pantalla/interactiu. El Makefile defineix
% \imprimible per generar una còpia pensada per imprimir: fons blanc,
% enllaços discrets i tests estàtics sense camps AcroForm.
\newif\ifimprimible
\ifdefined\imprimible\imprimibletrue\else\imprimiblefalse\fi
% Nom de la Part IV segons el mode: en paper no hi ha interactivitat.
\ifimprimible
  \newcommand{\NomTests}{Tests d'autoavaluació}
  \newcommand{\nomtests}{tests d'autoavaluació}
\else
  \newcommand{\NomTests}{Tests interactius}
  \newcommand{\nomtests}{tests interactius}
\fi

% Els fitxers compartits viuen ara a la subcarpeta common/ d'aquest mateix
% projecte, de manera que es troben amb ruta explícita, sense dependre de
% TEXINPUTS ni de cap configuració del sistema de compilació.
\usepackage[sobri]{common/aprendes}
\input{common/preamble.tex}
\input{common/macros.tex}
\input{caps/categories-macros.tex}
\input{caps/tikz-estils.tex}
\input{common/interactive-tests.tex}
\input{caps/tests-estatics.tex}
\input{caps/didactica.tex}

\makeindex[title={Índex alfabètic},columns=2,intoc]

\begin{document}
\ifimprimible\pagecolor{white}\else\pagecolor{fons!20}\fi
\hypermetadata

\frontmatter
\maketitle
\input{caps/credits}
\tableofcontents
\let\cleardoublepage\clearpage

\input{caps/prefaci}

\mainmatter
\input{caps/introduccio}
\input{caps/guia}

% ============================================================
% PART I — TEORIA
% ============================================================
\part{Teoria}
\setcounter{chapter}{0}
\input{caps/teoria/categories}
\input{caps/teoria/functors}
\input{caps/teoria/transformacions}
\ifpublicacioparcial\else
\input{caps/teoria/yoneda}
\input{caps/teoria/limits}
\input{caps/teoria/adjuncions}
\input{caps/teoria/cartesianes}
\input{caps/teoria/subobjectes}
\input{caps/teoria/topos}
\input{caps/teoria/pont}
\fi

% ============================================================
% PART II — PRÀCTICA
% Els capítols reprodueixen els noms i l'ordre de la part teòrica.
% ============================================================
\part{Pràctica}
\setcounter{chapter}{0}
\input{caps/practica/categories}
\input{caps/practica/functors}
\input{caps/practica/transformacions}
\ifpublicacioparcial\else
\input{caps/practica/yoneda}
\input{caps/practica/limits}
\input{caps/practica/adjuncions}
\input{caps/practica/cartesianes}
\input{caps/practica/subobjectes}
\input{caps/practica/topos}
\fi

% ============================================================
% PART III — EXERCICIS PROPOSATS
% ============================================================
\part{Exercicis proposats}
\setcounter{chapter}{0}
\input{caps/exercicis/categories}
\input{caps/exercicis/functors}
\input{caps/exercicis/transformacions}
\ifpublicacioparcial\else
\input{caps/exercicis/yoneda}
\input{caps/exercicis/limits}
\input{caps/exercicis/adjuncions}
\input{caps/exercicis/cartesianes}
\input{caps/exercicis/subobjectes}
\input{caps/exercicis/topos}
\fi

% ============================================================
% PART IV — TESTS INTERACTIUS
% ============================================================
\part{\NomTests}
\setcounter{chapter}{0}
% En pantalla, un sol entorn Form per a tots els tests: un únic diccionari
% AcroForm. En mode imprimible les macros dels tests són estàtiques i no es
% crea cap formulari.
\ifimprimible\else\begin{Form}\fi
\input{caps/tests/categories}
\input{caps/tests/functors}
\input{caps/tests/transformacions}
\ifpublicacioparcial\else
\input{caps/tests/yoneda}
\input{caps/tests/limits}
\input{caps/tests/adjuncions}
\input{caps/tests/cartesianes}
\input{caps/tests/subobjectes}
\input{caps/tests/topos}
\fi
\ifimprimible\else\end{Form}\fi

% ============================================================
% APÈNDIXS NUMERATS
% \appendix ha d'anar dins \mainmatter (abans de \backmatter):
% a la classe book, \backmatter desactiva la numeració de capítols,
% de manera que els capítols d'apèndix es numeren A, B, ... aquí.
% ============================================================
\appendix
\part{Apèndixs}
\input{caps/apendix-axiomatica-relacional}
\input{caps/apendix-grafs}
\input{caps/apendix-mida}
\input{caps/solucions-tests}
\ifpublicacioparcial\else
\input{caps/apendix-contraexemples}
\fi

\backmatter
\input{caps/bibliografia}
\input{caps/glossari}
\input{caps/notacio-mida}
\printindex

\end{document}
"
%
%   (b) descomentant la línia \publicacioparcialtrue de més avall.
%
% Recordeu recompilar amb XeLaTeX + makeindex + XeLaTeX x2.
% ============================================================
\newif\ifpublicacioparcial
\ifdefined\fascicle\publicacioparcialtrue\else\publicacioparcialfalse\fi
% \publicacioparcialtrue % <- descomenteu per forçar el fascicle sense \fascicle

% Mode de sortida. Per defecte: pantalla/interactiu. El Makefile defineix
% \imprimible per generar una còpia pensada per imprimir: fons blanc,
% enllaços discrets i tests estàtics sense camps AcroForm.
\newif\ifimprimible
\ifdefined\imprimible\imprimibletrue\else\imprimiblefalse\fi
% Nom de la Part IV segons el mode: en paper no hi ha interactivitat.
\ifimprimible
  \newcommand{\NomTests}{Tests d'autoavaluació}
  \newcommand{\nomtests}{tests d'autoavaluació}
\else
  \newcommand{\NomTests}{Tests interactius}
  \newcommand{\nomtests}{tests interactius}
\fi

% Els fitxers compartits viuen ara a la subcarpeta common/ d'aquest mateix
% projecte, de manera que es troben amb ruta explícita, sense dependre de
% TEXINPUTS ni de cap configuració del sistema de compilació.
\usepackage[sobri]{common/aprendes}
% ============================================================
%  Preàmbul compartit — Aprendes Universitat
%  Compartit pels llibres d'Aprendes Universitat.
%  Les dades pròpies de cada llibre (títol, autor i metadades PDF)
%  van al metadata.tex local de cada projecte, carregat ABANS
%  d'aquest fitxer. La macro \hypermetadata es crida DESPRÉS,
%  quan hyperref ja ha estat carregat.
% ============================================================

\usepackage{eurosym}
\usepackage{iftex}
\ifpdftex
  \usepackage[T1]{fontenc}
  \usepackage[utf8]{inputenc}
  \usepackage{lmodern}
\else
  % XeLaTeX / LuaLaTeX: Unicode natiu, res d'inputenc
  \usepackage{fontspec}
  % Carrega Latin Modern per nom de fitxer, que XeTeX/LuaTeX pot
  % resoldre via kpathsea fins i tot si fontconfig no l'ha registrat.
  \setmainfont{lmroman10-regular.otf}[
    BoldFont=lmroman10-bold.otf,
    ItalicFont=lmroman10-italic.otf,
    BoldItalicFont=lmroman10-bolditalic.otf
  ]
  \setsansfont{lmsans10-regular.otf}[
    BoldFont=lmsans10-bold.otf,
    ItalicFont=lmsans10-oblique.otf,
    BoldItalicFont=lmsans10-boldoblique.otf
  ]
  \setmonofont{lmmono10-regular.otf}[ItalicFont=lmmono10-italic.otf]
\fi
% microtype: protrusion i expansió completes només amb pdfLaTeX.
% Amb XeLaTeX/LuaLaTeX només hi ha protrusion i generaria avisos de
% «slot number»; per això només l'activem sota pdfTeX.
\ifpdftex
  \usepackage{microtype}
\fi
\usepackage[catalan, provide=*]{babel}
\usepackage[autostyle=true,babel=true]{csquotes}
% csquotes no porta estil català: el declarem perquè \enquote{...}
% doni guillemets «...» (exteriors) i cometes altes “...” (interiors),
% independentment de la versió de csquotes instal·lada.
\DeclareQuoteStyle{catalan}
  {\guillemotleft}{\guillemotright}
  {\quotedblbase}{\textquotedblleft}
\DeclareQuoteAlias{catalan}{catala}
\usepackage{amsmath,mathtools}% amssymb, amsthm: via aprendes.sty
\usepackage{geometry}
\usepackage{import}
\usepackage{graphicx}
\usepackage{tabularx}
\usepackage{xltabular}% tabularx que pagina (taula de notació)
\usepackage{tikz}
\usetikzlibrary{positioning,arrows.meta,matrix}
\usepackage[all]{xy}% diagrames de grafs (llaços i arestes múltiples) a l'apèndix de grafs
\usepackage{float}
\usepackage[section]{placeins}% FloatBarrier automàtic: cap figura no travessa una secció
\usepackage{pgfplots}
\pgfplotsset{compat=1.18}
\usepackage{imakeidx}
\usepackage{booktabs}

% La bibliografia es configura localment en cada projecte quan hi ha referències.

\setcounter{MaxMatrixCols}{10}

\graphicspath{{img/}}
\geometry{hmargin={3cm,3cm},vmargin={2.5cm,2.5cm},offset=0pt,footskip=1cm,headsep=1cm}


\setlist{itemsep=0.25em, topsep=0.25em}
\numberwithin{equation}{section}
\setcounter{secnumdepth}{3}
\setcounter{tocdepth}{2}

% Configuració dels paràgrafs
\usepackage{indentfirst}
\setlength{\parindent}{1.5em}
\setlength{\parskip}{0.3em}

% ============================================================
%  hyperref: ÚLTIM paquet, com recomana el seu manual.
%  Els colors dels enllaços es fixen aquí. Les metadades
%  pròpies del llibre ja s'han definit abans en metadata.tex;
%  la macro \hypermetadata s'executa després de carregar aquest fitxer.
% ============================================================
\usepackage[pdflang=ca-ES,pdfdisplaydoctitle=true]{hyperref}
% Color dels enllaços: el blau pissarra de la paleta d'aprendes,
% una mica enfosquit. Prou distint del vermell corporatiu dels
% títols, exemples i exercicis, i coherent amb la resta de colors.
\colorlet{EnllacBlau}{ResultatBlau!90!black}
\ifimprimible
  \hypersetup{
    colorlinks=true,
    linktoc=page,
    linkcolor=black,
    urlcolor=black,
    citecolor=black,
    filecolor=black,
    pdfborder={0 0 0},
    hypertexnames=false,
  }
\else
  \hypersetup{
    colorlinks=true,
    linktoc=page,
    linkcolor=EnllacBlau,
    urlcolor=EnllacBlau,
    citecolor=EnllacBlau,
    filecolor=EnllacBlau,
    hypertexnames=false,
  }
\fi
\usepackage{bookmark}% millora els marcadors del PDF; ha d'anar després d'hyperref

% Macros matemàtiques i auxiliars compartides
\newcommand{\Qcb}[1]{#1}
\newcommand{\func}[1]{\operatorname{#1}}
\newcommand{\dbigcup}{\displaystyle\bigcup}
\newcommand{\dbigcap}{\displaystyle\bigcap}
\newcommand{\FRAME}[8]{%
\begin{center}
#5\par\smallskip
\includegraphics[width=#2,height=#3,keepaspectratio]{img/#7}
\end{center}}

\newcommand{\R}{\mathbb{R}}
\newcommand{\N}{\mathbb{N}}
\newcommand{\Z}{\mathbb{Z}}
\newcommand{\Q}{\mathbb{Q}}
\newcommand{\Natpos}{\mathbb{N}^{*}}
\newcommand{\Zstar}{\mathbb{Z}^{*}}
\newcommand{\mcd}{\operatorname{mcd}}
\newcommand{\mcm}{\operatorname{mcm}}
\newcommand{\abs}[1]{\left\lvert #1\right\rvert}
\newcommand{\classe}[1]{\left[#1\right]}
\newcommand{\divides}{\mid}
\newcommand{\notdivides}{\nmid}
\newcommand{\decimalperiodic}[2]{#1,\overline{#2}}
\DeclareMathOperator{\sgn}{sgn} 
% ------------------------------------------------------------
% Macros pròpies de teoria de categories
% ------------------------------------------------------------
\newcommand{\cat}[1]{\mathbf{#1}}          % noms de categories: \cat{Set}
\newcommand{\id}{\operatorname{id}}         % identitat
\providecommand{\Hom}{\operatorname{Hom}}   % conjunts de morfismes
\newcommand{\op}{^{\mathrm{op}}}            % categoria oposada
\newcommand{\Obj}{\operatorname{Obj}}       % col·lecció d'objectes
\newcommand{\comp}{\circ}                    % composició
\DeclareMathOperator{\dom}{dom}
\DeclareMathOperator{\cod}{cod}
\DeclareMathOperator{\Ob}{Ob}
\DeclareMathOperator{\Mor}{Mor}
\DeclareMathOperator{\Aut}{Aut}
\DeclareMathOperator{\Free}{Free}
\DeclareMathOperator{\Nat}{Nat}       % transformacions naturals com a homconjunt
\DeclareMathOperator{\Sub}{Sub}       % preordre de subobjectes
\DeclareMathOperator{\ev}{ev}         % morfisme d'avaluació (exponencials)
\DeclareMathOperator{\Imatge}{Im}     % imatge d'un morfisme (factorització imatge)
% Claudàtors semàntics [[-]] sense dependre de stmaryrd:
\providecommand{\llbracket}{\mathopen{[\![}}
\providecommand{\rrbracket}{\mathclose{]\!]}}
\newcommand{\yon}{\mathrm{y}}          % immersió de Yoneda
\DeclareMathOperator*{\colimop}{colim} % colímit (amb límits sota)

% ============================================================
%  Macros pròpies de la teoria de categories, locals al volum.
%
%  El fitxer compartit ../common/macros.tex barreja notació
%  genèrica (\R, \N, \mcd, decimals periòdics...) amb la de
%  categories. Aquest fitxer recull, en un sol lloc i documentada,
%  només la notació que aquest llibre fa servir, de manera que el
%  volum no depèn de l'organització del fitxer compartit.
%
%  S'usa \providecommand perquè, si macros.tex ja ha definit una
%  macro, no es produeixi cap col·lisió: la definició compartida i
%  la local són idèntiques. Si algun dia es depura macros.tex
%  traient-ne la part de categories, aquest fitxer la garanteix.
% ============================================================

\providecommand{\cat}[1]{\mathbf{#1}}          % noms de categories: \cat{Set}
\providecommand{\id}{\operatorname{id}}         % identitat
\providecommand{\Hom}{\operatorname{Hom}}       % conjunts de morfismes
\providecommand{\op}{^{\mathrm{op}}}            % categoria oposada
\providecommand{\Obj}{\operatorname{Obj}}       % col·lecció d'objectes
\providecommand{\comp}{\circ}                    % composició
\providecommand{\yon}{\mathrm{y}}                % immersió de Yoneda

\providecommand{\dom}{\operatorname{dom}}
\providecommand{\cod}{\operatorname{cod}}

% Predicats relacionals de l'apèndix d'axiomàtica relacional (Dom, Cod, Com,
% Id). S'escriuen amb versaletes matemàtiques per distingir-los clarament dels
% operadors funcionals \dom, \cod: aquí són relacions primitives, no funcions.
\providecommand{\Dom}{\mathsf{Dom}}
\providecommand{\Cod}{\mathsf{Cod}}
\providecommand{\Com}{\mathsf{Com}}
\providecommand{\Id}{\mathsf{Id}}
\providecommand{\Free}{\operatorname{Free}}
\providecommand{\Nat}{\operatorname{Nat}}        % transformacions naturals com a homconjunt
\providecommand{\Sub}{\operatorname{Sub}}        % preordre de subobjectes
\providecommand{\ev}{\operatorname{ev}}          % morfisme d'avaluació (exponencials)
\providecommand{\Imatge}{\operatorname{Im}}      % imatge d'un morfisme

% Claudàtors semàntics [[-]] sense dependre de stmaryrd:
\providecommand{\llbracket}{\mathopen{[\![}}
\providecommand{\rrbracket}{\mathclose{]\!]}}

% colimit amb subíndex a sota (com \lim), només si encara no existeix:
\makeatletter
\@ifundefined{colimop}{\DeclareMathOperator*{\colimop}{colim}}{}
\makeatother

% Remissions a capítols i elements que no són a l'edició parcial (capítols 1--3).
% A l'edició completa fan servir \ref; a la parcial escriuen el número o un text
% alternatiu, per no deixar referències «??».
\newcommand{\refcap}[2]{\ifpublicacioparcial #2\else\ref{#1}\fi}
\newcommand{\refalt}[2]{\ifpublicacioparcial #2\else\ref{#1}\fi}

% ============================================================
%  Estils TikZ comuns per als diagrames del llibre.
%
%  Els diagrames de teoria de categories repeteixen molt codi de
%  puntes de fletxa i de matrius de nodes. Aquests estils permeten
%  escriure els diagrames nous de manera compacta i canviar
%  l'estètica de tots des d'un sol lloc:
%
%     \begin{tikzpicture}
%       \matrix (m) [cat matrix] { A & B \\ C & D \\ };
%       \draw[cat] (m-1-1) -- node[above] {$f$} (m-1-2);
%     \end{tikzpicture}
%
%  Els diagrames antics, escrits amb -{Stealth[scale=1.1]} i amb
%  els paràmetres de matriu explícits, continuen funcionant sense
%  canvis; aquests estils són per a diagrames nous i per a una
%  homogeneïtzació progressiva.
% ============================================================

\tikzset{
  % Fletxa estàndard d'un morfisme.
  cat/.style={-{Stealth[scale=1.1]}},
  % Fletxa de monomorfisme (cua d'ham).
  cat mono/.style={{Hooks[right]}-{Stealth[scale=1.1]}},
  % Fletxa d'epimorfisme (doble punta).
  cat epi/.style={-{Stealth[scale=1.1]Stealth[scale=1.1]}},
  % Fletxa de transformació natural (doble).
  cat nat/.style={-{Implies},double,double distance=1.5pt},
  % Matriu de nodes amb la separació habitual del llibre.
  cat matrix/.style={matrix of math nodes, row sep=2.6em, column sep=3.2em},
}

% Tests interactius en PDF
\newcounter{testquestion}
\newcommand{\testprefix}{}
\newcommand{\testkeys}{}
\newcommand{\iniciTest}[2]{%
  \renewcommand{\testprefix}{#1}%
  \renewcommand{\testkeys}{#2}%
  \setcounter{testquestion}{0}%
  \gdef\testjustificacions{}%
  \label{test:#1}%
  \begin{center}
  \fbox{\begin{minipage}{0.9\textwidth}
  \ifimprimible
    \textbf{Instruccions.} Marca una sola resposta per pregunta. Comprova les respostes i les justificacions a l'apèndix de solucions.
  \else
    \textbf{Instruccions.} Marca una sola resposta per pregunta. En acabar, prem el bot\'{o} \textbf{Corregir}. El camp \textbf{Nota} mostrar\`{a} el resultat.\par\medskip
    \PushButton[name=#1corregir,width=3.2cm,height=0.8cm,bordercolor={0 0 1},onclick={var p='#1';var k='#2'.split(',');var opts=['a','b','c','d'];var s=0;for(var i=0;i<k.length;i++){var ans='';var n=0;for(var j=0;j<opts.length;j++){var q=this.getField(p+'q'+(i+1)+opts[j]);if(q && q.value!='Off'){ans=opts[j];n++;}}var f=this.getField(p+'f'+(i+1));if(n==1 && ans==k[i]){s++;if(f)f.value='Correcte';}else{if(f)f.value=(n==0?'Sense respondre. Correcta: '+k[i]:(n>1?'Marca una sola resposta. Correcta: '+k[i]:'Incorrecte. Correcta: '+k[i]));}}this.getField(p+'nota').value=s+' / '+k.length;app.alert('Nota: '+s+' / '+k.length);}]{Corregir}\quad
    \PushButton[name=#1netejar,width=3.2cm,height=0.8cm,bordercolor={1 0 0},onclick={var keys='#2'.split(',');var opts=['a','b','c','d'];this.getField('#1nota').value='';for(var i=1;i<=keys.length;i++){var f=this.getField('#1f'+i);if(f)f.value='';for(var j=0;j<opts.length;j++){var q=this.getField('#1q'+i+opts[j]);if(q)q.value='Off';}}}]{Netejar}\quad
    \textbf{Nota: }\TextField[name=#1nota,readonly=true,width=4cm,bordercolor={0 0 0}]{}
  \fi
  \end{minipage}}
  \end{center}
  \bigskip
}
\newcommand{\pregunta}{\stepcounter{testquestion}\item}
% Justificació breu de la resposta correcta d'una pregunta. No es mostra
% al costat de les opcions (per no donar pistes de longitud); s'acumula i
% \clauestatica la imprimeix després de la clau de respostes.
\newcommand{\testjustificacions}{}
\makeatletter
\newcommand{\justificacio}[1]{%
  \edef\@justnum{\thetestquestion}%
  \expandafter\g@addto@macro\expandafter\testjustificacions\expandafter{%
    \expandafter\justitem\expandafter{\@justnum}{#1}}%
}
\makeatother
\newcommand{\justitem}[2]{\item[\textbf{#1.}] #2}
\long\def\opcio#1#2{\item \ifimprimible$\square$\else\CheckBox[name=\testprefix{}q\thetestquestion#1,width=1.1em,height=1.1em,bordercolor={0 0 0},onclick={var opts=['a','b','c','d'];for(var j=0;j<opts.length;j++){if(opts[j]!='#1'){var q=this.getField('\testprefix{}q\thetestquestion'+opts[j]);if(q)q.value='Off';}}}]{}\fi\quad #2}
\newcommand{\feedback}{\ifimprimible\par\medskip\else\par\smallskip\textbf{Correcci\'{o}: }\TextField[name=\testprefix{}f\thetestquestion,readonly=true,width=9cm,bordercolor={0.6 0.6 0.6}]{}\par\medskip\fi}
% ---- Validació estructural de la clau (C7) ----
% \finalitzaTest compara el nombre de preguntes realment plantejades
% (\thetestquestion) amb el nombre de respostes de la clau, i emet un
% avís de compilació si no coincideixen. Cal cridar-la després de
% \end{Form}. Fem servir \@for, que recorre llistes separades per comes
% de manera segura, sense bucles \ifx artesanals.
\newcounter{testkeycount}
\makeatletter
\newcommand{\finalitzaTest}{%
  \setcounter{testkeycount}{0}%
  \expandafter\@for\expandafter\next\expandafter:\expandafter=\testkeys\do{\stepcounter{testkeycount}}%
  \ifnum\value{testquestion}=\value{testkeycount}\else
    \PackageWarning{interactive-tests}{El test '\testprefix' te \thetestquestion\space preguntes pero la clau en te \thetestkeycount\space respostes. Reviseu la crida a \string\iniciTest.}%
  \fi
}
\makeatother

% ============================================================
%  Claus estàtiques dels tests, reunides al final del llibre.
%
%  Els tests del llibre es corregeixen amb JavaScript d'AcroForm,
%  que no funciona en molts visors integrats, navegadors ni
%  aplicacions mòbils. Per a aquests casos, i per a l'ús en paper,
%  cada test té una clau de respostes i unes justificacions. Perquè
%  no apareguin just després de les preguntes, cosa que restaria
%  valor a l'autoavaluació, no s'imprimeixen al lloc del test:
%
%   - \clauestatica (cridada just abans de \finalitzaTest) desa la
%     clau i les justificacions del test en curs a \totesclaus i
%     imprimeix només una remissió a l'apèndix de solucions;
%   - \imprimeixsolucionstests, a caps/solucions-tests.tex, les
%     imprimeix totes al final del llibre, cadascuna amb un enllaç de
%     tornada al test.
% ============================================================

\makeatletter
\newcounter{clauestindex}
\gdef\totesclaus{}
\newcommand{\clauestatica}{%
  \begingroup
    \edef\@tmpclau{\noexpand\blocsoluciotest{\testprefix}{\thechapter}{\testkeys}{\unexpanded\expandafter{\testjustificacions}}}%
    \expandafter\g@addto@macro\expandafter\totesclaus\expandafter{\@tmpclau}%
  \endgroup
  \par\medskip
  \begin{center}
  \small\itshape
  \ifimprimible
  Les respostes i les justificacions d'aquest test són a l'apèndix \hyperref[cap:solucions-tests]{\textup{\enquote{Solucions dels tests}}}, pàgina~\pageref{sol:\testprefix}. Consulteu-les després d'haver respost tot el test.
  \else
  Si el vostre lector de PDF no admet el botó \textup{\textbf{Corregir}}, trobareu les respostes i les justificacions d'aquest test a l'apèndix \hyperref[cap:solucions-tests]{\textup{\enquote{Solucions dels tests}}}, pàgina~\pageref{sol:\testprefix}.
  \fi
  \end{center}
  \par\medskip
}
% #1 prefix del test, #2 número de capítol, #3 clau, #4 justificacions
\newcommand{\blocsoluciotest}[4]{%
  \section*{Test del capítol #2}\label{sol:#1}%
  \addcontentsline{toc}{section}{Test del capítol #2}%
  \noindent\textbf{Respostes.}\quad
  \setcounter{clauestindex}{0}%
  \@for\next:=#3\do{%
    \stepcounter{clauestindex}%
    \mbox{\textbf{\theclauestindex.}~\next}\hspace{1.2em plus 0.4em}%
  }%
  \par
  \def\@tmpjust{#4}%
  \ifx\@tmpjust\@empty\else
    \medskip
    \noindent\textbf{Justificacions.}
    \begin{list}{}{\setlength{\leftmargin}{2.2em}\setlength{\labelwidth}{1.8em}\setlength{\itemsep}{0.2em}\small}
    #4
    \end{list}
  \fi
  \par\smallskip
  \noindent{\small\hyperref[test:#1]{Tornar al test} (pàgina~\pageref{test:#1}).}
  \par\bigskip
}
\newcommand{\imprimeixsolucionstests}{\totesclaus}
\makeatother

% ============================================================
%  Entorns didàctics locals del volum de categories.
%  Definits aquí (no a l'estil compartit) per no acoblar-hi
%  el projecte. Reutilitzen tcolorbox, etoolbox i la paleta
%  ja carregats per aprendes.sty en mode «sobri».
%
%    \begin{objectius} ... \end{objectius}   -> bloc d'obertura
%    \begin{resumcap}   ... \end{resumcap}   -> bloc de tancament
%
%  Dins d'aquests blocs s'hi poden usar els encapçalaments
%  \apartat{...} per separar prerequisits/objectius/etc.
% ============================================================

% Color propi per als blocs didàctics: verd pissarra derivat de la
% paleta existent (cexer = verd), enfosquit per contrastar amb el fons.
\colorlet{DidacticVerd}{cexer!55!black}

% Estils de bloc: mateixa gramàtica visual que blocenunciat/blocobservacio.
\tcbset{
  blocdidactic/.style={
    enhanced, breakable, sharp corners,
    frame hidden, boxrule=0pt,
    colback=DidacticVerd!6,
    borderline west={2.0pt}{0pt}{DidacticVerd},
    left=8pt, right=4pt, top=5pt, bottom=5pt,
    before skip=1.0em, after skip=1.0em,
  },
}

% Encapçalament intern d'un apartat del bloc (Prerequisits, Objectius...).
\newcommand{\apartat}[1]{\par\smallskip{\color{DidacticVerd}\ifpdftex\scshape\else\sffamily\bfseries\fi #1}\par\nobreak\smallskip}

% Bloc d'obertura del capítol.
\newenvironment{objectius}{%
  \begin{tcolorbox}[blocdidactic]%
  \ignorespaces
}{\end{tcolorbox}}

% Bloc de tancament del capítol.
\newenvironment{resumcap}{%
  \begin{tcolorbox}[blocdidactic]%
  {\color{DidacticVerd}\ifpdftex\scshape\else\sffamily\bfseries\fi S\'{\i}ntesi del cap\'{\i}tol.}\hspace{0.6em}%
  \ignorespaces
}{\end{tcolorbox}}


\makeindex[title={Índex alfabètic},columns=2,intoc]

\begin{document}
\ifimprimible\pagecolor{white}\else\pagecolor{fons!20}\fi
\hypermetadata

\frontmatter
\maketitle
% Pàgina de crèdits (revers de la portada).
\thispagestyle{empty}
\vspace*{\fill}
\begin{flushleft}
\small
\textbf{\llibretitol}\\
\llibreautor\\[0.6em]
Versió \llibreversio\ · \llibredata\\
\llibreweb\\[1.2em]

\textbf{Llicència.} © \llibreany\ \llibreautor. Aquesta obra està subjecta a la llicència Creative Commons Reconeixement-CompartirIgual 4.0 Internacional (CC BY-SA 4.0): \url{https://creativecommons.org/licenses/by-sa/4.0/deed.ca}. Podeu copiar-la, redistribuir-la i adaptar-la, també amb finalitat comercial, sempre que en reconegueu l'autoria i distribuïu les obres derivades amb la mateixa llicència.\\[1em]

\textbf{Com citar-lo.} \llibreautor, \emph{\llibretitol}, versió \llibreversio, \llibreany. Projecte aprendes\,/\,logcat, \llibreweb.\\[1em]

\textbf{Errates i suggeriments.} Es poden comunicar obrint una incidència (\emph{issue}) al repositori del projecte, \llibreerrates, indicant la versió del llibre i la pàgina. Les correccions s'incorporen a les versions següents.
\end{flushleft}
\clearpage

\tableofcontents
\let\cleardoublepage\clearpage

\chapter*{Prefaci}
\addcontentsline{toc}{chapter}{Prefaci}

Aquest llibre forma part d'\textbf{aprendes}, una col·lecció de materials
d'estudi que comparteixen una mateixa manera de treballar: cada tema es
presenta en una part de \emph{teoria}, es reprèn en una part de
\emph{pràctica} amb problemes resolts, es consolida amb \emph{exercicis
proposats} i es comprova amb \emph{\nomtests}. Qui ja conegui
altres volums de la col·lecció hi reconeixerà l'estructura; qui s'hi
acosti per primera vegada la trobarà explicada a la guia de lectura.

El volum té un propòsit precís. És el primer d'un projecte en dues
etapes: aquí es construeix el llenguatge de la teoria de categories
---des de les definicions bàsiques fins als classificadors de subobjectes
i els topos--- amb el detall que necessitarà un volum posterior dedicat
a la lògica categòrica. Per això la lògica hi és present des del
començament, com a font d'exemples i com a direcció del recorregut,
tot i que el seu desenvolupament sistemàtic queda per a aquella segona
etapa.

\ifpublicacioparcial
Aquesta és una edició parcial, que conté els tres primers temes
---categories, functors i transformacions naturals--- en les quatre parts,
juntament amb els apèndixs i els materials de consulta. Els capítols
restants s'incorporaran en l'edició completa.
\else
Aquesta edició conté el volum complet: nou temes desenvolupats en les
quatre parts i un capítol final que fa de pont cap a la lògica categòrica.
\fi

Abans del primer capítol, la \emph{Introducció} explica per què val la
pena aprendre aquest llenguatge i la \emph{Guia de lectura} proposa com
fer-ho. Es poden llegir en aquest ordre o tornar-hi quan convingui.

\medskip
\noindent\textbf{Prerequisits.} El text pressuposa la maduresa matemàtica habitual d'un primer curs universitari de matemàtiques, informàtica o lògica: treballar amb conjunts i aplicacions, seguir una demostració elemental i llegir notació algebraica bàsica. No es pressuposa haver estudiat teoria de categories, topologia, àlgebra abstracta avançada ni lògica categòrica. Quan un exemple necessita vocabulari addicional, s'introdueix localment o es remet a un apèndix.

\medskip
\noindent\textbf{Errates i suggeriments.} Són benvinguts. Es poden comunicar obrint una incidència al repositori del projecte, \llibreerrates, indicant la versió (\llibreversio) i la pàgina.

% ------------------------------------------------------------
% Espai reservat per a l'autor (descomenteu i completeu):
%
% \medskip
% \noindent\textbf{Agraïments.} ...
%
% (Les errates i els suggeriments: vegeu la pàgina de crèdits i el text següent.)
% ------------------------------------------------------------


\mainmatter
\chapter*{\texorpdfstring{Per què aprendre\\ teoria de categories?}{Per què aprendre teoria de categories?}}
\addcontentsline{toc}{chapter}{Introducció}
\markboth{Introducció}{Introducció}

\begingroup
\renewcommand{\thesection}{\arabic{section}}
\setcounter{section}{0}
\setlength{\emergencystretch}{1.5em}

Hi ha construccions matemàtiques que aprenem diverses vegades, amb noms i
representacions diferents. Aparellar dues dades, imposar simultàniament dues
condicions, trobar un màxim comú divisor o introduir una conjunció semblen
operacions allunyades. I, tanmateix, totes poden respondre a una mateixa
pregunta: com podem reunir dues exigències en una sola dada, de manera que
qualsevol altra manera de satisfer-les passi per aquesta?

La teoria de categories ensenya a reconèixer aquestes formes comunes. No ho
fa declarant que els conjunts, els nombres i les proposicions són la mateixa
cosa. Ho fa distingint, en cada context, quines relacions són rellevants i
com es poden compondre. El que unifica no és la naturalesa dels objectes,
sinó el paper que tenen dins d'un sistema de relacions.

Aquest llibre proposa aprendre aquest llenguatge des dels seus fonaments,
amb quatre famílies d'exemples que ens acompanyaran al llarg del recorregut:
els conjunts, els preordres, els grafs i la lògica. La finalitat és doble:
adquirir les eines bàsiques de la teoria de categories i comprendre per què
aquestes eines preparen una lectura estructural de la lògica. Abans de
començar amb les definicions, convé veure què ens permeten preguntar.

\section{Una estructura també es coneix per les seves relacions}
\label{sec:mirar-fora}

Una manera habitual d'estudiar un objecte consisteix a mirar-ne l'interior:
els elements d'un conjunt, l'operació d'un grup o les parts d'un graf. La
perspectiva categòrica no elimina aquesta mirada, però la complementa amb
una altra: quines transformacions podem fer entre aquest objecte i els
altres? Quines hi arriben, quines en surten i com s'encadenen?
Podem resumir aquest canvi amb una consigna que ens acompanyarà al llarg
del llibre: \emph{relaciona, no individualitza}. En lloc de fixar un
objecte per allò que és aïlladament, el situem pel paper que té dins
d'una xarxa de transformacions.

Entre conjunts tenim aplicacions; entre grups, homomorfismes; entre
preordres, aplicacions monòtones. Aquestes transformacions tenen una
propietat comuna. Si podem passar d'$A$ a $B$ i de $B$ a $C$, podem fer els
dos passos seguits:
\[
  A\xrightarrow{f}B\xrightarrow{g}C,
  \qquad g\comp f\colon A\longrightarrow C.
\]
La composició és associativa: quan tres passos es poden encadenar, agrupar
els dos primers o els dos darrers no canvia el resultat. A més, cada objecte
té una transformació identitat, que el deixa com és i no altera cap altra
transformació quan s'hi compon.

Una \emph{categoria}\index{categoria} reté justament aquestes dades: objectes, fletxes entre
objectes, composició i identitats, amb les lleis que acabem de descriure.
La paraula \emph{fletxa}, o \emph{morfisme}, és més general que la paraula
\emph{funció}. En la categoria lliure\index{categoria!lliure} sobre un graf dirigit, els objectes
són els vèrtexs i les fletxes són els camins: compondre és concatenar-los,
i la identitat és el camí de longitud zero. Un monoide també es pot veure
com una categoria: té un sol objecte, els elements del monoide en són les
fletxes i la multiplicació en dona la composició.

Ni tan sols els mateixos objectes fixen per si sols la categoria. Amb
conjunts i aplicacions obtenim $\cat{Set}$; amb conjunts i relacions,
$\cat{Rel}$\index{Rel@$\cat{Rel}$}. En aquest segon cas, compondre $R\subseteq A\times B$ amb
$S\subseteq B\times C$ significa que
\[
  a\,(S\comp R)\,c
  \quad\Longleftrightarrow\quad
  \exists b\in B\,\bigl(aRb\land bSc\bigr).
\]
La fletxa composta registra que hi ha algun terme intermedi que connecta
els dos extrems. Ja trobem aquí una operació lògica, la quantificació
existencial, dins d'una operació sobre fletxes.

Per tant, la primera pregunta categòrica no és només \enquote{quins objectes tenim?},
sinó \enquote{en quina categoria els estudiem?}. Els nombres enters positius,
ordenats per $\leq$ o per divisibilitat, ofereixen dues categories diferents:
en cadascuna, una fletxa expressa una relació diferent. Escollir les fletxes
és escollir quina part de la matemàtica volem fer visible.

\section{Del que un objecte conté al paper que compleix}

Considerem el producte cartesià de dos conjunts, $A\times B$. El podem
descriure enumerant-ne els elements: són els parells $(a,b)$ amb $a\in A$
i $b\in B$. Però també el podem descriure sense començar pels parells.
Hi ha dues projeccions,
\[
  \pi_A\colon A\times B\longrightarrow A,
  \qquad \pi_B\colon A\times B\longrightarrow B.
\]
Si des d'un mateix conjunt $X$ tenim una aplicació $f\colon X\to A$ i una
aplicació $g\colon X\to B$, hi ha exactament una aplicació cap a $A\times B$
que reuneix aquestes dues dades: $x\mapsto(f(x),g(x))$.

El punt important és que aquesta segona descripció només parla de fletxes.
Té sentit, doncs, en una categoria qualsevol. Un \emph{producte}\index{producte} d'$A$ i
$B$ és un objecte $P$ amb projeccions $\pi_A\colon P\to A$ i
$\pi_B\colon P\to B$ tal que, per a tot objecte $X$ i totes les fletxes
$f\colon X\to A$ i $g\colon X\to B$, existeix una única fletxa
$u\colon X\to P$ amb
\[
  \pi_A\comp u=f,
  \qquad \pi_B\comp u=g.
\]
Aquestes equacions es fan visibles en el diagrama
\[
\begin{tikzpicture}[baseline=(m.center)]
  \matrix (m) [matrix of math nodes,column sep=3.4em,row sep=3em] {
       & X & \\
     A & P & B \\
  };
  \draw[-{Stealth}] (m-1-2) -- node[above left] {$f$} (m-2-1);
  \draw[-{Stealth}] (m-1-2) -- node[above right] {$g$} (m-2-3);
  \draw[-{Stealth},dashed] (m-1-2) -- node[right] {$u$} (m-2-2);
  \draw[-{Stealth}] (m-2-2) -- node[below] {$\pi_A$} (m-2-1);
  \draw[-{Stealth}] (m-2-2) -- node[below] {$\pi_B$} (m-2-3);
\end{tikzpicture}
\]
La fletxa discontínua és la que la propietat ens permet construir.
Dir que el diagrama \emph{commuta} és dir que els camins que comparem,
amb els mateixos extrems, donen la mateixa fletxa.

En un preordre, posem una única fletxa $x\to y$ quan $x\leq y$, i cap quan
aquesta relació no es compleix. El producte de $a$ i $b$ és aleshores el seu
ínfim: un objecte per sota de tots dos, i tan gran com sigui possible entre
els que tenen aquesta propietat. La mateixa condició produeix les
construccions següents, quan existeixen:

\begin{center}
\renewcommand{\arraystretch}{1.2}
\begin{tabularx}{\linewidth}{@{}>{\raggedright\arraybackslash}X
                                  >{\raggedright\arraybackslash}p{0.32\linewidth}@{}}
\toprule
\textbf{Context i fletxes} & \textbf{Producte} \\
\midrule
Conjunts amb aplicacions & Producte cartesià \\
Subconjunts d'un conjunt fix amb inclusions & Intersecció \\
Enters positius amb divisibilitat & Màxim comú divisor \\
Fórmules amb deduïbilitat & Conjunció \\
\bottomrule
\end{tabularx}
\end{center}

La lectura lògica és especialment reveladora. Fixem un sistema deductiu
proposicional amb les regles usuals de la conjunció i posem una fletxa
$p\to q$ quan $q$ és deduïble de $p$. Les projeccions del producte són les
deduccions de $p$ i de $q$ a partir de $p\land q$. Si des de $r$ podem
deduir totes dues fórmules, podem deduir-ne la conjunció:
\[
  r\vdash_T p\land q
  \quad\Longleftrightarrow\quad
  \bigl(r\vdash_T p\ \text{i}\ r\vdash_T q\bigr).
\]
La unicitat de la fletxa és automàtica perquè només registrem si una
deducció existeix, no quina prova la realitza. El producte expressa aquí
una regla lògica amb el mateix patró que expressa l'aparellament de dades.

Això és una \emph{propietat universal}\index{propietat universal}: una caracterització mitjançant
l'existència i la unicitat de fletxes que satisfan unes condicions.
La paraula \enquote{universal} indica que la condició val per a qualsevol objecte
$X$ i qualsevol parell de fletxes adequat, no només per a uns exemples
triats. Bona part del llibre consisteix a reconèixer aquest patró en
construccions successivament més riques.

\section{La mateixa estructura no vol dir el mateix objecte}

Una propietat universal no ens obliga a triar una representació concreta.
Podem construir dues realitzacions diferents d'un producte i obtenir dos
objectes que no siguin literalment iguals. La propietat universal, però,
força un \emph{isomorfisme}\index{isomorfisme} entre ells compatible amb les projeccions.

Una fletxa $f\colon A\to B$ és un isomorfisme si hi ha una fletxa
$g\colon B\to A$ que desfà el canvi en tots dos sentits:
\[
  g\comp f=\id_A,
  \qquad f\comp g=\id_B.
\]
Escrivim aleshores $A\cong B$. El sentit d'aquest símbol depèn de la
categoria: a $\cat{Set}$ parlem de bijeccions; entre grups, d'isomorfismes
de grups; entre espais topològics amb aplicacions contínues, d'homeomorfismes.
No n'hi ha prou de poder desfer una aplicació entre conjunts si la inversa
no respecta l'estructura que hem fixat.

En la categoria prima de la deduïbilitat, dues fórmules són isomorfes quan
són interdeduïbles. Així, $p\land q$ i $q\land p$ poden tenir sintaxis
diferents i el mateix paper deductiu. El llenguatge categòric no les
declara iguals: precisa en quin sentit la diferència de presentació no
altera l'estructura que estem estudiant.

Cal conservar també les dades que acompanyen l'objecte. L'isomorfisme
\emph{únic} entre dos productes és el que respecta les projeccions;
l'objecte subjacent, sense aquestes dades, pot tenir altres automorfismes.
Aprendre a distingir l'objecte de la construcció completa és part de
l'aprenentatge de les propietats universals.

Un altre canvi de perspectiva consisteix a invertir sistemàticament les
fletxes. La categoria oposada\index{categoria!oposada} $\cat{C}\op$ té els mateixos
objectes que $\cat{C}$, però una fletxa $A\to B$ de $\cat{C}$ hi apareix
com una fletxa $B\to A$. No afirmem que la fletxa original tingui inversa:
canviem el context en què la llegim.

En invertir les fletxes de la propietat universal del producte obtenim
el \emph{coproducte}. A $\cat{Set}$ és la unió disjunta; entre subconjunts
ordenats per inclusió és la unió; en el context deductiu apropiat és la
disjunció. De manera semblant, el dual d'un objecte terminal, que rep una
única fletxa des de qualsevol objecte, és un objecte inicial, que n'envia
una única a cadascun.

La \emph{dualitat}\index{dualitat} és més que una simetria de dibuixos. Quan un argument
està formulat categòricament, en podem obtenir l'argument dual invertint
també les composicions i les hipòtesis. Aprenem així a buscar, al costat
d'una pregunta, la pregunta que es veu des de l'altra direcció.

\section{Comparar contextos i comparar les comparacions}

Si les categories descriuen contextos, com podem passar d'un context a
un altre? Un \emph{functor}\index{functor} $F\colon\cat{C}\to\cat{D}$ assigna objectes
a objectes i fletxes a fletxes, preservant els extrems, les identitats i
la composició:
\[
  F(\id_A)=\id_{F(A)},
  \qquad F(g\comp f)=F(g)\comp F(f).
\]
Un functor és, en aquest sentit precís, una traducció estructurada.
Per exemple, podem passar d'un grup al seu conjunt subjacent i d'un
homomorfisme a l'aplicació que el realitza. També podem passar d'un graf
a la categoria dels seus camins.

La precisió importa: tot functor preserva les lleis categòriques, però
no tot functor preserva productes, límits o qualsevol altra estructura
addicional. El llibre ensenyarà a preguntar què preserva cada traducció,
què oblida i què permet recuperar.

Hi ha encara un altre nivell de comparació. Donats dos functors
$F,G\colon\cat{C}\to\cat{D}$, una \emph{transformació natural}\index{transformació natural}
$\eta\colon F\Rightarrow G$ dona una fletxa
$\eta_A\colon F(A)\to G(A)$ per a cada objecte $A$, amb la condició
que aquestes fletxes siguin compatibles amb tots els morfismes de
$\cat{C}$. Per a $f\colon A\to B$, això vol dir
\[
  G(f)\comp\eta_A=\eta_B\comp F(f).
\]
No és una col·lecció de comparacions independents: totes s'han d'ajustar
a una mateixa xarxa de relacions.

Un exemple elemental és la diagonal dels conjunts,
$\delta_A\colon A\to A\times A$, definida per $a\mapsto(a,a)$.
Copiar una dada i després aplicar $f$ a les dues còpies dona el mateix
resultat que aplicar $f$ primer i copiar després:
\[
\begin{tikzpicture}[baseline=(m.center)]
  \matrix (m) [matrix of math nodes,column sep=4em,row sep=3em] {
    A & B \\
    A\times A & B\times B \\
  };
  \draw[-{Stealth}] (m-1-1) -- node[above] {$f$} (m-1-2);
  \draw[-{Stealth}] (m-1-1) -- node[left] {$\delta_A$} (m-2-1);
  \draw[-{Stealth}] (m-1-2) -- node[right] {$\delta_B$} (m-2-2);
  \draw[-{Stealth}] (m-2-1) -- node[below] {$f\times f$} (m-2-2);
\end{tikzpicture}
\]
La naturalitat expressa que l'operació de copiar és compatible amb
qualsevol aplicació, no amb una tria particular de noms per als elements.
En un exemple lògic del llibre, un quadrat d'aquesta mena expressarà
que substituir variables i després interpretar una fórmula dona el mateix
resultat que interpretar-la i transportar-ne el significat.

Aquest nivell de coherència permet formular també l'\emph{equivalència de
categories}\index{equivalència de categories}. Dues categories $\cat{C}$ i $\cat{D}$ són equivalents si hi
ha functors $F\colon\cat{C}\to\cat{D}$ i $G\colon\cat{D}\to\cat{C}$
tals que anar i tornar és naturalment isomorf a no haver canviat de context:
\[
  G\comp F\cong\id_{\cat{C}},
  \qquad F\comp G\cong\id_{\cat{D}}.
\]
No demanem una igualtat literal de presentacions. Demanem que el canvi
conservi l'estructura categòrica essencial i que la recuperació sigui
coherent. L'isomorfisme compara objectes dins d'una categoria;
l'equivalència compara categories senceres.

\section{Posar a prova la mirada des de fora}

Aquesta secció avança idees que el llibre introduirà pas a pas en els
capítols centrals. No cal retenir-ne ara els detalls: n'hi ha prou de
veure cap a on apunten i com continuen la perspectiva de les seccions
anteriors.

La idea de conèixer un objecte per les seves relacions es pot formular
amb exactitud. Per a un objecte $A$ d'una categoria localment petita,
$\Hom(X,A)$ és el conjunt de fletxes de $X$ cap a $A$. Podem pensar cada
objecte $X$ com una forma de posar $A$ a prova, i cada fletxa $X\to A$
com una dada d'$A$ observada amb aquella forma.

A $\cat{Set}$, prendre $X$ amb un sol element recupera els elements
ordinaris d'$A$. En altres categories, però, aquesta única mena de prova
pot resultar insuficient. La perspectiva general considera \emph{tots}
els objectes $X$ i, a més, registra com una prova canvia quan precomponem
amb una fletxa $X'\to X$.

Aquesta informació s'organitza en l'homfunctor\index{homfunctor} $\Hom(-,A)$.
El lema de Yoneda\index{Yoneda!lema} donarà forma precisa a una intuïció anunciada des
del començament: els objectes queden determinats, llevat d'isomorfisme,
per aquest sistema de relacions. No n'hi ha prou de saber quantes
fletxes hi arriben; cal conservar la manera com es transformen i es
componen. Dos homfunctors d'aquest tipus naturalment isomorfs provenen
d'objectes isomorfs.

La \emph{representabilitat} aplica aquesta idea a les construccions.
En el cas del producte, donar una fletxa $X\to A\times B$ equival a
donar dues fletxes amb domini $X$:
\[
  \Hom(X,A\times B)
  \cong \Hom(X,A)\times\Hom(X,B),
\]
de manera compatible amb els canvis de $X$. El producte representa,
amb un sol objecte, l'assignació que a cada $X$ associa aquestes parelles.
Límits i colímits reuniran moltes propietats universals sota un mateix
patró; no seran una llista de receptes sense connexió.

Les \emph{adjuncions}\index{adjunció} ens permetran fer un pas més i relacionar
construccions senceres. Si $F\colon\cat{C}\to\cat{D}$ i
$G\colon\cat{D}\to\cat{C}$ són adjunts, amb $F$ per l'esquerra de
$G$, hi ha una correspondència natural
\[
  \Hom_{\cat{D}}(F(A),B)
  \cong \Hom_{\cat{C}}(A,G(B)).
\]
Una mateixa dada es pot descriure a banda i banda del canvi de context.
Per exemple, donar un functor des de la categoria lliure d'un graf cap
a una categoria és equivalent a donar un morfisme del graf inicial cap
al graf subjacent d'aquella categoria. Els camins compostos no s'han
d'escollir de nou: la composició en determina la imatge.

En un altre exemple, donar una aplicació $X\times A\to B$ equival a
donar una aplicació $X\to B^A$, on $B^A$ és el conjunt d'aplicacions
d'$A$ a $B$. És la \emph{currificació}: una funció de dues variables
es pot descriure com una funció que retorna funcions. El tancament
cartesià generalitza aquesta correspondència i ens porta cap a la
semàntica del càlcul lambda simplement tipat.

\section{La lògica com a fil conductor, no com a drecera}

La lògica apareix en aquest llibre des del primer capítol, però a
diferents nivells. El més elemental és la categoria prima de la
deduïbilitat: les fórmules són objectes i una fletxa registra que una
fórmula es pot deduir d'una altra. Aquesta lectura ajuda a reconèixer
la conjunció com a producte i la interdeduïbilitat com a isomorfisme.

Conservar les proves mateixes és una tasca més exigent. Cal decidir
quan dues proves representen la mateixa fletxa i com la composició
respecta aquesta identificació. No podem passar de l'existència d'una
deducció a una categoria de proves sense afegir aquestes precisions.
La interpretació de tipus i termes en categories cartesianes tancades
ens permetrà començar a veure la relació entre proves, programes i
morfismes, sense confondre aquests nivells.

Més endavant, els \emph{subobjectes}\index{subobjecte} oferiran una altra lectura lògica.
A $\cat{Set}$, un predicat sobre $A$ es pot descriure pel subconjunt
d'elements que el satisfan. Si $f\colon X\to A$, substituir $f(x)$
en el predicat correspon a prendre'n l'antiimatge. El \emph{pullback}\index{pullback}
generalitza aquesta operació, i la reindexació de subobjectes expressa
la substitució sense haver de parlar dels elements d'una categoria
arbitrària.

Les imatges i les adjuncions recuperen llavors l'existencial que havíem
trobat en compondre relacions. Projectar un predicat sobre $X\times A$
cap a $X$ significa retenir els $x$ per als quals hi ha algun $a$ que
el satisfà. En una categoria regular, aquesta construcció dona un
adjunt per l'esquerra de la reindexació. L'adjunt per la dreta, que
interpretaria el quantificador universal, requereix estructura
addicional: no es dedueix de la regularitat per si sola.

Finalment, el \emph{classificador de subobjectes} generalitza una altra
operació familiar. Un subconjunt $S\subseteq A$ es pot descriure per
la seva funció característica $A\to\{0,1\}$. En un topos elemental,
un objecte $\Omega$ fa aquest paper: els morfismes $A\to\Omega$
classifiquen els subobjectes d'$A$. Aquí es retroben dues línies del
llibre, la representabilitat i la lectura dels subobjectes com a
predicats. No hem d'esperar, però, que $\Omega$ es comporti sempre
com els dos valors de veritat clàssics de $\cat{Set}$.

Aquesta és la destinació del volum: arribar al llindar de la lògica
categòrica amb les eines necessàries per comprendre'n el llenguatge.
El desenvolupament sistemàtic de la sintaxi, la semàntica, la correcció,
la completesa i les condicions de coherència queda reservat al volum
següent. El pont final n'explica la continuació; no substitueix aquest
desenvolupament. La lògica dona direcció al recorregut, però no ens
eximeix d'aprendre les construccions categòriques amb les seves
hipòtesis i demostracions.

\section{Què vol dir haver après aquest llenguatge}

Aprendre categories no consisteix a abandonar els exemples concrets
tan aviat com tenim una definició general. Consisteix a poder fer el
camí d'anada i tornada: reconèixer una estructura en un exemple,
formular-la amb fletxes i tornar a interpretar-la en un altre context.
Una bona abstracció ens ha de permetre veure més, no només dir menys.

Al llarg del llibre adquirirem l'hàbit de preguntar-nos:
\begin{itemize}
  \item Quins són els objectes i les fletxes, i què significa compondre-les?
  \item Quina propietat expressa el diagrama que tenim davant?
  \item Quines dades caracteritzen una construcció i quines en són
        només una realització particular?
  \item Què queda determinat llevat d'isomorfisme, i quines
        compatibilitats fan únic aquest isomorfisme?
  \item Què preserva el functor amb què canviem de context?
  \item Quin és l'enunciat dual i sota quines hipòtesis té sentit?
  \item Quina lectura lògica té la construcció, si el context la permet?
\end{itemize}

Els primers capítols fixaran el vocabulari; els següents ensenyaran
a organitzar-lo amb propietats universals, representabilitat, límits
i adjuncions; els darrers mostraran com aquestes eines es combinen
en el tancament cartesià, els subobjectes i els topos. La guia de
lectura que acompanya el llibre orienta sobre com treballar aquest
recorregut amb la teoria, la pràctica, els exercicis i els tests.

El resultat que busquem no és disposar d'una paraula nova per a cada
construcció coneguda. És poder reconèixer, davant d'una situació nova,
quines relacions en governen l'estructura. Quan una intersecció, una
conjunció i un producte deixen de ser semblances vagues i esdevenen
manifestacions d'una mateixa propietat, el llenguatge categòric
comença a fer la seva feina.

\endgroup

\chapter*{\texorpdfstring{Guia de lectura i itineraris\\ d'aprenentatge}{Guia de lectura i itineraris d'aprenentatge}}
\addcontentsline{toc}{chapter}{Guia de lectura}
\markboth{Guia de lectura}{Guia de lectura}

\begingroup
\renewcommand{\thesection}{\arabic{section}}
\setcounter{section}{0}
\setlength{\emergencystretch}{1.5em}

Aquest capítol és un mapa per a la primera lectura. No cal aprendre
ara els noms de totes les construccions que hi apareixen: les trobarem
definides al seu lloc. La seva funció és mostrar què aporta cada etapa,
com es relacionen les quatre parts del llibre i com podem comprovar
que l'estudi ens està fent avançar.

La introducció explica \emph{per què} val la pena aprendre categories.
Aquí ens ocuparem de \emph{com} fer-ho. Podeu llegir aquesta guia
abans de començar i tornar-hi en acabar un bloc de capítols, quan
els noms del mapa ja s'hagin convertit en construccions familiars.


\section{El punt de partida i la destinació}

El llibre no pressuposa coneixements previs de teoria de categories.
Sí que pressuposa una familiaritat elemental amb el llenguatge
matemàtic: conjunts, aplicacions, composició, relacions i
demostracions senzilles. Per començar, convé poder fer quatre coses:
\begin{itemize}
  \item distingir el domini i el codomini d'una aplicació i calcular
        una composició;
  \item treballar amb productes cartesians, subconjunts i antiimatges;
  \item llegir quantificadors i separar una afirmació d'existència
        d'una afirmació d'unicitat;
  \item distingir una implicació, el seu recíproc i una equivalència,
        i entendre el paper d'un contraexemple.
\end{itemize}
Si alguna d'aquestes operacions és poc familiar, és millor dedicar-hi
una breu revisió que atribuir-ne la dificultat a l'abstracció categòrica.

No cal dominar prèviament tots els àmbits que proporcionen exemples.
Els grups, els espais vectorials i la topologia enriqueixen la lectura,
però, si encara no els coneixeu, podeu començar amb conjunts, preordres
i grafs. Els exemples del dual i del bidual demanen àlgebra lineal;
els que utilitzen continuïtat demanen nocions de topologia. Es poden
reprendre en una segona lectura sense abandonar el fil general.
La lògica aporta exemples des del començament, però no cal conèixer
ja la lògica categòrica per seguir-los.

La destinació és una comprensió inicial de les categories, els
functors i la naturalitat, seguida de les propietats universals,
Yoneda, els límits i les adjuncions. Aquest bagatge permet arribar
al tancament cartesià, als subobjectes i a una primera introducció
als topos. El volum prepara l'estudi posterior de la lògica categòrica;
no pretén desenvolupar-ne aquí una teoria completa.

\section{Quatre parts per a un mateix aprenentatge}

El llibre separa físicament la \textbf{Teoria}, la \textbf{Pràctica},
els \textbf{Exercicis proposats} i els \textbf{\NomTests}.
Aquesta separació facilita la consulta, però no obliga a estudiar
primer tota la teoria i només després tota la pràctica. Per a una
primera lectura, és més útil treballar \emph{per temes}: llegir un
capítol teòric, estudiar els problemes resolts corresponents,
intentar exercicis i comprovar la comprensió amb el test del mateix
tema.

\begin{center}
\renewcommand{\arraystretch}{1.2}
\begin{tabularx}{\linewidth}{@{}>{\raggedright\arraybackslash}p{0.25\linewidth}
                                  >{\raggedright\arraybackslash}X@{}}
\toprule
\textbf{Part} & \textbf{La pregunta que ajuda a respondre} \\
\midrule
Teoria & Quina és la definició, per què té sentit i què se'n pot deduir? \\
Pràctica & Com s'utilitza en un problema i com s'escriu l'argument? \\
Exercicis proposats & Puc reconèixer i aplicar la idea sense una solució completa? \\
\NomTests & Distingeixo el concepte de les confusions que s'hi assemblen? \\
\bottomrule
\end{tabularx}
\end{center}

Els nou primers temes tenen capítols corresponents en les quatre
parts. La numeració torna a començar en cada part; per això, quan
cerqueu material d'un tema, fixeu-vos tant en el número com en la
part i el títol. El capítol teòric desè, \emph{Pont cap a la lògica
categòrica}, és un epíleg d'orientació i no té un capítol paral·lel
de pràctica, exercicis o tests.

Els blocs d'obertura dels capítols teòrics indiquen prerequisits
i objectius. Llegiu-los com una invitació a activar coneixements
anteriors. Els blocs de síntesi del final recullen les idees centrals,
els errors típics i lectures recomanades. Serveixen per revisar,
però no substitueixen els exemples i les demostracions del capítol.

\section{El mapa conceptual dels capítols}
\label{sec:mapa-lectura}

\ifpublicacioparcial
\noindent\textbf{Nota sobre aquesta edició.} Aquesta versió parcial
només inclou els capítols 1--3. El mapa descriu, tanmateix, els deu
capítols del volum complet, perquè pugueu situar el que estudieu
dins del recorregut sencer.
\fi

El recorregut es pot agrupar en quatre etapes. Els capítols 1--3
construeixen el llenguatge i els nivells de comparació; els 4--6
organitzen les construccions mitjançant propietats universals;
els 7--9 combinen aquestes eines en estructures amb lectura lògica;
el 10 mostra la continuació. Aquesta divisió és una ajuda per orientar-se,
no una frontera rígida: les idees de les primeres etapes reapareixen
constantment en les següents.

\subsection*{1. Categories: aprendre la gramàtica de les fletxes}

El primer capítol introdueix objectes, morfismes, identitats i
composició. També ensenya a llegir diagrames, construeix exemples
i diferencia isomorfismes, monomorfismes, epimorfismes, seccions
i retraccions. La categoria oposada i la dualitat permeten
reutilitzar definicions i arguments en l'altra direcció.

És un capítol extens perquè fixa els hàbits que necessitarem després.
No convé voler memoritzar tots els exemples alhora. Escolliu-ne uns
quants i pregunteu sempre què són els objectes, què són les fletxes
i com es componen. Les construccions de categories noves i les
precaucions de mida amplien aquest vocabulari.

\noindent\textbf{Comprovació de progrés.} Podeu verificar els axiomes
en un exemple concret i explicar per què una fletxa és, o no és,
un isomorfisme? Podeu escriure l'equació d'un quadrat commutatiu
sense invertir l'ordre de la composició?

\subsection*{2. Functors: passar d'un context a un altre}

El segon capítol puja un nivell: no compara només objectes d'una
categoria, sinó categories entre elles. Els functors actuen sobre
objectes i morfismes i preserven identitats i composició. El graf
subjacent i la categoria lliure permeten veure una parella de
construccions que reapareixerà com a adjunció.

Cal dedicar una atenció especial als homfunctors. $\Hom(A,-)$
actua per postcomposició; $\Hom(-,B)$, per precomposició i amb
inversió de sentit. Aquesta variància no és un detall de notació:
és la base de la representabilitat i de Yoneda. Les nocions de
functor ple i fidel descriuen l'acció sobre cada homconjunt.

\noindent\textbf{Comprovació de progrés.} Per definir un functor,
podeu donar-ne l'acció sobre les fletxes, no només sobre els objectes?
Podeu explicar què fa $\Hom(-,B)$ amb una fletxa $X'\to X$?

\subsection*{3. Transformacions naturals i equivalències: exigir coherència}

El tercer capítol compara functors. La naturalitat exigeix que
una família de morfismes sigui compatible amb totes les fletxes
del domini. Els exemples de la interpretació de fórmules i del
bidual mostren per què aquesta compatibilitat és matemàticament
significativa, no només una convenció gràfica.

Les categories de functors permeten tractar aquestes comparacions
com a morfismes. L'equivalència de categories precisa quan dos
contextos descriuen la mateixa estructura essencial sense tenir
la mateixa presentació. La composició horitzontal afegeix un
segon tipus de composició que cal distingir de la vertical.

\noindent\textbf{Comprovació de progrés.} Podeu dibuixar el quadrat
de naturalitat d'una família proposada i comprovar-lo per a una
fletxa arbitrària? Podeu distingir un isomorfisme d'objectes,
un isomorfisme natural i una equivalència de categories?

\subsection*{4. Propietats universals, representabilitat i Yoneda:
caracteritzar pel paper}

Aquí canvia el centre de gravetat. En lloc de construir un objecte
i estudiar-ne després les propietats, busquem un objecte que
resolgui una necessitat expressada amb fletxes. Els objectes
inicials i terminals ofereixen els primers casos; les categories
coma ajuden a reunir dades i compatibilitats en un mateix context.

La representabilitat expressa aquestes caracteritzacions amb
homfunctors. Yoneda fa precisa la idea de conèixer un objecte
pel sistema de morfismes que hi arriben. És normal que aquesta
etapa demani més d'una lectura: cal coordinar categories, functors,
transformacions i elements d'homconjunts.

\noindent\textbf{Comprovació de progrés.} Podeu demostrar que dos
objectes terminals són isomorfs i indicar on s'utilitza la unicitat?
En una representació, podeu distingir el functor representat,
l'objecte representant i la dada universal? Podeu seguir la
bijecció de Yoneda a partir de l'avaluació en la identitat?

\subsection*{5. Límits i colímits: reunir construccions en un patró}

Un diagrama fixa una configuració d'objectes i morfismes; un con
és una manera compatible de relacionar-s'hi. El límit és el con
universal. Segons la forma del diagrama, obtenim productes,
equalitzadors, pullbacks o l'objecte terminal. Els colímits en
són la versió dual.

El pullback mereix una atenció especial, perquè reapareixerà
en la reindexació i els subobjectes. També cal separar dues
preguntes: si una categoria té un límit i si un functor el
preserva. No són la mateixa afirmació.

\noindent\textbf{Comprovació de progrés.} Podeu calcular un pullback
de conjunts i comprovar-ne la propietat universal? Podeu
distingir el diagrama, el con i el seu vèrtex, i explicar
què vol dir que un functor preserva aquest límit?

\subsection*{6. Adjuncions: dues descripcions de la mateixa dada}

Les adjuncions relacionen functors mitjançant una bijecció
natural d'hom\-con\-junts. El capítol comença amb preordres,
on la correspondència es llegeix com una equivalència de
desigualtats. Després introdueix unitat, counitat i identitats
triangulars, i mostra com les adjuncions expliquen la preservació
de límits i colímits.

Per a una primera aproximació, escolliu una adjunció i seguiu-la
en les dues descripcions: per bijeccions d'homconjunts i per
unitat i counitat. Les construccions lliures i la currificació
connecten aquest capítol amb exemples ja coneguts.

\noindent\textbf{Comprovació de progrés.} Podeu dir quins morfismes
es corresponen en una adjunció concreta i construir-ne els
transposats? Podeu recordar, amb una justificació conceptual,
que els adjunts per la dreta preserven límits i els de
l'esquerra preserven colímits?

\subsection*{7. Categories cartesianes tancades: funcions dins del context}

Els productes i els exponencials permeten parlar d'aparellament,
avaluació i currificació dins d'una categoria. El capítol
caracteritza el tancament cartesià i n'explica la relació
amb el càlcul lambda simplement tipat: els tipus es llegeixen
com a objectes i els termes, en context, com a morfismes.

El conjunt de totes les aplicacions $B\to C$ ofereix l'exemple
més familiar d'exponencial. Els preordres permeten reconèixer
la implicació com a adjunt de la conjunció. Els contraexemples
recorden que no tota categoria té aquesta estructura.

\noindent\textbf{Comprovació de progrés.} Podeu passar d'una fletxa
$A\times B\to C$ a la seva currificació $A\to C^B$ i tornar-ne
amb l'avaluació? Podeu explicar per què tenir productes no
garanteix tenir exponencials?

\subsection*{8. Subobjectes, imatges i factoritzacions: de les parts als predicats}

Un subobjecte no és simplement un subconjunt escrit amb un
altre nom: és una classe de monomorfismes amb un mateix
codomini, identificats de manera compatible. La reindexació
per pullback generalitza l'antiimatge i ofereix una lectura
estructural de la substitució.

Les interseccions expressen conjuncions de predicats; les
imatges i la regularitat permeten construir un adjunt
existencial de la reindexació. Aquesta etapa obliga a
controlar les hipòtesis: la regularitat no garanteix per
si sola unions de subobjectes ni un quantificador universal.

\noindent\textbf{Comprovació de progrés.} Podeu calcular una
reindexació a $\cat{Set}$ i interpretar-la com a substitució?
Podeu descriure la imatge d'un predicat sota una projecció
i explicar la relació $\exists_f\dashv f^*$?

\subsection*{9. Classificadors de subobjectes i topos: retrobar la representabilitat}

El classificador generalitza la funció característica d'un
subconjunt. La correspondència entre subobjectes d'$A$ i
morfismes $A\to\Omega$ recupera la representabilitat del
capítol quart. El topos elemental combina límits finits,
exponencials i classificador de subobjectes.

El capítol presenta $\cat{Set}$ i les categories de prefeixos
com a exemples. L'objectiu és comprendre les definicions
i la connexió entre les peces, no dominar encara la
teoria de feixos o la lògica interna dels topos.

\noindent\textbf{Comprovació de progrés.} Podeu reconstruir un
subconjunt com l'antiimatge del valor vertader per la
seva funció característica? Podeu enunciar les tres
condicions d'un topos i explicar què aporta cadascuna?

\subsection*{10. Pont cap a la lògica categòrica: saber què queda preparat}

L'epíleg situa les eines adquirides en un programa més ampli:
contextos, substitució, predicats, quantificadors, teories
i models. El mapa de fragments lògics i l'exemple de la
categoria sintàctica orienten sobre què es prolonga al
volum següent.

Llegiu-lo com una mirada retrospectiva i una invitació a
continuar, no com una exigència de dominar de cop tots els
temes anunciats. La sintaxi general, la correcció, la
completesa i les condicions de coherència pertanyen a
l'estudi posterior.

\noindent\textbf{Comprovació de progrés.} Quan l'epíleg parla de
models com a functors o de quantificadors com a adjunts,
podeu identificar les construccions d'aquest volum que
donen sentit a aquestes expressions?

\section{Com treballar un capítol}

Una lectura activa alterna comprensió, reconstrucció i aplicació.
Podeu utilitzar el cicle següent, adaptant-ne el ritme a la
dificultat del tema.

\begin{enumerate}
  \item \textbf{Situar la pregunta.} Llegiu l'obertura i els
        objectius del capítol. Escriviu en una frase quin
        problema intenta resoldre i quines idees anteriors
        necessita.
  \item \textbf{Desmuntar la definició.} Separeu les dades
        que s'han de donar, les operacions que s'hi afegeixen
        i les propietats que han de satisfer. Anoteu el
        domini i el codomini de cada fletxa.
  \item \textbf{Ancorar-la en exemples.} Treballeu un cas
        amb conjunts i un altre amb preordres o grafs.
        Quan pertoqui, afegiu-hi la lectura lògica. Busqueu
        també un exemple que no compleixi la definició.
  \item \textbf{Reconstruir una demostració.} Tanqueu el
        text i intenteu recuperar un argument curt. Després
        compareu-lo amb el llibre: on s'usa cada hipòtesi?
        On apareix l'existència i on la unicitat?
  \item \textbf{Fer servir la pràctica amb pausa.} Llegiu
        l'enunciat d'un problema resolt i proveu de fer-ne
        el primer pas abans de mirar la solució. Un cop
        llegida, torneu a escriure l'argument sense copiar-lo.
  \item \textbf{Intentar un exercici nou.} Comenceu pels
        de consolidació. Si us encalleu, identifiqueu el
        punt exacte i consulteu la indicació, no només
        una altra solució que s'hi assembli.
  \item \textbf{Tancar el cicle.} Feu el test, reviseu els
        errors i rellegiu la síntesi. Escriviu què podeu
        fer ara que no podíeu fer abans i què ha de quedar
        anotat per a una segona lectura.
\end{enumerate}

No cal reconstruir immediatament totes les proves llargues.
Però convé entendre'n l'estratègia i ser capaç de reproduir
arguments curts que reapareixen: provar una naturalitat,
usar la unicitat d'una propietat universal o construir
una fletxa amb el domini i el codomini correctes.

\subsection*{Una pauta per llegir diagrames}

Al llarg del llibre, quan una propietat categòrica té una expressió
natural amb fletxes, la presentem alhora amb un diagrama i amb les
equacions de composició que aquest expressa. Els diagrames formen
part del llenguatge matemàtic, no són una il·lustració opcional.
Aprofiteu aquesta doble presentació: passar del dibuix a les equacions
i de les equacions al dibuix és una destresa que cal practicar.

Davant d'un diagrama, no comenceu per recordar-ne el nom.
Seguiu aquesta pauta:
\begin{enumerate}
  \item Identifiqueu els objectes i els extrems de cada fletxa.
  \item Distingiu les dades donades de la fletxa que s'ha de construir.
  \item Escriviu els camins que tenen el mateix origen i el mateix destí.
  \item Traduïu-ne la commutativitat en equacions de composició.
  \item Si hi ha una propietat universal, digueu sobre quines
        dades es quantifica i què es demana que sigui únic.
\end{enumerate}

Una fletxa discontínua en un diagrama universal no es construeix
perquè el dibuix ho suggereixi: l'existència es justifica amb
una hipòtesi o una propietat. Igualment, un quadrat commutatiu
no és automàticament un pullback. Cal comprovar-ne també la
propietat universal.

\section{Els exemples recurrents com a punts de suport}

Canviar d'exemple no ha de significar tornar a començar.
Les mateixes famílies permeten veure cares diferents d'una idea.

\begin{description}
  \item[Conjunts.] Permeten calcular explícitament amb elements,
        productes, antiimatges, funcions i imatges. Són un
        primer laboratori, però no autoritzen a usar elements
        dins d'una categoria arbitrària.
  \item[Preordres.] Simplifiquen els homconjunts: hi ha com a
        màxim una fletxa entre dos objectes. Això converteix
        propietats universals en condicions d'ordre i adjuncions
        en equivalències de desigualtats. La unicitat automàtica
        és una ajuda, però també pot amagar una dificultat
        que reapareix en categories no primes.
  \item[Grafs.] Fan visible la diferència entre tenir fletxes
        i disposar de composició. Cal distingir la categoria
        de grafs, on les fletxes són homomorfismes de grafs,
        de la categoria lliure d'un graf, on són camins.
        La construcció lliure connecta els primers capítols
        amb les propietats universals i les adjuncions.
  \item[Lògica.] Permet interpretar deducció, conjunció,
        substitució i quantificació en els contextos adequats.
        Cal separar sempre la categoria prima que registra
        deduïbilitat de la perspectiva que conserva les proves,
        i també dels subobjectes que representen predicats.
\end{description}

Un quadern de correspondències pot ajudar: per a cada nova
construcció, escriviu què significa en dues d'aquestes famílies,
quina propietat comparteixen i quines hipòtesis necessita
cada realització. Si no existeix en un dels exemples, anoteu-ho:
aquest límit també forma part de la comprensió.

\section{Itineraris possibles}

Els itineraris següents són propostes de treball, no permisos
per ignorar les dependències. L'ordre del llibre continua sent
el recorregut més segur per a una primera lectura completa.

\subsection*{Una primera entrada: els tres capítols inicials}

Treballeu categories, functors i transformacions naturals amb
els seus materials paral·lels de pràctica, exercicis i tests.
La fita és poder moure-us entre els tres nivells: objectes i
morfismes, categories i functors, functors i transformacions.
No cal dominar encara tot el desplegament de la composició
horitzontal, però sí distingir-la de la vertical.

\ifpublicacioparcial
Aquesta edició parcial conté precisament aquests tres temes en
les quatre parts, a més dels materials complementaris. No és una
introducció completa als límits, les adjuncions o els topos;
les remissions als capítols absents s'han d'entendre com a
indicacions per continuar amb el volum complet.
\fi

\subsection*{Un recorregut complet d'autoaprenentatge}

Seguiu els nou temes amb el cicle de treball per capítols i
llegiu després el pont final. Cada cop que arribeu a una
definició general, torneu a un exemple anterior. Després
de Yoneda, reviseu els homfunctors; després de les adjuncions,
reviseu la construcció lliure; després del classificador,
reviseu la representabilitat. Aquestes tornades no són
repeticions inútils: fan visible la unitat del recorregut.

\subsection*{Un curs semestral}

Per a un curs d'un semestre, una selecció natural és treballar els capítols 1--6 com a nucli obligatori i triar, segons l'orientació del curs, els capítols 7--9. La teoria s'ha de combinar cada setmana amb una selecció de problemes resolts i exercicis proposats; els tests poden servir d'autoavaluació, no de substitut d'una prova escrita. El capítol 10 funciona bé com a sessió final de síntesi i d'obertura cap a la lògica categòrica.

\subsection*{Una lectura per al professorat}

Per preparar docència, convé començar per aquesta guia i pels blocs d'objectius i d'errors típics de cada capítol. Les parts de Pràctica permeten seleccionar exemples per a classe, mentre que els Exercicis proposats ofereixen material de treball autònom. Els quatre fils recurrents ---conjunts, preordres, grafs i lògica--- es poden modular segons el perfil de l'alumnat, però és recomanable conservar-ne almenys dos en paral·lel per evitar que les definicions quedin lligades a una sola categoria concreta.

\subsection*{Una lectura orientada a la lògica}

Manteniu el recorregut general, però doneu més pes a la
deduïbilitat del capítol 1 i als exemples d'interpretació
i substitució dels capítols 2 i 3. Als capítols 4--6,
consolideu representabilitat, pullbacks i adjuncions;
després, treballeu amb especial cura el tancament cartesià,
els subobjectes i el classificador.

No convé saltar directament a l'epíleg: expressions com
\enquote{el quantificador és un adjunt} orienten, però només es
comprenen quan sabem entre quines categories o preordres
actua l'adjunció i què afirma la seva correspondència.
Les equivalències lògiques no substitueixen la comprovació
de les hipòtesis categòriques.

\subsection*{Una lectura orientada als tipus i a la computació}

Feu servir els grafs i les categories lliures per entendre
la composició. Després de consolidar functors i naturalitat,
poseu l'accent en propietats universals, productes,
adjuncions, exponencials i currificació. El capítol 7
reuneix aquestes eines en la semàntica de tipus i termes.

Els capítols 8 i 9 poden constituir una segona etapa
si el primer objectiu és el càlcul lambda simplement tipat.
Caldrà reprendre'ls per arribar als predicats, els
classificadors i la perspectiva lògica més àmplia.

\section{Convencions tipogràfiques i de lectura}

Els noms de categories s'escriuen amb tipografia matemàtica, com $\cat{Set}$, $\cat{Graph}$ o $\cat{Cat}$. Els termes definits i les idees que cal retenir apareixen destacats dins dels entorns corresponents; el color reforça la jerarquia visual, però l'etiqueta textual ---\enquote{Definició}, \enquote{Teorema}, \enquote{Exemple}, \enquote{Exercici}, etc.--- n'indica sempre la funció. Les fletxes discontínues dels diagrames assenyalen habitualment el morfisme que una propietat universal afirma que existeix; el text adjacent n'explicita el paper. Els símbols amb més d'un ús, especialment $f^{*}$, es desambigüen pel context i es recullen a la taula de notació.

\section{Exercicis i tests: què comproven i què no}

Els exercicis proposats utilitzen una estrella,
$\star$, per assenyalar consolidació, i dues,
$\star\star$, per assenyalar ampliació. Comenceu pels
primers, però no interpreteu les marques com una mesura
exacta del temps que necessitareu. Una dificultat breu
de variància pot resultar més important que un càlcul llarg.

Les indicacions són una ajuda per reprendre el treball.
Abans de consultar-les, escriviu les dades, la conclusió
i almenys un intent. Després d'utilitzar una indicació,
torneu a fer l'exercici sense mirar-la. L'objectiu no és
haver reconegut una solució, sinó poder construir un argument.

Els tests tenen una única resposta correcta per pregunta.
\ifimprimible
En aquesta edició per imprimir, les respostes i les
justificacions de cada test es reuneixen al final del llibre
(apèndix \enquote{Solucions dels tests}). L'edició de pantalla
ofereix, a més, correcció automàtica en els lectors de PDF
compatibles.
\else
La correcció automàtica necessita un lector de PDF compatible
amb formularis AcroForm i JavaScript. Si el lector no
l'admet, podeu respondre igualment les preguntes i
consultar les respostes i les justificacions de cada test,
reunides al final del llibre (apèndix \enquote{Solucions dels
tests}), també usables en paper.
\fi
Consulteu-les només després d'haver respost tot el test.

Una resposta encertada no demostra per si sola que sapigueu
resoldre un problema. Per aprofitar el test, justifiqueu
per què l'opció triada és correcta i per què les altres
fallen. Si heu dubtat entre dues opcions, convertiu el
dubte en una pregunta per revisar a la teoria o a la
pràctica. Aquesta explicació és més útil que la puntuació
obtinguda sense justificació.

\section{Algunes dificultats que convé reconèixer aviat}

Hi ha errors que travessen diversos capítols. Identificar-los
permet corregir la manera de treballar, no només una resposta.

\begin{description}
  \item[Compondre sense comprovar els tipus.]
        Abans d'escriure $g\comp f$, verifiqueu que el
        codomini de $f$ és el domini de $g$. L'associativitat
        no fa componibles fletxes amb extrems incompatibles.
  \item[Confondre un exemple amb la definició general.]
        A $\cat{Set}$, mono i epi corresponen a injectivitat
        i exhaustivitat. En altres categories, les definicions
        són condicions de cancel·lació. Un mono que també
        és epi no és sempre un isomorfisme.
  \item[Tractar la naturalitat com una etiqueta.]
        Una família de fletxes no és natural només perquè
        s'hagi definit amb la mateixa fórmula. Cal escriure
        el quadrat i comprovar-lo per a tots els morfismes.
  \item[Oblidar les dades universals.]
        El producte inclou les projeccions; el límit inclou
        el con; l'exponencial inclou l'avaluació. La unicitat
        de l'isomorfisme es refereix a les compatibilitats
        amb aquestes dades, no a l'objecte nu.
  \item[Canviar l'orientació sense controlar-la.]
        La contravariància, la dualitat i la reindexació
        inverteixen direccions de maneres que cal precisar.
        Escriviu dominis i codominis abans de fer servir
        una regla recordada de memòria.
  \item[Afegir estructura que no s'ha suposat.]
        Tenir productes no implica tenir exponencials;
        la regularitat no dona per si sola disjuncions
        ni quantificació universal; un topos no és
        necessàriament booleà. Rellegiu les hipòtesis.
\end{description}

Si una demostració sembla bloquejada, pregunteu-vos primer
si el problema és de definició, de tipus, de variància o
d'hipòtesis. Moltes dificultats aparentment abstractes
es resolen quan precisem una d'aquestes quatre coses.

\section{Quan consultar els apèndixs i els materials finals}

Els apèndixs no constitueixen una etapa que s'hagi de
completar obligatòriament abans de començar la teoria.
Es poden consultar quan la lectura ho demani.

\begin{description}
  \item[Apèndix A: axiomàtica relacional.]
        Ofereix una presentació en què les fletxes són
        els individus primitius i els objectes es
        reconstrueixen a partir de les identitats.
        També hi trobareu la formalització del principi
        de dualitat, en versió sintàctica i semàntica.
        És una ampliació adequada després del primer
        capítol, sobretot per a qui s'interessi per
        la formalització lògica. Demana familiaritat
        amb lògica de primer ordre amb igualtat.
  \item[Apèndix B: grafs.]
        Aclareix les convencions, els homomorfismes,
        els camins i el pas a categories. Consulteu-lo
        quan treballar amb camins, llaços o arestes
        múltiples us plantegi dubtes en els primers
        capítols.
  \item[Apèndix C: convenció de mida.]
        Explica el marc de conjunts i classes del llibre.
        És especialment útil quan apareixen
        $\cat{Cat}$, categories de functors i
        categories de prefeixos. Cal distingir
        \enquote{petita} de \enquote{localment petita}: la segona
        condició garanteix homconjunts, no que tots
        els objectes formin un conjunt.
  \item[Apèndix D: solucions dels tests.]
        Respostes i justificacions. Consulteu-les només
        després d'haver respost tot el test.
\ifpublicacioparcial\else
  \item[Apèndix E: quasiexemples i contraexemples.]
        Recull estructures que fallen exactament un axioma
        (categoria, functor, transformació natural), la
        jerarquia entre igualtat, isomorfisme, isomorfisme
        natural i equivalència, i models de la lògica que
        queden fora dels topos. La secció sobre la
        jerarquia és útil des del capítol~\ref{cap:transformacions};
        la de la lògica, després del capítol~\ref{cap:topos}.
\fi
\end{description}

La taula de notació serveix per recuperar el significat
i la primera aparició dels símbols. El glossari
català--anglès ajuda a llegir la bibliografia internacional,
sobretot en termes com \emph{pullback}, \emph{slice},
\emph{presheaf} i \emph{whiskering}. L'índex alfabètic
és una eina per retrobar conceptes quan reapareixen en
un context nou.

Les lectures recomanades al final dels capítols tenen
una funció de contrast i ampliació. En una primera
lectura, és millor consultar una font alternativa per
aclarir una pregunta concreta que obrir alhora molts
tractaments generals del mateix tema.

\section{Una manera de mesurar l'aprenentatge}

No mesureu el progrés només pel nombre de pàgines llegides.
Una fita més fiable és poder fer, sense el text davant,
quatre operacions amb una idea nova:
\begin{enumerate}
  \item enunciar-ne la definició amb les hipòtesis correctes;
  \item interpretar-la en almenys dos exemples;
  \item utilitzar-la en un argument o un càlcul senzill;
  \item explicar què falla en un contraexemple o quan
        es retira una hipòtesi.
\end{enumerate}

No cal satisfer aquestes quatre condicions amb la mateixa
fluïdesa des del primer dia. És raonable avançar amb
una dificultat ben identificada i tornar-hi després.
El que convé evitar és substituir la comprensió per
la familiaritat visual amb els símbols.

En acabar el llibre, el guany més important serà poder
relacionar les peces: veure una propietat universal
darrere d'una construcció, una naturalitat darrere
d'una compatibilitat, una adjunció darrere d'una
correspondència i una operació sobre subobjectes
darrere d'una lectura lògica. Aquesta capacitat de
passar entre exemples i estructura és el criteri
que dona unitat a tot el recorregut.

\endgroup


% ============================================================
% PART I — TEORIA
% ============================================================
\part{Teoria}
\setcounter{chapter}{0}
\chapter{Categories}
\label{cap:categories}

\begin{objectius}
\apartat{Prerequisits} Llenguatge elemental de conjunts i aplicacions: domini, codomini i composició; producte cartesià, subconjunts i relacions. Lectura de quantificadors elementals i distinció entre existència i unicitat. Nocions bàsiques d'implicació, equivalència i contraexemple. Els exemples amb grups, espais vectorials o espais topològics enriqueixen la lectura, però no són imprescindibles. No es pressuposa cap coneixement previ de teoria de categories.
\apartat{Objectius} En acabar el capítol, el lector hauria de poder:
\begin{itemize}
  \item enunciar la definició de categoria i verificar-ne els axiomes en exemples concrets;
  \item llegir i escriure diagrames commutatius i traduir-los a igualtats de composicions;
  \item reconèixer la mateixa estructura categòrica en conjunts, preordres, monoides, grafs, relacions i sistemes deductius;
  \item definir i utilitzar isomorfismes, monomorfismes, epimorfismes, seccions i retraccions, distingint-los de les nocions conjuntistes corresponents;
  \item construir i interpretar la categoria oposada, categories producte, subcategories i categories sobre i sota un objecte en casos senzills;
  \item aplicar el principi de dualitat a definicions i resultats elementals;
  \item descriure la categoria lliure generada per un graf finit senzill;
  \item distingir categories petites i localment petites.
\end{itemize}
\end{objectius}

\section{Cap al concepte de categoria}

Quan comencem a estudiar matemàtiques, aprenem a treballar amb molts tipus d'\emph{estructures}: ordres, grups, espais vectorials, espais topològics, etc. Aviat notem que aquestes estructures rarament apareixen soles. Al costat de cadascuna hi ha sempre una classe de transformacions que les relacionen respectant la seva manera de comportar-se: entre ordres, les aplicacions monòtones; entre grups, els homomorfismes; entre espais vectorials, les aplicacions lineals; entre espais topològics, les aplicacions contínues. Aquestes transformacions admissibles les anomenarem, en general, \emph{morfismes} o \emph{fletxes}.

La noció de categoria no decideix per si sola quins morfismes són admissibles: ho fixem nosaltres en cada context, d'acord amb l'estructura que volem respectar. El que fa la teoria de categories és abstreure el patró comú a tots aquests contextos: treballar amb objectes, amb transformacions admissibles entre ells i amb la manera com aquestes transformacions es componen. Cada categoria concreta enregistra aquesta tria de morfismes i, a partir d'aquí, només en reté la composició i les identitats.

Els exemples concrets que acabem d'esmentar tenen tots un conjunt subjacent. En aquests casos, preservar l'estructura significa començar per una aplicació entre els conjunts subjacents i imposar-hi les condicions adequades. No és, però, un requisit de la noció de categoria: més endavant trobarem categories els objectes de les quals no són conjunts amb estructura. Com hem dit a la secció~\ref{sec:mirar-fora} de la introducció, la teoria de categories desplaça la visió clàssica, puntual i conjuntista, de les aplicacions cap a una visió abstracta com a fletxes.

Recordem breument la notació de fletxa per a les aplicacions entre conjunts. Una aplicació $f$ amb domini $X$ i codomini $Y$ es denota $f\colon X\to Y$, o bé com a fletxa $X\xrightarrow{f}Y$. L'operació fonamental sobre les aplicacions és la \emph{composició}: si $f\colon X\to Y$ i $g\colon Y\to Z$, es defineix $g\comp f\colon X\to Z$ per
\[
(g\comp f)(x)\coloneqq g(f(x)).
\]
Perquè la composició estigui definida, el codomini de $f$ ha de coincidir amb el domini de $g$; amb fletxes, $X\xrightarrow{f}Y\xrightarrow{g}Z$. A més, per a cada conjunt $X$ hi ha l'aplicació \emph{identitat} sobre $X$,
\[
\id_X\colon X\to X,\qquad \id_X(x)\coloneqq x.
\]
Aquestes operacions es regeixen per la llei d'associativitat i les lleis de la unitat: per a $f\colon X\to Y$, $g\colon Y\to Z$ i $h\colon Z\to W$,
\[
(h\comp g)\comp f=h\comp(g\comp f),\qquad f\comp\id_X=f=\id_Y\comp f.
\]
Observem que aquestes equacions es formulen purament en termes de les operacions sobre aplicacions, sense cap referència als elements dels conjunts $X$, $Y$, $Z$, $W$.

Això no és cap ingenuïtat, sinó una evidència sobre les estructures que hem mencionat: els homomorfismes entre grups, les aplicacions monòtones entre preordres, les aplicacions lineals entre espais vectorials i les aplicacions contínues entre espais topològics satisfan les condicions que acabem d'escriure.

Aquesta és la idea que orientarà tot el que segueix: estudiar un objecte no pel que és per dins, sinó pel paper que fa dins del sistema de fletxes que el connecten amb els altres. És el que ens permetrà de parlar-ne amb un mateix llenguatge i, fins i tot, estendre'l a altres tipus d'objectes no necessàriament basats en la noció d'aplicació, com ara grafs, sistemes deductius o relacions entre conjunts.

\section{Definició de categoria}

\begin{definicio}
\label{def:categoria}
\index{categoria}Una \textbf{categoria} $\cat{C}$ consta de les dades següents:
\begin{enumerate}
\item una col·lecció d'\textbf{objectes}\index{objecte} $A,B,C,\dots$, que denotem $\Ob(\cat{C})$;
\item una col·lecció de \textbf{morfismes}\index{morfisme} $f,g,h,\dots$, que denotem $\Mor(\cat{C})$; cada morfisme $f$ té associats un únic \textbf{domini}\index{domini} $A$ i un únic \textbf{codomini}\index{codomini} $B$, i ho escrivim com $A=\dom(f)$ i $B=\cod(f)$; en aquest cas també escrivim $f\colon A\to B$ o com a fletxa $A\xrightarrow{f}B$;
\item per a cada objecte $A$, un morfisme \textbf{identitat}\index{identitat} $\id_A\colon A\to A$;
\item per a cada parella de morfismes componibles $f\colon A\to B$ i $g\colon B\to C$, la \textbf{composició}\index{composicio@composició} $g\comp f\colon A\to C$.
\end{enumerate}
Aquestes dades han de satisfer els dos axiomes següents, sempre que les composicions indicades tinguin sentit:
\begin{enumerate}
\item \textbf{Associativitat.} $h\comp(g\comp f)=(h\comp g)\comp f$, per a $f\colon A\to B$, $g\colon B\to C$, $h\colon C\to D$.
\item \textbf{Identitats.} $\id_B\comp f=f$ i $f\comp\id_A=f$, per a tot $f\colon A\to B$.
\end{enumerate}
La secció~\ref{sec:diagrames} dona la lectura diagramàtica d'aquests axiomes.
\end{definicio}

Parlem aquí de \emph{col·leccions} i no de \emph{conjunts}, deliberadament: en alguns exemples els objectes, o els morfismes, són massa nombrosos per formar un conjunt. Així, tant $\Ob(\cat{C})$ com $\Mor(\cat{C})$ són en general col·leccions, no necessàriament conjunts. Precisarem aquesta qüestió a la secció~\ref{sec:mida}. Sovint és més còmode agrupar els morfismes segons el seu domini i codomini: quan la col·lecció de morfismes de $A$ a $B$ és un conjunt, se sol denotar
\[
\Hom_{\cat{C}}(A,B)
\]
o simplement $\Hom(A,B)$ quan la categoria és clara i l'anomenem \emph{homconjunt} de $A$ i $B$. Com que cada morfisme té un únic domini i un únic codomini, $\Mor(\cat{C})$ sempre es reparteix segons el parell $(A,B)$ de domini i codomini, en parts disjuntes dos a dos;\footnote{Aquesta disjunció és, de fet, equivalent a la unicitat del domini i el codomini. A l'axiomàtica relacional de l'apèndix~\ref{cap:axiomatica-relacional} en fa el paper d'un axioma.} el que no està garantit és que cada part sigui un \emph{conjunt}. Quan ho és per a tots els parells ---propietat que rep el nom de \emph{petitesa local} (secció~\ref{sec:mida})---, aquestes parts són els homconjunts $\Hom(A,B)$ i $\Mor(\cat{C})$ n'és la unió. Al llarg del llibre totes les categories que fem servir són localment petites, de manera que aquesta subtilesa no ens afectarà.

\begin{observacio}
Fem algunes precisions sobre la definició.
\begin{enumerate}
\item Els morfismes també s'anomenen \textbf{fletxes}, i sovint farem servir les dues paraules indistintament.
\item L'axioma d'associativitat fa que puguem escriure $h\comp g\comp f$ sense parèntesis, sabent que el resultat no depèn de la manera d'agrupar-los.
\item Els axiomes d'identitat diuen que $\id_A$ es comporta com un neutre per la composició. Veurem tot seguit que aquesta propietat determina $\id_A$ de manera única.
\item Aquesta definició axiomàtica pren com a termes no definits els objectes i els morfismes. Una categoria és qualsevol estructura que contingui aquestes dues menes de dades i que satisfaci les condicions i propietats de la definició.\footnotemark
\end{enumerate}
\end{observacio}
\footnotetext{De fet, es pot prescindir d'una de les dues menes de dades: com que cada objecte queda determinat per la seva identitat, una categoria petita es pot descriure en un llenguatge amb una sola mena d'individus ---les fletxes--- i tres relacions primitives de domini, codomini i composició, amb els objectes reconstruïts com les fletxes identitat. L'apèndix~\ref{cap:axiomatica-relacional} desenvolupa aquesta presentació relacional i en demostra l'equivalència amb la definició d'aquí.}

\begin{proposicio}
La fletxa identitat d'un objecte és única.
\end{proposicio}

\begin{prova}
Suposem que $e_A\colon A\to A$ i $e'_A\colon A\to A$ són dues fletxes identitat de $A$. Com que $e_A$ és una identitat,
\[
e_A\comp e'_A=e'_A.
\]
D'altra banda, com que $e'_A$ és una identitat,
\[
e_A\comp e'_A=e_A.
\]
Per tant $e_A=e'_A$.
\end{prova}

\section{Fletxes, camins i diagrames commutatius}\label{sec:diagrames}\index{diagrama commutatiu}

La definició anterior es pot llegir algebraicament, mitjançant equacions, però també gràficament. Aquesta segona lectura serà una part essencial del llenguatge categòric.

Si $f\colon A\to B$ i $g\colon B\to C$ són componibles, la composició es pot representar amb el triangle
\[
\begin{tikzpicture}[baseline=(m.center)]
  \matrix (m) [matrix of math nodes, column sep=3.4em, row sep=2.7em] {
    A & B \\
      & C \\
  };
  \draw[-{Stealth[scale=1.1]}] (m-1-1) -- node[above] {$f$} (m-1-2);
  \draw[-{Stealth[scale=1.1]}] (m-1-2) -- node[right] {$g$} (m-2-2);
  \draw[-{Stealth[scale=1.1]}] (m-1-1) -- node[below left] {$g\comp f$} (m-2-2);
\end{tikzpicture}
\]
La fletxa directa de $A$ a $C$ és precisament el resultat de seguir el camí $A\to B\to C$.

Més generalment, un diagrama pot contenir diferents camins dirigits amb el mateix origen i el mateix destí. Per exemple,
\[
\begin{tikzpicture}[baseline=(m.center)]
  \matrix (m) [matrix of math nodes, row sep=2.7em, column sep=3.5em] {
    A & B \\
    C & D \\
  };
  \draw[-{Stealth[scale=1.1]}] (m-1-1) -- node[above] {$f$} (m-1-2);
  \draw[-{Stealth[scale=1.1]}] (m-1-1) -- node[left] {$u$} (m-2-1);
  \draw[-{Stealth[scale=1.1]}] (m-1-2) -- node[right] {$v$} (m-2-2);
  \draw[-{Stealth[scale=1.1]}] (m-2-1) -- node[below] {$g$} (m-2-2);
\end{tikzpicture}
\]
té dos camins de $A$ a $D$: primer $f$ i després $v$, o primer $u$ i després $g$.

\begin{definicio}
Un diagrama com l'anterior \textbf{commuta} si els dos camins determinen el mateix morfisme, és a dir, si
\[
v\comp f=g\comp u.
\]
En general, direm que un diagrama commuta quan tots els camins dirigits amb el mateix origen i el mateix destí que cal comparar donen la mateixa composició.
\end{definicio}

\begin{observacio}
Un dibuix de fletxes no és automàticament un diagrama commutatiu. El dibuix especifica objectes i morfismes; afirmar que \emph{commuta} afegeix una o més igualtats entre composicions de fletxes. Per això convé aprendre a passar en tots dos sentits:
\[
\text{diagrama}\quad\longleftrightarrow\quad\text{equacions entre fletxes i composicions}.
\]
\end{observacio}

L'associativitat es pot llegir com un diagrama commutatiu. En una cadena
\[
A\xrightarrow{f}B\xrightarrow{g}C\xrightarrow{h}D
\]
podem agrupar les composicions de dues maneres, $(h\comp g)\comp f$ o $h\comp(g\comp f)$. El diagrama següent ho mostra com un paral·lelogram de vèrtexs $A,B,D,C$ i costats $f$, $h\comp g$, $h$ i $g\comp f$; la diagonal interior és $g$.
\[
\begin{tikzpicture}[>={Stealth[scale=1.1]}]
  \node (A) at (-3,2) {$A$};
  \node (B) at (0,2)  {$B$};
  \node (C) at (0,0)  {$C$};
  \node (D) at (3,0)  {$D$};
  \draw[->] (A) -- node[above] {$f$} (B);
  \draw[->] (B) -- node[right] {$h\comp g$} (D);
  \draw[->] (C) -- node[below] {$h$} (D);
  \draw[->] (A) -- node[below left=-2pt] {$g\comp f$} (C);
  \draw[->] (B) -- node[left] {$g$} (C);
\end{tikzpicture}
\]
Recórrer $A\to B\to C\to D$ agrupant primer $g\comp f$ i després component amb $h$, o agrupant primer $h\comp g$ i precomponent amb $f$, dona la mateixa fletxa $A\to D$. Això és exactament l'axioma d'associativitat:
\[
(h\comp g)\comp f=h\comp(g\comp f).
\]

Les identitats també tenen una lectura diagramàtica. Si $f\colon A\to B$, inserir $\id_A$ abans de $f$ o $\id_B$ després de $f$ no canvia el camí:
\[
\begin{tikzpicture}[baseline=(m.center)]
  \matrix (m) [matrix of math nodes, row sep=2.6em, column sep=2.6em] {
    A & A \\
      & B \\
  };
  \draw[-{Stealth[scale=1.1]}] (m-1-1) -- node[above] {\(\id_A\)} (m-1-2);
  \draw[-{Stealth[scale=1.1]}] (m-1-2) -- node[right] {\(f\)} (m-2-2);
  \draw[-{Stealth[scale=1.1]}] (m-1-1) -- node[below left] {\(f\)} (m-2-2);
\end{tikzpicture}
\qquad\qquad
\begin{tikzpicture}[baseline=(m.center)]
  \matrix (m) [matrix of math nodes, row sep=2.6em, column sep=2.6em] {
    A & B \\
      & B \\
  };
  \draw[-{Stealth[scale=1.1]}] (m-1-1) -- node[above] {\(f\)} (m-1-2);
  \draw[-{Stealth[scale=1.1]}] (m-1-2) -- node[right] {\(\id_B\)} (m-2-2);
  \draw[-{Stealth[scale=1.1]}] (m-1-1) -- node[below left] {\(f\)} (m-2-2);
\end{tikzpicture}
\]
Els dos triangles commutatius expressen exactament els axiomes d'identitat:
\[
f\comp\id_A=f,\qquad \id_B\comp f=f.
\]
A partir d'ara seguirem aquesta norma: quan una propietat categòrica tingui una expressió diagramàtica natural, en donarem alhora, al mateix lloc, l'equació i el diagrama. Llegir l'una i dibuixar l'altre, i a l'inrevés, és una destresa que el llibre demana constantment.

\section{Exemples de categories}

La força de la definició es veu de seguida en la varietat d'exemples que hi encaixen. En cadascun cal identificar els objectes, els morfismes, les identitats i la composició. Comencem pels exemples més petits, que malgrat la seva simplicitat apareixeran sovint com a \emph{categories índex} (la definició~\refalt{def:diagrama}{de diagrama del capítol~5}) per construir diagrames.

\subsection{Categories finites}\label{subsec:finites}\index{categoria!buida}\index{categoria!unitat}\index{categoria!finita}

Per descomptat, els objectes d'una categoria no cal que siguin conjunts. Considerem alguns exemples molt senzills que es poden descriure per complet. En cada cas identifiquem els objectes, els morfismes, la composició i les identitats; l'associativitat i les lleis d'identitat es comproven directament, ja que hi ha molt poques composicions a considerar.
\begin{enumerate}
\item La \emph{categoria buida} $\cat{0}$ té $\Ob(\cat{0})=\varnothing$ i cap morfisme. No hi ha composició ni identitats que definir, i els axiomes es compleixen trivialment perquè no hi ha res a comprovar; és la categoria \enquote{sense res}.
\item La \emph{categoria unitat} $\cat{1}$, també anomenada categoria \emph{terminal}, té $\Ob(\cat{1})=\{\ast\}$ i $\Hom(\ast,\ast)=\{\id_\ast\}$. L'única composició és $\id_\ast\comp\id_\ast=\id_\ast$, i els axiomes es redueixen a aquesta igualtat:
\[
\begin{tikzpicture}[baseline=(a.base)]
  \node (a) {$\ast$};
  \draw[-{Stealth[scale=1.1]}] (a) to[out=40,in=-40,looseness=8] node[right] {$\id_\ast$} (a);
\end{tikzpicture}
\]
\item La categoria $\cat{2}$ té $\Ob(\cat{2})=\{\ast,\star\}$ i una única fletxa no trivial $u\colon\ast\to\star$, a més de les dues identitats:
\[
\begin{aligned}
\Hom(\ast,\ast)&=\{\id_\ast\}, &
\Hom(\star,\star)&=\{\id_\star\},\\
\Hom(\ast,\star)&=\{u\}, &
\Hom(\star,\ast)&=\varnothing.
\end{aligned}
\]
Les úniques composicions són les que involucren identitats, $u\comp\id_\ast=u=\id_\star\comp u$, de manera que els axiomes es compleixen automàticament. A partir d'ara, com és costum, no dibuixem les identitats:
\[
\begin{tikzpicture}[baseline=(a.base)]
  \node (a) at (0,0) {$\ast$};
  \node (b) at (1.8,0) {$\star$};
  \draw[-{Stealth[scale=1.1]}] (a) -- node[above] {$u$} (b);
\end{tikzpicture}
\]
\item La categoria amb \textbf{dues fletxes paral·leles} té $\Ob=\{\ast,\star\}$ i dues fletxes no trivials $f,g\colon\ast\to\star$ amb el mateix domini i el mateix codomini:
\[
\Hom(\ast,\star)=\{f,g\},\qquad
\Hom(\star,\ast)=\varnothing,
\]
i als homconjunts $\Hom(\ast,\ast)$ i $\Hom(\star,\star)$ només hi ha les identitats. Com que cap parell d'aquestes fletxes no és componible, de nou les úniques composicions són amb identitats:
\[
\begin{tikzpicture}[baseline=(a.base)]
  \node (a) at (0,0) {$\ast$};
  \node (b) at (1.8,0) {$\star$};
  \draw[-{Stealth[scale=1.1]}] ([yshift=2.2pt]a.east) -- node[above] {$f$} ([yshift=2.2pt]b.west);
  \draw[-{Stealth[scale=1.1]}] ([yshift=-2.2pt]a.east) -- node[below] {$g$} ([yshift=-2.2pt]b.west);
\end{tikzpicture}
\]
Aquesta darrera reapareixerà com a categoria índex dels \emph{equalitzadors} (la definició~\refalt{def:equalitzador}{d'equalitzador del capítol~5}).
\item La categoria $\cat{3}$ té $\Ob(\cat{3})=\{\ast,\star,\bullet\}$ i tres fletxes no trivials, $f\colon\ast\to\star$, $g\colon\star\to\bullet$ i $h\colon\ast\to\bullet$:
\[
\Hom(\ast,\star)=\{f\},\qquad
\Hom(\star,\bullet)=\{g\},\qquad
\Hom(\ast,\bullet)=\{h\};
\]
els homconjunts d'un objecte a si mateix contenen només la identitat, i els homconjunts en sentit contrari són buits. La composició obliga a incloure la composició de les dues fletxes consecutives: l'única composició no trivial és $g\comp f=h$. Amb aquesta composició (i les identitats) els axiomes se satisfan:
\[
\begin{tikzpicture}[baseline=(m.center)]
  \matrix (m) [matrix of math nodes, row sep=2.6em, column sep=2.6em] {
    \ast & \star \\
         & \bullet \\
  };
  \draw[-{Stealth[scale=1.1]}] (m-1-1) -- node[above] {$f$} (m-1-2);
  \draw[-{Stealth[scale=1.1]}] (m-1-2) -- node[right] {$g$} (m-2-2);
  \draw[-{Stealth[scale=1.1]}] (m-1-1) -- node[below left] {$h=g\comp f$} (m-2-2);
\end{tikzpicture}
\]
\end{enumerate}

\begin{observacio}
Afegir fletxes no és gratuït: la composició obliga a incloure totes les composicions. Dues fletxes de sentits oposats $f\colon\ast\to\star$ i $g\colon\star\to\ast$
\[
\begin{tikzpicture}[baseline=(a.base)]
  \node (a) at (0,0) {$\ast$};
  \node (b) at (2.2,0) {$\star$};
  \draw[-{Stealth[scale=1.1]}] ([yshift=2.6pt]a.east) -- node[above] {$f$} ([yshift=2.6pt]b.west);
  \draw[-{Stealth[scale=1.1]}] ([yshift=-2.6pt]b.west) -- node[below] {$g$} ([yshift=-2.6pt]a.east);
\end{tikzpicture}
\]
generen $g\comp f$, $f\comp g$, $g\comp f\comp g$, i així indefinidament; la categoria resultant és infinita. Per obtenir-ne una de finita cal imposar equacions entre composicions ---per exemple $g\comp f=\id_\ast$ i $f\comp g=\id_\star$, cosa que fa de $f$ un isomorfisme (la definició~\ref{def:isomorfisme}). En canvi, dues fletxes \emph{paral·leles} $f,g\colon\ast\to\star$ no són componibles entre elles, i la categoria és finita sense afegir res.
\end{observacio}

\subsection{Categories de conjunts}\label{subsec:conjunts}

La categoria $\cat{Set}$ té:
\begin{itemize}
\item objectes: $\Ob(\cat{Set})$ és la col·lecció de tots els conjunts;
\item morfismes: $\Hom_{\cat{Set}}(A,B)$ és el conjunt de totes les aplicacions $A\to B$;
\item composició: la composició habitual d'aplicacions, $(g\comp f)(x)=g(f(x))$;
\item identitats: l'aplicació identitat, $\id_A(x)=x$.
\end{itemize}
L'associativitat i les lleis d'identitat són propietats ben conegudes de les aplicacions, de manera que $\cat{Set}$ és una categoria. Serà una categoria de referència: quan un concepte categòric nou ens desconcerti, mirar què vol dir a $\cat{Set}$ sol ser un bon primer pas.

Si ens restringim als conjunts finits i a les aplicacions entre ells, obtenim la categoria $\cat{FinSet}$: $\Ob(\cat{FinSet})$ és la col·lecció dels conjunts finits, $\Hom_{\cat{FinSet}}(A,B)=\Hom_{\cat{Set}}(A,B)$, i la composició i les identitats són les mateixes que a $\cat{Set}$. Com que la composició d'aplicacions entre conjunts finits torna a ser una aplicació entre conjunts finits, els axiomes se satisfan igual que a $\cat{Set}$.

Tot conjunt $S$ es pot veure també com una categoria, que continuem escrivint $S$:\index{categoria!discreta} $\Ob(S)=S$ i els únics morfismes són les identitats, una per objecte,
\[
\Hom_S(x,y)=\begin{cases}\{\id_x\} & \text{si } x=y,\\[2pt] \varnothing & \text{altrament.}\end{cases}
\]
La composició és forçada ($\id_x\comp\id_x=\id_x$) i els axiomes es compleixen trivialment. Una categoria en què els únics morfismes són les identitats s'anomena \textbf{categoria discreta}. En una categoria discreta no hi ha cap informació més enllà de la col·lecció d'objectes.

\begin{observacio}
Una categoria no queda determinada pels seus objectes: cal dir també quines són les fletxes. Amb només dos objectes, per exemple, hem trobat ja tres categories ben diferents:
\begin{itemize}
\item la \textbf{categoria discreta} de dos punts, només amb les identitats i cap fletxa no trivial;
\item la categoria $\cat{2}$, amb una única fletxa no trivial d'un objecte a l'altre;
\item la categoria amb \textbf{dues fletxes paral·leles}, amb dues fletxes no trivials entre els mateixos objectes.
\end{itemize}
Totes tres tenen els mateixos dos objectes, però són categories diferents perquè difereixen en les fletxes. En particular, $\cat{2}$ \emph{no} és la categoria discreta de dos punts: la presència de la fletxa no trivial les distingeix.
\end{observacio}

\subsection{Estructures elementals com a categories}\label{subsec:elementals}

Les estructures matemàtiques més simples sovint són categories. Per exemple, un monoide no s'assembla a una categoria d'un sol objecte: és exactament una categoria d'un sol objecte. Recollim aquesta idea en una llista d'identificacions, i tot seguit les justifiquem:
\begin{itemize}
\item un conjunt és una categoria discreta (ja vist a la subsecció anterior);
\item un monoide és una categoria d'un sol objecte;
\item un grup és una categoria d'un sol objecte en què tot morfisme és un isomorfisme;
\item un preordre és una categoria prima;
\item un ordre parcial és una categoria prima en què els únics isomorfismes són les identitats;
\item un sistema deductiu és una categoria prima.
\end{itemize}
En tots els casos, els objectes i els morfismes deixen de ser conjunts i aplicacions: passen a ser els elements i les relacions internes de l'estructura.

\paragraph{Preordres.}\index{preordre}
Un \emph{preordre} és un parell $(P,\le)$ on $P$ és un conjunt i $\le$ una relació binària sobre $P$ reflexiva ($x\le x$) i transitiva (si $x\le y$ i $y\le z$, aleshores $x\le z$). Per exemple, $(\mathbb{R},\le)$, o els subconjunts d'un conjunt ordenats per la inclusió $\subseteq$.

Tot preordre $(P,\le)$ \textbf{és} una categoria, que continuem escrivint $P$:
\begin{itemize}
\item objectes: $\Ob(P)=P$, els elements del preordre;
\item morfismes: hi ha exactament una fletxa $x\to y$ quan $x\le y$,
\[
\Hom_P(x,y)=\begin{cases}\{(x,y)\} & \text{si } x\le y,\\[2pt] \varnothing & \text{altrament;}\end{cases}
\]
\item composició: $(y,z)\comp(x,y)=(x,z)$, que existeix per transitivitat;
\item identitats: $\id_x=(x,x)$, que existeix per reflexivitat.
\end{itemize}
\[
\begin{tikzpicture}[baseline=(m.center)]
  \matrix (m) [matrix of math nodes, row sep=2.6em, column sep=2.6em] {
    x & y \\
      & z \\
  };
  \draw[-{Stealth[scale=1.1]}] (m-1-1) -- node[above] {$(x,y)$} (m-1-2);
  \draw[-{Stealth[scale=1.1]}] (m-1-2) -- node[right] {$(y,z)$} (m-2-2);
  \draw[-{Stealth[scale=1.1]}] (m-1-1) -- node[below left] {$(x,z)$} (m-2-2);
\end{tikzpicture}
\]
Com que cada homconjunt té com a màxim un element, l'associativitat i les lleis d'identitat se satisfan automàticament. Una categoria en què $|\Hom(x,y)|\le 1$ per a tot parell d'objectes s'anomena \textbf{prima}\index{categoria!prima} (en anglès, \emph{thin category}).

Si el preordre és a més antisimètric ($x\le y$ i $y\le x$ impliquen $x=y$), és un \emph{ordre parcial}. Categòricament, l'antisimetria diu que els únics isomorfismes (la definició~\ref{def:isomorfisme}) són les identitats: si $x\to y$ i $y\to x$ existeixen totes dues, aleshores $x=y$.

\paragraph{Monoides.}\index{monoide}
Un \emph{monoide} és una terna $(M,\ast,e)$ on $\ast$ és una operació associativa sobre $M$ amb element neutre $e$. Per exemple, $(\mathbb{N},+,0)$, o les paraules sobre un alfabet amb la concatenació i la paraula buida.

Tot monoide \textbf{és} una categoria amb un sol objecte $\ast$, que escrivim $\cat{B}M$.\footnote{La notació $\cat{B}M$ prové de la topologia, on $BM$ designa l'\emph{espai classificador} del monoide (o grup) $M$; la categoria $\cat{B}M$ n'és el model categòric.} Concretament:
\begin{itemize}
\item objectes: $\Ob(\cat{B}M)=\{\ast\}$;
\item morfismes: l'únic homconjunt és tot el monoide, $\Hom_{\cat{B}M}(\ast,\ast)=M$;
\item composició: l'operació del monoide, $g\comp f=g\ast f$, on l'element de la dreta actua primer;
\item identitat: $\id_\ast=e$.
\end{itemize}
Els axiomes de categoria són exactament els de monoide. Els morfismes, aquí, no són aplicacions: són els elements de $M$.

\paragraph{Grups.}\index{grup}
Un \emph{grup} és una terna $(G,\ast,e)$ que és un monoide en què, a més, tot element $a$ té un invers $a^{-1}$ amb $a\ast a^{-1}=e=a^{-1}\ast a$. Per exemple, $(\mathbb{Z},+,0)$, o les permutacions d'un conjunt amb la composició.

Un grup \textbf{és} una categoria d'un sol objecte en què tot morfisme és un isomorfisme. Com que un grup és en particular un monoide, la categoria és la de l'apartat anterior, que escrivim igualment $\cat{B}G$: $\Ob(\cat{B}G)=\{\ast\}$, $\Hom_{\cat{B}G}(\ast,\ast)=G$, $g\comp f=g\ast f$ i $\id_\ast=e$. En efecte, l'invers de l'element $a$ és exactament el morfisme invers de la fletxa corresponent. La teoria de grups és, doncs, la part invertible de la teoria de monoides. Cada element del grup és un bucle sobre l'únic objecte, i la identitat $e$ n'és el bucle destacat:
\[
\begin{tikzpicture}[baseline=(o.base)]
  \node (o) {$\ast$};
  % pètals genèrics (elements del grup)
  \foreach \ang in {30,150,210,270,330}{
    \draw[-{Stealth[scale=0.9]},gray] (o) to[out=\ang+14,in=\ang-14,looseness=16] (o);
  }
  % pètal destacat: la identitat e
  \draw[-{Stealth[scale=1.0]},thick] (o) to[out=104,in=76,looseness=16] (o);
  \node at (0,1.32) {$e$};
\end{tikzpicture}
\]

\paragraph{Sistemes deductius.}\index{deduibilitat@deduïbilitat}\index{Prov@$\cat{Prov}_T$}
Fixem un \emph{sistema deductiu} $T$ i considerem la seva relació de deduïbilitat $\vdash_T$ entre fórmules, reflexiva ($p\vdash_T p$) i transitiva (de $p\vdash_T q$ i $q\vdash_T r$ se segueix $p\vdash_T r$). Aquesta relació determina una categoria, que escrivim $\cat{Prov}_T$ per distingir-la del sistema $T$ del qual prové:
\begin{itemize}
\item objectes: $\Ob(\cat{Prov}_T)$ és el conjunt de les fórmules $p,q,r,\dots$ de $T$;
\item morfismes: hi ha exactament una fletxa $p\to q$ quan $p\vdash_T q$,
\[
\Hom_{\cat{Prov}_T}(p,q)=\begin{cases}\{(p,q)\} & \text{si } p\vdash_T q,\\[2pt] \varnothing & \text{altrament;}\end{cases}
\]
\item composició: $(q,r)\comp(p,q)=(p,r)$, que existeix per transitivitat;
\item identitats: $\id_p=(p,p)$, que existeix per reflexivitat.
\end{itemize}
\[
\begin{tikzpicture}[baseline=(m.center)]
  \matrix (m) [matrix of math nodes, row sep=2.6em, column sep=2.6em] {
    p & q \\
      & r \\
  };
  \draw[-{Stealth[scale=1.1]}] (m-1-1) -- node[above] {$p\vdash_T q$} (m-1-2);
  \draw[-{Stealth[scale=1.1]}] (m-1-2) -- node[right] {$q\vdash_T r$} (m-2-2);
  \draw[-{Stealth[scale=1.1]}] (m-1-1) -- node[below left] {$p\vdash_T r$} (m-2-2);
\end{tikzpicture}
\]
Com el preordre, $\cat{Prov}_T$ és una categoria prima. A diferència de l'ordre parcial, però, la deduïbilitat no acostuma a ser antisimètrica: dues fórmules diferents poden ser interdeduïbles ($p\vdash_T q$ i $q\vdash_T p$ amb $p\neq q$), com $p\land q$ i $q\land p$. Categòricament, això vol dir que hi ha isomorfismes no trivials: fórmules distintes però isomorfes. Aquesta categoria només registra \emph{si} $q$ es dedueix de $p$, no \emph{de quina manera} es dedueix; l'observació següent presenta la versió que distingeix les proves.

\begin{observacio}[Les proves mateixes com a fletxes]\index{prova!com a morfisme}
Hi ha una interpretació lògica més fina en què una prova concreta
\[
\pi\colon p\longrightarrow q
\]
és ella mateixa un morfisme, i dues proves $\pi,\rho\colon p\to q$ poden ser morfismes diferents. La composició representa aleshores l'encadenament de proves i la identitat, la prova trivial d'una fórmula a partir d'ella mateixa. Passem, doncs, d'una sola fletxa (registrar \emph{si} hi ha deducció) a moltes fletxes (registrar \emph{quina} prova).

Per obtenir una categoria de proves completament definida cal fixar el sistema deductiu i precisar també quan dues derivacions es consideren iguals ---per exemple, mitjançant una equivalència adequada entre proves. No necessitem encara aquesta teoria; de moment distingirem entre la categoria prima $\cat{Prov}_T$ de la \emph{deduïbilitat} i aquesta interpretació més rica de les \emph{proves com a fletxes}. Hi tornarem quan el llenguatge categòric ens permeti fer-ho amb més precisió.
\end{observacio}

\subsection{Categories de conjunts estructurats}

Ara considerem les categories els objectes de les quals són conjunts amb alguna estructura, i els morfismes, les aplicacions que preserven aquesta estructura. En tots els casos la composició és la d'aplicacions i les identitats són les aplicacions identitat, de manera que l'associativitat i les lleis d'identitat s'hereten de les de $\cat{Set}$; només cal comprovar, en cada cas, que la composició d'aplicacions que respecten l'estructura torna a respectar-la, i que la identitat també ho fa.

\paragraph{Grups.}
La categoria $\cat{Grp}$ té:
\begin{itemize}
\item objectes: $\Ob(\cat{Grp})$ és la col·lecció de tots els grups;
\item morfismes: $\Hom_{\cat{Grp}}(G,H)$ és el conjunt dels homomorfismes de grups $G\to H$;
\item composició: la composició d'aplicacions;
\item identitats: l'aplicació identitat $\id_{G}$.
\end{itemize}
Un \emph{homomorfisme de grups} $h\colon G\to H$ és una aplicació entre els conjunts subjacents que preserva l'operació: $h(a\ast b)=h(a)\ast h(b)$ per a tot $a,b$ (d'on se segueix que preserva també el neutre i els inversos). Si $h\colon G\to H$ i $k\colon H\to K$ són homomorfismes, la composició també ho és, ja que
\[
k(h(a\ast b))=k(h(a)\ast h(b))=k(h(a))\ast k(h(b));
\]
i l'aplicació identitat preserva trivialment l'operació. Com que la composició i les identitats són les d'aplicacions entre conjunts, l'associativitat i les lleis d'identitat s'hereten de $\cat{Set}$.

\paragraph{Espais vectorials.}\index{espai vectorial}
Fixem un cos $\mathbb{K}$. La categoria $\cat{Vect}_{\mathbb K}$ té:
\begin{itemize}
\item objectes: $\Ob(\cat{Vect}_{\mathbb K})$ és la col·lecció dels espais vectorials sobre $\mathbb{K}$;
\item morfismes: $\Hom_{\mathbb K}(V,W)$ és el conjunt de les aplicacions lineals $V\to W$;
\item composició: la composició d'aplicacions;
\item identitats: l'aplicació identitat $\id_V$.
\end{itemize}
Un \emph{espai vectorial} sobre $\mathbb{K}$ és un conjunt $V$ amb una suma commutativa i associativa $V\times V\to V$ ---amb element neutre i oposats--- i una multiplicació per escalars $\mathbb{K}\times V\to V$ compatible amb les operacions. Una \emph{aplicació lineal} $f\colon V\to W$ és una aplicació entre els conjunts subjacents que preserva l'estructura: $f(u+v)=f(u)+f(v)$ i $f(\lambda v)=\lambda f(v)$. Si $f\colon V\to W$ i $g\colon W\to U$ són lineals, la composició també ho és, ja que
\[
g(f(u+v))=g(f(u))+g(f(v))\qquad\text{i}\qquad g(f(\lambda v))=\lambda\,g(f(v));
\]
i la identitat és lineal. Per tant la composició i les identitats són les d'aplicacions entre conjunts, i l'associativitat i les lleis d'identitat s'hereten de $\cat{Set}$.

\paragraph{Espais topològics.}\index{espai topològic}
La categoria $\cat{Top}$ té:
\begin{itemize}
\item objectes: $\Ob(\cat{Top})$ és la col·lecció de tots els espais topològics;
\item morfismes: $\Hom_{\cat{Top}}(X,Y)$ és el conjunt de les aplicacions contínues $X\to Y$;
\item composició: la composició d'aplicacions;
\item identitats: l'aplicació identitat $\id_{X}$.
\end{itemize}
Un \emph{espai topològic} és un parell $(X,\mathcal{T}_X)$ on $X$ és un conjunt i $\mathcal{T}_X$ una família de subconjunts de $X$ (els \emph{oberts}) tal que:
\begin{itemize}
\item $\varnothing, X\in\mathcal{T}_X$;
\item si $U,V\in\mathcal{T}_X$, aleshores $U\cap V\in\mathcal{T}_X$;
\item si $\{U_i\}_{i\in I}$ és una família qualsevol en $\mathcal{T}_X$, aleshores $\bigcup_{i\in I}U_i\in\mathcal{T}_X$.
\end{itemize}
Una \emph{aplicació contínua} $f\colon(X,\mathcal{T}_X)\to(Y,\mathcal{T}_Y)$ és una aplicació $f\colon X\to Y$ que preserva l'estructura en el sentit següent: $f^{-1}(U)\in\mathcal{T}_X$ per a tot $U\in\mathcal{T}_Y$. Si $f\colon X\to Y$ i $g\colon Y\to Z$ són contínues, la composició també ho és, ja que $(g\comp f)^{-1}(U)=f^{-1}(g^{-1}(U))$ és obert quan $U$ ho és; i la identitat és contínua, perquè $\id^{-1}(U)=U$. Per tant, un cop més, la composició i les identitats són les d'aplicacions entre conjunts, i l'associativitat i les lleis d'identitat s'hereten de $\cat{Set}$.

\paragraph{Monoides.}
La categoria $\cat{Mon}$ té:
\begin{itemize}
\item objectes: $\Ob(\cat{Mon})$ és la col·lecció de tots els monoides;
\item morfismes: $\Hom_{\cat{Mon}}(M,N)$ és el conjunt dels homomorfismes de monoides $M\to N$;
\item composició: la composició d'aplicacions;
\item identitats: l'aplicació identitat $\id_{M}$.
\end{itemize}
Un \emph{homomorfisme de monoides} $h\colon M\to N$ és una aplicació entre els conjunts subjacents que preserva l'operació i el neutre: $h(a\ast b)=h(a)\ast h(b)$ i $h(e_M)=e_N$. Si $h\colon M\to N$ i $k\colon N\to P$ són homomorfismes, la composició també ho és, ja que
\[
k(h(a\ast b))=k(h(a)\ast h(b))=k(h(a))\ast k(h(b))\qquad\text{i}\qquad k(h(e_M))=k(e_N)=e_P;
\]
i la identitat preserva trivialment l'operació i el neutre. Com que la composició i les identitats són les d'aplicacions entre conjunts, l'associativitat i les lleis d'identitat s'hereten de $\cat{Set}$.

\paragraph{Ordres parcials.}
La categoria $\cat{Pos}$ té:
\begin{itemize}
\item objectes: $\Ob(\cat{Pos})$ és la col·lecció de tots els ordres parcials;
\item morfismes: $\Hom_{\cat{Pos}}(P,Q)$ és el conjunt de les aplicacions monòtones $P\to Q$;
\item composició: la composició d'aplicacions;
\item identitats: l'aplicació identitat $\id_{P}$.
\end{itemize}
Una \emph{aplicació monòtona} $f\colon P\to Q$ és una aplicació entre els conjunts subjacents que preserva l'ordre: si $x\le y$, aleshores $f(x)\le f(y)$. Si $f\colon P\to Q$ i $g\colon Q\to R$ són monòtones, la composició també ho és, ja que
\[
x\le y \ \Longrightarrow\ f(x)\le f(y) \ \Longrightarrow\ g(f(x))\le g(f(y));
\]
i la identitat preserva trivialment l'ordre. Un cop més, la composició i les identitats són les d'aplicacions entre conjunts, i les lleis s'hereten de $\cat{Set}$.

Convé no confondre les dues mirades: un monoide concret \emph{és} una categoria (d'un sol objecte) i un ordre parcial concret \emph{és} una categoria (prima), mentre que $\cat{Mon}$ i $\cat{Pos}$ tenen \emph{tots} els monoides i \emph{tots} els ordres parcials per objectes, respectivament.

\begin{observacio}[Categories com a contextos matemàtics]
Cadascuna d'aquestes categories és l'\emph{àmbit} on es desenvolupa una teoria matemàtica: $\cat{Top}$ és el context de la topologia general, $\cat{Grp}$ el de la teoria de grups, $\cat{Vect}_{\mathbb K}$ el de l'àlgebra lineal, i així successivament. Aquí, la categoria la formen tots els objectes d'una mena; a les estructures elementals, en canvi, un sol objecte estructurat ja és una categoria.
\end{observacio}

\subsection{Conjunts i relacions}\index{Rel@$\cat{Rel}$}

Com hem avançat a la introducció (secció~\ref{sec:mirar-fora}), els conjunts admeten una segona categoria, on els morfismes no són aplicacions sinó \emph{relacions}. La categoria $\cat{Rel}$ té:
\begin{itemize}
\item objectes: $\Ob(\cat{Rel})$ és la col·lecció de tots els conjunts;
\item morfismes: $\Hom_{\cat{Rel}}(X,Y)$ és el conjunt de les relacions $R\subseteq X\times Y$ (amb la precisió sobre els extrems de l'observació~\ref{obs:extrems-implicits});
\item composició: per a $R\subseteq X\times Y$ i $S\subseteq Y\times Z$,
\[
S\comp R=\{(x,z)\mid \exists y\in Y.\ (x,y)\in R \ \text{i}\ (y,z)\in S\}\subseteq X\times Z;
\]
\item identitats: la relació diagonal $\id_X=\{(x,x)\mid x\in X\}$.
\end{itemize}
Els axiomes es comproven directament a partir de la definició. Per a l'associativitat, un parell $(x,w)$ pertany a les dues maneres d'agrupar tres relacions encadenades exactament quan hi ha termes intermedis que connecten $x$ amb $w$ passant per les tres relacions; com que la condició no depèn de l'agrupament, les dues composicions coincideixen. Per a les identitats, compondre amb la diagonal no afegeix cap pas efectiu, perquè obliga el terme intermedi a coincidir amb l'extrem. La verificació completa és el primer exercici resolt del capítol~1 de la part de Pràctica (exercici~\ref{exr:pr-rel}). Compondre relacions equival a aplicar un quantificador existencial sobre el terme intermedi: aquest és el primer pont amb la lògica que assenyalàvem a la introducció. A diferència del que passa a $\cat{Set}$, un morfisme $X\to Y$ de $\cat{Rel}$ no ha d'assignar necessàriament a cada element de $X$ un únic element de $Y$: pot relacionar-lo amb diversos elements de $Y$, o amb cap.

\begin{observacio}[Extrems implícits]\label{obs:extrems-implicits}
En rigor, un morfisme $X\to Y$ de $\cat{Rel}$ no és una relació $R$ tota sola, sinó la relació \emph{junt amb} els seus extrems: la terna $(X,Y,R)$ amb $R\subseteq X\times Y$. Aquesta precisió és necessària perquè els homconjunts siguin disjunts, tal com exigeix la definició de categoria. Per exemple, la relació buida $R=\varnothing$ està continguda en $X\times Y$ per a qualsevol parell $(X,Y)$; si no registréssim els extrems, un mateix objecte $\varnothing$ seria alhora morfisme entre infinites parelles d'objectes, amb domini i codomini indeterminats. Registrant $(X,Y,R)$, cada morfisme té un únic domini $X$ i un únic codomini $Y$. Escriurem simplement $R$ quan els extrems se sobreentenguin; és la mateixa convenció d'\emph{extrems implícits} que ja hem fet servir en codificar la fletxa d'un preordre o de $\cat{Prov}_T$ com el parell $(x,y)$, i que retrobarem tot seguit a les categories sobre i sota un objecte. No canvia cap de les construccions: només n'explicita la compatibilitat amb la definició inicial de categoria.
\end{observacio}

\subsection{Categories de grafs}\label{subsec:grafs}

\begin{definicio}\index{graf dirigit}
Un \textbf{graf dirigit} $G$ consta d'un conjunt $V_G$ de \textbf{vèrtexs}, un conjunt $E_G$ d'\textbf{arestes} i dues aplicacions
\[
s_G,t_G\colon E_G\longrightarrow V_G,
\]
que assignen a cada aresta el seu \textbf{origen} i el seu \textbf{destí}, respectivament.
\end{definicio}

\begin{definicio}\index{morfisme de grafs}
Un \textbf{morfisme de grafs} $F\colon G\to H$ és una parella d'aplicacions
\[
F_V\colon V_G\to V_H,
\qquad
F_E\colon E_G\to E_H,
\]
que preserven origen i destí:
\[
F_V\comp s_G=s_H\comp F_E,
\qquad
F_V\comp t_G=t_H\comp F_E.
\]
Equivalentment, han de commutar els dos quadrats següents:

\begin{center}
\begin{tikzpicture}[baseline=(m.center)]
  \matrix (m) [matrix of math nodes, row sep=2.5em, column sep=3.2em] {
    E_G & V_G \\
    E_H & V_H \\
  };
  \draw[-{Stealth[scale=1.0]}] (m-1-1) -- node[above] {$s_G$} (m-1-2);
  \draw[-{Stealth[scale=1.0]}] (m-1-1) -- node[left] {$F_E$} (m-2-1);
  \draw[-{Stealth[scale=1.0]}] (m-1-2) -- node[right] {$F_V$} (m-2-2);
  \draw[-{Stealth[scale=1.0]}] (m-2-1) -- node[below] {$s_H$} (m-2-2);
\end{tikzpicture}
\hspace{2.2em}
\begin{tikzpicture}[baseline=(n.center)]
  \matrix (n) [matrix of math nodes, row sep=2.5em, column sep=3.2em] {
    E_G & V_G \\
    E_H & V_H \\
  };
  \draw[-{Stealth[scale=1.0]}] (n-1-1) -- node[above] {$t_G$} (n-1-2);
  \draw[-{Stealth[scale=1.0]}] (n-1-1) -- node[left] {$F_E$} (n-2-1);
  \draw[-{Stealth[scale=1.0]}] (n-1-2) -- node[right] {$F_V$} (n-2-2);
  \draw[-{Stealth[scale=1.0]}] (n-2-1) -- node[below] {$t_H$} (n-2-2);
\end{tikzpicture}
\end{center}
\end{definicio}

Els grafs dirigits i els morfismes de grafs formen una categoria, que denotarem $\cat{Graph}$:\index{Graph@$\cat{Graph}$}
\begin{itemize}
\item objectes: $\Ob(\cat{Graph})$ és la col·lecció de tots els grafs dirigits;
\item morfismes: $\Hom_{\cat{Graph}}(G,H)$ és el conjunt dels morfismes de grafs $G\to H$;
\item composició: component a component; per a $F\colon G\to H$ i $F'\colon H\to K$,
\[
F'\comp F=(F'_V\comp F_V,\ F'_E\comp F_E);
\]
\item identitats: $\id_G=(\id_{V_G},\id_{E_G})$.
\end{itemize}
La composició torna a ser un morfisme de grafs, perquè les equacions es conserven:
\[
(F'_V\comp F_V)\comp s_G=F'_V\comp s_H\comp F_E=s_K\comp(F'_E\comp F_E),
\]
i anàlogament per a $t$. L'associativitat i les lleis d'identitat es compleixen component a component, perquè són les de $\cat{Set}$.

Aquest exemple és especialment rellevant perquè els grafs són l'esquelet combinatori dels diagrames: tenen vèrtexs i arestes dirigides, però encara no incorporen per si sols identitats, composicions ni axiomes categòrics.

Per fixar idees, vet aquí un graf dirigit concret amb tres vèrtexs $V_G=\{a,b,c\}$ i tres arestes:
\[
\begin{tikzpicture}[baseline=(m.center)]
  \matrix (m) [matrix of math nodes, row sep=2.6em, column sep=2.6em] {
    a & b \\
      & c \\
  };
  \draw[-{Stealth[scale=1.1]}] (m-1-1) -- node[above] {$e_1$} (m-1-2);
  \draw[-{Stealth[scale=1.1]}] (m-1-2) -- node[right] {$e_2$} (m-2-2);
  \draw[-{Stealth[scale=1.1]}] (m-1-1) -- node[below left] {$e_3$} (m-2-2);
\end{tikzpicture}
\]
Aquí $e_1$ va de $a$ a $b$, $e_2$ de $b$ a $c$ i $e_3$ de $a$ a $c$. A diferència d'una categoria, un graf no demana que hi hagi cap aresta que \enquote{representi la composició} d'$e_1$ i $e_2$: $e_3$ hi és perquè l'hem posada, no perquè cap axioma l'obligui. Aquesta és exactament la diferència entre un graf i una categoria: la categoria afegeix identitats, composició i els axiomes que les regeixen. Per a un desenvolupament complet dels grafs, amb nombrosos exemples i els homomorfismes de grafs, vegeu l'apèndix~\ref{cap:grafs}.

\subsection{Resum: objectes i morfismes dels exemples}

La taula següent recull, per a cada exemple d'aquesta secció, quines dades constitueixen els objectes i quines els morfismes.

\begin{center}
\footnotesize
\setlength{\tabcolsep}{6pt}
\begin{tabularx}{\linewidth}{@{}>{\raggedright\arraybackslash}X>{\raggedright\arraybackslash}X>{\raggedright\arraybackslash}X@{}}
\toprule
Categoria & Objectes & Morfismes \\
\midrule
$\cat{0},\cat{1},\cat{2},\cat{3}$ & marcadors abstractes & fletxes formals \\
$\cat{Set}$, $\cat{FinSet}$ & conjunts & aplicacions \\
categoria discreta & marcadors abstractes & només identitats \\
$\cat{Grp}$, $\cat{Vect}_{\mathbb K}$, $\cat{Top}$, $\cat{Mon}$, $\cat{Pos}$ & conjunts amb estructura & aplicacions que preserven l'estructura \\
$\cat{B}M$ (un monoide) & un únic objecte $\ast$ & elements de $M$ \\
un preordre $P$; $\cat{Prov}_T$ & elements de $P$; fórmules de $T$ & $x\le y$; $p\vdash_T q$ \\
$\cat{Rel}$ & conjunts & relacions \\
$\cat{Graph}$ & grafs $(V,E,s,t)$ & parelles compatibles d'aplicacions \\
\bottomrule
\end{tabularx}
\end{center}

Aquesta taula fa visible que la definició de categoria no pressuposa res sobre la naturalesa dels objectes ni dels morfismes: no cal que els objectes siguin conjunts, ni que els morfismes siguin aplicacions. En una categoria discreta, per exemple, les identitats són morfismes \emph{formals}: presentar-les com a aplicacions pressuposaria una realització concreta que no forma part de la definició. I a $\cat{B}M$ els morfismes són els elements de $M$, que en general tampoc no són aplicacions. La categoria $\cat{Rel}$ ho remarca des de l'altre costat: els seus objectes són conjunts, però els seus morfismes ---les relacions--- no són funcions.

Convé no confondre aquesta descripció \emph{estructural} de les dades amb la seva eventual \emph{codificació} conjuntista. En el marc de fonaments que adoptem, qualsevol d'aquestes dades ---inclosos els marcadors abstractes--- es pot representar per conjunts; però aquesta codificació és una tria externa, no una part de la definició de categoria. Per això la taula descriu què \emph{són} les dades i no si \enquote{són conjunts} o \enquote{són funcions}: aquesta darrera pregunta barreja la naturalesa estructural amb una representació particular.

\section{Isomorfismes}\index{isomorfisme}

Entre tots els morfismes d'una categoria, n'hi ha uns d'especialment importants: els que tenen invers. Capturen, en el llenguatge abstracte de les fletxes, la idea que dos objectes tenen la mateixa estructura des del punt de vista de la categoria.

\begin{definicio}\label{def:isomorfisme}
En una categoria $\cat{C}$, un morfisme $f\colon A\to B$ és un \textbf{isomorfisme} si existeix un morfisme $g\colon B\to A$ tal que
\[
g\comp f=\id_A
\qquad\text{i}\qquad
f\comp g=\id_B.
\]
El morfisme $g$ s'anomena \textbf{invers} de $f$. Si existeix un isomorfisme entre $A$ i $B$, diem que $A$ i $B$ són \textbf{isomorfs}, i escrivim $A\cong B$.
\end{definicio}

La situació es caracteritza per dues equacions, $g\comp f=\id_A$ i $f\comp g=\id_B$, que es llegeixen com dos triangles commutatius:
\[
\begin{tikzpicture}[baseline=(m.center)]
  \matrix (m) [matrix of math nodes, row sep=2.6em, column sep=2.6em] {
    A & B \\
    A & \\
  };
  \draw[-{Stealth[scale=1.1]}] (m-1-1) -- node[above] {$f$} (m-1-2);
  \draw[-{Stealth[scale=1.1]}] (m-1-2) -- node[right] {$g$} (m-2-1);
  \draw[-{Stealth[scale=1.1]}] (m-1-1) -- node[left] {$\id_A$} (m-2-1);
\end{tikzpicture}
\qquad\qquad
\begin{tikzpicture}[baseline=(m.center)]
  \matrix (m) [matrix of math nodes, row sep=2.6em, column sep=2.6em] {
    B & A \\
    B & \\
  };
  \draw[-{Stealth[scale=1.1]}] (m-1-1) -- node[above] {$g$} (m-1-2);
  \draw[-{Stealth[scale=1.1]}] (m-1-2) -- node[right] {$f$} (m-2-1);
  \draw[-{Stealth[scale=1.1]}] (m-1-1) -- node[left] {$\id_B$} (m-2-1);
\end{tikzpicture}
\]
Cada triangle diu que anar d'un objecte a l'altre amb $f$ o $g$ i tornar és el mateix que seguir la identitat.

\begin{proposicio}
L'invers d'un isomorfisme, si existeix, és únic.
\end{proposicio}

\begin{prova}
Suposem que $f\colon A\to B$ té dos inversos, $g$ i $g'$. Aleshores
\[
g=g\comp\id_B=g\comp(f\comp g')=(g\comp f)\comp g'=\id_A\comp g'=g'.
\]
Per tant $g=g'$. Com que l'invers és únic, el podem denotar sense ambigüitat $f^{-1}$.
\end{prova}

\begin{proposicio}
Sigui $\cat{C}$ una categoria.
\begin{enumerate}
\item Per a cada objecte $A$, la identitat $\id_A$ és un isomorfisme, amb $\id_A^{-1}=\id_A$.
\item Si $f$ és un isomorfisme, aleshores $f^{-1}$ també ho és, i $(f^{-1})^{-1}=f$.
\item Si $f\colon A\to B$ i $g\colon B\to C$ són isomorfismes, aleshores $g\comp f$ també ho és, i
\[
(g\comp f)^{-1}=f^{-1}\comp g^{-1}.
\]
\end{enumerate}
\end{proposicio}

\begin{prova}
(1) De $\id_A\comp\id_A=\id_A$ se segueix que $\id_A$ és el seu propi invers. (2) Les equacions $g\comp f=\id_A$ i $f\comp g=\id_B$ que fan de $g=f^{-1}$ l'invers de $f$ diuen també que $f$ és l'invers de $g$; per tant $f^{-1}$ és un isomorfisme i $(f^{-1})^{-1}=f$. (3) Comprovem que $f^{-1}\comp g^{-1}$ és l'invers de $g\comp f$:
\[
(g\comp f)\comp(f^{-1}\comp g^{-1})=g\comp(f\comp f^{-1})\comp g^{-1}=g\comp\id_B\comp g^{-1}=g\comp g^{-1}=\id_C,
\]
i anàlogament $(f^{-1}\comp g^{-1})\comp(g\comp f)=\id_A$.
\end{prova}

\begin{corollari}
Donada una categoria $\cat{C}$, definim sobre $\Ob(\cat{C})$ la relació $\cong_{\cat{C}}$ per: $A\cong_{\cat{C}}B$ si i només si existeix un isomorfisme $f\colon A\to B$. Aleshores $\cong_{\cat{C}}$ és una relació d'equivalència.
\end{corollari}

\begin{prova}
És reflexiva: $\id_A$ és un isomorfisme per (1), de manera que $A\cong_{\cat{C}}A$. És simètrica: si $f$ és un isomorfisme de $A$ a $B$, aleshores $f^{-1}$ ho és de $B$ a $A$ per (2), de manera que $B\cong_{\cat{C}}A$. És transitiva: si $f\colon A\to B$ i $g\colon B\to C$ són isomorfismes, aleshores $g\comp f$ també ho és per (3), de manera que $A\cong_{\cat{C}}C$.
\end{prova}

\begin{observacio}[Igualtat i isomorfisme]\label{obs:igualtat-isomorfisme}\index{isomorfisme!i igualtat}
Convé fixar des d'ara una convenció que el llibre manté sempre. Escrivim $A=B$ només quan $A$ i $B$ són \emph{el mateix} objecte, i $A\cong B$ quan hi ha un isomorfisme entre ells. Quan un isomorfisme és distingit ---no depèn de cap tria---, l'anomenem \textbf{canònic}, però continua sent un isomorfisme i no una igualtat (en veurem un primer cas en el monoide lliure, exemple~\ref{ex:monoide-lliure-cat}).

La relació $\cong_{\cat{C}}$ permet parlar d'una \emph{igualtat estructural relativa a $\cat{C}$}: dos objectes isomorfs a $\cat{C}$ no es poden distingir per cap propietat formulada només amb els morfismes, la composició i les identitats de $\cat{C}$, perquè l'isomorfisme la transporta d'un a l'altre (ho precisarem a l'observació~\ref{obs:invariancia}). Diem, doncs, que tenen la mateixa estructura \emph{des del punt de vista de $\cat{C}$}, i per això sovint la pregunta rellevant no és si $A$ i $B$ són iguals, sinó si són isomorfs a $\cat{C}$.

Aquesta igualtat no és absoluta: depèn dels morfismes que la categoria considera admissibles. Per exemple, el conjunt $\{0,1\}$ amb l'ordre discret i el mateix conjunt amb l'ordre $0<1$ no són isomorfs a $\cat{Pos}$, perquè cap bijecció monòtona no té inversa monòtona, però els seus conjunts subjacents sí que són isomorfs a $\cat{Set}$. Canviar de categoria canvia, doncs, quins objectes compten com a estructuralment iguals. Aquesta és una de les idees centrals de la teoria de categories: la forma d'un objecte sempre es mira dins d'una categoria determinada.
La jerarquia completa ---igualtat, isomorfisme, isomorfisme natural i equivalència---, amb un contraexemple per a cada pas, es tracta a l'apèndix~\refalt{cap:contraexemples}{de contraexemples del llibre complet}.
\end{observacio}

\begin{exemple}
A $\cat{Set}$, els isomorfismes són exactament les aplicacions bijectives. En efecte, una aplicació té inversa per la composició si i només si és injectiva i exhaustiva.
\end{exemple}

\begin{exemple}[Isomorfismes en les categories de conjunts estructurats]
A $\cat{Grp}$, $\cat{Vect}_{\mathbb K}$ i $\cat{Mon}$, els isomorfismes són les aplicacions bijectives que preserven l'estructura i la inversa de les quals també la preserva; recuperem així les nocions habituals d'isomorfisme de grups, d'espais vectorials i de monoides. A $\cat{Top}$, en canvi, un isomorfisme és un homeomorfisme\index{homeomorfisme}: una aplicació contínua i bijectiva amb inversa contínua. Aquí es veu que la bijectivitat no basta: hi ha aplicacions contínues i bijectives que no són isomorfismes, perquè la seva inversa no és contínua. La definició categòrica capta la noció correcta sense haver-la d'enunciar cas per cas.
\end{exemple}

\begin{exemple}
En un preordre, un isomorfisme entre $x$ i $y$ vol dir que hi ha fletxes $x\to y$ i $y\to x$, és a dir,
\[
x\le y
\qquad\text{i}\qquad
y\le x.
\]
En un ordre parcial això obliga $x=y$. En un preordre general, en canvi, poden existir objectes isomorfs diferents: són exactament els elements equivalents per la relació $x\sim y$ si $x\le y$ i $y\le x$.
\end{exemple}

\begin{exemple}[Isomorfisme i equivalència lògica]\label{ex:iso-equivalencia-logica}
A la categoria prima $\cat{Prov}_T$ d'un sistema deductiu $T$, dues fórmules $p$ i $q$ són isomorfes exactament quan
\[
p\vdash_T q
\qquad\text{i}\qquad
q\vdash_T p.
\]
Per tant, l'isomorfisme categòric es tradueix aquí en \textbf{interdeduïbilitat}. Aquest és un primer exemple de com una noció categòrica adquireix una lectura lògica natural.
\end{exemple}

\begin{observacio}
L'exemple dels preordres mostra per què la definició categòrica és la bona. Es podria pensar a definir un isomorfisme com un morfisme \enquote{bijectiu}, però això requeriria parlar dels elements dels objectes, i no sempre té sentit. A més, fins i tot quan en té, pot donar la noció equivocada: a $\cat{Pos}$ existeixen aplicacions monòtones bijectives que no són isomorfismes perquè la seva inversa no és monòtona. La definició mitjançant composició i identitats evita aquest parany.
\end{observacio}

\begin{observacio}[Propietats invariants per isomorfisme]\label{obs:invariancia}
Sigui $\cat{C}$ una categoria. Una propietat dels objectes de $\cat{C}$ és \textbf{invariant per isomorfisme} a $\cat{C}$ si, sempre que un objecte la té, tot objecte isomorf a ell a $\cat{C}$ també la té. Com la relació $\cong_{\cat{C}}$ de l'observació anterior, aquesta noció és sempre relativa a una categoria determinada. Una propietat definida només a partir de l'estructura categòrica de $\cat{C}$ ---els seus objectes, morfismes, composició i identitats, sense fer servir la igualtat entre objectes--- ha de ser invariant per isomorfisme a $\cat{C}$, perquè un isomorfisme transporta qualsevol afirmació d'aquesta mena d'un objecte a l'altre. Per això la invariància per isomorfisme és un criteri bàsic per reconèixer quan una propietat depèn de la forma de l'objecte en la categoria i no de la seva presentació concreta.

Aquesta noció dona un mètode útil. Per provar que dos objectes \emph{són} isomorfs, n'hi ha prou d'exhibir un isomorfisme entre ells. Provar que \emph{no} ho són sol ser més difícil, perquè caldria descartar tots els isomorfismes possibles; la via pràctica és trobar una propietat invariant per isomorfisme a la categoria en qüestió que un dels objectes tingui i l'altre no. Per exemple, a $\cat{Pos}$, siguin $A$ la cadena $a<b<c$ i $B$ l'ordre amb $x<y$ i $x<z$ però $y,z$ incomparables. Comprovar directament que no són isomorfs exigiria descartar totes les bijeccions monòtones amb inversa monòtona; en canvi, n'hi ha prou d'observar que \emph{estar totalment ordenat} és una propietat invariant per isomorfisme, que $A$ té i $B$ no. Per tant $A$ i $B$ no són isomorfs.

La implicació inversa no és automàtica: una propietat pot ser invariant per isomorfisme sense haver estat formulada internament en el llenguatge categòric. L'exemple anterior n'és una mostra: \emph{estar totalment ordenat} és invariant per isomorfisme a $\cat{Pos}$, però l'hem definit parlant dels elements de l'ordre, no de fletxes.
\end{observacio}

\begin{observacio}
La noció d'isomorfisme és autènticament categòrica: es formula només amb composició i identitats, de manera que té sentit a \emph{qualsevol} categoria, sense referència als elements dels objectes. Això la fa més abstracta i més general que les nocions d'isomorfisme pròpies de cada context ---bijeccions, isomorfismes de grups, homeomorfismes---, que en són casos particulars.

Aquesta generalitat es fa visible en un cas que ja coneixem. Havíem identificat els monoides amb les categories d'un sol objecte; ara podem identificar els grups amb exactament aquelles categories d'un sol objecte en què tota fletxa és un isomorfisme. Portant la mateixa idea a un nombre qualsevol d'objectes s'obté una noció d'importància considerable en les matemàtiques actuals: un \emph{grupoide} és una categoria en què tot morfisme és un isomorfisme.\index{grupoide}

El grupoide unifica dos casos que ja hem trobat. Un grup és un grupoide amb un sol objecte, com acabem de veure. I un preordre en què tot morfisme és un isomorfisme és, precisament, un en què $x\le y$ implica $y\le x$: una relació simètrica, és a dir, una \emph{relació d'equivalència}. Així, una relació d'equivalència no és res més que un grupoide prim: les classes d'isomorfisme dels seus objectes són les classes d'equivalència.
\end{observacio}

\begin{definicio}\index{endomorfisme}\index{automorfisme}
Un morfisme $f\colon A\to A$, amb domini i codomini iguals, s'anomena \textbf{endomorfisme} de $A$. Un endomorfisme que a més és un isomorfisme s'anomena \textbf{automorfisme} de $A$.
\end{definicio}

Els automorfismes d'un objecte $A$, amb la composició, formen un grup, el \textbf{grup d'automorfismes} de $A$, que escrivim $\Aut(A)$.\index{Aut@$\Aut$} Torna a aparèixer així la identificació d'abans: $\Aut(A)$ és el grup associat a la categoria d'un sol objecte que té els automorfismes de $A$ com a fletxes.

\section{Classes especials de morfismes}

Els isomorfismes són els morfismes invertibles: tenen un invers per tots dos costats. Sovint, però, un morfisme té invers només per un costat, o satisfà una propietat de cancel·lació més feble que la invertibilitat. Aquestes nocions ---monomorfismes, epimorfismes, seccions i retraccions--- són bàsiques. A $\cat{Set}$, els monomorfismes i els epimorfismes coincideixen amb les aplicacions injectives i exhaustives; com veurem, però, aquesta coincidència és pròpia de $\cat{Set}$ i d'algunes categories concretes, no una propietat de les categories en general.

Recordem que, a $\cat{Set}$, una aplicació $f$ és \emph{injectiva} si $f(a)=f(a')$ implica $a=a'$, i \emph{exhaustiva} si tot element del codomini té alguna antiimatge. Aquestes definicions parlen d'elements; les versions categòriques que segueixen les reformulen només amb fletxes.

\subsection{Monomorfismes i epimorfismes}\label{subsec:mono-epi}

\begin{definicio}\index{monomorfisme}\index{epimorfisme}
\label{def:mono-epi}
Un morfisme $f\colon A\to B$ és un \textbf{monomorfisme} (o \textbf{mono}), i escrivim $f\colon A\rightarrowtail B$, si es pot cancel·lar per l'esquerra: per a tota parella de morfismes $g,h\colon C\to A$,
\[
\begin{tikzpicture}[baseline=(m.center)]
  \matrix (m) [matrix of math nodes, column sep=3.4em] { C & A & B \\ };
  \draw[-{Stealth[scale=1.0]}] ([yshift=2.6pt]m-1-1.east) -- node[above] {$g$} ([yshift=2.6pt]m-1-2.west);
  \draw[-{Stealth[scale=1.0]}] ([yshift=-2.6pt]m-1-1.east) -- node[below] {$h$} ([yshift=-2.6pt]m-1-2.west);
  \draw[-{Stealth[scale=1.1]}] (m-1-2) -- node[above] {$f$} (m-1-3);
\end{tikzpicture}
\qquad
f\comp g=f\comp h\ \Longrightarrow\ g=h.
\]
Un morfisme $f\colon A\to B$ és un \textbf{epimorfisme} (o \textbf{epi}), i escrivim $f\colon A\twoheadrightarrow B$, si es pot cancel·lar per la dreta: per a tota parella $g,h\colon B\to C$,
\[
\begin{tikzpicture}[baseline=(m.center)]
  \matrix (m) [matrix of math nodes, column sep=3.4em] { A & B & C \\ };
  \draw[-{Stealth[scale=1.1]}] (m-1-1) -- node[above] {$f$} (m-1-2);
  \draw[-{Stealth[scale=1.0]}] ([yshift=2.6pt]m-1-2.east) -- node[above] {$g$} ([yshift=2.6pt]m-1-3.west);
  \draw[-{Stealth[scale=1.0]}] ([yshift=-2.6pt]m-1-2.east) -- node[below] {$h$} ([yshift=-2.6pt]m-1-3.west);
\end{tikzpicture}
\qquad
g\comp f=h\comp f\ \Longrightarrow\ g=h.
\]
\end{definicio}

\begin{observacio}[Com cal llegir aquests diagrames]\label{obs:llegir-mono-epi}
Els diagrames de la definició~\ref{def:mono-epi} \emph{no} afirmen que res commuti: són esquemes de prova, i s'han de llegir per a cada parella possible de fletxes $g,h$. Al diagrama del monomorfisme hi ha dos camins de $C$ a $B$, \enquote{primer $g$, després $f$} i \enquote{primer $h$, després $f$}, i hi ha també les dues fletxes paral·leles $g,h$ de $C$ a $A$. Recordem que un diagrama amb dues fletxes paral·leles commuta exactament quan són iguals. La definició diu, doncs:
\begin{quote}
si els dos camins de $C$ a $B$ coincideixen, aleshores les fletxes paral·leles ja coincideixen.
\end{quote}
Amb el vocabulari de la secció~\ref{sec:diagrames}: l'antecedent no és que el diagrama commuti, sinó només que commutin els camins que acaben a $B$; la conclusió és que el diagrama commuta del tot. Un monomorfisme és, doncs, una fletxa per a la qual \emph{commutar després de $f$ implica commutar}.
Llegit en sentit contrari: si $g\neq h$, aleshores $f\comp g\neq f\comp h$. Un monomorfisme \emph{no confon} mai dues fletxes diferents que arriben a $A$; després de compondre-les amb $f$, continuen sent diferents.

Cal anar amb compte amb una lectura massa ràpida. Que $f$ sigui mono no impedeix que hi hagi parelles $g\neq h\colon C\to A$; el diagrama amb dues fletxes paral·leles diferents existeix perfectament. El que exclou és que, sent diferents, es facin iguals en compondre-les amb $f$.

El diagrama de l'epimorfisme es llegeix de la mateixa manera, amb les fletxes de prova a l'altre costat: si $g\neq h\colon B\to C$, aleshores $g\comp f\neq h\comp f$. Dues fletxes que surten de $B$ i coincideixen després de $f$ ja eren iguals; $f$ arriba prou a $B$ perquè cap diferència entre $g$ i $h$ no se li escapi. A $\cat{Set}$, prenent com a fletxes de prova les aplicacions des d'un conjunt d'un sol element, que són els elements de $A$, la lectura del monomorfisme es converteix exactament en la injectivitat, com veurem tot seguit.
\end{observacio}

\begin{definicio}\index{bimorfisme}
Un morfisme que és alhora un monomorfisme i un epimorfisme s'anomena \textbf{bimorfisme}.
\end{definicio}

\begin{exemple}
A $\cat{Set}$, els monomorfismes són exactament les aplicacions \textbf{injectives} i els epimorfismes, les \textbf{exhaustives}. Vegem el cas dels monomorfismes: si $f$ és injectiva i $f\comp g=f\comp h$, aleshores per a cada $c$ tenim $f(g(c))=f(h(c))$, d'on $g(c)=h(c)$; recíprocament, si $f$ no és injectiva, hi ha $a\neq a'$ amb $f(a)=f(a')$, i les dues aplicacions constants $g,h\colon\{\ast\}\to A$ amb valors $a$ i $a'$ compleixen $f\comp g=f\comp h$ però $g\neq h$.

Per als epimorfismes, sigui $f\colon A\to B$. Si $f$ és exhaustiva i $g\comp f=h\comp f$, aleshores per a cada $b\in B$ existeix $a$ amb $f(a)=b$, i $g(b)=g(f(a))=h(f(a))=h(b)$, d'on $g=h$. Recíprocament, si $f$ no és exhaustiva, hi ha $b_0\in B$ sense antiimatge, és a dir, $b_0\notin f(A)$. Definim $g,h\colon B\to\{0,1\}$ a \emph{tot} $B$ per
\[
g(b)=0 \text{ per a tot } b\in B,
\qquad
h(b)=\begin{cases}1 & \text{si } b=b_0,\\ 0 & \text{si } b\neq b_0.\end{cases}
\]
Aleshores $g$ i $h$ coincideixen a tot $B$ tret de $b_0$, de manera que $g\neq h$. Però per a tot $a\in A$ tenim $f(a)\neq b_0$, ja que $b_0\notin f(A)$; per tant $g(f(a))=0=h(f(a))$, és a dir, $g\comp f=h\comp f$. Això mostra que $f$ no és un epimorfisme.

Convé notar que cap dels dos arguments no tria una \emph{secció} de $f$: per als monomorfismes es comparen dues seleccions d'elements des d'un conjunt unitari, i per als epimorfismes es construeixen explícitament dues aplicacions cap a $\{0,1\}$. En particular, la caracterització \enquote{epimorfisme $=$ exhaustiva} a $\cat{Set}$ no fa servir l'axioma de l'elecció, i no s'ha de confondre amb l'enunciat, ben diferent, que tota aplicació exhaustiva \emph{admet una secció} ---aquest sí equivalent a l'axioma de l'elecció (vegeu més avall).
\end{exemple}

\begin{proposicio}
Tot isomorfisme és un bimorfisme: si $f$ és un isomorfisme, aleshores és alhora un monomorfisme i un epimorfisme.
\end{proposicio}

\begin{prova}
Sigui $f\colon A\to B$ un isomorfisme amb invers $f^{-1}$. Si $f\comp g=f\comp h$, aleshores
\[
g=f^{-1}\comp(f\comp g)=f^{-1}\comp(f\comp h)=h,
\]
de manera que $f$ és mono. Si $g\comp f=h\comp f$, aleshores
\[
g=(g\comp f)\comp f^{-1}=(h\comp f)\comp f^{-1}=h,
\]
de manera que $f$ és epi.
\end{prova}

\begin{observacio}
A $\cat{Set}$, un morfisme és un isomorfisme si i només si és alhora mono i epi, és a dir, si i només si, com a aplicació, és injectiu i exhaustiu, o sigui \textbf{bijectiu}. Dit d'una altra manera, a $\cat{Set}$ els bimorfismes són exactament els isomorfismes. El recíproc de la proposició anterior, però, no és cert: en general un bimorfisme no és un isomorfisme. L'única fletxa no trivial de la categoria $\cat{2}$ és un bimorfisme ---per la simple raó que gairebé no hi ha morfismes amb què fer les cancel·lacions---, però no té invers, de manera que no és un isomorfisme. Un exemple més substancial es dona a $\cat{Top}$: una bijecció contínua és un bimorfisme, però no és un isomorfisme si la seva inversa no és contínua.
\end{observacio}

\begin{observacio}[Atenció: raonar amb elements]\label{obs:raonar-elements}\index{element!perill de raonar amb}
A $\cat{Set}$ podem raonar amb elements per dos motius: cada element $a\in A$ és una fletxa $\{\ast\}\to A$, i dues aplicacions $A\to B$ són iguals si coincideixen sobre tots els elements. En una categoria qualsevol, cap d'aquestes dues coses no està garantida.
\begin{itemize}
\item \emph{Pot ser que no hi hagi elements.} En un preordre o en un monoide vist com a categoria, una fletxa no envia res enlloc, i preguntar si és \enquote{injectiva} no té sentit. Mono i epi, en canvi, sempre tenen sentit, perquè es defineixen per composició (exercici~\ref{exr:pr-prima-mono-epi} de la Pràctica).
\item \emph{L'aplicació subjacent pot enganyar.} A $\cat{Top}$, una bijecció contínua és mono i epi sense ser un isomorfisme (observació anterior). A $\cat{Mon}$, la inclusió $\iota\colon(\mathbb N,+,0)\hookrightarrow(\mathbb Z,+,0)$ és un epimorfisme, tot i que no és exhaustiva. Si dos homomorfismes $g,h\colon\mathbb Z\to M$ coincideixen sobre $\mathbb N$, aleshores $g(-1)$ i $h(-1)$ són tots dos inversos de $g(1)=h(1)$, i en un monoide l'invers d'un element, si existeix, és únic. Per tant, $g(-1)=h(-1)$ i, com que tot enter és suma d'un natural i de còpies de $-1$, $g=h$.
\item \emph{Els \enquote{punts} poden no separar fletxes.} A $\cat{Grp}$, el grup trivial $1$ fa el paper de $\{\ast\}$, però de $1$ a qualsevol grup només hi ha un homomorfisme. Les fletxes $1\to G$ no distingeixen, doncs, cap parell d'homomorfismes diferents.
\end{itemize}
Per això, en una categoria general els arguments es fan amb fletxes i composicions, i reservem els adjectius \emph{injectiu}, \emph{exhaustiu} i \emph{bijectiu} per a les aplicacions, incloses les aplicacions subjacents dels morfismes d'una categoria concreta. La definició de monomorfisme és, de fet, una \enquote{injectivitat} respecte de \emph{totes} les fletxes de prova $C\to A$, i no només respecte dels elements.
\end{observacio}

\subsection{Seccions i retraccions}\label{subsec:seccions-retraccions}

\begin{definicio}\index{secció}\index{retracció}\index{invers!per un costat}
Considerem un morfisme $f\colon A\to B$. Diem que un morfisme $s\colon B\to A$ és una \textbf{secció} de $f$ quan es compleix $f\comp s=\id_B$. Això equival a dir que el diagrama següent és commutatiu:
\[
\begin{tikzpicture}[baseline=(m.center)]
  \matrix (m) [matrix of math nodes, row sep=2.6em, column sep=2.6em] {
    B & A \\
    B & \\
  };
  \draw[-{Stealth[scale=1.1]}] (m-1-1) -- node[above] {$s$} (m-1-2);
  \draw[-{Stealth[scale=1.1]}] (m-1-2) -- node[right] {$f$} (m-2-1);
  \draw[-{Stealth[scale=1.1]}] (m-1-1) -- node[left] {$\id_B$} (m-2-1);
\end{tikzpicture}
\]
També diem en aquest cas que $s$ és un \emph{invers per la dreta} de $f$.

Simètricament, diem que un morfisme $r\colon B\to A$ és una \textbf{retracció} de $f$ quan es compleix $r\comp f=\id_A$. Això equival a dir que el diagrama següent és commutatiu:
\[
\begin{tikzpicture}[baseline=(m.center)]
  \matrix (m) [matrix of math nodes, row sep=2.6em, column sep=2.6em] {
    A & B \\
    A & \\
  };
  \draw[-{Stealth[scale=1.1]}] (m-1-1) -- node[above] {$f$} (m-1-2);
  \draw[-{Stealth[scale=1.1]}] (m-1-2) -- node[right] {$r$} (m-2-1);
  \draw[-{Stealth[scale=1.1]}] (m-1-1) -- node[left] {$\id_A$} (m-2-1);
\end{tikzpicture}
\]
També diem en aquest cas que $r$ és un \emph{invers per l'esquerra} de $f$.
\end{definicio}

Una conseqüència òbvia d'aquestes dues definicions és que quan es compleix, per exemple, $f\comp s=\id_B$, se'n segueix alhora que $s$ és una secció de $f$ i que $f$ és una retracció de $s$; el mateix succeeix amb $r\comp f=\id_A$, on $r$ és una retracció de $f$ i $f$ és una secció de $r$.

Un morfisme pot tenir cap, una o diverses seccions, i igualment cap, una o diverses retraccions. Vegem-ne exemples a $\cat{Set}$:
\begin{itemize}
\item L'única aplicació $f\colon\emptyset\to\{*\}$ no té cap secció: no existeix cap aplicació $\{*\}\to\emptyset$.
\item La injecció $f\colon\{0\}\hookrightarrow\{0,1\}$ associada a la inclusió $\{0\}\subseteq\{0,1\}$ té una única retracció, l'única aplicació $\{0,1\}\to\{0\}$; en canvi no té cap secció.
\item L'aplicació $f\colon\{0,1\}\to\{*\}$ té dues seccions distintes, $*\mapsto 0$ i $*\mapsto 1$: cada tria d'una antiimatge en dona una.
\end{itemize}
Tenir secció o retracció és, doncs, una propietat especial de $f$, i quan en té, en general no és única.

\begin{proposicio}\label{prop:escindit-mono-epi}
Sigui $f\colon A\to B$. Si $f$ té alguna secció, aleshores $f$ és un epimorfisme. Si $f$ té alguna retracció, aleshores $f$ és un monomorfisme.
\end{proposicio}

\begin{prova}
Suposem que $f$ té una retracció $r$, amb $r\comp f=\id_A$, i siguin $g,h\colon C\to A$ amb $f\comp g=f\comp h$. Component per l'esquerra amb $r$,
\[
g=\id_A\comp g=(r\comp f)\comp g=r\comp(f\comp g)=r\comp(f\comp h)=(r\comp f)\comp h=h,
\]
de manera que $f$ és mono. El cas de la secció és anàleg: si $f$ té secció $s$, amb $f\comp s=\id_B$, i $g,h\colon B\to C$ amb $g\comp f=h\comp f$, aleshores component per la dreta amb $s$,
\[
g=g\comp\id_B=g\comp(f\comp s)=(g\comp f)\comp s=(h\comp f)\comp s=h\comp(f\comp s)=h\comp\id_B=h,
\]
de manera que $f$ és epi.
\end{prova}

\begin{definicio}\index{morfisme!escindit}\index{epimorfisme!escindit}\index{monomorfisme!escindit}
Diem que un morfisme $f$ queda \textbf{escindit} quan té una secció o una retracció. En concret, si $f$ té una secció, aleshores $f$ és un epimorfisme i se'n diu \textbf{epimorfisme escindit}; si $f$ té una retracció, aleshores $f$ és un monomorfisme i se'n diu \textbf{monomorfisme escindit}.
\end{definicio}

\begin{observacio}
El recíproc no val en general: hi ha monomorfismes que no tenen cap retracció (i per tant no són escindits) i epimorfismes que no tenen cap secció. En efecte, a $\cat{Set}$ tota aplicació injectiva amb domini no buit té alguna retracció i tota exhaustiva té alguna secció ---aquesta darrera afirmació és, de fet, equivalent a l'\emph{axioma de l'elecció}---, però en altres categories la situació és diferent.

Finalment, un morfisme $f$ és un \textbf{isomorfisme} si i només si té alhora alguna secció i alguna retracció. En aquest cas la secció i la retracció són forçosament \emph{úniques} i coincideixen: si $f\comp s=\id_B$ i $r\comp f=\id_A$, aleshores
\[
r=r\comp\id_B=r\comp(f\comp s)=(r\comp f)\comp s=\id_A\comp s=s.
\]
En canvi, la unicitat d’una secció o d’una retracció, considerada
separadament, no caracteritza els isomorfismes en una categoria general; per exemple, la injecció
\[
 i\colon\{0\}\hookrightarrow\{0,1\}
\]
té una única retracció $r\colon\{0,1\}\to\{0\}$, però no és un isomorfisme. No és una subtilesa avançada, sinó una contradicció amb un exemple de la mateixa subsecció.
\end{observacio}

\section{Construccions sobre categories}

A partir de categories donades se'n poden construir de noves. Aquestes construccions seran útils constantment, i algunes ---com la categoria oposada--- són fonamentals.

\subsection{La categoria oposada}\index{categoria!oposada}

\begin{definicio}
Donada una categoria $\cat{C}$, la seva \textbf{categoria oposada} $\cat{C}\op$ té:
\begin{itemize}
\item objectes: els mateixos que $\cat{C}$, $\Ob(\cat{C}\op)=\Ob(\cat{C})$;
\item morfismes: per a cada morfisme $f\colon A\to B$ de $\cat{C}$, un morfisme $f\op\colon B\to A$ de $\cat{C}\op$, de manera que
\[
\Hom_{\cat{C}\op}(B,A)=\{f\op\mid f\in\Hom_{\cat{C}}(A,B)\};
\]
\item composició: s'obté invertint l'ordre; per a $f\colon A\to B$ i $g\colon B\to C$ de $\cat{C}$,
\[
f\op\comp g\op=(g\comp f)\op\colon C\to A;
\]
\item identitats: la identitat de $A$ a $\cat{C}\op$ és $(\id_A)\op$.
\end{itemize}
És a dir, $\cat{C}\op$ s'obté invertint el sentit de totes les fletxes.
\end{definicio}

Visualment, la composició a $\cat{C}$ i la corresponent a $\cat{C}\op$ inverteixen l'ordre:
\[
\begin{array}{c@{\qquad\qquad}c}
\cat{C}: & \cat{C}\op: \\[0.6em]
\begin{tikzpicture}[baseline=(m.center)]
  \matrix (m) [matrix of math nodes, row sep=2.6em, column sep=2.6em] {
    A & B \\
      & C \\
  };
  \draw[-{Stealth[scale=1.1]}] (m-1-1) -- node[above] {$f$} (m-1-2);
  \draw[-{Stealth[scale=1.1]}] (m-1-2) -- node[right] {$g$} (m-2-2);
  \draw[-{Stealth[scale=1.1]}] (m-1-1) -- node[below left] {$g\comp f$} (m-2-2);
\end{tikzpicture}
&
\begin{tikzpicture}[baseline=(m.center)]
  \matrix (m) [matrix of math nodes, row sep=2.6em, column sep=2.6em] {
    A & B \\
      & C \\
  };
  \draw[-{Stealth[scale=1.1]}] (m-1-2) -- node[above] {$f\op$} (m-1-1);
  \draw[-{Stealth[scale=1.1]}] (m-2-2) -- node[right] {$g\op$} (m-1-2);
  \draw[-{Stealth[scale=1.1]}] (m-2-2) -- node[below left] {$f\op\comp g\op$} (m-1-1);
\end{tikzpicture}
\end{array}
\]
L'associativitat a $\cat{C}\op$ es dedueix de la de $\cat{C}$ invertint l'ordre, i les lleis d'identitat, de les de $\cat{C}$: per exemple, $f\op\comp(\id_B)\op=(\id_B\comp f)\op=f\op$. La categoria oposada és la base de la \textbf{dualitat}\index{dualitat}, que desenvolupem a la secció~\ref{sec:dualitat}: tota afirmació formulada purament en el llenguatge categòric dona, invertint les fletxes, una afirmació dual, i aquesta correspondència es pot fer del tot precisa.

\begin{observacio}
Aplicar dues vegades l'oposada retorna la categoria de partida:
\[
(\cat{C}\op)\op=\cat{C}.
\]
Invertir les fletxes i tornar-les a invertir les deixa com estaven.
\end{observacio}

Vegem com queda l'oposada en dos casos. Quan un preordre $(P,\le)$ es considera com una categoria ---amb un morfisme $x\to y$ quan $x\le y$---, la seva oposada té un morfisme $x\to y$ exactament quan $y\le x$. I quan un monoide $(M,\ast,e)$ es considera com una categoria, la seva oposada té els mateixos morfismes ---els elements de $M$---, però la composició de $a$ i $b$ és $b\ast a$ en lloc de $a\ast b$.

\subsection{La categoria producte}\label{subsec:categoria-producte}\index{categoria!producte}

\begin{definicio}
Donades dues categories $\cat{C}$ i $\cat{D}$, la \textbf{categoria producte} $\cat{C}\times\cat{D}$ té:
\begin{itemize}
\item objectes: $\Ob(\cat{C}\times\cat{D})$ és la col·lecció de les parelles $(A,X)$ amb $A\in\Ob(\cat{C})$ i $X\in\Ob(\cat{D})$;
\item morfismes: $\Hom_{\cat{C}\times\cat{D}}\bigl((A,X),(B,Y)\bigr)=\Hom_{\cat{C}}(A,B)\times\Hom_{\cat{D}}(X,Y)$, és a dir, les parelles $(f,g)$ amb $f\colon A\to B$ a $\cat{C}$ i $g\colon X\to Y$ a $\cat{D}$;
\item composició: component a component,
\[
(f',g')\comp(f,g)=(f'\comp f,\ g'\comp g);
\]
\item identitats: $\id_{(A,X)}=(\id_A,\id_X)$.
\end{itemize} Els axiomes de categoria per a $\cat{C}\times\cat{D}$ se segueixen dels de $\cat{C}$ i $\cat{D}$ aplicats a cada component.
\end{definicio}

\begin{exemple}
Prenem la categoria $\cat{2}$ de la subsecció~\ref{subsec:finites} ---dos objectes i una única fletxa no trivial---, i n'escrivim els objectes com $0$ i $1$, amb la fletxa no trivial $0\to 1$. La categoria producte $\cat{2}\times\cat{2}$ té quatre objectes, que es disposen de manera natural en un quadrat:
\[
\begin{tikzpicture}[baseline=(m.center)]
  \matrix (m) [matrix of math nodes, row sep=2.6em, column sep=3.4em] {
    (0,0) & (0,1) \\
    (1,0) & (1,1) \\
  };
  \draw[-{Stealth[scale=1.1]}] (m-1-1) -- (m-1-2);
  \draw[-{Stealth[scale=1.1]}] (m-1-1) -- (m-2-1);
  \draw[-{Stealth[scale=1.1]}] (m-1-2) -- (m-2-2);
  \draw[-{Stealth[scale=1.1]}] (m-2-1) -- (m-2-2);
  \draw[-{Stealth[scale=1.1]}] (m-1-1) -- (m-2-2);
\end{tikzpicture}
\]
Cada fletxa és una parella de fletxes de $\cat{2}$. Com que $\cat{2}$ és prima, també ho és $\cat{2}\times\cat{2}$: entre dos objectes hi ha com a màxim un morfisme. Per això els dos camins del quadrat de $(0,0)$ a $(1,1)$ tenen necessàriament la mateixa composició, que és la diagonal; el quadrat commuta.
\end{exemple}

\begin{observacio}[Dos usos de \enquote{producte}]\label{obs:dos-productes}\index{producte!d'objectes i de categories}
No s'han de confondre dues construccions que porten el mateix nom. La categoria producte $\cat{C}\times\cat{D}$ és una \emph{categoria} nova, construïda a partir de dues categories. En canvi, el producte $A\times B$ de dos objectes ---que la introducció fa servir com a primer exemple de propietat universal i que avançarem a la subsecció~\ref{subsec:producte-coproducte}--- és un \emph{objecte} d'una mateixa categoria, caracteritzat per una propietat universal. Les dues nocions estan relacionades: quan disposem de la categoria $\cat{Cat}$ de les categories petites (capítol~\ref{cap:functors}), si $\cat{C}$ i $\cat{D}$ són petites, la categoria $\cat{C}\times\cat{D}$ es podrà veure com el producte dels objectes $\cat{C}$ i $\cat{D}$ dins de $\cat{Cat}$. Però no són la mateixa construcció: l'una fabrica una categoria a partir de dues, i l'altra viu dins d'una sola categoria.
\end{observacio}

\subsection{Subcategories}\label{subsec:subcategories}\index{subcategoria}

Una altra manera d'obtenir noves categories a partir de categories antigues és restringint-ne els objectes o les fletxes.

\begin{definicio}
Una \textbf{subcategoria} $\cat{S}$ d'una categoria $\cat{C}$ ve donada per una subcol·lecció d'objectes $\Ob(\cat{S})\subseteq\Ob(\cat{C})$ i, per a cada parella $A,B\in\Ob(\cat{S})$, una subcol·lecció de morfismes $\Hom_{\cat{S}}(A,B)\subseteq\Hom_{\cat{C}}(A,B)$, de manera que:
\begin{itemize}
\item la identitat $\id_A$ pertany a $\Hom_{\cat{S}}(A,A)$ per a cada $A\in\Ob(\cat{S})$;
\item la composició de dos morfismes de $\cat{S}$ componibles pertany a $\cat{S}$.
\end{itemize}
Amb aquestes condicions, $\cat{S}$ és una categoria per dret propi, amb la composició i les identitats heretades de $\cat{C}$.
\end{definicio}

Dos casos particulars, en cert sentit oposats, són especialment útils. Es diu que $\cat{S}$ és una \textbf{subcategoria plena}\index{subcategoria!plena} quan $\Hom_{\cat{S}}(A,B)=\Hom_{\cat{C}}(A,B)$ per a tot $A,B\in\Ob(\cat{S})$; una subcategoria plena queda determinada, doncs, només per la seva col·lecció d'objectes. Simètricament, es diu que $\cat{S}$ és una \textbf{subcategoria ampla}\index{subcategoria!ampla} quan $\Ob(\cat{S})=\Ob(\cat{C})$, encara que només contingui alguns dels morfismes. En resum: una subcategoria plena restringeix els objectes i conserva tots els morfismes entre ells; una subcategoria ampla conserva tots els objectes i en restringeix els morfismes.

\begin{exemple}
$\cat{FinSet}$ és una subcategoria plena de $\cat{Set}$: pren com a objectes els conjunts finits i, entre dos d'ells, totes les aplicacions. En canvi, la subcategoria de $\cat{Set}$ que té tots els conjunts per objectes però només les aplicacions \emph{injectives} per morfismes és una subcategoria ampla que no és plena: conserva tots els objectes, però deixa fora les aplicacions que no són injectives.
\end{exemple}

\subsection{Categories sobre i sota un objecte}\index{categoria!sobre un objecte}\index{categoria!sota un objecte}

Fixat un objecte $A$ d'una categoria $\cat{C}$, se'n pot construir una categoria nova els objectes de la qual són els morfismes que apunten cap a $A$.

\begin{definicio}\label{def:slice}
Sigui $\cat{C}$ una categoria i $A$ un objecte. La \textbf{categoria de $\cat{C}$ sobre $A$} (en anglès, \emph{slice category}), que escrivim $\cat{C}/A$, té:
\begin{itemize}
\item objectes: $\Ob(\cat{C}/A)$ és la col·lecció dels morfismes $x\colon X\to A$ de $\cat{C}$ amb codomini $A$;
\item morfismes: $\Hom_{\cat{C}/A}(x,y)$, per a $x\colon X\to A$ i $y\colon Y\to A$, és el conjunt dels morfismes $h\colon X\to Y$ de $\cat{C}$ que fan commutar el triangle
\[
\begin{tikzpicture}[baseline=(m.center)]
  \matrix (m) [matrix of math nodes, row sep=2.4em, column sep=2.4em] {
    X & & Y \\
      & A & \\
  };
  \draw[-{Stealth[scale=1.0]}] (m-1-1) -- node[above] {$h$} (m-1-3);
  \draw[-{Stealth[scale=1.0]}] (m-1-1) -- node[below left] {$x$} (m-2-2);
  \draw[-{Stealth[scale=1.0]}] (m-1-3) -- node[below right] {$y$} (m-2-2);
\end{tikzpicture}
\qquad\text{és a dir, } y\comp h=x;
\]
\item composició: la de $\cat{C}$; si $h\in\Hom_{\cat{C}/A}(x,y)$ i $k\in\Hom_{\cat{C}/A}(y,z)$, aleshores $k\comp h\in\Hom_{\cat{C}/A}(x,z)$, perquè $z\comp(k\comp h)=(z\comp k)\comp h=y\comp h=x$;
\item identitats: la identitat de $x\colon X\to A$ és $\id_X$, que fa commutar el triangle perquè $x\comp\id_X=x$.
\end{itemize}
Els axiomes s'hereten de $\cat{C}$.
\end{definicio}

Anàlogament es defineix la \textbf{categoria de $\cat{C}$ sota $A$} (en anglès, \emph{coslice category}), que escrivim $A/\cat{C}$:
\begin{itemize}
\item objectes: $\Ob(A/\cat{C})$ és la col·lecció dels morfismes $x\colon A\to X$ de $\cat{C}$ amb domini $A$;
\item morfismes: $\Hom_{A/\cat{C}}(x,y)$, per a $x\colon A\to X$ i $y\colon A\to Y$, és el conjunt dels morfismes $h\colon X\to Y$ de $\cat{C}$ amb $h\comp x=y$;
\item composició i identitats: les de $\cat{C}$, com a $\cat{C}/A$.
\end{itemize}
Un morfisme de $A/\cat{C}$ és, doncs, un $h$ que fa commutar el triangle sota $A$:
\[
\begin{tikzpicture}[baseline=(m.center)]
  \matrix (m) [matrix of math nodes, row sep=2.4em, column sep=2.4em] {
      & A & \\
    X & & Y \\
  };
  \draw[-{Stealth[scale=1.0]}] (m-1-2) -- node[above left] {$x$} (m-2-1);
  \draw[-{Stealth[scale=1.0]}] (m-1-2) -- node[above right] {$y$} (m-2-3);
  \draw[-{Stealth[scale=1.0]}] (m-2-1) -- node[below] {$h$} (m-2-3);
\end{tikzpicture}
\]
Aquí també convé fer explícits els extrems, com a l'observació~\ref{obs:extrems-implicits}. Un morfisme de $\cat{C}/A$ de $x$ a $y$ no és el morfisme $h\colon X\to Y$ de $\cat{C}$ tot sol, sinó $h$ \emph{junt amb} els seus extrems $x$ i $y$ a $\cat{C}/A$: la terna $(x,y,h)$. La raó és la mateixa: un mateix $h\colon X\to Y$ de $\cat{C}$ pot fer commutar el triangle per a objectes distints de $\cat{C}/A$ ---per exemple, si $X$ és el domini de dos objectes diferents $x,x'\colon X\to A$---, i sense registrar els extrems tindria dominis i codominis indeterminats a $\cat{C}/A$. Escriurem $h$ quan els extrems se sobreentenguin. Anàlogament per a $A/\cat{C}$.

\begin{exemple}
A $\cat{Set}/A$, un objecte és una aplicació $x\colon X\to A$. Les fibres $x^{-1}(a)$, per a $a\in A$, permeten interpretar aquest objecte com una família de conjunts $\bigl(x^{-1}(a)\bigr)_{a\in A}$ indexada per $A$. Recíprocament, una família de conjunts indexada per $A$ dona, mitjançant la unió disjunta, una aplicació cap a $A$. Aquestes dues construccions no identifiquen literalment les dades ---la unió disjunta demana triar-ne una representació---, però sí que estableixen la correspondència adequada entre totes dues descripcions. La farem precisa als capítols~\ref{cap:functors} i~\ref{cap:transformacions}, quan disposem del llenguatge dels functors i de les equivalències de categories.
\end{exemple}

\begin{exemple}\index{conjunt apuntat}
Prenem $A=\{\ast\}$, un conjunt d'un sol element. Un objecte de $\{\ast\}/\cat{Set}$ és una aplicació $\{\ast\}\to X$, és a dir, l'elecció d'un element distingit de $X$, el \textbf{punt base}. Un morfisme és una aplicació que respecta el punt base. Així, $\{\ast\}/\cat{Set}$ és la categoria dels \textbf{conjunts apuntats} (en anglès, \emph{pointed sets}).
\end{exemple}

\section{Principi de dualitat}\label{sec:dualitat}\index{dualitat!principi de}

La dualitat no és només una llista de parelles de conceptes. És un principi que transforma definicions, enunciats i demostracions, i la categoria oposada permet formular-lo amb precisió. Aquesta secció en fa explícit el mecanisme; el retrobarem en tractar objectes inicials i terminals, productes i coproductes, productes fibrats i coproductes fibrats (en anglès, \emph{pullbacks} i \emph{pushouts}), i límits i colímits.

\subsection{Dualitzar un enunciat}\index{dualitat!d'un enunciat}

El principi s'aplica a afirmacions escrites \emph{exclusivament} en el llenguatge de les categories: hi poden intervenir objectes i morfismes, el domini i el codomini de cada morfisme, identitats, composicions i igualtats entre composicions, combinats amb connectors lògics i amb quantificadors sobre objectes i morfismes. Les definicions de monomorfisme, d'isomorfisme o de secció són d'aquesta mena. En canvi, \enquote{$f$ és injectiva} no ho és, perquè parla dels elements d'un conjunt.

Donada una afirmació $\varphi$ d'aquest tipus, n'obtenim la \textbf{dual} $\varphi\op$ \emph{invertint totes les fletxes}:
\begin{itemize}
  \item cada morfisme $f\colon A\to B$ s'inverteix: a la dual es llegeix com una fletxa $B\to A$; per tant, s'intercanvien domini i codomini;
  \item cada composició $g\comp f$ passa a ser $f\comp g$, és a dir, s'inverteix l'ordre de composició;
  \item les identitats, les igualtats, els connectors i els quantificadors es mantenen.
\end{itemize}
Dualitzar dues vegades retorna l'afirmació de partida: $(\varphi\op)\op=\varphi$.

La dual té una lectura immediata: $\varphi\op$ diu de $\cat{C}$ exactament el que $\varphi$ diu de $\cat{C}\op$. En efecte, un morfisme $A\to B$ de $\cat{C}\op$ és un morfisme $B\to A$ de $\cat{C}$, les identitats són les mateixes i la composició de $\cat{C}\op$ és la de $\cat{C}$ en l'ordre invers, $g\op\comp f\op=(f\comp g)\op$. Per tant, per a tota categoria $\cat{C}$,
\[
\varphi \text{ es compleix a } \cat{C}\op
\quad\Longleftrightarrow\quad
\varphi\op \text{ es compleix a } \cat{C}.
\]
D'aquí surt el principi.

\begin{teorema}[Principi de dualitat]\label{teo:dualitat}\index{dualitat!principi de}
Sigui $\varphi$ una afirmació expressable en el llenguatge de les categories. Si $\varphi$ és vàlida en totes les categories, la seva dual $\varphi\op$ també és vàlida en totes les categories.
\end{teorema}

\begin{prova}
Sigui $\cat{C}$ una categoria qualsevol. Com que $\cat{C}\op$ també és una categoria, per hipòtesi $\varphi$ es compleix a $\cat{C}\op$; per l'equivalència anterior, $\varphi\op$ es compleix a $\cat{C}$.
\end{prova}

La manera habitual d'usar el principi és aquesta: s'aplica un teorema conegut a $\cat{C}\op$ i se'n tradueix la conclusió al llenguatge de $\cat{C}$. No cal repetir la demostració; \textbf{la demostració dual és la demostració original amb totes les fletxes invertides}.

Aquesta formulació operativa és la que farem servir al llarg del llibre. Se'n pot donar també una versió completament formal: fixar un llenguatge de primer ordre per a la teoria de categories, definir la dualització com una transformació de fórmules i comprovar que els axiomes en queden invariants. Aquesta formalització, amb una versió sintàctica (sobre deduccions) i una de semàntica (sobre models), és a l'apèndix~\ref{cap:axiomatica-relacional} (secció~\ref{sec:ar-dualitat}); no és necessària per seguir la resta del llibre.

\subsection{Primer exemple: de monomorfisme a epimorfisme}

Recordem (subsecció~\ref{subsec:mono-epi}) que $m\colon A\to B$ és un \textbf{monomorfisme} si, per a tot parell $u,v\colon X\to A$,
\[
\begin{tikzpicture}[baseline=(m.center)]
  \matrix (m) [matrix of math nodes, column sep=3.2em] { X & A & B \\ };
  \draw[-{Stealth[scale=1.0]}] ([yshift=2.6pt]m-1-1.east) -- node[above] {$u$} ([yshift=2.6pt]m-1-2.west);
  \draw[-{Stealth[scale=1.0]}] ([yshift=-2.6pt]m-1-1.east) -- node[below] {$v$} ([yshift=-2.6pt]m-1-2.west);
  \draw[-{Stealth[scale=1.1]}] (m-1-2) -- node[above] {$m$} (m-1-3);
\end{tikzpicture}
\qquad
m\comp u = m\comp v \implies u=v.
\]
Formem-ne el dual: intercanviem domini i codomini i invertim l'ordre de cada composició. L'enunciat dual diu que, per a tot parell $u,v\colon A\to X$,
\[
\begin{tikzpicture}[baseline=(m.center)]
  \matrix (m) [matrix of math nodes, column sep=3.2em] { B & A & X \\ };
  \draw[-{Stealth[scale=1.1]}] (m-1-1) -- node[above] {$m$} (m-1-2);
  \draw[-{Stealth[scale=1.0]}] ([yshift=2.6pt]m-1-2.east) -- node[above] {$u$} ([yshift=2.6pt]m-1-3.west);
  \draw[-{Stealth[scale=1.0]}] ([yshift=-2.6pt]m-1-2.east) -- node[below] {$v$} ([yshift=-2.6pt]m-1-3.west);
\end{tikzpicture}
\qquad
u\comp m = v\comp m \implies u=v,
\]
que és exactament la definició d'\textbf{epimorfisme}. El dual de \enquote{ser un monomorfisme} és, doncs, \enquote{ser un epimorfisme}.

Ara podem dualitzar teoremes sense repetir-ne la demostració. Per exemple, en qualsevol categoria, si $A\xrightarrow{f}B\xrightarrow{g}C$ són monomorfismes, la composició $g\comp f$ també ho és. Pel principi de dualitat (teorema~\ref{teo:dualitat}), l'enunciat dual val igualment:
\begin{quote}
en qualsevol categoria, la composició de dos epimorfismes és un epimorfisme.
\end{quote}
No n'hem hagut de fer cap demostració nova. L'exercici~\ref{exr:pr-dual-mono-epi} de la Pràctica fa aquest pas en detall i assenyala exactament què s'inverteix en passar a $\cat{C}\op$.

\subsection{Segon exemple: de secció a retracció}

Recordem (subsecció~\ref{subsec:seccions-retraccions}) que $s\colon B\to A$ és una \textbf{secció} de $f\colon A\to B$ quan $f\comp s=\id_B$. Dualitzem aquest enunciat: $f$ passa a ser una fletxa $B\to A$, $s$ una fletxa $A\to B$, i la igualtat $f\comp s=\id_B$ es converteix en $s\comp f=\id_B$. En diagrames, invertim les tres fletxes del triangle:
\[
\begin{tikzpicture}[baseline=(m.center)]
  \matrix (m) [matrix of math nodes, row sep=2.6em, column sep=2.6em] {
    B & A \\
    B & \\
  };
  \draw[-{Stealth[scale=1.1]}] (m-1-1) -- node[above] {$s$} (m-1-2);
  \draw[-{Stealth[scale=1.1]}] (m-1-2) -- node[right] {$f$} (m-2-1);
  \draw[-{Stealth[scale=1.1]}] (m-1-1) -- node[left] {$\id_B$} (m-2-1);
\end{tikzpicture}
\quad f\comp s=\id_B
\qquad\qquad
\begin{tikzpicture}[baseline=(m.center)]
  \matrix (m) [matrix of math nodes, row sep=2.6em, column sep=2.6em] {
    B & A \\
    B & \\
  };
  \draw[-{Stealth[scale=1.1]}] (m-1-2) -- node[above] {$s$} (m-1-1);
  \draw[-{Stealth[scale=1.1]}] (m-2-1) -- node[right] {$f$} (m-1-2);
  \draw[-{Stealth[scale=1.1]}] (m-2-1) -- node[left] {$\id_B$} (m-1-1);
\end{tikzpicture}
\quad s\comp f=\id_B
\]
La igualtat de la dreta diu exactament que $s$ és una \textbf{retracció} de $f$. El dual de \enquote{tenir una secció} és, doncs, \enquote{tenir una retracció}, i el d'epimorfisme escindit és monomorfisme escindit.

Amb aquest diccionari es poden dualitzar també els resultats sobre seccions i retraccions; els exercicis proposats del capítol ho apliquen a la proposició~\ref{prop:escindit-mono-epi}.

\subsection{Tercer exemple: del producte al coproducte}\label{subsec:producte-coproducte}

A tall d'avançament ---els definirem en rigor al capítol de límits---, considerem el producte i el coproducte de dos objectes, que il·lustren la dualització d'un enunciat amb un quantificador d'existència única.

Un \emph{producte} de $A$ i $B$ és un objecte $A\times B$ amb dos morfismes $\pi_1\colon A\times B\to A$ i $\pi_2\colon A\times B\to B$ tals que, per a tot $X$ i tot parell $f\colon X\to A$, $g\colon X\to B$, existeix un únic $u\colon X\to A\times B$ amb $\pi_1\comp u=f$ i $\pi_2\comp u=g$:
\[
\begin{tikzpicture}[baseline=(m.center)]
  \matrix (m) [matrix of math nodes, row sep=2.6em, column sep=2.8em] {
     & X & \\
    A & A\times B & B \\
  };
  \draw[-{Stealth[scale=1.0]}] (m-1-2) -- node[above left] {$f$} (m-2-1);
  \draw[-{Stealth[scale=1.0]}] (m-1-2) -- node[above right] {$g$} (m-2-3);
  \draw[-{Stealth[scale=1.0]},dashed] (m-1-2) -- node[right] {$u$} (m-2-2);
  \draw[-{Stealth[scale=1.0]}] (m-2-2) -- node[below] {$\pi_1$} (m-2-1);
  \draw[-{Stealth[scale=1.0]}] (m-2-2) -- node[below] {$\pi_2$} (m-2-3);
\end{tikzpicture}
\]
En el diagrama seguim una convenció que farem servir cada cop que dibuixem una propietat universal: la \textbf{fletxa discontínua} marca el morfisme l'existència del qual afirma la propietat ---aquí el mediador $u$, un cop donades la resta de dades: l'objecte $X$ i les fletxes $f$ i $g$---, mentre que les fletxes contínues són les dades de partida. Igualment, escriurem $\exists!$ per abreujar \enquote{existeix un únic}, com acabem de fer en enunciar la propietat.

Invertint totes les fletxes obtenim un objecte $A+B$ amb dos morfismes $\iota_1\colon A\to A+B$ i $\iota_2\colon B\to A+B$ tals que, per a tot $X$ i tot parell $f\colon A\to X$, $g\colon B\to X$, existeix un únic $u\colon A+B\to X$ amb $u\comp\iota_1=f$ i $u\comp\iota_2=g$:
\[
\begin{tikzpicture}[baseline=(m.center)]
  \matrix (m) [matrix of math nodes, row sep=2.6em, column sep=2.8em] {
    A & A+B & B \\
     & X & \\
  };
  \draw[-{Stealth[scale=1.0]}] (m-1-1) -- node[above] {$\iota_1$} (m-1-2);
  \draw[-{Stealth[scale=1.0]}] (m-1-3) -- node[above] {$\iota_2$} (m-1-2);
  \draw[-{Stealth[scale=1.0]}] (m-1-1) -- node[below left] {$f$} (m-2-2);
  \draw[-{Stealth[scale=1.0]}] (m-1-3) -- node[below right] {$g$} (m-2-2);
  \draw[-{Stealth[scale=1.0]},dashed] (m-1-2) -- node[right] {$u$} (m-2-2);
\end{tikzpicture}
\]
Aquesta és la definició del \emph{coproducte} $A+B$. Els quantificadors $\forall$ i $\exists!$ i la conjunció de les dues equacions es mantenen; el que s'inverteix és el domini i el codomini de cada morfisme i l'ordre de cada composició.

\paragraph{Un cas que cal vigilar.} La construcció \enquote{categoria sobre $A$} es dualitza en la construcció \enquote{categoria sota $A$} (definició~\ref{def:slice}). Ara bé, en dualitzar aquesta construcció cal dualitzar també la categoria ambient. La fórmula precisa és
\[
(\cat{C}/A)\op\;\cong\;A/(\cat{C}\op),
\qquad
(A/\cat{C})\op\;\cong\;(\cat{C}\op)/A.
\]
No s'ha de llegir, doncs, que $(\cat{C}/A)\op\cong A/\cat{C}$ mantenint $\cat{C}$ intacta: això és fals en general. Per exemple, sigui $\cat{C}=\cat{2}=(0\to 1)$ i prenguem com a $A$ l'\emph{objecte} $1$ de $\cat{2}$ (no la categoria $\cat{1}$). La categoria $\cat{2}/1$ té dos objectes (les dues fletxes amb codomini $1$: la fletxa $0\to 1$ i $\id_1$), mentre que $1/\cat{2}$ en té un de sol ($\id_1$, l'única fletxa amb domini $1$); per tant $(\cat{2}/1)\op$ \emph{no} és isomorfa a $1/\cat{2}$. El que sí que val és $(\cat{2}/1)\op\cong 1/(\cat{2}\op)$: a $\cat{2}\op$ hi ha dues fletxes amb domini $1$, la identitat i la fletxa $1\to 0$. En tots els casos, la mateixa regla ---invertir les fletxes--- transforma cada noció en la seva dual, però la regla actua sobre \emph{tot} l'enunciat, ambient inclòs.

\subsection{Diccionari bàsic de dualitats}

Recollim les parelles duals. Les primeres ja han aparegut en aquest capítol; la resta les trobarem més endavant, i la taula creixerà a mesura que apareguin nous conceptes.
\begin{center}
\begin{tabular}{ll}
\hline
\textbf{Concepte} & \textbf{Concepte dual}\\
\hline
\multicolumn{2}{l}{\emph{Ja vistos en aquest capítol}}\\
monomorfisme & epimorfisme\\
secció & retracció\\
producte & coproducte\\
categoria sobre $A$ \ ($\cat{C}/A$) & categoria sota $A$ \ ($A/\cat{C}$)\\
\hline
\multicolumn{2}{l}{\emph{Més endavant}}\\
objecte inicial & objecte terminal\\
equalitzador & coequalitzador\\
producte fibrat & coproducte fibrat\\
con & cocon\\
límit & colímit\\
\hline
\end{tabular}
\end{center}
Algunes nocions són \textbf{autoduals}\index{dualitat!noció autodual}: la identitat, l'isomorfisme i la commutativitat d'un diagrama continuen sent nocions del mateix tipus després de dualitzar-les.

La taula recull les parelles duals com a \emph{nocions}: cada concepte i el seu dual. Quan aquestes nocions es realitzen en categories concretes, però, la dualitat s'ha d'aplicar a tot l'enunciat, ambient inclòs. La parella \enquote{sobre $A$ / sota $A$} n'és l'exemple a vigilar: com hem vist, l'enunciat correcte és $(\cat{C}/A)\op\cong A/(\cat{C}\op)$, no $(\cat{C}/A)\op\cong A/\cat{C}$ amb $\cat{C}$ intacta.

\subsection{Abast i precaucions}

El principi només s'aplica a afirmacions expressables en el llenguatge de les categories. Una afirmació que faci servir informació externa ---elements d'un conjunt, oberts d'un espai o una construcció concreta--- no té dual fins que no es reformula exclusivament amb objectes i morfismes. I quan l'afirmació es refereix a construccions fetes a partir d'una categoria, cal dualitzar també la categoria ambient, com acabem de veure amb les categories sobre i sota un objecte.

El principi de dualitat no permet concloure, per si sol, que una categoria concreta sigui equivalent a la seva oposada. El principi relaciona afirmacions sobre $\cat{C}$ amb les afirmacions duals sobre $\cat{C}\op$; no proporciona en general cap equivalència $\cat{C}\simeq\cat{C}\op$. Algunes categories sí que ho són ---per exemple, $\cat{2}=(0\to 1)$ és isomorfa a $\cat{2}\op$ intercanviant-ne els dos objectes---, però això s'ha de comprovar en cada cas. En canvi, $\cat{Set}$ no s'identifica amb $\cat{Set}\op$, encara que els teoremes generals sobre productes tinguin teoremes duals sobre coproductes.

Finalment, una precaució sobre la lògica. Els connectors i els quantificadors amb què escrivim una afirmació $\varphi$ pertanyen al llenguatge amb què \emph{parlem} de la categoria, i la dualitat els deixa intactes. Una altra cosa és el que passa quan els objectes de la categoria són ells mateixos fórmules, com a la categoria $\cat{Prov}_T$ de la subsecció~\ref{subsec:elementals}. Allà, \emph{quan existeixen}, la conjunció és un producte, la disjunció un coproducte, el cert un objecte terminal i el fals un objecte inicial, i la dualitat de fletxes els intercanvia. Però un preordre de deduïbilitat qualsevol no garanteix que aquests connectors existeixin: en el preordre discret de dues fórmules sense cap deducció entre elles, aquestes dues fórmules no tenen ni producte ni coproducte. I que la dualitat de fletxes coincideixi exactament amb la dualitat de De~Morgan exigeix, a més, un complement involutiu com el de les àlgebres booleanes. Tornarem sobre aquesta lògica interna més endavant, amb el llenguatge adequat.

\section{Propietats duals dels morfismes}\label{sec:propietats-duals}

El principi de dualitat estalvia feina: n'hi ha prou de demostrar les propietats dels monomorfismes; les dels epimorfismes se n'obtenen invertint les fletxes, sense cap demostració nova. Ho il·lustrem amb les propietats bàsiques.

\subsection{Propietats dels monomorfismes}

\begin{proposicio}\label{prop:mono-comp}
La composició de dos monomorfismes és un monomorfisme: si $f\colon A\to B$ i $g\colon B\to C$ són monomorfismes, aleshores $g\comp f$ ho és.
\end{proposicio}

\begin{prova}
Siguin $u,v\colon X\to A$ amb $(g\comp f)\comp u=(g\comp f)\comp v$, és a dir $g\comp(f\comp u)=g\comp(f\comp v)$. Com que $g$ és mono, $f\comp u=f\comp v$; i com que $f$ és mono, $u=v$.
\[
\begin{tikzpicture}[baseline=(m.center)]
  \matrix (m) [matrix of math nodes, column sep=3em] { X & A & B & C \\ };
  \draw[-{Stealth[scale=1.0]}] ([yshift=2.6pt]m-1-1.east) -- node[above] {$u$} ([yshift=2.6pt]m-1-2.west);
  \draw[-{Stealth[scale=1.0]}] ([yshift=-2.6pt]m-1-1.east) -- node[below] {$v$} ([yshift=-2.6pt]m-1-2.west);
  \draw[-{Stealth[scale=1.1]}] (m-1-2) -- node[above] {$f$} (m-1-3);
  \draw[-{Stealth[scale=1.1]}] (m-1-3) -- node[above] {$g$} (m-1-4);
\end{tikzpicture}
\]
\end{prova}

\begin{proposicio}\label{prop:mono-factor}
Si la composició $g\comp f$ és un monomorfisme, aleshores $f$ és un monomorfisme.
\end{proposicio}

\begin{prova}
Siguin $u,v\colon X\to A$ amb $f\comp u=f\comp v$. Component amb $g$, $(g\comp f)\comp u=(g\comp f)\comp v$; com que $g\comp f$ és mono, $u=v$. (No cal cap hipòtesi sobre $g$.)
\end{prova}

\begin{proposicio}\label{prop:iso-retr-mono}
Tot monomorfisme escindit és un monomorfisme, i tot isomorfisme és un monomorfisme escindit. Per tant,
\[
\text{isomorfisme} \;\Longrightarrow\; \text{monomorfisme escindit} \;\Longrightarrow\; \text{monomorfisme}.
\]
\end{proposicio}

\begin{prova}
La primera implicació és la proposició~\ref{prop:escindit-mono-epi}: si $f$ té una retracció, aleshores $f$ és un monomorfisme. Per a la segona, si $f$ és un isomorfisme, el seu invers $f^{-1}$ satisfà $f^{-1}\comp f=\id_A$, de manera que $f^{-1}$ és una retracció de $f$; per tant $f$ té retracció i és un monomorfisme escindit. La cadena d'implicacions se'n segueix.
\end{prova}

\subsection{Les propietats duals dels epimorfismes}

Aplicant el principi de dualitat a les proposicions anteriors ---invertint les fletxes--- obtenim, sense cap demostració nova, les propietats corresponents dels epimorfismes.

\begin{proposicio}[dual de la proposició~\ref{prop:mono-comp}]\label{prop:epi-comp}
La composició de dos epimorfismes és un epimorfisme.
\end{proposicio}

\begin{proposicio}[dual de la proposició~\ref{prop:mono-factor}]\label{prop:epi-factor}
Si la composició $g\comp f$ és un epimorfisme, aleshores $g$ és un epimorfisme.
\end{proposicio}

\begin{proposicio}[dual de la proposició~\ref{prop:iso-retr-mono}]\label{prop:iso-sec-epi}
Tot epimorfisme escindit és un epimorfisme, i tot isomorfisme és un epimorfisme escindit. Per tant,
\[
\text{isomorfisme} \;\Longrightarrow\; \text{epimorfisme escindit} \;\Longrightarrow\; \text{epimorfisme}.
\]
\end{proposicio}

Observem l'asimetria que revela la dualitat: a la proposició~\ref{prop:mono-factor}, d'un mono $g\comp f$ hereta el factor \emph{de la dreta}, $f$; a la seva dual~\ref{prop:epi-factor}, d'un epi hereta el factor \emph{de l'esquerra}, $g$.

\subsection{Condicions necessàries i suficients}

Les implicacions anteriors \emph{no} es reverteixen en general, i és aquí on convé anar amb compte. La cadena vàlida a qualsevol categoria és
\[
\text{iso} \;\Longrightarrow\; \text{mono escindit} \;\overset{(1)}{\Longrightarrow}\; \text{mono},
\]
i cap de les dues fletxes es pot invertir en general:
\begin{itemize}
\item[(1)] \textbf{mono $\not\Rightarrow$ mono escindit.} A $\cat{Grp}$, l'homomorfisme $\mathbb{Z}\xrightarrow{\,\cdot 2\,}\mathbb{Z}$ és un monomorfisme, però no té cap retracció: no és escindit.
\item[(2)] \textbf{mono escindit $\not\Rightarrow$ iso.} Un monomorfisme escindit no ha de ser exhaustiu en cap sentit: a $\cat{Set}$, la injecció $\{0\}\hookrightarrow\{0,1\}$ té retracció (l'única aplicació $\{0,1\}\to\{0\}$), és per tant un monomorfisme escindit, però no és un isomorfisme.
\end{itemize}
Dualment, $\text{iso}\Rightarrow\text{epi escindit}\Rightarrow\text{epi}$, amb les mateixes dues fletxes no reversibles. Convé no confondre \enquote{mono escindit $+$ epi} amb \enquote{mono $+$ epi}: un \textbf{bimorfisme} (alhora mono i epi) no és necessàriament un isomorfisme ---l'única fletxa no trivial de $\cat{2}$ ho és sense tenir invers, i a $\cat{Top}$ una bijecció contínua amb inversa no contínua també---, mentre que un morfisme alhora mono escindit i epi sí que ho és, com recull la caracterització següent:
\[
f \text{ és iso} \iff f \text{ és mono escindit i epi} \iff f \text{ és epi escindit i mono}.
\]
Les demostracions d'aquestes equivalències, i la de la cadena d'implicacions anterior, es proposen a les parts de pràctica i d'exercicis.

\section{Grafs i categories lliures}\label{sec:grafs-lliures}\index{categoria!lliure}

Els grafs i les categories estan estretament relacionats, però no són el mateix. Un graf dirigit només especifica vèrtexs i arestes. Una categoria, a més, disposa d'identitats, permet compondre fletxes i exigeix que aquestes composicions satisfacin els axiomes. De fet, per a una categoria petita, si oblidem la composició i quines fletxes són identitats, els objectes i morfismes formen un graf dirigit: és el seu \emph{graf subjacent}. Les fletxes identitat no desapareixen ---continuen sent arestes del graf, un bucle a cada vèrtex---, només que ja no estan marcades com a identitats. Cal que la categoria sigui petita, perquè els vèrtexs i les arestes d'un graf formen conjunts.

Hi ha una manera natural de fer el camí contrari: generar una categoria a partir d'un graf dirigit prenent com a morfismes tots els camins possibles.

\begin{definicio}\label{def:free}
Sigui $G$ un graf dirigit, amb conjunt de vèrtexs $V_G$, conjunt d'arestes $E_G$ i aplicacions origen i destí $s_G,t_G\colon E_G\to V_G$. Per a $A,B\in V_G$, un \textbf{camí}\index{camí} de $A$ a $B$ de longitud $n\ge 1$ és una successió $(e_1,\dots,e_n)$ d'arestes amb
\[
s_G(e_1)=A,\qquad t_G(e_i)=s_G(e_{i+1})\ \ (1\le i<n),\qquad t_G(e_n)=B.
\]
Per a cada vèrtex $A$ hi ha, a més, el \textbf{camí buit} $()_A$ de $A$ a $A$, de longitud zero.

La \textbf{categoria lliure generada pel graf} $G$, que escrivim $\Free(G)$,\index{Free@$\Free$}\index{categoria!lliure} té:
\begin{itemize}
\item objectes: $\Ob(\Free(G))=V_G$, els vèrtexs de $G$;
\item morfismes: $\Hom_{\Free(G)}(A,B)$ és el conjunt dels camins de $A$ a $B$;
\item composició: la \textbf{concatenació} en ordre de recorregut; per a $p=(a_1,\dots,a_r)\in\Hom(A,B)$ i $q=(b_1,\dots,b_s)\in\Hom(B,C)$,
\[
q\comp p=(a_1,\dots,a_r,b_1,\dots,b_s),
\]
on $p$, que surt de $A$, es recorre primer, d'acord amb la convenció $q\comp p$; concatenar amb un camí buit deixa el camí igual;
\item identitats: $\id_A=()_A$.
\end{itemize}
\end{definicio}

Per exemple, a partir del graf
\[
A\xrightarrow{f}B\xrightarrow{g}C
\]
la categoria lliure conté, a més de les identitats i dels camins $(f)$ i $(g)$ de longitud~1, que escrivim simplement $f$ i $g$, la fletxa composta $g\comp f=(f,g)\colon A\to C$:
\[
\begin{tikzpicture}[baseline=(m.center)]
  \matrix (m) [matrix of math nodes, column sep=3.5em] { A & B & C \\ };
  \draw[-{Stealth[scale=1.05]}] (m-1-1) -- node[above] {$f$} (m-1-2);
  \draw[-{Stealth[scale=1.05]}] (m-1-2) -- node[above] {$g$} (m-1-3);
  \draw[-{Stealth[scale=1.0]}] (m-1-1) to[bend right=25] node[below] {$g\comp f$} (m-1-3);
\end{tikzpicture}
\]
Els axiomes de categoria se satisfan. La composició és associativa, perquè $r\comp(q\comp p)$ i $(r\comp q)\comp p$ són totes dues la successió de totes les arestes de $p$, $q$ i $r$, en aquest ordre; i el camí buit és neutre, perquè concatenar-lo no afegeix cap aresta: $p\comp()_A=p=()_B\comp p$ per a tot camí $p$ de $A$ a $B$. Com que $V_G$ i $E_G$ són conjunts, també ho són els conjunts de camins, i $\Free(G)$ és una categoria petita. L'apèndix~\ref{cap:grafs} desenvolupa aquesta construcció amb més exemples. El nom \emph{lliure} quedarà justificat més endavant, quan disposem del llenguatge dels functors i puguem formular la propietat universal d'aquesta construcció.

\begin{exemple}[El monoide lliure i la convenció de composició]\label{ex:monoide-lliure-cat}
Prenem el cas d'un graf amb \textbf{un sol vèrtex} $\ast$ i conjunt d'arestes $E$. Aleshores totes les arestes són bucles ---comencen i acaben a $\ast$--- i cap parell no deixa d'estar encadenat, de manera que un camí és simplement una seqüència finita d'arestes, és a dir, una \emph{paraula} sobre $E$. Escrivim $E^*$ per al conjunt de totes aquestes paraules, incloent-hi la \emph{paraula buida} $\varepsilon$, de longitud zero.

Fixem tres convencions, que s'han de llegir juntes.
\begin{itemize}
\item \emph{Lectura de les paraules.} Escrivim les lletres en l'\emph{ordre de recorregut}: la paraula $a_1a_2\cdots a_n$ és el camí $(a_1,a_2,\dots,a_n)$, que recorre primer $a_1$, després $a_2$, i així successivament.
\item \emph{Concatenació.} Per a $u=a_1\cdots a_m$ i $v=b_1\cdots b_n$, la concatenació és la juxtaposició $uv=a_1\cdots a_m b_1\cdots b_n$: primer les lletres de $u$ i després les de $v$. Amb aquesta operació i la paraula buida, $(E^*,\text{concatenació},\varepsilon)$ és un monoide, el \textbf{monoide lliure}\index{monoide!lliure} sobre $E$: $(uv)w$ i $u(vw)$ són totes dues la paraula que juxtaposa $u$, $v$ i $w$ en aquest ordre, i $\varepsilon u=u=u\varepsilon$.
\item \emph{Composició.} A $\Free(G)$, la composició és la de la definició: $v\comp u$ vol dir \enquote{primer $u$, després $v$}, i el camí que en resulta recorre primer les lletres de $u$. Per tant,
\[
v\comp u=uv.
\]
\end{itemize}
Amb dues lletres $a,b\in E$, tot cap en una línia:
\[
b\comp a=ab\quad(\text{primer }a,\text{ després }b),
\qquad\qquad
a\comp b=ba\quad(\text{primer }b,\text{ després }a).
\]
L'ordre de la composició és, doncs, el \emph{contrari} del de la juxtaposició: el símbol de la dreta a $b\comp a$ és el que s'escriu primer a la paraula.

Quina categoria hem obtingut? Té un sol objecte, $\Hom_{\Free(G)}(\ast,\ast)=E^*$ i $\id_\ast=\varepsilon$. Però, amb la convenció de $\cat{B}M$, en què $g\comp f=g\cdot f$, la composició $v\comp u=uv$ correspon al producte \emph{invertit}. Si anomenem \textbf{monoide oposat}\index{monoide!oposat} $M\op$ d'un monoide $M$ el que té els mateixos elements i el mateix neutre i l'operació $m\cdot_{\mathrm{op}} n=n\cdot m$ ---de manera que $\cat{B}(M\op)=(\cat{B}M)\op$---, el que tenim és la \emph{igualtat}
\[
\Free(G)=\cat{B}\bigl((E^*)\op\bigr)=\bigl(\cat{B}(E^*)\bigr)\op,
\]
i no $\cat{B}(E^*)$. La diferència és només d'ordre d'escriptura. L'aplicació que inverteix l'ordre de les lletres, $a_1\cdots a_n\mapsto a_n\cdots a_1$, compleix $(uv)^{\mathrm{rev}}=v^{\mathrm{rev}}u^{\mathrm{rev}}$, i és per tant un isomorfisme de monoides $E^*\cong(E^*)\op$. Dona, doncs, un isomorfisme canònic de categories $\cat{B}(E^*)\cong\Free(G)$. Si haguéssim escrit les paraules en l'ordre de la composició ($a_n\cdots a_1$ per al camí que comença per $a_1$), hauríem obtingut literalment $\cat{B}(E^*)$. Hem preferit l'ordre de recorregut, que és el natural per als camins. En aquest sentit, i llevat d'aquest isomorfisme, el monoide lliure és el cas particular de la categoria lliure per a un graf amb un sol vèrtex.
\end{exemple}

\begin{exemple}
La categoria lliure sobre el graf amb dos vèrtexs $A,B$ i una sola aresta $A\to B$ és la categoria finita $\mathbf{2}$. Si, en canvi, el graf té dues arestes en sentits oposats, $e\colon A\to B$ i $f\colon B\to A$, la categoria lliure té infinits morfismes: $e$, $f$, $f\comp e$, $e\comp f$, $e\comp f\comp e$, i així successivament, perquè cap relació no identifica camins diferents. Així, \emph{un graf finit no genera necessàriament una categoria lliure finita}.
\end{exemple}

Aquesta construcció, que genera l'estructura més econòmica a partir d'unes dades sense imposar cap relació que no hi vingui forçada, és un exemple del patró \emph{lliure}, que reapareix per a monoides, grups i moltes altres estructures i que, al capítol~\refcap{cap:adjuncions}{6}, es formalitza com una \emph{adjunció} entre un functor lliure i un functor oblit.

\section{Categories segons la mida}\index{categoria!petita}\index{categoria!localment petita}
\label{sec:mida}

Hem parlat de \enquote{col·lecció} d'objectes en lloc de \enquote{conjunt}, i ara podem precisar per què. En alguns exemples, els objectes són massa nombrosos per formar un conjunt.

\begin{observacio}[La paradoxa de Russell]
La categoria $\cat{Set}$ té per objectes \emph{tots} els conjunts. Ara bé, no pot existir el conjunt de tots els conjunts. En efecte, si $U$ fos aquest conjunt, en podríem separar el subconjunt $R=\{x\in U\mid x\notin x\}$ dels conjunts que no es pertanyen a si mateixos, i preguntar-nos si $R\in R$: si $R\in R$, la definició força $R\notin R$; i si $R\notin R$, la definició força $R\in R$. Totes dues possibilitats són contradictòries. Aquesta és la \textbf{paradoxa de Russell}, i mostra que la col·lecció de tots els conjunts és massa gran per ser un conjunt.
\end{observacio}

Per això, en parlar dels objectes d'una categoria, diem \enquote{col·lecció} i no \enquote{conjunt}: pot ser massa gran. Segons la mida d'aquestes col·leccions introduïm els conceptes següents.

\begin{definicio}
Una categoria és \textbf{petita} si la seva col·lecció d'objectes i la seva col·lecció de morfismes són tots dos conjunts. En cas contrari es diu \textbf{gran}\index{categoria!gran}. I és \textbf{localment petita} si, per a cada parella d'objectes $A,B$, la col·lecció de morfismes $\Hom(A,B)$ és un conjunt, encara que la col·lecció de tots els objectes pugui ser massa gran per ser-ho.
\end{definicio}

Convé no llegir aquestes tres nocions com una partició en tres classes disjuntes. La partició real és \emph{petita} contra \emph{gran}. La petitesa local és una condició transversal: tota categoria petita és localment petita ---si tots els morfismes formen un conjunt, també en formen un els de cada parella $(A,B)$---, i una categoria gran pot ser localment petita o no. Així,
\[
\text{petita}\;\Longrightarrow\;\text{localment petita},
\]
i el recíproc és fals: $\cat{Set}$ és localment petita però no petita. Que no és petita ho mostra la paradoxa de Russell; que és localment petita es comprova identificant cada aplicació amb el seu graf (exercici~\ref{exr:pr-set-localment-petita} de la Pràctica). Les tres possibilitats efectives són, doncs, petita (per tant localment petita), gran localment petita i gran no localment petita.

Moltes de les categories habituals que utilitzem ---$\cat{Set}$, $\cat{Grp}$, $\cat{Top}$, $\cat{Graph}$ i altres--- són localment petites però no petites: tenen massa objectes per formar un conjunt, però entre dos objectes qualssevol hi ha només un conjunt de morfismes. En canvi, cada preordre concret sobre un conjunt i cada categoria finita concreta són categories petites. És important no confondre aquestes categories petites amb les categories grans que les apleguen, com $\cat{Pos}$, la categoria de tots els ordres parcials, o la categoria de totes les categories finites: en tots dos casos hi ha una classe pròpia de conjunts subjacents. Per la mateixa raó és gran $\cat{FinSet}$, tot i que cadascun dels seus objectes és un conjunt finit. Al llarg del llibre suposarem, tret que diguem el contrari, que les categories amb què treballem són localment petites; això garanteix que sempre podem parlar del conjunt $\Hom(A,B)$, cosa que serà essencial més endavant.\footnote{Aquí prenem com a base la teoria de conjunts clàssica, amb classes pròpies; l'apèndix~\ref{cap:mida} fixa aquest marc, i \cite[§1.8]{awodey}, \cite[§1.1]{riehl}, \cite[§3.2]{leinster} i \cite[cap.~1]{perrone} en donen presentacions introductòries. També és possible prendre el mateix llenguatge estructural com a fonament, sense una teoria de conjunts prèvia: l'apèndix~\ref{cap:axiomatica-relacional} n'ofereix un tast, i \cite[cap.~3]{leinster} en fa una introducció.}

\begin{resumcap}
\apartat{Idees centrals}
\begin{enumerate}
  \item Una categoria organitza objectes i morfismes mitjançant identitats i una composició associativa. El que importa no és què són els objectes, sinó com es relacionen a través de les fletxes. Els diagrames commutatius representen igualtats entre composicions.
  \item Una mateixa forma categòrica apareix en contextos molt diferents: conjunts i aplicacions, preordres, monoides, grafs, relacions i sistemes deductius.
  \item Els isomorfismes expressen quan dos objectes tenen la mateixa estructura dins la categoria; sovint, la noció estructuralment rellevant és l'isomorfisme més que no pas la igualtat estricta. Monomorfismes i epimorfismes es defineixen per cancel·lació. A $\cat{Set}$, i en algunes categories concretes, coincideixen amb les aplicacions injectives i exhaustives, però aquesta coincidència no és vàlida en general: en una categoria prima, per exemple, tot morfisme és mono i epi.
  \item Seccions i retraccions són inverses unilaterals. Un morfisme és un monomorfisme escindit si i només si té una retracció, i és un epimorfisme escindit si i només si té una secció. Tot monomorfisme escindit és mono i tot epimorfisme escindit és epi, però les implicacions inverses no valen en general:
  \[
  \text{iso}\Longrightarrow\text{mono escindit}\Longrightarrow\text{mono},
  \qquad
  \text{iso}\Longrightarrow\text{epi escindit}\Longrightarrow\text{epi}.
  \]
  \item A partir d'una categoria en podem construir d'altres: l'oposada, productes de categories, subcategories i categories sobre i sota un objecte.
  \item El \textbf{principi de dualitat} ---invertir totes les fletxes--- converteix sistemàticament cada definició i cada resultat en el seu dual, amb la composició en l'ordre invers.
  \item Un graf dirigit no té, per si sol, composició ni identitats; la categoria lliure hi afegeix formalment tots els camins i les identitats necessàries.
  \item Les qüestions de mida obliguen a distingir categories petites de categories localment petites. Tota categoria petita és localment petita, però moltes categories habituals, com $\cat{Set}$, són localment petites sense ser petites.
\end{enumerate}
\apartat{Errors típics} Creure que \enquote{mono $+$ epi $\Rightarrow$ iso} (fals a $\cat{Top}$ i a $\cat{2}$: un bimorfisme no és sempre un isomorfisme); confondre un monomorfisme amb un monomorfisme escindit; confondre la igualtat d'objectes amb l'isomorfisme, encara que sigui canònic; aplicar els adjectius \enquote{injectiu} i \enquote{exhaustiu} a morfismes d'una categoria general, en lloc de reservar-los per a aplicacions; raonar amb elements en una categoria on els objectes no en tenen o on no separen les fletxes; escriure la composició en l'ordre equivocat en dualitzar; i oblidar la hipòtesi de petitesa local en parlar d'homconjunts.
\apartat{Lectura recomanada} \cite[cap.~1, §2.1 i §3.1]{awodey} per a la introducció paral·lela, els monomorfismes i epimorfismes i el principi de dualitat; \cite[§1.1]{leinster} per a una presentació breu i moderna; \cite[cap.~I i II.1--II.3, II.7]{maclane} per a la formulació clàssica, inclosos la dualitat, les categories producte i les categories lliures; \cite[§1.3]{barrwells} per als grafs, en una presentació orientada a la informàtica; \cite[parts~I--II, sessions~2--10]{lawvereschanuel} per a una introducció encara més elemental als morfismes, els isomorfismes i les seccions i retraccions; i \cite[§1.1.5 i §1.2]{perrone-notes} per a contraexemples elementals d'isomorfismes, monomorfismes i epimorfismes, escindits o no. Per a la presentació relacional de les categories petites, l'apèndix~\ref{cap:axiomatica-relacional}.
\end{resumcap}

\chapter{Functors}
\label{cap:functors}

\begin{objectius}
\apartat{Prerequisits} El capítol~\ref{cap:categories} complet, especialment la definició de categoria, els diagrames commutatius, els isomorfismes, la categoria oposada i la dualitat, les categories producte i les categories sobre i sota un objecte, els grafs i les categories lliures, i la distinció entre categories petites i localment petites. Alguns exemples opcionals utilitzen grups, espais vectorials i l'espai dual $V^*=\Hom_{\mathbb K}(V,\mathbb K)$, i un exemple de functor oblit esmenta els espais topològics; cap d'ells no és imprescindible per seguir el capítol.
\apartat{Objectius} En acabar el capítol, el lector hauria de poder:
\begin{itemize}
  \item definir un functor especificant-ne l'acció sobre objectes i morfismes, i comprovar-ne la functorialitat;
  \item compondre functors i treballar amb la categoria $\cat{Cat}$ de les categories petites;
  \item reconèixer i construir exemples de functors en conjunts, preordres, lògica i grafs, inclosos el graf subjacent i la categoria lliure, i la relació entre $\Free$ i $U$;
  \item distingir functors covariants i contravariants i controlar-ne la variància;
  \item manejar els homfunctors $\Hom(A,-)$ i $\Hom(-,B)$ i el bihomfunctor $\Hom(-,-)$;
  \item distingir functors fidels, plens i essencialment exhaustius;
  \item separar amb precisió les propietats que un functor preserva de les que reflecteix o pot deixar de preservar.
\end{itemize}
\end{objectius}

\section{Cap al concepte de functor}

Un cop fixada la noció de categoria, la pregunta natural és quina mena de correspondència relaciona dues categories. Igual que entre grups mirem els homomorfismes ---les aplicacions que respecten l'estructura de grup--- i entre espais topològics mirem les aplicacions contínues ---les que respecten l'estructura topològica---, entre categories mirarem les correspondències que respecten l'estructura categòrica: objectes, morfismes, identitats i composició. S'anomenen \emph{functors}.

Fixem-nos que una categoria té dos nivells de dades: els objectes i els morfismes. Un functor haurà d'actuar sobre tots dos alhora, i fer-ho de manera compatible: ha d'enviar cada fletxa de la categoria $\cat{C}$ a una fletxa de la categoria $\cat{D}$ que vagi entre els objectes on van a parar el seu domini i el seu codomini, i ha de preservar la manera com les fletxes es componen i quines fletxes són identitats. Aquesta és exactament la idea que fa que els functors siguin, a la vegada, els morfismes d'una categoria de categories.

Abans de formalitzar-ho, val la pena veure un exemple concret que surt d'una operació ben familiar amb conjunts. Treballem a $\cat{Set}$ i fixem un conjunt $S$. A cada conjunt $A$ li podem associar el conjunt de totes les aplicacions de $S$ cap a $A$,
\[
F(A)=A^S=\{\,g\mid g\colon S\to A\,\}.
\]
Això assigna un objecte de $\cat{Set}$ a cada objecte de $\cat{Set}$; però un functor també ha d'actuar sobre les fletxes. I, de fet, tota aplicació $f\colon A\to B$ n'indueix una entre les potències,
\[
F(f)\colon A^S\longrightarrow B^S,
\qquad
F(f)(g)=f\comp g,
\]
és a dir, postcompondre amb $f$: si $g\colon S\to A$ és un element de $A^S$, aleshores $f\comp g\colon S\to B$ és un element de $B^S$, tal com mostra el triangle
\[
\begin{tikzpicture}[baseline=(m.center)]
  \matrix (m) [matrix of math nodes, column sep=3.4em, row sep=2.2em] {
    S & A \\
      & B \\
  };
  \draw[-{Stealth[scale=1.1]}] (m-1-1) -- node[above] {$g$} (m-1-2);
  \draw[-{Stealth[scale=1.1]}] (m-1-2) -- node[right] {$f$} (m-2-2);
  \draw[-{Stealth[scale=1.1]}] (m-1-1) -- node[below left] {$f\comp g$} (m-2-2);
\end{tikzpicture}
\]
Igual que a la definició, $F$ transporta cada fletxa de $\cat{Set}$ a una fletxa entre les imatges dels seus extrems:
\[
\begin{tikzpicture}[baseline=(m.center)]
  \matrix (m) [matrix of math nodes, row sep=2.6em] {
    A \\
    B \\
  };
  \draw[-{Stealth[scale=1.1]}] (m-1-1) -- node[right] {$f$} (m-2-1);
\end{tikzpicture}
\qquad\stackrel{F}{\longmapsto}\qquad
\begin{tikzpicture}[baseline=(m.center)]
  \matrix (m) [matrix of math nodes, row sep=2.6em] {
    A^S \\
    B^S \\
  };
  \draw[-{Stealth[scale=1.1]}] (m-1-1) -- node[right] {$F(f)$} (m-2-1);
\end{tikzpicture}
\]
on cada $g\in A^S$ és una aplicació $g\colon S\to A$ i $F(f)(g)=f\comp g$.
Aquesta assignació compleix just les dues condicions que esperaríem d'un functor. D'una banda, preserva la composició: si $A\xrightarrow{f}B\xrightarrow{h}C$, aleshores, per a tot $g\in A^S$,
\[
\begin{aligned}
F(h\comp f)(g)&=(h\comp f)\comp g\\
&=h\comp(f\comp g)=F(h)\bigl(F(f)(g)\bigr),
\end{aligned}
\]
de manera que
\[
F(h\comp f)=F(h)\comp F(f);
\]
l'única cosa que hem fet servir és l'associativitat de la composició. De l'altra, preserva les identitats: $F(\id_A)(g)=\id_A\comp g=g$, és a dir $F(\id_A)=\id_{A^S}$. Tenim, doncs, una manera sistemàtica de transformar objectes \emph{i} fletxes de $\cat{Set}$ que respecta composició i identitats: exactament el que a la secció següent anomenarem un functor. (Aquest exemple reapareixerà: $A^S$ no és res més que $\Hom_{\cat{Set}}(S,A)$, i el retrobarem com a homfunctor covariant a la definició~\ref{def:homfunctors}.)

\section{Definició de functor}

\begin{definicio}\index{functor}
\label{def:functor}
Siguin $\cat{C}$ i $\cat{D}$ dues categories. Un \textbf{functor} $F\colon\cat{C}\to\cat{D}$ assigna:
\begin{itemize}
\item a cada objecte $A$ de $\cat{C}$, un objecte $F(A)$ de $\cat{D}$;
\item a cada morfisme $f\colon A\to B$ de $\cat{C}$, un morfisme $F(f)\colon F(A)\to F(B)$ de $\cat{D}$,
\end{itemize}
de manera que es compleixin les dues condicions següents:
\begin{enumerate}
\item \textbf{Preservació d'identitats.} Per a cada objecte $A$ de $\cat{C}$,
\[
F(\id_A)=\id_{F(A)}.
\]
\item \textbf{Preservació de la composició.} Per a cada parella de morfismes componibles $f\colon A\to B$ i $g\colon B\to C$ de $\cat{C}$,
\[
F(g\comp f)=F(g)\comp F(f).
\]
\end{enumerate}
\end{definicio}

Observem que la condició sobre morfismes ja imposa que $F$ \emph{respecti} dominis i codominis: si $f\colon A\to B$, aleshores $F(f)$ ha d'anar de $F(A)$ a $F(B)$. La preservació de la composició es pot llegir diagramàticament: un triangle commutatiu de $\cat{C}$ es transforma en un triangle commutatiu de $\cat{D}$.
\[
\begin{tikzpicture}[baseline=(m.center)]
  \matrix (m) [matrix of math nodes, column sep=3.2em, row sep=2.6em] {
    A & B \\
      & C \\
  };
  \draw[-{Stealth[scale=1.1]}] (m-1-1) -- node[above] {$f$} (m-1-2);
  \draw[-{Stealth[scale=1.1]}] (m-1-2) -- node[right] {$g$} (m-2-2);
  \draw[-{Stealth[scale=1.1]}] (m-1-1) -- node[below left] {$g\comp f$} (m-2-2);
\end{tikzpicture}
\qquad\stackrel{F}{\longmapsto}\qquad
\begin{tikzpicture}[baseline=(m.center)]
  \matrix (m) [matrix of math nodes, column sep=3.6em, row sep=2.6em] {
    F(A) & F(B) \\
         & F(C) \\
  };
  \draw[-{Stealth[scale=1.1]}] (m-1-1) -- node[above] {$F(f)$} (m-1-2);
  \draw[-{Stealth[scale=1.1]}] (m-1-2) -- node[right] {$F(g)$} (m-2-2);
  \draw[-{Stealth[scale=1.1]}] (m-1-1) -- node[below left] {$F(g\comp f)$} (m-2-2);
\end{tikzpicture}
\]

\begin{observacio}[Primer els tipus, després les identitats]\label{obs:comprovar-tipus}\index{functor!comprovació de tipus}
Per comprovar que una assignació és un functor convé seguir sempre el mateix ordre, i no passar a les equacions fins que els tipus no quadren.
\begin{enumerate}
\item \emph{Objectes.} Cada $F(A)$ ha de ser un objecte de $\cat{D}$.
\item \emph{Tipus dels morfismes.} Per a cada $f\colon A\to B$, cal que $F(f)$ sigui un morfisme de $\cat{D}$ \emph{i} que vagi de $F(A)$ a $F(B)$. Aquí és on fallen més sovint les construccions candidates.
\item \emph{Identitats i composició.} Només aleshores té sentit comprovar $F(\id_A)=\id_{F(A)}$ i $F(g\comp f)=F(g)\comp F(f)$. El pas~2 garanteix que els dos membres de cada equació tenen el mateix tipus ($F(A)\to F(A)$ i $F(A)\to F(C)$) i que la composició $F(g)\comp F(f)$ existeix, perquè el codomini de $F(f)$ i el domini de $F(g)$ són tots dos $F(B)$.
\end{enumerate}
Un exemple mostra per què l'ordre importa. Si a cada conjunt $X$ li assignem $\mathcal P(X)$ i a cada aplicació $f\colon X\to Y$ l'antiimatge $B\mapsto f^{-1}(B)$, el pas~2 ja falla: l'antiimatge va de $\mathcal P(Y)$ a $\mathcal P(X)$, no de $\mathcal P(X)$ a $\mathcal P(Y)$. No cal, doncs, comprovar cap equació per saber que això no és un functor covariant $\cat{Set}\to\cat{Set}$. Sí que és, en canvi, un functor \emph{contravariant}, la noció que introduirem a la secció~\ref{sec:variancia}, per a la qual el pas~2 demana precisament $F(f)\colon F(B)\to F(A)$ (exemple~\ref{ex:parts-dues-variancies}). Les solucions de la Pràctica segueixen aquest ordre; vegeu, per exemple, l'exercici resolt~\ref{pr-hom-cov}.
\end{observacio}

\begin{observacio}[Els functors preserven els diagrames commutatius]\index{diagrama commutatiu!preservat per functors}
Una conseqüència directa de la preservació de la composició és que un functor envia diagrames commutatius a diagrames commutatius. En efecte, dir que un diagrama commuta és dir que certes composicions de fletxes coincideixen; si dos camins amb el mateix origen i el mateix destí donen la mateixa composició a $\cat{C}$, aleshores, aplicant $F$ i usant que preserva la composició, les seves imatges també coincideixen a $\cat{D}$. Per això podem aplicar un functor a tot un diagrama i seguir raonant amb la seva imatge: les igualtats entre camins es conserven. Ho comprovarem en detall, per al cas d'un quadrat, a l'exercici~\ref{prac:functor-quadrat} del capítol~\ref{cap:practica-functors} de Pràctica.
\end{observacio}

\subsection{Els functors preserven els isomorfismes}\label{sec:preserven-isos}

Com que els functors preserven tota l'estructura categòrica, en particular preserven les equacions entre composicions. Una conseqüència immediata és que preserven els isomorfismes.

\begin{proposicio}\label{prop:functor-preserva-iso}\index{functor!preserva isomorfismes}
Si $F\colon\cat{C}\to\cat{D}$ és un functor i $f\colon A\to B$ és un isomorfisme a $\cat{C}$, aleshores $F(f)$ és un isomorfisme a $\cat{D}$, i
\[
\bigl(F(f)\bigr)^{-1}=F(f^{-1}).
\]
\end{proposicio}

\begin{prova}
Sigui $g=f^{-1}$, de manera que $g\comp f=\id_A$ i $f\comp g=\id_B$. Aplicant $F$ i usant que preserva composició i identitats,
\[
F(g)\comp F(f)=F(g\comp f)=F(\id_A)=\id_{F(A)},
\]
\[
F(f)\comp F(g)=F(f\comp g)=F(\id_B)=\id_{F(B)}.
\]
Per tant $F(g)$ és un invers de $F(f)$, i pel resultat d'unicitat de l'invers, $\bigl(F(f)\bigr)^{-1}=F(g)=F(f^{-1})$.
\end{prova}

\begin{observacio}
Val la pena veure per què és així, més enllà del càlcul de la prova. Dir que $f$ és un isomorfisme és dir que existeix $g$ amb $g\comp f=\id_A$ i $f\comp g=\id_B$: dues equacions entre composicions, és a dir, dos diagrames commutatius. Com que els functors preserven els diagrames commutatius, aquestes dues equacions es transporten a $\cat{D}$ en aplicar-hi $F$, i $F(g)$ resulta ser l'invers de $F(f)$.
\end{observacio}

\section{Exemples de functors}\label{sec:primers-exemples}

La força de la definició, com en el cas de les categories, es veu en la varietat d'exemples que hi encaixen. Recuperem els exemples de referència del capítol anterior.

\begin{exemple}[Functor identitat]\index{functor!identitat}
Per a tota categoria $\cat{C}$ hi ha el \textbf{functor identitat} $\id_{\cat{C}}\colon\cat{C}\to\cat{C}$, que deixa fixos tots els objectes i tots els morfismes. Preserva trivialment identitats i composició.
\end{exemple}

\begin{exemple}[Functor constant]\label{ex:functor-constant}\index{functor!constant}
Siguin $\cat{C}$ i $\cat{D}$ dues categories i $D$ un objecte de $\cat{D}$. El \textbf{functor constant} amb valor $D$,
\[
\Delta_D\colon\cat{C}\to\cat{D},
\]
envia tot objecte $A$ de $\cat{C}$ a $D$ i tot morfisme $f$ de $\cat{C}$ a la identitat $\id_D$. És un functor: $\Delta_D(\id_A)=\id_D=\id_{\Delta_D(A)}$, i $\Delta_D(g\comp f)=\id_D=\id_D\comp\id_D=\Delta_D(g)\comp\Delta_D(f)$. Quan $\cat{D}=\cat{Set}$ i $D$ és un conjunt d'un sol element $\{\ast\}$, escrivim $\Delta_{\{\ast\}}$ per al functor constant que envia cada objecte a aquest conjunt d'un sol element.
\end{exemple}

\begin{exemple}[Functors oblit]\index{functor!oblit}
Moltes categories d'objectes amb estructura $\cat{C}$ tenen un \textbf{functor oblit} $U\colon\cat{C}\to\cat{Set}$, que oblida l'estructura i reté només el conjunt subjacent:
\[
\cat{Grp}\to\cat{Set},
\qquad
\cat{Top}\to\cat{Set},
\qquad
\cat{Vect}_{\mathbb K}\to\cat{Set}.
\]
En cada cas, $U$ envia un objecte al seu conjunt subjacent i un morfisme ---un homomorfisme, una aplicació contínua, una aplicació lineal--- a la mateixa aplicació, vista simplement com a aplicació de conjunts. La identitat es preserva perquè l'aplicació identitat continua sent la identitat, i la composició es preserva perquè és la mateixa composició d'aplicacions. El functor oblit \enquote{no fa res} als morfismes; el seu contingut és que oblida l'estructura dels objectes.
\end{exemple}

\begin{exemple}[Functors des de les categories finites]\index{functor!des d'una categoria finita}
Les categories finites $\cat{0}$, $\cat{1}$ i $\cat{2}$ (subsecció~\ref{subsec:finites}) donen, com a dominis, functors amb una lectura molt directa. Un functor $\cat{1}\to\cat{D}$, on $\cat{1}$ té un únic objecte i cap morfisme no trivial, no és res més que la tria d'un objecte de $\cat{D}$: la imatge de l'únic objecte. Un functor $\cat{2}\to\cat{D}$, on $\cat{2}=(0\to 1)$ ---reanomenem els dos objectes de $\cat{2}$ com $0$ i $1$, i escrivim $(0\to 1)$ per indicar la seva única fletxa no trivial---, és la tria d'un morfisme de $\cat{D}$: cal dir on van $0$ i $1$ ---dos objectes--- i on va l'única fletxa no trivial ---un morfisme entre les seves imatges---, i la preservació d'identitats i composició no imposa res més. Finalment, des de la categoria buida $\cat{0}$ hi ha un únic functor $\cat{0}\to\cat{D}$, el functor buit, sigui quina sigui $\cat{D}$. Aquests tres fets ---objectes, morfismes i \enquote{res}--- reapareixeran contínuament: seleccionar un objecte o un morfisme d'una categoria és el mateix que donar un functor des de $\cat{1}$ o des de $\cat{2}$.
\end{exemple}

\begin{exemple}[Functors entre categories discretes]\index{functor!entre categories discretes}\index{categoria!discreta}
Un conjunt es pot veure com una categoria discreta, en què els únics morfismes són les identitats (subsecció~\ref{subsec:conjunts}). Si $S$ i $T$ són dos conjunts vistos així, un functor $F\colon S\to T$ només ha de dir on va cada objecte, perquè els únics morfismes són les identitats i tot functor les preserva automàticament. Per tant un functor entre categories discretes és exactament una aplicació de conjunts $S\to T$: sense estructura a preservar, la noció de functor es redueix a la d'aplicació.
\end{exemple}

\begin{exemple}[Functors entre preordres]\index{functor!entre preordres}
Siguin $(P,\leq)$ i $(Q,\leq)$ dos preordres, vistos com a categories. Un functor $F\colon P\to Q$ assigna a cada element $x\in P$ un element $F(x)\in Q$ i, com que hi ha un morfisme $x\to y$ exactament quan $x\leq y$, la condició sobre morfismes diu que
\[
x\leq y
\quad\Longrightarrow\quad
F(x)\leq F(y).
\]
És a dir, un functor entre preordres és exactament una aplicació monòtona. Les condicions d'identitat i composició es compleixen automàticament, perquè en un preordre entre dos objectes hi ha com a màxim una fletxa i no hi ha res més a comprovar.
\end{exemple}

\begin{exemple}[Functors des d'un monoide]\label{ex:functors-monoide}\index{functor!des d'un monoide}
Recordem que un monoide $(M,\cdot,1)$ es pot veure com una categoria d'un sol objecte $\ast$, que escrivim $\cat{B}M$. Un functor $F\colon\cat{B}M\to\cat{D}$ tria un objecte $D=F(\ast)$ de $\cat{D}$ i envia cada element $m\in M$ ---és a dir, cada morfisme $\ast\to\ast$--- a un morfisme $F(m)\colon D\to D$, de manera que
\[
F(1)=\id_D
\qquad\text{i}\qquad
F(m\cdot n)=F(m)\comp F(n).
\]
Aquestes són exactament les condicions perquè $F$ sigui un \textbf{homomorfisme de monoides} de $M$ cap al monoide dels endomorfismes de $D$,
\[
\mathrm{End}_{\cat{D}}(D)=\Hom_{\cat{D}}(D,D).
\]
En particular, quan $\cat{D}=\cat{Set}$, un functor $\cat{B}M\to\cat{Set}$ és una \textbf{acció} del monoide $M$ sobre un conjunt. Vegem per què. Un tal functor queda determinat per un conjunt $F(\ast)=S$ i, per a cada element $m\in M$, una funció $F(m)\colon S\to S$. Les dues condicions de functor ---preservar identitats i composició--- diuen exactament que, per a tot $x\in S$ i tots $m,n\in M$,
\[
F(1)(x)=x,
\qquad
F(m\cdot n)(x)=F(m)\bigl(F(n)(x)\bigr),
\]
Si escrivim $m\cdot x$ en lloc de $F(m)(x)$, aquestes igualtats esdevenen les condicions habituals d'una \emph{acció per l'esquerra} del monoide sobre $S$:
\[
1\cdot x=x,
\qquad
(m\cdot n)\cdot x=m\cdot(n\cdot x).
\]
\end{exemple}

\begin{exemple}[Functors des d'un grup: accions i representacions]\index{functor!des d'un grup}\index{accio@acció!de grup}\index{representacio@representació lineal}
Quan el monoide és un grup $G$, la categoria $\cat{B}G$ és un grupoide: tot morfisme és un isomorfisme. Un functor $F\colon\cat{B}G\to\cat{D}$ és aleshores un homomorfisme de $G$ cap al grup dels automorfismes de l'objecte $D=F(\ast)$,
\[
F\colon G\longrightarrow \Aut_{\cat{D}}(D),
\]
perquè un functor preserva isomorfismes (proposició~\ref{prop:functor-preserva-iso}) i, per tant, cada element de $G$ va a parar a un automorfisme de $D$. Dos casos són especialment importants:
\begin{itemize}
\item Un functor $\cat{B}G\to\cat{Set}$ és una \textbf{acció} del grup $G$ sobre un conjunt ---un \emph{$G$-conjunt}. El desplegament és el mateix que per a un monoide (exemple~\ref{ex:functors-monoide}), amb una novetat: com que ara cada $a\in G$ té un invers $a^{-1}$, la funció $F(a)\colon S\to S$ és bijectiva, amb inversa $F(a^{-1})$. Cada element del grup actua, doncs, com una \emph{permutació} de $S$. Un exemple típic són les simetries d'una figura, que actuen permutant-ne els vèrtexs.
\item Un functor $\cat{B}G\to\cat{Vect}_{\mathbb K}$ és una \textbf{representació lineal} de $G$ sobre el cos $\mathbb{K}$. El desplegament torna a ser el mateix, ara amb un espai vectorial $F(\ast)=V$ en lloc d'un conjunt: cada element $a\in G$ dona una aplicació \emph{lineal} invertible $F(a)\colon V\to V$ ---un automorfisme de $V$, amb inversa $F(a^{-1})$---, i el producte del grup es correspon amb la composició. Escrivint $a\cdot v$ per $F(a)(v)$, la novetat respecte d'una acció ordinària és que cada element actua respectant les operacions de l'espai:
\[
a\cdot(\lambda v+\mu w)=\lambda\,(a\cdot v)+\mu\,(a\cdot w).
\]
En dimensió finita, fixada una base, cada $F(a)$ és una matriu invertible: una representació \enquote{realitza} els elements abstractes del grup com a matrius concretes, cosa que permet estudiar el grup amb les eines de l'àlgebra lineal.
\end{itemize}
Així, dues teories aparentment allunyades ---les accions i representacions lineals d'un grup--- són, categòricament, el mateix tipus d'objecte: functors amb domini $\cat{B}G$, que només es distingeixen per la categoria d'arribada.
\end{exemple}

\begin{exemple}[Deduïbilitat i functors lògics]\label{ex:deduibilitat-functors}\index{functor!i deduïbilitat}\index{substitucio@substitució}
Siguin $T$ i $T'$ dos sistemes deductius, i $\cat{Prov}_T$ i $\cat{Prov}_{T'}$ les seves categories primes de deduïbilitat (subsecció~\ref{subsec:elementals}). Una assignació de fórmules $\varphi\colon A\mapsto \varphi(A)$ que preservi la deduïbilitat,
\[
A\vdash_T B
\quad\Longrightarrow\quad
\varphi(A)\vdash_{T'}\varphi(B),
\]
és exactament un functor $F\colon\cat{Prov}_T\to\cat{Prov}_{T'}$, amb $F(A)=\varphi(A)$: la fletxa que expressa $A\vdash_T B$ va a parar a la fletxa que expressa $\varphi(A)\vdash_{T'}\varphi(B)$. Com que aquestes categories són primes, no cal comprovar res més: la preservació de la deduïbilitat garanteix automàticament la preservació d'identitats i composicions. Aquest és el reflex categòric del fet que una traducció entre teories que respecti la deduïbilitat indueix un functor entre les categories associades.

Un cas concret i molt natural d'aquesta situació és el d'una \textbf{substitució} de variables proposicionals. Sigui $T_{\mathrm{prop}}$ el sistema deductiu de la lògica proposicional i $\cat{Prov}_{T_{\mathrm{prop}}}$ la seva categoria de deduïbilitat, amb les fórmules per objectes i una fletxa $A\to B$ exactament quan $A\vdash B$ (aquí escrivim $\vdash$ per $\vdash_{T_{\mathrm{prop}}}$). Considerem la substitució
\[
\sigma\colon\quad p\mapsto q,\qquad q\mapsto p\lor r,
\]
estesa a totes les fórmules respectant els connectors ($\land,\lor,\neg,\to$). Per exemple, $\sigma$ transforma $p\to(q\land p)$ en $q\to((p\lor r)\land q)$. El punt clau és que una substitució no sols transforma fórmules, sinó que també preserva la \emph{deduïbilitat}: si $A\vdash B$, aleshores $\sigma(A)\vdash\sigma(B)$, perquè qualsevol prova concreta de $B$ a partir de $A$ es tradueix, aplicant $\sigma$ a cada fórmula, en una derivació de $\sigma(B)$ a partir de $\sigma(A)$. Així, de la fletxa $p\land q\to p$ (perquè $p\land q\vdash p$) n'obtenim la fletxa $q\land(p\lor r)\to q$. Per tant $\sigma$ indueix un functor
\[
F_\sigma\colon\cat{Prov}_{T_{\mathrm{prop}}}\longrightarrow\cat{Prov}_{T_{\mathrm{prop}}},
\qquad
A\mapsto\sigma(A),
\qquad
(A\vdash B)\mapsto(\sigma(A)\vdash\sigma(B)),
\]
que és, precisament, el cas $T=T'=T_{\mathrm{prop}}$ de la construcció anterior (aquí amb la mateixa categoria de sortida i d'arribada).
\end{exemple}

\begin{observacio}
Si en comptes de la deduïbilitat considerem la \emph{categoria de proves} de $T$ ---els mateixos objectes, però ara amb les derivacions, mòdul equivalència, com a morfismes i el tall com a composició---, un functor cap a la categoria de proves de $T'$ demana també una traducció de cada derivació que en respecti l'encadenament i les identitats. A diferència del cas anterior, aquestes condicions ja no són automàtiques, perquè la categoria de proves no és prima.
\end{observacio}

\begin{exemple}[Functors oblit sobre i sota un objecte]\index{functor!oblit}\index{categoria!sobre un objecte}\index{categoria!sota un objecte}
Les categories de $\cat{C}$ sobre i sota un objecte $A$, $\cat{C}/A$ i $A/\cat{C}$ (definició~\ref{def:slice}), tenen un functor oblit cap a $\cat{C}$, un per a cadascuna. El functor
\[
U_A\colon\cat{C}/A\to\cat{C}
\]
envia cada objecte $x\colon X\to A$ al seu domini $X$ i cada morfisme (triangle) $h$ a ell mateix; oblida el morfisme distingit cap a $A$ i reté només el domini. Dualment, $A/\cat{C}\to\cat{C}$ envia cada objecte $x\colon A\to X$ al seu codomini $X$ i cada morfisme (triangle) a ell mateix; oblida el morfisme distingit des de $A$ i reté només el codomini.

Vegem-ho en un cas concret. Amb $\cat{2}=(0\to 1)$ i l'objecte $A=1$, els objectes de $\cat{2}/1$ són les fletxes de $\cat{2}$ amb codomini $1$: la identitat $\id_1$ i l'única fletxa no trivial $u\colon 0\to 1$. L'únic morfisme no trivial va de $u$ a $\id_1$, de manera que $\cat{2}/1$ és, de nou, una còpia de $\cat{2}$, i $U_1\colon\cat{2}/1\to\cat{2}$ envia $u\mapsto 0$ i $\id_1\mapsto 1$. Simètricament, sota l'objecte $0$, els objectes de $0/\cat{2}$ són les fletxes amb domini $0$ ---la identitat $\id_0$ i $u\colon 0\to 1$---, i el functor oblit $0/\cat{2}\to\cat{2}$ envia $\id_0\mapsto 0$ i $u\mapsto 1$.
\end{exemple}

\begin{exemple}[Projeccions i functor diagonal]\label{ex:projeccions-diagonal}\index{functor!de projecció}\index{functor!diagonal}
La categoria producte $\cat{C}\times\cat{D}$ (subsecció~\ref{subsec:categoria-producte}) té dos functors de \textbf{projecció}
\[
\pi_1\colon\cat{C}\times\cat{D}\to\cat{C},\qquad \pi_2\colon\cat{C}\times\cat{D}\to\cat{D},
\]
que es queden amb una de les components: $\pi_1(A,X)=A$ i $\pi_1(f,g)=f$, i anàlogament per a $\pi_2$. Són functors perquè la composició i les identitats de $\cat{C}\times\cat{D}$ es defineixen component a component:
\[
\begin{aligned}
\pi_1\bigl((f',g')\comp(f,g)\bigr)
 &= \pi_1(f'\comp f,\,g'\comp g)\\
 &= f'\comp f
  = \pi_1(f',g')\comp\pi_1(f,g),\\
\pi_1(\id_A,\id_X)&=\id_A,
\end{aligned}
\]
i el càlcul per a $\pi_2$ és el mateix sobre la segona component.
En sentit contrari, cada categoria $\cat{C}$ té un functor \textbf{diagonal}
\[
\Delta\colon\cat{C}\to\cat{C}\times\cat{C},\qquad \Delta(A)=(A,A),\qquad \Delta(f)=(f,f),
\]
que duplica objectes i morfismes; la functorialitat és de nou immediata, perquè $\Delta(g\comp f)=(g\comp f,\,g\comp f)=(g,g)\comp(f,f)$ i $\Delta(\id_A)=(\id_A,\id_A)$. Les projeccions i la diagonal encaixen en el diagrama commutatiu
\[
\begin{tikzpicture}[baseline=(m.center)]
  \matrix (m) [matrix of math nodes, row sep=2.6em, column sep=3.4em] {
    & \cat{C} & \\
    \cat{C} & \cat{C}\times\cat{C} & \cat{C} \\
  };
  \draw[-{Stealth[scale=1.1]}] (m-1-2) -- node[right] {$\Delta$} (m-2-2);
  \draw[-{Stealth[scale=1.1]}] (m-2-2) -- node[below] {$\pi_1$} (m-2-1);
  \draw[-{Stealth[scale=1.1]}] (m-2-2) -- node[below] {$\pi_2$} (m-2-3);
  \draw[-{Stealth[scale=1.1]}] (m-1-2) -- node[above left] {$\id_{\cat{C}}$} (m-2-1);
  \draw[-{Stealth[scale=1.1]}] (m-1-2) -- node[above right] {$\id_{\cat{C}}$} (m-2-3);
\end{tikzpicture}
\qquad
\pi_1\comp\Delta=\id_{\cat{C}}=\pi_2\comp\Delta,
\]
on la composició de functors es fa com a la secció següent: les dues igualtats es comproven directament, ja que $\pi_1(\Delta(A))=\pi_1(A,A)=A$ i $\pi_1(\Delta(f))=\pi_1(f,f)=f$, i igual per a $\pi_2$. La notació concorda amb la del functor constant $\Delta_D$ de l'exemple~\ref{ex:functor-constant}: més endavant, el functor diagonal s'estendrà a un functor que envia cada objecte $A$ a un diagrama constant de valor $A$, del qual $(A,A)$ és el cas de dues còpies. Són exemples elementals, però reapareixeran sovint: les projeccions fan de $\cat{C}\times\cat{D}$ el producte de $\cat{C}$ i $\cat{D}$ dins de $\cat{Cat}$ quan totes dues són petites (observació~\ref{obs:dos-productes}), i el functor diagonal és l'ingredient bàsic per tractar productes, límits i adjuncions en capítols posteriors.
\end{exemple}

\section{Composició de functors i la categoria \texorpdfstring{$\cat{Cat}$}{Cat}}\index{Cat@$\cat{Cat}$}

Els functors es poden compondre, i la composició torna a ser un functor. Això és el que permet organitzar les categories i els functors en una nova categoria.

\begin{proposicio}
Siguin $F\colon\cat{C}\to\cat{D}$ i $G\colon\cat{D}\to\cat{E}$ dos functors. La seva \textbf{composició} $G\comp F$, definida per
\[
(G\comp F)(A)=G(F(A))
\qquad\text{i}\qquad
(G\comp F)(f)=G(F(f)),
\]
és un functor $\cat{C}\to\cat{E}$.
\end{proposicio}

\begin{prova}
Cal comprovar les dues condicions. Per a les identitats, i usant primer que $F$ i després que $G$ les preserven,
\[
(G\comp F)(\id_A)=G(F(\id_A))=G(\id_{F(A)})=\id_{G(F(A))}=\id_{(G\comp F)(A)}.
\]
Per a la composició, siguin $f\colon A\to B$ i $g\colon B\to C$ a $\cat{C}$. Aleshores
\[
\begin{aligned}
(G\comp F)(g\comp f)
&=G\big(F(g\comp f)\big)
=G\big(F(g)\comp F(f)\big)\\
&=G(F(g))\comp G(F(f))
=(G\comp F)(g)\comp(G\comp F)(f).
\end{aligned}
\]
Per tant $G\comp F$ és un functor.
\end{prova}

La composició de functors és associativa i el functor identitat actua com a element neutre, exactament pels mateixos motius que a $\cat{Set}$. En efecte, un functor $F\colon\cat{C}\to\cat{D}$ consta de dues aplicacions ---una sobre objectes i una sobre morfismes---, i la composició de functors no és res més que la composició d'aquestes aplicacions component a component. Per tant, l'associativitat i les lleis d'identitat es redueixen a les propietats corresponents de la composició d'aplicacions, que ja coneixem. Això ens permet donar la definició següent.

\begin{definicio}\index{Cat@$\cat{Cat}$}
La categoria $\cat{Cat}$ té per objectes les categories petites i per morfismes els functors entre elles. La composició és la composició de functors i la identitat de cada categoria és el seu functor identitat.
\end{definicio}

\begin{observacio}[Qüestions de mida]\label{obs:mida-cat}
La categoria $\cat{Cat}$ té per objectes les categories \emph{petites}. Aquesta restricció la fa \emph{localment petita}: si $\cat{C}$ és petita, $\Ob(\cat{C})$ i $\Mor(\cat{C})$ són conjunts, i doncs, donades dues categories petites $\cat{C}$ i $\cat{D}$, el producte $\Ob(\cat{D})^{\Ob(\cat{C})}\times \Mor(\cat{D})^{\Mor(\cat{C})}$ és un conjunt; com que
\[
\Hom_{\cat{Cat}}(\cat{C},\cat{D})\subseteq \Ob(\cat{D})^{\Ob(\cat{C})}\times \Mor(\cat{D})^{\Mor(\cat{C})},
\]
també ho és $\Hom_{\cat{Cat}}(\cat{C},\cat{D})$. La mateixa $\cat{Cat}$, en canvi, és \emph{gran}: les categories petites formen una classe pròpia, no un conjunt.

Categories d'ús constant com $\cat{Set}$, $\cat{Grp}$ o $\cat{Top}$ són grans i, per tant, no són objectes de $\cat{Cat}$. Això no impedeix treballar amb functors entre elles: un functor oblit $\cat{Grp}\to\cat{Set}$, o el functor lliure $\Free\colon\cat{Graph}\to\cat{Cat}$, és simplement una assignació sobre objectes i una altra sobre morfismes que compleixen les condicions functorials, i la seva legitimitat no depèn de cap categoria que tingui les categories grans per objectes.

Al llarg del text, doncs, $\cat{Cat}$ designarà sempre la categoria de les categories \emph{petites}, i les categories grans, com $\cat{Set}$, no en són objectes.\footnotemark
\end{observacio}
\footnotetext{Aquesta afirmació és relativa al marc de mida adoptat (apèndix~\ref{cap:mida}). Amb una jerarquia d'universos de Grothendieck, una categoria gran respecte d'un univers pot ser petita respecte d'un univers més gran, i aleshores és objecte de la categoria de categories d'aquest nivell; vegeu \cite[§1.1]{riehl}.}%

Un cop els functors són els morfismes de $\cat{Cat}$, podem aplicar a les categories les mateixes nocions categòriques que ja coneixem. En particular, un isomorfisme a $\cat{Cat}$ és un functor que té un invers estricte.

\begin{definicio}[Isomorfisme de categories]\label{def:iso-categories}\index{isomorfisme!de categories}\index{categoria!isomorfa}
Dues categories $\cat{C}$ i $\cat{D}$ són \textbf{isomorfes}, i escrivim $\cat{C}\cong\cat{D}$, si hi ha functors $F\colon\cat{C}\to\cat{D}$ i $G\colon\cat{D}\to\cat{C}$ inversos l'un de l'altre, és a dir, amb $G\comp F=\id_{\cat{C}}$ i $F\comp G=\id_{\cat{D}}$; un tal $F$ és un \textbf{isomorfisme de categories}. Per a categories petites no és res més que un isomorfisme (definició~\ref{def:isomorfisme}) a $\cat{Cat}$, però la definició val igualment per a categories grans, que no són objectes de $\cat{Cat}$.
\end{definicio}

\section{Functors i grafs}\label{sec:functors-grafs}\index{functor!i grafs}\index{morfisme de grafs}

Val la pena comparar els functors amb els \emph{morfismes de grafs}, que vam introduir en definir $\cat{Graph}$ (subsecció~\ref{subsec:grafs}). Un morfisme de grafs $\alpha\colon G\to H$ és una parella d'aplicacions $\alpha_V\colon V_G\to V_H$ i $\alpha_E\colon E_G\to E_H$ compatibles amb les dues aplicacions d'origen i destí $s,t\colon E\to V$ de cada graf; res més, perquè un graf no té composició ni identitats. Un functor entre categories petites demana exactament aquestes mateixes dues dades ---una assignació sobre objectes i una sobre morfismes, compatibles amb domini i codomini--- però hi afegeix dues condicions que un graf no pot ni formular: la preservació de la composició i de les identitats. Un morfisme de grafs és, doncs, la part «sense composició ni identitats» d'un functor. Aquesta relació es concreta en dues construccions en sentits contraris entre $\cat{Graph}$ i $\cat{Cat}$, que estudiem tot seguit.

\subsection{El graf subjacent d'una categoria}

Tota categoria petita $\cat{C}$ té un \textbf{graf subjacent}: el graf que té per vèrtexs els objectes de $\cat{C}$ i per arestes \emph{tots} els seus morfismes, amb les aplicacions origen i destí donades pel domini i el codomini. En passar a aquest graf, oblidem quins morfismes són identitats i com es componen, però no eliminem les arestes corresponents: les fletxes identitat hi resten com a arestes, una a cada vèrtex, tot i que ja no marcades com a identitats.

Aquesta assignació és functorial. Donat un functor $F\colon\cat{C}\to\cat{D}$, obtenim un morfisme entre els grafs subjacents que actua, sobre vèrtexs, com $F$ sobre els objectes, i sobre arestes, com $F$ sobre els morfismes. És un morfisme de grafs ben definit precisament perquè $F$ respecta domini i codomini: una aresta $f\colon A\to B$ va a $F(f)\colon F(A)\to F(B)$, una aresta entre les imatges dels extrems. Per construir-lo n'hi ha prou amb això; la preservació de la composició i de les identitats no s'hi fa servir, i per això diem que «se n'oblida». Tenim, així, un functor oblit
\[
U\colon\cat{Cat}\to\cat{Graph},
\]
que envia cada categoria petita al seu graf subjacent $U(\cat{C})$ i cada functor $F$ al morfisme de grafs $U(F)$ que acabem de descriure. El recíproc del que fa $U$ és fals: un morfisme entre els grafs subjacents no respecta, en general, ni la composició ni les identitats, de manera que un functor és estrictament més que un morfisme de grafs. Vegem-ho amb un bucle. Sigui $M=(\{1,e\},\cdot,1)$ amb $e\cdot e=e$, vist com a categoria d'un sol objecte $\ast$; el functor oblit en dona el graf subjacent:
\[
\begin{tikzpicture}[baseline=(m.center)]
  \matrix (m) [cat matrix] { \ast \\ };
  \draw[cat] (m-1-1) to[out=55,in=125,looseness=9] node[above] {\(\id_\ast\)} (m-1-1);
  \draw[cat] (m-1-1) to[out=-125,in=-55,looseness=9] node[below] {\(e\)} (m-1-1);
\end{tikzpicture}
\qquad\stackrel{U}{\longmapsto}\qquad
\begin{tikzpicture}[baseline=(m.center)]
  \matrix (m) [cat matrix] { \ast \\ };
  \draw[cat] (m-1-1) to[out=55,in=125,looseness=9] node[above] {\(\id_\ast\)} (m-1-1);
  \draw[cat] (m-1-1) to[out=-125,in=-55,looseness=9] node[below] {\(e\)} (m-1-1);
\end{tikzpicture}
\]
Sobre aquest graf ---que conserva les dues arestes $\id_\ast$ i $e$, però on $U$ ha oblidat tota la composició: les igualtats $\id_\ast\comp\id_\ast=\id_\ast$, $\id_\ast\comp e=e$, $e\comp\id_\ast=e$ i $e\comp e=e$ (les tres que involucren $\id_\ast$ són el que en feia la identitat)--- considerem l'assignació $\alpha$ que fixa el vèrtex i envia tots dos bucles a $e$:
\[
\begin{tikzpicture}[baseline=(m.center)]
  \matrix (m) [cat matrix] { \ast \\ };
  \draw[cat] (m-1-1) to[out=55,in=125,looseness=9] node[above] {\(\id_\ast\)} (m-1-1);
  \draw[cat] (m-1-1) to[out=-125,in=-55,looseness=9] node[below] {\(e\)} (m-1-1);
\end{tikzpicture}
\qquad\stackrel{\alpha}{\longmapsto}\qquad
\begin{tikzpicture}[baseline=(m.center)]
  \matrix (m) [cat matrix] { \ast \\ };
  \draw[cat] (m-1-1) to[out=55,in=125,looseness=9] node[above] {\(e\)} (m-1-1);
\end{tikzpicture}
\]
Com a morfisme de grafs, $\alpha$ és vàlid: respecta origen i destí, perquè tots els bucles van del vèrtex a si mateix. Però no prové de cap functor: si fos $\alpha=U(F)$, el functor $F$ enviaria $\id_\ast$ a $e$, i cap functor no pot fer-ho ---ha d'enviar $\id_\ast$ a una identitat, i $e$ no ho és---. Els morfismes de grafs són, doncs, més abundants que els functors. Aquest mateix contraexemple, des de l'angle de la definició de functor, es reprèn a l'exercici~\ref{ex:identitat-necessaria}.

\subsection{La categoria lliure sobre un graf}

En sentit contrari, la construcció de la categoria lliure generada per un graf, $\Free(G)$, que vam definir a la secció~\ref{sec:grafs-lliures}, també és functorial. Recordem-ne l'acció sobre objectes: $\Free(G)$ té per objectes els vèrtexs de $G$ i per morfismes els camins d'arestes, amb la concatenació com a composició i el camí buit com a identitat de cada vèrtex.

Vegem l'acció sobre morfismes. Sigui $\alpha\colon G\to H$ un morfisme de grafs, amb components $\alpha_V\colon V_G\to V_H$ i $\alpha_E\colon E_G\to E_H$. Definim un functor
\[
\Free(\alpha)\colon\Free(G)\to\Free(H)
\]
que, sobre objectes, actua com $\alpha_V$, i sobre morfismes transforma cada camí aresta a aresta: el camí $(f_1,\dots,f_k)$ de $G$ va al camí $(\alpha_E(f_1),\dots,\alpha_E(f_k))$ de $H$, i el camí buit en $v$, $()_v$, va al camí buit en $\alpha_V(v)$. L'assignació està ben definida ---com que $\alpha$ preserva origen i destí, arestes encadenades van a arestes encadenades, i el resultat és de nou un camí--- i és un functor: transforma concatenacions en concatenacions i camins buits en camins buits, és a dir, composicions en composicions i identitats en identitats.

Finalment, $\Free$ respecta la identitat i la composició de $\cat{Graph}$: $\Free(\id_G)=\id_{\Free(G)}$ i $\Free(\beta\comp\alpha)=\Free(\beta)\comp\Free(\alpha)$, ja que aplicar $\alpha$ i després $\beta$ aresta a aresta és el mateix que aplicar $\beta\comp\alpha$. Tenim, doncs, un functor
\[
\Free\colon\cat{Graph}\to\cat{Cat}.
\]

Vegem-ho amb el mateix graf d'un bucle de l'exemple anterior. Sigui $G$ el graf amb un sol vèrtex $\ast$ i una sola aresta $a$, un bucle. Els camins de $G$ són les paraules en la lletra $a$ ---$\varepsilon, a, aa, aaa,\dots$---, de manera que $\Free(G)$ és la categoria d'un sol objecte els morfismes de la qual són aquestes paraules, amb la concatenació per composició i la paraula buida $\varepsilon$ per identitat:
\[
\begin{tikzpicture}[baseline=(m.center)]
  \matrix (m) [cat matrix] { \ast \\ };
  \draw[cat] (m-1-1) to[out=55,in=125,looseness=9] node[above] {\(a\)} (m-1-1);
\end{tikzpicture}
\qquad\stackrel{\Free}{\longmapsto}\qquad
\begin{tikzpicture}[baseline=(m.center)]
  \matrix (m) [cat matrix] { \ast \\ };
  \draw[cat] (m-1-1) to[out=120,in=175,looseness=12] node[left] {\(\varepsilon\)} (m-1-1);
  \draw[cat] (m-1-1) to[out=65,in=115,looseness=12] node[above] {\(a\)} (m-1-1);
  \draw[cat] (m-1-1) to[out=5,in=60,looseness=12] node[right] {\(aa\)} (m-1-1);
  \draw[cat] (m-1-1) to[out=-55,in=-5,looseness=12] node[right] {\(\dots\)} (m-1-1);
  \draw[cat] (m-1-1) to[out=-125,in=-70,looseness=12] node[below] {\(a^n\)} (m-1-1);
\end{tikzpicture}
\]
Comptant la paraula buida com a $a^0$, els morfismes són exactament les potències $a^n$ amb $n\geq 0$, i concatenar-les suma els exponents: $a^m\comp a^n=a^{m+n}$. Aquest monoide de morfismes és $(\mathbb{N},+,0)$, de manera que $\Free(G)=\cat{B}\mathbb{N}$. Mentre que $U$ oblidava estructura, $\Free$ en crea: d'una sola aresta en surten infinits morfismes, un per cada longitud de camí.

Cal notar que $\Free$ i $U$ són functors \emph{entre categories grans} ($\cat{Graph}$ i $\cat{Cat}$ no són objectes de la $\cat{Cat}$ de les categories petites); les tractem com a categories localment petites, no com a fletxes de $\cat{Cat}$.

\subsection{La relació entre \texorpdfstring{$\Free$}{Free} i \texorpdfstring{$U$}{U}}\label{subsec:free-u}

Les dues construccions van en sentits contraris, però no són inverses l'una de l'altra. El functor $U$ oblida quins morfismes són identitats i com es componen; $\Free$ afegeix formalment identitats i composicions al graf, sense imposar cap altra relació. Quan partim del graf subjacent d'una categoria $\cat{C}$, la composició lliure crea camins formals nous, de manera que $\Free(U(\cat{C}))$ no és $\cat{C}$ en general; però hi ha un functor
\[
\Free(U(\cat{C}))\to\cat{C}
\]
que avalua cada camí formal mitjançant la composició de $\cat{C}$. Per exemple, prenem per $\cat{C}$ el monoide $M=(\{1,e\},\cdot,1)$ amb $e\cdot e=e$, vist com a categoria d'un sol objecte $\ast$. El seu graf subjacent té un vèrtex i dos bucles $\id_\ast$ i $e$, de manera que $\Free(U(\cat{C}))$ té per morfismes \emph{totes} les paraules en aquestes dues lletres ---infinites---, amb la paraula buida $\varepsilon$ per identitat. El functor $\Free(U(\cat{C}))\to\cat{C}$ avalua cada paraula multiplicant-la a $M$: $\varepsilon\mapsto\id_\ast$, $ee\mapsto e\comp e=e$, $\id_\ast e\mapsto e$, i així successivament. Col·lapsa, doncs, la infinitat de camins formals sobre els dos únics morfismes de $\cat{C}$, recuperant les composicions que $\Free$ havia descartat. Aquest functor il·lustra una propietat que caracteritza $\Free(G)$. Sigui $\cat{D}$ una categoria petita. Definir un functor $\Free(G)\to\cat{D}$ equival a triar les imatges dels vèrtexs \emph{i} de les arestes de $G$, respectant origen i destí ---és a dir, a donar un morfisme de grafs $G\to U(\cat{D})$---, i aquesta tria s'estén de manera única a $\Free(G)$. Cal comptar totes dues assignacions: un functor ha d'enviar cada objecte de $\Free(G)$ ---cada vèrtex de $G$--- a un objecte de $\cat{D}$, també els vèrtexs aïllats, que no són extrem de cap aresta. Per exemple, si $G$ té un únic vèrtex i cap aresta, no hi ha cap imatge d'aresta per triar, però un functor $\Free(G)\to\cat{D}$ encara exigeix escollir un objecte de $\cat{D}$; és precisament la component sobre vèrtexs del morfisme de grafs $G\to U(\cat{D})$, que en aquest cas es redueix a triar aquest objecte. La formulació via $G\to U(\cat{D})$ recull sempre les dues dades, perquè un morfisme de grafs consta d'una aplicació sobre vèrtexs i una sobre arestes. Aquesta correspondència és una bijecció
\[
\Hom_{\cat{Cat}}\big(\Free(G),\cat{D}\big)\;\cong\;\Hom_{\cat{Graph}}\big(G,U(\cat{D})\big).
\]
Aquesta és la formulació precisa de \enquote{lliure}. La retrobarem com a exemple al capítol~\refcap{cap:yoneda}{4}, i en tota la seva generalitat com l'\emph{adjunció} $\Free\dashv U$ al capítol~\refcap{cap:adjuncions}{6}.

\begin{observacio}
Reconèixer els functors com a fletxes d'una categoria és el primer pas per poder plantejar-nos si són isomorfismes, si es componen o si hi ha dualitat. Dins $\cat{Cat}$, per exemple, una aplicació monòtona $P\to Q$ entre preordres petits és una fletxa; en canvi, $\Free$ i $U$ són functors entre categories grans, i per això no són morfismes de $\cat{Cat}$; això no els fa menys legítims, sinó que simplement no els agrupem com a fletxes de cap categoria de categories.
\end{observacio}

\section{Functors covariants i contravariants}\label{sec:variancia}

Fins ara els functors preserven el sentit de les fletxes. Sovint, però, apareixen assignacions que \emph{inverteixen} el sentit: enviar $f\colon A\to B$ a un morfisme $F(B)\to F(A)$. Parlem aleshores de \textbf{functors contravariants}, que s'expliquen millor mitjançant la categoria oposada.

\begin{definicio}\index{functor!contravariant}\index{functor!covariant}
Un \textbf{functor contravariant} de $\cat{C}$ a $\cat{D}$ és un functor
\[
F\colon\cat{C}\op\to\cat{D}.
\]
Per contrast, els functors de la definició original s'anomenen \textbf{covariants}.
\end{definicio}

Desplegant la definició de functor sobre $\cat{C}\op$, un functor contravariant assigna a cada objecte $A$ un objecte $F(A)$ i a cada morfisme $f\colon A\to B$ de $\cat{C}$ un morfisme en sentit invertit
\[
F(f)\colon F(B)\to F(A),
\]
de manera que $F(\id_A)=\id_{F(A)}$ i $F(g\comp f)=F(f)\comp F(g)$. La inversió de l'ordre en aquesta última igualtat no és cap condició nova: és la condició de functor covariant sobre $\cat{C}\op$, un cop recordem que compondre a $\cat{C}\op$ és compondre a $\cat{C}$ en l'ordre invers. Aquest és un primer exemple de la potència de la categoria oposada: permet tractar covariància i contravariància amb una sola definició.

Diagramàticament, un triangle commutatiu de $\cat{C}$ es transforma en un triangle commutatiu de $\cat{D}$ amb les fletxes imatge invertides:
\[
\begin{tikzpicture}[baseline=(m.center)]
  \matrix (m) [matrix of math nodes, column sep=3.2em, row sep=2.6em] {
    A & B \\
      & C \\
  };
  \draw[-{Stealth[scale=1.1]}] (m-1-1) -- node[above] {$f$} (m-1-2);
  \draw[-{Stealth[scale=1.1]}] (m-1-2) -- node[right] {$g$} (m-2-2);
  \draw[-{Stealth[scale=1.1]}] (m-1-1) -- node[below left] {$g\comp f$} (m-2-2);
\end{tikzpicture}
\qquad\stackrel{F}{\longmapsto}\qquad
\begin{tikzpicture}[baseline=(m.center)]
  \matrix (m) [matrix of math nodes, column sep=3.6em, row sep=2.6em] {
    F(A) & F(B) \\
         & F(C) \\
  };
  \draw[{Stealth[scale=1.1]}-] (m-1-1) -- node[above] {$F(f)$} (m-1-2);
  \draw[{Stealth[scale=1.1]}-] (m-1-2) -- node[right] {$F(g)$} (m-2-2);
  \draw[{Stealth[scale=1.1]}-] (m-1-1) -- node[below left] {$F(g\comp f)$} (m-2-2);
\end{tikzpicture}
\]

\begin{definicio}[Functor oposat]\label{def:functor-oposat}\index{functor!oposat}
Tot functor $F\colon\cat{C}\to\cat{D}$ indueix un \textbf{functor oposat}
\[
F\op\colon\cat{C}\op\to\cat{D}\op,
\]
que actua igual que $F$ sobre objectes i sobre morfismes, però llegit entre les categories oposades. Com que compondre a $\cat{C}\op$ i a $\cat{D}\op$ és compondre a $\cat{C}$ i a $\cat{D}$ en l'ordre invers, les dues condicions de functor es transporten sense canvis.
\end{definicio}

\begin{observacio}
El pas a l'oposat és, de fet, functorial: assigna a cada categoria petita $\cat{C}$ la seva oposada $\cat{C}\op$ i a cada functor $F$ el functor oposat $F\op$, i respecta composicions i identitats, de manera que
\[
(-)\op\colon\cat{Cat}\to\cat{Cat}
\]
és un functor ---covariant: envia $F\colon\cat{C}\to\cat{D}$ a $F\op\colon\cat{C}\op\to\cat{D}\op$, en el mateix sentit. A més, aplicar-lo dues vegades torna al punt de partida, tant sobre objectes com sobre morfismes: $(\cat{C}\op)\op=\cat{C}$ i $(F\op)\op=F$. Dit altrament, $(-)\op\comp(-)\op=\id_{\cat{Cat}}$: el functor $(-)\op$ és la seva pròpia inversa, és a dir, una \textbf{involució}. En particular és un \textbf{isomorfisme de categories} de $\cat{Cat}$ en ella mateixa ---un endofunctor amb functor invers, que a més coincideix amb ell mateix. És la manifestació, al nivell de $\cat{Cat}$, de l'autodualitat que ja hem trobat objecte a objecte amb el principi de dualitat. Convé recordar, això sí, que tot plegat viu sobre categories \emph{petites}, les úniques que són objectes de $\cat{Cat}$; per a una categoria gran, l'oposada segueix tenint sentit, però no com a imatge d'una fletxa de $\cat{Cat}$.
\end{observacio}

\begin{proposicio}[Variància d'una composició]\label{prop:variancia-composicio}\index{functor!variància d'una composició}\index{regla dels signes (variància)}
Siguin $F$ i $G$ dos functors, cadascun covariant o contravariant, amb el codomini de $F$ i el domini de $G$ iguals com a categoria (a banda de la variància). L'assignació sobre objectes $A\mapsto G(F(A))$ és de nou un functor, i la seva variància segueix la \textbf{regla dels signes}: la composició és covariant si $F$ i $G$ tenen la mateixa variància, i contravariant si en tenen de diferents. En particular, la composició de dos functors contravariants és covariant.
\end{proposicio}

\begin{prova}
Un functor contravariant $\cat{C}\to\cat{D}$ és, per definició, un functor covariant $\cat{C}\op\to\cat{D}$, i la composició de functors torna a ser un functor (ho hem vist en construir $\cat{Cat}$). Només cal fer encaixar dominis i codominis, i d'això se n'ocupa el functor oposat.

Si $F$ i $G$ són tots dos covariants, $G\comp F\colon\cat{C}\to\cat{E}$ és directament la composició, covariant. Si tots dos són contravariants, escrits com $F\colon\cat{C}\op\to\cat{D}$ i $G\colon\cat{D}\op\to\cat{E}$, prenem el functor oposat $F\op\colon\cat{C}\to\cat{D}\op$ (definició~\ref{def:functor-oposat}), que actua igual que $F$ sobre objectes; aleshores
\[
G\comp F\op\colon\cat{C}\to\cat{E}
\]
és una composició de functors ---per tant un functor--- covariant, i sobre objectes fa $A\mapsto G(F(A))$. Els dos casos mixtos són idèntics en mètode: un ajust anàleg deixa el domini en $\cat{C}\op$ i la composició resulta contravariant. La variància només depèn, doncs, de la paritat del nombre de factors contravariants.
\end{prova}

\subsection{Els homfunctors}\label{subsec:homfunctors}

L'exemple més important de functors covariants i contravariants ---i el que sostindrà bona part del que ve després, fins al lema de Yoneda--- surt de mirar els morfismes d'una categoria cap a un objecte fix, o des d'un objecte fix.

\begin{definicio}[Els homfunctors]\label{def:homfunctors}\label{ex:homfunctors}\index{functor!Hom}\index{homfunctor}
Sigui $\cat{C}$ una categoria localment petita. Fixat un objecte $A$ de $\cat{C}$, definim \textbf{homfunctor covariant}\index{homfunctor!covariant}\index{functor!Hom!covariant}
\[
\Hom(A,-)\colon\cat{C}\to\cat{Set}
\]
com el functor que envia cada objecte $X$ al conjunt $\Hom(A,X)$ i cada morfisme $f\colon X\to Y$ a l'aplicació
\[
f_*\colon\Hom(A,X)\to\Hom(A,Y),
\qquad
\varphi\mapsto f\comp\varphi,
\]
la \textbf{postcomposició} amb $f$. Fixat un objecte $B$, definim \textbf{homfunctor contravariant}\index{homfunctor!contravariant}\index{functor!Hom!contravariant}
\[
\Hom(-,B)\colon\cat{C}\op\to\cat{Set}
\]
com el functor que envia cada objecte $X$ al conjunt $\Hom(X,B)$ i cada morfisme $f\colon X\to Y$ a l'aplicació
\[
f^*\colon\Hom(Y,B)\to\Hom(X,B),
\qquad
\psi\mapsto\psi\comp f,
\]
la \textbf{precomposició} amb $f$.
\end{definicio}

\begin{notacio}\label{not:hom-morfismes}\index{Hom@$\Hom$!aplicat a un morfisme}
Com per a qualsevol functor $F$, que envia un morfisme $f$ a $F(f)$, escriurem també
\[
\Hom(A,f)\coloneqq f_*
\qquad\text{i}\qquad
\Hom(f,B)\coloneqq f^*
\]
per a les imatges de $f\colon X\to Y$ pels homfunctors $\Hom(A,-)$ i $\Hom(-,B)$, respectivament.
\end{notacio}

En el cas covariant, un element $\varphi\colon A\to X$ es perllonga amb $f$ per obtenir $f\comp\varphi\colon A\to Y$; en el contravariant, un element $\psi\colon Y\to B$ es precedeix de $f$ per obtenir $\psi\comp f\colon X\to B$:
\[
\begin{tikzpicture}[baseline=(m.center)]
  \matrix (m) [matrix of math nodes, column sep=3.2em, row sep=2.6em] {
    A & X \\
      & Y \\
  };
  \draw[-{Stealth[scale=1.1]}] (m-1-1) -- node[above] {$\varphi$} (m-1-2);
  \draw[-{Stealth[scale=1.1]}] (m-1-2) -- node[right] {$f$} (m-2-2);
  \draw[-{Stealth[scale=1.1]}] (m-1-1) -- node[below left] {$f\comp\varphi$} (m-2-2);
\end{tikzpicture}
\qquad\qquad
\begin{tikzpicture}[baseline=(m.center)]
  \matrix (m) [matrix of math nodes, column sep=3.2em, row sep=2.6em] {
    X & Y \\
      & B \\
  };
  \draw[-{Stealth[scale=1.1]}] (m-1-1) -- node[above] {$f$} (m-1-2);
  \draw[-{Stealth[scale=1.1]}] (m-1-2) -- node[right] {$\psi$} (m-2-2);
  \draw[-{Stealth[scale=1.1]}] (m-1-1) -- node[below left] {$\psi\comp f$} (m-2-2);
\end{tikzpicture}
\]
Observem que, en el cas covariant, $f$ actua sobre el \emph{destí} de les fletxes (les allarga per l'extrem final), mentre que en el contravariant actua sobre l'\emph{origen} (les allarga per l'extrem inicial), i això és el que inverteix el sentit del functor. La functorialitat és conseqüència directa de l'associativitat de la composició i de les lleis de les identitats. El que importa és l'ordre que en resulta:
\[
(g\comp f)_*=g_*\comp f_*,
\qquad\qquad
(g\comp f)^*=f^*\comp g^*.
\]
En el cas covariant l'ordre es conserva; en el contravariant s'inverteix, perquè $f$ actua per l'extrem inicial de les fletxes. Les comprovacions detallades són els exercicis resolts~\ref{pr-hom-cov} i~\ref{pr-hom-contra} de la Pràctica. Aquests dos functors són centrals en teoria de categories i els retrobarem en el lema de Yoneda (capítol~\refcap{cap:yoneda}{4}).

Els homfunctors traslladen els isomorfismes de $\cat{C}$ a bijeccions de conjunts. Això és un cas particular de la proposició~\ref{prop:functor-preserva-iso}, però val la pena escriure'l explícitament, perquè és la forma en què l'usarem sovint.

\begin{corollari}[Els homfunctors transformen isos en bijeccions]\label{prop:homfunctor-iso}\index{homfunctor!i isomorfismes}
Sigui $\cat{C}$ localment petita i sigui $f\colon X\to Y$ un isomorfisme de $\cat{C}$. Aleshores, per a tot objecte $A$, la postcomposició
\[
f_*\colon\Hom(A,X)\to\Hom(A,Y),\qquad\varphi\mapsto f\comp\varphi,
\]
és una bijecció, amb inversa $(f^{-1})_*$; i, per a tot objecte $B$, la precomposició
\[
f^*\colon\Hom(Y,B)\to\Hom(X,B),\qquad\psi\mapsto\psi\comp f,
\]
és una bijecció, amb inversa $(f^{-1})^*$.
\end{corollari}

\begin{prova}
Un isomorfisme de $\cat{Set}$ no és res més que una bijecció. Apliquem, doncs, la proposició~\ref{prop:functor-preserva-iso} a cada homfunctor. Per al covariant $\Hom(A,-)$, la imatge de l'isomorfisme $f$ és $f_*$, que resulta ser un isomorfisme de $\cat{Set}$ amb invers la imatge de $f^{-1}$, és a dir $(f^{-1})_*$. Per al contravariant $\Hom(-,B)$ ---un functor sobre $\cat{C}\op$, on $f$ segueix sent un isomorfisme--- la imatge de $f$ és $f^*$, que és igualment un isomorfisme de $\cat{Set}$ amb invers $(f^{-1})^*$. En tots dos casos, ser isomorfisme a $\cat{Set}$ vol dir ser bijecció.
\end{prova}

Dit d'una altra manera: si $X\cong Y$, aleshores $\Hom(A,X)\cong\Hom(A,Y)$ i $\Hom(Y,B)\cong\Hom(X,B)$ com a conjunts, per a qualssevol $A$ i $B$. Els objectes isomorfs són, doncs, indistingibles vistos a través de qualsevol homfunctor: una primera pista de la idea que un objecte queda determinat pels morfismes que hi arriben o que en surten, que el lema de Yoneda portarà fins al final.

Els functors contravariants cap a $\cat{Set}$, com l'homfunctor $\Hom(-,B)$, apareixen tan sovint que reben un nom propi.

\begin{definicio}[Prefeix]\label{def:prefeix}\index{prefeix@prefeix}
Sigui $\cat{C}$ una categoria. Un \textbf{prefeix} sobre $\cat{C}$ és un functor contravariant de $\cat{C}$ a $\cat{Set}$, és a dir, un functor
\[
P\colon\cat{C}\op\to\cat{Set}.
\]
Explícitament, $P$ assigna a cada objecte $A$ de $\cat{C}$ un conjunt $P(A)$ i a cada morfisme $f\colon A\to B$ una aplicació $P(f)\colon P(B)\to P(A)$, de manera que $P(\id_A)=\id_{P(A)}$ i $P(g\comp f)=P(f)\comp P(g)$.
\end{definicio}

Si $\cat{C}$ és localment petita, cada homfunctor contravariant $\Hom(-,B)$ és un prefeix sobre $\cat{C}$, i el functor de parts de l'exemple següent n'és un sobre $\cat{Set}$. Els prefeixos tornaran a aparèixer als capítols~\refcap{cap:yoneda}{4} i~\refcap{cap:topos}{9}.

Així, \emph{prefeix sobre $\cat{C}$} i \emph{functor contravariant de $\cat{C}$ a $\cat{Set}$} són sinònims: el nom no afegeix cap condició, ni tan sols de mida. A partir d'ara, per a un functor contravariant cap a $\cat{Set}$ direm sempre \emph{prefeix},\footnote{A la literatura també es parla de \emph{prefeixos sobre $\cat{C}$ amb valors a $\cat{D}$} (per exemple, prefeixos de grups o d'espais vectorials) per als functors $\cat{C}\op\to\cat{D}$. En aquest text, \emph{prefeix} vol dir sempre amb valors a $\cat{Set}$.} i n'escriurem la signatura $P\colon\cat{C}\op\to\cat{Set}$ quan convingui. Reservarem l'expressió \emph{functor contravariant} per als que arriben a altres categories, com el functor dual de l'exemple~\ref{ex:homfunctor-dual}. Els functors \emph{covariants} $\cat{C}\to\cat{Set}$ no reben nom propi en aquest text: són, simplement, els prefeixos sobre $\cat{C}\op$.

\begin{exemple}[El functor de parts i l'homfunctor $\Hom(-,\{0,1\})$]\label{ex:functor-parts}\index{functor!de parts}\index{homfunctor!i functor de parts}
Comencem dins de $\cat{Set}$, prenent com a objecte fixat el conjunt de dos elements $\{0,1\}$. Cada subconjunt $A\subseteq X$ es correspon bijectivament amb la seva \textbf{funció característica} $\chi_A\colon X\to\{0,1\}$, que val $1$ sobre $A$ i $0$ fora; recíprocament, cada aplicació $X\to\{0,1\}$ és la característica del subconjunt on val $1$. Les funcions característiques donen, doncs, per a cada conjunt $X$, una bijecció canònica
\[
\mathcal P(X)\;\cong\;\Hom_{\cat{Set}}(X,\{0,1\}).
\]
Els dos conjunts no són iguals ---un conté subconjunts, l'altre aplicacions---, però la bijecció no depèn de cap tria.

Vegem que, sota aquestes bijeccions, la precomposició correspon a l'antiimatge. Per a $f\colon X\to Y$, l'homfunctor envia una aplicació $\chi\colon Y\to\{0,1\}$ a $\chi\comp f\colon X\to\{0,1\}$; i si $\chi=\chi_B$ és la característica de $B\subseteq Y$, aleshores $\chi_B\comp f$ val $1$ exactament sobre els $x$ amb $f(x)\in B$, és a dir, és la característica de l'antiimatge $f^{-1}(B)$. Així, la precomposició
\[
\chi\longmapsto\chi\comp f
\]
és precisament l'aplicació d'antiimatge
\[
f^{-1}\colon\mathcal P(Y)\to\mathcal P(X),
\qquad
B\mapsto\{x\in X : f(x)\in B\}.
\]
Aquí $f^{-1}$ denota l'aplicació d'antiimatge, no l'invers de $f$, que pot no existir. Les antiimatges defineixen el \textbf{functor de parts} contravariant
\[
\mathcal P\colon\cat{Set}\op\to\cat{Set},\qquad X\mapsto\mathcal P(X),\qquad f\mapsto f^{-1},
\]
i la seva contravariància, $(g\comp f)^{-1}=f^{-1}\comp g^{-1}$, es pot transportar, a través de les bijeccions anteriors, des de la de l'homfunctor $\Hom(-,\{0,1\})$; la comprovació directa és l'exercici resolt~\ref{pr-parts-contra} de la Pràctica.

Per tant, el functor de parts contravariant i l'homfunctor $\Hom_{\cat{Set}}(-,\{0,1\})$ descriuen la mateixa construcció fins a una identificació canònica: provar cada conjunt amb un objecte distingit, el conjunt $\{0,1\}$ de valors de veritat. No són literalment el mateix functor. El llenguatge de les transformacions naturals (capítol~\ref{cap:transformacions}) permetrà expressar aquesta afirmació amb precisió: tots dos functors seran \emph{naturalment isomorfs} (exemple~\ref{ex:isos-naturals-cap2}).

Convé un avís sobre l'altra variància, i val la pena justificar-lo. Un homfunctor covariant $\Hom_{\cat{Set}}(\{0,1\},-)$ envia cada conjunt $X$ al conjunt d'aplicacions $\{0,1\}\to X$. Una tal aplicació queda completament determinada per la parella d'imatges $(x_0,x_1)$ amb $x_0$ la imatge de $0$ i $x_1$ la de $1$, i tota parella d'elements de $X$ prové d'una aplicació; per tant
\[
\Hom_{\cat{Set}}(\{0,1\},X)\;\cong\;X\times X,
\]
i sobre morfismes la postcomposició amb $f\colon X\to Y$ actua component a component, $(x_0,x_1)\mapsto(f(x_0),f(x_1))$. Aquest és, doncs, el functor \enquote{parell ordenat}, ben diferent del functor de parts: $\Hom(\{0,1\},X)\cong X\times X$ té sempre $\lvert X\rvert^2$ elements (quan $X$ és finit), mentre que $\mathcal P(X)$ en té $2^{\lvert X\rvert}$. La correspondència canònica entre el functor de parts i un homfunctor és, per tant, un fenomen exclusiu de la cara contravariant: $\mathcal P$ correspon a $\Hom(-,\{0,1\})$, no a $\Hom(\{0,1\},-)$. La raó de fons és que a $\mathcal P(X)$ el conjunt $\{0,1\}$ fa de \emph{destí} (cada subconjunt és una aplicació \emph{cap a} $\{0,1\}$), i fixar el destí dona un homfunctor contravariant; posar-lo com a \emph{origen} dona una construcció covariant, però tota una altra. La versió covariant genuïna del functor de parts existeix ---envia $f$ a la imatge directa $A\mapsto\{f(x):x\in A\}$--- i és functorial, però no correspon a cap homfunctor; quan calgui distingir-la de $\mathcal P$ l'escriurem $\mathcal P_!$.
\end{exemple}

\begin{exemple}[Les dues variàncies del functor de parts]\label{ex:parts-dues-variancies}\index{functor!de parts!covariant i contravariant}
Posem en paral·lel les dues maneres de fer functorial l'assignació $X\mapsto\mathcal P(X)$. Totes dues coincideixen sobre objectes i difereixen en tot el que fan amb una aplicació $f\colon X\to Y$:
\[
\renewcommand{\arraystretch}{1.3}
\begin{array}{l|l|l}
 & \mathcal P_!\colon\cat{Set}\to\cat{Set}\ \text{(covariant)} & \mathcal P\colon\cat{Set}\op\to\cat{Set}\ \text{(contravariant)}\\ \hline
\text{acció sobre } f & \text{imatge directa } A\mapsto f(A) & \text{antiimatge } B\mapsto f^{-1}(B)\\
\text{tipus} & \mathcal P_!(f)\colon\mathcal P(X)\to\mathcal P(Y) & \mathcal P(f)\colon\mathcal P(Y)\to\mathcal P(X)\\
\text{identitat} & \id_X(A)=A & \id_X^{-1}(B)=B\\
\text{composició} & (g\comp f)(A)=g\bigl(f(A)\bigr) & (g\comp f)^{-1}(C)=f^{-1}\bigl(g^{-1}(C)\bigr)\\
\text{llei functorial} & \mathcal P_!(g\comp f)=\mathcal P_!(g)\comp\mathcal P_!(f) & \mathcal P(g\comp f)=\mathcal P(f)\comp\mathcal P(g)
\end{array}
\]
Llegida per files, la taula segueix l'ordre de l'observació~\ref{obs:comprovar-tipus}: primer el tipus, que ja fixa la variància, i després les dues equacions. Les comprovacions completes són els exercicis resolts~\ref{pr-parts-cov} i~\ref{pr-parts-contra} de la Pràctica. Observem que a la columna contravariant l'ordre de $f$ i $g$ s'inverteix: és el mateix fenomen que a $(g\comp f)^*=f^*\comp g^*$ per a l'homfunctor contravariant.

Les dues columnes no són independents. Per a $A\subseteq X$ i $B\subseteq Y$,
\[
f(A)\subseteq B\quad\Longleftrightarrow\quad A\subseteq f^{-1}(B),
\]
perquè totes dues condicions diuen que $f(a)\in B$ per a tot $a\in A$. Aquesta equivalència entre les dues accions, una en cada sentit, és el primer exemple d'una \emph{adjunció}; la retrobarem al capítol~\refcap{cap:adjuncions}{6} com a $\exists_f\dashv f^{*}$, l'origen categòric del quantificador existencial (exemple~\refalt{ex:quantificadors-set}{d'imatge, antiimatge i quantificadors del capítol~6}).
\end{exemple}

En aquest exemple l'objecte fixat feia de \emph{destí}; fixant-lo com a \emph{origen} obtenim homfunctors covariants amb una lectura igual de concreta. Ho veiem a $\cat{Graph}$, on l'estructura dels objectes ja no és la d'un simple conjunt.

\begin{exemple}[Vèrtexs i arestes com a homfunctors]\label{ex:homfunctor-grafs}\index{homfunctor!a Graph}\index{graf dirigit!vèrtexs i arestes com a homfunctors}
Treballem a $\cat{Graph}$, amb la notació de la subsecció~\ref{subsec:grafs}: un graf $G$ té vèrtexs $V_G$, arestes $E_G$ i aplicacions origen i destí $s_G,t_G\colon E_G\to V_G$. Triem dos grafs \enquote{patró} ben senzills:
\[
\mathbf V=\bigl(\{\ast\},\ \varnothing\bigr),
\qquad
\mathbf E=\bigl(\{0,1\},\ \{e\}\bigr)\ \text{amb } s(e)=0,\ t(e)=1,
\]
és a dir, $\mathbf V$ és un sol vèrtex sense arestes i $\mathbf E$ és una sola aresta entre dos vèrtexs.

Un morfisme de grafs $\mathbf V\to G$ no té cap aresta on decidir res, i sobre l'únic vèrtex tria un vèrtex de $G$; per tant queda determinat per aquesta tria, i n'hi ha un per cada vèrtex:
\[
\Hom_{\cat{Graph}}(\mathbf V,G)\;\cong\;V_G.
\]
Un morfisme de grafs $\mathbf E\to G$ tria una aresta $\varepsilon$ de $G$ ---la imatge de $e$--- i, forçosament, els seus dos extrems han d'anar a $s_G(\varepsilon)$ i $t_G(\varepsilon)$, de manera que la tria de l'aresta ja determina tot el morfisme; n'hi ha un per cada aresta:
\[
\Hom_{\cat{Graph}}(\mathbf E,G)\;\cong\;E_G.
\]
Així, els homfunctors covariants $\Hom_{\cat{Graph}}(\mathbf V,-)$ i $\Hom_{\cat{Graph}}(\mathbf E,-)$ corresponen canònicament, respectivament, al functor \enquote{conjunt de vèrtexs} i al functor \enquote{conjunt d'arestes} $\cat{Graph}\to\cat{Set}$. Sobre un morfisme $F=(F_V,F_E)\colon G\to H$, el functor $\Hom_{\cat{Graph}}(\mathbf V,-)$ actua per la postcomposició $F_*\colon\Hom_{\cat{Graph}}(\mathbf V,G)\to\Hom_{\cat{Graph}}(\mathbf V,H)$, $v\mapsto F\comp v$ (definició~\ref{def:homfunctors}). Si $v$ tria el vèrtex $v(\ast)$, aleshores $F\comp v$ tria el vèrtex $F_V(v(\ast))$; per tant, llegida a través de la bijecció $\Hom_{\cat{Graph}}(\mathbf V,G)\cong V_G$, la postcomposició $F_*$ és exactament $F_V\colon V_G\to V_H$. Anàlogament, per a $\Hom_{\cat{Graph}}(\mathbf E,-)$, la postcomposició $F_*$ envia el morfisme que tria l'aresta $\varepsilon$ al que tria $F_E(\varepsilon)$, i correspon a $F_E\colon E_G\to E_H$.

El que hem trobat és que dos functors oblit familiars ---prendre el conjunt de vèrtexs i prendre el conjunt d'arestes--- corresponen canònicament als homfunctors covariants $\Hom_{\cat{Graph}}(\mathbf V,-)$ i $\Hom_{\cat{Graph}}(\mathbf E,-)$, per a una tria adient de l'objecte d'origen; al capítol~\ref{cap:transformacions} reconeixerem aquestes identificacions com a isomorfismes naturals. És el mateix fenomen que fa que, a $\cat{Set}$, el functor identitat correspongui canònicament a $\Hom_{\cat{Set}}(\{\ast\},-)$: provar un objecte amb els morfismes des d'un patró petit en llegeix una part de l'estructura. Que un functor cap a $\cat{Set}$ es pugui reconstruir així, com un $\Hom(A,-)$ per a cert objecte $A$, és una propietat prou important perquè hi tornem amb nom propi ---els \emph{functors representables}--- i és el punt de partida del lema de Yoneda (capítol~\refcap{cap:yoneda}{4}); de moment n'hi ha prou de retenir que aquests functors oblit són homfunctors disfressats, llevat d'una identificació canònica.
\end{exemple}

Als exemples anteriors la categoria d'arribada era $\cat{Set}$. Quan els homconjunts tenen \emph{ells mateixos} estructura, l'homfunctor la recull i aterra en una categoria més rica; el cas paradigmàtic és el dels espais vectorials.

\begin{exemple}[El dual d'un espai vectorial com a homfunctor]\label{ex:homfunctor-dual}\index{homfunctor!i espai dual}\index{espai dual!com a homfunctor}\index{aplicació dual}
Fixem un cos $\mathbb K$ i treballem a $\cat{Vect}_{\mathbb K}$. Aquí cada homconjunt $\Hom_{\mathbb K}(V,W)$ no és un conjunt qualsevol: és ell mateix un $\mathbb K$-espai vectorial, amb la suma d'aplicacions lineals i el producte per escalars. Prenent com a segon objecte el cos $\mathbb K$ vist com a espai vectorial sobre si mateix, l'homconjunt
\[
\Hom_{\mathbb K}(V,\mathbb K)=V^*
\]
és exactament l'\textbf{espai dual} de $V$: les formes lineals sobre $V$. Com a conjunt, $V^*$ és el valor en $V$ de l'homfunctor contravariant de la definició~\ref{def:homfunctors}, especialitzat a $B=\mathbb K$:
\[
\Hom_{\mathbb K}(-,\mathbb K)\colon\cat{Vect}_{\mathbb K}\op\to\cat{Set}.
\]
Però aquests homconjunts tenen canònicament estructura d'espai vectorial, amb la suma i el producte per escalars punt a punt, i la precomposició amb una aplicació lineal és lineal (ho veurem tot seguit). Per tant, la mateixa assignació s'\emph{aixeca} a un functor amb valors a $\cat{Vect}_{\mathbb K}$, el \textbf{functor dual}
\[
(-)^*=\Hom_{\mathbb K}(-,\mathbb K)\colon\cat{Vect}_{\mathbb K}\op\to\cat{Vect}_{\mathbb K},
\qquad V\mapsto V^*.
\]
És el mateix patró que el functor de parts, amb $\mathbb K$ en lloc de $\{0,1\}$ com a objecte distingit. Sobre morfismes, la precomposició $\psi\mapsto\psi\comp f$ és precisament l'\textbf{aplicació dual} (o transposada) d'una aplicació lineal $f\colon V\to W$:
\[
f^*\colon W^*\to V^*,\qquad \psi\mapsto\psi\comp f.
\]
La contravariància, $(g\comp f)^*=f^*\comp g^*$, és la regla familiar de la transposició; si triem bases i representem les aplicacions per matrius, $f^*$ correspon a la matriu transposada i la inversió de l'ordre és la fórmula $(AB)^{\mathsf T}=B^{\mathsf T}A^{\mathsf T}$.

Cada $f^*$ és, en efecte, una aplicació lineal, $f^*(\psi+\lambda\psi')=(\psi+\lambda\psi')\comp f=f^*(\psi)+\lambda f^*(\psi')$, i això és el que permet l'aixecament: l'homfunctor $\cat{Set}$-valuat i el functor dual tenen la mateixa acció sobre objectes i morfismes, però el segon recorda, a més, l'estructura lineal dels valors. El mateix passa amb el covariant: $\Hom_{\mathbb K}(V,-)\colon\cat{Vect}_{\mathbb K}\to\cat{Set}$ s'aixeca a un functor
\[
\Hom_{\mathbb K}(V,-)\colon\cat{Vect}_{\mathbb K}\to\cat{Vect}_{\mathbb K}.
\]
Quan parlem d'homfunctors sense cap més precisió, ens referim sempre a la versió amb valors a $\cat{Set}$ de la definició~\ref{def:homfunctors}; és la que intervindrà en la representabilitat.

El corol·lari~\ref{prop:homfunctor-iso} es llegeix aquí com un fet conegut de l'àlgebra lineal: si $f\colon V\to W$ és un isomorfisme, aleshores la transposada $f^*\colon W^*\to V^*$ també ho és, amb inversa $(f^{-1})^*$. En particular, espais isomorfs tenen duals isomorfs.

No seguim, de moment, cap a la comparació entre $V$ i el seu dual, ni cap al bidual (o doble dual) $V^{**}$: aquesta és una qüestió de \emph{naturalitat}, que demana el llenguatge de les transformacions naturals. La reprendrem exactament on la deixem, a l'exemple~\ref{ex:bidual} del capítol~\ref{cap:transformacions}.
\end{exemple}


\subsection{El bihomfunctor}\label{subsec:bihomfunctor}

Les dues variables dels homfunctors es poden tractar alhora. Fixem-nos que $\Hom$ pren \emph{dos} objectes i en dona un conjunt, de manera contravariant en el primer i covariant en el segon; això és exactament un functor sobre un producte de categories.

\begin{definicio}[El bihomfunctor]\label{def:bihomfunctor}\index{homfunctor!bihomfunctor}\index{Hom@$\Hom$!bihomfunctor}
Sigui $\cat{C}$ localment petita. Definim \textbf{bihomfunctor} (o \textbf{bifunctor Hom})
\[
\Hom_{\cat{C}}(-,-)\colon\cat{C}\op\times\cat{C}\to\cat{Set}
\]
com el functor que envia una parella d'objectes $(A,B)$ al conjunt $\Hom(A,B)$, i una parella de morfismes $(f\colon A'\to A,\ g\colon B\to B')$ a l'aplicació
\[
\Hom(f,g)\colon\Hom(A,B)\to\Hom(A',B'),\qquad h\mapsto g\comp h\comp f.
\]
\end{definicio}

La fórmula $h\mapsto g\comp h\comp f$ combina la precomposició amb $f$ (contravariant, en la primera variable) i la postcomposició amb $g$ (covariant, en la segona): un element $h\colon A\to B$ es prolonga per l'origen amb $f$ i pel destí amb $g$.
\[
\begin{tikzpicture}[baseline=(m.center)]
  \matrix (m) [matrix of math nodes, column sep=3em] {
    A' & A & B & B' \\
  };
  \draw[-{Stealth[scale=1.1]}] (m-1-1) -- node[above] {$f$} (m-1-2);
  \draw[-{Stealth[scale=1.1]}] (m-1-2) -- node[above] {$h$} (m-1-3);
  \draw[-{Stealth[scale=1.1]}] (m-1-3) -- node[above] {$g$} (m-1-4);
  \draw[-{Stealth[scale=1.1]}] (m-1-1) to[bend right=22] node[below] {$g\comp h\comp f$} (m-1-4);
\end{tikzpicture}
\]
Fixant la primera variable en $A$ recuperem l'homfunctor covariant $\Hom(A,-)$; fixant la segona en $B$ recuperem el contravariant $\Hom(-,B)$. També sobre morfismes les notacions són coherents: prenent una de les dues components igual a una identitat,
\[
\Hom(\id_A,g)=g_*=\Hom(A,g),
\qquad
\Hom(f,\id_B)=f^*=\Hom(f,B),
\]
de manera que la notació~\ref{not:hom-morfismes} és un cas particular de la del bihomfunctor. Ara bé, la functorialitat en cada variable per separat no basta: cal que les dues accions siguin compatibles, és a dir, que $\Hom(-,-)$ preservi la composició i les identitats de $\cat{C}\op\times\cat{C}$. També això és conseqüència de l'associativitat i de les identitats: aplicar primer $(f,g)$ i després $(f',g')$ envia $h$ a $g'\comp g\comp h\comp f\comp f'$, que és l'acció de la parella composta $(f\comp f',\,g'\comp g)$, amb la primera component en l'ordre invertit, com correspon a $\cat{C}\op$. La comprovació detallada és l'exercici resolt~\ref{pr-bihom} de la Pràctica.

\begin{exemple}[El bihomfunctor de $\cat{Set}$]\label{ex:bihom-set}\index{homfunctor!bihomfunctor!a Set}
A $\cat{Set}$, el bihomfunctor envia una parella de conjunts $(X,Y)$ al conjunt d'aplicacions
\[
\Hom_{\cat{Set}}(X,Y)=Y^X,
\]
amb la mateixa notació exponencial de l'inici del capítol. L'acció sobre una parella de morfismes $(f\colon X'\to X,\ g\colon Y\to Y')$ és el \enquote{canvi de variables} a l'entrada i a la sortida: donada una aplicació $h\colon X\to Y$, primer la reindexem pel domini amb $f$ i després en traslladem els valors amb $g$,
\[
\Hom(f,g)\colon Y^X\to (Y')^{X'},\qquad h\mapsto g\comp h\comp f.
\]
Fixant el primer conjunt en un punt $\{\ast\}$ recuperem $\Hom(\{\ast\},Y)\cong Y$ ---els elements de $Y$--- amb l'acció covariant habitual per postcomposició; fixant el segon en $\{0,1\}$ recuperem $\Hom(X,\{0,1\})\cong\mathcal P(X)$, és a dir, el functor de parts de l'exemple~\ref{ex:functor-parts}, llevat de la identificació canònica. El bihomfunctor és, doncs, la construcció que engloba tots dos alhora.
\end{exemple}

\begin{exemple}[El bihomfunctor de $\cat{Vect}_{\mathbb K}$]\label{ex:bihom-vect}\index{homfunctor!bihomfunctor!a Vect}
El bihomfunctor ordinari de $\cat{Vect}_{\mathbb K}$ és
\[
\Hom(-,-)\colon\cat{Vect}_{\mathbb K}\op\times\cat{Vect}_{\mathbb K}\to\cat{Set},
\]
que envia una parella d'espais $(V,W)$ al conjunt $\Hom_{\mathbb K}(V,W)$ de les aplicacions lineals de $V$ a $W$. Però aquest conjunt és canònicament un espai vectorial, i sobre una parella $(f\colon V'\to V,\ g\colon W\to W')$ l'acció $h\mapsto g\comp h\comp f$ combina la precomposició amb $f$ i la postcomposició amb $g$, totes dues lineals. Per això el bihomfunctor admet un aixecament canònic a $\cat{Vect}_{\mathbb K}$:
\[
\Hom_{\mathbb K}(-,-)\colon\cat{Vect}_{\mathbb K}\op\times\cat{Vect}_{\mathbb K}\to\cat{Vect}_{\mathbb K}.
\]
Fixant la segona variable en el cos, aquest aixecament dona el functor dual $(-)^*$ de l'exemple~\ref{ex:homfunctor-dual}. Fixant la primera, $\Hom_{\mathbb K}(\mathbb K,-)$ correspon canònicament al functor identitat. En efecte, una aplicació lineal $\mathbb K\to W$ queda determinada per la imatge de $1$, i això dona una bijecció lineal canònica
\[
\Hom_{\mathbb K}(\mathbb K,W)\;\cong\;W;
\]
al capítol~\ref{cap:transformacions} la reconeixerem com un isomorfisme natural. Les dues cares del dual i la identitat es reuneixen, així, en un sol functor de dues variables.
\end{exemple}

\section{Functors fidels i plens}

% CONVENCIÓ D'ESTIL (aplicada a tot el llibre): s'eviten les nominalitzacions
% «fidelitat» i «plenitud» en la prosa expositiva; es prediquen com «ser fidel» /
% «ser ple». Propagat a teoria, exercicis i a la prosa de practica/tests.
% Excepcions conservades a propòsit: les etiquetes de pas en demostracions
% «(fidelitat)» / «per fidelitat» (caps/practica/*, caps/tests/*), i l'ús ordinari
% «per fidelitat a la bibliografia» del glossari (caps/glossari.tex).

Un functor actua sobre els morfismes; és natural preguntar-se com és aquesta acció entre cada parella d'objectes. Fixats $A,B$ a $\cat{C}$, un functor $F\colon\cat{C}\to\cat{D}$ indueix una aplicació
\[
F_{A,B}\colon\Hom_{\cat{C}}(A,B)\to\Hom_{\cat{D}}(F(A),F(B)),
\qquad
f\mapsto F(f).
\]
Les propietats d'aquesta aplicació donen lloc a una classificació bàsica.

\begin{definicio}\index{functor!fidel}\index{functor!ple}
\label{def:ple-fidel}
Un functor $F\colon\cat{C}\to\cat{D}$ és:
\begin{itemize}
\item \textbf{fidel} si cada aplicació $F_{A,B}$ és injectiva; és a dir, si $F(f)=F(g)$ amb $f,g\colon A\to B$ implica $f=g$;
\item \textbf{ple} si cada aplicació $F_{A,B}$ és exhaustiva; és a dir, si tot morfisme $F(A)\to F(B)$ de $\cat{D}$ és de la forma $F(f)$ per a algun $f\colon A\to B$ de $\cat{C}$;
\item \textbf{ple i fidel} si cada $F_{A,B}$ és bijectiva.
\end{itemize}
\end{definicio}

\begin{observacio}
Cal no confondre aquestes condicions amb la injectivitat o l'exhaustivitat \emph{sobre objectes}. Un functor pot ser ple i fidel sense ser injectiu sobre objectes: ser ple i fidel només afecta els morfismes entre cada parella fixada d'objectes. Un functor ple i fidel identifica $\Hom_{\cat{C}}(A,B)$ amb $\Hom_{\cat{D}}(F(A),F(B))$, però pot enviar objectes diferents a un mateix objecte.
\end{observacio}

\begin{observacio}
Aquestes propietats es transporten al functor oposat. Recordem (definició~\ref{def:functor-oposat}) que el functor $F\op\colon\cat{C}\op\to\cat{D}\op$ actua igual que $F$ sobre objectes i morfismes. Cal no confondre el \emph{functor} $F\op$ amb l'\emph{aplicació entre homconjunts} que indueix. Fixats dos objectes $A,B$, que $F$ sigui ple i fidel es mesura amb l'aplicació
\[
F_{A,B}\colon\Hom_{\cat{C}}(A,B)\to\Hom_{\cat{D}}(F(A),F(B)),\qquad f\mapsto F(f),
\]
i les de $F\op$, amb l'aplicació anàloga $(F\op)_{B,A}$ ---amb la parella d'objectes en l'ordre invers, perquè a $\cat{C}\op$ una fletxa $A\to B$ és una fletxa $B\to A$ de $\cat{C}$---
\[
(F\op)_{B,A}\colon\Hom_{\cat{C}\op}(B,A)\to\Hom_{\cat{D}\op}(F(B),F(A)).
\]
Ara bé, per definició d'oposada $\Hom_{\cat{C}\op}(B,A)=\Hom_{\cat{C}}(A,B)$ i $\Hom_{\cat{D}\op}(F(B),F(A))=\Hom_{\cat{D}}(F(A),F(B))$; i sobre aquests conjunts $(F\op)_{B,A}$ envia cada $f$ a $F(f)$, exactament com $F_{A,B}$. Les dues aplicacions són, doncs, literalment la mateixa. Com que \enquote{ser injectiva} i \enquote{ser exhaustiva} són propietats d'aquesta aplicació, $F$ és fidel (respectivament ple) si i només si $F\op$ ho és.
\end{observacio}

\begin{exemple}
El functor oblit $U\colon\cat{Grp}\to\cat{Set}$ és fidel: dos homomorfismes diferents segueixen sent aplicacions diferents, de manera que $U$ no identifica morfismes. En canvi no és ple: entre dos grups hi ha aplicacions de conjunts que no són homomorfismes, i aquestes no provenen de cap morfisme de $\cat{Grp}$. Aquesta és la situació típica dels functors oblit: retenen prou informació per distingir morfismes, però obren la porta a fletxes noves que no respecten l'estructura.
\end{exemple}

Les subcategories, introduïdes elementalment a la subsecció~\ref{subsec:subcategories}, tenen una lectura natural en el llenguatge dels functors. Si $\cat{S}$ és una subcategoria de $\cat{C}$, el \textbf{functor inclusió} $\iota\colon\cat{S}\hookrightarrow\cat{C}$ ---que envia cada objecte i cada morfisme de $\cat{S}$ a ell mateix a $\cat{C}$--- sempre és \textbf{fidel}, perquè no identifica mai dos morfismes diferents. Les dues menes de subcategoria que vam distingir ---la \textbf{plena} i l'\textbf{ampla}--- es tradueixen, cadascuna, en una propietat d'aquest functor. D'una banda, $\cat{S}$ és una \textbf{subcategoria plena} exactament quan $\iota$ és \textbf{ple}: que $\iota$ sigui ple vol dir que entre dos objectes de $\cat{S}$ hi ha \emph{tots} els morfismes que hi ha a $\cat{C}$, que és justament la definició de subcategoria plena. De l'altra, $\cat{S}$ és una \textbf{subcategoria ampla} exactament quan l'acció de $\iota$ \emph{sobre els objectes} és exhaustiva, és a dir, quan $\Ob(\cat{S})=\Ob(\cat{C})$: la subcategoria conté tots els objectes de $\cat{C}$. (Parlem aquí de l'aplicació $\Ob(\cat{S})\to\Ob(\cat{C})$ induïda per $\iota$, no del functor mateix; per al functor inclusió aquesta aplicació és sempre injectiva, i ser ampla equival que sigui, a més, exhaustiva.)\index{subcategoria!plena}\index{subcategoria!ampla}\index{functor!inclusió}

\begin{exemple}\label{ex:inclusio-finset-plena}
El \textbf{functor inclusió} $\cat{FinSet}\hookrightarrow\cat{Set}$ dels conjunts finits és ple i fidel: entre dos conjunts finits hi ha exactament les mateixes aplicacions que a $\cat{Set}$, de manera que un cop triats els objectes no descartem cap morfisme. En canvi, el functor inclusió de la subcategoria de $\cat{Set}$ que conserva tots els conjunts com a objectes però només les aplicacions \emph{injectives} com a morfismes és fidel, però \emph{no} ple: entre dos conjunts hi ha aplicacions no injectives, que per tant no provenen de cap morfisme de la subcategoria. El mateix passa amb la subcategoria ampla de les aplicacions \emph{exhaustives}: el seu functor inclusió és fidel però no ple, perquè deixa fora les aplicacions que no són exhaustives.
\end{exemple}

\begin{proposicio}\label{prop:plenfidel-reflecteix-iso}\index{functor!ple i fidel}
Sigui $F\colon\cat{C}\to\cat{D}$ un functor ple i fidel, i sigui $f\colon A\to B$ un morfisme de $\cat{C}$. Aleshores
\[
f\ \text{és un isomorfisme a }\cat{C}
\quad\Longleftrightarrow\quad
F(f)\ \text{és un isomorfisme a }\cat{D}.
\]
\end{proposicio}

\begin{prova}
La implicació cap a la dreta val per a qualsevol functor: és la proposició~\ref{prop:functor-preserva-iso}, i no fa servir que $F$ sigui ple ni fidel. Vegem el recíproc, que sí que ho necessita. Suposem que $F(f)\colon F(A)\to F(B)$ té un invers $h\colon F(B)\to F(A)$ a $\cat{D}$. Com que $F$ és ple, existeix $g\colon B\to A$ a $\cat{C}$ amb $F(g)=h$. Aleshores
\[
F(g\comp f)=F(g)\comp F(f)=h\comp F(f)=\id_{F(A)}=F(\id_A),
\]
i com que $F$ és fidel, $g\comp f=\id_A$. Anàlogament $f\comp g=\id_B$. Per tant $f$ és un isomorfisme amb invers $g$.
\end{prova}

\begin{observacio}
Les dues direccions tenen estatus diferent. La directa ---preservar isomorfismes--- val per a tot functor; el que caracteritza els plens i fidels és el recíproc: \emph{reflecteixen} isomorfismes, i no només els preserven. Un functor ple i fidel reprodueix bijectivament els conjunts de morfismes entre els objectes de la seva imatge, però això \emph{no} implica que sigui injectiu sobre objectes: pot enviar objectes diferents a un mateix objecte. Aquesta situació ---morfismes identificats bijectivament, objectes no necessàriament--- conduirà més endavant a la noció d'\emph{equivalència de categories}.
\end{observacio}

Ser ple i fidel descriu com actua un functor sobre els morfismes. Una tercera condició, sobre els objectes, completa la terna que caracteritzarà les equivalències.

\begin{definicio}[Functor essencialment exhaustiu]\label{def:ess-exhaustiu}\index{functor!essencialment exhaustiu}
Un functor $F\colon\cat{C}\to\cat{D}$ és \textbf{essencialment exhaustiu} si tot objecte de $\cat{D}$ és isomorf a la imatge d'algun objecte de $\cat{C}$:
\[
\forall D\in\cat{D},\ \exists A\in\cat{C}\ \text{tal que}\ F(A)\cong D.
\]
\end{definicio}

\begin{observacio}
La condició demana \emph{isomorfisme}, no igualtat: no cal que tot objecte de $\cat{D}$ sigui exactament una imatge, només que ho sigui llevat d'isomorfisme. Un functor alhora ple, fidel i essencialment exhaustiu recull bijectivament els morfismes i, a més, arriba a tots els objectes llevat d'isomorfisme; aquestes són precisament les condicions que, un cop disposem de les transformacions naturals, caracteritzaran les \emph{equivalències de categories}.
\end{observacio}

\begin{exemple}[Un functor essencialment exhaustiu que no és exhaustiu sobre objectes]\label{ex:ess-exhaustiu-finset}\index{functor!essencialment exhaustiu}
Considerem la subcategoria plena $\cat{S}$ de $\cat{FinSet}$ que té per objectes només els conjunts
\[
\underline{0}=\varnothing,\qquad \underline{n}=\{1,\dots,n\}\quad(n\geq 1)
\]
---un de sol de cada mida, en lloc de tots els conjunts finits---, i el seu functor inclusió $\iota\colon\cat{S}\hookrightarrow\cat{FinSet}$. Aquest functor \emph{és} essencialment exhaustiu: tot conjunt finit $X$ té algun cardinal $n$, i aleshores és isomorf (bijectiu) a $\underline{n}=\iota(\underline{n})$. En canvi, $\iota$ no és exhaustiu sobre objectes: el conjunt $\{a,b\}$, posem per cas, no és \emph{igual} a cap $\underline{n}$, tot i ser isomorf a $\underline{2}$. Aquest és, doncs, l'exemple prototípic de la diferència entre \enquote{arribar a tots els objectes} i \enquote{arribar-hi llevat d'isomorfisme}: $\iota$ deixa fora molts objectes com a imatges exactes, però cap classe d'isomorfisme. L'exercici resolt~\ref{pr-esquelet-finset} de la Pràctica comprova, a més, que $\iota$ és ple i fidel.
\end{exemple}

\begin{exemple}[Un functor que no és essencialment exhaustiu]\index{functor!essencialment exhaustiu}
En canvi, la inclusió $\cat{FinSet}\hookrightarrow\cat{Set}$ \emph{no} és essencialment exhaustiva: un conjunt infinit no és isomorf a cap conjunt finit, de manera que cap objecte finit de la imatge no hi pot arribar. Aquí falla, doncs, la tercera condició, tot i que aquest functor sí que és ple i fidel (exemple~\ref{ex:inclusio-finset-plena}). Ser ple i fidel no diu res sobre quins objectes s'assoleixen: són condicions independents.
\end{exemple}

\subsection{Què conserva un functor: un resum}\label{subsec:que-conserva}

Convé aclarir, abans de tancar el capítol, quines propietats respecta un functor \emph{qualsevol} i quines en demanen més. Un functor transporta qualsevol configuració finita de morfismes i les equacions de composició que aquests morfismes satisfan. Per això preserva, per exemple, isomorfismes, seccions i retraccions: la propietat està testimoniada per morfismes de la categoria d'origen ---l'invers, la secció o la retracció--- i per una equació com $g\comp f=\id$, i el functor transporta tant el testimoni com l'equació. En canvi, una propietat com ser mono o epi quantifica sobre totes les fletxes possibles d'un cert tipus (\enquote{per a tota parella $g,h$}), incloses fletxes de la categoria d'arribada que poden no provenir de la categoria d'origen; i és aquesta quantificació la que un functor no pot garantir.

\begin{itemize}
\item Un functor \emph{qualsevol} preserva composició i identitats, diagrames commutatius, isomorfismes, i seccions i retraccions (és a dir, els epimorfismes i monomorfismes \emph{escindits}). En tots els casos, la propietat està testimoniada per morfismes i equacions de la categoria d'origen, que el functor transporta.
\item Ja un functor \emph{fidel} ---sense necessitat de ser ple--- \emph{reflecteix} monomorfismes i epimorfismes: si $F(f)$ és mono, aleshores $f$ ja ho era, perquè de $f\comp g=f\comp h$ se'n dedueix $F(f)\comp F(g)=F(f)\comp F(h)$, d'on $F(g)=F(h)$ i, com que $F$ és fidel, $g=h$; l'argument per als epimorfismes és dual.
\item Un functor \emph{ple i fidel} reflecteix, a més, isomorfismes, seccions i retraccions: si la imatge d'un morfisme és un isomorfisme, una secció o una retracció, el morfisme original ja ho era. Aquí calen totes dues hipòtesis: ser ple permet aixecar l'invers ---unilateral o bilateral--- que \emph{a priori} només viu al codomini, i ser fidel permet reflectir l'equació que el defineix. Recupera, doncs, aquestes nocions a banda i banda.
\item Al capítol~\ref{cap:transformacions} veurem que un functor \emph{ple, fidel i essencialment exhaustiu} determina una \emph{equivalència de categories} (en general, amb l'ajut de l'axioma de l'elecció). Els capítols posteriors precisaran, per a cada construcció, quines propietats es preserven i es reflecteixen sota equivalència.

\end{itemize}

En canvi, un functor qualsevol \emph{no} preserva, en general, ni els monomorfismes ni els epimorfismes ni les propietats universals. Vegem-ho amb els epimorfismes.

\begin{exemple}[Un functor que no preserva epimorfismes]\index{functor!no preserva epimorfismes}
Prenem el monoide additiu $(\mathbb{N},+,0)$ i el monoide $(\{0,1\},\vee,0)$, on $\vee$ és el màxim (o disjunció), tots dos vistos com a categories d'un sol objecte, $\cat{B}\mathbb{N}$ i $\cat{B}M$. L'aplicació $h\colon\mathbb{N}\to\{0,1\}$ amb $h(0)=0$ i $h(n)=1$ per a $n\geq 1$ és un homomorfisme de monoides, i defineix per tant un functor $F\colon\cat{B}\mathbb{N}\to\cat{B}M$ (exemple~\ref{ex:functors-monoide}).

A $\cat{B}\mathbb{N}$, el morfisme $1$ és un epimorfisme: la cancel·lació per la dreta a $\cat{B}\mathbb{N}$ és la cancel·lació de la suma, i a $(\mathbb{N},+)$ es pot cancel·lar. La seva imatge $F(1)=1$, en canvi, \emph{no} és un epimorfisme a $\cat{B}M$, perquè $1\vee 0=1=1\vee 1$ amb $0\neq 1$: la cancel·lació falla. Així, $F$ envia un epimorfisme a un morfisme que no ho és.

L'arrel del fenomen és la que hem assenyalat: \enquote{ser epi} quantifica sobre totes les parelles $g,h$ de la categoria d'arribada, i el functor no controla les parelles noves que apareixen a $\cat{B}M$. Un epimorfisme \emph{escindit}, en canvi ---testimoniat per una secció $s$ amb $f\comp s=\id$---, sí que es preservaria, perquè el functor transporta la secció i l'equació.
\end{exemple}

\begin{resumcap}
\apartat{Idees centrals}
\begin{enumerate}
  \item Un functor preserva dominis, codominis, composició i identitats.
  \item Les categories petites i els functors formen una categoria, $\cat{Cat}$, que és gran però localment petita; els functors entre categories grans són igualment legítims, tot i que no són fletxes de $\cat{Cat}$.
  \item El graf subjacent i la categoria lliure formen una parella de functors en sentits contraris, $\Free\colon\cat{Graph}\to\cat{Cat}$ i $U\colon\cat{Cat}\to\cat{Graph}$. No són inversos l'un de l'altre, però estan lligats per l'avaluació $\Free(U(\cat{C}))\to\cat{C}$; aquesta relació reapareixerà més endavant com a adjunció.
  \item Un functor contravariant és un functor des de l'oposada. Per comprovar qualsevol functor, primer es verifiquen els tipus ($F(f)\colon F(A)\to F(B)$, o $F(B)\to F(A)$ en el cas contravariant) i després les identitats; el functor de parts té una versió de cada variància.
  \item Ser fidel o ser ple és una propietat de les aplicacions $F_{A,B}$ entre homconjunts, i no diu res de l'acció sobre objectes; afegint-hi l'exhaustivitat essencial, que sí que parla dels objectes (llevat d'isomorfisme), s'obtenen les equivalències.
  \item L'homfunctor covariant $\Hom(A,-)$ i el contravariant $\Hom(-,B)$ transformen isos en bijeccions, i el bihomfunctor $\Hom(-,-)$ els unifica en un functor de dues variables; tots tres són fonamentals i els retrobarem al lema de Yoneda.
\end{enumerate}
\apartat{Errors típics} Confondre \enquote{isomorfisme de categories} amb \enquote{equivalència}; creure que \enquote{fidel} vol dir injectiu sobre objectes o que \enquote{ple} vol dir exhaustiu sobre objectes, quan totes dues nocions parlen només dels homconjunts; comprovar les equacions functorials abans d'haver comprovat que $F(f)$ té el tipus correcte; oblidar la inversió de sentit en un functor contravariant; barrejar $\cat{C}/A$ (sobre $A$) amb $A/\cat{C}$ (sota $A$); i confondre \emph{preservar} amb \emph{reflectir} una propietat: tot functor preserva els isomorfismes; els functors plens i fidels també els reflecteixen, tot i que aquesta condició és suficient i no necessària. Un functor qualsevol no preserva, en general, ni els monomorfismes ni els epimorfismes.
\apartat{Lectura recomanada} \cite[part~III, sessions~11--17]{lawvereschanuel} per a exemples molt elementals d'aplicacions que preserven estructura entre grafs i entre sistemes dinàmics; \cite[§1.3 i §§7.1--7.2]{awodey} per a la definició de functor, la categoria de categories i l'estructura representable; \cite[I.3, II.2 i II.5]{maclane} per a la formulació clàssica, inclosos els functors contravariants; \cite[§1.2]{leinster} per als functors i les seves propietats, inclosos els contravariants i els plens i fidels; \cite[cap.~1]{riehl} per als homfunctors i el bifunctor $\Hom$; \cite[§§3.1--3.3]{barrwells} per a més exemples de functors, en particular les accions de monoides i els functors entre grafs, i per a la classificació dels functors en fidels, plens i altres tipus; \cite[§1.3 i §§1.5.1--1.5.3]{perrone-notes} per a una segona explicació de la functorialitat, amb una secció sobre què \emph{no} és un functor i una altra sobre functors, monos i epis.
\end{resumcap}

\chapter[Transformacions naturals]{\texorpdfstring{Transformacions naturals\\ i equivalències}{Transformacions naturals i equivalències}}
\label{cap:transformacions}

\begin{objectius}
\apartat{Prerequisits} Categories i diagrames commutatius (cap.~\ref{cap:categories}); functors, composició de functors, functors covariants i contravariants, i les nocions de functor ple, fidel i essencialment exhaustiu (cap.~\ref{cap:functors}). Els homfunctors intervenen en alguns exemples i exercicis, i seran importants en els capítols següents, però no són necessaris per a la definició de naturalitat. Alguns exemples opcionals fan servir una familiaritat bàsica amb espais vectorials, duals i biduals, i l'exemple d'interpretació de fórmules de primer ordre, marcat com a ampliació, pressuposa la semàntica d'aquesta lògica.
\apartat{Objectius} En acabar el capítol, el lector hauria de poder:
\begin{itemize}
  \item definir una transformació natural i comprovar-ne la naturalitat mitjançant quadrats commutatius;
  \item compondre transformacions verticalment i interpretar-les com a morfismes de la categoria de functors $\cat{D}^{\cat{C}}$;
  \item reconèixer un isomorfisme natural i distingir-lo d'una família d'isomorfismes que no sigui natural;
  \item treballar amb la composició horitzontal i distingir-la de la vertical;
  \item definir una equivalència de categories mitjançant un quasiinvers, i caracteritzar-la com un functor ple, fidel i essencialment exhaustiu (amb l'axioma de l'elecció per a la implicació recíproca).
\end{itemize}
\end{objectius}

\section{Cap al concepte de transformació natural}\label{sec:cap-transformacio}

Al capítol anterior vam veure que els functors són els morfismes entre categories: relacionen dues categories respectant-ne tota l'estructura. Ara bé, sovint tenim \emph{dos} functors $F,G\colon\cat{C}\to\cat{D}$ que van de la mateixa categoria a la mateixa categoria, i volem comparar-los. La pregunta natural és, doncs, si hi ha una noció de \emph{morfisme entre functors}.

Aquesta idea ---comparar functors de manera coherent--- va ser històricament una de les motivacions originals per introduir les categories. De fet, es diu sovint que les categories es van definir per poder definir els functors, i els functors per poder definir les transformacions naturals; totes tres nocions apareixen alhora, per primera vegada, a l'article fundacional d'Eilenberg i Mac~Lane~\cite{eilenbergmaclane}, que precisament les motiva a partir de la relació entre un espai vectorial i el seu bidual. El mot \enquote{natural} té aquí un sentit tècnic precís que farem explícit: una comparació és natural quan fa commutar els quadrats corresponents a tots els morfismes. Sovint s'expressa dient que \enquote{no depèn de cap elecció arbitrària}; convé retenir, però, que es tracta d'una intuïció orientadora i no d'una definició: una construcció pot haver requerit eleccions i, tanmateix, resultar natural. La condició precisa és sempre la commutativitat dels quadrats.

Abans de formalitzar-ho, val la pena veure un exemple que fa palès aquest \enquote{principi de coherència}, aquest cop entre \emph{sintaxi} i \emph{semàntica}. Treballem amb la lògica proposicional. Prenem com a categoria de base $\cat{FinSet}$, els conjunts finits amb les aplicacions entre ells (subsecció~\ref{subsec:conjunts}); ara en pensem els objectes com a conjunts de variables proposicionals, i els morfismes $f\colon X\to Y$ com a \emph{canvis de variables}, que a cada variable de $X$ n'assignen una de $Y$.

El primer functor, el \emph{sintàctic}, associa a cada conjunt de variables el conjunt de les fórmules que s'hi poden escriure:
\[
F(X)=\{\text{fórmules proposicionals amb variables de }X\},
\]
construïdes a partir de les variables de $X$ amb les connectives habituals. Sobre morfismes, un canvi de variables $f\colon X\to Y$ indueix l'aplicació de \emph{substitució}
\[
F(f)\colon F(X)\longrightarrow F(Y),
\qquad
F(f)(\varphi)=\varphi[f],
\]
on $\varphi[f]$ és la fórmula obtinguda reemplaçant dins $\varphi$ cada variable $x$ per la variable $f(x)$. Aquesta substitució és la mateixa operació que ja vam trobar a l'exemple~\ref{ex:deduibilitat-functors}, on una substitució de variables proposicionals induïa un functor entre categories de deduïbilitat perquè transformava deduccions; ara la mirem d'una altra manera, com l'acció sobre morfismes d'un functor $F\colon\cat{FinSet}\to\cat{Set}$. Formalment, $F(f)$ es defineix per recursió sobre l'estructura de $\varphi$: sobre una variable, $F(f)(x)=f(x)$, i $F(f)$ commuta amb les connectives.

El segon functor, el \emph{semàntic}, associa a cada conjunt de variables el conjunt de les funcions de veritat corresponents:
\[
G(X)=\{\text{funcions de veritat sobre }X\},
\]
on una \emph{funció de veritat} sobre $X$ és una aplicació
\[
\varphi\colon \{0,1\}^{X}\longrightarrow\{0,1\}
\]
que assigna un valor de veritat a cada valoració de les variables, entenent per \emph{valoració} un element $v\in\{0,1\}^{X}$, és a dir una aplicació $v\colon X\to\{0,1\}$ que fixa el valor de cada variable. Dit altrament, $G(X)=\{0,1\}^{(\{0,1\}^{X})}$. Sobre morfismes, un canvi de variables $f\colon X\to Y$ actua per \emph{precomposició en les valoracions}: tota valoració $w\colon Y\to\{0,1\}$ es restringeix a la valoració $w\comp f\colon X\to\{0,1\}$, i definim
\[
G(f)\colon G(X)\longrightarrow G(Y),
\qquad
G(f)(\theta)(w)=\theta(w\comp f),
\]
és a dir, $G(f)(\theta)=\theta\comp(-\comp f)$.

Tots dos són functors covariants $\cat{FinSet}\to\cat{Set}$. La comprovació és rutinària ---per recursió sobre les fórmules en el cas de $F$, i per l'associativitat de la composició en el de $G$--- i es fa a l'exercici resolt~\ref{pr-sintaxi-semantica} de la Pràctica.

A cada conjunt de variables $X$ li associem l'aplicació que interpreta cada fórmula com la seva funció de veritat (la seva taula de veritat):
\[
\eta_X\colon F(X)\longrightarrow G(X),
\qquad
\eta_X(\varphi)(v)=
\begin{cases}
1 & \text{si }v\models\varphi,\\
0 & \text{si }v\not\models\varphi,
\end{cases}
\]
on $v\models\varphi$ vol dir que la valoració $v$ satisfà la fórmula $\varphi$.

El fet clau és que aquesta interpretació és \emph{compatible amb el canvi de variables}: per a tot morfisme $f\colon X\to Y$ de $\cat{FinSet}$, el quadrat
\[
\begin{tikzpicture}[baseline=(m.center)]
  \matrix (m) [matrix of math nodes, row sep=2.8em, column sep=3.6em] {
    F(X) & F(Y) \\
    G(X) & G(Y) \\
  };
  \draw[-{Stealth[scale=1.1]}] (m-1-1) -- node[above] {$F(f)$} (m-1-2);
  \draw[-{Stealth[scale=1.1]}] (m-1-1) -- node[left] {$\eta_X$} (m-2-1);
  \draw[-{Stealth[scale=1.1]}] (m-1-2) -- node[right] {$\eta_Y$} (m-2-2);
  \draw[-{Stealth[scale=1.1]}] (m-2-1) -- node[below] {$G(f)$} (m-2-2);
\end{tikzpicture}
\]
commuta, és a dir $G(f)\comp\eta_X=\eta_Y\comp F(f)$. Avaluats sobre una fórmula $\varphi\in F(X)$ i una valoració $w\colon Y\to\{0,1\}$, els dos camins del quadrat diuen si $w\models\varphi[f]$ i si $(w\comp f)\models\varphi$, respectivament. Que coincideixin és exactament el \emph{lema de substitució} de la lògica proposicional,
\[
(w\comp f)\models\varphi
\quad\Longleftrightarrow\quad
w\models\varphi[f],
\]
que es demostra per recursió sobre $\varphi$; el càlcul complet és a l'exercici resolt~\ref{pr-sintaxi-semantica}.

En termes informals: substituir primer i interpretar després dona la mateixa funció de veritat que interpretar primer i transportar després. Per exemple, de $\varphi=p\land q$ i el canvi de variables $p\mapsto r,\ q\mapsto s$ obtenim $r\land s$, i tant se val si calculem la taula de veritat abans o després de substituir: el resultat és la mateixa funció (la que val $1$ només quan $r$ i $s$ són totes dues certes). Si el canvi identifica variables, diguem $p\mapsto r,\ q\mapsto r$, aleshores $\varphi$ es transforma en $r\land r$, amb taula de veritat la de $r$; també aquí interpretar i transportar coincideixen. Aquesta compatibilitat, per a tot canvi de variables alhora, és el que a la secció següent anomenarem la \textbf{naturalitat} de la família $(\eta_X)$: \emph{substituir variables i després calcular la funció de veritat és el mateix que calcular-la primer i transportar-la després al llarg del canvi de variables}.

Val la pena notar una diferència amb els exemples \enquote{d'identificació canònica} que veurem després (com el bidual d'un espai vectorial): aquí $\eta_X$ \emph{no} és injectiva ---dues fórmules diferents poden tenir la mateixa taula de veritat, com $p\lor q$ i $q\lor p$---, de manera que no estableix cap isomorfisme entre sintaxi i semàntica. El que expressa la naturalitat no és una identificació, sinó que la interpretació és compatible amb \emph{qualsevol} substitució o canvi de variables, incloses les substitucions que identifiquen variables. La mateixa idea s'estén, amb més maquinària, a la lògica de primer ordre, on les fórmules s'interpreten com els predicats que defineixen sobre una estructura fixada; ho veurem desplegat com a exemple formal a l'exemple~\ref{ex:interpretacio-primer-ordre} de la secció~\ref{sec:exemples-transf}.

\begin{observacio}[Per què estudiar lògica categòrica]
Un lector amb formació lògica pot objectar que, al capdavall, la feina l'ha feta el lema de substitució, i preguntar-se què hi guanya el llenguatge categòric. La resposta assenyala un canvi de punt de vista. Aquí hem hagut de definir la substitució per recursió i comprovar-ne la compatibilitat a mà perquè encara no disposem de cap maquinària; però la lògica categòrica reformula la substitució com una operació \emph{estructural} ---la composició, o més en general la reindexació al llarg d'un morfisme de la categoria de contextos--- i no com una manipulació sintàctica de fórmules. En aquest marc, la interpretació és un functor (o un morfisme entre estructures més riques) que \emph{preserva} aquella operació, i naturalitats com la d'aquest exemple deixen de demostrar-se cas per cas: s'obtenen com a corol·lari que la interpretació respecta l'estructura. La inducció sobre fórmules es fa aleshores una sola vegada i en abstracte, en establir que la categoria sintàctica té la propietat universal que la caracteritza. És el gir que fa que la lògica moderna no es defineixi per les seves peculiaritats sintàctiques, sinó per les estructures que preserva; el nostre lema de substitució n'és el cas visible més senzill.
\end{observacio}

\section{Definició de transformació natural}

Siguin $F,G\colon\cat{C}\to\cat{D}$ dos functors. Volem, per a cada objecte $A$ de $\cat{C}$, un morfisme de $\cat{D}$ que vagi de $F(A)$ a $G(A)$; i volem que aquests morfismes siguin compatibles amb l'acció de $F$ i $G$ sobre les fletxes.

\begin{definicio}\index{transformació natural}
\label{def:transf-natural}
Una \textbf{transformació natural} $\eta\colon F\Rightarrow G$ entre dos functors $F,G\colon\cat{C}\to\cat{D}$ és una \textbf{família de morfismes de $\cat{D}$}
\[
\eta=\bigl(\eta_A\bigr)_{A\in\Ob(\cat{C})},
\qquad
\eta_A\colon F(A)\to G(A),
\]
indexada pels objectes de $\cat{C}$ ---un morfisme $\eta_A$ per a cada objecte $A$, de $F(A)$ a $G(A)$, anomenat el \textbf{component} de $\eta$ a $A$---, subjecta a la condició següent: per a cada morfisme $f\colon A\to B$ de $\cat{C}$, el quadrat
\[
\begin{tikzpicture}[baseline=(m.center)]
  \matrix (m) [matrix of math nodes, row sep=2.8em, column sep=3.4em] {
    F(A) & F(B) \\
    G(A) & G(B) \\
  };
  \draw[-{Stealth[scale=1.1]}] (m-1-1) -- node[above] {$F(f)$} (m-1-2);
  \draw[-{Stealth[scale=1.1]}] (m-1-1) -- node[left] {$\eta_A$} (m-2-1);
  \draw[-{Stealth[scale=1.1]}] (m-1-2) -- node[right] {$\eta_B$} (m-2-2);
  \draw[-{Stealth[scale=1.1]}] (m-2-1) -- node[below] {$G(f)$} (m-2-2);
\end{tikzpicture}
\]
commuta; és a dir,
\[
G(f)\comp\eta_A=\eta_B\comp F(f).
\]
Aquesta igualtat s'anomena \textbf{condició de naturalitat}, i el quadrat, \textbf{quadrat de naturalitat}.
\end{definicio}

El quadrat de naturalitat exhibeix una transformació natural morfisme a morfisme, i és la manera pràctica de comprovar-la. Sovint, però, interessa representar tota la transformació $\eta$ d'un sol traç, com un objecte que \emph{relaciona} dos functors paral·lels. Dibuixarem, doncs, una transformació natural \emph{diagramàticament} d'aquesta manera: els dos functors $F,G\colon\cat{C}\to\cat{D}$ es representen com dues fletxes de $\cat{C}$ cap a $\cat{D}$, i la transformació natural com una fletxa doble $\Rightarrow$ que va de l'una cap a l'altra:
\[
\begin{tikzpicture}[baseline=(A.base)]
  \node (A) at (0,0) {$\cat{C}$};
  \node (B) at (2.9,0) {$\cat{D}$};
  \draw[cat] (A) to[bend left=26] node[above] {$F$} (B);
  \draw[cat] (A) to[bend right=26] node[below] {$G$} (B);
  \draw[cat nat] (1.45,0.2) -- node[right=1.5pt] {$\eta$} (1.45,-0.2);
\end{tikzpicture}
\]
És la primera vegada que apareix una fletxa doble entre fletxes: reflecteix que $\eta$ viu \enquote{un nivell més amunt} que els morfismes ordinaris. Aquest punt de vista ---objectes, functors i, damunt, transformacions naturals--- és el que més endavant resumirem dient que les categories formen una \emph{2-categoria}, i és la notació natural per a la composició horitzontal i per a la composició d'una transformació natural amb un functor que veurem al final del capítol.

Convé precisar bé què exigeix la condició de naturalitat i què no: les components $\eta_A$ es trien un per a cada objecte, però no lliurement, ja que han d'encaixar tots alhora amb l'acció dels functors sobre els morfismes. Ara bé, la naturalitat només demana que les components escollides siguin compatibles amb tots els morfismes; no demana que siguin l'única elecció possible. És una condició de compatibilitat, no d'unicitat ni d'absència de paràmetres, com mostra l'exemple següent.

Sobre els espais vectorials reals, prenem com a functors $F=G=\id_{\cat{Vect}_{\mathbb R}}$, el functor identitat de la categoria $\cat{Vect}_{\mathbb R}$, i busquem una transformació natural $\eta\colon\id\Rightarrow\id$. Fixem un escalar $\lambda\in\mathbb{R}$ i prenem, per a cada espai $V$, la component
\[
\eta_V=\lambda\,\id_V\colon V\to V,
\]
que és una aplicació lineal ---la multiplicació per l'escalar $\lambda$---, i per tant un morfisme de $\cat{Vect}_{\mathbb R}$, amb el tipus correcte perquè $F(V)=G(V)=V$. En efecte, aquesta família és una transformació natural: per a tota aplicació lineal $T\colon V\to W$ el quadrat de naturalitat és
\[
\begin{tikzpicture}[baseline=(m.center)]
  \matrix (m) [matrix of math nodes, row sep=2.8em, column sep=3.8em] {
    V & W \\
    V & W \\
  };
  \draw[-{Stealth[scale=1.1]}] (m-1-1) -- node[above] {$T$} (m-1-2);
  \draw[-{Stealth[scale=1.1]}] (m-1-1) -- node[left] {$\lambda\,\id_V$} (m-2-1);
  \draw[-{Stealth[scale=1.1]}] (m-1-2) -- node[right] {$\lambda\,\id_W$} (m-2-2);
  \draw[-{Stealth[scale=1.1]}] (m-2-1) -- node[below] {$T$} (m-2-2);
\end{tikzpicture}
\]
i commuta, perquè
\[
T\comp(\lambda\,\id_V)=\lambda\,T=(\lambda\,\id_W)\comp T,
\]
on les dues fletxes horitzontals són $T$ perquè el functor identitat actua sobre $T$ deixant-lo igual. Quan $\lambda\neq 0$, cada component $\eta_V=\lambda\,\id_V$ és a més un isomorfisme d'espais vectorials, amb invers $\lambda^{-1}\id_V$; això farà de $\eta$ un exemple del que més endavant anomenarem un \emph{isomorfisme natural}. Tenim, doncs, tota una família de transformacions naturals, una per cada valor de $\lambda$: la naturalitat es compleix per a qualsevol $\lambda$ fixat, tot i que la construcció depèn d'aquesta elecció. Naturalitat no vol dir, per tant, \enquote{sense eleccions}, sinó \enquote{compatible amb tots els morfismes}.


\begin{exemple}[Una família que \emph{no} és natural]\label{ex:no-natural}\index{transformació natural!contraexemple}
És igual d'instructiu veure una família de components que \emph{falla} la condició de naturalitat. Treballem a $\cat{Set}$ amb $F=G=\id_{\cat{Set}}$, de manera que una transformació $\id\Rightarrow\id$ hauria de donar, per a cada conjunt $A$, una aplicació $\eta_A\colon A\to A$ que commuti amb \emph{totes} les aplicacions. Definim una família completament explícita, sense cap tria:
\[
\eta_A=\id_A\quad\text{si } A\neq\{0,1\},
\qquad\qquad
\eta_{\{0,1\}}=\tau,\quad \tau(0)=1,\ \tau(1)=0.
\]
És a dir, totes les components són identitats excepte la del conjunt $\{0,1\}$, que és la transposició dels dos elements. És una família de morfismes $\eta_A\colon A\to A$ indexada per \emph{tots} els conjunts, i cada component és fins i tot una bijecció.

Vegem què demana el quadrat de naturalitat per a una aplicació $f\colon A\to B$. Com que $F=G=\id_{\cat{Set}}$, es té $F(f)=G(f)=f$, i la condició és $f\comp\eta_A=\eta_B\comp f$. Prenem $A=B=\{0,1\}$ i l'aplicació constant de valor $0$, $f(0)=f(1)=0$:
\[
\begin{tikzpicture}[baseline=(m.center)]
  \matrix (m) [matrix of math nodes, row sep=2.8em, column sep=3.8em] {
    \{0,1\} & \{0,1\} \\
    \{0,1\} & \{0,1\} \\
  };
  \draw[-{Stealth[scale=1.1]}] (m-1-1) -- node[above] {$f$} (m-1-2);
  \draw[-{Stealth[scale=1.1]}] (m-1-1) -- node[left] {$\tau$} (m-2-1);
  \draw[-{Stealth[scale=1.1]}] (m-1-2) -- node[right] {$\tau$} (m-2-2);
  \draw[-{Stealth[scale=1.1]}] (m-2-1) -- node[below] {$f$} (m-2-2);
\end{tikzpicture}
\qquad
f\comp\tau=f,\qquad \tau\comp f=\text{constant de valor }1.
\]
El camí de l'esquerra i a baix dona $f\comp\tau$, que continua sent constant de valor $0$; el de dalt i a la dreta dona $\tau\comp f$, que és constant de valor $1$. Els dos camins difereixen, i la família $(\eta_A)$ \emph{no} és una transformació natural.

La lliçó és que tenir una regla objecte a objecte, fins i tot formada per bijeccions, no basta. La component a $\{0,1\}$ s'ha definit sense tenir en compte les aplicacions que entren i surten d'aquest conjunt, i una sola d'aquestes aplicacions n'hi ha prou per trencar la compatibilitat. La naturalitat és exactament aquesta compatibilitat amb \emph{totes} les fletxes.
\end{exemple}

\section{Exemples de transformacions naturals}\label{sec:exemples-transf}

\begin{exemple}[Transformació identitat]\index{transformació natural!identitat}
Per a tot functor $F\colon\cat{C}\to\cat{D}$ hi ha la \textbf{transformació natural identitat} $\id_F\colon F\Rightarrow F$, amb components $(\id_F)_A=\id_{F(A)}$. El quadrat de naturalitat commuta trivialment, perquè els dos camins són $F(f)$.
\end{exemple}

\begin{exemple}[Preordres]\index{transformació natural!entre preordres}
Siguin $P,Q$ dos preordres i $F,G\colon P\to Q$ dos functors, és a dir, dues aplicacions monòtones. Una transformació natural $\eta\colon F\Rightarrow G$ dona, per a cada $x\in P$, un morfisme $\eta_x\colon F(x)\to G(x)$; però a $Q$ hi ha una fletxa $F(x)\to G(x)$ exactament quan $F(x)\leq G(x)$. La condició de naturalitat es compleix automàticament, perquè entre dos objectes d'un preordre hi ha com a màxim una fletxa. Per tant,
\[
\text{existeix }\eta\colon F\Rightarrow G
\quad\Longleftrightarrow\quad
F(x)\leq G(x)\ \text{per a tot }x.
\]
En el llenguatge de l'ordre, una transformació natural entre aplicacions monòtones és exactament la relació \enquote{$F$ és puntualment menor o igual que $G$}.
\end{exemple}

\begin{exemple}[Ampliació: interpretació de fórmules de primer ordre]\label{ex:interpretacio-primer-ordre}\index{transformació natural!interpretació lògica}\index{lema de substitució}
Reprenem el fil de la lògica que hem obert a la secció~\ref{sec:cap-transformacio} (\enquote{Cap al concepte de transformació natural}), ara en primer ordre i amb la naturalitat plantejada com a exemple formal. Fixem un llenguatge de primer ordre i una estructura $M$. L'exemple pressuposa la semàntica de la lògica de primer ordre (estructures, satisfacció, variables lliures i lligades) i es pot ometre en una primera lectura.

Convé fixar d'entrada les convencions que fan functorial la construcció sintàctica, en lloc d'adoptar-les a posteriori. Disposem d'una reserva infinita de variables; un \emph{context} és un conjunt finit $X$ de variables d'aquesta reserva, i escrivim $\cat{Ctx}$ per a la subcategoria plena de $\cat{FinSet}$ que té per objectes els contextos i per morfismes totes les aplicacions entre ells (com que la reserva és infinita, $\cat{Ctx}$ és equivalent a $\cat{FinSet}$, però no farem servir aquesta equivalència); treballem amb fórmules \emph{mòdul reanomenament de variables lligades} ($\alpha$-equivalència), i totes les substitucions s'entenen \emph{lliures de captura}. Sense aquestes convencions les lleis functorials de sota no se sostenen: en $\varphi=\exists y\,R(x,y)$, substituir la variable lliure $x$ per $y$ no pot donar ingènuament $\exists y\,R(y,y)$; cal reanomenar abans la variable lligada i obtenir, per exemple, $\exists z\,R(y,z)$.

Amb aquestes convencions considerem dues construccions. Per una banda,
\[
F(X)=\{\text{fórmules amb variables lliures entre }X\},
\]
enteses com a classes d'$\alpha$-equivalència; un canvi de variables $f\colon X\to Y$ hi actua per substitució, $F(f)(\varphi)=\varphi[f]$, reemplaçant cada variable lliure $x$ per $f(x)$. Per l'altra,
\[
G(X)=\{\text{predicats definibles sobre }X\},
\]
on un \emph{predicat definible} és un subconjunt de $M^X$ (el conjunt d'assignacions $s\colon X\to M$) de la forma
\[
\llbracket\varphi\rrbracket=\{\,s\colon X\to M\mid M\models\varphi[s]\,\},
\]
és a dir, l'\emph{extensió} d'alguna fórmula. El canvi de variables $f$ actua sobre $G$ transportant assignacions: com que $f$ indueix la precomposició $M^{f}\colon M^{Y}\to M^{X}$, $s\mapsto s\comp f$, posem
\[
G(f)(S)=\bigl(M^{f}\bigr)^{-1}(S)=\{\,s\in M^{Y}\mid s\comp f\in S\,\},\qquad S\subseteq M^{X}.
\]
Totes dues són functors covariants $\cat{Ctx}\to\cat{Set}$: per a $F$ es comprova que la substitució respecta identitats i composició\footnotemark; per a $G$, la functorialitat és la de la precomposició $M^{(-)}$ seguida de preimatge.

A cada context $X$ li associem ara l'aplicació que interpreta cada fórmula com la seva extensió,
\[
\eta_X\colon F(X)\to G(X),
\qquad
\eta_X(\varphi)=\llbracket\varphi\rrbracket.
\]
Aquesta $\eta_X$ està ben definida sobre classes d'$\alpha$-equivalència, perquè dues fórmules que només difereixen en el nom de les variables lligades tenen la mateixa extensió: la semàntica és $\alpha$-invariant, i per això el pas al quocient no altera res. La condició de naturalitat, per a un canvi de variables $f\colon X\to Y$, és que interpretar la fórmula substituïda coincideixi amb transportar l'extensió de la fórmula original:
\[
\llbracket\varphi[f]\rrbracket=\{\,s\colon Y\to M\mid s\comp f\in\llbracket\varphi\rrbracket\,\}.
\]
Aquesta igualtat és exactament el \textbf{lema de substitució} de la teoria de models ---el significat d'una fórmula després de canviar-ne les variables es calcula transportant el significat original--- i fa doble feina: garanteix alhora que $G(f)$ envia predicats definibles a predicats definibles (la preimatge d'un definible és definible) i que el quadrat de naturalitat commuta. Per tant $\eta\colon F\Rightarrow G$ és una transformació natural, i expressa que la interpretació és compatible amb la substitució.

Com en el cas proposicional de la secció~\ref{sec:cap-transformacio}, $\eta_X$ no és injectiva ---dues fórmules lògicament equivalents defineixen el mateix predicat---, de manera que la naturalitat no és aquí una identificació, sinó que expressa que la interpretació és compatible amb els canvis de variables, inclosos els reanomenaments i les identificacions. Convé, finalment, no confondre dues nocions que aquí conviuen: en una estructura $M$ fixada, dues fórmules poden tenir la mateixa extensió sense ser lògicament equivalents en \emph{totes} les estructures; la relació d'\emph{equivalència semàntica relativa a $M$} és, per tant, més grollera que l'\emph{equivalència lògica} global.
\end{exemple}
\footnotetext{És on es fan servir les convencions anteriors. Quan $f$ no és injectiva, identifica variables lliures, i és la substitució lliure de captura, definida sobre classes d'$\alpha$-equivalència, la que fa que $\varphi[\id]=\varphi$ i $\varphi[g\comp f]=(\varphi[f])[g]$ valguin com a \emph{igualtats}, no només mòdul reanomenament ad hoc.}%

\begin{exemple}[Bidual d'espais vectorials]\label{ex:bidual}\index{transformació natural!bidual}
Reprenem el fil que vam deixar obert a l'exemple~\ref{ex:homfunctor-dual} (\enquote{El dual d'un espai vectorial com a homfunctor}), on vam construir el \textbf{functor dual}
\[
(-)^*=\Hom_{\mathbb K}(-,\mathbb K)\colon\cat{Vect}_{\mathbb K}\op\to\cat{Vect}_{\mathbb K},\qquad V\mapsto V^*,
\]
contravariant, que sobre un morfisme $T\colon V\to W$ dona l'aplicació dual $T^*\colon W^*\to V^*$, $\psi\mapsto\psi\comp T$, amb la regla $(g\comp f)^*=f^*\comp g^*$. Vam ajornar deliberadament la comparació de $V$ amb el bidual fins a tenir el llenguatge de les transformacions naturals; ja el tenim.

El \textbf{functor bidual} (o doble dual) s'obté per composició del dual amb ell mateix. Com que $(-)^*$ és contravariant ---un functor amb domini $\cat{Vect}_{\mathbb K}\op$---, per compondre'l amb ell mateix cal fer encaixar dominis i codominis amb l'oposat, exactament com a la proposició~\ref{prop:variancia-composicio}: es pren el functor oposat $\bigl((-)^*\bigr)\op\colon\cat{Vect}_{\mathbb K}\to\cat{Vect}_{\mathbb K}\op$ ---que actua igual sobre objectes--- i s'hi torna a aplicar el dual,
\[
(-)^{**}=(-)^*\comp\bigl((-)^*\bigr)\op\colon\cat{Vect}_{\mathbb K}\to\cat{Vect}_{\mathbb K}.
\]
Per la regla dels signes, la composició de dos functors contravariants és \emph{covariant}. Desplegada, actua sobre objectes com
\[
V\mapsto V^{**}=(V^*)^*,
\]
i sobre un morfisme $T\colon V\to W$ com el dual aplicat dues vegades,
\[
T^{**}=(T^*)^*\colon V^{**}\to W^{**},\qquad \rho\mapsto\rho\comp T^*,
\]
amb la regla $(g\comp f)^{**}=g^{**}\comp f^{**}$: invertir l'ordre dos cops el restaura, i per això el bidual és un endofunctor covariant de $\cat{Vect}_{\mathbb K}$, paral·lel a la identitat.

L'efecte sobre els morfismes ho resumeix el diagrama següent, on el dual inverteix el sentit de la fletxa i el bidual, aplicant-lo dues vegades, el recupera:
\[
\begin{tikzpicture}[baseline=(m.center)]
  \matrix (m) [matrix of math nodes, column sep=3.2em, row sep=3.0em] {
    {} & V & W & {} \\
    {} & V^* & W^* & {} \\
    {} & V^{**} & W^{**} & {} \\
  };
  \draw[cat] (m-1-2) -- node[above] {$T$} (m-1-3);
  \draw[cat] (m-2-3) -- node[above] {$T^*$} (m-2-2);
  \draw[cat] (m-3-2) -- node[above] {$T^{**}$} (m-3-3);
  \draw[cat,dashed] (m-1-1) -- node[left] {$(-)^*$} (m-2-1);
  \draw[cat,dashed] (m-2-1) -- node[left] {$(-)^*$} (m-3-1);
  \draw[cat,dashed] (m-1-4) -- node[right] {$(-)^{**}$} (m-3-4);
\end{tikzpicture}
\]
Les fletxes contínues són morfismes de $\cat{Vect}_{\mathbb K}$; les puntejades, l'acció del functor. La fila del mig apunta en sentit contrari ---el dual $(-)^*$ és contravariant---, i la de baix en recupera el sentit ---dues inversions es cancel·len---: el bidual $(-)^{**}$, composició del dual amb ell mateix, és covariant, paral·lel a la identitat. Ara sí que té sentit comparar-los amb una transformació natural. Tots dos ---la identitat i el bidual--- són endofunctors covariants de $\cat{Vect}_{\mathbb K}$, de manera que la comparació és una transformació natural $\eta$ entre functors paral·lels:
\[
\begin{tikzpicture}[baseline=(A.base)]
  \node (A) at (0,0) {$\cat{Vect}_{\mathbb K}$};
  \node (B) at (3.6,0) {$\cat{Vect}_{\mathbb K}$};
  \draw[cat] (A) to[bend left=26] node[above] {$\id_{\cat{Vect}_{\mathbb K}}$} (B);
  \draw[cat] (A) to[bend right=26] node[below] {$(-)^{**}$} (B);
  \draw[cat nat] (1.8,0.2) -- node[right=1.5pt] {$\eta$} (1.8,-0.2);
\end{tikzpicture}
\]
El component de $\eta$ a cada espai $V$ és l'\textbf{avaluació}, que envia un vector $v$ a la forma \enquote{avaluar en $v$}. La transformació natural és, doncs, la família
\[
\eta=\bigl(\eta_V\bigr)_{V\in\Ob(\cat{Vect}_{\mathbb K})},
\]
i cada component es desglossa nivell a nivell, amb la imatge de $v$ com una aplicació de $V^*$ en $\mathbb K$:
\[
\begin{array}{r c l}
\eta_V\colon V & \longrightarrow & V^{**}\\[6pt]
v & \longmapsto & \eta_V(v)\colon
  \begin{array}[t]{r c l}
    V^* & \longrightarrow & \mathbb K\\[4pt]
    \varphi & \longmapsto & \varphi(v).
  \end{array}
\end{array}
\]
Cada $\eta_V$ és lineal i està ben definida. La condició de naturalitat, per a una aplicació lineal $T\colon V\to W$, és la commutativitat del quadrat
\[
\begin{tikzpicture}[baseline=(m.center)]
  \matrix (m) [matrix of math nodes, row sep=2.8em, column sep=3.6em] {
    V & W \\
    V^{**} & W^{**} \\
  };
  \draw[-{Stealth[scale=1.1]}] (m-1-1) -- node[above] {$T$} (m-1-2);
  \draw[-{Stealth[scale=1.1]}] (m-1-1) -- node[left] {$\eta_V$} (m-2-1);
  \draw[-{Stealth[scale=1.1]}] (m-1-2) -- node[right] {$\eta_W$} (m-2-2);
  \draw[-{Stealth[scale=1.1]}] (m-2-1) -- node[below] {$T^{**}$} (m-2-2);
\end{tikzpicture}
\]
és a dir, $T^{**}\comp\eta_V=\eta_W\comp T$: prendre l'avaluació i després transportar-la pel bidual coincideix amb aplicar $T$ i després avaluar. Aquí la fletxa horitzontal superior és $T$ perquè el functor identitat actua sobre $T$ deixant-lo igual. En aquest cas, la naturalitat es pot llegir dient que la identificació de $V$ amb la seva imatge dins $V^{**}$ no fa servir cap elecció de base: la construcció $v\mapsto(\varphi\mapsto\varphi(v))$ és intrínseca. Aquesta lectura ---\enquote{independent de la base}--- és una interpretació d'aquest exemple concret, no la definició de naturalitat, que continua sent la commutativitat dels quadrats. En restringir-se als espais de dimensió finita, cada $\eta_V$ és un isomorfisme, cosa que reprendrem en parlar dels isomorfismes naturals. En canvi, la comparació de $V$ amb el \emph{dual} $V^*$ ---i no amb el bidual--- no es pot fer natural; ho precisarem a l'observació~\ref{obs:iso-no-natural}, un cop definits els isomorfismes naturals.
\end{exemple}

\section{Categoria de functors}\label{sec:categoria-functors}

Els functors entre dues categories fixes, amb les transformacions naturals com a morfismes, formen al seu torn una categoria. Aquesta secció en desenvolupa la construcció: la composició que la fa possible, la categoria $\cat{D}^{\cat{C}}$ mateixa, un exemple fonamental ---els grafs--- i els seus isomorfismes, que resulten ser els isomorfismes naturals; en tanquem l'estudi amb l'equivalència de categories.

Les transformacions naturals no només relacionen functors: també es poden \emph{compondre} entre elles. La composició bàsica apila una transformació damunt d'una altra, i és la que després convertirà els functors entre dues categories en una categoria, i la que dona sentit a parlar de la \emph{inversa} d'un isomorfisme natural.

\begin{definicio}[Composició vertical]\index{transformació natural!composició vertical}
Donades dues transformacions naturals $\eta\colon F\Rightarrow G$ i $\theta\colon G\Rightarrow H$ entre functors $\cat{C}\to\cat{D}$, la seva \textbf{composició vertical} $\theta\comp\eta\colon F\Rightarrow H$ té per components
\[
(\theta\comp\eta)_A=\theta_A\comp\eta_A.
\]
El nom \emph{vertical} ve d'aquesta representació: les dues transformacions s'apilen ---$\eta$ a dalt i $\theta$ a sota--- entre els tres functors paral·lels $F,G,H\colon\cat{C}\to\cat{D}$, i la composició $\theta\comp\eta$ les recorre de dalt a baix.
\[
\begin{tikzpicture}[baseline=(A.base)]
  \node (A) at (0,0) {$\cat{C}$};
  \node (B) at (3.4,0) {$\cat{D}$};
  \draw[cat] (A) to[bend left=44] node[above] {$F$} (B);
  \draw[cat] (A) -- (B);
  \draw[cat] (A) to[bend right=44] node[below] {$H$} (B);
  \node[fill=white,inner sep=1.5pt] at (1.7,0) {$G$};
  \draw[cat nat] (1.7,0.52) -- node[right=1.5pt] {$\eta$} (1.7,0.27);
  \draw[cat nat] (1.7,-0.27) -- node[right=1.5pt] {$\theta$} (1.7,-0.52);
\end{tikzpicture}
\]
El mateix nom es reconeix a escala de components: fixat un morfisme $f\colon A\to B$, les components de $\eta$ i de $\theta$ són fletxes \emph{verticals}, i els seus quadrats de naturalitat s'apilen l'un damunt de l'altre ---compartint la fila central $G(A)\to G(B)$---, de manera que $\theta\comp\eta$ en compon les columnes:
\[
\begin{tikzpicture}[baseline=(m.center)]
  \matrix (m) [matrix of math nodes, row sep=3em, column sep=3.6em] {
    F(A) & F(B) \\
    G(A) & G(B) \\
    H(A) & H(B) \\
  };
  \draw[cat] (m-1-1) -- node[above] {$F(f)$} (m-1-2);
  \draw[cat] (m-2-1) -- node[above] {$G(f)$} (m-2-2);
  \draw[cat] (m-3-1) -- node[below] {$H(f)$} (m-3-2);
  \draw[cat] (m-1-1) -- node[left] {$\eta_A$} (m-2-1);
  \draw[cat] (m-2-1) -- node[left] {$\theta_A$} (m-3-1);
  \draw[cat] (m-1-2) -- node[right] {$\eta_B$} (m-2-2);
  \draw[cat] (m-2-2) -- node[right] {$\theta_B$} (m-3-2);
\end{tikzpicture}
\]
Aïllant-ne el rectangle exterior, la composició $\theta\comp\eta\colon F\Rightarrow H$ té el seu propi quadrat de naturalitat (a l'esquerra), i es representa com una única cel·la de $F$ cap a $H$ (a la dreta):
\[
\begin{tikzpicture}[baseline=(m.center)]
  \matrix (m) [matrix of math nodes, row sep=2.8em, column sep=3.6em] {
    F(A) & F(B) \\
    H(A) & H(B) \\
  };
  \draw[cat] (m-1-1) -- node[above] {$F(f)$} (m-1-2);
  \draw[cat] (m-1-1) -- node[left] {$(\theta\comp\eta)_A$} (m-2-1);
  \draw[cat] (m-1-2) -- node[right] {$(\theta\comp\eta)_B$} (m-2-2);
  \draw[cat] (m-2-1) -- node[below] {$H(f)$} (m-2-2);
\end{tikzpicture}
\qquad\qquad
\begin{tikzpicture}[baseline=(A.base)]
  \node (A) at (0,0) {$\cat{C}$};
  \node (B) at (3.7,0) {$\cat{D}$};
  \draw[cat] (A) to[bend left=26] node[above] {$F$} (B);
  \draw[cat] (A) to[bend right=26] node[below] {$H$} (B);
  \draw[cat nat] (1.85,0.2) -- node[right=1.5pt] {$\theta\comp\eta$} (1.85,-0.2);
\end{tikzpicture}
\]
\end{definicio}

\begin{proposicio}
La composició vertical torna a ser una transformació natural, és associativa, i la transformació identitat $\id_F$ n'és l'element neutre.
\end{proposicio}

\begin{prova}
Per a la naturalitat n'hi ha prou d'observar que, al diagrama de components de la definició, tots dos quadrats commuten (naturalitat de $\eta$ i de $\theta$); per tant també commuta el rectangle exterior, les columnes del qual són $(\theta\comp\eta)_A$ i $(\theta\comp\eta)_B$, i això és el quadrat de naturalitat de $\theta\comp\eta$ ---equivalentment, $H(f)\comp\theta_A\comp\eta_A=\theta_B\comp\eta_B\comp F(f)$, encadenant els dos quadrats. L'associativitat i la propietat de la identitat es redueixen a les propietats corresponents de la composició a $\cat{D}$, aplicades component a component.
\end{prova}

La composició vertical de transformacions naturals, introduïda més amunt, organitza els functors entre dues categories en una nova categoria.

\begin{definicio}\index{categoria!de functors}\index{DC@$\cat{D}^{\cat{C}}$}
Sigui $\cat{C}$ una categoria petita i $\cat{D}$ una categoria qualsevol. La \textbf{categoria de functors} $\cat{D}^{\cat{C}}$ (també escrita $[\cat{C},\cat{D}]$ o $\operatorname{Fun}(\cat{C},\cat{D})$) té:
\begin{itemize}
\item per objectes, els functors $\cat{C}\to\cat{D}$;
\item per morfismes, les transformacions naturals entre ells;
\item per composició, la composició vertical, i per identitats, les transformacions identitat.
\end{itemize}
\end{definicio}

\begin{observacio}[Qüestions de mida]\label{obs:mida-functors}\index{transformació natural!Nat@$\operatorname{Nat}$}
Convé dir per què la definició demana $\cat{C}$ petita. La raó de fons és la mateixa de l'observació~\ref{obs:mida-cat}, on vam haver de restringir $\cat{Cat}$ a les categories \emph{petites}. Si $\cat{C}$ fos gran, les dades que defineixen un functor $\cat{C}\to\cat{D}$ podrien formar una classe pròpia ---l'assignació recorreria una classe pròpia d'objectes---, i les classes pròpies no poden ser elements de cap col·lecció; aleshores els functors $\cat{C}\to\cat{D}$ no es poden reunir com els \emph{objectes} d'una categoria.\footnotemark Amb $\cat{C}$ petita, en canvi, aquestes dades formen un conjunt i la construcció no té cap problema. Això no impedeix parlar \emph{individualment} de functors i transformacions naturals amb domini gran ---ho fem contínuament, com amb $\cat{Set}\to\cat{Set}$---; el que no fem és agrupar-los tots en una categoria de functors.

Amb $\cat{C}$ petita, $\cat{D}^{\cat{C}}$ pot ser encara gran, però és \textbf{localment petita} tan bon punt $\cat{D}$ ho és: una transformació natural $F\Rightarrow G$ queda determinada pels seus components $\eta_A$, un per a cada objecte $A$ de $\cat{C}$, de manera que la col·lecció de totes les transformacions naturals $F\Rightarrow G$, que denotem $\operatorname{Nat}(F,G)$, és un subconjunt del producte
\[
\operatorname{Nat}(F,G)\;\subseteq\;\prod_{A\in\Ob(\cat{C})}\Hom_{\cat{D}}(F(A),G(A)),
\]
indexat per $\Ob(\cat{C})$ ---un conjunt, precisament perquè $\cat{C}$ és petita---, pel mateix argument que feia $\cat{Cat}$ localment petita. Com que aquest producte és un conjunt, $\operatorname{Nat}(F,G)$ també ho és, i és el conjunt de morfismes de $F$ a $G$ a la categoria de functors:
\[
\operatorname{Nat}(F,G)=\Hom_{\cat{D}^{\cat{C}}}(F,G).
\]
\end{observacio}
\footnotetext{Això depèn del marc de mida, que aquí és el de classes i conjunts de l'apèndix~\ref{cap:mida}. Amb universos de Grothendieck, la categoria de functors d'una categoria gran es pot formar en un univers superior; vegeu \cite[§1.1]{riehl}.}%

\index{graf!com a categoria de functors}\index{categoria!de functors!grafs}
Un dels exemples més il·luminatius de categoria de functors mostra que un objecte aparentment aliè a les categories ---un graf dirigit--- \emph{és}, en realitat, un functor, i que els seus morfismes són transformacions naturals. Aquest exemple reprèn el fil dels grafs que hem seguit des del primer capítol i és, a més, el primer indici que categories de functors del tipus $\cat{Set}^{\cat{C}}$ codifiquen estructures matemàtiques familiars.

La categoria índex $\cat{P}$ no és nova: és la categoria de \textbf{dues fletxes paral·leles} de la subsecció~\ref{subsec:finites}, amb dos objectes i dues fletxes no trivials amb el mateix domini i el mateix codomini. Els seus dos objectes són, en si, abstractes; els rebategem $E$ i $V$ ---i les fletxes, $s$ i $t$--- per anticipar la lectura com a grafs. Com que una aresta apunta cap als seus extrems, el domini de les fletxes és l'objecte de les \emph{arestes}, $E$, i el codomini el dels \emph{vèrtexs}, $V$, amb $s,t\colon E\to V$:
\[
\begin{tikzpicture}[baseline=(E.base)]
  \node (E) at (0,0) {$E$};
  \node (V) at (2.4,0) {$V$};
  \draw[-{Stealth[scale=1.1]}] ([yshift=3pt]E.east) -- node[above] {$s$} ([yshift=3pt]V.west);
  \draw[-{Stealth[scale=1.1]}] ([yshift=-3pt]E.east) -- node[below] {$t$} ([yshift=-3pt]V.west);
\end{tikzpicture}
\]

\begin{proposicio}\index{graf!com a categoria de functors}
La categoria de functors $\cat{Set}^{\cat{P}}$ és isomorfa a la categoria $\cat{Graph}$ dels grafs dirigits i els morfismes de grafs.
\end{proposicio}

\begin{prova}
Un functor $G\colon\cat{P}\to\cat{Set}$ consisteix en dos conjunts $G(E)$ i $G(V)$ i dues aplicacions $G(s),G(t)\colon G(E)\to G(V)$. Interpretant $G(E)$ com el conjunt d'\textbf{arestes}, $G(V)$ com el conjunt de \textbf{vèrtexs} i $G(s),G(t)$ com les aplicacions \enquote{origen} i \enquote{destí}, aquesta dada és \emph{exactament} un graf dirigit. Recíprocament, tot graf dirigit determina un tal functor. Els objectes de $\cat{Set}^{\cat{P}}$ ---que són functors $\cat{P}\to\cat{Set}$--- es corresponen, doncs, bijectivament amb els grafs dirigits.

Vegem ara els morfismes. Una transformació natural $\varphi\colon G\Rightarrow H$ entre dos d'aquests functors té dos components, $\varphi_E\colon G(E)\to H(E)$ i $\varphi_V\colon G(V)\to H(V)$, i la condició de naturalitat per a les fletxes $s$ i $t$ dona els dos quadrats
\[
\begin{tikzpicture}[baseline=(m.center)]
  \matrix (m) [matrix of math nodes, row sep=2.4em, column sep=3em] {
    G(E) & G(V) \\
    H(E) & H(V) \\
  };
  \draw[-{Stealth[scale=1.1]}] (m-1-1) -- node[above] {$G(s)$} (m-1-2);
  \draw[-{Stealth[scale=1.1]}] (m-1-1) -- node[left] {$\varphi_E$} (m-2-1);
  \draw[-{Stealth[scale=1.1]}] (m-1-2) -- node[right] {$\varphi_V$} (m-2-2);
  \draw[-{Stealth[scale=1.1]}] (m-2-1) -- node[below] {$H(s)$} (m-2-2);
\end{tikzpicture}
\qquad\text{i}\qquad
\begin{tikzpicture}[baseline=(n.center)]
  \matrix (n) [matrix of math nodes, row sep=2.4em, column sep=3em] {
    G(E) & G(V) \\
    H(E) & H(V) \\
  };
  \draw[-{Stealth[scale=1.1]}] (n-1-1) -- node[above] {$G(t)$} (n-1-2);
  \draw[-{Stealth[scale=1.1]}] (n-1-1) -- node[left] {$\varphi_E$} (n-2-1);
  \draw[-{Stealth[scale=1.1]}] (n-1-2) -- node[right] {$\varphi_V$} (n-2-2);
  \draw[-{Stealth[scale=1.1]}] (n-2-1) -- node[below] {$H(t)$} (n-2-2);
\end{tikzpicture}
\]
Aquests quadrats diuen que $\varphi_E$ i $\varphi_V$ respecten l'origen i el destí de cada aresta; és a dir, són precisament la definició de \textbf{morfisme de grafs}. Per tant els morfismes de $\cat{Set}^{\cat{P}}$ es corresponen amb els morfismes de grafs. Aquesta correspondència, sobre objectes i sobre morfismes, és un parell de functors inversos l'un de l'altre, i ho és \emph{estrictament}: les seves composicions són les identitats \emph{exactes} ---no només isomorfismes naturals---, perquè un functor $\cat{P}\to\cat{Set}$ no és res més que els seus valors sobre $E$, $V$, $s$ i $t$; anar del functor al graf i tornar dona el mateix functor, i del graf al functor i tornar, el mateix graf. És, doncs, un isomorfisme de categories en el sentit estricte de la definició~\ref{def:iso-categories} ---més fort que l'equivalència que veurem tot seguit---, i $\cat{Set}^{\cat{P}}\cong\cat{Graph}$. El caràcter estricte d'aquest isomorfisme descansa sobre la codificació que hem triat: en identificar un graf amb les dades $(E,V,s,t)$ ---exactament les que porta un functor $\cat{P}\to\cat{Set}$---, l'anada i la tornada retornen les \emph{mateixes} dades. Amb altres codificacions habituals dels grafs, la correspondència continua sent un isomorfisme \emph{canònic} de categories, però no una igualtat literal de dades; en cap cas \enquote{exactes} no s'ha de llegir com una igualtat fundacional independent de la codificació.
\end{prova}

\begin{observacio}
Aquest exemple té una lectura de fons: un cop decidim representar els grafs com a functors $\cat{P}\to\cat{Set}$, la representació \emph{justifica i recupera canònicament} la definició habitual de morfisme de grafs, perquè els morfismes de $\cat{Set}^{\cat{P}}$ són, per força, les transformacions naturals. És un criteri de coherència que sovint es fa servir en teoria de categories: quan una estructura s'expressa com a categoria de functors, la noció de morfisme que se'n dedueix hauria de coincidir amb la que s'havia adoptat independentment. Les categories del tipus $\cat{Set}^{\cat{C}}$, que codifiquen estructures com els grafs, inclouen com a cas particular les categories de prefeixos, que presentem tot seguit i que seran fonamentals per al lema de Yoneda i, més endavant, per a la semàntica categòrica.
\end{observacio}

\begin{exemple}[La categoria de prefeixos]\label{ex:categoria-prefeixos}\index{prefeix@prefeix!categoria de}\index{prefeix@prefeix!morfisme de}
Sigui $\cat{C}$ una categoria petita. Els prefeixos sobre $\cat{C}$ (definició~\ref{def:prefeix}) són els objectes de la categoria de functors
\[
\widehat{\cat{C}}\coloneqq\cat{Set}^{\cat{C}\op},
\]
que anomenem \textbf{categoria de prefeixos} sobre $\cat{C}$. Els seus morfismes, els \textbf{morfismes de prefeixos} $\sigma\colon P\Rightarrow Q$, són les transformacions naturals: famílies d'aplicacions $\sigma_A\colon P(A)\to Q(A)$ tals que, per a cada morfisme $f\colon A\to B$ de $\cat{C}$,
\[
\sigma_A\comp P(f)=Q(f)\comp\sigma_B\colon P(B)\to Q(A),
\]
que és el quadrat de naturalitat escrit per a prefeixos. Com que $\cat{Set}$ és localment petita, $\widehat{\cat{C}}$ també ho és, per l'observació sobre qüestions de mida. La categoria $\cat{Set}^{\cat{P}}$ dels grafs n'és un exemple: com que un functor $\cat{P}\to\cat{Set}$ és un functor $(\cat{P}\op)\op\to\cat{Set}$, els grafs són els prefeixos sobre $\cat{P}\op$.
\end{exemple}

\begin{observacio}[La forma dels grafs és, ella mateixa, un graf]\index{categoria!lliure}
La categoria índex $\cat{P}$ té una lectura circular ben curiosa. El seu \emph{graf subjacent} ---els dos objectes com a vèrtexs i les dues fletxes generadores $s,t$ com a arestes--- és exactament $\bullet\rightrightarrows\bullet$, i $\cat{P}$ n'és la \textbf{categoria lliure} (secció~\ref{sec:grafs-lliures}): com que les dues fletxes paral·leles no es poden encadenar ---de l'objecte d'arribada no en surt cap---, no hi ha camins de longitud $\geq 2$, i els únics morfismes de $\cat{P}$ són $\id$, $\id$, $s$ i $t$. És a dir, $\cat{P}=\Free(\bullet\rightrightarrows\bullet)$. Així, la \enquote{forma} amb què definim els grafs ---el graf subjacent de $\cat{P}$--- és, ella mateixa, un graf, de fet un objecte de $\cat{Graph}=\cat{Set}^{\cat{P}}$: el patró es tanca sobre si mateix.
\end{observacio}

Ara que les transformacions naturals formen una categoria $\cat{D}^{\cat{C}}$, en podem estudiar els isomorfismes.

\begin{definicio}\index{isomorfisme natural}
Una transformació natural $\eta\colon F\Rightarrow G$ és un \textbf{isomorfisme natural} si admet una \emph{inversa}: una transformació natural $\eta^{-1}\colon G\Rightarrow F$ amb $\eta^{-1}\comp\eta=\id_F$ i $\eta\comp\eta^{-1}=\id_G$ (composició vertical). En aquest cas escrivim $F\cong G$ i diem que $F$ i $G$ són \textbf{naturalment isomorfs}.
\end{definicio}

Quan $\cat{C}$ és petita, aquesta condició és, literalment, la d'un isomorfisme a la categoria de functors $\cat{D}^{\cat{C}}$. En efecte, els morfismes de $\cat{D}^{\cat{C}}$ són les transformacions naturals i la seva composició és la vertical; per tant, que $\eta$ admeti una inversa $\eta^{-1}$ com a transformació natural ---amb $\eta^{-1}\comp\eta=\id_F$ i $\eta\comp\eta^{-1}=\id_G$--- no és res més que tenir un invers a $\cat{D}^{\cat{C}}$. Els isomorfismes naturals coincideixen, doncs, amb els isomorfismes de $\cat{D}^{\cat{C}}$: la noció d'isomorfisme del primer capítol (definició~\ref{def:isomorfisme}), un nivell més amunt.

El fet clau és que aquesta condició \enquote{global} ---tenir inversa com a transformació natural--- es pot comprovar \enquote{punt a punt}, sobre les components:

\begin{proposicio}[Caracterització]\label{prop:iso-natural-components}
Siguin $F,G\colon\cat{C}\to\cat{D}$ dos functors. Una transformació natural $\eta\colon F\Rightarrow G$ és un isomorfisme natural si, i només si, cada component $\eta_A$ (per a $A$ objecte de $\cat{C}$) és un isomorfisme a $\cat{D}$.
\end{proposicio}

\begin{prova}
Si $\eta$ té inversa $\eta^{-1}$, avaluant les igualtats $\eta^{-1}\comp\eta=\id_F$ i $\eta\comp\eta^{-1}=\id_G$ a cada objecte $A$ ---i recordant que la composició vertical és component a component--- obtenim $(\eta^{-1})_A\comp\eta_A=\id_{F(A)}$ i $\eta_A\comp(\eta^{-1})_A=\id_{G(A)}$; per tant $\eta_A$ és un isomorfisme a $\cat{D}$, amb invers $(\eta^{-1})_A$.

Recíprocament, suposem que cada $\eta_A$ és un isomorfisme a $\cat{D}$, amb invers $\eta_A^{-1}$. N'hi ha prou de veure que aquests inversos són les components d'una transformació natural $\eta^{-1}\colon G\Rightarrow F$, perquè aleshores, component a component, és la inversa de $\eta$. Partim del quadrat de naturalitat de $\eta$ per a un morfisme $f\colon A\to B$,
\[
\begin{tikzpicture}[baseline=(m.center)]
  \matrix (m) [matrix of math nodes, row sep=2.6em, column sep=3.4em] {
    F(A) & F(B) \\
    G(A) & G(B) \\
  };
  \draw[-{Stealth[scale=1.1]}] (m-1-1) -- node[above] {$F(f)$} (m-1-2);
  \draw[-{Stealth[scale=1.1]}] (m-1-1) -- node[left] {$\eta_A$} (m-2-1);
  \draw[-{Stealth[scale=1.1]}] (m-1-2) -- node[right] {$\eta_B$} (m-2-2);
  \draw[-{Stealth[scale=1.1]}] (m-2-1) -- node[below] {$G(f)$} (m-2-2);
\end{tikzpicture}
\]
és a dir, $G(f)\comp\eta_A=\eta_B\comp F(f)$, i component per l'esquerra amb $\eta_B^{-1}$ i per la dreta amb $\eta_A^{-1}$,
\[
\eta_B^{-1}\comp G(f)=\eta_B^{-1}\comp G(f)\comp\eta_A\comp\eta_A^{-1}=\eta_B^{-1}\comp\eta_B\comp F(f)\comp\eta_A^{-1}=F(f)\comp\eta_A^{-1},
\]
que és la condició de naturalitat per a $\eta^{-1}$.
\end{prova}

\begin{observacio}\label{obs:iso-no-natural}
La distinció entre isomorfisme i isomorfisme \emph{natural} és fonamental. Que $F(A)$ i $G(A)$ siguin isomorfs per a cada $A$ no basta: cal que els isomorfismes ---un per objecte, les components $\eta_A$--- es puguin triar \emph{alhora} de manera compatible amb tots els morfismes. Aquesta compatibilitat no es refereix a triar més components, sinó a la \emph{naturalitat}: per a cada morfisme $f\colon A\to B$ de $\cat{C}$, el quadrat de $\eta$ ha de commutar ($G(f)\comp\eta_A=\eta_B\comp F(f)$), lligant les components dels dos extrems.

Un exemple mínim ho mostra sense cap tecnicisme. Prenem un grup de dos elements, $\{e,g\}$ amb $g^2=e$ (on $e$ és la identitat), vist com a categoria d'un sol objecte $\cat{B}(\{e,g\})$, i dos functors $F,G\colon\cat{B}(\{e,g\})\to\cat{Set}$. Tots dos envien l'únic objecte $\ast$ al mateix conjunt $\{0,1\}$, però l'element $g$ actua de manera diferent: en $F$ actua per la identitat, i en $G$ per la transposició $\tau$ que intercanvia $0$ i $1$. Puntualment no hi ha res a distingir: $F(\ast)=G(\ast)=\{0,1\}$ són el mateix conjunt, i per tant \emph{isomorfs}. Tanmateix, no hi ha cap isomorfisme natural $F\cong G$. Una transformació natural $\eta\colon F\Rightarrow G$ tindria una única component $\eta_\ast\colon\{0,1\}\to\{0,1\}$, i la naturalitat respecte del morfisme $g$ exigiria
\[
G(g)\comp\eta_\ast=\eta_\ast\comp F(g),
\qquad\text{és a dir}\qquad
\tau\comp\eta_\ast=\eta_\ast.
\]
Però $\tau$ no té punts fixos, de manera que cap aplicació $\eta_\ast$ pot complir $\tau\comp\eta_\ast=\eta_\ast$ (aplicant a qualsevol $x$ obtindríem $\tau(\eta_\ast(x))=\eta_\ast(x)$, impossible). No hi ha, doncs, cap transformació natural $F\Rightarrow G$, i menys encara un isomorfisme natural: la igualtat puntual dels objectes no es pot fer compatible amb l'acció del morfisme. Val a dir que aquí obtenim un contraexemple més fort que el necessari: per descartar un isomorfisme \emph{natural} n'hi hauria prou de veure que cap component \emph{bijectiva} compleix la naturalitat, mentre que l'argument anterior descarta \emph{totes} les components, bijectives o no. La manca d'isomorfisme natural en surt, doncs, reforçada.

El cas clàssic del dual d'un espai vectorial il·lustra la mateixa idea, però demana una precaució addicional. Sovint es diu que un espai de dimensió finita \enquote{és isomorf al seu dual $V\cong V^*$, però no naturalment}. Cal formular-ho amb cura, perquè la identitat és un functor \emph{covariant} $\cat{Vect}_{\mathbb K}\to\cat{Vect}_{\mathbb K}$ mentre que el dual $(-)^*$ és \emph{contravariant}, $\cat{Vect}_{\mathbb K}\op\to\cat{Vect}_{\mathbb K}$: no són dos functors paral·lels als quals es pugui aplicar directament la definició de transformació natural d'aquest capítol. La manera correcta de precisar-ho és restringir-se als \emph{isomorfismes} lineals i comparar la identitat amb l'acció covariant $T\mapsto(T^{-1})^*$; aleshores es comprova que cap família d'isomorfismes $V\cong V^*$ és compatible amb tots els isomorfismes lineals. Aquest desenvolupament es fa amb detall a l'exercici~\ref{ex:dual-no-natural} de la part de pràctica. En canvi, l'isomorfisme $V\cong V^{**}$ amb el bidual sí que està ben plantejat sobre \emph{totes} les aplicacions lineals ---identitat i bidual són tots dos covariants--- i és natural.
\end{observacio}

\begin{exemple}[Les identificacions canòniques del capítol anterior]\label{ex:isos-naturals-cap2}\index{isomorfisme natural!functor de parts}
Al capítol~\ref{cap:functors} vam trobar diverses identificacions canòniques entre un functor i un homfunctor, i vam anunciar que les reconeixeríem com a isomorfismes naturals. La proposició~\ref{prop:iso-natural-components} ho fa immediat: en cada cas, les components són bijeccions (o isomorfismes lineals), i només cal comprovar la naturalitat.

\emph{El functor de parts.} Els dos functors contravariants de l'exemple~\ref{ex:functor-parts} són
\[
\mathcal P,\ \Hom_{\cat{Set}}(-,\{0,1\})\colon\cat{Set}\op\longrightarrow\cat{Set}.
\] Les funcions característiques donen, per a cada conjunt $X$, la bijecció
\[
\chi_X\colon\mathcal P(X)\to\Hom_{\cat{Set}}(X,\{0,1\}),\qquad B\mapsto\chi_B.
\]
Un morfisme de $\cat{Set}\op$ de $Y$ a $X$ és una aplicació $f\colon X\to Y$, que $\mathcal P$ envia a l'antiimatge $f^{-1}$ i l'homfunctor a la precomposició $-\comp f$. El quadrat de naturalitat és
\[
\begin{tikzpicture}[baseline=(m.center)]
  \matrix (m) [matrix of math nodes, row sep=2.8em, column sep=4.2em] {
    \mathcal P(Y) & \mathcal P(X) \\
    \Hom(Y,\{0,1\}) & \Hom(X,\{0,1\}) \\
  };
  \draw[-{Stealth[scale=1.1]}] (m-1-1) -- node[above] {$f^{-1}$} (m-1-2);
  \draw[-{Stealth[scale=1.1]}] (m-1-1) -- node[left] {$\chi_Y$} (m-2-1);
  \draw[-{Stealth[scale=1.1]}] (m-1-2) -- node[right] {$\chi_X$} (m-2-2);
  \draw[-{Stealth[scale=1.1]}] (m-2-1) -- node[below] {$-\comp f$} (m-2-2);
\end{tikzpicture}
\qquad
\chi_B\comp f=\chi_{f^{-1}(B)}\quad(B\subseteq Y),
\]
i commuta, perquè per a tot $x\in X$ es té $(\chi_B\comp f)(x)=1$ si i només si $f(x)\in B$, és a dir, si i només si $x\in f^{-1}(B)$. Per tant $\chi\colon\mathcal P\Rightarrow\Hom_{\cat{Set}}(-,\{0,1\})$ és un isomorfisme natural, i ara podem escriure amb tota precisió
\[
\mathcal P\;\cong\;\Hom_{\cat{Set}}(-,\{0,1\}).
\]

\emph{Els altres casos.} El mateix esquema val per a les altres identificacions del capítol anterior; n'indiquem les components i l'equació de naturalitat, que es comprova directament.
\begin{itemize}
  \item A $\cat{Set}$, $\id_{\cat{Set}}\cong\Hom_{\cat{Set}}(\{\ast\},-)$: la component envia $x\in X$ a l'aplicació $\bar x\colon\ast\mapsto x$, i la naturalitat diu que $f\comp\bar x=\overline{f(x)}$.
  \item A $\cat{Graph}$, el functor de vèrtexs és naturalment isomorf a $\Hom_{\cat{Graph}}(\mathbf V,-)$ i el d'arestes a $\Hom_{\cat{Graph}}(\mathbf E,-)$ (exemple~\ref{ex:homfunctor-grafs}): la component envia cada vèrtex, o cada aresta, al morfisme que el tria, i la naturalitat diu que compondre aquest morfisme amb un morfisme de grafs $F$ tria la imatge per $F$.
  \item A $\cat{Vect}_{\mathbb K}$, el functor identitat és naturalment isomorf a l'aixecament
\[
\Hom_{\mathbb K}(\mathbb K,-)\colon\cat{Vect}_{\mathbb K}\longrightarrow\cat{Vect}_{\mathbb K}.
\]
La component envia $w\in W$ a l'aplicació lineal $\lambda\mapsto\lambda w$, amb inversa $\varphi\mapsto\varphi(1)$, i la naturalitat diu que $T\comp(\lambda\mapsto\lambda w)=(\lambda\mapsto\lambda\,T(w))$ per a tota aplicació lineal $T$.
\end{itemize}
Aquests exemples tornaran a aparèixer al capítol del lema de Yoneda, on direm que aquests functors són \emph{representables}.
\end{exemple}

\subsection{Equivalència de categories}\label{sec:equivalencia}

Quan volem dir que dues categories són \enquote{essencialment la mateixa} ---amb la mateixa estructura categòrica, distingint els objectes només llevat d'isomorfisme---, la igualtat estricta és massa rígida, i l'isomorfisme de categories, que demana functors inversos ($G\comp F=\id_{\cat{C}}$ i $F\comp G=\id_{\cat{D}}$), també ho és sovint. La noció adequada relaxa aquestes igualtats a isomorfismes naturals: en lloc d'un invers demanarem un \textbf{quasiinvers}, un functor $G$ amb $G\comp F\cong\id_{\cat{C}}$ i $F\comp G\cong\id_{\cat{D}}$.

\begin{definicio}\index{equivalència de categories}
\label{def:equivalencia}
Una \textbf{equivalència de categories} entre $\cat{C}$ i $\cat{D}$ és una parella de functors
\[
F\colon\cat{C}\to\cat{D},
\qquad
G\colon\cat{D}\to\cat{C},
\]
juntament amb isomorfismes naturals
\[
G\comp F\cong\id_{\cat{C}}
\qquad\text{i}\qquad
F\comp G\cong\id_{\cat{D}}.
\]
Quan existeix una tal equivalència, diem que $\cat{C}$ i $\cat{D}$ són \textbf{equivalents} i escrivim $\cat{C}\simeq\cat{D}$. El functor $G$ s'anomena un \textbf{quasiinvers} de $F$. Quan calgui donar nom als isomorfismes naturals, els orientarem sempre així:
\[
\eta\colon\id_{\cat{C}}\Rightarrow G\comp F,
\qquad
\varepsilon\colon F\comp G\Rightarrow\id_{\cat{D}}.
\]
\end{definicio}

\begin{observacio}
La diferència amb l'isomorfisme de categories és que aquí no demanem $G\comp F=\id_{\cat{C}}$, sinó només $G\comp F\cong\id_{\cat{C}}$: en tornar de $\cat{C}$ a $\cat{C}$ no recuperem exactament cada objecte, sinó un objecte \emph{naturalment isomorf}. Aquesta flexibilitat és la que fa que l'equivalència sigui la noció adequada de comparació entre categories quan, com passa sovint, només ens interessen els objectes llevat d'isomorfisme.
\end{observacio}

El resultat central caracteritza les equivalències mitjançant les tres propietats dels functors que vam introduir al capítol anterior: ser ple i fidel (la definició~\ref{def:ple-fidel}) i l'exhaustivitat essencial (la definició~\ref{def:ess-exhaustiu}).

\begin{teorema}\index{equivalència de categories!caracterització}
\label{teo:caract-equivalencia}
Un functor $F\colon\cat{C}\to\cat{D}$ forma part d'una equivalència de categories si i només si és \textbf{ple}, \textbf{fidel} i \textbf{essencialment exhaustiu}.
\end{teorema}

\begin{prova}
\emph{($\Rightarrow$)} Suposem que $F$ forma part d'una equivalència, amb quasiinvers $G\colon\cat{D}\to\cat{C}$ i isomorfismes naturals $\eta\colon\id_{\cat{C}}\Rightarrow G\comp F$ i $\varepsilon\colon F\comp G\Rightarrow\id_{\cat{D}}$. Com que les components de $\eta$ i de $\varepsilon$ són invertibles, els quadrats de naturalitat es poden llegir com a fórmules: per a $f\colon A\to B$ a $\cat{C}$ i $g\colon D\to D'$ a $\cat{D}$,
\[
G(F(f))=\eta_B\comp f\comp\eta_A^{-1},
\qquad\qquad
F(G(g))=\varepsilon_{D'}^{-1}\comp g\comp\varepsilon_D.
\tag{$\ast$}
\]

\emph{$F$ és fidel.} Siguin $f,f'\colon A\to B$ amb $F(f)=F(f')$. Aleshores $G(F(f))=G(F(f'))$ i, per~$(\ast)$, $\eta_B\comp f\comp\eta_A^{-1}=\eta_B\comp f'\comp\eta_A^{-1}$. Component amb $\eta_B^{-1}$ per l'esquerra i amb $\eta_A$ per la dreta, $f=f'$.

\emph{$G$ és fidel.} És el mateix argument amb la segona fórmula de~$(\ast)$: si $G(g)=G(g')$ per a $g,g'\colon D\to D'$, aleshores $\varepsilon_{D'}^{-1}\comp g\comp\varepsilon_D=\varepsilon_{D'}^{-1}\comp g'\comp\varepsilon_D$, i per tant $g=g'$.

\emph{$F$ és ple.} Sigui $u\colon F(A)\to F(B)$ un morfisme qualsevol de $\cat{D}$, i posem
\[
f=\eta_B^{-1}\comp G(u)\comp\eta_A\colon A\to B.
\]
Per~$(\ast)$, $G(F(f))=\eta_B\comp\bigl(\eta_B^{-1}\comp G(u)\comp\eta_A\bigr)\comp\eta_A^{-1}=G(u)$, i com que $G$ és fidel, $F(f)=u$. Per concloure $F(f)=u$ no es pot fer servir que $F$ és fidel, perquè encara no sabem que $u$ sigui la imatge de cap morfisme; és que $G$ sigui fidel el que ho permet.

\emph{$F$ és essencialment exhaustiu.} Per a cada objecte $D$ de $\cat{D}$, la component $\varepsilon_D\colon F(G(D))\to D$ és un isomorfisme, de manera que $D$ és isomorf a l'objecte $F(G(D))$ de la imatge de $F$.

\emph{($\Leftarrow$)} Suposem que $F$ és ple, fidel i essencialment exhaustiu. Construïm un quasiinvers en quatre passos.

\emph{Pas 1: $G$ sobre objectes.} Per l'exhaustivitat essencial, per a cada objecte $D$ de $\cat{D}$ hi ha algun objecte $C$ de $\cat{C}$ amb $F(C)\cong D$. Triem, per a cada $D$, un tal objecte, que anomenem $G(D)$, i un isomorfisme $\varepsilon_D\colon F(G(D))\to D$.

\emph{Pas 2: $G$ sobre morfismes.} Sigui $g\colon D\to D'$. El morfisme $\varepsilon_{D'}^{-1}\comp g\comp\varepsilon_D$ va de $F(G(D))$ a $F(G(D'))$. Com que $F$ és ple, prové d'algun morfisme $G(D)\to G(D')$, i com que $F$ és fidel, d'un de sol. Definim $G(g)$ com aquest únic morfisme; és a dir, $G(g)$ queda caracteritzat per
\[
F(G(g))=\varepsilon_{D'}^{-1}\comp g\comp\varepsilon_D.
\tag{$\ast\ast$}
\]

\emph{Pas 3: $G$ és un functor.} Per a la identitat, $(\ast\ast)$ dona $F(G(\id_D))=\varepsilon_D^{-1}\comp\varepsilon_D=\id_{F(G(D))}=F(\id_{G(D)})$, i, com que $F$ és fidel, $G(\id_D)=\id_{G(D)}$. Per a la composició, siguin $g\colon D\to D'$ i $g'\colon D'\to D''$. Aplicant $(\ast\ast)$ dues vegades,
\[
\begin{aligned}
F\bigl(G(g')\comp G(g)\bigr)&=F(G(g'))\comp F(G(g))
=\bigl(\varepsilon_{D''}^{-1}\comp g'\comp\varepsilon_{D'}\bigr)\comp\bigl(\varepsilon_{D'}^{-1}\comp g\comp\varepsilon_D\bigr)\\
&=\varepsilon_{D''}^{-1}\comp(g'\comp g)\comp\varepsilon_D=F\bigl(G(g'\comp g)\bigr),
\end{aligned}
\]
i, com que $F$ és fidel, $G(g')\comp G(g)=G(g'\comp g)$.

\emph{Pas 4: els isomorfismes naturals.} Component $(\ast\ast)$ amb $\varepsilon_{D'}$ per l'esquerra s'obté
\[
\varepsilon_{D'}\comp F(G(g))=g\comp\varepsilon_D,
\]
que és el quadrat de naturalitat de $\varepsilon\colon F\comp G\Rightarrow\id_{\cat{D}}$ per a $g$. Com que cada $\varepsilon_D$ és un isomorfisme, $\varepsilon$ és un isomorfisme natural.

Per construir $\eta$, sigui $A$ un objecte de $\cat{C}$. Com que $F$ és ple i fidel, hi ha un únic $\eta_A\colon A\to G(F(A))$ amb $F(\eta_A)=\varepsilon_{F(A)}^{-1}$, i un únic $\zeta_A\colon G(F(A))\to A$ amb $F(\zeta_A)=\varepsilon_{F(A)}$. Aleshores
\[
F(\zeta_A\comp\eta_A)=\varepsilon_{F(A)}\comp\varepsilon_{F(A)}^{-1}=F(\id_A),
\qquad
F(\eta_A\comp\zeta_A)=\varepsilon_{F(A)}^{-1}\comp\varepsilon_{F(A)}=F(\id_{G(F(A))}),
\]
i, com que $F$ és fidel, $\zeta_A$ és l'invers de $\eta_A$: cada $\eta_A$ és un isomorfisme. Queda la naturalitat: per a $f\colon A\to B$, cal veure $\eta_B\comp f=G(F(f))\comp\eta_A$. Apliquem $F$ als dos membres. D'una banda, $F(\eta_B\comp f)=\varepsilon_{F(B)}^{-1}\comp F(f)$. De l'altra, aplicant $(\ast\ast)$ al morfisme $g=F(f)$,
\[
F\bigl(G(F(f))\comp\eta_A\bigr)=\bigl(\varepsilon_{F(B)}^{-1}\comp F(f)\comp\varepsilon_{F(A)}\bigr)\comp\varepsilon_{F(A)}^{-1}=\varepsilon_{F(B)}^{-1}\comp F(f).
\]
Les dues imatges coincideixen, i, com que $F$ és fidel, $\eta_B\comp f=G(F(f))\comp\eta_A$. Així $\eta\colon\id_{\cat{C}}\Rightarrow G\comp F$ és un isomorfisme natural, i $G$ és un quasiinvers de $F$.
\end{prova}

\begin{observacio}[Sobre l'elecció]
A la direcció ($\Leftarrow$), el pas~1 tria simultàniament un objecte $G(D)$ i un isomorfisme $\varepsilon_D$ per a cada $D$, i per tant fa servir l'\emph{axioma de l'elecció}.\footnotemark{} Tota la resta de la construcció és \emph{determinada} per aquesta tria: el pas~2 no tria res, perquè, com que $F$ és ple i fidel, $G(g)$ és únic, i $\eta$ també queda fixat. En casos concrets sovint es pot donar un quasiinvers canònic sense recórrer a l'elecció. La direcció ($\Rightarrow$), en canvi, no tria res, i tampoc no fa servir cap relació entre $\eta$ i $\varepsilon$ més enllà de la naturalitat i la invertibilitat.
\end{observacio}
\footnotetext{Si $\cat{D}$ és gran, la família de tries està indexada per una classe pròpia, i cal una forma d'elecció més forta que l'habitual, l'\emph{elecció global}, disponible per exemple a la teoria de classes NBG (apèndix~\ref{cap:mida}). Vegeu \cite[§1.5]{riehl}.}%

\begin{definicio}[Esquelet]\index{categoria!esquelètica}\index{esquelet}\label{def:esquelet}
Una categoria és \textbf{esquelètica} si dos objectes isomorfs qualssevol són \emph{iguals}; és a dir, si cada classe d'isomorfisme té un únic objecte. Un \textbf{esquelet} d'una categoria $\cat{C}$ és una subcategoria plena que conté \emph{exactament} un objecte de cada classe d'isomorfisme de $\cat{C}$; un esquelet és sempre una categoria esquelètica.
\end{definicio}

\begin{observacio}
Convé distingir dues nocions properes. Una subcategoria plena que conté \emph{almenys} un objecte de cada classe d'isomorfisme ja és equivalent a $\cat{C}$ (com veurem tot seguit), però pot tenir diversos representants isomorfs d'una mateixa classe. Un \emph{esquelet} n'és el cas ajustat: en conté \emph{exactament} un, de manera que no queden objectes isomorfs distints. Triar exactament un representant de cada classe d'isomorfisme és, però, una elecció simultània. Amb l'\emph{axioma de l'elecció}, tota categoria \emph{petita} admet un esquelet, i tots els seus esquelets són equivalents entre si. Per a una categoria gran, seleccionar simultàniament un representant de cada classe d'isomorfisme pot requerir un principi d'elecció global, segons el marc fundacional adoptat.\footnotemark És la mateixa elecció que fa servir la construcció d'un quasiinvers (observació anterior).
\end{observacio}
\footnotetext{En una categoria gran, les classes d'isomorfisme formen una classe pròpia, i cal la mateixa forma d'elecció global que a l'observació anterior (apèndix~\ref{cap:mida}); la petitesa local només controla els homconjunts i no resol aquesta selecció.}%

\begin{exemple}
Tota categoria $\cat{C}$ és equivalent a qualsevol subcategoria plena que contingui almenys un objecte de cada classe d'isomorfisme de $\cat{C}$. En efecte, el functor inclusió d'aquesta subcategoria és ple i fidel per ser una subcategoria plena, i essencialment exhaustiu per contenir un representant de cada classe. Per exemple, la categoria $\cat{FinSet}$ de tots els conjunts finits és equivalent a la seva subcategoria plena formada pels conjunts $\underline{0}=\varnothing$ i $\underline{n}=\{1,\dots,n\}$ per a $n\geq 1$ (exemple~\ref{ex:ess-exhaustiu-finset}): no cal treballar amb tots els conjunts finits, sinó només amb un de cada mida. Aquesta subcategoria $\{\underline{n}\mid n\geq 0\}$ és, de fet, un \emph{esquelet} de $\cat{FinSet}$, perquè conté exactament un conjunt de cada cardinalitat finita. Convé notar que $\cat{FinSet}$ és una categoria \emph{gran} ---els seus objectes són massa nombrosos per formar un conjunt, tal com passa amb $\cat{Set}$--- mentre que l'esquelet $\{\underline{n}\mid n\geq 0\}$ és \emph{petit}, ja que té un conjunt (de fet, numerable) d'objectes. Una categoria s'anomena \textbf{essencialment petita}\index{categoria!essencialment petita} si és equivalent a una categoria petita. L'equivalència $\cat{FinSet}\simeq\{\underline{n}\mid n\geq 0\}$ mostra, doncs, que $\cat{FinSet}$ és essencialment petita.
\end{exemple}

\begin{observacio}
Aquest exemple mostra la utilitat pràctica de l'equivalència: permet substituir una categoria gran i redundant per una de més petita i manejable, sense perdre cap informació categòrica. És el reflex, a nivell de categories, del fet que en teoria de categories els objectes només importen llevat d'isomorfisme.
\end{observacio}

\begin{exemple}[Matrius i espais vectorials: equivalència sense igualtat ni isomorfisme]\label{ex:mat-finvect}\index{equivalència de categories!matrius i espais vectorials}\index{Mat@$\cat{Mat}_{\mathbb K}$}
Fixem un cos $\mathbb K$. La categoria $\cat{Mat}_{\mathbb K}$ té per objectes els nombres naturals $n\geq 0$. Els morfismes $m\to n$ són les matrius $n\times m$ amb entrades a $\mathbb K$, la composició és el producte, $B\comp A=BA$, i la identitat de $n$ és la matriu identitat $I_n$; els axiomes són l'associativitat del producte de matrius i la neutralitat de $I_n$. Sigui $\cat{FinVect}_{\mathbb K}$ la subcategoria plena de $\cat{Vect}_{\mathbb K}$ formada pels espais de dimensió finita. El functor
\[
F\colon\cat{Mat}_{\mathbb K}\to\cat{FinVect}_{\mathbb K},\qquad n\mapsto\mathbb K^n,\qquad A\mapsto(x\mapsto Ax),
\]
té els tipus correctes: una matriu $n\times m$ dona una aplicació lineal $\mathbb K^m\to\mathbb K^n$. A més, $F(BA)=F(B)\comp F(A)$ i $F(I_n)=\id_{\mathbb K^n}$. És fidel i ple, perquè tota aplicació lineal $\mathbb K^m\to\mathbb K^n$ és la multiplicació per una única matriu, la que té per columnes les imatges de la base canònica. És essencialment exhaustiu, perquè tot espai de dimensió $n$ és isomorf a $\mathbb K^n$. Pel teorema~\ref{teo:caract-equivalencia}, $F$ és una equivalència: $\cat{Mat}_{\mathbb K}\simeq\cat{FinVect}_{\mathbb K}$. Aquest exemple permet veure els tres nivells de comparació.
\begin{itemize}
\item \emph{No és una igualtat.} Són dues categories diferents: els objectes de l'una són nombres i els de l'altra, espais vectorials.
\item \emph{No és un isomorfisme de categories.} A $\cat{Mat}_{\mathbb K}$, dos objectes isomorfs són iguals, perquè una matriu invertible és quadrada: la categoria és esquelètica. A $\cat{FinVect}_{\mathbb K}$, en canvi, $\mathbb K^2$ i l'espai dels polinomis de grau $\leq 1$ són objectes diferents i isomorfs. Un isomorfisme de categories és bijectiu sobre objectes i preserva i reflecteix els isomorfismes, de manera que conservaria la propietat de ser esquelètica. Per tant, no hi ha cap isomorfisme de categories entre $\cat{Mat}_{\mathbb K}$ i $\cat{FinVect}_{\mathbb K}$, ni aquest $F$ ni cap altre.
\item \emph{El quasiinvers només torna llevat d'isomorfisme.} El pas~1 de la demostració del teorema consisteix aquí a \emph{triar una base} de cada espai $V$, és a dir, un isomorfisme $\varepsilon_V\colon\mathbb K^{\dim V}\to V$. Aleshores $G(V)=\dim V$, i $G(g)$ és la matriu de $g$ en les bases triades. Si per als espais $\mathbb K^n$ triem la base canònica, es té literalment $G\comp F=\id_{\cat{Mat}_{\mathbb K}}$. En canvi, $F(G(V))=\mathbb K^{\dim V}$, que no és $V$ tret que $V$ sigui algun $\mathbb K^n$. Així, $F\comp G$ \emph{no} és igual a $\id_{\cat{FinVect}_{\mathbb K}}$, sinó només isomorf a través de $\varepsilon$. La condició de naturalitat $g\comp\varepsilon_V=\varepsilon_W\comp F(G(g))$ és la coneguda fórmula que relaciona una aplicació lineal amb la seva matriu en unes bases.
\end{itemize}
Un espai vectorial qualsevol no té cap base privilegiada; per això el quasiinvers depèn d'una tria, i el que s'obté és una equivalència i no un isomorfisme. Per a l'àlgebra lineal, la conclusió és la que ja es fa servir a la pràctica: treballar amb matrius no perd cap informació sobre les aplicacions lineals entre espais de dimensió finita, sempre que es recordi que la identificació depèn de les bases.
\end{exemple}

Quines propietats es conserven per equivalència i quines no, i com es demostra que dues categories \emph{no} són equivalents, es tracta a l'apèndix~\refalt{cap:contraexemples}{de contraexemples del llibre complet}, junt amb la relació entre equivalència i esquelets.


\section{Composició horitzontal}\index{transformació natural!composició horitzontal}

La composició vertical apila transformacions entre functors \emph{amb el mateix domini i codomini}. Hi ha una segona manera de compondre, que combina transformacions naturals amb functors i, més generalment, transformacions naturals \enquote{de costat}.

Comencem pel cas mixt, en què componem una transformació natural amb un functor\footnote{Segons el costat pel qual s'afegeix el functor, es parla de composició \enquote{per l'esquerra} i \enquote{per la dreta}; convé advertir que aquesta assignació d'\enquote{esquerra} i \enquote{dreta} no és uniforme a la bibliografia. Aquí la fixem per la posició del functor en la notació: $K\eta$ (functor a l'esquerra) és la composició per l'esquerra, i $\eta H$ (functor a la dreta), la composició per la dreta.}.

\begin{definicio}[Composició amb un functor]\index{transformació natural!composició amb un functor}
Sigui $\eta\colon F\Rightarrow G$ una transformació natural entre functors $F,G\colon\cat{C}\to\cat{D}$.
\begin{itemize}
\item Si $K\colon\cat{D}\to\cat{E}$ és un functor, la \textbf{composició per l'esquerra} $K\eta\colon K\comp F\Rightarrow K\comp G$ té per components $(K\eta)_A=K(\eta_A)$.
\item Si $H\colon\cat{B}\to\cat{C}$ és un functor, la \textbf{composició per la dreta} $\eta H\colon F\comp H\Rightarrow G\comp H$ té per components $(\eta H)_B=\eta_{H(B)}$.
\end{itemize}
\end{definicio}

\begin{proposicio}
Les dues construccions anteriors són transformacions naturals.
\end{proposicio}

\begin{prova}
Per a $K\eta$, apliquem el functor $K$ al quadrat de naturalitat de $\eta$ per a un morfisme $f\colon A\to B$: de $G(f)\comp\eta_A=\eta_B\comp F(f)$, i com que $K$ preserva la composició,
\[
K(G(f))\comp K(\eta_A)=K(\eta_B)\comp K(F(f)),
\]
que és la condició de naturalitat per a $K\eta$ respecte de $K\comp F$ i $K\comp G$. Per a $\eta H$, la naturalitat respecte d'un morfisme $b\colon B\to B'$ de $\cat{B}$ és el quadrat de naturalitat de $\eta$ per al morfisme $H(b)$ de $\cat{C}$.
\end{prova}

Amb aquestes dues operacions podem definir la composició horitzontal general.

\begin{definicio}[Composició horitzontal]\index{transformació natural!composició horitzontal}
Siguin $\eta\colon F\Rightarrow G$ (amb $F,G\colon\cat{C}\to\cat{D}$) i $\theta\colon K\Rightarrow L$ (amb $K,L\colon\cat{D}\to\cat{E}$). La seva \textbf{composició horitzontal} $\theta\ast\eta\colon K\comp F\Rightarrow L\comp G$ té per components
\[
(\theta\ast\eta)_A=\theta_{G(A)}\comp K(\eta_A)=L(\eta_A)\comp\theta_{F(A)}.
\]
\end{definicio}

Diagramàticament, la composició horitzontal juxtaposa les dues transformacions, l'una a continuació de l'altra:
\[
\begin{tikzpicture}[baseline=(A.base)]
  \node (A) at (0,0) {$\cat{C}$};
  \node (B) at (2.9,0) {$\cat{D}$};
  \node (Cc) at (5.8,0) {$\cat{E}$};
  \draw[cat] (A) to[bend left=26] node[above] {$F$} (B);
  \draw[cat] (A) to[bend right=26] node[below] {$G$} (B);
  \draw[cat nat] (1.45,0.2) -- node[right=1.5pt] {$\eta$} (1.45,-0.2);
  \draw[cat] (B) to[bend left=26] node[above] {$K$} (Cc);
  \draw[cat] (B) to[bend right=26] node[below] {$L$} (Cc);
  \draw[cat nat] (4.35,0.2) -- node[right=1.5pt] {$\theta$} (4.35,-0.2);
\end{tikzpicture}
\]

\begin{observacio}\label{obs:comp-horitzontal}
Les dues expressions de $(\theta\ast\eta)_A$ coincideixen precisament pel quadrat de naturalitat de $\theta$ aplicat al morfisme $\eta_A\colon F(A)\to G(A)$ de $\cat{D}$:
\[
\begin{tikzpicture}[baseline=(m.center)]
  \matrix (m) [matrix of math nodes, row sep=2.8em, column sep=4.2em] {
    K(F(A)) & K(G(A)) \\
    L(F(A)) & L(G(A)) \\
  };
  \draw[-{Stealth[scale=1.1]}] (m-1-1) -- node[above] {$K(\eta_A)$} (m-1-2);
  \draw[-{Stealth[scale=1.1]}] (m-1-1) -- node[left] {$\theta_{F(A)}$} (m-2-1);
  \draw[-{Stealth[scale=1.1]}] (m-1-2) -- node[right] {$\theta_{G(A)}$} (m-2-2);
  \draw[-{Stealth[scale=1.1]}] (m-2-1) -- node[below] {$L(\eta_A)$} (m-2-2);
\end{tikzpicture}
\qquad
\theta_{G(A)}\comp K(\eta_A)=L(\eta_A)\comp\theta_{F(A)}.
\]
La composició horitzontal es pot llegir, doncs, component una transformació natural amb un functor a cada costat, en qualsevol dels dos ordres:
\[
\theta\ast\eta=(\theta G)\comp(K\eta)=(L\eta)\comp(\theta F).
\]
La composició vertical i l'horitzontal, juntes, satisfan una llei d'intercanvi que fa de $\cat{Cat}$ una \emph{2-categoria}: en una graella de transformacions naturals com la següent, compondre primer en vertical i després en horitzontal dona el mateix que fer-ho en l'ordre contrari,
\[
(\beta'\comp\alpha')\ast(\beta\comp\alpha)=(\beta'\ast\beta)\comp(\alpha'\ast\alpha).
\]
\[
\begin{tikzpicture}[baseline=(A.base)]
  \node (A) at (0,0) {$\cat{C}$};
  \node (B) at (3.4,0) {$\cat{D}$};
  \node (Cc) at (6.8,0) {$\cat{E}$};
  \draw[cat] (A) to[bend left=44] node[above] {$F$} (B);
  \draw[cat] (A) -- (B);
  \draw[cat] (A) to[bend right=44] node[below] {$H$} (B);
  \node[fill=white,inner sep=1.5pt] at (1.7,0) {$G$};
  \draw[cat nat] (1.7,0.52) -- node[right=1.5pt] {$\alpha$} (1.7,0.27);
  \draw[cat nat] (1.7,-0.27) -- node[right=1.5pt] {$\beta$} (1.7,-0.52);
  \draw[cat] (B) to[bend left=44] node[above] {$F'$} (Cc);
  \draw[cat] (B) -- (Cc);
  \draw[cat] (B) to[bend right=44] node[below] {$H'$} (Cc);
  \node[fill=white,inner sep=1.5pt] at (5.1,0) {$G'$};
  \draw[cat nat] (5.1,0.52) -- node[right=1.5pt] {$\alpha'$} (5.1,0.27);
  \draw[cat nat] (5.1,-0.27) -- node[right=1.5pt] {$\beta'$} (5.1,-0.52);
\end{tikzpicture}
\]
No en necessitarem els detalls, però és el marc en què les transformacions naturals adquireixen tota la seva potència.
\end{observacio}

\begin{resumcap}
\apartat{Idees centrals}
\begin{enumerate}
  \item Una transformació natural és una família de morfismes compatible amb \emph{totes} les fletxes de la categoria de partida: cada quadrat de naturalitat commuta.
  \item Les transformacions naturals es componen verticalment, component a component; amb aquesta composició, els functors $\cat{C}\to\cat{D}$ i les transformacions naturals formen la categoria de functors $\cat{D}^{\cat{C}}$.
  \item Un isomorfisme natural és una transformació natural amb totes les components invertibles. La naturalitat no s'hi pot ometre: una família d'isomorfismes, un per a cada objecte, que no sigui natural no és un isomorfisme natural.
  \item Una equivalència de categories és un functor amb un quasiinvers, llevat d'isomorfismes naturals. En el marc d'elecció adoptat al llibre (apèndix~\ref{cap:mida}), es caracteritza equivalentment per ser un functor ple, fidel i essencialment exhaustiu (teorema~\ref{teo:caract-equivalencia}, amb demostració completa). No és una igualtat ni, en general, un isomorfisme de categories: $\cat{Mat}_{\mathbb K}\simeq\cat{FinVect}_{\mathbb K}$, però no són isomorfes.
  \item La composició horitzontal és un segon mode de composició de transformacions naturals, compatible amb la vertical a través de la llei d'intercanvi.
\end{enumerate}
\apartat{Exemple clau} Els grafs dirigits són els objectes d'una categoria de functors: $\cat{Graph}\cong\cat{Set}^{\cat{P}}$, on els morfismes de grafs són exactament les transformacions naturals. Convé recordar l'orientació: $\cat{P}$ té dues fletxes $s,t\colon E\to V$, de manera que un graf és un functor \emph{covariant} $\cat{P}\to\cat{Set}$; no s'ha de confondre amb un prefeix sobre $\cat{P}$.
\apartat{Errors típics} Oblidar comprovar la naturalitat (que el quadrat commuti per a \emph{tot} morfisme); confondre isomorfisme natural amb igualtat de functors; creure que una equivalència conserva propietats com \enquote{tenir un sol objecte} o \enquote{ser esquelètica}; confondre una equivalència amb un isomorfisme de categories, o esperar que $F\comp G$ sigui literalment la identitat; i barrejar composició vertical i horitzontal de transformacions.
\apartat{Lectura recomanada} \cite[§§7.4--7.10]{awodey} per a la naturalitat, les categories de functors i les equivalències; \cite[§§1.4--1.5]{riehl}; \cite[§1.3]{leinster}, que inclou també les equivalències; \cite[I.4, II.4--II.5 i IV.4]{maclane} per a la formulació clàssica de la naturalitat, de les categories de functors amb la composició horitzontal i de les equivalències; \cite[§§4.2--4.4]{barrwells} per a la naturalitat amb exemples de grafs i llistes i per al càlcul de Godement, que és la composició horitzontal; \cite[§1.4 i §1.5.4]{perrone-notes} per a exemples de naturalitat, una secció sobre què \emph{no} és natural i les equivalències.
\end{resumcap}

\ifpublicacioparcial\else
\chapter[Representabilitat i Yoneda]{\texorpdfstring{Propietats universals,\\ representabilitat i Yoneda}{Propietats universals, representabilitat i Yoneda}}
\label{cap:yoneda}

\begin{objectius}
\apartat{Prerequisits} Categories sobre i sota un objecte i dualitat (cap.~\ref{cap:categories}); homfunctors, bifunctor $\Hom$, prefeixos i la parella $\Free$, $U$ (cap.~\ref{cap:functors}); transformacions i isomorfismes naturals, i categories de functors, en particular la categoria de prefeixos (cap.~\ref{cap:transformacions}). No es pressuposa cap coneixement previ de representabilitat ni del lema de Yoneda.
\apartat{Objectius} En acabar el capítol, el lector hauria de poder:
\begin{itemize}
  \item definir objectes inicials i terminals i demostrar-ne la unicitat corresponent;
  \item construir i llegir categories coma en exemples senzills, i interpretar-hi dades universals com a objectes inicials o terminals;
  \item definir una representació mitjançant un element universal i traduir-la a un isomorfisme natural amb un homfunctor;
  \item distingir el functor representat, l'objecte representant, l'element universal i la representació completa;
  \item definir la immersió de Yoneda, enunciar i demostrar el lema de Yoneda i seguir-ne explícitament la bijecció;
  \item deduir que la immersió de Yoneda és plena i fidel i que un objecte queda determinat, llevat d'isomorfisme, pel seu functor representable.
\end{itemize}
\end{objectius}

\section{Cap a les propietats universals}

Amb les transformacions naturals vam completar, al capítol~\ref{cap:transformacions}, una torre de tres nivells: objectes i morfismes dins d'una categoria, functors entre categories, i transformacions naturals entre functors ---l'estructura que, amb les composicions vertical i horitzontal, fa de $\cat{Cat}$ una 2-categoria. És el marc en què es formulen les construccions més potents de la teoria, i les primeres construccions que estudiarem són les propietats universals.

Als capítols anteriors han aparegut, sempre com a exemples i sense tractar-los en
detall, objectes i construccions definits no per la seva constitució interna sinó
per la manera com es relacionen amb tots els altres. El producte de dos objectes
$A$ i $B$ (subsecció~\ref{subsec:producte-coproducte}) es defineix per una condició
d'aquest tipus: per a cada objecte $X$ i cada parell de morfismes $f\colon X\to A$
i $g\colon X\to B$ existeix un \emph{únic} morfisme $u\colon X\to A\times B$ amb
$\pi_1\comp u=f$ i $\pi_2\comp u=g$, és a dir, que fa commutar el diagrama
\[
\begin{tikzpicture}[baseline=(m.center)]
  \matrix (m) [matrix of math nodes, row sep=2.6em, column sep=2.8em] {
     & X & \\
    A & A\times B & B \\
  };
  \draw[-{Stealth[scale=1.0]}] (m-1-2) -- node[above left] {$f$} (m-2-1);
  \draw[-{Stealth[scale=1.0]}] (m-1-2) -- node[above right] {$g$} (m-2-3);
  \draw[-{Stealth[scale=1.0]},dashed] (m-1-2) -- node[right] {$u$} (m-2-2);
  \draw[-{Stealth[scale=1.0]}] (m-2-2) -- node[below] {$\pi_1$} (m-2-1);
  \draw[-{Stealth[scale=1.0]}] (m-2-2) -- node[below] {$\pi_2$} (m-2-3);
\end{tikzpicture}
\]
El coproducte es defineix per la condició dual, amb les fletxes girades. La categoria lliure $\Free(G)$ la satisfà també: un functor
$\Free(G)\to\cat{D}$ queda determinat, de manera única, per un morfisme de grafs
$G\to U(\cat{D})$ (subsecció~\ref{subsec:free-u}). I els objectes inicial i
terminal, que fins ara només hem anomenat (secció~\ref{sec:dualitat}), en són el
cas més senzill. Aquestes condicions, que quantifiquen sobre \emph{totes} les
fletxes d'una categoria, s'anomenen \textbf{propietats
universals}\index{propietat universal}, i són, potser, la contribució més
característica de la teoria de categories. La seva virtut, que encara no hem demostrat en cap cas, és que determinen la \emph{dada universal} ---l'objecte juntament amb les fletxes que intervenen en la condició--- llevat d'un únic isomorfisme compatible. En particular, l'objecte subjacent queda determinat llevat d'isomorfisme, però pot admetre altres isomorfismes que no respectin la dada universal. Abans de donar-ne la forma general, les estudiem amb detall en dos
casos: el més simple, el dels objectes terminals i inicials, que definim ara, i el
de la construcció lliure, que ja coneixem.

\subsection{El cas més simple: objectes terminals i inicials}

El cas extrem, i el més fàcil de dibuixar, és el d'un objecte al qual s'arriba d'una
\emph{única} manera des de qualsevol altre, o del qual es \emph{surt} d'una única
manera cap a qualsevol altre. Diem que un objecte $T$ és \textbf{terminal} si per a
cada objecte $X$ hi ha exactament una fletxa $X\to T$, i que un objecte $I$ és
\textbf{inicial} si per a cada $X$ n'hi ha exactament una $I\to X$:
\[
\begin{tikzpicture}[baseline=(m.center)]
  \matrix (m) [matrix of math nodes, column sep=3.2em] { X & T \\ };
  \draw[-{Stealth[scale=1.1]}, dashed] (m-1-1) -- node[above] {$\exists!$} (m-1-2);
\end{tikzpicture}
\qquad\qquad
\begin{tikzpicture}[baseline=(m.center)]
  \matrix (m) [matrix of math nodes, column sep=3.2em] { I & X \\ };
  \draw[-{Stealth[scale=1.1]}, dashed] (m-1-1) -- node[above] {$\exists!$} (m-1-2);
\end{tikzpicture}
\]
En aquests diagrames seguim la convenció que ja vam introduir en presentar el producte
(subsecció~\ref{subsec:producte-coproducte}): la \textbf{fletxa discontínua} assenyala
el morfisme l'existència del qual afirma la propietat universal, i $\exists!$ abreuja
\enquote{existeix un únic}; com que aquí el morfisme encara no té nom, hi escrivim
directament $\exists!$. Cada diagrama es llegeix, doncs: per a tot $X$ existeix un únic
morfisme cap a $T$ (respectivament, des de $I$).

A $\cat{Set}$, per exemple, un objecte terminal és qualsevol conjunt d'un sol element i
un objecte inicial és el conjunt buit. El que distingeix aquests objectes no és cap
propietat dels seus elements, sinó la fletxa: la condició els fixa completament, fins
al punt que dos objectes terminals qualssevol són isomorfs \emph{d'una única manera},
com demostrarem tot seguit (proposició~\ref{prop:unicitat-terminal}). Aquesta és, en
miniatura, una propietat universal: una condició sobre \emph{totes} les fletxes cap a
un objecte ---o des d'ell--- que el determina llevat d'un únic isomorfisme. En aquest cas, l'isomorfisme és únic fins i tot com a isomorfisme d'objectes, perquè no hi ha cap dada addicional a preservar: entre dos objectes terminals només hi ha una fletxa en cada sentit. Escrivim
encara $T$ i $I$, i no $1$ i $0$, a propòsit: la notació numèrica dona per fet que
\enquote{l'} objecte terminal està ben determinat, cosa que només ens autoritzem a
afirmar un cop provada aquesta unicitat. A la secció següent la demostrem i, aleshores
sí, batejarem $T$ com a $1$ i $I$ com a $0$.

\subsection{Un exemple ja conegut: la construcció lliure}\label{subsec:construccio-lliure}

Un exemple menys trivial ja ha aparegut: la categoria lliure\index{categoria!lliure}
$\Free(G)$ generada per un graf (seccions~\ref{sec:grafs-lliures}
i~\ref{sec:functors-grafs}), i el seu cas d'un sol vèrtex, el monoide
lliure\index{monoide!lliure}. Allà vam deixar pendent de justificar el nom
\enquote{lliure}; ara podem llegir-lo com una propietat universal.

La propietat posa en joc dues categories diferents, i convé tenir present en
quina treballem a cada moment. D'una banda hi ha $\cat{Graph}$, on viuen els
grafs i els morfismes de grafs; de l'altra, $\cat{Cat}$, on viuen les categories
petites i els functors entre elles. Els functors del
capítol~\ref{cap:functors}, $\Free\colon\cat{Graph}\to\cat{Cat}$ i
$U\colon\cat{Cat}\to\cat{Graph}$, fan de pont entre tots dos nivells; no són
fletxes de cap d'aquestes categories ---són functors entre categories grans---,
sinó la manera de passar de l'una a l'altra. Així, $\Free(G)$ és un objecte de
$\cat{Cat}$, mentre que $U(\Free(G))$, el seu graf subjacent, és un objecte de
$\cat{Graph}$.

La clau és una \textbf{inserció dels generadors}, el morfisme de grafs
\[
\eta_G\colon G\to U(\Free(G)),
\]
que deixa fixos els vèrtexs i envia cada aresta al camí de longitud~1
corresponent. Aquesta inserció té la propietat següent: per a tota categoria
petita $\cat{D}$ i tot morfisme de grafs $f\colon G\to U(\cat{D})$ ---una fletxa
de $\cat{Graph}$, que tria un objecte de $\cat{D}$ per a cada vèrtex i un
morfisme de $\cat{D}$ per a cada aresta, sense cap exigència sobre
composicions---, existeix un \emph{únic} functor $\bar f\colon\Free(G)\to\cat{D}$
---una fletxa de $\cat{Cat}$--- que estén $f$. Com que $f$ i $\bar f$ no viuen a
la mateixa categoria, \enquote{estendre} no es pot expressar comparant-los
directament: primer apliquem $U$ a $\bar f$, cosa que en dona el morfisme de
grafs $U(\bar f)\colon U(\Free(G))\to U(\cat{D})$, i aleshores la condició és la
igualtat $U(\bar f)\comp\eta_G=f$ entre morfismes de $\cat{Graph}$.
\[
\begin{tikzpicture}[baseline=(m.center)]
  \matrix (m) [matrix of math nodes, column sep=3.4em, row sep=2.8em] {
    G & U(\Free(G)) \\
      & U(\cat{D}) \\
  };
  \draw[-{Stealth[scale=1.1]}] (m-1-1) -- node[above] {$\eta_G$} (m-1-2);
  \draw[-{Stealth[scale=1.1]}] (m-1-1) -- node[below left] {$f$} (m-2-2);
  \draw[-{Stealth[scale=1.1]}] (m-1-2) -- node[right] {$U(\bar f)$} (m-2-2);
  \node[anchor=south, font=\small] at ([yshift=0.6em]m.north) {a $\cat{Graph}$};
\end{tikzpicture}
\qquad\qquad
\begin{tikzpicture}[baseline=(m.center)]
  \matrix (m) [matrix of math nodes, row sep=2.8em] {
    \Free(G) \\
    \cat{D} \\
  };
  \draw[-{Stealth[scale=1.1]}, dashed] (m-1-1) -- node[right] {$\exists!\,\bar f$} (m-2-1);
  \node[anchor=south, font=\small] at ([yshift=0.6em]m.north) {a $\cat{Cat}$};
\end{tikzpicture}
\]
El triangle de l'esquerra commuta a $\cat{Graph}$. La fletxa discontínua, la
que la propietat afirma que existeix i és única, és $\bar f$, i viu a $\cat{Cat}$;
la seva imatge per $U$ és la que tanca el triangle.
Aquesta és, precisament, la correspondència que vam establir al
capítol~\ref{cap:functors}: definir un functor \emph{des de} $\Free(G)$ no demana res
més que dir on van els generadors. La construcció $\Free(G)$ no s'ha \emph{definit} per
aquesta propietat ---s'ha bastit amb camins---, però la propietat la
\emph{caracteritza}: qualsevol altra categoria petita $\cat{F}$ dotada d'un morfisme de grafs
$G\to U(\cat{F})$ amb la mateixa condició seria isomorfa a $\Free(G)$, i d'una única manera compatible amb la
inserció. En això consisteix una propietat universal. Dit en el llenguatge que introduirem a la secció de representabilitat, $\Free(G)$ representa el functor
\[
\cat{D}\longmapsto\Hom_{\cat{Graph}}\bigl(G,U(\cat{D})\bigr),
\]
i $\eta_G\in\Hom_{\cat{Graph}}\bigl(G,U(\Free(G))\bigr)$ n'és l'\emph{element universal}. La formulació completa d'aquesta correspondència com a adjunció $\Free\dashv U$ arribarà al capítol d'adjuncions.

En aquest capítol precisem què és una propietat universal, la lliguem amb els
homfunctors $\Hom(A,-)$ i $\Hom(-,B)$ que vam introduir, i culminem amb el
\emph{lema de Yoneda}, que fa explícita la idea que un objecte queda determinat pels
morfismes que hi arriben o que en surten.

\section{Objectes inicials i terminals}

\begin{definicio}\index{objecte!inicial}\index{objecte!terminal}
En una categoria $\cat{C}$:
\begin{itemize}
\item un objecte $I$ és \textbf{inicial} si per a cada objecte $X$ hi ha \emph{exactament un} morfisme $I\to X$;
\item un objecte $T$ és \textbf{terminal} si per a cada objecte $X$ hi ha \emph{exactament un} morfisme $X\to T$.
\end{itemize}
Les dues nocions són duals: un objecte inicial a $\cat{C}$ és un objecte terminal a $\cat{C}\op$.
\end{definicio}

\begin{exemple}
A $\cat{Set}$, l'objecte inicial és el conjunt buit $\varnothing$ ---hi ha una única aplicació $\varnothing\to X$, la buida--- i l'objecte terminal és qualsevol conjunt d'un sol element $\{\ast\}$ ---hi ha una única aplicació $X\to\{\ast\}$. A $\cat{Grp}$, en canvi, el grup trivial és alhora inicial i terminal (un tal objecte s'anomena \emph{zero}). En un preordre vist com a categoria, un objecte inicial és un \emph{mínim global}, un element $m$ amb $m\le x$ per a tot $x$, i un de terminal, un \emph{màxim global}. No s'han de confondre amb els elements minimals o maximals, que només demanen que no hi hagi cap element estrictament per sota o per sobre.
\end{exemple}

La propietat clau, model de totes les propietats universals, és la unicitat.

\begin{proposicio}\label{prop:unicitat-terminal}\index{objecte!unicitat llevat d'isomorfisme}
Si una categoria té dos objectes terminals, aleshores són isomorfs, i l'isomorfisme entre ells és únic. El mateix val per als objectes inicials.
\end{proposicio}

\begin{prova}
Siguin $T$ i $T'$ dos objectes terminals. Com que $T'$ és terminal, hi ha un únic morfisme $f\colon T\to T'$; com que $T$ és terminal, hi ha un únic morfisme $g\colon T'\to T$. La composició $g\comp f\colon T\to T$ és un morfisme de $T$ a $T$; però $T$ és terminal, de manera que l'únic morfisme $T\to T$ és $\id_T$, i per tant $g\comp f=\id_T$. Anàlogament $f\comp g=\id_{T'}$. Així $f$ és un isomorfisme, i és únic perquè ho és el morfisme $T\to T'$. El cas inicial és dual.
\end{prova}

Provada la unicitat llevat d'un únic isomorfisme, ja té sentit parlar de \enquote{l'} objecte terminal i \enquote{l'} objecte inicial. D'ara endavant els escrivim, respectivament, $1$ i $0$. Aquesta notació és estàndard i no pressuposa cap finitud: prové del fet que l'objecte terminal és el producte buit ---la unitat de $\times$--- i l'inicial el coproducte buit ---el zero de $+$---, com veurem al capítol~\ref{cap:limits}.

\begin{observacio}
Aquesta demostració és el patró que retrobarem sempre, però convé enunciar-lo amb la precisió justa. Dues \emph{dades universals} del mateix tipus ---dos objectes terminals, dos productes $(P,p_1,p_2)$ i $(Q,q_1,q_2)$, dues representacions--- són isomorfes mitjançant un \emph{únic} isomorfisme compatible amb \emph{tota} la dada universal (les projeccions, la representació, etc.). Els objectes \emph{subjacents} són, per tant, isomorfs; però l'isomorfisme entre els objectes \emph{nus}, un cop oblidada la dada, no ha de ser únic: un producte $P$ pot admetre automorfismes no trivials que no respecten les projeccions. En el cas dels objectes terminals i inicials no hi ha dada addicional a preservar, i per això l'isomorfisme és únic sense més (proposició~\ref{prop:unicitat-terminal}). Amb aquesta cautela parlem de \enquote{l'} objecte terminal, \enquote{el} producte, etc.: la unicitat \emph{compatible amb la dada} és el que dona sentit a caracteritzar objectes per la seva propietat universal en lloc de per una construcció concreta.
\end{observacio}

\section{Categories coma}\label{sec:coma}\index{categoria!coma}

Abans de continuar val la pena introduir una construcció que unifica diverses coses que hem trobat de passada ---les categories sobre i sota un objecte, i ben aviat els cons sobre un diagrama--- i que permet dir amb precisió què vol dir que un objecte és \enquote{universal}: n'hi haurà prou que sigui \emph{inicial} o \emph{terminal} en una categoria adequada.

\begin{definicio}[Categoria coma]\index{categoria!coma}
\label{def:coma}
Siguin $S\colon\cat{A}\to\cat{C}$ i $T\colon\cat{B}\to\cat{C}$ dos functors amb el mateix codomini. La \textbf{categoria coma} $(S\downarrow T)$ té:
\begin{itemize}
  \item per \textbf{objectes}, les ternes $(A,B,f)$ amb $A$ objecte de $\cat{A}$, $B$ objecte de $\cat{B}$ i $f\colon S(A)\to T(B)$ un morfisme de $\cat{C}$;
  \item per \textbf{morfismes} $(A,B,f)\to(A',B',f')$, les parelles $(a,b)$ amb $a\colon A\to A'$ a $\cat{A}$ i $b\colon B\to B'$ a $\cat{B}$ tals que el quadrat
  \[
  \begin{tikzpicture}[baseline=(m.center)]
    \matrix (m) [matrix of math nodes, column sep=3em, row sep=2.4em] {
      S(A) & S(A') \\
      T(B) & T(B') \\
    };
    \draw[-{Stealth[scale=1.1]}] (m-1-1) -- node[above] {$S(a)$} (m-1-2);
    \draw[-{Stealth[scale=1.1]}] (m-1-1) -- node[left] {$f$} (m-2-1);
    \draw[-{Stealth[scale=1.1]}] (m-1-2) -- node[right] {$f'$} (m-2-2);
    \draw[-{Stealth[scale=1.1]}] (m-2-1) -- node[below] {$T(b)$} (m-2-2);
  \end{tikzpicture}
  \]
  commuta a $\cat{C}$, és a dir, $T(b)\comp f=f'\comp S(a)$.
\end{itemize}
La composició i les identitats es fan component a component.
\end{definicio}

\begin{observacio}[Sobre el nom i la notació]\index{categoria!coma!notació}
El concepte i el nom es deuen a Lawvere, que a la seva tesi doctoral (1963)
escrivia aquesta categoria $(S,T)$, per analogia amb la notació $(A,B)$ dels
homconjunts: els objectes de la categoria coma són, en efecte, morfismes de
$S(A)$ a $T(B)$, una mena d'\enquote{homconjunt entre functors}. D'aquí ve el
nom de \emph{categoria coma}, que ha sobreviscut tot i que la coma, usada com a
operador, es prestava a confusió i ha estat substituïda. La notació
$(S\downarrow T)$, que és la de Mac Lane~\cite{maclane} i la que seguim, conserva
els parèntesis de l'original, que a més delimiten els operands quan són
expressions compostes. La fletxa es pot llegir com un dibuix dels objectes: al
quadrat anterior, cada objecte $f$ és una fletxa que \emph{baixa} de la fila de
$S$ a la fila de $T$. També es troben les notacions $S\downarrow T$, sense
parèntesis, i $(S/T)$, que generalitza la notació $\cat{C}/A$ de les categories
sobre un objecte (exemple següent).
\end{observacio}

La notació més freqüent fixa un dels functors en un objecte. Si $\cat{1}$ és la categoria terminal (un objecte i la seva identitat) i $C$ és un objecte de $\cat{C}$, escrivim $C$ també per al functor $\cat{1}\to\cat{C}$ que el tria. Aleshores:

\begin{exemple}[Categories sobre i sota un objecte com a categories coma]\index{categoria!sobre un objecte}\index{categoria!sota un objecte}\label{ex:slice-coma}
Prenem $S=\id_{\cat{C}}$ i $T=C\colon\cat{1}\to\cat{C}$. Un objecte de $(\id_{\cat{C}}\downarrow C)$ és una terna $(A,\ast,f)$ amb $f\colon A\to C$, és a dir, simplement un morfisme cap a $C$. Un morfisme $(A,\ast,f)\to(A',\ast,f')$ és una parella $(a,\id_\ast)$ amb $a\colon A\to A'$ tal que $f'\comp a=f$ (el quadrat de la definició, amb $T(\id_\ast)=\id_C$), és a dir, un triangle commutatiu sobre $C$:
\[
\begin{tikzpicture}[baseline=(m.center)]
  \matrix (m) [matrix of math nodes, column sep=1.6em, row sep=2.6em] {
    A & & A' \\
      & C &  \\
  };
  \draw[-{Stealth[scale=1.1]}] (m-1-1) -- node[above] {$a$} (m-1-3);
  \draw[-{Stealth[scale=1.1]}] (m-1-1) -- node[below left] {$f$} (m-2-2);
  \draw[-{Stealth[scale=1.1]}] (m-1-3) -- node[below right] {$f'$} (m-2-2);
\end{tikzpicture}
\]
Recuperem exactament la categoria $\cat{C}/C$ de $\cat{C}$ sobre $C$. Dualment, $(C\downarrow\id_{\cat{C}})$ és la categoria $C/\cat{C}$ de $\cat{C}$ sota $C$, la dels morfismes que \emph{surten} de $C$. Les categories sobre i sota un objecte són, doncs, els casos més simples de categoria coma.
\end{exemple}

\begin{exemple}[Objectes universals com a inicials o terminals]\label{ex:universal-coma}
Fem precís, en el cas més senzill, l'eslògan que motiva aquesta secció. Fixem
un objecte $C$ d'una categoria localment petita $\cat{C}$ i considerem el
prefeix $\Hom(-,C)\colon\cat{C}\op\to\cat{Set}$ (definició~\ref{def:homfunctors}),
que envia cada objecte $A$ a $\Hom(A,C)$ i cada morfisme $h\colon A\to A'$ a la
precomposició, que amb la notació~\ref{not:hom-morfismes} escrivim
$\Hom(h,C)=h^*$:
\[
\Hom(h,C)\colon\Hom(A',C)\to\Hom(A,C),\qquad x'\mapsto x'\comp h.
\]
La categoria $\cat{C}/C=(\id_{\cat{C}}\downarrow C)$ de l'exemple anterior es
descriu només amb aquest functor. Identificant cada terna $(A,\ast,x)$ amb la
parella $(A,x)$:
\begin{itemize}
  \item un \textbf{objecte} és una parella $(A,x)$ amb $A$ objecte de $\cat{C}$ i
  $x\in\Hom(A,C)$;
  \item un \textbf{morfisme} $(A,x)\to(A',x')$ és un morfisme $h\colon A\to A'$
  de $\cat{C}$ tal que
  \[
    \Hom(h,C)(x')=x,
    \qquad\text{és a dir,}\qquad
    x'\comp h=x,
  \]
  que és el triangle commutatiu sobre $C$:
  \[
  \begin{tikzpicture}[baseline=(m.center)]
    \matrix (m) [matrix of math nodes, column sep=1.6em, row sep=2.6em] {
      A & & A' \\
        & C &  \\
    };
    \draw[-{Stealth[scale=1.1]}] (m-1-1) -- node[above] {$h$} (m-1-3);
    \draw[-{Stealth[scale=1.1]}] (m-1-1) -- node[below left] {$x$} (m-2-2);
    \draw[-{Stealth[scale=1.1]}] (m-1-3) -- node[below right] {$x'$} (m-2-2);
  \end{tikzpicture}
  \]
\end{itemize}

\emph{Afirmació: $(C,\id_C)$ és un objecte terminal de $\cat{C}/C$.} Sigui
$(A,x)$ un objecte qualsevol. Un morfisme $(A,x)\to(C,\id_C)$ és un
$h\colon A\to C$ amb $\Hom(h,C)(\id_C)=x$. Com que
$\Hom(h,C)(\id_C)=\id_C\comp h=h$, la condició diu exactament $h=x$. Per tant
existeix un morfisme $(A,x)\to(C,\id_C)$, i només un: $h=x$.

Llegit en termes del functor, això diu que per a cada objecte $A$ l'aplicació
\[
\Hom(A,C)\longrightarrow\Hom(A,C),\qquad h\longmapsto\Hom(h,C)(\id_C),
\]
és bijectiva (de fet, és la identitat): tot element $x$ de qualsevol $\Hom(A,C)$
s'obté d'un \emph{únic} element distingit, $\id_C\in\Hom(C,C)$, aplicant-hi
$\Hom(h,C)$ per a un \emph{únic} $h$. Aquesta és precisament la forma de la
definició~\ref{def:representable} de la secció següent: $(C,\id_C)$ és una
representació de $\Hom(-,C)$, amb element universal $\id_C$. Allà veurem
(observació~\ref{obs:representacio-inicial}) que, per a qualsevol prefeix $P$
sobre $\cat{C}$ (definició~\ref{def:prefeix}), ser una representació de $P$ equival a ser un
objecte inicial d'una certa categoria coma construïda a partir de $P$; per a
$P=\Hom(-,C)$ aquesta categoria és $(\cat{C}/C)\op$, i ser-hi inicial és ser
terminal a $\cat{C}/C$.

Dualment, amb el homfunctor covariant $\Hom(C,-)$, $(C,\id_C)$ és un objecte
inicial de $C/\cat{C}$: un morfisme $(C,\id_C)\to(A,x)$ és un $h\colon C\to A$
amb $\Hom(C,h)(\id_C)=h\comp\id_C=x$, és a dir, $h=x$.

Un cas menys trivial és la construcció lliure de la
subsecció~\ref{subsec:construccio-lliure}. Sigui $G$ un graf, vist com a functor
$G\colon\cat{1}\to\cat{Graph}$, i sigui $U\colon\cat{Cat}\to\cat{Graph}$ el
functor oblit. Un objecte de la categoria coma $(G\downarrow U)$ és una parella
$(\cat{D},f)$ formada per una categoria petita $\cat{D}$ i un morfisme de grafs
$f\colon G\to U(\cat{D})$; un morfisme $(\cat{D},f)\to(\cat{D}',f')$ és un
functor $F\colon\cat{D}\to\cat{D}'$ amb $U(F)\comp f=f'$. Que
$(\Free(G),\eta_G)$ sigui un objecte inicial de $(G\downarrow U)$ vol dir que
per a cada $(\cat{D},f)$ hi ha un únic functor $\bar f\colon\Free(G)\to\cat{D}$
amb $U(\bar f)\comp\eta_G=f$: és, literalment, la propietat universal que hi
vam enunciar.

Aquests casos fan precís l'eslògan que unifica les propietats universals:
\emph{ser universal és ser inicial o terminal en una categoria coma adient}. Ho
retrobarem al capítol següent: un límit és un objecte terminal de la categoria dels
cons, que és una categoria coma, i un colímit, un objecte inicial de la dels cocons
(observació~\ref{obs:cons-coma}).
\end{exemple}

\section{Propietats universals com a representabilitat}

Podem reformular aquestes idees amb el llenguatge dels functors representables. Recordem que, en una categoria localment petita, cada objecte $A$ determina el prefeix $\Hom(-,A)\colon\cat{C}\op\to\cat{Set}$ i el functor covariant $\Hom(A,-)\colon\cat{C}\to\cat{Set}$.

Observem primer com es llegeixen les nocions anteriors. Que $1$ sigui terminal vol dir que $\Hom(X,1)$ té exactament un element per a cada $X$; és a dir, que el functor $\Hom(-,1)$ és \emph{naturalment isomorf} al prefeix constant $\Delta_{\{\ast\}}\colon\cat{C}\op\to\cat{Set}$ (exemple~\ref{ex:functor-constant}). Les components de l'isomorfisme són les úniques bijeccions $\Hom(X,1)\to\{\ast\}$, i la naturalitat és automàtica, perquè tots els quadrats acaben en un conjunt d'un sol element. Dualment, $0$ és inicial quan $\Hom(0,-)$ és naturalment isomorf al functor constant $\Delta_{\{\ast\}}\colon\cat{C}\to\cat{Set}$. En cap dels dos casos no és una igualtat literal: els conjunts d'un sol element $\Hom(X,1)$, i també els $\Hom(0,X)$, són conjunts diferents per a objectes $X$ diferents, i no s'han identificat per definició amb $\{\ast\}$.

Una forma general de les propietats universals que estudiarem és la representabilitat d'un functor cap a $\cat{Set}$: un objecte, juntament amb una dada universal, satisfà una propietat d'aquesta mena quan \emph{representa} el functor que recull les dades del problema.

\begin{definicio}[Representació]\index{propietat universal}\index{functor!representable}\index{element universal}
\label{def:representable}
Sigui $\cat{C}$ una categoria localment petita.
\begin{enumerate}
\item Sigui $F\colon\cat{C}\to\cat{Set}$ un functor covariant. Una \textbf{representació} de $F$ és una parella $(A,u)$ formada per un objecte $A$ de $\cat{C}$ i un element $u\in F(A)$, l'\textbf{element universal}, tal que per a cada objecte $X$ i cada $x\in F(X)$ existeix un únic morfisme $f\colon A\to X$ amb
\[
F(f)(u)=x.
\]
\item Sigui $P$ un prefeix sobre $\cat{C}$ (definició~\ref{def:prefeix}). Una \textbf{representació} de $P$ és una parella $(A,u)$ formada per un objecte $A$ de $\cat{C}$ i un element $u\in P(A)$, l'\textbf{element universal}, tal que per a cada objecte $X$ i cada $x\in P(X)$ existeix un únic morfisme $f\colon X\to A$ amb
\[
P(f)(u)=x.
\]
\end{enumerate}
En tots dos casos es diu que el functor és \textbf{representable}, que $A$ n'és l'\textbf{objecte representant} i que $(A,u)$ n'és una representació.
\end{definicio}

Amb la noció d'isomorfisme natural (capítol~\ref{cap:transformacions}), la representabilitat es pot formular sense esmentar l'element universal.

\begin{proposicio}[Representabilitat com a isomorfisme natural]\label{prop:representable-iso}\index{functor!representable}
Sigui $\cat{C}$ una categoria localment petita.
\begin{enumerate}
\item Un functor $F\colon\cat{C}\to\cat{Set}$ és representable si i només si és naturalment isomorf al functor $\Hom_{\cat{C}}(S,-)\colon\cat{C}\to\cat{Set}$ per a algun objecte $S$ de $\cat{C}$. En aquest cas, $S$ és un objecte representant de $F$.
\item Un prefeix $P\colon\cat{C}\op\to\cat{Set}$ és representable si i només si és naturalment isomorf al prefeix $\Hom_{\cat{C}}(-,S)\colon\cat{C}\op\to\cat{Set}$ per a algun objecte $S$ de $\cat{C}$. De nou, en aquest cas $S$ és un objecte representant de $P$.
\end{enumerate}
Més precisament, les representacions $(S,u)$ es corresponen amb els isomorfismes naturals
\[
\alpha\colon\Hom(S,-)\Rightarrow F
\qquad\text{o bé}\qquad
\alpha\colon\Hom(-,S)\Rightarrow P,
\]
segons el cas, mitjançant $u=\alpha_S(\id_S)$ i $\alpha_X(f)=F(f)(u)$, respectivament $\alpha_X(f)=P(f)(u)$.
\end{proposicio}

\begin{prova}
Fem primer el cas~(1). Els quadrats de naturalitat segueixen el criteri del capítol~\ref{cap:transformacions}: l'acció dels functors sobre un morfisme, en horitzontal (a dalt l'homfunctor, a baix $F$), i les components de la transformació, en vertical. Al costat de cada quadrat hi ha el camí que hi fa un element: les fletxes $\mapsto$ indiquen on va a parar cada element.

\emph{D'una representació a un isomorfisme natural.} Sigui $(S,u)$ una representació de $F$ i definim $\alpha_X\colon\Hom(S,X)\to F(X)$ per $\alpha_X(f)=F(f)(u)$. Per veure que $\alpha\colon\Hom(S,-)\Rightarrow F$ és natural, cal que per a cada $g\colon X\to Y$ commuti el quadrat següent, on l'acció de l'homfunctor sobre $g$ és la postcomposició $g_*=\Hom(S,g)$ (notació~\ref{not:hom-morfismes}):
\[
\begin{tikzpicture}[baseline=(m.center)]
  \matrix (m) [matrix of math nodes, row sep=2.8em, column sep=3.6em] {
    \Hom(S,X) & \Hom(S,Y) \\
    F(X) & F(Y) \\
  };
  \draw[-{Stealth[scale=1.1]}] (m-1-1) -- node[above] {$g_*$} (m-1-2);
  \draw[-{Stealth[scale=1.1]}] (m-1-1) -- node[left] {$\alpha_X$} (m-2-1);
  \draw[-{Stealth[scale=1.1]}] (m-1-2) -- node[right] {$\alpha_Y$} (m-2-2);
  \draw[-{Stealth[scale=1.1]}] (m-2-1) -- node[below] {$F(g)$} (m-2-2);
\end{tikzpicture}
\qquad\qquad
\begin{tikzpicture}[baseline=(m.center)]
  \matrix (m) [matrix of math nodes, row sep=2.8em, column sep=2.4em] {
    f & g\comp f \\
    F(f)(u) & F(g)\bigl(F(f)(u)\bigr) \\
  };
  \draw[{Bar[width=4pt]}-{Stealth[scale=1.1]}] (m-1-1) -- (m-1-2);
  \draw[{Bar[width=4pt]}-{Stealth[scale=1.1]}] (m-1-1) -- (m-2-1);
  \draw[{Bar[width=4pt]}-{Stealth[scale=1.1]}] (m-1-2) -- (m-2-2);
  \draw[{Bar[width=4pt]}-{Stealth[scale=1.1]}] (m-2-1) -- (m-2-2);
\end{tikzpicture}
\]
Seguint un element $f$ pels dos camins, per dalt i després avall s'obté $\alpha_Y(g\comp f)=F(g\comp f)(u)$, i avall i després per baix s'obté $F(g)\bigl(F(f)(u)\bigr)$. Totes dues coincideixen perquè $F$ preserva la composició:
\[
\alpha_Y(g\comp f)=F(g\comp f)(u)=F(g)\bigl(F(f)(u)\bigr)=F(g)\bigl(\alpha_X(f)\bigr).
\]
La condició de representació diu exactament que cada $\alpha_X$ és bijectiva, i per la proposició~\ref{prop:iso-natural-components} $\alpha$ és un isomorfisme natural.

\emph{D'un isomorfisme natural a una representació.} Recíprocament, sigui $\alpha\colon\Hom(S,-)\Rightarrow F$ un isomorfisme natural i posem $u=\alpha_S(\id_S)\in F(S)$. Per a cada $f\colon S\to X$, escrivim el quadrat de naturalitat de $\alpha$ respecte de $f$ i hi seguim l'element $\id_S$:
\[
\begin{tikzpicture}[baseline=(m.center)]
  \matrix (m) [matrix of math nodes, row sep=2.8em, column sep=3.6em] {
    \Hom(S,S) & \Hom(S,X) \\
    F(S) & F(X) \\
  };
  \draw[-{Stealth[scale=1.1]}] (m-1-1) -- node[above] {$f_*$} (m-1-2);
  \draw[-{Stealth[scale=1.1]}] (m-1-1) -- node[left] {$\alpha_S$} (m-2-1);
  \draw[-{Stealth[scale=1.1]}] (m-1-2) -- node[right] {$\alpha_X$} (m-2-2);
  \draw[-{Stealth[scale=1.1]}] (m-2-1) -- node[below] {$F(f)$} (m-2-2);
\end{tikzpicture}
\qquad\qquad
\begin{tikzpicture}[baseline=(m.center)]
  \matrix (m) [matrix of math nodes, row sep=2.8em, column sep=2.4em] {
    \id_S & f \\
    u & F(f)(u) \\
  };
  \draw[{Bar[width=4pt]}-{Stealth[scale=1.1]}] (m-1-1) -- (m-1-2);
  \draw[{Bar[width=4pt]}-{Stealth[scale=1.1]}] (m-1-1) -- (m-2-1);
  \draw[{Bar[width=4pt]}-{Stealth[scale=1.1]}] (m-1-2) -- (m-2-2);
  \draw[{Bar[width=4pt]}-{Stealth[scale=1.1]}] (m-2-1) -- (m-2-2);
\end{tikzpicture}
\]
Per dalt, $\id_S$ va a $f\comp\id_S=f$ i després a $\alpha_X(f)$; avall, va a $u$ i després a $F(f)(u)$. La commutativitat del quadrat dona, doncs,
\[
\alpha_X(f)=\alpha_X(f\comp\id_S)=F(f)\bigl(\alpha_S(\id_S)\bigr)=F(f)(u).
\]
Això diu que tota la transformació $\alpha$ queda determinada pel sol element $u$. Com que $\alpha_X$ és bijectiva (proposició~\ref{prop:iso-natural-components}), per a cada $x\in F(X)$ hi ha un únic $f\colon S\to X$ amb $F(f)(u)=x$: $(S,u)$ és una representació de $F$.

Les dues construccions són inverses l'una de l'altra. Partint de $u$, la transformació $\alpha$ torna $\alpha_S(\id_S)=F(\id_S)(u)=u$; i partint de $\alpha$, la igualtat anterior diu que $\alpha$ és precisament la transformació $f\mapsto F(f)(u)$ construïda a partir de $u=\alpha_S(\id_S)$.

El cas~(2) és el cas~(1) aplicat a $\cat{C}\op$, però val la pena veure com canvia el diagrama clau. Ara $\alpha\colon\Hom(-,S)\Rightarrow P$, els morfismes arriben a $S$ en lloc de sortir-ne, i l'acció de l'homfunctor sobre un morfisme és la precomposició $f^*=\Hom(f,S)$. Per a cada $f\colon X\to S$, el quadrat de naturalitat respecte de $f$, amb l'element $\id_S$, és
\[
\begin{tikzpicture}[baseline=(m.center)]
  \matrix (m) [matrix of math nodes, row sep=2.8em, column sep=3.6em] {
    \Hom(S,S) & \Hom(X,S) \\
    P(S) & P(X) \\
  };
  \draw[-{Stealth[scale=1.1]}] (m-1-1) -- node[above] {$f^*$} (m-1-2);
  \draw[-{Stealth[scale=1.1]}] (m-1-1) -- node[left] {$\alpha_S$} (m-2-1);
  \draw[-{Stealth[scale=1.1]}] (m-1-2) -- node[right] {$\alpha_X$} (m-2-2);
  \draw[-{Stealth[scale=1.1]}] (m-2-1) -- node[below] {$P(f)$} (m-2-2);
\end{tikzpicture}
\qquad\qquad
\begin{tikzpicture}[baseline=(m.center)]
  \matrix (m) [matrix of math nodes, row sep=2.8em, column sep=2.4em] {
    \id_S & f \\
    u & P(f)(u) \\
  };
  \draw[{Bar[width=4pt]}-{Stealth[scale=1.1]}] (m-1-1) -- (m-1-2);
  \draw[{Bar[width=4pt]}-{Stealth[scale=1.1]}] (m-1-1) -- (m-2-1);
  \draw[{Bar[width=4pt]}-{Stealth[scale=1.1]}] (m-1-2) -- (m-2-2);
  \draw[{Bar[width=4pt]}-{Stealth[scale=1.1]}] (m-2-1) -- (m-2-2);
\end{tikzpicture}
\]
on ara, per dalt, $\id_S$ va a $\id_S\comp f=f$. Com que els functors són contravariants, les fletxes horitzontals van de l'objecte $S$ a l'objecte $X$, en sentit contrari a $f\colon X\to S$. Se'n dedueix $\alpha_X(f)=P(f)(u)$, i la resta de l'argument és el mateix.
\end{prova}

Les dues nocions són duals l'una de l'altra. Un prefeix sobre $\cat{C}$ és un functor covariant $\cat{C}\op\to\cat{Set}$, i com que un morfisme $X\to A$ de $\cat{C}$ és un morfisme $A\to X$ de $\cat{C}\op$, una representació d'un prefeix en el sentit~(2) és exactament una representació en el sentit~(1) del functor covariant corresponent sobre $\cat{C}\op$. Per això, tot el que demostrem per a un dels dos casos val, per dualitat, per a l'altre. Per exemple, $0$ és inicial si i només si $(0,\ast)$ és una representació del functor covariant $\Delta_{\{\ast\}}$, i $1$ és terminal si i només si $(1,\ast)$ és una representació del prefeix $\Delta_{\{\ast\}}$.

\begin{observacio}[Sondes i pantalles]\label{obs:sondes-pantalles}\index{functor!representable!interpretació}
Els functors representables admeten una lectura intuïtiva. El functor representable $\Hom_{\cat{C}}(S,-)\colon\cat{C}\to\cat{Set}$ pren un objecte $X$ i en dona el conjunt $\Hom_{\cat{C}}(S,X)$ de fletxes de $S$ a $X$. Les característiques de $X$ que extreu són exactament totes les maneres possibles de fer arribar $S$ dins de $X$, com si $S$ fos una \emph{sonda} amb què explorem l'estructura de $X$. Dir que un functor $F$ és representable per $S$ és dir que el que $F$ mesura de cada objecte és precisament el que en veu aquesta sonda.

Dualment, el prefeix representable $\Hom_{\cat{C}}(-,S)\colon\cat{C}\op\to\cat{Set}$ pren un objecte $X$ i en dona el conjunt $\Hom_{\cat{C}}(X,S)$ de fletxes de $X$ a $S$. Les característiques de $X$ que extreu són totes les maneres possibles d'enviar $X$ a $S$, o totes les observacions de $X$ amb valors a $S$: podem pensar $S$ com una \emph{pantalla} sobre la qual $X$ es projecta de moltes maneres diferents.
\end{observacio}

Vegem-ne exemples. En tots, l'isomorfisme natural es pot comprovar directament o, equivalentment, n'hi ha prou de donar l'element universal (proposició~\ref{prop:representable-iso}).

\begin{exemple}[Vèrtexs i arestes d'un graf]\label{ex:rep-grafs}\index{functor!representable!a Graph}
A l'exemple~\ref{ex:homfunctor-grafs} vam veure que els functors $\cat{Graph}\to\cat{Set}$ \enquote{conjunt de vèrtexs} i \enquote{conjunt d'arestes} són naturalment isomorfs a $\Hom_{\cat{Graph}}(\mathbf V,-)$ i a $\Hom_{\cat{Graph}}(\mathbf E,-)$. Són, doncs, representables, amb objectes representants el graf $\mathbf V$ d'un sol vèrtex i el graf $\mathbf E$ d'una sola aresta. Els elements universals són l'únic vèrtex $\ast$ de $\mathbf V$ i l'única aresta $e$ de $\mathbf E$: tot vèrtex d'un graf $G$ és la imatge de $\ast$ per un únic morfisme $\mathbf V\to G$, i tota aresta és la imatge de $e$ per un únic morfisme $\mathbf E\to G$. Són sondes en el sentit literal: $\mathbf V$ detecta els vèrtexs i $\mathbf E$ detecta les arestes.
\end{exemple}

\begin{exemple}[Functors oblit de $\cat{Grp}$, $\cat{Mon}$ i $\cat{Vect}_{\mathbb K}$]\label{ex:rep-grp}\index{functor!representable!oblit}
El functor oblit $U\colon\cat{Grp}\to\cat{Set}$ és representable, amb objecte representant el grup additiu $\mathbb Z$ i element universal $1\in U(\mathbb Z)$. En efecte, un homomorfisme $\varphi\colon\mathbb Z\to G$ queda determinat per $g=\varphi(1)$, perquè
\[
\varphi(n)=g^n\qquad\text{per a tot } n\in\mathbb Z;
\]
i, recíprocament, per a cada $g\in G$ l'aplicació $n\mapsto g^n$ és un homomorfisme. Per tant, $\varphi\mapsto\varphi(1)$ és una bijecció
\[
\Hom_{\cat{Grp}}(\mathbb Z,G)\;\cong\;U(G),
\]
natural en $G$, perquè per a un homomorfisme $\psi\colon G\to H$ es té $(\psi\comp\varphi)(1)=\psi(\varphi(1))$.

El mateix argument val per als altres dos models algebraics del capítol~\ref{cap:categories}. El functor oblit $\cat{Mon}\to\cat{Set}$ és representable pel monoide $(\mathbb N,+,0)$ amb element universal $1$: un homomorfisme de monoides $h\colon\mathbb N\to M$ queda determinat per $m=h(1)$, perquè $h(n)=m\ast\cdots\ast m$ ($n$ vegades, i $h(0)$ el neutre), i per a cada $m$ aquesta fórmula defineix un homomorfisme. El functor oblit $\cat{Vect}_{\mathbb K}\to\cat{Set}$ és representable per $\mathbb K$, vist com a espai vectorial sobre si mateix, amb element universal $1$: una aplicació lineal $T\colon\mathbb K\to V$ queda determinada per $v=T(1)$, perquè $T(\lambda)=\lambda v$, i per a cada $v\in V$ l'aplicació $\lambda\mapsto\lambda v$ és lineal.

En els tres casos l'objecte representant és l'estructura \enquote{generada lliurement per un sol element}, i com a sonda veu exactament els elements de cada objecte.
\end{exemple}

\begin{exemple}[El punt a $\cat{Set}$]\label{ex:rep-punt}
El functor identitat de $\cat{Set}$ és representable per un conjunt d'un sol element $\{\ast\}$, amb element universal $\ast$: una aplicació $\{\ast\}\to X$ és el mateix que un element de $X$, i $\Hom_{\cat{Set}}(\{\ast\},X)\cong X$ naturalment en $X$. El punt és la sonda que veu els elements d'un conjunt.
\end{exemple}

\begin{exemple}[Subconjunts i valors de veritat]\label{ex:rep-parts}\index{functor!representable!de parts}
El functor de parts $\mathcal P\colon\cat{Set}\op\to\cat{Set}$ és un prefeix representable: a l'exemple~\ref{ex:functor-parts} vam veure que les funcions característiques donen bijeccions canòniques $\mathcal P(X)\cong\Hom_{\cat{Set}}(X,\{0,1\})$ sota les quals l'antiimatge correspon a la precomposició. En el llenguatge del capítol~\ref{cap:transformacions}, això diu exactament que $\mathcal P\cong\Hom_{\cat{Set}}(-,\{0,1\})$ naturalment: el quadrat de naturalitat és la igualtat $\chi_B\comp f=\chi_{f^{-1}(B)}$ (exemple~\ref{ex:isos-naturals-cap2}). L'objecte representant és el conjunt $\{0,1\}$ de valors de veritat, i l'element universal és $\{1\}\in\mathcal P(\{0,1\})$, el subconjunt dels valors \enquote{cert}: per a cada $B\subseteq X$, l'única aplicació $\chi\colon X\to\{0,1\}$ amb $\chi^{-1}(\{1\})=B$ és la funció característica $\chi_B$.

Aquí $\{0,1\}$ és la pantalla. Un subconjunt $B\subseteq X$ és el mateix que un predicat sobre $X$, i el predicat és una observació de $X$ amb valors de veritat: $\chi_B(x)=1$ vol dir que $x$ compleix el predicat. Que els subconjunts es puguin classificar així, per morfismes cap a un objecte fix de valors de veritat, és una idea que tornarà amb força al capítol~\ref{cap:topos}, on l'objecte $\{0,1\}$ es generalitza en el \emph{classificador de subobjectes} $\Omega$ (definició~\ref{def:classificador}).
\end{exemple}

\begin{exemple}[El producte com a representació]\label{ex:rep-producte}\index{producte!com a representació}
Siguin $A$ i $B$ objectes d'una categoria localment petita $\cat{C}$. L'assignació
\[
X\longmapsto\Hom(X,A)\times\Hom(X,B),
\qquad
f\longmapsto\bigl((a,b)\mapsto(a\comp f,\,b\comp f)\bigr),
\]
és un prefeix sobre $\cat{C}$: preserva identitats i, per a $f\colon X\to Y$ i $g\colon Y\to Z$, $(a\comp g\comp f,\,b\comp g\comp f)$ s'obté aplicant primer $g$ i després $f$. Una representació d'aquest prefeix és un objecte $Q$ amb un element universal $(p_1,p_2)\in\Hom(Q,A)\times\Hom(Q,B)$ tal que, per a cada $X$ i cada parella $(a,b)$, hi ha un únic $u\colon X\to Q$ amb $p_1\comp u=a$ i $p_2\comp u=b$. És exactament la definició de producte de la subsecció~\ref{subsec:producte-coproducte}: el producte de $A$ i $B$ existeix si i només si aquest prefeix és representable, i un producte és una representació seva.

Aquest exemple té també una lectura lògica. A la categoria prima $\cat{Prov}_T$ associada a un sistema deductiu $T$ (capítol~\ref{cap:categories}), entre dues fórmules hi ha com a màxim un morfisme, i n'hi ha un exactament quan hi ha deduïbilitat $p\vdash_T q$. Per a dues fórmules $p$ i $q$, el prefeix
\[
r\longmapsto\Hom_{\cat{Prov}_T}(r,p)\times\Hom_{\cat{Prov}_T}(r,q)
\]
és, doncs, no buit exactament quan $r\vdash_T p$ i $r\vdash_T q$. Una representació és una fórmula $c$ amb $c\vdash_T p$ i $c\vdash_T q$ (l'element universal) tal que tota fórmula $r$ que dedueix $p$ i $q$ dedueix $c$; la unicitat del morfisme és automàtica. Per transitivitat, això equival a
\[
r\vdash_T c
\quad\Longleftrightarrow\quad
r\vdash_T p\ \text{i}\ r\vdash_T q,
\]
que són les regles d'introducció i d'eliminació de la conjunció. La conjunció $p\wedge q$, quan existeix, és doncs una representació d'aquest prefeix: satisfà la propietat universal del producte, llegida en lògica, tal com anunciava la secció~\ref{sec:dualitat}.
\end{exemple}

\begin{exemple}[Functors que no són representables]\label{ex:no-representable}
Si $\cat{C}$ té algun objecte, el functor constant $\Delta_\varnothing\colon\cat{C}\to\cat{Set}$ no és representable: per a tot objecte $S$, el conjunt $\Hom(S,S)$ conté $\id_S$ i no és buit, mentre que $\Delta_\varnothing(S)=\varnothing$, de manera que no hi pot haver cap bijecció $\Hom(S,S)\to\Delta_\varnothing(S)$. El mateix argument val per al prefeix $\Delta_\varnothing$. Això explica, per exemple, que en una categoria discreta amb dos objectes $A\neq B$ (només hi ha identitats) no hi hagi producte de $A$ i $B$: el prefeix de l'exemple~\ref{ex:rep-producte} hi val $\Hom(X,A)\times\Hom(X,B)=\varnothing$ per a tot $X$, perquè cap objecte té fletxes cap a tots dos, i és, doncs, el prefeix $\Delta_\varnothing$.
\end{exemple}

\begin{observacio}[Representacions com a objectes inicials]\label{obs:representacio-inicial}
Això fa precís l'eslògan de la secció~\ref{sec:coma}. Sigui
$P$ un prefeix sobre $\cat{C}$ i sigui
$\{\ast\}\colon\cat{1}\to\cat{Set}$ el functor que tria un conjunt d'un sol
element. Un objecte de la categoria coma $(\{\ast\}\downarrow P)$ és una terna
$(\ast,A,x)$ amb $x\colon\{\ast\}\to P(A)$, és a dir, un objecte $A$ de
$\cat{C}$ junt amb un element $x\in P(A)$; l'escrivim $(A,x)$. Un morfisme
$(A,x)\to(A',x')$ és un morfisme $A\to A'$ de $\cat{C}\op$ ---és a dir, un
$h\colon A'\to A$ de $\cat{C}$--- que fa commutar el quadrat de la
definició~\ref{def:coma}, condició que aquí diu $P(h)(x)=x'$. Per tant,
$(A,u)$ és un objecte inicial de $(\{\ast\}\downarrow P)$ si i només si per a
cada $(X,x)$ hi ha un únic $h\colon X\to A$ amb $P(h)(u)=x$: exactament la
condició de representació. Per a $P=\Hom(-,C)$, la categoria
$(\{\ast\}\downarrow P)$ és $(\cat{C}/C)\op$, i retrobem
l'exemple~\ref{ex:universal-coma}.
\end{observacio}

\begin{observacio}[La categoria d'elements]\label{obs:categoria-elements}\index{categoria!d'elements}
A la bibliografia és habitual presentar la mateixa construcció amb les fletxes en el sentit de $\cat{C}$. Per a un prefeix $P\colon\cat{C}\op\to\cat{Set}$, la \textbf{categoria d'elements} $\textstyle\int P$ té per objectes les parelles $(A,x)$ amb $x\in P(A)$, i per morfismes $f\colon(A,x)\to(B,y)$ els morfismes $f\colon A\to B$ de $\cat{C}$ amb $P(f)(y)=x$. Té, doncs, les mateixes dades que la categoria coma de l'observació anterior, amb les fletxes girades:
\[
\textstyle\int P\;\cong\;(\{\ast\}\downarrow P)\op.
\]
Per tant, una representació $(C,u)$ de $P$ correspon a un objecte \emph{terminal} de $\int P$: per a cada $(A,x)$ hi ha un únic $f\colon A\to C$ amb $P(f)(u)=x$. És la formulació que fa servir, per exemple, \cite[cap.~2]{riehl}.
\end{observacio}

\begin{observacio}
Cal distingir dos nivells d'unicitat. L'\textbf{objecte} representant $A$ és únic llevat d'isomorfisme, però no necessàriament d'un únic isomorfisme. En canvi, la \textbf{representació} $(A,u)$ ---amb la dada de l'element universal--- és única llevat d'un \emph{únic} isomorfisme compatible amb els elements universals: si $(A,u)$ i $(B,v)$ representen $P$, hi ha un únic isomorfisme $\varphi\colon A\to B$ amb $P(\varphi)(v)=u$. Aquesta distinció, entre la unicitat de l'objecte i la de la dada universal, és essencial i la retrobarem en tractar límits, adjuncions i objectes classificadors. Que dos functors representables $\Hom(-,A)$ i $\Hom(-,B)$ siguin naturalment isomorfs equivaldrà, pel lema de Yoneda, a $A\cong B$; la resta del capítol ho fa precís.
\end{observacio}

\section{La immersió de Yoneda}

Fixem una categoria \textbf{petita} $\cat{C}$. La noció de prefeix, i la de transformació natural entre prefeixos, té sentit sobre qualsevol categoria (definició~\ref{def:prefeix}). Els prefeixos representables $\Hom(-,A)$, en canvi, només són prefeixos si cada $\Hom(X,A)$ és un conjunt, és a dir, si $\cat{C}$ és localment petita, que és la hipòtesi amb què els hem fet servir a la secció de representabilitat. Aquí demanem més: que $\cat{C}$ sigui petita. Aquesta restricció afecta la categoria que aplega els prefeixos. Si $\cat{C}$ no és petita, encara que sigui localment petita, cada prefeix recorre una classe pròpia d'objectes i les transformacions naturals entre dos prefeixos poden formar també una classe pròpia. Aleshores no es pot formar $\cat{Set}^{\cat{C}\op}$ com a categoria (observació~\ref{obs:mida-functors}). Amb $\cat{C}$ petita, en canvi, la categoria de prefeixos $\widehat{\cat{C}}$ i tots els seus homconjunts queden ben definits (exemple~\ref{ex:categoria-prefeixos}).\footnote{Amb universos de Grothendieck, la construcció es pot fer també per a $\cat{C}$ només localment petita, formant $\widehat{\cat{C}}$ en un univers superior; vegeu l'apèndix~\ref{cap:mida} i \cite[§1.1]{riehl}.} A cada objecte $A$ li hem associat el prefeix representable $\Hom(-,A)\colon\cat{C}\op\to\cat{Set}$. Aquesta assignació és, ella mateixa, functorial: dona un functor de $\cat{C}$ cap a la categoria de prefeixos $\widehat{\cat{C}}=\cat{Set}^{\cat{C}\op}$.

\begin{definicio}[Immersió de Yoneda]\index{Yoneda!immersió}
\label{def:immersio-yoneda}
La \textbf{immersió de Yoneda} (en anglès, \emph{Yoneda embedding}) és el functor
\[
\yon\colon\cat{C}\to\cat{Set}^{\cat{C}\op},
\qquad
\yon(A)=\Hom(-,A),
\]
que envia cada objecte $A$ al functor representable $\Hom(-,A)$, i cada morfisme $f\colon A\to B$ a la transformació natural $\yon(f)\colon\Hom(-,A)\Rightarrow\Hom(-,B)$ de components la postcomposició amb $f$,
\[
\yon(f)_X=\Hom(X,f)\colon\Hom(X,A)\to\Hom(X,B),
\qquad
\varphi\mapsto f\comp\varphi.
\]
D'acord amb la notació dels homfunctors, escriurem també $\Hom(-,f)$\index{Hom@$\Hom$!$\Hom(-,f)$} per a aquesta transformació natural $\yon(f)$.
\end{definicio}

\begin{observacio}
Cal comprovar que $\yon(f)$ és realment natural i que $\yon$ preserva composició i identitats, però tot es redueix, un cop més, a l'associativitat i a les lleis de la identitat de $\cat{C}$. El contingut de la immersió de Yoneda és que \emph{tota} categoria petita s'immergeix en la seva categoria de prefeixos, que és una categoria molt rica: com veurem al capítol següent, té tots els límits i colímits \emph{petits}, calculats puntualment. El teorema següent en justifica el nom: $\yon$ és una immersió plena i fidel.
\end{observacio}

\section{El lema de Yoneda}

El resultat central relaciona les transformacions naturals que surten d'un functor representable amb els elements del functor de destí.

\begin{teorema}[Lema de Yoneda]\index{Yoneda!lema}
\label{teo:yoneda}
Sigui $\cat{C}$ una categoria \textbf{petita}, $A$ un objecte de $\cat{C}$ i $P$ un prefeix sobre $\cat{C}$. Aleshores hi ha una bijecció
\[
\Nat\big(\Hom(-,A),\,P\big)\;\cong\;P(A),
\]
natural \emph{contravariantment} en $A\in\cat{C}$ i \emph{covariantment} en $P\in\widehat{\cat{C}}$. Explícitament, la bijecció envia una transformació natural $\eta$ a l'element $\eta_A(\id_A)\in P(A)$.\footnotemark
\end{teorema}
\footnotetext{Si $\cat{C}$ és només localment petita, la mateixa construcció classifica les transformacions naturals que surten del representable mitjançant els elements de $P(A)$, sense haver de formar la categoria $\widehat{\cat{C}}$ de tots els prefeixos. En aquest cas, $\Nat(\Hom(-,A),P)$ s'ha d'entendre com aquesta col·lecció parametritzada pels elements de $P(A)$: en el marc de classes de l'apèndix~\ref{cap:mida}, una transformació amb domini gran és una família indexada per una classe, i no esdevé per això element d'un conjunt. Per parlar-ne literalment com d'un homconjunt dins d'una categoria de prefeixos cal fixar una convenció de mida addicional, per exemple amb universos; vegeu \cite[§1.1]{riehl}.}%

\begin{prova}
Descrivim les dues aplicacions inverses.

\textbf{De transformacions a elements.} A cada $\eta\colon\Hom(-,A)\Rightarrow P$ li assignem
\[
\Phi(\eta)=\eta_A(\id_A)\in P(A).
\]
Té sentit perquè $\id_A\in\Hom(A,A)$ i $\eta_A\colon\Hom(A,A)\to P(A)$.

\textbf{D'elements a transformacions.} A cada $x\in P(A)$ li assignem la transformació $\Psi(x)$ de components
\[
\Psi(x)_X\colon\Hom(X,A)\to P(X),
\qquad
\varphi\mapsto P(\varphi)(x).
\]
Té sentit perquè $\varphi\colon X\to A$ dona, per contravariància, $P(\varphi)\colon P(A)\to P(X)$, i $x\in P(A)$. Cal veure que $\Psi(x)$ és natural: per a un morfisme $g\colon Y\to X$, el quadrat de naturalitat és $P(g)\comp\Psi(x)_X=\Psi(x)_Y\comp g^*$, que sobre $\varphi\in\Hom(X,A)$ dona
\[
P(g)\big(P(\varphi)(x)\big)=P(\varphi\comp g)(x)=\Psi(x)_Y(\varphi\comp g)=\Psi(x)_Y\big(g^*(\varphi)\big),
\]
on hem usat la functorialitat $P(g)\comp P(\varphi)=P(\varphi\comp g)$.

\textbf{Són inverses.} D'una banda,
\[
\Phi(\Psi(x))=\Psi(x)_A(\id_A)=P(\id_A)(x)=x.
\]
De l'altra, donada $\eta$, cal veure que $\Psi(\eta_A(\id_A))=\eta$; és a dir, que per a tot $X$ i tot $\varphi\colon X\to A$,
\[
P(\varphi)\big(\eta_A(\id_A)\big)=\eta_X(\varphi).
\]
Això és exactament el quadrat de naturalitat de $\eta$ per al morfisme $\varphi\colon X\to A$, aplicat a $\id_A\in\Hom(A,A)$: en efecte, $\varphi^*(\id_A)=\id_A\comp\varphi=\varphi$, de manera que
\[
\eta_X(\varphi)=\eta_X(\varphi^*(\id_A))=P(\varphi)(\eta_A(\id_A)).
\]
Les dues assignacions són, doncs, inverses. Queda comprovar que la bijecció és \emph{natural} en les dues variables: és part de l'enunciat, no un afegit, i les dues verificacions són breus. Escrivim $\Phi_{A,P}\colon\Nat(\Hom(-,A),P)\to P(A)$ per exhibir-ne la dependència.

\emph{Naturalitat en $P$} (covariant). Sigui $\sigma\colon P\Rightarrow Q$ una transformació natural de prefeixos, amb postcomposició $(\sigma\comp-)\colon\Nat(\Hom(-,A),P)\to\Nat(\Hom(-,A),Q)$. Per a $\eta\colon\Hom(-,A)\Rightarrow P$,
\[
\Phi_{A,Q}(\sigma\comp\eta)=(\sigma\comp\eta)_A(\id_A)=\sigma_A\big(\eta_A(\id_A)\big)=\sigma_A\big(\Phi_{A,P}(\eta)\big),
\]
és a dir $\Phi_{A,Q}\comp(\sigma\comp-)=\sigma_A\comp\Phi_{A,P}$.

\emph{Naturalitat en $A$} (contravariant). Sigui $a\colon A'\to A$ un morfisme de $\cat{C}$. La immersió dona $\yon(a)=\Hom(-,a)\colon\Hom(-,A')\Rightarrow\Hom(-,A)$, i la precomposició indueix $(-\comp\yon(a))\colon\Nat(\Hom(-,A),P)\to\Nat(\Hom(-,A'),P)$. Per a $\eta\colon\Hom(-,A)\Rightarrow P$, com que $\yon(a)_{A'}(\id_{A'})=a\comp\id_{A'}=a$,
\[
\Phi_{A',P}\big(\eta\comp\yon(a)\big)=\big(\eta\comp\yon(a)\big)_{A'}(\id_{A'})=\eta_{A'}(a)=P(a)\big(\eta_A(\id_A)\big)=P(a)\big(\Phi_{A,P}(\eta)\big),
\]
on la penúltima igualtat és la identitat de naturalitat de $\eta$ demostrada més amunt (la que dona $\Psi(\eta_A(\id_A))=\eta$), aplicada amb $X=A'$ i $\varphi=a$; és a dir $\Phi_{A',P}\comp(-\comp\yon(a))=P(a)\comp\Phi_{A,P}$. Aquestes dues identitats completen el lema tal com s'ha enunciat.
\end{prova}

\begin{observacio}[Taula de control: variància i tipus]\label{obs:taula-yoneda}\index{Yoneda!taula de variància i tipus}
Els errors més freqüents d'aquest capítol són errors de tipus: una fletxa que va en el sentit equivocat o una composició que no existeix. Abans de qualsevol càlcul convé escriure el domini i el codomini de cada component, com al capítol~\ref{cap:functors} (observació~\ref{obs:comprovar-tipus}). La taula recull els casos que hem fet servir, amb $f\colon A\to B$, $u\colon X'\to X$, $v\colon X\to Y$, $a\colon A'\to A$ i $\sigma\colon P\Rightarrow Q$.
\begin{center}
\footnotesize
\renewcommand{\arraystretch}{1.35}
\setlength{\tabcolsep}{4pt}
\begin{tabularx}{\linewidth}{@{}>{\raggedright\arraybackslash}p{0.22\linewidth}>{\raggedright\arraybackslash}p{0.24\linewidth}>{\raggedright\arraybackslash}X@{}}
\toprule
\textbf{Expressió} & \textbf{Què és i variància} & \textbf{Tipus i acció}\\
\midrule
$\Hom(-,A)$ & prefeix $\cat{C}\op\to\cat{Set}$ (contravariant) & $X\mapsto\Hom(X,A)$; \ $\Hom(u,A)=u^*\colon\Hom(X,A)\to\Hom(X',A)$, $\psi\mapsto\psi\comp u$\\
$\Hom(A,-)$ & functor $\cat{C}\to\cat{Set}$ (covariant) & $X\mapsto\Hom(A,X)$; \ $\Hom(A,v)=v_*\colon\Hom(A,X)\to\Hom(A,Y)$, $\varphi\mapsto v\comp\varphi$\\
$\Hom(-,f)=\yon(f)$ & transformació natural $\Hom(-,A)\Rightarrow\Hom(-,B)$, en el sentit de $f$ (no és un functor) & component en $X$: $\Hom(X,f)\colon\Hom(X,A)\to\Hom(X,B)$, $\varphi\mapsto f\comp\varphi$\\
$\Hom(f,-)$ & transformació natural $\Hom(B,-)\Rightarrow\Hom(A,-)$, en sentit \emph{contrari} a $f$ & component en $X$: $\Hom(f,X)\colon\Hom(B,X)\to\Hom(A,X)$, $\psi\mapsto\psi\comp f$\\
$\yon$ & functor $\cat{C}\to\cat{Set}^{\cat{C}\op}$ (covariant) & $A\mapsto\Hom(-,A)$; \ $f\mapsto\Hom(-,f)\colon\yon(A)\Rightarrow\yon(B)$\\
\midrule
naturalitat de $\Hom(-,f)$ en $X$ & quadrat per a $u\colon X'\to X$ & $\Hom(X',f)\comp\Hom(u,A)=\Hom(u,B)\comp\Hom(X,f)\colon\Hom(X,A)\to\Hom(X',B)$; \ tots dos envien $\varphi$ a $f\comp\varphi\comp u$\\
lema de Yoneda, naturalitat en $P$ & covariant, per a $\sigma\colon P\Rightarrow Q$ & $\Phi_{A,Q}\comp(\sigma\comp-)=\sigma_A\comp\Phi_{A,P}\colon\Nat(\yon(A),P)\to Q(A)$\\
lema de Yoneda, naturalitat en $A$ & contravariant, per a $a\colon A'\to A$ & $\Phi_{A',P}\comp(-\comp\yon(a))=P(a)\comp\Phi_{A,P}\colon\Nat(\yon(A),P)\to P(A')$\\
\bottomrule
\end{tabularx}
\end{center}
La regla que hi ha darrere és sempre la mateixa. A $\Hom(X,Y)$, un morfisme que actua sobre la \emph{segona} variable ho fa per postcomposició i conserva el sentit; un que actua sobre la \emph{primera} ho fa per precomposició i l'inverteix. Per això $\Hom(-,f)$ va en el sentit de $f$ i $\Hom(f,-)$ en el contrari, i per això la naturalitat del lema és covariant en $P$ i contravariant en $A$.
\end{observacio}

\section{Conseqüències}

La primera conseqüència és que la immersió de Yoneda mereix el seu nom.

\begin{corollari}[La immersió de Yoneda és plena i fidel]\index{Yoneda!immersió plena i fidel}
Per a tota parella d'objectes $A,B$ de $\cat{C}$, la immersió $\yon$ dona una bijecció
\[
\Hom(A,B)\;\cong\;\Nat\big(\Hom(-,A),\,\Hom(-,B)\big).
\]
En particular, $\yon\colon\cat{C}\to\cat{Set}^{\cat{C}\op}$ és ple i fidel.
\end{corollari}

\begin{prova}
Apliquem el lema de Yoneda amb $P=\Hom(-,B)$. Aleshores
\[
\Nat\big(\Hom(-,A),\,\Hom(-,B)\big)\cong \Hom(-,B)(A)=\Hom(A,B),
\]
i seguint la bijecció es comprova que és precisament $f\mapsto\yon(f)$. Per tant $\yon$ és bijectiu sobre cada conjunt de morfismes, és a dir, ple i fidel.
\end{prova}

\begin{corollari}\index{Yoneda!objectes determinats per morfismes}
Dos objectes $A$ i $B$ de $\cat{C}$ són isomorfs si i només si els seus functors representables ho són:
\[
A\cong B
\quad\Longleftrightarrow\quad
\Hom(-,A)\cong\Hom(-,B).
\]
\end{corollari}

\begin{prova}
Un functor ple i fidel reflecteix els isomorfismes (capítol dels functors): si $\yon(A)\cong\yon(B)$, aleshores $A\cong B$. El recíproc és que tot functor preserva isomorfismes. Aplicat a $\yon$, dona l'equivalència.
\end{prova}

\begin{observacio}
Aquest darrer corol·lari és la formulació precisa de l'eslògan del primer capítol: \emph{un objecte queda determinat, llevat d'isomorfisme, per la xarxa de morfismes que hi arriben}. Conèixer $\Hom(-,A)$ ---és a dir, saber com és el conjunt de morfismes $X\to A$ per a cada $X$, de manera compatible amb la composició--- equival a conèixer $A$ llevat d'isomorfisme. En particular, l'objecte que representa un functor representable és únic llevat d'isomorfisme, cosa que unifica totes les unicitats de les propietats universals que hem trobat.
\end{observacio}

\begin{observacio}
Hi ha una versió covariant de tot plegat, obtinguda treballant amb $\cat{C}\op$ en lloc de $\cat{C}$: dona una immersió $\cat{C}\op\to\cat{Set}^{\cat{C}}$, també plena i fidel, i el lema $\Nat(\Hom(A,-),Q)\cong Q(A)$ per a functors covariants $Q\colon\cat{C}\to\cat{Set}$. Totes dues versions s'usen contínuament; la contravariant, amb prefeixos, és la que enllaça amb la teoria de feixos i la lògica categòrica que estudiarem més endavant.
\end{observacio}

\begin{resumcap}
\apartat{Idees centrals}
\begin{enumerate}
  \item Les dades universals es caracteritzen per l'existència i la unicitat de morfismes; sovint es poden veure com a objectes inicials o terminals en una categoria coma adequada.
  \item Una representació d'un functor és un objecte representant juntament amb un element universal; equivalentment, és un isomorfisme natural amb un homfunctor.
  \item El lema de Yoneda dona una bijecció natural
  \[
  \Nat\bigl(\Hom(-,A),P\bigr)\;\cong\;P(A),
  \]
  determinada per l'avaluació a $\id_A$.
  \item La immersió de Yoneda és plena i fidel; per tant, els morfismes entre objectes es recuperen exactament com a transformacions naturals entre els seus prefeixos representables.
  \item En particular, un objecte queda determinat, llevat d'isomorfisme, pel sistema de morfismes que hi arriben; la representació completa és única llevat d'un únic isomorfisme compatible amb la dada universal.
\end{enumerate}
\apartat{Recordatori operatiu} Les dues direccions de la bijecció de Yoneda són
\[
\eta\longmapsto\eta_A(\id_A),
\qquad\qquad
x\longmapsto\bigl(f\mapsto P(f)(x)\bigr).
\]
\apartat{Errors típics} Dir \enquote{únic llevat d'isomorfisme únic} de l'objecte nu quan això només val per a la representació amb la dada universal; usar la propietat universal del producte abans d'haver-la definit; i confondre el functor representat amb l'objecte representant.
\apartat{Lectura recomanada} \cite[§2.2 i §§8.1--8.4]{awodey} per als objectes inicials i terminals i per a la immersió, el lema i les aplicacions de Yoneda; \cite[§§4.1--4.3]{leinster}; \cite[cap.~2]{riehl}, d'arquitectura molt propera a la d'aquest capítol; \cite[II.6 i III.1--III.2]{maclane} per a les categories coma, les fletxes universals i el lema de Yoneda; i \cite[cap.~2]{perrone-notes}, per a una segona explicació del lema de Yoneda, amb molts casos particulars i exemples de matemàtica elemental.
\end{resumcap}

\chapter[Límits i colímits]{Límits i colímits}
\label{cap:limits}

\begin{objectius}
\apartat{Prerequisits} Categories oposades i dualitat, i categories sobre un objecte (cap.~\ref{cap:categories}); functors, functors constants i el functor diagonal (cap.~\ref{cap:functors}); transformacions naturals i categories de functors (cap.~\ref{cap:transformacions}); categories coma, propietats universals, representabilitat i lema de Yoneda, especialment la unicitat llevat d'isomorfisme (cap.~\ref{cap:yoneda}).
\apartat{Objectius} En acabar el capítol, el lector hauria de poder:
\begin{itemize}
  \item interpretar un diagrama com un functor i un con com una transformació natural;
  \item definir límit i colímit com a cons i cocons universals;
  \item reconèixer productes, equalitzadors, pullbacks i l'objecte terminal com a límits, i els seus duals com a colímits;
  \item calcular aquests casos bàsics a $\cat{Set}$ i interpretar-los en preordres;
  \item fer servir els pullbacks per descriure la reindexació;
  \item enunciar la caracterització dels límits finits mitjançant objecte terminal, productes binaris i equalitzadors, o bé objecte terminal i pullbacks;
  \item descriure els límits i colímits petits a $\cat{Set}$ i explicar com es calculen punt a punt els límits i colímits en categories de functors;
  \item distingir preservació i reflexió de límits;
  \item explicar per què els representables covariants preserven límits.
\end{itemize}
\end{objectius}

\section{Cap a la noció de límit}

El capítol anterior ha caracteritzat diversos objectes per propietats universals: l'objecte terminal, l'inicial, i de passada el producte i el coproducte. Tots responen a un mateix patró: un objecte que es relaciona de manera òptima amb tots els altres. I tots són únics llevat d'isomorfisme. En aquest capítol unifiquem aquest patró en la noció general de \textbf{límit} d'un diagrama, i la seva dual, el \textbf{colímit}.

La idea és senzilla. Un diagrama és una col·lecció d'objectes i morfismes amb una forma determinada. Un \emph{con} sobre el diagrama és un objecte que projecta cap a tots els seus vèrtexs de manera compatible. El \emph{límit} és el con òptim: aquell pel qual tots els altres factoritzen de manera única. Segons la forma del diagrama, aquesta recepta produeix productes, equalitzadors, pullbacks, objectes terminals i molt més, tots com a casos particulars d'una sola definició. La representabilitat i el lema de Yoneda en seran l'eina conceptual: en el llenguatge del capítol anterior, un límit és un objecte que representa el prefeix que envia cada objecte $X$ al conjunt dels cons de vèrtex $X$ sobre el diagrama.

\section{Diagrames i cons}

Formalitzem la idea de \enquote{diagrama d'una forma donada} amb el llenguatge dels functors, que ja tenim a mà.

\begin{definicio}[Diagrama]\label{def:diagrama}\index{diagrama!de forma $\cat{I}$}\index{categoria!índex}
Sigui $\cat{I}$ una categoria petita, anomenada \textbf{categoria índex}. Un \textbf{diagrama de forma $\cat{I}$} en una categoria $\cat{C}$ és un functor $D\colon\cat{I}\to\cat{C}$. Escrivim $D_i$ per a l'objecte $D(i)$ i, per a un morfisme $u\colon i\to j$ de $\cat{I}$, $D_u\colon D_i\to D_j$ per a la fletxa corresponent.
\end{definicio}

\begin{observacio}
La categoria índex $\cat{I}$ fixa la \emph{forma} del diagrama, no el seu contingut. Amb $\cat{I}$ discreta de dos objectes obtenim una parella d'objectes sense fletxes; amb $\cat{I}$ la categoria de dues fletxes paral·leles, un parell de morfismes amb el mateix domini i codomini; amb $\cat{I}$ buida, el diagrama buit. Cada forma donarà lloc a un tipus de límit.
\end{observacio}

\begin{definicio}[Con]\index{con}
Sigui $D\colon\cat{I}\to\cat{C}$ un diagrama. Un \textbf{con} sobre $D$ amb \textbf{vèrtex} $A$ és una família de morfismes
\[
\alpha_i\colon A\to D_i,\qquad i\in\cat{I},
\]
tal que, per a cada morfisme $u\colon i\to j$ de $\cat{I}$, el triangle commuta:
\[
\begin{tikzpicture}[baseline=(m.center)]
  \matrix (m) [matrix of math nodes, row sep=2.4em, column sep=3.0em] {
     & A & \\
    D_i & & D_j \\
  };
  \draw[-{Stealth[scale=1.1]}] (m-1-2) -- node[above left] {$\alpha_i$} (m-2-1);
  \draw[-{Stealth[scale=1.1]}] (m-1-2) -- node[above right] {$\alpha_j$} (m-2-3);
  \draw[-{Stealth[scale=1.1]}] (m-2-1) -- node[below] {$D_u$} (m-2-3);
\end{tikzpicture}
\qquad
D_u\comp\alpha_i=\alpha_j.
\]
\end{definicio}

\begin{observacio}\label{obs:con-transformacio}
Un con és exactament una transformació natural $\alpha\colon\Delta_A\Rightarrow D$, on $\Delta_A\colon\cat{I}\to\cat{C}$ és el functor constant amb valor $A$ (exemple~\ref{ex:functor-constant}). En efecte, els components $\alpha_i\colon A\to D_i$ i la condició de naturalitat $D_u\comp\alpha_i=\alpha_j$ són precisament la definició de con. Aquest és el primer lloc on la maquinària de les transformacions naturals rendeix directament.
\end{observacio}

\begin{definicio}[Morfisme de cons]
Un \textbf{morfisme de cons} de $(A,\alpha)$ a $(B,\beta)$, tots dos sobre $D$, és un morfisme $f\colon A\to B$ tal que $\beta_i\comp f=\alpha_i$ per a tot $i$. Els cons sobre $D$ i els seus morfismes formen una categoria, que denotem $\mathrm{Con}(D)$.
\end{definicio}

\begin{observacio}[Els cons són objectes d'una categoria coma]\label{obs:cons-coma}\index{con}\index{categoria!coma}\index{functor!diagonal}
Els functors constants $\Delta_A$ s'apleguen en un sol functor, el \textbf{functor diagonal}
\[
\Delta\colon\cat{C}\longrightarrow\cat{C}^{\cat{I}},
\]
que envia cada objecte $A$ al functor constant $\Delta_A$ i cada morfisme $f\colon A\to B$ a la transformació natural $\Delta_f\colon\Delta_A\Rightarrow\Delta_B$ de components $(\Delta_f)_i=f$; la naturalitat és trivial, perquè $\Delta_A$ i $\Delta_B$ envien tots els morfismes de $\cat{I}$ a identitats, i la functorialitat es comprova component a component. Amb aquest functor, la categoria de cons és una categoria coma (secció~\ref{sec:coma}):
\[
\mathrm{Con}(D)=(\Delta\downarrow D).
\]
En efecte, prenent $S=\Delta$ i $T=D\colon\cat{1}\to\cat{C}^{\cat{I}}$, un objecte de $(\Delta\downarrow D)$ és una terna $(A,\ast,\alpha)$ amb $A$ objecte de $\cat{C}$ i $\alpha\colon\Delta_A\Rightarrow D$ un morfisme de $\cat{C}^{\cat{I}}$, que identifiquem amb la parella $(A,\alpha)$, com a l'exemple~\ref{ex:slice-coma}; és a dir, un con sobre $D$ amb vèrtex $A$ (observació~\ref{obs:con-transformacio}). Un morfisme $(A,\alpha)\to(B,\beta)$ és una parella $(f,\id_\ast)$ amb $f\colon A\to B$ tal que el quadrat de la definició de categoria coma commuta, és a dir, $\beta\comp\Delta_f=\alpha$; component a component, $\beta_i\comp f=\alpha_i$ per a tot $i$, que és exactament la condició de morfisme de cons.

Dualment, un cocon sobre $D$ amb vèrtex $A$ (definició~\ref{def:cocon}) és una transformació natural $D\Rightarrow\Delta_A$, i els cocons formen la categoria coma $(D\downarrow\Delta)$. Així, el límit és un objecte terminal de $(\Delta\downarrow D)$ i el colímit un objecte inicial de $(D\downarrow\Delta)$: és l'eslògan de l'exemple~\ref{ex:universal-coma}, \emph{ser universal és ser inicial o terminal en una categoria coma adient}. La categoria coma dona, doncs, el llenguatge comú per a les categories sobre i sota un objecte, els objectes universals del capítol~\ref{cap:yoneda} i els límits i colímits.
\end{observacio}

\section{Límits}

\begin{definicio}[Límit]\index{límit}
\label{def:limit}
Un \textbf{límit} del diagrama $D\colon\cat{I}\to\cat{C}$ és un objecte terminal de la categoria de cons $\mathrm{Con}(D)$. Explícitament, és un con $(L,\lambda)$ tal que, per a tot con $(A,\alpha)$ sobre $D$, hi ha un \emph{únic} morfisme $f\colon A\to L$ amb
\[
\lambda_i\comp f=\alpha_i\qquad\text{per a tot }i.
\]
Escrivim $L=\varprojlim D$ o $L=\lim D$.
\end{definicio}

\begin{observacio}
Un límit és, doncs, el con \enquote{universal}: tots els altres cons factoritzen a través d'ell de manera única. Com que és un objecte terminal de $\mathrm{Con}(D)$, el con límit és \textbf{únic llevat d'un únic isomorfisme} de cons, és a dir, compatible amb les projeccions, per la unicitat dels objectes terminals (proposició~\ref{prop:unicitat-terminal}); el vèrtex sol és únic llevat d'isomorfisme. Direm que una categoria \textbf{té límits de forma $\cat{I}$} quan tot diagrama de forma $\cat{I}$ hi té límit.
\end{observacio}

La propietat universal del límit es pot formular, en el llenguatge del capítol~\ref{cap:yoneda}, com una representabilitat. Suposem $\cat{I}$ petita i $\cat{C}$ localment petita, de manera que, per a cada objecte $A$, els cons sobre $D$ amb vèrtex $A$ formen un conjunt
\[
\mathrm{Con}(A,D)=\Nat(\Delta_A,D).
\]
Un morfisme $g\colon A'\to A$ transforma cada con $\alpha$ amb vèrtex $A$ en el con $\alpha\comp\Delta_g$ amb vèrtex $A'$, de components $\alpha_i\comp g$; obtenim així un prefeix
\[
\mathrm{Con}(-,D)\colon\cat{C}\op\longrightarrow\cat{Set}.
\]

\begin{proposicio}[Els límits com a representacions]\label{prop:limit-representable}\index{límit!com a representació}
Un con $(L,\lambda)$ sobre $D$ és un límit si i només si les aplicacions
\[
\Hom(A,L)\longrightarrow\mathrm{Con}(A,D),
\qquad
f\longmapsto\lambda\comp\Delta_f=(\lambda_i\comp f)_i,
\]
formen un isomorfisme natural $\Hom(-,L)\cong\mathrm{Con}(-,D)$. Dit d'una altra manera, un límit de $D$ és una representació $(L,\lambda)$ del prefeix $\mathrm{Con}(-,D)$, amb element universal $\lambda\in\mathrm{Con}(L,D)$.
\end{proposicio}

\begin{prova}
Les aplicacions són naturals en $A$: per a $g\colon A'\to A$, enviar primer $f$ a $(\lambda_i\comp f)_i$ i precompondre després amb $g$ dona $(\lambda_i\comp f\comp g)_i$, que és la imatge de $f\comp g$. Una transformació natural és un isomorfisme natural exactament quan totes les seves components són bijectives, i que l'aplicació corresponent a $A$ sigui bijectiva diu exactament que, per a tot con $\alpha$ amb vèrtex $A$, hi ha un únic $f\colon A\to L$ amb $\lambda_i\comp f=\alpha_i$ per a tot $i$. Això és la definició de límit. L'element universal és la imatge de $\id_L$, que és $\lambda$.
\end{prova}

Aquesta formulació explica d'un sol cop tres fets del capítol. Un límit és una propietat universal en el sentit del capítol~\ref{cap:yoneda}; és, per tant, únic llevat d'un únic isomorfisme compatible amb la dada universal, que aquí és el con $\lambda$; i la mateixa idea mostrarà que els representables preserven límits (proposició~\ref{prop:representables-limits}).

Recorrem ara les formes particulars, cadascuna amb la seva categoria índex.

\subsection{Producte}\index{producte}

Prenem $\cat{I}$ la categoria \textbf{discreta} amb dos objectes $1,2$ (sense fletxes no trivials). Un diagrama és simplement una parella d'objectes $A,B$. Un con amb vèrtex $X$ és una parella de morfismes $f\colon X\to A$, $g\colon X\to B$, sense cap condició (no hi ha fletxes a $\cat{I}$ que imposin res). El límit és el \textbf{producte}.

\begin{definicio}\index{producte}
\label{def:producte}
El \textbf{producte} de dos objectes $A,B$ és un objecte $A\times B$ amb dues projeccions $\pi_1\colon A\times B\to A$ i $\pi_2\colon A\times B\to B$ tals que, per a tot objecte $X$ i tota parella de morfismes $f\colon X\to A$, $g\colon X\to B$, hi ha un únic morfisme $\langle f,g\rangle\colon X\to A\times B$ amb
\[
\pi_1\comp\langle f,g\rangle=f,\qquad \pi_2\comp\langle f,g\rangle=g.
\]
\[
\begin{tikzpicture}[baseline=(m.center)]
  \matrix (m) [matrix of math nodes, column sep=2.8em, row sep=2.6em] {
      & X & \\
    A & A\times B & B \\
  };
  \draw[-{Stealth[scale=1.1]}] (m-1-2) -- node[above left] {$f$} (m-2-1);
  \draw[-{Stealth[scale=1.1]}] (m-1-2) -- node[above right] {$g$} (m-2-3);
  \draw[-{Stealth[scale=1.1]},dashed] (m-1-2) -- node[right] {$\langle f,g\rangle$} (m-2-2);
  \draw[-{Stealth[scale=1.1]}] (m-2-2) -- node[below] {$\pi_1$} (m-2-1);
  \draw[-{Stealth[scale=1.1]}] (m-2-2) -- node[below] {$\pi_2$} (m-2-3);
\end{tikzpicture}
\]
\end{definicio}

\begin{exemple}
A $\cat{Set}$, el producte és el producte cartesià $A\times B=\{(a,b)\mid a\in A,\ b\in B\}$ amb les projeccions habituals. A $\cat{Grp}$ és el producte directe de grups; a $\cat{Top}$, l'espai producte amb la topologia producte; en un preordre, l'ínfim dels dos elements, quan existeix. En cada cas, la propietat universal recupera la construcció coneguda.
\end{exemple}

\subsection{Equalitzador}\index{equalitzador}

Prenem $\cat{I}$ la categoria amb dues fletxes paral·leles $\bullet\rightrightarrows\bullet$. Un diagrama és un parell de morfismes $f,g\colon A\to B$. Un con amb vèrtex $X$ és un morfisme $e\colon X\to A$ amb $f\comp e=g\comp e$ (la condició prové de les dues fletxes). El límit és l'\textbf{equalitzador}.

\begin{definicio}\index{equalitzador}
\label{def:equalitzador}
L'\textbf{equalitzador} de $f,g\colon A\to B$ és un objecte $E$ amb un morfisme $e\colon E\to A$ tal que $f\comp e=g\comp e$ i que és universal amb aquesta propietat: per a tot $h\colon X\to A$ amb $f\comp h=g\comp h$, hi ha un únic $k\colon X\to E$ amb $e\comp k=h$.
\[
\begin{tikzpicture}[baseline=(m.center)]
  \matrix (m) [matrix of math nodes, column sep=3em, row sep=2.6em] {
    E & A & B \\
    X & & \\
  };
  \draw[-{Stealth[scale=1.1]}] (m-1-1) -- node[above] {$e$} (m-1-2);
  \draw[-{Stealth[scale=1.1]}] ([yshift=2pt]m-1-2.east) -- node[above] {$f$} ([yshift=2pt]m-1-3.west);
  \draw[-{Stealth[scale=1.1]}] ([yshift=-2pt]m-1-2.east) -- node[below] {$g$} ([yshift=-2pt]m-1-3.west);
  \draw[-{Stealth[scale=1.1]}] (m-2-1) -- node[left] {$h$} (m-1-2);
  \draw[-{Stealth[scale=1.1]},dashed] (m-2-1) -- node[left] {$k$} (m-1-1);
\end{tikzpicture}
\]
\end{definicio}

\begin{exemple}
A $\cat{Set}$, l'equalitzador de $f,g\colon A\to B$ és el subconjunt $E=\{a\in A\mid f(a)=g(a)\}$ amb la injecció $E\hookrightarrow A$. En general, els equalitzadors representen els subobjectes definits per una equació. Tot equalitzador és un monomorfisme.
\end{exemple}

\subsection{Pullback}\index{pullback}

Prenem $\cat{I}$ la categoria amb forma de racó, $\bullet\to\bullet\leftarrow\bullet$. Un diagrama és un parell de morfismes amb codomini comú, $f\colon A\to C$ i $g\colon B\to C$. El límit és el \textbf{pullback}.

\begin{definicio}\index{pullback}
\label{def:pullback}
El \textbf{pullback} de $f\colon A\to C$ i $g\colon B\to C$ és un objecte $P$ amb morfismes $p_1\colon P\to A$ i $p_2\colon P\to B$ tals que $f\comp p_1=g\comp p_2$ i que és universal: per a tot $X$ amb $r_1\colon X\to A$, $r_2\colon X\to B$ i $f\comp r_1=g\comp r_2$, hi ha un únic $h\colon X\to P$ amb $p_1\comp h=r_1$ i $p_2\comp h=r_2$.
\[
\begin{tikzpicture}[baseline=(m.center)]
  \matrix (m) [matrix of math nodes, column sep=3em, row sep=2.6em] {
    P & B \\
    A & C \\
  };
  \draw[-{Stealth[scale=1.1]}] (m-1-1) -- node[above] {$p_2$} (m-1-2);
  \draw[-{Stealth[scale=1.1]}] (m-1-1) -- node[left] {$p_1$} (m-2-1);
  \draw[-{Stealth[scale=1.1]}] (m-1-2) -- node[right] {$g$} (m-2-2);
  \draw[-{Stealth[scale=1.1]}] (m-2-1) -- node[below] {$f$} (m-2-2);
  \node at ([xshift=0.75em,yshift=-0.75em]m-1-1) {$\scriptstyle\lrcorner$};
\end{tikzpicture}
\]
\end{definicio}

\begin{exemple}
A $\cat{Set}$, el pullback és $P=\{(a,b)\in A\times B\mid f(a)=g(b)\}$ amb les projeccions. Quan $g\colon B\hookrightarrow C$ és la inclusió d'un subconjunt $B\subseteq C$, el pullback recupera l'antiimatge $f^{-1}(B)$, mitjançant la bijecció $(a,b)\mapsto a$. En aquest cas, doncs, el pullback categoritza l'operació d'antiimatge. En general, el pullback també s'anomena \emph{producte fibrat} i s'escriu $A\times_C B$.
\end{exemple}

\begin{observacio}[Reindexació per pullback]\index{reindexació}\index{functor!de reindexació}
\label{obs:reindexacio}
Ara que disposem de pullbacks podem definir la reindexació entre categories sobre un objecte. Suposem que $\cat{C}$ té tots els pullbacks i sigui $f\colon A\to B$ un morfisme. Definim un functor
\[
f^{*}\colon\cat{C}/B\longrightarrow\cat{C}/A
\]
de la manera següent: a un objecte $y\colon Y\to B$ de $\cat{C}/B$ li assignem el pullback de $y$ al llarg de $f$, és a dir, l'objecte $f^{*}Y=A\times_B Y$ amb la seva projecció $p_1\colon A\times_B Y\to A$, que és un objecte de $\cat{C}/A$:
\[
\begin{tikzpicture}[baseline=(m.center)]
  \matrix (m) [matrix of math nodes, column sep=3em, row sep=2.6em] {
    A\times_B Y & Y \\
    A & B \\
  };
  \draw[-{Stealth[scale=1.1]}] (m-1-1) -- node[above] {$p_2$} (m-1-2);
  \draw[-{Stealth[scale=1.1]}] (m-1-1) -- node[left] {$p_1=f^{*}y$} (m-2-1);
  \draw[-{Stealth[scale=1.1]}] (m-1-2) -- node[right] {$y$} (m-2-2);
  \draw[-{Stealth[scale=1.1]}] (m-2-1) -- node[below] {$f$} (m-2-2);
  \node at ([xshift=0.9em,yshift=-0.8em]m-1-1) {$\scriptstyle\lrcorner$};
\end{tikzpicture}
\]
Sobre morfismes, la propietat universal del pullback assigna a cada triangle de $\cat{C}/B$ un únic triangle de $\cat{C}/A$. Un cop \textbf{triats} els pullbacks que defineixen $f^{*}$, obtenim un functor ordinari $f^{*}\colon\cat{C}/B\to\cat{C}/A$. Per a morfismes composables $A\xrightarrow{f}B\xrightarrow{g}C$, els functors $(g\comp f)^{*}$ i $f^{*}\comp g^{*}$ no són literalment iguals per a una tria arbitrària de pullbacks, però són canònicament isomorfs. Aquesta coherència s'organitza, en un llenguatge més avançat, mitjançant la noció de \emph{pseudofunctor}, que no necessitarem aquí. En passar als preordres de subobjectes (capítol~\ref{cap:subobjectes}), aquests isomorfismes queden identificats i la reindexació és estrictament functorial:
\[
(g\comp f)^{*}=f^{*}\comp g^{*},\qquad (\id_A)^{*}=\id_{\Sub(A)}.
\]
A $\cat{Set}$, si interpretem $y\colon Y\to B$ com la família de fibres $\{y^{-1}(b)\}_{b\in B}$, aleshores $f^{*}Y$ s'interpreta com la família reindexada $\{y^{-1}(f(a))\}_{a\in A}$: exactament la \textbf{substitució} al llarg de $f$. Aquesta és la lectura lògica que explotarem en tractar els subobjectes, on $f^{*}$ es restringeix a un morfisme de preordres $\Sub(B)\to\Sub(A)$ i s'interpreta com la substitució en un predicat.
\end{observacio}

\subsection{Objecte terminal com a límit}

Prenem $\cat{I}=\cat{0}$, la categoria buida. L'únic diagrama de forma $\cat{0}$ és el diagrama buit. Un con sobre el diagrama buit és, simplement, un objecte $X$ (sense cap component, perquè no hi ha vèrtexs). El límit ---el con terminal--- és, doncs, l'\textbf{objecte terminal} de $\cat{C}$. Així, l'objecte terminal és el límit del diagrama buit, i queda unificat amb els altres.

\section{Colímits}

Tota la teoria anterior es dualitza girant les fletxes. Un \textbf{colímit} d'un diagrama $D\colon\cat{I}\to\cat{C}$ és un límit del diagrama oposat $D\op\colon\cat{I}\op\to\cat{C}\op$; equivalentment, és un objecte inicial a la categoria de \textbf{cocons} sobre $D$.

\begin{definicio}[Cocon i colímit]\label{def:cocon}\index{cocon}\index{colímit}
Un \textbf{cocon} sobre $D$ amb vèrtex $A$ és una família $\alpha_i\colon D_i\to A$ tal que, per a cada $u\colon i\to j$, commuta el triangle dual
\[
\begin{tikzpicture}[baseline=(m.center)]
  \matrix (m) [matrix of math nodes, row sep=2.4em, column sep=3.0em] {
    D_i & & D_j \\
     & A & \\
  };
  \draw[-{Stealth[scale=1.1]}] (m-1-1) -- node[below left] {$\alpha_i$} (m-2-2);
  \draw[-{Stealth[scale=1.1]}] (m-1-3) -- node[below right] {$\alpha_j$} (m-2-2);
  \draw[-{Stealth[scale=1.1]}] (m-1-1) -- node[above] {$D_u$} (m-1-3);
\end{tikzpicture}
\qquad
\alpha_j\comp D_u=\alpha_i.
\]
Un \textbf{colímit} és un cocon $(L,\lambda)$ universal: per a tot cocon $(A,\alpha)$ hi ha un únic $f\colon L\to A$ amb $f\comp\lambda_i=\alpha_i$. Escrivim $L=\colimop D$ o $L=\varinjlim D$.
\end{definicio}

Dualitzant cada límit obtenim el colímit corresponent:

\begin{itemize}
\item El \textbf{coproducte} $A+B$ (dual del producte) té injeccions $\iota_1\colon A\to A+B$, $\iota_2\colon B\to A+B$ universals. A $\cat{Set}$ és la unió disjunta; a $\cat{Grp}$, el producte lliure; en un preordre, el suprem.
\item El \textbf{coequalitzador} de $f,g\colon A\to B$ (dual de l'equalitzador) és un $q\colon B\to Q$ universal amb $q\comp f=q\comp g$. A $\cat{Set}$ és el quocient de $B$ per la relació d'equivalència generada per $f(a)\sim g(a)$. Tot coequalitzador és un epimorfisme.
\item El \textbf{pushout} de $f\colon C\to A$ i $g\colon C\to B$ (dual del pullback) és un $Q$ universal amb el quadrat commutatiu. A $\cat{Set}$ és la unió de $A$ i $B$ enganxada al llarg de $C$, i s'escriu $A+_C B$.
\item L'\textbf{objecte inicial} (dual del terminal) és el colímit del diagrama buit.
\end{itemize}

\begin{observacio}
No cal demostrar res de nou per als colímits: cada enunciat sobre límits a $\cat{C}$ es tradueix en un enunciat sobre colímits a $\cat{C}\op$, i viceversa. Aquesta és la dualitat en acció, aplicada de manera sistemàtica. En endavant, quan enunciem un resultat per a límits, el dual per a colímits s'entén sobreentès.
\end{observacio}

\section{Límits finits}

\begin{definicio}\index{límit!finit}\index{categoria!finitament completa}
Un límit es diu \textbf{finit} quan la seva categoria índex $\cat{I}$ és finita, és a dir, quan té un nombre finit d'objectes \emph{i} un nombre finit de morfismes. (No n'hi ha prou que els objectes siguin finits: una categoria d'un sol objecte pot tenir infinits endomorfismes.) Una categoria és \textbf{finitament completa} (o \textbf{lex}) si té tots els límits finits; és \textbf{completa} si té tots els límits de forma petita.
\end{definicio}

El resultat clau redueix tots els límits finits a dos casos bàsics: productes finits i equalitzadors. Això és molt útil, perquè per comprovar que una categoria és finitament completa n'hi ha prou de verificar dues coses.

\begin{teorema}\index{límit!finit!construcció}\label{teo:limits-finits}
Per a una categoria $\cat{C}$, són equivalents:
\begin{enumerate}
\item $\cat{C}$ té tots els límits finits;
\item $\cat{C}$ té objecte terminal, productes binaris i equalitzadors;
\item $\cat{C}$ té objecte terminal i pullbacks.
\end{enumerate}
\end{teorema}

\begin{prova}
\emph{(1) implica (2) i (3).} L'objecte terminal, el producte binari, l'equalitzador i el pullback són límits de diagrames de formes finites: la categoria buida, la discreta de dos objectes, $\bullet\rightrightarrows\bullet$ i $\bullet\to\bullet\leftarrow\bullet$.

\emph{(2) implica (1).} Primer, $\cat{C}$ té productes de qualsevol família finita. El de la família buida és l'objecte terminal, i per inducció, si $P_n$ és un producte de $A_1,\dots,A_n$ amb projeccions $\pi_k$, aleshores $P_n\times A_{n+1}$, amb projeccions $\pi_k\comp p_1$ ($k\le n$) i $p_2$, és un producte de $A_1,\dots,A_{n+1}$. En efecte, donades $f_k\colon X\to A_k$, hi ha una única $f\colon X\to P_n$ amb $\pi_k\comp f=f_k$ per a $k\le n$, i una única $X\to P_n\times A_{n+1}$ amb components $f$ i $f_{n+1}$; tota fletxa amb les mateixes composicions amb les $n+1$ projeccions té necessàriament aquestes dues components.

Sigui ara $D\colon\cat{I}\to\cat{C}$ un diagrama finit, i escrivim $\cat{I}_0$ i $\cat{I}_1$ per als conjunts d'objectes i de morfismes de $\cat{I}$. Considerem els productes finits
\[
P=\prod_{i\in\cat{I}_0}D_i
\qquad\text{i}\qquad
Q=\prod_{u\in\cat{I}_1}D_{\mathrm{cod}(u)},
\]
amb projeccions $\pi_i\colon P\to D_i$ i $q_u\colon Q\to D_{\mathrm{cod}(u)}$, i les dues fletxes $s,t\colon P\to Q$ determinades, per a cada $u\colon i\to j$, per
\[
q_u\comp s=\pi_j,
\qquad\qquad
q_u\comp t=D_u\comp\pi_i.
\]
Sigui $e\colon L\to P$ un equalitzador de $s$ i $t$, i posem $\lambda_i=\pi_i\comp e\colon L\to D_i$. Vegem que $(L,\lambda)$ és un límit de $D$.

\emph{És un con.} Per a $u\colon i\to j$,
\[
D_u\comp\lambda_i=D_u\comp\pi_i\comp e=q_u\comp t\comp e=q_u\comp s\comp e=\pi_j\comp e=\lambda_j.
\]

\emph{És universal.} Sigui $(A,\alpha)$ un con sobre $D$, i sigui $a\colon A\to P$ l'única fletxa amb $\pi_i\comp a=\alpha_i$ per a tot $i$. Per a cada $u\colon i\to j$,
\[
q_u\comp s\comp a=\pi_j\comp a=\alpha_j,
\qquad
q_u\comp t\comp a=D_u\comp\pi_i\comp a=D_u\comp\alpha_i=\alpha_j,
\]
de manera que $s\comp a$ i $t\comp a$ tenen les mateixes components i, per la unicitat del producte $Q$, $s\comp a=t\comp a$. Per la propietat de l'equalitzador hi ha una única $h\colon A\to L$ amb $e\comp h=a$, i aleshores $\lambda_i\comp h=\pi_i\comp e\comp h=\alpha_i$. Si $h'\colon A\to L$ també compleix $\lambda_i\comp h'=\alpha_i$ per a tot $i$, aleshores $\pi_i\comp(e\comp h')=\alpha_i$, d'on $e\comp h'=a$ per la unicitat del producte $P$, i $h'=h$ per la de l'equalitzador. El diagrama buit hi queda inclòs: aleshores $P$ i $Q$ són productes buits, és a dir, objectes terminals, i $L$ és terminal.

\emph{(3) implica (2).} Sigui $1$ l'objecte terminal i, per a cada objecte $A$, $!_A\colon A\to 1$ l'única fletxa. \emph{Productes binaris}: un pullback de $!_A$ i $!_B$ és un producte de $A$ i $B$. En efecte, per a qualsevol parell $f\colon X\to A$, $g\colon X\to B$, la condició de con $!_A\comp f=\,!_B\comp g$ es compleix automàticament, perquè totes dues fletxes van a $1$; per tant, la propietat universal del pullback diu exactament la del producte.

\emph{Equalitzadors}: siguin $f,g\colon A\to B$. Ja sabem que hi ha productes binaris; formem el pullback de
\[
\langle\id_A,f\rangle,\ \langle\id_A,g\rangle\colon A\to A\times B,
\]
amb projeccions $p_1,p_2\colon P\to A$, de manera que $\langle\id_A,f\rangle\comp p_1=\langle\id_A,g\rangle\comp p_2$. Component amb la primera projecció del producte, $p_1=p_2$; anomenem $e$ aquesta fletxa. Component amb la segona, $f\comp e=g\comp e$. Si $h\colon X\to A$ compleix $f\comp h=g\comp h$, aleshores
\[
\langle\id_A,f\rangle\comp h=\langle h,f\comp h\rangle=\langle h,g\comp h\rangle=\langle\id_A,g\rangle\comp h,
\]
de manera que $(h,h)$ és un con sobre el pullback, i hi ha una única $k\colon X\to P$ amb $p_1\comp k=h=p_2\comp k$, és a dir, amb $e\comp k=h$. Si $k'$ també compleix $e\comp k'=h$, aleshores $p_1\comp k'=p_2\comp k'=h$ i $k'=k$. Per tant $e$ és un equalitzador de $f$ i $g$.

Amb les implicacions (3)$\Rightarrow$(2)$\Rightarrow$(1)$\Rightarrow$(3), les tres condicions són equivalents.
\end{prova}

\begin{observacio}
Dualment, una categoria és \textbf{finitament cocompleta} si té objecte inicial, coproductes binaris i coequalitzadors, o equivalentment objecte inicial i pushouts. Per a la lògica categòrica, els límits finits ---i molt especialment els \emph{pullbacks}--- són centrals: serveixen per interpretar la substitució i les operacions sobre subobjectes que veurem més endavant.
\end{observacio}

\begin{exemple}[Límits en un preordre i en la deduïbilitat]\label{ex:limits-preordre}\index{límit!en un preordre}
Sigui $P$ un preordre vist com a categoria i $D\colon\cat{I}\to P$ un diagrama. Com que entre dos objectes hi ha com a màxim una fletxa, un con amb vèrtex $x$ diu simplement que $x\le D_i$ per a tot $i$, i totes les condicions de commutativitat es compleixen automàticament. Per tant, quan existeixen,
\[
\varprojlim D=\bigwedge_{i\in\cat{I}_0}D_i,
\qquad\qquad
\varinjlim D=\bigvee_{i\in\cat{I}_0}D_i:
\]
el límit és l'ínfim dels objectes del diagrama, i el colímit, el suprem; les fletxes del diagrama no hi intervenen. En particular, l'equalitzador de dues fletxes paral·leles $x\rightrightarrows y$ és $x$ mateix (les dues fletxes coincideixen), i el pullback de $a\to c\leftarrow b$ és l'ínfim $a\wedge b$, que no depèn de $c$. Així, un preordre té tots els límits finits si i només si té element màxim (l'ínfim de la família buida, que és l'objecte terminal) i ínfims binaris.

A la categoria prima $\cat{Prov}_T$ d'un sistema deductiu $T$, això diu que, quan el sistema disposa de $\top$ i de conjunció, $\top$ és un objecte terminal ($p\vdash_T\top$ per a tota fórmula $p$) i $p\land q$ és un producte de $p$ i $q$ (exemple~\ref{ex:rep-producte}). Aquestes dues operacions donen, doncs, tots els límits finits en el context deductiu: la part \enquote{conjuntiva} de la lògica és exactament la dels límits finits.
\end{exemple}

\section{Límits a \texorpdfstring{$\cat{Set}$}{Set} i en categories de functors}

\begin{exemple}[Límits a $\cat{Set}$]\label{ex:limits-set}\index{límit!a Set@a $\cat{Set}$}
Sigui $D\colon\cat{I}\to\cat{Set}$ un diagrama amb $\cat{I}$ petita. El conjunt de les famílies compatibles
\[
L=\Bigl\{(x_i)_{i}\in\textstyle\prod_{i\in\cat{I}_0}D_i\ \Bigm|\ D_u(x_i)=x_j\ \text{per a tot } u\colon i\to j\Bigr\},
\]
amb les projeccions $\lambda_i\colon L\to D_i$, és un límit de $D$. En efecte, un con $(A,\alpha)$ envia cada $a\in A$ a la família $(\alpha_i(a))_i$, que és compatible per la condició de con; l'aplicació $a\mapsto(\alpha_i(a))_i$ és l'única $f\colon A\to L$ amb $\lambda_i\comp f=\alpha_i$. En particular, $\cat{Set}$ té tots els límits petits. Dualment, el colímit és el quocient de la unió disjunta $\coprod_i D_i$, els elements de la qual són parelles $(i,x)$ amb $x\in D_i$, per la relació d'equivalència generada per
\[
(i,x)\sim\bigl(j,D_u(x)\bigr)\qquad(u\colon i\to j);
\]
per tant, $\cat{Set}$ té també tots els colímits petits. A $\cat{Set}$, doncs, els límits imposen compatibilitats i els colímits enganxen dades identificant-les.
\end{exemple}

Els límits en categories de functors es calculen a partir dels límits de la categoria d'arribada, objecte per objecte.

\begin{proposicio}[Límits punt a punt]\label{prop:limits-punt-a-punt}\index{límit!en categories de functors}\index{categoria!de functors!límits}
Siguin $\cat{A}$ i $\cat{I}$ categories petites. Si $\cat{C}$ té límits de forma $\cat{I}$, aleshores la categoria de functors $\cat{C}^{\cat{A}}$ també en té, i es calculen \emph{punt a punt}: per a un diagrama $F\colon\cat{I}\to\cat{C}^{\cat{A}}$,
\[
\Bigl(\varprojlim F\Bigr)(a)\;\cong\;\varprojlim_{i\in\cat{I}}F(i)(a)
\qquad(a\in\cat{A}).
\]
Dualment, si $\cat{C}$ té colímits de forma $\cat{I}$, els colímits a $\cat{C}^{\cat{A}}$ es calculen punt a punt. Equivalentment, per a cada objecte $a$ de $\cat{A}$, el functor d'avaluació
\[
\mathrm{ev}_a\colon\cat{C}^{\cat{A}}\longrightarrow\cat{C},\qquad F\longmapsto F(a),
\]
preserva aquests límits i colímits.
\end{proposicio}
Vegeu també \cite[§3.3]{riehl} i, en el context de les categories de diagrames, \cite[§§8.5--8.6]{awodey}.

\begin{prova}
Recordem que un objecte de $\cat{C}^{\cat{A}}$ és un functor $\cat{A}\to\cat{C}$, i que un morfisme és una transformació natural, amb la composició calculada component a component. Així, el diagrama $F$ assigna a cada objecte $i$ de $\cat{I}$ un functor $F(i)\colon\cat{A}\to\cat{C}$, i a cada $u\colon i\to j$ una transformació natural $F(u)\colon F(i)\Rightarrow F(j)$, de components $F(u)_a$. Per a cada objecte $a$ de $\cat{A}$, el diagrama avaluat
\[
F_a=\mathrm{ev}_a\comp F\colon\cat{I}\to\cat{C},\qquad i\mapsto F(i)(a),\quad u\mapsto F(u)_a,
\]
és un functor, perquè la composició i les identitats de $\cat{C}^{\cat{A}}$ es calculen component a component.

\emph{Pas 1: el límit sobre objectes.} Per a cada objecte $a$ de $\cat{A}$, triem un límit $\bigl(L(a),(\lambda_i(a))_i\bigr)$ de $F_a$; per tant, $F(u)_a\comp\lambda_i(a)=\lambda_j(a)$ per a cada $u\colon i\to j$.

\emph{Pas 2: el límit sobre morfismes.} Sigui $v\colon a\to a'$ un morfisme de $\cat{A}$. La família $\bigl(F(i)(v)\comp\lambda_i(a)\bigr)_i$, de fletxes $L(a)\to F(i)(a')$, és un con sobre $F_{a'}$: per a $u\colon i\to j$, la naturalitat de $F(u)$ i la condició de con de $\lambda(a)$ donen
\[
F(u)_{a'}\comp F(i)(v)\comp\lambda_i(a)=F(j)(v)\comp F(u)_a\comp\lambda_i(a)=F(j)(v)\comp\lambda_j(a).
\]
Per la universalitat de $L(a')$, hi ha una única $L(v)\colon L(a)\to L(a')$ tal que
\[
\lambda_i(a')\comp L(v)=F(i)(v)\comp\lambda_i(a)\qquad\text{per a tot } i.
\tag{$\dagger$}
\]

\emph{Pas 3: $L$ és un functor.} Fem servir que dues fletxes cap a un límit que tenen les mateixes composicions amb totes les projeccions són iguals, per la unicitat de la factorització. Per a $v=\id_a$, tant $L(\id_a)$ com $\id_{L(a)}$ compleixen $(\dagger)$, perquè $F(i)(\id_a)=\id$; per tant són iguals. Per a $v\colon a\to a'$ i $w\colon a'\to a''$, aplicant $(\dagger)$ dues vegades i la functorialitat de $F(i)$,
\[
\begin{aligned}
\lambda_i(a'')\comp L(w)\comp L(v)&=F(i)(w)\comp\lambda_i(a')\comp L(v)\\
&=F(i)(w)\comp F(i)(v)\comp\lambda_i(a)=F(i)(w\comp v)\comp\lambda_i(a),
\end{aligned}
\]
de manera que $L(w)\comp L(v)$ compleix la condició $(\dagger)$ que caracteritza $L(w\comp v)$, i $L(w)\comp L(v)=L(w\comp v)$.

\emph{Pas 4: un con a $\cat{C}^{\cat{A}}$.} La condició $(\dagger)$ és exactament el quadrat de naturalitat de la família $\lambda_i=(\lambda_i(a))_a$; per tant, cada $\lambda_i$ és una transformació natural $L\Rightarrow F(i)$. A més, $F(u)\comp\lambda_i=\lambda_j$ a $\cat{C}^{\cat{A}}$, perquè aquesta igualtat de transformacions naturals es comprova component a component i ho diu el pas~1. Així, $(L,\lambda)$ és un con sobre $F$.

\emph{Pas 5: universalitat.} Sigui $(G,\alpha)$ un con sobre $F$, amb $\alpha_i\colon G\Rightarrow F(i)$ i $F(u)\comp\alpha_i=\alpha_j$. Avaluant en $a$, la família $\bigl((\alpha_i)_a\bigr)_i$ és un con sobre $F_a$ amb vèrtex $G(a)$, de manera que hi ha una única $\theta_a\colon G(a)\to L(a)$ amb $\lambda_i(a)\comp\theta_a=(\alpha_i)_a$ per a tot $i$. La família $\theta$ és natural: per a $v\colon a\to a'$, per $(\dagger)$ i per la naturalitat de $\alpha_i$,
\[
\lambda_i(a')\comp L(v)\comp\theta_a=F(i)(v)\comp(\alpha_i)_a=(\alpha_i)_{a'}\comp G(v)=\lambda_i(a')\comp\theta_{a'}\comp G(v)
\]
per a tot $i$, i per tant $L(v)\comp\theta_a=\theta_{a'}\comp G(v)$. Així $\theta\colon G\Rightarrow L$ és un morfisme de $\cat{C}^{\cat{A}}$ amb $\lambda_i\comp\theta=\alpha_i$. És l'únic: si $\theta'$ també ho compleix, a cada $a$ es té $\lambda_i(a)\comp\theta'_a=(\alpha_i)_a$, i la unicitat de la factorització a través de $L(a)$ dona $\theta'_a=\theta_a$. Per tant $(L,\lambda)$ és un límit de $F$.

\emph{Pas 6: el càlcul punt a punt.} Per construcció, $\mathrm{ev}_a(L,\lambda)=\bigl(L(a),\lambda(a)\bigr)$ és un límit de $F_a$. Si $(L',\lambda')$ és un altre límit de $F$, hi ha un isomorfisme $\varphi\colon L'\Rightarrow L$ amb $\lambda_i\comp\varphi=\lambda'_i$, i la seva component $\varphi_a$ és un isomorfisme amb $\lambda_i(a)\comp\varphi_a=\lambda'_i(a)$. Un con isomorf a un con límit, per un isomorfisme compatible amb les projeccions, també és un límit, perquè les factoritzacions a través de l'un i de l'altre es corresponen component amb $\varphi_a$ o amb $\varphi_a^{-1}$. Per tant $\bigl(L'(a),\lambda'(a)\bigr)$ és un límit de $F_a$: $\mathrm{ev}_a$ preserva els límits de forma $\cat{I}$, i $\bigl(\varprojlim F\bigr)(a)\cong\varprojlim_i F(i)(a)$.

L'argument per als colímits és el mateix amb totes les fletxes invertides: cocons en lloc de cons, i la propietat universal del colímit en lloc de la del límit.
\end{prova}

\begin{exemple}[Un mateix producte a $\cat{Set}$, a $\cat{Graph}$ i a $\cat{Set}^{\cat{A}}$]\label{ex:producte-tres-contextos}\index{producte!a Set, Graph i categories de functors}\index{objecte terminal!a Graph}
Fixem un sol diagrama abstracte: la parella d'objectes $(X,T)$, sobre la categoria discreta de dos objectes, juntament amb el diagrama buit, que dona l'objecte terminal. Calculem-ne el límit en tres llocs.

\emph{A $\cat{Set}$.} El producte és $X\times T$, i l'objecte terminal, un conjunt d'un sol element $\{\ast\}$. Amb $T=\{\ast\}$, l'aplicació $(x,\ast)\mapsto x$ és una bijecció $X\times\{\ast\}\cong X$. A més, les fletxes $\{\ast\}\to X$ són exactament els elements de $X$.

\emph{A $\cat{Graph}$.} Recordem que $\cat{Graph}\cong\cat{Set}^{\cat{P}}$, on $\cat{P}$ té dos objectes $E,V$ i dues fletxes $s,t\colon E\to V$. Per la proposició anterior, el producte es calcula component a component,
\[
V_{G\times H}=V_G\times V_H,
\qquad
E_{G\times H}=E_G\times E_H,
\]
amb origen i destí també component a component: l'aresta $(e,e')$ va del vèrtex $(s_G(e),s_H(e'))$ al vèrtex $(t_G(e),t_H(e'))$. Les projeccions són els morfismes de grafs que retenen una de les dues components. Prenem ara com a $T$ el candidat ingenu a \enquote{punt}, el graf $\mathbf V$ d'un sol vèrtex i cap aresta. La fórmula dona
\[
V_{G\times\mathbf V}=V_G\times\{\ast\}\cong V_G,
\qquad
E_{G\times\mathbf V}=E_G\times\varnothing=\varnothing:
\]
el producte amb $\mathbf V$ \emph{esborra totes les arestes}. Per tant, $\mathbf V$ no és l'objecte terminal de $\cat{Graph}$; de fet, d'un graf amb alguna aresta no hi ha cap morfisme cap a $\mathbf V$, perquè l'aresta no té on anar. L'objecte terminal es calcula també punt a punt, com el límit del diagrama buit a cada objecte de $\cat{P}$: un vèrtex \emph{i} una aresta, és a dir, el graf $\mathbf L$ d'un sol \emph{bucle}. Amb $T=\mathbf L$, sí que $G\times\mathbf L\cong G$. I els morfismes $\mathbf L\to G$ ja no són els vèrtexs de $G$, sinó els seus bucles: un graf sense bucles no té cap \enquote{punt} en sentit categòric, encara que tingui molts vèrtexs (compareu-ho amb l'observació~\ref{obs:raonar-elements}).

\emph{A una categoria de functors $\cat{Set}^{\cat{A}}$.} La recepta és la mateixa per a qualsevol categoria petita $\cat{A}$: $(F\times G)(a)=F(a)\times G(a)$ i $(F\times G)(v)=F(v)\times G(v)$, i l'objecte terminal és el functor constant de valor $\{\ast\}$. Els dos casos anteriors en són casos particulars: $\cat{Set}\cong\cat{Set}^{\cat{1}}$, i $\cat{Graph}\cong\cat{Set}^{\cat{P}}$. Ara s'entén la sorpresa del bucle: el terminal de $\cat{Set}^{\cat{P}}$ ha de ser un punt \emph{a cada objecte} de $\cat{P}$, també a $E$, i això obliga a tenir una aresta.

La forma del diagrama i la recepta del càlcul no canvien d'un context a l'altre. El que canvia és quins objectes fan cada paper, i la intuició de $\cat{Set}$ només es pot traslladar si es fa objecte per objecte de la categoria índex.
\end{exemple}

Com que $\cat{Set}$ té tots els límits i colímits petits (exemple~\ref{ex:limits-set}), la proposició mostra que també els té la categoria de prefeixos $\cat{Set}^{\cat{C}\op}$ de qualsevol categoria petita $\cat{C}$, tal com s'anunciava al capítol~\ref{cap:yoneda}. Aquest fet serà essencial en els capítols de subobjectes i de topos.

\section{Preservació i reflexió de límits}

Quan apliquem un functor a un diagrama i al seu límit, la imatge és un con sobre el diagrama imatge, però no té per què ser-ne el límit. Els functors que respecten límits reben un nom.

\begin{definicio}\index{functor!preserva límits}\index{functor!reflecteix límits}
Sigui $F\colon\cat{C}\to\cat{D}$ un functor.
\begin{itemize}
\item $F$ \textbf{preserva} el límit $(L,\lambda)$ d'un diagrama $D\colon\cat{I}\to\cat{C}$ si $(F(L),F\lambda)$ és un límit de $F\comp D$ a $\cat{D}$.
\item $F$ \textbf{reflecteix} límits si, per a tot diagrama $D\colon\cat{I}\to\cat{C}$ i tot con $(L,\lambda)$ sobre $D$, quan $(F(L),F\lambda)$ és un límit de $F\comp D$, aleshores $(L,\lambda)$ ja era un límit de $D$.
\end{itemize}
Es diu que $F$ \textbf{preserva els límits finits} si preserva el límit de tot diagrama finit que en tingui, i anàlogament per a altres classes.
\end{definicio}

\begin{observacio}
La reflexió es formula sempre \emph{respecte d'un con i un diagrama concrets}: cal partir d'un con $(L,\lambda)$ sobre $D$ la imatge del qual és limitant. Un functor ple i fidel permet aleshores aixecar i detectar les factoritzacions úniques ---perquè estableix una bijecció sobre els homconjunts---, però això no vol dir que \enquote{creï} límits del no-res: si a $\cat{C}$ no hi ha cap con sobre $D$ que projecti sobre el límit imatge, no hi ha res a reflectir. La distinció entre \emph{reflectir} (comparar un con donat) i \emph{crear} (produir-ne un de nou) és important i convé tenir-la present sempre que un functor interactuï amb els límits. La noció formal de crear límits, útil sobretot per als functors oblit de categories algebraiques, no serà necessària aquí; n'hi ha prou de retenir que és més forta que la mera reflexió d'un con ja donat.
\end{observacio}

\begin{exemple}
El functor oblit $U\colon\cat{Grp}\to\cat{Set}$ preserva tots els límits petits. En efecte, si $D\colon\cat{I}\to\cat{Grp}$ és un diagrama petit, el límit dels conjunts subjacents és el conjunt de famílies compatibles $(x_i)_i$ (exemple~\ref{ex:limits-set}). Les operacions de grup, definides component a component, conserven aquestes compatibilitats, perquè cada $D_u$ és un homomorfisme; aquest conjunt adquireix així una estructura de grup per a la qual les projeccions són homomorfismes i que satisfà la propietat universal del límit a $\cat{Grp}$. En particular, això recupera els productes i els equalitzadors habituals. En canvi, $U$ \emph{no} preserva tots els colímits: ja falla per als coproductes, perquè el conjunt subjacent d'un coproducte de grups (el producte lliure) no és, en general, la unió disjunta dels conjunts subjacents. Aquesta asimetria és típica dels functors oblit i anticipa el teorema del capítol següent: els functors oblit solen ser adjunts per la dreta, i els adjunts per la dreta preserven tots els límits; no hi ha cap afirmació anàloga que els obligui a preservar tots els colímits.
\end{exemple}

\begin{proposicio}\label{prop:representables-limits}\index{functor!representable!preserva límits}
Per a tot objecte $A$ d'una categoria localment petita, el functor representable $\Hom(A,-)\colon\cat{C}\to\cat{Set}$ preserva tots els límits que existeixen a $\cat{C}$.
\end{proposicio}

\begin{prova}
Sigui $(L,\lambda)$ un límit d'un diagrama $D\colon\cat{I}\to\cat{C}$. Escrivim $H=\Hom(A,-)$, de manera que $H(X)=\Hom(A,X)$ i $H(v)=v_*\colon\varphi\mapsto v\comp\varphi$.

\emph{El con imatge.} La família $H\lambda=\bigl((\lambda_i)_*\bigr)_i$, amb $(\lambda_i)_*\colon\Hom(A,L)\to\Hom(A,D_i)$, és un con sobre $H\comp D$: per a $u\colon i\to j$, la functorialitat de $H$ dona $(D_u)_*\comp(\lambda_i)_*=(D_u\comp\lambda_i)_*=(\lambda_j)_*$.

\emph{Existència de la factorització.} Sigui $(S,\sigma)$ un con sobre $H\comp D$ a $\cat{Set}$. És a dir, $S$ és un conjunt i $\sigma_i\colon S\to\Hom(A,D_i)$ són aplicacions amb $(D_u)_*\comp\sigma_i=\sigma_j$. Avaluada en un element $s\in S$, aquesta condició diu $D_u\comp\sigma_i(s)=\sigma_j(s)$: la família $\bigl(\sigma_i(s)\bigr)_i$ és un con sobre $D$ amb vèrtex $A$. Per la propietat universal de $L$, hi ha una única fletxa $h(s)\colon A\to L$ amb $\lambda_i\comp h(s)=\sigma_i(s)$ per a tot $i$. Això defineix una aplicació $h\colon S\to\Hom(A,L)$ amb $(\lambda_i)_*\comp h=\sigma_i$ per a tot $i$.

\emph{Unicitat.} Si $h'\colon S\to\Hom(A,L)$ també compleix $(\lambda_i)_*\comp h'=\sigma_i$, aleshores, per a cada $s$, $\lambda_i\comp h'(s)=\sigma_i(s)$ per a tot $i$, i la unicitat de la factorització a través de $L$ dona $h'(s)=h(s)$. Per tant $h'=h$.

Així $\bigl(\Hom(A,L),H\lambda\bigr)$ és un límit de $H\comp D$, i $\Hom(A,-)$ preserva el límit $(L,\lambda)$.
\end{prova}

\begin{observacio}[La prova en una línia]
La demostració anterior és la lectura element a element de la cadena de bijeccions
\[
\Hom\bigl(A,\varprojlim D\bigr)\;\cong\;\mathrm{Con}(A,D)\;\cong\;\varprojlim\nolimits_i\Hom(A,D_i).
\]
La primera bijecció és la proposició~\ref{prop:limit-representable}. La segona és la descripció dels límits a $\cat{Set}$ com a famílies compatibles (exemple~\ref{ex:limits-set}): un element de $\varprojlim_i\Hom(A,D_i)$ és una família $(\alpha_i)_i$ amb $D_u\comp\alpha_i=\alpha_j$, és a dir, un con. La composta envia $f$ a $(\lambda_i\comp f)_i$, que és l'acció del con imatge.
\end{observacio}

\begin{observacio}
Aquest fet ---els representables covariants preserven límits--- és una peça clau que retrobarem. Dualment, els representables contravariants $\Hom(-,B)$ envien colímits a límits: com que $\Hom_{\cat{C}}(-,B)=\Hom_{\cat{C}\op}(B,-)$ i un colímit de $\cat{C}$ és un límit de $\cat{C}\op$, n'hi ha prou d'aplicar la proposició a $\cat{C}\op$. Per exemple, $\Hom(A+A',B)\cong\Hom(A,B)\times\Hom(A',B)$. Combinat amb el lema de Yoneda, permet raonar sobre límits component a component, i és la base tècnica de molts arguments de la teoria.
\end{observacio}

\begin{resumcap}
\apartat{Idees centrals}
\begin{enumerate}
  \item Un límit és el con universal sobre un diagrama; un colímit, el cocon universal, és a dir, el límit del diagrama corresponent a la categoria oposada.
  \item El límit d'un diagrama $D$ representa el prefeix de cons: $\Hom(-,\varprojlim D)\cong\mathrm{Con}(-,D)$. Per tant, el con límit és únic llevat d'un únic isomorfisme compatible amb les projeccions, i el vèrtex sol, llevat d'isomorfisme.
  \item Productes, equalitzadors, pullbacks i l'objecte terminal són límits; els seus duals són colímits. Tenir tots els límits finits equival a tenir objecte terminal, productes binaris i equalitzadors; equivalentment, objecte terminal i pullbacks.
  \item $\cat{Set}$ té tots els límits i colímits petits (famílies compatibles i quocients de sumes), i les categories de functors els calculen punt a punt quan existeixen a la categoria d'arribada.
  \item Els pullbacks indueixen la reindexació $f^{*}\colon\cat{C}/B\to\cat{C}/A$, que a $\cat{Set}$ reindexa fibres i anticipa la substitució de predicats.
  \item Els representables $\Hom(A,-)$ preserven límits.
\end{enumerate}
\apartat{Errors típics} Confondre el con (la família de projeccions) amb el seu vèrtex (l'objecte); creure que \enquote{tenir límits} depèn de la tria d'un límit concret, quan la propietat universal el fixa llevat d'isomorfisme; oblidar que un colímit no és un límit \emph{de la mateixa} categoria, sinó de l'oposada; suposar que tot functor preserva límits (només ho fan certs functors, com els representables covariants); i confondre \enquote{$F$ preserva un límit}, que parteix d'un límit de $\cat{C}$ i en demana la imatge, amb \enquote{$F$ reflecteix límits}, que parteix d'un con de $\cat{C}$ amb imatge limitant i en conclou que ja era un límit.
\apartat{Lectura recomanada} \cite[cap.~5]{awodey} per a l'exposició paral·lela a aquesta; \cite[III.3--III.4 i cap.~V]{maclane} per al tractament clàssic de límits com a cons universals; \cite[cap.~3]{riehl} per a exemples addicionals, en particular §3.3, sobre preservació, reflexió i creació de límits, i §3.4, sobre la naturalesa representable dels límits; \cite[cap.~9]{barrwells} per a la intuïció del límit com a subconjunt d'un producte tallat per equacions, i del colímit com a quocient d'una suma; \cite[cap.~3]{perrone-notes} per al prefeix de cons i el cas dels preordres; \cite[cap.~5]{leinster} per a una presentació breu, amb la interacció entre functors i límits; i \cite[part~IV, sessions~19--25]{lawvereschanuel}, per a objectes terminals, productes i sumes, amb els productes de grafs, en una presentació molt elemental.
\end{resumcap}

\chapter[Adjuncions]{Adjuncions}
\label{cap:adjuncions}

\begin{objectius}
\apartat{Prerequisits} Functors, preordres vistos com a categories i la categoria lliure d'un graf (cap.~\ref{cap:functors}); transformacions i isomorfismes naturals, i equivalències de categories (cap.~\ref{cap:transformacions}); representabilitat i propietats universals (cap.~\ref{cap:yoneda}); límits i colímits, i la reindexació per pullback (cap.~\ref{cap:limits}).
\apartat{Objectius} En acabar el capítol, el lector hauria de poder:
\begin{itemize}
  \item definir una adjunció mitjançant una bijecció natural d'homconjunts i mitjançant unitat i counitat;
  \item calcular transposats i reconèixer la propietat universal de la unitat;
  \item verificar les identitats triangulars;
  \item reconèixer adjuncions en preordres i en construccions lliure--oblit;
  \item interpretar $\colimop_{\cat{I}}\dashv\Delta\dashv\varprojlim_{\cat{I}}$ i la currificació com a adjuncions;
  \item interpretar a $\cat{Set}$ la cadena $\exists_f\dashv f^{*}\dashv\forall_f$ com a prototip de la quantificació;
  \item demostrar la unicitat dels adjunts llevat d'isomorfisme natural;
  \item aplicar que els adjunts per la dreta preserven límits i els adjunts per l'esquerra, colímits.
\end{itemize}
\end{objectius}

\section{Cap a la noció d'adjunció}

Al llarg del llibre hem trobat ja la parella formada pel functor categoria lliure $\Free\colon\cat{Graph}\to\cat{Cat}$ i el functor graf subjacent $U\colon\cat{Cat}\to\cat{Graph}$. Un altre patró clàssic, que veurem ara en una versió elemental, és el del monoide lliure i el functor oblit $U\colon\cat{Mon}\to\cat{Set}$. En tots dos casos, un functor afegeix estructura de la manera més lliure possible i l'altre l'oblida, i tots dos estan lligats per una correspondència precisa. Aquesta relació s'anomena \textbf{adjunció}, i és un dels conceptes centrals de la teoria de categories.

Les adjuncions explicaran també un patró que vam observar en estudiar els límits: els representables i molts functors oblit preserven límits de manera sistemàtica, mentre que moltes construccions lliures preserven colímits. El teorema d'aquest capítol explicarà aquests fenòmens quan els functors en qüestió siguin adjunts. I seran un pont directe cap a la lògica: els quantificadors existencial i universal es revelaran com a adjunts.

La idea es pot copsar primer en el cas més simple, el dels preordres, on una adjunció és una parella d'aplicacions monòtones lligades per una equivalència de desigualtats. Després la generalitzarem a categories qualssevol, substituint les desigualtats per bijeccions entre conjunts de morfismes.

\section{Adjuncions entre preordres}

Siguin $(P,\leq)$ i $(Q,\leq)$ dos preordres, vistos com a categories, i $f\colon P\to Q$, $g\colon Q\to P$ dues aplicacions monòtones.

\begin{definicio}\index{adjunció!entre preordres}
Diem que $f$ és \textbf{adjunta per l'esquerra} de $g$ (i $g$ \textbf{adjunta per la dreta} de $f$), i escrivim $f\dashv g$, quan per a tot $x\in P$ i tot $y\in Q$,
\[
f(x)\leq y
\quad\Longleftrightarrow\quad
x\leq g(y).
\]
\end{definicio}

\begin{observacio}[Unitat i counitat en un preordre]\index{connexió de Galois}
Si $f\dashv g$, prenent $y=f(x)$ i $x=g(y)$ a la definició s'obté
\[
x\leq g(f(x)),
\qquad
f(g(y))\leq y
\]
per a tot $x\in P$ i tot $y\in Q$. Recíprocament, si $f$ i $g$ són monòtones i compleixen aquestes dues desigualtats, aleshores $f\dashv g$: si $f(x)\leq y$, aleshores $x\leq g(f(x))\leq g(y)$, i si $x\leq g(y)$, aleshores $f(x)\leq f(g(y))\leq y$. Aquestes dues desigualtats són exactament la \emph{unitat} i la \emph{counitat} de l'adjunció, que veurem en general més endavant; en una categoria prima, les identitats triangulars no aporten cap igualtat nova, perquè entre dos objectes hi ha com a màxim una fletxa. Una adjunció entre preordres s'anomena sovint \textbf{connexió de Galois}.
\end{observacio}

\begin{exemple}[Conjunció i implicació]\index{adjunció!conjunció i implicació}
Sigui $\cat{Prov}_T$ el preordre de fórmules d'un sistema proposicional $T$ amb les regles usuals de $\wedge$ i $\to$, i fixem una fórmula $A$. Les operacions $A\wedge-$ i $A\to-$ són monòtones, i per a totes les fórmules $B$ i $C$
\[
A\wedge B\vdash_T C
\quad\Longleftrightarrow\quad
B\vdash_T A\to C,
\]
de manera que
\[
A\wedge-\;\dashv\;A\to-.
\]
Aquesta adjunció codifica les regles de la implicació; és la forma categòrica del teorema de deducció. És un primer indici que alguns connectius lògics es caracteritzen mitjançant adjuncions, idea que reapareixerà de manera sistemàtica en la lògica categòrica.
\end{exemple}

\begin{exemple}[Imatge, antiimatge i quantificadors]\label{ex:quantificadors-set}\index{adjunció!quantificadors}\index{quantificador!com a adjunt}
Sigui $f\colon X\to Y$ una aplicació, i considerem els preordres $\mathcal P(X)$ i $\mathcal P(Y)$ de les parts, ordenats per inclusió. L'antiimatge
\[
f^{*}(T)=f^{-1}(T),\qquad T\subseteq Y,
\]
és monòtona, i té adjunts a tots dos costats: per a $S\subseteq X$, posem
\[
\exists_f(S)=f[S]=\{f(x)\mid x\in S\},
\qquad
\forall_f(S)=\{y\in Y\mid f^{-1}(\{y\})\subseteq S\}.
\]
Aleshores, per a tot $S\subseteq X$ i tot $T\subseteq Y$,
\[
\exists_f(S)\subseteq T
\quad\Longleftrightarrow\quad
S\subseteq f^{*}(T),
\qquad\qquad
f^{*}(T)\subseteq S
\quad\Longleftrightarrow\quad
T\subseteq\forall_f(S).
\]
La primera equivalència diu que les imatges dels elements de $S$ són a $T$ exactament quan tots els elements de $S$ van a parar a $T$. Per a la segona: si $f^{-1}(T)\subseteq S$ i $y\in T$, tot $x$ amb $f(x)=y$ és a $f^{-1}(T)$ i, per tant, a $S$, de manera que $y\in\forall_f(S)$; i recíprocament, si $T\subseteq\forall_f(S)$ i $f(x)\in T$, aleshores $x\in f^{-1}(\{f(x)\})\subseteq S$. Per tant,
\[
\exists_f\;\dashv\;f^{*}\;\dashv\;\forall_f.
\]

El nom dels adjunts s'explica quan $f$ és una projecció $\pi\colon X\times A\to X$. Un subconjunt $S\subseteq X\times A$ és un predicat en dues variables, i
\[
\begin{aligned}
\exists_\pi(S)&=\{x\mid \text{existeix }a\in A\text{ amb }(x,a)\in S\},\\
\forall_\pi(S)&=\{x\mid (x,a)\in S\text{ per a tot }a\in A\}:
\end{aligned}
\]
són exactament les lectures conjuntistes de la quantificació existencial i universal sobre la variable d'$A$, mentre que $\pi^{*}$ afegeix una variable muda. Al capítol de subobjectes veurem en quines categories existeixen aquests dos adjunts; aquí només afirmem el cas de $\cat{Set}$.
\end{exemple}

\begin{exemple}[Interior topològic]\index{adjunció!interior topològic}
Sigui $X$ un espai topològic, $\mathcal P(X)$ el preordre de les parts per inclusió i $\mathcal O(X)$ el dels oberts. La injecció $i\colon\mathcal O(X)\to\mathcal P(X)$ té com a adjunt per la dreta l'\textbf{interior} $\mathrm{int}\colon\mathcal P(X)\to\mathcal O(X)$:
\[
i(U)\subseteq S
\quad\Longleftrightarrow\quad
U\subseteq\mathrm{int}(S),
\]
per a tot obert $U$ i tota part $S$. És a dir, l'interior és el més gran obert contingut en $S$.
\end{exemple}

\section{Adjuncions entre categories}

La definició general substitueix l'equivalència de desigualtats per una \textbf{bijecció entre conjunts de morfismes}, natural en totes dues variables.

\begin{definicio}[Adjunció]\index{adjunció}
\label{def:adjuncio}
Siguin $F\colon\cat{C}\to\cat{D}$ i $G\colon\cat{D}\to\cat{C}$ dos functors entre categories localment petites. Una \textbf{adjunció} $F\dashv G$ és una bijecció
\[
\theta_{A,B}\colon\Hom_{\cat{D}}(F(A),B)\;\xrightarrow{\ \cong\ }\;\Hom_{\cat{C}}(A,G(B)),
\]
per a cada $A$ de $\cat{C}$ i $B$ de $\cat{D}$, \textbf{natural} en $A$ i en $B$. Es diu que $F$ és l'\textbf{adjunt per l'esquerra} i $G$ l'\textbf{adjunt per la dreta}.
\end{definicio}

\begin{observacio}
La naturalitat vol dir que $\theta$ és un isomorfisme natural entre els dos functors
\[
\Hom_{\cat{D}}(F(-),-)\quad\text{i}\quad\Hom_{\cat{C}}(-,G(-)),
\]
tots dos de $\cat{C}\op\times\cat{D}$ cap a $\cat{Set}$. Explícitament, per a $a\colon A'\to A$ a $\cat{C}$, $b\colon B\to B'$ a $\cat{D}$ i $\varphi\colon F(A)\to B$, el pas per $\theta$ commuta amb la postcomposició i amb la precomposició:
\[
\theta_{A,B'}(b\comp\varphi)=G(b)\comp\theta_{A,B}(\varphi),
\qquad\qquad
\theta_{A',B}\bigl(\varphi\comp F(a)\bigr)=\theta_{A,B}(\varphi)\comp a.
\]
En paraules: transposar i després compondre amb $G(b)$ o amb $a$ és el mateix que compondre primer amb $b$ o amb $F(a)$ i després transposar. Són els dos quadrats de naturalitat
\[
\begin{tikzpicture}[baseline=(m.center)]
  \matrix (m) [matrix of math nodes, row sep=2.8em, column sep=3.6em] {
    \Hom(F(A),B) & \Hom(A,G(B)) \\
    \Hom(F(A),B') & \Hom(A,G(B')) \\
  };
  \draw[-{Stealth[scale=1.1]}] (m-1-1) -- node[above] {$\theta_{A,B}$} (m-1-2);
  \draw[-{Stealth[scale=1.1]}] (m-1-1) -- node[left] {$b\comp-$} (m-2-1);
  \draw[-{Stealth[scale=1.1]}] (m-1-2) -- node[right] {$G(b)\comp-$} (m-2-2);
  \draw[-{Stealth[scale=1.1]}] (m-2-1) -- node[below] {$\theta_{A,B'}$} (m-2-2);
\end{tikzpicture}
\]
i
\[
\begin{tikzpicture}[baseline=(m.center)]
  \matrix (m) [matrix of math nodes, row sep=2.8em, column sep=3.6em] {
    \Hom(F(A),B) & \Hom(A,G(B)) \\
    \Hom(F(A'),B) & \Hom(A',G(B)) \\
  };
  \draw[-{Stealth[scale=1.1]}] (m-1-1) -- node[above] {$\theta_{A,B}$} (m-1-2);
  \draw[-{Stealth[scale=1.1]}] (m-1-1) -- node[left] {$-\comp F(a)$} (m-2-1);
  \draw[-{Stealth[scale=1.1]}] (m-1-2) -- node[right] {$-\comp a$} (m-2-2);
  \draw[-{Stealth[scale=1.1]}] (m-2-1) -- node[below] {$\theta_{A',B}$} (m-2-2);
\end{tikzpicture}
\]
on els vèrtexs són conjunts de morfismes i les fletxes verticals actuen per composició. Els morfismes que es corresponen sota $\theta$ s'anomenen \textbf{transposats} l'un de l'altre: si $\varphi\colon F(A)\to B$, el seu transposat és $\theta(\varphi)\colon A\to G(B)$.
\end{observacio}

\begin{observacio}
El cas dels preordres n'és un cas particular: si $P,Q$ són preordres, un conjunt de morfismes $\Hom(x,y)$ té com a molt un element, de manera que la bijecció $\Hom(f(x),y)\cong\Hom(x,g(y))$ es redueix a l'equivalència $f(x)\leq y\Leftrightarrow x\leq g(y)$. Així, la definició general estén literalment la dels preordres.
\end{observacio}

\section{Unitat i counitat}

Una adjunció es pot descriure, de manera equivalent, sense esmentar la bijecció, mitjançant dues transformacions naturals distingides.

\begin{definicio}[Unitat i counitat]\index{adjunció!unitat}\index{adjunció!counitat}
\label{def:unitat-counitat}
Donada una adjunció $F\dashv G$ amb bijecció $\theta$:
\begin{itemize}
\item la \textbf{unitat} és la transformació natural $\eta\colon\id_{\cat{C}}\Rightarrow G\comp F$ de components
\[
\eta_A=\theta_{A,F(A)}(\id_{F(A)})\colon A\to G(F(A));
\]
\item la \textbf{counitat} és la transformació natural $\varepsilon\colon F\comp G\Rightarrow\id_{\cat{D}}$ de components
\[
\varepsilon_B=\theta_{G(B),B}^{-1}(\id_{G(B)})\colon F(G(B))\to B.
\]
\end{itemize}
\end{definicio}

\begin{observacio}[De la unitat i la counitat a la bijecció]\label{obs:adj-transposats}
La unitat i la counitat permeten recuperar la bijecció. Donat $\varphi\colon F(A)\to B$, el seu transposat és
\[
\theta(\varphi)=G(\varphi)\comp\eta_A\colon A\to G(B),
\]
i dualment, donat $\psi\colon A\to G(B)$, el seu transposat invers és
\[
\theta^{-1}(\psi)=\varepsilon_B\comp F(\psi)\colon F(A)\to B.
\]
\end{observacio}

\begin{definicio}[Fletxa universal]\label{def:fletxa-universal}\index{fletxa universal}\index{adjunció!fletxa universal}
Sigui $G\colon\cat{D}\to\cat{C}$ un functor i $A$ un objecte de $\cat{C}$. Una \textbf{fletxa universal} de $A$ cap a $G$ és una parella $(R,\eta_A)$ formada per un objecte $R$ de $\cat{D}$ i un morfisme $\eta_A\colon A\to G(R)$ tal que tot $\psi\colon A\to G(B)$ factoritza de manera única com $\psi=G(\varphi)\comp\eta_A$, amb $\varphi\colon R\to B$.
\end{definicio}

En una adjunció $F\dashv G$, cada component de la unitat és una fletxa universal, amb $R=F(A)$ (exercici~\ref{exr:adj-unitat-universal} dels Exercicis):
\[
\begin{tikzpicture}[baseline=(m.center)]
  \matrix (m) [matrix of math nodes, column sep=4em, row sep=2.7em] {
    A & G(F(A)) \\
      & G(B) \\
  };
  \draw[-{Stealth[scale=1.1]}] (m-1-1) -- node[above] {$\eta_A$} (m-1-2);
  \draw[-{Stealth[scale=1.1]}] (m-1-1) -- node[below left] {$\psi$} (m-2-2);
  \draw[-{Stealth[scale=1.1]},dashed] (m-1-2) -- node[right] {$G(\varphi)$} (m-2-2);
\end{tikzpicture}
\qquad
\psi=G(\varphi)\comp\eta_A,\quad\varphi\text{ únic}.
\]
Aquesta és la formulació de l'adjunció que connecta directament amb els objectes lliures.

\begin{exemple}[El monoide lliure]\index{adjunció!lliure-oblit}
Sigui $U\colon\cat{Mon}\to\cat{Set}$ el functor oblit dels monoides i $F\colon\cat{Set}\to\cat{Mon}$ el functor que envia un conjunt $X$ al \textbf{monoide lliure} $X^*$ de les paraules finites en $X$, amb la concatenació. Aleshores $F\dashv U$: un homomorfisme de monoides $X^*\to M$ està determinat de manera única pels valors sobre les lletres, és a dir, per una aplicació $X\to U(M)$. Aquesta és exactament la bijecció
\[
\Hom_{\cat{Mon}}(X^*,M)\cong\Hom_{\cat{Set}}(X,U(M)).
\]
La unitat $\eta_X\colon X\to U(X^*)$ envia cada element a la paraula d'una lletra; és la fletxa universal que exhibeix la propietat universal del monoide lliure. Aquest patró ---\emph{lliure $\dashv$ oblit}--- es repeteix per a grups, espais vectorials, categories i moltes altres estructures.
\end{exemple}

\begin{exemple}[La categoria lliure d'un graf: $\Free\dashv U$]\label{ex:free-u-adjuncio}\index{adjunció!categoria lliure}
Els capítols anteriors ja han preparat aquest exemple. El functor graf subjacent $U\colon\cat{Cat}\to\cat{Graph}$ i el functor categoria lliure $\Free\colon\cat{Graph}\to\cat{Cat}$ satisfan la bijecció
\[
\Hom_{\cat{Cat}}\bigl(\Free(G),\cat{D}\bigr)\;\cong\;\Hom_{\cat{Graph}}\bigl(G,U(\cat{D})\bigr),
\]
natural en $G$ i en $\cat{D}$, que el capítol~\ref{cap:yoneda} va obtenir com la representabilitat del functor $\cat{D}\mapsto\Hom_{\cat{Graph}}(G,U(\cat{D}))$. És a dir, $\Free\dashv U$. El que aporta ara el llenguatge de les adjuncions és identificar-ne les dades:
\begin{itemize}
  \item la \emph{unitat} $\eta_G\colon G\to U(\Free(G))$ és la inserció dels generadors: cada vèrtex va a ell mateix i cada aresta, al camí de longitud~$1$ que la recorre;
  \item la \emph{counitat} $\varepsilon_{\cat{C}}\colon\Free(U(\cat{C}))\to\cat{C}$ \emph{avalua} els camins formals: envia un camí $(f_1,\dots,f_n)$ de fletxes de $\cat{C}$, recorregut en aquest ordre, a la composta $f_n\comp\cdots\comp f_1$, i el camí buit en un objecte, a la seva identitat.
\end{itemize}
\end{exemple}

\section{Identitats triangulars}

La unitat i la counitat no són independents: satisfan dues equacions que, de fet, caracteritzen les adjuncions.

\begin{proposicio}[Identitats triangulars]\label{prop:triangulars}\index{adjunció!identitats triangulars}
Per a una adjunció $F\dashv G$ amb unitat $\eta$ i counitat $\varepsilon$, es compleixen les \textbf{identitats triangulars}:
\[
\varepsilon_{F(A)}\comp F(\eta_A)=\id_{F(A)},
\qquad
G(\varepsilon_B)\comp\eta_{G(B)}=\id_{G(B)},
\]
és a dir, commuten els dos triangles
\[
\begin{tikzpicture}[baseline=(m.center)]
  \matrix (m) [matrix of math nodes, column sep=3.6em, row sep=2.7em] {
    F(A) & F(G(F(A))) \\
         & F(A) \\
  };
  \draw[-{Stealth[scale=1.1]}] (m-1-1) -- node[above] {$F(\eta_A)$} (m-1-2);
  \draw[-{Stealth[scale=1.1]}] (m-1-1) -- node[below left] {$\id_{F(A)}$} (m-2-2);
  \draw[-{Stealth[scale=1.1]}] (m-1-2) -- node[right] {$\varepsilon_{F(A)}$} (m-2-2);
\end{tikzpicture}
\qquad\qquad
\begin{tikzpicture}[baseline=(m.center)]
  \matrix (m) [matrix of math nodes, column sep=3.6em, row sep=2.7em] {
    G(B) & G(F(G(B))) \\
         & G(B) \\
  };
  \draw[-{Stealth[scale=1.1]}] (m-1-1) -- node[above] {$\eta_{G(B)}$} (m-1-2);
  \draw[-{Stealth[scale=1.1]}] (m-1-1) -- node[below left] {$\id_{G(B)}$} (m-2-2);
  \draw[-{Stealth[scale=1.1]}] (m-1-2) -- node[right] {$G(\varepsilon_B)$} (m-2-2);
\end{tikzpicture}
\]
Recíprocament, dues transformacions naturals $\eta\colon\id\Rightarrow G F$ i $\varepsilon\colon F G\Rightarrow\id$ que les satisfacin defineixen una adjunció $F\dashv G$.
\end{proposicio}

\begin{prova}
Comprovem la primera identitat. Pel càlcul del transposat, $\theta^{-1}(\psi)=\varepsilon_B\comp F(\psi)$; prenent $B=F(A)$ i $\psi=\eta_A=\theta(\id_{F(A)})$,
\[
\varepsilon_{F(A)}\comp F(\eta_A)=\theta^{-1}(\eta_A)=\theta^{-1}(\theta(\id_{F(A)}))=\id_{F(A)}.
\]
La segona identitat és dual, usant $\theta(\varphi)=G(\varphi)\comp\eta_A$ amb $A=G(B)$ i $\varphi=\varepsilon_B$. Per al recíproc, es defineix la bijecció per $\theta(\varphi)=G(\varphi)\comp\eta_A$ i $\theta^{-1}(\psi)=\varepsilon_B\comp F(\psi)$; les identitats triangulars garanteixen que són inverses l'una de l'altra, i la naturalitat de $\eta,\varepsilon$ dona la de $\theta$.
\end{prova}

\begin{observacio}
Les identitats triangulars reben aquest nom perquè expressen la commutativitat de dos triangles que involucren $F$, $G$, $\eta$ i $\varepsilon$. Donen una definició d'adjunció \emph{purament diagramàtica}, sense referència a conjunts de morfismes, la qual cosa és sovint la més còmoda per treballar i l'única disponible en contextos on no hi ha homconjunts (per exemple, en 2-categories abstractes).
\end{observacio}

\begin{observacio}[Tres presentacions d'una adjunció]\label{obs:tres-presentacions}\index{adjunció!tres presentacions}
Hem trobat tres maneres de donar una adjunció entre $\cat{C}$ i $\cat{D}$. Convé no barrejar-les, perquè no parteixen de les mateixes dades.
\begin{center}
\footnotesize
\renewcommand{\arraystretch}{1.3}
\begin{tabularx}{\linewidth}{@{}>{\raggedright\arraybackslash}p{0.2\linewidth}>{\raggedright\arraybackslash}X>{\raggedright\arraybackslash}X@{}}
\toprule
\textbf{Presentació} & \textbf{Dades} & \textbf{Condicions}\\
\midrule
bijecció d'homconjunts (definició~\ref{def:adjuncio}) & functors $F$ i $G$; bijeccions $\theta_{A,B}\colon\Hom(F(A),B)\to\Hom(A,G(B))$ & naturalitat en $A$ i en $B$\\
unitat i counitat (proposició~\ref{prop:triangulars}) & functors $F$ i $G$; transformacions naturals $\eta\colon\id\Rightarrow GF$ i $\varepsilon\colon FG\Rightarrow\id$ & identitats triangulars\\
fletxes universals (definició~\ref{def:fletxa-universal}) & només el functor $G$; per a cada $A$, un objecte $F(A)$ i una fletxa $\eta_A\colon A\to G(F(A))$ & cada $\eta_A$ és universal; $F$ no es dona sobre morfismes, sinó que la propietat universal el determina\\
\bottomrule
\end{tabularx}
\end{center}
Els passos d'una presentació a una altra són aquests:
\begin{itemize}
\item De la primera a la segona: $\eta$ i $\varepsilon$ són els transposats de les identitats (definició~\ref{def:unitat-counitat}).
\item De la segona a la primera: les fórmules de l'observació~\ref{obs:adj-transposats}, que són inverses gràcies a les identitats triangulars (exercici~\ref{exr:adj-triangulars-reciproc} dels Exercicis).
\item De la primera a la tercera: la unitat és universal (exercici~\ref{exr:adj-unitat-universal}).
\item De la tercera a la primera: la propietat universal defineix $F$ sobre morfismes i la bijecció $\varphi\mapsto G(\varphi)\comp\eta_A$ (exercici~\ref{exr:adj-fletxes-universals}).
\end{itemize}
La tercera presentació és la més econòmica, perquè no cal conèixer $F$ per endavant; per això és la que s'usa per \emph{construir} adjunts, com els objectes lliures. La segona és la més còmoda per calcular, i la primera és la que dona directament els teoremes de preservació de límits de la secció següent.
\end{observacio}

\begin{observacio}[Adjuncions i equivalències]\index{equivalència de categories!adjunta}
Al capítol~\ref{cap:transformacions} vam definir una equivalència mitjançant un quasiinvers $G$ i dos isomorfismes naturals orientats com $\eta\colon\id_{\cat{C}}\Rightarrow G\comp F$ i $\varepsilon\colon F\comp G\Rightarrow\id_{\cat{D}}$ (definició~\ref{def:equivalencia}), la mateixa orientació que la unitat i la counitat d'una adjunció $F\dashv G$. La coincidència no és casual. Tota equivalència es pot equipar amb una estructura d'\textbf{equivalència adjunta}: modificant, si cal, un dels dos isomorfismes naturals, es pot aconseguir que $\eta$ i $\varepsilon$ satisfacin les identitats triangulars, i aleshores $F\dashv G$ amb unitat i counitat isomorfismes naturals. Recíprocament, una adjunció en què la unitat i la counitat són isomorfismes naturals és una equivalència. Les equivalències són, doncs, les adjuncions més simètriques possibles.
\end{observacio}

\section{Unicitat dels adjunts}

\begin{proposicio}\index{adjunció!unicitat}
Els adjunts són únics llevat d'isomorfisme natural. Si $F\dashv G$ i $F\dashv G'$, aleshores $G\cong G'$; i si $F\dashv G$ i $F'\dashv G$, aleshores $F\cong F'$.
\end{proposicio}

\begin{prova}
Si $F\dashv G$ i $F\dashv G'$, aleshores per a tot $A,B$,
\[
\Hom(A,G(B))\cong\Hom(F(A),B)\cong\Hom(A,G'(B)),
\]
naturalment en $A$. Pel corol·lari del lema de Yoneda ---dos objectes amb functors representables naturalment isomorfs són isomorfs---, $G(B)\cong G'(B)$ naturalment en $B$, és a dir, $G\cong G'$. El cas de $F$ és dual.
\end{prova}

\section{Els adjunts i els límits}

El resultat més útil sobre adjuncions diu com interactuen amb límits i colímits. És la raó de fons per la qual molts functors oblit preserven els productes i, en general, els límits.

\begin{teorema}[Els adjunts per la dreta preserven límits]\index{adjunció!preservació de límits}
\label{teo:adjunts-preserven-limits}
Si $G\colon\cat{D}\to\cat{C}$ és un adjunt per la dreta, aleshores $G$ preserva tots els límits petits que existeixen a $\cat{D}$. Dualment, un adjunt per l'esquerra preserva tots els colímits petits que existeixen.
\end{teorema}

\begin{prova}
Sigui $F\colon\cat{C}\to\cat{D}$ l'adjunt per l'esquerra de $G$, amb bijecció $\theta$ (definició~\ref{def:adjuncio}). Farem servir les dues equacions de naturalitat de $\theta$: per a $b\colon B\to B'$ a $\cat{D}$ i $a\colon A'\to A$ a $\cat{C}$,
\[
\theta(b\comp\varphi)=G(b)\comp\theta(\varphi),
\qquad\qquad
\theta\bigl(\varphi\comp F(a)\bigr)=\theta(\varphi)\comp a.
\]

\emph{Límits.} Sigui $(L,\lambda)$ un límit d'un diagrama $D\colon\cat{I}\to\cat{D}$. La família $G\lambda=(G(\lambda_i))_i$ és un con sobre $G\comp D$, perquè $G(D_u)\comp G(\lambda_i)=G(D_u\comp\lambda_i)=G(\lambda_j)$ per a cada $u\colon i\to j$. Vegem que és universal. Sigui $(A,\alpha)$ un con sobre $G\comp D$, amb $\alpha_i\colon A\to G(D_i)$ i $G(D_u)\comp\alpha_i=\alpha_j$, i siguin $\varphi_i=\theta^{-1}(\alpha_i)\colon F(A)\to D_i$ els transposats. Per la primera equació de naturalitat,
\[
\theta(D_u\comp\varphi_i)=G(D_u)\comp\alpha_i=\alpha_j=\theta(\varphi_j),
\]
i com que $\theta$ és bijectiva, $D_u\comp\varphi_i=\varphi_j$: les $\varphi_i$ formen un con sobre $D$ amb vèrtex $F(A)$. Per tant hi ha una única $h\colon F(A)\to L$ amb $\lambda_i\comp h=\varphi_i$ per a tot $i$. Posem $k=\theta(h)\colon A\to G(L)$. De nou per la primera equació,
\[
G(\lambda_i)\comp k=\theta(\lambda_i\comp h)=\theta(\varphi_i)=\alpha_i,
\]
de manera que $k$ factoritza el con $\alpha$. És l'única factorització: si $k'\colon A\to G(L)$ compleix $G(\lambda_i)\comp k'=\alpha_i$ per a tot $i$, aleshores $h'=\theta^{-1}(k')$ compleix $\theta(\lambda_i\comp h')=G(\lambda_i)\comp k'=\alpha_i=\theta(\varphi_i)$, d'on $\lambda_i\comp h'=\varphi_i$; per la unicitat de $h$, $h'=h$, i $k'=\theta(h')=k$. Així $(G(L),G\lambda)$ és un límit de $G\comp D$.

\emph{Colímits.} Sigui ara $(K,\kappa)$ un colímit d'un diagrama $E\colon\cat{I}\to\cat{C}$. La família $F\kappa=(F(\kappa_i))_i$ és un cocon sobre $F\comp E$, perquè $F(\kappa_j)\comp F(E_u)=F(\kappa_j\comp E_u)=F(\kappa_i)$. Sigui $(B,\beta)$ un cocon sobre $F\comp E$, amb $\beta_i\colon F(E_i)\to B$ i $\beta_j\comp F(E_u)=\beta_i$, i siguin $\psi_i=\theta(\beta_i)\colon E_i\to G(B)$. Per la segona equació de naturalitat,
\[
\psi_j\comp E_u=\theta\bigl(\beta_j\comp F(E_u)\bigr)=\theta(\beta_i)=\psi_i,
\]
de manera que les $\psi_i$ formen un cocon sobre $E$, i hi ha una única $m\colon K\to G(B)$ amb $m\comp\kappa_i=\psi_i$. Posem $n=\theta^{-1}(m)\colon F(K)\to B$. Aleshores $\theta\bigl(n\comp F(\kappa_i)\bigr)=m\comp\kappa_i=\psi_i=\theta(\beta_i)$, és a dir, $n\comp F(\kappa_i)=\beta_i$. Si $n'$ també compleix $n'\comp F(\kappa_i)=\beta_i$, aleshores $\theta(n')\comp\kappa_i=\theta(\beta_i)=\psi_i$, d'on $\theta(n')=m$ i $n'=n$. Així $(F(K),F\kappa)$ és un colímit de $F\comp E$.
\end{prova}

\begin{observacio}[La demostració en termes de representables]
La demostració anterior és la versió \enquote{element a element} d'una cadena de bijeccions. Per a tot objecte $A$ de $\cat{C}$, usant l'adjunció i que $\Hom(F(A),-)$ preserva límits (proposició~\ref{prop:representables-limits}),
\[
\begin{aligned}
\Hom(A,G(L))&\cong\Hom(F(A),L) &&\text{(adjunció)}\\
&\cong\varprojlim\nolimits_i\Hom(F(A),D_i) &&\text{($\Hom(F(A),-)$ preserva límits)}\\
&\cong\varprojlim\nolimits_i\Hom(A,G(D_i)) &&\text{(adjunció, natural en $D_i$)}\\
&\cong\mathrm{Con}(A,G\comp D) &&\text{(límits a $\cat{Set}$)},
\end{aligned}
\]
naturalment en $A$; la bijecció composta envia $f\colon A\to G(L)$ a $(G(\lambda_i)\comp f)_i$. Llegida així, l'afirmació diu que $G(L)$ representa el prefeix dels cons sobre $G\comp D$, amb element universal $G\lambda$ (proposició~\ref{prop:limit-representable}).
\end{observacio}

\begin{observacio}[Què vol dir exactament \enquote{preservar}]\label{obs:preservar-canonic}\index{morfisme de comparació}
El teorema no afirma cap igualtat $G(\varprojlim D)=\varprojlim(G\comp D)$: els límits només estan determinats llevat d'isomorfisme, i una igualtat no tindria sentit. Tampoc no afirma només que \emph{existeixi} algun isomorfisme entre aquests dos objectes. Afirma el que diu la definició del capítol~\ref{cap:limits}: el con imatge $(G(L),G\lambda)$ \emph{és un con límit} de $G\comp D$.

Si hem triat un límit $\bigl(\varprojlim(G\comp D),\pi\bigr)$ de $G\comp D$, la condició es pot expressar amb un sol morfisme. El con $G\lambda$ factoritza de manera única a través del límit triat, i dona l'únic \textbf{morfisme de comparació}
\[
c\colon G(L)\to\varprojlim(G\comp D),\qquad \pi_i\comp c=G(\lambda_i)\ \text{per a tot } i.
\]
Aleshores $G$ preserva el límit si i només si $c$ és un isomorfisme. En efecte, si $(G(L),G\lambda)$ és un límit, l'únic isomorfisme compatible amb les projeccions entre els dos límits de $G\comp D$ compleix $\pi_i\comp c=G(\lambda_i)$, i per tant és $c$. Recíprocament, si $c$ és invertible, el con $G\lambda$ és isomorf al con límit $\pi$ per un isomorfisme compatible amb les projeccions, i per tant és un límit. En aquest sentit precís escrivim $G(\varprojlim D)\cong\varprojlim(G\comp D)$: l'isomorfisme és el \emph{canònic} $c$, no un isomorfisme qualsevol. A més, la noció no depèn del límit triat de $D$, perquè dos límits de $D$ són isomorfs per un únic isomorfisme compatible amb les projeccions, i $G$ el transporta.

La diferència importa. El functor $X\mapsto X+X$ de $\cat{Set}$ en si mateix no preserva productes binaris. El morfisme de comparació $(X\times Y)+(X\times Y)\to(X+X)\times(Y+Y)$ envia cada parella de la còpia $k$ a la parella de còpies $(k,k)$, i no arriba als elements en què les dues còpies són diferents. Tanmateix, per a $X=Y=\mathbb N$ tots dos conjunts són numerables, i per tant \emph{hi ha} una bijecció entre ells. Dualment, $F$ preserva el colímit $(K,\kappa)$ de $E$ quan el morfisme de comparació $\colimop(F\comp E)\to F(K)$ és un isomorfisme.
\end{observacio}

\begin{observacio}
Aquest teorema es recorda amb les sigles angleses \enquote{RAPL} (\emph{right adjoints preserve limits}) i \enquote{LAPC} (\emph{left adjoints preserve colimits}). Explica de cop molts fets. Els functors oblit que són adjunts per la dreta preserven límits, però això no implica que preservin colímits; en molts exemples algebraics no els preserven. Dualment, els functors lliures, com a adjunts per l'esquerra, preserven colímits, però no hi ha cap teorema general que els obligui a preservar límits. El teorema dona una sola direcció i no prohibeix l'altra: un functor amb adjunts a tots dos costats, com l'oblit $\cat{Top}\to\cat{Set}$, preserva tant límits com colímits. És també un criteri negatiu útil: un functor que \emph{no} preserva límits no pot ser un adjunt per la dreta.
\end{observacio}

\subsection{El límit i el colímit com a adjunts del functor diagonal}\index{límit!com a adjunt}\index{colímit!com a adjunt}

Hi ha una segona manera, més sintètica, en què límits i colímits són adjuncions ---no ja \emph{preservats} per adjunts, sinó ells mateixos adjunts. Fixem una categoria d'índexs $\cat{I}$ i considerem el \textbf{functor diagonal}
\[
\Delta\colon\cat{C}\longrightarrow\cat{C}^{\cat{I}},
\]
que envia un objecte $A$ al diagrama constant de valor $A$ (i un morfisme a la transformació natural constant). Prendre el límit i prendre el colímit d'un diagrama són operacions $\cat{C}^{\cat{I}}\to\cat{C}$, i resulten ser, respectivament, l'adjunt per la dreta i per l'esquerra de $\Delta$.

\begin{teorema}[$\colimop\dashv\Delta\dashv\varprojlim$]\index{adjunció!límit-diagonal}
\label{teo:lim-adjunt}
Sigui $\cat{I}$ una categoria petita i suposem que $\cat{C}$ té tots els límits i colímits de forma $\cat{I}$. Fixem, per a cada diagrama de forma $\cat{I}$, un límit i un colímit; quan la família de diagrames és pròpiament gran, fem servir la convenció d'elecció global de l'apèndix~\ref{cap:mida}. La propietat universal estén aquestes tries, de manera única, a functors $\varprojlim_{\cat{I}}$ i $\colimop_{\cat{I}}$ (pas~2 de la demostració). Aleshores el functor diagonal $\Delta\colon\cat{C}\to\cat{C}^{\cat{I}}$ té adjunt per l'esquerra $\colimop_{\cat{I}}$ i adjunt per la dreta $\varprojlim_{\cat{I}}$:
\[
\colimop_{\cat{I}}\;\dashv\;\Delta\;\dashv\;\varprojlim_{\cat{I}}.
\]
\end{teorema}

\begin{prova}
Recordem que $\Delta$ envia $A$ al diagrama constant $\Delta A$, amb $(\Delta A)_i=A$ i $(\Delta A)_u=\id_A$, i un morfisme $a\colon A'\to A$ a la transformació natural $\Delta a$ de components $a$. Escrivim $(\varprojlim D,\lambda^D)$ per al límit triat de cada diagrama $D$.

\emph{Pas 1: un morfisme $\Delta A\Rightarrow D$ és un con.} Una transformació natural $\alpha\colon\Delta A\Rightarrow D$ té components $\alpha_i\colon A\to D_i$, i el seu quadrat de naturalitat per a $u\colon i\to j$ diu $D_u\comp\alpha_i=\alpha_j\comp\id_A=\alpha_j$. És exactament la condició de con. Per tant, $\Hom_{\cat{C}^{\cat{I}}}(\Delta A,D)$ és el conjunt dels cons sobre $D$ amb vèrtex $A$.

\emph{Pas 2: $\varprojlim$ és un functor.} Sigui $\beta\colon D\Rightarrow D'$ un morfisme de $\cat{C}^{\cat{I}}$. La família $\beta_i\comp\lambda^D_i$ és un con sobre $D'$, perquè $D'_u\comp\beta_i\comp\lambda^D_i=\beta_j\comp D_u\comp\lambda^D_i=\beta_j\comp\lambda^D_j$. Definim $\varprojlim\beta$ com l'única fletxa amb
\[
\lambda^{D'}_i\comp\varprojlim\beta=\beta_i\comp\lambda^D_i\qquad\text{per a tot } i.
\]
La identitat i $\varprojlim\id_D$ compleixen totes dues aquesta condició per a $\beta=\id_D$. Per a $\beta'\colon D'\Rightarrow D''$, la composició $\varprojlim\beta'\comp\varprojlim\beta$ la compleix per a $\beta'\comp\beta$. Per la unicitat de la factorització, $\varprojlim$ preserva identitats i composició.

\emph{Pas 3: la bijecció.} Definim $\theta_{A,D}\colon\Hom_{\cat{C}^{\cat{I}}}(\Delta A,D)\to\Hom_{\cat{C}}(A,\varprojlim D)$ enviant cada con $\alpha$ a l'única fletxa $\theta(\alpha)$ amb $\lambda^D_i\comp\theta(\alpha)=\alpha_i$ per a tot $i$. Per la propietat universal del límit, $\theta$ és bijectiva, amb inversa $h\mapsto(\lambda^D_i\comp h)_i=\lambda^D\comp\Delta h$.

\emph{Pas 4: naturalitat.} Per a $a\colon A'\to A$, les fletxes $\theta(\alpha\comp\Delta a)$ i $\theta(\alpha)\comp a$ compleixen totes dues $\lambda^D_i\comp(-)=\alpha_i\comp a$, i per tant coincideixen: és la naturalitat en $A$. Per a $\beta\colon D\Rightarrow D'$,
\[
\lambda^{D'}_i\comp\varprojlim\beta\comp\theta(\alpha)=\beta_i\comp\lambda^D_i\comp\theta(\alpha)=\beta_i\comp\alpha_i=\lambda^{D'}_i\comp\theta(\beta\comp\alpha),
\]
d'on $\varprojlim\beta\comp\theta(\alpha)=\theta(\beta\comp\alpha)$: és la naturalitat en $D$. Així $\Delta\dashv\varprojlim_{\cat{I}}$.

\emph{L'adjunció de l'esquerra} és la duala, pas per pas. Una transformació natural $\gamma\colon D\Rightarrow\Delta A$ té components $\gamma_i\colon D_i\to A$ amb $\gamma_j\comp D_u=\gamma_i$, és a dir, és un cocon. Si $(\colimop D,\kappa^D)$ és el colímit triat, $\colimop\beta$ és l'única fletxa amb $\colimop\beta\comp\kappa^D_i=\kappa^{D'}_i\comp\beta_i$, i és functorial per la mateixa raó. La bijecció $\Hom_{\cat{C}}(\colimop D,A)\to\Hom_{\cat{C}^{\cat{I}}}(D,\Delta A)$ és $h\mapsto(h\comp\kappa^D_i)_i$, bijectiva per la propietat universal del colímit. És natural en $A$, perquè $(a\comp h)\comp\kappa^D_i=a\comp(h\comp\kappa^D_i)$, i en $D$, perquè $h\comp\colimop\beta\comp\kappa^D_i=(h\comp\kappa^{D'}_i)\comp\beta_i$. Així $\colimop_{\cat{I}}\dashv\Delta$.
\end{prova}

\begin{observacio}
Aquesta fórmula és una de les síntesis més netes del llibre: recull en una sola línia tot el capítol de límits. En l'adjunció $\Delta\dashv\varprojlim_{\cat{I}}$, la \emph{counitat} en un diagrama $D$,
\[
\varepsilon_D\colon\Delta\bigl(\varprojlim\nolimits_{\cat{I}}D\bigr)\Rightarrow D,
\]
és exactament el con límit de $D$, és a dir, la família de projeccions; la \emph{unitat} en un objecte $A$ és el morfisme $A\to\varprojlim_{\cat{I}}\Delta A$. Dualment, en l'adjunció $\colimop_{\cat{I}}\dashv\Delta$, la \emph{unitat}
\[
\eta_D\colon D\Rightarrow\Delta\bigl(\colimop\nolimits_{\cat{I}}D\bigr)
\]
és el cocon colímit, la família d'injeccions, mentre que la \emph{counitat} és el morfisme $\colimop_{\cat{I}}\Delta A\to A$. Identificar quina de les dues transformacions és el con o el cocon universal és un bon exercici per fixar la diferència entre unitat i counitat. I la lectura és profètica: quan al capítol de subobjectes vegem l'existencial i l'universal com a adjunts de la reindexació $f^{*}$, estarem veient una versió \enquote{fibrada} d'aquesta mateixa idea ---$f^{*}$ fa el paper de $\Delta$, i els quantificadors, el de $\colimop$ i $\varprojlim$.
\end{observacio}

\section{Exponencials i currificació}

Recuperem ara la relació que vam ajornar en tractar les categories de functors: la que lliga el producte amb els objectes de functors.

\begin{proposicio}\label{prop:cat-currificacio}\index{adjunció!producte-exponencial}\index{Cat@$\cat{Cat}$!exponencials}
Per a categories petites $\cat{A},\cat{B},\cat{C}$, hi ha un isomorfisme natural
\[
\Hom_{\cat{Cat}}(\cat{A}\times\cat{B},\cat{C})\;\cong\;\Hom_{\cat{Cat}}(\cat{A},\cat{C}^{\cat{B}}).
\]
És a dir, el functor $-\times\cat{B}$ és adjunt per l'esquerra del functor $(-)^{\cat{B}}$:
\[
-\times\cat{B}\;\dashv\;(-)^{\cat{B}}.
\]
\end{proposicio}

\begin{prova}
Primer, els dos costats són functors $\cat{Cat}\to\cat{Cat}$. El functor $-\times\cat{B}$ envia $P\colon\cat{A}'\to\cat{A}$ a $P\times\id_{\cat{B}}$. El functor $(-)^{\cat{B}}$ envia $K\colon\cat{C}\to\cat{C}'$ a la composició amb $K$, $\cat{C}^{\cat{B}}\to\cat{C}'^{\cat{B}}$, que fa $H\mapsto K\comp H$ i $\eta\mapsto K\eta$ (capítol~\ref{cap:transformacions}). Si $\cat{B}$ i $\cat{C}$ són petites, $\cat{C}^{\cat{B}}$ també ho és. A $\cat{A}\times\cat{B}$ escrivim els morfismes com a parelles, que es componen component a component, i fem servir la descomposició
\[
(f,g)=(f,\id_{b'})\comp(\id_a,g)=(\id_{a'},g)\comp(f,\id_b)\qquad\text{per a } f\colon a\to a',\ g\colon b\to b'.
\tag{$\ast$}
\]

\emph{Currificació.} Sigui $F\colon\cat{A}\times\cat{B}\to\cat{C}$. Per a cada objecte $a$, $\Lambda F(a)=F(a,-)$ és el functor $\cat{B}\to\cat{C}$ que envia $b\mapsto F(a,b)$ i $g\mapsto F(\id_a,g)$. És un functor perquè $F$ ho és i $(\id_a,g')\comp(\id_a,g)=(\id_a,g'\comp g)$. Per a $f\colon a\to a'$, $\Lambda F(f)\colon F(a,-)\Rightarrow F(a',-)$ té components $F(f,\id_b)$. És natural perquè, aplicant $F$ a $(\ast)$, $F(\id_{a'},g)\comp F(f,\id_b)=F(f,g)=F(f,\id_{b'})\comp F(\id_a,g)$. Finalment, $\Lambda F\colon\cat{A}\to\cat{C}^{\cat{B}}$ és un functor: les components de $\Lambda F(\id_a)$ són identitats, i les de $\Lambda F(f'\comp f)$ són $F(f'\comp f,\id_b)=F(f',\id_b)\comp F(f,\id_b)$, que són les de la composició vertical $\Lambda F(f')\comp\Lambda F(f)$.

\emph{Descurrificació.} Sigui $H\colon\cat{A}\to\cat{C}^{\cat{B}}$. Definim $\bar H\colon\cat{A}\times\cat{B}\to\cat{C}$ per
\[
\bar H(a,b)=H(a)(b),\qquad \bar H(f,g)=H(f)_{b'}\comp H(a)(g)=H(a')(g)\comp H(f)_b,
\]
on la segona igualtat és la naturalitat de $H(f)$. Les identitats es preserven perquè $H(\id_a)$ és la transformació identitat i $H(a)$ és un functor. Per a $(f',g')\colon(a',b')\to(a'',b'')$,
\[
\begin{aligned}
\bar H(f',g')\comp\bar H(f,g)
&=H(f')_{b''}\comp H(a')(g')\comp H(f)_{b'}\comp H(a)(g)\\
&=H(f')_{b''}\comp H(f)_{b''}\comp H(a)(g')\comp H(a)(g)\\
&=H(f'\comp f)_{b''}\comp H(a)(g'\comp g),
\end{aligned}
\]
on hem fet servir la naturalitat de $H(f)$ per a $g'$, que $H$ preserva la composició (vertical, component a component) i que $H(a)$ és un functor. Això és $\bar H\bigl((f',g')\comp(f,g)\bigr)$.

\emph{Són inverses.} Per a $F$: $\overline{\Lambda F}(a,b)=F(a,b)$ i $\overline{\Lambda F}(f,g)=F(f,\id_{b'})\comp F(\id_a,g)=F(f,g)$, per $(\ast)$. Per a $H$: $\Lambda\bar H(a)$ envia $g$ a $\bar H(\id_a,g)=H(a)(g)$, de manera que $\Lambda\bar H(a)=H(a)$; i les components de $\Lambda\bar H(f)$ són $\bar H(f,\id_b)=H(f)_b$, de manera que $\Lambda\bar H(f)=H(f)$.

\emph{Naturalitat.} Per a $P\colon\cat{A}'\to\cat{A}$, el functor $\Lambda\bigl(F\comp(P\times\id_{\cat{B}})\bigr)$ envia $a'$ a $F(P(a'),-)$ i $f'$ a la transformació de components $F(P(f'),\id_b)$; és a dir, és $\Lambda F\comp P$. Per a $K\colon\cat{C}\to\cat{C}'$, $\Lambda(K\comp F)$ envia $a$ a $K\comp F(a,-)$ i $f$ a la transformació de components $K(F(f,\id_b))$, que és $K\,\Lambda F(f)$; és a dir, $\Lambda(K\comp F)=K^{\cat{B}}\comp\Lambda F$, on $K^{\cat{B}}$ és la imatge de $K$ pel functor $(-)^{\cat{B}}$. Aquestes dues igualtats són la naturalitat de $\Lambda$ en $\cat{A}$ i en $\cat{C}$. Per tant, $-\times\cat{B}\dashv(-)^{\cat{B}}$.
\end{prova}

\begin{observacio}
Aquesta adjunció ---\emph{producte $\dashv$ exponencial}--- és el prototip de les \textbf{categories cartesianes tancades}, que estudiarem al capítol següent i que són el marc categòric del càlcul lambda simplement tipat. La \enquote{currificació} és exactament la que apareix en programació funcional: transformar una funció de dues variables en una funció que retorna funcions.
\end{observacio}

\begin{resumcap}
\apartat{Idees centrals}
\begin{enumerate}
  \item Una adjunció $F\dashv G$ és una bijecció natural $\Hom(F(A),B)\cong\Hom(A,G(B))$, i dona els transposats.
  \item Equival a donar una unitat i una counitat que satisfan les identitats triangulars; la unitat expressa una propietat universal, i dualment la counitat.
  \item Exemples clau: les connexions de Galois entre preordres, $\Free\dashv U$ i, a $\cat{Set}$, la cadena $\exists_f\dashv f^{*}\dashv\forall_f$.
  \item $\colimop_{\cat{I}}\dashv\Delta\dashv\varprojlim_{\cat{I}}$ i $-\times B\dashv(-)^{B}$ mostren que moltes propietats universals són adjuncions.
  \item Els adjunts són únics llevat d'isomorfisme natural; els adjunts per la dreta preserven límits, i els adjunts per l'esquerra, colímits.
\end{enumerate}
\apartat{Errors típics} Intercanviar quin adjunt preserva límits i quin colímits; creure que un adjunt per la dreta \emph{no pot} preservar colímits, quan el teorema només dona una direcció; oblidar la naturalitat de la bijecció d'homconjunts; confondre unitat i counitat (per exemple, a $\Delta\dashv\varprojlim$ el con límit és la counitat, no la unitat); i tractar les mnemotècniques (\enquote{RAPL\,/\,LAPC}) com a terminologia establerta en lloc d'ajudes informals.
\apartat{Lectura recomanada} \cite[cap.~9]{awodey}, especialment §§9.4--9.6, per als adjunts d'ordre, els quantificadors com a adjunts i RAPL; \cite[cap.~2]{leinster} per a la unitat, la counitat i les fletxes universals; \cite[§§4.1--4.5]{riehl} per al càlcul d'adjuncions i la preservació de límits; \cite[cap.~4]{perrone-notes} per als exemples lliure--oblit, les connexions de Galois i l'adjunció entre categories i grafs; \cite[cap.~13]{barrwells} per a una exposició orientada a exemples; i \cite[cap.~IV]{maclane} per al tractament clàssic.
\end{resumcap}

\chapter[Cartesianes tancades]{\texorpdfstring{Categories cartesianes\\ tancades}{Categories cartesianes tancades}}
\label{cap:cartesianes}

\begin{objectius}
\apartat{Prerequisits} Representabilitat i lema de Yoneda (cap.~\ref{cap:yoneda}); productes i límits finits (cap.~\ref{cap:limits}); adjuncions i currificació (cap.~\ref{cap:adjuncions}). No cal cap coneixement previ de càlcul lambda: el capítol n'introdueix la semàntica necessària.
\apartat{Objectius} En acabar el capítol, el lector hauria de poder:
\begin{itemize}
  \item definir l'objecte exponencial per la seva propietat universal, com a representació del functor $A\mapsto\Hom(A\times B,C)$, i escriure les dues lleis de la currificació;
  \item reconèixer una categoria cartesiana tancada i l'exponencial com a adjunt per la dreta de $-\times B$;
  \item caracteritzar els ordres parcials cartesians tancats com a semireticles implicatius amb màxim, i relacionar-los amb les àlgebres de Heyting;
  \item situar $\cat{Set}$, els preordres, les categories de prefeixos, $\cat{Graph}$ i $\cat{Top}$ respecte del tancament cartesià;
  \item interpretar el càlcul lambda simplement tipat en una CCC, incloses les lleis $\beta$ i $\eta$ i les regles estructurals.
\end{itemize}
\end{objectius}

En tractar les adjuncions vam trobar, al final del capítol, l'adjunció entre el producte i l'exponencial a $\cat{Cat}$: $-\times\cat{B}\dashv(-)^{\cat{B}}$. Aquella era una manifestació externa d'un fenomen que ara volem fer \emph{intern} a una categoria qualsevol. Una categoria cartesiana en què, per a cada objecte $B$, el functor $-\times B$ té adjunt per la dreta s'anomena \textbf{cartesiana tancada}, i és, exactament, el marc on es pot interpretar el \textbf{càlcul lambda simplement tipat}. Aquest capítol construeix aquesta noció pas a pas, la il·lustra en els exemples fonamentals i n'explica la lectura com a semàntica de tipus i termes. És el primer pont explícit del llibre cap a la lògica categòrica.

Al capítol anterior, les adjuncions van unificar construccions aparentment diverses: objectes lliures, interiors i clausures, connectius lògics i exponencials. El patró comú, una bijecció natural d'homconjunts o, equivalentment, una unitat i una counitat que satisfan les identitats triangulars, reapareixerà com a columna vertebral de tot el que segueix. En aquest capítol el trobem en l'adjunció entre el producte i l'exponencial; més endavant, en les categories amb l'estructura suficient, veurem que els \textbf{quantificadors} $\exists$ i $\forall$ es modelen, respectivament, com els adjunts per l'esquerra i per la dreta dels functors de reindexació entre subobjectes: és la formulació categòrica de la relació entre quantificació i substitució.

\section{Categories cartesianes}

Recordem que un \textbf{producte finit} és el límit d'un diagrama amb un nombre finit d'objectes i cap morfisme no trivial. El cas de zero objectes és l'objecte terminal; el de dos, el producte binari. Una categoria té \textbf{tots els productes finits} si té objecte terminal i producte binari de qualsevol parella, ja que els productes de tres o més objectes s'obtenen iterant.

\begin{definicio}[Categoria cartesiana]\index{categoria!cartesiana}
\label{def:cartesiana}
Una categoria $\cat{C}$ és \textbf{cartesiana} si té objecte terminal $1$ i producte binari $A\times B$ de qualsevol parella d'objectes. Equivalentment, si té tots els productes finits.
\end{definicio}

En una categoria cartesiana disposem, per a cada objecte $A$, del \textbf{morfisme diagonal}
\[
\delta_A=\langle\id_A,\id_A\rangle\colon A\to A\times A,
\]
l'únic morfisme les dues projeccions del qual són la identitat, i del morfisme únic $!_A\colon A\to 1$. Escrivim la diagonal amb minúscula per no confondre-la amb el functor constant $\Delta_A$ (exemple~\ref{ex:functor-constant}). Aquestes dades no són estructura addicional: queden \emph{determinades} per les propietats universals.

Els productes escollits determinen també la functorialitat del producte. Donats $f\colon A\to A'$ i $g\colon B\to B'$, posem
\[
f\times g=\langle f\comp\pi_1,\ g\comp\pi_2\rangle\colon A\times B\to A'\times B'.
\]
En particular, per a $B$ fix, l'assignació $A\mapsto A\times B$ defineix el functor $-\times B$, tal com precisem ara.

\begin{proposicio}[El producte per un objecte és functorial]\index{producte!functorialitat}
\label{prop:producte-functor}
Sigui $\cat{C}$ cartesiana i $B$ un objecte fix. L'assignació $A\mapsto A\times B$ s'estén a un functor
\[
-\times B\colon\cat{C}\to\cat{C},
\]
que envia $f\colon A\to A'$ al morfisme $f\times\id_B=\langle f\comp\pi_1,\pi_2\rangle\colon A\times B\to A'\times B$.
\end{proposicio}

\begin{prova}
Cal comprovar que $-\times B$ respecta identitats i composicions. Per a la identitat, $\id_A\times\id_B=\langle\pi_1,\pi_2\rangle=\id_{A\times B}$ per la unicitat de la propietat universal del producte. Per a la composició, siguin $f\colon A\to A'$ i $g\colon A'\to A''$. Tant $(g\comp f)\times\id_B$ com $(g\times\id_B)\comp(f\times\id_B)$ tenen com a projecció primera $g\comp f\comp\pi_1$ i com a projecció segona $\pi_2$; per unicitat, coincideixen.
\end{prova}

\section{Objectes exponencials}

En una categoria cartesiana, la pregunta clau és si el functor $-\times B$ té adjunt per la dreta. Per a $B$ i $C$ fixos, l'exponencial $C^B$, quan existeix, representa el functor
\[
A\longmapsto\Hom(A\times B,C)
\]
(capítol~\ref{cap:yoneda}); quan existeix per a tot $C$, defineix l'adjunt per la dreta de $-\times B$.

\begin{definicio}[Objecte exponencial]\index{objecte!exponencial}\index{exponencial}
\label{def:exponencial}
Siguin $B,C$ objectes d'una categoria cartesiana $\cat{C}$. Un \textbf{objecte exponencial} $C^{B}$ és un objecte dotat d'un morfisme d'\textbf{avaluació}
\[
\ev\colon C^{B}\times B\to C
\]
amb la propietat universal següent: per a tot objecte $A$ i tot morfisme $f\colon A\times B\to C$, existeix un únic morfisme $\lambda f\colon A\to C^{B}$, la \textbf{transposada} o \textbf{currificació} de $f$, tal que el triangle
\[
\begin{tikzpicture}[baseline=(m.center)]
  \matrix (m) [matrix of math nodes, column sep=3.4em, row sep=2.6em] {
    A\times B & C^{B}\times B \\
              & C \\
  };
  \draw[-{Stealth[scale=1.1]}] (m-1-1) -- node[above] {$\lambda f\times\id_B$} (m-1-2);
  \draw[-{Stealth[scale=1.1]}] (m-1-1) -- node[below left] {$f$} (m-2-2);
  \draw[-{Stealth[scale=1.1]}] (m-1-2) -- node[right] {$\ev$} (m-2-2);
\end{tikzpicture}
\]
commuta, és a dir, $\ev\comp(\lambda f\times\id_B)=f$.
\end{definicio}

L'avaluació formalitza la idea d'\enquote{aplicar una funció al seu argument}, i la currificació, la de \enquote{fixar un paràmetre}. Com tota representació, l'exponencial és únic llevat d'un únic isomorfisme compatible amb l'avaluació: si $(E,\ev)$ i $(E',\ev')$ són exponencials de $C$ per $B$, hi ha un únic isomorfisme $u\colon E\to E'$ amb $\ev'\comp(u\times\id_B)=\ev$.

\begin{observacio}[Les dues lleis de la currificació]\label{obs:lleis-currificacio}\index{currificació!lleis}
La correspondència $f\mapsto\lambda f$ és una bijecció entre $\Hom(A\times B,C)$ i $\Hom(A,C^B)$, amb inversa $g\mapsto\ev\comp(g\times\id_B)$. Les dues composicions donen identitats:
\[
\ev\comp(\lambda f\times\id_B)=f,
\qquad\qquad
\lambda\bigl(\ev\comp(g\times\id_B)\bigr)=g,
\]
per a tot $f\colon A\times B\to C$ i tot $g\colon A\to C^B$. La primera és la condició de la definició. La segona en surt per unicitat: $g$ és un morfisme $A\to C^B$ que fa commutar el triangle per a $f=\ev\comp(g\times\id_B)$, i l'únic que ho fa és $\lambda f$. Aquestes dues lleis són la versió categòrica de les lleis $\beta$ i $\eta$ del càlcul lambda, que veurem més endavant; la naturalitat de la bijecció la converteix en una adjunció.
\end{observacio}

\begin{proposicio}[L'exponencial és un adjunt per la dreta]\index{adjunció!producte-exponencial}
\label{prop:adjuncio-exp}
Si per a un objecte fix $B$ existeix l'exponencial $C^{B}$ per a tot $C$, aleshores l'assignació $C\mapsto C^{B}$ s'estén a un functor $(-)^{B}\colon\cat{C}\to\cat{C}$ i hi ha un isomorfisme natural
\[
\Hom(A\times B,\,C)\;\cong\;\Hom(A,\,C^{B}),
\]
natural en $A$ i en $C$. És a dir, $-\times B\dashv(-)^{B}$.
\end{proposicio}

\begin{prova}
La currificació $f\mapsto\lambda f$ és la bijecció buscada (observació~\ref{obs:lleis-currificacio}). La naturalitat en $A$ es comprova precomponent amb $a\colon A'\to A$ i usant la functorialitat de $-\times B$ (la proposició~\ref{prop:producte-functor}); la naturalitat en $C$ defineix, de fet, l'acció de $(-)^{B}$ sobre morfismes: donat $c\colon C\to C'$, es posa $c^{B}=\lambda\bigl(c\comp\ev\bigr)\colon C^{B}\to C'^{B}$. Que això respecta identitats i composicions se segueix, un cop més, de la unicitat de la transposada.
\end{prova}

\begin{notacio}
Escrivim indistintament $C^{B}$ o $[B,C]$ per a l'exponencial. La transposada de $f$ es denota $\lambda f$ o $\widehat f$; la seva inversa, $\ev\comp(g\times\id_B)$, es denota $\widetilde g$ i s'anomena \textbf{descurrificació} de $g$.
\end{notacio}

\section{La definició de categoria cartesiana tancada}

\begin{definicio}[Categoria cartesiana tancada]\index{categoria!cartesiana tancada}\index{CCC}
\label{def:ccc}
Una categoria $\cat{C}$ és \textbf{cartesiana tancada} (abreujat \textbf{CCC}, de l'anglès \emph{cartesian closed category}) si és cartesiana i, a més, per a cada objecte $B$, el functor $-\times B\colon\cat{C}\to\cat{C}$ té adjunt per la dreta $(-)^{B}$. Explícitament, $\cat{C}$ té:
\begin{enumerate}
\item un objecte terminal $1$;
\item productes binaris $A\times B$;
\item exponencials $C^{B}$ amb avaluació i currificació.
\end{enumerate}
\end{definicio}

\begin{observacio}[L'exponencial com a bifunctor]\label{obs:exponencial-bifunctor}\index{exponencial!bifunctor}
En una CCC, els exponencials s'organitzen en un bifunctor
\[
(-)^{(-)}\colon\cat{C}\op\times\cat{C}\longrightarrow\cat{C},
\qquad (B,C)\longmapsto C^{B},
\]
contravariant en $B$ i covariant en $C$, i la bijecció
\[
\Hom(A\times B,C)\;\cong\;\Hom(A,C^{B})
\]
és natural en les tres variables $A$, $B$ i $C$. L'acció en $C$ és la de la proposició~\ref{prop:adjuncio-exp}; la construcció de l'acció en $B$, mitjançant la transposada de $\ev\comp(\id\times f)$, es proposa als Exercicis.
\end{observacio}

La condició es pot resumir dient que en una CCC les tres operacions \enquote{no fer res} (terminal), \enquote{aparellar} (producte) i \enquote{formar espais de funcions} (exponencial) hi són totes i estan lligades per adjuncions. Aquesta és exactament l'estructura que necessita un llenguatge de tipus amb producte de tipus i tipus funció.

\begin{observacio}
La tancadura és una propietat \emph{objecte a objecte}: cal l'exponencial per a \emph{cada} base $B$. Una categoria pot tenir alguns exponencials i no ser tancada. La demanda uniforme ---un adjunt per la dreta de $-\times B$ per a tot $B$--- és el que fa que la currificació estigui sempre disponible.
\end{observacio}

\section{Exemples fonamentals}

\begin{exemple}[$\cat{Set}$]\index{Set@$\cat{Set}$!cartesiana tancada}
La categoria $\cat{Set}$ és cartesiana tancada. L'objecte terminal és qualsevol conjunt d'un element; el producte és el producte cartesià; i l'exponencial $C^{B}$ és el conjunt de \emph{totes} les aplicacions $B\to C$, amb l'avaluació $\ev(g,b)=g(b)$. La currificació $\lambda f\colon A\to C^{B}$ d'una aplicació $f\colon A\times B\to C$ és $\lambda f(a)=(b\mapsto f(a,b))$: exactament la currificació de la programació funcional. Aquest és el model prototípic i el que dona nom a tota la construcció.
\end{exemple}

\begin{exemple}[$\cat{Cat}$: transformacions naturals com a functors]\label{ex:cat-ccc}\index{Cat@$\cat{Cat}$!cartesiana tancada}\index{avaluació!a Cat@a $\cat{Cat}$}
La categoria $\cat{Cat}$ de les categories petites és cartesiana tancada: el terminal és $\cat{1}$, el producte és la categoria producte i l'exponencial de $\cat{C}$ per $\cat{B}$ és la categoria de functors $\cat{C}^{\cat{B}}$ (proposició~\ref{prop:cat-currificacio}). Fem explícits l'avaluació i la currificació, i després els llegim en un cas concret.

\emph{Avaluació.} El functor $\ev\colon\cat{C}^{\cat{B}}\times\cat{B}\to\cat{C}$ és
\[
\ev(H,b)=H(b),\qquad \ev(\eta,g)=\eta_{b'}\comp H(g)=H'(g)\comp\eta_b,
\]
per a $\eta\colon H\Rightarrow H'$ i $g\colon b\to b'$; la segona igualtat és la naturalitat de $\eta$. És la descurrificació de $\id_{\cat{C}^{\cat{B}}}$ a la demostració de la proposició~\ref{prop:cat-currificacio}, i d'allà en surt que és un functor.

\emph{Currificació i llei $\beta$.} La currificació de $F\colon\cat{A}\times\cat{B}\to\cat{C}$ és el functor $\Lambda F\colon\cat{A}\to\cat{C}^{\cat{B}}$ d'aquella demostració. El triangle de la definició~\ref{def:exponencial} commuta: sobre objectes, $\ev\bigl(\Lambda F(a),b\bigr)=F(a,b)$; sobre un morfisme $(f,g)$,
\[
\ev\bigl(\Lambda F(f),g\bigr)=\Lambda F(f)_{b'}\comp\Lambda F(a)(g)=F(f,\id_{b'})\comp F(\id_a,g)=F(f,g).
\]
Així, $\ev\comp(\Lambda F\times\id_{\cat{B}})=F$: avaluar la currificació torna la funció de dues variables.

\emph{Un cas concret: $\cat{B}=\cat{2}$.} La categoria $\cat{C}^{\cat{2}}$ és la categoria de fletxes de $\cat{C}$: els objectes són els morfismes de $\cat{C}$ i els morfismes, els quadrats commutatius (Pràctica del capítol~\ref{cap:transformacions}). L'avaluació envia $(h,0)$ al domini de $h$, $(h,1)$ al codomini i $(h,u)$ a $h$ mateix, on $u\colon0\to1$. Un functor $\Phi\colon\cat{A}\times\cat{2}\to\cat{C}$ consisteix en tres dades:
\begin{itemize}
\item dos functors $F=\Phi(-,0)$ i $G=\Phi(-,1)$ de $\cat{A}$ a $\cat{C}$;
\item una família $\alpha_a=\Phi(\id_a,u)\colon F(a)\to G(a)$;
\item la naturalitat $G(f)\comp\alpha_a=\alpha_{a'}\comp F(f)$, que no és cap condició nova: és $\Phi$ aplicat a les dues descomposicions de $(f,u)$.
\end{itemize}
És a dir, un functor $\cat{A}\times\cat{2}\to\cat{C}$ \emph{és} una transformació natural $\alpha\colon F\Rightarrow G$. La seva currificació $\Lambda\Phi\colon\cat{A}\to\cat{C}^{\cat{2}}$ envia $a$ a la fletxa $\alpha_a$ i $f$ al quadrat de naturalitat de $f$. L'avaluació recupera les tres peces: $\ev(\Lambda\Phi(a),0)=F(a)$, $\ev(\Lambda\Phi(a),1)=G(a)$ i $\ev(\Lambda\Phi(a),u)=\alpha_a$. Tenim, doncs, tres lectures d'una mateixa dada,
\[
\alpha\colon F\Rightarrow G\qquad\longleftrightarrow\qquad\cat{A}\times\cat{2}\to\cat{C}\qquad\longleftrightarrow\qquad\cat{A}\to\cat{C}^{\cat{2}},
\]
i el pas de la segona a la tercera és la currificació, amb l'avaluació com a inversa.
\end{exemple}

\begin{exemple}[Preordres cartesians tancats]\index{semireticle implicatiu}\index{àlgebra de Heyting}\index{Heyting}
Un preordre amb ínfims finits, vist com a categoria, és cartesià: el terminal és l'element màxim $\top$ i el producte és l'ínfim $x\wedge y$. Té exponencials si i només si, per a cada $b$, l'operació $-\wedge b$ té adjunt per la dreta; aquest adjunt és la \textbf{implicació} $b\Rightarrow-$, caracteritzada per
\[
a\wedge b\leq c\quad\Longleftrightarrow\quad a\leq(b\Rightarrow c).
\]
Un ordre parcial amb ínfims finits és, doncs, cartesià, i és cartesià tancat si i només si és un \textbf{semireticle implicatiu amb màxim}: un $\wedge$-semireticle amb element màxim $\top$ i una operació d'implicació que satisfà l'equivalència anterior (també anomenat \emph{semireticle de Brouwer}). Per a un preordre, la mateixa caracterització val després de passar al quocient per l'equivalència
\[
a\sim b\quad\Longleftrightarrow\quad a\leq b\ \text{i}\ b\leq a,
\]
perquè en la terminologia algebraica habitual un semireticle és un ordre parcial. Aquesta estructura captura la part $(\top,\wedge,\Rightarrow)$ de la lògica.

Convé no confondre aquesta estructura amb una àlgebra de Heyting. Una \textbf{àlgebra de Heyting}\index{àlgebra de Heyting} és un \emph{reticle distributiu acotat} ---amb ínfims finits $\top,\wedge$ i \emph{suprems finits} $\bot,\vee$--- proveït d'una implicació que satisfà l'equivalència anterior. El tancament cartesià per si sol només dona la part multiplicativa: no garanteix ni els joins $\vee$ ni el mínim $\bot$. Reservem, doncs, el nom \emph{àlgebra de Heyting} per al cas amb suprems finits; un semireticle implicatiu amb màxim que a més tingui suprems finits (i sigui distributiu, cosa automàtica en presència de la implicació) és una àlgebra de Heyting. Aquest exemple és decisiu per a la lògica: la implicació intuïcionista \emph{és} l'exponencial, i el tancament cartesià és la traducció algebraica de disposar del connectiu $\to$, mentre que la disjunció $\vee$ correspon als suprems.
\end{exemple}

\begin{observacio}[Atenció: cartesià tancat no vol dir àlgebra de Heyting]\label{obs:ccc-no-heyting}\index{àlgebra de Heyting!i preordres cartesians tancats}
En un preordre, ser cartesià tancat dona $\top$, $\wedge$ i $\Rightarrow$, i res més. Un exemple mínim mostra què pot faltar. Sigui $\mathbb Z\cup\{+\infty\}$ amb l'ordre usual. És una cadena amb màxim $+\infty$, amb ínfims $a\wedge b=\min(a,b)$, i té implicació
\[
b\Rightarrow c=\begin{cases}+\infty&\text{si } b\leq c,\\ c&\text{si } b>c.\end{cases}
\]
En efecte, si $b\leq c$ tenim $a\wedge b\leq c$ per a tot $a$; si $b>c$, tenim $\min(a,b)\leq c$ si i només si $a\leq c$. Per tant és un ordre parcial cartesià tancat. Però no té element mínim, de manera que no és una àlgebra de Heyting. Aquí no existeixen ni $\bot$, ni la negació $\neg b=(b\Rightarrow\bot)$, ni el principi \emph{ex falso}; no hi ha, doncs, cap interpretació de la falsedat. En altres exemples pot faltar el suprem binari $\vee$, i amb ell la disjunció.

La relació exacta és aquesta: una àlgebra de Heyting és un ordre parcial cartesià tancat que, a més, té coproductes finits ($\bot$ i $\vee$). La distributivitat ve aleshores gratis, perquè $-\wedge b$, com a adjunt per l'esquerra, preserva els suprems (capítol~\ref{cap:adjuncions}). L'estructura de Heyting també condiciona els morfismes: un morfisme d'àlgebres de Heyting ha de preservar $\bot$ i $\vee$, cosa que no demana un functor que només preservi l'estructura cartesiana tancada. En lògica, el tancament cartesià modela el fragment $(\top,\wedge,\to)$; la falsedat i la disjunció són estructura addicional.
\end{observacio}

\begin{exemple}[Categories de prefeixos]\index{prefeix@prefeix!cartesiana tancada}
Per a qualsevol categoria petita $\cat{C}$, la categoria de prefeixos $\widehat{\cat{C}}=\cat{Set}^{\cat{C}\op}$ (exemple~\ref{ex:categoria-prefeixos}) és cartesiana tancada. Els productes i el terminal es calculen \emph{objecte a objecte} a $\cat{Set}$ (proposició~\ref{prop:limits-punt-a-punt}), i l'exponencial es defineix, mitjançant el lema de Yoneda, per $Q^{P}(A)=\Nat(\Hom(-,A)\times P,\,Q)$. Aquestes categories seran el terreny natural dels topos de prefeixos.
\end{exemple}

\begin{exemple}[Grafs]\index{graf dirigit!cartesiana tancada}
Al capítol~\ref{cap:transformacions} vam veure que $\cat{Graph}\cong\cat{Set}^{\cat{P}}$, on $\cat{P}$ té dos objectes $E,V$ i dues fletxes $s,t\colon E\to V$. Com que $\cat{Set}^{\cat{P}}=\cat{Set}^{(\cat{P}\op)\op}$, els grafs són exactament els prefeixos sobre $\cat{P}\op$: un prefeix $G\colon(\cat{P}\op)\op\to\cat{Set}$ és la dada de dos conjunts, de vèrtexs i d'arestes, i de les aplicacions origen i destí. Per tant, com a cas particular de l'exemple anterior, $\cat{Graph}$ és cartesiana tancada. No calcularem aquí l'exponencial de dos grafs; n'hi ha prou de saber que existeix i que es pot obtenir amb la fórmula de Yoneda de l'exemple anterior.
\end{exemple}

\begin{observacio}[Un no-exemple instructiu]
La categoria $\cat{Top}$ dels espais topològics \emph{no} és cartesiana tancada: no sempre existeix una topologia sobre l'espai de funcions contínues que faci de l'avaluació un morfisme universal. Aquesta mancança va motivar històricament la introducció de categories \enquote{convenients} d'espais topològics, com la dels espais Hausdorff compactament generats, amb les construccions categòriques adequades, que sí que és cartesiana tancada. El tancament cartesià és, doncs, una propietat delicada i valuosa, no automàtica.
\end{observacio}

\section[Semàntica del càlcul lambda simplement tipat]{Semàntica del càlcul lambda\\ simplement tipat}

El motiu pel qual les CCC són centrals en lògica categòrica és la seva correspondència amb el \textbf{càlcul lambda simplement tipat}. Descrivim la traducció a nivell introductori; és una de les formes de l'anomenada \emph{correspondència de Curry--Howard--Lambek}.

\subsection{Tipus com a objectes}

Fixem un conjunt de tipus base. Els \textbf{tipus} es construeixen amb dues operacions: el \textbf{producte de tipus} $\sigma\times\tau$ i el \textbf{tipus funció} $\sigma\to\tau$, més un tipus unitat $\mathbf{1}$. En una CCC $\cat{C}$, la interpretació $\llbracket-\rrbracket$ assigna:
\[
\llbracket\mathbf{1}\rrbracket=1,\qquad
\llbracket\sigma\times\tau\rrbracket=\llbracket\sigma\rrbracket\times\llbracket\tau\rrbracket,\qquad
\llbracket\sigma\to\tau\rrbracket=\llbracket\tau\rrbracket^{\llbracket\sigma\rrbracket}.
\]
Cada tipus esdevé un objecte; el tipus funció esdevé un exponencial.

\subsection{Termes com a morfismes}

Un \textbf{context} $\Gamma=x_1:\sigma_1,\dots,x_n:\sigma_n$ s'interpreta, per recursió, com un producte:
\[
\llbracket\varnothing\rrbracket=1,
\qquad
\llbracket\Gamma,x:\sigma\rrbracket=\llbracket\Gamma\rrbracket\times\llbracket\sigma\rrbracket.
\]
Un \textbf{terme} $\Gamma\vdash t:\tau$, amb variables lliures en el context $\Gamma$, s'interpreta com un morfisme
\[
\llbracket t\rrbracket\colon\llbracket\Gamma\rrbracket\longrightarrow\llbracket\tau\rrbracket.
\]
Les regles de formació de termes es corresponen amb les operacions de la CCC:
\begin{itemize}
\item una \textbf{variable} $x_i$ s'interpreta amb la projecció $\pi_i$;
\item una \textbf{parella} $\langle t,u\rangle$, amb el morfisme aparellament $\langle\llbracket t\rrbracket,\llbracket u\rrbracket\rangle$ cap a un producte;
\item una \textbf{aplicació} $t\,u$, amb la composició de l'avaluació i l'aparellament de $\llbracket t\rrbracket$ i $\llbracket u\rrbracket$;
\item una \textbf{abstracció} $\lambda x{:}\sigma.\,t$, amb la currificació $\lambda\llbracket t\rrbracket$: si $\Gamma,x:\sigma\vdash t:\tau$, aleshores $\llbracket t\rrbracket\colon\llbracket\Gamma\rrbracket\times\llbracket\sigma\rrbracket\to\llbracket\tau\rrbracket$, i la seva transposada és un morfisme $\llbracket\Gamma\rrbracket\to\llbracket\tau\rrbracket^{\llbracket\sigma\rrbracket}$, exactament el tipus de $\lambda x{:}\sigma.\,t$.
\end{itemize}
La correspondència entre abstracció i currificació és el cor de la traducció: lligar una variable és transposar per l'adjunció, i aplicar una funció és avaluar. Les regles de conversió del càlcul es tornen igualtats de morfismes:
\[
(\beta)\quad(\lambda x.\,t)\,u = t[u/x]
\qquad\text{esdevé}\qquad
\ev\comp\langle\lambda\llbracket t\rrbracket,\llbracket u\rrbracket\rangle=\llbracket t\rrbracket\comp\langle\id,\llbracket u\rrbracket\rangle,
\]
que no és literalment l'equació de la definició~\ref{def:exponencial}, sinó que \emph{se'n dedueix}. Si $f=\llbracket t\rrbracket\colon\llbracket\Gamma\rrbracket\times\llbracket\sigma\rrbracket\to\llbracket\tau\rrbracket$, la primera llei de la currificació (observació~\ref{obs:lleis-currificacio}) dona $\ev\comp(\lambda f\times\id)=f$, i precomponent amb $\langle\id_{\llbracket\Gamma\rrbracket},\llbracket u\rrbracket\rangle\colon\llbracket\Gamma\rrbracket\to\llbracket\Gamma\rrbracket\times\llbracket\sigma\rrbracket$ s'obté
\[
\ev\comp\langle\lambda f,\llbracket u\rrbracket\rangle
=\ev\comp(\lambda f\times\id)\comp\langle\id,\llbracket u\rrbracket\rangle
=f\comp\langle\id,\llbracket u\rrbracket\rangle,
\]
que és la interpretació de la substitució $t[u/x]$. La regla $(\eta)$, $\lambda x.\,(t\,x)=t$, es correspon amb la segona llei, és a dir, amb la unicitat de la transposada.

\begin{observacio}[L'estructura cartesiana i les regles estructurals]\index{regles estructurals}
Els morfismes canònics del producte interpreten les \emph{regles estructurals} de la lògica i del càlcul. Una projecció $\llbracket\Gamma\rrbracket\times\llbracket\sigma\rrbracket\to\llbracket\Gamma\rrbracket$ permet \emph{oblidar} una variable (debilitament); la diagonal $\delta\colon\llbracket\sigma\rrbracket\to\llbracket\sigma\rrbracket\times\llbracket\sigma\rrbracket$ permet \emph{reutilitzar-la} (contracció); i l'isomorfisme de simetria $\llbracket\sigma\rrbracket\times\llbracket\tau\rrbracket\cong\llbracket\tau\rrbracket\times\llbracket\sigma\rrbracket$ permet \emph{intercanviar-ne} l'ordre (intercanvi). És per això que la diagonal, introduïda al començament del capítol, no és un detall aïllat: l'estructura \emph{cartesiana} és exactament la que fa admissibles aquestes tres regles.
\end{observacio}

\begin{observacio}[Curry--Howard--Lambek]\index{Curry--Howard--Lambek}
La correspondència es pot fer precisa: després de fixar una noció adequada de teoria del càlcul lambda simplement tipat, una estructura cartesiana tancada escollida i els morfismes que la preserven, s'obté una \emph{equivalència} de categories entre les dues presentacions. En aquest sentit, tipus, termes i equacions del càlcul lambda corresponen a objectes, morfismes i igualtats de morfismes en una CCC. Aquesta és la primera de les grans traduccions \enquote{lògica $=$ categoria} que estructuraran l'estudi posterior de la lògica categòrica; les següents afegiran, sobre aquesta base, els subobjectes per als predicats i la reindexació per a la quantificació.
\end{observacio}

\begin{resumcap}
\apartat{Idees centrals}
\begin{enumerate}
  \item L'objecte exponencial $C^{B}$ representa el functor $A\mapsto\Hom(A\times B,C)$; la currificació és la bijecció corresponent, amb les dues lleis $\ev\comp(\lambda f\times\id_B)=f$ i $\lambda(\ev\comp(g\times\id_B))=g$.
  \item Una categoria cartesiana tancada té productes finits i tots els exponencials; els exponencials s'organitzen en un bifunctor contravariant en l'exponent i covariant en la base.
  \item Un ordre parcial cartesià tancat és un semireticle implicatiu amb màxim; per a un preordre, aquesta descripció val després de passar al quocient per l'equivalència $a\sim b$. Si, a més, hi ha suprems finits, s'obté una àlgebra de Heyting; sense ells, no (per exemple, $\mathbb Z\cup\{+\infty\}$ no té $\bot$).
  \item $\cat{Set}$, $\cat{Cat}$, les categories de prefeixos i, en particular, $\cat{Graph}$ són cartesianes tancades. A $\cat{Cat}$, l'avaluació i la currificació identifiquen les transformacions naturals amb els functors $\cat{A}\times\cat{2}\to\cat{C}$ i $\cat{A}\to\cat{C}^{\cat{2}}$; $\cat{Top}$ no ho és en general.
  \item El càlcul lambda simplement tipat s'interpreta en tota CCC: els contextos com a productes, l'abstracció com a currificació, i les lleis $\beta$ i $\eta$ com a conseqüències de les dues lleis de la currificació.
\end{enumerate}
\apartat{Errors típics} Identificar tota categoria cartesiana tancada amb una àlgebra de Heyting, fins i tot quan és un ordre parcial (poden faltar $\bot$ i els joins, i amb ells la negació i la disjunció); atribuir la functorialitat del producte a la diagonal en lloc de a la seva propietat universal; i tractar la correspondència de Curry--Howard--Lambek com una igualtat en lloc d'una equivalència amb estructura triada.
\apartat{Lectura recomanada} \cite[cap.~6]{awodey} per a exponencials i CCC; \cite{lambekscott} per a la correspondència amb el càlcul lambda; \cite[cap.~1]{maclanemoerdijk} per als exemples de prefeixos; \cite[cap.~6]{barrwells}, que presenta les CCC i el càlcul lambda tipat i calcula explícitament l'exponencial de dos grafs (exemple~6.1.12); i \cite[§4.3]{riehl} per a la formulació de l'exponencial com a adjunt; \cite[§11]{streicher}, per als exponencials en categories de prefeixos (§11.1) i la semàntica del càlcul lambda tipat (§11.2); i \cite[article~V, sessions~30--31]{lawvereschanuel}, per a una introducció molt elemental als objectes de morfismes.
\end{resumcap}

\chapter[Subobjectes, imatges i factoritzacions]{\texorpdfstring{Subobjectes, imatges\\ i factoritzacions}{Subobjectes, imatges i factoritzacions}}
\label{cap:subobjectes}

\begin{objectius}
\apartat{Prerequisits} Monomorfismes i epimorfismes (cap.~\ref{cap:categories}); pullbacks, coequalitzadors, límits finits i reindexació per pullback (cap.~\ref{cap:limits}); adjuncions, especialment entre preordres (cap.~\ref{cap:adjuncions}).
\apartat{Objectius} En acabar el capítol, el lector hauria de poder:
\begin{itemize}
  \item definir els subobjectes i entendre $\Sub(A)$ com un ordre parcial de predicats sobre $A$;
  \item interpretar la reindexació $f^{*}$ com a substitució i organitzar-la, en una categoria ben potenciada, en un functor contravariant $\Sub\colon\cat{C}\op\to\cat{Poset}$;
  \item interpretar la intersecció de subobjectes com a pullback;
  \item definir epimorfisme regular, parella nucli i factorització imatge;
  \item reconèixer una categoria regular;
  \item demostrar l'adjunció $\exists_f\dashv f^{*}$.
\end{itemize}
\end{objectius}

El capítol anterior ens ha donat la semàntica dels tipus i els termes. Ara construïm la peça complementària: la interpretació de les \textbf{proposicions} i la \textbf{substitució}. La idea directriu és senzilla i profunda: un predicat sobre un objecte $A$ ---una propietat que els \enquote{elements} de $A$ poden tenir o no--- es representa categòricament per un \textbf{subobjecte} de $A$. A $\cat{Set}$, un subconjunt es representa per la seva inclusió, que és injectiva; en una categoria arbitrària, el paper de les injeccions el fan els monomorfismes. Els subobjectes d'un objecte formen un ordre parcial $\Sub(A)$, i la reindexació per pullback, que ja vam construir, hi actua com la substitució.

A la segona meitat del capítol construïm la \textbf{factorització imatge} d'un morfisme. A $\cat{Set}$, tota aplicació es factoritza com una aplicació exhaustiva sobre la seva imatge seguida de la inclusió d'aquesta imatge; en una categoria arbitrària, aquesta idea es reformula mitjançant epimorfismes regulars i monomorfismes. Després definim les \textbf{categories regulars} com aquelles on aquesta factorització és estable per canvi de base, i n'obtenim una conseqüència notable: l'\textbf{existencial} apareix com a \textbf{adjunt per l'esquerra} de la reindexació. L'universal $\forall_f$ requereix estructura addicional i es tractarà al volum de lògica categòrica. Amb això, el volum tanca la infraestructura categòrica que la lògica farà servir.

\section{Subobjectes}

Volem capturar la idea de \enquote{subconjunt} de manera purament categòrica. A $\cat{Set}$, un subconjunt $S\subseteq A$ es pot presentar per la seva injecció $S\hookrightarrow A$, que és un monomorfisme. Però dos monomorfismes diferents poden tenir la mateixa imatge: una injecció i una reparametrització bijectiva del domini defineixen \enquote{el mateix} subconjunt. Cal, doncs, identificar monomorfismes que difereixen per un isomorfisme del domini.

\begin{definicio}[Subobjecte]\index{subobjecte}
\label{def:subobjecte}
Sigui $A$ un objecte d'una categoria $\cat{C}$. Considerem els monomorfismes amb codomini $A$. Diem que dos monomorfismes $m\colon S\rightarrowtail A$ i $m'\colon S'\rightarrowtail A$ són \textbf{equivalents}, i escrivim $m\sim m'$, si existeix un isomorfisme $\varphi\colon S\to S'$ amb $m'\comp\varphi=m$:
\[
\begin{tikzpicture}[baseline=(m.center)]
  \matrix (m) [matrix of math nodes, column sep=2.8em, row sep=2.0em] {
    S & & S' \\
      & A & \\
  };
  \draw[-{Stealth[scale=1.1]}] (m-1-1) -- node[above] {$\varphi$} node[below] {$\sim$} (m-1-3);
  \draw[-{Stealth[scale=1.1]}] (m-1-1) -- node[below left] {$m$} (m-2-2);
  \draw[-{Stealth[scale=1.1]}] (m-1-3) -- node[below right] {$m'$} (m-2-2);
\end{tikzpicture}
\]
Un \textbf{subobjecte} de $A$ és una classe d'equivalència de monomorfismes per aquesta relació. Escrivim $[m]$ per al subobjecte representat per $m$.
\end{definicio}

\begin{observacio}
La relació $\sim$ és d'equivalència: la reflexivitat usa $\varphi=\id$, la simetria usa $\varphi^{-1}$ (que existeix perquè $\varphi$ és isomorfisme), i la transitivitat, la composició d'isomorfismes. Cal notar que l'isomorfisme $\varphi$, si existeix, és \emph{únic}: de $m'\comp\varphi=m=m'\comp\varphi'$ i el fet que $m'$ és mono se segueix $\varphi=\varphi'$. Els subobjectes estan, doncs, ben definits sense ambigüitat.
\end{observacio}

\begin{exemple}
A $\cat{Set}$, els subobjectes de $A$ es corresponen exactament amb els \textbf{subconjunts} de $A$. Un monomorfisme $m\colon S\rightarrowtail A$ és una aplicació injectiva, i la seva classe d'equivalència queda determinada per la imatge $m(S)\subseteq A$: dos monos són equivalents si i només si tenen la mateixa imatge. Així, $\Sub(A)\cong\mathcal{P}(A)$, el conjunt de parts de $A$.
\end{exemple}

\section{L'ordre de subobjectes}

Els subobjectes d'un objecte no formen només un conjunt: hi ha una noció natural d'inclusió entre ells.

\begin{definicio}[Ordre entre subobjectes]\index{subobjecte!ordre}\index{Sub@$\Sub(A)$}
\label{def:ordre-sub}
Donats dos monomorfismes $m\colon S\rightarrowtail A$ i $n\colon T\rightarrowtail A$, diem que $[m]\leq[n]$ si $m$ \textbf{es factoritza a través de} $n$, és a dir, si existeix un morfisme $h\colon S\to T$ amb $n\comp h=m$:
\[
\begin{tikzpicture}[baseline=(m.center)]
  \matrix (m) [matrix of math nodes, column sep=2.8em, row sep=2.0em] {
    S & & T \\
      & A & \\
  };
  \draw[-{Stealth[scale=1.1]}] (m-1-1) -- node[above] {$h$} (m-1-3);
  \draw[-{Stealth[scale=1.1]}] (m-1-1) -- node[below left] {$m$} (m-2-2);
  \draw[-{Stealth[scale=1.1]}] (m-1-3) -- node[below right] {$n$} (m-2-2);
\end{tikzpicture}
\]
\end{definicio}

\begin{proposicio}[$\Sub(A)$ és un ordre parcial]
\label{prop:sub-preordre}
La relació $\leq$ de la definició~\ref{def:ordre-sub} està ben definida sobre classes d'equivalència i és reflexiva i transitiva. A més, és antisimètrica: $[m]\leq[n]$ i $[n]\leq[m]$ impliquen $[m]=[n]$. Per tant, $\Sub(A)$ és un \textbf{ordre parcial}. (Sobre els monomorfismes mateixos, abans de passar a classes d'equivalència, la factorització només defineix un preordre.)
\end{proposicio}

\begin{prova}
El morfisme de factorització $h$, si existeix, és únic, perquè $n$ és mono: de $n\comp h=m=n\comp h'$ se segueix $h=h'$. A més, $h$ és automàticament un monomorfisme: si $h\comp u=h\comp v$, aleshores $m\comp u=n\comp h\comp u=n\comp h\comp v=m\comp v$, i com que $m$ és mono, $u=v$. La bona definició sobre classes és immediata: substituir $m$ o $n$ per equivalents compon $h$ amb isomorfismes. La reflexivitat usa $h=\id$; la transitivitat, la composició de factoritzacions.

Per a l'antisimetria, suposem $[m]\leq[n]$ i $[n]\leq[m]$, amb factoritzacions $h\colon S\to T$ i $k\colon T\to S$. Aleshores $n\comp h=m$ i $m\comp k=n$, d'on $n\comp h\comp k=m\comp k=n$; com que $n$ és mono, $h\comp k=\id_T$. Simètricament, $k\comp h=\id_S$. Per tant $h$ és un isomorfisme amb $n\comp h=m$, és a dir, $m\sim n$ i $[m]=[n]$.
\end{prova}

\begin{observacio}[Categories ben potenciades]\index{categoria!ben potenciada}\index{ben potenciada}
La petitesa local garanteix que cada $\Hom(A,B)$ és un conjunt, però \emph{no} que la col·lecció $\Sub(A)$ ho sigui: en principi podria ser massa gran. Direm que una categoria és \textbf{ben potenciada} si, per a cada objecte $A$, els subobjectes de $A$ formen un conjunt. En el marc de classes que adoptem, la formulació precisa demana un conjunt de monomorfismes representants, perquè cada classe d'equivalència pot ser una classe pròpia (apèndix~\ref{cap:mida}). Aquesta és la hipòtesi que garanteix que $\Sub(A)$ és un conjunt de debò i que, més endavant, es pugui reunir en un prefeix
\[
\Sub\colon\cat{C}\op\longrightarrow\cat{Set}.
\]
Totes les categories que trobarem seran ben potenciades; a $\cat{Set}$, per exemple, $\Sub(A)\cong\mathcal{P}(A)$ és un conjunt. De fet, l'existència d'un classificador de subobjectes en una categoria localment petita ja implica aquesta condició.
\end{observacio}

\begin{observacio}[Predicats com a subobjectes]
Aquesta és la lectura lògica fonamental. Pensem un objecte $A$ com un \textbf{tipus} i un subobjecte $[m]\in\Sub(A)$ com un \textbf{predicat} sobre aquell tipus: la propietat \enquote{pertànyer a $S$}. L'ordre $[m]\leq[n]$ és la \textbf{implicació} entre predicats: tot element que satisfà $m$ satisfà $n$. Veurem que $\Sub(A)$ té sovint estructura de reticle ---amb intersecció i unió de subobjectes--- i fins i tot d'àlgebra de Heyting, cosa que en fa l'àlgebra de proposicions sobre $A$.
\end{observacio}

\section{Reindexació de subobjectes}

Al capítol de límits vam construir, per a un morfisme $f\colon A\to B$ en una categoria amb pullbacks, el functor de reindexació $f^{*}\colon\cat{C}/B\to\cat{C}/A$ (observació~\ref{obs:reindexacio}). La clau ara és que aquesta reindexació \emph{preserva els monomorfismes}, de manera que es restringeix a una operació entre ordres parcials de subobjectes.

\begin{proposicio}[El pullback preserva monomorfismes]
\label{prop:pullback-mono}
En un quadrat cartesià
\[
\begin{tikzpicture}[baseline=(m.center)]
  \matrix (m) [matrix of math nodes, column sep=3em, row sep=2.4em] {
    P & S \\
    A & B \\
  };
  \draw[-{Stealth[scale=1.1]}] (m-1-1) -- node[above] {$f'$} (m-1-2);
  \draw[-{Stealth[scale=1.1]}] (m-1-1) -- node[left] {$m'$} (m-2-1);
  \draw[-{Stealth[scale=1.1]}] (m-1-2) -- node[right] {$m$} (m-2-2);
  \draw[-{Stealth[scale=1.1]}] (m-2-1) -- node[below] {$f$} (m-2-2);
  \node at ([xshift=0.8em,yshift=-0.7em]m-1-1) {$\scriptstyle\lrcorner$};
\end{tikzpicture}
\qquad
m\comp f'=f\comp m',
\]
si $m$ és un monomorfisme, aleshores el seu pullback $m'\colon P\to A$ al llarg de $f$ també ho és.
\end{proposicio}

\begin{prova}
Siguin $u,v\colon Z\to P$ amb $m'\comp u=m'\comp v$. Per la commutativitat del quadrat,
\[
m\comp f'\comp u=f\comp m'\comp u=f\comp m'\comp v=m\comp f'\comp v,
\]
i, com que $m$ és mono, $f'\comp u=f'\comp v$. Així, $u$ i $v$ tenen les mateixes composicions amb les dues projeccions del pullback, $m'$ i $f'$, i per la unicitat de la propietat universal del pullback, $u=v$. Per tant $m'$ és un monomorfisme.
\end{prova}

\begin{observacio}
La conseqüència que ens interessa és immediata: com que la reindexació $f^{*}$ envia monomorfismes a monomorfismes, es restringeix a una operació ben definida entre els ordres parcials de subobjectes, que és la que ara definim. A $\cat{Set}$, aquesta operació és l'antiimatge $S\mapsto f^{-1}(S)$, tal com vam veure en introduir el pullback.
\end{observacio}

\begin{definicio}[Functor de reindexació sobre subobjectes]\index{reindexació!de subobjectes}\index{substitució}
\label{def:reindexacio-sub}
Sigui $f\colon A\to B$ un morfisme en una categoria $\cat{C}$ amb pullbacks. La \textbf{reindexació} indueix una aplicació
\[
f^{*}\colon\Sub(B)\longrightarrow\Sub(A)
\]
que envia un subobjecte $[m]$, amb $m\colon S\rightarrowtail B$, al subobjecte $[m']$ del seu pullback al llarg de $f$. Per la proposició~\ref{prop:pullback-mono}, $m'$ és mono, de manera que $f^{*}[m]$ és un subobjecte de $A$ ben definit.
\end{definicio}

\begin{proposicio}[La reindexació és monòtona]
\label{prop:reindexacio-monotona}
L'aplicació $f^{*}\colon\Sub(B)\to\Sub(A)$ preserva l'ordre: si $[m]\leq[n]$ a $\Sub(B)$, aleshores $f^{*}[m]\leq f^{*}[n]$ a $\Sub(A)$. És a dir, $f^{*}$ és una aplicació monòtona.
\end{proposicio}

\begin{prova}
Suposem $[m]\leq[n]$, amb $m\colon S\rightarrowtail B$, $n\colon T\rightarrowtail B$ i $h\colon S\to T$ tal que $n\comp h=m$. Siguin $(P,m',f_m)$ un pullback de $m$ al llarg de $f$ i $(Q,n',f_n)$ un de $n$, amb $f\comp m'=m\comp f_m$ i $f\comp n'=n\comp f_n$. La parella $m'\colon P\to A$, $h\comp f_m\colon P\to T$ és un con per al pullback de $n$, perquè
\[
f\comp m'=m\comp f_m=n\comp h\comp f_m.
\]
Per tant hi ha un únic $k\colon P\to Q$ amb $n'\comp k=m'$ i $f_n\comp k=h\comp f_m$. La primera igualtat diu que $[m']\leq[n']$. Observem que l'argument val per a representants i pullbacks \emph{qualssevol}.
\end{prova}

\begin{observacio}[$f^{*}$ està ben definida]\label{obs:reindexacio-ben-definida}
La definició~\ref{def:reindexacio-sub} tria un representant $m$ de la classe i un pullback. El resultat no depèn de cap de les dues tries. Si $m\sim m_1$, aleshores $[m]\leq[m_1]$ i $[m_1]\leq[m]$, i la demostració anterior, aplicada en tots dos sentits, dona que els pullbacks corresponents són equivalents. El cas de dos pullbacks del mateix $m$ és el cas particular $m_1=m$.
\end{observacio}

\begin{proposicio}[El functor de subobjectes]\label{prop:functor-sub}\index{Sub@$\Sub$!functor}
Si $\cat{C}$ és ben potenciada i té pullbacks, les reindexacions defineixen un functor
\[
\Sub\colon\cat{C}\op\longrightarrow\cat{Poset},
\]
que envia cada objecte $A$ a l'ordre parcial $\Sub(A)$ i cada morfisme $f\colon A\to B$ a l'aplicació monòtona $f^{*}\colon\Sub(B)\to\Sub(A)$. És a dir, $\id_A^{*}=\id_{\Sub(A)}$ i $(g\comp f)^{*}=f^{*}\comp g^{*}$, com a igualtats d'aplicacions entre ordres parcials.
\end{proposicio}

\begin{prova}
Com que $\cat{C}$ és ben potenciada, cada $\Sub(A)$ és un conjunt parcialment ordenat, i cada $f^{*}$ és una aplicació monòtona (proposició~\ref{prop:reindexacio-monotona}) ben definida (observació~\ref{obs:reindexacio-ben-definida}). Per aquesta mateixa observació, per calcular $f^{*}[m]$ podem fer servir \emph{qualsevol} pullback de $m$ al llarg de $f$. Queden les dues lleis functorials.

\emph{Identitats.} Sigui $m\colon S\rightarrowtail A$. El quadrat format per $m$ als dos costats verticals i per les identitats $\id_S$ i $\id_A$ als horitzontals és un pullback de $m$ al llarg de $\id_A$. En efecte, si $x\colon Z\to A$ i $y\colon Z\to S$ compleixen $\id_A\comp x=m\comp y$, l'única fletxa $z\colon Z\to S$ amb $m\comp z=x$ i $\id_S\comp z=y$ és $z=y$. Per tant $\id_A^{*}[m]=[m]$.

\emph{Composició.} Siguin $f\colon A\to B$, $g\colon B\to C$ i $m\colon S\rightarrowtail C$. Prenem un pullback $(Q,m_1,g')$ de $m$ al llarg de $g$, amb $g\comp m_1=m\comp g'$, de manera que $g^{*}[m]=[m_1]$. Prenem després un pullback $(P,m_2,f')$ de $m_1$ al llarg de $f$, amb $f\comp m_2=m_1\comp f'$, de manera que $f^{*}(g^{*}[m])=[m_2]$. Vegem que $(P,m_2,g'\comp f')$ és un pullback de $m$ al llarg de $g\comp f$; així $(g\comp f)^{*}[m]=[m_2]$ i les dues classes coincideixen. Primer, el quadrat commuta: $(g\comp f)\comp m_2=g\comp m_1\comp f'=m\comp(g'\comp f')$. Siguin ara $x\colon Z\to A$ i $y\colon Z\to S$ amb $g\comp f\comp x=m\comp y$.
\begin{itemize}
\item \emph{Existència.} La parella $(f\comp x,\,y)$ és un con per al primer pullback, de manera que hi ha un únic $q\colon Z\to Q$ amb $m_1\comp q=f\comp x$ i $g'\comp q=y$. Aleshores $(x,q)$ és un con per al segon pullback, i hi ha un únic $p\colon Z\to P$ amb $m_2\comp p=x$ i $f'\comp p=q$. Per tant, $m_2\comp p=x$ i $(g'\comp f')\comp p=y$.
\item \emph{Unicitat.} Si $p'\colon Z\to P$ també compleix $m_2\comp p'=x$ i $g'\comp f'\comp p'=y$, aleshores $f'\comp p'$ compleix $m_1\comp(f'\comp p')=f\comp m_2\comp p'=f\comp x$ i $g'\comp(f'\comp p')=y$. Per la unicitat al primer pullback, $f'\comp p'=q$. Per la unicitat al segon, $p'=p$.
\end{itemize}
Aquest argument és el lema d'enganxament de pullbacks, en la direcció que aquí cal (exercici resolt corresponent de la Pràctica del capítol~\ref{cap:limits}). Per tant $(g\comp f)^{*}=f^{*}\comp g^{*}$.

Les dues lleis són \emph{igualtats} d'aplicacions entre ordres parcials, i no només isomorfismes. A nivell de monomorfismes, el pullback al llarg de $g\comp f$ i el pullback iterat només coincideixen llevat d'isomorfisme. En passar a classes, però, aquest isomorfisme queda absorbit.
\end{prova}

\begin{observacio}
Aquí no cal cap noció de pseudofunctor. La coherència \enquote{llevat d'isomorfisme} que vam trobar en reindexar les categories tall $\cat{C}/B$ (observació~\ref{obs:reindexacio}) desapareix en passar a classes de subobjectes. Aquest functor és precisament el que el capítol següent voldrà representar.
\end{observacio}

\begin{observacio}[La substitució]
Aquesta és la interpretació que anunciàvem. Si pensem un subobjecte $[m]\in\Sub(B)$ com un predicat $P(y)$ sobre la variable $y:B$, i $f\colon A\to B$ com un \enquote{terme} $f(x)$ amb $x:A$, aleshores $f^{*}[m]\in\Sub(A)$ és el predicat $P(f(x))$: la \textbf{substitució} de $y$ per $f(x)$ dins de $P$. A $\cat{Set}$, on $[m]$ és un subconjunt $S\subseteq B$, la reindexació és l'antiimatge:
\[
f^{*}(S)=f^{-1}(S)=\{x\in A\mid f(x)\in S\}.
\]
La monotonia de $f^{*}$ diu que la substitució respecta la implicació: si $P\Rightarrow Q$, aleshores $P(f)\Rightarrow Q(f)$.
\end{observacio}

\begin{exemple}[Subgrafs]\label{ex:subgrafs}\index{graf dirigit!subobjectes}
A $\cat{Graph}$, els monomorfismes són els morfismes de grafs injectius sobre vèrtexs i sobre arestes (vegeu els Exercicis). Un subobjecte d'un graf $G$ es pot representar, doncs, per un \textbf{subgraf}: un subconjunt de vèrtexs $V'\subseteq V_G$ i un subconjunt d'arestes $E'\subseteq E_G$ tancat per origen i destí, és a dir, tal que l'origen i el destí de tota aresta de $E'$ són a $V'$. La reindexació al llarg d'un morfisme de grafs $F\colon G'\to G$ es calcula component a component, per antiimatge de vèrtexs i d'arestes:
\[
F^{*}(V',E')=\bigl(F_V^{-1}(V'),\ F_E^{-1}(E')\bigr),
\]
que torna a ser un subgraf de $G'$, perquè $F$ preserva origen i destí. L'identificació \enquote{predicat $=$ subobjecte} no és, doncs, exclusiva dels conjunts: un predicat sobre un graf és un subgraf, i la substitució és l'antiimatge component a component.
\end{exemple}

\section{Intersecció de subobjectes}

L'ordre parcial $\Sub(A)$ té més estructura que l'ordre: quan la categoria té pullbacks, admet \textbf{ínfims binaris}, que interpreten la conjunció de predicats.

\begin{proposicio}[La intersecció és un pullback]\index{subobjecte!intersecció}
\label{prop:interseccio}
Siguin $m\colon S\rightarrowtail A$ i $n\colon T\rightarrowtail A$ dos subobjectes de $A$. El seu \textbf{pullback} sobre $A$,
\[
\begin{tikzpicture}[baseline=(m.center)]
  \matrix (m) [matrix of math nodes, column sep=3em, row sep=2.4em] {
    S\cap T & T \\
    S & A \\
  };
  \draw[-{Stealth[scale=1.1]}] (m-1-1) -- node[above] {$q$} (m-1-2);
  \draw[-{Stealth[scale=1.1]}] (m-1-1) -- node[left] {$p$} (m-2-1);
  \draw[-{Stealth[scale=1.1]}] (m-1-2) -- node[right] {$n$} (m-2-2);
  \draw[-{Stealth[scale=1.1]}] (m-2-1) -- node[below] {$m$} (m-2-2);
  \node at ([xshift=1em,yshift=-0.7em]m-1-1) {$\scriptstyle\lrcorner$};
\end{tikzpicture}
\qquad
m\comp p=n\comp q,
\]
amb el morfisme induït $m\comp p=n\comp q\colon S\cap T\to A$, que és un monomorfisme, representa l'\textbf{ínfim} $[m]\wedge[n]$ a $\Sub(A)$. Escrivim també $S\cap T=S\times_A T$.
\end{proposicio}

\begin{prova}
El morfisme $m\comp p\colon S\cap T\to A$ és mono per ser composició de monos (la projecció $p$ del pullback és mono per la proposició~\ref{prop:pullback-mono}, i $m$ és mono). Per construcció $[m\wedge n]\leq[m]$ i $[m\wedge n]\leq[n]$, ja que hi ha les projeccions cap a $S$ i $T$. I si $[k]\leq[m]$ i $[k]\leq[n]$ per a un tercer subobjecte $k\colon R\rightarrowtail A$, les factoritzacions $R\to S$ i $R\to T$ són compatibles sobre $A$, de manera que la propietat universal del pullback dona una única factorització $R\to S\cap T$; per tant $[k]\leq[m]\wedge[n]$. Això caracteritza l'ínfim.
\end{prova}

\begin{observacio}
En termes lògics, $[m]\wedge[n]$ és la \textbf{conjunció} dels predicats: l'element pertany a $S\cap T$ si i només si pertany a $S$ \emph{i} a $T$. A $\cat{Set}$ recuperem la intersecció de subconjunts. En qualsevol categoria, el subobjecte \textbf{màxim} de $\Sub(A)$ és $[\id_A]$ (el predicat sempre cert), com es comprova a la Pràctica; i quan hi ha objecte inicial estricte, el subobjecte \textbf{mínim} és la injecció de l'objecte inicial (el predicat sempre fals). Amb intersecció i aquests extrems, $\Sub(A)$ és un \emph{semireticle amb màxim}. Les \textbf{unions} ---la disjunció--- exigeixen estructura addicional: la regularitat proporciona imatges, existencials i interseccions, però no, en general, suprems finits estables a les fibres $\Sub(A)$. Aquests apareixen, per exemple, en les \textbf{categories coherents}, que es tractaran al volum de lògica categòrica.
\end{observacio}

\section{Epimorfismes regulars i imatges}

Per definir la imatge d'un morfisme, ens cal la classe adequada d'epimorfismes: els que sorgeixen com a quocients.

\begin{definicio}[Epimorfisme regular]\index{epimorfisme!regular}
\label{def:epi-regular}
Un morfisme $e\colon A\to B$ és un \textbf{epimorfisme regular} si és el coequalitzador d'alguna parella de morfismes $u,v\colon R\to A$.
\end{definicio}

\begin{definicio}[Parella nucli]\index{parella nucli}
\label{def:parella-nucli}
La \textbf{parella nucli} d'un morfisme $f\colon A\to B$ és el pullback de $f$ amb ell mateix,
\[
\begin{tikzpicture}[baseline=(q.center)]
  \matrix (q) [matrix of math nodes, column sep=3em, row sep=2.4em] {
    R & A \\
    A & B \\
  };
  \draw[-{Stealth[scale=1.1]}] (q-1-1) -- node[above] {$p_2$} (q-1-2);
  \draw[-{Stealth[scale=1.1]}] (q-1-1) -- node[left] {$p_1$} (q-2-1);
  \draw[-{Stealth[scale=1.1]}] (q-1-2) -- node[right] {$f$} (q-2-2);
  \draw[-{Stealth[scale=1.1]}] (q-2-1) -- node[below] {$f$} (q-2-2);
  \node at ([xshift=0.8em,yshift=-0.7em]q-1-1) {$\scriptstyle\lrcorner$};
\end{tikzpicture}
\qquad
f\comp p_1=f\comp p_2;
\]
els seus dos morfismes $p_1,p_2\colon R\to A$ recullen les parelles d'elements que $f$ identifica. A $\cat{Set}$, $R=\{(a,a')\mid f(a)=f(a')\}$.
\end{definicio}

\begin{lema}\label{lema:epi-regular-nucli}\index{epimorfisme!regular}\index{parella nucli}
En una categoria amb pullbacks, si $e\colon A\to B$ és un epimorfisme regular, aleshores $e$ és el coequalitzador de la seva parella nucli.
\end{lema}

\begin{prova}
Suposem que $e$ és el coequalitzador de $u,v\colon K\rightrightarrows A$, i sigui $p_1,p_2\colon R\rightrightarrows A$ la parella nucli de $e$. Com que $e\comp u=e\comp v$, la propietat universal del pullback dona $k\colon K\to R$ amb $p_1\comp k=u$ i $p_2\comp k=v$. Si $h\colon A\to Z$ satisfà $h\comp p_1=h\comp p_2$, aleshores
\[
h\comp u=h\comp p_1\comp k=h\comp p_2\comp k=h\comp v,
\]
i com que $e$ és el coequalitzador de $u$ i $v$, hi ha una única $\bar h\colon B\to Z$ amb $\bar h\comp e=h$. Com que, a més, $e\comp p_1=e\comp p_2$, això és exactament la propietat universal del coequalitzador de $p_1$ i $p_2$.
\end{prova}

\begin{definicio}[Factorització imatge]\index{imatge}\index{factorització!imatge}
\label{def:imatge}
Una \textbf{factorització imatge} de $f\colon A\to B$ és una descomposició $f=m\comp e$ amb $A\xrightarrow{\,e\,}I\xrightarrow{\,m\,}B$, on $e$ és un epimorfisme regular i $m$ un monomorfisme. El subobjecte $[m]\in\Sub(B)$ és la \textbf{imatge} de $f$, denotada $\Imatge(f)$.
\end{definicio}

\begin{proposicio}[La imatge és el mínim subobjecte]
\label{prop:imatge-unica}
Quan existeix, la factorització imatge és única llevat d'isomorfisme canònic, i $\Imatge(f)$ és el \emph{menor} subobjecte de $B$ a través del qual $f$ es factoritza: si $f$ es factoritza per un mono $n\colon T\rightarrowtail B$, aleshores $\Imatge(f)\leq[n]$.
\end{proposicio}

\begin{prova}
Sigui $f=m\comp e$ una factorització imatge. Com que $m$ és mono, qualsevol parella nucli de $e$ és també una parella nucli de $f$ (perquè $m\comp e\comp u=m\comp e\comp v$ si i només si $e\comp u=e\comp v$); en particular, les parelles nucli escollides per a $e$ i per a $f$ són canònicament isomorfes. Pel lema~\ref{lema:epi-regular-nucli}, $e$ és el coequalitzador de la seva parella nucli $p_1,p_2$. Si $f=m'\comp e'$ és una altra factorització amb $m'$ mono, aleshores $e'$ iguala les branques d'aquesta parella nucli, perquè $m'\comp e'\comp p_1=f\comp p_1=f\comp p_2=m'\comp e'\comp p_2$ i $m'$ és mono, de manera que $e'$ es factoritza per $e$ mitjançant un únic morfisme; simètricament s'obté l'invers, i les dues composicions són identitats perquè $e$ és epi. Això dona l'isomorfisme canònic. La minimalitat en segueix: si $f=n\comp g$ amb $n$ mono, aleshores $g$ iguala la parella nucli i es factoritza per $e$, d'on $\Imatge(f)=[m]\leq[n]$.
\end{prova}

\begin{observacio}
A $\cat{Set}$, $\Imatge(f)$ és el subconjunt imatge $f(A)\subseteq B$, amb $e$ l'aplicació exhaustiva $A\to f(A)$ i $m$ la injecció. Recuperem la noció usual d'imatge.
\end{observacio}

\section{Categories regulars i estabilitat}

Perquè la teoria d'imatges sigui utilitzable, cal que les imatges es comportin bé sota el canvi de base.

\begin{definicio}[Categoria regular]\index{categoria!regular}
\label{def:regular}
Una categoria $\cat{C}$ és \textbf{regular} si té tots els límits finits, tota fletxa admet una factorització imatge (epi regular seguit de mono), i els epimorfismes regulars són \textbf{estables per pullback}.
\end{definicio}

\begin{exemple}
Són regulars $\cat{Set}$, tota categoria d'àlgebres d'una teoria algebraica (grups, anells, mòduls\ldots), tota categoria abeliana i tot topos elemental. La categoria $\cat{Top}$ \emph{no} és regular, perquè les imatges no hi són estables per pullback.
\end{exemple}

\begin{proposicio}[Estabilitat de la imatge]
\label{prop:imatge-estable}
En una categoria regular, la formació d'imatges commuta amb el pullback: si es forma el pullback de $f\colon A\to B$ al llarg de $g\colon B'\to B$, aleshores $\Imatge(f')=g^{*}\bigl(\Imatge(f)\bigr)$, on $f'$ és el morfisme reindexat.
\end{proposicio}

\begin{prova}
Factoritzem $f=m\comp e$. En prendre el pullback al llarg de $g$, l'epimorfisme regular $e$ dona un epi regular $e'$ (per l'estabilitat) i el mono $m$ dona un mono $m'$ (proposició~\ref{prop:pullback-mono}). Així $f'=m'\comp e'$ és una factorització epi regular\,--\,mono, i per la unicitat de la imatge, $m'$ representa $g^{*}(\Imatge(f))$.
\end{prova}

\section{L'existencial com a adjunt de la reindexació}

De tot això en resulta que la reindexació té un adjunt per l'esquerra, que reprèn el fil del capítol d'adjuncions.

\begin{teorema}[$\exists_f\dashv f^{*}$]\index{quantificador!existencial}\index{adjunció!existencial}
\label{teo:existencial-adjunt}
Sigui $\cat{C}$ una categoria regular i $f\colon A\to B$. L'operació
\[
\exists_f[s]:=\Imatge(f\comp s)\in\Sub(B)
\]
és adjunta per l'esquerra de la reindexació $f^{*}\colon\Sub(B)\to\Sub(A)$; és a dir,
\[
\exists_f[s]\leq[t]\quad\Longleftrightarrow\quad[s]\leq f^{*}[t].
\]
\end{teorema}

\begin{observacio}
L'operació $\exists_f$ està definida sobre classes de subobjectes. Si $s\sim s'$, és a dir, $s'=s\comp\varphi$ amb $\varphi$ isomorfisme, aleshores $f\comp s'=(f\comp s)\comp\varphi$, i les factoritzacions imatge de $f\comp s$ i de $f\comp s'$ donen el mateix subobjecte de $B$. És exactament la unicitat de la imatge llevat d'isomorfisme (proposició~\ref{prop:imatge-unica}) el que fa que $\exists_f$ estigui ben definida.
\end{observacio}

\begin{prova}
Si $[s]\leq f^{*}[t]$, aleshores $S$ es factoritza a través del pullback $f^{*}T$, i component-lo amb la projecció cap a $T$ s'obté que $f\comp s$ es factoritza per $t$; com que $\Imatge(f\comp s)$ és el mínim subobjecte amb aquesta propietat, $\exists_f[s]\leq[t]$. Recíprocament, si $\exists_f[s]\leq[t]$, aleshores $f\comp s=t\comp k$ per a algun $k$, i aquesta igualtat exhibeix $S$ com un con sobre el diagrama del pullback $f^{*}T$, de manera que $[s]\leq f^{*}[t]$.
\end{prova}

\begin{observacio}[L'universal com a estructura addicional]\index{quantificador!universal}\index{adjunció!quantificadors}
En una categoria regular, la reindexació $f^{*}$ té sempre l'adjunt per l'esquerra $\exists_f$ que acabem de construir, però \emph{no} té necessàriament un adjunt per la dreta. Quan cada $f^{*}$ admet a més un adjunt per la dreta, el denotem $\forall_f$ i queda caracteritzat per
\[
f^{*}[t]\leq[s]\quad\Longleftrightarrow\quad[t]\leq\forall_f[s].
\]
Aquesta és estructura addicional, que la regularitat no garanteix: la regularitat només dona el costat existencial. Apareix, per exemple, en les categories localment cartesianes tancades amb les hipòtesis de mida adequades i, en particular, en tot topos elemental. La construcció de $\forall_f$ i el seu paper lògic es tracten al volum de lògica categòrica; no s'obté simplement dualitzant la construcció de la imatge dins de la mateixa fibra de subobjectes.
\end{observacio}

\begin{observacio}[Compatibilitat amb la substitució]\index{Beck--Chevalley}
En una categoria regular, l'estabilitat de les imatges per pullback dona una compatibilitat entre l'existencial i la substitució. Per a un quadrat cartesià
\[
\begin{tikzpicture}[baseline=(q.center)]
  \matrix (q) [matrix of math nodes, column sep=3em, row sep=2.4em] {
    A' & A \\
    B' & B \\
  };
  \draw[-{Stealth[scale=1.1]}] (q-1-1) -- node[above] {$h$} (q-1-2);
  \draw[-{Stealth[scale=1.1]}] (q-1-1) -- node[left] {$f'$} (q-2-1);
  \draw[-{Stealth[scale=1.1]}] (q-1-2) -- node[right] {$f$} (q-2-2);
  \draw[-{Stealth[scale=1.1]}] (q-2-1) -- node[below] {$g$} (q-2-2);
  \node at ([xshift=0.8em,yshift=-0.7em]q-1-1) {$\scriptstyle\lrcorner$};
\end{tikzpicture}
\]
es té la \textbf{condició de Beck--Chevalley}
\[
g^{*}\comp\exists_f=\exists_{f'}\comp h^{*}\colon\Sub(A)\to\Sub(B'):
\]
quantificar existencialment i després substituir és el mateix que substituir i després quantificar. No ho demostrem aquí; és un primer exemple del tipus de coherència que la lògica categòrica exigeix.
\end{observacio}

\begin{observacio}[Lectura lògica]
Aquesta adjunció té una lectura lògica il·lustrativa. Si $f$ és la projecció $\pi\colon X\times Y\to X$ que oblida la variable $Y$, i pensem un subobjecte de $X\times Y$ com un predicat $s(x,y)$, aleshores
\[
\exists_\pi[s]=\{x\mid\exists y,\ s(x,y)\}.
\]
L'existencial apareix, doncs, com a adjunt per l'esquerra de la substitució. Quan la categoria té també els adjunts per la dreta, l'universal es llegeix anàlogament com $\forall_\pi[s]=\{x\mid\forall y,\ s(x,y)\}$. El desenvolupament d'aquesta idea com a \emph{sistema lògic} ---amb el seu càlcul, la interpretació de l'existencial per la factorització imatge i les condicions de coherència pertinents--- és objecte de l'estudi de la lògica categòrica, i queda fora de l'abast d'aquest volum.
\end{observacio}

\begin{resumcap}
\apartat{Idees centrals}
\begin{enumerate}
  \item Un subobjecte és una classe d'equivalència de monomorfismes amb el mateix codomini; $\Sub(A)$ és un ordre parcial, l'àlgebra de predicats sobre $A$. A $\cat{Set}$ són els subconjunts, i a $\cat{Graph}$, els subgrafs.
  \item La reindexació $f^{*}$, construïda per pullback, interpreta la substitució; en una categoria ben potenciada amb pullbacks dona un functor $\Sub\colon\cat{C}\op\to\cat{Poset}$, amb igualtats estrictes.
  \item La intersecció de subobjectes és el seu pullback i interpreta la conjunció; el màxim $[\id_A]$ interpreta el predicat cert.
  \item Un epimorfisme regular és el coequalitzador de la seva parella nucli; la factorització imatge (epi regular seguit de mono) és única llevat d'isomorfisme, i en una categoria regular existeix i és estable per pullback.
  \item En una categoria regular, $f^{*}$ té un adjunt per l'esquerra $\exists_f$, que interpreta la quantificació existencial; l'universal $\forall_f$ és estructura addicional que la regularitat no garanteix.
\end{enumerate}
\apartat{Errors típics} Creure que la regularitat dona unions de subobjectes (calen categories coherents) o l'adjunt dret $\forall_f$; confondre els epimorfismes amb els epimorfismes regulars (a $\cat{Top}$, els primers són les aplicacions contínues exhaustives i els segons, les aplicacions quocient); confondre la reindexació $f^{*}$ amb la imatge $\exists_f$; i usar \enquote{imatge}, \enquote{buit} o \enquote{elements} en una categoria arbitrària sense declarar les hipòtesis.
\apartat{Lectura recomanada} \cite[cap.~5 i~\S9.5]{awodey} per a subobjectes, imatges i l'adjunció existencial; \cite[§§2.8, 2.10, 9.4 i 9.8]{barrwells} per a subobjectes, epimorfismes regulars, categories regulars i sistemes de factorització; \cite[V.7]{maclane} per al tractament clàssic dels subobjectes; \cite[A1.3]{johnstone} i \cite{borceux} per a categories regulars amb més profunditat; \cite{jacobs} per a la lectura fibracional de la reindexació.
\end{resumcap}

\chapter[Classificadors de subobjectes i topos]{\texorpdfstring{Classificadors de subobjectes\\ i introducció als topos}{Classificadors de subobjectes i introducció als topos}}
\label{cap:topos}

\begin{objectius}
\apartat{Prerequisits} Subobjectes i el functor $\Sub$ (cap.~\ref{cap:subobjectes}); representabilitat i Yoneda (cap.~\ref{cap:yoneda}); tancament cartesià (cap.~\ref{cap:cartesianes}).
\apartat{Objectius} En acabar el capítol, el lector hauria de poder:
\begin{itemize}
  \item definir el classificador de subobjectes $\Omega$ i el morfisme característic d'un subobjecte;
  \item expressar la pertinença d'un element generalitzat com una igualtat de morfismes cap a $\Omega$;
  \item demostrar que tenir classificador equival a la representabilitat del functor $\Sub$;
  \item enunciar la definició de topos elemental i reconèixer $\Omega^{E}$ com a objecte de parts;
  \item calcular morfismes característics a $\cat{Set}$ i a $\cat{Graph}$, i explicar per què la representabilitat de $\Sub$ converteix els predicats en morfismes i la substitució en composició;
  \item reconèixer com a topos $\cat{Set}$, $\cat{FinSet}$, les categories de prefeixos sobre una categoria petita, amb el classificador de sedassos, i $\cat{Graph}$, i distingir els topos elementals de les categories de prefeixos;
  \item explicar en quin sentit $\Omega$ és un objecte de valors de veritat, i per què això no fa booleà tot topos.
\end{itemize}
\end{objectius}

Al capítol anterior hem vist que els subobjectes d'un objecte $A$ formen un ordre parcial $\Sub(A)$, que podem llegir com l'ordre de predicats sobre $A$, i que la reindexació el fa functorial: en una categoria ben potenciada amb pullbacks obtenim un functor $\Sub\colon\cat{C}\op\to\cat{Poset}$ (proposició~\ref{prop:functor-sub}). Sorgeix una pregunta natural: hi ha un objecte que \emph{representi} aquest functor? És a dir, un objecte $\Omega$ tal que donar un subobjecte de $A$ equivalgui a donar un morfisme $A\to\Omega$? Quan aquest objecte existeix, s'anomena \textbf{classificador de subobjectes}, i les categories que en tenen ---juntament amb límits finits i exponencials--- són els \textbf{topos elementals}. Aquest capítol n'introdueix la definició i els exemples bàsics. No en desenvolupem la lògica interna: aquesta pertany a l'estudi de la lògica categòrica, per al qual aquest volum és la preparació.

\section{El classificador de subobjectes}

A $\cat{Set}$, un subconjunt $U\subseteq E$ es descriu per la seva \textbf{funció característica} $\chi_U\colon E\to\{0,1\}$, que val $1$ als elements de $U$ i $0$ a la resta. El conjunt $\{0,1\}$ actua com un objecte de \enquote{valors de veritat}, i triar un subconjunt equival a triar una funció cap a ell. La noció següent generalitza aquesta idea a una categoria arbitrària.

\begin{definicio}[Classificador de subobjectes]\index{classificador de subobjectes}\index{Omega@$\Omega$}
\label{def:classificador}
Sigui $\cat{E}$ una categoria amb tots els límits finits. Un \textbf{classificador de subobjectes} és un objecte $\Omega$ juntament amb un morfisme $\mathsf{true}\colon 1\to\Omega$ des de l'objecte terminal, amb la propietat universal següent: per a cada objecte $E$ i cada subobjecte $m\colon U\rightarrowtail E$, existeix un \emph{únic} morfisme $\chi_m\colon E\to\Omega$, el \textbf{morfisme característic} del subobjecte, que fa del quadrat
\[
\begin{tikzpicture}[baseline=(m.center)]
  \matrix (m) [matrix of math nodes, column sep=3em, row sep=2.4em] {
    U & 1 \\
    E & \Omega \\
  };
  \draw[-{Stealth[scale=1.1]}] (m-1-1) -- node[above] {$!$} (m-1-2);
  \draw[-{Stealth[scale=1.1]}] (m-1-1) -- node[left] {$m$} (m-2-1);
  \draw[-{Stealth[scale=1.1]}] (m-1-2) -- node[right] {$\mathsf{true}$} (m-2-2);
  \draw[-{Stealth[scale=1.1]}] (m-2-1) -- node[below] {$\chi_m$} (m-2-2);
  \node at ([xshift=0.8em,yshift=-0.7em]m-1-1) {$\scriptstyle\lrcorner$};
\end{tikzpicture}
\]
un pullback. És a dir, $U$ es recupera com el pullback de $\mathsf{true}$ al llarg de $\chi_m$. A $\cat{Set}$, el morfisme característic és la funció característica usual.
\end{definicio}

La condició diu que cada subobjecte de $E$ té un i només un morfisme característic, i que el subobjecte es reconstrueix com l'antiimatge del \enquote{punt vertader} $\mathsf{true}$. A $\cat{Set}$, $\Omega=\{0,1\}$ i $\mathsf{true}\colon 1\to\{0,1\}$ tria l'element $1$; la funció característica és la de sempre, i el pullback de $\mathsf{true}$ al llarg de $\chi_U$ recupera exactament $U$.

\begin{exemple}[Un morfisme característic calculat a $\cat{Set}$]\label{ex:caracteristic-set}\index{morfisme característic!a Set@a $\cat{Set}$}
Sigui $E=\{1,2,3,4\}$ i $U=\{2,4\}$, el predicat \enquote{$x$ és parell}, amb la injecció $m\colon U\hookrightarrow E$. El morfisme característic és
\[
\chi_U\colon E\to\{0,1\},\qquad 1\mapsto0,\quad 2\mapsto1,\quad 3\mapsto0,\quad 4\mapsto1.
\]
Comprovem les tres afirmacions de la definició sobre aquest cas.
\begin{itemize}
\item \emph{El quadrat commuta.} $\chi_U\comp m$ envia $2$ i $4$ a $1$, que és $\mathsf{true}\comp\,!_U$.
\item \emph{És un pullback.} El pullback de $\mathsf{true}$ al llarg de $\chi_U$ és $\{x\in E\mid\chi_U(x)=1\}=\{2,4\}$, amb la inclusió (el pullback a $\cat{Set}$, capítol~\ref{cap:limits}): recupera exactament $U$.
\item \emph{Unicitat.} Una altra aplicació $\chi\colon E\to\{0,1\}$ amb el mateix pullback ha de valer $1$ exactament a $\{2,4\}$, i per tant és $\chi_U$. Per exemple, l'aplicació que val $1$ a $\{2,3,4\}$ fa commutar el quadrat, però el seu pullback és $\{2,3,4\}$, no $U$: la commutativitat sola no basta.
\end{itemize}
La pertinença per elements, que formalitzarem tot seguit, hi és visible: l'element $x\colon\{\ast\}\to E$ amb $x(\ast)=4$ compleix $\chi_U\comp x=\mathsf{true}$, i el que tria $3$, no. La substitució també: si $g\colon\{a,b,c\}\to E$ és $g(a)=1$, $g(b)=2$, $g(c)=4$, aleshores $\chi_U\comp g$ envia $a\mapsto0$, $b\mapsto1$ i $c\mapsto1$, i classifica $g^{-1}(U)=\{b,c\}$. \emph{Substituir} en un predicat és \emph{precompondre} amb el seu morfisme característic.
\end{exemple}

\begin{observacio}[Pertinença per elements generalitzats]\index{element generalitzat!pertinença}
Sigui $m\colon U\rightarrowtail E$ un subobjecte amb morfisme característic $\chi_m\colon E\to\Omega$, i sigui $x\colon X\to E$ un element generalitzat. Aleshores
\[
x\ \text{factoritza a través de } m
\quad\Longleftrightarrow\quad
\chi_m\comp x=\mathsf{true}\comp\,!_X.
\]
En efecte, si $x=m\comp y$, aleshores $\chi_m\comp x=\chi_m\comp m\comp y=\mathsf{true}\comp\,!_U\comp y=\mathsf{true}\comp\,!_X$; i recíprocament, si $\chi_m\comp x=\mathsf{true}\comp\,!_X$, la propietat universal del pullback dona una factorització (única) de $x$ a través de $m$. És la lectura precisa, sense elements ordinaris, de
\[
x\in U\quad\Longleftrightarrow\quad\chi_U(x)=\mathsf{true}:
\]
el classificador converteix la pertinença a un subobjecte en una igualtat de morfismes cap a $\Omega$.
\end{observacio}

\begin{proposicio}[Classificador i representabilitat]
\label{prop:omega-representa}
Suposem $\cat{E}$ localment petita, ben potenciada i amb límits finits, de manera que el functor de subobjectes pren valors a $\cat{Set}$. Aleshores les dades d'un classificador de subobjectes $(\Omega,\mathsf{true})$ són equivalents a una representació del functor de subobjectes:
\[
\Sub_{\cat{E}}(E)\;\cong\;\Hom_{\cat{E}}(E,\Omega),
\]
natural en $E$. En conseqüència, el classificador, si existeix, és únic llevat d'un únic isomorfisme compatible amb $\mathsf{true}$.
\end{proposicio}

\begin{prova}
\emph{Del classificador a la representació.} Enviem un subobjecte $[m]$ al seu morfisme característic $\chi_m$. L'aplicació és injectiva perquè $\chi_m$ determina $[m]$ com el pullback de $\mathsf{true}$ al llarg de $\chi_m$, i és exhaustiva perquè, per a tot $\chi\colon E\to\Omega$, el pullback de $\mathsf{true}$ al llarg de $\chi$ és un subobjecte el morfisme característic del qual és $\chi$, per la unicitat de la definició~\ref{def:classificador}. La naturalitat en $E$ és la compatibilitat amb la reindexació: reindexar un subobjecte al llarg de $g\colon E'\to E$ correspon a precompondre la característica amb $g$, pel lema d'enganxament de pullbacks.

\emph{De la representació al classificador.} Suposem que $\Sub_{\cat{E}}\cong\Hom_{\cat{E}}(-,\Omega)$ naturalment, i sigui $t\colon T\rightarrowtail\Omega$ un representant de l'element universal, és a dir, del subobjecte de $\Omega$ que correspon a $\id_\Omega$. Pel lema de Yoneda, el subobjecte que correspon a un morfisme $\chi\colon E\to\Omega$ és la reindexació $\chi^{*}[t]$: tot subobjecte és el pullback de $t$ al llarg d'un únic morfisme. Vegem que $T$ és terminal. Per a cada objecte $A$, el subobjecte màxim $[\id_A]$ és el pullback de $t$ al llarg d'un únic $\chi\colon A\to\Omega$, i el pullback dona un morfisme $A\to T$. Si $g\colon A\to T$ és un morfisme qualsevol, el pullback de $t$ al llarg de $t\comp g$ és $[\id_A]$, perquè $t$ és mono; per unicitat, $t\comp g=\chi$, i com que $t$ és mono, $g$ queda determinat. Així hi ha exactament un morfisme $A\to T$, i $T\cong 1$. Identificant $T$ amb $1$, el mono $t$ és un morfisme $\mathsf{true}\colon 1\to\Omega$ que satisfà la definició~\ref{def:classificador}.

La unicitat del classificador és el principi de Yoneda aplicat a dos objectes que representen el mateix functor, amb l'element universal $\mathsf{true}$ com a dada universal.
\end{prova}

\begin{observacio}[Per què la representabilitat és un pas conceptual]\label{obs:sub-representable}
L'equivalència de la proposició no és només una reformulació. El functor $\Sub$ és una dada \emph{dispersa}: un ordre parcial per a cada objecte, i una reindexació per a cada morfisme. La representabilitat el concentra en un sol objecte $\Omega$ i en un sol subobjecte universal, $\mathsf{true}\colon 1\rightarrowtail\Omega$. Aquest canvi té tres conseqüències:
\begin{itemize}
\item \emph{Els predicats esdevenen morfismes.} Un predicat sobre $E$ ja no és una classe d'equivalència de monomorfismes, sinó una fletxa $E\to\Omega$, que es pot compondre, aparellar i currificar com qualsevol altra.
\item \emph{La substitució esdevé composició.} La naturalitat de la bijecció diu que reindexar al llarg de $g$ és precompondre amb $g$, tal com hem vist a l'exemple~\ref{ex:caracteristic-set}.
\item \emph{Tot subobjecte és un cas d'un de sol.} Cada subobjecte és el pullback de $\mathsf{true}$ al llarg del seu morfisme característic. Les operacions lògiques es poden definir, doncs, una sola vegada sobre $\Omega$, com a morfismes $\Omega\times\Omega\to\Omega$, i transportar-se a tots els objectes (exercici resolt~\ref{pr-conjuncio-omega} de la Pràctica).
\end{itemize}
És el mateix pas que va de $\Hom(A\times B,C)$ a l'exponencial $C^B$: una família de conjunts esdevé un objecte de la categoria. Per això el classificador és el que permet que la lògica \enquote{visqui dins} de la categoria. Pel lema de Yoneda, a més, el classificador es pot calcular sondejant-lo amb objectes representables, com farem a $\cat{Graph}$.
\end{observacio}

\section{Topos elementals}

\begin{definicio}[Topos elemental]\index{topos}\index{topos!elemental}
\label{def:topos}
Un \textbf{topos elemental} és una categoria $\cat{E}$ que compleix:
\begin{enumerate}
\item $\cat{E}$ té tots els límits finits;
\item $\cat{E}$ té un classificador de subobjectes $\Omega$;
\item $\cat{E}$ és cartesiana tancada (té tots els exponencials).
\end{enumerate}
\end{definicio}

Aquesta definició, compacta, resulta ser extraordinàriament rica. Es pot demostrar, per exemple, que tot topos té també tots els colímits finits, i que, si $\cat{E}$ és un topos, també ho és cada categoria $\cat{E}/A$ de $\cat{E}$ sobre un objecte $A$. No demostrem aquí aquestes propietats; les recollim per il·lustrar que la combinació \enquote{límits finits $+$ classificador $+$ exponencials} és molt més potent del que sembla a primera vista.

\begin{observacio}[Topos elemental i categoria de prefeixos]\label{obs:elemental-prefeixos}\index{topos!elemental i de prefeixos}
Convé no confondre la definició amb el seu exemple principal. Un topos \emph{elemental} es defineix per tres condicions expressables en el llenguatge de les categories, i només demana límits \emph{finits}. Les categories de prefeixos $\cat{Set}^{\cat{C}\op}$ en són una classe important, però tenen molta més estructura: tots els límits i colímits \emph{petits}, i un classificador que es pot escriure explícitament amb sedassos. No tot topos elemental és d'aquesta mena. La categoria $\cat{FinSet}$ dels conjunts finits és un topos: té límits finits, l'exponencial de dos conjunts finits és finit, i $\{0,1\}$ hi és un classificador. Però no té coproductes infinits: no hi ha cap coproducte d'una família numerable de còpies de $1$, perquè un tal objecte hauria d'admetre una aplicació cap a $\{0,1\}$ per a cada successió de zeros i uns, i un conjunt finit només n'admet un nombre finit. En canvi, tota categoria de prefeixos té tots els colímits petits, i una equivalència de categories els conserva. Per tant, $\cat{FinSet}$ no és equivalent a cap categoria de prefeixos. Entre totes dues nocions hi ha els \emph{topos de Grothendieck} (categories de feixos), que aquí no tractem.
\end{observacio}

\begin{observacio}[Objecte de parts]
En un topos, l'exponencial $\Omega^{E}$ mereix un nom propi: és l'\textbf{objecte de parts} $\mathcal{P}(E)$, l'anàleg categòric del conjunt de parts. En efecte, per la propietat de l'exponencial i la del classificador,
\[
\Hom(A,\Omega^{E})\cong\Hom(A\times E,\Omega)\cong\Sub(A\times E),
\]
de manera que els morfismes cap a $\Omega^{E}$ classifiquen relacions entre $A$ i $E$. En el cas $A=1$, com que $1\times E\cong E$, s'obté
\[
\Hom(1,\Omega^{E})\cong\Sub(E):
\]
els punts globals de $\Omega^{E}$ són exactament els subobjectes de $E$, que és el que justifica el nom d'objecte de parts. Aquesta és la porta d'entrada a la lògica d'ordre superior interna d'un topos, que és tema de l'estudi posterior.
\end{observacio}

\section{Exemples}

\begin{exemple}[Conjunts]
$\cat{Set}$ és un topos. Té límits finits, és cartesià tancat (l'exponencial és el conjunt de funcions) i té classificador $\Omega=\{0,1\}$ amb $\mathsf{true}$ l'element $1$. És el topos prototípic, i tota la teoria es pot llegir com una generalització d'aquest cas.
\end{exemple}

\begin{exemple}[Prefeixos]\index{prefeix@prefeix!topos de}\index{sedàs}
Per a qualsevol categoria petita $\cat{C}$, la categoria de prefeixos $\widehat{\cat{C}}=\cat{Set}^{\cat{C}\op}$ (exemple~\ref{ex:categoria-prefeixos}) és un topos. Ja sabíem que té límits finits, calculats objecte a objecte (proposició~\ref{prop:limits-punt-a-punt}), i que és cartesiana tancada; el que hi falta és el classificador. Aquest es construeix amb els \textbf{sedassos}. Un sedàs sobre un objecte $C$ és un conjunt $S$ de morfismes amb codomini $C$ tal que
\[
f\colon D\to C,\quad f\in S,\quad g\colon X\to D
\qquad\Longrightarrow\qquad
f\comp g\in S.
\]
El prefeix $\Omega$ envia cada objecte $C$ al conjunt $\Omega(C)$ dels sedassos sobre $C$, i cada morfisme $h\colon D\to C$ a l'aplicació
\[
\Omega(h)\colon\Omega(C)\to\Omega(D),
\qquad
S\longmapsto h^{*}S=\{g\colon X\to D\mid h\comp g\in S\},
\]
que torna a donar un sedàs. El punt distingit $\mathsf{true}\colon 1\to\Omega$ tria a cada $C$ el sedàs màxim, format per tots els morfismes cap a $C$.

Si $U\rightarrowtail F$ és un subprefeix, el seu morfisme característic envia $x\in F(C)$ al sedàs dels morfismes que porten $x$ dins de $U$:
\[
\chi_U(C)(x)=\{f\colon D\to C\mid F(f)(x)\in U(D)\}.
\]
Es comprova que aquest conjunt és un sedàs, que les aplicacions $\chi_U(C)$ són naturals en $C$, i que $U$ és el pullback de $\mathsf{true}$ al llarg de $\chi_U$. Aquest exemple cobreix una quantitat enorme de casos i és el pont cap a la teoria de feixos.
\end{exemple}

\begin{exemple}[El topos dels grafs: el classificador calculat]\label{ex:omega-graph}\index{graf dirigit!topos}\index{classificador de subobjectes!a Graph@a $\cat{Graph}$}
Al capítol~\ref{cap:transformacions} vam veure que $\cat{Graph}\cong\cat{Set}^{\cat{P}}$, i com que $\cat{Set}^{\cat{P}}=\cat{Set}^{(\cat{P}\op)\op}$, és la categoria de prefeixos sobre $\cat{P}\op$. Per l'exemple anterior, $\cat{Graph}$ és, doncs, un topos. Calculem-ne el classificador.

\emph{Els vèrtexs i les arestes de $\Omega$.} Els dos representables són el graf $\mathbf V$ d'un sol vèrtex i el graf $\mathbf E$ d'una sola aresta $e\colon s\to t$ (exemple~\ref{ex:homfunctor-grafs}): els vèrtexs d'un graf $G$ corresponen als morfismes $\mathbf V\to G$, i les arestes, als morfismes $\mathbf E\to G$. Aplicant-ho a $\Omega$ i fent servir $\Hom(-,\Omega)\cong\Sub$, obtenim
\[
V_\Omega\cong\Sub(\mathbf V),\qquad E_\Omega\cong\Sub(\mathbf E).
\]
El graf $\mathbf V$ té dos subgrafs, el buit i ell mateix; així, $\Omega$ té dos vèrtexs, que anomenem $0$ (\enquote{el vèrtex no hi és}) i $1$ (\enquote{el vèrtex hi és}). El graf $\mathbf E$ té cinc subgrafs (exemple~\ref{ex:subgrafs}): $\varnothing$, $\{s\}$, $\{t\}$, $\{s,t\}$ i $\{s,t,e\}$. Així, $\Omega$ té cinc arestes.

\emph{Origen i destí.} L'origen d'una aresta $K\subseteq\mathbf E$ de $\Omega$ s'obté reindexant $K$ al llarg del morfisme $\mathbf V\to\mathbf E$ que tria $s$; és a dir, val $1$ si $s\in K$ i $0$ altrament. El destí, igual amb $t$. Anomenant les arestes per la informació que codifiquen, queda
\[
\renewcommand{\arraystretch}{1.2}
\begin{array}{c|c|l}
\text{aresta} & \text{subgraf de }\mathbf E & \text{de}\to\text{a}\\ \hline
\bot & \varnothing & 0\to0\\
\sigma & \{s\} & 1\to0\\
\tau & \{t\} & 0\to1\\
\partial & \{s,t\} & 1\to1\\
\top & \{s,t,e\} & 1\to1
\end{array}
\]
El graf $\Omega$ té, doncs, dos vèrtexs, un bucle a $0$, una aresta en cada sentit entre $0$ i $1$, i \emph{dos} bucles a $1$ (vegeu també \cite[exemple~15.4.3]{barrwells}). El morfisme $\mathsf{true}\colon\mathbf L\to\Omega$, definit sobre el graf terminal d'un bucle (exemple~\ref{ex:producte-tres-contextos}), envia el vèrtex a $1$ i el bucle a $\top$.

\emph{El morfisme característic d'un subgraf.} Si $H\subseteq G$ és un subgraf, $\chi_H$ envia cada vèrtex a $1$ o a $0$ segons que sigui o no a $H$. Cada aresta $x\colon v\to w$ va a l'aresta de $\Omega$ que registra si $v$, $w$ i $x$ són a $H$: a $\top$ si $x\in H$; a $\partial$ si $v,w\in H$ però $x\notin H$; i a $\sigma$, $\tau$ o $\bot$ segons quins extrems hi siguin. Per exemple, sigui $G$ el cicle de vèrtexs $a,b,c$ i arestes $e\colon a\to b$, $f\colon b\to c$, $k\colon c\to a$, i sigui $H$ el subgraf amb vèrtexs $a,b$ i aresta $e$. Aleshores
\[
\chi_H\colon\quad a,b\mapsto1,\quad c\mapsto0,\qquad e\mapsto\top,\quad f\mapsto\sigma,\quad k\mapsto\tau.
\]
És un morfisme de grafs, perquè $\sigma\colon1\to0$ i $f\colon b\to c$ tenen extrems compatibles, i el mateix per a les altres arestes. A més, el pullback de $\mathsf{true}$ al llarg de $\chi_H$, que es calcula component a component, té per vèrtexs els que van a $1$, $\{a,b\}$, i per arestes les que van a $\top$, $\{e\}$: és exactament $H$. La unicitat també es veu directament. Els vèrtexs de $H$ han d'anar a $1$ i els altres a $0$. Una aresta ha d'anar a una aresta de $\Omega$ amb els extrems ja determinats, i a $\top$ exactament si és a $H$; amb aquestes condicions, en cada cas hi ha una sola opció.

\emph{Més de dos valors de veritat.} La diferència entre $\partial$ i $\top$ és el que fa interessant aquest topos. Al graf $\mathbf E$, el subgraf $\{s\}$ i el subgraf $\{t\}$ són disjunts, i $\{t\}$ és el més gran dels subgrafs disjunts de $\{s\}$, perquè tot subgraf que contingui l'aresta ha de contenir $s$. Però la seva unió és $\{s,t\}$, que no és tot $\mathbf E$: hi falta l'aresta, que és justament el que distingeix $\partial$ de $\top$. A $\Sub(\mathbf E)$, doncs, un predicat i la seva negació no cobreixen tot l'objecte, i el principi del tercer exclòs falla. El classificador d'aquest topos és un graf amb estructura pròpia, no un conjunt $\{0,1\}$. La mateixa categoria que ens ha servit per introduir objectes, functors, Yoneda i límits és, així, un exemple genuí de topos no booleà.
\end{exemple}

\begin{observacio}[$\Omega$ com a objecte de valors de veritat]\label{obs:omega-veritat}\index{Omega@$\Omega$!valors de veritat}\index{topos!booleà}
La noció de topos va aparèixer primer a l'escola de Grothendieck de geometria algebraica, com a generalització de la d'espai topològic. Un dels seus aspectes més fascinants, però, és la relació amb la lògica. Per la representabilitat de $\Sub$, un morfisme $\varphi\colon E\to\Omega$ \emph{és} un predicat sobre $E$, la seva \enquote{extensió} és el subobjecte que classifica, i $\mathsf{true}$ fa el paper del valor \enquote{cert}. En aquest sentit precís ---i només en aquest--- $\Omega$ és un objecte de valors de veritat.

Això no vol dir que $\Omega$ tingui dos valors, ni que la lògica sigui clàssica. En tot topos cada $\Sub(E)$ és una àlgebra de Heyting (nota d'ampliació següent), de manera que la lògica interna és, en general, intuïcionista. Un topos s'anomena \textbf{booleà} quan totes aquestes àlgebres de Heyting són àlgebres de Boole, és a dir, quan per a tot subobjecte $U$ es compleix $U\vee\neg U=\top$. $\cat{Set}$ i $\cat{FinSet}$ són booleans; $\cat{Graph}$ no ho és, com hem vist a l'exemple~\ref{ex:omega-graph}. Fins i tot els \enquote{punts} de $\Omega$, els valors de veritat globals $1\to\Omega$, poden ser més de dos. I, a la inversa, tenir exactament dos valors globals no fa booleà un topos: els valors globals només compten $\Hom(1,\Omega)\cong\Sub(1)$, mentre que la booleanitat és una condició sobre tots els $\Sub(E)$. Els prefeixos sobre el monoide $(\mathbb N,+)$ tenen dos valors globals i no són booleans (vegeu els Exercicis per a tots dos fenòmens). Aquesta lectura permet interpretar la lògica ---de primer ordre i d'ordre superior--- dins de qualsevol topos. El seu desenvolupament és, precisament, l'objecte de l'estudi de la lògica categòrica.
\end{observacio}

\begin{nota}[Ampliació: $\Sub(E)$ és una àlgebra de Heyting]\label{nota:sub-heyting}\index{àlgebra de Heyting!subobjectes d'un topos}
L'afirmació de l'observació anterior és un teorema de la teoria de topos que aquest volum \emph{no demostra}; la fem servir només per situar la qüestió de la booleanitat. L'enunciat és aquest: en un topos elemental, cada ordre parcial $\Sub(E)$ té màxim, mínim, ínfims i suprems binaris i una implicació $\Rightarrow$ adjunta per la dreta de $-\wedge U$, i les reindexacions $f^{*}$ preserven tota aquesta estructura.

Algunes peces ja les tenim. El màxim és $[\id_E]$, i els ínfims són pullbacks (capítol~\ref{cap:subobjectes}). Les altres requereixen resultats que no hem provat. El mínim i els suprems es construeixen amb l'objecte inicial, els coproductes i les imatges; la seva existència fa servir que tot topos té colímits finits i és una categoria regular. La implicació s'obté amb el classificador i els exponencials. A $\cat{Graph}$, les dues construccions es poden verificar a mà sobre els subgrafs, cosa que justifica independentment l'exemple~\ref{ex:omega-graph}. Per a les demostracions, vegeu \cite[cap.~IV]{maclanemoerdijk} i \cite{johnstone}; el desenvolupament correspon a l'estudi de la lògica categòrica.
\end{nota}

\begin{resumcap}
\apartat{Idees centrals}
\begin{enumerate}
  \item El classificador $\Omega$ és un objecte amb $\mathsf{true}\colon 1\to\Omega$ tal que cada subobjecte té un únic morfisme característic, del qual és el pullback; la pertinença d'un element generalitzat esdevé una igualtat de morfismes cap a $\Omega$.
  \item Tenir classificador equival a la representabilitat del functor $\Sub$, i $\Omega$ n'és l'objecte representant: els predicats esdevenen morfismes $E\to\Omega$ i la substitució, precomposició.
  \item Un topos elemental té límits finits, exponencials i classificador; en un topos, $\Omega^{E}$ és l'objecte de parts de $E$.
  \item $\cat{Set}$ és el topos prototípic, amb $\Omega=\{0,1\}$.
  \item Tota categoria de prefeixos sobre una categoria petita és un topos, amb $\Omega$ el prefeix de sedassos; en particular, $\cat{Graph}$ és un topos, amb un classificador de dos vèrtexs i cinc arestes. No tot topos elemental és de prefeixos: $\cat{FinSet}$ no ho és.
  \item $\Omega$ és un objecte de valors de veritat perquè classifica predicats, però la lògica d'un topos és en general intuïcionista: $\cat{Graph}$ no és booleà.
\end{enumerate}
\apartat{Errors típics} Oblidar la hipòtesi de petitesa local\,/\,bona potenciació en escriure $\Sub\colon\cat{C}\op\to\cat{Set}$; confondre topos elemental amb categoria de prefeixos; confondre el classificador amb un conjunt de dos elements (només a $\cat{Set}$ el model prototípic és $\{0,1\}$; en un topos general, $\Omega$ és un objecte de la categoria i pot tenir una estructura molt més rica); i suposar que la lògica interna de qualsevol topos és clàssica: en general és intuïcionista, i la booleanitat és una propietat addicional.
\apartat{Lectura recomanada} \cite[cap.~8]{awodey}, especialment §8.8, per a classificadors, topos de prefeixos i la connexió amb Yoneda; \cite[cap.~15]{barrwells} per a una presentació alternativa i l'exemple dels grafs; \cite[cap.~1--2]{maclanemoerdijk} per continuar cap a la teoria de topos i feixos; \cite[§§12 i 14]{streicher} per als topos elementals i exercicis amb topos de prefeixos; \cite[article~VI, sessions~32--33]{lawvereschanuel} per a una presentació molt elemental dels subobjectes, la veritat i els topos; \cite{johnstone} com a obra de consulta avançada.
\end{resumcap}

\chapter[Pont cap a la lògica categòrica]{\texorpdfstring{Pont cap a la lògica\\ categòrica}{Pont cap a la lògica categòrica}}
\label{cap:pont}

\begin{objectius}
\apartat{Prerequisits} Tot el volum, especialment subobjectes, reindexació i categories regulars (cap.~\ref{cap:subobjectes}) i els classificadors (cap.~\ref{cap:topos}).
\apartat{Funció del capítol} Aquest capítol és un epíleg sense demostracions ni exercicis. No desenvolupa una teoria nova, però introdueix la terminologia mínima necessària per dibuixar el mapa que porta de les construccions d'aquest volum a la lògica categòrica. Serveix d'orientació per a qui vulgui continuar cap al volum següent.
\end{objectius}

Aquest volum tanca aquí. Els capítols anteriors han construït, pas a pas, el llenguatge categòric bàsic necessari per al programa que anunciem: categories i morfismes, functors, transformacions naturals i equivalències, representabilitat i el lema de Yoneda, límits i colímits, adjuncions, tancament cartesià, subobjectes, imatges i factoritzacions, i els classificadors de subobjectes. Aquest epíleg mostra què d'aquest bagatge es prolonga en l'estudi de la \textbf{lògica categòrica}, i com: cada eina apresa aquí té una contrapartida lògica precisa a l'altra banda del pont.

\section{El que aquest volum deixa preparat}

Convé fer inventari del que ja tenim, perquè és exactament la infraestructura que la lògica categòrica pressuposa:
\begin{itemize}
\item \textbf{Categories, functors i naturalitat}: el llenguatge bàsic en què s'expressa tota la resta.
\item \textbf{La categoria prima de la deduïbilitat}: la primera aparició, al capítol inicial, de la idea que una lògica es pot veure com una categoria.
\item \textbf{Representabilitat i Yoneda}: la manera d'identificar un objecte per com es relaciona amb tots els altres; el classificador de subobjectes n'és un cas.
\item \textbf{Límits finits, i molt especialment els pullbacks}: la base per interpretar contextos, substitució i relacions.
\item \textbf{Adjuncions}: el patró que reapareixerà en operacions lògiques centrals, especialment en la implicació i en els quantificadors existencial i universal.
\item \textbf{Tancament cartesià i càlcul lambda}: la semàntica dels tipus i els termes.
\item \textbf{Subobjectes i reindexació}: la interpretació dels predicats i la substitució.
\item \textbf{Imatges i factoritzacions}: la construcció que, a l'altra banda, interpretarà la quantificació existencial.
\item \textbf{Classificadors de subobjectes i topos}: el marc on la lògica esdevé interna a la categoria.
\end{itemize}

Amb tot això, el salt cap a la lògica categòrica deixa de ser un salt: és una continuació natural.

\section{El mapa: de regular a topos}\label{sec:mapa}\index{categoria!regular}\index{categoria!coherent}\index{categoria!de Heyting}

Abans de descriure el recorregut, val la pena fixar un mapa que ordena, d'una vegada, quina estructura categòrica interpreta quin fragment lògic. És la classificació que dona sentit retrospectiu a les distincions que hem anat marcant ---per exemple, per què la regularitat basta per a l'existencial però no per a la disjunció ni per a l'universal.

La idea directriu és que cada nivell \emph{afegeix} una peça d'estructura a les fibres de subobjectes $\Sub(A)$, i que cada peça interpreta un connectiu o quantificador:
\begin{center}
\renewcommand{\arraystretch}{1.3}
\noindent\begin{tabularx}{\linewidth}{@{}>{\raggedright\arraybackslash}X >{\raggedright\arraybackslash}X >{\raggedright\arraybackslash}X@{}}
\hline
\textbf{Estructura categòrica} & \textbf{Nova estructura a $\Sub(A)$} & \textbf{Fragment lògic}\\
\hline
categoria amb límits finits & $\top$, $\wedge$ (interseccions) & conjunció, veritat\\
categoria \textbf{regular} & imatges estables: $\exists_f$ & $+\ \exists$\\
categoria \textbf{coherent} & joins finits estables: $\bot$, $\vee$ & $+\ \vee,\ \bot$\\
categoria \textbf{de Heyting} & implicació $\Rightarrow$; $\forall_f$ & $+\ \Rightarrow,\ \forall$: lògica intuïcionista de primer ordre\\
\textbf{topos} elemental & classificador $\Omega$, objectes de parts, exponencials & lògica intuïcionista d'ordre superior\\
\hline
\end{tabularx}
\renewcommand{\arraystretch}{1.0}
\end{center}

La lectura de la taula, de dalt a baix, és acumulativa i respon les preguntes que aquest volum ha deixat obertes deliberadament:
\begin{itemize}
  \item \textbf{Regular} ($\top,\wedge,\exists$). La factorització imatge estable dona l'existencial com a adjunt per l'esquerra de la reindexació. És fins on arriba el capítol de subobjectes: per això hi vam construir només $\exists_f$, i no $\forall_f$ ni la disjunció.
  \item \textbf{Coherent} ($+\vee,\bot$). S'hi afegeixen \emph{joins finits estables} a les fibres, que interpreten la disjunció i la falsedat. Aquesta és l'estructura que a la teoria de subobjectes anunciàvem com a necessària per a les \enquote{unions}, i que la regularitat sola no proporciona.
  \item \textbf{De Heyting} ($+\Rightarrow,\forall$). Les fibres esdevenen àlgebres de Heyting ---amb implicació, adjunta per la dreta de la conjunció--- i cada reindexació $f^{*}$ adquireix un adjunt per la dreta $\forall_f$. Apareixen la implicació i el quantificador universal.
  \item \textbf{Topos} elemental. El classificador $\Omega$ i els exponencials clouen el quadre amb una lògica intuïcionista d'ordre superior, on fins i tot els predicats sobre predicats tenen extensió.
\end{itemize}

\index{categoria!localment cartesiana tancada}En paral·lel, la semàntica dels \emph{tipus} segueix una altra línia. Les categories cartesianes tancades interpreten el càlcul lambda simplement tipat (capítol~\ref{cap:cartesianes}), mentre que les categories \textbf{localment cartesianes tancades} proporcionen productes dependents: per a cada $f\colon A\to B$, la reindexació $f^{*}\colon\cat{C}/B\to\cat{C}/A$ admet un adjunt per la dreta $\Pi_f$. Aquesta via és especialment important per a la teoria de tipus dependents. No és un esglaó de la cadena regular--coherent--Heyting, però totes dues vies es troben en els topos elementals, que són localment cartesians tancats i tenen lògica de Heyting als subobjectes.

Aquest mapa no és matèria d'aquest volum (no en demostrem les equivalències), però és la brúixola que orienta tot el que segueix: cada fila és, de fet, un capítol del programa de la lògica categòrica.

\section{El que continua a l'altra banda}

L'estudi de la lògica categòrica pren aquestes eines i les fa servir per construir una \textbf{semàntica de la lògica} de complexitat creixent. Sense entrar-hi ---perquè és l'objecte d'un estudi a part---, en dibuixem el recorregut, perquè el lector vegi on desemboca cada cosa que ha après:

\begin{description}
\item[Lògica d'equacions.] El punt de partida: teories algebraiques a l'estil de Lawvere, on la categoria sintàctica té productes finits i els models són functors que els preserven. És la lògica que ja s'endevina en la universalitat dels productes finits: els contextos es combinen per producte i les operacions esdevenen morfismes entre potències finites d'un objecte. El tancament cartesià pertany a una via posterior, la del càlcul lambda i la implicació.

\item[Lògica de límits finits.] El pont immediat: teories amb diverses menes, operacions parcialment definides i axiomes d'existència i unicitat. La categoria sintàctica passa a tenir \emph{límits finits}, no només productes; és el marc de les teories essencialment algebraiques.

\item[Lògica regular.] Fórmules amb igualtat, conjunció i quantificació existencial, escrites com a seqüents en context. Aquí és on les \textbf{imatges regulars} construïdes en aquest volum interpreten l'existencial, i on la factorització imatge interpreta l'eliminació de testimonis. El resultat central és que, per a una teoria regular $T$ i una categoria regular $\cat{C}$, els models de $T$ en $\cat{C}$ es descriuen mitjançant functors regulars des de la categoria sintàctica regular $\cat{C}_T$ cap a $\cat{C}$; amb la noció adequada de morfisme de models, aquesta correspondència s'organitza categòricament.

\item[Lògica coherent i geomètrica.] S'hi afegeixen disjuncions finites (coherent) i disjuncions indexades per conjunts (geomètrica). És el marc que condueix als feixos i, de nou, als topos.

\item[Lògica intuïcionista de primer ordre.] Apareixen la implicació i el quantificador universal. La idea organitzadora és que les fibres de subobjectes formen \textbf{àlgebres de Heyting}: la implicació és l'adjunt dret de la conjunció, i $\forall$ és l'adjunt dret de la substitució. Els adjunts de la reindexació que aquí hem vist com a exemple es converteixen allà en els quantificadors del sistema.

\item[Doctrines i fibracions lògiques.] La síntesi: contextos, substitució, predicats i quantificadors s'organitzen en una \emph{doctrina indexada}. Les condicions de coherència ---Beck--Chevalley i Frobenius--- expressen que substituir i quantificar interactuen correctament. És el llenguatge que unifica tots els casos anteriors.

\item[Teories classificadores i llenguatge intern.] El punt culminant: la categoria o el topos classificador d'una teoria, amb la seva propietat universal, i el llenguatge intern que permet raonar dins d'una categoria com si els seus objectes fossin conjunts, sempre que les regles emprades corresponguin a l'estructura disponible.

\item[Lògica clàssica i booleanització.] Com a epíleg, què s'afegeix en imposar el tercer exclòs, quines construccions deixen de ser constructives i com, en teoria de topos, es passa a la booleanització d'un topos. La lògica clàssica apareix així com una especialització, no com el rerefons silenciós de tota la teoria.
\end{description}

\begin{nota}[Ampliació: estructures sense conjunt subjacent]
Al capítol~\ref{cap:categories} hem partit d'estructures que tenen un conjunt subjacent. No totes les estructures matemàtiques són conjunts amb una estructura afegida, i el programa anterior ho fa visible. A cada teoria geomètrica li correspon un \emph{topos classificador}, la categoria de punts del qual és equivalent a la categoria de models de la teoria a $\cat{Set}$ \cite{caramello}. Però aquest \enquote{espai} pot ser no trivial i no tenir cap punt. Per exemple, la teoria d'una aplicació exhaustiva $\mathbb N\to\mathbb R$ no té cap model a $\cat{Set}$, pel teorema de Cantor; tanmateix, el seu \emph{locale} classificador és no trivial, i el topos de feixos sobre aquest locale, que n'és el topos classificador, no té cap punt \cite{blechschmidt}. Aquestes situacions queden fora d'aquest volum; les esmentem perquè mostren que la perspectiva de les fletxes no depèn de disposar d'elements.
\end{nota}

\section{Un exemple: la categoria sintàctica d'una teoria}\label{sec:sintactica}\index{categoria!sintàctica}

Fins ara la correspondència lògica--categòrica ha aparegut repartida en exemples solts. Val la pena reunir-la en un únic exemple pont, complet i concret, que il·lustra de cop l'eslògan \enquote{una teoria es codifica en una categoria sintàctica}. Prenem la teoria algebraica més simple no trivial: una \textbf{mena} $s$ i una operació binària $m\colon s\times s\to s$, sense cap axioma (la teoria dels \emph{magmes}). Construïm pas a pas la seva categoria sintàctica $\cat{Syn}$.

\begin{enumerate}
  \item \textbf{Contextos com a objectes.} Un \emph{context} és una llista finita de variables de la mena $s$, que escrivim $[x_1,\dots,x_n]$ i abreugem $s^n$. Els objectes de $\cat{Syn}$ són els contextos, un per a cada $n\geq 0$; el context buit $s^0$ és l'objecte terminal.
  \item \textbf{Termes com a morfismes.} Un morfisme $s^n\to s^m$ és una llista de $m$ termes $(t_1,\dots,t_m)$, cadascun construït amb les variables $x_1,\dots,x_n$ i l'operació $m$ (per exemple, $m(x_1,m(x_2,x_1))$). La composició és la \textbf{substitució}: components d'un morfisme dins d'un altre. La identitat de $s^n$ és la llista de variables $(x_1,\dots,x_n)$.
  \item \textbf{Productes com a concatenació de contextos.} El producte $s^n\times s^m$ és el context $s^{n+m}$: juxtaposar dos contextos és formar-ne el producte, i les projeccions obliden les variables sobrants. Així $\cat{Syn}$ té tots els productes finits, i l'operació $m$ del llenguatge és, ara, un morfisme distingit $s\times s\to s$ de la categoria.
  \item \textbf{Equacions com a igualtats de morfismes.} Si la teoria tingués axiomes ---posem l'associativitat de $m$---, s'imposarien identificant els dos morfismes $s^3\to s$ corresponents a $m(m(x_1,x_2),x_3)$ i $m(x_1,m(x_2,x_3))$. Una equació entre termes \emph{és}, literalment, una igualtat entre morfismes de $\cat{Syn}$.
  \item \textbf{Models com a functors.} Un model de la teoria en una categoria $\cat{C}$ amb productes finits ---per exemple, un magma ordinari a $\cat{Set}$--- és exactament un functor $\cat{Syn}\to\cat{C}$ que preserva els productes finits: tria un objecte per a $s$ i un morfisme per a $m$, respectant tota la sintaxi. Els homomorfismes de models són transformacions naturals. Aquesta és la forma precisa de \enquote{els models són functors}.
  \item \textbf{Fórmules i existencial.} Per passar de la teoria algebraica anterior a la lògica regular cal enriquir també la categoria sintàctica: la categoria $\cat{Syn}$ dels passos anteriors no conté automàticament els subobjectes que representen totes les fórmules. En la \textbf{categoria sintàctica regular} $\cat{C}_T$ d'una teoria amb predicats, una fórmula $\varphi(\bar x)$ en context determina, llevat d'equivalència demostrable, un subobjecte
  \[
  [\bar x\mid\varphi]\rightarrowtail[\bar x\mid\top].
  \]
  La substitució al llarg d'un morfisme de contextos s'interpreta per \textbf{reindexació}, és a dir, per pullback. En particular, si $\pi\colon s^{n+1}\to s^n$ és la projecció que oblida la darrera variable, $\pi^{*}$ és el debilitament, que considera un predicat sobre $s^n$ en el context ampliat. Si $P\rightarrowtail s^{n+1}$ representa $\varphi(\bar x,y)$, la \textbf{imatge} del compost
  \[
  P\rightarrowtail s^{n+1}\xrightarrow{\ \pi\ }s^n
  \]
  representa $\exists y\,\varphi(\bar x,y)$. Així, l'existencial és la imatge d'un predicat al llarg d'una projecció o, equivalentment, l'adjunt per l'esquerra de la reindexació (capítol~\ref{cap:subobjectes}).
  \item \textbf{Un càlcul complet en un model.} Fem concret tot l'anterior en un sol model, el magma $(\mathbb N,+)$, és a dir, el functor $M\colon\cat{Syn}\to\cat{Set}$ amb $M(s)=\mathbb N$ i $M(m)=+$. Prenem la fórmula $\varphi(x,y)$: $m(y,y)=x$, en el context $[x,y]$.
  \begin{itemize}
    \item \emph{Una equació és un equalitzador.} Els dos termes $m(y,y)$ i $x$ són morfismes $s^2\to s$, interpretats com les aplicacions $(x,y)\mapsto y+y$ i $(x,y)\mapsto x$ de $\mathbb N^2$ a $\mathbb N$. La fórmula s'interpreta com el seu equalitzador,
    \[
    P=\{(x,y)\in\mathbb N^2\mid y+y=x\}\rightarrowtail\mathbb N^2.
    \]
    Equivalentment, $P$ és el pullback de la diagonal $\delta_{\mathbb N}\colon\mathbb N\rightarrowtail\mathbb N\times\mathbb N$ al llarg del parell format pels dos termes: la igualtat és la diagonal.
    \item \emph{L'existencial és una imatge.} La projecció $\pi\colon\mathbb N^2\to\mathbb N$ oblida $y$. La imatge de $P\rightarrowtail\mathbb N^2\to\mathbb N$ és $\{x\mid\exists y\ y+y=x\}$, el conjunt dels nombres parells. Així, $\exists y\,(m(y,y)=x)$ s'interpreta com el predicat \enquote{$x$ és parell}, és a dir, com $\exists_\pi[P]$.
    \item \emph{El predicat és un morfisme cap a $\Omega$.} El seu morfisme característic $\mathbb N\to\{0,1\}$ val $1$ exactament als parells, com a l'exemple~\ref{ex:caracteristic-set}.
    \item \emph{La substitució és un pullback.} Substituir $x$ pel terme $m(z,m(z,z))$, en el context $[z]$, dona la fórmula $\exists y\,(m(y,y)=m(z,m(z,z)))$. Categòricament, és reindexar el predicat al llarg del morfisme $s\to s$ definit pel terme, interpretat com $z\mapsto 3z$. Per tant, s'interpreta com $\{z\mid 3z\text{ és parell}\}$, que és el conjunt dels parells. En canvi, substituir pel terme $m(z,z)$ dona $\{z\mid 2z\text{ és parell}\}=\mathbb N$, el subobjecte màxim: la fórmula $\exists y\,(m(y,y)=m(z,z))$ és demostrable, amb el testimoni $y=z$, i el model ho reflecteix.
  \end{itemize}
  Un sol exemple recorre així les peces del volum: termes com a morfismes, equacions com a equalitzadors i igualtat com a diagonal (capítols~\ref{cap:limits} i~\ref{cap:cartesianes}), existencial com a imatge (capítol~\ref{cap:subobjectes}), predicats com a morfismes cap a $\Omega$ (capítol~\ref{cap:topos}) i substitució com a pullback.
\end{enumerate}

Amb aquest exemple, cada peça del llibre troba el seu lloc: els contextos són objectes, els termes són morfismes, la substitució és composició i reindexació, els productes són concatenacions, les equacions són igualtats de fletxes, els models són functors que preserven estructura, i els quantificadors, quan existeixen, són adjunts de la reindexació. La categoria sintàctica és el punt on la sintaxi lògica i la semàntica categòrica es toquen; construir-la en general, amb l'estructura exacta que cada fragment lògic requereix, és una de les primeres tasques del volum següent.

\section{Cloenda}

El fil que recorre tot aquest programa és un de sol, i ja el coneixem: \emph{traduir entre lògica i categories}. En el marc categòric adequat a cada fragment lògic, un predicat s'interpreta com un subobjecte; la substitució, com un pullback; la conjunció, com una intersecció; l'existencial, com una imatge; la implicació i l'universal, mitjançant adjunts; una teoria es codifica en una categoria sintàctica i els seus models es descriuen per functors que preserven l'estructura pertinent. Les teories algebraiques, regulars, coherents o geomètriques tenen categories sintàctiques amb l'estructura corresponent, i per això teories sintàcticament diferents poden tenir semàntiques categòriques equivalents, o fins i tot el mateix topos classificador. Cadascuna d'aquestes traduccions descansa sobre una construcció que aquest volum ha desenvolupat amb detall.

No hem demostrat cap teorema de la lògica categòrica en aquestes pàgines, i era deliberat: la seva demostració pertany al volum següent. El que sí que hem fet és deixar el terreny preparat, de manera que quan aquells teoremes arribin, cap de les seves peces categòriques resulti desconeguda. Si el lector, en trobar més endavant que \enquote{l'existencial és un adjunt per l'esquerra de la substitució, estable per canvi de base}, reconeix en aquesta frase la factorització imatge, l'adjunció i el pullback que ha estudiat aquí, aleshores aquest volum haurà complert la seva funció.

\section{Mapa del volum següent}\label{sec:mapa-volum-seguent}

La secció~\ref{sec:mapa} ordenava la continuació per fragments lògics. Per acabar, l'ordenem per temes: per a cadascun dels cinc temes centrals del volum de lògica categòrica, què en queda establert aquí i què s'hi afegirà.
\begin{center}
\footnotesize
\renewcommand{\arraystretch}{1.3}
\noindent\begin{tabularx}{\linewidth}{@{}>{\raggedright\arraybackslash}p{0.2\linewidth}>{\raggedright\arraybackslash}X>{\raggedright\arraybackslash}X@{}}
\toprule
\textbf{Tema} & \textbf{En aquest volum} & \textbf{Al volum següent}\\
\midrule
substitució i reindexació & $f^{*}$ per pullback a $\cat{C}/B$ (capítol~\ref{cap:limits}) i a $\Sub(B)$; el functor $\Sub\colon\cat{C}\op\to\cat{Poset}$ (capítol~\ref{cap:subobjectes}) & la substitució de termes en fórmules com a reindexació en una doctrina; la condició de Beck--Chevalley\\
connectius sobre subobjectes & $\top$ i $\wedge$ com a pullbacks (capítol~\ref{cap:subobjectes}); $\Rightarrow$ als preordres cartesians tancats (capítol~\ref{cap:cartesianes}); $\wedge$ sobre $\Omega$ (capítol~\ref{cap:topos}) & $\bot$ i $\vee$ estables (categories coherents); $\Rightarrow$ a les fibres (categories de Heyting)\\
quantificadors com a adjunts & $\exists_f\dashv f^{*}\dashv\forall_f$ a $\cat{Set}$ (capítol~\ref{cap:adjuncions}); $\exists_f$ com a imatge en categories regulars (capítol~\ref{cap:subobjectes}) & $\forall_f$ en categories de Heyting; la condició de Frobenius; la quantificació al llarg de projeccions de contextos\\
igualtat & la diagonal $\delta_A$ (capítol~\ref{cap:cartesianes}); la parella nucli, que és el pullback de $\delta_B$ al llarg de $f\times f$ (capítol~\ref{cap:subobjectes}); les equacions com a equalitzadors (secció~\ref{sec:sintactica}) & la igualtat com a $\exists_{\delta_A}(\top)$, on $\exists_{\delta_A}$ és l'adjunt per l'esquerra de la contracció $\delta_A^{*}$; les regles de la igualtat com a propietats d'aquesta adjunció\\
semàntica interna & $\Omega$, la pertinença per elements generalitzats i l'objecte de parts (capítol~\ref{cap:topos}); l'exemple de la secció~\ref{sec:sintactica} & el llenguatge intern d'un topos, la semàntica de Kripke--Joyal, la lògica d'ordre superior, i la correcció i la completesa\\
\bottomrule
\end{tabularx}
\end{center}
Cap fila de la columna central no demana res que el lector no hagi vist en aquest volum. La columna de la dreta és el programa del volum següent.

\begin{resumcap}
\apartat{El fil en una frase} En el marc categòric adequat a cada fragment lògic, un predicat s'interpreta com un subobjecte; la substitució, com un pullback; la conjunció, com una intersecció; l'existencial, com una imatge; la implicació i l'universal, mitjançant adjunts; una teoria es codifica en una categoria sintàctica i els seus models es descriuen per functors que preserven l'estructura pertinent.
\apartat{Cap on continua} La lògica pròpiament dita ---sintaxi, correcció i completesa, i les condicions de coherència Beck--Chevalley i Frobenius--- es reserva per al volum de lògica categòrica, del qual aquest text és la preparació. La secció~\ref{sec:mapa-volum-seguent} en dona el mapa per temes: substitució, connectius, quantificadors, igualtat i semàntica interna.
\apartat{Lectura recomanada} \cite{lambekscott} i \cite{jacobs} per a la lògica categòrica i la teoria de tipus; \cite[§13]{streicher} per a la lògica interna dels topos; \cite{makkaireyes} i \cite{pitts} per a la lògica de primer ordre en categories regulars, coherents i de Heyting, que és el punt de partida del volum de lògica; \cite{maclanemoerdijk} per a la lògica geomètrica i els topos; \cite{johnstone} com a obra de consulta. Com a direccions complementàries: \cite{hott}, per a la teoria homotòpica de tipus, on la igualtat entre dos elements pot tenir estructura pròpia; i l'nLab \cite{nlab-classificador,nlab-locale}, com a recurs de consulta sobre topos classificadors i locales.
\end{resumcap}

\fi

% ============================================================
% PART II — PRÀCTICA
% Els capítols reprodueixen els noms i l'ordre de la part teòrica.
% ============================================================
\part{Pràctica}
\setcounter{chapter}{0}
\chapter{Categories}

En aquest capítol resolem amb detall una selecció d'exercicis corresponents al capítol \emph{Categories} de la part de Teoria. L'objectiu no és només arribar al resultat, sinó mostrar \emph{com} es raona en el llenguatge categòric: comprovar axiomes, seguir camins, traduir diagrames a equacions i reconèixer una mateixa estructura en conjunts, preordres, lògica i grafs. Convé intentar cada exercici abans de llegir-ne la solució.

\section{Definició de categoria}

\begin{exercici}\label{exr:pr-rel}
Comprova que els conjunts amb les relacions formen una categoria $\cat{Rel}$, on un morfisme $A\to B$ és una relació $R\subseteq A\times B$, la identitat de $A$ és la relació d'igualtat, i la composició de $R\subseteq A\times B$ amb $S\subseteq B\times C$ és
\[
S\comp R=\{(a,c)\in A\times C\mid \exists\,b\in B\text{ tal que }(a,b)\in R\text{ i }(b,c)\in S\}.
\]
\end{exercici}

\begin{solucio}
Hem de comprovar els dos axiomes. Comencem per les \textbf{identitats}. Sigui $\id_A=\{(a,a)\mid a\in A\}$ la relació d'igualtat sobre $A$. Per a tota relació $R\subseteq A\times B$,
\[
R\comp\id_A
=\{(a,b)\mid \exists\,a'\ (a,a')\in\id_A\text{ i }(a',b)\in R\}
=R,
\]
perquè $(a,a')\in\id_A$ obliga $a'=a$. Anàlogament $\id_B\comp R=R$.

Per a l'\textbf{associativitat}, siguin $R\subseteq A\times B$, $S\subseteq B\times C$ i $T\subseteq C\times D$. Un parell $(a,d)$ pertany a $T\comp(S\comp R)$ si i només si existeixen $b\in B$ i $c\in C$ tals que
\[
(a,b)\in R,
\qquad
(b,c)\in S,
\qquad
(c,d)\in T.
\]
La mateixa condició caracteritza $(T\comp S)\comp R$. Per tant
\[
T\comp(S\comp R)=(T\comp S)\comp R,
\]
i $\cat{Rel}$ és una categoria.
\end{solucio}

\section{Fletxes, camins i diagrames commutatius}

\begin{exercici}
Llegeix el diagrama següent i escriu l'equació que expressa que commuta:
\[
\begin{tikzpicture}[baseline=(m.center)]
  \matrix (m) [matrix of math nodes, row sep=2.7em, column sep=3.6em] {
    A & B \\
    C & D \\
  };
  \draw[-{Stealth[scale=1.1]}] (m-1-1) -- node[above] {$f$} (m-1-2);
  \draw[-{Stealth[scale=1.1]}] (m-1-1) -- node[left] {$u$} (m-2-1);
  \draw[-{Stealth[scale=1.1]}] (m-1-2) -- node[right] {$v$} (m-2-2);
  \draw[-{Stealth[scale=1.1]}] (m-2-1) -- node[below] {$g$} (m-2-2);
\end{tikzpicture}
\]
Explica també per què l'ordre de les composicions és el que escrius.
\end{exercici}

\begin{solucio}
Hi ha dos camins dirigits de $A$ a $D$. El camí superior-dret és
\[
A\xrightarrow{f}B\xrightarrow{v}D,
\]
i per tant determina $v\comp f$. El camí esquerre-inferior és
\[
A\xrightarrow{u}C\xrightarrow{g}D,
\]
i determina $g\comp u$. El quadrat commuta exactament quan
\[
\boxed{v\comp f=g\comp u}.
\]
L'ordre de cada composició es llegeix de dreta a esquerra: en $v\comp f$ seguim primer $f$ i després $v$.
\end{solucio}

\begin{exercici}
Considera a $\cat{Set}$ quatre còpies de $\mathbb Z$ i les aplicacions
\[
f(n)=n+1,
\qquad
u(n)=2n,
\qquad
v(n)=2n,
\qquad
g(n)=n+2.
\]
Comprova que el quadrat
\[
\begin{tikzpicture}[baseline=(m.center)]
  \matrix (m) [matrix of math nodes, row sep=2.8em, column sep=4.2em] {
    \mathbb Z & \mathbb Z \\
    \mathbb Z & \mathbb Z \\
  };
  \draw[-{Stealth[scale=1.05]}] (m-1-1) -- node[above] {$f$} (m-1-2);
  \draw[-{Stealth[scale=1.05]}] (m-1-1) -- node[left] {$u$} (m-2-1);
  \draw[-{Stealth[scale=1.05]}] (m-1-2) -- node[right] {$v$} (m-2-2);
  \draw[-{Stealth[scale=1.05]}] (m-2-1) -- node[below] {$g$} (m-2-2);
\end{tikzpicture}
\]
commuta.
\end{exercici}

\begin{solucio}
Cal comprovar que $v\comp f=g\comp u$. Per a qualsevol $n\in\mathbb Z$,
\[
(v\comp f)(n)=v(n+1)=2n+2,
\]
mentre que
\[
(g\comp u)(n)=g(2n)=2n+2.
\]
Les dues aplicacions coincideixen en tots els enters. Per tant
\[
v\comp f=g\comp u,
\]
i el quadrat commuta. Aquest exemple mostra com una afirmació sobre funcions es pot llegir simultàniament com una igualtat algebraica i com la commutativitat d'un diagrama.
\end{solucio}

\section{Primers exemples}

\begin{exercici}
Sigui $T$ un sistema deductiu amb una relació de deduïbilitat $\vdash_T$ reflexiva i transitiva. Comprova amb detall que les fórmules, amb una única fletxa $A\to B$ quan $A\vdash_T B$, formen una categoria prima $\cat{Prov}_T$.
\end{exercici}

\begin{solucio}
Els \textbf{objectes} són les fórmules. Entre dues fórmules $A$ i $B$ posem un únic morfisme $A\to B$ exactament quan $A\vdash_T B$.

La \textbf{identitat} existeix perquè la deduïbilitat és reflexiva:
\[
A\vdash_T A.
\]
Per tant hi ha una fletxa $\id_A\colon A\to A$.

La \textbf{composició} existeix perquè la deduïbilitat és transitiva. Si
\[
A\vdash_T B
\qquad\text{i}\qquad
B\vdash_T C,
\]
aleshores
\[
A\vdash_T C.
\]
Diagramàticament,
\[
\begin{tikzpicture}[baseline=(m.center)]
  \matrix (m) [matrix of math nodes, column sep=3.4em] { A & B & C \\ };
  \draw[-{Stealth[scale=1.05]}] (m-1-1) -- (m-1-2);
  \draw[-{Stealth[scale=1.05]}] (m-1-2) -- (m-1-3);
  \draw[-{Stealth[scale=1.0]}] (m-1-1) to[bend right=25] (m-1-3);
\end{tikzpicture}
\]

Finalment, entre dues fórmules hi ha com a màxim una fletxa. Això fa que les igualtats exigides pels axiomes d'associativitat i identitat siguin automàtiques sempre que les composicions existeixin. Per tant $\cat{Prov}_T$ és una categoria prima.

Cal remarcar què hem oblidat: si hi ha dues demostracions diferents de $B$ a partir d'$A$, aquesta categoria les identifica en una sola fletxa. Només conserva la relació de deduïbilitat.
\end{solucio}

\begin{exercici}
Siguin els grafs dirigits $G$ i $H$ següents:
\[
G:\quad 0\xrightarrow{a}1,
\qquad
H:\quad
\begin{tikzpicture}[baseline=(m.center)]
  \matrix (m) [matrix of math nodes, row sep=2.0em, column sep=2.8em] {
    & y \\
    x & z \\
  };
  \draw[-{Stealth[scale=1.0]}] (m-2-1) -- node[above left] {$r$} (m-1-2);
  \draw[-{Stealth[scale=1.0]}] (m-2-1) -- node[below] {$s$} (m-2-2);
\end{tikzpicture}
\]
Defineix $F_V(0)=x$, $F_V(1)=y$ i $F_E(a)=r$. Comprova que això defineix un morfisme de grafs $F\colon G\to H$. Explica per què substituir $F_E(a)=r$ per $F_E(a)=s$, mantenint $F_V$, no defineix un morfisme de grafs.
\end{exercici}

\begin{solucio}
En $G$ tenim
\[
s_G(a)=0,
\qquad
t_G(a)=1,
\]
i en $H$,
\[
s_H(r)=x,
\qquad
t_H(r)=y.
\]
Per tant
\[
F_V(s_G(a))=F_V(0)=x=s_H(r)=s_H(F_E(a))
\]
i
\[
F_V(t_G(a))=F_V(1)=y=t_H(r)=t_H(F_E(a)).
\]
Així es compleixen les dues equacions
\[
F_V\comp s_G=s_H\comp F_E,
\qquad
F_V\comp t_G=t_H\comp F_E,
\]
i $F$ és un morfisme de grafs.

Si poséssim $F_E(a)=s$, l'origen encara seria correcte perquè $s_H(s)=x$, però el destí fallaria:
\[
t_H(s)=z\neq y=F_V(1).
\]
Per tant el quadrat corresponent al destí no commutaria i no tindríem un morfisme de grafs.
\end{solucio}

\section{Isomorfismes}

\begin{exercici}
En una categoria qualsevol, comprova que si $f\colon A\to B$ i $g\colon B\to C$ tenen inversos, aleshores $g\comp f$ també en té, i
\[
(g\comp f)^{-1}=f^{-1}\comp g^{-1}.
\]
\end{exercici}

\begin{solucio}
La situació és
\[
A\xrightarrow{f}B\xrightarrow{g}C.
\]
Per recórrer aquest camí en sentit contrari hem d'invertir també l'ordre: primer $g^{-1}$ i després $f^{-1}$. Proposem, doncs, $f^{-1}\comp g^{-1}$ com a invers.

D'una banda,
\[
(f^{-1}\comp g^{-1})\comp(g\comp f)
=f^{-1}\comp(g^{-1}\comp g)\comp f
=f^{-1}\comp f
=\id_A,
\]
i de l'altra,
\[
(g\comp f)\comp(f^{-1}\comp g^{-1})
=g\comp(f\comp f^{-1})\comp g^{-1}
=g\comp g^{-1}
=\id_C.
\]
Per tant $(g\comp f)^{-1}=f^{-1}\comp g^{-1}$.
\end{solucio}

\begin{exercici}
Comprova que a $\cat{Rel}$ una relació $R\colon A\to B$ és un isomorfisme si i només si és el graf d'una bijecció de $A$ a $B$.
\end{exercici}

\begin{solucio}
Si $f\colon A\to B$ és una bijecció i $R$ és el seu graf, el graf de $f^{-1}$ és una relació $B\to A$ i les dues composicions són les relacions d'igualtat. Per tant $R$ és un isomorfisme.

Recíprocament, suposem que $R$ té un invers $S\colon B\to A$, és a dir,
\[
S\comp R=\id_A,
\qquad
R\comp S=\id_B.
\]
Per a cada $a\in A$, com $(a,a)\in S\comp R$, existeix algun $b\in B$ amb $(a,b)\in R$ i $(b,a)\in S$. Així cada $a$ es relaciona amb almenys un $b$.

Si $(a,b)\in R$ i $(a,b')\in R$, triem $b_0$ amb $(a,b_0)\in R$ i $(b_0,a)\in S$. Llavors
\[
(b_0,b)\in R\comp S=\id_B,
\qquad
(b_0,b')\in R\comp S=\id_B,
\]
i per tant $b=b_0=b'$. Així cada $a$ es relaciona amb un únic $b$, i $R$ és el graf d'una aplicació $f\colon A\to B$.

Aplicat a $S$, el mateix argument dona una aplicació $g\colon B\to A$. Les igualtats $S\comp R=\id_A$ i $R\comp S=\id_B$ es converteixen en
\[
g\comp f=\id_A,
\qquad
f\comp g=\id_B.
\]
Per tant $f$ és bijectiva i $R$ n'és el graf.
\end{solucio}

\begin{exercici}
Sigui $\cat{C}$ una categoria. Comprova que la relació \enquote{ésser isomorf} entre objectes de $\cat{C}$ és una relació d'equivalència.
\end{exercici}

\begin{solucio}
\textbf{Reflexivitat}: $\id_A$ és el seu propi invers, de manera que $A\cong A$.

\textbf{Simetria}: si $f\colon A\to B$ és un isomorfisme amb invers $f^{-1}\colon B\to A$, aleshores $f^{-1}$ és també un isomorfisme, amb invers $f$; per tant $B\cong A$.

\textbf{Transitivitat}: si $f\colon A\to B$ i $g\colon B\to C$ són isomorfismes, la composició $g\comp f$ també ho és pel primer exercici de la secció; per tant $A\cong C$.
\end{solucio}

\begin{exercici}
A la categoria prima $\cat{Prov}_T$, comprova que dues fórmules $A$ i $B$ són isomorfes si i només si són interdeduïbles:
\[
A\vdash_T B
\qquad\text{i}\qquad
B\vdash_T A.
\]
\end{exercici}

\begin{solucio}
Si $A\cong B$, per definició existeixen fletxes
\[
A\longrightarrow B
\qquad\text{i}\qquad
B\longrightarrow A.
\]
A $\cat{Prov}_T$ això significa exactament $A\vdash_T B$ i $B\vdash_T A$.

Recíprocament, si es donen aquestes dues deduïbilitats, existeixen les dues fletxes. Com que $\cat{Prov}_T$ és prima, qualsevol endomorfisme $A\to A$ és necessàriament l'única fletxa $\id_A$, i igualment per a $B$. Per tant les dues composicions són les identitats i les fletxes són inverses. Així $A\cong B$.
\end{solucio}

\section{Classes especials de morfismes}

\begin{exercici}
Demostra que un morfisme $f$ és un isomorfisme si i només si és un monomorfisme escindit (té retracció) i un epimorfisme.
\end{exercici}

\begin{solucio}
Si $f$ és un isomorfisme, el seu invers és alhora retracció i secció, de manera que $f$ és mono escindit; i tot isomorfisme és epi. Recíprocament, suposem que $f\colon A\to B$ té una retracció $r$ (amb $r\comp f=\id_A$) i que és epi. Partint de $r\comp f=\id_A$ i component amb $f$ per l'esquerra,
\[
(f\comp r)\comp f = f\comp(r\comp f) = f\comp\id_A = f = \id_B\comp f.
\]
Com que $f$ és epi, es pot cancel·lar el factor $f$ de la dreta, d'on
\[
f\comp r=\id_B.
\]
Juntament amb $r\comp f=\id_A$, això mostra que $r$ és també invers per la dreta de $f$: $f$ és un isomorfisme amb invers $r=f^{-1}$.
\end{solucio}

\begin{exercici}
A $\cat{Set}$, considera $f\colon\{0,1\}\to\{\ast\}$, l'única aplicació, i $g\colon\{\ast\}\to\{0,1\}$ amb $g(\ast)=0$. Comprova quina de les dues composicions $f\comp g$ i $g\comp f$ és una identitat, digues quin dels dos morfismes és una secció i quin una retracció (i de qui), i determina quin dels dos és un monomorfisme i quin un epimorfisme.
\end{exercici}

\begin{solucio}
Calculem les dues composicions. Per una banda, $f\comp g\colon\{\ast\}\to\{\ast\}$ compleix $(f\comp g)(\ast)=f(0)=\ast$, de manera que $f\comp g=\id_{\{\ast\}}$. Per l'altra, $g\comp f\colon\{0,1\}\to\{0,1\}$ envia tant $0$ com $1$ a $0$, de manera que $g\comp f\neq\id_{\{0,1\}}$ (per exemple, $(g\comp f)(1)=0\neq 1$).

Així, l'única igualtat que val és $f\comp g=\id_{\{\ast\}}$. Llegida des de cada morfisme, aquesta equació diu dues coses alhora: $g$ és una \emph{secció} de $f$ (un invers per la dreta de $f$), i $f$ és una \emph{retracció} de $g$ (un invers per l'esquerra de $g$). Convé insistir que \enquote{secció} i \enquote{retracció} són papers relatius a l'altre morfisme, no propietats absolutes: aquí $g$ no és \enquote{una secció} en abstracte, sinó una secció \emph{de $f$}.

Pel que fa a mono i epi: com que $f$ té una secció, $f$ és un epimorfisme; com a aplicació, és exhaustiva, d'acord amb la caracterització dels epimorfismes de $\cat{Set}$. Com que $g$ té una retracció, $g$ és un monomorfisme; com a aplicació, és injectiva. En canvi, $f$ no és mono, perquè l'aplicació $f$ no és injectiva ($f(0)=f(1)$), i $g$ no és epi, perquè l'aplicació $g$ no és exhaustiva ($1\notin g(\{\ast\})$). Cap dels dos és un isomorfisme: ho serien només si $g\comp f$ fos també una identitat, cosa que hem vist que no passa.
\end{solucio}

\begin{exercici}\label{exr:pr-prima-mono-epi}
Sigui $\cat{C}$ una categoria prima, és a dir, una categoria en què entre dos objectes hi ha com a màxim un morfisme. Comprova que tot morfisme de $\cat{C}$ és alhora un monomorfisme i un epimorfisme. Explica per què aquest fet mostra que, en una categoria arbitrària, \enquote{mono} i \enquote{epi} no es poden interpretar com \enquote{injectiu} i \enquote{exhaustiu}.
\end{exercici}

\begin{solucio}
Sigui $f\colon A\to B$ un morfisme de $\cat{C}$. Per veure que és un \textbf{monomorfisme}, siguin $g,h\colon X\to A$ amb $f\comp g=f\comp h$:
\[
\begin{tikzpicture}[baseline=(m.center)]
  \matrix (m) [matrix of math nodes, column sep=3.4em] { X & A & B \\ };
  \draw[-{Stealth[scale=1.0]}] ([yshift=2.6pt]m-1-1.east) -- node[above] {$g$} ([yshift=2.6pt]m-1-2.west);
  \draw[-{Stealth[scale=1.0]}] ([yshift=-2.6pt]m-1-1.east) -- node[below] {$h$} ([yshift=-2.6pt]m-1-2.west);
  \draw[-{Stealth[scale=1.1]}] (m-1-2) -- node[above] {$f$} (m-1-3);
\end{tikzpicture}
\qquad
f\comp g=f\comp h\ \Longrightarrow\ g=h.
\]
Els morfismes $g$ i $h$ tenen el mateix domini $X$ i el mateix codomini $A$: tots dos pertanyen a $\Hom(X,A)$. Com que la categoria és prima, aquest homconjunt té com a màxim un element, i per tant $g=h$. Observem que ni tan sols hem fet servir la hipòtesi $f\comp g=f\comp h$: la cancel·lació es compleix perquè no hi ha res a cancel·lar.

Per veure que és un \textbf{epimorfisme}, siguin $g,h\colon B\to Y$ amb $g\comp f=h\comp f$:
\[
\begin{tikzpicture}[baseline=(m.center)]
  \matrix (m) [matrix of math nodes, column sep=3.4em] { A & B & Y \\ };
  \draw[-{Stealth[scale=1.1]}] (m-1-1) -- node[above] {$f$} (m-1-2);
  \draw[-{Stealth[scale=1.0]}] ([yshift=2.6pt]m-1-2.east) -- node[above] {$g$} ([yshift=2.6pt]m-1-3.west);
  \draw[-{Stealth[scale=1.0]}] ([yshift=-2.6pt]m-1-2.east) -- node[below] {$h$} ([yshift=-2.6pt]m-1-3.west);
\end{tikzpicture}
\qquad
g\comp f=h\comp f\ \Longrightarrow\ g=h.
\]
Ara $g,h\in\Hom(B,Y)$, i el mateix argument dona $g=h$. També ho podem obtenir per dualitat: la categoria oposada $\cat{C}\op$ també és prima, perquè $\Hom_{\cat{C}\op}(B,A)$ està en bijecció amb $\Hom_{\cat{C}}(A,B)$; per tant $f\op$ és mono a $\cat{C}\op$, és a dir, $f$ és epi a $\cat{C}$ (teorema~\ref{teo:dualitat}).

\emph{Per què això impedeix llegir mono com a injectiu i epi com a exhaustiu.} Hi ha dues raons, una més radical que l'altra.

Primer, en una categoria prima els morfismes no són, en general, aplicacions. En un preordre, la fletxa $x\to y$ només registra que $x\le y$; no envia cap element enlloc, i preguntar si és \enquote{injectiva} no té sentit. Mono i epi, en canvi, sempre en tenen, perquè es defineixen només amb la composició.

Segon, fins i tot si volguéssim mantenir l'analogia, entraria en contradicció. A $\cat{Set}$, una aplicació injectiva i exhaustiva és bijectiva, és a dir, un isomorfisme. Si \enquote{mono $=$ injectiu} i \enquote{epi $=$ exhaustiu} valguessin en general, tot morfisme d'una categoria prima seria un isomorfisme. Però a $\cat{2}=(0\to 1)$ la fletxa $u\colon 0\to 1$ és mono i epi pel que acabem de demostrar, i no té invers, perquè no hi ha cap fletxa $1\to 0$. En un preordre qualsevol passa el mateix: la fletxa $x\to y$ és sempre mono i epi, i només és un isomorfisme quan també $y\le x$.

La conclusió és que \enquote{mono $=$ injectiu} i \enquote{epi $=$ exhaustiu} són teoremes sobre $\cat{Set}$, no definicions. Les definicions són les de cancel·lació, i en una categoria prima es compleixen de la manera més trivial possible.
\end{solucio}

\begin{exercici}\label{pr-mono-no-escindit}
Compara les aplicacions injectives $i\colon\varnothing\to\{\ast\}$ i $j\colon\{0\}\hookrightarrow\{0,1\}$ de $\cat{Set}$. Totes dues són monomorfismes? Quina d'elles és un monomorfisme escindit?
\end{exercici}

\begin{solucio}
Totes dues són injectives i, per tant, monomorfismes de $\cat{Set}$.

La inclusió $j$ té una retracció: l'aplicació $r\colon\{0,1\}\to\{0\}$ constant compleix $r\comp j=\id_{\{0\}}$. És, doncs, un monomorfisme escindit.

L'aplicació $i$, en canvi, no en té cap: una retracció seria una aplicació $\{\ast\}\to\varnothing$, i no n'hi ha cap, perquè $\ast$ no té on anar. Per tant, $i$ és un monomorfisme que no és escindit. L'exemple mostra que l'afirmació \enquote{tota aplicació injectiva té una inversa per l'esquerra} només és certa per a dominis no buits; i aquest contraexemple no depèn de l'axioma de l'elecció, a diferència del cas dual de les seccions.
\end{solucio}

\section{Construccions sobre categories}

\begin{exercici}
Descriu explícitament la categoria producte $\cat{2}\times\cat{3}$, on $\cat{2}$ és la categoria associada a l'ordre $0<1$ i $\cat{3}$ la categoria associada a $0<1<2$. Quants objectes i quants morfismes té? Com es componen?
\end{exercici}

\begin{solucio}
Els objectes són les parelles $(i,j)$ amb $i\in\{0,1\}$ i $j\in\{0,1,2\}$. N'hi ha sis i es poden disposar en una graella:
\[
\begin{tikzpicture}[baseline=(m.center)]
  \matrix (m) [matrix of math nodes, row sep=2.4em, column sep=2.8em] {
    (0,0) & (0,1) & (0,2) \\
    (1,0) & (1,1) & (1,2) \\
  };
  \draw[-{Stealth[scale=1.05]}] (m-1-1) -- (m-1-2);
  \draw[-{Stealth[scale=1.05]}] (m-1-2) -- (m-1-3);
  \draw[-{Stealth[scale=1.05]}] (m-2-1) -- (m-2-2);
  \draw[-{Stealth[scale=1.05]}] (m-2-2) -- (m-2-3);
  \draw[-{Stealth[scale=1.05]}] (m-1-1) -- (m-2-1);
  \draw[-{Stealth[scale=1.05]}] (m-1-2) -- (m-2-2);
  \draw[-{Stealth[scale=1.05]}] (m-1-3) -- (m-2-3);
\end{tikzpicture}
\]
Només hi hem dibuixat les fletxes bàsiques; les composicions produeixen també fletxes més llargues i diagonals.

Com que $\cat{2}$ i $\cat{3}$ són primes, entre $(i,j)$ i $(i',j')$ hi ha un únic morfisme exactament quan
\[
i\le i'
\qquad\text{i}\qquad
j\le j'.
\]
$\cat{2}$ té tres morfismes i $\cat{3}$ en té sis. Com que un morfisme del producte és una parella, hi ha
\[
3\cdot6=18
\]
morfismes. La composició es fa component a component. Tots els quadrats de la graella commuten perquè el producte és també una categoria prima.
\end{solucio}

\begin{exercici}
Descriu explícitament la categoria oposada $\cat{P}\op$ quan $\cat{P}$ és un preordre $(P,\leq)$ vist com a categoria. Quina relació d'ordre li correspon?
\end{exercici}

\begin{solucio}
A $\cat{P}$ hi ha un morfisme $x\to y$ exactament quan $x\leq y$. A $\cat{P}\op$ les fletxes s'inverteixen:
\[
x\longrightarrow y\text{ a }\cat{P}\op
\quad\Longleftrightarrow\quad
y\longrightarrow x\text{ a }\cat{P}
\quad\Longleftrightarrow\quad
y\leq x.
\]
Per tant $\cat{P}\op$ és el preordre invers $(P,\geq)$. En aquest exemple, invertir les fletxes és literalment invertir la desigualtat.
\end{solucio}

\begin{exercici}
Sigui $\cat{2}=(0\to 1)$ la categoria amb dos objectes i una fletxa no trivial. Prenent com a ambient la mateixa $\cat{2}$, descriu explícitament la categoria sobre l'objecte $1$, $\cat{2}/1$, i la categoria sota l'objecte $1$, $1/\cat{2}$, comptant-ne objectes i morfismes. Comprova després que $(\cat{2}/1)\op$ \emph{no} és isomorfa a $1/\cat{2}$, i identifica amb quina categoria sí que ho és.
\end{exercici}

\begin{solucio}
Escrivim la fletxa no trivial de $\cat{2}$ com $u\colon 0\to 1$.

\emph{La categoria $\cat{2}/1$} (sobre $1$). Els seus objectes són els morfismes de $\cat{2}$ amb codomini $1$: la fletxa $u\colon 0\to 1$ i la identitat $\id_1\colon 1\to 1$. Són, doncs, \emph{dos} objectes. Un morfisme de $u$ a $\id_1$ és una fletxa $h$ de $\cat{2}$ amb $\id_1\comp h=u$, és a dir $h=u$; i les identitats de cada objecte completen els morfismes. La categoria $\cat{2}/1$ té la forma d'una fletxa entre dos objectes: és isomorfa a $\cat{2}$.

\emph{La categoria $1/\cat{2}$} (sota $1$). Els seus objectes són els morfismes de $\cat{2}$ amb domini $1$: només $\id_1$, perquè de $1$ no en surt cap altra fletxa. És, doncs, \emph{un} sol objecte, amb només la seva identitat: la categoria $\cat{1}$.

\emph{El contrast.} $\cat{2}/1$ té dos objectes i $1/\cat{2}$ en té un. Per tant $(\cat{2}/1)\op$ té dos objectes i \emph{no} pot ser isomorfa a $1/\cat{2}$. El que sí que val és la fórmula general de dualitat
\[
(\cat{2}/1)\op\;\cong\;1/(\cat{2}\op).
\]
En efecte, $\cat{2}\op$ és la fletxa invertida $1\to 0$; sota l'objecte $1$ a $\cat{2}\op$ hi ha els morfismes amb domini $1$, que són $\id_1$ i la fletxa $1\to 0$: \emph{dos} objectes, com $(\cat{2}/1)\op$. La moral: en dualitzar la construcció \enquote{sobre} en \enquote{sota} cal dualitzar també la categoria ambient, no deixar-la intacta.
\end{solucio}

\section{Principi de dualitat}

\begin{exercici}\label{exr:pr-dual-mono-epi}
Demostra que la composició de dos monomorfismes és un monomorfisme. A continuació, sense repetir el càlcul, dedueix per dualitat que la composició de dos epimorfismes és un epimorfisme. Indica exactament què s'inverteix en passar a $\cat{C}\op$ i què es manté.
\end{exercici}

\begin{solucio}
\emph{Primera part: el cas dels monomorfismes.} Siguin $f\colon A\to B$ i $g\colon B\to C$ monomorfismes, i siguin $u,v\colon X\to A$ amb $(g\comp f)\comp u=(g\comp f)\comp v$:
\[
\begin{tikzpicture}[baseline=(m.center)]
  \matrix (m) [matrix of math nodes, column sep=3em] { X & A & B & C \\ };
  \draw[-{Stealth[scale=1.0]}] ([yshift=2.6pt]m-1-1.east) -- node[above] {$u$} ([yshift=2.6pt]m-1-2.west);
  \draw[-{Stealth[scale=1.0]}] ([yshift=-2.6pt]m-1-1.east) -- node[below] {$v$} ([yshift=-2.6pt]m-1-2.west);
  \draw[-{Stealth[scale=1.1]}] (m-1-2) -- node[above] {$f$} (m-1-3);
  \draw[-{Stealth[scale=1.1]}] (m-1-3) -- node[above] {$g$} (m-1-4);
\end{tikzpicture}
\qquad
(g\comp f)\comp u=(g\comp f)\comp v\ \Longrightarrow\ u=v.
\]
Per associativitat, la hipòtesi és $g\comp(f\comp u)=g\comp(f\comp v)$. Com que $g$ és mono, en cancel·lem el factor de l'esquerra i obtenim $f\comp u=f\comp v$. Com que $f$ és mono, tornem a cancel·lar i obtenim $u=v$. Per tant $g\comp f$ és mono. Observem l'ordre: primer es cancel·la $g$, el factor \emph{exterior}, i després $f$.

\emph{Segona part: dualització.} Procedim en tres passos.

\textbf{1. L'enunciat és categòric.} \enquote{Si $f\colon A\to B$ i $g\colon B\to C$ són monomorfismes, aleshores $g\comp f$ és un monomorfisme} només parla d'objectes, morfismes, composició i igualtats entre composicions (la definició de mono és d'aquesta mena). Per tant té dual, i com que l'hem demostrat en una categoria qualsevol, el principi de dualitat (teorema~\ref{teo:dualitat}) garanteix que el dual també és vàlid en tota categoria.

\textbf{2. Què s'inverteix.} Recorrem l'enunciat element per element:
\begin{center}
\begin{tabular}{@{}ll@{}}
\toprule
A $\cat{C}$ & A l'enunciat dual \\
\midrule
$f\colon A\to B$, \quad $g\colon B\to C$ & $f\colon B\to A$, \quad $g\colon C\to B$ \\
la composta $g\comp f\colon A\to C$ & la composta $f\comp g\colon C\to A$ \\
fletxes de prova $u,v\colon X\to A$ & fletxes de prova $u,v\colon A\to X$ \\
mono: $m\comp u=m\comp v\Rightarrow u=v$ & epi: $u\comp m=v\comp m\Rightarrow u=v$ \\
\bottomrule
\end{tabular}
\end{center}
Es mantenen els quantificadors (\enquote{per a tot $u,v$}), la implicació, la igualtat i l'associativitat, que és autodual. L'enunciat dual diu, doncs: si $f\colon B\to A$ i $g\colon C\to B$ són epimorfismes, aleshores $f\comp g\colon C\to A$ és un epimorfisme. Rebatejant les fletxes perquè la cadena es llegeixi d'esquerra a dreta, $p\colon A\to B$ i $q\colon B\to C$, això és: \emph{si $p$ i $q$ són epimorfismes, $q\comp p$ és un epimorfisme.}

\textbf{3. Per què és vàlid.} És l'argument del principi, fet explícit. Siguin $p\colon A\to B$ i $q\colon B\to C$ epimorfismes de $\cat{C}$. A $\cat{C}\op$, les fletxes $p\op\colon B\to A$ i $q\op\colon C\to B$ són monomorfismes, perquè ser epi a $\cat{C}$ és ser mono a $\cat{C}\op$. Per la primera part, aplicada a $\cat{C}\op$, la composta $p\op\comp q\op\colon C\to A$ és mono a $\cat{C}\op$. Però $p\op\comp q\op=(q\comp p)\op$, de manera que $(q\comp p)\op$ és mono a $\cat{C}\op$, és a dir, $q\comp p$ és epi a $\cat{C}$.

\emph{Comprovació: la prova dual.} No l'hem necessitat, però val la pena escriure-la per veure que és la primera llegida a l'inrevés. Siguin $u,v\colon C\to X$ amb $u\comp(q\comp p)=v\comp(q\comp p)$:
\[
\begin{tikzpicture}[baseline=(m.center)]
  \matrix (m) [matrix of math nodes, column sep=3em] { A & B & C & X \\ };
  \draw[-{Stealth[scale=1.1]}] (m-1-1) -- node[above] {$p$} (m-1-2);
  \draw[-{Stealth[scale=1.1]}] (m-1-2) -- node[above] {$q$} (m-1-3);
  \draw[-{Stealth[scale=1.0]}] ([yshift=2.6pt]m-1-3.east) -- node[above] {$u$} ([yshift=2.6pt]m-1-4.west);
  \draw[-{Stealth[scale=1.0]}] ([yshift=-2.6pt]m-1-3.east) -- node[below] {$v$} ([yshift=-2.6pt]m-1-4.west);
\end{tikzpicture}
\qquad
u\comp(q\comp p)=v\comp(q\comp p)\ \Longrightarrow\ u=v.
\]
Per associativitat, $(u\comp q)\comp p=(v\comp q)\comp p$; com que $p$ és epi, $u\comp q=v\comp q$; com que $q$ és epi, $u=v$. Les mateixes passes, amb cada composició en l'ordre invers: ara es cancel·la primer $p$, el factor de la \emph{dreta}, que és la fletxa per on comença el camí.

\emph{Error típic.} En dualitzar és fàcil invertir les fletxes però oblidar d'invertir l'ordre de la composició, i escriure \enquote{$p\comp q$ és epi}. Si $p\colon A\to B$ i $q\colon B\to C$, aquesta composta ni tan sols està definida, tret que $C=A$. Comprovar els dominis i els codominis de cada composta és la millor manera de detectar l'error.
\end{solucio}

\section{Grafs i categories lliures}

\begin{exercici}
Considera el graf amb un sol vèrtex i dues arestes que són bucles, $a$ i $b$. Descriu la categoria lliure que genera. Quants morfismes té?
\end{exercici}

\begin{solucio}
La categoria lliure té un sol objecte. Els morfismes són tots els camins finits:
\[
\varepsilon,
\quad a,
\quad b,
\quad aa,
\quad ab,
\quad ba,
\quad bb,
\quad aaa,
\quad\ldots
\]
on $\varepsilon$ és el camí buit i actua com a identitat. La composició és la concatenació de paraules en ordre de recorregut: $b\comp a=ab$. Com que cap relació no identifica camins diferents, hi ha una infinitat numerable de morfismes. Llevat de l'isomorfisme canònic que inverteix l'ordre de les lletres (exemple~\ref{ex:monoide-lliure-cat}), és el monoide lliure sobre l'alfabet $\{a,b\}$.
\end{solucio}

\section{Qüestions de mida}

\begin{exercici}
En una categoria petita, l'oposada $\cat{C}\op$ també és petita. Comprova-ho, i comprova a més que $\cat{C}$ i $\cat{C}\op$ tenen exactament el mateix nombre de morfismes.
\end{exercici}

\begin{solucio}
Els objectes de $\cat{C}\op$ són els mateixos que els de $\cat{C}$. Cada morfisme $f\colon A\to B$ de $\cat{C}$ correspon al mateix morfisme vist com $f\colon B\to A$ a $\cat{C}\op$. Això dona una bijecció entre les col·leccions de morfismes de totes dues categories.

Per tant, si els objectes i els morfismes de $\cat{C}$ formen conjunts, també ho fan els de $\cat{C}\op$. En particular, $\cat{C}$ és petita si i només si ho és $\cat{C}\op$.
\end{solucio}

\begin{exercici}\label{exr:pr-set-localment-petita}
Comprova que $\cat{Set}$ és localment petita: per a conjunts $A$ i $B$, mostra que la col·lecció de les aplicacions $A\to B$ és un conjunt, identificant cada aplicació amb el seu graf, que és un subconjunt d'$A\times B$.
\end{exercici}

\begin{solucio}
Una aplicació $f\colon A\to B$ queda determinada pel seu \textbf{graf}
\[
\Gamma_f=\{(a,f(a))\mid a\in A\}\subseteq A\times B.
\]
Un subconjunt $\Gamma\subseteq A\times B$ és el graf d'alguna aplicació $A\to B$ exactament quan és \emph{funcional}: per a cada $a\in A$ hi ha un únic $b\in B$ amb $(a,b)\in\Gamma$. Així, $f\mapsto\Gamma_f$ és una bijecció entre les aplicacions $A\to B$ i els subconjunts funcionals d'$A\times B$; la inversa envia $\Gamma$ a l'aplicació que assigna a cada $a$ l'únic $b$ amb $(a,b)\in\Gamma$.

Vegem ara que aquests subconjunts funcionals formen un conjunt. Com que $A$ i $B$ són conjunts, el producte $A\times B$ és un conjunt, i també ho és el seu conjunt de parts $\mathcal P(A\times B)$. La condició \enquote{ser funcional} és una propietat ben definida dels elements de $\mathcal P(A\times B)$, de manera que, per l'axioma de separació,
\[
\{\Gamma\in\mathcal P(A\times B)\mid \forall a\in A\ \exists!\,b\in B\ (a,b)\in\Gamma\}
\]
és un conjunt. Per la bijecció anterior, $\Hom_{\cat{Set}}(A,B)$ és un conjunt.

Una precisió, la mateixa de l'observació~\ref{obs:extrems-implicits} de la Teoria per a $\cat{Rel}$: el graf no determina el codomini, perquè $\Gamma_f\subseteq A\times B$ també és un subconjunt d'$A\times B'$ per a tot $B'\supseteq B$. Per això, en rigor, un morfisme de $\cat{Set}$ és la terna $(A,B,\Gamma_f)$. Això no canvia la conclusió: les ternes $(A,B,\Gamma)$ amb $\Gamma$ funcional formen el conjunt $\{A\}\times\{B\}\times\{\Gamma\in\mathcal P(A\times B)\mid \Gamma\text{ funcional}\}$, que està en bijecció amb l'anterior.

\emph{Abast.} El mateix argument explica per què $\Hom(A,B)$ és un conjunt en els exemples habituals. A $\cat{Rel}$, directament, $\Hom_{\cat{Rel}}(X,Y)$ és, llevat dels extrems, $\mathcal P(X\times Y)$. A les categories de conjunts estructurats ($\cat{Grp}$, $\cat{Top}$, $\cat{Vect}_{\mathbb K}$, $\cat{Mon}$, $\cat{Pos}$), els morfismes són aplicacions entre els conjunts subjacents que compleixen una condició addicional; formen, doncs, un subconjunt de l'homconjunt de $\cat{Set}$ corresponent, de nou per separació. En canvi, la col·lecció de \emph{tots} els objectes de $\cat{Set}$ no és un conjunt (paradoxa de Russell): $\cat{Set}$ és localment petita sense ser petita.
\end{solucio}

\chapter{Functors}\label{cap:practica-functors}

En aquest capítol resolem amb detall una selecció d'exercicis sobre \emph{functors}. Com al capítol anterior, l'objectiu no és només arribar al resultat, sinó mostrar com es comprova que una assignació és un functor, com es reconeixen els functors coneguts sota formes noves i com es distingeixen les propietats que un functor preserva de les que no. Convé intentar cada exercici abans de llegir-ne la solució.

\section{Definició de functor}

\begin{exercici}\label{prac:functor-quadrat}
Sigui $F\colon\cat{C}\to\cat{D}$ un functor i suposem que a $\cat{C}$ tenim un quadrat commutatiu, és a dir, quatre morfismes $f\colon A\to B$, $u\colon A\to C$, $v\colon B\to D$ i $g\colon C\to D$ amb
\[
v\comp f=g\comp u.
\]
Mostra que la seva imatge per $F$ torna a ser un quadrat commutatiu:
\[
F(v)\comp F(f)=F(g)\comp F(u).
\]
\end{exercici}

\begin{solucio}
El quadrat de partida i la seva imatge són:
\[
\begin{tikzpicture}[baseline=(m.center)]
  \matrix (m) [matrix of math nodes, row sep=2.6em, column sep=3em] {
    A & B \\
    C & D \\
  };
  \draw[-{Stealth[scale=1.1]}] (m-1-1) -- node[above] {$f$} (m-1-2);
  \draw[-{Stealth[scale=1.1]}] (m-1-1) -- node[left] {$u$} (m-2-1);
  \draw[-{Stealth[scale=1.1]}] (m-1-2) -- node[right] {$v$} (m-2-2);
  \draw[-{Stealth[scale=1.1]}] (m-2-1) -- node[below] {$g$} (m-2-2);
\end{tikzpicture}
\qquad\stackrel{F}{\longmapsto}\qquad
\begin{tikzpicture}[baseline=(m.center)]
  \matrix (m) [matrix of math nodes, row sep=2.6em, column sep=3.4em] {
    F(A) & F(B) \\
    F(C) & F(D) \\
  };
  \draw[-{Stealth[scale=1.1]}] (m-1-1) -- node[above] {$F(f)$} (m-1-2);
  \draw[-{Stealth[scale=1.1]}] (m-1-1) -- node[left] {$F(u)$} (m-2-1);
  \draw[-{Stealth[scale=1.1]}] (m-1-2) -- node[right] {$F(v)$} (m-2-2);
  \draw[-{Stealth[scale=1.1]}] (m-2-1) -- node[below] {$F(g)$} (m-2-2);
\end{tikzpicture}
\]
Partim de la hipòtesi $v\comp f=g\comp u$. Apliquem $F$ a aquesta igualtat; com que $F$ és una aplicació ben definida sobre els morfismes, morfismes iguals tenen imatges iguals:
\[
F(v\comp f)=F(g\comp u).
\]
Ara fem servir que $F$ preserva la composició a cada costat:
\[
F(v)\comp F(f)=F(v\comp f)=F(g\comp u)=F(g)\comp F(u).
\]
Per tant el quadrat imatge commuta. Aquest és el contingut de la idea que \emph{els functors envien diagrames commutatius a diagrames commutatius}: n'hi ha prou que ho comprovem sobre les composicions, perquè un diagrama commuta exactament quan certes composicions coincideixen.
\end{solucio}

\begin{exercici}\label{pr-parts-cov}
Comprova que l'assignació que envia cada conjunt $X$ al conjunt de parts $\mathcal P(X)$ i cada aplicació $f\colon X\to Y$ a la \textbf{imatge directa}
\[
\mathcal P_!(f)\colon\mathcal P(X)\to\mathcal P(Y),
\qquad
A\mapsto f(A)=\{f(a)\mid a\in A\},
\]
és un functor covariant
\[
\mathcal P_!\colon\cat{Set}\to\cat{Set},
\]
el \textbf{functor de parts per imatge directa}. (Reservem $\mathcal P$, sense subíndex, per al functor de parts contravariant de la Teoria, que actua per antiimatge.)
\end{exercici}

\begin{solucio}
Hem de comprovar la preservació d'identitats i de composició.

Per a la \textbf{identitat}, si $f=\id_X$ aleshores per a tot $A\subseteq X$,
\[
\mathcal P_!(\id_X)(A)=\{\id_X(a)\mid a\in A\}=A,
\]
de manera que $\mathcal P_!(\id_X)=\id_{\mathcal P(X)}$.

Per a la \textbf{composició}, siguin $f\colon X\to Y$ i $g\colon Y\to Z$. Per a tot $A\subseteq X$,
\[
\begin{aligned}
\mathcal P_!(g\comp f)(A)&=(g\comp f)(A)=g\big(f(A)\big)\\
&=\mathcal P_!(g)\big(\mathcal P_!(f)(A)\big)=\big(\mathcal P_!(g)\comp\mathcal P_!(f)\big)(A).
\end{aligned}
\]
La igualtat central, $(g\comp f)(A)=g(f(A))$, és la propietat coneguda de la imatge directa. Per tant $\mathcal P_!$ és un functor covariant. Té els mateixos objectes que $\mathcal P$, però actua sobre els morfismes en sentit contrari i per una operació diferent: són dos functors distints.
\end{solucio}

\begin{exercici}\label{pr-hom-cov}
Sigui $\cat{C}$ una categoria localment petita i $A$ un objecte fix. Comprova que l'assignació que envia cada objecte $X$ al conjunt $\Hom(A,X)$ i cada morfisme $f\colon X\to Y$ a l'aplicació
\[
f_*\colon\Hom(A,X)\to\Hom(A,Y),
\qquad
\varphi\mapsto f\comp\varphi,
\]
és un functor covariant $\Hom(A,-)\colon\cat{C}\to\cat{Set}$ (suposant $\cat{C}$ localment petita).
\end{exercici}

\begin{solucio}
Primer cal veure que $f_*$ està ben definida: si $\varphi\colon A\to X$, aleshores $f\comp\varphi\colon A\to Y$, de manera que $f_*(\varphi)\in\Hom(A,Y)$.

Per a la \textbf{identitat}, prenent $f=\id_X$,
\[
(\id_X)_*(\varphi)=\id_X\comp\varphi=\varphi,
\]
de manera que $(\id_X)_*=\id_{\Hom(A,X)}$.

Per a la \textbf{composició}, siguin $f\colon X\to Y$ i $g\colon Y\to Z$. Per a tot $\varphi\in\Hom(A,X)$, usant l'associativitat de $\cat{C}$,
\[
(g\comp f)_*(\varphi)=(g\comp f)\comp\varphi=g\comp(f\comp\varphi)=g_*\big(f_*(\varphi)\big).
\]
Per tant $(g\comp f)_*=g_*\comp f_*$ i $\Hom(A,-)$ és un functor covariant.
\end{solucio}

\section{Exemples de functors}

\begin{exercici}
Siguin $T$ i $T'$ dos sistemes deductius sobre el mateix conjunt de fórmules, amb $T'$ una \textbf{extensió} de $T$ ---és a dir, tot el que és deduïble a $T$ també ho és a $T'$: si $A\vdash_T B$, aleshores $A\vdash_{T'} B$. Comprova que la identitat sobre les fórmules indueix un functor
\[
J\colon\cat{Prov}_T\to\cat{Prov}_{T'},
\]
i explica en quins termes es podria formular que $J$ sigui ple.
\end{exercici}

\begin{solucio}
Recordem que $\cat{Prov}_T$ té les fórmules per objectes i una única fletxa $A\to B$ exactament quan $A\vdash_T B$; és una categoria prima. Definim $J$ com la identitat sobre els objectes, $J(A)=A$.

Sobre els morfismes, una fletxa $A\to B$ de $\cat{Prov}_T$ existeix exactament quan $A\vdash_T B$; la fletxa registra, doncs, el fet de la deduïbilitat, no una prova concreta. Com que $T'$ és una extensió de $T$, tenim $A\vdash_{T'} B$, de manera que hi ha una fletxa $J(A)\to J(B)$ a $\cat{Prov}_{T'}$, i és única perquè $\cat{Prov}_{T'}$ és prima. Això defineix $J$ sobre morfismes.

Les condicions de functor se satisfan automàticament perquè totes dues categories són primes: entre dos objectes hi ha com a màxim una fletxa, així que la preservació d'identitats i de composicions no imposa res més enllà de l'existència de les fletxes imatge, que ja hem comprovat.

Quant a \textbf{ser ple}: $J$ seria ple si, per a cada parella de fórmules $A,B$, tota fletxa $J(A)\to J(B)$ a $\cat{Prov}_{T'}$ provingués d'una fletxa $A\to B$ a $\cat{Prov}_T$. És a dir, si
\[
A\vdash_{T'} B\quad\Longrightarrow\quad A\vdash_T B.
\]
Dit d'una altra manera, $J$ és ple exactament quan l'extensió $T'$ \emph{no afegeix cap conseqüència nova} entre les fórmules considerades. Quan $T'$ demostra estrictament més relacions que $T$, el functor $J$ és fidel ---de fet, totes dues categories són primes, així que $J$ és sempre fidel--- però no ple.
\end{solucio}

\begin{exercici}
Comprova amb detall que un functor entre dos preordres $(P,\leq)$ i $(Q,\leq)$, vistos com a categories, és exactament una aplicació monòtona.
\end{exercici}

\begin{solucio}
Un functor $F\colon P\to Q$ dona, en primer lloc, una aplicació $x\mapsto F(x)$ dels objectes, és a dir, dels elements de $P$ als de $Q$.

La condició sobre morfismes diu que si hi ha una fletxa $x\to y$ a $P$, aleshores hi ha una fletxa $F(x)\to F(y)$ a $Q$. Com que en un preordre hi ha una fletxa $x\to y$ exactament quan $x\leq y$, això equival a
\[
x\leq y
\quad\Longrightarrow\quad
F(x)\leq F(y),
\]
que és precisament la definició d'aplicació monòtona.

Recíprocament, tota aplicació monòtona defineix un functor: envia l'única fletxa $x\to y$ (quan $x\leq y$) a l'única fletxa $F(x)\to F(y)$. Les condicions de preservació d'identitats i composició es compleixen automàticament, perquè entre dos objectes hi ha com a màxim una fletxa i totes les igualtats entre morfismes són trivials. Així, els functors $P\to Q$ i les aplicacions monòtones $P\to Q$ són el mateix.
\end{solucio}

\begin{exercici}
Sigui $M$ un monoide vist com a categoria $\cat{B}M$ amb un sol objecte $\ast$. Comprova que un functor $F\colon\cat{B}M\to\cat{Set}$ és el mateix que una acció de $M$ sobre un conjunt, és a dir, un conjunt $S$ juntament amb una operació
\[
M\times S\to S,
\qquad
(m,s)\mapsto m\cdot s,
\]
tal que $1_M\cdot s=s$ i $(mn)\cdot s=m\cdot(n\cdot s)$.
\end{exercici}
\begin{solucio}
Un functor $F\colon\cat{B}M\to\cat{Set}$ tria un conjunt $S=F(\ast)$ i envia cada element $m\in M$ ---cada morfisme $\ast\to\ast$--- a una aplicació $F(m)\colon S\to S$. Definim l'acció per
\[
m\cdot s=F(m)(s).
\]

La preservació de la identitat diu $F(1_M)=\id_S$, és a dir, $1_M\cdot s=s$ per a tot $s$. La preservació de la composició diu $F(mn)=F(m)\comp F(n)$, és a dir,
\[
(mn)\cdot s=F(mn)(s)=F(m)\big(F(n)(s)\big)=m\cdot(n\cdot s).
\]
Aquestes són exactament les dues lleis d'una acció per l'esquerra. Recíprocament, tota acció defineix, per la mateixa fórmula, un functor. Per tant els functors $\cat{B}M\to\cat{Set}$ i les accions de $M$ coincideixen.
\end{solucio}

\begin{exercici}
Comprova que la construcció de la categoria lliure és functorial sobre les fletxes: donat un morfisme de grafs $\alpha\colon G\to H$, l'assignació que envia cada camí d'arestes de $G$ a la seqüència de les seves imatges és un functor $\Free(\alpha)\colon\Free(G)\to\Free(H)$.
\end{exercici}

\begin{solucio}
Un camí a $\Free(G)$ és una seqüència finita d'arestes encadenades $a_1a_2\cdots a_n$. Definim
\[
\Free(\alpha)(a_1a_2\cdots a_n)=\alpha_E(a_1)\,\alpha_E(a_2)\cdots\alpha_E(a_n),
\]
i sobre els objectes $\Free(\alpha)=\alpha_V$. Cal comprovar que aquesta seqüència és realment un camí a $H$: com que $\alpha$ preserva origen i destí, si $a_i$ acaba on comença $a_{i+1}$, aleshores $\alpha_E(a_i)$ acaba on comença $\alpha_E(a_{i+1})$, de manera que les imatges també s'encadenen.

La \textbf{identitat} de cada vèrtex és el camí buit, que s'envia al camí buit del vèrtex imatge: $\Free(\alpha)$ preserva identitats. La \textbf{composició} és la concatenació. Siguin $p=(a_1,\dots,a_r)$ i $q=(b_1,\dots,b_s)$ dos camins componibles, amb el final de $p$ igual a l'inici de $q$; la seva composició és el camí que recorre primer $p$ i després $q$,
\[
q\comp p=(a_1,\dots,a_r,b_1,\dots,b_s).
\]
Aplicar $\alpha_E$ terme a terme i després concatenar dona el mateix que concatenar i després aplicar $\alpha_E$ terme a terme, de manera que
\[
\Free(\alpha)(q\comp p)=\Free(\alpha)(q)\comp\Free(\alpha)(p).
\]
Per tant $\Free(\alpha)$ és un functor.
\end{solucio}

\section{Composició de functors i la categoria \texorpdfstring{$\cat{Cat}$}{Cat}}

\begin{exercici}
Comprova que la composició de dos functors $F\colon\cat{C}\to\cat{D}$ i $G\colon\cat{D}\to\cat{E}$, seguida d'un tercer $H\colon\cat{E}\to\cat{F}$, és associativa:
\[
H\comp(G\comp F)=(H\comp G)\comp F.
\]
\end{exercici}

\begin{solucio}
Dues igualtats de functors es comproven sobre objectes i sobre morfismes. Sobre un objecte $A$ de $\cat{C}$,
\[
\big(H\comp(G\comp F)\big)(A)=H\big(G(F(A))\big)=\big((H\comp G)\comp F\big)(A),
\]
i sobre un morfisme $f$,
\[
\big(H\comp(G\comp F)\big)(f)=H\big(G(F(f))\big)=\big((H\comp G)\comp F\big)(f).
\]
En tots dos casos la igualtat es redueix a l'associativitat de l'aplicació successiva, primer $F$, després $G$, després $H$. Per tant la composició de functors és associativa, i amb els functors identitat com a neutres, això confirma que $\cat{Cat}$ és una categoria.
\end{solucio}

\begin{exercici}
Fixat un objecte $A$ d'una categoria $\cat{C}$, descriu el functor oblit $U_A\colon\cat{C}/A\to\cat{C}$ de la categoria de $\cat{C}$ sobre $A$ cap a $\cat{C}$, i comprova que és un functor.
\end{exercici}

\begin{solucio}
Recordem que un objecte de $\cat{C}/A$ és un morfisme $x\colon X\to A$ de $\cat{C}$, i un morfisme de $x\colon X\to A$ cap a $y\colon Y\to A$ és un morfisme $h\colon X\to Y$ de $\cat{C}$ tal que $y\comp h=x$ (el triangle commuta). Definim $U_A$ així: sobre objectes, envia $x\colon X\to A$ al seu domini $X$; sobre morfismes, envia el morfisme $h$ del triangle al mateix $h\colon X\to Y$ vist a $\cat{C}$, oblidant la condició $y\comp h=x$.

Comprovem que és un functor. \textbf{Identitats:} la identitat de l'objecte $x\colon X\to A$ a $\cat{C}/A$ és $\id_X$ (que compleix $x\comp\id_X=x$), i $U_A(\id_X)=\id_X=\id_{U_A(x)}$. \textbf{Composició:} si $h\colon x\to y$ i $k\colon y\to z$ són morfismes componibles a $\cat{C}/A$, la seva composició és $k\comp h$ com a morfisme de $\cat{C}$; per tant $U_A(k\comp h)=k\comp h=U_A(k)\comp U_A(h)$. Així $U_A$ preserva identitats i composició, i és un functor. Intuïtivament, $U_A$ \enquote{oblida} el morfisme cap a $A$ i reté només el domini i les fletxes entre dominis.
\end{solucio}

\section{Functors i isomorfismes}

\begin{exercici}
Dona un exemple de functor i d'un morfisme \emph{no} isomorf la imatge del qual \emph{sí} que és un isomorfisme. Això mostra que els functors preserven isomorfismes, però no els reflecteixen en general.
\end{exercici}

\begin{solucio}
Considerem el functor oblit $U\colon\cat{Pos}\to\cat{Set}$, que envia cada conjunt parcialment ordenat al seu conjunt subjacent i cada aplicació monòtona a ella mateixa. Siguin $P$ el conjunt $\{0,1\}$ amb l'ordre discret (cap parella d'elements diferents no és comparable) i $Q$ el mateix conjunt amb l'ordre $0<1$, i sigui
\[
f\colon P\to Q
\]
la identitat de conjunts. És monòtona, perquè a $P$ les úniques relacions són $0\le 0$ i $1\le 1$, que es conserven a $Q$. La inversa $Q\to P$, en canvi, no ho és: a $Q$ tenim $0\le 1$, però a $P$ no. Per tant $f$ \textbf{no és un isomorfisme} a $\cat{Pos}$, perquè no té cap invers dins de la categoria.

Ara bé, la seva imatge $U(f)$ és l'aplicació identitat del conjunt $\{0,1\}$, que \textbf{sí que és un isomorfisme} a $\cat{Set}$.

\emph{Un exemple anàleg, per a qui conegui la topologia.} El functor oblit $U\colon\cat{Top}\to\cat{Set}$ es comporta igual. Si $X$ és un conjunt amb dues topologies $\tau$ i $\tau'$, amb $\tau$ estrictament més fina que $\tau'$, la identitat $(X,\tau)\to(X,\tau')$ és contínua, però la seva inversa no ho és; no és, doncs, un isomorfisme a $\cat{Top}$, mentre que la seva imatge a $\cat{Set}$ és la identitat de $X$.

En tots dos casos, $U$ no reflecteix isomorfismes: la seva imatge pot ser invertible sense que l'original ho sigui. Els functors plens i fidels sí que reflecteixen isomorfismes, encara que aquesta condició no és necessària: hi ha functors que els reflecteixen sense ser plens. Per exemple, el functor oblit $\cat{Grp}\to\cat{Set}$ reflecteix isomorfismes ---un homomorfisme bijectiu de grups és un isomorfisme---, malgrat no ser ple.
\end{solucio}

\begin{exercici}\label{pr-2-set-mono-epi}
A la categoria $\cat{2}=(0\xrightarrow{u}1)$, la fletxa $u$ és mono i epi. Construeix un functor $\cat{2}\to\cat{Set}$ que no l'enviï a un monomorfisme, i un altre que no l'enviï a un epimorfisme.
\end{exercici}

\begin{solucio}
Un functor $\cat{2}\to\cat{Set}$ queda determinat per dos conjunts i una aplicació entre ells, imatges de $0$, $1$ i $u$: les úniques composicions de $\cat{2}$ involucren identitats, de manera que la functorialitat només demana enviar identitats a identitats.
\begin{itemize}
  \item $F(0)=\{0,1\}$, $F(1)=\{\ast\}$ i $F(u)$ l'única aplicació $\{0,1\}\to\{\ast\}$. No és injectiva, i per tant $F(u)$ no és mono.
  \item $G(0)=\{\ast\}$, $G(1)=\{0,1\}$ i $G(u)$ l'aplicació de valor $0$. No és exhaustiva, i per tant $G(u)$ no és epi.
\end{itemize}
Un functor, doncs, no preserva en general ni els monomorfismes ni els epimorfismes. A $\cat{2}$, que $u$ sigui mono i epi és conseqüència que la categoria és prima; a $\cat{Set}$ hi ha més fletxes amb què comparar, i aquesta trivialitat es perd.
\end{solucio}

\section{Functors covariants i contravariants}

\begin{exercici}\label{pr-parts-contra}
Comprova que el functor de parts contravariant $\mathcal P\colon\cat{Set}\op\to\cat{Set}$, que envia $f\colon X\to Y$ a l'antiimatge
\[
f^{-1}\colon\mathcal P(Y)\to\mathcal P(X),
\qquad
B\mapsto f^{-1}(B),
\]
és realment un prefeix sobre $\cat{Set}$ (definició~\ref{def:prefeix}).
\end{exercici}

\begin{solucio}
Cal comprovar la preservació d'identitats i la inversió de l'ordre en la composició.

Per a la \textbf{identitat}, $(\id_X)^{-1}(B)=B$ per a tot $B\subseteq X$, de manera que $(\id_X)^{-1}=\id_{\mathcal P(X)}$.

Per a la \textbf{composició}, siguin $f\colon X\to Y$ i $g\colon Y\to Z$. Per a tot $C\subseteq Z$,
\[
(g\comp f)^{-1}(C)=\{x\mid g(f(x))\in C\}=\{x\mid f(x)\in g^{-1}(C)\}=f^{-1}\big(g^{-1}(C)\big).
\]
Per tant
\[
(g\comp f)^{-1}=f^{-1}\comp g^{-1},
\]
amb l'ordre invertit. Aquesta és exactament la condició de prefeix, és a dir, de functor $\cat{Set}\op\to\cat{Set}$.
\end{solucio}

\begin{exercici}\label{pr-hom-contra}
Sigui $B$ un objecte fix d'una categoria localment petita $\cat{C}$. Comprova que $\Hom(-,B)$ és un prefeix sobre $\cat{C}$, és a dir, un functor $\Hom(-,B)\colon\cat{C}\op\to\cat{Set}$.
\end{exercici}

\begin{solucio}
A cada objecte $X$ li assignem el conjunt $\Hom(X,B)$. A cada morfisme $f\colon X\to Y$ li assignem l'aplicació de \textbf{precomposició}
\[
f^*\colon\Hom(Y,B)\to\Hom(X,B),
\qquad
\psi\mapsto\psi\comp f,
\]
que està ben definida perquè si $\psi\colon Y\to B$ aleshores $\psi\comp f\colon X\to B$. Observem que $f^*$ inverteix el sentit: parteix dels morfismes des de $Y$ i arriba als morfismes des de $X$.

Per a la \textbf{identitat}, $(\id_X)^*(\psi)=\psi\comp\id_X=\psi$. Per a la \textbf{composició}, siguin $f\colon X\to Y$ i $g\colon Y\to Z$; per a tot $\psi\in\Hom(Z,B)$,
\[
(g\comp f)^*(\psi)=\psi\comp(g\comp f)=(\psi\comp g)\comp f=f^*\big(g^*(\psi)\big).
\]
Per tant $(g\comp f)^*=f^*\comp g^*$, amb l'ordre invertit, i $\Hom(-,B)$ és contravariant.
\end{solucio}

\begin{exercici}\label{pr-bihom}
Sigui $\cat{C}$ localment petita. Comprova que el bifunctor
\[
\Hom(-,-)\colon\cat{C}\op\times\cat{C}\to\cat{Set},\qquad \Hom(f,g)(h)=g\comp h\comp f,
\]
respecta la composició simultània en les dues variables i les identitats. És a dir, comprova que és realment un functor, no només en cada variable per separat.
\end{exercici}

\begin{solucio}
Siguin $f\colon A'\to A$, $f'\colon A''\to A'$, $g\colon B\to B'$ i $g'\colon B'\to B''$ morfismes de $\cat{C}$. Recordem que la composició a $\cat{C}\op\times\cat{C}$ actua component a component, i que a la primera component (la de $\cat{C}\op$) l'ordre és l'invers del de $\cat{C}$. \textbf{Composició.} Per a tot $h\colon A\to B$,
\[
\begin{aligned}
\Hom(f',g')\big(\Hom(f,g)(h)\big)
&=g'\comp(g\comp h\comp f)\comp f'\\
&=(g'\comp g)\comp h\comp(f\comp f')\\
&=\Hom(f\comp f',\,g'\comp g)(h),
\end{aligned}
\]
on el pas central és només associativitat. El resultat $\Hom(f\comp f',g'\comp g)$ és precisament $\Hom$ aplicat a la composició de la parella $(f',g')$ amb $(f,g)$ a $\cat{C}\op\times\cat{C}$: a la segona component obtenim $g'\comp g$, i a la primera $f\comp f'$ (l'ordre invertit, coherent amb la contravariància). \textbf{Identitats.} Per a la parella identitat $(\id_A,\id_B)$,
\[
\Hom(\id_A,\id_B)(h)=\id_B\comp h\comp\id_A=h,
\]
és a dir $\Hom(\id_A,\id_B)=\id_{\Hom(A,B)}$. Per tant $\Hom(-,-)$ preserva composició i identitats de $\cat{C}\op\times\cat{C}$, i és un functor de ple dret.
\end{solucio}

\section{Functors fidels i plens}

\begin{exercici}
Comprova que el functor oblit $U\colon\cat{Pos}\to\cat{Set}$, que envia cada ordre parcial al seu conjunt subjacent i cada aplicació monòtona a ella mateixa, és fidel però no ple.
\end{exercici}

\begin{solucio}
\textbf{Fidel.} Siguin $f,g\colon P\to Q$ dues aplicacions monòtones amb $U(f)=U(g)$. Això vol dir que $f$ i $g$ coincideixen com a aplicacions entre els conjunts subjacents, és a dir, en tot element de $P$. Una aplicació monòtona no és res més que la seva acció sobre els elements (amb la condició afegida de preservar l'ordre), de manera que $f=g$. Així cada aplicació $U_{P,Q}\colon\Hom_{\cat{Pos}}(P,Q)\to\Hom_{\cat{Set}}(U(P),U(Q))$ és injectiva i, per tant, $U$ és fidel. (Sobre objectes, en canvi, $U$ no és injectiu: ordres diferents sobre un mateix conjunt tenen la mateixa imatge. Ser fidel no hi té res a veure.)

\textbf{No ple.} Prenem l'ordre parcial $P=\{0,1\}$ amb $0<1$ i el mateix $Q=P$. L'aplicació $\sigma\colon\{0,1\}\to\{0,1\}$ que intercanvia els dos elements ($\sigma(0)=1$, $\sigma(1)=0$) és una aplicació de conjunts perfectament vàlida, i pertany a $\Hom_{\cat{Set}}(U P,U Q)$. Però no és monòtona: de $0<1$ es passaria a $\sigma(0)=1$ i $\sigma(1)=0$, és a dir $1\le 0$, que és fals. Per tant $\sigma$ no és de la forma $U(f)$ per a cap aplicació monòtona $f$, l'aplicació $U_{P,Q}$ no és exhaustiva i $U$ no és ple.

\emph{Per a qui conegui els grups.} El functor oblit $U\colon\cat{Grp}\to\cat{Set}$ es comporta igual: és fidel perquè un homomorfisme queda determinat per la seva acció sobre els elements, i no és ple perquè hi ha aplicacions entre els conjunts subjacents que no són homomorfismes, com $x\mapsto x+1$ de $\mathbb Z$ en $\mathbb Z$, que no envia $0$ a $0$.
\end{solucio}

\begin{exercici}\label{pr-esquelet-finset}
Sigui $\cat{S}$ la subcategoria plena de $\cat{FinSet}$ que té per objectes els conjunts $\underline{0}=\varnothing$ i $\underline{n}=\{1,\dots,n\}$ per a $n\geq 1$. Comprova que la inclusió
\[
\iota\colon\cat{S}\hookrightarrow\cat{FinSet}
\]
és plena, fidel i essencialment exhaustiva, però no exhaustiva sobre objectes en sentit estricte.
\end{exercici}

\begin{solucio}
\textbf{Fidel.} La inclusió envia cada morfisme a ell mateix: per a objectes $\underline{m},\underline{n}$ de $\cat{S}$, l'aplicació $\iota_{\underline{m},\underline{n}}\colon\Hom_{\cat{S}}(\underline{m},\underline{n})\to\Hom_{\cat{FinSet}}(\underline{m},\underline{n})$ és $h\mapsto h$, que és injectiva.

\textbf{Plena.} Com que $\cat{S}$ és una subcategoria \emph{plena}, conté tots els morfismes de $\cat{FinSet}$ entre dos objectes seus: $\Hom_{\cat{S}}(\underline{m},\underline{n})=\Hom_{\cat{FinSet}}(\underline{m},\underline{n})$, el conjunt de totes les aplicacions $\underline{m}\to\underline{n}$. L'aplicació $\iota_{\underline{m},\underline{n}}$ és, doncs, la identitat d'aquest conjunt, i en particular exhaustiva.

\textbf{Essencialment exhaustiva.} Sigui $X$ un conjunt finit, i sigui $n$ el seu nombre d'elements. Per definició de cardinal finit, hi ha una bijecció $X\to\underline{n}$ (per a $n=0$, $X=\varnothing=\underline{0}$ i n'hi ha prou amb la identitat). Una bijecció és un isomorfisme de $\cat{FinSet}$, de manera que $X\cong\iota(\underline{n})$. Tot objecte de $\cat{FinSet}$ és, doncs, isomorf a una imatge de $\iota$.

\textbf{No exhaustiva sobre objectes.} El conjunt $\{2\}$ és un objecte de $\cat{FinSet}$, però no és cap dels $\underline{n}$: té un sol element, i l'únic $\underline{n}$ amb un element és $\underline{1}=\{1\}\neq\{2\}$. No és, per tant, igual a cap imatge de $\iota$, tot i ser isomorf a $\underline{1}$.

\emph{Per què és important.} Aquest functor reuneix les tres propietats ---ple, fidel i essencialment exhaustiu--- que al capítol~\ref{cap:transformacions} caracteritzaran les equivalències de categories, i alhora mostra que no cal arribar a tots els objectes en sentit estricte. Observem també que per a cada $X$ hem triat \emph{una} bijecció $X\to\underline{n}$; per construir un functor en sentit invers caldria triar-ne una per a cada conjunt finit alhora, i aquest és el punt on el capítol~\ref{cap:transformacions} recorrerà a l'axioma de l'elecció.
\end{solucio}


\chapter[Transformacions naturals]{\texorpdfstring{Transformacions naturals\\ i equivalències}{Transformacions naturals i equivalències}}

En aquest capítol resolem amb detall exercicis sobre \emph{transformacions naturals}, \emph{isomorfismes naturals} i \emph{equivalències de categories}. L'èmfasi és a comprovar la condició de naturalitat en casos concrets, a distingir un isomorfisme d'un isomorfisme natural, i a reconèixer una equivalència mitjançant les tres propietats que la caracteritzen. Convé intentar cada exercici abans de llegir-ne la solució.

\section{Definició de transformació natural}

\begin{exercici}
Siguin $F,G\colon\cat{C}\to\cat{D}$ dos functors i suposem que, per a cada objecte $A$, tenim un morfisme $\eta_A\colon F(A)\to G(A)$. Comprova que, per verificar que $\eta$ és una transformació natural, n'hi ha prou de comprovar la condició de naturalitat sobre un conjunt de morfismes que \emph{generi} $\cat{C}$ per composició, i que en el cas d'una categoria lliure $\Free(G_0)$ això vol dir comprovar-la només sobre les arestes del graf $G_0$.
\end{exercici}

\begin{solucio}
La condició de naturalitat per a un morfisme componible diu $G(f)\comp\eta_A=\eta_B\comp F(f)$. Suposem que val per a dos morfismes componibles $f\colon A\to B$ i $g\colon B\to C$. Aleshores, per a la composició $g\comp f$,
\[
G(g\comp f)\comp\eta_A=G(g)\comp G(f)\comp\eta_A=G(g)\comp\eta_B\comp F(f)=\eta_C\comp F(g)\comp F(f)=\eta_C\comp F(g\comp f),
\]
on hem fet servir la functorialitat de $G$ i $F$ i les condicions de naturalitat per a $f$ i per a $g$. Per tant, si la naturalitat val per a un conjunt de morfismes, val per a totes les seves composicions. També val trivialment per a les identitats, perquè $G(\id_A)\comp\eta_A=\eta_A=\eta_A\comp F(\id_A)$.

En una categoria lliure $\Free(G_0)$, tot morfisme és un camí, és a dir, una composició d'arestes del graf $G_0$. Per l'argument anterior, doncs, la naturalitat sobre tots els morfismes es dedueix de la naturalitat sobre les arestes. Això redueix molt la feina de comprovació.
\end{solucio}

\begin{exercici}
Considera la categoria $\cat{2}=(0\to 1)$ i dos functors $F,G\colon\cat{2}\to\cat{D}$. Descriu explícitament què és una transformació natural $\eta\colon F\Rightarrow G$.
\end{exercici}

\begin{solucio}
Un functor $F\colon\cat{2}\to\cat{D}$ és el mateix que un morfisme de $\cat{D}$: tria dos objectes $F(0),F(1)$ i un morfisme $F(u)\colon F(0)\to F(1)$, on $u\colon 0\to 1$ és l'única fletxa no trivial. Igualment $G$ dona un morfisme $G(u)\colon G(0)\to G(1)$.

Una transformació natural $\eta\colon F\Rightarrow G$ consisteix en dos components $\eta_0\colon F(0)\to G(0)$ i $\eta_1\colon F(1)\to G(1)$, subjectes a la condició de naturalitat per a $u$:
\[
G(u)\comp\eta_0=\eta_1\comp F(u).
\]
És a dir, una transformació natural entre dos morfismes de $\cat{D}$ és exactament un \textbf{quadrat commutatiu}
\[
\begin{tikzpicture}[baseline=(m.center)]
  \matrix (m) [matrix of math nodes, row sep=2.6em, column sep=3.2em] {
    F(0) & F(1) \\
    G(0) & G(1) \\
  };
  \draw[-{Stealth[scale=1.1]}] (m-1-1) -- node[above] {$F(u)$} (m-1-2);
  \draw[-{Stealth[scale=1.1]}] (m-1-1) -- node[left] {$\eta_0$} (m-2-1);
  \draw[-{Stealth[scale=1.1]}] (m-1-2) -- node[right] {$\eta_1$} (m-2-2);
  \draw[-{Stealth[scale=1.1]}] (m-2-1) -- node[below] {$G(u)$} (m-2-2);
\end{tikzpicture}
\]
Aquest exemple explica per què els quadrats commutatius apareixen tan sovint: són les transformacions naturals més simples possibles.
\end{solucio}

\section{Exemples de transformacions naturals}

\begin{exercici}\label{pr-sintaxi-semantica}
Considerem els functors sintàctic i semàntic $F,G\colon\cat{FinSet}\to\cat{Set}$ i la interpretació $\eta_X\colon F(X)\to G(X)$ de la secció~\ref{sec:cap-transformacio} de la Teoria. Comprova que $F$ i $G$ són functors i que $\eta$ és una transformació natural $F\Rightarrow G$.
\end{exercici}

\begin{solucio}
\textbf{$F$ és un functor.} Recordem que $F(f)(\varphi)=\varphi[f]$ es defineix per recursió: $F(f)(x)=f(x)$ sobre una variable, i $F(f)$ commuta amb les connectives, per exemple $F(f)(\varphi\land\psi)=F(f)(\varphi)\land F(f)(\psi)$. Comprovem les dues igualtats per recursió sobre $\varphi$. Per a la identitat, sobre una variable $F(\id_X)(x)=x$, i si $F(\id_X)$ deixa fixes $\varphi$ i $\psi$, també deixa fixa $\varphi\land\psi$ (i anàlogament per a les altres connectives); per tant $F(\id_X)=\id_{F(X)}$. Per a la composició, siguin $f\colon X\to Y$ i $g\colon Y\to Z$. Sobre una variable,
\[
F(g\comp f)(x)=g(f(x))=F(g)\bigl(F(f)(x)\bigr),
\]
i si la igualtat $F(g\comp f)=F(g)\comp F(f)$ val per a $\varphi$ i $\psi$, val per a $\varphi\land\psi$, perquè les tres aplicacions commuten amb $\land$. Així, substituir per $f$ i després per $g$ és substituir per $g\comp f$.

\textbf{$G$ és un functor.} Recordem que $G(f)(\theta)(w)=\theta(w\comp f)$ per a $\theta\in G(X)$ i $w\colon Y\to\{0,1\}$. Per a la identitat, $G(\id_X)(\theta)(v)=\theta(v\comp\id_X)=\theta(v)$. Per a la composició, si $f\colon X\to Y$, $g\colon Y\to Z$ i $u\colon Z\to\{0,1\}$,
\[
\begin{aligned}
G(g\comp f)(\theta)(u)
&=\theta\bigl(u\comp(g\comp f)\bigr)
=\theta\bigl((u\comp g)\comp f\bigr)\\
&=G(f)(\theta)(u\comp g)
=\bigl(G(g)\comp G(f)\bigr)(\theta)(u),
\end{aligned}
\]
on el pas clau és l'associativitat de la composició.

\textbf{Lema de substitució.} Per a tota $\varphi\in F(X)$ i tota $w\colon Y\to\{0,1\}$,
\[
(w\comp f)\models\varphi
\quad\Longleftrightarrow\quad
w\models\varphi[f].
\]
Per recursió sobre $\varphi$. Si $\varphi$ és una variable $x$, els dos costats diuen que $w(f(x))=1$, perquè $(w\comp f)(x)=w(f(x))$ i $x[f]=f(x)$. Si $\varphi=\psi\land\chi$, aleshores $(w\comp f)\models\varphi$ si i només si $(w\comp f)\models\psi$ i $(w\comp f)\models\chi$; per hipòtesi de recursió, si i només si $w\models\psi[f]$ i $w\models\chi[f]$, és a dir, si i només si $w\models\psi[f]\land\chi[f]=\varphi[f]$. Les altres connectives són anàlogues.

\textbf{Naturalitat.} Avaluem els dos camins del quadrat de naturalitat sobre $\varphi\in F(X)$ i $w\colon Y\to\{0,1\}$:
\[
\begin{aligned}
\bigl(\eta_Y(F(f)(\varphi))\bigr)(w)&=\eta_Y(\varphi[f])(w)=[\,w\models\varphi[f]\,],\\
\bigl(G(f)(\eta_X(\varphi))\bigr)(w)&=\eta_X(\varphi)(w\comp f)=[\,(w\comp f)\models\varphi\,],
\end{aligned}
\]
on $[\,\cdot\,]$ val $1$ o $0$ segons que la condició sigui certa o falsa. Pel lema de substitució, les dues expressions coincideixen per a tota $w$, de manera que $\eta_Y\comp F(f)=G(f)\comp\eta_X$ per a tota $f$, i $\eta$ és natural.
\end{solucio}

\begin{exercici}
Comprova en detall que l'avaluació $\eta_V\colon V\to V^{**}$, $\eta_V(v)(\varphi)=\varphi(v)$, és una transformació natural de la identitat cap al functor bidual (o doble dual), sobre la categoria dels espais vectorials sobre un cos $\mathbb K$ i les aplicacions lineals.
\end{exercici}

\begin{solucio}
Recordem primer com actua el bidual sobre morfismes. Donada una aplicació lineal $T\colon V\to W$, el seu dual $T^*\colon W^*\to V^*$ envia $\psi\mapsto\psi\comp T$, i el bidual $T^{**}\colon V^{**}\to W^{**}$ envia $\Phi\mapsto\Phi\comp T^*$.

La condició de naturalitat per a $T\colon V\to W$ és el quadrat
\[
T^{**}\comp\eta_V=\eta_W\comp T.
\]
Comprovem que tots dos costats coincideixen aplicats a un vector $v\in V$ i avaluats en una forma $\psi\in W^*$. D'una banda,
\[
\begin{aligned}
\big(T^{**}(\eta_V(v))\big)(\psi)
&=\big(\eta_V(v)\comp T^*\big)(\psi)
=\eta_V(v)\big(T^*(\psi)\big)\\
&=\eta_V(v)(\psi\comp T)
=(\psi\comp T)(v)
=\psi(T(v)).
\end{aligned}
\]
De l'altra,
\[
\big(\eta_W(T(v))\big)(\psi)=\psi\big(T(v)\big).
\]
Els dos resultats coincideixen per a tot $v$ i tot $\psi$, de manera que $T^{**}\comp\eta_V=\eta_W\comp T$. Per tant $\eta$ és natural. Cap pas no ha requerit triar una base: aquesta és la manifestació concreta de la naturalitat.
\end{solucio}

\begin{exercici}
Siguin $P$ i $Q$ preordres i $F,G,H\colon P\to Q$ aplicacions monòtones. Interpreta les transformacions naturals $F\Rightarrow G$ i comprova que la composició vertical correspon a la transitivitat de la desigualtat puntual.
\end{exercici}

\begin{solucio}
Com hem vist a la teoria, hi ha una transformació natural $F\Rightarrow G$ si i només si $F(x)\leq G(x)$ per a tot $x\in P$; i quan existeix, és única, perquè entre dos objectes d'un preordre hi ha com a màxim una fletxa.

Si tenim $\eta\colon F\Rightarrow G$ i $\theta\colon G\Rightarrow H$, això vol dir $F(x)\leq G(x)$ i $G(x)\leq H(x)$ per a tot $x$. La composició vertical $\theta\comp\eta\colon F\Rightarrow H$ existeix perquè, per transitivitat, $F(x)\leq H(x)$. Així, la composició vertical de transformacions naturals entre aplicacions monòtones és exactament la transitivitat de la relació \enquote{puntualment menor o igual}. La categoria de functors $Q^{P}$ resulta ser, en aquest cas, un preordre: el preordre de les aplicacions monòtones $P\to Q$ amb l'ordre puntual.
\end{solucio}

\section{Isomorfismes naturals}

\begin{exercici}[Ampliació, per a lectors amb àlgebra lineal]\label{ex:dual-no-natural}
Comprova que tot espai vectorial real $V$ de dimensió finita és isomorf al seu dual $V^*$, però que aquesta identificació no és \emph{canònica}: no existeix cap família d'isomorfismes $\eta_V\colon V\to V^*$ (per a $V$ espai vectorial real de dimensió finita) que sigui invariant sota tots els isomorfismes lineals, en el sentit que, per a tot isomorfisme $T\colon V\to W$, es compleixi $\eta_W\comp T=(T^{-1})^*\comp\eta_V$.
\end{exercici}

\begin{solucio}
Si $V$ té dimensió finita $n$, aleshores $V^*$ també té dimensió $n$, i dos espais reals de la mateixa dimensió finita són isomorfs. Un isomorfisme concret s'obté triant una base $\{e_1,\dots,e_n\}$ de $V$ i enviant $e_i$ a la base dual $\{e_i^*\}$. L'isomorfisme depèn de la base triada.

Cal notar primer que la comparació s'ha de plantejar amb cura, perquè la identitat és un functor covariant $\cat{Vect}_{\mathbb R}\to\cat{Vect}_{\mathbb R}$ i el dual és \emph{contravariant}, és a dir, un functor $\cat{Vect}_{\mathbb R}\op\to\cat{Vect}_{\mathbb R}$; entre functors de variància diferent no té sentit la noció de quadrat de naturalitat que hem definit. Per això formulem la invariància com a compatibilitat amb els \emph{isomorfismes} lineals, on el dual $(T^{-1})^*$ i $T$ van en el mateix sentit.

Que no existeixi una tal família es pot veure amb l'escalament. Sigui $V\neq 0$ i prenem l'homotècia $T=2\,\id_V$, que és un isomorfisme lineal (aquí fem servir que treballem sobre $\mathbb R$, on $2\neq 0$). Aleshores $T^{-1}=\tfrac12\,\id_V$ i $(T^{-1})^*=\tfrac12\,\id_{V^*}$. La condició d'invariància demanaria
\[
\eta_V\comp(2\,\id_V)=\bigl(\tfrac12\,\id_{V^*}\bigr)\comp\eta_V,
\]
és a dir, $2\,\eta_V=\tfrac12\,\eta_V$, o equivalentment $\tfrac32\,\eta_V=0$. Com que $\tfrac32\neq 0$, això força $\eta_V=0$, que no és un isomorfisme. Per tant no hi ha cap identificació $V\cong V^*$ invariant sota tots els isomorfismes lineals.

\begin{observacio}
El punt clau de l'argument és triar un escalar $\lambda$ amb $\lambda\neq\lambda^{-1}$, és a dir $\lambda^2\neq 1$: de la igualtat $\lambda\,\eta_V=\lambda^{-1}\eta_V$ només se'n dedueix $\eta_V=0$ quan $\lambda-\lambda^{-1}\neq 0$. No n'hi ha prou, doncs, amb $\lambda\neq 0,1$: sobre $\mathbb R$ l'escalar $\lambda=-1$ compliria $\lambda\neq 0,1$ però $\lambda^2=1$, i no serviria. Per això prenem $\lambda=2$, amb $\lambda^2=4\neq 1$. Sobre cossos petits pot no existir cap escalar amb $\lambda^2\neq 1$ (per exemple a $\mathbb F_2$ i $\mathbb F_3$ tots els no nuls compleixen $\lambda^2=1$), i l'argument de l'escalament s'ha d'ajustar; per això aquí ens restringim a $\mathbb R$.
\end{observacio}

Aquest contrast amb el bidual ---que, restringit als espais de dimensió finita, sí que és naturalment isomorf a la identitat--- és l'exemple canònic de la diferència entre \enquote{isomorf} i \enquote{naturalment isomorf}.
\end{solucio}

\begin{exercici}
Comprova que la relació \enquote{ser naturalment isomorf} és una relació d'equivalència sobre els functors $\cat{C}\to\cat{D}$.
\end{exercici}

\begin{solucio}
Cal veure reflexivitat, simetria i transitivitat.

\textbf{Reflexivitat}: la transformació identitat $\id_F\colon F\Rightarrow F$ té components $\id_{F(A)}$, que són isomorfismes; per tant $F\cong F$.

\textbf{Simetria}: si $\eta\colon F\Rightarrow G$ és un isomorfisme natural, hem vist a la teoria que els inversos $\eta_A^{-1}$ formen una transformació natural $\eta^{-1}\colon G\Rightarrow F$, també amb components isomorfismes; per tant $G\cong F$.

\textbf{Transitivitat}: si $\eta\colon F\Rightarrow G$ i $\theta\colon G\Rightarrow H$ són isomorfismes naturals, la composició vertical $\theta\comp\eta$ té components $\theta_A\comp\eta_A$, composició de dos isomorfismes, que és un isomorfisme. Per tant $F\cong H$.

De fet, aquesta és simplement la relació d'isomorfisme entre objectes de la categoria de functors $\cat{D}^{\cat{C}}$, que ja sabem que és una relació d'equivalència.
\end{solucio}

\section{La categoria de functors}

\begin{exercici}
Descriu la categoria de functors $\cat{D}^{\cat{1}}$, on $\cat{1}$ és la categoria amb un sol objecte i només la identitat.
\end{exercici}

\begin{solucio}
Un functor $F\colon\cat{1}\to\cat{D}$ tria un únic objecte $F(\ast)=D$ de $\cat{D}$ (i envia l'única identitat a $\id_D$). Per tant, els objectes de $\cat{D}^{\cat{1}}$ es corresponen amb els objectes de $\cat{D}$.

Una transformació natural $\eta\colon F\Rightarrow G$ entre dos tals functors té un únic component $\eta_\ast\colon F(\ast)\to G(\ast)$, és a dir, un morfisme $D\to D'$ de $\cat{D}$; la condició de naturalitat per a l'única fletxa (la identitat) es compleix trivialment. Per tant, els morfismes de $\cat{D}^{\cat{1}}$ es corresponen amb els morfismes de $\cat{D}$.

Concloem que $\cat{D}^{\cat{1}}$ és \emph{isomorfa} a $\cat{D}$: prendre functors des de la categoria terminal $\cat{1}$ recupera la categoria original.
\end{solucio}

\begin{exercici}
Comprova que un functor $F\colon\cat{2}\to\cat{D}$ és un objecte de la categoria de functors $\cat{D}^{\cat{2}}$, i que aquesta categoria és la \textbf{categoria de fletxes} de $\cat{D}$: els seus objectes són els morfismes de $\cat{D}$ i els seus morfismes són els quadrats commutatius.
\end{exercici}

\begin{solucio}
Ja hem vist que un functor $\cat{2}\to\cat{D}$ és un morfisme de $\cat{D}$, i que una transformació natural entre dos d'aquests functors és un quadrat commutatiu. Per tant, els objectes de $\cat{D}^{\cat{2}}$ són els morfismes de $\cat{D}$, i els morfismes de $\cat{D}^{\cat{2}}$ són els quadrats commutatius, amb la composició vertical corresponent a apilar quadrats. Aquesta categoria s'anomena \textbf{categoria de fletxes} de $\cat{D}$ i es denota sovint $\cat{D}^{\to}$. És un primer exemple de com les categories de functors produeixen construccions útils a partir de categories índex molt simples.
\end{solucio}

\section{Equivalència de categories}

\begin{exercici}
Comprova que la categoria dels conjunts finits és equivalent a la seva subcategoria plena formada pels conjunts $\underline{0}=\varnothing$ i $\underline{n}=\{1,\dots,n\}$ per a $n\geq 1$.
\end{exercici}

\begin{solucio}
Sigui $\cat{S}$ aquesta subcategoria plena i $\iota\colon\cat{S}\hookrightarrow\cat{FinSet}$ la inclusió. A l'exercici resolt~\ref{pr-esquelet-finset} del capítol~\ref{cap:functors} vam comprovar que $\iota$ és ple, fidel i essencialment exhaustiu. Pel teorema~\ref{teo:caract-equivalencia}, $\iota$ forma part d'una equivalència, i per tant $\cat{FinSet}\simeq\cat{S}$.

Per construir un quasiinvers $G\colon\cat{FinSet}\to\cat{S}$, la demostració del teorema tria, per a cada conjunt finit $X$ de $n$ elements, una bijecció $\varepsilon_X\colon\underline{n}\to X$, i posa $G(X)=\underline{n}$. Com que $\cat{FinSet}$ és gran, aquesta tria simultània fa servir l'axioma de l'elecció en la forma que explica l'observació que segueix el teorema. Intuïtivament: per fer teoria de categories amb conjunts finits, n'hi ha prou de considerar un conjunt de cada mida.
\end{solucio}

\begin{exercici}\label{pr-equivalencia-preserva}
Comprova que una equivalència de categories preserva monomorfismes, isomorfismes i objectes terminals. Explica per què, en canvi, \emph{no} preserva propietats estrictes, com ara la de tenir exactament un objecte.
\end{exercici}

\begin{solucio}
Sigui $F\colon\cat{C}\to\cat{D}$ part d'una equivalència, amb quasiinvers $G$ i isomorfismes naturals $\eta\colon\id_{\cat{C}}\Rightarrow G\comp F$ i $\varepsilon\colon F\comp G\Rightarrow\id_{\cat{D}}$, orientats com a la definició~\ref{def:equivalencia}. Tots dos són isomorfismes component a component, amb inversos $\eta^{-1}$ i $\varepsilon^{-1}$.

\textbf{Isomorfismes}: com que $F$ és un functor, preserva isomorfismes (proposició~\ref{prop:functor-preserva-iso}); i com que és ple i fidel ---per ser part d'una equivalència---, també els reflecteix. Així, $f$ és un isomorfisme a $\cat{C}$ si i només si $F(f)$ ho és a $\cat{D}$.

\textbf{Monomorfismes}: sigui $f\colon A\to B$ mono a $\cat{C}$; vegem que $F(f)\colon F(A)\to F(B)$ és mono a $\cat{D}$. Siguin $u,v\colon Z\to F(A)$ dos morfismes de $\cat{D}$ amb
\[
F(f)\comp u=F(f)\comp v.
\]
Apliquem el quasiinvers $G$ i fem servir la seva functorialitat:
\[
G(F(f))\comp G(u)=G\bigl(F(f)\comp u\bigr)=G\bigl(F(f)\comp v\bigr)=G(F(f))\comp G(v).
\]
Ara traslladem aquesta igualtat de $G(F(A))$ a $A$ mitjançant la naturalitat de $\eta$. Per a $f\colon A\to B$, el quadrat de naturalitat de $\eta$ dona
\[
G(F(f))\comp\eta_A=\eta_B\comp f,
\]
i, en ser $\eta_A,\eta_B$ isomorfismes, $G(F(f))=\eta_B\comp f\comp\eta_A^{-1}$. Substituint a la igualtat anterior i cancel·lant l'isomorfisme $\eta_B$ per l'esquerra,
\[
f\comp\bigl(\eta_A^{-1}\comp G(u)\bigr)=f\comp\bigl(\eta_A^{-1}\comp G(v)\bigr).
\]
Com que $f$ és mono a $\cat{C}$, en resulta $\eta_A^{-1}\comp G(u)=\eta_A^{-1}\comp G(v)$, i cancel·lant l'isomorfisme $\eta_A^{-1}$ obtenim $G(u)=G(v)$. Finalment, $G$ és fidel ---per ser part d'una equivalència---, de manera que $u=v$. Per tant $F(f)$ és mono.

\textbf{Objectes terminals}: sigui $1$ terminal a $\cat{C}$; vegem que $F(1)$ és terminal a $\cat{D}$. Sigui $D$ un objecte qualsevol de $\cat{D}$.

\emph{Existència}: com que $1$ és terminal, hi ha un únic morfisme $!\colon G(D)\to 1$ a $\cat{C}$. Aleshores el morfisme
\[
D\xrightarrow{\ \varepsilon_D^{-1}\ }F(G(D))\xrightarrow{\ F(!)\ }F(1)
\]
és un morfisme $D\to F(1)$, cosa que en garanteix l'existència.

\emph{Unicitat}: siguin $h,h'\colon D\to F(1)$ dos morfismes. Aplicant $G$, obtenim $G(h),G(h')\colon G(D)\to G(F(1))$. Component amb l'isomorfisme $\eta_1^{-1}\colon G(F(1))\to 1$, els morfismes $\eta_1^{-1}\comp G(h)$ i $\eta_1^{-1}\comp G(h')$ van de $G(D)$ a $1$ i, per tant, coincideixen per la unicitat cap a l'objecte terminal $1$. Com que $\eta_1^{-1}$ és invertible, resulta $G(h)=G(h')$, i com que $G$ és fidel, $h=h'$. Per tant $F(1)$ és terminal.

\textbf{Per què no tota propietat}: una equivalència \emph{no} preserva les propietats que depenen de la \emph{igualtat} d'objectes. Per exemple, \enquote{tenir exactament un objecte} o \enquote{ser esquelètica} no es conserven: la categoria amb un sol objecte i la seva versió amb dos objectes isomorfs són equivalents, però la primera té un objecte i la segona en té dos.

Convé distingir, a més, dues menes de nocions entre les que hem comprovat. Ser isomorfisme i ser monomorfisme són \emph{propietats de morfismes}, definides per equacions o per cancel·lació; l'equivalència les preserva, i de la mateixa manera es comprova que preserva els epimorfismes. Ser objecte terminal és, en canvi, una \emph{propietat universal} d'un objecte, que el caracteritza només llevat d'isomorfisme; el que hem vist és que la imatge d'un terminal torna a ser terminal. Més endavant trobarem altres construccions universals ---objectes inicials, límits i colímits--- i en comprovarem el comportament sota equivalència cas per cas.
\end{solucio}


\section{Composició horitzontal}

\begin{exercici}\label{pr-comp-horitzontal}
Siguin $\eta\colon F\Rightarrow G$, amb $F,G\colon\cat{C}\to\cat{D}$, i $\theta\colon K\Rightarrow L$, amb $K,L\colon\cat{D}\to\cat{E}$, com a la definició de composició horitzontal. Escriu els tipus de
\[
\theta_{G(A)}\comp K(\eta_A)
\qquad\text{i}\qquad
L(\eta_A)\comp\theta_{F(A)},
\]
explica per què coincideixen per la naturalitat de $\theta$, i verifica el quadrat de naturalitat de $\theta\ast\eta\colon K\comp F\Rightarrow L\comp G$.
\end{exercici}

\begin{solucio}
\textbf{Els tipus.} El morfisme $\eta_A\colon F(A)\to G(A)$ viu a $\cat{D}$, i $K$ el porta a $K(\eta_A)\colon K(F(A))\to K(G(A))$ a $\cat{E}$. La component de $\theta$ a l'objecte $G(A)$ de $\cat{D}$ és $\theta_{G(A)}\colon K(G(A))\to L(G(A))$. Per tant
\[
\theta_{G(A)}\comp K(\eta_A)\colon K(F(A))\longrightarrow L(G(A)).
\]
Anàlogament, $\theta_{F(A)}\colon K(F(A))\to L(F(A))$ i $L(\eta_A)\colon L(F(A))\to L(G(A))$, de manera que
\[
L(\eta_A)\comp\theta_{F(A)}\colon K(F(A))\longrightarrow L(G(A)).
\]
Tots dos morfismes tenen, doncs, el mateix origen i el mateix destí, i té sentit preguntar-se si coincideixen.

\textbf{Per què coincideixen.} Són els dos camins del quadrat de naturalitat de $\theta$ per al morfisme $\eta_A\colon F(A)\to G(A)$ de $\cat{D}$, dibuixat a l'observació~\ref{obs:comp-horitzontal} de la Teoria:
\[
\theta_{G(A)}\comp K(\eta_A)=L(\eta_A)\comp\theta_{F(A)}.
\]
La naturalitat de $\theta$ diu que aquest quadrat commuta per a \emph{tot} morfisme de $\cat{D}$; en particular, per a $\eta_A$. La component $(\theta\ast\eta)_A$ està, doncs, ben definida, sigui quina sigui l'expressió que fem servir.

\textbf{Naturalitat de $\theta\ast\eta$.} Sigui $f\colon A\to B$ un morfisme de $\cat{C}$. Cal veure que
\[
\begin{tikzpicture}[baseline=(m.center)]
  \matrix (m) [matrix of math nodes, row sep=2.8em, column sep=4.6em] {
    K(F(A)) & K(F(B)) \\
    L(G(A)) & L(G(B)) \\
  };
  \draw[-{Stealth[scale=1.1]}] (m-1-1) -- node[above] {$K(F(f))$} (m-1-2);
  \draw[-{Stealth[scale=1.1]}] (m-1-1) -- node[left] {$(\theta\ast\eta)_A$} (m-2-1);
  \draw[-{Stealth[scale=1.1]}] (m-1-2) -- node[right] {$(\theta\ast\eta)_B$} (m-2-2);
  \draw[-{Stealth[scale=1.1]}] (m-2-1) -- node[below] {$L(G(f))$} (m-2-2);
\end{tikzpicture}
\qquad
(\theta\ast\eta)_B\comp K(F(f))=L(G(f))\comp(\theta\ast\eta)_A
\]
commuta. Fent servir l'expressió $\theta_{G(-)}\comp K(\eta_-)$ per a les components,
\[
\begin{aligned}
(\theta\ast\eta)_B\comp K(F(f))
&=\theta_{G(B)}\comp K(\eta_B)\comp K(F(f))
=\theta_{G(B)}\comp K\bigl(\eta_B\comp F(f)\bigr)\\
&=\theta_{G(B)}\comp K\bigl(G(f)\comp\eta_A\bigr)
=\theta_{G(B)}\comp K(G(f))\comp K(\eta_A)\\
&=L(G(f))\comp\theta_{G(A)}\comp K(\eta_A)
=L(G(f))\comp(\theta\ast\eta)_A.
\end{aligned}
\]
Hem fet servir, per ordre: la functorialitat de $K$, la naturalitat de $\eta$ per a $f$, de nou la functorialitat de $K$, i la naturalitat de $\theta$ per al morfisme $G(f)$ de $\cat{D}$. Cada naturalitat s'utilitza una sola vegada, i cadascuna al seu nivell: la de $\eta$ dins de $\cat{D}$, i la de $\theta$ dins de $\cat{E}$.
\end{solucio}


\ifpublicacioparcial\else
\chapter[Representabilitat i Yoneda]{\texorpdfstring{Propietats universals,\\ representabilitat i Yoneda}{Propietats universals, representabilitat i Yoneda}}

Els exercicis d'aquest capítol treballen les propietats universals a través de la representabilitat i el lema de Yoneda. L'objectiu és guanyar fluïdesa amb els functors representables, amb el càlcul de transformacions naturals que en surten, i amb l'argument ---repetitiu, un cop entès--- que dedueix unicitat llevat d'isomorfisme.

\section{Objectes inicials i terminals}

\begin{exercici}\label{pr-producte-terminal}
Siguin $A$ i $B$ dos objectes d'una categoria $\cat{C}$. Un \textbf{con} sobre $A$ i $B$ és una terna $(X,f,g)$ amb $f\colon X\to A$ i $g\colon X\to B$, i un morfisme de cons $(X,f,g)\to(X',f',g')$ és un morfisme $h\colon X\to X'$ de $\cat{C}$ amb $f'\comp h=f$ i $g'\comp h=g$. Comprova que els cons sobre $A$ i $B$ formen una categoria $\mathrm{Con}(A,B)$, i que un con $(P,p_1,p_2)$ és un producte de $A$ i $B$ si i només si és un objecte terminal de $\mathrm{Con}(A,B)$. Interpreta el resultat en la categoria d'un preordre $(Q,\le)$.
\end{exercici}

\begin{solucio}
\textbf{És una categoria.} La identitat de $(X,f,g)$ és $\id_X$, que compleix $f\comp\id_X=f$ i $g\comp\id_X=g$. Si $h\colon(X,f,g)\to(X',f',g')$ i $h'\colon(X',f',g')\to(X'',f'',g'')$ són morfismes de cons, la composta $h'\comp h$ també ho és, perquè
\[
f''\comp(h'\comp h)=(f''\comp h')\comp h=f'\comp h=f,
\]
i igual amb $g$. L'associativitat i les lleis d'identitat s'hereten de $\cat{C}$, perquè la composició és la de $\cat{C}$.

\textbf{Producte i objecte terminal.} Un morfisme de cons $h\colon(X,f,g)\to(P,p_1,p_2)$ és exactament un morfisme $h\colon X\to P$ que fa commutar els dos triangles,
\[
\begin{tikzpicture}[baseline=(m.center)]
  \matrix (m) [matrix of math nodes, row sep=2.6em, column sep=2.8em] {
     & X & \\
    A & P & B \\
  };
  \draw[-{Stealth[scale=1.0]}] (m-1-2) -- node[above left] {$f$} (m-2-1);
  \draw[-{Stealth[scale=1.0]}] (m-1-2) -- node[above right] {$g$} (m-2-3);
  \draw[-{Stealth[scale=1.0]},dashed] (m-1-2) -- node[right] {$h$} (m-2-2);
  \draw[-{Stealth[scale=1.0]}] (m-2-2) -- node[below] {$p_1$} (m-2-1);
  \draw[-{Stealth[scale=1.0]}] (m-2-2) -- node[below] {$p_2$} (m-2-3);
\end{tikzpicture}
\qquad
p_1\comp h=f,\quad p_2\comp h=g.
\]
Que $(P,p_1,p_2)$ sigui terminal a $\mathrm{Con}(A,B)$ diu que, per a cada con $(X,f,g)$, hi ha exactament un morfisme de cons cap a $(P,p_1,p_2)$, és a dir, exactament un $h$ amb $p_1\comp h=f$ i $p_2\comp h=g$. Això és, paraula per paraula, la definició de producte (subsecció~\ref{subsec:producte-coproducte}). Les dues afirmacions són, doncs, la mateixa.

Com a conseqüència, la unicitat dels productes es dedueix de la dels objectes terminals (proposició~\ref{prop:unicitat-terminal}): dos productes de $A$ i $B$ són isomorfs per un únic isomorfisme \emph{de cons}, és a dir, compatible amb les projeccions. L'objecte $P$ tot sol, en canvi, pot tenir altres automorfismes que no respectin $p_1$ i $p_2$; per exemple, a $\cat{Set}$, l'intercanvi $(a,b)\mapsto(b,a)$ de $A\times A$.

\textbf{En un preordre.} A la categoria d'un preordre $(Q,\le)$, un con sobre $a$ i $b$ és un element $x$ amb $x\le a$ i $x\le b$, és a dir, una cota inferior comuna, i entre dos cons hi ha com a màxim un morfisme, que existeix quan $x\le x'$. Un objecte terminal de $\mathrm{Con}(a,b)$ és, doncs, una cota inferior $p$ tal que tota altra cota inferior $x$ compleix $x\le p$: la \emph{màxima cota inferior}, o ínfim, de $a$ i $b$. Els productes en un preordre són els ínfims, i en un ordre parcial són únics; en un preordre, únics llevat de la relació $p\le p'\le p$.
\end{solucio}

\begin{exercici}
Comprova que un objecte $Z$ és alhora inicial i terminal en una categoria si i només si, per a cada objecte $X$, hi ha exactament un morfisme $Z\to X$ i exactament un morfisme $X\to Z$. Un tal objecte s'anomena \emph{objecte zero}. Comprova que a $\cat{Pos}$ no n'hi ha cap i, per a qui conegui els grups, que a $\cat{Grp}$ el grup trivial n'és un.
\end{exercici}

\begin{solucio}
La primera afirmació és la conjunció literal de les dues definicions: $Z$ inicial dona el morfisme únic $Z\to X$, i $Z$ terminal, el morfisme únic $X\to Z$.

A $\cat{Pos}$, l'objecte inicial és el conjunt ordenat buit i el terminal és el d'un sol element. Un objecte zero hauria de ser isomorf a tots dos, i aquests dos no són isomorfs, perquè no hi ha cap aplicació $\{\ast\}\to\varnothing$. Per tant, $\cat{Pos}$ no té objecte zero.

Per al grup trivial $1=\{e\}$: donat un grup $G$, hi ha un únic homomorfisme $1\to G$ (l'element neutre ha d'anar al neutre), i un únic homomorfisme $G\to 1$ (tot va al neutre). Per tant $1$ és inicial i terminal a $\cat{Grp}$, és a dir, un objecte zero.
\end{solucio}

\section{Categories coma}

\begin{exercici}\label{pr-coma-lliure}
Sigui $G$ un graf, i sigui $U\colon\cat{Cat}\to\cat{Graph}$ el functor graf subjacent, on $\cat{Cat}$ té per objectes les categories petites. Descriu explícitament la categoria coma $(G\downarrow U)$: els seus objectes, els seus morfismes, la composició i les identitats. Demostra que la parella $(\Free(G),\eta_G)$, on $\eta_G\colon G\to U(\Free(G))$ és la inserció dels generadors, és un objecte inicial de $(G\downarrow U)$.
\end{exercici}

\begin{solucio}
\textbf{Objectes i morfismes.} Un objecte de $(G\downarrow U)$ és una parella $(\cat{D},\varphi)$ formada per una categoria petita $\cat{D}$ i un morfisme de grafs $\varphi\colon G\to U(\cat{D})$. Un morfisme $(\cat{D},\varphi)\to(\cat{D}',\varphi')$ és un functor $H\colon\cat{D}\to\cat{D}'$ que fa commutar el triangle
\[
\begin{tikzpicture}[baseline=(m.center)]
  \matrix (m) [matrix of math nodes, row sep=2.6em, column sep=2.4em] {
     & G & \\
    U(\cat{D}) & & U(\cat{D}') \\
  };
  \draw[-{Stealth[scale=1.0]}] (m-1-2) -- node[above left] {$\varphi$} (m-2-1);
  \draw[-{Stealth[scale=1.0]}] (m-1-2) -- node[above right] {$\varphi'$} (m-2-3);
  \draw[-{Stealth[scale=1.0]}] (m-2-1) -- node[below] {$U(H)$} (m-2-3);
\end{tikzpicture}
\qquad
U(H)\comp\varphi=\varphi'.
\]

\textbf{És una categoria.} La identitat de $(\cat{D},\varphi)$ és $\id_{\cat{D}}$, perquè $U(\id_{\cat{D}})\comp\varphi=\varphi$. Si $H\colon(\cat{D},\varphi)\to(\cat{D}',\varphi')$ i $H'\colon(\cat{D}',\varphi')\to(\cat{D}'',\varphi'')$, la composta $H'\comp H$ també és un morfisme, per la functorialitat de $U$:
\[
U(H'\comp H)\comp\varphi=U(H')\comp U(H)\comp\varphi=U(H')\comp\varphi'=\varphi''.
\]
L'associativitat i les lleis d'identitat s'hereten de $\cat{Cat}$.

\textbf{$(\Free(G),\eta_G)$ és inicial.} Sigui $(\cat{D},\varphi)$ un objecte qualsevol. Un morfisme $(\Free(G),\eta_G)\to(\cat{D},\varphi)$ és un functor $\bar\varphi\colon\Free(G)\to\cat{D}$ amb $U(\bar\varphi)\comp\eta_G=\varphi$, és a dir, que coincideix amb $\varphi$ sobre els vèrtexs i sobre els camins de longitud~$1$.
\begin{itemize}
  \item \emph{Existència.} Definim $\bar\varphi$ sobre els objectes com $\varphi$, sobre el camí buit en un vèrtex $x$ com $\id_{\varphi(x)}$, i sobre un camí $(a_1,\dots,a_n)$ recorregut en aquest ordre com la composta $\varphi(a_n)\comp\cdots\comp\varphi(a_1)$. Com que la composició de $\Free(G)$ és la concatenació, $\bar\varphi$ respecta composicions i identitats.
  \item \emph{Unicitat.} Tot camí és la composta de les seves arestes, i tot camí buit és una identitat. Un functor amb $U(\bar\varphi)\comp\eta_G=\varphi$ queda, per tant, determinat sobre tots els morfismes per la seva acció sobre les arestes, que està fixada per $\varphi$.
\end{itemize}
Així, des de $(\Free(G),\eta_G)$ hi ha exactament un morfisme cap a cada objecte de $(G\downarrow U)$: és un objecte inicial. Observem que no ha calgut la noció d'adjunció; és aquesta propietat la que, al capítol d'adjuncions, farà de $\Free$ un adjunt per l'esquerra de $U$.
\end{solucio}

\section{Propietats universals com a representabilitat}

\begin{exercici}
Comprova que un objecte $1$ és terminal si i només si el functor $\Hom(-,1)$ és naturalment isomorf al functor constant $\Delta_{\{\ast\}}\colon\cat{C}\op\to\cat{Set}$ (exemple~\ref{ex:functor-constant}).
\end{exercici}

\begin{solucio}
Suposem $1$ terminal. Aleshores, per a cada objecte $X$, el conjunt $\Hom(X,1)$ té exactament un element, de manera que hi ha una bijecció $\Hom(X,1)\cong\{\ast\}$. Aquestes bijeccions són automàticament naturals: l'única aplicació possible cap a un conjunt d'un punt fa commutar qualsevol quadrat. Per tant $\Hom(-,1)$ és naturalment isomorf al functor constant $\Delta_{\{\ast\}}$.

Recíprocament, si $\Hom(-,1)\cong\Delta_{\{\ast\}}$, aleshores cada $\Hom(X,1)$ té exactament un element, és a dir, hi ha exactament un morfisme $X\to 1$ per a cada $X$, i $1$ és terminal.
\end{solucio}

\begin{exercici}
Sigui $\cat{C}$ una categoria localment petita amb un objecte terminal $1$. Comprova que el \textbf{functor de punts globals}
\[
\Hom(1,-)\colon\cat{C}\longrightarrow\cat{Set}
\]
envia cada objecte $X$ a $\Hom(1,X)$. Aquests morfismes s'anomenen \emph{punts globals} de $X$. Calcula aquest conjunt a $\cat{Set}$.
\end{exercici}

\begin{solucio}
Per definició, $\Hom(1,-)$ envia $X$ a $\Hom(1,X)$, el conjunt de morfismes des de l'objecte terminal, és a dir, dels punts globals de $X$. A $\cat{Set}$, l'objecte terminal és un conjunt d'un sol element $\{\ast\}$, i una aplicació $\{\ast\}\to X$ equival a triar un element de $X$: la imatge de $\ast$. Per tant $\Hom(1,X)\cong X$ de manera natural, i a $\cat{Set}$ els punts globals recuperen exactament els elements. En general, però, $\Hom(1,-)$ pot perdre informació: hi ha categories on objectes diferents tenen els mateixos punts globals.
\end{solucio}

\section{El lema de Yoneda}

\begin{exercici}
Sigui $\cat{C}$ localment petita i $A$ un objecte. Aplicant el lema de Yoneda amb $P=\Hom(-,A)$, comprova que les úniques transformacions naturals $\Hom(-,A)\Rightarrow\Hom(-,A)$ que hi ha es corresponen amb els endomorfismes de $A$, i que la identitat correspon a $\id_A$.
\end{exercici}

\begin{solucio}
El lema de Yoneda dona una bijecció
\[
\Nat\big(\Hom(-,A),\Hom(-,A)\big)\cong\Hom(-,A)(A)=\Hom(A,A),
\]
que envia $\eta$ a $\eta_A(\id_A)$. Així, cada transformació natural de $\Hom(-,A)$ en ell mateix correspon a un endomorfisme de $A$. La transformació identitat $\id_{\Hom(-,A)}$ té components la identitat de cada $\Hom(X,A)$; el seu valor a $\id_A$ és $\id_A$, de manera que correspon a l'endomorfisme identitat. Aquest càlcul és el nucli de la demostració que la immersió de Yoneda és plena i fidel.
\end{solucio}

\begin{exercici}\label{pr-yoneda-parts}
Sigui $A$ un conjunt i sigui $\mathcal P\colon\cat{Set}\op\to\cat{Set}$ el functor de parts contravariant. Descriu explícitament, mitjançant la bijecció del lema de Yoneda, les transformacions naturals $\Hom_{\cat{Set}}(-,A)\Rightarrow\mathcal P$. Quina transformació correspon a un subconjunt $B\subseteq A$? Comprova directament que la bijecció recupera $B$, i calcula-la en el cas $A=\{0,1\}$, $B=\{1\}$, sobre l'aplicació $\varphi\colon\{a,b,c\}\to A$ amb $\varphi(a)=\varphi(c)=1$ i $\varphi(b)=0$.
\end{exercici}

\begin{solucio}
El lema de Yoneda dona una bijecció
\[
\Phi\colon\Nat\bigl(\Hom(-,A),\mathcal P\bigr)\longrightarrow\mathcal P(A),
\qquad
\Phi(\alpha)=\alpha_A(\id_A),
\]
amb inversa $\Psi$, que envia un element $x\in\mathcal P(A)$ a la transformació $\Psi(x)_X(\varphi)=\mathcal P(\varphi)(x)$. Aquí un element de $\mathcal P(A)$ és un subconjunt $B\subseteq A$, i $\mathcal P(\varphi)$ és l'antiimatge. Per tant, la transformació natural que correspon a $B$ és
\[
\Psi(B)_X\colon\Hom(X,A)\to\mathcal P(X),
\qquad
\varphi\longmapsto\varphi^{-1}(B).
\]
És a dir: tota assignació natural d'aquest tipus s'obté per antiimatge d'un únic subconjunt fix de $A$, i subconjunts diferents donen transformacions naturals diferents.

\textbf{La bijecció recupera $B$.} En un sentit, $\Phi(\Psi(B))=\Psi(B)_A(\id_A)=\id_A^{-1}(B)=B$. En l'altre, si $\alpha$ és una transformació natural qualsevol i $B=\alpha_A(\id_A)$, la naturalitat de $\alpha$ aplicada a $\varphi\colon X\to A$ i avaluada sobre $\id_A$ dona
\[
\alpha_X(\varphi)=\alpha_X(\id_A\comp\varphi)=\mathcal P(\varphi)\bigl(\alpha_A(\id_A)\bigr)=\varphi^{-1}(B),
\]
de manera que $\alpha=\Psi(B)$. La naturalitat de $\Psi(B)$ es veu també directament: per a $g\colon Y\to X$, $(\varphi\comp g)^{-1}(B)=g^{-1}\bigl(\varphi^{-1}(B)\bigr)$.

\textbf{Un càlcul.} Amb $A=\{0,1\}$ i $B=\{1\}$, la component a $X=\{a,b,c\}$ envia $\varphi$ a $\varphi^{-1}(\{1\})=\{a,c\}$. Observem que, sota l'isomorfisme natural $\mathcal P\cong\Hom(-,\{0,1\})$ de l'exemple~\ref{ex:isos-naturals-cap2}, el subconjunt $B$ correspon a la seva funció característica $\chi_B\colon A\to\{0,1\}$, i la transformació $\varphi\mapsto\varphi^{-1}(B)$ correspon a $\varphi\mapsto\chi_B\comp\varphi$: és el cas $P=\Hom(-,\{0,1\})$ del lema, en què les transformacions naturals $\Hom(-,A)\Rightarrow\Hom(-,\{0,1\})$ són les postcomposicions amb morfismes $A\to\{0,1\}$.

\textbf{Lectura lògica.} Un subconjunt $B\subseteq A$ és un \emph{predicat} sobre $A$, i $\varphi^{-1}(B)$ és el predicat sobre $X$ que s'obté substituint $\varphi$: $x\in\varphi^{-1}(B)$ si i només si $\varphi(x)\in B$. El lema de Yoneda diu, doncs, que les maneres naturals de produir predicats a partir d'aplicacions cap a $A$ són exactament les substitucions en un predicat fix sobre $A$. És la mateixa operació de substitució que, al capítol~\ref{cap:transformacions}, feia natural la interpretació de les fórmules.
\end{solucio}

\begin{exercici}\label{pr-yoneda-grafs}
Siguin $V,E\colon\cat{Graph}\to\cat{Set}$ els functors de vèrtexs i d'arestes, siguin $\mathbf V$ el graf d'un sol vèrtex sense arestes i $\mathbf E$ el graf $0\xrightarrow{e}1$. Calcula, amb la versió covariant del lema de Yoneda, $\Nat(E,V)$ i $\Nat(V,E)$, i interpreta les transformacions que hi surtin.
\end{exercici}

\begin{solucio}
Per l'exemple~\ref{ex:isos-naturals-cap2}, $V\cong\Hom_{\cat{Graph}}(\mathbf V,-)$ i $E\cong\Hom_{\cat{Graph}}(\mathbf E,-)$. La versió covariant del lema dona, per a qualsevol functor $Q\colon\cat{Graph}\to\cat{Set}$, una bijecció entre les transformacions naturals que surten de $\Hom(A,-)$ i els elements de $Q(A)$. Per tant,
\[
\Nat(E,V)\;\cong\;V(\mathbf E)=\{0,1\},
\qquad
\Nat(V,E)\;\cong\;E(\mathbf V)=\varnothing.
\]
Observem la inversió: el primer càlcul avalua el functor d'\emph{arribada}, $V$, en l'objecte que \emph{representa} el functor de sortida, $E$.

\textbf{Les dues transformacions $E\Rightarrow V$.} L'element $0\in V(\mathbf E)$ correspon a la transformació que, a cada graf $G$, envia una aresta a una aplicació $E_G\to V_G$ calculada així: l'aresta $\varepsilon$ correspon al morfisme $\mathbf E\to G$ que la tria, i aquest envia el vèrtex $0$ a l'origen de $\varepsilon$. És, doncs, l'aplicació origen $s_G$. Anàlogament, l'element $1$ dona l'aplicació destí $t_G$. La naturalitat de $s$ i de $t$ és exactament la condició que un morfisme de grafs preservi origen i destí. Encara que en un graf on totes les arestes són bucles $s_G=t_G$, les dues transformacions són diferents, perquè a $\mathbf E$ les components difereixen.

\textbf{Cap transformació $V\Rightarrow E$.} També es veu directament: una transformació natural $V\Rightarrow E$ hauria de tenir una component $V_{\mathbf V}=\{\ast\}\to E_{\mathbf V}=\varnothing$, que no existeix. No hi ha cap manera natural de triar una aresta per a cada vèrtex en tots els grafs alhora.
\end{solucio}

\section{Conseqüències}

\begin{exercici}\label{pr-yoneda-preordre}
Sigui $(Q,\le)$ un preordre vist com a categoria. Descriu el prefeix representable $\Hom(-,a)$ i calcula, amb el lema de Yoneda, les transformacions naturals $\Hom(-,a)\Rightarrow\Hom(-,b)$. Dedueix que
\[
a\le b\quad\Longleftrightarrow\quad{\downarrow}a\subseteq{\downarrow}b,
\qquad\text{on } {\downarrow}a=\{x\in Q\mid x\le a\},
\]
i interpreta el resultat a la categoria $\cat{Prov}_T$ d'un sistema deductiu.
\end{exercici}

\begin{solucio}
\textbf{El representable.} Com que entre dos objectes hi ha com a màxim una fletxa, $\Hom(x,a)$ és un conjunt d'un sol element si $x\le a$, i és buit altrament. El prefeix $\Hom(-,a)$ és, doncs, la funció característica del conjunt ${\downarrow}a$ dels elements que estan per sota de $a$: \enquote{val} en $x$ exactament quan $x\in{\downarrow}a$.

\textbf{Les transformacions naturals.} Pel lema de Yoneda (o per la immersió plena i fidel),
\[
\Nat\bigl(\Hom(-,a),\Hom(-,b)\bigr)\;\cong\;\Hom(a,b),
\]
que té un element si $a\le b$ i cap altrament. Ho podem veure també directament: una transformació natural té una component $\Hom(x,a)\to\Hom(x,b)$ per a cada $x$, i una aplicació d'un conjunt d'un sol element en un conjunt buit no existeix. Hi ha, doncs, una transformació natural exactament quan $x\le a$ implica $x\le b$ per a tot $x$, és a dir, quan ${\downarrow}a\subseteq{\downarrow}b$; i aleshores és única, i la naturalitat és automàtica, perquè tots els valors tenen com a màxim un element.

\textbf{La conclusió.} Comparant les dues descripcions, $a\le b$ si i només si ${\downarrow}a\subseteq{\downarrow}b$. La implicació d'esquerra a dreta és la transitivitat; la de dreta a esquerra s'obté prenent $x=a$, que és exactament l'avaluació a $\id_a$ del lema de Yoneda. Com a cas particular, $a\cong b$ si i només si ${\downarrow}a={\downarrow}b$: dos elements són isomorfs quan tenen els mateixos elements per sota.

\textbf{A $\cat{Prov}_T$.} Aquí ${\downarrow}p$ és el conjunt de les fórmules que dedueixen $p$, i el resultat diu que $p\vdash_T q$ si i només si tota fórmula que dedueix $p$ dedueix també $q$. Una fórmula queda determinada, llevat d'interdeduïbilitat, per les fórmules de les quals es dedueix.
\end{solucio}

\begin{exercici}\label{pr-yoneda-monoide}
Sigui $M$ un monoide i $\cat{B}M$ la categoria d'un sol objecte $\ast$ associada, amb composició $g\comp f=g\cdot f$ (escrivim $\cdot$ per a l'operació del monoide i $e$ per al neutre, que és $\id_\ast$). Comprova que un prefeix $X\colon(\cat{B}M)\op\to\cat{Set}$ és el mateix que un conjunt amb una \emph{acció per la dreta} de $M$, i descriu el prefeix representable $\Hom(-,\ast)$. Aplica el lema de Yoneda per calcular $\Nat\bigl(\Hom(-,\ast),X\bigr)$ i, en particular, $\Nat\bigl(\Hom(-,\ast),\Hom(-,\ast)\bigr)$.
\end{exercici}

\begin{solucio}
\textbf{Prefeixos i accions.} Un prefeix $X$ dona un conjunt $X(\ast)$, que escrivim $X$, i per a cada $g\in M$ una aplicació $X(g)\colon X\to X$. Escrivim $x\cdot g=X(g)(x)$. La functorialitat contravariant, $X(g\comp f)=X(f)\comp X(g)$ i $X(e)=\id$, diu exactament
\[
(x\cdot g)\cdot f=x\cdot(g\cdot f),
\qquad
x\cdot e=x:
\]
és una acció per la dreta. Una transformació natural entre dos prefeixos és una aplicació $\eta\colon X\to Y$ amb $\eta(x\cdot g)=\eta(x)\cdot g$, és a dir, una aplicació \emph{equivariant}.

\textbf{El representable.} $\Hom(\ast,\ast)=M$, i un morfisme $g$ actua per precomposició: $m\mapsto m\comp g=m\cdot g$. El prefeix representable és, doncs, $M$ amb l'acció per la dreta per multiplicació, l'\emph{acció regular}.

\textbf{Yoneda.} El lema dona una bijecció $\Nat\bigl(\Hom(-,\ast),X\bigr)\cong X$: l'aplicació equivariant $\eta\colon M\to X$ correspon a $\eta(e)$, i recíprocament $x\in X$ dona l'aplicació $m\mapsto x\cdot m$. Una aplicació equivariant des de l'acció regular queda determinada pel lloc on envia l'element neutre.

Per a $X=M$ regular, les aplicacions equivariants $M\to M$ són exactament les multiplicacions per l'esquerra $\lambda_a\colon m\mapsto a\cdot m$, una per a cada $a\in M$. Com que $\lambda_a\comp\lambda_b=\lambda_{a\cdot b}$ i $\lambda_e=\id$, la immersió de Yoneda identifica $M$ amb aquest monoide d'aplicacions de $M$ en si mateix, i el fet que sigui fidel diu que $a\mapsto\lambda_a$ és injectiva. Obtenim el \textbf{teorema de Cayley} per a monoides: tot monoide és isomorf a un monoide d'aplicacions d'un conjunt en si mateix, amb la composició.
\end{solucio}

\chapter[Límits i colímits]{Límits i colímits}

Els exercicis d'aquest capítol treballen les propietats universals dels límits i colímits concrets, i l'argument ---sempre el mateix--- que dedueix unicitat i factoritzacions. Convé dibuixar cada diagrama abans de començar.

\section{Diagrames i cons}

\begin{exercici}
Comprova que un con sobre un diagrama $D\colon\cat{I}\to\cat{C}$ amb vèrtex $A$ és el mateix que una transformació natural $\Delta_A\Rightarrow D$, on $\Delta_A$ és el functor constant amb valor $A$.
\end{exercici}

\begin{solucio}
El functor constant $\Delta_A$ envia cada objecte $i$ de $\cat{I}$ a $A$ i cada morfisme a $\id_A$. Una transformació natural $\alpha\colon\Delta_A\Rightarrow D$ té components $\alpha_i\colon\Delta_A(i)=A\to D_i$, i la condició de naturalitat per a un morfisme $u\colon i\to j$ diu
\[
D_u\comp\alpha_i=\alpha_j\comp\Delta_A(u)=\alpha_j\comp\id_A=\alpha_j.
\]
Aquesta és exactament la condició de con. Recíprocament, tota família amb aquesta propietat defineix una transformació natural. Així, els cons sobre $D$ amb vèrtex $A$ i les transformacions naturals $\Delta_A\Rightarrow D$ coincideixen.
\end{solucio}

\section{Productes i equalitzadors}

\begin{exercici}
Comprova que a $\cat{Set}$ el producte categòric coincideix amb el producte cartesià, verificant la propietat universal.
\end{exercici}

\begin{solucio}
Sigui $A\times B=\{(a,b)\mid a\in A,\ b\in B\}$ amb $\pi_1(a,b)=a$ i $\pi_2(a,b)=b$. Donats $f\colon X\to A$ i $g\colon X\to B$, definim $\langle f,g\rangle\colon X\to A\times B$ per $\langle f,g\rangle(x)=(f(x),g(x))$. Aleshores $\pi_1\comp\langle f,g\rangle=f$ i $\pi_2\comp\langle f,g\rangle=g$. I és l'únic possible: si $h\colon X\to A\times B$ compleix $\pi_1\comp h=f$ i $\pi_2\comp h=g$, aleshores $h(x)=(\pi_1(h(x)),\pi_2(h(x)))=(f(x),g(x))=\langle f,g\rangle(x)$. Per tant el producte cartesià satisfà la propietat universal del producte.
\end{solucio}

\begin{exercici}
Comprova que tot equalitzador és un monomorfisme.
\end{exercici}

\begin{solucio}
Sigui $e\colon E\to A$ l'equalitzador de $f,g\colon A\to B$, i siguin $r,s\colon Z\to E$ amb $e\comp r=e\comp s$. Volem $r=s$. Sigui $h=e\comp r=e\comp s\colon Z\to A$. Aleshores $f\comp h=f\comp e\comp r=g\comp e\comp r=g\comp h$, de manera que $h$ satisfà la condició que defineix els cons de l'equalitzador. Per la propietat universal, hi ha un \emph{únic} $k\colon Z\to E$ amb $e\comp k=h$. Però tant $r$ com $s$ compleixen aquesta equació; per unicitat, $r=s$. Així $e$ és un monomorfisme.
\end{solucio}

\section{Pullbacks}

\begin{exercici}
Comprova que a $\cat{Set}$ el pullback de $f\colon A\to C$ i $g\colon B\to C$ és $P=\{(a,b)\in A\times B\mid f(a)=g(b)\}$, i que quan $g$ és la injecció d'un subconjunt $B\subseteq C$ el pullback recupera l'antiimatge $f^{-1}(B)$.
\end{exercici}

\begin{solucio}
Amb $P=\{(a,b)\mid f(a)=g(b)\}$ i les projeccions $p_1,p_2$, es compleix $f\comp p_1=g\comp p_2$ per construcció. Donats $r_1\colon X\to A$ i $r_2\colon X\to B$ amb $f\comp r_1=g\comp r_2$, l'aplicació $h(x)=(r_1(x),r_2(x))$ té imatge dins $P$ (perquè $f(r_1(x))=g(r_2(x))$) i és l'únic morfisme amb $p_1\comp h=r_1$, $p_2\comp h=r_2$. Això és la propietat universal.

Si $B\subseteq C$ i $g$ és la injecció, la condició $f(a)=g(b)=b$ força $b=f(a)$ amb $f(a)\in B$; per tant $P\cong\{a\in A\mid f(a)\in B\}=f^{-1}(B)$. El pullback recupera l'antiimatge.
\end{solucio}

\begin{exercici}[Lema d'enganxament de pullbacks]
Considera el diagrama de dos quadrats adossats en una categoria arbitrària:
\[
\begin{tikzpicture}[baseline=(m.center)]
  \matrix (m) [matrix of math nodes, column sep=3.2em, row sep=2.6em] {
    A & B & C \\
    A' & B' & C' \\
  };
  \draw[-{Stealth[scale=1.1]}] (m-1-1) -- node[above] {$u$} (m-1-2);
  \draw[-{Stealth[scale=1.1]}] (m-1-2) -- node[above] {$v$} (m-1-3);
  \draw[-{Stealth[scale=1.1]}] (m-2-1) -- node[below] {$u'$} (m-2-2);
  \draw[-{Stealth[scale=1.1]}] (m-2-2) -- node[below] {$v'$} (m-2-3);
  \draw[-{Stealth[scale=1.1]}] (m-1-1) -- node[left] {$a$} (m-2-1);
  \draw[-{Stealth[scale=1.1]}] (m-1-2) -- node[left] {$b$} (m-2-2);
  \draw[-{Stealth[scale=1.1]}] (m-1-3) -- node[right] {$c$} (m-2-3);
\end{tikzpicture}
\]
amb tots dos quadrats commutatius. Suposem que el \textbf{quadrat dret} (sobre $B,C,B',C'$) és un pullback. Comprova que el \textbf{quadrat esquerre} és un pullback si i només si el \textbf{rectangle exterior} (sobre $A,C,A',C'$) ho és.
\end{exercici}

\begin{solucio}
Treballem en una categoria arbitrària, usant només la propietat universal.

\textbf{($\Rightarrow$) Si el quadrat esquerre és pullback, l'exterior ho és.} Sigui $T$ un objecte amb morfismes $p\colon T\to A'$ i $q\colon T\to C$ compatibles per al rectangle exterior, és a dir amb $c\comp q=v'\comp u'\comp p$. Com que el quadrat dret és un pullback i els morfismes $q\colon T\to C$ i $u'\comp p\colon T\to B'$ satisfan $v'\comp(u'\comp p)=c\comp q$, hi ha un únic $r\colon T\to B$ amb $v\comp r=q$ i $b\comp r=u'\comp p$. Ara $r$ i $p$ són compatibles per al quadrat esquerre, perquè $b\comp r=u'\comp p$; com que aquest és un pullback, hi ha un únic $t\colon T\to A$ amb $u\comp t=r$ i $a\comp t=p$. Aleshores $(v\comp u)\comp t=v\comp r=q$ i $a\comp t=p$, de manera que $t$ factoritza el con exterior, i és únic per les unicitats successives.

\textbf{($\Leftarrow$) Si el rectangle exterior és pullback, l'esquerre ho és.} Sigui $T$ amb $p\colon T\to A'$ i $r\colon T\to B$ compatibles per al quadrat esquerre, és a dir amb $b\comp r=u'\comp p$. Posem $q=v\comp r\colon T\to C$. Aleshores $p$ i $q$ són compatibles per a l'exterior, perquè
\[
c\comp q=c\comp v\comp r=v'\comp b\comp r=v'\comp u'\comp p.
\]
Com que l'exterior és un pullback, hi ha un únic $t\colon T\to A$ amb $(v\comp u)\comp t=q$ i $a\comp t=p$. Queda veure que $u\comp t=r$: component amb $v$ i amb $b$,
\[
v\comp(u\comp t)=q=v\comp r,\qquad b\comp(u\comp t)=u'\comp a\comp t=u'\comp p=b\comp r,
\]
i com que el quadrat dret és un pullback, aquestes dues igualtats forcen $u\comp t=r$. Així $t$ és la factorització única per al quadrat esquerre.

En tots dos sentits, l'ingredient és la propietat universal del quadrat dret, que transfereix les factoritzacions entre el rectangle i el quadrat esquerre. Aquest \emph{lema d'enganxament} és una de les eines més usades per raonar amb pullbacks, i el retrobarem en tractar la reindexació entre subobjectes.
\end{solucio}

\begin{exercici}\label{pr-reindexacio-set}
Sigui $f\colon A\to B$ una aplicació i $y\colon Y\to B$ un objecte de $\cat{Set}/B$. Calcula explícitament la reindexació $f^{*}Y$ (observació~\ref{obs:reindexacio}) i comprova que la fibra de $f^{*}Y\to A$ sobre cada $a\in A$ està en bijecció canònica amb la fibra $y^{-1}(f(a))$.
\end{exercici}

\begin{solucio}
Pel càlcul del pullback a $\cat{Set}$ de l'exercici anterior,
\[
f^{*}Y=A\times_B Y=\{(a,z)\in A\times Y\mid f(a)=y(z)\},
\]
amb la projecció $p_1\colon(a,z)\mapsto a$ com a objecte de $\cat{Set}/A$. La fibra de $p_1$ sobre $a$ és
\[
p_1^{-1}(a)=\{(a,z)\mid y(z)=f(a)\},
\]
i l'aplicació $(a,z)\mapsto z$ en dona una bijecció canònica amb
\[
\{z\in Y\mid y(z)=f(a)\}=y^{-1}(f(a)).
\]
Així, si interpretem $y\colon Y\to B$ com la família de fibres $(Y_b)_{b\in B}$, amb $Y_b=y^{-1}(b)$, la reindexació al llarg de $f$ s'interpreta com la família $(Y_{f(a)})_{a\in A}$: cada índex $a$ rep la fibra de l'índex $f(a)$. És la \emph{substitució} al llarg de $f$, que reapareixerà com a operació sobre predicats en el capítol de subobjectes.
\end{solucio}

\begin{exercici}\label{pr-provt-limits-finits}
A la categoria prima $\cat{Prov}_T$, suposem que el sistema deductiu $T$ té una fórmula $\top$ i una conjunció $\land$ amb les regles habituals. Comprova que $\top$ és un objecte terminal i que $p\land q$ és un producte de $p$ i $q$. Dedueix que $\cat{Prov}_T$ té tots els límits finits.
\end{exercici}

\begin{solucio}
\textbf{Terminal.} Per a tota fórmula $r$ es té $r\vdash_T\top$, i com que $\cat{Prov}_T$ és prima, aquesta fletxa és única. Per tant $\top$ és terminal.

\textbf{Producte.} Les regles d'eliminació de la conjunció donen les projeccions
\[
p\land q\vdash_T p,
\qquad
p\land q\vdash_T q.
\]
Si $r\vdash_T p$ i $r\vdash_T q$, la regla d'introducció dona $r\vdash_T p\land q$, que és la fletxa mediadora; els triangles commuten automàticament i la fletxa és única, perquè la categoria és prima. Per tant $p\land q$, amb les dues eliminacions, és un producte de $p$ i $q$.

\textbf{Límits finits.} En una categoria prima, dues fletxes paral·leles coincideixen, de manera que l'equalitzador de $x\rightrightarrows y$ és $\id_x$ i sempre existeix. Pel teorema de caracterització dels límits finits (teorema~\ref{teo:limits-finits}), objecte terminal, productes binaris i equalitzadors donen tots els límits finits. És la versió resolta de la lectura lògica de l'exemple~\ref{ex:limits-preordre}.
\end{solucio}

\section{Colímits}

\begin{exercici}
Descriu el coproducte a $\cat{Set}$ i comprova'n la propietat universal.
\end{exercici}

\begin{solucio}
El coproducte de $A$ i $B$ a $\cat{Set}$ és la unió disjunta $A+B=(\{1\}\times A)\cup(\{2\}\times B)$, amb les injeccions $\iota_1(a)=(1,a)$ i $\iota_2(b)=(2,b)$. Donats $f\colon A\to X$ i $g\colon B\to X$, l'aplicació $[f,g]\colon A+B\to X$ definida per $[f,g](1,a)=f(a)$ i $[f,g](2,b)=g(b)$ és l'únic morfisme amb $[f,g]\comp\iota_1=f$ i $[f,g]\comp\iota_2=g$. Aquesta és la definició per casos, i és exactament la propietat universal del coproducte, dual de la del producte.
\end{solucio}

\begin{exercici}
Comprova que a $\cat{Set}$ el coequalitzador de $f,g\colon A\to B$ és el quocient $B/{\sim}$ per la relació d'equivalència generada per $f(a)\sim g(a)$.
\end{exercici}

\begin{solucio}
Sigui $\sim$ la menor relació d'equivalència amb $f(a)\sim g(a)$ per a tot $a$, i $q\colon B\to B/{\sim}$ la projecció al quocient. Aleshores $q\comp f=q\comp g$, perquè $f(a)$ i $g(a)$ tenen la mateixa classe. Si $s\colon B\to S$ compleix $s\comp f=s\comp g$, aleshores $s$ és constant sobre cada classe d'equivalència, de manera que factoritza de manera única com $s=r\comp q$ amb $r([b])=s(b)$ ben definida. Aquesta és la propietat universal del coequalitzador; el quocient hi encaixa exactament.
\end{solucio}

\section{Preservació i reflexió}

\begin{exercici}
Comprova que el functor oblit $U\colon\cat{Grp}\to\cat{Set}$ preserva els productes, però no preserva tots els coproductes.
\end{exercici}

\begin{solucio}
Per als productes: el grup producte $G\times H$ té com a conjunt subjacent el producte cartesià dels conjunts subjacents, i les projeccions són les aplicacions projecció; per tant $U(G\times H)=U(G)\times U(H)$ amb les projeccions correctes, i $U$ preserva el producte.

Per als coproductes: el coproducte a $\cat{Grp}$ és el \emph{producte lliure} $G*H$. N'hi ha prou d'un contraexemple, i el més senzill és prendre dos grups trivials $1=\{e\}$. El seu producte lliure és de nou trivial, $1*1\cong 1$, de manera que $U(1*1)$ té un sol element. En canvi, el coproducte a $\cat{Set}$ dels conjunts subjacents, $U(1)+U(1)$, és la unió disjunta de dos conjunts d'un element i en té dos. El cocon imatge, amb vèrtex $U(1*1)$, no pot ser, doncs, un coproducte a $\cat{Set}$, i $U$ no preserva aquest coproducte.

Aquesta asimetria anticipa el teorema del capítol següent: $U$ és un adjunt per la dreta, i els adjunts per la dreta preserven tots els límits; no hi ha cap afirmació anàloga que els obligui a preservar tots els colímits.
\end{solucio}

\begin{exercici}
Comprova que el functor representable $\Hom(A,-)\colon\cat{Set}\to\cat{Set}$ preserva productes, és a dir, $\Hom(A,X\times Y)\cong\Hom(A,X)\times\Hom(A,Y)$.
\end{exercici}

\begin{solucio}
Per la propietat universal del producte, l'aplicació $h\mapsto(\pi_1\comp h,\ \pi_2\comp h)$ és una bijecció
\[
\Hom(A,X\times Y)\cong\Hom(A,X)\times\Hom(A,Y),
\]
natural en $X$ i $Y$ (i també en $A$, si es considera la variació contravariant en aquesta variable). Aquesta bijecció és exactament la induïda pel con imatge $\bigl(\Hom(A,X\times Y),\Hom(A,\pi_1),\Hom(A,\pi_2)\bigr)$, de manera que aquest con és un producte a $\cat{Set}$ i $\Hom(A,-)$ preserva el producte binari. El mateix argument, amb la propietat universal general del límit, mostra que $\Hom(A,-)$ preserva tots els límits (proposició~\ref{prop:representables-limits}).
\end{solucio}

\chapter[Adjuncions]{Adjuncions}

Els exercicis d'aquest capítol treballen les adjuncions des dels tres punts de vista equivalents ---bijecció d'homconjunts, unitat/counitat, identitats triangulars--- i n'apliquen les conseqüències sobre límits. Convé identificar sempre qui és l'adjunt per l'esquerra i qui per la dreta.

\section{Adjuncions entre preordres}

\begin{exercici}
Comprova que al preordre $\cat{Prov}_T$ d'un sistema proposicional amb les regles usuals de $\wedge$ i $\to$, fixada una fórmula $A$, l'operació $B\mapsto A\wedge B$ és adjunta per l'esquerra de $B\mapsto A\to B$.
\end{exercici}

\begin{solucio}
Hem de veure que $A\wedge B\vdash_T C$ si i només si $B\vdash_T A\to C$. 

Suposem $A\wedge B\vdash_T C$. Aleshores, assumint $B$, volem deduir $A\to C$: assumim també $A$; tenim $A$ i $B$, per tant $A\wedge B$, d'on $C$; descarregant la hipòtesi $A$ obtenim $A\to C$. Així $B\vdash_T A\to C$.

Recíprocament, suposem $B\vdash_T A\to C$. Assumint $A\wedge B$, en tenim $A$ i $B$; de $B$ obtenim $A\to C$, i amb $A$, per modus ponens, $C$. Així $A\wedge B\vdash_T C$.

Les dues implicacions donen l'equivalència, és a dir, $(A\wedge-)\dashv(A\to-)$. Aquesta adjunció és la lectura categòrica de la regla de la implicació.
\end{solucio}

\begin{exercici}\label{pr-quantificadors-calcul}
Sigui $f\colon X\to Y$ l'aplicació de $X=\{1,2,3,4\}$ en $Y=\{a,b,c\}$ amb $f(1)=f(2)=a$, $f(3)=b$ i $f(4)=b$, i sigui $S=\{1,2,3\}$. Calcula $\exists_f(S)$ i $\forall_f(S)$ (exemple~\ref{ex:quantificadors-set}). Després, fent servir només la cadena $\exists_f\dashv f^{*}\dashv\forall_f$, justifica que $\exists_f$ preserva unions, $\forall_f$ preserva interseccions i $f^{*}$ preserva totes dues.
\end{exercici}

\begin{solucio}
\textbf{Càlcul.} La imatge és $\exists_f(S)=\{a,b\}$. Per a $\forall_f(S)$, mirem les fibres: $f^{-1}(\{a\})=\{1,2\}\subseteq S$, de manera que $a\in\forall_f(S)$; $f^{-1}(\{b\})=\{3,4\}\not\subseteq S$, perquè $4\notin S$; i $f^{-1}(\{c\})=\varnothing\subseteq S$. Per tant
\[
\forall_f(S)=\{a,c\}.
\]
Observem dos fenòmens: $c$ hi és perquè la seva fibra és buida (la quantificació universal sobre un conjunt buit és certa), i $b$ no hi és tot i que $3\in S$.

\textbf{Preservació.} En un preordre de parts, les unions són els coproductes (suprems) i les interseccions, els productes (ínfims). Pel teorema que els adjunts per l'esquerra preserven colímits i els adjunts per la dreta preserven límits:
\begin{itemize}
  \item $\exists_f$, adjunt per l'esquerra de $f^{*}$, preserva unions: $f[S\cup S']=f[S]\cup f[S']$;
  \item $\forall_f$, adjunt per la dreta de $f^{*}$, preserva interseccions;
  \item $f^{*}$, que és adjunt per la dreta de $\exists_f$ i adjunt per l'esquerra de $\forall_f$, preserva totes dues.
\end{itemize}
En canvi, $\exists_f$ no preserva en general les interseccions: amb l'aplicació de l'enunciat, $\exists_f(\{1\}\cap\{2\})=\varnothing$, però $\exists_f(\{1\})\cap\exists_f(\{2\})=\{a\}$. Lògicament: l'existencial commuta amb la disjunció, però no amb la conjunció.
\end{solucio}

\section{Adjuncions entre categories}

\begin{exercici}\label{pr-disc-ob}
Sigui $\mathrm{Ob}\colon\cat{Cat}\to\cat{Set}$ el functor que envia cada categoria petita al seu conjunt d'objectes, i $\mathrm{Disc}\colon\cat{Set}\to\cat{Cat}$ el que envia un conjunt $X$ a la categoria discreta amb objectes $X$. Comprova que $\mathrm{Disc}\dashv\mathrm{Ob}$ i descriu-ne la unitat i la counitat.
\end{exercici}

\begin{solucio}
Un functor $H\colon\mathrm{Disc}(X)\to\cat{D}$ ha d'enviar cada identitat a una identitat, i $\mathrm{Disc}(X)$ no té cap altre morfisme; queda, doncs, determinat per una aplicació $X\to\mathrm{Ob}(\cat{D})$ entre els objectes, i qualsevol aplicació dona un functor. Això és una bijecció
\[
\Hom_{\cat{Cat}}\bigl(\mathrm{Disc}(X),\cat{D}\bigr)\;\cong\;\Hom_{\cat{Set}}\bigl(X,\mathrm{Ob}(\cat{D})\bigr),
\]
natural en $X$ i en $\cat{D}$, perquè les dues bandes es componen de la mateixa manera sobre els objectes. Per tant $\mathrm{Disc}\dashv\mathrm{Ob}$.

La \emph{unitat} $\eta_X\colon X\to\mathrm{Ob}(\mathrm{Disc}(X))=X$ és la identitat: és un isomorfisme. La \emph{counitat} $\varepsilon_{\cat{D}}\colon\mathrm{Disc}(\mathrm{Ob}(\cat{D}))\to\cat{D}$ és el functor que inclou els objectes de $\cat{D}$ amb les seves identitats, i en general no és un isomorfisme, perquè oblida tots els morfismes no identitat.

\emph{Observació.} $\mathrm{Ob}$ té també un adjunt per la dreta: el functor que envia $X$ a la categoria \emph{indiscreta}, amb exactament un morfisme entre cada parella d'objectes. Per tant $\mathrm{Ob}$ preserva límits i colímits, d'acord amb el fet que els límits i colímits de categories petites es calculen, sobre els objectes, com a $\cat{Set}$.
\end{solucio}

\section{Unitat i counitat}

\begin{exercici}
Per a l'adjunció lliure-oblit $F\dashv U$ dels monoides, descriu la unitat $\eta_X\colon X\to U(X^*)$ i la counitat $\varepsilon_M\colon (U(M))^*\to M$, i comprova que la unitat expressa la propietat universal del monoide lliure.
\end{exercici}

\begin{solucio}
La unitat $\eta_X\colon X\to U(X^*)$ envia cada element $x$ a la paraula d'una sola lletra $x$. La counitat $\varepsilon_M\colon (U(M))^*\to M$ envia una paraula d'elements de $M$ al seu producte a $M$: $\varepsilon_M(m_1\cdots m_n)=m_1\cdots m_n$ (producte calculat a $M$).

La propietat universal de la unitat diu que per a tota aplicació $\varphi\colon X\to U(M)$ hi ha un únic homomorfisme $\bar\varphi\colon X^*\to M$ amb $U(\bar\varphi)\comp\eta_X=\varphi$. Això és exactament la definició del monoide lliure: un homomorfisme des de $X^*$ queda determinat, de manera única, pels valors sobre els generadors $X$.
\end{solucio}

\begin{exercici}
Comprova que, en una adjunció $F\dashv G$, el transposat de $\varphi\colon F(A)\to B$ és $G(\varphi)\comp\eta_A$, i el transposat invers de $\psi\colon A\to G(B)$ és $\varepsilon_B\comp F(\psi)$.
\end{exercici}

\begin{solucio}
Per naturalitat de la bijecció $\theta$ en la segona variable, aplicada al morfisme $\varphi\colon F(A)\to B$ i partint de $\id_{F(A)}$,
\[
\theta(\varphi)=\theta(\varphi\comp\id_{F(A)})=G(\varphi)\comp\theta(\id_{F(A)})=G(\varphi)\comp\eta_A,
\]
on hem usat que $\theta(\id_{F(A)})=\eta_A$ per definició de la unitat. Dualment, per naturalitat en la primera variable i la definició de la counitat $\varepsilon_B=\theta^{-1}(\id_{G(B)})$,
\[
\theta^{-1}(\psi)=\theta^{-1}(\id_{G(B)}\comp\psi)=\theta^{-1}(\id_{G(B)})\comp F(\psi)=\varepsilon_B\comp F(\psi).
\]
Aquestes fórmules permeten passar de la bijecció a la unitat/counitat i viceversa.
\end{solucio}

\section{Identitats triangulars}

\begin{exercici}\label{pr-diagonal-producte}
Sigui $\cat{C}$ una categoria amb productes binaris, i siguin $\Delta\colon\cat{C}\to\cat{C}\times\cat{C}$ el functor diagonal, $\Delta(A)=(A,A)$, i $\times\colon\cat{C}\times\cat{C}\to\cat{C}$ el functor producte, $(X,Y)\mapsto X\times Y$. Comprova que $\Delta\dashv\times$, descriu-ne la unitat i la counitat, i verifica directament les dues identitats triangulars.
\end{exercici}

\begin{solucio}
\textbf{L'adjunció.} Un morfisme $\Delta(A)\to(X,Y)$ a $\cat{C}\times\cat{C}$ és una parella $(f\colon A\to X,\ g\colon A\to Y)$, i la propietat universal del producte la posa en bijecció amb els morfismes $\langle f,g\rangle\colon A\to X\times Y$:
\[
\Hom_{\cat{C}\times\cat{C}}\bigl((A,A),(X,Y)\bigr)\;\cong\;\Hom_{\cat{C}}(A,X\times Y).
\]
La bijecció és natural en $A$ i en $(X,Y)$, perquè $\langle f,g\rangle\comp a=\langle f\comp a,g\comp a\rangle$ i $(u\times v)\comp\langle f,g\rangle=\langle u\comp f,v\comp g\rangle$. Per tant $\Delta\dashv\times$.

\textbf{Unitat i counitat.} La unitat és el transposat de $\id_{(A,A)}$, és a dir, el morfisme diagonal
\[
\eta_A=\delta_A=\langle\id_A,\id_A\rangle\colon A\to A\times A,
\]
i la counitat és el transposat invers de $\id_{X\times Y}$, la parella de projeccions
\[
\varepsilon_{(X,Y)}=(\pi_1,\pi_2)\colon(X\times Y,\ X\times Y)\to(X,Y).
\]

\textbf{Primera identitat.} $\varepsilon_{\Delta(A)}\comp\Delta(\eta_A)$ és la parella
\[
(\pi_1\comp\delta_A,\ \pi_2\comp\delta_A)=(\id_A,\id_A)=\id_{\Delta(A)}:
\]
copiar $A$ i projectar a cada còpia no canvia res.

\textbf{Segona identitat.} $(\pi_1\times\pi_2)\comp\delta_{X\times Y}$ és el morfisme $X\times Y\to X\times Y$ de components $\pi_1$ i $\pi_2$, és a dir,
\[
\langle\pi_1,\pi_2\rangle=\id_{X\times Y}:
\]
duplicar una parella i quedar-se amb la primera component de la primera còpia i la segona de la segona retorna la parella original.

Observem que la unitat $\delta_A$ no és un isomorfisme (a $\cat{Set}$, només ho és si $A$ té com a màxim un element): les identitats triangulars no demanen que ho sigui.
\end{solucio}

\begin{exercici}\label{pr-free-u-triangulars}
Per a l'adjunció $\Free\dashv U$ entre $\cat{Graph}$ i $\cat{Cat}$ (exemple~\ref{ex:free-u-adjuncio}), comprova directament les dues identitats triangulars
\[
\varepsilon_{\Free(G)}\comp\Free(\eta_G)=\id_{\Free(G)},
\qquad
U(\varepsilon_{\cat{C}})\comp\eta_{U(\cat{C})}=\id_{U(\cat{C})},
\]
i explica què diu cadascuna en termes de camins.
\end{exercici}

\begin{solucio}
\textbf{Primera identitat.} El functor $\Free(\eta_G)\colon\Free(G)\to\Free(U(\Free(G)))$ envia un camí $(a_1,\dots,a_n)$ de $G$ al camí formal $\bigl((a_1),\dots,(a_n)\bigr)$, les arestes del qual són camins de longitud~$1$ de $\Free(G)$. La counitat $\varepsilon_{\Free(G)}$ avalua aquest camí formal component les seves arestes a $\Free(G)$, és a dir, concatenant:
\[
\varepsilon_{\Free(G)}\Bigl(\bigl((a_1),\dots,(a_n)\bigr)\Bigr)=(a_n)\comp\cdots\comp(a_1)=(a_1,\dots,a_n).
\]
El camí buit va al camí buit. Per tant, la composta és la identitat de $\Free(G)$: \emph{trencar un camí en arestes i tornar-lo a compondre no canvia res}.

\textbf{Segona identitat.} La unitat $\eta_{U(\cat{C})}$ envia una fletxa $f$ de $\cat{C}$, vista com a aresta del graf $U(\cat{C})$, al camí $(f)$ de longitud~$1$, i cada objecte a ell mateix. El morfisme de grafs $U(\varepsilon_{\cat{C}})$ avalua aquest camí, que té una sola fletxa, i en dona $f$. Per tant, la composta és la identitat de $U(\cat{C})$: \emph{avaluar el camí d'una sola fletxa retorna la fletxa}.

Les dues identitats expressen, doncs, que la inserció dels generadors i l'avaluació de camins són inverses en el sentit precís que demanen les identitats triangulars, sense que $\eta$ ni $\varepsilon$ siguin isomorfismes: $\Free(U(\cat{C}))$ té en general molts més morfismes que $\cat{C}$, perquè hi ha camins diferents amb la mateixa composta.
\end{solucio}

\section{Els adjunts i els límits}

\begin{exercici}\label{pr-rapl-mon}
Comprova, fent servir que els adjunts per la dreta preserven límits, que el functor oblit $U\colon\cat{Mon}\to\cat{Set}$ preserva productes. Comprova també que $U$ no preserva tots els coproductes, i explica què en dedueixes.
\end{exercici}

\begin{solucio}
El functor monoide lliure $F\colon\cat{Set}\to\cat{Mon}$, $X\mapsto X^*$, és adjunt per l'esquerra de $U$ (exemple del monoide lliure de la Teoria). Per tant $U$ és un adjunt per la dreta, i pel teorema \enquote{els adjunts per la dreta preserven límits}, $U$ preserva tots els límits petits, en particular els productes: el conjunt subjacent d'un producte de monoides és el producte dels conjunts subjacents.

El teorema, però, no diu res dels colímits, i de fet $U$ no preserva tots els coproductes. Sigui $1=\{e\}$ el monoide trivial. És un objecte inicial de $\cat{Mon}$, perquè l'únic homomorfisme $1\to M$ envia $e$ al neutre, i el coproducte d'un objecte amb un objecte inicial és l'objecte mateix; per tant, el coproducte $1+1$ a $\cat{Mon}$ és $1$, amb un sol element. En canvi, el coproducte a $\cat{Set}$ dels conjunts subjacents, $U(1)+U(1)$, en té dos, i el cocon imatge no és un coproducte. Com a conseqüència, $U$ no pot ser un adjunt per l'esquerra, perquè els adjunts per l'esquerra preserven colímits.

\emph{Ampliació, per a qui conegui els grups.} El mateix passa amb $U\colon\cat{Grp}\to\cat{Set}$. El seu adjunt per l'esquerra és el functor \emph{grup lliure}, que envia un conjunt $X$ al grup de les paraules reduïdes en les lletres de $X$ i els seus inversos formals; per tant, $U$ preserva tots els límits. El grup trivial també és inicial a $\cat{Grp}$, el coproducte (el \emph{producte lliure}) de dos grups trivials és trivial, $1*1\cong 1$, i $U$ tampoc no preserva aquest coproducte.
\end{solucio}

\begin{exercici}
Comprova que si un functor $G$ \emph{no} preserva algun límit petit, aleshores $G$ no pot tenir cap adjunt per l'esquerra.
\end{exercici}

\begin{solucio}
És el contrarecíproc del teorema. Si $G$ tingués un adjunt per l'esquerra $F$, amb $F\dashv G$, aleshores $G$ seria un adjunt per la dreta i, pel teorema, preservaria tots els límits petits que existeixen. Com que per hipòtesi $G$ no en preserva algun, no pot tenir adjunt per l'esquerra. Aquest és un criteri negatiu molt pràctic: per exemple, si un functor oblit no preserva algun coproducte, això impedeix que sigui adjunt per l'esquerra, però no impedeix que sigui adjunt per la dreta.
\end{solucio}

\section{Exponencials}

\begin{exercici}
Comprova la bijecció
\[
\Hom_{\cat{Set}}(A\times B,C)\cong\Hom_{\cat{Set}}(A,C^B)
\]
i identifica-la amb la currificació.
\end{exercici}

\begin{solucio}
A $\cat{Set}$, $C^B$ és el conjunt de les aplicacions $B\to C$. Donada $f\colon A\times B\to C$, definim $\tilde f\colon A\to C^B$ per $\tilde f(a)=(b\mapsto f(a,b))$; és a dir, $\tilde f(a)$ és la funció que fixa la primera coordenada en $a$. Recíprocament, donada $g\colon A\to C^B$, definim $\hat g\colon A\times B\to C$ per $\hat g(a,b)=g(a)(b)$. Aquestes assignacions són inverses l'una de l'altra i naturals en totes les variables. La correspondència $f\mapsto\tilde f$ és exactament la \textbf{currificació}: transformar una funció de dues variables en una funció d'una variable que retorna funcions. Aquesta bijecció expressa l'adjunció $-\times B\dashv(-)^B$.
\end{solucio}

\chapter[Cartesianes tancades]{\texorpdfstring{Categories cartesianes\\ tancades}{Categories cartesianes tancades}}

Els exercicis treballen tres aspectes del tancament cartesià: productes i diagonals, la propietat universal de l'exponencial i la interpretació de tipus i lògica. Convé tenir sempre present la bijecció
\[
\Hom(A\times B,C)\cong\Hom(A,C^{B}).
\]

\section{Productes i diagonals}

\begin{exercici}
En una categoria cartesiana, comprova que la diagonal $\delta_A\colon A\to A\times A$ satisfà
\[
\pi_1\comp\delta_A=\pi_2\comp\delta_A=\id_A
\]
i que és l'únic morfisme amb aquesta propietat.
\end{exercici}

\begin{solucio}
Per la propietat universal del producte $A\times A$, un morfisme $A\to A\times A$ està determinat per les seves dues composicions amb les projeccions. Definim $\delta_A$ com l'únic morfisme amb $\pi_1\comp\delta_A=\id_A$ i $\pi_2\comp\delta_A=\id_A$; l'existència i la unicitat són exactament el contingut de la propietat universal aplicada a la parella $(\id_A,\id_A)$. Qualsevol altre morfisme $d\colon A\to A\times A$ amb $\pi_1\comp d=\pi_2\comp d=\id_A$ compleix la mateixa condició universal i, per unicitat, $d=\delta_A$.
\end{solucio}

\section{La propietat universal de l'exponencial}

\begin{exercici}
A $\cat{Set}$, descriu explícitament l'avaluació, la currificació i la descurrificació per a l'exponencial $C^{B}$, i verifica l'equació $\ev\comp(\lambda f\times\id_B)=f$.
\end{exercici}

\begin{solucio}
A $\cat{Set}$, $C^{B}$ és el conjunt d'aplicacions $B\to C$. L'avaluació és $\ev\colon C^{B}\times B\to C$, $\ev(g,b)=g(b)$. Donada $f\colon A\times B\to C$, la currificació és $\lambda f\colon A\to C^{B}$ amb $\lambda f(a)=(b\mapsto f(a,b))$. Comprovem l'equació: per a $(a,b)\in A\times B$,
\[
\bigl(\ev\comp(\lambda f\times\id_B)\bigr)(a,b)=\ev\bigl(\lambda f(a),b\bigr)=\lambda f(a)(b)=f(a,b).
\]
Per tant $\ev\comp(\lambda f\times\id_B)=f$. La descurrificació de $g\colon A\to C^{B}$ és $\widetilde g(a,b)=g(a)(b)$, i és inversa de la currificació.
\end{solucio}

\begin{exercici}
Demostra que en qualsevol CCC hi ha un isomorfisme natural $C^{1}\cong C$ i que $1^{B}\cong 1$.
\end{exercici}

\begin{solucio}
Per a $C^{1}$: per la propietat universal, $\Hom(A,C^{1})\cong\Hom(A\times 1,C)$. Com que $A\times 1\cong A$ (l'objecte terminal és neutre per al producte), $\Hom(A\times 1,C)\cong\Hom(A,C)$, natural en $A$. Pel lema de Yoneda, $C^{1}\cong C$.

Per a $1^{B}$: $\Hom(A,1^{B})\cong\Hom(A\times B,1)$, que té exactament un element per ser $1$ terminal. Per tant $\Hom(A,1^{B})$ és un conjunt d'un punt per a tot $A$, natural en $A$, i pel lema de Yoneda $1^{B}\cong 1$. Aquestes són les versions categòriques de les identitats $c^{1}=c$ i $1^{b}=1$.
\end{solucio}

\begin{exercici}
Comprova la \enquote{llei exponencial} $C^{A\times B}\cong(C^{B})^{A}$ en una CCC.
\end{exercici}

\begin{solucio}
Fem servir la propietat universal repetidament i el lema de Yoneda. Per a qualsevol $X$,
\[
\Hom(X,C^{A\times B})\cong\Hom(X\times(A\times B),C)\cong\Hom((X\times A)\times B,C),
\]
usant l'associativitat del producte. Aplicant ara l'adjunció respecte de $B$ i després respecte de $A$,
\[
\Hom((X\times A)\times B,C)\cong\Hom(X\times A,C^{B})\cong\Hom(X,(C^{B})^{A}).
\]
Totes les bijeccions són naturals en $X$; pel lema de Yoneda, $C^{A\times B}\cong(C^{B})^{A}$. És la llei $c^{a\cdot b}=(c^{b})^{a}$, i correspon, en el càlcul lambda, al fet que una funció de dos arguments és una funció que retorna una funció (currificació iterada).
\end{solucio}

\begin{exercici}\label{pr-distributivitat}
Sigui $\cat{C}$ una CCC amb coproductes binaris i objecte inicial $0$. Demostra que el producte es distribueix sobre el coproducte:
\[
A\times(B+C)\;\cong\;(A\times B)+(A\times C),
\qquad
A\times 0\;\cong\;0.
\]
Interpreta el resultat a $\cat{Set}$ i en una àlgebra de Heyting.
\end{exercici}

\begin{solucio}
Pel tancament cartesià, $-\times A\dashv(-)^{A}$, i el producte és commutatiu llevat d'isomorfisme natural; per tant $A\times-$ també és un adjunt per l'esquerra. Els adjunts per l'esquerra preserven colímits, en particular els coproductes binaris i l'objecte inicial, que és el colímit del diagrama buit. Explícitament, el morfisme canònic
\[
[\id_A\times\iota_1,\ \id_A\times\iota_2]\colon(A\times B)+(A\times C)\longrightarrow A\times(B+C)
\]
és un isomorfisme, i $A\times 0$ és inicial, de manera que $A\times 0\cong 0$.

\emph{A $\cat{Set}$}, la primera bijecció envia una parella $(a,x)$ amb $x$ a la component $B$ o a la component $C$ de la unió disjunta al lloc corresponent; la segona diu que $A\times\varnothing=\varnothing$.

\emph{En una àlgebra de Heyting}, vista com a preordre cartesià tancat amb suprems finits, el coproducte és el suprem i l'inicial és $\bot$. Obtenim
\[
a\wedge(b\vee c)=(a\wedge b)\vee(a\wedge c),
\qquad
a\wedge\bot=\bot:
\]
tota àlgebra de Heyting és un reticle distributiu, i la distributivitat és conseqüència de l'existència de la implicació. Lògicament, la conjunció es distribueix sobre la disjunció.
\end{solucio}

\section{Heyting i lògica}

\begin{exercici}
Comprova que en una àlgebra de Heyting l'exponencial $b\Rightarrow c$ satisfà $a\wedge b\leq c\Leftrightarrow a\leq(b\Rightarrow c)$, i deriva-hi que $b\wedge(b\Rightarrow c)\leq c$ (\emph{modus ponens}).
\end{exercici}

\begin{solucio}
La condició $a\wedge b\leq c\Leftrightarrow a\leq(b\Rightarrow c)$ és, precisament, l'adjunció $-\wedge b\dashv(b\Rightarrow-)$ escrita en el llenguatge dels preordres: és la definició d'exponencial en un preordre cartesià, on $\Hom$ té com a molt un element i la bijecció d'homconjunts es redueix a una equivalència de desigualtats. Per obtenir el \emph{modus ponens}, prenem $a=(b\Rightarrow c)$. Aleshores el membre dret $a\leq(b\Rightarrow c)$ es compleix trivialment (és $b\Rightarrow c\leq b\Rightarrow c$), de manera que el membre esquerre també: $(b\Rightarrow c)\wedge b\leq c$. Aquesta és la counitat de l'adjunció, i és la forma algebraica de la regla que de $b$ i $b\to c$ se'n segueix $c$.
\end{solucio}

\begin{exercici}
En una àlgebra de Heyting, demostra que $a\leq(b\Rightarrow a)$ per a tot $a,b$, i interpreta-ho lògicament.
\end{exercici}

\begin{solucio}
Per l'adjunció, $a\leq(b\Rightarrow a)$ equival a $a\wedge b\leq a$, que és cert perquè $a\wedge b$ és un ínfim i, per tant, $a\wedge b\leq a$. Lògicament, correspon a l'axioma $a\to(b\to a)$: una proposició verdadera se segueix de qualsevol hipòtesi. És l'esquelet del combinador $\mathsf{K}$ del càlcul lambda, la funció constant.
\end{solucio}

\section{Semàntica de termes}

\begin{exercici}
Interpreta el terme $x:\sigma,\,y:\tau\vdash\langle y,x\rangle:\tau\times\sigma$ com a morfisme en una CCC i identifica'l amb el morfisme de commutació del producte.
\end{exercici}

\begin{solucio}
El context $x:\sigma,y:\tau$ s'interpreta com l'objecte $\llbracket\sigma\rrbracket\times\llbracket\tau\rrbracket$. Les variables $x$ i $y$ són les projeccions $\pi_1$ i $\pi_2$. La parella $\langle y,x\rangle$ s'interpreta amb l'aparellament de les interpretacions de $y$ i $x$, és a dir,
\[
\llbracket\langle y,x\rangle\rrbracket=\langle\pi_2,\pi_1\rangle\colon\llbracket\sigma\rrbracket\times\llbracket\tau\rrbracket\to\llbracket\tau\rrbracket\times\llbracket\sigma\rrbracket.
\]
Aquest és exactament el morfisme de \textbf{commutació} (o \emph{swap}) del producte, que intercanvia els dos factors. La seva inversa és $\langle\pi_2,\pi_1\rangle$ en sentit contrari, cosa que expressa que $\sigma\times\tau\cong\tau\times\sigma$.
\end{solucio}

\begin{exercici}
Interpreta l'abstracció $x:\sigma\vdash(\lambda y{:}\tau.\,x):\tau\to\sigma$ i comprova que és la currificació d'una projecció.
\end{exercici}

\begin{solucio}
El cos del terme, $x:\sigma,\,y:\tau\vdash x:\sigma$, s'interpreta amb la projecció $\pi_1\colon\llbracket\sigma\rrbracket\times\llbracket\tau\rrbracket\to\llbracket\sigma\rrbracket$. L'abstracció respecte de $y$ es tradueix per la currificació:
\[
\llbracket\lambda y.\,x\rrbracket=\lambda(\pi_1)\colon\llbracket\sigma\rrbracket\to\llbracket\sigma\rrbracket^{\llbracket\tau\rrbracket}.
\]
És a dir, la funció que envia cada $x$ a la funció constant de valor $x$: el combinador $\mathsf{K}$, la versió proposicional del qual és la desigualtat $a\leq(b\Rightarrow a)$ de l'exercici de Heyting anterior.
\end{solucio}

\chapter[Subobjectes, imatges i factoritzacions]{\texorpdfstring{Subobjectes, imatges\\ i factoritzacions}{Subobjectes, imatges i factoritzacions}}

Els exercicis d'aquest capítol fixen les tres idees centrals: els subobjectes com a classes de monomorfismes, l'ordre $\Sub(A)$ com a àlgebra de predicats, i la reindexació $f^{*}$ com a substitució. Molts arguments es fan \enquote{sense elements}, usant només propietats universals; convé acostumar-s'hi, perquè és el que permet traslladar-los a categories arbitràries.

\section{Subobjectes i el seu ordre}

\begin{exercici}
Comprova que a $\cat{Set}$ dos monomorfismes $m\colon S\rightarrowtail A$ i $n\colon T\rightarrowtail A$ són equivalents si i només si tenen la mateixa imatge, i dedueix que $\Sub(A)\cong\mathcal{P}(A)$.
\end{exercici}

\begin{solucio}
Si $m\sim n$, hi ha un isomorfisme $\varphi\colon S\to T$ amb $n\comp\varphi=m$; aleshores $m(S)=n(\varphi(S))=n(T)$, de manera que tenen la mateixa imatge. Recíprocament, suposem $m(S)=n(T)=:I$. Siguin $\bar m\colon S\to I$ i $\bar n\colon T\to I$ les \emph{corestriccions} de $m$ i $n$ a $I$, és a dir, les mateixes aplicacions amb codomini $I$: $\bar m(s)=m(s)$ i $\bar n(t)=n(t)$. Com que $m$ i $n$ són injectives i tenen imatge $I$, $\bar m$ i $\bar n$ són bijeccions. La composició $\varphi=\bar n^{-1}\comp\bar m\colon S\to T$ és, doncs, una bijecció, i compleix $n\comp\varphi=m$, perquè per a tot $s\in S$, $n(\varphi(s))=\bar n(\bar n^{-1}(\bar m(s)))=m(s)$. Així $m\sim n$. Per tant la imatge determina la classe, i l'assignació $[m]\mapsto m(S)$ és una bijecció entre $\Sub(A)$ i el conjunt de subconjunts $\mathcal{P}(A)$; a més respecta l'ordre, ja que $[m]\leq[n]$ equival a $m(S)\subseteq n(T)$.
\end{solucio}

\begin{exercici}
Demostra que en qualsevol categoria l'objecte $\id_A\colon A\to A$ és el subobjecte \textbf{màxim} de $\Sub(A)$.
\end{exercici}

\begin{solucio}
La identitat és un monomorfisme (és un isomorfisme). Per a qualsevol altre subobjecte $m\colon S\rightarrowtail A$, el mateix $m$ és una factorització de $m$ a través de $\id_A$, ja que $\id_A\comp m=m$. Per tant $[m]\leq[\id_A]$ per a tot $m$, cosa que fa de $[\id_A]$ el màxim. Lògicament, és el predicat \enquote{sempre cert}: tot element de $A$ el satisfà.
\end{solucio}

\begin{exercici}
Comprova que si $m\colon S\rightarrowtail A$ és un subobjecte i $\varphi\colon S'\to S$ és un isomorfisme, aleshores $m\comp\varphi$ representa el mateix subobjecte. Per què això justifica treballar amb classes i no amb monos concrets?
\end{exercici}

\begin{solucio}
Per definició de la relació d'equivalència, $m\comp\varphi\sim m$ mitjançant l'isomorfisme $\varphi$: $m\comp\varphi=m\comp\varphi$ amb $m\comp\varphi$ factoritzant. Formalment, $[m\comp\varphi]=[m]$ perquè $\varphi$ és un isomorfisme del domini que fa commutar el triangle. Això justifica treballar amb classes: la teoria dels subobjectes ha de ser invariant per reparametritzacions del domini, ja que \enquote{el subconjunt} no depèn de com l'indexem. Prendre classes elimina aquesta ambigüitat sense perdre informació geomètrica.
\end{solucio}

\section{Reindexació}

\begin{exercici}
Sigui $f\colon A\to B$ a $\cat{Set}$ i $S\subseteq B$. Comprova que $f^{*}(S)=f^{-1}(S)$ i que $f^{*}$ preserva interseccions: $f^{*}(S\cap S')=f^{*}(S)\cap f^{*}(S')$.
\end{exercici}

\begin{solucio}
El pullback de la injecció $S\hookrightarrow B$ al llarg de $f$ és $\{(a,s)\in A\times S\mid f(a)=s\}$, que s'identifica amb $\{a\in A\mid f(a)\in S\}=f^{-1}(S)$, amb la projecció cap a $A$ com a injecció. Per a les interseccions, $x\in f^{-1}(S\cap S')$ si i només si $f(x)\in S\cap S'$, és a dir, $f(x)\in S$ i $f(x)\in S'$, això és $x\in f^{-1}(S)\cap f^{-1}(S')$. La reindexació preserva, doncs, la conjunció; és una manifestació que la substitució commuta amb $\wedge$.
\end{solucio}


\begin{exercici}
Comprova que si $f\colon A\to B$ és un epimorfisme escindit (té secció $s$ amb $f\comp s=\id_B$), aleshores $f^{*}$ és injectiva sobre subobjectes.
\end{exercici}

\begin{solucio}
Si $f\comp s=\id_B$, aplicant la functorialitat de la reindexació (proposició~\ref{prop:functor-sub}), $s^{*}\comp f^{*}=(f\comp s)^{*}=(\id_B)^{*}=\id_{\Sub(B)}$. Per tant $f^{*}$ té inversa per l'esquerra $s^{*}$, cosa que la fa injectiva: si $f^{*}[m]=f^{*}[n]$, aplicant $s^{*}$ obtenim $[m]=[n]$. Com que $f$ admet una secció, els subobjectes de $B$ es poden recuperar després de reindexar-los al llarg de $f$; per això $f^{*}$ és injectiva. (No vol dir que $f$ no identifiqui elements: una epi escindida pot identificar-ne molts.)
\end{solucio}

\section{Interseccions i lògica}

\begin{exercici}
Verifica amb detall que el pullback de dos subobjectes $m\colon S\rightarrowtail A$ i $n\colon T\rightarrowtail A$ és el seu ínfim a $\Sub(A)$, escrivint les dues direccions de la propietat universal.
\end{exercici}

\begin{solucio}
Sigui $P=S\times_A T$ el pullback, amb projeccions $p\colon P\to S$ i $q\colon P\to T$ i injecció $m\comp p=n\comp q\colon P\rightarrowtail A$ (mono per ser composició de monos). \emph{Fita inferior}: $[P]\leq[m]$ via $p$, ja que $m\comp p$ és la injecció i factoritza per $m$; anàlogament $[P]\leq[n]$ via $q$. \emph{Màxima fita inferior}: sigui $k\colon R\rightarrowtail A$ amb $[k]\leq[m]$ i $[k]\leq[n]$, amb factoritzacions $a\colon R\to S$ ($m\comp a=k$) i $b\colon R\to T$ ($n\comp b=k$). Aleshores $m\comp a=k=n\comp b$, de manera que la parella $(a,b)$ és compatible amb el pullback i indueix un únic $u\colon R\to P$ amb $p\comp u=a$ i $q\comp u=b$. Aquest $u$ factoritza $k$ a través de la injecció de $P$: $(m\comp p)\comp u=m\comp a=k$. Per tant $[k]\leq[P]$, i $[P]$ és l'ínfim $[m]\wedge[n]$.
\end{solucio}

\section{Imatges i factoritzacions}

\begin{exercici}
Comprova que a $\cat{Set}$ tota aplicació $f\colon A\to B$ té factorització imatge amb $I=f(A)$, i identifica l'epimorfisme regular i el monomorfisme.
\end{exercici}

\begin{solucio}
Prenem $I=f(A)\subseteq B$. L'aplicació $e\colon A\to I$, $e(a)=f(a)$, és exhaustiva, i tota aplicació exhaustiva a $\cat{Set}$ és un epimorfisme regular: és el coequalitzador de la seva parella nucli $\{(a,a')\mid f(a)=f(a')\}$. La injecció $m\colon I\hookrightarrow B$ és injectiva, per tant un monomorfisme. Aleshores $f=m\comp e$ és la factorització imatge, i $\Imatge(f)=[m]$ és el subconjunt $f(A)$. Recuperem la noció usual d'imatge d'una aplicació.
\end{solucio}

\begin{exercici}
Demostra que si $f=m\comp e$ és una factorització imatge i $f$ és ell mateix un monomorfisme, aleshores $e$ és un isomorfisme.
\end{exercici}

\begin{solucio}
Si $f=m\comp e$ és mono i $m$ és mono, aleshores $e$ és mono: de $e\comp u=e\comp v$ se segueix $m\comp e\comp u=m\comp e\comp v$, és a dir $f\comp u=f\comp v$, d'on $u=v$. Però $e$ també és un epimorfisme regular, en particular un epimorfisme. Un morfisme que és alhora mono i epi regular és un isomorfisme: essent el coequalitzador d'una parella que ja iguala (perquè $e$ mono força la parella nucli a ser la diagonal), $e$ coequalitza la parella trivial i és, doncs, invertible. Per tant $\Imatge(f)=[f]$: la imatge d'un mono és ell mateix.
\end{solucio}

\begin{exercici}
Comprova que la imatge és el menor subobjecte pel qual es factoritza $f$: si $f=n\comp g$ amb $n\colon T\rightarrowtail B$ mono, aleshores $\Imatge(f)\leq[n]$.
\end{exercici}

\begin{solucio}
Sigui $f=m\comp e$ la factorització imatge i $f=n\comp g$ una factorització qualsevol a través del mono $n$. Aleshores $n\comp g=m\comp e$. Com que $e$ és el coequalitzador de la parella nucli de $f$ i $n$ és mono, $g$ iguala les dues branques de la parella nucli, de manera que es factoritza per $e$: hi ha $h$ amb $h\comp e=g$. Aleshores $n\comp h\comp e=n\comp g=f=m\comp e$, i com que $e$ és epi, $n\comp h=m$. Aquesta igualtat és precisament una factorització de $m$ a través de $n$, és a dir $\Imatge(f)=[m]\leq[n]$.
\end{solucio}

\section{Estabilitat i regularitat}

\begin{exercici}
A $\cat{Set}$, comprova directament que les imatges són estables per pullback: si es reindexà $f\colon A\to B$ al llarg de $g\colon B'\to B$, aleshores $\Imatge(f')=g^{-1}(\Imatge(f))$.
\end{exercici}

\begin{solucio}
El pullback de $f$ al llarg de $g$ és $A'=\{(b',a)\mid g(b')=f(a)\}$ amb $f'(b',a)=b'$. La seva imatge és $\{b'\in B'\mid\exists a,\ g(b')=f(a)\}=\{b'\mid g(b')\in f(A)\}=g^{-1}(f(A))=g^{-1}(\Imatge(f))$. Per tant $\Imatge(f')=g^{*}(\Imatge(f))$, que és l'enunciat d'estabilitat. Aquesta comprovació explícita és el model de la propietat abstracta que defineix les categories regulars.
\end{solucio}

\begin{exercici}
Demostra que en una categoria regular els epimorfismes regulars són tancats per composició: si $f\colon A\to B$ i $g\colon B\to C$ són epimorfismes regulars, aleshores $g\comp f$ també ho és.
\end{exercici}

\begin{solucio}
Siguin $f\colon A\to B$ i $g\colon B\to C$ epimorfismes regulars, i factoritzem la composició com a imatge, $g\comp f=m\comp e$, amb $e\colon A\to I$ epi regular i $m\colon I\rightarrowtail C$ mono. Sigui $p_1,p_2\colon R\rightrightarrows A$ la parella nucli de $f$. Com que $f\comp p_1=f\comp p_2$,
\[
m\comp e\comp p_1=g\comp f\comp p_1=g\comp f\comp p_2=m\comp e\comp p_2,
\]
i, com que $m$ és mono, $e\comp p_1=e\comp p_2$. Pel lema~\ref{lema:epi-regular-nucli}, $f$ és el coequalitzador de la seva parella nucli; per tant, existeix $k\colon B\to I$ amb $e=k\comp f$. Aleshores
\[
m\comp k\comp f=m\comp e=g\comp f,
\]
i com que $f$ és epi, $m\comp k=g$: el morfisme $g$ es factoritza a través de $m$. Però $g=\id_C\comp g$ és una factorització imatge de $g$, de manera que $\Imatge(g)=[\id_C]$, i per la minimalitat de la imatge (proposició~\ref{prop:imatge-unica}), $[\id_C]\leq[m]$. Com que $[\id_C]$ és el màxim de $\Sub(C)$, $[m]=[\id_C]$ i $m$ és un isomorfisme. Per tant $g\comp f=m\comp e$ és un epi regular seguit d'un isomorfisme, que torna a ser un epi regular.

A $\cat{Set}$, tot plegat es redueix al fet que la composició d'aplicacions exhaustives és exhaustiva.
\end{solucio}

\section{L'existencial}

\begin{exercici}
Verifica les dues direccions de l'adjunció $\exists_f\dashv f^{*}$ a $\cat{Set}$, on $\exists_f(S)=f(S)$ i $f^{*}(T)=f^{-1}(T)$.
\end{exercici}

\begin{solucio}
Siguin $S\subseteq A$ i $T\subseteq B$. Cal veure $f(S)\subseteq T\Leftrightarrow S\subseteq f^{-1}(T)$.

\emph{($\Rightarrow$)} Si $f(S)\subseteq T$ i $a\in S$, aleshores $f(a)\in f(S)\subseteq T$, així que $a\in f^{-1}(T)$; per tant $S\subseteq f^{-1}(T)$.

\emph{($\Leftarrow$)} Si $S\subseteq f^{-1}(T)$ i $b\in f(S)$, aleshores $b=f(a)$ amb $a\in S\subseteq f^{-1}(T)$, d'on $b=f(a)\in T$; per tant $f(S)\subseteq T$.

Aquesta és la forma conjuntista de l'adjunció imatge directa\,--\,antiimatge, i és exactament la regla d'introducció/eliminació de l'existencial quan $f$ és una projecció.
\end{solucio}

\begin{exercici}
Sigui $\pi\colon X\times Y\to X$ la projecció. Descriu $\exists_\pi[s]$ i $\forall_\pi[s]$ per a un subobjecte $[s]\subseteq X\times Y$ a $\cat{Set}$, i relaciona-ho amb els quantificadors.
\end{exercici}

\begin{solucio}
Sigui $S\subseteq X\times Y$ el subconjunt corresponent a $[s]$. Aleshores
\[
\begin{aligned}
\exists_\pi(S)&=\{x\in X\mid \exists y\in Y:\ (x,y)\in S\},\\
\forall_\pi(S)&=\{x\in X\mid \forall y\in Y:\ (x,y)\in S\}.
\end{aligned}
\]
El primer és la imatge de $S$ per $\pi$ (l'existencial); el segon, el conjunt de fibres senceres contingudes en $S$ (l'universal). Són els casos particulars, per a una projecció, de la cadena $\exists_f\dashv f^{*}\dashv\forall_f$ de l'exemple~\ref{ex:quantificadors-set} del capítol d'adjuncions: l'adjunció $\exists_\pi\dashv\pi^{*}$ es comprova com a l'exercici anterior, i $\pi^{*}\dashv\forall_\pi$ amb un argument anàleg, no per dualitat. Així, en una categoria regular tenim l'existencial; l'universal requereix, a més, adjunts per la dreta, que apareixen en contextos més rics (categories de Heyting).
\end{solucio}


\chapter[Classificadors de subobjectes i topos]{\texorpdfstring{Classificadors de subobjectes i introducció als topos}{Classificadors de subobjectes i introducció als topos}}

Els exercicis d'aquest capítol són de familiarització amb el classificador de subobjectes i la definició de topos. No entren en la lògica interna, que pertany a l'estudi posterior.

\section{El classificador}

\begin{exercici}
Comprova amb detall que a $\cat{Set}$ l'objecte $\Omega=\{0,1\}$ amb $\mathsf{true}\colon 1\to\{0,1\}$, $\mathsf{true}(\ast)=1$, és un classificador de subobjectes.
\end{exercici}

\begin{solucio}
Sigui $U\subseteq E$ un subconjunt, amb injecció $m\colon U\hookrightarrow E$. Definim $\chi_U\colon E\to\{0,1\}$ per $\chi_U(x)=1$ si $x\in U$ i $\chi_U(x)=0$ altrament. El pullback de $\mathsf{true}\colon 1\to\{0,1\}$ al llarg de $\chi_U$ és $\{x\in E\mid\chi_U(x)=1\}=U$, amb la injecció com a projecció; per tant el quadrat és cartesià. Per a la unicitat: si $\chi\colon E\to\{0,1\}$ fa el quadrat cartesià, aleshores $\chi^{-1}(1)=U$, cosa que força $\chi=\chi_U$. Així $\{0,1\}$ classifica els subobjectes de qualsevol conjunt.
\end{solucio}

\begin{exercici}\label{pr-ordre-pertinenca}
En una categoria amb classificador de subobjectes, siguin $m\colon U\rightarrowtail E$ i $n\colon V\rightarrowtail E$ dos subobjectes, amb morfismes característics $\chi_m$ i $\chi_n$. Demostra que
\[
[m]\leq[n]
\quad\Longleftrightarrow\quad
\chi_n\comp m=\mathsf{true}\comp\,!_U,
\]
és a dir, que $U$ està contingut en $V$ si i només si tots els \enquote{elements} de $U$ pertanyen a $V$. Dedueix que $\chi_m=\chi_n$ si i només si $[m]=[n]$.
\end{exercici}

\begin{solucio}
Per definició de l'ordre, $[m]\leq[n]$ vol dir que $m$ factoritza a través de $n$. Aplicant la caracterització de la pertinença per elements generalitzats de la Teoria a l'element generalitzat $m\colon U\to E$, això passa si i només si $\chi_n\comp m=\mathsf{true}\comp\,!_U$.

Per a la segona part, si $[m]=[n]$, els dos subobjectes són el mateix i tenen el mateix morfisme característic, per la unicitat de la definició. Recíprocament, si $\chi_m=\chi_n$, aleshores $\chi_n\comp m=\chi_m\comp m=\mathsf{true}\comp\,!_U$, perquè el quadrat de $m$ commuta, i per la primera part $[m]\leq[n]$; simètricament, $[n]\leq[m]$, i per l'antisimetria de $\Sub(E)$, $[m]=[n]$. Així, el morfisme característic determina el subobjecte, que és la injectivitat de la bijecció $\Sub(E)\cong\Hom(E,\Omega)$ vista ara en termes de pertinença.
\end{solucio}

\begin{exercici}\label{pr-mono-equalitzador}
En una categoria amb classificador de subobjectes, sigui $m\colon U\rightarrowtail E$ un subobjecte amb morfisme característic $\chi_m$. Demostra que $m$ és un equalitzador de $\chi_m$ i $\mathsf{true}\comp\,!_E$. Dedueix que un morfisme que és alhora mono i epi és un isomorfisme.
\end{exercici}

\begin{solucio}
\textbf{Equalitza.} Pel quadrat del classificador, $\chi_m\comp m=\mathsf{true}\comp\,!_U=\mathsf{true}\comp\,!_E\comp m$, perquè $!_E\comp m$ i $!_U$ són fletxes cap al terminal i, per tant, coincideixen.

\textbf{És universal.} Sigui $x\colon X\to E$ amb $\chi_m\comp x=\mathsf{true}\comp\,!_E\comp x$. El membre dret és $\mathsf{true}\comp\,!_X$, de manera que, per la pertinença per elements generalitzats de la Teoria, $x$ factoritza a través de $m$. La factorització és única perquè $m$ és mono. Així $m$ és un equalitzador.

\textbf{Mono i epi implica iso.} Suposem, a més, que $m$ és epi. De $\chi_m\comp m=(\mathsf{true}\comp\,!_E)\comp m$ i la cancel·lació per la dreta, $\chi_m=\mathsf{true}\comp\,!_E$. Aleshores $\id_E$ també equalitza les dues fletxes, i per la propietat de l'equalitzador hi ha $h\colon E\to U$ amb $m\comp h=\id_E$. A més, $m\comp(h\comp m)=m=m\comp\id_U$, i com que $m$ és mono, $h\comp m=\id_U$. Per tant $m$ és un isomorfisme.

Al capítol~\ref{cap:categories} vam veure que, en general, un bimorfisme no és un isomorfisme: a $\cat{Top}$, per exemple, una bijecció contínua pot no ser un homeomorfisme. Una categoria amb classificador no admet aquest fenomen. Als Exercicis es proposa treure'n la conseqüència per a $\cat{Top}$.
\end{solucio}

\section{Topos}

\begin{exercici}
Comprova que $\cat{Set}$ satisfà les tres condicions de la definició de topos.
\end{exercici}

\begin{solucio}
\emph{Límits finits}: $\cat{Set}$ té terminal (un punt), productes (producte cartesià) i equalitzadors (subconjunts definits per igualtats), i per tant tots els límits finits. \emph{Classificador}: l'exercici anterior mostra que $\{0,1\}$ n'és un. \emph{Exponencials}: l'exponencial $C^{B}$ és el conjunt d'aplicacions $B\to C$, i vam veure al capítol de tancament cartesià que $\cat{Set}$ és cartesiana tancada. Per tant $\cat{Set}$ és un topos.
\end{solucio}

\begin{exercici}
Demostra que en un topos l'objecte de parts $\mathcal{P}(E)=\Omega^{E}$ satisfà $\Hom(A,\mathcal{P}(E))\cong\Sub(A\times E)$.
\end{exercici}

\begin{solucio}
Component les dues propietats universals disponibles en un topos:
\[
\Hom(A,\Omega^{E})\cong\Hom(A\times E,\Omega)\cong\Sub(A\times E),
\]
on la primera bijecció és la de l'exponencial (currificació) i la segona és la representabilitat del classificador (\ref{prop:omega-representa}). Totes dues són naturals en $A$. Així, els morfismes cap a $\Omega^{E}$ es corresponen amb subobjectes de $A\times E$, és a dir, amb relacions entre $A$ i $E$. A $\cat{Set}$, això recupera que $\mathcal{P}(E)$ és el conjunt de parts i que una aplicació $A\to\mathcal{P}(E)$ és una relació entre $A$ i $E$.
\end{solucio}

\section{Prefeixos}

\begin{exercici}\label{pr-sedassos}
Sigui $\cat{C}$ una categoria petita, $F$ un prefeix sobre $\cat{C}$ i $U\rightarrowtail F$ un subprefeix, és a dir, $U(C)\subseteq F(C)$ per a cada $C$, amb $F(f)$ enviant $U(C)$ dins de $U(D)$ per a cada $f\colon D\to C$. Per a $x\in F(C)$, posem
\[
\chi_U(C)(x)=\{f\colon D\to C\mid F(f)(x)\in U(D)\}.
\]
Comprova que $\chi_U(C)(x)$ és un sedàs sobre $C$, que les aplicacions $\chi_U(C)$ formen una transformació natural $\chi_U\colon F\Rightarrow\Omega$, i que $U$ és el pullback de $\mathsf{true}$ al llarg de $\chi_U$.
\end{exercici}

\begin{solucio}
\textbf{És un sedàs.} Si $f\colon D\to C$ pertany a $\chi_U(C)(x)$ i $g\colon X\to D$, aleshores $F(f\comp g)(x)=F(g)\bigl(F(f)(x)\bigr)$, i com que $F(f)(x)\in U(D)$ i $U$ és un subprefeix, $F(g)$ l'envia dins de $U(X)$. Per tant $f\comp g$ pertany a $\chi_U(C)(x)$.

\textbf{Naturalitat.} Per a $h\colon D\to C$ i $x\in F(C)$, cal veure que $\chi_U(D)(F(h)(x))=h^{*}\bigl(\chi_U(C)(x)\bigr)$. Per a $g\colon X\to D$,
\[
\begin{aligned}
g\in\chi_U(D)\bigl(F(h)(x)\bigr)
&\iff F(g)\bigl(F(h)(x)\bigr)\in U(X)\\
&\iff F(h\comp g)(x)\in U(X)\\
&\iff h\comp g\in\chi_U(C)(x),
\end{aligned}
\]
i l'última condició és la definició de $g\in h^{*}\bigl(\chi_U(C)(x)\bigr)$.

\textbf{Pullback.} Com que els límits de prefeixos es calculen objecte a objecte, n'hi ha prou de veure que, per a cada $C$, $U(C)$ és el conjunt dels $x\in F(C)$ amb $\chi_U(C)(x)$ igual al sedàs màxim. Si $x\in U(C)$, tot $f\colon D\to C$ compleix $F(f)(x)\in U(D)$, de manera que el sedàs és màxim. Recíprocament, si el sedàs és màxim, conté $\id_C$, i aleshores $x=F(\id_C)(x)\in U(C)$. La unicitat de $\chi_U$ es comprova de la mateixa manera, mirant quins morfismes porten $x$ dins de $U$.

A $\cat{Set}$, vist com a prefeixos sobre la categoria d'un sol objecte i només la identitat, els sedassos sobre l'únic objecte són $\varnothing$ i $\{\id\}$, i es recupera $\Omega=\{0,1\}$.
\end{solucio}

\begin{exercici}\label{pr-conjuncio-omega}
En un topos, sigui $\wedge\colon\Omega\times\Omega\to\Omega$ el morfisme característic del subobjecte $\langle\mathsf{true},\mathsf{true}\rangle\colon 1\rightarrowtail\Omega\times\Omega$. Demostra que, per a dos subobjectes $U,V$ de $E$,
\[
\chi_{U\wedge V}=\wedge\comp\langle\chi_U,\chi_V\rangle,
\]
on $U\wedge V$ és la intersecció. Interpreta-ho a $\cat{Set}$.
\end{exercici}

\begin{solucio}
Per l'exercici~\ref{pr-ordre-pertinenca}, n'hi ha prou de veure que tots dos morfismes classifiquen el mateix subobjecte, és a dir, que un element generalitzat $x\colon X\to E$ hi pertany en un cas si i només si hi pertany en l'altre. Per la pertinença per elements generalitzats:
\begin{itemize}
  \item $x$ factoritza a través de $U\wedge V$ si i només si factoritza a través de $U$ i de $V$, perquè la intersecció és el pullback;
  \item és a dir, si i només si $\chi_U\comp x=\mathsf{true}\comp\,!_X$ i $\chi_V\comp x=\mathsf{true}\comp\,!_X$;
  \item és a dir, si i només si $\langle\chi_U,\chi_V\rangle\comp x=\langle\mathsf{true},\mathsf{true}\rangle\comp\,!_X$, que és la pertinença de $\langle\chi_U,\chi_V\rangle\comp x$ al subobjecte $\langle\mathsf{true},\mathsf{true}\rangle$;
  \item és a dir, si i només si $\wedge\comp\langle\chi_U,\chi_V\rangle\comp x=\mathsf{true}\comp\,!_X$.
\end{itemize}
Sigui $W$ el subobjecte classificat per $\wedge\comp\langle\chi_U,\chi_V\rangle$. Els dos subobjectes $U\wedge V$ i $W$ tenen, doncs, els mateixos elements generalitzats. Prenent com a $x$ un representant de $U\wedge V$, obtenim $U\wedge V\leq W$, i prenent-ne un de $W$, $W\leq U\wedge V$. Per tant són iguals, i per la unicitat del morfisme característic, $\chi_{U\wedge V}=\wedge\comp\langle\chi_U,\chi_V\rangle$.

A $\cat{Set}$, $\wedge\colon\{0,1\}\times\{0,1\}\to\{0,1\}$ val $1$ exactament en $(1,1)$: és la taula de veritat de la conjunció, i la fórmula diu que la funció característica d'una intersecció és el producte, o el mínim, de les funcions característiques.
\end{solucio}

\begin{exercici}\label{pr-reindexacio-graph}
Sigui $G$ el cicle de l'exemple~\ref{ex:omega-graph} de la Teoria, amb vèrtexs $a,b,c$, arestes $e\colon a\to b$, $f\colon b\to c$, $k\colon c\to a$, i el subgraf $H$ de vèrtexs $a,b$ i aresta $e$. Sigui $G'$ el cicle de longitud sis, amb vèrtexs $p_0,\dots,p_5$ i arestes $u_i\colon p_i\to p_{i+1}$ (índexs mòdul $6$), i sigui $F\colon G'\to G$ el morfisme que \enquote{dona dues voltes}: $p_i\mapsto a,b,c$ i $u_i\mapsto e,f,k$ segons que $i$ sigui $0,1,2$ mòdul $3$. Calcula $F^{*}H$ i comprova que $\chi_{F^{*}H}=\chi_H\comp F$.
\end{exercici}

\begin{solucio}
\textbf{La reindexació.} Es calcula per antiimatge, component a component (exemple de subgrafs del capítol~\ref{cap:subobjectes}). Els vèrtexs de $F^{*}H$ són els que van a $a$ o a $b$, és a dir, $p_0,p_1,p_3,p_4$. Les arestes són les que van a $e$, és a dir, $u_0$ i $u_3$.

\textbf{La precomposició.} Amb la taula de $\chi_H$ de la Teoria ($a,b\mapsto1$, $c\mapsto0$, $e\mapsto\top$, $f\mapsto\sigma$, $k\mapsto\tau$), el morfisme $\chi_H\comp F$ envia $p_0,p_1,p_3,p_4\mapsto1$ i $p_2,p_5\mapsto0$, i
\[
u_0,u_3\mapsto\top,\qquad u_1,u_4\mapsto\sigma,\qquad u_2,u_5\mapsto\tau.
\]

\textbf{El morfisme característic de $F^{*}H$.} Apliquem directament la regla de la Teoria a $F^{*}H\subseteq G'$. Les arestes $u_0$ i $u_3$ són al subgraf, i van a $\top$. Les arestes $u_1\colon p_1\to p_2$ i $u_4\colon p_4\to p_5$ tenen l'origen dins i el destí fora, i van a $\sigma$. Les arestes $u_2\colon p_2\to p_3$ i $u_5\colon p_5\to p_0$ tenen l'origen fora i el destí dins, i van a $\tau$. Sobre els vèrtexs coincideix també amb $\chi_H\comp F$. Així $\chi_{F^{*}H}=\chi_H\comp F$: substituir és precompondre, ara en un topos que no és $\cat{Set}$, i al llarg d'un morfisme que no és injectiu.
\end{solucio}

\section{Connectius i igualtat}

\begin{exercici}\label{pr-igualtat-omega}
En una categoria amb productes binaris i classificador, sigui $\delta_A=\langle\id_A,\id_A\rangle\colon A\to A\times A$. Comprova que $\delta_A$ és un monomorfisme i que, per a dos elements generalitzats $x,y\colon X\to A$,
\[
\chi_{\delta_A}\comp\langle x,y\rangle=\mathsf{true}\comp\,!_X
\quad\Longleftrightarrow\quad
x=y.
\]
Descriu $\chi_{\delta_A}$ a $\cat{Set}$.
\end{exercici}

\begin{solucio}
\textbf{És mono.} La projecció $\pi_1$ compleix $\pi_1\comp\delta_A=\id_A$, de manera que $\delta_A$ té una retracció i és un monomorfisme escindit.

\textbf{Pertinença a la diagonal.} Per la pertinença per elements generalitzats de la Teoria, la igualtat de l'esquerra equival a dir que $\langle x,y\rangle$ factoritza a través de $\delta_A$, és a dir, que hi ha $z\colon X\to A$ amb $\langle z,z\rangle=\langle x,y\rangle$. Component amb les projeccions, això vol dir $x=z=y$. Recíprocament, si $x=y$, n'hi ha prou de prendre $z=x$.

\textbf{A $\cat{Set}$.} $\chi_{\delta_A}(a,b)=1$ si $a=b$ i $0$ altrament: és la delta de Kronecker. En termes lògics, $\chi_{\delta_A}$ és el predicat d'igualtat \enquote{$x=y$} sobre $A$, i la diagonal n'és l'extensió. És la fila \enquote{igualtat} del mapa del capítol~\ref{cap:pont}.
\end{solucio}

\begin{exercici}\label{pr-disjuncio-negacio-set}
A $\cat{Set}$, amb $\Omega=\{0,1\}$, sigui ${\vee}\colon\Omega\times\Omega\to\Omega$ el morfisme característic del subconjunt $\{(0,1),(1,0),(1,1)\}$, i sigui $\neg\colon\Omega\to\Omega$ el del subconjunt $\{0\}$. Comprova que, per a subconjunts $U,V\subseteq E$,
\[
\chi_{U\cup V}={\vee}\comp\langle\chi_U,\chi_V\rangle,
\qquad
\chi_{E\setminus U}=\neg\comp\chi_U.
\]
Explica per què la segona igualtat no té un anàleg directe a $\cat{Graph}$.
\end{exercici}

\begin{solucio}
Per a $x\in E$, $({\vee}\comp\langle\chi_U,\chi_V\rangle)(x)=1$ si i només si $(\chi_U(x),\chi_V(x))$ és en $\{(0,1),(1,0),(1,1)\}$, és a dir, si i només si $x\in U$ o $x\in V$. Les dues aplicacions valen, doncs, $1$ exactament a $U\cup V$, i coincideixen. Anàlogament, $(\neg\comp\chi_U)(x)=1$ si i només si $\chi_U(x)=0$, és a dir, si i només si $x\notin U$. Els morfismes $\vee$ i $\neg$ són les taules de veritat de la disjunció i de la negació, i les dues fórmules diuen que la unió i el complement es calculen \enquote{punt a punt} sobre els valors de veritat, com la intersecció a l'exercici~\ref{pr-conjuncio-omega}.

A $\cat{Graph}$, el complement de vèrtexs i d'arestes d'un subgraf no és, en general, un subgraf: pot contenir una aresta sense un dels seus extrems. El millor substitut és el subgraf més gran disjunt de $U$, però, com hem vist a l'exemple~\ref{ex:omega-graph} de la Teoria, la seva unió amb $U$ pot no ser tot el graf. No hi ha, doncs, cap morfisme $\Omega\to\Omega$ que faci el paper de $\neg$ amb la propietat $U\cup\neg U=E$.
\end{solucio}

\fi

% ============================================================
% PART III — EXERCICIS PROPOSATS
% ============================================================
\part{Exercicis proposats}
\setcounter{chapter}{0}
\chapter{Categories}

Els exercicis següents completen els resolts al capítol corresponent de la part de Pràctica. Cada enunciat va acompanyat d'una indicació breu per si el lector s'encalla; convé, però, intentar-lo primer sense mirar-la. La marca \enquote{$\star$} assenyala un exercici de consolidació i \enquote{$\star\star$} un d'ampliació.

\section{Categories i axiomes}

\begin{exercici}
Comprova que un conjunt qualsevol $X$ dona lloc a una categoria \textbf{discreta}: els objectes són els elements de $X$ i els únics morfismes són les identitats. Verifica els axiomes. \hfill{\normalfont$\star$}
\begin{indicacio}
No hi ha cap morfisme entre objectes diferents; l'única composició possible en cada objecte és $\id_x\comp\id_x$.
\end{indicacio}
\end{exercici}

\begin{exercici}
Sigui $\cat{C}$ una categoria. Comprova que fixar un objecte $A$ i considerar només les fletxes $A\to A$, amb la composició de $\cat{C}$, dona un \textbf{monoide}. Quin és el seu element neutre? \hfill{\normalfont$\star$}
\begin{indicacio}
Els morfismes $A\to A$ es poden compondre entre ells i el resultat torna a anar d'$A$ a $A$. El neutre és $\id_A$.
\end{indicacio}
\end{exercici}

\begin{exercici}[Ampliació, per a lectors amb topologia]
Comprova que els espais topològics amb les aplicacions contínues formen una categoria $\cat{Top}$. Cal usar dues propietats de les aplicacions contínues: quines? \hfill{\normalfont$\star$}
\begin{indicacio}
La composició de contínues és contínua, i la identitat és contínua.
\end{indicacio}
\end{exercici}

\section{Fletxes, camins i diagrames commutatius}

\begin{exercici}
Escriu l'equació que expressa la commutativitat del triangle
\[
\begin{tikzpicture}[baseline=(m.center)]
  \matrix (m) [matrix of math nodes, column sep=3.3em, row sep=2.6em] {
    A & B \\
      & C \\
  };
  \draw[-{Stealth[scale=1.05]}] (m-1-1) -- node[above] {$f$} (m-1-2);
  \draw[-{Stealth[scale=1.05]}] (m-1-2) -- node[right] {$g$} (m-2-2);
  \draw[-{Stealth[scale=1.05]}] (m-1-1) -- node[below left] {$h$} (m-2-2);
\end{tikzpicture}
\]
i explica quins són els dos camins que es comparen. \hfill{\normalfont$\star$}
\begin{indicacio}
Tots dos camins comencen a $A$ i acaben a $C$; un passa per $B$ i l'altre és la fletxa directa.
\end{indicacio}
\end{exercici}

\begin{exercici}
Representa mitjançant un quadrat commutatiu l'equació
\[
q\comp u=v\comp p,
\]
on $u\colon A\to B$, $q\colon B\to D$, $p\colon A\to C$ i $v\colon C\to D$. \hfill{\normalfont$\star$}
\begin{indicacio}
Col·loca $A$ a l'angle superior esquerre i $D$ a l'inferior dret. Els dos membres de l'equació són dos camins d'$A$ a $D$.
\end{indicacio}
\end{exercici}

\begin{exercici}
Dona un exemple de quatre aplicacions entre conjunts que formin un quadrat que \emph{no} commuti. Indica explícitament un element sobre el qual les dues composicions donin resultats diferents. \hfill{\normalfont$\star$}
\begin{indicacio}
Pots treballar amb conjunts molt petits, per exemple $\{0,1\}$, i comparar els dos camins sobre l'element $0$.
\end{indicacio}
\end{exercici}

\section{Primers exemples}

\begin{exercici}
En un sistema deductiu $T$, suposem
\[
A\vdash_T B,
\qquad
B\vdash_T C,
\qquad
C\vdash_T A.
\]
Dibuixa el diagrama corresponent a $\cat{Prov}_T$ i comprova que $A$, $B$ i $C$ són isomorfs dos a dos. \hfill{\normalfont$\star$}
\begin{indicacio}
La transitivitat dona també $A\vdash_T C$, $B\vdash_T A$ i $C\vdash_T B$. En una categoria prima, dues fletxes en sentits oposats ja donen un isomorfisme.
\end{indicacio}
\end{exercici}

\begin{exercici}
Comprova directament que la composició de dos morfismes de grafs és un morfisme de grafs. Si $F\colon G\to H$ i $K\colon H\to L$, verifica les equacions d'origen i destí per a $K\comp F$. \hfill{\normalfont$\star$}
\begin{indicacio}
Per a l'origen, parteix de
\[
F_V\comp s_G=s_H\comp F_E,
\qquad
K_V\comp s_H=s_L\comp K_E,
\]
i compon les igualtats adequadament. Fes el mateix amb $t$.
\end{indicacio}
\end{exercici}

\begin{exercici}
Explica amb precisió quina informació té un graf dirigit que encara no té l'estructura d'una categoria. Després indica què afegeix la categoria lliure sobre aquest graf. \hfill{\normalfont$\star$}
\begin{indicacio}
Un graf té vèrtexs, arestes, origen i destí. Una categoria necessita, a més, identitats, composició i axiomes; la categoria lliure els genera mitjançant camins.
\end{indicacio}
\end{exercici}

\section{Isomorfismes}

\begin{exercici}
Comprova que a la categoria $\cat{Mon}$ dels monoides i els homomorfismes, un isomorfisme és exactament un homomorfisme bijectiu. Compara-ho amb el que pot passar a $\cat{Pos}$. \hfill{\normalfont$\star\star$}
\begin{indicacio}
Comprova que la inversa d'un homomorfisme bijectiu de monoides preserva l'operació i el neutre, de manera que torna a ser un homomorfisme. Per a $\cat{Pos}$, en canvi, la inversa d'una aplicació monòtona bijectiva pot no ser monòtona. Fes el contrast explícit amb un exemple: pren el conjunt $\{0,1\}$ amb l'ordre \emph{discret} (cap dels dos elements comparable amb l'altre) com a domini, i el mateix conjunt amb l'ordre $0<1$ com a codomini. La identitat subjacent és una aplicació monòtona (trivialment, perquè al domini no hi ha cap parella comparable) i bijectiva, però la seva inversa ---la identitat de l'ordre $0<1$ cap a l'ordre discret--- no és monòtona, ja que hauria d'enviar $0<1$ a dos elements incomparables. Per tant no és un isomorfisme a $\cat{Pos}$, tot i ser una bijecció monòtona.
\end{indicacio}
\end{exercici}

\begin{exercici}
Comprova que si $g\comp f$ és un isomorfisme, no cal que ho siguin ni $f$ ni $g$. Busca'n un exemple senzill a $\cat{Set}$ i representa'l amb fletxes. \hfill{\normalfont$\star$}
\begin{indicacio}
Pensa en una injecció seguida d'una projecció entre conjunts de mides diferents.
\end{indicacio}
\end{exercici}

\section{Classes especials de morfismes}

\begin{exercici}
Comprova que tot isomorfisme té secció. Després, prenent $A=\{0,1\}$, $B=\{\ast\}$ i $f$ l'única aplicació $A\to B$, dona una secció $s\colon B\to A$ explícita i comprova que $f$ no és isomorfisme. Observa que en aquest cas $f$ té \emph{dues} seccions distintes: quines? \hfill{\normalfont$\star$}
\begin{indicacio}
Tria $s(\ast)=0$ o $s(\ast)=1$. En tots dos casos $f\comp s=\id_B$, però $s\comp f\neq\id_A$. Les dues tries donen les dues seccions.
\end{indicacio}
\end{exercici}

\begin{exercici}
A la categoria $\cat{Set}$, exhibeix morfismes que il·lustrin els tres casos possibles pel que fa a l'existència de seccions: un morfisme sense cap secció, un amb una única secció i un amb més d'una. \hfill{\normalfont$\star\star$}
\begin{indicacio}
Per a \emph{cap} secció, cerca una aplicació no exhaustiva, com $\emptyset\to\{\ast\}$ o una injecció no exhaustiva. Per a \emph{una única}, una aplicació exhaustiva on cada element del codomini tingui una sola antiimatge (una bijecció). Per a \emph{diverses}, una aplicació exhaustiva amb alguna fibra de més d'un element: cada elecció d'antiimatges dona una secció distinta. Recorda que a $\cat{Set}$ triar una secció d'una aplicació exhaustiva és triar un element a cada fibra.
\end{indicacio}
\end{exercici}

\begin{exercici}
Demostra la caracterització dual de l'exercici resolt a la part de pràctica: un morfisme $f$ és un isomorfisme si i només si és un epimorfisme escindit (té secció) i un monomorfisme. \hfill{\normalfont$\star$}
\begin{indicacio}
És el dual de \enquote{mono escindit i epi $\Rightarrow$ iso}: parteix d'una secció $s$ amb $f\comp s=\id_B$ i fes servir que $f$ és mono per veure que $s\comp f=\id_A$.
\end{indicacio}
\end{exercici}

\begin{exercici}
Sigui $\cat{C}$ una categoria prima. Per l'exercici resolt~\ref{exr:pr-prima-mono-epi} de la Pràctica, tot morfisme de $\cat{C}$ és un bimorfisme. Demostra que un morfisme $f\colon A\to B$ de $\cat{C}$ és un isomorfisme si i només si existeix algun morfisme $B\to A$. Aplica-ho a la fletxa $u\colon 0\to 1$ de $\cat{2}$ i a dues fórmules interdeduïbles i distintes de $\cat{Prov}_T$. \hfill{\normalfont$\star$}
\begin{indicacio}
Si $g\colon B\to A$ existeix, $g\comp f$ i $\id_A$ pertanyen tots dos a $\Hom(A,A)$, que té com a màxim un element; raona igual amb $f\comp g$ i $\id_B$. A $\cat{2}$ no hi ha cap fletxa $1\to 0$, de manera que $u$ és un bimorfisme que no és iso. A $\cat{Prov}_T$, en canvi, la fletxa que expressa $p\vdash_T q$ és un isomorfisme que no és una identitat quan també $q\vdash_T p$.
\end{indicacio}
\end{exercici}

\begin{exercici}
Estudia quin factor hereta la propietat en una composició. Demostra que si $g\comp f$ és un monomorfisme, aleshores $f$ és un monomorfisme; i que, dualment, si $g\comp f$ és un epimorfisme, aleshores $g$ és un epimorfisme. Comprova amb un contraexemple que, en el primer cas, $g$ no cal que sigui mono, i en el segon, $f$ no cal que sigui epi. \hfill{\normalfont$\star\star$}
\begin{indicacio}
Per al mono: de $f\comp u=f\comp v$, compon per l'esquerra amb $g$ per obtenir $(g\comp f)\comp u=(g\comp f)\comp v$ i aplica que $g\comp f$ és mono. El cas de l'epi és el dual: compon per la dreta. Observa l'asimetria: d'un mono $g\comp f$ hereta el factor de la \emph{dreta} ($f$), i d'un epi, el de l'\emph{esquerra} ($g$). Per als contraexemples, pensa en una injecció i una projecció a $\cat{Set}$.
\end{indicacio}
\end{exercici}

\begin{exercici}
Sigui $M=\{1,e\}$ el monoide amb $e\cdot e=e$ i $1$ neutre, i sigui $\cat{B}M$ la categoria d'un sol objecte associada. Decideix si $e$ és un monomorfisme, si és un epimorfisme i si té secció o retracció. \hfill{\normalfont$\star\star$}
\begin{indicacio}
Compara $e\comp 1$ amb $e\comp e$, i $1\comp e$ amb $e\comp e$. Per a les inverses unilaterals, només hi ha dos candidats.
\end{indicacio}
\end{exercici}

\begin{exercici}
A $\cat{Set}$, demostra que tota aplicació exhaustiva $f\colon A\to B$ és un epimorfisme sense triar cap secció de $f$. Explica després per què l'enunciat \enquote{tota aplicació exhaustiva té una secció} no és una conseqüència d'aquest fet, sinó un enunciat diferent, equivalent a l'axioma de l'elecció. \hfill{\normalfont$\star\star$}
\begin{indicacio}
Si $g\comp f=h\comp f$ i $y\in B$, n'hi ha prou de tenir \emph{algun} $x$ amb $f(x)=y$ per concloure $g(y)=h(y)$; no cal triar-ne un per a cada $y$ alhora. Una secció, en canvi, és precisament una tria simultània d'una antiimatge per a cada element de $B$.
\end{indicacio}
\end{exercici}

\section{Construccions sobre categories}

\begin{exercici}
Descriu la categoria producte $\cat{3}\times\cat{3}$, on $\cat{3}$ és la categoria associada a l'ordre $0<1<2$. Quants objectes i quants morfismes té? Dibuixa les fletxes bàsiques i explica quines altres fletxes apareixen per composició. \hfill{\normalfont$\star\star$}
\begin{indicacio}
$\cat{3}$ té tres objectes i sis morfismes. En el producte, els objectes i els morfismes són parelles.
\end{indicacio}
\end{exercici}

\begin{exercici}
Descriu explícitament la categoria $\cat{3}/2$ sobre l'objecte $2$ de $\cat{3}$, on $\cat{3}$ és la categoria associada a l'ordre $0<1<2$. Enumera'n els objectes i els morfismes i identifica, si pots, una categoria finita isomorfa a aquesta categoria sobre un objecte. \hfill{\normalfont$\star$}
\begin{indicacio}
Els objectes de $\cat{3}/2$ són les fletxes de $\cat{3}$ amb codomini $2$, i un morfisme entre dos d'aquests objectes és una fletxa de $\cat{3}$ que fa commutar el triangle corresponent. Com que $\cat{3}$ és prima, tots els triangles commuten. Observa també que de cada objecte de $\cat{3}$ surt exactament una fletxa cap a $2$.
\end{indicacio}
\end{exercici}

\begin{exercici}
Compara directament $(\cat{C}\times\cat{D})\op$ i $\cat{C}\op\times\cat{D}\op$. Comprova que tenen els mateixos objectes i que els seus morfismes, identitats i composicions s'identifiquen component a component. \hfill{\normalfont$\star$}
\begin{indicacio}
Un morfisme $(A,X)\to(B,Y)$ a $(\cat{C}\times\cat{D})\op$ és, per definició, un morfisme $(B,Y)\to(A,X)$ a $\cat{C}\times\cat{D}$. Escriu què significa això en cadascuna de les dues components.
\end{indicacio}
\end{exercici}

\begin{exercici}
Per a cadascuna de les subcategories següents de $\cat{Set}$, digues si és \textbf{plena}, \textbf{ampla}, totes dues coses o cap: (a)~$\cat{FinSet}$, amb els conjunts finits per objectes i totes les aplicacions entre ells; (b)~la subcategoria amb tots els conjunts per objectes i només les aplicacions injectives; (c)~la subcategoria amb tots els conjunts per objectes i totes les aplicacions (és a dir, $\cat{Set}$ mateixa); (d)~la subcategoria amb els conjunts finits per objectes i només les aplicacions injectives entre ells. \hfill{\normalfont$\star$}
\begin{indicacio}
Recorda que \emph{plena} vol dir que es conserven \emph{tots} els morfismes entre els objectes escollits (la subcategoria queda determinada només pels objectes), i \emph{ampla} vol dir que es conserven \emph{tots} els objectes. Comprova cada cas contra aquestes dues condicions per separat: poden fallar independentment. En particular, mira si (d) és plena o ampla, i observa que una subcategoria pot no ser ni l'una ni l'altra.
\end{indicacio}
\end{exercici}

\section{Principi de dualitat}

\begin{exercici}
Escriu el dual de l'afirmació següent i justifica'l sense fer una nova demostració completa:
\begin{quote}
Si $f\colon A\to B$ té una retracció, aleshores $f$ és un monomorfisme.
\end{quote}
\hfill{\normalfont$\star$}
\begin{indicacio}
Passa a la categoria oposada i identifica quina noció és dual de \enquote{retracció} i quina és dual de \enquote{monomorfisme}. Recorda que, en dualitzar, s'inverteixen tant el sentit de les fletxes com l'ordre de les composicions.
\end{indicacio}
\end{exercici}

\section{Grafs i categories lliures}

\begin{exercici}
Classifica, \textbf{fins a isomorfisme}, els grafs dirigits que generen una categoria lliure amb exactament \textbf{quatre} morfismes. Recorda que les identitats compten com a morfismes. \hfill{\normalfont$\star\star$}
\begin{indicacio}
Cada vèrtex aporta una identitat i cada aresta un camí de longitud~$1$: quantes arestes i quants vèrtexs pot tenir, com a màxim, el graf? Què passaria si hi hagués un cicle dirigit? Després, separa els casos segons el nombre de vèrtexs.
\end{indicacio}
\end{exercici}

\begin{exercici}
Considera el graf
\[
A\xrightarrow{f}B\xrightarrow{g}C.
\]
Llista tots els morfismes de la categoria lliure que genera, dibuixa també la fletxa composta i comprova que hi ha exactament sis morfismes. \hfill{\normalfont$\star$}
\begin{indicacio}
Compta els tres camins buits, les dues arestes $f$ i $g$, i el camí compost de longitud dos.
\end{indicacio}
\end{exercici}

\begin{exercici}
Per què un graf amb un cicle, per exemple una aresta $a\colon A\to A$, genera sempre una categoria lliure amb infinits morfismes? \hfill{\normalfont$\star$}
\begin{indicacio}
El cicle es pot recórrer tantes vegades com es vulgui: $a,a^2,a^3,\ldots$, i cada repetició dona un camí diferent.
\end{indicacio}
\end{exercici}

\section{Qüestions de mida}

\begin{exercici}
Comprova que un preordre, vist com a categoria, és una categoria \textbf{petita} sempre que el conjunt subjacent $P$ sigui un conjunt. Quants morfismes hi ha, com a màxim, entre dos objectes? \hfill{\normalfont$\star$}
\begin{indicacio}
Entre dos objectes hi ha zero o un morfisme; la col·lecció total de morfismes es pot identificar amb un subconjunt de $P\times P$.
\end{indicacio}
\end{exercici}

\begin{exercici}
Explica per què el fet que tots els objectes d'una categoria siguin finits no implica que la categoria sigui petita. Aplica-ho a $\cat{FinSet}$, la categoria que té per objectes tots els conjunts finits i per morfismes totes les aplicacions entre ells. \hfill{\normalfont$\star$}
\begin{indicacio}
La mida de cada objecte i la mida de la col·lecció de tots els objectes són qüestions diferents. Per a cada conjunt $x$, el conjunt $\{x\}$ és finit: quants objectes té, doncs, $\cat{FinSet}$? Compara-ho amb la paradoxa de Russell.
\end{indicacio}
\end{exercici}

\begin{exercici}
Explica per què $\cat{Graph}$ és localment petita: donats dos grafs dirigits $G$ i $H$, per què els morfismes de grafs $G\to H$ formen un conjunt? \hfill{\normalfont$\star$}
\begin{indicacio}
Un morfisme de grafs és una parella d'aplicacions $V_G\to V_H$ i $E_G\to E_H$ que satisfà dues condicions. Aquestes parelles formen un subconjunt d'un producte de dos conjunts d'aplicacions.
\end{indicacio}
\end{exercici}

\chapter{Functors}

Els exercicis següents completen els resolts al capítol corresponent de la part de Pràctica. Cada enunciat va acompanyat d'una indicació breu per si el lector s'encalla; convé, però, intentar-lo primer sense mirar-la. La marca \enquote{$\star$} assenyala un exercici de consolidació i \enquote{$\star\star$} un d'ampliació.

\section{Definició de functor}

\begin{exercici}
Comprova que el functor identitat $\id_{\cat{C}}\colon\cat{C}\to\cat{C}$ i, per a cada objecte $D$ d'una categoria $\cat{D}$, el functor constant $\Delta_D\colon\cat{C}\to\cat{D}$ ---que envia tot objecte a $D$ i tot morfisme a $\id_D$ (exemple~\ref{ex:functor-constant})--- són functors. \hfill{\normalfont$\star$}
\begin{indicacio}
Per al functor constant, comprova que $\Delta_D(\id_A)=\id_D$ i que $\Delta_D(g\comp f)=\id_D=\id_D\comp\id_D=\Delta_D(g)\comp\Delta_D(f)$.
\end{indicacio}
\end{exercici}

\begin{exercici}\label{ex:identitat-necessaria}
Considera una assignació que respecti dominis, codominis i composició, però de la qual \emph{no} se suposi que preservi identitats. Comprova, amb un contraexemple senzill, que la preservació de la identitat no se'n dedueix, i conclou que la condició $F(\id_A)=\id_{F(A)}$ és realment necessària a la definició de functor. \hfill{\normalfont$\star$}
\begin{indicacio}
Pren el monoide de dos elements $\{1,e\}$ amb $e\comp e=e$ i $e\comp m=m\comp e=e$ per a tot $m$, vist com a categoria d'un sol objecte. L'assignació que envia cada morfisme $m$ a $e$ respecta la composició, perquè $F(m\comp n)=e=e\comp e=F(m)\comp F(n)$, però envia la identitat $1$ a $e\neq 1$. Per tant no preserva identitats malgrat respectar la composició.
\end{indicacio}
\end{exercici}

\section{Exemples de functors}

\begin{exercici}
Comprova que hi ha un functor $V\colon\cat{Graph}\to\cat{Set}$ que envia cada graf $G$ al seu conjunt de vèrtexs $V_G$ i cada morfisme de grafs $F=(F_V,F_E)\colon G\to H$ a la seva component sobre vèrtexs $F_V\colon V_G\to V_H$. Quines propietats de la composició i de les identitats de $\cat{Graph}$ en garanteixen la functorialitat? Fes el mateix per al functor d'arestes $E\colon\cat{Graph}\to\cat{Set}$. \hfill{\normalfont$\star$}
\begin{indicacio}
La identitat d'un graf té per component sobre vèrtexs la identitat de $V_G$, i la composició de morfismes de grafs es fa component a component: $(F'\comp F)_V=F'_V\comp F_V$. Quan hagis estudiat els homfunctors (subsecció~\ref{subsec:homfunctors}), compara $V$ amb $\Hom_{\cat{Graph}}(\mathbf V,-)$, on $\mathbf V$ és el graf d'un sol vèrtex.
\end{indicacio}
\end{exercici}

\begin{exercici}[Ampliació, per a lectors amb topologia]
Comprova que hi ha un functor oblit $U\colon\cat{Top}\to\cat{Set}$ i descriu explícitament què fa sobre objectes i sobre morfismes. Quines propietats de les aplicacions contínues en garanteixen la functorialitat? \hfill{\normalfont$\star$}
\begin{indicacio}
$U$ envia cada espai al seu conjunt de punts i cada aplicació contínua a la mateixa aplicació. Cal que la identitat sigui contínua i que la composició de contínues sigui contínua.
\end{indicacio}
\end{exercici}

\begin{exercici}[Ampliació, per a lectors amb àlgebra]
Sigui $G$ un grup vist com a categoria $\cat{B}G$ amb un sol objecte. Comprova que un functor $\cat{B}G\to\cat{Set}$ és el mateix que una \textbf{acció} de $G$ sobre un conjunt, i que la imatge de cada element del grup és necessàriament una \emph{bijecció}. \hfill{\normalfont$\star$}
\begin{indicacio}
Segueix el cas del monoide. La diferència és que en un grup cada element té invers; usa que un functor preserva isomorfismes i que tot morfisme de $\cat{B}G$ és un isomorfisme.
\end{indicacio}
\end{exercici}

\begin{exercici}
Descriu un functor $\cat{2}\to\cat{C}$, on $\cat{2}$ és la categoria associada a l'ordre $0<1$. Comprova que triar un tal functor equival a triar un morfisme de $\cat{C}$. \hfill{\normalfont$\star$}
\begin{indicacio}
El functor ha d'enviar l'única fletxa no trivial $0\to1$ a un morfisme de $\cat{C}$; els dos objectes es determinen com el seu domini i el seu codomini.
\end{indicacio}
\end{exercici}

\section{Composició de functors i la categoria \texorpdfstring{$\cat{Cat}$}{Cat}}

\begin{exercici}
Comprova que un isomorfisme a $\cat{Cat}$ ---és a dir, un functor $F\colon\cat{C}\to\cat{D}$ amb un functor invers $G$ tal que $G\comp F=\id_{\cat{C}}$ i $F\comp G=\id_{\cat{D}}$--- és bijectiu sobre objectes i sobre morfismes. \hfill{\normalfont$\star\star$}
\begin{indicacio}
De $G\comp F=\id_{\cat{C}}$ i $F\comp G=\id_{\cat{D}}$ dedueix que $F$ té inversa com a aplicació sobre objectes i com a aplicació sobre morfismes.
\end{indicacio}
\end{exercici}

\begin{exercici}
Siguin $F\colon\cat{E}\to\cat{C}$ i $G\colon\cat{E}\to\cat{D}$ dos functors. Defineix un functor $\langle F,G\rangle\colon\cat{E}\to\cat{C}\times\cat{D}$ per $\langle F,G\rangle(E)=(F(E),G(E))$ i $\langle F,G\rangle(h)=(F(h),G(h))$. Comprova que és un functor, que
\[
\pi_1\comp\langle F,G\rangle=F,\qquad \pi_2\comp\langle F,G\rangle=G,
\]
on $\pi_1,\pi_2$ són les projeccions de l'exemple~\ref{ex:projeccions-diagonal}, i que és l'únic functor $\cat{E}\to\cat{C}\times\cat{D}$ amb aquestes dues propietats. Quin functor s'obté amb $\cat{E}=\cat{C}=\cat{D}$ i $F=G=\id_{\cat{C}}$? \hfill{\normalfont$\star$}
\begin{indicacio}
La functorialitat es comprova component a component, com per a $\Delta$. Per a la unicitat, si $H\colon\cat{E}\to\cat{C}\times\cat{D}$ compleix les dues igualtats, escriu $H(E)$ i $H(h)$ com a parelles i mira què en diuen $\pi_1$ i $\pi_2$. És la propietat universal del producte, ara a nivell de categories.
\end{indicacio}
\end{exercici}

\section{Functors i isomorfismes}

\begin{exercici}
Comprova que si $F\colon\cat{C}\to\cat{D}$ és un functor i $A\cong B$ a $\cat{C}$, aleshores $F(A)\cong F(B)$ a $\cat{D}$. És a dir, els functors preserven la relació d'isomorfisme entre objectes. \hfill{\normalfont$\star$}
\begin{indicacio}
Aplica $F$ a un isomorfisme $A\to B$ i usa que la imatge d'un isomorfisme és un isomorfisme.
\end{indicacio}
\end{exercici}

\begin{exercici}
Dona un exemple de dos objectes no isomorfs que un functor envia a objectes isomorfs. Això mostra que els functors no reflecteixen, en general, la relació d'isomorfisme entre objectes. \hfill{\normalfont$\star$}
\begin{indicacio}
Un functor constant envia tots els objectes al mateix objecte, i per tant tots a objectes isomorfs entre ells.
\end{indicacio}
\end{exercici}

\section{Functors covariants i contravariants}

\begin{exercici}
Donat un functor $F\colon\cat{C}\to\cat{D}$, comprova que la mateixa regla sobre objectes i morfismes defineix un functor
\[
F\op\colon\cat{C}\op\to\cat{D}\op,
\]
el \textbf{functor oposat} de $F$. Comprova que $F$ és fidel (respectivament ple) si i només si $F\op$ ho és. \hfill{\normalfont$\star$}
\begin{indicacio}
Sobre objectes, $F\op=F$; sobre morfismes, $F\op(f)=F(f)$, però ara llegit entre els mateixos objectes a les categories oposades. Compondre a $\cat{C}\op$ i a $\cat{D}\op$ és compondre a $\cat{C}$ i a $\cat{D}$ en l'ordre invers, de manera que les dues condicions de functor es transporten. Els conjunts de morfismes són els mateixos, així que injectivitat i exhaustivitat es conserven.
\end{indicacio}
\end{exercici}

\begin{exercici}
Fixat un objecte $A$ d'una categoria localment petita, comprova que $\Hom(A,-)$ envia isomorfismes a bijeccions: si $f\colon X\to Y$ és un isomorfisme, aleshores
\[
f_*\colon\Hom(A,X)\to\Hom(A,Y)
\]
és una bijecció. \hfill{\normalfont$\star$}
\begin{indicacio}
$\Hom(A,-)$ és un functor i, per tant, preserva isomorfismes; a $\cat{Set}$ els isomorfismes són les bijeccions. Alternativament, comprova directament que $(f^{-1})_*$ és la inversa de $f_*$.
\end{indicacio}
\end{exercici}

\section{Functors fidels i plens}

\begin{exercici}
Comprova que la composició de dos functors fidels és fidel, i que la composició de dos functors plens és plena. Dedueix que la composició de functors plens i fidels és plena i fidel. \hfill{\normalfont$\star$}
\begin{indicacio}
Treballa amb les aplicacions $F_{A,B}$ sobre conjunts de morfismes: la composició de dues injeccions és injectiva, i la de dues exhaustives, exhaustiva.
\end{indicacio}
\end{exercici}

\begin{exercici}
Demostra que un functor fidel reflecteix monomorfismes i epimorfismes: si $F$ és fidel i $F(f)$ és mono (respectivament epi), aleshores $f$ és mono (respectivament epi). Indica on fas servir que $F$ és fidel, i comprova que no cal que $F$ sigui ple. \hfill{\normalfont$\star$}
\begin{indicacio}
Parteix d'una igualtat de composicions a la categoria d'origen, aplica-hi el functor i cancel·la a la d'arribada.
\end{indicacio}
\end{exercici}

\begin{exercici}
Comprova que el functor diagonal $\Delta\colon\cat{C}\to\cat{C}\times\cat{C}$ de l'exemple~\ref{ex:projeccions-diagonal} és sempre fidel, i que és ple si i només si $\cat{C}$ és prima. \hfill{\normalfont$\star$}
\begin{indicacio}
Sobre cada homconjunt, $\Delta$ actua com $f\mapsto(f,f)$. Quines parelles $(f,g)$ de $\cat{C}\times\cat{C}$ provenen d'una sola fletxa?
\end{indicacio}
\end{exercici}

\begin{exercici}
Comprova que el functor oblit $U\colon\cat{Cat}\to\cat{Graph}$, que envia cada categoria petita al seu graf subjacent, és fidel però \emph{no} ple. Mostra també que, per a tot graf $G$, hi ha un morfisme canònic de grafs $\eta_G\colon G\to U(\Free(G))$. \hfill{\normalfont$\star$}
\begin{indicacio}
Quant a ser fidel, un functor queda determinat per la seva acció sobre objectes i morfismes, que és justament el que reté el graf subjacent. Per veure que no és ple, observa que un morfisme de grafs entre els grafs subjacents no ha de respectar la composició, i per tant no prové necessàriament d'un functor. Per al morfisme canònic $\eta_G$, envia cada vèrtex a ell mateix i cada aresta de $G$ al camí de longitud~$1$ corresponent dins $\Free(G)$.
\end{indicacio}
\end{exercici}

\begin{exercici}[Ampliació]
Caracteritza exactament per a quins grafs $G$ el morfisme canònic $\eta_G\colon G\to U(\Free(G))$ de l'exercici anterior és un isomorfisme de grafs. \hfill{\normalfont$\star\star$}
\begin{indicacio}
Compara les arestes de $U(\Free(G))$ amb les imatges de les arestes de $G$. Quins camins de $\Free(G)$ existeixen sempre, sigui quin sigui $G$, i no tenen longitud~$1$? Tingues en compte també el cas extrem.
\end{indicacio}
\end{exercici}

\begin{exercici}
Comprova que un functor ple i fidel \textbf{reflecteix els isomorfismes entre objectes}, en el sentit següent:
\[
F(A)\cong F(B)\quad\Longrightarrow\quad A\cong B.
\]
Suposant a més que $\cat{C}$ i $\cat{D}$ són petites, dedueix que $F$ indueix una aplicació injectiva entre les classes d'isomorfisme d'objectes de $\cat{C}$ i les de $\cat{D}$. \hfill{\normalfont$\star\star$}
\begin{indicacio}
Pren un isomorfisme $F(A)\to F(B)$. Com que $F$ és ple, prové d'un morfisme $A\to B$; per la propietat que els functors plens i fidels reflecteixen isomorfismes, aquest morfisme és un isomorfisme, de manera que $A\cong B$. La implicació $F(A)\cong F(B)\Rightarrow A\cong B$ val sempre; la hipòtesi de petitesa només s'usa a la darrera part, per poder parlar del conjunt de classes d'isomorfisme: en una categoria petita els objectes formen un conjunt, i el quocient per la relació $\cong$ també, de manera que \enquote{l'aplicació entre classes d'isomorfisme} té sentit literal. En una categoria gran caldria una convenció addicional per manejar aquestes classes, que aquí evitem.
\end{indicacio}
\end{exercici}

\begin{exercici}
Considera l'assignació $\cat{C}\to\cat{C}\op$ que deixa fixos els objectes i envia cada morfisme $f$ a ell mateix, vist ara en sentit contrari. Comprova que \emph{no} és, en general, un functor covariant, però sí un functor contravariant ple i fidel. Quin és aquest functor? \hfill{\normalfont$\star\star$}
\begin{indicacio}
Sobre morfismes, aquesta assignació inverteix l'ordre de la composició; és el functor identitat vist com a functor contravariant $\cat{C}\to\cat{C}\op$. La bijecció sobre cada conjunt de morfismes és evident.
\end{indicacio}
\end{exercici}

\begin{exercici}
Comprova que si $F\colon\cat{C}\to\cat{D}$ és ple i fidel, i $\cat{C}$ és localment petita, aleshores $F$ estableix una bijecció
\[
\Hom_{\cat{C}}(A,B)\;\xrightarrow{\;\sim\;}\;\Hom_{\cat{D}}(F(A),F(B))
\]
per a cada parella d'objectes. Explica per què això no obliga $F$ a ser injectiu sobre objectes. \hfill{\normalfont$\star\star$}
\begin{indicacio}
La bijecció és la definició mateixa de ple i fidel. Per a la segona part, busca una categoria amb dos objectes isomorfs i exactament un morfisme entre cada parella ordenada, i compara-la amb $\cat{1}$.
\end{indicacio}
\end{exercici}

\begin{exercici}
Distingeix l'exhaustivitat \emph{estricta} sobre objectes (tot objecte del codomini és exactament una imatge) de l'\emph{essencial} (tot objecte del codomini és isomorf a una imatge). Sigui $\cat{I}$ la categoria amb dos objectes $a,b$ i exactament un morfisme entre cada parella ordenada d'objectes, i considera el functor $E\colon\cat{1}\to\cat{I}$ que envia l'únic objecte $\ast$ a $a$. Comprova que $E$ és essencialment exhaustiu però no exhaustiu sobre objectes. \hfill{\normalfont$\star\star$}
\begin{indicacio}
Observa que $E$ arriba estrictament només a $a$, però que $b\cong a$.
\end{indicacio}
\end{exercici}

\begin{exercici}
Sigui $T$ la lògica proposicional clàssica. Comprova que la negació $A\mapsto\neg A$ defineix un functor contravariant
\[
N\colon\cat{Prov}_T\op\to\cat{Prov}_T,
\]
i que, per tant, la contraposició ---de $A\vdash_T B$ es dedueix $\neg B\vdash_T\neg A$--- és exactament la functorialitat de $N$. Comprova que $N$ és fidel i que és ple. Explica quina llei lògica cal perquè $N$ sigui ple i per què, en lògica intuïcionista, $N$ continua sent un functor però deixa de ser ple. \hfill{\normalfont$\star\star$}
\begin{indicacio}
Les dues categories són primes, de manera que només cal comprovar que les fletxes imatge existeixen. Que $N$ sigui ple demana que de $\neg B\vdash_T\neg A$ se'n segueixi $A\vdash_T B$. Per veure que pot fallar en lògica intuïcionista, busca una fórmula que no sigui estable per doble negació.
\end{indicacio}
\end{exercici}

\chapter[Transformacions naturals]{\texorpdfstring{Transformacions naturals\\ i equivalències}{Transformacions naturals i equivalències}}

Els exercicis següents completen els resolts a la part de Pràctica. Cada enunciat va acompanyat d'una indicació breu; convé intentar-lo primer sense mirar-la. La marca \enquote{$\star$} assenyala un exercici de consolidació i \enquote{$\star\star$} un d'ampliació.

\section{Definició de transformació natural}

\begin{exercici}
Siguin $F,G,H\colon\cat{C}\to\cat{D}$ tres functors i $\eta\colon F\Rightarrow G$, $\theta\colon G\Rightarrow H$ dues transformacions naturals. Comprova directament, a partir de la definició, que la composició vertical $\theta\comp\eta$ satisfà la condició de naturalitat. \hfill{\normalfont$\star$}
\begin{indicacio}
Encadena els quadrats de naturalitat de $\eta$ i de $\theta$: $H(f)\comp\theta_A\comp\eta_A=\theta_B\comp G(f)\comp\eta_A=\theta_B\comp\eta_B\comp F(f)$.
\end{indicacio}
\end{exercici}

\begin{exercici}
Sigui $\eta\colon F\Rightarrow G$ una transformació natural i $K\colon\cat{D}\to\cat{E}$ un functor. Comprova que els morfismes $K(\eta_A)$ són les components d'una transformació natural $K\eta\colon K\comp F\Rightarrow K\comp G$ (composició d'una transformació natural amb un functor per l'esquerra). \hfill{\normalfont$\star$}
\begin{indicacio}
Aplica $K$ al quadrat de naturalitat de $\eta$ i usa la functorialitat de $K$: de $G(f)\comp\eta_A=\eta_B\comp F(f)$ dedueix $K(G(f))\comp K(\eta_A)=K(\eta_B)\comp K(F(f))$.
\end{indicacio}
\end{exercici}

\begin{exercici}
Siguin $\eta\colon F\Rightarrow G$, amb $F,G\colon\cat{C}\to\cat{D}$, i $\theta\colon K\Rightarrow L$, amb $K,L\colon\cat{D}\to\cat{E}$. Comprova que la composició horitzontal amb una transformació identitat es redueix a la composició amb un functor:
\[
\theta\ast\id_F=\theta F,
\qquad\qquad
\id_K\ast\eta=K\eta.
\]
Dedueix que la composició horitzontal generalitza les dues composicions d'una transformació natural amb un functor. \hfill{\normalfont$\star$}
\begin{indicacio}
Calcula les components amb la definició, fent servir l'expressió de $(\theta\ast\eta)_A$ que més convingui a cada cas, i recorda que un functor envia identitats a identitats.
\end{indicacio}
\end{exercici}

\begin{exercici}
Siguin $\alpha\colon F\Rightarrow G$ i $\beta\colon G\Rightarrow H$, amb $F,G,H\colon\cat{C}\to\cat{D}$, i $\sigma\colon K\Rightarrow L$ i $\tau\colon L\Rightarrow M$, amb $K,L,M\colon\cat{D}\to\cat{E}$. Escriu el tipus de cada costat i demostra la \textbf{llei d'intercanvi}
\[
(\tau\comp\sigma)\ast(\beta\comp\alpha)=(\tau\ast\beta)\comp(\sigma\ast\alpha)\colon K\comp F\Rightarrow M\comp H.
\]
\hfill{\normalfont$\star\star$}
\begin{indicacio}
Calcula les components en un objecte $A$ de les dues bandes i fes servir la naturalitat de $\sigma$ aplicada al morfisme $\beta_A$.
\end{indicacio}
\end{exercici}

\section{Exemples de transformacions naturals}

\begin{exercici}
En l'exemple sintaxi/semàntica de la secció~\ref{sec:cap-transformacio}, escriu explícitament l'equació de naturalitat de la interpretació $\eta$ per a un canvi de variables $f\colon X\to Y$, i explica en paraules què diuen els dos costats. Comprova-la per a la fórmula $\varphi=p\lor\lnot q$, amb $X=\{p,q\}$, en dos casos:
\begin{enumerate}
  \item el reanomenament $f\colon p\mapsto r,\ q\mapsto s$, amb $Y=\{r,s\}$;
  \item la substitució $f\colon p\mapsto r,\ q\mapsto r$, amb $Y=\{r\}$, que identifica les dues variables.
\end{enumerate}
En el segon cas, $\varphi$ no és una tautologia, però $\varphi[f]$ sí que ho és. Explica per què això no contradiu la naturalitat. \hfill{\normalfont$\star$}
\begin{indicacio}
Avalua tots dos camins del quadrat sobre una valoració qualsevol $w\colon Y\to\{0,1\}$: l'un passa per $w\models\varphi[f]$ i l'altre per $(w\comp f)\models\varphi$. Al segon cas, observa quines valoracions de $X$ són de la forma $w\comp f$.
\end{indicacio}
\end{exercici}

\begin{exercici}
Sigui $\cat{C}$ una categoria localment petita i siguin $X,Y$ dos objectes seus. Comprova que cada morfisme $h\colon Y\to X$ indueix una transformació natural entre els homfunctors covariants
\[
\Hom(X,-)\Rightarrow\Hom(Y,-).
\]
\hfill{\normalfont$\star$}
\begin{indicacio}
Defineix la component a $A$ per precomposició amb $h$ i usa l'associativitat per comprovar la naturalitat.
\end{indicacio}
\end{exercici}

\begin{exercici}
Considera el functor \enquote{llista} $L\colon\cat{Set}\to\cat{Set}$, que envia cada conjunt $X$ al conjunt $L(X)$ de les llistes finites d'elements de $X$. Comprova que l'aplicació que envia cada element $x$ a la llista d'un sol element $[x]$ defineix una transformació natural $\id_{\cat{Set}}\Rightarrow L$. \hfill{\normalfont$\star\star$}
\begin{indicacio}
El component a $X$ és $\eta_X\colon X\to L(X)$, $x\mapsto[x]$. La naturalitat respecte d'una aplicació $f\colon X\to Y$ diu que aplicar $f$ i després formar la llista d'un element coincideix amb formar la llista i després aplicar $f$ terme a terme.
\end{indicacio}
\end{exercici}

\section{Isomorfismes naturals}

\begin{exercici}
Considera el functor d'\textbf{intercanvi} $\sigma_{\cat{C},\cat{D}}\colon\cat{C}\times\cat{D}\to\cat{D}\times\cat{C}$ definit, sobre objectes i sobre morfismes, per
\[
\sigma_{\cat{C},\cat{D}}(A,B)=(B,A),
\qquad
\sigma_{\cat{C},\cat{D}}(f,g)=(g,f).
\]
Comprova que és un functor i que la seva composició amb l'intercanvi en sentit contrari,
\[
\sigma_{\cat{D},\cat{C}}\comp\sigma_{\cat{C},\cat{D}}=\id_{\cat{C}\times\cat{D}},
\]
és \emph{estrictament} la identitat (i igualment en l'altre sentit). Dedueix que $\sigma_{\cat{C},\cat{D}}$ és un \emph{isomorfisme de categories}, no només una equivalència, i explica per què aquí no cal passar per un isomorfisme natural com $\id\cong\sigma\comp\sigma$. \hfill{\normalfont$\star$}
\begin{indicacio}
Comprova la functorialitat component a component i calcula $\sigma_{\cat{D},\cat{C}}\comp\sigma_{\cat{C},\cat{D}}$ sobre objectes i sobre morfismes. Per a l'última pregunta, compara la definició d'isomorfisme de categories amb la d'equivalència.
\end{indicacio}
\end{exercici}

\begin{exercici}
Sigui $\cat{C}$ una categoria amb \emph{almenys un objecte}. Comprova que dos functors constants $\Delta_D,\Delta_{D'}\colon\cat{C}\to\cat{D}$ (amb valors $D$ i $D'$) són naturalment isomorfs si i només si $D\cong D'$ a $\cat{D}$. Explica què falla si $\cat{C}$ és buida. \hfill{\normalfont$\star\star$}
\begin{indicacio}
Per a ($\Rightarrow$), mira què és una component en un objecte qualsevol de $\cat{C}$. Per a ($\Leftarrow$), prova amb una família constant i escriu-ne el quadrat de naturalitat. Per al cas buit, pensa quantes transformacions naturals hi ha entre dos functors amb domini buit.
\end{indicacio}
\end{exercici}

\begin{exercici}
Sigui $\mathbb{Z}/2\mathbb{Z}=\{1,g\}$ (amb $g^2=1$) vist com a categoria d'un sol objecte $\cat{B}(\mathbb{Z}/2\mathbb{Z})$, i considera dos functors $F,G\colon\cat{B}(\mathbb{Z}/2\mathbb{Z})\to\cat{Set}$ que envien l'únic objecte al conjunt $\{0,1,2\}$. En $F$, el generador $g$ actua per la identitat; en $G$, actua per l'aplicació $\sigma$ que intercanvia $0$ i $1$ i deixa fix el $2$.
\begin{enumerate}
  \item Troba totes les transformacions naturals $F\Rightarrow G$.
  \item Dedueix que $F$ i $G$ no són naturalment isomorfs, tot i que $F(\ast)=G(\ast)$.
  \item Hi ha transformacions naturals $G\Rightarrow F$? N'hi ha alguna de bijectiva?
\end{enumerate}
Compara el resultat amb l'exemple de l'observació~\ref{obs:iso-no-natural}, on no hi havia \emph{cap} transformació natural. \hfill{\normalfont$\star\star$}
\begin{indicacio}
Escriu la condició de naturalitat per al generador $g$ i estudia si pot existir una component bijectiva que la satisfaci.
\end{indicacio}
\end{exercici}

\section{La categoria de functors}

\begin{exercici}
Comprova que si $\cat{C}$ és una categoria petita i $\cat{D}$ és localment petita, aleshores la categoria de functors $\cat{D}^{\cat{C}}$ és localment petita. \hfill{\normalfont$\star$}
\begin{indicacio}
Una transformació natural està determinada pels seus components, un per a cada objecte de $\cat{C}$; com que $\cat{C}$ és petita i cada $\Hom_{\cat{D}}(F(A),G(A))$ és un conjunt, la col·lecció de transformacions naturals $F\Rightarrow G$ és un subconjunt d'un producte de conjunts indexat per un conjunt.
\end{indicacio}
\end{exercici}

\begin{exercici}
Descriu la categoria de functors $\cat{Set}^{\cat{2}}$ i comprova que els seus isomorfismes són els quadrats commutatius amb les dues fletxes verticals bijectives. \hfill{\normalfont$\star\star$}
\begin{indicacio}
Un objecte és una aplicació $f\colon A\to B$; un morfisme $(u,v)$ és un quadrat commutatiu. Un isomorfisme a la categoria de fletxes és un quadrat amb $u$ i $v$ bijectives, ja que la inversa s'obté invertint-les.
\end{indicacio}
\end{exercici}

\section{Equivalència de categories}

\begin{exercici}
Comprova que \enquote{ser equivalents} és una relació d'equivalència entre categories. És a dir, comprova que tota categoria és equivalent a ella mateixa, que la relació és simètrica i que la composició de dues equivalències és una equivalència. \hfill{\normalfont$\star$}
\begin{indicacio}
Per a la reflexivitat, pren el functor identitat. Per a la simetria, intercanvia els papers de $F$ i $G$. Per a la transitivitat, comprova que la composició de functors plens, fidels i essencialment exhaustius torna a tenir les tres propietats, o compon directament els functors i els isomorfismes naturals.
\end{indicacio}
\end{exercici}

\begin{exercici}
Sigui $\cat{C}$ una categoria en què tot objecte és isomorf a un objecte fix $D_0$ (és a dir, hi ha un sol objecte llevat d'isomorfisme; \emph{no} se suposa que tots els morfismes siguin isomorfismes). Comprova que $\cat{C}$ és equivalent al monoide $\mathrm{End}(D_0)=\Hom(D_0,D_0)$ vist com a categoria d'un sol objecte. \hfill{\normalfont$\star\star$}
\begin{indicacio}
Considera la subcategoria plena de $\cat{C}$ que té $D_0$ com a únic objecte, i fes servir la caracterització de les equivalències. Per què no cal cap hipòtesi sobre la invertibilitat dels morfismes?
\end{indicacio}
\end{exercici}

\begin{exercici}
Dona un exemple de dues categories que siguin \emph{equivalents} però no \emph{isomorfes}, i explica per què no són isomorfes. \hfill{\normalfont$\star\star$}
\begin{indicacio}
Un isomorfisme de categories indueix una bijecció entre els objectes, mentre que una equivalència no. Busca dues categories equivalents amb un nombre diferent d'objectes; n'hi ha un exemple amb categories molt petites.
\end{indicacio}
\end{exercici}

\ifpublicacioparcial\else
\chapter[Representabilitat i Yoneda]{\texorpdfstring{Propietats universals,\\ representabilitat i Yoneda}{Propietats universals, representabilitat i Yoneda}}

Els exercicis següents completen els resolts a la part de Pràctica. Cada enunciat va acompanyat d'una indicació breu. La marca \enquote{$\star$} assenyala un exercici de consolidació i \enquote{$\star\star$} un d'ampliació.

\section{Objectes inicials i terminals}

\begin{exercici}
Comprova que si una categoria té objecte inicial, aquest és únic llevat d'isomorfisme únic, repetint l'argument fet per als objectes terminals. \hfill{\normalfont$\star$}
\begin{indicacio}
Dualitza l'argument dels objectes terminals: compara les composicions de les úniques fletxes entre dos objectes inicials amb les identitats.
\end{indicacio}
\end{exercici}

\begin{exercici}
Determina els objectes inicial i terminal a $\cat{Set}$, a $\cat{Pos}$ i en un preordre $(P,\leq)$ vist com a categoria. \hfill{\normalfont$\star$}
\begin{indicacio}
Per a cada categoria, pregunta't quins objectes reben exactament una fletxa des de cada objecte, i quins n'envien exactament una cap a cada objecte. En un preordre, tradueix-ho en termes de l'ordre.
\end{indicacio}
\end{exercici}

\begin{exercici}[Ampliació, per a lectors amb àlgebra]
Comprova que la categoria $\cat{Field}$ dels cossos i homomorfismes de cossos \emph{no} té objecte inicial ni terminal. \hfill{\normalfont$\star$}
\begin{indicacio}
Un homomorfisme de cossos preserva $0$, $1$ i les operacions, i és sempre injectiu. No hi ha cap cos que rebi un morfisme de tots els altres (pensa en característiques diferents), ni cap que n'enviï a tots.
\end{indicacio}
\end{exercici}

\section{Propietats universals com a representabilitat}

\begin{exercici}
Comprova que un objecte $0$ és inicial si i només si el functor covariant $\Hom(0,-)$ és naturalment isomorf al functor constant $\Delta_{\{\ast\}}\colon\cat{C}\to\cat{Set}$ (exemple~\ref{ex:functor-constant}). \hfill{\normalfont$\star$}
\begin{indicacio}
Compara el nombre d'elements de cada $\Hom(0,X)$ amb el de $\{\ast\}$, i pregunta't quines condicions de naturalitat queden per comprovar.
\end{indicacio}
\end{exercici}

\begin{exercici}
Sigui $\cat{C}$ localment petita amb objecte terminal $1$. Comprova que $\Hom(1,-)\colon\cat{C}\to\cat{Set}$ és un functor i demostra que envia l'objecte terminal a un conjunt terminal. Mostra amb un exemple que, tot i això, $\Hom(1,-)$ pot no distingir objectes no isomorfs: pot enviar-los a conjunts amb el mateix nombre d'elements, o fins i tot al mateix conjunt. \hfill{\normalfont$\star$}
\begin{indicacio}
Calcula $\Hom(1,1)$. Per a l'exemple, pensa en un preordre amb un element màxim i dos elements incomparables per sota.
\end{indicacio}
\end{exercici}

\begin{exercici}
Considera el prefeix de parts $\mathcal P\colon\cat{Set}\op\to\cat{Set}$ i les dues parelles candidates $(\{0,1\},\{1\})$ i $(\{0,1\},\varnothing)$. Quina de les dues és una representació de $\mathcal P$? Justifica per què l'altra no ho és. \hfill{\normalfont$\star$}
\begin{indicacio}
Per a la segona parella, intenta obtenir el subconjunt $\{\ast\}$ d'un conjunt d'un sol element com una antiimatge de $\varnothing$.
\end{indicacio}
\end{exercici}

\section{El lema de Yoneda}

\begin{exercici}
Escriu l'enunciat covariant del lema de Yoneda: per a un functor covariant $Q\colon\cat{C}\to\cat{Set}$ i un objecte $A$, hi ha una bijecció natural
\[
\Nat\bigl(\Hom(A,-),Q\bigr)\;\cong\;Q(A).
\]
Identifica la bijecció. \hfill{\normalfont$\star\star$}
\begin{indicacio}
Reprodueix la demostració del cas contravariant a $\cat{C}\op$ i avalua la transformació en $\id_A$.
\end{indicacio}
\end{exercici}

\begin{exercici}
Fes servir el lema de Yoneda per calcular
\[
\Nat\bigl(\Hom(-,A),\Hom(-,B)\bigr).
\]
Comprova que la bijecció resultant és precisament l'acció de la immersió de Yoneda sobre els morfismes $A\to B$. \hfill{\normalfont$\star\star$}
\begin{indicacio}
Aplica Yoneda al prefeix representable $P=\Hom(-,B)$ i identifica què correspon a un morfisme $A\to B$.
\end{indicacio}
\end{exercici}

\begin{exercici}
Siguin $(A,u)$ i $(B,v)$ dues representacions d'un mateix prefeix $P$. Demostra que hi ha un únic isomorfisme $f\colon A\to B$ amb $P(f)(v)=u$. Explica per què això no diu que $A$ tingui un únic automorfisme. \hfill{\normalfont$\star$}
\begin{indicacio}
Fes servir la universalitat de $v$ per obtenir $f$, i la de $u$ per obtenir $g\colon B\to A$ amb $P(g)(u)=v$. Aplica després la unicitat a $g\comp f$ i a $\id_A$.
\end{indicacio}
\end{exercici}

\begin{exercici}
Comprova que si un prefeix $P$ és representable, aleshores l'element $x\in P(A)$ que correspon a la identitat sota l'isomorfisme $\Hom(-,A)\cong P$ té la propietat universal següent: per a cada objecte $X$ i cada $y\in P(X)$, hi ha un únic morfisme $\varphi\colon X\to A$ amb $P(\varphi)(x)=y$. (Aquest $x$ s'anomena \emph{element universal}.) \hfill{\normalfont$\star$}
\begin{indicacio}
Escriu la naturalitat de l'isomorfisme $\Hom(-,A)\cong P$ per a un morfisme $\varphi\colon X\to A$, avaluada en $\id_A$.
\end{indicacio}
\end{exercici}

\section{Conseqüències}

\begin{exercici}
Sigui $\cat{C}$ una categoria petita. Comprova que la immersió de Yoneda $\yon\colon\cat{C}\to\cat{Set}^{\cat{C}\op}$ \emph{no} és, en general, essencialment exhaustiva, tot i ser plena i fidel. Quins prefeixos són de la forma $\Hom(-,A)$? \hfill{\normalfont$\star\star$}
\begin{indicacio}
Busca una propietat que tot prefeix representable $\Hom(-,A)$ compleix necessàriament en l'objecte $A$, i un prefeix que no la compleixi.
\end{indicacio}
\end{exercici}

\begin{exercici}
Dedueix del corol·lari de Yoneda que si dos objectes tenen conjunts de morfismes \enquote{naturalment iguals} des de tot objecte, aleshores són isomorfs; i explica per què això formalitza l'eslògan \enquote{un objecte es coneix per les fletxes que hi arriben}. \hfill{\normalfont$\star\star$}
\begin{indicacio}
Tradueix \enquote{naturalment iguals des de tot objecte} en termes de prefeixos representables.
\end{indicacio}
\end{exercici}

\begin{exercici}
Enuncia i comprova la versió dual del corol·lari: $A\cong B$ si i només si $\Hom(A,-)\cong\Hom(B,-)$ com a functors covariants. \hfill{\normalfont$\star\star$}
\begin{indicacio}
Aplica el corol·lari a $\cat{C}\op$, o repeteix l'argument amb la immersió covariant $\cat{C}\op\to\cat{Set}^{\cat{C}}$, que també és plena i fidel.
\end{indicacio}
\end{exercici}

\chapter[Límits i colímits]{Límits i colímits}

Els exercicis següents completen els resolts a la part de Pràctica. Cada enunciat va acompanyat d'una indicació breu. La marca \enquote{$\star$} assenyala un exercici de consolidació i \enquote{$\star\star$} un d'ampliació.

\section{Diagrames i cons}

\begin{exercici}
Descriu la categoria índex $\cat{I}$ i el diagrama corresponent per a cadascun d'aquests límits: producte binari, equalitzador, pullback i objecte terminal. \hfill{\normalfont$\star$}
\begin{indicacio}
Identifica, per a cada propietat universal, la forma mínima de diagrama que la descriu: quants objectes hi intervenen i quines fletxes cal que commutin.
\end{indicacio}
\end{exercici}

\begin{exercici}
Comprova que un morfisme de cons és un cas particular de morfisme a la categoria $\mathrm{Con}(D)$, i que la composició de morfismes de cons torna a ser-ne un. \hfill{\normalfont$\star$}
\begin{indicacio}
Si $f\colon A\to B$ i $g\colon B\to C$ satisfan $\beta_i\comp f=\alpha_i$ i $\gamma_i\comp g=\beta_i$, aleshores $\gamma_i\comp(g\comp f)=\alpha_i$.
\end{indicacio}
\end{exercici}

\section{Límits}

\begin{exercici}
Comprova que en un preordre vist com a categoria, el producte de dos elements és el seu ínfim i el coproducte és el seu suprem, quan existeixen. \hfill{\normalfont$\star$}
\begin{indicacio}
Tradueix la propietat universal del producte a desigualtats.
\end{indicacio}
\end{exercici}

\begin{exercici}
Comprova que dos límits d'un mateix diagrama són isomorfs per un únic isomorfisme compatible amb totes les projeccions. \hfill{\normalfont$\star$}
\begin{indicacio}
Treballa a la categoria $\mathrm{Con}(D)$ i fes servir la unicitat dels objectes terminals.
\end{indicacio}
\end{exercici}

\section{Productes i equalitzadors}

\begin{exercici}
Comprova, llevat d'isomorfisme, la commutativitat i l'associativitat del producte:
\[
A\times B\cong B\times A,\quad (A\times B)\times C\cong A\times(B\times C).
\] \hfill{\normalfont$\star$}
\begin{indicacio}
Fes servir la unicitat d'un objecte definit per una mateixa propietat universal.
\end{indicacio}
\end{exercici}

\begin{exercici}
Comprova que si $\cat{C}$ té objecte terminal $1$, aleshores $A\times 1\cong A$ per a tot objecte $A$. \hfill{\normalfont$\star$}
\begin{indicacio}
Compara les propietats universals de $A\times 1$ i de $A$.
\end{indicacio}
\end{exercici}

\begin{exercici}
Comprova que a $\cat{Grp}$ l'equalitzador de $f,g\colon G\to H$ és el subgrup $\{x\in G\mid f(x)=g(x)\}$. \hfill{\normalfont$\star$}
\begin{indicacio}
Comprova que aquest subconjunt és un subgrup i que la injecció satisfà la propietat universal de l'equalitzador; el functor oblit cap a $\cat{Set}$ ajuda a veure-ho.
\end{indicacio}
\end{exercici}

\section{Pullbacks}

\begin{exercici}
Comprova que en un pullback, si $g\colon B\to C$ és un monomorfisme, aleshores $p_1\colon P\to A$ també ho és. (Es diu que els monomorfismes són \emph{estables per pullback}.) \hfill{\normalfont$\star\star$}
\begin{indicacio}
Usa la propietat universal del pullback i la cancel·lació de $g$. Aquesta estabilitat serà essencial per a la teoria de subobjectes.
\end{indicacio}
\end{exercici}

\begin{exercici}
Comprova que un quadrat amb $\id$ als dos costats,
\[
\begin{array}{ccc}
A & \xrightarrow{\;\id\;} & A\\[-2pt]
\downarrow{\scriptstyle\id} & & \downarrow{\scriptstyle f}\\[-2pt]
A & \xrightarrow{\;f\;} & B
\end{array}
\]
és un pullback si i només si $f$ és un monomorfisme. \hfill{\normalfont$\star\star$}
\begin{indicacio}
El pullback d'aquest quadrat a $\cat{Set}$ és $\{(a,a')\mid f(a)=f(a')\}$; que la projecció sigui iso equival a la injectivitat de $f$. En general, tradueix-ho amb la propietat universal.
\end{indicacio}
\end{exercici}

\section{Colímits}

\begin{exercici}
Comprova que tot coequalitzador és un epimorfisme, dualitzant la demostració que tot equalitzador és un monomorfisme. \hfill{\normalfont$\star$}
\begin{indicacio}
Dualitza l'argument de la Pràctica que prova que tot equalitzador és mono.
\end{indicacio}
\end{exercici}

\begin{exercici}
Comprova que el pushout de $f\colon C\to A$ i $g\colon C\to B$ a $\cat{Set}$ és $(A+B)/{\sim}$, on $\sim$ enganxa $f(c)$ amb $g(c)$ per a cada $c\in C$. \hfill{\normalfont$\star$}
\begin{indicacio}
És el dual del pullback. Pren la unió disjunta $A+B$ i quocienta per la relació generada per $\iota_1(f(c))\sim\iota_2(g(c))$; comprova la propietat universal.
\end{indicacio}
\end{exercici}

\section{Límits finits}

\begin{exercici}
Sigui $\cat{C}$ una categoria amb productes binaris i pullbacks, i siguin $f,g\colon A\to B$. La demostració del teorema~\ref{teo:limits-finits} obté l'equalitzador de $f$ i $g$ com un pullback de dues fletxes $A\to A\times B$. Comprova que també és un pullback de la diagonal $\langle\id_B,\id_B\rangle\colon B\to B\times B$ al llarg de $\langle f,g\rangle\colon A\to B\times B$. \hfill{\normalfont$\star\star$}
\begin{indicacio}
Escriu què és un con sobre aquest pullback i tradueix la condició de con component a component, amb les dues projeccions de $B\times B$.
\end{indicacio}
\end{exercici}

\begin{exercici}
Comprova directament que una categoria amb productes binaris i equalitzadors té tots els pullbacks, amb una construcció més econòmica que la general de la demostració del teorema~\ref{teo:limits-finits}. \hfill{\normalfont$\star$}
\begin{indicacio}
Expressa el pullback com un equalitzador de dues fletxes que surten d'un producte.
\end{indicacio}
\end{exercici}

\section{Límits en categories de functors}

\begin{exercici}
Sigui $\mathbf E$ el graf amb dos vèrtexs $0,1$ i una sola aresta $e\colon 0\to 1$. Fent servir que $\cat{Graph}\cong\cat{Set}^{\cat{P}}$ i que els límits i colímits s'hi calculen punt a punt (proposició~\ref{prop:limits-punt-a-punt}), descriu el producte $\mathbf E\times\mathbf E$ i el coproducte $\mathbf E+\mathbf E$: els seus vèrtexs, les seves arestes, i l'origen i el destí de cada aresta. Dibuixa'ls. És $\mathbf E\times\mathbf E$ un quadrat? \hfill{\normalfont$\star$}
\begin{indicacio}
Calcula per separat el conjunt de vèrtexs i el d'arestes, i defineix origen i destí component a component.
\end{indicacio}
\end{exercici}

\section{Preservació i reflexió}

\begin{exercici}
Comprova que tot functor ple i fidel reflecteix límits. Dualment, reflecteix colímits. \hfill{\normalfont$\star\star$}
\begin{indicacio}
Aixeca l'existència de la factorització amb la plenitud i prova'n la unicitat amb la fidelitat.
\end{indicacio}
\end{exercici}

\begin{exercici}
Comprova que el prefeix representable
\[
\Hom(-,B)\colon\cat{C}\op\to\cat{Set}
\]
envia colímits de $\cat{C}$ a límits de $\cat{Set}$. \hfill{\normalfont$\star\star$}
\begin{indicacio}
És el dual del fet que $\Hom(A,-)$ preserva límits: un cocon amb vèrtex $B$ sobre $D$ és un con sobre $\Hom(D_-,B)$, i la universalitat es correspon.
\end{indicacio}
\end{exercici}

\chapter[Adjuncions]{Adjuncions}

Els exercicis següents completen els resolts a la part de Pràctica. Cada enunciat va acompanyat d'una indicació breu. La marca \enquote{$\star$} assenyala un exercici de consolidació i \enquote{$\star\star$} un d'ampliació.

\section{Adjuncions entre preordres}

\begin{exercici}
Comprova que en una adjunció entre preordres $f\dashv g$, l'aplicació $f$ preserva els suprems que existeixin i $g$ preserva els ínfims. \hfill{\normalfont$\star$}
\begin{indicacio}
És el cas particular, per a preordres, que els adjunts per l'esquerra preserven colímits (suprems) i els adjunts per la dreta preserven límits (ínfims). Demostra-ho directament amb la definició $f(x)\leq y\Leftrightarrow x\leq g(y)$.
\end{indicacio}
\end{exercici}

\begin{exercici}
Sigui $f\colon P\to Q$ monòtona entre \textbf{reticles complets}, és a dir, preordres en què tot subconjunt té suprem i ínfim. Comprova que $f$ té adjunt per la dreta si i només si preserva tots els suprems. \hfill{\normalfont$\star\star$}
\begin{indicacio}
Si $f$ preserva suprems, defineix $g(y)$ com el suprem d'un conjunt adequat d'elements de $P$.
\end{indicacio}
\end{exercici}

\section{Adjuncions entre categories}

\begin{exercici}\label{exr:adj-unitat-universal}
Comprova que si $F\dashv G$, aleshores per a cada objecte $A$ el morfisme $\eta_A\colon A\to G(F(A))$ és universal: tot $\psi\colon A\to G(B)$ factoritza de manera única com $G(\varphi)\comp\eta_A$. \hfill{\normalfont$\star\star$}
\begin{indicacio}
Fes servir la bijecció $\theta$ i la fórmula del transposat en termes de la unitat.
\end{indicacio}
\end{exercici}

\begin{exercici}
Sigui $\cat{C}$ una categoria amb coproductes binaris i $\Delta\colon\cat{C}\to\cat{C}\times\cat{C}$ el functor diagonal, $A\mapsto(A,A)$. Dualitzant l'exercici resolt de la Pràctica sobre $\Delta\dashv\times$, comprova que el coproducte és adjunt per l'esquerra de $\Delta$, $(+)\dashv\Delta$, i descriu-ne la unitat i la counitat. \hfill{\normalfont$\star\star$}
\begin{indicacio}
Fes servir la propietat universal del coproducte; la unitat i la counitat són les duals de la diagonal i les projeccions.
\end{indicacio}
\end{exercici}

\section{Unitat i counitat}

\begin{exercici}\label{exr:adj-fletxes-universals}
Sigui $G\colon\cat{D}\to\cat{C}$ un functor. Suposem que, per a cada objecte $A$ de $\cat{C}$, hem triat un objecte $F(A)$ de $\cat{D}$ i una fletxa universal $\eta_A\colon A\to G(F(A))$, en el sentit que tot $\psi\colon A\to G(B)$ factoritza de manera única com $G(\varphi)\comp\eta_A$. Demostra que aquestes dades determinen de manera única l'acció de $F$ sobre morfismes i defineixen un adjunt per l'esquerra $F\dashv G$. \hfill{\normalfont$\star\star$}
\begin{indicacio}
Defineix $F$ sobre objectes i sobre morfismes amb la propietat universal, i fes servir la unicitat per obtenir la functorialitat.
\end{indicacio}
\end{exercici}

\begin{exercici}
Per a l'adjunció $\Free\dashv U$ entre $\cat{Graph}$ i $\cat{Cat}$ (exemple~\ref{ex:free-u-adjuncio}), comprova que la counitat
\[
\varepsilon_{\cat{C}}\colon\Free(U(\cat{C}))\longrightarrow\cat{C}
\]
és la identitat sobre els objectes i és un functor ple. Mostra amb un exemple que, en general, no és fidel. \hfill{\normalfont$\star$}
\begin{indicacio}
Per veure que és ple, pensa en els camins de longitud~$1$. Per veure que no és fidel, busca dos camins diferents amb la mateixa composta.
\end{indicacio}
\end{exercici}

\section{Identitats triangulars}

\begin{exercici}
Comprova la segona identitat triangular,
\[
G(\varepsilon_B)\comp\eta_{G(B)}=\id_{G(B)},
\]
dualitzant la demostració de la primera. \hfill{\normalfont$\star$}
\begin{indicacio}
Dualitza la demostració de la Teoria, fent servir la fórmula del transposat en termes de la unitat.
\end{indicacio}
\end{exercici}

\begin{exercici}\label{exr:adj-triangulars-reciproc}
Comprova que dues transformacions naturals $\eta\colon\id\Rightarrow GF$ i $\varepsilon\colon FG\Rightarrow\id$ que satisfan les identitats triangulars defineixen una adjunció $F\dashv G$. \hfill{\normalfont$\star\star$}
\begin{indicacio}
Defineix els dos transposats amb $\eta$ i $\varepsilon$, i fes servir les identitats triangulars per veure que són inversos.
\end{indicacio}
\end{exercici}

\begin{exercici}
Siguin $F\colon\cat{C}\to\cat{D}$, $G\colon\cat{D}\to\cat{C}$, $F'\colon\cat{D}\to\cat{E}$ i $G'\colon\cat{E}\to\cat{D}$ functors entre categories localment petites, amb $F\dashv G$ i $F'\dashv G'$. Demostra que les adjuncions es componen:
\[
F'\comp F\;\dashv\;G\comp G'.
\]
Descriu la unitat de l'adjunció composta en termes de les unitats $\eta$ i $\eta'$. \hfill{\normalfont$\star$}
\begin{indicacio}
Encadena les dues bijeccions d'homconjunts, passant per $\cat{D}$, i comprova que la composta és natural en les dues variables. Per a la unitat, transposa la identitat de $F'(F(A))$.
\end{indicacio}
\end{exercici}

\begin{exercici}
Sigui $F\dashv G$ una adjunció amb unitat $\eta$ i counitat $\varepsilon$. Demostra que, si $\eta$ i $\varepsilon$ són isomorfismes naturals, aleshores $F$ és una equivalència de categories amb quasiinvers $G$ (definició~\ref{def:equivalencia}), i que en aquest cas $G$ també és adjunt per l'esquerra de $F$. \hfill{\normalfont$\star$}
\begin{indicacio}
Compara les dades amb la definició d'equivalència. Per a la segona part, inverteix $\eta$ i $\varepsilon$ i comprova les identitats triangulars de l'adjunció $G\dashv F$ a partir de les de $F\dashv G$.
\end{indicacio}
\end{exercici}

\section{Els adjunts i els límits}

\begin{exercici}
Comprova que el prefeix de parts $\mathcal P\colon\cat{Set}\op\to\cat{Set}$ envia coproductes a productes: $\mathcal P(A+B)\cong\mathcal P(A)\times\mathcal P(B)$. \hfill{\normalfont$\star$}
\begin{indicacio}
Descriu un subconjunt de $A+B$ a partir de les seves interseccions amb les dues components.
\end{indicacio}
\end{exercici}

\begin{exercici}[Ampliació, per a lectors amb topologia]
Sigui $U\colon\cat{Top}\to\cat{Set}$ el functor oblit. Comprova que té \emph{dos} adjunts, un per l'esquerra i un per la dreta, i digues quins són. \hfill{\normalfont$\star$}
\begin{indicacio}
Pensa en les dues topologies extremes sobre un conjunt.
\end{indicacio}
\end{exercici}

\section{Exponencials}

\begin{exercici}
Comprova que a $\cat{Set}$ es compleix $C^{A+B}\cong C^A\times C^B$ i $(C\times D)^A\cong C^A\times D^A$, i interpreta-ho en termes d'adjuncions. \hfill{\normalfont$\star$}
\begin{indicacio}
Per a la primera, fes servir que $\Hom(A+B,C)\cong\Hom(A,C)\times\Hom(B,C)$, és a dir, que el representable contravariant $\Hom(-,C)$ envia coproductes a productes. Per a la segona, fes servir que $(-)^A$ és adjunt per la dreta de $-\times A$ i, per tant, preserva productes.
\end{indicacio}
\end{exercici}

\begin{exercici}
Comprova que la currificació $\Hom(A\times B,C)\cong\Hom(A,C^B)$ és natural en les tres variables $A,B,C$. \hfill{\normalfont$\star\star$}
\begin{indicacio}
Comprova la commutativitat amb la precomposició en $A$, la postcomposició en $C$ i l'acció en $B$, aplicant les definicions de $\tilde f$ i $\hat g$.
\end{indicacio}
\end{exercici}

\chapter[Cartesianes tancades]{\texorpdfstring{Categories cartesianes\\ tancades}{Categories cartesianes tancades}}

Els exercicis següents completen els resolts a la part de Pràctica. Cada enunciat va acompanyat d'una indicació breu. La marca \enquote{$\star$} assenyala un exercici de consolidació i \enquote{$\star\star$} un d'ampliació, sensiblement més exigent.

\section{Productes finits}

\begin{exercici}
Comprova que en una categoria cartesiana el producte és associatiu i commutatiu llevat d'isomorfisme: $(A\times B)\times C\cong A\times(B\times C)$ i $A\times B\cong B\times A$. \hfill{\normalfont$\star$}
\begin{indicacio}
Construeix els morfismes en tots dos sentits amb aparellaments de projeccions i comprova que en són inversos usant la unicitat de la propietat universal. No cal calcular amb elements.
\end{indicacio}
\end{exercici}

\begin{exercici}
Demostra que si $\cat{C}$ és cartesiana, aleshores $\cat{C}\op$ té coproductes finits, i descriu-los. \hfill{\normalfont$\star$}
\begin{indicacio}
El producte a $\cat{C}$ és el coproducte a $\cat{C}\op$ per dualització de la propietat universal; el terminal esdevé inicial.
\end{indicacio}
\end{exercici}

\section{Exponencials}

\begin{exercici}
En una CCC, construeix un morfisme canònic \enquote{composició} $C^{B}\times B^{A}\to C^{A}$ i comprova que és associatiu en el sentit adient. \hfill{\normalfont$\star\star$}
\begin{indicacio}
Currifica el morfisme $C^{B}\times B^{A}\times A\to C$ obtingut avaluant primer sobre $A$ i després sobre $B$. L'associativitat es comprova amb la unicitat de la transposada.
\end{indicacio}
\end{exercici}

\begin{exercici}
Comprova que l'assignació $B\mapsto C^{B}$ és \emph{contravariant}: un morfisme $f\colon B'\to B$ indueix $C^{f}\colon C^{B}\to C^{B'}$. \hfill{\normalfont$\star$}
\begin{indicacio}
Currifica
\[
C^{B}\times B'\xrightarrow{\id\times f}C^{B}\times B\xrightarrow{\ev}C.
\]
Això mostra la contravariància en l'exponent $B$; en canvi, per a $B$ fix, el functor $(-)^{B}$ és covariant en $C$.
\end{indicacio}
\end{exercici}

\begin{exercici}
En una CCC amb coproductes, demostra la llei $C^{A+B}\cong C^{A}\times C^{B}$. \hfill{\normalfont$\star\star$}
\begin{indicacio}
Fes servir que $X\times-$ és un adjunt per l'esquerra i, per tant, preserva coproductes (exercici resolt~\ref{pr-distributivitat} de la Pràctica). Després aplica la propietat universal del coproducte i el lema de Yoneda.
\end{indicacio}
\end{exercici}

\begin{exercici}
Sabem que $\cat{Graph}$ és cartesiana tancada, perquè és una categoria de prefeixos. Fent servir la fórmula de l'exponencial de prefeixos i el lema de Yoneda, demostra que els \emph{vèrtexs} de l'exponencial $H^{G}$ de dos grafs corresponen a totes les aplicacions $V_G\to V_H$ entre els conjunts de vèrtexs, sense cap condició sobre les arestes. \hfill{\normalfont$\star\star$}
\begin{indicacio}
Els vèrtexs de $H^{G}$ es corresponen canònicament amb els morfismes de grafs $\mathbf V\to H^{G}$, on $\mathbf V$ és el graf d'un sol vèrtex. Transposa'ls i calcula el graf $\mathbf V\times G$.
\end{indicacio}
\end{exercici}

\begin{exercici}
Sigui $\cat{2}=(0\to1)$. Fent servir l'exemple de $\cat{Cat}$ de la Teoria amb $\cat{A}=\cat{2}$, comprova que un functor $\cat{2}\times\cat{2}\to\cat{C}$ és el mateix que un quadrat commutatiu de $\cat{C}$, i que la seva currificació $\cat{2}\to\cat{C}^{\cat{2}}$ és el mateix que un morfisme de la categoria de fletxes. Què fa l'avaluació sobre aquestes dades? \hfill{\normalfont$\star$}
\begin{indicacio}
Escriu els quatre objectes i les cinc fletxes no identitat de $\cat{2}\times\cat{2}$. Aplica la lectura \enquote{transformació natural $=$ functor $\cat{A}\times\cat{2}\to\cat{C}$} amb $\cat{A}=\cat{2}$.
\end{indicacio}
\end{exercici}

\section{Models i lògica}

\begin{exercici}
Comprova que tota àlgebra de Boole és una àlgebra de Heyting en què $b\Rightarrow c=\neg b\vee c$, i per tant un preordre cartesià tancat. \hfill{\normalfont$\star$}
\begin{indicacio}
Verifica l'adjunció $a\wedge b\leq c\Leftrightarrow a\leq\neg b\vee c$ usant les lleis booleanes. La lògica clàssica és el cas booleà de la intuïcionista.
\end{indicacio}
\end{exercici}

\begin{exercici}
Sigui $\cat{C}$ una CCC. Demostra que el conjunt de morfismes $1\to C^{B}$ està en bijecció amb el de morfismes $B\to C$ (els \enquote{punts} de l'exponencial són les funcions). \hfill{\normalfont$\star$}
\begin{indicacio}
És l'adjunció amb $A=1$ i la identificació $1\times B\cong B$.
\end{indicacio}
\end{exercici}

\begin{exercici}
Interpreta el combinador $\mathsf{S}$, de tipus $(\sigma\to\tau\to\rho)\to(\sigma\to\tau)\to\sigma\to\rho$, com a morfisme d'una CCC. \hfill{\normalfont$\star\star$}
\begin{indicacio}
Descurrifica successivament fins a un morfisme amb domini un producte de tres exponencials i $\llbracket\sigma\rrbracket$, i construeix-lo amb avaluacions i la diagonal. Compara'l amb la demostració intuïcionista de l'esquema d'axioma corresponent.
\end{indicacio}
\end{exercici}

\begin{exercici}
Al preordre cartesià tancat $\mathbb Z\cup\{+\infty\}$ de la Teoria, comprova el \emph{modus ponens} $b\wedge(b\Rightarrow c)\leq c$ i la llei de currificació $b\Rightarrow(c\Rightarrow d)=(b\wedge c)\Rightarrow d$, calculant els dos costats per casos. \hfill{\normalfont$\star$}
\begin{indicacio}
Fes servir la fórmula explícita de la implicació en una cadena, i distingeix els casos segons l'ordre relatiu de $b$, $c$ i $d$.
\end{indicacio}
\end{exercici}

\begin{exercici}
Afegeix a $\mathbb Z\cup\{+\infty\}$ un element mínim $-\infty$. Comprova que s'obté una àlgebra de Heyting, calcula la pseudocomplementació $\neg b=(b\Rightarrow-\infty)$ i decideix si és una àlgebra de Boole. \hfill{\normalfont$\star\star$}
\begin{indicacio}
En una cadena, els suprems finits són màxims. Per a la segona part, calcula $\neg\neg b$ i $b\vee\neg b$ per a un $b$ que no sigui un dels extrems.
\end{indicacio}
\end{exercici}

\chapter[Subobjectes, imatges i factoritzacions]{\texorpdfstring{Subobjectes, imatges\\ i factoritzacions}{Subobjectes, imatges i factoritzacions}}

Els exercicis següents completen els resolts a la part de Pràctica. Cada enunciat va acompanyat d'una indicació breu. La marca \enquote{$\star$} assenyala un exercici de consolidació i \enquote{$\star\star$} un d'ampliació.

\section{Subobjectes}

\begin{exercici}
Descriu $\Sub(A)$ quan $A$ és un objecte d'un preordre vist com a categoria. \hfill{\normalfont$\star$}
\begin{indicacio}
En un preordre hi ha com a molt un morfisme entre dos objectes, i tots són monos. Quins subobjectes té, doncs, $A$? Compara amb el conjunt d'elements $x\leq A$: quan donen dos d'aquests elements el mateix subobjecte? Què passa si el preordre és un ordre parcial?
\end{indicacio}
\end{exercici}

\begin{exercici}
Comprova que en una categoria amb objecte inicial estricte $0$, la injecció $0\rightarrowtail A$ és el subobjecte mínim de $\Sub(A)$. \hfill{\normalfont$\star$}
\begin{indicacio}
\enquote{Estricte} vol dir que tot morfisme cap a $0$ és un isomorfisme. Comprova primer que $0\to A$ és mono, i després que es factoritza a través de qualsevol subobjecte.
\end{indicacio}
\end{exercici}

\begin{exercici}[Ampliació, per a lectors amb àlgebra]
Demostra que a $\cat{Grp}$ els subobjectes d'un grup $G$ es corresponen exactament amb els seus \textbf{subgrups}. \hfill{\normalfont$\star\star$}
\begin{indicacio}
Els monomorfismes a $\cat{Grp}$ són els homomorfismes injectius. Comprova que dos monos amb codomini $G$ són equivalents si i només si tenen la mateixa imatge, i que les imatges de monos són exactament els subgrups.
\end{indicacio}
\end{exercici}

\begin{exercici}
Comprova que un morfisme de grafs $F\colon G'\to G$ és un monomorfisme de $\cat{Graph}$ si i només si les aplicacions $F_V$ i $F_E$ són injectives. Dedueix la descripció dels subobjectes d'un graf com a subgrafs de l'exemple~\ref{ex:subgrafs}. \hfill{\normalfont$\star$}
\begin{indicacio}
Per a una direcció, prova els morfismes que surten del graf d'un sol vèrtex i del graf d'una sola aresta.
\end{indicacio}
\end{exercici}

\section{Reindexació}

\begin{exercici}
Sigui $f\colon A\to B$. Comprova que $f^{*}$ envia el subobjecte màxim de $\Sub(B)$ al màxim de $\Sub(A)$. \hfill{\normalfont$\star$}
\begin{indicacio}
El pullback de $\id_B$ al llarg de $f$ és un mono equivalent a $\id_A$. Interpreta-ho: substituir dins del predicat \enquote{sempre cert} torna a donar \enquote{sempre cert}.
\end{indicacio}
\end{exercici}

\begin{exercici}
Treballa a $\cat{Set}$. Comprova que si $m\colon S\rightarrowtail B$ és un subobjecte (un subconjunt $S\subseteq B$) i $f\colon A\to B$, aleshores $f^{*}[m]$ és el subobjecte \emph{mínim} (buit) de $A$ exactament quan la imatge de $f$ no talla $S$, és a dir, quan $f(A)\cap S=\varnothing$. \hfill{\normalfont$\star$}
\begin{indicacio}
A $\cat{Set}$, $f^{-1}(S)=\varnothing$ si i només si cap element de la imatge de $f$ és a $S$. Formula-ho amb pullbacks: el pullback de $m$ al llarg de $f$ és l'objecte inicial exactament en aquest cas.
\end{indicacio}
\end{exercici}

\section{Interseccions i estructura}

\begin{exercici}
Demostra que la reindexació $f^{*}$ preserva interseccions: $f^{*}([m]\wedge[n])=f^{*}[m]\wedge f^{*}[n]$. \hfill{\normalfont$\star\star$}
\begin{indicacio}
Tant l'ínfim com la reindexació són pullbacks. Usa que els pullbacks es componen: el pullback d'un pullback sobre $B$ al llarg de $f$ es reorganitza en la intersecció dels pullbacks. És un exercici de manipulació de quadrats cartesians.
\end{indicacio}
\end{exercici}

\begin{exercici}
A $\cat{Set}$, el functor de subobjectes $\Sub\colon\cat{Set}\op\to\cat{Poset}$ (proposició~\ref{prop:functor-sub}) i el functor de parts contravariant $\mathcal P$, amb $\mathcal P(A)$ ordenat per inclusió, envien cada conjunt a ordres parcials isomorfs. Comprova que aquests isomorfismes formen un isomorfisme natural $\Sub\cong\mathcal P$. \hfill{\normalfont$\star$}
\begin{indicacio}
Fes servir la imatge d'un monomorfisme per definir les components, i compara la reindexació amb l'antiimatge.
\end{indicacio}
\end{exercici}


\section{Factorització imatge}

\begin{exercici}[Ampliació, per a lectors amb àlgebra]
Comprova que a $\cat{Grp}$ la factorització imatge d'un homomorfisme $f\colon G\to H$ és $G\twoheadrightarrow \Imatge(f)\hookrightarrow H$, on $\Imatge(f)$ és el subgrup imatge. \hfill{\normalfont$\star$}
\begin{indicacio}
Els epimorfismes regulars a $\cat{Grp}$ són les projeccions al quocient $G\to G/N$; la primera part de $f$ és $G\to G/\ker f\cong\Imatge(f)$. Fes servir el primer teorema d'isomorfia.
\end{indicacio}
\end{exercici}

\begin{exercici}
Demostra que en qualsevol categoria amb factoritzacions imatge, la imatge d'una composició satisfà $\Imatge(g\comp f)\leq\Imatge(g)$. \hfill{\normalfont$\star$}
\begin{indicacio}
$g\comp f$ es factoritza a través de $\Imatge(g)$ (via la factorització de $g$), i la imatge és el menor subobjecte pel qual això passa.
\end{indicacio}
\end{exercici}

\begin{exercici}
Comprova que un epimorfisme regular que és alhora monomorfisme és un isomorfisme, i dedueix que en una categoria regular la factorització imatge d'un mono és trivial. \hfill{\normalfont$\star$}
\begin{indicacio}
Un epi regular és el coequalitzador de la seva parella nucli; si a més és mono, la parella nucli és la diagonal, i el coequalitzador d'una parella trivial és una identitat llevat d'isomorfisme.
\end{indicacio}
\end{exercici}

\begin{exercici}
Siguin $H$ i $K$ dos subgrafs d'un graf $G$, i $F\colon G\to G''$ un morfisme de grafs. Descriu la intersecció $H\wedge K$ a $\Sub(G)$ i la imatge $\exists_F(H)$ a $\Sub(G'')$, i comprova que totes dues es calculen component a component, sobre vèrtexs i sobre arestes. \hfill{\normalfont$\star$}
\begin{indicacio}
Fes servir l'exemple~\ref{ex:subgrafs} i comprova que les operacions component a component donen subconjunts tancats per origen i destí.
\end{indicacio}
\end{exercici}

\section{Regularitat}

\begin{exercici}[Ampliació, per a lectors amb àlgebra homològica]
Demostra que tota categoria abeliana és regular. \hfill{\normalfont$\star\star$}
\begin{indicacio}
En una categoria abeliana tot morfisme es factoritza com epi (sobre la coimatge) seguit de mono (la imatge), i coimatge i imatge coincideixen. L'estabilitat per pullback se segueix de l'exactitud. Pots limitar-te a esbossar per què cada condició de la definició es compleix.
\end{indicacio}
\end{exercici}

\begin{exercici}[Ampliació, per a lectors amb topologia]
Exhibeix una aplicació quocient a $\cat{Top}$ el pullback de la qual no és una aplicació quocient, i dedueix-ne que $\cat{Top}$ no és regular. \hfill{\normalfont$\star\star$}
\begin{indicacio}
A $\cat{Top}$, els epimorfismes regulars són les aplicacions quocient (no s'han de confondre amb els epimorfismes, que són les aplicacions contínues exhaustives). Busca una aplicació quocient el pullback de la qual, al llarg d'una aplicació adequada, deixa de ser quocient.
\end{indicacio}
\end{exercici}

\section{Quantificació}

\begin{exercici}
Demostra que $\exists_f$ és monòtona. Suposa, a més, que $\Sub(A)$ té un mínim $\bot_A$ (el predicat \enquote{fals}); demostra, fent servir l'adjunció $\exists_f\dashv f^{*}$, que $\exists_f(\bot_A)$ és el mínim de $\Sub(B)$. \hfill{\normalfont$\star$}
\begin{indicacio}
Per a la monotonia, recorda que els adjunts entre preordres són monòtons. Per al mínim, escriu la condició de l'adjunció per a $\exists_f(\bot_A)\leq[t]$.
\end{indicacio}
\end{exercici}

\begin{exercici}
Usant que $\exists_f$ és adjunt per l'esquerra, demostra que preserva unions (suprems) de subobjectes quan existeixen. \hfill{\normalfont$\star\star$}
\begin{indicacio}
Els adjunts per l'esquerra preserven tots els colímits que existeixin; en un preordre, els suprems són els colímits. Escriu la cadena de bicondicionals de l'adjunció aplicada a un suprem.
\end{indicacio}
\end{exercici}


\chapter[Classificadors de subobjectes i topos]{\texorpdfstring{Classificadors de subobjectes i introducció als topos}{Classificadors de subobjectes i introducció als topos}}

Els exercicis següents completen els resolts a la part de Pràctica. La marca \enquote{$\star$} assenyala un exercici de consolidació i \enquote{$\star\star$} un d'ampliació.

\section{Classificadors}

\begin{exercici}
Comprova que a $\cat{Set}$ el morfisme característic $\chi_U$ d'un subconjunt determina i és determinat per $U$, de manera que $\Sub(E)\cong\{0,1\}^{E}$. \hfill{\normalfont$\star$}
\begin{indicacio}
Associa a cada subconjunt la seva funció característica i comprova que l'aplicació és una bijecció amb inversa $\chi\mapsto\chi^{-1}(1)$.
\end{indicacio}
\end{exercici}

\begin{exercici}
En un topos, sigui $m\colon U\rightarrowtail E$ un subobjecte i $g\colon E'\to E$ un morfisme. Demostra que el morfisme característic de la reindexació $g^{*}m$ és
\[
\chi_{g^{*}m}=\chi_m\comp g.
\]
Interpreta-ho com la compatibilitat entre la representació $\Sub\cong\Hom(-,\Omega)$ i la substitució. \hfill{\normalfont$\star$}
\begin{indicacio}
Compara els dos pullbacks de $\mathsf{true}$ i fes servir el lema d'enganxament de pullbacks.
\end{indicacio}
\end{exercici}

\begin{exercici}
Sigui $E=\{1,\dots,6\}$ i $U=\{3,6\}$, el predicat \enquote{$x$ és múltiple de $3$}. Calcula $\chi_U\colon E\to\{0,1\}$. Sigui $g\colon\{a,b,c,d\}\to E$ l'aplicació amb $g(a)=3$, $g(b)=4$, $g(c)=6$ i $g(d)=3$. Calcula $\chi_U\comp g$ i identifica el subconjunt que classifica. \hfill{\normalfont$\star$}
\begin{indicacio}
Segueix l'exemple de la Teoria. Com que $g$ no és injectiva, fixa't que dos elements diferents poden tenir la mateixa imatge i que això no afecta la reindexació.
\end{indicacio}
\end{exercici}

\begin{exercici}
Comprova que $\mathsf{true}\colon1\to\Omega$ és sempre un monomorfisme, i que $\Hom(1,\Omega)\cong\Sub(1)$. Quants valors de veritat globals té $\cat{Set}$? \hfill{\normalfont$\star$}
\begin{indicacio}
Per a la primera part, compara dues fletxes qualssevol $X\to1$. Per a la segona, particularitza la representabilitat de $\Sub$.
\end{indicacio}
\end{exercici}

\section{Topos}

\begin{exercici}
Sigui $\mathbf 2=(0\to 1)$. La categoria de fletxes $\cat{Set}^{\mathbf 2}$ té per objectes les aplicacions $X_0\to X_1$, i és una categoria de prefeixos sobre $\mathbf 2\op$, per tant un topos. Fent servir la construcció amb sedassos:
\begin{enumerate}
  \item calcula els conjunts $\Omega(0)$ i $\Omega(1)$;
  \item calcula l'aplicació de restricció entre ells, de manera que $\Omega$ és un objecte de $\cat{Set}^{\mathbf 2}$, és a dir, una fletxa;
  \item determina els elements globals $1\to\Omega$, que són els \enquote{valors de veritat globals}.
\end{enumerate}
\hfill{\normalfont$\star\star$}
\begin{indicacio}
Compta els sedassos sobre cada objecte de $\mathbf 2\op$. Observa que $\Omega$ no és un conjunt de valors, sinó una fletxa entre dos conjunts.
\end{indicacio}
\end{exercici}

\begin{exercici}
Demostra que un topos que és alhora una categoria prima és equivalent a la categoria terminal $\cat{1}$. \hfill{\normalfont$\star\star$}
\begin{indicacio}
En una categoria prima, tot morfisme és mono. Aplica-ho al morfisme $A\to 1$ i fes servir que $\Sub(1)\cong\Hom(1,\Omega)$ té com a màxim un element.
\end{indicacio}
\end{exercici}

\begin{exercici}[Ampliació]
Fent servir el classificador de $\cat{Graph}$ calculat a la Teoria, determina els valors de veritat globals, és a dir, els morfismes $\mathbf L\to\Omega$ des del graf terminal d'un bucle. Comprova que n'hi ha tres, identifica cadascun amb un subgraf de $\mathbf L$ i dedueix que $\Omega$ no és isomorf a $1+1$. \hfill{\normalfont$\star\star$}
\begin{indicacio}
Un morfisme $\mathbf L\to\Omega$ tria un bucle de $\Omega$. Contrasta el resultat amb $\Sub(\mathbf L)\cong\Hom(\mathbf L,\Omega)$ i compta els morfismes $\mathbf L\to\mathbf L+\mathbf L$.
\end{indicacio}
\end{exercici}

\begin{exercici}
Dedueix de l'exercici resolt~\ref{pr-mono-equalitzador} de la Pràctica que $\cat{Top}$ no té classificador de subobjectes, i per tant no és un topos. \hfill{\normalfont$\star$}
\begin{indicacio}
A $\cat{Top}$, els monomorfismes són les aplicacions contínues injectives i els epimorfismes, les exhaustives. Busca una bijecció contínua entre dos espais de dos punts que no sigui un homeomorfisme.
\end{indicacio}
\end{exercici}

\begin{exercici}
Sigui $\cat{B}\mathbb N$ el monoide additiu $(\mathbb N,+,0)$ vist com a categoria d'un sol objecte $\ast$. Descriu tots els sedassos sobre $\ast$ i dedueix que $\Omega(\ast)$ està en bijecció amb $\mathbb N\cup\{\infty\}$. Descriu l'acció $\Omega(h)$ de cada $h\in\mathbb N$. Comprova que només hi ha dos valors de veritat globals, però que $\Sub(\yon(\ast))$ no és una àlgebra de Boole. \hfill{\normalfont$\star\star$}
\begin{indicacio}
Un sedàs no buit és tancat per sumar qualsevol natural; mira'n l'element mínim. Un valor global és un sedàs fix per totes les accions. Per a la darrera part, fes servir que $\Sub(\yon(\ast))\cong\Omega(\ast)$ i busca un sedàs $S$ tal que el més gran dels sedassos disjunts de $S$, unit amb $S$, no doni el sedàs màxim.
\end{indicacio}
\end{exercici}

\section{Grafs i lògica}

\begin{exercici}
Sigui $G$ el graf amb vèrtexs $a,b$ i arestes $\ell\colon a\to a$, $e\colon a\to b$ i $d\colon b\to a$, i sigui $H$ el subgraf format pel vèrtex $a$ i el bucle $\ell$. Fent servir la descripció de $\Omega$ de la Teoria, calcula $\chi_H\colon G\to\Omega$ i comprova que el pullback de $\mathsf{true}$ al llarg de $\chi_H$ és $H$. \hfill{\normalfont$\star$}
\begin{indicacio}
Per a cada aresta, mira quins dels seus dos extrems són a $H$ i si ella mateixa hi és, i tria l'aresta de $\Omega$ que registra aquesta informació.
\end{indicacio}
\end{exercici}

\begin{exercici}
A $\Sub(\mathbf E)$, on $\mathbf E$ és el graf d'una sola aresta $e\colon s\to t$, sigui $\neg K$ el més gran dels subgrafs disjunts de $K$. Comprova que existeix, calcula $\neg K$ per als cinc subgrafs de $\mathbf E$ i determina per a quins es compleix $\neg\neg K=K$. Què en dedueixes sobre la doble negació? \hfill{\normalfont$\star\star$}
\begin{indicacio}
Un subgraf que conté l'aresta ha de contenir els seus dos extrems. Fes una taula amb $K$, $\neg K$ i $\neg\neg K$.
\end{indicacio}
\end{exercici}

\fi

% ============================================================
% PART IV — TESTS INTERACTIUS
% ============================================================
\part{\NomTests}
\setcounter{chapter}{0}
% En pantalla, un sol entorn Form per a tots els tests: un únic diccionari
% AcroForm. En mode imprimible les macros dels tests són estàtiques i no es
% crea cap formulari.
\ifimprimible\else\begin{Form}\fi
\chapter{Categories}

Aquest test recull els conceptes essencials del capítol. Cada pregunta té una única resposta correcta.

\iniciTest{cat}{c,b,a,d,c,d,a,b,a,d,b,c}

\begin{enumerate}

% --- Definició i composició ---
\pregunta Un lector afirma: \enquote{Com que la composició d'una categoria és associativa, l'expressió $h\comp g\comp f$ està ben definida per a qualssevol tres morfismes $f,g,h$.} On falla?
\begin{enumerate}
  \opcio{a}{No falla: l'associativitat garanteix que sempre es pot compondre.}
  \opcio{b}{Falla perquè l'associativitat només val en categories petites.}
  \opcio{c}{Falla perquè la composició és \emph{parcial}.}
  \opcio{d}{Falla perquè $h\comp g\comp f$ mai no està definit sense parèntesis.}
\end{enumerate}
\feedback
\justificacio{$h\comp g\comp f$ només té sentit si $\cod(f)=\dom(g)$ i $\cod(g)=\dom(h)$; l'associativitat només diu que, quan les dues agrupacions existeixen, coincideixen.}

% --- Diagrames commutatius ---

\pregunta Considera el quadrat format per $f\colon A\to B$, $u\colon A\to C$, $v\colon B\to D$ i $g\colon C\to D$. Quina equació expressa que el quadrat commuta?
\begin{enumerate}
  \opcio{a}{$f\comp v=u\comp g$, perquè es llegeix en l'ordre en què es recorren les fletxes.}
  \opcio{b}{$v\comp f=g\comp u$.}
  \opcio{c}{$v\comp u=f\comp g$.}
  \opcio{d}{$f=g$ i $u=v$.}
\end{enumerate}
\feedback
\justificacio{Els dos camins de $A$ a $D$ són \enquote{primer $f$, després $v$} i \enquote{primer $u$, després $g$}, i la composició s'escriu de dreta a esquerra.}

% --- Exemples: lògica ---

\pregunta A la categoria $\cat{Prov}_T$ d'un sistema deductiu $T$, siguin $p$ i $q$ dues fórmules \emph{distintes} tals que $p\vdash_T q$ i $q\vdash_T p$. Què en podem dir?
\begin{enumerate}
  \opcio{a}{Són objectes isomorfs però no iguals: $\cat{Prov}_T$ és prima, però la deduïbilitat no és antisimètrica.}
  \opcio{b}{Han de ser iguals, perquè en una categoria prima dos objectes isomorfs coincideixen.}
  \opcio{c}{Hi ha exactament dues fletxes de $p$ a $q$, una per a cada deducció.}
  \opcio{d}{No hi ha cap fletxa entre elles, perquè són fórmules diferents.}
\end{enumerate}
\feedback
\justificacio{Hi ha fletxes en tots dos sentits, i com que la categoria és prima, les composicions són identitats: $p\cong q$. Però $p$ i $q$ són fórmules diferents. Una categoria prima és un preordre, no necessàriament un ordre parcial, i aquí es veu la diferència entre igualtat i isomorfisme.}

% --- Exemples: categories primes ---

\pregunta Quina de les categories següents \emph{no} és prima?
\begin{enumerate}
  \opcio{a}{La categoria d'un preordre $(P,\le)$.}
  \opcio{b}{La categoria $\cat{Prov}_T$ d'un sistema deductiu.}
  \opcio{c}{La categoria $\cat{2}=(0\to 1)$.}
  \opcio{d}{La categoria $\cat{B}M$ del monoide $(\mathbb N,+,0)$.}
\end{enumerate}
\feedback
\justificacio{Té un sol objecte, però $\Hom(\ast,\ast)=\mathbb N$ té més d'un element.}

% --- Isomorfismes ---

\pregunta Sigui $D$ el conjunt $\{0,1\}$ amb l'ordre discret i sigui $C$ el mateix conjunt amb l'ordre $0<1$. L'aplicació identitat subjacent $f\colon D\to C$ és bijectiva i monòtona. És un isomorfisme a $\cat{Pos}$?
\begin{enumerate}
  \opcio{a}{Sí, perquè qualsevol aplicació bijectiva és un isomorfisme.}
  \opcio{b}{Sí, perquè una aplicació monòtona bijectiva té sempre inversa monòtona.}
  \opcio{c}{No, perquè la inversa subjacent $C\to D$ no és monòtona.}
  \opcio{d}{No, perquè $f$ no és monòtona.}
\end{enumerate}
\feedback
\justificacio{Per tant no és un morfisme de $\cat{Pos}$, i $f$ no té invers en la categoria.}

% --- Mono/epi: cancel·lació ---

\pregunta L'aplicació $i\colon\varnothing\to\{\ast\}$ de $\cat{Set}$ és:
\begin{enumerate}
  \opcio{a}{un isomorfisme.}
  \opcio{b}{un epimorfisme amb secció.}
  \opcio{c}{ni mono ni epi.}
  \opcio{d}{un monomorfisme sense retracció.}
\end{enumerate}
\feedback
\justificacio{És injectiva i, per tant, un monomorfisme de $\cat{Set}$; però no té retracció, perquè no hi ha cap aplicació $\{\ast\}\to\varnothing$.}

\pregunta Un morfisme $f\colon A\to B$ té una \emph{retracció} $r$ (amb $r\comp f=\id_A$). Quina conclusió és correcta i per què?
\begin{enumerate}
  \opcio{a}{$f$ és un monomorfisme: si $f\comp g=f\comp h$, component amb $r$ per l'esquerra dona $g=h$.}
  \opcio{b}{$f$ és un isomorfisme, perquè té invers per un costat.}
  \opcio{c}{$f$ és un epimorfisme, perquè $r$ el cancel·la per la dreta.}
  \opcio{d}{No se'n pot dir res sense saber la categoria.}
\end{enumerate}
\feedback
\justificacio{Si $f\comp g=f\comp h$, aleshores $g=r\comp f\comp g=r\comp f\comp h=h$. Tenir invers per un costat no fa $f$ isomorfisme: a $\cat{Set}$, la inclusió $\{0\}\hookrightarrow\{0,1\}$ té retracció però no és bijectiva.}

% --- Dualitat i oposada ---

\pregunta Considera la categoria oposada $\cat{C}\op$. Un lector diu: \enquote{Si $f$ és un monomorfisme a $\cat{C}$, aleshores $f$ és un monomorfisme a $\cat{C}\op$.} És correcte?
\begin{enumerate}
  \opcio{a}{Sí, perquè $\cat{C}$ i $\cat{C}\op$ tenen els mateixos morfismes.}
  \opcio{b}{No, perquè un mono de $\cat{C}$ és un epi de $\cat{C}\op$.}
  \opcio{c}{No, perquè a $\cat{C}\op$ no hi ha monomorfismes.}
  \opcio{d}{Sí, perquè els monomorfismes són autoduals.}
\end{enumerate}
\feedback
\justificacio{En passar a l'oposada, la cancel·lació per l'esquerra es converteix en cancel·lació per la dreta, igual que una retracció passa a ser una secció.}

% --- Construccions de categories ---

\pregunta Sigui $A$ un objecte d'una categoria $\cat{C}$. Què és un morfisme de $x\colon X\to A$ a $y\colon Y\to A$ a la categoria $\cat{C}/A$ sobre $A$?
\begin{enumerate}
  \opcio{a}{Un morfisme $h\colon X\to Y$ de $\cat{C}$ tal que $y\comp h=x$.}
  \opcio{b}{Qualsevol morfisme $h\colon X\to Y$ de $\cat{C}$, sense cap condició.}
  \opcio{c}{Un morfisme $h\colon Y\to X$ tal que $x\comp h=y$.}
  \opcio{d}{Un morfisme $A\to A$.}
\end{enumerate}
\feedback
\justificacio{És a dir, que fa commutar el triangle sobre $A$.}

% --- Construccions de categories: categoria producte ---

\pregunta Siguin $\cat{C}$ i $\cat{D}$ dues categories. Què és un morfisme de $(A,X)$ a $(B,Y)$ a la categoria producte $\cat{C}\times\cat{D}$?
\begin{enumerate}
  \opcio{a}{Un morfisme $A\times X\to B\times Y$ d'una de les dues categories.}
  \opcio{b}{Un morfisme $A\to B$ de $\cat{C}$ o bé un morfisme $X\to Y$ de $\cat{D}$.}
  \opcio{c}{L'objecte producte $A\times B$, caracteritzat per la seva propietat universal dins de $\cat{C}$.}
  \opcio{d}{Una parella $(f,g)$ amb $f\colon A\to B$ a $\cat{C}$ i $g\colon X\to Y$ a $\cat{D}$.}
\end{enumerate}
\feedback
\justificacio{La composició es fa component a component. No s'ha de confondre amb el producte d'objectes dins d'una sola categoria.}

% --- Categoria lliure ---

\pregunta Considera el graf amb un sol vèrtex $\ast$ i una sola aresta $a\colon\ast\to\ast$. Quants morfismes té la categoria lliure que genera?
\begin{enumerate}
  \opcio{a}{Dos: la identitat i $a$.}
  \opcio{b}{Infinits, un per a cada longitud de camí.}
  \opcio{c}{Un de sol, perquè el graf té un sol vèrtex.}
  \opcio{d}{Cap, perquè $a$ no es pot compondre amb si mateixa.}
\end{enumerate}
\feedback
\justificacio{Són $\id_\ast,\ a,\ a\comp a,\ a\comp a\comp a,\dots$; és la categoria $\cat{B}M$ del monoide $(\mathbb N,+,0)$.}

% --- Mida ---

\pregunta Un lector afirma: \enquote{La categoria $\cat{Set}$ és petita, perquè els seus morfismes són simplement aplicacions.} On s'equivoca?
\begin{enumerate}
  \opcio{a}{No s'equivoca: $\cat{Set}$ és petita.}
  \opcio{b}{S'equivoca perquè $\cat{Set}$ no té identitats.}
  \opcio{c}{S'equivoca perquè $\cat{Set}$ no és petita.}
  \opcio{d}{S'equivoca perquè $\cat{Set}$ ni tan sols és una categoria.}
\end{enumerate}
\feedback
\justificacio{És \emph{localment petita} (cada $\Hom(A,B)$ és un conjunt), però la col·lecció de tots els conjunts és una classe pròpia, no un conjunt.}

\end{enumerate}
\clauestatica
\finalitzaTest

\chapter{Functors}

Aquest test recull els conceptes essencials del capítol. Cada pregunta té una única resposta correcta.

\iniciTest{fun}{a,b,d,c,a,c,d,a,b,c,b}

\begin{enumerate}

\pregunta Un lector afirma: \enquote{Per definir un functor $F$ n'hi ha prou de dir on va cada objecte; l'acció sobre els morfismes queda determinada.} On falla?
\begin{enumerate}
  \opcio{a}{Falla perquè l'acció sobre morfismes no queda determinada per la dels objectes.}
  \opcio{b}{Falla només quan $\cat{C}$ és gran.}
  \opcio{c}{No falla: els objectes determinen el functor.}
  \opcio{d}{Falla perquè els functors no actuen sobre objectes.}
\end{enumerate}
\feedback
\justificacio{Un functor és \emph{dues} assignacions, sobre objectes i sobre morfismes, que han de respectar dominis, codominis, composició i identitats.}

\pregunta Un lector proposa un \enquote{functor} $\cat{B}M\to\cat{B}M$ (amb $M=\{1,e\}$, $e\comp e=e$) que fixa l'objecte i envia el bucle $1$ a $e$ i el bucle $e$ a $e$. És un functor?
\begin{enumerate}
  \opcio{a}{Sí, perquè respecta dominis i codominis.}
  \opcio{b}{No, perquè no envia la identitat a la identitat.}
  \opcio{c}{Sí, perquè $e$ és idempotent.}
  \opcio{d}{No, perquè cap assignació entre monoides és functorial.}
\end{enumerate}
\feedback
\justificacio{La identitat $1$ de $\cat{B}M$ va a parar a $e\neq 1$; és un morfisme de grafs, però no un functor.}

\pregunta Es pot afirmar que \emph{tot} functor preserva els monomorfismes?
\begin{enumerate}
  \opcio{a}{Sí, perquè preserva tota l'estructura categòrica.}
  \opcio{b}{Sí, perquè preserva la composició.}
  \opcio{c}{No, perquè cap functor preserva mai els monomorfismes.}
  \opcio{d}{No en general, tot i que alguns functors sí que els preserven.}
\end{enumerate}
\feedback
\justificacio{Un functor preserva el que està testimoniat per morfismes i equacions (isomorfismes, seccions, retraccions), però ser mono es defineix per cancel·lació sobre \emph{totes} les fletxes de la categoria d'arribada, incloses les que no provenen de la categoria d'origen.}

\pregunta Sigui $\cat{S}$ la subcategoria plena de $\cat{FinSet}$ que té per objectes els conjunts $\underline{0}=\varnothing$ i $\underline{n}=\{1,\dots,n\}$ per a $n\geq 1$, un de cada cardinal finit, i sigui $\iota\colon\cat{S}\hookrightarrow\cat{FinSet}$ la inclusió. Quina afirmació és correcta?
\begin{enumerate}
  \opcio{a}{$\iota$ és exhaustiu sobre objectes en sentit estricte.}
  \opcio{b}{$\iota$ no és ni ple ni fidel.}
  \opcio{c}{$\iota$ és essencialment exhaustiu però no exhaustiu sobre objectes.}
  \opcio{d}{$\iota$ és un isomorfisme de categories.}
\end{enumerate}
\feedback
\justificacio{Tot conjunt finit és isomorf a algun $\underline{n}$, però en general no és igual a cap d'ells.}

\pregunta Sigui $F$ un functor \emph{ple i fidel} i suposem que $F(f)$ és un isomorfisme a $\cat{D}$. Quina conclusió és correcta, i sobre quin fet reposa?
\begin{enumerate}
  \opcio{a}{Que $f$ ja era un isomorfisme a $\cat{C}$.}
  \opcio{b}{Que no es pot concloure que $f$ sigui un isomorfisme sense suposar, a més, que $F$ és exhaustiu sobre objectes.}
  \opcio{c}{Res: ser ple i fidel no diu res sobre isomorfismes.}
  \opcio{d}{Que $F$ és una equivalència.}
\end{enumerate}
\feedback
\justificacio{Com que $F$ és ple, l'invers de $F(f)$ prové d'algun $g$; com que $F$ és fidel, les identitats $g\comp f=\id$ i $f\comp g=\id$ es dedueixen de les seves imatges. Els functors plens i fidels \emph{reflecteixen} isomorfismes.}

\pregunta Siguin $f\colon X\to Y$ i $h\colon Y\to B$. Quina és l'acció de l'homfunctor contravariant $\Hom(-,B)$ sobre $f$, aplicada a $h$?
\begin{enumerate}
  \opcio{a}{$h\mapsto f\comp h$, amb valors a $\Hom(X,B)$.}
  \opcio{b}{$h\mapsto f\comp h$, amb valors a $\Hom(Y,B)$.}
  \opcio{c}{$h\mapsto h\comp f$, amb valors a $\Hom(X,B)$.}
  \opcio{d}{$h\mapsto h\comp f$, amb valors a $\Hom(Y,B)$.}
\end{enumerate}
\feedback
\justificacio{La precomposició amb $f$ va de $\Hom(Y,B)$ a $\Hom(X,B)$: inverteix el sentit de $f$.}

\pregunta Un lector diu: \enquote{El functor oblit $U\colon\cat{Grp}\to\cat{Set}$ és ple, perquè envia cada homomorfisme a la seva aplicació subjacent.} On s'equivoca?
\begin{enumerate}
  \opcio{a}{No s'equivoca: $U$ és ple.}
  \opcio{b}{S'equivoca perquè $U$ no és ni tan sols un functor.}
  \opcio{c}{S'equivoca perquè $U$ no actua sobre morfismes.}
  \opcio{d}{S'equivoca: $U$ és fidel, però no ple.}
\end{enumerate}
\feedback
\justificacio{Homomorfismes diferents donen aplicacions diferents; però entre dos grups hi ha aplicacions de conjunts que no són homomorfismes i que, per tant, no provenen de cap morfisme de $\cat{Grp}$.}

\pregunta Un lector afirma: \enquote{Com que $\cat{Cat}$ té totes les categories petites per objectes, $\cat{Cat}$ és ella mateixa una categoria petita.} És correcte?
\begin{enumerate}
  \opcio{a}{No, perquè $\cat{Cat}$ és gran, encara que localment petita.}
  \opcio{b}{Sí, perquè els functors formen conjunts.}
  \opcio{c}{Sí: una categoria de categories petites és petita.}
  \opcio{d}{No, perquè $\cat{Cat}$ no té identitats.}
\end{enumerate}
\feedback
\justificacio{La col·lecció de totes les categories petites és una classe pròpia, com la de tots els conjunts.}

\pregunta Donats un graf $G$ i una categoria petita $\cat{D}$, definir un functor $\Free(G)\to\cat{D}$ des de la categoria lliure de $G$ equival a:
\begin{enumerate}
  \opcio{a}{triar només les imatges de les arestes de $G$.}
  \opcio{b}{donar un morfisme de grafs $G\to U(\cat{D})$.}
  \opcio{c}{triar un objecte de $\cat{D}$ per a cada camí de $G$ de manera independent.}
  \opcio{d}{donar una equivalència entre $\Free(G)$ i $\cat{D}$.}
\end{enumerate}
\feedback
\justificacio{Cal triar les imatges dels vèrtexs \emph{i} de les arestes, respectant origen i destí; aquesta tria s'estén de manera única als camins.}

\pregunta Sobre el bifunctor $\Hom(-,-)\colon\cat{C}\op\times\cat{C}\to\cat{Set}$, amb $\Hom(f,g)(h)=g\comp h\comp f$, quina descripció de la variància és correcta?
\begin{enumerate}
  \opcio{a}{Covariant en totes dues variables.}
  \opcio{b}{Contravariant en totes dues variables.}
  \opcio{c}{Contravariant en la primera variable i covariant en la segona.}
  \opcio{d}{Covariant en la primera i contravariant en la segona.}
\end{enumerate}
\feedback
\justificacio{Precomposició amb $f$ en la primera variable i postcomposició amb $g$ en la segona.}

% --- Fil lògic ---

\pregunta Siguin $T\subseteq T'$ dos sistemes deductius sobre les mateixes fórmules, i sigui
\[
J\colon\cat{Prov}_T\to\cat{Prov}_{T'}
\]
el functor identitat sobre fórmules. Quan és $J$ ple?
\begin{enumerate}
  \opcio{a}{Sempre, perquè és la identitat sobre objectes.}
  \opcio{b}{Exactament quan $T'$ no introdueix cap conseqüència nova $A\vdash_{T'}B$ entre les fórmules considerades.}
  \opcio{c}{Només quan $T=T'$ com a conjunts d'axiomes.}
  \opcio{d}{Mai, perquè les categories de deduïbilitat són primes.}
\end{enumerate}
\feedback
\justificacio{Totes dues categories són primes, de manera que $J$ és ple si i només si cada fletxa $A\to B$ de $\cat{Prov}_{T'}$ prové d'una de $\cat{Prov}_T$, és a dir, si $A\vdash_{T'}B$ implica $A\vdash_T B$. Ser la identitat sobre objectes (a) no diu res sobre els morfismes.}

\end{enumerate}
\clauestatica
\finalitzaTest

\chapter[Transformacions naturals]{\texorpdfstring{Transformacions naturals\\ i equivalències}{Transformacions naturals i equivalències}}

Aquest test recull els conceptes essencials del capítol. Cada pregunta té una única resposta correcta.

\iniciTest{trn}{c,b,c,d,a,b,c,b,a,d,a,d}

\begin{enumerate}

\pregunta A $\cat{Set}$, prenem $F=G=\id_{\cat{Set}}$ i definim $\eta_A=\id_A$ per a tot conjunt $A\neq\{0,1\}$, i $\eta_{\{0,1\}}=\tau$, la transposició $\tau(0)=1$, $\tau(1)=0$. Aquesta família, \emph{és} una transformació natural $\id\Rightarrow\id$?
\begin{enumerate}
  \opcio{a}{Sí, perquè cada component $\eta_A$ és una bijecció.}
  \opcio{b}{Sí, perquè $\eta$ coincideix amb la identitat a tots els conjunts excepte un.}
  \opcio{c}{No, perquè una aplicació constant trenca el quadrat de naturalitat.}
  \opcio{d}{No, perquè la transposició $\tau$ no és un morfisme de $\cat{Set}$.}
\end{enumerate}
\feedback
\justificacio{Per a $f\colon\{0,1\}\to\{0,1\}$ constant de valor $0$, es té $f\comp\tau=f$, mentre que $\tau\comp f$ és constant de valor $1$.}

\pregunta Un lector proposa aquest argument: \enquote{Si $F(A)\cong G(A)$ per a cada objecte $A$, aleshores $F$ i $G$ són naturalment isomorfs, perquè n'hi ha prou d'agafar un isomorfisme $\eta_A\colon F(A)\to G(A)$ per a cada $A$.} On falla?
\begin{enumerate}
  \opcio{a}{No falla: l'argument és correcte.}
  \opcio{b}{Falla perquè els $\eta_A$ triats poden no complir el quadrat de naturalitat.}
  \opcio{c}{Falla perquè un isomorfisme natural exigeix que $F=G$.}
  \opcio{d}{Falla perquè $F(A)$ i $G(A)$ no poden ser mai isomorfs si $F\neq G$.}
\end{enumerate}
\feedback
\justificacio{L'existència d'isomorfismes objecte a objecte no garanteix que siguin \emph{compatibles} amb l'acció dels functors sobre els morfismes.}

\pregunta Sigui $\mathbb{Z}/2\mathbb{Z}=\{1,g\}$ vist com a categoria d'un objecte, i $F,G\colon\cat{B}(\mathbb{Z}/2\mathbb{Z})\to\cat{Set}$ que envien l'objecte a $\{0,1\}$, on $g$ actua per la identitat en $F$ i per la transposició $\tau$ en $G$. Quantes transformacions naturals $\eta\colon F\Rightarrow G$ hi ha?
\begin{enumerate}
  \opcio{a}{Exactament una, l'única component possible $\eta_\ast=\id$.}
  \opcio{b}{Dues: la identitat i la transposició.}
  \opcio{c}{Cap.}
  \opcio{d}{Infinites, una per cada aplicació $\{0,1\}\to\{0,1\}$.}
\end{enumerate}
\feedback
\justificacio{La naturalitat respecte de $g$ exigiria $\tau\comp\eta_\ast=\eta_\ast$, impossible perquè $\tau$ no té punts fixos.}

\pregunta Què és un morfisme de $f\colon A\to B$ a $g\colon A'\to B'$ a la categoria de functors $\cat{D}^{\cat{2}}$, on $\cat{2}=(0\to 1)$?
\begin{enumerate}
  \opcio{a}{Un morfisme qualsevol $A\to A'$.}
  \opcio{b}{Un functor $\cat{2}\to\cat{D}$.}
  \opcio{c}{Una parella qualsevol de morfismes $u\colon A\to A'$ i $v\colon B\to B'$.}
  \opcio{d}{Una parella $u\colon A\to A'$, $v\colon B\to B'$ amb $v\comp f=g\comp u$.}
\end{enumerate}
\feedback
\justificacio{És la condició de naturalitat per a l'única fletxa no trivial de $\cat{2}$.}

\pregunta En l'exemple sintaxi/semàntica, amb el functor sintàctic $F$, el semàntic $G$ i la interpretació $\eta$, què expressa la igualtat $G(f)\comp\eta_X=\eta_Y\comp F(f)$ per a un canvi de variables $f\colon X\to Y$?
\begin{enumerate}
  \opcio{a}{Que substituir variables i després interpretar dona la mateixa funció de veritat que interpretar i després transportar la semàntica al llarg del canvi de variables.}
  \opcio{b}{Que $\eta_X$ és una bijecció entre fórmules i funcions de veritat.}
  \opcio{c}{Que dues fórmules amb la mateixa taula de veritat són iguals.}
  \opcio{d}{Que el canvi de variables $f$ ha de ser bijectiu.}
\end{enumerate}
\feedback
\justificacio{És el quadrat de naturalitat de $\eta$ per a $f$: els dos camins de $F(X)$ a $G(Y)$ coincideixen. En termes lògics, la interpretació és compatible amb la substitució. Ni (b) ni (c) són certes, perquè fórmules diferents poden tenir la mateixa taula de veritat.}

\pregunta Un estudiant escriu que la composició vertical de $\eta\colon F\Rightarrow G$ i $\theta\colon G\Rightarrow H$ té components $(\theta\comp\eta)_A=\eta_A\comp\theta_A$. Per què és incorrecte, i quina és la forma bona?
\begin{enumerate}
  \opcio{a}{És correcte tal com està.}
  \opcio{b}{L'ordre està invertit: la composta és $\theta_A\comp\eta_A$.}
  \opcio{c}{Cap dels dos ordres té sentit; la composició vertical no existeix.}
  \opcio{d}{Tots dos ordres donen el mateix resultat, així que és igual.}
\end{enumerate}
\feedback
\justificacio{$\eta_A\colon F(A)\to G(A)$ i $\theta_A\colon G(A)\to H(A)$ només es componen en aquest ordre, i la composta va de $F(A)$ a $H(A)$.}

\pregunta Sigui $L\colon\cat{Set}\to\cat{Set}$ el functor de llistes finites, amb $L(f)$ aplicant $f$ a cada entrada, i sigui $\eta_X(x)=[x]$ la llista d'un sol element. Per a $f\colon X\to Y$, quina és l'equació de naturalitat ben tipada?
\begin{enumerate}
  \opcio{a}{$f\comp\eta_X=\eta_Y\comp L(f)$.}
  \opcio{b}{$\eta_X\comp L(f)=f\comp\eta_Y$.}
  \opcio{c}{$L(f)\comp\eta_X=\eta_Y\comp f$.}
  \opcio{d}{$\eta_Y=L(f)$.}
\end{enumerate}
\feedback
\justificacio{Tots dos costats envien $x$ a $[f(x)]$.}

\pregunta Un functor $F\colon\cat{C}\to\cat{D}$ és ple, fidel i essencialment exhaustiu. Un lector conclou: \enquote{Aleshores $F$ és bijectiu sobre els objectes.} És correcte?
\begin{enumerate}
  \opcio{a}{Sí: essencialment exhaustiu vol dir exhaustiu sobre objectes, i ple i fidel vol dir injectiu sobre objectes.}
  \opcio{b}{No: $F$ pot no ser bijectiu sobre objectes.}
  \opcio{c}{No, perquè cap functor és mai bijectiu sobre objectes.}
  \opcio{d}{Sí, sempre que $\cat{C}$ i $\cat{D}$ siguin petites.}
\end{enumerate}
\feedback
\justificacio{Essencialment exhaustiu només demana que tot objecte de $\cat{D}$ sigui \emph{isomorf} a una imatge, i ser ple i fidel afecta els morfismes; per exemple, la inclusió d'un esquelet no és bijectiva sobre objectes.}

\pregunta A la construcció del quasiinvers $G$ d'una equivalència, per a cada objecte $D$ es tria un objecte $G(D)$ i un isomorfisme $\varepsilon_D\colon F(G(D))\to D$, i sobre morfismes es defineix $G(v)$ com l'\emph{única} fletxa amb $F(G(v))=\varepsilon_{D'}^{-1}\comp v\comp\varepsilon_D$. Quina propietat de $F$ garanteix aquesta unicitat i existència?
\begin{enumerate}
  \opcio{a}{Que $F$ sigui ple i fidel.}
  \opcio{b}{Només ser fidel.}
  \opcio{c}{L'exhaustivitat essencial.}
  \opcio{d}{Que $\cat{D}$ sigui petita.}
\end{enumerate}
\feedback
\justificacio{Ser ple dona l'existència de $G(v)$, i ser fidel, la unicitat.}

\pregunta Siguin $\eta\colon F\Rightarrow G$, amb $F,G\colon\cat{C}\to\cat{D}$, i $\theta\colon K\Rightarrow L$, amb $K,L\colon\cat{D}\to\cat{E}$. Quina és la component $(\theta\ast\eta)_A$ de la composició horitzontal?
\begin{enumerate}
  \opcio{a}{$\theta_A\comp\eta_A$.}
  \opcio{b}{$K(\eta_A)\comp\theta_{G(A)}$.}
  \opcio{c}{$\eta_{K(A)}\comp\theta_A$.}
  \opcio{d}{$\theta_{G(A)}\comp K(\eta_A)$.}
\end{enumerate}
\feedback
\justificacio{Coincideix amb $L(\eta_A)\comp\theta_{F(A)}$ per la naturalitat de $\theta$ aplicada a $\eta_A$.}

\pregunta Siguin $\cat{1}$ la categoria terminal i $\cat{I}$ la categoria amb dos objectes i exactament un morfisme entre cada parella ordenada d'objectes. Quina afirmació és correcta?
\begin{enumerate}
  \opcio{a}{Són equivalents però no isomorfes.}
  \opcio{b}{Són isomorfes, perquè tots els objectes de $\cat{I}$ són isomorfs entre si.}
  \opcio{c}{No són equivalents, perquè tenen un nombre diferent d'objectes.}
  \opcio{d}{Són isomorfes, perquè totes dues són categories primes.}
\end{enumerate}
\feedback
\justificacio{La inclusió $\cat{1}\to\cat{I}$ és plena, fidel i essencialment exhaustiva, però cap functor entre elles no pot ser bijectiu sobre objectes.}

\pregunta Sigui $F\colon\cat{Mat}_{\mathbb K}\to\cat{FinVect}_{\mathbb K}$, $n\mapsto\mathbb K^n$, el functor de l'exemple de la Teoria. Quina afirmació és correcta?
\begin{enumerate}
  \opcio{a}{Les dues categories són iguals, perquè tot espai de dimensió $n$ \emph{és} $\mathbb K^n$.}
  \opcio{b}{Són isomorfes però no iguals: $F$ té un invers estricte.}
  \opcio{c}{Són equivalents i, per tant, isomorfes.}
  \opcio{d}{Són equivalents però no isomorfes; amb una tria adequada de bases, $G\comp F=\id$ literalment, però $F\comp G$ només és isomorf a la identitat.}
\end{enumerate}
\feedback
\justificacio{$F$ és ple, fidel i essencialment exhaustiu, i per tant una equivalència. No hi ha cap isomorfisme de categories, perquè $\cat{Mat}_{\mathbb K}$ és esquelètica i $\cat{FinVect}_{\mathbb K}$ no ho és. El quasiinvers tria una base de cada espai: $F(G(V))=\mathbb K^{\dim V}$, que és isomorf a $V$ però en general no hi és igual. Un espai de dimensió $n$ és isomorf a $\mathbb K^n$, no igual.}

\end{enumerate}
\clauestatica
\finalitzaTest

\ifpublicacioparcial\else
\chapter[Representabilitat i Yoneda]{\texorpdfstring{Propietats universals,\\ representabilitat i Yoneda}{Propietats universals, representabilitat i Yoneda}}

Aquest test recull els conceptes essencials del capítol. Cada pregunta té una única resposta correcta.

\iniciTest{yon}{a,c,b,d,a,d,b,c,a,c,b,d,c,d,b}

\begin{enumerate}

\pregunta Quan és \textbf{terminal} un objecte $1$ d'una categoria?
\begin{enumerate}
  \opcio{a}{Quan hi ha exactament un morfisme $X\to 1$ per a cada $X$.}
  \opcio{b}{Quan hi ha exactament un morfisme $1\to X$ per a cada $X$.}
  \opcio{c}{Quan no hi ha cap morfisme cap a ell.}
  \opcio{d}{Quan és isomorf a tots els altres objectes.}
\end{enumerate}
\feedback
\justificacio{És la definició. L'opció (b) és la d'objecte inicial, i (d) no té per què complir-se: ser terminal és una propietat universal, no un isomorfisme amb tots els objectes.}

\pregunta Siguin $S\colon\cat{A}\to\cat{C}$ i $T\colon\cat{B}\to\cat{C}$. Una parella $a\colon A\to A'$, $b\colon B\to B'$ és un morfisme $(A,B,u)\to(A',B',u')$ de la categoria coma $(S\downarrow T)$ quan:
\begin{enumerate}
  \opcio{a}{$S(a)\comp u=u'\comp T(b)$.}
  \opcio{b}{$T(b)\comp u'=u\comp S(a)$.}
  \opcio{c}{$T(b)\comp u=u'\comp S(a)$.}
  \opcio{d}{$S(a)=T(b)$.}
\end{enumerate}
\feedback
\justificacio{Comprova els tipus: $u\colon S(A)\to T(B)$ i $u'\colon S(A')\to T(B')$. L'únic quadrat ben tipat de $S(A)$ a $T(B')$ és $T(b)\comp u=u'\comp S(a)$. A (a) i (b) les composicions no existeixen en general.}

\pregunta En la formulació per representabilitat desenvolupada en aquest capítol, una \textbf{propietat universal} caracteritza una dada...
\begin{enumerate}
  \opcio{a}{per la seva constitució interna.}
  \opcio{b}{per la manera com es relaciona amb tots els altres objectes, mitjançant la representabilitat d'un functor cap a $\cat{Set}$.}
  \opcio{c}{per la seva cardinalitat.}
  \opcio{d}{només a $\cat{Set}$.}
\end{enumerate}
\feedback
\justificacio{Una propietat universal no descriu què \emph{és} la dada, sinó com es relaciona amb tota la resta: la dada és un objecte representant, amb el seu element universal.}

\pregunta Que $1$ sigui terminal equival a dir que:
\begin{enumerate}
  \opcio{a}{$\Hom(1,-)$ és constant.}
  \opcio{b}{$\Hom(-,1)$ no és functor.}
  \opcio{c}{$\Hom(-,1)$ és buit.}
  \opcio{d}{$\Hom(-,1)$ és isomorf al functor constant amb valor un punt.}
\end{enumerate}
\feedback
\justificacio{Ser terminal vol dir que cada $\Hom(X,1)$ té exactament un element. Aquestes bijeccions amb $\{\ast\}$ formen un isomorfisme natural, perquè tots els quadrats acaben en un conjunt d'un element.}

\pregunta Quan es diu que un prefeix $P\colon\cat{C}\op\to\cat{Set}$ és \textbf{representable}?
\begin{enumerate}
  \opcio{a}{Quan és isomorf a $\Hom(-,A)$ per a algun objecte $A$.}
  \opcio{b}{Quan és constant.}
  \opcio{c}{Quan és ple i fidel.}
  \opcio{d}{Quan té un sol element a cada conjunt.}
\end{enumerate}
\feedback
\justificacio{És la definició, en la forma de la proposició sobre representabilitat i isomorfisme natural. Ser constant o tenir un element a cada conjunt són casos molt particulars; \enquote{ple i fidel} és una propietat de functors entre categories, no de prefeixos.}

\pregunta Sigui $P\colon\cat{C}\op\to\cat{Set}$ un prefeix i $x\in P(A)$. Quina és la component en $X$ de la transformació natural $\Hom(-,A)\Rightarrow P$ que el lema de Yoneda associa a $x$, avaluada en $f\colon X\to A$?
\begin{enumerate}
  \opcio{a}{$P(f^{-1})(x)$.}
  \opcio{b}{$x$, sense canviar de conjunt.}
  \opcio{c}{$P(\id_A)(f)$.}
  \opcio{d}{$P(f)(x)$.}
\end{enumerate}
\feedback
\justificacio{Per contravariància, $P(f)\colon P(A)\to P(X)$; no cal que $f$ tingui invers.}

\pregunta Sigui $f\colon A\to B$ un morfisme. Les quatre descripcions següents de la component en $X$ de $\yon(f)=\Hom(-,f)$ semblen raonables. Quina és l'única en què el domini, el codomini i la composició tenen els tipus correctes?
\begin{enumerate}
  \opcio{a}{$\Hom(X,A)\to\Hom(X,B)$, $\varphi\mapsto\varphi\comp f$.}
  \opcio{b}{$\Hom(X,A)\to\Hom(X,B)$, $\varphi\mapsto f\comp\varphi$.}
  \opcio{c}{$\Hom(A,X)\to\Hom(B,X)$, $\varphi\mapsto\varphi\comp f$.}
  \opcio{d}{$\Hom(X,B)\to\Hom(X,A)$, $\varphi\mapsto\varphi\comp f$.}
\end{enumerate}
\feedback
\justificacio{Si $\varphi\colon X\to A$, l'única composició possible amb $f\colon A\to B$ és $f\comp\varphi\colon X\to B$. A (a), $\varphi\comp f$ demanaria $B=X$. A (c), amb $\varphi\colon A\to X$, $\varphi\comp f$ demanaria $B=A$: és una versió mal tipada de $\Hom(f,X)$, que va de $\Hom(B,X)$ a $\Hom(A,X)$. A (d), amb $\varphi\colon X\to B$, $\varphi\comp f$ demanaria de nou $B=X$.}

\pregunta Sigui $\mathcal P$ el prefeix de parts, $A=\{0,1\}$ i $B=\{1\}$. Quin subconjunt assigna la transformació natural $\Hom(-,A)\Rightarrow\mathcal P$ determinada per $B$ a l'aplicació $f\colon\{a,b,c\}\to A$ amb $f(a)=f(c)=1$ i $f(b)=0$?
\begin{enumerate}
  \opcio{a}{$\{b\}$.}
  \opcio{b}{$\{a,b,c\}$.}
  \opcio{c}{$\{a,c\}$.}
  \opcio{d}{$\{1\}$.}
\end{enumerate}
\feedback
\justificacio{Pel lema de Yoneda, la transformació determinada per $B\in\mathcal P(A)$ envia $f$ a $\mathcal P(f)(B)=f^{-1}(B)=\{a,c\}$. L'opció (d) és $B$ mateix, que viu a $\mathcal P(A)$ i no a $\mathcal P(\{a,b,c\})$.}

\pregunta Sigui $\Phi_{A,P}\colon\Nat(\yon(A),P)\to P(A)$, $\eta\mapsto\eta_A(\id_A)$, la bijecció del lema de Yoneda, i sigui $a\colon A'\to A$. Quina de les igualtats següents expressa la naturalitat en $A$ amb tots els tipus correctes?
\begin{enumerate}
  \opcio{a}{$\Phi_{A',P}\bigl(\eta\comp\yon(a)\bigr)=P(a)\bigl(\Phi_{A,P}(\eta)\bigr)$, com a elements de $P(A')$.}
  \opcio{b}{$\Phi_{A,P}\bigl(\eta\comp\yon(a)\bigr)=P(a)\bigl(\Phi_{A',P}(\eta)\bigr)$, com a elements de $P(A)$.}
  \opcio{c}{$\Phi_{A',P}\bigl(\yon(a)\comp\eta\bigr)=P(a)\bigl(\Phi_{A,P}(\eta)\bigr)$, com a elements de $P(A')$.}
  \opcio{d}{$\Phi_{A',P}\bigl(\eta\comp\yon(a)\bigr)=P(a)^{-1}\bigl(\Phi_{A,P}(\eta)\bigr)$, com a elements de $P(A')$.}
\end{enumerate}
\feedback
\justificacio{Tenim $\eta\colon\yon(A)\Rightarrow P$ i $\yon(a)\colon\yon(A')\Rightarrow\yon(A)$. L'única composició possible és, doncs, $\eta\comp\yon(a)$, que va de $\yon(A')$ a $P$, i s'hi aplica $\Phi_{A',P}$. Per contravariància, $P(a)\colon P(A)\to P(A')$. A (b), $\Phi_{A,P}$ no s'aplica a una transformació que surt de $\yon(A')$. A (c), $\yon(a)\comp\eta$ no és componible. A (d), $P(a)$ no té per què ser invertible.}

\pregunta A la demostració del lema, a quina dada s'envia una transformació natural $\eta\colon\Hom(-,A)\Rightarrow P$?
\begin{enumerate}
  \opcio{a}{A $\eta_A$.}
  \opcio{b}{A tot el functor $P$.}
  \opcio{c}{A l'element $\eta_A(\id_A)\in P(A)$.}
  \opcio{d}{A la identitat de $P$.}
\end{enumerate}
\feedback
\justificacio{L'aplicació del lema és $\eta\mapsto\eta_A(\id_A)$. L'opció (a) és tota la component, no un element de $P(A)$; és l'element el que determina $\eta$.}

\pregunta Una conseqüència directa del lema de Yoneda és que la immersió $\yon$ és:
\begin{enumerate}
  \opcio{a}{constant.}
  \opcio{b}{plena i fidel.}
  \opcio{c}{essencialment exhaustiva.}
  \opcio{d}{un isomorfisme de categories.}
\end{enumerate}
\feedback
\justificacio{Amb $P=\Hom(-,B)$, el lema dona $\Nat(\Hom(-,A),\Hom(-,B))\cong\Hom(A,B)$, i la bijecció és $f\mapsto\yon(f)$. No és essencialment exhaustiva: hi ha prefeixos que no són representables.}

\pregunta Segons el corol·lari de Yoneda, $A\cong B$ si i només si:
\begin{enumerate}
  \opcio{a}{$A$ i $B$ tenen la mateixa cardinalitat.}
  \opcio{b}{$A=B$.}
  \opcio{c}{hi ha un functor entre ells.}
  \opcio{d}{$\Hom(-,A)\cong\Hom(-,B)$ com a prefeixos.}
\end{enumerate}
\feedback
\justificacio{Un functor ple i fidel reflecteix isomorfismes, i tot functor els preserva; aplicat a $\yon$, dona l'equivalència. La igualtat (b) és massa forta, i la cardinalitat (a) no té sentit en una categoria general.}

\pregunta Siguin $(P,p_1,p_2)$ i $(Q,q_1,q_2)$ dos productes dels mateixos objectes. Què és únic?
\begin{enumerate}
  \opcio{a}{Res: poden ser literalment iguals o no.}
  \opcio{b}{L'isomorfisme entre els objectes $P$ i $Q$, sense cap condició.}
  \opcio{c}{L'isomorfisme $u\colon P\to Q$ que compleix $q_i\comp u=p_i$ per a $i=1,2$.}
  \opcio{d}{L'automorfisme de $P$, que només pot ser la identitat.}
\end{enumerate}
\feedback
\justificacio{La propietat universal fa que l'isomorfisme sigui únic \emph{entre els compatibles amb les projeccions}. Sense aquesta condició, pot haver-n'hi molts: a $\cat{Set}$, per exemple, qualsevol bijecció de $P$ en $Q$.}

\pregunta Siguin $S\colon\cat{A}\to\cat{C}$ i $T\colon\cat{B}\to\cat{C}$ dos functors. Un objecte de la categoria coma $(S\downarrow T)$ és:
\begin{enumerate}
  \opcio{a}{només un objecte de $\cat{C}$.}
  \opcio{b}{una transformació natural $S\Rightarrow T$.}
  \opcio{c}{una parella d'objectes isomorfs.}
  \opcio{d}{una terna $(A,B,f)$ amb $A$ de $\cat{A}$, $B$ de $\cat{B}$ i $f\colon S(A)\to T(B)$.}
\end{enumerate}
\feedback
\justificacio{És la definició: un objecte de $\cat{A}$, un de $\cat{B}$ i una fletxa de $S(A)$ a $T(B)$ a $\cat{C}$. Les categories sobre i sota un objecte en són casos particulars.}

\pregunta Si $(A,u)$ representa un prefeix $P\colon\cat{C}\op\to\cat{Set}$, què expressa la universalitat de l'element $u\in P(A)$?
\begin{enumerate}
  \opcio{a}{Que $u$ és l'únic element de $P(A)$.}
  \opcio{b}{Que per a cada objecte $X$ i cada $x\in P(X)$ hi ha un únic $f\colon X\to A$ amb $P(f)(u)=x$.}
  \opcio{c}{Que $A$ és un objecte terminal.}
  \opcio{d}{Que $P$ és un functor constant.}
\end{enumerate}
\feedback
\justificacio{És la definició de representació d'un prefeix: cada element de cada $P(X)$ s'obté de $u$ per un únic morfisme $X\to A$. Equival a l'isomorfisme natural $\Hom(-,A)\cong P$ que envia $f$ a $P(f)(u)$.}

\end{enumerate}
\clauestatica
\finalitzaTest

\chapter[Límits i colímits]{Límits i colímits}

Aquest test recull els conceptes essencials del capítol. Cada pregunta té una única resposta correcta.

\iniciTest{lim}{b,c,a,d,d,b,c,a,d,c,a,c,b,b,a}

\begin{enumerate}

\pregunta Què és un \textbf{diagrama} de forma $\cat{I}$ en una categoria $\cat{C}$?
\begin{enumerate}
  \opcio{a}{Un objecte de $\cat{C}$.}
  \opcio{b}{Un functor $D\colon\cat{I}\to\cat{C}$, on $\cat{I}$ és una categoria índex petita.}
  \opcio{c}{Una transformació natural.}
  \opcio{d}{Un morfisme de $\cat{I}$.}
\end{enumerate}
\feedback
\justificacio{És la definició: un diagrama és un functor, i la categoria índex en dona la forma. Un objecte o un morfisme sols no recullen les compatibilitats entre les fletxes.}

\pregunta Un \textbf{con} sobre $D$ amb vèrtex $A$ és el mateix que:
\begin{enumerate}
  \opcio{a}{un objecte terminal.}
  \opcio{b}{un límit.}
  \opcio{c}{una transformació natural $\Delta_A\Rightarrow D$ des del functor constant.}
  \opcio{d}{un functor $\cat{C}\to\cat{I}$.}
\end{enumerate}
\feedback
\justificacio{Les components de $\alpha\colon\Delta_A\Rightarrow D$ són fletxes $A\to D_i$, i la naturalitat per a $u\colon i\to j$ diu $D_u\comp\alpha_i=\alpha_j$, que és exactament la condició de con.}

\pregunta Un \textbf{límit} de $D$ és:
\begin{enumerate}
  \opcio{a}{un objecte terminal de la categoria de cons $\mathrm{Con}(D)$.}
  \opcio{b}{qualsevol con sobre $D$.}
  \opcio{c}{un objecte inicial de $\cat{C}$.}
  \opcio{d}{el producte de tots els objectes.}
\end{enumerate}
\feedback
\justificacio{Un con límit és un con pel qual factoritza de manera única qualsevol altre con: és a dir, un objecte terminal de $\mathrm{Con}(D)$. D'aquí ve la unicitat llevat d'un únic isomorfisme compatible amb les projeccions.}

\pregunta Quina categoria índex dona lloc al \textbf{producte} de dos objectes?
\begin{enumerate}
  \opcio{a}{La categoria buida.}
  \opcio{b}{Dues fletxes paral·leles.}
  \opcio{c}{La forma de racó.}
  \opcio{d}{La categoria discreta amb dos objectes.}
\end{enumerate}
\feedback
\justificacio{Un con sobre un parell d'objectes, sense cap fletxa entre ells, és una parella de fletxes sense cap condició: la propietat del producte. Dues fletxes paral·leles donen l'equalitzador, i el racó, el pullback.}

\pregunta A $\cat{Set}$, l'\textbf{equalitzador} de $f,g\colon A\to B$ és:
\begin{enumerate}
  \opcio{a}{el producte cartesià $A\times B$.}
  \opcio{b}{la unió disjunta.}
  \opcio{c}{el quocient de $B$.}
  \opcio{d}{el subconjunt $\{a\in A\mid f(a)=g(a)\}$.}
\end{enumerate}
\feedback
\justificacio{Una fletxa $h\colon X\to A$ amb $f\comp h=g\comp h$ pren valors dins d'aquest subconjunt, i hi factoritza de manera única. L'opció (c) és el coequalitzador, el concepte dual.}

\pregunta El \textbf{pullback} de $f\colon A\to C$ i $g\colon B\to C$ a $\cat{Set}$ és:
\begin{enumerate}
  \opcio{a}{$A\times B$ sencer.}
  \opcio{b}{$\{(a,b)\in A\times B\mid f(a)=g(b)\}$.}
  \opcio{c}{la unió $A\cup B$.}
  \opcio{d}{el quocient $A/B$.}
\end{enumerate}
\feedback
\justificacio{Una parella $x_1\colon X\to A$, $x_2\colon X\to B$ amb $f\comp x_1=g\comp x_2$ factoritza de manera única per $x\mapsto(x_1(x),x_2(x))$, que pren valors en aquest subconjunt. Tot $A\times B$ només és el pullback quan $C$ té un sol element.}

\pregunta L'\textbf{objecte terminal} és el límit de quin diagrama?
\begin{enumerate}
  \opcio{a}{El diagrama de dues fletxes paral·leles.}
  \opcio{b}{El diagrama d'un sol objecte.}
  \opcio{c}{El diagrama buit.}
  \opcio{d}{El diagrama de racó.}
\end{enumerate}
\feedback
\justificacio{Un con sobre el diagrama buit és només un vèrtex, sense projeccions. Un con límit és, doncs, un objecte al qual arriba una única fletxa des de cada objecte.}

\pregunta Un \textbf{colímit} d'un diagrama $D$ és:
\begin{enumerate}
  \opcio{a}{un objecte inicial a la categoria de cocons (un límit a la categoria oposada).}
  \opcio{b}{un altre nom per al límit.}
  \opcio{c}{sempre l'objecte terminal.}
  \opcio{d}{un con.}
\end{enumerate}
\feedback
\justificacio{És la noció dual: un cocon universal, és a dir, un objecte inicial de la categoria de cocons, o equivalentment un límit del diagrama corresponent a $\cat{C}\op$. No és un límit de la mateixa categoria.}

\pregunta Si $\cat{C}$ té límits de forma $\cat{I}$ i $\cat{A}$ és petita, com es calcula el límit d'un diagrama $F\colon\cat{I}\to\cat{C}^{\cat{A}}$?
\begin{enumerate}
  \opcio{a}{Prenent el límit només sobre els objectes terminals de $\cat{A}$.}
  \opcio{b}{Aplicant el functor de parts a cada $F(i)$.}
  \opcio{c}{No en té: les categories de functors no tenen límits.}
  \opcio{d}{Punt a punt: $\bigl(\varprojlim F\bigr)(a)\cong\varprojlim_i F(i)(a)$.}
\end{enumerate}
\feedback
\justificacio{Els límits a $\cat{C}^{\cat{A}}$ es calculen objecte per objecte a $\cat{C}$; equivalentment, cada functor d'avaluació $\mathrm{ev}_a$ els preserva.}

\pregunta Quina és la caracterització dels \textbf{límits finits}?
\begin{enumerate}
  \opcio{a}{N'hi ha prou amb tenir coproductes.}
  \opcio{b}{Cal tenir tots els límits de forma petita.}
  \opcio{c}{Equivalen a tenir objecte terminal, productes binaris i equalitzadors (o terminal i pullbacks).}
  \opcio{d}{No es poden construir a partir de casos més simples.}
\end{enumerate}
\feedback
\justificacio{És el teorema de caracterització: el límit d'un diagrama finit es construeix com un equalitzador de dues fletxes entre productes finits. Amb terminal i pullbacks s'obtenen productes i equalitzadors.}

\pregunta Quan \textbf{preserva} un functor $F$ el límit d'un diagrama $D$?
\begin{enumerate}
  \opcio{a}{Quan $(F(L),F\lambda)$ és un límit de $F\comp D$.}
  \opcio{b}{Sempre, per definició de functor.}
  \opcio{c}{Quan $F$ és constant.}
  \opcio{d}{Quan $F$ no té adjunts.}
\end{enumerate}
\feedback
\justificacio{És la definició de preservació. No tot functor preserva límits: per exemple, l'oblit $\cat{Grp}\to\cat{Set}$ no preserva els coproductes, que són límits a l'oposada.}

\pregunta El functor oblit $U\colon\cat{Grp}\to\cat{Set}$:
\begin{enumerate}
  \opcio{a}{no preserva cap límit ni cap colímit.}
  \opcio{b}{preserva tots els colímits, però no tots els límits.}
  \opcio{c}{preserva tots els límits, però no tots els colímits; per exemple, no preserva tots els coproductes.}
  \opcio{d}{preserva tots els límits i tots els colímits.}
\end{enumerate}
\feedback
\justificacio{El conjunt subjacent d'un límit de grups és el límit dels conjunts subjacents (famílies compatibles). En canvi, el coproducte de grups és el producte lliure, que no és la unió disjunta dels conjunts subjacents. Més endavant es veurà que $U$ és un adjunt per la dreta.}

\pregunta En una categoria localment petita $\cat{C}$, el functor representable
\[
\Hom(A,-)\colon\cat{C}\to\cat{Set}
\]
sempre:
\begin{enumerate}
  \opcio{a}{no és un functor.}
  \opcio{b}{preserva els límits que existeixen.}
  \opcio{c}{envia límits a colímits.}
  \opcio{d}{és constant.}
\end{enumerate}
\feedback
\justificacio{Les fletxes $A\to\varprojlim D$ es corresponen bijectivament amb els cons amb vèrtex $A$, que són les famílies compatibles de $\Hom(A,D_i)$, és a dir, els elements de $\varprojlim\Hom(A,D_i)$. L'opció (c) correspon al representable contravariant: $\Hom(-,B)$ envia colímits a límits.}

\pregunta A $\cat{Set}$, siguin $f\colon A\to B$ i $y\colon Y\to B$. Què fa la reindexació $f^{*}$?
\begin{enumerate}
  \opcio{a}{Envia $Y\to B$ al coproducte $A+Y$.}
  \opcio{b}{Envia $Y\to B$ al pullback $A\times_B Y\to A$.}
  \opcio{c}{Envia $Y\to B$ al coequalitzador de $f$ i $y$.}
  \opcio{d}{Només està definida quan $f$ és un isomorfisme.}
\end{enumerate}
\feedback
\justificacio{La fibra de $A\times_B Y\to A$ sobre $a$ està en bijecció amb la fibra $y^{-1}(f(a))$: la reindexació és la substitució al llarg de $f$.}

\pregunta En el diagrama següent, quina propietat universal fa de $P$ el \textbf{pullback} de $f$ i $g$?
\[
\begin{tikzpicture}[baseline=(m.center)]
  \matrix (m) [matrix of math nodes, row sep=2.6em, column sep=2.6em]
    { P & A \\ B & C \\ };
  \draw[-{Stealth[scale=1.1]}] (m-1-1) -- node[above] {$p_1$} (m-1-2);
  \draw[-{Stealth[scale=1.1]}] (m-1-1) -- node[left] {$p_2$} (m-2-1);
  \draw[-{Stealth[scale=1.1]}] (m-1-2) -- node[right] {$f$} (m-2-2);
  \draw[-{Stealth[scale=1.1]}] (m-2-1) -- node[below] {$g$} (m-2-2);
\end{tikzpicture}
\]
\begin{enumerate}
  \opcio{a}{Que per a tot objecte $X$ amb fletxes $x_1\colon X\to A$, $x_2\colon X\to B$ tals que $f\comp x_1=g\comp x_2$, hi ha un únic $u\colon X\to P$ amb $p_1\comp u=x_1$ i $p_2\comp u=x_2$.}
  \opcio{b}{Que $f\comp p_1=g\comp p_2$ i res més.}
  \opcio{c}{Que $P=A\times B$ sempre.}
  \opcio{d}{Que $f$ i $g$ són isomorfismes.}
\end{enumerate}
\feedback
\justificacio{És la propietat universal del pullback: el quadrat commuta i qualsevol altre quadrat commutatiu sobre $f$ i $g$ hi factoritza de manera única. L'opció (b) només diu que el quadrat commuta, i això no basta.}

\end{enumerate}
\clauestatica
\finalitzaTest

\chapter[Adjuncions]{Adjuncions}

Aquest test recull els conceptes essencials del capítol. Cada pregunta té una única resposta correcta.

\iniciTest{adj}{b,c,b,d,a,b,c,d,b,c,a,d,a,a}

\begin{enumerate}

\pregunta Què vol dir que $f\dashv g$ entre dos preordres?
\begin{enumerate}
  \opcio{a}{Que $f=g$.}
  \opcio{b}{Que $f(x)\leq y$ si i només si $x\leq g(y)$, per a tot $x,y$.}
  \opcio{c}{Que $f$ i $g$ són inverses.}
  \opcio{d}{Que $f$ i $g$ són constants.}
\end{enumerate}
\feedback
\justificacio{És la definició d'adjunció entre preordres: $f$ és adjunt per l'esquerra de $g$. No cal que siguin inverses: $f(x)\leq y\Leftrightarrow x\leq g(y)$ és més feble que $g\comp f=\id$ i $f\comp g=\id$.}

\pregunta Fixada una fórmula $A$, quina parella d'operacions forma una adjunció al preordre $\cat{Prov}_T$ d'un sistema proposicional?
\begin{enumerate}
  \opcio{a}{La negació amb ella mateixa.}
  \opcio{b}{La disjunció amb la conjunció.}
  \opcio{c}{$A\wedge-$ (esquerra) amb $A\to-$ (dreta).}
  \opcio{d}{La implicació amb la negació.}
\end{enumerate}
\feedback
\justificacio{$A\wedge B\vdash_T C$ si i només si $B\vdash_T A\to C$: és la regla de la deducció (introducció i eliminació de la implicació). Les altres parelles no satisfan cap equivalència d'aquesta forma.}

\pregunta Una \textbf{adjunció} $F\dashv G$ entre categories és:
\begin{enumerate}
  \opcio{a}{una igualtat de functors.}
  \opcio{b}{una bijecció natural $\Hom(F(A),B)\cong\Hom(A,G(B))$.}
  \opcio{c}{un isomorfisme de categories.}
  \opcio{d}{una transformació natural qualsevol.}
\end{enumerate}
\feedback
\justificacio{És la definició: una família de bijeccions $\theta_{A,B}$, natural en $A$ i en $B$. Un isomorfisme de categories és un cas molt particular, i una transformació natural qualsevol no relaciona homconjunts.}

\pregunta Sigui $\cat{I}$ una categoria petita i suposem que $\cat{C}$ té tots els límits i colímits de forma $\cat{I}$. Quina és la relació entre el functor diagonal $\Delta\colon\cat{C}\to\cat{C}^{\cat{I}}$, el límit i el colímit?
\begin{enumerate}
  \opcio{a}{$\varprojlim_{\cat{I}}\dashv\Delta\dashv\colimop_{\cat{I}}$.}
  \opcio{b}{$\Delta\dashv\colimop_{\cat{I}}$ i $\Delta\dashv\varprojlim_{\cat{I}}$.}
  \opcio{c}{El límit i el colímit són tots dos adjunts per l'esquerra de $\Delta$.}
  \opcio{d}{$\colimop_{\cat{I}}\dashv\Delta\dashv\varprojlim_{\cat{I}}$.}
\end{enumerate}
\feedback
\justificacio{Un morfisme $\Delta A\Rightarrow D$ és un con, que correspon a una fletxa $A\to\varprojlim D$; un morfisme $D\Rightarrow\Delta A$ és un cocon, que correspon a una fletxa $\colimop D\to A$.}

\pregunta La \textbf{unitat} d'una adjunció $F\dashv G$ és una transformació natural:
\begin{enumerate}
  \opcio{a}{$\id\Rightarrow G F$.}
  \opcio{b}{$F\Rightarrow G$.}
  \opcio{c}{$F G\Rightarrow\id$.}
  \opcio{d}{$G\Rightarrow F$.}
\end{enumerate}
\feedback
\justificacio{$\eta_A=\theta(\id_{F(A)})\colon A\to G(F(A))$. L'opció (c) és la counitat. Les opcions (b) i (d) no tenen sentit en general, perquè $F$ i $G$ van en sentits contraris i no comparteixen domini.}

\pregunta La \textbf{counitat} d'una adjunció $F\dashv G$ és:
\begin{enumerate}
  \opcio{a}{$\id\Rightarrow G F$.}
  \opcio{b}{$F G\Rightarrow\id$.}
  \opcio{c}{una igualtat.}
  \opcio{d}{un objecte terminal.}
\end{enumerate}
\feedback
\justificacio{$\varepsilon_B=\theta^{-1}(\id_{G(B)})\colon F(G(B))\to B$. L'opció (a) és la unitat. La counitat és una transformació natural, no una igualtat ni un objecte.}

\pregunta El transposat de $\varphi\colon F(A)\to B$ sota l'adjunció és:
\begin{enumerate}
  \opcio{a}{$\varepsilon_B\comp F(\varphi)$.}
  \opcio{b}{$\varphi$ mateix.}
  \opcio{c}{$G(\varphi)\comp\eta_A$.}
  \opcio{d}{$\eta_A\comp G(\varphi)$.}
\end{enumerate}
\feedback
\justificacio{Comprova els tipus: amb $\varphi\colon F(A)\to B$, $G(\varphi)\colon G(F(A))\to G(B)$, i compondre amb $\eta_A\colon A\to G(F(A))$ dona $A\to G(B)$. L'opció (a) és la fórmula del transposat \emph{invers}, per a $\psi\colon A\to G(B)$; l'opció (d) no és componible.}

\pregunta Les \textbf{identitats triangulars} relacionen:
\begin{enumerate}
  \opcio{a}{dos objectes terminals.}
  \opcio{b}{dos límits.}
  \opcio{c}{dues categories isomorfes.}
  \opcio{d}{la unitat i la counitat, i caracteritzen l'adjunció.}
\end{enumerate}
\feedback
\justificacio{Són $\varepsilon_{F(A)}\comp F(\eta_A)=\id_{F(A)}$ i $G(\varepsilon_B)\comp\eta_{G(B)}=\id_{G(B)}$. A més, dues transformacions naturals que les satisfan defineixen una adjunció: és la segona presentació.}

\pregunta Sigui $U\colon\cat{Mon}\to\cat{Set}$ el functor oblit. Quin fet mostra que $U$ no preserva tots els colímits?
\begin{enumerate}
  \opcio{a}{No preserva productes.}
  \opcio{b}{El monoide trivial és inicial a $\cat{Mon}$, però el seu conjunt subjacent no és $\varnothing$.}
  \opcio{c}{$\cat{Mon}$ no té objecte inicial.}
  \opcio{d}{Cap adjunt per la dreta no pot preservar colímits.}
\end{enumerate}
\feedback
\justificacio{L'objecte inicial és el colímit del diagrama buit. L'objecte inicial de $\cat{Set}$ és $\varnothing$, però $U$ envia l'inicial de $\cat{Mon}$ a un conjunt d'un element; per tant, no preserva aquest colímit.}

\pregunta Els adjunts són únics:
\begin{enumerate}
  \opcio{a}{de manera estricta (igualtat).}
  \opcio{b}{mai.}
  \opcio{c}{llevat d'isomorfisme natural.}
  \opcio{d}{només en preordres.}
\end{enumerate}
\feedback
\justificacio{Si $F\dashv G$ i $F\dashv G'$, aleshores $\Hom(-,G(B))\cong\Hom(-,G'(B))$ naturalment, i Yoneda dona $G(B)\cong G'(B)$, naturalment en $B$. No hi ha igualtat estricta: qualsevol functor isomorf a un adjunt també ho és.}

\pregunta Què diu el teorema \enquote{RAPL}?
\begin{enumerate}
  \opcio{a}{Que els adjunts per la dreta preserven límits.}
  \opcio{b}{Que tot functor és adjunt.}
  \opcio{c}{Que els límits no existeixen.}
  \opcio{d}{Que els adjunts per l'esquerra preserven límits.}
\end{enumerate}
\feedback
\justificacio{RAPL abreuja \emph{right adjoints preserve limits}. L'opció (d) confon els costats: els adjunts per l'esquerra preserven colímits.}

\pregunta Per què el functor graf subjacent $U\colon\cat{Cat}\to\cat{Graph}$ preserva límits?
\begin{enumerate}
  \opcio{a}{Perquè és ple.}
  \opcio{b}{Perquè és adjunt per l'esquerra.}
  \opcio{c}{Perquè tots els functors preserven límits.}
  \opcio{d}{Perquè és adjunt per la dreta de $\Free$.}
\end{enumerate}
\feedback
\justificacio{$\Free\dashv U$, i els adjunts per la dreta preserven límits.}

\pregunta L'adjunció \enquote{producte-exponencial} $-\times B\dashv(-)^B$ correspon a:
\begin{enumerate}
  \opcio{a}{la currificació $\Hom(A\times B,C)\cong\Hom(A,C^B)$.}
  \opcio{b}{la dualitat.}
  \opcio{c}{la negació.}
  \opcio{d}{la composició de functors.}
\end{enumerate}
\feedback
\justificacio{Fixat $B$, una fletxa $A\times B\to C$ correspon a una fletxa $A\to C^B$, naturalment en $A$ i en $C$. És exactament la bijecció de l'adjunció $-\times B\dashv(-)^B$.}

\pregunta Si $f^{*}$ admet adjunts a tots dos costats, quin paper tenen $\exists_f$ i $\forall_f$?
\begin{enumerate}
  \opcio{a}{Com l'adjunt per l'esquerra i l'adjunt per la dreta de $f^{*}$, respectivament.}
  \opcio{b}{Com objectes terminals.}
  \opcio{c}{Com functors constants.}
  \opcio{d}{Com isomorfismes de categories.}
\end{enumerate}
\feedback
\justificacio{A $\cat{Set}$, $\exists_f$ és la imatge i $\forall_f(S)$ el conjunt dels $y$ amb fibra continguda en $S$; per a una projecció, són la quantificació existencial i la universal.}

\end{enumerate}
\clauestatica
\finalitzaTest

\chapter[Cartesianes tancades]{\texorpdfstring{Categories cartesianes\\ tancades}{Categories cartesianes tancades}}

Aquest test recull els conceptes essencials del capítol. Cada pregunta té una única resposta correcta.

\iniciTest{cct}{a,c,d,b,c,a,d,a,b,c,d,b}

\begin{enumerate}

\pregunta Una categoria és \textbf{cartesiana} quan té:
\begin{enumerate}
  \opcio{a}{objecte terminal i productes binaris (és a dir, tots els productes finits).}
  \opcio{b}{tots els colímits.}
  \opcio{c}{només un objecte inicial.}
  \opcio{d}{exponencials però no productes.}
\end{enumerate}
\feedback
\justificacio{El terminal és el producte buit, i els productes de qualsevol família finita s'obtenen per inducció a partir dels binaris. Els colímits i els exponencials no formen part de la definició.}

\pregunta Quina bijecció natural expressa la propietat universal de l'exponencial $C^{B}$?
\begin{enumerate}
  \opcio{a}{$\Hom(A,C)\cong\Hom(B,C)$.}
  \opcio{b}{$\Hom(A+B,C)\cong\Hom(A,C)\times\Hom(B,C)$.}
  \opcio{c}{$\Hom(A\times B,C)\cong\Hom(A,C^{B})$.}
  \opcio{d}{$\Hom(A,C\times B)\cong\Hom(A,C^{B})$.}
\end{enumerate}
\feedback
\justificacio{És la currificació $f\mapsto\lambda f$, natural en $A$ i en $C$: expressa l'adjunció $-\times B\dashv(-)^{B}$.}

\pregunta En una CCC, quin morfisme indueix $f\colon B'\to B$ entre exponencials amb el mateix $C$?
\begin{enumerate}
  \opcio{a}{$C^{B'}\to C^{B}$, perquè l'exponencial és covariant en totes dues variables.}
  \opcio{b}{Cap: l'exponencial no depèn functorialment de $B$.}
  \opcio{c}{$B^{C}\to B'^{C}$.}
  \opcio{d}{$C^{B}\to C^{B'}$: l'exponencial és contravariant en l'exponent.}
\end{enumerate}
\feedback
\justificacio{$C^{f}$ és la transposada de $C^{B}\times B'\xrightarrow{\id\times f}C^{B}\times B\xrightarrow{\ev}C$.}

\pregunta La \textbf{currificació} $\lambda f$ d'un morfisme $f\colon A\times B\to C$ és l'únic morfisme que compleix:
\begin{enumerate}
  \opcio{a}{$\lambda f\comp\ev=f$.}
  \opcio{b}{$\ev\comp(\lambda f\times\id_B)=f$.}
  \opcio{c}{$\lambda f=f$.}
  \opcio{d}{$\lambda f\comp\pi_1=f$.}
\end{enumerate}
\feedback
\justificacio{És el triangle de la definició: currificar i després avaluar torna $f$. L'opció (a) no és ni tan sols componible, perquè $\ev$ surt de $C^B\times B$ i $\lambda f$ hi arriba només des de $A$.}

\pregunta Una categoria és \textbf{cartesiana tancada} quan és cartesiana i, a més, per a cada objecte $B$:
\begin{enumerate}
  \opcio{a}{$B$ és terminal.}
  \opcio{b}{$-\times B$ és un isomorfisme.}
  \opcio{c}{el functor $-\times B$ té adjunt per la dreta $(-)^{B}$.}
  \opcio{d}{$-\times B$ té adjunt per l'esquerra.}
\end{enumerate}
\feedback
\justificacio{Tancar cartesianament vol dir que la currificació és sempre disponible: per a cada $B$, $-\times B\dashv(-)^B$. Un adjunt per l'esquerra de $-\times B$ seria una altra cosa, i no existeix en general.}

\pregunta A $\cat{Set}$, l'exponencial $C^{B}$ és:
\begin{enumerate}
  \opcio{a}{el conjunt de \emph{totes} les aplicacions $B\to C$.}
  \opcio{b}{el producte cartesià $C\times B$.}
  \opcio{c}{la unió disjunta $C+B$.}
  \opcio{d}{el conjunt de bijeccions $B\to C$.}
\end{enumerate}
\feedback
\justificacio{L'avaluació és $\ev(g,b)=g(b)$, i la currificació de $f\colon A\times B\to C$ és $a\mapsto(b\mapsto f(a,b))$. Cal prendre totes les aplicacions, no només les bijeccions, perquè $\lambda f(a)$ pot ser qualsevol aplicació.}

\pregunta En una \textbf{àlgebra de Heyting}, l'objecte exponencial $b\Rightarrow c$ és:
\begin{enumerate}
  \opcio{a}{el suprem $b\vee c$.}
  \opcio{b}{la negació $\neg b$.}
  \opcio{c}{l'ínfim $b\wedge c$.}
  \opcio{d}{l'adjunt per la dreta de $-\wedge b$, caracteritzat per $a\wedge b\leq c\Leftrightarrow a\leq(b\Rightarrow c)$.}
\end{enumerate}
\feedback
\justificacio{En un preordre, l'exponencial és l'adjunt per la dreta del producte, que és l'ínfim. En una àlgebra de Boole, a més, $b\Rightarrow c=\neg b\vee c$, però aquesta fórmula no val en una àlgebra de Heyting qualsevol.}

\pregunta Quina categoria \emph{no} és, en general, cartesiana tancada?
\begin{enumerate}
  \opcio{a}{$\cat{Top}$, els espais topològics.}
  \opcio{b}{una àlgebra de Heyting.}
  \opcio{c}{$\cat{Set}$.}
  \opcio{d}{una categoria de prefeixos $\cat{Set}^{\cat{C}\op}$.}
\end{enumerate}
\feedback
\justificacio{A $\cat{Top}$, l'avaluació no sempre té un espai de funcions contínues universal. $\cat{Set}$, les àlgebres de Heyting com a preordres i les categories de prefeixos sí que són cartesianes tancades.}

\pregunta En la semàntica del càlcul lambda simplement tipat, el \textbf{tipus funció} $\sigma\to\tau$ s'interpreta amb:
\begin{enumerate}
  \opcio{a}{el producte $\llbracket\sigma\rrbracket\times\llbracket\tau\rrbracket$.}
  \opcio{b}{l'exponencial $\llbracket\tau\rrbracket^{\llbracket\sigma\rrbracket}$.}
  \opcio{c}{l'objecte terminal.}
  \opcio{d}{el coproducte.}
\end{enumerate}
\feedback
\justificacio{Un terme de tipus $\sigma\to\tau$ és una funció, i en una CCC les funcions de $\llbracket\sigma\rrbracket$ a $\llbracket\tau\rrbracket$ formen l'exponencial. El producte interpreta, en canvi, el tipus parella $\sigma\times\tau$.}

\pregunta En aquesta semàntica, l'\textbf{abstracció} $\lambda x.\,t$ es correspon amb:
\begin{enumerate}
  \opcio{a}{una projecció.}
  \opcio{b}{la diagonal.}
  \opcio{c}{la currificació de la interpretació de $t$.}
  \opcio{d}{l'objecte terminal.}
\end{enumerate}
\feedback
\justificacio{Si $t$ té tipus $\tau$ en el context $\Gamma,x{:}\sigma$, s'interpreta com un morfisme $\llbracket\Gamma\rrbracket\times\llbracket\sigma\rrbracket\to\llbracket\tau\rrbracket$. La seva currificació és un morfisme $\llbracket\Gamma\rrbracket\to\llbracket\tau\rrbracket^{\llbracket\sigma\rrbracket}$, que interpreta $\lambda x.\,t$.}

\pregunta La regla de conversió $(\beta)$ del càlcul lambda es correspon categòricament amb:
\begin{enumerate}
  \opcio{a}{la unicitat de l'objecte terminal.}
  \opcio{b}{l'associativitat del producte.}
  \opcio{c}{la commutativitat del coproducte.}
  \opcio{d}{l'equació de l'avaluació $\ev\comp(\lambda f\times\id_B)=f$, de la qual es dedueix.}
\end{enumerate}
\feedback
\justificacio{Aplicar una abstracció a un argument és avaluar una currificació, i l'equació $\ev\comp(\lambda f\times\id)=f$ diu que el resultat és $f$, que és el que afirma la regla $(\beta)$. La regla $(\eta)$ correspon a la segona llei, $\lambda(\ev\comp(g\times\id))=g$.}

\pregunta La correspondència de \textbf{Curry--Howard--Lambek} relaciona:
\begin{enumerate}
  \opcio{a}{grups amb anells.}
  \opcio{b}{el càlcul lambda simplement tipat amb les categories cartesianes tancades.}
  \opcio{c}{límits amb colímits.}
  \opcio{d}{monoides amb preordres.}
\end{enumerate}
\feedback
\justificacio{Els tipus s'interpreten com a objectes, els termes com a morfismes i les conversions com a igualtats. En la lectura lògica, les proposicions del fragment $(\top,\wedge,\to)$ són els tipus, i les demostracions, els termes.}

\end{enumerate}
\clauestatica
\finalitzaTest

\chapter[Subobjectes, imatges i factoritzacions]{\texorpdfstring{Subobjectes, imatges\\ i factoritzacions}{Subobjectes, imatges i factoritzacions}}

Aquest test recull els conceptes essencials del capítol. Cada pregunta té una única resposta correcta.

\iniciTest{sub}{c,b,d,b,c,d,a,c,a,d,c,b,d,a,b,a,c}

\begin{enumerate}

\pregunta Un \textbf{subobjecte} d'un objecte $A$ és:
\begin{enumerate}
  \opcio{a}{qualsevol morfisme cap a $A$.}
  \opcio{b}{un element de $A$.}
  \opcio{c}{una classe d'equivalència de monomorfismes amb codomini $A$.}
  \opcio{d}{un epimorfisme des de $A$.}
\end{enumerate}
\feedback
\justificacio{Dos monomorfismes que difereixen en un isomorfisme dels dominis representen el mateix subobjecte. Un morfisme qualsevol no basta: cal que sigui mono, el paper que a $\cat{Set}$ fan les injeccions.}

\pregunta Sigui $H$ un subgraf d'un graf $G$ i $F\colon G'\to G$ un morfisme de grafs. Què és la reindexació $F^{*}H$?
\begin{enumerate}
  \opcio{a}{La imatge de $H$ per $F$.}
  \opcio{b}{El subgraf de $G'$ format per les antiimatges dels vèrtexs i de les arestes de $H$.}
  \opcio{c}{El graf $H$ mateix, vist dins de $G'$.}
  \opcio{d}{El producte de grafs $H\times G'$.}
\end{enumerate}
\feedback
\justificacio{La reindexació es calcula component a component: $F^{*}(V',E')=(F_V^{-1}(V'),F_E^{-1}(E'))$, que és tancat per origen i destí perquè $F$ els preserva.}

\pregunta A $\cat{Set}$, l'ordre parcial $\Sub(A)$ és isomorf a:
\begin{enumerate}
  \opcio{a}{el conjunt $A$.}
  \opcio{b}{el monoide de les aplicacions $A\to A$.}
  \opcio{c}{el grup de permutacions de $A$.}
  \opcio{d}{el conjunt de parts $\mathcal{P}(A)$, ordenat per inclusió.}
\end{enumerate}
\feedback
\justificacio{Dos monomorfismes cap a $A$ són equivalents si i només si tenen la mateixa imatge, i la imatge determina la classe. L'ordre correspon a la inclusió de les imatges.}

\pregunta La \textbf{parella nucli} d'una fletxa $f\colon A\to B$, quan existeix, s'obté com:
\begin{enumerate}
  \opcio{a}{l'equalitzador de $f$ i $\id_A$.}
  \opcio{b}{el pullback de $f$ amb ella mateixa.}
  \opcio{c}{el coequalitzador de $f$ amb ella mateixa.}
  \opcio{d}{la factorització imatge de $f$.}
\end{enumerate}
\feedback
\justificacio{És el parell de projeccions $p_1,p_2\colon R\to A$ del pullback de $f$ amb $f$; a $\cat{Set}$, $R=\{(a,a')\mid f(a)=f(a')\}$. Quan existeix la parella nucli, un epimorfisme regular n'és el coequalitzador.}

\pregunta Lògicament, l'ordre $[m]\leq[n]$ a $\Sub(A)$ s'interpreta com:
\begin{enumerate}
  \opcio{a}{la negació del predicat.}
  \opcio{b}{la igualtat de predicats.}
  \opcio{c}{la implicació entre predicats sobre $A$.}
  \opcio{d}{la disjunció de predicats.}
\end{enumerate}
\feedback
\justificacio{$[m]\leq[n]$ vol dir que $m$ factoritza a través de $n$: tot el que satisfà el predicat $m$ satisfà el predicat $n$. És l'ordre de la deduïbilitat entre predicats sobre $A$.}

\pregunta El \textbf{pullback} d'un monomorfisme al llarg de qualsevol morfisme és:
\begin{enumerate}
  \opcio{a}{un epimorfisme.}
  \opcio{b}{un isomorfisme.}
  \opcio{c}{mai un monomorfisme.}
  \opcio{d}{sempre un monomorfisme.}
\end{enumerate}
\feedback
\justificacio{Si $m'\comp u=m'\comp v$ al pullback, aleshores $u$ i $v$ tenen la mateixa composició amb totes dues projeccions (la segona, per la cancel·lació de $m$), i per la unicitat del pullback $u=v$.}

\pregunta El \textbf{functor de reindexació} $f^{*}\colon\Sub(B)\to\Sub(A)$ associat a $f\colon A\to B$ envia un subobjecte al seu:
\begin{enumerate}
  \opcio{a}{pullback al llarg de $f$.}
  \opcio{b}{coproducte amb $A$.}
  \opcio{c}{producte amb $B$.}
  \opcio{d}{imatge directa.}
\end{enumerate}
\feedback
\justificacio{És la definició: el pullback d'un mono és mono, i en resulta una aplicació monòtona ben definida entre classes. La imatge directa (d) correspon a l'existencial, que va en sentit contrari.}

\pregunta A $\cat{Set}$, la reindexació $f^{*}(S)$ d'un subconjunt $S\subseteq B$ és:
\begin{enumerate}
  \opcio{a}{la imatge directa $f(S)$.}
  \opcio{b}{el complement de $S$.}
  \opcio{c}{l'antiimatge $f^{-1}(S)$.}
  \opcio{d}{el producte $S\times A$.}
\end{enumerate}
\feedback
\justificacio{El pullback de la inclusió $S\hookrightarrow B$ al llarg de $f$ és $\{a\in A\mid f(a)\in S\}$. La imatge directa (a) és l'existencial $\exists_f$, l'adjunt per l'esquerra de $f^{*}$.}

\pregunta Lògicament, la reindexació $f^{*}$ interpreta:
\begin{enumerate}
  \opcio{a}{la substitució d'una variable per un terme.}
  \opcio{b}{la quantificació existencial.}
  \opcio{c}{la negació.}
  \opcio{d}{la conjunció.}
\end{enumerate}
\feedback
\justificacio{Si $\varphi(y)$ és un predicat sobre $B$ i $f$ interpreta un terme, $f^{*}\varphi$ és $\varphi(f(x))$: substituir és fer el pullback. La quantificació existencial correspon a la imatge.}

\pregunta La \textbf{intersecció} de dos subobjectes $[m]\wedge[n]$ a $\Sub(A)$ es calcula com:
\begin{enumerate}
  \opcio{a}{el coproducte de $S$ i $T$.}
  \opcio{b}{la unió de les imatges.}
  \opcio{c}{l'exponencial $T^{S}$.}
  \opcio{d}{el pullback de $m$ i $n$ sobre $A$.}
\end{enumerate}
\feedback
\justificacio{Una fletxa factoritza a través de $m$ i de $n$ si i només si factoritza a través del seu pullback. Per tant, el pullback és el major subobjecte per sota de tots dos, i interpreta la conjunció.}

\pregunta En una categoria regular, considerem un quadrat cartesià amb $h\colon A'\to A$, $f'\colon A'\to B'$, $f\colon A\to B$ i $g\colon B'\to B$, de manera que $f\comp h=g\comp f'$. Quina igualtat és la condició de Beck--Chevalley, amb els tipus correctes?
\begin{enumerate}
  \opcio{a}{$\exists_f\comp g^{*}=h^{*}\comp\exists_{f'}$.}
  \opcio{b}{$h^{*}\comp\exists_f=\exists_{f'}\comp g^{*}\colon\Sub(B)\to\Sub(A')$.}
  \opcio{c}{$g^{*}\comp\exists_f=\exists_{f'}\comp h^{*}\colon\Sub(A)\to\Sub(B')$.}
  \opcio{d}{$g^{*}\comp\exists_f=\exists_{f'}\comp h^{*}\colon\Sub(B)\to\Sub(A')$.}
\end{enumerate}
\feedback
\justificacio{Partint de $\Sub(A)$, hi ha dos camins fins a $\Sub(B')$: quantificar al llarg de $f$ i substituir al llarg de $g$, o substituir al llarg de $h$ i quantificar al llarg de $f'$. La condició diu que coincideixen: quantificar i substituir commuten. A (a) i (b) les composicions no existeixen ($\exists_f$ surt de $\Sub(A)$, no de $\Sub(B)$), i a (d) el domini i el codomini estan intercanviats.}

\pregunta Segons el capítol, en una categoria \textbf{regular} la reindexació $f^{*}$ té garantit:
\begin{enumerate}
  \opcio{a}{cap adjunt.}
  \opcio{b}{un adjunt per l'esquerra $\exists_f$ (l'universal $\forall_f$ requereix estructura addicional).}
  \opcio{c}{un adjunt per la dreta $\forall_f$, però no per l'esquerra.}
  \opcio{d}{adjunts per tots dos costats, sempre.}
\end{enumerate}
\feedback
\justificacio{L'existencial és la imatge, i la regularitat en garanteix l'existència i l'estabilitat. L'adjunt per la dreta $\forall_f$ és estructura addicional, que tenen, per exemple, els topos.}

\pregunta Una \textbf{factorització imatge} de $f$ descompon $f$ com:
\begin{enumerate}
  \opcio{a}{un mono seguit d'un epi qualsevol.}
  \opcio{b}{dos isomorfismes.}
  \opcio{c}{dos monomorfismes.}
  \opcio{d}{un epimorfisme regular seguit d'un monomorfisme.}
\end{enumerate}
\feedback
\justificacio{A $\cat{Set}$, l'epi regular és l'aplicació exhaustiva sobre la imatge, i el mono és la inclusió de la imatge. La factorització és única llevat d'un únic isomorfisme compatible.}

\pregunta La \textbf{imatge} $\Imatge(f)$ es caracteritza com:
\begin{enumerate}
  \opcio{a}{el menor subobjecte del codomini a través del qual $f$ es factoritza.}
  \opcio{b}{el major subobjecte del domini.}
  \opcio{c}{el nucli de $f$.}
  \opcio{d}{el coproducte de domini i codomini.}
\end{enumerate}
\feedback
\justificacio{La imatge és un subobjecte pel qual factoritza $f$, i és per sota de qualsevol altre que tingui aquesta propietat. A $\cat{Set}$ és el subconjunt $f(A)\subseteq B$.}

\pregunta Una categoria és \textbf{regular} si té límits finits, factoritzacions imatge i, a més:
\begin{enumerate}
  \opcio{a}{tots els colímits.}
  \opcio{b}{els epimorfismes regulars són estables per pullback.}
  \opcio{c}{és cartesiana tancada.}
  \opcio{d}{tots els objectes són projectius.}
\end{enumerate}
\feedback
\justificacio{És la condició que fa estables les imatges per reindexació i, per tant, l'existencial compatible amb la substitució. Els colímits, el tancament cartesià i els objectes projectius no formen part de la definició.}

\pregunta En una categoria regular, la \textbf{quantificació existencial} $\exists_f[s]$ es defineix com:
\begin{enumerate}
  \opcio{a}{la imatge $\Imatge(f\comp s)$.}
  \opcio{b}{el pullback de $[s]$.}
  \opcio{c}{el complement de $[s]$.}
  \opcio{d}{l'exponencial de $[s]$.}
\end{enumerate}
\feedback
\justificacio{Es fa la composició $f\comp s$ i se'n pren la imatge a $B$; a $\cat{Set}$, $f(S)$. Aquest operador és l'adjunt per l'esquerra de $f^{*}$: $\exists_f[s]\leq[t]$ si i només si $[s]\leq f^{*}[t]$.}

\pregunta Sigui $\cat{C}$ una categoria regular i $f\colon A\to B$ un morfisme. Quina d'aquestes afirmacions \emph{no} es pot deduir només de la regularitat?
\begin{enumerate}
  \opcio{a}{Tota fletxa té una factorització imatge (única llevat d'iso).}
  \opcio{b}{La reindexació $f^{*}\colon\Sub(B)\to\Sub(A)$ té un adjunt per l'esquerra $\exists_f$.}
  \opcio{c}{Cada $\Sub(A)$ té unions binàries, i $f^{*}$ les preserva.}
  \opcio{d}{Cada $\Sub(A)$ té ínfims binaris, calculats per pullback, i $f^{*}$ els preserva.}
\end{enumerate}
\feedback
\justificacio{Les unions de subobjectes estables per reindexació no formen part de la definició de regularitat: són l'estructura de les categories coherents. El mateix passa amb l'adjunt per la dreta $\forall_f$. En canvi, (a) i (b) són part del que dona la regularitat, i (d) també val: els ínfims són pullbacks, i $f^{*}$ els preserva perquè és adjunt per la dreta de $\exists_f$.}

\end{enumerate}
\clauestatica
\finalitzaTest

\chapter[Classificadors de subobjectes i topos]{\texorpdfstring{Classificadors de subobjectes i introducció als topos}{Classificadors de subobjectes i introducció als topos}}

Aquest test recull els conceptes essencials del capítol. Cada pregunta té una única resposta correcta.

\iniciTest{top}{c,b,d,a,c,d,b,a,a,b,c,d,b}

\begin{enumerate}

\pregunta Sigui $m\colon U\rightarrowtail E$ un subobjecte classificat per $\chi_m\colon E\to\Omega$. Un element generalitzat $x\colon X\to E$ factoritza a través de $m$ exactament quan:
\begin{enumerate}
  \opcio{a}{$\chi_m\comp x$ és un isomorfisme.}
  \opcio{b}{$x$ és un monomorfisme.}
  \opcio{c}{$\chi_m\comp x=\mathsf{true}\comp\,!_X$.}
  \opcio{d}{$X$ és l'objecte terminal.}
\end{enumerate}
\feedback
\justificacio{És la propietat universal del quadrat cartesià de $m$: la pertinença d'un element generalitzat esdevé una igualtat de morfismes cap a $\Omega$.}

\pregunta A la categoria de prefeixos sobre $\cat{C}$, siguin $h\colon D\to C$ i $S$ un sedàs sobre $C$. Quin sedàs sobre $D$ és $\Omega(h)(S)$?
\begin{enumerate}
  \opcio{a}{$\{h\comp g\mid g\in S\}$.}
  \opcio{b}{$\{g\colon X\to D\mid h\comp g\in S\}$.}
  \opcio{c}{$S$ mateix.}
  \opcio{d}{El sedàs màxim sobre $D$.}
\end{enumerate}
\feedback
\justificacio{$\Omega$ actua per reindexació de sedassos: $\Omega(h)(S)=h^{*}S$, i aquest conjunt torna a ser un sedàs.}

\pregunta La propietat universal del classificador diu que cada subobjecte $U\rightarrowtail E$ determina:
\begin{enumerate}
  \opcio{a}{un isomorfisme $E\cong\Omega$.}
  \opcio{b}{cap morfisme.}
  \opcio{c}{diversos morfismes $E\to\Omega$.}
  \opcio{d}{un únic morfisme $\chi\colon E\to\Omega$ que fa un quadrat cartesià amb $\mathsf{true}$.}
\end{enumerate}
\feedback
\justificacio{És la definició de classificador: existència i unicitat de $\chi$, amb $U$ com a pullback de $\mathsf{true}$ al llarg de $\chi$. Que el quadrat només commuti no basta per a la unicitat.}

\pregunta L'existència d'un classificador de subobjectes equival a dir que el functor $\Sub(-)$ és:
\begin{enumerate}
  \opcio{a}{representable, amb objecte representant $\Omega$.}
  \opcio{b}{constant.}
  \opcio{c}{no functorial.}
  \opcio{d}{un isomorfisme.}
\end{enumerate}
\feedback
\justificacio{La bijecció $\Sub(E)\cong\Hom(E,\Omega)$, natural en $E$, és exactament una representació de $\Sub$, amb element universal $\mathsf{true}$.}

\pregunta El classificador de subobjectes, quan existeix, és únic:
\begin{enumerate}
  \opcio{a}{només a $\cat{Set}$.}
  \opcio{b}{mai.}
  \opcio{c}{llevat d'un únic isomorfisme compatible amb $\mathsf{true}$ (per Yoneda).}
  \opcio{d}{llevat d'homotopia.}
\end{enumerate}
\feedback
\justificacio{Dos objectes que representen el mateix functor són isomorfs per un únic isomorfisme compatible amb els elements universals: és el principi de Yoneda aplicat a $\Sub$.}

\pregunta Un \textbf{topos elemental} és una categoria amb límits finits, exponencials i:
\begin{enumerate}
  \opcio{a}{tots els colímits infinits.}
  \opcio{b}{un objecte inicial estricte.}
  \opcio{c}{una estructura de grup.}
  \opcio{d}{un classificador de subobjectes.}
\end{enumerate}
\feedback
\justificacio{Límits finits, exponencials i classificador de subobjectes: no es demanen colímits infinits. $\cat{FinSet}$, que no en té, és un topos.}

\pregunta Quina d'aquestes categories és un topos?
\begin{enumerate}
  \opcio{a}{la categoria dels grups.}
  \opcio{b}{$\cat{Set}$.}
  \opcio{c}{$\cat{Top}$.}
  \opcio{d}{un monoide no trivial vist com a categoria.}
\end{enumerate}
\feedback
\justificacio{$\cat{Set}$ té límits finits, exponencials i el classificador $\{0,1\}$. $\cat{Top}$ no és cartesiana tancada, i un monoide no trivial vist com a categoria no té ni tan sols tots els límits finits.}

\pregunta Per a $\cat{C}$ petita, la categoria de prefeixos $\cat{Set}^{\cat{C}\op}$ és:
\begin{enumerate}
  \opcio{a}{sempre un topos.}
  \opcio{b}{un topos només si $\cat{C}$ és un grup.}
  \opcio{c}{mai un topos.}
  \opcio{d}{cartesiana tancada però sense classificador.}
\end{enumerate}
\feedback
\justificacio{Té límits finits i exponencials (capítols de límits i de cartesianes tancades) i el classificador de sedassos. No cal cap condició sobre $\cat{C}$ més enllà de ser petita.}

\pregunta A $\cat{Graph}$, sigui $x\colon v\to w$ una aresta d'un graf $G$ i $H$ un subgraf que conté els vèrtexs $v$ i $w$ però no l'aresta $x$. A quina aresta del classificador $\Omega$ envia $\chi_H$ l'aresta $x$?
\begin{enumerate}
  \opcio{a}{Al bucle $\partial$ del vèrtex $1$, que registra \enquote{extrems sí, aresta no}.}
  \opcio{b}{Al bucle $\top$ del vèrtex $1$, perquè els dos extrems són a $H$.}
  \opcio{c}{A l'aresta $\sigma\colon1\to0$.}
  \opcio{d}{Enlloc: $\chi_H$ només actua sobre els vèrtexs.}
\end{enumerate}
\feedback
\justificacio{Com que $v,w\in H$, la imatge de $x$ ha d'anar de $1$ a $1$, i n'hi ha dues: $\partial$ i $\top$. Si fos $\top$, $x$ seria al pullback de $\mathsf{true}$, que ha de ser exactament $H$; per tant, va a $\partial$. L'opció (c) tindria el destí equivocat, i un morfisme de grafs actua també sobre les arestes.}

\pregunta En un topos, l'objecte de parts $\mathcal P(E)=\Omega^{E}$ satisfà, per a tot objecte $A$,
\[
\Hom(A,\mathcal P(E))\;\cong\;\Sub(A\times E).
\]
Quina interpretació expressa aquesta bijecció?
\begin{enumerate}
  \opcio{a}{Que tot objecte és isomorf al seu objecte de parts.}
  \opcio{b}{Que els morfismes cap a $\mathcal P(E)$ classifiquen les relacions entre $A$ i $E$.}
  \opcio{c}{Que $\mathcal P(E)$ és l'objecte terminal.}
  \opcio{d}{Que $E$ té exactament dos subobjectes.}
\end{enumerate}
\feedback
\justificacio{Combina la currificació $\Hom(A,\Omega^{E})\cong\Hom(A\times E,\Omega)$ amb la representabilitat de $\Sub$; per a $A=1$ recupera $\Hom(1,\mathcal P(E))\cong\Sub(E)$.}

\pregunta Quina afirmació sobre la categoria $\cat{FinSet}$ dels conjunts finits és correcta?
\begin{enumerate}
  \opcio{a}{No és un topos, perquè no té coproductes infinits.}
  \opcio{b}{És equivalent a $\cat{Set}^{\cat{C}\op}$ per a alguna categoria petita $\cat{C}$.}
  \opcio{c}{És un topos elemental que no és equivalent a cap categoria de prefeixos.}
  \opcio{d}{Només és un topos si s'hi afegeixen els conjunts infinits.}
\end{enumerate}
\feedback
\justificacio{Té límits finits, exponencials finits i el classificador $\{0,1\}$, de manera que és un topos. Tota categoria de prefeixos té colímits petits, i $\cat{FinSet}$ no té coproductes infinits; com que una equivalència conserva els colímits, no és equivalent a cap. Un topos elemental no demana colímits infinits.}

\pregunta A $\cat{Graph}$, $\Omega$ té dos vèrtexs i cinc arestes. Quin fet mostra que $\cat{Graph}$ no és un topos booleà?
\begin{enumerate}
  \opcio{a}{Que $\Omega$ no és un conjunt.}
  \opcio{b}{Que $\Omega$ té més de dos vèrtexs.}
  \opcio{c}{Que $\Hom(\mathbf L,\Omega)$ té exactament dos elements.}
  \opcio{d}{Que $\{t\}$ és el més gran dels subgrafs de $\mathbf E$ disjunts de $\{s\}$, però $\{s\}\cup\{t\}$ no és tot $\mathbf E$.}
\end{enumerate}
\feedback
\justificacio{Per a $U=\{s\}$, el seu \enquote{complement} $\neg U=\{t\}$ no cobreix, juntament amb $U$, tot $\mathbf E$: hi falta l'aresta. Falla, doncs, $U\vee\neg U=\top$. L'opció (b) és falsa ($\Omega$ té dos vèrtexs), i també la (c) (n'hi ha tres).}

\pregunta Un topos té exactament dos morfismes $1\to\Omega$. Què se'n pot concloure?
\begin{enumerate}
  \opcio{a}{Que és booleà, perquè $\Omega$ té exactament dos valors de veritat.}
  \opcio{b}{Res sobre la booleanitat: els morfismes $1\to\Omega$ només compten $\Sub(1)$, mentre que ser booleà és una condició sobre tots els $\Sub(E)$.}
  \opcio{c}{Que $\Omega\cong1+1$.}
  \opcio{d}{Que el topos és equivalent a $\cat{Set}$.}
\end{enumerate}
\feedback
\justificacio{$\Hom(1,\Omega)\cong\Sub(1)$. Els prefeixos sobre el monoide $(\mathbb N,+)$ tenen exactament dos valors globals, però els subobjectes del representable formen una cadena $\mathbb N\cup\{\infty\}$, que no és una àlgebra de Boole, i aquest $\Omega$ no és $1+1$. Tenir dos valors globals i ser booleà són coses diferents.}

\end{enumerate}
\clauestatica
\finalitzaTest

\fi
\ifimprimible\else\end{Form}\fi

% ============================================================
% APÈNDIXS NUMERATS
% \appendix ha d'anar dins \mainmatter (abans de \backmatter):
% a la classe book, \backmatter desactiva la numeració de capítols,
% de manera que els capítols d'apèndix es numeren A, B, ... aquí.
% ============================================================
\appendix
\part{Apèndixs}
\chapter[Axiomàtica relacional]{Una axiomàtica relacional de les categories petites}
\label{cap:axiomatica-relacional}

\begin{objectius}
\apartat{Prerequisits} La definició de categoria, la categoria oposada i la formulació operativa del principi de dualitat (capítol~\ref{cap:categories}, secció~\ref{sec:dualitat}). De lògica de primer ordre amb igualtat: la sintaxi (fórmules, sentències, variables lliures i lligades, substitució), la semàntica (estructures, satisfacció i models), la noció de deducció en un càlcul de Hilbert i les demostracions per inducció sobre fórmules. El teorema de completesa només s'esmenta. Les referències a la lògica lliure de la secció~\ref{sec:ar-relacional-funcional} són informatives i no es pressuposen.
\apartat{Objectius} En acabar l'apèndix, el lector hauria de poder:
\begin{itemize}
  \item presentar una categoria petita en un llenguatge de primer ordre \emph{amb una sola mena d'individus}, on tot individu és una fletxa i $\Dom$, $\Cod$ i $\Com$ són les úniques relacions primitives;
  \item distingir el llenguatge formal, els axiomes i el metallenguatge, i reconèixer quines nocions són primitives i quines en són derivades;
  \item deduir dels axiomes l'estabilitat dels extrems, la unicitat de la composició i la derivabilitat del tercer axioma;
  \item reconstruir els objectes com les fletxes identitat i justificar el pas a la notació funcional usual;
  \item explicar en quin sentit la presentació relacional és equivalent a la definició usual, fins a isomorfisme canònic;
  \item comparar la presentació relacional amb una presentació funcional amb operacions parcials;
  \item definir la traducció dual de fórmules, comprovar l'autodualitat dels axiomes i demostrar el principi de dualitat en les versions sintàctica i semàntica, relacionant-lo amb la versió operativa del capítol~\ref{cap:categories}.
\end{itemize}
\end{objectius}

La definició habitual d'una categoria (definició~\ref{def:categoria}) introdueix objectes, fletxes, identitats i composició. En aquest apèndix adoptem un punt de vista més auster: el llenguatge categòric primitiu només parla de fletxes. Allò que habitualment anomenem objectes es reconstrueix després com una propietat especial d'algunes fletxes. És, en un sentit ben literal, el programa de la introducció ---\emph{relaciona, no individualitza} (secció~\ref{sec:mirar-fora})--- dut fins a la seva conclusió sintàctica: si el que compta d'un objecte és el seu paper dins el sistema de fletxes, aleshores no cal introduir-lo com a dada primitiva.

L'objectiu no és substituir la presentació usual, que és molt més còmoda per al treball matemàtic ordinari. L'objectiu és separar amb precisió tres nivells: el \emph{llenguatge formal} de la teoria; els \emph{axiomes} formulats en aquest llenguatge; i el \emph{metallenguatge}, on interpretem les fórmules i introduïm notació derivada. Aquesta separació permet veure quines nocions són primitives i quines són conseqüència dels axiomes. En particular, la funcionalitat de la composició no s'incorpora a la sintaxi: es demostrarà com un teorema.

\begin{nota}
La idea de descriure una categoria \enquote{només amb fletxes}, amb els objectes identificats amb les identitats, és clàssica: apareix ja en Mac~Lane~\cite{maclane}, es formula amb una relació ternària de composició en la teoria elemental de categories de Lawvere~\cite{lawvere-ccfm}, i s'estudia sistemàticament, amb operacions parcials, en Scott~\cite{scott-identity} i en Freyd i Scedrov~\cite{freydscedrov}. Diversos d'aquests sistemes han estat comparats i verificats formalment per Benzmüller i Scott~\cite{BS2020,BSAFP}. La contribució d'aquest apèndix no és, doncs, la idea, sinó una presentació concreta: una axiomàtica de primer ordre \emph{purament relacional} ---fins i tot el domini i el codomini són relacions--- en què la funcionalitat de la composició i el comportament dels seus extrems esdevenen \emph{teoremes}.
\end{nota}

\section{El llenguatge formal}
\label{sec:ar-llenguatge}

Treballarem en lògica clàssica de primer ordre amb igualtat i amb una sola mena d'individus. La signatura no lògica és
\[
\mathcal L_{\mathrm{Cat}}
=
\{\Dom^{2},\Cod^{2},\Com^{3}\}.
\]
No hi ha constants ni símbols de funció. Per tant, els termes del llenguatge són només variables individuals $f,g,h,d,c,e,\dots$, i totes recorren un mateix univers.

Les fórmules atòmiques no lògiques són de les formes $\Dom(f,d)$, $\Cod(f,c)$ i $\Com(f,g,h)$. La seva lectura metalingüística serà:
\begin{align*}
\Dom(f,d)&:\quad \text{\enquote{el domini de $f$ és $d$}},\\
\Cod(f,c)&:\quad \text{\enquote{el codomini de $f$ és $c$}},\\
\Com(f,g,h)&:\quad \text{\enquote{primer $f$, després $g$, i la composta és $h$}}.
\end{align*}
Així, quan la notació funcional ja estigui justificada, $\Com(f,g,h)$ correspondrà metalingüísticament a $h=g\comp f$.

La paraula \emph{fletxa} pertany, en aquest punt, al metallenguatge: no hi ha cap predicat primitiu $\operatorname{Fletxa}(f)$. Tots els individus del model s'interpreten com a fletxes. Una estructura per a $\mathcal L_{\mathrm{Cat}}$ és, doncs,
\[
\mathfrak M=(M,\Dom^{\mathfrak M},\Cod^{\mathfrak M},\Com^{\mathfrak M}),
\]
on $M$ és l'univers, $\Dom^{\mathfrak M},\Cod^{\mathfrak M}\subseteq M^2$ i $\Com^{\mathfrak M}\subseteq M^3$.

\begin{observacio}[metateòrica]
\label{obs:ar-metateorica}
Amb la semàntica estàndard de primer ordre, $M$ és un conjunt \emph{no buit}; per tant, els models en sentit estricte corresponen a categories petites no buides. La categoria buida queda exclosa per la convenció habitual que l'univers d'una estructura de primer ordre no és buit; si es vol incloure-la, cal permetre explícitament estructures amb univers buit. Les categories grans, d'altra banda, exigeixen una convenció metateòrica addicional ---classes pròpies o universos de Grothendieck, en la línia de l'apèndix~\ref{cap:mida}---, però això no modifica el llenguatge categòric primitiu anterior.
\end{observacio}

\section{Els axiomes}
\label{sec:ar-axiomes}

Presentem primer el sistema en cinc blocs conceptuals. A la secció~\ref{sec:ar-redundancia} veurem que, amb una formulació prou forta de l'associativitat, el tercer bloc és derivable dels altres; la base final del sistema és, doncs, \eqref{ax:A1D}--\eqref{ax:A1C}, \eqref{ax:A2}, \eqref{ax:A4E}--\eqref{ax:A4U} i \eqref{ax:A5D}--\eqref{ax:A5C}.

\paragraph{Axioma 1: existència i unicitat de domini i codomini.}
Tota fletxa té exactament un domini i exactament un codomini:
\begin{align}
\forall f\,\exists!d\,\Dom(f,d),
\tag{A1D}\label{ax:A1D}\\
\forall f\,\exists!c\,\Cod(f,c).
\tag{A1C}\label{ax:A1C}
\end{align}
Com és habitual, $\exists!$ és només una abreviació de la lògica de primer ordre amb igualtat.

\paragraph{Axioma 2: existència de la composició.}
Dues fletxes es poden compondre exactament quan el codomini de la primera coincideix amb el domini de la segona:
\begin{equation}
\forall f\,\forall g\,
\left[
\exists e\bigl(\Cod(f,e)\land\Dom(g,e)\bigr)
\leftrightarrow
\exists h\,\Com(f,g,h)
\right].
\tag{A2}\label{ax:A2}
\end{equation}
Aquest axioma afirma existència, però no unicitat. La unicitat de la composta serà un teorema (secció~\ref{sec:ar-unicitat}).

\paragraph{Axioma 3: extrems de la composta.}
Si $h$ és una composta de $f$ amb $g$, el domini de $h$ és el de $f$ i el codomini de $h$ és el de $g$:
\begin{equation}
\begin{split}
\forall f,g,h,d,c\,\bigl[
&\Com(f,g,h)\land\Dom(f,d)\land\Cod(g,c)\\
&\longrightarrow
\Dom(h,d)\land\Cod(h,c)
\bigr].
\end{split}
\tag{A3}\label{ax:A3}
\end{equation}

\paragraph{Axioma 4: associativitat forta.}
Com que $\Com$ és una relació i la composició és parcial, la mera igualtat de dues composicions iterades ja existents no expressa tota la força de l'associativitat usual. Adoptarem una formulació en dues parts. Primer, si existeixen les dues composicions consecutives $g\comp f$ i $h\comp g$, existeixen també les dues composicions iterades:
\begin{equation}
\begin{split}
\forall f,g,h,u,v\,\bigl[
&\Com(f,g,u)\land\Com(g,h,v)\\
&\longrightarrow
\exists p\,\exists q\,
\bigl(\Com(u,h,p)\land\Com(f,v,q)\bigr)
\bigr].
\end{split}
\tag{A4E}\label{ax:A4E}
\end{equation}
Segon, qualssevol resultats obtinguts per les dues associacions coincideixen:
\begin{equation}
\begin{split}
\forall f,g,h,u,v,p,q\,\bigl[
&\Com(f,g,u)\land\Com(g,h,v)\land\\
&\Com(u,h,p)\land\Com(f,v,q)
\longrightarrow p=q
\bigr].
\end{split}
\tag{A4U}\label{ax:A4U}
\end{equation}
Els dos enunciats, \eqref{ax:A4E} i \eqref{ax:A4U}, formen el bloc d'associativitat (A4). El primer controla l'existència de les composicions iterades; el segon n'expressa la coherència. Un cop demostrada la funcionalitat de $\Com$ (secció~\ref{sec:ar-unicitat}), aquest bloc es llegeix com la llei usual d'un semigrup parcial: si $g\comp f$ i $h\comp g$ existeixen, aleshores existeixen $h\comp(g\comp f)$ i $(h\comp g)\comp f$, i són iguals.

\paragraph{Axioma 5: neutralitat.}
El domini és neutre per la dreta i el codomini és neutre per l'esquerra:
\begin{align}
\forall f\,\forall d\,
\bigl(\Dom(f,d)\rightarrow\Com(d,f,f)\bigr),
\tag{A5D}\label{ax:A5D}\\
\forall f\,\forall c\,
\bigl(\Cod(f,c)\rightarrow\Com(f,c,f)\bigr).
\tag{A5C}\label{ax:A5C}
\end{align}
Metalingüísticament, un cop recuperada la notació funcional, aquestes fórmules diran $f\comp d=f$ i $c\comp f=f$.

\section{Estabilitat dels extrems}
\label{sec:ar-estabilitat}

El primer fet rellevant és que qualsevol fletxa que apareix com a domini o codomini és estable sota les dues relacions d'extrem.

\begin{proposicio}[estabilitat dels extrems]
\label{prop:ar-estabilitat}
Si $\Dom(f,d)$, aleshores $\Dom(d,d)\land\Cod(d,d)$. Si $\Cod(f,c)$, aleshores $\Dom(c,c)\land\Cod(c,c)$.
\end{proposicio}

\begin{prova}
Suposem $\Dom(f,d)$. Per \eqref{ax:A5D}, $\Com(d,f,f)$. Per la direcció de dreta a esquerra de \eqref{ax:A2}, l'existència d'aquesta composició implica que existeix $e$ tal que $\Cod(d,e)\land\Dom(f,e)$. Com que també $\Dom(f,d)$, la unicitat del domini de $f$, donada per \eqref{ax:A1D}, implica $e=d$; per tant $\Cod(d,d)$.

Ara \eqref{ax:A5C}, aplicada a $d$, dona $\Com(d,d,d)$. Novament per \eqref{ax:A2}, existeix $e$ tal que $\Cod(d,e)\land\Dom(d,e)$. Ja sabem que $\Cod(d,d)$; per unicitat del codomini de $d$, $e=d$, i per tant $\Dom(d,d)$. Així, $\Dom(f,d)\longrightarrow \Dom(d,d)\land\Cod(d,d)$.

El cas del codomini és dual. Si $\Cod(f,c)$, llavors \eqref{ax:A5C} dona $\Com(f,c,f)$. Per \eqref{ax:A2} existeix $e$ amb $\Cod(f,e)\land\Dom(c,e)$; la unicitat del codomini de $f$ força $e=c$, de manera que $\Dom(c,c)$. Aleshores \eqref{ax:A5D} dona $\Com(c,c,c)$, i un nou ús de \eqref{ax:A2} i de la unicitat del domini de $c$ proporciona $\Cod(c,c)$.
\end{prova}

\begin{observacio}
La demostració només utilitza (A1), (A2) i (A5). No utilitza ni l'axioma dels extrems de la composta (A3) ni l'associativitat (A4).
\end{observacio}

\section{Unicitat de la composició}
\label{sec:ar-unicitat}

La relació $\Com$ no s'ha postulat com a funcional en el tercer argument. Aquesta funcionalitat es dedueix de la componibilitat, la neutralitat i l'associativitat.

\begin{proposicio}[funcionalitat de la composició]
\label{prop:ar-unicitat}
Per a totes les fletxes $f,g,h,k$,
\[
\Com(f,g,h)\land\Com(f,g,k)\longrightarrow h=k.
\]
\end{proposicio}

\begin{prova}
Suposem $\Com(f,g,h)$ i $\Com(f,g,k)$. Per \eqref{ax:A2}, existeix $e$ tal que $\Cod(f,e)\land\Dom(g,e)$. Per neutralitat, $\Com(f,e,f)$ i $\Com(e,g,g)$. Apliquem ara \eqref{ax:A4U} a la terna $f,e,g$. Prenent $u=f$, $v=g$, $p=h$, $q=k$, les quatre premisses de \eqref{ax:A4U} són precisament $\Com(f,e,f)$, $\Com(e,g,g)$, $\Com(f,g,h)$ i $\Com(f,g,k)$. Per tant, $h=k$. La demostració només fa servir (A2), (A5) i \eqref{ax:A4U}; en particular, no fa servir (A3).
\end{prova}

Així, $\Com$ és el graf d'una operació binària parcial: per a cada parella $(f,g)$ hi ha com a màxim una fletxa $h$ tal que $\Com(f,g,h)$, i (A2) determina exactament quan aquesta fletxa existeix.

\section{Identitats i objectes com a nocions derivades}
\label{sec:ar-identitats}

Definim, com a abreviació del llenguatge,
\begin{equation}
\Id(e)
\;:\Longleftrightarrow\;
\Dom(e,e)\land\Cod(e,e).
\label{def:ar-Id}
\end{equation}
El símbol $\Id$ no s'afegeix a la signatura primitiva: és només una abreviació definicional. La proposició d'estabilitat dona immediatament $\Dom(f,d)\rightarrow\Id(d)$ i $\Cod(f,c)\rightarrow\Id(c)$: tot domini i tot codomini és una fletxa identitària.

També podem introduir $\Obj(e):\Longleftrightarrow\Id(e)$. Així, un \emph{objecte} no és una segona mena d'individu: és una fletxa que té per domini i codomini ella mateixa. En particular, sota els axiomes anteriors,
\[
\Dom(e,e)
\;\Longleftrightarrow\;
\Cod(e,e)
\;\Longleftrightarrow\;
\Id(e),
\]
perquè qualsevol de les dues primeres propietats, aplicada a la proposició~\ref{prop:ar-estabilitat}, implica l'altra.

\section{El tercer axioma és derivable}
\label{sec:ar-redundancia}

La forma forta de l'associativitat permet demostrar l'axioma (A3) a partir de (A1), (A2), (A4) i (A5).

\begin{proposicio}[redundància de (A3)]
\label{prop:ar-redundancia}
Suposem (A1), (A2), (A4) i (A5). Si $\Com(f,g,h)$, $\Dom(f,d)$ i $\Cod(g,c)$, aleshores $\Dom(h,d)$ i $\Cod(h,c)$.
\end{proposicio}

\begin{prova}
Les proposicions~\ref{prop:ar-estabilitat} i~\ref{prop:ar-unicitat}, que farem servir, no depenen de (A3); l'argument no és, doncs, circular. Comencem pel domini. De $\Dom(f,d)$ i \eqref{ax:A5D} obtenim $\Com(d,f,f)$. Juntament amb $\Com(f,g,h)$, l'axioma d'existència associativa \eqref{ax:A4E}, aplicat a la terna $d,f,g$, proporciona $p,q$ tals que $\Com(f,g,p)$ i $\Com(d,h,q)$. Per \eqref{ax:A4U}, aquestes dues composicions iterades tenen el mateix resultat, és a dir, $p=q$. D'altra banda, per la unicitat de la composició (proposició~\ref{prop:ar-unicitat}), $\Com(f,g,h)$ i $\Com(f,g,p)$ impliquen $p=h$. Per tant $q=h$, i tenim $\Com(d,h,h)$. Per \eqref{ax:A2}, existeix $e$ amb $\Cod(d,e)\land\Dom(h,e)$. Com que $d$ és domini de $f$, la proposició~\ref{prop:ar-estabilitat} dona $\Cod(d,d)$. La unicitat del codomini de $d$ implica $e=d$, i obtenim $\Dom(h,d)$.

Per al codomini, de $\Cod(g,c)$ i \eqref{ax:A5C} obtenim $\Com(g,c,g)$. Aplicada a $f,g,c$, l'associativitat forta proporciona composicions iterades $\Com(h,c,p)$ i $\Com(f,g,q)$ amb $p=q$. La unicitat de la composició de $f$ amb $g$ dona $q=h$, i per tant $\Com(h,c,h)$. Per \eqref{ax:A2}, existeix $e$ amb $\Cod(h,e)\land\Dom(c,e)$. Com que $c$ és codomini de $g$, l'estabilitat dona $\Dom(c,c)$. Per unicitat del domini de $c$, $e=c$, i així $\Cod(h,c)$.
\end{prova}

\begin{corollari}
\label{cor:ar-equivalencia-sistemes}
Amb l'associativitat forta adoptada aquí, el sistema de cinc blocs $(A1),(A2),(A3),(A4),(A5)$ és equivalent al sistema reduït $(A1),(A2),(A4),(A5)$. El tercer axioma original es pot conservar en una primera presentació per la seva transparència conceptual, però no és necessari com a axioma independent.
\end{corollari}

No afirmem que els quatre blocs restants siguin independents entre si: la redundància de (A3) no estableix la minimalitat del sistema reduït. L'anàlisi d'independència de sistemes funcionals afins ha estat estudiada i formalitzada per Benzmüller i Scott~\cite{BS2020,BSAFP}, però la transferència literal d'aquells resultats al llenguatge purament relacional adoptat aquí requeriria una comparació formal addicional; una via natural seria construir models finits que separessin cadascun dels quatre blocs.

\begin{observacio}[la forma lògica de (A4) és essencial]
Aquesta reducció depèn de la forma lògica de (A4). Si l'associativitat es formula només com la igualtat de dues composicions iterades \emph{suposant que totes dues ja existeixen}, no es controla l'existència de les composicions iterades i la demostració anterior deixa de funcionar. En una teoria de composició parcial, la forma lògica de l'associativitat és, doncs, part essencial de l'axiomàtica.
\end{observacio}

\section{Recuperació de la notació funcional}
\label{sec:ar-notacio-funcional}

L'axioma (A1) permet introduir en el metallenguatge les notacions $d(f)$ i $c(f)$ per a les úniques fletxes que satisfan, respectivament, $\Dom(f,d(f))$ i $\Cod(f,c(f))$. Aquestes expressions no són termes del llenguatge primitiu $\mathcal L_{\mathrm{Cat}}$: són notació derivada, legitimada per un teorema d'existència i unicitat; en termes lògics, passem a una extensió definicional per descripcions úniques. De la mateixa manera, si $f$ i $g$ són componibles, (A2) dona l'existència d'una $h$ tal que $\Com(f,g,h)$, i la proposició~\ref{prop:ar-unicitat} en dona la unicitat; escrivim, doncs, $g\comp f$ per designar aquesta única fletxa.

Amb aquesta notació, (A2) es llegeix
\[
g\comp f\ \text{existeix}
\quad\Longleftrightarrow\quad
c(f)=d(g),
\]
i la proposició~\ref{prop:ar-redundancia} dona $d(g\comp f)=d(f)$ i $c(g\comp f)=c(g)$.

Si $A$ i $B$ són fletxes identitàries, la notació usual $f\colon A\longrightarrow B$ abrevia metalingüísticament $\Dom(f,A)\land\Cod(f,B)$; les condicions $\Id(A)$ i $\Id(B)$ són automàtiques per estabilitat. També podem escriure, metalingüísticament, $\Hom(A,B)=\{f\in M: \Dom(f,A)\land\Cod(f,B)\}$. El símbol $\Hom$ no forma part tampoc del llenguatge primitiu.

\section{Equivalència amb la definició usual}
\label{sec:ar-equivalencia}

La presentació relacional no defineix una estructura matemàtica nova: recupera la noció usual de categoria petita (definició~\ref{def:categoria}), fins a l'isomorfisme canònic que identifica cada objecte amb la seva identitat.

Abans de les demostracions, val la pena desplegar un model concret i petit, per veure com les relacions $\Dom$, $\Cod$ i $\Com$ codifiquen una categoria familiar i com els objectes hi apareixen com a fletxes identitat.

\begin{exemple}[Un model amb tres objectes i una cadena de morfismes]\label{ex:ar-model-3}
Considerem la categoria $\cat{3}$ (amb els objectes i les fletxes reanomenats), formada per una cadena de dues fletxes, $A\xrightarrow{u}B\xrightarrow{v}C$, amb la seva composta $w=v\comp u\colon A\to C$:
\[
\begin{tikzpicture}[baseline=(A.base)]
  \node (A) at (0,0) {$A$};
  \node (B) at (2.2,0) {$B$};
  \node (C) at (4.4,0) {$C$};
  \draw[-{Stealth[scale=1.1]}] (A) -- node[above] {$u$} (B);
  \draw[-{Stealth[scale=1.1]}] (B) -- node[above] {$v$} (C);
  \draw[-{Stealth[scale=1.1]}] (A) to[out=-35,in=-145] node[below] {$w=v\comp u$} (C);
\end{tikzpicture}
\]
El model relacional corresponent té per univers el conjunt de \emph{totes} les fletxes, incloses les tres identitats:
\[
M=\{\,A,\ B,\ C,\ u,\ v,\ w\,\},
\]
on $A,B,C$ són alhora els objectes i les seves fletxes identitat. Les relacions unàries derivades i els extrems de cada fletxa són:
\[
\renewcommand{\arraystretch}{1.2}
\begin{array}{c|c|c|c}
f & \Dom(f,-) & \Cod(f,-) & \Id(f) \\ \hline
A & A & A & \text{sí} \\
B & B & B & \text{sí} \\
C & C & C & \text{sí} \\
u & A & B & \text{no} \\
v & B & C & \text{no} \\
w & A & C & \text{no}
\end{array}
\]
Els objectes són exactament les fletxes $e$ amb $\Id(e)$, és a dir $A,B,C$: no calen com a dada primitiva, sinó que es reconeixen per la propietat $\Dom(e,e)\land\Cod(e,e)$. Les úniques parelles componibles $(f,g)$ ---les que compleixen $c(f)=d(g)$--- donen la taula de $\Com$, on recordem que $\Com(f,g,h)$ es llegeix \enquote{primer $f$, després $g$, amb composta $h=g\comp f$}:
\[
\renewcommand{\arraystretch}{1.2}
\begin{array}{ll}
\Com(A,A,A) & \text{(identitat de }A) \\
\Com(B,B,B) & \text{(identitat de }B) \\
\Com(C,C,C) & \text{(identitat de }C) \\
\Com(A,u,u),\ \Com(u,B,u) & (u\comp\id_A=u=\id_B\comp u) \\
\Com(B,v,v),\ \Com(v,C,v) & (v\comp\id_B=v=\id_C\comp v) \\
\Com(A,w,w),\ \Com(w,C,w) & (w\comp\id_A=w=\id_C\comp w) \\
\Com(u,v,w) & (\text{la composició no trivial }v\comp u=w)
\end{array}
\]
Val la pena llegir l'última línia amb l'ordre de $\Com$ ben present: $\Com(u,v,w)$ diu que en recórrer \emph{primer} $u$ i \emph{després} $v$ s'obté $w$, és a dir $w=v\comp u$. No hi ha cap altra parella componible; en particular, $u$ i $w$ no es componen en cap ordre, perquè $c(u)=B\neq A=d(w)$ i $c(w)=C\neq A=d(u)$. Comprovem els axiomes. (A1) és la taula d'extrems. (A2) es comprova parella per parella: les deu ternes de la taula de $\Com$ corresponen exactament a les deu parelles $(f,g)$ amb $c(f)=d(g)$. (A5) és el contingut de les sis primeres línies, que donen $\Com(d(f),f,f)$ i $\Com(f,c(f),f)$ per a cada $f$. Per a (A4), observem que una terna encadenada $(f,g,h)$, amb $\Com(f,g,u)$ i $\Com(g,h,v)$, conté sempre alguna identitat, perquè la cadena més llarga de fletxes no identitat és $u,v$, de longitud dos. Si la identitat és $f$, aleshores $u=g$ i $f$ és el domini de $v$, de manera que totes dues associacions donen $v$: $\Com(u,h,v)$ i $\Com(f,v,v)$. Els casos en què la identitat és $g$ o $h$ són anàlegs, i totes dues associacions donen la composta de les altres dues fletxes. Així es compleixen \eqref{ax:A4E} i \eqref{ax:A4U}. Reconstruint la categoria a partir d'aquest model (com farà el teorema~\ref{th:ar-model-a-categoria}) es recupera la categoria de partida, en què ara cada objecte és la seva pròpia fletxa identitat.
\end{exemple}

\begin{teorema}
\label{th:ar-model-a-categoria}
Tot model del sistema reduït $(A1),(A2),(A4),(A5)$ amb univers no buit determina una categoria petita no buida en el sentit de la definició~\ref{def:categoria}.
\end{teorema}

\begin{prova}
Sigui $\mathfrak M$ un model. Prenem com a fletxes els elements de $M$ i com a objectes $\Obj(\mathfrak M)=\{e\in M:\Id(e)\}$. Per (A1), cada fletxa $f$ té un domini únic $d(f)$ i un codomini únic $c(f)$; per estabilitat, tots dos són objectes. Si $c(f)=d(g)$, (A2) proporciona una composta i la proposició~\ref{prop:ar-unicitat} en proporciona una de sola, que denotem $g\comp f$; la proposició~\ref{prop:ar-redundancia} assegura que $d(g\comp f)=d(f)$ i $c(g\comp f)=c(g)$. Per a un objecte $A$, la seva fletxa identitat és $A$ mateix. Aquesta fletxa va de $A$ a $A$, perquè $\Dom(A,A)$ i $\Cod(A,A)$. Si $f\colon A\to B$, (A5) dona $\Com(A,f,f)$ i $\Com(f,B,f)$, és a dir, $f\comp\id_A=f=\id_B\comp f$. Per a l'associativitat, siguin $f\colon A\to B$, $g\colon B\to C$ i $h\colon C\to D$. Aplicant \eqref{ax:A4E} a $\Com(f,g,g\comp f)$ i $\Com(g,h,h\comp g)$, existeixen $p,q$ amb $\Com(g\comp f,h,p)$ i $\Com(f,h\comp g,q)$; per la unicitat de la composició, $p=h\comp(g\comp f)$ i $q=(h\comp g)\comp f$, i \eqref{ax:A4U} dona $p=q$. Així obtenim una categoria, petita perquè $M$ és un conjunt.
\end{prova}

\begin{teorema}[recíproc]
\label{th:ar-categoria-a-model}
Tota categoria petita no buida determina un model del sistema relacional.
\end{teorema}

\begin{prova}
Sigui $\cat{C}$ una categoria petita no buida. Prenem com a univers $M=\operatorname{Mor}(\cat{C})$. Si $f\colon A\to B$, definim
\[
\Dom(f,d)\Longleftrightarrow d=\id_A,
\qquad
\Cod(f,c)\Longleftrightarrow c=\id_B,
\]
i
i, per a $f\colon A\to B$ i $g\colon B'\to C$,
\[
\Com(f,g,h)\Longleftrightarrow \bigl(B=B'\ \text{i}\ h=g\comp f\bigr),
\]
on la notació de la dreta pertany ara al metallenguatge de la categoria usual. L'univers és no buit, perquè $\cat{C}$ té algun objecte i, per tant, alguna identitat. Comprovem els axiomes. (A1) és immediat. Per a (A2), $\Cod(f,e)\land\Dom(g,e)$ vol dir $e=\id_B=\id_{B'}$, i això passa si i només si $B=B'$, perquè $\id_B$ determina $B$ com a domini; i $B=B'$ és exactament la condició perquè existeixi $h$ amb $\Com(f,g,h)$. (A3) diu que $g\comp f\colon A\to C$. Per a (A4), $\Com(f,g,u)$ i $\Com(g,h,v)$ volen dir que $f,g,h$ són consecutives, amb $u=g\comp f$ i $v=h\comp g$; aleshores existeixen $h\comp u$ i $v\comp f$, que són iguals per l'associativitat de $\cat{C}$. Finalment, (A5) són les lleis d'identitat $f\comp\id_A=f$ i $\id_B\comp f=f$.
\end{prova}

\begin{observacio}
En passar d'una categoria usual $\cat{C}$ al model relacional, els individus són els morfismes de $\cat{C}$; en reconstruir després una categoria a partir del model, els seus objectes són les fletxes identitat $\id_A$, no els objectes originals $A$. La categoria reconstruïda és canònicament isomorfa a $\cat{C}$ ---la bijecció d'objectes és $A\mapsto\id_A$--- però en general no n'és literalment la mateixa estructura conjuntista. Per aquest motiu diem que les dues presentacions descriuen la mateixa estructura categòrica \emph{fins a isomorfisme canònic}, no que en siguin idèntiques. Una afirmació metateòrica més forta ---una equivalència definicional o una bi-interpretabilitat--- requeriria definir explícitament les traduccions entre les dues teories i comparar-les; no la formulem aquí.

\end{observacio}

Així, la diferència entre les dues presentacions és una diferència de \emph{llenguatge primitiu}, no de contingut categòric.

\section{Relacional i funcional}
\label{sec:ar-relacional-funcional}

Els predicats $\Dom$ i $\Cod$, sota (A1), són els grafs de dues funcions unàries totals. La relació $\Com$, sota (A2) i la proposició~\ref{prop:ar-unicitat}, és el graf d'una funció binària parcial. Aquesta observació explica dues maneres naturals de formalitzar les categories.

\paragraph{Presentació relacional.}
En lògica de primer ordre estàndard podem prendre com a primitives $\Dom(f,d)$, $\Cod(f,c)$ i $\Com(f,g,h)$. L'existència s'expressa mitjançant quantificadors i la funcionalitat s'axiomatitza o es demostra. És la via d'aquest apèndix.

\paragraph{Presentació funcional.}
Podem prendre $\dom$ i $\cod$ com a operadors unaris. El problema és la composició: com que no tota parella de fletxes és componible, una operació binària de composició no és total sobre $M\times M$. En lògica de primer ordre estàndard, en canvi, els símbols de funció s'interpreten com a funcions totals. Cal, doncs, ampliar el marc lògic o codificar explícitament la parcialitat.

Aquest és el camí seguit, en un formalisme diferent, per Benzmüller i Scott~\cite{BS2020,BSAFP}. Treballen amb una lògica lliure embeguda en Isabelle/HOL, un predicat $E$ d'existència o definició, i tres operacions parcials $\dom$, $\cod$ i composició. En el seu \emph{AxiomsSet5}, atribuït a Scott, els axiomes principals tenen la forma
\begin{align*}
E(\dom x)&\to E(x),\tag{S1}\\
E(\cod y)&\to E(y),\tag{S2}\\
E(x\cdot y)&\leftrightarrow \dom x\approx\cod y,\tag{S3}\\
x\cdot(y\cdot z)&\simeq (x\cdot y)\cdot z,\tag{S4}\\
(\cod y)\cdot y&\simeq y,\tag{S5}\\
x\cdot(\dom x)&\simeq x,\tag{S6}
\end{align*}
amb la composició escrita en ordre funcional ($x\cdot y$ és \enquote{primer $y$, després $x$}), on $\simeq$ és la \emph{igualtat de Kleene} (dos termes són iguals si tots dos són indefinits, o bé tots dos definits i iguals) i $\approx$ és la \emph{identitat existent} (tots dos termes definits i iguals). En aquesta presentació, la funcionalitat de $\dom$, $\cod$ i de la composició queda incorporada en els símbols funcionals; la parcialitat es gestiona mitjançant la lògica lliure. Els mateixos autors formalitzen també equivalències entre diversos sistemes axiomàtics i el pas, per Skolemització, de formulacions existencials cap a formulacions funcionals.

\begin{nota}
Els símbols tipogràfics per a aquestes dues igualtats varien entre fonts; aquí fixem $\simeq$ per a la igualtat de Kleene i $\approx$ per a la identitat existent, que no coincideixen necessàriament amb la notació de l'obra citada.
\end{nota}

La correspondència conceptual amb la presentació relacional és:
\[
\begin{array}{c|c}
\text{presentació relacional} & \text{presentació funcional lliure}\\
\hline
\text{(A1): domini i codomini} & \dom,\cod\ \text{i } S1,S2\\
\text{(A2): componibilitat} & S3\\
\text{(A3): extrems de la composta} & \text{derivable en sistemes reduïts}\\
\text{(A4): associativitat} & S4\\
\text{(A5): neutralitat} & S5,S6
\end{array}
\]
La correspondència no és literal fórmula per fórmula, perquè els dos llenguatges distribueixen la informació de manera diferent.

\section{El principi de dualitat}
\label{sec:ar-dualitat}\index{dualitat!principi de}

El capítol~\ref{cap:categories} formula el principi de dualitat de manera operativa (teorema~\ref{teo:dualitat}): una afirmació expressable en el llenguatge de les categories es dualitza invertint les fletxes, i, si és vàlida en totes les categories, la seva dual també ho és, perquè l'afirmació original es pot aplicar a la categoria oposada. La presentació relacional permet donar-ne una versió completament formal, en dues formes: una de \emph{sintàctica}, sobre deduccions, i una de \emph{semàntica}, sobre models. En aquesta secció escrivim $\mathrm{CT}$ per al conjunt d'axiomes \eqref{ax:A1D}--\eqref{ax:A5C}.

\paragraph{La traducció dual.}
Per a cada fórmula $\varphi$ de $\mathcal L_{\mathrm{Cat}}$ definim la seva \textbf{dual}\index{dualitat!d'un enunciat} $\varphi\op$ substituint cada fórmula atòmica no lògica segons
\begin{gather*}
\Dom(f,d)\longmapsto\Cod(f,d),\qquad
\Cod(f,c)\longmapsto\Dom(f,c),\\
\Com(f,g,h)\longmapsto\Com(g,f,h),
\end{gather*}
i deixant intactes les igualtats, els connectors i els quantificadors. Com que $\Com(f,g,h)$ es llegeix \enquote{primer $f$, després $g$, amb composta $h$}, intercanviar-ne els dos primers arguments és exactament invertir l'ordre de composició. La traducció és una involució, $(\varphi\op)\op=\varphi$, i commuta amb la substitució de variables. Un cop recuperada la notació funcional (secció~\ref{sec:ar-notacio-funcional}), aquestes substitucions són les del capítol~\ref{cap:categories}: $d(f)$ i $c(f)$ s'intercanvien i $g\comp f$ passa a $f\comp g$. Observem que la dualitat no depèn de tractar la composició com un símbol funcional total: opera directament sobre la relació ternària.

\begin{observacio}
Els connectors i quantificadors de $\varphi$ pertanyen al llenguatge $\mathcal L_{\mathrm{Cat}}$, amb què \emph{parlem} d'una categoria, i la dualitat els deixa intactes. No s'han de confondre amb la lògica que una categoria pot codificar \emph{internament}, quan els seus objectes són fórmules i els seus morfismes deduccions. En aquest cas, la dualitat de fletxes pot intercanviar la conjunció i la disjunció del sistema codificat, però el $\land$ i el $\lor$ de $\mathcal L_{\mathrm{Cat}}$ es mantenen.
\end{observacio}

\paragraph{Autodualitat dels axiomes.}

\begin{lema}\label{lema:ar-autodual}
Per a cada axioma $\alpha$ de $\mathrm{CT}$, la fórmula $\alpha\op$ és lògicament equivalent a un axioma de $\mathrm{CT}$. Concretament, \eqref{ax:A1D} i \eqref{ax:A1C} es transformen l'un en l'altre, igual que \eqref{ax:A5D} i \eqref{ax:A5C}; i cadascun dels axiomes \eqref{ax:A2}, \eqref{ax:A3}, \eqref{ax:A4E} i \eqref{ax:A4U} es transforma en una variant d'ell mateix, que només en difereix pel nom de les variables lligades, per l'ordre dels quantificadors i dels membres de les conjuncions i, en el cas de \eqref{ax:A4U}, per la simetria de la igualtat.
\end{lema}

\begin{prova}
Els casos d'\eqref{ax:A1D} i \eqref{ax:A5D} són immediats: $(\forall f\,\exists!d\,\Dom(f,d))\op$ és $\forall f\,\exists!d\,\Cod(f,d)$, que és \eqref{ax:A1C}, i $(\forall f\,\forall d\,(\Dom(f,d)\to\Com(d,f,f)))\op$ és $\forall f\,\forall d\,(\Cod(f,d)\to\Com(f,d,f))$, que és \eqref{ax:A5C}; per involució, els casos d'\eqref{ax:A1C} i \eqref{ax:A5C} són els recíprocs.

La dual d'\eqref{ax:A2} és
\[
\forall f\,\forall g\,\bigl[\exists e\bigl(\Dom(f,e)\land\Cod(g,e)\bigr)\leftrightarrow\exists h\,\Com(g,f,h)\bigr],
\]
i, intercanviant els noms de les variables lligades $f$ i $g$, s'obté \eqref{ax:A2} llevat de l'ordre dels quantificadors i dels membres de la conjunció. La dual d'\eqref{ax:A3} és
\[
\forall f,g,h,d,c\,\bigl[\Com(g,f,h)\land\Cod(f,d)\land\Dom(g,c)\longrightarrow\Cod(h,d)\land\Dom(h,c)\bigr],
\]
i l'intercanvi $f\leftrightarrow g$, $d\leftrightarrow c$ la torna en \eqref{ax:A3}. La dual d'\eqref{ax:A4E} és
\[
\forall f,g,h,u,v\,\bigl[\Com(g,f,u)\land\Com(h,g,v)\longrightarrow\exists p\,\exists q\,\bigl(\Com(h,u,p)\land\Com(v,f,q)\bigr)\bigr];
\]
amb l'intercanvi $f\leftrightarrow h$, $u\leftrightarrow v$, $p\leftrightarrow q$ esdevé
\[
\forall f,g,h,u,v\,\bigl[\Com(g,h,v)\land\Com(f,g,u)\longrightarrow\exists q\,\exists p\,\bigl(\Com(f,v,q)\land\Com(u,h,p)\bigr)\bigr],
\]
que és \eqref{ax:A4E}. El mateix reanomenament transforma la dual d'\eqref{ax:A4U} en \eqref{ax:A4U}, amb la conclusió escrita $q=p$.
\end{prova}

Intuïtivament, el lema diu que la definició de categoria no distingeix entre una fletxa i la mateixa fletxa recorreguda en sentit contrari. És l'únic ingredient específic de la teoria que necessiten els dos teoremes següents.

\paragraph{Dualitat sintàctica.}

\begin{teorema}[Dualitat sintàctica]\label{teo:ar-dualitat-sint}
Per a tota sentència $\varphi$ de $\mathcal L_{\mathrm{Cat}}$, si $\mathrm{CT}\vdash\varphi$, aleshores $\mathrm{CT}\vdash\varphi\op$.
\end{teorema}

\begin{prova}
Fixem un càlcul de Hilbert estàndard per a la lògica de primer ordre amb igualtat. La traducció dual només reanomena les relacions primitives i permuta els arguments de $\Com$; commuta amb la substitució, conserva les variables lliures i no toca ni els connectors, ni els quantificadors, ni la igualtat. Per tant, transforma cada axioma lògic (inclosos els de la igualtat) en un axioma lògic, i cada aplicació d'una regla d'inferència en una aplicació de la mateixa regla; les restriccions sobre variables de la regla de generalització es conserven, perquè les variables lliures no canvien. Aplicada fórmula a fórmula a una deducció de $\varphi$ a partir de $\mathrm{CT}$, en dona una de $\varphi\op$ a partir de $\mathrm{CT}\op=\{\alpha\op:\alpha\in\mathrm{CT}\}$. Pel lema~\ref{lema:ar-autodual}, cada $\alpha\op$ es dedueix de $\mathrm{CT}$; per tant, $\mathrm{CT}\vdash\varphi\op$.
\end{prova}

\paragraph{Dualitat semàntica.}
Per a una estructura $\mathfrak M=(M,\Dom^{\mathfrak M},\Cod^{\mathfrak M},\Com^{\mathfrak M})$ definim l'\textbf{estructura oposada}
\[
\mathfrak M\op=\bigl(M,\ \Cod^{\mathfrak M},\ \Dom^{\mathfrak M},\ \{(g,f,h):(f,g,h)\in\Com^{\mathfrak M}\}\bigr),
\]
amb el mateix univers, els papers de domini i codomini intercanviats i la composició en l'ordre invers. Si $\mathfrak M$ és el model associat a una categoria petita $\cat{C}$ pel teorema~\ref{th:ar-categoria-a-model}, aleshores l'aplicació $f\mapsto f\op$ és un isomorfisme de $\mathfrak M\op$ en el model associat a $\cat{C}\op$: per a $f\colon A\to B$, la relació $\Dom^{\mathfrak M\op}(f,\id_B)$, que és $\Cod^{\mathfrak M}(f,\id_B)$, correspon a $\Dom(f\op,(\id_B)\op)$ a $\cat{C}\op$, i $\Com^{\mathfrak M\op}(f,g,h)$, que és $\Com^{\mathfrak M}(g,f,h)$, és a dir $h=f\comp g$, correspon a $h\op=g\op\comp f\op$, és a dir, a $\Com(f\op,g\op,h\op)$ a $\cat{C}\op$. Si s'identifica cada $f\op$ amb $f$, els dos models coincideixen.

\begin{lema}\label{lema:ar-oposada}
Per a tota fórmula $\varphi$ de $\mathcal L_{\mathrm{Cat}}$ i tota assignació $s$ de les seves variables lliures,
\[
\mathfrak M\op\models\varphi[s]
\quad\Longleftrightarrow\quad
\mathfrak M\models\varphi\op[s].
\]
\end{lema}

\begin{prova}
Per inducció sobre $\varphi$. Per a les fórmules atòmiques és la definició de $\mathfrak M\op$: per exemple, $\mathfrak M\op\models\Com(f,g,h)[s]$ si i només si $(s(g),s(f),s(h))\in\Com^{\mathfrak M}$, és a dir, si i només si $\mathfrak M\models\Com(g,f,h)[s]$. Les igualtats s'interpreten igual a totes dues estructures, i els passos dels connectors i dels quantificadors són immediats perquè l'univers és el mateix.
\end{prova}

\begin{teorema}[Dualitat semàntica]\label{teo:ar-dualitat-sem}
Si $\mathfrak M$ és un model de $\mathrm{CT}$, també ho és $\mathfrak M\op$. En conseqüència, per a tota sentència $\varphi$ de $\mathcal L_{\mathrm{Cat}}$, si $\varphi$ es compleix en tots els models de $\mathrm{CT}$ ---que, pels teoremes~\ref{th:ar-model-a-categoria} i~\ref{th:ar-categoria-a-model}, corresponen a les categories petites no buides---, també s'hi compleix $\varphi\op$.
\end{teorema}

\begin{prova}
Si $\mathfrak M\models\mathrm{CT}$ i $\alpha\in\mathrm{CT}$, pel lema~\ref{lema:ar-autodual} la fórmula $\alpha\op$ és equivalent a un axioma, de manera que $\mathfrak M\models\alpha\op$, i pel lema~\ref{lema:ar-oposada}, $\mathfrak M\op\models\alpha$. Per a la segona afirmació, sigui $\mathfrak M$ un model de $\mathrm{CT}$. Com que $\mathfrak M\op$ també ho és, $\mathfrak M\op\models\varphi$, i el lema~\ref{lema:ar-oposada} dona $\mathfrak M\models\varphi\op$.
\end{prova}

Aquesta demostració és la versió formal de la del capítol~\ref{cap:categories}: allà es fa servir que $\cat{C}\op$ és una categoria; aquí, que $\mathfrak M\op$ és un model. Pel teorema de completesa, $\mathrm{CT}\vdash\varphi$ equival a $\mathrm{CT}\models\varphi$, i els dos teoremes són, doncs, dues cares d'un mateix fet. Les restriccions sobre la mida de l'observació~\ref{obs:ar-metateorica} es mantenen: la versió semàntica parla de categories petites no buides, i l'extensió a categories grans demana la mateixa convenció metateòrica addicional.

\paragraph{Abast.}
Tots dos teoremes parlen de fórmules de $\mathcal L_{\mathrm{Cat}}$. Una afirmació que mencioni elements d'un conjunt, una categoria concreta o qualsevol estructura externa no és una fórmula d'aquest llenguatge i no té dual fins que no s'hi reformula. Per això el capítol~\ref{cap:categories} enuncia el principi per a afirmacions \emph{expressables en el llenguatge de les categories}, i no per a \enquote{qualsevol afirmació sobre categories}. Els enunciats que involucren diverses categories i functors entre elles demanen un llenguatge més ric; el mecanisme és el mateix, però cal dualitzar també les categories implicades, com mostra l'exemple de les categories sobre i sota un objecte.

\paragraph{Un mecanisme general.}
El que hem fet ---una involució del llenguatge que deixa invariant el conjunt d'axiomes--- no és exclusiu de la teoria de categories. La lògica clàssica en dona una altra instància: en el llenguatge de les àlgebres booleanes, la involució que intercanvia $\wedge$ amb $\vee$ i el cert amb el fals deixa els axiomes invariants, i produeix la dualitat de De~Morgan. Un exemple encara més antic és la \textbf{dualitat projectiva}\index{dualitat!projectiva}: en el pla projectiu, intercanviar \emph{punt} per \emph{recta} ---i \enquote{passar per} per \enquote{tallar-se en}--- deixa els axiomes invariants, de manera que cada teorema té un teorema dual, com el de Desargues o la parella Pascal--Brianchon. En tots els casos actua el mateix mecanisme, aplicat a llenguatges diferents. Quan una categoria codifica lògica, la dualitat de fletxes intercanvia conjunció i disjunció si existeixen; que aquesta dualitat coincideixi exactament amb la de De~Morgan requereix, a més, un complement involutiu, i és una hipòtesi addicional sobre la categoria, no una conseqüència automàtica de la dualitat categòrica.

\section{Una observació computacional}
\label{sec:ar-computacional}

La presentació relacional és especialment transparent per a l'axiomàtica, però no és necessàriament la més còmoda per a una implementació. Un programa prefereix, habitualment, operacions que retornin valors, $\dom\colon M\to M$ i $\cod\colon M\to M$, i una composició parcial. Aquesta darrera es pot representar, per exemple, com $\operatorname{comp}\colon M\times M\longrightarrow M_{\bot}$ o mitjançant un tipus opcional que distingeixi entre una composta existent i una parella no componible.

No hi ha contradicció entre les dues perspectives. La teoria relacional demostra primer existència i unicitat; aquests resultats justifiquen després el pas a una interfície funcional. El millor llenguatge per axiomatitzar una estructura no ha de coincidir amb el millor llenguatge per implementar-la.

\section{Connexió amb una lectura estructural}
\label{sec:ar-estructural}

La presentació anterior respon a un criteri estructural molt simple: no introduir com a primitives entitats que poden reconstruir-se a partir de les relacions fonamentals de la teoria. En la presentació usual es parla de dos tipus de dades, objectes i fletxes; en la relacional elemental només hi ha una mena d'individus, i els objectes apareixen com les fletxes que satisfan $\Dom(e,e)\land\Cod(e,e)$. L'objecte no és, doncs, una entitat afegida al costat de les fletxes, sinó una posició estructural determinada per les relacions de domini i codomini ---exactament el gir que la introducció anunciava en preferir \emph{relacionar} a \emph{individualitzar} (secció~\ref{sec:mirar-fora}).

Aquesta perspectiva és compatible amb una lectura estructural de la teoria de conjunts: les realitzacions conjuntistes concretes poden quedar al metallenguatge, mentre que el llenguatge de la teoria categòrica reté només les relacions que expressen l'estructura que volem estudiar. Això no elimina la teoria de conjunts del metallenguatge ---la semàntica estàndard de primer ordre continua sent conjuntista, i el marc de mida el fixa l'apèndix corresponent--- però evita confondre una realització externa amb les nocions primitives de la teoria objecte.

La lliçó metodològica és, doncs, doble. D'una banda, la formulació relacional mostra que una categoria es pot reconstruir a partir d'un univers de fletxes i de tres relacions primitives. De l'altra, un cop feta la reconstrucció, no hi ha cap inconvenient a recuperar la notació funcional i la terminologia usual, que són molt més eficients per desenvolupar la teoria ---i que són, de fet, les que la resta d'aquest llibre empra.

\chapter{Grafs}\label{cap:grafs}

Aquest apèndix desenvolupa amb detall la noció de graf dirigit, els seus homomorfismes i els camins, i mostra com generar una categoria a partir d'un graf. El material amplia, amb molts exemples concrets, la subsecció~\ref{subsec:grafs} del capítol de categories.

\section{Grafs}

Qualsevol sistema de relacions, com les rutes de transport o les connexions d'aerolínies, es pot dibuixar com una col·lecció de punts i fletxes. En aquesta etapa no ens importa què són els punts, sinó com estan connectats. Aquesta idea es modela amb el concepte de graf. En general, un graf consta de punts (o objectes), anomenats \emph{vèrtexs}, alguns dels quals estan units per \emph{arestes}; aquí considerarem que els grafs són \emph{dirigits}, és a dir, que cada aresta apunta en una direcció, d'un punt a un altre, i per això fem servir fletxes per connectar punts. Entre dos punts hi pot haver moltes fletxes, o cap, i fins i tot pot haver-hi fletxes d'un punt a si mateix. Sovint la manera més fàcil de descriure un graf petit és dibuixar-lo, com a la figura següent.
\begin{equation}\label{graf:1}
\xymatrix{
1 \ar@(ul,ur)[]^a \ar@/^/[dr]^b
& & 2 \ar@/^/[dl]_c \ar@/_/[dr]^d \ar[ll]_e
& & 3 \ar[ll]_f \ar@(ur,dr)[]^g \\
& 4 \ar@/^/[ul]^h \ar[rr]^k
& & 5 \ar[ur]_l & 6
}
\end{equation}
El tipus de graf que tractem aquí és una versió del graf dirigit que sovint s'anomena \emph{multigraf dirigit amb bucles}, encara que farem servir la paraula \enquote{graf} per abreujar. En donem la definició formal.

\begin{definicio}\index{graf dirigit}
Un \textbf{graf} $G$ comprèn:
\begin{itemize}
\item un conjunt $V_G$ de \textbf{vèrtexs};
\item un conjunt $E_G$ d'\textbf{arestes};
\item dues aplicacions $s_G\colon E_G\to V_G$ i $t_G\colon E_G\to V_G$, anomenades \textbf{origen} i \textbf{destí}. Si $s_G(f)=a$ i $t_G(f)=b$, escrivim $a\xrightarrow{f}b$ o $f\colon a\to b$, i diem que $a$ és el vèrtex d'origen i $b$ el de destí.
\end{itemize}
La notació
\[
G=\xymatrix{
E_G \ar@<0.6ex>[r]^{s_G} \ar@<-0.6ex>[r]_{t_G} & V_G
}
\]
indica que $G$ és un graf amb vèrtexs a $V_G$ i arestes a $E_G$.
\end{definicio}

\begin{exemple}
Considerem el graf $G$ de la figura següent:
\[
\xymatrix{
	& \bullet_{a} \ar@<1ex>[dl]^f \ar@<1ex>[dr]^g &  \\
	\bullet_{b} \ar@(ul,dl)_l \ar@<1ex>[ur]^h & & \bullet_{c}  \ar@<1ex>[ll]^k \ar@(ur,dr)^m \\
}
\]
Especifica quins són els seus elements.
\end{exemple}
\begin{solucio}
Es té $V_G=\{a,b,c\}$ i $E_G=\{f,g,h,k,l,m\}$, i les dues aplicacions són
\[
    \begin{array}{c|c|c}
        E_G & s_G & t_G \\ \hline
        f & a & b  \\
        g & a & c \\
        h & b & a \\
        k & c & b \\
        l & b & b \\
        m & c & c
    \end{array}
\]
\end{solucio}

\begin{exemple}
Defineix el graf $G$ perquè la seva representació gràfica sigui \eqref{graf:1}.
\end{exemple}
\begin{solucio}
Prenem $V_G=\{1,2,3,4,5,6\}$, $E_G=\{a,b,c,d,e,f,g,h,k,l\}$ i les aplicacions
\[
    \begin{array}{c|c|c}
        E_G & s_G & t_G \\ \hline
        a & 1 & 1  \\
        b & 1 & 4 \\
        c & 2 & 4 \\
        d & 2 & 5 \\
        e & 2 & 1 \\
        f & 3 & 2 \\
        g & 3 & 3 \\
        h & 4 & 1 \\
        k & 4 & 5 \\
        l & 5 & 3
    \end{array}
\]
que reprodueixen la figura \eqref{graf:1}.
\end{solucio}

\begin{exemple}\leavevmode
\begin{enumerate}
\item Dibuixa el graf corresponent a les dades següents:
\[
    \begin{array}{c|c|c}
        E_G & s_G & t_G \\ \hline
        f & 2 & 3 \\
        g & 2 & 3 \\
        h & 2 & 3 \\
        k & 4 & 3 \\
        l & 6 & 3 \\
        m & 6 & 6
    \end{array}
    \hspace{2cm} V_G=\{1,2,3,4,5,6\}
\]
\item Escriu la taula corresponent al graf $H$ següent:
\[
\xymatrix{
	1 \ar[r]^{a} & 2 \ar[d]_{g} \ar@/^2ex/[r]^-{b} \ar[r]^-{c} & 3 \ar[r]^-{e} \ar@/^2ex/[l]^-{d} & 4 \\
	5 \ar@(ul,dl)_l & 6 \ar[l]_-{h} \ar[r]^-{k} & 7 \ar[ur]^-{f}
}
\]
\end{enumerate}
\end{exemple}
\begin{solucio}\leavevmode
\begin{enumerate}
\item
\[
\xymatrix{
	4 \ar[r]^-{k} & 3 & 6 \ar[l]_-{l} \ar@(ur,dr)\\
	1 & 2 \ar@/^2ex/[u]^-{f} \ar[u]^-{g} \ar@/_2ex/[u]_{h} & 5
}
\]
\item
\[
    \begin{array}{c|c|c}
        E_H & s_H & t_H \\ \hline
        a & 1 & 2 \\
        b & 2 & 3 \\
        c & 2 & 3 \\
        d & 3 & 2 \\
        e & 3 & 4 \\
        f & 7 & 4 \\
        g & 2 & 6 \\
        h & 6 & 5 \\
        k & 6 & 7 \\
        l & 5 & 5
    \end{array}
    \hspace{2cm} V_H=\{1,2,3,4,5,6,7\}
\]
\end{enumerate}
\end{solucio}

\begin{exemple}
En una empresa petita hi ha tres departaments: Administració, Comptabilitat i Atenció al client. Durant el funcionament diari, els documents poden circular diverses vegades entre els mateixos departaments, i alguns departaments processen documents interns sense enviar-los enlloc. En un dia es donen les situacions següents:
\begin{itemize}
\item Dos documents passen d'Administració a Comptabilitat.
\item Un document passa de Comptabilitat a Administració.
\item Tres documents passen de Comptabilitat a Atenció al client.
\item Un document passa d'Atenció al client a Comptabilitat.
\item Dos documents són processats internament a Administració.
\item Un document és processat internament a Comptabilitat.
\item Un document és processat internament a Atenció al client.
\end{itemize}
Modela aquesta situació com un graf $G$, especificant-ne els elements, interpretant-ne les arestes i representant-lo gràficament.
\end{exemple}
\begin{solucio}
El graf $G$ té com a conjunt de vèrtexs $V_G=\{a,b,c\}$, on $a,b,c$ representen Administració, Comptabilitat i Atenció al client, respectivament. El conjunt d'arestes és
\[
\begin{aligned}
E_G = \{
& a\xrightarrow{f_1}b,\ a\xrightarrow{f_2}b, \\
& b\xrightarrow{g}a, \\
& b\xrightarrow{h_1}c,\ b\xrightarrow{h_2}c,\ b\xrightarrow{h_3}c, \\
& c\xrightarrow{k}b, \\
& a\xrightarrow{j_1}a,\ a\xrightarrow{j_2}a, \\
& b\xrightarrow{i}b, \\
& c\xrightarrow{l}c
\},
\end{aligned}
\]
on els subíndexs indiquen arestes diferents entre els mateixos vèrtexs (diversos documents amb el mateix recorregut) i els bucles $a\xrightarrow{j}a$, $b\xrightarrow{i}b$ i $c\xrightarrow{l}c$ representen processament intern. La representació gràfica de $G$ és
\[
\xymatrix@R=7ex@C=8ex{
	c \ar@(ul,dl)_l \ar[r]^-{k} & b \ar@/^1.5ex/[l]^-{h_2} \ar@/_2.5ex/[l]_{h_1} \ar@/_5ex/[l]_{h_3}  \ar@(dr,ur)_i \ar@/^4ex/[d]^-{g} \\
	& a \ar@/^4ex/[u]^-{f_1} \ar@/^1ex/[u]_(.45){f_2}  \ar@(dl,dr)_{j_{1}}  \ar@(dr,ur)_{j_{2}}
}
\]
Aquest graf modela el flux real de documents dins l'empresa en un dia: les direccions indiquen el sentit del flux, el nombre d'arestes entre dos departaments indica la quantitat de documents, i els bucles indiquen activitat interna.
\end{solucio}

\begin{exercici}
Defineix el graf $G$ perquè la seva representació gràfica sigui la següent:
\[
\xymatrix{
	\bullet \ar@<1ex>[r]  & \bullet \ar@<1ex>[l] \ar@(ul,ur)  & \bullet \ar[l] \ar@(ur,dr) \ar[dl] & \\
	\bullet \ar[ur] \ar[r]  & \bullet  & \bullet
}
\]
\end{exercici}

\begin{exercici}
Dibuixa el graf corresponent a les dades següents:
\[
    \begin{array}{c|c|c}
        E_G & s_G & t_G \\ \hline
        f & v & w \\
        g & w & x \\
        h & w & x \\
        i & y & y \\
        j & y & z \\
        k & z & y
    \end{array}
    \hspace{2cm} V_G=\{v,w,x,y,z\}
\]
\end{exercici}

\begin{exercici}
En un centre logístic hi ha quatre seccions: Recepció, Magatzem, Preparació de comandes i Expedició. Els paquets es mouen entre les seccions al llarg del dia, i alguns processos es fan dins d'una mateixa secció. En una jornada s'observen els fluxos següents:
\begin{itemize}
\item Dos paquets passen de Recepció a Magatzem.
\item Un paquet passa de Magatzem a Recepció.
\item Tres paquets passen de Magatzem a Preparació de comandes.
\item Dos paquets passen de Preparació de comandes a Expedició.
\item Un paquet passa d'Expedició a Magatzem.
\item Dos paquets són reprocessats dins de Magatzem.
\item Un paquet és revisat dins de Preparació de comandes.
\item Un paquet és verificat dins d'Expedició.
\end{itemize}
Modela aquesta situació mitjançant un graf. Especifica'n els elements i representa'l gràficament.
\end{exercici}

\section{Homomorfismes de grafs}\label{sec:grafs-homomorfismes}

L'estructura d'un graf
\[
G=\xymatrix{
E_G \ar@<0.6ex>[r]^{s_G} \ar@<-0.6ex>[r]_{t_G} & V_G
}
\]
consisteix en dos conjunts $V_G$ i $E_G$ i dues aplicacions $s_G$ i $t_G$. Per relacionar dos grafs, hem de poder relacionar adequadament els seus conjunts i les seves aplicacions.

\begin{definicio}\index{morfisme de grafs}
Considerem dos grafs
\[
G=\xymatrix{
E_G \ar@<0.6ex>[r]^{s_G} \ar@<-0.6ex>[r]_{t_G} & V_G
}
\qquad\text{i}\qquad
H=\xymatrix{
E_H \ar@<0.6ex>[r]^{s_H} \ar@<-0.6ex>[r]_{t_H} & V_H
}
\]
Un \textbf{homomorfisme} $F$ de $G$ a $H$, que escrivim $F\colon G\to H$, és un parell d'aplicacions $F_V\colon V_G\to V_H$ i $F_E\colon E_G\to E_H$ tals que els diagrames següents commuten:
\begin{equation}\label{grafs:hom}
\xymatrix{
	E_G \ar[d]_-{F_E} \ar[r]^-{s_G} & V_G \ar[d]^-{F_V}\\
	E_H \ar[r]_-{s_H} & V_H
}
\hspace{2cm}
\xymatrix{
	E_G \ar[d]_-{F_E} \ar[r]^-{t_G} & V_G \ar[d]^-{F_V}\\
	E_H \ar[r]_-{t_H} & V_H
}
\end{equation}
\end{definicio}

Dit breument, un homomorfisme \emph{preserva} l'origen (diagrama de l'esquerra),
\[
F_V\comp s_G=s_H\comp F_E,
\]
i el destí (diagrama de la dreta),
\[
F_V\comp t_G=t_H\comp F_E.
\]

\begin{exemple}
Considerem els dos grafs següents:
\[
G:\ 
\xymatrix{
	1 \ar[d]_{h} \ar[r]^{f} & 2 \ar[d]^{g} \ar@(ur,ul)_{k} \\
	4 & 3
}
\qquad\text{i}\qquad
H:\ 
\xymatrix{
	b \ar@(ur,ul)_{y} \ar@(dr,ur)_{z} \\
	a \ar[u]_{x} \ar[r]_{v} & c
}
\]
Defineix un homomorfisme $F$ entre ells.
\end{exemple}
\begin{solucio}
Com que $k$ és un bucle i l'únic vèrtex de $H$ amb bucles és $b$, necessàriament $F_V(2)=b$. Com que $g$ surt de $2$, la seva imatge ha de sortir de $b$, i les úniques arestes que surten de $b$ són els bucles $y,z$; per tant $g$ també ha d'anar a un bucle i $F_V(3)=b$. Aquesta part és forçada; la resta és una elecció. Per construir un exemple concret, triem $F_V(1)=a$ i $F_V(4)=c$, i posem $F_E(k)=F_E(g)=y$, $F_E(f)=x$ i $F_E(h)=v$. Les taules següents comproven que aquesta tria preserva origen i destí.

La distinció entre una condició necessària i una elecció és important: que $F_V(2)=F_V(3)=b$ estigui forçat no vol dir que l'homomorfisme sigui únic. De fet, una enumeració finita directa de les aplicacions compatibles dona $24$ homomorfismes entre aquests dos grafs, dels quals $8$ tenen $F_V(1)=a$ i $16$ tenen $F_V(1)=b$.

Comprovem que les condicions \eqref{grafs:hom} es compleixen. Les taules dels dos grafs són
\[
    \begin{array}{c|c|c}
        E_G & s_G & t_G \\ \hline
        f & 1 & 2  \\
        g & 2 & 3 \\
        h & 1 & 4 \\
        k & 2 & 2 
    \end{array}
    \qquad
    \begin{array}{c|c|c}
        E_H & s_H & t_H \\ \hline
        v & a & c  \\
        x & a & b \\
        y & b & b \\
        z & b & b 
    \end{array}
\]
i les dues aplicacions de l'homomorfisme són
\[
F_V =
\begin{array}{c|c}
    V_G & V_H \\ \hline
     1 & a \\
     2 & b \\
     3 & b \\
     4 & c
\end{array}
\qquad
F_E =
\begin{array}{c|c}
   E_G  & E_H \\ \hline
   f & x \\
   g & y \\
   h & v \\
   k & y
\end{array}
\]
Per exemple,
\[
\begin{aligned}
(F_V\comp s_G)(f)&=F_V(s_G(f))\\
&=F_V(1)=a,\\
(s_H\comp F_E)(f)&=s_H(F_E(f))\\
&=s_H(x)=a,
\end{aligned}
\]
de manera que $(F_V\comp s_G)(f)=(s_H\comp F_E)(f)$; deixem la resta d'arestes com a exercici. Anàlogament,
\[
\begin{aligned}
(F_V\comp t_G)(h)&=F_V(t_G(h))\\
&=F_V(4)=c,\\
(t_H\comp F_E)(h)&=t_H(F_E(h))\\
&=t_H(v)=c,
\end{aligned}
\]
i per tant $(F_V\comp t_G)(h)=(t_H\comp F_E)(h)$.
\end{solucio}

\begin{exercici}
Considerem els grafs
\[
G:\ 
\xymatrix{
    1 \ar[r]^-{a} & 2 \ar[r]^-{b}& 3
}
\qquad
H:\ 
\xymatrix{
    4 \ar@<0.6ex>[r]^{d} \ar@<-0.6ex>[r]_{c} & 5 \ar@(ur,dr)^{e}
}
\]
Troba un homomorfisme $F$ de $G$ a $H$ i comprova les igualtats \eqref{grafs:hom}.
\end{exercici}

\begin{exercici}
Considerem el graf $G$ definit per
\[
    \begin{array}{c|c|c}
        E_G & s_G & t_G \\ \hline
        a & 1 & 2 \\
        b & 2 & 3 \\
        c & 1 & 4 \\
        d & 1 & 4 \\
        e & 5 & 6 
    \end{array}
    \hspace{2cm} V_G=\{1,2,3,4,5,6\}
\]
i el graf $H$ representat per
\[
\xymatrix{
	m \ar[d]_{y} \ar@<0.6ex>[r]^{w} & n \ar[l]^{x} \\
	p  & q \ar@(dr,ur)_{z}
}
\]
Es demana:
\begin{enumerate}
\item representar gràficament el graf $G$;
\item trobar un homomorfisme $F$ de $G$ a $H$ i comprovar les igualtats \eqref{grafs:hom}.
\end{enumerate}
\end{exercici}

\section{Camins en un graf}

Sigui $G$ un graf. Un \emph{camí} a $G$ és qualsevol successió d'arestes tal que el destí d'una aresta sigui l'origen de la següent. Això inclou successions de longitud~1, que són elements de $E_G$, i successions de longitud~0, que comencen i acaben al mateix vèrtex sense seguir cap aresta.

\begin{definicio}
Sigui $n\geq 1$. En un graf
\[
G=\xymatrix{
E_G \ar@<0.6ex>[r]^{s_G} \ar@<-0.6ex>[r]_{t_G} & V_G
}
\]
un \textbf{camí de longitud} $n$ del vèrtex $a$ al vèrtex $b$ és una successió d'arestes $(f_1,f_2,\dots,f_n)$, escrita en \emph{ordre de recorregut} ---primer es recorre $f_1$---, tal que
\begin{enumerate}
\item[(i)] $s_G(f_1)=a$;
\item[(ii)] $t_G(f_i)=s_G(f_{i+1})$ per a $i=1,\dots,n-1$;
\item[(iii)] $t_G(f_n)=b$.
\end{enumerate}
El conjunt de camins de longitud $n$ s'escriu $C_n$. En particular, $C_2$ és el conjunt de parells d'arestes $(f,g)$ amb $t_G(f)=s_G(g)$, anomenats parells d'\emph{arestes componibles}.
\end{definicio}

Per conveni, per a cada vèrtex $a\in V_G$ hi ha un únic camí de longitud~0 d'$a$ a si mateix, que denotem $()_a$ i anomenem \emph{camí buit} d'$a$; és el cas $n=0$, exclòs de la definició anterior perquè aleshores no hi ha cap aresta $f_1,\dots,f_n$. Si dibuixem un camí
\[
\xymatrix{
\cdot \ar[r]^{f_{1}} & \cdot \ar[r]^-{f_{2}} & \cdots \ar[r]^{f_{n-1}} & \cdot \ar[r]^{f_{n}} & \cdot
}
\]
l'aresta $f_1$ queda a l'esquerra i es recorre primer; $f_n$ queda a la dreta i es recorre última. Si $f\in E_G$, aleshores $(f)$ és un camí de longitud~1, és a dir, $(f)\in C_1$.

\begin{exemple}\label{ex:grafs-camins-G}
Considerem el graf $G$ de l'exemple d'homomorfisme de la secció anterior:
\[
G:\ 
\xymatrix{
	1 \ar[d]_{h} \ar[r]^{f} & 2 \ar[d]^{g} \ar@(ur,ul)_{k} \\
	4 & 3
}
\]
Escriu els camins de longitud $0$, $1$, $2$ i $3$ de $G$, i descriu tots els camins d'un vèrtex a un altre.
\end{exemple}
\begin{solucio}
Recordem la taula de $G$: $f\colon 1\to 2$, $g\colon 2\to 3$, $h\colon 1\to 4$ i el bucle $k\colon 2\to 2$.

\emph{Longitud $0$.} Hi ha un camí buit per a cada vèrtex: $C_0=\{()_1,\ ()_2,\ ()_3,\ ()_4\}$.

\emph{Longitud $1$.} Són les arestes: $C_1=\{(f),\ (g),\ (h),\ (k)\}$.

\emph{Longitud $2$.} Busquem els parells $(u,v)$ amb $t_G(u)=s_G(v)$. Com que $t_G(f)=t_G(k)=2$ i les arestes que surten de $2$ són $g$ i $k$, mentre que de $3$ i de $4$ no en surt cap,
\[
C_2=\{(f,g),\ (f,k),\ (k,g),\ (k,k)\}.
\]

\emph{Longitud $3$.} Cada camí de longitud $2$ que acaba a $2$ es pot allargar amb $g$ o amb $k$; els que acaben a $3$ no es poden allargar:
\[
C_3=\{(f,k,g),\ (f,k,k),\ (k,k,g),\ (k,k,k)\}.
\]

\emph{Descripció general.} Escrivim $k^n$ per a la successió $(k,\dots,k)$ de $n$ bucles, amb $k^0$ buida. Els camins d'un vèrtex a un altre són exactament:
\[
\begin{array}{c|l}
\text{d'origen a destí} & \text{camins} \\ \hline
1\to 1,\ 3\to 3,\ 4\to 4 & \text{només el camí buit} \\
1\to 2 & (f,k^n),\ n\geq 0 \\
1\to 3 & (f,k^n,g),\ n\geq 0 \\
1\to 4 & (h) \\
2\to 2 & k^n,\ n\geq 0 \ (\text{per a } n=0,\ \text{el camí buit } ()_2) \\
2\to 3 & (k^n,g),\ n\geq 0
\end{array}
\]
i no n'hi ha cap més. El bucle $k$ fa que el nombre de camins sigui infinit, però per a cada longitud $n\geq 2$ n'hi ha exactament quatre.
\end{solucio}

\begin{exemple}\label{ex:grafs-camins-H}
Fem el mateix amb el graf $H$ del mateix exemple:
\[
H:\ 
\xymatrix{
	b \ar@(ur,ul)_{y} \ar@(dr,ur)_{z} \\
	a \ar[u]_{x} \ar[r]_{v} & c
}
\]
Escriu els camins de longitud $0$, $1$ i $2$ de $H$, i compta quants camins de longitud $n$ hi ha.
\end{exemple}
\begin{solucio}
La taula de $H$ és $x\colon a\to b$, $v\colon a\to c$ i els bucles $y,z\colon b\to b$.

\emph{Longitud $0$ i $1$.} $C_0=\{()_a,\ ()_b,\ ()_c\}$ i $C_1=\{(x),\ (v),\ (y),\ (z)\}$.

\emph{Longitud $2$.} De $c$ no surt cap aresta, de manera que $v$ no es pot allargar. Després de $x$, $y$ o $z$ som a $b$, i podem continuar amb $y$ o amb $z$:
\[
C_2=\{(x,y),\ (x,z),\ (y,y),\ (y,z),\ (z,y),\ (z,z)\}.
\]

\emph{Recompte.} Un camí de longitud $n\geq 1$ de $b$ a $b$ és una successió qualsevol de $n$ bucles, triant a cada pas entre $y$ i $z$: n'hi ha $2^n$. Un camí de longitud $n\geq 2$ que surt d'$a$ ha de començar per $x$ i continuar amb $n-1$ bucles: n'hi ha $2^{n-1}$. L'únic camí que va d'$a$ a $c$ és $(v)$. Per tant, per a $n\geq 2$ hi ha
\[
2^n+2^{n-1}=3\cdot 2^{n-1}
\]
camins de longitud $n$; per exemple, $6$ de longitud $2$, d'acord amb la llista anterior.

Comparat amb l'exemple anterior, un sol bucle produeix un nombre constant de camins per a cada longitud, mentre que dos bucles al mateix vèrtex en produeixen un nombre que creix exponencialment.
\end{solucio}

\section{De grafs a categories}

Podem passar d'un graf a una categoria afegint-hi composició i identitats. La construcció següent és fonamental: mostra que tot graf genera una categoria de manera lliure.

Per a qualsevol graf $G$ hi ha una categoria $\Free(G)$ els objectes de la qual són els vèrtexs de $G$ i els morfismes de la qual són els camins de $G$. Un camí d'$a$ a $b$ és un morfisme $a\to b$. La composició s'obté concatenant camins, en \emph{ordre de recorregut}: si $p=(a_1,\dots,a_r)$ és un camí d'$A$ a $B$ i $q=(b_1,\dots,b_s)$ un camí de $B$ a $C$, la seva composició és el camí d'$A$ a $C$
\[
q\comp p=(a_1,\dots,a_r,b_1,\dots,b_s),
\]
on, d'acord amb la convenció categòrica $q\comp p$, el camí $p$ ---el que surt d'$A$--- es recorre primer. Per a cada vèrtex $A$, la identitat $\id_A$ és el camí buit $()_A$.

Comprovem que $\Free(G)$ és una categoria. La composició és associativa perquè la concatenació de successions ho és: en concatenar tres camins componibles, el resultat és la successió de totes les seves arestes en ordre, amb independència de l'ordre en què s'agrupin. El camí buit és neutre, ja que concatenar-lo no afegeix cap aresta: $(a_1,\dots,a_r)\comp()_A=(a_1,\dots,a_r)=()_B\comp(a_1,\dots,a_r)$ per a tot camí d'$A$ a $B$. Per tant, $\Free(G)$ satisfà els axiomes de categoria. S'anomena \textbf{categoria lliure generada pel graf} $G$, o també \textbf{categoria de camins} de $G$.\index{categoria!lliure}\index{categoria!de camins}

\begin{exemple}\label{ex:grafs-free-cadena}
Especifica la categoria $\Free(G)$ generada pel graf
\[
G:\ 
\xymatrix{
    1 \ar[r]^-{a} & 2 \ar[r]^-{b}& 3
}
\]
de l'exercici d'homomorfismes de la secció~\ref{sec:grafs-homomorfismes}.
\end{exemple}
\begin{solucio}
Els objectes són els vèrtexs $1$, $2$ i $3$. Els morfismes són els camins, i com que el graf no té cicles, n'hi ha un nombre finit:
\[
\begin{array}{c|l}
\Hom(1,1)=\{()_1\} & \Hom(1,2)=\{(a)\} \\
\Hom(2,2)=\{()_2\} & \Hom(2,3)=\{(b)\} \\
\Hom(3,3)=\{()_3\} & \Hom(1,3)=\{(a,b)\}
\end{array}
\]
i tots els altres homconjunts ($\Hom(2,1)$, $\Hom(3,1)$, $\Hom(3,2)$) són buits. En total, sis morfismes: les tres identitats, les dues arestes i el camí de longitud~$2$.

L'única composició que no involucra identitats és la de $(a)$ i $(b)$:
\[
\xymatrix{
1 \ar[r]^-{(a)} \ar[dr]_-{(a,b)} & 2 \ar[d]^-{(b)} \\
& 3
}
\qquad\qquad
(b)\comp(a)=(a,b).
\]
Les composicions amb identitats són les que imposen els axiomes, perquè els camins buits són neutres per a la concatenació.

Com que entre dos objectes hi ha com a màxim un morfisme, $\Free(G)$ és una categoria prima. De fet, és la categoria associada a l'ordre $1<2<3$ ---hi ha un morfisme $i\to j$ exactament quan $i\leq j$---, isomorfa, doncs, a la categoria $\cat{3}$ de l'ordre $0<1<2$ que apareix al capítol~\ref{cap:categories}.
\end{solucio}

\begin{exemple}\label{ex:grafs-free-bucle}
Especifica la categoria $\Free(H)$ generada pel graf
\[
H:\ 
\xymatrix{
    4 \ar@<0.6ex>[r]^{d} \ar@<-0.6ex>[r]_{c} & 5 \ar@(ur,dr)^{e}
}
\]
del mateix exercici.
\end{exemple}
\begin{solucio}
Els objectes són $4$ i $5$. Escrivim $e^n$ per al camí $(e,\dots,e)$ de $n$ bucles, amb $e^0=()_5$. Els camins de $H$ donen els homconjunts següents:
\[
\begin{aligned}
\Hom(4,4)&=\{()_4\},\\
\Hom(5,5)&=\{e^n : n\geq 0\},\\
\Hom(4,5)&=\{(c,e^n) : n\geq 0\}\cup\{(d,e^n) : n\geq 0\},\\
\Hom(5,4)&=\varnothing.
\end{aligned}
\]
De $5$ a $4$ no hi ha cap camí, perquè cap aresta surt de $5$ cap a $4$.

La composició és la concatenació. Per a bucles, $e^m\comp e^n=e^{m+n}$, i per als camins de $4$ a $5$,
\[
\xymatrix{
4 \ar[r]^-{(c,e^n)} \ar[dr]_-{(c,e^{n+m})} & 5 \ar[d]^-{e^m} \\
& 5
}
\qquad\qquad
e^m\comp(c,e^n)=(c,e^{n+m}),
\]
i igual amb $d$ en lloc de $c$.

Hi ha dues observacions a fer. Primera: $\Hom(5,5)$ amb la composició és el monoide $(\mathbb N,+,0)$, via $e^n\mapsto n$; la subcategoria plena sobre l'objecte $5$ és, doncs, la categoria d'un sol objecte d'aquest monoide. Segona: les arestes paral·leles $c$ i $d$ donen morfismes diferents, $(c)\neq(d)$, i també $(c,e^n)\neq(d,e^n)$ per a tot $n$. La categoria lliure no identifica mai dos camins diferents: no imposa cap relació que no vingui forçada pels axiomes.
\end{solucio}

\chapter{Convenció de mida}
\label{cap:mida}

Al llarg del llibre hem parlat de categories \emph{petites}, \emph{localment petites} i de \emph{classes pròpies} sense fixar del tot el marc conjuntista. Aquest apèndix el fa explícit, perquè el volum de lògica categòrica ---on apareixen categories de models, doctrines i fibracions--- l'exigirà.

\medskip
\noindent\textbf{Marc adoptat: classes i conjunts.} Treballem en una teoria de conjunts amb \textbf{classes pròpies}, a l'estil de von Neumann--Bernays--Gödel (NBG). Hi ha dos nivells de col·leccions: els \emph{conjunts}, que poden ser elements d'altres col·leccions, i les \emph{classes pròpies}, que no. Amb això:
\begin{itemize}
  \item una categoria és \textbf{petita} si la seva col·lecció d'objectes i la de morfismes són totes dues \emph{conjunts};
  \item és \textbf{localment petita} si tots els homconjunts $\Hom(A,B)$ són conjunts, tot i que la col·lecció d'objectes pot ser una classe pròpia;
  \item és \textbf{ben potenciada} si, per a cada objecte $A$, els subobjectes de $A$ formen un conjunt.
\end{itemize}
L'última condició demana una precisió. Un \emph{subobjecte} de $A$ no és un monomorfisme concret cap a $A$, sinó una classe d'equivalència de monomorfismes. Dos monomorfismes $m\colon S\rightarrowtail A$ i $n\colon T\rightarrowtail A$ són equivalents quan cadascun factoritza a través de l'altre, és a dir, quan difereixen en un isomorfisme entre els dominis. La distinció és essencial, i en el marc de classes cal formular-la amb cura. A $\cat{Set}$, els monomorfismes cap a un conjunt donat formen una classe pròpia, perquè cada subconjunt té moltes còpies isomorfes; i cada classe d'equivalència és, ella mateixa, una classe pròpia. Com que les classes pròpies no poden ser elements de cap col·lecció, \enquote{el conjunt de les classes d'equivalència} no té sentit literal. La definició precisa és, doncs, la següent: una categoria és \textbf{ben potenciada} si, per a cada objecte $A$, hi ha un \emph{conjunt} $R_A$ de monomorfismes cap a $A$ tal que tot monomorfisme cap a $A$ és equivalent a algun element de $R_A$. Aleshores es defineix $\Sub(A)$ com el conjunt quocient $R_A/{\sim}$, i $[m]$ designa la classe dels representants equivalents a $m$. Un altre conjunt de representants dona un quocient en bijecció canònica amb aquest, i la bijecció respecta l'ordre; per això el resultat no depèn de la tria de $R_A$. A $\cat{Set}$ es pot prendre com a $R_A$ el conjunt de les inclusions dels subconjunts de $A$, i s'obté $\Sub(A)\cong\mathcal P(A)$. La resta del llibre parla de \enquote{classes de monomorfismes} en aquest sentit.

Les categories $\cat{Set}$, $\cat{Grp}$, $\cat{Top}$, $\cat{Graph}$ i $\cat{Cat}$ són localment petites però no petites; cada preordre concret sobre un conjunt i cada categoria finita concreta són petites. Això no s'ha de confondre amb la petitesa de la categoria de \emph{tots} els preordres o de \emph{totes} les categories finites, que ---en tenir una classe pròpia de conjunts subjacents--- són grans llevat que es fixi una convenció de mida o s'esculli un conjunt de representants.

Quan alguna construcció demani triar simultàniament un objecte (o un morfisme) per a cada objecte d'una categoria \emph{gran} ---per exemple, en construir un invers d'una equivalència---, la família de tries està indexada per una classe pròpia, i l'axioma de l'elecció ordinari, que val per a famílies indexades per un conjunt, no hi arriba. En aquests casos assumim la forma \textbf{global} de l'elecció (disponible, per exemple, a NBG amb l'axioma d'elecció global). Ho indicarem allà on calgui.

\medskip
\noindent\textbf{L'alternativa dels universos.} Una segona opció, molt usada en geometria algebraica i teoria de topos, fixa un o més \textbf{universos de Grothendieck} $\mathcal{U}$ ---conjunts prou grans i tancats per les operacions habituals--- i parla de categories $\mathcal{U}$-petites i $\mathcal{U}$-grans de manera \emph{relativa} a $\mathcal{U}$. Aquesta convenció té l'avantatge que \enquote{la categoria de tots els $\mathcal{U}$-conjunts} és, ella mateixa, un objecte d'un univers més gran, cosa que simplifica certes construccions. En aquest volum \emph{no} l'adoptem, per no carregar una introducció amb la teoria d'universos; però el lector que continuï cap a la teoria de feixos la retrobarà. Les dues opcions són maneres alternatives de \emph{gestionar la mida} més que no pas fonaments idèntics: a grans trets, el paper de les \enquote{classes pròpies} de NBG el fa aquí la distinció entre $\mathcal{U}$-petit i $\mathcal{U}$-gran relativa a un univers $\mathcal{U}$ fixat. Fer precisa aquesta correspondència exigeix especificar l'univers i les hipòtesis pertinents; no cal, però, que això sigui un obstacle per a una primera lectura.

\medskip
\noindent Fixar aquesta convenció des d'ara evita l'ambigüitat que, altrament, es faria notar quan apareguin categories de models i functors classificadors: en tot moment, cada construcció d'aquest llibre declara si viu en un conjunt o en una classe pròpia.

\medskip
\noindent\textbf{Per saber-ne més.} Les qüestions de mida es tracten, al nivell d'aquest llibre, a \cite[§1.8]{awodey}, \cite[§1.1]{riehl} ---on l'observació~1.1.5 presenta la convenció d'universos---, \cite[§3.2]{leinster} i \cite[cap.~1]{perrone}. El capítol~3 de \cite{leinster} presenta, a més, la perspectiva en què el mateix llenguatge estructural fa de fonament, sense una teoria de conjunts prèvia; en són les fonts \cite{lawvere-etcs} i, com a introducció informal, \cite{leinster-rethinking}. Per a una comparació detallada de les diferents maneres de gestionar la mida ---classes, universos i altres alternatives---, adequada un cop se'n té una certa experiència, vegeu \cite{shulman-mida}.

\chapter{Solucions dels tests}
\label{cap:solucions-tests}

\ifimprimible
Aquest apèndix reuneix les respostes i les justificacions dels tests. En aquesta edició per imprimir, són la manera de corregir-los.
\else
Aquest apèndix reuneix les respostes i les justificacions dels tests interactius. Serveixen per als lectors de PDF que no admeten la correcció automàtica i per a l'ús en paper.
\fi Es recomana consultar-les només després d'haver respost tot el test: cada justificació explica per què la resposta correcta ho és i, sovint, on falla el distractor més temptador.

\imprimeixsolucionstests

\ifpublicacioparcial\else
\chapter[Quasiexemples i contraexemples]{Quasiexemples\\ i contraexemples}
\label{cap:contraexemples}

Una definició matemàtica s'entén de veritat quan se'n coneixen tant els exemples com els \emph{quasiexemples}: estructures que s'hi assemblen molt, que apareixen de manera natural, i que tanmateix fallen en un punt precís. Aquest apèndix recull quasiexemples i contraexemples per a les nocions centrals del llibre: categoria, functor, transformació natural i equivalència. Acaba amb models de la lògica que queden fora de l'explicació per topos dels capítols~\ref{cap:subobjectes}--\ref{cap:pont}.

La secció~\ref{sec:igualtat-equivalencia}, sobre igualtat, isomorfisme, isomorfisme natural i equivalència, és la més important. Confondre aquests nivells és un error freqüent, i greu, perquè porta a demostrar coses falses o a formular propietats que no tenen sentit categòric.

El nivell és el del llibre, i els exemples s'han triat perquè només demanin els coneixements que el llibre ja pressuposa. Tota afirmació que es pot justificar amb aquests coneixements es demostra. Quan un comentari remet a teories que el llibre no desenvolupa, es presenta de manera informativa i es diu explícitament.

% =====================================================================
\section{Què no és una categoria}
\label{sec:no-categoria}
% =====================================================================

\subsection{Les condicions que es poden trencar}

La definició~\ref{def:categoria} demana unes dades i dos axiomes. Per poder dir, exemple a exemple, quina condició falla, les desglossem en cinc:
\begin{enumerate}
\item[(C1)] \textbf{Extrems.} Cada morfisme $f$ té un \emph{únic} domini i un \emph{únic} codomini.
\item[(C2)] \textbf{Composició.} Per a cada parella componible $f\colon A\to B$, $g\colon B\to C$ hi ha un morfisme $g\comp f\colon A\to C$, determinat sense ambigüitat.
\item[(C3)] \textbf{Associativitat.} $h\comp(g\comp f)=(h\comp g)\comp f$.
\item[(C4)] \textbf{Existència d'identitats.} Per a cada objecte $A$ hi ha un morfisme $\id_A\colon A\to A$.
\item[(C5)] \textbf{Lleis d'identitat.} $f\comp\id_A=f=\id_B\comp f$ per a tot $f\colon A\to B$.
\end{enumerate}
Les igualtats de (C3) i (C5) són \emph{igualtats}: no isomorfismes, ni homotopies, ni cap mena d'equivalència. Gairebé tots els exemples d'aquesta secció fallen exactament aquí. Recordem també que la demostració de la unicitat de la identitat (capítol~\ref{cap:categories}) no fa servir l'associativitat: només les lleis d'identitat. Ho aprofitarem en parlar dels magmes.

\subsection{L'exemple central: els camins en un espai}

Aquest exemple fa servir la noció de continuïtat. Sigui $X$ un espai topològic. Proposem:
\begin{itemize}
\item \emph{objectes}: els punts de $X$;
\item \emph{morfismes}: els camins, és a dir, les aplicacions contínues $\gamma\colon[0,1]\to X$; escrivim $\gamma\colon x\to y$ si $\gamma(0)=x$ i $\gamma(1)=y$;
\item \emph{composició}: la concatenació. Si $\gamma\colon x\to y$ i $\delta\colon y\to z$, definim $\delta\ast\gamma\colon x\to z$ per
\[
(\delta\ast\gamma)(t)=
\begin{cases}
\gamma(2t), & 0\le t\le\tfrac12,\\
\delta(2t-1), & \tfrac12\le t\le1;
\end{cases}
\]
\item \emph{candidats a identitat}: els camins constants $c_x(t)=x$.
\end{itemize}
Escrivim $\delta\ast\gamma$ (primer $\gamma$, després $\delta$) per coherència amb $g\comp f$. Les dues branques coincideixen a $t=\frac12$, perquè $\gamma(1)=y=\delta(0)$, i per tant $\delta\ast\gamma$ és contínua. Així (C1) i (C2) es compleixen. El problema és a (C3) i (C5).

\paragraph{Fallada de l'associativitat.} Siguin $\gamma\colon x\to y$, $\delta\colon y\to z$ i $\varepsilon\colon z\to w$. Un càlcul directe dona
\[
\bigl((\varepsilon\ast\delta)\ast\gamma\bigr)(t)=
\begin{cases}
\gamma(2t), & t\in[0,\frac12],\\
\delta(4t-2), & t\in[\frac12,\frac34],\\
\varepsilon(4t-3), & t\in[\frac34,1],
\end{cases}
\qquad
\bigl(\varepsilon\ast(\delta\ast\gamma)\bigr)(t)=
\begin{cases}
\gamma(4t), & t\in[0,\frac14],\\
\delta(4t-1), & t\in[\frac14,\frac12],\\
\varepsilon(2t-1), & t\in[\frac12,1].
\end{cases}
\]
Tots dos recorren els mateixos tres trams en el mateix ordre, però a velocitats diferents.

\begin{exemple}[Tres segments a la recta]\label{ex:camins-recta}
A $X=\mathbb R$, prenem $\gamma(t)=t$, $\delta(t)=1+t$ i $\varepsilon(t)=2+t$, de manera que $\gamma\colon0\to1$, $\delta\colon1\to2$ i $\varepsilon\colon2\to3$. Avaluant a $t=\frac12$,
\[
\bigl((\varepsilon\ast\delta)\ast\gamma\bigr)(\tfrac12)=\gamma(1)=1,
\qquad
\bigl(\varepsilon\ast(\delta\ast\gamma)\bigr)(\tfrac12)=\varepsilon(0)=2.
\]
Els dos camins van de $0$ a $3$ i tenen la mateixa imatge, l'interval $[0,3]$, però són diferents: falla (C3).
\end{exemple}

\paragraph{Fallada de les identitats.} El fracàs és més radical que el de l'associativitat:

\begin{proposicio}\label{prop:camins-sense-identitat}
Sigui $\gamma\colon x\to y$ un camí injectiu. Aleshores cap camí $e\colon y\to y$ compleix $e\ast\gamma=\gamma$. Anàlogament, si $\gamma\colon y\to z$ és injectiu, cap camí $e\colon y\to y$ compleix $\gamma\ast e=\gamma$.
\end{proposicio}

\begin{prova}
Per a qualsevol $e$, $(e\ast\gamma)(\frac14)=\gamma(\frac12)$, que és diferent de $\gamma(\frac14)$ perquè $\gamma$ és injectiu. Per a l'altra afirmació, $(\gamma\ast e)(\frac34)=\gamma(\frac12)\neq\gamma(\frac34)$.
\end{prova}

A $X=\mathbb R$, entre dos punts diferents qualssevol hi ha camins injectius. Per tant, no és que s'hagin triat malament els camins constants: (C4)--(C5) no es poden satisfer de cap manera amb aquesta composició.

\paragraph{Però fallen de manera controlada.} Les igualtats que fallen valen \emph{llevat d'una reparametrització}, i les reparametritzacions són inofensives llevat d'homotopia. Dos camins $\alpha,\beta\colon x\to y$ són \emph{homòtops relativament als extrems}, i escrivim $\alpha\simeq\beta$, si hi ha una aplicació contínua $H\colon[0,1]\times[0,1]\to X$ amb $H(t,0)=\alpha(t)$, $H(t,1)=\beta(t)$, $H(0,s)=x$ i $H(1,s)=y$.

\begin{lema}[Reparametrització]\label{lema:reparametritzacio}
Sigui $\alpha\colon[0,1]\to X$ un camí i $\varphi\colon[0,1]\to[0,1]$ contínua amb $\varphi(0)=0$ i $\varphi(1)=1$. Aleshores $\alpha\comp\varphi\simeq\alpha$.
\end{lema}

\begin{prova}
L'aplicació $H(t,s)=\alpha\bigl((1-s)\varphi(t)+st\bigr)$ és contínua i ben definida, perquè $(1-s)\varphi(t)+st\in[0,1]$. Val $\alpha\comp\varphi$ per a $s=0$ i $\alpha$ per a $s=1$, i $H(0,s)=\alpha(0)$, $H(1,s)=\alpha(1)$ per a tot $s$.
\end{prova}

\begin{proposicio}\label{prop:associador-camins}
\begin{enumerate}
\item[(a)] $(\varepsilon\ast\delta)\ast\gamma=\bigl(\varepsilon\ast(\delta\ast\gamma)\bigr)\comp\varphi$, on $\varphi(t)=t/2$ a $[0,\frac12]$, $\varphi(t)=t-\frac14$ a $[\frac12,\frac34]$ i $\varphi(t)=2t-1$ a $[\frac34,1]$.
\item[(b)] $c_y\ast\gamma=\gamma\comp\psi$ amb $\psi(t)=\min(2t,1)$, i $\gamma\ast c_x=\gamma\comp\psi'$ amb $\psi'(t)=\max(2t-1,0)$.
\end{enumerate}
En conseqüència, $(\varepsilon\ast\delta)\ast\gamma\simeq\varepsilon\ast(\delta\ast\gamma)$, $c_y\ast\gamma\simeq\gamma$ i $\gamma\ast c_x\simeq\gamma$.
\end{proposicio}

\begin{prova}
Les igualtats (a) i (b) es comproven tram a tram amb les fórmules anteriors; per exemple, a $[\frac12,\frac34]$, $\varphi(t)=t-\frac14\in[\frac14,\frac12]$ i $\delta\bigl(4(t-\frac14)-1\bigr)=\delta(4t-2)$. Les funcions $\varphi$, $\psi$ i $\psi'$ són contínues i fixen $0$ i $1$, i el lema~\ref{lema:reparametritzacio} dona les homotopies.
\end{prova}

\paragraph{Primer remei: passar al quocient.} Si $\gamma\simeq\gamma'$ i $\delta\simeq\delta'$, concatenant les homotopies per a cada $s$ fix s'obté $\delta\ast\gamma\simeq\delta'\ast\gamma'$. Per tant, la concatenació indueix una composició ben definida entre classes, $[\delta]\comp[\gamma]:=[\delta\ast\gamma]$, i la proposició~\ref{prop:associador-camins} esdevé exactament (C3) i (C5) entre classes, amb $\id_x=[c_x]$. S'obté una categoria, el \textbf{grupoide fonamental} $\Pi_1(X)$\index{grupoide fonamental}: els objectes són els punts de $X$ i $\Hom(x,y)$ és el conjunt de classes d'homotopia de camins de $x$ a $y$. És un grupoide perquè el camí recorregut a l'inrevés, $\bar\gamma(t)=\gamma(1-t)$, és un invers de $[\gamma]$ (la comprovació és una homotopia explícita que no reproduïm). L'homconjunt $\Hom(x,x)$ és el que en topologia s'anomena \emph{grup fonamental} de $X$ en el punt $x$ (informatiu).

\paragraph{Segon remei: canviar els morfismes.} L'origen del problema és que tots els camins viuen en l'interval $[0,1]$, i cada concatenació ha de comprimir el temps a la meitat. Si cada camí té la seva durada, el problema desapareix. Un \textbf{camí de Moore}\index{camí de Moore} és un parell $(r,\gamma)$ amb $r\ge0$ i $\gamma\colon[0,r]\to X$ contínua, i es concatena sumant les durades:
\[
(s,\delta)\comp(r,\gamma)=(r+s,\ \delta\star\gamma),\qquad
(\delta\star\gamma)(t)=
\begin{cases}
\gamma(t), & t\in[0,r],\\
\delta(t-r), & t\in[r,r+s].
\end{cases}
\]

\begin{proposicio}
Els punts de $X$ i els camins de Moore formen una categoria, amb identitats $\id_x=(0,c_x)$.
\end{proposicio}

\begin{prova}
Les dues maneres d'agrupar tres camins $(r,\gamma)$, $(s,\delta)$ i $(u,\varepsilon)$ donen el camí de durada $r+s+u$ que val $\gamma(t)$ a $[0,r]$, $\delta(t-r)$ a $[r,r+s]$ i $\varepsilon(t-r-s)$ a $[r+s,r+s+u]$: \emph{la mateixa funció}. Les lleis d'identitat també són literals, perquè un camí de durada $0$ no ocupa temps: $(r,\gamma)\comp(0,c_x)$ té durada $r$ i val $\gamma(t-0)=\gamma(t)$.
\end{prova}

Tenim, doncs, les dues estratègies que reapareixen constantment en matemàtiques: \emph{quocientar}, és a dir, oblidar la informació que impedeix l'estrictesa (com a $\Pi_1(X)$), i \emph{estrictificar}, és a dir, substituir la composició per una altra que conservi la informació que interessa i en la qual les lleis valguin literalment (com amb els camins de Moore). Aquí no es pot parlar d'equivalència de categories, perquè l'estructura de partida no és una categoria. Els camins d'arestes d'un graf es concatenen sense problemes, amb el camí buit com a identitat ---és la categoria lliure de la definició~\ref{def:free}--- per la mateixa raó: la longitud s'hi suma i no es renormalitza mai.

\paragraph{Tercera via: fer teoria de la feblesa (informatiu).} La proposició~\ref{prop:associador-camins} no diu només que les dues associacions són homòtopes: en dona una homotopia \emph{concreta}, l'\emph{associador}. Amb quatre camins hi ha cinc maneres de posar els parèntesis, i els associadors les connecten formant un pentàgon. Si s'exigeix que els dos camins del pentàgon donin \emph{exactament} el mateix resultat, s'obté l'estructura d'una \emph{bicategoria}. Si només es demana que coincideixin llevat d'una nova homotopia, i així successivament a tots els nivells, s'arriba a les $\infty$-categories. Cap de les dues no es tracta en aquest llibre.

\subsection{Un sol objecte: els magmes}

Un \textbf{magma}\index{magma} és un conjunt $M$ amb una operació binària qualsevol, $(a,b)\mapsto ab$; un \emph{semigrup} és un magma associatiu, i un monoide, un semigrup amb element neutre. Com al capítol~\ref{cap:categories} per als monoides, a tot magma li podem associar una estructura $\cat{B}M$ amb un sol objecte $\ast$, un morfisme $\ast\to\ast$ per a cada $a\in M$ i la composició $a\comp b=ab$. Les condicions (C1) i (C2) es compleixen sempre, i $\cat{B}M$ és una categoria si i només si $M$ és un monoide. Cada magma no associatiu, o sense neutre, és un quasiexemple de categoria amb un objecte.

\begin{exemple}[La resta d'enters]
$M=(\mathbb Z,-)$. No és associatiu: $(1-1)-1=-1$, però $1-(1-1)=1$. Té neutre per la dreta, $a-0=a$, però no per l'esquerra: $e-a=a$ per a tot $a$ obligaria $e=2a$ per a tot $a$. Pel raonament de la unicitat de la identitat, un neutre bilateral hauria de coincidir amb $0$, que no ho és. Així, (C5) es compleix a mitges.
\end{exemple}

\begin{exemple}[Pedra, paper, tisora]\label{ex:ppt}
$M=\{\mathrm{Pe},\mathrm{Pa},\mathrm{Ti}\}$, on $xy$ és el guanyador de la partida entre $x$ i $y$, i $xx=x$:
\[
\begin{array}{c|ccc}
 & \mathrm{Pe} & \mathrm{Pa} & \mathrm{Ti}\\\hline
\mathrm{Pe} & \mathrm{Pe} & \mathrm{Pa} & \mathrm{Pe}\\
\mathrm{Pa} & \mathrm{Pa} & \mathrm{Pa} & \mathrm{Ti}\\
\mathrm{Ti} & \mathrm{Pe} & \mathrm{Ti} & \mathrm{Ti}
\end{array}
\]
És commutatiu i idempotent, però $(\mathrm{Pe}\,\mathrm{Pa})\,\mathrm{Ti}=\mathrm{Pa}\,\mathrm{Ti}=\mathrm{Ti}$ i $\mathrm{Pe}\,(\mathrm{Pa}\,\mathrm{Ti})=\mathrm{Pe}\,\mathrm{Ti}=\mathrm{Pe}$. Tampoc no té neutre: $\mathrm{Pe}\,e=\mathrm{Pe}$ exigeix $e\in\{\mathrm{Pe},\mathrm{Ti}\}$, $\mathrm{Pa}\,e=\mathrm{Pa}$ exigeix $e\in\{\mathrm{Pe},\mathrm{Pa}\}$ i $\mathrm{Ti}\,e=\mathrm{Ti}$ exigeix $e\in\{\mathrm{Pa},\mathrm{Ti}\}$, i la intersecció és buida. És un model finit, comprovable a mà, on fallen alhora (C3), (C4) i (C5).
\end{exemple}

\begin{exemple}[Altres magmes naturals]
\begin{itemize}
\item \emph{El punt mitjà}: $\mathbb R$ amb $ab=\frac{a+b}2$. Per exemple, $(0\cdot0)\cdot4=2$ i $0\cdot(0\cdot4)=1$; no hi ha neutre, perquè $a\cdot e=a$ força $e=a$.
\item \emph{El producte vectorial}: $\mathbb R^3$ amb $a\times b$. $(\mathbf i\times\mathbf i)\times\mathbf j=\mathbf 0$, però $\mathbf i\times(\mathbf i\times\mathbf j)=\mathbf i\times\mathbf k=-\mathbf j$. Satisfà, en lloc de l'associativitat, la identitat de Jacobi: és l'exemple bàsic d'àlgebra de Lie, una estructura perfectament bona que no és una categoria d'un objecte.
\item \emph{La coma flotant} (com a il·lustració): la suma de l'aritmètica de coma flotant d'un ordinador no és, en general, associativa. En doble precisió, una avaluació típica dona $(0.1+0.2)+0.3=0.6000000000000001$, mentre que $0.1+(0.2+0.3)=0.6$. Per convertir-ho en un magma literal caldria precisar el conjunt de valors i el tractament dels casos excepcionals (desbordaments, zeros amb signe...), cosa que no fem.
\item \emph{Associativitat sense neutre}: $(2\mathbb Z,\cdot)$ és un semigrup, però l'únic candidat a neutre seria $1\notin2\mathbb Z$. Aquesta fallada es repara de franc, afegint formalment un neutre.
\end{itemize}
\end{exemple}

\paragraph{Quocientar no sempre és un bon remei.} Davant d'un magma no associatiu podríem intentar el truc dels camins: quocientar per la relació d'equivalència més petita, compatible amb l'operació, que la faci associativa. Però el resultat pot ser catastròfic.

\begin{proposicio}
El quocient associatiu del magma pedra-paper-tisora té un sol element.
\end{proposicio}

\begin{prova}
En el quocient, $\mathrm{Ti}=(\mathrm{Pe}\,\mathrm{Pa})\,\mathrm{Ti}=\mathrm{Pe}\,(\mathrm{Pa}\,\mathrm{Ti})=\mathrm{Pe}$ i $\mathrm{Pe}=(\mathrm{Pa}\,\mathrm{Ti})\,\mathrm{Pe}=\mathrm{Pa}\,(\mathrm{Ti}\,\mathrm{Pe})=\mathrm{Pa}$.
\end{prova}

La diferència amb els camins és significativa. Allà la manca d'associativitat era \emph{coherent}, testimoniada per homotopies canòniques, i el quocient conservava la informació essencial. Aquí és arbitrària, i quocientar ho esborra tot. La jerarquia de les estructures que hem vist és aquesta:
\[
\renewcommand{\arraystretch}{1.2}
\begin{array}{lll}
\text{un sol objecte} & \text{diversos objectes} & \text{condicions}\\\hline
\text{magma} & \text{magmoide} & \text{(C1), (C2)}\\
\text{semigrup} & \text{semicategoria} & \text{(C1)--(C3)}\\
\text{monoide} & \text{categoria} & \text{(C1)--(C5)}\\
\text{grup} & \text{grupoide} & \text{(C1)--(C5) i inversos}
\end{array}
\]

\subsection{Altres maneres de fallar}

\begin{exemple}[Ordre estricte: una semicategoria]
Objectes: els nombres reals; un únic morfisme $x\to y$ si $x<y$, i cap altrament. La composició existeix per transitivitat i és associativa, perquè entre dos objectes hi ha com a molt un morfisme. Però no hi ha cap morfisme $x\to x$: falla (C4). Substituint $<$ per $\le$ s'obté una categoria.
\end{exemple}

\begin{exemple}[Grafs: no hi ha composició]
Un graf dirigit té vèrtexs i arestes, però dues arestes consecutives $x\to y\to z$ no tenen per què donar una aresta $x\to z$: falla (C2). El remei és la categoria lliure del graf (definició~\ref{def:free}).
\end{exemple}

\begin{exemple}[Aplicacions lineals no nul·les]
Objectes: els espais vectorials reals no nuls; morfismes: les aplicacions lineals \emph{no nul·les}. Les identitats hi són, i la composició és associativa quan està definida, però no sempre ho està: a $\mathbb R^2$, les projeccions $p(x,y)=(x,0)$ i $q(x,y)=(0,y)$ són no nul·les i $q\comp p=0$. Falla (C2).
\end{exemple}

\begin{exemple}[Funcions identificades amb el seu graf]
Si una funció és només el seu graf, $\{(1,1)\}$ és alhora una funció $\{1\}\to\{1\}$ i una funció $\{1\}\to\{1,2\}$: el codomini no queda determinat i falla (C1). En particular, \enquote{ser exhaustiva} deixa de ser una propietat de la funció. És el mateix fenomen que l'observació~\ref{obs:extrems-implicits} assenyala per a $\cat{Rel}$, i el remei és el mateix: el domini i el codomini formen part de les dades del morfisme.
\end{exemple}

\begin{exemple}[Un quocient que no és compatible amb la composició]
Al monoide de les funcions contínues $\mathbb R\to\mathbb R$ amb la composició, identifiquem $f\sim g$ si $f(0)=g(0)$, i intentem definir $[g]\comp[f]=[g\comp f]$. Amb $f(x)=x+1$, $g(x)=x$ i $g'(x)=2x$ tenim $g\sim g'$, però $(g\comp f)(0)=1\neq2=(g'\comp f)(0)$. La composició no està ben definida i falla (C2). Per això, en construir $\Pi_1(X)$, calia comprovar que l'homotopia és compatible amb la concatenació: no tot quocient d'una categoria és una categoria.
\end{exemple}

\subsection{Falsos contraexemples}

Per equilibrar, convé recordar estructures que \emph{semblen} fallar algun axioma però que són categories genuïnes. Ensenyen que els morfismes no han de ser funcions, ni estar definits a tot arreu.

\begin{exemple}[Funcions parcials]
Objectes: els conjunts; morfismes $f\colon A\rightharpoonup B$: les funcions definides en un subconjunt $D_f\subseteq A$. La composició $g\comp f$ és definida en $\{a\in D_f\mid f(a)\in D_g\}$. Les dues maneres d'agrupar $h$, $g$ i $f$ tenen el mateix domini de definició, $\{a\in D_f\mid f(a)\in D_g,\ g(f(a))\in D_h\}$, i hi prenen el mateix valor; les identitats són les funcions totals identitat. Obtenim la categoria $\cat{Par}$: que un morfisme no estigui definit a tot arreu no impedeix res.
\end{exemple}

Les relacions (la categoria $\cat{Rel}$ del capítol~\ref{cap:categories}) i les matrius ($\cat{Mat}_{\mathbb K}$, exemple~\ref{ex:mat-finvect}) són dos altres casos: els morfismes no són funcions, i en el segon cas ni tan sols són transformacions de res, sinó taules de nombres.

\subsection{Resum}

\[
\renewcommand{\arraystretch}{1.15}
\begin{array}{lccccl}
\text{Exemple} & \text{(C1)} & \text{(C2)} & \text{(C3)} & \text{(C4--5)} & \text{Remei}\\\hline
\text{camins sobre }[0,1] & \text{sí} & \text{sí} & \text{no} & \text{no} & \Pi_1(X)\text{ o camins de Moore}\\
(\mathbb Z,-) & \text{sí} & \text{sí} & \text{no} & \text{a mitges} & \text{---}\\
\text{pedra-paper-tisora} & \text{sí} & \text{sí} & \text{no} & \text{no} & \text{(el quocient és trivial)}\\
(2\mathbb Z,\cdot) & \text{sí} & \text{sí} & \text{sí} & \text{no} & \text{afegir un neutre}\\
(\mathbb R,<) & \text{sí} & \text{sí} & \text{sí} & \text{no} & (\mathbb R,\le)\\
\text{graf dirigit} & \text{sí} & \text{no} & \text{---} & \text{---} & \text{categoria lliure}\\
\text{lineals no nul·les} & \text{sí} & \text{no} & \text{sí} & \text{sí} & \text{admetre el }0\\
\text{funcions com a grafs} & \text{no} & \text{sí} & \text{sí} & \text{sí} & \text{morfismes amb extrems}
\end{array}
\]
Els axiomes de categoria no són convencions arbitràries: n'hi ha prou d'agafar l'estructura més natural possible ---els camins en un espai--- per veure'ls fallar. Quan fallen, hi ha tres actituds possibles: quocientar, estrictificar o acceptar la feblesa i fer-ne teoria.

% =====================================================================
\section{Què no és un functor}
\label{sec:no-functor}
% =====================================================================

Un functor ha de superar dues proves, en aquest ordre (observació~\ref{obs:comprovar-tipus}): els tipus i, després, les equacions $F(\id_A)=\id_{F(A)}$ i $F(g\comp f)=F(g)\comp F(f)$. Els gairebé-functors fallen en una de les dues.

\paragraph{Fallada de tipus.} El capítol~\ref{cap:functors} ja en dona l'exemple bàsic: l'antiimatge $f\mapsto f^{-1}$ no pot ser un functor covariant de $\cat{Set}$ en $\cat{Set}$, perquè va en sentit contrari; sí que ho és com a functor contravariant (exemple~\ref{ex:parts-dues-variancies}). Un altre exemple: l'assignació que envia cada grup $G$ al seu grup d'automorfismes $\mathrm{Aut}(G)$ no té cap manera natural d'actuar sobre els homomorfismes. Un homomorfisme $f\colon G\to H$ no indueix cap aplicació raonable $\mathrm{Aut}(G)\to\mathrm{Aut}(H)$; només ho fa quan $f$ és un isomorfisme, que és la manera de convertir-la en un functor sobre la subcategoria dels isomorfismes.

\paragraph{Una assignació sobre objectes que no admet cap functor.} El cas més instructiu és el d'un objecte que sembla \enquote{canònic} però no és functorial de cap manera.

\begin{proposicio}[El centre d'un grup no és functorial]\label{prop:centre-no-functor}
No hi ha cap functor $Z\colon\cat{Grp}\to\cat{Grp}$ que enviï cada grup $G$ al seu centre $Z(G)=\{z\in G\mid zg=gz \text{ per a tot } g\}$, sigui quina sigui la seva acció sobre els homomorfismes.
\end{proposicio}

\begin{prova}
Sigui $C_2=\{1,c\}$ el grup de dos elements i $S_3$ el grup de permutacions de tres elements. Considerem la inclusió $i\colon C_2\to S_3$, que envia $c$ a la transposició $(1\,2)$, i el signe $s\colon S_3\to C_2$, que envia les transposicions a $c$. Aleshores $s\comp i=\id_{C_2}$. Com que $C_2$ és commutatiu, $Z(C_2)=C_2$. En canvi, $Z(S_3)=\{1\}$: les tres transposicions no commuten entre si (per exemple, $(1\,2)(1\,3)=(1\,3\,2)$, mentre que $(1\,3)(1\,2)=(1\,2\,3)$), i les dues permutacions cícliques $(1\,2\,3)$ i $(1\,3\,2)$ no commuten amb $(1\,2)$, perquè $(1\,2)(1\,2\,3)(1\,2)=(1\,3\,2)$. Si $Z$ fos un functor,
\[
\id_{C_2}=Z(\id_{C_2})=Z(s\comp i)=Z(s)\comp Z(i),
\]
amb $Z(i)\colon C_2\to\{1\}$ i $Z(s)\colon\{1\}\to C_2$. Però una composició que passa per un grup d'un element és constant, i no pot ser la identitat d'un grup de dos elements.
\end{prova}

L'argument és un patró general: si $s\comp i=\id$, aleshores $F(s)\comp F(i)=\id$, de manera que tot functor converteix les retraccions en retraccions. Una assignació que envia un objecte a un de \enquote{massa petit} per factoritzar-hi una identitat no pot ser functorial. En canvi, assignacions semblants que només fan servir propietats preservades pels homomorfismes sí que són functors (vegeu els exercicis).

\paragraph{Fallada de les equacions.} Considerem el monoide de les funcions derivables $\mathbb R\to\mathbb R$ amb la composició, i l'assignació $f\mapsto f'$ cap al mateix monoide. Els tipus són correctes, però les equacions fallen. La identitat va a la funció constant $1$, que no és la identitat. I, amb $f(x)=g(x)=2x$, tenim $(g\comp f)'=4$ (constant), mentre que $g'\comp f'=2$ (constant). La regla de la cadena explica què cal canviar:

\begin{exemple}[La regla de la cadena és una llei functorial]\label{ex:regla-cadena}
Sigui $\cat{Der}$ la categoria que té per objectes els nombres reals i per morfismes $a\to b$ les funcions derivables $f\colon\mathbb R\to\mathbb R$ amb $f(a)=b$, amb la composició de funcions. Sigui $\cat{B}(\mathbb R,\cdot)$ el monoide multiplicatiu dels reals vist com a categoria d'un objecte. L'assignació
\[
D\colon\cat{Der}\to\cat{B}(\mathbb R,\cdot),\qquad (f\colon a\to b)\ \longmapsto\ f'(a)
\]
és un functor. Les identitats van a $1$, perquè $\id'(a)=1$, i la regla de la cadena, $(g\comp f)'(a)=g'(f(a))\cdot f'(a)$, diu exactament que $D(g\comp f)=D(g)\comp D(f)$, ja que $f(a)$ és el domini de $g$. El que fallava era el lloc on es mira la derivada: el functor no és $f\mapsto f'$, sinó $f\mapsto f'$ \emph{al punt de partida}.
\end{exemple}

\paragraph{Functors llevat d'isomorfisme (informatiu).} Hi ha assignacions que són functorials només llevat d'isomorfisme, exactament com els camins eren associatius llevat d'homotopia. L'exemple més important del llibre és la reindexació de l'observació~\ref{obs:reindexacio}. Per a cada $f\colon A\to B$, el pullback dona $f^{*}\colon\cat{C}/B\to\cat{C}/A$, però, un cop triats els pullbacks, $(g\comp f)^{*}$ i $f^{*}\comp g^{*}$ no són iguals, sinó canònicament isomorfs. A $\cat{Set}$, amb els pullbacks triats com a conjunts de parelles, $(g\comp f)^{*}$ envia $p\colon E\to C$ al conjunt de les parelles $(a,e)$ amb $g(f(a))=p(e)$, mentre que $f^{*}g^{*}$ l'envia al conjunt de les parelles $\bigl(a,(b,e)\bigr)$ amb $f(a)=b$ i $g(b)=p(e)$: són conjunts diferents, en bijecció canònica. Una tal assignació s'anomena \emph{pseudofunctor}. Al capítol~\ref{cap:subobjectes} el problema desapareix perquè $\Sub$ treballa amb classes de monomorfismes, i l'isomorfisme canònic queda absorbit (proposició~\ref{prop:functor-sub}): és el primer remei de la secció anterior, el quocient.

% =====================================================================
\section{Què no és una transformació natural}
\label{sec:no-natural}
% =====================================================================

El llibre ja dona els contraexemples essencials, i aquí només els ordenem. Una família de components $\eta_A\colon F(A)\to G(A)$ pot deixar de ser natural per tres motius diferents:
\begin{enumerate}
\item \emph{Les components s'han triat objecte per objecte}, sense mirar els morfismes. L'exemple~\ref{ex:no-natural} té totes les components bijectives i, tanmateix, un sol morfisme $\{0,1\}\to\{0,1\}$ trenca el quadrat.
\item \emph{La naturalitat només val per a una part dels morfismes.} L'isomorfisme entre un espai de dimensió finita i el seu dual depèn d'una base, i ni tan sols es pot plantejar com a transformació natural entre functors paral·lels, perquè la identitat és covariant i el dual contravariant. Només s'hi pot comparar restringint-se als isomorfismes, i fins i tot així no és natural (observació~\ref{obs:iso-no-natural} i exercici resolt~\ref{ex:dual-no-natural} de la Pràctica del capítol~\ref{cap:transformacions}). En canvi, $V\cong V^{**}$ és natural (exemple~\ref{ex:bidual}).
\item \emph{No hi ha cap component possible.} Amb el grup de dos elements actuant trivialment i per la transposició sobre $\{0,1\}$, no hi ha cap transformació natural entre els dos functors, ni tan sols una que no sigui un isomorfisme (observació~\ref{obs:iso-no-natural}).
\end{enumerate}
La lliçó comuna és la de la secció següent: tenir objectes isomorfs \emph{un per un} no és el mateix que tenir functors isomorfs.

% =====================================================================
\section{Igualtat, isomorfisme, isomorfisme natural i equivalència}
\label{sec:igualtat-equivalencia}
% =====================================================================

\subsection{La jerarquia}

Per a cada mena d'entitat hi ha una noció \emph{estricta} de \enquote{ser el mateix}, la igualtat, i una noció \emph{estructural}, que és la que respecta la teoria:
\[
\renewcommand{\arraystretch}{1.25}
\begin{array}{lll}
\text{Què comparem} & \text{Noció estricta} & \text{Noció estructural}\\\hline
\text{dos objectes d'una categoria }\cat{C} & A=B & A\cong B\ \text{(isomorfisme a }\cat{C})\\
\text{dos functors }F,G\colon\cat{C}\to\cat{D} & F=G & F\cong G\ \text{(isomorfisme natural)}\\
\text{dues categories} & \cat{C}=\cat{D} & \cat{C}\simeq\cat{D}\ \text{(equivalència)}
\end{array}
\]
Entre les categories hi ha, a més, un nivell intermedi: l'isomorfisme de categories $\cat{C}\cong\cat{D}$ (definició~\ref{def:iso-categories}). Les tres files tenen la mateixa forma. Un isomorfisme natural és un isomorfisme a la categoria de functors $\cat{D}^{\cat{C}}$, i per tant la segona fila és un cas de la primera. L'equivalència és el que s'obté quan, a la definició d'isomorfisme de categories, se substitueixen les igualtats $G\comp F=\id$ i $F\comp G=\id$ per isomorfismes naturals $G\comp F\cong\id$ i $F\comp G\cong\id$, perquè la manera correcta de comparar functors és l'isomorfisme natural. Cada nivell fa servir la noció estructural del nivell anterior.

Totes les implicacions possibles entre aquestes nocions van en un sol sentit:
\[
\begin{gathered}
A=B\ \Rightarrow\ A\cong B,\\
F=G\ \Rightarrow\ F\cong G\ \Rightarrow\ F(A)\cong G(A)\ \text{per a tot }A,\\
\cat{C}=\cat{D}\ \Rightarrow\ \cat{C}\cong\cat{D}\ \Rightarrow\ \cat{C}\simeq\cat{D},
\end{gathered}
\]
i cap no és reversible. Vegem-ne un contraexemple per a cada pas.

\subsection{Un contraexemple per a cada pas}

\begin{itemize}
\item \emph{Isomorfs però no iguals.} A $\cat{Set}$, $\{0,1\}\cong\{a,b\}$. A $\cat{Prov}_T$, $p\land q\cong q\land p$, que són fórmules diferents (exemple~\ref{ex:iso-equivalencia-logica}).
\item \emph{Naturalment isomorfs però no iguals.} A la categoria dels espais vectorials de dimensió finita, la identitat i el bidual són naturalment isomorfs (exemple~\ref{ex:bidual}), però no són iguals: $V^{**}$ no és $V$, sinó un espai format per funcionals sobre $V^{*}$.
\item \emph{Isomorfs objecte per objecte, però no naturalment isomorfs.} Els dos functors del grup de dos elements en $\cat{Set}$ de l'observació~\ref{obs:iso-no-natural} envien l'únic objecte al mateix conjunt $\{0,1\}$, però no són naturalment isomorfs.
\item \emph{Isomorfes però no iguals.} Una categoria i una còpia seva amb els objectes reanomenats.
\item \emph{Equivalents però no isomorfes.} $\cat{Mat}_{\mathbb K}\simeq\cat{FinVect}_{\mathbb K}$ (exemple~\ref{ex:mat-finvect}), o la categoria $\cat{1}$ i la categoria amb dos objectes isomorfs i un sol morfisme entre cada parell d'objectes.
\end{itemize}

\subsection{La invariància estructural}

El principi que governa tot el llibre és aquest: \emph{una propietat formulada només amb l'estructura ---composició, identitats, propietats universals---, sense afirmar mai que dos objectes són iguals, es transporta per la noció estructural corresponent.} Per als objectes ho diu l'observació~\ref{obs:invariancia}. Per a les categories, la peça clau és que les equivalències conserven les propietats universals.

\begin{teorema}\label{teo:equivalencia-limits}
Sigui $F\colon\cat{C}\to\cat{D}$ una equivalència. Si $(L,\lambda)$ és un límit d'un diagrama $D\colon\cat{I}\to\cat{C}$, aleshores $(F(L),F\lambda)$ és un límit de $F\comp D$. Dualment, $F$ preserva els colímits.
\end{teorema}

\begin{prova}
Pel teorema~\ref{teo:caract-equivalencia}, $F$ és ple, fidel i essencialment exhaustiu. La família $F\lambda=(F(\lambda_i))_i$ és un con sobre $F\comp D$, perquè $F$ preserva la composició. Sigui $(X,\alpha)$ un con sobre $F\comp D$. Com que $F$ és essencialment exhaustiu, hi ha un objecte $C$ de $\cat{C}$ i un isomorfisme $e\colon F(C)\to X$. Com que $F$ és ple i fidel, per a cada $i$ hi ha un únic $\beta_i\colon C\to D_i$ amb $F(\beta_i)=\alpha_i\comp e$. Aquests $\beta_i$ formen un con: per a $u\colon i\to j$,
\[
F(D_u\comp\beta_i)=F(D_u)\comp\alpha_i\comp e=\alpha_j\comp e=F(\beta_j),
\]
i, com que $F$ és fidel, $D_u\comp\beta_i=\beta_j$. Per la propietat del límit, hi ha un únic $h\colon C\to L$ amb $\lambda_i\comp h=\beta_i$. Posem $k=F(h)\comp e^{-1}\colon X\to F(L)$. Aleshores $F(\lambda_i)\comp k=F(\beta_i)\comp e^{-1}=\alpha_i$: $k$ factoritza el con.

Per a la unicitat, sigui $k'\colon X\to F(L)$ amb $F(\lambda_i)\comp k'=\alpha_i$ per a tot $i$. Com que $F$ és ple, $k'\comp e=F(h')$ per a algun $h'\colon C\to L$, i aleshores $F(\lambda_i\comp h')=\alpha_i\comp e=F(\beta_i)$. Com que $F$ és fidel, $\lambda_i\comp h'=\beta_i$, i per la unicitat del límit, $h'=h$. Per tant, $k'=F(h)\comp e^{-1}=k$. Els colímits es tracten igual, invertint totes les fletxes.
\end{prova}

D'aquest teorema, de l'exercici resolt~\ref{pr-equivalencia-preserva} de la Pràctica del capítol~\ref{cap:transformacions} (que una equivalència preserva i reflecteix isomorfismes, monomorfismes i epimorfismes) i del fet que el quasiinvers també és una equivalència, se segueix que les propietats de la columna esquerra de la taula següent les té $\cat{C}$ si i només si les té qualsevol categoria equivalent a $\cat{C}$. Per a les dues últimes (cartesiana tancada, topos), el mateix argument s'aplica a les propietats universals de l'exponencial i del classificador; no el reproduïm.
\begin{center}
\small
\renewcommand{\arraystretch}{1.2}
\begin{tabularx}{\linewidth}{@{}>{\raggedright\arraybackslash}X>{\raggedright\arraybackslash}X@{}}
\toprule
\textbf{Invariants per equivalència} & \textbf{No invariants per equivalència}\\
\midrule
tenir objecte terminal o inicial & el nombre d'objectes\\
tenir productes, equalitzadors, límits finits & tenir un sol objecte\\
ser un preordre (com a molt un morfisme entre dos objectes) & ser un ordre parcial\\
ser un grupoide (tot morfisme és invertible) & ser esquelètica\\
que tot morfisme mono i epi sigui un isomorfisme & que dos objectes concrets siguin iguals\\
ser cartesiana tancada; ser un topos & ser petita (\enquote{essencialment petita} sí que és invariant)\\
\bottomrule
\end{tabularx}
\end{center}
Les propietats de la columna de la dreta fan servir, d'una manera o altra, la igualtat d'objectes, i per això no són categòriques. Per exemple, $\cat{1}$ és un ordre parcial, i la categoria amb dos objectes isomorfs, que li és equivalent, no ho és. Quan una definició depèn de la igualtat d'objectes, convé sospitar-ne.

\subsection{Equivalència i esquelets}

El resultat següent fa veure que l'equivalència és exactament l'isomorfisme \enquote{llevat de les còpies isomorfes}.

\begin{lema}\label{lema:esqueletiques}
Una equivalència entre dues categories esquelètiques és un isomorfisme de categories.
\end{lema}

\begin{prova}
Sigui $F\colon\cat{C}\to\cat{D}$ una equivalència entre categories esquelètiques; és ple, fidel i essencialment exhaustiu. \emph{És injectiu sobre objectes}: si $F(A)=F(B)$, la identitat de $F(A)$ és $F(u)$ per a algun $u\colon A\to B$, perquè $F$ és ple, i $u$ és un isomorfisme, perquè un functor ple i fidel reflecteix els isomorfismes (proposició~\ref{prop:plenfidel-reflecteix-iso}). Com que $\cat{C}$ és esquelètica, $A=B$. \emph{És exhaustiu sobre objectes}: per a cada $D$ hi ha $C$ amb $D\cong F(C)$, i com que $\cat{D}$ és esquelètica, $D=F(C)$. A més, $F$ és bijectiu sobre cada homconjunt, perquè és ple i fidel. Per tant, $F$ és bijectiu sobre objectes i sobre morfismes, i l'assignació inversa és un functor (preserva identitats i composicions perquè $F$ les preserva). És, doncs, un isomorfisme de categories.
\end{prova}

\begin{corollari}\label{cor:equivalencia-esquelets}
Dues categories són equivalents si i només si tenen esquelets isomorfs.
\end{corollari}

\begin{prova}
En el marc d'elecció del llibre, tota categoria té esquelets (definició~\ref{def:esquelet}), i és equivalent a qualsevol d'ells: la inclusió de l'esquelet és plena i fidel per ser una subcategoria plena, i essencialment exhaustiva perquè l'esquelet conté un objecte de cada classe d'isomorfisme. Les equivalències es componen, perquè la composició de functors plens, fidels i essencialment exhaustius ho torna a ser. Si $\cat{C}\simeq\cat{D}$, els seus esquelets són equivalents, i pel lema~\ref{lema:esqueletiques} són isomorfs. Recíprocament, si els esquelets són isomorfs, $\cat{C}\simeq\mathrm{esquelet}(\cat{C})\cong\mathrm{esquelet}(\cat{D})\simeq\cat{D}$.
\end{prova}

Per als preordres, l'esquelet és l'ordre parcial que s'obté identificant els elements equivalents ($x\le y$ i $y\le x$): dos preordres són equivalents si i només si aquests ordres parcials són isomorfs.

\subsection{Com es demostra que dues categories no són equivalents}

Per demostrar que $\cat{C}\not\simeq\cat{D}$ no es poden revisar tots els functors possibles. La via és buscar una propietat invariant per equivalència que l'una tingui i l'altra no.

\begin{exemple}[$\cat{Set}\not\simeq\cat{Set}\op$]
Diem que un objecte inicial $0$ és \emph{estricte} si tot morfisme $X\to0$ és un isomorfisme. Tenir un objecte inicial estricte és invariant per equivalència. Si $F\colon\cat{C}\to\cat{D}$ és una equivalència, $F(0)$ és inicial (teorema~\ref{teo:equivalencia-limits}). Donat un morfisme $k\colon Y\to F(0)$, prenem un isomorfisme $e\colon F(C)\to Y$; com que $F$ és ple, $k\comp e=F(u)$ per a algun $u\colon C\to0$, que és un isomorfisme, i aleshores $F(u)$ i $k=F(u)\comp e^{-1}$ també ho són. A $\cat{Set}$, l'objecte inicial $\varnothing$ és estricte: només hi ha aplicacions $X\to\varnothing$ si $X=\varnothing$, i aleshores és la identitat. A $\cat{Set}\op$, l'objecte inicial és un conjunt d'un element $\{\ast\}$, i un morfisme $X\to\{\ast\}$ de $\cat{Set}\op$ és una aplicació $\{\ast\}\to X$ de $\cat{Set}$. Amb $X=\{0,1\}$ no és bijectiva, i per tant no és un isomorfisme. Així, $\cat{Set}$ no és equivalent a la seva oposada.
\end{exemple}

\begin{exemple}[$\cat{Grp}\not\simeq\cat{Ab}$]
La propietat \enquote{per a dos objectes qualssevol, el producte i el coproducte són isomorfs} és invariant per equivalència, perquè una equivalència preserva productes, coproductes i isomorfismes. 
\emph{A $\cat{Ab}$ es compleix.} Siguin $A$ i $B$ grups abelians, escrits additivament. El producte $A\times B$, amb les inclusions $a\mapsto(a,0)$ i $b\mapsto(0,b)$, és també un coproducte. Donats homomorfismes $f\colon A\to G$ i $g\colon B\to G$ cap a un grup abelià $G$, l'aplicació $(a,b)\mapsto f(a)+g(b)$ és un homomorfisme, precisament perquè $G$ és commutatiu, i restringida a les inclusions dona $f$ i $g$. És l'única, perquè $(a,b)=(a,0)+(0,b)$.

\emph{A $\cat{Grp}$ no es compleix.} Sigui $C_2=\{1,c\}$ i suposem que el coproducte $P=C_2+C_2$ és isomorf al producte $C_2\times C_2$, que és commutatiu. Els homomorfismes $C_2\to S_3$ que envien $c$ a $(1\,2)$ i a $(1\,3)$ indueixen un homomorfisme $P\to S_3$ la imatge del qual conté $(1\,2)$ i $(1\,3)$. Però la imatge d'un grup commutatiu és commutativa, i $(1\,2)$ i $(1\,3)$ no commuten. Per tant, $C_2+C_2\not\cong C_2\times C_2$.
\end{exemple}

\begin{exemple}[$\cat{FinSet}\not\simeq\cat{Set}$]
\enquote{Tota família numerable d'objectes té coproducte} és invariant per equivalència. $\cat{Set}$ la compleix i $\cat{FinSet}$ no (observació~\ref{obs:elemental-prefeixos}).
\end{exemple}

\subsection{Què té a veure tot això amb la lògica}

A la categoria $\cat{Prov}_T$ d'un sistema deductiu, els tres nivells tenen una lectura lògica precisa:
\begin{itemize}
\item la \emph{igualtat} de dues fórmules és la igualtat sintàctica: $p\land q$ i $q\land p$ són fórmules diferents;
\item l'\emph{isomorfisme} és la interdeduïbilitat, $p\vdash_T q$ i $q\vdash_T p$ (exemple~\ref{ex:iso-equivalencia-logica});
\item l'\emph{esquelet} de $\cat{Prov}_T$ és isomorf a l'ordre parcial de les classes de fórmules interdeduïbles, que en lògica s'anomena \textbf{àlgebra de Lindenbaum--Tarski}\index{àlgebra de Lindenbaum--Tarski} de $T$. Pel corol·lari~\ref{cor:equivalencia-esquelets}, $\cat{Prov}_T$ és equivalent, però en general no isomorf, a aquesta àlgebra.
\end{itemize}
En conseqüència, qualsevol propietat categòrica de $\cat{Prov}_T$ és una propietat de l'àlgebra de Lindenbaum--Tarski: tenir productes (la conjunció), ser cartesiana tancada (la implicació), ser una àlgebra de Heyting o de Boole. La lògica estudia les teories \emph{llevat d'interdeduïbilitat}, i la teoria de categories hi aporta el llenguatge que fa aquesta actitud sistemàtica: el que compta d'una teoria és la seva estructura, no la tria particular de fórmules.

\subsection{Errors típics}

\begin{itemize}
\item Escriure $A=B$ quan només es té $A\cong B$, encara que l'isomorfisme sigui canònic.
\item Deduir $F\cong G$ de $F(A)\cong G(A)$ per a cada $A$, sense comprovar la naturalitat.
\item Creure que una equivalència és un isomorfisme de categories, o que $G\comp F$ ha de ser literalment la identitat.
\item Demostrar que dues categories no són equivalents amb un argument que fa servir la igualtat d'objectes, com comptar-ne el nombre.
\item Definir una propietat d'una categoria que depèn d'objectes concrets, i després aplicar-la a una categoria equivalent.
\end{itemize}

% =====================================================================
\section{La lògica fora dels topos}
\label{sec:fora-topos}
% =====================================================================

Al llarg dels capítols~\ref{cap:cartesianes}--\ref{cap:pont}, la lògica s'ha llegit d'una manera determinada, que comença amb el tancament cartesià: una proposició sobre $E$ és un subobjecte de $E$, la conjunció és un pullback, la implicació és l'exponencial (l'adjunt per la dreta de $-\wedge U$), els quantificadors són adjunts de la reindexació, i en un topos tot això és la lògica interna de $\Omega$. Hi ha lògiques i models perfectament legítims que no encaixen en aquest esquema. Els presentem de manera sobretot informativa, però en cada cas diem amb precisió quina condició falla i, quan és fàcil, ho comprovem.

\subsection{Categories que no són cartesianes tancades}
\label{subsec:no-ccc}

Abans de mirar què hi ha entre les categories cartesianes tancades i els topos, convé veure que moltes categories molt naturals no són ni tan sols cartesianes tancades. Per demostrar-ho, n'hi ha prou de trobar una propietat que tota categoria cartesiana tancada tingui i que la categoria en qüestió no tingui. La més útil és la que dona el teorema de preservació de colímits pels adjunts per l'esquerra (capítol~\ref{cap:adjuncions}): en una categoria cartesiana tancada, $-\times B$ és adjunt per l'esquerra de $(-)^B$, i per tant preserva tots els colímits que existeixin. En particular, si $0$ és inicial, $0\times B$ també ho és.

Un \textbf{objecte zero}\index{objecte zero} és un objecte que és alhora inicial i terminal.

\begin{proposicio}\label{prop:zero-no-ccc}
Si una categoria cartesiana tancada té un objecte zero $0$, aleshores tots els seus objectes són isomorfs a $0$.
\end{proposicio}

\begin{prova}
Sigui $B$ un objecte. Com que $0$ és inicial i $-\times B$ preserva l'objecte inicial, $0\times B$ és inicial, i per tant $0\times B\cong0$. D'altra banda, com que $0$ és terminal, la projecció $\pi_2\colon0\times B\to B$ és un isomorfisme. En efecte, sigui $!_B\colon B\to0$ l'únic morfisme cap al terminal, i $s=\langle !_B,\id_B\rangle\colon B\to0\times B$. Aleshores $\pi_2\comp s=\id_B$, i
\[
s\comp\pi_2=\langle !_B\comp\pi_2,\ \pi_2\rangle=\langle\pi_1,\pi_2\rangle=\id_{0\times B},
\]
perquè $!_B\comp\pi_2$ i $\pi_1$ són dos morfismes $0\times B\to0$, que han de coincidir perquè $0$ és terminal. Per tant, $B\cong0\times B\cong0$.
\end{prova}

\begin{corollari}\label{cor:zero-no-ccc}
Les categories $\cat{Vect}_{\mathbb K}$, $\cat{Ab}$, $\cat{Grp}$, $\cat{Mon}$ i $\cat{Rel}$ no són cartesianes tancades.
\end{corollari}

\begin{prova}
Totes tenen un objecte zero: l'espai $\{0\}$; el grup, o el monoide, d'un sol element, perquè de l'element neutre surt un únic homomorfisme cap a qualsevol altre grup o monoide i cap a l'estructura trivial només n'hi ha un; i, a $\cat{Rel}$, el conjunt buit, perquè l'única relació $\varnothing\to X$ i l'única relació $X\to\varnothing$ són la relació buida. Totes tenen, però, objectes que no són isomorfs al zero: $\mathbb K$, el grup $\mathbb Z$, el monoide $(\mathbb N,+)$ o qualsevol conjunt no buit. Per la proposició~\ref{prop:zero-no-ccc}, cap no és cartesiana tancada.
\end{prova}

\paragraph{Preordres.} En un preordre amb ínfims finits, ser cartesià tancat vol dir tenir una implicació $b\Rightarrow c$ amb $a\wedge b\le c$ si i només si $a\le b\Rightarrow c$ (capítol~\ref{cap:cartesianes}). Si, a més, té suprems binaris, aleshores és distributiu, perquè $-\wedge b$ és adjunt per l'esquerra i preserva els suprems. Per tant, \emph{un reticle no distributiu no és cartesià tancat}. El reticle dels subespais de $\mathbb R^2$ que veurem tot seguit n'és un exemple.

\paragraph{Categories que ni tan sols són cartesianes.} Un monoide $M$ no trivial, vist com a categoria $\cat{B}M$, no té objecte terminal: l'únic objecte $\ast$ ho seria només si $\Hom(\ast,\ast)=M$ tingués un sol element. Un preordre amb dos elements incomparables que no tinguin cap fita inferior comuna no té el producte d'aquests dos elements. Cap d'aquestes categories no pot ser cartesiana tancada, perquè no té productes finits.

\paragraph{Els espais topològics (informatiu).} $\cat{Top}$ tampoc no és cartesiana tancada: hi ha quocients d'espais que el producte amb un espai fix no conserva. La demostració demana més topologia de la que pressuposa el llibre, i no la fem.

\subsection{Cartesianes tancades que no són topos}

La lògica de la implicació necessita només el tancament cartesià, i això és molt menys que un topos.

\begin{exemple}[$\cat{Cat}$ no és un topos]
La categoria $\cat{Cat}$ és cartesiana tancada (exemple~\ref{ex:cat-ccc}), però no té classificador de subobjectes. Sigui $\cat{2}$ la categoria $0\to1$, i $\cat{I}$ la categoria amb dos objectes i, a més de les identitats, dos morfismes $u\colon0\to1$ i $u^{-1}\colon1\to0$ inversos l'un de l'altre. La inclusió $i\colon\cat{2}\to\cat{I}$ és:
\begin{itemize}
\item \emph{mono}: és injectiva sobre objectes i morfismes, de manera que si $i\comp H=i\comp K$, aleshores $H=K$;
\item \emph{epi}: si $H,K\colon\cat{I}\to\cat{C}$ coincideixen sobre $u$, aleshores també coincideixen sobre $u^{-1}$, perquè un functor envia l'invers de $u$ a l'invers de la imatge de $u$, i l'invers és únic;
\item \emph{no és un isomorfisme}: $u^{-1}$ no és a la imatge.
\end{itemize}
En una categoria amb classificador, tot morfisme mono i epi és un isomorfisme (exercici resolt~\ref{pr-mono-equalitzador} de la Pràctica del capítol~\ref{cap:topos}). Per tant, $\cat{Cat}$ no en té cap.
\end{exemple}

L'ordre parcial $\mathbb Z\cup\{+\infty\}$ de l'observació~\ref{obs:ccc-no-heyting} és un altre cas: té conjunció i implicació, però no té falsedat.

\subsection{Lògiques on falla la distributivitat}

En una àlgebra de Heyting, $a\wedge(b\vee c)=(a\wedge b)\vee(a\wedge c)$, perquè $a\wedge-$ és adjunt per l'esquerra i preserva els suprems (capítol~\ref{cap:cartesianes}). Com que els subobjectes d'un objecte d'un topos formen una àlgebra de Heyting (nota~\ref{nota:sub-heyting}), cap reticle no distributiu pot ser-ne un.

\begin{exemple}[La lògica quàntica]
Siguin els subespais vectorials de $\mathbb R^2$, ordenats per inclusió: $\{0\}$, les rectes que passen per l'origen i $\mathbb R^2$. L'ínfim és la intersecció i el suprem és la suma. Si $a$, $b$ i $c$ són tres rectes diferents,
\[
a\wedge(b\vee c)=a\wedge\mathbb R^2=a,
\qquad
(a\wedge b)\vee(a\wedge c)=\{0\}\vee\{0\}=\{0\}.
\]
La distributivitat falla. En mecànica quàntica, les propietats observables d'un sistema es modelen amb reticles d'aquest tipus (de subespais tancats), i la \enquote{lògica quàntica} de Birkhoff i von Neumann és la lògica d'aquests reticles no distributius (informatiu). Amb aquestes operacions, aquest reticle no pot ser el reticle de subobjectes de cap objecte d'un topos.
\end{exemple}

\subsection{Lògiques amb una altra conjunció}

\begin{exemple}[La lògica de Łukasiewicz]
Els valors de veritat són els nombres de $[0,1]$, amb la conjunció forta i la implicació
\[
x\otimes y=\max(0,\,x+y-1),\qquad y\to z=\min(1,\,1-y+z).
\]
La implicació és adjunta d'aquesta conjunció, i no de l'ínfim. En efecte, si $z\ge0$, $\max(0,x+y-1)\le z$ si i només si $x\le1-y+z$, és a dir, $x\le y\to z$. En canvi, amb l'ínfim falla: per a $y=0.5$, $z=0.4$ i $x=0.6$, tenim $\min(x,y)=0.5>0.4$, però $x=0.6\le0.9=y\to z$. A més, $\otimes$ no és idempotent: $0.5\otimes0.5=0\neq0.5$. Lògicament, això vol dir que de $A$ no se'n dedueix $A\otimes A$: una hipòtesi no es pot fer servir dues vegades. Categòricament, $([0,1],\otimes,\to)$ és un preordre \emph{monoidal tancat} que no és cartesià tancat respecte de $\otimes$. Aquesta idea ---hipòtesis que són recursos i no es poden duplicar--- és la de la \emph{lògica lineal} de Girard; els seus fragments més simples es modelen en categories monoidals tancades, i la resta demana estructura addicional (informatiu).
\end{exemple}

\subsection{Lògiques on la contradicció no és el fals}

\begin{exemple}[Els tancats d'un espai]
A $\mathbb R$, els subconjunts tancats, ordenats per inclusió, formen un reticle distributiu on el \enquote{complement} natural d'un tancat $A$ és l'adherència del seu complementari, $\neg A=\overline{\mathbb R\setminus A}$. Per a $A=[0,1]$,
\[
A\wedge\neg A=[0,1]\cap\bigl((-\infty,0]\cup[1,\infty)\bigr)=\{0,1\}\neq\varnothing:
\]
una proposició i la seva negació són certes alhora a la frontera. Aquesta estructura és la duala d'una àlgebra de Heyting (una àlgebra de \emph{co-Heyting}), i serveix de model a lògiques \emph{paraconsistents}, on una contradicció no implica qualsevol cosa. Amb aquestes operacions, no pot ser el reticle de subobjectes de cap objecte d'un topos, perquè hi falla $A\wedge\neg A=\bot$.
\end{exemple}

\subsection{Lògiques amb modalitats}

Una \emph{lògica modal} afegeix operadors com \enquote{necessàriament} ($\Box$) i \enquote{possiblement} ($\Diamond$). En un model de Kripke, hi ha un conjunt de mons $W$ i una relació d'accessibilitat $R\subseteq W\times W$, i per a $A\subseteq W$ es defineix
\[
\Box A=\{x\mid \text{per a tot } y,\ xRy\Rightarrow y\in A\},
\qquad
\Diamond A=\{x\mid \text{hi ha } y \text{ amb } xRy \text{ i } y\in A\}.
\]
Aquestes modalitats són quantificadors \emph{al llarg de la relació} $R$, com els $\forall_f$ i $\exists_f$ del capítol~\ref{cap:subobjectes} ho eren al llarg d'una aplicació. Per veure-ho, sigui $\exists_R(B)=\{y\mid \text{hi ha } x\in B \text{ amb } xRy\}$, el conjunt dels mons accessibles des de $B$. Aleshores, per a $A,B\subseteq W$,
\[
\exists_R(B)\subseteq A\iff B\subseteq\Box A,
\]
és a dir, $\exists_R\dashv\Box$ entre els ordres parcials $(\mathcal P(W),\subseteq)$. En efecte, totes dues condicions diuen el mateix: per a tot $x\in B$ i tot $y$ amb $xRy$, es té $y\in A$. La modalitat $\Diamond$ és el mateix tipus de quantificador, però al llarg de la relació inversa: $\Diamond A=\{x\mid\text{hi ha } y\in A \text{ amb } yR^{-1}x\}$. Les modalitats no són, però, connectius de la lògica interna d'un topos, sinó estructura afegida (informatiu).

\subsection{Resum}

\begin{center}
\small
\renewcommand{\arraystretch}{1.2}
\begin{tabularx}{\linewidth}{@{}>{\raggedright\arraybackslash}p{0.24\linewidth}>{\raggedright\arraybackslash}X>{\raggedright\arraybackslash}X@{}}
\toprule
\textbf{Model} & \textbf{Estructura} & \textbf{Què falla respecte d'un topos}\\
\midrule
$\cat{Vect}_{\mathbb K}$, $\cat{Ab}$, $\cat{Grp}$, $\cat{Rel}$ & amb objecte zero & no són cartesianes tancades\\
$\cat{Cat}$ & cartesiana tancada & no hi ha classificador\\
$\mathbb Z\cup\{+\infty\}$ & preordre cartesià tancat & no hi ha fals\\
subespais de $\mathbb R^2$ & reticle & la distributivitat, i per tant el tancament cartesià\\
Łukasiewicz & preordre monoidal tancat & la implicació no és adjunta de $\wedge$\\
tancats de $\mathbb R$ & co-Heyting & $A\wedge\neg A\neq\bot$\\
models de Kripke & àlgebra de Boole amb modalitats & estructura afegida\\
\bottomrule
\end{tabularx}
\end{center}
Cada fila marca un camí que la lògica categòrica pot seguir més enllà dels topos: les estructures monoidals tancades per a les lògiques de recursos, els reticles no distributius per a la lògica quàntica, els operadors modals com a adjunts al llarg de relacions. El llibre s'ha centrat en l'esquema dels topos perquè és el que unifica més bé la lògica intuïcionista i la teoria de conjunts. Aquestes alternatives mostren que és una tria, no l'únic camí possible.

% =====================================================================
\section{Exercicis}
% =====================================================================

\begin{exercici}
Demostra que $c_y\ast\gamma=\gamma$ si i només si $\gamma$ és constant. \hfill{\normalfont$\star$}
\begin{indicacio}
La igualtat a $[\frac12,1]$ força $\gamma=y$ en aquest interval, i a $[0,\frac12]$ diu $\gamma(t)=\gamma(2t)$. Itera-ho i fes servir la continuïtat a $0$.
\end{indicacio}
\end{exercici}

\begin{exercici}
Comprova que, a la categoria dels camins de Moore, un camí $(r,\gamma)$ amb $r>0$ no té cap invers: no hi ha cap $(s,\delta)$ amb $(s,\delta)\comp(r,\gamma)=(0,c_x)$. \hfill{\normalfont$\star$}
\begin{indicacio}
Mira la durada de la composició.
\end{indicacio}
\end{exercici}

\begin{exercici}
Troba un magma de dos elements que no sigui associatiu, i un altre que sigui associatiu però no tingui neutre. \hfill{\normalfont$\star$}
\begin{indicacio}
Hi ha només setze operacions possibles en un conjunt de dos elements. Per al segon, prova amb una operació que no depengui del segon argument.
\end{indicacio}
\end{exercici}

\begin{exercici}
Considera els nombres reals com a objectes, i com a morfismes $x\to y$ els reals $d\ge|y-x|$, amb la composició $d'\comp d=d+d'$. És una categoria? Què canvia si s'exigeix $d=|y-x|$? \hfill{\normalfont$\star\star$}
\begin{indicacio}
Per al primer cas, fes servir la desigualtat triangular. Per al segon, mira què passa amb $x<y>z$.
\end{indicacio}
\end{exercici}

\begin{exercici}
Sigui $T(G)=\{g\in G\mid g^2=1\}$. Comprova que $G\mapsto T(G)$, amb $f\mapsto f|_{T(G)}$, és un functor $\cat{Grp}\to\cat{Set}$, i explica per què l'argument de la proposició~\ref{prop:centre-no-functor} no s'hi aplica. \hfill{\normalfont$\star$}
\begin{indicacio}
Mira primer els tipus: si $g^2=1$, què es pot dir de $f(g)^2$? Per a la segona part, calcula $T(S_3)$.
\end{indicacio}
\end{exercici}

\begin{exercici}
Demostra que un functor bijectiu sobre objectes i sobre morfismes és un isomorfisme de categories. \hfill{\normalfont$\star$}
\begin{indicacio}
Defineix l'invers sobre objectes i morfismes, i comprova que preserva la composició aplicant-hi el functor original.
\end{indicacio}
\end{exercici}

\begin{exercici}
Demostra que \enquote{ser un grupoide} és invariant per equivalència, i que \enquote{ser un ordre parcial} no ho és. \hfill{\normalfont$\star$}
\begin{indicacio}
Per a la primera part, fes servir que una equivalència és plena i fidel i que reflecteix els isomorfismes.
\end{indicacio}
\end{exercici}

\begin{exercici}
La categoria $\cat{Rel}$ és isomorfa a la seva oposada, enviant cada relació $R$ a la relació inversa. Fes-ho servir, junt amb l'exemple $\cat{Set}\not\simeq\cat{Set}\op$, per demostrar que $\cat{Set}\not\simeq\cat{Rel}$. \hfill{\normalfont$\star\star$}
\begin{indicacio}
Si $F\colon\cat{C}\to\cat{D}$ és una equivalència, $F\op\colon\cat{C}\op\to\cat{D}\op$ també ho és.
\end{indicacio}
\end{exercici}

\begin{exercici}
Sigui $\cat{Set}_\ast$ la categoria dels conjunts amb un punt distingit, amb les aplicacions que preserven el punt distingit. Comprova que té un objecte zero i dedueix que no és cartesiana tancada. \hfill{\normalfont$\star$}
\begin{indicacio}
Pensa en el conjunt d'un sol element. Per a la segona part, aplica la proposició~\ref{prop:zero-no-ccc}.
\end{indicacio}
\end{exercici}

\begin{exercici}
Comprova la fallada de la distributivitat al reticle dels subespais de $\mathbb R^2$ amb les rectes $y=0$, $x=0$ i $y=x$. Té aquest reticle un \enquote{complement} per a cada recta? \hfill{\normalfont$\star$}
\begin{indicacio}
Un complement de $a$ és un $b$ amb $a\wedge b=\{0\}$ i $a\vee b=\mathbb R^2$. És únic?
\end{indicacio}
\end{exercici}

\begin{exercici}
A $\mathbb R$, calcula $A\wedge\neg A$ per al tancat $A=[0,1]\cup\{2\}$, amb $\neg A=\overline{\mathbb R\setminus A}$. \hfill{\normalfont$\star$}
\begin{indicacio}
El resultat és la frontera de $A$.
\end{indicacio}
\end{exercici}

\fi

\backmatter
\chapter*{Bibliografia}
\addcontentsline{toc}{chapter}{Bibliografia}
\markboth{Bibliografia}{Bibliografia}

Les referències següents cobreixen tres funcions: introduccions al mateix nivell que aquest text, obres de consulta i fonts per a la transició cap a la lògica categòrica, i les fonts de la terminologia matemàtica catalana emprada.

\begin{thebibliography}{99}

\bibitem{eilenbergmaclane}
S.~Eilenberg i S.~Mac~Lane,
\emph{General theory of natural equivalences},
Transactions of the American Mathematical Society~\textbf{58} (1945), 231--294.
\newblock (Article fundacional de la teoria de categories, on es defineixen per primera vegada les nocions de categoria, functor i transformació natural. \url{https://doi.org/10.1090/S0002-9947-1945-0013131-6}.)

\bibitem{awodey}
S.~Awodey,
\emph{Category Theory},
2a edició, Oxford Logic Guides~52, Oxford University Press, 2010.
\newblock (Text introductori de referència; guia principal per a l'ordre i el nivell d'aquest llibre, i per als capítols sobre categories cartesianes tancades i lògica. Les remissions del llibre es refereixen a aquesta edició.)

\bibitem{maclane}
S.~Mac~Lane,
\emph{Categories for the Working Mathematician},
2a edició, Graduate Texts in Mathematics~5, Springer, 1998. \url{https://doi.org/10.1007/978-1-4757-4721-8}.
\newblock (Obra de consulta clàssica sobre teoria de categories. Les remissions del llibre es refereixen a aquesta edició.)

\bibitem{leinster}
T.~Leinster,
\emph{Basic Category Theory},
Cambridge Studies in Advanced Mathematics~143, Cambridge University Press, 2014. Versió en accés obert: \url{https://arxiv.org/abs/1612.09375}.
\newblock (Introducció alternativa, breu i moderna. La versió oberta es va publicar a arXiv el desembre de 2016 per acord amb Cambridge University Press; incorpora les correccions de la llista d'errates de l'edició impresa i en conserva la numeració de capítols i seccions, que és la que segueixen les remissions d'aquest llibre. Consulta: octubre de 2026.)

\bibitem{riehl}
E.~Riehl,
\emph{Category Theory in Context},
Aurora: Dover Modern Math Originals, Dover, 2016.
\newblock (Consulta complementària, amb molts exemples. Les remissions d'aquest llibre, per capítols i seccions, es refereixen a l'edició de 2016. La pàgina de l'autora ofereix una versió en accés obert, \url{https://emilyriehl.github.io/files/context.pdf}, que s'actualitza i pot no coincidir amb l'edició de 2016 en la paginació. La pàgina de l'autora distingeix actualment la primera edició de la versió en preparació de la segona edició; les remissions d'aquest llibre es refereixen a la primera edició de 2016. Consulta: octubre de 2026.)

\bibitem{barrwells}
M.~Barr i C.~Wells,
\emph{Category Theory for Computing Science},
3a edició, Les Publications CRM, Montréal, 1999. Reimprès a \emph{Reprints in Theory and Applications of Categories}~22 (2012). \url{https://www.tac.mta.ca/tac/reprints/articles/22/tr22abs.html}.
\newblock (Introducció orientada a la informàtica i la lògica, amb sistemes deductius i llenguatges funcionals com a exemples de categories. El portal de TAC n'ofereix també versions corregides posteriors a la reimpressió de 2012; les remissions d'aquest llibre es refereixen a la reimpressió de 2012; per a pàgines exactes cal indicar la versió consultada. Consulta: octubre de 2026.)

\bibitem{perrone}
P.~Perrone,
\emph{Starting Category Theory},
World Scientific, Singapur, 2024. \url{https://doi.org/10.1142/13670}.
\newblock (Introducció detallada i accessible, amb molts exemples de matemàtica bàsica.)

\bibitem{perrone-notes}
P.~Perrone,
\emph{Notes on Category Theory with examples from basic mathematics},
notes en accés obert, arXiv:1912.10642, versió~7 (abril de 2024; text no revisat, darrera actualització de febrer de 2021), \url{https://arxiv.org/abs/1912.10642}.
\newblock (Material previ al llibre \cite{perrone}, que el va ampliar i reorganitzar; les remissions a aquestes notes no es corresponen necessàriament amb els capítols del llibre. Consulta: octubre de 2026.)

\bibitem{lawvereschanuel}
F.~W.~Lawvere i S.~H.~Schanuel,
\emph{Conceptual Mathematics: A First Introduction to Categories},
2a edició, Cambridge University Press, 2009.
\newblock (Introducció molt elemental, amb conjunts, grafs i sistemes dinàmics com a exemples recurrents. Les remissions es fan per parts, articles i sessions, i s'han comprovat amb la traducció castellana de la segona edició, \emph{Matemáticas conceptuales. Una primera introducción a categorías}, de F.~Marmolejo, que en conserva la numeració.)

\bibitem{streicher}
T.~Streicher,
\emph{Introduction to Category Theory and Categorical Logic},
notes de curs, Technische Universität Darmstadt, 2003--2004.
\newblock (Introducció a la teoria de categories orientada a la lògica categòrica; referència per als capítols de categories cartesianes tancades i topos i per al pont cap al volum de lògica. Les remissions es refereixen a la versió en línia de les notes indicada a la referència.)

\bibitem{borceux}
F.~Borceux,
\emph{Handbook of Categorical Algebra},
3 volums, Encyclopedia of Mathematics and its Applications, Cambridge University Press, 1994.
\newblock (Referència exhaustiva per a límits, adjuncions i exactitud.)

\bibitem{maclanemoerdijk}
S.~Mac~Lane i I.~Moerdijk,
\emph{Sheaves in Geometry and Logic: A First Introduction to Topos Theory},
Universitext, Springer, 1992.
\newblock (Pont cap a la teoria de topos i la lògica geomètrica. Es cita l'edició estàndard impresa de 1992.)

\bibitem{lambekscott}
J.~Lambek i P.~J.~Scott,
\emph{Introduction to Higher-Order Categorical Logic},
Cambridge Studies in Advanced Mathematics~7, Cambridge University Press, 1986.
\newblock (Correspondència entre lògica d'ordre superior, càlcul lambda i categories cartesianes tancades.)

\bibitem{makkaireyes}
M.~Makkai i G.~E.~Reyes,
\emph{First Order Categorical Logic},
Lecture Notes in Mathematics~611, Springer, 1977.
\newblock (Tractament clàssic de la lògica de primer ordre en categories regulars, coherents i de Heyting, i dels topos classificadors; referència per al volum de lògica categòrica.)

\bibitem{pitts}
A.~M.~Pitts,
\emph{Categorical Logic},
dins S.~Abramsky, D.~M.~Gabbay i T.~S.~E.~Maibaum (eds.), \emph{Handbook of Logic in Computer Science}, vol.~5, Oxford University Press, 2000, pàg.~39--128.
\newblock (Introducció sistemàtica a la semàntica categòrica de la lògica, des de les teories equacionals fins a la lògica d'ordre superior; referència per al volum de lògica categòrica.)

\bibitem{jacobs}
B.~Jacobs,
\emph{Categorical Logic and Type Theory},
Studies in Logic and the Foundations of Mathematics~141, Elsevier, 1999.
\newblock (Fibracions, hiperdoctrines i teoria de tipus des del punt de vista categòric.)

\bibitem{johnstone}
P.~T.~Johnstone,
\emph{Sketches of an Elephant: A Topos Theory Compendium},
Oxford Logic Guides 43--44, Oxford University Press, 2002.
\newblock (Obra avançada de consulta sobre teoria de topos.)

\bibitem{lawvere-ccfm}
F.~W.~Lawvere,
\emph{The Category of Categories as a Foundation for Mathematics},
en \emph{Proceedings of the Conference on Categorical Algebra (La Jolla, 1965)}, Springer, 1966, pàg.~1--20.
\newblock (Teoria elemental de categories en primer ordre; precedent de la presentació relacional de l'apèndix~\ref{cap:axiomatica-relacional}.)

\bibitem{lawvere-etcs}
F.~W.~Lawvere,
\emph{An Elementary Theory of the Category of Sets},
Proceedings of the National Academy of Sciences of the USA~\textbf{52} (1964), 1506--1511. Versió llarga, amb comentaris de C.~McLarty i de l'autor, a \emph{Reprints in Theory and Applications of Categories}~11 (2005): \url{http://www.tac.mta.ca/tac/reprints/articles/11/tr11abs.html}.
\newblock (Axiomàtica de la teoria de conjunts formulada en termes de la categoria de conjunts: la teoria de categories com a fonament.)

\bibitem{leinster-rethinking}
T.~Leinster,
\emph{Rethinking Set Theory},
American Mathematical Monthly~\textbf{121} (2014), núm.~5, 403--415. Versió en accés obert: \url{https://arxiv.org/abs/1212.6543}.
\newblock (Presentació informal i breu de la teoria de conjunts estructural de Lawvere.)

\bibitem{shulman-mida}
M.~A.~Shulman,
\emph{Set Theory for Category Theory},
prepublicació, 2008. \url{https://arxiv.org/abs/0810.1279}.
\newblock (Comparació de les diferents maneres de tractar la mida en teoria de categories: classes, universos i altres alternatives.)

\bibitem{nlab-classificador}
nLab,
\emph{classifying topos},
\url{https://ncatlab.org/nlab/show/classifying+topos} (consultat el 29 de setembre de 2026).
\newblock (Wiki de referència en teoria de categories; entrada sobre el topos classificador d'una teoria geomètrica.)

\bibitem{nlab-locale}
nLab,
\emph{locale},
\url{https://ncatlab.org/nlab/show/locale} (consultat el 29 de setembre de 2026).
\newblock (Hi apareix l'exemple del locale de les aplicacions exhaustives $\mathbb N\to\mathbb R$, no trivial i sense punts.)

\bibitem{blechschmidt}
I.~Blechschmidt,
\emph{Generalized Spaces for Constructive Algebra},
prepublicació, 2020. \url{https://arxiv.org/abs/2012.13850}.
\newblock (Presentació accessible dels espais generalitzats sense punts, amb l'exemple de les aplicacions exhaustives $\mathbb N\to\mathbb R$.)

\bibitem{caramello}
O.~Caramello,
\emph{Theories, Sites, Toposes: Relating and Studying Mathematical Theories through Topos-Theoretic \enquote{Bridges}},
Oxford University Press, 2018.
\newblock (Referència sobre els toposos classificadors des del punt de vista de la lògica. Primera edició impresa de 2018; la versió electrònica es va publicar el desembre de 2017.)

\bibitem{hott}
The Univalent Foundations Program,
\emph{Homotopy Type Theory: Univalent Foundations of Mathematics},
Institute for Advanced Study, 2013. Accés obert: \url{https://homotopytypetheory.org/book/}.
\newblock (Introducció a la teoria homotòpica de tipus, llegible des de la lògica.)

\bibitem{scott-identity}
D.~S.~Scott,
\emph{Identity and Existence in Intuitionistic Logic},
en M.~Fourman, C.~Mulvey i D.~Scott (eds.), \emph{Applications of Sheaves}, Lecture Notes in Mathematics~753, Springer, 1979, pàg.~660--696.
\newblock (Operacions parcials, existència i identitat; rerefons de les axiomàtiques \enquote{només amb fletxes}.)

\bibitem{freydscedrov}
P.~J.~Freyd i A.~Scedrov,
\emph{Categories, Allegories},
North-Holland Mathematical Library~39, North-Holland, 1990.
\newblock (Presentació sistemàtica sense objectes primitius.)

\bibitem{BS2020}
C.~Benzmüller i D.~S.~Scott,
\emph{Automating Free Logic in HOL, with an Experimental Application in Category Theory},
Journal of Automated Reasoning~64, 2020, pàg.~53--72.
\newblock (Formalització en Isabelle/HOL de diverses axiomàtiques de categories amb lògica lliure.)

\bibitem{BSAFP}
C.~Benzmüller i D.~S.~Scott,
\emph{Axiom Systems for Category Theory in Free Logic},
Archive of Formal Proofs, 2018. \url{https://isa-afp.org/entries/AxiomaticCategoryTheory.html}.
\newblock (Comparació i verificació d'equivalències entre sistemes de Mac~Lane, Scott i Freyd--Scedrov.)

\bibitem{diec}
Institut d'Estudis Catalans,
\emph{Diccionari de la llengua catalana},
2a edició, Barcelona. \url{https://dlc.iec.cat}.
\newblock (Font normativa de la terminologia general catalana.)

\bibitem{termcat}
Universitat Politècnica de Catalunya; Enciclopèdia Catalana,
\emph{Diccionari de matemàtiques i estadística},
1a ed., Barcelona, 2002; 2a ed. en línia, Institut d'Estudis Catalans, 2021. \url{https://cit.iec.cat/DME/}.
\newblock (Font de la terminologia matemàtica catalana: functor, morfisme, preordre, etc.)

\end{thebibliography}

\chapter*{Glossari català--anglès}
\addcontentsline{toc}{chapter}{Glossari català--anglès}
\markboth{Glossari català--anglès}{Glossari català--anglès}

Aquest glossari recull sobretot termes conflictius, no transparents o especialment útils per llegir bibliografia internacional, i no pretén ser un vocabulari exhaustiu de teoria de categories. S'hi dona la forma anglesa i, quan escau, la notació emprada.

\medskip
\noindent\textbf{Criteri d'anglicismes.} Diversos termes s'han manllevat de l'anglès per manca d'una forma catalana establerta o per fidelitat a la bibliografia (\emph{pullback}, \emph{pushout}). En la seva primera aparició es dona, quan existeix, el terme català al costat de l'anglès i de la notació; després se n'usa una sola forma. La terminologia general segueix el \emph{Diccionari de la llengua catalana} de l'Institut d'Estudis Catalans \cite{diec}, i la matemàtica, el \emph{Diccionari de matemàtiques i estadística} del TERMCAT \cite{termcat}.

\medskip

\renewcommand{\arraystretch}{1.25}
\begin{xltabular}{\linewidth}{@{}>{\raggedright\arraybackslash}X >{\raggedright\arraybackslash}X l@{}}
\hline
\textbf{Català} & \textbf{Anglès} & \textbf{Notació}\\
\hline
\endfirsthead
\hline
\textbf{Català} & \textbf{Anglès} & \textbf{Notació}\\
\hline
\endhead
\hline
\endfoot
\endlastfoot
categoria oposada & opposite category & $\cat{C}\op$\\
categoria prima & thin category & \\
categoria sobre un objecte & slice category & $\cat{C}/A$\\
categoria sota un objecte & coslice category & $A/\cat{C}$\\
homconjunt & hom-set & $\Hom(A,B)$\\
functor representable & representable functor & p.\,ex.\ $\Hom(-,A)$\\
immersió de Yoneda & Yoneda embedding & \\
bihomfunctor (o bifunctor Hom) & hom functor, Hom bifunctor & $\Hom(-,-)$\\
component (d'una transformació natural; femení) & component & $\eta_A$\\
antiimatge & preimage, inverse image & $f^{-1}(B)$\\
exhaustiu / essencialment exhaustiu & surjective / essentially surjective & \\
\multicolumn{3}{@{}p{\linewidth}@{}}{\footnotesize En aquest llibre, \emph{exhaustiu} és el terme sistemàtic per a \emph{surjective}; \emph{essencialment exhaustiu} tradueix \emph{essentially surjective}.}\\[0.35em]
transformació natural & natural transformation & $\eta\colon F\Rightarrow G$\\
equivalència de categories & equivalence of categories & $\cat{C}\simeq\cat{D}$\\
monomorfisme / epimorfisme & monomorphism / epimorphism & mono / epi\\
objecte inicial / terminal & initial / terminal object & $0$ / $1$ (quan s'adopta aquesta notació)\\
functor ple / fidel & full / faithful functor & \\
imatge essencial & essential image & \\
categoria d'elements & category of elements & $\textstyle\int P$\\
categoria essencialment petita & essentially small category & \\
quasiinvers (d'una equivalència) & quasi-inverse / pseudo-inverse & \\
prefeix & presheaf & $\cat{Set}^{\cat{C}\op}$\\
producte fibrat (pullback) & pullback & \\
suma amalgamada (pushout) & pushout & \\
con & cone & \\
cocon & cocone & \\
límit / colímit & limit / colimit & $\varprojlim$ / $\varinjlim$\\
unitat / counitat & unit / counit & $\eta$ / $\varepsilon$\\
composició horitzontal & horizontal composition & $\theta\ast\eta$\\
composició amb un functor per l'esquerra & (left) whiskering & $K\eta$\\
composició amb un functor per la dreta & (right) whiskering & $\eta H$\\
epimorfisme regular & regular epimorphism & \\
parella nucli & kernel pair & \\
factorització imatge & image factorization & \\
reindexació & reindexing / base change & $f^{*}$\\
subobjecte & subobject & $\Sub(A)$\\
classificador de subobjectes & subobject classifier & $\Omega$\\
sedàs & sieve & \\
semireticle implicatiu & implicative meet-semilattice & \\
categoria cartesiana tancada & cartesian closed category & \\
categoria regular & regular category & \\
\hline
\end{xltabular}
\renewcommand{\arraystretch}{1.0}

\chapter*{Taula de notació}
\addcontentsline{toc}{chapter}{Taula de notació}
\markboth{Taula de notació}{Taula de notació}

Recollim la notació principal del llibre, amb el seu significat i el lloc de la primera aparició o definició. Les referències remeten al capítol o a la secció corresponent.

\medskip
\renewcommand{\arraystretch}{1.25}
\begin{xltabular}{\linewidth}{@{}l >{\raggedright\arraybackslash}X l@{}}
\hline
\textbf{Símbol} & \textbf{Significat} & \textbf{Primera aparició}\\
\hline
\endfirsthead
\hline
\textbf{Símbol} & \textbf{Significat} & \textbf{Primera aparició}\\
\hline
\endhead
\hline
\endfoot
\endlastfoot
$\cat{C},\cat{D},\dots$ & categories & cap.~\ref{cap:categories}\\
$\Hom(A,B)$, $\cat{C}(A,B)$ & homconjunt de morfismes $A\to B$ & cap.~\ref{cap:categories}\\
$\id_A$ & morfisme identitat de $A$ & cap.~\ref{cap:categories}\\
$g\comp f$ & composició ($f$ seguit de $g$) & cap.~\ref{cap:categories}\\
$\cat{C}\op$ & categoria oposada & cap.~\ref{cap:categories}\\
$f\op$ & morfisme dual de $f$ a $\cat{C}\op$ & \S\ref{sec:dualitat}\\
$\cat{Set},\cat{Grp},\cat{Top},\cat{Graph}$ & categories concretes habituals & cap.~\ref{cap:categories}\\
$\cat{Cat}$ & categoria de les categories petites & cap.~\ref{cap:functors}\\
$\cat{C}/A$, $A/\cat{C}$ & categoria sobre, sota un objecte & cap.~\ref{cap:functors}\\
$(S\downarrow T)$ & categoria coma & \S\refalt{sec:coma}{del cap.~4}\\
$F,G\colon\cat{C}\to\cat{D}$ & functors & cap.~\ref{cap:functors}\\
$\Hom(A,-)$, $\Hom(-,B)$ & homfunctors covariant i contravariant & def.~\ref{def:homfunctors}\\
$\Hom(A,f)=f_*$ & postcomposició amb $f$ (acció de $\Hom(A,-)$ sobre $f$) & not.~\ref{not:hom-morfismes}\\
$\Hom(f,B)=f^*$ & precomposició amb $f$ (acció de $\Hom(-,B)$ sobre $f$). \emph{Atenció:} el mateix símbol designa la reindexació (vegeu més avall) & not.~\ref{not:hom-morfismes}\\
$\widehat{\cat{C}}=\cat{Set}^{\cat{C}\op}$ & categoria de prefeixos sobre $\cat{C}$ & ex.~\ref{ex:categoria-prefeixos}\\
$\Hom(-,-)$, $\Hom(f,g)$ & bihomfunctor i la seva acció $h\mapsto g\comp h\comp f$ & def.~\ref{def:bihomfunctor}\\
$\alpha\colon F\Rightarrow G$ & transformació natural & cap.~\ref{cap:transformacions}\\
$\cat{D}^{\cat{C}}$, $[\cat{C},\cat{D}]$ & categoria de functors & cap.~\ref{cap:transformacions}\\
$\Hom(-,A)$ & exemple de functor representable (contravariant) & cap.~\refcap{cap:yoneda}{4}\\
$\yon$, $\Hom(-,f)$ & immersió de Yoneda i la seva acció $\yon(f)$ sobre $f$ & def.~\refalt{def:immersio-yoneda}{del cap.~4}\\
$\Delta_A$ & functor constant amb valor $A$ & ex.~\ref{ex:functor-constant}\\
$\Delta$ & functor diagonal $\cat{C}\to\cat{C}^{\cat{I}}$ & obs.~\refalt{obs:cons-coma}{del cap.~5}\\
$\delta_A$ & morfisme diagonal $\langle\id_A,\id_A\rangle\colon A\to A\times A$ & cap.~\refcap{cap:cartesianes}{7}\\
$\varprojlim D$, $\colimop D$ & límit, colímit d'un diagrama & cap.~\refcap{cap:limits}{5}\\
$A\times B$, $A+B$ & producte, coproducte & cap.~\refcap{cap:limits}{5}\\
$F\dashv G$ & $F$ és adjunt per l'esquerra de $G$ & cap.~\refcap{cap:adjuncions}{6}\\
$\eta$, $\varepsilon$ & unitat, counitat d'una adjunció & cap.~\refcap{cap:adjuncions}{6}\\
$C^{B}$, $\lambda f$ & exponencial, currificació & cap.~\refcap{cap:cartesianes}{7}\\
$\Sub(A)$ & ordre parcial de subobjectes de $A$ & cap.~\refcap{cap:subobjectes}{8}\\
$f^{*}$ & reindexació per pullback (substitució), a $\cat{C}/B$ i a $\Sub(B)$; a $\cat{Set}$, l'antiimatge $f^{-1}$. És el mateix símbol que la precomposició; el context (homconjunts o subobjectes i fibres) decideix. A $\cat{Set}$ totes dues lectures coincideixen a través de $\Sub(B)\cong\Hom(B,\{0,1\})$ & obs.~\refalt{obs:reindexacio}{del cap.~5}; cap.~\refcap{cap:subobjectes}{8}\\
$\Imatge(f)$ & imatge d'un morfisme & cap.~\refcap{cap:subobjectes}{8}\\
$\exists_f$, $\forall_f$ & quantificadors com a adjunts de $f^{*}$ & cap.~\refcap{cap:subobjectes}{8}\\
$\Omega$, $\mathsf{true}$ & classificador de subobjectes & cap.~\refcap{cap:topos}{9}\\
\hline
\end{xltabular}
\renewcommand{\arraystretch}{1.0}

\printindex

\end{document}
"
%
%   (b) descomentant la línia \publicacioparcialtrue de més avall.
%
% Recordeu recompilar amb XeLaTeX + makeindex + XeLaTeX x2.
% ============================================================
\newif\ifpublicacioparcial
\ifdefined\fascicle\publicacioparcialtrue\else\publicacioparcialfalse\fi
% \publicacioparcialtrue % <- descomenteu per forçar el fascicle sense \fascicle

% Mode de sortida. Per defecte: pantalla/interactiu. El Makefile defineix
% \imprimible per generar una còpia pensada per imprimir: fons blanc,
% enllaços discrets i tests estàtics sense camps AcroForm.
\newif\ifimprimible
\ifdefined\imprimible\imprimibletrue\else\imprimiblefalse\fi
% Nom de la Part IV segons el mode: en paper no hi ha interactivitat.
\ifimprimible
  \newcommand{\NomTests}{Tests d'autoavaluació}
  \newcommand{\nomtests}{tests d'autoavaluació}
\else
  \newcommand{\NomTests}{Tests interactius}
  \newcommand{\nomtests}{tests interactius}
\fi

% Els fitxers compartits viuen ara a la subcarpeta common/ d'aquest mateix
% projecte, de manera que es troben amb ruta explícita, sense dependre de
% TEXINPUTS ni de cap configuració del sistema de compilació.
\usepackage[sobri]{common/aprendes}
% ============================================================
%  Preàmbul compartit — Aprendes Universitat
%  Compartit pels llibres d'Aprendes Universitat.
%  Les dades pròpies de cada llibre (títol, autor i metadades PDF)
%  van al metadata.tex local de cada projecte, carregat ABANS
%  d'aquest fitxer. La macro \hypermetadata es crida DESPRÉS,
%  quan hyperref ja ha estat carregat.
% ============================================================

\usepackage{eurosym}
\usepackage{iftex}
\ifpdftex
  \usepackage[T1]{fontenc}
  \usepackage[utf8]{inputenc}
  \usepackage{lmodern}
\else
  % XeLaTeX / LuaLaTeX: Unicode natiu, res d'inputenc
  \usepackage{fontspec}
  % Carrega Latin Modern per nom de fitxer, que XeTeX/LuaTeX pot
  % resoldre via kpathsea fins i tot si fontconfig no l'ha registrat.
  \setmainfont{lmroman10-regular.otf}[
    BoldFont=lmroman10-bold.otf,
    ItalicFont=lmroman10-italic.otf,
    BoldItalicFont=lmroman10-bolditalic.otf
  ]
  \setsansfont{lmsans10-regular.otf}[
    BoldFont=lmsans10-bold.otf,
    ItalicFont=lmsans10-oblique.otf,
    BoldItalicFont=lmsans10-boldoblique.otf
  ]
  \setmonofont{lmmono10-regular.otf}[ItalicFont=lmmono10-italic.otf]
\fi
% microtype: protrusion i expansió completes només amb pdfLaTeX.
% Amb XeLaTeX/LuaLaTeX només hi ha protrusion i generaria avisos de
% «slot number»; per això només l'activem sota pdfTeX.
\ifpdftex
  \usepackage{microtype}
\fi
\usepackage[catalan, provide=*]{babel}
\usepackage[autostyle=true,babel=true]{csquotes}
% csquotes no porta estil català: el declarem perquè \enquote{...}
% doni guillemets «...» (exteriors) i cometes altes “...” (interiors),
% independentment de la versió de csquotes instal·lada.
\DeclareQuoteStyle{catalan}
  {\guillemotleft}{\guillemotright}
  {\quotedblbase}{\textquotedblleft}
\DeclareQuoteAlias{catalan}{catala}
\usepackage{amsmath,mathtools}% amssymb, amsthm: via aprendes.sty
\usepackage{geometry}
\usepackage{import}
\usepackage{graphicx}
\usepackage{tabularx}
\usepackage{xltabular}% tabularx que pagina (taula de notació)
\usepackage{tikz}
\usetikzlibrary{positioning,arrows.meta,matrix}
\usepackage[all]{xy}% diagrames de grafs (llaços i arestes múltiples) a l'apèndix de grafs
\usepackage{float}
\usepackage[section]{placeins}% FloatBarrier automàtic: cap figura no travessa una secció
\usepackage{pgfplots}
\pgfplotsset{compat=1.18}
\usepackage{imakeidx}
\usepackage{booktabs}

% La bibliografia es configura localment en cada projecte quan hi ha referències.

\setcounter{MaxMatrixCols}{10}

\graphicspath{{img/}}
\geometry{hmargin={3cm,3cm},vmargin={2.5cm,2.5cm},offset=0pt,footskip=1cm,headsep=1cm}


\setlist{itemsep=0.25em, topsep=0.25em}
\numberwithin{equation}{section}
\setcounter{secnumdepth}{3}
\setcounter{tocdepth}{2}

% Configuració dels paràgrafs
\usepackage{indentfirst}
\setlength{\parindent}{1.5em}
\setlength{\parskip}{0.3em}

% ============================================================
%  hyperref: ÚLTIM paquet, com recomana el seu manual.
%  Els colors dels enllaços es fixen aquí. Les metadades
%  pròpies del llibre ja s'han definit abans en metadata.tex;
%  la macro \hypermetadata s'executa després de carregar aquest fitxer.
% ============================================================
\usepackage[pdflang=ca-ES,pdfdisplaydoctitle=true]{hyperref}
% Color dels enllaços: el blau pissarra de la paleta d'aprendes,
% una mica enfosquit. Prou distint del vermell corporatiu dels
% títols, exemples i exercicis, i coherent amb la resta de colors.
\colorlet{EnllacBlau}{ResultatBlau!90!black}
\ifimprimible
  \hypersetup{
    colorlinks=true,
    linktoc=page,
    linkcolor=black,
    urlcolor=black,
    citecolor=black,
    filecolor=black,
    pdfborder={0 0 0},
    hypertexnames=false,
  }
\else
  \hypersetup{
    colorlinks=true,
    linktoc=page,
    linkcolor=EnllacBlau,
    urlcolor=EnllacBlau,
    citecolor=EnllacBlau,
    filecolor=EnllacBlau,
    hypertexnames=false,
  }
\fi
\usepackage{bookmark}% millora els marcadors del PDF; ha d'anar després d'hyperref

% Macros matemàtiques i auxiliars compartides
\newcommand{\Qcb}[1]{#1}
\newcommand{\func}[1]{\operatorname{#1}}
\newcommand{\dbigcup}{\displaystyle\bigcup}
\newcommand{\dbigcap}{\displaystyle\bigcap}
\newcommand{\FRAME}[8]{%
\begin{center}
#5\par\smallskip
\includegraphics[width=#2,height=#3,keepaspectratio]{img/#7}
\end{center}}

\newcommand{\R}{\mathbb{R}}
\newcommand{\N}{\mathbb{N}}
\newcommand{\Z}{\mathbb{Z}}
\newcommand{\Q}{\mathbb{Q}}
\newcommand{\Natpos}{\mathbb{N}^{*}}
\newcommand{\Zstar}{\mathbb{Z}^{*}}
\newcommand{\mcd}{\operatorname{mcd}}
\newcommand{\mcm}{\operatorname{mcm}}
\newcommand{\abs}[1]{\left\lvert #1\right\rvert}
\newcommand{\classe}[1]{\left[#1\right]}
\newcommand{\divides}{\mid}
\newcommand{\notdivides}{\nmid}
\newcommand{\decimalperiodic}[2]{#1,\overline{#2}}
\DeclareMathOperator{\sgn}{sgn} 
% ------------------------------------------------------------
% Macros pròpies de teoria de categories
% ------------------------------------------------------------
\newcommand{\cat}[1]{\mathbf{#1}}          % noms de categories: \cat{Set}
\newcommand{\id}{\operatorname{id}}         % identitat
\providecommand{\Hom}{\operatorname{Hom}}   % conjunts de morfismes
\newcommand{\op}{^{\mathrm{op}}}            % categoria oposada
\newcommand{\Obj}{\operatorname{Obj}}       % col·lecció d'objectes
\newcommand{\comp}{\circ}                    % composició
\DeclareMathOperator{\dom}{dom}
\DeclareMathOperator{\cod}{cod}
\DeclareMathOperator{\Ob}{Ob}
\DeclareMathOperator{\Mor}{Mor}
\DeclareMathOperator{\Aut}{Aut}
\DeclareMathOperator{\Free}{Free}
\DeclareMathOperator{\Nat}{Nat}       % transformacions naturals com a homconjunt
\DeclareMathOperator{\Sub}{Sub}       % preordre de subobjectes
\DeclareMathOperator{\ev}{ev}         % morfisme d'avaluació (exponencials)
\DeclareMathOperator{\Imatge}{Im}     % imatge d'un morfisme (factorització imatge)
% Claudàtors semàntics [[-]] sense dependre de stmaryrd:
\providecommand{\llbracket}{\mathopen{[\![}}
\providecommand{\rrbracket}{\mathclose{]\!]}}
\newcommand{\yon}{\mathrm{y}}          % immersió de Yoneda
\DeclareMathOperator*{\colimop}{colim} % colímit (amb límits sota)

% ============================================================
%  Macros pròpies de la teoria de categories, locals al volum.
%
%  El fitxer compartit ../common/macros.tex barreja notació
%  genèrica (\R, \N, \mcd, decimals periòdics...) amb la de
%  categories. Aquest fitxer recull, en un sol lloc i documentada,
%  només la notació que aquest llibre fa servir, de manera que el
%  volum no depèn de l'organització del fitxer compartit.
%
%  S'usa \providecommand perquè, si macros.tex ja ha definit una
%  macro, no es produeixi cap col·lisió: la definició compartida i
%  la local són idèntiques. Si algun dia es depura macros.tex
%  traient-ne la part de categories, aquest fitxer la garanteix.
% ============================================================

\providecommand{\cat}[1]{\mathbf{#1}}          % noms de categories: \cat{Set}
\providecommand{\id}{\operatorname{id}}         % identitat
\providecommand{\Hom}{\operatorname{Hom}}       % conjunts de morfismes
\providecommand{\op}{^{\mathrm{op}}}            % categoria oposada
\providecommand{\Obj}{\operatorname{Obj}}       % col·lecció d'objectes
\providecommand{\comp}{\circ}                    % composició
\providecommand{\yon}{\mathrm{y}}                % immersió de Yoneda

\providecommand{\dom}{\operatorname{dom}}
\providecommand{\cod}{\operatorname{cod}}

% Predicats relacionals de l'apèndix d'axiomàtica relacional (Dom, Cod, Com,
% Id). S'escriuen amb versaletes matemàtiques per distingir-los clarament dels
% operadors funcionals \dom, \cod: aquí són relacions primitives, no funcions.
\providecommand{\Dom}{\mathsf{Dom}}
\providecommand{\Cod}{\mathsf{Cod}}
\providecommand{\Com}{\mathsf{Com}}
\providecommand{\Id}{\mathsf{Id}}
\providecommand{\Free}{\operatorname{Free}}
\providecommand{\Nat}{\operatorname{Nat}}        % transformacions naturals com a homconjunt
\providecommand{\Sub}{\operatorname{Sub}}        % preordre de subobjectes
\providecommand{\ev}{\operatorname{ev}}          % morfisme d'avaluació (exponencials)
\providecommand{\Imatge}{\operatorname{Im}}      % imatge d'un morfisme

% Claudàtors semàntics [[-]] sense dependre de stmaryrd:
\providecommand{\llbracket}{\mathopen{[\![}}
\providecommand{\rrbracket}{\mathclose{]\!]}}

% colimit amb subíndex a sota (com \lim), només si encara no existeix:
\makeatletter
\@ifundefined{colimop}{\DeclareMathOperator*{\colimop}{colim}}{}
\makeatother

% Remissions a capítols i elements que no són a l'edició parcial (capítols 1--3).
% A l'edició completa fan servir \ref; a la parcial escriuen el número o un text
% alternatiu, per no deixar referències «??».
\newcommand{\refcap}[2]{\ifpublicacioparcial #2\else\ref{#1}\fi}
\newcommand{\refalt}[2]{\ifpublicacioparcial #2\else\ref{#1}\fi}

% ============================================================
%  Estils TikZ comuns per als diagrames del llibre.
%
%  Els diagrames de teoria de categories repeteixen molt codi de
%  puntes de fletxa i de matrius de nodes. Aquests estils permeten
%  escriure els diagrames nous de manera compacta i canviar
%  l'estètica de tots des d'un sol lloc:
%
%     \begin{tikzpicture}
%       \matrix (m) [cat matrix] { A & B \\ C & D \\ };
%       \draw[cat] (m-1-1) -- node[above] {$f$} (m-1-2);
%     \end{tikzpicture}
%
%  Els diagrames antics, escrits amb -{Stealth[scale=1.1]} i amb
%  els paràmetres de matriu explícits, continuen funcionant sense
%  canvis; aquests estils són per a diagrames nous i per a una
%  homogeneïtzació progressiva.
% ============================================================

\tikzset{
  % Fletxa estàndard d'un morfisme.
  cat/.style={-{Stealth[scale=1.1]}},
  % Fletxa de monomorfisme (cua d'ham).
  cat mono/.style={{Hooks[right]}-{Stealth[scale=1.1]}},
  % Fletxa d'epimorfisme (doble punta).
  cat epi/.style={-{Stealth[scale=1.1]Stealth[scale=1.1]}},
  % Fletxa de transformació natural (doble).
  cat nat/.style={-{Implies},double,double distance=1.5pt},
  % Matriu de nodes amb la separació habitual del llibre.
  cat matrix/.style={matrix of math nodes, row sep=2.6em, column sep=3.2em},
}

% Tests interactius en PDF
\newcounter{testquestion}
\newcommand{\testprefix}{}
\newcommand{\testkeys}{}
\newcommand{\iniciTest}[2]{%
  \renewcommand{\testprefix}{#1}%
  \renewcommand{\testkeys}{#2}%
  \setcounter{testquestion}{0}%
  \gdef\testjustificacions{}%
  \label{test:#1}%
  \begin{center}
  \fbox{\begin{minipage}{0.9\textwidth}
  \ifimprimible
    \textbf{Instruccions.} Marca una sola resposta per pregunta. Comprova les respostes i les justificacions a l'apèndix de solucions.
  \else
    \textbf{Instruccions.} Marca una sola resposta per pregunta. En acabar, prem el bot\'{o} \textbf{Corregir}. El camp \textbf{Nota} mostrar\`{a} el resultat.\par\medskip
    \PushButton[name=#1corregir,width=3.2cm,height=0.8cm,bordercolor={0 0 1},onclick={var p='#1';var k='#2'.split(',');var opts=['a','b','c','d'];var s=0;for(var i=0;i<k.length;i++){var ans='';var n=0;for(var j=0;j<opts.length;j++){var q=this.getField(p+'q'+(i+1)+opts[j]);if(q && q.value!='Off'){ans=opts[j];n++;}}var f=this.getField(p+'f'+(i+1));if(n==1 && ans==k[i]){s++;if(f)f.value='Correcte';}else{if(f)f.value=(n==0?'Sense respondre. Correcta: '+k[i]:(n>1?'Marca una sola resposta. Correcta: '+k[i]:'Incorrecte. Correcta: '+k[i]));}}this.getField(p+'nota').value=s+' / '+k.length;app.alert('Nota: '+s+' / '+k.length);}]{Corregir}\quad
    \PushButton[name=#1netejar,width=3.2cm,height=0.8cm,bordercolor={1 0 0},onclick={var keys='#2'.split(',');var opts=['a','b','c','d'];this.getField('#1nota').value='';for(var i=1;i<=keys.length;i++){var f=this.getField('#1f'+i);if(f)f.value='';for(var j=0;j<opts.length;j++){var q=this.getField('#1q'+i+opts[j]);if(q)q.value='Off';}}}]{Netejar}\quad
    \textbf{Nota: }\TextField[name=#1nota,readonly=true,width=4cm,bordercolor={0 0 0}]{}
  \fi
  \end{minipage}}
  \end{center}
  \bigskip
}
\newcommand{\pregunta}{\stepcounter{testquestion}\item}
% Justificació breu de la resposta correcta d'una pregunta. No es mostra
% al costat de les opcions (per no donar pistes de longitud); s'acumula i
% \clauestatica la imprimeix després de la clau de respostes.
\newcommand{\testjustificacions}{}
\makeatletter
\newcommand{\justificacio}[1]{%
  \edef\@justnum{\thetestquestion}%
  \expandafter\g@addto@macro\expandafter\testjustificacions\expandafter{%
    \expandafter\justitem\expandafter{\@justnum}{#1}}%
}
\makeatother
\newcommand{\justitem}[2]{\item[\textbf{#1.}] #2}
\long\def\opcio#1#2{\item \ifimprimible$\square$\else\CheckBox[name=\testprefix{}q\thetestquestion#1,width=1.1em,height=1.1em,bordercolor={0 0 0},onclick={var opts=['a','b','c','d'];for(var j=0;j<opts.length;j++){if(opts[j]!='#1'){var q=this.getField('\testprefix{}q\thetestquestion'+opts[j]);if(q)q.value='Off';}}}]{}\fi\quad #2}
\newcommand{\feedback}{\ifimprimible\par\medskip\else\par\smallskip\textbf{Correcci\'{o}: }\TextField[name=\testprefix{}f\thetestquestion,readonly=true,width=9cm,bordercolor={0.6 0.6 0.6}]{}\par\medskip\fi}
% ---- Validació estructural de la clau (C7) ----
% \finalitzaTest compara el nombre de preguntes realment plantejades
% (\thetestquestion) amb el nombre de respostes de la clau, i emet un
% avís de compilació si no coincideixen. Cal cridar-la després de
% \end{Form}. Fem servir \@for, que recorre llistes separades per comes
% de manera segura, sense bucles \ifx artesanals.
\newcounter{testkeycount}
\makeatletter
\newcommand{\finalitzaTest}{%
  \setcounter{testkeycount}{0}%
  \expandafter\@for\expandafter\next\expandafter:\expandafter=\testkeys\do{\stepcounter{testkeycount}}%
  \ifnum\value{testquestion}=\value{testkeycount}\else
    \PackageWarning{interactive-tests}{El test '\testprefix' te \thetestquestion\space preguntes pero la clau en te \thetestkeycount\space respostes. Reviseu la crida a \string\iniciTest.}%
  \fi
}
\makeatother

% ============================================================
%  Claus estàtiques dels tests, reunides al final del llibre.
%
%  Els tests del llibre es corregeixen amb JavaScript d'AcroForm,
%  que no funciona en molts visors integrats, navegadors ni
%  aplicacions mòbils. Per a aquests casos, i per a l'ús en paper,
%  cada test té una clau de respostes i unes justificacions. Perquè
%  no apareguin just després de les preguntes, cosa que restaria
%  valor a l'autoavaluació, no s'imprimeixen al lloc del test:
%
%   - \clauestatica (cridada just abans de \finalitzaTest) desa la
%     clau i les justificacions del test en curs a \totesclaus i
%     imprimeix només una remissió a l'apèndix de solucions;
%   - \imprimeixsolucionstests, a caps/solucions-tests.tex, les
%     imprimeix totes al final del llibre, cadascuna amb un enllaç de
%     tornada al test.
% ============================================================

\makeatletter
\newcounter{clauestindex}
\gdef\totesclaus{}
\newcommand{\clauestatica}{%
  \begingroup
    \edef\@tmpclau{\noexpand\blocsoluciotest{\testprefix}{\thechapter}{\testkeys}{\unexpanded\expandafter{\testjustificacions}}}%
    \expandafter\g@addto@macro\expandafter\totesclaus\expandafter{\@tmpclau}%
  \endgroup
  \par\medskip
  \begin{center}
  \small\itshape
  \ifimprimible
  Les respostes i les justificacions d'aquest test són a l'apèndix \hyperref[cap:solucions-tests]{\textup{\enquote{Solucions dels tests}}}, pàgina~\pageref{sol:\testprefix}. Consulteu-les després d'haver respost tot el test.
  \else
  Si el vostre lector de PDF no admet el botó \textup{\textbf{Corregir}}, trobareu les respostes i les justificacions d'aquest test a l'apèndix \hyperref[cap:solucions-tests]{\textup{\enquote{Solucions dels tests}}}, pàgina~\pageref{sol:\testprefix}.
  \fi
  \end{center}
  \par\medskip
}
% #1 prefix del test, #2 número de capítol, #3 clau, #4 justificacions
\newcommand{\blocsoluciotest}[4]{%
  \section*{Test del capítol #2}\label{sol:#1}%
  \addcontentsline{toc}{section}{Test del capítol #2}%
  \noindent\textbf{Respostes.}\quad
  \setcounter{clauestindex}{0}%
  \@for\next:=#3\do{%
    \stepcounter{clauestindex}%
    \mbox{\textbf{\theclauestindex.}~\next}\hspace{1.2em plus 0.4em}%
  }%
  \par
  \def\@tmpjust{#4}%
  \ifx\@tmpjust\@empty\else
    \medskip
    \noindent\textbf{Justificacions.}
    \begin{list}{}{\setlength{\leftmargin}{2.2em}\setlength{\labelwidth}{1.8em}\setlength{\itemsep}{0.2em}\small}
    #4
    \end{list}
  \fi
  \par\smallskip
  \noindent{\small\hyperref[test:#1]{Tornar al test} (pàgina~\pageref{test:#1}).}
  \par\bigskip
}
\newcommand{\imprimeixsolucionstests}{\totesclaus}
\makeatother

% ============================================================
%  Entorns didàctics locals del volum de categories.
%  Definits aquí (no a l'estil compartit) per no acoblar-hi
%  el projecte. Reutilitzen tcolorbox, etoolbox i la paleta
%  ja carregats per aprendes.sty en mode «sobri».
%
%    \begin{objectius} ... \end{objectius}   -> bloc d'obertura
%    \begin{resumcap}   ... \end{resumcap}   -> bloc de tancament
%
%  Dins d'aquests blocs s'hi poden usar els encapçalaments
%  \apartat{...} per separar prerequisits/objectius/etc.
% ============================================================

% Color propi per als blocs didàctics: verd pissarra derivat de la
% paleta existent (cexer = verd), enfosquit per contrastar amb el fons.
\colorlet{DidacticVerd}{cexer!55!black}

% Estils de bloc: mateixa gramàtica visual que blocenunciat/blocobservacio.
\tcbset{
  blocdidactic/.style={
    enhanced, breakable, sharp corners,
    frame hidden, boxrule=0pt,
    colback=DidacticVerd!6,
    borderline west={2.0pt}{0pt}{DidacticVerd},
    left=8pt, right=4pt, top=5pt, bottom=5pt,
    before skip=1.0em, after skip=1.0em,
  },
}

% Encapçalament intern d'un apartat del bloc (Prerequisits, Objectius...).
\newcommand{\apartat}[1]{\par\smallskip{\color{DidacticVerd}\ifpdftex\scshape\else\sffamily\bfseries\fi #1}\par\nobreak\smallskip}

% Bloc d'obertura del capítol.
\newenvironment{objectius}{%
  \begin{tcolorbox}[blocdidactic]%
  \ignorespaces
}{\end{tcolorbox}}

% Bloc de tancament del capítol.
\newenvironment{resumcap}{%
  \begin{tcolorbox}[blocdidactic]%
  {\color{DidacticVerd}\ifpdftex\scshape\else\sffamily\bfseries\fi S\'{\i}ntesi del cap\'{\i}tol.}\hspace{0.6em}%
  \ignorespaces
}{\end{tcolorbox}}


\makeindex[title={Índex alfabètic},columns=2,intoc]

\begin{document}
\ifimprimible\pagecolor{white}\else\pagecolor{fons!20}\fi
\hypermetadata

\frontmatter
\maketitle
% Pàgina de crèdits (revers de la portada).
\thispagestyle{empty}
\vspace*{\fill}
\begin{flushleft}
\small
\textbf{\llibretitol}\\
\llibreautor\\[0.6em]
Versió \llibreversio\ · \llibredata\\
\llibreweb\\[1.2em]

\textbf{Llicència.} © \llibreany\ \llibreautor. Aquesta obra està subjecta a la llicència Creative Commons Reconeixement-CompartirIgual 4.0 Internacional (CC BY-SA 4.0): \url{https://creativecommons.org/licenses/by-sa/4.0/deed.ca}. Podeu copiar-la, redistribuir-la i adaptar-la, també amb finalitat comercial, sempre que en reconegueu l'autoria i distribuïu les obres derivades amb la mateixa llicència.\\[1em]

\textbf{Com citar-lo.} \llibreautor, \emph{\llibretitol}, versió \llibreversio, \llibreany. Projecte aprendes\,/\,logcat, \llibreweb.\\[1em]

\textbf{Errates i suggeriments.} Es poden comunicar obrint una incidència (\emph{issue}) al repositori del projecte, \llibreerrates, indicant la versió del llibre i la pàgina. Les correccions s'incorporen a les versions següents.
\end{flushleft}
\clearpage

\tableofcontents
\let\cleardoublepage\clearpage

\chapter*{Prefaci}
\addcontentsline{toc}{chapter}{Prefaci}

Aquest llibre forma part d'\textbf{aprendes}, una col·lecció de materials
d'estudi que comparteixen una mateixa manera de treballar: cada tema es
presenta en una part de \emph{teoria}, es reprèn en una part de
\emph{pràctica} amb problemes resolts, es consolida amb \emph{exercicis
proposats} i es comprova amb \emph{\nomtests}. Qui ja conegui
altres volums de la col·lecció hi reconeixerà l'estructura; qui s'hi
acosti per primera vegada la trobarà explicada a la guia de lectura.

El volum té un propòsit precís. És el primer d'un projecte en dues
etapes: aquí es construeix el llenguatge de la teoria de categories
---des de les definicions bàsiques fins als classificadors de subobjectes
i els topos--- amb el detall que necessitarà un volum posterior dedicat
a la lògica categòrica. Per això la lògica hi és present des del
començament, com a font d'exemples i com a direcció del recorregut,
tot i que el seu desenvolupament sistemàtic queda per a aquella segona
etapa.

\ifpublicacioparcial
Aquesta és una edició parcial, que conté els tres primers temes
---categories, functors i transformacions naturals--- en les quatre parts,
juntament amb els apèndixs i els materials de consulta. Els capítols
restants s'incorporaran en l'edició completa.
\else
Aquesta edició conté el volum complet: nou temes desenvolupats en les
quatre parts i un capítol final que fa de pont cap a la lògica categòrica.
\fi

Abans del primer capítol, la \emph{Introducció} explica per què val la
pena aprendre aquest llenguatge i la \emph{Guia de lectura} proposa com
fer-ho. Es poden llegir en aquest ordre o tornar-hi quan convingui.

\medskip
\noindent\textbf{Prerequisits.} El text pressuposa la maduresa matemàtica habitual d'un primer curs universitari de matemàtiques, informàtica o lògica: treballar amb conjunts i aplicacions, seguir una demostració elemental i llegir notació algebraica bàsica. No es pressuposa haver estudiat teoria de categories, topologia, àlgebra abstracta avançada ni lògica categòrica. Quan un exemple necessita vocabulari addicional, s'introdueix localment o es remet a un apèndix.

\medskip
\noindent\textbf{Errates i suggeriments.} Són benvinguts. Es poden comunicar obrint una incidència al repositori del projecte, \llibreerrates, indicant la versió (\llibreversio) i la pàgina.

% ------------------------------------------------------------
% Espai reservat per a l'autor (descomenteu i completeu):
%
% \medskip
% \noindent\textbf{Agraïments.} ...
%
% (Les errates i els suggeriments: vegeu la pàgina de crèdits i el text següent.)
% ------------------------------------------------------------


\mainmatter
\chapter*{\texorpdfstring{Per què aprendre\\ teoria de categories?}{Per què aprendre teoria de categories?}}
\addcontentsline{toc}{chapter}{Introducció}
\markboth{Introducció}{Introducció}

\begingroup
\renewcommand{\thesection}{\arabic{section}}
\setcounter{section}{0}
\setlength{\emergencystretch}{1.5em}

Hi ha construccions matemàtiques que aprenem diverses vegades, amb noms i
representacions diferents. Aparellar dues dades, imposar simultàniament dues
condicions, trobar un màxim comú divisor o introduir una conjunció semblen
operacions allunyades. I, tanmateix, totes poden respondre a una mateixa
pregunta: com podem reunir dues exigències en una sola dada, de manera que
qualsevol altra manera de satisfer-les passi per aquesta?

La teoria de categories ensenya a reconèixer aquestes formes comunes. No ho
fa declarant que els conjunts, els nombres i les proposicions són la mateixa
cosa. Ho fa distingint, en cada context, quines relacions són rellevants i
com es poden compondre. El que unifica no és la naturalesa dels objectes,
sinó el paper que tenen dins d'un sistema de relacions.

Aquest llibre proposa aprendre aquest llenguatge des dels seus fonaments,
amb quatre famílies d'exemples que ens acompanyaran al llarg del recorregut:
els conjunts, els preordres, els grafs i la lògica. La finalitat és doble:
adquirir les eines bàsiques de la teoria de categories i comprendre per què
aquestes eines preparen una lectura estructural de la lògica. Abans de
començar amb les definicions, convé veure què ens permeten preguntar.

\section{Una estructura també es coneix per les seves relacions}
\label{sec:mirar-fora}

Una manera habitual d'estudiar un objecte consisteix a mirar-ne l'interior:
els elements d'un conjunt, l'operació d'un grup o les parts d'un graf. La
perspectiva categòrica no elimina aquesta mirada, però la complementa amb
una altra: quines transformacions podem fer entre aquest objecte i els
altres? Quines hi arriben, quines en surten i com s'encadenen?
Podem resumir aquest canvi amb una consigna que ens acompanyarà al llarg
del llibre: \emph{relaciona, no individualitza}. En lloc de fixar un
objecte per allò que és aïlladament, el situem pel paper que té dins
d'una xarxa de transformacions.

Entre conjunts tenim aplicacions; entre grups, homomorfismes; entre
preordres, aplicacions monòtones. Aquestes transformacions tenen una
propietat comuna. Si podem passar d'$A$ a $B$ i de $B$ a $C$, podem fer els
dos passos seguits:
\[
  A\xrightarrow{f}B\xrightarrow{g}C,
  \qquad g\comp f\colon A\longrightarrow C.
\]
La composició és associativa: quan tres passos es poden encadenar, agrupar
els dos primers o els dos darrers no canvia el resultat. A més, cada objecte
té una transformació identitat, que el deixa com és i no altera cap altra
transformació quan s'hi compon.

Una \emph{categoria}\index{categoria} reté justament aquestes dades: objectes, fletxes entre
objectes, composició i identitats, amb les lleis que acabem de descriure.
La paraula \emph{fletxa}, o \emph{morfisme}, és més general que la paraula
\emph{funció}. En la categoria lliure\index{categoria!lliure} sobre un graf dirigit, els objectes
són els vèrtexs i les fletxes són els camins: compondre és concatenar-los,
i la identitat és el camí de longitud zero. Un monoide també es pot veure
com una categoria: té un sol objecte, els elements del monoide en són les
fletxes i la multiplicació en dona la composició.

Ni tan sols els mateixos objectes fixen per si sols la categoria. Amb
conjunts i aplicacions obtenim $\cat{Set}$; amb conjunts i relacions,
$\cat{Rel}$\index{Rel@$\cat{Rel}$}. En aquest segon cas, compondre $R\subseteq A\times B$ amb
$S\subseteq B\times C$ significa que
\[
  a\,(S\comp R)\,c
  \quad\Longleftrightarrow\quad
  \exists b\in B\,\bigl(aRb\land bSc\bigr).
\]
La fletxa composta registra que hi ha algun terme intermedi que connecta
els dos extrems. Ja trobem aquí una operació lògica, la quantificació
existencial, dins d'una operació sobre fletxes.

Per tant, la primera pregunta categòrica no és només \enquote{quins objectes tenim?},
sinó \enquote{en quina categoria els estudiem?}. Els nombres enters positius,
ordenats per $\leq$ o per divisibilitat, ofereixen dues categories diferents:
en cadascuna, una fletxa expressa una relació diferent. Escollir les fletxes
és escollir quina part de la matemàtica volem fer visible.

\section{Del que un objecte conté al paper que compleix}

Considerem el producte cartesià de dos conjunts, $A\times B$. El podem
descriure enumerant-ne els elements: són els parells $(a,b)$ amb $a\in A$
i $b\in B$. Però també el podem descriure sense començar pels parells.
Hi ha dues projeccions,
\[
  \pi_A\colon A\times B\longrightarrow A,
  \qquad \pi_B\colon A\times B\longrightarrow B.
\]
Si des d'un mateix conjunt $X$ tenim una aplicació $f\colon X\to A$ i una
aplicació $g\colon X\to B$, hi ha exactament una aplicació cap a $A\times B$
que reuneix aquestes dues dades: $x\mapsto(f(x),g(x))$.

El punt important és que aquesta segona descripció només parla de fletxes.
Té sentit, doncs, en una categoria qualsevol. Un \emph{producte}\index{producte} d'$A$ i
$B$ és un objecte $P$ amb projeccions $\pi_A\colon P\to A$ i
$\pi_B\colon P\to B$ tal que, per a tot objecte $X$ i totes les fletxes
$f\colon X\to A$ i $g\colon X\to B$, existeix una única fletxa
$u\colon X\to P$ amb
\[
  \pi_A\comp u=f,
  \qquad \pi_B\comp u=g.
\]
Aquestes equacions es fan visibles en el diagrama
\[
\begin{tikzpicture}[baseline=(m.center)]
  \matrix (m) [matrix of math nodes,column sep=3.4em,row sep=3em] {
       & X & \\
     A & P & B \\
  };
  \draw[-{Stealth}] (m-1-2) -- node[above left] {$f$} (m-2-1);
  \draw[-{Stealth}] (m-1-2) -- node[above right] {$g$} (m-2-3);
  \draw[-{Stealth},dashed] (m-1-2) -- node[right] {$u$} (m-2-2);
  \draw[-{Stealth}] (m-2-2) -- node[below] {$\pi_A$} (m-2-1);
  \draw[-{Stealth}] (m-2-2) -- node[below] {$\pi_B$} (m-2-3);
\end{tikzpicture}
\]
La fletxa discontínua és la que la propietat ens permet construir.
Dir que el diagrama \emph{commuta} és dir que els camins que comparem,
amb els mateixos extrems, donen la mateixa fletxa.

En un preordre, posem una única fletxa $x\to y$ quan $x\leq y$, i cap quan
aquesta relació no es compleix. El producte de $a$ i $b$ és aleshores el seu
ínfim: un objecte per sota de tots dos, i tan gran com sigui possible entre
els que tenen aquesta propietat. La mateixa condició produeix les
construccions següents, quan existeixen:

\begin{center}
\renewcommand{\arraystretch}{1.2}
\begin{tabularx}{\linewidth}{@{}>{\raggedright\arraybackslash}X
                                  >{\raggedright\arraybackslash}p{0.32\linewidth}@{}}
\toprule
\textbf{Context i fletxes} & \textbf{Producte} \\
\midrule
Conjunts amb aplicacions & Producte cartesià \\
Subconjunts d'un conjunt fix amb inclusions & Intersecció \\
Enters positius amb divisibilitat & Màxim comú divisor \\
Fórmules amb deduïbilitat & Conjunció \\
\bottomrule
\end{tabularx}
\end{center}

La lectura lògica és especialment reveladora. Fixem un sistema deductiu
proposicional amb les regles usuals de la conjunció i posem una fletxa
$p\to q$ quan $q$ és deduïble de $p$. Les projeccions del producte són les
deduccions de $p$ i de $q$ a partir de $p\land q$. Si des de $r$ podem
deduir totes dues fórmules, podem deduir-ne la conjunció:
\[
  r\vdash_T p\land q
  \quad\Longleftrightarrow\quad
  \bigl(r\vdash_T p\ \text{i}\ r\vdash_T q\bigr).
\]
La unicitat de la fletxa és automàtica perquè només registrem si una
deducció existeix, no quina prova la realitza. El producte expressa aquí
una regla lògica amb el mateix patró que expressa l'aparellament de dades.

Això és una \emph{propietat universal}\index{propietat universal}: una caracterització mitjançant
l'existència i la unicitat de fletxes que satisfan unes condicions.
La paraula \enquote{universal} indica que la condició val per a qualsevol objecte
$X$ i qualsevol parell de fletxes adequat, no només per a uns exemples
triats. Bona part del llibre consisteix a reconèixer aquest patró en
construccions successivament més riques.

\section{La mateixa estructura no vol dir el mateix objecte}

Una propietat universal no ens obliga a triar una representació concreta.
Podem construir dues realitzacions diferents d'un producte i obtenir dos
objectes que no siguin literalment iguals. La propietat universal, però,
força un \emph{isomorfisme}\index{isomorfisme} entre ells compatible amb les projeccions.

Una fletxa $f\colon A\to B$ és un isomorfisme si hi ha una fletxa
$g\colon B\to A$ que desfà el canvi en tots dos sentits:
\[
  g\comp f=\id_A,
  \qquad f\comp g=\id_B.
\]
Escrivim aleshores $A\cong B$. El sentit d'aquest símbol depèn de la
categoria: a $\cat{Set}$ parlem de bijeccions; entre grups, d'isomorfismes
de grups; entre espais topològics amb aplicacions contínues, d'homeomorfismes.
No n'hi ha prou de poder desfer una aplicació entre conjunts si la inversa
no respecta l'estructura que hem fixat.

En la categoria prima de la deduïbilitat, dues fórmules són isomorfes quan
són interdeduïbles. Així, $p\land q$ i $q\land p$ poden tenir sintaxis
diferents i el mateix paper deductiu. El llenguatge categòric no les
declara iguals: precisa en quin sentit la diferència de presentació no
altera l'estructura que estem estudiant.

Cal conservar també les dades que acompanyen l'objecte. L'isomorfisme
\emph{únic} entre dos productes és el que respecta les projeccions;
l'objecte subjacent, sense aquestes dades, pot tenir altres automorfismes.
Aprendre a distingir l'objecte de la construcció completa és part de
l'aprenentatge de les propietats universals.

Un altre canvi de perspectiva consisteix a invertir sistemàticament les
fletxes. La categoria oposada\index{categoria!oposada} $\cat{C}\op$ té els mateixos
objectes que $\cat{C}$, però una fletxa $A\to B$ de $\cat{C}$ hi apareix
com una fletxa $B\to A$. No afirmem que la fletxa original tingui inversa:
canviem el context en què la llegim.

En invertir les fletxes de la propietat universal del producte obtenim
el \emph{coproducte}. A $\cat{Set}$ és la unió disjunta; entre subconjunts
ordenats per inclusió és la unió; en el context deductiu apropiat és la
disjunció. De manera semblant, el dual d'un objecte terminal, que rep una
única fletxa des de qualsevol objecte, és un objecte inicial, que n'envia
una única a cadascun.

La \emph{dualitat}\index{dualitat} és més que una simetria de dibuixos. Quan un argument
està formulat categòricament, en podem obtenir l'argument dual invertint
també les composicions i les hipòtesis. Aprenem així a buscar, al costat
d'una pregunta, la pregunta que es veu des de l'altra direcció.

\section{Comparar contextos i comparar les comparacions}

Si les categories descriuen contextos, com podem passar d'un context a
un altre? Un \emph{functor}\index{functor} $F\colon\cat{C}\to\cat{D}$ assigna objectes
a objectes i fletxes a fletxes, preservant els extrems, les identitats i
la composició:
\[
  F(\id_A)=\id_{F(A)},
  \qquad F(g\comp f)=F(g)\comp F(f).
\]
Un functor és, en aquest sentit precís, una traducció estructurada.
Per exemple, podem passar d'un grup al seu conjunt subjacent i d'un
homomorfisme a l'aplicació que el realitza. També podem passar d'un graf
a la categoria dels seus camins.

La precisió importa: tot functor preserva les lleis categòriques, però
no tot functor preserva productes, límits o qualsevol altra estructura
addicional. El llibre ensenyarà a preguntar què preserva cada traducció,
què oblida i què permet recuperar.

Hi ha encara un altre nivell de comparació. Donats dos functors
$F,G\colon\cat{C}\to\cat{D}$, una \emph{transformació natural}\index{transformació natural}
$\eta\colon F\Rightarrow G$ dona una fletxa
$\eta_A\colon F(A)\to G(A)$ per a cada objecte $A$, amb la condició
que aquestes fletxes siguin compatibles amb tots els morfismes de
$\cat{C}$. Per a $f\colon A\to B$, això vol dir
\[
  G(f)\comp\eta_A=\eta_B\comp F(f).
\]
No és una col·lecció de comparacions independents: totes s'han d'ajustar
a una mateixa xarxa de relacions.

Un exemple elemental és la diagonal dels conjunts,
$\delta_A\colon A\to A\times A$, definida per $a\mapsto(a,a)$.
Copiar una dada i després aplicar $f$ a les dues còpies dona el mateix
resultat que aplicar $f$ primer i copiar després:
\[
\begin{tikzpicture}[baseline=(m.center)]
  \matrix (m) [matrix of math nodes,column sep=4em,row sep=3em] {
    A & B \\
    A\times A & B\times B \\
  };
  \draw[-{Stealth}] (m-1-1) -- node[above] {$f$} (m-1-2);
  \draw[-{Stealth}] (m-1-1) -- node[left] {$\delta_A$} (m-2-1);
  \draw[-{Stealth}] (m-1-2) -- node[right] {$\delta_B$} (m-2-2);
  \draw[-{Stealth}] (m-2-1) -- node[below] {$f\times f$} (m-2-2);
\end{tikzpicture}
\]
La naturalitat expressa que l'operació de copiar és compatible amb
qualsevol aplicació, no amb una tria particular de noms per als elements.
En un exemple lògic del llibre, un quadrat d'aquesta mena expressarà
que substituir variables i després interpretar una fórmula dona el mateix
resultat que interpretar-la i transportar-ne el significat.

Aquest nivell de coherència permet formular també l'\emph{equivalència de
categories}\index{equivalència de categories}. Dues categories $\cat{C}$ i $\cat{D}$ són equivalents si hi
ha functors $F\colon\cat{C}\to\cat{D}$ i $G\colon\cat{D}\to\cat{C}$
tals que anar i tornar és naturalment isomorf a no haver canviat de context:
\[
  G\comp F\cong\id_{\cat{C}},
  \qquad F\comp G\cong\id_{\cat{D}}.
\]
No demanem una igualtat literal de presentacions. Demanem que el canvi
conservi l'estructura categòrica essencial i que la recuperació sigui
coherent. L'isomorfisme compara objectes dins d'una categoria;
l'equivalència compara categories senceres.

\section{Posar a prova la mirada des de fora}

Aquesta secció avança idees que el llibre introduirà pas a pas en els
capítols centrals. No cal retenir-ne ara els detalls: n'hi ha prou de
veure cap a on apunten i com continuen la perspectiva de les seccions
anteriors.

La idea de conèixer un objecte per les seves relacions es pot formular
amb exactitud. Per a un objecte $A$ d'una categoria localment petita,
$\Hom(X,A)$ és el conjunt de fletxes de $X$ cap a $A$. Podem pensar cada
objecte $X$ com una forma de posar $A$ a prova, i cada fletxa $X\to A$
com una dada d'$A$ observada amb aquella forma.

A $\cat{Set}$, prendre $X$ amb un sol element recupera els elements
ordinaris d'$A$. En altres categories, però, aquesta única mena de prova
pot resultar insuficient. La perspectiva general considera \emph{tots}
els objectes $X$ i, a més, registra com una prova canvia quan precomponem
amb una fletxa $X'\to X$.

Aquesta informació s'organitza en l'homfunctor\index{homfunctor} $\Hom(-,A)$.
El lema de Yoneda\index{Yoneda!lema} donarà forma precisa a una intuïció anunciada des
del començament: els objectes queden determinats, llevat d'isomorfisme,
per aquest sistema de relacions. No n'hi ha prou de saber quantes
fletxes hi arriben; cal conservar la manera com es transformen i es
componen. Dos homfunctors d'aquest tipus naturalment isomorfs provenen
d'objectes isomorfs.

La \emph{representabilitat} aplica aquesta idea a les construccions.
En el cas del producte, donar una fletxa $X\to A\times B$ equival a
donar dues fletxes amb domini $X$:
\[
  \Hom(X,A\times B)
  \cong \Hom(X,A)\times\Hom(X,B),
\]
de manera compatible amb els canvis de $X$. El producte representa,
amb un sol objecte, l'assignació que a cada $X$ associa aquestes parelles.
Límits i colímits reuniran moltes propietats universals sota un mateix
patró; no seran una llista de receptes sense connexió.

Les \emph{adjuncions}\index{adjunció} ens permetran fer un pas més i relacionar
construccions senceres. Si $F\colon\cat{C}\to\cat{D}$ i
$G\colon\cat{D}\to\cat{C}$ són adjunts, amb $F$ per l'esquerra de
$G$, hi ha una correspondència natural
\[
  \Hom_{\cat{D}}(F(A),B)
  \cong \Hom_{\cat{C}}(A,G(B)).
\]
Una mateixa dada es pot descriure a banda i banda del canvi de context.
Per exemple, donar un functor des de la categoria lliure d'un graf cap
a una categoria és equivalent a donar un morfisme del graf inicial cap
al graf subjacent d'aquella categoria. Els camins compostos no s'han
d'escollir de nou: la composició en determina la imatge.

En un altre exemple, donar una aplicació $X\times A\to B$ equival a
donar una aplicació $X\to B^A$, on $B^A$ és el conjunt d'aplicacions
d'$A$ a $B$. És la \emph{currificació}: una funció de dues variables
es pot descriure com una funció que retorna funcions. El tancament
cartesià generalitza aquesta correspondència i ens porta cap a la
semàntica del càlcul lambda simplement tipat.

\section{La lògica com a fil conductor, no com a drecera}

La lògica apareix en aquest llibre des del primer capítol, però a
diferents nivells. El més elemental és la categoria prima de la
deduïbilitat: les fórmules són objectes i una fletxa registra que una
fórmula es pot deduir d'una altra. Aquesta lectura ajuda a reconèixer
la conjunció com a producte i la interdeduïbilitat com a isomorfisme.

Conservar les proves mateixes és una tasca més exigent. Cal decidir
quan dues proves representen la mateixa fletxa i com la composició
respecta aquesta identificació. No podem passar de l'existència d'una
deducció a una categoria de proves sense afegir aquestes precisions.
La interpretació de tipus i termes en categories cartesianes tancades
ens permetrà començar a veure la relació entre proves, programes i
morfismes, sense confondre aquests nivells.

Més endavant, els \emph{subobjectes}\index{subobjecte} oferiran una altra lectura lògica.
A $\cat{Set}$, un predicat sobre $A$ es pot descriure pel subconjunt
d'elements que el satisfan. Si $f\colon X\to A$, substituir $f(x)$
en el predicat correspon a prendre'n l'antiimatge. El \emph{pullback}\index{pullback}
generalitza aquesta operació, i la reindexació de subobjectes expressa
la substitució sense haver de parlar dels elements d'una categoria
arbitrària.

Les imatges i les adjuncions recuperen llavors l'existencial que havíem
trobat en compondre relacions. Projectar un predicat sobre $X\times A$
cap a $X$ significa retenir els $x$ per als quals hi ha algun $a$ que
el satisfà. En una categoria regular, aquesta construcció dona un
adjunt per l'esquerra de la reindexació. L'adjunt per la dreta, que
interpretaria el quantificador universal, requereix estructura
addicional: no es dedueix de la regularitat per si sola.

Finalment, el \emph{classificador de subobjectes} generalitza una altra
operació familiar. Un subconjunt $S\subseteq A$ es pot descriure per
la seva funció característica $A\to\{0,1\}$. En un topos elemental,
un objecte $\Omega$ fa aquest paper: els morfismes $A\to\Omega$
classifiquen els subobjectes d'$A$. Aquí es retroben dues línies del
llibre, la representabilitat i la lectura dels subobjectes com a
predicats. No hem d'esperar, però, que $\Omega$ es comporti sempre
com els dos valors de veritat clàssics de $\cat{Set}$.

Aquesta és la destinació del volum: arribar al llindar de la lògica
categòrica amb les eines necessàries per comprendre'n el llenguatge.
El desenvolupament sistemàtic de la sintaxi, la semàntica, la correcció,
la completesa i les condicions de coherència queda reservat al volum
següent. El pont final n'explica la continuació; no substitueix aquest
desenvolupament. La lògica dona direcció al recorregut, però no ens
eximeix d'aprendre les construccions categòriques amb les seves
hipòtesis i demostracions.

\section{Què vol dir haver après aquest llenguatge}

Aprendre categories no consisteix a abandonar els exemples concrets
tan aviat com tenim una definició general. Consisteix a poder fer el
camí d'anada i tornada: reconèixer una estructura en un exemple,
formular-la amb fletxes i tornar a interpretar-la en un altre context.
Una bona abstracció ens ha de permetre veure més, no només dir menys.

Al llarg del llibre adquirirem l'hàbit de preguntar-nos:
\begin{itemize}
  \item Quins són els objectes i les fletxes, i què significa compondre-les?
  \item Quina propietat expressa el diagrama que tenim davant?
  \item Quines dades caracteritzen una construcció i quines en són
        només una realització particular?
  \item Què queda determinat llevat d'isomorfisme, i quines
        compatibilitats fan únic aquest isomorfisme?
  \item Què preserva el functor amb què canviem de context?
  \item Quin és l'enunciat dual i sota quines hipòtesis té sentit?
  \item Quina lectura lògica té la construcció, si el context la permet?
\end{itemize}

Els primers capítols fixaran el vocabulari; els següents ensenyaran
a organitzar-lo amb propietats universals, representabilitat, límits
i adjuncions; els darrers mostraran com aquestes eines es combinen
en el tancament cartesià, els subobjectes i els topos. La guia de
lectura que acompanya el llibre orienta sobre com treballar aquest
recorregut amb la teoria, la pràctica, els exercicis i els tests.

El resultat que busquem no és disposar d'una paraula nova per a cada
construcció coneguda. És poder reconèixer, davant d'una situació nova,
quines relacions en governen l'estructura. Quan una intersecció, una
conjunció i un producte deixen de ser semblances vagues i esdevenen
manifestacions d'una mateixa propietat, el llenguatge categòric
comença a fer la seva feina.

\endgroup

\chapter*{\texorpdfstring{Guia de lectura i itineraris\\ d'aprenentatge}{Guia de lectura i itineraris d'aprenentatge}}
\addcontentsline{toc}{chapter}{Guia de lectura}
\markboth{Guia de lectura}{Guia de lectura}

\begingroup
\renewcommand{\thesection}{\arabic{section}}
\setcounter{section}{0}
\setlength{\emergencystretch}{1.5em}

Aquest capítol és un mapa per a la primera lectura. No cal aprendre
ara els noms de totes les construccions que hi apareixen: les trobarem
definides al seu lloc. La seva funció és mostrar què aporta cada etapa,
com es relacionen les quatre parts del llibre i com podem comprovar
que l'estudi ens està fent avançar.

La introducció explica \emph{per què} val la pena aprendre categories.
Aquí ens ocuparem de \emph{com} fer-ho. Podeu llegir aquesta guia
abans de començar i tornar-hi en acabar un bloc de capítols, quan
els noms del mapa ja s'hagin convertit en construccions familiars.


\section{El punt de partida i la destinació}

El llibre no pressuposa coneixements previs de teoria de categories.
Sí que pressuposa una familiaritat elemental amb el llenguatge
matemàtic: conjunts, aplicacions, composició, relacions i
demostracions senzilles. Per començar, convé poder fer quatre coses:
\begin{itemize}
  \item distingir el domini i el codomini d'una aplicació i calcular
        una composició;
  \item treballar amb productes cartesians, subconjunts i antiimatges;
  \item llegir quantificadors i separar una afirmació d'existència
        d'una afirmació d'unicitat;
  \item distingir una implicació, el seu recíproc i una equivalència,
        i entendre el paper d'un contraexemple.
\end{itemize}
Si alguna d'aquestes operacions és poc familiar, és millor dedicar-hi
una breu revisió que atribuir-ne la dificultat a l'abstracció categòrica.

No cal dominar prèviament tots els àmbits que proporcionen exemples.
Els grups, els espais vectorials i la topologia enriqueixen la lectura,
però, si encara no els coneixeu, podeu començar amb conjunts, preordres
i grafs. Els exemples del dual i del bidual demanen àlgebra lineal;
els que utilitzen continuïtat demanen nocions de topologia. Es poden
reprendre en una segona lectura sense abandonar el fil general.
La lògica aporta exemples des del començament, però no cal conèixer
ja la lògica categòrica per seguir-los.

La destinació és una comprensió inicial de les categories, els
functors i la naturalitat, seguida de les propietats universals,
Yoneda, els límits i les adjuncions. Aquest bagatge permet arribar
al tancament cartesià, als subobjectes i a una primera introducció
als topos. El volum prepara l'estudi posterior de la lògica categòrica;
no pretén desenvolupar-ne aquí una teoria completa.

\section{Quatre parts per a un mateix aprenentatge}

El llibre separa físicament la \textbf{Teoria}, la \textbf{Pràctica},
els \textbf{Exercicis proposats} i els \textbf{\NomTests}.
Aquesta separació facilita la consulta, però no obliga a estudiar
primer tota la teoria i només després tota la pràctica. Per a una
primera lectura, és més útil treballar \emph{per temes}: llegir un
capítol teòric, estudiar els problemes resolts corresponents,
intentar exercicis i comprovar la comprensió amb el test del mateix
tema.

\begin{center}
\renewcommand{\arraystretch}{1.2}
\begin{tabularx}{\linewidth}{@{}>{\raggedright\arraybackslash}p{0.25\linewidth}
                                  >{\raggedright\arraybackslash}X@{}}
\toprule
\textbf{Part} & \textbf{La pregunta que ajuda a respondre} \\
\midrule
Teoria & Quina és la definició, per què té sentit i què se'n pot deduir? \\
Pràctica & Com s'utilitza en un problema i com s'escriu l'argument? \\
Exercicis proposats & Puc reconèixer i aplicar la idea sense una solució completa? \\
\NomTests & Distingeixo el concepte de les confusions que s'hi assemblen? \\
\bottomrule
\end{tabularx}
\end{center}

Els nou primers temes tenen capítols corresponents en les quatre
parts. La numeració torna a començar en cada part; per això, quan
cerqueu material d'un tema, fixeu-vos tant en el número com en la
part i el títol. El capítol teòric desè, \emph{Pont cap a la lògica
categòrica}, és un epíleg d'orientació i no té un capítol paral·lel
de pràctica, exercicis o tests.

Els blocs d'obertura dels capítols teòrics indiquen prerequisits
i objectius. Llegiu-los com una invitació a activar coneixements
anteriors. Els blocs de síntesi del final recullen les idees centrals,
els errors típics i lectures recomanades. Serveixen per revisar,
però no substitueixen els exemples i les demostracions del capítol.

\section{El mapa conceptual dels capítols}
\label{sec:mapa-lectura}

\ifpublicacioparcial
\noindent\textbf{Nota sobre aquesta edició.} Aquesta versió parcial
només inclou els capítols 1--3. El mapa descriu, tanmateix, els deu
capítols del volum complet, perquè pugueu situar el que estudieu
dins del recorregut sencer.
\fi

El recorregut es pot agrupar en quatre etapes. Els capítols 1--3
construeixen el llenguatge i els nivells de comparació; els 4--6
organitzen les construccions mitjançant propietats universals;
els 7--9 combinen aquestes eines en estructures amb lectura lògica;
el 10 mostra la continuació. Aquesta divisió és una ajuda per orientar-se,
no una frontera rígida: les idees de les primeres etapes reapareixen
constantment en les següents.

\subsection*{1. Categories: aprendre la gramàtica de les fletxes}

El primer capítol introdueix objectes, morfismes, identitats i
composició. També ensenya a llegir diagrames, construeix exemples
i diferencia isomorfismes, monomorfismes, epimorfismes, seccions
i retraccions. La categoria oposada i la dualitat permeten
reutilitzar definicions i arguments en l'altra direcció.

És un capítol extens perquè fixa els hàbits que necessitarem després.
No convé voler memoritzar tots els exemples alhora. Escolliu-ne uns
quants i pregunteu sempre què són els objectes, què són les fletxes
i com es componen. Les construccions de categories noves i les
precaucions de mida amplien aquest vocabulari.

\noindent\textbf{Comprovació de progrés.} Podeu verificar els axiomes
en un exemple concret i explicar per què una fletxa és, o no és,
un isomorfisme? Podeu escriure l'equació d'un quadrat commutatiu
sense invertir l'ordre de la composició?

\subsection*{2. Functors: passar d'un context a un altre}

El segon capítol puja un nivell: no compara només objectes d'una
categoria, sinó categories entre elles. Els functors actuen sobre
objectes i morfismes i preserven identitats i composició. El graf
subjacent i la categoria lliure permeten veure una parella de
construccions que reapareixerà com a adjunció.

Cal dedicar una atenció especial als homfunctors. $\Hom(A,-)$
actua per postcomposició; $\Hom(-,B)$, per precomposició i amb
inversió de sentit. Aquesta variància no és un detall de notació:
és la base de la representabilitat i de Yoneda. Les nocions de
functor ple i fidel descriuen l'acció sobre cada homconjunt.

\noindent\textbf{Comprovació de progrés.} Per definir un functor,
podeu donar-ne l'acció sobre les fletxes, no només sobre els objectes?
Podeu explicar què fa $\Hom(-,B)$ amb una fletxa $X'\to X$?

\subsection*{3. Transformacions naturals i equivalències: exigir coherència}

El tercer capítol compara functors. La naturalitat exigeix que
una família de morfismes sigui compatible amb totes les fletxes
del domini. Els exemples de la interpretació de fórmules i del
bidual mostren per què aquesta compatibilitat és matemàticament
significativa, no només una convenció gràfica.

Les categories de functors permeten tractar aquestes comparacions
com a morfismes. L'equivalència de categories precisa quan dos
contextos descriuen la mateixa estructura essencial sense tenir
la mateixa presentació. La composició horitzontal afegeix un
segon tipus de composició que cal distingir de la vertical.

\noindent\textbf{Comprovació de progrés.} Podeu dibuixar el quadrat
de naturalitat d'una família proposada i comprovar-lo per a una
fletxa arbitrària? Podeu distingir un isomorfisme d'objectes,
un isomorfisme natural i una equivalència de categories?

\subsection*{4. Propietats universals, representabilitat i Yoneda:
caracteritzar pel paper}

Aquí canvia el centre de gravetat. En lloc de construir un objecte
i estudiar-ne després les propietats, busquem un objecte que
resolgui una necessitat expressada amb fletxes. Els objectes
inicials i terminals ofereixen els primers casos; les categories
coma ajuden a reunir dades i compatibilitats en un mateix context.

La representabilitat expressa aquestes caracteritzacions amb
homfunctors. Yoneda fa precisa la idea de conèixer un objecte
pel sistema de morfismes que hi arriben. És normal que aquesta
etapa demani més d'una lectura: cal coordinar categories, functors,
transformacions i elements d'homconjunts.

\noindent\textbf{Comprovació de progrés.} Podeu demostrar que dos
objectes terminals són isomorfs i indicar on s'utilitza la unicitat?
En una representació, podeu distingir el functor representat,
l'objecte representant i la dada universal? Podeu seguir la
bijecció de Yoneda a partir de l'avaluació en la identitat?

\subsection*{5. Límits i colímits: reunir construccions en un patró}

Un diagrama fixa una configuració d'objectes i morfismes; un con
és una manera compatible de relacionar-s'hi. El límit és el con
universal. Segons la forma del diagrama, obtenim productes,
equalitzadors, pullbacks o l'objecte terminal. Els colímits en
són la versió dual.

El pullback mereix una atenció especial, perquè reapareixerà
en la reindexació i els subobjectes. També cal separar dues
preguntes: si una categoria té un límit i si un functor el
preserva. No són la mateixa afirmació.

\noindent\textbf{Comprovació de progrés.} Podeu calcular un pullback
de conjunts i comprovar-ne la propietat universal? Podeu
distingir el diagrama, el con i el seu vèrtex, i explicar
què vol dir que un functor preserva aquest límit?

\subsection*{6. Adjuncions: dues descripcions de la mateixa dada}

Les adjuncions relacionen functors mitjançant una bijecció
natural d'hom\-con\-junts. El capítol comença amb preordres,
on la correspondència es llegeix com una equivalència de
desigualtats. Després introdueix unitat, counitat i identitats
triangulars, i mostra com les adjuncions expliquen la preservació
de límits i colímits.

Per a una primera aproximació, escolliu una adjunció i seguiu-la
en les dues descripcions: per bijeccions d'homconjunts i per
unitat i counitat. Les construccions lliures i la currificació
connecten aquest capítol amb exemples ja coneguts.

\noindent\textbf{Comprovació de progrés.} Podeu dir quins morfismes
es corresponen en una adjunció concreta i construir-ne els
transposats? Podeu recordar, amb una justificació conceptual,
que els adjunts per la dreta preserven límits i els de
l'esquerra preserven colímits?

\subsection*{7. Categories cartesianes tancades: funcions dins del context}

Els productes i els exponencials permeten parlar d'aparellament,
avaluació i currificació dins d'una categoria. El capítol
caracteritza el tancament cartesià i n'explica la relació
amb el càlcul lambda simplement tipat: els tipus es llegeixen
com a objectes i els termes, en context, com a morfismes.

El conjunt de totes les aplicacions $B\to C$ ofereix l'exemple
més familiar d'exponencial. Els preordres permeten reconèixer
la implicació com a adjunt de la conjunció. Els contraexemples
recorden que no tota categoria té aquesta estructura.

\noindent\textbf{Comprovació de progrés.} Podeu passar d'una fletxa
$A\times B\to C$ a la seva currificació $A\to C^B$ i tornar-ne
amb l'avaluació? Podeu explicar per què tenir productes no
garanteix tenir exponencials?

\subsection*{8. Subobjectes, imatges i factoritzacions: de les parts als predicats}

Un subobjecte no és simplement un subconjunt escrit amb un
altre nom: és una classe de monomorfismes amb un mateix
codomini, identificats de manera compatible. La reindexació
per pullback generalitza l'antiimatge i ofereix una lectura
estructural de la substitució.

Les interseccions expressen conjuncions de predicats; les
imatges i la regularitat permeten construir un adjunt
existencial de la reindexació. Aquesta etapa obliga a
controlar les hipòtesis: la regularitat no garanteix per
si sola unions de subobjectes ni un quantificador universal.

\noindent\textbf{Comprovació de progrés.} Podeu calcular una
reindexació a $\cat{Set}$ i interpretar-la com a substitució?
Podeu descriure la imatge d'un predicat sota una projecció
i explicar la relació $\exists_f\dashv f^*$?

\subsection*{9. Classificadors de subobjectes i topos: retrobar la representabilitat}

El classificador generalitza la funció característica d'un
subconjunt. La correspondència entre subobjectes d'$A$ i
morfismes $A\to\Omega$ recupera la representabilitat del
capítol quart. El topos elemental combina límits finits,
exponencials i classificador de subobjectes.

El capítol presenta $\cat{Set}$ i les categories de prefeixos
com a exemples. L'objectiu és comprendre les definicions
i la connexió entre les peces, no dominar encara la
teoria de feixos o la lògica interna dels topos.

\noindent\textbf{Comprovació de progrés.} Podeu reconstruir un
subconjunt com l'antiimatge del valor vertader per la
seva funció característica? Podeu enunciar les tres
condicions d'un topos i explicar què aporta cadascuna?

\subsection*{10. Pont cap a la lògica categòrica: saber què queda preparat}

L'epíleg situa les eines adquirides en un programa més ampli:
contextos, substitució, predicats, quantificadors, teories
i models. El mapa de fragments lògics i l'exemple de la
categoria sintàctica orienten sobre què es prolonga al
volum següent.

Llegiu-lo com una mirada retrospectiva i una invitació a
continuar, no com una exigència de dominar de cop tots els
temes anunciats. La sintaxi general, la correcció, la
completesa i les condicions de coherència pertanyen a
l'estudi posterior.

\noindent\textbf{Comprovació de progrés.} Quan l'epíleg parla de
models com a functors o de quantificadors com a adjunts,
podeu identificar les construccions d'aquest volum que
donen sentit a aquestes expressions?

\section{Com treballar un capítol}

Una lectura activa alterna comprensió, reconstrucció i aplicació.
Podeu utilitzar el cicle següent, adaptant-ne el ritme a la
dificultat del tema.

\begin{enumerate}
  \item \textbf{Situar la pregunta.} Llegiu l'obertura i els
        objectius del capítol. Escriviu en una frase quin
        problema intenta resoldre i quines idees anteriors
        necessita.
  \item \textbf{Desmuntar la definició.} Separeu les dades
        que s'han de donar, les operacions que s'hi afegeixen
        i les propietats que han de satisfer. Anoteu el
        domini i el codomini de cada fletxa.
  \item \textbf{Ancorar-la en exemples.} Treballeu un cas
        amb conjunts i un altre amb preordres o grafs.
        Quan pertoqui, afegiu-hi la lectura lògica. Busqueu
        també un exemple que no compleixi la definició.
  \item \textbf{Reconstruir una demostració.} Tanqueu el
        text i intenteu recuperar un argument curt. Després
        compareu-lo amb el llibre: on s'usa cada hipòtesi?
        On apareix l'existència i on la unicitat?
  \item \textbf{Fer servir la pràctica amb pausa.} Llegiu
        l'enunciat d'un problema resolt i proveu de fer-ne
        el primer pas abans de mirar la solució. Un cop
        llegida, torneu a escriure l'argument sense copiar-lo.
  \item \textbf{Intentar un exercici nou.} Comenceu pels
        de consolidació. Si us encalleu, identifiqueu el
        punt exacte i consulteu la indicació, no només
        una altra solució que s'hi assembli.
  \item \textbf{Tancar el cicle.} Feu el test, reviseu els
        errors i rellegiu la síntesi. Escriviu què podeu
        fer ara que no podíeu fer abans i què ha de quedar
        anotat per a una segona lectura.
\end{enumerate}

No cal reconstruir immediatament totes les proves llargues.
Però convé entendre'n l'estratègia i ser capaç de reproduir
arguments curts que reapareixen: provar una naturalitat,
usar la unicitat d'una propietat universal o construir
una fletxa amb el domini i el codomini correctes.

\subsection*{Una pauta per llegir diagrames}

Al llarg del llibre, quan una propietat categòrica té una expressió
natural amb fletxes, la presentem alhora amb un diagrama i amb les
equacions de composició que aquest expressa. Els diagrames formen
part del llenguatge matemàtic, no són una il·lustració opcional.
Aprofiteu aquesta doble presentació: passar del dibuix a les equacions
i de les equacions al dibuix és una destresa que cal practicar.

Davant d'un diagrama, no comenceu per recordar-ne el nom.
Seguiu aquesta pauta:
\begin{enumerate}
  \item Identifiqueu els objectes i els extrems de cada fletxa.
  \item Distingiu les dades donades de la fletxa que s'ha de construir.
  \item Escriviu els camins que tenen el mateix origen i el mateix destí.
  \item Traduïu-ne la commutativitat en equacions de composició.
  \item Si hi ha una propietat universal, digueu sobre quines
        dades es quantifica i què es demana que sigui únic.
\end{enumerate}

Una fletxa discontínua en un diagrama universal no es construeix
perquè el dibuix ho suggereixi: l'existència es justifica amb
una hipòtesi o una propietat. Igualment, un quadrat commutatiu
no és automàticament un pullback. Cal comprovar-ne també la
propietat universal.

\section{Els exemples recurrents com a punts de suport}

Canviar d'exemple no ha de significar tornar a començar.
Les mateixes famílies permeten veure cares diferents d'una idea.

\begin{description}
  \item[Conjunts.] Permeten calcular explícitament amb elements,
        productes, antiimatges, funcions i imatges. Són un
        primer laboratori, però no autoritzen a usar elements
        dins d'una categoria arbitrària.
  \item[Preordres.] Simplifiquen els homconjunts: hi ha com a
        màxim una fletxa entre dos objectes. Això converteix
        propietats universals en condicions d'ordre i adjuncions
        en equivalències de desigualtats. La unicitat automàtica
        és una ajuda, però també pot amagar una dificultat
        que reapareix en categories no primes.
  \item[Grafs.] Fan visible la diferència entre tenir fletxes
        i disposar de composició. Cal distingir la categoria
        de grafs, on les fletxes són homomorfismes de grafs,
        de la categoria lliure d'un graf, on són camins.
        La construcció lliure connecta els primers capítols
        amb les propietats universals i les adjuncions.
  \item[Lògica.] Permet interpretar deducció, conjunció,
        substitució i quantificació en els contextos adequats.
        Cal separar sempre la categoria prima que registra
        deduïbilitat de la perspectiva que conserva les proves,
        i també dels subobjectes que representen predicats.
\end{description}

Un quadern de correspondències pot ajudar: per a cada nova
construcció, escriviu què significa en dues d'aquestes famílies,
quina propietat comparteixen i quines hipòtesis necessita
cada realització. Si no existeix en un dels exemples, anoteu-ho:
aquest límit també forma part de la comprensió.

\section{Itineraris possibles}

Els itineraris següents són propostes de treball, no permisos
per ignorar les dependències. L'ordre del llibre continua sent
el recorregut més segur per a una primera lectura completa.

\subsection*{Una primera entrada: els tres capítols inicials}

Treballeu categories, functors i transformacions naturals amb
els seus materials paral·lels de pràctica, exercicis i tests.
La fita és poder moure-us entre els tres nivells: objectes i
morfismes, categories i functors, functors i transformacions.
No cal dominar encara tot el desplegament de la composició
horitzontal, però sí distingir-la de la vertical.

\ifpublicacioparcial
Aquesta edició parcial conté precisament aquests tres temes en
les quatre parts, a més dels materials complementaris. No és una
introducció completa als límits, les adjuncions o els topos;
les remissions als capítols absents s'han d'entendre com a
indicacions per continuar amb el volum complet.
\fi

\subsection*{Un recorregut complet d'autoaprenentatge}

Seguiu els nou temes amb el cicle de treball per capítols i
llegiu després el pont final. Cada cop que arribeu a una
definició general, torneu a un exemple anterior. Després
de Yoneda, reviseu els homfunctors; després de les adjuncions,
reviseu la construcció lliure; després del classificador,
reviseu la representabilitat. Aquestes tornades no són
repeticions inútils: fan visible la unitat del recorregut.

\subsection*{Un curs semestral}

Per a un curs d'un semestre, una selecció natural és treballar els capítols 1--6 com a nucli obligatori i triar, segons l'orientació del curs, els capítols 7--9. La teoria s'ha de combinar cada setmana amb una selecció de problemes resolts i exercicis proposats; els tests poden servir d'autoavaluació, no de substitut d'una prova escrita. El capítol 10 funciona bé com a sessió final de síntesi i d'obertura cap a la lògica categòrica.

\subsection*{Una lectura per al professorat}

Per preparar docència, convé començar per aquesta guia i pels blocs d'objectius i d'errors típics de cada capítol. Les parts de Pràctica permeten seleccionar exemples per a classe, mentre que els Exercicis proposats ofereixen material de treball autònom. Els quatre fils recurrents ---conjunts, preordres, grafs i lògica--- es poden modular segons el perfil de l'alumnat, però és recomanable conservar-ne almenys dos en paral·lel per evitar que les definicions quedin lligades a una sola categoria concreta.

\subsection*{Una lectura orientada a la lògica}

Manteniu el recorregut general, però doneu més pes a la
deduïbilitat del capítol 1 i als exemples d'interpretació
i substitució dels capítols 2 i 3. Als capítols 4--6,
consolideu representabilitat, pullbacks i adjuncions;
després, treballeu amb especial cura el tancament cartesià,
els subobjectes i el classificador.

No convé saltar directament a l'epíleg: expressions com
\enquote{el quantificador és un adjunt} orienten, però només es
comprenen quan sabem entre quines categories o preordres
actua l'adjunció i què afirma la seva correspondència.
Les equivalències lògiques no substitueixen la comprovació
de les hipòtesis categòriques.

\subsection*{Una lectura orientada als tipus i a la computació}

Feu servir els grafs i les categories lliures per entendre
la composició. Després de consolidar functors i naturalitat,
poseu l'accent en propietats universals, productes,
adjuncions, exponencials i currificació. El capítol 7
reuneix aquestes eines en la semàntica de tipus i termes.

Els capítols 8 i 9 poden constituir una segona etapa
si el primer objectiu és el càlcul lambda simplement tipat.
Caldrà reprendre'ls per arribar als predicats, els
classificadors i la perspectiva lògica més àmplia.

\section{Convencions tipogràfiques i de lectura}

Els noms de categories s'escriuen amb tipografia matemàtica, com $\cat{Set}$, $\cat{Graph}$ o $\cat{Cat}$. Els termes definits i les idees que cal retenir apareixen destacats dins dels entorns corresponents; el color reforça la jerarquia visual, però l'etiqueta textual ---\enquote{Definició}, \enquote{Teorema}, \enquote{Exemple}, \enquote{Exercici}, etc.--- n'indica sempre la funció. Les fletxes discontínues dels diagrames assenyalen habitualment el morfisme que una propietat universal afirma que existeix; el text adjacent n'explicita el paper. Els símbols amb més d'un ús, especialment $f^{*}$, es desambigüen pel context i es recullen a la taula de notació.

\section{Exercicis i tests: què comproven i què no}

Els exercicis proposats utilitzen una estrella,
$\star$, per assenyalar consolidació, i dues,
$\star\star$, per assenyalar ampliació. Comenceu pels
primers, però no interpreteu les marques com una mesura
exacta del temps que necessitareu. Una dificultat breu
de variància pot resultar més important que un càlcul llarg.

Les indicacions són una ajuda per reprendre el treball.
Abans de consultar-les, escriviu les dades, la conclusió
i almenys un intent. Després d'utilitzar una indicació,
torneu a fer l'exercici sense mirar-la. L'objectiu no és
haver reconegut una solució, sinó poder construir un argument.

Els tests tenen una única resposta correcta per pregunta.
\ifimprimible
En aquesta edició per imprimir, les respostes i les
justificacions de cada test es reuneixen al final del llibre
(apèndix \enquote{Solucions dels tests}). L'edició de pantalla
ofereix, a més, correcció automàtica en els lectors de PDF
compatibles.
\else
La correcció automàtica necessita un lector de PDF compatible
amb formularis AcroForm i JavaScript. Si el lector no
l'admet, podeu respondre igualment les preguntes i
consultar les respostes i les justificacions de cada test,
reunides al final del llibre (apèndix \enquote{Solucions dels
tests}), també usables en paper.
\fi
Consulteu-les només després d'haver respost tot el test.

Una resposta encertada no demostra per si sola que sapigueu
resoldre un problema. Per aprofitar el test, justifiqueu
per què l'opció triada és correcta i per què les altres
fallen. Si heu dubtat entre dues opcions, convertiu el
dubte en una pregunta per revisar a la teoria o a la
pràctica. Aquesta explicació és més útil que la puntuació
obtinguda sense justificació.

\section{Algunes dificultats que convé reconèixer aviat}

Hi ha errors que travessen diversos capítols. Identificar-los
permet corregir la manera de treballar, no només una resposta.

\begin{description}
  \item[Compondre sense comprovar els tipus.]
        Abans d'escriure $g\comp f$, verifiqueu que el
        codomini de $f$ és el domini de $g$. L'associativitat
        no fa componibles fletxes amb extrems incompatibles.
  \item[Confondre un exemple amb la definició general.]
        A $\cat{Set}$, mono i epi corresponen a injectivitat
        i exhaustivitat. En altres categories, les definicions
        són condicions de cancel·lació. Un mono que també
        és epi no és sempre un isomorfisme.
  \item[Tractar la naturalitat com una etiqueta.]
        Una família de fletxes no és natural només perquè
        s'hagi definit amb la mateixa fórmula. Cal escriure
        el quadrat i comprovar-lo per a tots els morfismes.
  \item[Oblidar les dades universals.]
        El producte inclou les projeccions; el límit inclou
        el con; l'exponencial inclou l'avaluació. La unicitat
        de l'isomorfisme es refereix a les compatibilitats
        amb aquestes dades, no a l'objecte nu.
  \item[Canviar l'orientació sense controlar-la.]
        La contravariància, la dualitat i la reindexació
        inverteixen direccions de maneres que cal precisar.
        Escriviu dominis i codominis abans de fer servir
        una regla recordada de memòria.
  \item[Afegir estructura que no s'ha suposat.]
        Tenir productes no implica tenir exponencials;
        la regularitat no dona per si sola disjuncions
        ni quantificació universal; un topos no és
        necessàriament booleà. Rellegiu les hipòtesis.
\end{description}

Si una demostració sembla bloquejada, pregunteu-vos primer
si el problema és de definició, de tipus, de variància o
d'hipòtesis. Moltes dificultats aparentment abstractes
es resolen quan precisem una d'aquestes quatre coses.

\section{Quan consultar els apèndixs i els materials finals}

Els apèndixs no constitueixen una etapa que s'hagi de
completar obligatòriament abans de començar la teoria.
Es poden consultar quan la lectura ho demani.

\begin{description}
  \item[Apèndix A: axiomàtica relacional.]
        Ofereix una presentació en què les fletxes són
        els individus primitius i els objectes es
        reconstrueixen a partir de les identitats.
        També hi trobareu la formalització del principi
        de dualitat, en versió sintàctica i semàntica.
        És una ampliació adequada després del primer
        capítol, sobretot per a qui s'interessi per
        la formalització lògica. Demana familiaritat
        amb lògica de primer ordre amb igualtat.
  \item[Apèndix B: grafs.]
        Aclareix les convencions, els homomorfismes,
        els camins i el pas a categories. Consulteu-lo
        quan treballar amb camins, llaços o arestes
        múltiples us plantegi dubtes en els primers
        capítols.
  \item[Apèndix C: convenció de mida.]
        Explica el marc de conjunts i classes del llibre.
        És especialment útil quan apareixen
        $\cat{Cat}$, categories de functors i
        categories de prefeixos. Cal distingir
        \enquote{petita} de \enquote{localment petita}: la segona
        condició garanteix homconjunts, no que tots
        els objectes formin un conjunt.
  \item[Apèndix D: solucions dels tests.]
        Respostes i justificacions. Consulteu-les només
        després d'haver respost tot el test.
\ifpublicacioparcial\else
  \item[Apèndix E: quasiexemples i contraexemples.]
        Recull estructures que fallen exactament un axioma
        (categoria, functor, transformació natural), la
        jerarquia entre igualtat, isomorfisme, isomorfisme
        natural i equivalència, i models de la lògica que
        queden fora dels topos. La secció sobre la
        jerarquia és útil des del capítol~\ref{cap:transformacions};
        la de la lògica, després del capítol~\ref{cap:topos}.
\fi
\end{description}

La taula de notació serveix per recuperar el significat
i la primera aparició dels símbols. El glossari
català--anglès ajuda a llegir la bibliografia internacional,
sobretot en termes com \emph{pullback}, \emph{slice},
\emph{presheaf} i \emph{whiskering}. L'índex alfabètic
és una eina per retrobar conceptes quan reapareixen en
un context nou.

Les lectures recomanades al final dels capítols tenen
una funció de contrast i ampliació. En una primera
lectura, és millor consultar una font alternativa per
aclarir una pregunta concreta que obrir alhora molts
tractaments generals del mateix tema.

\section{Una manera de mesurar l'aprenentatge}

No mesureu el progrés només pel nombre de pàgines llegides.
Una fita més fiable és poder fer, sense el text davant,
quatre operacions amb una idea nova:
\begin{enumerate}
  \item enunciar-ne la definició amb les hipòtesis correctes;
  \item interpretar-la en almenys dos exemples;
  \item utilitzar-la en un argument o un càlcul senzill;
  \item explicar què falla en un contraexemple o quan
        es retira una hipòtesi.
\end{enumerate}

No cal satisfer aquestes quatre condicions amb la mateixa
fluïdesa des del primer dia. És raonable avançar amb
una dificultat ben identificada i tornar-hi després.
El que convé evitar és substituir la comprensió per
la familiaritat visual amb els símbols.

En acabar el llibre, el guany més important serà poder
relacionar les peces: veure una propietat universal
darrere d'una construcció, una naturalitat darrere
d'una compatibilitat, una adjunció darrere d'una
correspondència i una operació sobre subobjectes
darrere d'una lectura lògica. Aquesta capacitat de
passar entre exemples i estructura és el criteri
que dona unitat a tot el recorregut.

\endgroup


% ============================================================
% PART I — TEORIA
% ============================================================
\part{Teoria}
\setcounter{chapter}{0}
\chapter{Categories}
\label{cap:categories}

\begin{objectius}
\apartat{Prerequisits} Llenguatge elemental de conjunts i aplicacions: domini, codomini i composició; producte cartesià, subconjunts i relacions. Lectura de quantificadors elementals i distinció entre existència i unicitat. Nocions bàsiques d'implicació, equivalència i contraexemple. Els exemples amb grups, espais vectorials o espais topològics enriqueixen la lectura, però no són imprescindibles. No es pressuposa cap coneixement previ de teoria de categories.
\apartat{Objectius} En acabar el capítol, el lector hauria de poder:
\begin{itemize}
  \item enunciar la definició de categoria i verificar-ne els axiomes en exemples concrets;
  \item llegir i escriure diagrames commutatius i traduir-los a igualtats de composicions;
  \item reconèixer la mateixa estructura categòrica en conjunts, preordres, monoides, grafs, relacions i sistemes deductius;
  \item definir i utilitzar isomorfismes, monomorfismes, epimorfismes, seccions i retraccions, distingint-los de les nocions conjuntistes corresponents;
  \item construir i interpretar la categoria oposada, categories producte, subcategories i categories sobre i sota un objecte en casos senzills;
  \item aplicar el principi de dualitat a definicions i resultats elementals;
  \item descriure la categoria lliure generada per un graf finit senzill;
  \item distingir categories petites i localment petites.
\end{itemize}
\end{objectius}

\section{Cap al concepte de categoria}

Quan comencem a estudiar matemàtiques, aprenem a treballar amb molts tipus d'\emph{estructures}: ordres, grups, espais vectorials, espais topològics, etc. Aviat notem que aquestes estructures rarament apareixen soles. Al costat de cadascuna hi ha sempre una classe de transformacions que les relacionen respectant la seva manera de comportar-se: entre ordres, les aplicacions monòtones; entre grups, els homomorfismes; entre espais vectorials, les aplicacions lineals; entre espais topològics, les aplicacions contínues. Aquestes transformacions admissibles les anomenarem, en general, \emph{morfismes} o \emph{fletxes}.

La noció de categoria no decideix per si sola quins morfismes són admissibles: ho fixem nosaltres en cada context, d'acord amb l'estructura que volem respectar. El que fa la teoria de categories és abstreure el patró comú a tots aquests contextos: treballar amb objectes, amb transformacions admissibles entre ells i amb la manera com aquestes transformacions es componen. Cada categoria concreta enregistra aquesta tria de morfismes i, a partir d'aquí, només en reté la composició i les identitats.

Els exemples concrets que acabem d'esmentar tenen tots un conjunt subjacent. En aquests casos, preservar l'estructura significa començar per una aplicació entre els conjunts subjacents i imposar-hi les condicions adequades. No és, però, un requisit de la noció de categoria: més endavant trobarem categories els objectes de les quals no són conjunts amb estructura. Com hem dit a la secció~\ref{sec:mirar-fora} de la introducció, la teoria de categories desplaça la visió clàssica, puntual i conjuntista, de les aplicacions cap a una visió abstracta com a fletxes.

Recordem breument la notació de fletxa per a les aplicacions entre conjunts. Una aplicació $f$ amb domini $X$ i codomini $Y$ es denota $f\colon X\to Y$, o bé com a fletxa $X\xrightarrow{f}Y$. L'operació fonamental sobre les aplicacions és la \emph{composició}: si $f\colon X\to Y$ i $g\colon Y\to Z$, es defineix $g\comp f\colon X\to Z$ per
\[
(g\comp f)(x)\coloneqq g(f(x)).
\]
Perquè la composició estigui definida, el codomini de $f$ ha de coincidir amb el domini de $g$; amb fletxes, $X\xrightarrow{f}Y\xrightarrow{g}Z$. A més, per a cada conjunt $X$ hi ha l'aplicació \emph{identitat} sobre $X$,
\[
\id_X\colon X\to X,\qquad \id_X(x)\coloneqq x.
\]
Aquestes operacions es regeixen per la llei d'associativitat i les lleis de la unitat: per a $f\colon X\to Y$, $g\colon Y\to Z$ i $h\colon Z\to W$,
\[
(h\comp g)\comp f=h\comp(g\comp f),\qquad f\comp\id_X=f=\id_Y\comp f.
\]
Observem que aquestes equacions es formulen purament en termes de les operacions sobre aplicacions, sense cap referència als elements dels conjunts $X$, $Y$, $Z$, $W$.

Això no és cap ingenuïtat, sinó una evidència sobre les estructures que hem mencionat: els homomorfismes entre grups, les aplicacions monòtones entre preordres, les aplicacions lineals entre espais vectorials i les aplicacions contínues entre espais topològics satisfan les condicions que acabem d'escriure.

Aquesta és la idea que orientarà tot el que segueix: estudiar un objecte no pel que és per dins, sinó pel paper que fa dins del sistema de fletxes que el connecten amb els altres. És el que ens permetrà de parlar-ne amb un mateix llenguatge i, fins i tot, estendre'l a altres tipus d'objectes no necessàriament basats en la noció d'aplicació, com ara grafs, sistemes deductius o relacions entre conjunts.

\section{Definició de categoria}

\begin{definicio}
\label{def:categoria}
\index{categoria}Una \textbf{categoria} $\cat{C}$ consta de les dades següents:
\begin{enumerate}
\item una col·lecció d'\textbf{objectes}\index{objecte} $A,B,C,\dots$, que denotem $\Ob(\cat{C})$;
\item una col·lecció de \textbf{morfismes}\index{morfisme} $f,g,h,\dots$, que denotem $\Mor(\cat{C})$; cada morfisme $f$ té associats un únic \textbf{domini}\index{domini} $A$ i un únic \textbf{codomini}\index{codomini} $B$, i ho escrivim com $A=\dom(f)$ i $B=\cod(f)$; en aquest cas també escrivim $f\colon A\to B$ o com a fletxa $A\xrightarrow{f}B$;
\item per a cada objecte $A$, un morfisme \textbf{identitat}\index{identitat} $\id_A\colon A\to A$;
\item per a cada parella de morfismes componibles $f\colon A\to B$ i $g\colon B\to C$, la \textbf{composició}\index{composicio@composició} $g\comp f\colon A\to C$.
\end{enumerate}
Aquestes dades han de satisfer els dos axiomes següents, sempre que les composicions indicades tinguin sentit:
\begin{enumerate}
\item \textbf{Associativitat.} $h\comp(g\comp f)=(h\comp g)\comp f$, per a $f\colon A\to B$, $g\colon B\to C$, $h\colon C\to D$.
\item \textbf{Identitats.} $\id_B\comp f=f$ i $f\comp\id_A=f$, per a tot $f\colon A\to B$.
\end{enumerate}
La secció~\ref{sec:diagrames} dona la lectura diagramàtica d'aquests axiomes.
\end{definicio}

Parlem aquí de \emph{col·leccions} i no de \emph{conjunts}, deliberadament: en alguns exemples els objectes, o els morfismes, són massa nombrosos per formar un conjunt. Així, tant $\Ob(\cat{C})$ com $\Mor(\cat{C})$ són en general col·leccions, no necessàriament conjunts. Precisarem aquesta qüestió a la secció~\ref{sec:mida}. Sovint és més còmode agrupar els morfismes segons el seu domini i codomini: quan la col·lecció de morfismes de $A$ a $B$ és un conjunt, se sol denotar
\[
\Hom_{\cat{C}}(A,B)
\]
o simplement $\Hom(A,B)$ quan la categoria és clara i l'anomenem \emph{homconjunt} de $A$ i $B$. Com que cada morfisme té un únic domini i un únic codomini, $\Mor(\cat{C})$ sempre es reparteix segons el parell $(A,B)$ de domini i codomini, en parts disjuntes dos a dos;\footnote{Aquesta disjunció és, de fet, equivalent a la unicitat del domini i el codomini. A l'axiomàtica relacional de l'apèndix~\ref{cap:axiomatica-relacional} en fa el paper d'un axioma.} el que no està garantit és que cada part sigui un \emph{conjunt}. Quan ho és per a tots els parells ---propietat que rep el nom de \emph{petitesa local} (secció~\ref{sec:mida})---, aquestes parts són els homconjunts $\Hom(A,B)$ i $\Mor(\cat{C})$ n'és la unió. Al llarg del llibre totes les categories que fem servir són localment petites, de manera que aquesta subtilesa no ens afectarà.

\begin{observacio}
Fem algunes precisions sobre la definició.
\begin{enumerate}
\item Els morfismes també s'anomenen \textbf{fletxes}, i sovint farem servir les dues paraules indistintament.
\item L'axioma d'associativitat fa que puguem escriure $h\comp g\comp f$ sense parèntesis, sabent que el resultat no depèn de la manera d'agrupar-los.
\item Els axiomes d'identitat diuen que $\id_A$ es comporta com un neutre per la composició. Veurem tot seguit que aquesta propietat determina $\id_A$ de manera única.
\item Aquesta definició axiomàtica pren com a termes no definits els objectes i els morfismes. Una categoria és qualsevol estructura que contingui aquestes dues menes de dades i que satisfaci les condicions i propietats de la definició.\footnotemark
\end{enumerate}
\end{observacio}
\footnotetext{De fet, es pot prescindir d'una de les dues menes de dades: com que cada objecte queda determinat per la seva identitat, una categoria petita es pot descriure en un llenguatge amb una sola mena d'individus ---les fletxes--- i tres relacions primitives de domini, codomini i composició, amb els objectes reconstruïts com les fletxes identitat. L'apèndix~\ref{cap:axiomatica-relacional} desenvolupa aquesta presentació relacional i en demostra l'equivalència amb la definició d'aquí.}

\begin{proposicio}
La fletxa identitat d'un objecte és única.
\end{proposicio}

\begin{prova}
Suposem que $e_A\colon A\to A$ i $e'_A\colon A\to A$ són dues fletxes identitat de $A$. Com que $e_A$ és una identitat,
\[
e_A\comp e'_A=e'_A.
\]
D'altra banda, com que $e'_A$ és una identitat,
\[
e_A\comp e'_A=e_A.
\]
Per tant $e_A=e'_A$.
\end{prova}

\section{Fletxes, camins i diagrames commutatius}\label{sec:diagrames}\index{diagrama commutatiu}

La definició anterior es pot llegir algebraicament, mitjançant equacions, però també gràficament. Aquesta segona lectura serà una part essencial del llenguatge categòric.

Si $f\colon A\to B$ i $g\colon B\to C$ són componibles, la composició es pot representar amb el triangle
\[
\begin{tikzpicture}[baseline=(m.center)]
  \matrix (m) [matrix of math nodes, column sep=3.4em, row sep=2.7em] {
    A & B \\
      & C \\
  };
  \draw[-{Stealth[scale=1.1]}] (m-1-1) -- node[above] {$f$} (m-1-2);
  \draw[-{Stealth[scale=1.1]}] (m-1-2) -- node[right] {$g$} (m-2-2);
  \draw[-{Stealth[scale=1.1]}] (m-1-1) -- node[below left] {$g\comp f$} (m-2-2);
\end{tikzpicture}
\]
La fletxa directa de $A$ a $C$ és precisament el resultat de seguir el camí $A\to B\to C$.

Més generalment, un diagrama pot contenir diferents camins dirigits amb el mateix origen i el mateix destí. Per exemple,
\[
\begin{tikzpicture}[baseline=(m.center)]
  \matrix (m) [matrix of math nodes, row sep=2.7em, column sep=3.5em] {
    A & B \\
    C & D \\
  };
  \draw[-{Stealth[scale=1.1]}] (m-1-1) -- node[above] {$f$} (m-1-2);
  \draw[-{Stealth[scale=1.1]}] (m-1-1) -- node[left] {$u$} (m-2-1);
  \draw[-{Stealth[scale=1.1]}] (m-1-2) -- node[right] {$v$} (m-2-2);
  \draw[-{Stealth[scale=1.1]}] (m-2-1) -- node[below] {$g$} (m-2-2);
\end{tikzpicture}
\]
té dos camins de $A$ a $D$: primer $f$ i després $v$, o primer $u$ i després $g$.

\begin{definicio}
Un diagrama com l'anterior \textbf{commuta} si els dos camins determinen el mateix morfisme, és a dir, si
\[
v\comp f=g\comp u.
\]
En general, direm que un diagrama commuta quan tots els camins dirigits amb el mateix origen i el mateix destí que cal comparar donen la mateixa composició.
\end{definicio}

\begin{observacio}
Un dibuix de fletxes no és automàticament un diagrama commutatiu. El dibuix especifica objectes i morfismes; afirmar que \emph{commuta} afegeix una o més igualtats entre composicions de fletxes. Per això convé aprendre a passar en tots dos sentits:
\[
\text{diagrama}\quad\longleftrightarrow\quad\text{equacions entre fletxes i composicions}.
\]
\end{observacio}

L'associativitat es pot llegir com un diagrama commutatiu. En una cadena
\[
A\xrightarrow{f}B\xrightarrow{g}C\xrightarrow{h}D
\]
podem agrupar les composicions de dues maneres, $(h\comp g)\comp f$ o $h\comp(g\comp f)$. El diagrama següent ho mostra com un paral·lelogram de vèrtexs $A,B,D,C$ i costats $f$, $h\comp g$, $h$ i $g\comp f$; la diagonal interior és $g$.
\[
\begin{tikzpicture}[>={Stealth[scale=1.1]}]
  \node (A) at (-3,2) {$A$};
  \node (B) at (0,2)  {$B$};
  \node (C) at (0,0)  {$C$};
  \node (D) at (3,0)  {$D$};
  \draw[->] (A) -- node[above] {$f$} (B);
  \draw[->] (B) -- node[right] {$h\comp g$} (D);
  \draw[->] (C) -- node[below] {$h$} (D);
  \draw[->] (A) -- node[below left=-2pt] {$g\comp f$} (C);
  \draw[->] (B) -- node[left] {$g$} (C);
\end{tikzpicture}
\]
Recórrer $A\to B\to C\to D$ agrupant primer $g\comp f$ i després component amb $h$, o agrupant primer $h\comp g$ i precomponent amb $f$, dona la mateixa fletxa $A\to D$. Això és exactament l'axioma d'associativitat:
\[
(h\comp g)\comp f=h\comp(g\comp f).
\]

Les identitats també tenen una lectura diagramàtica. Si $f\colon A\to B$, inserir $\id_A$ abans de $f$ o $\id_B$ després de $f$ no canvia el camí:
\[
\begin{tikzpicture}[baseline=(m.center)]
  \matrix (m) [matrix of math nodes, row sep=2.6em, column sep=2.6em] {
    A & A \\
      & B \\
  };
  \draw[-{Stealth[scale=1.1]}] (m-1-1) -- node[above] {\(\id_A\)} (m-1-2);
  \draw[-{Stealth[scale=1.1]}] (m-1-2) -- node[right] {\(f\)} (m-2-2);
  \draw[-{Stealth[scale=1.1]}] (m-1-1) -- node[below left] {\(f\)} (m-2-2);
\end{tikzpicture}
\qquad\qquad
\begin{tikzpicture}[baseline=(m.center)]
  \matrix (m) [matrix of math nodes, row sep=2.6em, column sep=2.6em] {
    A & B \\
      & B \\
  };
  \draw[-{Stealth[scale=1.1]}] (m-1-1) -- node[above] {\(f\)} (m-1-2);
  \draw[-{Stealth[scale=1.1]}] (m-1-2) -- node[right] {\(\id_B\)} (m-2-2);
  \draw[-{Stealth[scale=1.1]}] (m-1-1) -- node[below left] {\(f\)} (m-2-2);
\end{tikzpicture}
\]
Els dos triangles commutatius expressen exactament els axiomes d'identitat:
\[
f\comp\id_A=f,\qquad \id_B\comp f=f.
\]
A partir d'ara seguirem aquesta norma: quan una propietat categòrica tingui una expressió diagramàtica natural, en donarem alhora, al mateix lloc, l'equació i el diagrama. Llegir l'una i dibuixar l'altre, i a l'inrevés, és una destresa que el llibre demana constantment.

\section{Exemples de categories}

La força de la definició es veu de seguida en la varietat d'exemples que hi encaixen. En cadascun cal identificar els objectes, els morfismes, les identitats i la composició. Comencem pels exemples més petits, que malgrat la seva simplicitat apareixeran sovint com a \emph{categories índex} (la definició~\refalt{def:diagrama}{de diagrama del capítol~5}) per construir diagrames.

\subsection{Categories finites}\label{subsec:finites}\index{categoria!buida}\index{categoria!unitat}\index{categoria!finita}

Per descomptat, els objectes d'una categoria no cal que siguin conjunts. Considerem alguns exemples molt senzills que es poden descriure per complet. En cada cas identifiquem els objectes, els morfismes, la composició i les identitats; l'associativitat i les lleis d'identitat es comproven directament, ja que hi ha molt poques composicions a considerar.
\begin{enumerate}
\item La \emph{categoria buida} $\cat{0}$ té $\Ob(\cat{0})=\varnothing$ i cap morfisme. No hi ha composició ni identitats que definir, i els axiomes es compleixen trivialment perquè no hi ha res a comprovar; és la categoria \enquote{sense res}.
\item La \emph{categoria unitat} $\cat{1}$, també anomenada categoria \emph{terminal}, té $\Ob(\cat{1})=\{\ast\}$ i $\Hom(\ast,\ast)=\{\id_\ast\}$. L'única composició és $\id_\ast\comp\id_\ast=\id_\ast$, i els axiomes es redueixen a aquesta igualtat:
\[
\begin{tikzpicture}[baseline=(a.base)]
  \node (a) {$\ast$};
  \draw[-{Stealth[scale=1.1]}] (a) to[out=40,in=-40,looseness=8] node[right] {$\id_\ast$} (a);
\end{tikzpicture}
\]
\item La categoria $\cat{2}$ té $\Ob(\cat{2})=\{\ast,\star\}$ i una única fletxa no trivial $u\colon\ast\to\star$, a més de les dues identitats:
\[
\begin{aligned}
\Hom(\ast,\ast)&=\{\id_\ast\}, &
\Hom(\star,\star)&=\{\id_\star\},\\
\Hom(\ast,\star)&=\{u\}, &
\Hom(\star,\ast)&=\varnothing.
\end{aligned}
\]
Les úniques composicions són les que involucren identitats, $u\comp\id_\ast=u=\id_\star\comp u$, de manera que els axiomes es compleixen automàticament. A partir d'ara, com és costum, no dibuixem les identitats:
\[
\begin{tikzpicture}[baseline=(a.base)]
  \node (a) at (0,0) {$\ast$};
  \node (b) at (1.8,0) {$\star$};
  \draw[-{Stealth[scale=1.1]}] (a) -- node[above] {$u$} (b);
\end{tikzpicture}
\]
\item La categoria amb \textbf{dues fletxes paral·leles} té $\Ob=\{\ast,\star\}$ i dues fletxes no trivials $f,g\colon\ast\to\star$ amb el mateix domini i el mateix codomini:
\[
\Hom(\ast,\star)=\{f,g\},\qquad
\Hom(\star,\ast)=\varnothing,
\]
i als homconjunts $\Hom(\ast,\ast)$ i $\Hom(\star,\star)$ només hi ha les identitats. Com que cap parell d'aquestes fletxes no és componible, de nou les úniques composicions són amb identitats:
\[
\begin{tikzpicture}[baseline=(a.base)]
  \node (a) at (0,0) {$\ast$};
  \node (b) at (1.8,0) {$\star$};
  \draw[-{Stealth[scale=1.1]}] ([yshift=2.2pt]a.east) -- node[above] {$f$} ([yshift=2.2pt]b.west);
  \draw[-{Stealth[scale=1.1]}] ([yshift=-2.2pt]a.east) -- node[below] {$g$} ([yshift=-2.2pt]b.west);
\end{tikzpicture}
\]
Aquesta darrera reapareixerà com a categoria índex dels \emph{equalitzadors} (la definició~\refalt{def:equalitzador}{d'equalitzador del capítol~5}).
\item La categoria $\cat{3}$ té $\Ob(\cat{3})=\{\ast,\star,\bullet\}$ i tres fletxes no trivials, $f\colon\ast\to\star$, $g\colon\star\to\bullet$ i $h\colon\ast\to\bullet$:
\[
\Hom(\ast,\star)=\{f\},\qquad
\Hom(\star,\bullet)=\{g\},\qquad
\Hom(\ast,\bullet)=\{h\};
\]
els homconjunts d'un objecte a si mateix contenen només la identitat, i els homconjunts en sentit contrari són buits. La composició obliga a incloure la composició de les dues fletxes consecutives: l'única composició no trivial és $g\comp f=h$. Amb aquesta composició (i les identitats) els axiomes se satisfan:
\[
\begin{tikzpicture}[baseline=(m.center)]
  \matrix (m) [matrix of math nodes, row sep=2.6em, column sep=2.6em] {
    \ast & \star \\
         & \bullet \\
  };
  \draw[-{Stealth[scale=1.1]}] (m-1-1) -- node[above] {$f$} (m-1-2);
  \draw[-{Stealth[scale=1.1]}] (m-1-2) -- node[right] {$g$} (m-2-2);
  \draw[-{Stealth[scale=1.1]}] (m-1-1) -- node[below left] {$h=g\comp f$} (m-2-2);
\end{tikzpicture}
\]
\end{enumerate}

\begin{observacio}
Afegir fletxes no és gratuït: la composició obliga a incloure totes les composicions. Dues fletxes de sentits oposats $f\colon\ast\to\star$ i $g\colon\star\to\ast$
\[
\begin{tikzpicture}[baseline=(a.base)]
  \node (a) at (0,0) {$\ast$};
  \node (b) at (2.2,0) {$\star$};
  \draw[-{Stealth[scale=1.1]}] ([yshift=2.6pt]a.east) -- node[above] {$f$} ([yshift=2.6pt]b.west);
  \draw[-{Stealth[scale=1.1]}] ([yshift=-2.6pt]b.west) -- node[below] {$g$} ([yshift=-2.6pt]a.east);
\end{tikzpicture}
\]
generen $g\comp f$, $f\comp g$, $g\comp f\comp g$, i així indefinidament; la categoria resultant és infinita. Per obtenir-ne una de finita cal imposar equacions entre composicions ---per exemple $g\comp f=\id_\ast$ i $f\comp g=\id_\star$, cosa que fa de $f$ un isomorfisme (la definició~\ref{def:isomorfisme}). En canvi, dues fletxes \emph{paral·leles} $f,g\colon\ast\to\star$ no són componibles entre elles, i la categoria és finita sense afegir res.
\end{observacio}

\subsection{Categories de conjunts}\label{subsec:conjunts}

La categoria $\cat{Set}$ té:
\begin{itemize}
\item objectes: $\Ob(\cat{Set})$ és la col·lecció de tots els conjunts;
\item morfismes: $\Hom_{\cat{Set}}(A,B)$ és el conjunt de totes les aplicacions $A\to B$;
\item composició: la composició habitual d'aplicacions, $(g\comp f)(x)=g(f(x))$;
\item identitats: l'aplicació identitat, $\id_A(x)=x$.
\end{itemize}
L'associativitat i les lleis d'identitat són propietats ben conegudes de les aplicacions, de manera que $\cat{Set}$ és una categoria. Serà una categoria de referència: quan un concepte categòric nou ens desconcerti, mirar què vol dir a $\cat{Set}$ sol ser un bon primer pas.

Si ens restringim als conjunts finits i a les aplicacions entre ells, obtenim la categoria $\cat{FinSet}$: $\Ob(\cat{FinSet})$ és la col·lecció dels conjunts finits, $\Hom_{\cat{FinSet}}(A,B)=\Hom_{\cat{Set}}(A,B)$, i la composició i les identitats són les mateixes que a $\cat{Set}$. Com que la composició d'aplicacions entre conjunts finits torna a ser una aplicació entre conjunts finits, els axiomes se satisfan igual que a $\cat{Set}$.

Tot conjunt $S$ es pot veure també com una categoria, que continuem escrivint $S$:\index{categoria!discreta} $\Ob(S)=S$ i els únics morfismes són les identitats, una per objecte,
\[
\Hom_S(x,y)=\begin{cases}\{\id_x\} & \text{si } x=y,\\[2pt] \varnothing & \text{altrament.}\end{cases}
\]
La composició és forçada ($\id_x\comp\id_x=\id_x$) i els axiomes es compleixen trivialment. Una categoria en què els únics morfismes són les identitats s'anomena \textbf{categoria discreta}. En una categoria discreta no hi ha cap informació més enllà de la col·lecció d'objectes.

\begin{observacio}
Una categoria no queda determinada pels seus objectes: cal dir també quines són les fletxes. Amb només dos objectes, per exemple, hem trobat ja tres categories ben diferents:
\begin{itemize}
\item la \textbf{categoria discreta} de dos punts, només amb les identitats i cap fletxa no trivial;
\item la categoria $\cat{2}$, amb una única fletxa no trivial d'un objecte a l'altre;
\item la categoria amb \textbf{dues fletxes paral·leles}, amb dues fletxes no trivials entre els mateixos objectes.
\end{itemize}
Totes tres tenen els mateixos dos objectes, però són categories diferents perquè difereixen en les fletxes. En particular, $\cat{2}$ \emph{no} és la categoria discreta de dos punts: la presència de la fletxa no trivial les distingeix.
\end{observacio}

\subsection{Estructures elementals com a categories}\label{subsec:elementals}

Les estructures matemàtiques més simples sovint són categories. Per exemple, un monoide no s'assembla a una categoria d'un sol objecte: és exactament una categoria d'un sol objecte. Recollim aquesta idea en una llista d'identificacions, i tot seguit les justifiquem:
\begin{itemize}
\item un conjunt és una categoria discreta (ja vist a la subsecció anterior);
\item un monoide és una categoria d'un sol objecte;
\item un grup és una categoria d'un sol objecte en què tot morfisme és un isomorfisme;
\item un preordre és una categoria prima;
\item un ordre parcial és una categoria prima en què els únics isomorfismes són les identitats;
\item un sistema deductiu és una categoria prima.
\end{itemize}
En tots els casos, els objectes i els morfismes deixen de ser conjunts i aplicacions: passen a ser els elements i les relacions internes de l'estructura.

\paragraph{Preordres.}\index{preordre}
Un \emph{preordre} és un parell $(P,\le)$ on $P$ és un conjunt i $\le$ una relació binària sobre $P$ reflexiva ($x\le x$) i transitiva (si $x\le y$ i $y\le z$, aleshores $x\le z$). Per exemple, $(\mathbb{R},\le)$, o els subconjunts d'un conjunt ordenats per la inclusió $\subseteq$.

Tot preordre $(P,\le)$ \textbf{és} una categoria, que continuem escrivint $P$:
\begin{itemize}
\item objectes: $\Ob(P)=P$, els elements del preordre;
\item morfismes: hi ha exactament una fletxa $x\to y$ quan $x\le y$,
\[
\Hom_P(x,y)=\begin{cases}\{(x,y)\} & \text{si } x\le y,\\[2pt] \varnothing & \text{altrament;}\end{cases}
\]
\item composició: $(y,z)\comp(x,y)=(x,z)$, que existeix per transitivitat;
\item identitats: $\id_x=(x,x)$, que existeix per reflexivitat.
\end{itemize}
\[
\begin{tikzpicture}[baseline=(m.center)]
  \matrix (m) [matrix of math nodes, row sep=2.6em, column sep=2.6em] {
    x & y \\
      & z \\
  };
  \draw[-{Stealth[scale=1.1]}] (m-1-1) -- node[above] {$(x,y)$} (m-1-2);
  \draw[-{Stealth[scale=1.1]}] (m-1-2) -- node[right] {$(y,z)$} (m-2-2);
  \draw[-{Stealth[scale=1.1]}] (m-1-1) -- node[below left] {$(x,z)$} (m-2-2);
\end{tikzpicture}
\]
Com que cada homconjunt té com a màxim un element, l'associativitat i les lleis d'identitat se satisfan automàticament. Una categoria en què $|\Hom(x,y)|\le 1$ per a tot parell d'objectes s'anomena \textbf{prima}\index{categoria!prima} (en anglès, \emph{thin category}).

Si el preordre és a més antisimètric ($x\le y$ i $y\le x$ impliquen $x=y$), és un \emph{ordre parcial}. Categòricament, l'antisimetria diu que els únics isomorfismes (la definició~\ref{def:isomorfisme}) són les identitats: si $x\to y$ i $y\to x$ existeixen totes dues, aleshores $x=y$.

\paragraph{Monoides.}\index{monoide}
Un \emph{monoide} és una terna $(M,\ast,e)$ on $\ast$ és una operació associativa sobre $M$ amb element neutre $e$. Per exemple, $(\mathbb{N},+,0)$, o les paraules sobre un alfabet amb la concatenació i la paraula buida.

Tot monoide \textbf{és} una categoria amb un sol objecte $\ast$, que escrivim $\cat{B}M$.\footnote{La notació $\cat{B}M$ prové de la topologia, on $BM$ designa l'\emph{espai classificador} del monoide (o grup) $M$; la categoria $\cat{B}M$ n'és el model categòric.} Concretament:
\begin{itemize}
\item objectes: $\Ob(\cat{B}M)=\{\ast\}$;
\item morfismes: l'únic homconjunt és tot el monoide, $\Hom_{\cat{B}M}(\ast,\ast)=M$;
\item composició: l'operació del monoide, $g\comp f=g\ast f$, on l'element de la dreta actua primer;
\item identitat: $\id_\ast=e$.
\end{itemize}
Els axiomes de categoria són exactament els de monoide. Els morfismes, aquí, no són aplicacions: són els elements de $M$.

\paragraph{Grups.}\index{grup}
Un \emph{grup} és una terna $(G,\ast,e)$ que és un monoide en què, a més, tot element $a$ té un invers $a^{-1}$ amb $a\ast a^{-1}=e=a^{-1}\ast a$. Per exemple, $(\mathbb{Z},+,0)$, o les permutacions d'un conjunt amb la composició.

Un grup \textbf{és} una categoria d'un sol objecte en què tot morfisme és un isomorfisme. Com que un grup és en particular un monoide, la categoria és la de l'apartat anterior, que escrivim igualment $\cat{B}G$: $\Ob(\cat{B}G)=\{\ast\}$, $\Hom_{\cat{B}G}(\ast,\ast)=G$, $g\comp f=g\ast f$ i $\id_\ast=e$. En efecte, l'invers de l'element $a$ és exactament el morfisme invers de la fletxa corresponent. La teoria de grups és, doncs, la part invertible de la teoria de monoides. Cada element del grup és un bucle sobre l'únic objecte, i la identitat $e$ n'és el bucle destacat:
\[
\begin{tikzpicture}[baseline=(o.base)]
  \node (o) {$\ast$};
  % pètals genèrics (elements del grup)
  \foreach \ang in {30,150,210,270,330}{
    \draw[-{Stealth[scale=0.9]},gray] (o) to[out=\ang+14,in=\ang-14,looseness=16] (o);
  }
  % pètal destacat: la identitat e
  \draw[-{Stealth[scale=1.0]},thick] (o) to[out=104,in=76,looseness=16] (o);
  \node at (0,1.32) {$e$};
\end{tikzpicture}
\]

\paragraph{Sistemes deductius.}\index{deduibilitat@deduïbilitat}\index{Prov@$\cat{Prov}_T$}
Fixem un \emph{sistema deductiu} $T$ i considerem la seva relació de deduïbilitat $\vdash_T$ entre fórmules, reflexiva ($p\vdash_T p$) i transitiva (de $p\vdash_T q$ i $q\vdash_T r$ se segueix $p\vdash_T r$). Aquesta relació determina una categoria, que escrivim $\cat{Prov}_T$ per distingir-la del sistema $T$ del qual prové:
\begin{itemize}
\item objectes: $\Ob(\cat{Prov}_T)$ és el conjunt de les fórmules $p,q,r,\dots$ de $T$;
\item morfismes: hi ha exactament una fletxa $p\to q$ quan $p\vdash_T q$,
\[
\Hom_{\cat{Prov}_T}(p,q)=\begin{cases}\{(p,q)\} & \text{si } p\vdash_T q,\\[2pt] \varnothing & \text{altrament;}\end{cases}
\]
\item composició: $(q,r)\comp(p,q)=(p,r)$, que existeix per transitivitat;
\item identitats: $\id_p=(p,p)$, que existeix per reflexivitat.
\end{itemize}
\[
\begin{tikzpicture}[baseline=(m.center)]
  \matrix (m) [matrix of math nodes, row sep=2.6em, column sep=2.6em] {
    p & q \\
      & r \\
  };
  \draw[-{Stealth[scale=1.1]}] (m-1-1) -- node[above] {$p\vdash_T q$} (m-1-2);
  \draw[-{Stealth[scale=1.1]}] (m-1-2) -- node[right] {$q\vdash_T r$} (m-2-2);
  \draw[-{Stealth[scale=1.1]}] (m-1-1) -- node[below left] {$p\vdash_T r$} (m-2-2);
\end{tikzpicture}
\]
Com el preordre, $\cat{Prov}_T$ és una categoria prima. A diferència de l'ordre parcial, però, la deduïbilitat no acostuma a ser antisimètrica: dues fórmules diferents poden ser interdeduïbles ($p\vdash_T q$ i $q\vdash_T p$ amb $p\neq q$), com $p\land q$ i $q\land p$. Categòricament, això vol dir que hi ha isomorfismes no trivials: fórmules distintes però isomorfes. Aquesta categoria només registra \emph{si} $q$ es dedueix de $p$, no \emph{de quina manera} es dedueix; l'observació següent presenta la versió que distingeix les proves.

\begin{observacio}[Les proves mateixes com a fletxes]\index{prova!com a morfisme}
Hi ha una interpretació lògica més fina en què una prova concreta
\[
\pi\colon p\longrightarrow q
\]
és ella mateixa un morfisme, i dues proves $\pi,\rho\colon p\to q$ poden ser morfismes diferents. La composició representa aleshores l'encadenament de proves i la identitat, la prova trivial d'una fórmula a partir d'ella mateixa. Passem, doncs, d'una sola fletxa (registrar \emph{si} hi ha deducció) a moltes fletxes (registrar \emph{quina} prova).

Per obtenir una categoria de proves completament definida cal fixar el sistema deductiu i precisar també quan dues derivacions es consideren iguals ---per exemple, mitjançant una equivalència adequada entre proves. No necessitem encara aquesta teoria; de moment distingirem entre la categoria prima $\cat{Prov}_T$ de la \emph{deduïbilitat} i aquesta interpretació més rica de les \emph{proves com a fletxes}. Hi tornarem quan el llenguatge categòric ens permeti fer-ho amb més precisió.
\end{observacio}

\subsection{Categories de conjunts estructurats}

Ara considerem les categories els objectes de les quals són conjunts amb alguna estructura, i els morfismes, les aplicacions que preserven aquesta estructura. En tots els casos la composició és la d'aplicacions i les identitats són les aplicacions identitat, de manera que l'associativitat i les lleis d'identitat s'hereten de les de $\cat{Set}$; només cal comprovar, en cada cas, que la composició d'aplicacions que respecten l'estructura torna a respectar-la, i que la identitat també ho fa.

\paragraph{Grups.}
La categoria $\cat{Grp}$ té:
\begin{itemize}
\item objectes: $\Ob(\cat{Grp})$ és la col·lecció de tots els grups;
\item morfismes: $\Hom_{\cat{Grp}}(G,H)$ és el conjunt dels homomorfismes de grups $G\to H$;
\item composició: la composició d'aplicacions;
\item identitats: l'aplicació identitat $\id_{G}$.
\end{itemize}
Un \emph{homomorfisme de grups} $h\colon G\to H$ és una aplicació entre els conjunts subjacents que preserva l'operació: $h(a\ast b)=h(a)\ast h(b)$ per a tot $a,b$ (d'on se segueix que preserva també el neutre i els inversos). Si $h\colon G\to H$ i $k\colon H\to K$ són homomorfismes, la composició també ho és, ja que
\[
k(h(a\ast b))=k(h(a)\ast h(b))=k(h(a))\ast k(h(b));
\]
i l'aplicació identitat preserva trivialment l'operació. Com que la composició i les identitats són les d'aplicacions entre conjunts, l'associativitat i les lleis d'identitat s'hereten de $\cat{Set}$.

\paragraph{Espais vectorials.}\index{espai vectorial}
Fixem un cos $\mathbb{K}$. La categoria $\cat{Vect}_{\mathbb K}$ té:
\begin{itemize}
\item objectes: $\Ob(\cat{Vect}_{\mathbb K})$ és la col·lecció dels espais vectorials sobre $\mathbb{K}$;
\item morfismes: $\Hom_{\mathbb K}(V,W)$ és el conjunt de les aplicacions lineals $V\to W$;
\item composició: la composició d'aplicacions;
\item identitats: l'aplicació identitat $\id_V$.
\end{itemize}
Un \emph{espai vectorial} sobre $\mathbb{K}$ és un conjunt $V$ amb una suma commutativa i associativa $V\times V\to V$ ---amb element neutre i oposats--- i una multiplicació per escalars $\mathbb{K}\times V\to V$ compatible amb les operacions. Una \emph{aplicació lineal} $f\colon V\to W$ és una aplicació entre els conjunts subjacents que preserva l'estructura: $f(u+v)=f(u)+f(v)$ i $f(\lambda v)=\lambda f(v)$. Si $f\colon V\to W$ i $g\colon W\to U$ són lineals, la composició també ho és, ja que
\[
g(f(u+v))=g(f(u))+g(f(v))\qquad\text{i}\qquad g(f(\lambda v))=\lambda\,g(f(v));
\]
i la identitat és lineal. Per tant la composició i les identitats són les d'aplicacions entre conjunts, i l'associativitat i les lleis d'identitat s'hereten de $\cat{Set}$.

\paragraph{Espais topològics.}\index{espai topològic}
La categoria $\cat{Top}$ té:
\begin{itemize}
\item objectes: $\Ob(\cat{Top})$ és la col·lecció de tots els espais topològics;
\item morfismes: $\Hom_{\cat{Top}}(X,Y)$ és el conjunt de les aplicacions contínues $X\to Y$;
\item composició: la composició d'aplicacions;
\item identitats: l'aplicació identitat $\id_{X}$.
\end{itemize}
Un \emph{espai topològic} és un parell $(X,\mathcal{T}_X)$ on $X$ és un conjunt i $\mathcal{T}_X$ una família de subconjunts de $X$ (els \emph{oberts}) tal que:
\begin{itemize}
\item $\varnothing, X\in\mathcal{T}_X$;
\item si $U,V\in\mathcal{T}_X$, aleshores $U\cap V\in\mathcal{T}_X$;
\item si $\{U_i\}_{i\in I}$ és una família qualsevol en $\mathcal{T}_X$, aleshores $\bigcup_{i\in I}U_i\in\mathcal{T}_X$.
\end{itemize}
Una \emph{aplicació contínua} $f\colon(X,\mathcal{T}_X)\to(Y,\mathcal{T}_Y)$ és una aplicació $f\colon X\to Y$ que preserva l'estructura en el sentit següent: $f^{-1}(U)\in\mathcal{T}_X$ per a tot $U\in\mathcal{T}_Y$. Si $f\colon X\to Y$ i $g\colon Y\to Z$ són contínues, la composició també ho és, ja que $(g\comp f)^{-1}(U)=f^{-1}(g^{-1}(U))$ és obert quan $U$ ho és; i la identitat és contínua, perquè $\id^{-1}(U)=U$. Per tant, un cop més, la composició i les identitats són les d'aplicacions entre conjunts, i l'associativitat i les lleis d'identitat s'hereten de $\cat{Set}$.

\paragraph{Monoides.}
La categoria $\cat{Mon}$ té:
\begin{itemize}
\item objectes: $\Ob(\cat{Mon})$ és la col·lecció de tots els monoides;
\item morfismes: $\Hom_{\cat{Mon}}(M,N)$ és el conjunt dels homomorfismes de monoides $M\to N$;
\item composició: la composició d'aplicacions;
\item identitats: l'aplicació identitat $\id_{M}$.
\end{itemize}
Un \emph{homomorfisme de monoides} $h\colon M\to N$ és una aplicació entre els conjunts subjacents que preserva l'operació i el neutre: $h(a\ast b)=h(a)\ast h(b)$ i $h(e_M)=e_N$. Si $h\colon M\to N$ i $k\colon N\to P$ són homomorfismes, la composició també ho és, ja que
\[
k(h(a\ast b))=k(h(a)\ast h(b))=k(h(a))\ast k(h(b))\qquad\text{i}\qquad k(h(e_M))=k(e_N)=e_P;
\]
i la identitat preserva trivialment l'operació i el neutre. Com que la composició i les identitats són les d'aplicacions entre conjunts, l'associativitat i les lleis d'identitat s'hereten de $\cat{Set}$.

\paragraph{Ordres parcials.}
La categoria $\cat{Pos}$ té:
\begin{itemize}
\item objectes: $\Ob(\cat{Pos})$ és la col·lecció de tots els ordres parcials;
\item morfismes: $\Hom_{\cat{Pos}}(P,Q)$ és el conjunt de les aplicacions monòtones $P\to Q$;
\item composició: la composició d'aplicacions;
\item identitats: l'aplicació identitat $\id_{P}$.
\end{itemize}
Una \emph{aplicació monòtona} $f\colon P\to Q$ és una aplicació entre els conjunts subjacents que preserva l'ordre: si $x\le y$, aleshores $f(x)\le f(y)$. Si $f\colon P\to Q$ i $g\colon Q\to R$ són monòtones, la composició també ho és, ja que
\[
x\le y \ \Longrightarrow\ f(x)\le f(y) \ \Longrightarrow\ g(f(x))\le g(f(y));
\]
i la identitat preserva trivialment l'ordre. Un cop més, la composició i les identitats són les d'aplicacions entre conjunts, i les lleis s'hereten de $\cat{Set}$.

Convé no confondre les dues mirades: un monoide concret \emph{és} una categoria (d'un sol objecte) i un ordre parcial concret \emph{és} una categoria (prima), mentre que $\cat{Mon}$ i $\cat{Pos}$ tenen \emph{tots} els monoides i \emph{tots} els ordres parcials per objectes, respectivament.

\begin{observacio}[Categories com a contextos matemàtics]
Cadascuna d'aquestes categories és l'\emph{àmbit} on es desenvolupa una teoria matemàtica: $\cat{Top}$ és el context de la topologia general, $\cat{Grp}$ el de la teoria de grups, $\cat{Vect}_{\mathbb K}$ el de l'àlgebra lineal, i així successivament. Aquí, la categoria la formen tots els objectes d'una mena; a les estructures elementals, en canvi, un sol objecte estructurat ja és una categoria.
\end{observacio}

\subsection{Conjunts i relacions}\index{Rel@$\cat{Rel}$}

Com hem avançat a la introducció (secció~\ref{sec:mirar-fora}), els conjunts admeten una segona categoria, on els morfismes no són aplicacions sinó \emph{relacions}. La categoria $\cat{Rel}$ té:
\begin{itemize}
\item objectes: $\Ob(\cat{Rel})$ és la col·lecció de tots els conjunts;
\item morfismes: $\Hom_{\cat{Rel}}(X,Y)$ és el conjunt de les relacions $R\subseteq X\times Y$ (amb la precisió sobre els extrems de l'observació~\ref{obs:extrems-implicits});
\item composició: per a $R\subseteq X\times Y$ i $S\subseteq Y\times Z$,
\[
S\comp R=\{(x,z)\mid \exists y\in Y.\ (x,y)\in R \ \text{i}\ (y,z)\in S\}\subseteq X\times Z;
\]
\item identitats: la relació diagonal $\id_X=\{(x,x)\mid x\in X\}$.
\end{itemize}
Els axiomes es comproven directament a partir de la definició. Per a l'associativitat, un parell $(x,w)$ pertany a les dues maneres d'agrupar tres relacions encadenades exactament quan hi ha termes intermedis que connecten $x$ amb $w$ passant per les tres relacions; com que la condició no depèn de l'agrupament, les dues composicions coincideixen. Per a les identitats, compondre amb la diagonal no afegeix cap pas efectiu, perquè obliga el terme intermedi a coincidir amb l'extrem. La verificació completa és el primer exercici resolt del capítol~1 de la part de Pràctica (exercici~\ref{exr:pr-rel}). Compondre relacions equival a aplicar un quantificador existencial sobre el terme intermedi: aquest és el primer pont amb la lògica que assenyalàvem a la introducció. A diferència del que passa a $\cat{Set}$, un morfisme $X\to Y$ de $\cat{Rel}$ no ha d'assignar necessàriament a cada element de $X$ un únic element de $Y$: pot relacionar-lo amb diversos elements de $Y$, o amb cap.

\begin{observacio}[Extrems implícits]\label{obs:extrems-implicits}
En rigor, un morfisme $X\to Y$ de $\cat{Rel}$ no és una relació $R$ tota sola, sinó la relació \emph{junt amb} els seus extrems: la terna $(X,Y,R)$ amb $R\subseteq X\times Y$. Aquesta precisió és necessària perquè els homconjunts siguin disjunts, tal com exigeix la definició de categoria. Per exemple, la relació buida $R=\varnothing$ està continguda en $X\times Y$ per a qualsevol parell $(X,Y)$; si no registréssim els extrems, un mateix objecte $\varnothing$ seria alhora morfisme entre infinites parelles d'objectes, amb domini i codomini indeterminats. Registrant $(X,Y,R)$, cada morfisme té un únic domini $X$ i un únic codomini $Y$. Escriurem simplement $R$ quan els extrems se sobreentenguin; és la mateixa convenció d'\emph{extrems implícits} que ja hem fet servir en codificar la fletxa d'un preordre o de $\cat{Prov}_T$ com el parell $(x,y)$, i que retrobarem tot seguit a les categories sobre i sota un objecte. No canvia cap de les construccions: només n'explicita la compatibilitat amb la definició inicial de categoria.
\end{observacio}

\subsection{Categories de grafs}\label{subsec:grafs}

\begin{definicio}\index{graf dirigit}
Un \textbf{graf dirigit} $G$ consta d'un conjunt $V_G$ de \textbf{vèrtexs}, un conjunt $E_G$ d'\textbf{arestes} i dues aplicacions
\[
s_G,t_G\colon E_G\longrightarrow V_G,
\]
que assignen a cada aresta el seu \textbf{origen} i el seu \textbf{destí}, respectivament.
\end{definicio}

\begin{definicio}\index{morfisme de grafs}
Un \textbf{morfisme de grafs} $F\colon G\to H$ és una parella d'aplicacions
\[
F_V\colon V_G\to V_H,
\qquad
F_E\colon E_G\to E_H,
\]
que preserven origen i destí:
\[
F_V\comp s_G=s_H\comp F_E,
\qquad
F_V\comp t_G=t_H\comp F_E.
\]
Equivalentment, han de commutar els dos quadrats següents:

\begin{center}
\begin{tikzpicture}[baseline=(m.center)]
  \matrix (m) [matrix of math nodes, row sep=2.5em, column sep=3.2em] {
    E_G & V_G \\
    E_H & V_H \\
  };
  \draw[-{Stealth[scale=1.0]}] (m-1-1) -- node[above] {$s_G$} (m-1-2);
  \draw[-{Stealth[scale=1.0]}] (m-1-1) -- node[left] {$F_E$} (m-2-1);
  \draw[-{Stealth[scale=1.0]}] (m-1-2) -- node[right] {$F_V$} (m-2-2);
  \draw[-{Stealth[scale=1.0]}] (m-2-1) -- node[below] {$s_H$} (m-2-2);
\end{tikzpicture}
\hspace{2.2em}
\begin{tikzpicture}[baseline=(n.center)]
  \matrix (n) [matrix of math nodes, row sep=2.5em, column sep=3.2em] {
    E_G & V_G \\
    E_H & V_H \\
  };
  \draw[-{Stealth[scale=1.0]}] (n-1-1) -- node[above] {$t_G$} (n-1-2);
  \draw[-{Stealth[scale=1.0]}] (n-1-1) -- node[left] {$F_E$} (n-2-1);
  \draw[-{Stealth[scale=1.0]}] (n-1-2) -- node[right] {$F_V$} (n-2-2);
  \draw[-{Stealth[scale=1.0]}] (n-2-1) -- node[below] {$t_H$} (n-2-2);
\end{tikzpicture}
\end{center}
\end{definicio}

Els grafs dirigits i els morfismes de grafs formen una categoria, que denotarem $\cat{Graph}$:\index{Graph@$\cat{Graph}$}
\begin{itemize}
\item objectes: $\Ob(\cat{Graph})$ és la col·lecció de tots els grafs dirigits;
\item morfismes: $\Hom_{\cat{Graph}}(G,H)$ és el conjunt dels morfismes de grafs $G\to H$;
\item composició: component a component; per a $F\colon G\to H$ i $F'\colon H\to K$,
\[
F'\comp F=(F'_V\comp F_V,\ F'_E\comp F_E);
\]
\item identitats: $\id_G=(\id_{V_G},\id_{E_G})$.
\end{itemize}
La composició torna a ser un morfisme de grafs, perquè les equacions es conserven:
\[
(F'_V\comp F_V)\comp s_G=F'_V\comp s_H\comp F_E=s_K\comp(F'_E\comp F_E),
\]
i anàlogament per a $t$. L'associativitat i les lleis d'identitat es compleixen component a component, perquè són les de $\cat{Set}$.

Aquest exemple és especialment rellevant perquè els grafs són l'esquelet combinatori dels diagrames: tenen vèrtexs i arestes dirigides, però encara no incorporen per si sols identitats, composicions ni axiomes categòrics.

Per fixar idees, vet aquí un graf dirigit concret amb tres vèrtexs $V_G=\{a,b,c\}$ i tres arestes:
\[
\begin{tikzpicture}[baseline=(m.center)]
  \matrix (m) [matrix of math nodes, row sep=2.6em, column sep=2.6em] {
    a & b \\
      & c \\
  };
  \draw[-{Stealth[scale=1.1]}] (m-1-1) -- node[above] {$e_1$} (m-1-2);
  \draw[-{Stealth[scale=1.1]}] (m-1-2) -- node[right] {$e_2$} (m-2-2);
  \draw[-{Stealth[scale=1.1]}] (m-1-1) -- node[below left] {$e_3$} (m-2-2);
\end{tikzpicture}
\]
Aquí $e_1$ va de $a$ a $b$, $e_2$ de $b$ a $c$ i $e_3$ de $a$ a $c$. A diferència d'una categoria, un graf no demana que hi hagi cap aresta que \enquote{representi la composició} d'$e_1$ i $e_2$: $e_3$ hi és perquè l'hem posada, no perquè cap axioma l'obligui. Aquesta és exactament la diferència entre un graf i una categoria: la categoria afegeix identitats, composició i els axiomes que les regeixen. Per a un desenvolupament complet dels grafs, amb nombrosos exemples i els homomorfismes de grafs, vegeu l'apèndix~\ref{cap:grafs}.

\subsection{Resum: objectes i morfismes dels exemples}

La taula següent recull, per a cada exemple d'aquesta secció, quines dades constitueixen els objectes i quines els morfismes.

\begin{center}
\footnotesize
\setlength{\tabcolsep}{6pt}
\begin{tabularx}{\linewidth}{@{}>{\raggedright\arraybackslash}X>{\raggedright\arraybackslash}X>{\raggedright\arraybackslash}X@{}}
\toprule
Categoria & Objectes & Morfismes \\
\midrule
$\cat{0},\cat{1},\cat{2},\cat{3}$ & marcadors abstractes & fletxes formals \\
$\cat{Set}$, $\cat{FinSet}$ & conjunts & aplicacions \\
categoria discreta & marcadors abstractes & només identitats \\
$\cat{Grp}$, $\cat{Vect}_{\mathbb K}$, $\cat{Top}$, $\cat{Mon}$, $\cat{Pos}$ & conjunts amb estructura & aplicacions que preserven l'estructura \\
$\cat{B}M$ (un monoide) & un únic objecte $\ast$ & elements de $M$ \\
un preordre $P$; $\cat{Prov}_T$ & elements de $P$; fórmules de $T$ & $x\le y$; $p\vdash_T q$ \\
$\cat{Rel}$ & conjunts & relacions \\
$\cat{Graph}$ & grafs $(V,E,s,t)$ & parelles compatibles d'aplicacions \\
\bottomrule
\end{tabularx}
\end{center}

Aquesta taula fa visible que la definició de categoria no pressuposa res sobre la naturalesa dels objectes ni dels morfismes: no cal que els objectes siguin conjunts, ni que els morfismes siguin aplicacions. En una categoria discreta, per exemple, les identitats són morfismes \emph{formals}: presentar-les com a aplicacions pressuposaria una realització concreta que no forma part de la definició. I a $\cat{B}M$ els morfismes són els elements de $M$, que en general tampoc no són aplicacions. La categoria $\cat{Rel}$ ho remarca des de l'altre costat: els seus objectes són conjunts, però els seus morfismes ---les relacions--- no són funcions.

Convé no confondre aquesta descripció \emph{estructural} de les dades amb la seva eventual \emph{codificació} conjuntista. En el marc de fonaments que adoptem, qualsevol d'aquestes dades ---inclosos els marcadors abstractes--- es pot representar per conjunts; però aquesta codificació és una tria externa, no una part de la definició de categoria. Per això la taula descriu què \emph{són} les dades i no si \enquote{són conjunts} o \enquote{són funcions}: aquesta darrera pregunta barreja la naturalesa estructural amb una representació particular.

\section{Isomorfismes}\index{isomorfisme}

Entre tots els morfismes d'una categoria, n'hi ha uns d'especialment importants: els que tenen invers. Capturen, en el llenguatge abstracte de les fletxes, la idea que dos objectes tenen la mateixa estructura des del punt de vista de la categoria.

\begin{definicio}\label{def:isomorfisme}
En una categoria $\cat{C}$, un morfisme $f\colon A\to B$ és un \textbf{isomorfisme} si existeix un morfisme $g\colon B\to A$ tal que
\[
g\comp f=\id_A
\qquad\text{i}\qquad
f\comp g=\id_B.
\]
El morfisme $g$ s'anomena \textbf{invers} de $f$. Si existeix un isomorfisme entre $A$ i $B$, diem que $A$ i $B$ són \textbf{isomorfs}, i escrivim $A\cong B$.
\end{definicio}

La situació es caracteritza per dues equacions, $g\comp f=\id_A$ i $f\comp g=\id_B$, que es llegeixen com dos triangles commutatius:
\[
\begin{tikzpicture}[baseline=(m.center)]
  \matrix (m) [matrix of math nodes, row sep=2.6em, column sep=2.6em] {
    A & B \\
    A & \\
  };
  \draw[-{Stealth[scale=1.1]}] (m-1-1) -- node[above] {$f$} (m-1-2);
  \draw[-{Stealth[scale=1.1]}] (m-1-2) -- node[right] {$g$} (m-2-1);
  \draw[-{Stealth[scale=1.1]}] (m-1-1) -- node[left] {$\id_A$} (m-2-1);
\end{tikzpicture}
\qquad\qquad
\begin{tikzpicture}[baseline=(m.center)]
  \matrix (m) [matrix of math nodes, row sep=2.6em, column sep=2.6em] {
    B & A \\
    B & \\
  };
  \draw[-{Stealth[scale=1.1]}] (m-1-1) -- node[above] {$g$} (m-1-2);
  \draw[-{Stealth[scale=1.1]}] (m-1-2) -- node[right] {$f$} (m-2-1);
  \draw[-{Stealth[scale=1.1]}] (m-1-1) -- node[left] {$\id_B$} (m-2-1);
\end{tikzpicture}
\]
Cada triangle diu que anar d'un objecte a l'altre amb $f$ o $g$ i tornar és el mateix que seguir la identitat.

\begin{proposicio}
L'invers d'un isomorfisme, si existeix, és únic.
\end{proposicio}

\begin{prova}
Suposem que $f\colon A\to B$ té dos inversos, $g$ i $g'$. Aleshores
\[
g=g\comp\id_B=g\comp(f\comp g')=(g\comp f)\comp g'=\id_A\comp g'=g'.
\]
Per tant $g=g'$. Com que l'invers és únic, el podem denotar sense ambigüitat $f^{-1}$.
\end{prova}

\begin{proposicio}
Sigui $\cat{C}$ una categoria.
\begin{enumerate}
\item Per a cada objecte $A$, la identitat $\id_A$ és un isomorfisme, amb $\id_A^{-1}=\id_A$.
\item Si $f$ és un isomorfisme, aleshores $f^{-1}$ també ho és, i $(f^{-1})^{-1}=f$.
\item Si $f\colon A\to B$ i $g\colon B\to C$ són isomorfismes, aleshores $g\comp f$ també ho és, i
\[
(g\comp f)^{-1}=f^{-1}\comp g^{-1}.
\]
\end{enumerate}
\end{proposicio}

\begin{prova}
(1) De $\id_A\comp\id_A=\id_A$ se segueix que $\id_A$ és el seu propi invers. (2) Les equacions $g\comp f=\id_A$ i $f\comp g=\id_B$ que fan de $g=f^{-1}$ l'invers de $f$ diuen també que $f$ és l'invers de $g$; per tant $f^{-1}$ és un isomorfisme i $(f^{-1})^{-1}=f$. (3) Comprovem que $f^{-1}\comp g^{-1}$ és l'invers de $g\comp f$:
\[
(g\comp f)\comp(f^{-1}\comp g^{-1})=g\comp(f\comp f^{-1})\comp g^{-1}=g\comp\id_B\comp g^{-1}=g\comp g^{-1}=\id_C,
\]
i anàlogament $(f^{-1}\comp g^{-1})\comp(g\comp f)=\id_A$.
\end{prova}

\begin{corollari}
Donada una categoria $\cat{C}$, definim sobre $\Ob(\cat{C})$ la relació $\cong_{\cat{C}}$ per: $A\cong_{\cat{C}}B$ si i només si existeix un isomorfisme $f\colon A\to B$. Aleshores $\cong_{\cat{C}}$ és una relació d'equivalència.
\end{corollari}

\begin{prova}
És reflexiva: $\id_A$ és un isomorfisme per (1), de manera que $A\cong_{\cat{C}}A$. És simètrica: si $f$ és un isomorfisme de $A$ a $B$, aleshores $f^{-1}$ ho és de $B$ a $A$ per (2), de manera que $B\cong_{\cat{C}}A$. És transitiva: si $f\colon A\to B$ i $g\colon B\to C$ són isomorfismes, aleshores $g\comp f$ també ho és per (3), de manera que $A\cong_{\cat{C}}C$.
\end{prova}

\begin{observacio}[Igualtat i isomorfisme]\label{obs:igualtat-isomorfisme}\index{isomorfisme!i igualtat}
Convé fixar des d'ara una convenció que el llibre manté sempre. Escrivim $A=B$ només quan $A$ i $B$ són \emph{el mateix} objecte, i $A\cong B$ quan hi ha un isomorfisme entre ells. Quan un isomorfisme és distingit ---no depèn de cap tria---, l'anomenem \textbf{canònic}, però continua sent un isomorfisme i no una igualtat (en veurem un primer cas en el monoide lliure, exemple~\ref{ex:monoide-lliure-cat}).

La relació $\cong_{\cat{C}}$ permet parlar d'una \emph{igualtat estructural relativa a $\cat{C}$}: dos objectes isomorfs a $\cat{C}$ no es poden distingir per cap propietat formulada només amb els morfismes, la composició i les identitats de $\cat{C}$, perquè l'isomorfisme la transporta d'un a l'altre (ho precisarem a l'observació~\ref{obs:invariancia}). Diem, doncs, que tenen la mateixa estructura \emph{des del punt de vista de $\cat{C}$}, i per això sovint la pregunta rellevant no és si $A$ i $B$ són iguals, sinó si són isomorfs a $\cat{C}$.

Aquesta igualtat no és absoluta: depèn dels morfismes que la categoria considera admissibles. Per exemple, el conjunt $\{0,1\}$ amb l'ordre discret i el mateix conjunt amb l'ordre $0<1$ no són isomorfs a $\cat{Pos}$, perquè cap bijecció monòtona no té inversa monòtona, però els seus conjunts subjacents sí que són isomorfs a $\cat{Set}$. Canviar de categoria canvia, doncs, quins objectes compten com a estructuralment iguals. Aquesta és una de les idees centrals de la teoria de categories: la forma d'un objecte sempre es mira dins d'una categoria determinada.
La jerarquia completa ---igualtat, isomorfisme, isomorfisme natural i equivalència---, amb un contraexemple per a cada pas, es tracta a l'apèndix~\refalt{cap:contraexemples}{de contraexemples del llibre complet}.
\end{observacio}

\begin{exemple}
A $\cat{Set}$, els isomorfismes són exactament les aplicacions bijectives. En efecte, una aplicació té inversa per la composició si i només si és injectiva i exhaustiva.
\end{exemple}

\begin{exemple}[Isomorfismes en les categories de conjunts estructurats]
A $\cat{Grp}$, $\cat{Vect}_{\mathbb K}$ i $\cat{Mon}$, els isomorfismes són les aplicacions bijectives que preserven l'estructura i la inversa de les quals també la preserva; recuperem així les nocions habituals d'isomorfisme de grups, d'espais vectorials i de monoides. A $\cat{Top}$, en canvi, un isomorfisme és un homeomorfisme\index{homeomorfisme}: una aplicació contínua i bijectiva amb inversa contínua. Aquí es veu que la bijectivitat no basta: hi ha aplicacions contínues i bijectives que no són isomorfismes, perquè la seva inversa no és contínua. La definició categòrica capta la noció correcta sense haver-la d'enunciar cas per cas.
\end{exemple}

\begin{exemple}
En un preordre, un isomorfisme entre $x$ i $y$ vol dir que hi ha fletxes $x\to y$ i $y\to x$, és a dir,
\[
x\le y
\qquad\text{i}\qquad
y\le x.
\]
En un ordre parcial això obliga $x=y$. En un preordre general, en canvi, poden existir objectes isomorfs diferents: són exactament els elements equivalents per la relació $x\sim y$ si $x\le y$ i $y\le x$.
\end{exemple}

\begin{exemple}[Isomorfisme i equivalència lògica]\label{ex:iso-equivalencia-logica}
A la categoria prima $\cat{Prov}_T$ d'un sistema deductiu $T$, dues fórmules $p$ i $q$ són isomorfes exactament quan
\[
p\vdash_T q
\qquad\text{i}\qquad
q\vdash_T p.
\]
Per tant, l'isomorfisme categòric es tradueix aquí en \textbf{interdeduïbilitat}. Aquest és un primer exemple de com una noció categòrica adquireix una lectura lògica natural.
\end{exemple}

\begin{observacio}
L'exemple dels preordres mostra per què la definició categòrica és la bona. Es podria pensar a definir un isomorfisme com un morfisme \enquote{bijectiu}, però això requeriria parlar dels elements dels objectes, i no sempre té sentit. A més, fins i tot quan en té, pot donar la noció equivocada: a $\cat{Pos}$ existeixen aplicacions monòtones bijectives que no són isomorfismes perquè la seva inversa no és monòtona. La definició mitjançant composició i identitats evita aquest parany.
\end{observacio}

\begin{observacio}[Propietats invariants per isomorfisme]\label{obs:invariancia}
Sigui $\cat{C}$ una categoria. Una propietat dels objectes de $\cat{C}$ és \textbf{invariant per isomorfisme} a $\cat{C}$ si, sempre que un objecte la té, tot objecte isomorf a ell a $\cat{C}$ també la té. Com la relació $\cong_{\cat{C}}$ de l'observació anterior, aquesta noció és sempre relativa a una categoria determinada. Una propietat definida només a partir de l'estructura categòrica de $\cat{C}$ ---els seus objectes, morfismes, composició i identitats, sense fer servir la igualtat entre objectes--- ha de ser invariant per isomorfisme a $\cat{C}$, perquè un isomorfisme transporta qualsevol afirmació d'aquesta mena d'un objecte a l'altre. Per això la invariància per isomorfisme és un criteri bàsic per reconèixer quan una propietat depèn de la forma de l'objecte en la categoria i no de la seva presentació concreta.

Aquesta noció dona un mètode útil. Per provar que dos objectes \emph{són} isomorfs, n'hi ha prou d'exhibir un isomorfisme entre ells. Provar que \emph{no} ho són sol ser més difícil, perquè caldria descartar tots els isomorfismes possibles; la via pràctica és trobar una propietat invariant per isomorfisme a la categoria en qüestió que un dels objectes tingui i l'altre no. Per exemple, a $\cat{Pos}$, siguin $A$ la cadena $a<b<c$ i $B$ l'ordre amb $x<y$ i $x<z$ però $y,z$ incomparables. Comprovar directament que no són isomorfs exigiria descartar totes les bijeccions monòtones amb inversa monòtona; en canvi, n'hi ha prou d'observar que \emph{estar totalment ordenat} és una propietat invariant per isomorfisme, que $A$ té i $B$ no. Per tant $A$ i $B$ no són isomorfs.

La implicació inversa no és automàtica: una propietat pot ser invariant per isomorfisme sense haver estat formulada internament en el llenguatge categòric. L'exemple anterior n'és una mostra: \emph{estar totalment ordenat} és invariant per isomorfisme a $\cat{Pos}$, però l'hem definit parlant dels elements de l'ordre, no de fletxes.
\end{observacio}

\begin{observacio}
La noció d'isomorfisme és autènticament categòrica: es formula només amb composició i identitats, de manera que té sentit a \emph{qualsevol} categoria, sense referència als elements dels objectes. Això la fa més abstracta i més general que les nocions d'isomorfisme pròpies de cada context ---bijeccions, isomorfismes de grups, homeomorfismes---, que en són casos particulars.

Aquesta generalitat es fa visible en un cas que ja coneixem. Havíem identificat els monoides amb les categories d'un sol objecte; ara podem identificar els grups amb exactament aquelles categories d'un sol objecte en què tota fletxa és un isomorfisme. Portant la mateixa idea a un nombre qualsevol d'objectes s'obté una noció d'importància considerable en les matemàtiques actuals: un \emph{grupoide} és una categoria en què tot morfisme és un isomorfisme.\index{grupoide}

El grupoide unifica dos casos que ja hem trobat. Un grup és un grupoide amb un sol objecte, com acabem de veure. I un preordre en què tot morfisme és un isomorfisme és, precisament, un en què $x\le y$ implica $y\le x$: una relació simètrica, és a dir, una \emph{relació d'equivalència}. Així, una relació d'equivalència no és res més que un grupoide prim: les classes d'isomorfisme dels seus objectes són les classes d'equivalència.
\end{observacio}

\begin{definicio}\index{endomorfisme}\index{automorfisme}
Un morfisme $f\colon A\to A$, amb domini i codomini iguals, s'anomena \textbf{endomorfisme} de $A$. Un endomorfisme que a més és un isomorfisme s'anomena \textbf{automorfisme} de $A$.
\end{definicio}

Els automorfismes d'un objecte $A$, amb la composició, formen un grup, el \textbf{grup d'automorfismes} de $A$, que escrivim $\Aut(A)$.\index{Aut@$\Aut$} Torna a aparèixer així la identificació d'abans: $\Aut(A)$ és el grup associat a la categoria d'un sol objecte que té els automorfismes de $A$ com a fletxes.

\section{Classes especials de morfismes}

Els isomorfismes són els morfismes invertibles: tenen un invers per tots dos costats. Sovint, però, un morfisme té invers només per un costat, o satisfà una propietat de cancel·lació més feble que la invertibilitat. Aquestes nocions ---monomorfismes, epimorfismes, seccions i retraccions--- són bàsiques. A $\cat{Set}$, els monomorfismes i els epimorfismes coincideixen amb les aplicacions injectives i exhaustives; com veurem, però, aquesta coincidència és pròpia de $\cat{Set}$ i d'algunes categories concretes, no una propietat de les categories en general.

Recordem que, a $\cat{Set}$, una aplicació $f$ és \emph{injectiva} si $f(a)=f(a')$ implica $a=a'$, i \emph{exhaustiva} si tot element del codomini té alguna antiimatge. Aquestes definicions parlen d'elements; les versions categòriques que segueixen les reformulen només amb fletxes.

\subsection{Monomorfismes i epimorfismes}\label{subsec:mono-epi}

\begin{definicio}\index{monomorfisme}\index{epimorfisme}
\label{def:mono-epi}
Un morfisme $f\colon A\to B$ és un \textbf{monomorfisme} (o \textbf{mono}), i escrivim $f\colon A\rightarrowtail B$, si es pot cancel·lar per l'esquerra: per a tota parella de morfismes $g,h\colon C\to A$,
\[
\begin{tikzpicture}[baseline=(m.center)]
  \matrix (m) [matrix of math nodes, column sep=3.4em] { C & A & B \\ };
  \draw[-{Stealth[scale=1.0]}] ([yshift=2.6pt]m-1-1.east) -- node[above] {$g$} ([yshift=2.6pt]m-1-2.west);
  \draw[-{Stealth[scale=1.0]}] ([yshift=-2.6pt]m-1-1.east) -- node[below] {$h$} ([yshift=-2.6pt]m-1-2.west);
  \draw[-{Stealth[scale=1.1]}] (m-1-2) -- node[above] {$f$} (m-1-3);
\end{tikzpicture}
\qquad
f\comp g=f\comp h\ \Longrightarrow\ g=h.
\]
Un morfisme $f\colon A\to B$ és un \textbf{epimorfisme} (o \textbf{epi}), i escrivim $f\colon A\twoheadrightarrow B$, si es pot cancel·lar per la dreta: per a tota parella $g,h\colon B\to C$,
\[
\begin{tikzpicture}[baseline=(m.center)]
  \matrix (m) [matrix of math nodes, column sep=3.4em] { A & B & C \\ };
  \draw[-{Stealth[scale=1.1]}] (m-1-1) -- node[above] {$f$} (m-1-2);
  \draw[-{Stealth[scale=1.0]}] ([yshift=2.6pt]m-1-2.east) -- node[above] {$g$} ([yshift=2.6pt]m-1-3.west);
  \draw[-{Stealth[scale=1.0]}] ([yshift=-2.6pt]m-1-2.east) -- node[below] {$h$} ([yshift=-2.6pt]m-1-3.west);
\end{tikzpicture}
\qquad
g\comp f=h\comp f\ \Longrightarrow\ g=h.
\]
\end{definicio}

\begin{observacio}[Com cal llegir aquests diagrames]\label{obs:llegir-mono-epi}
Els diagrames de la definició~\ref{def:mono-epi} \emph{no} afirmen que res commuti: són esquemes de prova, i s'han de llegir per a cada parella possible de fletxes $g,h$. Al diagrama del monomorfisme hi ha dos camins de $C$ a $B$, \enquote{primer $g$, després $f$} i \enquote{primer $h$, després $f$}, i hi ha també les dues fletxes paral·leles $g,h$ de $C$ a $A$. Recordem que un diagrama amb dues fletxes paral·leles commuta exactament quan són iguals. La definició diu, doncs:
\begin{quote}
si els dos camins de $C$ a $B$ coincideixen, aleshores les fletxes paral·leles ja coincideixen.
\end{quote}
Amb el vocabulari de la secció~\ref{sec:diagrames}: l'antecedent no és que el diagrama commuti, sinó només que commutin els camins que acaben a $B$; la conclusió és que el diagrama commuta del tot. Un monomorfisme és, doncs, una fletxa per a la qual \emph{commutar després de $f$ implica commutar}.
Llegit en sentit contrari: si $g\neq h$, aleshores $f\comp g\neq f\comp h$. Un monomorfisme \emph{no confon} mai dues fletxes diferents que arriben a $A$; després de compondre-les amb $f$, continuen sent diferents.

Cal anar amb compte amb una lectura massa ràpida. Que $f$ sigui mono no impedeix que hi hagi parelles $g\neq h\colon C\to A$; el diagrama amb dues fletxes paral·leles diferents existeix perfectament. El que exclou és que, sent diferents, es facin iguals en compondre-les amb $f$.

El diagrama de l'epimorfisme es llegeix de la mateixa manera, amb les fletxes de prova a l'altre costat: si $g\neq h\colon B\to C$, aleshores $g\comp f\neq h\comp f$. Dues fletxes que surten de $B$ i coincideixen després de $f$ ja eren iguals; $f$ arriba prou a $B$ perquè cap diferència entre $g$ i $h$ no se li escapi. A $\cat{Set}$, prenent com a fletxes de prova les aplicacions des d'un conjunt d'un sol element, que són els elements de $A$, la lectura del monomorfisme es converteix exactament en la injectivitat, com veurem tot seguit.
\end{observacio}

\begin{definicio}\index{bimorfisme}
Un morfisme que és alhora un monomorfisme i un epimorfisme s'anomena \textbf{bimorfisme}.
\end{definicio}

\begin{exemple}
A $\cat{Set}$, els monomorfismes són exactament les aplicacions \textbf{injectives} i els epimorfismes, les \textbf{exhaustives}. Vegem el cas dels monomorfismes: si $f$ és injectiva i $f\comp g=f\comp h$, aleshores per a cada $c$ tenim $f(g(c))=f(h(c))$, d'on $g(c)=h(c)$; recíprocament, si $f$ no és injectiva, hi ha $a\neq a'$ amb $f(a)=f(a')$, i les dues aplicacions constants $g,h\colon\{\ast\}\to A$ amb valors $a$ i $a'$ compleixen $f\comp g=f\comp h$ però $g\neq h$.

Per als epimorfismes, sigui $f\colon A\to B$. Si $f$ és exhaustiva i $g\comp f=h\comp f$, aleshores per a cada $b\in B$ existeix $a$ amb $f(a)=b$, i $g(b)=g(f(a))=h(f(a))=h(b)$, d'on $g=h$. Recíprocament, si $f$ no és exhaustiva, hi ha $b_0\in B$ sense antiimatge, és a dir, $b_0\notin f(A)$. Definim $g,h\colon B\to\{0,1\}$ a \emph{tot} $B$ per
\[
g(b)=0 \text{ per a tot } b\in B,
\qquad
h(b)=\begin{cases}1 & \text{si } b=b_0,\\ 0 & \text{si } b\neq b_0.\end{cases}
\]
Aleshores $g$ i $h$ coincideixen a tot $B$ tret de $b_0$, de manera que $g\neq h$. Però per a tot $a\in A$ tenim $f(a)\neq b_0$, ja que $b_0\notin f(A)$; per tant $g(f(a))=0=h(f(a))$, és a dir, $g\comp f=h\comp f$. Això mostra que $f$ no és un epimorfisme.

Convé notar que cap dels dos arguments no tria una \emph{secció} de $f$: per als monomorfismes es comparen dues seleccions d'elements des d'un conjunt unitari, i per als epimorfismes es construeixen explícitament dues aplicacions cap a $\{0,1\}$. En particular, la caracterització \enquote{epimorfisme $=$ exhaustiva} a $\cat{Set}$ no fa servir l'axioma de l'elecció, i no s'ha de confondre amb l'enunciat, ben diferent, que tota aplicació exhaustiva \emph{admet una secció} ---aquest sí equivalent a l'axioma de l'elecció (vegeu més avall).
\end{exemple}

\begin{proposicio}
Tot isomorfisme és un bimorfisme: si $f$ és un isomorfisme, aleshores és alhora un monomorfisme i un epimorfisme.
\end{proposicio}

\begin{prova}
Sigui $f\colon A\to B$ un isomorfisme amb invers $f^{-1}$. Si $f\comp g=f\comp h$, aleshores
\[
g=f^{-1}\comp(f\comp g)=f^{-1}\comp(f\comp h)=h,
\]
de manera que $f$ és mono. Si $g\comp f=h\comp f$, aleshores
\[
g=(g\comp f)\comp f^{-1}=(h\comp f)\comp f^{-1}=h,
\]
de manera que $f$ és epi.
\end{prova}

\begin{observacio}
A $\cat{Set}$, un morfisme és un isomorfisme si i només si és alhora mono i epi, és a dir, si i només si, com a aplicació, és injectiu i exhaustiu, o sigui \textbf{bijectiu}. Dit d'una altra manera, a $\cat{Set}$ els bimorfismes són exactament els isomorfismes. El recíproc de la proposició anterior, però, no és cert: en general un bimorfisme no és un isomorfisme. L'única fletxa no trivial de la categoria $\cat{2}$ és un bimorfisme ---per la simple raó que gairebé no hi ha morfismes amb què fer les cancel·lacions---, però no té invers, de manera que no és un isomorfisme. Un exemple més substancial es dona a $\cat{Top}$: una bijecció contínua és un bimorfisme, però no és un isomorfisme si la seva inversa no és contínua.
\end{observacio}

\begin{observacio}[Atenció: raonar amb elements]\label{obs:raonar-elements}\index{element!perill de raonar amb}
A $\cat{Set}$ podem raonar amb elements per dos motius: cada element $a\in A$ és una fletxa $\{\ast\}\to A$, i dues aplicacions $A\to B$ són iguals si coincideixen sobre tots els elements. En una categoria qualsevol, cap d'aquestes dues coses no està garantida.
\begin{itemize}
\item \emph{Pot ser que no hi hagi elements.} En un preordre o en un monoide vist com a categoria, una fletxa no envia res enlloc, i preguntar si és \enquote{injectiva} no té sentit. Mono i epi, en canvi, sempre tenen sentit, perquè es defineixen per composició (exercici~\ref{exr:pr-prima-mono-epi} de la Pràctica).
\item \emph{L'aplicació subjacent pot enganyar.} A $\cat{Top}$, una bijecció contínua és mono i epi sense ser un isomorfisme (observació anterior). A $\cat{Mon}$, la inclusió $\iota\colon(\mathbb N,+,0)\hookrightarrow(\mathbb Z,+,0)$ és un epimorfisme, tot i que no és exhaustiva. Si dos homomorfismes $g,h\colon\mathbb Z\to M$ coincideixen sobre $\mathbb N$, aleshores $g(-1)$ i $h(-1)$ són tots dos inversos de $g(1)=h(1)$, i en un monoide l'invers d'un element, si existeix, és únic. Per tant, $g(-1)=h(-1)$ i, com que tot enter és suma d'un natural i de còpies de $-1$, $g=h$.
\item \emph{Els \enquote{punts} poden no separar fletxes.} A $\cat{Grp}$, el grup trivial $1$ fa el paper de $\{\ast\}$, però de $1$ a qualsevol grup només hi ha un homomorfisme. Les fletxes $1\to G$ no distingeixen, doncs, cap parell d'homomorfismes diferents.
\end{itemize}
Per això, en una categoria general els arguments es fan amb fletxes i composicions, i reservem els adjectius \emph{injectiu}, \emph{exhaustiu} i \emph{bijectiu} per a les aplicacions, incloses les aplicacions subjacents dels morfismes d'una categoria concreta. La definició de monomorfisme és, de fet, una \enquote{injectivitat} respecte de \emph{totes} les fletxes de prova $C\to A$, i no només respecte dels elements.
\end{observacio}

\subsection{Seccions i retraccions}\label{subsec:seccions-retraccions}

\begin{definicio}\index{secció}\index{retracció}\index{invers!per un costat}
Considerem un morfisme $f\colon A\to B$. Diem que un morfisme $s\colon B\to A$ és una \textbf{secció} de $f$ quan es compleix $f\comp s=\id_B$. Això equival a dir que el diagrama següent és commutatiu:
\[
\begin{tikzpicture}[baseline=(m.center)]
  \matrix (m) [matrix of math nodes, row sep=2.6em, column sep=2.6em] {
    B & A \\
    B & \\
  };
  \draw[-{Stealth[scale=1.1]}] (m-1-1) -- node[above] {$s$} (m-1-2);
  \draw[-{Stealth[scale=1.1]}] (m-1-2) -- node[right] {$f$} (m-2-1);
  \draw[-{Stealth[scale=1.1]}] (m-1-1) -- node[left] {$\id_B$} (m-2-1);
\end{tikzpicture}
\]
També diem en aquest cas que $s$ és un \emph{invers per la dreta} de $f$.

Simètricament, diem que un morfisme $r\colon B\to A$ és una \textbf{retracció} de $f$ quan es compleix $r\comp f=\id_A$. Això equival a dir que el diagrama següent és commutatiu:
\[
\begin{tikzpicture}[baseline=(m.center)]
  \matrix (m) [matrix of math nodes, row sep=2.6em, column sep=2.6em] {
    A & B \\
    A & \\
  };
  \draw[-{Stealth[scale=1.1]}] (m-1-1) -- node[above] {$f$} (m-1-2);
  \draw[-{Stealth[scale=1.1]}] (m-1-2) -- node[right] {$r$} (m-2-1);
  \draw[-{Stealth[scale=1.1]}] (m-1-1) -- node[left] {$\id_A$} (m-2-1);
\end{tikzpicture}
\]
També diem en aquest cas que $r$ és un \emph{invers per l'esquerra} de $f$.
\end{definicio}

Una conseqüència òbvia d'aquestes dues definicions és que quan es compleix, per exemple, $f\comp s=\id_B$, se'n segueix alhora que $s$ és una secció de $f$ i que $f$ és una retracció de $s$; el mateix succeeix amb $r\comp f=\id_A$, on $r$ és una retracció de $f$ i $f$ és una secció de $r$.

Un morfisme pot tenir cap, una o diverses seccions, i igualment cap, una o diverses retraccions. Vegem-ne exemples a $\cat{Set}$:
\begin{itemize}
\item L'única aplicació $f\colon\emptyset\to\{*\}$ no té cap secció: no existeix cap aplicació $\{*\}\to\emptyset$.
\item La injecció $f\colon\{0\}\hookrightarrow\{0,1\}$ associada a la inclusió $\{0\}\subseteq\{0,1\}$ té una única retracció, l'única aplicació $\{0,1\}\to\{0\}$; en canvi no té cap secció.
\item L'aplicació $f\colon\{0,1\}\to\{*\}$ té dues seccions distintes, $*\mapsto 0$ i $*\mapsto 1$: cada tria d'una antiimatge en dona una.
\end{itemize}
Tenir secció o retracció és, doncs, una propietat especial de $f$, i quan en té, en general no és única.

\begin{proposicio}\label{prop:escindit-mono-epi}
Sigui $f\colon A\to B$. Si $f$ té alguna secció, aleshores $f$ és un epimorfisme. Si $f$ té alguna retracció, aleshores $f$ és un monomorfisme.
\end{proposicio}

\begin{prova}
Suposem que $f$ té una retracció $r$, amb $r\comp f=\id_A$, i siguin $g,h\colon C\to A$ amb $f\comp g=f\comp h$. Component per l'esquerra amb $r$,
\[
g=\id_A\comp g=(r\comp f)\comp g=r\comp(f\comp g)=r\comp(f\comp h)=(r\comp f)\comp h=h,
\]
de manera que $f$ és mono. El cas de la secció és anàleg: si $f$ té secció $s$, amb $f\comp s=\id_B$, i $g,h\colon B\to C$ amb $g\comp f=h\comp f$, aleshores component per la dreta amb $s$,
\[
g=g\comp\id_B=g\comp(f\comp s)=(g\comp f)\comp s=(h\comp f)\comp s=h\comp(f\comp s)=h\comp\id_B=h,
\]
de manera que $f$ és epi.
\end{prova}

\begin{definicio}\index{morfisme!escindit}\index{epimorfisme!escindit}\index{monomorfisme!escindit}
Diem que un morfisme $f$ queda \textbf{escindit} quan té una secció o una retracció. En concret, si $f$ té una secció, aleshores $f$ és un epimorfisme i se'n diu \textbf{epimorfisme escindit}; si $f$ té una retracció, aleshores $f$ és un monomorfisme i se'n diu \textbf{monomorfisme escindit}.
\end{definicio}

\begin{observacio}
El recíproc no val en general: hi ha monomorfismes que no tenen cap retracció (i per tant no són escindits) i epimorfismes que no tenen cap secció. En efecte, a $\cat{Set}$ tota aplicació injectiva amb domini no buit té alguna retracció i tota exhaustiva té alguna secció ---aquesta darrera afirmació és, de fet, equivalent a l'\emph{axioma de l'elecció}---, però en altres categories la situació és diferent.

Finalment, un morfisme $f$ és un \textbf{isomorfisme} si i només si té alhora alguna secció i alguna retracció. En aquest cas la secció i la retracció són forçosament \emph{úniques} i coincideixen: si $f\comp s=\id_B$ i $r\comp f=\id_A$, aleshores
\[
r=r\comp\id_B=r\comp(f\comp s)=(r\comp f)\comp s=\id_A\comp s=s.
\]
En canvi, la unicitat d’una secció o d’una retracció, considerada
separadament, no caracteritza els isomorfismes en una categoria general; per exemple, la injecció
\[
 i\colon\{0\}\hookrightarrow\{0,1\}
\]
té una única retracció $r\colon\{0,1\}\to\{0\}$, però no és un isomorfisme. No és una subtilesa avançada, sinó una contradicció amb un exemple de la mateixa subsecció.
\end{observacio}

\section{Construccions sobre categories}

A partir de categories donades se'n poden construir de noves. Aquestes construccions seran útils constantment, i algunes ---com la categoria oposada--- són fonamentals.

\subsection{La categoria oposada}\index{categoria!oposada}

\begin{definicio}
Donada una categoria $\cat{C}$, la seva \textbf{categoria oposada} $\cat{C}\op$ té:
\begin{itemize}
\item objectes: els mateixos que $\cat{C}$, $\Ob(\cat{C}\op)=\Ob(\cat{C})$;
\item morfismes: per a cada morfisme $f\colon A\to B$ de $\cat{C}$, un morfisme $f\op\colon B\to A$ de $\cat{C}\op$, de manera que
\[
\Hom_{\cat{C}\op}(B,A)=\{f\op\mid f\in\Hom_{\cat{C}}(A,B)\};
\]
\item composició: s'obté invertint l'ordre; per a $f\colon A\to B$ i $g\colon B\to C$ de $\cat{C}$,
\[
f\op\comp g\op=(g\comp f)\op\colon C\to A;
\]
\item identitats: la identitat de $A$ a $\cat{C}\op$ és $(\id_A)\op$.
\end{itemize}
És a dir, $\cat{C}\op$ s'obté invertint el sentit de totes les fletxes.
\end{definicio}

Visualment, la composició a $\cat{C}$ i la corresponent a $\cat{C}\op$ inverteixen l'ordre:
\[
\begin{array}{c@{\qquad\qquad}c}
\cat{C}: & \cat{C}\op: \\[0.6em]
\begin{tikzpicture}[baseline=(m.center)]
  \matrix (m) [matrix of math nodes, row sep=2.6em, column sep=2.6em] {
    A & B \\
      & C \\
  };
  \draw[-{Stealth[scale=1.1]}] (m-1-1) -- node[above] {$f$} (m-1-2);
  \draw[-{Stealth[scale=1.1]}] (m-1-2) -- node[right] {$g$} (m-2-2);
  \draw[-{Stealth[scale=1.1]}] (m-1-1) -- node[below left] {$g\comp f$} (m-2-2);
\end{tikzpicture}
&
\begin{tikzpicture}[baseline=(m.center)]
  \matrix (m) [matrix of math nodes, row sep=2.6em, column sep=2.6em] {
    A & B \\
      & C \\
  };
  \draw[-{Stealth[scale=1.1]}] (m-1-2) -- node[above] {$f\op$} (m-1-1);
  \draw[-{Stealth[scale=1.1]}] (m-2-2) -- node[right] {$g\op$} (m-1-2);
  \draw[-{Stealth[scale=1.1]}] (m-2-2) -- node[below left] {$f\op\comp g\op$} (m-1-1);
\end{tikzpicture}
\end{array}
\]
L'associativitat a $\cat{C}\op$ es dedueix de la de $\cat{C}$ invertint l'ordre, i les lleis d'identitat, de les de $\cat{C}$: per exemple, $f\op\comp(\id_B)\op=(\id_B\comp f)\op=f\op$. La categoria oposada és la base de la \textbf{dualitat}\index{dualitat}, que desenvolupem a la secció~\ref{sec:dualitat}: tota afirmació formulada purament en el llenguatge categòric dona, invertint les fletxes, una afirmació dual, i aquesta correspondència es pot fer del tot precisa.

\begin{observacio}
Aplicar dues vegades l'oposada retorna la categoria de partida:
\[
(\cat{C}\op)\op=\cat{C}.
\]
Invertir les fletxes i tornar-les a invertir les deixa com estaven.
\end{observacio}

Vegem com queda l'oposada en dos casos. Quan un preordre $(P,\le)$ es considera com una categoria ---amb un morfisme $x\to y$ quan $x\le y$---, la seva oposada té un morfisme $x\to y$ exactament quan $y\le x$. I quan un monoide $(M,\ast,e)$ es considera com una categoria, la seva oposada té els mateixos morfismes ---els elements de $M$---, però la composició de $a$ i $b$ és $b\ast a$ en lloc de $a\ast b$.

\subsection{La categoria producte}\label{subsec:categoria-producte}\index{categoria!producte}

\begin{definicio}
Donades dues categories $\cat{C}$ i $\cat{D}$, la \textbf{categoria producte} $\cat{C}\times\cat{D}$ té:
\begin{itemize}
\item objectes: $\Ob(\cat{C}\times\cat{D})$ és la col·lecció de les parelles $(A,X)$ amb $A\in\Ob(\cat{C})$ i $X\in\Ob(\cat{D})$;
\item morfismes: $\Hom_{\cat{C}\times\cat{D}}\bigl((A,X),(B,Y)\bigr)=\Hom_{\cat{C}}(A,B)\times\Hom_{\cat{D}}(X,Y)$, és a dir, les parelles $(f,g)$ amb $f\colon A\to B$ a $\cat{C}$ i $g\colon X\to Y$ a $\cat{D}$;
\item composició: component a component,
\[
(f',g')\comp(f,g)=(f'\comp f,\ g'\comp g);
\]
\item identitats: $\id_{(A,X)}=(\id_A,\id_X)$.
\end{itemize} Els axiomes de categoria per a $\cat{C}\times\cat{D}$ se segueixen dels de $\cat{C}$ i $\cat{D}$ aplicats a cada component.
\end{definicio}

\begin{exemple}
Prenem la categoria $\cat{2}$ de la subsecció~\ref{subsec:finites} ---dos objectes i una única fletxa no trivial---, i n'escrivim els objectes com $0$ i $1$, amb la fletxa no trivial $0\to 1$. La categoria producte $\cat{2}\times\cat{2}$ té quatre objectes, que es disposen de manera natural en un quadrat:
\[
\begin{tikzpicture}[baseline=(m.center)]
  \matrix (m) [matrix of math nodes, row sep=2.6em, column sep=3.4em] {
    (0,0) & (0,1) \\
    (1,0) & (1,1) \\
  };
  \draw[-{Stealth[scale=1.1]}] (m-1-1) -- (m-1-2);
  \draw[-{Stealth[scale=1.1]}] (m-1-1) -- (m-2-1);
  \draw[-{Stealth[scale=1.1]}] (m-1-2) -- (m-2-2);
  \draw[-{Stealth[scale=1.1]}] (m-2-1) -- (m-2-2);
  \draw[-{Stealth[scale=1.1]}] (m-1-1) -- (m-2-2);
\end{tikzpicture}
\]
Cada fletxa és una parella de fletxes de $\cat{2}$. Com que $\cat{2}$ és prima, també ho és $\cat{2}\times\cat{2}$: entre dos objectes hi ha com a màxim un morfisme. Per això els dos camins del quadrat de $(0,0)$ a $(1,1)$ tenen necessàriament la mateixa composició, que és la diagonal; el quadrat commuta.
\end{exemple}

\begin{observacio}[Dos usos de \enquote{producte}]\label{obs:dos-productes}\index{producte!d'objectes i de categories}
No s'han de confondre dues construccions que porten el mateix nom. La categoria producte $\cat{C}\times\cat{D}$ és una \emph{categoria} nova, construïda a partir de dues categories. En canvi, el producte $A\times B$ de dos objectes ---que la introducció fa servir com a primer exemple de propietat universal i que avançarem a la subsecció~\ref{subsec:producte-coproducte}--- és un \emph{objecte} d'una mateixa categoria, caracteritzat per una propietat universal. Les dues nocions estan relacionades: quan disposem de la categoria $\cat{Cat}$ de les categories petites (capítol~\ref{cap:functors}), si $\cat{C}$ i $\cat{D}$ són petites, la categoria $\cat{C}\times\cat{D}$ es podrà veure com el producte dels objectes $\cat{C}$ i $\cat{D}$ dins de $\cat{Cat}$. Però no són la mateixa construcció: l'una fabrica una categoria a partir de dues, i l'altra viu dins d'una sola categoria.
\end{observacio}

\subsection{Subcategories}\label{subsec:subcategories}\index{subcategoria}

Una altra manera d'obtenir noves categories a partir de categories antigues és restringint-ne els objectes o les fletxes.

\begin{definicio}
Una \textbf{subcategoria} $\cat{S}$ d'una categoria $\cat{C}$ ve donada per una subcol·lecció d'objectes $\Ob(\cat{S})\subseteq\Ob(\cat{C})$ i, per a cada parella $A,B\in\Ob(\cat{S})$, una subcol·lecció de morfismes $\Hom_{\cat{S}}(A,B)\subseteq\Hom_{\cat{C}}(A,B)$, de manera que:
\begin{itemize}
\item la identitat $\id_A$ pertany a $\Hom_{\cat{S}}(A,A)$ per a cada $A\in\Ob(\cat{S})$;
\item la composició de dos morfismes de $\cat{S}$ componibles pertany a $\cat{S}$.
\end{itemize}
Amb aquestes condicions, $\cat{S}$ és una categoria per dret propi, amb la composició i les identitats heretades de $\cat{C}$.
\end{definicio}

Dos casos particulars, en cert sentit oposats, són especialment útils. Es diu que $\cat{S}$ és una \textbf{subcategoria plena}\index{subcategoria!plena} quan $\Hom_{\cat{S}}(A,B)=\Hom_{\cat{C}}(A,B)$ per a tot $A,B\in\Ob(\cat{S})$; una subcategoria plena queda determinada, doncs, només per la seva col·lecció d'objectes. Simètricament, es diu que $\cat{S}$ és una \textbf{subcategoria ampla}\index{subcategoria!ampla} quan $\Ob(\cat{S})=\Ob(\cat{C})$, encara que només contingui alguns dels morfismes. En resum: una subcategoria plena restringeix els objectes i conserva tots els morfismes entre ells; una subcategoria ampla conserva tots els objectes i en restringeix els morfismes.

\begin{exemple}
$\cat{FinSet}$ és una subcategoria plena de $\cat{Set}$: pren com a objectes els conjunts finits i, entre dos d'ells, totes les aplicacions. En canvi, la subcategoria de $\cat{Set}$ que té tots els conjunts per objectes però només les aplicacions \emph{injectives} per morfismes és una subcategoria ampla que no és plena: conserva tots els objectes, però deixa fora les aplicacions que no són injectives.
\end{exemple}

\subsection{Categories sobre i sota un objecte}\index{categoria!sobre un objecte}\index{categoria!sota un objecte}

Fixat un objecte $A$ d'una categoria $\cat{C}$, se'n pot construir una categoria nova els objectes de la qual són els morfismes que apunten cap a $A$.

\begin{definicio}\label{def:slice}
Sigui $\cat{C}$ una categoria i $A$ un objecte. La \textbf{categoria de $\cat{C}$ sobre $A$} (en anglès, \emph{slice category}), que escrivim $\cat{C}/A$, té:
\begin{itemize}
\item objectes: $\Ob(\cat{C}/A)$ és la col·lecció dels morfismes $x\colon X\to A$ de $\cat{C}$ amb codomini $A$;
\item morfismes: $\Hom_{\cat{C}/A}(x,y)$, per a $x\colon X\to A$ i $y\colon Y\to A$, és el conjunt dels morfismes $h\colon X\to Y$ de $\cat{C}$ que fan commutar el triangle
\[
\begin{tikzpicture}[baseline=(m.center)]
  \matrix (m) [matrix of math nodes, row sep=2.4em, column sep=2.4em] {
    X & & Y \\
      & A & \\
  };
  \draw[-{Stealth[scale=1.0]}] (m-1-1) -- node[above] {$h$} (m-1-3);
  \draw[-{Stealth[scale=1.0]}] (m-1-1) -- node[below left] {$x$} (m-2-2);
  \draw[-{Stealth[scale=1.0]}] (m-1-3) -- node[below right] {$y$} (m-2-2);
\end{tikzpicture}
\qquad\text{és a dir, } y\comp h=x;
\]
\item composició: la de $\cat{C}$; si $h\in\Hom_{\cat{C}/A}(x,y)$ i $k\in\Hom_{\cat{C}/A}(y,z)$, aleshores $k\comp h\in\Hom_{\cat{C}/A}(x,z)$, perquè $z\comp(k\comp h)=(z\comp k)\comp h=y\comp h=x$;
\item identitats: la identitat de $x\colon X\to A$ és $\id_X$, que fa commutar el triangle perquè $x\comp\id_X=x$.
\end{itemize}
Els axiomes s'hereten de $\cat{C}$.
\end{definicio}

Anàlogament es defineix la \textbf{categoria de $\cat{C}$ sota $A$} (en anglès, \emph{coslice category}), que escrivim $A/\cat{C}$:
\begin{itemize}
\item objectes: $\Ob(A/\cat{C})$ és la col·lecció dels morfismes $x\colon A\to X$ de $\cat{C}$ amb domini $A$;
\item morfismes: $\Hom_{A/\cat{C}}(x,y)$, per a $x\colon A\to X$ i $y\colon A\to Y$, és el conjunt dels morfismes $h\colon X\to Y$ de $\cat{C}$ amb $h\comp x=y$;
\item composició i identitats: les de $\cat{C}$, com a $\cat{C}/A$.
\end{itemize}
Un morfisme de $A/\cat{C}$ és, doncs, un $h$ que fa commutar el triangle sota $A$:
\[
\begin{tikzpicture}[baseline=(m.center)]
  \matrix (m) [matrix of math nodes, row sep=2.4em, column sep=2.4em] {
      & A & \\
    X & & Y \\
  };
  \draw[-{Stealth[scale=1.0]}] (m-1-2) -- node[above left] {$x$} (m-2-1);
  \draw[-{Stealth[scale=1.0]}] (m-1-2) -- node[above right] {$y$} (m-2-3);
  \draw[-{Stealth[scale=1.0]}] (m-2-1) -- node[below] {$h$} (m-2-3);
\end{tikzpicture}
\]
Aquí també convé fer explícits els extrems, com a l'observació~\ref{obs:extrems-implicits}. Un morfisme de $\cat{C}/A$ de $x$ a $y$ no és el morfisme $h\colon X\to Y$ de $\cat{C}$ tot sol, sinó $h$ \emph{junt amb} els seus extrems $x$ i $y$ a $\cat{C}/A$: la terna $(x,y,h)$. La raó és la mateixa: un mateix $h\colon X\to Y$ de $\cat{C}$ pot fer commutar el triangle per a objectes distints de $\cat{C}/A$ ---per exemple, si $X$ és el domini de dos objectes diferents $x,x'\colon X\to A$---, i sense registrar els extrems tindria dominis i codominis indeterminats a $\cat{C}/A$. Escriurem $h$ quan els extrems se sobreentenguin. Anàlogament per a $A/\cat{C}$.

\begin{exemple}
A $\cat{Set}/A$, un objecte és una aplicació $x\colon X\to A$. Les fibres $x^{-1}(a)$, per a $a\in A$, permeten interpretar aquest objecte com una família de conjunts $\bigl(x^{-1}(a)\bigr)_{a\in A}$ indexada per $A$. Recíprocament, una família de conjunts indexada per $A$ dona, mitjançant la unió disjunta, una aplicació cap a $A$. Aquestes dues construccions no identifiquen literalment les dades ---la unió disjunta demana triar-ne una representació---, però sí que estableixen la correspondència adequada entre totes dues descripcions. La farem precisa als capítols~\ref{cap:functors} i~\ref{cap:transformacions}, quan disposem del llenguatge dels functors i de les equivalències de categories.
\end{exemple}

\begin{exemple}\index{conjunt apuntat}
Prenem $A=\{\ast\}$, un conjunt d'un sol element. Un objecte de $\{\ast\}/\cat{Set}$ és una aplicació $\{\ast\}\to X$, és a dir, l'elecció d'un element distingit de $X$, el \textbf{punt base}. Un morfisme és una aplicació que respecta el punt base. Així, $\{\ast\}/\cat{Set}$ és la categoria dels \textbf{conjunts apuntats} (en anglès, \emph{pointed sets}).
\end{exemple}

\section{Principi de dualitat}\label{sec:dualitat}\index{dualitat!principi de}

La dualitat no és només una llista de parelles de conceptes. És un principi que transforma definicions, enunciats i demostracions, i la categoria oposada permet formular-lo amb precisió. Aquesta secció en fa explícit el mecanisme; el retrobarem en tractar objectes inicials i terminals, productes i coproductes, productes fibrats i coproductes fibrats (en anglès, \emph{pullbacks} i \emph{pushouts}), i límits i colímits.

\subsection{Dualitzar un enunciat}\index{dualitat!d'un enunciat}

El principi s'aplica a afirmacions escrites \emph{exclusivament} en el llenguatge de les categories: hi poden intervenir objectes i morfismes, el domini i el codomini de cada morfisme, identitats, composicions i igualtats entre composicions, combinats amb connectors lògics i amb quantificadors sobre objectes i morfismes. Les definicions de monomorfisme, d'isomorfisme o de secció són d'aquesta mena. En canvi, \enquote{$f$ és injectiva} no ho és, perquè parla dels elements d'un conjunt.

Donada una afirmació $\varphi$ d'aquest tipus, n'obtenim la \textbf{dual} $\varphi\op$ \emph{invertint totes les fletxes}:
\begin{itemize}
  \item cada morfisme $f\colon A\to B$ s'inverteix: a la dual es llegeix com una fletxa $B\to A$; per tant, s'intercanvien domini i codomini;
  \item cada composició $g\comp f$ passa a ser $f\comp g$, és a dir, s'inverteix l'ordre de composició;
  \item les identitats, les igualtats, els connectors i els quantificadors es mantenen.
\end{itemize}
Dualitzar dues vegades retorna l'afirmació de partida: $(\varphi\op)\op=\varphi$.

La dual té una lectura immediata: $\varphi\op$ diu de $\cat{C}$ exactament el que $\varphi$ diu de $\cat{C}\op$. En efecte, un morfisme $A\to B$ de $\cat{C}\op$ és un morfisme $B\to A$ de $\cat{C}$, les identitats són les mateixes i la composició de $\cat{C}\op$ és la de $\cat{C}$ en l'ordre invers, $g\op\comp f\op=(f\comp g)\op$. Per tant, per a tota categoria $\cat{C}$,
\[
\varphi \text{ es compleix a } \cat{C}\op
\quad\Longleftrightarrow\quad
\varphi\op \text{ es compleix a } \cat{C}.
\]
D'aquí surt el principi.

\begin{teorema}[Principi de dualitat]\label{teo:dualitat}\index{dualitat!principi de}
Sigui $\varphi$ una afirmació expressable en el llenguatge de les categories. Si $\varphi$ és vàlida en totes les categories, la seva dual $\varphi\op$ també és vàlida en totes les categories.
\end{teorema}

\begin{prova}
Sigui $\cat{C}$ una categoria qualsevol. Com que $\cat{C}\op$ també és una categoria, per hipòtesi $\varphi$ es compleix a $\cat{C}\op$; per l'equivalència anterior, $\varphi\op$ es compleix a $\cat{C}$.
\end{prova}

La manera habitual d'usar el principi és aquesta: s'aplica un teorema conegut a $\cat{C}\op$ i se'n tradueix la conclusió al llenguatge de $\cat{C}$. No cal repetir la demostració; \textbf{la demostració dual és la demostració original amb totes les fletxes invertides}.

Aquesta formulació operativa és la que farem servir al llarg del llibre. Se'n pot donar també una versió completament formal: fixar un llenguatge de primer ordre per a la teoria de categories, definir la dualització com una transformació de fórmules i comprovar que els axiomes en queden invariants. Aquesta formalització, amb una versió sintàctica (sobre deduccions) i una de semàntica (sobre models), és a l'apèndix~\ref{cap:axiomatica-relacional} (secció~\ref{sec:ar-dualitat}); no és necessària per seguir la resta del llibre.

\subsection{Primer exemple: de monomorfisme a epimorfisme}

Recordem (subsecció~\ref{subsec:mono-epi}) que $m\colon A\to B$ és un \textbf{monomorfisme} si, per a tot parell $u,v\colon X\to A$,
\[
\begin{tikzpicture}[baseline=(m.center)]
  \matrix (m) [matrix of math nodes, column sep=3.2em] { X & A & B \\ };
  \draw[-{Stealth[scale=1.0]}] ([yshift=2.6pt]m-1-1.east) -- node[above] {$u$} ([yshift=2.6pt]m-1-2.west);
  \draw[-{Stealth[scale=1.0]}] ([yshift=-2.6pt]m-1-1.east) -- node[below] {$v$} ([yshift=-2.6pt]m-1-2.west);
  \draw[-{Stealth[scale=1.1]}] (m-1-2) -- node[above] {$m$} (m-1-3);
\end{tikzpicture}
\qquad
m\comp u = m\comp v \implies u=v.
\]
Formem-ne el dual: intercanviem domini i codomini i invertim l'ordre de cada composició. L'enunciat dual diu que, per a tot parell $u,v\colon A\to X$,
\[
\begin{tikzpicture}[baseline=(m.center)]
  \matrix (m) [matrix of math nodes, column sep=3.2em] { B & A & X \\ };
  \draw[-{Stealth[scale=1.1]}] (m-1-1) -- node[above] {$m$} (m-1-2);
  \draw[-{Stealth[scale=1.0]}] ([yshift=2.6pt]m-1-2.east) -- node[above] {$u$} ([yshift=2.6pt]m-1-3.west);
  \draw[-{Stealth[scale=1.0]}] ([yshift=-2.6pt]m-1-2.east) -- node[below] {$v$} ([yshift=-2.6pt]m-1-3.west);
\end{tikzpicture}
\qquad
u\comp m = v\comp m \implies u=v,
\]
que és exactament la definició d'\textbf{epimorfisme}. El dual de \enquote{ser un monomorfisme} és, doncs, \enquote{ser un epimorfisme}.

Ara podem dualitzar teoremes sense repetir-ne la demostració. Per exemple, en qualsevol categoria, si $A\xrightarrow{f}B\xrightarrow{g}C$ són monomorfismes, la composició $g\comp f$ també ho és. Pel principi de dualitat (teorema~\ref{teo:dualitat}), l'enunciat dual val igualment:
\begin{quote}
en qualsevol categoria, la composició de dos epimorfismes és un epimorfisme.
\end{quote}
No n'hem hagut de fer cap demostració nova. L'exercici~\ref{exr:pr-dual-mono-epi} de la Pràctica fa aquest pas en detall i assenyala exactament què s'inverteix en passar a $\cat{C}\op$.

\subsection{Segon exemple: de secció a retracció}

Recordem (subsecció~\ref{subsec:seccions-retraccions}) que $s\colon B\to A$ és una \textbf{secció} de $f\colon A\to B$ quan $f\comp s=\id_B$. Dualitzem aquest enunciat: $f$ passa a ser una fletxa $B\to A$, $s$ una fletxa $A\to B$, i la igualtat $f\comp s=\id_B$ es converteix en $s\comp f=\id_B$. En diagrames, invertim les tres fletxes del triangle:
\[
\begin{tikzpicture}[baseline=(m.center)]
  \matrix (m) [matrix of math nodes, row sep=2.6em, column sep=2.6em] {
    B & A \\
    B & \\
  };
  \draw[-{Stealth[scale=1.1]}] (m-1-1) -- node[above] {$s$} (m-1-2);
  \draw[-{Stealth[scale=1.1]}] (m-1-2) -- node[right] {$f$} (m-2-1);
  \draw[-{Stealth[scale=1.1]}] (m-1-1) -- node[left] {$\id_B$} (m-2-1);
\end{tikzpicture}
\quad f\comp s=\id_B
\qquad\qquad
\begin{tikzpicture}[baseline=(m.center)]
  \matrix (m) [matrix of math nodes, row sep=2.6em, column sep=2.6em] {
    B & A \\
    B & \\
  };
  \draw[-{Stealth[scale=1.1]}] (m-1-2) -- node[above] {$s$} (m-1-1);
  \draw[-{Stealth[scale=1.1]}] (m-2-1) -- node[right] {$f$} (m-1-2);
  \draw[-{Stealth[scale=1.1]}] (m-2-1) -- node[left] {$\id_B$} (m-1-1);
\end{tikzpicture}
\quad s\comp f=\id_B
\]
La igualtat de la dreta diu exactament que $s$ és una \textbf{retracció} de $f$. El dual de \enquote{tenir una secció} és, doncs, \enquote{tenir una retracció}, i el d'epimorfisme escindit és monomorfisme escindit.

Amb aquest diccionari es poden dualitzar també els resultats sobre seccions i retraccions; els exercicis proposats del capítol ho apliquen a la proposició~\ref{prop:escindit-mono-epi}.

\subsection{Tercer exemple: del producte al coproducte}\label{subsec:producte-coproducte}

A tall d'avançament ---els definirem en rigor al capítol de límits---, considerem el producte i el coproducte de dos objectes, que il·lustren la dualització d'un enunciat amb un quantificador d'existència única.

Un \emph{producte} de $A$ i $B$ és un objecte $A\times B$ amb dos morfismes $\pi_1\colon A\times B\to A$ i $\pi_2\colon A\times B\to B$ tals que, per a tot $X$ i tot parell $f\colon X\to A$, $g\colon X\to B$, existeix un únic $u\colon X\to A\times B$ amb $\pi_1\comp u=f$ i $\pi_2\comp u=g$:
\[
\begin{tikzpicture}[baseline=(m.center)]
  \matrix (m) [matrix of math nodes, row sep=2.6em, column sep=2.8em] {
     & X & \\
    A & A\times B & B \\
  };
  \draw[-{Stealth[scale=1.0]}] (m-1-2) -- node[above left] {$f$} (m-2-1);
  \draw[-{Stealth[scale=1.0]}] (m-1-2) -- node[above right] {$g$} (m-2-3);
  \draw[-{Stealth[scale=1.0]},dashed] (m-1-2) -- node[right] {$u$} (m-2-2);
  \draw[-{Stealth[scale=1.0]}] (m-2-2) -- node[below] {$\pi_1$} (m-2-1);
  \draw[-{Stealth[scale=1.0]}] (m-2-2) -- node[below] {$\pi_2$} (m-2-3);
\end{tikzpicture}
\]
En el diagrama seguim una convenció que farem servir cada cop que dibuixem una propietat universal: la \textbf{fletxa discontínua} marca el morfisme l'existència del qual afirma la propietat ---aquí el mediador $u$, un cop donades la resta de dades: l'objecte $X$ i les fletxes $f$ i $g$---, mentre que les fletxes contínues són les dades de partida. Igualment, escriurem $\exists!$ per abreujar \enquote{existeix un únic}, com acabem de fer en enunciar la propietat.

Invertint totes les fletxes obtenim un objecte $A+B$ amb dos morfismes $\iota_1\colon A\to A+B$ i $\iota_2\colon B\to A+B$ tals que, per a tot $X$ i tot parell $f\colon A\to X$, $g\colon B\to X$, existeix un únic $u\colon A+B\to X$ amb $u\comp\iota_1=f$ i $u\comp\iota_2=g$:
\[
\begin{tikzpicture}[baseline=(m.center)]
  \matrix (m) [matrix of math nodes, row sep=2.6em, column sep=2.8em] {
    A & A+B & B \\
     & X & \\
  };
  \draw[-{Stealth[scale=1.0]}] (m-1-1) -- node[above] {$\iota_1$} (m-1-2);
  \draw[-{Stealth[scale=1.0]}] (m-1-3) -- node[above] {$\iota_2$} (m-1-2);
  \draw[-{Stealth[scale=1.0]}] (m-1-1) -- node[below left] {$f$} (m-2-2);
  \draw[-{Stealth[scale=1.0]}] (m-1-3) -- node[below right] {$g$} (m-2-2);
  \draw[-{Stealth[scale=1.0]},dashed] (m-1-2) -- node[right] {$u$} (m-2-2);
\end{tikzpicture}
\]
Aquesta és la definició del \emph{coproducte} $A+B$. Els quantificadors $\forall$ i $\exists!$ i la conjunció de les dues equacions es mantenen; el que s'inverteix és el domini i el codomini de cada morfisme i l'ordre de cada composició.

\paragraph{Un cas que cal vigilar.} La construcció \enquote{categoria sobre $A$} es dualitza en la construcció \enquote{categoria sota $A$} (definició~\ref{def:slice}). Ara bé, en dualitzar aquesta construcció cal dualitzar també la categoria ambient. La fórmula precisa és
\[
(\cat{C}/A)\op\;\cong\;A/(\cat{C}\op),
\qquad
(A/\cat{C})\op\;\cong\;(\cat{C}\op)/A.
\]
No s'ha de llegir, doncs, que $(\cat{C}/A)\op\cong A/\cat{C}$ mantenint $\cat{C}$ intacta: això és fals en general. Per exemple, sigui $\cat{C}=\cat{2}=(0\to 1)$ i prenguem com a $A$ l'\emph{objecte} $1$ de $\cat{2}$ (no la categoria $\cat{1}$). La categoria $\cat{2}/1$ té dos objectes (les dues fletxes amb codomini $1$: la fletxa $0\to 1$ i $\id_1$), mentre que $1/\cat{2}$ en té un de sol ($\id_1$, l'única fletxa amb domini $1$); per tant $(\cat{2}/1)\op$ \emph{no} és isomorfa a $1/\cat{2}$. El que sí que val és $(\cat{2}/1)\op\cong 1/(\cat{2}\op)$: a $\cat{2}\op$ hi ha dues fletxes amb domini $1$, la identitat i la fletxa $1\to 0$. En tots els casos, la mateixa regla ---invertir les fletxes--- transforma cada noció en la seva dual, però la regla actua sobre \emph{tot} l'enunciat, ambient inclòs.

\subsection{Diccionari bàsic de dualitats}

Recollim les parelles duals. Les primeres ja han aparegut en aquest capítol; la resta les trobarem més endavant, i la taula creixerà a mesura que apareguin nous conceptes.
\begin{center}
\begin{tabular}{ll}
\hline
\textbf{Concepte} & \textbf{Concepte dual}\\
\hline
\multicolumn{2}{l}{\emph{Ja vistos en aquest capítol}}\\
monomorfisme & epimorfisme\\
secció & retracció\\
producte & coproducte\\
categoria sobre $A$ \ ($\cat{C}/A$) & categoria sota $A$ \ ($A/\cat{C}$)\\
\hline
\multicolumn{2}{l}{\emph{Més endavant}}\\
objecte inicial & objecte terminal\\
equalitzador & coequalitzador\\
producte fibrat & coproducte fibrat\\
con & cocon\\
límit & colímit\\
\hline
\end{tabular}
\end{center}
Algunes nocions són \textbf{autoduals}\index{dualitat!noció autodual}: la identitat, l'isomorfisme i la commutativitat d'un diagrama continuen sent nocions del mateix tipus després de dualitzar-les.

La taula recull les parelles duals com a \emph{nocions}: cada concepte i el seu dual. Quan aquestes nocions es realitzen en categories concretes, però, la dualitat s'ha d'aplicar a tot l'enunciat, ambient inclòs. La parella \enquote{sobre $A$ / sota $A$} n'és l'exemple a vigilar: com hem vist, l'enunciat correcte és $(\cat{C}/A)\op\cong A/(\cat{C}\op)$, no $(\cat{C}/A)\op\cong A/\cat{C}$ amb $\cat{C}$ intacta.

\subsection{Abast i precaucions}

El principi només s'aplica a afirmacions expressables en el llenguatge de les categories. Una afirmació que faci servir informació externa ---elements d'un conjunt, oberts d'un espai o una construcció concreta--- no té dual fins que no es reformula exclusivament amb objectes i morfismes. I quan l'afirmació es refereix a construccions fetes a partir d'una categoria, cal dualitzar també la categoria ambient, com acabem de veure amb les categories sobre i sota un objecte.

El principi de dualitat no permet concloure, per si sol, que una categoria concreta sigui equivalent a la seva oposada. El principi relaciona afirmacions sobre $\cat{C}$ amb les afirmacions duals sobre $\cat{C}\op$; no proporciona en general cap equivalència $\cat{C}\simeq\cat{C}\op$. Algunes categories sí que ho són ---per exemple, $\cat{2}=(0\to 1)$ és isomorfa a $\cat{2}\op$ intercanviant-ne els dos objectes---, però això s'ha de comprovar en cada cas. En canvi, $\cat{Set}$ no s'identifica amb $\cat{Set}\op$, encara que els teoremes generals sobre productes tinguin teoremes duals sobre coproductes.

Finalment, una precaució sobre la lògica. Els connectors i els quantificadors amb què escrivim una afirmació $\varphi$ pertanyen al llenguatge amb què \emph{parlem} de la categoria, i la dualitat els deixa intactes. Una altra cosa és el que passa quan els objectes de la categoria són ells mateixos fórmules, com a la categoria $\cat{Prov}_T$ de la subsecció~\ref{subsec:elementals}. Allà, \emph{quan existeixen}, la conjunció és un producte, la disjunció un coproducte, el cert un objecte terminal i el fals un objecte inicial, i la dualitat de fletxes els intercanvia. Però un preordre de deduïbilitat qualsevol no garanteix que aquests connectors existeixin: en el preordre discret de dues fórmules sense cap deducció entre elles, aquestes dues fórmules no tenen ni producte ni coproducte. I que la dualitat de fletxes coincideixi exactament amb la dualitat de De~Morgan exigeix, a més, un complement involutiu com el de les àlgebres booleanes. Tornarem sobre aquesta lògica interna més endavant, amb el llenguatge adequat.

\section{Propietats duals dels morfismes}\label{sec:propietats-duals}

El principi de dualitat estalvia feina: n'hi ha prou de demostrar les propietats dels monomorfismes; les dels epimorfismes se n'obtenen invertint les fletxes, sense cap demostració nova. Ho il·lustrem amb les propietats bàsiques.

\subsection{Propietats dels monomorfismes}

\begin{proposicio}\label{prop:mono-comp}
La composició de dos monomorfismes és un monomorfisme: si $f\colon A\to B$ i $g\colon B\to C$ són monomorfismes, aleshores $g\comp f$ ho és.
\end{proposicio}

\begin{prova}
Siguin $u,v\colon X\to A$ amb $(g\comp f)\comp u=(g\comp f)\comp v$, és a dir $g\comp(f\comp u)=g\comp(f\comp v)$. Com que $g$ és mono, $f\comp u=f\comp v$; i com que $f$ és mono, $u=v$.
\[
\begin{tikzpicture}[baseline=(m.center)]
  \matrix (m) [matrix of math nodes, column sep=3em] { X & A & B & C \\ };
  \draw[-{Stealth[scale=1.0]}] ([yshift=2.6pt]m-1-1.east) -- node[above] {$u$} ([yshift=2.6pt]m-1-2.west);
  \draw[-{Stealth[scale=1.0]}] ([yshift=-2.6pt]m-1-1.east) -- node[below] {$v$} ([yshift=-2.6pt]m-1-2.west);
  \draw[-{Stealth[scale=1.1]}] (m-1-2) -- node[above] {$f$} (m-1-3);
  \draw[-{Stealth[scale=1.1]}] (m-1-3) -- node[above] {$g$} (m-1-4);
\end{tikzpicture}
\]
\end{prova}

\begin{proposicio}\label{prop:mono-factor}
Si la composició $g\comp f$ és un monomorfisme, aleshores $f$ és un monomorfisme.
\end{proposicio}

\begin{prova}
Siguin $u,v\colon X\to A$ amb $f\comp u=f\comp v$. Component amb $g$, $(g\comp f)\comp u=(g\comp f)\comp v$; com que $g\comp f$ és mono, $u=v$. (No cal cap hipòtesi sobre $g$.)
\end{prova}

\begin{proposicio}\label{prop:iso-retr-mono}
Tot monomorfisme escindit és un monomorfisme, i tot isomorfisme és un monomorfisme escindit. Per tant,
\[
\text{isomorfisme} \;\Longrightarrow\; \text{monomorfisme escindit} \;\Longrightarrow\; \text{monomorfisme}.
\]
\end{proposicio}

\begin{prova}
La primera implicació és la proposició~\ref{prop:escindit-mono-epi}: si $f$ té una retracció, aleshores $f$ és un monomorfisme. Per a la segona, si $f$ és un isomorfisme, el seu invers $f^{-1}$ satisfà $f^{-1}\comp f=\id_A$, de manera que $f^{-1}$ és una retracció de $f$; per tant $f$ té retracció i és un monomorfisme escindit. La cadena d'implicacions se'n segueix.
\end{prova}

\subsection{Les propietats duals dels epimorfismes}

Aplicant el principi de dualitat a les proposicions anteriors ---invertint les fletxes--- obtenim, sense cap demostració nova, les propietats corresponents dels epimorfismes.

\begin{proposicio}[dual de la proposició~\ref{prop:mono-comp}]\label{prop:epi-comp}
La composició de dos epimorfismes és un epimorfisme.
\end{proposicio}

\begin{proposicio}[dual de la proposició~\ref{prop:mono-factor}]\label{prop:epi-factor}
Si la composició $g\comp f$ és un epimorfisme, aleshores $g$ és un epimorfisme.
\end{proposicio}

\begin{proposicio}[dual de la proposició~\ref{prop:iso-retr-mono}]\label{prop:iso-sec-epi}
Tot epimorfisme escindit és un epimorfisme, i tot isomorfisme és un epimorfisme escindit. Per tant,
\[
\text{isomorfisme} \;\Longrightarrow\; \text{epimorfisme escindit} \;\Longrightarrow\; \text{epimorfisme}.
\]
\end{proposicio}

Observem l'asimetria que revela la dualitat: a la proposició~\ref{prop:mono-factor}, d'un mono $g\comp f$ hereta el factor \emph{de la dreta}, $f$; a la seva dual~\ref{prop:epi-factor}, d'un epi hereta el factor \emph{de l'esquerra}, $g$.

\subsection{Condicions necessàries i suficients}

Les implicacions anteriors \emph{no} es reverteixen en general, i és aquí on convé anar amb compte. La cadena vàlida a qualsevol categoria és
\[
\text{iso} \;\Longrightarrow\; \text{mono escindit} \;\overset{(1)}{\Longrightarrow}\; \text{mono},
\]
i cap de les dues fletxes es pot invertir en general:
\begin{itemize}
\item[(1)] \textbf{mono $\not\Rightarrow$ mono escindit.} A $\cat{Grp}$, l'homomorfisme $\mathbb{Z}\xrightarrow{\,\cdot 2\,}\mathbb{Z}$ és un monomorfisme, però no té cap retracció: no és escindit.
\item[(2)] \textbf{mono escindit $\not\Rightarrow$ iso.} Un monomorfisme escindit no ha de ser exhaustiu en cap sentit: a $\cat{Set}$, la injecció $\{0\}\hookrightarrow\{0,1\}$ té retracció (l'única aplicació $\{0,1\}\to\{0\}$), és per tant un monomorfisme escindit, però no és un isomorfisme.
\end{itemize}
Dualment, $\text{iso}\Rightarrow\text{epi escindit}\Rightarrow\text{epi}$, amb les mateixes dues fletxes no reversibles. Convé no confondre \enquote{mono escindit $+$ epi} amb \enquote{mono $+$ epi}: un \textbf{bimorfisme} (alhora mono i epi) no és necessàriament un isomorfisme ---l'única fletxa no trivial de $\cat{2}$ ho és sense tenir invers, i a $\cat{Top}$ una bijecció contínua amb inversa no contínua també---, mentre que un morfisme alhora mono escindit i epi sí que ho és, com recull la caracterització següent:
\[
f \text{ és iso} \iff f \text{ és mono escindit i epi} \iff f \text{ és epi escindit i mono}.
\]
Les demostracions d'aquestes equivalències, i la de la cadena d'implicacions anterior, es proposen a les parts de pràctica i d'exercicis.

\section{Grafs i categories lliures}\label{sec:grafs-lliures}\index{categoria!lliure}

Els grafs i les categories estan estretament relacionats, però no són el mateix. Un graf dirigit només especifica vèrtexs i arestes. Una categoria, a més, disposa d'identitats, permet compondre fletxes i exigeix que aquestes composicions satisfacin els axiomes. De fet, per a una categoria petita, si oblidem la composició i quines fletxes són identitats, els objectes i morfismes formen un graf dirigit: és el seu \emph{graf subjacent}. Les fletxes identitat no desapareixen ---continuen sent arestes del graf, un bucle a cada vèrtex---, només que ja no estan marcades com a identitats. Cal que la categoria sigui petita, perquè els vèrtexs i les arestes d'un graf formen conjunts.

Hi ha una manera natural de fer el camí contrari: generar una categoria a partir d'un graf dirigit prenent com a morfismes tots els camins possibles.

\begin{definicio}\label{def:free}
Sigui $G$ un graf dirigit, amb conjunt de vèrtexs $V_G$, conjunt d'arestes $E_G$ i aplicacions origen i destí $s_G,t_G\colon E_G\to V_G$. Per a $A,B\in V_G$, un \textbf{camí}\index{camí} de $A$ a $B$ de longitud $n\ge 1$ és una successió $(e_1,\dots,e_n)$ d'arestes amb
\[
s_G(e_1)=A,\qquad t_G(e_i)=s_G(e_{i+1})\ \ (1\le i<n),\qquad t_G(e_n)=B.
\]
Per a cada vèrtex $A$ hi ha, a més, el \textbf{camí buit} $()_A$ de $A$ a $A$, de longitud zero.

La \textbf{categoria lliure generada pel graf} $G$, que escrivim $\Free(G)$,\index{Free@$\Free$}\index{categoria!lliure} té:
\begin{itemize}
\item objectes: $\Ob(\Free(G))=V_G$, els vèrtexs de $G$;
\item morfismes: $\Hom_{\Free(G)}(A,B)$ és el conjunt dels camins de $A$ a $B$;
\item composició: la \textbf{concatenació} en ordre de recorregut; per a $p=(a_1,\dots,a_r)\in\Hom(A,B)$ i $q=(b_1,\dots,b_s)\in\Hom(B,C)$,
\[
q\comp p=(a_1,\dots,a_r,b_1,\dots,b_s),
\]
on $p$, que surt de $A$, es recorre primer, d'acord amb la convenció $q\comp p$; concatenar amb un camí buit deixa el camí igual;
\item identitats: $\id_A=()_A$.
\end{itemize}
\end{definicio}

Per exemple, a partir del graf
\[
A\xrightarrow{f}B\xrightarrow{g}C
\]
la categoria lliure conté, a més de les identitats i dels camins $(f)$ i $(g)$ de longitud~1, que escrivim simplement $f$ i $g$, la fletxa composta $g\comp f=(f,g)\colon A\to C$:
\[
\begin{tikzpicture}[baseline=(m.center)]
  \matrix (m) [matrix of math nodes, column sep=3.5em] { A & B & C \\ };
  \draw[-{Stealth[scale=1.05]}] (m-1-1) -- node[above] {$f$} (m-1-2);
  \draw[-{Stealth[scale=1.05]}] (m-1-2) -- node[above] {$g$} (m-1-3);
  \draw[-{Stealth[scale=1.0]}] (m-1-1) to[bend right=25] node[below] {$g\comp f$} (m-1-3);
\end{tikzpicture}
\]
Els axiomes de categoria se satisfan. La composició és associativa, perquè $r\comp(q\comp p)$ i $(r\comp q)\comp p$ són totes dues la successió de totes les arestes de $p$, $q$ i $r$, en aquest ordre; i el camí buit és neutre, perquè concatenar-lo no afegeix cap aresta: $p\comp()_A=p=()_B\comp p$ per a tot camí $p$ de $A$ a $B$. Com que $V_G$ i $E_G$ són conjunts, també ho són els conjunts de camins, i $\Free(G)$ és una categoria petita. L'apèndix~\ref{cap:grafs} desenvolupa aquesta construcció amb més exemples. El nom \emph{lliure} quedarà justificat més endavant, quan disposem del llenguatge dels functors i puguem formular la propietat universal d'aquesta construcció.

\begin{exemple}[El monoide lliure i la convenció de composició]\label{ex:monoide-lliure-cat}
Prenem el cas d'un graf amb \textbf{un sol vèrtex} $\ast$ i conjunt d'arestes $E$. Aleshores totes les arestes són bucles ---comencen i acaben a $\ast$--- i cap parell no deixa d'estar encadenat, de manera que un camí és simplement una seqüència finita d'arestes, és a dir, una \emph{paraula} sobre $E$. Escrivim $E^*$ per al conjunt de totes aquestes paraules, incloent-hi la \emph{paraula buida} $\varepsilon$, de longitud zero.

Fixem tres convencions, que s'han de llegir juntes.
\begin{itemize}
\item \emph{Lectura de les paraules.} Escrivim les lletres en l'\emph{ordre de recorregut}: la paraula $a_1a_2\cdots a_n$ és el camí $(a_1,a_2,\dots,a_n)$, que recorre primer $a_1$, després $a_2$, i així successivament.
\item \emph{Concatenació.} Per a $u=a_1\cdots a_m$ i $v=b_1\cdots b_n$, la concatenació és la juxtaposició $uv=a_1\cdots a_m b_1\cdots b_n$: primer les lletres de $u$ i després les de $v$. Amb aquesta operació i la paraula buida, $(E^*,\text{concatenació},\varepsilon)$ és un monoide, el \textbf{monoide lliure}\index{monoide!lliure} sobre $E$: $(uv)w$ i $u(vw)$ són totes dues la paraula que juxtaposa $u$, $v$ i $w$ en aquest ordre, i $\varepsilon u=u=u\varepsilon$.
\item \emph{Composició.} A $\Free(G)$, la composició és la de la definició: $v\comp u$ vol dir \enquote{primer $u$, després $v$}, i el camí que en resulta recorre primer les lletres de $u$. Per tant,
\[
v\comp u=uv.
\]
\end{itemize}
Amb dues lletres $a,b\in E$, tot cap en una línia:
\[
b\comp a=ab\quad(\text{primer }a,\text{ després }b),
\qquad\qquad
a\comp b=ba\quad(\text{primer }b,\text{ després }a).
\]
L'ordre de la composició és, doncs, el \emph{contrari} del de la juxtaposició: el símbol de la dreta a $b\comp a$ és el que s'escriu primer a la paraula.

Quina categoria hem obtingut? Té un sol objecte, $\Hom_{\Free(G)}(\ast,\ast)=E^*$ i $\id_\ast=\varepsilon$. Però, amb la convenció de $\cat{B}M$, en què $g\comp f=g\cdot f$, la composició $v\comp u=uv$ correspon al producte \emph{invertit}. Si anomenem \textbf{monoide oposat}\index{monoide!oposat} $M\op$ d'un monoide $M$ el que té els mateixos elements i el mateix neutre i l'operació $m\cdot_{\mathrm{op}} n=n\cdot m$ ---de manera que $\cat{B}(M\op)=(\cat{B}M)\op$---, el que tenim és la \emph{igualtat}
\[
\Free(G)=\cat{B}\bigl((E^*)\op\bigr)=\bigl(\cat{B}(E^*)\bigr)\op,
\]
i no $\cat{B}(E^*)$. La diferència és només d'ordre d'escriptura. L'aplicació que inverteix l'ordre de les lletres, $a_1\cdots a_n\mapsto a_n\cdots a_1$, compleix $(uv)^{\mathrm{rev}}=v^{\mathrm{rev}}u^{\mathrm{rev}}$, i és per tant un isomorfisme de monoides $E^*\cong(E^*)\op$. Dona, doncs, un isomorfisme canònic de categories $\cat{B}(E^*)\cong\Free(G)$. Si haguéssim escrit les paraules en l'ordre de la composició ($a_n\cdots a_1$ per al camí que comença per $a_1$), hauríem obtingut literalment $\cat{B}(E^*)$. Hem preferit l'ordre de recorregut, que és el natural per als camins. En aquest sentit, i llevat d'aquest isomorfisme, el monoide lliure és el cas particular de la categoria lliure per a un graf amb un sol vèrtex.
\end{exemple}

\begin{exemple}
La categoria lliure sobre el graf amb dos vèrtexs $A,B$ i una sola aresta $A\to B$ és la categoria finita $\mathbf{2}$. Si, en canvi, el graf té dues arestes en sentits oposats, $e\colon A\to B$ i $f\colon B\to A$, la categoria lliure té infinits morfismes: $e$, $f$, $f\comp e$, $e\comp f$, $e\comp f\comp e$, i així successivament, perquè cap relació no identifica camins diferents. Així, \emph{un graf finit no genera necessàriament una categoria lliure finita}.
\end{exemple}

Aquesta construcció, que genera l'estructura més econòmica a partir d'unes dades sense imposar cap relació que no hi vingui forçada, és un exemple del patró \emph{lliure}, que reapareix per a monoides, grups i moltes altres estructures i que, al capítol~\refcap{cap:adjuncions}{6}, es formalitza com una \emph{adjunció} entre un functor lliure i un functor oblit.

\section{Categories segons la mida}\index{categoria!petita}\index{categoria!localment petita}
\label{sec:mida}

Hem parlat de \enquote{col·lecció} d'objectes en lloc de \enquote{conjunt}, i ara podem precisar per què. En alguns exemples, els objectes són massa nombrosos per formar un conjunt.

\begin{observacio}[La paradoxa de Russell]
La categoria $\cat{Set}$ té per objectes \emph{tots} els conjunts. Ara bé, no pot existir el conjunt de tots els conjunts. En efecte, si $U$ fos aquest conjunt, en podríem separar el subconjunt $R=\{x\in U\mid x\notin x\}$ dels conjunts que no es pertanyen a si mateixos, i preguntar-nos si $R\in R$: si $R\in R$, la definició força $R\notin R$; i si $R\notin R$, la definició força $R\in R$. Totes dues possibilitats són contradictòries. Aquesta és la \textbf{paradoxa de Russell}, i mostra que la col·lecció de tots els conjunts és massa gran per ser un conjunt.
\end{observacio}

Per això, en parlar dels objectes d'una categoria, diem \enquote{col·lecció} i no \enquote{conjunt}: pot ser massa gran. Segons la mida d'aquestes col·leccions introduïm els conceptes següents.

\begin{definicio}
Una categoria és \textbf{petita} si la seva col·lecció d'objectes i la seva col·lecció de morfismes són tots dos conjunts. En cas contrari es diu \textbf{gran}\index{categoria!gran}. I és \textbf{localment petita} si, per a cada parella d'objectes $A,B$, la col·lecció de morfismes $\Hom(A,B)$ és un conjunt, encara que la col·lecció de tots els objectes pugui ser massa gran per ser-ho.
\end{definicio}

Convé no llegir aquestes tres nocions com una partició en tres classes disjuntes. La partició real és \emph{petita} contra \emph{gran}. La petitesa local és una condició transversal: tota categoria petita és localment petita ---si tots els morfismes formen un conjunt, també en formen un els de cada parella $(A,B)$---, i una categoria gran pot ser localment petita o no. Així,
\[
\text{petita}\;\Longrightarrow\;\text{localment petita},
\]
i el recíproc és fals: $\cat{Set}$ és localment petita però no petita. Que no és petita ho mostra la paradoxa de Russell; que és localment petita es comprova identificant cada aplicació amb el seu graf (exercici~\ref{exr:pr-set-localment-petita} de la Pràctica). Les tres possibilitats efectives són, doncs, petita (per tant localment petita), gran localment petita i gran no localment petita.

Moltes de les categories habituals que utilitzem ---$\cat{Set}$, $\cat{Grp}$, $\cat{Top}$, $\cat{Graph}$ i altres--- són localment petites però no petites: tenen massa objectes per formar un conjunt, però entre dos objectes qualssevol hi ha només un conjunt de morfismes. En canvi, cada preordre concret sobre un conjunt i cada categoria finita concreta són categories petites. És important no confondre aquestes categories petites amb les categories grans que les apleguen, com $\cat{Pos}$, la categoria de tots els ordres parcials, o la categoria de totes les categories finites: en tots dos casos hi ha una classe pròpia de conjunts subjacents. Per la mateixa raó és gran $\cat{FinSet}$, tot i que cadascun dels seus objectes és un conjunt finit. Al llarg del llibre suposarem, tret que diguem el contrari, que les categories amb què treballem són localment petites; això garanteix que sempre podem parlar del conjunt $\Hom(A,B)$, cosa que serà essencial més endavant.\footnote{Aquí prenem com a base la teoria de conjunts clàssica, amb classes pròpies; l'apèndix~\ref{cap:mida} fixa aquest marc, i \cite[§1.8]{awodey}, \cite[§1.1]{riehl}, \cite[§3.2]{leinster} i \cite[cap.~1]{perrone} en donen presentacions introductòries. També és possible prendre el mateix llenguatge estructural com a fonament, sense una teoria de conjunts prèvia: l'apèndix~\ref{cap:axiomatica-relacional} n'ofereix un tast, i \cite[cap.~3]{leinster} en fa una introducció.}

\begin{resumcap}
\apartat{Idees centrals}
\begin{enumerate}
  \item Una categoria organitza objectes i morfismes mitjançant identitats i una composició associativa. El que importa no és què són els objectes, sinó com es relacionen a través de les fletxes. Els diagrames commutatius representen igualtats entre composicions.
  \item Una mateixa forma categòrica apareix en contextos molt diferents: conjunts i aplicacions, preordres, monoides, grafs, relacions i sistemes deductius.
  \item Els isomorfismes expressen quan dos objectes tenen la mateixa estructura dins la categoria; sovint, la noció estructuralment rellevant és l'isomorfisme més que no pas la igualtat estricta. Monomorfismes i epimorfismes es defineixen per cancel·lació. A $\cat{Set}$, i en algunes categories concretes, coincideixen amb les aplicacions injectives i exhaustives, però aquesta coincidència no és vàlida en general: en una categoria prima, per exemple, tot morfisme és mono i epi.
  \item Seccions i retraccions són inverses unilaterals. Un morfisme és un monomorfisme escindit si i només si té una retracció, i és un epimorfisme escindit si i només si té una secció. Tot monomorfisme escindit és mono i tot epimorfisme escindit és epi, però les implicacions inverses no valen en general:
  \[
  \text{iso}\Longrightarrow\text{mono escindit}\Longrightarrow\text{mono},
  \qquad
  \text{iso}\Longrightarrow\text{epi escindit}\Longrightarrow\text{epi}.
  \]
  \item A partir d'una categoria en podem construir d'altres: l'oposada, productes de categories, subcategories i categories sobre i sota un objecte.
  \item El \textbf{principi de dualitat} ---invertir totes les fletxes--- converteix sistemàticament cada definició i cada resultat en el seu dual, amb la composició en l'ordre invers.
  \item Un graf dirigit no té, per si sol, composició ni identitats; la categoria lliure hi afegeix formalment tots els camins i les identitats necessàries.
  \item Les qüestions de mida obliguen a distingir categories petites de categories localment petites. Tota categoria petita és localment petita, però moltes categories habituals, com $\cat{Set}$, són localment petites sense ser petites.
\end{enumerate}
\apartat{Errors típics} Creure que \enquote{mono $+$ epi $\Rightarrow$ iso} (fals a $\cat{Top}$ i a $\cat{2}$: un bimorfisme no és sempre un isomorfisme); confondre un monomorfisme amb un monomorfisme escindit; confondre la igualtat d'objectes amb l'isomorfisme, encara que sigui canònic; aplicar els adjectius \enquote{injectiu} i \enquote{exhaustiu} a morfismes d'una categoria general, en lloc de reservar-los per a aplicacions; raonar amb elements en una categoria on els objectes no en tenen o on no separen les fletxes; escriure la composició en l'ordre equivocat en dualitzar; i oblidar la hipòtesi de petitesa local en parlar d'homconjunts.
\apartat{Lectura recomanada} \cite[cap.~1, §2.1 i §3.1]{awodey} per a la introducció paral·lela, els monomorfismes i epimorfismes i el principi de dualitat; \cite[§1.1]{leinster} per a una presentació breu i moderna; \cite[cap.~I i II.1--II.3, II.7]{maclane} per a la formulació clàssica, inclosos la dualitat, les categories producte i les categories lliures; \cite[§1.3]{barrwells} per als grafs, en una presentació orientada a la informàtica; \cite[parts~I--II, sessions~2--10]{lawvereschanuel} per a una introducció encara més elemental als morfismes, els isomorfismes i les seccions i retraccions; i \cite[§1.1.5 i §1.2]{perrone-notes} per a contraexemples elementals d'isomorfismes, monomorfismes i epimorfismes, escindits o no. Per a la presentació relacional de les categories petites, l'apèndix~\ref{cap:axiomatica-relacional}.
\end{resumcap}

\chapter{Functors}
\label{cap:functors}

\begin{objectius}
\apartat{Prerequisits} El capítol~\ref{cap:categories} complet, especialment la definició de categoria, els diagrames commutatius, els isomorfismes, la categoria oposada i la dualitat, les categories producte i les categories sobre i sota un objecte, els grafs i les categories lliures, i la distinció entre categories petites i localment petites. Alguns exemples opcionals utilitzen grups, espais vectorials i l'espai dual $V^*=\Hom_{\mathbb K}(V,\mathbb K)$, i un exemple de functor oblit esmenta els espais topològics; cap d'ells no és imprescindible per seguir el capítol.
\apartat{Objectius} En acabar el capítol, el lector hauria de poder:
\begin{itemize}
  \item definir un functor especificant-ne l'acció sobre objectes i morfismes, i comprovar-ne la functorialitat;
  \item compondre functors i treballar amb la categoria $\cat{Cat}$ de les categories petites;
  \item reconèixer i construir exemples de functors en conjunts, preordres, lògica i grafs, inclosos el graf subjacent i la categoria lliure, i la relació entre $\Free$ i $U$;
  \item distingir functors covariants i contravariants i controlar-ne la variància;
  \item manejar els homfunctors $\Hom(A,-)$ i $\Hom(-,B)$ i el bihomfunctor $\Hom(-,-)$;
  \item distingir functors fidels, plens i essencialment exhaustius;
  \item separar amb precisió les propietats que un functor preserva de les que reflecteix o pot deixar de preservar.
\end{itemize}
\end{objectius}

\section{Cap al concepte de functor}

Un cop fixada la noció de categoria, la pregunta natural és quina mena de correspondència relaciona dues categories. Igual que entre grups mirem els homomorfismes ---les aplicacions que respecten l'estructura de grup--- i entre espais topològics mirem les aplicacions contínues ---les que respecten l'estructura topològica---, entre categories mirarem les correspondències que respecten l'estructura categòrica: objectes, morfismes, identitats i composició. S'anomenen \emph{functors}.

Fixem-nos que una categoria té dos nivells de dades: els objectes i els morfismes. Un functor haurà d'actuar sobre tots dos alhora, i fer-ho de manera compatible: ha d'enviar cada fletxa de la categoria $\cat{C}$ a una fletxa de la categoria $\cat{D}$ que vagi entre els objectes on van a parar el seu domini i el seu codomini, i ha de preservar la manera com les fletxes es componen i quines fletxes són identitats. Aquesta és exactament la idea que fa que els functors siguin, a la vegada, els morfismes d'una categoria de categories.

Abans de formalitzar-ho, val la pena veure un exemple concret que surt d'una operació ben familiar amb conjunts. Treballem a $\cat{Set}$ i fixem un conjunt $S$. A cada conjunt $A$ li podem associar el conjunt de totes les aplicacions de $S$ cap a $A$,
\[
F(A)=A^S=\{\,g\mid g\colon S\to A\,\}.
\]
Això assigna un objecte de $\cat{Set}$ a cada objecte de $\cat{Set}$; però un functor també ha d'actuar sobre les fletxes. I, de fet, tota aplicació $f\colon A\to B$ n'indueix una entre les potències,
\[
F(f)\colon A^S\longrightarrow B^S,
\qquad
F(f)(g)=f\comp g,
\]
és a dir, postcompondre amb $f$: si $g\colon S\to A$ és un element de $A^S$, aleshores $f\comp g\colon S\to B$ és un element de $B^S$, tal com mostra el triangle
\[
\begin{tikzpicture}[baseline=(m.center)]
  \matrix (m) [matrix of math nodes, column sep=3.4em, row sep=2.2em] {
    S & A \\
      & B \\
  };
  \draw[-{Stealth[scale=1.1]}] (m-1-1) -- node[above] {$g$} (m-1-2);
  \draw[-{Stealth[scale=1.1]}] (m-1-2) -- node[right] {$f$} (m-2-2);
  \draw[-{Stealth[scale=1.1]}] (m-1-1) -- node[below left] {$f\comp g$} (m-2-2);
\end{tikzpicture}
\]
Igual que a la definició, $F$ transporta cada fletxa de $\cat{Set}$ a una fletxa entre les imatges dels seus extrems:
\[
\begin{tikzpicture}[baseline=(m.center)]
  \matrix (m) [matrix of math nodes, row sep=2.6em] {
    A \\
    B \\
  };
  \draw[-{Stealth[scale=1.1]}] (m-1-1) -- node[right] {$f$} (m-2-1);
\end{tikzpicture}
\qquad\stackrel{F}{\longmapsto}\qquad
\begin{tikzpicture}[baseline=(m.center)]
  \matrix (m) [matrix of math nodes, row sep=2.6em] {
    A^S \\
    B^S \\
  };
  \draw[-{Stealth[scale=1.1]}] (m-1-1) -- node[right] {$F(f)$} (m-2-1);
\end{tikzpicture}
\]
on cada $g\in A^S$ és una aplicació $g\colon S\to A$ i $F(f)(g)=f\comp g$.
Aquesta assignació compleix just les dues condicions que esperaríem d'un functor. D'una banda, preserva la composició: si $A\xrightarrow{f}B\xrightarrow{h}C$, aleshores, per a tot $g\in A^S$,
\[
\begin{aligned}
F(h\comp f)(g)&=(h\comp f)\comp g\\
&=h\comp(f\comp g)=F(h)\bigl(F(f)(g)\bigr),
\end{aligned}
\]
de manera que
\[
F(h\comp f)=F(h)\comp F(f);
\]
l'única cosa que hem fet servir és l'associativitat de la composició. De l'altra, preserva les identitats: $F(\id_A)(g)=\id_A\comp g=g$, és a dir $F(\id_A)=\id_{A^S}$. Tenim, doncs, una manera sistemàtica de transformar objectes \emph{i} fletxes de $\cat{Set}$ que respecta composició i identitats: exactament el que a la secció següent anomenarem un functor. (Aquest exemple reapareixerà: $A^S$ no és res més que $\Hom_{\cat{Set}}(S,A)$, i el retrobarem com a homfunctor covariant a la definició~\ref{def:homfunctors}.)

\section{Definició de functor}

\begin{definicio}\index{functor}
\label{def:functor}
Siguin $\cat{C}$ i $\cat{D}$ dues categories. Un \textbf{functor} $F\colon\cat{C}\to\cat{D}$ assigna:
\begin{itemize}
\item a cada objecte $A$ de $\cat{C}$, un objecte $F(A)$ de $\cat{D}$;
\item a cada morfisme $f\colon A\to B$ de $\cat{C}$, un morfisme $F(f)\colon F(A)\to F(B)$ de $\cat{D}$,
\end{itemize}
de manera que es compleixin les dues condicions següents:
\begin{enumerate}
\item \textbf{Preservació d'identitats.} Per a cada objecte $A$ de $\cat{C}$,
\[
F(\id_A)=\id_{F(A)}.
\]
\item \textbf{Preservació de la composició.} Per a cada parella de morfismes componibles $f\colon A\to B$ i $g\colon B\to C$ de $\cat{C}$,
\[
F(g\comp f)=F(g)\comp F(f).
\]
\end{enumerate}
\end{definicio}

Observem que la condició sobre morfismes ja imposa que $F$ \emph{respecti} dominis i codominis: si $f\colon A\to B$, aleshores $F(f)$ ha d'anar de $F(A)$ a $F(B)$. La preservació de la composició es pot llegir diagramàticament: un triangle commutatiu de $\cat{C}$ es transforma en un triangle commutatiu de $\cat{D}$.
\[
\begin{tikzpicture}[baseline=(m.center)]
  \matrix (m) [matrix of math nodes, column sep=3.2em, row sep=2.6em] {
    A & B \\
      & C \\
  };
  \draw[-{Stealth[scale=1.1]}] (m-1-1) -- node[above] {$f$} (m-1-2);
  \draw[-{Stealth[scale=1.1]}] (m-1-2) -- node[right] {$g$} (m-2-2);
  \draw[-{Stealth[scale=1.1]}] (m-1-1) -- node[below left] {$g\comp f$} (m-2-2);
\end{tikzpicture}
\qquad\stackrel{F}{\longmapsto}\qquad
\begin{tikzpicture}[baseline=(m.center)]
  \matrix (m) [matrix of math nodes, column sep=3.6em, row sep=2.6em] {
    F(A) & F(B) \\
         & F(C) \\
  };
  \draw[-{Stealth[scale=1.1]}] (m-1-1) -- node[above] {$F(f)$} (m-1-2);
  \draw[-{Stealth[scale=1.1]}] (m-1-2) -- node[right] {$F(g)$} (m-2-2);
  \draw[-{Stealth[scale=1.1]}] (m-1-1) -- node[below left] {$F(g\comp f)$} (m-2-2);
\end{tikzpicture}
\]

\begin{observacio}[Primer els tipus, després les identitats]\label{obs:comprovar-tipus}\index{functor!comprovació de tipus}
Per comprovar que una assignació és un functor convé seguir sempre el mateix ordre, i no passar a les equacions fins que els tipus no quadren.
\begin{enumerate}
\item \emph{Objectes.} Cada $F(A)$ ha de ser un objecte de $\cat{D}$.
\item \emph{Tipus dels morfismes.} Per a cada $f\colon A\to B$, cal que $F(f)$ sigui un morfisme de $\cat{D}$ \emph{i} que vagi de $F(A)$ a $F(B)$. Aquí és on fallen més sovint les construccions candidates.
\item \emph{Identitats i composició.} Només aleshores té sentit comprovar $F(\id_A)=\id_{F(A)}$ i $F(g\comp f)=F(g)\comp F(f)$. El pas~2 garanteix que els dos membres de cada equació tenen el mateix tipus ($F(A)\to F(A)$ i $F(A)\to F(C)$) i que la composició $F(g)\comp F(f)$ existeix, perquè el codomini de $F(f)$ i el domini de $F(g)$ són tots dos $F(B)$.
\end{enumerate}
Un exemple mostra per què l'ordre importa. Si a cada conjunt $X$ li assignem $\mathcal P(X)$ i a cada aplicació $f\colon X\to Y$ l'antiimatge $B\mapsto f^{-1}(B)$, el pas~2 ja falla: l'antiimatge va de $\mathcal P(Y)$ a $\mathcal P(X)$, no de $\mathcal P(X)$ a $\mathcal P(Y)$. No cal, doncs, comprovar cap equació per saber que això no és un functor covariant $\cat{Set}\to\cat{Set}$. Sí que és, en canvi, un functor \emph{contravariant}, la noció que introduirem a la secció~\ref{sec:variancia}, per a la qual el pas~2 demana precisament $F(f)\colon F(B)\to F(A)$ (exemple~\ref{ex:parts-dues-variancies}). Les solucions de la Pràctica segueixen aquest ordre; vegeu, per exemple, l'exercici resolt~\ref{pr-hom-cov}.
\end{observacio}

\begin{observacio}[Els functors preserven els diagrames commutatius]\index{diagrama commutatiu!preservat per functors}
Una conseqüència directa de la preservació de la composició és que un functor envia diagrames commutatius a diagrames commutatius. En efecte, dir que un diagrama commuta és dir que certes composicions de fletxes coincideixen; si dos camins amb el mateix origen i el mateix destí donen la mateixa composició a $\cat{C}$, aleshores, aplicant $F$ i usant que preserva la composició, les seves imatges també coincideixen a $\cat{D}$. Per això podem aplicar un functor a tot un diagrama i seguir raonant amb la seva imatge: les igualtats entre camins es conserven. Ho comprovarem en detall, per al cas d'un quadrat, a l'exercici~\ref{prac:functor-quadrat} del capítol~\ref{cap:practica-functors} de Pràctica.
\end{observacio}

\subsection{Els functors preserven els isomorfismes}\label{sec:preserven-isos}

Com que els functors preserven tota l'estructura categòrica, en particular preserven les equacions entre composicions. Una conseqüència immediata és que preserven els isomorfismes.

\begin{proposicio}\label{prop:functor-preserva-iso}\index{functor!preserva isomorfismes}
Si $F\colon\cat{C}\to\cat{D}$ és un functor i $f\colon A\to B$ és un isomorfisme a $\cat{C}$, aleshores $F(f)$ és un isomorfisme a $\cat{D}$, i
\[
\bigl(F(f)\bigr)^{-1}=F(f^{-1}).
\]
\end{proposicio}

\begin{prova}
Sigui $g=f^{-1}$, de manera que $g\comp f=\id_A$ i $f\comp g=\id_B$. Aplicant $F$ i usant que preserva composició i identitats,
\[
F(g)\comp F(f)=F(g\comp f)=F(\id_A)=\id_{F(A)},
\]
\[
F(f)\comp F(g)=F(f\comp g)=F(\id_B)=\id_{F(B)}.
\]
Per tant $F(g)$ és un invers de $F(f)$, i pel resultat d'unicitat de l'invers, $\bigl(F(f)\bigr)^{-1}=F(g)=F(f^{-1})$.
\end{prova}

\begin{observacio}
Val la pena veure per què és així, més enllà del càlcul de la prova. Dir que $f$ és un isomorfisme és dir que existeix $g$ amb $g\comp f=\id_A$ i $f\comp g=\id_B$: dues equacions entre composicions, és a dir, dos diagrames commutatius. Com que els functors preserven els diagrames commutatius, aquestes dues equacions es transporten a $\cat{D}$ en aplicar-hi $F$, i $F(g)$ resulta ser l'invers de $F(f)$.
\end{observacio}

\section{Exemples de functors}\label{sec:primers-exemples}

La força de la definició, com en el cas de les categories, es veu en la varietat d'exemples que hi encaixen. Recuperem els exemples de referència del capítol anterior.

\begin{exemple}[Functor identitat]\index{functor!identitat}
Per a tota categoria $\cat{C}$ hi ha el \textbf{functor identitat} $\id_{\cat{C}}\colon\cat{C}\to\cat{C}$, que deixa fixos tots els objectes i tots els morfismes. Preserva trivialment identitats i composició.
\end{exemple}

\begin{exemple}[Functor constant]\label{ex:functor-constant}\index{functor!constant}
Siguin $\cat{C}$ i $\cat{D}$ dues categories i $D$ un objecte de $\cat{D}$. El \textbf{functor constant} amb valor $D$,
\[
\Delta_D\colon\cat{C}\to\cat{D},
\]
envia tot objecte $A$ de $\cat{C}$ a $D$ i tot morfisme $f$ de $\cat{C}$ a la identitat $\id_D$. És un functor: $\Delta_D(\id_A)=\id_D=\id_{\Delta_D(A)}$, i $\Delta_D(g\comp f)=\id_D=\id_D\comp\id_D=\Delta_D(g)\comp\Delta_D(f)$. Quan $\cat{D}=\cat{Set}$ i $D$ és un conjunt d'un sol element $\{\ast\}$, escrivim $\Delta_{\{\ast\}}$ per al functor constant que envia cada objecte a aquest conjunt d'un sol element.
\end{exemple}

\begin{exemple}[Functors oblit]\index{functor!oblit}
Moltes categories d'objectes amb estructura $\cat{C}$ tenen un \textbf{functor oblit} $U\colon\cat{C}\to\cat{Set}$, que oblida l'estructura i reté només el conjunt subjacent:
\[
\cat{Grp}\to\cat{Set},
\qquad
\cat{Top}\to\cat{Set},
\qquad
\cat{Vect}_{\mathbb K}\to\cat{Set}.
\]
En cada cas, $U$ envia un objecte al seu conjunt subjacent i un morfisme ---un homomorfisme, una aplicació contínua, una aplicació lineal--- a la mateixa aplicació, vista simplement com a aplicació de conjunts. La identitat es preserva perquè l'aplicació identitat continua sent la identitat, i la composició es preserva perquè és la mateixa composició d'aplicacions. El functor oblit \enquote{no fa res} als morfismes; el seu contingut és que oblida l'estructura dels objectes.
\end{exemple}

\begin{exemple}[Functors des de les categories finites]\index{functor!des d'una categoria finita}
Les categories finites $\cat{0}$, $\cat{1}$ i $\cat{2}$ (subsecció~\ref{subsec:finites}) donen, com a dominis, functors amb una lectura molt directa. Un functor $\cat{1}\to\cat{D}$, on $\cat{1}$ té un únic objecte i cap morfisme no trivial, no és res més que la tria d'un objecte de $\cat{D}$: la imatge de l'únic objecte. Un functor $\cat{2}\to\cat{D}$, on $\cat{2}=(0\to 1)$ ---reanomenem els dos objectes de $\cat{2}$ com $0$ i $1$, i escrivim $(0\to 1)$ per indicar la seva única fletxa no trivial---, és la tria d'un morfisme de $\cat{D}$: cal dir on van $0$ i $1$ ---dos objectes--- i on va l'única fletxa no trivial ---un morfisme entre les seves imatges---, i la preservació d'identitats i composició no imposa res més. Finalment, des de la categoria buida $\cat{0}$ hi ha un únic functor $\cat{0}\to\cat{D}$, el functor buit, sigui quina sigui $\cat{D}$. Aquests tres fets ---objectes, morfismes i \enquote{res}--- reapareixeran contínuament: seleccionar un objecte o un morfisme d'una categoria és el mateix que donar un functor des de $\cat{1}$ o des de $\cat{2}$.
\end{exemple}

\begin{exemple}[Functors entre categories discretes]\index{functor!entre categories discretes}\index{categoria!discreta}
Un conjunt es pot veure com una categoria discreta, en què els únics morfismes són les identitats (subsecció~\ref{subsec:conjunts}). Si $S$ i $T$ són dos conjunts vistos així, un functor $F\colon S\to T$ només ha de dir on va cada objecte, perquè els únics morfismes són les identitats i tot functor les preserva automàticament. Per tant un functor entre categories discretes és exactament una aplicació de conjunts $S\to T$: sense estructura a preservar, la noció de functor es redueix a la d'aplicació.
\end{exemple}

\begin{exemple}[Functors entre preordres]\index{functor!entre preordres}
Siguin $(P,\leq)$ i $(Q,\leq)$ dos preordres, vistos com a categories. Un functor $F\colon P\to Q$ assigna a cada element $x\in P$ un element $F(x)\in Q$ i, com que hi ha un morfisme $x\to y$ exactament quan $x\leq y$, la condició sobre morfismes diu que
\[
x\leq y
\quad\Longrightarrow\quad
F(x)\leq F(y).
\]
És a dir, un functor entre preordres és exactament una aplicació monòtona. Les condicions d'identitat i composició es compleixen automàticament, perquè en un preordre entre dos objectes hi ha com a màxim una fletxa i no hi ha res més a comprovar.
\end{exemple}

\begin{exemple}[Functors des d'un monoide]\label{ex:functors-monoide}\index{functor!des d'un monoide}
Recordem que un monoide $(M,\cdot,1)$ es pot veure com una categoria d'un sol objecte $\ast$, que escrivim $\cat{B}M$. Un functor $F\colon\cat{B}M\to\cat{D}$ tria un objecte $D=F(\ast)$ de $\cat{D}$ i envia cada element $m\in M$ ---és a dir, cada morfisme $\ast\to\ast$--- a un morfisme $F(m)\colon D\to D$, de manera que
\[
F(1)=\id_D
\qquad\text{i}\qquad
F(m\cdot n)=F(m)\comp F(n).
\]
Aquestes són exactament les condicions perquè $F$ sigui un \textbf{homomorfisme de monoides} de $M$ cap al monoide dels endomorfismes de $D$,
\[
\mathrm{End}_{\cat{D}}(D)=\Hom_{\cat{D}}(D,D).
\]
En particular, quan $\cat{D}=\cat{Set}$, un functor $\cat{B}M\to\cat{Set}$ és una \textbf{acció} del monoide $M$ sobre un conjunt. Vegem per què. Un tal functor queda determinat per un conjunt $F(\ast)=S$ i, per a cada element $m\in M$, una funció $F(m)\colon S\to S$. Les dues condicions de functor ---preservar identitats i composició--- diuen exactament que, per a tot $x\in S$ i tots $m,n\in M$,
\[
F(1)(x)=x,
\qquad
F(m\cdot n)(x)=F(m)\bigl(F(n)(x)\bigr),
\]
Si escrivim $m\cdot x$ en lloc de $F(m)(x)$, aquestes igualtats esdevenen les condicions habituals d'una \emph{acció per l'esquerra} del monoide sobre $S$:
\[
1\cdot x=x,
\qquad
(m\cdot n)\cdot x=m\cdot(n\cdot x).
\]
\end{exemple}

\begin{exemple}[Functors des d'un grup: accions i representacions]\index{functor!des d'un grup}\index{accio@acció!de grup}\index{representacio@representació lineal}
Quan el monoide és un grup $G$, la categoria $\cat{B}G$ és un grupoide: tot morfisme és un isomorfisme. Un functor $F\colon\cat{B}G\to\cat{D}$ és aleshores un homomorfisme de $G$ cap al grup dels automorfismes de l'objecte $D=F(\ast)$,
\[
F\colon G\longrightarrow \Aut_{\cat{D}}(D),
\]
perquè un functor preserva isomorfismes (proposició~\ref{prop:functor-preserva-iso}) i, per tant, cada element de $G$ va a parar a un automorfisme de $D$. Dos casos són especialment importants:
\begin{itemize}
\item Un functor $\cat{B}G\to\cat{Set}$ és una \textbf{acció} del grup $G$ sobre un conjunt ---un \emph{$G$-conjunt}. El desplegament és el mateix que per a un monoide (exemple~\ref{ex:functors-monoide}), amb una novetat: com que ara cada $a\in G$ té un invers $a^{-1}$, la funció $F(a)\colon S\to S$ és bijectiva, amb inversa $F(a^{-1})$. Cada element del grup actua, doncs, com una \emph{permutació} de $S$. Un exemple típic són les simetries d'una figura, que actuen permutant-ne els vèrtexs.
\item Un functor $\cat{B}G\to\cat{Vect}_{\mathbb K}$ és una \textbf{representació lineal} de $G$ sobre el cos $\mathbb{K}$. El desplegament torna a ser el mateix, ara amb un espai vectorial $F(\ast)=V$ en lloc d'un conjunt: cada element $a\in G$ dona una aplicació \emph{lineal} invertible $F(a)\colon V\to V$ ---un automorfisme de $V$, amb inversa $F(a^{-1})$---, i el producte del grup es correspon amb la composició. Escrivint $a\cdot v$ per $F(a)(v)$, la novetat respecte d'una acció ordinària és que cada element actua respectant les operacions de l'espai:
\[
a\cdot(\lambda v+\mu w)=\lambda\,(a\cdot v)+\mu\,(a\cdot w).
\]
En dimensió finita, fixada una base, cada $F(a)$ és una matriu invertible: una representació \enquote{realitza} els elements abstractes del grup com a matrius concretes, cosa que permet estudiar el grup amb les eines de l'àlgebra lineal.
\end{itemize}
Així, dues teories aparentment allunyades ---les accions i representacions lineals d'un grup--- són, categòricament, el mateix tipus d'objecte: functors amb domini $\cat{B}G$, que només es distingeixen per la categoria d'arribada.
\end{exemple}

\begin{exemple}[Deduïbilitat i functors lògics]\label{ex:deduibilitat-functors}\index{functor!i deduïbilitat}\index{substitucio@substitució}
Siguin $T$ i $T'$ dos sistemes deductius, i $\cat{Prov}_T$ i $\cat{Prov}_{T'}$ les seves categories primes de deduïbilitat (subsecció~\ref{subsec:elementals}). Una assignació de fórmules $\varphi\colon A\mapsto \varphi(A)$ que preservi la deduïbilitat,
\[
A\vdash_T B
\quad\Longrightarrow\quad
\varphi(A)\vdash_{T'}\varphi(B),
\]
és exactament un functor $F\colon\cat{Prov}_T\to\cat{Prov}_{T'}$, amb $F(A)=\varphi(A)$: la fletxa que expressa $A\vdash_T B$ va a parar a la fletxa que expressa $\varphi(A)\vdash_{T'}\varphi(B)$. Com que aquestes categories són primes, no cal comprovar res més: la preservació de la deduïbilitat garanteix automàticament la preservació d'identitats i composicions. Aquest és el reflex categòric del fet que una traducció entre teories que respecti la deduïbilitat indueix un functor entre les categories associades.

Un cas concret i molt natural d'aquesta situació és el d'una \textbf{substitució} de variables proposicionals. Sigui $T_{\mathrm{prop}}$ el sistema deductiu de la lògica proposicional i $\cat{Prov}_{T_{\mathrm{prop}}}$ la seva categoria de deduïbilitat, amb les fórmules per objectes i una fletxa $A\to B$ exactament quan $A\vdash B$ (aquí escrivim $\vdash$ per $\vdash_{T_{\mathrm{prop}}}$). Considerem la substitució
\[
\sigma\colon\quad p\mapsto q,\qquad q\mapsto p\lor r,
\]
estesa a totes les fórmules respectant els connectors ($\land,\lor,\neg,\to$). Per exemple, $\sigma$ transforma $p\to(q\land p)$ en $q\to((p\lor r)\land q)$. El punt clau és que una substitució no sols transforma fórmules, sinó que també preserva la \emph{deduïbilitat}: si $A\vdash B$, aleshores $\sigma(A)\vdash\sigma(B)$, perquè qualsevol prova concreta de $B$ a partir de $A$ es tradueix, aplicant $\sigma$ a cada fórmula, en una derivació de $\sigma(B)$ a partir de $\sigma(A)$. Així, de la fletxa $p\land q\to p$ (perquè $p\land q\vdash p$) n'obtenim la fletxa $q\land(p\lor r)\to q$. Per tant $\sigma$ indueix un functor
\[
F_\sigma\colon\cat{Prov}_{T_{\mathrm{prop}}}\longrightarrow\cat{Prov}_{T_{\mathrm{prop}}},
\qquad
A\mapsto\sigma(A),
\qquad
(A\vdash B)\mapsto(\sigma(A)\vdash\sigma(B)),
\]
que és, precisament, el cas $T=T'=T_{\mathrm{prop}}$ de la construcció anterior (aquí amb la mateixa categoria de sortida i d'arribada).
\end{exemple}

\begin{observacio}
Si en comptes de la deduïbilitat considerem la \emph{categoria de proves} de $T$ ---els mateixos objectes, però ara amb les derivacions, mòdul equivalència, com a morfismes i el tall com a composició---, un functor cap a la categoria de proves de $T'$ demana també una traducció de cada derivació que en respecti l'encadenament i les identitats. A diferència del cas anterior, aquestes condicions ja no són automàtiques, perquè la categoria de proves no és prima.
\end{observacio}

\begin{exemple}[Functors oblit sobre i sota un objecte]\index{functor!oblit}\index{categoria!sobre un objecte}\index{categoria!sota un objecte}
Les categories de $\cat{C}$ sobre i sota un objecte $A$, $\cat{C}/A$ i $A/\cat{C}$ (definició~\ref{def:slice}), tenen un functor oblit cap a $\cat{C}$, un per a cadascuna. El functor
\[
U_A\colon\cat{C}/A\to\cat{C}
\]
envia cada objecte $x\colon X\to A$ al seu domini $X$ i cada morfisme (triangle) $h$ a ell mateix; oblida el morfisme distingit cap a $A$ i reté només el domini. Dualment, $A/\cat{C}\to\cat{C}$ envia cada objecte $x\colon A\to X$ al seu codomini $X$ i cada morfisme (triangle) a ell mateix; oblida el morfisme distingit des de $A$ i reté només el codomini.

Vegem-ho en un cas concret. Amb $\cat{2}=(0\to 1)$ i l'objecte $A=1$, els objectes de $\cat{2}/1$ són les fletxes de $\cat{2}$ amb codomini $1$: la identitat $\id_1$ i l'única fletxa no trivial $u\colon 0\to 1$. L'únic morfisme no trivial va de $u$ a $\id_1$, de manera que $\cat{2}/1$ és, de nou, una còpia de $\cat{2}$, i $U_1\colon\cat{2}/1\to\cat{2}$ envia $u\mapsto 0$ i $\id_1\mapsto 1$. Simètricament, sota l'objecte $0$, els objectes de $0/\cat{2}$ són les fletxes amb domini $0$ ---la identitat $\id_0$ i $u\colon 0\to 1$---, i el functor oblit $0/\cat{2}\to\cat{2}$ envia $\id_0\mapsto 0$ i $u\mapsto 1$.
\end{exemple}

\begin{exemple}[Projeccions i functor diagonal]\label{ex:projeccions-diagonal}\index{functor!de projecció}\index{functor!diagonal}
La categoria producte $\cat{C}\times\cat{D}$ (subsecció~\ref{subsec:categoria-producte}) té dos functors de \textbf{projecció}
\[
\pi_1\colon\cat{C}\times\cat{D}\to\cat{C},\qquad \pi_2\colon\cat{C}\times\cat{D}\to\cat{D},
\]
que es queden amb una de les components: $\pi_1(A,X)=A$ i $\pi_1(f,g)=f$, i anàlogament per a $\pi_2$. Són functors perquè la composició i les identitats de $\cat{C}\times\cat{D}$ es defineixen component a component:
\[
\begin{aligned}
\pi_1\bigl((f',g')\comp(f,g)\bigr)
 &= \pi_1(f'\comp f,\,g'\comp g)\\
 &= f'\comp f
  = \pi_1(f',g')\comp\pi_1(f,g),\\
\pi_1(\id_A,\id_X)&=\id_A,
\end{aligned}
\]
i el càlcul per a $\pi_2$ és el mateix sobre la segona component.
En sentit contrari, cada categoria $\cat{C}$ té un functor \textbf{diagonal}
\[
\Delta\colon\cat{C}\to\cat{C}\times\cat{C},\qquad \Delta(A)=(A,A),\qquad \Delta(f)=(f,f),
\]
que duplica objectes i morfismes; la functorialitat és de nou immediata, perquè $\Delta(g\comp f)=(g\comp f,\,g\comp f)=(g,g)\comp(f,f)$ i $\Delta(\id_A)=(\id_A,\id_A)$. Les projeccions i la diagonal encaixen en el diagrama commutatiu
\[
\begin{tikzpicture}[baseline=(m.center)]
  \matrix (m) [matrix of math nodes, row sep=2.6em, column sep=3.4em] {
    & \cat{C} & \\
    \cat{C} & \cat{C}\times\cat{C} & \cat{C} \\
  };
  \draw[-{Stealth[scale=1.1]}] (m-1-2) -- node[right] {$\Delta$} (m-2-2);
  \draw[-{Stealth[scale=1.1]}] (m-2-2) -- node[below] {$\pi_1$} (m-2-1);
  \draw[-{Stealth[scale=1.1]}] (m-2-2) -- node[below] {$\pi_2$} (m-2-3);
  \draw[-{Stealth[scale=1.1]}] (m-1-2) -- node[above left] {$\id_{\cat{C}}$} (m-2-1);
  \draw[-{Stealth[scale=1.1]}] (m-1-2) -- node[above right] {$\id_{\cat{C}}$} (m-2-3);
\end{tikzpicture}
\qquad
\pi_1\comp\Delta=\id_{\cat{C}}=\pi_2\comp\Delta,
\]
on la composició de functors es fa com a la secció següent: les dues igualtats es comproven directament, ja que $\pi_1(\Delta(A))=\pi_1(A,A)=A$ i $\pi_1(\Delta(f))=\pi_1(f,f)=f$, i igual per a $\pi_2$. La notació concorda amb la del functor constant $\Delta_D$ de l'exemple~\ref{ex:functor-constant}: més endavant, el functor diagonal s'estendrà a un functor que envia cada objecte $A$ a un diagrama constant de valor $A$, del qual $(A,A)$ és el cas de dues còpies. Són exemples elementals, però reapareixeran sovint: les projeccions fan de $\cat{C}\times\cat{D}$ el producte de $\cat{C}$ i $\cat{D}$ dins de $\cat{Cat}$ quan totes dues són petites (observació~\ref{obs:dos-productes}), i el functor diagonal és l'ingredient bàsic per tractar productes, límits i adjuncions en capítols posteriors.
\end{exemple}

\section{Composició de functors i la categoria \texorpdfstring{$\cat{Cat}$}{Cat}}\index{Cat@$\cat{Cat}$}

Els functors es poden compondre, i la composició torna a ser un functor. Això és el que permet organitzar les categories i els functors en una nova categoria.

\begin{proposicio}
Siguin $F\colon\cat{C}\to\cat{D}$ i $G\colon\cat{D}\to\cat{E}$ dos functors. La seva \textbf{composició} $G\comp F$, definida per
\[
(G\comp F)(A)=G(F(A))
\qquad\text{i}\qquad
(G\comp F)(f)=G(F(f)),
\]
és un functor $\cat{C}\to\cat{E}$.
\end{proposicio}

\begin{prova}
Cal comprovar les dues condicions. Per a les identitats, i usant primer que $F$ i després que $G$ les preserven,
\[
(G\comp F)(\id_A)=G(F(\id_A))=G(\id_{F(A)})=\id_{G(F(A))}=\id_{(G\comp F)(A)}.
\]
Per a la composició, siguin $f\colon A\to B$ i $g\colon B\to C$ a $\cat{C}$. Aleshores
\[
\begin{aligned}
(G\comp F)(g\comp f)
&=G\big(F(g\comp f)\big)
=G\big(F(g)\comp F(f)\big)\\
&=G(F(g))\comp G(F(f))
=(G\comp F)(g)\comp(G\comp F)(f).
\end{aligned}
\]
Per tant $G\comp F$ és un functor.
\end{prova}

La composició de functors és associativa i el functor identitat actua com a element neutre, exactament pels mateixos motius que a $\cat{Set}$. En efecte, un functor $F\colon\cat{C}\to\cat{D}$ consta de dues aplicacions ---una sobre objectes i una sobre morfismes---, i la composició de functors no és res més que la composició d'aquestes aplicacions component a component. Per tant, l'associativitat i les lleis d'identitat es redueixen a les propietats corresponents de la composició d'aplicacions, que ja coneixem. Això ens permet donar la definició següent.

\begin{definicio}\index{Cat@$\cat{Cat}$}
La categoria $\cat{Cat}$ té per objectes les categories petites i per morfismes els functors entre elles. La composició és la composició de functors i la identitat de cada categoria és el seu functor identitat.
\end{definicio}

\begin{observacio}[Qüestions de mida]\label{obs:mida-cat}
La categoria $\cat{Cat}$ té per objectes les categories \emph{petites}. Aquesta restricció la fa \emph{localment petita}: si $\cat{C}$ és petita, $\Ob(\cat{C})$ i $\Mor(\cat{C})$ són conjunts, i doncs, donades dues categories petites $\cat{C}$ i $\cat{D}$, el producte $\Ob(\cat{D})^{\Ob(\cat{C})}\times \Mor(\cat{D})^{\Mor(\cat{C})}$ és un conjunt; com que
\[
\Hom_{\cat{Cat}}(\cat{C},\cat{D})\subseteq \Ob(\cat{D})^{\Ob(\cat{C})}\times \Mor(\cat{D})^{\Mor(\cat{C})},
\]
també ho és $\Hom_{\cat{Cat}}(\cat{C},\cat{D})$. La mateixa $\cat{Cat}$, en canvi, és \emph{gran}: les categories petites formen una classe pròpia, no un conjunt.

Categories d'ús constant com $\cat{Set}$, $\cat{Grp}$ o $\cat{Top}$ són grans i, per tant, no són objectes de $\cat{Cat}$. Això no impedeix treballar amb functors entre elles: un functor oblit $\cat{Grp}\to\cat{Set}$, o el functor lliure $\Free\colon\cat{Graph}\to\cat{Cat}$, és simplement una assignació sobre objectes i una altra sobre morfismes que compleixen les condicions functorials, i la seva legitimitat no depèn de cap categoria que tingui les categories grans per objectes.

Al llarg del text, doncs, $\cat{Cat}$ designarà sempre la categoria de les categories \emph{petites}, i les categories grans, com $\cat{Set}$, no en són objectes.\footnotemark
\end{observacio}
\footnotetext{Aquesta afirmació és relativa al marc de mida adoptat (apèndix~\ref{cap:mida}). Amb una jerarquia d'universos de Grothendieck, una categoria gran respecte d'un univers pot ser petita respecte d'un univers més gran, i aleshores és objecte de la categoria de categories d'aquest nivell; vegeu \cite[§1.1]{riehl}.}%

Un cop els functors són els morfismes de $\cat{Cat}$, podem aplicar a les categories les mateixes nocions categòriques que ja coneixem. En particular, un isomorfisme a $\cat{Cat}$ és un functor que té un invers estricte.

\begin{definicio}[Isomorfisme de categories]\label{def:iso-categories}\index{isomorfisme!de categories}\index{categoria!isomorfa}
Dues categories $\cat{C}$ i $\cat{D}$ són \textbf{isomorfes}, i escrivim $\cat{C}\cong\cat{D}$, si hi ha functors $F\colon\cat{C}\to\cat{D}$ i $G\colon\cat{D}\to\cat{C}$ inversos l'un de l'altre, és a dir, amb $G\comp F=\id_{\cat{C}}$ i $F\comp G=\id_{\cat{D}}$; un tal $F$ és un \textbf{isomorfisme de categories}. Per a categories petites no és res més que un isomorfisme (definició~\ref{def:isomorfisme}) a $\cat{Cat}$, però la definició val igualment per a categories grans, que no són objectes de $\cat{Cat}$.
\end{definicio}

\section{Functors i grafs}\label{sec:functors-grafs}\index{functor!i grafs}\index{morfisme de grafs}

Val la pena comparar els functors amb els \emph{morfismes de grafs}, que vam introduir en definir $\cat{Graph}$ (subsecció~\ref{subsec:grafs}). Un morfisme de grafs $\alpha\colon G\to H$ és una parella d'aplicacions $\alpha_V\colon V_G\to V_H$ i $\alpha_E\colon E_G\to E_H$ compatibles amb les dues aplicacions d'origen i destí $s,t\colon E\to V$ de cada graf; res més, perquè un graf no té composició ni identitats. Un functor entre categories petites demana exactament aquestes mateixes dues dades ---una assignació sobre objectes i una sobre morfismes, compatibles amb domini i codomini--- però hi afegeix dues condicions que un graf no pot ni formular: la preservació de la composició i de les identitats. Un morfisme de grafs és, doncs, la part «sense composició ni identitats» d'un functor. Aquesta relació es concreta en dues construccions en sentits contraris entre $\cat{Graph}$ i $\cat{Cat}$, que estudiem tot seguit.

\subsection{El graf subjacent d'una categoria}

Tota categoria petita $\cat{C}$ té un \textbf{graf subjacent}: el graf que té per vèrtexs els objectes de $\cat{C}$ i per arestes \emph{tots} els seus morfismes, amb les aplicacions origen i destí donades pel domini i el codomini. En passar a aquest graf, oblidem quins morfismes són identitats i com es componen, però no eliminem les arestes corresponents: les fletxes identitat hi resten com a arestes, una a cada vèrtex, tot i que ja no marcades com a identitats.

Aquesta assignació és functorial. Donat un functor $F\colon\cat{C}\to\cat{D}$, obtenim un morfisme entre els grafs subjacents que actua, sobre vèrtexs, com $F$ sobre els objectes, i sobre arestes, com $F$ sobre els morfismes. És un morfisme de grafs ben definit precisament perquè $F$ respecta domini i codomini: una aresta $f\colon A\to B$ va a $F(f)\colon F(A)\to F(B)$, una aresta entre les imatges dels extrems. Per construir-lo n'hi ha prou amb això; la preservació de la composició i de les identitats no s'hi fa servir, i per això diem que «se n'oblida». Tenim, així, un functor oblit
\[
U\colon\cat{Cat}\to\cat{Graph},
\]
que envia cada categoria petita al seu graf subjacent $U(\cat{C})$ i cada functor $F$ al morfisme de grafs $U(F)$ que acabem de descriure. El recíproc del que fa $U$ és fals: un morfisme entre els grafs subjacents no respecta, en general, ni la composició ni les identitats, de manera que un functor és estrictament més que un morfisme de grafs. Vegem-ho amb un bucle. Sigui $M=(\{1,e\},\cdot,1)$ amb $e\cdot e=e$, vist com a categoria d'un sol objecte $\ast$; el functor oblit en dona el graf subjacent:
\[
\begin{tikzpicture}[baseline=(m.center)]
  \matrix (m) [cat matrix] { \ast \\ };
  \draw[cat] (m-1-1) to[out=55,in=125,looseness=9] node[above] {\(\id_\ast\)} (m-1-1);
  \draw[cat] (m-1-1) to[out=-125,in=-55,looseness=9] node[below] {\(e\)} (m-1-1);
\end{tikzpicture}
\qquad\stackrel{U}{\longmapsto}\qquad
\begin{tikzpicture}[baseline=(m.center)]
  \matrix (m) [cat matrix] { \ast \\ };
  \draw[cat] (m-1-1) to[out=55,in=125,looseness=9] node[above] {\(\id_\ast\)} (m-1-1);
  \draw[cat] (m-1-1) to[out=-125,in=-55,looseness=9] node[below] {\(e\)} (m-1-1);
\end{tikzpicture}
\]
Sobre aquest graf ---que conserva les dues arestes $\id_\ast$ i $e$, però on $U$ ha oblidat tota la composició: les igualtats $\id_\ast\comp\id_\ast=\id_\ast$, $\id_\ast\comp e=e$, $e\comp\id_\ast=e$ i $e\comp e=e$ (les tres que involucren $\id_\ast$ són el que en feia la identitat)--- considerem l'assignació $\alpha$ que fixa el vèrtex i envia tots dos bucles a $e$:
\[
\begin{tikzpicture}[baseline=(m.center)]
  \matrix (m) [cat matrix] { \ast \\ };
  \draw[cat] (m-1-1) to[out=55,in=125,looseness=9] node[above] {\(\id_\ast\)} (m-1-1);
  \draw[cat] (m-1-1) to[out=-125,in=-55,looseness=9] node[below] {\(e\)} (m-1-1);
\end{tikzpicture}
\qquad\stackrel{\alpha}{\longmapsto}\qquad
\begin{tikzpicture}[baseline=(m.center)]
  \matrix (m) [cat matrix] { \ast \\ };
  \draw[cat] (m-1-1) to[out=55,in=125,looseness=9] node[above] {\(e\)} (m-1-1);
\end{tikzpicture}
\]
Com a morfisme de grafs, $\alpha$ és vàlid: respecta origen i destí, perquè tots els bucles van del vèrtex a si mateix. Però no prové de cap functor: si fos $\alpha=U(F)$, el functor $F$ enviaria $\id_\ast$ a $e$, i cap functor no pot fer-ho ---ha d'enviar $\id_\ast$ a una identitat, i $e$ no ho és---. Els morfismes de grafs són, doncs, més abundants que els functors. Aquest mateix contraexemple, des de l'angle de la definició de functor, es reprèn a l'exercici~\ref{ex:identitat-necessaria}.

\subsection{La categoria lliure sobre un graf}

En sentit contrari, la construcció de la categoria lliure generada per un graf, $\Free(G)$, que vam definir a la secció~\ref{sec:grafs-lliures}, també és functorial. Recordem-ne l'acció sobre objectes: $\Free(G)$ té per objectes els vèrtexs de $G$ i per morfismes els camins d'arestes, amb la concatenació com a composició i el camí buit com a identitat de cada vèrtex.

Vegem l'acció sobre morfismes. Sigui $\alpha\colon G\to H$ un morfisme de grafs, amb components $\alpha_V\colon V_G\to V_H$ i $\alpha_E\colon E_G\to E_H$. Definim un functor
\[
\Free(\alpha)\colon\Free(G)\to\Free(H)
\]
que, sobre objectes, actua com $\alpha_V$, i sobre morfismes transforma cada camí aresta a aresta: el camí $(f_1,\dots,f_k)$ de $G$ va al camí $(\alpha_E(f_1),\dots,\alpha_E(f_k))$ de $H$, i el camí buit en $v$, $()_v$, va al camí buit en $\alpha_V(v)$. L'assignació està ben definida ---com que $\alpha$ preserva origen i destí, arestes encadenades van a arestes encadenades, i el resultat és de nou un camí--- i és un functor: transforma concatenacions en concatenacions i camins buits en camins buits, és a dir, composicions en composicions i identitats en identitats.

Finalment, $\Free$ respecta la identitat i la composició de $\cat{Graph}$: $\Free(\id_G)=\id_{\Free(G)}$ i $\Free(\beta\comp\alpha)=\Free(\beta)\comp\Free(\alpha)$, ja que aplicar $\alpha$ i després $\beta$ aresta a aresta és el mateix que aplicar $\beta\comp\alpha$. Tenim, doncs, un functor
\[
\Free\colon\cat{Graph}\to\cat{Cat}.
\]

Vegem-ho amb el mateix graf d'un bucle de l'exemple anterior. Sigui $G$ el graf amb un sol vèrtex $\ast$ i una sola aresta $a$, un bucle. Els camins de $G$ són les paraules en la lletra $a$ ---$\varepsilon, a, aa, aaa,\dots$---, de manera que $\Free(G)$ és la categoria d'un sol objecte els morfismes de la qual són aquestes paraules, amb la concatenació per composició i la paraula buida $\varepsilon$ per identitat:
\[
\begin{tikzpicture}[baseline=(m.center)]
  \matrix (m) [cat matrix] { \ast \\ };
  \draw[cat] (m-1-1) to[out=55,in=125,looseness=9] node[above] {\(a\)} (m-1-1);
\end{tikzpicture}
\qquad\stackrel{\Free}{\longmapsto}\qquad
\begin{tikzpicture}[baseline=(m.center)]
  \matrix (m) [cat matrix] { \ast \\ };
  \draw[cat] (m-1-1) to[out=120,in=175,looseness=12] node[left] {\(\varepsilon\)} (m-1-1);
  \draw[cat] (m-1-1) to[out=65,in=115,looseness=12] node[above] {\(a\)} (m-1-1);
  \draw[cat] (m-1-1) to[out=5,in=60,looseness=12] node[right] {\(aa\)} (m-1-1);
  \draw[cat] (m-1-1) to[out=-55,in=-5,looseness=12] node[right] {\(\dots\)} (m-1-1);
  \draw[cat] (m-1-1) to[out=-125,in=-70,looseness=12] node[below] {\(a^n\)} (m-1-1);
\end{tikzpicture}
\]
Comptant la paraula buida com a $a^0$, els morfismes són exactament les potències $a^n$ amb $n\geq 0$, i concatenar-les suma els exponents: $a^m\comp a^n=a^{m+n}$. Aquest monoide de morfismes és $(\mathbb{N},+,0)$, de manera que $\Free(G)=\cat{B}\mathbb{N}$. Mentre que $U$ oblidava estructura, $\Free$ en crea: d'una sola aresta en surten infinits morfismes, un per cada longitud de camí.

Cal notar que $\Free$ i $U$ són functors \emph{entre categories grans} ($\cat{Graph}$ i $\cat{Cat}$ no són objectes de la $\cat{Cat}$ de les categories petites); les tractem com a categories localment petites, no com a fletxes de $\cat{Cat}$.

\subsection{La relació entre \texorpdfstring{$\Free$}{Free} i \texorpdfstring{$U$}{U}}\label{subsec:free-u}

Les dues construccions van en sentits contraris, però no són inverses l'una de l'altra. El functor $U$ oblida quins morfismes són identitats i com es componen; $\Free$ afegeix formalment identitats i composicions al graf, sense imposar cap altra relació. Quan partim del graf subjacent d'una categoria $\cat{C}$, la composició lliure crea camins formals nous, de manera que $\Free(U(\cat{C}))$ no és $\cat{C}$ en general; però hi ha un functor
\[
\Free(U(\cat{C}))\to\cat{C}
\]
que avalua cada camí formal mitjançant la composició de $\cat{C}$. Per exemple, prenem per $\cat{C}$ el monoide $M=(\{1,e\},\cdot,1)$ amb $e\cdot e=e$, vist com a categoria d'un sol objecte $\ast$. El seu graf subjacent té un vèrtex i dos bucles $\id_\ast$ i $e$, de manera que $\Free(U(\cat{C}))$ té per morfismes \emph{totes} les paraules en aquestes dues lletres ---infinites---, amb la paraula buida $\varepsilon$ per identitat. El functor $\Free(U(\cat{C}))\to\cat{C}$ avalua cada paraula multiplicant-la a $M$: $\varepsilon\mapsto\id_\ast$, $ee\mapsto e\comp e=e$, $\id_\ast e\mapsto e$, i així successivament. Col·lapsa, doncs, la infinitat de camins formals sobre els dos únics morfismes de $\cat{C}$, recuperant les composicions que $\Free$ havia descartat. Aquest functor il·lustra una propietat que caracteritza $\Free(G)$. Sigui $\cat{D}$ una categoria petita. Definir un functor $\Free(G)\to\cat{D}$ equival a triar les imatges dels vèrtexs \emph{i} de les arestes de $G$, respectant origen i destí ---és a dir, a donar un morfisme de grafs $G\to U(\cat{D})$---, i aquesta tria s'estén de manera única a $\Free(G)$. Cal comptar totes dues assignacions: un functor ha d'enviar cada objecte de $\Free(G)$ ---cada vèrtex de $G$--- a un objecte de $\cat{D}$, també els vèrtexs aïllats, que no són extrem de cap aresta. Per exemple, si $G$ té un únic vèrtex i cap aresta, no hi ha cap imatge d'aresta per triar, però un functor $\Free(G)\to\cat{D}$ encara exigeix escollir un objecte de $\cat{D}$; és precisament la component sobre vèrtexs del morfisme de grafs $G\to U(\cat{D})$, que en aquest cas es redueix a triar aquest objecte. La formulació via $G\to U(\cat{D})$ recull sempre les dues dades, perquè un morfisme de grafs consta d'una aplicació sobre vèrtexs i una sobre arestes. Aquesta correspondència és una bijecció
\[
\Hom_{\cat{Cat}}\big(\Free(G),\cat{D}\big)\;\cong\;\Hom_{\cat{Graph}}\big(G,U(\cat{D})\big).
\]
Aquesta és la formulació precisa de \enquote{lliure}. La retrobarem com a exemple al capítol~\refcap{cap:yoneda}{4}, i en tota la seva generalitat com l'\emph{adjunció} $\Free\dashv U$ al capítol~\refcap{cap:adjuncions}{6}.

\begin{observacio}
Reconèixer els functors com a fletxes d'una categoria és el primer pas per poder plantejar-nos si són isomorfismes, si es componen o si hi ha dualitat. Dins $\cat{Cat}$, per exemple, una aplicació monòtona $P\to Q$ entre preordres petits és una fletxa; en canvi, $\Free$ i $U$ són functors entre categories grans, i per això no són morfismes de $\cat{Cat}$; això no els fa menys legítims, sinó que simplement no els agrupem com a fletxes de cap categoria de categories.
\end{observacio}

\section{Functors covariants i contravariants}\label{sec:variancia}

Fins ara els functors preserven el sentit de les fletxes. Sovint, però, apareixen assignacions que \emph{inverteixen} el sentit: enviar $f\colon A\to B$ a un morfisme $F(B)\to F(A)$. Parlem aleshores de \textbf{functors contravariants}, que s'expliquen millor mitjançant la categoria oposada.

\begin{definicio}\index{functor!contravariant}\index{functor!covariant}
Un \textbf{functor contravariant} de $\cat{C}$ a $\cat{D}$ és un functor
\[
F\colon\cat{C}\op\to\cat{D}.
\]
Per contrast, els functors de la definició original s'anomenen \textbf{covariants}.
\end{definicio}

Desplegant la definició de functor sobre $\cat{C}\op$, un functor contravariant assigna a cada objecte $A$ un objecte $F(A)$ i a cada morfisme $f\colon A\to B$ de $\cat{C}$ un morfisme en sentit invertit
\[
F(f)\colon F(B)\to F(A),
\]
de manera que $F(\id_A)=\id_{F(A)}$ i $F(g\comp f)=F(f)\comp F(g)$. La inversió de l'ordre en aquesta última igualtat no és cap condició nova: és la condició de functor covariant sobre $\cat{C}\op$, un cop recordem que compondre a $\cat{C}\op$ és compondre a $\cat{C}$ en l'ordre invers. Aquest és un primer exemple de la potència de la categoria oposada: permet tractar covariància i contravariància amb una sola definició.

Diagramàticament, un triangle commutatiu de $\cat{C}$ es transforma en un triangle commutatiu de $\cat{D}$ amb les fletxes imatge invertides:
\[
\begin{tikzpicture}[baseline=(m.center)]
  \matrix (m) [matrix of math nodes, column sep=3.2em, row sep=2.6em] {
    A & B \\
      & C \\
  };
  \draw[-{Stealth[scale=1.1]}] (m-1-1) -- node[above] {$f$} (m-1-2);
  \draw[-{Stealth[scale=1.1]}] (m-1-2) -- node[right] {$g$} (m-2-2);
  \draw[-{Stealth[scale=1.1]}] (m-1-1) -- node[below left] {$g\comp f$} (m-2-2);
\end{tikzpicture}
\qquad\stackrel{F}{\longmapsto}\qquad
\begin{tikzpicture}[baseline=(m.center)]
  \matrix (m) [matrix of math nodes, column sep=3.6em, row sep=2.6em] {
    F(A) & F(B) \\
         & F(C) \\
  };
  \draw[{Stealth[scale=1.1]}-] (m-1-1) -- node[above] {$F(f)$} (m-1-2);
  \draw[{Stealth[scale=1.1]}-] (m-1-2) -- node[right] {$F(g)$} (m-2-2);
  \draw[{Stealth[scale=1.1]}-] (m-1-1) -- node[below left] {$F(g\comp f)$} (m-2-2);
\end{tikzpicture}
\]

\begin{definicio}[Functor oposat]\label{def:functor-oposat}\index{functor!oposat}
Tot functor $F\colon\cat{C}\to\cat{D}$ indueix un \textbf{functor oposat}
\[
F\op\colon\cat{C}\op\to\cat{D}\op,
\]
que actua igual que $F$ sobre objectes i sobre morfismes, però llegit entre les categories oposades. Com que compondre a $\cat{C}\op$ i a $\cat{D}\op$ és compondre a $\cat{C}$ i a $\cat{D}$ en l'ordre invers, les dues condicions de functor es transporten sense canvis.
\end{definicio}

\begin{observacio}
El pas a l'oposat és, de fet, functorial: assigna a cada categoria petita $\cat{C}$ la seva oposada $\cat{C}\op$ i a cada functor $F$ el functor oposat $F\op$, i respecta composicions i identitats, de manera que
\[
(-)\op\colon\cat{Cat}\to\cat{Cat}
\]
és un functor ---covariant: envia $F\colon\cat{C}\to\cat{D}$ a $F\op\colon\cat{C}\op\to\cat{D}\op$, en el mateix sentit. A més, aplicar-lo dues vegades torna al punt de partida, tant sobre objectes com sobre morfismes: $(\cat{C}\op)\op=\cat{C}$ i $(F\op)\op=F$. Dit altrament, $(-)\op\comp(-)\op=\id_{\cat{Cat}}$: el functor $(-)\op$ és la seva pròpia inversa, és a dir, una \textbf{involució}. En particular és un \textbf{isomorfisme de categories} de $\cat{Cat}$ en ella mateixa ---un endofunctor amb functor invers, que a més coincideix amb ell mateix. És la manifestació, al nivell de $\cat{Cat}$, de l'autodualitat que ja hem trobat objecte a objecte amb el principi de dualitat. Convé recordar, això sí, que tot plegat viu sobre categories \emph{petites}, les úniques que són objectes de $\cat{Cat}$; per a una categoria gran, l'oposada segueix tenint sentit, però no com a imatge d'una fletxa de $\cat{Cat}$.
\end{observacio}

\begin{proposicio}[Variància d'una composició]\label{prop:variancia-composicio}\index{functor!variància d'una composició}\index{regla dels signes (variància)}
Siguin $F$ i $G$ dos functors, cadascun covariant o contravariant, amb el codomini de $F$ i el domini de $G$ iguals com a categoria (a banda de la variància). L'assignació sobre objectes $A\mapsto G(F(A))$ és de nou un functor, i la seva variància segueix la \textbf{regla dels signes}: la composició és covariant si $F$ i $G$ tenen la mateixa variància, i contravariant si en tenen de diferents. En particular, la composició de dos functors contravariants és covariant.
\end{proposicio}

\begin{prova}
Un functor contravariant $\cat{C}\to\cat{D}$ és, per definició, un functor covariant $\cat{C}\op\to\cat{D}$, i la composició de functors torna a ser un functor (ho hem vist en construir $\cat{Cat}$). Només cal fer encaixar dominis i codominis, i d'això se n'ocupa el functor oposat.

Si $F$ i $G$ són tots dos covariants, $G\comp F\colon\cat{C}\to\cat{E}$ és directament la composició, covariant. Si tots dos són contravariants, escrits com $F\colon\cat{C}\op\to\cat{D}$ i $G\colon\cat{D}\op\to\cat{E}$, prenem el functor oposat $F\op\colon\cat{C}\to\cat{D}\op$ (definició~\ref{def:functor-oposat}), que actua igual que $F$ sobre objectes; aleshores
\[
G\comp F\op\colon\cat{C}\to\cat{E}
\]
és una composició de functors ---per tant un functor--- covariant, i sobre objectes fa $A\mapsto G(F(A))$. Els dos casos mixtos són idèntics en mètode: un ajust anàleg deixa el domini en $\cat{C}\op$ i la composició resulta contravariant. La variància només depèn, doncs, de la paritat del nombre de factors contravariants.
\end{prova}

\subsection{Els homfunctors}\label{subsec:homfunctors}

L'exemple més important de functors covariants i contravariants ---i el que sostindrà bona part del que ve després, fins al lema de Yoneda--- surt de mirar els morfismes d'una categoria cap a un objecte fix, o des d'un objecte fix.

\begin{definicio}[Els homfunctors]\label{def:homfunctors}\label{ex:homfunctors}\index{functor!Hom}\index{homfunctor}
Sigui $\cat{C}$ una categoria localment petita. Fixat un objecte $A$ de $\cat{C}$, definim \textbf{homfunctor covariant}\index{homfunctor!covariant}\index{functor!Hom!covariant}
\[
\Hom(A,-)\colon\cat{C}\to\cat{Set}
\]
com el functor que envia cada objecte $X$ al conjunt $\Hom(A,X)$ i cada morfisme $f\colon X\to Y$ a l'aplicació
\[
f_*\colon\Hom(A,X)\to\Hom(A,Y),
\qquad
\varphi\mapsto f\comp\varphi,
\]
la \textbf{postcomposició} amb $f$. Fixat un objecte $B$, definim \textbf{homfunctor contravariant}\index{homfunctor!contravariant}\index{functor!Hom!contravariant}
\[
\Hom(-,B)\colon\cat{C}\op\to\cat{Set}
\]
com el functor que envia cada objecte $X$ al conjunt $\Hom(X,B)$ i cada morfisme $f\colon X\to Y$ a l'aplicació
\[
f^*\colon\Hom(Y,B)\to\Hom(X,B),
\qquad
\psi\mapsto\psi\comp f,
\]
la \textbf{precomposició} amb $f$.
\end{definicio}

\begin{notacio}\label{not:hom-morfismes}\index{Hom@$\Hom$!aplicat a un morfisme}
Com per a qualsevol functor $F$, que envia un morfisme $f$ a $F(f)$, escriurem també
\[
\Hom(A,f)\coloneqq f_*
\qquad\text{i}\qquad
\Hom(f,B)\coloneqq f^*
\]
per a les imatges de $f\colon X\to Y$ pels homfunctors $\Hom(A,-)$ i $\Hom(-,B)$, respectivament.
\end{notacio}

En el cas covariant, un element $\varphi\colon A\to X$ es perllonga amb $f$ per obtenir $f\comp\varphi\colon A\to Y$; en el contravariant, un element $\psi\colon Y\to B$ es precedeix de $f$ per obtenir $\psi\comp f\colon X\to B$:
\[
\begin{tikzpicture}[baseline=(m.center)]
  \matrix (m) [matrix of math nodes, column sep=3.2em, row sep=2.6em] {
    A & X \\
      & Y \\
  };
  \draw[-{Stealth[scale=1.1]}] (m-1-1) -- node[above] {$\varphi$} (m-1-2);
  \draw[-{Stealth[scale=1.1]}] (m-1-2) -- node[right] {$f$} (m-2-2);
  \draw[-{Stealth[scale=1.1]}] (m-1-1) -- node[below left] {$f\comp\varphi$} (m-2-2);
\end{tikzpicture}
\qquad\qquad
\begin{tikzpicture}[baseline=(m.center)]
  \matrix (m) [matrix of math nodes, column sep=3.2em, row sep=2.6em] {
    X & Y \\
      & B \\
  };
  \draw[-{Stealth[scale=1.1]}] (m-1-1) -- node[above] {$f$} (m-1-2);
  \draw[-{Stealth[scale=1.1]}] (m-1-2) -- node[right] {$\psi$} (m-2-2);
  \draw[-{Stealth[scale=1.1]}] (m-1-1) -- node[below left] {$\psi\comp f$} (m-2-2);
\end{tikzpicture}
\]
Observem que, en el cas covariant, $f$ actua sobre el \emph{destí} de les fletxes (les allarga per l'extrem final), mentre que en el contravariant actua sobre l'\emph{origen} (les allarga per l'extrem inicial), i això és el que inverteix el sentit del functor. La functorialitat és conseqüència directa de l'associativitat de la composició i de les lleis de les identitats. El que importa és l'ordre que en resulta:
\[
(g\comp f)_*=g_*\comp f_*,
\qquad\qquad
(g\comp f)^*=f^*\comp g^*.
\]
En el cas covariant l'ordre es conserva; en el contravariant s'inverteix, perquè $f$ actua per l'extrem inicial de les fletxes. Les comprovacions detallades són els exercicis resolts~\ref{pr-hom-cov} i~\ref{pr-hom-contra} de la Pràctica. Aquests dos functors són centrals en teoria de categories i els retrobarem en el lema de Yoneda (capítol~\refcap{cap:yoneda}{4}).

Els homfunctors traslladen els isomorfismes de $\cat{C}$ a bijeccions de conjunts. Això és un cas particular de la proposició~\ref{prop:functor-preserva-iso}, però val la pena escriure'l explícitament, perquè és la forma en què l'usarem sovint.

\begin{corollari}[Els homfunctors transformen isos en bijeccions]\label{prop:homfunctor-iso}\index{homfunctor!i isomorfismes}
Sigui $\cat{C}$ localment petita i sigui $f\colon X\to Y$ un isomorfisme de $\cat{C}$. Aleshores, per a tot objecte $A$, la postcomposició
\[
f_*\colon\Hom(A,X)\to\Hom(A,Y),\qquad\varphi\mapsto f\comp\varphi,
\]
és una bijecció, amb inversa $(f^{-1})_*$; i, per a tot objecte $B$, la precomposició
\[
f^*\colon\Hom(Y,B)\to\Hom(X,B),\qquad\psi\mapsto\psi\comp f,
\]
és una bijecció, amb inversa $(f^{-1})^*$.
\end{corollari}

\begin{prova}
Un isomorfisme de $\cat{Set}$ no és res més que una bijecció. Apliquem, doncs, la proposició~\ref{prop:functor-preserva-iso} a cada homfunctor. Per al covariant $\Hom(A,-)$, la imatge de l'isomorfisme $f$ és $f_*$, que resulta ser un isomorfisme de $\cat{Set}$ amb invers la imatge de $f^{-1}$, és a dir $(f^{-1})_*$. Per al contravariant $\Hom(-,B)$ ---un functor sobre $\cat{C}\op$, on $f$ segueix sent un isomorfisme--- la imatge de $f$ és $f^*$, que és igualment un isomorfisme de $\cat{Set}$ amb invers $(f^{-1})^*$. En tots dos casos, ser isomorfisme a $\cat{Set}$ vol dir ser bijecció.
\end{prova}

Dit d'una altra manera: si $X\cong Y$, aleshores $\Hom(A,X)\cong\Hom(A,Y)$ i $\Hom(Y,B)\cong\Hom(X,B)$ com a conjunts, per a qualssevol $A$ i $B$. Els objectes isomorfs són, doncs, indistingibles vistos a través de qualsevol homfunctor: una primera pista de la idea que un objecte queda determinat pels morfismes que hi arriben o que en surten, que el lema de Yoneda portarà fins al final.

Els functors contravariants cap a $\cat{Set}$, com l'homfunctor $\Hom(-,B)$, apareixen tan sovint que reben un nom propi.

\begin{definicio}[Prefeix]\label{def:prefeix}\index{prefeix@prefeix}
Sigui $\cat{C}$ una categoria. Un \textbf{prefeix} sobre $\cat{C}$ és un functor contravariant de $\cat{C}$ a $\cat{Set}$, és a dir, un functor
\[
P\colon\cat{C}\op\to\cat{Set}.
\]
Explícitament, $P$ assigna a cada objecte $A$ de $\cat{C}$ un conjunt $P(A)$ i a cada morfisme $f\colon A\to B$ una aplicació $P(f)\colon P(B)\to P(A)$, de manera que $P(\id_A)=\id_{P(A)}$ i $P(g\comp f)=P(f)\comp P(g)$.
\end{definicio}

Si $\cat{C}$ és localment petita, cada homfunctor contravariant $\Hom(-,B)$ és un prefeix sobre $\cat{C}$, i el functor de parts de l'exemple següent n'és un sobre $\cat{Set}$. Els prefeixos tornaran a aparèixer als capítols~\refcap{cap:yoneda}{4} i~\refcap{cap:topos}{9}.

Així, \emph{prefeix sobre $\cat{C}$} i \emph{functor contravariant de $\cat{C}$ a $\cat{Set}$} són sinònims: el nom no afegeix cap condició, ni tan sols de mida. A partir d'ara, per a un functor contravariant cap a $\cat{Set}$ direm sempre \emph{prefeix},\footnote{A la literatura també es parla de \emph{prefeixos sobre $\cat{C}$ amb valors a $\cat{D}$} (per exemple, prefeixos de grups o d'espais vectorials) per als functors $\cat{C}\op\to\cat{D}$. En aquest text, \emph{prefeix} vol dir sempre amb valors a $\cat{Set}$.} i n'escriurem la signatura $P\colon\cat{C}\op\to\cat{Set}$ quan convingui. Reservarem l'expressió \emph{functor contravariant} per als que arriben a altres categories, com el functor dual de l'exemple~\ref{ex:homfunctor-dual}. Els functors \emph{covariants} $\cat{C}\to\cat{Set}$ no reben nom propi en aquest text: són, simplement, els prefeixos sobre $\cat{C}\op$.

\begin{exemple}[El functor de parts i l'homfunctor $\Hom(-,\{0,1\})$]\label{ex:functor-parts}\index{functor!de parts}\index{homfunctor!i functor de parts}
Comencem dins de $\cat{Set}$, prenent com a objecte fixat el conjunt de dos elements $\{0,1\}$. Cada subconjunt $A\subseteq X$ es correspon bijectivament amb la seva \textbf{funció característica} $\chi_A\colon X\to\{0,1\}$, que val $1$ sobre $A$ i $0$ fora; recíprocament, cada aplicació $X\to\{0,1\}$ és la característica del subconjunt on val $1$. Les funcions característiques donen, doncs, per a cada conjunt $X$, una bijecció canònica
\[
\mathcal P(X)\;\cong\;\Hom_{\cat{Set}}(X,\{0,1\}).
\]
Els dos conjunts no són iguals ---un conté subconjunts, l'altre aplicacions---, però la bijecció no depèn de cap tria.

Vegem que, sota aquestes bijeccions, la precomposició correspon a l'antiimatge. Per a $f\colon X\to Y$, l'homfunctor envia una aplicació $\chi\colon Y\to\{0,1\}$ a $\chi\comp f\colon X\to\{0,1\}$; i si $\chi=\chi_B$ és la característica de $B\subseteq Y$, aleshores $\chi_B\comp f$ val $1$ exactament sobre els $x$ amb $f(x)\in B$, és a dir, és la característica de l'antiimatge $f^{-1}(B)$. Així, la precomposició
\[
\chi\longmapsto\chi\comp f
\]
és precisament l'aplicació d'antiimatge
\[
f^{-1}\colon\mathcal P(Y)\to\mathcal P(X),
\qquad
B\mapsto\{x\in X : f(x)\in B\}.
\]
Aquí $f^{-1}$ denota l'aplicació d'antiimatge, no l'invers de $f$, que pot no existir. Les antiimatges defineixen el \textbf{functor de parts} contravariant
\[
\mathcal P\colon\cat{Set}\op\to\cat{Set},\qquad X\mapsto\mathcal P(X),\qquad f\mapsto f^{-1},
\]
i la seva contravariància, $(g\comp f)^{-1}=f^{-1}\comp g^{-1}$, es pot transportar, a través de les bijeccions anteriors, des de la de l'homfunctor $\Hom(-,\{0,1\})$; la comprovació directa és l'exercici resolt~\ref{pr-parts-contra} de la Pràctica.

Per tant, el functor de parts contravariant i l'homfunctor $\Hom_{\cat{Set}}(-,\{0,1\})$ descriuen la mateixa construcció fins a una identificació canònica: provar cada conjunt amb un objecte distingit, el conjunt $\{0,1\}$ de valors de veritat. No són literalment el mateix functor. El llenguatge de les transformacions naturals (capítol~\ref{cap:transformacions}) permetrà expressar aquesta afirmació amb precisió: tots dos functors seran \emph{naturalment isomorfs} (exemple~\ref{ex:isos-naturals-cap2}).

Convé un avís sobre l'altra variància, i val la pena justificar-lo. Un homfunctor covariant $\Hom_{\cat{Set}}(\{0,1\},-)$ envia cada conjunt $X$ al conjunt d'aplicacions $\{0,1\}\to X$. Una tal aplicació queda completament determinada per la parella d'imatges $(x_0,x_1)$ amb $x_0$ la imatge de $0$ i $x_1$ la de $1$, i tota parella d'elements de $X$ prové d'una aplicació; per tant
\[
\Hom_{\cat{Set}}(\{0,1\},X)\;\cong\;X\times X,
\]
i sobre morfismes la postcomposició amb $f\colon X\to Y$ actua component a component, $(x_0,x_1)\mapsto(f(x_0),f(x_1))$. Aquest és, doncs, el functor \enquote{parell ordenat}, ben diferent del functor de parts: $\Hom(\{0,1\},X)\cong X\times X$ té sempre $\lvert X\rvert^2$ elements (quan $X$ és finit), mentre que $\mathcal P(X)$ en té $2^{\lvert X\rvert}$. La correspondència canònica entre el functor de parts i un homfunctor és, per tant, un fenomen exclusiu de la cara contravariant: $\mathcal P$ correspon a $\Hom(-,\{0,1\})$, no a $\Hom(\{0,1\},-)$. La raó de fons és que a $\mathcal P(X)$ el conjunt $\{0,1\}$ fa de \emph{destí} (cada subconjunt és una aplicació \emph{cap a} $\{0,1\}$), i fixar el destí dona un homfunctor contravariant; posar-lo com a \emph{origen} dona una construcció covariant, però tota una altra. La versió covariant genuïna del functor de parts existeix ---envia $f$ a la imatge directa $A\mapsto\{f(x):x\in A\}$--- i és functorial, però no correspon a cap homfunctor; quan calgui distingir-la de $\mathcal P$ l'escriurem $\mathcal P_!$.
\end{exemple}

\begin{exemple}[Les dues variàncies del functor de parts]\label{ex:parts-dues-variancies}\index{functor!de parts!covariant i contravariant}
Posem en paral·lel les dues maneres de fer functorial l'assignació $X\mapsto\mathcal P(X)$. Totes dues coincideixen sobre objectes i difereixen en tot el que fan amb una aplicació $f\colon X\to Y$:
\[
\renewcommand{\arraystretch}{1.3}
\begin{array}{l|l|l}
 & \mathcal P_!\colon\cat{Set}\to\cat{Set}\ \text{(covariant)} & \mathcal P\colon\cat{Set}\op\to\cat{Set}\ \text{(contravariant)}\\ \hline
\text{acció sobre } f & \text{imatge directa } A\mapsto f(A) & \text{antiimatge } B\mapsto f^{-1}(B)\\
\text{tipus} & \mathcal P_!(f)\colon\mathcal P(X)\to\mathcal P(Y) & \mathcal P(f)\colon\mathcal P(Y)\to\mathcal P(X)\\
\text{identitat} & \id_X(A)=A & \id_X^{-1}(B)=B\\
\text{composició} & (g\comp f)(A)=g\bigl(f(A)\bigr) & (g\comp f)^{-1}(C)=f^{-1}\bigl(g^{-1}(C)\bigr)\\
\text{llei functorial} & \mathcal P_!(g\comp f)=\mathcal P_!(g)\comp\mathcal P_!(f) & \mathcal P(g\comp f)=\mathcal P(f)\comp\mathcal P(g)
\end{array}
\]
Llegida per files, la taula segueix l'ordre de l'observació~\ref{obs:comprovar-tipus}: primer el tipus, que ja fixa la variància, i després les dues equacions. Les comprovacions completes són els exercicis resolts~\ref{pr-parts-cov} i~\ref{pr-parts-contra} de la Pràctica. Observem que a la columna contravariant l'ordre de $f$ i $g$ s'inverteix: és el mateix fenomen que a $(g\comp f)^*=f^*\comp g^*$ per a l'homfunctor contravariant.

Les dues columnes no són independents. Per a $A\subseteq X$ i $B\subseteq Y$,
\[
f(A)\subseteq B\quad\Longleftrightarrow\quad A\subseteq f^{-1}(B),
\]
perquè totes dues condicions diuen que $f(a)\in B$ per a tot $a\in A$. Aquesta equivalència entre les dues accions, una en cada sentit, és el primer exemple d'una \emph{adjunció}; la retrobarem al capítol~\refcap{cap:adjuncions}{6} com a $\exists_f\dashv f^{*}$, l'origen categòric del quantificador existencial (exemple~\refalt{ex:quantificadors-set}{d'imatge, antiimatge i quantificadors del capítol~6}).
\end{exemple}

En aquest exemple l'objecte fixat feia de \emph{destí}; fixant-lo com a \emph{origen} obtenim homfunctors covariants amb una lectura igual de concreta. Ho veiem a $\cat{Graph}$, on l'estructura dels objectes ja no és la d'un simple conjunt.

\begin{exemple}[Vèrtexs i arestes com a homfunctors]\label{ex:homfunctor-grafs}\index{homfunctor!a Graph}\index{graf dirigit!vèrtexs i arestes com a homfunctors}
Treballem a $\cat{Graph}$, amb la notació de la subsecció~\ref{subsec:grafs}: un graf $G$ té vèrtexs $V_G$, arestes $E_G$ i aplicacions origen i destí $s_G,t_G\colon E_G\to V_G$. Triem dos grafs \enquote{patró} ben senzills:
\[
\mathbf V=\bigl(\{\ast\},\ \varnothing\bigr),
\qquad
\mathbf E=\bigl(\{0,1\},\ \{e\}\bigr)\ \text{amb } s(e)=0,\ t(e)=1,
\]
és a dir, $\mathbf V$ és un sol vèrtex sense arestes i $\mathbf E$ és una sola aresta entre dos vèrtexs.

Un morfisme de grafs $\mathbf V\to G$ no té cap aresta on decidir res, i sobre l'únic vèrtex tria un vèrtex de $G$; per tant queda determinat per aquesta tria, i n'hi ha un per cada vèrtex:
\[
\Hom_{\cat{Graph}}(\mathbf V,G)\;\cong\;V_G.
\]
Un morfisme de grafs $\mathbf E\to G$ tria una aresta $\varepsilon$ de $G$ ---la imatge de $e$--- i, forçosament, els seus dos extrems han d'anar a $s_G(\varepsilon)$ i $t_G(\varepsilon)$, de manera que la tria de l'aresta ja determina tot el morfisme; n'hi ha un per cada aresta:
\[
\Hom_{\cat{Graph}}(\mathbf E,G)\;\cong\;E_G.
\]
Així, els homfunctors covariants $\Hom_{\cat{Graph}}(\mathbf V,-)$ i $\Hom_{\cat{Graph}}(\mathbf E,-)$ corresponen canònicament, respectivament, al functor \enquote{conjunt de vèrtexs} i al functor \enquote{conjunt d'arestes} $\cat{Graph}\to\cat{Set}$. Sobre un morfisme $F=(F_V,F_E)\colon G\to H$, el functor $\Hom_{\cat{Graph}}(\mathbf V,-)$ actua per la postcomposició $F_*\colon\Hom_{\cat{Graph}}(\mathbf V,G)\to\Hom_{\cat{Graph}}(\mathbf V,H)$, $v\mapsto F\comp v$ (definició~\ref{def:homfunctors}). Si $v$ tria el vèrtex $v(\ast)$, aleshores $F\comp v$ tria el vèrtex $F_V(v(\ast))$; per tant, llegida a través de la bijecció $\Hom_{\cat{Graph}}(\mathbf V,G)\cong V_G$, la postcomposició $F_*$ és exactament $F_V\colon V_G\to V_H$. Anàlogament, per a $\Hom_{\cat{Graph}}(\mathbf E,-)$, la postcomposició $F_*$ envia el morfisme que tria l'aresta $\varepsilon$ al que tria $F_E(\varepsilon)$, i correspon a $F_E\colon E_G\to E_H$.

El que hem trobat és que dos functors oblit familiars ---prendre el conjunt de vèrtexs i prendre el conjunt d'arestes--- corresponen canònicament als homfunctors covariants $\Hom_{\cat{Graph}}(\mathbf V,-)$ i $\Hom_{\cat{Graph}}(\mathbf E,-)$, per a una tria adient de l'objecte d'origen; al capítol~\ref{cap:transformacions} reconeixerem aquestes identificacions com a isomorfismes naturals. És el mateix fenomen que fa que, a $\cat{Set}$, el functor identitat correspongui canònicament a $\Hom_{\cat{Set}}(\{\ast\},-)$: provar un objecte amb els morfismes des d'un patró petit en llegeix una part de l'estructura. Que un functor cap a $\cat{Set}$ es pugui reconstruir així, com un $\Hom(A,-)$ per a cert objecte $A$, és una propietat prou important perquè hi tornem amb nom propi ---els \emph{functors representables}--- i és el punt de partida del lema de Yoneda (capítol~\refcap{cap:yoneda}{4}); de moment n'hi ha prou de retenir que aquests functors oblit són homfunctors disfressats, llevat d'una identificació canònica.
\end{exemple}

Als exemples anteriors la categoria d'arribada era $\cat{Set}$. Quan els homconjunts tenen \emph{ells mateixos} estructura, l'homfunctor la recull i aterra en una categoria més rica; el cas paradigmàtic és el dels espais vectorials.

\begin{exemple}[El dual d'un espai vectorial com a homfunctor]\label{ex:homfunctor-dual}\index{homfunctor!i espai dual}\index{espai dual!com a homfunctor}\index{aplicació dual}
Fixem un cos $\mathbb K$ i treballem a $\cat{Vect}_{\mathbb K}$. Aquí cada homconjunt $\Hom_{\mathbb K}(V,W)$ no és un conjunt qualsevol: és ell mateix un $\mathbb K$-espai vectorial, amb la suma d'aplicacions lineals i el producte per escalars. Prenent com a segon objecte el cos $\mathbb K$ vist com a espai vectorial sobre si mateix, l'homconjunt
\[
\Hom_{\mathbb K}(V,\mathbb K)=V^*
\]
és exactament l'\textbf{espai dual} de $V$: les formes lineals sobre $V$. Com a conjunt, $V^*$ és el valor en $V$ de l'homfunctor contravariant de la definició~\ref{def:homfunctors}, especialitzat a $B=\mathbb K$:
\[
\Hom_{\mathbb K}(-,\mathbb K)\colon\cat{Vect}_{\mathbb K}\op\to\cat{Set}.
\]
Però aquests homconjunts tenen canònicament estructura d'espai vectorial, amb la suma i el producte per escalars punt a punt, i la precomposició amb una aplicació lineal és lineal (ho veurem tot seguit). Per tant, la mateixa assignació s'\emph{aixeca} a un functor amb valors a $\cat{Vect}_{\mathbb K}$, el \textbf{functor dual}
\[
(-)^*=\Hom_{\mathbb K}(-,\mathbb K)\colon\cat{Vect}_{\mathbb K}\op\to\cat{Vect}_{\mathbb K},
\qquad V\mapsto V^*.
\]
És el mateix patró que el functor de parts, amb $\mathbb K$ en lloc de $\{0,1\}$ com a objecte distingit. Sobre morfismes, la precomposició $\psi\mapsto\psi\comp f$ és precisament l'\textbf{aplicació dual} (o transposada) d'una aplicació lineal $f\colon V\to W$:
\[
f^*\colon W^*\to V^*,\qquad \psi\mapsto\psi\comp f.
\]
La contravariància, $(g\comp f)^*=f^*\comp g^*$, és la regla familiar de la transposició; si triem bases i representem les aplicacions per matrius, $f^*$ correspon a la matriu transposada i la inversió de l'ordre és la fórmula $(AB)^{\mathsf T}=B^{\mathsf T}A^{\mathsf T}$.

Cada $f^*$ és, en efecte, una aplicació lineal, $f^*(\psi+\lambda\psi')=(\psi+\lambda\psi')\comp f=f^*(\psi)+\lambda f^*(\psi')$, i això és el que permet l'aixecament: l'homfunctor $\cat{Set}$-valuat i el functor dual tenen la mateixa acció sobre objectes i morfismes, però el segon recorda, a més, l'estructura lineal dels valors. El mateix passa amb el covariant: $\Hom_{\mathbb K}(V,-)\colon\cat{Vect}_{\mathbb K}\to\cat{Set}$ s'aixeca a un functor
\[
\Hom_{\mathbb K}(V,-)\colon\cat{Vect}_{\mathbb K}\to\cat{Vect}_{\mathbb K}.
\]
Quan parlem d'homfunctors sense cap més precisió, ens referim sempre a la versió amb valors a $\cat{Set}$ de la definició~\ref{def:homfunctors}; és la que intervindrà en la representabilitat.

El corol·lari~\ref{prop:homfunctor-iso} es llegeix aquí com un fet conegut de l'àlgebra lineal: si $f\colon V\to W$ és un isomorfisme, aleshores la transposada $f^*\colon W^*\to V^*$ també ho és, amb inversa $(f^{-1})^*$. En particular, espais isomorfs tenen duals isomorfs.

No seguim, de moment, cap a la comparació entre $V$ i el seu dual, ni cap al bidual (o doble dual) $V^{**}$: aquesta és una qüestió de \emph{naturalitat}, que demana el llenguatge de les transformacions naturals. La reprendrem exactament on la deixem, a l'exemple~\ref{ex:bidual} del capítol~\ref{cap:transformacions}.
\end{exemple}


\subsection{El bihomfunctor}\label{subsec:bihomfunctor}

Les dues variables dels homfunctors es poden tractar alhora. Fixem-nos que $\Hom$ pren \emph{dos} objectes i en dona un conjunt, de manera contravariant en el primer i covariant en el segon; això és exactament un functor sobre un producte de categories.

\begin{definicio}[El bihomfunctor]\label{def:bihomfunctor}\index{homfunctor!bihomfunctor}\index{Hom@$\Hom$!bihomfunctor}
Sigui $\cat{C}$ localment petita. Definim \textbf{bihomfunctor} (o \textbf{bifunctor Hom})
\[
\Hom_{\cat{C}}(-,-)\colon\cat{C}\op\times\cat{C}\to\cat{Set}
\]
com el functor que envia una parella d'objectes $(A,B)$ al conjunt $\Hom(A,B)$, i una parella de morfismes $(f\colon A'\to A,\ g\colon B\to B')$ a l'aplicació
\[
\Hom(f,g)\colon\Hom(A,B)\to\Hom(A',B'),\qquad h\mapsto g\comp h\comp f.
\]
\end{definicio}

La fórmula $h\mapsto g\comp h\comp f$ combina la precomposició amb $f$ (contravariant, en la primera variable) i la postcomposició amb $g$ (covariant, en la segona): un element $h\colon A\to B$ es prolonga per l'origen amb $f$ i pel destí amb $g$.
\[
\begin{tikzpicture}[baseline=(m.center)]
  \matrix (m) [matrix of math nodes, column sep=3em] {
    A' & A & B & B' \\
  };
  \draw[-{Stealth[scale=1.1]}] (m-1-1) -- node[above] {$f$} (m-1-2);
  \draw[-{Stealth[scale=1.1]}] (m-1-2) -- node[above] {$h$} (m-1-3);
  \draw[-{Stealth[scale=1.1]}] (m-1-3) -- node[above] {$g$} (m-1-4);
  \draw[-{Stealth[scale=1.1]}] (m-1-1) to[bend right=22] node[below] {$g\comp h\comp f$} (m-1-4);
\end{tikzpicture}
\]
Fixant la primera variable en $A$ recuperem l'homfunctor covariant $\Hom(A,-)$; fixant la segona en $B$ recuperem el contravariant $\Hom(-,B)$. També sobre morfismes les notacions són coherents: prenent una de les dues components igual a una identitat,
\[
\Hom(\id_A,g)=g_*=\Hom(A,g),
\qquad
\Hom(f,\id_B)=f^*=\Hom(f,B),
\]
de manera que la notació~\ref{not:hom-morfismes} és un cas particular de la del bihomfunctor. Ara bé, la functorialitat en cada variable per separat no basta: cal que les dues accions siguin compatibles, és a dir, que $\Hom(-,-)$ preservi la composició i les identitats de $\cat{C}\op\times\cat{C}$. També això és conseqüència de l'associativitat i de les identitats: aplicar primer $(f,g)$ i després $(f',g')$ envia $h$ a $g'\comp g\comp h\comp f\comp f'$, que és l'acció de la parella composta $(f\comp f',\,g'\comp g)$, amb la primera component en l'ordre invertit, com correspon a $\cat{C}\op$. La comprovació detallada és l'exercici resolt~\ref{pr-bihom} de la Pràctica.

\begin{exemple}[El bihomfunctor de $\cat{Set}$]\label{ex:bihom-set}\index{homfunctor!bihomfunctor!a Set}
A $\cat{Set}$, el bihomfunctor envia una parella de conjunts $(X,Y)$ al conjunt d'aplicacions
\[
\Hom_{\cat{Set}}(X,Y)=Y^X,
\]
amb la mateixa notació exponencial de l'inici del capítol. L'acció sobre una parella de morfismes $(f\colon X'\to X,\ g\colon Y\to Y')$ és el \enquote{canvi de variables} a l'entrada i a la sortida: donada una aplicació $h\colon X\to Y$, primer la reindexem pel domini amb $f$ i després en traslladem els valors amb $g$,
\[
\Hom(f,g)\colon Y^X\to (Y')^{X'},\qquad h\mapsto g\comp h\comp f.
\]
Fixant el primer conjunt en un punt $\{\ast\}$ recuperem $\Hom(\{\ast\},Y)\cong Y$ ---els elements de $Y$--- amb l'acció covariant habitual per postcomposició; fixant el segon en $\{0,1\}$ recuperem $\Hom(X,\{0,1\})\cong\mathcal P(X)$, és a dir, el functor de parts de l'exemple~\ref{ex:functor-parts}, llevat de la identificació canònica. El bihomfunctor és, doncs, la construcció que engloba tots dos alhora.
\end{exemple}

\begin{exemple}[El bihomfunctor de $\cat{Vect}_{\mathbb K}$]\label{ex:bihom-vect}\index{homfunctor!bihomfunctor!a Vect}
El bihomfunctor ordinari de $\cat{Vect}_{\mathbb K}$ és
\[
\Hom(-,-)\colon\cat{Vect}_{\mathbb K}\op\times\cat{Vect}_{\mathbb K}\to\cat{Set},
\]
que envia una parella d'espais $(V,W)$ al conjunt $\Hom_{\mathbb K}(V,W)$ de les aplicacions lineals de $V$ a $W$. Però aquest conjunt és canònicament un espai vectorial, i sobre una parella $(f\colon V'\to V,\ g\colon W\to W')$ l'acció $h\mapsto g\comp h\comp f$ combina la precomposició amb $f$ i la postcomposició amb $g$, totes dues lineals. Per això el bihomfunctor admet un aixecament canònic a $\cat{Vect}_{\mathbb K}$:
\[
\Hom_{\mathbb K}(-,-)\colon\cat{Vect}_{\mathbb K}\op\times\cat{Vect}_{\mathbb K}\to\cat{Vect}_{\mathbb K}.
\]
Fixant la segona variable en el cos, aquest aixecament dona el functor dual $(-)^*$ de l'exemple~\ref{ex:homfunctor-dual}. Fixant la primera, $\Hom_{\mathbb K}(\mathbb K,-)$ correspon canònicament al functor identitat. En efecte, una aplicació lineal $\mathbb K\to W$ queda determinada per la imatge de $1$, i això dona una bijecció lineal canònica
\[
\Hom_{\mathbb K}(\mathbb K,W)\;\cong\;W;
\]
al capítol~\ref{cap:transformacions} la reconeixerem com un isomorfisme natural. Les dues cares del dual i la identitat es reuneixen, així, en un sol functor de dues variables.
\end{exemple}

\section{Functors fidels i plens}

% CONVENCIÓ D'ESTIL (aplicada a tot el llibre): s'eviten les nominalitzacions
% «fidelitat» i «plenitud» en la prosa expositiva; es prediquen com «ser fidel» /
% «ser ple». Propagat a teoria, exercicis i a la prosa de practica/tests.
% Excepcions conservades a propòsit: les etiquetes de pas en demostracions
% «(fidelitat)» / «per fidelitat» (caps/practica/*, caps/tests/*), i l'ús ordinari
% «per fidelitat a la bibliografia» del glossari (caps/glossari.tex).

Un functor actua sobre els morfismes; és natural preguntar-se com és aquesta acció entre cada parella d'objectes. Fixats $A,B$ a $\cat{C}$, un functor $F\colon\cat{C}\to\cat{D}$ indueix una aplicació
\[
F_{A,B}\colon\Hom_{\cat{C}}(A,B)\to\Hom_{\cat{D}}(F(A),F(B)),
\qquad
f\mapsto F(f).
\]
Les propietats d'aquesta aplicació donen lloc a una classificació bàsica.

\begin{definicio}\index{functor!fidel}\index{functor!ple}
\label{def:ple-fidel}
Un functor $F\colon\cat{C}\to\cat{D}$ és:
\begin{itemize}
\item \textbf{fidel} si cada aplicació $F_{A,B}$ és injectiva; és a dir, si $F(f)=F(g)$ amb $f,g\colon A\to B$ implica $f=g$;
\item \textbf{ple} si cada aplicació $F_{A,B}$ és exhaustiva; és a dir, si tot morfisme $F(A)\to F(B)$ de $\cat{D}$ és de la forma $F(f)$ per a algun $f\colon A\to B$ de $\cat{C}$;
\item \textbf{ple i fidel} si cada $F_{A,B}$ és bijectiva.
\end{itemize}
\end{definicio}

\begin{observacio}
Cal no confondre aquestes condicions amb la injectivitat o l'exhaustivitat \emph{sobre objectes}. Un functor pot ser ple i fidel sense ser injectiu sobre objectes: ser ple i fidel només afecta els morfismes entre cada parella fixada d'objectes. Un functor ple i fidel identifica $\Hom_{\cat{C}}(A,B)$ amb $\Hom_{\cat{D}}(F(A),F(B))$, però pot enviar objectes diferents a un mateix objecte.
\end{observacio}

\begin{observacio}
Aquestes propietats es transporten al functor oposat. Recordem (definició~\ref{def:functor-oposat}) que el functor $F\op\colon\cat{C}\op\to\cat{D}\op$ actua igual que $F$ sobre objectes i morfismes. Cal no confondre el \emph{functor} $F\op$ amb l'\emph{aplicació entre homconjunts} que indueix. Fixats dos objectes $A,B$, que $F$ sigui ple i fidel es mesura amb l'aplicació
\[
F_{A,B}\colon\Hom_{\cat{C}}(A,B)\to\Hom_{\cat{D}}(F(A),F(B)),\qquad f\mapsto F(f),
\]
i les de $F\op$, amb l'aplicació anàloga $(F\op)_{B,A}$ ---amb la parella d'objectes en l'ordre invers, perquè a $\cat{C}\op$ una fletxa $A\to B$ és una fletxa $B\to A$ de $\cat{C}$---
\[
(F\op)_{B,A}\colon\Hom_{\cat{C}\op}(B,A)\to\Hom_{\cat{D}\op}(F(B),F(A)).
\]
Ara bé, per definició d'oposada $\Hom_{\cat{C}\op}(B,A)=\Hom_{\cat{C}}(A,B)$ i $\Hom_{\cat{D}\op}(F(B),F(A))=\Hom_{\cat{D}}(F(A),F(B))$; i sobre aquests conjunts $(F\op)_{B,A}$ envia cada $f$ a $F(f)$, exactament com $F_{A,B}$. Les dues aplicacions són, doncs, literalment la mateixa. Com que \enquote{ser injectiva} i \enquote{ser exhaustiva} són propietats d'aquesta aplicació, $F$ és fidel (respectivament ple) si i només si $F\op$ ho és.
\end{observacio}

\begin{exemple}
El functor oblit $U\colon\cat{Grp}\to\cat{Set}$ és fidel: dos homomorfismes diferents segueixen sent aplicacions diferents, de manera que $U$ no identifica morfismes. En canvi no és ple: entre dos grups hi ha aplicacions de conjunts que no són homomorfismes, i aquestes no provenen de cap morfisme de $\cat{Grp}$. Aquesta és la situació típica dels functors oblit: retenen prou informació per distingir morfismes, però obren la porta a fletxes noves que no respecten l'estructura.
\end{exemple}

Les subcategories, introduïdes elementalment a la subsecció~\ref{subsec:subcategories}, tenen una lectura natural en el llenguatge dels functors. Si $\cat{S}$ és una subcategoria de $\cat{C}$, el \textbf{functor inclusió} $\iota\colon\cat{S}\hookrightarrow\cat{C}$ ---que envia cada objecte i cada morfisme de $\cat{S}$ a ell mateix a $\cat{C}$--- sempre és \textbf{fidel}, perquè no identifica mai dos morfismes diferents. Les dues menes de subcategoria que vam distingir ---la \textbf{plena} i l'\textbf{ampla}--- es tradueixen, cadascuna, en una propietat d'aquest functor. D'una banda, $\cat{S}$ és una \textbf{subcategoria plena} exactament quan $\iota$ és \textbf{ple}: que $\iota$ sigui ple vol dir que entre dos objectes de $\cat{S}$ hi ha \emph{tots} els morfismes que hi ha a $\cat{C}$, que és justament la definició de subcategoria plena. De l'altra, $\cat{S}$ és una \textbf{subcategoria ampla} exactament quan l'acció de $\iota$ \emph{sobre els objectes} és exhaustiva, és a dir, quan $\Ob(\cat{S})=\Ob(\cat{C})$: la subcategoria conté tots els objectes de $\cat{C}$. (Parlem aquí de l'aplicació $\Ob(\cat{S})\to\Ob(\cat{C})$ induïda per $\iota$, no del functor mateix; per al functor inclusió aquesta aplicació és sempre injectiva, i ser ampla equival que sigui, a més, exhaustiva.)\index{subcategoria!plena}\index{subcategoria!ampla}\index{functor!inclusió}

\begin{exemple}\label{ex:inclusio-finset-plena}
El \textbf{functor inclusió} $\cat{FinSet}\hookrightarrow\cat{Set}$ dels conjunts finits és ple i fidel: entre dos conjunts finits hi ha exactament les mateixes aplicacions que a $\cat{Set}$, de manera que un cop triats els objectes no descartem cap morfisme. En canvi, el functor inclusió de la subcategoria de $\cat{Set}$ que conserva tots els conjunts com a objectes però només les aplicacions \emph{injectives} com a morfismes és fidel, però \emph{no} ple: entre dos conjunts hi ha aplicacions no injectives, que per tant no provenen de cap morfisme de la subcategoria. El mateix passa amb la subcategoria ampla de les aplicacions \emph{exhaustives}: el seu functor inclusió és fidel però no ple, perquè deixa fora les aplicacions que no són exhaustives.
\end{exemple}

\begin{proposicio}\label{prop:plenfidel-reflecteix-iso}\index{functor!ple i fidel}
Sigui $F\colon\cat{C}\to\cat{D}$ un functor ple i fidel, i sigui $f\colon A\to B$ un morfisme de $\cat{C}$. Aleshores
\[
f\ \text{és un isomorfisme a }\cat{C}
\quad\Longleftrightarrow\quad
F(f)\ \text{és un isomorfisme a }\cat{D}.
\]
\end{proposicio}

\begin{prova}
La implicació cap a la dreta val per a qualsevol functor: és la proposició~\ref{prop:functor-preserva-iso}, i no fa servir que $F$ sigui ple ni fidel. Vegem el recíproc, que sí que ho necessita. Suposem que $F(f)\colon F(A)\to F(B)$ té un invers $h\colon F(B)\to F(A)$ a $\cat{D}$. Com que $F$ és ple, existeix $g\colon B\to A$ a $\cat{C}$ amb $F(g)=h$. Aleshores
\[
F(g\comp f)=F(g)\comp F(f)=h\comp F(f)=\id_{F(A)}=F(\id_A),
\]
i com que $F$ és fidel, $g\comp f=\id_A$. Anàlogament $f\comp g=\id_B$. Per tant $f$ és un isomorfisme amb invers $g$.
\end{prova}

\begin{observacio}
Les dues direccions tenen estatus diferent. La directa ---preservar isomorfismes--- val per a tot functor; el que caracteritza els plens i fidels és el recíproc: \emph{reflecteixen} isomorfismes, i no només els preserven. Un functor ple i fidel reprodueix bijectivament els conjunts de morfismes entre els objectes de la seva imatge, però això \emph{no} implica que sigui injectiu sobre objectes: pot enviar objectes diferents a un mateix objecte. Aquesta situació ---morfismes identificats bijectivament, objectes no necessàriament--- conduirà més endavant a la noció d'\emph{equivalència de categories}.
\end{observacio}

Ser ple i fidel descriu com actua un functor sobre els morfismes. Una tercera condició, sobre els objectes, completa la terna que caracteritzarà les equivalències.

\begin{definicio}[Functor essencialment exhaustiu]\label{def:ess-exhaustiu}\index{functor!essencialment exhaustiu}
Un functor $F\colon\cat{C}\to\cat{D}$ és \textbf{essencialment exhaustiu} si tot objecte de $\cat{D}$ és isomorf a la imatge d'algun objecte de $\cat{C}$:
\[
\forall D\in\cat{D},\ \exists A\in\cat{C}\ \text{tal que}\ F(A)\cong D.
\]
\end{definicio}

\begin{observacio}
La condició demana \emph{isomorfisme}, no igualtat: no cal que tot objecte de $\cat{D}$ sigui exactament una imatge, només que ho sigui llevat d'isomorfisme. Un functor alhora ple, fidel i essencialment exhaustiu recull bijectivament els morfismes i, a més, arriba a tots els objectes llevat d'isomorfisme; aquestes són precisament les condicions que, un cop disposem de les transformacions naturals, caracteritzaran les \emph{equivalències de categories}.
\end{observacio}

\begin{exemple}[Un functor essencialment exhaustiu que no és exhaustiu sobre objectes]\label{ex:ess-exhaustiu-finset}\index{functor!essencialment exhaustiu}
Considerem la subcategoria plena $\cat{S}$ de $\cat{FinSet}$ que té per objectes només els conjunts
\[
\underline{0}=\varnothing,\qquad \underline{n}=\{1,\dots,n\}\quad(n\geq 1)
\]
---un de sol de cada mida, en lloc de tots els conjunts finits---, i el seu functor inclusió $\iota\colon\cat{S}\hookrightarrow\cat{FinSet}$. Aquest functor \emph{és} essencialment exhaustiu: tot conjunt finit $X$ té algun cardinal $n$, i aleshores és isomorf (bijectiu) a $\underline{n}=\iota(\underline{n})$. En canvi, $\iota$ no és exhaustiu sobre objectes: el conjunt $\{a,b\}$, posem per cas, no és \emph{igual} a cap $\underline{n}$, tot i ser isomorf a $\underline{2}$. Aquest és, doncs, l'exemple prototípic de la diferència entre \enquote{arribar a tots els objectes} i \enquote{arribar-hi llevat d'isomorfisme}: $\iota$ deixa fora molts objectes com a imatges exactes, però cap classe d'isomorfisme. L'exercici resolt~\ref{pr-esquelet-finset} de la Pràctica comprova, a més, que $\iota$ és ple i fidel.
\end{exemple}

\begin{exemple}[Un functor que no és essencialment exhaustiu]\index{functor!essencialment exhaustiu}
En canvi, la inclusió $\cat{FinSet}\hookrightarrow\cat{Set}$ \emph{no} és essencialment exhaustiva: un conjunt infinit no és isomorf a cap conjunt finit, de manera que cap objecte finit de la imatge no hi pot arribar. Aquí falla, doncs, la tercera condició, tot i que aquest functor sí que és ple i fidel (exemple~\ref{ex:inclusio-finset-plena}). Ser ple i fidel no diu res sobre quins objectes s'assoleixen: són condicions independents.
\end{exemple}

\subsection{Què conserva un functor: un resum}\label{subsec:que-conserva}

Convé aclarir, abans de tancar el capítol, quines propietats respecta un functor \emph{qualsevol} i quines en demanen més. Un functor transporta qualsevol configuració finita de morfismes i les equacions de composició que aquests morfismes satisfan. Per això preserva, per exemple, isomorfismes, seccions i retraccions: la propietat està testimoniada per morfismes de la categoria d'origen ---l'invers, la secció o la retracció--- i per una equació com $g\comp f=\id$, i el functor transporta tant el testimoni com l'equació. En canvi, una propietat com ser mono o epi quantifica sobre totes les fletxes possibles d'un cert tipus (\enquote{per a tota parella $g,h$}), incloses fletxes de la categoria d'arribada que poden no provenir de la categoria d'origen; i és aquesta quantificació la que un functor no pot garantir.

\begin{itemize}
\item Un functor \emph{qualsevol} preserva composició i identitats, diagrames commutatius, isomorfismes, i seccions i retraccions (és a dir, els epimorfismes i monomorfismes \emph{escindits}). En tots els casos, la propietat està testimoniada per morfismes i equacions de la categoria d'origen, que el functor transporta.
\item Ja un functor \emph{fidel} ---sense necessitat de ser ple--- \emph{reflecteix} monomorfismes i epimorfismes: si $F(f)$ és mono, aleshores $f$ ja ho era, perquè de $f\comp g=f\comp h$ se'n dedueix $F(f)\comp F(g)=F(f)\comp F(h)$, d'on $F(g)=F(h)$ i, com que $F$ és fidel, $g=h$; l'argument per als epimorfismes és dual.
\item Un functor \emph{ple i fidel} reflecteix, a més, isomorfismes, seccions i retraccions: si la imatge d'un morfisme és un isomorfisme, una secció o una retracció, el morfisme original ja ho era. Aquí calen totes dues hipòtesis: ser ple permet aixecar l'invers ---unilateral o bilateral--- que \emph{a priori} només viu al codomini, i ser fidel permet reflectir l'equació que el defineix. Recupera, doncs, aquestes nocions a banda i banda.
\item Al capítol~\ref{cap:transformacions} veurem que un functor \emph{ple, fidel i essencialment exhaustiu} determina una \emph{equivalència de categories} (en general, amb l'ajut de l'axioma de l'elecció). Els capítols posteriors precisaran, per a cada construcció, quines propietats es preserven i es reflecteixen sota equivalència.

\end{itemize}

En canvi, un functor qualsevol \emph{no} preserva, en general, ni els monomorfismes ni els epimorfismes ni les propietats universals. Vegem-ho amb els epimorfismes.

\begin{exemple}[Un functor que no preserva epimorfismes]\index{functor!no preserva epimorfismes}
Prenem el monoide additiu $(\mathbb{N},+,0)$ i el monoide $(\{0,1\},\vee,0)$, on $\vee$ és el màxim (o disjunció), tots dos vistos com a categories d'un sol objecte, $\cat{B}\mathbb{N}$ i $\cat{B}M$. L'aplicació $h\colon\mathbb{N}\to\{0,1\}$ amb $h(0)=0$ i $h(n)=1$ per a $n\geq 1$ és un homomorfisme de monoides, i defineix per tant un functor $F\colon\cat{B}\mathbb{N}\to\cat{B}M$ (exemple~\ref{ex:functors-monoide}).

A $\cat{B}\mathbb{N}$, el morfisme $1$ és un epimorfisme: la cancel·lació per la dreta a $\cat{B}\mathbb{N}$ és la cancel·lació de la suma, i a $(\mathbb{N},+)$ es pot cancel·lar. La seva imatge $F(1)=1$, en canvi, \emph{no} és un epimorfisme a $\cat{B}M$, perquè $1\vee 0=1=1\vee 1$ amb $0\neq 1$: la cancel·lació falla. Així, $F$ envia un epimorfisme a un morfisme que no ho és.

L'arrel del fenomen és la que hem assenyalat: \enquote{ser epi} quantifica sobre totes les parelles $g,h$ de la categoria d'arribada, i el functor no controla les parelles noves que apareixen a $\cat{B}M$. Un epimorfisme \emph{escindit}, en canvi ---testimoniat per una secció $s$ amb $f\comp s=\id$---, sí que es preservaria, perquè el functor transporta la secció i l'equació.
\end{exemple}

\begin{resumcap}
\apartat{Idees centrals}
\begin{enumerate}
  \item Un functor preserva dominis, codominis, composició i identitats.
  \item Les categories petites i els functors formen una categoria, $\cat{Cat}$, que és gran però localment petita; els functors entre categories grans són igualment legítims, tot i que no són fletxes de $\cat{Cat}$.
  \item El graf subjacent i la categoria lliure formen una parella de functors en sentits contraris, $\Free\colon\cat{Graph}\to\cat{Cat}$ i $U\colon\cat{Cat}\to\cat{Graph}$. No són inversos l'un de l'altre, però estan lligats per l'avaluació $\Free(U(\cat{C}))\to\cat{C}$; aquesta relació reapareixerà més endavant com a adjunció.
  \item Un functor contravariant és un functor des de l'oposada. Per comprovar qualsevol functor, primer es verifiquen els tipus ($F(f)\colon F(A)\to F(B)$, o $F(B)\to F(A)$ en el cas contravariant) i després les identitats; el functor de parts té una versió de cada variància.
  \item Ser fidel o ser ple és una propietat de les aplicacions $F_{A,B}$ entre homconjunts, i no diu res de l'acció sobre objectes; afegint-hi l'exhaustivitat essencial, que sí que parla dels objectes (llevat d'isomorfisme), s'obtenen les equivalències.
  \item L'homfunctor covariant $\Hom(A,-)$ i el contravariant $\Hom(-,B)$ transformen isos en bijeccions, i el bihomfunctor $\Hom(-,-)$ els unifica en un functor de dues variables; tots tres són fonamentals i els retrobarem al lema de Yoneda.
\end{enumerate}
\apartat{Errors típics} Confondre \enquote{isomorfisme de categories} amb \enquote{equivalència}; creure que \enquote{fidel} vol dir injectiu sobre objectes o que \enquote{ple} vol dir exhaustiu sobre objectes, quan totes dues nocions parlen només dels homconjunts; comprovar les equacions functorials abans d'haver comprovat que $F(f)$ té el tipus correcte; oblidar la inversió de sentit en un functor contravariant; barrejar $\cat{C}/A$ (sobre $A$) amb $A/\cat{C}$ (sota $A$); i confondre \emph{preservar} amb \emph{reflectir} una propietat: tot functor preserva els isomorfismes; els functors plens i fidels també els reflecteixen, tot i que aquesta condició és suficient i no necessària. Un functor qualsevol no preserva, en general, ni els monomorfismes ni els epimorfismes.
\apartat{Lectura recomanada} \cite[part~III, sessions~11--17]{lawvereschanuel} per a exemples molt elementals d'aplicacions que preserven estructura entre grafs i entre sistemes dinàmics; \cite[§1.3 i §§7.1--7.2]{awodey} per a la definició de functor, la categoria de categories i l'estructura representable; \cite[I.3, II.2 i II.5]{maclane} per a la formulació clàssica, inclosos els functors contravariants; \cite[§1.2]{leinster} per als functors i les seves propietats, inclosos els contravariants i els plens i fidels; \cite[cap.~1]{riehl} per als homfunctors i el bifunctor $\Hom$; \cite[§§3.1--3.3]{barrwells} per a més exemples de functors, en particular les accions de monoides i els functors entre grafs, i per a la classificació dels functors en fidels, plens i altres tipus; \cite[§1.3 i §§1.5.1--1.5.3]{perrone-notes} per a una segona explicació de la functorialitat, amb una secció sobre què \emph{no} és un functor i una altra sobre functors, monos i epis.
\end{resumcap}

\chapter[Transformacions naturals]{\texorpdfstring{Transformacions naturals\\ i equivalències}{Transformacions naturals i equivalències}}
\label{cap:transformacions}

\begin{objectius}
\apartat{Prerequisits} Categories i diagrames commutatius (cap.~\ref{cap:categories}); functors, composició de functors, functors covariants i contravariants, i les nocions de functor ple, fidel i essencialment exhaustiu (cap.~\ref{cap:functors}). Els homfunctors intervenen en alguns exemples i exercicis, i seran importants en els capítols següents, però no són necessaris per a la definició de naturalitat. Alguns exemples opcionals fan servir una familiaritat bàsica amb espais vectorials, duals i biduals, i l'exemple d'interpretació de fórmules de primer ordre, marcat com a ampliació, pressuposa la semàntica d'aquesta lògica.
\apartat{Objectius} En acabar el capítol, el lector hauria de poder:
\begin{itemize}
  \item definir una transformació natural i comprovar-ne la naturalitat mitjançant quadrats commutatius;
  \item compondre transformacions verticalment i interpretar-les com a morfismes de la categoria de functors $\cat{D}^{\cat{C}}$;
  \item reconèixer un isomorfisme natural i distingir-lo d'una família d'isomorfismes que no sigui natural;
  \item treballar amb la composició horitzontal i distingir-la de la vertical;
  \item definir una equivalència de categories mitjançant un quasiinvers, i caracteritzar-la com un functor ple, fidel i essencialment exhaustiu (amb l'axioma de l'elecció per a la implicació recíproca).
\end{itemize}
\end{objectius}

\section{Cap al concepte de transformació natural}\label{sec:cap-transformacio}

Al capítol anterior vam veure que els functors són els morfismes entre categories: relacionen dues categories respectant-ne tota l'estructura. Ara bé, sovint tenim \emph{dos} functors $F,G\colon\cat{C}\to\cat{D}$ que van de la mateixa categoria a la mateixa categoria, i volem comparar-los. La pregunta natural és, doncs, si hi ha una noció de \emph{morfisme entre functors}.

Aquesta idea ---comparar functors de manera coherent--- va ser històricament una de les motivacions originals per introduir les categories. De fet, es diu sovint que les categories es van definir per poder definir els functors, i els functors per poder definir les transformacions naturals; totes tres nocions apareixen alhora, per primera vegada, a l'article fundacional d'Eilenberg i Mac~Lane~\cite{eilenbergmaclane}, que precisament les motiva a partir de la relació entre un espai vectorial i el seu bidual. El mot \enquote{natural} té aquí un sentit tècnic precís que farem explícit: una comparació és natural quan fa commutar els quadrats corresponents a tots els morfismes. Sovint s'expressa dient que \enquote{no depèn de cap elecció arbitrària}; convé retenir, però, que es tracta d'una intuïció orientadora i no d'una definició: una construcció pot haver requerit eleccions i, tanmateix, resultar natural. La condició precisa és sempre la commutativitat dels quadrats.

Abans de formalitzar-ho, val la pena veure un exemple que fa palès aquest \enquote{principi de coherència}, aquest cop entre \emph{sintaxi} i \emph{semàntica}. Treballem amb la lògica proposicional. Prenem com a categoria de base $\cat{FinSet}$, els conjunts finits amb les aplicacions entre ells (subsecció~\ref{subsec:conjunts}); ara en pensem els objectes com a conjunts de variables proposicionals, i els morfismes $f\colon X\to Y$ com a \emph{canvis de variables}, que a cada variable de $X$ n'assignen una de $Y$.

El primer functor, el \emph{sintàctic}, associa a cada conjunt de variables el conjunt de les fórmules que s'hi poden escriure:
\[
F(X)=\{\text{fórmules proposicionals amb variables de }X\},
\]
construïdes a partir de les variables de $X$ amb les connectives habituals. Sobre morfismes, un canvi de variables $f\colon X\to Y$ indueix l'aplicació de \emph{substitució}
\[
F(f)\colon F(X)\longrightarrow F(Y),
\qquad
F(f)(\varphi)=\varphi[f],
\]
on $\varphi[f]$ és la fórmula obtinguda reemplaçant dins $\varphi$ cada variable $x$ per la variable $f(x)$. Aquesta substitució és la mateixa operació que ja vam trobar a l'exemple~\ref{ex:deduibilitat-functors}, on una substitució de variables proposicionals induïa un functor entre categories de deduïbilitat perquè transformava deduccions; ara la mirem d'una altra manera, com l'acció sobre morfismes d'un functor $F\colon\cat{FinSet}\to\cat{Set}$. Formalment, $F(f)$ es defineix per recursió sobre l'estructura de $\varphi$: sobre una variable, $F(f)(x)=f(x)$, i $F(f)$ commuta amb les connectives.

El segon functor, el \emph{semàntic}, associa a cada conjunt de variables el conjunt de les funcions de veritat corresponents:
\[
G(X)=\{\text{funcions de veritat sobre }X\},
\]
on una \emph{funció de veritat} sobre $X$ és una aplicació
\[
\varphi\colon \{0,1\}^{X}\longrightarrow\{0,1\}
\]
que assigna un valor de veritat a cada valoració de les variables, entenent per \emph{valoració} un element $v\in\{0,1\}^{X}$, és a dir una aplicació $v\colon X\to\{0,1\}$ que fixa el valor de cada variable. Dit altrament, $G(X)=\{0,1\}^{(\{0,1\}^{X})}$. Sobre morfismes, un canvi de variables $f\colon X\to Y$ actua per \emph{precomposició en les valoracions}: tota valoració $w\colon Y\to\{0,1\}$ es restringeix a la valoració $w\comp f\colon X\to\{0,1\}$, i definim
\[
G(f)\colon G(X)\longrightarrow G(Y),
\qquad
G(f)(\theta)(w)=\theta(w\comp f),
\]
és a dir, $G(f)(\theta)=\theta\comp(-\comp f)$.

Tots dos són functors covariants $\cat{FinSet}\to\cat{Set}$. La comprovació és rutinària ---per recursió sobre les fórmules en el cas de $F$, i per l'associativitat de la composició en el de $G$--- i es fa a l'exercici resolt~\ref{pr-sintaxi-semantica} de la Pràctica.

A cada conjunt de variables $X$ li associem l'aplicació que interpreta cada fórmula com la seva funció de veritat (la seva taula de veritat):
\[
\eta_X\colon F(X)\longrightarrow G(X),
\qquad
\eta_X(\varphi)(v)=
\begin{cases}
1 & \text{si }v\models\varphi,\\
0 & \text{si }v\not\models\varphi,
\end{cases}
\]
on $v\models\varphi$ vol dir que la valoració $v$ satisfà la fórmula $\varphi$.

El fet clau és que aquesta interpretació és \emph{compatible amb el canvi de variables}: per a tot morfisme $f\colon X\to Y$ de $\cat{FinSet}$, el quadrat
\[
\begin{tikzpicture}[baseline=(m.center)]
  \matrix (m) [matrix of math nodes, row sep=2.8em, column sep=3.6em] {
    F(X) & F(Y) \\
    G(X) & G(Y) \\
  };
  \draw[-{Stealth[scale=1.1]}] (m-1-1) -- node[above] {$F(f)$} (m-1-2);
  \draw[-{Stealth[scale=1.1]}] (m-1-1) -- node[left] {$\eta_X$} (m-2-1);
  \draw[-{Stealth[scale=1.1]}] (m-1-2) -- node[right] {$\eta_Y$} (m-2-2);
  \draw[-{Stealth[scale=1.1]}] (m-2-1) -- node[below] {$G(f)$} (m-2-2);
\end{tikzpicture}
\]
commuta, és a dir $G(f)\comp\eta_X=\eta_Y\comp F(f)$. Avaluats sobre una fórmula $\varphi\in F(X)$ i una valoració $w\colon Y\to\{0,1\}$, els dos camins del quadrat diuen si $w\models\varphi[f]$ i si $(w\comp f)\models\varphi$, respectivament. Que coincideixin és exactament el \emph{lema de substitució} de la lògica proposicional,
\[
(w\comp f)\models\varphi
\quad\Longleftrightarrow\quad
w\models\varphi[f],
\]
que es demostra per recursió sobre $\varphi$; el càlcul complet és a l'exercici resolt~\ref{pr-sintaxi-semantica}.

En termes informals: substituir primer i interpretar després dona la mateixa funció de veritat que interpretar primer i transportar després. Per exemple, de $\varphi=p\land q$ i el canvi de variables $p\mapsto r,\ q\mapsto s$ obtenim $r\land s$, i tant se val si calculem la taula de veritat abans o després de substituir: el resultat és la mateixa funció (la que val $1$ només quan $r$ i $s$ són totes dues certes). Si el canvi identifica variables, diguem $p\mapsto r,\ q\mapsto r$, aleshores $\varphi$ es transforma en $r\land r$, amb taula de veritat la de $r$; també aquí interpretar i transportar coincideixen. Aquesta compatibilitat, per a tot canvi de variables alhora, és el que a la secció següent anomenarem la \textbf{naturalitat} de la família $(\eta_X)$: \emph{substituir variables i després calcular la funció de veritat és el mateix que calcular-la primer i transportar-la després al llarg del canvi de variables}.

Val la pena notar una diferència amb els exemples \enquote{d'identificació canònica} que veurem després (com el bidual d'un espai vectorial): aquí $\eta_X$ \emph{no} és injectiva ---dues fórmules diferents poden tenir la mateixa taula de veritat, com $p\lor q$ i $q\lor p$---, de manera que no estableix cap isomorfisme entre sintaxi i semàntica. El que expressa la naturalitat no és una identificació, sinó que la interpretació és compatible amb \emph{qualsevol} substitució o canvi de variables, incloses les substitucions que identifiquen variables. La mateixa idea s'estén, amb més maquinària, a la lògica de primer ordre, on les fórmules s'interpreten com els predicats que defineixen sobre una estructura fixada; ho veurem desplegat com a exemple formal a l'exemple~\ref{ex:interpretacio-primer-ordre} de la secció~\ref{sec:exemples-transf}.

\begin{observacio}[Per què estudiar lògica categòrica]
Un lector amb formació lògica pot objectar que, al capdavall, la feina l'ha feta el lema de substitució, i preguntar-se què hi guanya el llenguatge categòric. La resposta assenyala un canvi de punt de vista. Aquí hem hagut de definir la substitució per recursió i comprovar-ne la compatibilitat a mà perquè encara no disposem de cap maquinària; però la lògica categòrica reformula la substitució com una operació \emph{estructural} ---la composició, o més en general la reindexació al llarg d'un morfisme de la categoria de contextos--- i no com una manipulació sintàctica de fórmules. En aquest marc, la interpretació és un functor (o un morfisme entre estructures més riques) que \emph{preserva} aquella operació, i naturalitats com la d'aquest exemple deixen de demostrar-se cas per cas: s'obtenen com a corol·lari que la interpretació respecta l'estructura. La inducció sobre fórmules es fa aleshores una sola vegada i en abstracte, en establir que la categoria sintàctica té la propietat universal que la caracteritza. És el gir que fa que la lògica moderna no es defineixi per les seves peculiaritats sintàctiques, sinó per les estructures que preserva; el nostre lema de substitució n'és el cas visible més senzill.
\end{observacio}

\section{Definició de transformació natural}

Siguin $F,G\colon\cat{C}\to\cat{D}$ dos functors. Volem, per a cada objecte $A$ de $\cat{C}$, un morfisme de $\cat{D}$ que vagi de $F(A)$ a $G(A)$; i volem que aquests morfismes siguin compatibles amb l'acció de $F$ i $G$ sobre les fletxes.

\begin{definicio}\index{transformació natural}
\label{def:transf-natural}
Una \textbf{transformació natural} $\eta\colon F\Rightarrow G$ entre dos functors $F,G\colon\cat{C}\to\cat{D}$ és una \textbf{família de morfismes de $\cat{D}$}
\[
\eta=\bigl(\eta_A\bigr)_{A\in\Ob(\cat{C})},
\qquad
\eta_A\colon F(A)\to G(A),
\]
indexada pels objectes de $\cat{C}$ ---un morfisme $\eta_A$ per a cada objecte $A$, de $F(A)$ a $G(A)$, anomenat el \textbf{component} de $\eta$ a $A$---, subjecta a la condició següent: per a cada morfisme $f\colon A\to B$ de $\cat{C}$, el quadrat
\[
\begin{tikzpicture}[baseline=(m.center)]
  \matrix (m) [matrix of math nodes, row sep=2.8em, column sep=3.4em] {
    F(A) & F(B) \\
    G(A) & G(B) \\
  };
  \draw[-{Stealth[scale=1.1]}] (m-1-1) -- node[above] {$F(f)$} (m-1-2);
  \draw[-{Stealth[scale=1.1]}] (m-1-1) -- node[left] {$\eta_A$} (m-2-1);
  \draw[-{Stealth[scale=1.1]}] (m-1-2) -- node[right] {$\eta_B$} (m-2-2);
  \draw[-{Stealth[scale=1.1]}] (m-2-1) -- node[below] {$G(f)$} (m-2-2);
\end{tikzpicture}
\]
commuta; és a dir,
\[
G(f)\comp\eta_A=\eta_B\comp F(f).
\]
Aquesta igualtat s'anomena \textbf{condició de naturalitat}, i el quadrat, \textbf{quadrat de naturalitat}.
\end{definicio}

El quadrat de naturalitat exhibeix una transformació natural morfisme a morfisme, i és la manera pràctica de comprovar-la. Sovint, però, interessa representar tota la transformació $\eta$ d'un sol traç, com un objecte que \emph{relaciona} dos functors paral·lels. Dibuixarem, doncs, una transformació natural \emph{diagramàticament} d'aquesta manera: els dos functors $F,G\colon\cat{C}\to\cat{D}$ es representen com dues fletxes de $\cat{C}$ cap a $\cat{D}$, i la transformació natural com una fletxa doble $\Rightarrow$ que va de l'una cap a l'altra:
\[
\begin{tikzpicture}[baseline=(A.base)]
  \node (A) at (0,0) {$\cat{C}$};
  \node (B) at (2.9,0) {$\cat{D}$};
  \draw[cat] (A) to[bend left=26] node[above] {$F$} (B);
  \draw[cat] (A) to[bend right=26] node[below] {$G$} (B);
  \draw[cat nat] (1.45,0.2) -- node[right=1.5pt] {$\eta$} (1.45,-0.2);
\end{tikzpicture}
\]
És la primera vegada que apareix una fletxa doble entre fletxes: reflecteix que $\eta$ viu \enquote{un nivell més amunt} que els morfismes ordinaris. Aquest punt de vista ---objectes, functors i, damunt, transformacions naturals--- és el que més endavant resumirem dient que les categories formen una \emph{2-categoria}, i és la notació natural per a la composició horitzontal i per a la composició d'una transformació natural amb un functor que veurem al final del capítol.

Convé precisar bé què exigeix la condició de naturalitat i què no: les components $\eta_A$ es trien un per a cada objecte, però no lliurement, ja que han d'encaixar tots alhora amb l'acció dels functors sobre els morfismes. Ara bé, la naturalitat només demana que les components escollides siguin compatibles amb tots els morfismes; no demana que siguin l'única elecció possible. És una condició de compatibilitat, no d'unicitat ni d'absència de paràmetres, com mostra l'exemple següent.

Sobre els espais vectorials reals, prenem com a functors $F=G=\id_{\cat{Vect}_{\mathbb R}}$, el functor identitat de la categoria $\cat{Vect}_{\mathbb R}$, i busquem una transformació natural $\eta\colon\id\Rightarrow\id$. Fixem un escalar $\lambda\in\mathbb{R}$ i prenem, per a cada espai $V$, la component
\[
\eta_V=\lambda\,\id_V\colon V\to V,
\]
que és una aplicació lineal ---la multiplicació per l'escalar $\lambda$---, i per tant un morfisme de $\cat{Vect}_{\mathbb R}$, amb el tipus correcte perquè $F(V)=G(V)=V$. En efecte, aquesta família és una transformació natural: per a tota aplicació lineal $T\colon V\to W$ el quadrat de naturalitat és
\[
\begin{tikzpicture}[baseline=(m.center)]
  \matrix (m) [matrix of math nodes, row sep=2.8em, column sep=3.8em] {
    V & W \\
    V & W \\
  };
  \draw[-{Stealth[scale=1.1]}] (m-1-1) -- node[above] {$T$} (m-1-2);
  \draw[-{Stealth[scale=1.1]}] (m-1-1) -- node[left] {$\lambda\,\id_V$} (m-2-1);
  \draw[-{Stealth[scale=1.1]}] (m-1-2) -- node[right] {$\lambda\,\id_W$} (m-2-2);
  \draw[-{Stealth[scale=1.1]}] (m-2-1) -- node[below] {$T$} (m-2-2);
\end{tikzpicture}
\]
i commuta, perquè
\[
T\comp(\lambda\,\id_V)=\lambda\,T=(\lambda\,\id_W)\comp T,
\]
on les dues fletxes horitzontals són $T$ perquè el functor identitat actua sobre $T$ deixant-lo igual. Quan $\lambda\neq 0$, cada component $\eta_V=\lambda\,\id_V$ és a més un isomorfisme d'espais vectorials, amb invers $\lambda^{-1}\id_V$; això farà de $\eta$ un exemple del que més endavant anomenarem un \emph{isomorfisme natural}. Tenim, doncs, tota una família de transformacions naturals, una per cada valor de $\lambda$: la naturalitat es compleix per a qualsevol $\lambda$ fixat, tot i que la construcció depèn d'aquesta elecció. Naturalitat no vol dir, per tant, \enquote{sense eleccions}, sinó \enquote{compatible amb tots els morfismes}.


\begin{exemple}[Una família que \emph{no} és natural]\label{ex:no-natural}\index{transformació natural!contraexemple}
És igual d'instructiu veure una família de components que \emph{falla} la condició de naturalitat. Treballem a $\cat{Set}$ amb $F=G=\id_{\cat{Set}}$, de manera que una transformació $\id\Rightarrow\id$ hauria de donar, per a cada conjunt $A$, una aplicació $\eta_A\colon A\to A$ que commuti amb \emph{totes} les aplicacions. Definim una família completament explícita, sense cap tria:
\[
\eta_A=\id_A\quad\text{si } A\neq\{0,1\},
\qquad\qquad
\eta_{\{0,1\}}=\tau,\quad \tau(0)=1,\ \tau(1)=0.
\]
És a dir, totes les components són identitats excepte la del conjunt $\{0,1\}$, que és la transposició dels dos elements. És una família de morfismes $\eta_A\colon A\to A$ indexada per \emph{tots} els conjunts, i cada component és fins i tot una bijecció.

Vegem què demana el quadrat de naturalitat per a una aplicació $f\colon A\to B$. Com que $F=G=\id_{\cat{Set}}$, es té $F(f)=G(f)=f$, i la condició és $f\comp\eta_A=\eta_B\comp f$. Prenem $A=B=\{0,1\}$ i l'aplicació constant de valor $0$, $f(0)=f(1)=0$:
\[
\begin{tikzpicture}[baseline=(m.center)]
  \matrix (m) [matrix of math nodes, row sep=2.8em, column sep=3.8em] {
    \{0,1\} & \{0,1\} \\
    \{0,1\} & \{0,1\} \\
  };
  \draw[-{Stealth[scale=1.1]}] (m-1-1) -- node[above] {$f$} (m-1-2);
  \draw[-{Stealth[scale=1.1]}] (m-1-1) -- node[left] {$\tau$} (m-2-1);
  \draw[-{Stealth[scale=1.1]}] (m-1-2) -- node[right] {$\tau$} (m-2-2);
  \draw[-{Stealth[scale=1.1]}] (m-2-1) -- node[below] {$f$} (m-2-2);
\end{tikzpicture}
\qquad
f\comp\tau=f,\qquad \tau\comp f=\text{constant de valor }1.
\]
El camí de l'esquerra i a baix dona $f\comp\tau$, que continua sent constant de valor $0$; el de dalt i a la dreta dona $\tau\comp f$, que és constant de valor $1$. Els dos camins difereixen, i la família $(\eta_A)$ \emph{no} és una transformació natural.

La lliçó és que tenir una regla objecte a objecte, fins i tot formada per bijeccions, no basta. La component a $\{0,1\}$ s'ha definit sense tenir en compte les aplicacions que entren i surten d'aquest conjunt, i una sola d'aquestes aplicacions n'hi ha prou per trencar la compatibilitat. La naturalitat és exactament aquesta compatibilitat amb \emph{totes} les fletxes.
\end{exemple}

\section{Exemples de transformacions naturals}\label{sec:exemples-transf}

\begin{exemple}[Transformació identitat]\index{transformació natural!identitat}
Per a tot functor $F\colon\cat{C}\to\cat{D}$ hi ha la \textbf{transformació natural identitat} $\id_F\colon F\Rightarrow F$, amb components $(\id_F)_A=\id_{F(A)}$. El quadrat de naturalitat commuta trivialment, perquè els dos camins són $F(f)$.
\end{exemple}

\begin{exemple}[Preordres]\index{transformació natural!entre preordres}
Siguin $P,Q$ dos preordres i $F,G\colon P\to Q$ dos functors, és a dir, dues aplicacions monòtones. Una transformació natural $\eta\colon F\Rightarrow G$ dona, per a cada $x\in P$, un morfisme $\eta_x\colon F(x)\to G(x)$; però a $Q$ hi ha una fletxa $F(x)\to G(x)$ exactament quan $F(x)\leq G(x)$. La condició de naturalitat es compleix automàticament, perquè entre dos objectes d'un preordre hi ha com a màxim una fletxa. Per tant,
\[
\text{existeix }\eta\colon F\Rightarrow G
\quad\Longleftrightarrow\quad
F(x)\leq G(x)\ \text{per a tot }x.
\]
En el llenguatge de l'ordre, una transformació natural entre aplicacions monòtones és exactament la relació \enquote{$F$ és puntualment menor o igual que $G$}.
\end{exemple}

\begin{exemple}[Ampliació: interpretació de fórmules de primer ordre]\label{ex:interpretacio-primer-ordre}\index{transformació natural!interpretació lògica}\index{lema de substitució}
Reprenem el fil de la lògica que hem obert a la secció~\ref{sec:cap-transformacio} (\enquote{Cap al concepte de transformació natural}), ara en primer ordre i amb la naturalitat plantejada com a exemple formal. Fixem un llenguatge de primer ordre i una estructura $M$. L'exemple pressuposa la semàntica de la lògica de primer ordre (estructures, satisfacció, variables lliures i lligades) i es pot ometre en una primera lectura.

Convé fixar d'entrada les convencions que fan functorial la construcció sintàctica, en lloc d'adoptar-les a posteriori. Disposem d'una reserva infinita de variables; un \emph{context} és un conjunt finit $X$ de variables d'aquesta reserva, i escrivim $\cat{Ctx}$ per a la subcategoria plena de $\cat{FinSet}$ que té per objectes els contextos i per morfismes totes les aplicacions entre ells (com que la reserva és infinita, $\cat{Ctx}$ és equivalent a $\cat{FinSet}$, però no farem servir aquesta equivalència); treballem amb fórmules \emph{mòdul reanomenament de variables lligades} ($\alpha$-equivalència), i totes les substitucions s'entenen \emph{lliures de captura}. Sense aquestes convencions les lleis functorials de sota no se sostenen: en $\varphi=\exists y\,R(x,y)$, substituir la variable lliure $x$ per $y$ no pot donar ingènuament $\exists y\,R(y,y)$; cal reanomenar abans la variable lligada i obtenir, per exemple, $\exists z\,R(y,z)$.

Amb aquestes convencions considerem dues construccions. Per una banda,
\[
F(X)=\{\text{fórmules amb variables lliures entre }X\},
\]
enteses com a classes d'$\alpha$-equivalència; un canvi de variables $f\colon X\to Y$ hi actua per substitució, $F(f)(\varphi)=\varphi[f]$, reemplaçant cada variable lliure $x$ per $f(x)$. Per l'altra,
\[
G(X)=\{\text{predicats definibles sobre }X\},
\]
on un \emph{predicat definible} és un subconjunt de $M^X$ (el conjunt d'assignacions $s\colon X\to M$) de la forma
\[
\llbracket\varphi\rrbracket=\{\,s\colon X\to M\mid M\models\varphi[s]\,\},
\]
és a dir, l'\emph{extensió} d'alguna fórmula. El canvi de variables $f$ actua sobre $G$ transportant assignacions: com que $f$ indueix la precomposició $M^{f}\colon M^{Y}\to M^{X}$, $s\mapsto s\comp f$, posem
\[
G(f)(S)=\bigl(M^{f}\bigr)^{-1}(S)=\{\,s\in M^{Y}\mid s\comp f\in S\,\},\qquad S\subseteq M^{X}.
\]
Totes dues són functors covariants $\cat{Ctx}\to\cat{Set}$: per a $F$ es comprova que la substitució respecta identitats i composició\footnotemark; per a $G$, la functorialitat és la de la precomposició $M^{(-)}$ seguida de preimatge.

A cada context $X$ li associem ara l'aplicació que interpreta cada fórmula com la seva extensió,
\[
\eta_X\colon F(X)\to G(X),
\qquad
\eta_X(\varphi)=\llbracket\varphi\rrbracket.
\]
Aquesta $\eta_X$ està ben definida sobre classes d'$\alpha$-equivalència, perquè dues fórmules que només difereixen en el nom de les variables lligades tenen la mateixa extensió: la semàntica és $\alpha$-invariant, i per això el pas al quocient no altera res. La condició de naturalitat, per a un canvi de variables $f\colon X\to Y$, és que interpretar la fórmula substituïda coincideixi amb transportar l'extensió de la fórmula original:
\[
\llbracket\varphi[f]\rrbracket=\{\,s\colon Y\to M\mid s\comp f\in\llbracket\varphi\rrbracket\,\}.
\]
Aquesta igualtat és exactament el \textbf{lema de substitució} de la teoria de models ---el significat d'una fórmula després de canviar-ne les variables es calcula transportant el significat original--- i fa doble feina: garanteix alhora que $G(f)$ envia predicats definibles a predicats definibles (la preimatge d'un definible és definible) i que el quadrat de naturalitat commuta. Per tant $\eta\colon F\Rightarrow G$ és una transformació natural, i expressa que la interpretació és compatible amb la substitució.

Com en el cas proposicional de la secció~\ref{sec:cap-transformacio}, $\eta_X$ no és injectiva ---dues fórmules lògicament equivalents defineixen el mateix predicat---, de manera que la naturalitat no és aquí una identificació, sinó que expressa que la interpretació és compatible amb els canvis de variables, inclosos els reanomenaments i les identificacions. Convé, finalment, no confondre dues nocions que aquí conviuen: en una estructura $M$ fixada, dues fórmules poden tenir la mateixa extensió sense ser lògicament equivalents en \emph{totes} les estructures; la relació d'\emph{equivalència semàntica relativa a $M$} és, per tant, més grollera que l'\emph{equivalència lògica} global.
\end{exemple}
\footnotetext{És on es fan servir les convencions anteriors. Quan $f$ no és injectiva, identifica variables lliures, i és la substitució lliure de captura, definida sobre classes d'$\alpha$-equivalència, la que fa que $\varphi[\id]=\varphi$ i $\varphi[g\comp f]=(\varphi[f])[g]$ valguin com a \emph{igualtats}, no només mòdul reanomenament ad hoc.}%

\begin{exemple}[Bidual d'espais vectorials]\label{ex:bidual}\index{transformació natural!bidual}
Reprenem el fil que vam deixar obert a l'exemple~\ref{ex:homfunctor-dual} (\enquote{El dual d'un espai vectorial com a homfunctor}), on vam construir el \textbf{functor dual}
\[
(-)^*=\Hom_{\mathbb K}(-,\mathbb K)\colon\cat{Vect}_{\mathbb K}\op\to\cat{Vect}_{\mathbb K},\qquad V\mapsto V^*,
\]
contravariant, que sobre un morfisme $T\colon V\to W$ dona l'aplicació dual $T^*\colon W^*\to V^*$, $\psi\mapsto\psi\comp T$, amb la regla $(g\comp f)^*=f^*\comp g^*$. Vam ajornar deliberadament la comparació de $V$ amb el bidual fins a tenir el llenguatge de les transformacions naturals; ja el tenim.

El \textbf{functor bidual} (o doble dual) s'obté per composició del dual amb ell mateix. Com que $(-)^*$ és contravariant ---un functor amb domini $\cat{Vect}_{\mathbb K}\op$---, per compondre'l amb ell mateix cal fer encaixar dominis i codominis amb l'oposat, exactament com a la proposició~\ref{prop:variancia-composicio}: es pren el functor oposat $\bigl((-)^*\bigr)\op\colon\cat{Vect}_{\mathbb K}\to\cat{Vect}_{\mathbb K}\op$ ---que actua igual sobre objectes--- i s'hi torna a aplicar el dual,
\[
(-)^{**}=(-)^*\comp\bigl((-)^*\bigr)\op\colon\cat{Vect}_{\mathbb K}\to\cat{Vect}_{\mathbb K}.
\]
Per la regla dels signes, la composició de dos functors contravariants és \emph{covariant}. Desplegada, actua sobre objectes com
\[
V\mapsto V^{**}=(V^*)^*,
\]
i sobre un morfisme $T\colon V\to W$ com el dual aplicat dues vegades,
\[
T^{**}=(T^*)^*\colon V^{**}\to W^{**},\qquad \rho\mapsto\rho\comp T^*,
\]
amb la regla $(g\comp f)^{**}=g^{**}\comp f^{**}$: invertir l'ordre dos cops el restaura, i per això el bidual és un endofunctor covariant de $\cat{Vect}_{\mathbb K}$, paral·lel a la identitat.

L'efecte sobre els morfismes ho resumeix el diagrama següent, on el dual inverteix el sentit de la fletxa i el bidual, aplicant-lo dues vegades, el recupera:
\[
\begin{tikzpicture}[baseline=(m.center)]
  \matrix (m) [matrix of math nodes, column sep=3.2em, row sep=3.0em] {
    {} & V & W & {} \\
    {} & V^* & W^* & {} \\
    {} & V^{**} & W^{**} & {} \\
  };
  \draw[cat] (m-1-2) -- node[above] {$T$} (m-1-3);
  \draw[cat] (m-2-3) -- node[above] {$T^*$} (m-2-2);
  \draw[cat] (m-3-2) -- node[above] {$T^{**}$} (m-3-3);
  \draw[cat,dashed] (m-1-1) -- node[left] {$(-)^*$} (m-2-1);
  \draw[cat,dashed] (m-2-1) -- node[left] {$(-)^*$} (m-3-1);
  \draw[cat,dashed] (m-1-4) -- node[right] {$(-)^{**}$} (m-3-4);
\end{tikzpicture}
\]
Les fletxes contínues són morfismes de $\cat{Vect}_{\mathbb K}$; les puntejades, l'acció del functor. La fila del mig apunta en sentit contrari ---el dual $(-)^*$ és contravariant---, i la de baix en recupera el sentit ---dues inversions es cancel·len---: el bidual $(-)^{**}$, composició del dual amb ell mateix, és covariant, paral·lel a la identitat. Ara sí que té sentit comparar-los amb una transformació natural. Tots dos ---la identitat i el bidual--- són endofunctors covariants de $\cat{Vect}_{\mathbb K}$, de manera que la comparació és una transformació natural $\eta$ entre functors paral·lels:
\[
\begin{tikzpicture}[baseline=(A.base)]
  \node (A) at (0,0) {$\cat{Vect}_{\mathbb K}$};
  \node (B) at (3.6,0) {$\cat{Vect}_{\mathbb K}$};
  \draw[cat] (A) to[bend left=26] node[above] {$\id_{\cat{Vect}_{\mathbb K}}$} (B);
  \draw[cat] (A) to[bend right=26] node[below] {$(-)^{**}$} (B);
  \draw[cat nat] (1.8,0.2) -- node[right=1.5pt] {$\eta$} (1.8,-0.2);
\end{tikzpicture}
\]
El component de $\eta$ a cada espai $V$ és l'\textbf{avaluació}, que envia un vector $v$ a la forma \enquote{avaluar en $v$}. La transformació natural és, doncs, la família
\[
\eta=\bigl(\eta_V\bigr)_{V\in\Ob(\cat{Vect}_{\mathbb K})},
\]
i cada component es desglossa nivell a nivell, amb la imatge de $v$ com una aplicació de $V^*$ en $\mathbb K$:
\[
\begin{array}{r c l}
\eta_V\colon V & \longrightarrow & V^{**}\\[6pt]
v & \longmapsto & \eta_V(v)\colon
  \begin{array}[t]{r c l}
    V^* & \longrightarrow & \mathbb K\\[4pt]
    \varphi & \longmapsto & \varphi(v).
  \end{array}
\end{array}
\]
Cada $\eta_V$ és lineal i està ben definida. La condició de naturalitat, per a una aplicació lineal $T\colon V\to W$, és la commutativitat del quadrat
\[
\begin{tikzpicture}[baseline=(m.center)]
  \matrix (m) [matrix of math nodes, row sep=2.8em, column sep=3.6em] {
    V & W \\
    V^{**} & W^{**} \\
  };
  \draw[-{Stealth[scale=1.1]}] (m-1-1) -- node[above] {$T$} (m-1-2);
  \draw[-{Stealth[scale=1.1]}] (m-1-1) -- node[left] {$\eta_V$} (m-2-1);
  \draw[-{Stealth[scale=1.1]}] (m-1-2) -- node[right] {$\eta_W$} (m-2-2);
  \draw[-{Stealth[scale=1.1]}] (m-2-1) -- node[below] {$T^{**}$} (m-2-2);
\end{tikzpicture}
\]
és a dir, $T^{**}\comp\eta_V=\eta_W\comp T$: prendre l'avaluació i després transportar-la pel bidual coincideix amb aplicar $T$ i després avaluar. Aquí la fletxa horitzontal superior és $T$ perquè el functor identitat actua sobre $T$ deixant-lo igual. En aquest cas, la naturalitat es pot llegir dient que la identificació de $V$ amb la seva imatge dins $V^{**}$ no fa servir cap elecció de base: la construcció $v\mapsto(\varphi\mapsto\varphi(v))$ és intrínseca. Aquesta lectura ---\enquote{independent de la base}--- és una interpretació d'aquest exemple concret, no la definició de naturalitat, que continua sent la commutativitat dels quadrats. En restringir-se als espais de dimensió finita, cada $\eta_V$ és un isomorfisme, cosa que reprendrem en parlar dels isomorfismes naturals. En canvi, la comparació de $V$ amb el \emph{dual} $V^*$ ---i no amb el bidual--- no es pot fer natural; ho precisarem a l'observació~\ref{obs:iso-no-natural}, un cop definits els isomorfismes naturals.
\end{exemple}

\section{Categoria de functors}\label{sec:categoria-functors}

Els functors entre dues categories fixes, amb les transformacions naturals com a morfismes, formen al seu torn una categoria. Aquesta secció en desenvolupa la construcció: la composició que la fa possible, la categoria $\cat{D}^{\cat{C}}$ mateixa, un exemple fonamental ---els grafs--- i els seus isomorfismes, que resulten ser els isomorfismes naturals; en tanquem l'estudi amb l'equivalència de categories.

Les transformacions naturals no només relacionen functors: també es poden \emph{compondre} entre elles. La composició bàsica apila una transformació damunt d'una altra, i és la que després convertirà els functors entre dues categories en una categoria, i la que dona sentit a parlar de la \emph{inversa} d'un isomorfisme natural.

\begin{definicio}[Composició vertical]\index{transformació natural!composició vertical}
Donades dues transformacions naturals $\eta\colon F\Rightarrow G$ i $\theta\colon G\Rightarrow H$ entre functors $\cat{C}\to\cat{D}$, la seva \textbf{composició vertical} $\theta\comp\eta\colon F\Rightarrow H$ té per components
\[
(\theta\comp\eta)_A=\theta_A\comp\eta_A.
\]
El nom \emph{vertical} ve d'aquesta representació: les dues transformacions s'apilen ---$\eta$ a dalt i $\theta$ a sota--- entre els tres functors paral·lels $F,G,H\colon\cat{C}\to\cat{D}$, i la composició $\theta\comp\eta$ les recorre de dalt a baix.
\[
\begin{tikzpicture}[baseline=(A.base)]
  \node (A) at (0,0) {$\cat{C}$};
  \node (B) at (3.4,0) {$\cat{D}$};
  \draw[cat] (A) to[bend left=44] node[above] {$F$} (B);
  \draw[cat] (A) -- (B);
  \draw[cat] (A) to[bend right=44] node[below] {$H$} (B);
  \node[fill=white,inner sep=1.5pt] at (1.7,0) {$G$};
  \draw[cat nat] (1.7,0.52) -- node[right=1.5pt] {$\eta$} (1.7,0.27);
  \draw[cat nat] (1.7,-0.27) -- node[right=1.5pt] {$\theta$} (1.7,-0.52);
\end{tikzpicture}
\]
El mateix nom es reconeix a escala de components: fixat un morfisme $f\colon A\to B$, les components de $\eta$ i de $\theta$ són fletxes \emph{verticals}, i els seus quadrats de naturalitat s'apilen l'un damunt de l'altre ---compartint la fila central $G(A)\to G(B)$---, de manera que $\theta\comp\eta$ en compon les columnes:
\[
\begin{tikzpicture}[baseline=(m.center)]
  \matrix (m) [matrix of math nodes, row sep=3em, column sep=3.6em] {
    F(A) & F(B) \\
    G(A) & G(B) \\
    H(A) & H(B) \\
  };
  \draw[cat] (m-1-1) -- node[above] {$F(f)$} (m-1-2);
  \draw[cat] (m-2-1) -- node[above] {$G(f)$} (m-2-2);
  \draw[cat] (m-3-1) -- node[below] {$H(f)$} (m-3-2);
  \draw[cat] (m-1-1) -- node[left] {$\eta_A$} (m-2-1);
  \draw[cat] (m-2-1) -- node[left] {$\theta_A$} (m-3-1);
  \draw[cat] (m-1-2) -- node[right] {$\eta_B$} (m-2-2);
  \draw[cat] (m-2-2) -- node[right] {$\theta_B$} (m-3-2);
\end{tikzpicture}
\]
Aïllant-ne el rectangle exterior, la composició $\theta\comp\eta\colon F\Rightarrow H$ té el seu propi quadrat de naturalitat (a l'esquerra), i es representa com una única cel·la de $F$ cap a $H$ (a la dreta):
\[
\begin{tikzpicture}[baseline=(m.center)]
  \matrix (m) [matrix of math nodes, row sep=2.8em, column sep=3.6em] {
    F(A) & F(B) \\
    H(A) & H(B) \\
  };
  \draw[cat] (m-1-1) -- node[above] {$F(f)$} (m-1-2);
  \draw[cat] (m-1-1) -- node[left] {$(\theta\comp\eta)_A$} (m-2-1);
  \draw[cat] (m-1-2) -- node[right] {$(\theta\comp\eta)_B$} (m-2-2);
  \draw[cat] (m-2-1) -- node[below] {$H(f)$} (m-2-2);
\end{tikzpicture}
\qquad\qquad
\begin{tikzpicture}[baseline=(A.base)]
  \node (A) at (0,0) {$\cat{C}$};
  \node (B) at (3.7,0) {$\cat{D}$};
  \draw[cat] (A) to[bend left=26] node[above] {$F$} (B);
  \draw[cat] (A) to[bend right=26] node[below] {$H$} (B);
  \draw[cat nat] (1.85,0.2) -- node[right=1.5pt] {$\theta\comp\eta$} (1.85,-0.2);
\end{tikzpicture}
\]
\end{definicio}

\begin{proposicio}
La composició vertical torna a ser una transformació natural, és associativa, i la transformació identitat $\id_F$ n'és l'element neutre.
\end{proposicio}

\begin{prova}
Per a la naturalitat n'hi ha prou d'observar que, al diagrama de components de la definició, tots dos quadrats commuten (naturalitat de $\eta$ i de $\theta$); per tant també commuta el rectangle exterior, les columnes del qual són $(\theta\comp\eta)_A$ i $(\theta\comp\eta)_B$, i això és el quadrat de naturalitat de $\theta\comp\eta$ ---equivalentment, $H(f)\comp\theta_A\comp\eta_A=\theta_B\comp\eta_B\comp F(f)$, encadenant els dos quadrats. L'associativitat i la propietat de la identitat es redueixen a les propietats corresponents de la composició a $\cat{D}$, aplicades component a component.
\end{prova}

La composició vertical de transformacions naturals, introduïda més amunt, organitza els functors entre dues categories en una nova categoria.

\begin{definicio}\index{categoria!de functors}\index{DC@$\cat{D}^{\cat{C}}$}
Sigui $\cat{C}$ una categoria petita i $\cat{D}$ una categoria qualsevol. La \textbf{categoria de functors} $\cat{D}^{\cat{C}}$ (també escrita $[\cat{C},\cat{D}]$ o $\operatorname{Fun}(\cat{C},\cat{D})$) té:
\begin{itemize}
\item per objectes, els functors $\cat{C}\to\cat{D}$;
\item per morfismes, les transformacions naturals entre ells;
\item per composició, la composició vertical, i per identitats, les transformacions identitat.
\end{itemize}
\end{definicio}

\begin{observacio}[Qüestions de mida]\label{obs:mida-functors}\index{transformació natural!Nat@$\operatorname{Nat}$}
Convé dir per què la definició demana $\cat{C}$ petita. La raó de fons és la mateixa de l'observació~\ref{obs:mida-cat}, on vam haver de restringir $\cat{Cat}$ a les categories \emph{petites}. Si $\cat{C}$ fos gran, les dades que defineixen un functor $\cat{C}\to\cat{D}$ podrien formar una classe pròpia ---l'assignació recorreria una classe pròpia d'objectes---, i les classes pròpies no poden ser elements de cap col·lecció; aleshores els functors $\cat{C}\to\cat{D}$ no es poden reunir com els \emph{objectes} d'una categoria.\footnotemark Amb $\cat{C}$ petita, en canvi, aquestes dades formen un conjunt i la construcció no té cap problema. Això no impedeix parlar \emph{individualment} de functors i transformacions naturals amb domini gran ---ho fem contínuament, com amb $\cat{Set}\to\cat{Set}$---; el que no fem és agrupar-los tots en una categoria de functors.

Amb $\cat{C}$ petita, $\cat{D}^{\cat{C}}$ pot ser encara gran, però és \textbf{localment petita} tan bon punt $\cat{D}$ ho és: una transformació natural $F\Rightarrow G$ queda determinada pels seus components $\eta_A$, un per a cada objecte $A$ de $\cat{C}$, de manera que la col·lecció de totes les transformacions naturals $F\Rightarrow G$, que denotem $\operatorname{Nat}(F,G)$, és un subconjunt del producte
\[
\operatorname{Nat}(F,G)\;\subseteq\;\prod_{A\in\Ob(\cat{C})}\Hom_{\cat{D}}(F(A),G(A)),
\]
indexat per $\Ob(\cat{C})$ ---un conjunt, precisament perquè $\cat{C}$ és petita---, pel mateix argument que feia $\cat{Cat}$ localment petita. Com que aquest producte és un conjunt, $\operatorname{Nat}(F,G)$ també ho és, i és el conjunt de morfismes de $F$ a $G$ a la categoria de functors:
\[
\operatorname{Nat}(F,G)=\Hom_{\cat{D}^{\cat{C}}}(F,G).
\]
\end{observacio}
\footnotetext{Això depèn del marc de mida, que aquí és el de classes i conjunts de l'apèndix~\ref{cap:mida}. Amb universos de Grothendieck, la categoria de functors d'una categoria gran es pot formar en un univers superior; vegeu \cite[§1.1]{riehl}.}%

\index{graf!com a categoria de functors}\index{categoria!de functors!grafs}
Un dels exemples més il·luminatius de categoria de functors mostra que un objecte aparentment aliè a les categories ---un graf dirigit--- \emph{és}, en realitat, un functor, i que els seus morfismes són transformacions naturals. Aquest exemple reprèn el fil dels grafs que hem seguit des del primer capítol i és, a més, el primer indici que categories de functors del tipus $\cat{Set}^{\cat{C}}$ codifiquen estructures matemàtiques familiars.

La categoria índex $\cat{P}$ no és nova: és la categoria de \textbf{dues fletxes paral·leles} de la subsecció~\ref{subsec:finites}, amb dos objectes i dues fletxes no trivials amb el mateix domini i el mateix codomini. Els seus dos objectes són, en si, abstractes; els rebategem $E$ i $V$ ---i les fletxes, $s$ i $t$--- per anticipar la lectura com a grafs. Com que una aresta apunta cap als seus extrems, el domini de les fletxes és l'objecte de les \emph{arestes}, $E$, i el codomini el dels \emph{vèrtexs}, $V$, amb $s,t\colon E\to V$:
\[
\begin{tikzpicture}[baseline=(E.base)]
  \node (E) at (0,0) {$E$};
  \node (V) at (2.4,0) {$V$};
  \draw[-{Stealth[scale=1.1]}] ([yshift=3pt]E.east) -- node[above] {$s$} ([yshift=3pt]V.west);
  \draw[-{Stealth[scale=1.1]}] ([yshift=-3pt]E.east) -- node[below] {$t$} ([yshift=-3pt]V.west);
\end{tikzpicture}
\]

\begin{proposicio}\index{graf!com a categoria de functors}
La categoria de functors $\cat{Set}^{\cat{P}}$ és isomorfa a la categoria $\cat{Graph}$ dels grafs dirigits i els morfismes de grafs.
\end{proposicio}

\begin{prova}
Un functor $G\colon\cat{P}\to\cat{Set}$ consisteix en dos conjunts $G(E)$ i $G(V)$ i dues aplicacions $G(s),G(t)\colon G(E)\to G(V)$. Interpretant $G(E)$ com el conjunt d'\textbf{arestes}, $G(V)$ com el conjunt de \textbf{vèrtexs} i $G(s),G(t)$ com les aplicacions \enquote{origen} i \enquote{destí}, aquesta dada és \emph{exactament} un graf dirigit. Recíprocament, tot graf dirigit determina un tal functor. Els objectes de $\cat{Set}^{\cat{P}}$ ---que són functors $\cat{P}\to\cat{Set}$--- es corresponen, doncs, bijectivament amb els grafs dirigits.

Vegem ara els morfismes. Una transformació natural $\varphi\colon G\Rightarrow H$ entre dos d'aquests functors té dos components, $\varphi_E\colon G(E)\to H(E)$ i $\varphi_V\colon G(V)\to H(V)$, i la condició de naturalitat per a les fletxes $s$ i $t$ dona els dos quadrats
\[
\begin{tikzpicture}[baseline=(m.center)]
  \matrix (m) [matrix of math nodes, row sep=2.4em, column sep=3em] {
    G(E) & G(V) \\
    H(E) & H(V) \\
  };
  \draw[-{Stealth[scale=1.1]}] (m-1-1) -- node[above] {$G(s)$} (m-1-2);
  \draw[-{Stealth[scale=1.1]}] (m-1-1) -- node[left] {$\varphi_E$} (m-2-1);
  \draw[-{Stealth[scale=1.1]}] (m-1-2) -- node[right] {$\varphi_V$} (m-2-2);
  \draw[-{Stealth[scale=1.1]}] (m-2-1) -- node[below] {$H(s)$} (m-2-2);
\end{tikzpicture}
\qquad\text{i}\qquad
\begin{tikzpicture}[baseline=(n.center)]
  \matrix (n) [matrix of math nodes, row sep=2.4em, column sep=3em] {
    G(E) & G(V) \\
    H(E) & H(V) \\
  };
  \draw[-{Stealth[scale=1.1]}] (n-1-1) -- node[above] {$G(t)$} (n-1-2);
  \draw[-{Stealth[scale=1.1]}] (n-1-1) -- node[left] {$\varphi_E$} (n-2-1);
  \draw[-{Stealth[scale=1.1]}] (n-1-2) -- node[right] {$\varphi_V$} (n-2-2);
  \draw[-{Stealth[scale=1.1]}] (n-2-1) -- node[below] {$H(t)$} (n-2-2);
\end{tikzpicture}
\]
Aquests quadrats diuen que $\varphi_E$ i $\varphi_V$ respecten l'origen i el destí de cada aresta; és a dir, són precisament la definició de \textbf{morfisme de grafs}. Per tant els morfismes de $\cat{Set}^{\cat{P}}$ es corresponen amb els morfismes de grafs. Aquesta correspondència, sobre objectes i sobre morfismes, és un parell de functors inversos l'un de l'altre, i ho és \emph{estrictament}: les seves composicions són les identitats \emph{exactes} ---no només isomorfismes naturals---, perquè un functor $\cat{P}\to\cat{Set}$ no és res més que els seus valors sobre $E$, $V$, $s$ i $t$; anar del functor al graf i tornar dona el mateix functor, i del graf al functor i tornar, el mateix graf. És, doncs, un isomorfisme de categories en el sentit estricte de la definició~\ref{def:iso-categories} ---més fort que l'equivalència que veurem tot seguit---, i $\cat{Set}^{\cat{P}}\cong\cat{Graph}$. El caràcter estricte d'aquest isomorfisme descansa sobre la codificació que hem triat: en identificar un graf amb les dades $(E,V,s,t)$ ---exactament les que porta un functor $\cat{P}\to\cat{Set}$---, l'anada i la tornada retornen les \emph{mateixes} dades. Amb altres codificacions habituals dels grafs, la correspondència continua sent un isomorfisme \emph{canònic} de categories, però no una igualtat literal de dades; en cap cas \enquote{exactes} no s'ha de llegir com una igualtat fundacional independent de la codificació.
\end{prova}

\begin{observacio}
Aquest exemple té una lectura de fons: un cop decidim representar els grafs com a functors $\cat{P}\to\cat{Set}$, la representació \emph{justifica i recupera canònicament} la definició habitual de morfisme de grafs, perquè els morfismes de $\cat{Set}^{\cat{P}}$ són, per força, les transformacions naturals. És un criteri de coherència que sovint es fa servir en teoria de categories: quan una estructura s'expressa com a categoria de functors, la noció de morfisme que se'n dedueix hauria de coincidir amb la que s'havia adoptat independentment. Les categories del tipus $\cat{Set}^{\cat{C}}$, que codifiquen estructures com els grafs, inclouen com a cas particular les categories de prefeixos, que presentem tot seguit i que seran fonamentals per al lema de Yoneda i, més endavant, per a la semàntica categòrica.
\end{observacio}

\begin{exemple}[La categoria de prefeixos]\label{ex:categoria-prefeixos}\index{prefeix@prefeix!categoria de}\index{prefeix@prefeix!morfisme de}
Sigui $\cat{C}$ una categoria petita. Els prefeixos sobre $\cat{C}$ (definició~\ref{def:prefeix}) són els objectes de la categoria de functors
\[
\widehat{\cat{C}}\coloneqq\cat{Set}^{\cat{C}\op},
\]
que anomenem \textbf{categoria de prefeixos} sobre $\cat{C}$. Els seus morfismes, els \textbf{morfismes de prefeixos} $\sigma\colon P\Rightarrow Q$, són les transformacions naturals: famílies d'aplicacions $\sigma_A\colon P(A)\to Q(A)$ tals que, per a cada morfisme $f\colon A\to B$ de $\cat{C}$,
\[
\sigma_A\comp P(f)=Q(f)\comp\sigma_B\colon P(B)\to Q(A),
\]
que és el quadrat de naturalitat escrit per a prefeixos. Com que $\cat{Set}$ és localment petita, $\widehat{\cat{C}}$ també ho és, per l'observació sobre qüestions de mida. La categoria $\cat{Set}^{\cat{P}}$ dels grafs n'és un exemple: com que un functor $\cat{P}\to\cat{Set}$ és un functor $(\cat{P}\op)\op\to\cat{Set}$, els grafs són els prefeixos sobre $\cat{P}\op$.
\end{exemple}

\begin{observacio}[La forma dels grafs és, ella mateixa, un graf]\index{categoria!lliure}
La categoria índex $\cat{P}$ té una lectura circular ben curiosa. El seu \emph{graf subjacent} ---els dos objectes com a vèrtexs i les dues fletxes generadores $s,t$ com a arestes--- és exactament $\bullet\rightrightarrows\bullet$, i $\cat{P}$ n'és la \textbf{categoria lliure} (secció~\ref{sec:grafs-lliures}): com que les dues fletxes paral·leles no es poden encadenar ---de l'objecte d'arribada no en surt cap---, no hi ha camins de longitud $\geq 2$, i els únics morfismes de $\cat{P}$ són $\id$, $\id$, $s$ i $t$. És a dir, $\cat{P}=\Free(\bullet\rightrightarrows\bullet)$. Així, la \enquote{forma} amb què definim els grafs ---el graf subjacent de $\cat{P}$--- és, ella mateixa, un graf, de fet un objecte de $\cat{Graph}=\cat{Set}^{\cat{P}}$: el patró es tanca sobre si mateix.
\end{observacio}

Ara que les transformacions naturals formen una categoria $\cat{D}^{\cat{C}}$, en podem estudiar els isomorfismes.

\begin{definicio}\index{isomorfisme natural}
Una transformació natural $\eta\colon F\Rightarrow G$ és un \textbf{isomorfisme natural} si admet una \emph{inversa}: una transformació natural $\eta^{-1}\colon G\Rightarrow F$ amb $\eta^{-1}\comp\eta=\id_F$ i $\eta\comp\eta^{-1}=\id_G$ (composició vertical). En aquest cas escrivim $F\cong G$ i diem que $F$ i $G$ són \textbf{naturalment isomorfs}.
\end{definicio}

Quan $\cat{C}$ és petita, aquesta condició és, literalment, la d'un isomorfisme a la categoria de functors $\cat{D}^{\cat{C}}$. En efecte, els morfismes de $\cat{D}^{\cat{C}}$ són les transformacions naturals i la seva composició és la vertical; per tant, que $\eta$ admeti una inversa $\eta^{-1}$ com a transformació natural ---amb $\eta^{-1}\comp\eta=\id_F$ i $\eta\comp\eta^{-1}=\id_G$--- no és res més que tenir un invers a $\cat{D}^{\cat{C}}$. Els isomorfismes naturals coincideixen, doncs, amb els isomorfismes de $\cat{D}^{\cat{C}}$: la noció d'isomorfisme del primer capítol (definició~\ref{def:isomorfisme}), un nivell més amunt.

El fet clau és que aquesta condició \enquote{global} ---tenir inversa com a transformació natural--- es pot comprovar \enquote{punt a punt}, sobre les components:

\begin{proposicio}[Caracterització]\label{prop:iso-natural-components}
Siguin $F,G\colon\cat{C}\to\cat{D}$ dos functors. Una transformació natural $\eta\colon F\Rightarrow G$ és un isomorfisme natural si, i només si, cada component $\eta_A$ (per a $A$ objecte de $\cat{C}$) és un isomorfisme a $\cat{D}$.
\end{proposicio}

\begin{prova}
Si $\eta$ té inversa $\eta^{-1}$, avaluant les igualtats $\eta^{-1}\comp\eta=\id_F$ i $\eta\comp\eta^{-1}=\id_G$ a cada objecte $A$ ---i recordant que la composició vertical és component a component--- obtenim $(\eta^{-1})_A\comp\eta_A=\id_{F(A)}$ i $\eta_A\comp(\eta^{-1})_A=\id_{G(A)}$; per tant $\eta_A$ és un isomorfisme a $\cat{D}$, amb invers $(\eta^{-1})_A$.

Recíprocament, suposem que cada $\eta_A$ és un isomorfisme a $\cat{D}$, amb invers $\eta_A^{-1}$. N'hi ha prou de veure que aquests inversos són les components d'una transformació natural $\eta^{-1}\colon G\Rightarrow F$, perquè aleshores, component a component, és la inversa de $\eta$. Partim del quadrat de naturalitat de $\eta$ per a un morfisme $f\colon A\to B$,
\[
\begin{tikzpicture}[baseline=(m.center)]
  \matrix (m) [matrix of math nodes, row sep=2.6em, column sep=3.4em] {
    F(A) & F(B) \\
    G(A) & G(B) \\
  };
  \draw[-{Stealth[scale=1.1]}] (m-1-1) -- node[above] {$F(f)$} (m-1-2);
  \draw[-{Stealth[scale=1.1]}] (m-1-1) -- node[left] {$\eta_A$} (m-2-1);
  \draw[-{Stealth[scale=1.1]}] (m-1-2) -- node[right] {$\eta_B$} (m-2-2);
  \draw[-{Stealth[scale=1.1]}] (m-2-1) -- node[below] {$G(f)$} (m-2-2);
\end{tikzpicture}
\]
és a dir, $G(f)\comp\eta_A=\eta_B\comp F(f)$, i component per l'esquerra amb $\eta_B^{-1}$ i per la dreta amb $\eta_A^{-1}$,
\[
\eta_B^{-1}\comp G(f)=\eta_B^{-1}\comp G(f)\comp\eta_A\comp\eta_A^{-1}=\eta_B^{-1}\comp\eta_B\comp F(f)\comp\eta_A^{-1}=F(f)\comp\eta_A^{-1},
\]
que és la condició de naturalitat per a $\eta^{-1}$.
\end{prova}

\begin{observacio}\label{obs:iso-no-natural}
La distinció entre isomorfisme i isomorfisme \emph{natural} és fonamental. Que $F(A)$ i $G(A)$ siguin isomorfs per a cada $A$ no basta: cal que els isomorfismes ---un per objecte, les components $\eta_A$--- es puguin triar \emph{alhora} de manera compatible amb tots els morfismes. Aquesta compatibilitat no es refereix a triar més components, sinó a la \emph{naturalitat}: per a cada morfisme $f\colon A\to B$ de $\cat{C}$, el quadrat de $\eta$ ha de commutar ($G(f)\comp\eta_A=\eta_B\comp F(f)$), lligant les components dels dos extrems.

Un exemple mínim ho mostra sense cap tecnicisme. Prenem un grup de dos elements, $\{e,g\}$ amb $g^2=e$ (on $e$ és la identitat), vist com a categoria d'un sol objecte $\cat{B}(\{e,g\})$, i dos functors $F,G\colon\cat{B}(\{e,g\})\to\cat{Set}$. Tots dos envien l'únic objecte $\ast$ al mateix conjunt $\{0,1\}$, però l'element $g$ actua de manera diferent: en $F$ actua per la identitat, i en $G$ per la transposició $\tau$ que intercanvia $0$ i $1$. Puntualment no hi ha res a distingir: $F(\ast)=G(\ast)=\{0,1\}$ són el mateix conjunt, i per tant \emph{isomorfs}. Tanmateix, no hi ha cap isomorfisme natural $F\cong G$. Una transformació natural $\eta\colon F\Rightarrow G$ tindria una única component $\eta_\ast\colon\{0,1\}\to\{0,1\}$, i la naturalitat respecte del morfisme $g$ exigiria
\[
G(g)\comp\eta_\ast=\eta_\ast\comp F(g),
\qquad\text{és a dir}\qquad
\tau\comp\eta_\ast=\eta_\ast.
\]
Però $\tau$ no té punts fixos, de manera que cap aplicació $\eta_\ast$ pot complir $\tau\comp\eta_\ast=\eta_\ast$ (aplicant a qualsevol $x$ obtindríem $\tau(\eta_\ast(x))=\eta_\ast(x)$, impossible). No hi ha, doncs, cap transformació natural $F\Rightarrow G$, i menys encara un isomorfisme natural: la igualtat puntual dels objectes no es pot fer compatible amb l'acció del morfisme. Val a dir que aquí obtenim un contraexemple més fort que el necessari: per descartar un isomorfisme \emph{natural} n'hi hauria prou de veure que cap component \emph{bijectiva} compleix la naturalitat, mentre que l'argument anterior descarta \emph{totes} les components, bijectives o no. La manca d'isomorfisme natural en surt, doncs, reforçada.

El cas clàssic del dual d'un espai vectorial il·lustra la mateixa idea, però demana una precaució addicional. Sovint es diu que un espai de dimensió finita \enquote{és isomorf al seu dual $V\cong V^*$, però no naturalment}. Cal formular-ho amb cura, perquè la identitat és un functor \emph{covariant} $\cat{Vect}_{\mathbb K}\to\cat{Vect}_{\mathbb K}$ mentre que el dual $(-)^*$ és \emph{contravariant}, $\cat{Vect}_{\mathbb K}\op\to\cat{Vect}_{\mathbb K}$: no són dos functors paral·lels als quals es pugui aplicar directament la definició de transformació natural d'aquest capítol. La manera correcta de precisar-ho és restringir-se als \emph{isomorfismes} lineals i comparar la identitat amb l'acció covariant $T\mapsto(T^{-1})^*$; aleshores es comprova que cap família d'isomorfismes $V\cong V^*$ és compatible amb tots els isomorfismes lineals. Aquest desenvolupament es fa amb detall a l'exercici~\ref{ex:dual-no-natural} de la part de pràctica. En canvi, l'isomorfisme $V\cong V^{**}$ amb el bidual sí que està ben plantejat sobre \emph{totes} les aplicacions lineals ---identitat i bidual són tots dos covariants--- i és natural.
\end{observacio}

\begin{exemple}[Les identificacions canòniques del capítol anterior]\label{ex:isos-naturals-cap2}\index{isomorfisme natural!functor de parts}
Al capítol~\ref{cap:functors} vam trobar diverses identificacions canòniques entre un functor i un homfunctor, i vam anunciar que les reconeixeríem com a isomorfismes naturals. La proposició~\ref{prop:iso-natural-components} ho fa immediat: en cada cas, les components són bijeccions (o isomorfismes lineals), i només cal comprovar la naturalitat.

\emph{El functor de parts.} Els dos functors contravariants de l'exemple~\ref{ex:functor-parts} són
\[
\mathcal P,\ \Hom_{\cat{Set}}(-,\{0,1\})\colon\cat{Set}\op\longrightarrow\cat{Set}.
\] Les funcions característiques donen, per a cada conjunt $X$, la bijecció
\[
\chi_X\colon\mathcal P(X)\to\Hom_{\cat{Set}}(X,\{0,1\}),\qquad B\mapsto\chi_B.
\]
Un morfisme de $\cat{Set}\op$ de $Y$ a $X$ és una aplicació $f\colon X\to Y$, que $\mathcal P$ envia a l'antiimatge $f^{-1}$ i l'homfunctor a la precomposició $-\comp f$. El quadrat de naturalitat és
\[
\begin{tikzpicture}[baseline=(m.center)]
  \matrix (m) [matrix of math nodes, row sep=2.8em, column sep=4.2em] {
    \mathcal P(Y) & \mathcal P(X) \\
    \Hom(Y,\{0,1\}) & \Hom(X,\{0,1\}) \\
  };
  \draw[-{Stealth[scale=1.1]}] (m-1-1) -- node[above] {$f^{-1}$} (m-1-2);
  \draw[-{Stealth[scale=1.1]}] (m-1-1) -- node[left] {$\chi_Y$} (m-2-1);
  \draw[-{Stealth[scale=1.1]}] (m-1-2) -- node[right] {$\chi_X$} (m-2-2);
  \draw[-{Stealth[scale=1.1]}] (m-2-1) -- node[below] {$-\comp f$} (m-2-2);
\end{tikzpicture}
\qquad
\chi_B\comp f=\chi_{f^{-1}(B)}\quad(B\subseteq Y),
\]
i commuta, perquè per a tot $x\in X$ es té $(\chi_B\comp f)(x)=1$ si i només si $f(x)\in B$, és a dir, si i només si $x\in f^{-1}(B)$. Per tant $\chi\colon\mathcal P\Rightarrow\Hom_{\cat{Set}}(-,\{0,1\})$ és un isomorfisme natural, i ara podem escriure amb tota precisió
\[
\mathcal P\;\cong\;\Hom_{\cat{Set}}(-,\{0,1\}).
\]

\emph{Els altres casos.} El mateix esquema val per a les altres identificacions del capítol anterior; n'indiquem les components i l'equació de naturalitat, que es comprova directament.
\begin{itemize}
  \item A $\cat{Set}$, $\id_{\cat{Set}}\cong\Hom_{\cat{Set}}(\{\ast\},-)$: la component envia $x\in X$ a l'aplicació $\bar x\colon\ast\mapsto x$, i la naturalitat diu que $f\comp\bar x=\overline{f(x)}$.
  \item A $\cat{Graph}$, el functor de vèrtexs és naturalment isomorf a $\Hom_{\cat{Graph}}(\mathbf V,-)$ i el d'arestes a $\Hom_{\cat{Graph}}(\mathbf E,-)$ (exemple~\ref{ex:homfunctor-grafs}): la component envia cada vèrtex, o cada aresta, al morfisme que el tria, i la naturalitat diu que compondre aquest morfisme amb un morfisme de grafs $F$ tria la imatge per $F$.
  \item A $\cat{Vect}_{\mathbb K}$, el functor identitat és naturalment isomorf a l'aixecament
\[
\Hom_{\mathbb K}(\mathbb K,-)\colon\cat{Vect}_{\mathbb K}\longrightarrow\cat{Vect}_{\mathbb K}.
\]
La component envia $w\in W$ a l'aplicació lineal $\lambda\mapsto\lambda w$, amb inversa $\varphi\mapsto\varphi(1)$, i la naturalitat diu que $T\comp(\lambda\mapsto\lambda w)=(\lambda\mapsto\lambda\,T(w))$ per a tota aplicació lineal $T$.
\end{itemize}
Aquests exemples tornaran a aparèixer al capítol del lema de Yoneda, on direm que aquests functors són \emph{representables}.
\end{exemple}

\subsection{Equivalència de categories}\label{sec:equivalencia}

Quan volem dir que dues categories són \enquote{essencialment la mateixa} ---amb la mateixa estructura categòrica, distingint els objectes només llevat d'isomorfisme---, la igualtat estricta és massa rígida, i l'isomorfisme de categories, que demana functors inversos ($G\comp F=\id_{\cat{C}}$ i $F\comp G=\id_{\cat{D}}$), també ho és sovint. La noció adequada relaxa aquestes igualtats a isomorfismes naturals: en lloc d'un invers demanarem un \textbf{quasiinvers}, un functor $G$ amb $G\comp F\cong\id_{\cat{C}}$ i $F\comp G\cong\id_{\cat{D}}$.

\begin{definicio}\index{equivalència de categories}
\label{def:equivalencia}
Una \textbf{equivalència de categories} entre $\cat{C}$ i $\cat{D}$ és una parella de functors
\[
F\colon\cat{C}\to\cat{D},
\qquad
G\colon\cat{D}\to\cat{C},
\]
juntament amb isomorfismes naturals
\[
G\comp F\cong\id_{\cat{C}}
\qquad\text{i}\qquad
F\comp G\cong\id_{\cat{D}}.
\]
Quan existeix una tal equivalència, diem que $\cat{C}$ i $\cat{D}$ són \textbf{equivalents} i escrivim $\cat{C}\simeq\cat{D}$. El functor $G$ s'anomena un \textbf{quasiinvers} de $F$. Quan calgui donar nom als isomorfismes naturals, els orientarem sempre així:
\[
\eta\colon\id_{\cat{C}}\Rightarrow G\comp F,
\qquad
\varepsilon\colon F\comp G\Rightarrow\id_{\cat{D}}.
\]
\end{definicio}

\begin{observacio}
La diferència amb l'isomorfisme de categories és que aquí no demanem $G\comp F=\id_{\cat{C}}$, sinó només $G\comp F\cong\id_{\cat{C}}$: en tornar de $\cat{C}$ a $\cat{C}$ no recuperem exactament cada objecte, sinó un objecte \emph{naturalment isomorf}. Aquesta flexibilitat és la que fa que l'equivalència sigui la noció adequada de comparació entre categories quan, com passa sovint, només ens interessen els objectes llevat d'isomorfisme.
\end{observacio}

El resultat central caracteritza les equivalències mitjançant les tres propietats dels functors que vam introduir al capítol anterior: ser ple i fidel (la definició~\ref{def:ple-fidel}) i l'exhaustivitat essencial (la definició~\ref{def:ess-exhaustiu}).

\begin{teorema}\index{equivalència de categories!caracterització}
\label{teo:caract-equivalencia}
Un functor $F\colon\cat{C}\to\cat{D}$ forma part d'una equivalència de categories si i només si és \textbf{ple}, \textbf{fidel} i \textbf{essencialment exhaustiu}.
\end{teorema}

\begin{prova}
\emph{($\Rightarrow$)} Suposem que $F$ forma part d'una equivalència, amb quasiinvers $G\colon\cat{D}\to\cat{C}$ i isomorfismes naturals $\eta\colon\id_{\cat{C}}\Rightarrow G\comp F$ i $\varepsilon\colon F\comp G\Rightarrow\id_{\cat{D}}$. Com que les components de $\eta$ i de $\varepsilon$ són invertibles, els quadrats de naturalitat es poden llegir com a fórmules: per a $f\colon A\to B$ a $\cat{C}$ i $g\colon D\to D'$ a $\cat{D}$,
\[
G(F(f))=\eta_B\comp f\comp\eta_A^{-1},
\qquad\qquad
F(G(g))=\varepsilon_{D'}^{-1}\comp g\comp\varepsilon_D.
\tag{$\ast$}
\]

\emph{$F$ és fidel.} Siguin $f,f'\colon A\to B$ amb $F(f)=F(f')$. Aleshores $G(F(f))=G(F(f'))$ i, per~$(\ast)$, $\eta_B\comp f\comp\eta_A^{-1}=\eta_B\comp f'\comp\eta_A^{-1}$. Component amb $\eta_B^{-1}$ per l'esquerra i amb $\eta_A$ per la dreta, $f=f'$.

\emph{$G$ és fidel.} És el mateix argument amb la segona fórmula de~$(\ast)$: si $G(g)=G(g')$ per a $g,g'\colon D\to D'$, aleshores $\varepsilon_{D'}^{-1}\comp g\comp\varepsilon_D=\varepsilon_{D'}^{-1}\comp g'\comp\varepsilon_D$, i per tant $g=g'$.

\emph{$F$ és ple.} Sigui $u\colon F(A)\to F(B)$ un morfisme qualsevol de $\cat{D}$, i posem
\[
f=\eta_B^{-1}\comp G(u)\comp\eta_A\colon A\to B.
\]
Per~$(\ast)$, $G(F(f))=\eta_B\comp\bigl(\eta_B^{-1}\comp G(u)\comp\eta_A\bigr)\comp\eta_A^{-1}=G(u)$, i com que $G$ és fidel, $F(f)=u$. Per concloure $F(f)=u$ no es pot fer servir que $F$ és fidel, perquè encara no sabem que $u$ sigui la imatge de cap morfisme; és que $G$ sigui fidel el que ho permet.

\emph{$F$ és essencialment exhaustiu.} Per a cada objecte $D$ de $\cat{D}$, la component $\varepsilon_D\colon F(G(D))\to D$ és un isomorfisme, de manera que $D$ és isomorf a l'objecte $F(G(D))$ de la imatge de $F$.

\emph{($\Leftarrow$)} Suposem que $F$ és ple, fidel i essencialment exhaustiu. Construïm un quasiinvers en quatre passos.

\emph{Pas 1: $G$ sobre objectes.} Per l'exhaustivitat essencial, per a cada objecte $D$ de $\cat{D}$ hi ha algun objecte $C$ de $\cat{C}$ amb $F(C)\cong D$. Triem, per a cada $D$, un tal objecte, que anomenem $G(D)$, i un isomorfisme $\varepsilon_D\colon F(G(D))\to D$.

\emph{Pas 2: $G$ sobre morfismes.} Sigui $g\colon D\to D'$. El morfisme $\varepsilon_{D'}^{-1}\comp g\comp\varepsilon_D$ va de $F(G(D))$ a $F(G(D'))$. Com que $F$ és ple, prové d'algun morfisme $G(D)\to G(D')$, i com que $F$ és fidel, d'un de sol. Definim $G(g)$ com aquest únic morfisme; és a dir, $G(g)$ queda caracteritzat per
\[
F(G(g))=\varepsilon_{D'}^{-1}\comp g\comp\varepsilon_D.
\tag{$\ast\ast$}
\]

\emph{Pas 3: $G$ és un functor.} Per a la identitat, $(\ast\ast)$ dona $F(G(\id_D))=\varepsilon_D^{-1}\comp\varepsilon_D=\id_{F(G(D))}=F(\id_{G(D)})$, i, com que $F$ és fidel, $G(\id_D)=\id_{G(D)}$. Per a la composició, siguin $g\colon D\to D'$ i $g'\colon D'\to D''$. Aplicant $(\ast\ast)$ dues vegades,
\[
\begin{aligned}
F\bigl(G(g')\comp G(g)\bigr)&=F(G(g'))\comp F(G(g))
=\bigl(\varepsilon_{D''}^{-1}\comp g'\comp\varepsilon_{D'}\bigr)\comp\bigl(\varepsilon_{D'}^{-1}\comp g\comp\varepsilon_D\bigr)\\
&=\varepsilon_{D''}^{-1}\comp(g'\comp g)\comp\varepsilon_D=F\bigl(G(g'\comp g)\bigr),
\end{aligned}
\]
i, com que $F$ és fidel, $G(g')\comp G(g)=G(g'\comp g)$.

\emph{Pas 4: els isomorfismes naturals.} Component $(\ast\ast)$ amb $\varepsilon_{D'}$ per l'esquerra s'obté
\[
\varepsilon_{D'}\comp F(G(g))=g\comp\varepsilon_D,
\]
que és el quadrat de naturalitat de $\varepsilon\colon F\comp G\Rightarrow\id_{\cat{D}}$ per a $g$. Com que cada $\varepsilon_D$ és un isomorfisme, $\varepsilon$ és un isomorfisme natural.

Per construir $\eta$, sigui $A$ un objecte de $\cat{C}$. Com que $F$ és ple i fidel, hi ha un únic $\eta_A\colon A\to G(F(A))$ amb $F(\eta_A)=\varepsilon_{F(A)}^{-1}$, i un únic $\zeta_A\colon G(F(A))\to A$ amb $F(\zeta_A)=\varepsilon_{F(A)}$. Aleshores
\[
F(\zeta_A\comp\eta_A)=\varepsilon_{F(A)}\comp\varepsilon_{F(A)}^{-1}=F(\id_A),
\qquad
F(\eta_A\comp\zeta_A)=\varepsilon_{F(A)}^{-1}\comp\varepsilon_{F(A)}=F(\id_{G(F(A))}),
\]
i, com que $F$ és fidel, $\zeta_A$ és l'invers de $\eta_A$: cada $\eta_A$ és un isomorfisme. Queda la naturalitat: per a $f\colon A\to B$, cal veure $\eta_B\comp f=G(F(f))\comp\eta_A$. Apliquem $F$ als dos membres. D'una banda, $F(\eta_B\comp f)=\varepsilon_{F(B)}^{-1}\comp F(f)$. De l'altra, aplicant $(\ast\ast)$ al morfisme $g=F(f)$,
\[
F\bigl(G(F(f))\comp\eta_A\bigr)=\bigl(\varepsilon_{F(B)}^{-1}\comp F(f)\comp\varepsilon_{F(A)}\bigr)\comp\varepsilon_{F(A)}^{-1}=\varepsilon_{F(B)}^{-1}\comp F(f).
\]
Les dues imatges coincideixen, i, com que $F$ és fidel, $\eta_B\comp f=G(F(f))\comp\eta_A$. Així $\eta\colon\id_{\cat{C}}\Rightarrow G\comp F$ és un isomorfisme natural, i $G$ és un quasiinvers de $F$.
\end{prova}

\begin{observacio}[Sobre l'elecció]
A la direcció ($\Leftarrow$), el pas~1 tria simultàniament un objecte $G(D)$ i un isomorfisme $\varepsilon_D$ per a cada $D$, i per tant fa servir l'\emph{axioma de l'elecció}.\footnotemark{} Tota la resta de la construcció és \emph{determinada} per aquesta tria: el pas~2 no tria res, perquè, com que $F$ és ple i fidel, $G(g)$ és únic, i $\eta$ també queda fixat. En casos concrets sovint es pot donar un quasiinvers canònic sense recórrer a l'elecció. La direcció ($\Rightarrow$), en canvi, no tria res, i tampoc no fa servir cap relació entre $\eta$ i $\varepsilon$ més enllà de la naturalitat i la invertibilitat.
\end{observacio}
\footnotetext{Si $\cat{D}$ és gran, la família de tries està indexada per una classe pròpia, i cal una forma d'elecció més forta que l'habitual, l'\emph{elecció global}, disponible per exemple a la teoria de classes NBG (apèndix~\ref{cap:mida}). Vegeu \cite[§1.5]{riehl}.}%

\begin{definicio}[Esquelet]\index{categoria!esquelètica}\index{esquelet}\label{def:esquelet}
Una categoria és \textbf{esquelètica} si dos objectes isomorfs qualssevol són \emph{iguals}; és a dir, si cada classe d'isomorfisme té un únic objecte. Un \textbf{esquelet} d'una categoria $\cat{C}$ és una subcategoria plena que conté \emph{exactament} un objecte de cada classe d'isomorfisme de $\cat{C}$; un esquelet és sempre una categoria esquelètica.
\end{definicio}

\begin{observacio}
Convé distingir dues nocions properes. Una subcategoria plena que conté \emph{almenys} un objecte de cada classe d'isomorfisme ja és equivalent a $\cat{C}$ (com veurem tot seguit), però pot tenir diversos representants isomorfs d'una mateixa classe. Un \emph{esquelet} n'és el cas ajustat: en conté \emph{exactament} un, de manera que no queden objectes isomorfs distints. Triar exactament un representant de cada classe d'isomorfisme és, però, una elecció simultània. Amb l'\emph{axioma de l'elecció}, tota categoria \emph{petita} admet un esquelet, i tots els seus esquelets són equivalents entre si. Per a una categoria gran, seleccionar simultàniament un representant de cada classe d'isomorfisme pot requerir un principi d'elecció global, segons el marc fundacional adoptat.\footnotemark És la mateixa elecció que fa servir la construcció d'un quasiinvers (observació anterior).
\end{observacio}
\footnotetext{En una categoria gran, les classes d'isomorfisme formen una classe pròpia, i cal la mateixa forma d'elecció global que a l'observació anterior (apèndix~\ref{cap:mida}); la petitesa local només controla els homconjunts i no resol aquesta selecció.}%

\begin{exemple}
Tota categoria $\cat{C}$ és equivalent a qualsevol subcategoria plena que contingui almenys un objecte de cada classe d'isomorfisme de $\cat{C}$. En efecte, el functor inclusió d'aquesta subcategoria és ple i fidel per ser una subcategoria plena, i essencialment exhaustiu per contenir un representant de cada classe. Per exemple, la categoria $\cat{FinSet}$ de tots els conjunts finits és equivalent a la seva subcategoria plena formada pels conjunts $\underline{0}=\varnothing$ i $\underline{n}=\{1,\dots,n\}$ per a $n\geq 1$ (exemple~\ref{ex:ess-exhaustiu-finset}): no cal treballar amb tots els conjunts finits, sinó només amb un de cada mida. Aquesta subcategoria $\{\underline{n}\mid n\geq 0\}$ és, de fet, un \emph{esquelet} de $\cat{FinSet}$, perquè conté exactament un conjunt de cada cardinalitat finita. Convé notar que $\cat{FinSet}$ és una categoria \emph{gran} ---els seus objectes són massa nombrosos per formar un conjunt, tal com passa amb $\cat{Set}$--- mentre que l'esquelet $\{\underline{n}\mid n\geq 0\}$ és \emph{petit}, ja que té un conjunt (de fet, numerable) d'objectes. Una categoria s'anomena \textbf{essencialment petita}\index{categoria!essencialment petita} si és equivalent a una categoria petita. L'equivalència $\cat{FinSet}\simeq\{\underline{n}\mid n\geq 0\}$ mostra, doncs, que $\cat{FinSet}$ és essencialment petita.
\end{exemple}

\begin{observacio}
Aquest exemple mostra la utilitat pràctica de l'equivalència: permet substituir una categoria gran i redundant per una de més petita i manejable, sense perdre cap informació categòrica. És el reflex, a nivell de categories, del fet que en teoria de categories els objectes només importen llevat d'isomorfisme.
\end{observacio}

\begin{exemple}[Matrius i espais vectorials: equivalència sense igualtat ni isomorfisme]\label{ex:mat-finvect}\index{equivalència de categories!matrius i espais vectorials}\index{Mat@$\cat{Mat}_{\mathbb K}$}
Fixem un cos $\mathbb K$. La categoria $\cat{Mat}_{\mathbb K}$ té per objectes els nombres naturals $n\geq 0$. Els morfismes $m\to n$ són les matrius $n\times m$ amb entrades a $\mathbb K$, la composició és el producte, $B\comp A=BA$, i la identitat de $n$ és la matriu identitat $I_n$; els axiomes són l'associativitat del producte de matrius i la neutralitat de $I_n$. Sigui $\cat{FinVect}_{\mathbb K}$ la subcategoria plena de $\cat{Vect}_{\mathbb K}$ formada pels espais de dimensió finita. El functor
\[
F\colon\cat{Mat}_{\mathbb K}\to\cat{FinVect}_{\mathbb K},\qquad n\mapsto\mathbb K^n,\qquad A\mapsto(x\mapsto Ax),
\]
té els tipus correctes: una matriu $n\times m$ dona una aplicació lineal $\mathbb K^m\to\mathbb K^n$. A més, $F(BA)=F(B)\comp F(A)$ i $F(I_n)=\id_{\mathbb K^n}$. És fidel i ple, perquè tota aplicació lineal $\mathbb K^m\to\mathbb K^n$ és la multiplicació per una única matriu, la que té per columnes les imatges de la base canònica. És essencialment exhaustiu, perquè tot espai de dimensió $n$ és isomorf a $\mathbb K^n$. Pel teorema~\ref{teo:caract-equivalencia}, $F$ és una equivalència: $\cat{Mat}_{\mathbb K}\simeq\cat{FinVect}_{\mathbb K}$. Aquest exemple permet veure els tres nivells de comparació.
\begin{itemize}
\item \emph{No és una igualtat.} Són dues categories diferents: els objectes de l'una són nombres i els de l'altra, espais vectorials.
\item \emph{No és un isomorfisme de categories.} A $\cat{Mat}_{\mathbb K}$, dos objectes isomorfs són iguals, perquè una matriu invertible és quadrada: la categoria és esquelètica. A $\cat{FinVect}_{\mathbb K}$, en canvi, $\mathbb K^2$ i l'espai dels polinomis de grau $\leq 1$ són objectes diferents i isomorfs. Un isomorfisme de categories és bijectiu sobre objectes i preserva i reflecteix els isomorfismes, de manera que conservaria la propietat de ser esquelètica. Per tant, no hi ha cap isomorfisme de categories entre $\cat{Mat}_{\mathbb K}$ i $\cat{FinVect}_{\mathbb K}$, ni aquest $F$ ni cap altre.
\item \emph{El quasiinvers només torna llevat d'isomorfisme.} El pas~1 de la demostració del teorema consisteix aquí a \emph{triar una base} de cada espai $V$, és a dir, un isomorfisme $\varepsilon_V\colon\mathbb K^{\dim V}\to V$. Aleshores $G(V)=\dim V$, i $G(g)$ és la matriu de $g$ en les bases triades. Si per als espais $\mathbb K^n$ triem la base canònica, es té literalment $G\comp F=\id_{\cat{Mat}_{\mathbb K}}$. En canvi, $F(G(V))=\mathbb K^{\dim V}$, que no és $V$ tret que $V$ sigui algun $\mathbb K^n$. Així, $F\comp G$ \emph{no} és igual a $\id_{\cat{FinVect}_{\mathbb K}}$, sinó només isomorf a través de $\varepsilon$. La condició de naturalitat $g\comp\varepsilon_V=\varepsilon_W\comp F(G(g))$ és la coneguda fórmula que relaciona una aplicació lineal amb la seva matriu en unes bases.
\end{itemize}
Un espai vectorial qualsevol no té cap base privilegiada; per això el quasiinvers depèn d'una tria, i el que s'obté és una equivalència i no un isomorfisme. Per a l'àlgebra lineal, la conclusió és la que ja es fa servir a la pràctica: treballar amb matrius no perd cap informació sobre les aplicacions lineals entre espais de dimensió finita, sempre que es recordi que la identificació depèn de les bases.
\end{exemple}

Quines propietats es conserven per equivalència i quines no, i com es demostra que dues categories \emph{no} són equivalents, es tracta a l'apèndix~\refalt{cap:contraexemples}{de contraexemples del llibre complet}, junt amb la relació entre equivalència i esquelets.


\section{Composició horitzontal}\index{transformació natural!composició horitzontal}

La composició vertical apila transformacions entre functors \emph{amb el mateix domini i codomini}. Hi ha una segona manera de compondre, que combina transformacions naturals amb functors i, més generalment, transformacions naturals \enquote{de costat}.

Comencem pel cas mixt, en què componem una transformació natural amb un functor\footnote{Segons el costat pel qual s'afegeix el functor, es parla de composició \enquote{per l'esquerra} i \enquote{per la dreta}; convé advertir que aquesta assignació d'\enquote{esquerra} i \enquote{dreta} no és uniforme a la bibliografia. Aquí la fixem per la posició del functor en la notació: $K\eta$ (functor a l'esquerra) és la composició per l'esquerra, i $\eta H$ (functor a la dreta), la composició per la dreta.}.

\begin{definicio}[Composició amb un functor]\index{transformació natural!composició amb un functor}
Sigui $\eta\colon F\Rightarrow G$ una transformació natural entre functors $F,G\colon\cat{C}\to\cat{D}$.
\begin{itemize}
\item Si $K\colon\cat{D}\to\cat{E}$ és un functor, la \textbf{composició per l'esquerra} $K\eta\colon K\comp F\Rightarrow K\comp G$ té per components $(K\eta)_A=K(\eta_A)$.
\item Si $H\colon\cat{B}\to\cat{C}$ és un functor, la \textbf{composició per la dreta} $\eta H\colon F\comp H\Rightarrow G\comp H$ té per components $(\eta H)_B=\eta_{H(B)}$.
\end{itemize}
\end{definicio}

\begin{proposicio}
Les dues construccions anteriors són transformacions naturals.
\end{proposicio}

\begin{prova}
Per a $K\eta$, apliquem el functor $K$ al quadrat de naturalitat de $\eta$ per a un morfisme $f\colon A\to B$: de $G(f)\comp\eta_A=\eta_B\comp F(f)$, i com que $K$ preserva la composició,
\[
K(G(f))\comp K(\eta_A)=K(\eta_B)\comp K(F(f)),
\]
que és la condició de naturalitat per a $K\eta$ respecte de $K\comp F$ i $K\comp G$. Per a $\eta H$, la naturalitat respecte d'un morfisme $b\colon B\to B'$ de $\cat{B}$ és el quadrat de naturalitat de $\eta$ per al morfisme $H(b)$ de $\cat{C}$.
\end{prova}

Amb aquestes dues operacions podem definir la composició horitzontal general.

\begin{definicio}[Composició horitzontal]\index{transformació natural!composició horitzontal}
Siguin $\eta\colon F\Rightarrow G$ (amb $F,G\colon\cat{C}\to\cat{D}$) i $\theta\colon K\Rightarrow L$ (amb $K,L\colon\cat{D}\to\cat{E}$). La seva \textbf{composició horitzontal} $\theta\ast\eta\colon K\comp F\Rightarrow L\comp G$ té per components
\[
(\theta\ast\eta)_A=\theta_{G(A)}\comp K(\eta_A)=L(\eta_A)\comp\theta_{F(A)}.
\]
\end{definicio}

Diagramàticament, la composició horitzontal juxtaposa les dues transformacions, l'una a continuació de l'altra:
\[
\begin{tikzpicture}[baseline=(A.base)]
  \node (A) at (0,0) {$\cat{C}$};
  \node (B) at (2.9,0) {$\cat{D}$};
  \node (Cc) at (5.8,0) {$\cat{E}$};
  \draw[cat] (A) to[bend left=26] node[above] {$F$} (B);
  \draw[cat] (A) to[bend right=26] node[below] {$G$} (B);
  \draw[cat nat] (1.45,0.2) -- node[right=1.5pt] {$\eta$} (1.45,-0.2);
  \draw[cat] (B) to[bend left=26] node[above] {$K$} (Cc);
  \draw[cat] (B) to[bend right=26] node[below] {$L$} (Cc);
  \draw[cat nat] (4.35,0.2) -- node[right=1.5pt] {$\theta$} (4.35,-0.2);
\end{tikzpicture}
\]

\begin{observacio}\label{obs:comp-horitzontal}
Les dues expressions de $(\theta\ast\eta)_A$ coincideixen precisament pel quadrat de naturalitat de $\theta$ aplicat al morfisme $\eta_A\colon F(A)\to G(A)$ de $\cat{D}$:
\[
\begin{tikzpicture}[baseline=(m.center)]
  \matrix (m) [matrix of math nodes, row sep=2.8em, column sep=4.2em] {
    K(F(A)) & K(G(A)) \\
    L(F(A)) & L(G(A)) \\
  };
  \draw[-{Stealth[scale=1.1]}] (m-1-1) -- node[above] {$K(\eta_A)$} (m-1-2);
  \draw[-{Stealth[scale=1.1]}] (m-1-1) -- node[left] {$\theta_{F(A)}$} (m-2-1);
  \draw[-{Stealth[scale=1.1]}] (m-1-2) -- node[right] {$\theta_{G(A)}$} (m-2-2);
  \draw[-{Stealth[scale=1.1]}] (m-2-1) -- node[below] {$L(\eta_A)$} (m-2-2);
\end{tikzpicture}
\qquad
\theta_{G(A)}\comp K(\eta_A)=L(\eta_A)\comp\theta_{F(A)}.
\]
La composició horitzontal es pot llegir, doncs, component una transformació natural amb un functor a cada costat, en qualsevol dels dos ordres:
\[
\theta\ast\eta=(\theta G)\comp(K\eta)=(L\eta)\comp(\theta F).
\]
La composició vertical i l'horitzontal, juntes, satisfan una llei d'intercanvi que fa de $\cat{Cat}$ una \emph{2-categoria}: en una graella de transformacions naturals com la següent, compondre primer en vertical i després en horitzontal dona el mateix que fer-ho en l'ordre contrari,
\[
(\beta'\comp\alpha')\ast(\beta\comp\alpha)=(\beta'\ast\beta)\comp(\alpha'\ast\alpha).
\]
\[
\begin{tikzpicture}[baseline=(A.base)]
  \node (A) at (0,0) {$\cat{C}$};
  \node (B) at (3.4,0) {$\cat{D}$};
  \node (Cc) at (6.8,0) {$\cat{E}$};
  \draw[cat] (A) to[bend left=44] node[above] {$F$} (B);
  \draw[cat] (A) -- (B);
  \draw[cat] (A) to[bend right=44] node[below] {$H$} (B);
  \node[fill=white,inner sep=1.5pt] at (1.7,0) {$G$};
  \draw[cat nat] (1.7,0.52) -- node[right=1.5pt] {$\alpha$} (1.7,0.27);
  \draw[cat nat] (1.7,-0.27) -- node[right=1.5pt] {$\beta$} (1.7,-0.52);
  \draw[cat] (B) to[bend left=44] node[above] {$F'$} (Cc);
  \draw[cat] (B) -- (Cc);
  \draw[cat] (B) to[bend right=44] node[below] {$H'$} (Cc);
  \node[fill=white,inner sep=1.5pt] at (5.1,0) {$G'$};
  \draw[cat nat] (5.1,0.52) -- node[right=1.5pt] {$\alpha'$} (5.1,0.27);
  \draw[cat nat] (5.1,-0.27) -- node[right=1.5pt] {$\beta'$} (5.1,-0.52);
\end{tikzpicture}
\]
No en necessitarem els detalls, però és el marc en què les transformacions naturals adquireixen tota la seva potència.
\end{observacio}

\begin{resumcap}
\apartat{Idees centrals}
\begin{enumerate}
  \item Una transformació natural és una família de morfismes compatible amb \emph{totes} les fletxes de la categoria de partida: cada quadrat de naturalitat commuta.
  \item Les transformacions naturals es componen verticalment, component a component; amb aquesta composició, els functors $\cat{C}\to\cat{D}$ i les transformacions naturals formen la categoria de functors $\cat{D}^{\cat{C}}$.
  \item Un isomorfisme natural és una transformació natural amb totes les components invertibles. La naturalitat no s'hi pot ometre: una família d'isomorfismes, un per a cada objecte, que no sigui natural no és un isomorfisme natural.
  \item Una equivalència de categories és un functor amb un quasiinvers, llevat d'isomorfismes naturals. En el marc d'elecció adoptat al llibre (apèndix~\ref{cap:mida}), es caracteritza equivalentment per ser un functor ple, fidel i essencialment exhaustiu (teorema~\ref{teo:caract-equivalencia}, amb demostració completa). No és una igualtat ni, en general, un isomorfisme de categories: $\cat{Mat}_{\mathbb K}\simeq\cat{FinVect}_{\mathbb K}$, però no són isomorfes.
  \item La composició horitzontal és un segon mode de composició de transformacions naturals, compatible amb la vertical a través de la llei d'intercanvi.
\end{enumerate}
\apartat{Exemple clau} Els grafs dirigits són els objectes d'una categoria de functors: $\cat{Graph}\cong\cat{Set}^{\cat{P}}$, on els morfismes de grafs són exactament les transformacions naturals. Convé recordar l'orientació: $\cat{P}$ té dues fletxes $s,t\colon E\to V$, de manera que un graf és un functor \emph{covariant} $\cat{P}\to\cat{Set}$; no s'ha de confondre amb un prefeix sobre $\cat{P}$.
\apartat{Errors típics} Oblidar comprovar la naturalitat (que el quadrat commuti per a \emph{tot} morfisme); confondre isomorfisme natural amb igualtat de functors; creure que una equivalència conserva propietats com \enquote{tenir un sol objecte} o \enquote{ser esquelètica}; confondre una equivalència amb un isomorfisme de categories, o esperar que $F\comp G$ sigui literalment la identitat; i barrejar composició vertical i horitzontal de transformacions.
\apartat{Lectura recomanada} \cite[§§7.4--7.10]{awodey} per a la naturalitat, les categories de functors i les equivalències; \cite[§§1.4--1.5]{riehl}; \cite[§1.3]{leinster}, que inclou també les equivalències; \cite[I.4, II.4--II.5 i IV.4]{maclane} per a la formulació clàssica de la naturalitat, de les categories de functors amb la composició horitzontal i de les equivalències; \cite[§§4.2--4.4]{barrwells} per a la naturalitat amb exemples de grafs i llistes i per al càlcul de Godement, que és la composició horitzontal; \cite[§1.4 i §1.5.4]{perrone-notes} per a exemples de naturalitat, una secció sobre què \emph{no} és natural i les equivalències.
\end{resumcap}

\ifpublicacioparcial\else
\chapter[Representabilitat i Yoneda]{\texorpdfstring{Propietats universals,\\ representabilitat i Yoneda}{Propietats universals, representabilitat i Yoneda}}
\label{cap:yoneda}

\begin{objectius}
\apartat{Prerequisits} Categories sobre i sota un objecte i dualitat (cap.~\ref{cap:categories}); homfunctors, bifunctor $\Hom$, prefeixos i la parella $\Free$, $U$ (cap.~\ref{cap:functors}); transformacions i isomorfismes naturals, i categories de functors, en particular la categoria de prefeixos (cap.~\ref{cap:transformacions}). No es pressuposa cap coneixement previ de representabilitat ni del lema de Yoneda.
\apartat{Objectius} En acabar el capítol, el lector hauria de poder:
\begin{itemize}
  \item definir objectes inicials i terminals i demostrar-ne la unicitat corresponent;
  \item construir i llegir categories coma en exemples senzills, i interpretar-hi dades universals com a objectes inicials o terminals;
  \item definir una representació mitjançant un element universal i traduir-la a un isomorfisme natural amb un homfunctor;
  \item distingir el functor representat, l'objecte representant, l'element universal i la representació completa;
  \item definir la immersió de Yoneda, enunciar i demostrar el lema de Yoneda i seguir-ne explícitament la bijecció;
  \item deduir que la immersió de Yoneda és plena i fidel i que un objecte queda determinat, llevat d'isomorfisme, pel seu functor representable.
\end{itemize}
\end{objectius}

\section{Cap a les propietats universals}

Amb les transformacions naturals vam completar, al capítol~\ref{cap:transformacions}, una torre de tres nivells: objectes i morfismes dins d'una categoria, functors entre categories, i transformacions naturals entre functors ---l'estructura que, amb les composicions vertical i horitzontal, fa de $\cat{Cat}$ una 2-categoria. És el marc en què es formulen les construccions més potents de la teoria, i les primeres construccions que estudiarem són les propietats universals.

Als capítols anteriors han aparegut, sempre com a exemples i sense tractar-los en
detall, objectes i construccions definits no per la seva constitució interna sinó
per la manera com es relacionen amb tots els altres. El producte de dos objectes
$A$ i $B$ (subsecció~\ref{subsec:producte-coproducte}) es defineix per una condició
d'aquest tipus: per a cada objecte $X$ i cada parell de morfismes $f\colon X\to A$
i $g\colon X\to B$ existeix un \emph{únic} morfisme $u\colon X\to A\times B$ amb
$\pi_1\comp u=f$ i $\pi_2\comp u=g$, és a dir, que fa commutar el diagrama
\[
\begin{tikzpicture}[baseline=(m.center)]
  \matrix (m) [matrix of math nodes, row sep=2.6em, column sep=2.8em] {
     & X & \\
    A & A\times B & B \\
  };
  \draw[-{Stealth[scale=1.0]}] (m-1-2) -- node[above left] {$f$} (m-2-1);
  \draw[-{Stealth[scale=1.0]}] (m-1-2) -- node[above right] {$g$} (m-2-3);
  \draw[-{Stealth[scale=1.0]},dashed] (m-1-2) -- node[right] {$u$} (m-2-2);
  \draw[-{Stealth[scale=1.0]}] (m-2-2) -- node[below] {$\pi_1$} (m-2-1);
  \draw[-{Stealth[scale=1.0]}] (m-2-2) -- node[below] {$\pi_2$} (m-2-3);
\end{tikzpicture}
\]
El coproducte es defineix per la condició dual, amb les fletxes girades. La categoria lliure $\Free(G)$ la satisfà també: un functor
$\Free(G)\to\cat{D}$ queda determinat, de manera única, per un morfisme de grafs
$G\to U(\cat{D})$ (subsecció~\ref{subsec:free-u}). I els objectes inicial i
terminal, que fins ara només hem anomenat (secció~\ref{sec:dualitat}), en són el
cas més senzill. Aquestes condicions, que quantifiquen sobre \emph{totes} les
fletxes d'una categoria, s'anomenen \textbf{propietats
universals}\index{propietat universal}, i són, potser, la contribució més
característica de la teoria de categories. La seva virtut, que encara no hem demostrat en cap cas, és que determinen la \emph{dada universal} ---l'objecte juntament amb les fletxes que intervenen en la condició--- llevat d'un únic isomorfisme compatible. En particular, l'objecte subjacent queda determinat llevat d'isomorfisme, però pot admetre altres isomorfismes que no respectin la dada universal. Abans de donar-ne la forma general, les estudiem amb detall en dos
casos: el més simple, el dels objectes terminals i inicials, que definim ara, i el
de la construcció lliure, que ja coneixem.

\subsection{El cas més simple: objectes terminals i inicials}

El cas extrem, i el més fàcil de dibuixar, és el d'un objecte al qual s'arriba d'una
\emph{única} manera des de qualsevol altre, o del qual es \emph{surt} d'una única
manera cap a qualsevol altre. Diem que un objecte $T$ és \textbf{terminal} si per a
cada objecte $X$ hi ha exactament una fletxa $X\to T$, i que un objecte $I$ és
\textbf{inicial} si per a cada $X$ n'hi ha exactament una $I\to X$:
\[
\begin{tikzpicture}[baseline=(m.center)]
  \matrix (m) [matrix of math nodes, column sep=3.2em] { X & T \\ };
  \draw[-{Stealth[scale=1.1]}, dashed] (m-1-1) -- node[above] {$\exists!$} (m-1-2);
\end{tikzpicture}
\qquad\qquad
\begin{tikzpicture}[baseline=(m.center)]
  \matrix (m) [matrix of math nodes, column sep=3.2em] { I & X \\ };
  \draw[-{Stealth[scale=1.1]}, dashed] (m-1-1) -- node[above] {$\exists!$} (m-1-2);
\end{tikzpicture}
\]
En aquests diagrames seguim la convenció que ja vam introduir en presentar el producte
(subsecció~\ref{subsec:producte-coproducte}): la \textbf{fletxa discontínua} assenyala
el morfisme l'existència del qual afirma la propietat universal, i $\exists!$ abreuja
\enquote{existeix un únic}; com que aquí el morfisme encara no té nom, hi escrivim
directament $\exists!$. Cada diagrama es llegeix, doncs: per a tot $X$ existeix un únic
morfisme cap a $T$ (respectivament, des de $I$).

A $\cat{Set}$, per exemple, un objecte terminal és qualsevol conjunt d'un sol element i
un objecte inicial és el conjunt buit. El que distingeix aquests objectes no és cap
propietat dels seus elements, sinó la fletxa: la condició els fixa completament, fins
al punt que dos objectes terminals qualssevol són isomorfs \emph{d'una única manera},
com demostrarem tot seguit (proposició~\ref{prop:unicitat-terminal}). Aquesta és, en
miniatura, una propietat universal: una condició sobre \emph{totes} les fletxes cap a
un objecte ---o des d'ell--- que el determina llevat d'un únic isomorfisme. En aquest cas, l'isomorfisme és únic fins i tot com a isomorfisme d'objectes, perquè no hi ha cap dada addicional a preservar: entre dos objectes terminals només hi ha una fletxa en cada sentit. Escrivim
encara $T$ i $I$, i no $1$ i $0$, a propòsit: la notació numèrica dona per fet que
\enquote{l'} objecte terminal està ben determinat, cosa que només ens autoritzem a
afirmar un cop provada aquesta unicitat. A la secció següent la demostrem i, aleshores
sí, batejarem $T$ com a $1$ i $I$ com a $0$.

\subsection{Un exemple ja conegut: la construcció lliure}\label{subsec:construccio-lliure}

Un exemple menys trivial ja ha aparegut: la categoria lliure\index{categoria!lliure}
$\Free(G)$ generada per un graf (seccions~\ref{sec:grafs-lliures}
i~\ref{sec:functors-grafs}), i el seu cas d'un sol vèrtex, el monoide
lliure\index{monoide!lliure}. Allà vam deixar pendent de justificar el nom
\enquote{lliure}; ara podem llegir-lo com una propietat universal.

La propietat posa en joc dues categories diferents, i convé tenir present en
quina treballem a cada moment. D'una banda hi ha $\cat{Graph}$, on viuen els
grafs i els morfismes de grafs; de l'altra, $\cat{Cat}$, on viuen les categories
petites i els functors entre elles. Els functors del
capítol~\ref{cap:functors}, $\Free\colon\cat{Graph}\to\cat{Cat}$ i
$U\colon\cat{Cat}\to\cat{Graph}$, fan de pont entre tots dos nivells; no són
fletxes de cap d'aquestes categories ---són functors entre categories grans---,
sinó la manera de passar de l'una a l'altra. Així, $\Free(G)$ és un objecte de
$\cat{Cat}$, mentre que $U(\Free(G))$, el seu graf subjacent, és un objecte de
$\cat{Graph}$.

La clau és una \textbf{inserció dels generadors}, el morfisme de grafs
\[
\eta_G\colon G\to U(\Free(G)),
\]
que deixa fixos els vèrtexs i envia cada aresta al camí de longitud~1
corresponent. Aquesta inserció té la propietat següent: per a tota categoria
petita $\cat{D}$ i tot morfisme de grafs $f\colon G\to U(\cat{D})$ ---una fletxa
de $\cat{Graph}$, que tria un objecte de $\cat{D}$ per a cada vèrtex i un
morfisme de $\cat{D}$ per a cada aresta, sense cap exigència sobre
composicions---, existeix un \emph{únic} functor $\bar f\colon\Free(G)\to\cat{D}$
---una fletxa de $\cat{Cat}$--- que estén $f$. Com que $f$ i $\bar f$ no viuen a
la mateixa categoria, \enquote{estendre} no es pot expressar comparant-los
directament: primer apliquem $U$ a $\bar f$, cosa que en dona el morfisme de
grafs $U(\bar f)\colon U(\Free(G))\to U(\cat{D})$, i aleshores la condició és la
igualtat $U(\bar f)\comp\eta_G=f$ entre morfismes de $\cat{Graph}$.
\[
\begin{tikzpicture}[baseline=(m.center)]
  \matrix (m) [matrix of math nodes, column sep=3.4em, row sep=2.8em] {
    G & U(\Free(G)) \\
      & U(\cat{D}) \\
  };
  \draw[-{Stealth[scale=1.1]}] (m-1-1) -- node[above] {$\eta_G$} (m-1-2);
  \draw[-{Stealth[scale=1.1]}] (m-1-1) -- node[below left] {$f$} (m-2-2);
  \draw[-{Stealth[scale=1.1]}] (m-1-2) -- node[right] {$U(\bar f)$} (m-2-2);
  \node[anchor=south, font=\small] at ([yshift=0.6em]m.north) {a $\cat{Graph}$};
\end{tikzpicture}
\qquad\qquad
\begin{tikzpicture}[baseline=(m.center)]
  \matrix (m) [matrix of math nodes, row sep=2.8em] {
    \Free(G) \\
    \cat{D} \\
  };
  \draw[-{Stealth[scale=1.1]}, dashed] (m-1-1) -- node[right] {$\exists!\,\bar f$} (m-2-1);
  \node[anchor=south, font=\small] at ([yshift=0.6em]m.north) {a $\cat{Cat}$};
\end{tikzpicture}
\]
El triangle de l'esquerra commuta a $\cat{Graph}$. La fletxa discontínua, la
que la propietat afirma que existeix i és única, és $\bar f$, i viu a $\cat{Cat}$;
la seva imatge per $U$ és la que tanca el triangle.
Aquesta és, precisament, la correspondència que vam establir al
capítol~\ref{cap:functors}: definir un functor \emph{des de} $\Free(G)$ no demana res
més que dir on van els generadors. La construcció $\Free(G)$ no s'ha \emph{definit} per
aquesta propietat ---s'ha bastit amb camins---, però la propietat la
\emph{caracteritza}: qualsevol altra categoria petita $\cat{F}$ dotada d'un morfisme de grafs
$G\to U(\cat{F})$ amb la mateixa condició seria isomorfa a $\Free(G)$, i d'una única manera compatible amb la
inserció. En això consisteix una propietat universal. Dit en el llenguatge que introduirem a la secció de representabilitat, $\Free(G)$ representa el functor
\[
\cat{D}\longmapsto\Hom_{\cat{Graph}}\bigl(G,U(\cat{D})\bigr),
\]
i $\eta_G\in\Hom_{\cat{Graph}}\bigl(G,U(\Free(G))\bigr)$ n'és l'\emph{element universal}. La formulació completa d'aquesta correspondència com a adjunció $\Free\dashv U$ arribarà al capítol d'adjuncions.

En aquest capítol precisem què és una propietat universal, la lliguem amb els
homfunctors $\Hom(A,-)$ i $\Hom(-,B)$ que vam introduir, i culminem amb el
\emph{lema de Yoneda}, que fa explícita la idea que un objecte queda determinat pels
morfismes que hi arriben o que en surten.

\section{Objectes inicials i terminals}

\begin{definicio}\index{objecte!inicial}\index{objecte!terminal}
En una categoria $\cat{C}$:
\begin{itemize}
\item un objecte $I$ és \textbf{inicial} si per a cada objecte $X$ hi ha \emph{exactament un} morfisme $I\to X$;
\item un objecte $T$ és \textbf{terminal} si per a cada objecte $X$ hi ha \emph{exactament un} morfisme $X\to T$.
\end{itemize}
Les dues nocions són duals: un objecte inicial a $\cat{C}$ és un objecte terminal a $\cat{C}\op$.
\end{definicio}

\begin{exemple}
A $\cat{Set}$, l'objecte inicial és el conjunt buit $\varnothing$ ---hi ha una única aplicació $\varnothing\to X$, la buida--- i l'objecte terminal és qualsevol conjunt d'un sol element $\{\ast\}$ ---hi ha una única aplicació $X\to\{\ast\}$. A $\cat{Grp}$, en canvi, el grup trivial és alhora inicial i terminal (un tal objecte s'anomena \emph{zero}). En un preordre vist com a categoria, un objecte inicial és un \emph{mínim global}, un element $m$ amb $m\le x$ per a tot $x$, i un de terminal, un \emph{màxim global}. No s'han de confondre amb els elements minimals o maximals, que només demanen que no hi hagi cap element estrictament per sota o per sobre.
\end{exemple}

La propietat clau, model de totes les propietats universals, és la unicitat.

\begin{proposicio}\label{prop:unicitat-terminal}\index{objecte!unicitat llevat d'isomorfisme}
Si una categoria té dos objectes terminals, aleshores són isomorfs, i l'isomorfisme entre ells és únic. El mateix val per als objectes inicials.
\end{proposicio}

\begin{prova}
Siguin $T$ i $T'$ dos objectes terminals. Com que $T'$ és terminal, hi ha un únic morfisme $f\colon T\to T'$; com que $T$ és terminal, hi ha un únic morfisme $g\colon T'\to T$. La composició $g\comp f\colon T\to T$ és un morfisme de $T$ a $T$; però $T$ és terminal, de manera que l'únic morfisme $T\to T$ és $\id_T$, i per tant $g\comp f=\id_T$. Anàlogament $f\comp g=\id_{T'}$. Així $f$ és un isomorfisme, i és únic perquè ho és el morfisme $T\to T'$. El cas inicial és dual.
\end{prova}

Provada la unicitat llevat d'un únic isomorfisme, ja té sentit parlar de \enquote{l'} objecte terminal i \enquote{l'} objecte inicial. D'ara endavant els escrivim, respectivament, $1$ i $0$. Aquesta notació és estàndard i no pressuposa cap finitud: prové del fet que l'objecte terminal és el producte buit ---la unitat de $\times$--- i l'inicial el coproducte buit ---el zero de $+$---, com veurem al capítol~\ref{cap:limits}.

\begin{observacio}
Aquesta demostració és el patró que retrobarem sempre, però convé enunciar-lo amb la precisió justa. Dues \emph{dades universals} del mateix tipus ---dos objectes terminals, dos productes $(P,p_1,p_2)$ i $(Q,q_1,q_2)$, dues representacions--- són isomorfes mitjançant un \emph{únic} isomorfisme compatible amb \emph{tota} la dada universal (les projeccions, la representació, etc.). Els objectes \emph{subjacents} són, per tant, isomorfs; però l'isomorfisme entre els objectes \emph{nus}, un cop oblidada la dada, no ha de ser únic: un producte $P$ pot admetre automorfismes no trivials que no respecten les projeccions. En el cas dels objectes terminals i inicials no hi ha dada addicional a preservar, i per això l'isomorfisme és únic sense més (proposició~\ref{prop:unicitat-terminal}). Amb aquesta cautela parlem de \enquote{l'} objecte terminal, \enquote{el} producte, etc.: la unicitat \emph{compatible amb la dada} és el que dona sentit a caracteritzar objectes per la seva propietat universal en lloc de per una construcció concreta.
\end{observacio}

\section{Categories coma}\label{sec:coma}\index{categoria!coma}

Abans de continuar val la pena introduir una construcció que unifica diverses coses que hem trobat de passada ---les categories sobre i sota un objecte, i ben aviat els cons sobre un diagrama--- i que permet dir amb precisió què vol dir que un objecte és \enquote{universal}: n'hi haurà prou que sigui \emph{inicial} o \emph{terminal} en una categoria adequada.

\begin{definicio}[Categoria coma]\index{categoria!coma}
\label{def:coma}
Siguin $S\colon\cat{A}\to\cat{C}$ i $T\colon\cat{B}\to\cat{C}$ dos functors amb el mateix codomini. La \textbf{categoria coma} $(S\downarrow T)$ té:
\begin{itemize}
  \item per \textbf{objectes}, les ternes $(A,B,f)$ amb $A$ objecte de $\cat{A}$, $B$ objecte de $\cat{B}$ i $f\colon S(A)\to T(B)$ un morfisme de $\cat{C}$;
  \item per \textbf{morfismes} $(A,B,f)\to(A',B',f')$, les parelles $(a,b)$ amb $a\colon A\to A'$ a $\cat{A}$ i $b\colon B\to B'$ a $\cat{B}$ tals que el quadrat
  \[
  \begin{tikzpicture}[baseline=(m.center)]
    \matrix (m) [matrix of math nodes, column sep=3em, row sep=2.4em] {
      S(A) & S(A') \\
      T(B) & T(B') \\
    };
    \draw[-{Stealth[scale=1.1]}] (m-1-1) -- node[above] {$S(a)$} (m-1-2);
    \draw[-{Stealth[scale=1.1]}] (m-1-1) -- node[left] {$f$} (m-2-1);
    \draw[-{Stealth[scale=1.1]}] (m-1-2) -- node[right] {$f'$} (m-2-2);
    \draw[-{Stealth[scale=1.1]}] (m-2-1) -- node[below] {$T(b)$} (m-2-2);
  \end{tikzpicture}
  \]
  commuta a $\cat{C}$, és a dir, $T(b)\comp f=f'\comp S(a)$.
\end{itemize}
La composició i les identitats es fan component a component.
\end{definicio}

\begin{observacio}[Sobre el nom i la notació]\index{categoria!coma!notació}
El concepte i el nom es deuen a Lawvere, que a la seva tesi doctoral (1963)
escrivia aquesta categoria $(S,T)$, per analogia amb la notació $(A,B)$ dels
homconjunts: els objectes de la categoria coma són, en efecte, morfismes de
$S(A)$ a $T(B)$, una mena d'\enquote{homconjunt entre functors}. D'aquí ve el
nom de \emph{categoria coma}, que ha sobreviscut tot i que la coma, usada com a
operador, es prestava a confusió i ha estat substituïda. La notació
$(S\downarrow T)$, que és la de Mac Lane~\cite{maclane} i la que seguim, conserva
els parèntesis de l'original, que a més delimiten els operands quan són
expressions compostes. La fletxa es pot llegir com un dibuix dels objectes: al
quadrat anterior, cada objecte $f$ és una fletxa que \emph{baixa} de la fila de
$S$ a la fila de $T$. També es troben les notacions $S\downarrow T$, sense
parèntesis, i $(S/T)$, que generalitza la notació $\cat{C}/A$ de les categories
sobre un objecte (exemple següent).
\end{observacio}

La notació més freqüent fixa un dels functors en un objecte. Si $\cat{1}$ és la categoria terminal (un objecte i la seva identitat) i $C$ és un objecte de $\cat{C}$, escrivim $C$ també per al functor $\cat{1}\to\cat{C}$ que el tria. Aleshores:

\begin{exemple}[Categories sobre i sota un objecte com a categories coma]\index{categoria!sobre un objecte}\index{categoria!sota un objecte}\label{ex:slice-coma}
Prenem $S=\id_{\cat{C}}$ i $T=C\colon\cat{1}\to\cat{C}$. Un objecte de $(\id_{\cat{C}}\downarrow C)$ és una terna $(A,\ast,f)$ amb $f\colon A\to C$, és a dir, simplement un morfisme cap a $C$. Un morfisme $(A,\ast,f)\to(A',\ast,f')$ és una parella $(a,\id_\ast)$ amb $a\colon A\to A'$ tal que $f'\comp a=f$ (el quadrat de la definició, amb $T(\id_\ast)=\id_C$), és a dir, un triangle commutatiu sobre $C$:
\[
\begin{tikzpicture}[baseline=(m.center)]
  \matrix (m) [matrix of math nodes, column sep=1.6em, row sep=2.6em] {
    A & & A' \\
      & C &  \\
  };
  \draw[-{Stealth[scale=1.1]}] (m-1-1) -- node[above] {$a$} (m-1-3);
  \draw[-{Stealth[scale=1.1]}] (m-1-1) -- node[below left] {$f$} (m-2-2);
  \draw[-{Stealth[scale=1.1]}] (m-1-3) -- node[below right] {$f'$} (m-2-2);
\end{tikzpicture}
\]
Recuperem exactament la categoria $\cat{C}/C$ de $\cat{C}$ sobre $C$. Dualment, $(C\downarrow\id_{\cat{C}})$ és la categoria $C/\cat{C}$ de $\cat{C}$ sota $C$, la dels morfismes que \emph{surten} de $C$. Les categories sobre i sota un objecte són, doncs, els casos més simples de categoria coma.
\end{exemple}

\begin{exemple}[Objectes universals com a inicials o terminals]\label{ex:universal-coma}
Fem precís, en el cas més senzill, l'eslògan que motiva aquesta secció. Fixem
un objecte $C$ d'una categoria localment petita $\cat{C}$ i considerem el
prefeix $\Hom(-,C)\colon\cat{C}\op\to\cat{Set}$ (definició~\ref{def:homfunctors}),
que envia cada objecte $A$ a $\Hom(A,C)$ i cada morfisme $h\colon A\to A'$ a la
precomposició, que amb la notació~\ref{not:hom-morfismes} escrivim
$\Hom(h,C)=h^*$:
\[
\Hom(h,C)\colon\Hom(A',C)\to\Hom(A,C),\qquad x'\mapsto x'\comp h.
\]
La categoria $\cat{C}/C=(\id_{\cat{C}}\downarrow C)$ de l'exemple anterior es
descriu només amb aquest functor. Identificant cada terna $(A,\ast,x)$ amb la
parella $(A,x)$:
\begin{itemize}
  \item un \textbf{objecte} és una parella $(A,x)$ amb $A$ objecte de $\cat{C}$ i
  $x\in\Hom(A,C)$;
  \item un \textbf{morfisme} $(A,x)\to(A',x')$ és un morfisme $h\colon A\to A'$
  de $\cat{C}$ tal que
  \[
    \Hom(h,C)(x')=x,
    \qquad\text{és a dir,}\qquad
    x'\comp h=x,
  \]
  que és el triangle commutatiu sobre $C$:
  \[
  \begin{tikzpicture}[baseline=(m.center)]
    \matrix (m) [matrix of math nodes, column sep=1.6em, row sep=2.6em] {
      A & & A' \\
        & C &  \\
    };
    \draw[-{Stealth[scale=1.1]}] (m-1-1) -- node[above] {$h$} (m-1-3);
    \draw[-{Stealth[scale=1.1]}] (m-1-1) -- node[below left] {$x$} (m-2-2);
    \draw[-{Stealth[scale=1.1]}] (m-1-3) -- node[below right] {$x'$} (m-2-2);
  \end{tikzpicture}
  \]
\end{itemize}

\emph{Afirmació: $(C,\id_C)$ és un objecte terminal de $\cat{C}/C$.} Sigui
$(A,x)$ un objecte qualsevol. Un morfisme $(A,x)\to(C,\id_C)$ és un
$h\colon A\to C$ amb $\Hom(h,C)(\id_C)=x$. Com que
$\Hom(h,C)(\id_C)=\id_C\comp h=h$, la condició diu exactament $h=x$. Per tant
existeix un morfisme $(A,x)\to(C,\id_C)$, i només un: $h=x$.

Llegit en termes del functor, això diu que per a cada objecte $A$ l'aplicació
\[
\Hom(A,C)\longrightarrow\Hom(A,C),\qquad h\longmapsto\Hom(h,C)(\id_C),
\]
és bijectiva (de fet, és la identitat): tot element $x$ de qualsevol $\Hom(A,C)$
s'obté d'un \emph{únic} element distingit, $\id_C\in\Hom(C,C)$, aplicant-hi
$\Hom(h,C)$ per a un \emph{únic} $h$. Aquesta és precisament la forma de la
definició~\ref{def:representable} de la secció següent: $(C,\id_C)$ és una
representació de $\Hom(-,C)$, amb element universal $\id_C$. Allà veurem
(observació~\ref{obs:representacio-inicial}) que, per a qualsevol prefeix $P$
sobre $\cat{C}$ (definició~\ref{def:prefeix}), ser una representació de $P$ equival a ser un
objecte inicial d'una certa categoria coma construïda a partir de $P$; per a
$P=\Hom(-,C)$ aquesta categoria és $(\cat{C}/C)\op$, i ser-hi inicial és ser
terminal a $\cat{C}/C$.

Dualment, amb el homfunctor covariant $\Hom(C,-)$, $(C,\id_C)$ és un objecte
inicial de $C/\cat{C}$: un morfisme $(C,\id_C)\to(A,x)$ és un $h\colon C\to A$
amb $\Hom(C,h)(\id_C)=h\comp\id_C=x$, és a dir, $h=x$.

Un cas menys trivial és la construcció lliure de la
subsecció~\ref{subsec:construccio-lliure}. Sigui $G$ un graf, vist com a functor
$G\colon\cat{1}\to\cat{Graph}$, i sigui $U\colon\cat{Cat}\to\cat{Graph}$ el
functor oblit. Un objecte de la categoria coma $(G\downarrow U)$ és una parella
$(\cat{D},f)$ formada per una categoria petita $\cat{D}$ i un morfisme de grafs
$f\colon G\to U(\cat{D})$; un morfisme $(\cat{D},f)\to(\cat{D}',f')$ és un
functor $F\colon\cat{D}\to\cat{D}'$ amb $U(F)\comp f=f'$. Que
$(\Free(G),\eta_G)$ sigui un objecte inicial de $(G\downarrow U)$ vol dir que
per a cada $(\cat{D},f)$ hi ha un únic functor $\bar f\colon\Free(G)\to\cat{D}$
amb $U(\bar f)\comp\eta_G=f$: és, literalment, la propietat universal que hi
vam enunciar.

Aquests casos fan precís l'eslògan que unifica les propietats universals:
\emph{ser universal és ser inicial o terminal en una categoria coma adient}. Ho
retrobarem al capítol següent: un límit és un objecte terminal de la categoria dels
cons, que és una categoria coma, i un colímit, un objecte inicial de la dels cocons
(observació~\ref{obs:cons-coma}).
\end{exemple}

\section{Propietats universals com a representabilitat}

Podem reformular aquestes idees amb el llenguatge dels functors representables. Recordem que, en una categoria localment petita, cada objecte $A$ determina el prefeix $\Hom(-,A)\colon\cat{C}\op\to\cat{Set}$ i el functor covariant $\Hom(A,-)\colon\cat{C}\to\cat{Set}$.

Observem primer com es llegeixen les nocions anteriors. Que $1$ sigui terminal vol dir que $\Hom(X,1)$ té exactament un element per a cada $X$; és a dir, que el functor $\Hom(-,1)$ és \emph{naturalment isomorf} al prefeix constant $\Delta_{\{\ast\}}\colon\cat{C}\op\to\cat{Set}$ (exemple~\ref{ex:functor-constant}). Les components de l'isomorfisme són les úniques bijeccions $\Hom(X,1)\to\{\ast\}$, i la naturalitat és automàtica, perquè tots els quadrats acaben en un conjunt d'un sol element. Dualment, $0$ és inicial quan $\Hom(0,-)$ és naturalment isomorf al functor constant $\Delta_{\{\ast\}}\colon\cat{C}\to\cat{Set}$. En cap dels dos casos no és una igualtat literal: els conjunts d'un sol element $\Hom(X,1)$, i també els $\Hom(0,X)$, són conjunts diferents per a objectes $X$ diferents, i no s'han identificat per definició amb $\{\ast\}$.

Una forma general de les propietats universals que estudiarem és la representabilitat d'un functor cap a $\cat{Set}$: un objecte, juntament amb una dada universal, satisfà una propietat d'aquesta mena quan \emph{representa} el functor que recull les dades del problema.

\begin{definicio}[Representació]\index{propietat universal}\index{functor!representable}\index{element universal}
\label{def:representable}
Sigui $\cat{C}$ una categoria localment petita.
\begin{enumerate}
\item Sigui $F\colon\cat{C}\to\cat{Set}$ un functor covariant. Una \textbf{representació} de $F$ és una parella $(A,u)$ formada per un objecte $A$ de $\cat{C}$ i un element $u\in F(A)$, l'\textbf{element universal}, tal que per a cada objecte $X$ i cada $x\in F(X)$ existeix un únic morfisme $f\colon A\to X$ amb
\[
F(f)(u)=x.
\]
\item Sigui $P$ un prefeix sobre $\cat{C}$ (definició~\ref{def:prefeix}). Una \textbf{representació} de $P$ és una parella $(A,u)$ formada per un objecte $A$ de $\cat{C}$ i un element $u\in P(A)$, l'\textbf{element universal}, tal que per a cada objecte $X$ i cada $x\in P(X)$ existeix un únic morfisme $f\colon X\to A$ amb
\[
P(f)(u)=x.
\]
\end{enumerate}
En tots dos casos es diu que el functor és \textbf{representable}, que $A$ n'és l'\textbf{objecte representant} i que $(A,u)$ n'és una representació.
\end{definicio}

Amb la noció d'isomorfisme natural (capítol~\ref{cap:transformacions}), la representabilitat es pot formular sense esmentar l'element universal.

\begin{proposicio}[Representabilitat com a isomorfisme natural]\label{prop:representable-iso}\index{functor!representable}
Sigui $\cat{C}$ una categoria localment petita.
\begin{enumerate}
\item Un functor $F\colon\cat{C}\to\cat{Set}$ és representable si i només si és naturalment isomorf al functor $\Hom_{\cat{C}}(S,-)\colon\cat{C}\to\cat{Set}$ per a algun objecte $S$ de $\cat{C}$. En aquest cas, $S$ és un objecte representant de $F$.
\item Un prefeix $P\colon\cat{C}\op\to\cat{Set}$ és representable si i només si és naturalment isomorf al prefeix $\Hom_{\cat{C}}(-,S)\colon\cat{C}\op\to\cat{Set}$ per a algun objecte $S$ de $\cat{C}$. De nou, en aquest cas $S$ és un objecte representant de $P$.
\end{enumerate}
Més precisament, les representacions $(S,u)$ es corresponen amb els isomorfismes naturals
\[
\alpha\colon\Hom(S,-)\Rightarrow F
\qquad\text{o bé}\qquad
\alpha\colon\Hom(-,S)\Rightarrow P,
\]
segons el cas, mitjançant $u=\alpha_S(\id_S)$ i $\alpha_X(f)=F(f)(u)$, respectivament $\alpha_X(f)=P(f)(u)$.
\end{proposicio}

\begin{prova}
Fem primer el cas~(1). Els quadrats de naturalitat segueixen el criteri del capítol~\ref{cap:transformacions}: l'acció dels functors sobre un morfisme, en horitzontal (a dalt l'homfunctor, a baix $F$), i les components de la transformació, en vertical. Al costat de cada quadrat hi ha el camí que hi fa un element: les fletxes $\mapsto$ indiquen on va a parar cada element.

\emph{D'una representació a un isomorfisme natural.} Sigui $(S,u)$ una representació de $F$ i definim $\alpha_X\colon\Hom(S,X)\to F(X)$ per $\alpha_X(f)=F(f)(u)$. Per veure que $\alpha\colon\Hom(S,-)\Rightarrow F$ és natural, cal que per a cada $g\colon X\to Y$ commuti el quadrat següent, on l'acció de l'homfunctor sobre $g$ és la postcomposició $g_*=\Hom(S,g)$ (notació~\ref{not:hom-morfismes}):
\[
\begin{tikzpicture}[baseline=(m.center)]
  \matrix (m) [matrix of math nodes, row sep=2.8em, column sep=3.6em] {
    \Hom(S,X) & \Hom(S,Y) \\
    F(X) & F(Y) \\
  };
  \draw[-{Stealth[scale=1.1]}] (m-1-1) -- node[above] {$g_*$} (m-1-2);
  \draw[-{Stealth[scale=1.1]}] (m-1-1) -- node[left] {$\alpha_X$} (m-2-1);
  \draw[-{Stealth[scale=1.1]}] (m-1-2) -- node[right] {$\alpha_Y$} (m-2-2);
  \draw[-{Stealth[scale=1.1]}] (m-2-1) -- node[below] {$F(g)$} (m-2-2);
\end{tikzpicture}
\qquad\qquad
\begin{tikzpicture}[baseline=(m.center)]
  \matrix (m) [matrix of math nodes, row sep=2.8em, column sep=2.4em] {
    f & g\comp f \\
    F(f)(u) & F(g)\bigl(F(f)(u)\bigr) \\
  };
  \draw[{Bar[width=4pt]}-{Stealth[scale=1.1]}] (m-1-1) -- (m-1-2);
  \draw[{Bar[width=4pt]}-{Stealth[scale=1.1]}] (m-1-1) -- (m-2-1);
  \draw[{Bar[width=4pt]}-{Stealth[scale=1.1]}] (m-1-2) -- (m-2-2);
  \draw[{Bar[width=4pt]}-{Stealth[scale=1.1]}] (m-2-1) -- (m-2-2);
\end{tikzpicture}
\]
Seguint un element $f$ pels dos camins, per dalt i després avall s'obté $\alpha_Y(g\comp f)=F(g\comp f)(u)$, i avall i després per baix s'obté $F(g)\bigl(F(f)(u)\bigr)$. Totes dues coincideixen perquè $F$ preserva la composició:
\[
\alpha_Y(g\comp f)=F(g\comp f)(u)=F(g)\bigl(F(f)(u)\bigr)=F(g)\bigl(\alpha_X(f)\bigr).
\]
La condició de representació diu exactament que cada $\alpha_X$ és bijectiva, i per la proposició~\ref{prop:iso-natural-components} $\alpha$ és un isomorfisme natural.

\emph{D'un isomorfisme natural a una representació.} Recíprocament, sigui $\alpha\colon\Hom(S,-)\Rightarrow F$ un isomorfisme natural i posem $u=\alpha_S(\id_S)\in F(S)$. Per a cada $f\colon S\to X$, escrivim el quadrat de naturalitat de $\alpha$ respecte de $f$ i hi seguim l'element $\id_S$:
\[
\begin{tikzpicture}[baseline=(m.center)]
  \matrix (m) [matrix of math nodes, row sep=2.8em, column sep=3.6em] {
    \Hom(S,S) & \Hom(S,X) \\
    F(S) & F(X) \\
  };
  \draw[-{Stealth[scale=1.1]}] (m-1-1) -- node[above] {$f_*$} (m-1-2);
  \draw[-{Stealth[scale=1.1]}] (m-1-1) -- node[left] {$\alpha_S$} (m-2-1);
  \draw[-{Stealth[scale=1.1]}] (m-1-2) -- node[right] {$\alpha_X$} (m-2-2);
  \draw[-{Stealth[scale=1.1]}] (m-2-1) -- node[below] {$F(f)$} (m-2-2);
\end{tikzpicture}
\qquad\qquad
\begin{tikzpicture}[baseline=(m.center)]
  \matrix (m) [matrix of math nodes, row sep=2.8em, column sep=2.4em] {
    \id_S & f \\
    u & F(f)(u) \\
  };
  \draw[{Bar[width=4pt]}-{Stealth[scale=1.1]}] (m-1-1) -- (m-1-2);
  \draw[{Bar[width=4pt]}-{Stealth[scale=1.1]}] (m-1-1) -- (m-2-1);
  \draw[{Bar[width=4pt]}-{Stealth[scale=1.1]}] (m-1-2) -- (m-2-2);
  \draw[{Bar[width=4pt]}-{Stealth[scale=1.1]}] (m-2-1) -- (m-2-2);
\end{tikzpicture}
\]
Per dalt, $\id_S$ va a $f\comp\id_S=f$ i després a $\alpha_X(f)$; avall, va a $u$ i després a $F(f)(u)$. La commutativitat del quadrat dona, doncs,
\[
\alpha_X(f)=\alpha_X(f\comp\id_S)=F(f)\bigl(\alpha_S(\id_S)\bigr)=F(f)(u).
\]
Això diu que tota la transformació $\alpha$ queda determinada pel sol element $u$. Com que $\alpha_X$ és bijectiva (proposició~\ref{prop:iso-natural-components}), per a cada $x\in F(X)$ hi ha un únic $f\colon S\to X$ amb $F(f)(u)=x$: $(S,u)$ és una representació de $F$.

Les dues construccions són inverses l'una de l'altra. Partint de $u$, la transformació $\alpha$ torna $\alpha_S(\id_S)=F(\id_S)(u)=u$; i partint de $\alpha$, la igualtat anterior diu que $\alpha$ és precisament la transformació $f\mapsto F(f)(u)$ construïda a partir de $u=\alpha_S(\id_S)$.

El cas~(2) és el cas~(1) aplicat a $\cat{C}\op$, però val la pena veure com canvia el diagrama clau. Ara $\alpha\colon\Hom(-,S)\Rightarrow P$, els morfismes arriben a $S$ en lloc de sortir-ne, i l'acció de l'homfunctor sobre un morfisme és la precomposició $f^*=\Hom(f,S)$. Per a cada $f\colon X\to S$, el quadrat de naturalitat respecte de $f$, amb l'element $\id_S$, és
\[
\begin{tikzpicture}[baseline=(m.center)]
  \matrix (m) [matrix of math nodes, row sep=2.8em, column sep=3.6em] {
    \Hom(S,S) & \Hom(X,S) \\
    P(S) & P(X) \\
  };
  \draw[-{Stealth[scale=1.1]}] (m-1-1) -- node[above] {$f^*$} (m-1-2);
  \draw[-{Stealth[scale=1.1]}] (m-1-1) -- node[left] {$\alpha_S$} (m-2-1);
  \draw[-{Stealth[scale=1.1]}] (m-1-2) -- node[right] {$\alpha_X$} (m-2-2);
  \draw[-{Stealth[scale=1.1]}] (m-2-1) -- node[below] {$P(f)$} (m-2-2);
\end{tikzpicture}
\qquad\qquad
\begin{tikzpicture}[baseline=(m.center)]
  \matrix (m) [matrix of math nodes, row sep=2.8em, column sep=2.4em] {
    \id_S & f \\
    u & P(f)(u) \\
  };
  \draw[{Bar[width=4pt]}-{Stealth[scale=1.1]}] (m-1-1) -- (m-1-2);
  \draw[{Bar[width=4pt]}-{Stealth[scale=1.1]}] (m-1-1) -- (m-2-1);
  \draw[{Bar[width=4pt]}-{Stealth[scale=1.1]}] (m-1-2) -- (m-2-2);
  \draw[{Bar[width=4pt]}-{Stealth[scale=1.1]}] (m-2-1) -- (m-2-2);
\end{tikzpicture}
\]
on ara, per dalt, $\id_S$ va a $\id_S\comp f=f$. Com que els functors són contravariants, les fletxes horitzontals van de l'objecte $S$ a l'objecte $X$, en sentit contrari a $f\colon X\to S$. Se'n dedueix $\alpha_X(f)=P(f)(u)$, i la resta de l'argument és el mateix.
\end{prova}

Les dues nocions són duals l'una de l'altra. Un prefeix sobre $\cat{C}$ és un functor covariant $\cat{C}\op\to\cat{Set}$, i com que un morfisme $X\to A$ de $\cat{C}$ és un morfisme $A\to X$ de $\cat{C}\op$, una representació d'un prefeix en el sentit~(2) és exactament una representació en el sentit~(1) del functor covariant corresponent sobre $\cat{C}\op$. Per això, tot el que demostrem per a un dels dos casos val, per dualitat, per a l'altre. Per exemple, $0$ és inicial si i només si $(0,\ast)$ és una representació del functor covariant $\Delta_{\{\ast\}}$, i $1$ és terminal si i només si $(1,\ast)$ és una representació del prefeix $\Delta_{\{\ast\}}$.

\begin{observacio}[Sondes i pantalles]\label{obs:sondes-pantalles}\index{functor!representable!interpretació}
Els functors representables admeten una lectura intuïtiva. El functor representable $\Hom_{\cat{C}}(S,-)\colon\cat{C}\to\cat{Set}$ pren un objecte $X$ i en dona el conjunt $\Hom_{\cat{C}}(S,X)$ de fletxes de $S$ a $X$. Les característiques de $X$ que extreu són exactament totes les maneres possibles de fer arribar $S$ dins de $X$, com si $S$ fos una \emph{sonda} amb què explorem l'estructura de $X$. Dir que un functor $F$ és representable per $S$ és dir que el que $F$ mesura de cada objecte és precisament el que en veu aquesta sonda.

Dualment, el prefeix representable $\Hom_{\cat{C}}(-,S)\colon\cat{C}\op\to\cat{Set}$ pren un objecte $X$ i en dona el conjunt $\Hom_{\cat{C}}(X,S)$ de fletxes de $X$ a $S$. Les característiques de $X$ que extreu són totes les maneres possibles d'enviar $X$ a $S$, o totes les observacions de $X$ amb valors a $S$: podem pensar $S$ com una \emph{pantalla} sobre la qual $X$ es projecta de moltes maneres diferents.
\end{observacio}

Vegem-ne exemples. En tots, l'isomorfisme natural es pot comprovar directament o, equivalentment, n'hi ha prou de donar l'element universal (proposició~\ref{prop:representable-iso}).

\begin{exemple}[Vèrtexs i arestes d'un graf]\label{ex:rep-grafs}\index{functor!representable!a Graph}
A l'exemple~\ref{ex:homfunctor-grafs} vam veure que els functors $\cat{Graph}\to\cat{Set}$ \enquote{conjunt de vèrtexs} i \enquote{conjunt d'arestes} són naturalment isomorfs a $\Hom_{\cat{Graph}}(\mathbf V,-)$ i a $\Hom_{\cat{Graph}}(\mathbf E,-)$. Són, doncs, representables, amb objectes representants el graf $\mathbf V$ d'un sol vèrtex i el graf $\mathbf E$ d'una sola aresta. Els elements universals són l'únic vèrtex $\ast$ de $\mathbf V$ i l'única aresta $e$ de $\mathbf E$: tot vèrtex d'un graf $G$ és la imatge de $\ast$ per un únic morfisme $\mathbf V\to G$, i tota aresta és la imatge de $e$ per un únic morfisme $\mathbf E\to G$. Són sondes en el sentit literal: $\mathbf V$ detecta els vèrtexs i $\mathbf E$ detecta les arestes.
\end{exemple}

\begin{exemple}[Functors oblit de $\cat{Grp}$, $\cat{Mon}$ i $\cat{Vect}_{\mathbb K}$]\label{ex:rep-grp}\index{functor!representable!oblit}
El functor oblit $U\colon\cat{Grp}\to\cat{Set}$ és representable, amb objecte representant el grup additiu $\mathbb Z$ i element universal $1\in U(\mathbb Z)$. En efecte, un homomorfisme $\varphi\colon\mathbb Z\to G$ queda determinat per $g=\varphi(1)$, perquè
\[
\varphi(n)=g^n\qquad\text{per a tot } n\in\mathbb Z;
\]
i, recíprocament, per a cada $g\in G$ l'aplicació $n\mapsto g^n$ és un homomorfisme. Per tant, $\varphi\mapsto\varphi(1)$ és una bijecció
\[
\Hom_{\cat{Grp}}(\mathbb Z,G)\;\cong\;U(G),
\]
natural en $G$, perquè per a un homomorfisme $\psi\colon G\to H$ es té $(\psi\comp\varphi)(1)=\psi(\varphi(1))$.

El mateix argument val per als altres dos models algebraics del capítol~\ref{cap:categories}. El functor oblit $\cat{Mon}\to\cat{Set}$ és representable pel monoide $(\mathbb N,+,0)$ amb element universal $1$: un homomorfisme de monoides $h\colon\mathbb N\to M$ queda determinat per $m=h(1)$, perquè $h(n)=m\ast\cdots\ast m$ ($n$ vegades, i $h(0)$ el neutre), i per a cada $m$ aquesta fórmula defineix un homomorfisme. El functor oblit $\cat{Vect}_{\mathbb K}\to\cat{Set}$ és representable per $\mathbb K$, vist com a espai vectorial sobre si mateix, amb element universal $1$: una aplicació lineal $T\colon\mathbb K\to V$ queda determinada per $v=T(1)$, perquè $T(\lambda)=\lambda v$, i per a cada $v\in V$ l'aplicació $\lambda\mapsto\lambda v$ és lineal.

En els tres casos l'objecte representant és l'estructura \enquote{generada lliurement per un sol element}, i com a sonda veu exactament els elements de cada objecte.
\end{exemple}

\begin{exemple}[El punt a $\cat{Set}$]\label{ex:rep-punt}
El functor identitat de $\cat{Set}$ és representable per un conjunt d'un sol element $\{\ast\}$, amb element universal $\ast$: una aplicació $\{\ast\}\to X$ és el mateix que un element de $X$, i $\Hom_{\cat{Set}}(\{\ast\},X)\cong X$ naturalment en $X$. El punt és la sonda que veu els elements d'un conjunt.
\end{exemple}

\begin{exemple}[Subconjunts i valors de veritat]\label{ex:rep-parts}\index{functor!representable!de parts}
El functor de parts $\mathcal P\colon\cat{Set}\op\to\cat{Set}$ és un prefeix representable: a l'exemple~\ref{ex:functor-parts} vam veure que les funcions característiques donen bijeccions canòniques $\mathcal P(X)\cong\Hom_{\cat{Set}}(X,\{0,1\})$ sota les quals l'antiimatge correspon a la precomposició. En el llenguatge del capítol~\ref{cap:transformacions}, això diu exactament que $\mathcal P\cong\Hom_{\cat{Set}}(-,\{0,1\})$ naturalment: el quadrat de naturalitat és la igualtat $\chi_B\comp f=\chi_{f^{-1}(B)}$ (exemple~\ref{ex:isos-naturals-cap2}). L'objecte representant és el conjunt $\{0,1\}$ de valors de veritat, i l'element universal és $\{1\}\in\mathcal P(\{0,1\})$, el subconjunt dels valors \enquote{cert}: per a cada $B\subseteq X$, l'única aplicació $\chi\colon X\to\{0,1\}$ amb $\chi^{-1}(\{1\})=B$ és la funció característica $\chi_B$.

Aquí $\{0,1\}$ és la pantalla. Un subconjunt $B\subseteq X$ és el mateix que un predicat sobre $X$, i el predicat és una observació de $X$ amb valors de veritat: $\chi_B(x)=1$ vol dir que $x$ compleix el predicat. Que els subconjunts es puguin classificar així, per morfismes cap a un objecte fix de valors de veritat, és una idea que tornarà amb força al capítol~\ref{cap:topos}, on l'objecte $\{0,1\}$ es generalitza en el \emph{classificador de subobjectes} $\Omega$ (definició~\ref{def:classificador}).
\end{exemple}

\begin{exemple}[El producte com a representació]\label{ex:rep-producte}\index{producte!com a representació}
Siguin $A$ i $B$ objectes d'una categoria localment petita $\cat{C}$. L'assignació
\[
X\longmapsto\Hom(X,A)\times\Hom(X,B),
\qquad
f\longmapsto\bigl((a,b)\mapsto(a\comp f,\,b\comp f)\bigr),
\]
és un prefeix sobre $\cat{C}$: preserva identitats i, per a $f\colon X\to Y$ i $g\colon Y\to Z$, $(a\comp g\comp f,\,b\comp g\comp f)$ s'obté aplicant primer $g$ i després $f$. Una representació d'aquest prefeix és un objecte $Q$ amb un element universal $(p_1,p_2)\in\Hom(Q,A)\times\Hom(Q,B)$ tal que, per a cada $X$ i cada parella $(a,b)$, hi ha un únic $u\colon X\to Q$ amb $p_1\comp u=a$ i $p_2\comp u=b$. És exactament la definició de producte de la subsecció~\ref{subsec:producte-coproducte}: el producte de $A$ i $B$ existeix si i només si aquest prefeix és representable, i un producte és una representació seva.

Aquest exemple té també una lectura lògica. A la categoria prima $\cat{Prov}_T$ associada a un sistema deductiu $T$ (capítol~\ref{cap:categories}), entre dues fórmules hi ha com a màxim un morfisme, i n'hi ha un exactament quan hi ha deduïbilitat $p\vdash_T q$. Per a dues fórmules $p$ i $q$, el prefeix
\[
r\longmapsto\Hom_{\cat{Prov}_T}(r,p)\times\Hom_{\cat{Prov}_T}(r,q)
\]
és, doncs, no buit exactament quan $r\vdash_T p$ i $r\vdash_T q$. Una representació és una fórmula $c$ amb $c\vdash_T p$ i $c\vdash_T q$ (l'element universal) tal que tota fórmula $r$ que dedueix $p$ i $q$ dedueix $c$; la unicitat del morfisme és automàtica. Per transitivitat, això equival a
\[
r\vdash_T c
\quad\Longleftrightarrow\quad
r\vdash_T p\ \text{i}\ r\vdash_T q,
\]
que són les regles d'introducció i d'eliminació de la conjunció. La conjunció $p\wedge q$, quan existeix, és doncs una representació d'aquest prefeix: satisfà la propietat universal del producte, llegida en lògica, tal com anunciava la secció~\ref{sec:dualitat}.
\end{exemple}

\begin{exemple}[Functors que no són representables]\label{ex:no-representable}
Si $\cat{C}$ té algun objecte, el functor constant $\Delta_\varnothing\colon\cat{C}\to\cat{Set}$ no és representable: per a tot objecte $S$, el conjunt $\Hom(S,S)$ conté $\id_S$ i no és buit, mentre que $\Delta_\varnothing(S)=\varnothing$, de manera que no hi pot haver cap bijecció $\Hom(S,S)\to\Delta_\varnothing(S)$. El mateix argument val per al prefeix $\Delta_\varnothing$. Això explica, per exemple, que en una categoria discreta amb dos objectes $A\neq B$ (només hi ha identitats) no hi hagi producte de $A$ i $B$: el prefeix de l'exemple~\ref{ex:rep-producte} hi val $\Hom(X,A)\times\Hom(X,B)=\varnothing$ per a tot $X$, perquè cap objecte té fletxes cap a tots dos, i és, doncs, el prefeix $\Delta_\varnothing$.
\end{exemple}

\begin{observacio}[Representacions com a objectes inicials]\label{obs:representacio-inicial}
Això fa precís l'eslògan de la secció~\ref{sec:coma}. Sigui
$P$ un prefeix sobre $\cat{C}$ i sigui
$\{\ast\}\colon\cat{1}\to\cat{Set}$ el functor que tria un conjunt d'un sol
element. Un objecte de la categoria coma $(\{\ast\}\downarrow P)$ és una terna
$(\ast,A,x)$ amb $x\colon\{\ast\}\to P(A)$, és a dir, un objecte $A$ de
$\cat{C}$ junt amb un element $x\in P(A)$; l'escrivim $(A,x)$. Un morfisme
$(A,x)\to(A',x')$ és un morfisme $A\to A'$ de $\cat{C}\op$ ---és a dir, un
$h\colon A'\to A$ de $\cat{C}$--- que fa commutar el quadrat de la
definició~\ref{def:coma}, condició que aquí diu $P(h)(x)=x'$. Per tant,
$(A,u)$ és un objecte inicial de $(\{\ast\}\downarrow P)$ si i només si per a
cada $(X,x)$ hi ha un únic $h\colon X\to A$ amb $P(h)(u)=x$: exactament la
condició de representació. Per a $P=\Hom(-,C)$, la categoria
$(\{\ast\}\downarrow P)$ és $(\cat{C}/C)\op$, i retrobem
l'exemple~\ref{ex:universal-coma}.
\end{observacio}

\begin{observacio}[La categoria d'elements]\label{obs:categoria-elements}\index{categoria!d'elements}
A la bibliografia és habitual presentar la mateixa construcció amb les fletxes en el sentit de $\cat{C}$. Per a un prefeix $P\colon\cat{C}\op\to\cat{Set}$, la \textbf{categoria d'elements} $\textstyle\int P$ té per objectes les parelles $(A,x)$ amb $x\in P(A)$, i per morfismes $f\colon(A,x)\to(B,y)$ els morfismes $f\colon A\to B$ de $\cat{C}$ amb $P(f)(y)=x$. Té, doncs, les mateixes dades que la categoria coma de l'observació anterior, amb les fletxes girades:
\[
\textstyle\int P\;\cong\;(\{\ast\}\downarrow P)\op.
\]
Per tant, una representació $(C,u)$ de $P$ correspon a un objecte \emph{terminal} de $\int P$: per a cada $(A,x)$ hi ha un únic $f\colon A\to C$ amb $P(f)(u)=x$. És la formulació que fa servir, per exemple, \cite[cap.~2]{riehl}.
\end{observacio}

\begin{observacio}
Cal distingir dos nivells d'unicitat. L'\textbf{objecte} representant $A$ és únic llevat d'isomorfisme, però no necessàriament d'un únic isomorfisme. En canvi, la \textbf{representació} $(A,u)$ ---amb la dada de l'element universal--- és única llevat d'un \emph{únic} isomorfisme compatible amb els elements universals: si $(A,u)$ i $(B,v)$ representen $P$, hi ha un únic isomorfisme $\varphi\colon A\to B$ amb $P(\varphi)(v)=u$. Aquesta distinció, entre la unicitat de l'objecte i la de la dada universal, és essencial i la retrobarem en tractar límits, adjuncions i objectes classificadors. Que dos functors representables $\Hom(-,A)$ i $\Hom(-,B)$ siguin naturalment isomorfs equivaldrà, pel lema de Yoneda, a $A\cong B$; la resta del capítol ho fa precís.
\end{observacio}

\section{La immersió de Yoneda}

Fixem una categoria \textbf{petita} $\cat{C}$. La noció de prefeix, i la de transformació natural entre prefeixos, té sentit sobre qualsevol categoria (definició~\ref{def:prefeix}). Els prefeixos representables $\Hom(-,A)$, en canvi, només són prefeixos si cada $\Hom(X,A)$ és un conjunt, és a dir, si $\cat{C}$ és localment petita, que és la hipòtesi amb què els hem fet servir a la secció de representabilitat. Aquí demanem més: que $\cat{C}$ sigui petita. Aquesta restricció afecta la categoria que aplega els prefeixos. Si $\cat{C}$ no és petita, encara que sigui localment petita, cada prefeix recorre una classe pròpia d'objectes i les transformacions naturals entre dos prefeixos poden formar també una classe pròpia. Aleshores no es pot formar $\cat{Set}^{\cat{C}\op}$ com a categoria (observació~\ref{obs:mida-functors}). Amb $\cat{C}$ petita, en canvi, la categoria de prefeixos $\widehat{\cat{C}}$ i tots els seus homconjunts queden ben definits (exemple~\ref{ex:categoria-prefeixos}).\footnote{Amb universos de Grothendieck, la construcció es pot fer també per a $\cat{C}$ només localment petita, formant $\widehat{\cat{C}}$ en un univers superior; vegeu l'apèndix~\ref{cap:mida} i \cite[§1.1]{riehl}.} A cada objecte $A$ li hem associat el prefeix representable $\Hom(-,A)\colon\cat{C}\op\to\cat{Set}$. Aquesta assignació és, ella mateixa, functorial: dona un functor de $\cat{C}$ cap a la categoria de prefeixos $\widehat{\cat{C}}=\cat{Set}^{\cat{C}\op}$.

\begin{definicio}[Immersió de Yoneda]\index{Yoneda!immersió}
\label{def:immersio-yoneda}
La \textbf{immersió de Yoneda} (en anglès, \emph{Yoneda embedding}) és el functor
\[
\yon\colon\cat{C}\to\cat{Set}^{\cat{C}\op},
\qquad
\yon(A)=\Hom(-,A),
\]
que envia cada objecte $A$ al functor representable $\Hom(-,A)$, i cada morfisme $f\colon A\to B$ a la transformació natural $\yon(f)\colon\Hom(-,A)\Rightarrow\Hom(-,B)$ de components la postcomposició amb $f$,
\[
\yon(f)_X=\Hom(X,f)\colon\Hom(X,A)\to\Hom(X,B),
\qquad
\varphi\mapsto f\comp\varphi.
\]
D'acord amb la notació dels homfunctors, escriurem també $\Hom(-,f)$\index{Hom@$\Hom$!$\Hom(-,f)$} per a aquesta transformació natural $\yon(f)$.
\end{definicio}

\begin{observacio}
Cal comprovar que $\yon(f)$ és realment natural i que $\yon$ preserva composició i identitats, però tot es redueix, un cop més, a l'associativitat i a les lleis de la identitat de $\cat{C}$. El contingut de la immersió de Yoneda és que \emph{tota} categoria petita s'immergeix en la seva categoria de prefeixos, que és una categoria molt rica: com veurem al capítol següent, té tots els límits i colímits \emph{petits}, calculats puntualment. El teorema següent en justifica el nom: $\yon$ és una immersió plena i fidel.
\end{observacio}

\section{El lema de Yoneda}

El resultat central relaciona les transformacions naturals que surten d'un functor representable amb els elements del functor de destí.

\begin{teorema}[Lema de Yoneda]\index{Yoneda!lema}
\label{teo:yoneda}
Sigui $\cat{C}$ una categoria \textbf{petita}, $A$ un objecte de $\cat{C}$ i $P$ un prefeix sobre $\cat{C}$. Aleshores hi ha una bijecció
\[
\Nat\big(\Hom(-,A),\,P\big)\;\cong\;P(A),
\]
natural \emph{contravariantment} en $A\in\cat{C}$ i \emph{covariantment} en $P\in\widehat{\cat{C}}$. Explícitament, la bijecció envia una transformació natural $\eta$ a l'element $\eta_A(\id_A)\in P(A)$.\footnotemark
\end{teorema}
\footnotetext{Si $\cat{C}$ és només localment petita, la mateixa construcció classifica les transformacions naturals que surten del representable mitjançant els elements de $P(A)$, sense haver de formar la categoria $\widehat{\cat{C}}$ de tots els prefeixos. En aquest cas, $\Nat(\Hom(-,A),P)$ s'ha d'entendre com aquesta col·lecció parametritzada pels elements de $P(A)$: en el marc de classes de l'apèndix~\ref{cap:mida}, una transformació amb domini gran és una família indexada per una classe, i no esdevé per això element d'un conjunt. Per parlar-ne literalment com d'un homconjunt dins d'una categoria de prefeixos cal fixar una convenció de mida addicional, per exemple amb universos; vegeu \cite[§1.1]{riehl}.}%

\begin{prova}
Descrivim les dues aplicacions inverses.

\textbf{De transformacions a elements.} A cada $\eta\colon\Hom(-,A)\Rightarrow P$ li assignem
\[
\Phi(\eta)=\eta_A(\id_A)\in P(A).
\]
Té sentit perquè $\id_A\in\Hom(A,A)$ i $\eta_A\colon\Hom(A,A)\to P(A)$.

\textbf{D'elements a transformacions.} A cada $x\in P(A)$ li assignem la transformació $\Psi(x)$ de components
\[
\Psi(x)_X\colon\Hom(X,A)\to P(X),
\qquad
\varphi\mapsto P(\varphi)(x).
\]
Té sentit perquè $\varphi\colon X\to A$ dona, per contravariància, $P(\varphi)\colon P(A)\to P(X)$, i $x\in P(A)$. Cal veure que $\Psi(x)$ és natural: per a un morfisme $g\colon Y\to X$, el quadrat de naturalitat és $P(g)\comp\Psi(x)_X=\Psi(x)_Y\comp g^*$, que sobre $\varphi\in\Hom(X,A)$ dona
\[
P(g)\big(P(\varphi)(x)\big)=P(\varphi\comp g)(x)=\Psi(x)_Y(\varphi\comp g)=\Psi(x)_Y\big(g^*(\varphi)\big),
\]
on hem usat la functorialitat $P(g)\comp P(\varphi)=P(\varphi\comp g)$.

\textbf{Són inverses.} D'una banda,
\[
\Phi(\Psi(x))=\Psi(x)_A(\id_A)=P(\id_A)(x)=x.
\]
De l'altra, donada $\eta$, cal veure que $\Psi(\eta_A(\id_A))=\eta$; és a dir, que per a tot $X$ i tot $\varphi\colon X\to A$,
\[
P(\varphi)\big(\eta_A(\id_A)\big)=\eta_X(\varphi).
\]
Això és exactament el quadrat de naturalitat de $\eta$ per al morfisme $\varphi\colon X\to A$, aplicat a $\id_A\in\Hom(A,A)$: en efecte, $\varphi^*(\id_A)=\id_A\comp\varphi=\varphi$, de manera que
\[
\eta_X(\varphi)=\eta_X(\varphi^*(\id_A))=P(\varphi)(\eta_A(\id_A)).
\]
Les dues assignacions són, doncs, inverses. Queda comprovar que la bijecció és \emph{natural} en les dues variables: és part de l'enunciat, no un afegit, i les dues verificacions són breus. Escrivim $\Phi_{A,P}\colon\Nat(\Hom(-,A),P)\to P(A)$ per exhibir-ne la dependència.

\emph{Naturalitat en $P$} (covariant). Sigui $\sigma\colon P\Rightarrow Q$ una transformació natural de prefeixos, amb postcomposició $(\sigma\comp-)\colon\Nat(\Hom(-,A),P)\to\Nat(\Hom(-,A),Q)$. Per a $\eta\colon\Hom(-,A)\Rightarrow P$,
\[
\Phi_{A,Q}(\sigma\comp\eta)=(\sigma\comp\eta)_A(\id_A)=\sigma_A\big(\eta_A(\id_A)\big)=\sigma_A\big(\Phi_{A,P}(\eta)\big),
\]
és a dir $\Phi_{A,Q}\comp(\sigma\comp-)=\sigma_A\comp\Phi_{A,P}$.

\emph{Naturalitat en $A$} (contravariant). Sigui $a\colon A'\to A$ un morfisme de $\cat{C}$. La immersió dona $\yon(a)=\Hom(-,a)\colon\Hom(-,A')\Rightarrow\Hom(-,A)$, i la precomposició indueix $(-\comp\yon(a))\colon\Nat(\Hom(-,A),P)\to\Nat(\Hom(-,A'),P)$. Per a $\eta\colon\Hom(-,A)\Rightarrow P$, com que $\yon(a)_{A'}(\id_{A'})=a\comp\id_{A'}=a$,
\[
\Phi_{A',P}\big(\eta\comp\yon(a)\big)=\big(\eta\comp\yon(a)\big)_{A'}(\id_{A'})=\eta_{A'}(a)=P(a)\big(\eta_A(\id_A)\big)=P(a)\big(\Phi_{A,P}(\eta)\big),
\]
on la penúltima igualtat és la identitat de naturalitat de $\eta$ demostrada més amunt (la que dona $\Psi(\eta_A(\id_A))=\eta$), aplicada amb $X=A'$ i $\varphi=a$; és a dir $\Phi_{A',P}\comp(-\comp\yon(a))=P(a)\comp\Phi_{A,P}$. Aquestes dues identitats completen el lema tal com s'ha enunciat.
\end{prova}

\begin{observacio}[Taula de control: variància i tipus]\label{obs:taula-yoneda}\index{Yoneda!taula de variància i tipus}
Els errors més freqüents d'aquest capítol són errors de tipus: una fletxa que va en el sentit equivocat o una composició que no existeix. Abans de qualsevol càlcul convé escriure el domini i el codomini de cada component, com al capítol~\ref{cap:functors} (observació~\ref{obs:comprovar-tipus}). La taula recull els casos que hem fet servir, amb $f\colon A\to B$, $u\colon X'\to X$, $v\colon X\to Y$, $a\colon A'\to A$ i $\sigma\colon P\Rightarrow Q$.
\begin{center}
\footnotesize
\renewcommand{\arraystretch}{1.35}
\setlength{\tabcolsep}{4pt}
\begin{tabularx}{\linewidth}{@{}>{\raggedright\arraybackslash}p{0.22\linewidth}>{\raggedright\arraybackslash}p{0.24\linewidth}>{\raggedright\arraybackslash}X@{}}
\toprule
\textbf{Expressió} & \textbf{Què és i variància} & \textbf{Tipus i acció}\\
\midrule
$\Hom(-,A)$ & prefeix $\cat{C}\op\to\cat{Set}$ (contravariant) & $X\mapsto\Hom(X,A)$; \ $\Hom(u,A)=u^*\colon\Hom(X,A)\to\Hom(X',A)$, $\psi\mapsto\psi\comp u$\\
$\Hom(A,-)$ & functor $\cat{C}\to\cat{Set}$ (covariant) & $X\mapsto\Hom(A,X)$; \ $\Hom(A,v)=v_*\colon\Hom(A,X)\to\Hom(A,Y)$, $\varphi\mapsto v\comp\varphi$\\
$\Hom(-,f)=\yon(f)$ & transformació natural $\Hom(-,A)\Rightarrow\Hom(-,B)$, en el sentit de $f$ (no és un functor) & component en $X$: $\Hom(X,f)\colon\Hom(X,A)\to\Hom(X,B)$, $\varphi\mapsto f\comp\varphi$\\
$\Hom(f,-)$ & transformació natural $\Hom(B,-)\Rightarrow\Hom(A,-)$, en sentit \emph{contrari} a $f$ & component en $X$: $\Hom(f,X)\colon\Hom(B,X)\to\Hom(A,X)$, $\psi\mapsto\psi\comp f$\\
$\yon$ & functor $\cat{C}\to\cat{Set}^{\cat{C}\op}$ (covariant) & $A\mapsto\Hom(-,A)$; \ $f\mapsto\Hom(-,f)\colon\yon(A)\Rightarrow\yon(B)$\\
\midrule
naturalitat de $\Hom(-,f)$ en $X$ & quadrat per a $u\colon X'\to X$ & $\Hom(X',f)\comp\Hom(u,A)=\Hom(u,B)\comp\Hom(X,f)\colon\Hom(X,A)\to\Hom(X',B)$; \ tots dos envien $\varphi$ a $f\comp\varphi\comp u$\\
lema de Yoneda, naturalitat en $P$ & covariant, per a $\sigma\colon P\Rightarrow Q$ & $\Phi_{A,Q}\comp(\sigma\comp-)=\sigma_A\comp\Phi_{A,P}\colon\Nat(\yon(A),P)\to Q(A)$\\
lema de Yoneda, naturalitat en $A$ & contravariant, per a $a\colon A'\to A$ & $\Phi_{A',P}\comp(-\comp\yon(a))=P(a)\comp\Phi_{A,P}\colon\Nat(\yon(A),P)\to P(A')$\\
\bottomrule
\end{tabularx}
\end{center}
La regla que hi ha darrere és sempre la mateixa. A $\Hom(X,Y)$, un morfisme que actua sobre la \emph{segona} variable ho fa per postcomposició i conserva el sentit; un que actua sobre la \emph{primera} ho fa per precomposició i l'inverteix. Per això $\Hom(-,f)$ va en el sentit de $f$ i $\Hom(f,-)$ en el contrari, i per això la naturalitat del lema és covariant en $P$ i contravariant en $A$.
\end{observacio}

\section{Conseqüències}

La primera conseqüència és que la immersió de Yoneda mereix el seu nom.

\begin{corollari}[La immersió de Yoneda és plena i fidel]\index{Yoneda!immersió plena i fidel}
Per a tota parella d'objectes $A,B$ de $\cat{C}$, la immersió $\yon$ dona una bijecció
\[
\Hom(A,B)\;\cong\;\Nat\big(\Hom(-,A),\,\Hom(-,B)\big).
\]
En particular, $\yon\colon\cat{C}\to\cat{Set}^{\cat{C}\op}$ és ple i fidel.
\end{corollari}

\begin{prova}
Apliquem el lema de Yoneda amb $P=\Hom(-,B)$. Aleshores
\[
\Nat\big(\Hom(-,A),\,\Hom(-,B)\big)\cong \Hom(-,B)(A)=\Hom(A,B),
\]
i seguint la bijecció es comprova que és precisament $f\mapsto\yon(f)$. Per tant $\yon$ és bijectiu sobre cada conjunt de morfismes, és a dir, ple i fidel.
\end{prova}

\begin{corollari}\index{Yoneda!objectes determinats per morfismes}
Dos objectes $A$ i $B$ de $\cat{C}$ són isomorfs si i només si els seus functors representables ho són:
\[
A\cong B
\quad\Longleftrightarrow\quad
\Hom(-,A)\cong\Hom(-,B).
\]
\end{corollari}

\begin{prova}
Un functor ple i fidel reflecteix els isomorfismes (capítol dels functors): si $\yon(A)\cong\yon(B)$, aleshores $A\cong B$. El recíproc és que tot functor preserva isomorfismes. Aplicat a $\yon$, dona l'equivalència.
\end{prova}

\begin{observacio}
Aquest darrer corol·lari és la formulació precisa de l'eslògan del primer capítol: \emph{un objecte queda determinat, llevat d'isomorfisme, per la xarxa de morfismes que hi arriben}. Conèixer $\Hom(-,A)$ ---és a dir, saber com és el conjunt de morfismes $X\to A$ per a cada $X$, de manera compatible amb la composició--- equival a conèixer $A$ llevat d'isomorfisme. En particular, l'objecte que representa un functor representable és únic llevat d'isomorfisme, cosa que unifica totes les unicitats de les propietats universals que hem trobat.
\end{observacio}

\begin{observacio}
Hi ha una versió covariant de tot plegat, obtinguda treballant amb $\cat{C}\op$ en lloc de $\cat{C}$: dona una immersió $\cat{C}\op\to\cat{Set}^{\cat{C}}$, també plena i fidel, i el lema $\Nat(\Hom(A,-),Q)\cong Q(A)$ per a functors covariants $Q\colon\cat{C}\to\cat{Set}$. Totes dues versions s'usen contínuament; la contravariant, amb prefeixos, és la que enllaça amb la teoria de feixos i la lògica categòrica que estudiarem més endavant.
\end{observacio}

\begin{resumcap}
\apartat{Idees centrals}
\begin{enumerate}
  \item Les dades universals es caracteritzen per l'existència i la unicitat de morfismes; sovint es poden veure com a objectes inicials o terminals en una categoria coma adequada.
  \item Una representació d'un functor és un objecte representant juntament amb un element universal; equivalentment, és un isomorfisme natural amb un homfunctor.
  \item El lema de Yoneda dona una bijecció natural
  \[
  \Nat\bigl(\Hom(-,A),P\bigr)\;\cong\;P(A),
  \]
  determinada per l'avaluació a $\id_A$.
  \item La immersió de Yoneda és plena i fidel; per tant, els morfismes entre objectes es recuperen exactament com a transformacions naturals entre els seus prefeixos representables.
  \item En particular, un objecte queda determinat, llevat d'isomorfisme, pel sistema de morfismes que hi arriben; la representació completa és única llevat d'un únic isomorfisme compatible amb la dada universal.
\end{enumerate}
\apartat{Recordatori operatiu} Les dues direccions de la bijecció de Yoneda són
\[
\eta\longmapsto\eta_A(\id_A),
\qquad\qquad
x\longmapsto\bigl(f\mapsto P(f)(x)\bigr).
\]
\apartat{Errors típics} Dir \enquote{únic llevat d'isomorfisme únic} de l'objecte nu quan això només val per a la representació amb la dada universal; usar la propietat universal del producte abans d'haver-la definit; i confondre el functor representat amb l'objecte representant.
\apartat{Lectura recomanada} \cite[§2.2 i §§8.1--8.4]{awodey} per als objectes inicials i terminals i per a la immersió, el lema i les aplicacions de Yoneda; \cite[§§4.1--4.3]{leinster}; \cite[cap.~2]{riehl}, d'arquitectura molt propera a la d'aquest capítol; \cite[II.6 i III.1--III.2]{maclane} per a les categories coma, les fletxes universals i el lema de Yoneda; i \cite[cap.~2]{perrone-notes}, per a una segona explicació del lema de Yoneda, amb molts casos particulars i exemples de matemàtica elemental.
\end{resumcap}

\chapter[Límits i colímits]{Límits i colímits}
\label{cap:limits}

\begin{objectius}
\apartat{Prerequisits} Categories oposades i dualitat, i categories sobre un objecte (cap.~\ref{cap:categories}); functors, functors constants i el functor diagonal (cap.~\ref{cap:functors}); transformacions naturals i categories de functors (cap.~\ref{cap:transformacions}); categories coma, propietats universals, representabilitat i lema de Yoneda, especialment la unicitat llevat d'isomorfisme (cap.~\ref{cap:yoneda}).
\apartat{Objectius} En acabar el capítol, el lector hauria de poder:
\begin{itemize}
  \item interpretar un diagrama com un functor i un con com una transformació natural;
  \item definir límit i colímit com a cons i cocons universals;
  \item reconèixer productes, equalitzadors, pullbacks i l'objecte terminal com a límits, i els seus duals com a colímits;
  \item calcular aquests casos bàsics a $\cat{Set}$ i interpretar-los en preordres;
  \item fer servir els pullbacks per descriure la reindexació;
  \item enunciar la caracterització dels límits finits mitjançant objecte terminal, productes binaris i equalitzadors, o bé objecte terminal i pullbacks;
  \item descriure els límits i colímits petits a $\cat{Set}$ i explicar com es calculen punt a punt els límits i colímits en categories de functors;
  \item distingir preservació i reflexió de límits;
  \item explicar per què els representables covariants preserven límits.
\end{itemize}
\end{objectius}

\section{Cap a la noció de límit}

El capítol anterior ha caracteritzat diversos objectes per propietats universals: l'objecte terminal, l'inicial, i de passada el producte i el coproducte. Tots responen a un mateix patró: un objecte que es relaciona de manera òptima amb tots els altres. I tots són únics llevat d'isomorfisme. En aquest capítol unifiquem aquest patró en la noció general de \textbf{límit} d'un diagrama, i la seva dual, el \textbf{colímit}.

La idea és senzilla. Un diagrama és una col·lecció d'objectes i morfismes amb una forma determinada. Un \emph{con} sobre el diagrama és un objecte que projecta cap a tots els seus vèrtexs de manera compatible. El \emph{límit} és el con òptim: aquell pel qual tots els altres factoritzen de manera única. Segons la forma del diagrama, aquesta recepta produeix productes, equalitzadors, pullbacks, objectes terminals i molt més, tots com a casos particulars d'una sola definició. La representabilitat i el lema de Yoneda en seran l'eina conceptual: en el llenguatge del capítol anterior, un límit és un objecte que representa el prefeix que envia cada objecte $X$ al conjunt dels cons de vèrtex $X$ sobre el diagrama.

\section{Diagrames i cons}

Formalitzem la idea de \enquote{diagrama d'una forma donada} amb el llenguatge dels functors, que ja tenim a mà.

\begin{definicio}[Diagrama]\label{def:diagrama}\index{diagrama!de forma $\cat{I}$}\index{categoria!índex}
Sigui $\cat{I}$ una categoria petita, anomenada \textbf{categoria índex}. Un \textbf{diagrama de forma $\cat{I}$} en una categoria $\cat{C}$ és un functor $D\colon\cat{I}\to\cat{C}$. Escrivim $D_i$ per a l'objecte $D(i)$ i, per a un morfisme $u\colon i\to j$ de $\cat{I}$, $D_u\colon D_i\to D_j$ per a la fletxa corresponent.
\end{definicio}

\begin{observacio}
La categoria índex $\cat{I}$ fixa la \emph{forma} del diagrama, no el seu contingut. Amb $\cat{I}$ discreta de dos objectes obtenim una parella d'objectes sense fletxes; amb $\cat{I}$ la categoria de dues fletxes paral·leles, un parell de morfismes amb el mateix domini i codomini; amb $\cat{I}$ buida, el diagrama buit. Cada forma donarà lloc a un tipus de límit.
\end{observacio}

\begin{definicio}[Con]\index{con}
Sigui $D\colon\cat{I}\to\cat{C}$ un diagrama. Un \textbf{con} sobre $D$ amb \textbf{vèrtex} $A$ és una família de morfismes
\[
\alpha_i\colon A\to D_i,\qquad i\in\cat{I},
\]
tal que, per a cada morfisme $u\colon i\to j$ de $\cat{I}$, el triangle commuta:
\[
\begin{tikzpicture}[baseline=(m.center)]
  \matrix (m) [matrix of math nodes, row sep=2.4em, column sep=3.0em] {
     & A & \\
    D_i & & D_j \\
  };
  \draw[-{Stealth[scale=1.1]}] (m-1-2) -- node[above left] {$\alpha_i$} (m-2-1);
  \draw[-{Stealth[scale=1.1]}] (m-1-2) -- node[above right] {$\alpha_j$} (m-2-3);
  \draw[-{Stealth[scale=1.1]}] (m-2-1) -- node[below] {$D_u$} (m-2-3);
\end{tikzpicture}
\qquad
D_u\comp\alpha_i=\alpha_j.
\]
\end{definicio}

\begin{observacio}\label{obs:con-transformacio}
Un con és exactament una transformació natural $\alpha\colon\Delta_A\Rightarrow D$, on $\Delta_A\colon\cat{I}\to\cat{C}$ és el functor constant amb valor $A$ (exemple~\ref{ex:functor-constant}). En efecte, els components $\alpha_i\colon A\to D_i$ i la condició de naturalitat $D_u\comp\alpha_i=\alpha_j$ són precisament la definició de con. Aquest és el primer lloc on la maquinària de les transformacions naturals rendeix directament.
\end{observacio}

\begin{definicio}[Morfisme de cons]
Un \textbf{morfisme de cons} de $(A,\alpha)$ a $(B,\beta)$, tots dos sobre $D$, és un morfisme $f\colon A\to B$ tal que $\beta_i\comp f=\alpha_i$ per a tot $i$. Els cons sobre $D$ i els seus morfismes formen una categoria, que denotem $\mathrm{Con}(D)$.
\end{definicio}

\begin{observacio}[Els cons són objectes d'una categoria coma]\label{obs:cons-coma}\index{con}\index{categoria!coma}\index{functor!diagonal}
Els functors constants $\Delta_A$ s'apleguen en un sol functor, el \textbf{functor diagonal}
\[
\Delta\colon\cat{C}\longrightarrow\cat{C}^{\cat{I}},
\]
que envia cada objecte $A$ al functor constant $\Delta_A$ i cada morfisme $f\colon A\to B$ a la transformació natural $\Delta_f\colon\Delta_A\Rightarrow\Delta_B$ de components $(\Delta_f)_i=f$; la naturalitat és trivial, perquè $\Delta_A$ i $\Delta_B$ envien tots els morfismes de $\cat{I}$ a identitats, i la functorialitat es comprova component a component. Amb aquest functor, la categoria de cons és una categoria coma (secció~\ref{sec:coma}):
\[
\mathrm{Con}(D)=(\Delta\downarrow D).
\]
En efecte, prenent $S=\Delta$ i $T=D\colon\cat{1}\to\cat{C}^{\cat{I}}$, un objecte de $(\Delta\downarrow D)$ és una terna $(A,\ast,\alpha)$ amb $A$ objecte de $\cat{C}$ i $\alpha\colon\Delta_A\Rightarrow D$ un morfisme de $\cat{C}^{\cat{I}}$, que identifiquem amb la parella $(A,\alpha)$, com a l'exemple~\ref{ex:slice-coma}; és a dir, un con sobre $D$ amb vèrtex $A$ (observació~\ref{obs:con-transformacio}). Un morfisme $(A,\alpha)\to(B,\beta)$ és una parella $(f,\id_\ast)$ amb $f\colon A\to B$ tal que el quadrat de la definició de categoria coma commuta, és a dir, $\beta\comp\Delta_f=\alpha$; component a component, $\beta_i\comp f=\alpha_i$ per a tot $i$, que és exactament la condició de morfisme de cons.

Dualment, un cocon sobre $D$ amb vèrtex $A$ (definició~\ref{def:cocon}) és una transformació natural $D\Rightarrow\Delta_A$, i els cocons formen la categoria coma $(D\downarrow\Delta)$. Així, el límit és un objecte terminal de $(\Delta\downarrow D)$ i el colímit un objecte inicial de $(D\downarrow\Delta)$: és l'eslògan de l'exemple~\ref{ex:universal-coma}, \emph{ser universal és ser inicial o terminal en una categoria coma adient}. La categoria coma dona, doncs, el llenguatge comú per a les categories sobre i sota un objecte, els objectes universals del capítol~\ref{cap:yoneda} i els límits i colímits.
\end{observacio}

\section{Límits}

\begin{definicio}[Límit]\index{límit}
\label{def:limit}
Un \textbf{límit} del diagrama $D\colon\cat{I}\to\cat{C}$ és un objecte terminal de la categoria de cons $\mathrm{Con}(D)$. Explícitament, és un con $(L,\lambda)$ tal que, per a tot con $(A,\alpha)$ sobre $D$, hi ha un \emph{únic} morfisme $f\colon A\to L$ amb
\[
\lambda_i\comp f=\alpha_i\qquad\text{per a tot }i.
\]
Escrivim $L=\varprojlim D$ o $L=\lim D$.
\end{definicio}

\begin{observacio}
Un límit és, doncs, el con \enquote{universal}: tots els altres cons factoritzen a través d'ell de manera única. Com que és un objecte terminal de $\mathrm{Con}(D)$, el con límit és \textbf{únic llevat d'un únic isomorfisme} de cons, és a dir, compatible amb les projeccions, per la unicitat dels objectes terminals (proposició~\ref{prop:unicitat-terminal}); el vèrtex sol és únic llevat d'isomorfisme. Direm que una categoria \textbf{té límits de forma $\cat{I}$} quan tot diagrama de forma $\cat{I}$ hi té límit.
\end{observacio}

La propietat universal del límit es pot formular, en el llenguatge del capítol~\ref{cap:yoneda}, com una representabilitat. Suposem $\cat{I}$ petita i $\cat{C}$ localment petita, de manera que, per a cada objecte $A$, els cons sobre $D$ amb vèrtex $A$ formen un conjunt
\[
\mathrm{Con}(A,D)=\Nat(\Delta_A,D).
\]
Un morfisme $g\colon A'\to A$ transforma cada con $\alpha$ amb vèrtex $A$ en el con $\alpha\comp\Delta_g$ amb vèrtex $A'$, de components $\alpha_i\comp g$; obtenim així un prefeix
\[
\mathrm{Con}(-,D)\colon\cat{C}\op\longrightarrow\cat{Set}.
\]

\begin{proposicio}[Els límits com a representacions]\label{prop:limit-representable}\index{límit!com a representació}
Un con $(L,\lambda)$ sobre $D$ és un límit si i només si les aplicacions
\[
\Hom(A,L)\longrightarrow\mathrm{Con}(A,D),
\qquad
f\longmapsto\lambda\comp\Delta_f=(\lambda_i\comp f)_i,
\]
formen un isomorfisme natural $\Hom(-,L)\cong\mathrm{Con}(-,D)$. Dit d'una altra manera, un límit de $D$ és una representació $(L,\lambda)$ del prefeix $\mathrm{Con}(-,D)$, amb element universal $\lambda\in\mathrm{Con}(L,D)$.
\end{proposicio}

\begin{prova}
Les aplicacions són naturals en $A$: per a $g\colon A'\to A$, enviar primer $f$ a $(\lambda_i\comp f)_i$ i precompondre després amb $g$ dona $(\lambda_i\comp f\comp g)_i$, que és la imatge de $f\comp g$. Una transformació natural és un isomorfisme natural exactament quan totes les seves components són bijectives, i que l'aplicació corresponent a $A$ sigui bijectiva diu exactament que, per a tot con $\alpha$ amb vèrtex $A$, hi ha un únic $f\colon A\to L$ amb $\lambda_i\comp f=\alpha_i$ per a tot $i$. Això és la definició de límit. L'element universal és la imatge de $\id_L$, que és $\lambda$.
\end{prova}

Aquesta formulació explica d'un sol cop tres fets del capítol. Un límit és una propietat universal en el sentit del capítol~\ref{cap:yoneda}; és, per tant, únic llevat d'un únic isomorfisme compatible amb la dada universal, que aquí és el con $\lambda$; i la mateixa idea mostrarà que els representables preserven límits (proposició~\ref{prop:representables-limits}).

Recorrem ara les formes particulars, cadascuna amb la seva categoria índex.

\subsection{Producte}\index{producte}

Prenem $\cat{I}$ la categoria \textbf{discreta} amb dos objectes $1,2$ (sense fletxes no trivials). Un diagrama és simplement una parella d'objectes $A,B$. Un con amb vèrtex $X$ és una parella de morfismes $f\colon X\to A$, $g\colon X\to B$, sense cap condició (no hi ha fletxes a $\cat{I}$ que imposin res). El límit és el \textbf{producte}.

\begin{definicio}\index{producte}
\label{def:producte}
El \textbf{producte} de dos objectes $A,B$ és un objecte $A\times B$ amb dues projeccions $\pi_1\colon A\times B\to A$ i $\pi_2\colon A\times B\to B$ tals que, per a tot objecte $X$ i tota parella de morfismes $f\colon X\to A$, $g\colon X\to B$, hi ha un únic morfisme $\langle f,g\rangle\colon X\to A\times B$ amb
\[
\pi_1\comp\langle f,g\rangle=f,\qquad \pi_2\comp\langle f,g\rangle=g.
\]
\[
\begin{tikzpicture}[baseline=(m.center)]
  \matrix (m) [matrix of math nodes, column sep=2.8em, row sep=2.6em] {
      & X & \\
    A & A\times B & B \\
  };
  \draw[-{Stealth[scale=1.1]}] (m-1-2) -- node[above left] {$f$} (m-2-1);
  \draw[-{Stealth[scale=1.1]}] (m-1-2) -- node[above right] {$g$} (m-2-3);
  \draw[-{Stealth[scale=1.1]},dashed] (m-1-2) -- node[right] {$\langle f,g\rangle$} (m-2-2);
  \draw[-{Stealth[scale=1.1]}] (m-2-2) -- node[below] {$\pi_1$} (m-2-1);
  \draw[-{Stealth[scale=1.1]}] (m-2-2) -- node[below] {$\pi_2$} (m-2-3);
\end{tikzpicture}
\]
\end{definicio}

\begin{exemple}
A $\cat{Set}$, el producte és el producte cartesià $A\times B=\{(a,b)\mid a\in A,\ b\in B\}$ amb les projeccions habituals. A $\cat{Grp}$ és el producte directe de grups; a $\cat{Top}$, l'espai producte amb la topologia producte; en un preordre, l'ínfim dels dos elements, quan existeix. En cada cas, la propietat universal recupera la construcció coneguda.
\end{exemple}

\subsection{Equalitzador}\index{equalitzador}

Prenem $\cat{I}$ la categoria amb dues fletxes paral·leles $\bullet\rightrightarrows\bullet$. Un diagrama és un parell de morfismes $f,g\colon A\to B$. Un con amb vèrtex $X$ és un morfisme $e\colon X\to A$ amb $f\comp e=g\comp e$ (la condició prové de les dues fletxes). El límit és l'\textbf{equalitzador}.

\begin{definicio}\index{equalitzador}
\label{def:equalitzador}
L'\textbf{equalitzador} de $f,g\colon A\to B$ és un objecte $E$ amb un morfisme $e\colon E\to A$ tal que $f\comp e=g\comp e$ i que és universal amb aquesta propietat: per a tot $h\colon X\to A$ amb $f\comp h=g\comp h$, hi ha un únic $k\colon X\to E$ amb $e\comp k=h$.
\[
\begin{tikzpicture}[baseline=(m.center)]
  \matrix (m) [matrix of math nodes, column sep=3em, row sep=2.6em] {
    E & A & B \\
    X & & \\
  };
  \draw[-{Stealth[scale=1.1]}] (m-1-1) -- node[above] {$e$} (m-1-2);
  \draw[-{Stealth[scale=1.1]}] ([yshift=2pt]m-1-2.east) -- node[above] {$f$} ([yshift=2pt]m-1-3.west);
  \draw[-{Stealth[scale=1.1]}] ([yshift=-2pt]m-1-2.east) -- node[below] {$g$} ([yshift=-2pt]m-1-3.west);
  \draw[-{Stealth[scale=1.1]}] (m-2-1) -- node[left] {$h$} (m-1-2);
  \draw[-{Stealth[scale=1.1]},dashed] (m-2-1) -- node[left] {$k$} (m-1-1);
\end{tikzpicture}
\]
\end{definicio}

\begin{exemple}
A $\cat{Set}$, l'equalitzador de $f,g\colon A\to B$ és el subconjunt $E=\{a\in A\mid f(a)=g(a)\}$ amb la injecció $E\hookrightarrow A$. En general, els equalitzadors representen els subobjectes definits per una equació. Tot equalitzador és un monomorfisme.
\end{exemple}

\subsection{Pullback}\index{pullback}

Prenem $\cat{I}$ la categoria amb forma de racó, $\bullet\to\bullet\leftarrow\bullet$. Un diagrama és un parell de morfismes amb codomini comú, $f\colon A\to C$ i $g\colon B\to C$. El límit és el \textbf{pullback}.

\begin{definicio}\index{pullback}
\label{def:pullback}
El \textbf{pullback} de $f\colon A\to C$ i $g\colon B\to C$ és un objecte $P$ amb morfismes $p_1\colon P\to A$ i $p_2\colon P\to B$ tals que $f\comp p_1=g\comp p_2$ i que és universal: per a tot $X$ amb $r_1\colon X\to A$, $r_2\colon X\to B$ i $f\comp r_1=g\comp r_2$, hi ha un únic $h\colon X\to P$ amb $p_1\comp h=r_1$ i $p_2\comp h=r_2$.
\[
\begin{tikzpicture}[baseline=(m.center)]
  \matrix (m) [matrix of math nodes, column sep=3em, row sep=2.6em] {
    P & B \\
    A & C \\
  };
  \draw[-{Stealth[scale=1.1]}] (m-1-1) -- node[above] {$p_2$} (m-1-2);
  \draw[-{Stealth[scale=1.1]}] (m-1-1) -- node[left] {$p_1$} (m-2-1);
  \draw[-{Stealth[scale=1.1]}] (m-1-2) -- node[right] {$g$} (m-2-2);
  \draw[-{Stealth[scale=1.1]}] (m-2-1) -- node[below] {$f$} (m-2-2);
  \node at ([xshift=0.75em,yshift=-0.75em]m-1-1) {$\scriptstyle\lrcorner$};
\end{tikzpicture}
\]
\end{definicio}

\begin{exemple}
A $\cat{Set}$, el pullback és $P=\{(a,b)\in A\times B\mid f(a)=g(b)\}$ amb les projeccions. Quan $g\colon B\hookrightarrow C$ és la inclusió d'un subconjunt $B\subseteq C$, el pullback recupera l'antiimatge $f^{-1}(B)$, mitjançant la bijecció $(a,b)\mapsto a$. En aquest cas, doncs, el pullback categoritza l'operació d'antiimatge. En general, el pullback també s'anomena \emph{producte fibrat} i s'escriu $A\times_C B$.
\end{exemple}

\begin{observacio}[Reindexació per pullback]\index{reindexació}\index{functor!de reindexació}
\label{obs:reindexacio}
Ara que disposem de pullbacks podem definir la reindexació entre categories sobre un objecte. Suposem que $\cat{C}$ té tots els pullbacks i sigui $f\colon A\to B$ un morfisme. Definim un functor
\[
f^{*}\colon\cat{C}/B\longrightarrow\cat{C}/A
\]
de la manera següent: a un objecte $y\colon Y\to B$ de $\cat{C}/B$ li assignem el pullback de $y$ al llarg de $f$, és a dir, l'objecte $f^{*}Y=A\times_B Y$ amb la seva projecció $p_1\colon A\times_B Y\to A$, que és un objecte de $\cat{C}/A$:
\[
\begin{tikzpicture}[baseline=(m.center)]
  \matrix (m) [matrix of math nodes, column sep=3em, row sep=2.6em] {
    A\times_B Y & Y \\
    A & B \\
  };
  \draw[-{Stealth[scale=1.1]}] (m-1-1) -- node[above] {$p_2$} (m-1-2);
  \draw[-{Stealth[scale=1.1]}] (m-1-1) -- node[left] {$p_1=f^{*}y$} (m-2-1);
  \draw[-{Stealth[scale=1.1]}] (m-1-2) -- node[right] {$y$} (m-2-2);
  \draw[-{Stealth[scale=1.1]}] (m-2-1) -- node[below] {$f$} (m-2-2);
  \node at ([xshift=0.9em,yshift=-0.8em]m-1-1) {$\scriptstyle\lrcorner$};
\end{tikzpicture}
\]
Sobre morfismes, la propietat universal del pullback assigna a cada triangle de $\cat{C}/B$ un únic triangle de $\cat{C}/A$. Un cop \textbf{triats} els pullbacks que defineixen $f^{*}$, obtenim un functor ordinari $f^{*}\colon\cat{C}/B\to\cat{C}/A$. Per a morfismes composables $A\xrightarrow{f}B\xrightarrow{g}C$, els functors $(g\comp f)^{*}$ i $f^{*}\comp g^{*}$ no són literalment iguals per a una tria arbitrària de pullbacks, però són canònicament isomorfs. Aquesta coherència s'organitza, en un llenguatge més avançat, mitjançant la noció de \emph{pseudofunctor}, que no necessitarem aquí. En passar als preordres de subobjectes (capítol~\ref{cap:subobjectes}), aquests isomorfismes queden identificats i la reindexació és estrictament functorial:
\[
(g\comp f)^{*}=f^{*}\comp g^{*},\qquad (\id_A)^{*}=\id_{\Sub(A)}.
\]
A $\cat{Set}$, si interpretem $y\colon Y\to B$ com la família de fibres $\{y^{-1}(b)\}_{b\in B}$, aleshores $f^{*}Y$ s'interpreta com la família reindexada $\{y^{-1}(f(a))\}_{a\in A}$: exactament la \textbf{substitució} al llarg de $f$. Aquesta és la lectura lògica que explotarem en tractar els subobjectes, on $f^{*}$ es restringeix a un morfisme de preordres $\Sub(B)\to\Sub(A)$ i s'interpreta com la substitució en un predicat.
\end{observacio}

\subsection{Objecte terminal com a límit}

Prenem $\cat{I}=\cat{0}$, la categoria buida. L'únic diagrama de forma $\cat{0}$ és el diagrama buit. Un con sobre el diagrama buit és, simplement, un objecte $X$ (sense cap component, perquè no hi ha vèrtexs). El límit ---el con terminal--- és, doncs, l'\textbf{objecte terminal} de $\cat{C}$. Així, l'objecte terminal és el límit del diagrama buit, i queda unificat amb els altres.

\section{Colímits}

Tota la teoria anterior es dualitza girant les fletxes. Un \textbf{colímit} d'un diagrama $D\colon\cat{I}\to\cat{C}$ és un límit del diagrama oposat $D\op\colon\cat{I}\op\to\cat{C}\op$; equivalentment, és un objecte inicial a la categoria de \textbf{cocons} sobre $D$.

\begin{definicio}[Cocon i colímit]\label{def:cocon}\index{cocon}\index{colímit}
Un \textbf{cocon} sobre $D$ amb vèrtex $A$ és una família $\alpha_i\colon D_i\to A$ tal que, per a cada $u\colon i\to j$, commuta el triangle dual
\[
\begin{tikzpicture}[baseline=(m.center)]
  \matrix (m) [matrix of math nodes, row sep=2.4em, column sep=3.0em] {
    D_i & & D_j \\
     & A & \\
  };
  \draw[-{Stealth[scale=1.1]}] (m-1-1) -- node[below left] {$\alpha_i$} (m-2-2);
  \draw[-{Stealth[scale=1.1]}] (m-1-3) -- node[below right] {$\alpha_j$} (m-2-2);
  \draw[-{Stealth[scale=1.1]}] (m-1-1) -- node[above] {$D_u$} (m-1-3);
\end{tikzpicture}
\qquad
\alpha_j\comp D_u=\alpha_i.
\]
Un \textbf{colímit} és un cocon $(L,\lambda)$ universal: per a tot cocon $(A,\alpha)$ hi ha un únic $f\colon L\to A$ amb $f\comp\lambda_i=\alpha_i$. Escrivim $L=\colimop D$ o $L=\varinjlim D$.
\end{definicio}

Dualitzant cada límit obtenim el colímit corresponent:

\begin{itemize}
\item El \textbf{coproducte} $A+B$ (dual del producte) té injeccions $\iota_1\colon A\to A+B$, $\iota_2\colon B\to A+B$ universals. A $\cat{Set}$ és la unió disjunta; a $\cat{Grp}$, el producte lliure; en un preordre, el suprem.
\item El \textbf{coequalitzador} de $f,g\colon A\to B$ (dual de l'equalitzador) és un $q\colon B\to Q$ universal amb $q\comp f=q\comp g$. A $\cat{Set}$ és el quocient de $B$ per la relació d'equivalència generada per $f(a)\sim g(a)$. Tot coequalitzador és un epimorfisme.
\item El \textbf{pushout} de $f\colon C\to A$ i $g\colon C\to B$ (dual del pullback) és un $Q$ universal amb el quadrat commutatiu. A $\cat{Set}$ és la unió de $A$ i $B$ enganxada al llarg de $C$, i s'escriu $A+_C B$.
\item L'\textbf{objecte inicial} (dual del terminal) és el colímit del diagrama buit.
\end{itemize}

\begin{observacio}
No cal demostrar res de nou per als colímits: cada enunciat sobre límits a $\cat{C}$ es tradueix en un enunciat sobre colímits a $\cat{C}\op$, i viceversa. Aquesta és la dualitat en acció, aplicada de manera sistemàtica. En endavant, quan enunciem un resultat per a límits, el dual per a colímits s'entén sobreentès.
\end{observacio}

\section{Límits finits}

\begin{definicio}\index{límit!finit}\index{categoria!finitament completa}
Un límit es diu \textbf{finit} quan la seva categoria índex $\cat{I}$ és finita, és a dir, quan té un nombre finit d'objectes \emph{i} un nombre finit de morfismes. (No n'hi ha prou que els objectes siguin finits: una categoria d'un sol objecte pot tenir infinits endomorfismes.) Una categoria és \textbf{finitament completa} (o \textbf{lex}) si té tots els límits finits; és \textbf{completa} si té tots els límits de forma petita.
\end{definicio}

El resultat clau redueix tots els límits finits a dos casos bàsics: productes finits i equalitzadors. Això és molt útil, perquè per comprovar que una categoria és finitament completa n'hi ha prou de verificar dues coses.

\begin{teorema}\index{límit!finit!construcció}\label{teo:limits-finits}
Per a una categoria $\cat{C}$, són equivalents:
\begin{enumerate}
\item $\cat{C}$ té tots els límits finits;
\item $\cat{C}$ té objecte terminal, productes binaris i equalitzadors;
\item $\cat{C}$ té objecte terminal i pullbacks.
\end{enumerate}
\end{teorema}

\begin{prova}
\emph{(1) implica (2) i (3).} L'objecte terminal, el producte binari, l'equalitzador i el pullback són límits de diagrames de formes finites: la categoria buida, la discreta de dos objectes, $\bullet\rightrightarrows\bullet$ i $\bullet\to\bullet\leftarrow\bullet$.

\emph{(2) implica (1).} Primer, $\cat{C}$ té productes de qualsevol família finita. El de la família buida és l'objecte terminal, i per inducció, si $P_n$ és un producte de $A_1,\dots,A_n$ amb projeccions $\pi_k$, aleshores $P_n\times A_{n+1}$, amb projeccions $\pi_k\comp p_1$ ($k\le n$) i $p_2$, és un producte de $A_1,\dots,A_{n+1}$. En efecte, donades $f_k\colon X\to A_k$, hi ha una única $f\colon X\to P_n$ amb $\pi_k\comp f=f_k$ per a $k\le n$, i una única $X\to P_n\times A_{n+1}$ amb components $f$ i $f_{n+1}$; tota fletxa amb les mateixes composicions amb les $n+1$ projeccions té necessàriament aquestes dues components.

Sigui ara $D\colon\cat{I}\to\cat{C}$ un diagrama finit, i escrivim $\cat{I}_0$ i $\cat{I}_1$ per als conjunts d'objectes i de morfismes de $\cat{I}$. Considerem els productes finits
\[
P=\prod_{i\in\cat{I}_0}D_i
\qquad\text{i}\qquad
Q=\prod_{u\in\cat{I}_1}D_{\mathrm{cod}(u)},
\]
amb projeccions $\pi_i\colon P\to D_i$ i $q_u\colon Q\to D_{\mathrm{cod}(u)}$, i les dues fletxes $s,t\colon P\to Q$ determinades, per a cada $u\colon i\to j$, per
\[
q_u\comp s=\pi_j,
\qquad\qquad
q_u\comp t=D_u\comp\pi_i.
\]
Sigui $e\colon L\to P$ un equalitzador de $s$ i $t$, i posem $\lambda_i=\pi_i\comp e\colon L\to D_i$. Vegem que $(L,\lambda)$ és un límit de $D$.

\emph{És un con.} Per a $u\colon i\to j$,
\[
D_u\comp\lambda_i=D_u\comp\pi_i\comp e=q_u\comp t\comp e=q_u\comp s\comp e=\pi_j\comp e=\lambda_j.
\]

\emph{És universal.} Sigui $(A,\alpha)$ un con sobre $D$, i sigui $a\colon A\to P$ l'única fletxa amb $\pi_i\comp a=\alpha_i$ per a tot $i$. Per a cada $u\colon i\to j$,
\[
q_u\comp s\comp a=\pi_j\comp a=\alpha_j,
\qquad
q_u\comp t\comp a=D_u\comp\pi_i\comp a=D_u\comp\alpha_i=\alpha_j,
\]
de manera que $s\comp a$ i $t\comp a$ tenen les mateixes components i, per la unicitat del producte $Q$, $s\comp a=t\comp a$. Per la propietat de l'equalitzador hi ha una única $h\colon A\to L$ amb $e\comp h=a$, i aleshores $\lambda_i\comp h=\pi_i\comp e\comp h=\alpha_i$. Si $h'\colon A\to L$ també compleix $\lambda_i\comp h'=\alpha_i$ per a tot $i$, aleshores $\pi_i\comp(e\comp h')=\alpha_i$, d'on $e\comp h'=a$ per la unicitat del producte $P$, i $h'=h$ per la de l'equalitzador. El diagrama buit hi queda inclòs: aleshores $P$ i $Q$ són productes buits, és a dir, objectes terminals, i $L$ és terminal.

\emph{(3) implica (2).} Sigui $1$ l'objecte terminal i, per a cada objecte $A$, $!_A\colon A\to 1$ l'única fletxa. \emph{Productes binaris}: un pullback de $!_A$ i $!_B$ és un producte de $A$ i $B$. En efecte, per a qualsevol parell $f\colon X\to A$, $g\colon X\to B$, la condició de con $!_A\comp f=\,!_B\comp g$ es compleix automàticament, perquè totes dues fletxes van a $1$; per tant, la propietat universal del pullback diu exactament la del producte.

\emph{Equalitzadors}: siguin $f,g\colon A\to B$. Ja sabem que hi ha productes binaris; formem el pullback de
\[
\langle\id_A,f\rangle,\ \langle\id_A,g\rangle\colon A\to A\times B,
\]
amb projeccions $p_1,p_2\colon P\to A$, de manera que $\langle\id_A,f\rangle\comp p_1=\langle\id_A,g\rangle\comp p_2$. Component amb la primera projecció del producte, $p_1=p_2$; anomenem $e$ aquesta fletxa. Component amb la segona, $f\comp e=g\comp e$. Si $h\colon X\to A$ compleix $f\comp h=g\comp h$, aleshores
\[
\langle\id_A,f\rangle\comp h=\langle h,f\comp h\rangle=\langle h,g\comp h\rangle=\langle\id_A,g\rangle\comp h,
\]
de manera que $(h,h)$ és un con sobre el pullback, i hi ha una única $k\colon X\to P$ amb $p_1\comp k=h=p_2\comp k$, és a dir, amb $e\comp k=h$. Si $k'$ també compleix $e\comp k'=h$, aleshores $p_1\comp k'=p_2\comp k'=h$ i $k'=k$. Per tant $e$ és un equalitzador de $f$ i $g$.

Amb les implicacions (3)$\Rightarrow$(2)$\Rightarrow$(1)$\Rightarrow$(3), les tres condicions són equivalents.
\end{prova}

\begin{observacio}
Dualment, una categoria és \textbf{finitament cocompleta} si té objecte inicial, coproductes binaris i coequalitzadors, o equivalentment objecte inicial i pushouts. Per a la lògica categòrica, els límits finits ---i molt especialment els \emph{pullbacks}--- són centrals: serveixen per interpretar la substitució i les operacions sobre subobjectes que veurem més endavant.
\end{observacio}

\begin{exemple}[Límits en un preordre i en la deduïbilitat]\label{ex:limits-preordre}\index{límit!en un preordre}
Sigui $P$ un preordre vist com a categoria i $D\colon\cat{I}\to P$ un diagrama. Com que entre dos objectes hi ha com a màxim una fletxa, un con amb vèrtex $x$ diu simplement que $x\le D_i$ per a tot $i$, i totes les condicions de commutativitat es compleixen automàticament. Per tant, quan existeixen,
\[
\varprojlim D=\bigwedge_{i\in\cat{I}_0}D_i,
\qquad\qquad
\varinjlim D=\bigvee_{i\in\cat{I}_0}D_i:
\]
el límit és l'ínfim dels objectes del diagrama, i el colímit, el suprem; les fletxes del diagrama no hi intervenen. En particular, l'equalitzador de dues fletxes paral·leles $x\rightrightarrows y$ és $x$ mateix (les dues fletxes coincideixen), i el pullback de $a\to c\leftarrow b$ és l'ínfim $a\wedge b$, que no depèn de $c$. Així, un preordre té tots els límits finits si i només si té element màxim (l'ínfim de la família buida, que és l'objecte terminal) i ínfims binaris.

A la categoria prima $\cat{Prov}_T$ d'un sistema deductiu $T$, això diu que, quan el sistema disposa de $\top$ i de conjunció, $\top$ és un objecte terminal ($p\vdash_T\top$ per a tota fórmula $p$) i $p\land q$ és un producte de $p$ i $q$ (exemple~\ref{ex:rep-producte}). Aquestes dues operacions donen, doncs, tots els límits finits en el context deductiu: la part \enquote{conjuntiva} de la lògica és exactament la dels límits finits.
\end{exemple}

\section{Límits a \texorpdfstring{$\cat{Set}$}{Set} i en categories de functors}

\begin{exemple}[Límits a $\cat{Set}$]\label{ex:limits-set}\index{límit!a Set@a $\cat{Set}$}
Sigui $D\colon\cat{I}\to\cat{Set}$ un diagrama amb $\cat{I}$ petita. El conjunt de les famílies compatibles
\[
L=\Bigl\{(x_i)_{i}\in\textstyle\prod_{i\in\cat{I}_0}D_i\ \Bigm|\ D_u(x_i)=x_j\ \text{per a tot } u\colon i\to j\Bigr\},
\]
amb les projeccions $\lambda_i\colon L\to D_i$, és un límit de $D$. En efecte, un con $(A,\alpha)$ envia cada $a\in A$ a la família $(\alpha_i(a))_i$, que és compatible per la condició de con; l'aplicació $a\mapsto(\alpha_i(a))_i$ és l'única $f\colon A\to L$ amb $\lambda_i\comp f=\alpha_i$. En particular, $\cat{Set}$ té tots els límits petits. Dualment, el colímit és el quocient de la unió disjunta $\coprod_i D_i$, els elements de la qual són parelles $(i,x)$ amb $x\in D_i$, per la relació d'equivalència generada per
\[
(i,x)\sim\bigl(j,D_u(x)\bigr)\qquad(u\colon i\to j);
\]
per tant, $\cat{Set}$ té també tots els colímits petits. A $\cat{Set}$, doncs, els límits imposen compatibilitats i els colímits enganxen dades identificant-les.
\end{exemple}

Els límits en categories de functors es calculen a partir dels límits de la categoria d'arribada, objecte per objecte.

\begin{proposicio}[Límits punt a punt]\label{prop:limits-punt-a-punt}\index{límit!en categories de functors}\index{categoria!de functors!límits}
Siguin $\cat{A}$ i $\cat{I}$ categories petites. Si $\cat{C}$ té límits de forma $\cat{I}$, aleshores la categoria de functors $\cat{C}^{\cat{A}}$ també en té, i es calculen \emph{punt a punt}: per a un diagrama $F\colon\cat{I}\to\cat{C}^{\cat{A}}$,
\[
\Bigl(\varprojlim F\Bigr)(a)\;\cong\;\varprojlim_{i\in\cat{I}}F(i)(a)
\qquad(a\in\cat{A}).
\]
Dualment, si $\cat{C}$ té colímits de forma $\cat{I}$, els colímits a $\cat{C}^{\cat{A}}$ es calculen punt a punt. Equivalentment, per a cada objecte $a$ de $\cat{A}$, el functor d'avaluació
\[
\mathrm{ev}_a\colon\cat{C}^{\cat{A}}\longrightarrow\cat{C},\qquad F\longmapsto F(a),
\]
preserva aquests límits i colímits.
\end{proposicio}
Vegeu també \cite[§3.3]{riehl} i, en el context de les categories de diagrames, \cite[§§8.5--8.6]{awodey}.

\begin{prova}
Recordem que un objecte de $\cat{C}^{\cat{A}}$ és un functor $\cat{A}\to\cat{C}$, i que un morfisme és una transformació natural, amb la composició calculada component a component. Així, el diagrama $F$ assigna a cada objecte $i$ de $\cat{I}$ un functor $F(i)\colon\cat{A}\to\cat{C}$, i a cada $u\colon i\to j$ una transformació natural $F(u)\colon F(i)\Rightarrow F(j)$, de components $F(u)_a$. Per a cada objecte $a$ de $\cat{A}$, el diagrama avaluat
\[
F_a=\mathrm{ev}_a\comp F\colon\cat{I}\to\cat{C},\qquad i\mapsto F(i)(a),\quad u\mapsto F(u)_a,
\]
és un functor, perquè la composició i les identitats de $\cat{C}^{\cat{A}}$ es calculen component a component.

\emph{Pas 1: el límit sobre objectes.} Per a cada objecte $a$ de $\cat{A}$, triem un límit $\bigl(L(a),(\lambda_i(a))_i\bigr)$ de $F_a$; per tant, $F(u)_a\comp\lambda_i(a)=\lambda_j(a)$ per a cada $u\colon i\to j$.

\emph{Pas 2: el límit sobre morfismes.} Sigui $v\colon a\to a'$ un morfisme de $\cat{A}$. La família $\bigl(F(i)(v)\comp\lambda_i(a)\bigr)_i$, de fletxes $L(a)\to F(i)(a')$, és un con sobre $F_{a'}$: per a $u\colon i\to j$, la naturalitat de $F(u)$ i la condició de con de $\lambda(a)$ donen
\[
F(u)_{a'}\comp F(i)(v)\comp\lambda_i(a)=F(j)(v)\comp F(u)_a\comp\lambda_i(a)=F(j)(v)\comp\lambda_j(a).
\]
Per la universalitat de $L(a')$, hi ha una única $L(v)\colon L(a)\to L(a')$ tal que
\[
\lambda_i(a')\comp L(v)=F(i)(v)\comp\lambda_i(a)\qquad\text{per a tot } i.
\tag{$\dagger$}
\]

\emph{Pas 3: $L$ és un functor.} Fem servir que dues fletxes cap a un límit que tenen les mateixes composicions amb totes les projeccions són iguals, per la unicitat de la factorització. Per a $v=\id_a$, tant $L(\id_a)$ com $\id_{L(a)}$ compleixen $(\dagger)$, perquè $F(i)(\id_a)=\id$; per tant són iguals. Per a $v\colon a\to a'$ i $w\colon a'\to a''$, aplicant $(\dagger)$ dues vegades i la functorialitat de $F(i)$,
\[
\begin{aligned}
\lambda_i(a'')\comp L(w)\comp L(v)&=F(i)(w)\comp\lambda_i(a')\comp L(v)\\
&=F(i)(w)\comp F(i)(v)\comp\lambda_i(a)=F(i)(w\comp v)\comp\lambda_i(a),
\end{aligned}
\]
de manera que $L(w)\comp L(v)$ compleix la condició $(\dagger)$ que caracteritza $L(w\comp v)$, i $L(w)\comp L(v)=L(w\comp v)$.

\emph{Pas 4: un con a $\cat{C}^{\cat{A}}$.} La condició $(\dagger)$ és exactament el quadrat de naturalitat de la família $\lambda_i=(\lambda_i(a))_a$; per tant, cada $\lambda_i$ és una transformació natural $L\Rightarrow F(i)$. A més, $F(u)\comp\lambda_i=\lambda_j$ a $\cat{C}^{\cat{A}}$, perquè aquesta igualtat de transformacions naturals es comprova component a component i ho diu el pas~1. Així, $(L,\lambda)$ és un con sobre $F$.

\emph{Pas 5: universalitat.} Sigui $(G,\alpha)$ un con sobre $F$, amb $\alpha_i\colon G\Rightarrow F(i)$ i $F(u)\comp\alpha_i=\alpha_j$. Avaluant en $a$, la família $\bigl((\alpha_i)_a\bigr)_i$ és un con sobre $F_a$ amb vèrtex $G(a)$, de manera que hi ha una única $\theta_a\colon G(a)\to L(a)$ amb $\lambda_i(a)\comp\theta_a=(\alpha_i)_a$ per a tot $i$. La família $\theta$ és natural: per a $v\colon a\to a'$, per $(\dagger)$ i per la naturalitat de $\alpha_i$,
\[
\lambda_i(a')\comp L(v)\comp\theta_a=F(i)(v)\comp(\alpha_i)_a=(\alpha_i)_{a'}\comp G(v)=\lambda_i(a')\comp\theta_{a'}\comp G(v)
\]
per a tot $i$, i per tant $L(v)\comp\theta_a=\theta_{a'}\comp G(v)$. Així $\theta\colon G\Rightarrow L$ és un morfisme de $\cat{C}^{\cat{A}}$ amb $\lambda_i\comp\theta=\alpha_i$. És l'únic: si $\theta'$ també ho compleix, a cada $a$ es té $\lambda_i(a)\comp\theta'_a=(\alpha_i)_a$, i la unicitat de la factorització a través de $L(a)$ dona $\theta'_a=\theta_a$. Per tant $(L,\lambda)$ és un límit de $F$.

\emph{Pas 6: el càlcul punt a punt.} Per construcció, $\mathrm{ev}_a(L,\lambda)=\bigl(L(a),\lambda(a)\bigr)$ és un límit de $F_a$. Si $(L',\lambda')$ és un altre límit de $F$, hi ha un isomorfisme $\varphi\colon L'\Rightarrow L$ amb $\lambda_i\comp\varphi=\lambda'_i$, i la seva component $\varphi_a$ és un isomorfisme amb $\lambda_i(a)\comp\varphi_a=\lambda'_i(a)$. Un con isomorf a un con límit, per un isomorfisme compatible amb les projeccions, també és un límit, perquè les factoritzacions a través de l'un i de l'altre es corresponen component amb $\varphi_a$ o amb $\varphi_a^{-1}$. Per tant $\bigl(L'(a),\lambda'(a)\bigr)$ és un límit de $F_a$: $\mathrm{ev}_a$ preserva els límits de forma $\cat{I}$, i $\bigl(\varprojlim F\bigr)(a)\cong\varprojlim_i F(i)(a)$.

L'argument per als colímits és el mateix amb totes les fletxes invertides: cocons en lloc de cons, i la propietat universal del colímit en lloc de la del límit.
\end{prova}

\begin{exemple}[Un mateix producte a $\cat{Set}$, a $\cat{Graph}$ i a $\cat{Set}^{\cat{A}}$]\label{ex:producte-tres-contextos}\index{producte!a Set, Graph i categories de functors}\index{objecte terminal!a Graph}
Fixem un sol diagrama abstracte: la parella d'objectes $(X,T)$, sobre la categoria discreta de dos objectes, juntament amb el diagrama buit, que dona l'objecte terminal. Calculem-ne el límit en tres llocs.

\emph{A $\cat{Set}$.} El producte és $X\times T$, i l'objecte terminal, un conjunt d'un sol element $\{\ast\}$. Amb $T=\{\ast\}$, l'aplicació $(x,\ast)\mapsto x$ és una bijecció $X\times\{\ast\}\cong X$. A més, les fletxes $\{\ast\}\to X$ són exactament els elements de $X$.

\emph{A $\cat{Graph}$.} Recordem que $\cat{Graph}\cong\cat{Set}^{\cat{P}}$, on $\cat{P}$ té dos objectes $E,V$ i dues fletxes $s,t\colon E\to V$. Per la proposició anterior, el producte es calcula component a component,
\[
V_{G\times H}=V_G\times V_H,
\qquad
E_{G\times H}=E_G\times E_H,
\]
amb origen i destí també component a component: l'aresta $(e,e')$ va del vèrtex $(s_G(e),s_H(e'))$ al vèrtex $(t_G(e),t_H(e'))$. Les projeccions són els morfismes de grafs que retenen una de les dues components. Prenem ara com a $T$ el candidat ingenu a \enquote{punt}, el graf $\mathbf V$ d'un sol vèrtex i cap aresta. La fórmula dona
\[
V_{G\times\mathbf V}=V_G\times\{\ast\}\cong V_G,
\qquad
E_{G\times\mathbf V}=E_G\times\varnothing=\varnothing:
\]
el producte amb $\mathbf V$ \emph{esborra totes les arestes}. Per tant, $\mathbf V$ no és l'objecte terminal de $\cat{Graph}$; de fet, d'un graf amb alguna aresta no hi ha cap morfisme cap a $\mathbf V$, perquè l'aresta no té on anar. L'objecte terminal es calcula també punt a punt, com el límit del diagrama buit a cada objecte de $\cat{P}$: un vèrtex \emph{i} una aresta, és a dir, el graf $\mathbf L$ d'un sol \emph{bucle}. Amb $T=\mathbf L$, sí que $G\times\mathbf L\cong G$. I els morfismes $\mathbf L\to G$ ja no són els vèrtexs de $G$, sinó els seus bucles: un graf sense bucles no té cap \enquote{punt} en sentit categòric, encara que tingui molts vèrtexs (compareu-ho amb l'observació~\ref{obs:raonar-elements}).

\emph{A una categoria de functors $\cat{Set}^{\cat{A}}$.} La recepta és la mateixa per a qualsevol categoria petita $\cat{A}$: $(F\times G)(a)=F(a)\times G(a)$ i $(F\times G)(v)=F(v)\times G(v)$, i l'objecte terminal és el functor constant de valor $\{\ast\}$. Els dos casos anteriors en són casos particulars: $\cat{Set}\cong\cat{Set}^{\cat{1}}$, i $\cat{Graph}\cong\cat{Set}^{\cat{P}}$. Ara s'entén la sorpresa del bucle: el terminal de $\cat{Set}^{\cat{P}}$ ha de ser un punt \emph{a cada objecte} de $\cat{P}$, també a $E$, i això obliga a tenir una aresta.

La forma del diagrama i la recepta del càlcul no canvien d'un context a l'altre. El que canvia és quins objectes fan cada paper, i la intuició de $\cat{Set}$ només es pot traslladar si es fa objecte per objecte de la categoria índex.
\end{exemple}

Com que $\cat{Set}$ té tots els límits i colímits petits (exemple~\ref{ex:limits-set}), la proposició mostra que també els té la categoria de prefeixos $\cat{Set}^{\cat{C}\op}$ de qualsevol categoria petita $\cat{C}$, tal com s'anunciava al capítol~\ref{cap:yoneda}. Aquest fet serà essencial en els capítols de subobjectes i de topos.

\section{Preservació i reflexió de límits}

Quan apliquem un functor a un diagrama i al seu límit, la imatge és un con sobre el diagrama imatge, però no té per què ser-ne el límit. Els functors que respecten límits reben un nom.

\begin{definicio}\index{functor!preserva límits}\index{functor!reflecteix límits}
Sigui $F\colon\cat{C}\to\cat{D}$ un functor.
\begin{itemize}
\item $F$ \textbf{preserva} el límit $(L,\lambda)$ d'un diagrama $D\colon\cat{I}\to\cat{C}$ si $(F(L),F\lambda)$ és un límit de $F\comp D$ a $\cat{D}$.
\item $F$ \textbf{reflecteix} límits si, per a tot diagrama $D\colon\cat{I}\to\cat{C}$ i tot con $(L,\lambda)$ sobre $D$, quan $(F(L),F\lambda)$ és un límit de $F\comp D$, aleshores $(L,\lambda)$ ja era un límit de $D$.
\end{itemize}
Es diu que $F$ \textbf{preserva els límits finits} si preserva el límit de tot diagrama finit que en tingui, i anàlogament per a altres classes.
\end{definicio}

\begin{observacio}
La reflexió es formula sempre \emph{respecte d'un con i un diagrama concrets}: cal partir d'un con $(L,\lambda)$ sobre $D$ la imatge del qual és limitant. Un functor ple i fidel permet aleshores aixecar i detectar les factoritzacions úniques ---perquè estableix una bijecció sobre els homconjunts---, però això no vol dir que \enquote{creï} límits del no-res: si a $\cat{C}$ no hi ha cap con sobre $D$ que projecti sobre el límit imatge, no hi ha res a reflectir. La distinció entre \emph{reflectir} (comparar un con donat) i \emph{crear} (produir-ne un de nou) és important i convé tenir-la present sempre que un functor interactuï amb els límits. La noció formal de crear límits, útil sobretot per als functors oblit de categories algebraiques, no serà necessària aquí; n'hi ha prou de retenir que és més forta que la mera reflexió d'un con ja donat.
\end{observacio}

\begin{exemple}
El functor oblit $U\colon\cat{Grp}\to\cat{Set}$ preserva tots els límits petits. En efecte, si $D\colon\cat{I}\to\cat{Grp}$ és un diagrama petit, el límit dels conjunts subjacents és el conjunt de famílies compatibles $(x_i)_i$ (exemple~\ref{ex:limits-set}). Les operacions de grup, definides component a component, conserven aquestes compatibilitats, perquè cada $D_u$ és un homomorfisme; aquest conjunt adquireix així una estructura de grup per a la qual les projeccions són homomorfismes i que satisfà la propietat universal del límit a $\cat{Grp}$. En particular, això recupera els productes i els equalitzadors habituals. En canvi, $U$ \emph{no} preserva tots els colímits: ja falla per als coproductes, perquè el conjunt subjacent d'un coproducte de grups (el producte lliure) no és, en general, la unió disjunta dels conjunts subjacents. Aquesta asimetria és típica dels functors oblit i anticipa el teorema del capítol següent: els functors oblit solen ser adjunts per la dreta, i els adjunts per la dreta preserven tots els límits; no hi ha cap afirmació anàloga que els obligui a preservar tots els colímits.
\end{exemple}

\begin{proposicio}\label{prop:representables-limits}\index{functor!representable!preserva límits}
Per a tot objecte $A$ d'una categoria localment petita, el functor representable $\Hom(A,-)\colon\cat{C}\to\cat{Set}$ preserva tots els límits que existeixen a $\cat{C}$.
\end{proposicio}

\begin{prova}
Sigui $(L,\lambda)$ un límit d'un diagrama $D\colon\cat{I}\to\cat{C}$. Escrivim $H=\Hom(A,-)$, de manera que $H(X)=\Hom(A,X)$ i $H(v)=v_*\colon\varphi\mapsto v\comp\varphi$.

\emph{El con imatge.} La família $H\lambda=\bigl((\lambda_i)_*\bigr)_i$, amb $(\lambda_i)_*\colon\Hom(A,L)\to\Hom(A,D_i)$, és un con sobre $H\comp D$: per a $u\colon i\to j$, la functorialitat de $H$ dona $(D_u)_*\comp(\lambda_i)_*=(D_u\comp\lambda_i)_*=(\lambda_j)_*$.

\emph{Existència de la factorització.} Sigui $(S,\sigma)$ un con sobre $H\comp D$ a $\cat{Set}$. És a dir, $S$ és un conjunt i $\sigma_i\colon S\to\Hom(A,D_i)$ són aplicacions amb $(D_u)_*\comp\sigma_i=\sigma_j$. Avaluada en un element $s\in S$, aquesta condició diu $D_u\comp\sigma_i(s)=\sigma_j(s)$: la família $\bigl(\sigma_i(s)\bigr)_i$ és un con sobre $D$ amb vèrtex $A$. Per la propietat universal de $L$, hi ha una única fletxa $h(s)\colon A\to L$ amb $\lambda_i\comp h(s)=\sigma_i(s)$ per a tot $i$. Això defineix una aplicació $h\colon S\to\Hom(A,L)$ amb $(\lambda_i)_*\comp h=\sigma_i$ per a tot $i$.

\emph{Unicitat.} Si $h'\colon S\to\Hom(A,L)$ també compleix $(\lambda_i)_*\comp h'=\sigma_i$, aleshores, per a cada $s$, $\lambda_i\comp h'(s)=\sigma_i(s)$ per a tot $i$, i la unicitat de la factorització a través de $L$ dona $h'(s)=h(s)$. Per tant $h'=h$.

Així $\bigl(\Hom(A,L),H\lambda\bigr)$ és un límit de $H\comp D$, i $\Hom(A,-)$ preserva el límit $(L,\lambda)$.
\end{prova}

\begin{observacio}[La prova en una línia]
La demostració anterior és la lectura element a element de la cadena de bijeccions
\[
\Hom\bigl(A,\varprojlim D\bigr)\;\cong\;\mathrm{Con}(A,D)\;\cong\;\varprojlim\nolimits_i\Hom(A,D_i).
\]
La primera bijecció és la proposició~\ref{prop:limit-representable}. La segona és la descripció dels límits a $\cat{Set}$ com a famílies compatibles (exemple~\ref{ex:limits-set}): un element de $\varprojlim_i\Hom(A,D_i)$ és una família $(\alpha_i)_i$ amb $D_u\comp\alpha_i=\alpha_j$, és a dir, un con. La composta envia $f$ a $(\lambda_i\comp f)_i$, que és l'acció del con imatge.
\end{observacio}

\begin{observacio}
Aquest fet ---els representables covariants preserven límits--- és una peça clau que retrobarem. Dualment, els representables contravariants $\Hom(-,B)$ envien colímits a límits: com que $\Hom_{\cat{C}}(-,B)=\Hom_{\cat{C}\op}(B,-)$ i un colímit de $\cat{C}$ és un límit de $\cat{C}\op$, n'hi ha prou d'aplicar la proposició a $\cat{C}\op$. Per exemple, $\Hom(A+A',B)\cong\Hom(A,B)\times\Hom(A',B)$. Combinat amb el lema de Yoneda, permet raonar sobre límits component a component, i és la base tècnica de molts arguments de la teoria.
\end{observacio}

\begin{resumcap}
\apartat{Idees centrals}
\begin{enumerate}
  \item Un límit és el con universal sobre un diagrama; un colímit, el cocon universal, és a dir, el límit del diagrama corresponent a la categoria oposada.
  \item El límit d'un diagrama $D$ representa el prefeix de cons: $\Hom(-,\varprojlim D)\cong\mathrm{Con}(-,D)$. Per tant, el con límit és únic llevat d'un únic isomorfisme compatible amb les projeccions, i el vèrtex sol, llevat d'isomorfisme.
  \item Productes, equalitzadors, pullbacks i l'objecte terminal són límits; els seus duals són colímits. Tenir tots els límits finits equival a tenir objecte terminal, productes binaris i equalitzadors; equivalentment, objecte terminal i pullbacks.
  \item $\cat{Set}$ té tots els límits i colímits petits (famílies compatibles i quocients de sumes), i les categories de functors els calculen punt a punt quan existeixen a la categoria d'arribada.
  \item Els pullbacks indueixen la reindexació $f^{*}\colon\cat{C}/B\to\cat{C}/A$, que a $\cat{Set}$ reindexa fibres i anticipa la substitució de predicats.
  \item Els representables $\Hom(A,-)$ preserven límits.
\end{enumerate}
\apartat{Errors típics} Confondre el con (la família de projeccions) amb el seu vèrtex (l'objecte); creure que \enquote{tenir límits} depèn de la tria d'un límit concret, quan la propietat universal el fixa llevat d'isomorfisme; oblidar que un colímit no és un límit \emph{de la mateixa} categoria, sinó de l'oposada; suposar que tot functor preserva límits (només ho fan certs functors, com els representables covariants); i confondre \enquote{$F$ preserva un límit}, que parteix d'un límit de $\cat{C}$ i en demana la imatge, amb \enquote{$F$ reflecteix límits}, que parteix d'un con de $\cat{C}$ amb imatge limitant i en conclou que ja era un límit.
\apartat{Lectura recomanada} \cite[cap.~5]{awodey} per a l'exposició paral·lela a aquesta; \cite[III.3--III.4 i cap.~V]{maclane} per al tractament clàssic de límits com a cons universals; \cite[cap.~3]{riehl} per a exemples addicionals, en particular §3.3, sobre preservació, reflexió i creació de límits, i §3.4, sobre la naturalesa representable dels límits; \cite[cap.~9]{barrwells} per a la intuïció del límit com a subconjunt d'un producte tallat per equacions, i del colímit com a quocient d'una suma; \cite[cap.~3]{perrone-notes} per al prefeix de cons i el cas dels preordres; \cite[cap.~5]{leinster} per a una presentació breu, amb la interacció entre functors i límits; i \cite[part~IV, sessions~19--25]{lawvereschanuel}, per a objectes terminals, productes i sumes, amb els productes de grafs, en una presentació molt elemental.
\end{resumcap}

\chapter[Adjuncions]{Adjuncions}
\label{cap:adjuncions}

\begin{objectius}
\apartat{Prerequisits} Functors, preordres vistos com a categories i la categoria lliure d'un graf (cap.~\ref{cap:functors}); transformacions i isomorfismes naturals, i equivalències de categories (cap.~\ref{cap:transformacions}); representabilitat i propietats universals (cap.~\ref{cap:yoneda}); límits i colímits, i la reindexació per pullback (cap.~\ref{cap:limits}).
\apartat{Objectius} En acabar el capítol, el lector hauria de poder:
\begin{itemize}
  \item definir una adjunció mitjançant una bijecció natural d'homconjunts i mitjançant unitat i counitat;
  \item calcular transposats i reconèixer la propietat universal de la unitat;
  \item verificar les identitats triangulars;
  \item reconèixer adjuncions en preordres i en construccions lliure--oblit;
  \item interpretar $\colimop_{\cat{I}}\dashv\Delta\dashv\varprojlim_{\cat{I}}$ i la currificació com a adjuncions;
  \item interpretar a $\cat{Set}$ la cadena $\exists_f\dashv f^{*}\dashv\forall_f$ com a prototip de la quantificació;
  \item demostrar la unicitat dels adjunts llevat d'isomorfisme natural;
  \item aplicar que els adjunts per la dreta preserven límits i els adjunts per l'esquerra, colímits.
\end{itemize}
\end{objectius}

\section{Cap a la noció d'adjunció}

Al llarg del llibre hem trobat ja la parella formada pel functor categoria lliure $\Free\colon\cat{Graph}\to\cat{Cat}$ i el functor graf subjacent $U\colon\cat{Cat}\to\cat{Graph}$. Un altre patró clàssic, que veurem ara en una versió elemental, és el del monoide lliure i el functor oblit $U\colon\cat{Mon}\to\cat{Set}$. En tots dos casos, un functor afegeix estructura de la manera més lliure possible i l'altre l'oblida, i tots dos estan lligats per una correspondència precisa. Aquesta relació s'anomena \textbf{adjunció}, i és un dels conceptes centrals de la teoria de categories.

Les adjuncions explicaran també un patró que vam observar en estudiar els límits: els representables i molts functors oblit preserven límits de manera sistemàtica, mentre que moltes construccions lliures preserven colímits. El teorema d'aquest capítol explicarà aquests fenòmens quan els functors en qüestió siguin adjunts. I seran un pont directe cap a la lògica: els quantificadors existencial i universal es revelaran com a adjunts.

La idea es pot copsar primer en el cas més simple, el dels preordres, on una adjunció és una parella d'aplicacions monòtones lligades per una equivalència de desigualtats. Després la generalitzarem a categories qualssevol, substituint les desigualtats per bijeccions entre conjunts de morfismes.

\section{Adjuncions entre preordres}

Siguin $(P,\leq)$ i $(Q,\leq)$ dos preordres, vistos com a categories, i $f\colon P\to Q$, $g\colon Q\to P$ dues aplicacions monòtones.

\begin{definicio}\index{adjunció!entre preordres}
Diem que $f$ és \textbf{adjunta per l'esquerra} de $g$ (i $g$ \textbf{adjunta per la dreta} de $f$), i escrivim $f\dashv g$, quan per a tot $x\in P$ i tot $y\in Q$,
\[
f(x)\leq y
\quad\Longleftrightarrow\quad
x\leq g(y).
\]
\end{definicio}

\begin{observacio}[Unitat i counitat en un preordre]\index{connexió de Galois}
Si $f\dashv g$, prenent $y=f(x)$ i $x=g(y)$ a la definició s'obté
\[
x\leq g(f(x)),
\qquad
f(g(y))\leq y
\]
per a tot $x\in P$ i tot $y\in Q$. Recíprocament, si $f$ i $g$ són monòtones i compleixen aquestes dues desigualtats, aleshores $f\dashv g$: si $f(x)\leq y$, aleshores $x\leq g(f(x))\leq g(y)$, i si $x\leq g(y)$, aleshores $f(x)\leq f(g(y))\leq y$. Aquestes dues desigualtats són exactament la \emph{unitat} i la \emph{counitat} de l'adjunció, que veurem en general més endavant; en una categoria prima, les identitats triangulars no aporten cap igualtat nova, perquè entre dos objectes hi ha com a màxim una fletxa. Una adjunció entre preordres s'anomena sovint \textbf{connexió de Galois}.
\end{observacio}

\begin{exemple}[Conjunció i implicació]\index{adjunció!conjunció i implicació}
Sigui $\cat{Prov}_T$ el preordre de fórmules d'un sistema proposicional $T$ amb les regles usuals de $\wedge$ i $\to$, i fixem una fórmula $A$. Les operacions $A\wedge-$ i $A\to-$ són monòtones, i per a totes les fórmules $B$ i $C$
\[
A\wedge B\vdash_T C
\quad\Longleftrightarrow\quad
B\vdash_T A\to C,
\]
de manera que
\[
A\wedge-\;\dashv\;A\to-.
\]
Aquesta adjunció codifica les regles de la implicació; és la forma categòrica del teorema de deducció. És un primer indici que alguns connectius lògics es caracteritzen mitjançant adjuncions, idea que reapareixerà de manera sistemàtica en la lògica categòrica.
\end{exemple}

\begin{exemple}[Imatge, antiimatge i quantificadors]\label{ex:quantificadors-set}\index{adjunció!quantificadors}\index{quantificador!com a adjunt}
Sigui $f\colon X\to Y$ una aplicació, i considerem els preordres $\mathcal P(X)$ i $\mathcal P(Y)$ de les parts, ordenats per inclusió. L'antiimatge
\[
f^{*}(T)=f^{-1}(T),\qquad T\subseteq Y,
\]
és monòtona, i té adjunts a tots dos costats: per a $S\subseteq X$, posem
\[
\exists_f(S)=f[S]=\{f(x)\mid x\in S\},
\qquad
\forall_f(S)=\{y\in Y\mid f^{-1}(\{y\})\subseteq S\}.
\]
Aleshores, per a tot $S\subseteq X$ i tot $T\subseteq Y$,
\[
\exists_f(S)\subseteq T
\quad\Longleftrightarrow\quad
S\subseteq f^{*}(T),
\qquad\qquad
f^{*}(T)\subseteq S
\quad\Longleftrightarrow\quad
T\subseteq\forall_f(S).
\]
La primera equivalència diu que les imatges dels elements de $S$ són a $T$ exactament quan tots els elements de $S$ van a parar a $T$. Per a la segona: si $f^{-1}(T)\subseteq S$ i $y\in T$, tot $x$ amb $f(x)=y$ és a $f^{-1}(T)$ i, per tant, a $S$, de manera que $y\in\forall_f(S)$; i recíprocament, si $T\subseteq\forall_f(S)$ i $f(x)\in T$, aleshores $x\in f^{-1}(\{f(x)\})\subseteq S$. Per tant,
\[
\exists_f\;\dashv\;f^{*}\;\dashv\;\forall_f.
\]

El nom dels adjunts s'explica quan $f$ és una projecció $\pi\colon X\times A\to X$. Un subconjunt $S\subseteq X\times A$ és un predicat en dues variables, i
\[
\begin{aligned}
\exists_\pi(S)&=\{x\mid \text{existeix }a\in A\text{ amb }(x,a)\in S\},\\
\forall_\pi(S)&=\{x\mid (x,a)\in S\text{ per a tot }a\in A\}:
\end{aligned}
\]
són exactament les lectures conjuntistes de la quantificació existencial i universal sobre la variable d'$A$, mentre que $\pi^{*}$ afegeix una variable muda. Al capítol de subobjectes veurem en quines categories existeixen aquests dos adjunts; aquí només afirmem el cas de $\cat{Set}$.
\end{exemple}

\begin{exemple}[Interior topològic]\index{adjunció!interior topològic}
Sigui $X$ un espai topològic, $\mathcal P(X)$ el preordre de les parts per inclusió i $\mathcal O(X)$ el dels oberts. La injecció $i\colon\mathcal O(X)\to\mathcal P(X)$ té com a adjunt per la dreta l'\textbf{interior} $\mathrm{int}\colon\mathcal P(X)\to\mathcal O(X)$:
\[
i(U)\subseteq S
\quad\Longleftrightarrow\quad
U\subseteq\mathrm{int}(S),
\]
per a tot obert $U$ i tota part $S$. És a dir, l'interior és el més gran obert contingut en $S$.
\end{exemple}

\section{Adjuncions entre categories}

La definició general substitueix l'equivalència de desigualtats per una \textbf{bijecció entre conjunts de morfismes}, natural en totes dues variables.

\begin{definicio}[Adjunció]\index{adjunció}
\label{def:adjuncio}
Siguin $F\colon\cat{C}\to\cat{D}$ i $G\colon\cat{D}\to\cat{C}$ dos functors entre categories localment petites. Una \textbf{adjunció} $F\dashv G$ és una bijecció
\[
\theta_{A,B}\colon\Hom_{\cat{D}}(F(A),B)\;\xrightarrow{\ \cong\ }\;\Hom_{\cat{C}}(A,G(B)),
\]
per a cada $A$ de $\cat{C}$ i $B$ de $\cat{D}$, \textbf{natural} en $A$ i en $B$. Es diu que $F$ és l'\textbf{adjunt per l'esquerra} i $G$ l'\textbf{adjunt per la dreta}.
\end{definicio}

\begin{observacio}
La naturalitat vol dir que $\theta$ és un isomorfisme natural entre els dos functors
\[
\Hom_{\cat{D}}(F(-),-)\quad\text{i}\quad\Hom_{\cat{C}}(-,G(-)),
\]
tots dos de $\cat{C}\op\times\cat{D}$ cap a $\cat{Set}$. Explícitament, per a $a\colon A'\to A$ a $\cat{C}$, $b\colon B\to B'$ a $\cat{D}$ i $\varphi\colon F(A)\to B$, el pas per $\theta$ commuta amb la postcomposició i amb la precomposició:
\[
\theta_{A,B'}(b\comp\varphi)=G(b)\comp\theta_{A,B}(\varphi),
\qquad\qquad
\theta_{A',B}\bigl(\varphi\comp F(a)\bigr)=\theta_{A,B}(\varphi)\comp a.
\]
En paraules: transposar i després compondre amb $G(b)$ o amb $a$ és el mateix que compondre primer amb $b$ o amb $F(a)$ i després transposar. Són els dos quadrats de naturalitat
\[
\begin{tikzpicture}[baseline=(m.center)]
  \matrix (m) [matrix of math nodes, row sep=2.8em, column sep=3.6em] {
    \Hom(F(A),B) & \Hom(A,G(B)) \\
    \Hom(F(A),B') & \Hom(A,G(B')) \\
  };
  \draw[-{Stealth[scale=1.1]}] (m-1-1) -- node[above] {$\theta_{A,B}$} (m-1-2);
  \draw[-{Stealth[scale=1.1]}] (m-1-1) -- node[left] {$b\comp-$} (m-2-1);
  \draw[-{Stealth[scale=1.1]}] (m-1-2) -- node[right] {$G(b)\comp-$} (m-2-2);
  \draw[-{Stealth[scale=1.1]}] (m-2-1) -- node[below] {$\theta_{A,B'}$} (m-2-2);
\end{tikzpicture}
\]
i
\[
\begin{tikzpicture}[baseline=(m.center)]
  \matrix (m) [matrix of math nodes, row sep=2.8em, column sep=3.6em] {
    \Hom(F(A),B) & \Hom(A,G(B)) \\
    \Hom(F(A'),B) & \Hom(A',G(B)) \\
  };
  \draw[-{Stealth[scale=1.1]}] (m-1-1) -- node[above] {$\theta_{A,B}$} (m-1-2);
  \draw[-{Stealth[scale=1.1]}] (m-1-1) -- node[left] {$-\comp F(a)$} (m-2-1);
  \draw[-{Stealth[scale=1.1]}] (m-1-2) -- node[right] {$-\comp a$} (m-2-2);
  \draw[-{Stealth[scale=1.1]}] (m-2-1) -- node[below] {$\theta_{A',B}$} (m-2-2);
\end{tikzpicture}
\]
on els vèrtexs són conjunts de morfismes i les fletxes verticals actuen per composició. Els morfismes que es corresponen sota $\theta$ s'anomenen \textbf{transposats} l'un de l'altre: si $\varphi\colon F(A)\to B$, el seu transposat és $\theta(\varphi)\colon A\to G(B)$.
\end{observacio}

\begin{observacio}
El cas dels preordres n'és un cas particular: si $P,Q$ són preordres, un conjunt de morfismes $\Hom(x,y)$ té com a molt un element, de manera que la bijecció $\Hom(f(x),y)\cong\Hom(x,g(y))$ es redueix a l'equivalència $f(x)\leq y\Leftrightarrow x\leq g(y)$. Així, la definició general estén literalment la dels preordres.
\end{observacio}

\section{Unitat i counitat}

Una adjunció es pot descriure, de manera equivalent, sense esmentar la bijecció, mitjançant dues transformacions naturals distingides.

\begin{definicio}[Unitat i counitat]\index{adjunció!unitat}\index{adjunció!counitat}
\label{def:unitat-counitat}
Donada una adjunció $F\dashv G$ amb bijecció $\theta$:
\begin{itemize}
\item la \textbf{unitat} és la transformació natural $\eta\colon\id_{\cat{C}}\Rightarrow G\comp F$ de components
\[
\eta_A=\theta_{A,F(A)}(\id_{F(A)})\colon A\to G(F(A));
\]
\item la \textbf{counitat} és la transformació natural $\varepsilon\colon F\comp G\Rightarrow\id_{\cat{D}}$ de components
\[
\varepsilon_B=\theta_{G(B),B}^{-1}(\id_{G(B)})\colon F(G(B))\to B.
\]
\end{itemize}
\end{definicio}

\begin{observacio}[De la unitat i la counitat a la bijecció]\label{obs:adj-transposats}
La unitat i la counitat permeten recuperar la bijecció. Donat $\varphi\colon F(A)\to B$, el seu transposat és
\[
\theta(\varphi)=G(\varphi)\comp\eta_A\colon A\to G(B),
\]
i dualment, donat $\psi\colon A\to G(B)$, el seu transposat invers és
\[
\theta^{-1}(\psi)=\varepsilon_B\comp F(\psi)\colon F(A)\to B.
\]
\end{observacio}

\begin{definicio}[Fletxa universal]\label{def:fletxa-universal}\index{fletxa universal}\index{adjunció!fletxa universal}
Sigui $G\colon\cat{D}\to\cat{C}$ un functor i $A$ un objecte de $\cat{C}$. Una \textbf{fletxa universal} de $A$ cap a $G$ és una parella $(R,\eta_A)$ formada per un objecte $R$ de $\cat{D}$ i un morfisme $\eta_A\colon A\to G(R)$ tal que tot $\psi\colon A\to G(B)$ factoritza de manera única com $\psi=G(\varphi)\comp\eta_A$, amb $\varphi\colon R\to B$.
\end{definicio}

En una adjunció $F\dashv G$, cada component de la unitat és una fletxa universal, amb $R=F(A)$ (exercici~\ref{exr:adj-unitat-universal} dels Exercicis):
\[
\begin{tikzpicture}[baseline=(m.center)]
  \matrix (m) [matrix of math nodes, column sep=4em, row sep=2.7em] {
    A & G(F(A)) \\
      & G(B) \\
  };
  \draw[-{Stealth[scale=1.1]}] (m-1-1) -- node[above] {$\eta_A$} (m-1-2);
  \draw[-{Stealth[scale=1.1]}] (m-1-1) -- node[below left] {$\psi$} (m-2-2);
  \draw[-{Stealth[scale=1.1]},dashed] (m-1-2) -- node[right] {$G(\varphi)$} (m-2-2);
\end{tikzpicture}
\qquad
\psi=G(\varphi)\comp\eta_A,\quad\varphi\text{ únic}.
\]
Aquesta és la formulació de l'adjunció que connecta directament amb els objectes lliures.

\begin{exemple}[El monoide lliure]\index{adjunció!lliure-oblit}
Sigui $U\colon\cat{Mon}\to\cat{Set}$ el functor oblit dels monoides i $F\colon\cat{Set}\to\cat{Mon}$ el functor que envia un conjunt $X$ al \textbf{monoide lliure} $X^*$ de les paraules finites en $X$, amb la concatenació. Aleshores $F\dashv U$: un homomorfisme de monoides $X^*\to M$ està determinat de manera única pels valors sobre les lletres, és a dir, per una aplicació $X\to U(M)$. Aquesta és exactament la bijecció
\[
\Hom_{\cat{Mon}}(X^*,M)\cong\Hom_{\cat{Set}}(X,U(M)).
\]
La unitat $\eta_X\colon X\to U(X^*)$ envia cada element a la paraula d'una lletra; és la fletxa universal que exhibeix la propietat universal del monoide lliure. Aquest patró ---\emph{lliure $\dashv$ oblit}--- es repeteix per a grups, espais vectorials, categories i moltes altres estructures.
\end{exemple}

\begin{exemple}[La categoria lliure d'un graf: $\Free\dashv U$]\label{ex:free-u-adjuncio}\index{adjunció!categoria lliure}
Els capítols anteriors ja han preparat aquest exemple. El functor graf subjacent $U\colon\cat{Cat}\to\cat{Graph}$ i el functor categoria lliure $\Free\colon\cat{Graph}\to\cat{Cat}$ satisfan la bijecció
\[
\Hom_{\cat{Cat}}\bigl(\Free(G),\cat{D}\bigr)\;\cong\;\Hom_{\cat{Graph}}\bigl(G,U(\cat{D})\bigr),
\]
natural en $G$ i en $\cat{D}$, que el capítol~\ref{cap:yoneda} va obtenir com la representabilitat del functor $\cat{D}\mapsto\Hom_{\cat{Graph}}(G,U(\cat{D}))$. És a dir, $\Free\dashv U$. El que aporta ara el llenguatge de les adjuncions és identificar-ne les dades:
\begin{itemize}
  \item la \emph{unitat} $\eta_G\colon G\to U(\Free(G))$ és la inserció dels generadors: cada vèrtex va a ell mateix i cada aresta, al camí de longitud~$1$ que la recorre;
  \item la \emph{counitat} $\varepsilon_{\cat{C}}\colon\Free(U(\cat{C}))\to\cat{C}$ \emph{avalua} els camins formals: envia un camí $(f_1,\dots,f_n)$ de fletxes de $\cat{C}$, recorregut en aquest ordre, a la composta $f_n\comp\cdots\comp f_1$, i el camí buit en un objecte, a la seva identitat.
\end{itemize}
\end{exemple}

\section{Identitats triangulars}

La unitat i la counitat no són independents: satisfan dues equacions que, de fet, caracteritzen les adjuncions.

\begin{proposicio}[Identitats triangulars]\label{prop:triangulars}\index{adjunció!identitats triangulars}
Per a una adjunció $F\dashv G$ amb unitat $\eta$ i counitat $\varepsilon$, es compleixen les \textbf{identitats triangulars}:
\[
\varepsilon_{F(A)}\comp F(\eta_A)=\id_{F(A)},
\qquad
G(\varepsilon_B)\comp\eta_{G(B)}=\id_{G(B)},
\]
és a dir, commuten els dos triangles
\[
\begin{tikzpicture}[baseline=(m.center)]
  \matrix (m) [matrix of math nodes, column sep=3.6em, row sep=2.7em] {
    F(A) & F(G(F(A))) \\
         & F(A) \\
  };
  \draw[-{Stealth[scale=1.1]}] (m-1-1) -- node[above] {$F(\eta_A)$} (m-1-2);
  \draw[-{Stealth[scale=1.1]}] (m-1-1) -- node[below left] {$\id_{F(A)}$} (m-2-2);
  \draw[-{Stealth[scale=1.1]}] (m-1-2) -- node[right] {$\varepsilon_{F(A)}$} (m-2-2);
\end{tikzpicture}
\qquad\qquad
\begin{tikzpicture}[baseline=(m.center)]
  \matrix (m) [matrix of math nodes, column sep=3.6em, row sep=2.7em] {
    G(B) & G(F(G(B))) \\
         & G(B) \\
  };
  \draw[-{Stealth[scale=1.1]}] (m-1-1) -- node[above] {$\eta_{G(B)}$} (m-1-2);
  \draw[-{Stealth[scale=1.1]}] (m-1-1) -- node[below left] {$\id_{G(B)}$} (m-2-2);
  \draw[-{Stealth[scale=1.1]}] (m-1-2) -- node[right] {$G(\varepsilon_B)$} (m-2-2);
\end{tikzpicture}
\]
Recíprocament, dues transformacions naturals $\eta\colon\id\Rightarrow G F$ i $\varepsilon\colon F G\Rightarrow\id$ que les satisfacin defineixen una adjunció $F\dashv G$.
\end{proposicio}

\begin{prova}
Comprovem la primera identitat. Pel càlcul del transposat, $\theta^{-1}(\psi)=\varepsilon_B\comp F(\psi)$; prenent $B=F(A)$ i $\psi=\eta_A=\theta(\id_{F(A)})$,
\[
\varepsilon_{F(A)}\comp F(\eta_A)=\theta^{-1}(\eta_A)=\theta^{-1}(\theta(\id_{F(A)}))=\id_{F(A)}.
\]
La segona identitat és dual, usant $\theta(\varphi)=G(\varphi)\comp\eta_A$ amb $A=G(B)$ i $\varphi=\varepsilon_B$. Per al recíproc, es defineix la bijecció per $\theta(\varphi)=G(\varphi)\comp\eta_A$ i $\theta^{-1}(\psi)=\varepsilon_B\comp F(\psi)$; les identitats triangulars garanteixen que són inverses l'una de l'altra, i la naturalitat de $\eta,\varepsilon$ dona la de $\theta$.
\end{prova}

\begin{observacio}
Les identitats triangulars reben aquest nom perquè expressen la commutativitat de dos triangles que involucren $F$, $G$, $\eta$ i $\varepsilon$. Donen una definició d'adjunció \emph{purament diagramàtica}, sense referència a conjunts de morfismes, la qual cosa és sovint la més còmoda per treballar i l'única disponible en contextos on no hi ha homconjunts (per exemple, en 2-categories abstractes).
\end{observacio}

\begin{observacio}[Tres presentacions d'una adjunció]\label{obs:tres-presentacions}\index{adjunció!tres presentacions}
Hem trobat tres maneres de donar una adjunció entre $\cat{C}$ i $\cat{D}$. Convé no barrejar-les, perquè no parteixen de les mateixes dades.
\begin{center}
\footnotesize
\renewcommand{\arraystretch}{1.3}
\begin{tabularx}{\linewidth}{@{}>{\raggedright\arraybackslash}p{0.2\linewidth}>{\raggedright\arraybackslash}X>{\raggedright\arraybackslash}X@{}}
\toprule
\textbf{Presentació} & \textbf{Dades} & \textbf{Condicions}\\
\midrule
bijecció d'homconjunts (definició~\ref{def:adjuncio}) & functors $F$ i $G$; bijeccions $\theta_{A,B}\colon\Hom(F(A),B)\to\Hom(A,G(B))$ & naturalitat en $A$ i en $B$\\
unitat i counitat (proposició~\ref{prop:triangulars}) & functors $F$ i $G$; transformacions naturals $\eta\colon\id\Rightarrow GF$ i $\varepsilon\colon FG\Rightarrow\id$ & identitats triangulars\\
fletxes universals (definició~\ref{def:fletxa-universal}) & només el functor $G$; per a cada $A$, un objecte $F(A)$ i una fletxa $\eta_A\colon A\to G(F(A))$ & cada $\eta_A$ és universal; $F$ no es dona sobre morfismes, sinó que la propietat universal el determina\\
\bottomrule
\end{tabularx}
\end{center}
Els passos d'una presentació a una altra són aquests:
\begin{itemize}
\item De la primera a la segona: $\eta$ i $\varepsilon$ són els transposats de les identitats (definició~\ref{def:unitat-counitat}).
\item De la segona a la primera: les fórmules de l'observació~\ref{obs:adj-transposats}, que són inverses gràcies a les identitats triangulars (exercici~\ref{exr:adj-triangulars-reciproc} dels Exercicis).
\item De la primera a la tercera: la unitat és universal (exercici~\ref{exr:adj-unitat-universal}).
\item De la tercera a la primera: la propietat universal defineix $F$ sobre morfismes i la bijecció $\varphi\mapsto G(\varphi)\comp\eta_A$ (exercici~\ref{exr:adj-fletxes-universals}).
\end{itemize}
La tercera presentació és la més econòmica, perquè no cal conèixer $F$ per endavant; per això és la que s'usa per \emph{construir} adjunts, com els objectes lliures. La segona és la més còmoda per calcular, i la primera és la que dona directament els teoremes de preservació de límits de la secció següent.
\end{observacio}

\begin{observacio}[Adjuncions i equivalències]\index{equivalència de categories!adjunta}
Al capítol~\ref{cap:transformacions} vam definir una equivalència mitjançant un quasiinvers $G$ i dos isomorfismes naturals orientats com $\eta\colon\id_{\cat{C}}\Rightarrow G\comp F$ i $\varepsilon\colon F\comp G\Rightarrow\id_{\cat{D}}$ (definició~\ref{def:equivalencia}), la mateixa orientació que la unitat i la counitat d'una adjunció $F\dashv G$. La coincidència no és casual. Tota equivalència es pot equipar amb una estructura d'\textbf{equivalència adjunta}: modificant, si cal, un dels dos isomorfismes naturals, es pot aconseguir que $\eta$ i $\varepsilon$ satisfacin les identitats triangulars, i aleshores $F\dashv G$ amb unitat i counitat isomorfismes naturals. Recíprocament, una adjunció en què la unitat i la counitat són isomorfismes naturals és una equivalència. Les equivalències són, doncs, les adjuncions més simètriques possibles.
\end{observacio}

\section{Unicitat dels adjunts}

\begin{proposicio}\index{adjunció!unicitat}
Els adjunts són únics llevat d'isomorfisme natural. Si $F\dashv G$ i $F\dashv G'$, aleshores $G\cong G'$; i si $F\dashv G$ i $F'\dashv G$, aleshores $F\cong F'$.
\end{proposicio}

\begin{prova}
Si $F\dashv G$ i $F\dashv G'$, aleshores per a tot $A,B$,
\[
\Hom(A,G(B))\cong\Hom(F(A),B)\cong\Hom(A,G'(B)),
\]
naturalment en $A$. Pel corol·lari del lema de Yoneda ---dos objectes amb functors representables naturalment isomorfs són isomorfs---, $G(B)\cong G'(B)$ naturalment en $B$, és a dir, $G\cong G'$. El cas de $F$ és dual.
\end{prova}

\section{Els adjunts i els límits}

El resultat més útil sobre adjuncions diu com interactuen amb límits i colímits. És la raó de fons per la qual molts functors oblit preserven els productes i, en general, els límits.

\begin{teorema}[Els adjunts per la dreta preserven límits]\index{adjunció!preservació de límits}
\label{teo:adjunts-preserven-limits}
Si $G\colon\cat{D}\to\cat{C}$ és un adjunt per la dreta, aleshores $G$ preserva tots els límits petits que existeixen a $\cat{D}$. Dualment, un adjunt per l'esquerra preserva tots els colímits petits que existeixen.
\end{teorema}

\begin{prova}
Sigui $F\colon\cat{C}\to\cat{D}$ l'adjunt per l'esquerra de $G$, amb bijecció $\theta$ (definició~\ref{def:adjuncio}). Farem servir les dues equacions de naturalitat de $\theta$: per a $b\colon B\to B'$ a $\cat{D}$ i $a\colon A'\to A$ a $\cat{C}$,
\[
\theta(b\comp\varphi)=G(b)\comp\theta(\varphi),
\qquad\qquad
\theta\bigl(\varphi\comp F(a)\bigr)=\theta(\varphi)\comp a.
\]

\emph{Límits.} Sigui $(L,\lambda)$ un límit d'un diagrama $D\colon\cat{I}\to\cat{D}$. La família $G\lambda=(G(\lambda_i))_i$ és un con sobre $G\comp D$, perquè $G(D_u)\comp G(\lambda_i)=G(D_u\comp\lambda_i)=G(\lambda_j)$ per a cada $u\colon i\to j$. Vegem que és universal. Sigui $(A,\alpha)$ un con sobre $G\comp D$, amb $\alpha_i\colon A\to G(D_i)$ i $G(D_u)\comp\alpha_i=\alpha_j$, i siguin $\varphi_i=\theta^{-1}(\alpha_i)\colon F(A)\to D_i$ els transposats. Per la primera equació de naturalitat,
\[
\theta(D_u\comp\varphi_i)=G(D_u)\comp\alpha_i=\alpha_j=\theta(\varphi_j),
\]
i com que $\theta$ és bijectiva, $D_u\comp\varphi_i=\varphi_j$: les $\varphi_i$ formen un con sobre $D$ amb vèrtex $F(A)$. Per tant hi ha una única $h\colon F(A)\to L$ amb $\lambda_i\comp h=\varphi_i$ per a tot $i$. Posem $k=\theta(h)\colon A\to G(L)$. De nou per la primera equació,
\[
G(\lambda_i)\comp k=\theta(\lambda_i\comp h)=\theta(\varphi_i)=\alpha_i,
\]
de manera que $k$ factoritza el con $\alpha$. És l'única factorització: si $k'\colon A\to G(L)$ compleix $G(\lambda_i)\comp k'=\alpha_i$ per a tot $i$, aleshores $h'=\theta^{-1}(k')$ compleix $\theta(\lambda_i\comp h')=G(\lambda_i)\comp k'=\alpha_i=\theta(\varphi_i)$, d'on $\lambda_i\comp h'=\varphi_i$; per la unicitat de $h$, $h'=h$, i $k'=\theta(h')=k$. Així $(G(L),G\lambda)$ és un límit de $G\comp D$.

\emph{Colímits.} Sigui ara $(K,\kappa)$ un colímit d'un diagrama $E\colon\cat{I}\to\cat{C}$. La família $F\kappa=(F(\kappa_i))_i$ és un cocon sobre $F\comp E$, perquè $F(\kappa_j)\comp F(E_u)=F(\kappa_j\comp E_u)=F(\kappa_i)$. Sigui $(B,\beta)$ un cocon sobre $F\comp E$, amb $\beta_i\colon F(E_i)\to B$ i $\beta_j\comp F(E_u)=\beta_i$, i siguin $\psi_i=\theta(\beta_i)\colon E_i\to G(B)$. Per la segona equació de naturalitat,
\[
\psi_j\comp E_u=\theta\bigl(\beta_j\comp F(E_u)\bigr)=\theta(\beta_i)=\psi_i,
\]
de manera que les $\psi_i$ formen un cocon sobre $E$, i hi ha una única $m\colon K\to G(B)$ amb $m\comp\kappa_i=\psi_i$. Posem $n=\theta^{-1}(m)\colon F(K)\to B$. Aleshores $\theta\bigl(n\comp F(\kappa_i)\bigr)=m\comp\kappa_i=\psi_i=\theta(\beta_i)$, és a dir, $n\comp F(\kappa_i)=\beta_i$. Si $n'$ també compleix $n'\comp F(\kappa_i)=\beta_i$, aleshores $\theta(n')\comp\kappa_i=\theta(\beta_i)=\psi_i$, d'on $\theta(n')=m$ i $n'=n$. Així $(F(K),F\kappa)$ és un colímit de $F\comp E$.
\end{prova}

\begin{observacio}[La demostració en termes de representables]
La demostració anterior és la versió \enquote{element a element} d'una cadena de bijeccions. Per a tot objecte $A$ de $\cat{C}$, usant l'adjunció i que $\Hom(F(A),-)$ preserva límits (proposició~\ref{prop:representables-limits}),
\[
\begin{aligned}
\Hom(A,G(L))&\cong\Hom(F(A),L) &&\text{(adjunció)}\\
&\cong\varprojlim\nolimits_i\Hom(F(A),D_i) &&\text{($\Hom(F(A),-)$ preserva límits)}\\
&\cong\varprojlim\nolimits_i\Hom(A,G(D_i)) &&\text{(adjunció, natural en $D_i$)}\\
&\cong\mathrm{Con}(A,G\comp D) &&\text{(límits a $\cat{Set}$)},
\end{aligned}
\]
naturalment en $A$; la bijecció composta envia $f\colon A\to G(L)$ a $(G(\lambda_i)\comp f)_i$. Llegida així, l'afirmació diu que $G(L)$ representa el prefeix dels cons sobre $G\comp D$, amb element universal $G\lambda$ (proposició~\ref{prop:limit-representable}).
\end{observacio}

\begin{observacio}[Què vol dir exactament \enquote{preservar}]\label{obs:preservar-canonic}\index{morfisme de comparació}
El teorema no afirma cap igualtat $G(\varprojlim D)=\varprojlim(G\comp D)$: els límits només estan determinats llevat d'isomorfisme, i una igualtat no tindria sentit. Tampoc no afirma només que \emph{existeixi} algun isomorfisme entre aquests dos objectes. Afirma el que diu la definició del capítol~\ref{cap:limits}: el con imatge $(G(L),G\lambda)$ \emph{és un con límit} de $G\comp D$.

Si hem triat un límit $\bigl(\varprojlim(G\comp D),\pi\bigr)$ de $G\comp D$, la condició es pot expressar amb un sol morfisme. El con $G\lambda$ factoritza de manera única a través del límit triat, i dona l'únic \textbf{morfisme de comparació}
\[
c\colon G(L)\to\varprojlim(G\comp D),\qquad \pi_i\comp c=G(\lambda_i)\ \text{per a tot } i.
\]
Aleshores $G$ preserva el límit si i només si $c$ és un isomorfisme. En efecte, si $(G(L),G\lambda)$ és un límit, l'únic isomorfisme compatible amb les projeccions entre els dos límits de $G\comp D$ compleix $\pi_i\comp c=G(\lambda_i)$, i per tant és $c$. Recíprocament, si $c$ és invertible, el con $G\lambda$ és isomorf al con límit $\pi$ per un isomorfisme compatible amb les projeccions, i per tant és un límit. En aquest sentit precís escrivim $G(\varprojlim D)\cong\varprojlim(G\comp D)$: l'isomorfisme és el \emph{canònic} $c$, no un isomorfisme qualsevol. A més, la noció no depèn del límit triat de $D$, perquè dos límits de $D$ són isomorfs per un únic isomorfisme compatible amb les projeccions, i $G$ el transporta.

La diferència importa. El functor $X\mapsto X+X$ de $\cat{Set}$ en si mateix no preserva productes binaris. El morfisme de comparació $(X\times Y)+(X\times Y)\to(X+X)\times(Y+Y)$ envia cada parella de la còpia $k$ a la parella de còpies $(k,k)$, i no arriba als elements en què les dues còpies són diferents. Tanmateix, per a $X=Y=\mathbb N$ tots dos conjunts són numerables, i per tant \emph{hi ha} una bijecció entre ells. Dualment, $F$ preserva el colímit $(K,\kappa)$ de $E$ quan el morfisme de comparació $\colimop(F\comp E)\to F(K)$ és un isomorfisme.
\end{observacio}

\begin{observacio}
Aquest teorema es recorda amb les sigles angleses \enquote{RAPL} (\emph{right adjoints preserve limits}) i \enquote{LAPC} (\emph{left adjoints preserve colimits}). Explica de cop molts fets. Els functors oblit que són adjunts per la dreta preserven límits, però això no implica que preservin colímits; en molts exemples algebraics no els preserven. Dualment, els functors lliures, com a adjunts per l'esquerra, preserven colímits, però no hi ha cap teorema general que els obligui a preservar límits. El teorema dona una sola direcció i no prohibeix l'altra: un functor amb adjunts a tots dos costats, com l'oblit $\cat{Top}\to\cat{Set}$, preserva tant límits com colímits. És també un criteri negatiu útil: un functor que \emph{no} preserva límits no pot ser un adjunt per la dreta.
\end{observacio}

\subsection{El límit i el colímit com a adjunts del functor diagonal}\index{límit!com a adjunt}\index{colímit!com a adjunt}

Hi ha una segona manera, més sintètica, en què límits i colímits són adjuncions ---no ja \emph{preservats} per adjunts, sinó ells mateixos adjunts. Fixem una categoria d'índexs $\cat{I}$ i considerem el \textbf{functor diagonal}
\[
\Delta\colon\cat{C}\longrightarrow\cat{C}^{\cat{I}},
\]
que envia un objecte $A$ al diagrama constant de valor $A$ (i un morfisme a la transformació natural constant). Prendre el límit i prendre el colímit d'un diagrama són operacions $\cat{C}^{\cat{I}}\to\cat{C}$, i resulten ser, respectivament, l'adjunt per la dreta i per l'esquerra de $\Delta$.

\begin{teorema}[$\colimop\dashv\Delta\dashv\varprojlim$]\index{adjunció!límit-diagonal}
\label{teo:lim-adjunt}
Sigui $\cat{I}$ una categoria petita i suposem que $\cat{C}$ té tots els límits i colímits de forma $\cat{I}$. Fixem, per a cada diagrama de forma $\cat{I}$, un límit i un colímit; quan la família de diagrames és pròpiament gran, fem servir la convenció d'elecció global de l'apèndix~\ref{cap:mida}. La propietat universal estén aquestes tries, de manera única, a functors $\varprojlim_{\cat{I}}$ i $\colimop_{\cat{I}}$ (pas~2 de la demostració). Aleshores el functor diagonal $\Delta\colon\cat{C}\to\cat{C}^{\cat{I}}$ té adjunt per l'esquerra $\colimop_{\cat{I}}$ i adjunt per la dreta $\varprojlim_{\cat{I}}$:
\[
\colimop_{\cat{I}}\;\dashv\;\Delta\;\dashv\;\varprojlim_{\cat{I}}.
\]
\end{teorema}

\begin{prova}
Recordem que $\Delta$ envia $A$ al diagrama constant $\Delta A$, amb $(\Delta A)_i=A$ i $(\Delta A)_u=\id_A$, i un morfisme $a\colon A'\to A$ a la transformació natural $\Delta a$ de components $a$. Escrivim $(\varprojlim D,\lambda^D)$ per al límit triat de cada diagrama $D$.

\emph{Pas 1: un morfisme $\Delta A\Rightarrow D$ és un con.} Una transformació natural $\alpha\colon\Delta A\Rightarrow D$ té components $\alpha_i\colon A\to D_i$, i el seu quadrat de naturalitat per a $u\colon i\to j$ diu $D_u\comp\alpha_i=\alpha_j\comp\id_A=\alpha_j$. És exactament la condició de con. Per tant, $\Hom_{\cat{C}^{\cat{I}}}(\Delta A,D)$ és el conjunt dels cons sobre $D$ amb vèrtex $A$.

\emph{Pas 2: $\varprojlim$ és un functor.} Sigui $\beta\colon D\Rightarrow D'$ un morfisme de $\cat{C}^{\cat{I}}$. La família $\beta_i\comp\lambda^D_i$ és un con sobre $D'$, perquè $D'_u\comp\beta_i\comp\lambda^D_i=\beta_j\comp D_u\comp\lambda^D_i=\beta_j\comp\lambda^D_j$. Definim $\varprojlim\beta$ com l'única fletxa amb
\[
\lambda^{D'}_i\comp\varprojlim\beta=\beta_i\comp\lambda^D_i\qquad\text{per a tot } i.
\]
La identitat i $\varprojlim\id_D$ compleixen totes dues aquesta condició per a $\beta=\id_D$. Per a $\beta'\colon D'\Rightarrow D''$, la composició $\varprojlim\beta'\comp\varprojlim\beta$ la compleix per a $\beta'\comp\beta$. Per la unicitat de la factorització, $\varprojlim$ preserva identitats i composició.

\emph{Pas 3: la bijecció.} Definim $\theta_{A,D}\colon\Hom_{\cat{C}^{\cat{I}}}(\Delta A,D)\to\Hom_{\cat{C}}(A,\varprojlim D)$ enviant cada con $\alpha$ a l'única fletxa $\theta(\alpha)$ amb $\lambda^D_i\comp\theta(\alpha)=\alpha_i$ per a tot $i$. Per la propietat universal del límit, $\theta$ és bijectiva, amb inversa $h\mapsto(\lambda^D_i\comp h)_i=\lambda^D\comp\Delta h$.

\emph{Pas 4: naturalitat.} Per a $a\colon A'\to A$, les fletxes $\theta(\alpha\comp\Delta a)$ i $\theta(\alpha)\comp a$ compleixen totes dues $\lambda^D_i\comp(-)=\alpha_i\comp a$, i per tant coincideixen: és la naturalitat en $A$. Per a $\beta\colon D\Rightarrow D'$,
\[
\lambda^{D'}_i\comp\varprojlim\beta\comp\theta(\alpha)=\beta_i\comp\lambda^D_i\comp\theta(\alpha)=\beta_i\comp\alpha_i=\lambda^{D'}_i\comp\theta(\beta\comp\alpha),
\]
d'on $\varprojlim\beta\comp\theta(\alpha)=\theta(\beta\comp\alpha)$: és la naturalitat en $D$. Així $\Delta\dashv\varprojlim_{\cat{I}}$.

\emph{L'adjunció de l'esquerra} és la duala, pas per pas. Una transformació natural $\gamma\colon D\Rightarrow\Delta A$ té components $\gamma_i\colon D_i\to A$ amb $\gamma_j\comp D_u=\gamma_i$, és a dir, és un cocon. Si $(\colimop D,\kappa^D)$ és el colímit triat, $\colimop\beta$ és l'única fletxa amb $\colimop\beta\comp\kappa^D_i=\kappa^{D'}_i\comp\beta_i$, i és functorial per la mateixa raó. La bijecció $\Hom_{\cat{C}}(\colimop D,A)\to\Hom_{\cat{C}^{\cat{I}}}(D,\Delta A)$ és $h\mapsto(h\comp\kappa^D_i)_i$, bijectiva per la propietat universal del colímit. És natural en $A$, perquè $(a\comp h)\comp\kappa^D_i=a\comp(h\comp\kappa^D_i)$, i en $D$, perquè $h\comp\colimop\beta\comp\kappa^D_i=(h\comp\kappa^{D'}_i)\comp\beta_i$. Així $\colimop_{\cat{I}}\dashv\Delta$.
\end{prova}

\begin{observacio}
Aquesta fórmula és una de les síntesis més netes del llibre: recull en una sola línia tot el capítol de límits. En l'adjunció $\Delta\dashv\varprojlim_{\cat{I}}$, la \emph{counitat} en un diagrama $D$,
\[
\varepsilon_D\colon\Delta\bigl(\varprojlim\nolimits_{\cat{I}}D\bigr)\Rightarrow D,
\]
és exactament el con límit de $D$, és a dir, la família de projeccions; la \emph{unitat} en un objecte $A$ és el morfisme $A\to\varprojlim_{\cat{I}}\Delta A$. Dualment, en l'adjunció $\colimop_{\cat{I}}\dashv\Delta$, la \emph{unitat}
\[
\eta_D\colon D\Rightarrow\Delta\bigl(\colimop\nolimits_{\cat{I}}D\bigr)
\]
és el cocon colímit, la família d'injeccions, mentre que la \emph{counitat} és el morfisme $\colimop_{\cat{I}}\Delta A\to A$. Identificar quina de les dues transformacions és el con o el cocon universal és un bon exercici per fixar la diferència entre unitat i counitat. I la lectura és profètica: quan al capítol de subobjectes vegem l'existencial i l'universal com a adjunts de la reindexació $f^{*}$, estarem veient una versió \enquote{fibrada} d'aquesta mateixa idea ---$f^{*}$ fa el paper de $\Delta$, i els quantificadors, el de $\colimop$ i $\varprojlim$.
\end{observacio}

\section{Exponencials i currificació}

Recuperem ara la relació que vam ajornar en tractar les categories de functors: la que lliga el producte amb els objectes de functors.

\begin{proposicio}\label{prop:cat-currificacio}\index{adjunció!producte-exponencial}\index{Cat@$\cat{Cat}$!exponencials}
Per a categories petites $\cat{A},\cat{B},\cat{C}$, hi ha un isomorfisme natural
\[
\Hom_{\cat{Cat}}(\cat{A}\times\cat{B},\cat{C})\;\cong\;\Hom_{\cat{Cat}}(\cat{A},\cat{C}^{\cat{B}}).
\]
És a dir, el functor $-\times\cat{B}$ és adjunt per l'esquerra del functor $(-)^{\cat{B}}$:
\[
-\times\cat{B}\;\dashv\;(-)^{\cat{B}}.
\]
\end{proposicio}

\begin{prova}
Primer, els dos costats són functors $\cat{Cat}\to\cat{Cat}$. El functor $-\times\cat{B}$ envia $P\colon\cat{A}'\to\cat{A}$ a $P\times\id_{\cat{B}}$. El functor $(-)^{\cat{B}}$ envia $K\colon\cat{C}\to\cat{C}'$ a la composició amb $K$, $\cat{C}^{\cat{B}}\to\cat{C}'^{\cat{B}}$, que fa $H\mapsto K\comp H$ i $\eta\mapsto K\eta$ (capítol~\ref{cap:transformacions}). Si $\cat{B}$ i $\cat{C}$ són petites, $\cat{C}^{\cat{B}}$ també ho és. A $\cat{A}\times\cat{B}$ escrivim els morfismes com a parelles, que es componen component a component, i fem servir la descomposició
\[
(f,g)=(f,\id_{b'})\comp(\id_a,g)=(\id_{a'},g)\comp(f,\id_b)\qquad\text{per a } f\colon a\to a',\ g\colon b\to b'.
\tag{$\ast$}
\]

\emph{Currificació.} Sigui $F\colon\cat{A}\times\cat{B}\to\cat{C}$. Per a cada objecte $a$, $\Lambda F(a)=F(a,-)$ és el functor $\cat{B}\to\cat{C}$ que envia $b\mapsto F(a,b)$ i $g\mapsto F(\id_a,g)$. És un functor perquè $F$ ho és i $(\id_a,g')\comp(\id_a,g)=(\id_a,g'\comp g)$. Per a $f\colon a\to a'$, $\Lambda F(f)\colon F(a,-)\Rightarrow F(a',-)$ té components $F(f,\id_b)$. És natural perquè, aplicant $F$ a $(\ast)$, $F(\id_{a'},g)\comp F(f,\id_b)=F(f,g)=F(f,\id_{b'})\comp F(\id_a,g)$. Finalment, $\Lambda F\colon\cat{A}\to\cat{C}^{\cat{B}}$ és un functor: les components de $\Lambda F(\id_a)$ són identitats, i les de $\Lambda F(f'\comp f)$ són $F(f'\comp f,\id_b)=F(f',\id_b)\comp F(f,\id_b)$, que són les de la composició vertical $\Lambda F(f')\comp\Lambda F(f)$.

\emph{Descurrificació.} Sigui $H\colon\cat{A}\to\cat{C}^{\cat{B}}$. Definim $\bar H\colon\cat{A}\times\cat{B}\to\cat{C}$ per
\[
\bar H(a,b)=H(a)(b),\qquad \bar H(f,g)=H(f)_{b'}\comp H(a)(g)=H(a')(g)\comp H(f)_b,
\]
on la segona igualtat és la naturalitat de $H(f)$. Les identitats es preserven perquè $H(\id_a)$ és la transformació identitat i $H(a)$ és un functor. Per a $(f',g')\colon(a',b')\to(a'',b'')$,
\[
\begin{aligned}
\bar H(f',g')\comp\bar H(f,g)
&=H(f')_{b''}\comp H(a')(g')\comp H(f)_{b'}\comp H(a)(g)\\
&=H(f')_{b''}\comp H(f)_{b''}\comp H(a)(g')\comp H(a)(g)\\
&=H(f'\comp f)_{b''}\comp H(a)(g'\comp g),
\end{aligned}
\]
on hem fet servir la naturalitat de $H(f)$ per a $g'$, que $H$ preserva la composició (vertical, component a component) i que $H(a)$ és un functor. Això és $\bar H\bigl((f',g')\comp(f,g)\bigr)$.

\emph{Són inverses.} Per a $F$: $\overline{\Lambda F}(a,b)=F(a,b)$ i $\overline{\Lambda F}(f,g)=F(f,\id_{b'})\comp F(\id_a,g)=F(f,g)$, per $(\ast)$. Per a $H$: $\Lambda\bar H(a)$ envia $g$ a $\bar H(\id_a,g)=H(a)(g)$, de manera que $\Lambda\bar H(a)=H(a)$; i les components de $\Lambda\bar H(f)$ són $\bar H(f,\id_b)=H(f)_b$, de manera que $\Lambda\bar H(f)=H(f)$.

\emph{Naturalitat.} Per a $P\colon\cat{A}'\to\cat{A}$, el functor $\Lambda\bigl(F\comp(P\times\id_{\cat{B}})\bigr)$ envia $a'$ a $F(P(a'),-)$ i $f'$ a la transformació de components $F(P(f'),\id_b)$; és a dir, és $\Lambda F\comp P$. Per a $K\colon\cat{C}\to\cat{C}'$, $\Lambda(K\comp F)$ envia $a$ a $K\comp F(a,-)$ i $f$ a la transformació de components $K(F(f,\id_b))$, que és $K\,\Lambda F(f)$; és a dir, $\Lambda(K\comp F)=K^{\cat{B}}\comp\Lambda F$, on $K^{\cat{B}}$ és la imatge de $K$ pel functor $(-)^{\cat{B}}$. Aquestes dues igualtats són la naturalitat de $\Lambda$ en $\cat{A}$ i en $\cat{C}$. Per tant, $-\times\cat{B}\dashv(-)^{\cat{B}}$.
\end{prova}

\begin{observacio}
Aquesta adjunció ---\emph{producte $\dashv$ exponencial}--- és el prototip de les \textbf{categories cartesianes tancades}, que estudiarem al capítol següent i que són el marc categòric del càlcul lambda simplement tipat. La \enquote{currificació} és exactament la que apareix en programació funcional: transformar una funció de dues variables en una funció que retorna funcions.
\end{observacio}

\begin{resumcap}
\apartat{Idees centrals}
\begin{enumerate}
  \item Una adjunció $F\dashv G$ és una bijecció natural $\Hom(F(A),B)\cong\Hom(A,G(B))$, i dona els transposats.
  \item Equival a donar una unitat i una counitat que satisfan les identitats triangulars; la unitat expressa una propietat universal, i dualment la counitat.
  \item Exemples clau: les connexions de Galois entre preordres, $\Free\dashv U$ i, a $\cat{Set}$, la cadena $\exists_f\dashv f^{*}\dashv\forall_f$.
  \item $\colimop_{\cat{I}}\dashv\Delta\dashv\varprojlim_{\cat{I}}$ i $-\times B\dashv(-)^{B}$ mostren que moltes propietats universals són adjuncions.
  \item Els adjunts són únics llevat d'isomorfisme natural; els adjunts per la dreta preserven límits, i els adjunts per l'esquerra, colímits.
\end{enumerate}
\apartat{Errors típics} Intercanviar quin adjunt preserva límits i quin colímits; creure que un adjunt per la dreta \emph{no pot} preservar colímits, quan el teorema només dona una direcció; oblidar la naturalitat de la bijecció d'homconjunts; confondre unitat i counitat (per exemple, a $\Delta\dashv\varprojlim$ el con límit és la counitat, no la unitat); i tractar les mnemotècniques (\enquote{RAPL\,/\,LAPC}) com a terminologia establerta en lloc d'ajudes informals.
\apartat{Lectura recomanada} \cite[cap.~9]{awodey}, especialment §§9.4--9.6, per als adjunts d'ordre, els quantificadors com a adjunts i RAPL; \cite[cap.~2]{leinster} per a la unitat, la counitat i les fletxes universals; \cite[§§4.1--4.5]{riehl} per al càlcul d'adjuncions i la preservació de límits; \cite[cap.~4]{perrone-notes} per als exemples lliure--oblit, les connexions de Galois i l'adjunció entre categories i grafs; \cite[cap.~13]{barrwells} per a una exposició orientada a exemples; i \cite[cap.~IV]{maclane} per al tractament clàssic.
\end{resumcap}

\chapter[Cartesianes tancades]{\texorpdfstring{Categories cartesianes\\ tancades}{Categories cartesianes tancades}}
\label{cap:cartesianes}

\begin{objectius}
\apartat{Prerequisits} Representabilitat i lema de Yoneda (cap.~\ref{cap:yoneda}); productes i límits finits (cap.~\ref{cap:limits}); adjuncions i currificació (cap.~\ref{cap:adjuncions}). No cal cap coneixement previ de càlcul lambda: el capítol n'introdueix la semàntica necessària.
\apartat{Objectius} En acabar el capítol, el lector hauria de poder:
\begin{itemize}
  \item definir l'objecte exponencial per la seva propietat universal, com a representació del functor $A\mapsto\Hom(A\times B,C)$, i escriure les dues lleis de la currificació;
  \item reconèixer una categoria cartesiana tancada i l'exponencial com a adjunt per la dreta de $-\times B$;
  \item caracteritzar els ordres parcials cartesians tancats com a semireticles implicatius amb màxim, i relacionar-los amb les àlgebres de Heyting;
  \item situar $\cat{Set}$, els preordres, les categories de prefeixos, $\cat{Graph}$ i $\cat{Top}$ respecte del tancament cartesià;
  \item interpretar el càlcul lambda simplement tipat en una CCC, incloses les lleis $\beta$ i $\eta$ i les regles estructurals.
\end{itemize}
\end{objectius}

En tractar les adjuncions vam trobar, al final del capítol, l'adjunció entre el producte i l'exponencial a $\cat{Cat}$: $-\times\cat{B}\dashv(-)^{\cat{B}}$. Aquella era una manifestació externa d'un fenomen que ara volem fer \emph{intern} a una categoria qualsevol. Una categoria cartesiana en què, per a cada objecte $B$, el functor $-\times B$ té adjunt per la dreta s'anomena \textbf{cartesiana tancada}, i és, exactament, el marc on es pot interpretar el \textbf{càlcul lambda simplement tipat}. Aquest capítol construeix aquesta noció pas a pas, la il·lustra en els exemples fonamentals i n'explica la lectura com a semàntica de tipus i termes. És el primer pont explícit del llibre cap a la lògica categòrica.

Al capítol anterior, les adjuncions van unificar construccions aparentment diverses: objectes lliures, interiors i clausures, connectius lògics i exponencials. El patró comú, una bijecció natural d'homconjunts o, equivalentment, una unitat i una counitat que satisfan les identitats triangulars, reapareixerà com a columna vertebral de tot el que segueix. En aquest capítol el trobem en l'adjunció entre el producte i l'exponencial; més endavant, en les categories amb l'estructura suficient, veurem que els \textbf{quantificadors} $\exists$ i $\forall$ es modelen, respectivament, com els adjunts per l'esquerra i per la dreta dels functors de reindexació entre subobjectes: és la formulació categòrica de la relació entre quantificació i substitució.

\section{Categories cartesianes}

Recordem que un \textbf{producte finit} és el límit d'un diagrama amb un nombre finit d'objectes i cap morfisme no trivial. El cas de zero objectes és l'objecte terminal; el de dos, el producte binari. Una categoria té \textbf{tots els productes finits} si té objecte terminal i producte binari de qualsevol parella, ja que els productes de tres o més objectes s'obtenen iterant.

\begin{definicio}[Categoria cartesiana]\index{categoria!cartesiana}
\label{def:cartesiana}
Una categoria $\cat{C}$ és \textbf{cartesiana} si té objecte terminal $1$ i producte binari $A\times B$ de qualsevol parella d'objectes. Equivalentment, si té tots els productes finits.
\end{definicio}

En una categoria cartesiana disposem, per a cada objecte $A$, del \textbf{morfisme diagonal}
\[
\delta_A=\langle\id_A,\id_A\rangle\colon A\to A\times A,
\]
l'únic morfisme les dues projeccions del qual són la identitat, i del morfisme únic $!_A\colon A\to 1$. Escrivim la diagonal amb minúscula per no confondre-la amb el functor constant $\Delta_A$ (exemple~\ref{ex:functor-constant}). Aquestes dades no són estructura addicional: queden \emph{determinades} per les propietats universals.

Els productes escollits determinen també la functorialitat del producte. Donats $f\colon A\to A'$ i $g\colon B\to B'$, posem
\[
f\times g=\langle f\comp\pi_1,\ g\comp\pi_2\rangle\colon A\times B\to A'\times B'.
\]
En particular, per a $B$ fix, l'assignació $A\mapsto A\times B$ defineix el functor $-\times B$, tal com precisem ara.

\begin{proposicio}[El producte per un objecte és functorial]\index{producte!functorialitat}
\label{prop:producte-functor}
Sigui $\cat{C}$ cartesiana i $B$ un objecte fix. L'assignació $A\mapsto A\times B$ s'estén a un functor
\[
-\times B\colon\cat{C}\to\cat{C},
\]
que envia $f\colon A\to A'$ al morfisme $f\times\id_B=\langle f\comp\pi_1,\pi_2\rangle\colon A\times B\to A'\times B$.
\end{proposicio}

\begin{prova}
Cal comprovar que $-\times B$ respecta identitats i composicions. Per a la identitat, $\id_A\times\id_B=\langle\pi_1,\pi_2\rangle=\id_{A\times B}$ per la unicitat de la propietat universal del producte. Per a la composició, siguin $f\colon A\to A'$ i $g\colon A'\to A''$. Tant $(g\comp f)\times\id_B$ com $(g\times\id_B)\comp(f\times\id_B)$ tenen com a projecció primera $g\comp f\comp\pi_1$ i com a projecció segona $\pi_2$; per unicitat, coincideixen.
\end{prova}

\section{Objectes exponencials}

En una categoria cartesiana, la pregunta clau és si el functor $-\times B$ té adjunt per la dreta. Per a $B$ i $C$ fixos, l'exponencial $C^B$, quan existeix, representa el functor
\[
A\longmapsto\Hom(A\times B,C)
\]
(capítol~\ref{cap:yoneda}); quan existeix per a tot $C$, defineix l'adjunt per la dreta de $-\times B$.

\begin{definicio}[Objecte exponencial]\index{objecte!exponencial}\index{exponencial}
\label{def:exponencial}
Siguin $B,C$ objectes d'una categoria cartesiana $\cat{C}$. Un \textbf{objecte exponencial} $C^{B}$ és un objecte dotat d'un morfisme d'\textbf{avaluació}
\[
\ev\colon C^{B}\times B\to C
\]
amb la propietat universal següent: per a tot objecte $A$ i tot morfisme $f\colon A\times B\to C$, existeix un únic morfisme $\lambda f\colon A\to C^{B}$, la \textbf{transposada} o \textbf{currificació} de $f$, tal que el triangle
\[
\begin{tikzpicture}[baseline=(m.center)]
  \matrix (m) [matrix of math nodes, column sep=3.4em, row sep=2.6em] {
    A\times B & C^{B}\times B \\
              & C \\
  };
  \draw[-{Stealth[scale=1.1]}] (m-1-1) -- node[above] {$\lambda f\times\id_B$} (m-1-2);
  \draw[-{Stealth[scale=1.1]}] (m-1-1) -- node[below left] {$f$} (m-2-2);
  \draw[-{Stealth[scale=1.1]}] (m-1-2) -- node[right] {$\ev$} (m-2-2);
\end{tikzpicture}
\]
commuta, és a dir, $\ev\comp(\lambda f\times\id_B)=f$.
\end{definicio}

L'avaluació formalitza la idea d'\enquote{aplicar una funció al seu argument}, i la currificació, la de \enquote{fixar un paràmetre}. Com tota representació, l'exponencial és únic llevat d'un únic isomorfisme compatible amb l'avaluació: si $(E,\ev)$ i $(E',\ev')$ són exponencials de $C$ per $B$, hi ha un únic isomorfisme $u\colon E\to E'$ amb $\ev'\comp(u\times\id_B)=\ev$.

\begin{observacio}[Les dues lleis de la currificació]\label{obs:lleis-currificacio}\index{currificació!lleis}
La correspondència $f\mapsto\lambda f$ és una bijecció entre $\Hom(A\times B,C)$ i $\Hom(A,C^B)$, amb inversa $g\mapsto\ev\comp(g\times\id_B)$. Les dues composicions donen identitats:
\[
\ev\comp(\lambda f\times\id_B)=f,
\qquad\qquad
\lambda\bigl(\ev\comp(g\times\id_B)\bigr)=g,
\]
per a tot $f\colon A\times B\to C$ i tot $g\colon A\to C^B$. La primera és la condició de la definició. La segona en surt per unicitat: $g$ és un morfisme $A\to C^B$ que fa commutar el triangle per a $f=\ev\comp(g\times\id_B)$, i l'únic que ho fa és $\lambda f$. Aquestes dues lleis són la versió categòrica de les lleis $\beta$ i $\eta$ del càlcul lambda, que veurem més endavant; la naturalitat de la bijecció la converteix en una adjunció.
\end{observacio}

\begin{proposicio}[L'exponencial és un adjunt per la dreta]\index{adjunció!producte-exponencial}
\label{prop:adjuncio-exp}
Si per a un objecte fix $B$ existeix l'exponencial $C^{B}$ per a tot $C$, aleshores l'assignació $C\mapsto C^{B}$ s'estén a un functor $(-)^{B}\colon\cat{C}\to\cat{C}$ i hi ha un isomorfisme natural
\[
\Hom(A\times B,\,C)\;\cong\;\Hom(A,\,C^{B}),
\]
natural en $A$ i en $C$. És a dir, $-\times B\dashv(-)^{B}$.
\end{proposicio}

\begin{prova}
La currificació $f\mapsto\lambda f$ és la bijecció buscada (observació~\ref{obs:lleis-currificacio}). La naturalitat en $A$ es comprova precomponent amb $a\colon A'\to A$ i usant la functorialitat de $-\times B$ (la proposició~\ref{prop:producte-functor}); la naturalitat en $C$ defineix, de fet, l'acció de $(-)^{B}$ sobre morfismes: donat $c\colon C\to C'$, es posa $c^{B}=\lambda\bigl(c\comp\ev\bigr)\colon C^{B}\to C'^{B}$. Que això respecta identitats i composicions se segueix, un cop més, de la unicitat de la transposada.
\end{prova}

\begin{notacio}
Escrivim indistintament $C^{B}$ o $[B,C]$ per a l'exponencial. La transposada de $f$ es denota $\lambda f$ o $\widehat f$; la seva inversa, $\ev\comp(g\times\id_B)$, es denota $\widetilde g$ i s'anomena \textbf{descurrificació} de $g$.
\end{notacio}

\section{La definició de categoria cartesiana tancada}

\begin{definicio}[Categoria cartesiana tancada]\index{categoria!cartesiana tancada}\index{CCC}
\label{def:ccc}
Una categoria $\cat{C}$ és \textbf{cartesiana tancada} (abreujat \textbf{CCC}, de l'anglès \emph{cartesian closed category}) si és cartesiana i, a més, per a cada objecte $B$, el functor $-\times B\colon\cat{C}\to\cat{C}$ té adjunt per la dreta $(-)^{B}$. Explícitament, $\cat{C}$ té:
\begin{enumerate}
\item un objecte terminal $1$;
\item productes binaris $A\times B$;
\item exponencials $C^{B}$ amb avaluació i currificació.
\end{enumerate}
\end{definicio}

\begin{observacio}[L'exponencial com a bifunctor]\label{obs:exponencial-bifunctor}\index{exponencial!bifunctor}
En una CCC, els exponencials s'organitzen en un bifunctor
\[
(-)^{(-)}\colon\cat{C}\op\times\cat{C}\longrightarrow\cat{C},
\qquad (B,C)\longmapsto C^{B},
\]
contravariant en $B$ i covariant en $C$, i la bijecció
\[
\Hom(A\times B,C)\;\cong\;\Hom(A,C^{B})
\]
és natural en les tres variables $A$, $B$ i $C$. L'acció en $C$ és la de la proposició~\ref{prop:adjuncio-exp}; la construcció de l'acció en $B$, mitjançant la transposada de $\ev\comp(\id\times f)$, es proposa als Exercicis.
\end{observacio}

La condició es pot resumir dient que en una CCC les tres operacions \enquote{no fer res} (terminal), \enquote{aparellar} (producte) i \enquote{formar espais de funcions} (exponencial) hi són totes i estan lligades per adjuncions. Aquesta és exactament l'estructura que necessita un llenguatge de tipus amb producte de tipus i tipus funció.

\begin{observacio}
La tancadura és una propietat \emph{objecte a objecte}: cal l'exponencial per a \emph{cada} base $B$. Una categoria pot tenir alguns exponencials i no ser tancada. La demanda uniforme ---un adjunt per la dreta de $-\times B$ per a tot $B$--- és el que fa que la currificació estigui sempre disponible.
\end{observacio}

\section{Exemples fonamentals}

\begin{exemple}[$\cat{Set}$]\index{Set@$\cat{Set}$!cartesiana tancada}
La categoria $\cat{Set}$ és cartesiana tancada. L'objecte terminal és qualsevol conjunt d'un element; el producte és el producte cartesià; i l'exponencial $C^{B}$ és el conjunt de \emph{totes} les aplicacions $B\to C$, amb l'avaluació $\ev(g,b)=g(b)$. La currificació $\lambda f\colon A\to C^{B}$ d'una aplicació $f\colon A\times B\to C$ és $\lambda f(a)=(b\mapsto f(a,b))$: exactament la currificació de la programació funcional. Aquest és el model prototípic i el que dona nom a tota la construcció.
\end{exemple}

\begin{exemple}[$\cat{Cat}$: transformacions naturals com a functors]\label{ex:cat-ccc}\index{Cat@$\cat{Cat}$!cartesiana tancada}\index{avaluació!a Cat@a $\cat{Cat}$}
La categoria $\cat{Cat}$ de les categories petites és cartesiana tancada: el terminal és $\cat{1}$, el producte és la categoria producte i l'exponencial de $\cat{C}$ per $\cat{B}$ és la categoria de functors $\cat{C}^{\cat{B}}$ (proposició~\ref{prop:cat-currificacio}). Fem explícits l'avaluació i la currificació, i després els llegim en un cas concret.

\emph{Avaluació.} El functor $\ev\colon\cat{C}^{\cat{B}}\times\cat{B}\to\cat{C}$ és
\[
\ev(H,b)=H(b),\qquad \ev(\eta,g)=\eta_{b'}\comp H(g)=H'(g)\comp\eta_b,
\]
per a $\eta\colon H\Rightarrow H'$ i $g\colon b\to b'$; la segona igualtat és la naturalitat de $\eta$. És la descurrificació de $\id_{\cat{C}^{\cat{B}}}$ a la demostració de la proposició~\ref{prop:cat-currificacio}, i d'allà en surt que és un functor.

\emph{Currificació i llei $\beta$.} La currificació de $F\colon\cat{A}\times\cat{B}\to\cat{C}$ és el functor $\Lambda F\colon\cat{A}\to\cat{C}^{\cat{B}}$ d'aquella demostració. El triangle de la definició~\ref{def:exponencial} commuta: sobre objectes, $\ev\bigl(\Lambda F(a),b\bigr)=F(a,b)$; sobre un morfisme $(f,g)$,
\[
\ev\bigl(\Lambda F(f),g\bigr)=\Lambda F(f)_{b'}\comp\Lambda F(a)(g)=F(f,\id_{b'})\comp F(\id_a,g)=F(f,g).
\]
Així, $\ev\comp(\Lambda F\times\id_{\cat{B}})=F$: avaluar la currificació torna la funció de dues variables.

\emph{Un cas concret: $\cat{B}=\cat{2}$.} La categoria $\cat{C}^{\cat{2}}$ és la categoria de fletxes de $\cat{C}$: els objectes són els morfismes de $\cat{C}$ i els morfismes, els quadrats commutatius (Pràctica del capítol~\ref{cap:transformacions}). L'avaluació envia $(h,0)$ al domini de $h$, $(h,1)$ al codomini i $(h,u)$ a $h$ mateix, on $u\colon0\to1$. Un functor $\Phi\colon\cat{A}\times\cat{2}\to\cat{C}$ consisteix en tres dades:
\begin{itemize}
\item dos functors $F=\Phi(-,0)$ i $G=\Phi(-,1)$ de $\cat{A}$ a $\cat{C}$;
\item una família $\alpha_a=\Phi(\id_a,u)\colon F(a)\to G(a)$;
\item la naturalitat $G(f)\comp\alpha_a=\alpha_{a'}\comp F(f)$, que no és cap condició nova: és $\Phi$ aplicat a les dues descomposicions de $(f,u)$.
\end{itemize}
És a dir, un functor $\cat{A}\times\cat{2}\to\cat{C}$ \emph{és} una transformació natural $\alpha\colon F\Rightarrow G$. La seva currificació $\Lambda\Phi\colon\cat{A}\to\cat{C}^{\cat{2}}$ envia $a$ a la fletxa $\alpha_a$ i $f$ al quadrat de naturalitat de $f$. L'avaluació recupera les tres peces: $\ev(\Lambda\Phi(a),0)=F(a)$, $\ev(\Lambda\Phi(a),1)=G(a)$ i $\ev(\Lambda\Phi(a),u)=\alpha_a$. Tenim, doncs, tres lectures d'una mateixa dada,
\[
\alpha\colon F\Rightarrow G\qquad\longleftrightarrow\qquad\cat{A}\times\cat{2}\to\cat{C}\qquad\longleftrightarrow\qquad\cat{A}\to\cat{C}^{\cat{2}},
\]
i el pas de la segona a la tercera és la currificació, amb l'avaluació com a inversa.
\end{exemple}

\begin{exemple}[Preordres cartesians tancats]\index{semireticle implicatiu}\index{àlgebra de Heyting}\index{Heyting}
Un preordre amb ínfims finits, vist com a categoria, és cartesià: el terminal és l'element màxim $\top$ i el producte és l'ínfim $x\wedge y$. Té exponencials si i només si, per a cada $b$, l'operació $-\wedge b$ té adjunt per la dreta; aquest adjunt és la \textbf{implicació} $b\Rightarrow-$, caracteritzada per
\[
a\wedge b\leq c\quad\Longleftrightarrow\quad a\leq(b\Rightarrow c).
\]
Un ordre parcial amb ínfims finits és, doncs, cartesià, i és cartesià tancat si i només si és un \textbf{semireticle implicatiu amb màxim}: un $\wedge$-semireticle amb element màxim $\top$ i una operació d'implicació que satisfà l'equivalència anterior (també anomenat \emph{semireticle de Brouwer}). Per a un preordre, la mateixa caracterització val després de passar al quocient per l'equivalència
\[
a\sim b\quad\Longleftrightarrow\quad a\leq b\ \text{i}\ b\leq a,
\]
perquè en la terminologia algebraica habitual un semireticle és un ordre parcial. Aquesta estructura captura la part $(\top,\wedge,\Rightarrow)$ de la lògica.

Convé no confondre aquesta estructura amb una àlgebra de Heyting. Una \textbf{àlgebra de Heyting}\index{àlgebra de Heyting} és un \emph{reticle distributiu acotat} ---amb ínfims finits $\top,\wedge$ i \emph{suprems finits} $\bot,\vee$--- proveït d'una implicació que satisfà l'equivalència anterior. El tancament cartesià per si sol només dona la part multiplicativa: no garanteix ni els joins $\vee$ ni el mínim $\bot$. Reservem, doncs, el nom \emph{àlgebra de Heyting} per al cas amb suprems finits; un semireticle implicatiu amb màxim que a més tingui suprems finits (i sigui distributiu, cosa automàtica en presència de la implicació) és una àlgebra de Heyting. Aquest exemple és decisiu per a la lògica: la implicació intuïcionista \emph{és} l'exponencial, i el tancament cartesià és la traducció algebraica de disposar del connectiu $\to$, mentre que la disjunció $\vee$ correspon als suprems.
\end{exemple}

\begin{observacio}[Atenció: cartesià tancat no vol dir àlgebra de Heyting]\label{obs:ccc-no-heyting}\index{àlgebra de Heyting!i preordres cartesians tancats}
En un preordre, ser cartesià tancat dona $\top$, $\wedge$ i $\Rightarrow$, i res més. Un exemple mínim mostra què pot faltar. Sigui $\mathbb Z\cup\{+\infty\}$ amb l'ordre usual. És una cadena amb màxim $+\infty$, amb ínfims $a\wedge b=\min(a,b)$, i té implicació
\[
b\Rightarrow c=\begin{cases}+\infty&\text{si } b\leq c,\\ c&\text{si } b>c.\end{cases}
\]
En efecte, si $b\leq c$ tenim $a\wedge b\leq c$ per a tot $a$; si $b>c$, tenim $\min(a,b)\leq c$ si i només si $a\leq c$. Per tant és un ordre parcial cartesià tancat. Però no té element mínim, de manera que no és una àlgebra de Heyting. Aquí no existeixen ni $\bot$, ni la negació $\neg b=(b\Rightarrow\bot)$, ni el principi \emph{ex falso}; no hi ha, doncs, cap interpretació de la falsedat. En altres exemples pot faltar el suprem binari $\vee$, i amb ell la disjunció.

La relació exacta és aquesta: una àlgebra de Heyting és un ordre parcial cartesià tancat que, a més, té coproductes finits ($\bot$ i $\vee$). La distributivitat ve aleshores gratis, perquè $-\wedge b$, com a adjunt per l'esquerra, preserva els suprems (capítol~\ref{cap:adjuncions}). L'estructura de Heyting també condiciona els morfismes: un morfisme d'àlgebres de Heyting ha de preservar $\bot$ i $\vee$, cosa que no demana un functor que només preservi l'estructura cartesiana tancada. En lògica, el tancament cartesià modela el fragment $(\top,\wedge,\to)$; la falsedat i la disjunció són estructura addicional.
\end{observacio}

\begin{exemple}[Categories de prefeixos]\index{prefeix@prefeix!cartesiana tancada}
Per a qualsevol categoria petita $\cat{C}$, la categoria de prefeixos $\widehat{\cat{C}}=\cat{Set}^{\cat{C}\op}$ (exemple~\ref{ex:categoria-prefeixos}) és cartesiana tancada. Els productes i el terminal es calculen \emph{objecte a objecte} a $\cat{Set}$ (proposició~\ref{prop:limits-punt-a-punt}), i l'exponencial es defineix, mitjançant el lema de Yoneda, per $Q^{P}(A)=\Nat(\Hom(-,A)\times P,\,Q)$. Aquestes categories seran el terreny natural dels topos de prefeixos.
\end{exemple}

\begin{exemple}[Grafs]\index{graf dirigit!cartesiana tancada}
Al capítol~\ref{cap:transformacions} vam veure que $\cat{Graph}\cong\cat{Set}^{\cat{P}}$, on $\cat{P}$ té dos objectes $E,V$ i dues fletxes $s,t\colon E\to V$. Com que $\cat{Set}^{\cat{P}}=\cat{Set}^{(\cat{P}\op)\op}$, els grafs són exactament els prefeixos sobre $\cat{P}\op$: un prefeix $G\colon(\cat{P}\op)\op\to\cat{Set}$ és la dada de dos conjunts, de vèrtexs i d'arestes, i de les aplicacions origen i destí. Per tant, com a cas particular de l'exemple anterior, $\cat{Graph}$ és cartesiana tancada. No calcularem aquí l'exponencial de dos grafs; n'hi ha prou de saber que existeix i que es pot obtenir amb la fórmula de Yoneda de l'exemple anterior.
\end{exemple}

\begin{observacio}[Un no-exemple instructiu]
La categoria $\cat{Top}$ dels espais topològics \emph{no} és cartesiana tancada: no sempre existeix una topologia sobre l'espai de funcions contínues que faci de l'avaluació un morfisme universal. Aquesta mancança va motivar històricament la introducció de categories \enquote{convenients} d'espais topològics, com la dels espais Hausdorff compactament generats, amb les construccions categòriques adequades, que sí que és cartesiana tancada. El tancament cartesià és, doncs, una propietat delicada i valuosa, no automàtica.
\end{observacio}

\section[Semàntica del càlcul lambda simplement tipat]{Semàntica del càlcul lambda\\ simplement tipat}

El motiu pel qual les CCC són centrals en lògica categòrica és la seva correspondència amb el \textbf{càlcul lambda simplement tipat}. Descrivim la traducció a nivell introductori; és una de les formes de l'anomenada \emph{correspondència de Curry--Howard--Lambek}.

\subsection{Tipus com a objectes}

Fixem un conjunt de tipus base. Els \textbf{tipus} es construeixen amb dues operacions: el \textbf{producte de tipus} $\sigma\times\tau$ i el \textbf{tipus funció} $\sigma\to\tau$, més un tipus unitat $\mathbf{1}$. En una CCC $\cat{C}$, la interpretació $\llbracket-\rrbracket$ assigna:
\[
\llbracket\mathbf{1}\rrbracket=1,\qquad
\llbracket\sigma\times\tau\rrbracket=\llbracket\sigma\rrbracket\times\llbracket\tau\rrbracket,\qquad
\llbracket\sigma\to\tau\rrbracket=\llbracket\tau\rrbracket^{\llbracket\sigma\rrbracket}.
\]
Cada tipus esdevé un objecte; el tipus funció esdevé un exponencial.

\subsection{Termes com a morfismes}

Un \textbf{context} $\Gamma=x_1:\sigma_1,\dots,x_n:\sigma_n$ s'interpreta, per recursió, com un producte:
\[
\llbracket\varnothing\rrbracket=1,
\qquad
\llbracket\Gamma,x:\sigma\rrbracket=\llbracket\Gamma\rrbracket\times\llbracket\sigma\rrbracket.
\]
Un \textbf{terme} $\Gamma\vdash t:\tau$, amb variables lliures en el context $\Gamma$, s'interpreta com un morfisme
\[
\llbracket t\rrbracket\colon\llbracket\Gamma\rrbracket\longrightarrow\llbracket\tau\rrbracket.
\]
Les regles de formació de termes es corresponen amb les operacions de la CCC:
\begin{itemize}
\item una \textbf{variable} $x_i$ s'interpreta amb la projecció $\pi_i$;
\item una \textbf{parella} $\langle t,u\rangle$, amb el morfisme aparellament $\langle\llbracket t\rrbracket,\llbracket u\rrbracket\rangle$ cap a un producte;
\item una \textbf{aplicació} $t\,u$, amb la composició de l'avaluació i l'aparellament de $\llbracket t\rrbracket$ i $\llbracket u\rrbracket$;
\item una \textbf{abstracció} $\lambda x{:}\sigma.\,t$, amb la currificació $\lambda\llbracket t\rrbracket$: si $\Gamma,x:\sigma\vdash t:\tau$, aleshores $\llbracket t\rrbracket\colon\llbracket\Gamma\rrbracket\times\llbracket\sigma\rrbracket\to\llbracket\tau\rrbracket$, i la seva transposada és un morfisme $\llbracket\Gamma\rrbracket\to\llbracket\tau\rrbracket^{\llbracket\sigma\rrbracket}$, exactament el tipus de $\lambda x{:}\sigma.\,t$.
\end{itemize}
La correspondència entre abstracció i currificació és el cor de la traducció: lligar una variable és transposar per l'adjunció, i aplicar una funció és avaluar. Les regles de conversió del càlcul es tornen igualtats de morfismes:
\[
(\beta)\quad(\lambda x.\,t)\,u = t[u/x]
\qquad\text{esdevé}\qquad
\ev\comp\langle\lambda\llbracket t\rrbracket,\llbracket u\rrbracket\rangle=\llbracket t\rrbracket\comp\langle\id,\llbracket u\rrbracket\rangle,
\]
que no és literalment l'equació de la definició~\ref{def:exponencial}, sinó que \emph{se'n dedueix}. Si $f=\llbracket t\rrbracket\colon\llbracket\Gamma\rrbracket\times\llbracket\sigma\rrbracket\to\llbracket\tau\rrbracket$, la primera llei de la currificació (observació~\ref{obs:lleis-currificacio}) dona $\ev\comp(\lambda f\times\id)=f$, i precomponent amb $\langle\id_{\llbracket\Gamma\rrbracket},\llbracket u\rrbracket\rangle\colon\llbracket\Gamma\rrbracket\to\llbracket\Gamma\rrbracket\times\llbracket\sigma\rrbracket$ s'obté
\[
\ev\comp\langle\lambda f,\llbracket u\rrbracket\rangle
=\ev\comp(\lambda f\times\id)\comp\langle\id,\llbracket u\rrbracket\rangle
=f\comp\langle\id,\llbracket u\rrbracket\rangle,
\]
que és la interpretació de la substitució $t[u/x]$. La regla $(\eta)$, $\lambda x.\,(t\,x)=t$, es correspon amb la segona llei, és a dir, amb la unicitat de la transposada.

\begin{observacio}[L'estructura cartesiana i les regles estructurals]\index{regles estructurals}
Els morfismes canònics del producte interpreten les \emph{regles estructurals} de la lògica i del càlcul. Una projecció $\llbracket\Gamma\rrbracket\times\llbracket\sigma\rrbracket\to\llbracket\Gamma\rrbracket$ permet \emph{oblidar} una variable (debilitament); la diagonal $\delta\colon\llbracket\sigma\rrbracket\to\llbracket\sigma\rrbracket\times\llbracket\sigma\rrbracket$ permet \emph{reutilitzar-la} (contracció); i l'isomorfisme de simetria $\llbracket\sigma\rrbracket\times\llbracket\tau\rrbracket\cong\llbracket\tau\rrbracket\times\llbracket\sigma\rrbracket$ permet \emph{intercanviar-ne} l'ordre (intercanvi). És per això que la diagonal, introduïda al començament del capítol, no és un detall aïllat: l'estructura \emph{cartesiana} és exactament la que fa admissibles aquestes tres regles.
\end{observacio}

\begin{observacio}[Curry--Howard--Lambek]\index{Curry--Howard--Lambek}
La correspondència es pot fer precisa: després de fixar una noció adequada de teoria del càlcul lambda simplement tipat, una estructura cartesiana tancada escollida i els morfismes que la preserven, s'obté una \emph{equivalència} de categories entre les dues presentacions. En aquest sentit, tipus, termes i equacions del càlcul lambda corresponen a objectes, morfismes i igualtats de morfismes en una CCC. Aquesta és la primera de les grans traduccions \enquote{lògica $=$ categoria} que estructuraran l'estudi posterior de la lògica categòrica; les següents afegiran, sobre aquesta base, els subobjectes per als predicats i la reindexació per a la quantificació.
\end{observacio}

\begin{resumcap}
\apartat{Idees centrals}
\begin{enumerate}
  \item L'objecte exponencial $C^{B}$ representa el functor $A\mapsto\Hom(A\times B,C)$; la currificació és la bijecció corresponent, amb les dues lleis $\ev\comp(\lambda f\times\id_B)=f$ i $\lambda(\ev\comp(g\times\id_B))=g$.
  \item Una categoria cartesiana tancada té productes finits i tots els exponencials; els exponencials s'organitzen en un bifunctor contravariant en l'exponent i covariant en la base.
  \item Un ordre parcial cartesià tancat és un semireticle implicatiu amb màxim; per a un preordre, aquesta descripció val després de passar al quocient per l'equivalència $a\sim b$. Si, a més, hi ha suprems finits, s'obté una àlgebra de Heyting; sense ells, no (per exemple, $\mathbb Z\cup\{+\infty\}$ no té $\bot$).
  \item $\cat{Set}$, $\cat{Cat}$, les categories de prefeixos i, en particular, $\cat{Graph}$ són cartesianes tancades. A $\cat{Cat}$, l'avaluació i la currificació identifiquen les transformacions naturals amb els functors $\cat{A}\times\cat{2}\to\cat{C}$ i $\cat{A}\to\cat{C}^{\cat{2}}$; $\cat{Top}$ no ho és en general.
  \item El càlcul lambda simplement tipat s'interpreta en tota CCC: els contextos com a productes, l'abstracció com a currificació, i les lleis $\beta$ i $\eta$ com a conseqüències de les dues lleis de la currificació.
\end{enumerate}
\apartat{Errors típics} Identificar tota categoria cartesiana tancada amb una àlgebra de Heyting, fins i tot quan és un ordre parcial (poden faltar $\bot$ i els joins, i amb ells la negació i la disjunció); atribuir la functorialitat del producte a la diagonal en lloc de a la seva propietat universal; i tractar la correspondència de Curry--Howard--Lambek com una igualtat en lloc d'una equivalència amb estructura triada.
\apartat{Lectura recomanada} \cite[cap.~6]{awodey} per a exponencials i CCC; \cite{lambekscott} per a la correspondència amb el càlcul lambda; \cite[cap.~1]{maclanemoerdijk} per als exemples de prefeixos; \cite[cap.~6]{barrwells}, que presenta les CCC i el càlcul lambda tipat i calcula explícitament l'exponencial de dos grafs (exemple~6.1.12); i \cite[§4.3]{riehl} per a la formulació de l'exponencial com a adjunt; \cite[§11]{streicher}, per als exponencials en categories de prefeixos (§11.1) i la semàntica del càlcul lambda tipat (§11.2); i \cite[article~V, sessions~30--31]{lawvereschanuel}, per a una introducció molt elemental als objectes de morfismes.
\end{resumcap}

\chapter[Subobjectes, imatges i factoritzacions]{\texorpdfstring{Subobjectes, imatges\\ i factoritzacions}{Subobjectes, imatges i factoritzacions}}
\label{cap:subobjectes}

\begin{objectius}
\apartat{Prerequisits} Monomorfismes i epimorfismes (cap.~\ref{cap:categories}); pullbacks, coequalitzadors, límits finits i reindexació per pullback (cap.~\ref{cap:limits}); adjuncions, especialment entre preordres (cap.~\ref{cap:adjuncions}).
\apartat{Objectius} En acabar el capítol, el lector hauria de poder:
\begin{itemize}
  \item definir els subobjectes i entendre $\Sub(A)$ com un ordre parcial de predicats sobre $A$;
  \item interpretar la reindexació $f^{*}$ com a substitució i organitzar-la, en una categoria ben potenciada, en un functor contravariant $\Sub\colon\cat{C}\op\to\cat{Poset}$;
  \item interpretar la intersecció de subobjectes com a pullback;
  \item definir epimorfisme regular, parella nucli i factorització imatge;
  \item reconèixer una categoria regular;
  \item demostrar l'adjunció $\exists_f\dashv f^{*}$.
\end{itemize}
\end{objectius}

El capítol anterior ens ha donat la semàntica dels tipus i els termes. Ara construïm la peça complementària: la interpretació de les \textbf{proposicions} i la \textbf{substitució}. La idea directriu és senzilla i profunda: un predicat sobre un objecte $A$ ---una propietat que els \enquote{elements} de $A$ poden tenir o no--- es representa categòricament per un \textbf{subobjecte} de $A$. A $\cat{Set}$, un subconjunt es representa per la seva inclusió, que és injectiva; en una categoria arbitrària, el paper de les injeccions el fan els monomorfismes. Els subobjectes d'un objecte formen un ordre parcial $\Sub(A)$, i la reindexació per pullback, que ja vam construir, hi actua com la substitució.

A la segona meitat del capítol construïm la \textbf{factorització imatge} d'un morfisme. A $\cat{Set}$, tota aplicació es factoritza com una aplicació exhaustiva sobre la seva imatge seguida de la inclusió d'aquesta imatge; en una categoria arbitrària, aquesta idea es reformula mitjançant epimorfismes regulars i monomorfismes. Després definim les \textbf{categories regulars} com aquelles on aquesta factorització és estable per canvi de base, i n'obtenim una conseqüència notable: l'\textbf{existencial} apareix com a \textbf{adjunt per l'esquerra} de la reindexació. L'universal $\forall_f$ requereix estructura addicional i es tractarà al volum de lògica categòrica. Amb això, el volum tanca la infraestructura categòrica que la lògica farà servir.

\section{Subobjectes}

Volem capturar la idea de \enquote{subconjunt} de manera purament categòrica. A $\cat{Set}$, un subconjunt $S\subseteq A$ es pot presentar per la seva injecció $S\hookrightarrow A$, que és un monomorfisme. Però dos monomorfismes diferents poden tenir la mateixa imatge: una injecció i una reparametrització bijectiva del domini defineixen \enquote{el mateix} subconjunt. Cal, doncs, identificar monomorfismes que difereixen per un isomorfisme del domini.

\begin{definicio}[Subobjecte]\index{subobjecte}
\label{def:subobjecte}
Sigui $A$ un objecte d'una categoria $\cat{C}$. Considerem els monomorfismes amb codomini $A$. Diem que dos monomorfismes $m\colon S\rightarrowtail A$ i $m'\colon S'\rightarrowtail A$ són \textbf{equivalents}, i escrivim $m\sim m'$, si existeix un isomorfisme $\varphi\colon S\to S'$ amb $m'\comp\varphi=m$:
\[
\begin{tikzpicture}[baseline=(m.center)]
  \matrix (m) [matrix of math nodes, column sep=2.8em, row sep=2.0em] {
    S & & S' \\
      & A & \\
  };
  \draw[-{Stealth[scale=1.1]}] (m-1-1) -- node[above] {$\varphi$} node[below] {$\sim$} (m-1-3);
  \draw[-{Stealth[scale=1.1]}] (m-1-1) -- node[below left] {$m$} (m-2-2);
  \draw[-{Stealth[scale=1.1]}] (m-1-3) -- node[below right] {$m'$} (m-2-2);
\end{tikzpicture}
\]
Un \textbf{subobjecte} de $A$ és una classe d'equivalència de monomorfismes per aquesta relació. Escrivim $[m]$ per al subobjecte representat per $m$.
\end{definicio}

\begin{observacio}
La relació $\sim$ és d'equivalència: la reflexivitat usa $\varphi=\id$, la simetria usa $\varphi^{-1}$ (que existeix perquè $\varphi$ és isomorfisme), i la transitivitat, la composició d'isomorfismes. Cal notar que l'isomorfisme $\varphi$, si existeix, és \emph{únic}: de $m'\comp\varphi=m=m'\comp\varphi'$ i el fet que $m'$ és mono se segueix $\varphi=\varphi'$. Els subobjectes estan, doncs, ben definits sense ambigüitat.
\end{observacio}

\begin{exemple}
A $\cat{Set}$, els subobjectes de $A$ es corresponen exactament amb els \textbf{subconjunts} de $A$. Un monomorfisme $m\colon S\rightarrowtail A$ és una aplicació injectiva, i la seva classe d'equivalència queda determinada per la imatge $m(S)\subseteq A$: dos monos són equivalents si i només si tenen la mateixa imatge. Així, $\Sub(A)\cong\mathcal{P}(A)$, el conjunt de parts de $A$.
\end{exemple}

\section{L'ordre de subobjectes}

Els subobjectes d'un objecte no formen només un conjunt: hi ha una noció natural d'inclusió entre ells.

\begin{definicio}[Ordre entre subobjectes]\index{subobjecte!ordre}\index{Sub@$\Sub(A)$}
\label{def:ordre-sub}
Donats dos monomorfismes $m\colon S\rightarrowtail A$ i $n\colon T\rightarrowtail A$, diem que $[m]\leq[n]$ si $m$ \textbf{es factoritza a través de} $n$, és a dir, si existeix un morfisme $h\colon S\to T$ amb $n\comp h=m$:
\[
\begin{tikzpicture}[baseline=(m.center)]
  \matrix (m) [matrix of math nodes, column sep=2.8em, row sep=2.0em] {
    S & & T \\
      & A & \\
  };
  \draw[-{Stealth[scale=1.1]}] (m-1-1) -- node[above] {$h$} (m-1-3);
  \draw[-{Stealth[scale=1.1]}] (m-1-1) -- node[below left] {$m$} (m-2-2);
  \draw[-{Stealth[scale=1.1]}] (m-1-3) -- node[below right] {$n$} (m-2-2);
\end{tikzpicture}
\]
\end{definicio}

\begin{proposicio}[$\Sub(A)$ és un ordre parcial]
\label{prop:sub-preordre}
La relació $\leq$ de la definició~\ref{def:ordre-sub} està ben definida sobre classes d'equivalència i és reflexiva i transitiva. A més, és antisimètrica: $[m]\leq[n]$ i $[n]\leq[m]$ impliquen $[m]=[n]$. Per tant, $\Sub(A)$ és un \textbf{ordre parcial}. (Sobre els monomorfismes mateixos, abans de passar a classes d'equivalència, la factorització només defineix un preordre.)
\end{proposicio}

\begin{prova}
El morfisme de factorització $h$, si existeix, és únic, perquè $n$ és mono: de $n\comp h=m=n\comp h'$ se segueix $h=h'$. A més, $h$ és automàticament un monomorfisme: si $h\comp u=h\comp v$, aleshores $m\comp u=n\comp h\comp u=n\comp h\comp v=m\comp v$, i com que $m$ és mono, $u=v$. La bona definició sobre classes és immediata: substituir $m$ o $n$ per equivalents compon $h$ amb isomorfismes. La reflexivitat usa $h=\id$; la transitivitat, la composició de factoritzacions.

Per a l'antisimetria, suposem $[m]\leq[n]$ i $[n]\leq[m]$, amb factoritzacions $h\colon S\to T$ i $k\colon T\to S$. Aleshores $n\comp h=m$ i $m\comp k=n$, d'on $n\comp h\comp k=m\comp k=n$; com que $n$ és mono, $h\comp k=\id_T$. Simètricament, $k\comp h=\id_S$. Per tant $h$ és un isomorfisme amb $n\comp h=m$, és a dir, $m\sim n$ i $[m]=[n]$.
\end{prova}

\begin{observacio}[Categories ben potenciades]\index{categoria!ben potenciada}\index{ben potenciada}
La petitesa local garanteix que cada $\Hom(A,B)$ és un conjunt, però \emph{no} que la col·lecció $\Sub(A)$ ho sigui: en principi podria ser massa gran. Direm que una categoria és \textbf{ben potenciada} si, per a cada objecte $A$, els subobjectes de $A$ formen un conjunt. En el marc de classes que adoptem, la formulació precisa demana un conjunt de monomorfismes representants, perquè cada classe d'equivalència pot ser una classe pròpia (apèndix~\ref{cap:mida}). Aquesta és la hipòtesi que garanteix que $\Sub(A)$ és un conjunt de debò i que, més endavant, es pugui reunir en un prefeix
\[
\Sub\colon\cat{C}\op\longrightarrow\cat{Set}.
\]
Totes les categories que trobarem seran ben potenciades; a $\cat{Set}$, per exemple, $\Sub(A)\cong\mathcal{P}(A)$ és un conjunt. De fet, l'existència d'un classificador de subobjectes en una categoria localment petita ja implica aquesta condició.
\end{observacio}

\begin{observacio}[Predicats com a subobjectes]
Aquesta és la lectura lògica fonamental. Pensem un objecte $A$ com un \textbf{tipus} i un subobjecte $[m]\in\Sub(A)$ com un \textbf{predicat} sobre aquell tipus: la propietat \enquote{pertànyer a $S$}. L'ordre $[m]\leq[n]$ és la \textbf{implicació} entre predicats: tot element que satisfà $m$ satisfà $n$. Veurem que $\Sub(A)$ té sovint estructura de reticle ---amb intersecció i unió de subobjectes--- i fins i tot d'àlgebra de Heyting, cosa que en fa l'àlgebra de proposicions sobre $A$.
\end{observacio}

\section{Reindexació de subobjectes}

Al capítol de límits vam construir, per a un morfisme $f\colon A\to B$ en una categoria amb pullbacks, el functor de reindexació $f^{*}\colon\cat{C}/B\to\cat{C}/A$ (observació~\ref{obs:reindexacio}). La clau ara és que aquesta reindexació \emph{preserva els monomorfismes}, de manera que es restringeix a una operació entre ordres parcials de subobjectes.

\begin{proposicio}[El pullback preserva monomorfismes]
\label{prop:pullback-mono}
En un quadrat cartesià
\[
\begin{tikzpicture}[baseline=(m.center)]
  \matrix (m) [matrix of math nodes, column sep=3em, row sep=2.4em] {
    P & S \\
    A & B \\
  };
  \draw[-{Stealth[scale=1.1]}] (m-1-1) -- node[above] {$f'$} (m-1-2);
  \draw[-{Stealth[scale=1.1]}] (m-1-1) -- node[left] {$m'$} (m-2-1);
  \draw[-{Stealth[scale=1.1]}] (m-1-2) -- node[right] {$m$} (m-2-2);
  \draw[-{Stealth[scale=1.1]}] (m-2-1) -- node[below] {$f$} (m-2-2);
  \node at ([xshift=0.8em,yshift=-0.7em]m-1-1) {$\scriptstyle\lrcorner$};
\end{tikzpicture}
\qquad
m\comp f'=f\comp m',
\]
si $m$ és un monomorfisme, aleshores el seu pullback $m'\colon P\to A$ al llarg de $f$ també ho és.
\end{proposicio}

\begin{prova}
Siguin $u,v\colon Z\to P$ amb $m'\comp u=m'\comp v$. Per la commutativitat del quadrat,
\[
m\comp f'\comp u=f\comp m'\comp u=f\comp m'\comp v=m\comp f'\comp v,
\]
i, com que $m$ és mono, $f'\comp u=f'\comp v$. Així, $u$ i $v$ tenen les mateixes composicions amb les dues projeccions del pullback, $m'$ i $f'$, i per la unicitat de la propietat universal del pullback, $u=v$. Per tant $m'$ és un monomorfisme.
\end{prova}

\begin{observacio}
La conseqüència que ens interessa és immediata: com que la reindexació $f^{*}$ envia monomorfismes a monomorfismes, es restringeix a una operació ben definida entre els ordres parcials de subobjectes, que és la que ara definim. A $\cat{Set}$, aquesta operació és l'antiimatge $S\mapsto f^{-1}(S)$, tal com vam veure en introduir el pullback.
\end{observacio}

\begin{definicio}[Functor de reindexació sobre subobjectes]\index{reindexació!de subobjectes}\index{substitució}
\label{def:reindexacio-sub}
Sigui $f\colon A\to B$ un morfisme en una categoria $\cat{C}$ amb pullbacks. La \textbf{reindexació} indueix una aplicació
\[
f^{*}\colon\Sub(B)\longrightarrow\Sub(A)
\]
que envia un subobjecte $[m]$, amb $m\colon S\rightarrowtail B$, al subobjecte $[m']$ del seu pullback al llarg de $f$. Per la proposició~\ref{prop:pullback-mono}, $m'$ és mono, de manera que $f^{*}[m]$ és un subobjecte de $A$ ben definit.
\end{definicio}

\begin{proposicio}[La reindexació és monòtona]
\label{prop:reindexacio-monotona}
L'aplicació $f^{*}\colon\Sub(B)\to\Sub(A)$ preserva l'ordre: si $[m]\leq[n]$ a $\Sub(B)$, aleshores $f^{*}[m]\leq f^{*}[n]$ a $\Sub(A)$. És a dir, $f^{*}$ és una aplicació monòtona.
\end{proposicio}

\begin{prova}
Suposem $[m]\leq[n]$, amb $m\colon S\rightarrowtail B$, $n\colon T\rightarrowtail B$ i $h\colon S\to T$ tal que $n\comp h=m$. Siguin $(P,m',f_m)$ un pullback de $m$ al llarg de $f$ i $(Q,n',f_n)$ un de $n$, amb $f\comp m'=m\comp f_m$ i $f\comp n'=n\comp f_n$. La parella $m'\colon P\to A$, $h\comp f_m\colon P\to T$ és un con per al pullback de $n$, perquè
\[
f\comp m'=m\comp f_m=n\comp h\comp f_m.
\]
Per tant hi ha un únic $k\colon P\to Q$ amb $n'\comp k=m'$ i $f_n\comp k=h\comp f_m$. La primera igualtat diu que $[m']\leq[n']$. Observem que l'argument val per a representants i pullbacks \emph{qualssevol}.
\end{prova}

\begin{observacio}[$f^{*}$ està ben definida]\label{obs:reindexacio-ben-definida}
La definició~\ref{def:reindexacio-sub} tria un representant $m$ de la classe i un pullback. El resultat no depèn de cap de les dues tries. Si $m\sim m_1$, aleshores $[m]\leq[m_1]$ i $[m_1]\leq[m]$, i la demostració anterior, aplicada en tots dos sentits, dona que els pullbacks corresponents són equivalents. El cas de dos pullbacks del mateix $m$ és el cas particular $m_1=m$.
\end{observacio}

\begin{proposicio}[El functor de subobjectes]\label{prop:functor-sub}\index{Sub@$\Sub$!functor}
Si $\cat{C}$ és ben potenciada i té pullbacks, les reindexacions defineixen un functor
\[
\Sub\colon\cat{C}\op\longrightarrow\cat{Poset},
\]
que envia cada objecte $A$ a l'ordre parcial $\Sub(A)$ i cada morfisme $f\colon A\to B$ a l'aplicació monòtona $f^{*}\colon\Sub(B)\to\Sub(A)$. És a dir, $\id_A^{*}=\id_{\Sub(A)}$ i $(g\comp f)^{*}=f^{*}\comp g^{*}$, com a igualtats d'aplicacions entre ordres parcials.
\end{proposicio}

\begin{prova}
Com que $\cat{C}$ és ben potenciada, cada $\Sub(A)$ és un conjunt parcialment ordenat, i cada $f^{*}$ és una aplicació monòtona (proposició~\ref{prop:reindexacio-monotona}) ben definida (observació~\ref{obs:reindexacio-ben-definida}). Per aquesta mateixa observació, per calcular $f^{*}[m]$ podem fer servir \emph{qualsevol} pullback de $m$ al llarg de $f$. Queden les dues lleis functorials.

\emph{Identitats.} Sigui $m\colon S\rightarrowtail A$. El quadrat format per $m$ als dos costats verticals i per les identitats $\id_S$ i $\id_A$ als horitzontals és un pullback de $m$ al llarg de $\id_A$. En efecte, si $x\colon Z\to A$ i $y\colon Z\to S$ compleixen $\id_A\comp x=m\comp y$, l'única fletxa $z\colon Z\to S$ amb $m\comp z=x$ i $\id_S\comp z=y$ és $z=y$. Per tant $\id_A^{*}[m]=[m]$.

\emph{Composició.} Siguin $f\colon A\to B$, $g\colon B\to C$ i $m\colon S\rightarrowtail C$. Prenem un pullback $(Q,m_1,g')$ de $m$ al llarg de $g$, amb $g\comp m_1=m\comp g'$, de manera que $g^{*}[m]=[m_1]$. Prenem després un pullback $(P,m_2,f')$ de $m_1$ al llarg de $f$, amb $f\comp m_2=m_1\comp f'$, de manera que $f^{*}(g^{*}[m])=[m_2]$. Vegem que $(P,m_2,g'\comp f')$ és un pullback de $m$ al llarg de $g\comp f$; així $(g\comp f)^{*}[m]=[m_2]$ i les dues classes coincideixen. Primer, el quadrat commuta: $(g\comp f)\comp m_2=g\comp m_1\comp f'=m\comp(g'\comp f')$. Siguin ara $x\colon Z\to A$ i $y\colon Z\to S$ amb $g\comp f\comp x=m\comp y$.
\begin{itemize}
\item \emph{Existència.} La parella $(f\comp x,\,y)$ és un con per al primer pullback, de manera que hi ha un únic $q\colon Z\to Q$ amb $m_1\comp q=f\comp x$ i $g'\comp q=y$. Aleshores $(x,q)$ és un con per al segon pullback, i hi ha un únic $p\colon Z\to P$ amb $m_2\comp p=x$ i $f'\comp p=q$. Per tant, $m_2\comp p=x$ i $(g'\comp f')\comp p=y$.
\item \emph{Unicitat.} Si $p'\colon Z\to P$ també compleix $m_2\comp p'=x$ i $g'\comp f'\comp p'=y$, aleshores $f'\comp p'$ compleix $m_1\comp(f'\comp p')=f\comp m_2\comp p'=f\comp x$ i $g'\comp(f'\comp p')=y$. Per la unicitat al primer pullback, $f'\comp p'=q$. Per la unicitat al segon, $p'=p$.
\end{itemize}
Aquest argument és el lema d'enganxament de pullbacks, en la direcció que aquí cal (exercici resolt corresponent de la Pràctica del capítol~\ref{cap:limits}). Per tant $(g\comp f)^{*}=f^{*}\comp g^{*}$.

Les dues lleis són \emph{igualtats} d'aplicacions entre ordres parcials, i no només isomorfismes. A nivell de monomorfismes, el pullback al llarg de $g\comp f$ i el pullback iterat només coincideixen llevat d'isomorfisme. En passar a classes, però, aquest isomorfisme queda absorbit.
\end{prova}

\begin{observacio}
Aquí no cal cap noció de pseudofunctor. La coherència \enquote{llevat d'isomorfisme} que vam trobar en reindexar les categories tall $\cat{C}/B$ (observació~\ref{obs:reindexacio}) desapareix en passar a classes de subobjectes. Aquest functor és precisament el que el capítol següent voldrà representar.
\end{observacio}

\begin{observacio}[La substitució]
Aquesta és la interpretació que anunciàvem. Si pensem un subobjecte $[m]\in\Sub(B)$ com un predicat $P(y)$ sobre la variable $y:B$, i $f\colon A\to B$ com un \enquote{terme} $f(x)$ amb $x:A$, aleshores $f^{*}[m]\in\Sub(A)$ és el predicat $P(f(x))$: la \textbf{substitució} de $y$ per $f(x)$ dins de $P$. A $\cat{Set}$, on $[m]$ és un subconjunt $S\subseteq B$, la reindexació és l'antiimatge:
\[
f^{*}(S)=f^{-1}(S)=\{x\in A\mid f(x)\in S\}.
\]
La monotonia de $f^{*}$ diu que la substitució respecta la implicació: si $P\Rightarrow Q$, aleshores $P(f)\Rightarrow Q(f)$.
\end{observacio}

\begin{exemple}[Subgrafs]\label{ex:subgrafs}\index{graf dirigit!subobjectes}
A $\cat{Graph}$, els monomorfismes són els morfismes de grafs injectius sobre vèrtexs i sobre arestes (vegeu els Exercicis). Un subobjecte d'un graf $G$ es pot representar, doncs, per un \textbf{subgraf}: un subconjunt de vèrtexs $V'\subseteq V_G$ i un subconjunt d'arestes $E'\subseteq E_G$ tancat per origen i destí, és a dir, tal que l'origen i el destí de tota aresta de $E'$ són a $V'$. La reindexació al llarg d'un morfisme de grafs $F\colon G'\to G$ es calcula component a component, per antiimatge de vèrtexs i d'arestes:
\[
F^{*}(V',E')=\bigl(F_V^{-1}(V'),\ F_E^{-1}(E')\bigr),
\]
que torna a ser un subgraf de $G'$, perquè $F$ preserva origen i destí. L'identificació \enquote{predicat $=$ subobjecte} no és, doncs, exclusiva dels conjunts: un predicat sobre un graf és un subgraf, i la substitució és l'antiimatge component a component.
\end{exemple}

\section{Intersecció de subobjectes}

L'ordre parcial $\Sub(A)$ té més estructura que l'ordre: quan la categoria té pullbacks, admet \textbf{ínfims binaris}, que interpreten la conjunció de predicats.

\begin{proposicio}[La intersecció és un pullback]\index{subobjecte!intersecció}
\label{prop:interseccio}
Siguin $m\colon S\rightarrowtail A$ i $n\colon T\rightarrowtail A$ dos subobjectes de $A$. El seu \textbf{pullback} sobre $A$,
\[
\begin{tikzpicture}[baseline=(m.center)]
  \matrix (m) [matrix of math nodes, column sep=3em, row sep=2.4em] {
    S\cap T & T \\
    S & A \\
  };
  \draw[-{Stealth[scale=1.1]}] (m-1-1) -- node[above] {$q$} (m-1-2);
  \draw[-{Stealth[scale=1.1]}] (m-1-1) -- node[left] {$p$} (m-2-1);
  \draw[-{Stealth[scale=1.1]}] (m-1-2) -- node[right] {$n$} (m-2-2);
  \draw[-{Stealth[scale=1.1]}] (m-2-1) -- node[below] {$m$} (m-2-2);
  \node at ([xshift=1em,yshift=-0.7em]m-1-1) {$\scriptstyle\lrcorner$};
\end{tikzpicture}
\qquad
m\comp p=n\comp q,
\]
amb el morfisme induït $m\comp p=n\comp q\colon S\cap T\to A$, que és un monomorfisme, representa l'\textbf{ínfim} $[m]\wedge[n]$ a $\Sub(A)$. Escrivim també $S\cap T=S\times_A T$.
\end{proposicio}

\begin{prova}
El morfisme $m\comp p\colon S\cap T\to A$ és mono per ser composició de monos (la projecció $p$ del pullback és mono per la proposició~\ref{prop:pullback-mono}, i $m$ és mono). Per construcció $[m\wedge n]\leq[m]$ i $[m\wedge n]\leq[n]$, ja que hi ha les projeccions cap a $S$ i $T$. I si $[k]\leq[m]$ i $[k]\leq[n]$ per a un tercer subobjecte $k\colon R\rightarrowtail A$, les factoritzacions $R\to S$ i $R\to T$ són compatibles sobre $A$, de manera que la propietat universal del pullback dona una única factorització $R\to S\cap T$; per tant $[k]\leq[m]\wedge[n]$. Això caracteritza l'ínfim.
\end{prova}

\begin{observacio}
En termes lògics, $[m]\wedge[n]$ és la \textbf{conjunció} dels predicats: l'element pertany a $S\cap T$ si i només si pertany a $S$ \emph{i} a $T$. A $\cat{Set}$ recuperem la intersecció de subconjunts. En qualsevol categoria, el subobjecte \textbf{màxim} de $\Sub(A)$ és $[\id_A]$ (el predicat sempre cert), com es comprova a la Pràctica; i quan hi ha objecte inicial estricte, el subobjecte \textbf{mínim} és la injecció de l'objecte inicial (el predicat sempre fals). Amb intersecció i aquests extrems, $\Sub(A)$ és un \emph{semireticle amb màxim}. Les \textbf{unions} ---la disjunció--- exigeixen estructura addicional: la regularitat proporciona imatges, existencials i interseccions, però no, en general, suprems finits estables a les fibres $\Sub(A)$. Aquests apareixen, per exemple, en les \textbf{categories coherents}, que es tractaran al volum de lògica categòrica.
\end{observacio}

\section{Epimorfismes regulars i imatges}

Per definir la imatge d'un morfisme, ens cal la classe adequada d'epimorfismes: els que sorgeixen com a quocients.

\begin{definicio}[Epimorfisme regular]\index{epimorfisme!regular}
\label{def:epi-regular}
Un morfisme $e\colon A\to B$ és un \textbf{epimorfisme regular} si és el coequalitzador d'alguna parella de morfismes $u,v\colon R\to A$.
\end{definicio}

\begin{definicio}[Parella nucli]\index{parella nucli}
\label{def:parella-nucli}
La \textbf{parella nucli} d'un morfisme $f\colon A\to B$ és el pullback de $f$ amb ell mateix,
\[
\begin{tikzpicture}[baseline=(q.center)]
  \matrix (q) [matrix of math nodes, column sep=3em, row sep=2.4em] {
    R & A \\
    A & B \\
  };
  \draw[-{Stealth[scale=1.1]}] (q-1-1) -- node[above] {$p_2$} (q-1-2);
  \draw[-{Stealth[scale=1.1]}] (q-1-1) -- node[left] {$p_1$} (q-2-1);
  \draw[-{Stealth[scale=1.1]}] (q-1-2) -- node[right] {$f$} (q-2-2);
  \draw[-{Stealth[scale=1.1]}] (q-2-1) -- node[below] {$f$} (q-2-2);
  \node at ([xshift=0.8em,yshift=-0.7em]q-1-1) {$\scriptstyle\lrcorner$};
\end{tikzpicture}
\qquad
f\comp p_1=f\comp p_2;
\]
els seus dos morfismes $p_1,p_2\colon R\to A$ recullen les parelles d'elements que $f$ identifica. A $\cat{Set}$, $R=\{(a,a')\mid f(a)=f(a')\}$.
\end{definicio}

\begin{lema}\label{lema:epi-regular-nucli}\index{epimorfisme!regular}\index{parella nucli}
En una categoria amb pullbacks, si $e\colon A\to B$ és un epimorfisme regular, aleshores $e$ és el coequalitzador de la seva parella nucli.
\end{lema}

\begin{prova}
Suposem que $e$ és el coequalitzador de $u,v\colon K\rightrightarrows A$, i sigui $p_1,p_2\colon R\rightrightarrows A$ la parella nucli de $e$. Com que $e\comp u=e\comp v$, la propietat universal del pullback dona $k\colon K\to R$ amb $p_1\comp k=u$ i $p_2\comp k=v$. Si $h\colon A\to Z$ satisfà $h\comp p_1=h\comp p_2$, aleshores
\[
h\comp u=h\comp p_1\comp k=h\comp p_2\comp k=h\comp v,
\]
i com que $e$ és el coequalitzador de $u$ i $v$, hi ha una única $\bar h\colon B\to Z$ amb $\bar h\comp e=h$. Com que, a més, $e\comp p_1=e\comp p_2$, això és exactament la propietat universal del coequalitzador de $p_1$ i $p_2$.
\end{prova}

\begin{definicio}[Factorització imatge]\index{imatge}\index{factorització!imatge}
\label{def:imatge}
Una \textbf{factorització imatge} de $f\colon A\to B$ és una descomposició $f=m\comp e$ amb $A\xrightarrow{\,e\,}I\xrightarrow{\,m\,}B$, on $e$ és un epimorfisme regular i $m$ un monomorfisme. El subobjecte $[m]\in\Sub(B)$ és la \textbf{imatge} de $f$, denotada $\Imatge(f)$.
\end{definicio}

\begin{proposicio}[La imatge és el mínim subobjecte]
\label{prop:imatge-unica}
Quan existeix, la factorització imatge és única llevat d'isomorfisme canònic, i $\Imatge(f)$ és el \emph{menor} subobjecte de $B$ a través del qual $f$ es factoritza: si $f$ es factoritza per un mono $n\colon T\rightarrowtail B$, aleshores $\Imatge(f)\leq[n]$.
\end{proposicio}

\begin{prova}
Sigui $f=m\comp e$ una factorització imatge. Com que $m$ és mono, qualsevol parella nucli de $e$ és també una parella nucli de $f$ (perquè $m\comp e\comp u=m\comp e\comp v$ si i només si $e\comp u=e\comp v$); en particular, les parelles nucli escollides per a $e$ i per a $f$ són canònicament isomorfes. Pel lema~\ref{lema:epi-regular-nucli}, $e$ és el coequalitzador de la seva parella nucli $p_1,p_2$. Si $f=m'\comp e'$ és una altra factorització amb $m'$ mono, aleshores $e'$ iguala les branques d'aquesta parella nucli, perquè $m'\comp e'\comp p_1=f\comp p_1=f\comp p_2=m'\comp e'\comp p_2$ i $m'$ és mono, de manera que $e'$ es factoritza per $e$ mitjançant un únic morfisme; simètricament s'obté l'invers, i les dues composicions són identitats perquè $e$ és epi. Això dona l'isomorfisme canònic. La minimalitat en segueix: si $f=n\comp g$ amb $n$ mono, aleshores $g$ iguala la parella nucli i es factoritza per $e$, d'on $\Imatge(f)=[m]\leq[n]$.
\end{prova}

\begin{observacio}
A $\cat{Set}$, $\Imatge(f)$ és el subconjunt imatge $f(A)\subseteq B$, amb $e$ l'aplicació exhaustiva $A\to f(A)$ i $m$ la injecció. Recuperem la noció usual d'imatge.
\end{observacio}

\section{Categories regulars i estabilitat}

Perquè la teoria d'imatges sigui utilitzable, cal que les imatges es comportin bé sota el canvi de base.

\begin{definicio}[Categoria regular]\index{categoria!regular}
\label{def:regular}
Una categoria $\cat{C}$ és \textbf{regular} si té tots els límits finits, tota fletxa admet una factorització imatge (epi regular seguit de mono), i els epimorfismes regulars són \textbf{estables per pullback}.
\end{definicio}

\begin{exemple}
Són regulars $\cat{Set}$, tota categoria d'àlgebres d'una teoria algebraica (grups, anells, mòduls\ldots), tota categoria abeliana i tot topos elemental. La categoria $\cat{Top}$ \emph{no} és regular, perquè les imatges no hi són estables per pullback.
\end{exemple}

\begin{proposicio}[Estabilitat de la imatge]
\label{prop:imatge-estable}
En una categoria regular, la formació d'imatges commuta amb el pullback: si es forma el pullback de $f\colon A\to B$ al llarg de $g\colon B'\to B$, aleshores $\Imatge(f')=g^{*}\bigl(\Imatge(f)\bigr)$, on $f'$ és el morfisme reindexat.
\end{proposicio}

\begin{prova}
Factoritzem $f=m\comp e$. En prendre el pullback al llarg de $g$, l'epimorfisme regular $e$ dona un epi regular $e'$ (per l'estabilitat) i el mono $m$ dona un mono $m'$ (proposició~\ref{prop:pullback-mono}). Així $f'=m'\comp e'$ és una factorització epi regular\,--\,mono, i per la unicitat de la imatge, $m'$ representa $g^{*}(\Imatge(f))$.
\end{prova}

\section{L'existencial com a adjunt de la reindexació}

De tot això en resulta que la reindexació té un adjunt per l'esquerra, que reprèn el fil del capítol d'adjuncions.

\begin{teorema}[$\exists_f\dashv f^{*}$]\index{quantificador!existencial}\index{adjunció!existencial}
\label{teo:existencial-adjunt}
Sigui $\cat{C}$ una categoria regular i $f\colon A\to B$. L'operació
\[
\exists_f[s]:=\Imatge(f\comp s)\in\Sub(B)
\]
és adjunta per l'esquerra de la reindexació $f^{*}\colon\Sub(B)\to\Sub(A)$; és a dir,
\[
\exists_f[s]\leq[t]\quad\Longleftrightarrow\quad[s]\leq f^{*}[t].
\]
\end{teorema}

\begin{observacio}
L'operació $\exists_f$ està definida sobre classes de subobjectes. Si $s\sim s'$, és a dir, $s'=s\comp\varphi$ amb $\varphi$ isomorfisme, aleshores $f\comp s'=(f\comp s)\comp\varphi$, i les factoritzacions imatge de $f\comp s$ i de $f\comp s'$ donen el mateix subobjecte de $B$. És exactament la unicitat de la imatge llevat d'isomorfisme (proposició~\ref{prop:imatge-unica}) el que fa que $\exists_f$ estigui ben definida.
\end{observacio}

\begin{prova}
Si $[s]\leq f^{*}[t]$, aleshores $S$ es factoritza a través del pullback $f^{*}T$, i component-lo amb la projecció cap a $T$ s'obté que $f\comp s$ es factoritza per $t$; com que $\Imatge(f\comp s)$ és el mínim subobjecte amb aquesta propietat, $\exists_f[s]\leq[t]$. Recíprocament, si $\exists_f[s]\leq[t]$, aleshores $f\comp s=t\comp k$ per a algun $k$, i aquesta igualtat exhibeix $S$ com un con sobre el diagrama del pullback $f^{*}T$, de manera que $[s]\leq f^{*}[t]$.
\end{prova}

\begin{observacio}[L'universal com a estructura addicional]\index{quantificador!universal}\index{adjunció!quantificadors}
En una categoria regular, la reindexació $f^{*}$ té sempre l'adjunt per l'esquerra $\exists_f$ que acabem de construir, però \emph{no} té necessàriament un adjunt per la dreta. Quan cada $f^{*}$ admet a més un adjunt per la dreta, el denotem $\forall_f$ i queda caracteritzat per
\[
f^{*}[t]\leq[s]\quad\Longleftrightarrow\quad[t]\leq\forall_f[s].
\]
Aquesta és estructura addicional, que la regularitat no garanteix: la regularitat només dona el costat existencial. Apareix, per exemple, en les categories localment cartesianes tancades amb les hipòtesis de mida adequades i, en particular, en tot topos elemental. La construcció de $\forall_f$ i el seu paper lògic es tracten al volum de lògica categòrica; no s'obté simplement dualitzant la construcció de la imatge dins de la mateixa fibra de subobjectes.
\end{observacio}

\begin{observacio}[Compatibilitat amb la substitució]\index{Beck--Chevalley}
En una categoria regular, l'estabilitat de les imatges per pullback dona una compatibilitat entre l'existencial i la substitució. Per a un quadrat cartesià
\[
\begin{tikzpicture}[baseline=(q.center)]
  \matrix (q) [matrix of math nodes, column sep=3em, row sep=2.4em] {
    A' & A \\
    B' & B \\
  };
  \draw[-{Stealth[scale=1.1]}] (q-1-1) -- node[above] {$h$} (q-1-2);
  \draw[-{Stealth[scale=1.1]}] (q-1-1) -- node[left] {$f'$} (q-2-1);
  \draw[-{Stealth[scale=1.1]}] (q-1-2) -- node[right] {$f$} (q-2-2);
  \draw[-{Stealth[scale=1.1]}] (q-2-1) -- node[below] {$g$} (q-2-2);
  \node at ([xshift=0.8em,yshift=-0.7em]q-1-1) {$\scriptstyle\lrcorner$};
\end{tikzpicture}
\]
es té la \textbf{condició de Beck--Chevalley}
\[
g^{*}\comp\exists_f=\exists_{f'}\comp h^{*}\colon\Sub(A)\to\Sub(B'):
\]
quantificar existencialment i després substituir és el mateix que substituir i després quantificar. No ho demostrem aquí; és un primer exemple del tipus de coherència que la lògica categòrica exigeix.
\end{observacio}

\begin{observacio}[Lectura lògica]
Aquesta adjunció té una lectura lògica il·lustrativa. Si $f$ és la projecció $\pi\colon X\times Y\to X$ que oblida la variable $Y$, i pensem un subobjecte de $X\times Y$ com un predicat $s(x,y)$, aleshores
\[
\exists_\pi[s]=\{x\mid\exists y,\ s(x,y)\}.
\]
L'existencial apareix, doncs, com a adjunt per l'esquerra de la substitució. Quan la categoria té també els adjunts per la dreta, l'universal es llegeix anàlogament com $\forall_\pi[s]=\{x\mid\forall y,\ s(x,y)\}$. El desenvolupament d'aquesta idea com a \emph{sistema lògic} ---amb el seu càlcul, la interpretació de l'existencial per la factorització imatge i les condicions de coherència pertinents--- és objecte de l'estudi de la lògica categòrica, i queda fora de l'abast d'aquest volum.
\end{observacio}

\begin{resumcap}
\apartat{Idees centrals}
\begin{enumerate}
  \item Un subobjecte és una classe d'equivalència de monomorfismes amb el mateix codomini; $\Sub(A)$ és un ordre parcial, l'àlgebra de predicats sobre $A$. A $\cat{Set}$ són els subconjunts, i a $\cat{Graph}$, els subgrafs.
  \item La reindexació $f^{*}$, construïda per pullback, interpreta la substitució; en una categoria ben potenciada amb pullbacks dona un functor $\Sub\colon\cat{C}\op\to\cat{Poset}$, amb igualtats estrictes.
  \item La intersecció de subobjectes és el seu pullback i interpreta la conjunció; el màxim $[\id_A]$ interpreta el predicat cert.
  \item Un epimorfisme regular és el coequalitzador de la seva parella nucli; la factorització imatge (epi regular seguit de mono) és única llevat d'isomorfisme, i en una categoria regular existeix i és estable per pullback.
  \item En una categoria regular, $f^{*}$ té un adjunt per l'esquerra $\exists_f$, que interpreta la quantificació existencial; l'universal $\forall_f$ és estructura addicional que la regularitat no garanteix.
\end{enumerate}
\apartat{Errors típics} Creure que la regularitat dona unions de subobjectes (calen categories coherents) o l'adjunt dret $\forall_f$; confondre els epimorfismes amb els epimorfismes regulars (a $\cat{Top}$, els primers són les aplicacions contínues exhaustives i els segons, les aplicacions quocient); confondre la reindexació $f^{*}$ amb la imatge $\exists_f$; i usar \enquote{imatge}, \enquote{buit} o \enquote{elements} en una categoria arbitrària sense declarar les hipòtesis.
\apartat{Lectura recomanada} \cite[cap.~5 i~\S9.5]{awodey} per a subobjectes, imatges i l'adjunció existencial; \cite[§§2.8, 2.10, 9.4 i 9.8]{barrwells} per a subobjectes, epimorfismes regulars, categories regulars i sistemes de factorització; \cite[V.7]{maclane} per al tractament clàssic dels subobjectes; \cite[A1.3]{johnstone} i \cite{borceux} per a categories regulars amb més profunditat; \cite{jacobs} per a la lectura fibracional de la reindexació.
\end{resumcap}

\chapter[Classificadors de subobjectes i topos]{\texorpdfstring{Classificadors de subobjectes\\ i introducció als topos}{Classificadors de subobjectes i introducció als topos}}
\label{cap:topos}

\begin{objectius}
\apartat{Prerequisits} Subobjectes i el functor $\Sub$ (cap.~\ref{cap:subobjectes}); representabilitat i Yoneda (cap.~\ref{cap:yoneda}); tancament cartesià (cap.~\ref{cap:cartesianes}).
\apartat{Objectius} En acabar el capítol, el lector hauria de poder:
\begin{itemize}
  \item definir el classificador de subobjectes $\Omega$ i el morfisme característic d'un subobjecte;
  \item expressar la pertinença d'un element generalitzat com una igualtat de morfismes cap a $\Omega$;
  \item demostrar que tenir classificador equival a la representabilitat del functor $\Sub$;
  \item enunciar la definició de topos elemental i reconèixer $\Omega^{E}$ com a objecte de parts;
  \item calcular morfismes característics a $\cat{Set}$ i a $\cat{Graph}$, i explicar per què la representabilitat de $\Sub$ converteix els predicats en morfismes i la substitució en composició;
  \item reconèixer com a topos $\cat{Set}$, $\cat{FinSet}$, les categories de prefeixos sobre una categoria petita, amb el classificador de sedassos, i $\cat{Graph}$, i distingir els topos elementals de les categories de prefeixos;
  \item explicar en quin sentit $\Omega$ és un objecte de valors de veritat, i per què això no fa booleà tot topos.
\end{itemize}
\end{objectius}

Al capítol anterior hem vist que els subobjectes d'un objecte $A$ formen un ordre parcial $\Sub(A)$, que podem llegir com l'ordre de predicats sobre $A$, i que la reindexació el fa functorial: en una categoria ben potenciada amb pullbacks obtenim un functor $\Sub\colon\cat{C}\op\to\cat{Poset}$ (proposició~\ref{prop:functor-sub}). Sorgeix una pregunta natural: hi ha un objecte que \emph{representi} aquest functor? És a dir, un objecte $\Omega$ tal que donar un subobjecte de $A$ equivalgui a donar un morfisme $A\to\Omega$? Quan aquest objecte existeix, s'anomena \textbf{classificador de subobjectes}, i les categories que en tenen ---juntament amb límits finits i exponencials--- són els \textbf{topos elementals}. Aquest capítol n'introdueix la definició i els exemples bàsics. No en desenvolupem la lògica interna: aquesta pertany a l'estudi de la lògica categòrica, per al qual aquest volum és la preparació.

\section{El classificador de subobjectes}

A $\cat{Set}$, un subconjunt $U\subseteq E$ es descriu per la seva \textbf{funció característica} $\chi_U\colon E\to\{0,1\}$, que val $1$ als elements de $U$ i $0$ a la resta. El conjunt $\{0,1\}$ actua com un objecte de \enquote{valors de veritat}, i triar un subconjunt equival a triar una funció cap a ell. La noció següent generalitza aquesta idea a una categoria arbitrària.

\begin{definicio}[Classificador de subobjectes]\index{classificador de subobjectes}\index{Omega@$\Omega$}
\label{def:classificador}
Sigui $\cat{E}$ una categoria amb tots els límits finits. Un \textbf{classificador de subobjectes} és un objecte $\Omega$ juntament amb un morfisme $\mathsf{true}\colon 1\to\Omega$ des de l'objecte terminal, amb la propietat universal següent: per a cada objecte $E$ i cada subobjecte $m\colon U\rightarrowtail E$, existeix un \emph{únic} morfisme $\chi_m\colon E\to\Omega$, el \textbf{morfisme característic} del subobjecte, que fa del quadrat
\[
\begin{tikzpicture}[baseline=(m.center)]
  \matrix (m) [matrix of math nodes, column sep=3em, row sep=2.4em] {
    U & 1 \\
    E & \Omega \\
  };
  \draw[-{Stealth[scale=1.1]}] (m-1-1) -- node[above] {$!$} (m-1-2);
  \draw[-{Stealth[scale=1.1]}] (m-1-1) -- node[left] {$m$} (m-2-1);
  \draw[-{Stealth[scale=1.1]}] (m-1-2) -- node[right] {$\mathsf{true}$} (m-2-2);
  \draw[-{Stealth[scale=1.1]}] (m-2-1) -- node[below] {$\chi_m$} (m-2-2);
  \node at ([xshift=0.8em,yshift=-0.7em]m-1-1) {$\scriptstyle\lrcorner$};
\end{tikzpicture}
\]
un pullback. És a dir, $U$ es recupera com el pullback de $\mathsf{true}$ al llarg de $\chi_m$. A $\cat{Set}$, el morfisme característic és la funció característica usual.
\end{definicio}

La condició diu que cada subobjecte de $E$ té un i només un morfisme característic, i que el subobjecte es reconstrueix com l'antiimatge del \enquote{punt vertader} $\mathsf{true}$. A $\cat{Set}$, $\Omega=\{0,1\}$ i $\mathsf{true}\colon 1\to\{0,1\}$ tria l'element $1$; la funció característica és la de sempre, i el pullback de $\mathsf{true}$ al llarg de $\chi_U$ recupera exactament $U$.

\begin{exemple}[Un morfisme característic calculat a $\cat{Set}$]\label{ex:caracteristic-set}\index{morfisme característic!a Set@a $\cat{Set}$}
Sigui $E=\{1,2,3,4\}$ i $U=\{2,4\}$, el predicat \enquote{$x$ és parell}, amb la injecció $m\colon U\hookrightarrow E$. El morfisme característic és
\[
\chi_U\colon E\to\{0,1\},\qquad 1\mapsto0,\quad 2\mapsto1,\quad 3\mapsto0,\quad 4\mapsto1.
\]
Comprovem les tres afirmacions de la definició sobre aquest cas.
\begin{itemize}
\item \emph{El quadrat commuta.} $\chi_U\comp m$ envia $2$ i $4$ a $1$, que és $\mathsf{true}\comp\,!_U$.
\item \emph{És un pullback.} El pullback de $\mathsf{true}$ al llarg de $\chi_U$ és $\{x\in E\mid\chi_U(x)=1\}=\{2,4\}$, amb la inclusió (el pullback a $\cat{Set}$, capítol~\ref{cap:limits}): recupera exactament $U$.
\item \emph{Unicitat.} Una altra aplicació $\chi\colon E\to\{0,1\}$ amb el mateix pullback ha de valer $1$ exactament a $\{2,4\}$, i per tant és $\chi_U$. Per exemple, l'aplicació que val $1$ a $\{2,3,4\}$ fa commutar el quadrat, però el seu pullback és $\{2,3,4\}$, no $U$: la commutativitat sola no basta.
\end{itemize}
La pertinença per elements, que formalitzarem tot seguit, hi és visible: l'element $x\colon\{\ast\}\to E$ amb $x(\ast)=4$ compleix $\chi_U\comp x=\mathsf{true}$, i el que tria $3$, no. La substitució també: si $g\colon\{a,b,c\}\to E$ és $g(a)=1$, $g(b)=2$, $g(c)=4$, aleshores $\chi_U\comp g$ envia $a\mapsto0$, $b\mapsto1$ i $c\mapsto1$, i classifica $g^{-1}(U)=\{b,c\}$. \emph{Substituir} en un predicat és \emph{precompondre} amb el seu morfisme característic.
\end{exemple}

\begin{observacio}[Pertinença per elements generalitzats]\index{element generalitzat!pertinença}
Sigui $m\colon U\rightarrowtail E$ un subobjecte amb morfisme característic $\chi_m\colon E\to\Omega$, i sigui $x\colon X\to E$ un element generalitzat. Aleshores
\[
x\ \text{factoritza a través de } m
\quad\Longleftrightarrow\quad
\chi_m\comp x=\mathsf{true}\comp\,!_X.
\]
En efecte, si $x=m\comp y$, aleshores $\chi_m\comp x=\chi_m\comp m\comp y=\mathsf{true}\comp\,!_U\comp y=\mathsf{true}\comp\,!_X$; i recíprocament, si $\chi_m\comp x=\mathsf{true}\comp\,!_X$, la propietat universal del pullback dona una factorització (única) de $x$ a través de $m$. És la lectura precisa, sense elements ordinaris, de
\[
x\in U\quad\Longleftrightarrow\quad\chi_U(x)=\mathsf{true}:
\]
el classificador converteix la pertinença a un subobjecte en una igualtat de morfismes cap a $\Omega$.
\end{observacio}

\begin{proposicio}[Classificador i representabilitat]
\label{prop:omega-representa}
Suposem $\cat{E}$ localment petita, ben potenciada i amb límits finits, de manera que el functor de subobjectes pren valors a $\cat{Set}$. Aleshores les dades d'un classificador de subobjectes $(\Omega,\mathsf{true})$ són equivalents a una representació del functor de subobjectes:
\[
\Sub_{\cat{E}}(E)\;\cong\;\Hom_{\cat{E}}(E,\Omega),
\]
natural en $E$. En conseqüència, el classificador, si existeix, és únic llevat d'un únic isomorfisme compatible amb $\mathsf{true}$.
\end{proposicio}

\begin{prova}
\emph{Del classificador a la representació.} Enviem un subobjecte $[m]$ al seu morfisme característic $\chi_m$. L'aplicació és injectiva perquè $\chi_m$ determina $[m]$ com el pullback de $\mathsf{true}$ al llarg de $\chi_m$, i és exhaustiva perquè, per a tot $\chi\colon E\to\Omega$, el pullback de $\mathsf{true}$ al llarg de $\chi$ és un subobjecte el morfisme característic del qual és $\chi$, per la unicitat de la definició~\ref{def:classificador}. La naturalitat en $E$ és la compatibilitat amb la reindexació: reindexar un subobjecte al llarg de $g\colon E'\to E$ correspon a precompondre la característica amb $g$, pel lema d'enganxament de pullbacks.

\emph{De la representació al classificador.} Suposem que $\Sub_{\cat{E}}\cong\Hom_{\cat{E}}(-,\Omega)$ naturalment, i sigui $t\colon T\rightarrowtail\Omega$ un representant de l'element universal, és a dir, del subobjecte de $\Omega$ que correspon a $\id_\Omega$. Pel lema de Yoneda, el subobjecte que correspon a un morfisme $\chi\colon E\to\Omega$ és la reindexació $\chi^{*}[t]$: tot subobjecte és el pullback de $t$ al llarg d'un únic morfisme. Vegem que $T$ és terminal. Per a cada objecte $A$, el subobjecte màxim $[\id_A]$ és el pullback de $t$ al llarg d'un únic $\chi\colon A\to\Omega$, i el pullback dona un morfisme $A\to T$. Si $g\colon A\to T$ és un morfisme qualsevol, el pullback de $t$ al llarg de $t\comp g$ és $[\id_A]$, perquè $t$ és mono; per unicitat, $t\comp g=\chi$, i com que $t$ és mono, $g$ queda determinat. Així hi ha exactament un morfisme $A\to T$, i $T\cong 1$. Identificant $T$ amb $1$, el mono $t$ és un morfisme $\mathsf{true}\colon 1\to\Omega$ que satisfà la definició~\ref{def:classificador}.

La unicitat del classificador és el principi de Yoneda aplicat a dos objectes que representen el mateix functor, amb l'element universal $\mathsf{true}$ com a dada universal.
\end{prova}

\begin{observacio}[Per què la representabilitat és un pas conceptual]\label{obs:sub-representable}
L'equivalència de la proposició no és només una reformulació. El functor $\Sub$ és una dada \emph{dispersa}: un ordre parcial per a cada objecte, i una reindexació per a cada morfisme. La representabilitat el concentra en un sol objecte $\Omega$ i en un sol subobjecte universal, $\mathsf{true}\colon 1\rightarrowtail\Omega$. Aquest canvi té tres conseqüències:
\begin{itemize}
\item \emph{Els predicats esdevenen morfismes.} Un predicat sobre $E$ ja no és una classe d'equivalència de monomorfismes, sinó una fletxa $E\to\Omega$, que es pot compondre, aparellar i currificar com qualsevol altra.
\item \emph{La substitució esdevé composició.} La naturalitat de la bijecció diu que reindexar al llarg de $g$ és precompondre amb $g$, tal com hem vist a l'exemple~\ref{ex:caracteristic-set}.
\item \emph{Tot subobjecte és un cas d'un de sol.} Cada subobjecte és el pullback de $\mathsf{true}$ al llarg del seu morfisme característic. Les operacions lògiques es poden definir, doncs, una sola vegada sobre $\Omega$, com a morfismes $\Omega\times\Omega\to\Omega$, i transportar-se a tots els objectes (exercici resolt~\ref{pr-conjuncio-omega} de la Pràctica).
\end{itemize}
És el mateix pas que va de $\Hom(A\times B,C)$ a l'exponencial $C^B$: una família de conjunts esdevé un objecte de la categoria. Per això el classificador és el que permet que la lògica \enquote{visqui dins} de la categoria. Pel lema de Yoneda, a més, el classificador es pot calcular sondejant-lo amb objectes representables, com farem a $\cat{Graph}$.
\end{observacio}

\section{Topos elementals}

\begin{definicio}[Topos elemental]\index{topos}\index{topos!elemental}
\label{def:topos}
Un \textbf{topos elemental} és una categoria $\cat{E}$ que compleix:
\begin{enumerate}
\item $\cat{E}$ té tots els límits finits;
\item $\cat{E}$ té un classificador de subobjectes $\Omega$;
\item $\cat{E}$ és cartesiana tancada (té tots els exponencials).
\end{enumerate}
\end{definicio}

Aquesta definició, compacta, resulta ser extraordinàriament rica. Es pot demostrar, per exemple, que tot topos té també tots els colímits finits, i que, si $\cat{E}$ és un topos, també ho és cada categoria $\cat{E}/A$ de $\cat{E}$ sobre un objecte $A$. No demostrem aquí aquestes propietats; les recollim per il·lustrar que la combinació \enquote{límits finits $+$ classificador $+$ exponencials} és molt més potent del que sembla a primera vista.

\begin{observacio}[Topos elemental i categoria de prefeixos]\label{obs:elemental-prefeixos}\index{topos!elemental i de prefeixos}
Convé no confondre la definició amb el seu exemple principal. Un topos \emph{elemental} es defineix per tres condicions expressables en el llenguatge de les categories, i només demana límits \emph{finits}. Les categories de prefeixos $\cat{Set}^{\cat{C}\op}$ en són una classe important, però tenen molta més estructura: tots els límits i colímits \emph{petits}, i un classificador que es pot escriure explícitament amb sedassos. No tot topos elemental és d'aquesta mena. La categoria $\cat{FinSet}$ dels conjunts finits és un topos: té límits finits, l'exponencial de dos conjunts finits és finit, i $\{0,1\}$ hi és un classificador. Però no té coproductes infinits: no hi ha cap coproducte d'una família numerable de còpies de $1$, perquè un tal objecte hauria d'admetre una aplicació cap a $\{0,1\}$ per a cada successió de zeros i uns, i un conjunt finit només n'admet un nombre finit. En canvi, tota categoria de prefeixos té tots els colímits petits, i una equivalència de categories els conserva. Per tant, $\cat{FinSet}$ no és equivalent a cap categoria de prefeixos. Entre totes dues nocions hi ha els \emph{topos de Grothendieck} (categories de feixos), que aquí no tractem.
\end{observacio}

\begin{observacio}[Objecte de parts]
En un topos, l'exponencial $\Omega^{E}$ mereix un nom propi: és l'\textbf{objecte de parts} $\mathcal{P}(E)$, l'anàleg categòric del conjunt de parts. En efecte, per la propietat de l'exponencial i la del classificador,
\[
\Hom(A,\Omega^{E})\cong\Hom(A\times E,\Omega)\cong\Sub(A\times E),
\]
de manera que els morfismes cap a $\Omega^{E}$ classifiquen relacions entre $A$ i $E$. En el cas $A=1$, com que $1\times E\cong E$, s'obté
\[
\Hom(1,\Omega^{E})\cong\Sub(E):
\]
els punts globals de $\Omega^{E}$ són exactament els subobjectes de $E$, que és el que justifica el nom d'objecte de parts. Aquesta és la porta d'entrada a la lògica d'ordre superior interna d'un topos, que és tema de l'estudi posterior.
\end{observacio}

\section{Exemples}

\begin{exemple}[Conjunts]
$\cat{Set}$ és un topos. Té límits finits, és cartesià tancat (l'exponencial és el conjunt de funcions) i té classificador $\Omega=\{0,1\}$ amb $\mathsf{true}$ l'element $1$. És el topos prototípic, i tota la teoria es pot llegir com una generalització d'aquest cas.
\end{exemple}

\begin{exemple}[Prefeixos]\index{prefeix@prefeix!topos de}\index{sedàs}
Per a qualsevol categoria petita $\cat{C}$, la categoria de prefeixos $\widehat{\cat{C}}=\cat{Set}^{\cat{C}\op}$ (exemple~\ref{ex:categoria-prefeixos}) és un topos. Ja sabíem que té límits finits, calculats objecte a objecte (proposició~\ref{prop:limits-punt-a-punt}), i que és cartesiana tancada; el que hi falta és el classificador. Aquest es construeix amb els \textbf{sedassos}. Un sedàs sobre un objecte $C$ és un conjunt $S$ de morfismes amb codomini $C$ tal que
\[
f\colon D\to C,\quad f\in S,\quad g\colon X\to D
\qquad\Longrightarrow\qquad
f\comp g\in S.
\]
El prefeix $\Omega$ envia cada objecte $C$ al conjunt $\Omega(C)$ dels sedassos sobre $C$, i cada morfisme $h\colon D\to C$ a l'aplicació
\[
\Omega(h)\colon\Omega(C)\to\Omega(D),
\qquad
S\longmapsto h^{*}S=\{g\colon X\to D\mid h\comp g\in S\},
\]
que torna a donar un sedàs. El punt distingit $\mathsf{true}\colon 1\to\Omega$ tria a cada $C$ el sedàs màxim, format per tots els morfismes cap a $C$.

Si $U\rightarrowtail F$ és un subprefeix, el seu morfisme característic envia $x\in F(C)$ al sedàs dels morfismes que porten $x$ dins de $U$:
\[
\chi_U(C)(x)=\{f\colon D\to C\mid F(f)(x)\in U(D)\}.
\]
Es comprova que aquest conjunt és un sedàs, que les aplicacions $\chi_U(C)$ són naturals en $C$, i que $U$ és el pullback de $\mathsf{true}$ al llarg de $\chi_U$. Aquest exemple cobreix una quantitat enorme de casos i és el pont cap a la teoria de feixos.
\end{exemple}

\begin{exemple}[El topos dels grafs: el classificador calculat]\label{ex:omega-graph}\index{graf dirigit!topos}\index{classificador de subobjectes!a Graph@a $\cat{Graph}$}
Al capítol~\ref{cap:transformacions} vam veure que $\cat{Graph}\cong\cat{Set}^{\cat{P}}$, i com que $\cat{Set}^{\cat{P}}=\cat{Set}^{(\cat{P}\op)\op}$, és la categoria de prefeixos sobre $\cat{P}\op$. Per l'exemple anterior, $\cat{Graph}$ és, doncs, un topos. Calculem-ne el classificador.

\emph{Els vèrtexs i les arestes de $\Omega$.} Els dos representables són el graf $\mathbf V$ d'un sol vèrtex i el graf $\mathbf E$ d'una sola aresta $e\colon s\to t$ (exemple~\ref{ex:homfunctor-grafs}): els vèrtexs d'un graf $G$ corresponen als morfismes $\mathbf V\to G$, i les arestes, als morfismes $\mathbf E\to G$. Aplicant-ho a $\Omega$ i fent servir $\Hom(-,\Omega)\cong\Sub$, obtenim
\[
V_\Omega\cong\Sub(\mathbf V),\qquad E_\Omega\cong\Sub(\mathbf E).
\]
El graf $\mathbf V$ té dos subgrafs, el buit i ell mateix; així, $\Omega$ té dos vèrtexs, que anomenem $0$ (\enquote{el vèrtex no hi és}) i $1$ (\enquote{el vèrtex hi és}). El graf $\mathbf E$ té cinc subgrafs (exemple~\ref{ex:subgrafs}): $\varnothing$, $\{s\}$, $\{t\}$, $\{s,t\}$ i $\{s,t,e\}$. Així, $\Omega$ té cinc arestes.

\emph{Origen i destí.} L'origen d'una aresta $K\subseteq\mathbf E$ de $\Omega$ s'obté reindexant $K$ al llarg del morfisme $\mathbf V\to\mathbf E$ que tria $s$; és a dir, val $1$ si $s\in K$ i $0$ altrament. El destí, igual amb $t$. Anomenant les arestes per la informació que codifiquen, queda
\[
\renewcommand{\arraystretch}{1.2}
\begin{array}{c|c|l}
\text{aresta} & \text{subgraf de }\mathbf E & \text{de}\to\text{a}\\ \hline
\bot & \varnothing & 0\to0\\
\sigma & \{s\} & 1\to0\\
\tau & \{t\} & 0\to1\\
\partial & \{s,t\} & 1\to1\\
\top & \{s,t,e\} & 1\to1
\end{array}
\]
El graf $\Omega$ té, doncs, dos vèrtexs, un bucle a $0$, una aresta en cada sentit entre $0$ i $1$, i \emph{dos} bucles a $1$ (vegeu també \cite[exemple~15.4.3]{barrwells}). El morfisme $\mathsf{true}\colon\mathbf L\to\Omega$, definit sobre el graf terminal d'un bucle (exemple~\ref{ex:producte-tres-contextos}), envia el vèrtex a $1$ i el bucle a $\top$.

\emph{El morfisme característic d'un subgraf.} Si $H\subseteq G$ és un subgraf, $\chi_H$ envia cada vèrtex a $1$ o a $0$ segons que sigui o no a $H$. Cada aresta $x\colon v\to w$ va a l'aresta de $\Omega$ que registra si $v$, $w$ i $x$ són a $H$: a $\top$ si $x\in H$; a $\partial$ si $v,w\in H$ però $x\notin H$; i a $\sigma$, $\tau$ o $\bot$ segons quins extrems hi siguin. Per exemple, sigui $G$ el cicle de vèrtexs $a,b,c$ i arestes $e\colon a\to b$, $f\colon b\to c$, $k\colon c\to a$, i sigui $H$ el subgraf amb vèrtexs $a,b$ i aresta $e$. Aleshores
\[
\chi_H\colon\quad a,b\mapsto1,\quad c\mapsto0,\qquad e\mapsto\top,\quad f\mapsto\sigma,\quad k\mapsto\tau.
\]
És un morfisme de grafs, perquè $\sigma\colon1\to0$ i $f\colon b\to c$ tenen extrems compatibles, i el mateix per a les altres arestes. A més, el pullback de $\mathsf{true}$ al llarg de $\chi_H$, que es calcula component a component, té per vèrtexs els que van a $1$, $\{a,b\}$, i per arestes les que van a $\top$, $\{e\}$: és exactament $H$. La unicitat també es veu directament. Els vèrtexs de $H$ han d'anar a $1$ i els altres a $0$. Una aresta ha d'anar a una aresta de $\Omega$ amb els extrems ja determinats, i a $\top$ exactament si és a $H$; amb aquestes condicions, en cada cas hi ha una sola opció.

\emph{Més de dos valors de veritat.} La diferència entre $\partial$ i $\top$ és el que fa interessant aquest topos. Al graf $\mathbf E$, el subgraf $\{s\}$ i el subgraf $\{t\}$ són disjunts, i $\{t\}$ és el més gran dels subgrafs disjunts de $\{s\}$, perquè tot subgraf que contingui l'aresta ha de contenir $s$. Però la seva unió és $\{s,t\}$, que no és tot $\mathbf E$: hi falta l'aresta, que és justament el que distingeix $\partial$ de $\top$. A $\Sub(\mathbf E)$, doncs, un predicat i la seva negació no cobreixen tot l'objecte, i el principi del tercer exclòs falla. El classificador d'aquest topos és un graf amb estructura pròpia, no un conjunt $\{0,1\}$. La mateixa categoria que ens ha servit per introduir objectes, functors, Yoneda i límits és, així, un exemple genuí de topos no booleà.
\end{exemple}

\begin{observacio}[$\Omega$ com a objecte de valors de veritat]\label{obs:omega-veritat}\index{Omega@$\Omega$!valors de veritat}\index{topos!booleà}
La noció de topos va aparèixer primer a l'escola de Grothendieck de geometria algebraica, com a generalització de la d'espai topològic. Un dels seus aspectes més fascinants, però, és la relació amb la lògica. Per la representabilitat de $\Sub$, un morfisme $\varphi\colon E\to\Omega$ \emph{és} un predicat sobre $E$, la seva \enquote{extensió} és el subobjecte que classifica, i $\mathsf{true}$ fa el paper del valor \enquote{cert}. En aquest sentit precís ---i només en aquest--- $\Omega$ és un objecte de valors de veritat.

Això no vol dir que $\Omega$ tingui dos valors, ni que la lògica sigui clàssica. En tot topos cada $\Sub(E)$ és una àlgebra de Heyting (nota d'ampliació següent), de manera que la lògica interna és, en general, intuïcionista. Un topos s'anomena \textbf{booleà} quan totes aquestes àlgebres de Heyting són àlgebres de Boole, és a dir, quan per a tot subobjecte $U$ es compleix $U\vee\neg U=\top$. $\cat{Set}$ i $\cat{FinSet}$ són booleans; $\cat{Graph}$ no ho és, com hem vist a l'exemple~\ref{ex:omega-graph}. Fins i tot els \enquote{punts} de $\Omega$, els valors de veritat globals $1\to\Omega$, poden ser més de dos. I, a la inversa, tenir exactament dos valors globals no fa booleà un topos: els valors globals només compten $\Hom(1,\Omega)\cong\Sub(1)$, mentre que la booleanitat és una condició sobre tots els $\Sub(E)$. Els prefeixos sobre el monoide $(\mathbb N,+)$ tenen dos valors globals i no són booleans (vegeu els Exercicis per a tots dos fenòmens). Aquesta lectura permet interpretar la lògica ---de primer ordre i d'ordre superior--- dins de qualsevol topos. El seu desenvolupament és, precisament, l'objecte de l'estudi de la lògica categòrica.
\end{observacio}

\begin{nota}[Ampliació: $\Sub(E)$ és una àlgebra de Heyting]\label{nota:sub-heyting}\index{àlgebra de Heyting!subobjectes d'un topos}
L'afirmació de l'observació anterior és un teorema de la teoria de topos que aquest volum \emph{no demostra}; la fem servir només per situar la qüestió de la booleanitat. L'enunciat és aquest: en un topos elemental, cada ordre parcial $\Sub(E)$ té màxim, mínim, ínfims i suprems binaris i una implicació $\Rightarrow$ adjunta per la dreta de $-\wedge U$, i les reindexacions $f^{*}$ preserven tota aquesta estructura.

Algunes peces ja les tenim. El màxim és $[\id_E]$, i els ínfims són pullbacks (capítol~\ref{cap:subobjectes}). Les altres requereixen resultats que no hem provat. El mínim i els suprems es construeixen amb l'objecte inicial, els coproductes i les imatges; la seva existència fa servir que tot topos té colímits finits i és una categoria regular. La implicació s'obté amb el classificador i els exponencials. A $\cat{Graph}$, les dues construccions es poden verificar a mà sobre els subgrafs, cosa que justifica independentment l'exemple~\ref{ex:omega-graph}. Per a les demostracions, vegeu \cite[cap.~IV]{maclanemoerdijk} i \cite{johnstone}; el desenvolupament correspon a l'estudi de la lògica categòrica.
\end{nota}

\begin{resumcap}
\apartat{Idees centrals}
\begin{enumerate}
  \item El classificador $\Omega$ és un objecte amb $\mathsf{true}\colon 1\to\Omega$ tal que cada subobjecte té un únic morfisme característic, del qual és el pullback; la pertinença d'un element generalitzat esdevé una igualtat de morfismes cap a $\Omega$.
  \item Tenir classificador equival a la representabilitat del functor $\Sub$, i $\Omega$ n'és l'objecte representant: els predicats esdevenen morfismes $E\to\Omega$ i la substitució, precomposició.
  \item Un topos elemental té límits finits, exponencials i classificador; en un topos, $\Omega^{E}$ és l'objecte de parts de $E$.
  \item $\cat{Set}$ és el topos prototípic, amb $\Omega=\{0,1\}$.
  \item Tota categoria de prefeixos sobre una categoria petita és un topos, amb $\Omega$ el prefeix de sedassos; en particular, $\cat{Graph}$ és un topos, amb un classificador de dos vèrtexs i cinc arestes. No tot topos elemental és de prefeixos: $\cat{FinSet}$ no ho és.
  \item $\Omega$ és un objecte de valors de veritat perquè classifica predicats, però la lògica d'un topos és en general intuïcionista: $\cat{Graph}$ no és booleà.
\end{enumerate}
\apartat{Errors típics} Oblidar la hipòtesi de petitesa local\,/\,bona potenciació en escriure $\Sub\colon\cat{C}\op\to\cat{Set}$; confondre topos elemental amb categoria de prefeixos; confondre el classificador amb un conjunt de dos elements (només a $\cat{Set}$ el model prototípic és $\{0,1\}$; en un topos general, $\Omega$ és un objecte de la categoria i pot tenir una estructura molt més rica); i suposar que la lògica interna de qualsevol topos és clàssica: en general és intuïcionista, i la booleanitat és una propietat addicional.
\apartat{Lectura recomanada} \cite[cap.~8]{awodey}, especialment §8.8, per a classificadors, topos de prefeixos i la connexió amb Yoneda; \cite[cap.~15]{barrwells} per a una presentació alternativa i l'exemple dels grafs; \cite[cap.~1--2]{maclanemoerdijk} per continuar cap a la teoria de topos i feixos; \cite[§§12 i 14]{streicher} per als topos elementals i exercicis amb topos de prefeixos; \cite[article~VI, sessions~32--33]{lawvereschanuel} per a una presentació molt elemental dels subobjectes, la veritat i els topos; \cite{johnstone} com a obra de consulta avançada.
\end{resumcap}

\chapter[Pont cap a la lògica categòrica]{\texorpdfstring{Pont cap a la lògica\\ categòrica}{Pont cap a la lògica categòrica}}
\label{cap:pont}

\begin{objectius}
\apartat{Prerequisits} Tot el volum, especialment subobjectes, reindexació i categories regulars (cap.~\ref{cap:subobjectes}) i els classificadors (cap.~\ref{cap:topos}).
\apartat{Funció del capítol} Aquest capítol és un epíleg sense demostracions ni exercicis. No desenvolupa una teoria nova, però introdueix la terminologia mínima necessària per dibuixar el mapa que porta de les construccions d'aquest volum a la lògica categòrica. Serveix d'orientació per a qui vulgui continuar cap al volum següent.
\end{objectius}

Aquest volum tanca aquí. Els capítols anteriors han construït, pas a pas, el llenguatge categòric bàsic necessari per al programa que anunciem: categories i morfismes, functors, transformacions naturals i equivalències, representabilitat i el lema de Yoneda, límits i colímits, adjuncions, tancament cartesià, subobjectes, imatges i factoritzacions, i els classificadors de subobjectes. Aquest epíleg mostra què d'aquest bagatge es prolonga en l'estudi de la \textbf{lògica categòrica}, i com: cada eina apresa aquí té una contrapartida lògica precisa a l'altra banda del pont.

\section{El que aquest volum deixa preparat}

Convé fer inventari del que ja tenim, perquè és exactament la infraestructura que la lògica categòrica pressuposa:
\begin{itemize}
\item \textbf{Categories, functors i naturalitat}: el llenguatge bàsic en què s'expressa tota la resta.
\item \textbf{La categoria prima de la deduïbilitat}: la primera aparició, al capítol inicial, de la idea que una lògica es pot veure com una categoria.
\item \textbf{Representabilitat i Yoneda}: la manera d'identificar un objecte per com es relaciona amb tots els altres; el classificador de subobjectes n'és un cas.
\item \textbf{Límits finits, i molt especialment els pullbacks}: la base per interpretar contextos, substitució i relacions.
\item \textbf{Adjuncions}: el patró que reapareixerà en operacions lògiques centrals, especialment en la implicació i en els quantificadors existencial i universal.
\item \textbf{Tancament cartesià i càlcul lambda}: la semàntica dels tipus i els termes.
\item \textbf{Subobjectes i reindexació}: la interpretació dels predicats i la substitució.
\item \textbf{Imatges i factoritzacions}: la construcció que, a l'altra banda, interpretarà la quantificació existencial.
\item \textbf{Classificadors de subobjectes i topos}: el marc on la lògica esdevé interna a la categoria.
\end{itemize}

Amb tot això, el salt cap a la lògica categòrica deixa de ser un salt: és una continuació natural.

\section{El mapa: de regular a topos}\label{sec:mapa}\index{categoria!regular}\index{categoria!coherent}\index{categoria!de Heyting}

Abans de descriure el recorregut, val la pena fixar un mapa que ordena, d'una vegada, quina estructura categòrica interpreta quin fragment lògic. És la classificació que dona sentit retrospectiu a les distincions que hem anat marcant ---per exemple, per què la regularitat basta per a l'existencial però no per a la disjunció ni per a l'universal.

La idea directriu és que cada nivell \emph{afegeix} una peça d'estructura a les fibres de subobjectes $\Sub(A)$, i que cada peça interpreta un connectiu o quantificador:
\begin{center}
\renewcommand{\arraystretch}{1.3}
\noindent\begin{tabularx}{\linewidth}{@{}>{\raggedright\arraybackslash}X >{\raggedright\arraybackslash}X >{\raggedright\arraybackslash}X@{}}
\hline
\textbf{Estructura categòrica} & \textbf{Nova estructura a $\Sub(A)$} & \textbf{Fragment lògic}\\
\hline
categoria amb límits finits & $\top$, $\wedge$ (interseccions) & conjunció, veritat\\
categoria \textbf{regular} & imatges estables: $\exists_f$ & $+\ \exists$\\
categoria \textbf{coherent} & joins finits estables: $\bot$, $\vee$ & $+\ \vee,\ \bot$\\
categoria \textbf{de Heyting} & implicació $\Rightarrow$; $\forall_f$ & $+\ \Rightarrow,\ \forall$: lògica intuïcionista de primer ordre\\
\textbf{topos} elemental & classificador $\Omega$, objectes de parts, exponencials & lògica intuïcionista d'ordre superior\\
\hline
\end{tabularx}
\renewcommand{\arraystretch}{1.0}
\end{center}

La lectura de la taula, de dalt a baix, és acumulativa i respon les preguntes que aquest volum ha deixat obertes deliberadament:
\begin{itemize}
  \item \textbf{Regular} ($\top,\wedge,\exists$). La factorització imatge estable dona l'existencial com a adjunt per l'esquerra de la reindexació. És fins on arriba el capítol de subobjectes: per això hi vam construir només $\exists_f$, i no $\forall_f$ ni la disjunció.
  \item \textbf{Coherent} ($+\vee,\bot$). S'hi afegeixen \emph{joins finits estables} a les fibres, que interpreten la disjunció i la falsedat. Aquesta és l'estructura que a la teoria de subobjectes anunciàvem com a necessària per a les \enquote{unions}, i que la regularitat sola no proporciona.
  \item \textbf{De Heyting} ($+\Rightarrow,\forall$). Les fibres esdevenen àlgebres de Heyting ---amb implicació, adjunta per la dreta de la conjunció--- i cada reindexació $f^{*}$ adquireix un adjunt per la dreta $\forall_f$. Apareixen la implicació i el quantificador universal.
  \item \textbf{Topos} elemental. El classificador $\Omega$ i els exponencials clouen el quadre amb una lògica intuïcionista d'ordre superior, on fins i tot els predicats sobre predicats tenen extensió.
\end{itemize}

\index{categoria!localment cartesiana tancada}En paral·lel, la semàntica dels \emph{tipus} segueix una altra línia. Les categories cartesianes tancades interpreten el càlcul lambda simplement tipat (capítol~\ref{cap:cartesianes}), mentre que les categories \textbf{localment cartesianes tancades} proporcionen productes dependents: per a cada $f\colon A\to B$, la reindexació $f^{*}\colon\cat{C}/B\to\cat{C}/A$ admet un adjunt per la dreta $\Pi_f$. Aquesta via és especialment important per a la teoria de tipus dependents. No és un esglaó de la cadena regular--coherent--Heyting, però totes dues vies es troben en els topos elementals, que són localment cartesians tancats i tenen lògica de Heyting als subobjectes.

Aquest mapa no és matèria d'aquest volum (no en demostrem les equivalències), però és la brúixola que orienta tot el que segueix: cada fila és, de fet, un capítol del programa de la lògica categòrica.

\section{El que continua a l'altra banda}

L'estudi de la lògica categòrica pren aquestes eines i les fa servir per construir una \textbf{semàntica de la lògica} de complexitat creixent. Sense entrar-hi ---perquè és l'objecte d'un estudi a part---, en dibuixem el recorregut, perquè el lector vegi on desemboca cada cosa que ha après:

\begin{description}
\item[Lògica d'equacions.] El punt de partida: teories algebraiques a l'estil de Lawvere, on la categoria sintàctica té productes finits i els models són functors que els preserven. És la lògica que ja s'endevina en la universalitat dels productes finits: els contextos es combinen per producte i les operacions esdevenen morfismes entre potències finites d'un objecte. El tancament cartesià pertany a una via posterior, la del càlcul lambda i la implicació.

\item[Lògica de límits finits.] El pont immediat: teories amb diverses menes, operacions parcialment definides i axiomes d'existència i unicitat. La categoria sintàctica passa a tenir \emph{límits finits}, no només productes; és el marc de les teories essencialment algebraiques.

\item[Lògica regular.] Fórmules amb igualtat, conjunció i quantificació existencial, escrites com a seqüents en context. Aquí és on les \textbf{imatges regulars} construïdes en aquest volum interpreten l'existencial, i on la factorització imatge interpreta l'eliminació de testimonis. El resultat central és que, per a una teoria regular $T$ i una categoria regular $\cat{C}$, els models de $T$ en $\cat{C}$ es descriuen mitjançant functors regulars des de la categoria sintàctica regular $\cat{C}_T$ cap a $\cat{C}$; amb la noció adequada de morfisme de models, aquesta correspondència s'organitza categòricament.

\item[Lògica coherent i geomètrica.] S'hi afegeixen disjuncions finites (coherent) i disjuncions indexades per conjunts (geomètrica). És el marc que condueix als feixos i, de nou, als topos.

\item[Lògica intuïcionista de primer ordre.] Apareixen la implicació i el quantificador universal. La idea organitzadora és que les fibres de subobjectes formen \textbf{àlgebres de Heyting}: la implicació és l'adjunt dret de la conjunció, i $\forall$ és l'adjunt dret de la substitució. Els adjunts de la reindexació que aquí hem vist com a exemple es converteixen allà en els quantificadors del sistema.

\item[Doctrines i fibracions lògiques.] La síntesi: contextos, substitució, predicats i quantificadors s'organitzen en una \emph{doctrina indexada}. Les condicions de coherència ---Beck--Chevalley i Frobenius--- expressen que substituir i quantificar interactuen correctament. És el llenguatge que unifica tots els casos anteriors.

\item[Teories classificadores i llenguatge intern.] El punt culminant: la categoria o el topos classificador d'una teoria, amb la seva propietat universal, i el llenguatge intern que permet raonar dins d'una categoria com si els seus objectes fossin conjunts, sempre que les regles emprades corresponguin a l'estructura disponible.

\item[Lògica clàssica i booleanització.] Com a epíleg, què s'afegeix en imposar el tercer exclòs, quines construccions deixen de ser constructives i com, en teoria de topos, es passa a la booleanització d'un topos. La lògica clàssica apareix així com una especialització, no com el rerefons silenciós de tota la teoria.
\end{description}

\begin{nota}[Ampliació: estructures sense conjunt subjacent]
Al capítol~\ref{cap:categories} hem partit d'estructures que tenen un conjunt subjacent. No totes les estructures matemàtiques són conjunts amb una estructura afegida, i el programa anterior ho fa visible. A cada teoria geomètrica li correspon un \emph{topos classificador}, la categoria de punts del qual és equivalent a la categoria de models de la teoria a $\cat{Set}$ \cite{caramello}. Però aquest \enquote{espai} pot ser no trivial i no tenir cap punt. Per exemple, la teoria d'una aplicació exhaustiva $\mathbb N\to\mathbb R$ no té cap model a $\cat{Set}$, pel teorema de Cantor; tanmateix, el seu \emph{locale} classificador és no trivial, i el topos de feixos sobre aquest locale, que n'és el topos classificador, no té cap punt \cite{blechschmidt}. Aquestes situacions queden fora d'aquest volum; les esmentem perquè mostren que la perspectiva de les fletxes no depèn de disposar d'elements.
\end{nota}

\section{Un exemple: la categoria sintàctica d'una teoria}\label{sec:sintactica}\index{categoria!sintàctica}

Fins ara la correspondència lògica--categòrica ha aparegut repartida en exemples solts. Val la pena reunir-la en un únic exemple pont, complet i concret, que il·lustra de cop l'eslògan \enquote{una teoria es codifica en una categoria sintàctica}. Prenem la teoria algebraica més simple no trivial: una \textbf{mena} $s$ i una operació binària $m\colon s\times s\to s$, sense cap axioma (la teoria dels \emph{magmes}). Construïm pas a pas la seva categoria sintàctica $\cat{Syn}$.

\begin{enumerate}
  \item \textbf{Contextos com a objectes.} Un \emph{context} és una llista finita de variables de la mena $s$, que escrivim $[x_1,\dots,x_n]$ i abreugem $s^n$. Els objectes de $\cat{Syn}$ són els contextos, un per a cada $n\geq 0$; el context buit $s^0$ és l'objecte terminal.
  \item \textbf{Termes com a morfismes.} Un morfisme $s^n\to s^m$ és una llista de $m$ termes $(t_1,\dots,t_m)$, cadascun construït amb les variables $x_1,\dots,x_n$ i l'operació $m$ (per exemple, $m(x_1,m(x_2,x_1))$). La composició és la \textbf{substitució}: components d'un morfisme dins d'un altre. La identitat de $s^n$ és la llista de variables $(x_1,\dots,x_n)$.
  \item \textbf{Productes com a concatenació de contextos.} El producte $s^n\times s^m$ és el context $s^{n+m}$: juxtaposar dos contextos és formar-ne el producte, i les projeccions obliden les variables sobrants. Així $\cat{Syn}$ té tots els productes finits, i l'operació $m$ del llenguatge és, ara, un morfisme distingit $s\times s\to s$ de la categoria.
  \item \textbf{Equacions com a igualtats de morfismes.} Si la teoria tingués axiomes ---posem l'associativitat de $m$---, s'imposarien identificant els dos morfismes $s^3\to s$ corresponents a $m(m(x_1,x_2),x_3)$ i $m(x_1,m(x_2,x_3))$. Una equació entre termes \emph{és}, literalment, una igualtat entre morfismes de $\cat{Syn}$.
  \item \textbf{Models com a functors.} Un model de la teoria en una categoria $\cat{C}$ amb productes finits ---per exemple, un magma ordinari a $\cat{Set}$--- és exactament un functor $\cat{Syn}\to\cat{C}$ que preserva els productes finits: tria un objecte per a $s$ i un morfisme per a $m$, respectant tota la sintaxi. Els homomorfismes de models són transformacions naturals. Aquesta és la forma precisa de \enquote{els models són functors}.
  \item \textbf{Fórmules i existencial.} Per passar de la teoria algebraica anterior a la lògica regular cal enriquir també la categoria sintàctica: la categoria $\cat{Syn}$ dels passos anteriors no conté automàticament els subobjectes que representen totes les fórmules. En la \textbf{categoria sintàctica regular} $\cat{C}_T$ d'una teoria amb predicats, una fórmula $\varphi(\bar x)$ en context determina, llevat d'equivalència demostrable, un subobjecte
  \[
  [\bar x\mid\varphi]\rightarrowtail[\bar x\mid\top].
  \]
  La substitució al llarg d'un morfisme de contextos s'interpreta per \textbf{reindexació}, és a dir, per pullback. En particular, si $\pi\colon s^{n+1}\to s^n$ és la projecció que oblida la darrera variable, $\pi^{*}$ és el debilitament, que considera un predicat sobre $s^n$ en el context ampliat. Si $P\rightarrowtail s^{n+1}$ representa $\varphi(\bar x,y)$, la \textbf{imatge} del compost
  \[
  P\rightarrowtail s^{n+1}\xrightarrow{\ \pi\ }s^n
  \]
  representa $\exists y\,\varphi(\bar x,y)$. Així, l'existencial és la imatge d'un predicat al llarg d'una projecció o, equivalentment, l'adjunt per l'esquerra de la reindexació (capítol~\ref{cap:subobjectes}).
  \item \textbf{Un càlcul complet en un model.} Fem concret tot l'anterior en un sol model, el magma $(\mathbb N,+)$, és a dir, el functor $M\colon\cat{Syn}\to\cat{Set}$ amb $M(s)=\mathbb N$ i $M(m)=+$. Prenem la fórmula $\varphi(x,y)$: $m(y,y)=x$, en el context $[x,y]$.
  \begin{itemize}
    \item \emph{Una equació és un equalitzador.} Els dos termes $m(y,y)$ i $x$ són morfismes $s^2\to s$, interpretats com les aplicacions $(x,y)\mapsto y+y$ i $(x,y)\mapsto x$ de $\mathbb N^2$ a $\mathbb N$. La fórmula s'interpreta com el seu equalitzador,
    \[
    P=\{(x,y)\in\mathbb N^2\mid y+y=x\}\rightarrowtail\mathbb N^2.
    \]
    Equivalentment, $P$ és el pullback de la diagonal $\delta_{\mathbb N}\colon\mathbb N\rightarrowtail\mathbb N\times\mathbb N$ al llarg del parell format pels dos termes: la igualtat és la diagonal.
    \item \emph{L'existencial és una imatge.} La projecció $\pi\colon\mathbb N^2\to\mathbb N$ oblida $y$. La imatge de $P\rightarrowtail\mathbb N^2\to\mathbb N$ és $\{x\mid\exists y\ y+y=x\}$, el conjunt dels nombres parells. Així, $\exists y\,(m(y,y)=x)$ s'interpreta com el predicat \enquote{$x$ és parell}, és a dir, com $\exists_\pi[P]$.
    \item \emph{El predicat és un morfisme cap a $\Omega$.} El seu morfisme característic $\mathbb N\to\{0,1\}$ val $1$ exactament als parells, com a l'exemple~\ref{ex:caracteristic-set}.
    \item \emph{La substitució és un pullback.} Substituir $x$ pel terme $m(z,m(z,z))$, en el context $[z]$, dona la fórmula $\exists y\,(m(y,y)=m(z,m(z,z)))$. Categòricament, és reindexar el predicat al llarg del morfisme $s\to s$ definit pel terme, interpretat com $z\mapsto 3z$. Per tant, s'interpreta com $\{z\mid 3z\text{ és parell}\}$, que és el conjunt dels parells. En canvi, substituir pel terme $m(z,z)$ dona $\{z\mid 2z\text{ és parell}\}=\mathbb N$, el subobjecte màxim: la fórmula $\exists y\,(m(y,y)=m(z,z))$ és demostrable, amb el testimoni $y=z$, i el model ho reflecteix.
  \end{itemize}
  Un sol exemple recorre així les peces del volum: termes com a morfismes, equacions com a equalitzadors i igualtat com a diagonal (capítols~\ref{cap:limits} i~\ref{cap:cartesianes}), existencial com a imatge (capítol~\ref{cap:subobjectes}), predicats com a morfismes cap a $\Omega$ (capítol~\ref{cap:topos}) i substitució com a pullback.
\end{enumerate}

Amb aquest exemple, cada peça del llibre troba el seu lloc: els contextos són objectes, els termes són morfismes, la substitució és composició i reindexació, els productes són concatenacions, les equacions són igualtats de fletxes, els models són functors que preserven estructura, i els quantificadors, quan existeixen, són adjunts de la reindexació. La categoria sintàctica és el punt on la sintaxi lògica i la semàntica categòrica es toquen; construir-la en general, amb l'estructura exacta que cada fragment lògic requereix, és una de les primeres tasques del volum següent.

\section{Cloenda}

El fil que recorre tot aquest programa és un de sol, i ja el coneixem: \emph{traduir entre lògica i categories}. En el marc categòric adequat a cada fragment lògic, un predicat s'interpreta com un subobjecte; la substitució, com un pullback; la conjunció, com una intersecció; l'existencial, com una imatge; la implicació i l'universal, mitjançant adjunts; una teoria es codifica en una categoria sintàctica i els seus models es descriuen per functors que preserven l'estructura pertinent. Les teories algebraiques, regulars, coherents o geomètriques tenen categories sintàctiques amb l'estructura corresponent, i per això teories sintàcticament diferents poden tenir semàntiques categòriques equivalents, o fins i tot el mateix topos classificador. Cadascuna d'aquestes traduccions descansa sobre una construcció que aquest volum ha desenvolupat amb detall.

No hem demostrat cap teorema de la lògica categòrica en aquestes pàgines, i era deliberat: la seva demostració pertany al volum següent. El que sí que hem fet és deixar el terreny preparat, de manera que quan aquells teoremes arribin, cap de les seves peces categòriques resulti desconeguda. Si el lector, en trobar més endavant que \enquote{l'existencial és un adjunt per l'esquerra de la substitució, estable per canvi de base}, reconeix en aquesta frase la factorització imatge, l'adjunció i el pullback que ha estudiat aquí, aleshores aquest volum haurà complert la seva funció.

\section{Mapa del volum següent}\label{sec:mapa-volum-seguent}

La secció~\ref{sec:mapa} ordenava la continuació per fragments lògics. Per acabar, l'ordenem per temes: per a cadascun dels cinc temes centrals del volum de lògica categòrica, què en queda establert aquí i què s'hi afegirà.
\begin{center}
\footnotesize
\renewcommand{\arraystretch}{1.3}
\noindent\begin{tabularx}{\linewidth}{@{}>{\raggedright\arraybackslash}p{0.2\linewidth}>{\raggedright\arraybackslash}X>{\raggedright\arraybackslash}X@{}}
\toprule
\textbf{Tema} & \textbf{En aquest volum} & \textbf{Al volum següent}\\
\midrule
substitució i reindexació & $f^{*}$ per pullback a $\cat{C}/B$ (capítol~\ref{cap:limits}) i a $\Sub(B)$; el functor $\Sub\colon\cat{C}\op\to\cat{Poset}$ (capítol~\ref{cap:subobjectes}) & la substitució de termes en fórmules com a reindexació en una doctrina; la condició de Beck--Chevalley\\
connectius sobre subobjectes & $\top$ i $\wedge$ com a pullbacks (capítol~\ref{cap:subobjectes}); $\Rightarrow$ als preordres cartesians tancats (capítol~\ref{cap:cartesianes}); $\wedge$ sobre $\Omega$ (capítol~\ref{cap:topos}) & $\bot$ i $\vee$ estables (categories coherents); $\Rightarrow$ a les fibres (categories de Heyting)\\
quantificadors com a adjunts & $\exists_f\dashv f^{*}\dashv\forall_f$ a $\cat{Set}$ (capítol~\ref{cap:adjuncions}); $\exists_f$ com a imatge en categories regulars (capítol~\ref{cap:subobjectes}) & $\forall_f$ en categories de Heyting; la condició de Frobenius; la quantificació al llarg de projeccions de contextos\\
igualtat & la diagonal $\delta_A$ (capítol~\ref{cap:cartesianes}); la parella nucli, que és el pullback de $\delta_B$ al llarg de $f\times f$ (capítol~\ref{cap:subobjectes}); les equacions com a equalitzadors (secció~\ref{sec:sintactica}) & la igualtat com a $\exists_{\delta_A}(\top)$, on $\exists_{\delta_A}$ és l'adjunt per l'esquerra de la contracció $\delta_A^{*}$; les regles de la igualtat com a propietats d'aquesta adjunció\\
semàntica interna & $\Omega$, la pertinença per elements generalitzats i l'objecte de parts (capítol~\ref{cap:topos}); l'exemple de la secció~\ref{sec:sintactica} & el llenguatge intern d'un topos, la semàntica de Kripke--Joyal, la lògica d'ordre superior, i la correcció i la completesa\\
\bottomrule
\end{tabularx}
\end{center}
Cap fila de la columna central no demana res que el lector no hagi vist en aquest volum. La columna de la dreta és el programa del volum següent.

\begin{resumcap}
\apartat{El fil en una frase} En el marc categòric adequat a cada fragment lògic, un predicat s'interpreta com un subobjecte; la substitució, com un pullback; la conjunció, com una intersecció; l'existencial, com una imatge; la implicació i l'universal, mitjançant adjunts; una teoria es codifica en una categoria sintàctica i els seus models es descriuen per functors que preserven l'estructura pertinent.
\apartat{Cap on continua} La lògica pròpiament dita ---sintaxi, correcció i completesa, i les condicions de coherència Beck--Chevalley i Frobenius--- es reserva per al volum de lògica categòrica, del qual aquest text és la preparació. La secció~\ref{sec:mapa-volum-seguent} en dona el mapa per temes: substitució, connectius, quantificadors, igualtat i semàntica interna.
\apartat{Lectura recomanada} \cite{lambekscott} i \cite{jacobs} per a la lògica categòrica i la teoria de tipus; \cite[§13]{streicher} per a la lògica interna dels topos; \cite{makkaireyes} i \cite{pitts} per a la lògica de primer ordre en categories regulars, coherents i de Heyting, que és el punt de partida del volum de lògica; \cite{maclanemoerdijk} per a la lògica geomètrica i els topos; \cite{johnstone} com a obra de consulta. Com a direccions complementàries: \cite{hott}, per a la teoria homotòpica de tipus, on la igualtat entre dos elements pot tenir estructura pròpia; i l'nLab \cite{nlab-classificador,nlab-locale}, com a recurs de consulta sobre topos classificadors i locales.
\end{resumcap}

\fi

% ============================================================
% PART II — PRÀCTICA
% Els capítols reprodueixen els noms i l'ordre de la part teòrica.
% ============================================================
\part{Pràctica}
\setcounter{chapter}{0}
\chapter{Categories}

En aquest capítol resolem amb detall una selecció d'exercicis corresponents al capítol \emph{Categories} de la part de Teoria. L'objectiu no és només arribar al resultat, sinó mostrar \emph{com} es raona en el llenguatge categòric: comprovar axiomes, seguir camins, traduir diagrames a equacions i reconèixer una mateixa estructura en conjunts, preordres, lògica i grafs. Convé intentar cada exercici abans de llegir-ne la solució.

\section{Definició de categoria}

\begin{exercici}\label{exr:pr-rel}
Comprova que els conjunts amb les relacions formen una categoria $\cat{Rel}$, on un morfisme $A\to B$ és una relació $R\subseteq A\times B$, la identitat de $A$ és la relació d'igualtat, i la composició de $R\subseteq A\times B$ amb $S\subseteq B\times C$ és
\[
S\comp R=\{(a,c)\in A\times C\mid \exists\,b\in B\text{ tal que }(a,b)\in R\text{ i }(b,c)\in S\}.
\]
\end{exercici}

\begin{solucio}
Hem de comprovar els dos axiomes. Comencem per les \textbf{identitats}. Sigui $\id_A=\{(a,a)\mid a\in A\}$ la relació d'igualtat sobre $A$. Per a tota relació $R\subseteq A\times B$,
\[
R\comp\id_A
=\{(a,b)\mid \exists\,a'\ (a,a')\in\id_A\text{ i }(a',b)\in R\}
=R,
\]
perquè $(a,a')\in\id_A$ obliga $a'=a$. Anàlogament $\id_B\comp R=R$.

Per a l'\textbf{associativitat}, siguin $R\subseteq A\times B$, $S\subseteq B\times C$ i $T\subseteq C\times D$. Un parell $(a,d)$ pertany a $T\comp(S\comp R)$ si i només si existeixen $b\in B$ i $c\in C$ tals que
\[
(a,b)\in R,
\qquad
(b,c)\in S,
\qquad
(c,d)\in T.
\]
La mateixa condició caracteritza $(T\comp S)\comp R$. Per tant
\[
T\comp(S\comp R)=(T\comp S)\comp R,
\]
i $\cat{Rel}$ és una categoria.
\end{solucio}

\section{Fletxes, camins i diagrames commutatius}

\begin{exercici}
Llegeix el diagrama següent i escriu l'equació que expressa que commuta:
\[
\begin{tikzpicture}[baseline=(m.center)]
  \matrix (m) [matrix of math nodes, row sep=2.7em, column sep=3.6em] {
    A & B \\
    C & D \\
  };
  \draw[-{Stealth[scale=1.1]}] (m-1-1) -- node[above] {$f$} (m-1-2);
  \draw[-{Stealth[scale=1.1]}] (m-1-1) -- node[left] {$u$} (m-2-1);
  \draw[-{Stealth[scale=1.1]}] (m-1-2) -- node[right] {$v$} (m-2-2);
  \draw[-{Stealth[scale=1.1]}] (m-2-1) -- node[below] {$g$} (m-2-2);
\end{tikzpicture}
\]
Explica també per què l'ordre de les composicions és el que escrius.
\end{exercici}

\begin{solucio}
Hi ha dos camins dirigits de $A$ a $D$. El camí superior-dret és
\[
A\xrightarrow{f}B\xrightarrow{v}D,
\]
i per tant determina $v\comp f$. El camí esquerre-inferior és
\[
A\xrightarrow{u}C\xrightarrow{g}D,
\]
i determina $g\comp u$. El quadrat commuta exactament quan
\[
\boxed{v\comp f=g\comp u}.
\]
L'ordre de cada composició es llegeix de dreta a esquerra: en $v\comp f$ seguim primer $f$ i després $v$.
\end{solucio}

\begin{exercici}
Considera a $\cat{Set}$ quatre còpies de $\mathbb Z$ i les aplicacions
\[
f(n)=n+1,
\qquad
u(n)=2n,
\qquad
v(n)=2n,
\qquad
g(n)=n+2.
\]
Comprova que el quadrat
\[
\begin{tikzpicture}[baseline=(m.center)]
  \matrix (m) [matrix of math nodes, row sep=2.8em, column sep=4.2em] {
    \mathbb Z & \mathbb Z \\
    \mathbb Z & \mathbb Z \\
  };
  \draw[-{Stealth[scale=1.05]}] (m-1-1) -- node[above] {$f$} (m-1-2);
  \draw[-{Stealth[scale=1.05]}] (m-1-1) -- node[left] {$u$} (m-2-1);
  \draw[-{Stealth[scale=1.05]}] (m-1-2) -- node[right] {$v$} (m-2-2);
  \draw[-{Stealth[scale=1.05]}] (m-2-1) -- node[below] {$g$} (m-2-2);
\end{tikzpicture}
\]
commuta.
\end{exercici}

\begin{solucio}
Cal comprovar que $v\comp f=g\comp u$. Per a qualsevol $n\in\mathbb Z$,
\[
(v\comp f)(n)=v(n+1)=2n+2,
\]
mentre que
\[
(g\comp u)(n)=g(2n)=2n+2.
\]
Les dues aplicacions coincideixen en tots els enters. Per tant
\[
v\comp f=g\comp u,
\]
i el quadrat commuta. Aquest exemple mostra com una afirmació sobre funcions es pot llegir simultàniament com una igualtat algebraica i com la commutativitat d'un diagrama.
\end{solucio}

\section{Primers exemples}

\begin{exercici}
Sigui $T$ un sistema deductiu amb una relació de deduïbilitat $\vdash_T$ reflexiva i transitiva. Comprova amb detall que les fórmules, amb una única fletxa $A\to B$ quan $A\vdash_T B$, formen una categoria prima $\cat{Prov}_T$.
\end{exercici}

\begin{solucio}
Els \textbf{objectes} són les fórmules. Entre dues fórmules $A$ i $B$ posem un únic morfisme $A\to B$ exactament quan $A\vdash_T B$.

La \textbf{identitat} existeix perquè la deduïbilitat és reflexiva:
\[
A\vdash_T A.
\]
Per tant hi ha una fletxa $\id_A\colon A\to A$.

La \textbf{composició} existeix perquè la deduïbilitat és transitiva. Si
\[
A\vdash_T B
\qquad\text{i}\qquad
B\vdash_T C,
\]
aleshores
\[
A\vdash_T C.
\]
Diagramàticament,
\[
\begin{tikzpicture}[baseline=(m.center)]
  \matrix (m) [matrix of math nodes, column sep=3.4em] { A & B & C \\ };
  \draw[-{Stealth[scale=1.05]}] (m-1-1) -- (m-1-2);
  \draw[-{Stealth[scale=1.05]}] (m-1-2) -- (m-1-3);
  \draw[-{Stealth[scale=1.0]}] (m-1-1) to[bend right=25] (m-1-3);
\end{tikzpicture}
\]

Finalment, entre dues fórmules hi ha com a màxim una fletxa. Això fa que les igualtats exigides pels axiomes d'associativitat i identitat siguin automàtiques sempre que les composicions existeixin. Per tant $\cat{Prov}_T$ és una categoria prima.

Cal remarcar què hem oblidat: si hi ha dues demostracions diferents de $B$ a partir d'$A$, aquesta categoria les identifica en una sola fletxa. Només conserva la relació de deduïbilitat.
\end{solucio}

\begin{exercici}
Siguin els grafs dirigits $G$ i $H$ següents:
\[
G:\quad 0\xrightarrow{a}1,
\qquad
H:\quad
\begin{tikzpicture}[baseline=(m.center)]
  \matrix (m) [matrix of math nodes, row sep=2.0em, column sep=2.8em] {
    & y \\
    x & z \\
  };
  \draw[-{Stealth[scale=1.0]}] (m-2-1) -- node[above left] {$r$} (m-1-2);
  \draw[-{Stealth[scale=1.0]}] (m-2-1) -- node[below] {$s$} (m-2-2);
\end{tikzpicture}
\]
Defineix $F_V(0)=x$, $F_V(1)=y$ i $F_E(a)=r$. Comprova que això defineix un morfisme de grafs $F\colon G\to H$. Explica per què substituir $F_E(a)=r$ per $F_E(a)=s$, mantenint $F_V$, no defineix un morfisme de grafs.
\end{exercici}

\begin{solucio}
En $G$ tenim
\[
s_G(a)=0,
\qquad
t_G(a)=1,
\]
i en $H$,
\[
s_H(r)=x,
\qquad
t_H(r)=y.
\]
Per tant
\[
F_V(s_G(a))=F_V(0)=x=s_H(r)=s_H(F_E(a))
\]
i
\[
F_V(t_G(a))=F_V(1)=y=t_H(r)=t_H(F_E(a)).
\]
Així es compleixen les dues equacions
\[
F_V\comp s_G=s_H\comp F_E,
\qquad
F_V\comp t_G=t_H\comp F_E,
\]
i $F$ és un morfisme de grafs.

Si poséssim $F_E(a)=s$, l'origen encara seria correcte perquè $s_H(s)=x$, però el destí fallaria:
\[
t_H(s)=z\neq y=F_V(1).
\]
Per tant el quadrat corresponent al destí no commutaria i no tindríem un morfisme de grafs.
\end{solucio}

\section{Isomorfismes}

\begin{exercici}
En una categoria qualsevol, comprova que si $f\colon A\to B$ i $g\colon B\to C$ tenen inversos, aleshores $g\comp f$ també en té, i
\[
(g\comp f)^{-1}=f^{-1}\comp g^{-1}.
\]
\end{exercici}

\begin{solucio}
La situació és
\[
A\xrightarrow{f}B\xrightarrow{g}C.
\]
Per recórrer aquest camí en sentit contrari hem d'invertir també l'ordre: primer $g^{-1}$ i després $f^{-1}$. Proposem, doncs, $f^{-1}\comp g^{-1}$ com a invers.

D'una banda,
\[
(f^{-1}\comp g^{-1})\comp(g\comp f)
=f^{-1}\comp(g^{-1}\comp g)\comp f
=f^{-1}\comp f
=\id_A,
\]
i de l'altra,
\[
(g\comp f)\comp(f^{-1}\comp g^{-1})
=g\comp(f\comp f^{-1})\comp g^{-1}
=g\comp g^{-1}
=\id_C.
\]
Per tant $(g\comp f)^{-1}=f^{-1}\comp g^{-1}$.
\end{solucio}

\begin{exercici}
Comprova que a $\cat{Rel}$ una relació $R\colon A\to B$ és un isomorfisme si i només si és el graf d'una bijecció de $A$ a $B$.
\end{exercici}

\begin{solucio}
Si $f\colon A\to B$ és una bijecció i $R$ és el seu graf, el graf de $f^{-1}$ és una relació $B\to A$ i les dues composicions són les relacions d'igualtat. Per tant $R$ és un isomorfisme.

Recíprocament, suposem que $R$ té un invers $S\colon B\to A$, és a dir,
\[
S\comp R=\id_A,
\qquad
R\comp S=\id_B.
\]
Per a cada $a\in A$, com $(a,a)\in S\comp R$, existeix algun $b\in B$ amb $(a,b)\in R$ i $(b,a)\in S$. Així cada $a$ es relaciona amb almenys un $b$.

Si $(a,b)\in R$ i $(a,b')\in R$, triem $b_0$ amb $(a,b_0)\in R$ i $(b_0,a)\in S$. Llavors
\[
(b_0,b)\in R\comp S=\id_B,
\qquad
(b_0,b')\in R\comp S=\id_B,
\]
i per tant $b=b_0=b'$. Així cada $a$ es relaciona amb un únic $b$, i $R$ és el graf d'una aplicació $f\colon A\to B$.

Aplicat a $S$, el mateix argument dona una aplicació $g\colon B\to A$. Les igualtats $S\comp R=\id_A$ i $R\comp S=\id_B$ es converteixen en
\[
g\comp f=\id_A,
\qquad
f\comp g=\id_B.
\]
Per tant $f$ és bijectiva i $R$ n'és el graf.
\end{solucio}

\begin{exercici}
Sigui $\cat{C}$ una categoria. Comprova que la relació \enquote{ésser isomorf} entre objectes de $\cat{C}$ és una relació d'equivalència.
\end{exercici}

\begin{solucio}
\textbf{Reflexivitat}: $\id_A$ és el seu propi invers, de manera que $A\cong A$.

\textbf{Simetria}: si $f\colon A\to B$ és un isomorfisme amb invers $f^{-1}\colon B\to A$, aleshores $f^{-1}$ és també un isomorfisme, amb invers $f$; per tant $B\cong A$.

\textbf{Transitivitat}: si $f\colon A\to B$ i $g\colon B\to C$ són isomorfismes, la composició $g\comp f$ també ho és pel primer exercici de la secció; per tant $A\cong C$.
\end{solucio}

\begin{exercici}
A la categoria prima $\cat{Prov}_T$, comprova que dues fórmules $A$ i $B$ són isomorfes si i només si són interdeduïbles:
\[
A\vdash_T B
\qquad\text{i}\qquad
B\vdash_T A.
\]
\end{exercici}

\begin{solucio}
Si $A\cong B$, per definició existeixen fletxes
\[
A\longrightarrow B
\qquad\text{i}\qquad
B\longrightarrow A.
\]
A $\cat{Prov}_T$ això significa exactament $A\vdash_T B$ i $B\vdash_T A$.

Recíprocament, si es donen aquestes dues deduïbilitats, existeixen les dues fletxes. Com que $\cat{Prov}_T$ és prima, qualsevol endomorfisme $A\to A$ és necessàriament l'única fletxa $\id_A$, i igualment per a $B$. Per tant les dues composicions són les identitats i les fletxes són inverses. Així $A\cong B$.
\end{solucio}

\section{Classes especials de morfismes}

\begin{exercici}
Demostra que un morfisme $f$ és un isomorfisme si i només si és un monomorfisme escindit (té retracció) i un epimorfisme.
\end{exercici}

\begin{solucio}
Si $f$ és un isomorfisme, el seu invers és alhora retracció i secció, de manera que $f$ és mono escindit; i tot isomorfisme és epi. Recíprocament, suposem que $f\colon A\to B$ té una retracció $r$ (amb $r\comp f=\id_A$) i que és epi. Partint de $r\comp f=\id_A$ i component amb $f$ per l'esquerra,
\[
(f\comp r)\comp f = f\comp(r\comp f) = f\comp\id_A = f = \id_B\comp f.
\]
Com que $f$ és epi, es pot cancel·lar el factor $f$ de la dreta, d'on
\[
f\comp r=\id_B.
\]
Juntament amb $r\comp f=\id_A$, això mostra que $r$ és també invers per la dreta de $f$: $f$ és un isomorfisme amb invers $r=f^{-1}$.
\end{solucio}

\begin{exercici}
A $\cat{Set}$, considera $f\colon\{0,1\}\to\{\ast\}$, l'única aplicació, i $g\colon\{\ast\}\to\{0,1\}$ amb $g(\ast)=0$. Comprova quina de les dues composicions $f\comp g$ i $g\comp f$ és una identitat, digues quin dels dos morfismes és una secció i quin una retracció (i de qui), i determina quin dels dos és un monomorfisme i quin un epimorfisme.
\end{exercici}

\begin{solucio}
Calculem les dues composicions. Per una banda, $f\comp g\colon\{\ast\}\to\{\ast\}$ compleix $(f\comp g)(\ast)=f(0)=\ast$, de manera que $f\comp g=\id_{\{\ast\}}$. Per l'altra, $g\comp f\colon\{0,1\}\to\{0,1\}$ envia tant $0$ com $1$ a $0$, de manera que $g\comp f\neq\id_{\{0,1\}}$ (per exemple, $(g\comp f)(1)=0\neq 1$).

Així, l'única igualtat que val és $f\comp g=\id_{\{\ast\}}$. Llegida des de cada morfisme, aquesta equació diu dues coses alhora: $g$ és una \emph{secció} de $f$ (un invers per la dreta de $f$), i $f$ és una \emph{retracció} de $g$ (un invers per l'esquerra de $g$). Convé insistir que \enquote{secció} i \enquote{retracció} són papers relatius a l'altre morfisme, no propietats absolutes: aquí $g$ no és \enquote{una secció} en abstracte, sinó una secció \emph{de $f$}.

Pel que fa a mono i epi: com que $f$ té una secció, $f$ és un epimorfisme; com a aplicació, és exhaustiva, d'acord amb la caracterització dels epimorfismes de $\cat{Set}$. Com que $g$ té una retracció, $g$ és un monomorfisme; com a aplicació, és injectiva. En canvi, $f$ no és mono, perquè l'aplicació $f$ no és injectiva ($f(0)=f(1)$), i $g$ no és epi, perquè l'aplicació $g$ no és exhaustiva ($1\notin g(\{\ast\})$). Cap dels dos és un isomorfisme: ho serien només si $g\comp f$ fos també una identitat, cosa que hem vist que no passa.
\end{solucio}

\begin{exercici}\label{exr:pr-prima-mono-epi}
Sigui $\cat{C}$ una categoria prima, és a dir, una categoria en què entre dos objectes hi ha com a màxim un morfisme. Comprova que tot morfisme de $\cat{C}$ és alhora un monomorfisme i un epimorfisme. Explica per què aquest fet mostra que, en una categoria arbitrària, \enquote{mono} i \enquote{epi} no es poden interpretar com \enquote{injectiu} i \enquote{exhaustiu}.
\end{exercici}

\begin{solucio}
Sigui $f\colon A\to B$ un morfisme de $\cat{C}$. Per veure que és un \textbf{monomorfisme}, siguin $g,h\colon X\to A$ amb $f\comp g=f\comp h$:
\[
\begin{tikzpicture}[baseline=(m.center)]
  \matrix (m) [matrix of math nodes, column sep=3.4em] { X & A & B \\ };
  \draw[-{Stealth[scale=1.0]}] ([yshift=2.6pt]m-1-1.east) -- node[above] {$g$} ([yshift=2.6pt]m-1-2.west);
  \draw[-{Stealth[scale=1.0]}] ([yshift=-2.6pt]m-1-1.east) -- node[below] {$h$} ([yshift=-2.6pt]m-1-2.west);
  \draw[-{Stealth[scale=1.1]}] (m-1-2) -- node[above] {$f$} (m-1-3);
\end{tikzpicture}
\qquad
f\comp g=f\comp h\ \Longrightarrow\ g=h.
\]
Els morfismes $g$ i $h$ tenen el mateix domini $X$ i el mateix codomini $A$: tots dos pertanyen a $\Hom(X,A)$. Com que la categoria és prima, aquest homconjunt té com a màxim un element, i per tant $g=h$. Observem que ni tan sols hem fet servir la hipòtesi $f\comp g=f\comp h$: la cancel·lació es compleix perquè no hi ha res a cancel·lar.

Per veure que és un \textbf{epimorfisme}, siguin $g,h\colon B\to Y$ amb $g\comp f=h\comp f$:
\[
\begin{tikzpicture}[baseline=(m.center)]
  \matrix (m) [matrix of math nodes, column sep=3.4em] { A & B & Y \\ };
  \draw[-{Stealth[scale=1.1]}] (m-1-1) -- node[above] {$f$} (m-1-2);
  \draw[-{Stealth[scale=1.0]}] ([yshift=2.6pt]m-1-2.east) -- node[above] {$g$} ([yshift=2.6pt]m-1-3.west);
  \draw[-{Stealth[scale=1.0]}] ([yshift=-2.6pt]m-1-2.east) -- node[below] {$h$} ([yshift=-2.6pt]m-1-3.west);
\end{tikzpicture}
\qquad
g\comp f=h\comp f\ \Longrightarrow\ g=h.
\]
Ara $g,h\in\Hom(B,Y)$, i el mateix argument dona $g=h$. També ho podem obtenir per dualitat: la categoria oposada $\cat{C}\op$ també és prima, perquè $\Hom_{\cat{C}\op}(B,A)$ està en bijecció amb $\Hom_{\cat{C}}(A,B)$; per tant $f\op$ és mono a $\cat{C}\op$, és a dir, $f$ és epi a $\cat{C}$ (teorema~\ref{teo:dualitat}).

\emph{Per què això impedeix llegir mono com a injectiu i epi com a exhaustiu.} Hi ha dues raons, una més radical que l'altra.

Primer, en una categoria prima els morfismes no són, en general, aplicacions. En un preordre, la fletxa $x\to y$ només registra que $x\le y$; no envia cap element enlloc, i preguntar si és \enquote{injectiva} no té sentit. Mono i epi, en canvi, sempre en tenen, perquè es defineixen només amb la composició.

Segon, fins i tot si volguéssim mantenir l'analogia, entraria en contradicció. A $\cat{Set}$, una aplicació injectiva i exhaustiva és bijectiva, és a dir, un isomorfisme. Si \enquote{mono $=$ injectiu} i \enquote{epi $=$ exhaustiu} valguessin en general, tot morfisme d'una categoria prima seria un isomorfisme. Però a $\cat{2}=(0\to 1)$ la fletxa $u\colon 0\to 1$ és mono i epi pel que acabem de demostrar, i no té invers, perquè no hi ha cap fletxa $1\to 0$. En un preordre qualsevol passa el mateix: la fletxa $x\to y$ és sempre mono i epi, i només és un isomorfisme quan també $y\le x$.

La conclusió és que \enquote{mono $=$ injectiu} i \enquote{epi $=$ exhaustiu} són teoremes sobre $\cat{Set}$, no definicions. Les definicions són les de cancel·lació, i en una categoria prima es compleixen de la manera més trivial possible.
\end{solucio}

\begin{exercici}\label{pr-mono-no-escindit}
Compara les aplicacions injectives $i\colon\varnothing\to\{\ast\}$ i $j\colon\{0\}\hookrightarrow\{0,1\}$ de $\cat{Set}$. Totes dues són monomorfismes? Quina d'elles és un monomorfisme escindit?
\end{exercici}

\begin{solucio}
Totes dues són injectives i, per tant, monomorfismes de $\cat{Set}$.

La inclusió $j$ té una retracció: l'aplicació $r\colon\{0,1\}\to\{0\}$ constant compleix $r\comp j=\id_{\{0\}}$. És, doncs, un monomorfisme escindit.

L'aplicació $i$, en canvi, no en té cap: una retracció seria una aplicació $\{\ast\}\to\varnothing$, i no n'hi ha cap, perquè $\ast$ no té on anar. Per tant, $i$ és un monomorfisme que no és escindit. L'exemple mostra que l'afirmació \enquote{tota aplicació injectiva té una inversa per l'esquerra} només és certa per a dominis no buits; i aquest contraexemple no depèn de l'axioma de l'elecció, a diferència del cas dual de les seccions.
\end{solucio}

\section{Construccions sobre categories}

\begin{exercici}
Descriu explícitament la categoria producte $\cat{2}\times\cat{3}$, on $\cat{2}$ és la categoria associada a l'ordre $0<1$ i $\cat{3}$ la categoria associada a $0<1<2$. Quants objectes i quants morfismes té? Com es componen?
\end{exercici}

\begin{solucio}
Els objectes són les parelles $(i,j)$ amb $i\in\{0,1\}$ i $j\in\{0,1,2\}$. N'hi ha sis i es poden disposar en una graella:
\[
\begin{tikzpicture}[baseline=(m.center)]
  \matrix (m) [matrix of math nodes, row sep=2.4em, column sep=2.8em] {
    (0,0) & (0,1) & (0,2) \\
    (1,0) & (1,1) & (1,2) \\
  };
  \draw[-{Stealth[scale=1.05]}] (m-1-1) -- (m-1-2);
  \draw[-{Stealth[scale=1.05]}] (m-1-2) -- (m-1-3);
  \draw[-{Stealth[scale=1.05]}] (m-2-1) -- (m-2-2);
  \draw[-{Stealth[scale=1.05]}] (m-2-2) -- (m-2-3);
  \draw[-{Stealth[scale=1.05]}] (m-1-1) -- (m-2-1);
  \draw[-{Stealth[scale=1.05]}] (m-1-2) -- (m-2-2);
  \draw[-{Stealth[scale=1.05]}] (m-1-3) -- (m-2-3);
\end{tikzpicture}
\]
Només hi hem dibuixat les fletxes bàsiques; les composicions produeixen també fletxes més llargues i diagonals.

Com que $\cat{2}$ i $\cat{3}$ són primes, entre $(i,j)$ i $(i',j')$ hi ha un únic morfisme exactament quan
\[
i\le i'
\qquad\text{i}\qquad
j\le j'.
\]
$\cat{2}$ té tres morfismes i $\cat{3}$ en té sis. Com que un morfisme del producte és una parella, hi ha
\[
3\cdot6=18
\]
morfismes. La composició es fa component a component. Tots els quadrats de la graella commuten perquè el producte és també una categoria prima.
\end{solucio}

\begin{exercici}
Descriu explícitament la categoria oposada $\cat{P}\op$ quan $\cat{P}$ és un preordre $(P,\leq)$ vist com a categoria. Quina relació d'ordre li correspon?
\end{exercici}

\begin{solucio}
A $\cat{P}$ hi ha un morfisme $x\to y$ exactament quan $x\leq y$. A $\cat{P}\op$ les fletxes s'inverteixen:
\[
x\longrightarrow y\text{ a }\cat{P}\op
\quad\Longleftrightarrow\quad
y\longrightarrow x\text{ a }\cat{P}
\quad\Longleftrightarrow\quad
y\leq x.
\]
Per tant $\cat{P}\op$ és el preordre invers $(P,\geq)$. En aquest exemple, invertir les fletxes és literalment invertir la desigualtat.
\end{solucio}

\begin{exercici}
Sigui $\cat{2}=(0\to 1)$ la categoria amb dos objectes i una fletxa no trivial. Prenent com a ambient la mateixa $\cat{2}$, descriu explícitament la categoria sobre l'objecte $1$, $\cat{2}/1$, i la categoria sota l'objecte $1$, $1/\cat{2}$, comptant-ne objectes i morfismes. Comprova després que $(\cat{2}/1)\op$ \emph{no} és isomorfa a $1/\cat{2}$, i identifica amb quina categoria sí que ho és.
\end{exercici}

\begin{solucio}
Escrivim la fletxa no trivial de $\cat{2}$ com $u\colon 0\to 1$.

\emph{La categoria $\cat{2}/1$} (sobre $1$). Els seus objectes són els morfismes de $\cat{2}$ amb codomini $1$: la fletxa $u\colon 0\to 1$ i la identitat $\id_1\colon 1\to 1$. Són, doncs, \emph{dos} objectes. Un morfisme de $u$ a $\id_1$ és una fletxa $h$ de $\cat{2}$ amb $\id_1\comp h=u$, és a dir $h=u$; i les identitats de cada objecte completen els morfismes. La categoria $\cat{2}/1$ té la forma d'una fletxa entre dos objectes: és isomorfa a $\cat{2}$.

\emph{La categoria $1/\cat{2}$} (sota $1$). Els seus objectes són els morfismes de $\cat{2}$ amb domini $1$: només $\id_1$, perquè de $1$ no en surt cap altra fletxa. És, doncs, \emph{un} sol objecte, amb només la seva identitat: la categoria $\cat{1}$.

\emph{El contrast.} $\cat{2}/1$ té dos objectes i $1/\cat{2}$ en té un. Per tant $(\cat{2}/1)\op$ té dos objectes i \emph{no} pot ser isomorfa a $1/\cat{2}$. El que sí que val és la fórmula general de dualitat
\[
(\cat{2}/1)\op\;\cong\;1/(\cat{2}\op).
\]
En efecte, $\cat{2}\op$ és la fletxa invertida $1\to 0$; sota l'objecte $1$ a $\cat{2}\op$ hi ha els morfismes amb domini $1$, que són $\id_1$ i la fletxa $1\to 0$: \emph{dos} objectes, com $(\cat{2}/1)\op$. La moral: en dualitzar la construcció \enquote{sobre} en \enquote{sota} cal dualitzar també la categoria ambient, no deixar-la intacta.
\end{solucio}

\section{Principi de dualitat}

\begin{exercici}\label{exr:pr-dual-mono-epi}
Demostra que la composició de dos monomorfismes és un monomorfisme. A continuació, sense repetir el càlcul, dedueix per dualitat que la composició de dos epimorfismes és un epimorfisme. Indica exactament què s'inverteix en passar a $\cat{C}\op$ i què es manté.
\end{exercici}

\begin{solucio}
\emph{Primera part: el cas dels monomorfismes.} Siguin $f\colon A\to B$ i $g\colon B\to C$ monomorfismes, i siguin $u,v\colon X\to A$ amb $(g\comp f)\comp u=(g\comp f)\comp v$:
\[
\begin{tikzpicture}[baseline=(m.center)]
  \matrix (m) [matrix of math nodes, column sep=3em] { X & A & B & C \\ };
  \draw[-{Stealth[scale=1.0]}] ([yshift=2.6pt]m-1-1.east) -- node[above] {$u$} ([yshift=2.6pt]m-1-2.west);
  \draw[-{Stealth[scale=1.0]}] ([yshift=-2.6pt]m-1-1.east) -- node[below] {$v$} ([yshift=-2.6pt]m-1-2.west);
  \draw[-{Stealth[scale=1.1]}] (m-1-2) -- node[above] {$f$} (m-1-3);
  \draw[-{Stealth[scale=1.1]}] (m-1-3) -- node[above] {$g$} (m-1-4);
\end{tikzpicture}
\qquad
(g\comp f)\comp u=(g\comp f)\comp v\ \Longrightarrow\ u=v.
\]
Per associativitat, la hipòtesi és $g\comp(f\comp u)=g\comp(f\comp v)$. Com que $g$ és mono, en cancel·lem el factor de l'esquerra i obtenim $f\comp u=f\comp v$. Com que $f$ és mono, tornem a cancel·lar i obtenim $u=v$. Per tant $g\comp f$ és mono. Observem l'ordre: primer es cancel·la $g$, el factor \emph{exterior}, i després $f$.

\emph{Segona part: dualització.} Procedim en tres passos.

\textbf{1. L'enunciat és categòric.} \enquote{Si $f\colon A\to B$ i $g\colon B\to C$ són monomorfismes, aleshores $g\comp f$ és un monomorfisme} només parla d'objectes, morfismes, composició i igualtats entre composicions (la definició de mono és d'aquesta mena). Per tant té dual, i com que l'hem demostrat en una categoria qualsevol, el principi de dualitat (teorema~\ref{teo:dualitat}) garanteix que el dual també és vàlid en tota categoria.

\textbf{2. Què s'inverteix.} Recorrem l'enunciat element per element:
\begin{center}
\begin{tabular}{@{}ll@{}}
\toprule
A $\cat{C}$ & A l'enunciat dual \\
\midrule
$f\colon A\to B$, \quad $g\colon B\to C$ & $f\colon B\to A$, \quad $g\colon C\to B$ \\
la composta $g\comp f\colon A\to C$ & la composta $f\comp g\colon C\to A$ \\
fletxes de prova $u,v\colon X\to A$ & fletxes de prova $u,v\colon A\to X$ \\
mono: $m\comp u=m\comp v\Rightarrow u=v$ & epi: $u\comp m=v\comp m\Rightarrow u=v$ \\
\bottomrule
\end{tabular}
\end{center}
Es mantenen els quantificadors (\enquote{per a tot $u,v$}), la implicació, la igualtat i l'associativitat, que és autodual. L'enunciat dual diu, doncs: si $f\colon B\to A$ i $g\colon C\to B$ són epimorfismes, aleshores $f\comp g\colon C\to A$ és un epimorfisme. Rebatejant les fletxes perquè la cadena es llegeixi d'esquerra a dreta, $p\colon A\to B$ i $q\colon B\to C$, això és: \emph{si $p$ i $q$ són epimorfismes, $q\comp p$ és un epimorfisme.}

\textbf{3. Per què és vàlid.} És l'argument del principi, fet explícit. Siguin $p\colon A\to B$ i $q\colon B\to C$ epimorfismes de $\cat{C}$. A $\cat{C}\op$, les fletxes $p\op\colon B\to A$ i $q\op\colon C\to B$ són monomorfismes, perquè ser epi a $\cat{C}$ és ser mono a $\cat{C}\op$. Per la primera part, aplicada a $\cat{C}\op$, la composta $p\op\comp q\op\colon C\to A$ és mono a $\cat{C}\op$. Però $p\op\comp q\op=(q\comp p)\op$, de manera que $(q\comp p)\op$ és mono a $\cat{C}\op$, és a dir, $q\comp p$ és epi a $\cat{C}$.

\emph{Comprovació: la prova dual.} No l'hem necessitat, però val la pena escriure-la per veure que és la primera llegida a l'inrevés. Siguin $u,v\colon C\to X$ amb $u\comp(q\comp p)=v\comp(q\comp p)$:
\[
\begin{tikzpicture}[baseline=(m.center)]
  \matrix (m) [matrix of math nodes, column sep=3em] { A & B & C & X \\ };
  \draw[-{Stealth[scale=1.1]}] (m-1-1) -- node[above] {$p$} (m-1-2);
  \draw[-{Stealth[scale=1.1]}] (m-1-2) -- node[above] {$q$} (m-1-3);
  \draw[-{Stealth[scale=1.0]}] ([yshift=2.6pt]m-1-3.east) -- node[above] {$u$} ([yshift=2.6pt]m-1-4.west);
  \draw[-{Stealth[scale=1.0]}] ([yshift=-2.6pt]m-1-3.east) -- node[below] {$v$} ([yshift=-2.6pt]m-1-4.west);
\end{tikzpicture}
\qquad
u\comp(q\comp p)=v\comp(q\comp p)\ \Longrightarrow\ u=v.
\]
Per associativitat, $(u\comp q)\comp p=(v\comp q)\comp p$; com que $p$ és epi, $u\comp q=v\comp q$; com que $q$ és epi, $u=v$. Les mateixes passes, amb cada composició en l'ordre invers: ara es cancel·la primer $p$, el factor de la \emph{dreta}, que és la fletxa per on comença el camí.

\emph{Error típic.} En dualitzar és fàcil invertir les fletxes però oblidar d'invertir l'ordre de la composició, i escriure \enquote{$p\comp q$ és epi}. Si $p\colon A\to B$ i $q\colon B\to C$, aquesta composta ni tan sols està definida, tret que $C=A$. Comprovar els dominis i els codominis de cada composta és la millor manera de detectar l'error.
\end{solucio}

\section{Grafs i categories lliures}

\begin{exercici}
Considera el graf amb un sol vèrtex i dues arestes que són bucles, $a$ i $b$. Descriu la categoria lliure que genera. Quants morfismes té?
\end{exercici}

\begin{solucio}
La categoria lliure té un sol objecte. Els morfismes són tots els camins finits:
\[
\varepsilon,
\quad a,
\quad b,
\quad aa,
\quad ab,
\quad ba,
\quad bb,
\quad aaa,
\quad\ldots
\]
on $\varepsilon$ és el camí buit i actua com a identitat. La composició és la concatenació de paraules en ordre de recorregut: $b\comp a=ab$. Com que cap relació no identifica camins diferents, hi ha una infinitat numerable de morfismes. Llevat de l'isomorfisme canònic que inverteix l'ordre de les lletres (exemple~\ref{ex:monoide-lliure-cat}), és el monoide lliure sobre l'alfabet $\{a,b\}$.
\end{solucio}

\section{Qüestions de mida}

\begin{exercici}
En una categoria petita, l'oposada $\cat{C}\op$ també és petita. Comprova-ho, i comprova a més que $\cat{C}$ i $\cat{C}\op$ tenen exactament el mateix nombre de morfismes.
\end{exercici}

\begin{solucio}
Els objectes de $\cat{C}\op$ són els mateixos que els de $\cat{C}$. Cada morfisme $f\colon A\to B$ de $\cat{C}$ correspon al mateix morfisme vist com $f\colon B\to A$ a $\cat{C}\op$. Això dona una bijecció entre les col·leccions de morfismes de totes dues categories.

Per tant, si els objectes i els morfismes de $\cat{C}$ formen conjunts, també ho fan els de $\cat{C}\op$. En particular, $\cat{C}$ és petita si i només si ho és $\cat{C}\op$.
\end{solucio}

\begin{exercici}\label{exr:pr-set-localment-petita}
Comprova que $\cat{Set}$ és localment petita: per a conjunts $A$ i $B$, mostra que la col·lecció de les aplicacions $A\to B$ és un conjunt, identificant cada aplicació amb el seu graf, que és un subconjunt d'$A\times B$.
\end{exercici}

\begin{solucio}
Una aplicació $f\colon A\to B$ queda determinada pel seu \textbf{graf}
\[
\Gamma_f=\{(a,f(a))\mid a\in A\}\subseteq A\times B.
\]
Un subconjunt $\Gamma\subseteq A\times B$ és el graf d'alguna aplicació $A\to B$ exactament quan és \emph{funcional}: per a cada $a\in A$ hi ha un únic $b\in B$ amb $(a,b)\in\Gamma$. Així, $f\mapsto\Gamma_f$ és una bijecció entre les aplicacions $A\to B$ i els subconjunts funcionals d'$A\times B$; la inversa envia $\Gamma$ a l'aplicació que assigna a cada $a$ l'únic $b$ amb $(a,b)\in\Gamma$.

Vegem ara que aquests subconjunts funcionals formen un conjunt. Com que $A$ i $B$ són conjunts, el producte $A\times B$ és un conjunt, i també ho és el seu conjunt de parts $\mathcal P(A\times B)$. La condició \enquote{ser funcional} és una propietat ben definida dels elements de $\mathcal P(A\times B)$, de manera que, per l'axioma de separació,
\[
\{\Gamma\in\mathcal P(A\times B)\mid \forall a\in A\ \exists!\,b\in B\ (a,b)\in\Gamma\}
\]
és un conjunt. Per la bijecció anterior, $\Hom_{\cat{Set}}(A,B)$ és un conjunt.

Una precisió, la mateixa de l'observació~\ref{obs:extrems-implicits} de la Teoria per a $\cat{Rel}$: el graf no determina el codomini, perquè $\Gamma_f\subseteq A\times B$ també és un subconjunt d'$A\times B'$ per a tot $B'\supseteq B$. Per això, en rigor, un morfisme de $\cat{Set}$ és la terna $(A,B,\Gamma_f)$. Això no canvia la conclusió: les ternes $(A,B,\Gamma)$ amb $\Gamma$ funcional formen el conjunt $\{A\}\times\{B\}\times\{\Gamma\in\mathcal P(A\times B)\mid \Gamma\text{ funcional}\}$, que està en bijecció amb l'anterior.

\emph{Abast.} El mateix argument explica per què $\Hom(A,B)$ és un conjunt en els exemples habituals. A $\cat{Rel}$, directament, $\Hom_{\cat{Rel}}(X,Y)$ és, llevat dels extrems, $\mathcal P(X\times Y)$. A les categories de conjunts estructurats ($\cat{Grp}$, $\cat{Top}$, $\cat{Vect}_{\mathbb K}$, $\cat{Mon}$, $\cat{Pos}$), els morfismes són aplicacions entre els conjunts subjacents que compleixen una condició addicional; formen, doncs, un subconjunt de l'homconjunt de $\cat{Set}$ corresponent, de nou per separació. En canvi, la col·lecció de \emph{tots} els objectes de $\cat{Set}$ no és un conjunt (paradoxa de Russell): $\cat{Set}$ és localment petita sense ser petita.
\end{solucio}

\chapter{Functors}\label{cap:practica-functors}

En aquest capítol resolem amb detall una selecció d'exercicis sobre \emph{functors}. Com al capítol anterior, l'objectiu no és només arribar al resultat, sinó mostrar com es comprova que una assignació és un functor, com es reconeixen els functors coneguts sota formes noves i com es distingeixen les propietats que un functor preserva de les que no. Convé intentar cada exercici abans de llegir-ne la solució.

\section{Definició de functor}

\begin{exercici}\label{prac:functor-quadrat}
Sigui $F\colon\cat{C}\to\cat{D}$ un functor i suposem que a $\cat{C}$ tenim un quadrat commutatiu, és a dir, quatre morfismes $f\colon A\to B$, $u\colon A\to C$, $v\colon B\to D$ i $g\colon C\to D$ amb
\[
v\comp f=g\comp u.
\]
Mostra que la seva imatge per $F$ torna a ser un quadrat commutatiu:
\[
F(v)\comp F(f)=F(g)\comp F(u).
\]
\end{exercici}

\begin{solucio}
El quadrat de partida i la seva imatge són:
\[
\begin{tikzpicture}[baseline=(m.center)]
  \matrix (m) [matrix of math nodes, row sep=2.6em, column sep=3em] {
    A & B \\
    C & D \\
  };
  \draw[-{Stealth[scale=1.1]}] (m-1-1) -- node[above] {$f$} (m-1-2);
  \draw[-{Stealth[scale=1.1]}] (m-1-1) -- node[left] {$u$} (m-2-1);
  \draw[-{Stealth[scale=1.1]}] (m-1-2) -- node[right] {$v$} (m-2-2);
  \draw[-{Stealth[scale=1.1]}] (m-2-1) -- node[below] {$g$} (m-2-2);
\end{tikzpicture}
\qquad\stackrel{F}{\longmapsto}\qquad
\begin{tikzpicture}[baseline=(m.center)]
  \matrix (m) [matrix of math nodes, row sep=2.6em, column sep=3.4em] {
    F(A) & F(B) \\
    F(C) & F(D) \\
  };
  \draw[-{Stealth[scale=1.1]}] (m-1-1) -- node[above] {$F(f)$} (m-1-2);
  \draw[-{Stealth[scale=1.1]}] (m-1-1) -- node[left] {$F(u)$} (m-2-1);
  \draw[-{Stealth[scale=1.1]}] (m-1-2) -- node[right] {$F(v)$} (m-2-2);
  \draw[-{Stealth[scale=1.1]}] (m-2-1) -- node[below] {$F(g)$} (m-2-2);
\end{tikzpicture}
\]
Partim de la hipòtesi $v\comp f=g\comp u$. Apliquem $F$ a aquesta igualtat; com que $F$ és una aplicació ben definida sobre els morfismes, morfismes iguals tenen imatges iguals:
\[
F(v\comp f)=F(g\comp u).
\]
Ara fem servir que $F$ preserva la composició a cada costat:
\[
F(v)\comp F(f)=F(v\comp f)=F(g\comp u)=F(g)\comp F(u).
\]
Per tant el quadrat imatge commuta. Aquest és el contingut de la idea que \emph{els functors envien diagrames commutatius a diagrames commutatius}: n'hi ha prou que ho comprovem sobre les composicions, perquè un diagrama commuta exactament quan certes composicions coincideixen.
\end{solucio}

\begin{exercici}\label{pr-parts-cov}
Comprova que l'assignació que envia cada conjunt $X$ al conjunt de parts $\mathcal P(X)$ i cada aplicació $f\colon X\to Y$ a la \textbf{imatge directa}
\[
\mathcal P_!(f)\colon\mathcal P(X)\to\mathcal P(Y),
\qquad
A\mapsto f(A)=\{f(a)\mid a\in A\},
\]
és un functor covariant
\[
\mathcal P_!\colon\cat{Set}\to\cat{Set},
\]
el \textbf{functor de parts per imatge directa}. (Reservem $\mathcal P$, sense subíndex, per al functor de parts contravariant de la Teoria, que actua per antiimatge.)
\end{exercici}

\begin{solucio}
Hem de comprovar la preservació d'identitats i de composició.

Per a la \textbf{identitat}, si $f=\id_X$ aleshores per a tot $A\subseteq X$,
\[
\mathcal P_!(\id_X)(A)=\{\id_X(a)\mid a\in A\}=A,
\]
de manera que $\mathcal P_!(\id_X)=\id_{\mathcal P(X)}$.

Per a la \textbf{composició}, siguin $f\colon X\to Y$ i $g\colon Y\to Z$. Per a tot $A\subseteq X$,
\[
\begin{aligned}
\mathcal P_!(g\comp f)(A)&=(g\comp f)(A)=g\big(f(A)\big)\\
&=\mathcal P_!(g)\big(\mathcal P_!(f)(A)\big)=\big(\mathcal P_!(g)\comp\mathcal P_!(f)\big)(A).
\end{aligned}
\]
La igualtat central, $(g\comp f)(A)=g(f(A))$, és la propietat coneguda de la imatge directa. Per tant $\mathcal P_!$ és un functor covariant. Té els mateixos objectes que $\mathcal P$, però actua sobre els morfismes en sentit contrari i per una operació diferent: són dos functors distints.
\end{solucio}

\begin{exercici}\label{pr-hom-cov}
Sigui $\cat{C}$ una categoria localment petita i $A$ un objecte fix. Comprova que l'assignació que envia cada objecte $X$ al conjunt $\Hom(A,X)$ i cada morfisme $f\colon X\to Y$ a l'aplicació
\[
f_*\colon\Hom(A,X)\to\Hom(A,Y),
\qquad
\varphi\mapsto f\comp\varphi,
\]
és un functor covariant $\Hom(A,-)\colon\cat{C}\to\cat{Set}$ (suposant $\cat{C}$ localment petita).
\end{exercici}

\begin{solucio}
Primer cal veure que $f_*$ està ben definida: si $\varphi\colon A\to X$, aleshores $f\comp\varphi\colon A\to Y$, de manera que $f_*(\varphi)\in\Hom(A,Y)$.

Per a la \textbf{identitat}, prenent $f=\id_X$,
\[
(\id_X)_*(\varphi)=\id_X\comp\varphi=\varphi,
\]
de manera que $(\id_X)_*=\id_{\Hom(A,X)}$.

Per a la \textbf{composició}, siguin $f\colon X\to Y$ i $g\colon Y\to Z$. Per a tot $\varphi\in\Hom(A,X)$, usant l'associativitat de $\cat{C}$,
\[
(g\comp f)_*(\varphi)=(g\comp f)\comp\varphi=g\comp(f\comp\varphi)=g_*\big(f_*(\varphi)\big).
\]
Per tant $(g\comp f)_*=g_*\comp f_*$ i $\Hom(A,-)$ és un functor covariant.
\end{solucio}

\section{Exemples de functors}

\begin{exercici}
Siguin $T$ i $T'$ dos sistemes deductius sobre el mateix conjunt de fórmules, amb $T'$ una \textbf{extensió} de $T$ ---és a dir, tot el que és deduïble a $T$ també ho és a $T'$: si $A\vdash_T B$, aleshores $A\vdash_{T'} B$. Comprova que la identitat sobre les fórmules indueix un functor
\[
J\colon\cat{Prov}_T\to\cat{Prov}_{T'},
\]
i explica en quins termes es podria formular que $J$ sigui ple.
\end{exercici}

\begin{solucio}
Recordem que $\cat{Prov}_T$ té les fórmules per objectes i una única fletxa $A\to B$ exactament quan $A\vdash_T B$; és una categoria prima. Definim $J$ com la identitat sobre els objectes, $J(A)=A$.

Sobre els morfismes, una fletxa $A\to B$ de $\cat{Prov}_T$ existeix exactament quan $A\vdash_T B$; la fletxa registra, doncs, el fet de la deduïbilitat, no una prova concreta. Com que $T'$ és una extensió de $T$, tenim $A\vdash_{T'} B$, de manera que hi ha una fletxa $J(A)\to J(B)$ a $\cat{Prov}_{T'}$, i és única perquè $\cat{Prov}_{T'}$ és prima. Això defineix $J$ sobre morfismes.

Les condicions de functor se satisfan automàticament perquè totes dues categories són primes: entre dos objectes hi ha com a màxim una fletxa, així que la preservació d'identitats i de composicions no imposa res més enllà de l'existència de les fletxes imatge, que ja hem comprovat.

Quant a \textbf{ser ple}: $J$ seria ple si, per a cada parella de fórmules $A,B$, tota fletxa $J(A)\to J(B)$ a $\cat{Prov}_{T'}$ provingués d'una fletxa $A\to B$ a $\cat{Prov}_T$. És a dir, si
\[
A\vdash_{T'} B\quad\Longrightarrow\quad A\vdash_T B.
\]
Dit d'una altra manera, $J$ és ple exactament quan l'extensió $T'$ \emph{no afegeix cap conseqüència nova} entre les fórmules considerades. Quan $T'$ demostra estrictament més relacions que $T$, el functor $J$ és fidel ---de fet, totes dues categories són primes, així que $J$ és sempre fidel--- però no ple.
\end{solucio}

\begin{exercici}
Comprova amb detall que un functor entre dos preordres $(P,\leq)$ i $(Q,\leq)$, vistos com a categories, és exactament una aplicació monòtona.
\end{exercici}

\begin{solucio}
Un functor $F\colon P\to Q$ dona, en primer lloc, una aplicació $x\mapsto F(x)$ dels objectes, és a dir, dels elements de $P$ als de $Q$.

La condició sobre morfismes diu que si hi ha una fletxa $x\to y$ a $P$, aleshores hi ha una fletxa $F(x)\to F(y)$ a $Q$. Com que en un preordre hi ha una fletxa $x\to y$ exactament quan $x\leq y$, això equival a
\[
x\leq y
\quad\Longrightarrow\quad
F(x)\leq F(y),
\]
que és precisament la definició d'aplicació monòtona.

Recíprocament, tota aplicació monòtona defineix un functor: envia l'única fletxa $x\to y$ (quan $x\leq y$) a l'única fletxa $F(x)\to F(y)$. Les condicions de preservació d'identitats i composició es compleixen automàticament, perquè entre dos objectes hi ha com a màxim una fletxa i totes les igualtats entre morfismes són trivials. Així, els functors $P\to Q$ i les aplicacions monòtones $P\to Q$ són el mateix.
\end{solucio}

\begin{exercici}
Sigui $M$ un monoide vist com a categoria $\cat{B}M$ amb un sol objecte $\ast$. Comprova que un functor $F\colon\cat{B}M\to\cat{Set}$ és el mateix que una acció de $M$ sobre un conjunt, és a dir, un conjunt $S$ juntament amb una operació
\[
M\times S\to S,
\qquad
(m,s)\mapsto m\cdot s,
\]
tal que $1_M\cdot s=s$ i $(mn)\cdot s=m\cdot(n\cdot s)$.
\end{exercici}
\begin{solucio}
Un functor $F\colon\cat{B}M\to\cat{Set}$ tria un conjunt $S=F(\ast)$ i envia cada element $m\in M$ ---cada morfisme $\ast\to\ast$--- a una aplicació $F(m)\colon S\to S$. Definim l'acció per
\[
m\cdot s=F(m)(s).
\]

La preservació de la identitat diu $F(1_M)=\id_S$, és a dir, $1_M\cdot s=s$ per a tot $s$. La preservació de la composició diu $F(mn)=F(m)\comp F(n)$, és a dir,
\[
(mn)\cdot s=F(mn)(s)=F(m)\big(F(n)(s)\big)=m\cdot(n\cdot s).
\]
Aquestes són exactament les dues lleis d'una acció per l'esquerra. Recíprocament, tota acció defineix, per la mateixa fórmula, un functor. Per tant els functors $\cat{B}M\to\cat{Set}$ i les accions de $M$ coincideixen.
\end{solucio}

\begin{exercici}
Comprova que la construcció de la categoria lliure és functorial sobre les fletxes: donat un morfisme de grafs $\alpha\colon G\to H$, l'assignació que envia cada camí d'arestes de $G$ a la seqüència de les seves imatges és un functor $\Free(\alpha)\colon\Free(G)\to\Free(H)$.
\end{exercici}

\begin{solucio}
Un camí a $\Free(G)$ és una seqüència finita d'arestes encadenades $a_1a_2\cdots a_n$. Definim
\[
\Free(\alpha)(a_1a_2\cdots a_n)=\alpha_E(a_1)\,\alpha_E(a_2)\cdots\alpha_E(a_n),
\]
i sobre els objectes $\Free(\alpha)=\alpha_V$. Cal comprovar que aquesta seqüència és realment un camí a $H$: com que $\alpha$ preserva origen i destí, si $a_i$ acaba on comença $a_{i+1}$, aleshores $\alpha_E(a_i)$ acaba on comença $\alpha_E(a_{i+1})$, de manera que les imatges també s'encadenen.

La \textbf{identitat} de cada vèrtex és el camí buit, que s'envia al camí buit del vèrtex imatge: $\Free(\alpha)$ preserva identitats. La \textbf{composició} és la concatenació. Siguin $p=(a_1,\dots,a_r)$ i $q=(b_1,\dots,b_s)$ dos camins componibles, amb el final de $p$ igual a l'inici de $q$; la seva composició és el camí que recorre primer $p$ i després $q$,
\[
q\comp p=(a_1,\dots,a_r,b_1,\dots,b_s).
\]
Aplicar $\alpha_E$ terme a terme i després concatenar dona el mateix que concatenar i després aplicar $\alpha_E$ terme a terme, de manera que
\[
\Free(\alpha)(q\comp p)=\Free(\alpha)(q)\comp\Free(\alpha)(p).
\]
Per tant $\Free(\alpha)$ és un functor.
\end{solucio}

\section{Composició de functors i la categoria \texorpdfstring{$\cat{Cat}$}{Cat}}

\begin{exercici}
Comprova que la composició de dos functors $F\colon\cat{C}\to\cat{D}$ i $G\colon\cat{D}\to\cat{E}$, seguida d'un tercer $H\colon\cat{E}\to\cat{F}$, és associativa:
\[
H\comp(G\comp F)=(H\comp G)\comp F.
\]
\end{exercici}

\begin{solucio}
Dues igualtats de functors es comproven sobre objectes i sobre morfismes. Sobre un objecte $A$ de $\cat{C}$,
\[
\big(H\comp(G\comp F)\big)(A)=H\big(G(F(A))\big)=\big((H\comp G)\comp F\big)(A),
\]
i sobre un morfisme $f$,
\[
\big(H\comp(G\comp F)\big)(f)=H\big(G(F(f))\big)=\big((H\comp G)\comp F\big)(f).
\]
En tots dos casos la igualtat es redueix a l'associativitat de l'aplicació successiva, primer $F$, després $G$, després $H$. Per tant la composició de functors és associativa, i amb els functors identitat com a neutres, això confirma que $\cat{Cat}$ és una categoria.
\end{solucio}

\begin{exercici}
Fixat un objecte $A$ d'una categoria $\cat{C}$, descriu el functor oblit $U_A\colon\cat{C}/A\to\cat{C}$ de la categoria de $\cat{C}$ sobre $A$ cap a $\cat{C}$, i comprova que és un functor.
\end{exercici}

\begin{solucio}
Recordem que un objecte de $\cat{C}/A$ és un morfisme $x\colon X\to A$ de $\cat{C}$, i un morfisme de $x\colon X\to A$ cap a $y\colon Y\to A$ és un morfisme $h\colon X\to Y$ de $\cat{C}$ tal que $y\comp h=x$ (el triangle commuta). Definim $U_A$ així: sobre objectes, envia $x\colon X\to A$ al seu domini $X$; sobre morfismes, envia el morfisme $h$ del triangle al mateix $h\colon X\to Y$ vist a $\cat{C}$, oblidant la condició $y\comp h=x$.

Comprovem que és un functor. \textbf{Identitats:} la identitat de l'objecte $x\colon X\to A$ a $\cat{C}/A$ és $\id_X$ (que compleix $x\comp\id_X=x$), i $U_A(\id_X)=\id_X=\id_{U_A(x)}$. \textbf{Composició:} si $h\colon x\to y$ i $k\colon y\to z$ són morfismes componibles a $\cat{C}/A$, la seva composició és $k\comp h$ com a morfisme de $\cat{C}$; per tant $U_A(k\comp h)=k\comp h=U_A(k)\comp U_A(h)$. Així $U_A$ preserva identitats i composició, i és un functor. Intuïtivament, $U_A$ \enquote{oblida} el morfisme cap a $A$ i reté només el domini i les fletxes entre dominis.
\end{solucio}

\section{Functors i isomorfismes}

\begin{exercici}
Dona un exemple de functor i d'un morfisme \emph{no} isomorf la imatge del qual \emph{sí} que és un isomorfisme. Això mostra que els functors preserven isomorfismes, però no els reflecteixen en general.
\end{exercici}

\begin{solucio}
Considerem el functor oblit $U\colon\cat{Pos}\to\cat{Set}$, que envia cada conjunt parcialment ordenat al seu conjunt subjacent i cada aplicació monòtona a ella mateixa. Siguin $P$ el conjunt $\{0,1\}$ amb l'ordre discret (cap parella d'elements diferents no és comparable) i $Q$ el mateix conjunt amb l'ordre $0<1$, i sigui
\[
f\colon P\to Q
\]
la identitat de conjunts. És monòtona, perquè a $P$ les úniques relacions són $0\le 0$ i $1\le 1$, que es conserven a $Q$. La inversa $Q\to P$, en canvi, no ho és: a $Q$ tenim $0\le 1$, però a $P$ no. Per tant $f$ \textbf{no és un isomorfisme} a $\cat{Pos}$, perquè no té cap invers dins de la categoria.

Ara bé, la seva imatge $U(f)$ és l'aplicació identitat del conjunt $\{0,1\}$, que \textbf{sí que és un isomorfisme} a $\cat{Set}$.

\emph{Un exemple anàleg, per a qui conegui la topologia.} El functor oblit $U\colon\cat{Top}\to\cat{Set}$ es comporta igual. Si $X$ és un conjunt amb dues topologies $\tau$ i $\tau'$, amb $\tau$ estrictament més fina que $\tau'$, la identitat $(X,\tau)\to(X,\tau')$ és contínua, però la seva inversa no ho és; no és, doncs, un isomorfisme a $\cat{Top}$, mentre que la seva imatge a $\cat{Set}$ és la identitat de $X$.

En tots dos casos, $U$ no reflecteix isomorfismes: la seva imatge pot ser invertible sense que l'original ho sigui. Els functors plens i fidels sí que reflecteixen isomorfismes, encara que aquesta condició no és necessària: hi ha functors que els reflecteixen sense ser plens. Per exemple, el functor oblit $\cat{Grp}\to\cat{Set}$ reflecteix isomorfismes ---un homomorfisme bijectiu de grups és un isomorfisme---, malgrat no ser ple.
\end{solucio}

\begin{exercici}\label{pr-2-set-mono-epi}
A la categoria $\cat{2}=(0\xrightarrow{u}1)$, la fletxa $u$ és mono i epi. Construeix un functor $\cat{2}\to\cat{Set}$ que no l'enviï a un monomorfisme, i un altre que no l'enviï a un epimorfisme.
\end{exercici}

\begin{solucio}
Un functor $\cat{2}\to\cat{Set}$ queda determinat per dos conjunts i una aplicació entre ells, imatges de $0$, $1$ i $u$: les úniques composicions de $\cat{2}$ involucren identitats, de manera que la functorialitat només demana enviar identitats a identitats.
\begin{itemize}
  \item $F(0)=\{0,1\}$, $F(1)=\{\ast\}$ i $F(u)$ l'única aplicació $\{0,1\}\to\{\ast\}$. No és injectiva, i per tant $F(u)$ no és mono.
  \item $G(0)=\{\ast\}$, $G(1)=\{0,1\}$ i $G(u)$ l'aplicació de valor $0$. No és exhaustiva, i per tant $G(u)$ no és epi.
\end{itemize}
Un functor, doncs, no preserva en general ni els monomorfismes ni els epimorfismes. A $\cat{2}$, que $u$ sigui mono i epi és conseqüència que la categoria és prima; a $\cat{Set}$ hi ha més fletxes amb què comparar, i aquesta trivialitat es perd.
\end{solucio}

\section{Functors covariants i contravariants}

\begin{exercici}\label{pr-parts-contra}
Comprova que el functor de parts contravariant $\mathcal P\colon\cat{Set}\op\to\cat{Set}$, que envia $f\colon X\to Y$ a l'antiimatge
\[
f^{-1}\colon\mathcal P(Y)\to\mathcal P(X),
\qquad
B\mapsto f^{-1}(B),
\]
és realment un prefeix sobre $\cat{Set}$ (definició~\ref{def:prefeix}).
\end{exercici}

\begin{solucio}
Cal comprovar la preservació d'identitats i la inversió de l'ordre en la composició.

Per a la \textbf{identitat}, $(\id_X)^{-1}(B)=B$ per a tot $B\subseteq X$, de manera que $(\id_X)^{-1}=\id_{\mathcal P(X)}$.

Per a la \textbf{composició}, siguin $f\colon X\to Y$ i $g\colon Y\to Z$. Per a tot $C\subseteq Z$,
\[
(g\comp f)^{-1}(C)=\{x\mid g(f(x))\in C\}=\{x\mid f(x)\in g^{-1}(C)\}=f^{-1}\big(g^{-1}(C)\big).
\]
Per tant
\[
(g\comp f)^{-1}=f^{-1}\comp g^{-1},
\]
amb l'ordre invertit. Aquesta és exactament la condició de prefeix, és a dir, de functor $\cat{Set}\op\to\cat{Set}$.
\end{solucio}

\begin{exercici}\label{pr-hom-contra}
Sigui $B$ un objecte fix d'una categoria localment petita $\cat{C}$. Comprova que $\Hom(-,B)$ és un prefeix sobre $\cat{C}$, és a dir, un functor $\Hom(-,B)\colon\cat{C}\op\to\cat{Set}$.
\end{exercici}

\begin{solucio}
A cada objecte $X$ li assignem el conjunt $\Hom(X,B)$. A cada morfisme $f\colon X\to Y$ li assignem l'aplicació de \textbf{precomposició}
\[
f^*\colon\Hom(Y,B)\to\Hom(X,B),
\qquad
\psi\mapsto\psi\comp f,
\]
que està ben definida perquè si $\psi\colon Y\to B$ aleshores $\psi\comp f\colon X\to B$. Observem que $f^*$ inverteix el sentit: parteix dels morfismes des de $Y$ i arriba als morfismes des de $X$.

Per a la \textbf{identitat}, $(\id_X)^*(\psi)=\psi\comp\id_X=\psi$. Per a la \textbf{composició}, siguin $f\colon X\to Y$ i $g\colon Y\to Z$; per a tot $\psi\in\Hom(Z,B)$,
\[
(g\comp f)^*(\psi)=\psi\comp(g\comp f)=(\psi\comp g)\comp f=f^*\big(g^*(\psi)\big).
\]
Per tant $(g\comp f)^*=f^*\comp g^*$, amb l'ordre invertit, i $\Hom(-,B)$ és contravariant.
\end{solucio}

\begin{exercici}\label{pr-bihom}
Sigui $\cat{C}$ localment petita. Comprova que el bifunctor
\[
\Hom(-,-)\colon\cat{C}\op\times\cat{C}\to\cat{Set},\qquad \Hom(f,g)(h)=g\comp h\comp f,
\]
respecta la composició simultània en les dues variables i les identitats. És a dir, comprova que és realment un functor, no només en cada variable per separat.
\end{exercici}

\begin{solucio}
Siguin $f\colon A'\to A$, $f'\colon A''\to A'$, $g\colon B\to B'$ i $g'\colon B'\to B''$ morfismes de $\cat{C}$. Recordem que la composició a $\cat{C}\op\times\cat{C}$ actua component a component, i que a la primera component (la de $\cat{C}\op$) l'ordre és l'invers del de $\cat{C}$. \textbf{Composició.} Per a tot $h\colon A\to B$,
\[
\begin{aligned}
\Hom(f',g')\big(\Hom(f,g)(h)\big)
&=g'\comp(g\comp h\comp f)\comp f'\\
&=(g'\comp g)\comp h\comp(f\comp f')\\
&=\Hom(f\comp f',\,g'\comp g)(h),
\end{aligned}
\]
on el pas central és només associativitat. El resultat $\Hom(f\comp f',g'\comp g)$ és precisament $\Hom$ aplicat a la composició de la parella $(f',g')$ amb $(f,g)$ a $\cat{C}\op\times\cat{C}$: a la segona component obtenim $g'\comp g$, i a la primera $f\comp f'$ (l'ordre invertit, coherent amb la contravariància). \textbf{Identitats.} Per a la parella identitat $(\id_A,\id_B)$,
\[
\Hom(\id_A,\id_B)(h)=\id_B\comp h\comp\id_A=h,
\]
és a dir $\Hom(\id_A,\id_B)=\id_{\Hom(A,B)}$. Per tant $\Hom(-,-)$ preserva composició i identitats de $\cat{C}\op\times\cat{C}$, i és un functor de ple dret.
\end{solucio}

\section{Functors fidels i plens}

\begin{exercici}
Comprova que el functor oblit $U\colon\cat{Pos}\to\cat{Set}$, que envia cada ordre parcial al seu conjunt subjacent i cada aplicació monòtona a ella mateixa, és fidel però no ple.
\end{exercici}

\begin{solucio}
\textbf{Fidel.} Siguin $f,g\colon P\to Q$ dues aplicacions monòtones amb $U(f)=U(g)$. Això vol dir que $f$ i $g$ coincideixen com a aplicacions entre els conjunts subjacents, és a dir, en tot element de $P$. Una aplicació monòtona no és res més que la seva acció sobre els elements (amb la condició afegida de preservar l'ordre), de manera que $f=g$. Així cada aplicació $U_{P,Q}\colon\Hom_{\cat{Pos}}(P,Q)\to\Hom_{\cat{Set}}(U(P),U(Q))$ és injectiva i, per tant, $U$ és fidel. (Sobre objectes, en canvi, $U$ no és injectiu: ordres diferents sobre un mateix conjunt tenen la mateixa imatge. Ser fidel no hi té res a veure.)

\textbf{No ple.} Prenem l'ordre parcial $P=\{0,1\}$ amb $0<1$ i el mateix $Q=P$. L'aplicació $\sigma\colon\{0,1\}\to\{0,1\}$ que intercanvia els dos elements ($\sigma(0)=1$, $\sigma(1)=0$) és una aplicació de conjunts perfectament vàlida, i pertany a $\Hom_{\cat{Set}}(U P,U Q)$. Però no és monòtona: de $0<1$ es passaria a $\sigma(0)=1$ i $\sigma(1)=0$, és a dir $1\le 0$, que és fals. Per tant $\sigma$ no és de la forma $U(f)$ per a cap aplicació monòtona $f$, l'aplicació $U_{P,Q}$ no és exhaustiva i $U$ no és ple.

\emph{Per a qui conegui els grups.} El functor oblit $U\colon\cat{Grp}\to\cat{Set}$ es comporta igual: és fidel perquè un homomorfisme queda determinat per la seva acció sobre els elements, i no és ple perquè hi ha aplicacions entre els conjunts subjacents que no són homomorfismes, com $x\mapsto x+1$ de $\mathbb Z$ en $\mathbb Z$, que no envia $0$ a $0$.
\end{solucio}

\begin{exercici}\label{pr-esquelet-finset}
Sigui $\cat{S}$ la subcategoria plena de $\cat{FinSet}$ que té per objectes els conjunts $\underline{0}=\varnothing$ i $\underline{n}=\{1,\dots,n\}$ per a $n\geq 1$. Comprova que la inclusió
\[
\iota\colon\cat{S}\hookrightarrow\cat{FinSet}
\]
és plena, fidel i essencialment exhaustiva, però no exhaustiva sobre objectes en sentit estricte.
\end{exercici}

\begin{solucio}
\textbf{Fidel.} La inclusió envia cada morfisme a ell mateix: per a objectes $\underline{m},\underline{n}$ de $\cat{S}$, l'aplicació $\iota_{\underline{m},\underline{n}}\colon\Hom_{\cat{S}}(\underline{m},\underline{n})\to\Hom_{\cat{FinSet}}(\underline{m},\underline{n})$ és $h\mapsto h$, que és injectiva.

\textbf{Plena.} Com que $\cat{S}$ és una subcategoria \emph{plena}, conté tots els morfismes de $\cat{FinSet}$ entre dos objectes seus: $\Hom_{\cat{S}}(\underline{m},\underline{n})=\Hom_{\cat{FinSet}}(\underline{m},\underline{n})$, el conjunt de totes les aplicacions $\underline{m}\to\underline{n}$. L'aplicació $\iota_{\underline{m},\underline{n}}$ és, doncs, la identitat d'aquest conjunt, i en particular exhaustiva.

\textbf{Essencialment exhaustiva.} Sigui $X$ un conjunt finit, i sigui $n$ el seu nombre d'elements. Per definició de cardinal finit, hi ha una bijecció $X\to\underline{n}$ (per a $n=0$, $X=\varnothing=\underline{0}$ i n'hi ha prou amb la identitat). Una bijecció és un isomorfisme de $\cat{FinSet}$, de manera que $X\cong\iota(\underline{n})$. Tot objecte de $\cat{FinSet}$ és, doncs, isomorf a una imatge de $\iota$.

\textbf{No exhaustiva sobre objectes.} El conjunt $\{2\}$ és un objecte de $\cat{FinSet}$, però no és cap dels $\underline{n}$: té un sol element, i l'únic $\underline{n}$ amb un element és $\underline{1}=\{1\}\neq\{2\}$. No és, per tant, igual a cap imatge de $\iota$, tot i ser isomorf a $\underline{1}$.

\emph{Per què és important.} Aquest functor reuneix les tres propietats ---ple, fidel i essencialment exhaustiu--- que al capítol~\ref{cap:transformacions} caracteritzaran les equivalències de categories, i alhora mostra que no cal arribar a tots els objectes en sentit estricte. Observem també que per a cada $X$ hem triat \emph{una} bijecció $X\to\underline{n}$; per construir un functor en sentit invers caldria triar-ne una per a cada conjunt finit alhora, i aquest és el punt on el capítol~\ref{cap:transformacions} recorrerà a l'axioma de l'elecció.
\end{solucio}


\chapter[Transformacions naturals]{\texorpdfstring{Transformacions naturals\\ i equivalències}{Transformacions naturals i equivalències}}

En aquest capítol resolem amb detall exercicis sobre \emph{transformacions naturals}, \emph{isomorfismes naturals} i \emph{equivalències de categories}. L'èmfasi és a comprovar la condició de naturalitat en casos concrets, a distingir un isomorfisme d'un isomorfisme natural, i a reconèixer una equivalència mitjançant les tres propietats que la caracteritzen. Convé intentar cada exercici abans de llegir-ne la solució.

\section{Definició de transformació natural}

\begin{exercici}
Siguin $F,G\colon\cat{C}\to\cat{D}$ dos functors i suposem que, per a cada objecte $A$, tenim un morfisme $\eta_A\colon F(A)\to G(A)$. Comprova que, per verificar que $\eta$ és una transformació natural, n'hi ha prou de comprovar la condició de naturalitat sobre un conjunt de morfismes que \emph{generi} $\cat{C}$ per composició, i que en el cas d'una categoria lliure $\Free(G_0)$ això vol dir comprovar-la només sobre les arestes del graf $G_0$.
\end{exercici}

\begin{solucio}
La condició de naturalitat per a un morfisme componible diu $G(f)\comp\eta_A=\eta_B\comp F(f)$. Suposem que val per a dos morfismes componibles $f\colon A\to B$ i $g\colon B\to C$. Aleshores, per a la composició $g\comp f$,
\[
G(g\comp f)\comp\eta_A=G(g)\comp G(f)\comp\eta_A=G(g)\comp\eta_B\comp F(f)=\eta_C\comp F(g)\comp F(f)=\eta_C\comp F(g\comp f),
\]
on hem fet servir la functorialitat de $G$ i $F$ i les condicions de naturalitat per a $f$ i per a $g$. Per tant, si la naturalitat val per a un conjunt de morfismes, val per a totes les seves composicions. També val trivialment per a les identitats, perquè $G(\id_A)\comp\eta_A=\eta_A=\eta_A\comp F(\id_A)$.

En una categoria lliure $\Free(G_0)$, tot morfisme és un camí, és a dir, una composició d'arestes del graf $G_0$. Per l'argument anterior, doncs, la naturalitat sobre tots els morfismes es dedueix de la naturalitat sobre les arestes. Això redueix molt la feina de comprovació.
\end{solucio}

\begin{exercici}
Considera la categoria $\cat{2}=(0\to 1)$ i dos functors $F,G\colon\cat{2}\to\cat{D}$. Descriu explícitament què és una transformació natural $\eta\colon F\Rightarrow G$.
\end{exercici}

\begin{solucio}
Un functor $F\colon\cat{2}\to\cat{D}$ és el mateix que un morfisme de $\cat{D}$: tria dos objectes $F(0),F(1)$ i un morfisme $F(u)\colon F(0)\to F(1)$, on $u\colon 0\to 1$ és l'única fletxa no trivial. Igualment $G$ dona un morfisme $G(u)\colon G(0)\to G(1)$.

Una transformació natural $\eta\colon F\Rightarrow G$ consisteix en dos components $\eta_0\colon F(0)\to G(0)$ i $\eta_1\colon F(1)\to G(1)$, subjectes a la condició de naturalitat per a $u$:
\[
G(u)\comp\eta_0=\eta_1\comp F(u).
\]
És a dir, una transformació natural entre dos morfismes de $\cat{D}$ és exactament un \textbf{quadrat commutatiu}
\[
\begin{tikzpicture}[baseline=(m.center)]
  \matrix (m) [matrix of math nodes, row sep=2.6em, column sep=3.2em] {
    F(0) & F(1) \\
    G(0) & G(1) \\
  };
  \draw[-{Stealth[scale=1.1]}] (m-1-1) -- node[above] {$F(u)$} (m-1-2);
  \draw[-{Stealth[scale=1.1]}] (m-1-1) -- node[left] {$\eta_0$} (m-2-1);
  \draw[-{Stealth[scale=1.1]}] (m-1-2) -- node[right] {$\eta_1$} (m-2-2);
  \draw[-{Stealth[scale=1.1]}] (m-2-1) -- node[below] {$G(u)$} (m-2-2);
\end{tikzpicture}
\]
Aquest exemple explica per què els quadrats commutatius apareixen tan sovint: són les transformacions naturals més simples possibles.
\end{solucio}

\section{Exemples de transformacions naturals}

\begin{exercici}\label{pr-sintaxi-semantica}
Considerem els functors sintàctic i semàntic $F,G\colon\cat{FinSet}\to\cat{Set}$ i la interpretació $\eta_X\colon F(X)\to G(X)$ de la secció~\ref{sec:cap-transformacio} de la Teoria. Comprova que $F$ i $G$ són functors i que $\eta$ és una transformació natural $F\Rightarrow G$.
\end{exercici}

\begin{solucio}
\textbf{$F$ és un functor.} Recordem que $F(f)(\varphi)=\varphi[f]$ es defineix per recursió: $F(f)(x)=f(x)$ sobre una variable, i $F(f)$ commuta amb les connectives, per exemple $F(f)(\varphi\land\psi)=F(f)(\varphi)\land F(f)(\psi)$. Comprovem les dues igualtats per recursió sobre $\varphi$. Per a la identitat, sobre una variable $F(\id_X)(x)=x$, i si $F(\id_X)$ deixa fixes $\varphi$ i $\psi$, també deixa fixa $\varphi\land\psi$ (i anàlogament per a les altres connectives); per tant $F(\id_X)=\id_{F(X)}$. Per a la composició, siguin $f\colon X\to Y$ i $g\colon Y\to Z$. Sobre una variable,
\[
F(g\comp f)(x)=g(f(x))=F(g)\bigl(F(f)(x)\bigr),
\]
i si la igualtat $F(g\comp f)=F(g)\comp F(f)$ val per a $\varphi$ i $\psi$, val per a $\varphi\land\psi$, perquè les tres aplicacions commuten amb $\land$. Així, substituir per $f$ i després per $g$ és substituir per $g\comp f$.

\textbf{$G$ és un functor.} Recordem que $G(f)(\theta)(w)=\theta(w\comp f)$ per a $\theta\in G(X)$ i $w\colon Y\to\{0,1\}$. Per a la identitat, $G(\id_X)(\theta)(v)=\theta(v\comp\id_X)=\theta(v)$. Per a la composició, si $f\colon X\to Y$, $g\colon Y\to Z$ i $u\colon Z\to\{0,1\}$,
\[
\begin{aligned}
G(g\comp f)(\theta)(u)
&=\theta\bigl(u\comp(g\comp f)\bigr)
=\theta\bigl((u\comp g)\comp f\bigr)\\
&=G(f)(\theta)(u\comp g)
=\bigl(G(g)\comp G(f)\bigr)(\theta)(u),
\end{aligned}
\]
on el pas clau és l'associativitat de la composició.

\textbf{Lema de substitució.} Per a tota $\varphi\in F(X)$ i tota $w\colon Y\to\{0,1\}$,
\[
(w\comp f)\models\varphi
\quad\Longleftrightarrow\quad
w\models\varphi[f].
\]
Per recursió sobre $\varphi$. Si $\varphi$ és una variable $x$, els dos costats diuen que $w(f(x))=1$, perquè $(w\comp f)(x)=w(f(x))$ i $x[f]=f(x)$. Si $\varphi=\psi\land\chi$, aleshores $(w\comp f)\models\varphi$ si i només si $(w\comp f)\models\psi$ i $(w\comp f)\models\chi$; per hipòtesi de recursió, si i només si $w\models\psi[f]$ i $w\models\chi[f]$, és a dir, si i només si $w\models\psi[f]\land\chi[f]=\varphi[f]$. Les altres connectives són anàlogues.

\textbf{Naturalitat.} Avaluem els dos camins del quadrat de naturalitat sobre $\varphi\in F(X)$ i $w\colon Y\to\{0,1\}$:
\[
\begin{aligned}
\bigl(\eta_Y(F(f)(\varphi))\bigr)(w)&=\eta_Y(\varphi[f])(w)=[\,w\models\varphi[f]\,],\\
\bigl(G(f)(\eta_X(\varphi))\bigr)(w)&=\eta_X(\varphi)(w\comp f)=[\,(w\comp f)\models\varphi\,],
\end{aligned}
\]
on $[\,\cdot\,]$ val $1$ o $0$ segons que la condició sigui certa o falsa. Pel lema de substitució, les dues expressions coincideixen per a tota $w$, de manera que $\eta_Y\comp F(f)=G(f)\comp\eta_X$ per a tota $f$, i $\eta$ és natural.
\end{solucio}

\begin{exercici}
Comprova en detall que l'avaluació $\eta_V\colon V\to V^{**}$, $\eta_V(v)(\varphi)=\varphi(v)$, és una transformació natural de la identitat cap al functor bidual (o doble dual), sobre la categoria dels espais vectorials sobre un cos $\mathbb K$ i les aplicacions lineals.
\end{exercici}

\begin{solucio}
Recordem primer com actua el bidual sobre morfismes. Donada una aplicació lineal $T\colon V\to W$, el seu dual $T^*\colon W^*\to V^*$ envia $\psi\mapsto\psi\comp T$, i el bidual $T^{**}\colon V^{**}\to W^{**}$ envia $\Phi\mapsto\Phi\comp T^*$.

La condició de naturalitat per a $T\colon V\to W$ és el quadrat
\[
T^{**}\comp\eta_V=\eta_W\comp T.
\]
Comprovem que tots dos costats coincideixen aplicats a un vector $v\in V$ i avaluats en una forma $\psi\in W^*$. D'una banda,
\[
\begin{aligned}
\big(T^{**}(\eta_V(v))\big)(\psi)
&=\big(\eta_V(v)\comp T^*\big)(\psi)
=\eta_V(v)\big(T^*(\psi)\big)\\
&=\eta_V(v)(\psi\comp T)
=(\psi\comp T)(v)
=\psi(T(v)).
\end{aligned}
\]
De l'altra,
\[
\big(\eta_W(T(v))\big)(\psi)=\psi\big(T(v)\big).
\]
Els dos resultats coincideixen per a tot $v$ i tot $\psi$, de manera que $T^{**}\comp\eta_V=\eta_W\comp T$. Per tant $\eta$ és natural. Cap pas no ha requerit triar una base: aquesta és la manifestació concreta de la naturalitat.
\end{solucio}

\begin{exercici}
Siguin $P$ i $Q$ preordres i $F,G,H\colon P\to Q$ aplicacions monòtones. Interpreta les transformacions naturals $F\Rightarrow G$ i comprova que la composició vertical correspon a la transitivitat de la desigualtat puntual.
\end{exercici}

\begin{solucio}
Com hem vist a la teoria, hi ha una transformació natural $F\Rightarrow G$ si i només si $F(x)\leq G(x)$ per a tot $x\in P$; i quan existeix, és única, perquè entre dos objectes d'un preordre hi ha com a màxim una fletxa.

Si tenim $\eta\colon F\Rightarrow G$ i $\theta\colon G\Rightarrow H$, això vol dir $F(x)\leq G(x)$ i $G(x)\leq H(x)$ per a tot $x$. La composició vertical $\theta\comp\eta\colon F\Rightarrow H$ existeix perquè, per transitivitat, $F(x)\leq H(x)$. Així, la composició vertical de transformacions naturals entre aplicacions monòtones és exactament la transitivitat de la relació \enquote{puntualment menor o igual}. La categoria de functors $Q^{P}$ resulta ser, en aquest cas, un preordre: el preordre de les aplicacions monòtones $P\to Q$ amb l'ordre puntual.
\end{solucio}

\section{Isomorfismes naturals}

\begin{exercici}[Ampliació, per a lectors amb àlgebra lineal]\label{ex:dual-no-natural}
Comprova que tot espai vectorial real $V$ de dimensió finita és isomorf al seu dual $V^*$, però que aquesta identificació no és \emph{canònica}: no existeix cap família d'isomorfismes $\eta_V\colon V\to V^*$ (per a $V$ espai vectorial real de dimensió finita) que sigui invariant sota tots els isomorfismes lineals, en el sentit que, per a tot isomorfisme $T\colon V\to W$, es compleixi $\eta_W\comp T=(T^{-1})^*\comp\eta_V$.
\end{exercici}

\begin{solucio}
Si $V$ té dimensió finita $n$, aleshores $V^*$ també té dimensió $n$, i dos espais reals de la mateixa dimensió finita són isomorfs. Un isomorfisme concret s'obté triant una base $\{e_1,\dots,e_n\}$ de $V$ i enviant $e_i$ a la base dual $\{e_i^*\}$. L'isomorfisme depèn de la base triada.

Cal notar primer que la comparació s'ha de plantejar amb cura, perquè la identitat és un functor covariant $\cat{Vect}_{\mathbb R}\to\cat{Vect}_{\mathbb R}$ i el dual és \emph{contravariant}, és a dir, un functor $\cat{Vect}_{\mathbb R}\op\to\cat{Vect}_{\mathbb R}$; entre functors de variància diferent no té sentit la noció de quadrat de naturalitat que hem definit. Per això formulem la invariància com a compatibilitat amb els \emph{isomorfismes} lineals, on el dual $(T^{-1})^*$ i $T$ van en el mateix sentit.

Que no existeixi una tal família es pot veure amb l'escalament. Sigui $V\neq 0$ i prenem l'homotècia $T=2\,\id_V$, que és un isomorfisme lineal (aquí fem servir que treballem sobre $\mathbb R$, on $2\neq 0$). Aleshores $T^{-1}=\tfrac12\,\id_V$ i $(T^{-1})^*=\tfrac12\,\id_{V^*}$. La condició d'invariància demanaria
\[
\eta_V\comp(2\,\id_V)=\bigl(\tfrac12\,\id_{V^*}\bigr)\comp\eta_V,
\]
és a dir, $2\,\eta_V=\tfrac12\,\eta_V$, o equivalentment $\tfrac32\,\eta_V=0$. Com que $\tfrac32\neq 0$, això força $\eta_V=0$, que no és un isomorfisme. Per tant no hi ha cap identificació $V\cong V^*$ invariant sota tots els isomorfismes lineals.

\begin{observacio}
El punt clau de l'argument és triar un escalar $\lambda$ amb $\lambda\neq\lambda^{-1}$, és a dir $\lambda^2\neq 1$: de la igualtat $\lambda\,\eta_V=\lambda^{-1}\eta_V$ només se'n dedueix $\eta_V=0$ quan $\lambda-\lambda^{-1}\neq 0$. No n'hi ha prou, doncs, amb $\lambda\neq 0,1$: sobre $\mathbb R$ l'escalar $\lambda=-1$ compliria $\lambda\neq 0,1$ però $\lambda^2=1$, i no serviria. Per això prenem $\lambda=2$, amb $\lambda^2=4\neq 1$. Sobre cossos petits pot no existir cap escalar amb $\lambda^2\neq 1$ (per exemple a $\mathbb F_2$ i $\mathbb F_3$ tots els no nuls compleixen $\lambda^2=1$), i l'argument de l'escalament s'ha d'ajustar; per això aquí ens restringim a $\mathbb R$.
\end{observacio}

Aquest contrast amb el bidual ---que, restringit als espais de dimensió finita, sí que és naturalment isomorf a la identitat--- és l'exemple canònic de la diferència entre \enquote{isomorf} i \enquote{naturalment isomorf}.
\end{solucio}

\begin{exercici}
Comprova que la relació \enquote{ser naturalment isomorf} és una relació d'equivalència sobre els functors $\cat{C}\to\cat{D}$.
\end{exercici}

\begin{solucio}
Cal veure reflexivitat, simetria i transitivitat.

\textbf{Reflexivitat}: la transformació identitat $\id_F\colon F\Rightarrow F$ té components $\id_{F(A)}$, que són isomorfismes; per tant $F\cong F$.

\textbf{Simetria}: si $\eta\colon F\Rightarrow G$ és un isomorfisme natural, hem vist a la teoria que els inversos $\eta_A^{-1}$ formen una transformació natural $\eta^{-1}\colon G\Rightarrow F$, també amb components isomorfismes; per tant $G\cong F$.

\textbf{Transitivitat}: si $\eta\colon F\Rightarrow G$ i $\theta\colon G\Rightarrow H$ són isomorfismes naturals, la composició vertical $\theta\comp\eta$ té components $\theta_A\comp\eta_A$, composició de dos isomorfismes, que és un isomorfisme. Per tant $F\cong H$.

De fet, aquesta és simplement la relació d'isomorfisme entre objectes de la categoria de functors $\cat{D}^{\cat{C}}$, que ja sabem que és una relació d'equivalència.
\end{solucio}

\section{La categoria de functors}

\begin{exercici}
Descriu la categoria de functors $\cat{D}^{\cat{1}}$, on $\cat{1}$ és la categoria amb un sol objecte i només la identitat.
\end{exercici}

\begin{solucio}
Un functor $F\colon\cat{1}\to\cat{D}$ tria un únic objecte $F(\ast)=D$ de $\cat{D}$ (i envia l'única identitat a $\id_D$). Per tant, els objectes de $\cat{D}^{\cat{1}}$ es corresponen amb els objectes de $\cat{D}$.

Una transformació natural $\eta\colon F\Rightarrow G$ entre dos tals functors té un únic component $\eta_\ast\colon F(\ast)\to G(\ast)$, és a dir, un morfisme $D\to D'$ de $\cat{D}$; la condició de naturalitat per a l'única fletxa (la identitat) es compleix trivialment. Per tant, els morfismes de $\cat{D}^{\cat{1}}$ es corresponen amb els morfismes de $\cat{D}$.

Concloem que $\cat{D}^{\cat{1}}$ és \emph{isomorfa} a $\cat{D}$: prendre functors des de la categoria terminal $\cat{1}$ recupera la categoria original.
\end{solucio}

\begin{exercici}
Comprova que un functor $F\colon\cat{2}\to\cat{D}$ és un objecte de la categoria de functors $\cat{D}^{\cat{2}}$, i que aquesta categoria és la \textbf{categoria de fletxes} de $\cat{D}$: els seus objectes són els morfismes de $\cat{D}$ i els seus morfismes són els quadrats commutatius.
\end{exercici}

\begin{solucio}
Ja hem vist que un functor $\cat{2}\to\cat{D}$ és un morfisme de $\cat{D}$, i que una transformació natural entre dos d'aquests functors és un quadrat commutatiu. Per tant, els objectes de $\cat{D}^{\cat{2}}$ són els morfismes de $\cat{D}$, i els morfismes de $\cat{D}^{\cat{2}}$ són els quadrats commutatius, amb la composició vertical corresponent a apilar quadrats. Aquesta categoria s'anomena \textbf{categoria de fletxes} de $\cat{D}$ i es denota sovint $\cat{D}^{\to}$. És un primer exemple de com les categories de functors produeixen construccions útils a partir de categories índex molt simples.
\end{solucio}

\section{Equivalència de categories}

\begin{exercici}
Comprova que la categoria dels conjunts finits és equivalent a la seva subcategoria plena formada pels conjunts $\underline{0}=\varnothing$ i $\underline{n}=\{1,\dots,n\}$ per a $n\geq 1$.
\end{exercici}

\begin{solucio}
Sigui $\cat{S}$ aquesta subcategoria plena i $\iota\colon\cat{S}\hookrightarrow\cat{FinSet}$ la inclusió. A l'exercici resolt~\ref{pr-esquelet-finset} del capítol~\ref{cap:functors} vam comprovar que $\iota$ és ple, fidel i essencialment exhaustiu. Pel teorema~\ref{teo:caract-equivalencia}, $\iota$ forma part d'una equivalència, i per tant $\cat{FinSet}\simeq\cat{S}$.

Per construir un quasiinvers $G\colon\cat{FinSet}\to\cat{S}$, la demostració del teorema tria, per a cada conjunt finit $X$ de $n$ elements, una bijecció $\varepsilon_X\colon\underline{n}\to X$, i posa $G(X)=\underline{n}$. Com que $\cat{FinSet}$ és gran, aquesta tria simultània fa servir l'axioma de l'elecció en la forma que explica l'observació que segueix el teorema. Intuïtivament: per fer teoria de categories amb conjunts finits, n'hi ha prou de considerar un conjunt de cada mida.
\end{solucio}

\begin{exercici}\label{pr-equivalencia-preserva}
Comprova que una equivalència de categories preserva monomorfismes, isomorfismes i objectes terminals. Explica per què, en canvi, \emph{no} preserva propietats estrictes, com ara la de tenir exactament un objecte.
\end{exercici}

\begin{solucio}
Sigui $F\colon\cat{C}\to\cat{D}$ part d'una equivalència, amb quasiinvers $G$ i isomorfismes naturals $\eta\colon\id_{\cat{C}}\Rightarrow G\comp F$ i $\varepsilon\colon F\comp G\Rightarrow\id_{\cat{D}}$, orientats com a la definició~\ref{def:equivalencia}. Tots dos són isomorfismes component a component, amb inversos $\eta^{-1}$ i $\varepsilon^{-1}$.

\textbf{Isomorfismes}: com que $F$ és un functor, preserva isomorfismes (proposició~\ref{prop:functor-preserva-iso}); i com que és ple i fidel ---per ser part d'una equivalència---, també els reflecteix. Així, $f$ és un isomorfisme a $\cat{C}$ si i només si $F(f)$ ho és a $\cat{D}$.

\textbf{Monomorfismes}: sigui $f\colon A\to B$ mono a $\cat{C}$; vegem que $F(f)\colon F(A)\to F(B)$ és mono a $\cat{D}$. Siguin $u,v\colon Z\to F(A)$ dos morfismes de $\cat{D}$ amb
\[
F(f)\comp u=F(f)\comp v.
\]
Apliquem el quasiinvers $G$ i fem servir la seva functorialitat:
\[
G(F(f))\comp G(u)=G\bigl(F(f)\comp u\bigr)=G\bigl(F(f)\comp v\bigr)=G(F(f))\comp G(v).
\]
Ara traslladem aquesta igualtat de $G(F(A))$ a $A$ mitjançant la naturalitat de $\eta$. Per a $f\colon A\to B$, el quadrat de naturalitat de $\eta$ dona
\[
G(F(f))\comp\eta_A=\eta_B\comp f,
\]
i, en ser $\eta_A,\eta_B$ isomorfismes, $G(F(f))=\eta_B\comp f\comp\eta_A^{-1}$. Substituint a la igualtat anterior i cancel·lant l'isomorfisme $\eta_B$ per l'esquerra,
\[
f\comp\bigl(\eta_A^{-1}\comp G(u)\bigr)=f\comp\bigl(\eta_A^{-1}\comp G(v)\bigr).
\]
Com que $f$ és mono a $\cat{C}$, en resulta $\eta_A^{-1}\comp G(u)=\eta_A^{-1}\comp G(v)$, i cancel·lant l'isomorfisme $\eta_A^{-1}$ obtenim $G(u)=G(v)$. Finalment, $G$ és fidel ---per ser part d'una equivalència---, de manera que $u=v$. Per tant $F(f)$ és mono.

\textbf{Objectes terminals}: sigui $1$ terminal a $\cat{C}$; vegem que $F(1)$ és terminal a $\cat{D}$. Sigui $D$ un objecte qualsevol de $\cat{D}$.

\emph{Existència}: com que $1$ és terminal, hi ha un únic morfisme $!\colon G(D)\to 1$ a $\cat{C}$. Aleshores el morfisme
\[
D\xrightarrow{\ \varepsilon_D^{-1}\ }F(G(D))\xrightarrow{\ F(!)\ }F(1)
\]
és un morfisme $D\to F(1)$, cosa que en garanteix l'existència.

\emph{Unicitat}: siguin $h,h'\colon D\to F(1)$ dos morfismes. Aplicant $G$, obtenim $G(h),G(h')\colon G(D)\to G(F(1))$. Component amb l'isomorfisme $\eta_1^{-1}\colon G(F(1))\to 1$, els morfismes $\eta_1^{-1}\comp G(h)$ i $\eta_1^{-1}\comp G(h')$ van de $G(D)$ a $1$ i, per tant, coincideixen per la unicitat cap a l'objecte terminal $1$. Com que $\eta_1^{-1}$ és invertible, resulta $G(h)=G(h')$, i com que $G$ és fidel, $h=h'$. Per tant $F(1)$ és terminal.

\textbf{Per què no tota propietat}: una equivalència \emph{no} preserva les propietats que depenen de la \emph{igualtat} d'objectes. Per exemple, \enquote{tenir exactament un objecte} o \enquote{ser esquelètica} no es conserven: la categoria amb un sol objecte i la seva versió amb dos objectes isomorfs són equivalents, però la primera té un objecte i la segona en té dos.

Convé distingir, a més, dues menes de nocions entre les que hem comprovat. Ser isomorfisme i ser monomorfisme són \emph{propietats de morfismes}, definides per equacions o per cancel·lació; l'equivalència les preserva, i de la mateixa manera es comprova que preserva els epimorfismes. Ser objecte terminal és, en canvi, una \emph{propietat universal} d'un objecte, que el caracteritza només llevat d'isomorfisme; el que hem vist és que la imatge d'un terminal torna a ser terminal. Més endavant trobarem altres construccions universals ---objectes inicials, límits i colímits--- i en comprovarem el comportament sota equivalència cas per cas.
\end{solucio}


\section{Composició horitzontal}

\begin{exercici}\label{pr-comp-horitzontal}
Siguin $\eta\colon F\Rightarrow G$, amb $F,G\colon\cat{C}\to\cat{D}$, i $\theta\colon K\Rightarrow L$, amb $K,L\colon\cat{D}\to\cat{E}$, com a la definició de composició horitzontal. Escriu els tipus de
\[
\theta_{G(A)}\comp K(\eta_A)
\qquad\text{i}\qquad
L(\eta_A)\comp\theta_{F(A)},
\]
explica per què coincideixen per la naturalitat de $\theta$, i verifica el quadrat de naturalitat de $\theta\ast\eta\colon K\comp F\Rightarrow L\comp G$.
\end{exercici}

\begin{solucio}
\textbf{Els tipus.} El morfisme $\eta_A\colon F(A)\to G(A)$ viu a $\cat{D}$, i $K$ el porta a $K(\eta_A)\colon K(F(A))\to K(G(A))$ a $\cat{E}$. La component de $\theta$ a l'objecte $G(A)$ de $\cat{D}$ és $\theta_{G(A)}\colon K(G(A))\to L(G(A))$. Per tant
\[
\theta_{G(A)}\comp K(\eta_A)\colon K(F(A))\longrightarrow L(G(A)).
\]
Anàlogament, $\theta_{F(A)}\colon K(F(A))\to L(F(A))$ i $L(\eta_A)\colon L(F(A))\to L(G(A))$, de manera que
\[
L(\eta_A)\comp\theta_{F(A)}\colon K(F(A))\longrightarrow L(G(A)).
\]
Tots dos morfismes tenen, doncs, el mateix origen i el mateix destí, i té sentit preguntar-se si coincideixen.

\textbf{Per què coincideixen.} Són els dos camins del quadrat de naturalitat de $\theta$ per al morfisme $\eta_A\colon F(A)\to G(A)$ de $\cat{D}$, dibuixat a l'observació~\ref{obs:comp-horitzontal} de la Teoria:
\[
\theta_{G(A)}\comp K(\eta_A)=L(\eta_A)\comp\theta_{F(A)}.
\]
La naturalitat de $\theta$ diu que aquest quadrat commuta per a \emph{tot} morfisme de $\cat{D}$; en particular, per a $\eta_A$. La component $(\theta\ast\eta)_A$ està, doncs, ben definida, sigui quina sigui l'expressió que fem servir.

\textbf{Naturalitat de $\theta\ast\eta$.} Sigui $f\colon A\to B$ un morfisme de $\cat{C}$. Cal veure que
\[
\begin{tikzpicture}[baseline=(m.center)]
  \matrix (m) [matrix of math nodes, row sep=2.8em, column sep=4.6em] {
    K(F(A)) & K(F(B)) \\
    L(G(A)) & L(G(B)) \\
  };
  \draw[-{Stealth[scale=1.1]}] (m-1-1) -- node[above] {$K(F(f))$} (m-1-2);
  \draw[-{Stealth[scale=1.1]}] (m-1-1) -- node[left] {$(\theta\ast\eta)_A$} (m-2-1);
  \draw[-{Stealth[scale=1.1]}] (m-1-2) -- node[right] {$(\theta\ast\eta)_B$} (m-2-2);
  \draw[-{Stealth[scale=1.1]}] (m-2-1) -- node[below] {$L(G(f))$} (m-2-2);
\end{tikzpicture}
\qquad
(\theta\ast\eta)_B\comp K(F(f))=L(G(f))\comp(\theta\ast\eta)_A
\]
commuta. Fent servir l'expressió $\theta_{G(-)}\comp K(\eta_-)$ per a les components,
\[
\begin{aligned}
(\theta\ast\eta)_B\comp K(F(f))
&=\theta_{G(B)}\comp K(\eta_B)\comp K(F(f))
=\theta_{G(B)}\comp K\bigl(\eta_B\comp F(f)\bigr)\\
&=\theta_{G(B)}\comp K\bigl(G(f)\comp\eta_A\bigr)
=\theta_{G(B)}\comp K(G(f))\comp K(\eta_A)\\
&=L(G(f))\comp\theta_{G(A)}\comp K(\eta_A)
=L(G(f))\comp(\theta\ast\eta)_A.
\end{aligned}
\]
Hem fet servir, per ordre: la functorialitat de $K$, la naturalitat de $\eta$ per a $f$, de nou la functorialitat de $K$, i la naturalitat de $\theta$ per al morfisme $G(f)$ de $\cat{D}$. Cada naturalitat s'utilitza una sola vegada, i cadascuna al seu nivell: la de $\eta$ dins de $\cat{D}$, i la de $\theta$ dins de $\cat{E}$.
\end{solucio}


\ifpublicacioparcial\else
\chapter[Representabilitat i Yoneda]{\texorpdfstring{Propietats universals,\\ representabilitat i Yoneda}{Propietats universals, representabilitat i Yoneda}}

Els exercicis d'aquest capítol treballen les propietats universals a través de la representabilitat i el lema de Yoneda. L'objectiu és guanyar fluïdesa amb els functors representables, amb el càlcul de transformacions naturals que en surten, i amb l'argument ---repetitiu, un cop entès--- que dedueix unicitat llevat d'isomorfisme.

\section{Objectes inicials i terminals}

\begin{exercici}\label{pr-producte-terminal}
Siguin $A$ i $B$ dos objectes d'una categoria $\cat{C}$. Un \textbf{con} sobre $A$ i $B$ és una terna $(X,f,g)$ amb $f\colon X\to A$ i $g\colon X\to B$, i un morfisme de cons $(X,f,g)\to(X',f',g')$ és un morfisme $h\colon X\to X'$ de $\cat{C}$ amb $f'\comp h=f$ i $g'\comp h=g$. Comprova que els cons sobre $A$ i $B$ formen una categoria $\mathrm{Con}(A,B)$, i que un con $(P,p_1,p_2)$ és un producte de $A$ i $B$ si i només si és un objecte terminal de $\mathrm{Con}(A,B)$. Interpreta el resultat en la categoria d'un preordre $(Q,\le)$.
\end{exercici}

\begin{solucio}
\textbf{És una categoria.} La identitat de $(X,f,g)$ és $\id_X$, que compleix $f\comp\id_X=f$ i $g\comp\id_X=g$. Si $h\colon(X,f,g)\to(X',f',g')$ i $h'\colon(X',f',g')\to(X'',f'',g'')$ són morfismes de cons, la composta $h'\comp h$ també ho és, perquè
\[
f''\comp(h'\comp h)=(f''\comp h')\comp h=f'\comp h=f,
\]
i igual amb $g$. L'associativitat i les lleis d'identitat s'hereten de $\cat{C}$, perquè la composició és la de $\cat{C}$.

\textbf{Producte i objecte terminal.} Un morfisme de cons $h\colon(X,f,g)\to(P,p_1,p_2)$ és exactament un morfisme $h\colon X\to P$ que fa commutar els dos triangles,
\[
\begin{tikzpicture}[baseline=(m.center)]
  \matrix (m) [matrix of math nodes, row sep=2.6em, column sep=2.8em] {
     & X & \\
    A & P & B \\
  };
  \draw[-{Stealth[scale=1.0]}] (m-1-2) -- node[above left] {$f$} (m-2-1);
  \draw[-{Stealth[scale=1.0]}] (m-1-2) -- node[above right] {$g$} (m-2-3);
  \draw[-{Stealth[scale=1.0]},dashed] (m-1-2) -- node[right] {$h$} (m-2-2);
  \draw[-{Stealth[scale=1.0]}] (m-2-2) -- node[below] {$p_1$} (m-2-1);
  \draw[-{Stealth[scale=1.0]}] (m-2-2) -- node[below] {$p_2$} (m-2-3);
\end{tikzpicture}
\qquad
p_1\comp h=f,\quad p_2\comp h=g.
\]
Que $(P,p_1,p_2)$ sigui terminal a $\mathrm{Con}(A,B)$ diu que, per a cada con $(X,f,g)$, hi ha exactament un morfisme de cons cap a $(P,p_1,p_2)$, és a dir, exactament un $h$ amb $p_1\comp h=f$ i $p_2\comp h=g$. Això és, paraula per paraula, la definició de producte (subsecció~\ref{subsec:producte-coproducte}). Les dues afirmacions són, doncs, la mateixa.

Com a conseqüència, la unicitat dels productes es dedueix de la dels objectes terminals (proposició~\ref{prop:unicitat-terminal}): dos productes de $A$ i $B$ són isomorfs per un únic isomorfisme \emph{de cons}, és a dir, compatible amb les projeccions. L'objecte $P$ tot sol, en canvi, pot tenir altres automorfismes que no respectin $p_1$ i $p_2$; per exemple, a $\cat{Set}$, l'intercanvi $(a,b)\mapsto(b,a)$ de $A\times A$.

\textbf{En un preordre.} A la categoria d'un preordre $(Q,\le)$, un con sobre $a$ i $b$ és un element $x$ amb $x\le a$ i $x\le b$, és a dir, una cota inferior comuna, i entre dos cons hi ha com a màxim un morfisme, que existeix quan $x\le x'$. Un objecte terminal de $\mathrm{Con}(a,b)$ és, doncs, una cota inferior $p$ tal que tota altra cota inferior $x$ compleix $x\le p$: la \emph{màxima cota inferior}, o ínfim, de $a$ i $b$. Els productes en un preordre són els ínfims, i en un ordre parcial són únics; en un preordre, únics llevat de la relació $p\le p'\le p$.
\end{solucio}

\begin{exercici}
Comprova que un objecte $Z$ és alhora inicial i terminal en una categoria si i només si, per a cada objecte $X$, hi ha exactament un morfisme $Z\to X$ i exactament un morfisme $X\to Z$. Un tal objecte s'anomena \emph{objecte zero}. Comprova que a $\cat{Pos}$ no n'hi ha cap i, per a qui conegui els grups, que a $\cat{Grp}$ el grup trivial n'és un.
\end{exercici}

\begin{solucio}
La primera afirmació és la conjunció literal de les dues definicions: $Z$ inicial dona el morfisme únic $Z\to X$, i $Z$ terminal, el morfisme únic $X\to Z$.

A $\cat{Pos}$, l'objecte inicial és el conjunt ordenat buit i el terminal és el d'un sol element. Un objecte zero hauria de ser isomorf a tots dos, i aquests dos no són isomorfs, perquè no hi ha cap aplicació $\{\ast\}\to\varnothing$. Per tant, $\cat{Pos}$ no té objecte zero.

Per al grup trivial $1=\{e\}$: donat un grup $G$, hi ha un únic homomorfisme $1\to G$ (l'element neutre ha d'anar al neutre), i un únic homomorfisme $G\to 1$ (tot va al neutre). Per tant $1$ és inicial i terminal a $\cat{Grp}$, és a dir, un objecte zero.
\end{solucio}

\section{Categories coma}

\begin{exercici}\label{pr-coma-lliure}
Sigui $G$ un graf, i sigui $U\colon\cat{Cat}\to\cat{Graph}$ el functor graf subjacent, on $\cat{Cat}$ té per objectes les categories petites. Descriu explícitament la categoria coma $(G\downarrow U)$: els seus objectes, els seus morfismes, la composició i les identitats. Demostra que la parella $(\Free(G),\eta_G)$, on $\eta_G\colon G\to U(\Free(G))$ és la inserció dels generadors, és un objecte inicial de $(G\downarrow U)$.
\end{exercici}

\begin{solucio}
\textbf{Objectes i morfismes.} Un objecte de $(G\downarrow U)$ és una parella $(\cat{D},\varphi)$ formada per una categoria petita $\cat{D}$ i un morfisme de grafs $\varphi\colon G\to U(\cat{D})$. Un morfisme $(\cat{D},\varphi)\to(\cat{D}',\varphi')$ és un functor $H\colon\cat{D}\to\cat{D}'$ que fa commutar el triangle
\[
\begin{tikzpicture}[baseline=(m.center)]
  \matrix (m) [matrix of math nodes, row sep=2.6em, column sep=2.4em] {
     & G & \\
    U(\cat{D}) & & U(\cat{D}') \\
  };
  \draw[-{Stealth[scale=1.0]}] (m-1-2) -- node[above left] {$\varphi$} (m-2-1);
  \draw[-{Stealth[scale=1.0]}] (m-1-2) -- node[above right] {$\varphi'$} (m-2-3);
  \draw[-{Stealth[scale=1.0]}] (m-2-1) -- node[below] {$U(H)$} (m-2-3);
\end{tikzpicture}
\qquad
U(H)\comp\varphi=\varphi'.
\]

\textbf{És una categoria.} La identitat de $(\cat{D},\varphi)$ és $\id_{\cat{D}}$, perquè $U(\id_{\cat{D}})\comp\varphi=\varphi$. Si $H\colon(\cat{D},\varphi)\to(\cat{D}',\varphi')$ i $H'\colon(\cat{D}',\varphi')\to(\cat{D}'',\varphi'')$, la composta $H'\comp H$ també és un morfisme, per la functorialitat de $U$:
\[
U(H'\comp H)\comp\varphi=U(H')\comp U(H)\comp\varphi=U(H')\comp\varphi'=\varphi''.
\]
L'associativitat i les lleis d'identitat s'hereten de $\cat{Cat}$.

\textbf{$(\Free(G),\eta_G)$ és inicial.} Sigui $(\cat{D},\varphi)$ un objecte qualsevol. Un morfisme $(\Free(G),\eta_G)\to(\cat{D},\varphi)$ és un functor $\bar\varphi\colon\Free(G)\to\cat{D}$ amb $U(\bar\varphi)\comp\eta_G=\varphi$, és a dir, que coincideix amb $\varphi$ sobre els vèrtexs i sobre els camins de longitud~$1$.
\begin{itemize}
  \item \emph{Existència.} Definim $\bar\varphi$ sobre els objectes com $\varphi$, sobre el camí buit en un vèrtex $x$ com $\id_{\varphi(x)}$, i sobre un camí $(a_1,\dots,a_n)$ recorregut en aquest ordre com la composta $\varphi(a_n)\comp\cdots\comp\varphi(a_1)$. Com que la composició de $\Free(G)$ és la concatenació, $\bar\varphi$ respecta composicions i identitats.
  \item \emph{Unicitat.} Tot camí és la composta de les seves arestes, i tot camí buit és una identitat. Un functor amb $U(\bar\varphi)\comp\eta_G=\varphi$ queda, per tant, determinat sobre tots els morfismes per la seva acció sobre les arestes, que està fixada per $\varphi$.
\end{itemize}
Així, des de $(\Free(G),\eta_G)$ hi ha exactament un morfisme cap a cada objecte de $(G\downarrow U)$: és un objecte inicial. Observem que no ha calgut la noció d'adjunció; és aquesta propietat la que, al capítol d'adjuncions, farà de $\Free$ un adjunt per l'esquerra de $U$.
\end{solucio}

\section{Propietats universals com a representabilitat}

\begin{exercici}
Comprova que un objecte $1$ és terminal si i només si el functor $\Hom(-,1)$ és naturalment isomorf al functor constant $\Delta_{\{\ast\}}\colon\cat{C}\op\to\cat{Set}$ (exemple~\ref{ex:functor-constant}).
\end{exercici}

\begin{solucio}
Suposem $1$ terminal. Aleshores, per a cada objecte $X$, el conjunt $\Hom(X,1)$ té exactament un element, de manera que hi ha una bijecció $\Hom(X,1)\cong\{\ast\}$. Aquestes bijeccions són automàticament naturals: l'única aplicació possible cap a un conjunt d'un punt fa commutar qualsevol quadrat. Per tant $\Hom(-,1)$ és naturalment isomorf al functor constant $\Delta_{\{\ast\}}$.

Recíprocament, si $\Hom(-,1)\cong\Delta_{\{\ast\}}$, aleshores cada $\Hom(X,1)$ té exactament un element, és a dir, hi ha exactament un morfisme $X\to 1$ per a cada $X$, i $1$ és terminal.
\end{solucio}

\begin{exercici}
Sigui $\cat{C}$ una categoria localment petita amb un objecte terminal $1$. Comprova que el \textbf{functor de punts globals}
\[
\Hom(1,-)\colon\cat{C}\longrightarrow\cat{Set}
\]
envia cada objecte $X$ a $\Hom(1,X)$. Aquests morfismes s'anomenen \emph{punts globals} de $X$. Calcula aquest conjunt a $\cat{Set}$.
\end{exercici}

\begin{solucio}
Per definició, $\Hom(1,-)$ envia $X$ a $\Hom(1,X)$, el conjunt de morfismes des de l'objecte terminal, és a dir, dels punts globals de $X$. A $\cat{Set}$, l'objecte terminal és un conjunt d'un sol element $\{\ast\}$, i una aplicació $\{\ast\}\to X$ equival a triar un element de $X$: la imatge de $\ast$. Per tant $\Hom(1,X)\cong X$ de manera natural, i a $\cat{Set}$ els punts globals recuperen exactament els elements. En general, però, $\Hom(1,-)$ pot perdre informació: hi ha categories on objectes diferents tenen els mateixos punts globals.
\end{solucio}

\section{El lema de Yoneda}

\begin{exercici}
Sigui $\cat{C}$ localment petita i $A$ un objecte. Aplicant el lema de Yoneda amb $P=\Hom(-,A)$, comprova que les úniques transformacions naturals $\Hom(-,A)\Rightarrow\Hom(-,A)$ que hi ha es corresponen amb els endomorfismes de $A$, i que la identitat correspon a $\id_A$.
\end{exercici}

\begin{solucio}
El lema de Yoneda dona una bijecció
\[
\Nat\big(\Hom(-,A),\Hom(-,A)\big)\cong\Hom(-,A)(A)=\Hom(A,A),
\]
que envia $\eta$ a $\eta_A(\id_A)$. Així, cada transformació natural de $\Hom(-,A)$ en ell mateix correspon a un endomorfisme de $A$. La transformació identitat $\id_{\Hom(-,A)}$ té components la identitat de cada $\Hom(X,A)$; el seu valor a $\id_A$ és $\id_A$, de manera que correspon a l'endomorfisme identitat. Aquest càlcul és el nucli de la demostració que la immersió de Yoneda és plena i fidel.
\end{solucio}

\begin{exercici}\label{pr-yoneda-parts}
Sigui $A$ un conjunt i sigui $\mathcal P\colon\cat{Set}\op\to\cat{Set}$ el functor de parts contravariant. Descriu explícitament, mitjançant la bijecció del lema de Yoneda, les transformacions naturals $\Hom_{\cat{Set}}(-,A)\Rightarrow\mathcal P$. Quina transformació correspon a un subconjunt $B\subseteq A$? Comprova directament que la bijecció recupera $B$, i calcula-la en el cas $A=\{0,1\}$, $B=\{1\}$, sobre l'aplicació $\varphi\colon\{a,b,c\}\to A$ amb $\varphi(a)=\varphi(c)=1$ i $\varphi(b)=0$.
\end{exercici}

\begin{solucio}
El lema de Yoneda dona una bijecció
\[
\Phi\colon\Nat\bigl(\Hom(-,A),\mathcal P\bigr)\longrightarrow\mathcal P(A),
\qquad
\Phi(\alpha)=\alpha_A(\id_A),
\]
amb inversa $\Psi$, que envia un element $x\in\mathcal P(A)$ a la transformació $\Psi(x)_X(\varphi)=\mathcal P(\varphi)(x)$. Aquí un element de $\mathcal P(A)$ és un subconjunt $B\subseteq A$, i $\mathcal P(\varphi)$ és l'antiimatge. Per tant, la transformació natural que correspon a $B$ és
\[
\Psi(B)_X\colon\Hom(X,A)\to\mathcal P(X),
\qquad
\varphi\longmapsto\varphi^{-1}(B).
\]
És a dir: tota assignació natural d'aquest tipus s'obté per antiimatge d'un únic subconjunt fix de $A$, i subconjunts diferents donen transformacions naturals diferents.

\textbf{La bijecció recupera $B$.} En un sentit, $\Phi(\Psi(B))=\Psi(B)_A(\id_A)=\id_A^{-1}(B)=B$. En l'altre, si $\alpha$ és una transformació natural qualsevol i $B=\alpha_A(\id_A)$, la naturalitat de $\alpha$ aplicada a $\varphi\colon X\to A$ i avaluada sobre $\id_A$ dona
\[
\alpha_X(\varphi)=\alpha_X(\id_A\comp\varphi)=\mathcal P(\varphi)\bigl(\alpha_A(\id_A)\bigr)=\varphi^{-1}(B),
\]
de manera que $\alpha=\Psi(B)$. La naturalitat de $\Psi(B)$ es veu també directament: per a $g\colon Y\to X$, $(\varphi\comp g)^{-1}(B)=g^{-1}\bigl(\varphi^{-1}(B)\bigr)$.

\textbf{Un càlcul.} Amb $A=\{0,1\}$ i $B=\{1\}$, la component a $X=\{a,b,c\}$ envia $\varphi$ a $\varphi^{-1}(\{1\})=\{a,c\}$. Observem que, sota l'isomorfisme natural $\mathcal P\cong\Hom(-,\{0,1\})$ de l'exemple~\ref{ex:isos-naturals-cap2}, el subconjunt $B$ correspon a la seva funció característica $\chi_B\colon A\to\{0,1\}$, i la transformació $\varphi\mapsto\varphi^{-1}(B)$ correspon a $\varphi\mapsto\chi_B\comp\varphi$: és el cas $P=\Hom(-,\{0,1\})$ del lema, en què les transformacions naturals $\Hom(-,A)\Rightarrow\Hom(-,\{0,1\})$ són les postcomposicions amb morfismes $A\to\{0,1\}$.

\textbf{Lectura lògica.} Un subconjunt $B\subseteq A$ és un \emph{predicat} sobre $A$, i $\varphi^{-1}(B)$ és el predicat sobre $X$ que s'obté substituint $\varphi$: $x\in\varphi^{-1}(B)$ si i només si $\varphi(x)\in B$. El lema de Yoneda diu, doncs, que les maneres naturals de produir predicats a partir d'aplicacions cap a $A$ són exactament les substitucions en un predicat fix sobre $A$. És la mateixa operació de substitució que, al capítol~\ref{cap:transformacions}, feia natural la interpretació de les fórmules.
\end{solucio}

\begin{exercici}\label{pr-yoneda-grafs}
Siguin $V,E\colon\cat{Graph}\to\cat{Set}$ els functors de vèrtexs i d'arestes, siguin $\mathbf V$ el graf d'un sol vèrtex sense arestes i $\mathbf E$ el graf $0\xrightarrow{e}1$. Calcula, amb la versió covariant del lema de Yoneda, $\Nat(E,V)$ i $\Nat(V,E)$, i interpreta les transformacions que hi surtin.
\end{exercici}

\begin{solucio}
Per l'exemple~\ref{ex:isos-naturals-cap2}, $V\cong\Hom_{\cat{Graph}}(\mathbf V,-)$ i $E\cong\Hom_{\cat{Graph}}(\mathbf E,-)$. La versió covariant del lema dona, per a qualsevol functor $Q\colon\cat{Graph}\to\cat{Set}$, una bijecció entre les transformacions naturals que surten de $\Hom(A,-)$ i els elements de $Q(A)$. Per tant,
\[
\Nat(E,V)\;\cong\;V(\mathbf E)=\{0,1\},
\qquad
\Nat(V,E)\;\cong\;E(\mathbf V)=\varnothing.
\]
Observem la inversió: el primer càlcul avalua el functor d'\emph{arribada}, $V$, en l'objecte que \emph{representa} el functor de sortida, $E$.

\textbf{Les dues transformacions $E\Rightarrow V$.} L'element $0\in V(\mathbf E)$ correspon a la transformació que, a cada graf $G$, envia una aresta a una aplicació $E_G\to V_G$ calculada així: l'aresta $\varepsilon$ correspon al morfisme $\mathbf E\to G$ que la tria, i aquest envia el vèrtex $0$ a l'origen de $\varepsilon$. És, doncs, l'aplicació origen $s_G$. Anàlogament, l'element $1$ dona l'aplicació destí $t_G$. La naturalitat de $s$ i de $t$ és exactament la condició que un morfisme de grafs preservi origen i destí. Encara que en un graf on totes les arestes són bucles $s_G=t_G$, les dues transformacions són diferents, perquè a $\mathbf E$ les components difereixen.

\textbf{Cap transformació $V\Rightarrow E$.} També es veu directament: una transformació natural $V\Rightarrow E$ hauria de tenir una component $V_{\mathbf V}=\{\ast\}\to E_{\mathbf V}=\varnothing$, que no existeix. No hi ha cap manera natural de triar una aresta per a cada vèrtex en tots els grafs alhora.
\end{solucio}

\section{Conseqüències}

\begin{exercici}\label{pr-yoneda-preordre}
Sigui $(Q,\le)$ un preordre vist com a categoria. Descriu el prefeix representable $\Hom(-,a)$ i calcula, amb el lema de Yoneda, les transformacions naturals $\Hom(-,a)\Rightarrow\Hom(-,b)$. Dedueix que
\[
a\le b\quad\Longleftrightarrow\quad{\downarrow}a\subseteq{\downarrow}b,
\qquad\text{on } {\downarrow}a=\{x\in Q\mid x\le a\},
\]
i interpreta el resultat a la categoria $\cat{Prov}_T$ d'un sistema deductiu.
\end{exercici}

\begin{solucio}
\textbf{El representable.} Com que entre dos objectes hi ha com a màxim una fletxa, $\Hom(x,a)$ és un conjunt d'un sol element si $x\le a$, i és buit altrament. El prefeix $\Hom(-,a)$ és, doncs, la funció característica del conjunt ${\downarrow}a$ dels elements que estan per sota de $a$: \enquote{val} en $x$ exactament quan $x\in{\downarrow}a$.

\textbf{Les transformacions naturals.} Pel lema de Yoneda (o per la immersió plena i fidel),
\[
\Nat\bigl(\Hom(-,a),\Hom(-,b)\bigr)\;\cong\;\Hom(a,b),
\]
que té un element si $a\le b$ i cap altrament. Ho podem veure també directament: una transformació natural té una component $\Hom(x,a)\to\Hom(x,b)$ per a cada $x$, i una aplicació d'un conjunt d'un sol element en un conjunt buit no existeix. Hi ha, doncs, una transformació natural exactament quan $x\le a$ implica $x\le b$ per a tot $x$, és a dir, quan ${\downarrow}a\subseteq{\downarrow}b$; i aleshores és única, i la naturalitat és automàtica, perquè tots els valors tenen com a màxim un element.

\textbf{La conclusió.} Comparant les dues descripcions, $a\le b$ si i només si ${\downarrow}a\subseteq{\downarrow}b$. La implicació d'esquerra a dreta és la transitivitat; la de dreta a esquerra s'obté prenent $x=a$, que és exactament l'avaluació a $\id_a$ del lema de Yoneda. Com a cas particular, $a\cong b$ si i només si ${\downarrow}a={\downarrow}b$: dos elements són isomorfs quan tenen els mateixos elements per sota.

\textbf{A $\cat{Prov}_T$.} Aquí ${\downarrow}p$ és el conjunt de les fórmules que dedueixen $p$, i el resultat diu que $p\vdash_T q$ si i només si tota fórmula que dedueix $p$ dedueix també $q$. Una fórmula queda determinada, llevat d'interdeduïbilitat, per les fórmules de les quals es dedueix.
\end{solucio}

\begin{exercici}\label{pr-yoneda-monoide}
Sigui $M$ un monoide i $\cat{B}M$ la categoria d'un sol objecte $\ast$ associada, amb composició $g\comp f=g\cdot f$ (escrivim $\cdot$ per a l'operació del monoide i $e$ per al neutre, que és $\id_\ast$). Comprova que un prefeix $X\colon(\cat{B}M)\op\to\cat{Set}$ és el mateix que un conjunt amb una \emph{acció per la dreta} de $M$, i descriu el prefeix representable $\Hom(-,\ast)$. Aplica el lema de Yoneda per calcular $\Nat\bigl(\Hom(-,\ast),X\bigr)$ i, en particular, $\Nat\bigl(\Hom(-,\ast),\Hom(-,\ast)\bigr)$.
\end{exercici}

\begin{solucio}
\textbf{Prefeixos i accions.} Un prefeix $X$ dona un conjunt $X(\ast)$, que escrivim $X$, i per a cada $g\in M$ una aplicació $X(g)\colon X\to X$. Escrivim $x\cdot g=X(g)(x)$. La functorialitat contravariant, $X(g\comp f)=X(f)\comp X(g)$ i $X(e)=\id$, diu exactament
\[
(x\cdot g)\cdot f=x\cdot(g\cdot f),
\qquad
x\cdot e=x:
\]
és una acció per la dreta. Una transformació natural entre dos prefeixos és una aplicació $\eta\colon X\to Y$ amb $\eta(x\cdot g)=\eta(x)\cdot g$, és a dir, una aplicació \emph{equivariant}.

\textbf{El representable.} $\Hom(\ast,\ast)=M$, i un morfisme $g$ actua per precomposició: $m\mapsto m\comp g=m\cdot g$. El prefeix representable és, doncs, $M$ amb l'acció per la dreta per multiplicació, l'\emph{acció regular}.

\textbf{Yoneda.} El lema dona una bijecció $\Nat\bigl(\Hom(-,\ast),X\bigr)\cong X$: l'aplicació equivariant $\eta\colon M\to X$ correspon a $\eta(e)$, i recíprocament $x\in X$ dona l'aplicació $m\mapsto x\cdot m$. Una aplicació equivariant des de l'acció regular queda determinada pel lloc on envia l'element neutre.

Per a $X=M$ regular, les aplicacions equivariants $M\to M$ són exactament les multiplicacions per l'esquerra $\lambda_a\colon m\mapsto a\cdot m$, una per a cada $a\in M$. Com que $\lambda_a\comp\lambda_b=\lambda_{a\cdot b}$ i $\lambda_e=\id$, la immersió de Yoneda identifica $M$ amb aquest monoide d'aplicacions de $M$ en si mateix, i el fet que sigui fidel diu que $a\mapsto\lambda_a$ és injectiva. Obtenim el \textbf{teorema de Cayley} per a monoides: tot monoide és isomorf a un monoide d'aplicacions d'un conjunt en si mateix, amb la composició.
\end{solucio}

\chapter[Límits i colímits]{Límits i colímits}

Els exercicis d'aquest capítol treballen les propietats universals dels límits i colímits concrets, i l'argument ---sempre el mateix--- que dedueix unicitat i factoritzacions. Convé dibuixar cada diagrama abans de començar.

\section{Diagrames i cons}

\begin{exercici}
Comprova que un con sobre un diagrama $D\colon\cat{I}\to\cat{C}$ amb vèrtex $A$ és el mateix que una transformació natural $\Delta_A\Rightarrow D$, on $\Delta_A$ és el functor constant amb valor $A$.
\end{exercici}

\begin{solucio}
El functor constant $\Delta_A$ envia cada objecte $i$ de $\cat{I}$ a $A$ i cada morfisme a $\id_A$. Una transformació natural $\alpha\colon\Delta_A\Rightarrow D$ té components $\alpha_i\colon\Delta_A(i)=A\to D_i$, i la condició de naturalitat per a un morfisme $u\colon i\to j$ diu
\[
D_u\comp\alpha_i=\alpha_j\comp\Delta_A(u)=\alpha_j\comp\id_A=\alpha_j.
\]
Aquesta és exactament la condició de con. Recíprocament, tota família amb aquesta propietat defineix una transformació natural. Així, els cons sobre $D$ amb vèrtex $A$ i les transformacions naturals $\Delta_A\Rightarrow D$ coincideixen.
\end{solucio}

\section{Productes i equalitzadors}

\begin{exercici}
Comprova que a $\cat{Set}$ el producte categòric coincideix amb el producte cartesià, verificant la propietat universal.
\end{exercici}

\begin{solucio}
Sigui $A\times B=\{(a,b)\mid a\in A,\ b\in B\}$ amb $\pi_1(a,b)=a$ i $\pi_2(a,b)=b$. Donats $f\colon X\to A$ i $g\colon X\to B$, definim $\langle f,g\rangle\colon X\to A\times B$ per $\langle f,g\rangle(x)=(f(x),g(x))$. Aleshores $\pi_1\comp\langle f,g\rangle=f$ i $\pi_2\comp\langle f,g\rangle=g$. I és l'únic possible: si $h\colon X\to A\times B$ compleix $\pi_1\comp h=f$ i $\pi_2\comp h=g$, aleshores $h(x)=(\pi_1(h(x)),\pi_2(h(x)))=(f(x),g(x))=\langle f,g\rangle(x)$. Per tant el producte cartesià satisfà la propietat universal del producte.
\end{solucio}

\begin{exercici}
Comprova que tot equalitzador és un monomorfisme.
\end{exercici}

\begin{solucio}
Sigui $e\colon E\to A$ l'equalitzador de $f,g\colon A\to B$, i siguin $r,s\colon Z\to E$ amb $e\comp r=e\comp s$. Volem $r=s$. Sigui $h=e\comp r=e\comp s\colon Z\to A$. Aleshores $f\comp h=f\comp e\comp r=g\comp e\comp r=g\comp h$, de manera que $h$ satisfà la condició que defineix els cons de l'equalitzador. Per la propietat universal, hi ha un \emph{únic} $k\colon Z\to E$ amb $e\comp k=h$. Però tant $r$ com $s$ compleixen aquesta equació; per unicitat, $r=s$. Així $e$ és un monomorfisme.
\end{solucio}

\section{Pullbacks}

\begin{exercici}
Comprova que a $\cat{Set}$ el pullback de $f\colon A\to C$ i $g\colon B\to C$ és $P=\{(a,b)\in A\times B\mid f(a)=g(b)\}$, i que quan $g$ és la injecció d'un subconjunt $B\subseteq C$ el pullback recupera l'antiimatge $f^{-1}(B)$.
\end{exercici}

\begin{solucio}
Amb $P=\{(a,b)\mid f(a)=g(b)\}$ i les projeccions $p_1,p_2$, es compleix $f\comp p_1=g\comp p_2$ per construcció. Donats $r_1\colon X\to A$ i $r_2\colon X\to B$ amb $f\comp r_1=g\comp r_2$, l'aplicació $h(x)=(r_1(x),r_2(x))$ té imatge dins $P$ (perquè $f(r_1(x))=g(r_2(x))$) i és l'únic morfisme amb $p_1\comp h=r_1$, $p_2\comp h=r_2$. Això és la propietat universal.

Si $B\subseteq C$ i $g$ és la injecció, la condició $f(a)=g(b)=b$ força $b=f(a)$ amb $f(a)\in B$; per tant $P\cong\{a\in A\mid f(a)\in B\}=f^{-1}(B)$. El pullback recupera l'antiimatge.
\end{solucio}

\begin{exercici}[Lema d'enganxament de pullbacks]
Considera el diagrama de dos quadrats adossats en una categoria arbitrària:
\[
\begin{tikzpicture}[baseline=(m.center)]
  \matrix (m) [matrix of math nodes, column sep=3.2em, row sep=2.6em] {
    A & B & C \\
    A' & B' & C' \\
  };
  \draw[-{Stealth[scale=1.1]}] (m-1-1) -- node[above] {$u$} (m-1-2);
  \draw[-{Stealth[scale=1.1]}] (m-1-2) -- node[above] {$v$} (m-1-3);
  \draw[-{Stealth[scale=1.1]}] (m-2-1) -- node[below] {$u'$} (m-2-2);
  \draw[-{Stealth[scale=1.1]}] (m-2-2) -- node[below] {$v'$} (m-2-3);
  \draw[-{Stealth[scale=1.1]}] (m-1-1) -- node[left] {$a$} (m-2-1);
  \draw[-{Stealth[scale=1.1]}] (m-1-2) -- node[left] {$b$} (m-2-2);
  \draw[-{Stealth[scale=1.1]}] (m-1-3) -- node[right] {$c$} (m-2-3);
\end{tikzpicture}
\]
amb tots dos quadrats commutatius. Suposem que el \textbf{quadrat dret} (sobre $B,C,B',C'$) és un pullback. Comprova que el \textbf{quadrat esquerre} és un pullback si i només si el \textbf{rectangle exterior} (sobre $A,C,A',C'$) ho és.
\end{exercici}

\begin{solucio}
Treballem en una categoria arbitrària, usant només la propietat universal.

\textbf{($\Rightarrow$) Si el quadrat esquerre és pullback, l'exterior ho és.} Sigui $T$ un objecte amb morfismes $p\colon T\to A'$ i $q\colon T\to C$ compatibles per al rectangle exterior, és a dir amb $c\comp q=v'\comp u'\comp p$. Com que el quadrat dret és un pullback i els morfismes $q\colon T\to C$ i $u'\comp p\colon T\to B'$ satisfan $v'\comp(u'\comp p)=c\comp q$, hi ha un únic $r\colon T\to B$ amb $v\comp r=q$ i $b\comp r=u'\comp p$. Ara $r$ i $p$ són compatibles per al quadrat esquerre, perquè $b\comp r=u'\comp p$; com que aquest és un pullback, hi ha un únic $t\colon T\to A$ amb $u\comp t=r$ i $a\comp t=p$. Aleshores $(v\comp u)\comp t=v\comp r=q$ i $a\comp t=p$, de manera que $t$ factoritza el con exterior, i és únic per les unicitats successives.

\textbf{($\Leftarrow$) Si el rectangle exterior és pullback, l'esquerre ho és.} Sigui $T$ amb $p\colon T\to A'$ i $r\colon T\to B$ compatibles per al quadrat esquerre, és a dir amb $b\comp r=u'\comp p$. Posem $q=v\comp r\colon T\to C$. Aleshores $p$ i $q$ són compatibles per a l'exterior, perquè
\[
c\comp q=c\comp v\comp r=v'\comp b\comp r=v'\comp u'\comp p.
\]
Com que l'exterior és un pullback, hi ha un únic $t\colon T\to A$ amb $(v\comp u)\comp t=q$ i $a\comp t=p$. Queda veure que $u\comp t=r$: component amb $v$ i amb $b$,
\[
v\comp(u\comp t)=q=v\comp r,\qquad b\comp(u\comp t)=u'\comp a\comp t=u'\comp p=b\comp r,
\]
i com que el quadrat dret és un pullback, aquestes dues igualtats forcen $u\comp t=r$. Així $t$ és la factorització única per al quadrat esquerre.

En tots dos sentits, l'ingredient és la propietat universal del quadrat dret, que transfereix les factoritzacions entre el rectangle i el quadrat esquerre. Aquest \emph{lema d'enganxament} és una de les eines més usades per raonar amb pullbacks, i el retrobarem en tractar la reindexació entre subobjectes.
\end{solucio}

\begin{exercici}\label{pr-reindexacio-set}
Sigui $f\colon A\to B$ una aplicació i $y\colon Y\to B$ un objecte de $\cat{Set}/B$. Calcula explícitament la reindexació $f^{*}Y$ (observació~\ref{obs:reindexacio}) i comprova que la fibra de $f^{*}Y\to A$ sobre cada $a\in A$ està en bijecció canònica amb la fibra $y^{-1}(f(a))$.
\end{exercici}

\begin{solucio}
Pel càlcul del pullback a $\cat{Set}$ de l'exercici anterior,
\[
f^{*}Y=A\times_B Y=\{(a,z)\in A\times Y\mid f(a)=y(z)\},
\]
amb la projecció $p_1\colon(a,z)\mapsto a$ com a objecte de $\cat{Set}/A$. La fibra de $p_1$ sobre $a$ és
\[
p_1^{-1}(a)=\{(a,z)\mid y(z)=f(a)\},
\]
i l'aplicació $(a,z)\mapsto z$ en dona una bijecció canònica amb
\[
\{z\in Y\mid y(z)=f(a)\}=y^{-1}(f(a)).
\]
Així, si interpretem $y\colon Y\to B$ com la família de fibres $(Y_b)_{b\in B}$, amb $Y_b=y^{-1}(b)$, la reindexació al llarg de $f$ s'interpreta com la família $(Y_{f(a)})_{a\in A}$: cada índex $a$ rep la fibra de l'índex $f(a)$. És la \emph{substitució} al llarg de $f$, que reapareixerà com a operació sobre predicats en el capítol de subobjectes.
\end{solucio}

\begin{exercici}\label{pr-provt-limits-finits}
A la categoria prima $\cat{Prov}_T$, suposem que el sistema deductiu $T$ té una fórmula $\top$ i una conjunció $\land$ amb les regles habituals. Comprova que $\top$ és un objecte terminal i que $p\land q$ és un producte de $p$ i $q$. Dedueix que $\cat{Prov}_T$ té tots els límits finits.
\end{exercici}

\begin{solucio}
\textbf{Terminal.} Per a tota fórmula $r$ es té $r\vdash_T\top$, i com que $\cat{Prov}_T$ és prima, aquesta fletxa és única. Per tant $\top$ és terminal.

\textbf{Producte.} Les regles d'eliminació de la conjunció donen les projeccions
\[
p\land q\vdash_T p,
\qquad
p\land q\vdash_T q.
\]
Si $r\vdash_T p$ i $r\vdash_T q$, la regla d'introducció dona $r\vdash_T p\land q$, que és la fletxa mediadora; els triangles commuten automàticament i la fletxa és única, perquè la categoria és prima. Per tant $p\land q$, amb les dues eliminacions, és un producte de $p$ i $q$.

\textbf{Límits finits.} En una categoria prima, dues fletxes paral·leles coincideixen, de manera que l'equalitzador de $x\rightrightarrows y$ és $\id_x$ i sempre existeix. Pel teorema de caracterització dels límits finits (teorema~\ref{teo:limits-finits}), objecte terminal, productes binaris i equalitzadors donen tots els límits finits. És la versió resolta de la lectura lògica de l'exemple~\ref{ex:limits-preordre}.
\end{solucio}

\section{Colímits}

\begin{exercici}
Descriu el coproducte a $\cat{Set}$ i comprova'n la propietat universal.
\end{exercici}

\begin{solucio}
El coproducte de $A$ i $B$ a $\cat{Set}$ és la unió disjunta $A+B=(\{1\}\times A)\cup(\{2\}\times B)$, amb les injeccions $\iota_1(a)=(1,a)$ i $\iota_2(b)=(2,b)$. Donats $f\colon A\to X$ i $g\colon B\to X$, l'aplicació $[f,g]\colon A+B\to X$ definida per $[f,g](1,a)=f(a)$ i $[f,g](2,b)=g(b)$ és l'únic morfisme amb $[f,g]\comp\iota_1=f$ i $[f,g]\comp\iota_2=g$. Aquesta és la definició per casos, i és exactament la propietat universal del coproducte, dual de la del producte.
\end{solucio}

\begin{exercici}
Comprova que a $\cat{Set}$ el coequalitzador de $f,g\colon A\to B$ és el quocient $B/{\sim}$ per la relació d'equivalència generada per $f(a)\sim g(a)$.
\end{exercici}

\begin{solucio}
Sigui $\sim$ la menor relació d'equivalència amb $f(a)\sim g(a)$ per a tot $a$, i $q\colon B\to B/{\sim}$ la projecció al quocient. Aleshores $q\comp f=q\comp g$, perquè $f(a)$ i $g(a)$ tenen la mateixa classe. Si $s\colon B\to S$ compleix $s\comp f=s\comp g$, aleshores $s$ és constant sobre cada classe d'equivalència, de manera que factoritza de manera única com $s=r\comp q$ amb $r([b])=s(b)$ ben definida. Aquesta és la propietat universal del coequalitzador; el quocient hi encaixa exactament.
\end{solucio}

\section{Preservació i reflexió}

\begin{exercici}
Comprova que el functor oblit $U\colon\cat{Grp}\to\cat{Set}$ preserva els productes, però no preserva tots els coproductes.
\end{exercici}

\begin{solucio}
Per als productes: el grup producte $G\times H$ té com a conjunt subjacent el producte cartesià dels conjunts subjacents, i les projeccions són les aplicacions projecció; per tant $U(G\times H)=U(G)\times U(H)$ amb les projeccions correctes, i $U$ preserva el producte.

Per als coproductes: el coproducte a $\cat{Grp}$ és el \emph{producte lliure} $G*H$. N'hi ha prou d'un contraexemple, i el més senzill és prendre dos grups trivials $1=\{e\}$. El seu producte lliure és de nou trivial, $1*1\cong 1$, de manera que $U(1*1)$ té un sol element. En canvi, el coproducte a $\cat{Set}$ dels conjunts subjacents, $U(1)+U(1)$, és la unió disjunta de dos conjunts d'un element i en té dos. El cocon imatge, amb vèrtex $U(1*1)$, no pot ser, doncs, un coproducte a $\cat{Set}$, i $U$ no preserva aquest coproducte.

Aquesta asimetria anticipa el teorema del capítol següent: $U$ és un adjunt per la dreta, i els adjunts per la dreta preserven tots els límits; no hi ha cap afirmació anàloga que els obligui a preservar tots els colímits.
\end{solucio}

\begin{exercici}
Comprova que el functor representable $\Hom(A,-)\colon\cat{Set}\to\cat{Set}$ preserva productes, és a dir, $\Hom(A,X\times Y)\cong\Hom(A,X)\times\Hom(A,Y)$.
\end{exercici}

\begin{solucio}
Per la propietat universal del producte, l'aplicació $h\mapsto(\pi_1\comp h,\ \pi_2\comp h)$ és una bijecció
\[
\Hom(A,X\times Y)\cong\Hom(A,X)\times\Hom(A,Y),
\]
natural en $X$ i $Y$ (i també en $A$, si es considera la variació contravariant en aquesta variable). Aquesta bijecció és exactament la induïda pel con imatge $\bigl(\Hom(A,X\times Y),\Hom(A,\pi_1),\Hom(A,\pi_2)\bigr)$, de manera que aquest con és un producte a $\cat{Set}$ i $\Hom(A,-)$ preserva el producte binari. El mateix argument, amb la propietat universal general del límit, mostra que $\Hom(A,-)$ preserva tots els límits (proposició~\ref{prop:representables-limits}).
\end{solucio}

\chapter[Adjuncions]{Adjuncions}

Els exercicis d'aquest capítol treballen les adjuncions des dels tres punts de vista equivalents ---bijecció d'homconjunts, unitat/counitat, identitats triangulars--- i n'apliquen les conseqüències sobre límits. Convé identificar sempre qui és l'adjunt per l'esquerra i qui per la dreta.

\section{Adjuncions entre preordres}

\begin{exercici}
Comprova que al preordre $\cat{Prov}_T$ d'un sistema proposicional amb les regles usuals de $\wedge$ i $\to$, fixada una fórmula $A$, l'operació $B\mapsto A\wedge B$ és adjunta per l'esquerra de $B\mapsto A\to B$.
\end{exercici}

\begin{solucio}
Hem de veure que $A\wedge B\vdash_T C$ si i només si $B\vdash_T A\to C$. 

Suposem $A\wedge B\vdash_T C$. Aleshores, assumint $B$, volem deduir $A\to C$: assumim també $A$; tenim $A$ i $B$, per tant $A\wedge B$, d'on $C$; descarregant la hipòtesi $A$ obtenim $A\to C$. Així $B\vdash_T A\to C$.

Recíprocament, suposem $B\vdash_T A\to C$. Assumint $A\wedge B$, en tenim $A$ i $B$; de $B$ obtenim $A\to C$, i amb $A$, per modus ponens, $C$. Així $A\wedge B\vdash_T C$.

Les dues implicacions donen l'equivalència, és a dir, $(A\wedge-)\dashv(A\to-)$. Aquesta adjunció és la lectura categòrica de la regla de la implicació.
\end{solucio}

\begin{exercici}\label{pr-quantificadors-calcul}
Sigui $f\colon X\to Y$ l'aplicació de $X=\{1,2,3,4\}$ en $Y=\{a,b,c\}$ amb $f(1)=f(2)=a$, $f(3)=b$ i $f(4)=b$, i sigui $S=\{1,2,3\}$. Calcula $\exists_f(S)$ i $\forall_f(S)$ (exemple~\ref{ex:quantificadors-set}). Després, fent servir només la cadena $\exists_f\dashv f^{*}\dashv\forall_f$, justifica que $\exists_f$ preserva unions, $\forall_f$ preserva interseccions i $f^{*}$ preserva totes dues.
\end{exercici}

\begin{solucio}
\textbf{Càlcul.} La imatge és $\exists_f(S)=\{a,b\}$. Per a $\forall_f(S)$, mirem les fibres: $f^{-1}(\{a\})=\{1,2\}\subseteq S$, de manera que $a\in\forall_f(S)$; $f^{-1}(\{b\})=\{3,4\}\not\subseteq S$, perquè $4\notin S$; i $f^{-1}(\{c\})=\varnothing\subseteq S$. Per tant
\[
\forall_f(S)=\{a,c\}.
\]
Observem dos fenòmens: $c$ hi és perquè la seva fibra és buida (la quantificació universal sobre un conjunt buit és certa), i $b$ no hi és tot i que $3\in S$.

\textbf{Preservació.} En un preordre de parts, les unions són els coproductes (suprems) i les interseccions, els productes (ínfims). Pel teorema que els adjunts per l'esquerra preserven colímits i els adjunts per la dreta preserven límits:
\begin{itemize}
  \item $\exists_f$, adjunt per l'esquerra de $f^{*}$, preserva unions: $f[S\cup S']=f[S]\cup f[S']$;
  \item $\forall_f$, adjunt per la dreta de $f^{*}$, preserva interseccions;
  \item $f^{*}$, que és adjunt per la dreta de $\exists_f$ i adjunt per l'esquerra de $\forall_f$, preserva totes dues.
\end{itemize}
En canvi, $\exists_f$ no preserva en general les interseccions: amb l'aplicació de l'enunciat, $\exists_f(\{1\}\cap\{2\})=\varnothing$, però $\exists_f(\{1\})\cap\exists_f(\{2\})=\{a\}$. Lògicament: l'existencial commuta amb la disjunció, però no amb la conjunció.
\end{solucio}

\section{Adjuncions entre categories}

\begin{exercici}\label{pr-disc-ob}
Sigui $\mathrm{Ob}\colon\cat{Cat}\to\cat{Set}$ el functor que envia cada categoria petita al seu conjunt d'objectes, i $\mathrm{Disc}\colon\cat{Set}\to\cat{Cat}$ el que envia un conjunt $X$ a la categoria discreta amb objectes $X$. Comprova que $\mathrm{Disc}\dashv\mathrm{Ob}$ i descriu-ne la unitat i la counitat.
\end{exercici}

\begin{solucio}
Un functor $H\colon\mathrm{Disc}(X)\to\cat{D}$ ha d'enviar cada identitat a una identitat, i $\mathrm{Disc}(X)$ no té cap altre morfisme; queda, doncs, determinat per una aplicació $X\to\mathrm{Ob}(\cat{D})$ entre els objectes, i qualsevol aplicació dona un functor. Això és una bijecció
\[
\Hom_{\cat{Cat}}\bigl(\mathrm{Disc}(X),\cat{D}\bigr)\;\cong\;\Hom_{\cat{Set}}\bigl(X,\mathrm{Ob}(\cat{D})\bigr),
\]
natural en $X$ i en $\cat{D}$, perquè les dues bandes es componen de la mateixa manera sobre els objectes. Per tant $\mathrm{Disc}\dashv\mathrm{Ob}$.

La \emph{unitat} $\eta_X\colon X\to\mathrm{Ob}(\mathrm{Disc}(X))=X$ és la identitat: és un isomorfisme. La \emph{counitat} $\varepsilon_{\cat{D}}\colon\mathrm{Disc}(\mathrm{Ob}(\cat{D}))\to\cat{D}$ és el functor que inclou els objectes de $\cat{D}$ amb les seves identitats, i en general no és un isomorfisme, perquè oblida tots els morfismes no identitat.

\emph{Observació.} $\mathrm{Ob}$ té també un adjunt per la dreta: el functor que envia $X$ a la categoria \emph{indiscreta}, amb exactament un morfisme entre cada parella d'objectes. Per tant $\mathrm{Ob}$ preserva límits i colímits, d'acord amb el fet que els límits i colímits de categories petites es calculen, sobre els objectes, com a $\cat{Set}$.
\end{solucio}

\section{Unitat i counitat}

\begin{exercici}
Per a l'adjunció lliure-oblit $F\dashv U$ dels monoides, descriu la unitat $\eta_X\colon X\to U(X^*)$ i la counitat $\varepsilon_M\colon (U(M))^*\to M$, i comprova que la unitat expressa la propietat universal del monoide lliure.
\end{exercici}

\begin{solucio}
La unitat $\eta_X\colon X\to U(X^*)$ envia cada element $x$ a la paraula d'una sola lletra $x$. La counitat $\varepsilon_M\colon (U(M))^*\to M$ envia una paraula d'elements de $M$ al seu producte a $M$: $\varepsilon_M(m_1\cdots m_n)=m_1\cdots m_n$ (producte calculat a $M$).

La propietat universal de la unitat diu que per a tota aplicació $\varphi\colon X\to U(M)$ hi ha un únic homomorfisme $\bar\varphi\colon X^*\to M$ amb $U(\bar\varphi)\comp\eta_X=\varphi$. Això és exactament la definició del monoide lliure: un homomorfisme des de $X^*$ queda determinat, de manera única, pels valors sobre els generadors $X$.
\end{solucio}

\begin{exercici}
Comprova que, en una adjunció $F\dashv G$, el transposat de $\varphi\colon F(A)\to B$ és $G(\varphi)\comp\eta_A$, i el transposat invers de $\psi\colon A\to G(B)$ és $\varepsilon_B\comp F(\psi)$.
\end{exercici}

\begin{solucio}
Per naturalitat de la bijecció $\theta$ en la segona variable, aplicada al morfisme $\varphi\colon F(A)\to B$ i partint de $\id_{F(A)}$,
\[
\theta(\varphi)=\theta(\varphi\comp\id_{F(A)})=G(\varphi)\comp\theta(\id_{F(A)})=G(\varphi)\comp\eta_A,
\]
on hem usat que $\theta(\id_{F(A)})=\eta_A$ per definició de la unitat. Dualment, per naturalitat en la primera variable i la definició de la counitat $\varepsilon_B=\theta^{-1}(\id_{G(B)})$,
\[
\theta^{-1}(\psi)=\theta^{-1}(\id_{G(B)}\comp\psi)=\theta^{-1}(\id_{G(B)})\comp F(\psi)=\varepsilon_B\comp F(\psi).
\]
Aquestes fórmules permeten passar de la bijecció a la unitat/counitat i viceversa.
\end{solucio}

\section{Identitats triangulars}

\begin{exercici}\label{pr-diagonal-producte}
Sigui $\cat{C}$ una categoria amb productes binaris, i siguin $\Delta\colon\cat{C}\to\cat{C}\times\cat{C}$ el functor diagonal, $\Delta(A)=(A,A)$, i $\times\colon\cat{C}\times\cat{C}\to\cat{C}$ el functor producte, $(X,Y)\mapsto X\times Y$. Comprova que $\Delta\dashv\times$, descriu-ne la unitat i la counitat, i verifica directament les dues identitats triangulars.
\end{exercici}

\begin{solucio}
\textbf{L'adjunció.} Un morfisme $\Delta(A)\to(X,Y)$ a $\cat{C}\times\cat{C}$ és una parella $(f\colon A\to X,\ g\colon A\to Y)$, i la propietat universal del producte la posa en bijecció amb els morfismes $\langle f,g\rangle\colon A\to X\times Y$:
\[
\Hom_{\cat{C}\times\cat{C}}\bigl((A,A),(X,Y)\bigr)\;\cong\;\Hom_{\cat{C}}(A,X\times Y).
\]
La bijecció és natural en $A$ i en $(X,Y)$, perquè $\langle f,g\rangle\comp a=\langle f\comp a,g\comp a\rangle$ i $(u\times v)\comp\langle f,g\rangle=\langle u\comp f,v\comp g\rangle$. Per tant $\Delta\dashv\times$.

\textbf{Unitat i counitat.} La unitat és el transposat de $\id_{(A,A)}$, és a dir, el morfisme diagonal
\[
\eta_A=\delta_A=\langle\id_A,\id_A\rangle\colon A\to A\times A,
\]
i la counitat és el transposat invers de $\id_{X\times Y}$, la parella de projeccions
\[
\varepsilon_{(X,Y)}=(\pi_1,\pi_2)\colon(X\times Y,\ X\times Y)\to(X,Y).
\]

\textbf{Primera identitat.} $\varepsilon_{\Delta(A)}\comp\Delta(\eta_A)$ és la parella
\[
(\pi_1\comp\delta_A,\ \pi_2\comp\delta_A)=(\id_A,\id_A)=\id_{\Delta(A)}:
\]
copiar $A$ i projectar a cada còpia no canvia res.

\textbf{Segona identitat.} $(\pi_1\times\pi_2)\comp\delta_{X\times Y}$ és el morfisme $X\times Y\to X\times Y$ de components $\pi_1$ i $\pi_2$, és a dir,
\[
\langle\pi_1,\pi_2\rangle=\id_{X\times Y}:
\]
duplicar una parella i quedar-se amb la primera component de la primera còpia i la segona de la segona retorna la parella original.

Observem que la unitat $\delta_A$ no és un isomorfisme (a $\cat{Set}$, només ho és si $A$ té com a màxim un element): les identitats triangulars no demanen que ho sigui.
\end{solucio}

\begin{exercici}\label{pr-free-u-triangulars}
Per a l'adjunció $\Free\dashv U$ entre $\cat{Graph}$ i $\cat{Cat}$ (exemple~\ref{ex:free-u-adjuncio}), comprova directament les dues identitats triangulars
\[
\varepsilon_{\Free(G)}\comp\Free(\eta_G)=\id_{\Free(G)},
\qquad
U(\varepsilon_{\cat{C}})\comp\eta_{U(\cat{C})}=\id_{U(\cat{C})},
\]
i explica què diu cadascuna en termes de camins.
\end{exercici}

\begin{solucio}
\textbf{Primera identitat.} El functor $\Free(\eta_G)\colon\Free(G)\to\Free(U(\Free(G)))$ envia un camí $(a_1,\dots,a_n)$ de $G$ al camí formal $\bigl((a_1),\dots,(a_n)\bigr)$, les arestes del qual són camins de longitud~$1$ de $\Free(G)$. La counitat $\varepsilon_{\Free(G)}$ avalua aquest camí formal component les seves arestes a $\Free(G)$, és a dir, concatenant:
\[
\varepsilon_{\Free(G)}\Bigl(\bigl((a_1),\dots,(a_n)\bigr)\Bigr)=(a_n)\comp\cdots\comp(a_1)=(a_1,\dots,a_n).
\]
El camí buit va al camí buit. Per tant, la composta és la identitat de $\Free(G)$: \emph{trencar un camí en arestes i tornar-lo a compondre no canvia res}.

\textbf{Segona identitat.} La unitat $\eta_{U(\cat{C})}$ envia una fletxa $f$ de $\cat{C}$, vista com a aresta del graf $U(\cat{C})$, al camí $(f)$ de longitud~$1$, i cada objecte a ell mateix. El morfisme de grafs $U(\varepsilon_{\cat{C}})$ avalua aquest camí, que té una sola fletxa, i en dona $f$. Per tant, la composta és la identitat de $U(\cat{C})$: \emph{avaluar el camí d'una sola fletxa retorna la fletxa}.

Les dues identitats expressen, doncs, que la inserció dels generadors i l'avaluació de camins són inverses en el sentit precís que demanen les identitats triangulars, sense que $\eta$ ni $\varepsilon$ siguin isomorfismes: $\Free(U(\cat{C}))$ té en general molts més morfismes que $\cat{C}$, perquè hi ha camins diferents amb la mateixa composta.
\end{solucio}

\section{Els adjunts i els límits}

\begin{exercici}\label{pr-rapl-mon}
Comprova, fent servir que els adjunts per la dreta preserven límits, que el functor oblit $U\colon\cat{Mon}\to\cat{Set}$ preserva productes. Comprova també que $U$ no preserva tots els coproductes, i explica què en dedueixes.
\end{exercici}

\begin{solucio}
El functor monoide lliure $F\colon\cat{Set}\to\cat{Mon}$, $X\mapsto X^*$, és adjunt per l'esquerra de $U$ (exemple del monoide lliure de la Teoria). Per tant $U$ és un adjunt per la dreta, i pel teorema \enquote{els adjunts per la dreta preserven límits}, $U$ preserva tots els límits petits, en particular els productes: el conjunt subjacent d'un producte de monoides és el producte dels conjunts subjacents.

El teorema, però, no diu res dels colímits, i de fet $U$ no preserva tots els coproductes. Sigui $1=\{e\}$ el monoide trivial. És un objecte inicial de $\cat{Mon}$, perquè l'únic homomorfisme $1\to M$ envia $e$ al neutre, i el coproducte d'un objecte amb un objecte inicial és l'objecte mateix; per tant, el coproducte $1+1$ a $\cat{Mon}$ és $1$, amb un sol element. En canvi, el coproducte a $\cat{Set}$ dels conjunts subjacents, $U(1)+U(1)$, en té dos, i el cocon imatge no és un coproducte. Com a conseqüència, $U$ no pot ser un adjunt per l'esquerra, perquè els adjunts per l'esquerra preserven colímits.

\emph{Ampliació, per a qui conegui els grups.} El mateix passa amb $U\colon\cat{Grp}\to\cat{Set}$. El seu adjunt per l'esquerra és el functor \emph{grup lliure}, que envia un conjunt $X$ al grup de les paraules reduïdes en les lletres de $X$ i els seus inversos formals; per tant, $U$ preserva tots els límits. El grup trivial també és inicial a $\cat{Grp}$, el coproducte (el \emph{producte lliure}) de dos grups trivials és trivial, $1*1\cong 1$, i $U$ tampoc no preserva aquest coproducte.
\end{solucio}

\begin{exercici}
Comprova que si un functor $G$ \emph{no} preserva algun límit petit, aleshores $G$ no pot tenir cap adjunt per l'esquerra.
\end{exercici}

\begin{solucio}
És el contrarecíproc del teorema. Si $G$ tingués un adjunt per l'esquerra $F$, amb $F\dashv G$, aleshores $G$ seria un adjunt per la dreta i, pel teorema, preservaria tots els límits petits que existeixen. Com que per hipòtesi $G$ no en preserva algun, no pot tenir adjunt per l'esquerra. Aquest és un criteri negatiu molt pràctic: per exemple, si un functor oblit no preserva algun coproducte, això impedeix que sigui adjunt per l'esquerra, però no impedeix que sigui adjunt per la dreta.
\end{solucio}

\section{Exponencials}

\begin{exercici}
Comprova la bijecció
\[
\Hom_{\cat{Set}}(A\times B,C)\cong\Hom_{\cat{Set}}(A,C^B)
\]
i identifica-la amb la currificació.
\end{exercici}

\begin{solucio}
A $\cat{Set}$, $C^B$ és el conjunt de les aplicacions $B\to C$. Donada $f\colon A\times B\to C$, definim $\tilde f\colon A\to C^B$ per $\tilde f(a)=(b\mapsto f(a,b))$; és a dir, $\tilde f(a)$ és la funció que fixa la primera coordenada en $a$. Recíprocament, donada $g\colon A\to C^B$, definim $\hat g\colon A\times B\to C$ per $\hat g(a,b)=g(a)(b)$. Aquestes assignacions són inverses l'una de l'altra i naturals en totes les variables. La correspondència $f\mapsto\tilde f$ és exactament la \textbf{currificació}: transformar una funció de dues variables en una funció d'una variable que retorna funcions. Aquesta bijecció expressa l'adjunció $-\times B\dashv(-)^B$.
\end{solucio}

\chapter[Cartesianes tancades]{\texorpdfstring{Categories cartesianes\\ tancades}{Categories cartesianes tancades}}

Els exercicis treballen tres aspectes del tancament cartesià: productes i diagonals, la propietat universal de l'exponencial i la interpretació de tipus i lògica. Convé tenir sempre present la bijecció
\[
\Hom(A\times B,C)\cong\Hom(A,C^{B}).
\]

\section{Productes i diagonals}

\begin{exercici}
En una categoria cartesiana, comprova que la diagonal $\delta_A\colon A\to A\times A$ satisfà
\[
\pi_1\comp\delta_A=\pi_2\comp\delta_A=\id_A
\]
i que és l'únic morfisme amb aquesta propietat.
\end{exercici}

\begin{solucio}
Per la propietat universal del producte $A\times A$, un morfisme $A\to A\times A$ està determinat per les seves dues composicions amb les projeccions. Definim $\delta_A$ com l'únic morfisme amb $\pi_1\comp\delta_A=\id_A$ i $\pi_2\comp\delta_A=\id_A$; l'existència i la unicitat són exactament el contingut de la propietat universal aplicada a la parella $(\id_A,\id_A)$. Qualsevol altre morfisme $d\colon A\to A\times A$ amb $\pi_1\comp d=\pi_2\comp d=\id_A$ compleix la mateixa condició universal i, per unicitat, $d=\delta_A$.
\end{solucio}

\section{La propietat universal de l'exponencial}

\begin{exercici}
A $\cat{Set}$, descriu explícitament l'avaluació, la currificació i la descurrificació per a l'exponencial $C^{B}$, i verifica l'equació $\ev\comp(\lambda f\times\id_B)=f$.
\end{exercici}

\begin{solucio}
A $\cat{Set}$, $C^{B}$ és el conjunt d'aplicacions $B\to C$. L'avaluació és $\ev\colon C^{B}\times B\to C$, $\ev(g,b)=g(b)$. Donada $f\colon A\times B\to C$, la currificació és $\lambda f\colon A\to C^{B}$ amb $\lambda f(a)=(b\mapsto f(a,b))$. Comprovem l'equació: per a $(a,b)\in A\times B$,
\[
\bigl(\ev\comp(\lambda f\times\id_B)\bigr)(a,b)=\ev\bigl(\lambda f(a),b\bigr)=\lambda f(a)(b)=f(a,b).
\]
Per tant $\ev\comp(\lambda f\times\id_B)=f$. La descurrificació de $g\colon A\to C^{B}$ és $\widetilde g(a,b)=g(a)(b)$, i és inversa de la currificació.
\end{solucio}

\begin{exercici}
Demostra que en qualsevol CCC hi ha un isomorfisme natural $C^{1}\cong C$ i que $1^{B}\cong 1$.
\end{exercici}

\begin{solucio}
Per a $C^{1}$: per la propietat universal, $\Hom(A,C^{1})\cong\Hom(A\times 1,C)$. Com que $A\times 1\cong A$ (l'objecte terminal és neutre per al producte), $\Hom(A\times 1,C)\cong\Hom(A,C)$, natural en $A$. Pel lema de Yoneda, $C^{1}\cong C$.

Per a $1^{B}$: $\Hom(A,1^{B})\cong\Hom(A\times B,1)$, que té exactament un element per ser $1$ terminal. Per tant $\Hom(A,1^{B})$ és un conjunt d'un punt per a tot $A$, natural en $A$, i pel lema de Yoneda $1^{B}\cong 1$. Aquestes són les versions categòriques de les identitats $c^{1}=c$ i $1^{b}=1$.
\end{solucio}

\begin{exercici}
Comprova la \enquote{llei exponencial} $C^{A\times B}\cong(C^{B})^{A}$ en una CCC.
\end{exercici}

\begin{solucio}
Fem servir la propietat universal repetidament i el lema de Yoneda. Per a qualsevol $X$,
\[
\Hom(X,C^{A\times B})\cong\Hom(X\times(A\times B),C)\cong\Hom((X\times A)\times B,C),
\]
usant l'associativitat del producte. Aplicant ara l'adjunció respecte de $B$ i després respecte de $A$,
\[
\Hom((X\times A)\times B,C)\cong\Hom(X\times A,C^{B})\cong\Hom(X,(C^{B})^{A}).
\]
Totes les bijeccions són naturals en $X$; pel lema de Yoneda, $C^{A\times B}\cong(C^{B})^{A}$. És la llei $c^{a\cdot b}=(c^{b})^{a}$, i correspon, en el càlcul lambda, al fet que una funció de dos arguments és una funció que retorna una funció (currificació iterada).
\end{solucio}

\begin{exercici}\label{pr-distributivitat}
Sigui $\cat{C}$ una CCC amb coproductes binaris i objecte inicial $0$. Demostra que el producte es distribueix sobre el coproducte:
\[
A\times(B+C)\;\cong\;(A\times B)+(A\times C),
\qquad
A\times 0\;\cong\;0.
\]
Interpreta el resultat a $\cat{Set}$ i en una àlgebra de Heyting.
\end{exercici}

\begin{solucio}
Pel tancament cartesià, $-\times A\dashv(-)^{A}$, i el producte és commutatiu llevat d'isomorfisme natural; per tant $A\times-$ també és un adjunt per l'esquerra. Els adjunts per l'esquerra preserven colímits, en particular els coproductes binaris i l'objecte inicial, que és el colímit del diagrama buit. Explícitament, el morfisme canònic
\[
[\id_A\times\iota_1,\ \id_A\times\iota_2]\colon(A\times B)+(A\times C)\longrightarrow A\times(B+C)
\]
és un isomorfisme, i $A\times 0$ és inicial, de manera que $A\times 0\cong 0$.

\emph{A $\cat{Set}$}, la primera bijecció envia una parella $(a,x)$ amb $x$ a la component $B$ o a la component $C$ de la unió disjunta al lloc corresponent; la segona diu que $A\times\varnothing=\varnothing$.

\emph{En una àlgebra de Heyting}, vista com a preordre cartesià tancat amb suprems finits, el coproducte és el suprem i l'inicial és $\bot$. Obtenim
\[
a\wedge(b\vee c)=(a\wedge b)\vee(a\wedge c),
\qquad
a\wedge\bot=\bot:
\]
tota àlgebra de Heyting és un reticle distributiu, i la distributivitat és conseqüència de l'existència de la implicació. Lògicament, la conjunció es distribueix sobre la disjunció.
\end{solucio}

\section{Heyting i lògica}

\begin{exercici}
Comprova que en una àlgebra de Heyting l'exponencial $b\Rightarrow c$ satisfà $a\wedge b\leq c\Leftrightarrow a\leq(b\Rightarrow c)$, i deriva-hi que $b\wedge(b\Rightarrow c)\leq c$ (\emph{modus ponens}).
\end{exercici}

\begin{solucio}
La condició $a\wedge b\leq c\Leftrightarrow a\leq(b\Rightarrow c)$ és, precisament, l'adjunció $-\wedge b\dashv(b\Rightarrow-)$ escrita en el llenguatge dels preordres: és la definició d'exponencial en un preordre cartesià, on $\Hom$ té com a molt un element i la bijecció d'homconjunts es redueix a una equivalència de desigualtats. Per obtenir el \emph{modus ponens}, prenem $a=(b\Rightarrow c)$. Aleshores el membre dret $a\leq(b\Rightarrow c)$ es compleix trivialment (és $b\Rightarrow c\leq b\Rightarrow c$), de manera que el membre esquerre també: $(b\Rightarrow c)\wedge b\leq c$. Aquesta és la counitat de l'adjunció, i és la forma algebraica de la regla que de $b$ i $b\to c$ se'n segueix $c$.
\end{solucio}

\begin{exercici}
En una àlgebra de Heyting, demostra que $a\leq(b\Rightarrow a)$ per a tot $a,b$, i interpreta-ho lògicament.
\end{exercici}

\begin{solucio}
Per l'adjunció, $a\leq(b\Rightarrow a)$ equival a $a\wedge b\leq a$, que és cert perquè $a\wedge b$ és un ínfim i, per tant, $a\wedge b\leq a$. Lògicament, correspon a l'axioma $a\to(b\to a)$: una proposició verdadera se segueix de qualsevol hipòtesi. És l'esquelet del combinador $\mathsf{K}$ del càlcul lambda, la funció constant.
\end{solucio}

\section{Semàntica de termes}

\begin{exercici}
Interpreta el terme $x:\sigma,\,y:\tau\vdash\langle y,x\rangle:\tau\times\sigma$ com a morfisme en una CCC i identifica'l amb el morfisme de commutació del producte.
\end{exercici}

\begin{solucio}
El context $x:\sigma,y:\tau$ s'interpreta com l'objecte $\llbracket\sigma\rrbracket\times\llbracket\tau\rrbracket$. Les variables $x$ i $y$ són les projeccions $\pi_1$ i $\pi_2$. La parella $\langle y,x\rangle$ s'interpreta amb l'aparellament de les interpretacions de $y$ i $x$, és a dir,
\[
\llbracket\langle y,x\rangle\rrbracket=\langle\pi_2,\pi_1\rangle\colon\llbracket\sigma\rrbracket\times\llbracket\tau\rrbracket\to\llbracket\tau\rrbracket\times\llbracket\sigma\rrbracket.
\]
Aquest és exactament el morfisme de \textbf{commutació} (o \emph{swap}) del producte, que intercanvia els dos factors. La seva inversa és $\langle\pi_2,\pi_1\rangle$ en sentit contrari, cosa que expressa que $\sigma\times\tau\cong\tau\times\sigma$.
\end{solucio}

\begin{exercici}
Interpreta l'abstracció $x:\sigma\vdash(\lambda y{:}\tau.\,x):\tau\to\sigma$ i comprova que és la currificació d'una projecció.
\end{exercici}

\begin{solucio}
El cos del terme, $x:\sigma,\,y:\tau\vdash x:\sigma$, s'interpreta amb la projecció $\pi_1\colon\llbracket\sigma\rrbracket\times\llbracket\tau\rrbracket\to\llbracket\sigma\rrbracket$. L'abstracció respecte de $y$ es tradueix per la currificació:
\[
\llbracket\lambda y.\,x\rrbracket=\lambda(\pi_1)\colon\llbracket\sigma\rrbracket\to\llbracket\sigma\rrbracket^{\llbracket\tau\rrbracket}.
\]
És a dir, la funció que envia cada $x$ a la funció constant de valor $x$: el combinador $\mathsf{K}$, la versió proposicional del qual és la desigualtat $a\leq(b\Rightarrow a)$ de l'exercici de Heyting anterior.
\end{solucio}

\chapter[Subobjectes, imatges i factoritzacions]{\texorpdfstring{Subobjectes, imatges\\ i factoritzacions}{Subobjectes, imatges i factoritzacions}}

Els exercicis d'aquest capítol fixen les tres idees centrals: els subobjectes com a classes de monomorfismes, l'ordre $\Sub(A)$ com a àlgebra de predicats, i la reindexació $f^{*}$ com a substitució. Molts arguments es fan \enquote{sense elements}, usant només propietats universals; convé acostumar-s'hi, perquè és el que permet traslladar-los a categories arbitràries.

\section{Subobjectes i el seu ordre}

\begin{exercici}
Comprova que a $\cat{Set}$ dos monomorfismes $m\colon S\rightarrowtail A$ i $n\colon T\rightarrowtail A$ són equivalents si i només si tenen la mateixa imatge, i dedueix que $\Sub(A)\cong\mathcal{P}(A)$.
\end{exercici}

\begin{solucio}
Si $m\sim n$, hi ha un isomorfisme $\varphi\colon S\to T$ amb $n\comp\varphi=m$; aleshores $m(S)=n(\varphi(S))=n(T)$, de manera que tenen la mateixa imatge. Recíprocament, suposem $m(S)=n(T)=:I$. Siguin $\bar m\colon S\to I$ i $\bar n\colon T\to I$ les \emph{corestriccions} de $m$ i $n$ a $I$, és a dir, les mateixes aplicacions amb codomini $I$: $\bar m(s)=m(s)$ i $\bar n(t)=n(t)$. Com que $m$ i $n$ són injectives i tenen imatge $I$, $\bar m$ i $\bar n$ són bijeccions. La composició $\varphi=\bar n^{-1}\comp\bar m\colon S\to T$ és, doncs, una bijecció, i compleix $n\comp\varphi=m$, perquè per a tot $s\in S$, $n(\varphi(s))=\bar n(\bar n^{-1}(\bar m(s)))=m(s)$. Així $m\sim n$. Per tant la imatge determina la classe, i l'assignació $[m]\mapsto m(S)$ és una bijecció entre $\Sub(A)$ i el conjunt de subconjunts $\mathcal{P}(A)$; a més respecta l'ordre, ja que $[m]\leq[n]$ equival a $m(S)\subseteq n(T)$.
\end{solucio}

\begin{exercici}
Demostra que en qualsevol categoria l'objecte $\id_A\colon A\to A$ és el subobjecte \textbf{màxim} de $\Sub(A)$.
\end{exercici}

\begin{solucio}
La identitat és un monomorfisme (és un isomorfisme). Per a qualsevol altre subobjecte $m\colon S\rightarrowtail A$, el mateix $m$ és una factorització de $m$ a través de $\id_A$, ja que $\id_A\comp m=m$. Per tant $[m]\leq[\id_A]$ per a tot $m$, cosa que fa de $[\id_A]$ el màxim. Lògicament, és el predicat \enquote{sempre cert}: tot element de $A$ el satisfà.
\end{solucio}

\begin{exercici}
Comprova que si $m\colon S\rightarrowtail A$ és un subobjecte i $\varphi\colon S'\to S$ és un isomorfisme, aleshores $m\comp\varphi$ representa el mateix subobjecte. Per què això justifica treballar amb classes i no amb monos concrets?
\end{exercici}

\begin{solucio}
Per definició de la relació d'equivalència, $m\comp\varphi\sim m$ mitjançant l'isomorfisme $\varphi$: $m\comp\varphi=m\comp\varphi$ amb $m\comp\varphi$ factoritzant. Formalment, $[m\comp\varphi]=[m]$ perquè $\varphi$ és un isomorfisme del domini que fa commutar el triangle. Això justifica treballar amb classes: la teoria dels subobjectes ha de ser invariant per reparametritzacions del domini, ja que \enquote{el subconjunt} no depèn de com l'indexem. Prendre classes elimina aquesta ambigüitat sense perdre informació geomètrica.
\end{solucio}

\section{Reindexació}

\begin{exercici}
Sigui $f\colon A\to B$ a $\cat{Set}$ i $S\subseteq B$. Comprova que $f^{*}(S)=f^{-1}(S)$ i que $f^{*}$ preserva interseccions: $f^{*}(S\cap S')=f^{*}(S)\cap f^{*}(S')$.
\end{exercici}

\begin{solucio}
El pullback de la injecció $S\hookrightarrow B$ al llarg de $f$ és $\{(a,s)\in A\times S\mid f(a)=s\}$, que s'identifica amb $\{a\in A\mid f(a)\in S\}=f^{-1}(S)$, amb la projecció cap a $A$ com a injecció. Per a les interseccions, $x\in f^{-1}(S\cap S')$ si i només si $f(x)\in S\cap S'$, és a dir, $f(x)\in S$ i $f(x)\in S'$, això és $x\in f^{-1}(S)\cap f^{-1}(S')$. La reindexació preserva, doncs, la conjunció; és una manifestació que la substitució commuta amb $\wedge$.
\end{solucio}


\begin{exercici}
Comprova que si $f\colon A\to B$ és un epimorfisme escindit (té secció $s$ amb $f\comp s=\id_B$), aleshores $f^{*}$ és injectiva sobre subobjectes.
\end{exercici}

\begin{solucio}
Si $f\comp s=\id_B$, aplicant la functorialitat de la reindexació (proposició~\ref{prop:functor-sub}), $s^{*}\comp f^{*}=(f\comp s)^{*}=(\id_B)^{*}=\id_{\Sub(B)}$. Per tant $f^{*}$ té inversa per l'esquerra $s^{*}$, cosa que la fa injectiva: si $f^{*}[m]=f^{*}[n]$, aplicant $s^{*}$ obtenim $[m]=[n]$. Com que $f$ admet una secció, els subobjectes de $B$ es poden recuperar després de reindexar-los al llarg de $f$; per això $f^{*}$ és injectiva. (No vol dir que $f$ no identifiqui elements: una epi escindida pot identificar-ne molts.)
\end{solucio}

\section{Interseccions i lògica}

\begin{exercici}
Verifica amb detall que el pullback de dos subobjectes $m\colon S\rightarrowtail A$ i $n\colon T\rightarrowtail A$ és el seu ínfim a $\Sub(A)$, escrivint les dues direccions de la propietat universal.
\end{exercici}

\begin{solucio}
Sigui $P=S\times_A T$ el pullback, amb projeccions $p\colon P\to S$ i $q\colon P\to T$ i injecció $m\comp p=n\comp q\colon P\rightarrowtail A$ (mono per ser composició de monos). \emph{Fita inferior}: $[P]\leq[m]$ via $p$, ja que $m\comp p$ és la injecció i factoritza per $m$; anàlogament $[P]\leq[n]$ via $q$. \emph{Màxima fita inferior}: sigui $k\colon R\rightarrowtail A$ amb $[k]\leq[m]$ i $[k]\leq[n]$, amb factoritzacions $a\colon R\to S$ ($m\comp a=k$) i $b\colon R\to T$ ($n\comp b=k$). Aleshores $m\comp a=k=n\comp b$, de manera que la parella $(a,b)$ és compatible amb el pullback i indueix un únic $u\colon R\to P$ amb $p\comp u=a$ i $q\comp u=b$. Aquest $u$ factoritza $k$ a través de la injecció de $P$: $(m\comp p)\comp u=m\comp a=k$. Per tant $[k]\leq[P]$, i $[P]$ és l'ínfim $[m]\wedge[n]$.
\end{solucio}

\section{Imatges i factoritzacions}

\begin{exercici}
Comprova que a $\cat{Set}$ tota aplicació $f\colon A\to B$ té factorització imatge amb $I=f(A)$, i identifica l'epimorfisme regular i el monomorfisme.
\end{exercici}

\begin{solucio}
Prenem $I=f(A)\subseteq B$. L'aplicació $e\colon A\to I$, $e(a)=f(a)$, és exhaustiva, i tota aplicació exhaustiva a $\cat{Set}$ és un epimorfisme regular: és el coequalitzador de la seva parella nucli $\{(a,a')\mid f(a)=f(a')\}$. La injecció $m\colon I\hookrightarrow B$ és injectiva, per tant un monomorfisme. Aleshores $f=m\comp e$ és la factorització imatge, i $\Imatge(f)=[m]$ és el subconjunt $f(A)$. Recuperem la noció usual d'imatge d'una aplicació.
\end{solucio}

\begin{exercici}
Demostra que si $f=m\comp e$ és una factorització imatge i $f$ és ell mateix un monomorfisme, aleshores $e$ és un isomorfisme.
\end{exercici}

\begin{solucio}
Si $f=m\comp e$ és mono i $m$ és mono, aleshores $e$ és mono: de $e\comp u=e\comp v$ se segueix $m\comp e\comp u=m\comp e\comp v$, és a dir $f\comp u=f\comp v$, d'on $u=v$. Però $e$ també és un epimorfisme regular, en particular un epimorfisme. Un morfisme que és alhora mono i epi regular és un isomorfisme: essent el coequalitzador d'una parella que ja iguala (perquè $e$ mono força la parella nucli a ser la diagonal), $e$ coequalitza la parella trivial i és, doncs, invertible. Per tant $\Imatge(f)=[f]$: la imatge d'un mono és ell mateix.
\end{solucio}

\begin{exercici}
Comprova que la imatge és el menor subobjecte pel qual es factoritza $f$: si $f=n\comp g$ amb $n\colon T\rightarrowtail B$ mono, aleshores $\Imatge(f)\leq[n]$.
\end{exercici}

\begin{solucio}
Sigui $f=m\comp e$ la factorització imatge i $f=n\comp g$ una factorització qualsevol a través del mono $n$. Aleshores $n\comp g=m\comp e$. Com que $e$ és el coequalitzador de la parella nucli de $f$ i $n$ és mono, $g$ iguala les dues branques de la parella nucli, de manera que es factoritza per $e$: hi ha $h$ amb $h\comp e=g$. Aleshores $n\comp h\comp e=n\comp g=f=m\comp e$, i com que $e$ és epi, $n\comp h=m$. Aquesta igualtat és precisament una factorització de $m$ a través de $n$, és a dir $\Imatge(f)=[m]\leq[n]$.
\end{solucio}

\section{Estabilitat i regularitat}

\begin{exercici}
A $\cat{Set}$, comprova directament que les imatges són estables per pullback: si es reindexà $f\colon A\to B$ al llarg de $g\colon B'\to B$, aleshores $\Imatge(f')=g^{-1}(\Imatge(f))$.
\end{exercici}

\begin{solucio}
El pullback de $f$ al llarg de $g$ és $A'=\{(b',a)\mid g(b')=f(a)\}$ amb $f'(b',a)=b'$. La seva imatge és $\{b'\in B'\mid\exists a,\ g(b')=f(a)\}=\{b'\mid g(b')\in f(A)\}=g^{-1}(f(A))=g^{-1}(\Imatge(f))$. Per tant $\Imatge(f')=g^{*}(\Imatge(f))$, que és l'enunciat d'estabilitat. Aquesta comprovació explícita és el model de la propietat abstracta que defineix les categories regulars.
\end{solucio}

\begin{exercici}
Demostra que en una categoria regular els epimorfismes regulars són tancats per composició: si $f\colon A\to B$ i $g\colon B\to C$ són epimorfismes regulars, aleshores $g\comp f$ també ho és.
\end{exercici}

\begin{solucio}
Siguin $f\colon A\to B$ i $g\colon B\to C$ epimorfismes regulars, i factoritzem la composició com a imatge, $g\comp f=m\comp e$, amb $e\colon A\to I$ epi regular i $m\colon I\rightarrowtail C$ mono. Sigui $p_1,p_2\colon R\rightrightarrows A$ la parella nucli de $f$. Com que $f\comp p_1=f\comp p_2$,
\[
m\comp e\comp p_1=g\comp f\comp p_1=g\comp f\comp p_2=m\comp e\comp p_2,
\]
i, com que $m$ és mono, $e\comp p_1=e\comp p_2$. Pel lema~\ref{lema:epi-regular-nucli}, $f$ és el coequalitzador de la seva parella nucli; per tant, existeix $k\colon B\to I$ amb $e=k\comp f$. Aleshores
\[
m\comp k\comp f=m\comp e=g\comp f,
\]
i com que $f$ és epi, $m\comp k=g$: el morfisme $g$ es factoritza a través de $m$. Però $g=\id_C\comp g$ és una factorització imatge de $g$, de manera que $\Imatge(g)=[\id_C]$, i per la minimalitat de la imatge (proposició~\ref{prop:imatge-unica}), $[\id_C]\leq[m]$. Com que $[\id_C]$ és el màxim de $\Sub(C)$, $[m]=[\id_C]$ i $m$ és un isomorfisme. Per tant $g\comp f=m\comp e$ és un epi regular seguit d'un isomorfisme, que torna a ser un epi regular.

A $\cat{Set}$, tot plegat es redueix al fet que la composició d'aplicacions exhaustives és exhaustiva.
\end{solucio}

\section{L'existencial}

\begin{exercici}
Verifica les dues direccions de l'adjunció $\exists_f\dashv f^{*}$ a $\cat{Set}$, on $\exists_f(S)=f(S)$ i $f^{*}(T)=f^{-1}(T)$.
\end{exercici}

\begin{solucio}
Siguin $S\subseteq A$ i $T\subseteq B$. Cal veure $f(S)\subseteq T\Leftrightarrow S\subseteq f^{-1}(T)$.

\emph{($\Rightarrow$)} Si $f(S)\subseteq T$ i $a\in S$, aleshores $f(a)\in f(S)\subseteq T$, així que $a\in f^{-1}(T)$; per tant $S\subseteq f^{-1}(T)$.

\emph{($\Leftarrow$)} Si $S\subseteq f^{-1}(T)$ i $b\in f(S)$, aleshores $b=f(a)$ amb $a\in S\subseteq f^{-1}(T)$, d'on $b=f(a)\in T$; per tant $f(S)\subseteq T$.

Aquesta és la forma conjuntista de l'adjunció imatge directa\,--\,antiimatge, i és exactament la regla d'introducció/eliminació de l'existencial quan $f$ és una projecció.
\end{solucio}

\begin{exercici}
Sigui $\pi\colon X\times Y\to X$ la projecció. Descriu $\exists_\pi[s]$ i $\forall_\pi[s]$ per a un subobjecte $[s]\subseteq X\times Y$ a $\cat{Set}$, i relaciona-ho amb els quantificadors.
\end{exercici}

\begin{solucio}
Sigui $S\subseteq X\times Y$ el subconjunt corresponent a $[s]$. Aleshores
\[
\begin{aligned}
\exists_\pi(S)&=\{x\in X\mid \exists y\in Y:\ (x,y)\in S\},\\
\forall_\pi(S)&=\{x\in X\mid \forall y\in Y:\ (x,y)\in S\}.
\end{aligned}
\]
El primer és la imatge de $S$ per $\pi$ (l'existencial); el segon, el conjunt de fibres senceres contingudes en $S$ (l'universal). Són els casos particulars, per a una projecció, de la cadena $\exists_f\dashv f^{*}\dashv\forall_f$ de l'exemple~\ref{ex:quantificadors-set} del capítol d'adjuncions: l'adjunció $\exists_\pi\dashv\pi^{*}$ es comprova com a l'exercici anterior, i $\pi^{*}\dashv\forall_\pi$ amb un argument anàleg, no per dualitat. Així, en una categoria regular tenim l'existencial; l'universal requereix, a més, adjunts per la dreta, que apareixen en contextos més rics (categories de Heyting).
\end{solucio}


\chapter[Classificadors de subobjectes i topos]{\texorpdfstring{Classificadors de subobjectes i introducció als topos}{Classificadors de subobjectes i introducció als topos}}

Els exercicis d'aquest capítol són de familiarització amb el classificador de subobjectes i la definició de topos. No entren en la lògica interna, que pertany a l'estudi posterior.

\section{El classificador}

\begin{exercici}
Comprova amb detall que a $\cat{Set}$ l'objecte $\Omega=\{0,1\}$ amb $\mathsf{true}\colon 1\to\{0,1\}$, $\mathsf{true}(\ast)=1$, és un classificador de subobjectes.
\end{exercici}

\begin{solucio}
Sigui $U\subseteq E$ un subconjunt, amb injecció $m\colon U\hookrightarrow E$. Definim $\chi_U\colon E\to\{0,1\}$ per $\chi_U(x)=1$ si $x\in U$ i $\chi_U(x)=0$ altrament. El pullback de $\mathsf{true}\colon 1\to\{0,1\}$ al llarg de $\chi_U$ és $\{x\in E\mid\chi_U(x)=1\}=U$, amb la injecció com a projecció; per tant el quadrat és cartesià. Per a la unicitat: si $\chi\colon E\to\{0,1\}$ fa el quadrat cartesià, aleshores $\chi^{-1}(1)=U$, cosa que força $\chi=\chi_U$. Així $\{0,1\}$ classifica els subobjectes de qualsevol conjunt.
\end{solucio}

\begin{exercici}\label{pr-ordre-pertinenca}
En una categoria amb classificador de subobjectes, siguin $m\colon U\rightarrowtail E$ i $n\colon V\rightarrowtail E$ dos subobjectes, amb morfismes característics $\chi_m$ i $\chi_n$. Demostra que
\[
[m]\leq[n]
\quad\Longleftrightarrow\quad
\chi_n\comp m=\mathsf{true}\comp\,!_U,
\]
és a dir, que $U$ està contingut en $V$ si i només si tots els \enquote{elements} de $U$ pertanyen a $V$. Dedueix que $\chi_m=\chi_n$ si i només si $[m]=[n]$.
\end{exercici}

\begin{solucio}
Per definició de l'ordre, $[m]\leq[n]$ vol dir que $m$ factoritza a través de $n$. Aplicant la caracterització de la pertinença per elements generalitzats de la Teoria a l'element generalitzat $m\colon U\to E$, això passa si i només si $\chi_n\comp m=\mathsf{true}\comp\,!_U$.

Per a la segona part, si $[m]=[n]$, els dos subobjectes són el mateix i tenen el mateix morfisme característic, per la unicitat de la definició. Recíprocament, si $\chi_m=\chi_n$, aleshores $\chi_n\comp m=\chi_m\comp m=\mathsf{true}\comp\,!_U$, perquè el quadrat de $m$ commuta, i per la primera part $[m]\leq[n]$; simètricament, $[n]\leq[m]$, i per l'antisimetria de $\Sub(E)$, $[m]=[n]$. Així, el morfisme característic determina el subobjecte, que és la injectivitat de la bijecció $\Sub(E)\cong\Hom(E,\Omega)$ vista ara en termes de pertinença.
\end{solucio}

\begin{exercici}\label{pr-mono-equalitzador}
En una categoria amb classificador de subobjectes, sigui $m\colon U\rightarrowtail E$ un subobjecte amb morfisme característic $\chi_m$. Demostra que $m$ és un equalitzador de $\chi_m$ i $\mathsf{true}\comp\,!_E$. Dedueix que un morfisme que és alhora mono i epi és un isomorfisme.
\end{exercici}

\begin{solucio}
\textbf{Equalitza.} Pel quadrat del classificador, $\chi_m\comp m=\mathsf{true}\comp\,!_U=\mathsf{true}\comp\,!_E\comp m$, perquè $!_E\comp m$ i $!_U$ són fletxes cap al terminal i, per tant, coincideixen.

\textbf{És universal.} Sigui $x\colon X\to E$ amb $\chi_m\comp x=\mathsf{true}\comp\,!_E\comp x$. El membre dret és $\mathsf{true}\comp\,!_X$, de manera que, per la pertinença per elements generalitzats de la Teoria, $x$ factoritza a través de $m$. La factorització és única perquè $m$ és mono. Així $m$ és un equalitzador.

\textbf{Mono i epi implica iso.} Suposem, a més, que $m$ és epi. De $\chi_m\comp m=(\mathsf{true}\comp\,!_E)\comp m$ i la cancel·lació per la dreta, $\chi_m=\mathsf{true}\comp\,!_E$. Aleshores $\id_E$ també equalitza les dues fletxes, i per la propietat de l'equalitzador hi ha $h\colon E\to U$ amb $m\comp h=\id_E$. A més, $m\comp(h\comp m)=m=m\comp\id_U$, i com que $m$ és mono, $h\comp m=\id_U$. Per tant $m$ és un isomorfisme.

Al capítol~\ref{cap:categories} vam veure que, en general, un bimorfisme no és un isomorfisme: a $\cat{Top}$, per exemple, una bijecció contínua pot no ser un homeomorfisme. Una categoria amb classificador no admet aquest fenomen. Als Exercicis es proposa treure'n la conseqüència per a $\cat{Top}$.
\end{solucio}

\section{Topos}

\begin{exercici}
Comprova que $\cat{Set}$ satisfà les tres condicions de la definició de topos.
\end{exercici}

\begin{solucio}
\emph{Límits finits}: $\cat{Set}$ té terminal (un punt), productes (producte cartesià) i equalitzadors (subconjunts definits per igualtats), i per tant tots els límits finits. \emph{Classificador}: l'exercici anterior mostra que $\{0,1\}$ n'és un. \emph{Exponencials}: l'exponencial $C^{B}$ és el conjunt d'aplicacions $B\to C$, i vam veure al capítol de tancament cartesià que $\cat{Set}$ és cartesiana tancada. Per tant $\cat{Set}$ és un topos.
\end{solucio}

\begin{exercici}
Demostra que en un topos l'objecte de parts $\mathcal{P}(E)=\Omega^{E}$ satisfà $\Hom(A,\mathcal{P}(E))\cong\Sub(A\times E)$.
\end{exercici}

\begin{solucio}
Component les dues propietats universals disponibles en un topos:
\[
\Hom(A,\Omega^{E})\cong\Hom(A\times E,\Omega)\cong\Sub(A\times E),
\]
on la primera bijecció és la de l'exponencial (currificació) i la segona és la representabilitat del classificador (\ref{prop:omega-representa}). Totes dues són naturals en $A$. Així, els morfismes cap a $\Omega^{E}$ es corresponen amb subobjectes de $A\times E$, és a dir, amb relacions entre $A$ i $E$. A $\cat{Set}$, això recupera que $\mathcal{P}(E)$ és el conjunt de parts i que una aplicació $A\to\mathcal{P}(E)$ és una relació entre $A$ i $E$.
\end{solucio}

\section{Prefeixos}

\begin{exercici}\label{pr-sedassos}
Sigui $\cat{C}$ una categoria petita, $F$ un prefeix sobre $\cat{C}$ i $U\rightarrowtail F$ un subprefeix, és a dir, $U(C)\subseteq F(C)$ per a cada $C$, amb $F(f)$ enviant $U(C)$ dins de $U(D)$ per a cada $f\colon D\to C$. Per a $x\in F(C)$, posem
\[
\chi_U(C)(x)=\{f\colon D\to C\mid F(f)(x)\in U(D)\}.
\]
Comprova que $\chi_U(C)(x)$ és un sedàs sobre $C$, que les aplicacions $\chi_U(C)$ formen una transformació natural $\chi_U\colon F\Rightarrow\Omega$, i que $U$ és el pullback de $\mathsf{true}$ al llarg de $\chi_U$.
\end{exercici}

\begin{solucio}
\textbf{És un sedàs.} Si $f\colon D\to C$ pertany a $\chi_U(C)(x)$ i $g\colon X\to D$, aleshores $F(f\comp g)(x)=F(g)\bigl(F(f)(x)\bigr)$, i com que $F(f)(x)\in U(D)$ i $U$ és un subprefeix, $F(g)$ l'envia dins de $U(X)$. Per tant $f\comp g$ pertany a $\chi_U(C)(x)$.

\textbf{Naturalitat.} Per a $h\colon D\to C$ i $x\in F(C)$, cal veure que $\chi_U(D)(F(h)(x))=h^{*}\bigl(\chi_U(C)(x)\bigr)$. Per a $g\colon X\to D$,
\[
\begin{aligned}
g\in\chi_U(D)\bigl(F(h)(x)\bigr)
&\iff F(g)\bigl(F(h)(x)\bigr)\in U(X)\\
&\iff F(h\comp g)(x)\in U(X)\\
&\iff h\comp g\in\chi_U(C)(x),
\end{aligned}
\]
i l'última condició és la definició de $g\in h^{*}\bigl(\chi_U(C)(x)\bigr)$.

\textbf{Pullback.} Com que els límits de prefeixos es calculen objecte a objecte, n'hi ha prou de veure que, per a cada $C$, $U(C)$ és el conjunt dels $x\in F(C)$ amb $\chi_U(C)(x)$ igual al sedàs màxim. Si $x\in U(C)$, tot $f\colon D\to C$ compleix $F(f)(x)\in U(D)$, de manera que el sedàs és màxim. Recíprocament, si el sedàs és màxim, conté $\id_C$, i aleshores $x=F(\id_C)(x)\in U(C)$. La unicitat de $\chi_U$ es comprova de la mateixa manera, mirant quins morfismes porten $x$ dins de $U$.

A $\cat{Set}$, vist com a prefeixos sobre la categoria d'un sol objecte i només la identitat, els sedassos sobre l'únic objecte són $\varnothing$ i $\{\id\}$, i es recupera $\Omega=\{0,1\}$.
\end{solucio}

\begin{exercici}\label{pr-conjuncio-omega}
En un topos, sigui $\wedge\colon\Omega\times\Omega\to\Omega$ el morfisme característic del subobjecte $\langle\mathsf{true},\mathsf{true}\rangle\colon 1\rightarrowtail\Omega\times\Omega$. Demostra que, per a dos subobjectes $U,V$ de $E$,
\[
\chi_{U\wedge V}=\wedge\comp\langle\chi_U,\chi_V\rangle,
\]
on $U\wedge V$ és la intersecció. Interpreta-ho a $\cat{Set}$.
\end{exercici}

\begin{solucio}
Per l'exercici~\ref{pr-ordre-pertinenca}, n'hi ha prou de veure que tots dos morfismes classifiquen el mateix subobjecte, és a dir, que un element generalitzat $x\colon X\to E$ hi pertany en un cas si i només si hi pertany en l'altre. Per la pertinença per elements generalitzats:
\begin{itemize}
  \item $x$ factoritza a través de $U\wedge V$ si i només si factoritza a través de $U$ i de $V$, perquè la intersecció és el pullback;
  \item és a dir, si i només si $\chi_U\comp x=\mathsf{true}\comp\,!_X$ i $\chi_V\comp x=\mathsf{true}\comp\,!_X$;
  \item és a dir, si i només si $\langle\chi_U,\chi_V\rangle\comp x=\langle\mathsf{true},\mathsf{true}\rangle\comp\,!_X$, que és la pertinença de $\langle\chi_U,\chi_V\rangle\comp x$ al subobjecte $\langle\mathsf{true},\mathsf{true}\rangle$;
  \item és a dir, si i només si $\wedge\comp\langle\chi_U,\chi_V\rangle\comp x=\mathsf{true}\comp\,!_X$.
\end{itemize}
Sigui $W$ el subobjecte classificat per $\wedge\comp\langle\chi_U,\chi_V\rangle$. Els dos subobjectes $U\wedge V$ i $W$ tenen, doncs, els mateixos elements generalitzats. Prenent com a $x$ un representant de $U\wedge V$, obtenim $U\wedge V\leq W$, i prenent-ne un de $W$, $W\leq U\wedge V$. Per tant són iguals, i per la unicitat del morfisme característic, $\chi_{U\wedge V}=\wedge\comp\langle\chi_U,\chi_V\rangle$.

A $\cat{Set}$, $\wedge\colon\{0,1\}\times\{0,1\}\to\{0,1\}$ val $1$ exactament en $(1,1)$: és la taula de veritat de la conjunció, i la fórmula diu que la funció característica d'una intersecció és el producte, o el mínim, de les funcions característiques.
\end{solucio}

\begin{exercici}\label{pr-reindexacio-graph}
Sigui $G$ el cicle de l'exemple~\ref{ex:omega-graph} de la Teoria, amb vèrtexs $a,b,c$, arestes $e\colon a\to b$, $f\colon b\to c$, $k\colon c\to a$, i el subgraf $H$ de vèrtexs $a,b$ i aresta $e$. Sigui $G'$ el cicle de longitud sis, amb vèrtexs $p_0,\dots,p_5$ i arestes $u_i\colon p_i\to p_{i+1}$ (índexs mòdul $6$), i sigui $F\colon G'\to G$ el morfisme que \enquote{dona dues voltes}: $p_i\mapsto a,b,c$ i $u_i\mapsto e,f,k$ segons que $i$ sigui $0,1,2$ mòdul $3$. Calcula $F^{*}H$ i comprova que $\chi_{F^{*}H}=\chi_H\comp F$.
\end{exercici}

\begin{solucio}
\textbf{La reindexació.} Es calcula per antiimatge, component a component (exemple de subgrafs del capítol~\ref{cap:subobjectes}). Els vèrtexs de $F^{*}H$ són els que van a $a$ o a $b$, és a dir, $p_0,p_1,p_3,p_4$. Les arestes són les que van a $e$, és a dir, $u_0$ i $u_3$.

\textbf{La precomposició.} Amb la taula de $\chi_H$ de la Teoria ($a,b\mapsto1$, $c\mapsto0$, $e\mapsto\top$, $f\mapsto\sigma$, $k\mapsto\tau$), el morfisme $\chi_H\comp F$ envia $p_0,p_1,p_3,p_4\mapsto1$ i $p_2,p_5\mapsto0$, i
\[
u_0,u_3\mapsto\top,\qquad u_1,u_4\mapsto\sigma,\qquad u_2,u_5\mapsto\tau.
\]

\textbf{El morfisme característic de $F^{*}H$.} Apliquem directament la regla de la Teoria a $F^{*}H\subseteq G'$. Les arestes $u_0$ i $u_3$ són al subgraf, i van a $\top$. Les arestes $u_1\colon p_1\to p_2$ i $u_4\colon p_4\to p_5$ tenen l'origen dins i el destí fora, i van a $\sigma$. Les arestes $u_2\colon p_2\to p_3$ i $u_5\colon p_5\to p_0$ tenen l'origen fora i el destí dins, i van a $\tau$. Sobre els vèrtexs coincideix també amb $\chi_H\comp F$. Així $\chi_{F^{*}H}=\chi_H\comp F$: substituir és precompondre, ara en un topos que no és $\cat{Set}$, i al llarg d'un morfisme que no és injectiu.
\end{solucio}

\section{Connectius i igualtat}

\begin{exercici}\label{pr-igualtat-omega}
En una categoria amb productes binaris i classificador, sigui $\delta_A=\langle\id_A,\id_A\rangle\colon A\to A\times A$. Comprova que $\delta_A$ és un monomorfisme i que, per a dos elements generalitzats $x,y\colon X\to A$,
\[
\chi_{\delta_A}\comp\langle x,y\rangle=\mathsf{true}\comp\,!_X
\quad\Longleftrightarrow\quad
x=y.
\]
Descriu $\chi_{\delta_A}$ a $\cat{Set}$.
\end{exercici}

\begin{solucio}
\textbf{És mono.} La projecció $\pi_1$ compleix $\pi_1\comp\delta_A=\id_A$, de manera que $\delta_A$ té una retracció i és un monomorfisme escindit.

\textbf{Pertinença a la diagonal.} Per la pertinença per elements generalitzats de la Teoria, la igualtat de l'esquerra equival a dir que $\langle x,y\rangle$ factoritza a través de $\delta_A$, és a dir, que hi ha $z\colon X\to A$ amb $\langle z,z\rangle=\langle x,y\rangle$. Component amb les projeccions, això vol dir $x=z=y$. Recíprocament, si $x=y$, n'hi ha prou de prendre $z=x$.

\textbf{A $\cat{Set}$.} $\chi_{\delta_A}(a,b)=1$ si $a=b$ i $0$ altrament: és la delta de Kronecker. En termes lògics, $\chi_{\delta_A}$ és el predicat d'igualtat \enquote{$x=y$} sobre $A$, i la diagonal n'és l'extensió. És la fila \enquote{igualtat} del mapa del capítol~\ref{cap:pont}.
\end{solucio}

\begin{exercici}\label{pr-disjuncio-negacio-set}
A $\cat{Set}$, amb $\Omega=\{0,1\}$, sigui ${\vee}\colon\Omega\times\Omega\to\Omega$ el morfisme característic del subconjunt $\{(0,1),(1,0),(1,1)\}$, i sigui $\neg\colon\Omega\to\Omega$ el del subconjunt $\{0\}$. Comprova que, per a subconjunts $U,V\subseteq E$,
\[
\chi_{U\cup V}={\vee}\comp\langle\chi_U,\chi_V\rangle,
\qquad
\chi_{E\setminus U}=\neg\comp\chi_U.
\]
Explica per què la segona igualtat no té un anàleg directe a $\cat{Graph}$.
\end{exercici}

\begin{solucio}
Per a $x\in E$, $({\vee}\comp\langle\chi_U,\chi_V\rangle)(x)=1$ si i només si $(\chi_U(x),\chi_V(x))$ és en $\{(0,1),(1,0),(1,1)\}$, és a dir, si i només si $x\in U$ o $x\in V$. Les dues aplicacions valen, doncs, $1$ exactament a $U\cup V$, i coincideixen. Anàlogament, $(\neg\comp\chi_U)(x)=1$ si i només si $\chi_U(x)=0$, és a dir, si i només si $x\notin U$. Els morfismes $\vee$ i $\neg$ són les taules de veritat de la disjunció i de la negació, i les dues fórmules diuen que la unió i el complement es calculen \enquote{punt a punt} sobre els valors de veritat, com la intersecció a l'exercici~\ref{pr-conjuncio-omega}.

A $\cat{Graph}$, el complement de vèrtexs i d'arestes d'un subgraf no és, en general, un subgraf: pot contenir una aresta sense un dels seus extrems. El millor substitut és el subgraf més gran disjunt de $U$, però, com hem vist a l'exemple~\ref{ex:omega-graph} de la Teoria, la seva unió amb $U$ pot no ser tot el graf. No hi ha, doncs, cap morfisme $\Omega\to\Omega$ que faci el paper de $\neg$ amb la propietat $U\cup\neg U=E$.
\end{solucio}

\fi

% ============================================================
% PART III — EXERCICIS PROPOSATS
% ============================================================
\part{Exercicis proposats}
\setcounter{chapter}{0}
\chapter{Categories}

Els exercicis següents completen els resolts al capítol corresponent de la part de Pràctica. Cada enunciat va acompanyat d'una indicació breu per si el lector s'encalla; convé, però, intentar-lo primer sense mirar-la. La marca \enquote{$\star$} assenyala un exercici de consolidació i \enquote{$\star\star$} un d'ampliació.

\section{Categories i axiomes}

\begin{exercici}
Comprova que un conjunt qualsevol $X$ dona lloc a una categoria \textbf{discreta}: els objectes són els elements de $X$ i els únics morfismes són les identitats. Verifica els axiomes. \hfill{\normalfont$\star$}
\begin{indicacio}
No hi ha cap morfisme entre objectes diferents; l'única composició possible en cada objecte és $\id_x\comp\id_x$.
\end{indicacio}
\end{exercici}

\begin{exercici}
Sigui $\cat{C}$ una categoria. Comprova que fixar un objecte $A$ i considerar només les fletxes $A\to A$, amb la composició de $\cat{C}$, dona un \textbf{monoide}. Quin és el seu element neutre? \hfill{\normalfont$\star$}
\begin{indicacio}
Els morfismes $A\to A$ es poden compondre entre ells i el resultat torna a anar d'$A$ a $A$. El neutre és $\id_A$.
\end{indicacio}
\end{exercici}

\begin{exercici}[Ampliació, per a lectors amb topologia]
Comprova que els espais topològics amb les aplicacions contínues formen una categoria $\cat{Top}$. Cal usar dues propietats de les aplicacions contínues: quines? \hfill{\normalfont$\star$}
\begin{indicacio}
La composició de contínues és contínua, i la identitat és contínua.
\end{indicacio}
\end{exercici}

\section{Fletxes, camins i diagrames commutatius}

\begin{exercici}
Escriu l'equació que expressa la commutativitat del triangle
\[
\begin{tikzpicture}[baseline=(m.center)]
  \matrix (m) [matrix of math nodes, column sep=3.3em, row sep=2.6em] {
    A & B \\
      & C \\
  };
  \draw[-{Stealth[scale=1.05]}] (m-1-1) -- node[above] {$f$} (m-1-2);
  \draw[-{Stealth[scale=1.05]}] (m-1-2) -- node[right] {$g$} (m-2-2);
  \draw[-{Stealth[scale=1.05]}] (m-1-1) -- node[below left] {$h$} (m-2-2);
\end{tikzpicture}
\]
i explica quins són els dos camins que es comparen. \hfill{\normalfont$\star$}
\begin{indicacio}
Tots dos camins comencen a $A$ i acaben a $C$; un passa per $B$ i l'altre és la fletxa directa.
\end{indicacio}
\end{exercici}

\begin{exercici}
Representa mitjançant un quadrat commutatiu l'equació
\[
q\comp u=v\comp p,
\]
on $u\colon A\to B$, $q\colon B\to D$, $p\colon A\to C$ i $v\colon C\to D$. \hfill{\normalfont$\star$}
\begin{indicacio}
Col·loca $A$ a l'angle superior esquerre i $D$ a l'inferior dret. Els dos membres de l'equació són dos camins d'$A$ a $D$.
\end{indicacio}
\end{exercici}

\begin{exercici}
Dona un exemple de quatre aplicacions entre conjunts que formin un quadrat que \emph{no} commuti. Indica explícitament un element sobre el qual les dues composicions donin resultats diferents. \hfill{\normalfont$\star$}
\begin{indicacio}
Pots treballar amb conjunts molt petits, per exemple $\{0,1\}$, i comparar els dos camins sobre l'element $0$.
\end{indicacio}
\end{exercici}

\section{Primers exemples}

\begin{exercici}
En un sistema deductiu $T$, suposem
\[
A\vdash_T B,
\qquad
B\vdash_T C,
\qquad
C\vdash_T A.
\]
Dibuixa el diagrama corresponent a $\cat{Prov}_T$ i comprova que $A$, $B$ i $C$ són isomorfs dos a dos. \hfill{\normalfont$\star$}
\begin{indicacio}
La transitivitat dona també $A\vdash_T C$, $B\vdash_T A$ i $C\vdash_T B$. En una categoria prima, dues fletxes en sentits oposats ja donen un isomorfisme.
\end{indicacio}
\end{exercici}

\begin{exercici}
Comprova directament que la composició de dos morfismes de grafs és un morfisme de grafs. Si $F\colon G\to H$ i $K\colon H\to L$, verifica les equacions d'origen i destí per a $K\comp F$. \hfill{\normalfont$\star$}
\begin{indicacio}
Per a l'origen, parteix de
\[
F_V\comp s_G=s_H\comp F_E,
\qquad
K_V\comp s_H=s_L\comp K_E,
\]
i compon les igualtats adequadament. Fes el mateix amb $t$.
\end{indicacio}
\end{exercici}

\begin{exercici}
Explica amb precisió quina informació té un graf dirigit que encara no té l'estructura d'una categoria. Després indica què afegeix la categoria lliure sobre aquest graf. \hfill{\normalfont$\star$}
\begin{indicacio}
Un graf té vèrtexs, arestes, origen i destí. Una categoria necessita, a més, identitats, composició i axiomes; la categoria lliure els genera mitjançant camins.
\end{indicacio}
\end{exercici}

\section{Isomorfismes}

\begin{exercici}
Comprova que a la categoria $\cat{Mon}$ dels monoides i els homomorfismes, un isomorfisme és exactament un homomorfisme bijectiu. Compara-ho amb el que pot passar a $\cat{Pos}$. \hfill{\normalfont$\star\star$}
\begin{indicacio}
Comprova que la inversa d'un homomorfisme bijectiu de monoides preserva l'operació i el neutre, de manera que torna a ser un homomorfisme. Per a $\cat{Pos}$, en canvi, la inversa d'una aplicació monòtona bijectiva pot no ser monòtona. Fes el contrast explícit amb un exemple: pren el conjunt $\{0,1\}$ amb l'ordre \emph{discret} (cap dels dos elements comparable amb l'altre) com a domini, i el mateix conjunt amb l'ordre $0<1$ com a codomini. La identitat subjacent és una aplicació monòtona (trivialment, perquè al domini no hi ha cap parella comparable) i bijectiva, però la seva inversa ---la identitat de l'ordre $0<1$ cap a l'ordre discret--- no és monòtona, ja que hauria d'enviar $0<1$ a dos elements incomparables. Per tant no és un isomorfisme a $\cat{Pos}$, tot i ser una bijecció monòtona.
\end{indicacio}
\end{exercici}

\begin{exercici}
Comprova que si $g\comp f$ és un isomorfisme, no cal que ho siguin ni $f$ ni $g$. Busca'n un exemple senzill a $\cat{Set}$ i representa'l amb fletxes. \hfill{\normalfont$\star$}
\begin{indicacio}
Pensa en una injecció seguida d'una projecció entre conjunts de mides diferents.
\end{indicacio}
\end{exercici}

\section{Classes especials de morfismes}

\begin{exercici}
Comprova que tot isomorfisme té secció. Després, prenent $A=\{0,1\}$, $B=\{\ast\}$ i $f$ l'única aplicació $A\to B$, dona una secció $s\colon B\to A$ explícita i comprova que $f$ no és isomorfisme. Observa que en aquest cas $f$ té \emph{dues} seccions distintes: quines? \hfill{\normalfont$\star$}
\begin{indicacio}
Tria $s(\ast)=0$ o $s(\ast)=1$. En tots dos casos $f\comp s=\id_B$, però $s\comp f\neq\id_A$. Les dues tries donen les dues seccions.
\end{indicacio}
\end{exercici}

\begin{exercici}
A la categoria $\cat{Set}$, exhibeix morfismes que il·lustrin els tres casos possibles pel que fa a l'existència de seccions: un morfisme sense cap secció, un amb una única secció i un amb més d'una. \hfill{\normalfont$\star\star$}
\begin{indicacio}
Per a \emph{cap} secció, cerca una aplicació no exhaustiva, com $\emptyset\to\{\ast\}$ o una injecció no exhaustiva. Per a \emph{una única}, una aplicació exhaustiva on cada element del codomini tingui una sola antiimatge (una bijecció). Per a \emph{diverses}, una aplicació exhaustiva amb alguna fibra de més d'un element: cada elecció d'antiimatges dona una secció distinta. Recorda que a $\cat{Set}$ triar una secció d'una aplicació exhaustiva és triar un element a cada fibra.
\end{indicacio}
\end{exercici}

\begin{exercici}
Demostra la caracterització dual de l'exercici resolt a la part de pràctica: un morfisme $f$ és un isomorfisme si i només si és un epimorfisme escindit (té secció) i un monomorfisme. \hfill{\normalfont$\star$}
\begin{indicacio}
És el dual de \enquote{mono escindit i epi $\Rightarrow$ iso}: parteix d'una secció $s$ amb $f\comp s=\id_B$ i fes servir que $f$ és mono per veure que $s\comp f=\id_A$.
\end{indicacio}
\end{exercici}

\begin{exercici}
Sigui $\cat{C}$ una categoria prima. Per l'exercici resolt~\ref{exr:pr-prima-mono-epi} de la Pràctica, tot morfisme de $\cat{C}$ és un bimorfisme. Demostra que un morfisme $f\colon A\to B$ de $\cat{C}$ és un isomorfisme si i només si existeix algun morfisme $B\to A$. Aplica-ho a la fletxa $u\colon 0\to 1$ de $\cat{2}$ i a dues fórmules interdeduïbles i distintes de $\cat{Prov}_T$. \hfill{\normalfont$\star$}
\begin{indicacio}
Si $g\colon B\to A$ existeix, $g\comp f$ i $\id_A$ pertanyen tots dos a $\Hom(A,A)$, que té com a màxim un element; raona igual amb $f\comp g$ i $\id_B$. A $\cat{2}$ no hi ha cap fletxa $1\to 0$, de manera que $u$ és un bimorfisme que no és iso. A $\cat{Prov}_T$, en canvi, la fletxa que expressa $p\vdash_T q$ és un isomorfisme que no és una identitat quan també $q\vdash_T p$.
\end{indicacio}
\end{exercici}

\begin{exercici}
Estudia quin factor hereta la propietat en una composició. Demostra que si $g\comp f$ és un monomorfisme, aleshores $f$ és un monomorfisme; i que, dualment, si $g\comp f$ és un epimorfisme, aleshores $g$ és un epimorfisme. Comprova amb un contraexemple que, en el primer cas, $g$ no cal que sigui mono, i en el segon, $f$ no cal que sigui epi. \hfill{\normalfont$\star\star$}
\begin{indicacio}
Per al mono: de $f\comp u=f\comp v$, compon per l'esquerra amb $g$ per obtenir $(g\comp f)\comp u=(g\comp f)\comp v$ i aplica que $g\comp f$ és mono. El cas de l'epi és el dual: compon per la dreta. Observa l'asimetria: d'un mono $g\comp f$ hereta el factor de la \emph{dreta} ($f$), i d'un epi, el de l'\emph{esquerra} ($g$). Per als contraexemples, pensa en una injecció i una projecció a $\cat{Set}$.
\end{indicacio}
\end{exercici}

\begin{exercici}
Sigui $M=\{1,e\}$ el monoide amb $e\cdot e=e$ i $1$ neutre, i sigui $\cat{B}M$ la categoria d'un sol objecte associada. Decideix si $e$ és un monomorfisme, si és un epimorfisme i si té secció o retracció. \hfill{\normalfont$\star\star$}
\begin{indicacio}
Compara $e\comp 1$ amb $e\comp e$, i $1\comp e$ amb $e\comp e$. Per a les inverses unilaterals, només hi ha dos candidats.
\end{indicacio}
\end{exercici}

\begin{exercici}
A $\cat{Set}$, demostra que tota aplicació exhaustiva $f\colon A\to B$ és un epimorfisme sense triar cap secció de $f$. Explica després per què l'enunciat \enquote{tota aplicació exhaustiva té una secció} no és una conseqüència d'aquest fet, sinó un enunciat diferent, equivalent a l'axioma de l'elecció. \hfill{\normalfont$\star\star$}
\begin{indicacio}
Si $g\comp f=h\comp f$ i $y\in B$, n'hi ha prou de tenir \emph{algun} $x$ amb $f(x)=y$ per concloure $g(y)=h(y)$; no cal triar-ne un per a cada $y$ alhora. Una secció, en canvi, és precisament una tria simultània d'una antiimatge per a cada element de $B$.
\end{indicacio}
\end{exercici}

\section{Construccions sobre categories}

\begin{exercici}
Descriu la categoria producte $\cat{3}\times\cat{3}$, on $\cat{3}$ és la categoria associada a l'ordre $0<1<2$. Quants objectes i quants morfismes té? Dibuixa les fletxes bàsiques i explica quines altres fletxes apareixen per composició. \hfill{\normalfont$\star\star$}
\begin{indicacio}
$\cat{3}$ té tres objectes i sis morfismes. En el producte, els objectes i els morfismes són parelles.
\end{indicacio}
\end{exercici}

\begin{exercici}
Descriu explícitament la categoria $\cat{3}/2$ sobre l'objecte $2$ de $\cat{3}$, on $\cat{3}$ és la categoria associada a l'ordre $0<1<2$. Enumera'n els objectes i els morfismes i identifica, si pots, una categoria finita isomorfa a aquesta categoria sobre un objecte. \hfill{\normalfont$\star$}
\begin{indicacio}
Els objectes de $\cat{3}/2$ són les fletxes de $\cat{3}$ amb codomini $2$, i un morfisme entre dos d'aquests objectes és una fletxa de $\cat{3}$ que fa commutar el triangle corresponent. Com que $\cat{3}$ és prima, tots els triangles commuten. Observa també que de cada objecte de $\cat{3}$ surt exactament una fletxa cap a $2$.
\end{indicacio}
\end{exercici}

\begin{exercici}
Compara directament $(\cat{C}\times\cat{D})\op$ i $\cat{C}\op\times\cat{D}\op$. Comprova que tenen els mateixos objectes i que els seus morfismes, identitats i composicions s'identifiquen component a component. \hfill{\normalfont$\star$}
\begin{indicacio}
Un morfisme $(A,X)\to(B,Y)$ a $(\cat{C}\times\cat{D})\op$ és, per definició, un morfisme $(B,Y)\to(A,X)$ a $\cat{C}\times\cat{D}$. Escriu què significa això en cadascuna de les dues components.
\end{indicacio}
\end{exercici}

\begin{exercici}
Per a cadascuna de les subcategories següents de $\cat{Set}$, digues si és \textbf{plena}, \textbf{ampla}, totes dues coses o cap: (a)~$\cat{FinSet}$, amb els conjunts finits per objectes i totes les aplicacions entre ells; (b)~la subcategoria amb tots els conjunts per objectes i només les aplicacions injectives; (c)~la subcategoria amb tots els conjunts per objectes i totes les aplicacions (és a dir, $\cat{Set}$ mateixa); (d)~la subcategoria amb els conjunts finits per objectes i només les aplicacions injectives entre ells. \hfill{\normalfont$\star$}
\begin{indicacio}
Recorda que \emph{plena} vol dir que es conserven \emph{tots} els morfismes entre els objectes escollits (la subcategoria queda determinada només pels objectes), i \emph{ampla} vol dir que es conserven \emph{tots} els objectes. Comprova cada cas contra aquestes dues condicions per separat: poden fallar independentment. En particular, mira si (d) és plena o ampla, i observa que una subcategoria pot no ser ni l'una ni l'altra.
\end{indicacio}
\end{exercici}

\section{Principi de dualitat}

\begin{exercici}
Escriu el dual de l'afirmació següent i justifica'l sense fer una nova demostració completa:
\begin{quote}
Si $f\colon A\to B$ té una retracció, aleshores $f$ és un monomorfisme.
\end{quote}
\hfill{\normalfont$\star$}
\begin{indicacio}
Passa a la categoria oposada i identifica quina noció és dual de \enquote{retracció} i quina és dual de \enquote{monomorfisme}. Recorda que, en dualitzar, s'inverteixen tant el sentit de les fletxes com l'ordre de les composicions.
\end{indicacio}
\end{exercici}

\section{Grafs i categories lliures}

\begin{exercici}
Classifica, \textbf{fins a isomorfisme}, els grafs dirigits que generen una categoria lliure amb exactament \textbf{quatre} morfismes. Recorda que les identitats compten com a morfismes. \hfill{\normalfont$\star\star$}
\begin{indicacio}
Cada vèrtex aporta una identitat i cada aresta un camí de longitud~$1$: quantes arestes i quants vèrtexs pot tenir, com a màxim, el graf? Què passaria si hi hagués un cicle dirigit? Després, separa els casos segons el nombre de vèrtexs.
\end{indicacio}
\end{exercici}

\begin{exercici}
Considera el graf
\[
A\xrightarrow{f}B\xrightarrow{g}C.
\]
Llista tots els morfismes de la categoria lliure que genera, dibuixa també la fletxa composta i comprova que hi ha exactament sis morfismes. \hfill{\normalfont$\star$}
\begin{indicacio}
Compta els tres camins buits, les dues arestes $f$ i $g$, i el camí compost de longitud dos.
\end{indicacio}
\end{exercici}

\begin{exercici}
Per què un graf amb un cicle, per exemple una aresta $a\colon A\to A$, genera sempre una categoria lliure amb infinits morfismes? \hfill{\normalfont$\star$}
\begin{indicacio}
El cicle es pot recórrer tantes vegades com es vulgui: $a,a^2,a^3,\ldots$, i cada repetició dona un camí diferent.
\end{indicacio}
\end{exercici}

\section{Qüestions de mida}

\begin{exercici}
Comprova que un preordre, vist com a categoria, és una categoria \textbf{petita} sempre que el conjunt subjacent $P$ sigui un conjunt. Quants morfismes hi ha, com a màxim, entre dos objectes? \hfill{\normalfont$\star$}
\begin{indicacio}
Entre dos objectes hi ha zero o un morfisme; la col·lecció total de morfismes es pot identificar amb un subconjunt de $P\times P$.
\end{indicacio}
\end{exercici}

\begin{exercici}
Explica per què el fet que tots els objectes d'una categoria siguin finits no implica que la categoria sigui petita. Aplica-ho a $\cat{FinSet}$, la categoria que té per objectes tots els conjunts finits i per morfismes totes les aplicacions entre ells. \hfill{\normalfont$\star$}
\begin{indicacio}
La mida de cada objecte i la mida de la col·lecció de tots els objectes són qüestions diferents. Per a cada conjunt $x$, el conjunt $\{x\}$ és finit: quants objectes té, doncs, $\cat{FinSet}$? Compara-ho amb la paradoxa de Russell.
\end{indicacio}
\end{exercici}

\begin{exercici}
Explica per què $\cat{Graph}$ és localment petita: donats dos grafs dirigits $G$ i $H$, per què els morfismes de grafs $G\to H$ formen un conjunt? \hfill{\normalfont$\star$}
\begin{indicacio}
Un morfisme de grafs és una parella d'aplicacions $V_G\to V_H$ i $E_G\to E_H$ que satisfà dues condicions. Aquestes parelles formen un subconjunt d'un producte de dos conjunts d'aplicacions.
\end{indicacio}
\end{exercici}

\chapter{Functors}

Els exercicis següents completen els resolts al capítol corresponent de la part de Pràctica. Cada enunciat va acompanyat d'una indicació breu per si el lector s'encalla; convé, però, intentar-lo primer sense mirar-la. La marca \enquote{$\star$} assenyala un exercici de consolidació i \enquote{$\star\star$} un d'ampliació.

\section{Definició de functor}

\begin{exercici}
Comprova que el functor identitat $\id_{\cat{C}}\colon\cat{C}\to\cat{C}$ i, per a cada objecte $D$ d'una categoria $\cat{D}$, el functor constant $\Delta_D\colon\cat{C}\to\cat{D}$ ---que envia tot objecte a $D$ i tot morfisme a $\id_D$ (exemple~\ref{ex:functor-constant})--- són functors. \hfill{\normalfont$\star$}
\begin{indicacio}
Per al functor constant, comprova que $\Delta_D(\id_A)=\id_D$ i que $\Delta_D(g\comp f)=\id_D=\id_D\comp\id_D=\Delta_D(g)\comp\Delta_D(f)$.
\end{indicacio}
\end{exercici}

\begin{exercici}\label{ex:identitat-necessaria}
Considera una assignació que respecti dominis, codominis i composició, però de la qual \emph{no} se suposi que preservi identitats. Comprova, amb un contraexemple senzill, que la preservació de la identitat no se'n dedueix, i conclou que la condició $F(\id_A)=\id_{F(A)}$ és realment necessària a la definició de functor. \hfill{\normalfont$\star$}
\begin{indicacio}
Pren el monoide de dos elements $\{1,e\}$ amb $e\comp e=e$ i $e\comp m=m\comp e=e$ per a tot $m$, vist com a categoria d'un sol objecte. L'assignació que envia cada morfisme $m$ a $e$ respecta la composició, perquè $F(m\comp n)=e=e\comp e=F(m)\comp F(n)$, però envia la identitat $1$ a $e\neq 1$. Per tant no preserva identitats malgrat respectar la composició.
\end{indicacio}
\end{exercici}

\section{Exemples de functors}

\begin{exercici}
Comprova que hi ha un functor $V\colon\cat{Graph}\to\cat{Set}$ que envia cada graf $G$ al seu conjunt de vèrtexs $V_G$ i cada morfisme de grafs $F=(F_V,F_E)\colon G\to H$ a la seva component sobre vèrtexs $F_V\colon V_G\to V_H$. Quines propietats de la composició i de les identitats de $\cat{Graph}$ en garanteixen la functorialitat? Fes el mateix per al functor d'arestes $E\colon\cat{Graph}\to\cat{Set}$. \hfill{\normalfont$\star$}
\begin{indicacio}
La identitat d'un graf té per component sobre vèrtexs la identitat de $V_G$, i la composició de morfismes de grafs es fa component a component: $(F'\comp F)_V=F'_V\comp F_V$. Quan hagis estudiat els homfunctors (subsecció~\ref{subsec:homfunctors}), compara $V$ amb $\Hom_{\cat{Graph}}(\mathbf V,-)$, on $\mathbf V$ és el graf d'un sol vèrtex.
\end{indicacio}
\end{exercici}

\begin{exercici}[Ampliació, per a lectors amb topologia]
Comprova que hi ha un functor oblit $U\colon\cat{Top}\to\cat{Set}$ i descriu explícitament què fa sobre objectes i sobre morfismes. Quines propietats de les aplicacions contínues en garanteixen la functorialitat? \hfill{\normalfont$\star$}
\begin{indicacio}
$U$ envia cada espai al seu conjunt de punts i cada aplicació contínua a la mateixa aplicació. Cal que la identitat sigui contínua i que la composició de contínues sigui contínua.
\end{indicacio}
\end{exercici}

\begin{exercici}[Ampliació, per a lectors amb àlgebra]
Sigui $G$ un grup vist com a categoria $\cat{B}G$ amb un sol objecte. Comprova que un functor $\cat{B}G\to\cat{Set}$ és el mateix que una \textbf{acció} de $G$ sobre un conjunt, i que la imatge de cada element del grup és necessàriament una \emph{bijecció}. \hfill{\normalfont$\star$}
\begin{indicacio}
Segueix el cas del monoide. La diferència és que en un grup cada element té invers; usa que un functor preserva isomorfismes i que tot morfisme de $\cat{B}G$ és un isomorfisme.
\end{indicacio}
\end{exercici}

\begin{exercici}
Descriu un functor $\cat{2}\to\cat{C}$, on $\cat{2}$ és la categoria associada a l'ordre $0<1$. Comprova que triar un tal functor equival a triar un morfisme de $\cat{C}$. \hfill{\normalfont$\star$}
\begin{indicacio}
El functor ha d'enviar l'única fletxa no trivial $0\to1$ a un morfisme de $\cat{C}$; els dos objectes es determinen com el seu domini i el seu codomini.
\end{indicacio}
\end{exercici}

\section{Composició de functors i la categoria \texorpdfstring{$\cat{Cat}$}{Cat}}

\begin{exercici}
Comprova que un isomorfisme a $\cat{Cat}$ ---és a dir, un functor $F\colon\cat{C}\to\cat{D}$ amb un functor invers $G$ tal que $G\comp F=\id_{\cat{C}}$ i $F\comp G=\id_{\cat{D}}$--- és bijectiu sobre objectes i sobre morfismes. \hfill{\normalfont$\star\star$}
\begin{indicacio}
De $G\comp F=\id_{\cat{C}}$ i $F\comp G=\id_{\cat{D}}$ dedueix que $F$ té inversa com a aplicació sobre objectes i com a aplicació sobre morfismes.
\end{indicacio}
\end{exercici}

\begin{exercici}
Siguin $F\colon\cat{E}\to\cat{C}$ i $G\colon\cat{E}\to\cat{D}$ dos functors. Defineix un functor $\langle F,G\rangle\colon\cat{E}\to\cat{C}\times\cat{D}$ per $\langle F,G\rangle(E)=(F(E),G(E))$ i $\langle F,G\rangle(h)=(F(h),G(h))$. Comprova que és un functor, que
\[
\pi_1\comp\langle F,G\rangle=F,\qquad \pi_2\comp\langle F,G\rangle=G,
\]
on $\pi_1,\pi_2$ són les projeccions de l'exemple~\ref{ex:projeccions-diagonal}, i que és l'únic functor $\cat{E}\to\cat{C}\times\cat{D}$ amb aquestes dues propietats. Quin functor s'obté amb $\cat{E}=\cat{C}=\cat{D}$ i $F=G=\id_{\cat{C}}$? \hfill{\normalfont$\star$}
\begin{indicacio}
La functorialitat es comprova component a component, com per a $\Delta$. Per a la unicitat, si $H\colon\cat{E}\to\cat{C}\times\cat{D}$ compleix les dues igualtats, escriu $H(E)$ i $H(h)$ com a parelles i mira què en diuen $\pi_1$ i $\pi_2$. És la propietat universal del producte, ara a nivell de categories.
\end{indicacio}
\end{exercici}

\section{Functors i isomorfismes}

\begin{exercici}
Comprova que si $F\colon\cat{C}\to\cat{D}$ és un functor i $A\cong B$ a $\cat{C}$, aleshores $F(A)\cong F(B)$ a $\cat{D}$. És a dir, els functors preserven la relació d'isomorfisme entre objectes. \hfill{\normalfont$\star$}
\begin{indicacio}
Aplica $F$ a un isomorfisme $A\to B$ i usa que la imatge d'un isomorfisme és un isomorfisme.
\end{indicacio}
\end{exercici}

\begin{exercici}
Dona un exemple de dos objectes no isomorfs que un functor envia a objectes isomorfs. Això mostra que els functors no reflecteixen, en general, la relació d'isomorfisme entre objectes. \hfill{\normalfont$\star$}
\begin{indicacio}
Un functor constant envia tots els objectes al mateix objecte, i per tant tots a objectes isomorfs entre ells.
\end{indicacio}
\end{exercici}

\section{Functors covariants i contravariants}

\begin{exercici}
Donat un functor $F\colon\cat{C}\to\cat{D}$, comprova que la mateixa regla sobre objectes i morfismes defineix un functor
\[
F\op\colon\cat{C}\op\to\cat{D}\op,
\]
el \textbf{functor oposat} de $F$. Comprova que $F$ és fidel (respectivament ple) si i només si $F\op$ ho és. \hfill{\normalfont$\star$}
\begin{indicacio}
Sobre objectes, $F\op=F$; sobre morfismes, $F\op(f)=F(f)$, però ara llegit entre els mateixos objectes a les categories oposades. Compondre a $\cat{C}\op$ i a $\cat{D}\op$ és compondre a $\cat{C}$ i a $\cat{D}$ en l'ordre invers, de manera que les dues condicions de functor es transporten. Els conjunts de morfismes són els mateixos, així que injectivitat i exhaustivitat es conserven.
\end{indicacio}
\end{exercici}

\begin{exercici}
Fixat un objecte $A$ d'una categoria localment petita, comprova que $\Hom(A,-)$ envia isomorfismes a bijeccions: si $f\colon X\to Y$ és un isomorfisme, aleshores
\[
f_*\colon\Hom(A,X)\to\Hom(A,Y)
\]
és una bijecció. \hfill{\normalfont$\star$}
\begin{indicacio}
$\Hom(A,-)$ és un functor i, per tant, preserva isomorfismes; a $\cat{Set}$ els isomorfismes són les bijeccions. Alternativament, comprova directament que $(f^{-1})_*$ és la inversa de $f_*$.
\end{indicacio}
\end{exercici}

\section{Functors fidels i plens}

\begin{exercici}
Comprova que la composició de dos functors fidels és fidel, i que la composició de dos functors plens és plena. Dedueix que la composició de functors plens i fidels és plena i fidel. \hfill{\normalfont$\star$}
\begin{indicacio}
Treballa amb les aplicacions $F_{A,B}$ sobre conjunts de morfismes: la composició de dues injeccions és injectiva, i la de dues exhaustives, exhaustiva.
\end{indicacio}
\end{exercici}

\begin{exercici}
Demostra que un functor fidel reflecteix monomorfismes i epimorfismes: si $F$ és fidel i $F(f)$ és mono (respectivament epi), aleshores $f$ és mono (respectivament epi). Indica on fas servir que $F$ és fidel, i comprova que no cal que $F$ sigui ple. \hfill{\normalfont$\star$}
\begin{indicacio}
Parteix d'una igualtat de composicions a la categoria d'origen, aplica-hi el functor i cancel·la a la d'arribada.
\end{indicacio}
\end{exercici}

\begin{exercici}
Comprova que el functor diagonal $\Delta\colon\cat{C}\to\cat{C}\times\cat{C}$ de l'exemple~\ref{ex:projeccions-diagonal} és sempre fidel, i que és ple si i només si $\cat{C}$ és prima. \hfill{\normalfont$\star$}
\begin{indicacio}
Sobre cada homconjunt, $\Delta$ actua com $f\mapsto(f,f)$. Quines parelles $(f,g)$ de $\cat{C}\times\cat{C}$ provenen d'una sola fletxa?
\end{indicacio}
\end{exercici}

\begin{exercici}
Comprova que el functor oblit $U\colon\cat{Cat}\to\cat{Graph}$, que envia cada categoria petita al seu graf subjacent, és fidel però \emph{no} ple. Mostra també que, per a tot graf $G$, hi ha un morfisme canònic de grafs $\eta_G\colon G\to U(\Free(G))$. \hfill{\normalfont$\star$}
\begin{indicacio}
Quant a ser fidel, un functor queda determinat per la seva acció sobre objectes i morfismes, que és justament el que reté el graf subjacent. Per veure que no és ple, observa que un morfisme de grafs entre els grafs subjacents no ha de respectar la composició, i per tant no prové necessàriament d'un functor. Per al morfisme canònic $\eta_G$, envia cada vèrtex a ell mateix i cada aresta de $G$ al camí de longitud~$1$ corresponent dins $\Free(G)$.
\end{indicacio}
\end{exercici}

\begin{exercici}[Ampliació]
Caracteritza exactament per a quins grafs $G$ el morfisme canònic $\eta_G\colon G\to U(\Free(G))$ de l'exercici anterior és un isomorfisme de grafs. \hfill{\normalfont$\star\star$}
\begin{indicacio}
Compara les arestes de $U(\Free(G))$ amb les imatges de les arestes de $G$. Quins camins de $\Free(G)$ existeixen sempre, sigui quin sigui $G$, i no tenen longitud~$1$? Tingues en compte també el cas extrem.
\end{indicacio}
\end{exercici}

\begin{exercici}
Comprova que un functor ple i fidel \textbf{reflecteix els isomorfismes entre objectes}, en el sentit següent:
\[
F(A)\cong F(B)\quad\Longrightarrow\quad A\cong B.
\]
Suposant a més que $\cat{C}$ i $\cat{D}$ són petites, dedueix que $F$ indueix una aplicació injectiva entre les classes d'isomorfisme d'objectes de $\cat{C}$ i les de $\cat{D}$. \hfill{\normalfont$\star\star$}
\begin{indicacio}
Pren un isomorfisme $F(A)\to F(B)$. Com que $F$ és ple, prové d'un morfisme $A\to B$; per la propietat que els functors plens i fidels reflecteixen isomorfismes, aquest morfisme és un isomorfisme, de manera que $A\cong B$. La implicació $F(A)\cong F(B)\Rightarrow A\cong B$ val sempre; la hipòtesi de petitesa només s'usa a la darrera part, per poder parlar del conjunt de classes d'isomorfisme: en una categoria petita els objectes formen un conjunt, i el quocient per la relació $\cong$ també, de manera que \enquote{l'aplicació entre classes d'isomorfisme} té sentit literal. En una categoria gran caldria una convenció addicional per manejar aquestes classes, que aquí evitem.
\end{indicacio}
\end{exercici}

\begin{exercici}
Considera l'assignació $\cat{C}\to\cat{C}\op$ que deixa fixos els objectes i envia cada morfisme $f$ a ell mateix, vist ara en sentit contrari. Comprova que \emph{no} és, en general, un functor covariant, però sí un functor contravariant ple i fidel. Quin és aquest functor? \hfill{\normalfont$\star\star$}
\begin{indicacio}
Sobre morfismes, aquesta assignació inverteix l'ordre de la composició; és el functor identitat vist com a functor contravariant $\cat{C}\to\cat{C}\op$. La bijecció sobre cada conjunt de morfismes és evident.
\end{indicacio}
\end{exercici}

\begin{exercici}
Comprova que si $F\colon\cat{C}\to\cat{D}$ és ple i fidel, i $\cat{C}$ és localment petita, aleshores $F$ estableix una bijecció
\[
\Hom_{\cat{C}}(A,B)\;\xrightarrow{\;\sim\;}\;\Hom_{\cat{D}}(F(A),F(B))
\]
per a cada parella d'objectes. Explica per què això no obliga $F$ a ser injectiu sobre objectes. \hfill{\normalfont$\star\star$}
\begin{indicacio}
La bijecció és la definició mateixa de ple i fidel. Per a la segona part, busca una categoria amb dos objectes isomorfs i exactament un morfisme entre cada parella ordenada, i compara-la amb $\cat{1}$.
\end{indicacio}
\end{exercici}

\begin{exercici}
Distingeix l'exhaustivitat \emph{estricta} sobre objectes (tot objecte del codomini és exactament una imatge) de l'\emph{essencial} (tot objecte del codomini és isomorf a una imatge). Sigui $\cat{I}$ la categoria amb dos objectes $a,b$ i exactament un morfisme entre cada parella ordenada d'objectes, i considera el functor $E\colon\cat{1}\to\cat{I}$ que envia l'únic objecte $\ast$ a $a$. Comprova que $E$ és essencialment exhaustiu però no exhaustiu sobre objectes. \hfill{\normalfont$\star\star$}
\begin{indicacio}
Observa que $E$ arriba estrictament només a $a$, però que $b\cong a$.
\end{indicacio}
\end{exercici}

\begin{exercici}
Sigui $T$ la lògica proposicional clàssica. Comprova que la negació $A\mapsto\neg A$ defineix un functor contravariant
\[
N\colon\cat{Prov}_T\op\to\cat{Prov}_T,
\]
i que, per tant, la contraposició ---de $A\vdash_T B$ es dedueix $\neg B\vdash_T\neg A$--- és exactament la functorialitat de $N$. Comprova que $N$ és fidel i que és ple. Explica quina llei lògica cal perquè $N$ sigui ple i per què, en lògica intuïcionista, $N$ continua sent un functor però deixa de ser ple. \hfill{\normalfont$\star\star$}
\begin{indicacio}
Les dues categories són primes, de manera que només cal comprovar que les fletxes imatge existeixen. Que $N$ sigui ple demana que de $\neg B\vdash_T\neg A$ se'n segueixi $A\vdash_T B$. Per veure que pot fallar en lògica intuïcionista, busca una fórmula que no sigui estable per doble negació.
\end{indicacio}
\end{exercici}

\chapter[Transformacions naturals]{\texorpdfstring{Transformacions naturals\\ i equivalències}{Transformacions naturals i equivalències}}

Els exercicis següents completen els resolts a la part de Pràctica. Cada enunciat va acompanyat d'una indicació breu; convé intentar-lo primer sense mirar-la. La marca \enquote{$\star$} assenyala un exercici de consolidació i \enquote{$\star\star$} un d'ampliació.

\section{Definició de transformació natural}

\begin{exercici}
Siguin $F,G,H\colon\cat{C}\to\cat{D}$ tres functors i $\eta\colon F\Rightarrow G$, $\theta\colon G\Rightarrow H$ dues transformacions naturals. Comprova directament, a partir de la definició, que la composició vertical $\theta\comp\eta$ satisfà la condició de naturalitat. \hfill{\normalfont$\star$}
\begin{indicacio}
Encadena els quadrats de naturalitat de $\eta$ i de $\theta$: $H(f)\comp\theta_A\comp\eta_A=\theta_B\comp G(f)\comp\eta_A=\theta_B\comp\eta_B\comp F(f)$.
\end{indicacio}
\end{exercici}

\begin{exercici}
Sigui $\eta\colon F\Rightarrow G$ una transformació natural i $K\colon\cat{D}\to\cat{E}$ un functor. Comprova que els morfismes $K(\eta_A)$ són les components d'una transformació natural $K\eta\colon K\comp F\Rightarrow K\comp G$ (composició d'una transformació natural amb un functor per l'esquerra). \hfill{\normalfont$\star$}
\begin{indicacio}
Aplica $K$ al quadrat de naturalitat de $\eta$ i usa la functorialitat de $K$: de $G(f)\comp\eta_A=\eta_B\comp F(f)$ dedueix $K(G(f))\comp K(\eta_A)=K(\eta_B)\comp K(F(f))$.
\end{indicacio}
\end{exercici}

\begin{exercici}
Siguin $\eta\colon F\Rightarrow G$, amb $F,G\colon\cat{C}\to\cat{D}$, i $\theta\colon K\Rightarrow L$, amb $K,L\colon\cat{D}\to\cat{E}$. Comprova que la composició horitzontal amb una transformació identitat es redueix a la composició amb un functor:
\[
\theta\ast\id_F=\theta F,
\qquad\qquad
\id_K\ast\eta=K\eta.
\]
Dedueix que la composició horitzontal generalitza les dues composicions d'una transformació natural amb un functor. \hfill{\normalfont$\star$}
\begin{indicacio}
Calcula les components amb la definició, fent servir l'expressió de $(\theta\ast\eta)_A$ que més convingui a cada cas, i recorda que un functor envia identitats a identitats.
\end{indicacio}
\end{exercici}

\begin{exercici}
Siguin $\alpha\colon F\Rightarrow G$ i $\beta\colon G\Rightarrow H$, amb $F,G,H\colon\cat{C}\to\cat{D}$, i $\sigma\colon K\Rightarrow L$ i $\tau\colon L\Rightarrow M$, amb $K,L,M\colon\cat{D}\to\cat{E}$. Escriu el tipus de cada costat i demostra la \textbf{llei d'intercanvi}
\[
(\tau\comp\sigma)\ast(\beta\comp\alpha)=(\tau\ast\beta)\comp(\sigma\ast\alpha)\colon K\comp F\Rightarrow M\comp H.
\]
\hfill{\normalfont$\star\star$}
\begin{indicacio}
Calcula les components en un objecte $A$ de les dues bandes i fes servir la naturalitat de $\sigma$ aplicada al morfisme $\beta_A$.
\end{indicacio}
\end{exercici}

\section{Exemples de transformacions naturals}

\begin{exercici}
En l'exemple sintaxi/semàntica de la secció~\ref{sec:cap-transformacio}, escriu explícitament l'equació de naturalitat de la interpretació $\eta$ per a un canvi de variables $f\colon X\to Y$, i explica en paraules què diuen els dos costats. Comprova-la per a la fórmula $\varphi=p\lor\lnot q$, amb $X=\{p,q\}$, en dos casos:
\begin{enumerate}
  \item el reanomenament $f\colon p\mapsto r,\ q\mapsto s$, amb $Y=\{r,s\}$;
  \item la substitució $f\colon p\mapsto r,\ q\mapsto r$, amb $Y=\{r\}$, que identifica les dues variables.
\end{enumerate}
En el segon cas, $\varphi$ no és una tautologia, però $\varphi[f]$ sí que ho és. Explica per què això no contradiu la naturalitat. \hfill{\normalfont$\star$}
\begin{indicacio}
Avalua tots dos camins del quadrat sobre una valoració qualsevol $w\colon Y\to\{0,1\}$: l'un passa per $w\models\varphi[f]$ i l'altre per $(w\comp f)\models\varphi$. Al segon cas, observa quines valoracions de $X$ són de la forma $w\comp f$.
\end{indicacio}
\end{exercici}

\begin{exercici}
Sigui $\cat{C}$ una categoria localment petita i siguin $X,Y$ dos objectes seus. Comprova que cada morfisme $h\colon Y\to X$ indueix una transformació natural entre els homfunctors covariants
\[
\Hom(X,-)\Rightarrow\Hom(Y,-).
\]
\hfill{\normalfont$\star$}
\begin{indicacio}
Defineix la component a $A$ per precomposició amb $h$ i usa l'associativitat per comprovar la naturalitat.
\end{indicacio}
\end{exercici}

\begin{exercici}
Considera el functor \enquote{llista} $L\colon\cat{Set}\to\cat{Set}$, que envia cada conjunt $X$ al conjunt $L(X)$ de les llistes finites d'elements de $X$. Comprova que l'aplicació que envia cada element $x$ a la llista d'un sol element $[x]$ defineix una transformació natural $\id_{\cat{Set}}\Rightarrow L$. \hfill{\normalfont$\star\star$}
\begin{indicacio}
El component a $X$ és $\eta_X\colon X\to L(X)$, $x\mapsto[x]$. La naturalitat respecte d'una aplicació $f\colon X\to Y$ diu que aplicar $f$ i després formar la llista d'un element coincideix amb formar la llista i després aplicar $f$ terme a terme.
\end{indicacio}
\end{exercici}

\section{Isomorfismes naturals}

\begin{exercici}
Considera el functor d'\textbf{intercanvi} $\sigma_{\cat{C},\cat{D}}\colon\cat{C}\times\cat{D}\to\cat{D}\times\cat{C}$ definit, sobre objectes i sobre morfismes, per
\[
\sigma_{\cat{C},\cat{D}}(A,B)=(B,A),
\qquad
\sigma_{\cat{C},\cat{D}}(f,g)=(g,f).
\]
Comprova que és un functor i que la seva composició amb l'intercanvi en sentit contrari,
\[
\sigma_{\cat{D},\cat{C}}\comp\sigma_{\cat{C},\cat{D}}=\id_{\cat{C}\times\cat{D}},
\]
és \emph{estrictament} la identitat (i igualment en l'altre sentit). Dedueix que $\sigma_{\cat{C},\cat{D}}$ és un \emph{isomorfisme de categories}, no només una equivalència, i explica per què aquí no cal passar per un isomorfisme natural com $\id\cong\sigma\comp\sigma$. \hfill{\normalfont$\star$}
\begin{indicacio}
Comprova la functorialitat component a component i calcula $\sigma_{\cat{D},\cat{C}}\comp\sigma_{\cat{C},\cat{D}}$ sobre objectes i sobre morfismes. Per a l'última pregunta, compara la definició d'isomorfisme de categories amb la d'equivalència.
\end{indicacio}
\end{exercici}

\begin{exercici}
Sigui $\cat{C}$ una categoria amb \emph{almenys un objecte}. Comprova que dos functors constants $\Delta_D,\Delta_{D'}\colon\cat{C}\to\cat{D}$ (amb valors $D$ i $D'$) són naturalment isomorfs si i només si $D\cong D'$ a $\cat{D}$. Explica què falla si $\cat{C}$ és buida. \hfill{\normalfont$\star\star$}
\begin{indicacio}
Per a ($\Rightarrow$), mira què és una component en un objecte qualsevol de $\cat{C}$. Per a ($\Leftarrow$), prova amb una família constant i escriu-ne el quadrat de naturalitat. Per al cas buit, pensa quantes transformacions naturals hi ha entre dos functors amb domini buit.
\end{indicacio}
\end{exercici}

\begin{exercici}
Sigui $\mathbb{Z}/2\mathbb{Z}=\{1,g\}$ (amb $g^2=1$) vist com a categoria d'un sol objecte $\cat{B}(\mathbb{Z}/2\mathbb{Z})$, i considera dos functors $F,G\colon\cat{B}(\mathbb{Z}/2\mathbb{Z})\to\cat{Set}$ que envien l'únic objecte al conjunt $\{0,1,2\}$. En $F$, el generador $g$ actua per la identitat; en $G$, actua per l'aplicació $\sigma$ que intercanvia $0$ i $1$ i deixa fix el $2$.
\begin{enumerate}
  \item Troba totes les transformacions naturals $F\Rightarrow G$.
  \item Dedueix que $F$ i $G$ no són naturalment isomorfs, tot i que $F(\ast)=G(\ast)$.
  \item Hi ha transformacions naturals $G\Rightarrow F$? N'hi ha alguna de bijectiva?
\end{enumerate}
Compara el resultat amb l'exemple de l'observació~\ref{obs:iso-no-natural}, on no hi havia \emph{cap} transformació natural. \hfill{\normalfont$\star\star$}
\begin{indicacio}
Escriu la condició de naturalitat per al generador $g$ i estudia si pot existir una component bijectiva que la satisfaci.
\end{indicacio}
\end{exercici}

\section{La categoria de functors}

\begin{exercici}
Comprova que si $\cat{C}$ és una categoria petita i $\cat{D}$ és localment petita, aleshores la categoria de functors $\cat{D}^{\cat{C}}$ és localment petita. \hfill{\normalfont$\star$}
\begin{indicacio}
Una transformació natural està determinada pels seus components, un per a cada objecte de $\cat{C}$; com que $\cat{C}$ és petita i cada $\Hom_{\cat{D}}(F(A),G(A))$ és un conjunt, la col·lecció de transformacions naturals $F\Rightarrow G$ és un subconjunt d'un producte de conjunts indexat per un conjunt.
\end{indicacio}
\end{exercici}

\begin{exercici}
Descriu la categoria de functors $\cat{Set}^{\cat{2}}$ i comprova que els seus isomorfismes són els quadrats commutatius amb les dues fletxes verticals bijectives. \hfill{\normalfont$\star\star$}
\begin{indicacio}
Un objecte és una aplicació $f\colon A\to B$; un morfisme $(u,v)$ és un quadrat commutatiu. Un isomorfisme a la categoria de fletxes és un quadrat amb $u$ i $v$ bijectives, ja que la inversa s'obté invertint-les.
\end{indicacio}
\end{exercici}

\section{Equivalència de categories}

\begin{exercici}
Comprova que \enquote{ser equivalents} és una relació d'equivalència entre categories. És a dir, comprova que tota categoria és equivalent a ella mateixa, que la relació és simètrica i que la composició de dues equivalències és una equivalència. \hfill{\normalfont$\star$}
\begin{indicacio}
Per a la reflexivitat, pren el functor identitat. Per a la simetria, intercanvia els papers de $F$ i $G$. Per a la transitivitat, comprova que la composició de functors plens, fidels i essencialment exhaustius torna a tenir les tres propietats, o compon directament els functors i els isomorfismes naturals.
\end{indicacio}
\end{exercici}

\begin{exercici}
Sigui $\cat{C}$ una categoria en què tot objecte és isomorf a un objecte fix $D_0$ (és a dir, hi ha un sol objecte llevat d'isomorfisme; \emph{no} se suposa que tots els morfismes siguin isomorfismes). Comprova que $\cat{C}$ és equivalent al monoide $\mathrm{End}(D_0)=\Hom(D_0,D_0)$ vist com a categoria d'un sol objecte. \hfill{\normalfont$\star\star$}
\begin{indicacio}
Considera la subcategoria plena de $\cat{C}$ que té $D_0$ com a únic objecte, i fes servir la caracterització de les equivalències. Per què no cal cap hipòtesi sobre la invertibilitat dels morfismes?
\end{indicacio}
\end{exercici}

\begin{exercici}
Dona un exemple de dues categories que siguin \emph{equivalents} però no \emph{isomorfes}, i explica per què no són isomorfes. \hfill{\normalfont$\star\star$}
\begin{indicacio}
Un isomorfisme de categories indueix una bijecció entre els objectes, mentre que una equivalència no. Busca dues categories equivalents amb un nombre diferent d'objectes; n'hi ha un exemple amb categories molt petites.
\end{indicacio}
\end{exercici}

\ifpublicacioparcial\else
\chapter[Representabilitat i Yoneda]{\texorpdfstring{Propietats universals,\\ representabilitat i Yoneda}{Propietats universals, representabilitat i Yoneda}}

Els exercicis següents completen els resolts a la part de Pràctica. Cada enunciat va acompanyat d'una indicació breu. La marca \enquote{$\star$} assenyala un exercici de consolidació i \enquote{$\star\star$} un d'ampliació.

\section{Objectes inicials i terminals}

\begin{exercici}
Comprova que si una categoria té objecte inicial, aquest és únic llevat d'isomorfisme únic, repetint l'argument fet per als objectes terminals. \hfill{\normalfont$\star$}
\begin{indicacio}
Dualitza l'argument dels objectes terminals: compara les composicions de les úniques fletxes entre dos objectes inicials amb les identitats.
\end{indicacio}
\end{exercici}

\begin{exercici}
Determina els objectes inicial i terminal a $\cat{Set}$, a $\cat{Pos}$ i en un preordre $(P,\leq)$ vist com a categoria. \hfill{\normalfont$\star$}
\begin{indicacio}
Per a cada categoria, pregunta't quins objectes reben exactament una fletxa des de cada objecte, i quins n'envien exactament una cap a cada objecte. En un preordre, tradueix-ho en termes de l'ordre.
\end{indicacio}
\end{exercici}

\begin{exercici}[Ampliació, per a lectors amb àlgebra]
Comprova que la categoria $\cat{Field}$ dels cossos i homomorfismes de cossos \emph{no} té objecte inicial ni terminal. \hfill{\normalfont$\star$}
\begin{indicacio}
Un homomorfisme de cossos preserva $0$, $1$ i les operacions, i és sempre injectiu. No hi ha cap cos que rebi un morfisme de tots els altres (pensa en característiques diferents), ni cap que n'enviï a tots.
\end{indicacio}
\end{exercici}

\section{Propietats universals com a representabilitat}

\begin{exercici}
Comprova que un objecte $0$ és inicial si i només si el functor covariant $\Hom(0,-)$ és naturalment isomorf al functor constant $\Delta_{\{\ast\}}\colon\cat{C}\to\cat{Set}$ (exemple~\ref{ex:functor-constant}). \hfill{\normalfont$\star$}
\begin{indicacio}
Compara el nombre d'elements de cada $\Hom(0,X)$ amb el de $\{\ast\}$, i pregunta't quines condicions de naturalitat queden per comprovar.
\end{indicacio}
\end{exercici}

\begin{exercici}
Sigui $\cat{C}$ localment petita amb objecte terminal $1$. Comprova que $\Hom(1,-)\colon\cat{C}\to\cat{Set}$ és un functor i demostra que envia l'objecte terminal a un conjunt terminal. Mostra amb un exemple que, tot i això, $\Hom(1,-)$ pot no distingir objectes no isomorfs: pot enviar-los a conjunts amb el mateix nombre d'elements, o fins i tot al mateix conjunt. \hfill{\normalfont$\star$}
\begin{indicacio}
Calcula $\Hom(1,1)$. Per a l'exemple, pensa en un preordre amb un element màxim i dos elements incomparables per sota.
\end{indicacio}
\end{exercici}

\begin{exercici}
Considera el prefeix de parts $\mathcal P\colon\cat{Set}\op\to\cat{Set}$ i les dues parelles candidates $(\{0,1\},\{1\})$ i $(\{0,1\},\varnothing)$. Quina de les dues és una representació de $\mathcal P$? Justifica per què l'altra no ho és. \hfill{\normalfont$\star$}
\begin{indicacio}
Per a la segona parella, intenta obtenir el subconjunt $\{\ast\}$ d'un conjunt d'un sol element com una antiimatge de $\varnothing$.
\end{indicacio}
\end{exercici}

\section{El lema de Yoneda}

\begin{exercici}
Escriu l'enunciat covariant del lema de Yoneda: per a un functor covariant $Q\colon\cat{C}\to\cat{Set}$ i un objecte $A$, hi ha una bijecció natural
\[
\Nat\bigl(\Hom(A,-),Q\bigr)\;\cong\;Q(A).
\]
Identifica la bijecció. \hfill{\normalfont$\star\star$}
\begin{indicacio}
Reprodueix la demostració del cas contravariant a $\cat{C}\op$ i avalua la transformació en $\id_A$.
\end{indicacio}
\end{exercici}

\begin{exercici}
Fes servir el lema de Yoneda per calcular
\[
\Nat\bigl(\Hom(-,A),\Hom(-,B)\bigr).
\]
Comprova que la bijecció resultant és precisament l'acció de la immersió de Yoneda sobre els morfismes $A\to B$. \hfill{\normalfont$\star\star$}
\begin{indicacio}
Aplica Yoneda al prefeix representable $P=\Hom(-,B)$ i identifica què correspon a un morfisme $A\to B$.
\end{indicacio}
\end{exercici}

\begin{exercici}
Siguin $(A,u)$ i $(B,v)$ dues representacions d'un mateix prefeix $P$. Demostra que hi ha un únic isomorfisme $f\colon A\to B$ amb $P(f)(v)=u$. Explica per què això no diu que $A$ tingui un únic automorfisme. \hfill{\normalfont$\star$}
\begin{indicacio}
Fes servir la universalitat de $v$ per obtenir $f$, i la de $u$ per obtenir $g\colon B\to A$ amb $P(g)(u)=v$. Aplica després la unicitat a $g\comp f$ i a $\id_A$.
\end{indicacio}
\end{exercici}

\begin{exercici}
Comprova que si un prefeix $P$ és representable, aleshores l'element $x\in P(A)$ que correspon a la identitat sota l'isomorfisme $\Hom(-,A)\cong P$ té la propietat universal següent: per a cada objecte $X$ i cada $y\in P(X)$, hi ha un únic morfisme $\varphi\colon X\to A$ amb $P(\varphi)(x)=y$. (Aquest $x$ s'anomena \emph{element universal}.) \hfill{\normalfont$\star$}
\begin{indicacio}
Escriu la naturalitat de l'isomorfisme $\Hom(-,A)\cong P$ per a un morfisme $\varphi\colon X\to A$, avaluada en $\id_A$.
\end{indicacio}
\end{exercici}

\section{Conseqüències}

\begin{exercici}
Sigui $\cat{C}$ una categoria petita. Comprova que la immersió de Yoneda $\yon\colon\cat{C}\to\cat{Set}^{\cat{C}\op}$ \emph{no} és, en general, essencialment exhaustiva, tot i ser plena i fidel. Quins prefeixos són de la forma $\Hom(-,A)$? \hfill{\normalfont$\star\star$}
\begin{indicacio}
Busca una propietat que tot prefeix representable $\Hom(-,A)$ compleix necessàriament en l'objecte $A$, i un prefeix que no la compleixi.
\end{indicacio}
\end{exercici}

\begin{exercici}
Dedueix del corol·lari de Yoneda que si dos objectes tenen conjunts de morfismes \enquote{naturalment iguals} des de tot objecte, aleshores són isomorfs; i explica per què això formalitza l'eslògan \enquote{un objecte es coneix per les fletxes que hi arriben}. \hfill{\normalfont$\star\star$}
\begin{indicacio}
Tradueix \enquote{naturalment iguals des de tot objecte} en termes de prefeixos representables.
\end{indicacio}
\end{exercici}

\begin{exercici}
Enuncia i comprova la versió dual del corol·lari: $A\cong B$ si i només si $\Hom(A,-)\cong\Hom(B,-)$ com a functors covariants. \hfill{\normalfont$\star\star$}
\begin{indicacio}
Aplica el corol·lari a $\cat{C}\op$, o repeteix l'argument amb la immersió covariant $\cat{C}\op\to\cat{Set}^{\cat{C}}$, que també és plena i fidel.
\end{indicacio}
\end{exercici}

\chapter[Límits i colímits]{Límits i colímits}

Els exercicis següents completen els resolts a la part de Pràctica. Cada enunciat va acompanyat d'una indicació breu. La marca \enquote{$\star$} assenyala un exercici de consolidació i \enquote{$\star\star$} un d'ampliació.

\section{Diagrames i cons}

\begin{exercici}
Descriu la categoria índex $\cat{I}$ i el diagrama corresponent per a cadascun d'aquests límits: producte binari, equalitzador, pullback i objecte terminal. \hfill{\normalfont$\star$}
\begin{indicacio}
Identifica, per a cada propietat universal, la forma mínima de diagrama que la descriu: quants objectes hi intervenen i quines fletxes cal que commutin.
\end{indicacio}
\end{exercici}

\begin{exercici}
Comprova que un morfisme de cons és un cas particular de morfisme a la categoria $\mathrm{Con}(D)$, i que la composició de morfismes de cons torna a ser-ne un. \hfill{\normalfont$\star$}
\begin{indicacio}
Si $f\colon A\to B$ i $g\colon B\to C$ satisfan $\beta_i\comp f=\alpha_i$ i $\gamma_i\comp g=\beta_i$, aleshores $\gamma_i\comp(g\comp f)=\alpha_i$.
\end{indicacio}
\end{exercici}

\section{Límits}

\begin{exercici}
Comprova que en un preordre vist com a categoria, el producte de dos elements és el seu ínfim i el coproducte és el seu suprem, quan existeixen. \hfill{\normalfont$\star$}
\begin{indicacio}
Tradueix la propietat universal del producte a desigualtats.
\end{indicacio}
\end{exercici}

\begin{exercici}
Comprova que dos límits d'un mateix diagrama són isomorfs per un únic isomorfisme compatible amb totes les projeccions. \hfill{\normalfont$\star$}
\begin{indicacio}
Treballa a la categoria $\mathrm{Con}(D)$ i fes servir la unicitat dels objectes terminals.
\end{indicacio}
\end{exercici}

\section{Productes i equalitzadors}

\begin{exercici}
Comprova, llevat d'isomorfisme, la commutativitat i l'associativitat del producte:
\[
A\times B\cong B\times A,\quad (A\times B)\times C\cong A\times(B\times C).
\] \hfill{\normalfont$\star$}
\begin{indicacio}
Fes servir la unicitat d'un objecte definit per una mateixa propietat universal.
\end{indicacio}
\end{exercici}

\begin{exercici}
Comprova que si $\cat{C}$ té objecte terminal $1$, aleshores $A\times 1\cong A$ per a tot objecte $A$. \hfill{\normalfont$\star$}
\begin{indicacio}
Compara les propietats universals de $A\times 1$ i de $A$.
\end{indicacio}
\end{exercici}

\begin{exercici}
Comprova que a $\cat{Grp}$ l'equalitzador de $f,g\colon G\to H$ és el subgrup $\{x\in G\mid f(x)=g(x)\}$. \hfill{\normalfont$\star$}
\begin{indicacio}
Comprova que aquest subconjunt és un subgrup i que la injecció satisfà la propietat universal de l'equalitzador; el functor oblit cap a $\cat{Set}$ ajuda a veure-ho.
\end{indicacio}
\end{exercici}

\section{Pullbacks}

\begin{exercici}
Comprova que en un pullback, si $g\colon B\to C$ és un monomorfisme, aleshores $p_1\colon P\to A$ també ho és. (Es diu que els monomorfismes són \emph{estables per pullback}.) \hfill{\normalfont$\star\star$}
\begin{indicacio}
Usa la propietat universal del pullback i la cancel·lació de $g$. Aquesta estabilitat serà essencial per a la teoria de subobjectes.
\end{indicacio}
\end{exercici}

\begin{exercici}
Comprova que un quadrat amb $\id$ als dos costats,
\[
\begin{array}{ccc}
A & \xrightarrow{\;\id\;} & A\\[-2pt]
\downarrow{\scriptstyle\id} & & \downarrow{\scriptstyle f}\\[-2pt]
A & \xrightarrow{\;f\;} & B
\end{array}
\]
és un pullback si i només si $f$ és un monomorfisme. \hfill{\normalfont$\star\star$}
\begin{indicacio}
El pullback d'aquest quadrat a $\cat{Set}$ és $\{(a,a')\mid f(a)=f(a')\}$; que la projecció sigui iso equival a la injectivitat de $f$. En general, tradueix-ho amb la propietat universal.
\end{indicacio}
\end{exercici}

\section{Colímits}

\begin{exercici}
Comprova que tot coequalitzador és un epimorfisme, dualitzant la demostració que tot equalitzador és un monomorfisme. \hfill{\normalfont$\star$}
\begin{indicacio}
Dualitza l'argument de la Pràctica que prova que tot equalitzador és mono.
\end{indicacio}
\end{exercici}

\begin{exercici}
Comprova que el pushout de $f\colon C\to A$ i $g\colon C\to B$ a $\cat{Set}$ és $(A+B)/{\sim}$, on $\sim$ enganxa $f(c)$ amb $g(c)$ per a cada $c\in C$. \hfill{\normalfont$\star$}
\begin{indicacio}
És el dual del pullback. Pren la unió disjunta $A+B$ i quocienta per la relació generada per $\iota_1(f(c))\sim\iota_2(g(c))$; comprova la propietat universal.
\end{indicacio}
\end{exercici}

\section{Límits finits}

\begin{exercici}
Sigui $\cat{C}$ una categoria amb productes binaris i pullbacks, i siguin $f,g\colon A\to B$. La demostració del teorema~\ref{teo:limits-finits} obté l'equalitzador de $f$ i $g$ com un pullback de dues fletxes $A\to A\times B$. Comprova que també és un pullback de la diagonal $\langle\id_B,\id_B\rangle\colon B\to B\times B$ al llarg de $\langle f,g\rangle\colon A\to B\times B$. \hfill{\normalfont$\star\star$}
\begin{indicacio}
Escriu què és un con sobre aquest pullback i tradueix la condició de con component a component, amb les dues projeccions de $B\times B$.
\end{indicacio}
\end{exercici}

\begin{exercici}
Comprova directament que una categoria amb productes binaris i equalitzadors té tots els pullbacks, amb una construcció més econòmica que la general de la demostració del teorema~\ref{teo:limits-finits}. \hfill{\normalfont$\star$}
\begin{indicacio}
Expressa el pullback com un equalitzador de dues fletxes que surten d'un producte.
\end{indicacio}
\end{exercici}

\section{Límits en categories de functors}

\begin{exercici}
Sigui $\mathbf E$ el graf amb dos vèrtexs $0,1$ i una sola aresta $e\colon 0\to 1$. Fent servir que $\cat{Graph}\cong\cat{Set}^{\cat{P}}$ i que els límits i colímits s'hi calculen punt a punt (proposició~\ref{prop:limits-punt-a-punt}), descriu el producte $\mathbf E\times\mathbf E$ i el coproducte $\mathbf E+\mathbf E$: els seus vèrtexs, les seves arestes, i l'origen i el destí de cada aresta. Dibuixa'ls. És $\mathbf E\times\mathbf E$ un quadrat? \hfill{\normalfont$\star$}
\begin{indicacio}
Calcula per separat el conjunt de vèrtexs i el d'arestes, i defineix origen i destí component a component.
\end{indicacio}
\end{exercici}

\section{Preservació i reflexió}

\begin{exercici}
Comprova que tot functor ple i fidel reflecteix límits. Dualment, reflecteix colímits. \hfill{\normalfont$\star\star$}
\begin{indicacio}
Aixeca l'existència de la factorització amb la plenitud i prova'n la unicitat amb la fidelitat.
\end{indicacio}
\end{exercici}

\begin{exercici}
Comprova que el prefeix representable
\[
\Hom(-,B)\colon\cat{C}\op\to\cat{Set}
\]
envia colímits de $\cat{C}$ a límits de $\cat{Set}$. \hfill{\normalfont$\star\star$}
\begin{indicacio}
És el dual del fet que $\Hom(A,-)$ preserva límits: un cocon amb vèrtex $B$ sobre $D$ és un con sobre $\Hom(D_-,B)$, i la universalitat es correspon.
\end{indicacio}
\end{exercici}

\chapter[Adjuncions]{Adjuncions}

Els exercicis següents completen els resolts a la part de Pràctica. Cada enunciat va acompanyat d'una indicació breu. La marca \enquote{$\star$} assenyala un exercici de consolidació i \enquote{$\star\star$} un d'ampliació.

\section{Adjuncions entre preordres}

\begin{exercici}
Comprova que en una adjunció entre preordres $f\dashv g$, l'aplicació $f$ preserva els suprems que existeixin i $g$ preserva els ínfims. \hfill{\normalfont$\star$}
\begin{indicacio}
És el cas particular, per a preordres, que els adjunts per l'esquerra preserven colímits (suprems) i els adjunts per la dreta preserven límits (ínfims). Demostra-ho directament amb la definició $f(x)\leq y\Leftrightarrow x\leq g(y)$.
\end{indicacio}
\end{exercici}

\begin{exercici}
Sigui $f\colon P\to Q$ monòtona entre \textbf{reticles complets}, és a dir, preordres en què tot subconjunt té suprem i ínfim. Comprova que $f$ té adjunt per la dreta si i només si preserva tots els suprems. \hfill{\normalfont$\star\star$}
\begin{indicacio}
Si $f$ preserva suprems, defineix $g(y)$ com el suprem d'un conjunt adequat d'elements de $P$.
\end{indicacio}
\end{exercici}

\section{Adjuncions entre categories}

\begin{exercici}\label{exr:adj-unitat-universal}
Comprova que si $F\dashv G$, aleshores per a cada objecte $A$ el morfisme $\eta_A\colon A\to G(F(A))$ és universal: tot $\psi\colon A\to G(B)$ factoritza de manera única com $G(\varphi)\comp\eta_A$. \hfill{\normalfont$\star\star$}
\begin{indicacio}
Fes servir la bijecció $\theta$ i la fórmula del transposat en termes de la unitat.
\end{indicacio}
\end{exercici}

\begin{exercici}
Sigui $\cat{C}$ una categoria amb coproductes binaris i $\Delta\colon\cat{C}\to\cat{C}\times\cat{C}$ el functor diagonal, $A\mapsto(A,A)$. Dualitzant l'exercici resolt de la Pràctica sobre $\Delta\dashv\times$, comprova que el coproducte és adjunt per l'esquerra de $\Delta$, $(+)\dashv\Delta$, i descriu-ne la unitat i la counitat. \hfill{\normalfont$\star\star$}
\begin{indicacio}
Fes servir la propietat universal del coproducte; la unitat i la counitat són les duals de la diagonal i les projeccions.
\end{indicacio}
\end{exercici}

\section{Unitat i counitat}

\begin{exercici}\label{exr:adj-fletxes-universals}
Sigui $G\colon\cat{D}\to\cat{C}$ un functor. Suposem que, per a cada objecte $A$ de $\cat{C}$, hem triat un objecte $F(A)$ de $\cat{D}$ i una fletxa universal $\eta_A\colon A\to G(F(A))$, en el sentit que tot $\psi\colon A\to G(B)$ factoritza de manera única com $G(\varphi)\comp\eta_A$. Demostra que aquestes dades determinen de manera única l'acció de $F$ sobre morfismes i defineixen un adjunt per l'esquerra $F\dashv G$. \hfill{\normalfont$\star\star$}
\begin{indicacio}
Defineix $F$ sobre objectes i sobre morfismes amb la propietat universal, i fes servir la unicitat per obtenir la functorialitat.
\end{indicacio}
\end{exercici}

\begin{exercici}
Per a l'adjunció $\Free\dashv U$ entre $\cat{Graph}$ i $\cat{Cat}$ (exemple~\ref{ex:free-u-adjuncio}), comprova que la counitat
\[
\varepsilon_{\cat{C}}\colon\Free(U(\cat{C}))\longrightarrow\cat{C}
\]
és la identitat sobre els objectes i és un functor ple. Mostra amb un exemple que, en general, no és fidel. \hfill{\normalfont$\star$}
\begin{indicacio}
Per veure que és ple, pensa en els camins de longitud~$1$. Per veure que no és fidel, busca dos camins diferents amb la mateixa composta.
\end{indicacio}
\end{exercici}

\section{Identitats triangulars}

\begin{exercici}
Comprova la segona identitat triangular,
\[
G(\varepsilon_B)\comp\eta_{G(B)}=\id_{G(B)},
\]
dualitzant la demostració de la primera. \hfill{\normalfont$\star$}
\begin{indicacio}
Dualitza la demostració de la Teoria, fent servir la fórmula del transposat en termes de la unitat.
\end{indicacio}
\end{exercici}

\begin{exercici}\label{exr:adj-triangulars-reciproc}
Comprova que dues transformacions naturals $\eta\colon\id\Rightarrow GF$ i $\varepsilon\colon FG\Rightarrow\id$ que satisfan les identitats triangulars defineixen una adjunció $F\dashv G$. \hfill{\normalfont$\star\star$}
\begin{indicacio}
Defineix els dos transposats amb $\eta$ i $\varepsilon$, i fes servir les identitats triangulars per veure que són inversos.
\end{indicacio}
\end{exercici}

\begin{exercici}
Siguin $F\colon\cat{C}\to\cat{D}$, $G\colon\cat{D}\to\cat{C}$, $F'\colon\cat{D}\to\cat{E}$ i $G'\colon\cat{E}\to\cat{D}$ functors entre categories localment petites, amb $F\dashv G$ i $F'\dashv G'$. Demostra que les adjuncions es componen:
\[
F'\comp F\;\dashv\;G\comp G'.
\]
Descriu la unitat de l'adjunció composta en termes de les unitats $\eta$ i $\eta'$. \hfill{\normalfont$\star$}
\begin{indicacio}
Encadena les dues bijeccions d'homconjunts, passant per $\cat{D}$, i comprova que la composta és natural en les dues variables. Per a la unitat, transposa la identitat de $F'(F(A))$.
\end{indicacio}
\end{exercici}

\begin{exercici}
Sigui $F\dashv G$ una adjunció amb unitat $\eta$ i counitat $\varepsilon$. Demostra que, si $\eta$ i $\varepsilon$ són isomorfismes naturals, aleshores $F$ és una equivalència de categories amb quasiinvers $G$ (definició~\ref{def:equivalencia}), i que en aquest cas $G$ també és adjunt per l'esquerra de $F$. \hfill{\normalfont$\star$}
\begin{indicacio}
Compara les dades amb la definició d'equivalència. Per a la segona part, inverteix $\eta$ i $\varepsilon$ i comprova les identitats triangulars de l'adjunció $G\dashv F$ a partir de les de $F\dashv G$.
\end{indicacio}
\end{exercici}

\section{Els adjunts i els límits}

\begin{exercici}
Comprova que el prefeix de parts $\mathcal P\colon\cat{Set}\op\to\cat{Set}$ envia coproductes a productes: $\mathcal P(A+B)\cong\mathcal P(A)\times\mathcal P(B)$. \hfill{\normalfont$\star$}
\begin{indicacio}
Descriu un subconjunt de $A+B$ a partir de les seves interseccions amb les dues components.
\end{indicacio}
\end{exercici}

\begin{exercici}[Ampliació, per a lectors amb topologia]
Sigui $U\colon\cat{Top}\to\cat{Set}$ el functor oblit. Comprova que té \emph{dos} adjunts, un per l'esquerra i un per la dreta, i digues quins són. \hfill{\normalfont$\star$}
\begin{indicacio}
Pensa en les dues topologies extremes sobre un conjunt.
\end{indicacio}
\end{exercici}

\section{Exponencials}

\begin{exercici}
Comprova que a $\cat{Set}$ es compleix $C^{A+B}\cong C^A\times C^B$ i $(C\times D)^A\cong C^A\times D^A$, i interpreta-ho en termes d'adjuncions. \hfill{\normalfont$\star$}
\begin{indicacio}
Per a la primera, fes servir que $\Hom(A+B,C)\cong\Hom(A,C)\times\Hom(B,C)$, és a dir, que el representable contravariant $\Hom(-,C)$ envia coproductes a productes. Per a la segona, fes servir que $(-)^A$ és adjunt per la dreta de $-\times A$ i, per tant, preserva productes.
\end{indicacio}
\end{exercici}

\begin{exercici}
Comprova que la currificació $\Hom(A\times B,C)\cong\Hom(A,C^B)$ és natural en les tres variables $A,B,C$. \hfill{\normalfont$\star\star$}
\begin{indicacio}
Comprova la commutativitat amb la precomposició en $A$, la postcomposició en $C$ i l'acció en $B$, aplicant les definicions de $\tilde f$ i $\hat g$.
\end{indicacio}
\end{exercici}

\chapter[Cartesianes tancades]{\texorpdfstring{Categories cartesianes\\ tancades}{Categories cartesianes tancades}}

Els exercicis següents completen els resolts a la part de Pràctica. Cada enunciat va acompanyat d'una indicació breu. La marca \enquote{$\star$} assenyala un exercici de consolidació i \enquote{$\star\star$} un d'ampliació, sensiblement més exigent.

\section{Productes finits}

\begin{exercici}
Comprova que en una categoria cartesiana el producte és associatiu i commutatiu llevat d'isomorfisme: $(A\times B)\times C\cong A\times(B\times C)$ i $A\times B\cong B\times A$. \hfill{\normalfont$\star$}
\begin{indicacio}
Construeix els morfismes en tots dos sentits amb aparellaments de projeccions i comprova que en són inversos usant la unicitat de la propietat universal. No cal calcular amb elements.
\end{indicacio}
\end{exercici}

\begin{exercici}
Demostra que si $\cat{C}$ és cartesiana, aleshores $\cat{C}\op$ té coproductes finits, i descriu-los. \hfill{\normalfont$\star$}
\begin{indicacio}
El producte a $\cat{C}$ és el coproducte a $\cat{C}\op$ per dualització de la propietat universal; el terminal esdevé inicial.
\end{indicacio}
\end{exercici}

\section{Exponencials}

\begin{exercici}
En una CCC, construeix un morfisme canònic \enquote{composició} $C^{B}\times B^{A}\to C^{A}$ i comprova que és associatiu en el sentit adient. \hfill{\normalfont$\star\star$}
\begin{indicacio}
Currifica el morfisme $C^{B}\times B^{A}\times A\to C$ obtingut avaluant primer sobre $A$ i després sobre $B$. L'associativitat es comprova amb la unicitat de la transposada.
\end{indicacio}
\end{exercici}

\begin{exercici}
Comprova que l'assignació $B\mapsto C^{B}$ és \emph{contravariant}: un morfisme $f\colon B'\to B$ indueix $C^{f}\colon C^{B}\to C^{B'}$. \hfill{\normalfont$\star$}
\begin{indicacio}
Currifica
\[
C^{B}\times B'\xrightarrow{\id\times f}C^{B}\times B\xrightarrow{\ev}C.
\]
Això mostra la contravariància en l'exponent $B$; en canvi, per a $B$ fix, el functor $(-)^{B}$ és covariant en $C$.
\end{indicacio}
\end{exercici}

\begin{exercici}
En una CCC amb coproductes, demostra la llei $C^{A+B}\cong C^{A}\times C^{B}$. \hfill{\normalfont$\star\star$}
\begin{indicacio}
Fes servir que $X\times-$ és un adjunt per l'esquerra i, per tant, preserva coproductes (exercici resolt~\ref{pr-distributivitat} de la Pràctica). Després aplica la propietat universal del coproducte i el lema de Yoneda.
\end{indicacio}
\end{exercici}

\begin{exercici}
Sabem que $\cat{Graph}$ és cartesiana tancada, perquè és una categoria de prefeixos. Fent servir la fórmula de l'exponencial de prefeixos i el lema de Yoneda, demostra que els \emph{vèrtexs} de l'exponencial $H^{G}$ de dos grafs corresponen a totes les aplicacions $V_G\to V_H$ entre els conjunts de vèrtexs, sense cap condició sobre les arestes. \hfill{\normalfont$\star\star$}
\begin{indicacio}
Els vèrtexs de $H^{G}$ es corresponen canònicament amb els morfismes de grafs $\mathbf V\to H^{G}$, on $\mathbf V$ és el graf d'un sol vèrtex. Transposa'ls i calcula el graf $\mathbf V\times G$.
\end{indicacio}
\end{exercici}

\begin{exercici}
Sigui $\cat{2}=(0\to1)$. Fent servir l'exemple de $\cat{Cat}$ de la Teoria amb $\cat{A}=\cat{2}$, comprova que un functor $\cat{2}\times\cat{2}\to\cat{C}$ és el mateix que un quadrat commutatiu de $\cat{C}$, i que la seva currificació $\cat{2}\to\cat{C}^{\cat{2}}$ és el mateix que un morfisme de la categoria de fletxes. Què fa l'avaluació sobre aquestes dades? \hfill{\normalfont$\star$}
\begin{indicacio}
Escriu els quatre objectes i les cinc fletxes no identitat de $\cat{2}\times\cat{2}$. Aplica la lectura \enquote{transformació natural $=$ functor $\cat{A}\times\cat{2}\to\cat{C}$} amb $\cat{A}=\cat{2}$.
\end{indicacio}
\end{exercici}

\section{Models i lògica}

\begin{exercici}
Comprova que tota àlgebra de Boole és una àlgebra de Heyting en què $b\Rightarrow c=\neg b\vee c$, i per tant un preordre cartesià tancat. \hfill{\normalfont$\star$}
\begin{indicacio}
Verifica l'adjunció $a\wedge b\leq c\Leftrightarrow a\leq\neg b\vee c$ usant les lleis booleanes. La lògica clàssica és el cas booleà de la intuïcionista.
\end{indicacio}
\end{exercici}

\begin{exercici}
Sigui $\cat{C}$ una CCC. Demostra que el conjunt de morfismes $1\to C^{B}$ està en bijecció amb el de morfismes $B\to C$ (els \enquote{punts} de l'exponencial són les funcions). \hfill{\normalfont$\star$}
\begin{indicacio}
És l'adjunció amb $A=1$ i la identificació $1\times B\cong B$.
\end{indicacio}
\end{exercici}

\begin{exercici}
Interpreta el combinador $\mathsf{S}$, de tipus $(\sigma\to\tau\to\rho)\to(\sigma\to\tau)\to\sigma\to\rho$, com a morfisme d'una CCC. \hfill{\normalfont$\star\star$}
\begin{indicacio}
Descurrifica successivament fins a un morfisme amb domini un producte de tres exponencials i $\llbracket\sigma\rrbracket$, i construeix-lo amb avaluacions i la diagonal. Compara'l amb la demostració intuïcionista de l'esquema d'axioma corresponent.
\end{indicacio}
\end{exercici}

\begin{exercici}
Al preordre cartesià tancat $\mathbb Z\cup\{+\infty\}$ de la Teoria, comprova el \emph{modus ponens} $b\wedge(b\Rightarrow c)\leq c$ i la llei de currificació $b\Rightarrow(c\Rightarrow d)=(b\wedge c)\Rightarrow d$, calculant els dos costats per casos. \hfill{\normalfont$\star$}
\begin{indicacio}
Fes servir la fórmula explícita de la implicació en una cadena, i distingeix els casos segons l'ordre relatiu de $b$, $c$ i $d$.
\end{indicacio}
\end{exercici}

\begin{exercici}
Afegeix a $\mathbb Z\cup\{+\infty\}$ un element mínim $-\infty$. Comprova que s'obté una àlgebra de Heyting, calcula la pseudocomplementació $\neg b=(b\Rightarrow-\infty)$ i decideix si és una àlgebra de Boole. \hfill{\normalfont$\star\star$}
\begin{indicacio}
En una cadena, els suprems finits són màxims. Per a la segona part, calcula $\neg\neg b$ i $b\vee\neg b$ per a un $b$ que no sigui un dels extrems.
\end{indicacio}
\end{exercici}

\chapter[Subobjectes, imatges i factoritzacions]{\texorpdfstring{Subobjectes, imatges\\ i factoritzacions}{Subobjectes, imatges i factoritzacions}}

Els exercicis següents completen els resolts a la part de Pràctica. Cada enunciat va acompanyat d'una indicació breu. La marca \enquote{$\star$} assenyala un exercici de consolidació i \enquote{$\star\star$} un d'ampliació.

\section{Subobjectes}

\begin{exercici}
Descriu $\Sub(A)$ quan $A$ és un objecte d'un preordre vist com a categoria. \hfill{\normalfont$\star$}
\begin{indicacio}
En un preordre hi ha com a molt un morfisme entre dos objectes, i tots són monos. Quins subobjectes té, doncs, $A$? Compara amb el conjunt d'elements $x\leq A$: quan donen dos d'aquests elements el mateix subobjecte? Què passa si el preordre és un ordre parcial?
\end{indicacio}
\end{exercici}

\begin{exercici}
Comprova que en una categoria amb objecte inicial estricte $0$, la injecció $0\rightarrowtail A$ és el subobjecte mínim de $\Sub(A)$. \hfill{\normalfont$\star$}
\begin{indicacio}
\enquote{Estricte} vol dir que tot morfisme cap a $0$ és un isomorfisme. Comprova primer que $0\to A$ és mono, i després que es factoritza a través de qualsevol subobjecte.
\end{indicacio}
\end{exercici}

\begin{exercici}[Ampliació, per a lectors amb àlgebra]
Demostra que a $\cat{Grp}$ els subobjectes d'un grup $G$ es corresponen exactament amb els seus \textbf{subgrups}. \hfill{\normalfont$\star\star$}
\begin{indicacio}
Els monomorfismes a $\cat{Grp}$ són els homomorfismes injectius. Comprova que dos monos amb codomini $G$ són equivalents si i només si tenen la mateixa imatge, i que les imatges de monos són exactament els subgrups.
\end{indicacio}
\end{exercici}

\begin{exercici}
Comprova que un morfisme de grafs $F\colon G'\to G$ és un monomorfisme de $\cat{Graph}$ si i només si les aplicacions $F_V$ i $F_E$ són injectives. Dedueix la descripció dels subobjectes d'un graf com a subgrafs de l'exemple~\ref{ex:subgrafs}. \hfill{\normalfont$\star$}
\begin{indicacio}
Per a una direcció, prova els morfismes que surten del graf d'un sol vèrtex i del graf d'una sola aresta.
\end{indicacio}
\end{exercici}

\section{Reindexació}

\begin{exercici}
Sigui $f\colon A\to B$. Comprova que $f^{*}$ envia el subobjecte màxim de $\Sub(B)$ al màxim de $\Sub(A)$. \hfill{\normalfont$\star$}
\begin{indicacio}
El pullback de $\id_B$ al llarg de $f$ és un mono equivalent a $\id_A$. Interpreta-ho: substituir dins del predicat \enquote{sempre cert} torna a donar \enquote{sempre cert}.
\end{indicacio}
\end{exercici}

\begin{exercici}
Treballa a $\cat{Set}$. Comprova que si $m\colon S\rightarrowtail B$ és un subobjecte (un subconjunt $S\subseteq B$) i $f\colon A\to B$, aleshores $f^{*}[m]$ és el subobjecte \emph{mínim} (buit) de $A$ exactament quan la imatge de $f$ no talla $S$, és a dir, quan $f(A)\cap S=\varnothing$. \hfill{\normalfont$\star$}
\begin{indicacio}
A $\cat{Set}$, $f^{-1}(S)=\varnothing$ si i només si cap element de la imatge de $f$ és a $S$. Formula-ho amb pullbacks: el pullback de $m$ al llarg de $f$ és l'objecte inicial exactament en aquest cas.
\end{indicacio}
\end{exercici}

\section{Interseccions i estructura}

\begin{exercici}
Demostra que la reindexació $f^{*}$ preserva interseccions: $f^{*}([m]\wedge[n])=f^{*}[m]\wedge f^{*}[n]$. \hfill{\normalfont$\star\star$}
\begin{indicacio}
Tant l'ínfim com la reindexació són pullbacks. Usa que els pullbacks es componen: el pullback d'un pullback sobre $B$ al llarg de $f$ es reorganitza en la intersecció dels pullbacks. És un exercici de manipulació de quadrats cartesians.
\end{indicacio}
\end{exercici}

\begin{exercici}
A $\cat{Set}$, el functor de subobjectes $\Sub\colon\cat{Set}\op\to\cat{Poset}$ (proposició~\ref{prop:functor-sub}) i el functor de parts contravariant $\mathcal P$, amb $\mathcal P(A)$ ordenat per inclusió, envien cada conjunt a ordres parcials isomorfs. Comprova que aquests isomorfismes formen un isomorfisme natural $\Sub\cong\mathcal P$. \hfill{\normalfont$\star$}
\begin{indicacio}
Fes servir la imatge d'un monomorfisme per definir les components, i compara la reindexació amb l'antiimatge.
\end{indicacio}
\end{exercici}


\section{Factorització imatge}

\begin{exercici}[Ampliació, per a lectors amb àlgebra]
Comprova que a $\cat{Grp}$ la factorització imatge d'un homomorfisme $f\colon G\to H$ és $G\twoheadrightarrow \Imatge(f)\hookrightarrow H$, on $\Imatge(f)$ és el subgrup imatge. \hfill{\normalfont$\star$}
\begin{indicacio}
Els epimorfismes regulars a $\cat{Grp}$ són les projeccions al quocient $G\to G/N$; la primera part de $f$ és $G\to G/\ker f\cong\Imatge(f)$. Fes servir el primer teorema d'isomorfia.
\end{indicacio}
\end{exercici}

\begin{exercici}
Demostra que en qualsevol categoria amb factoritzacions imatge, la imatge d'una composició satisfà $\Imatge(g\comp f)\leq\Imatge(g)$. \hfill{\normalfont$\star$}
\begin{indicacio}
$g\comp f$ es factoritza a través de $\Imatge(g)$ (via la factorització de $g$), i la imatge és el menor subobjecte pel qual això passa.
\end{indicacio}
\end{exercici}

\begin{exercici}
Comprova que un epimorfisme regular que és alhora monomorfisme és un isomorfisme, i dedueix que en una categoria regular la factorització imatge d'un mono és trivial. \hfill{\normalfont$\star$}
\begin{indicacio}
Un epi regular és el coequalitzador de la seva parella nucli; si a més és mono, la parella nucli és la diagonal, i el coequalitzador d'una parella trivial és una identitat llevat d'isomorfisme.
\end{indicacio}
\end{exercici}

\begin{exercici}
Siguin $H$ i $K$ dos subgrafs d'un graf $G$, i $F\colon G\to G''$ un morfisme de grafs. Descriu la intersecció $H\wedge K$ a $\Sub(G)$ i la imatge $\exists_F(H)$ a $\Sub(G'')$, i comprova que totes dues es calculen component a component, sobre vèrtexs i sobre arestes. \hfill{\normalfont$\star$}
\begin{indicacio}
Fes servir l'exemple~\ref{ex:subgrafs} i comprova que les operacions component a component donen subconjunts tancats per origen i destí.
\end{indicacio}
\end{exercici}

\section{Regularitat}

\begin{exercici}[Ampliació, per a lectors amb àlgebra homològica]
Demostra que tota categoria abeliana és regular. \hfill{\normalfont$\star\star$}
\begin{indicacio}
En una categoria abeliana tot morfisme es factoritza com epi (sobre la coimatge) seguit de mono (la imatge), i coimatge i imatge coincideixen. L'estabilitat per pullback se segueix de l'exactitud. Pots limitar-te a esbossar per què cada condició de la definició es compleix.
\end{indicacio}
\end{exercici}

\begin{exercici}[Ampliació, per a lectors amb topologia]
Exhibeix una aplicació quocient a $\cat{Top}$ el pullback de la qual no és una aplicació quocient, i dedueix-ne que $\cat{Top}$ no és regular. \hfill{\normalfont$\star\star$}
\begin{indicacio}
A $\cat{Top}$, els epimorfismes regulars són les aplicacions quocient (no s'han de confondre amb els epimorfismes, que són les aplicacions contínues exhaustives). Busca una aplicació quocient el pullback de la qual, al llarg d'una aplicació adequada, deixa de ser quocient.
\end{indicacio}
\end{exercici}

\section{Quantificació}

\begin{exercici}
Demostra que $\exists_f$ és monòtona. Suposa, a més, que $\Sub(A)$ té un mínim $\bot_A$ (el predicat \enquote{fals}); demostra, fent servir l'adjunció $\exists_f\dashv f^{*}$, que $\exists_f(\bot_A)$ és el mínim de $\Sub(B)$. \hfill{\normalfont$\star$}
\begin{indicacio}
Per a la monotonia, recorda que els adjunts entre preordres són monòtons. Per al mínim, escriu la condició de l'adjunció per a $\exists_f(\bot_A)\leq[t]$.
\end{indicacio}
\end{exercici}

\begin{exercici}
Usant que $\exists_f$ és adjunt per l'esquerra, demostra que preserva unions (suprems) de subobjectes quan existeixen. \hfill{\normalfont$\star\star$}
\begin{indicacio}
Els adjunts per l'esquerra preserven tots els colímits que existeixin; en un preordre, els suprems són els colímits. Escriu la cadena de bicondicionals de l'adjunció aplicada a un suprem.
\end{indicacio}
\end{exercici}


\chapter[Classificadors de subobjectes i topos]{\texorpdfstring{Classificadors de subobjectes i introducció als topos}{Classificadors de subobjectes i introducció als topos}}

Els exercicis següents completen els resolts a la part de Pràctica. La marca \enquote{$\star$} assenyala un exercici de consolidació i \enquote{$\star\star$} un d'ampliació.

\section{Classificadors}

\begin{exercici}
Comprova que a $\cat{Set}$ el morfisme característic $\chi_U$ d'un subconjunt determina i és determinat per $U$, de manera que $\Sub(E)\cong\{0,1\}^{E}$. \hfill{\normalfont$\star$}
\begin{indicacio}
Associa a cada subconjunt la seva funció característica i comprova que l'aplicació és una bijecció amb inversa $\chi\mapsto\chi^{-1}(1)$.
\end{indicacio}
\end{exercici}

\begin{exercici}
En un topos, sigui $m\colon U\rightarrowtail E$ un subobjecte i $g\colon E'\to E$ un morfisme. Demostra que el morfisme característic de la reindexació $g^{*}m$ és
\[
\chi_{g^{*}m}=\chi_m\comp g.
\]
Interpreta-ho com la compatibilitat entre la representació $\Sub\cong\Hom(-,\Omega)$ i la substitució. \hfill{\normalfont$\star$}
\begin{indicacio}
Compara els dos pullbacks de $\mathsf{true}$ i fes servir el lema d'enganxament de pullbacks.
\end{indicacio}
\end{exercici}

\begin{exercici}
Sigui $E=\{1,\dots,6\}$ i $U=\{3,6\}$, el predicat \enquote{$x$ és múltiple de $3$}. Calcula $\chi_U\colon E\to\{0,1\}$. Sigui $g\colon\{a,b,c,d\}\to E$ l'aplicació amb $g(a)=3$, $g(b)=4$, $g(c)=6$ i $g(d)=3$. Calcula $\chi_U\comp g$ i identifica el subconjunt que classifica. \hfill{\normalfont$\star$}
\begin{indicacio}
Segueix l'exemple de la Teoria. Com que $g$ no és injectiva, fixa't que dos elements diferents poden tenir la mateixa imatge i que això no afecta la reindexació.
\end{indicacio}
\end{exercici}

\begin{exercici}
Comprova que $\mathsf{true}\colon1\to\Omega$ és sempre un monomorfisme, i que $\Hom(1,\Omega)\cong\Sub(1)$. Quants valors de veritat globals té $\cat{Set}$? \hfill{\normalfont$\star$}
\begin{indicacio}
Per a la primera part, compara dues fletxes qualssevol $X\to1$. Per a la segona, particularitza la representabilitat de $\Sub$.
\end{indicacio}
\end{exercici}

\section{Topos}

\begin{exercici}
Sigui $\mathbf 2=(0\to 1)$. La categoria de fletxes $\cat{Set}^{\mathbf 2}$ té per objectes les aplicacions $X_0\to X_1$, i és una categoria de prefeixos sobre $\mathbf 2\op$, per tant un topos. Fent servir la construcció amb sedassos:
\begin{enumerate}
  \item calcula els conjunts $\Omega(0)$ i $\Omega(1)$;
  \item calcula l'aplicació de restricció entre ells, de manera que $\Omega$ és un objecte de $\cat{Set}^{\mathbf 2}$, és a dir, una fletxa;
  \item determina els elements globals $1\to\Omega$, que són els \enquote{valors de veritat globals}.
\end{enumerate}
\hfill{\normalfont$\star\star$}
\begin{indicacio}
Compta els sedassos sobre cada objecte de $\mathbf 2\op$. Observa que $\Omega$ no és un conjunt de valors, sinó una fletxa entre dos conjunts.
\end{indicacio}
\end{exercici}

\begin{exercici}
Demostra que un topos que és alhora una categoria prima és equivalent a la categoria terminal $\cat{1}$. \hfill{\normalfont$\star\star$}
\begin{indicacio}
En una categoria prima, tot morfisme és mono. Aplica-ho al morfisme $A\to 1$ i fes servir que $\Sub(1)\cong\Hom(1,\Omega)$ té com a màxim un element.
\end{indicacio}
\end{exercici}

\begin{exercici}[Ampliació]
Fent servir el classificador de $\cat{Graph}$ calculat a la Teoria, determina els valors de veritat globals, és a dir, els morfismes $\mathbf L\to\Omega$ des del graf terminal d'un bucle. Comprova que n'hi ha tres, identifica cadascun amb un subgraf de $\mathbf L$ i dedueix que $\Omega$ no és isomorf a $1+1$. \hfill{\normalfont$\star\star$}
\begin{indicacio}
Un morfisme $\mathbf L\to\Omega$ tria un bucle de $\Omega$. Contrasta el resultat amb $\Sub(\mathbf L)\cong\Hom(\mathbf L,\Omega)$ i compta els morfismes $\mathbf L\to\mathbf L+\mathbf L$.
\end{indicacio}
\end{exercici}

\begin{exercici}
Dedueix de l'exercici resolt~\ref{pr-mono-equalitzador} de la Pràctica que $\cat{Top}$ no té classificador de subobjectes, i per tant no és un topos. \hfill{\normalfont$\star$}
\begin{indicacio}
A $\cat{Top}$, els monomorfismes són les aplicacions contínues injectives i els epimorfismes, les exhaustives. Busca una bijecció contínua entre dos espais de dos punts que no sigui un homeomorfisme.
\end{indicacio}
\end{exercici}

\begin{exercici}
Sigui $\cat{B}\mathbb N$ el monoide additiu $(\mathbb N,+,0)$ vist com a categoria d'un sol objecte $\ast$. Descriu tots els sedassos sobre $\ast$ i dedueix que $\Omega(\ast)$ està en bijecció amb $\mathbb N\cup\{\infty\}$. Descriu l'acció $\Omega(h)$ de cada $h\in\mathbb N$. Comprova que només hi ha dos valors de veritat globals, però que $\Sub(\yon(\ast))$ no és una àlgebra de Boole. \hfill{\normalfont$\star\star$}
\begin{indicacio}
Un sedàs no buit és tancat per sumar qualsevol natural; mira'n l'element mínim. Un valor global és un sedàs fix per totes les accions. Per a la darrera part, fes servir que $\Sub(\yon(\ast))\cong\Omega(\ast)$ i busca un sedàs $S$ tal que el més gran dels sedassos disjunts de $S$, unit amb $S$, no doni el sedàs màxim.
\end{indicacio}
\end{exercici}

\section{Grafs i lògica}

\begin{exercici}
Sigui $G$ el graf amb vèrtexs $a,b$ i arestes $\ell\colon a\to a$, $e\colon a\to b$ i $d\colon b\to a$, i sigui $H$ el subgraf format pel vèrtex $a$ i el bucle $\ell$. Fent servir la descripció de $\Omega$ de la Teoria, calcula $\chi_H\colon G\to\Omega$ i comprova que el pullback de $\mathsf{true}$ al llarg de $\chi_H$ és $H$. \hfill{\normalfont$\star$}
\begin{indicacio}
Per a cada aresta, mira quins dels seus dos extrems són a $H$ i si ella mateixa hi és, i tria l'aresta de $\Omega$ que registra aquesta informació.
\end{indicacio}
\end{exercici}

\begin{exercici}
A $\Sub(\mathbf E)$, on $\mathbf E$ és el graf d'una sola aresta $e\colon s\to t$, sigui $\neg K$ el més gran dels subgrafs disjunts de $K$. Comprova que existeix, calcula $\neg K$ per als cinc subgrafs de $\mathbf E$ i determina per a quins es compleix $\neg\neg K=K$. Què en dedueixes sobre la doble negació? \hfill{\normalfont$\star\star$}
\begin{indicacio}
Un subgraf que conté l'aresta ha de contenir els seus dos extrems. Fes una taula amb $K$, $\neg K$ i $\neg\neg K$.
\end{indicacio}
\end{exercici}

\fi

% ============================================================
% PART IV — TESTS INTERACTIUS
% ============================================================
\part{\NomTests}
\setcounter{chapter}{0}
% En pantalla, un sol entorn Form per a tots els tests: un únic diccionari
% AcroForm. En mode imprimible les macros dels tests són estàtiques i no es
% crea cap formulari.
\ifimprimible\else\begin{Form}\fi
\chapter{Categories}

Aquest test recull els conceptes essencials del capítol. Cada pregunta té una única resposta correcta.

\iniciTest{cat}{c,b,a,d,c,d,a,b,a,d,b,c}

\begin{enumerate}

% --- Definició i composició ---
\pregunta Un lector afirma: \enquote{Com que la composició d'una categoria és associativa, l'expressió $h\comp g\comp f$ està ben definida per a qualssevol tres morfismes $f,g,h$.} On falla?
\begin{enumerate}
  \opcio{a}{No falla: l'associativitat garanteix que sempre es pot compondre.}
  \opcio{b}{Falla perquè l'associativitat només val en categories petites.}
  \opcio{c}{Falla perquè la composició és \emph{parcial}.}
  \opcio{d}{Falla perquè $h\comp g\comp f$ mai no està definit sense parèntesis.}
\end{enumerate}
\feedback
\justificacio{$h\comp g\comp f$ només té sentit si $\cod(f)=\dom(g)$ i $\cod(g)=\dom(h)$; l'associativitat només diu que, quan les dues agrupacions existeixen, coincideixen.}

% --- Diagrames commutatius ---

\pregunta Considera el quadrat format per $f\colon A\to B$, $u\colon A\to C$, $v\colon B\to D$ i $g\colon C\to D$. Quina equació expressa que el quadrat commuta?
\begin{enumerate}
  \opcio{a}{$f\comp v=u\comp g$, perquè es llegeix en l'ordre en què es recorren les fletxes.}
  \opcio{b}{$v\comp f=g\comp u$.}
  \opcio{c}{$v\comp u=f\comp g$.}
  \opcio{d}{$f=g$ i $u=v$.}
\end{enumerate}
\feedback
\justificacio{Els dos camins de $A$ a $D$ són \enquote{primer $f$, després $v$} i \enquote{primer $u$, després $g$}, i la composició s'escriu de dreta a esquerra.}

% --- Exemples: lògica ---

\pregunta A la categoria $\cat{Prov}_T$ d'un sistema deductiu $T$, siguin $p$ i $q$ dues fórmules \emph{distintes} tals que $p\vdash_T q$ i $q\vdash_T p$. Què en podem dir?
\begin{enumerate}
  \opcio{a}{Són objectes isomorfs però no iguals: $\cat{Prov}_T$ és prima, però la deduïbilitat no és antisimètrica.}
  \opcio{b}{Han de ser iguals, perquè en una categoria prima dos objectes isomorfs coincideixen.}
  \opcio{c}{Hi ha exactament dues fletxes de $p$ a $q$, una per a cada deducció.}
  \opcio{d}{No hi ha cap fletxa entre elles, perquè són fórmules diferents.}
\end{enumerate}
\feedback
\justificacio{Hi ha fletxes en tots dos sentits, i com que la categoria és prima, les composicions són identitats: $p\cong q$. Però $p$ i $q$ són fórmules diferents. Una categoria prima és un preordre, no necessàriament un ordre parcial, i aquí es veu la diferència entre igualtat i isomorfisme.}

% --- Exemples: categories primes ---

\pregunta Quina de les categories següents \emph{no} és prima?
\begin{enumerate}
  \opcio{a}{La categoria d'un preordre $(P,\le)$.}
  \opcio{b}{La categoria $\cat{Prov}_T$ d'un sistema deductiu.}
  \opcio{c}{La categoria $\cat{2}=(0\to 1)$.}
  \opcio{d}{La categoria $\cat{B}M$ del monoide $(\mathbb N,+,0)$.}
\end{enumerate}
\feedback
\justificacio{Té un sol objecte, però $\Hom(\ast,\ast)=\mathbb N$ té més d'un element.}

% --- Isomorfismes ---

\pregunta Sigui $D$ el conjunt $\{0,1\}$ amb l'ordre discret i sigui $C$ el mateix conjunt amb l'ordre $0<1$. L'aplicació identitat subjacent $f\colon D\to C$ és bijectiva i monòtona. És un isomorfisme a $\cat{Pos}$?
\begin{enumerate}
  \opcio{a}{Sí, perquè qualsevol aplicació bijectiva és un isomorfisme.}
  \opcio{b}{Sí, perquè una aplicació monòtona bijectiva té sempre inversa monòtona.}
  \opcio{c}{No, perquè la inversa subjacent $C\to D$ no és monòtona.}
  \opcio{d}{No, perquè $f$ no és monòtona.}
\end{enumerate}
\feedback
\justificacio{Per tant no és un morfisme de $\cat{Pos}$, i $f$ no té invers en la categoria.}

% --- Mono/epi: cancel·lació ---

\pregunta L'aplicació $i\colon\varnothing\to\{\ast\}$ de $\cat{Set}$ és:
\begin{enumerate}
  \opcio{a}{un isomorfisme.}
  \opcio{b}{un epimorfisme amb secció.}
  \opcio{c}{ni mono ni epi.}
  \opcio{d}{un monomorfisme sense retracció.}
\end{enumerate}
\feedback
\justificacio{És injectiva i, per tant, un monomorfisme de $\cat{Set}$; però no té retracció, perquè no hi ha cap aplicació $\{\ast\}\to\varnothing$.}

\pregunta Un morfisme $f\colon A\to B$ té una \emph{retracció} $r$ (amb $r\comp f=\id_A$). Quina conclusió és correcta i per què?
\begin{enumerate}
  \opcio{a}{$f$ és un monomorfisme: si $f\comp g=f\comp h$, component amb $r$ per l'esquerra dona $g=h$.}
  \opcio{b}{$f$ és un isomorfisme, perquè té invers per un costat.}
  \opcio{c}{$f$ és un epimorfisme, perquè $r$ el cancel·la per la dreta.}
  \opcio{d}{No se'n pot dir res sense saber la categoria.}
\end{enumerate}
\feedback
\justificacio{Si $f\comp g=f\comp h$, aleshores $g=r\comp f\comp g=r\comp f\comp h=h$. Tenir invers per un costat no fa $f$ isomorfisme: a $\cat{Set}$, la inclusió $\{0\}\hookrightarrow\{0,1\}$ té retracció però no és bijectiva.}

% --- Dualitat i oposada ---

\pregunta Considera la categoria oposada $\cat{C}\op$. Un lector diu: \enquote{Si $f$ és un monomorfisme a $\cat{C}$, aleshores $f$ és un monomorfisme a $\cat{C}\op$.} És correcte?
\begin{enumerate}
  \opcio{a}{Sí, perquè $\cat{C}$ i $\cat{C}\op$ tenen els mateixos morfismes.}
  \opcio{b}{No, perquè un mono de $\cat{C}$ és un epi de $\cat{C}\op$.}
  \opcio{c}{No, perquè a $\cat{C}\op$ no hi ha monomorfismes.}
  \opcio{d}{Sí, perquè els monomorfismes són autoduals.}
\end{enumerate}
\feedback
\justificacio{En passar a l'oposada, la cancel·lació per l'esquerra es converteix en cancel·lació per la dreta, igual que una retracció passa a ser una secció.}

% --- Construccions de categories ---

\pregunta Sigui $A$ un objecte d'una categoria $\cat{C}$. Què és un morfisme de $x\colon X\to A$ a $y\colon Y\to A$ a la categoria $\cat{C}/A$ sobre $A$?
\begin{enumerate}
  \opcio{a}{Un morfisme $h\colon X\to Y$ de $\cat{C}$ tal que $y\comp h=x$.}
  \opcio{b}{Qualsevol morfisme $h\colon X\to Y$ de $\cat{C}$, sense cap condició.}
  \opcio{c}{Un morfisme $h\colon Y\to X$ tal que $x\comp h=y$.}
  \opcio{d}{Un morfisme $A\to A$.}
\end{enumerate}
\feedback
\justificacio{És a dir, que fa commutar el triangle sobre $A$.}

% --- Construccions de categories: categoria producte ---

\pregunta Siguin $\cat{C}$ i $\cat{D}$ dues categories. Què és un morfisme de $(A,X)$ a $(B,Y)$ a la categoria producte $\cat{C}\times\cat{D}$?
\begin{enumerate}
  \opcio{a}{Un morfisme $A\times X\to B\times Y$ d'una de les dues categories.}
  \opcio{b}{Un morfisme $A\to B$ de $\cat{C}$ o bé un morfisme $X\to Y$ de $\cat{D}$.}
  \opcio{c}{L'objecte producte $A\times B$, caracteritzat per la seva propietat universal dins de $\cat{C}$.}
  \opcio{d}{Una parella $(f,g)$ amb $f\colon A\to B$ a $\cat{C}$ i $g\colon X\to Y$ a $\cat{D}$.}
\end{enumerate}
\feedback
\justificacio{La composició es fa component a component. No s'ha de confondre amb el producte d'objectes dins d'una sola categoria.}

% --- Categoria lliure ---

\pregunta Considera el graf amb un sol vèrtex $\ast$ i una sola aresta $a\colon\ast\to\ast$. Quants morfismes té la categoria lliure que genera?
\begin{enumerate}
  \opcio{a}{Dos: la identitat i $a$.}
  \opcio{b}{Infinits, un per a cada longitud de camí.}
  \opcio{c}{Un de sol, perquè el graf té un sol vèrtex.}
  \opcio{d}{Cap, perquè $a$ no es pot compondre amb si mateixa.}
\end{enumerate}
\feedback
\justificacio{Són $\id_\ast,\ a,\ a\comp a,\ a\comp a\comp a,\dots$; és la categoria $\cat{B}M$ del monoide $(\mathbb N,+,0)$.}

% --- Mida ---

\pregunta Un lector afirma: \enquote{La categoria $\cat{Set}$ és petita, perquè els seus morfismes són simplement aplicacions.} On s'equivoca?
\begin{enumerate}
  \opcio{a}{No s'equivoca: $\cat{Set}$ és petita.}
  \opcio{b}{S'equivoca perquè $\cat{Set}$ no té identitats.}
  \opcio{c}{S'equivoca perquè $\cat{Set}$ no és petita.}
  \opcio{d}{S'equivoca perquè $\cat{Set}$ ni tan sols és una categoria.}
\end{enumerate}
\feedback
\justificacio{És \emph{localment petita} (cada $\Hom(A,B)$ és un conjunt), però la col·lecció de tots els conjunts és una classe pròpia, no un conjunt.}

\end{enumerate}
\clauestatica
\finalitzaTest

\chapter{Functors}

Aquest test recull els conceptes essencials del capítol. Cada pregunta té una única resposta correcta.

\iniciTest{fun}{a,b,d,c,a,c,d,a,b,c,b}

\begin{enumerate}

\pregunta Un lector afirma: \enquote{Per definir un functor $F$ n'hi ha prou de dir on va cada objecte; l'acció sobre els morfismes queda determinada.} On falla?
\begin{enumerate}
  \opcio{a}{Falla perquè l'acció sobre morfismes no queda determinada per la dels objectes.}
  \opcio{b}{Falla només quan $\cat{C}$ és gran.}
  \opcio{c}{No falla: els objectes determinen el functor.}
  \opcio{d}{Falla perquè els functors no actuen sobre objectes.}
\end{enumerate}
\feedback
\justificacio{Un functor és \emph{dues} assignacions, sobre objectes i sobre morfismes, que han de respectar dominis, codominis, composició i identitats.}

\pregunta Un lector proposa un \enquote{functor} $\cat{B}M\to\cat{B}M$ (amb $M=\{1,e\}$, $e\comp e=e$) que fixa l'objecte i envia el bucle $1$ a $e$ i el bucle $e$ a $e$. És un functor?
\begin{enumerate}
  \opcio{a}{Sí, perquè respecta dominis i codominis.}
  \opcio{b}{No, perquè no envia la identitat a la identitat.}
  \opcio{c}{Sí, perquè $e$ és idempotent.}
  \opcio{d}{No, perquè cap assignació entre monoides és functorial.}
\end{enumerate}
\feedback
\justificacio{La identitat $1$ de $\cat{B}M$ va a parar a $e\neq 1$; és un morfisme de grafs, però no un functor.}

\pregunta Es pot afirmar que \emph{tot} functor preserva els monomorfismes?
\begin{enumerate}
  \opcio{a}{Sí, perquè preserva tota l'estructura categòrica.}
  \opcio{b}{Sí, perquè preserva la composició.}
  \opcio{c}{No, perquè cap functor preserva mai els monomorfismes.}
  \opcio{d}{No en general, tot i que alguns functors sí que els preserven.}
\end{enumerate}
\feedback
\justificacio{Un functor preserva el que està testimoniat per morfismes i equacions (isomorfismes, seccions, retraccions), però ser mono es defineix per cancel·lació sobre \emph{totes} les fletxes de la categoria d'arribada, incloses les que no provenen de la categoria d'origen.}

\pregunta Sigui $\cat{S}$ la subcategoria plena de $\cat{FinSet}$ que té per objectes els conjunts $\underline{0}=\varnothing$ i $\underline{n}=\{1,\dots,n\}$ per a $n\geq 1$, un de cada cardinal finit, i sigui $\iota\colon\cat{S}\hookrightarrow\cat{FinSet}$ la inclusió. Quina afirmació és correcta?
\begin{enumerate}
  \opcio{a}{$\iota$ és exhaustiu sobre objectes en sentit estricte.}
  \opcio{b}{$\iota$ no és ni ple ni fidel.}
  \opcio{c}{$\iota$ és essencialment exhaustiu però no exhaustiu sobre objectes.}
  \opcio{d}{$\iota$ és un isomorfisme de categories.}
\end{enumerate}
\feedback
\justificacio{Tot conjunt finit és isomorf a algun $\underline{n}$, però en general no és igual a cap d'ells.}

\pregunta Sigui $F$ un functor \emph{ple i fidel} i suposem que $F(f)$ és un isomorfisme a $\cat{D}$. Quina conclusió és correcta, i sobre quin fet reposa?
\begin{enumerate}
  \opcio{a}{Que $f$ ja era un isomorfisme a $\cat{C}$.}
  \opcio{b}{Que no es pot concloure que $f$ sigui un isomorfisme sense suposar, a més, que $F$ és exhaustiu sobre objectes.}
  \opcio{c}{Res: ser ple i fidel no diu res sobre isomorfismes.}
  \opcio{d}{Que $F$ és una equivalència.}
\end{enumerate}
\feedback
\justificacio{Com que $F$ és ple, l'invers de $F(f)$ prové d'algun $g$; com que $F$ és fidel, les identitats $g\comp f=\id$ i $f\comp g=\id$ es dedueixen de les seves imatges. Els functors plens i fidels \emph{reflecteixen} isomorfismes.}

\pregunta Siguin $f\colon X\to Y$ i $h\colon Y\to B$. Quina és l'acció de l'homfunctor contravariant $\Hom(-,B)$ sobre $f$, aplicada a $h$?
\begin{enumerate}
  \opcio{a}{$h\mapsto f\comp h$, amb valors a $\Hom(X,B)$.}
  \opcio{b}{$h\mapsto f\comp h$, amb valors a $\Hom(Y,B)$.}
  \opcio{c}{$h\mapsto h\comp f$, amb valors a $\Hom(X,B)$.}
  \opcio{d}{$h\mapsto h\comp f$, amb valors a $\Hom(Y,B)$.}
\end{enumerate}
\feedback
\justificacio{La precomposició amb $f$ va de $\Hom(Y,B)$ a $\Hom(X,B)$: inverteix el sentit de $f$.}

\pregunta Un lector diu: \enquote{El functor oblit $U\colon\cat{Grp}\to\cat{Set}$ és ple, perquè envia cada homomorfisme a la seva aplicació subjacent.} On s'equivoca?
\begin{enumerate}
  \opcio{a}{No s'equivoca: $U$ és ple.}
  \opcio{b}{S'equivoca perquè $U$ no és ni tan sols un functor.}
  \opcio{c}{S'equivoca perquè $U$ no actua sobre morfismes.}
  \opcio{d}{S'equivoca: $U$ és fidel, però no ple.}
\end{enumerate}
\feedback
\justificacio{Homomorfismes diferents donen aplicacions diferents; però entre dos grups hi ha aplicacions de conjunts que no són homomorfismes i que, per tant, no provenen de cap morfisme de $\cat{Grp}$.}

\pregunta Un lector afirma: \enquote{Com que $\cat{Cat}$ té totes les categories petites per objectes, $\cat{Cat}$ és ella mateixa una categoria petita.} És correcte?
\begin{enumerate}
  \opcio{a}{No, perquè $\cat{Cat}$ és gran, encara que localment petita.}
  \opcio{b}{Sí, perquè els functors formen conjunts.}
  \opcio{c}{Sí: una categoria de categories petites és petita.}
  \opcio{d}{No, perquè $\cat{Cat}$ no té identitats.}
\end{enumerate}
\feedback
\justificacio{La col·lecció de totes les categories petites és una classe pròpia, com la de tots els conjunts.}

\pregunta Donats un graf $G$ i una categoria petita $\cat{D}$, definir un functor $\Free(G)\to\cat{D}$ des de la categoria lliure de $G$ equival a:
\begin{enumerate}
  \opcio{a}{triar només les imatges de les arestes de $G$.}
  \opcio{b}{donar un morfisme de grafs $G\to U(\cat{D})$.}
  \opcio{c}{triar un objecte de $\cat{D}$ per a cada camí de $G$ de manera independent.}
  \opcio{d}{donar una equivalència entre $\Free(G)$ i $\cat{D}$.}
\end{enumerate}
\feedback
\justificacio{Cal triar les imatges dels vèrtexs \emph{i} de les arestes, respectant origen i destí; aquesta tria s'estén de manera única als camins.}

\pregunta Sobre el bifunctor $\Hom(-,-)\colon\cat{C}\op\times\cat{C}\to\cat{Set}$, amb $\Hom(f,g)(h)=g\comp h\comp f$, quina descripció de la variància és correcta?
\begin{enumerate}
  \opcio{a}{Covariant en totes dues variables.}
  \opcio{b}{Contravariant en totes dues variables.}
  \opcio{c}{Contravariant en la primera variable i covariant en la segona.}
  \opcio{d}{Covariant en la primera i contravariant en la segona.}
\end{enumerate}
\feedback
\justificacio{Precomposició amb $f$ en la primera variable i postcomposició amb $g$ en la segona.}

% --- Fil lògic ---

\pregunta Siguin $T\subseteq T'$ dos sistemes deductius sobre les mateixes fórmules, i sigui
\[
J\colon\cat{Prov}_T\to\cat{Prov}_{T'}
\]
el functor identitat sobre fórmules. Quan és $J$ ple?
\begin{enumerate}
  \opcio{a}{Sempre, perquè és la identitat sobre objectes.}
  \opcio{b}{Exactament quan $T'$ no introdueix cap conseqüència nova $A\vdash_{T'}B$ entre les fórmules considerades.}
  \opcio{c}{Només quan $T=T'$ com a conjunts d'axiomes.}
  \opcio{d}{Mai, perquè les categories de deduïbilitat són primes.}
\end{enumerate}
\feedback
\justificacio{Totes dues categories són primes, de manera que $J$ és ple si i només si cada fletxa $A\to B$ de $\cat{Prov}_{T'}$ prové d'una de $\cat{Prov}_T$, és a dir, si $A\vdash_{T'}B$ implica $A\vdash_T B$. Ser la identitat sobre objectes (a) no diu res sobre els morfismes.}

\end{enumerate}
\clauestatica
\finalitzaTest

\chapter[Transformacions naturals]{\texorpdfstring{Transformacions naturals\\ i equivalències}{Transformacions naturals i equivalències}}

Aquest test recull els conceptes essencials del capítol. Cada pregunta té una única resposta correcta.

\iniciTest{trn}{c,b,c,d,a,b,c,b,a,d,a,d}

\begin{enumerate}

\pregunta A $\cat{Set}$, prenem $F=G=\id_{\cat{Set}}$ i definim $\eta_A=\id_A$ per a tot conjunt $A\neq\{0,1\}$, i $\eta_{\{0,1\}}=\tau$, la transposició $\tau(0)=1$, $\tau(1)=0$. Aquesta família, \emph{és} una transformació natural $\id\Rightarrow\id$?
\begin{enumerate}
  \opcio{a}{Sí, perquè cada component $\eta_A$ és una bijecció.}
  \opcio{b}{Sí, perquè $\eta$ coincideix amb la identitat a tots els conjunts excepte un.}
  \opcio{c}{No, perquè una aplicació constant trenca el quadrat de naturalitat.}
  \opcio{d}{No, perquè la transposició $\tau$ no és un morfisme de $\cat{Set}$.}
\end{enumerate}
\feedback
\justificacio{Per a $f\colon\{0,1\}\to\{0,1\}$ constant de valor $0$, es té $f\comp\tau=f$, mentre que $\tau\comp f$ és constant de valor $1$.}

\pregunta Un lector proposa aquest argument: \enquote{Si $F(A)\cong G(A)$ per a cada objecte $A$, aleshores $F$ i $G$ són naturalment isomorfs, perquè n'hi ha prou d'agafar un isomorfisme $\eta_A\colon F(A)\to G(A)$ per a cada $A$.} On falla?
\begin{enumerate}
  \opcio{a}{No falla: l'argument és correcte.}
  \opcio{b}{Falla perquè els $\eta_A$ triats poden no complir el quadrat de naturalitat.}
  \opcio{c}{Falla perquè un isomorfisme natural exigeix que $F=G$.}
  \opcio{d}{Falla perquè $F(A)$ i $G(A)$ no poden ser mai isomorfs si $F\neq G$.}
\end{enumerate}
\feedback
\justificacio{L'existència d'isomorfismes objecte a objecte no garanteix que siguin \emph{compatibles} amb l'acció dels functors sobre els morfismes.}

\pregunta Sigui $\mathbb{Z}/2\mathbb{Z}=\{1,g\}$ vist com a categoria d'un objecte, i $F,G\colon\cat{B}(\mathbb{Z}/2\mathbb{Z})\to\cat{Set}$ que envien l'objecte a $\{0,1\}$, on $g$ actua per la identitat en $F$ i per la transposició $\tau$ en $G$. Quantes transformacions naturals $\eta\colon F\Rightarrow G$ hi ha?
\begin{enumerate}
  \opcio{a}{Exactament una, l'única component possible $\eta_\ast=\id$.}
  \opcio{b}{Dues: la identitat i la transposició.}
  \opcio{c}{Cap.}
  \opcio{d}{Infinites, una per cada aplicació $\{0,1\}\to\{0,1\}$.}
\end{enumerate}
\feedback
\justificacio{La naturalitat respecte de $g$ exigiria $\tau\comp\eta_\ast=\eta_\ast$, impossible perquè $\tau$ no té punts fixos.}

\pregunta Què és un morfisme de $f\colon A\to B$ a $g\colon A'\to B'$ a la categoria de functors $\cat{D}^{\cat{2}}$, on $\cat{2}=(0\to 1)$?
\begin{enumerate}
  \opcio{a}{Un morfisme qualsevol $A\to A'$.}
  \opcio{b}{Un functor $\cat{2}\to\cat{D}$.}
  \opcio{c}{Una parella qualsevol de morfismes $u\colon A\to A'$ i $v\colon B\to B'$.}
  \opcio{d}{Una parella $u\colon A\to A'$, $v\colon B\to B'$ amb $v\comp f=g\comp u$.}
\end{enumerate}
\feedback
\justificacio{És la condició de naturalitat per a l'única fletxa no trivial de $\cat{2}$.}

\pregunta En l'exemple sintaxi/semàntica, amb el functor sintàctic $F$, el semàntic $G$ i la interpretació $\eta$, què expressa la igualtat $G(f)\comp\eta_X=\eta_Y\comp F(f)$ per a un canvi de variables $f\colon X\to Y$?
\begin{enumerate}
  \opcio{a}{Que substituir variables i després interpretar dona la mateixa funció de veritat que interpretar i després transportar la semàntica al llarg del canvi de variables.}
  \opcio{b}{Que $\eta_X$ és una bijecció entre fórmules i funcions de veritat.}
  \opcio{c}{Que dues fórmules amb la mateixa taula de veritat són iguals.}
  \opcio{d}{Que el canvi de variables $f$ ha de ser bijectiu.}
\end{enumerate}
\feedback
\justificacio{És el quadrat de naturalitat de $\eta$ per a $f$: els dos camins de $F(X)$ a $G(Y)$ coincideixen. En termes lògics, la interpretació és compatible amb la substitució. Ni (b) ni (c) són certes, perquè fórmules diferents poden tenir la mateixa taula de veritat.}

\pregunta Un estudiant escriu que la composició vertical de $\eta\colon F\Rightarrow G$ i $\theta\colon G\Rightarrow H$ té components $(\theta\comp\eta)_A=\eta_A\comp\theta_A$. Per què és incorrecte, i quina és la forma bona?
\begin{enumerate}
  \opcio{a}{És correcte tal com està.}
  \opcio{b}{L'ordre està invertit: la composta és $\theta_A\comp\eta_A$.}
  \opcio{c}{Cap dels dos ordres té sentit; la composició vertical no existeix.}
  \opcio{d}{Tots dos ordres donen el mateix resultat, així que és igual.}
\end{enumerate}
\feedback
\justificacio{$\eta_A\colon F(A)\to G(A)$ i $\theta_A\colon G(A)\to H(A)$ només es componen en aquest ordre, i la composta va de $F(A)$ a $H(A)$.}

\pregunta Sigui $L\colon\cat{Set}\to\cat{Set}$ el functor de llistes finites, amb $L(f)$ aplicant $f$ a cada entrada, i sigui $\eta_X(x)=[x]$ la llista d'un sol element. Per a $f\colon X\to Y$, quina és l'equació de naturalitat ben tipada?
\begin{enumerate}
  \opcio{a}{$f\comp\eta_X=\eta_Y\comp L(f)$.}
  \opcio{b}{$\eta_X\comp L(f)=f\comp\eta_Y$.}
  \opcio{c}{$L(f)\comp\eta_X=\eta_Y\comp f$.}
  \opcio{d}{$\eta_Y=L(f)$.}
\end{enumerate}
\feedback
\justificacio{Tots dos costats envien $x$ a $[f(x)]$.}

\pregunta Un functor $F\colon\cat{C}\to\cat{D}$ és ple, fidel i essencialment exhaustiu. Un lector conclou: \enquote{Aleshores $F$ és bijectiu sobre els objectes.} És correcte?
\begin{enumerate}
  \opcio{a}{Sí: essencialment exhaustiu vol dir exhaustiu sobre objectes, i ple i fidel vol dir injectiu sobre objectes.}
  \opcio{b}{No: $F$ pot no ser bijectiu sobre objectes.}
  \opcio{c}{No, perquè cap functor és mai bijectiu sobre objectes.}
  \opcio{d}{Sí, sempre que $\cat{C}$ i $\cat{D}$ siguin petites.}
\end{enumerate}
\feedback
\justificacio{Essencialment exhaustiu només demana que tot objecte de $\cat{D}$ sigui \emph{isomorf} a una imatge, i ser ple i fidel afecta els morfismes; per exemple, la inclusió d'un esquelet no és bijectiva sobre objectes.}

\pregunta A la construcció del quasiinvers $G$ d'una equivalència, per a cada objecte $D$ es tria un objecte $G(D)$ i un isomorfisme $\varepsilon_D\colon F(G(D))\to D$, i sobre morfismes es defineix $G(v)$ com l'\emph{única} fletxa amb $F(G(v))=\varepsilon_{D'}^{-1}\comp v\comp\varepsilon_D$. Quina propietat de $F$ garanteix aquesta unicitat i existència?
\begin{enumerate}
  \opcio{a}{Que $F$ sigui ple i fidel.}
  \opcio{b}{Només ser fidel.}
  \opcio{c}{L'exhaustivitat essencial.}
  \opcio{d}{Que $\cat{D}$ sigui petita.}
\end{enumerate}
\feedback
\justificacio{Ser ple dona l'existència de $G(v)$, i ser fidel, la unicitat.}

\pregunta Siguin $\eta\colon F\Rightarrow G$, amb $F,G\colon\cat{C}\to\cat{D}$, i $\theta\colon K\Rightarrow L$, amb $K,L\colon\cat{D}\to\cat{E}$. Quina és la component $(\theta\ast\eta)_A$ de la composició horitzontal?
\begin{enumerate}
  \opcio{a}{$\theta_A\comp\eta_A$.}
  \opcio{b}{$K(\eta_A)\comp\theta_{G(A)}$.}
  \opcio{c}{$\eta_{K(A)}\comp\theta_A$.}
  \opcio{d}{$\theta_{G(A)}\comp K(\eta_A)$.}
\end{enumerate}
\feedback
\justificacio{Coincideix amb $L(\eta_A)\comp\theta_{F(A)}$ per la naturalitat de $\theta$ aplicada a $\eta_A$.}

\pregunta Siguin $\cat{1}$ la categoria terminal i $\cat{I}$ la categoria amb dos objectes i exactament un morfisme entre cada parella ordenada d'objectes. Quina afirmació és correcta?
\begin{enumerate}
  \opcio{a}{Són equivalents però no isomorfes.}
  \opcio{b}{Són isomorfes, perquè tots els objectes de $\cat{I}$ són isomorfs entre si.}
  \opcio{c}{No són equivalents, perquè tenen un nombre diferent d'objectes.}
  \opcio{d}{Són isomorfes, perquè totes dues són categories primes.}
\end{enumerate}
\feedback
\justificacio{La inclusió $\cat{1}\to\cat{I}$ és plena, fidel i essencialment exhaustiva, però cap functor entre elles no pot ser bijectiu sobre objectes.}

\pregunta Sigui $F\colon\cat{Mat}_{\mathbb K}\to\cat{FinVect}_{\mathbb K}$, $n\mapsto\mathbb K^n$, el functor de l'exemple de la Teoria. Quina afirmació és correcta?
\begin{enumerate}
  \opcio{a}{Les dues categories són iguals, perquè tot espai de dimensió $n$ \emph{és} $\mathbb K^n$.}
  \opcio{b}{Són isomorfes però no iguals: $F$ té un invers estricte.}
  \opcio{c}{Són equivalents i, per tant, isomorfes.}
  \opcio{d}{Són equivalents però no isomorfes; amb una tria adequada de bases, $G\comp F=\id$ literalment, però $F\comp G$ només és isomorf a la identitat.}
\end{enumerate}
\feedback
\justificacio{$F$ és ple, fidel i essencialment exhaustiu, i per tant una equivalència. No hi ha cap isomorfisme de categories, perquè $\cat{Mat}_{\mathbb K}$ és esquelètica i $\cat{FinVect}_{\mathbb K}$ no ho és. El quasiinvers tria una base de cada espai: $F(G(V))=\mathbb K^{\dim V}$, que és isomorf a $V$ però en general no hi és igual. Un espai de dimensió $n$ és isomorf a $\mathbb K^n$, no igual.}

\end{enumerate}
\clauestatica
\finalitzaTest

\ifpublicacioparcial\else
\chapter[Representabilitat i Yoneda]{\texorpdfstring{Propietats universals,\\ representabilitat i Yoneda}{Propietats universals, representabilitat i Yoneda}}

Aquest test recull els conceptes essencials del capítol. Cada pregunta té una única resposta correcta.

\iniciTest{yon}{a,c,b,d,a,d,b,c,a,c,b,d,c,d,b}

\begin{enumerate}

\pregunta Quan és \textbf{terminal} un objecte $1$ d'una categoria?
\begin{enumerate}
  \opcio{a}{Quan hi ha exactament un morfisme $X\to 1$ per a cada $X$.}
  \opcio{b}{Quan hi ha exactament un morfisme $1\to X$ per a cada $X$.}
  \opcio{c}{Quan no hi ha cap morfisme cap a ell.}
  \opcio{d}{Quan és isomorf a tots els altres objectes.}
\end{enumerate}
\feedback
\justificacio{És la definició. L'opció (b) és la d'objecte inicial, i (d) no té per què complir-se: ser terminal és una propietat universal, no un isomorfisme amb tots els objectes.}

\pregunta Siguin $S\colon\cat{A}\to\cat{C}$ i $T\colon\cat{B}\to\cat{C}$. Una parella $a\colon A\to A'$, $b\colon B\to B'$ és un morfisme $(A,B,u)\to(A',B',u')$ de la categoria coma $(S\downarrow T)$ quan:
\begin{enumerate}
  \opcio{a}{$S(a)\comp u=u'\comp T(b)$.}
  \opcio{b}{$T(b)\comp u'=u\comp S(a)$.}
  \opcio{c}{$T(b)\comp u=u'\comp S(a)$.}
  \opcio{d}{$S(a)=T(b)$.}
\end{enumerate}
\feedback
\justificacio{Comprova els tipus: $u\colon S(A)\to T(B)$ i $u'\colon S(A')\to T(B')$. L'únic quadrat ben tipat de $S(A)$ a $T(B')$ és $T(b)\comp u=u'\comp S(a)$. A (a) i (b) les composicions no existeixen en general.}

\pregunta En la formulació per representabilitat desenvolupada en aquest capítol, una \textbf{propietat universal} caracteritza una dada...
\begin{enumerate}
  \opcio{a}{per la seva constitució interna.}
  \opcio{b}{per la manera com es relaciona amb tots els altres objectes, mitjançant la representabilitat d'un functor cap a $\cat{Set}$.}
  \opcio{c}{per la seva cardinalitat.}
  \opcio{d}{només a $\cat{Set}$.}
\end{enumerate}
\feedback
\justificacio{Una propietat universal no descriu què \emph{és} la dada, sinó com es relaciona amb tota la resta: la dada és un objecte representant, amb el seu element universal.}

\pregunta Que $1$ sigui terminal equival a dir que:
\begin{enumerate}
  \opcio{a}{$\Hom(1,-)$ és constant.}
  \opcio{b}{$\Hom(-,1)$ no és functor.}
  \opcio{c}{$\Hom(-,1)$ és buit.}
  \opcio{d}{$\Hom(-,1)$ és isomorf al functor constant amb valor un punt.}
\end{enumerate}
\feedback
\justificacio{Ser terminal vol dir que cada $\Hom(X,1)$ té exactament un element. Aquestes bijeccions amb $\{\ast\}$ formen un isomorfisme natural, perquè tots els quadrats acaben en un conjunt d'un element.}

\pregunta Quan es diu que un prefeix $P\colon\cat{C}\op\to\cat{Set}$ és \textbf{representable}?
\begin{enumerate}
  \opcio{a}{Quan és isomorf a $\Hom(-,A)$ per a algun objecte $A$.}
  \opcio{b}{Quan és constant.}
  \opcio{c}{Quan és ple i fidel.}
  \opcio{d}{Quan té un sol element a cada conjunt.}
\end{enumerate}
\feedback
\justificacio{És la definició, en la forma de la proposició sobre representabilitat i isomorfisme natural. Ser constant o tenir un element a cada conjunt són casos molt particulars; \enquote{ple i fidel} és una propietat de functors entre categories, no de prefeixos.}

\pregunta Sigui $P\colon\cat{C}\op\to\cat{Set}$ un prefeix i $x\in P(A)$. Quina és la component en $X$ de la transformació natural $\Hom(-,A)\Rightarrow P$ que el lema de Yoneda associa a $x$, avaluada en $f\colon X\to A$?
\begin{enumerate}
  \opcio{a}{$P(f^{-1})(x)$.}
  \opcio{b}{$x$, sense canviar de conjunt.}
  \opcio{c}{$P(\id_A)(f)$.}
  \opcio{d}{$P(f)(x)$.}
\end{enumerate}
\feedback
\justificacio{Per contravariància, $P(f)\colon P(A)\to P(X)$; no cal que $f$ tingui invers.}

\pregunta Sigui $f\colon A\to B$ un morfisme. Les quatre descripcions següents de la component en $X$ de $\yon(f)=\Hom(-,f)$ semblen raonables. Quina és l'única en què el domini, el codomini i la composició tenen els tipus correctes?
\begin{enumerate}
  \opcio{a}{$\Hom(X,A)\to\Hom(X,B)$, $\varphi\mapsto\varphi\comp f$.}
  \opcio{b}{$\Hom(X,A)\to\Hom(X,B)$, $\varphi\mapsto f\comp\varphi$.}
  \opcio{c}{$\Hom(A,X)\to\Hom(B,X)$, $\varphi\mapsto\varphi\comp f$.}
  \opcio{d}{$\Hom(X,B)\to\Hom(X,A)$, $\varphi\mapsto\varphi\comp f$.}
\end{enumerate}
\feedback
\justificacio{Si $\varphi\colon X\to A$, l'única composició possible amb $f\colon A\to B$ és $f\comp\varphi\colon X\to B$. A (a), $\varphi\comp f$ demanaria $B=X$. A (c), amb $\varphi\colon A\to X$, $\varphi\comp f$ demanaria $B=A$: és una versió mal tipada de $\Hom(f,X)$, que va de $\Hom(B,X)$ a $\Hom(A,X)$. A (d), amb $\varphi\colon X\to B$, $\varphi\comp f$ demanaria de nou $B=X$.}

\pregunta Sigui $\mathcal P$ el prefeix de parts, $A=\{0,1\}$ i $B=\{1\}$. Quin subconjunt assigna la transformació natural $\Hom(-,A)\Rightarrow\mathcal P$ determinada per $B$ a l'aplicació $f\colon\{a,b,c\}\to A$ amb $f(a)=f(c)=1$ i $f(b)=0$?
\begin{enumerate}
  \opcio{a}{$\{b\}$.}
  \opcio{b}{$\{a,b,c\}$.}
  \opcio{c}{$\{a,c\}$.}
  \opcio{d}{$\{1\}$.}
\end{enumerate}
\feedback
\justificacio{Pel lema de Yoneda, la transformació determinada per $B\in\mathcal P(A)$ envia $f$ a $\mathcal P(f)(B)=f^{-1}(B)=\{a,c\}$. L'opció (d) és $B$ mateix, que viu a $\mathcal P(A)$ i no a $\mathcal P(\{a,b,c\})$.}

\pregunta Sigui $\Phi_{A,P}\colon\Nat(\yon(A),P)\to P(A)$, $\eta\mapsto\eta_A(\id_A)$, la bijecció del lema de Yoneda, i sigui $a\colon A'\to A$. Quina de les igualtats següents expressa la naturalitat en $A$ amb tots els tipus correctes?
\begin{enumerate}
  \opcio{a}{$\Phi_{A',P}\bigl(\eta\comp\yon(a)\bigr)=P(a)\bigl(\Phi_{A,P}(\eta)\bigr)$, com a elements de $P(A')$.}
  \opcio{b}{$\Phi_{A,P}\bigl(\eta\comp\yon(a)\bigr)=P(a)\bigl(\Phi_{A',P}(\eta)\bigr)$, com a elements de $P(A)$.}
  \opcio{c}{$\Phi_{A',P}\bigl(\yon(a)\comp\eta\bigr)=P(a)\bigl(\Phi_{A,P}(\eta)\bigr)$, com a elements de $P(A')$.}
  \opcio{d}{$\Phi_{A',P}\bigl(\eta\comp\yon(a)\bigr)=P(a)^{-1}\bigl(\Phi_{A,P}(\eta)\bigr)$, com a elements de $P(A')$.}
\end{enumerate}
\feedback
\justificacio{Tenim $\eta\colon\yon(A)\Rightarrow P$ i $\yon(a)\colon\yon(A')\Rightarrow\yon(A)$. L'única composició possible és, doncs, $\eta\comp\yon(a)$, que va de $\yon(A')$ a $P$, i s'hi aplica $\Phi_{A',P}$. Per contravariància, $P(a)\colon P(A)\to P(A')$. A (b), $\Phi_{A,P}$ no s'aplica a una transformació que surt de $\yon(A')$. A (c), $\yon(a)\comp\eta$ no és componible. A (d), $P(a)$ no té per què ser invertible.}

\pregunta A la demostració del lema, a quina dada s'envia una transformació natural $\eta\colon\Hom(-,A)\Rightarrow P$?
\begin{enumerate}
  \opcio{a}{A $\eta_A$.}
  \opcio{b}{A tot el functor $P$.}
  \opcio{c}{A l'element $\eta_A(\id_A)\in P(A)$.}
  \opcio{d}{A la identitat de $P$.}
\end{enumerate}
\feedback
\justificacio{L'aplicació del lema és $\eta\mapsto\eta_A(\id_A)$. L'opció (a) és tota la component, no un element de $P(A)$; és l'element el que determina $\eta$.}

\pregunta Una conseqüència directa del lema de Yoneda és que la immersió $\yon$ és:
\begin{enumerate}
  \opcio{a}{constant.}
  \opcio{b}{plena i fidel.}
  \opcio{c}{essencialment exhaustiva.}
  \opcio{d}{un isomorfisme de categories.}
\end{enumerate}
\feedback
\justificacio{Amb $P=\Hom(-,B)$, el lema dona $\Nat(\Hom(-,A),\Hom(-,B))\cong\Hom(A,B)$, i la bijecció és $f\mapsto\yon(f)$. No és essencialment exhaustiva: hi ha prefeixos que no són representables.}

\pregunta Segons el corol·lari de Yoneda, $A\cong B$ si i només si:
\begin{enumerate}
  \opcio{a}{$A$ i $B$ tenen la mateixa cardinalitat.}
  \opcio{b}{$A=B$.}
  \opcio{c}{hi ha un functor entre ells.}
  \opcio{d}{$\Hom(-,A)\cong\Hom(-,B)$ com a prefeixos.}
\end{enumerate}
\feedback
\justificacio{Un functor ple i fidel reflecteix isomorfismes, i tot functor els preserva; aplicat a $\yon$, dona l'equivalència. La igualtat (b) és massa forta, i la cardinalitat (a) no té sentit en una categoria general.}

\pregunta Siguin $(P,p_1,p_2)$ i $(Q,q_1,q_2)$ dos productes dels mateixos objectes. Què és únic?
\begin{enumerate}
  \opcio{a}{Res: poden ser literalment iguals o no.}
  \opcio{b}{L'isomorfisme entre els objectes $P$ i $Q$, sense cap condició.}
  \opcio{c}{L'isomorfisme $u\colon P\to Q$ que compleix $q_i\comp u=p_i$ per a $i=1,2$.}
  \opcio{d}{L'automorfisme de $P$, que només pot ser la identitat.}
\end{enumerate}
\feedback
\justificacio{La propietat universal fa que l'isomorfisme sigui únic \emph{entre els compatibles amb les projeccions}. Sense aquesta condició, pot haver-n'hi molts: a $\cat{Set}$, per exemple, qualsevol bijecció de $P$ en $Q$.}

\pregunta Siguin $S\colon\cat{A}\to\cat{C}$ i $T\colon\cat{B}\to\cat{C}$ dos functors. Un objecte de la categoria coma $(S\downarrow T)$ és:
\begin{enumerate}
  \opcio{a}{només un objecte de $\cat{C}$.}
  \opcio{b}{una transformació natural $S\Rightarrow T$.}
  \opcio{c}{una parella d'objectes isomorfs.}
  \opcio{d}{una terna $(A,B,f)$ amb $A$ de $\cat{A}$, $B$ de $\cat{B}$ i $f\colon S(A)\to T(B)$.}
\end{enumerate}
\feedback
\justificacio{És la definició: un objecte de $\cat{A}$, un de $\cat{B}$ i una fletxa de $S(A)$ a $T(B)$ a $\cat{C}$. Les categories sobre i sota un objecte en són casos particulars.}

\pregunta Si $(A,u)$ representa un prefeix $P\colon\cat{C}\op\to\cat{Set}$, què expressa la universalitat de l'element $u\in P(A)$?
\begin{enumerate}
  \opcio{a}{Que $u$ és l'únic element de $P(A)$.}
  \opcio{b}{Que per a cada objecte $X$ i cada $x\in P(X)$ hi ha un únic $f\colon X\to A$ amb $P(f)(u)=x$.}
  \opcio{c}{Que $A$ és un objecte terminal.}
  \opcio{d}{Que $P$ és un functor constant.}
\end{enumerate}
\feedback
\justificacio{És la definició de representació d'un prefeix: cada element de cada $P(X)$ s'obté de $u$ per un únic morfisme $X\to A$. Equival a l'isomorfisme natural $\Hom(-,A)\cong P$ que envia $f$ a $P(f)(u)$.}

\end{enumerate}
\clauestatica
\finalitzaTest

\chapter[Límits i colímits]{Límits i colímits}

Aquest test recull els conceptes essencials del capítol. Cada pregunta té una única resposta correcta.

\iniciTest{lim}{b,c,a,d,d,b,c,a,d,c,a,c,b,b,a}

\begin{enumerate}

\pregunta Què és un \textbf{diagrama} de forma $\cat{I}$ en una categoria $\cat{C}$?
\begin{enumerate}
  \opcio{a}{Un objecte de $\cat{C}$.}
  \opcio{b}{Un functor $D\colon\cat{I}\to\cat{C}$, on $\cat{I}$ és una categoria índex petita.}
  \opcio{c}{Una transformació natural.}
  \opcio{d}{Un morfisme de $\cat{I}$.}
\end{enumerate}
\feedback
\justificacio{És la definició: un diagrama és un functor, i la categoria índex en dona la forma. Un objecte o un morfisme sols no recullen les compatibilitats entre les fletxes.}

\pregunta Un \textbf{con} sobre $D$ amb vèrtex $A$ és el mateix que:
\begin{enumerate}
  \opcio{a}{un objecte terminal.}
  \opcio{b}{un límit.}
  \opcio{c}{una transformació natural $\Delta_A\Rightarrow D$ des del functor constant.}
  \opcio{d}{un functor $\cat{C}\to\cat{I}$.}
\end{enumerate}
\feedback
\justificacio{Les components de $\alpha\colon\Delta_A\Rightarrow D$ són fletxes $A\to D_i$, i la naturalitat per a $u\colon i\to j$ diu $D_u\comp\alpha_i=\alpha_j$, que és exactament la condició de con.}

\pregunta Un \textbf{límit} de $D$ és:
\begin{enumerate}
  \opcio{a}{un objecte terminal de la categoria de cons $\mathrm{Con}(D)$.}
  \opcio{b}{qualsevol con sobre $D$.}
  \opcio{c}{un objecte inicial de $\cat{C}$.}
  \opcio{d}{el producte de tots els objectes.}
\end{enumerate}
\feedback
\justificacio{Un con límit és un con pel qual factoritza de manera única qualsevol altre con: és a dir, un objecte terminal de $\mathrm{Con}(D)$. D'aquí ve la unicitat llevat d'un únic isomorfisme compatible amb les projeccions.}

\pregunta Quina categoria índex dona lloc al \textbf{producte} de dos objectes?
\begin{enumerate}
  \opcio{a}{La categoria buida.}
  \opcio{b}{Dues fletxes paral·leles.}
  \opcio{c}{La forma de racó.}
  \opcio{d}{La categoria discreta amb dos objectes.}
\end{enumerate}
\feedback
\justificacio{Un con sobre un parell d'objectes, sense cap fletxa entre ells, és una parella de fletxes sense cap condició: la propietat del producte. Dues fletxes paral·leles donen l'equalitzador, i el racó, el pullback.}

\pregunta A $\cat{Set}$, l'\textbf{equalitzador} de $f,g\colon A\to B$ és:
\begin{enumerate}
  \opcio{a}{el producte cartesià $A\times B$.}
  \opcio{b}{la unió disjunta.}
  \opcio{c}{el quocient de $B$.}
  \opcio{d}{el subconjunt $\{a\in A\mid f(a)=g(a)\}$.}
\end{enumerate}
\feedback
\justificacio{Una fletxa $h\colon X\to A$ amb $f\comp h=g\comp h$ pren valors dins d'aquest subconjunt, i hi factoritza de manera única. L'opció (c) és el coequalitzador, el concepte dual.}

\pregunta El \textbf{pullback} de $f\colon A\to C$ i $g\colon B\to C$ a $\cat{Set}$ és:
\begin{enumerate}
  \opcio{a}{$A\times B$ sencer.}
  \opcio{b}{$\{(a,b)\in A\times B\mid f(a)=g(b)\}$.}
  \opcio{c}{la unió $A\cup B$.}
  \opcio{d}{el quocient $A/B$.}
\end{enumerate}
\feedback
\justificacio{Una parella $x_1\colon X\to A$, $x_2\colon X\to B$ amb $f\comp x_1=g\comp x_2$ factoritza de manera única per $x\mapsto(x_1(x),x_2(x))$, que pren valors en aquest subconjunt. Tot $A\times B$ només és el pullback quan $C$ té un sol element.}

\pregunta L'\textbf{objecte terminal} és el límit de quin diagrama?
\begin{enumerate}
  \opcio{a}{El diagrama de dues fletxes paral·leles.}
  \opcio{b}{El diagrama d'un sol objecte.}
  \opcio{c}{El diagrama buit.}
  \opcio{d}{El diagrama de racó.}
\end{enumerate}
\feedback
\justificacio{Un con sobre el diagrama buit és només un vèrtex, sense projeccions. Un con límit és, doncs, un objecte al qual arriba una única fletxa des de cada objecte.}

\pregunta Un \textbf{colímit} d'un diagrama $D$ és:
\begin{enumerate}
  \opcio{a}{un objecte inicial a la categoria de cocons (un límit a la categoria oposada).}
  \opcio{b}{un altre nom per al límit.}
  \opcio{c}{sempre l'objecte terminal.}
  \opcio{d}{un con.}
\end{enumerate}
\feedback
\justificacio{És la noció dual: un cocon universal, és a dir, un objecte inicial de la categoria de cocons, o equivalentment un límit del diagrama corresponent a $\cat{C}\op$. No és un límit de la mateixa categoria.}

\pregunta Si $\cat{C}$ té límits de forma $\cat{I}$ i $\cat{A}$ és petita, com es calcula el límit d'un diagrama $F\colon\cat{I}\to\cat{C}^{\cat{A}}$?
\begin{enumerate}
  \opcio{a}{Prenent el límit només sobre els objectes terminals de $\cat{A}$.}
  \opcio{b}{Aplicant el functor de parts a cada $F(i)$.}
  \opcio{c}{No en té: les categories de functors no tenen límits.}
  \opcio{d}{Punt a punt: $\bigl(\varprojlim F\bigr)(a)\cong\varprojlim_i F(i)(a)$.}
\end{enumerate}
\feedback
\justificacio{Els límits a $\cat{C}^{\cat{A}}$ es calculen objecte per objecte a $\cat{C}$; equivalentment, cada functor d'avaluació $\mathrm{ev}_a$ els preserva.}

\pregunta Quina és la caracterització dels \textbf{límits finits}?
\begin{enumerate}
  \opcio{a}{N'hi ha prou amb tenir coproductes.}
  \opcio{b}{Cal tenir tots els límits de forma petita.}
  \opcio{c}{Equivalen a tenir objecte terminal, productes binaris i equalitzadors (o terminal i pullbacks).}
  \opcio{d}{No es poden construir a partir de casos més simples.}
\end{enumerate}
\feedback
\justificacio{És el teorema de caracterització: el límit d'un diagrama finit es construeix com un equalitzador de dues fletxes entre productes finits. Amb terminal i pullbacks s'obtenen productes i equalitzadors.}

\pregunta Quan \textbf{preserva} un functor $F$ el límit d'un diagrama $D$?
\begin{enumerate}
  \opcio{a}{Quan $(F(L),F\lambda)$ és un límit de $F\comp D$.}
  \opcio{b}{Sempre, per definició de functor.}
  \opcio{c}{Quan $F$ és constant.}
  \opcio{d}{Quan $F$ no té adjunts.}
\end{enumerate}
\feedback
\justificacio{És la definició de preservació. No tot functor preserva límits: per exemple, l'oblit $\cat{Grp}\to\cat{Set}$ no preserva els coproductes, que són límits a l'oposada.}

\pregunta El functor oblit $U\colon\cat{Grp}\to\cat{Set}$:
\begin{enumerate}
  \opcio{a}{no preserva cap límit ni cap colímit.}
  \opcio{b}{preserva tots els colímits, però no tots els límits.}
  \opcio{c}{preserva tots els límits, però no tots els colímits; per exemple, no preserva tots els coproductes.}
  \opcio{d}{preserva tots els límits i tots els colímits.}
\end{enumerate}
\feedback
\justificacio{El conjunt subjacent d'un límit de grups és el límit dels conjunts subjacents (famílies compatibles). En canvi, el coproducte de grups és el producte lliure, que no és la unió disjunta dels conjunts subjacents. Més endavant es veurà que $U$ és un adjunt per la dreta.}

\pregunta En una categoria localment petita $\cat{C}$, el functor representable
\[
\Hom(A,-)\colon\cat{C}\to\cat{Set}
\]
sempre:
\begin{enumerate}
  \opcio{a}{no és un functor.}
  \opcio{b}{preserva els límits que existeixen.}
  \opcio{c}{envia límits a colímits.}
  \opcio{d}{és constant.}
\end{enumerate}
\feedback
\justificacio{Les fletxes $A\to\varprojlim D$ es corresponen bijectivament amb els cons amb vèrtex $A$, que són les famílies compatibles de $\Hom(A,D_i)$, és a dir, els elements de $\varprojlim\Hom(A,D_i)$. L'opció (c) correspon al representable contravariant: $\Hom(-,B)$ envia colímits a límits.}

\pregunta A $\cat{Set}$, siguin $f\colon A\to B$ i $y\colon Y\to B$. Què fa la reindexació $f^{*}$?
\begin{enumerate}
  \opcio{a}{Envia $Y\to B$ al coproducte $A+Y$.}
  \opcio{b}{Envia $Y\to B$ al pullback $A\times_B Y\to A$.}
  \opcio{c}{Envia $Y\to B$ al coequalitzador de $f$ i $y$.}
  \opcio{d}{Només està definida quan $f$ és un isomorfisme.}
\end{enumerate}
\feedback
\justificacio{La fibra de $A\times_B Y\to A$ sobre $a$ està en bijecció amb la fibra $y^{-1}(f(a))$: la reindexació és la substitució al llarg de $f$.}

\pregunta En el diagrama següent, quina propietat universal fa de $P$ el \textbf{pullback} de $f$ i $g$?
\[
\begin{tikzpicture}[baseline=(m.center)]
  \matrix (m) [matrix of math nodes, row sep=2.6em, column sep=2.6em]
    { P & A \\ B & C \\ };
  \draw[-{Stealth[scale=1.1]}] (m-1-1) -- node[above] {$p_1$} (m-1-2);
  \draw[-{Stealth[scale=1.1]}] (m-1-1) -- node[left] {$p_2$} (m-2-1);
  \draw[-{Stealth[scale=1.1]}] (m-1-2) -- node[right] {$f$} (m-2-2);
  \draw[-{Stealth[scale=1.1]}] (m-2-1) -- node[below] {$g$} (m-2-2);
\end{tikzpicture}
\]
\begin{enumerate}
  \opcio{a}{Que per a tot objecte $X$ amb fletxes $x_1\colon X\to A$, $x_2\colon X\to B$ tals que $f\comp x_1=g\comp x_2$, hi ha un únic $u\colon X\to P$ amb $p_1\comp u=x_1$ i $p_2\comp u=x_2$.}
  \opcio{b}{Que $f\comp p_1=g\comp p_2$ i res més.}
  \opcio{c}{Que $P=A\times B$ sempre.}
  \opcio{d}{Que $f$ i $g$ són isomorfismes.}
\end{enumerate}
\feedback
\justificacio{És la propietat universal del pullback: el quadrat commuta i qualsevol altre quadrat commutatiu sobre $f$ i $g$ hi factoritza de manera única. L'opció (b) només diu que el quadrat commuta, i això no basta.}

\end{enumerate}
\clauestatica
\finalitzaTest

\chapter[Adjuncions]{Adjuncions}

Aquest test recull els conceptes essencials del capítol. Cada pregunta té una única resposta correcta.

\iniciTest{adj}{b,c,b,d,a,b,c,d,b,c,a,d,a,a}

\begin{enumerate}

\pregunta Què vol dir que $f\dashv g$ entre dos preordres?
\begin{enumerate}
  \opcio{a}{Que $f=g$.}
  \opcio{b}{Que $f(x)\leq y$ si i només si $x\leq g(y)$, per a tot $x,y$.}
  \opcio{c}{Que $f$ i $g$ són inverses.}
  \opcio{d}{Que $f$ i $g$ són constants.}
\end{enumerate}
\feedback
\justificacio{És la definició d'adjunció entre preordres: $f$ és adjunt per l'esquerra de $g$. No cal que siguin inverses: $f(x)\leq y\Leftrightarrow x\leq g(y)$ és més feble que $g\comp f=\id$ i $f\comp g=\id$.}

\pregunta Fixada una fórmula $A$, quina parella d'operacions forma una adjunció al preordre $\cat{Prov}_T$ d'un sistema proposicional?
\begin{enumerate}
  \opcio{a}{La negació amb ella mateixa.}
  \opcio{b}{La disjunció amb la conjunció.}
  \opcio{c}{$A\wedge-$ (esquerra) amb $A\to-$ (dreta).}
  \opcio{d}{La implicació amb la negació.}
\end{enumerate}
\feedback
\justificacio{$A\wedge B\vdash_T C$ si i només si $B\vdash_T A\to C$: és la regla de la deducció (introducció i eliminació de la implicació). Les altres parelles no satisfan cap equivalència d'aquesta forma.}

\pregunta Una \textbf{adjunció} $F\dashv G$ entre categories és:
\begin{enumerate}
  \opcio{a}{una igualtat de functors.}
  \opcio{b}{una bijecció natural $\Hom(F(A),B)\cong\Hom(A,G(B))$.}
  \opcio{c}{un isomorfisme de categories.}
  \opcio{d}{una transformació natural qualsevol.}
\end{enumerate}
\feedback
\justificacio{És la definició: una família de bijeccions $\theta_{A,B}$, natural en $A$ i en $B$. Un isomorfisme de categories és un cas molt particular, i una transformació natural qualsevol no relaciona homconjunts.}

\pregunta Sigui $\cat{I}$ una categoria petita i suposem que $\cat{C}$ té tots els límits i colímits de forma $\cat{I}$. Quina és la relació entre el functor diagonal $\Delta\colon\cat{C}\to\cat{C}^{\cat{I}}$, el límit i el colímit?
\begin{enumerate}
  \opcio{a}{$\varprojlim_{\cat{I}}\dashv\Delta\dashv\colimop_{\cat{I}}$.}
  \opcio{b}{$\Delta\dashv\colimop_{\cat{I}}$ i $\Delta\dashv\varprojlim_{\cat{I}}$.}
  \opcio{c}{El límit i el colímit són tots dos adjunts per l'esquerra de $\Delta$.}
  \opcio{d}{$\colimop_{\cat{I}}\dashv\Delta\dashv\varprojlim_{\cat{I}}$.}
\end{enumerate}
\feedback
\justificacio{Un morfisme $\Delta A\Rightarrow D$ és un con, que correspon a una fletxa $A\to\varprojlim D$; un morfisme $D\Rightarrow\Delta A$ és un cocon, que correspon a una fletxa $\colimop D\to A$.}

\pregunta La \textbf{unitat} d'una adjunció $F\dashv G$ és una transformació natural:
\begin{enumerate}
  \opcio{a}{$\id\Rightarrow G F$.}
  \opcio{b}{$F\Rightarrow G$.}
  \opcio{c}{$F G\Rightarrow\id$.}
  \opcio{d}{$G\Rightarrow F$.}
\end{enumerate}
\feedback
\justificacio{$\eta_A=\theta(\id_{F(A)})\colon A\to G(F(A))$. L'opció (c) és la counitat. Les opcions (b) i (d) no tenen sentit en general, perquè $F$ i $G$ van en sentits contraris i no comparteixen domini.}

\pregunta La \textbf{counitat} d'una adjunció $F\dashv G$ és:
\begin{enumerate}
  \opcio{a}{$\id\Rightarrow G F$.}
  \opcio{b}{$F G\Rightarrow\id$.}
  \opcio{c}{una igualtat.}
  \opcio{d}{un objecte terminal.}
\end{enumerate}
\feedback
\justificacio{$\varepsilon_B=\theta^{-1}(\id_{G(B)})\colon F(G(B))\to B$. L'opció (a) és la unitat. La counitat és una transformació natural, no una igualtat ni un objecte.}

\pregunta El transposat de $\varphi\colon F(A)\to B$ sota l'adjunció és:
\begin{enumerate}
  \opcio{a}{$\varepsilon_B\comp F(\varphi)$.}
  \opcio{b}{$\varphi$ mateix.}
  \opcio{c}{$G(\varphi)\comp\eta_A$.}
  \opcio{d}{$\eta_A\comp G(\varphi)$.}
\end{enumerate}
\feedback
\justificacio{Comprova els tipus: amb $\varphi\colon F(A)\to B$, $G(\varphi)\colon G(F(A))\to G(B)$, i compondre amb $\eta_A\colon A\to G(F(A))$ dona $A\to G(B)$. L'opció (a) és la fórmula del transposat \emph{invers}, per a $\psi\colon A\to G(B)$; l'opció (d) no és componible.}

\pregunta Les \textbf{identitats triangulars} relacionen:
\begin{enumerate}
  \opcio{a}{dos objectes terminals.}
  \opcio{b}{dos límits.}
  \opcio{c}{dues categories isomorfes.}
  \opcio{d}{la unitat i la counitat, i caracteritzen l'adjunció.}
\end{enumerate}
\feedback
\justificacio{Són $\varepsilon_{F(A)}\comp F(\eta_A)=\id_{F(A)}$ i $G(\varepsilon_B)\comp\eta_{G(B)}=\id_{G(B)}$. A més, dues transformacions naturals que les satisfan defineixen una adjunció: és la segona presentació.}

\pregunta Sigui $U\colon\cat{Mon}\to\cat{Set}$ el functor oblit. Quin fet mostra que $U$ no preserva tots els colímits?
\begin{enumerate}
  \opcio{a}{No preserva productes.}
  \opcio{b}{El monoide trivial és inicial a $\cat{Mon}$, però el seu conjunt subjacent no és $\varnothing$.}
  \opcio{c}{$\cat{Mon}$ no té objecte inicial.}
  \opcio{d}{Cap adjunt per la dreta no pot preservar colímits.}
\end{enumerate}
\feedback
\justificacio{L'objecte inicial és el colímit del diagrama buit. L'objecte inicial de $\cat{Set}$ és $\varnothing$, però $U$ envia l'inicial de $\cat{Mon}$ a un conjunt d'un element; per tant, no preserva aquest colímit.}

\pregunta Els adjunts són únics:
\begin{enumerate}
  \opcio{a}{de manera estricta (igualtat).}
  \opcio{b}{mai.}
  \opcio{c}{llevat d'isomorfisme natural.}
  \opcio{d}{només en preordres.}
\end{enumerate}
\feedback
\justificacio{Si $F\dashv G$ i $F\dashv G'$, aleshores $\Hom(-,G(B))\cong\Hom(-,G'(B))$ naturalment, i Yoneda dona $G(B)\cong G'(B)$, naturalment en $B$. No hi ha igualtat estricta: qualsevol functor isomorf a un adjunt també ho és.}

\pregunta Què diu el teorema \enquote{RAPL}?
\begin{enumerate}
  \opcio{a}{Que els adjunts per la dreta preserven límits.}
  \opcio{b}{Que tot functor és adjunt.}
  \opcio{c}{Que els límits no existeixen.}
  \opcio{d}{Que els adjunts per l'esquerra preserven límits.}
\end{enumerate}
\feedback
\justificacio{RAPL abreuja \emph{right adjoints preserve limits}. L'opció (d) confon els costats: els adjunts per l'esquerra preserven colímits.}

\pregunta Per què el functor graf subjacent $U\colon\cat{Cat}\to\cat{Graph}$ preserva límits?
\begin{enumerate}
  \opcio{a}{Perquè és ple.}
  \opcio{b}{Perquè és adjunt per l'esquerra.}
  \opcio{c}{Perquè tots els functors preserven límits.}
  \opcio{d}{Perquè és adjunt per la dreta de $\Free$.}
\end{enumerate}
\feedback
\justificacio{$\Free\dashv U$, i els adjunts per la dreta preserven límits.}

\pregunta L'adjunció \enquote{producte-exponencial} $-\times B\dashv(-)^B$ correspon a:
\begin{enumerate}
  \opcio{a}{la currificació $\Hom(A\times B,C)\cong\Hom(A,C^B)$.}
  \opcio{b}{la dualitat.}
  \opcio{c}{la negació.}
  \opcio{d}{la composició de functors.}
\end{enumerate}
\feedback
\justificacio{Fixat $B$, una fletxa $A\times B\to C$ correspon a una fletxa $A\to C^B$, naturalment en $A$ i en $C$. És exactament la bijecció de l'adjunció $-\times B\dashv(-)^B$.}

\pregunta Si $f^{*}$ admet adjunts a tots dos costats, quin paper tenen $\exists_f$ i $\forall_f$?
\begin{enumerate}
  \opcio{a}{Com l'adjunt per l'esquerra i l'adjunt per la dreta de $f^{*}$, respectivament.}
  \opcio{b}{Com objectes terminals.}
  \opcio{c}{Com functors constants.}
  \opcio{d}{Com isomorfismes de categories.}
\end{enumerate}
\feedback
\justificacio{A $\cat{Set}$, $\exists_f$ és la imatge i $\forall_f(S)$ el conjunt dels $y$ amb fibra continguda en $S$; per a una projecció, són la quantificació existencial i la universal.}

\end{enumerate}
\clauestatica
\finalitzaTest

\chapter[Cartesianes tancades]{\texorpdfstring{Categories cartesianes\\ tancades}{Categories cartesianes tancades}}

Aquest test recull els conceptes essencials del capítol. Cada pregunta té una única resposta correcta.

\iniciTest{cct}{a,c,d,b,c,a,d,a,b,c,d,b}

\begin{enumerate}

\pregunta Una categoria és \textbf{cartesiana} quan té:
\begin{enumerate}
  \opcio{a}{objecte terminal i productes binaris (és a dir, tots els productes finits).}
  \opcio{b}{tots els colímits.}
  \opcio{c}{només un objecte inicial.}
  \opcio{d}{exponencials però no productes.}
\end{enumerate}
\feedback
\justificacio{El terminal és el producte buit, i els productes de qualsevol família finita s'obtenen per inducció a partir dels binaris. Els colímits i els exponencials no formen part de la definició.}

\pregunta Quina bijecció natural expressa la propietat universal de l'exponencial $C^{B}$?
\begin{enumerate}
  \opcio{a}{$\Hom(A,C)\cong\Hom(B,C)$.}
  \opcio{b}{$\Hom(A+B,C)\cong\Hom(A,C)\times\Hom(B,C)$.}
  \opcio{c}{$\Hom(A\times B,C)\cong\Hom(A,C^{B})$.}
  \opcio{d}{$\Hom(A,C\times B)\cong\Hom(A,C^{B})$.}
\end{enumerate}
\feedback
\justificacio{És la currificació $f\mapsto\lambda f$, natural en $A$ i en $C$: expressa l'adjunció $-\times B\dashv(-)^{B}$.}

\pregunta En una CCC, quin morfisme indueix $f\colon B'\to B$ entre exponencials amb el mateix $C$?
\begin{enumerate}
  \opcio{a}{$C^{B'}\to C^{B}$, perquè l'exponencial és covariant en totes dues variables.}
  \opcio{b}{Cap: l'exponencial no depèn functorialment de $B$.}
  \opcio{c}{$B^{C}\to B'^{C}$.}
  \opcio{d}{$C^{B}\to C^{B'}$: l'exponencial és contravariant en l'exponent.}
\end{enumerate}
\feedback
\justificacio{$C^{f}$ és la transposada de $C^{B}\times B'\xrightarrow{\id\times f}C^{B}\times B\xrightarrow{\ev}C$.}

\pregunta La \textbf{currificació} $\lambda f$ d'un morfisme $f\colon A\times B\to C$ és l'únic morfisme que compleix:
\begin{enumerate}
  \opcio{a}{$\lambda f\comp\ev=f$.}
  \opcio{b}{$\ev\comp(\lambda f\times\id_B)=f$.}
  \opcio{c}{$\lambda f=f$.}
  \opcio{d}{$\lambda f\comp\pi_1=f$.}
\end{enumerate}
\feedback
\justificacio{És el triangle de la definició: currificar i després avaluar torna $f$. L'opció (a) no és ni tan sols componible, perquè $\ev$ surt de $C^B\times B$ i $\lambda f$ hi arriba només des de $A$.}

\pregunta Una categoria és \textbf{cartesiana tancada} quan és cartesiana i, a més, per a cada objecte $B$:
\begin{enumerate}
  \opcio{a}{$B$ és terminal.}
  \opcio{b}{$-\times B$ és un isomorfisme.}
  \opcio{c}{el functor $-\times B$ té adjunt per la dreta $(-)^{B}$.}
  \opcio{d}{$-\times B$ té adjunt per l'esquerra.}
\end{enumerate}
\feedback
\justificacio{Tancar cartesianament vol dir que la currificació és sempre disponible: per a cada $B$, $-\times B\dashv(-)^B$. Un adjunt per l'esquerra de $-\times B$ seria una altra cosa, i no existeix en general.}

\pregunta A $\cat{Set}$, l'exponencial $C^{B}$ és:
\begin{enumerate}
  \opcio{a}{el conjunt de \emph{totes} les aplicacions $B\to C$.}
  \opcio{b}{el producte cartesià $C\times B$.}
  \opcio{c}{la unió disjunta $C+B$.}
  \opcio{d}{el conjunt de bijeccions $B\to C$.}
\end{enumerate}
\feedback
\justificacio{L'avaluació és $\ev(g,b)=g(b)$, i la currificació de $f\colon A\times B\to C$ és $a\mapsto(b\mapsto f(a,b))$. Cal prendre totes les aplicacions, no només les bijeccions, perquè $\lambda f(a)$ pot ser qualsevol aplicació.}

\pregunta En una \textbf{àlgebra de Heyting}, l'objecte exponencial $b\Rightarrow c$ és:
\begin{enumerate}
  \opcio{a}{el suprem $b\vee c$.}
  \opcio{b}{la negació $\neg b$.}
  \opcio{c}{l'ínfim $b\wedge c$.}
  \opcio{d}{l'adjunt per la dreta de $-\wedge b$, caracteritzat per $a\wedge b\leq c\Leftrightarrow a\leq(b\Rightarrow c)$.}
\end{enumerate}
\feedback
\justificacio{En un preordre, l'exponencial és l'adjunt per la dreta del producte, que és l'ínfim. En una àlgebra de Boole, a més, $b\Rightarrow c=\neg b\vee c$, però aquesta fórmula no val en una àlgebra de Heyting qualsevol.}

\pregunta Quina categoria \emph{no} és, en general, cartesiana tancada?
\begin{enumerate}
  \opcio{a}{$\cat{Top}$, els espais topològics.}
  \opcio{b}{una àlgebra de Heyting.}
  \opcio{c}{$\cat{Set}$.}
  \opcio{d}{una categoria de prefeixos $\cat{Set}^{\cat{C}\op}$.}
\end{enumerate}
\feedback
\justificacio{A $\cat{Top}$, l'avaluació no sempre té un espai de funcions contínues universal. $\cat{Set}$, les àlgebres de Heyting com a preordres i les categories de prefeixos sí que són cartesianes tancades.}

\pregunta En la semàntica del càlcul lambda simplement tipat, el \textbf{tipus funció} $\sigma\to\tau$ s'interpreta amb:
\begin{enumerate}
  \opcio{a}{el producte $\llbracket\sigma\rrbracket\times\llbracket\tau\rrbracket$.}
  \opcio{b}{l'exponencial $\llbracket\tau\rrbracket^{\llbracket\sigma\rrbracket}$.}
  \opcio{c}{l'objecte terminal.}
  \opcio{d}{el coproducte.}
\end{enumerate}
\feedback
\justificacio{Un terme de tipus $\sigma\to\tau$ és una funció, i en una CCC les funcions de $\llbracket\sigma\rrbracket$ a $\llbracket\tau\rrbracket$ formen l'exponencial. El producte interpreta, en canvi, el tipus parella $\sigma\times\tau$.}

\pregunta En aquesta semàntica, l'\textbf{abstracció} $\lambda x.\,t$ es correspon amb:
\begin{enumerate}
  \opcio{a}{una projecció.}
  \opcio{b}{la diagonal.}
  \opcio{c}{la currificació de la interpretació de $t$.}
  \opcio{d}{l'objecte terminal.}
\end{enumerate}
\feedback
\justificacio{Si $t$ té tipus $\tau$ en el context $\Gamma,x{:}\sigma$, s'interpreta com un morfisme $\llbracket\Gamma\rrbracket\times\llbracket\sigma\rrbracket\to\llbracket\tau\rrbracket$. La seva currificació és un morfisme $\llbracket\Gamma\rrbracket\to\llbracket\tau\rrbracket^{\llbracket\sigma\rrbracket}$, que interpreta $\lambda x.\,t$.}

\pregunta La regla de conversió $(\beta)$ del càlcul lambda es correspon categòricament amb:
\begin{enumerate}
  \opcio{a}{la unicitat de l'objecte terminal.}
  \opcio{b}{l'associativitat del producte.}
  \opcio{c}{la commutativitat del coproducte.}
  \opcio{d}{l'equació de l'avaluació $\ev\comp(\lambda f\times\id_B)=f$, de la qual es dedueix.}
\end{enumerate}
\feedback
\justificacio{Aplicar una abstracció a un argument és avaluar una currificació, i l'equació $\ev\comp(\lambda f\times\id)=f$ diu que el resultat és $f$, que és el que afirma la regla $(\beta)$. La regla $(\eta)$ correspon a la segona llei, $\lambda(\ev\comp(g\times\id))=g$.}

\pregunta La correspondència de \textbf{Curry--Howard--Lambek} relaciona:
\begin{enumerate}
  \opcio{a}{grups amb anells.}
  \opcio{b}{el càlcul lambda simplement tipat amb les categories cartesianes tancades.}
  \opcio{c}{límits amb colímits.}
  \opcio{d}{monoides amb preordres.}
\end{enumerate}
\feedback
\justificacio{Els tipus s'interpreten com a objectes, els termes com a morfismes i les conversions com a igualtats. En la lectura lògica, les proposicions del fragment $(\top,\wedge,\to)$ són els tipus, i les demostracions, els termes.}

\end{enumerate}
\clauestatica
\finalitzaTest

\chapter[Subobjectes, imatges i factoritzacions]{\texorpdfstring{Subobjectes, imatges\\ i factoritzacions}{Subobjectes, imatges i factoritzacions}}

Aquest test recull els conceptes essencials del capítol. Cada pregunta té una única resposta correcta.

\iniciTest{sub}{c,b,d,b,c,d,a,c,a,d,c,b,d,a,b,a,c}

\begin{enumerate}

\pregunta Un \textbf{subobjecte} d'un objecte $A$ és:
\begin{enumerate}
  \opcio{a}{qualsevol morfisme cap a $A$.}
  \opcio{b}{un element de $A$.}
  \opcio{c}{una classe d'equivalència de monomorfismes amb codomini $A$.}
  \opcio{d}{un epimorfisme des de $A$.}
\end{enumerate}
\feedback
\justificacio{Dos monomorfismes que difereixen en un isomorfisme dels dominis representen el mateix subobjecte. Un morfisme qualsevol no basta: cal que sigui mono, el paper que a $\cat{Set}$ fan les injeccions.}

\pregunta Sigui $H$ un subgraf d'un graf $G$ i $F\colon G'\to G$ un morfisme de grafs. Què és la reindexació $F^{*}H$?
\begin{enumerate}
  \opcio{a}{La imatge de $H$ per $F$.}
  \opcio{b}{El subgraf de $G'$ format per les antiimatges dels vèrtexs i de les arestes de $H$.}
  \opcio{c}{El graf $H$ mateix, vist dins de $G'$.}
  \opcio{d}{El producte de grafs $H\times G'$.}
\end{enumerate}
\feedback
\justificacio{La reindexació es calcula component a component: $F^{*}(V',E')=(F_V^{-1}(V'),F_E^{-1}(E'))$, que és tancat per origen i destí perquè $F$ els preserva.}

\pregunta A $\cat{Set}$, l'ordre parcial $\Sub(A)$ és isomorf a:
\begin{enumerate}
  \opcio{a}{el conjunt $A$.}
  \opcio{b}{el monoide de les aplicacions $A\to A$.}
  \opcio{c}{el grup de permutacions de $A$.}
  \opcio{d}{el conjunt de parts $\mathcal{P}(A)$, ordenat per inclusió.}
\end{enumerate}
\feedback
\justificacio{Dos monomorfismes cap a $A$ són equivalents si i només si tenen la mateixa imatge, i la imatge determina la classe. L'ordre correspon a la inclusió de les imatges.}

\pregunta La \textbf{parella nucli} d'una fletxa $f\colon A\to B$, quan existeix, s'obté com:
\begin{enumerate}
  \opcio{a}{l'equalitzador de $f$ i $\id_A$.}
  \opcio{b}{el pullback de $f$ amb ella mateixa.}
  \opcio{c}{el coequalitzador de $f$ amb ella mateixa.}
  \opcio{d}{la factorització imatge de $f$.}
\end{enumerate}
\feedback
\justificacio{És el parell de projeccions $p_1,p_2\colon R\to A$ del pullback de $f$ amb $f$; a $\cat{Set}$, $R=\{(a,a')\mid f(a)=f(a')\}$. Quan existeix la parella nucli, un epimorfisme regular n'és el coequalitzador.}

\pregunta Lògicament, l'ordre $[m]\leq[n]$ a $\Sub(A)$ s'interpreta com:
\begin{enumerate}
  \opcio{a}{la negació del predicat.}
  \opcio{b}{la igualtat de predicats.}
  \opcio{c}{la implicació entre predicats sobre $A$.}
  \opcio{d}{la disjunció de predicats.}
\end{enumerate}
\feedback
\justificacio{$[m]\leq[n]$ vol dir que $m$ factoritza a través de $n$: tot el que satisfà el predicat $m$ satisfà el predicat $n$. És l'ordre de la deduïbilitat entre predicats sobre $A$.}

\pregunta El \textbf{pullback} d'un monomorfisme al llarg de qualsevol morfisme és:
\begin{enumerate}
  \opcio{a}{un epimorfisme.}
  \opcio{b}{un isomorfisme.}
  \opcio{c}{mai un monomorfisme.}
  \opcio{d}{sempre un monomorfisme.}
\end{enumerate}
\feedback
\justificacio{Si $m'\comp u=m'\comp v$ al pullback, aleshores $u$ i $v$ tenen la mateixa composició amb totes dues projeccions (la segona, per la cancel·lació de $m$), i per la unicitat del pullback $u=v$.}

\pregunta El \textbf{functor de reindexació} $f^{*}\colon\Sub(B)\to\Sub(A)$ associat a $f\colon A\to B$ envia un subobjecte al seu:
\begin{enumerate}
  \opcio{a}{pullback al llarg de $f$.}
  \opcio{b}{coproducte amb $A$.}
  \opcio{c}{producte amb $B$.}
  \opcio{d}{imatge directa.}
\end{enumerate}
\feedback
\justificacio{És la definició: el pullback d'un mono és mono, i en resulta una aplicació monòtona ben definida entre classes. La imatge directa (d) correspon a l'existencial, que va en sentit contrari.}

\pregunta A $\cat{Set}$, la reindexació $f^{*}(S)$ d'un subconjunt $S\subseteq B$ és:
\begin{enumerate}
  \opcio{a}{la imatge directa $f(S)$.}
  \opcio{b}{el complement de $S$.}
  \opcio{c}{l'antiimatge $f^{-1}(S)$.}
  \opcio{d}{el producte $S\times A$.}
\end{enumerate}
\feedback
\justificacio{El pullback de la inclusió $S\hookrightarrow B$ al llarg de $f$ és $\{a\in A\mid f(a)\in S\}$. La imatge directa (a) és l'existencial $\exists_f$, l'adjunt per l'esquerra de $f^{*}$.}

\pregunta Lògicament, la reindexació $f^{*}$ interpreta:
\begin{enumerate}
  \opcio{a}{la substitució d'una variable per un terme.}
  \opcio{b}{la quantificació existencial.}
  \opcio{c}{la negació.}
  \opcio{d}{la conjunció.}
\end{enumerate}
\feedback
\justificacio{Si $\varphi(y)$ és un predicat sobre $B$ i $f$ interpreta un terme, $f^{*}\varphi$ és $\varphi(f(x))$: substituir és fer el pullback. La quantificació existencial correspon a la imatge.}

\pregunta La \textbf{intersecció} de dos subobjectes $[m]\wedge[n]$ a $\Sub(A)$ es calcula com:
\begin{enumerate}
  \opcio{a}{el coproducte de $S$ i $T$.}
  \opcio{b}{la unió de les imatges.}
  \opcio{c}{l'exponencial $T^{S}$.}
  \opcio{d}{el pullback de $m$ i $n$ sobre $A$.}
\end{enumerate}
\feedback
\justificacio{Una fletxa factoritza a través de $m$ i de $n$ si i només si factoritza a través del seu pullback. Per tant, el pullback és el major subobjecte per sota de tots dos, i interpreta la conjunció.}

\pregunta En una categoria regular, considerem un quadrat cartesià amb $h\colon A'\to A$, $f'\colon A'\to B'$, $f\colon A\to B$ i $g\colon B'\to B$, de manera que $f\comp h=g\comp f'$. Quina igualtat és la condició de Beck--Chevalley, amb els tipus correctes?
\begin{enumerate}
  \opcio{a}{$\exists_f\comp g^{*}=h^{*}\comp\exists_{f'}$.}
  \opcio{b}{$h^{*}\comp\exists_f=\exists_{f'}\comp g^{*}\colon\Sub(B)\to\Sub(A')$.}
  \opcio{c}{$g^{*}\comp\exists_f=\exists_{f'}\comp h^{*}\colon\Sub(A)\to\Sub(B')$.}
  \opcio{d}{$g^{*}\comp\exists_f=\exists_{f'}\comp h^{*}\colon\Sub(B)\to\Sub(A')$.}
\end{enumerate}
\feedback
\justificacio{Partint de $\Sub(A)$, hi ha dos camins fins a $\Sub(B')$: quantificar al llarg de $f$ i substituir al llarg de $g$, o substituir al llarg de $h$ i quantificar al llarg de $f'$. La condició diu que coincideixen: quantificar i substituir commuten. A (a) i (b) les composicions no existeixen ($\exists_f$ surt de $\Sub(A)$, no de $\Sub(B)$), i a (d) el domini i el codomini estan intercanviats.}

\pregunta Segons el capítol, en una categoria \textbf{regular} la reindexació $f^{*}$ té garantit:
\begin{enumerate}
  \opcio{a}{cap adjunt.}
  \opcio{b}{un adjunt per l'esquerra $\exists_f$ (l'universal $\forall_f$ requereix estructura addicional).}
  \opcio{c}{un adjunt per la dreta $\forall_f$, però no per l'esquerra.}
  \opcio{d}{adjunts per tots dos costats, sempre.}
\end{enumerate}
\feedback
\justificacio{L'existencial és la imatge, i la regularitat en garanteix l'existència i l'estabilitat. L'adjunt per la dreta $\forall_f$ és estructura addicional, que tenen, per exemple, els topos.}

\pregunta Una \textbf{factorització imatge} de $f$ descompon $f$ com:
\begin{enumerate}
  \opcio{a}{un mono seguit d'un epi qualsevol.}
  \opcio{b}{dos isomorfismes.}
  \opcio{c}{dos monomorfismes.}
  \opcio{d}{un epimorfisme regular seguit d'un monomorfisme.}
\end{enumerate}
\feedback
\justificacio{A $\cat{Set}$, l'epi regular és l'aplicació exhaustiva sobre la imatge, i el mono és la inclusió de la imatge. La factorització és única llevat d'un únic isomorfisme compatible.}

\pregunta La \textbf{imatge} $\Imatge(f)$ es caracteritza com:
\begin{enumerate}
  \opcio{a}{el menor subobjecte del codomini a través del qual $f$ es factoritza.}
  \opcio{b}{el major subobjecte del domini.}
  \opcio{c}{el nucli de $f$.}
  \opcio{d}{el coproducte de domini i codomini.}
\end{enumerate}
\feedback
\justificacio{La imatge és un subobjecte pel qual factoritza $f$, i és per sota de qualsevol altre que tingui aquesta propietat. A $\cat{Set}$ és el subconjunt $f(A)\subseteq B$.}

\pregunta Una categoria és \textbf{regular} si té límits finits, factoritzacions imatge i, a més:
\begin{enumerate}
  \opcio{a}{tots els colímits.}
  \opcio{b}{els epimorfismes regulars són estables per pullback.}
  \opcio{c}{és cartesiana tancada.}
  \opcio{d}{tots els objectes són projectius.}
\end{enumerate}
\feedback
\justificacio{És la condició que fa estables les imatges per reindexació i, per tant, l'existencial compatible amb la substitució. Els colímits, el tancament cartesià i els objectes projectius no formen part de la definició.}

\pregunta En una categoria regular, la \textbf{quantificació existencial} $\exists_f[s]$ es defineix com:
\begin{enumerate}
  \opcio{a}{la imatge $\Imatge(f\comp s)$.}
  \opcio{b}{el pullback de $[s]$.}
  \opcio{c}{el complement de $[s]$.}
  \opcio{d}{l'exponencial de $[s]$.}
\end{enumerate}
\feedback
\justificacio{Es fa la composició $f\comp s$ i se'n pren la imatge a $B$; a $\cat{Set}$, $f(S)$. Aquest operador és l'adjunt per l'esquerra de $f^{*}$: $\exists_f[s]\leq[t]$ si i només si $[s]\leq f^{*}[t]$.}

\pregunta Sigui $\cat{C}$ una categoria regular i $f\colon A\to B$ un morfisme. Quina d'aquestes afirmacions \emph{no} es pot deduir només de la regularitat?
\begin{enumerate}
  \opcio{a}{Tota fletxa té una factorització imatge (única llevat d'iso).}
  \opcio{b}{La reindexació $f^{*}\colon\Sub(B)\to\Sub(A)$ té un adjunt per l'esquerra $\exists_f$.}
  \opcio{c}{Cada $\Sub(A)$ té unions binàries, i $f^{*}$ les preserva.}
  \opcio{d}{Cada $\Sub(A)$ té ínfims binaris, calculats per pullback, i $f^{*}$ els preserva.}
\end{enumerate}
\feedback
\justificacio{Les unions de subobjectes estables per reindexació no formen part de la definició de regularitat: són l'estructura de les categories coherents. El mateix passa amb l'adjunt per la dreta $\forall_f$. En canvi, (a) i (b) són part del que dona la regularitat, i (d) també val: els ínfims són pullbacks, i $f^{*}$ els preserva perquè és adjunt per la dreta de $\exists_f$.}

\end{enumerate}
\clauestatica
\finalitzaTest

\chapter[Classificadors de subobjectes i topos]{\texorpdfstring{Classificadors de subobjectes i introducció als topos}{Classificadors de subobjectes i introducció als topos}}

Aquest test recull els conceptes essencials del capítol. Cada pregunta té una única resposta correcta.

\iniciTest{top}{c,b,d,a,c,d,b,a,a,b,c,d,b}

\begin{enumerate}

\pregunta Sigui $m\colon U\rightarrowtail E$ un subobjecte classificat per $\chi_m\colon E\to\Omega$. Un element generalitzat $x\colon X\to E$ factoritza a través de $m$ exactament quan:
\begin{enumerate}
  \opcio{a}{$\chi_m\comp x$ és un isomorfisme.}
  \opcio{b}{$x$ és un monomorfisme.}
  \opcio{c}{$\chi_m\comp x=\mathsf{true}\comp\,!_X$.}
  \opcio{d}{$X$ és l'objecte terminal.}
\end{enumerate}
\feedback
\justificacio{És la propietat universal del quadrat cartesià de $m$: la pertinença d'un element generalitzat esdevé una igualtat de morfismes cap a $\Omega$.}

\pregunta A la categoria de prefeixos sobre $\cat{C}$, siguin $h\colon D\to C$ i $S$ un sedàs sobre $C$. Quin sedàs sobre $D$ és $\Omega(h)(S)$?
\begin{enumerate}
  \opcio{a}{$\{h\comp g\mid g\in S\}$.}
  \opcio{b}{$\{g\colon X\to D\mid h\comp g\in S\}$.}
  \opcio{c}{$S$ mateix.}
  \opcio{d}{El sedàs màxim sobre $D$.}
\end{enumerate}
\feedback
\justificacio{$\Omega$ actua per reindexació de sedassos: $\Omega(h)(S)=h^{*}S$, i aquest conjunt torna a ser un sedàs.}

\pregunta La propietat universal del classificador diu que cada subobjecte $U\rightarrowtail E$ determina:
\begin{enumerate}
  \opcio{a}{un isomorfisme $E\cong\Omega$.}
  \opcio{b}{cap morfisme.}
  \opcio{c}{diversos morfismes $E\to\Omega$.}
  \opcio{d}{un únic morfisme $\chi\colon E\to\Omega$ que fa un quadrat cartesià amb $\mathsf{true}$.}
\end{enumerate}
\feedback
\justificacio{És la definició de classificador: existència i unicitat de $\chi$, amb $U$ com a pullback de $\mathsf{true}$ al llarg de $\chi$. Que el quadrat només commuti no basta per a la unicitat.}

\pregunta L'existència d'un classificador de subobjectes equival a dir que el functor $\Sub(-)$ és:
\begin{enumerate}
  \opcio{a}{representable, amb objecte representant $\Omega$.}
  \opcio{b}{constant.}
  \opcio{c}{no functorial.}
  \opcio{d}{un isomorfisme.}
\end{enumerate}
\feedback
\justificacio{La bijecció $\Sub(E)\cong\Hom(E,\Omega)$, natural en $E$, és exactament una representació de $\Sub$, amb element universal $\mathsf{true}$.}

\pregunta El classificador de subobjectes, quan existeix, és únic:
\begin{enumerate}
  \opcio{a}{només a $\cat{Set}$.}
  \opcio{b}{mai.}
  \opcio{c}{llevat d'un únic isomorfisme compatible amb $\mathsf{true}$ (per Yoneda).}
  \opcio{d}{llevat d'homotopia.}
\end{enumerate}
\feedback
\justificacio{Dos objectes que representen el mateix functor són isomorfs per un únic isomorfisme compatible amb els elements universals: és el principi de Yoneda aplicat a $\Sub$.}

\pregunta Un \textbf{topos elemental} és una categoria amb límits finits, exponencials i:
\begin{enumerate}
  \opcio{a}{tots els colímits infinits.}
  \opcio{b}{un objecte inicial estricte.}
  \opcio{c}{una estructura de grup.}
  \opcio{d}{un classificador de subobjectes.}
\end{enumerate}
\feedback
\justificacio{Límits finits, exponencials i classificador de subobjectes: no es demanen colímits infinits. $\cat{FinSet}$, que no en té, és un topos.}

\pregunta Quina d'aquestes categories és un topos?
\begin{enumerate}
  \opcio{a}{la categoria dels grups.}
  \opcio{b}{$\cat{Set}$.}
  \opcio{c}{$\cat{Top}$.}
  \opcio{d}{un monoide no trivial vist com a categoria.}
\end{enumerate}
\feedback
\justificacio{$\cat{Set}$ té límits finits, exponencials i el classificador $\{0,1\}$. $\cat{Top}$ no és cartesiana tancada, i un monoide no trivial vist com a categoria no té ni tan sols tots els límits finits.}

\pregunta Per a $\cat{C}$ petita, la categoria de prefeixos $\cat{Set}^{\cat{C}\op}$ és:
\begin{enumerate}
  \opcio{a}{sempre un topos.}
  \opcio{b}{un topos només si $\cat{C}$ és un grup.}
  \opcio{c}{mai un topos.}
  \opcio{d}{cartesiana tancada però sense classificador.}
\end{enumerate}
\feedback
\justificacio{Té límits finits i exponencials (capítols de límits i de cartesianes tancades) i el classificador de sedassos. No cal cap condició sobre $\cat{C}$ més enllà de ser petita.}

\pregunta A $\cat{Graph}$, sigui $x\colon v\to w$ una aresta d'un graf $G$ i $H$ un subgraf que conté els vèrtexs $v$ i $w$ però no l'aresta $x$. A quina aresta del classificador $\Omega$ envia $\chi_H$ l'aresta $x$?
\begin{enumerate}
  \opcio{a}{Al bucle $\partial$ del vèrtex $1$, que registra \enquote{extrems sí, aresta no}.}
  \opcio{b}{Al bucle $\top$ del vèrtex $1$, perquè els dos extrems són a $H$.}
  \opcio{c}{A l'aresta $\sigma\colon1\to0$.}
  \opcio{d}{Enlloc: $\chi_H$ només actua sobre els vèrtexs.}
\end{enumerate}
\feedback
\justificacio{Com que $v,w\in H$, la imatge de $x$ ha d'anar de $1$ a $1$, i n'hi ha dues: $\partial$ i $\top$. Si fos $\top$, $x$ seria al pullback de $\mathsf{true}$, que ha de ser exactament $H$; per tant, va a $\partial$. L'opció (c) tindria el destí equivocat, i un morfisme de grafs actua també sobre les arestes.}

\pregunta En un topos, l'objecte de parts $\mathcal P(E)=\Omega^{E}$ satisfà, per a tot objecte $A$,
\[
\Hom(A,\mathcal P(E))\;\cong\;\Sub(A\times E).
\]
Quina interpretació expressa aquesta bijecció?
\begin{enumerate}
  \opcio{a}{Que tot objecte és isomorf al seu objecte de parts.}
  \opcio{b}{Que els morfismes cap a $\mathcal P(E)$ classifiquen les relacions entre $A$ i $E$.}
  \opcio{c}{Que $\mathcal P(E)$ és l'objecte terminal.}
  \opcio{d}{Que $E$ té exactament dos subobjectes.}
\end{enumerate}
\feedback
\justificacio{Combina la currificació $\Hom(A,\Omega^{E})\cong\Hom(A\times E,\Omega)$ amb la representabilitat de $\Sub$; per a $A=1$ recupera $\Hom(1,\mathcal P(E))\cong\Sub(E)$.}

\pregunta Quina afirmació sobre la categoria $\cat{FinSet}$ dels conjunts finits és correcta?
\begin{enumerate}
  \opcio{a}{No és un topos, perquè no té coproductes infinits.}
  \opcio{b}{És equivalent a $\cat{Set}^{\cat{C}\op}$ per a alguna categoria petita $\cat{C}$.}
  \opcio{c}{És un topos elemental que no és equivalent a cap categoria de prefeixos.}
  \opcio{d}{Només és un topos si s'hi afegeixen els conjunts infinits.}
\end{enumerate}
\feedback
\justificacio{Té límits finits, exponencials finits i el classificador $\{0,1\}$, de manera que és un topos. Tota categoria de prefeixos té colímits petits, i $\cat{FinSet}$ no té coproductes infinits; com que una equivalència conserva els colímits, no és equivalent a cap. Un topos elemental no demana colímits infinits.}

\pregunta A $\cat{Graph}$, $\Omega$ té dos vèrtexs i cinc arestes. Quin fet mostra que $\cat{Graph}$ no és un topos booleà?
\begin{enumerate}
  \opcio{a}{Que $\Omega$ no és un conjunt.}
  \opcio{b}{Que $\Omega$ té més de dos vèrtexs.}
  \opcio{c}{Que $\Hom(\mathbf L,\Omega)$ té exactament dos elements.}
  \opcio{d}{Que $\{t\}$ és el més gran dels subgrafs de $\mathbf E$ disjunts de $\{s\}$, però $\{s\}\cup\{t\}$ no és tot $\mathbf E$.}
\end{enumerate}
\feedback
\justificacio{Per a $U=\{s\}$, el seu \enquote{complement} $\neg U=\{t\}$ no cobreix, juntament amb $U$, tot $\mathbf E$: hi falta l'aresta. Falla, doncs, $U\vee\neg U=\top$. L'opció (b) és falsa ($\Omega$ té dos vèrtexs), i també la (c) (n'hi ha tres).}

\pregunta Un topos té exactament dos morfismes $1\to\Omega$. Què se'n pot concloure?
\begin{enumerate}
  \opcio{a}{Que és booleà, perquè $\Omega$ té exactament dos valors de veritat.}
  \opcio{b}{Res sobre la booleanitat: els morfismes $1\to\Omega$ només compten $\Sub(1)$, mentre que ser booleà és una condició sobre tots els $\Sub(E)$.}
  \opcio{c}{Que $\Omega\cong1+1$.}
  \opcio{d}{Que el topos és equivalent a $\cat{Set}$.}
\end{enumerate}
\feedback
\justificacio{$\Hom(1,\Omega)\cong\Sub(1)$. Els prefeixos sobre el monoide $(\mathbb N,+)$ tenen exactament dos valors globals, però els subobjectes del representable formen una cadena $\mathbb N\cup\{\infty\}$, que no és una àlgebra de Boole, i aquest $\Omega$ no és $1+1$. Tenir dos valors globals i ser booleà són coses diferents.}

\end{enumerate}
\clauestatica
\finalitzaTest

\fi
\ifimprimible\else\end{Form}\fi

% ============================================================
% APÈNDIXS NUMERATS
% \appendix ha d'anar dins \mainmatter (abans de \backmatter):
% a la classe book, \backmatter desactiva la numeració de capítols,
% de manera que els capítols d'apèndix es numeren A, B, ... aquí.
% ============================================================
\appendix
\part{Apèndixs}
\chapter[Axiomàtica relacional]{Una axiomàtica relacional de les categories petites}
\label{cap:axiomatica-relacional}

\begin{objectius}
\apartat{Prerequisits} La definició de categoria, la categoria oposada i la formulació operativa del principi de dualitat (capítol~\ref{cap:categories}, secció~\ref{sec:dualitat}). De lògica de primer ordre amb igualtat: la sintaxi (fórmules, sentències, variables lliures i lligades, substitució), la semàntica (estructures, satisfacció i models), la noció de deducció en un càlcul de Hilbert i les demostracions per inducció sobre fórmules. El teorema de completesa només s'esmenta. Les referències a la lògica lliure de la secció~\ref{sec:ar-relacional-funcional} són informatives i no es pressuposen.
\apartat{Objectius} En acabar l'apèndix, el lector hauria de poder:
\begin{itemize}
  \item presentar una categoria petita en un llenguatge de primer ordre \emph{amb una sola mena d'individus}, on tot individu és una fletxa i $\Dom$, $\Cod$ i $\Com$ són les úniques relacions primitives;
  \item distingir el llenguatge formal, els axiomes i el metallenguatge, i reconèixer quines nocions són primitives i quines en són derivades;
  \item deduir dels axiomes l'estabilitat dels extrems, la unicitat de la composició i la derivabilitat del tercer axioma;
  \item reconstruir els objectes com les fletxes identitat i justificar el pas a la notació funcional usual;
  \item explicar en quin sentit la presentació relacional és equivalent a la definició usual, fins a isomorfisme canònic;
  \item comparar la presentació relacional amb una presentació funcional amb operacions parcials;
  \item definir la traducció dual de fórmules, comprovar l'autodualitat dels axiomes i demostrar el principi de dualitat en les versions sintàctica i semàntica, relacionant-lo amb la versió operativa del capítol~\ref{cap:categories}.
\end{itemize}
\end{objectius}

La definició habitual d'una categoria (definició~\ref{def:categoria}) introdueix objectes, fletxes, identitats i composició. En aquest apèndix adoptem un punt de vista més auster: el llenguatge categòric primitiu només parla de fletxes. Allò que habitualment anomenem objectes es reconstrueix després com una propietat especial d'algunes fletxes. És, en un sentit ben literal, el programa de la introducció ---\emph{relaciona, no individualitza} (secció~\ref{sec:mirar-fora})--- dut fins a la seva conclusió sintàctica: si el que compta d'un objecte és el seu paper dins el sistema de fletxes, aleshores no cal introduir-lo com a dada primitiva.

L'objectiu no és substituir la presentació usual, que és molt més còmoda per al treball matemàtic ordinari. L'objectiu és separar amb precisió tres nivells: el \emph{llenguatge formal} de la teoria; els \emph{axiomes} formulats en aquest llenguatge; i el \emph{metallenguatge}, on interpretem les fórmules i introduïm notació derivada. Aquesta separació permet veure quines nocions són primitives i quines són conseqüència dels axiomes. En particular, la funcionalitat de la composició no s'incorpora a la sintaxi: es demostrarà com un teorema.

\begin{nota}
La idea de descriure una categoria \enquote{només amb fletxes}, amb els objectes identificats amb les identitats, és clàssica: apareix ja en Mac~Lane~\cite{maclane}, es formula amb una relació ternària de composició en la teoria elemental de categories de Lawvere~\cite{lawvere-ccfm}, i s'estudia sistemàticament, amb operacions parcials, en Scott~\cite{scott-identity} i en Freyd i Scedrov~\cite{freydscedrov}. Diversos d'aquests sistemes han estat comparats i verificats formalment per Benzmüller i Scott~\cite{BS2020,BSAFP}. La contribució d'aquest apèndix no és, doncs, la idea, sinó una presentació concreta: una axiomàtica de primer ordre \emph{purament relacional} ---fins i tot el domini i el codomini són relacions--- en què la funcionalitat de la composició i el comportament dels seus extrems esdevenen \emph{teoremes}.
\end{nota}

\section{El llenguatge formal}
\label{sec:ar-llenguatge}

Treballarem en lògica clàssica de primer ordre amb igualtat i amb una sola mena d'individus. La signatura no lògica és
\[
\mathcal L_{\mathrm{Cat}}
=
\{\Dom^{2},\Cod^{2},\Com^{3}\}.
\]
No hi ha constants ni símbols de funció. Per tant, els termes del llenguatge són només variables individuals $f,g,h,d,c,e,\dots$, i totes recorren un mateix univers.

Les fórmules atòmiques no lògiques són de les formes $\Dom(f,d)$, $\Cod(f,c)$ i $\Com(f,g,h)$. La seva lectura metalingüística serà:
\begin{align*}
\Dom(f,d)&:\quad \text{\enquote{el domini de $f$ és $d$}},\\
\Cod(f,c)&:\quad \text{\enquote{el codomini de $f$ és $c$}},\\
\Com(f,g,h)&:\quad \text{\enquote{primer $f$, després $g$, i la composta és $h$}}.
\end{align*}
Així, quan la notació funcional ja estigui justificada, $\Com(f,g,h)$ correspondrà metalingüísticament a $h=g\comp f$.

La paraula \emph{fletxa} pertany, en aquest punt, al metallenguatge: no hi ha cap predicat primitiu $\operatorname{Fletxa}(f)$. Tots els individus del model s'interpreten com a fletxes. Una estructura per a $\mathcal L_{\mathrm{Cat}}$ és, doncs,
\[
\mathfrak M=(M,\Dom^{\mathfrak M},\Cod^{\mathfrak M},\Com^{\mathfrak M}),
\]
on $M$ és l'univers, $\Dom^{\mathfrak M},\Cod^{\mathfrak M}\subseteq M^2$ i $\Com^{\mathfrak M}\subseteq M^3$.

\begin{observacio}[metateòrica]
\label{obs:ar-metateorica}
Amb la semàntica estàndard de primer ordre, $M$ és un conjunt \emph{no buit}; per tant, els models en sentit estricte corresponen a categories petites no buides. La categoria buida queda exclosa per la convenció habitual que l'univers d'una estructura de primer ordre no és buit; si es vol incloure-la, cal permetre explícitament estructures amb univers buit. Les categories grans, d'altra banda, exigeixen una convenció metateòrica addicional ---classes pròpies o universos de Grothendieck, en la línia de l'apèndix~\ref{cap:mida}---, però això no modifica el llenguatge categòric primitiu anterior.
\end{observacio}

\section{Els axiomes}
\label{sec:ar-axiomes}

Presentem primer el sistema en cinc blocs conceptuals. A la secció~\ref{sec:ar-redundancia} veurem que, amb una formulació prou forta de l'associativitat, el tercer bloc és derivable dels altres; la base final del sistema és, doncs, \eqref{ax:A1D}--\eqref{ax:A1C}, \eqref{ax:A2}, \eqref{ax:A4E}--\eqref{ax:A4U} i \eqref{ax:A5D}--\eqref{ax:A5C}.

\paragraph{Axioma 1: existència i unicitat de domini i codomini.}
Tota fletxa té exactament un domini i exactament un codomini:
\begin{align}
\forall f\,\exists!d\,\Dom(f,d),
\tag{A1D}\label{ax:A1D}\\
\forall f\,\exists!c\,\Cod(f,c).
\tag{A1C}\label{ax:A1C}
\end{align}
Com és habitual, $\exists!$ és només una abreviació de la lògica de primer ordre amb igualtat.

\paragraph{Axioma 2: existència de la composició.}
Dues fletxes es poden compondre exactament quan el codomini de la primera coincideix amb el domini de la segona:
\begin{equation}
\forall f\,\forall g\,
\left[
\exists e\bigl(\Cod(f,e)\land\Dom(g,e)\bigr)
\leftrightarrow
\exists h\,\Com(f,g,h)
\right].
\tag{A2}\label{ax:A2}
\end{equation}
Aquest axioma afirma existència, però no unicitat. La unicitat de la composta serà un teorema (secció~\ref{sec:ar-unicitat}).

\paragraph{Axioma 3: extrems de la composta.}
Si $h$ és una composta de $f$ amb $g$, el domini de $h$ és el de $f$ i el codomini de $h$ és el de $g$:
\begin{equation}
\begin{split}
\forall f,g,h,d,c\,\bigl[
&\Com(f,g,h)\land\Dom(f,d)\land\Cod(g,c)\\
&\longrightarrow
\Dom(h,d)\land\Cod(h,c)
\bigr].
\end{split}
\tag{A3}\label{ax:A3}
\end{equation}

\paragraph{Axioma 4: associativitat forta.}
Com que $\Com$ és una relació i la composició és parcial, la mera igualtat de dues composicions iterades ja existents no expressa tota la força de l'associativitat usual. Adoptarem una formulació en dues parts. Primer, si existeixen les dues composicions consecutives $g\comp f$ i $h\comp g$, existeixen també les dues composicions iterades:
\begin{equation}
\begin{split}
\forall f,g,h,u,v\,\bigl[
&\Com(f,g,u)\land\Com(g,h,v)\\
&\longrightarrow
\exists p\,\exists q\,
\bigl(\Com(u,h,p)\land\Com(f,v,q)\bigr)
\bigr].
\end{split}
\tag{A4E}\label{ax:A4E}
\end{equation}
Segon, qualssevol resultats obtinguts per les dues associacions coincideixen:
\begin{equation}
\begin{split}
\forall f,g,h,u,v,p,q\,\bigl[
&\Com(f,g,u)\land\Com(g,h,v)\land\\
&\Com(u,h,p)\land\Com(f,v,q)
\longrightarrow p=q
\bigr].
\end{split}
\tag{A4U}\label{ax:A4U}
\end{equation}
Els dos enunciats, \eqref{ax:A4E} i \eqref{ax:A4U}, formen el bloc d'associativitat (A4). El primer controla l'existència de les composicions iterades; el segon n'expressa la coherència. Un cop demostrada la funcionalitat de $\Com$ (secció~\ref{sec:ar-unicitat}), aquest bloc es llegeix com la llei usual d'un semigrup parcial: si $g\comp f$ i $h\comp g$ existeixen, aleshores existeixen $h\comp(g\comp f)$ i $(h\comp g)\comp f$, i són iguals.

\paragraph{Axioma 5: neutralitat.}
El domini és neutre per la dreta i el codomini és neutre per l'esquerra:
\begin{align}
\forall f\,\forall d\,
\bigl(\Dom(f,d)\rightarrow\Com(d,f,f)\bigr),
\tag{A5D}\label{ax:A5D}\\
\forall f\,\forall c\,
\bigl(\Cod(f,c)\rightarrow\Com(f,c,f)\bigr).
\tag{A5C}\label{ax:A5C}
\end{align}
Metalingüísticament, un cop recuperada la notació funcional, aquestes fórmules diran $f\comp d=f$ i $c\comp f=f$.

\section{Estabilitat dels extrems}
\label{sec:ar-estabilitat}

El primer fet rellevant és que qualsevol fletxa que apareix com a domini o codomini és estable sota les dues relacions d'extrem.

\begin{proposicio}[estabilitat dels extrems]
\label{prop:ar-estabilitat}
Si $\Dom(f,d)$, aleshores $\Dom(d,d)\land\Cod(d,d)$. Si $\Cod(f,c)$, aleshores $\Dom(c,c)\land\Cod(c,c)$.
\end{proposicio}

\begin{prova}
Suposem $\Dom(f,d)$. Per \eqref{ax:A5D}, $\Com(d,f,f)$. Per la direcció de dreta a esquerra de \eqref{ax:A2}, l'existència d'aquesta composició implica que existeix $e$ tal que $\Cod(d,e)\land\Dom(f,e)$. Com que també $\Dom(f,d)$, la unicitat del domini de $f$, donada per \eqref{ax:A1D}, implica $e=d$; per tant $\Cod(d,d)$.

Ara \eqref{ax:A5C}, aplicada a $d$, dona $\Com(d,d,d)$. Novament per \eqref{ax:A2}, existeix $e$ tal que $\Cod(d,e)\land\Dom(d,e)$. Ja sabem que $\Cod(d,d)$; per unicitat del codomini de $d$, $e=d$, i per tant $\Dom(d,d)$. Així, $\Dom(f,d)\longrightarrow \Dom(d,d)\land\Cod(d,d)$.

El cas del codomini és dual. Si $\Cod(f,c)$, llavors \eqref{ax:A5C} dona $\Com(f,c,f)$. Per \eqref{ax:A2} existeix $e$ amb $\Cod(f,e)\land\Dom(c,e)$; la unicitat del codomini de $f$ força $e=c$, de manera que $\Dom(c,c)$. Aleshores \eqref{ax:A5D} dona $\Com(c,c,c)$, i un nou ús de \eqref{ax:A2} i de la unicitat del domini de $c$ proporciona $\Cod(c,c)$.
\end{prova}

\begin{observacio}
La demostració només utilitza (A1), (A2) i (A5). No utilitza ni l'axioma dels extrems de la composta (A3) ni l'associativitat (A4).
\end{observacio}

\section{Unicitat de la composició}
\label{sec:ar-unicitat}

La relació $\Com$ no s'ha postulat com a funcional en el tercer argument. Aquesta funcionalitat es dedueix de la componibilitat, la neutralitat i l'associativitat.

\begin{proposicio}[funcionalitat de la composició]
\label{prop:ar-unicitat}
Per a totes les fletxes $f,g,h,k$,
\[
\Com(f,g,h)\land\Com(f,g,k)\longrightarrow h=k.
\]
\end{proposicio}

\begin{prova}
Suposem $\Com(f,g,h)$ i $\Com(f,g,k)$. Per \eqref{ax:A2}, existeix $e$ tal que $\Cod(f,e)\land\Dom(g,e)$. Per neutralitat, $\Com(f,e,f)$ i $\Com(e,g,g)$. Apliquem ara \eqref{ax:A4U} a la terna $f,e,g$. Prenent $u=f$, $v=g$, $p=h$, $q=k$, les quatre premisses de \eqref{ax:A4U} són precisament $\Com(f,e,f)$, $\Com(e,g,g)$, $\Com(f,g,h)$ i $\Com(f,g,k)$. Per tant, $h=k$. La demostració només fa servir (A2), (A5) i \eqref{ax:A4U}; en particular, no fa servir (A3).
\end{prova}

Així, $\Com$ és el graf d'una operació binària parcial: per a cada parella $(f,g)$ hi ha com a màxim una fletxa $h$ tal que $\Com(f,g,h)$, i (A2) determina exactament quan aquesta fletxa existeix.

\section{Identitats i objectes com a nocions derivades}
\label{sec:ar-identitats}

Definim, com a abreviació del llenguatge,
\begin{equation}
\Id(e)
\;:\Longleftrightarrow\;
\Dom(e,e)\land\Cod(e,e).
\label{def:ar-Id}
\end{equation}
El símbol $\Id$ no s'afegeix a la signatura primitiva: és només una abreviació definicional. La proposició d'estabilitat dona immediatament $\Dom(f,d)\rightarrow\Id(d)$ i $\Cod(f,c)\rightarrow\Id(c)$: tot domini i tot codomini és una fletxa identitària.

També podem introduir $\Obj(e):\Longleftrightarrow\Id(e)$. Així, un \emph{objecte} no és una segona mena d'individu: és una fletxa que té per domini i codomini ella mateixa. En particular, sota els axiomes anteriors,
\[
\Dom(e,e)
\;\Longleftrightarrow\;
\Cod(e,e)
\;\Longleftrightarrow\;
\Id(e),
\]
perquè qualsevol de les dues primeres propietats, aplicada a la proposició~\ref{prop:ar-estabilitat}, implica l'altra.

\section{El tercer axioma és derivable}
\label{sec:ar-redundancia}

La forma forta de l'associativitat permet demostrar l'axioma (A3) a partir de (A1), (A2), (A4) i (A5).

\begin{proposicio}[redundància de (A3)]
\label{prop:ar-redundancia}
Suposem (A1), (A2), (A4) i (A5). Si $\Com(f,g,h)$, $\Dom(f,d)$ i $\Cod(g,c)$, aleshores $\Dom(h,d)$ i $\Cod(h,c)$.
\end{proposicio}

\begin{prova}
Les proposicions~\ref{prop:ar-estabilitat} i~\ref{prop:ar-unicitat}, que farem servir, no depenen de (A3); l'argument no és, doncs, circular. Comencem pel domini. De $\Dom(f,d)$ i \eqref{ax:A5D} obtenim $\Com(d,f,f)$. Juntament amb $\Com(f,g,h)$, l'axioma d'existència associativa \eqref{ax:A4E}, aplicat a la terna $d,f,g$, proporciona $p,q$ tals que $\Com(f,g,p)$ i $\Com(d,h,q)$. Per \eqref{ax:A4U}, aquestes dues composicions iterades tenen el mateix resultat, és a dir, $p=q$. D'altra banda, per la unicitat de la composició (proposició~\ref{prop:ar-unicitat}), $\Com(f,g,h)$ i $\Com(f,g,p)$ impliquen $p=h$. Per tant $q=h$, i tenim $\Com(d,h,h)$. Per \eqref{ax:A2}, existeix $e$ amb $\Cod(d,e)\land\Dom(h,e)$. Com que $d$ és domini de $f$, la proposició~\ref{prop:ar-estabilitat} dona $\Cod(d,d)$. La unicitat del codomini de $d$ implica $e=d$, i obtenim $\Dom(h,d)$.

Per al codomini, de $\Cod(g,c)$ i \eqref{ax:A5C} obtenim $\Com(g,c,g)$. Aplicada a $f,g,c$, l'associativitat forta proporciona composicions iterades $\Com(h,c,p)$ i $\Com(f,g,q)$ amb $p=q$. La unicitat de la composició de $f$ amb $g$ dona $q=h$, i per tant $\Com(h,c,h)$. Per \eqref{ax:A2}, existeix $e$ amb $\Cod(h,e)\land\Dom(c,e)$. Com que $c$ és codomini de $g$, l'estabilitat dona $\Dom(c,c)$. Per unicitat del domini de $c$, $e=c$, i així $\Cod(h,c)$.
\end{prova}

\begin{corollari}
\label{cor:ar-equivalencia-sistemes}
Amb l'associativitat forta adoptada aquí, el sistema de cinc blocs $(A1),(A2),(A3),(A4),(A5)$ és equivalent al sistema reduït $(A1),(A2),(A4),(A5)$. El tercer axioma original es pot conservar en una primera presentació per la seva transparència conceptual, però no és necessari com a axioma independent.
\end{corollari}

No afirmem que els quatre blocs restants siguin independents entre si: la redundància de (A3) no estableix la minimalitat del sistema reduït. L'anàlisi d'independència de sistemes funcionals afins ha estat estudiada i formalitzada per Benzmüller i Scott~\cite{BS2020,BSAFP}, però la transferència literal d'aquells resultats al llenguatge purament relacional adoptat aquí requeriria una comparació formal addicional; una via natural seria construir models finits que separessin cadascun dels quatre blocs.

\begin{observacio}[la forma lògica de (A4) és essencial]
Aquesta reducció depèn de la forma lògica de (A4). Si l'associativitat es formula només com la igualtat de dues composicions iterades \emph{suposant que totes dues ja existeixen}, no es controla l'existència de les composicions iterades i la demostració anterior deixa de funcionar. En una teoria de composició parcial, la forma lògica de l'associativitat és, doncs, part essencial de l'axiomàtica.
\end{observacio}

\section{Recuperació de la notació funcional}
\label{sec:ar-notacio-funcional}

L'axioma (A1) permet introduir en el metallenguatge les notacions $d(f)$ i $c(f)$ per a les úniques fletxes que satisfan, respectivament, $\Dom(f,d(f))$ i $\Cod(f,c(f))$. Aquestes expressions no són termes del llenguatge primitiu $\mathcal L_{\mathrm{Cat}}$: són notació derivada, legitimada per un teorema d'existència i unicitat; en termes lògics, passem a una extensió definicional per descripcions úniques. De la mateixa manera, si $f$ i $g$ són componibles, (A2) dona l'existència d'una $h$ tal que $\Com(f,g,h)$, i la proposició~\ref{prop:ar-unicitat} en dona la unicitat; escrivim, doncs, $g\comp f$ per designar aquesta única fletxa.

Amb aquesta notació, (A2) es llegeix
\[
g\comp f\ \text{existeix}
\quad\Longleftrightarrow\quad
c(f)=d(g),
\]
i la proposició~\ref{prop:ar-redundancia} dona $d(g\comp f)=d(f)$ i $c(g\comp f)=c(g)$.

Si $A$ i $B$ són fletxes identitàries, la notació usual $f\colon A\longrightarrow B$ abrevia metalingüísticament $\Dom(f,A)\land\Cod(f,B)$; les condicions $\Id(A)$ i $\Id(B)$ són automàtiques per estabilitat. També podem escriure, metalingüísticament, $\Hom(A,B)=\{f\in M: \Dom(f,A)\land\Cod(f,B)\}$. El símbol $\Hom$ no forma part tampoc del llenguatge primitiu.

\section{Equivalència amb la definició usual}
\label{sec:ar-equivalencia}

La presentació relacional no defineix una estructura matemàtica nova: recupera la noció usual de categoria petita (definició~\ref{def:categoria}), fins a l'isomorfisme canònic que identifica cada objecte amb la seva identitat.

Abans de les demostracions, val la pena desplegar un model concret i petit, per veure com les relacions $\Dom$, $\Cod$ i $\Com$ codifiquen una categoria familiar i com els objectes hi apareixen com a fletxes identitat.

\begin{exemple}[Un model amb tres objectes i una cadena de morfismes]\label{ex:ar-model-3}
Considerem la categoria $\cat{3}$ (amb els objectes i les fletxes reanomenats), formada per una cadena de dues fletxes, $A\xrightarrow{u}B\xrightarrow{v}C$, amb la seva composta $w=v\comp u\colon A\to C$:
\[
\begin{tikzpicture}[baseline=(A.base)]
  \node (A) at (0,0) {$A$};
  \node (B) at (2.2,0) {$B$};
  \node (C) at (4.4,0) {$C$};
  \draw[-{Stealth[scale=1.1]}] (A) -- node[above] {$u$} (B);
  \draw[-{Stealth[scale=1.1]}] (B) -- node[above] {$v$} (C);
  \draw[-{Stealth[scale=1.1]}] (A) to[out=-35,in=-145] node[below] {$w=v\comp u$} (C);
\end{tikzpicture}
\]
El model relacional corresponent té per univers el conjunt de \emph{totes} les fletxes, incloses les tres identitats:
\[
M=\{\,A,\ B,\ C,\ u,\ v,\ w\,\},
\]
on $A,B,C$ són alhora els objectes i les seves fletxes identitat. Les relacions unàries derivades i els extrems de cada fletxa són:
\[
\renewcommand{\arraystretch}{1.2}
\begin{array}{c|c|c|c}
f & \Dom(f,-) & \Cod(f,-) & \Id(f) \\ \hline
A & A & A & \text{sí} \\
B & B & B & \text{sí} \\
C & C & C & \text{sí} \\
u & A & B & \text{no} \\
v & B & C & \text{no} \\
w & A & C & \text{no}
\end{array}
\]
Els objectes són exactament les fletxes $e$ amb $\Id(e)$, és a dir $A,B,C$: no calen com a dada primitiva, sinó que es reconeixen per la propietat $\Dom(e,e)\land\Cod(e,e)$. Les úniques parelles componibles $(f,g)$ ---les que compleixen $c(f)=d(g)$--- donen la taula de $\Com$, on recordem que $\Com(f,g,h)$ es llegeix \enquote{primer $f$, després $g$, amb composta $h=g\comp f$}:
\[
\renewcommand{\arraystretch}{1.2}
\begin{array}{ll}
\Com(A,A,A) & \text{(identitat de }A) \\
\Com(B,B,B) & \text{(identitat de }B) \\
\Com(C,C,C) & \text{(identitat de }C) \\
\Com(A,u,u),\ \Com(u,B,u) & (u\comp\id_A=u=\id_B\comp u) \\
\Com(B,v,v),\ \Com(v,C,v) & (v\comp\id_B=v=\id_C\comp v) \\
\Com(A,w,w),\ \Com(w,C,w) & (w\comp\id_A=w=\id_C\comp w) \\
\Com(u,v,w) & (\text{la composició no trivial }v\comp u=w)
\end{array}
\]
Val la pena llegir l'última línia amb l'ordre de $\Com$ ben present: $\Com(u,v,w)$ diu que en recórrer \emph{primer} $u$ i \emph{després} $v$ s'obté $w$, és a dir $w=v\comp u$. No hi ha cap altra parella componible; en particular, $u$ i $w$ no es componen en cap ordre, perquè $c(u)=B\neq A=d(w)$ i $c(w)=C\neq A=d(u)$. Comprovem els axiomes. (A1) és la taula d'extrems. (A2) es comprova parella per parella: les deu ternes de la taula de $\Com$ corresponen exactament a les deu parelles $(f,g)$ amb $c(f)=d(g)$. (A5) és el contingut de les sis primeres línies, que donen $\Com(d(f),f,f)$ i $\Com(f,c(f),f)$ per a cada $f$. Per a (A4), observem que una terna encadenada $(f,g,h)$, amb $\Com(f,g,u)$ i $\Com(g,h,v)$, conté sempre alguna identitat, perquè la cadena més llarga de fletxes no identitat és $u,v$, de longitud dos. Si la identitat és $f$, aleshores $u=g$ i $f$ és el domini de $v$, de manera que totes dues associacions donen $v$: $\Com(u,h,v)$ i $\Com(f,v,v)$. Els casos en què la identitat és $g$ o $h$ són anàlegs, i totes dues associacions donen la composta de les altres dues fletxes. Així es compleixen \eqref{ax:A4E} i \eqref{ax:A4U}. Reconstruint la categoria a partir d'aquest model (com farà el teorema~\ref{th:ar-model-a-categoria}) es recupera la categoria de partida, en què ara cada objecte és la seva pròpia fletxa identitat.
\end{exemple}

\begin{teorema}
\label{th:ar-model-a-categoria}
Tot model del sistema reduït $(A1),(A2),(A4),(A5)$ amb univers no buit determina una categoria petita no buida en el sentit de la definició~\ref{def:categoria}.
\end{teorema}

\begin{prova}
Sigui $\mathfrak M$ un model. Prenem com a fletxes els elements de $M$ i com a objectes $\Obj(\mathfrak M)=\{e\in M:\Id(e)\}$. Per (A1), cada fletxa $f$ té un domini únic $d(f)$ i un codomini únic $c(f)$; per estabilitat, tots dos són objectes. Si $c(f)=d(g)$, (A2) proporciona una composta i la proposició~\ref{prop:ar-unicitat} en proporciona una de sola, que denotem $g\comp f$; la proposició~\ref{prop:ar-redundancia} assegura que $d(g\comp f)=d(f)$ i $c(g\comp f)=c(g)$. Per a un objecte $A$, la seva fletxa identitat és $A$ mateix. Aquesta fletxa va de $A$ a $A$, perquè $\Dom(A,A)$ i $\Cod(A,A)$. Si $f\colon A\to B$, (A5) dona $\Com(A,f,f)$ i $\Com(f,B,f)$, és a dir, $f\comp\id_A=f=\id_B\comp f$. Per a l'associativitat, siguin $f\colon A\to B$, $g\colon B\to C$ i $h\colon C\to D$. Aplicant \eqref{ax:A4E} a $\Com(f,g,g\comp f)$ i $\Com(g,h,h\comp g)$, existeixen $p,q$ amb $\Com(g\comp f,h,p)$ i $\Com(f,h\comp g,q)$; per la unicitat de la composició, $p=h\comp(g\comp f)$ i $q=(h\comp g)\comp f$, i \eqref{ax:A4U} dona $p=q$. Així obtenim una categoria, petita perquè $M$ és un conjunt.
\end{prova}

\begin{teorema}[recíproc]
\label{th:ar-categoria-a-model}
Tota categoria petita no buida determina un model del sistema relacional.
\end{teorema}

\begin{prova}
Sigui $\cat{C}$ una categoria petita no buida. Prenem com a univers $M=\operatorname{Mor}(\cat{C})$. Si $f\colon A\to B$, definim
\[
\Dom(f,d)\Longleftrightarrow d=\id_A,
\qquad
\Cod(f,c)\Longleftrightarrow c=\id_B,
\]
i
i, per a $f\colon A\to B$ i $g\colon B'\to C$,
\[
\Com(f,g,h)\Longleftrightarrow \bigl(B=B'\ \text{i}\ h=g\comp f\bigr),
\]
on la notació de la dreta pertany ara al metallenguatge de la categoria usual. L'univers és no buit, perquè $\cat{C}$ té algun objecte i, per tant, alguna identitat. Comprovem els axiomes. (A1) és immediat. Per a (A2), $\Cod(f,e)\land\Dom(g,e)$ vol dir $e=\id_B=\id_{B'}$, i això passa si i només si $B=B'$, perquè $\id_B$ determina $B$ com a domini; i $B=B'$ és exactament la condició perquè existeixi $h$ amb $\Com(f,g,h)$. (A3) diu que $g\comp f\colon A\to C$. Per a (A4), $\Com(f,g,u)$ i $\Com(g,h,v)$ volen dir que $f,g,h$ són consecutives, amb $u=g\comp f$ i $v=h\comp g$; aleshores existeixen $h\comp u$ i $v\comp f$, que són iguals per l'associativitat de $\cat{C}$. Finalment, (A5) són les lleis d'identitat $f\comp\id_A=f$ i $\id_B\comp f=f$.
\end{prova}

\begin{observacio}
En passar d'una categoria usual $\cat{C}$ al model relacional, els individus són els morfismes de $\cat{C}$; en reconstruir després una categoria a partir del model, els seus objectes són les fletxes identitat $\id_A$, no els objectes originals $A$. La categoria reconstruïda és canònicament isomorfa a $\cat{C}$ ---la bijecció d'objectes és $A\mapsto\id_A$--- però en general no n'és literalment la mateixa estructura conjuntista. Per aquest motiu diem que les dues presentacions descriuen la mateixa estructura categòrica \emph{fins a isomorfisme canònic}, no que en siguin idèntiques. Una afirmació metateòrica més forta ---una equivalència definicional o una bi-interpretabilitat--- requeriria definir explícitament les traduccions entre les dues teories i comparar-les; no la formulem aquí.

\end{observacio}

Així, la diferència entre les dues presentacions és una diferència de \emph{llenguatge primitiu}, no de contingut categòric.

\section{Relacional i funcional}
\label{sec:ar-relacional-funcional}

Els predicats $\Dom$ i $\Cod$, sota (A1), són els grafs de dues funcions unàries totals. La relació $\Com$, sota (A2) i la proposició~\ref{prop:ar-unicitat}, és el graf d'una funció binària parcial. Aquesta observació explica dues maneres naturals de formalitzar les categories.

\paragraph{Presentació relacional.}
En lògica de primer ordre estàndard podem prendre com a primitives $\Dom(f,d)$, $\Cod(f,c)$ i $\Com(f,g,h)$. L'existència s'expressa mitjançant quantificadors i la funcionalitat s'axiomatitza o es demostra. És la via d'aquest apèndix.

\paragraph{Presentació funcional.}
Podem prendre $\dom$ i $\cod$ com a operadors unaris. El problema és la composició: com que no tota parella de fletxes és componible, una operació binària de composició no és total sobre $M\times M$. En lògica de primer ordre estàndard, en canvi, els símbols de funció s'interpreten com a funcions totals. Cal, doncs, ampliar el marc lògic o codificar explícitament la parcialitat.

Aquest és el camí seguit, en un formalisme diferent, per Benzmüller i Scott~\cite{BS2020,BSAFP}. Treballen amb una lògica lliure embeguda en Isabelle/HOL, un predicat $E$ d'existència o definició, i tres operacions parcials $\dom$, $\cod$ i composició. En el seu \emph{AxiomsSet5}, atribuït a Scott, els axiomes principals tenen la forma
\begin{align*}
E(\dom x)&\to E(x),\tag{S1}\\
E(\cod y)&\to E(y),\tag{S2}\\
E(x\cdot y)&\leftrightarrow \dom x\approx\cod y,\tag{S3}\\
x\cdot(y\cdot z)&\simeq (x\cdot y)\cdot z,\tag{S4}\\
(\cod y)\cdot y&\simeq y,\tag{S5}\\
x\cdot(\dom x)&\simeq x,\tag{S6}
\end{align*}
amb la composició escrita en ordre funcional ($x\cdot y$ és \enquote{primer $y$, després $x$}), on $\simeq$ és la \emph{igualtat de Kleene} (dos termes són iguals si tots dos són indefinits, o bé tots dos definits i iguals) i $\approx$ és la \emph{identitat existent} (tots dos termes definits i iguals). En aquesta presentació, la funcionalitat de $\dom$, $\cod$ i de la composició queda incorporada en els símbols funcionals; la parcialitat es gestiona mitjançant la lògica lliure. Els mateixos autors formalitzen també equivalències entre diversos sistemes axiomàtics i el pas, per Skolemització, de formulacions existencials cap a formulacions funcionals.

\begin{nota}
Els símbols tipogràfics per a aquestes dues igualtats varien entre fonts; aquí fixem $\simeq$ per a la igualtat de Kleene i $\approx$ per a la identitat existent, que no coincideixen necessàriament amb la notació de l'obra citada.
\end{nota}

La correspondència conceptual amb la presentació relacional és:
\[
\begin{array}{c|c}
\text{presentació relacional} & \text{presentació funcional lliure}\\
\hline
\text{(A1): domini i codomini} & \dom,\cod\ \text{i } S1,S2\\
\text{(A2): componibilitat} & S3\\
\text{(A3): extrems de la composta} & \text{derivable en sistemes reduïts}\\
\text{(A4): associativitat} & S4\\
\text{(A5): neutralitat} & S5,S6
\end{array}
\]
La correspondència no és literal fórmula per fórmula, perquè els dos llenguatges distribueixen la informació de manera diferent.

\section{El principi de dualitat}
\label{sec:ar-dualitat}\index{dualitat!principi de}

El capítol~\ref{cap:categories} formula el principi de dualitat de manera operativa (teorema~\ref{teo:dualitat}): una afirmació expressable en el llenguatge de les categories es dualitza invertint les fletxes, i, si és vàlida en totes les categories, la seva dual també ho és, perquè l'afirmació original es pot aplicar a la categoria oposada. La presentació relacional permet donar-ne una versió completament formal, en dues formes: una de \emph{sintàctica}, sobre deduccions, i una de \emph{semàntica}, sobre models. En aquesta secció escrivim $\mathrm{CT}$ per al conjunt d'axiomes \eqref{ax:A1D}--\eqref{ax:A5C}.

\paragraph{La traducció dual.}
Per a cada fórmula $\varphi$ de $\mathcal L_{\mathrm{Cat}}$ definim la seva \textbf{dual}\index{dualitat!d'un enunciat} $\varphi\op$ substituint cada fórmula atòmica no lògica segons
\begin{gather*}
\Dom(f,d)\longmapsto\Cod(f,d),\qquad
\Cod(f,c)\longmapsto\Dom(f,c),\\
\Com(f,g,h)\longmapsto\Com(g,f,h),
\end{gather*}
i deixant intactes les igualtats, els connectors i els quantificadors. Com que $\Com(f,g,h)$ es llegeix \enquote{primer $f$, després $g$, amb composta $h$}, intercanviar-ne els dos primers arguments és exactament invertir l'ordre de composició. La traducció és una involució, $(\varphi\op)\op=\varphi$, i commuta amb la substitució de variables. Un cop recuperada la notació funcional (secció~\ref{sec:ar-notacio-funcional}), aquestes substitucions són les del capítol~\ref{cap:categories}: $d(f)$ i $c(f)$ s'intercanvien i $g\comp f$ passa a $f\comp g$. Observem que la dualitat no depèn de tractar la composició com un símbol funcional total: opera directament sobre la relació ternària.

\begin{observacio}
Els connectors i quantificadors de $\varphi$ pertanyen al llenguatge $\mathcal L_{\mathrm{Cat}}$, amb què \emph{parlem} d'una categoria, i la dualitat els deixa intactes. No s'han de confondre amb la lògica que una categoria pot codificar \emph{internament}, quan els seus objectes són fórmules i els seus morfismes deduccions. En aquest cas, la dualitat de fletxes pot intercanviar la conjunció i la disjunció del sistema codificat, però el $\land$ i el $\lor$ de $\mathcal L_{\mathrm{Cat}}$ es mantenen.
\end{observacio}

\paragraph{Autodualitat dels axiomes.}

\begin{lema}\label{lema:ar-autodual}
Per a cada axioma $\alpha$ de $\mathrm{CT}$, la fórmula $\alpha\op$ és lògicament equivalent a un axioma de $\mathrm{CT}$. Concretament, \eqref{ax:A1D} i \eqref{ax:A1C} es transformen l'un en l'altre, igual que \eqref{ax:A5D} i \eqref{ax:A5C}; i cadascun dels axiomes \eqref{ax:A2}, \eqref{ax:A3}, \eqref{ax:A4E} i \eqref{ax:A4U} es transforma en una variant d'ell mateix, que només en difereix pel nom de les variables lligades, per l'ordre dels quantificadors i dels membres de les conjuncions i, en el cas de \eqref{ax:A4U}, per la simetria de la igualtat.
\end{lema}

\begin{prova}
Els casos d'\eqref{ax:A1D} i \eqref{ax:A5D} són immediats: $(\forall f\,\exists!d\,\Dom(f,d))\op$ és $\forall f\,\exists!d\,\Cod(f,d)$, que és \eqref{ax:A1C}, i $(\forall f\,\forall d\,(\Dom(f,d)\to\Com(d,f,f)))\op$ és $\forall f\,\forall d\,(\Cod(f,d)\to\Com(f,d,f))$, que és \eqref{ax:A5C}; per involució, els casos d'\eqref{ax:A1C} i \eqref{ax:A5C} són els recíprocs.

La dual d'\eqref{ax:A2} és
\[
\forall f\,\forall g\,\bigl[\exists e\bigl(\Dom(f,e)\land\Cod(g,e)\bigr)\leftrightarrow\exists h\,\Com(g,f,h)\bigr],
\]
i, intercanviant els noms de les variables lligades $f$ i $g$, s'obté \eqref{ax:A2} llevat de l'ordre dels quantificadors i dels membres de la conjunció. La dual d'\eqref{ax:A3} és
\[
\forall f,g,h,d,c\,\bigl[\Com(g,f,h)\land\Cod(f,d)\land\Dom(g,c)\longrightarrow\Cod(h,d)\land\Dom(h,c)\bigr],
\]
i l'intercanvi $f\leftrightarrow g$, $d\leftrightarrow c$ la torna en \eqref{ax:A3}. La dual d'\eqref{ax:A4E} és
\[
\forall f,g,h,u,v\,\bigl[\Com(g,f,u)\land\Com(h,g,v)\longrightarrow\exists p\,\exists q\,\bigl(\Com(h,u,p)\land\Com(v,f,q)\bigr)\bigr];
\]
amb l'intercanvi $f\leftrightarrow h$, $u\leftrightarrow v$, $p\leftrightarrow q$ esdevé
\[
\forall f,g,h,u,v\,\bigl[\Com(g,h,v)\land\Com(f,g,u)\longrightarrow\exists q\,\exists p\,\bigl(\Com(f,v,q)\land\Com(u,h,p)\bigr)\bigr],
\]
que és \eqref{ax:A4E}. El mateix reanomenament transforma la dual d'\eqref{ax:A4U} en \eqref{ax:A4U}, amb la conclusió escrita $q=p$.
\end{prova}

Intuïtivament, el lema diu que la definició de categoria no distingeix entre una fletxa i la mateixa fletxa recorreguda en sentit contrari. És l'únic ingredient específic de la teoria que necessiten els dos teoremes següents.

\paragraph{Dualitat sintàctica.}

\begin{teorema}[Dualitat sintàctica]\label{teo:ar-dualitat-sint}
Per a tota sentència $\varphi$ de $\mathcal L_{\mathrm{Cat}}$, si $\mathrm{CT}\vdash\varphi$, aleshores $\mathrm{CT}\vdash\varphi\op$.
\end{teorema}

\begin{prova}
Fixem un càlcul de Hilbert estàndard per a la lògica de primer ordre amb igualtat. La traducció dual només reanomena les relacions primitives i permuta els arguments de $\Com$; commuta amb la substitució, conserva les variables lliures i no toca ni els connectors, ni els quantificadors, ni la igualtat. Per tant, transforma cada axioma lògic (inclosos els de la igualtat) en un axioma lògic, i cada aplicació d'una regla d'inferència en una aplicació de la mateixa regla; les restriccions sobre variables de la regla de generalització es conserven, perquè les variables lliures no canvien. Aplicada fórmula a fórmula a una deducció de $\varphi$ a partir de $\mathrm{CT}$, en dona una de $\varphi\op$ a partir de $\mathrm{CT}\op=\{\alpha\op:\alpha\in\mathrm{CT}\}$. Pel lema~\ref{lema:ar-autodual}, cada $\alpha\op$ es dedueix de $\mathrm{CT}$; per tant, $\mathrm{CT}\vdash\varphi\op$.
\end{prova}

\paragraph{Dualitat semàntica.}
Per a una estructura $\mathfrak M=(M,\Dom^{\mathfrak M},\Cod^{\mathfrak M},\Com^{\mathfrak M})$ definim l'\textbf{estructura oposada}
\[
\mathfrak M\op=\bigl(M,\ \Cod^{\mathfrak M},\ \Dom^{\mathfrak M},\ \{(g,f,h):(f,g,h)\in\Com^{\mathfrak M}\}\bigr),
\]
amb el mateix univers, els papers de domini i codomini intercanviats i la composició en l'ordre invers. Si $\mathfrak M$ és el model associat a una categoria petita $\cat{C}$ pel teorema~\ref{th:ar-categoria-a-model}, aleshores l'aplicació $f\mapsto f\op$ és un isomorfisme de $\mathfrak M\op$ en el model associat a $\cat{C}\op$: per a $f\colon A\to B$, la relació $\Dom^{\mathfrak M\op}(f,\id_B)$, que és $\Cod^{\mathfrak M}(f,\id_B)$, correspon a $\Dom(f\op,(\id_B)\op)$ a $\cat{C}\op$, i $\Com^{\mathfrak M\op}(f,g,h)$, que és $\Com^{\mathfrak M}(g,f,h)$, és a dir $h=f\comp g$, correspon a $h\op=g\op\comp f\op$, és a dir, a $\Com(f\op,g\op,h\op)$ a $\cat{C}\op$. Si s'identifica cada $f\op$ amb $f$, els dos models coincideixen.

\begin{lema}\label{lema:ar-oposada}
Per a tota fórmula $\varphi$ de $\mathcal L_{\mathrm{Cat}}$ i tota assignació $s$ de les seves variables lliures,
\[
\mathfrak M\op\models\varphi[s]
\quad\Longleftrightarrow\quad
\mathfrak M\models\varphi\op[s].
\]
\end{lema}

\begin{prova}
Per inducció sobre $\varphi$. Per a les fórmules atòmiques és la definició de $\mathfrak M\op$: per exemple, $\mathfrak M\op\models\Com(f,g,h)[s]$ si i només si $(s(g),s(f),s(h))\in\Com^{\mathfrak M}$, és a dir, si i només si $\mathfrak M\models\Com(g,f,h)[s]$. Les igualtats s'interpreten igual a totes dues estructures, i els passos dels connectors i dels quantificadors són immediats perquè l'univers és el mateix.
\end{prova}

\begin{teorema}[Dualitat semàntica]\label{teo:ar-dualitat-sem}
Si $\mathfrak M$ és un model de $\mathrm{CT}$, també ho és $\mathfrak M\op$. En conseqüència, per a tota sentència $\varphi$ de $\mathcal L_{\mathrm{Cat}}$, si $\varphi$ es compleix en tots els models de $\mathrm{CT}$ ---que, pels teoremes~\ref{th:ar-model-a-categoria} i~\ref{th:ar-categoria-a-model}, corresponen a les categories petites no buides---, també s'hi compleix $\varphi\op$.
\end{teorema}

\begin{prova}
Si $\mathfrak M\models\mathrm{CT}$ i $\alpha\in\mathrm{CT}$, pel lema~\ref{lema:ar-autodual} la fórmula $\alpha\op$ és equivalent a un axioma, de manera que $\mathfrak M\models\alpha\op$, i pel lema~\ref{lema:ar-oposada}, $\mathfrak M\op\models\alpha$. Per a la segona afirmació, sigui $\mathfrak M$ un model de $\mathrm{CT}$. Com que $\mathfrak M\op$ també ho és, $\mathfrak M\op\models\varphi$, i el lema~\ref{lema:ar-oposada} dona $\mathfrak M\models\varphi\op$.
\end{prova}

Aquesta demostració és la versió formal de la del capítol~\ref{cap:categories}: allà es fa servir que $\cat{C}\op$ és una categoria; aquí, que $\mathfrak M\op$ és un model. Pel teorema de completesa, $\mathrm{CT}\vdash\varphi$ equival a $\mathrm{CT}\models\varphi$, i els dos teoremes són, doncs, dues cares d'un mateix fet. Les restriccions sobre la mida de l'observació~\ref{obs:ar-metateorica} es mantenen: la versió semàntica parla de categories petites no buides, i l'extensió a categories grans demana la mateixa convenció metateòrica addicional.

\paragraph{Abast.}
Tots dos teoremes parlen de fórmules de $\mathcal L_{\mathrm{Cat}}$. Una afirmació que mencioni elements d'un conjunt, una categoria concreta o qualsevol estructura externa no és una fórmula d'aquest llenguatge i no té dual fins que no s'hi reformula. Per això el capítol~\ref{cap:categories} enuncia el principi per a afirmacions \emph{expressables en el llenguatge de les categories}, i no per a \enquote{qualsevol afirmació sobre categories}. Els enunciats que involucren diverses categories i functors entre elles demanen un llenguatge més ric; el mecanisme és el mateix, però cal dualitzar també les categories implicades, com mostra l'exemple de les categories sobre i sota un objecte.

\paragraph{Un mecanisme general.}
El que hem fet ---una involució del llenguatge que deixa invariant el conjunt d'axiomes--- no és exclusiu de la teoria de categories. La lògica clàssica en dona una altra instància: en el llenguatge de les àlgebres booleanes, la involució que intercanvia $\wedge$ amb $\vee$ i el cert amb el fals deixa els axiomes invariants, i produeix la dualitat de De~Morgan. Un exemple encara més antic és la \textbf{dualitat projectiva}\index{dualitat!projectiva}: en el pla projectiu, intercanviar \emph{punt} per \emph{recta} ---i \enquote{passar per} per \enquote{tallar-se en}--- deixa els axiomes invariants, de manera que cada teorema té un teorema dual, com el de Desargues o la parella Pascal--Brianchon. En tots els casos actua el mateix mecanisme, aplicat a llenguatges diferents. Quan una categoria codifica lògica, la dualitat de fletxes intercanvia conjunció i disjunció si existeixen; que aquesta dualitat coincideixi exactament amb la de De~Morgan requereix, a més, un complement involutiu, i és una hipòtesi addicional sobre la categoria, no una conseqüència automàtica de la dualitat categòrica.

\section{Una observació computacional}
\label{sec:ar-computacional}

La presentació relacional és especialment transparent per a l'axiomàtica, però no és necessàriament la més còmoda per a una implementació. Un programa prefereix, habitualment, operacions que retornin valors, $\dom\colon M\to M$ i $\cod\colon M\to M$, i una composició parcial. Aquesta darrera es pot representar, per exemple, com $\operatorname{comp}\colon M\times M\longrightarrow M_{\bot}$ o mitjançant un tipus opcional que distingeixi entre una composta existent i una parella no componible.

No hi ha contradicció entre les dues perspectives. La teoria relacional demostra primer existència i unicitat; aquests resultats justifiquen després el pas a una interfície funcional. El millor llenguatge per axiomatitzar una estructura no ha de coincidir amb el millor llenguatge per implementar-la.

\section{Connexió amb una lectura estructural}
\label{sec:ar-estructural}

La presentació anterior respon a un criteri estructural molt simple: no introduir com a primitives entitats que poden reconstruir-se a partir de les relacions fonamentals de la teoria. En la presentació usual es parla de dos tipus de dades, objectes i fletxes; en la relacional elemental només hi ha una mena d'individus, i els objectes apareixen com les fletxes que satisfan $\Dom(e,e)\land\Cod(e,e)$. L'objecte no és, doncs, una entitat afegida al costat de les fletxes, sinó una posició estructural determinada per les relacions de domini i codomini ---exactament el gir que la introducció anunciava en preferir \emph{relacionar} a \emph{individualitzar} (secció~\ref{sec:mirar-fora}).

Aquesta perspectiva és compatible amb una lectura estructural de la teoria de conjunts: les realitzacions conjuntistes concretes poden quedar al metallenguatge, mentre que el llenguatge de la teoria categòrica reté només les relacions que expressen l'estructura que volem estudiar. Això no elimina la teoria de conjunts del metallenguatge ---la semàntica estàndard de primer ordre continua sent conjuntista, i el marc de mida el fixa l'apèndix corresponent--- però evita confondre una realització externa amb les nocions primitives de la teoria objecte.

La lliçó metodològica és, doncs, doble. D'una banda, la formulació relacional mostra que una categoria es pot reconstruir a partir d'un univers de fletxes i de tres relacions primitives. De l'altra, un cop feta la reconstrucció, no hi ha cap inconvenient a recuperar la notació funcional i la terminologia usual, que són molt més eficients per desenvolupar la teoria ---i que són, de fet, les que la resta d'aquest llibre empra.

\chapter{Grafs}\label{cap:grafs}

Aquest apèndix desenvolupa amb detall la noció de graf dirigit, els seus homomorfismes i els camins, i mostra com generar una categoria a partir d'un graf. El material amplia, amb molts exemples concrets, la subsecció~\ref{subsec:grafs} del capítol de categories.

\section{Grafs}

Qualsevol sistema de relacions, com les rutes de transport o les connexions d'aerolínies, es pot dibuixar com una col·lecció de punts i fletxes. En aquesta etapa no ens importa què són els punts, sinó com estan connectats. Aquesta idea es modela amb el concepte de graf. En general, un graf consta de punts (o objectes), anomenats \emph{vèrtexs}, alguns dels quals estan units per \emph{arestes}; aquí considerarem que els grafs són \emph{dirigits}, és a dir, que cada aresta apunta en una direcció, d'un punt a un altre, i per això fem servir fletxes per connectar punts. Entre dos punts hi pot haver moltes fletxes, o cap, i fins i tot pot haver-hi fletxes d'un punt a si mateix. Sovint la manera més fàcil de descriure un graf petit és dibuixar-lo, com a la figura següent.
\begin{equation}\label{graf:1}
\xymatrix{
1 \ar@(ul,ur)[]^a \ar@/^/[dr]^b
& & 2 \ar@/^/[dl]_c \ar@/_/[dr]^d \ar[ll]_e
& & 3 \ar[ll]_f \ar@(ur,dr)[]^g \\
& 4 \ar@/^/[ul]^h \ar[rr]^k
& & 5 \ar[ur]_l & 6
}
\end{equation}
El tipus de graf que tractem aquí és una versió del graf dirigit que sovint s'anomena \emph{multigraf dirigit amb bucles}, encara que farem servir la paraula \enquote{graf} per abreujar. En donem la definició formal.

\begin{definicio}\index{graf dirigit}
Un \textbf{graf} $G$ comprèn:
\begin{itemize}
\item un conjunt $V_G$ de \textbf{vèrtexs};
\item un conjunt $E_G$ d'\textbf{arestes};
\item dues aplicacions $s_G\colon E_G\to V_G$ i $t_G\colon E_G\to V_G$, anomenades \textbf{origen} i \textbf{destí}. Si $s_G(f)=a$ i $t_G(f)=b$, escrivim $a\xrightarrow{f}b$ o $f\colon a\to b$, i diem que $a$ és el vèrtex d'origen i $b$ el de destí.
\end{itemize}
La notació
\[
G=\xymatrix{
E_G \ar@<0.6ex>[r]^{s_G} \ar@<-0.6ex>[r]_{t_G} & V_G
}
\]
indica que $G$ és un graf amb vèrtexs a $V_G$ i arestes a $E_G$.
\end{definicio}

\begin{exemple}
Considerem el graf $G$ de la figura següent:
\[
\xymatrix{
	& \bullet_{a} \ar@<1ex>[dl]^f \ar@<1ex>[dr]^g &  \\
	\bullet_{b} \ar@(ul,dl)_l \ar@<1ex>[ur]^h & & \bullet_{c}  \ar@<1ex>[ll]^k \ar@(ur,dr)^m \\
}
\]
Especifica quins són els seus elements.
\end{exemple}
\begin{solucio}
Es té $V_G=\{a,b,c\}$ i $E_G=\{f,g,h,k,l,m\}$, i les dues aplicacions són
\[
    \begin{array}{c|c|c}
        E_G & s_G & t_G \\ \hline
        f & a & b  \\
        g & a & c \\
        h & b & a \\
        k & c & b \\
        l & b & b \\
        m & c & c
    \end{array}
\]
\end{solucio}

\begin{exemple}
Defineix el graf $G$ perquè la seva representació gràfica sigui \eqref{graf:1}.
\end{exemple}
\begin{solucio}
Prenem $V_G=\{1,2,3,4,5,6\}$, $E_G=\{a,b,c,d,e,f,g,h,k,l\}$ i les aplicacions
\[
    \begin{array}{c|c|c}
        E_G & s_G & t_G \\ \hline
        a & 1 & 1  \\
        b & 1 & 4 \\
        c & 2 & 4 \\
        d & 2 & 5 \\
        e & 2 & 1 \\
        f & 3 & 2 \\
        g & 3 & 3 \\
        h & 4 & 1 \\
        k & 4 & 5 \\
        l & 5 & 3
    \end{array}
\]
que reprodueixen la figura \eqref{graf:1}.
\end{solucio}

\begin{exemple}\leavevmode
\begin{enumerate}
\item Dibuixa el graf corresponent a les dades següents:
\[
    \begin{array}{c|c|c}
        E_G & s_G & t_G \\ \hline
        f & 2 & 3 \\
        g & 2 & 3 \\
        h & 2 & 3 \\
        k & 4 & 3 \\
        l & 6 & 3 \\
        m & 6 & 6
    \end{array}
    \hspace{2cm} V_G=\{1,2,3,4,5,6\}
\]
\item Escriu la taula corresponent al graf $H$ següent:
\[
\xymatrix{
	1 \ar[r]^{a} & 2 \ar[d]_{g} \ar@/^2ex/[r]^-{b} \ar[r]^-{c} & 3 \ar[r]^-{e} \ar@/^2ex/[l]^-{d} & 4 \\
	5 \ar@(ul,dl)_l & 6 \ar[l]_-{h} \ar[r]^-{k} & 7 \ar[ur]^-{f}
}
\]
\end{enumerate}
\end{exemple}
\begin{solucio}\leavevmode
\begin{enumerate}
\item
\[
\xymatrix{
	4 \ar[r]^-{k} & 3 & 6 \ar[l]_-{l} \ar@(ur,dr)\\
	1 & 2 \ar@/^2ex/[u]^-{f} \ar[u]^-{g} \ar@/_2ex/[u]_{h} & 5
}
\]
\item
\[
    \begin{array}{c|c|c}
        E_H & s_H & t_H \\ \hline
        a & 1 & 2 \\
        b & 2 & 3 \\
        c & 2 & 3 \\
        d & 3 & 2 \\
        e & 3 & 4 \\
        f & 7 & 4 \\
        g & 2 & 6 \\
        h & 6 & 5 \\
        k & 6 & 7 \\
        l & 5 & 5
    \end{array}
    \hspace{2cm} V_H=\{1,2,3,4,5,6,7\}
\]
\end{enumerate}
\end{solucio}

\begin{exemple}
En una empresa petita hi ha tres departaments: Administració, Comptabilitat i Atenció al client. Durant el funcionament diari, els documents poden circular diverses vegades entre els mateixos departaments, i alguns departaments processen documents interns sense enviar-los enlloc. En un dia es donen les situacions següents:
\begin{itemize}
\item Dos documents passen d'Administració a Comptabilitat.
\item Un document passa de Comptabilitat a Administració.
\item Tres documents passen de Comptabilitat a Atenció al client.
\item Un document passa d'Atenció al client a Comptabilitat.
\item Dos documents són processats internament a Administració.
\item Un document és processat internament a Comptabilitat.
\item Un document és processat internament a Atenció al client.
\end{itemize}
Modela aquesta situació com un graf $G$, especificant-ne els elements, interpretant-ne les arestes i representant-lo gràficament.
\end{exemple}
\begin{solucio}
El graf $G$ té com a conjunt de vèrtexs $V_G=\{a,b,c\}$, on $a,b,c$ representen Administració, Comptabilitat i Atenció al client, respectivament. El conjunt d'arestes és
\[
\begin{aligned}
E_G = \{
& a\xrightarrow{f_1}b,\ a\xrightarrow{f_2}b, \\
& b\xrightarrow{g}a, \\
& b\xrightarrow{h_1}c,\ b\xrightarrow{h_2}c,\ b\xrightarrow{h_3}c, \\
& c\xrightarrow{k}b, \\
& a\xrightarrow{j_1}a,\ a\xrightarrow{j_2}a, \\
& b\xrightarrow{i}b, \\
& c\xrightarrow{l}c
\},
\end{aligned}
\]
on els subíndexs indiquen arestes diferents entre els mateixos vèrtexs (diversos documents amb el mateix recorregut) i els bucles $a\xrightarrow{j}a$, $b\xrightarrow{i}b$ i $c\xrightarrow{l}c$ representen processament intern. La representació gràfica de $G$ és
\[
\xymatrix@R=7ex@C=8ex{
	c \ar@(ul,dl)_l \ar[r]^-{k} & b \ar@/^1.5ex/[l]^-{h_2} \ar@/_2.5ex/[l]_{h_1} \ar@/_5ex/[l]_{h_3}  \ar@(dr,ur)_i \ar@/^4ex/[d]^-{g} \\
	& a \ar@/^4ex/[u]^-{f_1} \ar@/^1ex/[u]_(.45){f_2}  \ar@(dl,dr)_{j_{1}}  \ar@(dr,ur)_{j_{2}}
}
\]
Aquest graf modela el flux real de documents dins l'empresa en un dia: les direccions indiquen el sentit del flux, el nombre d'arestes entre dos departaments indica la quantitat de documents, i els bucles indiquen activitat interna.
\end{solucio}

\begin{exercici}
Defineix el graf $G$ perquè la seva representació gràfica sigui la següent:
\[
\xymatrix{
	\bullet \ar@<1ex>[r]  & \bullet \ar@<1ex>[l] \ar@(ul,ur)  & \bullet \ar[l] \ar@(ur,dr) \ar[dl] & \\
	\bullet \ar[ur] \ar[r]  & \bullet  & \bullet
}
\]
\end{exercici}

\begin{exercici}
Dibuixa el graf corresponent a les dades següents:
\[
    \begin{array}{c|c|c}
        E_G & s_G & t_G \\ \hline
        f & v & w \\
        g & w & x \\
        h & w & x \\
        i & y & y \\
        j & y & z \\
        k & z & y
    \end{array}
    \hspace{2cm} V_G=\{v,w,x,y,z\}
\]
\end{exercici}

\begin{exercici}
En un centre logístic hi ha quatre seccions: Recepció, Magatzem, Preparació de comandes i Expedició. Els paquets es mouen entre les seccions al llarg del dia, i alguns processos es fan dins d'una mateixa secció. En una jornada s'observen els fluxos següents:
\begin{itemize}
\item Dos paquets passen de Recepció a Magatzem.
\item Un paquet passa de Magatzem a Recepció.
\item Tres paquets passen de Magatzem a Preparació de comandes.
\item Dos paquets passen de Preparació de comandes a Expedició.
\item Un paquet passa d'Expedició a Magatzem.
\item Dos paquets són reprocessats dins de Magatzem.
\item Un paquet és revisat dins de Preparació de comandes.
\item Un paquet és verificat dins d'Expedició.
\end{itemize}
Modela aquesta situació mitjançant un graf. Especifica'n els elements i representa'l gràficament.
\end{exercici}

\section{Homomorfismes de grafs}\label{sec:grafs-homomorfismes}

L'estructura d'un graf
\[
G=\xymatrix{
E_G \ar@<0.6ex>[r]^{s_G} \ar@<-0.6ex>[r]_{t_G} & V_G
}
\]
consisteix en dos conjunts $V_G$ i $E_G$ i dues aplicacions $s_G$ i $t_G$. Per relacionar dos grafs, hem de poder relacionar adequadament els seus conjunts i les seves aplicacions.

\begin{definicio}\index{morfisme de grafs}
Considerem dos grafs
\[
G=\xymatrix{
E_G \ar@<0.6ex>[r]^{s_G} \ar@<-0.6ex>[r]_{t_G} & V_G
}
\qquad\text{i}\qquad
H=\xymatrix{
E_H \ar@<0.6ex>[r]^{s_H} \ar@<-0.6ex>[r]_{t_H} & V_H
}
\]
Un \textbf{homomorfisme} $F$ de $G$ a $H$, que escrivim $F\colon G\to H$, és un parell d'aplicacions $F_V\colon V_G\to V_H$ i $F_E\colon E_G\to E_H$ tals que els diagrames següents commuten:
\begin{equation}\label{grafs:hom}
\xymatrix{
	E_G \ar[d]_-{F_E} \ar[r]^-{s_G} & V_G \ar[d]^-{F_V}\\
	E_H \ar[r]_-{s_H} & V_H
}
\hspace{2cm}
\xymatrix{
	E_G \ar[d]_-{F_E} \ar[r]^-{t_G} & V_G \ar[d]^-{F_V}\\
	E_H \ar[r]_-{t_H} & V_H
}
\end{equation}
\end{definicio}

Dit breument, un homomorfisme \emph{preserva} l'origen (diagrama de l'esquerra),
\[
F_V\comp s_G=s_H\comp F_E,
\]
i el destí (diagrama de la dreta),
\[
F_V\comp t_G=t_H\comp F_E.
\]

\begin{exemple}
Considerem els dos grafs següents:
\[
G:\ 
\xymatrix{
	1 \ar[d]_{h} \ar[r]^{f} & 2 \ar[d]^{g} \ar@(ur,ul)_{k} \\
	4 & 3
}
\qquad\text{i}\qquad
H:\ 
\xymatrix{
	b \ar@(ur,ul)_{y} \ar@(dr,ur)_{z} \\
	a \ar[u]_{x} \ar[r]_{v} & c
}
\]
Defineix un homomorfisme $F$ entre ells.
\end{exemple}
\begin{solucio}
Com que $k$ és un bucle i l'únic vèrtex de $H$ amb bucles és $b$, necessàriament $F_V(2)=b$. Com que $g$ surt de $2$, la seva imatge ha de sortir de $b$, i les úniques arestes que surten de $b$ són els bucles $y,z$; per tant $g$ també ha d'anar a un bucle i $F_V(3)=b$. Aquesta part és forçada; la resta és una elecció. Per construir un exemple concret, triem $F_V(1)=a$ i $F_V(4)=c$, i posem $F_E(k)=F_E(g)=y$, $F_E(f)=x$ i $F_E(h)=v$. Les taules següents comproven que aquesta tria preserva origen i destí.

La distinció entre una condició necessària i una elecció és important: que $F_V(2)=F_V(3)=b$ estigui forçat no vol dir que l'homomorfisme sigui únic. De fet, una enumeració finita directa de les aplicacions compatibles dona $24$ homomorfismes entre aquests dos grafs, dels quals $8$ tenen $F_V(1)=a$ i $16$ tenen $F_V(1)=b$.

Comprovem que les condicions \eqref{grafs:hom} es compleixen. Les taules dels dos grafs són
\[
    \begin{array}{c|c|c}
        E_G & s_G & t_G \\ \hline
        f & 1 & 2  \\
        g & 2 & 3 \\
        h & 1 & 4 \\
        k & 2 & 2 
    \end{array}
    \qquad
    \begin{array}{c|c|c}
        E_H & s_H & t_H \\ \hline
        v & a & c  \\
        x & a & b \\
        y & b & b \\
        z & b & b 
    \end{array}
\]
i les dues aplicacions de l'homomorfisme són
\[
F_V =
\begin{array}{c|c}
    V_G & V_H \\ \hline
     1 & a \\
     2 & b \\
     3 & b \\
     4 & c
\end{array}
\qquad
F_E =
\begin{array}{c|c}
   E_G  & E_H \\ \hline
   f & x \\
   g & y \\
   h & v \\
   k & y
\end{array}
\]
Per exemple,
\[
\begin{aligned}
(F_V\comp s_G)(f)&=F_V(s_G(f))\\
&=F_V(1)=a,\\
(s_H\comp F_E)(f)&=s_H(F_E(f))\\
&=s_H(x)=a,
\end{aligned}
\]
de manera que $(F_V\comp s_G)(f)=(s_H\comp F_E)(f)$; deixem la resta d'arestes com a exercici. Anàlogament,
\[
\begin{aligned}
(F_V\comp t_G)(h)&=F_V(t_G(h))\\
&=F_V(4)=c,\\
(t_H\comp F_E)(h)&=t_H(F_E(h))\\
&=t_H(v)=c,
\end{aligned}
\]
i per tant $(F_V\comp t_G)(h)=(t_H\comp F_E)(h)$.
\end{solucio}

\begin{exercici}
Considerem els grafs
\[
G:\ 
\xymatrix{
    1 \ar[r]^-{a} & 2 \ar[r]^-{b}& 3
}
\qquad
H:\ 
\xymatrix{
    4 \ar@<0.6ex>[r]^{d} \ar@<-0.6ex>[r]_{c} & 5 \ar@(ur,dr)^{e}
}
\]
Troba un homomorfisme $F$ de $G$ a $H$ i comprova les igualtats \eqref{grafs:hom}.
\end{exercici}

\begin{exercici}
Considerem el graf $G$ definit per
\[
    \begin{array}{c|c|c}
        E_G & s_G & t_G \\ \hline
        a & 1 & 2 \\
        b & 2 & 3 \\
        c & 1 & 4 \\
        d & 1 & 4 \\
        e & 5 & 6 
    \end{array}
    \hspace{2cm} V_G=\{1,2,3,4,5,6\}
\]
i el graf $H$ representat per
\[
\xymatrix{
	m \ar[d]_{y} \ar@<0.6ex>[r]^{w} & n \ar[l]^{x} \\
	p  & q \ar@(dr,ur)_{z}
}
\]
Es demana:
\begin{enumerate}
\item representar gràficament el graf $G$;
\item trobar un homomorfisme $F$ de $G$ a $H$ i comprovar les igualtats \eqref{grafs:hom}.
\end{enumerate}
\end{exercici}

\section{Camins en un graf}

Sigui $G$ un graf. Un \emph{camí} a $G$ és qualsevol successió d'arestes tal que el destí d'una aresta sigui l'origen de la següent. Això inclou successions de longitud~1, que són elements de $E_G$, i successions de longitud~0, que comencen i acaben al mateix vèrtex sense seguir cap aresta.

\begin{definicio}
Sigui $n\geq 1$. En un graf
\[
G=\xymatrix{
E_G \ar@<0.6ex>[r]^{s_G} \ar@<-0.6ex>[r]_{t_G} & V_G
}
\]
un \textbf{camí de longitud} $n$ del vèrtex $a$ al vèrtex $b$ és una successió d'arestes $(f_1,f_2,\dots,f_n)$, escrita en \emph{ordre de recorregut} ---primer es recorre $f_1$---, tal que
\begin{enumerate}
\item[(i)] $s_G(f_1)=a$;
\item[(ii)] $t_G(f_i)=s_G(f_{i+1})$ per a $i=1,\dots,n-1$;
\item[(iii)] $t_G(f_n)=b$.
\end{enumerate}
El conjunt de camins de longitud $n$ s'escriu $C_n$. En particular, $C_2$ és el conjunt de parells d'arestes $(f,g)$ amb $t_G(f)=s_G(g)$, anomenats parells d'\emph{arestes componibles}.
\end{definicio}

Per conveni, per a cada vèrtex $a\in V_G$ hi ha un únic camí de longitud~0 d'$a$ a si mateix, que denotem $()_a$ i anomenem \emph{camí buit} d'$a$; és el cas $n=0$, exclòs de la definició anterior perquè aleshores no hi ha cap aresta $f_1,\dots,f_n$. Si dibuixem un camí
\[
\xymatrix{
\cdot \ar[r]^{f_{1}} & \cdot \ar[r]^-{f_{2}} & \cdots \ar[r]^{f_{n-1}} & \cdot \ar[r]^{f_{n}} & \cdot
}
\]
l'aresta $f_1$ queda a l'esquerra i es recorre primer; $f_n$ queda a la dreta i es recorre última. Si $f\in E_G$, aleshores $(f)$ és un camí de longitud~1, és a dir, $(f)\in C_1$.

\begin{exemple}\label{ex:grafs-camins-G}
Considerem el graf $G$ de l'exemple d'homomorfisme de la secció anterior:
\[
G:\ 
\xymatrix{
	1 \ar[d]_{h} \ar[r]^{f} & 2 \ar[d]^{g} \ar@(ur,ul)_{k} \\
	4 & 3
}
\]
Escriu els camins de longitud $0$, $1$, $2$ i $3$ de $G$, i descriu tots els camins d'un vèrtex a un altre.
\end{exemple}
\begin{solucio}
Recordem la taula de $G$: $f\colon 1\to 2$, $g\colon 2\to 3$, $h\colon 1\to 4$ i el bucle $k\colon 2\to 2$.

\emph{Longitud $0$.} Hi ha un camí buit per a cada vèrtex: $C_0=\{()_1,\ ()_2,\ ()_3,\ ()_4\}$.

\emph{Longitud $1$.} Són les arestes: $C_1=\{(f),\ (g),\ (h),\ (k)\}$.

\emph{Longitud $2$.} Busquem els parells $(u,v)$ amb $t_G(u)=s_G(v)$. Com que $t_G(f)=t_G(k)=2$ i les arestes que surten de $2$ són $g$ i $k$, mentre que de $3$ i de $4$ no en surt cap,
\[
C_2=\{(f,g),\ (f,k),\ (k,g),\ (k,k)\}.
\]

\emph{Longitud $3$.} Cada camí de longitud $2$ que acaba a $2$ es pot allargar amb $g$ o amb $k$; els que acaben a $3$ no es poden allargar:
\[
C_3=\{(f,k,g),\ (f,k,k),\ (k,k,g),\ (k,k,k)\}.
\]

\emph{Descripció general.} Escrivim $k^n$ per a la successió $(k,\dots,k)$ de $n$ bucles, amb $k^0$ buida. Els camins d'un vèrtex a un altre són exactament:
\[
\begin{array}{c|l}
\text{d'origen a destí} & \text{camins} \\ \hline
1\to 1,\ 3\to 3,\ 4\to 4 & \text{només el camí buit} \\
1\to 2 & (f,k^n),\ n\geq 0 \\
1\to 3 & (f,k^n,g),\ n\geq 0 \\
1\to 4 & (h) \\
2\to 2 & k^n,\ n\geq 0 \ (\text{per a } n=0,\ \text{el camí buit } ()_2) \\
2\to 3 & (k^n,g),\ n\geq 0
\end{array}
\]
i no n'hi ha cap més. El bucle $k$ fa que el nombre de camins sigui infinit, però per a cada longitud $n\geq 2$ n'hi ha exactament quatre.
\end{solucio}

\begin{exemple}\label{ex:grafs-camins-H}
Fem el mateix amb el graf $H$ del mateix exemple:
\[
H:\ 
\xymatrix{
	b \ar@(ur,ul)_{y} \ar@(dr,ur)_{z} \\
	a \ar[u]_{x} \ar[r]_{v} & c
}
\]
Escriu els camins de longitud $0$, $1$ i $2$ de $H$, i compta quants camins de longitud $n$ hi ha.
\end{exemple}
\begin{solucio}
La taula de $H$ és $x\colon a\to b$, $v\colon a\to c$ i els bucles $y,z\colon b\to b$.

\emph{Longitud $0$ i $1$.} $C_0=\{()_a,\ ()_b,\ ()_c\}$ i $C_1=\{(x),\ (v),\ (y),\ (z)\}$.

\emph{Longitud $2$.} De $c$ no surt cap aresta, de manera que $v$ no es pot allargar. Després de $x$, $y$ o $z$ som a $b$, i podem continuar amb $y$ o amb $z$:
\[
C_2=\{(x,y),\ (x,z),\ (y,y),\ (y,z),\ (z,y),\ (z,z)\}.
\]

\emph{Recompte.} Un camí de longitud $n\geq 1$ de $b$ a $b$ és una successió qualsevol de $n$ bucles, triant a cada pas entre $y$ i $z$: n'hi ha $2^n$. Un camí de longitud $n\geq 2$ que surt d'$a$ ha de començar per $x$ i continuar amb $n-1$ bucles: n'hi ha $2^{n-1}$. L'únic camí que va d'$a$ a $c$ és $(v)$. Per tant, per a $n\geq 2$ hi ha
\[
2^n+2^{n-1}=3\cdot 2^{n-1}
\]
camins de longitud $n$; per exemple, $6$ de longitud $2$, d'acord amb la llista anterior.

Comparat amb l'exemple anterior, un sol bucle produeix un nombre constant de camins per a cada longitud, mentre que dos bucles al mateix vèrtex en produeixen un nombre que creix exponencialment.
\end{solucio}

\section{De grafs a categories}

Podem passar d'un graf a una categoria afegint-hi composició i identitats. La construcció següent és fonamental: mostra que tot graf genera una categoria de manera lliure.

Per a qualsevol graf $G$ hi ha una categoria $\Free(G)$ els objectes de la qual són els vèrtexs de $G$ i els morfismes de la qual són els camins de $G$. Un camí d'$a$ a $b$ és un morfisme $a\to b$. La composició s'obté concatenant camins, en \emph{ordre de recorregut}: si $p=(a_1,\dots,a_r)$ és un camí d'$A$ a $B$ i $q=(b_1,\dots,b_s)$ un camí de $B$ a $C$, la seva composició és el camí d'$A$ a $C$
\[
q\comp p=(a_1,\dots,a_r,b_1,\dots,b_s),
\]
on, d'acord amb la convenció categòrica $q\comp p$, el camí $p$ ---el que surt d'$A$--- es recorre primer. Per a cada vèrtex $A$, la identitat $\id_A$ és el camí buit $()_A$.

Comprovem que $\Free(G)$ és una categoria. La composició és associativa perquè la concatenació de successions ho és: en concatenar tres camins componibles, el resultat és la successió de totes les seves arestes en ordre, amb independència de l'ordre en què s'agrupin. El camí buit és neutre, ja que concatenar-lo no afegeix cap aresta: $(a_1,\dots,a_r)\comp()_A=(a_1,\dots,a_r)=()_B\comp(a_1,\dots,a_r)$ per a tot camí d'$A$ a $B$. Per tant, $\Free(G)$ satisfà els axiomes de categoria. S'anomena \textbf{categoria lliure generada pel graf} $G$, o també \textbf{categoria de camins} de $G$.\index{categoria!lliure}\index{categoria!de camins}

\begin{exemple}\label{ex:grafs-free-cadena}
Especifica la categoria $\Free(G)$ generada pel graf
\[
G:\ 
\xymatrix{
    1 \ar[r]^-{a} & 2 \ar[r]^-{b}& 3
}
\]
de l'exercici d'homomorfismes de la secció~\ref{sec:grafs-homomorfismes}.
\end{exemple}
\begin{solucio}
Els objectes són els vèrtexs $1$, $2$ i $3$. Els morfismes són els camins, i com que el graf no té cicles, n'hi ha un nombre finit:
\[
\begin{array}{c|l}
\Hom(1,1)=\{()_1\} & \Hom(1,2)=\{(a)\} \\
\Hom(2,2)=\{()_2\} & \Hom(2,3)=\{(b)\} \\
\Hom(3,3)=\{()_3\} & \Hom(1,3)=\{(a,b)\}
\end{array}
\]
i tots els altres homconjunts ($\Hom(2,1)$, $\Hom(3,1)$, $\Hom(3,2)$) són buits. En total, sis morfismes: les tres identitats, les dues arestes i el camí de longitud~$2$.

L'única composició que no involucra identitats és la de $(a)$ i $(b)$:
\[
\xymatrix{
1 \ar[r]^-{(a)} \ar[dr]_-{(a,b)} & 2 \ar[d]^-{(b)} \\
& 3
}
\qquad\qquad
(b)\comp(a)=(a,b).
\]
Les composicions amb identitats són les que imposen els axiomes, perquè els camins buits són neutres per a la concatenació.

Com que entre dos objectes hi ha com a màxim un morfisme, $\Free(G)$ és una categoria prima. De fet, és la categoria associada a l'ordre $1<2<3$ ---hi ha un morfisme $i\to j$ exactament quan $i\leq j$---, isomorfa, doncs, a la categoria $\cat{3}$ de l'ordre $0<1<2$ que apareix al capítol~\ref{cap:categories}.
\end{solucio}

\begin{exemple}\label{ex:grafs-free-bucle}
Especifica la categoria $\Free(H)$ generada pel graf
\[
H:\ 
\xymatrix{
    4 \ar@<0.6ex>[r]^{d} \ar@<-0.6ex>[r]_{c} & 5 \ar@(ur,dr)^{e}
}
\]
del mateix exercici.
\end{exemple}
\begin{solucio}
Els objectes són $4$ i $5$. Escrivim $e^n$ per al camí $(e,\dots,e)$ de $n$ bucles, amb $e^0=()_5$. Els camins de $H$ donen els homconjunts següents:
\[
\begin{aligned}
\Hom(4,4)&=\{()_4\},\\
\Hom(5,5)&=\{e^n : n\geq 0\},\\
\Hom(4,5)&=\{(c,e^n) : n\geq 0\}\cup\{(d,e^n) : n\geq 0\},\\
\Hom(5,4)&=\varnothing.
\end{aligned}
\]
De $5$ a $4$ no hi ha cap camí, perquè cap aresta surt de $5$ cap a $4$.

La composició és la concatenació. Per a bucles, $e^m\comp e^n=e^{m+n}$, i per als camins de $4$ a $5$,
\[
\xymatrix{
4 \ar[r]^-{(c,e^n)} \ar[dr]_-{(c,e^{n+m})} & 5 \ar[d]^-{e^m} \\
& 5
}
\qquad\qquad
e^m\comp(c,e^n)=(c,e^{n+m}),
\]
i igual amb $d$ en lloc de $c$.

Hi ha dues observacions a fer. Primera: $\Hom(5,5)$ amb la composició és el monoide $(\mathbb N,+,0)$, via $e^n\mapsto n$; la subcategoria plena sobre l'objecte $5$ és, doncs, la categoria d'un sol objecte d'aquest monoide. Segona: les arestes paral·leles $c$ i $d$ donen morfismes diferents, $(c)\neq(d)$, i també $(c,e^n)\neq(d,e^n)$ per a tot $n$. La categoria lliure no identifica mai dos camins diferents: no imposa cap relació que no vingui forçada pels axiomes.
\end{solucio}

\chapter{Convenció de mida}
\label{cap:mida}

Al llarg del llibre hem parlat de categories \emph{petites}, \emph{localment petites} i de \emph{classes pròpies} sense fixar del tot el marc conjuntista. Aquest apèndix el fa explícit, perquè el volum de lògica categòrica ---on apareixen categories de models, doctrines i fibracions--- l'exigirà.

\medskip
\noindent\textbf{Marc adoptat: classes i conjunts.} Treballem en una teoria de conjunts amb \textbf{classes pròpies}, a l'estil de von Neumann--Bernays--Gödel (NBG). Hi ha dos nivells de col·leccions: els \emph{conjunts}, que poden ser elements d'altres col·leccions, i les \emph{classes pròpies}, que no. Amb això:
\begin{itemize}
  \item una categoria és \textbf{petita} si la seva col·lecció d'objectes i la de morfismes són totes dues \emph{conjunts};
  \item és \textbf{localment petita} si tots els homconjunts $\Hom(A,B)$ són conjunts, tot i que la col·lecció d'objectes pot ser una classe pròpia;
  \item és \textbf{ben potenciada} si, per a cada objecte $A$, els subobjectes de $A$ formen un conjunt.
\end{itemize}
L'última condició demana una precisió. Un \emph{subobjecte} de $A$ no és un monomorfisme concret cap a $A$, sinó una classe d'equivalència de monomorfismes. Dos monomorfismes $m\colon S\rightarrowtail A$ i $n\colon T\rightarrowtail A$ són equivalents quan cadascun factoritza a través de l'altre, és a dir, quan difereixen en un isomorfisme entre els dominis. La distinció és essencial, i en el marc de classes cal formular-la amb cura. A $\cat{Set}$, els monomorfismes cap a un conjunt donat formen una classe pròpia, perquè cada subconjunt té moltes còpies isomorfes; i cada classe d'equivalència és, ella mateixa, una classe pròpia. Com que les classes pròpies no poden ser elements de cap col·lecció, \enquote{el conjunt de les classes d'equivalència} no té sentit literal. La definició precisa és, doncs, la següent: una categoria és \textbf{ben potenciada} si, per a cada objecte $A$, hi ha un \emph{conjunt} $R_A$ de monomorfismes cap a $A$ tal que tot monomorfisme cap a $A$ és equivalent a algun element de $R_A$. Aleshores es defineix $\Sub(A)$ com el conjunt quocient $R_A/{\sim}$, i $[m]$ designa la classe dels representants equivalents a $m$. Un altre conjunt de representants dona un quocient en bijecció canònica amb aquest, i la bijecció respecta l'ordre; per això el resultat no depèn de la tria de $R_A$. A $\cat{Set}$ es pot prendre com a $R_A$ el conjunt de les inclusions dels subconjunts de $A$, i s'obté $\Sub(A)\cong\mathcal P(A)$. La resta del llibre parla de \enquote{classes de monomorfismes} en aquest sentit.

Les categories $\cat{Set}$, $\cat{Grp}$, $\cat{Top}$, $\cat{Graph}$ i $\cat{Cat}$ són localment petites però no petites; cada preordre concret sobre un conjunt i cada categoria finita concreta són petites. Això no s'ha de confondre amb la petitesa de la categoria de \emph{tots} els preordres o de \emph{totes} les categories finites, que ---en tenir una classe pròpia de conjunts subjacents--- són grans llevat que es fixi una convenció de mida o s'esculli un conjunt de representants.

Quan alguna construcció demani triar simultàniament un objecte (o un morfisme) per a cada objecte d'una categoria \emph{gran} ---per exemple, en construir un invers d'una equivalència---, la família de tries està indexada per una classe pròpia, i l'axioma de l'elecció ordinari, que val per a famílies indexades per un conjunt, no hi arriba. En aquests casos assumim la forma \textbf{global} de l'elecció (disponible, per exemple, a NBG amb l'axioma d'elecció global). Ho indicarem allà on calgui.

\medskip
\noindent\textbf{L'alternativa dels universos.} Una segona opció, molt usada en geometria algebraica i teoria de topos, fixa un o més \textbf{universos de Grothendieck} $\mathcal{U}$ ---conjunts prou grans i tancats per les operacions habituals--- i parla de categories $\mathcal{U}$-petites i $\mathcal{U}$-grans de manera \emph{relativa} a $\mathcal{U}$. Aquesta convenció té l'avantatge que \enquote{la categoria de tots els $\mathcal{U}$-conjunts} és, ella mateixa, un objecte d'un univers més gran, cosa que simplifica certes construccions. En aquest volum \emph{no} l'adoptem, per no carregar una introducció amb la teoria d'universos; però el lector que continuï cap a la teoria de feixos la retrobarà. Les dues opcions són maneres alternatives de \emph{gestionar la mida} més que no pas fonaments idèntics: a grans trets, el paper de les \enquote{classes pròpies} de NBG el fa aquí la distinció entre $\mathcal{U}$-petit i $\mathcal{U}$-gran relativa a un univers $\mathcal{U}$ fixat. Fer precisa aquesta correspondència exigeix especificar l'univers i les hipòtesis pertinents; no cal, però, que això sigui un obstacle per a una primera lectura.

\medskip
\noindent Fixar aquesta convenció des d'ara evita l'ambigüitat que, altrament, es faria notar quan apareguin categories de models i functors classificadors: en tot moment, cada construcció d'aquest llibre declara si viu en un conjunt o en una classe pròpia.

\medskip
\noindent\textbf{Per saber-ne més.} Les qüestions de mida es tracten, al nivell d'aquest llibre, a \cite[§1.8]{awodey}, \cite[§1.1]{riehl} ---on l'observació~1.1.5 presenta la convenció d'universos---, \cite[§3.2]{leinster} i \cite[cap.~1]{perrone}. El capítol~3 de \cite{leinster} presenta, a més, la perspectiva en què el mateix llenguatge estructural fa de fonament, sense una teoria de conjunts prèvia; en són les fonts \cite{lawvere-etcs} i, com a introducció informal, \cite{leinster-rethinking}. Per a una comparació detallada de les diferents maneres de gestionar la mida ---classes, universos i altres alternatives---, adequada un cop se'n té una certa experiència, vegeu \cite{shulman-mida}.

\chapter{Solucions dels tests}
\label{cap:solucions-tests}

\ifimprimible
Aquest apèndix reuneix les respostes i les justificacions dels tests. En aquesta edició per imprimir, són la manera de corregir-los.
\else
Aquest apèndix reuneix les respostes i les justificacions dels tests interactius. Serveixen per als lectors de PDF que no admeten la correcció automàtica i per a l'ús en paper.
\fi Es recomana consultar-les només després d'haver respost tot el test: cada justificació explica per què la resposta correcta ho és i, sovint, on falla el distractor més temptador.

\imprimeixsolucionstests

\ifpublicacioparcial\else
\chapter[Quasiexemples i contraexemples]{Quasiexemples\\ i contraexemples}
\label{cap:contraexemples}

Una definició matemàtica s'entén de veritat quan se'n coneixen tant els exemples com els \emph{quasiexemples}: estructures que s'hi assemblen molt, que apareixen de manera natural, i que tanmateix fallen en un punt precís. Aquest apèndix recull quasiexemples i contraexemples per a les nocions centrals del llibre: categoria, functor, transformació natural i equivalència. Acaba amb models de la lògica que queden fora de l'explicació per topos dels capítols~\ref{cap:subobjectes}--\ref{cap:pont}.

La secció~\ref{sec:igualtat-equivalencia}, sobre igualtat, isomorfisme, isomorfisme natural i equivalència, és la més important. Confondre aquests nivells és un error freqüent, i greu, perquè porta a demostrar coses falses o a formular propietats que no tenen sentit categòric.

El nivell és el del llibre, i els exemples s'han triat perquè només demanin els coneixements que el llibre ja pressuposa. Tota afirmació que es pot justificar amb aquests coneixements es demostra. Quan un comentari remet a teories que el llibre no desenvolupa, es presenta de manera informativa i es diu explícitament.

% =====================================================================
\section{Què no és una categoria}
\label{sec:no-categoria}
% =====================================================================

\subsection{Les condicions que es poden trencar}

La definició~\ref{def:categoria} demana unes dades i dos axiomes. Per poder dir, exemple a exemple, quina condició falla, les desglossem en cinc:
\begin{enumerate}
\item[(C1)] \textbf{Extrems.} Cada morfisme $f$ té un \emph{únic} domini i un \emph{únic} codomini.
\item[(C2)] \textbf{Composició.} Per a cada parella componible $f\colon A\to B$, $g\colon B\to C$ hi ha un morfisme $g\comp f\colon A\to C$, determinat sense ambigüitat.
\item[(C3)] \textbf{Associativitat.} $h\comp(g\comp f)=(h\comp g)\comp f$.
\item[(C4)] \textbf{Existència d'identitats.} Per a cada objecte $A$ hi ha un morfisme $\id_A\colon A\to A$.
\item[(C5)] \textbf{Lleis d'identitat.} $f\comp\id_A=f=\id_B\comp f$ per a tot $f\colon A\to B$.
\end{enumerate}
Les igualtats de (C3) i (C5) són \emph{igualtats}: no isomorfismes, ni homotopies, ni cap mena d'equivalència. Gairebé tots els exemples d'aquesta secció fallen exactament aquí. Recordem també que la demostració de la unicitat de la identitat (capítol~\ref{cap:categories}) no fa servir l'associativitat: només les lleis d'identitat. Ho aprofitarem en parlar dels magmes.

\subsection{L'exemple central: els camins en un espai}

Aquest exemple fa servir la noció de continuïtat. Sigui $X$ un espai topològic. Proposem:
\begin{itemize}
\item \emph{objectes}: els punts de $X$;
\item \emph{morfismes}: els camins, és a dir, les aplicacions contínues $\gamma\colon[0,1]\to X$; escrivim $\gamma\colon x\to y$ si $\gamma(0)=x$ i $\gamma(1)=y$;
\item \emph{composició}: la concatenació. Si $\gamma\colon x\to y$ i $\delta\colon y\to z$, definim $\delta\ast\gamma\colon x\to z$ per
\[
(\delta\ast\gamma)(t)=
\begin{cases}
\gamma(2t), & 0\le t\le\tfrac12,\\
\delta(2t-1), & \tfrac12\le t\le1;
\end{cases}
\]
\item \emph{candidats a identitat}: els camins constants $c_x(t)=x$.
\end{itemize}
Escrivim $\delta\ast\gamma$ (primer $\gamma$, després $\delta$) per coherència amb $g\comp f$. Les dues branques coincideixen a $t=\frac12$, perquè $\gamma(1)=y=\delta(0)$, i per tant $\delta\ast\gamma$ és contínua. Així (C1) i (C2) es compleixen. El problema és a (C3) i (C5).

\paragraph{Fallada de l'associativitat.} Siguin $\gamma\colon x\to y$, $\delta\colon y\to z$ i $\varepsilon\colon z\to w$. Un càlcul directe dona
\[
\bigl((\varepsilon\ast\delta)\ast\gamma\bigr)(t)=
\begin{cases}
\gamma(2t), & t\in[0,\frac12],\\
\delta(4t-2), & t\in[\frac12,\frac34],\\
\varepsilon(4t-3), & t\in[\frac34,1],
\end{cases}
\qquad
\bigl(\varepsilon\ast(\delta\ast\gamma)\bigr)(t)=
\begin{cases}
\gamma(4t), & t\in[0,\frac14],\\
\delta(4t-1), & t\in[\frac14,\frac12],\\
\varepsilon(2t-1), & t\in[\frac12,1].
\end{cases}
\]
Tots dos recorren els mateixos tres trams en el mateix ordre, però a velocitats diferents.

\begin{exemple}[Tres segments a la recta]\label{ex:camins-recta}
A $X=\mathbb R$, prenem $\gamma(t)=t$, $\delta(t)=1+t$ i $\varepsilon(t)=2+t$, de manera que $\gamma\colon0\to1$, $\delta\colon1\to2$ i $\varepsilon\colon2\to3$. Avaluant a $t=\frac12$,
\[
\bigl((\varepsilon\ast\delta)\ast\gamma\bigr)(\tfrac12)=\gamma(1)=1,
\qquad
\bigl(\varepsilon\ast(\delta\ast\gamma)\bigr)(\tfrac12)=\varepsilon(0)=2.
\]
Els dos camins van de $0$ a $3$ i tenen la mateixa imatge, l'interval $[0,3]$, però són diferents: falla (C3).
\end{exemple}

\paragraph{Fallada de les identitats.} El fracàs és més radical que el de l'associativitat:

\begin{proposicio}\label{prop:camins-sense-identitat}
Sigui $\gamma\colon x\to y$ un camí injectiu. Aleshores cap camí $e\colon y\to y$ compleix $e\ast\gamma=\gamma$. Anàlogament, si $\gamma\colon y\to z$ és injectiu, cap camí $e\colon y\to y$ compleix $\gamma\ast e=\gamma$.
\end{proposicio}

\begin{prova}
Per a qualsevol $e$, $(e\ast\gamma)(\frac14)=\gamma(\frac12)$, que és diferent de $\gamma(\frac14)$ perquè $\gamma$ és injectiu. Per a l'altra afirmació, $(\gamma\ast e)(\frac34)=\gamma(\frac12)\neq\gamma(\frac34)$.
\end{prova}

A $X=\mathbb R$, entre dos punts diferents qualssevol hi ha camins injectius. Per tant, no és que s'hagin triat malament els camins constants: (C4)--(C5) no es poden satisfer de cap manera amb aquesta composició.

\paragraph{Però fallen de manera controlada.} Les igualtats que fallen valen \emph{llevat d'una reparametrització}, i les reparametritzacions són inofensives llevat d'homotopia. Dos camins $\alpha,\beta\colon x\to y$ són \emph{homòtops relativament als extrems}, i escrivim $\alpha\simeq\beta$, si hi ha una aplicació contínua $H\colon[0,1]\times[0,1]\to X$ amb $H(t,0)=\alpha(t)$, $H(t,1)=\beta(t)$, $H(0,s)=x$ i $H(1,s)=y$.

\begin{lema}[Reparametrització]\label{lema:reparametritzacio}
Sigui $\alpha\colon[0,1]\to X$ un camí i $\varphi\colon[0,1]\to[0,1]$ contínua amb $\varphi(0)=0$ i $\varphi(1)=1$. Aleshores $\alpha\comp\varphi\simeq\alpha$.
\end{lema}

\begin{prova}
L'aplicació $H(t,s)=\alpha\bigl((1-s)\varphi(t)+st\bigr)$ és contínua i ben definida, perquè $(1-s)\varphi(t)+st\in[0,1]$. Val $\alpha\comp\varphi$ per a $s=0$ i $\alpha$ per a $s=1$, i $H(0,s)=\alpha(0)$, $H(1,s)=\alpha(1)$ per a tot $s$.
\end{prova}

\begin{proposicio}\label{prop:associador-camins}
\begin{enumerate}
\item[(a)] $(\varepsilon\ast\delta)\ast\gamma=\bigl(\varepsilon\ast(\delta\ast\gamma)\bigr)\comp\varphi$, on $\varphi(t)=t/2$ a $[0,\frac12]$, $\varphi(t)=t-\frac14$ a $[\frac12,\frac34]$ i $\varphi(t)=2t-1$ a $[\frac34,1]$.
\item[(b)] $c_y\ast\gamma=\gamma\comp\psi$ amb $\psi(t)=\min(2t,1)$, i $\gamma\ast c_x=\gamma\comp\psi'$ amb $\psi'(t)=\max(2t-1,0)$.
\end{enumerate}
En conseqüència, $(\varepsilon\ast\delta)\ast\gamma\simeq\varepsilon\ast(\delta\ast\gamma)$, $c_y\ast\gamma\simeq\gamma$ i $\gamma\ast c_x\simeq\gamma$.
\end{proposicio}

\begin{prova}
Les igualtats (a) i (b) es comproven tram a tram amb les fórmules anteriors; per exemple, a $[\frac12,\frac34]$, $\varphi(t)=t-\frac14\in[\frac14,\frac12]$ i $\delta\bigl(4(t-\frac14)-1\bigr)=\delta(4t-2)$. Les funcions $\varphi$, $\psi$ i $\psi'$ són contínues i fixen $0$ i $1$, i el lema~\ref{lema:reparametritzacio} dona les homotopies.
\end{prova}

\paragraph{Primer remei: passar al quocient.} Si $\gamma\simeq\gamma'$ i $\delta\simeq\delta'$, concatenant les homotopies per a cada $s$ fix s'obté $\delta\ast\gamma\simeq\delta'\ast\gamma'$. Per tant, la concatenació indueix una composició ben definida entre classes, $[\delta]\comp[\gamma]:=[\delta\ast\gamma]$, i la proposició~\ref{prop:associador-camins} esdevé exactament (C3) i (C5) entre classes, amb $\id_x=[c_x]$. S'obté una categoria, el \textbf{grupoide fonamental} $\Pi_1(X)$\index{grupoide fonamental}: els objectes són els punts de $X$ i $\Hom(x,y)$ és el conjunt de classes d'homotopia de camins de $x$ a $y$. És un grupoide perquè el camí recorregut a l'inrevés, $\bar\gamma(t)=\gamma(1-t)$, és un invers de $[\gamma]$ (la comprovació és una homotopia explícita que no reproduïm). L'homconjunt $\Hom(x,x)$ és el que en topologia s'anomena \emph{grup fonamental} de $X$ en el punt $x$ (informatiu).

\paragraph{Segon remei: canviar els morfismes.} L'origen del problema és que tots els camins viuen en l'interval $[0,1]$, i cada concatenació ha de comprimir el temps a la meitat. Si cada camí té la seva durada, el problema desapareix. Un \textbf{camí de Moore}\index{camí de Moore} és un parell $(r,\gamma)$ amb $r\ge0$ i $\gamma\colon[0,r]\to X$ contínua, i es concatena sumant les durades:
\[
(s,\delta)\comp(r,\gamma)=(r+s,\ \delta\star\gamma),\qquad
(\delta\star\gamma)(t)=
\begin{cases}
\gamma(t), & t\in[0,r],\\
\delta(t-r), & t\in[r,r+s].
\end{cases}
\]

\begin{proposicio}
Els punts de $X$ i els camins de Moore formen una categoria, amb identitats $\id_x=(0,c_x)$.
\end{proposicio}

\begin{prova}
Les dues maneres d'agrupar tres camins $(r,\gamma)$, $(s,\delta)$ i $(u,\varepsilon)$ donen el camí de durada $r+s+u$ que val $\gamma(t)$ a $[0,r]$, $\delta(t-r)$ a $[r,r+s]$ i $\varepsilon(t-r-s)$ a $[r+s,r+s+u]$: \emph{la mateixa funció}. Les lleis d'identitat també són literals, perquè un camí de durada $0$ no ocupa temps: $(r,\gamma)\comp(0,c_x)$ té durada $r$ i val $\gamma(t-0)=\gamma(t)$.
\end{prova}

Tenim, doncs, les dues estratègies que reapareixen constantment en matemàtiques: \emph{quocientar}, és a dir, oblidar la informació que impedeix l'estrictesa (com a $\Pi_1(X)$), i \emph{estrictificar}, és a dir, substituir la composició per una altra que conservi la informació que interessa i en la qual les lleis valguin literalment (com amb els camins de Moore). Aquí no es pot parlar d'equivalència de categories, perquè l'estructura de partida no és una categoria. Els camins d'arestes d'un graf es concatenen sense problemes, amb el camí buit com a identitat ---és la categoria lliure de la definició~\ref{def:free}--- per la mateixa raó: la longitud s'hi suma i no es renormalitza mai.

\paragraph{Tercera via: fer teoria de la feblesa (informatiu).} La proposició~\ref{prop:associador-camins} no diu només que les dues associacions són homòtopes: en dona una homotopia \emph{concreta}, l'\emph{associador}. Amb quatre camins hi ha cinc maneres de posar els parèntesis, i els associadors les connecten formant un pentàgon. Si s'exigeix que els dos camins del pentàgon donin \emph{exactament} el mateix resultat, s'obté l'estructura d'una \emph{bicategoria}. Si només es demana que coincideixin llevat d'una nova homotopia, i així successivament a tots els nivells, s'arriba a les $\infty$-categories. Cap de les dues no es tracta en aquest llibre.

\subsection{Un sol objecte: els magmes}

Un \textbf{magma}\index{magma} és un conjunt $M$ amb una operació binària qualsevol, $(a,b)\mapsto ab$; un \emph{semigrup} és un magma associatiu, i un monoide, un semigrup amb element neutre. Com al capítol~\ref{cap:categories} per als monoides, a tot magma li podem associar una estructura $\cat{B}M$ amb un sol objecte $\ast$, un morfisme $\ast\to\ast$ per a cada $a\in M$ i la composició $a\comp b=ab$. Les condicions (C1) i (C2) es compleixen sempre, i $\cat{B}M$ és una categoria si i només si $M$ és un monoide. Cada magma no associatiu, o sense neutre, és un quasiexemple de categoria amb un objecte.

\begin{exemple}[La resta d'enters]
$M=(\mathbb Z,-)$. No és associatiu: $(1-1)-1=-1$, però $1-(1-1)=1$. Té neutre per la dreta, $a-0=a$, però no per l'esquerra: $e-a=a$ per a tot $a$ obligaria $e=2a$ per a tot $a$. Pel raonament de la unicitat de la identitat, un neutre bilateral hauria de coincidir amb $0$, que no ho és. Així, (C5) es compleix a mitges.
\end{exemple}

\begin{exemple}[Pedra, paper, tisora]\label{ex:ppt}
$M=\{\mathrm{Pe},\mathrm{Pa},\mathrm{Ti}\}$, on $xy$ és el guanyador de la partida entre $x$ i $y$, i $xx=x$:
\[
\begin{array}{c|ccc}
 & \mathrm{Pe} & \mathrm{Pa} & \mathrm{Ti}\\\hline
\mathrm{Pe} & \mathrm{Pe} & \mathrm{Pa} & \mathrm{Pe}\\
\mathrm{Pa} & \mathrm{Pa} & \mathrm{Pa} & \mathrm{Ti}\\
\mathrm{Ti} & \mathrm{Pe} & \mathrm{Ti} & \mathrm{Ti}
\end{array}
\]
És commutatiu i idempotent, però $(\mathrm{Pe}\,\mathrm{Pa})\,\mathrm{Ti}=\mathrm{Pa}\,\mathrm{Ti}=\mathrm{Ti}$ i $\mathrm{Pe}\,(\mathrm{Pa}\,\mathrm{Ti})=\mathrm{Pe}\,\mathrm{Ti}=\mathrm{Pe}$. Tampoc no té neutre: $\mathrm{Pe}\,e=\mathrm{Pe}$ exigeix $e\in\{\mathrm{Pe},\mathrm{Ti}\}$, $\mathrm{Pa}\,e=\mathrm{Pa}$ exigeix $e\in\{\mathrm{Pe},\mathrm{Pa}\}$ i $\mathrm{Ti}\,e=\mathrm{Ti}$ exigeix $e\in\{\mathrm{Pa},\mathrm{Ti}\}$, i la intersecció és buida. És un model finit, comprovable a mà, on fallen alhora (C3), (C4) i (C5).
\end{exemple}

\begin{exemple}[Altres magmes naturals]
\begin{itemize}
\item \emph{El punt mitjà}: $\mathbb R$ amb $ab=\frac{a+b}2$. Per exemple, $(0\cdot0)\cdot4=2$ i $0\cdot(0\cdot4)=1$; no hi ha neutre, perquè $a\cdot e=a$ força $e=a$.
\item \emph{El producte vectorial}: $\mathbb R^3$ amb $a\times b$. $(\mathbf i\times\mathbf i)\times\mathbf j=\mathbf 0$, però $\mathbf i\times(\mathbf i\times\mathbf j)=\mathbf i\times\mathbf k=-\mathbf j$. Satisfà, en lloc de l'associativitat, la identitat de Jacobi: és l'exemple bàsic d'àlgebra de Lie, una estructura perfectament bona que no és una categoria d'un objecte.
\item \emph{La coma flotant} (com a il·lustració): la suma de l'aritmètica de coma flotant d'un ordinador no és, en general, associativa. En doble precisió, una avaluació típica dona $(0.1+0.2)+0.3=0.6000000000000001$, mentre que $0.1+(0.2+0.3)=0.6$. Per convertir-ho en un magma literal caldria precisar el conjunt de valors i el tractament dels casos excepcionals (desbordaments, zeros amb signe...), cosa que no fem.
\item \emph{Associativitat sense neutre}: $(2\mathbb Z,\cdot)$ és un semigrup, però l'únic candidat a neutre seria $1\notin2\mathbb Z$. Aquesta fallada es repara de franc, afegint formalment un neutre.
\end{itemize}
\end{exemple}

\paragraph{Quocientar no sempre és un bon remei.} Davant d'un magma no associatiu podríem intentar el truc dels camins: quocientar per la relació d'equivalència més petita, compatible amb l'operació, que la faci associativa. Però el resultat pot ser catastròfic.

\begin{proposicio}
El quocient associatiu del magma pedra-paper-tisora té un sol element.
\end{proposicio}

\begin{prova}
En el quocient, $\mathrm{Ti}=(\mathrm{Pe}\,\mathrm{Pa})\,\mathrm{Ti}=\mathrm{Pe}\,(\mathrm{Pa}\,\mathrm{Ti})=\mathrm{Pe}$ i $\mathrm{Pe}=(\mathrm{Pa}\,\mathrm{Ti})\,\mathrm{Pe}=\mathrm{Pa}\,(\mathrm{Ti}\,\mathrm{Pe})=\mathrm{Pa}$.
\end{prova}

La diferència amb els camins és significativa. Allà la manca d'associativitat era \emph{coherent}, testimoniada per homotopies canòniques, i el quocient conservava la informació essencial. Aquí és arbitrària, i quocientar ho esborra tot. La jerarquia de les estructures que hem vist és aquesta:
\[
\renewcommand{\arraystretch}{1.2}
\begin{array}{lll}
\text{un sol objecte} & \text{diversos objectes} & \text{condicions}\\\hline
\text{magma} & \text{magmoide} & \text{(C1), (C2)}\\
\text{semigrup} & \text{semicategoria} & \text{(C1)--(C3)}\\
\text{monoide} & \text{categoria} & \text{(C1)--(C5)}\\
\text{grup} & \text{grupoide} & \text{(C1)--(C5) i inversos}
\end{array}
\]

\subsection{Altres maneres de fallar}

\begin{exemple}[Ordre estricte: una semicategoria]
Objectes: els nombres reals; un únic morfisme $x\to y$ si $x<y$, i cap altrament. La composició existeix per transitivitat i és associativa, perquè entre dos objectes hi ha com a molt un morfisme. Però no hi ha cap morfisme $x\to x$: falla (C4). Substituint $<$ per $\le$ s'obté una categoria.
\end{exemple}

\begin{exemple}[Grafs: no hi ha composició]
Un graf dirigit té vèrtexs i arestes, però dues arestes consecutives $x\to y\to z$ no tenen per què donar una aresta $x\to z$: falla (C2). El remei és la categoria lliure del graf (definició~\ref{def:free}).
\end{exemple}

\begin{exemple}[Aplicacions lineals no nul·les]
Objectes: els espais vectorials reals no nuls; morfismes: les aplicacions lineals \emph{no nul·les}. Les identitats hi són, i la composició és associativa quan està definida, però no sempre ho està: a $\mathbb R^2$, les projeccions $p(x,y)=(x,0)$ i $q(x,y)=(0,y)$ són no nul·les i $q\comp p=0$. Falla (C2).
\end{exemple}

\begin{exemple}[Funcions identificades amb el seu graf]
Si una funció és només el seu graf, $\{(1,1)\}$ és alhora una funció $\{1\}\to\{1\}$ i una funció $\{1\}\to\{1,2\}$: el codomini no queda determinat i falla (C1). En particular, \enquote{ser exhaustiva} deixa de ser una propietat de la funció. És el mateix fenomen que l'observació~\ref{obs:extrems-implicits} assenyala per a $\cat{Rel}$, i el remei és el mateix: el domini i el codomini formen part de les dades del morfisme.
\end{exemple}

\begin{exemple}[Un quocient que no és compatible amb la composició]
Al monoide de les funcions contínues $\mathbb R\to\mathbb R$ amb la composició, identifiquem $f\sim g$ si $f(0)=g(0)$, i intentem definir $[g]\comp[f]=[g\comp f]$. Amb $f(x)=x+1$, $g(x)=x$ i $g'(x)=2x$ tenim $g\sim g'$, però $(g\comp f)(0)=1\neq2=(g'\comp f)(0)$. La composició no està ben definida i falla (C2). Per això, en construir $\Pi_1(X)$, calia comprovar que l'homotopia és compatible amb la concatenació: no tot quocient d'una categoria és una categoria.
\end{exemple}

\subsection{Falsos contraexemples}

Per equilibrar, convé recordar estructures que \emph{semblen} fallar algun axioma però que són categories genuïnes. Ensenyen que els morfismes no han de ser funcions, ni estar definits a tot arreu.

\begin{exemple}[Funcions parcials]
Objectes: els conjunts; morfismes $f\colon A\rightharpoonup B$: les funcions definides en un subconjunt $D_f\subseteq A$. La composició $g\comp f$ és definida en $\{a\in D_f\mid f(a)\in D_g\}$. Les dues maneres d'agrupar $h$, $g$ i $f$ tenen el mateix domini de definició, $\{a\in D_f\mid f(a)\in D_g,\ g(f(a))\in D_h\}$, i hi prenen el mateix valor; les identitats són les funcions totals identitat. Obtenim la categoria $\cat{Par}$: que un morfisme no estigui definit a tot arreu no impedeix res.
\end{exemple}

Les relacions (la categoria $\cat{Rel}$ del capítol~\ref{cap:categories}) i les matrius ($\cat{Mat}_{\mathbb K}$, exemple~\ref{ex:mat-finvect}) són dos altres casos: els morfismes no són funcions, i en el segon cas ni tan sols són transformacions de res, sinó taules de nombres.

\subsection{Resum}

\[
\renewcommand{\arraystretch}{1.15}
\begin{array}{lccccl}
\text{Exemple} & \text{(C1)} & \text{(C2)} & \text{(C3)} & \text{(C4--5)} & \text{Remei}\\\hline
\text{camins sobre }[0,1] & \text{sí} & \text{sí} & \text{no} & \text{no} & \Pi_1(X)\text{ o camins de Moore}\\
(\mathbb Z,-) & \text{sí} & \text{sí} & \text{no} & \text{a mitges} & \text{---}\\
\text{pedra-paper-tisora} & \text{sí} & \text{sí} & \text{no} & \text{no} & \text{(el quocient és trivial)}\\
(2\mathbb Z,\cdot) & \text{sí} & \text{sí} & \text{sí} & \text{no} & \text{afegir un neutre}\\
(\mathbb R,<) & \text{sí} & \text{sí} & \text{sí} & \text{no} & (\mathbb R,\le)\\
\text{graf dirigit} & \text{sí} & \text{no} & \text{---} & \text{---} & \text{categoria lliure}\\
\text{lineals no nul·les} & \text{sí} & \text{no} & \text{sí} & \text{sí} & \text{admetre el }0\\
\text{funcions com a grafs} & \text{no} & \text{sí} & \text{sí} & \text{sí} & \text{morfismes amb extrems}
\end{array}
\]
Els axiomes de categoria no són convencions arbitràries: n'hi ha prou d'agafar l'estructura més natural possible ---els camins en un espai--- per veure'ls fallar. Quan fallen, hi ha tres actituds possibles: quocientar, estrictificar o acceptar la feblesa i fer-ne teoria.

% =====================================================================
\section{Què no és un functor}
\label{sec:no-functor}
% =====================================================================

Un functor ha de superar dues proves, en aquest ordre (observació~\ref{obs:comprovar-tipus}): els tipus i, després, les equacions $F(\id_A)=\id_{F(A)}$ i $F(g\comp f)=F(g)\comp F(f)$. Els gairebé-functors fallen en una de les dues.

\paragraph{Fallada de tipus.} El capítol~\ref{cap:functors} ja en dona l'exemple bàsic: l'antiimatge $f\mapsto f^{-1}$ no pot ser un functor covariant de $\cat{Set}$ en $\cat{Set}$, perquè va en sentit contrari; sí que ho és com a functor contravariant (exemple~\ref{ex:parts-dues-variancies}). Un altre exemple: l'assignació que envia cada grup $G$ al seu grup d'automorfismes $\mathrm{Aut}(G)$ no té cap manera natural d'actuar sobre els homomorfismes. Un homomorfisme $f\colon G\to H$ no indueix cap aplicació raonable $\mathrm{Aut}(G)\to\mathrm{Aut}(H)$; només ho fa quan $f$ és un isomorfisme, que és la manera de convertir-la en un functor sobre la subcategoria dels isomorfismes.

\paragraph{Una assignació sobre objectes que no admet cap functor.} El cas més instructiu és el d'un objecte que sembla \enquote{canònic} però no és functorial de cap manera.

\begin{proposicio}[El centre d'un grup no és functorial]\label{prop:centre-no-functor}
No hi ha cap functor $Z\colon\cat{Grp}\to\cat{Grp}$ que enviï cada grup $G$ al seu centre $Z(G)=\{z\in G\mid zg=gz \text{ per a tot } g\}$, sigui quina sigui la seva acció sobre els homomorfismes.
\end{proposicio}

\begin{prova}
Sigui $C_2=\{1,c\}$ el grup de dos elements i $S_3$ el grup de permutacions de tres elements. Considerem la inclusió $i\colon C_2\to S_3$, que envia $c$ a la transposició $(1\,2)$, i el signe $s\colon S_3\to C_2$, que envia les transposicions a $c$. Aleshores $s\comp i=\id_{C_2}$. Com que $C_2$ és commutatiu, $Z(C_2)=C_2$. En canvi, $Z(S_3)=\{1\}$: les tres transposicions no commuten entre si (per exemple, $(1\,2)(1\,3)=(1\,3\,2)$, mentre que $(1\,3)(1\,2)=(1\,2\,3)$), i les dues permutacions cícliques $(1\,2\,3)$ i $(1\,3\,2)$ no commuten amb $(1\,2)$, perquè $(1\,2)(1\,2\,3)(1\,2)=(1\,3\,2)$. Si $Z$ fos un functor,
\[
\id_{C_2}=Z(\id_{C_2})=Z(s\comp i)=Z(s)\comp Z(i),
\]
amb $Z(i)\colon C_2\to\{1\}$ i $Z(s)\colon\{1\}\to C_2$. Però una composició que passa per un grup d'un element és constant, i no pot ser la identitat d'un grup de dos elements.
\end{prova}

L'argument és un patró general: si $s\comp i=\id$, aleshores $F(s)\comp F(i)=\id$, de manera que tot functor converteix les retraccions en retraccions. Una assignació que envia un objecte a un de \enquote{massa petit} per factoritzar-hi una identitat no pot ser functorial. En canvi, assignacions semblants que només fan servir propietats preservades pels homomorfismes sí que són functors (vegeu els exercicis).

\paragraph{Fallada de les equacions.} Considerem el monoide de les funcions derivables $\mathbb R\to\mathbb R$ amb la composició, i l'assignació $f\mapsto f'$ cap al mateix monoide. Els tipus són correctes, però les equacions fallen. La identitat va a la funció constant $1$, que no és la identitat. I, amb $f(x)=g(x)=2x$, tenim $(g\comp f)'=4$ (constant), mentre que $g'\comp f'=2$ (constant). La regla de la cadena explica què cal canviar:

\begin{exemple}[La regla de la cadena és una llei functorial]\label{ex:regla-cadena}
Sigui $\cat{Der}$ la categoria que té per objectes els nombres reals i per morfismes $a\to b$ les funcions derivables $f\colon\mathbb R\to\mathbb R$ amb $f(a)=b$, amb la composició de funcions. Sigui $\cat{B}(\mathbb R,\cdot)$ el monoide multiplicatiu dels reals vist com a categoria d'un objecte. L'assignació
\[
D\colon\cat{Der}\to\cat{B}(\mathbb R,\cdot),\qquad (f\colon a\to b)\ \longmapsto\ f'(a)
\]
és un functor. Les identitats van a $1$, perquè $\id'(a)=1$, i la regla de la cadena, $(g\comp f)'(a)=g'(f(a))\cdot f'(a)$, diu exactament que $D(g\comp f)=D(g)\comp D(f)$, ja que $f(a)$ és el domini de $g$. El que fallava era el lloc on es mira la derivada: el functor no és $f\mapsto f'$, sinó $f\mapsto f'$ \emph{al punt de partida}.
\end{exemple}

\paragraph{Functors llevat d'isomorfisme (informatiu).} Hi ha assignacions que són functorials només llevat d'isomorfisme, exactament com els camins eren associatius llevat d'homotopia. L'exemple més important del llibre és la reindexació de l'observació~\ref{obs:reindexacio}. Per a cada $f\colon A\to B$, el pullback dona $f^{*}\colon\cat{C}/B\to\cat{C}/A$, però, un cop triats els pullbacks, $(g\comp f)^{*}$ i $f^{*}\comp g^{*}$ no són iguals, sinó canònicament isomorfs. A $\cat{Set}$, amb els pullbacks triats com a conjunts de parelles, $(g\comp f)^{*}$ envia $p\colon E\to C$ al conjunt de les parelles $(a,e)$ amb $g(f(a))=p(e)$, mentre que $f^{*}g^{*}$ l'envia al conjunt de les parelles $\bigl(a,(b,e)\bigr)$ amb $f(a)=b$ i $g(b)=p(e)$: són conjunts diferents, en bijecció canònica. Una tal assignació s'anomena \emph{pseudofunctor}. Al capítol~\ref{cap:subobjectes} el problema desapareix perquè $\Sub$ treballa amb classes de monomorfismes, i l'isomorfisme canònic queda absorbit (proposició~\ref{prop:functor-sub}): és el primer remei de la secció anterior, el quocient.

% =====================================================================
\section{Què no és una transformació natural}
\label{sec:no-natural}
% =====================================================================

El llibre ja dona els contraexemples essencials, i aquí només els ordenem. Una família de components $\eta_A\colon F(A)\to G(A)$ pot deixar de ser natural per tres motius diferents:
\begin{enumerate}
\item \emph{Les components s'han triat objecte per objecte}, sense mirar els morfismes. L'exemple~\ref{ex:no-natural} té totes les components bijectives i, tanmateix, un sol morfisme $\{0,1\}\to\{0,1\}$ trenca el quadrat.
\item \emph{La naturalitat només val per a una part dels morfismes.} L'isomorfisme entre un espai de dimensió finita i el seu dual depèn d'una base, i ni tan sols es pot plantejar com a transformació natural entre functors paral·lels, perquè la identitat és covariant i el dual contravariant. Només s'hi pot comparar restringint-se als isomorfismes, i fins i tot així no és natural (observació~\ref{obs:iso-no-natural} i exercici resolt~\ref{ex:dual-no-natural} de la Pràctica del capítol~\ref{cap:transformacions}). En canvi, $V\cong V^{**}$ és natural (exemple~\ref{ex:bidual}).
\item \emph{No hi ha cap component possible.} Amb el grup de dos elements actuant trivialment i per la transposició sobre $\{0,1\}$, no hi ha cap transformació natural entre els dos functors, ni tan sols una que no sigui un isomorfisme (observació~\ref{obs:iso-no-natural}).
\end{enumerate}
La lliçó comuna és la de la secció següent: tenir objectes isomorfs \emph{un per un} no és el mateix que tenir functors isomorfs.

% =====================================================================
\section{Igualtat, isomorfisme, isomorfisme natural i equivalència}
\label{sec:igualtat-equivalencia}
% =====================================================================

\subsection{La jerarquia}

Per a cada mena d'entitat hi ha una noció \emph{estricta} de \enquote{ser el mateix}, la igualtat, i una noció \emph{estructural}, que és la que respecta la teoria:
\[
\renewcommand{\arraystretch}{1.25}
\begin{array}{lll}
\text{Què comparem} & \text{Noció estricta} & \text{Noció estructural}\\\hline
\text{dos objectes d'una categoria }\cat{C} & A=B & A\cong B\ \text{(isomorfisme a }\cat{C})\\
\text{dos functors }F,G\colon\cat{C}\to\cat{D} & F=G & F\cong G\ \text{(isomorfisme natural)}\\
\text{dues categories} & \cat{C}=\cat{D} & \cat{C}\simeq\cat{D}\ \text{(equivalència)}
\end{array}
\]
Entre les categories hi ha, a més, un nivell intermedi: l'isomorfisme de categories $\cat{C}\cong\cat{D}$ (definició~\ref{def:iso-categories}). Les tres files tenen la mateixa forma. Un isomorfisme natural és un isomorfisme a la categoria de functors $\cat{D}^{\cat{C}}$, i per tant la segona fila és un cas de la primera. L'equivalència és el que s'obté quan, a la definició d'isomorfisme de categories, se substitueixen les igualtats $G\comp F=\id$ i $F\comp G=\id$ per isomorfismes naturals $G\comp F\cong\id$ i $F\comp G\cong\id$, perquè la manera correcta de comparar functors és l'isomorfisme natural. Cada nivell fa servir la noció estructural del nivell anterior.

Totes les implicacions possibles entre aquestes nocions van en un sol sentit:
\[
\begin{gathered}
A=B\ \Rightarrow\ A\cong B,\\
F=G\ \Rightarrow\ F\cong G\ \Rightarrow\ F(A)\cong G(A)\ \text{per a tot }A,\\
\cat{C}=\cat{D}\ \Rightarrow\ \cat{C}\cong\cat{D}\ \Rightarrow\ \cat{C}\simeq\cat{D},
\end{gathered}
\]
i cap no és reversible. Vegem-ne un contraexemple per a cada pas.

\subsection{Un contraexemple per a cada pas}

\begin{itemize}
\item \emph{Isomorfs però no iguals.} A $\cat{Set}$, $\{0,1\}\cong\{a,b\}$. A $\cat{Prov}_T$, $p\land q\cong q\land p$, que són fórmules diferents (exemple~\ref{ex:iso-equivalencia-logica}).
\item \emph{Naturalment isomorfs però no iguals.} A la categoria dels espais vectorials de dimensió finita, la identitat i el bidual són naturalment isomorfs (exemple~\ref{ex:bidual}), però no són iguals: $V^{**}$ no és $V$, sinó un espai format per funcionals sobre $V^{*}$.
\item \emph{Isomorfs objecte per objecte, però no naturalment isomorfs.} Els dos functors del grup de dos elements en $\cat{Set}$ de l'observació~\ref{obs:iso-no-natural} envien l'únic objecte al mateix conjunt $\{0,1\}$, però no són naturalment isomorfs.
\item \emph{Isomorfes però no iguals.} Una categoria i una còpia seva amb els objectes reanomenats.
\item \emph{Equivalents però no isomorfes.} $\cat{Mat}_{\mathbb K}\simeq\cat{FinVect}_{\mathbb K}$ (exemple~\ref{ex:mat-finvect}), o la categoria $\cat{1}$ i la categoria amb dos objectes isomorfs i un sol morfisme entre cada parell d'objectes.
\end{itemize}

\subsection{La invariància estructural}

El principi que governa tot el llibre és aquest: \emph{una propietat formulada només amb l'estructura ---composició, identitats, propietats universals---, sense afirmar mai que dos objectes són iguals, es transporta per la noció estructural corresponent.} Per als objectes ho diu l'observació~\ref{obs:invariancia}. Per a les categories, la peça clau és que les equivalències conserven les propietats universals.

\begin{teorema}\label{teo:equivalencia-limits}
Sigui $F\colon\cat{C}\to\cat{D}$ una equivalència. Si $(L,\lambda)$ és un límit d'un diagrama $D\colon\cat{I}\to\cat{C}$, aleshores $(F(L),F\lambda)$ és un límit de $F\comp D$. Dualment, $F$ preserva els colímits.
\end{teorema}

\begin{prova}
Pel teorema~\ref{teo:caract-equivalencia}, $F$ és ple, fidel i essencialment exhaustiu. La família $F\lambda=(F(\lambda_i))_i$ és un con sobre $F\comp D$, perquè $F$ preserva la composició. Sigui $(X,\alpha)$ un con sobre $F\comp D$. Com que $F$ és essencialment exhaustiu, hi ha un objecte $C$ de $\cat{C}$ i un isomorfisme $e\colon F(C)\to X$. Com que $F$ és ple i fidel, per a cada $i$ hi ha un únic $\beta_i\colon C\to D_i$ amb $F(\beta_i)=\alpha_i\comp e$. Aquests $\beta_i$ formen un con: per a $u\colon i\to j$,
\[
F(D_u\comp\beta_i)=F(D_u)\comp\alpha_i\comp e=\alpha_j\comp e=F(\beta_j),
\]
i, com que $F$ és fidel, $D_u\comp\beta_i=\beta_j$. Per la propietat del límit, hi ha un únic $h\colon C\to L$ amb $\lambda_i\comp h=\beta_i$. Posem $k=F(h)\comp e^{-1}\colon X\to F(L)$. Aleshores $F(\lambda_i)\comp k=F(\beta_i)\comp e^{-1}=\alpha_i$: $k$ factoritza el con.

Per a la unicitat, sigui $k'\colon X\to F(L)$ amb $F(\lambda_i)\comp k'=\alpha_i$ per a tot $i$. Com que $F$ és ple, $k'\comp e=F(h')$ per a algun $h'\colon C\to L$, i aleshores $F(\lambda_i\comp h')=\alpha_i\comp e=F(\beta_i)$. Com que $F$ és fidel, $\lambda_i\comp h'=\beta_i$, i per la unicitat del límit, $h'=h$. Per tant, $k'=F(h)\comp e^{-1}=k$. Els colímits es tracten igual, invertint totes les fletxes.
\end{prova}

D'aquest teorema, de l'exercici resolt~\ref{pr-equivalencia-preserva} de la Pràctica del capítol~\ref{cap:transformacions} (que una equivalència preserva i reflecteix isomorfismes, monomorfismes i epimorfismes) i del fet que el quasiinvers també és una equivalència, se segueix que les propietats de la columna esquerra de la taula següent les té $\cat{C}$ si i només si les té qualsevol categoria equivalent a $\cat{C}$. Per a les dues últimes (cartesiana tancada, topos), el mateix argument s'aplica a les propietats universals de l'exponencial i del classificador; no el reproduïm.
\begin{center}
\small
\renewcommand{\arraystretch}{1.2}
\begin{tabularx}{\linewidth}{@{}>{\raggedright\arraybackslash}X>{\raggedright\arraybackslash}X@{}}
\toprule
\textbf{Invariants per equivalència} & \textbf{No invariants per equivalència}\\
\midrule
tenir objecte terminal o inicial & el nombre d'objectes\\
tenir productes, equalitzadors, límits finits & tenir un sol objecte\\
ser un preordre (com a molt un morfisme entre dos objectes) & ser un ordre parcial\\
ser un grupoide (tot morfisme és invertible) & ser esquelètica\\
que tot morfisme mono i epi sigui un isomorfisme & que dos objectes concrets siguin iguals\\
ser cartesiana tancada; ser un topos & ser petita (\enquote{essencialment petita} sí que és invariant)\\
\bottomrule
\end{tabularx}
\end{center}
Les propietats de la columna de la dreta fan servir, d'una manera o altra, la igualtat d'objectes, i per això no són categòriques. Per exemple, $\cat{1}$ és un ordre parcial, i la categoria amb dos objectes isomorfs, que li és equivalent, no ho és. Quan una definició depèn de la igualtat d'objectes, convé sospitar-ne.

\subsection{Equivalència i esquelets}

El resultat següent fa veure que l'equivalència és exactament l'isomorfisme \enquote{llevat de les còpies isomorfes}.

\begin{lema}\label{lema:esqueletiques}
Una equivalència entre dues categories esquelètiques és un isomorfisme de categories.
\end{lema}

\begin{prova}
Sigui $F\colon\cat{C}\to\cat{D}$ una equivalència entre categories esquelètiques; és ple, fidel i essencialment exhaustiu. \emph{És injectiu sobre objectes}: si $F(A)=F(B)$, la identitat de $F(A)$ és $F(u)$ per a algun $u\colon A\to B$, perquè $F$ és ple, i $u$ és un isomorfisme, perquè un functor ple i fidel reflecteix els isomorfismes (proposició~\ref{prop:plenfidel-reflecteix-iso}). Com que $\cat{C}$ és esquelètica, $A=B$. \emph{És exhaustiu sobre objectes}: per a cada $D$ hi ha $C$ amb $D\cong F(C)$, i com que $\cat{D}$ és esquelètica, $D=F(C)$. A més, $F$ és bijectiu sobre cada homconjunt, perquè és ple i fidel. Per tant, $F$ és bijectiu sobre objectes i sobre morfismes, i l'assignació inversa és un functor (preserva identitats i composicions perquè $F$ les preserva). És, doncs, un isomorfisme de categories.
\end{prova}

\begin{corollari}\label{cor:equivalencia-esquelets}
Dues categories són equivalents si i només si tenen esquelets isomorfs.
\end{corollari}

\begin{prova}
En el marc d'elecció del llibre, tota categoria té esquelets (definició~\ref{def:esquelet}), i és equivalent a qualsevol d'ells: la inclusió de l'esquelet és plena i fidel per ser una subcategoria plena, i essencialment exhaustiva perquè l'esquelet conté un objecte de cada classe d'isomorfisme. Les equivalències es componen, perquè la composició de functors plens, fidels i essencialment exhaustius ho torna a ser. Si $\cat{C}\simeq\cat{D}$, els seus esquelets són equivalents, i pel lema~\ref{lema:esqueletiques} són isomorfs. Recíprocament, si els esquelets són isomorfs, $\cat{C}\simeq\mathrm{esquelet}(\cat{C})\cong\mathrm{esquelet}(\cat{D})\simeq\cat{D}$.
\end{prova}

Per als preordres, l'esquelet és l'ordre parcial que s'obté identificant els elements equivalents ($x\le y$ i $y\le x$): dos preordres són equivalents si i només si aquests ordres parcials són isomorfs.

\subsection{Com es demostra que dues categories no són equivalents}

Per demostrar que $\cat{C}\not\simeq\cat{D}$ no es poden revisar tots els functors possibles. La via és buscar una propietat invariant per equivalència que l'una tingui i l'altra no.

\begin{exemple}[$\cat{Set}\not\simeq\cat{Set}\op$]
Diem que un objecte inicial $0$ és \emph{estricte} si tot morfisme $X\to0$ és un isomorfisme. Tenir un objecte inicial estricte és invariant per equivalència. Si $F\colon\cat{C}\to\cat{D}$ és una equivalència, $F(0)$ és inicial (teorema~\ref{teo:equivalencia-limits}). Donat un morfisme $k\colon Y\to F(0)$, prenem un isomorfisme $e\colon F(C)\to Y$; com que $F$ és ple, $k\comp e=F(u)$ per a algun $u\colon C\to0$, que és un isomorfisme, i aleshores $F(u)$ i $k=F(u)\comp e^{-1}$ també ho són. A $\cat{Set}$, l'objecte inicial $\varnothing$ és estricte: només hi ha aplicacions $X\to\varnothing$ si $X=\varnothing$, i aleshores és la identitat. A $\cat{Set}\op$, l'objecte inicial és un conjunt d'un element $\{\ast\}$, i un morfisme $X\to\{\ast\}$ de $\cat{Set}\op$ és una aplicació $\{\ast\}\to X$ de $\cat{Set}$. Amb $X=\{0,1\}$ no és bijectiva, i per tant no és un isomorfisme. Així, $\cat{Set}$ no és equivalent a la seva oposada.
\end{exemple}

\begin{exemple}[$\cat{Grp}\not\simeq\cat{Ab}$]
La propietat \enquote{per a dos objectes qualssevol, el producte i el coproducte són isomorfs} és invariant per equivalència, perquè una equivalència preserva productes, coproductes i isomorfismes. 
\emph{A $\cat{Ab}$ es compleix.} Siguin $A$ i $B$ grups abelians, escrits additivament. El producte $A\times B$, amb les inclusions $a\mapsto(a,0)$ i $b\mapsto(0,b)$, és també un coproducte. Donats homomorfismes $f\colon A\to G$ i $g\colon B\to G$ cap a un grup abelià $G$, l'aplicació $(a,b)\mapsto f(a)+g(b)$ és un homomorfisme, precisament perquè $G$ és commutatiu, i restringida a les inclusions dona $f$ i $g$. És l'única, perquè $(a,b)=(a,0)+(0,b)$.

\emph{A $\cat{Grp}$ no es compleix.} Sigui $C_2=\{1,c\}$ i suposem que el coproducte $P=C_2+C_2$ és isomorf al producte $C_2\times C_2$, que és commutatiu. Els homomorfismes $C_2\to S_3$ que envien $c$ a $(1\,2)$ i a $(1\,3)$ indueixen un homomorfisme $P\to S_3$ la imatge del qual conté $(1\,2)$ i $(1\,3)$. Però la imatge d'un grup commutatiu és commutativa, i $(1\,2)$ i $(1\,3)$ no commuten. Per tant, $C_2+C_2\not\cong C_2\times C_2$.
\end{exemple}

\begin{exemple}[$\cat{FinSet}\not\simeq\cat{Set}$]
\enquote{Tota família numerable d'objectes té coproducte} és invariant per equivalència. $\cat{Set}$ la compleix i $\cat{FinSet}$ no (observació~\ref{obs:elemental-prefeixos}).
\end{exemple}

\subsection{Què té a veure tot això amb la lògica}

A la categoria $\cat{Prov}_T$ d'un sistema deductiu, els tres nivells tenen una lectura lògica precisa:
\begin{itemize}
\item la \emph{igualtat} de dues fórmules és la igualtat sintàctica: $p\land q$ i $q\land p$ són fórmules diferents;
\item l'\emph{isomorfisme} és la interdeduïbilitat, $p\vdash_T q$ i $q\vdash_T p$ (exemple~\ref{ex:iso-equivalencia-logica});
\item l'\emph{esquelet} de $\cat{Prov}_T$ és isomorf a l'ordre parcial de les classes de fórmules interdeduïbles, que en lògica s'anomena \textbf{àlgebra de Lindenbaum--Tarski}\index{àlgebra de Lindenbaum--Tarski} de $T$. Pel corol·lari~\ref{cor:equivalencia-esquelets}, $\cat{Prov}_T$ és equivalent, però en general no isomorf, a aquesta àlgebra.
\end{itemize}
En conseqüència, qualsevol propietat categòrica de $\cat{Prov}_T$ és una propietat de l'àlgebra de Lindenbaum--Tarski: tenir productes (la conjunció), ser cartesiana tancada (la implicació), ser una àlgebra de Heyting o de Boole. La lògica estudia les teories \emph{llevat d'interdeduïbilitat}, i la teoria de categories hi aporta el llenguatge que fa aquesta actitud sistemàtica: el que compta d'una teoria és la seva estructura, no la tria particular de fórmules.

\subsection{Errors típics}

\begin{itemize}
\item Escriure $A=B$ quan només es té $A\cong B$, encara que l'isomorfisme sigui canònic.
\item Deduir $F\cong G$ de $F(A)\cong G(A)$ per a cada $A$, sense comprovar la naturalitat.
\item Creure que una equivalència és un isomorfisme de categories, o que $G\comp F$ ha de ser literalment la identitat.
\item Demostrar que dues categories no són equivalents amb un argument que fa servir la igualtat d'objectes, com comptar-ne el nombre.
\item Definir una propietat d'una categoria que depèn d'objectes concrets, i després aplicar-la a una categoria equivalent.
\end{itemize}

% =====================================================================
\section{La lògica fora dels topos}
\label{sec:fora-topos}
% =====================================================================

Al llarg dels capítols~\ref{cap:cartesianes}--\ref{cap:pont}, la lògica s'ha llegit d'una manera determinada, que comença amb el tancament cartesià: una proposició sobre $E$ és un subobjecte de $E$, la conjunció és un pullback, la implicació és l'exponencial (l'adjunt per la dreta de $-\wedge U$), els quantificadors són adjunts de la reindexació, i en un topos tot això és la lògica interna de $\Omega$. Hi ha lògiques i models perfectament legítims que no encaixen en aquest esquema. Els presentem de manera sobretot informativa, però en cada cas diem amb precisió quina condició falla i, quan és fàcil, ho comprovem.

\subsection{Categories que no són cartesianes tancades}
\label{subsec:no-ccc}

Abans de mirar què hi ha entre les categories cartesianes tancades i els topos, convé veure que moltes categories molt naturals no són ni tan sols cartesianes tancades. Per demostrar-ho, n'hi ha prou de trobar una propietat que tota categoria cartesiana tancada tingui i que la categoria en qüestió no tingui. La més útil és la que dona el teorema de preservació de colímits pels adjunts per l'esquerra (capítol~\ref{cap:adjuncions}): en una categoria cartesiana tancada, $-\times B$ és adjunt per l'esquerra de $(-)^B$, i per tant preserva tots els colímits que existeixin. En particular, si $0$ és inicial, $0\times B$ també ho és.

Un \textbf{objecte zero}\index{objecte zero} és un objecte que és alhora inicial i terminal.

\begin{proposicio}\label{prop:zero-no-ccc}
Si una categoria cartesiana tancada té un objecte zero $0$, aleshores tots els seus objectes són isomorfs a $0$.
\end{proposicio}

\begin{prova}
Sigui $B$ un objecte. Com que $0$ és inicial i $-\times B$ preserva l'objecte inicial, $0\times B$ és inicial, i per tant $0\times B\cong0$. D'altra banda, com que $0$ és terminal, la projecció $\pi_2\colon0\times B\to B$ és un isomorfisme. En efecte, sigui $!_B\colon B\to0$ l'únic morfisme cap al terminal, i $s=\langle !_B,\id_B\rangle\colon B\to0\times B$. Aleshores $\pi_2\comp s=\id_B$, i
\[
s\comp\pi_2=\langle !_B\comp\pi_2,\ \pi_2\rangle=\langle\pi_1,\pi_2\rangle=\id_{0\times B},
\]
perquè $!_B\comp\pi_2$ i $\pi_1$ són dos morfismes $0\times B\to0$, que han de coincidir perquè $0$ és terminal. Per tant, $B\cong0\times B\cong0$.
\end{prova}

\begin{corollari}\label{cor:zero-no-ccc}
Les categories $\cat{Vect}_{\mathbb K}$, $\cat{Ab}$, $\cat{Grp}$, $\cat{Mon}$ i $\cat{Rel}$ no són cartesianes tancades.
\end{corollari}

\begin{prova}
Totes tenen un objecte zero: l'espai $\{0\}$; el grup, o el monoide, d'un sol element, perquè de l'element neutre surt un únic homomorfisme cap a qualsevol altre grup o monoide i cap a l'estructura trivial només n'hi ha un; i, a $\cat{Rel}$, el conjunt buit, perquè l'única relació $\varnothing\to X$ i l'única relació $X\to\varnothing$ són la relació buida. Totes tenen, però, objectes que no són isomorfs al zero: $\mathbb K$, el grup $\mathbb Z$, el monoide $(\mathbb N,+)$ o qualsevol conjunt no buit. Per la proposició~\ref{prop:zero-no-ccc}, cap no és cartesiana tancada.
\end{prova}

\paragraph{Preordres.} En un preordre amb ínfims finits, ser cartesià tancat vol dir tenir una implicació $b\Rightarrow c$ amb $a\wedge b\le c$ si i només si $a\le b\Rightarrow c$ (capítol~\ref{cap:cartesianes}). Si, a més, té suprems binaris, aleshores és distributiu, perquè $-\wedge b$ és adjunt per l'esquerra i preserva els suprems. Per tant, \emph{un reticle no distributiu no és cartesià tancat}. El reticle dels subespais de $\mathbb R^2$ que veurem tot seguit n'és un exemple.

\paragraph{Categories que ni tan sols són cartesianes.} Un monoide $M$ no trivial, vist com a categoria $\cat{B}M$, no té objecte terminal: l'únic objecte $\ast$ ho seria només si $\Hom(\ast,\ast)=M$ tingués un sol element. Un preordre amb dos elements incomparables que no tinguin cap fita inferior comuna no té el producte d'aquests dos elements. Cap d'aquestes categories no pot ser cartesiana tancada, perquè no té productes finits.

\paragraph{Els espais topològics (informatiu).} $\cat{Top}$ tampoc no és cartesiana tancada: hi ha quocients d'espais que el producte amb un espai fix no conserva. La demostració demana més topologia de la que pressuposa el llibre, i no la fem.

\subsection{Cartesianes tancades que no són topos}

La lògica de la implicació necessita només el tancament cartesià, i això és molt menys que un topos.

\begin{exemple}[$\cat{Cat}$ no és un topos]
La categoria $\cat{Cat}$ és cartesiana tancada (exemple~\ref{ex:cat-ccc}), però no té classificador de subobjectes. Sigui $\cat{2}$ la categoria $0\to1$, i $\cat{I}$ la categoria amb dos objectes i, a més de les identitats, dos morfismes $u\colon0\to1$ i $u^{-1}\colon1\to0$ inversos l'un de l'altre. La inclusió $i\colon\cat{2}\to\cat{I}$ és:
\begin{itemize}
\item \emph{mono}: és injectiva sobre objectes i morfismes, de manera que si $i\comp H=i\comp K$, aleshores $H=K$;
\item \emph{epi}: si $H,K\colon\cat{I}\to\cat{C}$ coincideixen sobre $u$, aleshores també coincideixen sobre $u^{-1}$, perquè un functor envia l'invers de $u$ a l'invers de la imatge de $u$, i l'invers és únic;
\item \emph{no és un isomorfisme}: $u^{-1}$ no és a la imatge.
\end{itemize}
En una categoria amb classificador, tot morfisme mono i epi és un isomorfisme (exercici resolt~\ref{pr-mono-equalitzador} de la Pràctica del capítol~\ref{cap:topos}). Per tant, $\cat{Cat}$ no en té cap.
\end{exemple}

L'ordre parcial $\mathbb Z\cup\{+\infty\}$ de l'observació~\ref{obs:ccc-no-heyting} és un altre cas: té conjunció i implicació, però no té falsedat.

\subsection{Lògiques on falla la distributivitat}

En una àlgebra de Heyting, $a\wedge(b\vee c)=(a\wedge b)\vee(a\wedge c)$, perquè $a\wedge-$ és adjunt per l'esquerra i preserva els suprems (capítol~\ref{cap:cartesianes}). Com que els subobjectes d'un objecte d'un topos formen una àlgebra de Heyting (nota~\ref{nota:sub-heyting}), cap reticle no distributiu pot ser-ne un.

\begin{exemple}[La lògica quàntica]
Siguin els subespais vectorials de $\mathbb R^2$, ordenats per inclusió: $\{0\}$, les rectes que passen per l'origen i $\mathbb R^2$. L'ínfim és la intersecció i el suprem és la suma. Si $a$, $b$ i $c$ són tres rectes diferents,
\[
a\wedge(b\vee c)=a\wedge\mathbb R^2=a,
\qquad
(a\wedge b)\vee(a\wedge c)=\{0\}\vee\{0\}=\{0\}.
\]
La distributivitat falla. En mecànica quàntica, les propietats observables d'un sistema es modelen amb reticles d'aquest tipus (de subespais tancats), i la \enquote{lògica quàntica} de Birkhoff i von Neumann és la lògica d'aquests reticles no distributius (informatiu). Amb aquestes operacions, aquest reticle no pot ser el reticle de subobjectes de cap objecte d'un topos.
\end{exemple}

\subsection{Lògiques amb una altra conjunció}

\begin{exemple}[La lògica de Łukasiewicz]
Els valors de veritat són els nombres de $[0,1]$, amb la conjunció forta i la implicació
\[
x\otimes y=\max(0,\,x+y-1),\qquad y\to z=\min(1,\,1-y+z).
\]
La implicació és adjunta d'aquesta conjunció, i no de l'ínfim. En efecte, si $z\ge0$, $\max(0,x+y-1)\le z$ si i només si $x\le1-y+z$, és a dir, $x\le y\to z$. En canvi, amb l'ínfim falla: per a $y=0.5$, $z=0.4$ i $x=0.6$, tenim $\min(x,y)=0.5>0.4$, però $x=0.6\le0.9=y\to z$. A més, $\otimes$ no és idempotent: $0.5\otimes0.5=0\neq0.5$. Lògicament, això vol dir que de $A$ no se'n dedueix $A\otimes A$: una hipòtesi no es pot fer servir dues vegades. Categòricament, $([0,1],\otimes,\to)$ és un preordre \emph{monoidal tancat} que no és cartesià tancat respecte de $\otimes$. Aquesta idea ---hipòtesis que són recursos i no es poden duplicar--- és la de la \emph{lògica lineal} de Girard; els seus fragments més simples es modelen en categories monoidals tancades, i la resta demana estructura addicional (informatiu).
\end{exemple}

\subsection{Lògiques on la contradicció no és el fals}

\begin{exemple}[Els tancats d'un espai]
A $\mathbb R$, els subconjunts tancats, ordenats per inclusió, formen un reticle distributiu on el \enquote{complement} natural d'un tancat $A$ és l'adherència del seu complementari, $\neg A=\overline{\mathbb R\setminus A}$. Per a $A=[0,1]$,
\[
A\wedge\neg A=[0,1]\cap\bigl((-\infty,0]\cup[1,\infty)\bigr)=\{0,1\}\neq\varnothing:
\]
una proposició i la seva negació són certes alhora a la frontera. Aquesta estructura és la duala d'una àlgebra de Heyting (una àlgebra de \emph{co-Heyting}), i serveix de model a lògiques \emph{paraconsistents}, on una contradicció no implica qualsevol cosa. Amb aquestes operacions, no pot ser el reticle de subobjectes de cap objecte d'un topos, perquè hi falla $A\wedge\neg A=\bot$.
\end{exemple}

\subsection{Lògiques amb modalitats}

Una \emph{lògica modal} afegeix operadors com \enquote{necessàriament} ($\Box$) i \enquote{possiblement} ($\Diamond$). En un model de Kripke, hi ha un conjunt de mons $W$ i una relació d'accessibilitat $R\subseteq W\times W$, i per a $A\subseteq W$ es defineix
\[
\Box A=\{x\mid \text{per a tot } y,\ xRy\Rightarrow y\in A\},
\qquad
\Diamond A=\{x\mid \text{hi ha } y \text{ amb } xRy \text{ i } y\in A\}.
\]
Aquestes modalitats són quantificadors \emph{al llarg de la relació} $R$, com els $\forall_f$ i $\exists_f$ del capítol~\ref{cap:subobjectes} ho eren al llarg d'una aplicació. Per veure-ho, sigui $\exists_R(B)=\{y\mid \text{hi ha } x\in B \text{ amb } xRy\}$, el conjunt dels mons accessibles des de $B$. Aleshores, per a $A,B\subseteq W$,
\[
\exists_R(B)\subseteq A\iff B\subseteq\Box A,
\]
és a dir, $\exists_R\dashv\Box$ entre els ordres parcials $(\mathcal P(W),\subseteq)$. En efecte, totes dues condicions diuen el mateix: per a tot $x\in B$ i tot $y$ amb $xRy$, es té $y\in A$. La modalitat $\Diamond$ és el mateix tipus de quantificador, però al llarg de la relació inversa: $\Diamond A=\{x\mid\text{hi ha } y\in A \text{ amb } yR^{-1}x\}$. Les modalitats no són, però, connectius de la lògica interna d'un topos, sinó estructura afegida (informatiu).

\subsection{Resum}

\begin{center}
\small
\renewcommand{\arraystretch}{1.2}
\begin{tabularx}{\linewidth}{@{}>{\raggedright\arraybackslash}p{0.24\linewidth}>{\raggedright\arraybackslash}X>{\raggedright\arraybackslash}X@{}}
\toprule
\textbf{Model} & \textbf{Estructura} & \textbf{Què falla respecte d'un topos}\\
\midrule
$\cat{Vect}_{\mathbb K}$, $\cat{Ab}$, $\cat{Grp}$, $\cat{Rel}$ & amb objecte zero & no són cartesianes tancades\\
$\cat{Cat}$ & cartesiana tancada & no hi ha classificador\\
$\mathbb Z\cup\{+\infty\}$ & preordre cartesià tancat & no hi ha fals\\
subespais de $\mathbb R^2$ & reticle & la distributivitat, i per tant el tancament cartesià\\
Łukasiewicz & preordre monoidal tancat & la implicació no és adjunta de $\wedge$\\
tancats de $\mathbb R$ & co-Heyting & $A\wedge\neg A\neq\bot$\\
models de Kripke & àlgebra de Boole amb modalitats & estructura afegida\\
\bottomrule
\end{tabularx}
\end{center}
Cada fila marca un camí que la lògica categòrica pot seguir més enllà dels topos: les estructures monoidals tancades per a les lògiques de recursos, els reticles no distributius per a la lògica quàntica, els operadors modals com a adjunts al llarg de relacions. El llibre s'ha centrat en l'esquema dels topos perquè és el que unifica més bé la lògica intuïcionista i la teoria de conjunts. Aquestes alternatives mostren que és una tria, no l'únic camí possible.

% =====================================================================
\section{Exercicis}
% =====================================================================

\begin{exercici}
Demostra que $c_y\ast\gamma=\gamma$ si i només si $\gamma$ és constant. \hfill{\normalfont$\star$}
\begin{indicacio}
La igualtat a $[\frac12,1]$ força $\gamma=y$ en aquest interval, i a $[0,\frac12]$ diu $\gamma(t)=\gamma(2t)$. Itera-ho i fes servir la continuïtat a $0$.
\end{indicacio}
\end{exercici}

\begin{exercici}
Comprova que, a la categoria dels camins de Moore, un camí $(r,\gamma)$ amb $r>0$ no té cap invers: no hi ha cap $(s,\delta)$ amb $(s,\delta)\comp(r,\gamma)=(0,c_x)$. \hfill{\normalfont$\star$}
\begin{indicacio}
Mira la durada de la composició.
\end{indicacio}
\end{exercici}

\begin{exercici}
Troba un magma de dos elements que no sigui associatiu, i un altre que sigui associatiu però no tingui neutre. \hfill{\normalfont$\star$}
\begin{indicacio}
Hi ha només setze operacions possibles en un conjunt de dos elements. Per al segon, prova amb una operació que no depengui del segon argument.
\end{indicacio}
\end{exercici}

\begin{exercici}
Considera els nombres reals com a objectes, i com a morfismes $x\to y$ els reals $d\ge|y-x|$, amb la composició $d'\comp d=d+d'$. És una categoria? Què canvia si s'exigeix $d=|y-x|$? \hfill{\normalfont$\star\star$}
\begin{indicacio}
Per al primer cas, fes servir la desigualtat triangular. Per al segon, mira què passa amb $x<y>z$.
\end{indicacio}
\end{exercici}

\begin{exercici}
Sigui $T(G)=\{g\in G\mid g^2=1\}$. Comprova que $G\mapsto T(G)$, amb $f\mapsto f|_{T(G)}$, és un functor $\cat{Grp}\to\cat{Set}$, i explica per què l'argument de la proposició~\ref{prop:centre-no-functor} no s'hi aplica. \hfill{\normalfont$\star$}
\begin{indicacio}
Mira primer els tipus: si $g^2=1$, què es pot dir de $f(g)^2$? Per a la segona part, calcula $T(S_3)$.
\end{indicacio}
\end{exercici}

\begin{exercici}
Demostra que un functor bijectiu sobre objectes i sobre morfismes és un isomorfisme de categories. \hfill{\normalfont$\star$}
\begin{indicacio}
Defineix l'invers sobre objectes i morfismes, i comprova que preserva la composició aplicant-hi el functor original.
\end{indicacio}
\end{exercici}

\begin{exercici}
Demostra que \enquote{ser un grupoide} és invariant per equivalència, i que \enquote{ser un ordre parcial} no ho és. \hfill{\normalfont$\star$}
\begin{indicacio}
Per a la primera part, fes servir que una equivalència és plena i fidel i que reflecteix els isomorfismes.
\end{indicacio}
\end{exercici}

\begin{exercici}
La categoria $\cat{Rel}$ és isomorfa a la seva oposada, enviant cada relació $R$ a la relació inversa. Fes-ho servir, junt amb l'exemple $\cat{Set}\not\simeq\cat{Set}\op$, per demostrar que $\cat{Set}\not\simeq\cat{Rel}$. \hfill{\normalfont$\star\star$}
\begin{indicacio}
Si $F\colon\cat{C}\to\cat{D}$ és una equivalència, $F\op\colon\cat{C}\op\to\cat{D}\op$ també ho és.
\end{indicacio}
\end{exercici}

\begin{exercici}
Sigui $\cat{Set}_\ast$ la categoria dels conjunts amb un punt distingit, amb les aplicacions que preserven el punt distingit. Comprova que té un objecte zero i dedueix que no és cartesiana tancada. \hfill{\normalfont$\star$}
\begin{indicacio}
Pensa en el conjunt d'un sol element. Per a la segona part, aplica la proposició~\ref{prop:zero-no-ccc}.
\end{indicacio}
\end{exercici}

\begin{exercici}
Comprova la fallada de la distributivitat al reticle dels subespais de $\mathbb R^2$ amb les rectes $y=0$, $x=0$ i $y=x$. Té aquest reticle un \enquote{complement} per a cada recta? \hfill{\normalfont$\star$}
\begin{indicacio}
Un complement de $a$ és un $b$ amb $a\wedge b=\{0\}$ i $a\vee b=\mathbb R^2$. És únic?
\end{indicacio}
\end{exercici}

\begin{exercici}
A $\mathbb R$, calcula $A\wedge\neg A$ per al tancat $A=[0,1]\cup\{2\}$, amb $\neg A=\overline{\mathbb R\setminus A}$. \hfill{\normalfont$\star$}
\begin{indicacio}
El resultat és la frontera de $A$.
\end{indicacio}
\end{exercici}

\fi

\backmatter
\chapter*{Bibliografia}
\addcontentsline{toc}{chapter}{Bibliografia}
\markboth{Bibliografia}{Bibliografia}

Les referències següents cobreixen tres funcions: introduccions al mateix nivell que aquest text, obres de consulta i fonts per a la transició cap a la lògica categòrica, i les fonts de la terminologia matemàtica catalana emprada.

\begin{thebibliography}{99}

\bibitem{eilenbergmaclane}
S.~Eilenberg i S.~Mac~Lane,
\emph{General theory of natural equivalences},
Transactions of the American Mathematical Society~\textbf{58} (1945), 231--294.
\newblock (Article fundacional de la teoria de categories, on es defineixen per primera vegada les nocions de categoria, functor i transformació natural. \url{https://doi.org/10.1090/S0002-9947-1945-0013131-6}.)

\bibitem{awodey}
S.~Awodey,
\emph{Category Theory},
2a edició, Oxford Logic Guides~52, Oxford University Press, 2010.
\newblock (Text introductori de referència; guia principal per a l'ordre i el nivell d'aquest llibre, i per als capítols sobre categories cartesianes tancades i lògica. Les remissions del llibre es refereixen a aquesta edició.)

\bibitem{maclane}
S.~Mac~Lane,
\emph{Categories for the Working Mathematician},
2a edició, Graduate Texts in Mathematics~5, Springer, 1998. \url{https://doi.org/10.1007/978-1-4757-4721-8}.
\newblock (Obra de consulta clàssica sobre teoria de categories. Les remissions del llibre es refereixen a aquesta edició.)

\bibitem{leinster}
T.~Leinster,
\emph{Basic Category Theory},
Cambridge Studies in Advanced Mathematics~143, Cambridge University Press, 2014. Versió en accés obert: \url{https://arxiv.org/abs/1612.09375}.
\newblock (Introducció alternativa, breu i moderna. La versió oberta es va publicar a arXiv el desembre de 2016 per acord amb Cambridge University Press; incorpora les correccions de la llista d'errates de l'edició impresa i en conserva la numeració de capítols i seccions, que és la que segueixen les remissions d'aquest llibre. Consulta: octubre de 2026.)

\bibitem{riehl}
E.~Riehl,
\emph{Category Theory in Context},
Aurora: Dover Modern Math Originals, Dover, 2016.
\newblock (Consulta complementària, amb molts exemples. Les remissions d'aquest llibre, per capítols i seccions, es refereixen a l'edició de 2016. La pàgina de l'autora ofereix una versió en accés obert, \url{https://emilyriehl.github.io/files/context.pdf}, que s'actualitza i pot no coincidir amb l'edició de 2016 en la paginació. La pàgina de l'autora distingeix actualment la primera edició de la versió en preparació de la segona edició; les remissions d'aquest llibre es refereixen a la primera edició de 2016. Consulta: octubre de 2026.)

\bibitem{barrwells}
M.~Barr i C.~Wells,
\emph{Category Theory for Computing Science},
3a edició, Les Publications CRM, Montréal, 1999. Reimprès a \emph{Reprints in Theory and Applications of Categories}~22 (2012). \url{https://www.tac.mta.ca/tac/reprints/articles/22/tr22abs.html}.
\newblock (Introducció orientada a la informàtica i la lògica, amb sistemes deductius i llenguatges funcionals com a exemples de categories. El portal de TAC n'ofereix també versions corregides posteriors a la reimpressió de 2012; les remissions d'aquest llibre es refereixen a la reimpressió de 2012; per a pàgines exactes cal indicar la versió consultada. Consulta: octubre de 2026.)

\bibitem{perrone}
P.~Perrone,
\emph{Starting Category Theory},
World Scientific, Singapur, 2024. \url{https://doi.org/10.1142/13670}.
\newblock (Introducció detallada i accessible, amb molts exemples de matemàtica bàsica.)

\bibitem{perrone-notes}
P.~Perrone,
\emph{Notes on Category Theory with examples from basic mathematics},
notes en accés obert, arXiv:1912.10642, versió~7 (abril de 2024; text no revisat, darrera actualització de febrer de 2021), \url{https://arxiv.org/abs/1912.10642}.
\newblock (Material previ al llibre \cite{perrone}, que el va ampliar i reorganitzar; les remissions a aquestes notes no es corresponen necessàriament amb els capítols del llibre. Consulta: octubre de 2026.)

\bibitem{lawvereschanuel}
F.~W.~Lawvere i S.~H.~Schanuel,
\emph{Conceptual Mathematics: A First Introduction to Categories},
2a edició, Cambridge University Press, 2009.
\newblock (Introducció molt elemental, amb conjunts, grafs i sistemes dinàmics com a exemples recurrents. Les remissions es fan per parts, articles i sessions, i s'han comprovat amb la traducció castellana de la segona edició, \emph{Matemáticas conceptuales. Una primera introducción a categorías}, de F.~Marmolejo, que en conserva la numeració.)

\bibitem{streicher}
T.~Streicher,
\emph{Introduction to Category Theory and Categorical Logic},
notes de curs, Technische Universität Darmstadt, 2003--2004.
\newblock (Introducció a la teoria de categories orientada a la lògica categòrica; referència per als capítols de categories cartesianes tancades i topos i per al pont cap al volum de lògica. Les remissions es refereixen a la versió en línia de les notes indicada a la referència.)

\bibitem{borceux}
F.~Borceux,
\emph{Handbook of Categorical Algebra},
3 volums, Encyclopedia of Mathematics and its Applications, Cambridge University Press, 1994.
\newblock (Referència exhaustiva per a límits, adjuncions i exactitud.)

\bibitem{maclanemoerdijk}
S.~Mac~Lane i I.~Moerdijk,
\emph{Sheaves in Geometry and Logic: A First Introduction to Topos Theory},
Universitext, Springer, 1992.
\newblock (Pont cap a la teoria de topos i la lògica geomètrica. Es cita l'edició estàndard impresa de 1992.)

\bibitem{lambekscott}
J.~Lambek i P.~J.~Scott,
\emph{Introduction to Higher-Order Categorical Logic},
Cambridge Studies in Advanced Mathematics~7, Cambridge University Press, 1986.
\newblock (Correspondència entre lògica d'ordre superior, càlcul lambda i categories cartesianes tancades.)

\bibitem{makkaireyes}
M.~Makkai i G.~E.~Reyes,
\emph{First Order Categorical Logic},
Lecture Notes in Mathematics~611, Springer, 1977.
\newblock (Tractament clàssic de la lògica de primer ordre en categories regulars, coherents i de Heyting, i dels topos classificadors; referència per al volum de lògica categòrica.)

\bibitem{pitts}
A.~M.~Pitts,
\emph{Categorical Logic},
dins S.~Abramsky, D.~M.~Gabbay i T.~S.~E.~Maibaum (eds.), \emph{Handbook of Logic in Computer Science}, vol.~5, Oxford University Press, 2000, pàg.~39--128.
\newblock (Introducció sistemàtica a la semàntica categòrica de la lògica, des de les teories equacionals fins a la lògica d'ordre superior; referència per al volum de lògica categòrica.)

\bibitem{jacobs}
B.~Jacobs,
\emph{Categorical Logic and Type Theory},
Studies in Logic and the Foundations of Mathematics~141, Elsevier, 1999.
\newblock (Fibracions, hiperdoctrines i teoria de tipus des del punt de vista categòric.)

\bibitem{johnstone}
P.~T.~Johnstone,
\emph{Sketches of an Elephant: A Topos Theory Compendium},
Oxford Logic Guides 43--44, Oxford University Press, 2002.
\newblock (Obra avançada de consulta sobre teoria de topos.)

\bibitem{lawvere-ccfm}
F.~W.~Lawvere,
\emph{The Category of Categories as a Foundation for Mathematics},
en \emph{Proceedings of the Conference on Categorical Algebra (La Jolla, 1965)}, Springer, 1966, pàg.~1--20.
\newblock (Teoria elemental de categories en primer ordre; precedent de la presentació relacional de l'apèndix~\ref{cap:axiomatica-relacional}.)

\bibitem{lawvere-etcs}
F.~W.~Lawvere,
\emph{An Elementary Theory of the Category of Sets},
Proceedings of the National Academy of Sciences of the USA~\textbf{52} (1964), 1506--1511. Versió llarga, amb comentaris de C.~McLarty i de l'autor, a \emph{Reprints in Theory and Applications of Categories}~11 (2005): \url{http://www.tac.mta.ca/tac/reprints/articles/11/tr11abs.html}.
\newblock (Axiomàtica de la teoria de conjunts formulada en termes de la categoria de conjunts: la teoria de categories com a fonament.)

\bibitem{leinster-rethinking}
T.~Leinster,
\emph{Rethinking Set Theory},
American Mathematical Monthly~\textbf{121} (2014), núm.~5, 403--415. Versió en accés obert: \url{https://arxiv.org/abs/1212.6543}.
\newblock (Presentació informal i breu de la teoria de conjunts estructural de Lawvere.)

\bibitem{shulman-mida}
M.~A.~Shulman,
\emph{Set Theory for Category Theory},
prepublicació, 2008. \url{https://arxiv.org/abs/0810.1279}.
\newblock (Comparació de les diferents maneres de tractar la mida en teoria de categories: classes, universos i altres alternatives.)

\bibitem{nlab-classificador}
nLab,
\emph{classifying topos},
\url{https://ncatlab.org/nlab/show/classifying+topos} (consultat el 29 de setembre de 2026).
\newblock (Wiki de referència en teoria de categories; entrada sobre el topos classificador d'una teoria geomètrica.)

\bibitem{nlab-locale}
nLab,
\emph{locale},
\url{https://ncatlab.org/nlab/show/locale} (consultat el 29 de setembre de 2026).
\newblock (Hi apareix l'exemple del locale de les aplicacions exhaustives $\mathbb N\to\mathbb R$, no trivial i sense punts.)

\bibitem{blechschmidt}
I.~Blechschmidt,
\emph{Generalized Spaces for Constructive Algebra},
prepublicació, 2020. \url{https://arxiv.org/abs/2012.13850}.
\newblock (Presentació accessible dels espais generalitzats sense punts, amb l'exemple de les aplicacions exhaustives $\mathbb N\to\mathbb R$.)

\bibitem{caramello}
O.~Caramello,
\emph{Theories, Sites, Toposes: Relating and Studying Mathematical Theories through Topos-Theoretic \enquote{Bridges}},
Oxford University Press, 2018.
\newblock (Referència sobre els toposos classificadors des del punt de vista de la lògica. Primera edició impresa de 2018; la versió electrònica es va publicar el desembre de 2017.)

\bibitem{hott}
The Univalent Foundations Program,
\emph{Homotopy Type Theory: Univalent Foundations of Mathematics},
Institute for Advanced Study, 2013. Accés obert: \url{https://homotopytypetheory.org/book/}.
\newblock (Introducció a la teoria homotòpica de tipus, llegible des de la lògica.)

\bibitem{scott-identity}
D.~S.~Scott,
\emph{Identity and Existence in Intuitionistic Logic},
en M.~Fourman, C.~Mulvey i D.~Scott (eds.), \emph{Applications of Sheaves}, Lecture Notes in Mathematics~753, Springer, 1979, pàg.~660--696.
\newblock (Operacions parcials, existència i identitat; rerefons de les axiomàtiques \enquote{només amb fletxes}.)

\bibitem{freydscedrov}
P.~J.~Freyd i A.~Scedrov,
\emph{Categories, Allegories},
North-Holland Mathematical Library~39, North-Holland, 1990.
\newblock (Presentació sistemàtica sense objectes primitius.)

\bibitem{BS2020}
C.~Benzmüller i D.~S.~Scott,
\emph{Automating Free Logic in HOL, with an Experimental Application in Category Theory},
Journal of Automated Reasoning~64, 2020, pàg.~53--72.
\newblock (Formalització en Isabelle/HOL de diverses axiomàtiques de categories amb lògica lliure.)

\bibitem{BSAFP}
C.~Benzmüller i D.~S.~Scott,
\emph{Axiom Systems for Category Theory in Free Logic},
Archive of Formal Proofs, 2018. \url{https://isa-afp.org/entries/AxiomaticCategoryTheory.html}.
\newblock (Comparació i verificació d'equivalències entre sistemes de Mac~Lane, Scott i Freyd--Scedrov.)

\bibitem{diec}
Institut d'Estudis Catalans,
\emph{Diccionari de la llengua catalana},
2a edició, Barcelona. \url{https://dlc.iec.cat}.
\newblock (Font normativa de la terminologia general catalana.)

\bibitem{termcat}
Universitat Politècnica de Catalunya; Enciclopèdia Catalana,
\emph{Diccionari de matemàtiques i estadística},
1a ed., Barcelona, 2002; 2a ed. en línia, Institut d'Estudis Catalans, 2021. \url{https://cit.iec.cat/DME/}.
\newblock (Font de la terminologia matemàtica catalana: functor, morfisme, preordre, etc.)

\end{thebibliography}

\chapter*{Glossari català--anglès}
\addcontentsline{toc}{chapter}{Glossari català--anglès}
\markboth{Glossari català--anglès}{Glossari català--anglès}

Aquest glossari recull sobretot termes conflictius, no transparents o especialment útils per llegir bibliografia internacional, i no pretén ser un vocabulari exhaustiu de teoria de categories. S'hi dona la forma anglesa i, quan escau, la notació emprada.

\medskip
\noindent\textbf{Criteri d'anglicismes.} Diversos termes s'han manllevat de l'anglès per manca d'una forma catalana establerta o per fidelitat a la bibliografia (\emph{pullback}, \emph{pushout}). En la seva primera aparició es dona, quan existeix, el terme català al costat de l'anglès i de la notació; després se n'usa una sola forma. La terminologia general segueix el \emph{Diccionari de la llengua catalana} de l'Institut d'Estudis Catalans \cite{diec}, i la matemàtica, el \emph{Diccionari de matemàtiques i estadística} del TERMCAT \cite{termcat}.

\medskip

\renewcommand{\arraystretch}{1.25}
\begin{xltabular}{\linewidth}{@{}>{\raggedright\arraybackslash}X >{\raggedright\arraybackslash}X l@{}}
\hline
\textbf{Català} & \textbf{Anglès} & \textbf{Notació}\\
\hline
\endfirsthead
\hline
\textbf{Català} & \textbf{Anglès} & \textbf{Notació}\\
\hline
\endhead
\hline
\endfoot
\endlastfoot
categoria oposada & opposite category & $\cat{C}\op$\\
categoria prima & thin category & \\
categoria sobre un objecte & slice category & $\cat{C}/A$\\
categoria sota un objecte & coslice category & $A/\cat{C}$\\
homconjunt & hom-set & $\Hom(A,B)$\\
functor representable & representable functor & p.\,ex.\ $\Hom(-,A)$\\
immersió de Yoneda & Yoneda embedding & \\
bihomfunctor (o bifunctor Hom) & hom functor, Hom bifunctor & $\Hom(-,-)$\\
component (d'una transformació natural; femení) & component & $\eta_A$\\
antiimatge & preimage, inverse image & $f^{-1}(B)$\\
exhaustiu / essencialment exhaustiu & surjective / essentially surjective & \\
\multicolumn{3}{@{}p{\linewidth}@{}}{\footnotesize En aquest llibre, \emph{exhaustiu} és el terme sistemàtic per a \emph{surjective}; \emph{essencialment exhaustiu} tradueix \emph{essentially surjective}.}\\[0.35em]
transformació natural & natural transformation & $\eta\colon F\Rightarrow G$\\
equivalència de categories & equivalence of categories & $\cat{C}\simeq\cat{D}$\\
monomorfisme / epimorfisme & monomorphism / epimorphism & mono / epi\\
objecte inicial / terminal & initial / terminal object & $0$ / $1$ (quan s'adopta aquesta notació)\\
functor ple / fidel & full / faithful functor & \\
imatge essencial & essential image & \\
categoria d'elements & category of elements & $\textstyle\int P$\\
categoria essencialment petita & essentially small category & \\
quasiinvers (d'una equivalència) & quasi-inverse / pseudo-inverse & \\
prefeix & presheaf & $\cat{Set}^{\cat{C}\op}$\\
producte fibrat (pullback) & pullback & \\
suma amalgamada (pushout) & pushout & \\
con & cone & \\
cocon & cocone & \\
límit / colímit & limit / colimit & $\varprojlim$ / $\varinjlim$\\
unitat / counitat & unit / counit & $\eta$ / $\varepsilon$\\
composició horitzontal & horizontal composition & $\theta\ast\eta$\\
composició amb un functor per l'esquerra & (left) whiskering & $K\eta$\\
composició amb un functor per la dreta & (right) whiskering & $\eta H$\\
epimorfisme regular & regular epimorphism & \\
parella nucli & kernel pair & \\
factorització imatge & image factorization & \\
reindexació & reindexing / base change & $f^{*}$\\
subobjecte & subobject & $\Sub(A)$\\
classificador de subobjectes & subobject classifier & $\Omega$\\
sedàs & sieve & \\
semireticle implicatiu & implicative meet-semilattice & \\
categoria cartesiana tancada & cartesian closed category & \\
categoria regular & regular category & \\
\hline
\end{xltabular}
\renewcommand{\arraystretch}{1.0}

\chapter*{Taula de notació}
\addcontentsline{toc}{chapter}{Taula de notació}
\markboth{Taula de notació}{Taula de notació}

Recollim la notació principal del llibre, amb el seu significat i el lloc de la primera aparició o definició. Les referències remeten al capítol o a la secció corresponent.

\medskip
\renewcommand{\arraystretch}{1.25}
\begin{xltabular}{\linewidth}{@{}l >{\raggedright\arraybackslash}X l@{}}
\hline
\textbf{Símbol} & \textbf{Significat} & \textbf{Primera aparició}\\
\hline
\endfirsthead
\hline
\textbf{Símbol} & \textbf{Significat} & \textbf{Primera aparició}\\
\hline
\endhead
\hline
\endfoot
\endlastfoot
$\cat{C},\cat{D},\dots$ & categories & cap.~\ref{cap:categories}\\
$\Hom(A,B)$, $\cat{C}(A,B)$ & homconjunt de morfismes $A\to B$ & cap.~\ref{cap:categories}\\
$\id_A$ & morfisme identitat de $A$ & cap.~\ref{cap:categories}\\
$g\comp f$ & composició ($f$ seguit de $g$) & cap.~\ref{cap:categories}\\
$\cat{C}\op$ & categoria oposada & cap.~\ref{cap:categories}\\
$f\op$ & morfisme dual de $f$ a $\cat{C}\op$ & \S\ref{sec:dualitat}\\
$\cat{Set},\cat{Grp},\cat{Top},\cat{Graph}$ & categories concretes habituals & cap.~\ref{cap:categories}\\
$\cat{Cat}$ & categoria de les categories petites & cap.~\ref{cap:functors}\\
$\cat{C}/A$, $A/\cat{C}$ & categoria sobre, sota un objecte & cap.~\ref{cap:functors}\\
$(S\downarrow T)$ & categoria coma & \S\refalt{sec:coma}{del cap.~4}\\
$F,G\colon\cat{C}\to\cat{D}$ & functors & cap.~\ref{cap:functors}\\
$\Hom(A,-)$, $\Hom(-,B)$ & homfunctors covariant i contravariant & def.~\ref{def:homfunctors}\\
$\Hom(A,f)=f_*$ & postcomposició amb $f$ (acció de $\Hom(A,-)$ sobre $f$) & not.~\ref{not:hom-morfismes}\\
$\Hom(f,B)=f^*$ & precomposició amb $f$ (acció de $\Hom(-,B)$ sobre $f$). \emph{Atenció:} el mateix símbol designa la reindexació (vegeu més avall) & not.~\ref{not:hom-morfismes}\\
$\widehat{\cat{C}}=\cat{Set}^{\cat{C}\op}$ & categoria de prefeixos sobre $\cat{C}$ & ex.~\ref{ex:categoria-prefeixos}\\
$\Hom(-,-)$, $\Hom(f,g)$ & bihomfunctor i la seva acció $h\mapsto g\comp h\comp f$ & def.~\ref{def:bihomfunctor}\\
$\alpha\colon F\Rightarrow G$ & transformació natural & cap.~\ref{cap:transformacions}\\
$\cat{D}^{\cat{C}}$, $[\cat{C},\cat{D}]$ & categoria de functors & cap.~\ref{cap:transformacions}\\
$\Hom(-,A)$ & exemple de functor representable (contravariant) & cap.~\refcap{cap:yoneda}{4}\\
$\yon$, $\Hom(-,f)$ & immersió de Yoneda i la seva acció $\yon(f)$ sobre $f$ & def.~\refalt{def:immersio-yoneda}{del cap.~4}\\
$\Delta_A$ & functor constant amb valor $A$ & ex.~\ref{ex:functor-constant}\\
$\Delta$ & functor diagonal $\cat{C}\to\cat{C}^{\cat{I}}$ & obs.~\refalt{obs:cons-coma}{del cap.~5}\\
$\delta_A$ & morfisme diagonal $\langle\id_A,\id_A\rangle\colon A\to A\times A$ & cap.~\refcap{cap:cartesianes}{7}\\
$\varprojlim D$, $\colimop D$ & límit, colímit d'un diagrama & cap.~\refcap{cap:limits}{5}\\
$A\times B$, $A+B$ & producte, coproducte & cap.~\refcap{cap:limits}{5}\\
$F\dashv G$ & $F$ és adjunt per l'esquerra de $G$ & cap.~\refcap{cap:adjuncions}{6}\\
$\eta$, $\varepsilon$ & unitat, counitat d'una adjunció & cap.~\refcap{cap:adjuncions}{6}\\
$C^{B}$, $\lambda f$ & exponencial, currificació & cap.~\refcap{cap:cartesianes}{7}\\
$\Sub(A)$ & ordre parcial de subobjectes de $A$ & cap.~\refcap{cap:subobjectes}{8}\\
$f^{*}$ & reindexació per pullback (substitució), a $\cat{C}/B$ i a $\Sub(B)$; a $\cat{Set}$, l'antiimatge $f^{-1}$. És el mateix símbol que la precomposició; el context (homconjunts o subobjectes i fibres) decideix. A $\cat{Set}$ totes dues lectures coincideixen a través de $\Sub(B)\cong\Hom(B,\{0,1\})$ & obs.~\refalt{obs:reindexacio}{del cap.~5}; cap.~\refcap{cap:subobjectes}{8}\\
$\Imatge(f)$ & imatge d'un morfisme & cap.~\refcap{cap:subobjectes}{8}\\
$\exists_f$, $\forall_f$ & quantificadors com a adjunts de $f^{*}$ & cap.~\refcap{cap:subobjectes}{8}\\
$\Omega$, $\mathsf{true}$ & classificador de subobjectes & cap.~\refcap{cap:topos}{9}\\
\hline
\end{xltabular}
\renewcommand{\arraystretch}{1.0}

\printindex

\end{document}
"
%
%   (b) descomentant la línia \publicacioparcialtrue de més avall.
%
% Recordeu recompilar amb XeLaTeX + makeindex + XeLaTeX x2.
% ============================================================
\newif\ifpublicacioparcial
\ifdefined\fascicle\publicacioparcialtrue\else\publicacioparcialfalse\fi
% \publicacioparcialtrue % <- descomenteu per forçar el fascicle sense \fascicle

% Mode de sortida. Per defecte: pantalla/interactiu. El Makefile defineix
% \imprimible per generar una còpia pensada per imprimir: fons blanc,
% enllaços discrets i tests estàtics sense camps AcroForm.
\newif\ifimprimible
\ifdefined\imprimible\imprimibletrue\else\imprimiblefalse\fi
% Nom de la Part IV segons el mode: en paper no hi ha interactivitat.
\ifimprimible
  \newcommand{\NomTests}{Tests d'autoavaluació}
  \newcommand{\nomtests}{tests d'autoavaluació}
\else
  \newcommand{\NomTests}{Tests interactius}
  \newcommand{\nomtests}{tests interactius}
\fi

% Els fitxers compartits viuen ara a la subcarpeta common/ d'aquest mateix
% projecte, de manera que es troben amb ruta explícita, sense dependre de
% TEXINPUTS ni de cap configuració del sistema de compilació.
\usepackage[sobri]{common/aprendes}
% ============================================================
%  Preàmbul compartit — Aprendes Universitat
%  Compartit pels llibres d'Aprendes Universitat.
%  Les dades pròpies de cada llibre (títol, autor i metadades PDF)
%  van al metadata.tex local de cada projecte, carregat ABANS
%  d'aquest fitxer. La macro \hypermetadata es crida DESPRÉS,
%  quan hyperref ja ha estat carregat.
% ============================================================

\usepackage{eurosym}
\usepackage{iftex}
\ifpdftex
  \usepackage[T1]{fontenc}
  \usepackage[utf8]{inputenc}
  \usepackage{lmodern}
\else
  % XeLaTeX / LuaLaTeX: Unicode natiu, res d'inputenc
  \usepackage{fontspec}
  % Carrega Latin Modern per nom de fitxer, que XeTeX/LuaTeX pot
  % resoldre via kpathsea fins i tot si fontconfig no l'ha registrat.
  \setmainfont{lmroman10-regular.otf}[
    BoldFont=lmroman10-bold.otf,
    ItalicFont=lmroman10-italic.otf,
    BoldItalicFont=lmroman10-bolditalic.otf
  ]
  \setsansfont{lmsans10-regular.otf}[
    BoldFont=lmsans10-bold.otf,
    ItalicFont=lmsans10-oblique.otf,
    BoldItalicFont=lmsans10-boldoblique.otf
  ]
  \setmonofont{lmmono10-regular.otf}[ItalicFont=lmmono10-italic.otf]
\fi
% microtype: protrusion i expansió completes només amb pdfLaTeX.
% Amb XeLaTeX/LuaLaTeX només hi ha protrusion i generaria avisos de
% «slot number»; per això només l'activem sota pdfTeX.
\ifpdftex
  \usepackage{microtype}
\fi
\usepackage[catalan, provide=*]{babel}
\usepackage[autostyle=true,babel=true]{csquotes}
% csquotes no porta estil català: el declarem perquè \enquote{...}
% doni guillemets «...» (exteriors) i cometes altes “...” (interiors),
% independentment de la versió de csquotes instal·lada.
\DeclareQuoteStyle{catalan}
  {\guillemotleft}{\guillemotright}
  {\quotedblbase}{\textquotedblleft}
\DeclareQuoteAlias{catalan}{catala}
\usepackage{amsmath,mathtools}% amssymb, amsthm: via aprendes.sty
\usepackage{geometry}
\usepackage{import}
\usepackage{graphicx}
\usepackage{tabularx}
\usepackage{xltabular}% tabularx que pagina (taula de notació)
\usepackage{tikz}
\usetikzlibrary{positioning,arrows.meta,matrix}
\usepackage[all]{xy}% diagrames de grafs (llaços i arestes múltiples) a l'apèndix de grafs
\usepackage{float}
\usepackage[section]{placeins}% FloatBarrier automàtic: cap figura no travessa una secció
\usepackage{pgfplots}
\pgfplotsset{compat=1.18}
\usepackage{imakeidx}
\usepackage{booktabs}

% La bibliografia es configura localment en cada projecte quan hi ha referències.

\setcounter{MaxMatrixCols}{10}

\graphicspath{{img/}}
\geometry{hmargin={3cm,3cm},vmargin={2.5cm,2.5cm},offset=0pt,footskip=1cm,headsep=1cm}


\setlist{itemsep=0.25em, topsep=0.25em}
\numberwithin{equation}{section}
\setcounter{secnumdepth}{3}
\setcounter{tocdepth}{2}

% Configuració dels paràgrafs
\usepackage{indentfirst}
\setlength{\parindent}{1.5em}
\setlength{\parskip}{0.3em}

% ============================================================
%  hyperref: ÚLTIM paquet, com recomana el seu manual.
%  Els colors dels enllaços es fixen aquí. Les metadades
%  pròpies del llibre ja s'han definit abans en metadata.tex;
%  la macro \hypermetadata s'executa després de carregar aquest fitxer.
% ============================================================
\usepackage[pdflang=ca-ES,pdfdisplaydoctitle=true]{hyperref}
% Color dels enllaços: el blau pissarra de la paleta d'aprendes,
% una mica enfosquit. Prou distint del vermell corporatiu dels
% títols, exemples i exercicis, i coherent amb la resta de colors.
\colorlet{EnllacBlau}{ResultatBlau!90!black}
\ifimprimible
  \hypersetup{
    colorlinks=true,
    linktoc=page,
    linkcolor=black,
    urlcolor=black,
    citecolor=black,
    filecolor=black,
    pdfborder={0 0 0},
    hypertexnames=false,
  }
\else
  \hypersetup{
    colorlinks=true,
    linktoc=page,
    linkcolor=EnllacBlau,
    urlcolor=EnllacBlau,
    citecolor=EnllacBlau,
    filecolor=EnllacBlau,
    hypertexnames=false,
  }
\fi
\usepackage{bookmark}% millora els marcadors del PDF; ha d'anar després d'hyperref

% Macros matemàtiques i auxiliars compartides
\newcommand{\Qcb}[1]{#1}
\newcommand{\func}[1]{\operatorname{#1}}
\newcommand{\dbigcup}{\displaystyle\bigcup}
\newcommand{\dbigcap}{\displaystyle\bigcap}
\newcommand{\FRAME}[8]{%
\begin{center}
#5\par\smallskip
\includegraphics[width=#2,height=#3,keepaspectratio]{img/#7}
\end{center}}

\newcommand{\R}{\mathbb{R}}
\newcommand{\N}{\mathbb{N}}
\newcommand{\Z}{\mathbb{Z}}
\newcommand{\Q}{\mathbb{Q}}
\newcommand{\Natpos}{\mathbb{N}^{*}}
\newcommand{\Zstar}{\mathbb{Z}^{*}}
\newcommand{\mcd}{\operatorname{mcd}}
\newcommand{\mcm}{\operatorname{mcm}}
\newcommand{\abs}[1]{\left\lvert #1\right\rvert}
\newcommand{\classe}[1]{\left[#1\right]}
\newcommand{\divides}{\mid}
\newcommand{\notdivides}{\nmid}
\newcommand{\decimalperiodic}[2]{#1,\overline{#2}}
\DeclareMathOperator{\sgn}{sgn} 
% ------------------------------------------------------------
% Macros pròpies de teoria de categories
% ------------------------------------------------------------
\newcommand{\cat}[1]{\mathbf{#1}}          % noms de categories: \cat{Set}
\newcommand{\id}{\operatorname{id}}         % identitat
\providecommand{\Hom}{\operatorname{Hom}}   % conjunts de morfismes
\newcommand{\op}{^{\mathrm{op}}}            % categoria oposada
\newcommand{\Obj}{\operatorname{Obj}}       % col·lecció d'objectes
\newcommand{\comp}{\circ}                    % composició
\DeclareMathOperator{\dom}{dom}
\DeclareMathOperator{\cod}{cod}
\DeclareMathOperator{\Ob}{Ob}
\DeclareMathOperator{\Mor}{Mor}
\DeclareMathOperator{\Aut}{Aut}
\DeclareMathOperator{\Free}{Free}
\DeclareMathOperator{\Nat}{Nat}       % transformacions naturals com a homconjunt
\DeclareMathOperator{\Sub}{Sub}       % preordre de subobjectes
\DeclareMathOperator{\ev}{ev}         % morfisme d'avaluació (exponencials)
\DeclareMathOperator{\Imatge}{Im}     % imatge d'un morfisme (factorització imatge)
% Claudàtors semàntics [[-]] sense dependre de stmaryrd:
\providecommand{\llbracket}{\mathopen{[\![}}
\providecommand{\rrbracket}{\mathclose{]\!]}}
\newcommand{\yon}{\mathrm{y}}          % immersió de Yoneda
\DeclareMathOperator*{\colimop}{colim} % colímit (amb límits sota)

% ============================================================
%  Macros pròpies de la teoria de categories, locals al volum.
%
%  El fitxer compartit ../common/macros.tex barreja notació
%  genèrica (\R, \N, \mcd, decimals periòdics...) amb la de
%  categories. Aquest fitxer recull, en un sol lloc i documentada,
%  només la notació que aquest llibre fa servir, de manera que el
%  volum no depèn de l'organització del fitxer compartit.
%
%  S'usa \providecommand perquè, si macros.tex ja ha definit una
%  macro, no es produeixi cap col·lisió: la definició compartida i
%  la local són idèntiques. Si algun dia es depura macros.tex
%  traient-ne la part de categories, aquest fitxer la garanteix.
% ============================================================

\providecommand{\cat}[1]{\mathbf{#1}}          % noms de categories: \cat{Set}
\providecommand{\id}{\operatorname{id}}         % identitat
\providecommand{\Hom}{\operatorname{Hom}}       % conjunts de morfismes
\providecommand{\op}{^{\mathrm{op}}}            % categoria oposada
\providecommand{\Obj}{\operatorname{Obj}}       % col·lecció d'objectes
\providecommand{\comp}{\circ}                    % composició
\providecommand{\yon}{\mathrm{y}}                % immersió de Yoneda

\providecommand{\dom}{\operatorname{dom}}
\providecommand{\cod}{\operatorname{cod}}

% Predicats relacionals de l'apèndix d'axiomàtica relacional (Dom, Cod, Com,
% Id). S'escriuen amb versaletes matemàtiques per distingir-los clarament dels
% operadors funcionals \dom, \cod: aquí són relacions primitives, no funcions.
\providecommand{\Dom}{\mathsf{Dom}}
\providecommand{\Cod}{\mathsf{Cod}}
\providecommand{\Com}{\mathsf{Com}}
\providecommand{\Id}{\mathsf{Id}}
\providecommand{\Free}{\operatorname{Free}}
\providecommand{\Nat}{\operatorname{Nat}}        % transformacions naturals com a homconjunt
\providecommand{\Sub}{\operatorname{Sub}}        % preordre de subobjectes
\providecommand{\ev}{\operatorname{ev}}          % morfisme d'avaluació (exponencials)
\providecommand{\Imatge}{\operatorname{Im}}      % imatge d'un morfisme

% Claudàtors semàntics [[-]] sense dependre de stmaryrd:
\providecommand{\llbracket}{\mathopen{[\![}}
\providecommand{\rrbracket}{\mathclose{]\!]}}

% colimit amb subíndex a sota (com \lim), només si encara no existeix:
\makeatletter
\@ifundefined{colimop}{\DeclareMathOperator*{\colimop}{colim}}{}
\makeatother

% Remissions a capítols i elements que no són a l'edició parcial (capítols 1--3).
% A l'edició completa fan servir \ref; a la parcial escriuen el número o un text
% alternatiu, per no deixar referències «??».
\newcommand{\refcap}[2]{\ifpublicacioparcial #2\else\ref{#1}\fi}
\newcommand{\refalt}[2]{\ifpublicacioparcial #2\else\ref{#1}\fi}

% ============================================================
%  Estils TikZ comuns per als diagrames del llibre.
%
%  Els diagrames de teoria de categories repeteixen molt codi de
%  puntes de fletxa i de matrius de nodes. Aquests estils permeten
%  escriure els diagrames nous de manera compacta i canviar
%  l'estètica de tots des d'un sol lloc:
%
%     \begin{tikzpicture}
%       \matrix (m) [cat matrix] { A & B \\ C & D \\ };
%       \draw[cat] (m-1-1) -- node[above] {$f$} (m-1-2);
%     \end{tikzpicture}
%
%  Els diagrames antics, escrits amb -{Stealth[scale=1.1]} i amb
%  els paràmetres de matriu explícits, continuen funcionant sense
%  canvis; aquests estils són per a diagrames nous i per a una
%  homogeneïtzació progressiva.
% ============================================================

\tikzset{
  % Fletxa estàndard d'un morfisme.
  cat/.style={-{Stealth[scale=1.1]}},
  % Fletxa de monomorfisme (cua d'ham).
  cat mono/.style={{Hooks[right]}-{Stealth[scale=1.1]}},
  % Fletxa d'epimorfisme (doble punta).
  cat epi/.style={-{Stealth[scale=1.1]Stealth[scale=1.1]}},
  % Fletxa de transformació natural (doble).
  cat nat/.style={-{Implies},double,double distance=1.5pt},
  % Matriu de nodes amb la separació habitual del llibre.
  cat matrix/.style={matrix of math nodes, row sep=2.6em, column sep=3.2em},
}

% Tests interactius en PDF
\newcounter{testquestion}
\newcommand{\testprefix}{}
\newcommand{\testkeys}{}
\newcommand{\iniciTest}[2]{%
  \renewcommand{\testprefix}{#1}%
  \renewcommand{\testkeys}{#2}%
  \setcounter{testquestion}{0}%
  \gdef\testjustificacions{}%
  \label{test:#1}%
  \begin{center}
  \fbox{\begin{minipage}{0.9\textwidth}
  \ifimprimible
    \textbf{Instruccions.} Marca una sola resposta per pregunta. Comprova les respostes i les justificacions a l'apèndix de solucions.
  \else
    \textbf{Instruccions.} Marca una sola resposta per pregunta. En acabar, prem el bot\'{o} \textbf{Corregir}. El camp \textbf{Nota} mostrar\`{a} el resultat.\par\medskip
    \PushButton[name=#1corregir,width=3.2cm,height=0.8cm,bordercolor={0 0 1},onclick={var p='#1';var k='#2'.split(',');var opts=['a','b','c','d'];var s=0;for(var i=0;i<k.length;i++){var ans='';var n=0;for(var j=0;j<opts.length;j++){var q=this.getField(p+'q'+(i+1)+opts[j]);if(q && q.value!='Off'){ans=opts[j];n++;}}var f=this.getField(p+'f'+(i+1));if(n==1 && ans==k[i]){s++;if(f)f.value='Correcte';}else{if(f)f.value=(n==0?'Sense respondre. Correcta: '+k[i]:(n>1?'Marca una sola resposta. Correcta: '+k[i]:'Incorrecte. Correcta: '+k[i]));}}this.getField(p+'nota').value=s+' / '+k.length;app.alert('Nota: '+s+' / '+k.length);}]{Corregir}\quad
    \PushButton[name=#1netejar,width=3.2cm,height=0.8cm,bordercolor={1 0 0},onclick={var keys='#2'.split(',');var opts=['a','b','c','d'];this.getField('#1nota').value='';for(var i=1;i<=keys.length;i++){var f=this.getField('#1f'+i);if(f)f.value='';for(var j=0;j<opts.length;j++){var q=this.getField('#1q'+i+opts[j]);if(q)q.value='Off';}}}]{Netejar}\quad
    \textbf{Nota: }\TextField[name=#1nota,readonly=true,width=4cm,bordercolor={0 0 0}]{}
  \fi
  \end{minipage}}
  \end{center}
  \bigskip
}
\newcommand{\pregunta}{\stepcounter{testquestion}\item}
% Justificació breu de la resposta correcta d'una pregunta. No es mostra
% al costat de les opcions (per no donar pistes de longitud); s'acumula i
% \clauestatica la imprimeix després de la clau de respostes.
\newcommand{\testjustificacions}{}
\makeatletter
\newcommand{\justificacio}[1]{%
  \edef\@justnum{\thetestquestion}%
  \expandafter\g@addto@macro\expandafter\testjustificacions\expandafter{%
    \expandafter\justitem\expandafter{\@justnum}{#1}}%
}
\makeatother
\newcommand{\justitem}[2]{\item[\textbf{#1.}] #2}
\long\def\opcio#1#2{\item \ifimprimible$\square$\else\CheckBox[name=\testprefix{}q\thetestquestion#1,width=1.1em,height=1.1em,bordercolor={0 0 0},onclick={var opts=['a','b','c','d'];for(var j=0;j<opts.length;j++){if(opts[j]!='#1'){var q=this.getField('\testprefix{}q\thetestquestion'+opts[j]);if(q)q.value='Off';}}}]{}\fi\quad #2}
\newcommand{\feedback}{\ifimprimible\par\medskip\else\par\smallskip\textbf{Correcci\'{o}: }\TextField[name=\testprefix{}f\thetestquestion,readonly=true,width=9cm,bordercolor={0.6 0.6 0.6}]{}\par\medskip\fi}
% ---- Validació estructural de la clau (C7) ----
% \finalitzaTest compara el nombre de preguntes realment plantejades
% (\thetestquestion) amb el nombre de respostes de la clau, i emet un
% avís de compilació si no coincideixen. Cal cridar-la després de
% \end{Form}. Fem servir \@for, que recorre llistes separades per comes
% de manera segura, sense bucles \ifx artesanals.
\newcounter{testkeycount}
\makeatletter
\newcommand{\finalitzaTest}{%
  \setcounter{testkeycount}{0}%
  \expandafter\@for\expandafter\next\expandafter:\expandafter=\testkeys\do{\stepcounter{testkeycount}}%
  \ifnum\value{testquestion}=\value{testkeycount}\else
    \PackageWarning{interactive-tests}{El test '\testprefix' te \thetestquestion\space preguntes pero la clau en te \thetestkeycount\space respostes. Reviseu la crida a \string\iniciTest.}%
  \fi
}
\makeatother

% ============================================================
%  Claus estàtiques dels tests, reunides al final del llibre.
%
%  Els tests del llibre es corregeixen amb JavaScript d'AcroForm,
%  que no funciona en molts visors integrats, navegadors ni
%  aplicacions mòbils. Per a aquests casos, i per a l'ús en paper,
%  cada test té una clau de respostes i unes justificacions. Perquè
%  no apareguin just després de les preguntes, cosa que restaria
%  valor a l'autoavaluació, no s'imprimeixen al lloc del test:
%
%   - \clauestatica (cridada just abans de \finalitzaTest) desa la
%     clau i les justificacions del test en curs a \totesclaus i
%     imprimeix només una remissió a l'apèndix de solucions;
%   - \imprimeixsolucionstests, a caps/solucions-tests.tex, les
%     imprimeix totes al final del llibre, cadascuna amb un enllaç de
%     tornada al test.
% ============================================================

\makeatletter
\newcounter{clauestindex}
\gdef\totesclaus{}
\newcommand{\clauestatica}{%
  \begingroup
    \edef\@tmpclau{\noexpand\blocsoluciotest{\testprefix}{\thechapter}{\testkeys}{\unexpanded\expandafter{\testjustificacions}}}%
    \expandafter\g@addto@macro\expandafter\totesclaus\expandafter{\@tmpclau}%
  \endgroup
  \par\medskip
  \begin{center}
  \small\itshape
  \ifimprimible
  Les respostes i les justificacions d'aquest test són a l'apèndix \hyperref[cap:solucions-tests]{\textup{\enquote{Solucions dels tests}}}, pàgina~\pageref{sol:\testprefix}. Consulteu-les després d'haver respost tot el test.
  \else
  Si el vostre lector de PDF no admet el botó \textup{\textbf{Corregir}}, trobareu les respostes i les justificacions d'aquest test a l'apèndix \hyperref[cap:solucions-tests]{\textup{\enquote{Solucions dels tests}}}, pàgina~\pageref{sol:\testprefix}.
  \fi
  \end{center}
  \par\medskip
}
% #1 prefix del test, #2 número de capítol, #3 clau, #4 justificacions
\newcommand{\blocsoluciotest}[4]{%
  \section*{Test del capítol #2}\label{sol:#1}%
  \addcontentsline{toc}{section}{Test del capítol #2}%
  \noindent\textbf{Respostes.}\quad
  \setcounter{clauestindex}{0}%
  \@for\next:=#3\do{%
    \stepcounter{clauestindex}%
    \mbox{\textbf{\theclauestindex.}~\next}\hspace{1.2em plus 0.4em}%
  }%
  \par
  \def\@tmpjust{#4}%
  \ifx\@tmpjust\@empty\else
    \medskip
    \noindent\textbf{Justificacions.}
    \begin{list}{}{\setlength{\leftmargin}{2.2em}\setlength{\labelwidth}{1.8em}\setlength{\itemsep}{0.2em}\small}
    #4
    \end{list}
  \fi
  \par\smallskip
  \noindent{\small\hyperref[test:#1]{Tornar al test} (pàgina~\pageref{test:#1}).}
  \par\bigskip
}
\newcommand{\imprimeixsolucionstests}{\totesclaus}
\makeatother

% ============================================================
%  Entorns didàctics locals del volum de categories.
%  Definits aquí (no a l'estil compartit) per no acoblar-hi
%  el projecte. Reutilitzen tcolorbox, etoolbox i la paleta
%  ja carregats per aprendes.sty en mode «sobri».
%
%    \begin{objectius} ... \end{objectius}   -> bloc d'obertura
%    \begin{resumcap}   ... \end{resumcap}   -> bloc de tancament
%
%  Dins d'aquests blocs s'hi poden usar els encapçalaments
%  \apartat{...} per separar prerequisits/objectius/etc.
% ============================================================

% Color propi per als blocs didàctics: verd pissarra derivat de la
% paleta existent (cexer = verd), enfosquit per contrastar amb el fons.
\colorlet{DidacticVerd}{cexer!55!black}

% Estils de bloc: mateixa gramàtica visual que blocenunciat/blocobservacio.
\tcbset{
  blocdidactic/.style={
    enhanced, breakable, sharp corners,
    frame hidden, boxrule=0pt,
    colback=DidacticVerd!6,
    borderline west={2.0pt}{0pt}{DidacticVerd},
    left=8pt, right=4pt, top=5pt, bottom=5pt,
    before skip=1.0em, after skip=1.0em,
  },
}

% Encapçalament intern d'un apartat del bloc (Prerequisits, Objectius...).
\newcommand{\apartat}[1]{\par\smallskip{\color{DidacticVerd}\ifpdftex\scshape\else\sffamily\bfseries\fi #1}\par\nobreak\smallskip}

% Bloc d'obertura del capítol.
\newenvironment{objectius}{%
  \begin{tcolorbox}[blocdidactic]%
  \ignorespaces
}{\end{tcolorbox}}

% Bloc de tancament del capítol.
\newenvironment{resumcap}{%
  \begin{tcolorbox}[blocdidactic]%
  {\color{DidacticVerd}\ifpdftex\scshape\else\sffamily\bfseries\fi S\'{\i}ntesi del cap\'{\i}tol.}\hspace{0.6em}%
  \ignorespaces
}{\end{tcolorbox}}


\makeindex[title={Índex alfabètic},columns=2,intoc]

\begin{document}
\ifimprimible\pagecolor{white}\else\pagecolor{fons!20}\fi
\hypermetadata

\frontmatter
\maketitle
% Pàgina de crèdits (revers de la portada).
\thispagestyle{empty}
\vspace*{\fill}
\begin{flushleft}
\small
\textbf{\llibretitol}\\
\llibreautor\\[0.6em]
Versió \llibreversio\ · \llibredata\\
\llibreweb\\[1.2em]

\textbf{Llicència.} © \llibreany\ \llibreautor. Aquesta obra està subjecta a la llicència Creative Commons Reconeixement-CompartirIgual 4.0 Internacional (CC BY-SA 4.0): \url{https://creativecommons.org/licenses/by-sa/4.0/deed.ca}. Podeu copiar-la, redistribuir-la i adaptar-la, també amb finalitat comercial, sempre que en reconegueu l'autoria i distribuïu les obres derivades amb la mateixa llicència.\\[1em]

\textbf{Com citar-lo.} \llibreautor, \emph{\llibretitol}, versió \llibreversio, \llibreany. Projecte aprendes\,/\,logcat, \llibreweb.\\[1em]

\textbf{Errates i suggeriments.} Es poden comunicar obrint una incidència (\emph{issue}) al repositori del projecte, \llibreerrates, indicant la versió del llibre i la pàgina. Les correccions s'incorporen a les versions següents.
\end{flushleft}
\clearpage

\tableofcontents
\let\cleardoublepage\clearpage

\chapter*{Prefaci}
\addcontentsline{toc}{chapter}{Prefaci}

Aquest llibre forma part d'\textbf{aprendes}, una col·lecció de materials
d'estudi que comparteixen una mateixa manera de treballar: cada tema es
presenta en una part de \emph{teoria}, es reprèn en una part de
\emph{pràctica} amb problemes resolts, es consolida amb \emph{exercicis
proposats} i es comprova amb \emph{\nomtests}. Qui ja conegui
altres volums de la col·lecció hi reconeixerà l'estructura; qui s'hi
acosti per primera vegada la trobarà explicada a la guia de lectura.

El volum té un propòsit precís. És el primer d'un projecte en dues
etapes: aquí es construeix el llenguatge de la teoria de categories
---des de les definicions bàsiques fins als classificadors de subobjectes
i els topos--- amb el detall que necessitarà un volum posterior dedicat
a la lògica categòrica. Per això la lògica hi és present des del
començament, com a font d'exemples i com a direcció del recorregut,
tot i que el seu desenvolupament sistemàtic queda per a aquella segona
etapa.

\ifpublicacioparcial
Aquesta és una edició parcial, que conté els tres primers temes
---categories, functors i transformacions naturals--- en les quatre parts,
juntament amb els apèndixs i els materials de consulta. Els capítols
restants s'incorporaran en l'edició completa.
\else
Aquesta edició conté el volum complet: nou temes desenvolupats en les
quatre parts i un capítol final que fa de pont cap a la lògica categòrica.
\fi

Abans del primer capítol, la \emph{Introducció} explica per què val la
pena aprendre aquest llenguatge i la \emph{Guia de lectura} proposa com
fer-ho. Es poden llegir en aquest ordre o tornar-hi quan convingui.

\medskip
\noindent\textbf{Prerequisits.} El text pressuposa la maduresa matemàtica habitual d'un primer curs universitari de matemàtiques, informàtica o lògica: treballar amb conjunts i aplicacions, seguir una demostració elemental i llegir notació algebraica bàsica. No es pressuposa haver estudiat teoria de categories, topologia, àlgebra abstracta avançada ni lògica categòrica. Quan un exemple necessita vocabulari addicional, s'introdueix localment o es remet a un apèndix.

\medskip
\noindent\textbf{Errates i suggeriments.} Són benvinguts. Es poden comunicar obrint una incidència al repositori del projecte, \llibreerrates, indicant la versió (\llibreversio) i la pàgina.

% ------------------------------------------------------------
% Espai reservat per a l'autor (descomenteu i completeu):
%
% \medskip
% \noindent\textbf{Agraïments.} ...
%
% (Les errates i els suggeriments: vegeu la pàgina de crèdits i el text següent.)
% ------------------------------------------------------------


\mainmatter
\chapter*{\texorpdfstring{Per què aprendre\\ teoria de categories?}{Per què aprendre teoria de categories?}}
\addcontentsline{toc}{chapter}{Introducció}
\markboth{Introducció}{Introducció}

\begingroup
\renewcommand{\thesection}{\arabic{section}}
\setcounter{section}{0}
\setlength{\emergencystretch}{1.5em}

Hi ha construccions matemàtiques que aprenem diverses vegades, amb noms i
representacions diferents. Aparellar dues dades, imposar simultàniament dues
condicions, trobar un màxim comú divisor o introduir una conjunció semblen
operacions allunyades. I, tanmateix, totes poden respondre a una mateixa
pregunta: com podem reunir dues exigències en una sola dada, de manera que
qualsevol altra manera de satisfer-les passi per aquesta?

La teoria de categories ensenya a reconèixer aquestes formes comunes. No ho
fa declarant que els conjunts, els nombres i les proposicions són la mateixa
cosa. Ho fa distingint, en cada context, quines relacions són rellevants i
com es poden compondre. El que unifica no és la naturalesa dels objectes,
sinó el paper que tenen dins d'un sistema de relacions.

Aquest llibre proposa aprendre aquest llenguatge des dels seus fonaments,
amb quatre famílies d'exemples que ens acompanyaran al llarg del recorregut:
els conjunts, els preordres, els grafs i la lògica. La finalitat és doble:
adquirir les eines bàsiques de la teoria de categories i comprendre per què
aquestes eines preparen una lectura estructural de la lògica. Abans de
començar amb les definicions, convé veure què ens permeten preguntar.

\section{Una estructura també es coneix per les seves relacions}
\label{sec:mirar-fora}

Una manera habitual d'estudiar un objecte consisteix a mirar-ne l'interior:
els elements d'un conjunt, l'operació d'un grup o les parts d'un graf. La
perspectiva categòrica no elimina aquesta mirada, però la complementa amb
una altra: quines transformacions podem fer entre aquest objecte i els
altres? Quines hi arriben, quines en surten i com s'encadenen?
Podem resumir aquest canvi amb una consigna que ens acompanyarà al llarg
del llibre: \emph{relaciona, no individualitza}. En lloc de fixar un
objecte per allò que és aïlladament, el situem pel paper que té dins
d'una xarxa de transformacions.

Entre conjunts tenim aplicacions; entre grups, homomorfismes; entre
preordres, aplicacions monòtones. Aquestes transformacions tenen una
propietat comuna. Si podem passar d'$A$ a $B$ i de $B$ a $C$, podem fer els
dos passos seguits:
\[
  A\xrightarrow{f}B\xrightarrow{g}C,
  \qquad g\comp f\colon A\longrightarrow C.
\]
La composició és associativa: quan tres passos es poden encadenar, agrupar
els dos primers o els dos darrers no canvia el resultat. A més, cada objecte
té una transformació identitat, que el deixa com és i no altera cap altra
transformació quan s'hi compon.

Una \emph{categoria}\index{categoria} reté justament aquestes dades: objectes, fletxes entre
objectes, composició i identitats, amb les lleis que acabem de descriure.
La paraula \emph{fletxa}, o \emph{morfisme}, és més general que la paraula
\emph{funció}. En la categoria lliure\index{categoria!lliure} sobre un graf dirigit, els objectes
són els vèrtexs i les fletxes són els camins: compondre és concatenar-los,
i la identitat és el camí de longitud zero. Un monoide també es pot veure
com una categoria: té un sol objecte, els elements del monoide en són les
fletxes i la multiplicació en dona la composició.

Ni tan sols els mateixos objectes fixen per si sols la categoria. Amb
conjunts i aplicacions obtenim $\cat{Set}$; amb conjunts i relacions,
$\cat{Rel}$\index{Rel@$\cat{Rel}$}. En aquest segon cas, compondre $R\subseteq A\times B$ amb
$S\subseteq B\times C$ significa que
\[
  a\,(S\comp R)\,c
  \quad\Longleftrightarrow\quad
  \exists b\in B\,\bigl(aRb\land bSc\bigr).
\]
La fletxa composta registra que hi ha algun terme intermedi que connecta
els dos extrems. Ja trobem aquí una operació lògica, la quantificació
existencial, dins d'una operació sobre fletxes.

Per tant, la primera pregunta categòrica no és només \enquote{quins objectes tenim?},
sinó \enquote{en quina categoria els estudiem?}. Els nombres enters positius,
ordenats per $\leq$ o per divisibilitat, ofereixen dues categories diferents:
en cadascuna, una fletxa expressa una relació diferent. Escollir les fletxes
és escollir quina part de la matemàtica volem fer visible.

\section{Del que un objecte conté al paper que compleix}

Considerem el producte cartesià de dos conjunts, $A\times B$. El podem
descriure enumerant-ne els elements: són els parells $(a,b)$ amb $a\in A$
i $b\in B$. Però també el podem descriure sense començar pels parells.
Hi ha dues projeccions,
\[
  \pi_A\colon A\times B\longrightarrow A,
  \qquad \pi_B\colon A\times B\longrightarrow B.
\]
Si des d'un mateix conjunt $X$ tenim una aplicació $f\colon X\to A$ i una
aplicació $g\colon X\to B$, hi ha exactament una aplicació cap a $A\times B$
que reuneix aquestes dues dades: $x\mapsto(f(x),g(x))$.

El punt important és que aquesta segona descripció només parla de fletxes.
Té sentit, doncs, en una categoria qualsevol. Un \emph{producte}\index{producte} d'$A$ i
$B$ és un objecte $P$ amb projeccions $\pi_A\colon P\to A$ i
$\pi_B\colon P\to B$ tal que, per a tot objecte $X$ i totes les fletxes
$f\colon X\to A$ i $g\colon X\to B$, existeix una única fletxa
$u\colon X\to P$ amb
\[
  \pi_A\comp u=f,
  \qquad \pi_B\comp u=g.
\]
Aquestes equacions es fan visibles en el diagrama
\[
\begin{tikzpicture}[baseline=(m.center)]
  \matrix (m) [matrix of math nodes,column sep=3.4em,row sep=3em] {
       & X & \\
     A & P & B \\
  };
  \draw[-{Stealth}] (m-1-2) -- node[above left] {$f$} (m-2-1);
  \draw[-{Stealth}] (m-1-2) -- node[above right] {$g$} (m-2-3);
  \draw[-{Stealth},dashed] (m-1-2) -- node[right] {$u$} (m-2-2);
  \draw[-{Stealth}] (m-2-2) -- node[below] {$\pi_A$} (m-2-1);
  \draw[-{Stealth}] (m-2-2) -- node[below] {$\pi_B$} (m-2-3);
\end{tikzpicture}
\]
La fletxa discontínua és la que la propietat ens permet construir.
Dir que el diagrama \emph{commuta} és dir que els camins que comparem,
amb els mateixos extrems, donen la mateixa fletxa.

En un preordre, posem una única fletxa $x\to y$ quan $x\leq y$, i cap quan
aquesta relació no es compleix. El producte de $a$ i $b$ és aleshores el seu
ínfim: un objecte per sota de tots dos, i tan gran com sigui possible entre
els que tenen aquesta propietat. La mateixa condició produeix les
construccions següents, quan existeixen:

\begin{center}
\renewcommand{\arraystretch}{1.2}
\begin{tabularx}{\linewidth}{@{}>{\raggedright\arraybackslash}X
                                  >{\raggedright\arraybackslash}p{0.32\linewidth}@{}}
\toprule
\textbf{Context i fletxes} & \textbf{Producte} \\
\midrule
Conjunts amb aplicacions & Producte cartesià \\
Subconjunts d'un conjunt fix amb inclusions & Intersecció \\
Enters positius amb divisibilitat & Màxim comú divisor \\
Fórmules amb deduïbilitat & Conjunció \\
\bottomrule
\end{tabularx}
\end{center}

La lectura lògica és especialment reveladora. Fixem un sistema deductiu
proposicional amb les regles usuals de la conjunció i posem una fletxa
$p\to q$ quan $q$ és deduïble de $p$. Les projeccions del producte són les
deduccions de $p$ i de $q$ a partir de $p\land q$. Si des de $r$ podem
deduir totes dues fórmules, podem deduir-ne la conjunció:
\[
  r\vdash_T p\land q
  \quad\Longleftrightarrow\quad
  \bigl(r\vdash_T p\ \text{i}\ r\vdash_T q\bigr).
\]
La unicitat de la fletxa és automàtica perquè només registrem si una
deducció existeix, no quina prova la realitza. El producte expressa aquí
una regla lògica amb el mateix patró que expressa l'aparellament de dades.

Això és una \emph{propietat universal}\index{propietat universal}: una caracterització mitjançant
l'existència i la unicitat de fletxes que satisfan unes condicions.
La paraula \enquote{universal} indica que la condició val per a qualsevol objecte
$X$ i qualsevol parell de fletxes adequat, no només per a uns exemples
triats. Bona part del llibre consisteix a reconèixer aquest patró en
construccions successivament més riques.

\section{La mateixa estructura no vol dir el mateix objecte}

Una propietat universal no ens obliga a triar una representació concreta.
Podem construir dues realitzacions diferents d'un producte i obtenir dos
objectes que no siguin literalment iguals. La propietat universal, però,
força un \emph{isomorfisme}\index{isomorfisme} entre ells compatible amb les projeccions.

Una fletxa $f\colon A\to B$ és un isomorfisme si hi ha una fletxa
$g\colon B\to A$ que desfà el canvi en tots dos sentits:
\[
  g\comp f=\id_A,
  \qquad f\comp g=\id_B.
\]
Escrivim aleshores $A\cong B$. El sentit d'aquest símbol depèn de la
categoria: a $\cat{Set}$ parlem de bijeccions; entre grups, d'isomorfismes
de grups; entre espais topològics amb aplicacions contínues, d'homeomorfismes.
No n'hi ha prou de poder desfer una aplicació entre conjunts si la inversa
no respecta l'estructura que hem fixat.

En la categoria prima de la deduïbilitat, dues fórmules són isomorfes quan
són interdeduïbles. Així, $p\land q$ i $q\land p$ poden tenir sintaxis
diferents i el mateix paper deductiu. El llenguatge categòric no les
declara iguals: precisa en quin sentit la diferència de presentació no
altera l'estructura que estem estudiant.

Cal conservar també les dades que acompanyen l'objecte. L'isomorfisme
\emph{únic} entre dos productes és el que respecta les projeccions;
l'objecte subjacent, sense aquestes dades, pot tenir altres automorfismes.
Aprendre a distingir l'objecte de la construcció completa és part de
l'aprenentatge de les propietats universals.

Un altre canvi de perspectiva consisteix a invertir sistemàticament les
fletxes. La categoria oposada\index{categoria!oposada} $\cat{C}\op$ té els mateixos
objectes que $\cat{C}$, però una fletxa $A\to B$ de $\cat{C}$ hi apareix
com una fletxa $B\to A$. No afirmem que la fletxa original tingui inversa:
canviem el context en què la llegim.

En invertir les fletxes de la propietat universal del producte obtenim
el \emph{coproducte}. A $\cat{Set}$ és la unió disjunta; entre subconjunts
ordenats per inclusió és la unió; en el context deductiu apropiat és la
disjunció. De manera semblant, el dual d'un objecte terminal, que rep una
única fletxa des de qualsevol objecte, és un objecte inicial, que n'envia
una única a cadascun.

La \emph{dualitat}\index{dualitat} és més que una simetria de dibuixos. Quan un argument
està formulat categòricament, en podem obtenir l'argument dual invertint
també les composicions i les hipòtesis. Aprenem així a buscar, al costat
d'una pregunta, la pregunta que es veu des de l'altra direcció.

\section{Comparar contextos i comparar les comparacions}

Si les categories descriuen contextos, com podem passar d'un context a
un altre? Un \emph{functor}\index{functor} $F\colon\cat{C}\to\cat{D}$ assigna objectes
a objectes i fletxes a fletxes, preservant els extrems, les identitats i
la composició:
\[
  F(\id_A)=\id_{F(A)},
  \qquad F(g\comp f)=F(g)\comp F(f).
\]
Un functor és, en aquest sentit precís, una traducció estructurada.
Per exemple, podem passar d'un grup al seu conjunt subjacent i d'un
homomorfisme a l'aplicació que el realitza. També podem passar d'un graf
a la categoria dels seus camins.

La precisió importa: tot functor preserva les lleis categòriques, però
no tot functor preserva productes, límits o qualsevol altra estructura
addicional. El llibre ensenyarà a preguntar què preserva cada traducció,
què oblida i què permet recuperar.

Hi ha encara un altre nivell de comparació. Donats dos functors
$F,G\colon\cat{C}\to\cat{D}$, una \emph{transformació natural}\index{transformació natural}
$\eta\colon F\Rightarrow G$ dona una fletxa
$\eta_A\colon F(A)\to G(A)$ per a cada objecte $A$, amb la condició
que aquestes fletxes siguin compatibles amb tots els morfismes de
$\cat{C}$. Per a $f\colon A\to B$, això vol dir
\[
  G(f)\comp\eta_A=\eta_B\comp F(f).
\]
No és una col·lecció de comparacions independents: totes s'han d'ajustar
a una mateixa xarxa de relacions.

Un exemple elemental és la diagonal dels conjunts,
$\delta_A\colon A\to A\times A$, definida per $a\mapsto(a,a)$.
Copiar una dada i després aplicar $f$ a les dues còpies dona el mateix
resultat que aplicar $f$ primer i copiar després:
\[
\begin{tikzpicture}[baseline=(m.center)]
  \matrix (m) [matrix of math nodes,column sep=4em,row sep=3em] {
    A & B \\
    A\times A & B\times B \\
  };
  \draw[-{Stealth}] (m-1-1) -- node[above] {$f$} (m-1-2);
  \draw[-{Stealth}] (m-1-1) -- node[left] {$\delta_A$} (m-2-1);
  \draw[-{Stealth}] (m-1-2) -- node[right] {$\delta_B$} (m-2-2);
  \draw[-{Stealth}] (m-2-1) -- node[below] {$f\times f$} (m-2-2);
\end{tikzpicture}
\]
La naturalitat expressa que l'operació de copiar és compatible amb
qualsevol aplicació, no amb una tria particular de noms per als elements.
En un exemple lògic del llibre, un quadrat d'aquesta mena expressarà
que substituir variables i després interpretar una fórmula dona el mateix
resultat que interpretar-la i transportar-ne el significat.

Aquest nivell de coherència permet formular també l'\emph{equivalència de
categories}\index{equivalència de categories}. Dues categories $\cat{C}$ i $\cat{D}$ són equivalents si hi
ha functors $F\colon\cat{C}\to\cat{D}$ i $G\colon\cat{D}\to\cat{C}$
tals que anar i tornar és naturalment isomorf a no haver canviat de context:
\[
  G\comp F\cong\id_{\cat{C}},
  \qquad F\comp G\cong\id_{\cat{D}}.
\]
No demanem una igualtat literal de presentacions. Demanem que el canvi
conservi l'estructura categòrica essencial i que la recuperació sigui
coherent. L'isomorfisme compara objectes dins d'una categoria;
l'equivalència compara categories senceres.

\section{Posar a prova la mirada des de fora}

Aquesta secció avança idees que el llibre introduirà pas a pas en els
capítols centrals. No cal retenir-ne ara els detalls: n'hi ha prou de
veure cap a on apunten i com continuen la perspectiva de les seccions
anteriors.

La idea de conèixer un objecte per les seves relacions es pot formular
amb exactitud. Per a un objecte $A$ d'una categoria localment petita,
$\Hom(X,A)$ és el conjunt de fletxes de $X$ cap a $A$. Podem pensar cada
objecte $X$ com una forma de posar $A$ a prova, i cada fletxa $X\to A$
com una dada d'$A$ observada amb aquella forma.

A $\cat{Set}$, prendre $X$ amb un sol element recupera els elements
ordinaris d'$A$. En altres categories, però, aquesta única mena de prova
pot resultar insuficient. La perspectiva general considera \emph{tots}
els objectes $X$ i, a més, registra com una prova canvia quan precomponem
amb una fletxa $X'\to X$.

Aquesta informació s'organitza en l'homfunctor\index{homfunctor} $\Hom(-,A)$.
El lema de Yoneda\index{Yoneda!lema} donarà forma precisa a una intuïció anunciada des
del començament: els objectes queden determinats, llevat d'isomorfisme,
per aquest sistema de relacions. No n'hi ha prou de saber quantes
fletxes hi arriben; cal conservar la manera com es transformen i es
componen. Dos homfunctors d'aquest tipus naturalment isomorfs provenen
d'objectes isomorfs.

La \emph{representabilitat} aplica aquesta idea a les construccions.
En el cas del producte, donar una fletxa $X\to A\times B$ equival a
donar dues fletxes amb domini $X$:
\[
  \Hom(X,A\times B)
  \cong \Hom(X,A)\times\Hom(X,B),
\]
de manera compatible amb els canvis de $X$. El producte representa,
amb un sol objecte, l'assignació que a cada $X$ associa aquestes parelles.
Límits i colímits reuniran moltes propietats universals sota un mateix
patró; no seran una llista de receptes sense connexió.

Les \emph{adjuncions}\index{adjunció} ens permetran fer un pas més i relacionar
construccions senceres. Si $F\colon\cat{C}\to\cat{D}$ i
$G\colon\cat{D}\to\cat{C}$ són adjunts, amb $F$ per l'esquerra de
$G$, hi ha una correspondència natural
\[
  \Hom_{\cat{D}}(F(A),B)
  \cong \Hom_{\cat{C}}(A,G(B)).
\]
Una mateixa dada es pot descriure a banda i banda del canvi de context.
Per exemple, donar un functor des de la categoria lliure d'un graf cap
a una categoria és equivalent a donar un morfisme del graf inicial cap
al graf subjacent d'aquella categoria. Els camins compostos no s'han
d'escollir de nou: la composició en determina la imatge.

En un altre exemple, donar una aplicació $X\times A\to B$ equival a
donar una aplicació $X\to B^A$, on $B^A$ és el conjunt d'aplicacions
d'$A$ a $B$. És la \emph{currificació}: una funció de dues variables
es pot descriure com una funció que retorna funcions. El tancament
cartesià generalitza aquesta correspondència i ens porta cap a la
semàntica del càlcul lambda simplement tipat.

\section{La lògica com a fil conductor, no com a drecera}

La lògica apareix en aquest llibre des del primer capítol, però a
diferents nivells. El més elemental és la categoria prima de la
deduïbilitat: les fórmules són objectes i una fletxa registra que una
fórmula es pot deduir d'una altra. Aquesta lectura ajuda a reconèixer
la conjunció com a producte i la interdeduïbilitat com a isomorfisme.

Conservar les proves mateixes és una tasca més exigent. Cal decidir
quan dues proves representen la mateixa fletxa i com la composició
respecta aquesta identificació. No podem passar de l'existència d'una
deducció a una categoria de proves sense afegir aquestes precisions.
La interpretació de tipus i termes en categories cartesianes tancades
ens permetrà començar a veure la relació entre proves, programes i
morfismes, sense confondre aquests nivells.

Més endavant, els \emph{subobjectes}\index{subobjecte} oferiran una altra lectura lògica.
A $\cat{Set}$, un predicat sobre $A$ es pot descriure pel subconjunt
d'elements que el satisfan. Si $f\colon X\to A$, substituir $f(x)$
en el predicat correspon a prendre'n l'antiimatge. El \emph{pullback}\index{pullback}
generalitza aquesta operació, i la reindexació de subobjectes expressa
la substitució sense haver de parlar dels elements d'una categoria
arbitrària.

Les imatges i les adjuncions recuperen llavors l'existencial que havíem
trobat en compondre relacions. Projectar un predicat sobre $X\times A$
cap a $X$ significa retenir els $x$ per als quals hi ha algun $a$ que
el satisfà. En una categoria regular, aquesta construcció dona un
adjunt per l'esquerra de la reindexació. L'adjunt per la dreta, que
interpretaria el quantificador universal, requereix estructura
addicional: no es dedueix de la regularitat per si sola.

Finalment, el \emph{classificador de subobjectes} generalitza una altra
operació familiar. Un subconjunt $S\subseteq A$ es pot descriure per
la seva funció característica $A\to\{0,1\}$. En un topos elemental,
un objecte $\Omega$ fa aquest paper: els morfismes $A\to\Omega$
classifiquen els subobjectes d'$A$. Aquí es retroben dues línies del
llibre, la representabilitat i la lectura dels subobjectes com a
predicats. No hem d'esperar, però, que $\Omega$ es comporti sempre
com els dos valors de veritat clàssics de $\cat{Set}$.

Aquesta és la destinació del volum: arribar al llindar de la lògica
categòrica amb les eines necessàries per comprendre'n el llenguatge.
El desenvolupament sistemàtic de la sintaxi, la semàntica, la correcció,
la completesa i les condicions de coherència queda reservat al volum
següent. El pont final n'explica la continuació; no substitueix aquest
desenvolupament. La lògica dona direcció al recorregut, però no ens
eximeix d'aprendre les construccions categòriques amb les seves
hipòtesis i demostracions.

\section{Què vol dir haver après aquest llenguatge}

Aprendre categories no consisteix a abandonar els exemples concrets
tan aviat com tenim una definició general. Consisteix a poder fer el
camí d'anada i tornada: reconèixer una estructura en un exemple,
formular-la amb fletxes i tornar a interpretar-la en un altre context.
Una bona abstracció ens ha de permetre veure més, no només dir menys.

Al llarg del llibre adquirirem l'hàbit de preguntar-nos:
\begin{itemize}
  \item Quins són els objectes i les fletxes, i què significa compondre-les?
  \item Quina propietat expressa el diagrama que tenim davant?
  \item Quines dades caracteritzen una construcció i quines en són
        només una realització particular?
  \item Què queda determinat llevat d'isomorfisme, i quines
        compatibilitats fan únic aquest isomorfisme?
  \item Què preserva el functor amb què canviem de context?
  \item Quin és l'enunciat dual i sota quines hipòtesis té sentit?
  \item Quina lectura lògica té la construcció, si el context la permet?
\end{itemize}

Els primers capítols fixaran el vocabulari; els següents ensenyaran
a organitzar-lo amb propietats universals, representabilitat, límits
i adjuncions; els darrers mostraran com aquestes eines es combinen
en el tancament cartesià, els subobjectes i els topos. La guia de
lectura que acompanya el llibre orienta sobre com treballar aquest
recorregut amb la teoria, la pràctica, els exercicis i els tests.

El resultat que busquem no és disposar d'una paraula nova per a cada
construcció coneguda. És poder reconèixer, davant d'una situació nova,
quines relacions en governen l'estructura. Quan una intersecció, una
conjunció i un producte deixen de ser semblances vagues i esdevenen
manifestacions d'una mateixa propietat, el llenguatge categòric
comença a fer la seva feina.

\endgroup

\chapter*{\texorpdfstring{Guia de lectura i itineraris\\ d'aprenentatge}{Guia de lectura i itineraris d'aprenentatge}}
\addcontentsline{toc}{chapter}{Guia de lectura}
\markboth{Guia de lectura}{Guia de lectura}

\begingroup
\renewcommand{\thesection}{\arabic{section}}
\setcounter{section}{0}
\setlength{\emergencystretch}{1.5em}

Aquest capítol és un mapa per a la primera lectura. No cal aprendre
ara els noms de totes les construccions que hi apareixen: les trobarem
definides al seu lloc. La seva funció és mostrar què aporta cada etapa,
com es relacionen les quatre parts del llibre i com podem comprovar
que l'estudi ens està fent avançar.

La introducció explica \emph{per què} val la pena aprendre categories.
Aquí ens ocuparem de \emph{com} fer-ho. Podeu llegir aquesta guia
abans de començar i tornar-hi en acabar un bloc de capítols, quan
els noms del mapa ja s'hagin convertit en construccions familiars.


\section{El punt de partida i la destinació}

El llibre no pressuposa coneixements previs de teoria de categories.
Sí que pressuposa una familiaritat elemental amb el llenguatge
matemàtic: conjunts, aplicacions, composició, relacions i
demostracions senzilles. Per començar, convé poder fer quatre coses:
\begin{itemize}
  \item distingir el domini i el codomini d'una aplicació i calcular
        una composició;
  \item treballar amb productes cartesians, subconjunts i antiimatges;
  \item llegir quantificadors i separar una afirmació d'existència
        d'una afirmació d'unicitat;
  \item distingir una implicació, el seu recíproc i una equivalència,
        i entendre el paper d'un contraexemple.
\end{itemize}
Si alguna d'aquestes operacions és poc familiar, és millor dedicar-hi
una breu revisió que atribuir-ne la dificultat a l'abstracció categòrica.

No cal dominar prèviament tots els àmbits que proporcionen exemples.
Els grups, els espais vectorials i la topologia enriqueixen la lectura,
però, si encara no els coneixeu, podeu començar amb conjunts, preordres
i grafs. Els exemples del dual i del bidual demanen àlgebra lineal;
els que utilitzen continuïtat demanen nocions de topologia. Es poden
reprendre en una segona lectura sense abandonar el fil general.
La lògica aporta exemples des del començament, però no cal conèixer
ja la lògica categòrica per seguir-los.

La destinació és una comprensió inicial de les categories, els
functors i la naturalitat, seguida de les propietats universals,
Yoneda, els límits i les adjuncions. Aquest bagatge permet arribar
al tancament cartesià, als subobjectes i a una primera introducció
als topos. El volum prepara l'estudi posterior de la lògica categòrica;
no pretén desenvolupar-ne aquí una teoria completa.

\section{Quatre parts per a un mateix aprenentatge}

El llibre separa físicament la \textbf{Teoria}, la \textbf{Pràctica},
els \textbf{Exercicis proposats} i els \textbf{\NomTests}.
Aquesta separació facilita la consulta, però no obliga a estudiar
primer tota la teoria i només després tota la pràctica. Per a una
primera lectura, és més útil treballar \emph{per temes}: llegir un
capítol teòric, estudiar els problemes resolts corresponents,
intentar exercicis i comprovar la comprensió amb el test del mateix
tema.

\begin{center}
\renewcommand{\arraystretch}{1.2}
\begin{tabularx}{\linewidth}{@{}>{\raggedright\arraybackslash}p{0.25\linewidth}
                                  >{\raggedright\arraybackslash}X@{}}
\toprule
\textbf{Part} & \textbf{La pregunta que ajuda a respondre} \\
\midrule
Teoria & Quina és la definició, per què té sentit i què se'n pot deduir? \\
Pràctica & Com s'utilitza en un problema i com s'escriu l'argument? \\
Exercicis proposats & Puc reconèixer i aplicar la idea sense una solució completa? \\
\NomTests & Distingeixo el concepte de les confusions que s'hi assemblen? \\
\bottomrule
\end{tabularx}
\end{center}

Els nou primers temes tenen capítols corresponents en les quatre
parts. La numeració torna a començar en cada part; per això, quan
cerqueu material d'un tema, fixeu-vos tant en el número com en la
part i el títol. El capítol teòric desè, \emph{Pont cap a la lògica
categòrica}, és un epíleg d'orientació i no té un capítol paral·lel
de pràctica, exercicis o tests.

Els blocs d'obertura dels capítols teòrics indiquen prerequisits
i objectius. Llegiu-los com una invitació a activar coneixements
anteriors. Els blocs de síntesi del final recullen les idees centrals,
els errors típics i lectures recomanades. Serveixen per revisar,
però no substitueixen els exemples i les demostracions del capítol.

\section{El mapa conceptual dels capítols}
\label{sec:mapa-lectura}

\ifpublicacioparcial
\noindent\textbf{Nota sobre aquesta edició.} Aquesta versió parcial
només inclou els capítols 1--3. El mapa descriu, tanmateix, els deu
capítols del volum complet, perquè pugueu situar el que estudieu
dins del recorregut sencer.
\fi

El recorregut es pot agrupar en quatre etapes. Els capítols 1--3
construeixen el llenguatge i els nivells de comparació; els 4--6
organitzen les construccions mitjançant propietats universals;
els 7--9 combinen aquestes eines en estructures amb lectura lògica;
el 10 mostra la continuació. Aquesta divisió és una ajuda per orientar-se,
no una frontera rígida: les idees de les primeres etapes reapareixen
constantment en les següents.

\subsection*{1. Categories: aprendre la gramàtica de les fletxes}

El primer capítol introdueix objectes, morfismes, identitats i
composició. També ensenya a llegir diagrames, construeix exemples
i diferencia isomorfismes, monomorfismes, epimorfismes, seccions
i retraccions. La categoria oposada i la dualitat permeten
reutilitzar definicions i arguments en l'altra direcció.

És un capítol extens perquè fixa els hàbits que necessitarem després.
No convé voler memoritzar tots els exemples alhora. Escolliu-ne uns
quants i pregunteu sempre què són els objectes, què són les fletxes
i com es componen. Les construccions de categories noves i les
precaucions de mida amplien aquest vocabulari.

\noindent\textbf{Comprovació de progrés.} Podeu verificar els axiomes
en un exemple concret i explicar per què una fletxa és, o no és,
un isomorfisme? Podeu escriure l'equació d'un quadrat commutatiu
sense invertir l'ordre de la composició?

\subsection*{2. Functors: passar d'un context a un altre}

El segon capítol puja un nivell: no compara només objectes d'una
categoria, sinó categories entre elles. Els functors actuen sobre
objectes i morfismes i preserven identitats i composició. El graf
subjacent i la categoria lliure permeten veure una parella de
construccions que reapareixerà com a adjunció.

Cal dedicar una atenció especial als homfunctors. $\Hom(A,-)$
actua per postcomposició; $\Hom(-,B)$, per precomposició i amb
inversió de sentit. Aquesta variància no és un detall de notació:
és la base de la representabilitat i de Yoneda. Les nocions de
functor ple i fidel descriuen l'acció sobre cada homconjunt.

\noindent\textbf{Comprovació de progrés.} Per definir un functor,
podeu donar-ne l'acció sobre les fletxes, no només sobre els objectes?
Podeu explicar què fa $\Hom(-,B)$ amb una fletxa $X'\to X$?

\subsection*{3. Transformacions naturals i equivalències: exigir coherència}

El tercer capítol compara functors. La naturalitat exigeix que
una família de morfismes sigui compatible amb totes les fletxes
del domini. Els exemples de la interpretació de fórmules i del
bidual mostren per què aquesta compatibilitat és matemàticament
significativa, no només una convenció gràfica.

Les categories de functors permeten tractar aquestes comparacions
com a morfismes. L'equivalència de categories precisa quan dos
contextos descriuen la mateixa estructura essencial sense tenir
la mateixa presentació. La composició horitzontal afegeix un
segon tipus de composició que cal distingir de la vertical.

\noindent\textbf{Comprovació de progrés.} Podeu dibuixar el quadrat
de naturalitat d'una família proposada i comprovar-lo per a una
fletxa arbitrària? Podeu distingir un isomorfisme d'objectes,
un isomorfisme natural i una equivalència de categories?

\subsection*{4. Propietats universals, representabilitat i Yoneda:
caracteritzar pel paper}

Aquí canvia el centre de gravetat. En lloc de construir un objecte
i estudiar-ne després les propietats, busquem un objecte que
resolgui una necessitat expressada amb fletxes. Els objectes
inicials i terminals ofereixen els primers casos; les categories
coma ajuden a reunir dades i compatibilitats en un mateix context.

La representabilitat expressa aquestes caracteritzacions amb
homfunctors. Yoneda fa precisa la idea de conèixer un objecte
pel sistema de morfismes que hi arriben. És normal que aquesta
etapa demani més d'una lectura: cal coordinar categories, functors,
transformacions i elements d'homconjunts.

\noindent\textbf{Comprovació de progrés.} Podeu demostrar que dos
objectes terminals són isomorfs i indicar on s'utilitza la unicitat?
En una representació, podeu distingir el functor representat,
l'objecte representant i la dada universal? Podeu seguir la
bijecció de Yoneda a partir de l'avaluació en la identitat?

\subsection*{5. Límits i colímits: reunir construccions en un patró}

Un diagrama fixa una configuració d'objectes i morfismes; un con
és una manera compatible de relacionar-s'hi. El límit és el con
universal. Segons la forma del diagrama, obtenim productes,
equalitzadors, pullbacks o l'objecte terminal. Els colímits en
són la versió dual.

El pullback mereix una atenció especial, perquè reapareixerà
en la reindexació i els subobjectes. També cal separar dues
preguntes: si una categoria té un límit i si un functor el
preserva. No són la mateixa afirmació.

\noindent\textbf{Comprovació de progrés.} Podeu calcular un pullback
de conjunts i comprovar-ne la propietat universal? Podeu
distingir el diagrama, el con i el seu vèrtex, i explicar
què vol dir que un functor preserva aquest límit?

\subsection*{6. Adjuncions: dues descripcions de la mateixa dada}

Les adjuncions relacionen functors mitjançant una bijecció
natural d'hom\-con\-junts. El capítol comença amb preordres,
on la correspondència es llegeix com una equivalència de
desigualtats. Després introdueix unitat, counitat i identitats
triangulars, i mostra com les adjuncions expliquen la preservació
de límits i colímits.

Per a una primera aproximació, escolliu una adjunció i seguiu-la
en les dues descripcions: per bijeccions d'homconjunts i per
unitat i counitat. Les construccions lliures i la currificació
connecten aquest capítol amb exemples ja coneguts.

\noindent\textbf{Comprovació de progrés.} Podeu dir quins morfismes
es corresponen en una adjunció concreta i construir-ne els
transposats? Podeu recordar, amb una justificació conceptual,
que els adjunts per la dreta preserven límits i els de
l'esquerra preserven colímits?

\subsection*{7. Categories cartesianes tancades: funcions dins del context}

Els productes i els exponencials permeten parlar d'aparellament,
avaluació i currificació dins d'una categoria. El capítol
caracteritza el tancament cartesià i n'explica la relació
amb el càlcul lambda simplement tipat: els tipus es llegeixen
com a objectes i els termes, en context, com a morfismes.

El conjunt de totes les aplicacions $B\to C$ ofereix l'exemple
més familiar d'exponencial. Els preordres permeten reconèixer
la implicació com a adjunt de la conjunció. Els contraexemples
recorden que no tota categoria té aquesta estructura.

\noindent\textbf{Comprovació de progrés.} Podeu passar d'una fletxa
$A\times B\to C$ a la seva currificació $A\to C^B$ i tornar-ne
amb l'avaluació? Podeu explicar per què tenir productes no
garanteix tenir exponencials?

\subsection*{8. Subobjectes, imatges i factoritzacions: de les parts als predicats}

Un subobjecte no és simplement un subconjunt escrit amb un
altre nom: és una classe de monomorfismes amb un mateix
codomini, identificats de manera compatible. La reindexació
per pullback generalitza l'antiimatge i ofereix una lectura
estructural de la substitució.

Les interseccions expressen conjuncions de predicats; les
imatges i la regularitat permeten construir un adjunt
existencial de la reindexació. Aquesta etapa obliga a
controlar les hipòtesis: la regularitat no garanteix per
si sola unions de subobjectes ni un quantificador universal.

\noindent\textbf{Comprovació de progrés.} Podeu calcular una
reindexació a $\cat{Set}$ i interpretar-la com a substitució?
Podeu descriure la imatge d'un predicat sota una projecció
i explicar la relació $\exists_f\dashv f^*$?

\subsection*{9. Classificadors de subobjectes i topos: retrobar la representabilitat}

El classificador generalitza la funció característica d'un
subconjunt. La correspondència entre subobjectes d'$A$ i
morfismes $A\to\Omega$ recupera la representabilitat del
capítol quart. El topos elemental combina límits finits,
exponencials i classificador de subobjectes.

El capítol presenta $\cat{Set}$ i les categories de prefeixos
com a exemples. L'objectiu és comprendre les definicions
i la connexió entre les peces, no dominar encara la
teoria de feixos o la lògica interna dels topos.

\noindent\textbf{Comprovació de progrés.} Podeu reconstruir un
subconjunt com l'antiimatge del valor vertader per la
seva funció característica? Podeu enunciar les tres
condicions d'un topos i explicar què aporta cadascuna?

\subsection*{10. Pont cap a la lògica categòrica: saber què queda preparat}

L'epíleg situa les eines adquirides en un programa més ampli:
contextos, substitució, predicats, quantificadors, teories
i models. El mapa de fragments lògics i l'exemple de la
categoria sintàctica orienten sobre què es prolonga al
volum següent.

Llegiu-lo com una mirada retrospectiva i una invitació a
continuar, no com una exigència de dominar de cop tots els
temes anunciats. La sintaxi general, la correcció, la
completesa i les condicions de coherència pertanyen a
l'estudi posterior.

\noindent\textbf{Comprovació de progrés.} Quan l'epíleg parla de
models com a functors o de quantificadors com a adjunts,
podeu identificar les construccions d'aquest volum que
donen sentit a aquestes expressions?

\section{Com treballar un capítol}

Una lectura activa alterna comprensió, reconstrucció i aplicació.
Podeu utilitzar el cicle següent, adaptant-ne el ritme a la
dificultat del tema.

\begin{enumerate}
  \item \textbf{Situar la pregunta.} Llegiu l'obertura i els
        objectius del capítol. Escriviu en una frase quin
        problema intenta resoldre i quines idees anteriors
        necessita.
  \item \textbf{Desmuntar la definició.} Separeu les dades
        que s'han de donar, les operacions que s'hi afegeixen
        i les propietats que han de satisfer. Anoteu el
        domini i el codomini de cada fletxa.
  \item \textbf{Ancorar-la en exemples.} Treballeu un cas
        amb conjunts i un altre amb preordres o grafs.
        Quan pertoqui, afegiu-hi la lectura lògica. Busqueu
        també un exemple que no compleixi la definició.
  \item \textbf{Reconstruir una demostració.} Tanqueu el
        text i intenteu recuperar un argument curt. Després
        compareu-lo amb el llibre: on s'usa cada hipòtesi?
        On apareix l'existència i on la unicitat?
  \item \textbf{Fer servir la pràctica amb pausa.} Llegiu
        l'enunciat d'un problema resolt i proveu de fer-ne
        el primer pas abans de mirar la solució. Un cop
        llegida, torneu a escriure l'argument sense copiar-lo.
  \item \textbf{Intentar un exercici nou.} Comenceu pels
        de consolidació. Si us encalleu, identifiqueu el
        punt exacte i consulteu la indicació, no només
        una altra solució que s'hi assembli.
  \item \textbf{Tancar el cicle.} Feu el test, reviseu els
        errors i rellegiu la síntesi. Escriviu què podeu
        fer ara que no podíeu fer abans i què ha de quedar
        anotat per a una segona lectura.
\end{enumerate}

No cal reconstruir immediatament totes les proves llargues.
Però convé entendre'n l'estratègia i ser capaç de reproduir
arguments curts que reapareixen: provar una naturalitat,
usar la unicitat d'una propietat universal o construir
una fletxa amb el domini i el codomini correctes.

\subsection*{Una pauta per llegir diagrames}

Al llarg del llibre, quan una propietat categòrica té una expressió
natural amb fletxes, la presentem alhora amb un diagrama i amb les
equacions de composició que aquest expressa. Els diagrames formen
part del llenguatge matemàtic, no són una il·lustració opcional.
Aprofiteu aquesta doble presentació: passar del dibuix a les equacions
i de les equacions al dibuix és una destresa que cal practicar.

Davant d'un diagrama, no comenceu per recordar-ne el nom.
Seguiu aquesta pauta:
\begin{enumerate}
  \item Identifiqueu els objectes i els extrems de cada fletxa.
  \item Distingiu les dades donades de la fletxa que s'ha de construir.
  \item Escriviu els camins que tenen el mateix origen i el mateix destí.
  \item Traduïu-ne la commutativitat en equacions de composició.
  \item Si hi ha una propietat universal, digueu sobre quines
        dades es quantifica i què es demana que sigui únic.
\end{enumerate}

Una fletxa discontínua en un diagrama universal no es construeix
perquè el dibuix ho suggereixi: l'existència es justifica amb
una hipòtesi o una propietat. Igualment, un quadrat commutatiu
no és automàticament un pullback. Cal comprovar-ne també la
propietat universal.

\section{Els exemples recurrents com a punts de suport}

Canviar d'exemple no ha de significar tornar a començar.
Les mateixes famílies permeten veure cares diferents d'una idea.

\begin{description}
  \item[Conjunts.] Permeten calcular explícitament amb elements,
        productes, antiimatges, funcions i imatges. Són un
        primer laboratori, però no autoritzen a usar elements
        dins d'una categoria arbitrària.
  \item[Preordres.] Simplifiquen els homconjunts: hi ha com a
        màxim una fletxa entre dos objectes. Això converteix
        propietats universals en condicions d'ordre i adjuncions
        en equivalències de desigualtats. La unicitat automàtica
        és una ajuda, però també pot amagar una dificultat
        que reapareix en categories no primes.
  \item[Grafs.] Fan visible la diferència entre tenir fletxes
        i disposar de composició. Cal distingir la categoria
        de grafs, on les fletxes són homomorfismes de grafs,
        de la categoria lliure d'un graf, on són camins.
        La construcció lliure connecta els primers capítols
        amb les propietats universals i les adjuncions.
  \item[Lògica.] Permet interpretar deducció, conjunció,
        substitució i quantificació en els contextos adequats.
        Cal separar sempre la categoria prima que registra
        deduïbilitat de la perspectiva que conserva les proves,
        i també dels subobjectes que representen predicats.
\end{description}

Un quadern de correspondències pot ajudar: per a cada nova
construcció, escriviu què significa en dues d'aquestes famílies,
quina propietat comparteixen i quines hipòtesis necessita
cada realització. Si no existeix en un dels exemples, anoteu-ho:
aquest límit també forma part de la comprensió.

\section{Itineraris possibles}

Els itineraris següents són propostes de treball, no permisos
per ignorar les dependències. L'ordre del llibre continua sent
el recorregut més segur per a una primera lectura completa.

\subsection*{Una primera entrada: els tres capítols inicials}

Treballeu categories, functors i transformacions naturals amb
els seus materials paral·lels de pràctica, exercicis i tests.
La fita és poder moure-us entre els tres nivells: objectes i
morfismes, categories i functors, functors i transformacions.
No cal dominar encara tot el desplegament de la composició
horitzontal, però sí distingir-la de la vertical.

\ifpublicacioparcial
Aquesta edició parcial conté precisament aquests tres temes en
les quatre parts, a més dels materials complementaris. No és una
introducció completa als límits, les adjuncions o els topos;
les remissions als capítols absents s'han d'entendre com a
indicacions per continuar amb el volum complet.
\fi

\subsection*{Un recorregut complet d'autoaprenentatge}

Seguiu els nou temes amb el cicle de treball per capítols i
llegiu després el pont final. Cada cop que arribeu a una
definició general, torneu a un exemple anterior. Després
de Yoneda, reviseu els homfunctors; després de les adjuncions,
reviseu la construcció lliure; després del classificador,
reviseu la representabilitat. Aquestes tornades no són
repeticions inútils: fan visible la unitat del recorregut.

\subsection*{Un curs semestral}

Per a un curs d'un semestre, una selecció natural és treballar els capítols 1--6 com a nucli obligatori i triar, segons l'orientació del curs, els capítols 7--9. La teoria s'ha de combinar cada setmana amb una selecció de problemes resolts i exercicis proposats; els tests poden servir d'autoavaluació, no de substitut d'una prova escrita. El capítol 10 funciona bé com a sessió final de síntesi i d'obertura cap a la lògica categòrica.

\subsection*{Una lectura per al professorat}

Per preparar docència, convé començar per aquesta guia i pels blocs d'objectius i d'errors típics de cada capítol. Les parts de Pràctica permeten seleccionar exemples per a classe, mentre que els Exercicis proposats ofereixen material de treball autònom. Els quatre fils recurrents ---conjunts, preordres, grafs i lògica--- es poden modular segons el perfil de l'alumnat, però és recomanable conservar-ne almenys dos en paral·lel per evitar que les definicions quedin lligades a una sola categoria concreta.

\subsection*{Una lectura orientada a la lògica}

Manteniu el recorregut general, però doneu més pes a la
deduïbilitat del capítol 1 i als exemples d'interpretació
i substitució dels capítols 2 i 3. Als capítols 4--6,
consolideu representabilitat, pullbacks i adjuncions;
després, treballeu amb especial cura el tancament cartesià,
els subobjectes i el classificador.

No convé saltar directament a l'epíleg: expressions com
\enquote{el quantificador és un adjunt} orienten, però només es
comprenen quan sabem entre quines categories o preordres
actua l'adjunció i què afirma la seva correspondència.
Les equivalències lògiques no substitueixen la comprovació
de les hipòtesis categòriques.

\subsection*{Una lectura orientada als tipus i a la computació}

Feu servir els grafs i les categories lliures per entendre
la composició. Després de consolidar functors i naturalitat,
poseu l'accent en propietats universals, productes,
adjuncions, exponencials i currificació. El capítol 7
reuneix aquestes eines en la semàntica de tipus i termes.

Els capítols 8 i 9 poden constituir una segona etapa
si el primer objectiu és el càlcul lambda simplement tipat.
Caldrà reprendre'ls per arribar als predicats, els
classificadors i la perspectiva lògica més àmplia.

\section{Convencions tipogràfiques i de lectura}

Els noms de categories s'escriuen amb tipografia matemàtica, com $\cat{Set}$, $\cat{Graph}$ o $\cat{Cat}$. Els termes definits i les idees que cal retenir apareixen destacats dins dels entorns corresponents; el color reforça la jerarquia visual, però l'etiqueta textual ---\enquote{Definició}, \enquote{Teorema}, \enquote{Exemple}, \enquote{Exercici}, etc.--- n'indica sempre la funció. Les fletxes discontínues dels diagrames assenyalen habitualment el morfisme que una propietat universal afirma que existeix; el text adjacent n'explicita el paper. Els símbols amb més d'un ús, especialment $f^{*}$, es desambigüen pel context i es recullen a la taula de notació.

\section{Exercicis i tests: què comproven i què no}

Els exercicis proposats utilitzen una estrella,
$\star$, per assenyalar consolidació, i dues,
$\star\star$, per assenyalar ampliació. Comenceu pels
primers, però no interpreteu les marques com una mesura
exacta del temps que necessitareu. Una dificultat breu
de variància pot resultar més important que un càlcul llarg.

Les indicacions són una ajuda per reprendre el treball.
Abans de consultar-les, escriviu les dades, la conclusió
i almenys un intent. Després d'utilitzar una indicació,
torneu a fer l'exercici sense mirar-la. L'objectiu no és
haver reconegut una solució, sinó poder construir un argument.

Els tests tenen una única resposta correcta per pregunta.
\ifimprimible
En aquesta edició per imprimir, les respostes i les
justificacions de cada test es reuneixen al final del llibre
(apèndix \enquote{Solucions dels tests}). L'edició de pantalla
ofereix, a més, correcció automàtica en els lectors de PDF
compatibles.
\else
La correcció automàtica necessita un lector de PDF compatible
amb formularis AcroForm i JavaScript. Si el lector no
l'admet, podeu respondre igualment les preguntes i
consultar les respostes i les justificacions de cada test,
reunides al final del llibre (apèndix \enquote{Solucions dels
tests}), també usables en paper.
\fi
Consulteu-les només després d'haver respost tot el test.

Una resposta encertada no demostra per si sola que sapigueu
resoldre un problema. Per aprofitar el test, justifiqueu
per què l'opció triada és correcta i per què les altres
fallen. Si heu dubtat entre dues opcions, convertiu el
dubte en una pregunta per revisar a la teoria o a la
pràctica. Aquesta explicació és més útil que la puntuació
obtinguda sense justificació.

\section{Algunes dificultats que convé reconèixer aviat}

Hi ha errors que travessen diversos capítols. Identificar-los
permet corregir la manera de treballar, no només una resposta.

\begin{description}
  \item[Compondre sense comprovar els tipus.]
        Abans d'escriure $g\comp f$, verifiqueu que el
        codomini de $f$ és el domini de $g$. L'associativitat
        no fa componibles fletxes amb extrems incompatibles.
  \item[Confondre un exemple amb la definició general.]
        A $\cat{Set}$, mono i epi corresponen a injectivitat
        i exhaustivitat. En altres categories, les definicions
        són condicions de cancel·lació. Un mono que també
        és epi no és sempre un isomorfisme.
  \item[Tractar la naturalitat com una etiqueta.]
        Una família de fletxes no és natural només perquè
        s'hagi definit amb la mateixa fórmula. Cal escriure
        el quadrat i comprovar-lo per a tots els morfismes.
  \item[Oblidar les dades universals.]
        El producte inclou les projeccions; el límit inclou
        el con; l'exponencial inclou l'avaluació. La unicitat
        de l'isomorfisme es refereix a les compatibilitats
        amb aquestes dades, no a l'objecte nu.
  \item[Canviar l'orientació sense controlar-la.]
        La contravariància, la dualitat i la reindexació
        inverteixen direccions de maneres que cal precisar.
        Escriviu dominis i codominis abans de fer servir
        una regla recordada de memòria.
  \item[Afegir estructura que no s'ha suposat.]
        Tenir productes no implica tenir exponencials;
        la regularitat no dona per si sola disjuncions
        ni quantificació universal; un topos no és
        necessàriament booleà. Rellegiu les hipòtesis.
\end{description}

Si una demostració sembla bloquejada, pregunteu-vos primer
si el problema és de definició, de tipus, de variància o
d'hipòtesis. Moltes dificultats aparentment abstractes
es resolen quan precisem una d'aquestes quatre coses.

\section{Quan consultar els apèndixs i els materials finals}

Els apèndixs no constitueixen una etapa que s'hagi de
completar obligatòriament abans de començar la teoria.
Es poden consultar quan la lectura ho demani.

\begin{description}
  \item[Apèndix A: axiomàtica relacional.]
        Ofereix una presentació en què les fletxes són
        els individus primitius i els objectes es
        reconstrueixen a partir de les identitats.
        També hi trobareu la formalització del principi
        de dualitat, en versió sintàctica i semàntica.
        És una ampliació adequada després del primer
        capítol, sobretot per a qui s'interessi per
        la formalització lògica. Demana familiaritat
        amb lògica de primer ordre amb igualtat.
  \item[Apèndix B: grafs.]
        Aclareix les convencions, els homomorfismes,
        els camins i el pas a categories. Consulteu-lo
        quan treballar amb camins, llaços o arestes
        múltiples us plantegi dubtes en els primers
        capítols.
  \item[Apèndix C: convenció de mida.]
        Explica el marc de conjunts i classes del llibre.
        És especialment útil quan apareixen
        $\cat{Cat}$, categories de functors i
        categories de prefeixos. Cal distingir
        \enquote{petita} de \enquote{localment petita}: la segona
        condició garanteix homconjunts, no que tots
        els objectes formin un conjunt.
  \item[Apèndix D: solucions dels tests.]
        Respostes i justificacions. Consulteu-les només
        després d'haver respost tot el test.
\ifpublicacioparcial\else
  \item[Apèndix E: quasiexemples i contraexemples.]
        Recull estructures que fallen exactament un axioma
        (categoria, functor, transformació natural), la
        jerarquia entre igualtat, isomorfisme, isomorfisme
        natural i equivalència, i models de la lògica que
        queden fora dels topos. La secció sobre la
        jerarquia és útil des del capítol~\ref{cap:transformacions};
        la de la lògica, després del capítol~\ref{cap:topos}.
\fi
\end{description}

La taula de notació serveix per recuperar el significat
i la primera aparició dels símbols. El glossari
català--anglès ajuda a llegir la bibliografia internacional,
sobretot en termes com \emph{pullback}, \emph{slice},
\emph{presheaf} i \emph{whiskering}. L'índex alfabètic
és una eina per retrobar conceptes quan reapareixen en
un context nou.

Les lectures recomanades al final dels capítols tenen
una funció de contrast i ampliació. En una primera
lectura, és millor consultar una font alternativa per
aclarir una pregunta concreta que obrir alhora molts
tractaments generals del mateix tema.

\section{Una manera de mesurar l'aprenentatge}

No mesureu el progrés només pel nombre de pàgines llegides.
Una fita més fiable és poder fer, sense el text davant,
quatre operacions amb una idea nova:
\begin{enumerate}
  \item enunciar-ne la definició amb les hipòtesis correctes;
  \item interpretar-la en almenys dos exemples;
  \item utilitzar-la en un argument o un càlcul senzill;
  \item explicar què falla en un contraexemple o quan
        es retira una hipòtesi.
\end{enumerate}

No cal satisfer aquestes quatre condicions amb la mateixa
fluïdesa des del primer dia. És raonable avançar amb
una dificultat ben identificada i tornar-hi després.
El que convé evitar és substituir la comprensió per
la familiaritat visual amb els símbols.

En acabar el llibre, el guany més important serà poder
relacionar les peces: veure una propietat universal
darrere d'una construcció, una naturalitat darrere
d'una compatibilitat, una adjunció darrere d'una
correspondència i una operació sobre subobjectes
darrere d'una lectura lògica. Aquesta capacitat de
passar entre exemples i estructura és el criteri
que dona unitat a tot el recorregut.

\endgroup


% ============================================================
% PART I — TEORIA
% ============================================================
\part{Teoria}
\setcounter{chapter}{0}
\chapter{Categories}
\label{cap:categories}

\begin{objectius}
\apartat{Prerequisits} Llenguatge elemental de conjunts i aplicacions: domini, codomini i composició; producte cartesià, subconjunts i relacions. Lectura de quantificadors elementals i distinció entre existència i unicitat. Nocions bàsiques d'implicació, equivalència i contraexemple. Els exemples amb grups, espais vectorials o espais topològics enriqueixen la lectura, però no són imprescindibles. No es pressuposa cap coneixement previ de teoria de categories.
\apartat{Objectius} En acabar el capítol, el lector hauria de poder:
\begin{itemize}
  \item enunciar la definició de categoria i verificar-ne els axiomes en exemples concrets;
  \item llegir i escriure diagrames commutatius i traduir-los a igualtats de composicions;
  \item reconèixer la mateixa estructura categòrica en conjunts, preordres, monoides, grafs, relacions i sistemes deductius;
  \item definir i utilitzar isomorfismes, monomorfismes, epimorfismes, seccions i retraccions, distingint-los de les nocions conjuntistes corresponents;
  \item construir i interpretar la categoria oposada, categories producte, subcategories i categories sobre i sota un objecte en casos senzills;
  \item aplicar el principi de dualitat a definicions i resultats elementals;
  \item descriure la categoria lliure generada per un graf finit senzill;
  \item distingir categories petites i localment petites.
\end{itemize}
\end{objectius}

\section{Cap al concepte de categoria}

Quan comencem a estudiar matemàtiques, aprenem a treballar amb molts tipus d'\emph{estructures}: ordres, grups, espais vectorials, espais topològics, etc. Aviat notem que aquestes estructures rarament apareixen soles. Al costat de cadascuna hi ha sempre una classe de transformacions que les relacionen respectant la seva manera de comportar-se: entre ordres, les aplicacions monòtones; entre grups, els homomorfismes; entre espais vectorials, les aplicacions lineals; entre espais topològics, les aplicacions contínues. Aquestes transformacions admissibles les anomenarem, en general, \emph{morfismes} o \emph{fletxes}.

La noció de categoria no decideix per si sola quins morfismes són admissibles: ho fixem nosaltres en cada context, d'acord amb l'estructura que volem respectar. El que fa la teoria de categories és abstreure el patró comú a tots aquests contextos: treballar amb objectes, amb transformacions admissibles entre ells i amb la manera com aquestes transformacions es componen. Cada categoria concreta enregistra aquesta tria de morfismes i, a partir d'aquí, només en reté la composició i les identitats.

Els exemples concrets que acabem d'esmentar tenen tots un conjunt subjacent. En aquests casos, preservar l'estructura significa començar per una aplicació entre els conjunts subjacents i imposar-hi les condicions adequades. No és, però, un requisit de la noció de categoria: més endavant trobarem categories els objectes de les quals no són conjunts amb estructura. Com hem dit a la secció~\ref{sec:mirar-fora} de la introducció, la teoria de categories desplaça la visió clàssica, puntual i conjuntista, de les aplicacions cap a una visió abstracta com a fletxes.

Recordem breument la notació de fletxa per a les aplicacions entre conjunts. Una aplicació $f$ amb domini $X$ i codomini $Y$ es denota $f\colon X\to Y$, o bé com a fletxa $X\xrightarrow{f}Y$. L'operació fonamental sobre les aplicacions és la \emph{composició}: si $f\colon X\to Y$ i $g\colon Y\to Z$, es defineix $g\comp f\colon X\to Z$ per
\[
(g\comp f)(x)\coloneqq g(f(x)).
\]
Perquè la composició estigui definida, el codomini de $f$ ha de coincidir amb el domini de $g$; amb fletxes, $X\xrightarrow{f}Y\xrightarrow{g}Z$. A més, per a cada conjunt $X$ hi ha l'aplicació \emph{identitat} sobre $X$,
\[
\id_X\colon X\to X,\qquad \id_X(x)\coloneqq x.
\]
Aquestes operacions es regeixen per la llei d'associativitat i les lleis de la unitat: per a $f\colon X\to Y$, $g\colon Y\to Z$ i $h\colon Z\to W$,
\[
(h\comp g)\comp f=h\comp(g\comp f),\qquad f\comp\id_X=f=\id_Y\comp f.
\]
Observem que aquestes equacions es formulen purament en termes de les operacions sobre aplicacions, sense cap referència als elements dels conjunts $X$, $Y$, $Z$, $W$.

Això no és cap ingenuïtat, sinó una evidència sobre les estructures que hem mencionat: els homomorfismes entre grups, les aplicacions monòtones entre preordres, les aplicacions lineals entre espais vectorials i les aplicacions contínues entre espais topològics satisfan les condicions que acabem d'escriure.

Aquesta és la idea que orientarà tot el que segueix: estudiar un objecte no pel que és per dins, sinó pel paper que fa dins del sistema de fletxes que el connecten amb els altres. És el que ens permetrà de parlar-ne amb un mateix llenguatge i, fins i tot, estendre'l a altres tipus d'objectes no necessàriament basats en la noció d'aplicació, com ara grafs, sistemes deductius o relacions entre conjunts.

\section{Definició de categoria}

\begin{definicio}
\label{def:categoria}
\index{categoria}Una \textbf{categoria} $\cat{C}$ consta de les dades següents:
\begin{enumerate}
\item una col·lecció d'\textbf{objectes}\index{objecte} $A,B,C,\dots$, que denotem $\Ob(\cat{C})$;
\item una col·lecció de \textbf{morfismes}\index{morfisme} $f,g,h,\dots$, que denotem $\Mor(\cat{C})$; cada morfisme $f$ té associats un únic \textbf{domini}\index{domini} $A$ i un únic \textbf{codomini}\index{codomini} $B$, i ho escrivim com $A=\dom(f)$ i $B=\cod(f)$; en aquest cas també escrivim $f\colon A\to B$ o com a fletxa $A\xrightarrow{f}B$;
\item per a cada objecte $A$, un morfisme \textbf{identitat}\index{identitat} $\id_A\colon A\to A$;
\item per a cada parella de morfismes componibles $f\colon A\to B$ i $g\colon B\to C$, la \textbf{composició}\index{composicio@composició} $g\comp f\colon A\to C$.
\end{enumerate}
Aquestes dades han de satisfer els dos axiomes següents, sempre que les composicions indicades tinguin sentit:
\begin{enumerate}
\item \textbf{Associativitat.} $h\comp(g\comp f)=(h\comp g)\comp f$, per a $f\colon A\to B$, $g\colon B\to C$, $h\colon C\to D$.
\item \textbf{Identitats.} $\id_B\comp f=f$ i $f\comp\id_A=f$, per a tot $f\colon A\to B$.
\end{enumerate}
La secció~\ref{sec:diagrames} dona la lectura diagramàtica d'aquests axiomes.
\end{definicio}

Parlem aquí de \emph{col·leccions} i no de \emph{conjunts}, deliberadament: en alguns exemples els objectes, o els morfismes, són massa nombrosos per formar un conjunt. Així, tant $\Ob(\cat{C})$ com $\Mor(\cat{C})$ són en general col·leccions, no necessàriament conjunts. Precisarem aquesta qüestió a la secció~\ref{sec:mida}. Sovint és més còmode agrupar els morfismes segons el seu domini i codomini: quan la col·lecció de morfismes de $A$ a $B$ és un conjunt, se sol denotar
\[
\Hom_{\cat{C}}(A,B)
\]
o simplement $\Hom(A,B)$ quan la categoria és clara i l'anomenem \emph{homconjunt} de $A$ i $B$. Com que cada morfisme té un únic domini i un únic codomini, $\Mor(\cat{C})$ sempre es reparteix segons el parell $(A,B)$ de domini i codomini, en parts disjuntes dos a dos;\footnote{Aquesta disjunció és, de fet, equivalent a la unicitat del domini i el codomini. A l'axiomàtica relacional de l'apèndix~\ref{cap:axiomatica-relacional} en fa el paper d'un axioma.} el que no està garantit és que cada part sigui un \emph{conjunt}. Quan ho és per a tots els parells ---propietat que rep el nom de \emph{petitesa local} (secció~\ref{sec:mida})---, aquestes parts són els homconjunts $\Hom(A,B)$ i $\Mor(\cat{C})$ n'és la unió. Al llarg del llibre totes les categories que fem servir són localment petites, de manera que aquesta subtilesa no ens afectarà.

\begin{observacio}
Fem algunes precisions sobre la definició.
\begin{enumerate}
\item Els morfismes també s'anomenen \textbf{fletxes}, i sovint farem servir les dues paraules indistintament.
\item L'axioma d'associativitat fa que puguem escriure $h\comp g\comp f$ sense parèntesis, sabent que el resultat no depèn de la manera d'agrupar-los.
\item Els axiomes d'identitat diuen que $\id_A$ es comporta com un neutre per la composició. Veurem tot seguit que aquesta propietat determina $\id_A$ de manera única.
\item Aquesta definició axiomàtica pren com a termes no definits els objectes i els morfismes. Una categoria és qualsevol estructura que contingui aquestes dues menes de dades i que satisfaci les condicions i propietats de la definició.\footnotemark
\end{enumerate}
\end{observacio}
\footnotetext{De fet, es pot prescindir d'una de les dues menes de dades: com que cada objecte queda determinat per la seva identitat, una categoria petita es pot descriure en un llenguatge amb una sola mena d'individus ---les fletxes--- i tres relacions primitives de domini, codomini i composició, amb els objectes reconstruïts com les fletxes identitat. L'apèndix~\ref{cap:axiomatica-relacional} desenvolupa aquesta presentació relacional i en demostra l'equivalència amb la definició d'aquí.}

\begin{proposicio}
La fletxa identitat d'un objecte és única.
\end{proposicio}

\begin{prova}
Suposem que $e_A\colon A\to A$ i $e'_A\colon A\to A$ són dues fletxes identitat de $A$. Com que $e_A$ és una identitat,
\[
e_A\comp e'_A=e'_A.
\]
D'altra banda, com que $e'_A$ és una identitat,
\[
e_A\comp e'_A=e_A.
\]
Per tant $e_A=e'_A$.
\end{prova}

\section{Fletxes, camins i diagrames commutatius}\label{sec:diagrames}\index{diagrama commutatiu}

La definició anterior es pot llegir algebraicament, mitjançant equacions, però també gràficament. Aquesta segona lectura serà una part essencial del llenguatge categòric.

Si $f\colon A\to B$ i $g\colon B\to C$ són componibles, la composició es pot representar amb el triangle
\[
\begin{tikzpicture}[baseline=(m.center)]
  \matrix (m) [matrix of math nodes, column sep=3.4em, row sep=2.7em] {
    A & B \\
      & C \\
  };
  \draw[-{Stealth[scale=1.1]}] (m-1-1) -- node[above] {$f$} (m-1-2);
  \draw[-{Stealth[scale=1.1]}] (m-1-2) -- node[right] {$g$} (m-2-2);
  \draw[-{Stealth[scale=1.1]}] (m-1-1) -- node[below left] {$g\comp f$} (m-2-2);
\end{tikzpicture}
\]
La fletxa directa de $A$ a $C$ és precisament el resultat de seguir el camí $A\to B\to C$.

Més generalment, un diagrama pot contenir diferents camins dirigits amb el mateix origen i el mateix destí. Per exemple,
\[
\begin{tikzpicture}[baseline=(m.center)]
  \matrix (m) [matrix of math nodes, row sep=2.7em, column sep=3.5em] {
    A & B \\
    C & D \\
  };
  \draw[-{Stealth[scale=1.1]}] (m-1-1) -- node[above] {$f$} (m-1-2);
  \draw[-{Stealth[scale=1.1]}] (m-1-1) -- node[left] {$u$} (m-2-1);
  \draw[-{Stealth[scale=1.1]}] (m-1-2) -- node[right] {$v$} (m-2-2);
  \draw[-{Stealth[scale=1.1]}] (m-2-1) -- node[below] {$g$} (m-2-2);
\end{tikzpicture}
\]
té dos camins de $A$ a $D$: primer $f$ i després $v$, o primer $u$ i després $g$.

\begin{definicio}
Un diagrama com l'anterior \textbf{commuta} si els dos camins determinen el mateix morfisme, és a dir, si
\[
v\comp f=g\comp u.
\]
En general, direm que un diagrama commuta quan tots els camins dirigits amb el mateix origen i el mateix destí que cal comparar donen la mateixa composició.
\end{definicio}

\begin{observacio}
Un dibuix de fletxes no és automàticament un diagrama commutatiu. El dibuix especifica objectes i morfismes; afirmar que \emph{commuta} afegeix una o més igualtats entre composicions de fletxes. Per això convé aprendre a passar en tots dos sentits:
\[
\text{diagrama}\quad\longleftrightarrow\quad\text{equacions entre fletxes i composicions}.
\]
\end{observacio}

L'associativitat es pot llegir com un diagrama commutatiu. En una cadena
\[
A\xrightarrow{f}B\xrightarrow{g}C\xrightarrow{h}D
\]
podem agrupar les composicions de dues maneres, $(h\comp g)\comp f$ o $h\comp(g\comp f)$. El diagrama següent ho mostra com un paral·lelogram de vèrtexs $A,B,D,C$ i costats $f$, $h\comp g$, $h$ i $g\comp f$; la diagonal interior és $g$.
\[
\begin{tikzpicture}[>={Stealth[scale=1.1]}]
  \node (A) at (-3,2) {$A$};
  \node (B) at (0,2)  {$B$};
  \node (C) at (0,0)  {$C$};
  \node (D) at (3,0)  {$D$};
  \draw[->] (A) -- node[above] {$f$} (B);
  \draw[->] (B) -- node[right] {$h\comp g$} (D);
  \draw[->] (C) -- node[below] {$h$} (D);
  \draw[->] (A) -- node[below left=-2pt] {$g\comp f$} (C);
  \draw[->] (B) -- node[left] {$g$} (C);
\end{tikzpicture}
\]
Recórrer $A\to B\to C\to D$ agrupant primer $g\comp f$ i després component amb $h$, o agrupant primer $h\comp g$ i precomponent amb $f$, dona la mateixa fletxa $A\to D$. Això és exactament l'axioma d'associativitat:
\[
(h\comp g)\comp f=h\comp(g\comp f).
\]

Les identitats també tenen una lectura diagramàtica. Si $f\colon A\to B$, inserir $\id_A$ abans de $f$ o $\id_B$ després de $f$ no canvia el camí:
\[
\begin{tikzpicture}[baseline=(m.center)]
  \matrix (m) [matrix of math nodes, row sep=2.6em, column sep=2.6em] {
    A & A \\
      & B \\
  };
  \draw[-{Stealth[scale=1.1]}] (m-1-1) -- node[above] {\(\id_A\)} (m-1-2);
  \draw[-{Stealth[scale=1.1]}] (m-1-2) -- node[right] {\(f\)} (m-2-2);
  \draw[-{Stealth[scale=1.1]}] (m-1-1) -- node[below left] {\(f\)} (m-2-2);
\end{tikzpicture}
\qquad\qquad
\begin{tikzpicture}[baseline=(m.center)]
  \matrix (m) [matrix of math nodes, row sep=2.6em, column sep=2.6em] {
    A & B \\
      & B \\
  };
  \draw[-{Stealth[scale=1.1]}] (m-1-1) -- node[above] {\(f\)} (m-1-2);
  \draw[-{Stealth[scale=1.1]}] (m-1-2) -- node[right] {\(\id_B\)} (m-2-2);
  \draw[-{Stealth[scale=1.1]}] (m-1-1) -- node[below left] {\(f\)} (m-2-2);
\end{tikzpicture}
\]
Els dos triangles commutatius expressen exactament els axiomes d'identitat:
\[
f\comp\id_A=f,\qquad \id_B\comp f=f.
\]
A partir d'ara seguirem aquesta norma: quan una propietat categòrica tingui una expressió diagramàtica natural, en donarem alhora, al mateix lloc, l'equació i el diagrama. Llegir l'una i dibuixar l'altre, i a l'inrevés, és una destresa que el llibre demana constantment.

\section{Exemples de categories}

La força de la definició es veu de seguida en la varietat d'exemples que hi encaixen. En cadascun cal identificar els objectes, els morfismes, les identitats i la composició. Comencem pels exemples més petits, que malgrat la seva simplicitat apareixeran sovint com a \emph{categories índex} (la definició~\refalt{def:diagrama}{de diagrama del capítol~5}) per construir diagrames.

\subsection{Categories finites}\label{subsec:finites}\index{categoria!buida}\index{categoria!unitat}\index{categoria!finita}

Per descomptat, els objectes d'una categoria no cal que siguin conjunts. Considerem alguns exemples molt senzills que es poden descriure per complet. En cada cas identifiquem els objectes, els morfismes, la composició i les identitats; l'associativitat i les lleis d'identitat es comproven directament, ja que hi ha molt poques composicions a considerar.
\begin{enumerate}
\item La \emph{categoria buida} $\cat{0}$ té $\Ob(\cat{0})=\varnothing$ i cap morfisme. No hi ha composició ni identitats que definir, i els axiomes es compleixen trivialment perquè no hi ha res a comprovar; és la categoria \enquote{sense res}.
\item La \emph{categoria unitat} $\cat{1}$, també anomenada categoria \emph{terminal}, té $\Ob(\cat{1})=\{\ast\}$ i $\Hom(\ast,\ast)=\{\id_\ast\}$. L'única composició és $\id_\ast\comp\id_\ast=\id_\ast$, i els axiomes es redueixen a aquesta igualtat:
\[
\begin{tikzpicture}[baseline=(a.base)]
  \node (a) {$\ast$};
  \draw[-{Stealth[scale=1.1]}] (a) to[out=40,in=-40,looseness=8] node[right] {$\id_\ast$} (a);
\end{tikzpicture}
\]
\item La categoria $\cat{2}$ té $\Ob(\cat{2})=\{\ast,\star\}$ i una única fletxa no trivial $u\colon\ast\to\star$, a més de les dues identitats:
\[
\begin{aligned}
\Hom(\ast,\ast)&=\{\id_\ast\}, &
\Hom(\star,\star)&=\{\id_\star\},\\
\Hom(\ast,\star)&=\{u\}, &
\Hom(\star,\ast)&=\varnothing.
\end{aligned}
\]
Les úniques composicions són les que involucren identitats, $u\comp\id_\ast=u=\id_\star\comp u$, de manera que els axiomes es compleixen automàticament. A partir d'ara, com és costum, no dibuixem les identitats:
\[
\begin{tikzpicture}[baseline=(a.base)]
  \node (a) at (0,0) {$\ast$};
  \node (b) at (1.8,0) {$\star$};
  \draw[-{Stealth[scale=1.1]}] (a) -- node[above] {$u$} (b);
\end{tikzpicture}
\]
\item La categoria amb \textbf{dues fletxes paral·leles} té $\Ob=\{\ast,\star\}$ i dues fletxes no trivials $f,g\colon\ast\to\star$ amb el mateix domini i el mateix codomini:
\[
\Hom(\ast,\star)=\{f,g\},\qquad
\Hom(\star,\ast)=\varnothing,
\]
i als homconjunts $\Hom(\ast,\ast)$ i $\Hom(\star,\star)$ només hi ha les identitats. Com que cap parell d'aquestes fletxes no és componible, de nou les úniques composicions són amb identitats:
\[
\begin{tikzpicture}[baseline=(a.base)]
  \node (a) at (0,0) {$\ast$};
  \node (b) at (1.8,0) {$\star$};
  \draw[-{Stealth[scale=1.1]}] ([yshift=2.2pt]a.east) -- node[above] {$f$} ([yshift=2.2pt]b.west);
  \draw[-{Stealth[scale=1.1]}] ([yshift=-2.2pt]a.east) -- node[below] {$g$} ([yshift=-2.2pt]b.west);
\end{tikzpicture}
\]
Aquesta darrera reapareixerà com a categoria índex dels \emph{equalitzadors} (la definició~\refalt{def:equalitzador}{d'equalitzador del capítol~5}).
\item La categoria $\cat{3}$ té $\Ob(\cat{3})=\{\ast,\star,\bullet\}$ i tres fletxes no trivials, $f\colon\ast\to\star$, $g\colon\star\to\bullet$ i $h\colon\ast\to\bullet$:
\[
\Hom(\ast,\star)=\{f\},\qquad
\Hom(\star,\bullet)=\{g\},\qquad
\Hom(\ast,\bullet)=\{h\};
\]
els homconjunts d'un objecte a si mateix contenen només la identitat, i els homconjunts en sentit contrari són buits. La composició obliga a incloure la composició de les dues fletxes consecutives: l'única composició no trivial és $g\comp f=h$. Amb aquesta composició (i les identitats) els axiomes se satisfan:
\[
\begin{tikzpicture}[baseline=(m.center)]
  \matrix (m) [matrix of math nodes, row sep=2.6em, column sep=2.6em] {
    \ast & \star \\
         & \bullet \\
  };
  \draw[-{Stealth[scale=1.1]}] (m-1-1) -- node[above] {$f$} (m-1-2);
  \draw[-{Stealth[scale=1.1]}] (m-1-2) -- node[right] {$g$} (m-2-2);
  \draw[-{Stealth[scale=1.1]}] (m-1-1) -- node[below left] {$h=g\comp f$} (m-2-2);
\end{tikzpicture}
\]
\end{enumerate}

\begin{observacio}
Afegir fletxes no és gratuït: la composició obliga a incloure totes les composicions. Dues fletxes de sentits oposats $f\colon\ast\to\star$ i $g\colon\star\to\ast$
\[
\begin{tikzpicture}[baseline=(a.base)]
  \node (a) at (0,0) {$\ast$};
  \node (b) at (2.2,0) {$\star$};
  \draw[-{Stealth[scale=1.1]}] ([yshift=2.6pt]a.east) -- node[above] {$f$} ([yshift=2.6pt]b.west);
  \draw[-{Stealth[scale=1.1]}] ([yshift=-2.6pt]b.west) -- node[below] {$g$} ([yshift=-2.6pt]a.east);
\end{tikzpicture}
\]
generen $g\comp f$, $f\comp g$, $g\comp f\comp g$, i així indefinidament; la categoria resultant és infinita. Per obtenir-ne una de finita cal imposar equacions entre composicions ---per exemple $g\comp f=\id_\ast$ i $f\comp g=\id_\star$, cosa que fa de $f$ un isomorfisme (la definició~\ref{def:isomorfisme}). En canvi, dues fletxes \emph{paral·leles} $f,g\colon\ast\to\star$ no són componibles entre elles, i la categoria és finita sense afegir res.
\end{observacio}

\subsection{Categories de conjunts}\label{subsec:conjunts}

La categoria $\cat{Set}$ té:
\begin{itemize}
\item objectes: $\Ob(\cat{Set})$ és la col·lecció de tots els conjunts;
\item morfismes: $\Hom_{\cat{Set}}(A,B)$ és el conjunt de totes les aplicacions $A\to B$;
\item composició: la composició habitual d'aplicacions, $(g\comp f)(x)=g(f(x))$;
\item identitats: l'aplicació identitat, $\id_A(x)=x$.
\end{itemize}
L'associativitat i les lleis d'identitat són propietats ben conegudes de les aplicacions, de manera que $\cat{Set}$ és una categoria. Serà una categoria de referència: quan un concepte categòric nou ens desconcerti, mirar què vol dir a $\cat{Set}$ sol ser un bon primer pas.

Si ens restringim als conjunts finits i a les aplicacions entre ells, obtenim la categoria $\cat{FinSet}$: $\Ob(\cat{FinSet})$ és la col·lecció dels conjunts finits, $\Hom_{\cat{FinSet}}(A,B)=\Hom_{\cat{Set}}(A,B)$, i la composició i les identitats són les mateixes que a $\cat{Set}$. Com que la composició d'aplicacions entre conjunts finits torna a ser una aplicació entre conjunts finits, els axiomes se satisfan igual que a $\cat{Set}$.

Tot conjunt $S$ es pot veure també com una categoria, que continuem escrivint $S$:\index{categoria!discreta} $\Ob(S)=S$ i els únics morfismes són les identitats, una per objecte,
\[
\Hom_S(x,y)=\begin{cases}\{\id_x\} & \text{si } x=y,\\[2pt] \varnothing & \text{altrament.}\end{cases}
\]
La composició és forçada ($\id_x\comp\id_x=\id_x$) i els axiomes es compleixen trivialment. Una categoria en què els únics morfismes són les identitats s'anomena \textbf{categoria discreta}. En una categoria discreta no hi ha cap informació més enllà de la col·lecció d'objectes.

\begin{observacio}
Una categoria no queda determinada pels seus objectes: cal dir també quines són les fletxes. Amb només dos objectes, per exemple, hem trobat ja tres categories ben diferents:
\begin{itemize}
\item la \textbf{categoria discreta} de dos punts, només amb les identitats i cap fletxa no trivial;
\item la categoria $\cat{2}$, amb una única fletxa no trivial d'un objecte a l'altre;
\item la categoria amb \textbf{dues fletxes paral·leles}, amb dues fletxes no trivials entre els mateixos objectes.
\end{itemize}
Totes tres tenen els mateixos dos objectes, però són categories diferents perquè difereixen en les fletxes. En particular, $\cat{2}$ \emph{no} és la categoria discreta de dos punts: la presència de la fletxa no trivial les distingeix.
\end{observacio}

\subsection{Estructures elementals com a categories}\label{subsec:elementals}

Les estructures matemàtiques més simples sovint són categories. Per exemple, un monoide no s'assembla a una categoria d'un sol objecte: és exactament una categoria d'un sol objecte. Recollim aquesta idea en una llista d'identificacions, i tot seguit les justifiquem:
\begin{itemize}
\item un conjunt és una categoria discreta (ja vist a la subsecció anterior);
\item un monoide és una categoria d'un sol objecte;
\item un grup és una categoria d'un sol objecte en què tot morfisme és un isomorfisme;
\item un preordre és una categoria prima;
\item un ordre parcial és una categoria prima en què els únics isomorfismes són les identitats;
\item un sistema deductiu és una categoria prima.
\end{itemize}
En tots els casos, els objectes i els morfismes deixen de ser conjunts i aplicacions: passen a ser els elements i les relacions internes de l'estructura.

\paragraph{Preordres.}\index{preordre}
Un \emph{preordre} és un parell $(P,\le)$ on $P$ és un conjunt i $\le$ una relació binària sobre $P$ reflexiva ($x\le x$) i transitiva (si $x\le y$ i $y\le z$, aleshores $x\le z$). Per exemple, $(\mathbb{R},\le)$, o els subconjunts d'un conjunt ordenats per la inclusió $\subseteq$.

Tot preordre $(P,\le)$ \textbf{és} una categoria, que continuem escrivint $P$:
\begin{itemize}
\item objectes: $\Ob(P)=P$, els elements del preordre;
\item morfismes: hi ha exactament una fletxa $x\to y$ quan $x\le y$,
\[
\Hom_P(x,y)=\begin{cases}\{(x,y)\} & \text{si } x\le y,\\[2pt] \varnothing & \text{altrament;}\end{cases}
\]
\item composició: $(y,z)\comp(x,y)=(x,z)$, que existeix per transitivitat;
\item identitats: $\id_x=(x,x)$, que existeix per reflexivitat.
\end{itemize}
\[
\begin{tikzpicture}[baseline=(m.center)]
  \matrix (m) [matrix of math nodes, row sep=2.6em, column sep=2.6em] {
    x & y \\
      & z \\
  };
  \draw[-{Stealth[scale=1.1]}] (m-1-1) -- node[above] {$(x,y)$} (m-1-2);
  \draw[-{Stealth[scale=1.1]}] (m-1-2) -- node[right] {$(y,z)$} (m-2-2);
  \draw[-{Stealth[scale=1.1]}] (m-1-1) -- node[below left] {$(x,z)$} (m-2-2);
\end{tikzpicture}
\]
Com que cada homconjunt té com a màxim un element, l'associativitat i les lleis d'identitat se satisfan automàticament. Una categoria en què $|\Hom(x,y)|\le 1$ per a tot parell d'objectes s'anomena \textbf{prima}\index{categoria!prima} (en anglès, \emph{thin category}).

Si el preordre és a més antisimètric ($x\le y$ i $y\le x$ impliquen $x=y$), és un \emph{ordre parcial}. Categòricament, l'antisimetria diu que els únics isomorfismes (la definició~\ref{def:isomorfisme}) són les identitats: si $x\to y$ i $y\to x$ existeixen totes dues, aleshores $x=y$.

\paragraph{Monoides.}\index{monoide}
Un \emph{monoide} és una terna $(M,\ast,e)$ on $\ast$ és una operació associativa sobre $M$ amb element neutre $e$. Per exemple, $(\mathbb{N},+,0)$, o les paraules sobre un alfabet amb la concatenació i la paraula buida.

Tot monoide \textbf{és} una categoria amb un sol objecte $\ast$, que escrivim $\cat{B}M$.\footnote{La notació $\cat{B}M$ prové de la topologia, on $BM$ designa l'\emph{espai classificador} del monoide (o grup) $M$; la categoria $\cat{B}M$ n'és el model categòric.} Concretament:
\begin{itemize}
\item objectes: $\Ob(\cat{B}M)=\{\ast\}$;
\item morfismes: l'únic homconjunt és tot el monoide, $\Hom_{\cat{B}M}(\ast,\ast)=M$;
\item composició: l'operació del monoide, $g\comp f=g\ast f$, on l'element de la dreta actua primer;
\item identitat: $\id_\ast=e$.
\end{itemize}
Els axiomes de categoria són exactament els de monoide. Els morfismes, aquí, no són aplicacions: són els elements de $M$.

\paragraph{Grups.}\index{grup}
Un \emph{grup} és una terna $(G,\ast,e)$ que és un monoide en què, a més, tot element $a$ té un invers $a^{-1}$ amb $a\ast a^{-1}=e=a^{-1}\ast a$. Per exemple, $(\mathbb{Z},+,0)$, o les permutacions d'un conjunt amb la composició.

Un grup \textbf{és} una categoria d'un sol objecte en què tot morfisme és un isomorfisme. Com que un grup és en particular un monoide, la categoria és la de l'apartat anterior, que escrivim igualment $\cat{B}G$: $\Ob(\cat{B}G)=\{\ast\}$, $\Hom_{\cat{B}G}(\ast,\ast)=G$, $g\comp f=g\ast f$ i $\id_\ast=e$. En efecte, l'invers de l'element $a$ és exactament el morfisme invers de la fletxa corresponent. La teoria de grups és, doncs, la part invertible de la teoria de monoides. Cada element del grup és un bucle sobre l'únic objecte, i la identitat $e$ n'és el bucle destacat:
\[
\begin{tikzpicture}[baseline=(o.base)]
  \node (o) {$\ast$};
  % pètals genèrics (elements del grup)
  \foreach \ang in {30,150,210,270,330}{
    \draw[-{Stealth[scale=0.9]},gray] (o) to[out=\ang+14,in=\ang-14,looseness=16] (o);
  }
  % pètal destacat: la identitat e
  \draw[-{Stealth[scale=1.0]},thick] (o) to[out=104,in=76,looseness=16] (o);
  \node at (0,1.32) {$e$};
\end{tikzpicture}
\]

\paragraph{Sistemes deductius.}\index{deduibilitat@deduïbilitat}\index{Prov@$\cat{Prov}_T$}
Fixem un \emph{sistema deductiu} $T$ i considerem la seva relació de deduïbilitat $\vdash_T$ entre fórmules, reflexiva ($p\vdash_T p$) i transitiva (de $p\vdash_T q$ i $q\vdash_T r$ se segueix $p\vdash_T r$). Aquesta relació determina una categoria, que escrivim $\cat{Prov}_T$ per distingir-la del sistema $T$ del qual prové:
\begin{itemize}
\item objectes: $\Ob(\cat{Prov}_T)$ és el conjunt de les fórmules $p,q,r,\dots$ de $T$;
\item morfismes: hi ha exactament una fletxa $p\to q$ quan $p\vdash_T q$,
\[
\Hom_{\cat{Prov}_T}(p,q)=\begin{cases}\{(p,q)\} & \text{si } p\vdash_T q,\\[2pt] \varnothing & \text{altrament;}\end{cases}
\]
\item composició: $(q,r)\comp(p,q)=(p,r)$, que existeix per transitivitat;
\item identitats: $\id_p=(p,p)$, que existeix per reflexivitat.
\end{itemize}
\[
\begin{tikzpicture}[baseline=(m.center)]
  \matrix (m) [matrix of math nodes, row sep=2.6em, column sep=2.6em] {
    p & q \\
      & r \\
  };
  \draw[-{Stealth[scale=1.1]}] (m-1-1) -- node[above] {$p\vdash_T q$} (m-1-2);
  \draw[-{Stealth[scale=1.1]}] (m-1-2) -- node[right] {$q\vdash_T r$} (m-2-2);
  \draw[-{Stealth[scale=1.1]}] (m-1-1) -- node[below left] {$p\vdash_T r$} (m-2-2);
\end{tikzpicture}
\]
Com el preordre, $\cat{Prov}_T$ és una categoria prima. A diferència de l'ordre parcial, però, la deduïbilitat no acostuma a ser antisimètrica: dues fórmules diferents poden ser interdeduïbles ($p\vdash_T q$ i $q\vdash_T p$ amb $p\neq q$), com $p\land q$ i $q\land p$. Categòricament, això vol dir que hi ha isomorfismes no trivials: fórmules distintes però isomorfes. Aquesta categoria només registra \emph{si} $q$ es dedueix de $p$, no \emph{de quina manera} es dedueix; l'observació següent presenta la versió que distingeix les proves.

\begin{observacio}[Les proves mateixes com a fletxes]\index{prova!com a morfisme}
Hi ha una interpretació lògica més fina en què una prova concreta
\[
\pi\colon p\longrightarrow q
\]
és ella mateixa un morfisme, i dues proves $\pi,\rho\colon p\to q$ poden ser morfismes diferents. La composició representa aleshores l'encadenament de proves i la identitat, la prova trivial d'una fórmula a partir d'ella mateixa. Passem, doncs, d'una sola fletxa (registrar \emph{si} hi ha deducció) a moltes fletxes (registrar \emph{quina} prova).

Per obtenir una categoria de proves completament definida cal fixar el sistema deductiu i precisar també quan dues derivacions es consideren iguals ---per exemple, mitjançant una equivalència adequada entre proves. No necessitem encara aquesta teoria; de moment distingirem entre la categoria prima $\cat{Prov}_T$ de la \emph{deduïbilitat} i aquesta interpretació més rica de les \emph{proves com a fletxes}. Hi tornarem quan el llenguatge categòric ens permeti fer-ho amb més precisió.
\end{observacio}

\subsection{Categories de conjunts estructurats}

Ara considerem les categories els objectes de les quals són conjunts amb alguna estructura, i els morfismes, les aplicacions que preserven aquesta estructura. En tots els casos la composició és la d'aplicacions i les identitats són les aplicacions identitat, de manera que l'associativitat i les lleis d'identitat s'hereten de les de $\cat{Set}$; només cal comprovar, en cada cas, que la composició d'aplicacions que respecten l'estructura torna a respectar-la, i que la identitat també ho fa.

\paragraph{Grups.}
La categoria $\cat{Grp}$ té:
\begin{itemize}
\item objectes: $\Ob(\cat{Grp})$ és la col·lecció de tots els grups;
\item morfismes: $\Hom_{\cat{Grp}}(G,H)$ és el conjunt dels homomorfismes de grups $G\to H$;
\item composició: la composició d'aplicacions;
\item identitats: l'aplicació identitat $\id_{G}$.
\end{itemize}
Un \emph{homomorfisme de grups} $h\colon G\to H$ és una aplicació entre els conjunts subjacents que preserva l'operació: $h(a\ast b)=h(a)\ast h(b)$ per a tot $a,b$ (d'on se segueix que preserva també el neutre i els inversos). Si $h\colon G\to H$ i $k\colon H\to K$ són homomorfismes, la composició també ho és, ja que
\[
k(h(a\ast b))=k(h(a)\ast h(b))=k(h(a))\ast k(h(b));
\]
i l'aplicació identitat preserva trivialment l'operació. Com que la composició i les identitats són les d'aplicacions entre conjunts, l'associativitat i les lleis d'identitat s'hereten de $\cat{Set}$.

\paragraph{Espais vectorials.}\index{espai vectorial}
Fixem un cos $\mathbb{K}$. La categoria $\cat{Vect}_{\mathbb K}$ té:
\begin{itemize}
\item objectes: $\Ob(\cat{Vect}_{\mathbb K})$ és la col·lecció dels espais vectorials sobre $\mathbb{K}$;
\item morfismes: $\Hom_{\mathbb K}(V,W)$ és el conjunt de les aplicacions lineals $V\to W$;
\item composició: la composició d'aplicacions;
\item identitats: l'aplicació identitat $\id_V$.
\end{itemize}
Un \emph{espai vectorial} sobre $\mathbb{K}$ és un conjunt $V$ amb una suma commutativa i associativa $V\times V\to V$ ---amb element neutre i oposats--- i una multiplicació per escalars $\mathbb{K}\times V\to V$ compatible amb les operacions. Una \emph{aplicació lineal} $f\colon V\to W$ és una aplicació entre els conjunts subjacents que preserva l'estructura: $f(u+v)=f(u)+f(v)$ i $f(\lambda v)=\lambda f(v)$. Si $f\colon V\to W$ i $g\colon W\to U$ són lineals, la composició també ho és, ja que
\[
g(f(u+v))=g(f(u))+g(f(v))\qquad\text{i}\qquad g(f(\lambda v))=\lambda\,g(f(v));
\]
i la identitat és lineal. Per tant la composició i les identitats són les d'aplicacions entre conjunts, i l'associativitat i les lleis d'identitat s'hereten de $\cat{Set}$.

\paragraph{Espais topològics.}\index{espai topològic}
La categoria $\cat{Top}$ té:
\begin{itemize}
\item objectes: $\Ob(\cat{Top})$ és la col·lecció de tots els espais topològics;
\item morfismes: $\Hom_{\cat{Top}}(X,Y)$ és el conjunt de les aplicacions contínues $X\to Y$;
\item composició: la composició d'aplicacions;
\item identitats: l'aplicació identitat $\id_{X}$.
\end{itemize}
Un \emph{espai topològic} és un parell $(X,\mathcal{T}_X)$ on $X$ és un conjunt i $\mathcal{T}_X$ una família de subconjunts de $X$ (els \emph{oberts}) tal que:
\begin{itemize}
\item $\varnothing, X\in\mathcal{T}_X$;
\item si $U,V\in\mathcal{T}_X$, aleshores $U\cap V\in\mathcal{T}_X$;
\item si $\{U_i\}_{i\in I}$ és una família qualsevol en $\mathcal{T}_X$, aleshores $\bigcup_{i\in I}U_i\in\mathcal{T}_X$.
\end{itemize}
Una \emph{aplicació contínua} $f\colon(X,\mathcal{T}_X)\to(Y,\mathcal{T}_Y)$ és una aplicació $f\colon X\to Y$ que preserva l'estructura en el sentit següent: $f^{-1}(U)\in\mathcal{T}_X$ per a tot $U\in\mathcal{T}_Y$. Si $f\colon X\to Y$ i $g\colon Y\to Z$ són contínues, la composició també ho és, ja que $(g\comp f)^{-1}(U)=f^{-1}(g^{-1}(U))$ és obert quan $U$ ho és; i la identitat és contínua, perquè $\id^{-1}(U)=U$. Per tant, un cop més, la composició i les identitats són les d'aplicacions entre conjunts, i l'associativitat i les lleis d'identitat s'hereten de $\cat{Set}$.

\paragraph{Monoides.}
La categoria $\cat{Mon}$ té:
\begin{itemize}
\item objectes: $\Ob(\cat{Mon})$ és la col·lecció de tots els monoides;
\item morfismes: $\Hom_{\cat{Mon}}(M,N)$ és el conjunt dels homomorfismes de monoides $M\to N$;
\item composició: la composició d'aplicacions;
\item identitats: l'aplicació identitat $\id_{M}$.
\end{itemize}
Un \emph{homomorfisme de monoides} $h\colon M\to N$ és una aplicació entre els conjunts subjacents que preserva l'operació i el neutre: $h(a\ast b)=h(a)\ast h(b)$ i $h(e_M)=e_N$. Si $h\colon M\to N$ i $k\colon N\to P$ són homomorfismes, la composició també ho és, ja que
\[
k(h(a\ast b))=k(h(a)\ast h(b))=k(h(a))\ast k(h(b))\qquad\text{i}\qquad k(h(e_M))=k(e_N)=e_P;
\]
i la identitat preserva trivialment l'operació i el neutre. Com que la composició i les identitats són les d'aplicacions entre conjunts, l'associativitat i les lleis d'identitat s'hereten de $\cat{Set}$.

\paragraph{Ordres parcials.}
La categoria $\cat{Pos}$ té:
\begin{itemize}
\item objectes: $\Ob(\cat{Pos})$ és la col·lecció de tots els ordres parcials;
\item morfismes: $\Hom_{\cat{Pos}}(P,Q)$ és el conjunt de les aplicacions monòtones $P\to Q$;
\item composició: la composició d'aplicacions;
\item identitats: l'aplicació identitat $\id_{P}$.
\end{itemize}
Una \emph{aplicació monòtona} $f\colon P\to Q$ és una aplicació entre els conjunts subjacents que preserva l'ordre: si $x\le y$, aleshores $f(x)\le f(y)$. Si $f\colon P\to Q$ i $g\colon Q\to R$ són monòtones, la composició també ho és, ja que
\[
x\le y \ \Longrightarrow\ f(x)\le f(y) \ \Longrightarrow\ g(f(x))\le g(f(y));
\]
i la identitat preserva trivialment l'ordre. Un cop més, la composició i les identitats són les d'aplicacions entre conjunts, i les lleis s'hereten de $\cat{Set}$.

Convé no confondre les dues mirades: un monoide concret \emph{és} una categoria (d'un sol objecte) i un ordre parcial concret \emph{és} una categoria (prima), mentre que $\cat{Mon}$ i $\cat{Pos}$ tenen \emph{tots} els monoides i \emph{tots} els ordres parcials per objectes, respectivament.

\begin{observacio}[Categories com a contextos matemàtics]
Cadascuna d'aquestes categories és l'\emph{àmbit} on es desenvolupa una teoria matemàtica: $\cat{Top}$ és el context de la topologia general, $\cat{Grp}$ el de la teoria de grups, $\cat{Vect}_{\mathbb K}$ el de l'àlgebra lineal, i així successivament. Aquí, la categoria la formen tots els objectes d'una mena; a les estructures elementals, en canvi, un sol objecte estructurat ja és una categoria.
\end{observacio}

\subsection{Conjunts i relacions}\index{Rel@$\cat{Rel}$}

Com hem avançat a la introducció (secció~\ref{sec:mirar-fora}), els conjunts admeten una segona categoria, on els morfismes no són aplicacions sinó \emph{relacions}. La categoria $\cat{Rel}$ té:
\begin{itemize}
\item objectes: $\Ob(\cat{Rel})$ és la col·lecció de tots els conjunts;
\item morfismes: $\Hom_{\cat{Rel}}(X,Y)$ és el conjunt de les relacions $R\subseteq X\times Y$ (amb la precisió sobre els extrems de l'observació~\ref{obs:extrems-implicits});
\item composició: per a $R\subseteq X\times Y$ i $S\subseteq Y\times Z$,
\[
S\comp R=\{(x,z)\mid \exists y\in Y.\ (x,y)\in R \ \text{i}\ (y,z)\in S\}\subseteq X\times Z;
\]
\item identitats: la relació diagonal $\id_X=\{(x,x)\mid x\in X\}$.
\end{itemize}
Els axiomes es comproven directament a partir de la definició. Per a l'associativitat, un parell $(x,w)$ pertany a les dues maneres d'agrupar tres relacions encadenades exactament quan hi ha termes intermedis que connecten $x$ amb $w$ passant per les tres relacions; com que la condició no depèn de l'agrupament, les dues composicions coincideixen. Per a les identitats, compondre amb la diagonal no afegeix cap pas efectiu, perquè obliga el terme intermedi a coincidir amb l'extrem. La verificació completa és el primer exercici resolt del capítol~1 de la part de Pràctica (exercici~\ref{exr:pr-rel}). Compondre relacions equival a aplicar un quantificador existencial sobre el terme intermedi: aquest és el primer pont amb la lògica que assenyalàvem a la introducció. A diferència del que passa a $\cat{Set}$, un morfisme $X\to Y$ de $\cat{Rel}$ no ha d'assignar necessàriament a cada element de $X$ un únic element de $Y$: pot relacionar-lo amb diversos elements de $Y$, o amb cap.

\begin{observacio}[Extrems implícits]\label{obs:extrems-implicits}
En rigor, un morfisme $X\to Y$ de $\cat{Rel}$ no és una relació $R$ tota sola, sinó la relació \emph{junt amb} els seus extrems: la terna $(X,Y,R)$ amb $R\subseteq X\times Y$. Aquesta precisió és necessària perquè els homconjunts siguin disjunts, tal com exigeix la definició de categoria. Per exemple, la relació buida $R=\varnothing$ està continguda en $X\times Y$ per a qualsevol parell $(X,Y)$; si no registréssim els extrems, un mateix objecte $\varnothing$ seria alhora morfisme entre infinites parelles d'objectes, amb domini i codomini indeterminats. Registrant $(X,Y,R)$, cada morfisme té un únic domini $X$ i un únic codomini $Y$. Escriurem simplement $R$ quan els extrems se sobreentenguin; és la mateixa convenció d'\emph{extrems implícits} que ja hem fet servir en codificar la fletxa d'un preordre o de $\cat{Prov}_T$ com el parell $(x,y)$, i que retrobarem tot seguit a les categories sobre i sota un objecte. No canvia cap de les construccions: només n'explicita la compatibilitat amb la definició inicial de categoria.
\end{observacio}

\subsection{Categories de grafs}\label{subsec:grafs}

\begin{definicio}\index{graf dirigit}
Un \textbf{graf dirigit} $G$ consta d'un conjunt $V_G$ de \textbf{vèrtexs}, un conjunt $E_G$ d'\textbf{arestes} i dues aplicacions
\[
s_G,t_G\colon E_G\longrightarrow V_G,
\]
que assignen a cada aresta el seu \textbf{origen} i el seu \textbf{destí}, respectivament.
\end{definicio}

\begin{definicio}\index{morfisme de grafs}
Un \textbf{morfisme de grafs} $F\colon G\to H$ és una parella d'aplicacions
\[
F_V\colon V_G\to V_H,
\qquad
F_E\colon E_G\to E_H,
\]
que preserven origen i destí:
\[
F_V\comp s_G=s_H\comp F_E,
\qquad
F_V\comp t_G=t_H\comp F_E.
\]
Equivalentment, han de commutar els dos quadrats següents:

\begin{center}
\begin{tikzpicture}[baseline=(m.center)]
  \matrix (m) [matrix of math nodes, row sep=2.5em, column sep=3.2em] {
    E_G & V_G \\
    E_H & V_H \\
  };
  \draw[-{Stealth[scale=1.0]}] (m-1-1) -- node[above] {$s_G$} (m-1-2);
  \draw[-{Stealth[scale=1.0]}] (m-1-1) -- node[left] {$F_E$} (m-2-1);
  \draw[-{Stealth[scale=1.0]}] (m-1-2) -- node[right] {$F_V$} (m-2-2);
  \draw[-{Stealth[scale=1.0]}] (m-2-1) -- node[below] {$s_H$} (m-2-2);
\end{tikzpicture}
\hspace{2.2em}
\begin{tikzpicture}[baseline=(n.center)]
  \matrix (n) [matrix of math nodes, row sep=2.5em, column sep=3.2em] {
    E_G & V_G \\
    E_H & V_H \\
  };
  \draw[-{Stealth[scale=1.0]}] (n-1-1) -- node[above] {$t_G$} (n-1-2);
  \draw[-{Stealth[scale=1.0]}] (n-1-1) -- node[left] {$F_E$} (n-2-1);
  \draw[-{Stealth[scale=1.0]}] (n-1-2) -- node[right] {$F_V$} (n-2-2);
  \draw[-{Stealth[scale=1.0]}] (n-2-1) -- node[below] {$t_H$} (n-2-2);
\end{tikzpicture}
\end{center}
\end{definicio}

Els grafs dirigits i els morfismes de grafs formen una categoria, que denotarem $\cat{Graph}$:\index{Graph@$\cat{Graph}$}
\begin{itemize}
\item objectes: $\Ob(\cat{Graph})$ és la col·lecció de tots els grafs dirigits;
\item morfismes: $\Hom_{\cat{Graph}}(G,H)$ és el conjunt dels morfismes de grafs $G\to H$;
\item composició: component a component; per a $F\colon G\to H$ i $F'\colon H\to K$,
\[
F'\comp F=(F'_V\comp F_V,\ F'_E\comp F_E);
\]
\item identitats: $\id_G=(\id_{V_G},\id_{E_G})$.
\end{itemize}
La composició torna a ser un morfisme de grafs, perquè les equacions es conserven:
\[
(F'_V\comp F_V)\comp s_G=F'_V\comp s_H\comp F_E=s_K\comp(F'_E\comp F_E),
\]
i anàlogament per a $t$. L'associativitat i les lleis d'identitat es compleixen component a component, perquè són les de $\cat{Set}$.

Aquest exemple és especialment rellevant perquè els grafs són l'esquelet combinatori dels diagrames: tenen vèrtexs i arestes dirigides, però encara no incorporen per si sols identitats, composicions ni axiomes categòrics.

Per fixar idees, vet aquí un graf dirigit concret amb tres vèrtexs $V_G=\{a,b,c\}$ i tres arestes:
\[
\begin{tikzpicture}[baseline=(m.center)]
  \matrix (m) [matrix of math nodes, row sep=2.6em, column sep=2.6em] {
    a & b \\
      & c \\
  };
  \draw[-{Stealth[scale=1.1]}] (m-1-1) -- node[above] {$e_1$} (m-1-2);
  \draw[-{Stealth[scale=1.1]}] (m-1-2) -- node[right] {$e_2$} (m-2-2);
  \draw[-{Stealth[scale=1.1]}] (m-1-1) -- node[below left] {$e_3$} (m-2-2);
\end{tikzpicture}
\]
Aquí $e_1$ va de $a$ a $b$, $e_2$ de $b$ a $c$ i $e_3$ de $a$ a $c$. A diferència d'una categoria, un graf no demana que hi hagi cap aresta que \enquote{representi la composició} d'$e_1$ i $e_2$: $e_3$ hi és perquè l'hem posada, no perquè cap axioma l'obligui. Aquesta és exactament la diferència entre un graf i una categoria: la categoria afegeix identitats, composició i els axiomes que les regeixen. Per a un desenvolupament complet dels grafs, amb nombrosos exemples i els homomorfismes de grafs, vegeu l'apèndix~\ref{cap:grafs}.

\subsection{Resum: objectes i morfismes dels exemples}

La taula següent recull, per a cada exemple d'aquesta secció, quines dades constitueixen els objectes i quines els morfismes.

\begin{center}
\footnotesize
\setlength{\tabcolsep}{6pt}
\begin{tabularx}{\linewidth}{@{}>{\raggedright\arraybackslash}X>{\raggedright\arraybackslash}X>{\raggedright\arraybackslash}X@{}}
\toprule
Categoria & Objectes & Morfismes \\
\midrule
$\cat{0},\cat{1},\cat{2},\cat{3}$ & marcadors abstractes & fletxes formals \\
$\cat{Set}$, $\cat{FinSet}$ & conjunts & aplicacions \\
categoria discreta & marcadors abstractes & només identitats \\
$\cat{Grp}$, $\cat{Vect}_{\mathbb K}$, $\cat{Top}$, $\cat{Mon}$, $\cat{Pos}$ & conjunts amb estructura & aplicacions que preserven l'estructura \\
$\cat{B}M$ (un monoide) & un únic objecte $\ast$ & elements de $M$ \\
un preordre $P$; $\cat{Prov}_T$ & elements de $P$; fórmules de $T$ & $x\le y$; $p\vdash_T q$ \\
$\cat{Rel}$ & conjunts & relacions \\
$\cat{Graph}$ & grafs $(V,E,s,t)$ & parelles compatibles d'aplicacions \\
\bottomrule
\end{tabularx}
\end{center}

Aquesta taula fa visible que la definició de categoria no pressuposa res sobre la naturalesa dels objectes ni dels morfismes: no cal que els objectes siguin conjunts, ni que els morfismes siguin aplicacions. En una categoria discreta, per exemple, les identitats són morfismes \emph{formals}: presentar-les com a aplicacions pressuposaria una realització concreta que no forma part de la definició. I a $\cat{B}M$ els morfismes són els elements de $M$, que en general tampoc no són aplicacions. La categoria $\cat{Rel}$ ho remarca des de l'altre costat: els seus objectes són conjunts, però els seus morfismes ---les relacions--- no són funcions.

Convé no confondre aquesta descripció \emph{estructural} de les dades amb la seva eventual \emph{codificació} conjuntista. En el marc de fonaments que adoptem, qualsevol d'aquestes dades ---inclosos els marcadors abstractes--- es pot representar per conjunts; però aquesta codificació és una tria externa, no una part de la definició de categoria. Per això la taula descriu què \emph{són} les dades i no si \enquote{són conjunts} o \enquote{són funcions}: aquesta darrera pregunta barreja la naturalesa estructural amb una representació particular.

\section{Isomorfismes}\index{isomorfisme}

Entre tots els morfismes d'una categoria, n'hi ha uns d'especialment importants: els que tenen invers. Capturen, en el llenguatge abstracte de les fletxes, la idea que dos objectes tenen la mateixa estructura des del punt de vista de la categoria.

\begin{definicio}\label{def:isomorfisme}
En una categoria $\cat{C}$, un morfisme $f\colon A\to B$ és un \textbf{isomorfisme} si existeix un morfisme $g\colon B\to A$ tal que
\[
g\comp f=\id_A
\qquad\text{i}\qquad
f\comp g=\id_B.
\]
El morfisme $g$ s'anomena \textbf{invers} de $f$. Si existeix un isomorfisme entre $A$ i $B$, diem que $A$ i $B$ són \textbf{isomorfs}, i escrivim $A\cong B$.
\end{definicio}

La situació es caracteritza per dues equacions, $g\comp f=\id_A$ i $f\comp g=\id_B$, que es llegeixen com dos triangles commutatius:
\[
\begin{tikzpicture}[baseline=(m.center)]
  \matrix (m) [matrix of math nodes, row sep=2.6em, column sep=2.6em] {
    A & B \\
    A & \\
  };
  \draw[-{Stealth[scale=1.1]}] (m-1-1) -- node[above] {$f$} (m-1-2);
  \draw[-{Stealth[scale=1.1]}] (m-1-2) -- node[right] {$g$} (m-2-1);
  \draw[-{Stealth[scale=1.1]}] (m-1-1) -- node[left] {$\id_A$} (m-2-1);
\end{tikzpicture}
\qquad\qquad
\begin{tikzpicture}[baseline=(m.center)]
  \matrix (m) [matrix of math nodes, row sep=2.6em, column sep=2.6em] {
    B & A \\
    B & \\
  };
  \draw[-{Stealth[scale=1.1]}] (m-1-1) -- node[above] {$g$} (m-1-2);
  \draw[-{Stealth[scale=1.1]}] (m-1-2) -- node[right] {$f$} (m-2-1);
  \draw[-{Stealth[scale=1.1]}] (m-1-1) -- node[left] {$\id_B$} (m-2-1);
\end{tikzpicture}
\]
Cada triangle diu que anar d'un objecte a l'altre amb $f$ o $g$ i tornar és el mateix que seguir la identitat.

\begin{proposicio}
L'invers d'un isomorfisme, si existeix, és únic.
\end{proposicio}

\begin{prova}
Suposem que $f\colon A\to B$ té dos inversos, $g$ i $g'$. Aleshores
\[
g=g\comp\id_B=g\comp(f\comp g')=(g\comp f)\comp g'=\id_A\comp g'=g'.
\]
Per tant $g=g'$. Com que l'invers és únic, el podem denotar sense ambigüitat $f^{-1}$.
\end{prova}

\begin{proposicio}
Sigui $\cat{C}$ una categoria.
\begin{enumerate}
\item Per a cada objecte $A$, la identitat $\id_A$ és un isomorfisme, amb $\id_A^{-1}=\id_A$.
\item Si $f$ és un isomorfisme, aleshores $f^{-1}$ també ho és, i $(f^{-1})^{-1}=f$.
\item Si $f\colon A\to B$ i $g\colon B\to C$ són isomorfismes, aleshores $g\comp f$ també ho és, i
\[
(g\comp f)^{-1}=f^{-1}\comp g^{-1}.
\]
\end{enumerate}
\end{proposicio}

\begin{prova}
(1) De $\id_A\comp\id_A=\id_A$ se segueix que $\id_A$ és el seu propi invers. (2) Les equacions $g\comp f=\id_A$ i $f\comp g=\id_B$ que fan de $g=f^{-1}$ l'invers de $f$ diuen també que $f$ és l'invers de $g$; per tant $f^{-1}$ és un isomorfisme i $(f^{-1})^{-1}=f$. (3) Comprovem que $f^{-1}\comp g^{-1}$ és l'invers de $g\comp f$:
\[
(g\comp f)\comp(f^{-1}\comp g^{-1})=g\comp(f\comp f^{-1})\comp g^{-1}=g\comp\id_B\comp g^{-1}=g\comp g^{-1}=\id_C,
\]
i anàlogament $(f^{-1}\comp g^{-1})\comp(g\comp f)=\id_A$.
\end{prova}

\begin{corollari}
Donada una categoria $\cat{C}$, definim sobre $\Ob(\cat{C})$ la relació $\cong_{\cat{C}}$ per: $A\cong_{\cat{C}}B$ si i només si existeix un isomorfisme $f\colon A\to B$. Aleshores $\cong_{\cat{C}}$ és una relació d'equivalència.
\end{corollari}

\begin{prova}
És reflexiva: $\id_A$ és un isomorfisme per (1), de manera que $A\cong_{\cat{C}}A$. És simètrica: si $f$ és un isomorfisme de $A$ a $B$, aleshores $f^{-1}$ ho és de $B$ a $A$ per (2), de manera que $B\cong_{\cat{C}}A$. És transitiva: si $f\colon A\to B$ i $g\colon B\to C$ són isomorfismes, aleshores $g\comp f$ també ho és per (3), de manera que $A\cong_{\cat{C}}C$.
\end{prova}

\begin{observacio}[Igualtat i isomorfisme]\label{obs:igualtat-isomorfisme}\index{isomorfisme!i igualtat}
Convé fixar des d'ara una convenció que el llibre manté sempre. Escrivim $A=B$ només quan $A$ i $B$ són \emph{el mateix} objecte, i $A\cong B$ quan hi ha un isomorfisme entre ells. Quan un isomorfisme és distingit ---no depèn de cap tria---, l'anomenem \textbf{canònic}, però continua sent un isomorfisme i no una igualtat (en veurem un primer cas en el monoide lliure, exemple~\ref{ex:monoide-lliure-cat}).

La relació $\cong_{\cat{C}}$ permet parlar d'una \emph{igualtat estructural relativa a $\cat{C}$}: dos objectes isomorfs a $\cat{C}$ no es poden distingir per cap propietat formulada només amb els morfismes, la composició i les identitats de $\cat{C}$, perquè l'isomorfisme la transporta d'un a l'altre (ho precisarem a l'observació~\ref{obs:invariancia}). Diem, doncs, que tenen la mateixa estructura \emph{des del punt de vista de $\cat{C}$}, i per això sovint la pregunta rellevant no és si $A$ i $B$ són iguals, sinó si són isomorfs a $\cat{C}$.

Aquesta igualtat no és absoluta: depèn dels morfismes que la categoria considera admissibles. Per exemple, el conjunt $\{0,1\}$ amb l'ordre discret i el mateix conjunt amb l'ordre $0<1$ no són isomorfs a $\cat{Pos}$, perquè cap bijecció monòtona no té inversa monòtona, però els seus conjunts subjacents sí que són isomorfs a $\cat{Set}$. Canviar de categoria canvia, doncs, quins objectes compten com a estructuralment iguals. Aquesta és una de les idees centrals de la teoria de categories: la forma d'un objecte sempre es mira dins d'una categoria determinada.
La jerarquia completa ---igualtat, isomorfisme, isomorfisme natural i equivalència---, amb un contraexemple per a cada pas, es tracta a l'apèndix~\refalt{cap:contraexemples}{de contraexemples del llibre complet}.
\end{observacio}

\begin{exemple}
A $\cat{Set}$, els isomorfismes són exactament les aplicacions bijectives. En efecte, una aplicació té inversa per la composició si i només si és injectiva i exhaustiva.
\end{exemple}

\begin{exemple}[Isomorfismes en les categories de conjunts estructurats]
A $\cat{Grp}$, $\cat{Vect}_{\mathbb K}$ i $\cat{Mon}$, els isomorfismes són les aplicacions bijectives que preserven l'estructura i la inversa de les quals també la preserva; recuperem així les nocions habituals d'isomorfisme de grups, d'espais vectorials i de monoides. A $\cat{Top}$, en canvi, un isomorfisme és un homeomorfisme\index{homeomorfisme}: una aplicació contínua i bijectiva amb inversa contínua. Aquí es veu que la bijectivitat no basta: hi ha aplicacions contínues i bijectives que no són isomorfismes, perquè la seva inversa no és contínua. La definició categòrica capta la noció correcta sense haver-la d'enunciar cas per cas.
\end{exemple}

\begin{exemple}
En un preordre, un isomorfisme entre $x$ i $y$ vol dir que hi ha fletxes $x\to y$ i $y\to x$, és a dir,
\[
x\le y
\qquad\text{i}\qquad
y\le x.
\]
En un ordre parcial això obliga $x=y$. En un preordre general, en canvi, poden existir objectes isomorfs diferents: són exactament els elements equivalents per la relació $x\sim y$ si $x\le y$ i $y\le x$.
\end{exemple}

\begin{exemple}[Isomorfisme i equivalència lògica]\label{ex:iso-equivalencia-logica}
A la categoria prima $\cat{Prov}_T$ d'un sistema deductiu $T$, dues fórmules $p$ i $q$ són isomorfes exactament quan
\[
p\vdash_T q
\qquad\text{i}\qquad
q\vdash_T p.
\]
Per tant, l'isomorfisme categòric es tradueix aquí en \textbf{interdeduïbilitat}. Aquest és un primer exemple de com una noció categòrica adquireix una lectura lògica natural.
\end{exemple}

\begin{observacio}
L'exemple dels preordres mostra per què la definició categòrica és la bona. Es podria pensar a definir un isomorfisme com un morfisme \enquote{bijectiu}, però això requeriria parlar dels elements dels objectes, i no sempre té sentit. A més, fins i tot quan en té, pot donar la noció equivocada: a $\cat{Pos}$ existeixen aplicacions monòtones bijectives que no són isomorfismes perquè la seva inversa no és monòtona. La definició mitjançant composició i identitats evita aquest parany.
\end{observacio}

\begin{observacio}[Propietats invariants per isomorfisme]\label{obs:invariancia}
Sigui $\cat{C}$ una categoria. Una propietat dels objectes de $\cat{C}$ és \textbf{invariant per isomorfisme} a $\cat{C}$ si, sempre que un objecte la té, tot objecte isomorf a ell a $\cat{C}$ també la té. Com la relació $\cong_{\cat{C}}$ de l'observació anterior, aquesta noció és sempre relativa a una categoria determinada. Una propietat definida només a partir de l'estructura categòrica de $\cat{C}$ ---els seus objectes, morfismes, composició i identitats, sense fer servir la igualtat entre objectes--- ha de ser invariant per isomorfisme a $\cat{C}$, perquè un isomorfisme transporta qualsevol afirmació d'aquesta mena d'un objecte a l'altre. Per això la invariància per isomorfisme és un criteri bàsic per reconèixer quan una propietat depèn de la forma de l'objecte en la categoria i no de la seva presentació concreta.

Aquesta noció dona un mètode útil. Per provar que dos objectes \emph{són} isomorfs, n'hi ha prou d'exhibir un isomorfisme entre ells. Provar que \emph{no} ho són sol ser més difícil, perquè caldria descartar tots els isomorfismes possibles; la via pràctica és trobar una propietat invariant per isomorfisme a la categoria en qüestió que un dels objectes tingui i l'altre no. Per exemple, a $\cat{Pos}$, siguin $A$ la cadena $a<b<c$ i $B$ l'ordre amb $x<y$ i $x<z$ però $y,z$ incomparables. Comprovar directament que no són isomorfs exigiria descartar totes les bijeccions monòtones amb inversa monòtona; en canvi, n'hi ha prou d'observar que \emph{estar totalment ordenat} és una propietat invariant per isomorfisme, que $A$ té i $B$ no. Per tant $A$ i $B$ no són isomorfs.

La implicació inversa no és automàtica: una propietat pot ser invariant per isomorfisme sense haver estat formulada internament en el llenguatge categòric. L'exemple anterior n'és una mostra: \emph{estar totalment ordenat} és invariant per isomorfisme a $\cat{Pos}$, però l'hem definit parlant dels elements de l'ordre, no de fletxes.
\end{observacio}

\begin{observacio}
La noció d'isomorfisme és autènticament categòrica: es formula només amb composició i identitats, de manera que té sentit a \emph{qualsevol} categoria, sense referència als elements dels objectes. Això la fa més abstracta i més general que les nocions d'isomorfisme pròpies de cada context ---bijeccions, isomorfismes de grups, homeomorfismes---, que en són casos particulars.

Aquesta generalitat es fa visible en un cas que ja coneixem. Havíem identificat els monoides amb les categories d'un sol objecte; ara podem identificar els grups amb exactament aquelles categories d'un sol objecte en què tota fletxa és un isomorfisme. Portant la mateixa idea a un nombre qualsevol d'objectes s'obté una noció d'importància considerable en les matemàtiques actuals: un \emph{grupoide} és una categoria en què tot morfisme és un isomorfisme.\index{grupoide}

El grupoide unifica dos casos que ja hem trobat. Un grup és un grupoide amb un sol objecte, com acabem de veure. I un preordre en què tot morfisme és un isomorfisme és, precisament, un en què $x\le y$ implica $y\le x$: una relació simètrica, és a dir, una \emph{relació d'equivalència}. Així, una relació d'equivalència no és res més que un grupoide prim: les classes d'isomorfisme dels seus objectes són les classes d'equivalència.
\end{observacio}

\begin{definicio}\index{endomorfisme}\index{automorfisme}
Un morfisme $f\colon A\to A$, amb domini i codomini iguals, s'anomena \textbf{endomorfisme} de $A$. Un endomorfisme que a més és un isomorfisme s'anomena \textbf{automorfisme} de $A$.
\end{definicio}

Els automorfismes d'un objecte $A$, amb la composició, formen un grup, el \textbf{grup d'automorfismes} de $A$, que escrivim $\Aut(A)$.\index{Aut@$\Aut$} Torna a aparèixer així la identificació d'abans: $\Aut(A)$ és el grup associat a la categoria d'un sol objecte que té els automorfismes de $A$ com a fletxes.

\section{Classes especials de morfismes}

Els isomorfismes són els morfismes invertibles: tenen un invers per tots dos costats. Sovint, però, un morfisme té invers només per un costat, o satisfà una propietat de cancel·lació més feble que la invertibilitat. Aquestes nocions ---monomorfismes, epimorfismes, seccions i retraccions--- són bàsiques. A $\cat{Set}$, els monomorfismes i els epimorfismes coincideixen amb les aplicacions injectives i exhaustives; com veurem, però, aquesta coincidència és pròpia de $\cat{Set}$ i d'algunes categories concretes, no una propietat de les categories en general.

Recordem que, a $\cat{Set}$, una aplicació $f$ és \emph{injectiva} si $f(a)=f(a')$ implica $a=a'$, i \emph{exhaustiva} si tot element del codomini té alguna antiimatge. Aquestes definicions parlen d'elements; les versions categòriques que segueixen les reformulen només amb fletxes.

\subsection{Monomorfismes i epimorfismes}\label{subsec:mono-epi}

\begin{definicio}\index{monomorfisme}\index{epimorfisme}
\label{def:mono-epi}
Un morfisme $f\colon A\to B$ és un \textbf{monomorfisme} (o \textbf{mono}), i escrivim $f\colon A\rightarrowtail B$, si es pot cancel·lar per l'esquerra: per a tota parella de morfismes $g,h\colon C\to A$,
\[
\begin{tikzpicture}[baseline=(m.center)]
  \matrix (m) [matrix of math nodes, column sep=3.4em] { C & A & B \\ };
  \draw[-{Stealth[scale=1.0]}] ([yshift=2.6pt]m-1-1.east) -- node[above] {$g$} ([yshift=2.6pt]m-1-2.west);
  \draw[-{Stealth[scale=1.0]}] ([yshift=-2.6pt]m-1-1.east) -- node[below] {$h$} ([yshift=-2.6pt]m-1-2.west);
  \draw[-{Stealth[scale=1.1]}] (m-1-2) -- node[above] {$f$} (m-1-3);
\end{tikzpicture}
\qquad
f\comp g=f\comp h\ \Longrightarrow\ g=h.
\]
Un morfisme $f\colon A\to B$ és un \textbf{epimorfisme} (o \textbf{epi}), i escrivim $f\colon A\twoheadrightarrow B$, si es pot cancel·lar per la dreta: per a tota parella $g,h\colon B\to C$,
\[
\begin{tikzpicture}[baseline=(m.center)]
  \matrix (m) [matrix of math nodes, column sep=3.4em] { A & B & C \\ };
  \draw[-{Stealth[scale=1.1]}] (m-1-1) -- node[above] {$f$} (m-1-2);
  \draw[-{Stealth[scale=1.0]}] ([yshift=2.6pt]m-1-2.east) -- node[above] {$g$} ([yshift=2.6pt]m-1-3.west);
  \draw[-{Stealth[scale=1.0]}] ([yshift=-2.6pt]m-1-2.east) -- node[below] {$h$} ([yshift=-2.6pt]m-1-3.west);
\end{tikzpicture}
\qquad
g\comp f=h\comp f\ \Longrightarrow\ g=h.
\]
\end{definicio}

\begin{observacio}[Com cal llegir aquests diagrames]\label{obs:llegir-mono-epi}
Els diagrames de la definició~\ref{def:mono-epi} \emph{no} afirmen que res commuti: són esquemes de prova, i s'han de llegir per a cada parella possible de fletxes $g,h$. Al diagrama del monomorfisme hi ha dos camins de $C$ a $B$, \enquote{primer $g$, després $f$} i \enquote{primer $h$, després $f$}, i hi ha també les dues fletxes paral·leles $g,h$ de $C$ a $A$. Recordem que un diagrama amb dues fletxes paral·leles commuta exactament quan són iguals. La definició diu, doncs:
\begin{quote}
si els dos camins de $C$ a $B$ coincideixen, aleshores les fletxes paral·leles ja coincideixen.
\end{quote}
Amb el vocabulari de la secció~\ref{sec:diagrames}: l'antecedent no és que el diagrama commuti, sinó només que commutin els camins que acaben a $B$; la conclusió és que el diagrama commuta del tot. Un monomorfisme és, doncs, una fletxa per a la qual \emph{commutar després de $f$ implica commutar}.
Llegit en sentit contrari: si $g\neq h$, aleshores $f\comp g\neq f\comp h$. Un monomorfisme \emph{no confon} mai dues fletxes diferents que arriben a $A$; després de compondre-les amb $f$, continuen sent diferents.

Cal anar amb compte amb una lectura massa ràpida. Que $f$ sigui mono no impedeix que hi hagi parelles $g\neq h\colon C\to A$; el diagrama amb dues fletxes paral·leles diferents existeix perfectament. El que exclou és que, sent diferents, es facin iguals en compondre-les amb $f$.

El diagrama de l'epimorfisme es llegeix de la mateixa manera, amb les fletxes de prova a l'altre costat: si $g\neq h\colon B\to C$, aleshores $g\comp f\neq h\comp f$. Dues fletxes que surten de $B$ i coincideixen després de $f$ ja eren iguals; $f$ arriba prou a $B$ perquè cap diferència entre $g$ i $h$ no se li escapi. A $\cat{Set}$, prenent com a fletxes de prova les aplicacions des d'un conjunt d'un sol element, que són els elements de $A$, la lectura del monomorfisme es converteix exactament en la injectivitat, com veurem tot seguit.
\end{observacio}

\begin{definicio}\index{bimorfisme}
Un morfisme que és alhora un monomorfisme i un epimorfisme s'anomena \textbf{bimorfisme}.
\end{definicio}

\begin{exemple}
A $\cat{Set}$, els monomorfismes són exactament les aplicacions \textbf{injectives} i els epimorfismes, les \textbf{exhaustives}. Vegem el cas dels monomorfismes: si $f$ és injectiva i $f\comp g=f\comp h$, aleshores per a cada $c$ tenim $f(g(c))=f(h(c))$, d'on $g(c)=h(c)$; recíprocament, si $f$ no és injectiva, hi ha $a\neq a'$ amb $f(a)=f(a')$, i les dues aplicacions constants $g,h\colon\{\ast\}\to A$ amb valors $a$ i $a'$ compleixen $f\comp g=f\comp h$ però $g\neq h$.

Per als epimorfismes, sigui $f\colon A\to B$. Si $f$ és exhaustiva i $g\comp f=h\comp f$, aleshores per a cada $b\in B$ existeix $a$ amb $f(a)=b$, i $g(b)=g(f(a))=h(f(a))=h(b)$, d'on $g=h$. Recíprocament, si $f$ no és exhaustiva, hi ha $b_0\in B$ sense antiimatge, és a dir, $b_0\notin f(A)$. Definim $g,h\colon B\to\{0,1\}$ a \emph{tot} $B$ per
\[
g(b)=0 \text{ per a tot } b\in B,
\qquad
h(b)=\begin{cases}1 & \text{si } b=b_0,\\ 0 & \text{si } b\neq b_0.\end{cases}
\]
Aleshores $g$ i $h$ coincideixen a tot $B$ tret de $b_0$, de manera que $g\neq h$. Però per a tot $a\in A$ tenim $f(a)\neq b_0$, ja que $b_0\notin f(A)$; per tant $g(f(a))=0=h(f(a))$, és a dir, $g\comp f=h\comp f$. Això mostra que $f$ no és un epimorfisme.

Convé notar que cap dels dos arguments no tria una \emph{secció} de $f$: per als monomorfismes es comparen dues seleccions d'elements des d'un conjunt unitari, i per als epimorfismes es construeixen explícitament dues aplicacions cap a $\{0,1\}$. En particular, la caracterització \enquote{epimorfisme $=$ exhaustiva} a $\cat{Set}$ no fa servir l'axioma de l'elecció, i no s'ha de confondre amb l'enunciat, ben diferent, que tota aplicació exhaustiva \emph{admet una secció} ---aquest sí equivalent a l'axioma de l'elecció (vegeu més avall).
\end{exemple}

\begin{proposicio}
Tot isomorfisme és un bimorfisme: si $f$ és un isomorfisme, aleshores és alhora un monomorfisme i un epimorfisme.
\end{proposicio}

\begin{prova}
Sigui $f\colon A\to B$ un isomorfisme amb invers $f^{-1}$. Si $f\comp g=f\comp h$, aleshores
\[
g=f^{-1}\comp(f\comp g)=f^{-1}\comp(f\comp h)=h,
\]
de manera que $f$ és mono. Si $g\comp f=h\comp f$, aleshores
\[
g=(g\comp f)\comp f^{-1}=(h\comp f)\comp f^{-1}=h,
\]
de manera que $f$ és epi.
\end{prova}

\begin{observacio}
A $\cat{Set}$, un morfisme és un isomorfisme si i només si és alhora mono i epi, és a dir, si i només si, com a aplicació, és injectiu i exhaustiu, o sigui \textbf{bijectiu}. Dit d'una altra manera, a $\cat{Set}$ els bimorfismes són exactament els isomorfismes. El recíproc de la proposició anterior, però, no és cert: en general un bimorfisme no és un isomorfisme. L'única fletxa no trivial de la categoria $\cat{2}$ és un bimorfisme ---per la simple raó que gairebé no hi ha morfismes amb què fer les cancel·lacions---, però no té invers, de manera que no és un isomorfisme. Un exemple més substancial es dona a $\cat{Top}$: una bijecció contínua és un bimorfisme, però no és un isomorfisme si la seva inversa no és contínua.
\end{observacio}

\begin{observacio}[Atenció: raonar amb elements]\label{obs:raonar-elements}\index{element!perill de raonar amb}
A $\cat{Set}$ podem raonar amb elements per dos motius: cada element $a\in A$ és una fletxa $\{\ast\}\to A$, i dues aplicacions $A\to B$ són iguals si coincideixen sobre tots els elements. En una categoria qualsevol, cap d'aquestes dues coses no està garantida.
\begin{itemize}
\item \emph{Pot ser que no hi hagi elements.} En un preordre o en un monoide vist com a categoria, una fletxa no envia res enlloc, i preguntar si és \enquote{injectiva} no té sentit. Mono i epi, en canvi, sempre tenen sentit, perquè es defineixen per composició (exercici~\ref{exr:pr-prima-mono-epi} de la Pràctica).
\item \emph{L'aplicació subjacent pot enganyar.} A $\cat{Top}$, una bijecció contínua és mono i epi sense ser un isomorfisme (observació anterior). A $\cat{Mon}$, la inclusió $\iota\colon(\mathbb N,+,0)\hookrightarrow(\mathbb Z,+,0)$ és un epimorfisme, tot i que no és exhaustiva. Si dos homomorfismes $g,h\colon\mathbb Z\to M$ coincideixen sobre $\mathbb N$, aleshores $g(-1)$ i $h(-1)$ són tots dos inversos de $g(1)=h(1)$, i en un monoide l'invers d'un element, si existeix, és únic. Per tant, $g(-1)=h(-1)$ i, com que tot enter és suma d'un natural i de còpies de $-1$, $g=h$.
\item \emph{Els \enquote{punts} poden no separar fletxes.} A $\cat{Grp}$, el grup trivial $1$ fa el paper de $\{\ast\}$, però de $1$ a qualsevol grup només hi ha un homomorfisme. Les fletxes $1\to G$ no distingeixen, doncs, cap parell d'homomorfismes diferents.
\end{itemize}
Per això, en una categoria general els arguments es fan amb fletxes i composicions, i reservem els adjectius \emph{injectiu}, \emph{exhaustiu} i \emph{bijectiu} per a les aplicacions, incloses les aplicacions subjacents dels morfismes d'una categoria concreta. La definició de monomorfisme és, de fet, una \enquote{injectivitat} respecte de \emph{totes} les fletxes de prova $C\to A$, i no només respecte dels elements.
\end{observacio}

\subsection{Seccions i retraccions}\label{subsec:seccions-retraccions}

\begin{definicio}\index{secció}\index{retracció}\index{invers!per un costat}
Considerem un morfisme $f\colon A\to B$. Diem que un morfisme $s\colon B\to A$ és una \textbf{secció} de $f$ quan es compleix $f\comp s=\id_B$. Això equival a dir que el diagrama següent és commutatiu:
\[
\begin{tikzpicture}[baseline=(m.center)]
  \matrix (m) [matrix of math nodes, row sep=2.6em, column sep=2.6em] {
    B & A \\
    B & \\
  };
  \draw[-{Stealth[scale=1.1]}] (m-1-1) -- node[above] {$s$} (m-1-2);
  \draw[-{Stealth[scale=1.1]}] (m-1-2) -- node[right] {$f$} (m-2-1);
  \draw[-{Stealth[scale=1.1]}] (m-1-1) -- node[left] {$\id_B$} (m-2-1);
\end{tikzpicture}
\]
També diem en aquest cas que $s$ és un \emph{invers per la dreta} de $f$.

Simètricament, diem que un morfisme $r\colon B\to A$ és una \textbf{retracció} de $f$ quan es compleix $r\comp f=\id_A$. Això equival a dir que el diagrama següent és commutatiu:
\[
\begin{tikzpicture}[baseline=(m.center)]
  \matrix (m) [matrix of math nodes, row sep=2.6em, column sep=2.6em] {
    A & B \\
    A & \\
  };
  \draw[-{Stealth[scale=1.1]}] (m-1-1) -- node[above] {$f$} (m-1-2);
  \draw[-{Stealth[scale=1.1]}] (m-1-2) -- node[right] {$r$} (m-2-1);
  \draw[-{Stealth[scale=1.1]}] (m-1-1) -- node[left] {$\id_A$} (m-2-1);
\end{tikzpicture}
\]
També diem en aquest cas que $r$ és un \emph{invers per l'esquerra} de $f$.
\end{definicio}

Una conseqüència òbvia d'aquestes dues definicions és que quan es compleix, per exemple, $f\comp s=\id_B$, se'n segueix alhora que $s$ és una secció de $f$ i que $f$ és una retracció de $s$; el mateix succeeix amb $r\comp f=\id_A$, on $r$ és una retracció de $f$ i $f$ és una secció de $r$.

Un morfisme pot tenir cap, una o diverses seccions, i igualment cap, una o diverses retraccions. Vegem-ne exemples a $\cat{Set}$:
\begin{itemize}
\item L'única aplicació $f\colon\emptyset\to\{*\}$ no té cap secció: no existeix cap aplicació $\{*\}\to\emptyset$.
\item La injecció $f\colon\{0\}\hookrightarrow\{0,1\}$ associada a la inclusió $\{0\}\subseteq\{0,1\}$ té una única retracció, l'única aplicació $\{0,1\}\to\{0\}$; en canvi no té cap secció.
\item L'aplicació $f\colon\{0,1\}\to\{*\}$ té dues seccions distintes, $*\mapsto 0$ i $*\mapsto 1$: cada tria d'una antiimatge en dona una.
\end{itemize}
Tenir secció o retracció és, doncs, una propietat especial de $f$, i quan en té, en general no és única.

\begin{proposicio}\label{prop:escindit-mono-epi}
Sigui $f\colon A\to B$. Si $f$ té alguna secció, aleshores $f$ és un epimorfisme. Si $f$ té alguna retracció, aleshores $f$ és un monomorfisme.
\end{proposicio}

\begin{prova}
Suposem que $f$ té una retracció $r$, amb $r\comp f=\id_A$, i siguin $g,h\colon C\to A$ amb $f\comp g=f\comp h$. Component per l'esquerra amb $r$,
\[
g=\id_A\comp g=(r\comp f)\comp g=r\comp(f\comp g)=r\comp(f\comp h)=(r\comp f)\comp h=h,
\]
de manera que $f$ és mono. El cas de la secció és anàleg: si $f$ té secció $s$, amb $f\comp s=\id_B$, i $g,h\colon B\to C$ amb $g\comp f=h\comp f$, aleshores component per la dreta amb $s$,
\[
g=g\comp\id_B=g\comp(f\comp s)=(g\comp f)\comp s=(h\comp f)\comp s=h\comp(f\comp s)=h\comp\id_B=h,
\]
de manera que $f$ és epi.
\end{prova}

\begin{definicio}\index{morfisme!escindit}\index{epimorfisme!escindit}\index{monomorfisme!escindit}
Diem que un morfisme $f$ queda \textbf{escindit} quan té una secció o una retracció. En concret, si $f$ té una secció, aleshores $f$ és un epimorfisme i se'n diu \textbf{epimorfisme escindit}; si $f$ té una retracció, aleshores $f$ és un monomorfisme i se'n diu \textbf{monomorfisme escindit}.
\end{definicio}

\begin{observacio}
El recíproc no val en general: hi ha monomorfismes que no tenen cap retracció (i per tant no són escindits) i epimorfismes que no tenen cap secció. En efecte, a $\cat{Set}$ tota aplicació injectiva amb domini no buit té alguna retracció i tota exhaustiva té alguna secció ---aquesta darrera afirmació és, de fet, equivalent a l'\emph{axioma de l'elecció}---, però en altres categories la situació és diferent.

Finalment, un morfisme $f$ és un \textbf{isomorfisme} si i només si té alhora alguna secció i alguna retracció. En aquest cas la secció i la retracció són forçosament \emph{úniques} i coincideixen: si $f\comp s=\id_B$ i $r\comp f=\id_A$, aleshores
\[
r=r\comp\id_B=r\comp(f\comp s)=(r\comp f)\comp s=\id_A\comp s=s.
\]
En canvi, la unicitat d’una secció o d’una retracció, considerada
separadament, no caracteritza els isomorfismes en una categoria general; per exemple, la injecció
\[
 i\colon\{0\}\hookrightarrow\{0,1\}
\]
té una única retracció $r\colon\{0,1\}\to\{0\}$, però no és un isomorfisme. No és una subtilesa avançada, sinó una contradicció amb un exemple de la mateixa subsecció.
\end{observacio}

\section{Construccions sobre categories}

A partir de categories donades se'n poden construir de noves. Aquestes construccions seran útils constantment, i algunes ---com la categoria oposada--- són fonamentals.

\subsection{La categoria oposada}\index{categoria!oposada}

\begin{definicio}
Donada una categoria $\cat{C}$, la seva \textbf{categoria oposada} $\cat{C}\op$ té:
\begin{itemize}
\item objectes: els mateixos que $\cat{C}$, $\Ob(\cat{C}\op)=\Ob(\cat{C})$;
\item morfismes: per a cada morfisme $f\colon A\to B$ de $\cat{C}$, un morfisme $f\op\colon B\to A$ de $\cat{C}\op$, de manera que
\[
\Hom_{\cat{C}\op}(B,A)=\{f\op\mid f\in\Hom_{\cat{C}}(A,B)\};
\]
\item composició: s'obté invertint l'ordre; per a $f\colon A\to B$ i $g\colon B\to C$ de $\cat{C}$,
\[
f\op\comp g\op=(g\comp f)\op\colon C\to A;
\]
\item identitats: la identitat de $A$ a $\cat{C}\op$ és $(\id_A)\op$.
\end{itemize}
És a dir, $\cat{C}\op$ s'obté invertint el sentit de totes les fletxes.
\end{definicio}

Visualment, la composició a $\cat{C}$ i la corresponent a $\cat{C}\op$ inverteixen l'ordre:
\[
\begin{array}{c@{\qquad\qquad}c}
\cat{C}: & \cat{C}\op: \\[0.6em]
\begin{tikzpicture}[baseline=(m.center)]
  \matrix (m) [matrix of math nodes, row sep=2.6em, column sep=2.6em] {
    A & B \\
      & C \\
  };
  \draw[-{Stealth[scale=1.1]}] (m-1-1) -- node[above] {$f$} (m-1-2);
  \draw[-{Stealth[scale=1.1]}] (m-1-2) -- node[right] {$g$} (m-2-2);
  \draw[-{Stealth[scale=1.1]}] (m-1-1) -- node[below left] {$g\comp f$} (m-2-2);
\end{tikzpicture}
&
\begin{tikzpicture}[baseline=(m.center)]
  \matrix (m) [matrix of math nodes, row sep=2.6em, column sep=2.6em] {
    A & B \\
      & C \\
  };
  \draw[-{Stealth[scale=1.1]}] (m-1-2) -- node[above] {$f\op$} (m-1-1);
  \draw[-{Stealth[scale=1.1]}] (m-2-2) -- node[right] {$g\op$} (m-1-2);
  \draw[-{Stealth[scale=1.1]}] (m-2-2) -- node[below left] {$f\op\comp g\op$} (m-1-1);
\end{tikzpicture}
\end{array}
\]
L'associativitat a $\cat{C}\op$ es dedueix de la de $\cat{C}$ invertint l'ordre, i les lleis d'identitat, de les de $\cat{C}$: per exemple, $f\op\comp(\id_B)\op=(\id_B\comp f)\op=f\op$. La categoria oposada és la base de la \textbf{dualitat}\index{dualitat}, que desenvolupem a la secció~\ref{sec:dualitat}: tota afirmació formulada purament en el llenguatge categòric dona, invertint les fletxes, una afirmació dual, i aquesta correspondència es pot fer del tot precisa.

\begin{observacio}
Aplicar dues vegades l'oposada retorna la categoria de partida:
\[
(\cat{C}\op)\op=\cat{C}.
\]
Invertir les fletxes i tornar-les a invertir les deixa com estaven.
\end{observacio}

Vegem com queda l'oposada en dos casos. Quan un preordre $(P,\le)$ es considera com una categoria ---amb un morfisme $x\to y$ quan $x\le y$---, la seva oposada té un morfisme $x\to y$ exactament quan $y\le x$. I quan un monoide $(M,\ast,e)$ es considera com una categoria, la seva oposada té els mateixos morfismes ---els elements de $M$---, però la composició de $a$ i $b$ és $b\ast a$ en lloc de $a\ast b$.

\subsection{La categoria producte}\label{subsec:categoria-producte}\index{categoria!producte}

\begin{definicio}
Donades dues categories $\cat{C}$ i $\cat{D}$, la \textbf{categoria producte} $\cat{C}\times\cat{D}$ té:
\begin{itemize}
\item objectes: $\Ob(\cat{C}\times\cat{D})$ és la col·lecció de les parelles $(A,X)$ amb $A\in\Ob(\cat{C})$ i $X\in\Ob(\cat{D})$;
\item morfismes: $\Hom_{\cat{C}\times\cat{D}}\bigl((A,X),(B,Y)\bigr)=\Hom_{\cat{C}}(A,B)\times\Hom_{\cat{D}}(X,Y)$, és a dir, les parelles $(f,g)$ amb $f\colon A\to B$ a $\cat{C}$ i $g\colon X\to Y$ a $\cat{D}$;
\item composició: component a component,
\[
(f',g')\comp(f,g)=(f'\comp f,\ g'\comp g);
\]
\item identitats: $\id_{(A,X)}=(\id_A,\id_X)$.
\end{itemize} Els axiomes de categoria per a $\cat{C}\times\cat{D}$ se segueixen dels de $\cat{C}$ i $\cat{D}$ aplicats a cada component.
\end{definicio}

\begin{exemple}
Prenem la categoria $\cat{2}$ de la subsecció~\ref{subsec:finites} ---dos objectes i una única fletxa no trivial---, i n'escrivim els objectes com $0$ i $1$, amb la fletxa no trivial $0\to 1$. La categoria producte $\cat{2}\times\cat{2}$ té quatre objectes, que es disposen de manera natural en un quadrat:
\[
\begin{tikzpicture}[baseline=(m.center)]
  \matrix (m) [matrix of math nodes, row sep=2.6em, column sep=3.4em] {
    (0,0) & (0,1) \\
    (1,0) & (1,1) \\
  };
  \draw[-{Stealth[scale=1.1]}] (m-1-1) -- (m-1-2);
  \draw[-{Stealth[scale=1.1]}] (m-1-1) -- (m-2-1);
  \draw[-{Stealth[scale=1.1]}] (m-1-2) -- (m-2-2);
  \draw[-{Stealth[scale=1.1]}] (m-2-1) -- (m-2-2);
  \draw[-{Stealth[scale=1.1]}] (m-1-1) -- (m-2-2);
\end{tikzpicture}
\]
Cada fletxa és una parella de fletxes de $\cat{2}$. Com que $\cat{2}$ és prima, també ho és $\cat{2}\times\cat{2}$: entre dos objectes hi ha com a màxim un morfisme. Per això els dos camins del quadrat de $(0,0)$ a $(1,1)$ tenen necessàriament la mateixa composició, que és la diagonal; el quadrat commuta.
\end{exemple}

\begin{observacio}[Dos usos de \enquote{producte}]\label{obs:dos-productes}\index{producte!d'objectes i de categories}
No s'han de confondre dues construccions que porten el mateix nom. La categoria producte $\cat{C}\times\cat{D}$ és una \emph{categoria} nova, construïda a partir de dues categories. En canvi, el producte $A\times B$ de dos objectes ---que la introducció fa servir com a primer exemple de propietat universal i que avançarem a la subsecció~\ref{subsec:producte-coproducte}--- és un \emph{objecte} d'una mateixa categoria, caracteritzat per una propietat universal. Les dues nocions estan relacionades: quan disposem de la categoria $\cat{Cat}$ de les categories petites (capítol~\ref{cap:functors}), si $\cat{C}$ i $\cat{D}$ són petites, la categoria $\cat{C}\times\cat{D}$ es podrà veure com el producte dels objectes $\cat{C}$ i $\cat{D}$ dins de $\cat{Cat}$. Però no són la mateixa construcció: l'una fabrica una categoria a partir de dues, i l'altra viu dins d'una sola categoria.
\end{observacio}

\subsection{Subcategories}\label{subsec:subcategories}\index{subcategoria}

Una altra manera d'obtenir noves categories a partir de categories antigues és restringint-ne els objectes o les fletxes.

\begin{definicio}
Una \textbf{subcategoria} $\cat{S}$ d'una categoria $\cat{C}$ ve donada per una subcol·lecció d'objectes $\Ob(\cat{S})\subseteq\Ob(\cat{C})$ i, per a cada parella $A,B\in\Ob(\cat{S})$, una subcol·lecció de morfismes $\Hom_{\cat{S}}(A,B)\subseteq\Hom_{\cat{C}}(A,B)$, de manera que:
\begin{itemize}
\item la identitat $\id_A$ pertany a $\Hom_{\cat{S}}(A,A)$ per a cada $A\in\Ob(\cat{S})$;
\item la composició de dos morfismes de $\cat{S}$ componibles pertany a $\cat{S}$.
\end{itemize}
Amb aquestes condicions, $\cat{S}$ és una categoria per dret propi, amb la composició i les identitats heretades de $\cat{C}$.
\end{definicio}

Dos casos particulars, en cert sentit oposats, són especialment útils. Es diu que $\cat{S}$ és una \textbf{subcategoria plena}\index{subcategoria!plena} quan $\Hom_{\cat{S}}(A,B)=\Hom_{\cat{C}}(A,B)$ per a tot $A,B\in\Ob(\cat{S})$; una subcategoria plena queda determinada, doncs, només per la seva col·lecció d'objectes. Simètricament, es diu que $\cat{S}$ és una \textbf{subcategoria ampla}\index{subcategoria!ampla} quan $\Ob(\cat{S})=\Ob(\cat{C})$, encara que només contingui alguns dels morfismes. En resum: una subcategoria plena restringeix els objectes i conserva tots els morfismes entre ells; una subcategoria ampla conserva tots els objectes i en restringeix els morfismes.

\begin{exemple}
$\cat{FinSet}$ és una subcategoria plena de $\cat{Set}$: pren com a objectes els conjunts finits i, entre dos d'ells, totes les aplicacions. En canvi, la subcategoria de $\cat{Set}$ que té tots els conjunts per objectes però només les aplicacions \emph{injectives} per morfismes és una subcategoria ampla que no és plena: conserva tots els objectes, però deixa fora les aplicacions que no són injectives.
\end{exemple}

\subsection{Categories sobre i sota un objecte}\index{categoria!sobre un objecte}\index{categoria!sota un objecte}

Fixat un objecte $A$ d'una categoria $\cat{C}$, se'n pot construir una categoria nova els objectes de la qual són els morfismes que apunten cap a $A$.

\begin{definicio}\label{def:slice}
Sigui $\cat{C}$ una categoria i $A$ un objecte. La \textbf{categoria de $\cat{C}$ sobre $A$} (en anglès, \emph{slice category}), que escrivim $\cat{C}/A$, té:
\begin{itemize}
\item objectes: $\Ob(\cat{C}/A)$ és la col·lecció dels morfismes $x\colon X\to A$ de $\cat{C}$ amb codomini $A$;
\item morfismes: $\Hom_{\cat{C}/A}(x,y)$, per a $x\colon X\to A$ i $y\colon Y\to A$, és el conjunt dels morfismes $h\colon X\to Y$ de $\cat{C}$ que fan commutar el triangle
\[
\begin{tikzpicture}[baseline=(m.center)]
  \matrix (m) [matrix of math nodes, row sep=2.4em, column sep=2.4em] {
    X & & Y \\
      & A & \\
  };
  \draw[-{Stealth[scale=1.0]}] (m-1-1) -- node[above] {$h$} (m-1-3);
  \draw[-{Stealth[scale=1.0]}] (m-1-1) -- node[below left] {$x$} (m-2-2);
  \draw[-{Stealth[scale=1.0]}] (m-1-3) -- node[below right] {$y$} (m-2-2);
\end{tikzpicture}
\qquad\text{és a dir, } y\comp h=x;
\]
\item composició: la de $\cat{C}$; si $h\in\Hom_{\cat{C}/A}(x,y)$ i $k\in\Hom_{\cat{C}/A}(y,z)$, aleshores $k\comp h\in\Hom_{\cat{C}/A}(x,z)$, perquè $z\comp(k\comp h)=(z\comp k)\comp h=y\comp h=x$;
\item identitats: la identitat de $x\colon X\to A$ és $\id_X$, que fa commutar el triangle perquè $x\comp\id_X=x$.
\end{itemize}
Els axiomes s'hereten de $\cat{C}$.
\end{definicio}

Anàlogament es defineix la \textbf{categoria de $\cat{C}$ sota $A$} (en anglès, \emph{coslice category}), que escrivim $A/\cat{C}$:
\begin{itemize}
\item objectes: $\Ob(A/\cat{C})$ és la col·lecció dels morfismes $x\colon A\to X$ de $\cat{C}$ amb domini $A$;
\item morfismes: $\Hom_{A/\cat{C}}(x,y)$, per a $x\colon A\to X$ i $y\colon A\to Y$, és el conjunt dels morfismes $h\colon X\to Y$ de $\cat{C}$ amb $h\comp x=y$;
\item composició i identitats: les de $\cat{C}$, com a $\cat{C}/A$.
\end{itemize}
Un morfisme de $A/\cat{C}$ és, doncs, un $h$ que fa commutar el triangle sota $A$:
\[
\begin{tikzpicture}[baseline=(m.center)]
  \matrix (m) [matrix of math nodes, row sep=2.4em, column sep=2.4em] {
      & A & \\
    X & & Y \\
  };
  \draw[-{Stealth[scale=1.0]}] (m-1-2) -- node[above left] {$x$} (m-2-1);
  \draw[-{Stealth[scale=1.0]}] (m-1-2) -- node[above right] {$y$} (m-2-3);
  \draw[-{Stealth[scale=1.0]}] (m-2-1) -- node[below] {$h$} (m-2-3);
\end{tikzpicture}
\]
Aquí també convé fer explícits els extrems, com a l'observació~\ref{obs:extrems-implicits}. Un morfisme de $\cat{C}/A$ de $x$ a $y$ no és el morfisme $h\colon X\to Y$ de $\cat{C}$ tot sol, sinó $h$ \emph{junt amb} els seus extrems $x$ i $y$ a $\cat{C}/A$: la terna $(x,y,h)$. La raó és la mateixa: un mateix $h\colon X\to Y$ de $\cat{C}$ pot fer commutar el triangle per a objectes distints de $\cat{C}/A$ ---per exemple, si $X$ és el domini de dos objectes diferents $x,x'\colon X\to A$---, i sense registrar els extrems tindria dominis i codominis indeterminats a $\cat{C}/A$. Escriurem $h$ quan els extrems se sobreentenguin. Anàlogament per a $A/\cat{C}$.

\begin{exemple}
A $\cat{Set}/A$, un objecte és una aplicació $x\colon X\to A$. Les fibres $x^{-1}(a)$, per a $a\in A$, permeten interpretar aquest objecte com una família de conjunts $\bigl(x^{-1}(a)\bigr)_{a\in A}$ indexada per $A$. Recíprocament, una família de conjunts indexada per $A$ dona, mitjançant la unió disjunta, una aplicació cap a $A$. Aquestes dues construccions no identifiquen literalment les dades ---la unió disjunta demana triar-ne una representació---, però sí que estableixen la correspondència adequada entre totes dues descripcions. La farem precisa als capítols~\ref{cap:functors} i~\ref{cap:transformacions}, quan disposem del llenguatge dels functors i de les equivalències de categories.
\end{exemple}

\begin{exemple}\index{conjunt apuntat}
Prenem $A=\{\ast\}$, un conjunt d'un sol element. Un objecte de $\{\ast\}/\cat{Set}$ és una aplicació $\{\ast\}\to X$, és a dir, l'elecció d'un element distingit de $X$, el \textbf{punt base}. Un morfisme és una aplicació que respecta el punt base. Així, $\{\ast\}/\cat{Set}$ és la categoria dels \textbf{conjunts apuntats} (en anglès, \emph{pointed sets}).
\end{exemple}

\section{Principi de dualitat}\label{sec:dualitat}\index{dualitat!principi de}

La dualitat no és només una llista de parelles de conceptes. És un principi que transforma definicions, enunciats i demostracions, i la categoria oposada permet formular-lo amb precisió. Aquesta secció en fa explícit el mecanisme; el retrobarem en tractar objectes inicials i terminals, productes i coproductes, productes fibrats i coproductes fibrats (en anglès, \emph{pullbacks} i \emph{pushouts}), i límits i colímits.

\subsection{Dualitzar un enunciat}\index{dualitat!d'un enunciat}

El principi s'aplica a afirmacions escrites \emph{exclusivament} en el llenguatge de les categories: hi poden intervenir objectes i morfismes, el domini i el codomini de cada morfisme, identitats, composicions i igualtats entre composicions, combinats amb connectors lògics i amb quantificadors sobre objectes i morfismes. Les definicions de monomorfisme, d'isomorfisme o de secció són d'aquesta mena. En canvi, \enquote{$f$ és injectiva} no ho és, perquè parla dels elements d'un conjunt.

Donada una afirmació $\varphi$ d'aquest tipus, n'obtenim la \textbf{dual} $\varphi\op$ \emph{invertint totes les fletxes}:
\begin{itemize}
  \item cada morfisme $f\colon A\to B$ s'inverteix: a la dual es llegeix com una fletxa $B\to A$; per tant, s'intercanvien domini i codomini;
  \item cada composició $g\comp f$ passa a ser $f\comp g$, és a dir, s'inverteix l'ordre de composició;
  \item les identitats, les igualtats, els connectors i els quantificadors es mantenen.
\end{itemize}
Dualitzar dues vegades retorna l'afirmació de partida: $(\varphi\op)\op=\varphi$.

La dual té una lectura immediata: $\varphi\op$ diu de $\cat{C}$ exactament el que $\varphi$ diu de $\cat{C}\op$. En efecte, un morfisme $A\to B$ de $\cat{C}\op$ és un morfisme $B\to A$ de $\cat{C}$, les identitats són les mateixes i la composició de $\cat{C}\op$ és la de $\cat{C}$ en l'ordre invers, $g\op\comp f\op=(f\comp g)\op$. Per tant, per a tota categoria $\cat{C}$,
\[
\varphi \text{ es compleix a } \cat{C}\op
\quad\Longleftrightarrow\quad
\varphi\op \text{ es compleix a } \cat{C}.
\]
D'aquí surt el principi.

\begin{teorema}[Principi de dualitat]\label{teo:dualitat}\index{dualitat!principi de}
Sigui $\varphi$ una afirmació expressable en el llenguatge de les categories. Si $\varphi$ és vàlida en totes les categories, la seva dual $\varphi\op$ també és vàlida en totes les categories.
\end{teorema}

\begin{prova}
Sigui $\cat{C}$ una categoria qualsevol. Com que $\cat{C}\op$ també és una categoria, per hipòtesi $\varphi$ es compleix a $\cat{C}\op$; per l'equivalència anterior, $\varphi\op$ es compleix a $\cat{C}$.
\end{prova}

La manera habitual d'usar el principi és aquesta: s'aplica un teorema conegut a $\cat{C}\op$ i se'n tradueix la conclusió al llenguatge de $\cat{C}$. No cal repetir la demostració; \textbf{la demostració dual és la demostració original amb totes les fletxes invertides}.

Aquesta formulació operativa és la que farem servir al llarg del llibre. Se'n pot donar també una versió completament formal: fixar un llenguatge de primer ordre per a la teoria de categories, definir la dualització com una transformació de fórmules i comprovar que els axiomes en queden invariants. Aquesta formalització, amb una versió sintàctica (sobre deduccions) i una de semàntica (sobre models), és a l'apèndix~\ref{cap:axiomatica-relacional} (secció~\ref{sec:ar-dualitat}); no és necessària per seguir la resta del llibre.

\subsection{Primer exemple: de monomorfisme a epimorfisme}

Recordem (subsecció~\ref{subsec:mono-epi}) que $m\colon A\to B$ és un \textbf{monomorfisme} si, per a tot parell $u,v\colon X\to A$,
\[
\begin{tikzpicture}[baseline=(m.center)]
  \matrix (m) [matrix of math nodes, column sep=3.2em] { X & A & B \\ };
  \draw[-{Stealth[scale=1.0]}] ([yshift=2.6pt]m-1-1.east) -- node[above] {$u$} ([yshift=2.6pt]m-1-2.west);
  \draw[-{Stealth[scale=1.0]}] ([yshift=-2.6pt]m-1-1.east) -- node[below] {$v$} ([yshift=-2.6pt]m-1-2.west);
  \draw[-{Stealth[scale=1.1]}] (m-1-2) -- node[above] {$m$} (m-1-3);
\end{tikzpicture}
\qquad
m\comp u = m\comp v \implies u=v.
\]
Formem-ne el dual: intercanviem domini i codomini i invertim l'ordre de cada composició. L'enunciat dual diu que, per a tot parell $u,v\colon A\to X$,
\[
\begin{tikzpicture}[baseline=(m.center)]
  \matrix (m) [matrix of math nodes, column sep=3.2em] { B & A & X \\ };
  \draw[-{Stealth[scale=1.1]}] (m-1-1) -- node[above] {$m$} (m-1-2);
  \draw[-{Stealth[scale=1.0]}] ([yshift=2.6pt]m-1-2.east) -- node[above] {$u$} ([yshift=2.6pt]m-1-3.west);
  \draw[-{Stealth[scale=1.0]}] ([yshift=-2.6pt]m-1-2.east) -- node[below] {$v$} ([yshift=-2.6pt]m-1-3.west);
\end{tikzpicture}
\qquad
u\comp m = v\comp m \implies u=v,
\]
que és exactament la definició d'\textbf{epimorfisme}. El dual de \enquote{ser un monomorfisme} és, doncs, \enquote{ser un epimorfisme}.

Ara podem dualitzar teoremes sense repetir-ne la demostració. Per exemple, en qualsevol categoria, si $A\xrightarrow{f}B\xrightarrow{g}C$ són monomorfismes, la composició $g\comp f$ també ho és. Pel principi de dualitat (teorema~\ref{teo:dualitat}), l'enunciat dual val igualment:
\begin{quote}
en qualsevol categoria, la composició de dos epimorfismes és un epimorfisme.
\end{quote}
No n'hem hagut de fer cap demostració nova. L'exercici~\ref{exr:pr-dual-mono-epi} de la Pràctica fa aquest pas en detall i assenyala exactament què s'inverteix en passar a $\cat{C}\op$.

\subsection{Segon exemple: de secció a retracció}

Recordem (subsecció~\ref{subsec:seccions-retraccions}) que $s\colon B\to A$ és una \textbf{secció} de $f\colon A\to B$ quan $f\comp s=\id_B$. Dualitzem aquest enunciat: $f$ passa a ser una fletxa $B\to A$, $s$ una fletxa $A\to B$, i la igualtat $f\comp s=\id_B$ es converteix en $s\comp f=\id_B$. En diagrames, invertim les tres fletxes del triangle:
\[
\begin{tikzpicture}[baseline=(m.center)]
  \matrix (m) [matrix of math nodes, row sep=2.6em, column sep=2.6em] {
    B & A \\
    B & \\
  };
  \draw[-{Stealth[scale=1.1]}] (m-1-1) -- node[above] {$s$} (m-1-2);
  \draw[-{Stealth[scale=1.1]}] (m-1-2) -- node[right] {$f$} (m-2-1);
  \draw[-{Stealth[scale=1.1]}] (m-1-1) -- node[left] {$\id_B$} (m-2-1);
\end{tikzpicture}
\quad f\comp s=\id_B
\qquad\qquad
\begin{tikzpicture}[baseline=(m.center)]
  \matrix (m) [matrix of math nodes, row sep=2.6em, column sep=2.6em] {
    B & A \\
    B & \\
  };
  \draw[-{Stealth[scale=1.1]}] (m-1-2) -- node[above] {$s$} (m-1-1);
  \draw[-{Stealth[scale=1.1]}] (m-2-1) -- node[right] {$f$} (m-1-2);
  \draw[-{Stealth[scale=1.1]}] (m-2-1) -- node[left] {$\id_B$} (m-1-1);
\end{tikzpicture}
\quad s\comp f=\id_B
\]
La igualtat de la dreta diu exactament que $s$ és una \textbf{retracció} de $f$. El dual de \enquote{tenir una secció} és, doncs, \enquote{tenir una retracció}, i el d'epimorfisme escindit és monomorfisme escindit.

Amb aquest diccionari es poden dualitzar també els resultats sobre seccions i retraccions; els exercicis proposats del capítol ho apliquen a la proposició~\ref{prop:escindit-mono-epi}.

\subsection{Tercer exemple: del producte al coproducte}\label{subsec:producte-coproducte}

A tall d'avançament ---els definirem en rigor al capítol de límits---, considerem el producte i el coproducte de dos objectes, que il·lustren la dualització d'un enunciat amb un quantificador d'existència única.

Un \emph{producte} de $A$ i $B$ és un objecte $A\times B$ amb dos morfismes $\pi_1\colon A\times B\to A$ i $\pi_2\colon A\times B\to B$ tals que, per a tot $X$ i tot parell $f\colon X\to A$, $g\colon X\to B$, existeix un únic $u\colon X\to A\times B$ amb $\pi_1\comp u=f$ i $\pi_2\comp u=g$:
\[
\begin{tikzpicture}[baseline=(m.center)]
  \matrix (m) [matrix of math nodes, row sep=2.6em, column sep=2.8em] {
     & X & \\
    A & A\times B & B \\
  };
  \draw[-{Stealth[scale=1.0]}] (m-1-2) -- node[above left] {$f$} (m-2-1);
  \draw[-{Stealth[scale=1.0]}] (m-1-2) -- node[above right] {$g$} (m-2-3);
  \draw[-{Stealth[scale=1.0]},dashed] (m-1-2) -- node[right] {$u$} (m-2-2);
  \draw[-{Stealth[scale=1.0]}] (m-2-2) -- node[below] {$\pi_1$} (m-2-1);
  \draw[-{Stealth[scale=1.0]}] (m-2-2) -- node[below] {$\pi_2$} (m-2-3);
\end{tikzpicture}
\]
En el diagrama seguim una convenció que farem servir cada cop que dibuixem una propietat universal: la \textbf{fletxa discontínua} marca el morfisme l'existència del qual afirma la propietat ---aquí el mediador $u$, un cop donades la resta de dades: l'objecte $X$ i les fletxes $f$ i $g$---, mentre que les fletxes contínues són les dades de partida. Igualment, escriurem $\exists!$ per abreujar \enquote{existeix un únic}, com acabem de fer en enunciar la propietat.

Invertint totes les fletxes obtenim un objecte $A+B$ amb dos morfismes $\iota_1\colon A\to A+B$ i $\iota_2\colon B\to A+B$ tals que, per a tot $X$ i tot parell $f\colon A\to X$, $g\colon B\to X$, existeix un únic $u\colon A+B\to X$ amb $u\comp\iota_1=f$ i $u\comp\iota_2=g$:
\[
\begin{tikzpicture}[baseline=(m.center)]
  \matrix (m) [matrix of math nodes, row sep=2.6em, column sep=2.8em] {
    A & A+B & B \\
     & X & \\
  };
  \draw[-{Stealth[scale=1.0]}] (m-1-1) -- node[above] {$\iota_1$} (m-1-2);
  \draw[-{Stealth[scale=1.0]}] (m-1-3) -- node[above] {$\iota_2$} (m-1-2);
  \draw[-{Stealth[scale=1.0]}] (m-1-1) -- node[below left] {$f$} (m-2-2);
  \draw[-{Stealth[scale=1.0]}] (m-1-3) -- node[below right] {$g$} (m-2-2);
  \draw[-{Stealth[scale=1.0]},dashed] (m-1-2) -- node[right] {$u$} (m-2-2);
\end{tikzpicture}
\]
Aquesta és la definició del \emph{coproducte} $A+B$. Els quantificadors $\forall$ i $\exists!$ i la conjunció de les dues equacions es mantenen; el que s'inverteix és el domini i el codomini de cada morfisme i l'ordre de cada composició.

\paragraph{Un cas que cal vigilar.} La construcció \enquote{categoria sobre $A$} es dualitza en la construcció \enquote{categoria sota $A$} (definició~\ref{def:slice}). Ara bé, en dualitzar aquesta construcció cal dualitzar també la categoria ambient. La fórmula precisa és
\[
(\cat{C}/A)\op\;\cong\;A/(\cat{C}\op),
\qquad
(A/\cat{C})\op\;\cong\;(\cat{C}\op)/A.
\]
No s'ha de llegir, doncs, que $(\cat{C}/A)\op\cong A/\cat{C}$ mantenint $\cat{C}$ intacta: això és fals en general. Per exemple, sigui $\cat{C}=\cat{2}=(0\to 1)$ i prenguem com a $A$ l'\emph{objecte} $1$ de $\cat{2}$ (no la categoria $\cat{1}$). La categoria $\cat{2}/1$ té dos objectes (les dues fletxes amb codomini $1$: la fletxa $0\to 1$ i $\id_1$), mentre que $1/\cat{2}$ en té un de sol ($\id_1$, l'única fletxa amb domini $1$); per tant $(\cat{2}/1)\op$ \emph{no} és isomorfa a $1/\cat{2}$. El que sí que val és $(\cat{2}/1)\op\cong 1/(\cat{2}\op)$: a $\cat{2}\op$ hi ha dues fletxes amb domini $1$, la identitat i la fletxa $1\to 0$. En tots els casos, la mateixa regla ---invertir les fletxes--- transforma cada noció en la seva dual, però la regla actua sobre \emph{tot} l'enunciat, ambient inclòs.

\subsection{Diccionari bàsic de dualitats}

Recollim les parelles duals. Les primeres ja han aparegut en aquest capítol; la resta les trobarem més endavant, i la taula creixerà a mesura que apareguin nous conceptes.
\begin{center}
\begin{tabular}{ll}
\hline
\textbf{Concepte} & \textbf{Concepte dual}\\
\hline
\multicolumn{2}{l}{\emph{Ja vistos en aquest capítol}}\\
monomorfisme & epimorfisme\\
secció & retracció\\
producte & coproducte\\
categoria sobre $A$ \ ($\cat{C}/A$) & categoria sota $A$ \ ($A/\cat{C}$)\\
\hline
\multicolumn{2}{l}{\emph{Més endavant}}\\
objecte inicial & objecte terminal\\
equalitzador & coequalitzador\\
producte fibrat & coproducte fibrat\\
con & cocon\\
límit & colímit\\
\hline
\end{tabular}
\end{center}
Algunes nocions són \textbf{autoduals}\index{dualitat!noció autodual}: la identitat, l'isomorfisme i la commutativitat d'un diagrama continuen sent nocions del mateix tipus després de dualitzar-les.

La taula recull les parelles duals com a \emph{nocions}: cada concepte i el seu dual. Quan aquestes nocions es realitzen en categories concretes, però, la dualitat s'ha d'aplicar a tot l'enunciat, ambient inclòs. La parella \enquote{sobre $A$ / sota $A$} n'és l'exemple a vigilar: com hem vist, l'enunciat correcte és $(\cat{C}/A)\op\cong A/(\cat{C}\op)$, no $(\cat{C}/A)\op\cong A/\cat{C}$ amb $\cat{C}$ intacta.

\subsection{Abast i precaucions}

El principi només s'aplica a afirmacions expressables en el llenguatge de les categories. Una afirmació que faci servir informació externa ---elements d'un conjunt, oberts d'un espai o una construcció concreta--- no té dual fins que no es reformula exclusivament amb objectes i morfismes. I quan l'afirmació es refereix a construccions fetes a partir d'una categoria, cal dualitzar també la categoria ambient, com acabem de veure amb les categories sobre i sota un objecte.

El principi de dualitat no permet concloure, per si sol, que una categoria concreta sigui equivalent a la seva oposada. El principi relaciona afirmacions sobre $\cat{C}$ amb les afirmacions duals sobre $\cat{C}\op$; no proporciona en general cap equivalència $\cat{C}\simeq\cat{C}\op$. Algunes categories sí que ho són ---per exemple, $\cat{2}=(0\to 1)$ és isomorfa a $\cat{2}\op$ intercanviant-ne els dos objectes---, però això s'ha de comprovar en cada cas. En canvi, $\cat{Set}$ no s'identifica amb $\cat{Set}\op$, encara que els teoremes generals sobre productes tinguin teoremes duals sobre coproductes.

Finalment, una precaució sobre la lògica. Els connectors i els quantificadors amb què escrivim una afirmació $\varphi$ pertanyen al llenguatge amb què \emph{parlem} de la categoria, i la dualitat els deixa intactes. Una altra cosa és el que passa quan els objectes de la categoria són ells mateixos fórmules, com a la categoria $\cat{Prov}_T$ de la subsecció~\ref{subsec:elementals}. Allà, \emph{quan existeixen}, la conjunció és un producte, la disjunció un coproducte, el cert un objecte terminal i el fals un objecte inicial, i la dualitat de fletxes els intercanvia. Però un preordre de deduïbilitat qualsevol no garanteix que aquests connectors existeixin: en el preordre discret de dues fórmules sense cap deducció entre elles, aquestes dues fórmules no tenen ni producte ni coproducte. I que la dualitat de fletxes coincideixi exactament amb la dualitat de De~Morgan exigeix, a més, un complement involutiu com el de les àlgebres booleanes. Tornarem sobre aquesta lògica interna més endavant, amb el llenguatge adequat.

\section{Propietats duals dels morfismes}\label{sec:propietats-duals}

El principi de dualitat estalvia feina: n'hi ha prou de demostrar les propietats dels monomorfismes; les dels epimorfismes se n'obtenen invertint les fletxes, sense cap demostració nova. Ho il·lustrem amb les propietats bàsiques.

\subsection{Propietats dels monomorfismes}

\begin{proposicio}\label{prop:mono-comp}
La composició de dos monomorfismes és un monomorfisme: si $f\colon A\to B$ i $g\colon B\to C$ són monomorfismes, aleshores $g\comp f$ ho és.
\end{proposicio}

\begin{prova}
Siguin $u,v\colon X\to A$ amb $(g\comp f)\comp u=(g\comp f)\comp v$, és a dir $g\comp(f\comp u)=g\comp(f\comp v)$. Com que $g$ és mono, $f\comp u=f\comp v$; i com que $f$ és mono, $u=v$.
\[
\begin{tikzpicture}[baseline=(m.center)]
  \matrix (m) [matrix of math nodes, column sep=3em] { X & A & B & C \\ };
  \draw[-{Stealth[scale=1.0]}] ([yshift=2.6pt]m-1-1.east) -- node[above] {$u$} ([yshift=2.6pt]m-1-2.west);
  \draw[-{Stealth[scale=1.0]}] ([yshift=-2.6pt]m-1-1.east) -- node[below] {$v$} ([yshift=-2.6pt]m-1-2.west);
  \draw[-{Stealth[scale=1.1]}] (m-1-2) -- node[above] {$f$} (m-1-3);
  \draw[-{Stealth[scale=1.1]}] (m-1-3) -- node[above] {$g$} (m-1-4);
\end{tikzpicture}
\]
\end{prova}

\begin{proposicio}\label{prop:mono-factor}
Si la composició $g\comp f$ és un monomorfisme, aleshores $f$ és un monomorfisme.
\end{proposicio}

\begin{prova}
Siguin $u,v\colon X\to A$ amb $f\comp u=f\comp v$. Component amb $g$, $(g\comp f)\comp u=(g\comp f)\comp v$; com que $g\comp f$ és mono, $u=v$. (No cal cap hipòtesi sobre $g$.)
\end{prova}

\begin{proposicio}\label{prop:iso-retr-mono}
Tot monomorfisme escindit és un monomorfisme, i tot isomorfisme és un monomorfisme escindit. Per tant,
\[
\text{isomorfisme} \;\Longrightarrow\; \text{monomorfisme escindit} \;\Longrightarrow\; \text{monomorfisme}.
\]
\end{proposicio}

\begin{prova}
La primera implicació és la proposició~\ref{prop:escindit-mono-epi}: si $f$ té una retracció, aleshores $f$ és un monomorfisme. Per a la segona, si $f$ és un isomorfisme, el seu invers $f^{-1}$ satisfà $f^{-1}\comp f=\id_A$, de manera que $f^{-1}$ és una retracció de $f$; per tant $f$ té retracció i és un monomorfisme escindit. La cadena d'implicacions se'n segueix.
\end{prova}

\subsection{Les propietats duals dels epimorfismes}

Aplicant el principi de dualitat a les proposicions anteriors ---invertint les fletxes--- obtenim, sense cap demostració nova, les propietats corresponents dels epimorfismes.

\begin{proposicio}[dual de la proposició~\ref{prop:mono-comp}]\label{prop:epi-comp}
La composició de dos epimorfismes és un epimorfisme.
\end{proposicio}

\begin{proposicio}[dual de la proposició~\ref{prop:mono-factor}]\label{prop:epi-factor}
Si la composició $g\comp f$ és un epimorfisme, aleshores $g$ és un epimorfisme.
\end{proposicio}

\begin{proposicio}[dual de la proposició~\ref{prop:iso-retr-mono}]\label{prop:iso-sec-epi}
Tot epimorfisme escindit és un epimorfisme, i tot isomorfisme és un epimorfisme escindit. Per tant,
\[
\text{isomorfisme} \;\Longrightarrow\; \text{epimorfisme escindit} \;\Longrightarrow\; \text{epimorfisme}.
\]
\end{proposicio}

Observem l'asimetria que revela la dualitat: a la proposició~\ref{prop:mono-factor}, d'un mono $g\comp f$ hereta el factor \emph{de la dreta}, $f$; a la seva dual~\ref{prop:epi-factor}, d'un epi hereta el factor \emph{de l'esquerra}, $g$.

\subsection{Condicions necessàries i suficients}

Les implicacions anteriors \emph{no} es reverteixen en general, i és aquí on convé anar amb compte. La cadena vàlida a qualsevol categoria és
\[
\text{iso} \;\Longrightarrow\; \text{mono escindit} \;\overset{(1)}{\Longrightarrow}\; \text{mono},
\]
i cap de les dues fletxes es pot invertir en general:
\begin{itemize}
\item[(1)] \textbf{mono $\not\Rightarrow$ mono escindit.} A $\cat{Grp}$, l'homomorfisme $\mathbb{Z}\xrightarrow{\,\cdot 2\,}\mathbb{Z}$ és un monomorfisme, però no té cap retracció: no és escindit.
\item[(2)] \textbf{mono escindit $\not\Rightarrow$ iso.} Un monomorfisme escindit no ha de ser exhaustiu en cap sentit: a $\cat{Set}$, la injecció $\{0\}\hookrightarrow\{0,1\}$ té retracció (l'única aplicació $\{0,1\}\to\{0\}$), és per tant un monomorfisme escindit, però no és un isomorfisme.
\end{itemize}
Dualment, $\text{iso}\Rightarrow\text{epi escindit}\Rightarrow\text{epi}$, amb les mateixes dues fletxes no reversibles. Convé no confondre \enquote{mono escindit $+$ epi} amb \enquote{mono $+$ epi}: un \textbf{bimorfisme} (alhora mono i epi) no és necessàriament un isomorfisme ---l'única fletxa no trivial de $\cat{2}$ ho és sense tenir invers, i a $\cat{Top}$ una bijecció contínua amb inversa no contínua també---, mentre que un morfisme alhora mono escindit i epi sí que ho és, com recull la caracterització següent:
\[
f \text{ és iso} \iff f \text{ és mono escindit i epi} \iff f \text{ és epi escindit i mono}.
\]
Les demostracions d'aquestes equivalències, i la de la cadena d'implicacions anterior, es proposen a les parts de pràctica i d'exercicis.

\section{Grafs i categories lliures}\label{sec:grafs-lliures}\index{categoria!lliure}

Els grafs i les categories estan estretament relacionats, però no són el mateix. Un graf dirigit només especifica vèrtexs i arestes. Una categoria, a més, disposa d'identitats, permet compondre fletxes i exigeix que aquestes composicions satisfacin els axiomes. De fet, per a una categoria petita, si oblidem la composició i quines fletxes són identitats, els objectes i morfismes formen un graf dirigit: és el seu \emph{graf subjacent}. Les fletxes identitat no desapareixen ---continuen sent arestes del graf, un bucle a cada vèrtex---, només que ja no estan marcades com a identitats. Cal que la categoria sigui petita, perquè els vèrtexs i les arestes d'un graf formen conjunts.

Hi ha una manera natural de fer el camí contrari: generar una categoria a partir d'un graf dirigit prenent com a morfismes tots els camins possibles.

\begin{definicio}\label{def:free}
Sigui $G$ un graf dirigit, amb conjunt de vèrtexs $V_G$, conjunt d'arestes $E_G$ i aplicacions origen i destí $s_G,t_G\colon E_G\to V_G$. Per a $A,B\in V_G$, un \textbf{camí}\index{camí} de $A$ a $B$ de longitud $n\ge 1$ és una successió $(e_1,\dots,e_n)$ d'arestes amb
\[
s_G(e_1)=A,\qquad t_G(e_i)=s_G(e_{i+1})\ \ (1\le i<n),\qquad t_G(e_n)=B.
\]
Per a cada vèrtex $A$ hi ha, a més, el \textbf{camí buit} $()_A$ de $A$ a $A$, de longitud zero.

La \textbf{categoria lliure generada pel graf} $G$, que escrivim $\Free(G)$,\index{Free@$\Free$}\index{categoria!lliure} té:
\begin{itemize}
\item objectes: $\Ob(\Free(G))=V_G$, els vèrtexs de $G$;
\item morfismes: $\Hom_{\Free(G)}(A,B)$ és el conjunt dels camins de $A$ a $B$;
\item composició: la \textbf{concatenació} en ordre de recorregut; per a $p=(a_1,\dots,a_r)\in\Hom(A,B)$ i $q=(b_1,\dots,b_s)\in\Hom(B,C)$,
\[
q\comp p=(a_1,\dots,a_r,b_1,\dots,b_s),
\]
on $p$, que surt de $A$, es recorre primer, d'acord amb la convenció $q\comp p$; concatenar amb un camí buit deixa el camí igual;
\item identitats: $\id_A=()_A$.
\end{itemize}
\end{definicio}

Per exemple, a partir del graf
\[
A\xrightarrow{f}B\xrightarrow{g}C
\]
la categoria lliure conté, a més de les identitats i dels camins $(f)$ i $(g)$ de longitud~1, que escrivim simplement $f$ i $g$, la fletxa composta $g\comp f=(f,g)\colon A\to C$:
\[
\begin{tikzpicture}[baseline=(m.center)]
  \matrix (m) [matrix of math nodes, column sep=3.5em] { A & B & C \\ };
  \draw[-{Stealth[scale=1.05]}] (m-1-1) -- node[above] {$f$} (m-1-2);
  \draw[-{Stealth[scale=1.05]}] (m-1-2) -- node[above] {$g$} (m-1-3);
  \draw[-{Stealth[scale=1.0]}] (m-1-1) to[bend right=25] node[below] {$g\comp f$} (m-1-3);
\end{tikzpicture}
\]
Els axiomes de categoria se satisfan. La composició és associativa, perquè $r\comp(q\comp p)$ i $(r\comp q)\comp p$ són totes dues la successió de totes les arestes de $p$, $q$ i $r$, en aquest ordre; i el camí buit és neutre, perquè concatenar-lo no afegeix cap aresta: $p\comp()_A=p=()_B\comp p$ per a tot camí $p$ de $A$ a $B$. Com que $V_G$ i $E_G$ són conjunts, també ho són els conjunts de camins, i $\Free(G)$ és una categoria petita. L'apèndix~\ref{cap:grafs} desenvolupa aquesta construcció amb més exemples. El nom \emph{lliure} quedarà justificat més endavant, quan disposem del llenguatge dels functors i puguem formular la propietat universal d'aquesta construcció.

\begin{exemple}[El monoide lliure i la convenció de composició]\label{ex:monoide-lliure-cat}
Prenem el cas d'un graf amb \textbf{un sol vèrtex} $\ast$ i conjunt d'arestes $E$. Aleshores totes les arestes són bucles ---comencen i acaben a $\ast$--- i cap parell no deixa d'estar encadenat, de manera que un camí és simplement una seqüència finita d'arestes, és a dir, una \emph{paraula} sobre $E$. Escrivim $E^*$ per al conjunt de totes aquestes paraules, incloent-hi la \emph{paraula buida} $\varepsilon$, de longitud zero.

Fixem tres convencions, que s'han de llegir juntes.
\begin{itemize}
\item \emph{Lectura de les paraules.} Escrivim les lletres en l'\emph{ordre de recorregut}: la paraula $a_1a_2\cdots a_n$ és el camí $(a_1,a_2,\dots,a_n)$, que recorre primer $a_1$, després $a_2$, i així successivament.
\item \emph{Concatenació.} Per a $u=a_1\cdots a_m$ i $v=b_1\cdots b_n$, la concatenació és la juxtaposició $uv=a_1\cdots a_m b_1\cdots b_n$: primer les lletres de $u$ i després les de $v$. Amb aquesta operació i la paraula buida, $(E^*,\text{concatenació},\varepsilon)$ és un monoide, el \textbf{monoide lliure}\index{monoide!lliure} sobre $E$: $(uv)w$ i $u(vw)$ són totes dues la paraula que juxtaposa $u$, $v$ i $w$ en aquest ordre, i $\varepsilon u=u=u\varepsilon$.
\item \emph{Composició.} A $\Free(G)$, la composició és la de la definició: $v\comp u$ vol dir \enquote{primer $u$, després $v$}, i el camí que en resulta recorre primer les lletres de $u$. Per tant,
\[
v\comp u=uv.
\]
\end{itemize}
Amb dues lletres $a,b\in E$, tot cap en una línia:
\[
b\comp a=ab\quad(\text{primer }a,\text{ després }b),
\qquad\qquad
a\comp b=ba\quad(\text{primer }b,\text{ després }a).
\]
L'ordre de la composició és, doncs, el \emph{contrari} del de la juxtaposició: el símbol de la dreta a $b\comp a$ és el que s'escriu primer a la paraula.

Quina categoria hem obtingut? Té un sol objecte, $\Hom_{\Free(G)}(\ast,\ast)=E^*$ i $\id_\ast=\varepsilon$. Però, amb la convenció de $\cat{B}M$, en què $g\comp f=g\cdot f$, la composició $v\comp u=uv$ correspon al producte \emph{invertit}. Si anomenem \textbf{monoide oposat}\index{monoide!oposat} $M\op$ d'un monoide $M$ el que té els mateixos elements i el mateix neutre i l'operació $m\cdot_{\mathrm{op}} n=n\cdot m$ ---de manera que $\cat{B}(M\op)=(\cat{B}M)\op$---, el que tenim és la \emph{igualtat}
\[
\Free(G)=\cat{B}\bigl((E^*)\op\bigr)=\bigl(\cat{B}(E^*)\bigr)\op,
\]
i no $\cat{B}(E^*)$. La diferència és només d'ordre d'escriptura. L'aplicació que inverteix l'ordre de les lletres, $a_1\cdots a_n\mapsto a_n\cdots a_1$, compleix $(uv)^{\mathrm{rev}}=v^{\mathrm{rev}}u^{\mathrm{rev}}$, i és per tant un isomorfisme de monoides $E^*\cong(E^*)\op$. Dona, doncs, un isomorfisme canònic de categories $\cat{B}(E^*)\cong\Free(G)$. Si haguéssim escrit les paraules en l'ordre de la composició ($a_n\cdots a_1$ per al camí que comença per $a_1$), hauríem obtingut literalment $\cat{B}(E^*)$. Hem preferit l'ordre de recorregut, que és el natural per als camins. En aquest sentit, i llevat d'aquest isomorfisme, el monoide lliure és el cas particular de la categoria lliure per a un graf amb un sol vèrtex.
\end{exemple}

\begin{exemple}
La categoria lliure sobre el graf amb dos vèrtexs $A,B$ i una sola aresta $A\to B$ és la categoria finita $\mathbf{2}$. Si, en canvi, el graf té dues arestes en sentits oposats, $e\colon A\to B$ i $f\colon B\to A$, la categoria lliure té infinits morfismes: $e$, $f$, $f\comp e$, $e\comp f$, $e\comp f\comp e$, i així successivament, perquè cap relació no identifica camins diferents. Així, \emph{un graf finit no genera necessàriament una categoria lliure finita}.
\end{exemple}

Aquesta construcció, que genera l'estructura més econòmica a partir d'unes dades sense imposar cap relació que no hi vingui forçada, és un exemple del patró \emph{lliure}, que reapareix per a monoides, grups i moltes altres estructures i que, al capítol~\refcap{cap:adjuncions}{6}, es formalitza com una \emph{adjunció} entre un functor lliure i un functor oblit.

\section{Categories segons la mida}\index{categoria!petita}\index{categoria!localment petita}
\label{sec:mida}

Hem parlat de \enquote{col·lecció} d'objectes en lloc de \enquote{conjunt}, i ara podem precisar per què. En alguns exemples, els objectes són massa nombrosos per formar un conjunt.

\begin{observacio}[La paradoxa de Russell]
La categoria $\cat{Set}$ té per objectes \emph{tots} els conjunts. Ara bé, no pot existir el conjunt de tots els conjunts. En efecte, si $U$ fos aquest conjunt, en podríem separar el subconjunt $R=\{x\in U\mid x\notin x\}$ dels conjunts que no es pertanyen a si mateixos, i preguntar-nos si $R\in R$: si $R\in R$, la definició força $R\notin R$; i si $R\notin R$, la definició força $R\in R$. Totes dues possibilitats són contradictòries. Aquesta és la \textbf{paradoxa de Russell}, i mostra que la col·lecció de tots els conjunts és massa gran per ser un conjunt.
\end{observacio}

Per això, en parlar dels objectes d'una categoria, diem \enquote{col·lecció} i no \enquote{conjunt}: pot ser massa gran. Segons la mida d'aquestes col·leccions introduïm els conceptes següents.

\begin{definicio}
Una categoria és \textbf{petita} si la seva col·lecció d'objectes i la seva col·lecció de morfismes són tots dos conjunts. En cas contrari es diu \textbf{gran}\index{categoria!gran}. I és \textbf{localment petita} si, per a cada parella d'objectes $A,B$, la col·lecció de morfismes $\Hom(A,B)$ és un conjunt, encara que la col·lecció de tots els objectes pugui ser massa gran per ser-ho.
\end{definicio}

Convé no llegir aquestes tres nocions com una partició en tres classes disjuntes. La partició real és \emph{petita} contra \emph{gran}. La petitesa local és una condició transversal: tota categoria petita és localment petita ---si tots els morfismes formen un conjunt, també en formen un els de cada parella $(A,B)$---, i una categoria gran pot ser localment petita o no. Així,
\[
\text{petita}\;\Longrightarrow\;\text{localment petita},
\]
i el recíproc és fals: $\cat{Set}$ és localment petita però no petita. Que no és petita ho mostra la paradoxa de Russell; que és localment petita es comprova identificant cada aplicació amb el seu graf (exercici~\ref{exr:pr-set-localment-petita} de la Pràctica). Les tres possibilitats efectives són, doncs, petita (per tant localment petita), gran localment petita i gran no localment petita.

Moltes de les categories habituals que utilitzem ---$\cat{Set}$, $\cat{Grp}$, $\cat{Top}$, $\cat{Graph}$ i altres--- són localment petites però no petites: tenen massa objectes per formar un conjunt, però entre dos objectes qualssevol hi ha només un conjunt de morfismes. En canvi, cada preordre concret sobre un conjunt i cada categoria finita concreta són categories petites. És important no confondre aquestes categories petites amb les categories grans que les apleguen, com $\cat{Pos}$, la categoria de tots els ordres parcials, o la categoria de totes les categories finites: en tots dos casos hi ha una classe pròpia de conjunts subjacents. Per la mateixa raó és gran $\cat{FinSet}$, tot i que cadascun dels seus objectes és un conjunt finit. Al llarg del llibre suposarem, tret que diguem el contrari, que les categories amb què treballem són localment petites; això garanteix que sempre podem parlar del conjunt $\Hom(A,B)$, cosa que serà essencial més endavant.\footnote{Aquí prenem com a base la teoria de conjunts clàssica, amb classes pròpies; l'apèndix~\ref{cap:mida} fixa aquest marc, i \cite[§1.8]{awodey}, \cite[§1.1]{riehl}, \cite[§3.2]{leinster} i \cite[cap.~1]{perrone} en donen presentacions introductòries. També és possible prendre el mateix llenguatge estructural com a fonament, sense una teoria de conjunts prèvia: l'apèndix~\ref{cap:axiomatica-relacional} n'ofereix un tast, i \cite[cap.~3]{leinster} en fa una introducció.}

\begin{resumcap}
\apartat{Idees centrals}
\begin{enumerate}
  \item Una categoria organitza objectes i morfismes mitjançant identitats i una composició associativa. El que importa no és què són els objectes, sinó com es relacionen a través de les fletxes. Els diagrames commutatius representen igualtats entre composicions.
  \item Una mateixa forma categòrica apareix en contextos molt diferents: conjunts i aplicacions, preordres, monoides, grafs, relacions i sistemes deductius.
  \item Els isomorfismes expressen quan dos objectes tenen la mateixa estructura dins la categoria; sovint, la noció estructuralment rellevant és l'isomorfisme més que no pas la igualtat estricta. Monomorfismes i epimorfismes es defineixen per cancel·lació. A $\cat{Set}$, i en algunes categories concretes, coincideixen amb les aplicacions injectives i exhaustives, però aquesta coincidència no és vàlida en general: en una categoria prima, per exemple, tot morfisme és mono i epi.
  \item Seccions i retraccions són inverses unilaterals. Un morfisme és un monomorfisme escindit si i només si té una retracció, i és un epimorfisme escindit si i només si té una secció. Tot monomorfisme escindit és mono i tot epimorfisme escindit és epi, però les implicacions inverses no valen en general:
  \[
  \text{iso}\Longrightarrow\text{mono escindit}\Longrightarrow\text{mono},
  \qquad
  \text{iso}\Longrightarrow\text{epi escindit}\Longrightarrow\text{epi}.
  \]
  \item A partir d'una categoria en podem construir d'altres: l'oposada, productes de categories, subcategories i categories sobre i sota un objecte.
  \item El \textbf{principi de dualitat} ---invertir totes les fletxes--- converteix sistemàticament cada definició i cada resultat en el seu dual, amb la composició en l'ordre invers.
  \item Un graf dirigit no té, per si sol, composició ni identitats; la categoria lliure hi afegeix formalment tots els camins i les identitats necessàries.
  \item Les qüestions de mida obliguen a distingir categories petites de categories localment petites. Tota categoria petita és localment petita, però moltes categories habituals, com $\cat{Set}$, són localment petites sense ser petites.
\end{enumerate}
\apartat{Errors típics} Creure que \enquote{mono $+$ epi $\Rightarrow$ iso} (fals a $\cat{Top}$ i a $\cat{2}$: un bimorfisme no és sempre un isomorfisme); confondre un monomorfisme amb un monomorfisme escindit; confondre la igualtat d'objectes amb l'isomorfisme, encara que sigui canònic; aplicar els adjectius \enquote{injectiu} i \enquote{exhaustiu} a morfismes d'una categoria general, en lloc de reservar-los per a aplicacions; raonar amb elements en una categoria on els objectes no en tenen o on no separen les fletxes; escriure la composició en l'ordre equivocat en dualitzar; i oblidar la hipòtesi de petitesa local en parlar d'homconjunts.
\apartat{Lectura recomanada} \cite[cap.~1, §2.1 i §3.1]{awodey} per a la introducció paral·lela, els monomorfismes i epimorfismes i el principi de dualitat; \cite[§1.1]{leinster} per a una presentació breu i moderna; \cite[cap.~I i II.1--II.3, II.7]{maclane} per a la formulació clàssica, inclosos la dualitat, les categories producte i les categories lliures; \cite[§1.3]{barrwells} per als grafs, en una presentació orientada a la informàtica; \cite[parts~I--II, sessions~2--10]{lawvereschanuel} per a una introducció encara més elemental als morfismes, els isomorfismes i les seccions i retraccions; i \cite[§1.1.5 i §1.2]{perrone-notes} per a contraexemples elementals d'isomorfismes, monomorfismes i epimorfismes, escindits o no. Per a la presentació relacional de les categories petites, l'apèndix~\ref{cap:axiomatica-relacional}.
\end{resumcap}

\chapter{Functors}
\label{cap:functors}

\begin{objectius}
\apartat{Prerequisits} El capítol~\ref{cap:categories} complet, especialment la definició de categoria, els diagrames commutatius, els isomorfismes, la categoria oposada i la dualitat, les categories producte i les categories sobre i sota un objecte, els grafs i les categories lliures, i la distinció entre categories petites i localment petites. Alguns exemples opcionals utilitzen grups, espais vectorials i l'espai dual $V^*=\Hom_{\mathbb K}(V,\mathbb K)$, i un exemple de functor oblit esmenta els espais topològics; cap d'ells no és imprescindible per seguir el capítol.
\apartat{Objectius} En acabar el capítol, el lector hauria de poder:
\begin{itemize}
  \item definir un functor especificant-ne l'acció sobre objectes i morfismes, i comprovar-ne la functorialitat;
  \item compondre functors i treballar amb la categoria $\cat{Cat}$ de les categories petites;
  \item reconèixer i construir exemples de functors en conjunts, preordres, lògica i grafs, inclosos el graf subjacent i la categoria lliure, i la relació entre $\Free$ i $U$;
  \item distingir functors covariants i contravariants i controlar-ne la variància;
  \item manejar els homfunctors $\Hom(A,-)$ i $\Hom(-,B)$ i el bihomfunctor $\Hom(-,-)$;
  \item distingir functors fidels, plens i essencialment exhaustius;
  \item separar amb precisió les propietats que un functor preserva de les que reflecteix o pot deixar de preservar.
\end{itemize}
\end{objectius}

\section{Cap al concepte de functor}

Un cop fixada la noció de categoria, la pregunta natural és quina mena de correspondència relaciona dues categories. Igual que entre grups mirem els homomorfismes ---les aplicacions que respecten l'estructura de grup--- i entre espais topològics mirem les aplicacions contínues ---les que respecten l'estructura topològica---, entre categories mirarem les correspondències que respecten l'estructura categòrica: objectes, morfismes, identitats i composició. S'anomenen \emph{functors}.

Fixem-nos que una categoria té dos nivells de dades: els objectes i els morfismes. Un functor haurà d'actuar sobre tots dos alhora, i fer-ho de manera compatible: ha d'enviar cada fletxa de la categoria $\cat{C}$ a una fletxa de la categoria $\cat{D}$ que vagi entre els objectes on van a parar el seu domini i el seu codomini, i ha de preservar la manera com les fletxes es componen i quines fletxes són identitats. Aquesta és exactament la idea que fa que els functors siguin, a la vegada, els morfismes d'una categoria de categories.

Abans de formalitzar-ho, val la pena veure un exemple concret que surt d'una operació ben familiar amb conjunts. Treballem a $\cat{Set}$ i fixem un conjunt $S$. A cada conjunt $A$ li podem associar el conjunt de totes les aplicacions de $S$ cap a $A$,
\[
F(A)=A^S=\{\,g\mid g\colon S\to A\,\}.
\]
Això assigna un objecte de $\cat{Set}$ a cada objecte de $\cat{Set}$; però un functor també ha d'actuar sobre les fletxes. I, de fet, tota aplicació $f\colon A\to B$ n'indueix una entre les potències,
\[
F(f)\colon A^S\longrightarrow B^S,
\qquad
F(f)(g)=f\comp g,
\]
és a dir, postcompondre amb $f$: si $g\colon S\to A$ és un element de $A^S$, aleshores $f\comp g\colon S\to B$ és un element de $B^S$, tal com mostra el triangle
\[
\begin{tikzpicture}[baseline=(m.center)]
  \matrix (m) [matrix of math nodes, column sep=3.4em, row sep=2.2em] {
    S & A \\
      & B \\
  };
  \draw[-{Stealth[scale=1.1]}] (m-1-1) -- node[above] {$g$} (m-1-2);
  \draw[-{Stealth[scale=1.1]}] (m-1-2) -- node[right] {$f$} (m-2-2);
  \draw[-{Stealth[scale=1.1]}] (m-1-1) -- node[below left] {$f\comp g$} (m-2-2);
\end{tikzpicture}
\]
Igual que a la definició, $F$ transporta cada fletxa de $\cat{Set}$ a una fletxa entre les imatges dels seus extrems:
\[
\begin{tikzpicture}[baseline=(m.center)]
  \matrix (m) [matrix of math nodes, row sep=2.6em] {
    A \\
    B \\
  };
  \draw[-{Stealth[scale=1.1]}] (m-1-1) -- node[right] {$f$} (m-2-1);
\end{tikzpicture}
\qquad\stackrel{F}{\longmapsto}\qquad
\begin{tikzpicture}[baseline=(m.center)]
  \matrix (m) [matrix of math nodes, row sep=2.6em] {
    A^S \\
    B^S \\
  };
  \draw[-{Stealth[scale=1.1]}] (m-1-1) -- node[right] {$F(f)$} (m-2-1);
\end{tikzpicture}
\]
on cada $g\in A^S$ és una aplicació $g\colon S\to A$ i $F(f)(g)=f\comp g$.
Aquesta assignació compleix just les dues condicions que esperaríem d'un functor. D'una banda, preserva la composició: si $A\xrightarrow{f}B\xrightarrow{h}C$, aleshores, per a tot $g\in A^S$,
\[
\begin{aligned}
F(h\comp f)(g)&=(h\comp f)\comp g\\
&=h\comp(f\comp g)=F(h)\bigl(F(f)(g)\bigr),
\end{aligned}
\]
de manera que
\[
F(h\comp f)=F(h)\comp F(f);
\]
l'única cosa que hem fet servir és l'associativitat de la composició. De l'altra, preserva les identitats: $F(\id_A)(g)=\id_A\comp g=g$, és a dir $F(\id_A)=\id_{A^S}$. Tenim, doncs, una manera sistemàtica de transformar objectes \emph{i} fletxes de $\cat{Set}$ que respecta composició i identitats: exactament el que a la secció següent anomenarem un functor. (Aquest exemple reapareixerà: $A^S$ no és res més que $\Hom_{\cat{Set}}(S,A)$, i el retrobarem com a homfunctor covariant a la definició~\ref{def:homfunctors}.)

\section{Definició de functor}

\begin{definicio}\index{functor}
\label{def:functor}
Siguin $\cat{C}$ i $\cat{D}$ dues categories. Un \textbf{functor} $F\colon\cat{C}\to\cat{D}$ assigna:
\begin{itemize}
\item a cada objecte $A$ de $\cat{C}$, un objecte $F(A)$ de $\cat{D}$;
\item a cada morfisme $f\colon A\to B$ de $\cat{C}$, un morfisme $F(f)\colon F(A)\to F(B)$ de $\cat{D}$,
\end{itemize}
de manera que es compleixin les dues condicions següents:
\begin{enumerate}
\item \textbf{Preservació d'identitats.} Per a cada objecte $A$ de $\cat{C}$,
\[
F(\id_A)=\id_{F(A)}.
\]
\item \textbf{Preservació de la composició.} Per a cada parella de morfismes componibles $f\colon A\to B$ i $g\colon B\to C$ de $\cat{C}$,
\[
F(g\comp f)=F(g)\comp F(f).
\]
\end{enumerate}
\end{definicio}

Observem que la condició sobre morfismes ja imposa que $F$ \emph{respecti} dominis i codominis: si $f\colon A\to B$, aleshores $F(f)$ ha d'anar de $F(A)$ a $F(B)$. La preservació de la composició es pot llegir diagramàticament: un triangle commutatiu de $\cat{C}$ es transforma en un triangle commutatiu de $\cat{D}$.
\[
\begin{tikzpicture}[baseline=(m.center)]
  \matrix (m) [matrix of math nodes, column sep=3.2em, row sep=2.6em] {
    A & B \\
      & C \\
  };
  \draw[-{Stealth[scale=1.1]}] (m-1-1) -- node[above] {$f$} (m-1-2);
  \draw[-{Stealth[scale=1.1]}] (m-1-2) -- node[right] {$g$} (m-2-2);
  \draw[-{Stealth[scale=1.1]}] (m-1-1) -- node[below left] {$g\comp f$} (m-2-2);
\end{tikzpicture}
\qquad\stackrel{F}{\longmapsto}\qquad
\begin{tikzpicture}[baseline=(m.center)]
  \matrix (m) [matrix of math nodes, column sep=3.6em, row sep=2.6em] {
    F(A) & F(B) \\
         & F(C) \\
  };
  \draw[-{Stealth[scale=1.1]}] (m-1-1) -- node[above] {$F(f)$} (m-1-2);
  \draw[-{Stealth[scale=1.1]}] (m-1-2) -- node[right] {$F(g)$} (m-2-2);
  \draw[-{Stealth[scale=1.1]}] (m-1-1) -- node[below left] {$F(g\comp f)$} (m-2-2);
\end{tikzpicture}
\]

\begin{observacio}[Primer els tipus, després les identitats]\label{obs:comprovar-tipus}\index{functor!comprovació de tipus}
Per comprovar que una assignació és un functor convé seguir sempre el mateix ordre, i no passar a les equacions fins que els tipus no quadren.
\begin{enumerate}
\item \emph{Objectes.} Cada $F(A)$ ha de ser un objecte de $\cat{D}$.
\item \emph{Tipus dels morfismes.} Per a cada $f\colon A\to B$, cal que $F(f)$ sigui un morfisme de $\cat{D}$ \emph{i} que vagi de $F(A)$ a $F(B)$. Aquí és on fallen més sovint les construccions candidates.
\item \emph{Identitats i composició.} Només aleshores té sentit comprovar $F(\id_A)=\id_{F(A)}$ i $F(g\comp f)=F(g)\comp F(f)$. El pas~2 garanteix que els dos membres de cada equació tenen el mateix tipus ($F(A)\to F(A)$ i $F(A)\to F(C)$) i que la composició $F(g)\comp F(f)$ existeix, perquè el codomini de $F(f)$ i el domini de $F(g)$ són tots dos $F(B)$.
\end{enumerate}
Un exemple mostra per què l'ordre importa. Si a cada conjunt $X$ li assignem $\mathcal P(X)$ i a cada aplicació $f\colon X\to Y$ l'antiimatge $B\mapsto f^{-1}(B)$, el pas~2 ja falla: l'antiimatge va de $\mathcal P(Y)$ a $\mathcal P(X)$, no de $\mathcal P(X)$ a $\mathcal P(Y)$. No cal, doncs, comprovar cap equació per saber que això no és un functor covariant $\cat{Set}\to\cat{Set}$. Sí que és, en canvi, un functor \emph{contravariant}, la noció que introduirem a la secció~\ref{sec:variancia}, per a la qual el pas~2 demana precisament $F(f)\colon F(B)\to F(A)$ (exemple~\ref{ex:parts-dues-variancies}). Les solucions de la Pràctica segueixen aquest ordre; vegeu, per exemple, l'exercici resolt~\ref{pr-hom-cov}.
\end{observacio}

\begin{observacio}[Els functors preserven els diagrames commutatius]\index{diagrama commutatiu!preservat per functors}
Una conseqüència directa de la preservació de la composició és que un functor envia diagrames commutatius a diagrames commutatius. En efecte, dir que un diagrama commuta és dir que certes composicions de fletxes coincideixen; si dos camins amb el mateix origen i el mateix destí donen la mateixa composició a $\cat{C}$, aleshores, aplicant $F$ i usant que preserva la composició, les seves imatges també coincideixen a $\cat{D}$. Per això podem aplicar un functor a tot un diagrama i seguir raonant amb la seva imatge: les igualtats entre camins es conserven. Ho comprovarem en detall, per al cas d'un quadrat, a l'exercici~\ref{prac:functor-quadrat} del capítol~\ref{cap:practica-functors} de Pràctica.
\end{observacio}

\subsection{Els functors preserven els isomorfismes}\label{sec:preserven-isos}

Com que els functors preserven tota l'estructura categòrica, en particular preserven les equacions entre composicions. Una conseqüència immediata és que preserven els isomorfismes.

\begin{proposicio}\label{prop:functor-preserva-iso}\index{functor!preserva isomorfismes}
Si $F\colon\cat{C}\to\cat{D}$ és un functor i $f\colon A\to B$ és un isomorfisme a $\cat{C}$, aleshores $F(f)$ és un isomorfisme a $\cat{D}$, i
\[
\bigl(F(f)\bigr)^{-1}=F(f^{-1}).
\]
\end{proposicio}

\begin{prova}
Sigui $g=f^{-1}$, de manera que $g\comp f=\id_A$ i $f\comp g=\id_B$. Aplicant $F$ i usant que preserva composició i identitats,
\[
F(g)\comp F(f)=F(g\comp f)=F(\id_A)=\id_{F(A)},
\]
\[
F(f)\comp F(g)=F(f\comp g)=F(\id_B)=\id_{F(B)}.
\]
Per tant $F(g)$ és un invers de $F(f)$, i pel resultat d'unicitat de l'invers, $\bigl(F(f)\bigr)^{-1}=F(g)=F(f^{-1})$.
\end{prova}

\begin{observacio}
Val la pena veure per què és així, més enllà del càlcul de la prova. Dir que $f$ és un isomorfisme és dir que existeix $g$ amb $g\comp f=\id_A$ i $f\comp g=\id_B$: dues equacions entre composicions, és a dir, dos diagrames commutatius. Com que els functors preserven els diagrames commutatius, aquestes dues equacions es transporten a $\cat{D}$ en aplicar-hi $F$, i $F(g)$ resulta ser l'invers de $F(f)$.
\end{observacio}

\section{Exemples de functors}\label{sec:primers-exemples}

La força de la definició, com en el cas de les categories, es veu en la varietat d'exemples que hi encaixen. Recuperem els exemples de referència del capítol anterior.

\begin{exemple}[Functor identitat]\index{functor!identitat}
Per a tota categoria $\cat{C}$ hi ha el \textbf{functor identitat} $\id_{\cat{C}}\colon\cat{C}\to\cat{C}$, que deixa fixos tots els objectes i tots els morfismes. Preserva trivialment identitats i composició.
\end{exemple}

\begin{exemple}[Functor constant]\label{ex:functor-constant}\index{functor!constant}
Siguin $\cat{C}$ i $\cat{D}$ dues categories i $D$ un objecte de $\cat{D}$. El \textbf{functor constant} amb valor $D$,
\[
\Delta_D\colon\cat{C}\to\cat{D},
\]
envia tot objecte $A$ de $\cat{C}$ a $D$ i tot morfisme $f$ de $\cat{C}$ a la identitat $\id_D$. És un functor: $\Delta_D(\id_A)=\id_D=\id_{\Delta_D(A)}$, i $\Delta_D(g\comp f)=\id_D=\id_D\comp\id_D=\Delta_D(g)\comp\Delta_D(f)$. Quan $\cat{D}=\cat{Set}$ i $D$ és un conjunt d'un sol element $\{\ast\}$, escrivim $\Delta_{\{\ast\}}$ per al functor constant que envia cada objecte a aquest conjunt d'un sol element.
\end{exemple}

\begin{exemple}[Functors oblit]\index{functor!oblit}
Moltes categories d'objectes amb estructura $\cat{C}$ tenen un \textbf{functor oblit} $U\colon\cat{C}\to\cat{Set}$, que oblida l'estructura i reté només el conjunt subjacent:
\[
\cat{Grp}\to\cat{Set},
\qquad
\cat{Top}\to\cat{Set},
\qquad
\cat{Vect}_{\mathbb K}\to\cat{Set}.
\]
En cada cas, $U$ envia un objecte al seu conjunt subjacent i un morfisme ---un homomorfisme, una aplicació contínua, una aplicació lineal--- a la mateixa aplicació, vista simplement com a aplicació de conjunts. La identitat es preserva perquè l'aplicació identitat continua sent la identitat, i la composició es preserva perquè és la mateixa composició d'aplicacions. El functor oblit \enquote{no fa res} als morfismes; el seu contingut és que oblida l'estructura dels objectes.
\end{exemple}

\begin{exemple}[Functors des de les categories finites]\index{functor!des d'una categoria finita}
Les categories finites $\cat{0}$, $\cat{1}$ i $\cat{2}$ (subsecció~\ref{subsec:finites}) donen, com a dominis, functors amb una lectura molt directa. Un functor $\cat{1}\to\cat{D}$, on $\cat{1}$ té un únic objecte i cap morfisme no trivial, no és res més que la tria d'un objecte de $\cat{D}$: la imatge de l'únic objecte. Un functor $\cat{2}\to\cat{D}$, on $\cat{2}=(0\to 1)$ ---reanomenem els dos objectes de $\cat{2}$ com $0$ i $1$, i escrivim $(0\to 1)$ per indicar la seva única fletxa no trivial---, és la tria d'un morfisme de $\cat{D}$: cal dir on van $0$ i $1$ ---dos objectes--- i on va l'única fletxa no trivial ---un morfisme entre les seves imatges---, i la preservació d'identitats i composició no imposa res més. Finalment, des de la categoria buida $\cat{0}$ hi ha un únic functor $\cat{0}\to\cat{D}$, el functor buit, sigui quina sigui $\cat{D}$. Aquests tres fets ---objectes, morfismes i \enquote{res}--- reapareixeran contínuament: seleccionar un objecte o un morfisme d'una categoria és el mateix que donar un functor des de $\cat{1}$ o des de $\cat{2}$.
\end{exemple}

\begin{exemple}[Functors entre categories discretes]\index{functor!entre categories discretes}\index{categoria!discreta}
Un conjunt es pot veure com una categoria discreta, en què els únics morfismes són les identitats (subsecció~\ref{subsec:conjunts}). Si $S$ i $T$ són dos conjunts vistos així, un functor $F\colon S\to T$ només ha de dir on va cada objecte, perquè els únics morfismes són les identitats i tot functor les preserva automàticament. Per tant un functor entre categories discretes és exactament una aplicació de conjunts $S\to T$: sense estructura a preservar, la noció de functor es redueix a la d'aplicació.
\end{exemple}

\begin{exemple}[Functors entre preordres]\index{functor!entre preordres}
Siguin $(P,\leq)$ i $(Q,\leq)$ dos preordres, vistos com a categories. Un functor $F\colon P\to Q$ assigna a cada element $x\in P$ un element $F(x)\in Q$ i, com que hi ha un morfisme $x\to y$ exactament quan $x\leq y$, la condició sobre morfismes diu que
\[
x\leq y
\quad\Longrightarrow\quad
F(x)\leq F(y).
\]
És a dir, un functor entre preordres és exactament una aplicació monòtona. Les condicions d'identitat i composició es compleixen automàticament, perquè en un preordre entre dos objectes hi ha com a màxim una fletxa i no hi ha res més a comprovar.
\end{exemple}

\begin{exemple}[Functors des d'un monoide]\label{ex:functors-monoide}\index{functor!des d'un monoide}
Recordem que un monoide $(M,\cdot,1)$ es pot veure com una categoria d'un sol objecte $\ast$, que escrivim $\cat{B}M$. Un functor $F\colon\cat{B}M\to\cat{D}$ tria un objecte $D=F(\ast)$ de $\cat{D}$ i envia cada element $m\in M$ ---és a dir, cada morfisme $\ast\to\ast$--- a un morfisme $F(m)\colon D\to D$, de manera que
\[
F(1)=\id_D
\qquad\text{i}\qquad
F(m\cdot n)=F(m)\comp F(n).
\]
Aquestes són exactament les condicions perquè $F$ sigui un \textbf{homomorfisme de monoides} de $M$ cap al monoide dels endomorfismes de $D$,
\[
\mathrm{End}_{\cat{D}}(D)=\Hom_{\cat{D}}(D,D).
\]
En particular, quan $\cat{D}=\cat{Set}$, un functor $\cat{B}M\to\cat{Set}$ és una \textbf{acció} del monoide $M$ sobre un conjunt. Vegem per què. Un tal functor queda determinat per un conjunt $F(\ast)=S$ i, per a cada element $m\in M$, una funció $F(m)\colon S\to S$. Les dues condicions de functor ---preservar identitats i composició--- diuen exactament que, per a tot $x\in S$ i tots $m,n\in M$,
\[
F(1)(x)=x,
\qquad
F(m\cdot n)(x)=F(m)\bigl(F(n)(x)\bigr),
\]
Si escrivim $m\cdot x$ en lloc de $F(m)(x)$, aquestes igualtats esdevenen les condicions habituals d'una \emph{acció per l'esquerra} del monoide sobre $S$:
\[
1\cdot x=x,
\qquad
(m\cdot n)\cdot x=m\cdot(n\cdot x).
\]
\end{exemple}

\begin{exemple}[Functors des d'un grup: accions i representacions]\index{functor!des d'un grup}\index{accio@acció!de grup}\index{representacio@representació lineal}
Quan el monoide és un grup $G$, la categoria $\cat{B}G$ és un grupoide: tot morfisme és un isomorfisme. Un functor $F\colon\cat{B}G\to\cat{D}$ és aleshores un homomorfisme de $G$ cap al grup dels automorfismes de l'objecte $D=F(\ast)$,
\[
F\colon G\longrightarrow \Aut_{\cat{D}}(D),
\]
perquè un functor preserva isomorfismes (proposició~\ref{prop:functor-preserva-iso}) i, per tant, cada element de $G$ va a parar a un automorfisme de $D$. Dos casos són especialment importants:
\begin{itemize}
\item Un functor $\cat{B}G\to\cat{Set}$ és una \textbf{acció} del grup $G$ sobre un conjunt ---un \emph{$G$-conjunt}. El desplegament és el mateix que per a un monoide (exemple~\ref{ex:functors-monoide}), amb una novetat: com que ara cada $a\in G$ té un invers $a^{-1}$, la funció $F(a)\colon S\to S$ és bijectiva, amb inversa $F(a^{-1})$. Cada element del grup actua, doncs, com una \emph{permutació} de $S$. Un exemple típic són les simetries d'una figura, que actuen permutant-ne els vèrtexs.
\item Un functor $\cat{B}G\to\cat{Vect}_{\mathbb K}$ és una \textbf{representació lineal} de $G$ sobre el cos $\mathbb{K}$. El desplegament torna a ser el mateix, ara amb un espai vectorial $F(\ast)=V$ en lloc d'un conjunt: cada element $a\in G$ dona una aplicació \emph{lineal} invertible $F(a)\colon V\to V$ ---un automorfisme de $V$, amb inversa $F(a^{-1})$---, i el producte del grup es correspon amb la composició. Escrivint $a\cdot v$ per $F(a)(v)$, la novetat respecte d'una acció ordinària és que cada element actua respectant les operacions de l'espai:
\[
a\cdot(\lambda v+\mu w)=\lambda\,(a\cdot v)+\mu\,(a\cdot w).
\]
En dimensió finita, fixada una base, cada $F(a)$ és una matriu invertible: una representació \enquote{realitza} els elements abstractes del grup com a matrius concretes, cosa que permet estudiar el grup amb les eines de l'àlgebra lineal.
\end{itemize}
Així, dues teories aparentment allunyades ---les accions i representacions lineals d'un grup--- són, categòricament, el mateix tipus d'objecte: functors amb domini $\cat{B}G$, que només es distingeixen per la categoria d'arribada.
\end{exemple}

\begin{exemple}[Deduïbilitat i functors lògics]\label{ex:deduibilitat-functors}\index{functor!i deduïbilitat}\index{substitucio@substitució}
Siguin $T$ i $T'$ dos sistemes deductius, i $\cat{Prov}_T$ i $\cat{Prov}_{T'}$ les seves categories primes de deduïbilitat (subsecció~\ref{subsec:elementals}). Una assignació de fórmules $\varphi\colon A\mapsto \varphi(A)$ que preservi la deduïbilitat,
\[
A\vdash_T B
\quad\Longrightarrow\quad
\varphi(A)\vdash_{T'}\varphi(B),
\]
és exactament un functor $F\colon\cat{Prov}_T\to\cat{Prov}_{T'}$, amb $F(A)=\varphi(A)$: la fletxa que expressa $A\vdash_T B$ va a parar a la fletxa que expressa $\varphi(A)\vdash_{T'}\varphi(B)$. Com que aquestes categories són primes, no cal comprovar res més: la preservació de la deduïbilitat garanteix automàticament la preservació d'identitats i composicions. Aquest és el reflex categòric del fet que una traducció entre teories que respecti la deduïbilitat indueix un functor entre les categories associades.

Un cas concret i molt natural d'aquesta situació és el d'una \textbf{substitució} de variables proposicionals. Sigui $T_{\mathrm{prop}}$ el sistema deductiu de la lògica proposicional i $\cat{Prov}_{T_{\mathrm{prop}}}$ la seva categoria de deduïbilitat, amb les fórmules per objectes i una fletxa $A\to B$ exactament quan $A\vdash B$ (aquí escrivim $\vdash$ per $\vdash_{T_{\mathrm{prop}}}$). Considerem la substitució
\[
\sigma\colon\quad p\mapsto q,\qquad q\mapsto p\lor r,
\]
estesa a totes les fórmules respectant els connectors ($\land,\lor,\neg,\to$). Per exemple, $\sigma$ transforma $p\to(q\land p)$ en $q\to((p\lor r)\land q)$. El punt clau és que una substitució no sols transforma fórmules, sinó que també preserva la \emph{deduïbilitat}: si $A\vdash B$, aleshores $\sigma(A)\vdash\sigma(B)$, perquè qualsevol prova concreta de $B$ a partir de $A$ es tradueix, aplicant $\sigma$ a cada fórmula, en una derivació de $\sigma(B)$ a partir de $\sigma(A)$. Així, de la fletxa $p\land q\to p$ (perquè $p\land q\vdash p$) n'obtenim la fletxa $q\land(p\lor r)\to q$. Per tant $\sigma$ indueix un functor
\[
F_\sigma\colon\cat{Prov}_{T_{\mathrm{prop}}}\longrightarrow\cat{Prov}_{T_{\mathrm{prop}}},
\qquad
A\mapsto\sigma(A),
\qquad
(A\vdash B)\mapsto(\sigma(A)\vdash\sigma(B)),
\]
que és, precisament, el cas $T=T'=T_{\mathrm{prop}}$ de la construcció anterior (aquí amb la mateixa categoria de sortida i d'arribada).
\end{exemple}

\begin{observacio}
Si en comptes de la deduïbilitat considerem la \emph{categoria de proves} de $T$ ---els mateixos objectes, però ara amb les derivacions, mòdul equivalència, com a morfismes i el tall com a composició---, un functor cap a la categoria de proves de $T'$ demana també una traducció de cada derivació que en respecti l'encadenament i les identitats. A diferència del cas anterior, aquestes condicions ja no són automàtiques, perquè la categoria de proves no és prima.
\end{observacio}

\begin{exemple}[Functors oblit sobre i sota un objecte]\index{functor!oblit}\index{categoria!sobre un objecte}\index{categoria!sota un objecte}
Les categories de $\cat{C}$ sobre i sota un objecte $A$, $\cat{C}/A$ i $A/\cat{C}$ (definició~\ref{def:slice}), tenen un functor oblit cap a $\cat{C}$, un per a cadascuna. El functor
\[
U_A\colon\cat{C}/A\to\cat{C}
\]
envia cada objecte $x\colon X\to A$ al seu domini $X$ i cada morfisme (triangle) $h$ a ell mateix; oblida el morfisme distingit cap a $A$ i reté només el domini. Dualment, $A/\cat{C}\to\cat{C}$ envia cada objecte $x\colon A\to X$ al seu codomini $X$ i cada morfisme (triangle) a ell mateix; oblida el morfisme distingit des de $A$ i reté només el codomini.

Vegem-ho en un cas concret. Amb $\cat{2}=(0\to 1)$ i l'objecte $A=1$, els objectes de $\cat{2}/1$ són les fletxes de $\cat{2}$ amb codomini $1$: la identitat $\id_1$ i l'única fletxa no trivial $u\colon 0\to 1$. L'únic morfisme no trivial va de $u$ a $\id_1$, de manera que $\cat{2}/1$ és, de nou, una còpia de $\cat{2}$, i $U_1\colon\cat{2}/1\to\cat{2}$ envia $u\mapsto 0$ i $\id_1\mapsto 1$. Simètricament, sota l'objecte $0$, els objectes de $0/\cat{2}$ són les fletxes amb domini $0$ ---la identitat $\id_0$ i $u\colon 0\to 1$---, i el functor oblit $0/\cat{2}\to\cat{2}$ envia $\id_0\mapsto 0$ i $u\mapsto 1$.
\end{exemple}

\begin{exemple}[Projeccions i functor diagonal]\label{ex:projeccions-diagonal}\index{functor!de projecció}\index{functor!diagonal}
La categoria producte $\cat{C}\times\cat{D}$ (subsecció~\ref{subsec:categoria-producte}) té dos functors de \textbf{projecció}
\[
\pi_1\colon\cat{C}\times\cat{D}\to\cat{C},\qquad \pi_2\colon\cat{C}\times\cat{D}\to\cat{D},
\]
que es queden amb una de les components: $\pi_1(A,X)=A$ i $\pi_1(f,g)=f$, i anàlogament per a $\pi_2$. Són functors perquè la composició i les identitats de $\cat{C}\times\cat{D}$ es defineixen component a component:
\[
\begin{aligned}
\pi_1\bigl((f',g')\comp(f,g)\bigr)
 &= \pi_1(f'\comp f,\,g'\comp g)\\
 &= f'\comp f
  = \pi_1(f',g')\comp\pi_1(f,g),\\
\pi_1(\id_A,\id_X)&=\id_A,
\end{aligned}
\]
i el càlcul per a $\pi_2$ és el mateix sobre la segona component.
En sentit contrari, cada categoria $\cat{C}$ té un functor \textbf{diagonal}
\[
\Delta\colon\cat{C}\to\cat{C}\times\cat{C},\qquad \Delta(A)=(A,A),\qquad \Delta(f)=(f,f),
\]
que duplica objectes i morfismes; la functorialitat és de nou immediata, perquè $\Delta(g\comp f)=(g\comp f,\,g\comp f)=(g,g)\comp(f,f)$ i $\Delta(\id_A)=(\id_A,\id_A)$. Les projeccions i la diagonal encaixen en el diagrama commutatiu
\[
\begin{tikzpicture}[baseline=(m.center)]
  \matrix (m) [matrix of math nodes, row sep=2.6em, column sep=3.4em] {
    & \cat{C} & \\
    \cat{C} & \cat{C}\times\cat{C} & \cat{C} \\
  };
  \draw[-{Stealth[scale=1.1]}] (m-1-2) -- node[right] {$\Delta$} (m-2-2);
  \draw[-{Stealth[scale=1.1]}] (m-2-2) -- node[below] {$\pi_1$} (m-2-1);
  \draw[-{Stealth[scale=1.1]}] (m-2-2) -- node[below] {$\pi_2$} (m-2-3);
  \draw[-{Stealth[scale=1.1]}] (m-1-2) -- node[above left] {$\id_{\cat{C}}$} (m-2-1);
  \draw[-{Stealth[scale=1.1]}] (m-1-2) -- node[above right] {$\id_{\cat{C}}$} (m-2-3);
\end{tikzpicture}
\qquad
\pi_1\comp\Delta=\id_{\cat{C}}=\pi_2\comp\Delta,
\]
on la composició de functors es fa com a la secció següent: les dues igualtats es comproven directament, ja que $\pi_1(\Delta(A))=\pi_1(A,A)=A$ i $\pi_1(\Delta(f))=\pi_1(f,f)=f$, i igual per a $\pi_2$. La notació concorda amb la del functor constant $\Delta_D$ de l'exemple~\ref{ex:functor-constant}: més endavant, el functor diagonal s'estendrà a un functor que envia cada objecte $A$ a un diagrama constant de valor $A$, del qual $(A,A)$ és el cas de dues còpies. Són exemples elementals, però reapareixeran sovint: les projeccions fan de $\cat{C}\times\cat{D}$ el producte de $\cat{C}$ i $\cat{D}$ dins de $\cat{Cat}$ quan totes dues són petites (observació~\ref{obs:dos-productes}), i el functor diagonal és l'ingredient bàsic per tractar productes, límits i adjuncions en capítols posteriors.
\end{exemple}

\section{Composició de functors i la categoria \texorpdfstring{$\cat{Cat}$}{Cat}}\index{Cat@$\cat{Cat}$}

Els functors es poden compondre, i la composició torna a ser un functor. Això és el que permet organitzar les categories i els functors en una nova categoria.

\begin{proposicio}
Siguin $F\colon\cat{C}\to\cat{D}$ i $G\colon\cat{D}\to\cat{E}$ dos functors. La seva \textbf{composició} $G\comp F$, definida per
\[
(G\comp F)(A)=G(F(A))
\qquad\text{i}\qquad
(G\comp F)(f)=G(F(f)),
\]
és un functor $\cat{C}\to\cat{E}$.
\end{proposicio}

\begin{prova}
Cal comprovar les dues condicions. Per a les identitats, i usant primer que $F$ i després que $G$ les preserven,
\[
(G\comp F)(\id_A)=G(F(\id_A))=G(\id_{F(A)})=\id_{G(F(A))}=\id_{(G\comp F)(A)}.
\]
Per a la composició, siguin $f\colon A\to B$ i $g\colon B\to C$ a $\cat{C}$. Aleshores
\[
\begin{aligned}
(G\comp F)(g\comp f)
&=G\big(F(g\comp f)\big)
=G\big(F(g)\comp F(f)\big)\\
&=G(F(g))\comp G(F(f))
=(G\comp F)(g)\comp(G\comp F)(f).
\end{aligned}
\]
Per tant $G\comp F$ és un functor.
\end{prova}

La composició de functors és associativa i el functor identitat actua com a element neutre, exactament pels mateixos motius que a $\cat{Set}$. En efecte, un functor $F\colon\cat{C}\to\cat{D}$ consta de dues aplicacions ---una sobre objectes i una sobre morfismes---, i la composició de functors no és res més que la composició d'aquestes aplicacions component a component. Per tant, l'associativitat i les lleis d'identitat es redueixen a les propietats corresponents de la composició d'aplicacions, que ja coneixem. Això ens permet donar la definició següent.

\begin{definicio}\index{Cat@$\cat{Cat}$}
La categoria $\cat{Cat}$ té per objectes les categories petites i per morfismes els functors entre elles. La composició és la composició de functors i la identitat de cada categoria és el seu functor identitat.
\end{definicio}

\begin{observacio}[Qüestions de mida]\label{obs:mida-cat}
La categoria $\cat{Cat}$ té per objectes les categories \emph{petites}. Aquesta restricció la fa \emph{localment petita}: si $\cat{C}$ és petita, $\Ob(\cat{C})$ i $\Mor(\cat{C})$ són conjunts, i doncs, donades dues categories petites $\cat{C}$ i $\cat{D}$, el producte $\Ob(\cat{D})^{\Ob(\cat{C})}\times \Mor(\cat{D})^{\Mor(\cat{C})}$ és un conjunt; com que
\[
\Hom_{\cat{Cat}}(\cat{C},\cat{D})\subseteq \Ob(\cat{D})^{\Ob(\cat{C})}\times \Mor(\cat{D})^{\Mor(\cat{C})},
\]
també ho és $\Hom_{\cat{Cat}}(\cat{C},\cat{D})$. La mateixa $\cat{Cat}$, en canvi, és \emph{gran}: les categories petites formen una classe pròpia, no un conjunt.

Categories d'ús constant com $\cat{Set}$, $\cat{Grp}$ o $\cat{Top}$ són grans i, per tant, no són objectes de $\cat{Cat}$. Això no impedeix treballar amb functors entre elles: un functor oblit $\cat{Grp}\to\cat{Set}$, o el functor lliure $\Free\colon\cat{Graph}\to\cat{Cat}$, és simplement una assignació sobre objectes i una altra sobre morfismes que compleixen les condicions functorials, i la seva legitimitat no depèn de cap categoria que tingui les categories grans per objectes.

Al llarg del text, doncs, $\cat{Cat}$ designarà sempre la categoria de les categories \emph{petites}, i les categories grans, com $\cat{Set}$, no en són objectes.\footnotemark
\end{observacio}
\footnotetext{Aquesta afirmació és relativa al marc de mida adoptat (apèndix~\ref{cap:mida}). Amb una jerarquia d'universos de Grothendieck, una categoria gran respecte d'un univers pot ser petita respecte d'un univers més gran, i aleshores és objecte de la categoria de categories d'aquest nivell; vegeu \cite[§1.1]{riehl}.}%

Un cop els functors són els morfismes de $\cat{Cat}$, podem aplicar a les categories les mateixes nocions categòriques que ja coneixem. En particular, un isomorfisme a $\cat{Cat}$ és un functor que té un invers estricte.

\begin{definicio}[Isomorfisme de categories]\label{def:iso-categories}\index{isomorfisme!de categories}\index{categoria!isomorfa}
Dues categories $\cat{C}$ i $\cat{D}$ són \textbf{isomorfes}, i escrivim $\cat{C}\cong\cat{D}$, si hi ha functors $F\colon\cat{C}\to\cat{D}$ i $G\colon\cat{D}\to\cat{C}$ inversos l'un de l'altre, és a dir, amb $G\comp F=\id_{\cat{C}}$ i $F\comp G=\id_{\cat{D}}$; un tal $F$ és un \textbf{isomorfisme de categories}. Per a categories petites no és res més que un isomorfisme (definició~\ref{def:isomorfisme}) a $\cat{Cat}$, però la definició val igualment per a categories grans, que no són objectes de $\cat{Cat}$.
\end{definicio}

\section{Functors i grafs}\label{sec:functors-grafs}\index{functor!i grafs}\index{morfisme de grafs}

Val la pena comparar els functors amb els \emph{morfismes de grafs}, que vam introduir en definir $\cat{Graph}$ (subsecció~\ref{subsec:grafs}). Un morfisme de grafs $\alpha\colon G\to H$ és una parella d'aplicacions $\alpha_V\colon V_G\to V_H$ i $\alpha_E\colon E_G\to E_H$ compatibles amb les dues aplicacions d'origen i destí $s,t\colon E\to V$ de cada graf; res més, perquè un graf no té composició ni identitats. Un functor entre categories petites demana exactament aquestes mateixes dues dades ---una assignació sobre objectes i una sobre morfismes, compatibles amb domini i codomini--- però hi afegeix dues condicions que un graf no pot ni formular: la preservació de la composició i de les identitats. Un morfisme de grafs és, doncs, la part «sense composició ni identitats» d'un functor. Aquesta relació es concreta en dues construccions en sentits contraris entre $\cat{Graph}$ i $\cat{Cat}$, que estudiem tot seguit.

\subsection{El graf subjacent d'una categoria}

Tota categoria petita $\cat{C}$ té un \textbf{graf subjacent}: el graf que té per vèrtexs els objectes de $\cat{C}$ i per arestes \emph{tots} els seus morfismes, amb les aplicacions origen i destí donades pel domini i el codomini. En passar a aquest graf, oblidem quins morfismes són identitats i com es componen, però no eliminem les arestes corresponents: les fletxes identitat hi resten com a arestes, una a cada vèrtex, tot i que ja no marcades com a identitats.

Aquesta assignació és functorial. Donat un functor $F\colon\cat{C}\to\cat{D}$, obtenim un morfisme entre els grafs subjacents que actua, sobre vèrtexs, com $F$ sobre els objectes, i sobre arestes, com $F$ sobre els morfismes. És un morfisme de grafs ben definit precisament perquè $F$ respecta domini i codomini: una aresta $f\colon A\to B$ va a $F(f)\colon F(A)\to F(B)$, una aresta entre les imatges dels extrems. Per construir-lo n'hi ha prou amb això; la preservació de la composició i de les identitats no s'hi fa servir, i per això diem que «se n'oblida». Tenim, així, un functor oblit
\[
U\colon\cat{Cat}\to\cat{Graph},
\]
que envia cada categoria petita al seu graf subjacent $U(\cat{C})$ i cada functor $F$ al morfisme de grafs $U(F)$ que acabem de descriure. El recíproc del que fa $U$ és fals: un morfisme entre els grafs subjacents no respecta, en general, ni la composició ni les identitats, de manera que un functor és estrictament més que un morfisme de grafs. Vegem-ho amb un bucle. Sigui $M=(\{1,e\},\cdot,1)$ amb $e\cdot e=e$, vist com a categoria d'un sol objecte $\ast$; el functor oblit en dona el graf subjacent:
\[
\begin{tikzpicture}[baseline=(m.center)]
  \matrix (m) [cat matrix] { \ast \\ };
  \draw[cat] (m-1-1) to[out=55,in=125,looseness=9] node[above] {\(\id_\ast\)} (m-1-1);
  \draw[cat] (m-1-1) to[out=-125,in=-55,looseness=9] node[below] {\(e\)} (m-1-1);
\end{tikzpicture}
\qquad\stackrel{U}{\longmapsto}\qquad
\begin{tikzpicture}[baseline=(m.center)]
  \matrix (m) [cat matrix] { \ast \\ };
  \draw[cat] (m-1-1) to[out=55,in=125,looseness=9] node[above] {\(\id_\ast\)} (m-1-1);
  \draw[cat] (m-1-1) to[out=-125,in=-55,looseness=9] node[below] {\(e\)} (m-1-1);
\end{tikzpicture}
\]
Sobre aquest graf ---que conserva les dues arestes $\id_\ast$ i $e$, però on $U$ ha oblidat tota la composició: les igualtats $\id_\ast\comp\id_\ast=\id_\ast$, $\id_\ast\comp e=e$, $e\comp\id_\ast=e$ i $e\comp e=e$ (les tres que involucren $\id_\ast$ són el que en feia la identitat)--- considerem l'assignació $\alpha$ que fixa el vèrtex i envia tots dos bucles a $e$:
\[
\begin{tikzpicture}[baseline=(m.center)]
  \matrix (m) [cat matrix] { \ast \\ };
  \draw[cat] (m-1-1) to[out=55,in=125,looseness=9] node[above] {\(\id_\ast\)} (m-1-1);
  \draw[cat] (m-1-1) to[out=-125,in=-55,looseness=9] node[below] {\(e\)} (m-1-1);
\end{tikzpicture}
\qquad\stackrel{\alpha}{\longmapsto}\qquad
\begin{tikzpicture}[baseline=(m.center)]
  \matrix (m) [cat matrix] { \ast \\ };
  \draw[cat] (m-1-1) to[out=55,in=125,looseness=9] node[above] {\(e\)} (m-1-1);
\end{tikzpicture}
\]
Com a morfisme de grafs, $\alpha$ és vàlid: respecta origen i destí, perquè tots els bucles van del vèrtex a si mateix. Però no prové de cap functor: si fos $\alpha=U(F)$, el functor $F$ enviaria $\id_\ast$ a $e$, i cap functor no pot fer-ho ---ha d'enviar $\id_\ast$ a una identitat, i $e$ no ho és---. Els morfismes de grafs són, doncs, més abundants que els functors. Aquest mateix contraexemple, des de l'angle de la definició de functor, es reprèn a l'exercici~\ref{ex:identitat-necessaria}.

\subsection{La categoria lliure sobre un graf}

En sentit contrari, la construcció de la categoria lliure generada per un graf, $\Free(G)$, que vam definir a la secció~\ref{sec:grafs-lliures}, també és functorial. Recordem-ne l'acció sobre objectes: $\Free(G)$ té per objectes els vèrtexs de $G$ i per morfismes els camins d'arestes, amb la concatenació com a composició i el camí buit com a identitat de cada vèrtex.

Vegem l'acció sobre morfismes. Sigui $\alpha\colon G\to H$ un morfisme de grafs, amb components $\alpha_V\colon V_G\to V_H$ i $\alpha_E\colon E_G\to E_H$. Definim un functor
\[
\Free(\alpha)\colon\Free(G)\to\Free(H)
\]
que, sobre objectes, actua com $\alpha_V$, i sobre morfismes transforma cada camí aresta a aresta: el camí $(f_1,\dots,f_k)$ de $G$ va al camí $(\alpha_E(f_1),\dots,\alpha_E(f_k))$ de $H$, i el camí buit en $v$, $()_v$, va al camí buit en $\alpha_V(v)$. L'assignació està ben definida ---com que $\alpha$ preserva origen i destí, arestes encadenades van a arestes encadenades, i el resultat és de nou un camí--- i és un functor: transforma concatenacions en concatenacions i camins buits en camins buits, és a dir, composicions en composicions i identitats en identitats.

Finalment, $\Free$ respecta la identitat i la composició de $\cat{Graph}$: $\Free(\id_G)=\id_{\Free(G)}$ i $\Free(\beta\comp\alpha)=\Free(\beta)\comp\Free(\alpha)$, ja que aplicar $\alpha$ i després $\beta$ aresta a aresta és el mateix que aplicar $\beta\comp\alpha$. Tenim, doncs, un functor
\[
\Free\colon\cat{Graph}\to\cat{Cat}.
\]

Vegem-ho amb el mateix graf d'un bucle de l'exemple anterior. Sigui $G$ el graf amb un sol vèrtex $\ast$ i una sola aresta $a$, un bucle. Els camins de $G$ són les paraules en la lletra $a$ ---$\varepsilon, a, aa, aaa,\dots$---, de manera que $\Free(G)$ és la categoria d'un sol objecte els morfismes de la qual són aquestes paraules, amb la concatenació per composició i la paraula buida $\varepsilon$ per identitat:
\[
\begin{tikzpicture}[baseline=(m.center)]
  \matrix (m) [cat matrix] { \ast \\ };
  \draw[cat] (m-1-1) to[out=55,in=125,looseness=9] node[above] {\(a\)} (m-1-1);
\end{tikzpicture}
\qquad\stackrel{\Free}{\longmapsto}\qquad
\begin{tikzpicture}[baseline=(m.center)]
  \matrix (m) [cat matrix] { \ast \\ };
  \draw[cat] (m-1-1) to[out=120,in=175,looseness=12] node[left] {\(\varepsilon\)} (m-1-1);
  \draw[cat] (m-1-1) to[out=65,in=115,looseness=12] node[above] {\(a\)} (m-1-1);
  \draw[cat] (m-1-1) to[out=5,in=60,looseness=12] node[right] {\(aa\)} (m-1-1);
  \draw[cat] (m-1-1) to[out=-55,in=-5,looseness=12] node[right] {\(\dots\)} (m-1-1);
  \draw[cat] (m-1-1) to[out=-125,in=-70,looseness=12] node[below] {\(a^n\)} (m-1-1);
\end{tikzpicture}
\]
Comptant la paraula buida com a $a^0$, els morfismes són exactament les potències $a^n$ amb $n\geq 0$, i concatenar-les suma els exponents: $a^m\comp a^n=a^{m+n}$. Aquest monoide de morfismes és $(\mathbb{N},+,0)$, de manera que $\Free(G)=\cat{B}\mathbb{N}$. Mentre que $U$ oblidava estructura, $\Free$ en crea: d'una sola aresta en surten infinits morfismes, un per cada longitud de camí.

Cal notar que $\Free$ i $U$ són functors \emph{entre categories grans} ($\cat{Graph}$ i $\cat{Cat}$ no són objectes de la $\cat{Cat}$ de les categories petites); les tractem com a categories localment petites, no com a fletxes de $\cat{Cat}$.

\subsection{La relació entre \texorpdfstring{$\Free$}{Free} i \texorpdfstring{$U$}{U}}\label{subsec:free-u}

Les dues construccions van en sentits contraris, però no són inverses l'una de l'altra. El functor $U$ oblida quins morfismes són identitats i com es componen; $\Free$ afegeix formalment identitats i composicions al graf, sense imposar cap altra relació. Quan partim del graf subjacent d'una categoria $\cat{C}$, la composició lliure crea camins formals nous, de manera que $\Free(U(\cat{C}))$ no és $\cat{C}$ en general; però hi ha un functor
\[
\Free(U(\cat{C}))\to\cat{C}
\]
que avalua cada camí formal mitjançant la composició de $\cat{C}$. Per exemple, prenem per $\cat{C}$ el monoide $M=(\{1,e\},\cdot,1)$ amb $e\cdot e=e$, vist com a categoria d'un sol objecte $\ast$. El seu graf subjacent té un vèrtex i dos bucles $\id_\ast$ i $e$, de manera que $\Free(U(\cat{C}))$ té per morfismes \emph{totes} les paraules en aquestes dues lletres ---infinites---, amb la paraula buida $\varepsilon$ per identitat. El functor $\Free(U(\cat{C}))\to\cat{C}$ avalua cada paraula multiplicant-la a $M$: $\varepsilon\mapsto\id_\ast$, $ee\mapsto e\comp e=e$, $\id_\ast e\mapsto e$, i així successivament. Col·lapsa, doncs, la infinitat de camins formals sobre els dos únics morfismes de $\cat{C}$, recuperant les composicions que $\Free$ havia descartat. Aquest functor il·lustra una propietat que caracteritza $\Free(G)$. Sigui $\cat{D}$ una categoria petita. Definir un functor $\Free(G)\to\cat{D}$ equival a triar les imatges dels vèrtexs \emph{i} de les arestes de $G$, respectant origen i destí ---és a dir, a donar un morfisme de grafs $G\to U(\cat{D})$---, i aquesta tria s'estén de manera única a $\Free(G)$. Cal comptar totes dues assignacions: un functor ha d'enviar cada objecte de $\Free(G)$ ---cada vèrtex de $G$--- a un objecte de $\cat{D}$, també els vèrtexs aïllats, que no són extrem de cap aresta. Per exemple, si $G$ té un únic vèrtex i cap aresta, no hi ha cap imatge d'aresta per triar, però un functor $\Free(G)\to\cat{D}$ encara exigeix escollir un objecte de $\cat{D}$; és precisament la component sobre vèrtexs del morfisme de grafs $G\to U(\cat{D})$, que en aquest cas es redueix a triar aquest objecte. La formulació via $G\to U(\cat{D})$ recull sempre les dues dades, perquè un morfisme de grafs consta d'una aplicació sobre vèrtexs i una sobre arestes. Aquesta correspondència és una bijecció
\[
\Hom_{\cat{Cat}}\big(\Free(G),\cat{D}\big)\;\cong\;\Hom_{\cat{Graph}}\big(G,U(\cat{D})\big).
\]
Aquesta és la formulació precisa de \enquote{lliure}. La retrobarem com a exemple al capítol~\refcap{cap:yoneda}{4}, i en tota la seva generalitat com l'\emph{adjunció} $\Free\dashv U$ al capítol~\refcap{cap:adjuncions}{6}.

\begin{observacio}
Reconèixer els functors com a fletxes d'una categoria és el primer pas per poder plantejar-nos si són isomorfismes, si es componen o si hi ha dualitat. Dins $\cat{Cat}$, per exemple, una aplicació monòtona $P\to Q$ entre preordres petits és una fletxa; en canvi, $\Free$ i $U$ són functors entre categories grans, i per això no són morfismes de $\cat{Cat}$; això no els fa menys legítims, sinó que simplement no els agrupem com a fletxes de cap categoria de categories.
\end{observacio}

\section{Functors covariants i contravariants}\label{sec:variancia}

Fins ara els functors preserven el sentit de les fletxes. Sovint, però, apareixen assignacions que \emph{inverteixen} el sentit: enviar $f\colon A\to B$ a un morfisme $F(B)\to F(A)$. Parlem aleshores de \textbf{functors contravariants}, que s'expliquen millor mitjançant la categoria oposada.

\begin{definicio}\index{functor!contravariant}\index{functor!covariant}
Un \textbf{functor contravariant} de $\cat{C}$ a $\cat{D}$ és un functor
\[
F\colon\cat{C}\op\to\cat{D}.
\]
Per contrast, els functors de la definició original s'anomenen \textbf{covariants}.
\end{definicio}

Desplegant la definició de functor sobre $\cat{C}\op$, un functor contravariant assigna a cada objecte $A$ un objecte $F(A)$ i a cada morfisme $f\colon A\to B$ de $\cat{C}$ un morfisme en sentit invertit
\[
F(f)\colon F(B)\to F(A),
\]
de manera que $F(\id_A)=\id_{F(A)}$ i $F(g\comp f)=F(f)\comp F(g)$. La inversió de l'ordre en aquesta última igualtat no és cap condició nova: és la condició de functor covariant sobre $\cat{C}\op$, un cop recordem que compondre a $\cat{C}\op$ és compondre a $\cat{C}$ en l'ordre invers. Aquest és un primer exemple de la potència de la categoria oposada: permet tractar covariància i contravariància amb una sola definició.

Diagramàticament, un triangle commutatiu de $\cat{C}$ es transforma en un triangle commutatiu de $\cat{D}$ amb les fletxes imatge invertides:
\[
\begin{tikzpicture}[baseline=(m.center)]
  \matrix (m) [matrix of math nodes, column sep=3.2em, row sep=2.6em] {
    A & B \\
      & C \\
  };
  \draw[-{Stealth[scale=1.1]}] (m-1-1) -- node[above] {$f$} (m-1-2);
  \draw[-{Stealth[scale=1.1]}] (m-1-2) -- node[right] {$g$} (m-2-2);
  \draw[-{Stealth[scale=1.1]}] (m-1-1) -- node[below left] {$g\comp f$} (m-2-2);
\end{tikzpicture}
\qquad\stackrel{F}{\longmapsto}\qquad
\begin{tikzpicture}[baseline=(m.center)]
  \matrix (m) [matrix of math nodes, column sep=3.6em, row sep=2.6em] {
    F(A) & F(B) \\
         & F(C) \\
  };
  \draw[{Stealth[scale=1.1]}-] (m-1-1) -- node[above] {$F(f)$} (m-1-2);
  \draw[{Stealth[scale=1.1]}-] (m-1-2) -- node[right] {$F(g)$} (m-2-2);
  \draw[{Stealth[scale=1.1]}-] (m-1-1) -- node[below left] {$F(g\comp f)$} (m-2-2);
\end{tikzpicture}
\]

\begin{definicio}[Functor oposat]\label{def:functor-oposat}\index{functor!oposat}
Tot functor $F\colon\cat{C}\to\cat{D}$ indueix un \textbf{functor oposat}
\[
F\op\colon\cat{C}\op\to\cat{D}\op,
\]
que actua igual que $F$ sobre objectes i sobre morfismes, però llegit entre les categories oposades. Com que compondre a $\cat{C}\op$ i a $\cat{D}\op$ és compondre a $\cat{C}$ i a $\cat{D}$ en l'ordre invers, les dues condicions de functor es transporten sense canvis.
\end{definicio}

\begin{observacio}
El pas a l'oposat és, de fet, functorial: assigna a cada categoria petita $\cat{C}$ la seva oposada $\cat{C}\op$ i a cada functor $F$ el functor oposat $F\op$, i respecta composicions i identitats, de manera que
\[
(-)\op\colon\cat{Cat}\to\cat{Cat}
\]
és un functor ---covariant: envia $F\colon\cat{C}\to\cat{D}$ a $F\op\colon\cat{C}\op\to\cat{D}\op$, en el mateix sentit. A més, aplicar-lo dues vegades torna al punt de partida, tant sobre objectes com sobre morfismes: $(\cat{C}\op)\op=\cat{C}$ i $(F\op)\op=F$. Dit altrament, $(-)\op\comp(-)\op=\id_{\cat{Cat}}$: el functor $(-)\op$ és la seva pròpia inversa, és a dir, una \textbf{involució}. En particular és un \textbf{isomorfisme de categories} de $\cat{Cat}$ en ella mateixa ---un endofunctor amb functor invers, que a més coincideix amb ell mateix. És la manifestació, al nivell de $\cat{Cat}$, de l'autodualitat que ja hem trobat objecte a objecte amb el principi de dualitat. Convé recordar, això sí, que tot plegat viu sobre categories \emph{petites}, les úniques que són objectes de $\cat{Cat}$; per a una categoria gran, l'oposada segueix tenint sentit, però no com a imatge d'una fletxa de $\cat{Cat}$.
\end{observacio}

\begin{proposicio}[Variància d'una composició]\label{prop:variancia-composicio}\index{functor!variància d'una composició}\index{regla dels signes (variància)}
Siguin $F$ i $G$ dos functors, cadascun covariant o contravariant, amb el codomini de $F$ i el domini de $G$ iguals com a categoria (a banda de la variància). L'assignació sobre objectes $A\mapsto G(F(A))$ és de nou un functor, i la seva variància segueix la \textbf{regla dels signes}: la composició és covariant si $F$ i $G$ tenen la mateixa variància, i contravariant si en tenen de diferents. En particular, la composició de dos functors contravariants és covariant.
\end{proposicio}

\begin{prova}
Un functor contravariant $\cat{C}\to\cat{D}$ és, per definició, un functor covariant $\cat{C}\op\to\cat{D}$, i la composició de functors torna a ser un functor (ho hem vist en construir $\cat{Cat}$). Només cal fer encaixar dominis i codominis, i d'això se n'ocupa el functor oposat.

Si $F$ i $G$ són tots dos covariants, $G\comp F\colon\cat{C}\to\cat{E}$ és directament la composició, covariant. Si tots dos són contravariants, escrits com $F\colon\cat{C}\op\to\cat{D}$ i $G\colon\cat{D}\op\to\cat{E}$, prenem el functor oposat $F\op\colon\cat{C}\to\cat{D}\op$ (definició~\ref{def:functor-oposat}), que actua igual que $F$ sobre objectes; aleshores
\[
G\comp F\op\colon\cat{C}\to\cat{E}
\]
és una composició de functors ---per tant un functor--- covariant, i sobre objectes fa $A\mapsto G(F(A))$. Els dos casos mixtos són idèntics en mètode: un ajust anàleg deixa el domini en $\cat{C}\op$ i la composició resulta contravariant. La variància només depèn, doncs, de la paritat del nombre de factors contravariants.
\end{prova}

\subsection{Els homfunctors}\label{subsec:homfunctors}

L'exemple més important de functors covariants i contravariants ---i el que sostindrà bona part del que ve després, fins al lema de Yoneda--- surt de mirar els morfismes d'una categoria cap a un objecte fix, o des d'un objecte fix.

\begin{definicio}[Els homfunctors]\label{def:homfunctors}\label{ex:homfunctors}\index{functor!Hom}\index{homfunctor}
Sigui $\cat{C}$ una categoria localment petita. Fixat un objecte $A$ de $\cat{C}$, definim \textbf{homfunctor covariant}\index{homfunctor!covariant}\index{functor!Hom!covariant}
\[
\Hom(A,-)\colon\cat{C}\to\cat{Set}
\]
com el functor que envia cada objecte $X$ al conjunt $\Hom(A,X)$ i cada morfisme $f\colon X\to Y$ a l'aplicació
\[
f_*\colon\Hom(A,X)\to\Hom(A,Y),
\qquad
\varphi\mapsto f\comp\varphi,
\]
la \textbf{postcomposició} amb $f$. Fixat un objecte $B$, definim \textbf{homfunctor contravariant}\index{homfunctor!contravariant}\index{functor!Hom!contravariant}
\[
\Hom(-,B)\colon\cat{C}\op\to\cat{Set}
\]
com el functor que envia cada objecte $X$ al conjunt $\Hom(X,B)$ i cada morfisme $f\colon X\to Y$ a l'aplicació
\[
f^*\colon\Hom(Y,B)\to\Hom(X,B),
\qquad
\psi\mapsto\psi\comp f,
\]
la \textbf{precomposició} amb $f$.
\end{definicio}

\begin{notacio}\label{not:hom-morfismes}\index{Hom@$\Hom$!aplicat a un morfisme}
Com per a qualsevol functor $F$, que envia un morfisme $f$ a $F(f)$, escriurem també
\[
\Hom(A,f)\coloneqq f_*
\qquad\text{i}\qquad
\Hom(f,B)\coloneqq f^*
\]
per a les imatges de $f\colon X\to Y$ pels homfunctors $\Hom(A,-)$ i $\Hom(-,B)$, respectivament.
\end{notacio}

En el cas covariant, un element $\varphi\colon A\to X$ es perllonga amb $f$ per obtenir $f\comp\varphi\colon A\to Y$; en el contravariant, un element $\psi\colon Y\to B$ es precedeix de $f$ per obtenir $\psi\comp f\colon X\to B$:
\[
\begin{tikzpicture}[baseline=(m.center)]
  \matrix (m) [matrix of math nodes, column sep=3.2em, row sep=2.6em] {
    A & X \\
      & Y \\
  };
  \draw[-{Stealth[scale=1.1]}] (m-1-1) -- node[above] {$\varphi$} (m-1-2);
  \draw[-{Stealth[scale=1.1]}] (m-1-2) -- node[right] {$f$} (m-2-2);
  \draw[-{Stealth[scale=1.1]}] (m-1-1) -- node[below left] {$f\comp\varphi$} (m-2-2);
\end{tikzpicture}
\qquad\qquad
\begin{tikzpicture}[baseline=(m.center)]
  \matrix (m) [matrix of math nodes, column sep=3.2em, row sep=2.6em] {
    X & Y \\
      & B \\
  };
  \draw[-{Stealth[scale=1.1]}] (m-1-1) -- node[above] {$f$} (m-1-2);
  \draw[-{Stealth[scale=1.1]}] (m-1-2) -- node[right] {$\psi$} (m-2-2);
  \draw[-{Stealth[scale=1.1]}] (m-1-1) -- node[below left] {$\psi\comp f$} (m-2-2);
\end{tikzpicture}
\]
Observem que, en el cas covariant, $f$ actua sobre el \emph{destí} de les fletxes (les allarga per l'extrem final), mentre que en el contravariant actua sobre l'\emph{origen} (les allarga per l'extrem inicial), i això és el que inverteix el sentit del functor. La functorialitat és conseqüència directa de l'associativitat de la composició i de les lleis de les identitats. El que importa és l'ordre que en resulta:
\[
(g\comp f)_*=g_*\comp f_*,
\qquad\qquad
(g\comp f)^*=f^*\comp g^*.
\]
En el cas covariant l'ordre es conserva; en el contravariant s'inverteix, perquè $f$ actua per l'extrem inicial de les fletxes. Les comprovacions detallades són els exercicis resolts~\ref{pr-hom-cov} i~\ref{pr-hom-contra} de la Pràctica. Aquests dos functors són centrals en teoria de categories i els retrobarem en el lema de Yoneda (capítol~\refcap{cap:yoneda}{4}).

Els homfunctors traslladen els isomorfismes de $\cat{C}$ a bijeccions de conjunts. Això és un cas particular de la proposició~\ref{prop:functor-preserva-iso}, però val la pena escriure'l explícitament, perquè és la forma en què l'usarem sovint.

\begin{corollari}[Els homfunctors transformen isos en bijeccions]\label{prop:homfunctor-iso}\index{homfunctor!i isomorfismes}
Sigui $\cat{C}$ localment petita i sigui $f\colon X\to Y$ un isomorfisme de $\cat{C}$. Aleshores, per a tot objecte $A$, la postcomposició
\[
f_*\colon\Hom(A,X)\to\Hom(A,Y),\qquad\varphi\mapsto f\comp\varphi,
\]
és una bijecció, amb inversa $(f^{-1})_*$; i, per a tot objecte $B$, la precomposició
\[
f^*\colon\Hom(Y,B)\to\Hom(X,B),\qquad\psi\mapsto\psi\comp f,
\]
és una bijecció, amb inversa $(f^{-1})^*$.
\end{corollari}

\begin{prova}
Un isomorfisme de $\cat{Set}$ no és res més que una bijecció. Apliquem, doncs, la proposició~\ref{prop:functor-preserva-iso} a cada homfunctor. Per al covariant $\Hom(A,-)$, la imatge de l'isomorfisme $f$ és $f_*$, que resulta ser un isomorfisme de $\cat{Set}$ amb invers la imatge de $f^{-1}$, és a dir $(f^{-1})_*$. Per al contravariant $\Hom(-,B)$ ---un functor sobre $\cat{C}\op$, on $f$ segueix sent un isomorfisme--- la imatge de $f$ és $f^*$, que és igualment un isomorfisme de $\cat{Set}$ amb invers $(f^{-1})^*$. En tots dos casos, ser isomorfisme a $\cat{Set}$ vol dir ser bijecció.
\end{prova}

Dit d'una altra manera: si $X\cong Y$, aleshores $\Hom(A,X)\cong\Hom(A,Y)$ i $\Hom(Y,B)\cong\Hom(X,B)$ com a conjunts, per a qualssevol $A$ i $B$. Els objectes isomorfs són, doncs, indistingibles vistos a través de qualsevol homfunctor: una primera pista de la idea que un objecte queda determinat pels morfismes que hi arriben o que en surten, que el lema de Yoneda portarà fins al final.

Els functors contravariants cap a $\cat{Set}$, com l'homfunctor $\Hom(-,B)$, apareixen tan sovint que reben un nom propi.

\begin{definicio}[Prefeix]\label{def:prefeix}\index{prefeix@prefeix}
Sigui $\cat{C}$ una categoria. Un \textbf{prefeix} sobre $\cat{C}$ és un functor contravariant de $\cat{C}$ a $\cat{Set}$, és a dir, un functor
\[
P\colon\cat{C}\op\to\cat{Set}.
\]
Explícitament, $P$ assigna a cada objecte $A$ de $\cat{C}$ un conjunt $P(A)$ i a cada morfisme $f\colon A\to B$ una aplicació $P(f)\colon P(B)\to P(A)$, de manera que $P(\id_A)=\id_{P(A)}$ i $P(g\comp f)=P(f)\comp P(g)$.
\end{definicio}

Si $\cat{C}$ és localment petita, cada homfunctor contravariant $\Hom(-,B)$ és un prefeix sobre $\cat{C}$, i el functor de parts de l'exemple següent n'és un sobre $\cat{Set}$. Els prefeixos tornaran a aparèixer als capítols~\refcap{cap:yoneda}{4} i~\refcap{cap:topos}{9}.

Així, \emph{prefeix sobre $\cat{C}$} i \emph{functor contravariant de $\cat{C}$ a $\cat{Set}$} són sinònims: el nom no afegeix cap condició, ni tan sols de mida. A partir d'ara, per a un functor contravariant cap a $\cat{Set}$ direm sempre \emph{prefeix},\footnote{A la literatura també es parla de \emph{prefeixos sobre $\cat{C}$ amb valors a $\cat{D}$} (per exemple, prefeixos de grups o d'espais vectorials) per als functors $\cat{C}\op\to\cat{D}$. En aquest text, \emph{prefeix} vol dir sempre amb valors a $\cat{Set}$.} i n'escriurem la signatura $P\colon\cat{C}\op\to\cat{Set}$ quan convingui. Reservarem l'expressió \emph{functor contravariant} per als que arriben a altres categories, com el functor dual de l'exemple~\ref{ex:homfunctor-dual}. Els functors \emph{covariants} $\cat{C}\to\cat{Set}$ no reben nom propi en aquest text: són, simplement, els prefeixos sobre $\cat{C}\op$.

\begin{exemple}[El functor de parts i l'homfunctor $\Hom(-,\{0,1\})$]\label{ex:functor-parts}\index{functor!de parts}\index{homfunctor!i functor de parts}
Comencem dins de $\cat{Set}$, prenent com a objecte fixat el conjunt de dos elements $\{0,1\}$. Cada subconjunt $A\subseteq X$ es correspon bijectivament amb la seva \textbf{funció característica} $\chi_A\colon X\to\{0,1\}$, que val $1$ sobre $A$ i $0$ fora; recíprocament, cada aplicació $X\to\{0,1\}$ és la característica del subconjunt on val $1$. Les funcions característiques donen, doncs, per a cada conjunt $X$, una bijecció canònica
\[
\mathcal P(X)\;\cong\;\Hom_{\cat{Set}}(X,\{0,1\}).
\]
Els dos conjunts no són iguals ---un conté subconjunts, l'altre aplicacions---, però la bijecció no depèn de cap tria.

Vegem que, sota aquestes bijeccions, la precomposició correspon a l'antiimatge. Per a $f\colon X\to Y$, l'homfunctor envia una aplicació $\chi\colon Y\to\{0,1\}$ a $\chi\comp f\colon X\to\{0,1\}$; i si $\chi=\chi_B$ és la característica de $B\subseteq Y$, aleshores $\chi_B\comp f$ val $1$ exactament sobre els $x$ amb $f(x)\in B$, és a dir, és la característica de l'antiimatge $f^{-1}(B)$. Així, la precomposició
\[
\chi\longmapsto\chi\comp f
\]
és precisament l'aplicació d'antiimatge
\[
f^{-1}\colon\mathcal P(Y)\to\mathcal P(X),
\qquad
B\mapsto\{x\in X : f(x)\in B\}.
\]
Aquí $f^{-1}$ denota l'aplicació d'antiimatge, no l'invers de $f$, que pot no existir. Les antiimatges defineixen el \textbf{functor de parts} contravariant
\[
\mathcal P\colon\cat{Set}\op\to\cat{Set},\qquad X\mapsto\mathcal P(X),\qquad f\mapsto f^{-1},
\]
i la seva contravariància, $(g\comp f)^{-1}=f^{-1}\comp g^{-1}$, es pot transportar, a través de les bijeccions anteriors, des de la de l'homfunctor $\Hom(-,\{0,1\})$; la comprovació directa és l'exercici resolt~\ref{pr-parts-contra} de la Pràctica.

Per tant, el functor de parts contravariant i l'homfunctor $\Hom_{\cat{Set}}(-,\{0,1\})$ descriuen la mateixa construcció fins a una identificació canònica: provar cada conjunt amb un objecte distingit, el conjunt $\{0,1\}$ de valors de veritat. No són literalment el mateix functor. El llenguatge de les transformacions naturals (capítol~\ref{cap:transformacions}) permetrà expressar aquesta afirmació amb precisió: tots dos functors seran \emph{naturalment isomorfs} (exemple~\ref{ex:isos-naturals-cap2}).

Convé un avís sobre l'altra variància, i val la pena justificar-lo. Un homfunctor covariant $\Hom_{\cat{Set}}(\{0,1\},-)$ envia cada conjunt $X$ al conjunt d'aplicacions $\{0,1\}\to X$. Una tal aplicació queda completament determinada per la parella d'imatges $(x_0,x_1)$ amb $x_0$ la imatge de $0$ i $x_1$ la de $1$, i tota parella d'elements de $X$ prové d'una aplicació; per tant
\[
\Hom_{\cat{Set}}(\{0,1\},X)\;\cong\;X\times X,
\]
i sobre morfismes la postcomposició amb $f\colon X\to Y$ actua component a component, $(x_0,x_1)\mapsto(f(x_0),f(x_1))$. Aquest és, doncs, el functor \enquote{parell ordenat}, ben diferent del functor de parts: $\Hom(\{0,1\},X)\cong X\times X$ té sempre $\lvert X\rvert^2$ elements (quan $X$ és finit), mentre que $\mathcal P(X)$ en té $2^{\lvert X\rvert}$. La correspondència canònica entre el functor de parts i un homfunctor és, per tant, un fenomen exclusiu de la cara contravariant: $\mathcal P$ correspon a $\Hom(-,\{0,1\})$, no a $\Hom(\{0,1\},-)$. La raó de fons és que a $\mathcal P(X)$ el conjunt $\{0,1\}$ fa de \emph{destí} (cada subconjunt és una aplicació \emph{cap a} $\{0,1\}$), i fixar el destí dona un homfunctor contravariant; posar-lo com a \emph{origen} dona una construcció covariant, però tota una altra. La versió covariant genuïna del functor de parts existeix ---envia $f$ a la imatge directa $A\mapsto\{f(x):x\in A\}$--- i és functorial, però no correspon a cap homfunctor; quan calgui distingir-la de $\mathcal P$ l'escriurem $\mathcal P_!$.
\end{exemple}

\begin{exemple}[Les dues variàncies del functor de parts]\label{ex:parts-dues-variancies}\index{functor!de parts!covariant i contravariant}
Posem en paral·lel les dues maneres de fer functorial l'assignació $X\mapsto\mathcal P(X)$. Totes dues coincideixen sobre objectes i difereixen en tot el que fan amb una aplicació $f\colon X\to Y$:
\[
\renewcommand{\arraystretch}{1.3}
\begin{array}{l|l|l}
 & \mathcal P_!\colon\cat{Set}\to\cat{Set}\ \text{(covariant)} & \mathcal P\colon\cat{Set}\op\to\cat{Set}\ \text{(contravariant)}\\ \hline
\text{acció sobre } f & \text{imatge directa } A\mapsto f(A) & \text{antiimatge } B\mapsto f^{-1}(B)\\
\text{tipus} & \mathcal P_!(f)\colon\mathcal P(X)\to\mathcal P(Y) & \mathcal P(f)\colon\mathcal P(Y)\to\mathcal P(X)\\
\text{identitat} & \id_X(A)=A & \id_X^{-1}(B)=B\\
\text{composició} & (g\comp f)(A)=g\bigl(f(A)\bigr) & (g\comp f)^{-1}(C)=f^{-1}\bigl(g^{-1}(C)\bigr)\\
\text{llei functorial} & \mathcal P_!(g\comp f)=\mathcal P_!(g)\comp\mathcal P_!(f) & \mathcal P(g\comp f)=\mathcal P(f)\comp\mathcal P(g)
\end{array}
\]
Llegida per files, la taula segueix l'ordre de l'observació~\ref{obs:comprovar-tipus}: primer el tipus, que ja fixa la variància, i després les dues equacions. Les comprovacions completes són els exercicis resolts~\ref{pr-parts-cov} i~\ref{pr-parts-contra} de la Pràctica. Observem que a la columna contravariant l'ordre de $f$ i $g$ s'inverteix: és el mateix fenomen que a $(g\comp f)^*=f^*\comp g^*$ per a l'homfunctor contravariant.

Les dues columnes no són independents. Per a $A\subseteq X$ i $B\subseteq Y$,
\[
f(A)\subseteq B\quad\Longleftrightarrow\quad A\subseteq f^{-1}(B),
\]
perquè totes dues condicions diuen que $f(a)\in B$ per a tot $a\in A$. Aquesta equivalència entre les dues accions, una en cada sentit, és el primer exemple d'una \emph{adjunció}; la retrobarem al capítol~\refcap{cap:adjuncions}{6} com a $\exists_f\dashv f^{*}$, l'origen categòric del quantificador existencial (exemple~\refalt{ex:quantificadors-set}{d'imatge, antiimatge i quantificadors del capítol~6}).
\end{exemple}

En aquest exemple l'objecte fixat feia de \emph{destí}; fixant-lo com a \emph{origen} obtenim homfunctors covariants amb una lectura igual de concreta. Ho veiem a $\cat{Graph}$, on l'estructura dels objectes ja no és la d'un simple conjunt.

\begin{exemple}[Vèrtexs i arestes com a homfunctors]\label{ex:homfunctor-grafs}\index{homfunctor!a Graph}\index{graf dirigit!vèrtexs i arestes com a homfunctors}
Treballem a $\cat{Graph}$, amb la notació de la subsecció~\ref{subsec:grafs}: un graf $G$ té vèrtexs $V_G$, arestes $E_G$ i aplicacions origen i destí $s_G,t_G\colon E_G\to V_G$. Triem dos grafs \enquote{patró} ben senzills:
\[
\mathbf V=\bigl(\{\ast\},\ \varnothing\bigr),
\qquad
\mathbf E=\bigl(\{0,1\},\ \{e\}\bigr)\ \text{amb } s(e)=0,\ t(e)=1,
\]
és a dir, $\mathbf V$ és un sol vèrtex sense arestes i $\mathbf E$ és una sola aresta entre dos vèrtexs.

Un morfisme de grafs $\mathbf V\to G$ no té cap aresta on decidir res, i sobre l'únic vèrtex tria un vèrtex de $G$; per tant queda determinat per aquesta tria, i n'hi ha un per cada vèrtex:
\[
\Hom_{\cat{Graph}}(\mathbf V,G)\;\cong\;V_G.
\]
Un morfisme de grafs $\mathbf E\to G$ tria una aresta $\varepsilon$ de $G$ ---la imatge de $e$--- i, forçosament, els seus dos extrems han d'anar a $s_G(\varepsilon)$ i $t_G(\varepsilon)$, de manera que la tria de l'aresta ja determina tot el morfisme; n'hi ha un per cada aresta:
\[
\Hom_{\cat{Graph}}(\mathbf E,G)\;\cong\;E_G.
\]
Així, els homfunctors covariants $\Hom_{\cat{Graph}}(\mathbf V,-)$ i $\Hom_{\cat{Graph}}(\mathbf E,-)$ corresponen canònicament, respectivament, al functor \enquote{conjunt de vèrtexs} i al functor \enquote{conjunt d'arestes} $\cat{Graph}\to\cat{Set}$. Sobre un morfisme $F=(F_V,F_E)\colon G\to H$, el functor $\Hom_{\cat{Graph}}(\mathbf V,-)$ actua per la postcomposició $F_*\colon\Hom_{\cat{Graph}}(\mathbf V,G)\to\Hom_{\cat{Graph}}(\mathbf V,H)$, $v\mapsto F\comp v$ (definició~\ref{def:homfunctors}). Si $v$ tria el vèrtex $v(\ast)$, aleshores $F\comp v$ tria el vèrtex $F_V(v(\ast))$; per tant, llegida a través de la bijecció $\Hom_{\cat{Graph}}(\mathbf V,G)\cong V_G$, la postcomposició $F_*$ és exactament $F_V\colon V_G\to V_H$. Anàlogament, per a $\Hom_{\cat{Graph}}(\mathbf E,-)$, la postcomposició $F_*$ envia el morfisme que tria l'aresta $\varepsilon$ al que tria $F_E(\varepsilon)$, i correspon a $F_E\colon E_G\to E_H$.

El que hem trobat és que dos functors oblit familiars ---prendre el conjunt de vèrtexs i prendre el conjunt d'arestes--- corresponen canònicament als homfunctors covariants $\Hom_{\cat{Graph}}(\mathbf V,-)$ i $\Hom_{\cat{Graph}}(\mathbf E,-)$, per a una tria adient de l'objecte d'origen; al capítol~\ref{cap:transformacions} reconeixerem aquestes identificacions com a isomorfismes naturals. És el mateix fenomen que fa que, a $\cat{Set}$, el functor identitat correspongui canònicament a $\Hom_{\cat{Set}}(\{\ast\},-)$: provar un objecte amb els morfismes des d'un patró petit en llegeix una part de l'estructura. Que un functor cap a $\cat{Set}$ es pugui reconstruir així, com un $\Hom(A,-)$ per a cert objecte $A$, és una propietat prou important perquè hi tornem amb nom propi ---els \emph{functors representables}--- i és el punt de partida del lema de Yoneda (capítol~\refcap{cap:yoneda}{4}); de moment n'hi ha prou de retenir que aquests functors oblit són homfunctors disfressats, llevat d'una identificació canònica.
\end{exemple}

Als exemples anteriors la categoria d'arribada era $\cat{Set}$. Quan els homconjunts tenen \emph{ells mateixos} estructura, l'homfunctor la recull i aterra en una categoria més rica; el cas paradigmàtic és el dels espais vectorials.

\begin{exemple}[El dual d'un espai vectorial com a homfunctor]\label{ex:homfunctor-dual}\index{homfunctor!i espai dual}\index{espai dual!com a homfunctor}\index{aplicació dual}
Fixem un cos $\mathbb K$ i treballem a $\cat{Vect}_{\mathbb K}$. Aquí cada homconjunt $\Hom_{\mathbb K}(V,W)$ no és un conjunt qualsevol: és ell mateix un $\mathbb K$-espai vectorial, amb la suma d'aplicacions lineals i el producte per escalars. Prenent com a segon objecte el cos $\mathbb K$ vist com a espai vectorial sobre si mateix, l'homconjunt
\[
\Hom_{\mathbb K}(V,\mathbb K)=V^*
\]
és exactament l'\textbf{espai dual} de $V$: les formes lineals sobre $V$. Com a conjunt, $V^*$ és el valor en $V$ de l'homfunctor contravariant de la definició~\ref{def:homfunctors}, especialitzat a $B=\mathbb K$:
\[
\Hom_{\mathbb K}(-,\mathbb K)\colon\cat{Vect}_{\mathbb K}\op\to\cat{Set}.
\]
Però aquests homconjunts tenen canònicament estructura d'espai vectorial, amb la suma i el producte per escalars punt a punt, i la precomposició amb una aplicació lineal és lineal (ho veurem tot seguit). Per tant, la mateixa assignació s'\emph{aixeca} a un functor amb valors a $\cat{Vect}_{\mathbb K}$, el \textbf{functor dual}
\[
(-)^*=\Hom_{\mathbb K}(-,\mathbb K)\colon\cat{Vect}_{\mathbb K}\op\to\cat{Vect}_{\mathbb K},
\qquad V\mapsto V^*.
\]
És el mateix patró que el functor de parts, amb $\mathbb K$ en lloc de $\{0,1\}$ com a objecte distingit. Sobre morfismes, la precomposició $\psi\mapsto\psi\comp f$ és precisament l'\textbf{aplicació dual} (o transposada) d'una aplicació lineal $f\colon V\to W$:
\[
f^*\colon W^*\to V^*,\qquad \psi\mapsto\psi\comp f.
\]
La contravariància, $(g\comp f)^*=f^*\comp g^*$, és la regla familiar de la transposició; si triem bases i representem les aplicacions per matrius, $f^*$ correspon a la matriu transposada i la inversió de l'ordre és la fórmula $(AB)^{\mathsf T}=B^{\mathsf T}A^{\mathsf T}$.

Cada $f^*$ és, en efecte, una aplicació lineal, $f^*(\psi+\lambda\psi')=(\psi+\lambda\psi')\comp f=f^*(\psi)+\lambda f^*(\psi')$, i això és el que permet l'aixecament: l'homfunctor $\cat{Set}$-valuat i el functor dual tenen la mateixa acció sobre objectes i morfismes, però el segon recorda, a més, l'estructura lineal dels valors. El mateix passa amb el covariant: $\Hom_{\mathbb K}(V,-)\colon\cat{Vect}_{\mathbb K}\to\cat{Set}$ s'aixeca a un functor
\[
\Hom_{\mathbb K}(V,-)\colon\cat{Vect}_{\mathbb K}\to\cat{Vect}_{\mathbb K}.
\]
Quan parlem d'homfunctors sense cap més precisió, ens referim sempre a la versió amb valors a $\cat{Set}$ de la definició~\ref{def:homfunctors}; és la que intervindrà en la representabilitat.

El corol·lari~\ref{prop:homfunctor-iso} es llegeix aquí com un fet conegut de l'àlgebra lineal: si $f\colon V\to W$ és un isomorfisme, aleshores la transposada $f^*\colon W^*\to V^*$ també ho és, amb inversa $(f^{-1})^*$. En particular, espais isomorfs tenen duals isomorfs.

No seguim, de moment, cap a la comparació entre $V$ i el seu dual, ni cap al bidual (o doble dual) $V^{**}$: aquesta és una qüestió de \emph{naturalitat}, que demana el llenguatge de les transformacions naturals. La reprendrem exactament on la deixem, a l'exemple~\ref{ex:bidual} del capítol~\ref{cap:transformacions}.
\end{exemple}


\subsection{El bihomfunctor}\label{subsec:bihomfunctor}

Les dues variables dels homfunctors es poden tractar alhora. Fixem-nos que $\Hom$ pren \emph{dos} objectes i en dona un conjunt, de manera contravariant en el primer i covariant en el segon; això és exactament un functor sobre un producte de categories.

\begin{definicio}[El bihomfunctor]\label{def:bihomfunctor}\index{homfunctor!bihomfunctor}\index{Hom@$\Hom$!bihomfunctor}
Sigui $\cat{C}$ localment petita. Definim \textbf{bihomfunctor} (o \textbf{bifunctor Hom})
\[
\Hom_{\cat{C}}(-,-)\colon\cat{C}\op\times\cat{C}\to\cat{Set}
\]
com el functor que envia una parella d'objectes $(A,B)$ al conjunt $\Hom(A,B)$, i una parella de morfismes $(f\colon A'\to A,\ g\colon B\to B')$ a l'aplicació
\[
\Hom(f,g)\colon\Hom(A,B)\to\Hom(A',B'),\qquad h\mapsto g\comp h\comp f.
\]
\end{definicio}

La fórmula $h\mapsto g\comp h\comp f$ combina la precomposició amb $f$ (contravariant, en la primera variable) i la postcomposició amb $g$ (covariant, en la segona): un element $h\colon A\to B$ es prolonga per l'origen amb $f$ i pel destí amb $g$.
\[
\begin{tikzpicture}[baseline=(m.center)]
  \matrix (m) [matrix of math nodes, column sep=3em] {
    A' & A & B & B' \\
  };
  \draw[-{Stealth[scale=1.1]}] (m-1-1) -- node[above] {$f$} (m-1-2);
  \draw[-{Stealth[scale=1.1]}] (m-1-2) -- node[above] {$h$} (m-1-3);
  \draw[-{Stealth[scale=1.1]}] (m-1-3) -- node[above] {$g$} (m-1-4);
  \draw[-{Stealth[scale=1.1]}] (m-1-1) to[bend right=22] node[below] {$g\comp h\comp f$} (m-1-4);
\end{tikzpicture}
\]
Fixant la primera variable en $A$ recuperem l'homfunctor covariant $\Hom(A,-)$; fixant la segona en $B$ recuperem el contravariant $\Hom(-,B)$. També sobre morfismes les notacions són coherents: prenent una de les dues components igual a una identitat,
\[
\Hom(\id_A,g)=g_*=\Hom(A,g),
\qquad
\Hom(f,\id_B)=f^*=\Hom(f,B),
\]
de manera que la notació~\ref{not:hom-morfismes} és un cas particular de la del bihomfunctor. Ara bé, la functorialitat en cada variable per separat no basta: cal que les dues accions siguin compatibles, és a dir, que $\Hom(-,-)$ preservi la composició i les identitats de $\cat{C}\op\times\cat{C}$. També això és conseqüència de l'associativitat i de les identitats: aplicar primer $(f,g)$ i després $(f',g')$ envia $h$ a $g'\comp g\comp h\comp f\comp f'$, que és l'acció de la parella composta $(f\comp f',\,g'\comp g)$, amb la primera component en l'ordre invertit, com correspon a $\cat{C}\op$. La comprovació detallada és l'exercici resolt~\ref{pr-bihom} de la Pràctica.

\begin{exemple}[El bihomfunctor de $\cat{Set}$]\label{ex:bihom-set}\index{homfunctor!bihomfunctor!a Set}
A $\cat{Set}$, el bihomfunctor envia una parella de conjunts $(X,Y)$ al conjunt d'aplicacions
\[
\Hom_{\cat{Set}}(X,Y)=Y^X,
\]
amb la mateixa notació exponencial de l'inici del capítol. L'acció sobre una parella de morfismes $(f\colon X'\to X,\ g\colon Y\to Y')$ és el \enquote{canvi de variables} a l'entrada i a la sortida: donada una aplicació $h\colon X\to Y$, primer la reindexem pel domini amb $f$ i després en traslladem els valors amb $g$,
\[
\Hom(f,g)\colon Y^X\to (Y')^{X'},\qquad h\mapsto g\comp h\comp f.
\]
Fixant el primer conjunt en un punt $\{\ast\}$ recuperem $\Hom(\{\ast\},Y)\cong Y$ ---els elements de $Y$--- amb l'acció covariant habitual per postcomposició; fixant el segon en $\{0,1\}$ recuperem $\Hom(X,\{0,1\})\cong\mathcal P(X)$, és a dir, el functor de parts de l'exemple~\ref{ex:functor-parts}, llevat de la identificació canònica. El bihomfunctor és, doncs, la construcció que engloba tots dos alhora.
\end{exemple}

\begin{exemple}[El bihomfunctor de $\cat{Vect}_{\mathbb K}$]\label{ex:bihom-vect}\index{homfunctor!bihomfunctor!a Vect}
El bihomfunctor ordinari de $\cat{Vect}_{\mathbb K}$ és
\[
\Hom(-,-)\colon\cat{Vect}_{\mathbb K}\op\times\cat{Vect}_{\mathbb K}\to\cat{Set},
\]
que envia una parella d'espais $(V,W)$ al conjunt $\Hom_{\mathbb K}(V,W)$ de les aplicacions lineals de $V$ a $W$. Però aquest conjunt és canònicament un espai vectorial, i sobre una parella $(f\colon V'\to V,\ g\colon W\to W')$ l'acció $h\mapsto g\comp h\comp f$ combina la precomposició amb $f$ i la postcomposició amb $g$, totes dues lineals. Per això el bihomfunctor admet un aixecament canònic a $\cat{Vect}_{\mathbb K}$:
\[
\Hom_{\mathbb K}(-,-)\colon\cat{Vect}_{\mathbb K}\op\times\cat{Vect}_{\mathbb K}\to\cat{Vect}_{\mathbb K}.
\]
Fixant la segona variable en el cos, aquest aixecament dona el functor dual $(-)^*$ de l'exemple~\ref{ex:homfunctor-dual}. Fixant la primera, $\Hom_{\mathbb K}(\mathbb K,-)$ correspon canònicament al functor identitat. En efecte, una aplicació lineal $\mathbb K\to W$ queda determinada per la imatge de $1$, i això dona una bijecció lineal canònica
\[
\Hom_{\mathbb K}(\mathbb K,W)\;\cong\;W;
\]
al capítol~\ref{cap:transformacions} la reconeixerem com un isomorfisme natural. Les dues cares del dual i la identitat es reuneixen, així, en un sol functor de dues variables.
\end{exemple}

\section{Functors fidels i plens}

% CONVENCIÓ D'ESTIL (aplicada a tot el llibre): s'eviten les nominalitzacions
% «fidelitat» i «plenitud» en la prosa expositiva; es prediquen com «ser fidel» /
% «ser ple». Propagat a teoria, exercicis i a la prosa de practica/tests.
% Excepcions conservades a propòsit: les etiquetes de pas en demostracions
% «(fidelitat)» / «per fidelitat» (caps/practica/*, caps/tests/*), i l'ús ordinari
% «per fidelitat a la bibliografia» del glossari (caps/glossari.tex).

Un functor actua sobre els morfismes; és natural preguntar-se com és aquesta acció entre cada parella d'objectes. Fixats $A,B$ a $\cat{C}$, un functor $F\colon\cat{C}\to\cat{D}$ indueix una aplicació
\[
F_{A,B}\colon\Hom_{\cat{C}}(A,B)\to\Hom_{\cat{D}}(F(A),F(B)),
\qquad
f\mapsto F(f).
\]
Les propietats d'aquesta aplicació donen lloc a una classificació bàsica.

\begin{definicio}\index{functor!fidel}\index{functor!ple}
\label{def:ple-fidel}
Un functor $F\colon\cat{C}\to\cat{D}$ és:
\begin{itemize}
\item \textbf{fidel} si cada aplicació $F_{A,B}$ és injectiva; és a dir, si $F(f)=F(g)$ amb $f,g\colon A\to B$ implica $f=g$;
\item \textbf{ple} si cada aplicació $F_{A,B}$ és exhaustiva; és a dir, si tot morfisme $F(A)\to F(B)$ de $\cat{D}$ és de la forma $F(f)$ per a algun $f\colon A\to B$ de $\cat{C}$;
\item \textbf{ple i fidel} si cada $F_{A,B}$ és bijectiva.
\end{itemize}
\end{definicio}

\begin{observacio}
Cal no confondre aquestes condicions amb la injectivitat o l'exhaustivitat \emph{sobre objectes}. Un functor pot ser ple i fidel sense ser injectiu sobre objectes: ser ple i fidel només afecta els morfismes entre cada parella fixada d'objectes. Un functor ple i fidel identifica $\Hom_{\cat{C}}(A,B)$ amb $\Hom_{\cat{D}}(F(A),F(B))$, però pot enviar objectes diferents a un mateix objecte.
\end{observacio}

\begin{observacio}
Aquestes propietats es transporten al functor oposat. Recordem (definició~\ref{def:functor-oposat}) que el functor $F\op\colon\cat{C}\op\to\cat{D}\op$ actua igual que $F$ sobre objectes i morfismes. Cal no confondre el \emph{functor} $F\op$ amb l'\emph{aplicació entre homconjunts} que indueix. Fixats dos objectes $A,B$, que $F$ sigui ple i fidel es mesura amb l'aplicació
\[
F_{A,B}\colon\Hom_{\cat{C}}(A,B)\to\Hom_{\cat{D}}(F(A),F(B)),\qquad f\mapsto F(f),
\]
i les de $F\op$, amb l'aplicació anàloga $(F\op)_{B,A}$ ---amb la parella d'objectes en l'ordre invers, perquè a $\cat{C}\op$ una fletxa $A\to B$ és una fletxa $B\to A$ de $\cat{C}$---
\[
(F\op)_{B,A}\colon\Hom_{\cat{C}\op}(B,A)\to\Hom_{\cat{D}\op}(F(B),F(A)).
\]
Ara bé, per definició d'oposada $\Hom_{\cat{C}\op}(B,A)=\Hom_{\cat{C}}(A,B)$ i $\Hom_{\cat{D}\op}(F(B),F(A))=\Hom_{\cat{D}}(F(A),F(B))$; i sobre aquests conjunts $(F\op)_{B,A}$ envia cada $f$ a $F(f)$, exactament com $F_{A,B}$. Les dues aplicacions són, doncs, literalment la mateixa. Com que \enquote{ser injectiva} i \enquote{ser exhaustiva} són propietats d'aquesta aplicació, $F$ és fidel (respectivament ple) si i només si $F\op$ ho és.
\end{observacio}

\begin{exemple}
El functor oblit $U\colon\cat{Grp}\to\cat{Set}$ és fidel: dos homomorfismes diferents segueixen sent aplicacions diferents, de manera que $U$ no identifica morfismes. En canvi no és ple: entre dos grups hi ha aplicacions de conjunts que no són homomorfismes, i aquestes no provenen de cap morfisme de $\cat{Grp}$. Aquesta és la situació típica dels functors oblit: retenen prou informació per distingir morfismes, però obren la porta a fletxes noves que no respecten l'estructura.
\end{exemple}

Les subcategories, introduïdes elementalment a la subsecció~\ref{subsec:subcategories}, tenen una lectura natural en el llenguatge dels functors. Si $\cat{S}$ és una subcategoria de $\cat{C}$, el \textbf{functor inclusió} $\iota\colon\cat{S}\hookrightarrow\cat{C}$ ---que envia cada objecte i cada morfisme de $\cat{S}$ a ell mateix a $\cat{C}$--- sempre és \textbf{fidel}, perquè no identifica mai dos morfismes diferents. Les dues menes de subcategoria que vam distingir ---la \textbf{plena} i l'\textbf{ampla}--- es tradueixen, cadascuna, en una propietat d'aquest functor. D'una banda, $\cat{S}$ és una \textbf{subcategoria plena} exactament quan $\iota$ és \textbf{ple}: que $\iota$ sigui ple vol dir que entre dos objectes de $\cat{S}$ hi ha \emph{tots} els morfismes que hi ha a $\cat{C}$, que és justament la definició de subcategoria plena. De l'altra, $\cat{S}$ és una \textbf{subcategoria ampla} exactament quan l'acció de $\iota$ \emph{sobre els objectes} és exhaustiva, és a dir, quan $\Ob(\cat{S})=\Ob(\cat{C})$: la subcategoria conté tots els objectes de $\cat{C}$. (Parlem aquí de l'aplicació $\Ob(\cat{S})\to\Ob(\cat{C})$ induïda per $\iota$, no del functor mateix; per al functor inclusió aquesta aplicació és sempre injectiva, i ser ampla equival que sigui, a més, exhaustiva.)\index{subcategoria!plena}\index{subcategoria!ampla}\index{functor!inclusió}

\begin{exemple}\label{ex:inclusio-finset-plena}
El \textbf{functor inclusió} $\cat{FinSet}\hookrightarrow\cat{Set}$ dels conjunts finits és ple i fidel: entre dos conjunts finits hi ha exactament les mateixes aplicacions que a $\cat{Set}$, de manera que un cop triats els objectes no descartem cap morfisme. En canvi, el functor inclusió de la subcategoria de $\cat{Set}$ que conserva tots els conjunts com a objectes però només les aplicacions \emph{injectives} com a morfismes és fidel, però \emph{no} ple: entre dos conjunts hi ha aplicacions no injectives, que per tant no provenen de cap morfisme de la subcategoria. El mateix passa amb la subcategoria ampla de les aplicacions \emph{exhaustives}: el seu functor inclusió és fidel però no ple, perquè deixa fora les aplicacions que no són exhaustives.
\end{exemple}

\begin{proposicio}\label{prop:plenfidel-reflecteix-iso}\index{functor!ple i fidel}
Sigui $F\colon\cat{C}\to\cat{D}$ un functor ple i fidel, i sigui $f\colon A\to B$ un morfisme de $\cat{C}$. Aleshores
\[
f\ \text{és un isomorfisme a }\cat{C}
\quad\Longleftrightarrow\quad
F(f)\ \text{és un isomorfisme a }\cat{D}.
\]
\end{proposicio}

\begin{prova}
La implicació cap a la dreta val per a qualsevol functor: és la proposició~\ref{prop:functor-preserva-iso}, i no fa servir que $F$ sigui ple ni fidel. Vegem el recíproc, que sí que ho necessita. Suposem que $F(f)\colon F(A)\to F(B)$ té un invers $h\colon F(B)\to F(A)$ a $\cat{D}$. Com que $F$ és ple, existeix $g\colon B\to A$ a $\cat{C}$ amb $F(g)=h$. Aleshores
\[
F(g\comp f)=F(g)\comp F(f)=h\comp F(f)=\id_{F(A)}=F(\id_A),
\]
i com que $F$ és fidel, $g\comp f=\id_A$. Anàlogament $f\comp g=\id_B$. Per tant $f$ és un isomorfisme amb invers $g$.
\end{prova}

\begin{observacio}
Les dues direccions tenen estatus diferent. La directa ---preservar isomorfismes--- val per a tot functor; el que caracteritza els plens i fidels és el recíproc: \emph{reflecteixen} isomorfismes, i no només els preserven. Un functor ple i fidel reprodueix bijectivament els conjunts de morfismes entre els objectes de la seva imatge, però això \emph{no} implica que sigui injectiu sobre objectes: pot enviar objectes diferents a un mateix objecte. Aquesta situació ---morfismes identificats bijectivament, objectes no necessàriament--- conduirà més endavant a la noció d'\emph{equivalència de categories}.
\end{observacio}

Ser ple i fidel descriu com actua un functor sobre els morfismes. Una tercera condició, sobre els objectes, completa la terna que caracteritzarà les equivalències.

\begin{definicio}[Functor essencialment exhaustiu]\label{def:ess-exhaustiu}\index{functor!essencialment exhaustiu}
Un functor $F\colon\cat{C}\to\cat{D}$ és \textbf{essencialment exhaustiu} si tot objecte de $\cat{D}$ és isomorf a la imatge d'algun objecte de $\cat{C}$:
\[
\forall D\in\cat{D},\ \exists A\in\cat{C}\ \text{tal que}\ F(A)\cong D.
\]
\end{definicio}

\begin{observacio}
La condició demana \emph{isomorfisme}, no igualtat: no cal que tot objecte de $\cat{D}$ sigui exactament una imatge, només que ho sigui llevat d'isomorfisme. Un functor alhora ple, fidel i essencialment exhaustiu recull bijectivament els morfismes i, a més, arriba a tots els objectes llevat d'isomorfisme; aquestes són precisament les condicions que, un cop disposem de les transformacions naturals, caracteritzaran les \emph{equivalències de categories}.
\end{observacio}

\begin{exemple}[Un functor essencialment exhaustiu que no és exhaustiu sobre objectes]\label{ex:ess-exhaustiu-finset}\index{functor!essencialment exhaustiu}
Considerem la subcategoria plena $\cat{S}$ de $\cat{FinSet}$ que té per objectes només els conjunts
\[
\underline{0}=\varnothing,\qquad \underline{n}=\{1,\dots,n\}\quad(n\geq 1)
\]
---un de sol de cada mida, en lloc de tots els conjunts finits---, i el seu functor inclusió $\iota\colon\cat{S}\hookrightarrow\cat{FinSet}$. Aquest functor \emph{és} essencialment exhaustiu: tot conjunt finit $X$ té algun cardinal $n$, i aleshores és isomorf (bijectiu) a $\underline{n}=\iota(\underline{n})$. En canvi, $\iota$ no és exhaustiu sobre objectes: el conjunt $\{a,b\}$, posem per cas, no és \emph{igual} a cap $\underline{n}$, tot i ser isomorf a $\underline{2}$. Aquest és, doncs, l'exemple prototípic de la diferència entre \enquote{arribar a tots els objectes} i \enquote{arribar-hi llevat d'isomorfisme}: $\iota$ deixa fora molts objectes com a imatges exactes, però cap classe d'isomorfisme. L'exercici resolt~\ref{pr-esquelet-finset} de la Pràctica comprova, a més, que $\iota$ és ple i fidel.
\end{exemple}

\begin{exemple}[Un functor que no és essencialment exhaustiu]\index{functor!essencialment exhaustiu}
En canvi, la inclusió $\cat{FinSet}\hookrightarrow\cat{Set}$ \emph{no} és essencialment exhaustiva: un conjunt infinit no és isomorf a cap conjunt finit, de manera que cap objecte finit de la imatge no hi pot arribar. Aquí falla, doncs, la tercera condició, tot i que aquest functor sí que és ple i fidel (exemple~\ref{ex:inclusio-finset-plena}). Ser ple i fidel no diu res sobre quins objectes s'assoleixen: són condicions independents.
\end{exemple}

\subsection{Què conserva un functor: un resum}\label{subsec:que-conserva}

Convé aclarir, abans de tancar el capítol, quines propietats respecta un functor \emph{qualsevol} i quines en demanen més. Un functor transporta qualsevol configuració finita de morfismes i les equacions de composició que aquests morfismes satisfan. Per això preserva, per exemple, isomorfismes, seccions i retraccions: la propietat està testimoniada per morfismes de la categoria d'origen ---l'invers, la secció o la retracció--- i per una equació com $g\comp f=\id$, i el functor transporta tant el testimoni com l'equació. En canvi, una propietat com ser mono o epi quantifica sobre totes les fletxes possibles d'un cert tipus (\enquote{per a tota parella $g,h$}), incloses fletxes de la categoria d'arribada que poden no provenir de la categoria d'origen; i és aquesta quantificació la que un functor no pot garantir.

\begin{itemize}
\item Un functor \emph{qualsevol} preserva composició i identitats, diagrames commutatius, isomorfismes, i seccions i retraccions (és a dir, els epimorfismes i monomorfismes \emph{escindits}). En tots els casos, la propietat està testimoniada per morfismes i equacions de la categoria d'origen, que el functor transporta.
\item Ja un functor \emph{fidel} ---sense necessitat de ser ple--- \emph{reflecteix} monomorfismes i epimorfismes: si $F(f)$ és mono, aleshores $f$ ja ho era, perquè de $f\comp g=f\comp h$ se'n dedueix $F(f)\comp F(g)=F(f)\comp F(h)$, d'on $F(g)=F(h)$ i, com que $F$ és fidel, $g=h$; l'argument per als epimorfismes és dual.
\item Un functor \emph{ple i fidel} reflecteix, a més, isomorfismes, seccions i retraccions: si la imatge d'un morfisme és un isomorfisme, una secció o una retracció, el morfisme original ja ho era. Aquí calen totes dues hipòtesis: ser ple permet aixecar l'invers ---unilateral o bilateral--- que \emph{a priori} només viu al codomini, i ser fidel permet reflectir l'equació que el defineix. Recupera, doncs, aquestes nocions a banda i banda.
\item Al capítol~\ref{cap:transformacions} veurem que un functor \emph{ple, fidel i essencialment exhaustiu} determina una \emph{equivalència de categories} (en general, amb l'ajut de l'axioma de l'elecció). Els capítols posteriors precisaran, per a cada construcció, quines propietats es preserven i es reflecteixen sota equivalència.

\end{itemize}

En canvi, un functor qualsevol \emph{no} preserva, en general, ni els monomorfismes ni els epimorfismes ni les propietats universals. Vegem-ho amb els epimorfismes.

\begin{exemple}[Un functor que no preserva epimorfismes]\index{functor!no preserva epimorfismes}
Prenem el monoide additiu $(\mathbb{N},+,0)$ i el monoide $(\{0,1\},\vee,0)$, on $\vee$ és el màxim (o disjunció), tots dos vistos com a categories d'un sol objecte, $\cat{B}\mathbb{N}$ i $\cat{B}M$. L'aplicació $h\colon\mathbb{N}\to\{0,1\}$ amb $h(0)=0$ i $h(n)=1$ per a $n\geq 1$ és un homomorfisme de monoides, i defineix per tant un functor $F\colon\cat{B}\mathbb{N}\to\cat{B}M$ (exemple~\ref{ex:functors-monoide}).

A $\cat{B}\mathbb{N}$, el morfisme $1$ és un epimorfisme: la cancel·lació per la dreta a $\cat{B}\mathbb{N}$ és la cancel·lació de la suma, i a $(\mathbb{N},+)$ es pot cancel·lar. La seva imatge $F(1)=1$, en canvi, \emph{no} és un epimorfisme a $\cat{B}M$, perquè $1\vee 0=1=1\vee 1$ amb $0\neq 1$: la cancel·lació falla. Així, $F$ envia un epimorfisme a un morfisme que no ho és.

L'arrel del fenomen és la que hem assenyalat: \enquote{ser epi} quantifica sobre totes les parelles $g,h$ de la categoria d'arribada, i el functor no controla les parelles noves que apareixen a $\cat{B}M$. Un epimorfisme \emph{escindit}, en canvi ---testimoniat per una secció $s$ amb $f\comp s=\id$---, sí que es preservaria, perquè el functor transporta la secció i l'equació.
\end{exemple}

\begin{resumcap}
\apartat{Idees centrals}
\begin{enumerate}
  \item Un functor preserva dominis, codominis, composició i identitats.
  \item Les categories petites i els functors formen una categoria, $\cat{Cat}$, que és gran però localment petita; els functors entre categories grans són igualment legítims, tot i que no són fletxes de $\cat{Cat}$.
  \item El graf subjacent i la categoria lliure formen una parella de functors en sentits contraris, $\Free\colon\cat{Graph}\to\cat{Cat}$ i $U\colon\cat{Cat}\to\cat{Graph}$. No són inversos l'un de l'altre, però estan lligats per l'avaluació $\Free(U(\cat{C}))\to\cat{C}$; aquesta relació reapareixerà més endavant com a adjunció.
  \item Un functor contravariant és un functor des de l'oposada. Per comprovar qualsevol functor, primer es verifiquen els tipus ($F(f)\colon F(A)\to F(B)$, o $F(B)\to F(A)$ en el cas contravariant) i després les identitats; el functor de parts té una versió de cada variància.
  \item Ser fidel o ser ple és una propietat de les aplicacions $F_{A,B}$ entre homconjunts, i no diu res de l'acció sobre objectes; afegint-hi l'exhaustivitat essencial, que sí que parla dels objectes (llevat d'isomorfisme), s'obtenen les equivalències.
  \item L'homfunctor covariant $\Hom(A,-)$ i el contravariant $\Hom(-,B)$ transformen isos en bijeccions, i el bihomfunctor $\Hom(-,-)$ els unifica en un functor de dues variables; tots tres són fonamentals i els retrobarem al lema de Yoneda.
\end{enumerate}
\apartat{Errors típics} Confondre \enquote{isomorfisme de categories} amb \enquote{equivalència}; creure que \enquote{fidel} vol dir injectiu sobre objectes o que \enquote{ple} vol dir exhaustiu sobre objectes, quan totes dues nocions parlen només dels homconjunts; comprovar les equacions functorials abans d'haver comprovat que $F(f)$ té el tipus correcte; oblidar la inversió de sentit en un functor contravariant; barrejar $\cat{C}/A$ (sobre $A$) amb $A/\cat{C}$ (sota $A$); i confondre \emph{preservar} amb \emph{reflectir} una propietat: tot functor preserva els isomorfismes; els functors plens i fidels també els reflecteixen, tot i que aquesta condició és suficient i no necessària. Un functor qualsevol no preserva, en general, ni els monomorfismes ni els epimorfismes.
\apartat{Lectura recomanada} \cite[part~III, sessions~11--17]{lawvereschanuel} per a exemples molt elementals d'aplicacions que preserven estructura entre grafs i entre sistemes dinàmics; \cite[§1.3 i §§7.1--7.2]{awodey} per a la definició de functor, la categoria de categories i l'estructura representable; \cite[I.3, II.2 i II.5]{maclane} per a la formulació clàssica, inclosos els functors contravariants; \cite[§1.2]{leinster} per als functors i les seves propietats, inclosos els contravariants i els plens i fidels; \cite[cap.~1]{riehl} per als homfunctors i el bifunctor $\Hom$; \cite[§§3.1--3.3]{barrwells} per a més exemples de functors, en particular les accions de monoides i els functors entre grafs, i per a la classificació dels functors en fidels, plens i altres tipus; \cite[§1.3 i §§1.5.1--1.5.3]{perrone-notes} per a una segona explicació de la functorialitat, amb una secció sobre què \emph{no} és un functor i una altra sobre functors, monos i epis.
\end{resumcap}

\chapter[Transformacions naturals]{\texorpdfstring{Transformacions naturals\\ i equivalències}{Transformacions naturals i equivalències}}
\label{cap:transformacions}

\begin{objectius}
\apartat{Prerequisits} Categories i diagrames commutatius (cap.~\ref{cap:categories}); functors, composició de functors, functors covariants i contravariants, i les nocions de functor ple, fidel i essencialment exhaustiu (cap.~\ref{cap:functors}). Els homfunctors intervenen en alguns exemples i exercicis, i seran importants en els capítols següents, però no són necessaris per a la definició de naturalitat. Alguns exemples opcionals fan servir una familiaritat bàsica amb espais vectorials, duals i biduals, i l'exemple d'interpretació de fórmules de primer ordre, marcat com a ampliació, pressuposa la semàntica d'aquesta lògica.
\apartat{Objectius} En acabar el capítol, el lector hauria de poder:
\begin{itemize}
  \item definir una transformació natural i comprovar-ne la naturalitat mitjançant quadrats commutatius;
  \item compondre transformacions verticalment i interpretar-les com a morfismes de la categoria de functors $\cat{D}^{\cat{C}}$;
  \item reconèixer un isomorfisme natural i distingir-lo d'una família d'isomorfismes que no sigui natural;
  \item treballar amb la composició horitzontal i distingir-la de la vertical;
  \item definir una equivalència de categories mitjançant un quasiinvers, i caracteritzar-la com un functor ple, fidel i essencialment exhaustiu (amb l'axioma de l'elecció per a la implicació recíproca).
\end{itemize}
\end{objectius}

\section{Cap al concepte de transformació natural}\label{sec:cap-transformacio}

Al capítol anterior vam veure que els functors són els morfismes entre categories: relacionen dues categories respectant-ne tota l'estructura. Ara bé, sovint tenim \emph{dos} functors $F,G\colon\cat{C}\to\cat{D}$ que van de la mateixa categoria a la mateixa categoria, i volem comparar-los. La pregunta natural és, doncs, si hi ha una noció de \emph{morfisme entre functors}.

Aquesta idea ---comparar functors de manera coherent--- va ser històricament una de les motivacions originals per introduir les categories. De fet, es diu sovint que les categories es van definir per poder definir els functors, i els functors per poder definir les transformacions naturals; totes tres nocions apareixen alhora, per primera vegada, a l'article fundacional d'Eilenberg i Mac~Lane~\cite{eilenbergmaclane}, que precisament les motiva a partir de la relació entre un espai vectorial i el seu bidual. El mot \enquote{natural} té aquí un sentit tècnic precís que farem explícit: una comparació és natural quan fa commutar els quadrats corresponents a tots els morfismes. Sovint s'expressa dient que \enquote{no depèn de cap elecció arbitrària}; convé retenir, però, que es tracta d'una intuïció orientadora i no d'una definició: una construcció pot haver requerit eleccions i, tanmateix, resultar natural. La condició precisa és sempre la commutativitat dels quadrats.

Abans de formalitzar-ho, val la pena veure un exemple que fa palès aquest \enquote{principi de coherència}, aquest cop entre \emph{sintaxi} i \emph{semàntica}. Treballem amb la lògica proposicional. Prenem com a categoria de base $\cat{FinSet}$, els conjunts finits amb les aplicacions entre ells (subsecció~\ref{subsec:conjunts}); ara en pensem els objectes com a conjunts de variables proposicionals, i els morfismes $f\colon X\to Y$ com a \emph{canvis de variables}, que a cada variable de $X$ n'assignen una de $Y$.

El primer functor, el \emph{sintàctic}, associa a cada conjunt de variables el conjunt de les fórmules que s'hi poden escriure:
\[
F(X)=\{\text{fórmules proposicionals amb variables de }X\},
\]
construïdes a partir de les variables de $X$ amb les connectives habituals. Sobre morfismes, un canvi de variables $f\colon X\to Y$ indueix l'aplicació de \emph{substitució}
\[
F(f)\colon F(X)\longrightarrow F(Y),
\qquad
F(f)(\varphi)=\varphi[f],
\]
on $\varphi[f]$ és la fórmula obtinguda reemplaçant dins $\varphi$ cada variable $x$ per la variable $f(x)$. Aquesta substitució és la mateixa operació que ja vam trobar a l'exemple~\ref{ex:deduibilitat-functors}, on una substitució de variables proposicionals induïa un functor entre categories de deduïbilitat perquè transformava deduccions; ara la mirem d'una altra manera, com l'acció sobre morfismes d'un functor $F\colon\cat{FinSet}\to\cat{Set}$. Formalment, $F(f)$ es defineix per recursió sobre l'estructura de $\varphi$: sobre una variable, $F(f)(x)=f(x)$, i $F(f)$ commuta amb les connectives.

El segon functor, el \emph{semàntic}, associa a cada conjunt de variables el conjunt de les funcions de veritat corresponents:
\[
G(X)=\{\text{funcions de veritat sobre }X\},
\]
on una \emph{funció de veritat} sobre $X$ és una aplicació
\[
\varphi\colon \{0,1\}^{X}\longrightarrow\{0,1\}
\]
que assigna un valor de veritat a cada valoració de les variables, entenent per \emph{valoració} un element $v\in\{0,1\}^{X}$, és a dir una aplicació $v\colon X\to\{0,1\}$ que fixa el valor de cada variable. Dit altrament, $G(X)=\{0,1\}^{(\{0,1\}^{X})}$. Sobre morfismes, un canvi de variables $f\colon X\to Y$ actua per \emph{precomposició en les valoracions}: tota valoració $w\colon Y\to\{0,1\}$ es restringeix a la valoració $w\comp f\colon X\to\{0,1\}$, i definim
\[
G(f)\colon G(X)\longrightarrow G(Y),
\qquad
G(f)(\theta)(w)=\theta(w\comp f),
\]
és a dir, $G(f)(\theta)=\theta\comp(-\comp f)$.

Tots dos són functors covariants $\cat{FinSet}\to\cat{Set}$. La comprovació és rutinària ---per recursió sobre les fórmules en el cas de $F$, i per l'associativitat de la composició en el de $G$--- i es fa a l'exercici resolt~\ref{pr-sintaxi-semantica} de la Pràctica.

A cada conjunt de variables $X$ li associem l'aplicació que interpreta cada fórmula com la seva funció de veritat (la seva taula de veritat):
\[
\eta_X\colon F(X)\longrightarrow G(X),
\qquad
\eta_X(\varphi)(v)=
\begin{cases}
1 & \text{si }v\models\varphi,\\
0 & \text{si }v\not\models\varphi,
\end{cases}
\]
on $v\models\varphi$ vol dir que la valoració $v$ satisfà la fórmula $\varphi$.

El fet clau és que aquesta interpretació és \emph{compatible amb el canvi de variables}: per a tot morfisme $f\colon X\to Y$ de $\cat{FinSet}$, el quadrat
\[
\begin{tikzpicture}[baseline=(m.center)]
  \matrix (m) [matrix of math nodes, row sep=2.8em, column sep=3.6em] {
    F(X) & F(Y) \\
    G(X) & G(Y) \\
  };
  \draw[-{Stealth[scale=1.1]}] (m-1-1) -- node[above] {$F(f)$} (m-1-2);
  \draw[-{Stealth[scale=1.1]}] (m-1-1) -- node[left] {$\eta_X$} (m-2-1);
  \draw[-{Stealth[scale=1.1]}] (m-1-2) -- node[right] {$\eta_Y$} (m-2-2);
  \draw[-{Stealth[scale=1.1]}] (m-2-1) -- node[below] {$G(f)$} (m-2-2);
\end{tikzpicture}
\]
commuta, és a dir $G(f)\comp\eta_X=\eta_Y\comp F(f)$. Avaluats sobre una fórmula $\varphi\in F(X)$ i una valoració $w\colon Y\to\{0,1\}$, els dos camins del quadrat diuen si $w\models\varphi[f]$ i si $(w\comp f)\models\varphi$, respectivament. Que coincideixin és exactament el \emph{lema de substitució} de la lògica proposicional,
\[
(w\comp f)\models\varphi
\quad\Longleftrightarrow\quad
w\models\varphi[f],
\]
que es demostra per recursió sobre $\varphi$; el càlcul complet és a l'exercici resolt~\ref{pr-sintaxi-semantica}.

En termes informals: substituir primer i interpretar després dona la mateixa funció de veritat que interpretar primer i transportar després. Per exemple, de $\varphi=p\land q$ i el canvi de variables $p\mapsto r,\ q\mapsto s$ obtenim $r\land s$, i tant se val si calculem la taula de veritat abans o després de substituir: el resultat és la mateixa funció (la que val $1$ només quan $r$ i $s$ són totes dues certes). Si el canvi identifica variables, diguem $p\mapsto r,\ q\mapsto r$, aleshores $\varphi$ es transforma en $r\land r$, amb taula de veritat la de $r$; també aquí interpretar i transportar coincideixen. Aquesta compatibilitat, per a tot canvi de variables alhora, és el que a la secció següent anomenarem la \textbf{naturalitat} de la família $(\eta_X)$: \emph{substituir variables i després calcular la funció de veritat és el mateix que calcular-la primer i transportar-la després al llarg del canvi de variables}.

Val la pena notar una diferència amb els exemples \enquote{d'identificació canònica} que veurem després (com el bidual d'un espai vectorial): aquí $\eta_X$ \emph{no} és injectiva ---dues fórmules diferents poden tenir la mateixa taula de veritat, com $p\lor q$ i $q\lor p$---, de manera que no estableix cap isomorfisme entre sintaxi i semàntica. El que expressa la naturalitat no és una identificació, sinó que la interpretació és compatible amb \emph{qualsevol} substitució o canvi de variables, incloses les substitucions que identifiquen variables. La mateixa idea s'estén, amb més maquinària, a la lògica de primer ordre, on les fórmules s'interpreten com els predicats que defineixen sobre una estructura fixada; ho veurem desplegat com a exemple formal a l'exemple~\ref{ex:interpretacio-primer-ordre} de la secció~\ref{sec:exemples-transf}.

\begin{observacio}[Per què estudiar lògica categòrica]
Un lector amb formació lògica pot objectar que, al capdavall, la feina l'ha feta el lema de substitució, i preguntar-se què hi guanya el llenguatge categòric. La resposta assenyala un canvi de punt de vista. Aquí hem hagut de definir la substitució per recursió i comprovar-ne la compatibilitat a mà perquè encara no disposem de cap maquinària; però la lògica categòrica reformula la substitució com una operació \emph{estructural} ---la composició, o més en general la reindexació al llarg d'un morfisme de la categoria de contextos--- i no com una manipulació sintàctica de fórmules. En aquest marc, la interpretació és un functor (o un morfisme entre estructures més riques) que \emph{preserva} aquella operació, i naturalitats com la d'aquest exemple deixen de demostrar-se cas per cas: s'obtenen com a corol·lari que la interpretació respecta l'estructura. La inducció sobre fórmules es fa aleshores una sola vegada i en abstracte, en establir que la categoria sintàctica té la propietat universal que la caracteritza. És el gir que fa que la lògica moderna no es defineixi per les seves peculiaritats sintàctiques, sinó per les estructures que preserva; el nostre lema de substitució n'és el cas visible més senzill.
\end{observacio}

\section{Definició de transformació natural}

Siguin $F,G\colon\cat{C}\to\cat{D}$ dos functors. Volem, per a cada objecte $A$ de $\cat{C}$, un morfisme de $\cat{D}$ que vagi de $F(A)$ a $G(A)$; i volem que aquests morfismes siguin compatibles amb l'acció de $F$ i $G$ sobre les fletxes.

\begin{definicio}\index{transformació natural}
\label{def:transf-natural}
Una \textbf{transformació natural} $\eta\colon F\Rightarrow G$ entre dos functors $F,G\colon\cat{C}\to\cat{D}$ és una \textbf{família de morfismes de $\cat{D}$}
\[
\eta=\bigl(\eta_A\bigr)_{A\in\Ob(\cat{C})},
\qquad
\eta_A\colon F(A)\to G(A),
\]
indexada pels objectes de $\cat{C}$ ---un morfisme $\eta_A$ per a cada objecte $A$, de $F(A)$ a $G(A)$, anomenat el \textbf{component} de $\eta$ a $A$---, subjecta a la condició següent: per a cada morfisme $f\colon A\to B$ de $\cat{C}$, el quadrat
\[
\begin{tikzpicture}[baseline=(m.center)]
  \matrix (m) [matrix of math nodes, row sep=2.8em, column sep=3.4em] {
    F(A) & F(B) \\
    G(A) & G(B) \\
  };
  \draw[-{Stealth[scale=1.1]}] (m-1-1) -- node[above] {$F(f)$} (m-1-2);
  \draw[-{Stealth[scale=1.1]}] (m-1-1) -- node[left] {$\eta_A$} (m-2-1);
  \draw[-{Stealth[scale=1.1]}] (m-1-2) -- node[right] {$\eta_B$} (m-2-2);
  \draw[-{Stealth[scale=1.1]}] (m-2-1) -- node[below] {$G(f)$} (m-2-2);
\end{tikzpicture}
\]
commuta; és a dir,
\[
G(f)\comp\eta_A=\eta_B\comp F(f).
\]
Aquesta igualtat s'anomena \textbf{condició de naturalitat}, i el quadrat, \textbf{quadrat de naturalitat}.
\end{definicio}

El quadrat de naturalitat exhibeix una transformació natural morfisme a morfisme, i és la manera pràctica de comprovar-la. Sovint, però, interessa representar tota la transformació $\eta$ d'un sol traç, com un objecte que \emph{relaciona} dos functors paral·lels. Dibuixarem, doncs, una transformació natural \emph{diagramàticament} d'aquesta manera: els dos functors $F,G\colon\cat{C}\to\cat{D}$ es representen com dues fletxes de $\cat{C}$ cap a $\cat{D}$, i la transformació natural com una fletxa doble $\Rightarrow$ que va de l'una cap a l'altra:
\[
\begin{tikzpicture}[baseline=(A.base)]
  \node (A) at (0,0) {$\cat{C}$};
  \node (B) at (2.9,0) {$\cat{D}$};
  \draw[cat] (A) to[bend left=26] node[above] {$F$} (B);
  \draw[cat] (A) to[bend right=26] node[below] {$G$} (B);
  \draw[cat nat] (1.45,0.2) -- node[right=1.5pt] {$\eta$} (1.45,-0.2);
\end{tikzpicture}
\]
És la primera vegada que apareix una fletxa doble entre fletxes: reflecteix que $\eta$ viu \enquote{un nivell més amunt} que els morfismes ordinaris. Aquest punt de vista ---objectes, functors i, damunt, transformacions naturals--- és el que més endavant resumirem dient que les categories formen una \emph{2-categoria}, i és la notació natural per a la composició horitzontal i per a la composició d'una transformació natural amb un functor que veurem al final del capítol.

Convé precisar bé què exigeix la condició de naturalitat i què no: les components $\eta_A$ es trien un per a cada objecte, però no lliurement, ja que han d'encaixar tots alhora amb l'acció dels functors sobre els morfismes. Ara bé, la naturalitat només demana que les components escollides siguin compatibles amb tots els morfismes; no demana que siguin l'única elecció possible. És una condició de compatibilitat, no d'unicitat ni d'absència de paràmetres, com mostra l'exemple següent.

Sobre els espais vectorials reals, prenem com a functors $F=G=\id_{\cat{Vect}_{\mathbb R}}$, el functor identitat de la categoria $\cat{Vect}_{\mathbb R}$, i busquem una transformació natural $\eta\colon\id\Rightarrow\id$. Fixem un escalar $\lambda\in\mathbb{R}$ i prenem, per a cada espai $V$, la component
\[
\eta_V=\lambda\,\id_V\colon V\to V,
\]
que és una aplicació lineal ---la multiplicació per l'escalar $\lambda$---, i per tant un morfisme de $\cat{Vect}_{\mathbb R}$, amb el tipus correcte perquè $F(V)=G(V)=V$. En efecte, aquesta família és una transformació natural: per a tota aplicació lineal $T\colon V\to W$ el quadrat de naturalitat és
\[
\begin{tikzpicture}[baseline=(m.center)]
  \matrix (m) [matrix of math nodes, row sep=2.8em, column sep=3.8em] {
    V & W \\
    V & W \\
  };
  \draw[-{Stealth[scale=1.1]}] (m-1-1) -- node[above] {$T$} (m-1-2);
  \draw[-{Stealth[scale=1.1]}] (m-1-1) -- node[left] {$\lambda\,\id_V$} (m-2-1);
  \draw[-{Stealth[scale=1.1]}] (m-1-2) -- node[right] {$\lambda\,\id_W$} (m-2-2);
  \draw[-{Stealth[scale=1.1]}] (m-2-1) -- node[below] {$T$} (m-2-2);
\end{tikzpicture}
\]
i commuta, perquè
\[
T\comp(\lambda\,\id_V)=\lambda\,T=(\lambda\,\id_W)\comp T,
\]
on les dues fletxes horitzontals són $T$ perquè el functor identitat actua sobre $T$ deixant-lo igual. Quan $\lambda\neq 0$, cada component $\eta_V=\lambda\,\id_V$ és a més un isomorfisme d'espais vectorials, amb invers $\lambda^{-1}\id_V$; això farà de $\eta$ un exemple del que més endavant anomenarem un \emph{isomorfisme natural}. Tenim, doncs, tota una família de transformacions naturals, una per cada valor de $\lambda$: la naturalitat es compleix per a qualsevol $\lambda$ fixat, tot i que la construcció depèn d'aquesta elecció. Naturalitat no vol dir, per tant, \enquote{sense eleccions}, sinó \enquote{compatible amb tots els morfismes}.


\begin{exemple}[Una família que \emph{no} és natural]\label{ex:no-natural}\index{transformació natural!contraexemple}
És igual d'instructiu veure una família de components que \emph{falla} la condició de naturalitat. Treballem a $\cat{Set}$ amb $F=G=\id_{\cat{Set}}$, de manera que una transformació $\id\Rightarrow\id$ hauria de donar, per a cada conjunt $A$, una aplicació $\eta_A\colon A\to A$ que commuti amb \emph{totes} les aplicacions. Definim una família completament explícita, sense cap tria:
\[
\eta_A=\id_A\quad\text{si } A\neq\{0,1\},
\qquad\qquad
\eta_{\{0,1\}}=\tau,\quad \tau(0)=1,\ \tau(1)=0.
\]
És a dir, totes les components són identitats excepte la del conjunt $\{0,1\}$, que és la transposició dels dos elements. És una família de morfismes $\eta_A\colon A\to A$ indexada per \emph{tots} els conjunts, i cada component és fins i tot una bijecció.

Vegem què demana el quadrat de naturalitat per a una aplicació $f\colon A\to B$. Com que $F=G=\id_{\cat{Set}}$, es té $F(f)=G(f)=f$, i la condició és $f\comp\eta_A=\eta_B\comp f$. Prenem $A=B=\{0,1\}$ i l'aplicació constant de valor $0$, $f(0)=f(1)=0$:
\[
\begin{tikzpicture}[baseline=(m.center)]
  \matrix (m) [matrix of math nodes, row sep=2.8em, column sep=3.8em] {
    \{0,1\} & \{0,1\} \\
    \{0,1\} & \{0,1\} \\
  };
  \draw[-{Stealth[scale=1.1]}] (m-1-1) -- node[above] {$f$} (m-1-2);
  \draw[-{Stealth[scale=1.1]}] (m-1-1) -- node[left] {$\tau$} (m-2-1);
  \draw[-{Stealth[scale=1.1]}] (m-1-2) -- node[right] {$\tau$} (m-2-2);
  \draw[-{Stealth[scale=1.1]}] (m-2-1) -- node[below] {$f$} (m-2-2);
\end{tikzpicture}
\qquad
f\comp\tau=f,\qquad \tau\comp f=\text{constant de valor }1.
\]
El camí de l'esquerra i a baix dona $f\comp\tau$, que continua sent constant de valor $0$; el de dalt i a la dreta dona $\tau\comp f$, que és constant de valor $1$. Els dos camins difereixen, i la família $(\eta_A)$ \emph{no} és una transformació natural.

La lliçó és que tenir una regla objecte a objecte, fins i tot formada per bijeccions, no basta. La component a $\{0,1\}$ s'ha definit sense tenir en compte les aplicacions que entren i surten d'aquest conjunt, i una sola d'aquestes aplicacions n'hi ha prou per trencar la compatibilitat. La naturalitat és exactament aquesta compatibilitat amb \emph{totes} les fletxes.
\end{exemple}

\section{Exemples de transformacions naturals}\label{sec:exemples-transf}

\begin{exemple}[Transformació identitat]\index{transformació natural!identitat}
Per a tot functor $F\colon\cat{C}\to\cat{D}$ hi ha la \textbf{transformació natural identitat} $\id_F\colon F\Rightarrow F$, amb components $(\id_F)_A=\id_{F(A)}$. El quadrat de naturalitat commuta trivialment, perquè els dos camins són $F(f)$.
\end{exemple}

\begin{exemple}[Preordres]\index{transformació natural!entre preordres}
Siguin $P,Q$ dos preordres i $F,G\colon P\to Q$ dos functors, és a dir, dues aplicacions monòtones. Una transformació natural $\eta\colon F\Rightarrow G$ dona, per a cada $x\in P$, un morfisme $\eta_x\colon F(x)\to G(x)$; però a $Q$ hi ha una fletxa $F(x)\to G(x)$ exactament quan $F(x)\leq G(x)$. La condició de naturalitat es compleix automàticament, perquè entre dos objectes d'un preordre hi ha com a màxim una fletxa. Per tant,
\[
\text{existeix }\eta\colon F\Rightarrow G
\quad\Longleftrightarrow\quad
F(x)\leq G(x)\ \text{per a tot }x.
\]
En el llenguatge de l'ordre, una transformació natural entre aplicacions monòtones és exactament la relació \enquote{$F$ és puntualment menor o igual que $G$}.
\end{exemple}

\begin{exemple}[Ampliació: interpretació de fórmules de primer ordre]\label{ex:interpretacio-primer-ordre}\index{transformació natural!interpretació lògica}\index{lema de substitució}
Reprenem el fil de la lògica que hem obert a la secció~\ref{sec:cap-transformacio} (\enquote{Cap al concepte de transformació natural}), ara en primer ordre i amb la naturalitat plantejada com a exemple formal. Fixem un llenguatge de primer ordre i una estructura $M$. L'exemple pressuposa la semàntica de la lògica de primer ordre (estructures, satisfacció, variables lliures i lligades) i es pot ometre en una primera lectura.

Convé fixar d'entrada les convencions que fan functorial la construcció sintàctica, en lloc d'adoptar-les a posteriori. Disposem d'una reserva infinita de variables; un \emph{context} és un conjunt finit $X$ de variables d'aquesta reserva, i escrivim $\cat{Ctx}$ per a la subcategoria plena de $\cat{FinSet}$ que té per objectes els contextos i per morfismes totes les aplicacions entre ells (com que la reserva és infinita, $\cat{Ctx}$ és equivalent a $\cat{FinSet}$, però no farem servir aquesta equivalència); treballem amb fórmules \emph{mòdul reanomenament de variables lligades} ($\alpha$-equivalència), i totes les substitucions s'entenen \emph{lliures de captura}. Sense aquestes convencions les lleis functorials de sota no se sostenen: en $\varphi=\exists y\,R(x,y)$, substituir la variable lliure $x$ per $y$ no pot donar ingènuament $\exists y\,R(y,y)$; cal reanomenar abans la variable lligada i obtenir, per exemple, $\exists z\,R(y,z)$.

Amb aquestes convencions considerem dues construccions. Per una banda,
\[
F(X)=\{\text{fórmules amb variables lliures entre }X\},
\]
enteses com a classes d'$\alpha$-equivalència; un canvi de variables $f\colon X\to Y$ hi actua per substitució, $F(f)(\varphi)=\varphi[f]$, reemplaçant cada variable lliure $x$ per $f(x)$. Per l'altra,
\[
G(X)=\{\text{predicats definibles sobre }X\},
\]
on un \emph{predicat definible} és un subconjunt de $M^X$ (el conjunt d'assignacions $s\colon X\to M$) de la forma
\[
\llbracket\varphi\rrbracket=\{\,s\colon X\to M\mid M\models\varphi[s]\,\},
\]
és a dir, l'\emph{extensió} d'alguna fórmula. El canvi de variables $f$ actua sobre $G$ transportant assignacions: com que $f$ indueix la precomposició $M^{f}\colon M^{Y}\to M^{X}$, $s\mapsto s\comp f$, posem
\[
G(f)(S)=\bigl(M^{f}\bigr)^{-1}(S)=\{\,s\in M^{Y}\mid s\comp f\in S\,\},\qquad S\subseteq M^{X}.
\]
Totes dues són functors covariants $\cat{Ctx}\to\cat{Set}$: per a $F$ es comprova que la substitució respecta identitats i composició\footnotemark; per a $G$, la functorialitat és la de la precomposició $M^{(-)}$ seguida de preimatge.

A cada context $X$ li associem ara l'aplicació que interpreta cada fórmula com la seva extensió,
\[
\eta_X\colon F(X)\to G(X),
\qquad
\eta_X(\varphi)=\llbracket\varphi\rrbracket.
\]
Aquesta $\eta_X$ està ben definida sobre classes d'$\alpha$-equivalència, perquè dues fórmules que només difereixen en el nom de les variables lligades tenen la mateixa extensió: la semàntica és $\alpha$-invariant, i per això el pas al quocient no altera res. La condició de naturalitat, per a un canvi de variables $f\colon X\to Y$, és que interpretar la fórmula substituïda coincideixi amb transportar l'extensió de la fórmula original:
\[
\llbracket\varphi[f]\rrbracket=\{\,s\colon Y\to M\mid s\comp f\in\llbracket\varphi\rrbracket\,\}.
\]
Aquesta igualtat és exactament el \textbf{lema de substitució} de la teoria de models ---el significat d'una fórmula després de canviar-ne les variables es calcula transportant el significat original--- i fa doble feina: garanteix alhora que $G(f)$ envia predicats definibles a predicats definibles (la preimatge d'un definible és definible) i que el quadrat de naturalitat commuta. Per tant $\eta\colon F\Rightarrow G$ és una transformació natural, i expressa que la interpretació és compatible amb la substitució.

Com en el cas proposicional de la secció~\ref{sec:cap-transformacio}, $\eta_X$ no és injectiva ---dues fórmules lògicament equivalents defineixen el mateix predicat---, de manera que la naturalitat no és aquí una identificació, sinó que expressa que la interpretació és compatible amb els canvis de variables, inclosos els reanomenaments i les identificacions. Convé, finalment, no confondre dues nocions que aquí conviuen: en una estructura $M$ fixada, dues fórmules poden tenir la mateixa extensió sense ser lògicament equivalents en \emph{totes} les estructures; la relació d'\emph{equivalència semàntica relativa a $M$} és, per tant, més grollera que l'\emph{equivalència lògica} global.
\end{exemple}
\footnotetext{És on es fan servir les convencions anteriors. Quan $f$ no és injectiva, identifica variables lliures, i és la substitució lliure de captura, definida sobre classes d'$\alpha$-equivalència, la que fa que $\varphi[\id]=\varphi$ i $\varphi[g\comp f]=(\varphi[f])[g]$ valguin com a \emph{igualtats}, no només mòdul reanomenament ad hoc.}%

\begin{exemple}[Bidual d'espais vectorials]\label{ex:bidual}\index{transformació natural!bidual}
Reprenem el fil que vam deixar obert a l'exemple~\ref{ex:homfunctor-dual} (\enquote{El dual d'un espai vectorial com a homfunctor}), on vam construir el \textbf{functor dual}
\[
(-)^*=\Hom_{\mathbb K}(-,\mathbb K)\colon\cat{Vect}_{\mathbb K}\op\to\cat{Vect}_{\mathbb K},\qquad V\mapsto V^*,
\]
contravariant, que sobre un morfisme $T\colon V\to W$ dona l'aplicació dual $T^*\colon W^*\to V^*$, $\psi\mapsto\psi\comp T$, amb la regla $(g\comp f)^*=f^*\comp g^*$. Vam ajornar deliberadament la comparació de $V$ amb el bidual fins a tenir el llenguatge de les transformacions naturals; ja el tenim.

El \textbf{functor bidual} (o doble dual) s'obté per composició del dual amb ell mateix. Com que $(-)^*$ és contravariant ---un functor amb domini $\cat{Vect}_{\mathbb K}\op$---, per compondre'l amb ell mateix cal fer encaixar dominis i codominis amb l'oposat, exactament com a la proposició~\ref{prop:variancia-composicio}: es pren el functor oposat $\bigl((-)^*\bigr)\op\colon\cat{Vect}_{\mathbb K}\to\cat{Vect}_{\mathbb K}\op$ ---que actua igual sobre objectes--- i s'hi torna a aplicar el dual,
\[
(-)^{**}=(-)^*\comp\bigl((-)^*\bigr)\op\colon\cat{Vect}_{\mathbb K}\to\cat{Vect}_{\mathbb K}.
\]
Per la regla dels signes, la composició de dos functors contravariants és \emph{covariant}. Desplegada, actua sobre objectes com
\[
V\mapsto V^{**}=(V^*)^*,
\]
i sobre un morfisme $T\colon V\to W$ com el dual aplicat dues vegades,
\[
T^{**}=(T^*)^*\colon V^{**}\to W^{**},\qquad \rho\mapsto\rho\comp T^*,
\]
amb la regla $(g\comp f)^{**}=g^{**}\comp f^{**}$: invertir l'ordre dos cops el restaura, i per això el bidual és un endofunctor covariant de $\cat{Vect}_{\mathbb K}$, paral·lel a la identitat.

L'efecte sobre els morfismes ho resumeix el diagrama següent, on el dual inverteix el sentit de la fletxa i el bidual, aplicant-lo dues vegades, el recupera:
\[
\begin{tikzpicture}[baseline=(m.center)]
  \matrix (m) [matrix of math nodes, column sep=3.2em, row sep=3.0em] {
    {} & V & W & {} \\
    {} & V^* & W^* & {} \\
    {} & V^{**} & W^{**} & {} \\
  };
  \draw[cat] (m-1-2) -- node[above] {$T$} (m-1-3);
  \draw[cat] (m-2-3) -- node[above] {$T^*$} (m-2-2);
  \draw[cat] (m-3-2) -- node[above] {$T^{**}$} (m-3-3);
  \draw[cat,dashed] (m-1-1) -- node[left] {$(-)^*$} (m-2-1);
  \draw[cat,dashed] (m-2-1) -- node[left] {$(-)^*$} (m-3-1);
  \draw[cat,dashed] (m-1-4) -- node[right] {$(-)^{**}$} (m-3-4);
\end{tikzpicture}
\]
Les fletxes contínues són morfismes de $\cat{Vect}_{\mathbb K}$; les puntejades, l'acció del functor. La fila del mig apunta en sentit contrari ---el dual $(-)^*$ és contravariant---, i la de baix en recupera el sentit ---dues inversions es cancel·len---: el bidual $(-)^{**}$, composició del dual amb ell mateix, és covariant, paral·lel a la identitat. Ara sí que té sentit comparar-los amb una transformació natural. Tots dos ---la identitat i el bidual--- són endofunctors covariants de $\cat{Vect}_{\mathbb K}$, de manera que la comparació és una transformació natural $\eta$ entre functors paral·lels:
\[
\begin{tikzpicture}[baseline=(A.base)]
  \node (A) at (0,0) {$\cat{Vect}_{\mathbb K}$};
  \node (B) at (3.6,0) {$\cat{Vect}_{\mathbb K}$};
  \draw[cat] (A) to[bend left=26] node[above] {$\id_{\cat{Vect}_{\mathbb K}}$} (B);
  \draw[cat] (A) to[bend right=26] node[below] {$(-)^{**}$} (B);
  \draw[cat nat] (1.8,0.2) -- node[right=1.5pt] {$\eta$} (1.8,-0.2);
\end{tikzpicture}
\]
El component de $\eta$ a cada espai $V$ és l'\textbf{avaluació}, que envia un vector $v$ a la forma \enquote{avaluar en $v$}. La transformació natural és, doncs, la família
\[
\eta=\bigl(\eta_V\bigr)_{V\in\Ob(\cat{Vect}_{\mathbb K})},
\]
i cada component es desglossa nivell a nivell, amb la imatge de $v$ com una aplicació de $V^*$ en $\mathbb K$:
\[
\begin{array}{r c l}
\eta_V\colon V & \longrightarrow & V^{**}\\[6pt]
v & \longmapsto & \eta_V(v)\colon
  \begin{array}[t]{r c l}
    V^* & \longrightarrow & \mathbb K\\[4pt]
    \varphi & \longmapsto & \varphi(v).
  \end{array}
\end{array}
\]
Cada $\eta_V$ és lineal i està ben definida. La condició de naturalitat, per a una aplicació lineal $T\colon V\to W$, és la commutativitat del quadrat
\[
\begin{tikzpicture}[baseline=(m.center)]
  \matrix (m) [matrix of math nodes, row sep=2.8em, column sep=3.6em] {
    V & W \\
    V^{**} & W^{**} \\
  };
  \draw[-{Stealth[scale=1.1]}] (m-1-1) -- node[above] {$T$} (m-1-2);
  \draw[-{Stealth[scale=1.1]}] (m-1-1) -- node[left] {$\eta_V$} (m-2-1);
  \draw[-{Stealth[scale=1.1]}] (m-1-2) -- node[right] {$\eta_W$} (m-2-2);
  \draw[-{Stealth[scale=1.1]}] (m-2-1) -- node[below] {$T^{**}$} (m-2-2);
\end{tikzpicture}
\]
és a dir, $T^{**}\comp\eta_V=\eta_W\comp T$: prendre l'avaluació i després transportar-la pel bidual coincideix amb aplicar $T$ i després avaluar. Aquí la fletxa horitzontal superior és $T$ perquè el functor identitat actua sobre $T$ deixant-lo igual. En aquest cas, la naturalitat es pot llegir dient que la identificació de $V$ amb la seva imatge dins $V^{**}$ no fa servir cap elecció de base: la construcció $v\mapsto(\varphi\mapsto\varphi(v))$ és intrínseca. Aquesta lectura ---\enquote{independent de la base}--- és una interpretació d'aquest exemple concret, no la definició de naturalitat, que continua sent la commutativitat dels quadrats. En restringir-se als espais de dimensió finita, cada $\eta_V$ és un isomorfisme, cosa que reprendrem en parlar dels isomorfismes naturals. En canvi, la comparació de $V$ amb el \emph{dual} $V^*$ ---i no amb el bidual--- no es pot fer natural; ho precisarem a l'observació~\ref{obs:iso-no-natural}, un cop definits els isomorfismes naturals.
\end{exemple}

\section{Categoria de functors}\label{sec:categoria-functors}

Els functors entre dues categories fixes, amb les transformacions naturals com a morfismes, formen al seu torn una categoria. Aquesta secció en desenvolupa la construcció: la composició que la fa possible, la categoria $\cat{D}^{\cat{C}}$ mateixa, un exemple fonamental ---els grafs--- i els seus isomorfismes, que resulten ser els isomorfismes naturals; en tanquem l'estudi amb l'equivalència de categories.

Les transformacions naturals no només relacionen functors: també es poden \emph{compondre} entre elles. La composició bàsica apila una transformació damunt d'una altra, i és la que després convertirà els functors entre dues categories en una categoria, i la que dona sentit a parlar de la \emph{inversa} d'un isomorfisme natural.

\begin{definicio}[Composició vertical]\index{transformació natural!composició vertical}
Donades dues transformacions naturals $\eta\colon F\Rightarrow G$ i $\theta\colon G\Rightarrow H$ entre functors $\cat{C}\to\cat{D}$, la seva \textbf{composició vertical} $\theta\comp\eta\colon F\Rightarrow H$ té per components
\[
(\theta\comp\eta)_A=\theta_A\comp\eta_A.
\]
El nom \emph{vertical} ve d'aquesta representació: les dues transformacions s'apilen ---$\eta$ a dalt i $\theta$ a sota--- entre els tres functors paral·lels $F,G,H\colon\cat{C}\to\cat{D}$, i la composició $\theta\comp\eta$ les recorre de dalt a baix.
\[
\begin{tikzpicture}[baseline=(A.base)]
  \node (A) at (0,0) {$\cat{C}$};
  \node (B) at (3.4,0) {$\cat{D}$};
  \draw[cat] (A) to[bend left=44] node[above] {$F$} (B);
  \draw[cat] (A) -- (B);
  \draw[cat] (A) to[bend right=44] node[below] {$H$} (B);
  \node[fill=white,inner sep=1.5pt] at (1.7,0) {$G$};
  \draw[cat nat] (1.7,0.52) -- node[right=1.5pt] {$\eta$} (1.7,0.27);
  \draw[cat nat] (1.7,-0.27) -- node[right=1.5pt] {$\theta$} (1.7,-0.52);
\end{tikzpicture}
\]
El mateix nom es reconeix a escala de components: fixat un morfisme $f\colon A\to B$, les components de $\eta$ i de $\theta$ són fletxes \emph{verticals}, i els seus quadrats de naturalitat s'apilen l'un damunt de l'altre ---compartint la fila central $G(A)\to G(B)$---, de manera que $\theta\comp\eta$ en compon les columnes:
\[
\begin{tikzpicture}[baseline=(m.center)]
  \matrix (m) [matrix of math nodes, row sep=3em, column sep=3.6em] {
    F(A) & F(B) \\
    G(A) & G(B) \\
    H(A) & H(B) \\
  };
  \draw[cat] (m-1-1) -- node[above] {$F(f)$} (m-1-2);
  \draw[cat] (m-2-1) -- node[above] {$G(f)$} (m-2-2);
  \draw[cat] (m-3-1) -- node[below] {$H(f)$} (m-3-2);
  \draw[cat] (m-1-1) -- node[left] {$\eta_A$} (m-2-1);
  \draw[cat] (m-2-1) -- node[left] {$\theta_A$} (m-3-1);
  \draw[cat] (m-1-2) -- node[right] {$\eta_B$} (m-2-2);
  \draw[cat] (m-2-2) -- node[right] {$\theta_B$} (m-3-2);
\end{tikzpicture}
\]
Aïllant-ne el rectangle exterior, la composició $\theta\comp\eta\colon F\Rightarrow H$ té el seu propi quadrat de naturalitat (a l'esquerra), i es representa com una única cel·la de $F$ cap a $H$ (a la dreta):
\[
\begin{tikzpicture}[baseline=(m.center)]
  \matrix (m) [matrix of math nodes, row sep=2.8em, column sep=3.6em] {
    F(A) & F(B) \\
    H(A) & H(B) \\
  };
  \draw[cat] (m-1-1) -- node[above] {$F(f)$} (m-1-2);
  \draw[cat] (m-1-1) -- node[left] {$(\theta\comp\eta)_A$} (m-2-1);
  \draw[cat] (m-1-2) -- node[right] {$(\theta\comp\eta)_B$} (m-2-2);
  \draw[cat] (m-2-1) -- node[below] {$H(f)$} (m-2-2);
\end{tikzpicture}
\qquad\qquad
\begin{tikzpicture}[baseline=(A.base)]
  \node (A) at (0,0) {$\cat{C}$};
  \node (B) at (3.7,0) {$\cat{D}$};
  \draw[cat] (A) to[bend left=26] node[above] {$F$} (B);
  \draw[cat] (A) to[bend right=26] node[below] {$H$} (B);
  \draw[cat nat] (1.85,0.2) -- node[right=1.5pt] {$\theta\comp\eta$} (1.85,-0.2);
\end{tikzpicture}
\]
\end{definicio}

\begin{proposicio}
La composició vertical torna a ser una transformació natural, és associativa, i la transformació identitat $\id_F$ n'és l'element neutre.
\end{proposicio}

\begin{prova}
Per a la naturalitat n'hi ha prou d'observar que, al diagrama de components de la definició, tots dos quadrats commuten (naturalitat de $\eta$ i de $\theta$); per tant també commuta el rectangle exterior, les columnes del qual són $(\theta\comp\eta)_A$ i $(\theta\comp\eta)_B$, i això és el quadrat de naturalitat de $\theta\comp\eta$ ---equivalentment, $H(f)\comp\theta_A\comp\eta_A=\theta_B\comp\eta_B\comp F(f)$, encadenant els dos quadrats. L'associativitat i la propietat de la identitat es redueixen a les propietats corresponents de la composició a $\cat{D}$, aplicades component a component.
\end{prova}

La composició vertical de transformacions naturals, introduïda més amunt, organitza els functors entre dues categories en una nova categoria.

\begin{definicio}\index{categoria!de functors}\index{DC@$\cat{D}^{\cat{C}}$}
Sigui $\cat{C}$ una categoria petita i $\cat{D}$ una categoria qualsevol. La \textbf{categoria de functors} $\cat{D}^{\cat{C}}$ (també escrita $[\cat{C},\cat{D}]$ o $\operatorname{Fun}(\cat{C},\cat{D})$) té:
\begin{itemize}
\item per objectes, els functors $\cat{C}\to\cat{D}$;
\item per morfismes, les transformacions naturals entre ells;
\item per composició, la composició vertical, i per identitats, les transformacions identitat.
\end{itemize}
\end{definicio}

\begin{observacio}[Qüestions de mida]\label{obs:mida-functors}\index{transformació natural!Nat@$\operatorname{Nat}$}
Convé dir per què la definició demana $\cat{C}$ petita. La raó de fons és la mateixa de l'observació~\ref{obs:mida-cat}, on vam haver de restringir $\cat{Cat}$ a les categories \emph{petites}. Si $\cat{C}$ fos gran, les dades que defineixen un functor $\cat{C}\to\cat{D}$ podrien formar una classe pròpia ---l'assignació recorreria una classe pròpia d'objectes---, i les classes pròpies no poden ser elements de cap col·lecció; aleshores els functors $\cat{C}\to\cat{D}$ no es poden reunir com els \emph{objectes} d'una categoria.\footnotemark Amb $\cat{C}$ petita, en canvi, aquestes dades formen un conjunt i la construcció no té cap problema. Això no impedeix parlar \emph{individualment} de functors i transformacions naturals amb domini gran ---ho fem contínuament, com amb $\cat{Set}\to\cat{Set}$---; el que no fem és agrupar-los tots en una categoria de functors.

Amb $\cat{C}$ petita, $\cat{D}^{\cat{C}}$ pot ser encara gran, però és \textbf{localment petita} tan bon punt $\cat{D}$ ho és: una transformació natural $F\Rightarrow G$ queda determinada pels seus components $\eta_A$, un per a cada objecte $A$ de $\cat{C}$, de manera que la col·lecció de totes les transformacions naturals $F\Rightarrow G$, que denotem $\operatorname{Nat}(F,G)$, és un subconjunt del producte
\[
\operatorname{Nat}(F,G)\;\subseteq\;\prod_{A\in\Ob(\cat{C})}\Hom_{\cat{D}}(F(A),G(A)),
\]
indexat per $\Ob(\cat{C})$ ---un conjunt, precisament perquè $\cat{C}$ és petita---, pel mateix argument que feia $\cat{Cat}$ localment petita. Com que aquest producte és un conjunt, $\operatorname{Nat}(F,G)$ també ho és, i és el conjunt de morfismes de $F$ a $G$ a la categoria de functors:
\[
\operatorname{Nat}(F,G)=\Hom_{\cat{D}^{\cat{C}}}(F,G).
\]
\end{observacio}
\footnotetext{Això depèn del marc de mida, que aquí és el de classes i conjunts de l'apèndix~\ref{cap:mida}. Amb universos de Grothendieck, la categoria de functors d'una categoria gran es pot formar en un univers superior; vegeu \cite[§1.1]{riehl}.}%

\index{graf!com a categoria de functors}\index{categoria!de functors!grafs}
Un dels exemples més il·luminatius de categoria de functors mostra que un objecte aparentment aliè a les categories ---un graf dirigit--- \emph{és}, en realitat, un functor, i que els seus morfismes són transformacions naturals. Aquest exemple reprèn el fil dels grafs que hem seguit des del primer capítol i és, a més, el primer indici que categories de functors del tipus $\cat{Set}^{\cat{C}}$ codifiquen estructures matemàtiques familiars.

La categoria índex $\cat{P}$ no és nova: és la categoria de \textbf{dues fletxes paral·leles} de la subsecció~\ref{subsec:finites}, amb dos objectes i dues fletxes no trivials amb el mateix domini i el mateix codomini. Els seus dos objectes són, en si, abstractes; els rebategem $E$ i $V$ ---i les fletxes, $s$ i $t$--- per anticipar la lectura com a grafs. Com que una aresta apunta cap als seus extrems, el domini de les fletxes és l'objecte de les \emph{arestes}, $E$, i el codomini el dels \emph{vèrtexs}, $V$, amb $s,t\colon E\to V$:
\[
\begin{tikzpicture}[baseline=(E.base)]
  \node (E) at (0,0) {$E$};
  \node (V) at (2.4,0) {$V$};
  \draw[-{Stealth[scale=1.1]}] ([yshift=3pt]E.east) -- node[above] {$s$} ([yshift=3pt]V.west);
  \draw[-{Stealth[scale=1.1]}] ([yshift=-3pt]E.east) -- node[below] {$t$} ([yshift=-3pt]V.west);
\end{tikzpicture}
\]

\begin{proposicio}\index{graf!com a categoria de functors}
La categoria de functors $\cat{Set}^{\cat{P}}$ és isomorfa a la categoria $\cat{Graph}$ dels grafs dirigits i els morfismes de grafs.
\end{proposicio}

\begin{prova}
Un functor $G\colon\cat{P}\to\cat{Set}$ consisteix en dos conjunts $G(E)$ i $G(V)$ i dues aplicacions $G(s),G(t)\colon G(E)\to G(V)$. Interpretant $G(E)$ com el conjunt d'\textbf{arestes}, $G(V)$ com el conjunt de \textbf{vèrtexs} i $G(s),G(t)$ com les aplicacions \enquote{origen} i \enquote{destí}, aquesta dada és \emph{exactament} un graf dirigit. Recíprocament, tot graf dirigit determina un tal functor. Els objectes de $\cat{Set}^{\cat{P}}$ ---que són functors $\cat{P}\to\cat{Set}$--- es corresponen, doncs, bijectivament amb els grafs dirigits.

Vegem ara els morfismes. Una transformació natural $\varphi\colon G\Rightarrow H$ entre dos d'aquests functors té dos components, $\varphi_E\colon G(E)\to H(E)$ i $\varphi_V\colon G(V)\to H(V)$, i la condició de naturalitat per a les fletxes $s$ i $t$ dona els dos quadrats
\[
\begin{tikzpicture}[baseline=(m.center)]
  \matrix (m) [matrix of math nodes, row sep=2.4em, column sep=3em] {
    G(E) & G(V) \\
    H(E) & H(V) \\
  };
  \draw[-{Stealth[scale=1.1]}] (m-1-1) -- node[above] {$G(s)$} (m-1-2);
  \draw[-{Stealth[scale=1.1]}] (m-1-1) -- node[left] {$\varphi_E$} (m-2-1);
  \draw[-{Stealth[scale=1.1]}] (m-1-2) -- node[right] {$\varphi_V$} (m-2-2);
  \draw[-{Stealth[scale=1.1]}] (m-2-1) -- node[below] {$H(s)$} (m-2-2);
\end{tikzpicture}
\qquad\text{i}\qquad
\begin{tikzpicture}[baseline=(n.center)]
  \matrix (n) [matrix of math nodes, row sep=2.4em, column sep=3em] {
    G(E) & G(V) \\
    H(E) & H(V) \\
  };
  \draw[-{Stealth[scale=1.1]}] (n-1-1) -- node[above] {$G(t)$} (n-1-2);
  \draw[-{Stealth[scale=1.1]}] (n-1-1) -- node[left] {$\varphi_E$} (n-2-1);
  \draw[-{Stealth[scale=1.1]}] (n-1-2) -- node[right] {$\varphi_V$} (n-2-2);
  \draw[-{Stealth[scale=1.1]}] (n-2-1) -- node[below] {$H(t)$} (n-2-2);
\end{tikzpicture}
\]
Aquests quadrats diuen que $\varphi_E$ i $\varphi_V$ respecten l'origen i el destí de cada aresta; és a dir, són precisament la definició de \textbf{morfisme de grafs}. Per tant els morfismes de $\cat{Set}^{\cat{P}}$ es corresponen amb els morfismes de grafs. Aquesta correspondència, sobre objectes i sobre morfismes, és un parell de functors inversos l'un de l'altre, i ho és \emph{estrictament}: les seves composicions són les identitats \emph{exactes} ---no només isomorfismes naturals---, perquè un functor $\cat{P}\to\cat{Set}$ no és res més que els seus valors sobre $E$, $V$, $s$ i $t$; anar del functor al graf i tornar dona el mateix functor, i del graf al functor i tornar, el mateix graf. És, doncs, un isomorfisme de categories en el sentit estricte de la definició~\ref{def:iso-categories} ---més fort que l'equivalència que veurem tot seguit---, i $\cat{Set}^{\cat{P}}\cong\cat{Graph}$. El caràcter estricte d'aquest isomorfisme descansa sobre la codificació que hem triat: en identificar un graf amb les dades $(E,V,s,t)$ ---exactament les que porta un functor $\cat{P}\to\cat{Set}$---, l'anada i la tornada retornen les \emph{mateixes} dades. Amb altres codificacions habituals dels grafs, la correspondència continua sent un isomorfisme \emph{canònic} de categories, però no una igualtat literal de dades; en cap cas \enquote{exactes} no s'ha de llegir com una igualtat fundacional independent de la codificació.
\end{prova}

\begin{observacio}
Aquest exemple té una lectura de fons: un cop decidim representar els grafs com a functors $\cat{P}\to\cat{Set}$, la representació \emph{justifica i recupera canònicament} la definició habitual de morfisme de grafs, perquè els morfismes de $\cat{Set}^{\cat{P}}$ són, per força, les transformacions naturals. És un criteri de coherència que sovint es fa servir en teoria de categories: quan una estructura s'expressa com a categoria de functors, la noció de morfisme que se'n dedueix hauria de coincidir amb la que s'havia adoptat independentment. Les categories del tipus $\cat{Set}^{\cat{C}}$, que codifiquen estructures com els grafs, inclouen com a cas particular les categories de prefeixos, que presentem tot seguit i que seran fonamentals per al lema de Yoneda i, més endavant, per a la semàntica categòrica.
\end{observacio}

\begin{exemple}[La categoria de prefeixos]\label{ex:categoria-prefeixos}\index{prefeix@prefeix!categoria de}\index{prefeix@prefeix!morfisme de}
Sigui $\cat{C}$ una categoria petita. Els prefeixos sobre $\cat{C}$ (definició~\ref{def:prefeix}) són els objectes de la categoria de functors
\[
\widehat{\cat{C}}\coloneqq\cat{Set}^{\cat{C}\op},
\]
que anomenem \textbf{categoria de prefeixos} sobre $\cat{C}$. Els seus morfismes, els \textbf{morfismes de prefeixos} $\sigma\colon P\Rightarrow Q$, són les transformacions naturals: famílies d'aplicacions $\sigma_A\colon P(A)\to Q(A)$ tals que, per a cada morfisme $f\colon A\to B$ de $\cat{C}$,
\[
\sigma_A\comp P(f)=Q(f)\comp\sigma_B\colon P(B)\to Q(A),
\]
que és el quadrat de naturalitat escrit per a prefeixos. Com que $\cat{Set}$ és localment petita, $\widehat{\cat{C}}$ també ho és, per l'observació sobre qüestions de mida. La categoria $\cat{Set}^{\cat{P}}$ dels grafs n'és un exemple: com que un functor $\cat{P}\to\cat{Set}$ és un functor $(\cat{P}\op)\op\to\cat{Set}$, els grafs són els prefeixos sobre $\cat{P}\op$.
\end{exemple}

\begin{observacio}[La forma dels grafs és, ella mateixa, un graf]\index{categoria!lliure}
La categoria índex $\cat{P}$ té una lectura circular ben curiosa. El seu \emph{graf subjacent} ---els dos objectes com a vèrtexs i les dues fletxes generadores $s,t$ com a arestes--- és exactament $\bullet\rightrightarrows\bullet$, i $\cat{P}$ n'és la \textbf{categoria lliure} (secció~\ref{sec:grafs-lliures}): com que les dues fletxes paral·leles no es poden encadenar ---de l'objecte d'arribada no en surt cap---, no hi ha camins de longitud $\geq 2$, i els únics morfismes de $\cat{P}$ són $\id$, $\id$, $s$ i $t$. És a dir, $\cat{P}=\Free(\bullet\rightrightarrows\bullet)$. Així, la \enquote{forma} amb què definim els grafs ---el graf subjacent de $\cat{P}$--- és, ella mateixa, un graf, de fet un objecte de $\cat{Graph}=\cat{Set}^{\cat{P}}$: el patró es tanca sobre si mateix.
\end{observacio}

Ara que les transformacions naturals formen una categoria $\cat{D}^{\cat{C}}$, en podem estudiar els isomorfismes.

\begin{definicio}\index{isomorfisme natural}
Una transformació natural $\eta\colon F\Rightarrow G$ és un \textbf{isomorfisme natural} si admet una \emph{inversa}: una transformació natural $\eta^{-1}\colon G\Rightarrow F$ amb $\eta^{-1}\comp\eta=\id_F$ i $\eta\comp\eta^{-1}=\id_G$ (composició vertical). En aquest cas escrivim $F\cong G$ i diem que $F$ i $G$ són \textbf{naturalment isomorfs}.
\end{definicio}

Quan $\cat{C}$ és petita, aquesta condició és, literalment, la d'un isomorfisme a la categoria de functors $\cat{D}^{\cat{C}}$. En efecte, els morfismes de $\cat{D}^{\cat{C}}$ són les transformacions naturals i la seva composició és la vertical; per tant, que $\eta$ admeti una inversa $\eta^{-1}$ com a transformació natural ---amb $\eta^{-1}\comp\eta=\id_F$ i $\eta\comp\eta^{-1}=\id_G$--- no és res més que tenir un invers a $\cat{D}^{\cat{C}}$. Els isomorfismes naturals coincideixen, doncs, amb els isomorfismes de $\cat{D}^{\cat{C}}$: la noció d'isomorfisme del primer capítol (definició~\ref{def:isomorfisme}), un nivell més amunt.

El fet clau és que aquesta condició \enquote{global} ---tenir inversa com a transformació natural--- es pot comprovar \enquote{punt a punt}, sobre les components:

\begin{proposicio}[Caracterització]\label{prop:iso-natural-components}
Siguin $F,G\colon\cat{C}\to\cat{D}$ dos functors. Una transformació natural $\eta\colon F\Rightarrow G$ és un isomorfisme natural si, i només si, cada component $\eta_A$ (per a $A$ objecte de $\cat{C}$) és un isomorfisme a $\cat{D}$.
\end{proposicio}

\begin{prova}
Si $\eta$ té inversa $\eta^{-1}$, avaluant les igualtats $\eta^{-1}\comp\eta=\id_F$ i $\eta\comp\eta^{-1}=\id_G$ a cada objecte $A$ ---i recordant que la composició vertical és component a component--- obtenim $(\eta^{-1})_A\comp\eta_A=\id_{F(A)}$ i $\eta_A\comp(\eta^{-1})_A=\id_{G(A)}$; per tant $\eta_A$ és un isomorfisme a $\cat{D}$, amb invers $(\eta^{-1})_A$.

Recíprocament, suposem que cada $\eta_A$ és un isomorfisme a $\cat{D}$, amb invers $\eta_A^{-1}$. N'hi ha prou de veure que aquests inversos són les components d'una transformació natural $\eta^{-1}\colon G\Rightarrow F$, perquè aleshores, component a component, és la inversa de $\eta$. Partim del quadrat de naturalitat de $\eta$ per a un morfisme $f\colon A\to B$,
\[
\begin{tikzpicture}[baseline=(m.center)]
  \matrix (m) [matrix of math nodes, row sep=2.6em, column sep=3.4em] {
    F(A) & F(B) \\
    G(A) & G(B) \\
  };
  \draw[-{Stealth[scale=1.1]}] (m-1-1) -- node[above] {$F(f)$} (m-1-2);
  \draw[-{Stealth[scale=1.1]}] (m-1-1) -- node[left] {$\eta_A$} (m-2-1);
  \draw[-{Stealth[scale=1.1]}] (m-1-2) -- node[right] {$\eta_B$} (m-2-2);
  \draw[-{Stealth[scale=1.1]}] (m-2-1) -- node[below] {$G(f)$} (m-2-2);
\end{tikzpicture}
\]
és a dir, $G(f)\comp\eta_A=\eta_B\comp F(f)$, i component per l'esquerra amb $\eta_B^{-1}$ i per la dreta amb $\eta_A^{-1}$,
\[
\eta_B^{-1}\comp G(f)=\eta_B^{-1}\comp G(f)\comp\eta_A\comp\eta_A^{-1}=\eta_B^{-1}\comp\eta_B\comp F(f)\comp\eta_A^{-1}=F(f)\comp\eta_A^{-1},
\]
que és la condició de naturalitat per a $\eta^{-1}$.
\end{prova}

\begin{observacio}\label{obs:iso-no-natural}
La distinció entre isomorfisme i isomorfisme \emph{natural} és fonamental. Que $F(A)$ i $G(A)$ siguin isomorfs per a cada $A$ no basta: cal que els isomorfismes ---un per objecte, les components $\eta_A$--- es puguin triar \emph{alhora} de manera compatible amb tots els morfismes. Aquesta compatibilitat no es refereix a triar més components, sinó a la \emph{naturalitat}: per a cada morfisme $f\colon A\to B$ de $\cat{C}$, el quadrat de $\eta$ ha de commutar ($G(f)\comp\eta_A=\eta_B\comp F(f)$), lligant les components dels dos extrems.

Un exemple mínim ho mostra sense cap tecnicisme. Prenem un grup de dos elements, $\{e,g\}$ amb $g^2=e$ (on $e$ és la identitat), vist com a categoria d'un sol objecte $\cat{B}(\{e,g\})$, i dos functors $F,G\colon\cat{B}(\{e,g\})\to\cat{Set}$. Tots dos envien l'únic objecte $\ast$ al mateix conjunt $\{0,1\}$, però l'element $g$ actua de manera diferent: en $F$ actua per la identitat, i en $G$ per la transposició $\tau$ que intercanvia $0$ i $1$. Puntualment no hi ha res a distingir: $F(\ast)=G(\ast)=\{0,1\}$ són el mateix conjunt, i per tant \emph{isomorfs}. Tanmateix, no hi ha cap isomorfisme natural $F\cong G$. Una transformació natural $\eta\colon F\Rightarrow G$ tindria una única component $\eta_\ast\colon\{0,1\}\to\{0,1\}$, i la naturalitat respecte del morfisme $g$ exigiria
\[
G(g)\comp\eta_\ast=\eta_\ast\comp F(g),
\qquad\text{és a dir}\qquad
\tau\comp\eta_\ast=\eta_\ast.
\]
Però $\tau$ no té punts fixos, de manera que cap aplicació $\eta_\ast$ pot complir $\tau\comp\eta_\ast=\eta_\ast$ (aplicant a qualsevol $x$ obtindríem $\tau(\eta_\ast(x))=\eta_\ast(x)$, impossible). No hi ha, doncs, cap transformació natural $F\Rightarrow G$, i menys encara un isomorfisme natural: la igualtat puntual dels objectes no es pot fer compatible amb l'acció del morfisme. Val a dir que aquí obtenim un contraexemple més fort que el necessari: per descartar un isomorfisme \emph{natural} n'hi hauria prou de veure que cap component \emph{bijectiva} compleix la naturalitat, mentre que l'argument anterior descarta \emph{totes} les components, bijectives o no. La manca d'isomorfisme natural en surt, doncs, reforçada.

El cas clàssic del dual d'un espai vectorial il·lustra la mateixa idea, però demana una precaució addicional. Sovint es diu que un espai de dimensió finita \enquote{és isomorf al seu dual $V\cong V^*$, però no naturalment}. Cal formular-ho amb cura, perquè la identitat és un functor \emph{covariant} $\cat{Vect}_{\mathbb K}\to\cat{Vect}_{\mathbb K}$ mentre que el dual $(-)^*$ és \emph{contravariant}, $\cat{Vect}_{\mathbb K}\op\to\cat{Vect}_{\mathbb K}$: no són dos functors paral·lels als quals es pugui aplicar directament la definició de transformació natural d'aquest capítol. La manera correcta de precisar-ho és restringir-se als \emph{isomorfismes} lineals i comparar la identitat amb l'acció covariant $T\mapsto(T^{-1})^*$; aleshores es comprova que cap família d'isomorfismes $V\cong V^*$ és compatible amb tots els isomorfismes lineals. Aquest desenvolupament es fa amb detall a l'exercici~\ref{ex:dual-no-natural} de la part de pràctica. En canvi, l'isomorfisme $V\cong V^{**}$ amb el bidual sí que està ben plantejat sobre \emph{totes} les aplicacions lineals ---identitat i bidual són tots dos covariants--- i és natural.
\end{observacio}

\begin{exemple}[Les identificacions canòniques del capítol anterior]\label{ex:isos-naturals-cap2}\index{isomorfisme natural!functor de parts}
Al capítol~\ref{cap:functors} vam trobar diverses identificacions canòniques entre un functor i un homfunctor, i vam anunciar que les reconeixeríem com a isomorfismes naturals. La proposició~\ref{prop:iso-natural-components} ho fa immediat: en cada cas, les components són bijeccions (o isomorfismes lineals), i només cal comprovar la naturalitat.

\emph{El functor de parts.} Els dos functors contravariants de l'exemple~\ref{ex:functor-parts} són
\[
\mathcal P,\ \Hom_{\cat{Set}}(-,\{0,1\})\colon\cat{Set}\op\longrightarrow\cat{Set}.
\] Les funcions característiques donen, per a cada conjunt $X$, la bijecció
\[
\chi_X\colon\mathcal P(X)\to\Hom_{\cat{Set}}(X,\{0,1\}),\qquad B\mapsto\chi_B.
\]
Un morfisme de $\cat{Set}\op$ de $Y$ a $X$ és una aplicació $f\colon X\to Y$, que $\mathcal P$ envia a l'antiimatge $f^{-1}$ i l'homfunctor a la precomposició $-\comp f$. El quadrat de naturalitat és
\[
\begin{tikzpicture}[baseline=(m.center)]
  \matrix (m) [matrix of math nodes, row sep=2.8em, column sep=4.2em] {
    \mathcal P(Y) & \mathcal P(X) \\
    \Hom(Y,\{0,1\}) & \Hom(X,\{0,1\}) \\
  };
  \draw[-{Stealth[scale=1.1]}] (m-1-1) -- node[above] {$f^{-1}$} (m-1-2);
  \draw[-{Stealth[scale=1.1]}] (m-1-1) -- node[left] {$\chi_Y$} (m-2-1);
  \draw[-{Stealth[scale=1.1]}] (m-1-2) -- node[right] {$\chi_X$} (m-2-2);
  \draw[-{Stealth[scale=1.1]}] (m-2-1) -- node[below] {$-\comp f$} (m-2-2);
\end{tikzpicture}
\qquad
\chi_B\comp f=\chi_{f^{-1}(B)}\quad(B\subseteq Y),
\]
i commuta, perquè per a tot $x\in X$ es té $(\chi_B\comp f)(x)=1$ si i només si $f(x)\in B$, és a dir, si i només si $x\in f^{-1}(B)$. Per tant $\chi\colon\mathcal P\Rightarrow\Hom_{\cat{Set}}(-,\{0,1\})$ és un isomorfisme natural, i ara podem escriure amb tota precisió
\[
\mathcal P\;\cong\;\Hom_{\cat{Set}}(-,\{0,1\}).
\]

\emph{Els altres casos.} El mateix esquema val per a les altres identificacions del capítol anterior; n'indiquem les components i l'equació de naturalitat, que es comprova directament.
\begin{itemize}
  \item A $\cat{Set}$, $\id_{\cat{Set}}\cong\Hom_{\cat{Set}}(\{\ast\},-)$: la component envia $x\in X$ a l'aplicació $\bar x\colon\ast\mapsto x$, i la naturalitat diu que $f\comp\bar x=\overline{f(x)}$.
  \item A $\cat{Graph}$, el functor de vèrtexs és naturalment isomorf a $\Hom_{\cat{Graph}}(\mathbf V,-)$ i el d'arestes a $\Hom_{\cat{Graph}}(\mathbf E,-)$ (exemple~\ref{ex:homfunctor-grafs}): la component envia cada vèrtex, o cada aresta, al morfisme que el tria, i la naturalitat diu que compondre aquest morfisme amb un morfisme de grafs $F$ tria la imatge per $F$.
  \item A $\cat{Vect}_{\mathbb K}$, el functor identitat és naturalment isomorf a l'aixecament
\[
\Hom_{\mathbb K}(\mathbb K,-)\colon\cat{Vect}_{\mathbb K}\longrightarrow\cat{Vect}_{\mathbb K}.
\]
La component envia $w\in W$ a l'aplicació lineal $\lambda\mapsto\lambda w$, amb inversa $\varphi\mapsto\varphi(1)$, i la naturalitat diu que $T\comp(\lambda\mapsto\lambda w)=(\lambda\mapsto\lambda\,T(w))$ per a tota aplicació lineal $T$.
\end{itemize}
Aquests exemples tornaran a aparèixer al capítol del lema de Yoneda, on direm que aquests functors són \emph{representables}.
\end{exemple}

\subsection{Equivalència de categories}\label{sec:equivalencia}

Quan volem dir que dues categories són \enquote{essencialment la mateixa} ---amb la mateixa estructura categòrica, distingint els objectes només llevat d'isomorfisme---, la igualtat estricta és massa rígida, i l'isomorfisme de categories, que demana functors inversos ($G\comp F=\id_{\cat{C}}$ i $F\comp G=\id_{\cat{D}}$), també ho és sovint. La noció adequada relaxa aquestes igualtats a isomorfismes naturals: en lloc d'un invers demanarem un \textbf{quasiinvers}, un functor $G$ amb $G\comp F\cong\id_{\cat{C}}$ i $F\comp G\cong\id_{\cat{D}}$.

\begin{definicio}\index{equivalència de categories}
\label{def:equivalencia}
Una \textbf{equivalència de categories} entre $\cat{C}$ i $\cat{D}$ és una parella de functors
\[
F\colon\cat{C}\to\cat{D},
\qquad
G\colon\cat{D}\to\cat{C},
\]
juntament amb isomorfismes naturals
\[
G\comp F\cong\id_{\cat{C}}
\qquad\text{i}\qquad
F\comp G\cong\id_{\cat{D}}.
\]
Quan existeix una tal equivalència, diem que $\cat{C}$ i $\cat{D}$ són \textbf{equivalents} i escrivim $\cat{C}\simeq\cat{D}$. El functor $G$ s'anomena un \textbf{quasiinvers} de $F$. Quan calgui donar nom als isomorfismes naturals, els orientarem sempre així:
\[
\eta\colon\id_{\cat{C}}\Rightarrow G\comp F,
\qquad
\varepsilon\colon F\comp G\Rightarrow\id_{\cat{D}}.
\]
\end{definicio}

\begin{observacio}
La diferència amb l'isomorfisme de categories és que aquí no demanem $G\comp F=\id_{\cat{C}}$, sinó només $G\comp F\cong\id_{\cat{C}}$: en tornar de $\cat{C}$ a $\cat{C}$ no recuperem exactament cada objecte, sinó un objecte \emph{naturalment isomorf}. Aquesta flexibilitat és la que fa que l'equivalència sigui la noció adequada de comparació entre categories quan, com passa sovint, només ens interessen els objectes llevat d'isomorfisme.
\end{observacio}

El resultat central caracteritza les equivalències mitjançant les tres propietats dels functors que vam introduir al capítol anterior: ser ple i fidel (la definició~\ref{def:ple-fidel}) i l'exhaustivitat essencial (la definició~\ref{def:ess-exhaustiu}).

\begin{teorema}\index{equivalència de categories!caracterització}
\label{teo:caract-equivalencia}
Un functor $F\colon\cat{C}\to\cat{D}$ forma part d'una equivalència de categories si i només si és \textbf{ple}, \textbf{fidel} i \textbf{essencialment exhaustiu}.
\end{teorema}

\begin{prova}
\emph{($\Rightarrow$)} Suposem que $F$ forma part d'una equivalència, amb quasiinvers $G\colon\cat{D}\to\cat{C}$ i isomorfismes naturals $\eta\colon\id_{\cat{C}}\Rightarrow G\comp F$ i $\varepsilon\colon F\comp G\Rightarrow\id_{\cat{D}}$. Com que les components de $\eta$ i de $\varepsilon$ són invertibles, els quadrats de naturalitat es poden llegir com a fórmules: per a $f\colon A\to B$ a $\cat{C}$ i $g\colon D\to D'$ a $\cat{D}$,
\[
G(F(f))=\eta_B\comp f\comp\eta_A^{-1},
\qquad\qquad
F(G(g))=\varepsilon_{D'}^{-1}\comp g\comp\varepsilon_D.
\tag{$\ast$}
\]

\emph{$F$ és fidel.} Siguin $f,f'\colon A\to B$ amb $F(f)=F(f')$. Aleshores $G(F(f))=G(F(f'))$ i, per~$(\ast)$, $\eta_B\comp f\comp\eta_A^{-1}=\eta_B\comp f'\comp\eta_A^{-1}$. Component amb $\eta_B^{-1}$ per l'esquerra i amb $\eta_A$ per la dreta, $f=f'$.

\emph{$G$ és fidel.} És el mateix argument amb la segona fórmula de~$(\ast)$: si $G(g)=G(g')$ per a $g,g'\colon D\to D'$, aleshores $\varepsilon_{D'}^{-1}\comp g\comp\varepsilon_D=\varepsilon_{D'}^{-1}\comp g'\comp\varepsilon_D$, i per tant $g=g'$.

\emph{$F$ és ple.} Sigui $u\colon F(A)\to F(B)$ un morfisme qualsevol de $\cat{D}$, i posem
\[
f=\eta_B^{-1}\comp G(u)\comp\eta_A\colon A\to B.
\]
Per~$(\ast)$, $G(F(f))=\eta_B\comp\bigl(\eta_B^{-1}\comp G(u)\comp\eta_A\bigr)\comp\eta_A^{-1}=G(u)$, i com que $G$ és fidel, $F(f)=u$. Per concloure $F(f)=u$ no es pot fer servir que $F$ és fidel, perquè encara no sabem que $u$ sigui la imatge de cap morfisme; és que $G$ sigui fidel el que ho permet.

\emph{$F$ és essencialment exhaustiu.} Per a cada objecte $D$ de $\cat{D}$, la component $\varepsilon_D\colon F(G(D))\to D$ és un isomorfisme, de manera que $D$ és isomorf a l'objecte $F(G(D))$ de la imatge de $F$.

\emph{($\Leftarrow$)} Suposem que $F$ és ple, fidel i essencialment exhaustiu. Construïm un quasiinvers en quatre passos.

\emph{Pas 1: $G$ sobre objectes.} Per l'exhaustivitat essencial, per a cada objecte $D$ de $\cat{D}$ hi ha algun objecte $C$ de $\cat{C}$ amb $F(C)\cong D$. Triem, per a cada $D$, un tal objecte, que anomenem $G(D)$, i un isomorfisme $\varepsilon_D\colon F(G(D))\to D$.

\emph{Pas 2: $G$ sobre morfismes.} Sigui $g\colon D\to D'$. El morfisme $\varepsilon_{D'}^{-1}\comp g\comp\varepsilon_D$ va de $F(G(D))$ a $F(G(D'))$. Com que $F$ és ple, prové d'algun morfisme $G(D)\to G(D')$, i com que $F$ és fidel, d'un de sol. Definim $G(g)$ com aquest únic morfisme; és a dir, $G(g)$ queda caracteritzat per
\[
F(G(g))=\varepsilon_{D'}^{-1}\comp g\comp\varepsilon_D.
\tag{$\ast\ast$}
\]

\emph{Pas 3: $G$ és un functor.} Per a la identitat, $(\ast\ast)$ dona $F(G(\id_D))=\varepsilon_D^{-1}\comp\varepsilon_D=\id_{F(G(D))}=F(\id_{G(D)})$, i, com que $F$ és fidel, $G(\id_D)=\id_{G(D)}$. Per a la composició, siguin $g\colon D\to D'$ i $g'\colon D'\to D''$. Aplicant $(\ast\ast)$ dues vegades,
\[
\begin{aligned}
F\bigl(G(g')\comp G(g)\bigr)&=F(G(g'))\comp F(G(g))
=\bigl(\varepsilon_{D''}^{-1}\comp g'\comp\varepsilon_{D'}\bigr)\comp\bigl(\varepsilon_{D'}^{-1}\comp g\comp\varepsilon_D\bigr)\\
&=\varepsilon_{D''}^{-1}\comp(g'\comp g)\comp\varepsilon_D=F\bigl(G(g'\comp g)\bigr),
\end{aligned}
\]
i, com que $F$ és fidel, $G(g')\comp G(g)=G(g'\comp g)$.

\emph{Pas 4: els isomorfismes naturals.} Component $(\ast\ast)$ amb $\varepsilon_{D'}$ per l'esquerra s'obté
\[
\varepsilon_{D'}\comp F(G(g))=g\comp\varepsilon_D,
\]
que és el quadrat de naturalitat de $\varepsilon\colon F\comp G\Rightarrow\id_{\cat{D}}$ per a $g$. Com que cada $\varepsilon_D$ és un isomorfisme, $\varepsilon$ és un isomorfisme natural.

Per construir $\eta$, sigui $A$ un objecte de $\cat{C}$. Com que $F$ és ple i fidel, hi ha un únic $\eta_A\colon A\to G(F(A))$ amb $F(\eta_A)=\varepsilon_{F(A)}^{-1}$, i un únic $\zeta_A\colon G(F(A))\to A$ amb $F(\zeta_A)=\varepsilon_{F(A)}$. Aleshores
\[
F(\zeta_A\comp\eta_A)=\varepsilon_{F(A)}\comp\varepsilon_{F(A)}^{-1}=F(\id_A),
\qquad
F(\eta_A\comp\zeta_A)=\varepsilon_{F(A)}^{-1}\comp\varepsilon_{F(A)}=F(\id_{G(F(A))}),
\]
i, com que $F$ és fidel, $\zeta_A$ és l'invers de $\eta_A$: cada $\eta_A$ és un isomorfisme. Queda la naturalitat: per a $f\colon A\to B$, cal veure $\eta_B\comp f=G(F(f))\comp\eta_A$. Apliquem $F$ als dos membres. D'una banda, $F(\eta_B\comp f)=\varepsilon_{F(B)}^{-1}\comp F(f)$. De l'altra, aplicant $(\ast\ast)$ al morfisme $g=F(f)$,
\[
F\bigl(G(F(f))\comp\eta_A\bigr)=\bigl(\varepsilon_{F(B)}^{-1}\comp F(f)\comp\varepsilon_{F(A)}\bigr)\comp\varepsilon_{F(A)}^{-1}=\varepsilon_{F(B)}^{-1}\comp F(f).
\]
Les dues imatges coincideixen, i, com que $F$ és fidel, $\eta_B\comp f=G(F(f))\comp\eta_A$. Així $\eta\colon\id_{\cat{C}}\Rightarrow G\comp F$ és un isomorfisme natural, i $G$ és un quasiinvers de $F$.
\end{prova}

\begin{observacio}[Sobre l'elecció]
A la direcció ($\Leftarrow$), el pas~1 tria simultàniament un objecte $G(D)$ i un isomorfisme $\varepsilon_D$ per a cada $D$, i per tant fa servir l'\emph{axioma de l'elecció}.\footnotemark{} Tota la resta de la construcció és \emph{determinada} per aquesta tria: el pas~2 no tria res, perquè, com que $F$ és ple i fidel, $G(g)$ és únic, i $\eta$ també queda fixat. En casos concrets sovint es pot donar un quasiinvers canònic sense recórrer a l'elecció. La direcció ($\Rightarrow$), en canvi, no tria res, i tampoc no fa servir cap relació entre $\eta$ i $\varepsilon$ més enllà de la naturalitat i la invertibilitat.
\end{observacio}
\footnotetext{Si $\cat{D}$ és gran, la família de tries està indexada per una classe pròpia, i cal una forma d'elecció més forta que l'habitual, l'\emph{elecció global}, disponible per exemple a la teoria de classes NBG (apèndix~\ref{cap:mida}). Vegeu \cite[§1.5]{riehl}.}%

\begin{definicio}[Esquelet]\index{categoria!esquelètica}\index{esquelet}\label{def:esquelet}
Una categoria és \textbf{esquelètica} si dos objectes isomorfs qualssevol són \emph{iguals}; és a dir, si cada classe d'isomorfisme té un únic objecte. Un \textbf{esquelet} d'una categoria $\cat{C}$ és una subcategoria plena que conté \emph{exactament} un objecte de cada classe d'isomorfisme de $\cat{C}$; un esquelet és sempre una categoria esquelètica.
\end{definicio}

\begin{observacio}
Convé distingir dues nocions properes. Una subcategoria plena que conté \emph{almenys} un objecte de cada classe d'isomorfisme ja és equivalent a $\cat{C}$ (com veurem tot seguit), però pot tenir diversos representants isomorfs d'una mateixa classe. Un \emph{esquelet} n'és el cas ajustat: en conté \emph{exactament} un, de manera que no queden objectes isomorfs distints. Triar exactament un representant de cada classe d'isomorfisme és, però, una elecció simultània. Amb l'\emph{axioma de l'elecció}, tota categoria \emph{petita} admet un esquelet, i tots els seus esquelets són equivalents entre si. Per a una categoria gran, seleccionar simultàniament un representant de cada classe d'isomorfisme pot requerir un principi d'elecció global, segons el marc fundacional adoptat.\footnotemark És la mateixa elecció que fa servir la construcció d'un quasiinvers (observació anterior).
\end{observacio}
\footnotetext{En una categoria gran, les classes d'isomorfisme formen una classe pròpia, i cal la mateixa forma d'elecció global que a l'observació anterior (apèndix~\ref{cap:mida}); la petitesa local només controla els homconjunts i no resol aquesta selecció.}%

\begin{exemple}
Tota categoria $\cat{C}$ és equivalent a qualsevol subcategoria plena que contingui almenys un objecte de cada classe d'isomorfisme de $\cat{C}$. En efecte, el functor inclusió d'aquesta subcategoria és ple i fidel per ser una subcategoria plena, i essencialment exhaustiu per contenir un representant de cada classe. Per exemple, la categoria $\cat{FinSet}$ de tots els conjunts finits és equivalent a la seva subcategoria plena formada pels conjunts $\underline{0}=\varnothing$ i $\underline{n}=\{1,\dots,n\}$ per a $n\geq 1$ (exemple~\ref{ex:ess-exhaustiu-finset}): no cal treballar amb tots els conjunts finits, sinó només amb un de cada mida. Aquesta subcategoria $\{\underline{n}\mid n\geq 0\}$ és, de fet, un \emph{esquelet} de $\cat{FinSet}$, perquè conté exactament un conjunt de cada cardinalitat finita. Convé notar que $\cat{FinSet}$ és una categoria \emph{gran} ---els seus objectes són massa nombrosos per formar un conjunt, tal com passa amb $\cat{Set}$--- mentre que l'esquelet $\{\underline{n}\mid n\geq 0\}$ és \emph{petit}, ja que té un conjunt (de fet, numerable) d'objectes. Una categoria s'anomena \textbf{essencialment petita}\index{categoria!essencialment petita} si és equivalent a una categoria petita. L'equivalència $\cat{FinSet}\simeq\{\underline{n}\mid n\geq 0\}$ mostra, doncs, que $\cat{FinSet}$ és essencialment petita.
\end{exemple}

\begin{observacio}
Aquest exemple mostra la utilitat pràctica de l'equivalència: permet substituir una categoria gran i redundant per una de més petita i manejable, sense perdre cap informació categòrica. És el reflex, a nivell de categories, del fet que en teoria de categories els objectes només importen llevat d'isomorfisme.
\end{observacio}

\begin{exemple}[Matrius i espais vectorials: equivalència sense igualtat ni isomorfisme]\label{ex:mat-finvect}\index{equivalència de categories!matrius i espais vectorials}\index{Mat@$\cat{Mat}_{\mathbb K}$}
Fixem un cos $\mathbb K$. La categoria $\cat{Mat}_{\mathbb K}$ té per objectes els nombres naturals $n\geq 0$. Els morfismes $m\to n$ són les matrius $n\times m$ amb entrades a $\mathbb K$, la composició és el producte, $B\comp A=BA$, i la identitat de $n$ és la matriu identitat $I_n$; els axiomes són l'associativitat del producte de matrius i la neutralitat de $I_n$. Sigui $\cat{FinVect}_{\mathbb K}$ la subcategoria plena de $\cat{Vect}_{\mathbb K}$ formada pels espais de dimensió finita. El functor
\[
F\colon\cat{Mat}_{\mathbb K}\to\cat{FinVect}_{\mathbb K},\qquad n\mapsto\mathbb K^n,\qquad A\mapsto(x\mapsto Ax),
\]
té els tipus correctes: una matriu $n\times m$ dona una aplicació lineal $\mathbb K^m\to\mathbb K^n$. A més, $F(BA)=F(B)\comp F(A)$ i $F(I_n)=\id_{\mathbb K^n}$. És fidel i ple, perquè tota aplicació lineal $\mathbb K^m\to\mathbb K^n$ és la multiplicació per una única matriu, la que té per columnes les imatges de la base canònica. És essencialment exhaustiu, perquè tot espai de dimensió $n$ és isomorf a $\mathbb K^n$. Pel teorema~\ref{teo:caract-equivalencia}, $F$ és una equivalència: $\cat{Mat}_{\mathbb K}\simeq\cat{FinVect}_{\mathbb K}$. Aquest exemple permet veure els tres nivells de comparació.
\begin{itemize}
\item \emph{No és una igualtat.} Són dues categories diferents: els objectes de l'una són nombres i els de l'altra, espais vectorials.
\item \emph{No és un isomorfisme de categories.} A $\cat{Mat}_{\mathbb K}$, dos objectes isomorfs són iguals, perquè una matriu invertible és quadrada: la categoria és esquelètica. A $\cat{FinVect}_{\mathbb K}$, en canvi, $\mathbb K^2$ i l'espai dels polinomis de grau $\leq 1$ són objectes diferents i isomorfs. Un isomorfisme de categories és bijectiu sobre objectes i preserva i reflecteix els isomorfismes, de manera que conservaria la propietat de ser esquelètica. Per tant, no hi ha cap isomorfisme de categories entre $\cat{Mat}_{\mathbb K}$ i $\cat{FinVect}_{\mathbb K}$, ni aquest $F$ ni cap altre.
\item \emph{El quasiinvers només torna llevat d'isomorfisme.} El pas~1 de la demostració del teorema consisteix aquí a \emph{triar una base} de cada espai $V$, és a dir, un isomorfisme $\varepsilon_V\colon\mathbb K^{\dim V}\to V$. Aleshores $G(V)=\dim V$, i $G(g)$ és la matriu de $g$ en les bases triades. Si per als espais $\mathbb K^n$ triem la base canònica, es té literalment $G\comp F=\id_{\cat{Mat}_{\mathbb K}}$. En canvi, $F(G(V))=\mathbb K^{\dim V}$, que no és $V$ tret que $V$ sigui algun $\mathbb K^n$. Així, $F\comp G$ \emph{no} és igual a $\id_{\cat{FinVect}_{\mathbb K}}$, sinó només isomorf a través de $\varepsilon$. La condició de naturalitat $g\comp\varepsilon_V=\varepsilon_W\comp F(G(g))$ és la coneguda fórmula que relaciona una aplicació lineal amb la seva matriu en unes bases.
\end{itemize}
Un espai vectorial qualsevol no té cap base privilegiada; per això el quasiinvers depèn d'una tria, i el que s'obté és una equivalència i no un isomorfisme. Per a l'àlgebra lineal, la conclusió és la que ja es fa servir a la pràctica: treballar amb matrius no perd cap informació sobre les aplicacions lineals entre espais de dimensió finita, sempre que es recordi que la identificació depèn de les bases.
\end{exemple}

Quines propietats es conserven per equivalència i quines no, i com es demostra que dues categories \emph{no} són equivalents, es tracta a l'apèndix~\refalt{cap:contraexemples}{de contraexemples del llibre complet}, junt amb la relació entre equivalència i esquelets.


\section{Composició horitzontal}\index{transformació natural!composició horitzontal}

La composició vertical apila transformacions entre functors \emph{amb el mateix domini i codomini}. Hi ha una segona manera de compondre, que combina transformacions naturals amb functors i, més generalment, transformacions naturals \enquote{de costat}.

Comencem pel cas mixt, en què componem una transformació natural amb un functor\footnote{Segons el costat pel qual s'afegeix el functor, es parla de composició \enquote{per l'esquerra} i \enquote{per la dreta}; convé advertir que aquesta assignació d'\enquote{esquerra} i \enquote{dreta} no és uniforme a la bibliografia. Aquí la fixem per la posició del functor en la notació: $K\eta$ (functor a l'esquerra) és la composició per l'esquerra, i $\eta H$ (functor a la dreta), la composició per la dreta.}.

\begin{definicio}[Composició amb un functor]\index{transformació natural!composició amb un functor}
Sigui $\eta\colon F\Rightarrow G$ una transformació natural entre functors $F,G\colon\cat{C}\to\cat{D}$.
\begin{itemize}
\item Si $K\colon\cat{D}\to\cat{E}$ és un functor, la \textbf{composició per l'esquerra} $K\eta\colon K\comp F\Rightarrow K\comp G$ té per components $(K\eta)_A=K(\eta_A)$.
\item Si $H\colon\cat{B}\to\cat{C}$ és un functor, la \textbf{composició per la dreta} $\eta H\colon F\comp H\Rightarrow G\comp H$ té per components $(\eta H)_B=\eta_{H(B)}$.
\end{itemize}
\end{definicio}

\begin{proposicio}
Les dues construccions anteriors són transformacions naturals.
\end{proposicio}

\begin{prova}
Per a $K\eta$, apliquem el functor $K$ al quadrat de naturalitat de $\eta$ per a un morfisme $f\colon A\to B$: de $G(f)\comp\eta_A=\eta_B\comp F(f)$, i com que $K$ preserva la composició,
\[
K(G(f))\comp K(\eta_A)=K(\eta_B)\comp K(F(f)),
\]
que és la condició de naturalitat per a $K\eta$ respecte de $K\comp F$ i $K\comp G$. Per a $\eta H$, la naturalitat respecte d'un morfisme $b\colon B\to B'$ de $\cat{B}$ és el quadrat de naturalitat de $\eta$ per al morfisme $H(b)$ de $\cat{C}$.
\end{prova}

Amb aquestes dues operacions podem definir la composició horitzontal general.

\begin{definicio}[Composició horitzontal]\index{transformació natural!composició horitzontal}
Siguin $\eta\colon F\Rightarrow G$ (amb $F,G\colon\cat{C}\to\cat{D}$) i $\theta\colon K\Rightarrow L$ (amb $K,L\colon\cat{D}\to\cat{E}$). La seva \textbf{composició horitzontal} $\theta\ast\eta\colon K\comp F\Rightarrow L\comp G$ té per components
\[
(\theta\ast\eta)_A=\theta_{G(A)}\comp K(\eta_A)=L(\eta_A)\comp\theta_{F(A)}.
\]
\end{definicio}

Diagramàticament, la composició horitzontal juxtaposa les dues transformacions, l'una a continuació de l'altra:
\[
\begin{tikzpicture}[baseline=(A.base)]
  \node (A) at (0,0) {$\cat{C}$};
  \node (B) at (2.9,0) {$\cat{D}$};
  \node (Cc) at (5.8,0) {$\cat{E}$};
  \draw[cat] (A) to[bend left=26] node[above] {$F$} (B);
  \draw[cat] (A) to[bend right=26] node[below] {$G$} (B);
  \draw[cat nat] (1.45,0.2) -- node[right=1.5pt] {$\eta$} (1.45,-0.2);
  \draw[cat] (B) to[bend left=26] node[above] {$K$} (Cc);
  \draw[cat] (B) to[bend right=26] node[below] {$L$} (Cc);
  \draw[cat nat] (4.35,0.2) -- node[right=1.5pt] {$\theta$} (4.35,-0.2);
\end{tikzpicture}
\]

\begin{observacio}\label{obs:comp-horitzontal}
Les dues expressions de $(\theta\ast\eta)_A$ coincideixen precisament pel quadrat de naturalitat de $\theta$ aplicat al morfisme $\eta_A\colon F(A)\to G(A)$ de $\cat{D}$:
\[
\begin{tikzpicture}[baseline=(m.center)]
  \matrix (m) [matrix of math nodes, row sep=2.8em, column sep=4.2em] {
    K(F(A)) & K(G(A)) \\
    L(F(A)) & L(G(A)) \\
  };
  \draw[-{Stealth[scale=1.1]}] (m-1-1) -- node[above] {$K(\eta_A)$} (m-1-2);
  \draw[-{Stealth[scale=1.1]}] (m-1-1) -- node[left] {$\theta_{F(A)}$} (m-2-1);
  \draw[-{Stealth[scale=1.1]}] (m-1-2) -- node[right] {$\theta_{G(A)}$} (m-2-2);
  \draw[-{Stealth[scale=1.1]}] (m-2-1) -- node[below] {$L(\eta_A)$} (m-2-2);
\end{tikzpicture}
\qquad
\theta_{G(A)}\comp K(\eta_A)=L(\eta_A)\comp\theta_{F(A)}.
\]
La composició horitzontal es pot llegir, doncs, component una transformació natural amb un functor a cada costat, en qualsevol dels dos ordres:
\[
\theta\ast\eta=(\theta G)\comp(K\eta)=(L\eta)\comp(\theta F).
\]
La composició vertical i l'horitzontal, juntes, satisfan una llei d'intercanvi que fa de $\cat{Cat}$ una \emph{2-categoria}: en una graella de transformacions naturals com la següent, compondre primer en vertical i després en horitzontal dona el mateix que fer-ho en l'ordre contrari,
\[
(\beta'\comp\alpha')\ast(\beta\comp\alpha)=(\beta'\ast\beta)\comp(\alpha'\ast\alpha).
\]
\[
\begin{tikzpicture}[baseline=(A.base)]
  \node (A) at (0,0) {$\cat{C}$};
  \node (B) at (3.4,0) {$\cat{D}$};
  \node (Cc) at (6.8,0) {$\cat{E}$};
  \draw[cat] (A) to[bend left=44] node[above] {$F$} (B);
  \draw[cat] (A) -- (B);
  \draw[cat] (A) to[bend right=44] node[below] {$H$} (B);
  \node[fill=white,inner sep=1.5pt] at (1.7,0) {$G$};
  \draw[cat nat] (1.7,0.52) -- node[right=1.5pt] {$\alpha$} (1.7,0.27);
  \draw[cat nat] (1.7,-0.27) -- node[right=1.5pt] {$\beta$} (1.7,-0.52);
  \draw[cat] (B) to[bend left=44] node[above] {$F'$} (Cc);
  \draw[cat] (B) -- (Cc);
  \draw[cat] (B) to[bend right=44] node[below] {$H'$} (Cc);
  \node[fill=white,inner sep=1.5pt] at (5.1,0) {$G'$};
  \draw[cat nat] (5.1,0.52) -- node[right=1.5pt] {$\alpha'$} (5.1,0.27);
  \draw[cat nat] (5.1,-0.27) -- node[right=1.5pt] {$\beta'$} (5.1,-0.52);
\end{tikzpicture}
\]
No en necessitarem els detalls, però és el marc en què les transformacions naturals adquireixen tota la seva potència.
\end{observacio}

\begin{resumcap}
\apartat{Idees centrals}
\begin{enumerate}
  \item Una transformació natural és una família de morfismes compatible amb \emph{totes} les fletxes de la categoria de partida: cada quadrat de naturalitat commuta.
  \item Les transformacions naturals es componen verticalment, component a component; amb aquesta composició, els functors $\cat{C}\to\cat{D}$ i les transformacions naturals formen la categoria de functors $\cat{D}^{\cat{C}}$.
  \item Un isomorfisme natural és una transformació natural amb totes les components invertibles. La naturalitat no s'hi pot ometre: una família d'isomorfismes, un per a cada objecte, que no sigui natural no és un isomorfisme natural.
  \item Una equivalència de categories és un functor amb un quasiinvers, llevat d'isomorfismes naturals. En el marc d'elecció adoptat al llibre (apèndix~\ref{cap:mida}), es caracteritza equivalentment per ser un functor ple, fidel i essencialment exhaustiu (teorema~\ref{teo:caract-equivalencia}, amb demostració completa). No és una igualtat ni, en general, un isomorfisme de categories: $\cat{Mat}_{\mathbb K}\simeq\cat{FinVect}_{\mathbb K}$, però no són isomorfes.
  \item La composició horitzontal és un segon mode de composició de transformacions naturals, compatible amb la vertical a través de la llei d'intercanvi.
\end{enumerate}
\apartat{Exemple clau} Els grafs dirigits són els objectes d'una categoria de functors: $\cat{Graph}\cong\cat{Set}^{\cat{P}}$, on els morfismes de grafs són exactament les transformacions naturals. Convé recordar l'orientació: $\cat{P}$ té dues fletxes $s,t\colon E\to V$, de manera que un graf és un functor \emph{covariant} $\cat{P}\to\cat{Set}$; no s'ha de confondre amb un prefeix sobre $\cat{P}$.
\apartat{Errors típics} Oblidar comprovar la naturalitat (que el quadrat commuti per a \emph{tot} morfisme); confondre isomorfisme natural amb igualtat de functors; creure que una equivalència conserva propietats com \enquote{tenir un sol objecte} o \enquote{ser esquelètica}; confondre una equivalència amb un isomorfisme de categories, o esperar que $F\comp G$ sigui literalment la identitat; i barrejar composició vertical i horitzontal de transformacions.
\apartat{Lectura recomanada} \cite[§§7.4--7.10]{awodey} per a la naturalitat, les categories de functors i les equivalències; \cite[§§1.4--1.5]{riehl}; \cite[§1.3]{leinster}, que inclou també les equivalències; \cite[I.4, II.4--II.5 i IV.4]{maclane} per a la formulació clàssica de la naturalitat, de les categories de functors amb la composició horitzontal i de les equivalències; \cite[§§4.2--4.4]{barrwells} per a la naturalitat amb exemples de grafs i llistes i per al càlcul de Godement, que és la composició horitzontal; \cite[§1.4 i §1.5.4]{perrone-notes} per a exemples de naturalitat, una secció sobre què \emph{no} és natural i les equivalències.
\end{resumcap}

\ifpublicacioparcial\else
\chapter[Representabilitat i Yoneda]{\texorpdfstring{Propietats universals,\\ representabilitat i Yoneda}{Propietats universals, representabilitat i Yoneda}}
\label{cap:yoneda}

\begin{objectius}
\apartat{Prerequisits} Categories sobre i sota un objecte i dualitat (cap.~\ref{cap:categories}); homfunctors, bifunctor $\Hom$, prefeixos i la parella $\Free$, $U$ (cap.~\ref{cap:functors}); transformacions i isomorfismes naturals, i categories de functors, en particular la categoria de prefeixos (cap.~\ref{cap:transformacions}). No es pressuposa cap coneixement previ de representabilitat ni del lema de Yoneda.
\apartat{Objectius} En acabar el capítol, el lector hauria de poder:
\begin{itemize}
  \item definir objectes inicials i terminals i demostrar-ne la unicitat corresponent;
  \item construir i llegir categories coma en exemples senzills, i interpretar-hi dades universals com a objectes inicials o terminals;
  \item definir una representació mitjançant un element universal i traduir-la a un isomorfisme natural amb un homfunctor;
  \item distingir el functor representat, l'objecte representant, l'element universal i la representació completa;
  \item definir la immersió de Yoneda, enunciar i demostrar el lema de Yoneda i seguir-ne explícitament la bijecció;
  \item deduir que la immersió de Yoneda és plena i fidel i que un objecte queda determinat, llevat d'isomorfisme, pel seu functor representable.
\end{itemize}
\end{objectius}

\section{Cap a les propietats universals}

Amb les transformacions naturals vam completar, al capítol~\ref{cap:transformacions}, una torre de tres nivells: objectes i morfismes dins d'una categoria, functors entre categories, i transformacions naturals entre functors ---l'estructura que, amb les composicions vertical i horitzontal, fa de $\cat{Cat}$ una 2-categoria. És el marc en què es formulen les construccions més potents de la teoria, i les primeres construccions que estudiarem són les propietats universals.

Als capítols anteriors han aparegut, sempre com a exemples i sense tractar-los en
detall, objectes i construccions definits no per la seva constitució interna sinó
per la manera com es relacionen amb tots els altres. El producte de dos objectes
$A$ i $B$ (subsecció~\ref{subsec:producte-coproducte}) es defineix per una condició
d'aquest tipus: per a cada objecte $X$ i cada parell de morfismes $f\colon X\to A$
i $g\colon X\to B$ existeix un \emph{únic} morfisme $u\colon X\to A\times B$ amb
$\pi_1\comp u=f$ i $\pi_2\comp u=g$, és a dir, que fa commutar el diagrama
\[
\begin{tikzpicture}[baseline=(m.center)]
  \matrix (m) [matrix of math nodes, row sep=2.6em, column sep=2.8em] {
     & X & \\
    A & A\times B & B \\
  };
  \draw[-{Stealth[scale=1.0]}] (m-1-2) -- node[above left] {$f$} (m-2-1);
  \draw[-{Stealth[scale=1.0]}] (m-1-2) -- node[above right] {$g$} (m-2-3);
  \draw[-{Stealth[scale=1.0]},dashed] (m-1-2) -- node[right] {$u$} (m-2-2);
  \draw[-{Stealth[scale=1.0]}] (m-2-2) -- node[below] {$\pi_1$} (m-2-1);
  \draw[-{Stealth[scale=1.0]}] (m-2-2) -- node[below] {$\pi_2$} (m-2-3);
\end{tikzpicture}
\]
El coproducte es defineix per la condició dual, amb les fletxes girades. La categoria lliure $\Free(G)$ la satisfà també: un functor
$\Free(G)\to\cat{D}$ queda determinat, de manera única, per un morfisme de grafs
$G\to U(\cat{D})$ (subsecció~\ref{subsec:free-u}). I els objectes inicial i
terminal, que fins ara només hem anomenat (secció~\ref{sec:dualitat}), en són el
cas més senzill. Aquestes condicions, que quantifiquen sobre \emph{totes} les
fletxes d'una categoria, s'anomenen \textbf{propietats
universals}\index{propietat universal}, i són, potser, la contribució més
característica de la teoria de categories. La seva virtut, que encara no hem demostrat en cap cas, és que determinen la \emph{dada universal} ---l'objecte juntament amb les fletxes que intervenen en la condició--- llevat d'un únic isomorfisme compatible. En particular, l'objecte subjacent queda determinat llevat d'isomorfisme, però pot admetre altres isomorfismes que no respectin la dada universal. Abans de donar-ne la forma general, les estudiem amb detall en dos
casos: el més simple, el dels objectes terminals i inicials, que definim ara, i el
de la construcció lliure, que ja coneixem.

\subsection{El cas més simple: objectes terminals i inicials}

El cas extrem, i el més fàcil de dibuixar, és el d'un objecte al qual s'arriba d'una
\emph{única} manera des de qualsevol altre, o del qual es \emph{surt} d'una única
manera cap a qualsevol altre. Diem que un objecte $T$ és \textbf{terminal} si per a
cada objecte $X$ hi ha exactament una fletxa $X\to T$, i que un objecte $I$ és
\textbf{inicial} si per a cada $X$ n'hi ha exactament una $I\to X$:
\[
\begin{tikzpicture}[baseline=(m.center)]
  \matrix (m) [matrix of math nodes, column sep=3.2em] { X & T \\ };
  \draw[-{Stealth[scale=1.1]}, dashed] (m-1-1) -- node[above] {$\exists!$} (m-1-2);
\end{tikzpicture}
\qquad\qquad
\begin{tikzpicture}[baseline=(m.center)]
  \matrix (m) [matrix of math nodes, column sep=3.2em] { I & X \\ };
  \draw[-{Stealth[scale=1.1]}, dashed] (m-1-1) -- node[above] {$\exists!$} (m-1-2);
\end{tikzpicture}
\]
En aquests diagrames seguim la convenció que ja vam introduir en presentar el producte
(subsecció~\ref{subsec:producte-coproducte}): la \textbf{fletxa discontínua} assenyala
el morfisme l'existència del qual afirma la propietat universal, i $\exists!$ abreuja
\enquote{existeix un únic}; com que aquí el morfisme encara no té nom, hi escrivim
directament $\exists!$. Cada diagrama es llegeix, doncs: per a tot $X$ existeix un únic
morfisme cap a $T$ (respectivament, des de $I$).

A $\cat{Set}$, per exemple, un objecte terminal és qualsevol conjunt d'un sol element i
un objecte inicial és el conjunt buit. El que distingeix aquests objectes no és cap
propietat dels seus elements, sinó la fletxa: la condició els fixa completament, fins
al punt que dos objectes terminals qualssevol són isomorfs \emph{d'una única manera},
com demostrarem tot seguit (proposició~\ref{prop:unicitat-terminal}). Aquesta és, en
miniatura, una propietat universal: una condició sobre \emph{totes} les fletxes cap a
un objecte ---o des d'ell--- que el determina llevat d'un únic isomorfisme. En aquest cas, l'isomorfisme és únic fins i tot com a isomorfisme d'objectes, perquè no hi ha cap dada addicional a preservar: entre dos objectes terminals només hi ha una fletxa en cada sentit. Escrivim
encara $T$ i $I$, i no $1$ i $0$, a propòsit: la notació numèrica dona per fet que
\enquote{l'} objecte terminal està ben determinat, cosa que només ens autoritzem a
afirmar un cop provada aquesta unicitat. A la secció següent la demostrem i, aleshores
sí, batejarem $T$ com a $1$ i $I$ com a $0$.

\subsection{Un exemple ja conegut: la construcció lliure}\label{subsec:construccio-lliure}

Un exemple menys trivial ja ha aparegut: la categoria lliure\index{categoria!lliure}
$\Free(G)$ generada per un graf (seccions~\ref{sec:grafs-lliures}
i~\ref{sec:functors-grafs}), i el seu cas d'un sol vèrtex, el monoide
lliure\index{monoide!lliure}. Allà vam deixar pendent de justificar el nom
\enquote{lliure}; ara podem llegir-lo com una propietat universal.

La propietat posa en joc dues categories diferents, i convé tenir present en
quina treballem a cada moment. D'una banda hi ha $\cat{Graph}$, on viuen els
grafs i els morfismes de grafs; de l'altra, $\cat{Cat}$, on viuen les categories
petites i els functors entre elles. Els functors del
capítol~\ref{cap:functors}, $\Free\colon\cat{Graph}\to\cat{Cat}$ i
$U\colon\cat{Cat}\to\cat{Graph}$, fan de pont entre tots dos nivells; no són
fletxes de cap d'aquestes categories ---són functors entre categories grans---,
sinó la manera de passar de l'una a l'altra. Així, $\Free(G)$ és un objecte de
$\cat{Cat}$, mentre que $U(\Free(G))$, el seu graf subjacent, és un objecte de
$\cat{Graph}$.

La clau és una \textbf{inserció dels generadors}, el morfisme de grafs
\[
\eta_G\colon G\to U(\Free(G)),
\]
que deixa fixos els vèrtexs i envia cada aresta al camí de longitud~1
corresponent. Aquesta inserció té la propietat següent: per a tota categoria
petita $\cat{D}$ i tot morfisme de grafs $f\colon G\to U(\cat{D})$ ---una fletxa
de $\cat{Graph}$, que tria un objecte de $\cat{D}$ per a cada vèrtex i un
morfisme de $\cat{D}$ per a cada aresta, sense cap exigència sobre
composicions---, existeix un \emph{únic} functor $\bar f\colon\Free(G)\to\cat{D}$
---una fletxa de $\cat{Cat}$--- que estén $f$. Com que $f$ i $\bar f$ no viuen a
la mateixa categoria, \enquote{estendre} no es pot expressar comparant-los
directament: primer apliquem $U$ a $\bar f$, cosa que en dona el morfisme de
grafs $U(\bar f)\colon U(\Free(G))\to U(\cat{D})$, i aleshores la condició és la
igualtat $U(\bar f)\comp\eta_G=f$ entre morfismes de $\cat{Graph}$.
\[
\begin{tikzpicture}[baseline=(m.center)]
  \matrix (m) [matrix of math nodes, column sep=3.4em, row sep=2.8em] {
    G & U(\Free(G)) \\
      & U(\cat{D}) \\
  };
  \draw[-{Stealth[scale=1.1]}] (m-1-1) -- node[above] {$\eta_G$} (m-1-2);
  \draw[-{Stealth[scale=1.1]}] (m-1-1) -- node[below left] {$f$} (m-2-2);
  \draw[-{Stealth[scale=1.1]}] (m-1-2) -- node[right] {$U(\bar f)$} (m-2-2);
  \node[anchor=south, font=\small] at ([yshift=0.6em]m.north) {a $\cat{Graph}$};
\end{tikzpicture}
\qquad\qquad
\begin{tikzpicture}[baseline=(m.center)]
  \matrix (m) [matrix of math nodes, row sep=2.8em] {
    \Free(G) \\
    \cat{D} \\
  };
  \draw[-{Stealth[scale=1.1]}, dashed] (m-1-1) -- node[right] {$\exists!\,\bar f$} (m-2-1);
  \node[anchor=south, font=\small] at ([yshift=0.6em]m.north) {a $\cat{Cat}$};
\end{tikzpicture}
\]
El triangle de l'esquerra commuta a $\cat{Graph}$. La fletxa discontínua, la
que la propietat afirma que existeix i és única, és $\bar f$, i viu a $\cat{Cat}$;
la seva imatge per $U$ és la que tanca el triangle.
Aquesta és, precisament, la correspondència que vam establir al
capítol~\ref{cap:functors}: definir un functor \emph{des de} $\Free(G)$ no demana res
més que dir on van els generadors. La construcció $\Free(G)$ no s'ha \emph{definit} per
aquesta propietat ---s'ha bastit amb camins---, però la propietat la
\emph{caracteritza}: qualsevol altra categoria petita $\cat{F}$ dotada d'un morfisme de grafs
$G\to U(\cat{F})$ amb la mateixa condició seria isomorfa a $\Free(G)$, i d'una única manera compatible amb la
inserció. En això consisteix una propietat universal. Dit en el llenguatge que introduirem a la secció de representabilitat, $\Free(G)$ representa el functor
\[
\cat{D}\longmapsto\Hom_{\cat{Graph}}\bigl(G,U(\cat{D})\bigr),
\]
i $\eta_G\in\Hom_{\cat{Graph}}\bigl(G,U(\Free(G))\bigr)$ n'és l'\emph{element universal}. La formulació completa d'aquesta correspondència com a adjunció $\Free\dashv U$ arribarà al capítol d'adjuncions.

En aquest capítol precisem què és una propietat universal, la lliguem amb els
homfunctors $\Hom(A,-)$ i $\Hom(-,B)$ que vam introduir, i culminem amb el
\emph{lema de Yoneda}, que fa explícita la idea que un objecte queda determinat pels
morfismes que hi arriben o que en surten.

\section{Objectes inicials i terminals}

\begin{definicio}\index{objecte!inicial}\index{objecte!terminal}
En una categoria $\cat{C}$:
\begin{itemize}
\item un objecte $I$ és \textbf{inicial} si per a cada objecte $X$ hi ha \emph{exactament un} morfisme $I\to X$;
\item un objecte $T$ és \textbf{terminal} si per a cada objecte $X$ hi ha \emph{exactament un} morfisme $X\to T$.
\end{itemize}
Les dues nocions són duals: un objecte inicial a $\cat{C}$ és un objecte terminal a $\cat{C}\op$.
\end{definicio}

\begin{exemple}
A $\cat{Set}$, l'objecte inicial és el conjunt buit $\varnothing$ ---hi ha una única aplicació $\varnothing\to X$, la buida--- i l'objecte terminal és qualsevol conjunt d'un sol element $\{\ast\}$ ---hi ha una única aplicació $X\to\{\ast\}$. A $\cat{Grp}$, en canvi, el grup trivial és alhora inicial i terminal (un tal objecte s'anomena \emph{zero}). En un preordre vist com a categoria, un objecte inicial és un \emph{mínim global}, un element $m$ amb $m\le x$ per a tot $x$, i un de terminal, un \emph{màxim global}. No s'han de confondre amb els elements minimals o maximals, que només demanen que no hi hagi cap element estrictament per sota o per sobre.
\end{exemple}

La propietat clau, model de totes les propietats universals, és la unicitat.

\begin{proposicio}\label{prop:unicitat-terminal}\index{objecte!unicitat llevat d'isomorfisme}
Si una categoria té dos objectes terminals, aleshores són isomorfs, i l'isomorfisme entre ells és únic. El mateix val per als objectes inicials.
\end{proposicio}

\begin{prova}
Siguin $T$ i $T'$ dos objectes terminals. Com que $T'$ és terminal, hi ha un únic morfisme $f\colon T\to T'$; com que $T$ és terminal, hi ha un únic morfisme $g\colon T'\to T$. La composició $g\comp f\colon T\to T$ és un morfisme de $T$ a $T$; però $T$ és terminal, de manera que l'únic morfisme $T\to T$ és $\id_T$, i per tant $g\comp f=\id_T$. Anàlogament $f\comp g=\id_{T'}$. Així $f$ és un isomorfisme, i és únic perquè ho és el morfisme $T\to T'$. El cas inicial és dual.
\end{prova}

Provada la unicitat llevat d'un únic isomorfisme, ja té sentit parlar de \enquote{l'} objecte terminal i \enquote{l'} objecte inicial. D'ara endavant els escrivim, respectivament, $1$ i $0$. Aquesta notació és estàndard i no pressuposa cap finitud: prové del fet que l'objecte terminal és el producte buit ---la unitat de $\times$--- i l'inicial el coproducte buit ---el zero de $+$---, com veurem al capítol~\ref{cap:limits}.

\begin{observacio}
Aquesta demostració és el patró que retrobarem sempre, però convé enunciar-lo amb la precisió justa. Dues \emph{dades universals} del mateix tipus ---dos objectes terminals, dos productes $(P,p_1,p_2)$ i $(Q,q_1,q_2)$, dues representacions--- són isomorfes mitjançant un \emph{únic} isomorfisme compatible amb \emph{tota} la dada universal (les projeccions, la representació, etc.). Els objectes \emph{subjacents} són, per tant, isomorfs; però l'isomorfisme entre els objectes \emph{nus}, un cop oblidada la dada, no ha de ser únic: un producte $P$ pot admetre automorfismes no trivials que no respecten les projeccions. En el cas dels objectes terminals i inicials no hi ha dada addicional a preservar, i per això l'isomorfisme és únic sense més (proposició~\ref{prop:unicitat-terminal}). Amb aquesta cautela parlem de \enquote{l'} objecte terminal, \enquote{el} producte, etc.: la unicitat \emph{compatible amb la dada} és el que dona sentit a caracteritzar objectes per la seva propietat universal en lloc de per una construcció concreta.
\end{observacio}

\section{Categories coma}\label{sec:coma}\index{categoria!coma}

Abans de continuar val la pena introduir una construcció que unifica diverses coses que hem trobat de passada ---les categories sobre i sota un objecte, i ben aviat els cons sobre un diagrama--- i que permet dir amb precisió què vol dir que un objecte és \enquote{universal}: n'hi haurà prou que sigui \emph{inicial} o \emph{terminal} en una categoria adequada.

\begin{definicio}[Categoria coma]\index{categoria!coma}
\label{def:coma}
Siguin $S\colon\cat{A}\to\cat{C}$ i $T\colon\cat{B}\to\cat{C}$ dos functors amb el mateix codomini. La \textbf{categoria coma} $(S\downarrow T)$ té:
\begin{itemize}
  \item per \textbf{objectes}, les ternes $(A,B,f)$ amb $A$ objecte de $\cat{A}$, $B$ objecte de $\cat{B}$ i $f\colon S(A)\to T(B)$ un morfisme de $\cat{C}$;
  \item per \textbf{morfismes} $(A,B,f)\to(A',B',f')$, les parelles $(a,b)$ amb $a\colon A\to A'$ a $\cat{A}$ i $b\colon B\to B'$ a $\cat{B}$ tals que el quadrat
  \[
  \begin{tikzpicture}[baseline=(m.center)]
    \matrix (m) [matrix of math nodes, column sep=3em, row sep=2.4em] {
      S(A) & S(A') \\
      T(B) & T(B') \\
    };
    \draw[-{Stealth[scale=1.1]}] (m-1-1) -- node[above] {$S(a)$} (m-1-2);
    \draw[-{Stealth[scale=1.1]}] (m-1-1) -- node[left] {$f$} (m-2-1);
    \draw[-{Stealth[scale=1.1]}] (m-1-2) -- node[right] {$f'$} (m-2-2);
    \draw[-{Stealth[scale=1.1]}] (m-2-1) -- node[below] {$T(b)$} (m-2-2);
  \end{tikzpicture}
  \]
  commuta a $\cat{C}$, és a dir, $T(b)\comp f=f'\comp S(a)$.
\end{itemize}
La composició i les identitats es fan component a component.
\end{definicio}

\begin{observacio}[Sobre el nom i la notació]\index{categoria!coma!notació}
El concepte i el nom es deuen a Lawvere, que a la seva tesi doctoral (1963)
escrivia aquesta categoria $(S,T)$, per analogia amb la notació $(A,B)$ dels
homconjunts: els objectes de la categoria coma són, en efecte, morfismes de
$S(A)$ a $T(B)$, una mena d'\enquote{homconjunt entre functors}. D'aquí ve el
nom de \emph{categoria coma}, que ha sobreviscut tot i que la coma, usada com a
operador, es prestava a confusió i ha estat substituïda. La notació
$(S\downarrow T)$, que és la de Mac Lane~\cite{maclane} i la que seguim, conserva
els parèntesis de l'original, que a més delimiten els operands quan són
expressions compostes. La fletxa es pot llegir com un dibuix dels objectes: al
quadrat anterior, cada objecte $f$ és una fletxa que \emph{baixa} de la fila de
$S$ a la fila de $T$. També es troben les notacions $S\downarrow T$, sense
parèntesis, i $(S/T)$, que generalitza la notació $\cat{C}/A$ de les categories
sobre un objecte (exemple següent).
\end{observacio}

La notació més freqüent fixa un dels functors en un objecte. Si $\cat{1}$ és la categoria terminal (un objecte i la seva identitat) i $C$ és un objecte de $\cat{C}$, escrivim $C$ també per al functor $\cat{1}\to\cat{C}$ que el tria. Aleshores:

\begin{exemple}[Categories sobre i sota un objecte com a categories coma]\index{categoria!sobre un objecte}\index{categoria!sota un objecte}\label{ex:slice-coma}
Prenem $S=\id_{\cat{C}}$ i $T=C\colon\cat{1}\to\cat{C}$. Un objecte de $(\id_{\cat{C}}\downarrow C)$ és una terna $(A,\ast,f)$ amb $f\colon A\to C$, és a dir, simplement un morfisme cap a $C$. Un morfisme $(A,\ast,f)\to(A',\ast,f')$ és una parella $(a,\id_\ast)$ amb $a\colon A\to A'$ tal que $f'\comp a=f$ (el quadrat de la definició, amb $T(\id_\ast)=\id_C$), és a dir, un triangle commutatiu sobre $C$:
\[
\begin{tikzpicture}[baseline=(m.center)]
  \matrix (m) [matrix of math nodes, column sep=1.6em, row sep=2.6em] {
    A & & A' \\
      & C &  \\
  };
  \draw[-{Stealth[scale=1.1]}] (m-1-1) -- node[above] {$a$} (m-1-3);
  \draw[-{Stealth[scale=1.1]}] (m-1-1) -- node[below left] {$f$} (m-2-2);
  \draw[-{Stealth[scale=1.1]}] (m-1-3) -- node[below right] {$f'$} (m-2-2);
\end{tikzpicture}
\]
Recuperem exactament la categoria $\cat{C}/C$ de $\cat{C}$ sobre $C$. Dualment, $(C\downarrow\id_{\cat{C}})$ és la categoria $C/\cat{C}$ de $\cat{C}$ sota $C$, la dels morfismes que \emph{surten} de $C$. Les categories sobre i sota un objecte són, doncs, els casos més simples de categoria coma.
\end{exemple}

\begin{exemple}[Objectes universals com a inicials o terminals]\label{ex:universal-coma}
Fem precís, en el cas més senzill, l'eslògan que motiva aquesta secció. Fixem
un objecte $C$ d'una categoria localment petita $\cat{C}$ i considerem el
prefeix $\Hom(-,C)\colon\cat{C}\op\to\cat{Set}$ (definició~\ref{def:homfunctors}),
que envia cada objecte $A$ a $\Hom(A,C)$ i cada morfisme $h\colon A\to A'$ a la
precomposició, que amb la notació~\ref{not:hom-morfismes} escrivim
$\Hom(h,C)=h^*$:
\[
\Hom(h,C)\colon\Hom(A',C)\to\Hom(A,C),\qquad x'\mapsto x'\comp h.
\]
La categoria $\cat{C}/C=(\id_{\cat{C}}\downarrow C)$ de l'exemple anterior es
descriu només amb aquest functor. Identificant cada terna $(A,\ast,x)$ amb la
parella $(A,x)$:
\begin{itemize}
  \item un \textbf{objecte} és una parella $(A,x)$ amb $A$ objecte de $\cat{C}$ i
  $x\in\Hom(A,C)$;
  \item un \textbf{morfisme} $(A,x)\to(A',x')$ és un morfisme $h\colon A\to A'$
  de $\cat{C}$ tal que
  \[
    \Hom(h,C)(x')=x,
    \qquad\text{és a dir,}\qquad
    x'\comp h=x,
  \]
  que és el triangle commutatiu sobre $C$:
  \[
  \begin{tikzpicture}[baseline=(m.center)]
    \matrix (m) [matrix of math nodes, column sep=1.6em, row sep=2.6em] {
      A & & A' \\
        & C &  \\
    };
    \draw[-{Stealth[scale=1.1]}] (m-1-1) -- node[above] {$h$} (m-1-3);
    \draw[-{Stealth[scale=1.1]}] (m-1-1) -- node[below left] {$x$} (m-2-2);
    \draw[-{Stealth[scale=1.1]}] (m-1-3) -- node[below right] {$x'$} (m-2-2);
  \end{tikzpicture}
  \]
\end{itemize}

\emph{Afirmació: $(C,\id_C)$ és un objecte terminal de $\cat{C}/C$.} Sigui
$(A,x)$ un objecte qualsevol. Un morfisme $(A,x)\to(C,\id_C)$ és un
$h\colon A\to C$ amb $\Hom(h,C)(\id_C)=x$. Com que
$\Hom(h,C)(\id_C)=\id_C\comp h=h$, la condició diu exactament $h=x$. Per tant
existeix un morfisme $(A,x)\to(C,\id_C)$, i només un: $h=x$.

Llegit en termes del functor, això diu que per a cada objecte $A$ l'aplicació
\[
\Hom(A,C)\longrightarrow\Hom(A,C),\qquad h\longmapsto\Hom(h,C)(\id_C),
\]
és bijectiva (de fet, és la identitat): tot element $x$ de qualsevol $\Hom(A,C)$
s'obté d'un \emph{únic} element distingit, $\id_C\in\Hom(C,C)$, aplicant-hi
$\Hom(h,C)$ per a un \emph{únic} $h$. Aquesta és precisament la forma de la
definició~\ref{def:representable} de la secció següent: $(C,\id_C)$ és una
representació de $\Hom(-,C)$, amb element universal $\id_C$. Allà veurem
(observació~\ref{obs:representacio-inicial}) que, per a qualsevol prefeix $P$
sobre $\cat{C}$ (definició~\ref{def:prefeix}), ser una representació de $P$ equival a ser un
objecte inicial d'una certa categoria coma construïda a partir de $P$; per a
$P=\Hom(-,C)$ aquesta categoria és $(\cat{C}/C)\op$, i ser-hi inicial és ser
terminal a $\cat{C}/C$.

Dualment, amb el homfunctor covariant $\Hom(C,-)$, $(C,\id_C)$ és un objecte
inicial de $C/\cat{C}$: un morfisme $(C,\id_C)\to(A,x)$ és un $h\colon C\to A$
amb $\Hom(C,h)(\id_C)=h\comp\id_C=x$, és a dir, $h=x$.

Un cas menys trivial és la construcció lliure de la
subsecció~\ref{subsec:construccio-lliure}. Sigui $G$ un graf, vist com a functor
$G\colon\cat{1}\to\cat{Graph}$, i sigui $U\colon\cat{Cat}\to\cat{Graph}$ el
functor oblit. Un objecte de la categoria coma $(G\downarrow U)$ és una parella
$(\cat{D},f)$ formada per una categoria petita $\cat{D}$ i un morfisme de grafs
$f\colon G\to U(\cat{D})$; un morfisme $(\cat{D},f)\to(\cat{D}',f')$ és un
functor $F\colon\cat{D}\to\cat{D}'$ amb $U(F)\comp f=f'$. Que
$(\Free(G),\eta_G)$ sigui un objecte inicial de $(G\downarrow U)$ vol dir que
per a cada $(\cat{D},f)$ hi ha un únic functor $\bar f\colon\Free(G)\to\cat{D}$
amb $U(\bar f)\comp\eta_G=f$: és, literalment, la propietat universal que hi
vam enunciar.

Aquests casos fan precís l'eslògan que unifica les propietats universals:
\emph{ser universal és ser inicial o terminal en una categoria coma adient}. Ho
retrobarem al capítol següent: un límit és un objecte terminal de la categoria dels
cons, que és una categoria coma, i un colímit, un objecte inicial de la dels cocons
(observació~\ref{obs:cons-coma}).
\end{exemple}

\section{Propietats universals com a representabilitat}

Podem reformular aquestes idees amb el llenguatge dels functors representables. Recordem que, en una categoria localment petita, cada objecte $A$ determina el prefeix $\Hom(-,A)\colon\cat{C}\op\to\cat{Set}$ i el functor covariant $\Hom(A,-)\colon\cat{C}\to\cat{Set}$.

Observem primer com es llegeixen les nocions anteriors. Que $1$ sigui terminal vol dir que $\Hom(X,1)$ té exactament un element per a cada $X$; és a dir, que el functor $\Hom(-,1)$ és \emph{naturalment isomorf} al prefeix constant $\Delta_{\{\ast\}}\colon\cat{C}\op\to\cat{Set}$ (exemple~\ref{ex:functor-constant}). Les components de l'isomorfisme són les úniques bijeccions $\Hom(X,1)\to\{\ast\}$, i la naturalitat és automàtica, perquè tots els quadrats acaben en un conjunt d'un sol element. Dualment, $0$ és inicial quan $\Hom(0,-)$ és naturalment isomorf al functor constant $\Delta_{\{\ast\}}\colon\cat{C}\to\cat{Set}$. En cap dels dos casos no és una igualtat literal: els conjunts d'un sol element $\Hom(X,1)$, i també els $\Hom(0,X)$, són conjunts diferents per a objectes $X$ diferents, i no s'han identificat per definició amb $\{\ast\}$.

Una forma general de les propietats universals que estudiarem és la representabilitat d'un functor cap a $\cat{Set}$: un objecte, juntament amb una dada universal, satisfà una propietat d'aquesta mena quan \emph{representa} el functor que recull les dades del problema.

\begin{definicio}[Representació]\index{propietat universal}\index{functor!representable}\index{element universal}
\label{def:representable}
Sigui $\cat{C}$ una categoria localment petita.
\begin{enumerate}
\item Sigui $F\colon\cat{C}\to\cat{Set}$ un functor covariant. Una \textbf{representació} de $F$ és una parella $(A,u)$ formada per un objecte $A$ de $\cat{C}$ i un element $u\in F(A)$, l'\textbf{element universal}, tal que per a cada objecte $X$ i cada $x\in F(X)$ existeix un únic morfisme $f\colon A\to X$ amb
\[
F(f)(u)=x.
\]
\item Sigui $P$ un prefeix sobre $\cat{C}$ (definició~\ref{def:prefeix}). Una \textbf{representació} de $P$ és una parella $(A,u)$ formada per un objecte $A$ de $\cat{C}$ i un element $u\in P(A)$, l'\textbf{element universal}, tal que per a cada objecte $X$ i cada $x\in P(X)$ existeix un únic morfisme $f\colon X\to A$ amb
\[
P(f)(u)=x.
\]
\end{enumerate}
En tots dos casos es diu que el functor és \textbf{representable}, que $A$ n'és l'\textbf{objecte representant} i que $(A,u)$ n'és una representació.
\end{definicio}

Amb la noció d'isomorfisme natural (capítol~\ref{cap:transformacions}), la representabilitat es pot formular sense esmentar l'element universal.

\begin{proposicio}[Representabilitat com a isomorfisme natural]\label{prop:representable-iso}\index{functor!representable}
Sigui $\cat{C}$ una categoria localment petita.
\begin{enumerate}
\item Un functor $F\colon\cat{C}\to\cat{Set}$ és representable si i només si és naturalment isomorf al functor $\Hom_{\cat{C}}(S,-)\colon\cat{C}\to\cat{Set}$ per a algun objecte $S$ de $\cat{C}$. En aquest cas, $S$ és un objecte representant de $F$.
\item Un prefeix $P\colon\cat{C}\op\to\cat{Set}$ és representable si i només si és naturalment isomorf al prefeix $\Hom_{\cat{C}}(-,S)\colon\cat{C}\op\to\cat{Set}$ per a algun objecte $S$ de $\cat{C}$. De nou, en aquest cas $S$ és un objecte representant de $P$.
\end{enumerate}
Més precisament, les representacions $(S,u)$ es corresponen amb els isomorfismes naturals
\[
\alpha\colon\Hom(S,-)\Rightarrow F
\qquad\text{o bé}\qquad
\alpha\colon\Hom(-,S)\Rightarrow P,
\]
segons el cas, mitjançant $u=\alpha_S(\id_S)$ i $\alpha_X(f)=F(f)(u)$, respectivament $\alpha_X(f)=P(f)(u)$.
\end{proposicio}

\begin{prova}
Fem primer el cas~(1). Els quadrats de naturalitat segueixen el criteri del capítol~\ref{cap:transformacions}: l'acció dels functors sobre un morfisme, en horitzontal (a dalt l'homfunctor, a baix $F$), i les components de la transformació, en vertical. Al costat de cada quadrat hi ha el camí que hi fa un element: les fletxes $\mapsto$ indiquen on va a parar cada element.

\emph{D'una representació a un isomorfisme natural.} Sigui $(S,u)$ una representació de $F$ i definim $\alpha_X\colon\Hom(S,X)\to F(X)$ per $\alpha_X(f)=F(f)(u)$. Per veure que $\alpha\colon\Hom(S,-)\Rightarrow F$ és natural, cal que per a cada $g\colon X\to Y$ commuti el quadrat següent, on l'acció de l'homfunctor sobre $g$ és la postcomposició $g_*=\Hom(S,g)$ (notació~\ref{not:hom-morfismes}):
\[
\begin{tikzpicture}[baseline=(m.center)]
  \matrix (m) [matrix of math nodes, row sep=2.8em, column sep=3.6em] {
    \Hom(S,X) & \Hom(S,Y) \\
    F(X) & F(Y) \\
  };
  \draw[-{Stealth[scale=1.1]}] (m-1-1) -- node[above] {$g_*$} (m-1-2);
  \draw[-{Stealth[scale=1.1]}] (m-1-1) -- node[left] {$\alpha_X$} (m-2-1);
  \draw[-{Stealth[scale=1.1]}] (m-1-2) -- node[right] {$\alpha_Y$} (m-2-2);
  \draw[-{Stealth[scale=1.1]}] (m-2-1) -- node[below] {$F(g)$} (m-2-2);
\end{tikzpicture}
\qquad\qquad
\begin{tikzpicture}[baseline=(m.center)]
  \matrix (m) [matrix of math nodes, row sep=2.8em, column sep=2.4em] {
    f & g\comp f \\
    F(f)(u) & F(g)\bigl(F(f)(u)\bigr) \\
  };
  \draw[{Bar[width=4pt]}-{Stealth[scale=1.1]}] (m-1-1) -- (m-1-2);
  \draw[{Bar[width=4pt]}-{Stealth[scale=1.1]}] (m-1-1) -- (m-2-1);
  \draw[{Bar[width=4pt]}-{Stealth[scale=1.1]}] (m-1-2) -- (m-2-2);
  \draw[{Bar[width=4pt]}-{Stealth[scale=1.1]}] (m-2-1) -- (m-2-2);
\end{tikzpicture}
\]
Seguint un element $f$ pels dos camins, per dalt i després avall s'obté $\alpha_Y(g\comp f)=F(g\comp f)(u)$, i avall i després per baix s'obté $F(g)\bigl(F(f)(u)\bigr)$. Totes dues coincideixen perquè $F$ preserva la composició:
\[
\alpha_Y(g\comp f)=F(g\comp f)(u)=F(g)\bigl(F(f)(u)\bigr)=F(g)\bigl(\alpha_X(f)\bigr).
\]
La condició de representació diu exactament que cada $\alpha_X$ és bijectiva, i per la proposició~\ref{prop:iso-natural-components} $\alpha$ és un isomorfisme natural.

\emph{D'un isomorfisme natural a una representació.} Recíprocament, sigui $\alpha\colon\Hom(S,-)\Rightarrow F$ un isomorfisme natural i posem $u=\alpha_S(\id_S)\in F(S)$. Per a cada $f\colon S\to X$, escrivim el quadrat de naturalitat de $\alpha$ respecte de $f$ i hi seguim l'element $\id_S$:
\[
\begin{tikzpicture}[baseline=(m.center)]
  \matrix (m) [matrix of math nodes, row sep=2.8em, column sep=3.6em] {
    \Hom(S,S) & \Hom(S,X) \\
    F(S) & F(X) \\
  };
  \draw[-{Stealth[scale=1.1]}] (m-1-1) -- node[above] {$f_*$} (m-1-2);
  \draw[-{Stealth[scale=1.1]}] (m-1-1) -- node[left] {$\alpha_S$} (m-2-1);
  \draw[-{Stealth[scale=1.1]}] (m-1-2) -- node[right] {$\alpha_X$} (m-2-2);
  \draw[-{Stealth[scale=1.1]}] (m-2-1) -- node[below] {$F(f)$} (m-2-2);
\end{tikzpicture}
\qquad\qquad
\begin{tikzpicture}[baseline=(m.center)]
  \matrix (m) [matrix of math nodes, row sep=2.8em, column sep=2.4em] {
    \id_S & f \\
    u & F(f)(u) \\
  };
  \draw[{Bar[width=4pt]}-{Stealth[scale=1.1]}] (m-1-1) -- (m-1-2);
  \draw[{Bar[width=4pt]}-{Stealth[scale=1.1]}] (m-1-1) -- (m-2-1);
  \draw[{Bar[width=4pt]}-{Stealth[scale=1.1]}] (m-1-2) -- (m-2-2);
  \draw[{Bar[width=4pt]}-{Stealth[scale=1.1]}] (m-2-1) -- (m-2-2);
\end{tikzpicture}
\]
Per dalt, $\id_S$ va a $f\comp\id_S=f$ i després a $\alpha_X(f)$; avall, va a $u$ i després a $F(f)(u)$. La commutativitat del quadrat dona, doncs,
\[
\alpha_X(f)=\alpha_X(f\comp\id_S)=F(f)\bigl(\alpha_S(\id_S)\bigr)=F(f)(u).
\]
Això diu que tota la transformació $\alpha$ queda determinada pel sol element $u$. Com que $\alpha_X$ és bijectiva (proposició~\ref{prop:iso-natural-components}), per a cada $x\in F(X)$ hi ha un únic $f\colon S\to X$ amb $F(f)(u)=x$: $(S,u)$ és una representació de $F$.

Les dues construccions són inverses l'una de l'altra. Partint de $u$, la transformació $\alpha$ torna $\alpha_S(\id_S)=F(\id_S)(u)=u$; i partint de $\alpha$, la igualtat anterior diu que $\alpha$ és precisament la transformació $f\mapsto F(f)(u)$ construïda a partir de $u=\alpha_S(\id_S)$.

El cas~(2) és el cas~(1) aplicat a $\cat{C}\op$, però val la pena veure com canvia el diagrama clau. Ara $\alpha\colon\Hom(-,S)\Rightarrow P$, els morfismes arriben a $S$ en lloc de sortir-ne, i l'acció de l'homfunctor sobre un morfisme és la precomposició $f^*=\Hom(f,S)$. Per a cada $f\colon X\to S$, el quadrat de naturalitat respecte de $f$, amb l'element $\id_S$, és
\[
\begin{tikzpicture}[baseline=(m.center)]
  \matrix (m) [matrix of math nodes, row sep=2.8em, column sep=3.6em] {
    \Hom(S,S) & \Hom(X,S) \\
    P(S) & P(X) \\
  };
  \draw[-{Stealth[scale=1.1]}] (m-1-1) -- node[above] {$f^*$} (m-1-2);
  \draw[-{Stealth[scale=1.1]}] (m-1-1) -- node[left] {$\alpha_S$} (m-2-1);
  \draw[-{Stealth[scale=1.1]}] (m-1-2) -- node[right] {$\alpha_X$} (m-2-2);
  \draw[-{Stealth[scale=1.1]}] (m-2-1) -- node[below] {$P(f)$} (m-2-2);
\end{tikzpicture}
\qquad\qquad
\begin{tikzpicture}[baseline=(m.center)]
  \matrix (m) [matrix of math nodes, row sep=2.8em, column sep=2.4em] {
    \id_S & f \\
    u & P(f)(u) \\
  };
  \draw[{Bar[width=4pt]}-{Stealth[scale=1.1]}] (m-1-1) -- (m-1-2);
  \draw[{Bar[width=4pt]}-{Stealth[scale=1.1]}] (m-1-1) -- (m-2-1);
  \draw[{Bar[width=4pt]}-{Stealth[scale=1.1]}] (m-1-2) -- (m-2-2);
  \draw[{Bar[width=4pt]}-{Stealth[scale=1.1]}] (m-2-1) -- (m-2-2);
\end{tikzpicture}
\]
on ara, per dalt, $\id_S$ va a $\id_S\comp f=f$. Com que els functors són contravariants, les fletxes horitzontals van de l'objecte $S$ a l'objecte $X$, en sentit contrari a $f\colon X\to S$. Se'n dedueix $\alpha_X(f)=P(f)(u)$, i la resta de l'argument és el mateix.
\end{prova}

Les dues nocions són duals l'una de l'altra. Un prefeix sobre $\cat{C}$ és un functor covariant $\cat{C}\op\to\cat{Set}$, i com que un morfisme $X\to A$ de $\cat{C}$ és un morfisme $A\to X$ de $\cat{C}\op$, una representació d'un prefeix en el sentit~(2) és exactament una representació en el sentit~(1) del functor covariant corresponent sobre $\cat{C}\op$. Per això, tot el que demostrem per a un dels dos casos val, per dualitat, per a l'altre. Per exemple, $0$ és inicial si i només si $(0,\ast)$ és una representació del functor covariant $\Delta_{\{\ast\}}$, i $1$ és terminal si i només si $(1,\ast)$ és una representació del prefeix $\Delta_{\{\ast\}}$.

\begin{observacio}[Sondes i pantalles]\label{obs:sondes-pantalles}\index{functor!representable!interpretació}
Els functors representables admeten una lectura intuïtiva. El functor representable $\Hom_{\cat{C}}(S,-)\colon\cat{C}\to\cat{Set}$ pren un objecte $X$ i en dona el conjunt $\Hom_{\cat{C}}(S,X)$ de fletxes de $S$ a $X$. Les característiques de $X$ que extreu són exactament totes les maneres possibles de fer arribar $S$ dins de $X$, com si $S$ fos una \emph{sonda} amb què explorem l'estructura de $X$. Dir que un functor $F$ és representable per $S$ és dir que el que $F$ mesura de cada objecte és precisament el que en veu aquesta sonda.

Dualment, el prefeix representable $\Hom_{\cat{C}}(-,S)\colon\cat{C}\op\to\cat{Set}$ pren un objecte $X$ i en dona el conjunt $\Hom_{\cat{C}}(X,S)$ de fletxes de $X$ a $S$. Les característiques de $X$ que extreu són totes les maneres possibles d'enviar $X$ a $S$, o totes les observacions de $X$ amb valors a $S$: podem pensar $S$ com una \emph{pantalla} sobre la qual $X$ es projecta de moltes maneres diferents.
\end{observacio}

Vegem-ne exemples. En tots, l'isomorfisme natural es pot comprovar directament o, equivalentment, n'hi ha prou de donar l'element universal (proposició~\ref{prop:representable-iso}).

\begin{exemple}[Vèrtexs i arestes d'un graf]\label{ex:rep-grafs}\index{functor!representable!a Graph}
A l'exemple~\ref{ex:homfunctor-grafs} vam veure que els functors $\cat{Graph}\to\cat{Set}$ \enquote{conjunt de vèrtexs} i \enquote{conjunt d'arestes} són naturalment isomorfs a $\Hom_{\cat{Graph}}(\mathbf V,-)$ i a $\Hom_{\cat{Graph}}(\mathbf E,-)$. Són, doncs, representables, amb objectes representants el graf $\mathbf V$ d'un sol vèrtex i el graf $\mathbf E$ d'una sola aresta. Els elements universals són l'únic vèrtex $\ast$ de $\mathbf V$ i l'única aresta $e$ de $\mathbf E$: tot vèrtex d'un graf $G$ és la imatge de $\ast$ per un únic morfisme $\mathbf V\to G$, i tota aresta és la imatge de $e$ per un únic morfisme $\mathbf E\to G$. Són sondes en el sentit literal: $\mathbf V$ detecta els vèrtexs i $\mathbf E$ detecta les arestes.
\end{exemple}

\begin{exemple}[Functors oblit de $\cat{Grp}$, $\cat{Mon}$ i $\cat{Vect}_{\mathbb K}$]\label{ex:rep-grp}\index{functor!representable!oblit}
El functor oblit $U\colon\cat{Grp}\to\cat{Set}$ és representable, amb objecte representant el grup additiu $\mathbb Z$ i element universal $1\in U(\mathbb Z)$. En efecte, un homomorfisme $\varphi\colon\mathbb Z\to G$ queda determinat per $g=\varphi(1)$, perquè
\[
\varphi(n)=g^n\qquad\text{per a tot } n\in\mathbb Z;
\]
i, recíprocament, per a cada $g\in G$ l'aplicació $n\mapsto g^n$ és un homomorfisme. Per tant, $\varphi\mapsto\varphi(1)$ és una bijecció
\[
\Hom_{\cat{Grp}}(\mathbb Z,G)\;\cong\;U(G),
\]
natural en $G$, perquè per a un homomorfisme $\psi\colon G\to H$ es té $(\psi\comp\varphi)(1)=\psi(\varphi(1))$.

El mateix argument val per als altres dos models algebraics del capítol~\ref{cap:categories}. El functor oblit $\cat{Mon}\to\cat{Set}$ és representable pel monoide $(\mathbb N,+,0)$ amb element universal $1$: un homomorfisme de monoides $h\colon\mathbb N\to M$ queda determinat per $m=h(1)$, perquè $h(n)=m\ast\cdots\ast m$ ($n$ vegades, i $h(0)$ el neutre), i per a cada $m$ aquesta fórmula defineix un homomorfisme. El functor oblit $\cat{Vect}_{\mathbb K}\to\cat{Set}$ és representable per $\mathbb K$, vist com a espai vectorial sobre si mateix, amb element universal $1$: una aplicació lineal $T\colon\mathbb K\to V$ queda determinada per $v=T(1)$, perquè $T(\lambda)=\lambda v$, i per a cada $v\in V$ l'aplicació $\lambda\mapsto\lambda v$ és lineal.

En els tres casos l'objecte representant és l'estructura \enquote{generada lliurement per un sol element}, i com a sonda veu exactament els elements de cada objecte.
\end{exemple}

\begin{exemple}[El punt a $\cat{Set}$]\label{ex:rep-punt}
El functor identitat de $\cat{Set}$ és representable per un conjunt d'un sol element $\{\ast\}$, amb element universal $\ast$: una aplicació $\{\ast\}\to X$ és el mateix que un element de $X$, i $\Hom_{\cat{Set}}(\{\ast\},X)\cong X$ naturalment en $X$. El punt és la sonda que veu els elements d'un conjunt.
\end{exemple}

\begin{exemple}[Subconjunts i valors de veritat]\label{ex:rep-parts}\index{functor!representable!de parts}
El functor de parts $\mathcal P\colon\cat{Set}\op\to\cat{Set}$ és un prefeix representable: a l'exemple~\ref{ex:functor-parts} vam veure que les funcions característiques donen bijeccions canòniques $\mathcal P(X)\cong\Hom_{\cat{Set}}(X,\{0,1\})$ sota les quals l'antiimatge correspon a la precomposició. En el llenguatge del capítol~\ref{cap:transformacions}, això diu exactament que $\mathcal P\cong\Hom_{\cat{Set}}(-,\{0,1\})$ naturalment: el quadrat de naturalitat és la igualtat $\chi_B\comp f=\chi_{f^{-1}(B)}$ (exemple~\ref{ex:isos-naturals-cap2}). L'objecte representant és el conjunt $\{0,1\}$ de valors de veritat, i l'element universal és $\{1\}\in\mathcal P(\{0,1\})$, el subconjunt dels valors \enquote{cert}: per a cada $B\subseteq X$, l'única aplicació $\chi\colon X\to\{0,1\}$ amb $\chi^{-1}(\{1\})=B$ és la funció característica $\chi_B$.

Aquí $\{0,1\}$ és la pantalla. Un subconjunt $B\subseteq X$ és el mateix que un predicat sobre $X$, i el predicat és una observació de $X$ amb valors de veritat: $\chi_B(x)=1$ vol dir que $x$ compleix el predicat. Que els subconjunts es puguin classificar així, per morfismes cap a un objecte fix de valors de veritat, és una idea que tornarà amb força al capítol~\ref{cap:topos}, on l'objecte $\{0,1\}$ es generalitza en el \emph{classificador de subobjectes} $\Omega$ (definició~\ref{def:classificador}).
\end{exemple}

\begin{exemple}[El producte com a representació]\label{ex:rep-producte}\index{producte!com a representació}
Siguin $A$ i $B$ objectes d'una categoria localment petita $\cat{C}$. L'assignació
\[
X\longmapsto\Hom(X,A)\times\Hom(X,B),
\qquad
f\longmapsto\bigl((a,b)\mapsto(a\comp f,\,b\comp f)\bigr),
\]
és un prefeix sobre $\cat{C}$: preserva identitats i, per a $f\colon X\to Y$ i $g\colon Y\to Z$, $(a\comp g\comp f,\,b\comp g\comp f)$ s'obté aplicant primer $g$ i després $f$. Una representació d'aquest prefeix és un objecte $Q$ amb un element universal $(p_1,p_2)\in\Hom(Q,A)\times\Hom(Q,B)$ tal que, per a cada $X$ i cada parella $(a,b)$, hi ha un únic $u\colon X\to Q$ amb $p_1\comp u=a$ i $p_2\comp u=b$. És exactament la definició de producte de la subsecció~\ref{subsec:producte-coproducte}: el producte de $A$ i $B$ existeix si i només si aquest prefeix és representable, i un producte és una representació seva.

Aquest exemple té també una lectura lògica. A la categoria prima $\cat{Prov}_T$ associada a un sistema deductiu $T$ (capítol~\ref{cap:categories}), entre dues fórmules hi ha com a màxim un morfisme, i n'hi ha un exactament quan hi ha deduïbilitat $p\vdash_T q$. Per a dues fórmules $p$ i $q$, el prefeix
\[
r\longmapsto\Hom_{\cat{Prov}_T}(r,p)\times\Hom_{\cat{Prov}_T}(r,q)
\]
és, doncs, no buit exactament quan $r\vdash_T p$ i $r\vdash_T q$. Una representació és una fórmula $c$ amb $c\vdash_T p$ i $c\vdash_T q$ (l'element universal) tal que tota fórmula $r$ que dedueix $p$ i $q$ dedueix $c$; la unicitat del morfisme és automàtica. Per transitivitat, això equival a
\[
r\vdash_T c
\quad\Longleftrightarrow\quad
r\vdash_T p\ \text{i}\ r\vdash_T q,
\]
que són les regles d'introducció i d'eliminació de la conjunció. La conjunció $p\wedge q$, quan existeix, és doncs una representació d'aquest prefeix: satisfà la propietat universal del producte, llegida en lògica, tal com anunciava la secció~\ref{sec:dualitat}.
\end{exemple}

\begin{exemple}[Functors que no són representables]\label{ex:no-representable}
Si $\cat{C}$ té algun objecte, el functor constant $\Delta_\varnothing\colon\cat{C}\to\cat{Set}$ no és representable: per a tot objecte $S$, el conjunt $\Hom(S,S)$ conté $\id_S$ i no és buit, mentre que $\Delta_\varnothing(S)=\varnothing$, de manera que no hi pot haver cap bijecció $\Hom(S,S)\to\Delta_\varnothing(S)$. El mateix argument val per al prefeix $\Delta_\varnothing$. Això explica, per exemple, que en una categoria discreta amb dos objectes $A\neq B$ (només hi ha identitats) no hi hagi producte de $A$ i $B$: el prefeix de l'exemple~\ref{ex:rep-producte} hi val $\Hom(X,A)\times\Hom(X,B)=\varnothing$ per a tot $X$, perquè cap objecte té fletxes cap a tots dos, i és, doncs, el prefeix $\Delta_\varnothing$.
\end{exemple}

\begin{observacio}[Representacions com a objectes inicials]\label{obs:representacio-inicial}
Això fa precís l'eslògan de la secció~\ref{sec:coma}. Sigui
$P$ un prefeix sobre $\cat{C}$ i sigui
$\{\ast\}\colon\cat{1}\to\cat{Set}$ el functor que tria un conjunt d'un sol
element. Un objecte de la categoria coma $(\{\ast\}\downarrow P)$ és una terna
$(\ast,A,x)$ amb $x\colon\{\ast\}\to P(A)$, és a dir, un objecte $A$ de
$\cat{C}$ junt amb un element $x\in P(A)$; l'escrivim $(A,x)$. Un morfisme
$(A,x)\to(A',x')$ és un morfisme $A\to A'$ de $\cat{C}\op$ ---és a dir, un
$h\colon A'\to A$ de $\cat{C}$--- que fa commutar el quadrat de la
definició~\ref{def:coma}, condició que aquí diu $P(h)(x)=x'$. Per tant,
$(A,u)$ és un objecte inicial de $(\{\ast\}\downarrow P)$ si i només si per a
cada $(X,x)$ hi ha un únic $h\colon X\to A$ amb $P(h)(u)=x$: exactament la
condició de representació. Per a $P=\Hom(-,C)$, la categoria
$(\{\ast\}\downarrow P)$ és $(\cat{C}/C)\op$, i retrobem
l'exemple~\ref{ex:universal-coma}.
\end{observacio}

\begin{observacio}[La categoria d'elements]\label{obs:categoria-elements}\index{categoria!d'elements}
A la bibliografia és habitual presentar la mateixa construcció amb les fletxes en el sentit de $\cat{C}$. Per a un prefeix $P\colon\cat{C}\op\to\cat{Set}$, la \textbf{categoria d'elements} $\textstyle\int P$ té per objectes les parelles $(A,x)$ amb $x\in P(A)$, i per morfismes $f\colon(A,x)\to(B,y)$ els morfismes $f\colon A\to B$ de $\cat{C}$ amb $P(f)(y)=x$. Té, doncs, les mateixes dades que la categoria coma de l'observació anterior, amb les fletxes girades:
\[
\textstyle\int P\;\cong\;(\{\ast\}\downarrow P)\op.
\]
Per tant, una representació $(C,u)$ de $P$ correspon a un objecte \emph{terminal} de $\int P$: per a cada $(A,x)$ hi ha un únic $f\colon A\to C$ amb $P(f)(u)=x$. És la formulació que fa servir, per exemple, \cite[cap.~2]{riehl}.
\end{observacio}

\begin{observacio}
Cal distingir dos nivells d'unicitat. L'\textbf{objecte} representant $A$ és únic llevat d'isomorfisme, però no necessàriament d'un únic isomorfisme. En canvi, la \textbf{representació} $(A,u)$ ---amb la dada de l'element universal--- és única llevat d'un \emph{únic} isomorfisme compatible amb els elements universals: si $(A,u)$ i $(B,v)$ representen $P$, hi ha un únic isomorfisme $\varphi\colon A\to B$ amb $P(\varphi)(v)=u$. Aquesta distinció, entre la unicitat de l'objecte i la de la dada universal, és essencial i la retrobarem en tractar límits, adjuncions i objectes classificadors. Que dos functors representables $\Hom(-,A)$ i $\Hom(-,B)$ siguin naturalment isomorfs equivaldrà, pel lema de Yoneda, a $A\cong B$; la resta del capítol ho fa precís.
\end{observacio}

\section{La immersió de Yoneda}

Fixem una categoria \textbf{petita} $\cat{C}$. La noció de prefeix, i la de transformació natural entre prefeixos, té sentit sobre qualsevol categoria (definició~\ref{def:prefeix}). Els prefeixos representables $\Hom(-,A)$, en canvi, només són prefeixos si cada $\Hom(X,A)$ és un conjunt, és a dir, si $\cat{C}$ és localment petita, que és la hipòtesi amb què els hem fet servir a la secció de representabilitat. Aquí demanem més: que $\cat{C}$ sigui petita. Aquesta restricció afecta la categoria que aplega els prefeixos. Si $\cat{C}$ no és petita, encara que sigui localment petita, cada prefeix recorre una classe pròpia d'objectes i les transformacions naturals entre dos prefeixos poden formar també una classe pròpia. Aleshores no es pot formar $\cat{Set}^{\cat{C}\op}$ com a categoria (observació~\ref{obs:mida-functors}). Amb $\cat{C}$ petita, en canvi, la categoria de prefeixos $\widehat{\cat{C}}$ i tots els seus homconjunts queden ben definits (exemple~\ref{ex:categoria-prefeixos}).\footnote{Amb universos de Grothendieck, la construcció es pot fer també per a $\cat{C}$ només localment petita, formant $\widehat{\cat{C}}$ en un univers superior; vegeu l'apèndix~\ref{cap:mida} i \cite[§1.1]{riehl}.} A cada objecte $A$ li hem associat el prefeix representable $\Hom(-,A)\colon\cat{C}\op\to\cat{Set}$. Aquesta assignació és, ella mateixa, functorial: dona un functor de $\cat{C}$ cap a la categoria de prefeixos $\widehat{\cat{C}}=\cat{Set}^{\cat{C}\op}$.

\begin{definicio}[Immersió de Yoneda]\index{Yoneda!immersió}
\label{def:immersio-yoneda}
La \textbf{immersió de Yoneda} (en anglès, \emph{Yoneda embedding}) és el functor
\[
\yon\colon\cat{C}\to\cat{Set}^{\cat{C}\op},
\qquad
\yon(A)=\Hom(-,A),
\]
que envia cada objecte $A$ al functor representable $\Hom(-,A)$, i cada morfisme $f\colon A\to B$ a la transformació natural $\yon(f)\colon\Hom(-,A)\Rightarrow\Hom(-,B)$ de components la postcomposició amb $f$,
\[
\yon(f)_X=\Hom(X,f)\colon\Hom(X,A)\to\Hom(X,B),
\qquad
\varphi\mapsto f\comp\varphi.
\]
D'acord amb la notació dels homfunctors, escriurem també $\Hom(-,f)$\index{Hom@$\Hom$!$\Hom(-,f)$} per a aquesta transformació natural $\yon(f)$.
\end{definicio}

\begin{observacio}
Cal comprovar que $\yon(f)$ és realment natural i que $\yon$ preserva composició i identitats, però tot es redueix, un cop més, a l'associativitat i a les lleis de la identitat de $\cat{C}$. El contingut de la immersió de Yoneda és que \emph{tota} categoria petita s'immergeix en la seva categoria de prefeixos, que és una categoria molt rica: com veurem al capítol següent, té tots els límits i colímits \emph{petits}, calculats puntualment. El teorema següent en justifica el nom: $\yon$ és una immersió plena i fidel.
\end{observacio}

\section{El lema de Yoneda}

El resultat central relaciona les transformacions naturals que surten d'un functor representable amb els elements del functor de destí.

\begin{teorema}[Lema de Yoneda]\index{Yoneda!lema}
\label{teo:yoneda}
Sigui $\cat{C}$ una categoria \textbf{petita}, $A$ un objecte de $\cat{C}$ i $P$ un prefeix sobre $\cat{C}$. Aleshores hi ha una bijecció
\[
\Nat\big(\Hom(-,A),\,P\big)\;\cong\;P(A),
\]
natural \emph{contravariantment} en $A\in\cat{C}$ i \emph{covariantment} en $P\in\widehat{\cat{C}}$. Explícitament, la bijecció envia una transformació natural $\eta$ a l'element $\eta_A(\id_A)\in P(A)$.\footnotemark
\end{teorema}
\footnotetext{Si $\cat{C}$ és només localment petita, la mateixa construcció classifica les transformacions naturals que surten del representable mitjançant els elements de $P(A)$, sense haver de formar la categoria $\widehat{\cat{C}}$ de tots els prefeixos. En aquest cas, $\Nat(\Hom(-,A),P)$ s'ha d'entendre com aquesta col·lecció parametritzada pels elements de $P(A)$: en el marc de classes de l'apèndix~\ref{cap:mida}, una transformació amb domini gran és una família indexada per una classe, i no esdevé per això element d'un conjunt. Per parlar-ne literalment com d'un homconjunt dins d'una categoria de prefeixos cal fixar una convenció de mida addicional, per exemple amb universos; vegeu \cite[§1.1]{riehl}.}%

\begin{prova}
Descrivim les dues aplicacions inverses.

\textbf{De transformacions a elements.} A cada $\eta\colon\Hom(-,A)\Rightarrow P$ li assignem
\[
\Phi(\eta)=\eta_A(\id_A)\in P(A).
\]
Té sentit perquè $\id_A\in\Hom(A,A)$ i $\eta_A\colon\Hom(A,A)\to P(A)$.

\textbf{D'elements a transformacions.} A cada $x\in P(A)$ li assignem la transformació $\Psi(x)$ de components
\[
\Psi(x)_X\colon\Hom(X,A)\to P(X),
\qquad
\varphi\mapsto P(\varphi)(x).
\]
Té sentit perquè $\varphi\colon X\to A$ dona, per contravariància, $P(\varphi)\colon P(A)\to P(X)$, i $x\in P(A)$. Cal veure que $\Psi(x)$ és natural: per a un morfisme $g\colon Y\to X$, el quadrat de naturalitat és $P(g)\comp\Psi(x)_X=\Psi(x)_Y\comp g^*$, que sobre $\varphi\in\Hom(X,A)$ dona
\[
P(g)\big(P(\varphi)(x)\big)=P(\varphi\comp g)(x)=\Psi(x)_Y(\varphi\comp g)=\Psi(x)_Y\big(g^*(\varphi)\big),
\]
on hem usat la functorialitat $P(g)\comp P(\varphi)=P(\varphi\comp g)$.

\textbf{Són inverses.} D'una banda,
\[
\Phi(\Psi(x))=\Psi(x)_A(\id_A)=P(\id_A)(x)=x.
\]
De l'altra, donada $\eta$, cal veure que $\Psi(\eta_A(\id_A))=\eta$; és a dir, que per a tot $X$ i tot $\varphi\colon X\to A$,
\[
P(\varphi)\big(\eta_A(\id_A)\big)=\eta_X(\varphi).
\]
Això és exactament el quadrat de naturalitat de $\eta$ per al morfisme $\varphi\colon X\to A$, aplicat a $\id_A\in\Hom(A,A)$: en efecte, $\varphi^*(\id_A)=\id_A\comp\varphi=\varphi$, de manera que
\[
\eta_X(\varphi)=\eta_X(\varphi^*(\id_A))=P(\varphi)(\eta_A(\id_A)).
\]
Les dues assignacions són, doncs, inverses. Queda comprovar que la bijecció és \emph{natural} en les dues variables: és part de l'enunciat, no un afegit, i les dues verificacions són breus. Escrivim $\Phi_{A,P}\colon\Nat(\Hom(-,A),P)\to P(A)$ per exhibir-ne la dependència.

\emph{Naturalitat en $P$} (covariant). Sigui $\sigma\colon P\Rightarrow Q$ una transformació natural de prefeixos, amb postcomposició $(\sigma\comp-)\colon\Nat(\Hom(-,A),P)\to\Nat(\Hom(-,A),Q)$. Per a $\eta\colon\Hom(-,A)\Rightarrow P$,
\[
\Phi_{A,Q}(\sigma\comp\eta)=(\sigma\comp\eta)_A(\id_A)=\sigma_A\big(\eta_A(\id_A)\big)=\sigma_A\big(\Phi_{A,P}(\eta)\big),
\]
és a dir $\Phi_{A,Q}\comp(\sigma\comp-)=\sigma_A\comp\Phi_{A,P}$.

\emph{Naturalitat en $A$} (contravariant). Sigui $a\colon A'\to A$ un morfisme de $\cat{C}$. La immersió dona $\yon(a)=\Hom(-,a)\colon\Hom(-,A')\Rightarrow\Hom(-,A)$, i la precomposició indueix $(-\comp\yon(a))\colon\Nat(\Hom(-,A),P)\to\Nat(\Hom(-,A'),P)$. Per a $\eta\colon\Hom(-,A)\Rightarrow P$, com que $\yon(a)_{A'}(\id_{A'})=a\comp\id_{A'}=a$,
\[
\Phi_{A',P}\big(\eta\comp\yon(a)\big)=\big(\eta\comp\yon(a)\big)_{A'}(\id_{A'})=\eta_{A'}(a)=P(a)\big(\eta_A(\id_A)\big)=P(a)\big(\Phi_{A,P}(\eta)\big),
\]
on la penúltima igualtat és la identitat de naturalitat de $\eta$ demostrada més amunt (la que dona $\Psi(\eta_A(\id_A))=\eta$), aplicada amb $X=A'$ i $\varphi=a$; és a dir $\Phi_{A',P}\comp(-\comp\yon(a))=P(a)\comp\Phi_{A,P}$. Aquestes dues identitats completen el lema tal com s'ha enunciat.
\end{prova}

\begin{observacio}[Taula de control: variància i tipus]\label{obs:taula-yoneda}\index{Yoneda!taula de variància i tipus}
Els errors més freqüents d'aquest capítol són errors de tipus: una fletxa que va en el sentit equivocat o una composició que no existeix. Abans de qualsevol càlcul convé escriure el domini i el codomini de cada component, com al capítol~\ref{cap:functors} (observació~\ref{obs:comprovar-tipus}). La taula recull els casos que hem fet servir, amb $f\colon A\to B$, $u\colon X'\to X$, $v\colon X\to Y$, $a\colon A'\to A$ i $\sigma\colon P\Rightarrow Q$.
\begin{center}
\footnotesize
\renewcommand{\arraystretch}{1.35}
\setlength{\tabcolsep}{4pt}
\begin{tabularx}{\linewidth}{@{}>{\raggedright\arraybackslash}p{0.22\linewidth}>{\raggedright\arraybackslash}p{0.24\linewidth}>{\raggedright\arraybackslash}X@{}}
\toprule
\textbf{Expressió} & \textbf{Què és i variància} & \textbf{Tipus i acció}\\
\midrule
$\Hom(-,A)$ & prefeix $\cat{C}\op\to\cat{Set}$ (contravariant) & $X\mapsto\Hom(X,A)$; \ $\Hom(u,A)=u^*\colon\Hom(X,A)\to\Hom(X',A)$, $\psi\mapsto\psi\comp u$\\
$\Hom(A,-)$ & functor $\cat{C}\to\cat{Set}$ (covariant) & $X\mapsto\Hom(A,X)$; \ $\Hom(A,v)=v_*\colon\Hom(A,X)\to\Hom(A,Y)$, $\varphi\mapsto v\comp\varphi$\\
$\Hom(-,f)=\yon(f)$ & transformació natural $\Hom(-,A)\Rightarrow\Hom(-,B)$, en el sentit de $f$ (no és un functor) & component en $X$: $\Hom(X,f)\colon\Hom(X,A)\to\Hom(X,B)$, $\varphi\mapsto f\comp\varphi$\\
$\Hom(f,-)$ & transformació natural $\Hom(B,-)\Rightarrow\Hom(A,-)$, en sentit \emph{contrari} a $f$ & component en $X$: $\Hom(f,X)\colon\Hom(B,X)\to\Hom(A,X)$, $\psi\mapsto\psi\comp f$\\
$\yon$ & functor $\cat{C}\to\cat{Set}^{\cat{C}\op}$ (covariant) & $A\mapsto\Hom(-,A)$; \ $f\mapsto\Hom(-,f)\colon\yon(A)\Rightarrow\yon(B)$\\
\midrule
naturalitat de $\Hom(-,f)$ en $X$ & quadrat per a $u\colon X'\to X$ & $\Hom(X',f)\comp\Hom(u,A)=\Hom(u,B)\comp\Hom(X,f)\colon\Hom(X,A)\to\Hom(X',B)$; \ tots dos envien $\varphi$ a $f\comp\varphi\comp u$\\
lema de Yoneda, naturalitat en $P$ & covariant, per a $\sigma\colon P\Rightarrow Q$ & $\Phi_{A,Q}\comp(\sigma\comp-)=\sigma_A\comp\Phi_{A,P}\colon\Nat(\yon(A),P)\to Q(A)$\\
lema de Yoneda, naturalitat en $A$ & contravariant, per a $a\colon A'\to A$ & $\Phi_{A',P}\comp(-\comp\yon(a))=P(a)\comp\Phi_{A,P}\colon\Nat(\yon(A),P)\to P(A')$\\
\bottomrule
\end{tabularx}
\end{center}
La regla que hi ha darrere és sempre la mateixa. A $\Hom(X,Y)$, un morfisme que actua sobre la \emph{segona} variable ho fa per postcomposició i conserva el sentit; un que actua sobre la \emph{primera} ho fa per precomposició i l'inverteix. Per això $\Hom(-,f)$ va en el sentit de $f$ i $\Hom(f,-)$ en el contrari, i per això la naturalitat del lema és covariant en $P$ i contravariant en $A$.
\end{observacio}

\section{Conseqüències}

La primera conseqüència és que la immersió de Yoneda mereix el seu nom.

\begin{corollari}[La immersió de Yoneda és plena i fidel]\index{Yoneda!immersió plena i fidel}
Per a tota parella d'objectes $A,B$ de $\cat{C}$, la immersió $\yon$ dona una bijecció
\[
\Hom(A,B)\;\cong\;\Nat\big(\Hom(-,A),\,\Hom(-,B)\big).
\]
En particular, $\yon\colon\cat{C}\to\cat{Set}^{\cat{C}\op}$ és ple i fidel.
\end{corollari}

\begin{prova}
Apliquem el lema de Yoneda amb $P=\Hom(-,B)$. Aleshores
\[
\Nat\big(\Hom(-,A),\,\Hom(-,B)\big)\cong \Hom(-,B)(A)=\Hom(A,B),
\]
i seguint la bijecció es comprova que és precisament $f\mapsto\yon(f)$. Per tant $\yon$ és bijectiu sobre cada conjunt de morfismes, és a dir, ple i fidel.
\end{prova}

\begin{corollari}\index{Yoneda!objectes determinats per morfismes}
Dos objectes $A$ i $B$ de $\cat{C}$ són isomorfs si i només si els seus functors representables ho són:
\[
A\cong B
\quad\Longleftrightarrow\quad
\Hom(-,A)\cong\Hom(-,B).
\]
\end{corollari}

\begin{prova}
Un functor ple i fidel reflecteix els isomorfismes (capítol dels functors): si $\yon(A)\cong\yon(B)$, aleshores $A\cong B$. El recíproc és que tot functor preserva isomorfismes. Aplicat a $\yon$, dona l'equivalència.
\end{prova}

\begin{observacio}
Aquest darrer corol·lari és la formulació precisa de l'eslògan del primer capítol: \emph{un objecte queda determinat, llevat d'isomorfisme, per la xarxa de morfismes que hi arriben}. Conèixer $\Hom(-,A)$ ---és a dir, saber com és el conjunt de morfismes $X\to A$ per a cada $X$, de manera compatible amb la composició--- equival a conèixer $A$ llevat d'isomorfisme. En particular, l'objecte que representa un functor representable és únic llevat d'isomorfisme, cosa que unifica totes les unicitats de les propietats universals que hem trobat.
\end{observacio}

\begin{observacio}
Hi ha una versió covariant de tot plegat, obtinguda treballant amb $\cat{C}\op$ en lloc de $\cat{C}$: dona una immersió $\cat{C}\op\to\cat{Set}^{\cat{C}}$, també plena i fidel, i el lema $\Nat(\Hom(A,-),Q)\cong Q(A)$ per a functors covariants $Q\colon\cat{C}\to\cat{Set}$. Totes dues versions s'usen contínuament; la contravariant, amb prefeixos, és la que enllaça amb la teoria de feixos i la lògica categòrica que estudiarem més endavant.
\end{observacio}

\begin{resumcap}
\apartat{Idees centrals}
\begin{enumerate}
  \item Les dades universals es caracteritzen per l'existència i la unicitat de morfismes; sovint es poden veure com a objectes inicials o terminals en una categoria coma adequada.
  \item Una representació d'un functor és un objecte representant juntament amb un element universal; equivalentment, és un isomorfisme natural amb un homfunctor.
  \item El lema de Yoneda dona una bijecció natural
  \[
  \Nat\bigl(\Hom(-,A),P\bigr)\;\cong\;P(A),
  \]
  determinada per l'avaluació a $\id_A$.
  \item La immersió de Yoneda és plena i fidel; per tant, els morfismes entre objectes es recuperen exactament com a transformacions naturals entre els seus prefeixos representables.
  \item En particular, un objecte queda determinat, llevat d'isomorfisme, pel sistema de morfismes que hi arriben; la representació completa és única llevat d'un únic isomorfisme compatible amb la dada universal.
\end{enumerate}
\apartat{Recordatori operatiu} Les dues direccions de la bijecció de Yoneda són
\[
\eta\longmapsto\eta_A(\id_A),
\qquad\qquad
x\longmapsto\bigl(f\mapsto P(f)(x)\bigr).
\]
\apartat{Errors típics} Dir \enquote{únic llevat d'isomorfisme únic} de l'objecte nu quan això només val per a la representació amb la dada universal; usar la propietat universal del producte abans d'haver-la definit; i confondre el functor representat amb l'objecte representant.
\apartat{Lectura recomanada} \cite[§2.2 i §§8.1--8.4]{awodey} per als objectes inicials i terminals i per a la immersió, el lema i les aplicacions de Yoneda; \cite[§§4.1--4.3]{leinster}; \cite[cap.~2]{riehl}, d'arquitectura molt propera a la d'aquest capítol; \cite[II.6 i III.1--III.2]{maclane} per a les categories coma, les fletxes universals i el lema de Yoneda; i \cite[cap.~2]{perrone-notes}, per a una segona explicació del lema de Yoneda, amb molts casos particulars i exemples de matemàtica elemental.
\end{resumcap}

\chapter[Límits i colímits]{Límits i colímits}
\label{cap:limits}

\begin{objectius}
\apartat{Prerequisits} Categories oposades i dualitat, i categories sobre un objecte (cap.~\ref{cap:categories}); functors, functors constants i el functor diagonal (cap.~\ref{cap:functors}); transformacions naturals i categories de functors (cap.~\ref{cap:transformacions}); categories coma, propietats universals, representabilitat i lema de Yoneda, especialment la unicitat llevat d'isomorfisme (cap.~\ref{cap:yoneda}).
\apartat{Objectius} En acabar el capítol, el lector hauria de poder:
\begin{itemize}
  \item interpretar un diagrama com un functor i un con com una transformació natural;
  \item definir límit i colímit com a cons i cocons universals;
  \item reconèixer productes, equalitzadors, pullbacks i l'objecte terminal com a límits, i els seus duals com a colímits;
  \item calcular aquests casos bàsics a $\cat{Set}$ i interpretar-los en preordres;
  \item fer servir els pullbacks per descriure la reindexació;
  \item enunciar la caracterització dels límits finits mitjançant objecte terminal, productes binaris i equalitzadors, o bé objecte terminal i pullbacks;
  \item descriure els límits i colímits petits a $\cat{Set}$ i explicar com es calculen punt a punt els límits i colímits en categories de functors;
  \item distingir preservació i reflexió de límits;
  \item explicar per què els representables covariants preserven límits.
\end{itemize}
\end{objectius}

\section{Cap a la noció de límit}

El capítol anterior ha caracteritzat diversos objectes per propietats universals: l'objecte terminal, l'inicial, i de passada el producte i el coproducte. Tots responen a un mateix patró: un objecte que es relaciona de manera òptima amb tots els altres. I tots són únics llevat d'isomorfisme. En aquest capítol unifiquem aquest patró en la noció general de \textbf{límit} d'un diagrama, i la seva dual, el \textbf{colímit}.

La idea és senzilla. Un diagrama és una col·lecció d'objectes i morfismes amb una forma determinada. Un \emph{con} sobre el diagrama és un objecte que projecta cap a tots els seus vèrtexs de manera compatible. El \emph{límit} és el con òptim: aquell pel qual tots els altres factoritzen de manera única. Segons la forma del diagrama, aquesta recepta produeix productes, equalitzadors, pullbacks, objectes terminals i molt més, tots com a casos particulars d'una sola definició. La representabilitat i el lema de Yoneda en seran l'eina conceptual: en el llenguatge del capítol anterior, un límit és un objecte que representa el prefeix que envia cada objecte $X$ al conjunt dels cons de vèrtex $X$ sobre el diagrama.

\section{Diagrames i cons}

Formalitzem la idea de \enquote{diagrama d'una forma donada} amb el llenguatge dels functors, que ja tenim a mà.

\begin{definicio}[Diagrama]\label{def:diagrama}\index{diagrama!de forma $\cat{I}$}\index{categoria!índex}
Sigui $\cat{I}$ una categoria petita, anomenada \textbf{categoria índex}. Un \textbf{diagrama de forma $\cat{I}$} en una categoria $\cat{C}$ és un functor $D\colon\cat{I}\to\cat{C}$. Escrivim $D_i$ per a l'objecte $D(i)$ i, per a un morfisme $u\colon i\to j$ de $\cat{I}$, $D_u\colon D_i\to D_j$ per a la fletxa corresponent.
\end{definicio}

\begin{observacio}
La categoria índex $\cat{I}$ fixa la \emph{forma} del diagrama, no el seu contingut. Amb $\cat{I}$ discreta de dos objectes obtenim una parella d'objectes sense fletxes; amb $\cat{I}$ la categoria de dues fletxes paral·leles, un parell de morfismes amb el mateix domini i codomini; amb $\cat{I}$ buida, el diagrama buit. Cada forma donarà lloc a un tipus de límit.
\end{observacio}

\begin{definicio}[Con]\index{con}
Sigui $D\colon\cat{I}\to\cat{C}$ un diagrama. Un \textbf{con} sobre $D$ amb \textbf{vèrtex} $A$ és una família de morfismes
\[
\alpha_i\colon A\to D_i,\qquad i\in\cat{I},
\]
tal que, per a cada morfisme $u\colon i\to j$ de $\cat{I}$, el triangle commuta:
\[
\begin{tikzpicture}[baseline=(m.center)]
  \matrix (m) [matrix of math nodes, row sep=2.4em, column sep=3.0em] {
     & A & \\
    D_i & & D_j \\
  };
  \draw[-{Stealth[scale=1.1]}] (m-1-2) -- node[above left] {$\alpha_i$} (m-2-1);
  \draw[-{Stealth[scale=1.1]}] (m-1-2) -- node[above right] {$\alpha_j$} (m-2-3);
  \draw[-{Stealth[scale=1.1]}] (m-2-1) -- node[below] {$D_u$} (m-2-3);
\end{tikzpicture}
\qquad
D_u\comp\alpha_i=\alpha_j.
\]
\end{definicio}

\begin{observacio}\label{obs:con-transformacio}
Un con és exactament una transformació natural $\alpha\colon\Delta_A\Rightarrow D$, on $\Delta_A\colon\cat{I}\to\cat{C}$ és el functor constant amb valor $A$ (exemple~\ref{ex:functor-constant}). En efecte, els components $\alpha_i\colon A\to D_i$ i la condició de naturalitat $D_u\comp\alpha_i=\alpha_j$ són precisament la definició de con. Aquest és el primer lloc on la maquinària de les transformacions naturals rendeix directament.
\end{observacio}

\begin{definicio}[Morfisme de cons]
Un \textbf{morfisme de cons} de $(A,\alpha)$ a $(B,\beta)$, tots dos sobre $D$, és un morfisme $f\colon A\to B$ tal que $\beta_i\comp f=\alpha_i$ per a tot $i$. Els cons sobre $D$ i els seus morfismes formen una categoria, que denotem $\mathrm{Con}(D)$.
\end{definicio}

\begin{observacio}[Els cons són objectes d'una categoria coma]\label{obs:cons-coma}\index{con}\index{categoria!coma}\index{functor!diagonal}
Els functors constants $\Delta_A$ s'apleguen en un sol functor, el \textbf{functor diagonal}
\[
\Delta\colon\cat{C}\longrightarrow\cat{C}^{\cat{I}},
\]
que envia cada objecte $A$ al functor constant $\Delta_A$ i cada morfisme $f\colon A\to B$ a la transformació natural $\Delta_f\colon\Delta_A\Rightarrow\Delta_B$ de components $(\Delta_f)_i=f$; la naturalitat és trivial, perquè $\Delta_A$ i $\Delta_B$ envien tots els morfismes de $\cat{I}$ a identitats, i la functorialitat es comprova component a component. Amb aquest functor, la categoria de cons és una categoria coma (secció~\ref{sec:coma}):
\[
\mathrm{Con}(D)=(\Delta\downarrow D).
\]
En efecte, prenent $S=\Delta$ i $T=D\colon\cat{1}\to\cat{C}^{\cat{I}}$, un objecte de $(\Delta\downarrow D)$ és una terna $(A,\ast,\alpha)$ amb $A$ objecte de $\cat{C}$ i $\alpha\colon\Delta_A\Rightarrow D$ un morfisme de $\cat{C}^{\cat{I}}$, que identifiquem amb la parella $(A,\alpha)$, com a l'exemple~\ref{ex:slice-coma}; és a dir, un con sobre $D$ amb vèrtex $A$ (observació~\ref{obs:con-transformacio}). Un morfisme $(A,\alpha)\to(B,\beta)$ és una parella $(f,\id_\ast)$ amb $f\colon A\to B$ tal que el quadrat de la definició de categoria coma commuta, és a dir, $\beta\comp\Delta_f=\alpha$; component a component, $\beta_i\comp f=\alpha_i$ per a tot $i$, que és exactament la condició de morfisme de cons.

Dualment, un cocon sobre $D$ amb vèrtex $A$ (definició~\ref{def:cocon}) és una transformació natural $D\Rightarrow\Delta_A$, i els cocons formen la categoria coma $(D\downarrow\Delta)$. Així, el límit és un objecte terminal de $(\Delta\downarrow D)$ i el colímit un objecte inicial de $(D\downarrow\Delta)$: és l'eslògan de l'exemple~\ref{ex:universal-coma}, \emph{ser universal és ser inicial o terminal en una categoria coma adient}. La categoria coma dona, doncs, el llenguatge comú per a les categories sobre i sota un objecte, els objectes universals del capítol~\ref{cap:yoneda} i els límits i colímits.
\end{observacio}

\section{Límits}

\begin{definicio}[Límit]\index{límit}
\label{def:limit}
Un \textbf{límit} del diagrama $D\colon\cat{I}\to\cat{C}$ és un objecte terminal de la categoria de cons $\mathrm{Con}(D)$. Explícitament, és un con $(L,\lambda)$ tal que, per a tot con $(A,\alpha)$ sobre $D$, hi ha un \emph{únic} morfisme $f\colon A\to L$ amb
\[
\lambda_i\comp f=\alpha_i\qquad\text{per a tot }i.
\]
Escrivim $L=\varprojlim D$ o $L=\lim D$.
\end{definicio}

\begin{observacio}
Un límit és, doncs, el con \enquote{universal}: tots els altres cons factoritzen a través d'ell de manera única. Com que és un objecte terminal de $\mathrm{Con}(D)$, el con límit és \textbf{únic llevat d'un únic isomorfisme} de cons, és a dir, compatible amb les projeccions, per la unicitat dels objectes terminals (proposició~\ref{prop:unicitat-terminal}); el vèrtex sol és únic llevat d'isomorfisme. Direm que una categoria \textbf{té límits de forma $\cat{I}$} quan tot diagrama de forma $\cat{I}$ hi té límit.
\end{observacio}

La propietat universal del límit es pot formular, en el llenguatge del capítol~\ref{cap:yoneda}, com una representabilitat. Suposem $\cat{I}$ petita i $\cat{C}$ localment petita, de manera que, per a cada objecte $A$, els cons sobre $D$ amb vèrtex $A$ formen un conjunt
\[
\mathrm{Con}(A,D)=\Nat(\Delta_A,D).
\]
Un morfisme $g\colon A'\to A$ transforma cada con $\alpha$ amb vèrtex $A$ en el con $\alpha\comp\Delta_g$ amb vèrtex $A'$, de components $\alpha_i\comp g$; obtenim així un prefeix
\[
\mathrm{Con}(-,D)\colon\cat{C}\op\longrightarrow\cat{Set}.
\]

\begin{proposicio}[Els límits com a representacions]\label{prop:limit-representable}\index{límit!com a representació}
Un con $(L,\lambda)$ sobre $D$ és un límit si i només si les aplicacions
\[
\Hom(A,L)\longrightarrow\mathrm{Con}(A,D),
\qquad
f\longmapsto\lambda\comp\Delta_f=(\lambda_i\comp f)_i,
\]
formen un isomorfisme natural $\Hom(-,L)\cong\mathrm{Con}(-,D)$. Dit d'una altra manera, un límit de $D$ és una representació $(L,\lambda)$ del prefeix $\mathrm{Con}(-,D)$, amb element universal $\lambda\in\mathrm{Con}(L,D)$.
\end{proposicio}

\begin{prova}
Les aplicacions són naturals en $A$: per a $g\colon A'\to A$, enviar primer $f$ a $(\lambda_i\comp f)_i$ i precompondre després amb $g$ dona $(\lambda_i\comp f\comp g)_i$, que és la imatge de $f\comp g$. Una transformació natural és un isomorfisme natural exactament quan totes les seves components són bijectives, i que l'aplicació corresponent a $A$ sigui bijectiva diu exactament que, per a tot con $\alpha$ amb vèrtex $A$, hi ha un únic $f\colon A\to L$ amb $\lambda_i\comp f=\alpha_i$ per a tot $i$. Això és la definició de límit. L'element universal és la imatge de $\id_L$, que és $\lambda$.
\end{prova}

Aquesta formulació explica d'un sol cop tres fets del capítol. Un límit és una propietat universal en el sentit del capítol~\ref{cap:yoneda}; és, per tant, únic llevat d'un únic isomorfisme compatible amb la dada universal, que aquí és el con $\lambda$; i la mateixa idea mostrarà que els representables preserven límits (proposició~\ref{prop:representables-limits}).

Recorrem ara les formes particulars, cadascuna amb la seva categoria índex.

\subsection{Producte}\index{producte}

Prenem $\cat{I}$ la categoria \textbf{discreta} amb dos objectes $1,2$ (sense fletxes no trivials). Un diagrama és simplement una parella d'objectes $A,B$. Un con amb vèrtex $X$ és una parella de morfismes $f\colon X\to A$, $g\colon X\to B$, sense cap condició (no hi ha fletxes a $\cat{I}$ que imposin res). El límit és el \textbf{producte}.

\begin{definicio}\index{producte}
\label{def:producte}
El \textbf{producte} de dos objectes $A,B$ és un objecte $A\times B$ amb dues projeccions $\pi_1\colon A\times B\to A$ i $\pi_2\colon A\times B\to B$ tals que, per a tot objecte $X$ i tota parella de morfismes $f\colon X\to A$, $g\colon X\to B$, hi ha un únic morfisme $\langle f,g\rangle\colon X\to A\times B$ amb
\[
\pi_1\comp\langle f,g\rangle=f,\qquad \pi_2\comp\langle f,g\rangle=g.
\]
\[
\begin{tikzpicture}[baseline=(m.center)]
  \matrix (m) [matrix of math nodes, column sep=2.8em, row sep=2.6em] {
      & X & \\
    A & A\times B & B \\
  };
  \draw[-{Stealth[scale=1.1]}] (m-1-2) -- node[above left] {$f$} (m-2-1);
  \draw[-{Stealth[scale=1.1]}] (m-1-2) -- node[above right] {$g$} (m-2-3);
  \draw[-{Stealth[scale=1.1]},dashed] (m-1-2) -- node[right] {$\langle f,g\rangle$} (m-2-2);
  \draw[-{Stealth[scale=1.1]}] (m-2-2) -- node[below] {$\pi_1$} (m-2-1);
  \draw[-{Stealth[scale=1.1]}] (m-2-2) -- node[below] {$\pi_2$} (m-2-3);
\end{tikzpicture}
\]
\end{definicio}

\begin{exemple}
A $\cat{Set}$, el producte és el producte cartesià $A\times B=\{(a,b)\mid a\in A,\ b\in B\}$ amb les projeccions habituals. A $\cat{Grp}$ és el producte directe de grups; a $\cat{Top}$, l'espai producte amb la topologia producte; en un preordre, l'ínfim dels dos elements, quan existeix. En cada cas, la propietat universal recupera la construcció coneguda.
\end{exemple}

\subsection{Equalitzador}\index{equalitzador}

Prenem $\cat{I}$ la categoria amb dues fletxes paral·leles $\bullet\rightrightarrows\bullet$. Un diagrama és un parell de morfismes $f,g\colon A\to B$. Un con amb vèrtex $X$ és un morfisme $e\colon X\to A$ amb $f\comp e=g\comp e$ (la condició prové de les dues fletxes). El límit és l'\textbf{equalitzador}.

\begin{definicio}\index{equalitzador}
\label{def:equalitzador}
L'\textbf{equalitzador} de $f,g\colon A\to B$ és un objecte $E$ amb un morfisme $e\colon E\to A$ tal que $f\comp e=g\comp e$ i que és universal amb aquesta propietat: per a tot $h\colon X\to A$ amb $f\comp h=g\comp h$, hi ha un únic $k\colon X\to E$ amb $e\comp k=h$.
\[
\begin{tikzpicture}[baseline=(m.center)]
  \matrix (m) [matrix of math nodes, column sep=3em, row sep=2.6em] {
    E & A & B \\
    X & & \\
  };
  \draw[-{Stealth[scale=1.1]}] (m-1-1) -- node[above] {$e$} (m-1-2);
  \draw[-{Stealth[scale=1.1]}] ([yshift=2pt]m-1-2.east) -- node[above] {$f$} ([yshift=2pt]m-1-3.west);
  \draw[-{Stealth[scale=1.1]}] ([yshift=-2pt]m-1-2.east) -- node[below] {$g$} ([yshift=-2pt]m-1-3.west);
  \draw[-{Stealth[scale=1.1]}] (m-2-1) -- node[left] {$h$} (m-1-2);
  \draw[-{Stealth[scale=1.1]},dashed] (m-2-1) -- node[left] {$k$} (m-1-1);
\end{tikzpicture}
\]
\end{definicio}

\begin{exemple}
A $\cat{Set}$, l'equalitzador de $f,g\colon A\to B$ és el subconjunt $E=\{a\in A\mid f(a)=g(a)\}$ amb la injecció $E\hookrightarrow A$. En general, els equalitzadors representen els subobjectes definits per una equació. Tot equalitzador és un monomorfisme.
\end{exemple}

\subsection{Pullback}\index{pullback}

Prenem $\cat{I}$ la categoria amb forma de racó, $\bullet\to\bullet\leftarrow\bullet$. Un diagrama és un parell de morfismes amb codomini comú, $f\colon A\to C$ i $g\colon B\to C$. El límit és el \textbf{pullback}.

\begin{definicio}\index{pullback}
\label{def:pullback}
El \textbf{pullback} de $f\colon A\to C$ i $g\colon B\to C$ és un objecte $P$ amb morfismes $p_1\colon P\to A$ i $p_2\colon P\to B$ tals que $f\comp p_1=g\comp p_2$ i que és universal: per a tot $X$ amb $r_1\colon X\to A$, $r_2\colon X\to B$ i $f\comp r_1=g\comp r_2$, hi ha un únic $h\colon X\to P$ amb $p_1\comp h=r_1$ i $p_2\comp h=r_2$.
\[
\begin{tikzpicture}[baseline=(m.center)]
  \matrix (m) [matrix of math nodes, column sep=3em, row sep=2.6em] {
    P & B \\
    A & C \\
  };
  \draw[-{Stealth[scale=1.1]}] (m-1-1) -- node[above] {$p_2$} (m-1-2);
  \draw[-{Stealth[scale=1.1]}] (m-1-1) -- node[left] {$p_1$} (m-2-1);
  \draw[-{Stealth[scale=1.1]}] (m-1-2) -- node[right] {$g$} (m-2-2);
  \draw[-{Stealth[scale=1.1]}] (m-2-1) -- node[below] {$f$} (m-2-2);
  \node at ([xshift=0.75em,yshift=-0.75em]m-1-1) {$\scriptstyle\lrcorner$};
\end{tikzpicture}
\]
\end{definicio}

\begin{exemple}
A $\cat{Set}$, el pullback és $P=\{(a,b)\in A\times B\mid f(a)=g(b)\}$ amb les projeccions. Quan $g\colon B\hookrightarrow C$ és la inclusió d'un subconjunt $B\subseteq C$, el pullback recupera l'antiimatge $f^{-1}(B)$, mitjançant la bijecció $(a,b)\mapsto a$. En aquest cas, doncs, el pullback categoritza l'operació d'antiimatge. En general, el pullback també s'anomena \emph{producte fibrat} i s'escriu $A\times_C B$.
\end{exemple}

\begin{observacio}[Reindexació per pullback]\index{reindexació}\index{functor!de reindexació}
\label{obs:reindexacio}
Ara que disposem de pullbacks podem definir la reindexació entre categories sobre un objecte. Suposem que $\cat{C}$ té tots els pullbacks i sigui $f\colon A\to B$ un morfisme. Definim un functor
\[
f^{*}\colon\cat{C}/B\longrightarrow\cat{C}/A
\]
de la manera següent: a un objecte $y\colon Y\to B$ de $\cat{C}/B$ li assignem el pullback de $y$ al llarg de $f$, és a dir, l'objecte $f^{*}Y=A\times_B Y$ amb la seva projecció $p_1\colon A\times_B Y\to A$, que és un objecte de $\cat{C}/A$:
\[
\begin{tikzpicture}[baseline=(m.center)]
  \matrix (m) [matrix of math nodes, column sep=3em, row sep=2.6em] {
    A\times_B Y & Y \\
    A & B \\
  };
  \draw[-{Stealth[scale=1.1]}] (m-1-1) -- node[above] {$p_2$} (m-1-2);
  \draw[-{Stealth[scale=1.1]}] (m-1-1) -- node[left] {$p_1=f^{*}y$} (m-2-1);
  \draw[-{Stealth[scale=1.1]}] (m-1-2) -- node[right] {$y$} (m-2-2);
  \draw[-{Stealth[scale=1.1]}] (m-2-1) -- node[below] {$f$} (m-2-2);
  \node at ([xshift=0.9em,yshift=-0.8em]m-1-1) {$\scriptstyle\lrcorner$};
\end{tikzpicture}
\]
Sobre morfismes, la propietat universal del pullback assigna a cada triangle de $\cat{C}/B$ un únic triangle de $\cat{C}/A$. Un cop \textbf{triats} els pullbacks que defineixen $f^{*}$, obtenim un functor ordinari $f^{*}\colon\cat{C}/B\to\cat{C}/A$. Per a morfismes composables $A\xrightarrow{f}B\xrightarrow{g}C$, els functors $(g\comp f)^{*}$ i $f^{*}\comp g^{*}$ no són literalment iguals per a una tria arbitrària de pullbacks, però són canònicament isomorfs. Aquesta coherència s'organitza, en un llenguatge més avançat, mitjançant la noció de \emph{pseudofunctor}, que no necessitarem aquí. En passar als preordres de subobjectes (capítol~\ref{cap:subobjectes}), aquests isomorfismes queden identificats i la reindexació és estrictament functorial:
\[
(g\comp f)^{*}=f^{*}\comp g^{*},\qquad (\id_A)^{*}=\id_{\Sub(A)}.
\]
A $\cat{Set}$, si interpretem $y\colon Y\to B$ com la família de fibres $\{y^{-1}(b)\}_{b\in B}$, aleshores $f^{*}Y$ s'interpreta com la família reindexada $\{y^{-1}(f(a))\}_{a\in A}$: exactament la \textbf{substitució} al llarg de $f$. Aquesta és la lectura lògica que explotarem en tractar els subobjectes, on $f^{*}$ es restringeix a un morfisme de preordres $\Sub(B)\to\Sub(A)$ i s'interpreta com la substitució en un predicat.
\end{observacio}

\subsection{Objecte terminal com a límit}

Prenem $\cat{I}=\cat{0}$, la categoria buida. L'únic diagrama de forma $\cat{0}$ és el diagrama buit. Un con sobre el diagrama buit és, simplement, un objecte $X$ (sense cap component, perquè no hi ha vèrtexs). El límit ---el con terminal--- és, doncs, l'\textbf{objecte terminal} de $\cat{C}$. Així, l'objecte terminal és el límit del diagrama buit, i queda unificat amb els altres.

\section{Colímits}

Tota la teoria anterior es dualitza girant les fletxes. Un \textbf{colímit} d'un diagrama $D\colon\cat{I}\to\cat{C}$ és un límit del diagrama oposat $D\op\colon\cat{I}\op\to\cat{C}\op$; equivalentment, és un objecte inicial a la categoria de \textbf{cocons} sobre $D$.

\begin{definicio}[Cocon i colímit]\label{def:cocon}\index{cocon}\index{colímit}
Un \textbf{cocon} sobre $D$ amb vèrtex $A$ és una família $\alpha_i\colon D_i\to A$ tal que, per a cada $u\colon i\to j$, commuta el triangle dual
\[
\begin{tikzpicture}[baseline=(m.center)]
  \matrix (m) [matrix of math nodes, row sep=2.4em, column sep=3.0em] {
    D_i & & D_j \\
     & A & \\
  };
  \draw[-{Stealth[scale=1.1]}] (m-1-1) -- node[below left] {$\alpha_i$} (m-2-2);
  \draw[-{Stealth[scale=1.1]}] (m-1-3) -- node[below right] {$\alpha_j$} (m-2-2);
  \draw[-{Stealth[scale=1.1]}] (m-1-1) -- node[above] {$D_u$} (m-1-3);
\end{tikzpicture}
\qquad
\alpha_j\comp D_u=\alpha_i.
\]
Un \textbf{colímit} és un cocon $(L,\lambda)$ universal: per a tot cocon $(A,\alpha)$ hi ha un únic $f\colon L\to A$ amb $f\comp\lambda_i=\alpha_i$. Escrivim $L=\colimop D$ o $L=\varinjlim D$.
\end{definicio}

Dualitzant cada límit obtenim el colímit corresponent:

\begin{itemize}
\item El \textbf{coproducte} $A+B$ (dual del producte) té injeccions $\iota_1\colon A\to A+B$, $\iota_2\colon B\to A+B$ universals. A $\cat{Set}$ és la unió disjunta; a $\cat{Grp}$, el producte lliure; en un preordre, el suprem.
\item El \textbf{coequalitzador} de $f,g\colon A\to B$ (dual de l'equalitzador) és un $q\colon B\to Q$ universal amb $q\comp f=q\comp g$. A $\cat{Set}$ és el quocient de $B$ per la relació d'equivalència generada per $f(a)\sim g(a)$. Tot coequalitzador és un epimorfisme.
\item El \textbf{pushout} de $f\colon C\to A$ i $g\colon C\to B$ (dual del pullback) és un $Q$ universal amb el quadrat commutatiu. A $\cat{Set}$ és la unió de $A$ i $B$ enganxada al llarg de $C$, i s'escriu $A+_C B$.
\item L'\textbf{objecte inicial} (dual del terminal) és el colímit del diagrama buit.
\end{itemize}

\begin{observacio}
No cal demostrar res de nou per als colímits: cada enunciat sobre límits a $\cat{C}$ es tradueix en un enunciat sobre colímits a $\cat{C}\op$, i viceversa. Aquesta és la dualitat en acció, aplicada de manera sistemàtica. En endavant, quan enunciem un resultat per a límits, el dual per a colímits s'entén sobreentès.
\end{observacio}

\section{Límits finits}

\begin{definicio}\index{límit!finit}\index{categoria!finitament completa}
Un límit es diu \textbf{finit} quan la seva categoria índex $\cat{I}$ és finita, és a dir, quan té un nombre finit d'objectes \emph{i} un nombre finit de morfismes. (No n'hi ha prou que els objectes siguin finits: una categoria d'un sol objecte pot tenir infinits endomorfismes.) Una categoria és \textbf{finitament completa} (o \textbf{lex}) si té tots els límits finits; és \textbf{completa} si té tots els límits de forma petita.
\end{definicio}

El resultat clau redueix tots els límits finits a dos casos bàsics: productes finits i equalitzadors. Això és molt útil, perquè per comprovar que una categoria és finitament completa n'hi ha prou de verificar dues coses.

\begin{teorema}\index{límit!finit!construcció}\label{teo:limits-finits}
Per a una categoria $\cat{C}$, són equivalents:
\begin{enumerate}
\item $\cat{C}$ té tots els límits finits;
\item $\cat{C}$ té objecte terminal, productes binaris i equalitzadors;
\item $\cat{C}$ té objecte terminal i pullbacks.
\end{enumerate}
\end{teorema}

\begin{prova}
\emph{(1) implica (2) i (3).} L'objecte terminal, el producte binari, l'equalitzador i el pullback són límits de diagrames de formes finites: la categoria buida, la discreta de dos objectes, $\bullet\rightrightarrows\bullet$ i $\bullet\to\bullet\leftarrow\bullet$.

\emph{(2) implica (1).} Primer, $\cat{C}$ té productes de qualsevol família finita. El de la família buida és l'objecte terminal, i per inducció, si $P_n$ és un producte de $A_1,\dots,A_n$ amb projeccions $\pi_k$, aleshores $P_n\times A_{n+1}$, amb projeccions $\pi_k\comp p_1$ ($k\le n$) i $p_2$, és un producte de $A_1,\dots,A_{n+1}$. En efecte, donades $f_k\colon X\to A_k$, hi ha una única $f\colon X\to P_n$ amb $\pi_k\comp f=f_k$ per a $k\le n$, i una única $X\to P_n\times A_{n+1}$ amb components $f$ i $f_{n+1}$; tota fletxa amb les mateixes composicions amb les $n+1$ projeccions té necessàriament aquestes dues components.

Sigui ara $D\colon\cat{I}\to\cat{C}$ un diagrama finit, i escrivim $\cat{I}_0$ i $\cat{I}_1$ per als conjunts d'objectes i de morfismes de $\cat{I}$. Considerem els productes finits
\[
P=\prod_{i\in\cat{I}_0}D_i
\qquad\text{i}\qquad
Q=\prod_{u\in\cat{I}_1}D_{\mathrm{cod}(u)},
\]
amb projeccions $\pi_i\colon P\to D_i$ i $q_u\colon Q\to D_{\mathrm{cod}(u)}$, i les dues fletxes $s,t\colon P\to Q$ determinades, per a cada $u\colon i\to j$, per
\[
q_u\comp s=\pi_j,
\qquad\qquad
q_u\comp t=D_u\comp\pi_i.
\]
Sigui $e\colon L\to P$ un equalitzador de $s$ i $t$, i posem $\lambda_i=\pi_i\comp e\colon L\to D_i$. Vegem que $(L,\lambda)$ és un límit de $D$.

\emph{És un con.} Per a $u\colon i\to j$,
\[
D_u\comp\lambda_i=D_u\comp\pi_i\comp e=q_u\comp t\comp e=q_u\comp s\comp e=\pi_j\comp e=\lambda_j.
\]

\emph{És universal.} Sigui $(A,\alpha)$ un con sobre $D$, i sigui $a\colon A\to P$ l'única fletxa amb $\pi_i\comp a=\alpha_i$ per a tot $i$. Per a cada $u\colon i\to j$,
\[
q_u\comp s\comp a=\pi_j\comp a=\alpha_j,
\qquad
q_u\comp t\comp a=D_u\comp\pi_i\comp a=D_u\comp\alpha_i=\alpha_j,
\]
de manera que $s\comp a$ i $t\comp a$ tenen les mateixes components i, per la unicitat del producte $Q$, $s\comp a=t\comp a$. Per la propietat de l'equalitzador hi ha una única $h\colon A\to L$ amb $e\comp h=a$, i aleshores $\lambda_i\comp h=\pi_i\comp e\comp h=\alpha_i$. Si $h'\colon A\to L$ també compleix $\lambda_i\comp h'=\alpha_i$ per a tot $i$, aleshores $\pi_i\comp(e\comp h')=\alpha_i$, d'on $e\comp h'=a$ per la unicitat del producte $P$, i $h'=h$ per la de l'equalitzador. El diagrama buit hi queda inclòs: aleshores $P$ i $Q$ són productes buits, és a dir, objectes terminals, i $L$ és terminal.

\emph{(3) implica (2).} Sigui $1$ l'objecte terminal i, per a cada objecte $A$, $!_A\colon A\to 1$ l'única fletxa. \emph{Productes binaris}: un pullback de $!_A$ i $!_B$ és un producte de $A$ i $B$. En efecte, per a qualsevol parell $f\colon X\to A$, $g\colon X\to B$, la condició de con $!_A\comp f=\,!_B\comp g$ es compleix automàticament, perquè totes dues fletxes van a $1$; per tant, la propietat universal del pullback diu exactament la del producte.

\emph{Equalitzadors}: siguin $f,g\colon A\to B$. Ja sabem que hi ha productes binaris; formem el pullback de
\[
\langle\id_A,f\rangle,\ \langle\id_A,g\rangle\colon A\to A\times B,
\]
amb projeccions $p_1,p_2\colon P\to A$, de manera que $\langle\id_A,f\rangle\comp p_1=\langle\id_A,g\rangle\comp p_2$. Component amb la primera projecció del producte, $p_1=p_2$; anomenem $e$ aquesta fletxa. Component amb la segona, $f\comp e=g\comp e$. Si $h\colon X\to A$ compleix $f\comp h=g\comp h$, aleshores
\[
\langle\id_A,f\rangle\comp h=\langle h,f\comp h\rangle=\langle h,g\comp h\rangle=\langle\id_A,g\rangle\comp h,
\]
de manera que $(h,h)$ és un con sobre el pullback, i hi ha una única $k\colon X\to P$ amb $p_1\comp k=h=p_2\comp k$, és a dir, amb $e\comp k=h$. Si $k'$ també compleix $e\comp k'=h$, aleshores $p_1\comp k'=p_2\comp k'=h$ i $k'=k$. Per tant $e$ és un equalitzador de $f$ i $g$.

Amb les implicacions (3)$\Rightarrow$(2)$\Rightarrow$(1)$\Rightarrow$(3), les tres condicions són equivalents.
\end{prova}

\begin{observacio}
Dualment, una categoria és \textbf{finitament cocompleta} si té objecte inicial, coproductes binaris i coequalitzadors, o equivalentment objecte inicial i pushouts. Per a la lògica categòrica, els límits finits ---i molt especialment els \emph{pullbacks}--- són centrals: serveixen per interpretar la substitució i les operacions sobre subobjectes que veurem més endavant.
\end{observacio}

\begin{exemple}[Límits en un preordre i en la deduïbilitat]\label{ex:limits-preordre}\index{límit!en un preordre}
Sigui $P$ un preordre vist com a categoria i $D\colon\cat{I}\to P$ un diagrama. Com que entre dos objectes hi ha com a màxim una fletxa, un con amb vèrtex $x$ diu simplement que $x\le D_i$ per a tot $i$, i totes les condicions de commutativitat es compleixen automàticament. Per tant, quan existeixen,
\[
\varprojlim D=\bigwedge_{i\in\cat{I}_0}D_i,
\qquad\qquad
\varinjlim D=\bigvee_{i\in\cat{I}_0}D_i:
\]
el límit és l'ínfim dels objectes del diagrama, i el colímit, el suprem; les fletxes del diagrama no hi intervenen. En particular, l'equalitzador de dues fletxes paral·leles $x\rightrightarrows y$ és $x$ mateix (les dues fletxes coincideixen), i el pullback de $a\to c\leftarrow b$ és l'ínfim $a\wedge b$, que no depèn de $c$. Així, un preordre té tots els límits finits si i només si té element màxim (l'ínfim de la família buida, que és l'objecte terminal) i ínfims binaris.

A la categoria prima $\cat{Prov}_T$ d'un sistema deductiu $T$, això diu que, quan el sistema disposa de $\top$ i de conjunció, $\top$ és un objecte terminal ($p\vdash_T\top$ per a tota fórmula $p$) i $p\land q$ és un producte de $p$ i $q$ (exemple~\ref{ex:rep-producte}). Aquestes dues operacions donen, doncs, tots els límits finits en el context deductiu: la part \enquote{conjuntiva} de la lògica és exactament la dels límits finits.
\end{exemple}

\section{Límits a \texorpdfstring{$\cat{Set}$}{Set} i en categories de functors}

\begin{exemple}[Límits a $\cat{Set}$]\label{ex:limits-set}\index{límit!a Set@a $\cat{Set}$}
Sigui $D\colon\cat{I}\to\cat{Set}$ un diagrama amb $\cat{I}$ petita. El conjunt de les famílies compatibles
\[
L=\Bigl\{(x_i)_{i}\in\textstyle\prod_{i\in\cat{I}_0}D_i\ \Bigm|\ D_u(x_i)=x_j\ \text{per a tot } u\colon i\to j\Bigr\},
\]
amb les projeccions $\lambda_i\colon L\to D_i$, és un límit de $D$. En efecte, un con $(A,\alpha)$ envia cada $a\in A$ a la família $(\alpha_i(a))_i$, que és compatible per la condició de con; l'aplicació $a\mapsto(\alpha_i(a))_i$ és l'única $f\colon A\to L$ amb $\lambda_i\comp f=\alpha_i$. En particular, $\cat{Set}$ té tots els límits petits. Dualment, el colímit és el quocient de la unió disjunta $\coprod_i D_i$, els elements de la qual són parelles $(i,x)$ amb $x\in D_i$, per la relació d'equivalència generada per
\[
(i,x)\sim\bigl(j,D_u(x)\bigr)\qquad(u\colon i\to j);
\]
per tant, $\cat{Set}$ té també tots els colímits petits. A $\cat{Set}$, doncs, els límits imposen compatibilitats i els colímits enganxen dades identificant-les.
\end{exemple}

Els límits en categories de functors es calculen a partir dels límits de la categoria d'arribada, objecte per objecte.

\begin{proposicio}[Límits punt a punt]\label{prop:limits-punt-a-punt}\index{límit!en categories de functors}\index{categoria!de functors!límits}
Siguin $\cat{A}$ i $\cat{I}$ categories petites. Si $\cat{C}$ té límits de forma $\cat{I}$, aleshores la categoria de functors $\cat{C}^{\cat{A}}$ també en té, i es calculen \emph{punt a punt}: per a un diagrama $F\colon\cat{I}\to\cat{C}^{\cat{A}}$,
\[
\Bigl(\varprojlim F\Bigr)(a)\;\cong\;\varprojlim_{i\in\cat{I}}F(i)(a)
\qquad(a\in\cat{A}).
\]
Dualment, si $\cat{C}$ té colímits de forma $\cat{I}$, els colímits a $\cat{C}^{\cat{A}}$ es calculen punt a punt. Equivalentment, per a cada objecte $a$ de $\cat{A}$, el functor d'avaluació
\[
\mathrm{ev}_a\colon\cat{C}^{\cat{A}}\longrightarrow\cat{C},\qquad F\longmapsto F(a),
\]
preserva aquests límits i colímits.
\end{proposicio}
Vegeu també \cite[§3.3]{riehl} i, en el context de les categories de diagrames, \cite[§§8.5--8.6]{awodey}.

\begin{prova}
Recordem que un objecte de $\cat{C}^{\cat{A}}$ és un functor $\cat{A}\to\cat{C}$, i que un morfisme és una transformació natural, amb la composició calculada component a component. Així, el diagrama $F$ assigna a cada objecte $i$ de $\cat{I}$ un functor $F(i)\colon\cat{A}\to\cat{C}$, i a cada $u\colon i\to j$ una transformació natural $F(u)\colon F(i)\Rightarrow F(j)$, de components $F(u)_a$. Per a cada objecte $a$ de $\cat{A}$, el diagrama avaluat
\[
F_a=\mathrm{ev}_a\comp F\colon\cat{I}\to\cat{C},\qquad i\mapsto F(i)(a),\quad u\mapsto F(u)_a,
\]
és un functor, perquè la composició i les identitats de $\cat{C}^{\cat{A}}$ es calculen component a component.

\emph{Pas 1: el límit sobre objectes.} Per a cada objecte $a$ de $\cat{A}$, triem un límit $\bigl(L(a),(\lambda_i(a))_i\bigr)$ de $F_a$; per tant, $F(u)_a\comp\lambda_i(a)=\lambda_j(a)$ per a cada $u\colon i\to j$.

\emph{Pas 2: el límit sobre morfismes.} Sigui $v\colon a\to a'$ un morfisme de $\cat{A}$. La família $\bigl(F(i)(v)\comp\lambda_i(a)\bigr)_i$, de fletxes $L(a)\to F(i)(a')$, és un con sobre $F_{a'}$: per a $u\colon i\to j$, la naturalitat de $F(u)$ i la condició de con de $\lambda(a)$ donen
\[
F(u)_{a'}\comp F(i)(v)\comp\lambda_i(a)=F(j)(v)\comp F(u)_a\comp\lambda_i(a)=F(j)(v)\comp\lambda_j(a).
\]
Per la universalitat de $L(a')$, hi ha una única $L(v)\colon L(a)\to L(a')$ tal que
\[
\lambda_i(a')\comp L(v)=F(i)(v)\comp\lambda_i(a)\qquad\text{per a tot } i.
\tag{$\dagger$}
\]

\emph{Pas 3: $L$ és un functor.} Fem servir que dues fletxes cap a un límit que tenen les mateixes composicions amb totes les projeccions són iguals, per la unicitat de la factorització. Per a $v=\id_a$, tant $L(\id_a)$ com $\id_{L(a)}$ compleixen $(\dagger)$, perquè $F(i)(\id_a)=\id$; per tant són iguals. Per a $v\colon a\to a'$ i $w\colon a'\to a''$, aplicant $(\dagger)$ dues vegades i la functorialitat de $F(i)$,
\[
\begin{aligned}
\lambda_i(a'')\comp L(w)\comp L(v)&=F(i)(w)\comp\lambda_i(a')\comp L(v)\\
&=F(i)(w)\comp F(i)(v)\comp\lambda_i(a)=F(i)(w\comp v)\comp\lambda_i(a),
\end{aligned}
\]
de manera que $L(w)\comp L(v)$ compleix la condició $(\dagger)$ que caracteritza $L(w\comp v)$, i $L(w)\comp L(v)=L(w\comp v)$.

\emph{Pas 4: un con a $\cat{C}^{\cat{A}}$.} La condició $(\dagger)$ és exactament el quadrat de naturalitat de la família $\lambda_i=(\lambda_i(a))_a$; per tant, cada $\lambda_i$ és una transformació natural $L\Rightarrow F(i)$. A més, $F(u)\comp\lambda_i=\lambda_j$ a $\cat{C}^{\cat{A}}$, perquè aquesta igualtat de transformacions naturals es comprova component a component i ho diu el pas~1. Així, $(L,\lambda)$ és un con sobre $F$.

\emph{Pas 5: universalitat.} Sigui $(G,\alpha)$ un con sobre $F$, amb $\alpha_i\colon G\Rightarrow F(i)$ i $F(u)\comp\alpha_i=\alpha_j$. Avaluant en $a$, la família $\bigl((\alpha_i)_a\bigr)_i$ és un con sobre $F_a$ amb vèrtex $G(a)$, de manera que hi ha una única $\theta_a\colon G(a)\to L(a)$ amb $\lambda_i(a)\comp\theta_a=(\alpha_i)_a$ per a tot $i$. La família $\theta$ és natural: per a $v\colon a\to a'$, per $(\dagger)$ i per la naturalitat de $\alpha_i$,
\[
\lambda_i(a')\comp L(v)\comp\theta_a=F(i)(v)\comp(\alpha_i)_a=(\alpha_i)_{a'}\comp G(v)=\lambda_i(a')\comp\theta_{a'}\comp G(v)
\]
per a tot $i$, i per tant $L(v)\comp\theta_a=\theta_{a'}\comp G(v)$. Així $\theta\colon G\Rightarrow L$ és un morfisme de $\cat{C}^{\cat{A}}$ amb $\lambda_i\comp\theta=\alpha_i$. És l'únic: si $\theta'$ també ho compleix, a cada $a$ es té $\lambda_i(a)\comp\theta'_a=(\alpha_i)_a$, i la unicitat de la factorització a través de $L(a)$ dona $\theta'_a=\theta_a$. Per tant $(L,\lambda)$ és un límit de $F$.

\emph{Pas 6: el càlcul punt a punt.} Per construcció, $\mathrm{ev}_a(L,\lambda)=\bigl(L(a),\lambda(a)\bigr)$ és un límit de $F_a$. Si $(L',\lambda')$ és un altre límit de $F$, hi ha un isomorfisme $\varphi\colon L'\Rightarrow L$ amb $\lambda_i\comp\varphi=\lambda'_i$, i la seva component $\varphi_a$ és un isomorfisme amb $\lambda_i(a)\comp\varphi_a=\lambda'_i(a)$. Un con isomorf a un con límit, per un isomorfisme compatible amb les projeccions, també és un límit, perquè les factoritzacions a través de l'un i de l'altre es corresponen component amb $\varphi_a$ o amb $\varphi_a^{-1}$. Per tant $\bigl(L'(a),\lambda'(a)\bigr)$ és un límit de $F_a$: $\mathrm{ev}_a$ preserva els límits de forma $\cat{I}$, i $\bigl(\varprojlim F\bigr)(a)\cong\varprojlim_i F(i)(a)$.

L'argument per als colímits és el mateix amb totes les fletxes invertides: cocons en lloc de cons, i la propietat universal del colímit en lloc de la del límit.
\end{prova}

\begin{exemple}[Un mateix producte a $\cat{Set}$, a $\cat{Graph}$ i a $\cat{Set}^{\cat{A}}$]\label{ex:producte-tres-contextos}\index{producte!a Set, Graph i categories de functors}\index{objecte terminal!a Graph}
Fixem un sol diagrama abstracte: la parella d'objectes $(X,T)$, sobre la categoria discreta de dos objectes, juntament amb el diagrama buit, que dona l'objecte terminal. Calculem-ne el límit en tres llocs.

\emph{A $\cat{Set}$.} El producte és $X\times T$, i l'objecte terminal, un conjunt d'un sol element $\{\ast\}$. Amb $T=\{\ast\}$, l'aplicació $(x,\ast)\mapsto x$ és una bijecció $X\times\{\ast\}\cong X$. A més, les fletxes $\{\ast\}\to X$ són exactament els elements de $X$.

\emph{A $\cat{Graph}$.} Recordem que $\cat{Graph}\cong\cat{Set}^{\cat{P}}$, on $\cat{P}$ té dos objectes $E,V$ i dues fletxes $s,t\colon E\to V$. Per la proposició anterior, el producte es calcula component a component,
\[
V_{G\times H}=V_G\times V_H,
\qquad
E_{G\times H}=E_G\times E_H,
\]
amb origen i destí també component a component: l'aresta $(e,e')$ va del vèrtex $(s_G(e),s_H(e'))$ al vèrtex $(t_G(e),t_H(e'))$. Les projeccions són els morfismes de grafs que retenen una de les dues components. Prenem ara com a $T$ el candidat ingenu a \enquote{punt}, el graf $\mathbf V$ d'un sol vèrtex i cap aresta. La fórmula dona
\[
V_{G\times\mathbf V}=V_G\times\{\ast\}\cong V_G,
\qquad
E_{G\times\mathbf V}=E_G\times\varnothing=\varnothing:
\]
el producte amb $\mathbf V$ \emph{esborra totes les arestes}. Per tant, $\mathbf V$ no és l'objecte terminal de $\cat{Graph}$; de fet, d'un graf amb alguna aresta no hi ha cap morfisme cap a $\mathbf V$, perquè l'aresta no té on anar. L'objecte terminal es calcula també punt a punt, com el límit del diagrama buit a cada objecte de $\cat{P}$: un vèrtex \emph{i} una aresta, és a dir, el graf $\mathbf L$ d'un sol \emph{bucle}. Amb $T=\mathbf L$, sí que $G\times\mathbf L\cong G$. I els morfismes $\mathbf L\to G$ ja no són els vèrtexs de $G$, sinó els seus bucles: un graf sense bucles no té cap \enquote{punt} en sentit categòric, encara que tingui molts vèrtexs (compareu-ho amb l'observació~\ref{obs:raonar-elements}).

\emph{A una categoria de functors $\cat{Set}^{\cat{A}}$.} La recepta és la mateixa per a qualsevol categoria petita $\cat{A}$: $(F\times G)(a)=F(a)\times G(a)$ i $(F\times G)(v)=F(v)\times G(v)$, i l'objecte terminal és el functor constant de valor $\{\ast\}$. Els dos casos anteriors en són casos particulars: $\cat{Set}\cong\cat{Set}^{\cat{1}}$, i $\cat{Graph}\cong\cat{Set}^{\cat{P}}$. Ara s'entén la sorpresa del bucle: el terminal de $\cat{Set}^{\cat{P}}$ ha de ser un punt \emph{a cada objecte} de $\cat{P}$, també a $E$, i això obliga a tenir una aresta.

La forma del diagrama i la recepta del càlcul no canvien d'un context a l'altre. El que canvia és quins objectes fan cada paper, i la intuició de $\cat{Set}$ només es pot traslladar si es fa objecte per objecte de la categoria índex.
\end{exemple}

Com que $\cat{Set}$ té tots els límits i colímits petits (exemple~\ref{ex:limits-set}), la proposició mostra que també els té la categoria de prefeixos $\cat{Set}^{\cat{C}\op}$ de qualsevol categoria petita $\cat{C}$, tal com s'anunciava al capítol~\ref{cap:yoneda}. Aquest fet serà essencial en els capítols de subobjectes i de topos.

\section{Preservació i reflexió de límits}

Quan apliquem un functor a un diagrama i al seu límit, la imatge és un con sobre el diagrama imatge, però no té per què ser-ne el límit. Els functors que respecten límits reben un nom.

\begin{definicio}\index{functor!preserva límits}\index{functor!reflecteix límits}
Sigui $F\colon\cat{C}\to\cat{D}$ un functor.
\begin{itemize}
\item $F$ \textbf{preserva} el límit $(L,\lambda)$ d'un diagrama $D\colon\cat{I}\to\cat{C}$ si $(F(L),F\lambda)$ és un límit de $F\comp D$ a $\cat{D}$.
\item $F$ \textbf{reflecteix} límits si, per a tot diagrama $D\colon\cat{I}\to\cat{C}$ i tot con $(L,\lambda)$ sobre $D$, quan $(F(L),F\lambda)$ és un límit de $F\comp D$, aleshores $(L,\lambda)$ ja era un límit de $D$.
\end{itemize}
Es diu que $F$ \textbf{preserva els límits finits} si preserva el límit de tot diagrama finit que en tingui, i anàlogament per a altres classes.
\end{definicio}

\begin{observacio}
La reflexió es formula sempre \emph{respecte d'un con i un diagrama concrets}: cal partir d'un con $(L,\lambda)$ sobre $D$ la imatge del qual és limitant. Un functor ple i fidel permet aleshores aixecar i detectar les factoritzacions úniques ---perquè estableix una bijecció sobre els homconjunts---, però això no vol dir que \enquote{creï} límits del no-res: si a $\cat{C}$ no hi ha cap con sobre $D$ que projecti sobre el límit imatge, no hi ha res a reflectir. La distinció entre \emph{reflectir} (comparar un con donat) i \emph{crear} (produir-ne un de nou) és important i convé tenir-la present sempre que un functor interactuï amb els límits. La noció formal de crear límits, útil sobretot per als functors oblit de categories algebraiques, no serà necessària aquí; n'hi ha prou de retenir que és més forta que la mera reflexió d'un con ja donat.
\end{observacio}

\begin{exemple}
El functor oblit $U\colon\cat{Grp}\to\cat{Set}$ preserva tots els límits petits. En efecte, si $D\colon\cat{I}\to\cat{Grp}$ és un diagrama petit, el límit dels conjunts subjacents és el conjunt de famílies compatibles $(x_i)_i$ (exemple~\ref{ex:limits-set}). Les operacions de grup, definides component a component, conserven aquestes compatibilitats, perquè cada $D_u$ és un homomorfisme; aquest conjunt adquireix així una estructura de grup per a la qual les projeccions són homomorfismes i que satisfà la propietat universal del límit a $\cat{Grp}$. En particular, això recupera els productes i els equalitzadors habituals. En canvi, $U$ \emph{no} preserva tots els colímits: ja falla per als coproductes, perquè el conjunt subjacent d'un coproducte de grups (el producte lliure) no és, en general, la unió disjunta dels conjunts subjacents. Aquesta asimetria és típica dels functors oblit i anticipa el teorema del capítol següent: els functors oblit solen ser adjunts per la dreta, i els adjunts per la dreta preserven tots els límits; no hi ha cap afirmació anàloga que els obligui a preservar tots els colímits.
\end{exemple}

\begin{proposicio}\label{prop:representables-limits}\index{functor!representable!preserva límits}
Per a tot objecte $A$ d'una categoria localment petita, el functor representable $\Hom(A,-)\colon\cat{C}\to\cat{Set}$ preserva tots els límits que existeixen a $\cat{C}$.
\end{proposicio}

\begin{prova}
Sigui $(L,\lambda)$ un límit d'un diagrama $D\colon\cat{I}\to\cat{C}$. Escrivim $H=\Hom(A,-)$, de manera que $H(X)=\Hom(A,X)$ i $H(v)=v_*\colon\varphi\mapsto v\comp\varphi$.

\emph{El con imatge.} La família $H\lambda=\bigl((\lambda_i)_*\bigr)_i$, amb $(\lambda_i)_*\colon\Hom(A,L)\to\Hom(A,D_i)$, és un con sobre $H\comp D$: per a $u\colon i\to j$, la functorialitat de $H$ dona $(D_u)_*\comp(\lambda_i)_*=(D_u\comp\lambda_i)_*=(\lambda_j)_*$.

\emph{Existència de la factorització.} Sigui $(S,\sigma)$ un con sobre $H\comp D$ a $\cat{Set}$. És a dir, $S$ és un conjunt i $\sigma_i\colon S\to\Hom(A,D_i)$ són aplicacions amb $(D_u)_*\comp\sigma_i=\sigma_j$. Avaluada en un element $s\in S$, aquesta condició diu $D_u\comp\sigma_i(s)=\sigma_j(s)$: la família $\bigl(\sigma_i(s)\bigr)_i$ és un con sobre $D$ amb vèrtex $A$. Per la propietat universal de $L$, hi ha una única fletxa $h(s)\colon A\to L$ amb $\lambda_i\comp h(s)=\sigma_i(s)$ per a tot $i$. Això defineix una aplicació $h\colon S\to\Hom(A,L)$ amb $(\lambda_i)_*\comp h=\sigma_i$ per a tot $i$.

\emph{Unicitat.} Si $h'\colon S\to\Hom(A,L)$ també compleix $(\lambda_i)_*\comp h'=\sigma_i$, aleshores, per a cada $s$, $\lambda_i\comp h'(s)=\sigma_i(s)$ per a tot $i$, i la unicitat de la factorització a través de $L$ dona $h'(s)=h(s)$. Per tant $h'=h$.

Així $\bigl(\Hom(A,L),H\lambda\bigr)$ és un límit de $H\comp D$, i $\Hom(A,-)$ preserva el límit $(L,\lambda)$.
\end{prova}

\begin{observacio}[La prova en una línia]
La demostració anterior és la lectura element a element de la cadena de bijeccions
\[
\Hom\bigl(A,\varprojlim D\bigr)\;\cong\;\mathrm{Con}(A,D)\;\cong\;\varprojlim\nolimits_i\Hom(A,D_i).
\]
La primera bijecció és la proposició~\ref{prop:limit-representable}. La segona és la descripció dels límits a $\cat{Set}$ com a famílies compatibles (exemple~\ref{ex:limits-set}): un element de $\varprojlim_i\Hom(A,D_i)$ és una família $(\alpha_i)_i$ amb $D_u\comp\alpha_i=\alpha_j$, és a dir, un con. La composta envia $f$ a $(\lambda_i\comp f)_i$, que és l'acció del con imatge.
\end{observacio}

\begin{observacio}
Aquest fet ---els representables covariants preserven límits--- és una peça clau que retrobarem. Dualment, els representables contravariants $\Hom(-,B)$ envien colímits a límits: com que $\Hom_{\cat{C}}(-,B)=\Hom_{\cat{C}\op}(B,-)$ i un colímit de $\cat{C}$ és un límit de $\cat{C}\op$, n'hi ha prou d'aplicar la proposició a $\cat{C}\op$. Per exemple, $\Hom(A+A',B)\cong\Hom(A,B)\times\Hom(A',B)$. Combinat amb el lema de Yoneda, permet raonar sobre límits component a component, i és la base tècnica de molts arguments de la teoria.
\end{observacio}

\begin{resumcap}
\apartat{Idees centrals}
\begin{enumerate}
  \item Un límit és el con universal sobre un diagrama; un colímit, el cocon universal, és a dir, el límit del diagrama corresponent a la categoria oposada.
  \item El límit d'un diagrama $D$ representa el prefeix de cons: $\Hom(-,\varprojlim D)\cong\mathrm{Con}(-,D)$. Per tant, el con límit és únic llevat d'un únic isomorfisme compatible amb les projeccions, i el vèrtex sol, llevat d'isomorfisme.
  \item Productes, equalitzadors, pullbacks i l'objecte terminal són límits; els seus duals són colímits. Tenir tots els límits finits equival a tenir objecte terminal, productes binaris i equalitzadors; equivalentment, objecte terminal i pullbacks.
  \item $\cat{Set}$ té tots els límits i colímits petits (famílies compatibles i quocients de sumes), i les categories de functors els calculen punt a punt quan existeixen a la categoria d'arribada.
  \item Els pullbacks indueixen la reindexació $f^{*}\colon\cat{C}/B\to\cat{C}/A$, que a $\cat{Set}$ reindexa fibres i anticipa la substitució de predicats.
  \item Els representables $\Hom(A,-)$ preserven límits.
\end{enumerate}
\apartat{Errors típics} Confondre el con (la família de projeccions) amb el seu vèrtex (l'objecte); creure que \enquote{tenir límits} depèn de la tria d'un límit concret, quan la propietat universal el fixa llevat d'isomorfisme; oblidar que un colímit no és un límit \emph{de la mateixa} categoria, sinó de l'oposada; suposar que tot functor preserva límits (només ho fan certs functors, com els representables covariants); i confondre \enquote{$F$ preserva un límit}, que parteix d'un límit de $\cat{C}$ i en demana la imatge, amb \enquote{$F$ reflecteix límits}, que parteix d'un con de $\cat{C}$ amb imatge limitant i en conclou que ja era un límit.
\apartat{Lectura recomanada} \cite[cap.~5]{awodey} per a l'exposició paral·lela a aquesta; \cite[III.3--III.4 i cap.~V]{maclane} per al tractament clàssic de límits com a cons universals; \cite[cap.~3]{riehl} per a exemples addicionals, en particular §3.3, sobre preservació, reflexió i creació de límits, i §3.4, sobre la naturalesa representable dels límits; \cite[cap.~9]{barrwells} per a la intuïció del límit com a subconjunt d'un producte tallat per equacions, i del colímit com a quocient d'una suma; \cite[cap.~3]{perrone-notes} per al prefeix de cons i el cas dels preordres; \cite[cap.~5]{leinster} per a una presentació breu, amb la interacció entre functors i límits; i \cite[part~IV, sessions~19--25]{lawvereschanuel}, per a objectes terminals, productes i sumes, amb els productes de grafs, en una presentació molt elemental.
\end{resumcap}

\chapter[Adjuncions]{Adjuncions}
\label{cap:adjuncions}

\begin{objectius}
\apartat{Prerequisits} Functors, preordres vistos com a categories i la categoria lliure d'un graf (cap.~\ref{cap:functors}); transformacions i isomorfismes naturals, i equivalències de categories (cap.~\ref{cap:transformacions}); representabilitat i propietats universals (cap.~\ref{cap:yoneda}); límits i colímits, i la reindexació per pullback (cap.~\ref{cap:limits}).
\apartat{Objectius} En acabar el capítol, el lector hauria de poder:
\begin{itemize}
  \item definir una adjunció mitjançant una bijecció natural d'homconjunts i mitjançant unitat i counitat;
  \item calcular transposats i reconèixer la propietat universal de la unitat;
  \item verificar les identitats triangulars;
  \item reconèixer adjuncions en preordres i en construccions lliure--oblit;
  \item interpretar $\colimop_{\cat{I}}\dashv\Delta\dashv\varprojlim_{\cat{I}}$ i la currificació com a adjuncions;
  \item interpretar a $\cat{Set}$ la cadena $\exists_f\dashv f^{*}\dashv\forall_f$ com a prototip de la quantificació;
  \item demostrar la unicitat dels adjunts llevat d'isomorfisme natural;
  \item aplicar que els adjunts per la dreta preserven límits i els adjunts per l'esquerra, colímits.
\end{itemize}
\end{objectius}

\section{Cap a la noció d'adjunció}

Al llarg del llibre hem trobat ja la parella formada pel functor categoria lliure $\Free\colon\cat{Graph}\to\cat{Cat}$ i el functor graf subjacent $U\colon\cat{Cat}\to\cat{Graph}$. Un altre patró clàssic, que veurem ara en una versió elemental, és el del monoide lliure i el functor oblit $U\colon\cat{Mon}\to\cat{Set}$. En tots dos casos, un functor afegeix estructura de la manera més lliure possible i l'altre l'oblida, i tots dos estan lligats per una correspondència precisa. Aquesta relació s'anomena \textbf{adjunció}, i és un dels conceptes centrals de la teoria de categories.

Les adjuncions explicaran també un patró que vam observar en estudiar els límits: els representables i molts functors oblit preserven límits de manera sistemàtica, mentre que moltes construccions lliures preserven colímits. El teorema d'aquest capítol explicarà aquests fenòmens quan els functors en qüestió siguin adjunts. I seran un pont directe cap a la lògica: els quantificadors existencial i universal es revelaran com a adjunts.

La idea es pot copsar primer en el cas més simple, el dels preordres, on una adjunció és una parella d'aplicacions monòtones lligades per una equivalència de desigualtats. Després la generalitzarem a categories qualssevol, substituint les desigualtats per bijeccions entre conjunts de morfismes.

\section{Adjuncions entre preordres}

Siguin $(P,\leq)$ i $(Q,\leq)$ dos preordres, vistos com a categories, i $f\colon P\to Q$, $g\colon Q\to P$ dues aplicacions monòtones.

\begin{definicio}\index{adjunció!entre preordres}
Diem que $f$ és \textbf{adjunta per l'esquerra} de $g$ (i $g$ \textbf{adjunta per la dreta} de $f$), i escrivim $f\dashv g$, quan per a tot $x\in P$ i tot $y\in Q$,
\[
f(x)\leq y
\quad\Longleftrightarrow\quad
x\leq g(y).
\]
\end{definicio}

\begin{observacio}[Unitat i counitat en un preordre]\index{connexió de Galois}
Si $f\dashv g$, prenent $y=f(x)$ i $x=g(y)$ a la definició s'obté
\[
x\leq g(f(x)),
\qquad
f(g(y))\leq y
\]
per a tot $x\in P$ i tot $y\in Q$. Recíprocament, si $f$ i $g$ són monòtones i compleixen aquestes dues desigualtats, aleshores $f\dashv g$: si $f(x)\leq y$, aleshores $x\leq g(f(x))\leq g(y)$, i si $x\leq g(y)$, aleshores $f(x)\leq f(g(y))\leq y$. Aquestes dues desigualtats són exactament la \emph{unitat} i la \emph{counitat} de l'adjunció, que veurem en general més endavant; en una categoria prima, les identitats triangulars no aporten cap igualtat nova, perquè entre dos objectes hi ha com a màxim una fletxa. Una adjunció entre preordres s'anomena sovint \textbf{connexió de Galois}.
\end{observacio}

\begin{exemple}[Conjunció i implicació]\index{adjunció!conjunció i implicació}
Sigui $\cat{Prov}_T$ el preordre de fórmules d'un sistema proposicional $T$ amb les regles usuals de $\wedge$ i $\to$, i fixem una fórmula $A$. Les operacions $A\wedge-$ i $A\to-$ són monòtones, i per a totes les fórmules $B$ i $C$
\[
A\wedge B\vdash_T C
\quad\Longleftrightarrow\quad
B\vdash_T A\to C,
\]
de manera que
\[
A\wedge-\;\dashv\;A\to-.
\]
Aquesta adjunció codifica les regles de la implicació; és la forma categòrica del teorema de deducció. És un primer indici que alguns connectius lògics es caracteritzen mitjançant adjuncions, idea que reapareixerà de manera sistemàtica en la lògica categòrica.
\end{exemple}

\begin{exemple}[Imatge, antiimatge i quantificadors]\label{ex:quantificadors-set}\index{adjunció!quantificadors}\index{quantificador!com a adjunt}
Sigui $f\colon X\to Y$ una aplicació, i considerem els preordres $\mathcal P(X)$ i $\mathcal P(Y)$ de les parts, ordenats per inclusió. L'antiimatge
\[
f^{*}(T)=f^{-1}(T),\qquad T\subseteq Y,
\]
és monòtona, i té adjunts a tots dos costats: per a $S\subseteq X$, posem
\[
\exists_f(S)=f[S]=\{f(x)\mid x\in S\},
\qquad
\forall_f(S)=\{y\in Y\mid f^{-1}(\{y\})\subseteq S\}.
\]
Aleshores, per a tot $S\subseteq X$ i tot $T\subseteq Y$,
\[
\exists_f(S)\subseteq T
\quad\Longleftrightarrow\quad
S\subseteq f^{*}(T),
\qquad\qquad
f^{*}(T)\subseteq S
\quad\Longleftrightarrow\quad
T\subseteq\forall_f(S).
\]
La primera equivalència diu que les imatges dels elements de $S$ són a $T$ exactament quan tots els elements de $S$ van a parar a $T$. Per a la segona: si $f^{-1}(T)\subseteq S$ i $y\in T$, tot $x$ amb $f(x)=y$ és a $f^{-1}(T)$ i, per tant, a $S$, de manera que $y\in\forall_f(S)$; i recíprocament, si $T\subseteq\forall_f(S)$ i $f(x)\in T$, aleshores $x\in f^{-1}(\{f(x)\})\subseteq S$. Per tant,
\[
\exists_f\;\dashv\;f^{*}\;\dashv\;\forall_f.
\]

El nom dels adjunts s'explica quan $f$ és una projecció $\pi\colon X\times A\to X$. Un subconjunt $S\subseteq X\times A$ és un predicat en dues variables, i
\[
\begin{aligned}
\exists_\pi(S)&=\{x\mid \text{existeix }a\in A\text{ amb }(x,a)\in S\},\\
\forall_\pi(S)&=\{x\mid (x,a)\in S\text{ per a tot }a\in A\}:
\end{aligned}
\]
són exactament les lectures conjuntistes de la quantificació existencial i universal sobre la variable d'$A$, mentre que $\pi^{*}$ afegeix una variable muda. Al capítol de subobjectes veurem en quines categories existeixen aquests dos adjunts; aquí només afirmem el cas de $\cat{Set}$.
\end{exemple}

\begin{exemple}[Interior topològic]\index{adjunció!interior topològic}
Sigui $X$ un espai topològic, $\mathcal P(X)$ el preordre de les parts per inclusió i $\mathcal O(X)$ el dels oberts. La injecció $i\colon\mathcal O(X)\to\mathcal P(X)$ té com a adjunt per la dreta l'\textbf{interior} $\mathrm{int}\colon\mathcal P(X)\to\mathcal O(X)$:
\[
i(U)\subseteq S
\quad\Longleftrightarrow\quad
U\subseteq\mathrm{int}(S),
\]
per a tot obert $U$ i tota part $S$. És a dir, l'interior és el més gran obert contingut en $S$.
\end{exemple}

\section{Adjuncions entre categories}

La definició general substitueix l'equivalència de desigualtats per una \textbf{bijecció entre conjunts de morfismes}, natural en totes dues variables.

\begin{definicio}[Adjunció]\index{adjunció}
\label{def:adjuncio}
Siguin $F\colon\cat{C}\to\cat{D}$ i $G\colon\cat{D}\to\cat{C}$ dos functors entre categories localment petites. Una \textbf{adjunció} $F\dashv G$ és una bijecció
\[
\theta_{A,B}\colon\Hom_{\cat{D}}(F(A),B)\;\xrightarrow{\ \cong\ }\;\Hom_{\cat{C}}(A,G(B)),
\]
per a cada $A$ de $\cat{C}$ i $B$ de $\cat{D}$, \textbf{natural} en $A$ i en $B$. Es diu que $F$ és l'\textbf{adjunt per l'esquerra} i $G$ l'\textbf{adjunt per la dreta}.
\end{definicio}

\begin{observacio}
La naturalitat vol dir que $\theta$ és un isomorfisme natural entre els dos functors
\[
\Hom_{\cat{D}}(F(-),-)\quad\text{i}\quad\Hom_{\cat{C}}(-,G(-)),
\]
tots dos de $\cat{C}\op\times\cat{D}$ cap a $\cat{Set}$. Explícitament, per a $a\colon A'\to A$ a $\cat{C}$, $b\colon B\to B'$ a $\cat{D}$ i $\varphi\colon F(A)\to B$, el pas per $\theta$ commuta amb la postcomposició i amb la precomposició:
\[
\theta_{A,B'}(b\comp\varphi)=G(b)\comp\theta_{A,B}(\varphi),
\qquad\qquad
\theta_{A',B}\bigl(\varphi\comp F(a)\bigr)=\theta_{A,B}(\varphi)\comp a.
\]
En paraules: transposar i després compondre amb $G(b)$ o amb $a$ és el mateix que compondre primer amb $b$ o amb $F(a)$ i després transposar. Són els dos quadrats de naturalitat
\[
\begin{tikzpicture}[baseline=(m.center)]
  \matrix (m) [matrix of math nodes, row sep=2.8em, column sep=3.6em] {
    \Hom(F(A),B) & \Hom(A,G(B)) \\
    \Hom(F(A),B') & \Hom(A,G(B')) \\
  };
  \draw[-{Stealth[scale=1.1]}] (m-1-1) -- node[above] {$\theta_{A,B}$} (m-1-2);
  \draw[-{Stealth[scale=1.1]}] (m-1-1) -- node[left] {$b\comp-$} (m-2-1);
  \draw[-{Stealth[scale=1.1]}] (m-1-2) -- node[right] {$G(b)\comp-$} (m-2-2);
  \draw[-{Stealth[scale=1.1]}] (m-2-1) -- node[below] {$\theta_{A,B'}$} (m-2-2);
\end{tikzpicture}
\]
i
\[
\begin{tikzpicture}[baseline=(m.center)]
  \matrix (m) [matrix of math nodes, row sep=2.8em, column sep=3.6em] {
    \Hom(F(A),B) & \Hom(A,G(B)) \\
    \Hom(F(A'),B) & \Hom(A',G(B)) \\
  };
  \draw[-{Stealth[scale=1.1]}] (m-1-1) -- node[above] {$\theta_{A,B}$} (m-1-2);
  \draw[-{Stealth[scale=1.1]}] (m-1-1) -- node[left] {$-\comp F(a)$} (m-2-1);
  \draw[-{Stealth[scale=1.1]}] (m-1-2) -- node[right] {$-\comp a$} (m-2-2);
  \draw[-{Stealth[scale=1.1]}] (m-2-1) -- node[below] {$\theta_{A',B}$} (m-2-2);
\end{tikzpicture}
\]
on els vèrtexs són conjunts de morfismes i les fletxes verticals actuen per composició. Els morfismes que es corresponen sota $\theta$ s'anomenen \textbf{transposats} l'un de l'altre: si $\varphi\colon F(A)\to B$, el seu transposat és $\theta(\varphi)\colon A\to G(B)$.
\end{observacio}

\begin{observacio}
El cas dels preordres n'és un cas particular: si $P,Q$ són preordres, un conjunt de morfismes $\Hom(x,y)$ té com a molt un element, de manera que la bijecció $\Hom(f(x),y)\cong\Hom(x,g(y))$ es redueix a l'equivalència $f(x)\leq y\Leftrightarrow x\leq g(y)$. Així, la definició general estén literalment la dels preordres.
\end{observacio}

\section{Unitat i counitat}

Una adjunció es pot descriure, de manera equivalent, sense esmentar la bijecció, mitjançant dues transformacions naturals distingides.

\begin{definicio}[Unitat i counitat]\index{adjunció!unitat}\index{adjunció!counitat}
\label{def:unitat-counitat}
Donada una adjunció $F\dashv G$ amb bijecció $\theta$:
\begin{itemize}
\item la \textbf{unitat} és la transformació natural $\eta\colon\id_{\cat{C}}\Rightarrow G\comp F$ de components
\[
\eta_A=\theta_{A,F(A)}(\id_{F(A)})\colon A\to G(F(A));
\]
\item la \textbf{counitat} és la transformació natural $\varepsilon\colon F\comp G\Rightarrow\id_{\cat{D}}$ de components
\[
\varepsilon_B=\theta_{G(B),B}^{-1}(\id_{G(B)})\colon F(G(B))\to B.
\]
\end{itemize}
\end{definicio}

\begin{observacio}[De la unitat i la counitat a la bijecció]\label{obs:adj-transposats}
La unitat i la counitat permeten recuperar la bijecció. Donat $\varphi\colon F(A)\to B$, el seu transposat és
\[
\theta(\varphi)=G(\varphi)\comp\eta_A\colon A\to G(B),
\]
i dualment, donat $\psi\colon A\to G(B)$, el seu transposat invers és
\[
\theta^{-1}(\psi)=\varepsilon_B\comp F(\psi)\colon F(A)\to B.
\]
\end{observacio}

\begin{definicio}[Fletxa universal]\label{def:fletxa-universal}\index{fletxa universal}\index{adjunció!fletxa universal}
Sigui $G\colon\cat{D}\to\cat{C}$ un functor i $A$ un objecte de $\cat{C}$. Una \textbf{fletxa universal} de $A$ cap a $G$ és una parella $(R,\eta_A)$ formada per un objecte $R$ de $\cat{D}$ i un morfisme $\eta_A\colon A\to G(R)$ tal que tot $\psi\colon A\to G(B)$ factoritza de manera única com $\psi=G(\varphi)\comp\eta_A$, amb $\varphi\colon R\to B$.
\end{definicio}

En una adjunció $F\dashv G$, cada component de la unitat és una fletxa universal, amb $R=F(A)$ (exercici~\ref{exr:adj-unitat-universal} dels Exercicis):
\[
\begin{tikzpicture}[baseline=(m.center)]
  \matrix (m) [matrix of math nodes, column sep=4em, row sep=2.7em] {
    A & G(F(A)) \\
      & G(B) \\
  };
  \draw[-{Stealth[scale=1.1]}] (m-1-1) -- node[above] {$\eta_A$} (m-1-2);
  \draw[-{Stealth[scale=1.1]}] (m-1-1) -- node[below left] {$\psi$} (m-2-2);
  \draw[-{Stealth[scale=1.1]},dashed] (m-1-2) -- node[right] {$G(\varphi)$} (m-2-2);
\end{tikzpicture}
\qquad
\psi=G(\varphi)\comp\eta_A,\quad\varphi\text{ únic}.
\]
Aquesta és la formulació de l'adjunció que connecta directament amb els objectes lliures.

\begin{exemple}[El monoide lliure]\index{adjunció!lliure-oblit}
Sigui $U\colon\cat{Mon}\to\cat{Set}$ el functor oblit dels monoides i $F\colon\cat{Set}\to\cat{Mon}$ el functor que envia un conjunt $X$ al \textbf{monoide lliure} $X^*$ de les paraules finites en $X$, amb la concatenació. Aleshores $F\dashv U$: un homomorfisme de monoides $X^*\to M$ està determinat de manera única pels valors sobre les lletres, és a dir, per una aplicació $X\to U(M)$. Aquesta és exactament la bijecció
\[
\Hom_{\cat{Mon}}(X^*,M)\cong\Hom_{\cat{Set}}(X,U(M)).
\]
La unitat $\eta_X\colon X\to U(X^*)$ envia cada element a la paraula d'una lletra; és la fletxa universal que exhibeix la propietat universal del monoide lliure. Aquest patró ---\emph{lliure $\dashv$ oblit}--- es repeteix per a grups, espais vectorials, categories i moltes altres estructures.
\end{exemple}

\begin{exemple}[La categoria lliure d'un graf: $\Free\dashv U$]\label{ex:free-u-adjuncio}\index{adjunció!categoria lliure}
Els capítols anteriors ja han preparat aquest exemple. El functor graf subjacent $U\colon\cat{Cat}\to\cat{Graph}$ i el functor categoria lliure $\Free\colon\cat{Graph}\to\cat{Cat}$ satisfan la bijecció
\[
\Hom_{\cat{Cat}}\bigl(\Free(G),\cat{D}\bigr)\;\cong\;\Hom_{\cat{Graph}}\bigl(G,U(\cat{D})\bigr),
\]
natural en $G$ i en $\cat{D}$, que el capítol~\ref{cap:yoneda} va obtenir com la representabilitat del functor $\cat{D}\mapsto\Hom_{\cat{Graph}}(G,U(\cat{D}))$. És a dir, $\Free\dashv U$. El que aporta ara el llenguatge de les adjuncions és identificar-ne les dades:
\begin{itemize}
  \item la \emph{unitat} $\eta_G\colon G\to U(\Free(G))$ és la inserció dels generadors: cada vèrtex va a ell mateix i cada aresta, al camí de longitud~$1$ que la recorre;
  \item la \emph{counitat} $\varepsilon_{\cat{C}}\colon\Free(U(\cat{C}))\to\cat{C}$ \emph{avalua} els camins formals: envia un camí $(f_1,\dots,f_n)$ de fletxes de $\cat{C}$, recorregut en aquest ordre, a la composta $f_n\comp\cdots\comp f_1$, i el camí buit en un objecte, a la seva identitat.
\end{itemize}
\end{exemple}

\section{Identitats triangulars}

La unitat i la counitat no són independents: satisfan dues equacions que, de fet, caracteritzen les adjuncions.

\begin{proposicio}[Identitats triangulars]\label{prop:triangulars}\index{adjunció!identitats triangulars}
Per a una adjunció $F\dashv G$ amb unitat $\eta$ i counitat $\varepsilon$, es compleixen les \textbf{identitats triangulars}:
\[
\varepsilon_{F(A)}\comp F(\eta_A)=\id_{F(A)},
\qquad
G(\varepsilon_B)\comp\eta_{G(B)}=\id_{G(B)},
\]
és a dir, commuten els dos triangles
\[
\begin{tikzpicture}[baseline=(m.center)]
  \matrix (m) [matrix of math nodes, column sep=3.6em, row sep=2.7em] {
    F(A) & F(G(F(A))) \\
         & F(A) \\
  };
  \draw[-{Stealth[scale=1.1]}] (m-1-1) -- node[above] {$F(\eta_A)$} (m-1-2);
  \draw[-{Stealth[scale=1.1]}] (m-1-1) -- node[below left] {$\id_{F(A)}$} (m-2-2);
  \draw[-{Stealth[scale=1.1]}] (m-1-2) -- node[right] {$\varepsilon_{F(A)}$} (m-2-2);
\end{tikzpicture}
\qquad\qquad
\begin{tikzpicture}[baseline=(m.center)]
  \matrix (m) [matrix of math nodes, column sep=3.6em, row sep=2.7em] {
    G(B) & G(F(G(B))) \\
         & G(B) \\
  };
  \draw[-{Stealth[scale=1.1]}] (m-1-1) -- node[above] {$\eta_{G(B)}$} (m-1-2);
  \draw[-{Stealth[scale=1.1]}] (m-1-1) -- node[below left] {$\id_{G(B)}$} (m-2-2);
  \draw[-{Stealth[scale=1.1]}] (m-1-2) -- node[right] {$G(\varepsilon_B)$} (m-2-2);
\end{tikzpicture}
\]
Recíprocament, dues transformacions naturals $\eta\colon\id\Rightarrow G F$ i $\varepsilon\colon F G\Rightarrow\id$ que les satisfacin defineixen una adjunció $F\dashv G$.
\end{proposicio}

\begin{prova}
Comprovem la primera identitat. Pel càlcul del transposat, $\theta^{-1}(\psi)=\varepsilon_B\comp F(\psi)$; prenent $B=F(A)$ i $\psi=\eta_A=\theta(\id_{F(A)})$,
\[
\varepsilon_{F(A)}\comp F(\eta_A)=\theta^{-1}(\eta_A)=\theta^{-1}(\theta(\id_{F(A)}))=\id_{F(A)}.
\]
La segona identitat és dual, usant $\theta(\varphi)=G(\varphi)\comp\eta_A$ amb $A=G(B)$ i $\varphi=\varepsilon_B$. Per al recíproc, es defineix la bijecció per $\theta(\varphi)=G(\varphi)\comp\eta_A$ i $\theta^{-1}(\psi)=\varepsilon_B\comp F(\psi)$; les identitats triangulars garanteixen que són inverses l'una de l'altra, i la naturalitat de $\eta,\varepsilon$ dona la de $\theta$.
\end{prova}

\begin{observacio}
Les identitats triangulars reben aquest nom perquè expressen la commutativitat de dos triangles que involucren $F$, $G$, $\eta$ i $\varepsilon$. Donen una definició d'adjunció \emph{purament diagramàtica}, sense referència a conjunts de morfismes, la qual cosa és sovint la més còmoda per treballar i l'única disponible en contextos on no hi ha homconjunts (per exemple, en 2-categories abstractes).
\end{observacio}

\begin{observacio}[Tres presentacions d'una adjunció]\label{obs:tres-presentacions}\index{adjunció!tres presentacions}
Hem trobat tres maneres de donar una adjunció entre $\cat{C}$ i $\cat{D}$. Convé no barrejar-les, perquè no parteixen de les mateixes dades.
\begin{center}
\footnotesize
\renewcommand{\arraystretch}{1.3}
\begin{tabularx}{\linewidth}{@{}>{\raggedright\arraybackslash}p{0.2\linewidth}>{\raggedright\arraybackslash}X>{\raggedright\arraybackslash}X@{}}
\toprule
\textbf{Presentació} & \textbf{Dades} & \textbf{Condicions}\\
\midrule
bijecció d'homconjunts (definició~\ref{def:adjuncio}) & functors $F$ i $G$; bijeccions $\theta_{A,B}\colon\Hom(F(A),B)\to\Hom(A,G(B))$ & naturalitat en $A$ i en $B$\\
unitat i counitat (proposició~\ref{prop:triangulars}) & functors $F$ i $G$; transformacions naturals $\eta\colon\id\Rightarrow GF$ i $\varepsilon\colon FG\Rightarrow\id$ & identitats triangulars\\
fletxes universals (definició~\ref{def:fletxa-universal}) & només el functor $G$; per a cada $A$, un objecte $F(A)$ i una fletxa $\eta_A\colon A\to G(F(A))$ & cada $\eta_A$ és universal; $F$ no es dona sobre morfismes, sinó que la propietat universal el determina\\
\bottomrule
\end{tabularx}
\end{center}
Els passos d'una presentació a una altra són aquests:
\begin{itemize}
\item De la primera a la segona: $\eta$ i $\varepsilon$ són els transposats de les identitats (definició~\ref{def:unitat-counitat}).
\item De la segona a la primera: les fórmules de l'observació~\ref{obs:adj-transposats}, que són inverses gràcies a les identitats triangulars (exercici~\ref{exr:adj-triangulars-reciproc} dels Exercicis).
\item De la primera a la tercera: la unitat és universal (exercici~\ref{exr:adj-unitat-universal}).
\item De la tercera a la primera: la propietat universal defineix $F$ sobre morfismes i la bijecció $\varphi\mapsto G(\varphi)\comp\eta_A$ (exercici~\ref{exr:adj-fletxes-universals}).
\end{itemize}
La tercera presentació és la més econòmica, perquè no cal conèixer $F$ per endavant; per això és la que s'usa per \emph{construir} adjunts, com els objectes lliures. La segona és la més còmoda per calcular, i la primera és la que dona directament els teoremes de preservació de límits de la secció següent.
\end{observacio}

\begin{observacio}[Adjuncions i equivalències]\index{equivalència de categories!adjunta}
Al capítol~\ref{cap:transformacions} vam definir una equivalència mitjançant un quasiinvers $G$ i dos isomorfismes naturals orientats com $\eta\colon\id_{\cat{C}}\Rightarrow G\comp F$ i $\varepsilon\colon F\comp G\Rightarrow\id_{\cat{D}}$ (definició~\ref{def:equivalencia}), la mateixa orientació que la unitat i la counitat d'una adjunció $F\dashv G$. La coincidència no és casual. Tota equivalència es pot equipar amb una estructura d'\textbf{equivalència adjunta}: modificant, si cal, un dels dos isomorfismes naturals, es pot aconseguir que $\eta$ i $\varepsilon$ satisfacin les identitats triangulars, i aleshores $F\dashv G$ amb unitat i counitat isomorfismes naturals. Recíprocament, una adjunció en què la unitat i la counitat són isomorfismes naturals és una equivalència. Les equivalències són, doncs, les adjuncions més simètriques possibles.
\end{observacio}

\section{Unicitat dels adjunts}

\begin{proposicio}\index{adjunció!unicitat}
Els adjunts són únics llevat d'isomorfisme natural. Si $F\dashv G$ i $F\dashv G'$, aleshores $G\cong G'$; i si $F\dashv G$ i $F'\dashv G$, aleshores $F\cong F'$.
\end{proposicio}

\begin{prova}
Si $F\dashv G$ i $F\dashv G'$, aleshores per a tot $A,B$,
\[
\Hom(A,G(B))\cong\Hom(F(A),B)\cong\Hom(A,G'(B)),
\]
naturalment en $A$. Pel corol·lari del lema de Yoneda ---dos objectes amb functors representables naturalment isomorfs són isomorfs---, $G(B)\cong G'(B)$ naturalment en $B$, és a dir, $G\cong G'$. El cas de $F$ és dual.
\end{prova}

\section{Els adjunts i els límits}

El resultat més útil sobre adjuncions diu com interactuen amb límits i colímits. És la raó de fons per la qual molts functors oblit preserven els productes i, en general, els límits.

\begin{teorema}[Els adjunts per la dreta preserven límits]\index{adjunció!preservació de límits}
\label{teo:adjunts-preserven-limits}
Si $G\colon\cat{D}\to\cat{C}$ és un adjunt per la dreta, aleshores $G$ preserva tots els límits petits que existeixen a $\cat{D}$. Dualment, un adjunt per l'esquerra preserva tots els colímits petits que existeixen.
\end{teorema}

\begin{prova}
Sigui $F\colon\cat{C}\to\cat{D}$ l'adjunt per l'esquerra de $G$, amb bijecció $\theta$ (definició~\ref{def:adjuncio}). Farem servir les dues equacions de naturalitat de $\theta$: per a $b\colon B\to B'$ a $\cat{D}$ i $a\colon A'\to A$ a $\cat{C}$,
\[
\theta(b\comp\varphi)=G(b)\comp\theta(\varphi),
\qquad\qquad
\theta\bigl(\varphi\comp F(a)\bigr)=\theta(\varphi)\comp a.
\]

\emph{Límits.} Sigui $(L,\lambda)$ un límit d'un diagrama $D\colon\cat{I}\to\cat{D}$. La família $G\lambda=(G(\lambda_i))_i$ és un con sobre $G\comp D$, perquè $G(D_u)\comp G(\lambda_i)=G(D_u\comp\lambda_i)=G(\lambda_j)$ per a cada $u\colon i\to j$. Vegem que és universal. Sigui $(A,\alpha)$ un con sobre $G\comp D$, amb $\alpha_i\colon A\to G(D_i)$ i $G(D_u)\comp\alpha_i=\alpha_j$, i siguin $\varphi_i=\theta^{-1}(\alpha_i)\colon F(A)\to D_i$ els transposats. Per la primera equació de naturalitat,
\[
\theta(D_u\comp\varphi_i)=G(D_u)\comp\alpha_i=\alpha_j=\theta(\varphi_j),
\]
i com que $\theta$ és bijectiva, $D_u\comp\varphi_i=\varphi_j$: les $\varphi_i$ formen un con sobre $D$ amb vèrtex $F(A)$. Per tant hi ha una única $h\colon F(A)\to L$ amb $\lambda_i\comp h=\varphi_i$ per a tot $i$. Posem $k=\theta(h)\colon A\to G(L)$. De nou per la primera equació,
\[
G(\lambda_i)\comp k=\theta(\lambda_i\comp h)=\theta(\varphi_i)=\alpha_i,
\]
de manera que $k$ factoritza el con $\alpha$. És l'única factorització: si $k'\colon A\to G(L)$ compleix $G(\lambda_i)\comp k'=\alpha_i$ per a tot $i$, aleshores $h'=\theta^{-1}(k')$ compleix $\theta(\lambda_i\comp h')=G(\lambda_i)\comp k'=\alpha_i=\theta(\varphi_i)$, d'on $\lambda_i\comp h'=\varphi_i$; per la unicitat de $h$, $h'=h$, i $k'=\theta(h')=k$. Així $(G(L),G\lambda)$ és un límit de $G\comp D$.

\emph{Colímits.} Sigui ara $(K,\kappa)$ un colímit d'un diagrama $E\colon\cat{I}\to\cat{C}$. La família $F\kappa=(F(\kappa_i))_i$ és un cocon sobre $F\comp E$, perquè $F(\kappa_j)\comp F(E_u)=F(\kappa_j\comp E_u)=F(\kappa_i)$. Sigui $(B,\beta)$ un cocon sobre $F\comp E$, amb $\beta_i\colon F(E_i)\to B$ i $\beta_j\comp F(E_u)=\beta_i$, i siguin $\psi_i=\theta(\beta_i)\colon E_i\to G(B)$. Per la segona equació de naturalitat,
\[
\psi_j\comp E_u=\theta\bigl(\beta_j\comp F(E_u)\bigr)=\theta(\beta_i)=\psi_i,
\]
de manera que les $\psi_i$ formen un cocon sobre $E$, i hi ha una única $m\colon K\to G(B)$ amb $m\comp\kappa_i=\psi_i$. Posem $n=\theta^{-1}(m)\colon F(K)\to B$. Aleshores $\theta\bigl(n\comp F(\kappa_i)\bigr)=m\comp\kappa_i=\psi_i=\theta(\beta_i)$, és a dir, $n\comp F(\kappa_i)=\beta_i$. Si $n'$ també compleix $n'\comp F(\kappa_i)=\beta_i$, aleshores $\theta(n')\comp\kappa_i=\theta(\beta_i)=\psi_i$, d'on $\theta(n')=m$ i $n'=n$. Així $(F(K),F\kappa)$ és un colímit de $F\comp E$.
\end{prova}

\begin{observacio}[La demostració en termes de representables]
La demostració anterior és la versió \enquote{element a element} d'una cadena de bijeccions. Per a tot objecte $A$ de $\cat{C}$, usant l'adjunció i que $\Hom(F(A),-)$ preserva límits (proposició~\ref{prop:representables-limits}),
\[
\begin{aligned}
\Hom(A,G(L))&\cong\Hom(F(A),L) &&\text{(adjunció)}\\
&\cong\varprojlim\nolimits_i\Hom(F(A),D_i) &&\text{($\Hom(F(A),-)$ preserva límits)}\\
&\cong\varprojlim\nolimits_i\Hom(A,G(D_i)) &&\text{(adjunció, natural en $D_i$)}\\
&\cong\mathrm{Con}(A,G\comp D) &&\text{(límits a $\cat{Set}$)},
\end{aligned}
\]
naturalment en $A$; la bijecció composta envia $f\colon A\to G(L)$ a $(G(\lambda_i)\comp f)_i$. Llegida així, l'afirmació diu que $G(L)$ representa el prefeix dels cons sobre $G\comp D$, amb element universal $G\lambda$ (proposició~\ref{prop:limit-representable}).
\end{observacio}

\begin{observacio}[Què vol dir exactament \enquote{preservar}]\label{obs:preservar-canonic}\index{morfisme de comparació}
El teorema no afirma cap igualtat $G(\varprojlim D)=\varprojlim(G\comp D)$: els límits només estan determinats llevat d'isomorfisme, i una igualtat no tindria sentit. Tampoc no afirma només que \emph{existeixi} algun isomorfisme entre aquests dos objectes. Afirma el que diu la definició del capítol~\ref{cap:limits}: el con imatge $(G(L),G\lambda)$ \emph{és un con límit} de $G\comp D$.

Si hem triat un límit $\bigl(\varprojlim(G\comp D),\pi\bigr)$ de $G\comp D$, la condició es pot expressar amb un sol morfisme. El con $G\lambda$ factoritza de manera única a través del límit triat, i dona l'únic \textbf{morfisme de comparació}
\[
c\colon G(L)\to\varprojlim(G\comp D),\qquad \pi_i\comp c=G(\lambda_i)\ \text{per a tot } i.
\]
Aleshores $G$ preserva el límit si i només si $c$ és un isomorfisme. En efecte, si $(G(L),G\lambda)$ és un límit, l'únic isomorfisme compatible amb les projeccions entre els dos límits de $G\comp D$ compleix $\pi_i\comp c=G(\lambda_i)$, i per tant és $c$. Recíprocament, si $c$ és invertible, el con $G\lambda$ és isomorf al con límit $\pi$ per un isomorfisme compatible amb les projeccions, i per tant és un límit. En aquest sentit precís escrivim $G(\varprojlim D)\cong\varprojlim(G\comp D)$: l'isomorfisme és el \emph{canònic} $c$, no un isomorfisme qualsevol. A més, la noció no depèn del límit triat de $D$, perquè dos límits de $D$ són isomorfs per un únic isomorfisme compatible amb les projeccions, i $G$ el transporta.

La diferència importa. El functor $X\mapsto X+X$ de $\cat{Set}$ en si mateix no preserva productes binaris. El morfisme de comparació $(X\times Y)+(X\times Y)\to(X+X)\times(Y+Y)$ envia cada parella de la còpia $k$ a la parella de còpies $(k,k)$, i no arriba als elements en què les dues còpies són diferents. Tanmateix, per a $X=Y=\mathbb N$ tots dos conjunts són numerables, i per tant \emph{hi ha} una bijecció entre ells. Dualment, $F$ preserva el colímit $(K,\kappa)$ de $E$ quan el morfisme de comparació $\colimop(F\comp E)\to F(K)$ és un isomorfisme.
\end{observacio}

\begin{observacio}
Aquest teorema es recorda amb les sigles angleses \enquote{RAPL} (\emph{right adjoints preserve limits}) i \enquote{LAPC} (\emph{left adjoints preserve colimits}). Explica de cop molts fets. Els functors oblit que són adjunts per la dreta preserven límits, però això no implica que preservin colímits; en molts exemples algebraics no els preserven. Dualment, els functors lliures, com a adjunts per l'esquerra, preserven colímits, però no hi ha cap teorema general que els obligui a preservar límits. El teorema dona una sola direcció i no prohibeix l'altra: un functor amb adjunts a tots dos costats, com l'oblit $\cat{Top}\to\cat{Set}$, preserva tant límits com colímits. És també un criteri negatiu útil: un functor que \emph{no} preserva límits no pot ser un adjunt per la dreta.
\end{observacio}

\subsection{El límit i el colímit com a adjunts del functor diagonal}\index{límit!com a adjunt}\index{colímit!com a adjunt}

Hi ha una segona manera, més sintètica, en què límits i colímits són adjuncions ---no ja \emph{preservats} per adjunts, sinó ells mateixos adjunts. Fixem una categoria d'índexs $\cat{I}$ i considerem el \textbf{functor diagonal}
\[
\Delta\colon\cat{C}\longrightarrow\cat{C}^{\cat{I}},
\]
que envia un objecte $A$ al diagrama constant de valor $A$ (i un morfisme a la transformació natural constant). Prendre el límit i prendre el colímit d'un diagrama són operacions $\cat{C}^{\cat{I}}\to\cat{C}$, i resulten ser, respectivament, l'adjunt per la dreta i per l'esquerra de $\Delta$.

\begin{teorema}[$\colimop\dashv\Delta\dashv\varprojlim$]\index{adjunció!límit-diagonal}
\label{teo:lim-adjunt}
Sigui $\cat{I}$ una categoria petita i suposem que $\cat{C}$ té tots els límits i colímits de forma $\cat{I}$. Fixem, per a cada diagrama de forma $\cat{I}$, un límit i un colímit; quan la família de diagrames és pròpiament gran, fem servir la convenció d'elecció global de l'apèndix~\ref{cap:mida}. La propietat universal estén aquestes tries, de manera única, a functors $\varprojlim_{\cat{I}}$ i $\colimop_{\cat{I}}$ (pas~2 de la demostració). Aleshores el functor diagonal $\Delta\colon\cat{C}\to\cat{C}^{\cat{I}}$ té adjunt per l'esquerra $\colimop_{\cat{I}}$ i adjunt per la dreta $\varprojlim_{\cat{I}}$:
\[
\colimop_{\cat{I}}\;\dashv\;\Delta\;\dashv\;\varprojlim_{\cat{I}}.
\]
\end{teorema}

\begin{prova}
Recordem que $\Delta$ envia $A$ al diagrama constant $\Delta A$, amb $(\Delta A)_i=A$ i $(\Delta A)_u=\id_A$, i un morfisme $a\colon A'\to A$ a la transformació natural $\Delta a$ de components $a$. Escrivim $(\varprojlim D,\lambda^D)$ per al límit triat de cada diagrama $D$.

\emph{Pas 1: un morfisme $\Delta A\Rightarrow D$ és un con.} Una transformació natural $\alpha\colon\Delta A\Rightarrow D$ té components $\alpha_i\colon A\to D_i$, i el seu quadrat de naturalitat per a $u\colon i\to j$ diu $D_u\comp\alpha_i=\alpha_j\comp\id_A=\alpha_j$. És exactament la condició de con. Per tant, $\Hom_{\cat{C}^{\cat{I}}}(\Delta A,D)$ és el conjunt dels cons sobre $D$ amb vèrtex $A$.

\emph{Pas 2: $\varprojlim$ és un functor.} Sigui $\beta\colon D\Rightarrow D'$ un morfisme de $\cat{C}^{\cat{I}}$. La família $\beta_i\comp\lambda^D_i$ és un con sobre $D'$, perquè $D'_u\comp\beta_i\comp\lambda^D_i=\beta_j\comp D_u\comp\lambda^D_i=\beta_j\comp\lambda^D_j$. Definim $\varprojlim\beta$ com l'única fletxa amb
\[
\lambda^{D'}_i\comp\varprojlim\beta=\beta_i\comp\lambda^D_i\qquad\text{per a tot } i.
\]
La identitat i $\varprojlim\id_D$ compleixen totes dues aquesta condició per a $\beta=\id_D$. Per a $\beta'\colon D'\Rightarrow D''$, la composició $\varprojlim\beta'\comp\varprojlim\beta$ la compleix per a $\beta'\comp\beta$. Per la unicitat de la factorització, $\varprojlim$ preserva identitats i composició.

\emph{Pas 3: la bijecció.} Definim $\theta_{A,D}\colon\Hom_{\cat{C}^{\cat{I}}}(\Delta A,D)\to\Hom_{\cat{C}}(A,\varprojlim D)$ enviant cada con $\alpha$ a l'única fletxa $\theta(\alpha)$ amb $\lambda^D_i\comp\theta(\alpha)=\alpha_i$ per a tot $i$. Per la propietat universal del límit, $\theta$ és bijectiva, amb inversa $h\mapsto(\lambda^D_i\comp h)_i=\lambda^D\comp\Delta h$.

\emph{Pas 4: naturalitat.} Per a $a\colon A'\to A$, les fletxes $\theta(\alpha\comp\Delta a)$ i $\theta(\alpha)\comp a$ compleixen totes dues $\lambda^D_i\comp(-)=\alpha_i\comp a$, i per tant coincideixen: és la naturalitat en $A$. Per a $\beta\colon D\Rightarrow D'$,
\[
\lambda^{D'}_i\comp\varprojlim\beta\comp\theta(\alpha)=\beta_i\comp\lambda^D_i\comp\theta(\alpha)=\beta_i\comp\alpha_i=\lambda^{D'}_i\comp\theta(\beta\comp\alpha),
\]
d'on $\varprojlim\beta\comp\theta(\alpha)=\theta(\beta\comp\alpha)$: és la naturalitat en $D$. Així $\Delta\dashv\varprojlim_{\cat{I}}$.

\emph{L'adjunció de l'esquerra} és la duala, pas per pas. Una transformació natural $\gamma\colon D\Rightarrow\Delta A$ té components $\gamma_i\colon D_i\to A$ amb $\gamma_j\comp D_u=\gamma_i$, és a dir, és un cocon. Si $(\colimop D,\kappa^D)$ és el colímit triat, $\colimop\beta$ és l'única fletxa amb $\colimop\beta\comp\kappa^D_i=\kappa^{D'}_i\comp\beta_i$, i és functorial per la mateixa raó. La bijecció $\Hom_{\cat{C}}(\colimop D,A)\to\Hom_{\cat{C}^{\cat{I}}}(D,\Delta A)$ és $h\mapsto(h\comp\kappa^D_i)_i$, bijectiva per la propietat universal del colímit. És natural en $A$, perquè $(a\comp h)\comp\kappa^D_i=a\comp(h\comp\kappa^D_i)$, i en $D$, perquè $h\comp\colimop\beta\comp\kappa^D_i=(h\comp\kappa^{D'}_i)\comp\beta_i$. Així $\colimop_{\cat{I}}\dashv\Delta$.
\end{prova}

\begin{observacio}
Aquesta fórmula és una de les síntesis més netes del llibre: recull en una sola línia tot el capítol de límits. En l'adjunció $\Delta\dashv\varprojlim_{\cat{I}}$, la \emph{counitat} en un diagrama $D$,
\[
\varepsilon_D\colon\Delta\bigl(\varprojlim\nolimits_{\cat{I}}D\bigr)\Rightarrow D,
\]
és exactament el con límit de $D$, és a dir, la família de projeccions; la \emph{unitat} en un objecte $A$ és el morfisme $A\to\varprojlim_{\cat{I}}\Delta A$. Dualment, en l'adjunció $\colimop_{\cat{I}}\dashv\Delta$, la \emph{unitat}
\[
\eta_D\colon D\Rightarrow\Delta\bigl(\colimop\nolimits_{\cat{I}}D\bigr)
\]
és el cocon colímit, la família d'injeccions, mentre que la \emph{counitat} és el morfisme $\colimop_{\cat{I}}\Delta A\to A$. Identificar quina de les dues transformacions és el con o el cocon universal és un bon exercici per fixar la diferència entre unitat i counitat. I la lectura és profètica: quan al capítol de subobjectes vegem l'existencial i l'universal com a adjunts de la reindexació $f^{*}$, estarem veient una versió \enquote{fibrada} d'aquesta mateixa idea ---$f^{*}$ fa el paper de $\Delta$, i els quantificadors, el de $\colimop$ i $\varprojlim$.
\end{observacio}

\section{Exponencials i currificació}

Recuperem ara la relació que vam ajornar en tractar les categories de functors: la que lliga el producte amb els objectes de functors.

\begin{proposicio}\label{prop:cat-currificacio}\index{adjunció!producte-exponencial}\index{Cat@$\cat{Cat}$!exponencials}
Per a categories petites $\cat{A},\cat{B},\cat{C}$, hi ha un isomorfisme natural
\[
\Hom_{\cat{Cat}}(\cat{A}\times\cat{B},\cat{C})\;\cong\;\Hom_{\cat{Cat}}(\cat{A},\cat{C}^{\cat{B}}).
\]
És a dir, el functor $-\times\cat{B}$ és adjunt per l'esquerra del functor $(-)^{\cat{B}}$:
\[
-\times\cat{B}\;\dashv\;(-)^{\cat{B}}.
\]
\end{proposicio}

\begin{prova}
Primer, els dos costats són functors $\cat{Cat}\to\cat{Cat}$. El functor $-\times\cat{B}$ envia $P\colon\cat{A}'\to\cat{A}$ a $P\times\id_{\cat{B}}$. El functor $(-)^{\cat{B}}$ envia $K\colon\cat{C}\to\cat{C}'$ a la composició amb $K$, $\cat{C}^{\cat{B}}\to\cat{C}'^{\cat{B}}$, que fa $H\mapsto K\comp H$ i $\eta\mapsto K\eta$ (capítol~\ref{cap:transformacions}). Si $\cat{B}$ i $\cat{C}$ són petites, $\cat{C}^{\cat{B}}$ també ho és. A $\cat{A}\times\cat{B}$ escrivim els morfismes com a parelles, que es componen component a component, i fem servir la descomposició
\[
(f,g)=(f,\id_{b'})\comp(\id_a,g)=(\id_{a'},g)\comp(f,\id_b)\qquad\text{per a } f\colon a\to a',\ g\colon b\to b'.
\tag{$\ast$}
\]

\emph{Currificació.} Sigui $F\colon\cat{A}\times\cat{B}\to\cat{C}$. Per a cada objecte $a$, $\Lambda F(a)=F(a,-)$ és el functor $\cat{B}\to\cat{C}$ que envia $b\mapsto F(a,b)$ i $g\mapsto F(\id_a,g)$. És un functor perquè $F$ ho és i $(\id_a,g')\comp(\id_a,g)=(\id_a,g'\comp g)$. Per a $f\colon a\to a'$, $\Lambda F(f)\colon F(a,-)\Rightarrow F(a',-)$ té components $F(f,\id_b)$. És natural perquè, aplicant $F$ a $(\ast)$, $F(\id_{a'},g)\comp F(f,\id_b)=F(f,g)=F(f,\id_{b'})\comp F(\id_a,g)$. Finalment, $\Lambda F\colon\cat{A}\to\cat{C}^{\cat{B}}$ és un functor: les components de $\Lambda F(\id_a)$ són identitats, i les de $\Lambda F(f'\comp f)$ són $F(f'\comp f,\id_b)=F(f',\id_b)\comp F(f,\id_b)$, que són les de la composició vertical $\Lambda F(f')\comp\Lambda F(f)$.

\emph{Descurrificació.} Sigui $H\colon\cat{A}\to\cat{C}^{\cat{B}}$. Definim $\bar H\colon\cat{A}\times\cat{B}\to\cat{C}$ per
\[
\bar H(a,b)=H(a)(b),\qquad \bar H(f,g)=H(f)_{b'}\comp H(a)(g)=H(a')(g)\comp H(f)_b,
\]
on la segona igualtat és la naturalitat de $H(f)$. Les identitats es preserven perquè $H(\id_a)$ és la transformació identitat i $H(a)$ és un functor. Per a $(f',g')\colon(a',b')\to(a'',b'')$,
\[
\begin{aligned}
\bar H(f',g')\comp\bar H(f,g)
&=H(f')_{b''}\comp H(a')(g')\comp H(f)_{b'}\comp H(a)(g)\\
&=H(f')_{b''}\comp H(f)_{b''}\comp H(a)(g')\comp H(a)(g)\\
&=H(f'\comp f)_{b''}\comp H(a)(g'\comp g),
\end{aligned}
\]
on hem fet servir la naturalitat de $H(f)$ per a $g'$, que $H$ preserva la composició (vertical, component a component) i que $H(a)$ és un functor. Això és $\bar H\bigl((f',g')\comp(f,g)\bigr)$.

\emph{Són inverses.} Per a $F$: $\overline{\Lambda F}(a,b)=F(a,b)$ i $\overline{\Lambda F}(f,g)=F(f,\id_{b'})\comp F(\id_a,g)=F(f,g)$, per $(\ast)$. Per a $H$: $\Lambda\bar H(a)$ envia $g$ a $\bar H(\id_a,g)=H(a)(g)$, de manera que $\Lambda\bar H(a)=H(a)$; i les components de $\Lambda\bar H(f)$ són $\bar H(f,\id_b)=H(f)_b$, de manera que $\Lambda\bar H(f)=H(f)$.

\emph{Naturalitat.} Per a $P\colon\cat{A}'\to\cat{A}$, el functor $\Lambda\bigl(F\comp(P\times\id_{\cat{B}})\bigr)$ envia $a'$ a $F(P(a'),-)$ i $f'$ a la transformació de components $F(P(f'),\id_b)$; és a dir, és $\Lambda F\comp P$. Per a $K\colon\cat{C}\to\cat{C}'$, $\Lambda(K\comp F)$ envia $a$ a $K\comp F(a,-)$ i $f$ a la transformació de components $K(F(f,\id_b))$, que és $K\,\Lambda F(f)$; és a dir, $\Lambda(K\comp F)=K^{\cat{B}}\comp\Lambda F$, on $K^{\cat{B}}$ és la imatge de $K$ pel functor $(-)^{\cat{B}}$. Aquestes dues igualtats són la naturalitat de $\Lambda$ en $\cat{A}$ i en $\cat{C}$. Per tant, $-\times\cat{B}\dashv(-)^{\cat{B}}$.
\end{prova}

\begin{observacio}
Aquesta adjunció ---\emph{producte $\dashv$ exponencial}--- és el prototip de les \textbf{categories cartesianes tancades}, que estudiarem al capítol següent i que són el marc categòric del càlcul lambda simplement tipat. La \enquote{currificació} és exactament la que apareix en programació funcional: transformar una funció de dues variables en una funció que retorna funcions.
\end{observacio}

\begin{resumcap}
\apartat{Idees centrals}
\begin{enumerate}
  \item Una adjunció $F\dashv G$ és una bijecció natural $\Hom(F(A),B)\cong\Hom(A,G(B))$, i dona els transposats.
  \item Equival a donar una unitat i una counitat que satisfan les identitats triangulars; la unitat expressa una propietat universal, i dualment la counitat.
  \item Exemples clau: les connexions de Galois entre preordres, $\Free\dashv U$ i, a $\cat{Set}$, la cadena $\exists_f\dashv f^{*}\dashv\forall_f$.
  \item $\colimop_{\cat{I}}\dashv\Delta\dashv\varprojlim_{\cat{I}}$ i $-\times B\dashv(-)^{B}$ mostren que moltes propietats universals són adjuncions.
  \item Els adjunts són únics llevat d'isomorfisme natural; els adjunts per la dreta preserven límits, i els adjunts per l'esquerra, colímits.
\end{enumerate}
\apartat{Errors típics} Intercanviar quin adjunt preserva límits i quin colímits; creure que un adjunt per la dreta \emph{no pot} preservar colímits, quan el teorema només dona una direcció; oblidar la naturalitat de la bijecció d'homconjunts; confondre unitat i counitat (per exemple, a $\Delta\dashv\varprojlim$ el con límit és la counitat, no la unitat); i tractar les mnemotècniques (\enquote{RAPL\,/\,LAPC}) com a terminologia establerta en lloc d'ajudes informals.
\apartat{Lectura recomanada} \cite[cap.~9]{awodey}, especialment §§9.4--9.6, per als adjunts d'ordre, els quantificadors com a adjunts i RAPL; \cite[cap.~2]{leinster} per a la unitat, la counitat i les fletxes universals; \cite[§§4.1--4.5]{riehl} per al càlcul d'adjuncions i la preservació de límits; \cite[cap.~4]{perrone-notes} per als exemples lliure--oblit, les connexions de Galois i l'adjunció entre categories i grafs; \cite[cap.~13]{barrwells} per a una exposició orientada a exemples; i \cite[cap.~IV]{maclane} per al tractament clàssic.
\end{resumcap}

\chapter[Cartesianes tancades]{\texorpdfstring{Categories cartesianes\\ tancades}{Categories cartesianes tancades}}
\label{cap:cartesianes}

\begin{objectius}
\apartat{Prerequisits} Representabilitat i lema de Yoneda (cap.~\ref{cap:yoneda}); productes i límits finits (cap.~\ref{cap:limits}); adjuncions i currificació (cap.~\ref{cap:adjuncions}). No cal cap coneixement previ de càlcul lambda: el capítol n'introdueix la semàntica necessària.
\apartat{Objectius} En acabar el capítol, el lector hauria de poder:
\begin{itemize}
  \item definir l'objecte exponencial per la seva propietat universal, com a representació del functor $A\mapsto\Hom(A\times B,C)$, i escriure les dues lleis de la currificació;
  \item reconèixer una categoria cartesiana tancada i l'exponencial com a adjunt per la dreta de $-\times B$;
  \item caracteritzar els ordres parcials cartesians tancats com a semireticles implicatius amb màxim, i relacionar-los amb les àlgebres de Heyting;
  \item situar $\cat{Set}$, els preordres, les categories de prefeixos, $\cat{Graph}$ i $\cat{Top}$ respecte del tancament cartesià;
  \item interpretar el càlcul lambda simplement tipat en una CCC, incloses les lleis $\beta$ i $\eta$ i les regles estructurals.
\end{itemize}
\end{objectius}

En tractar les adjuncions vam trobar, al final del capítol, l'adjunció entre el producte i l'exponencial a $\cat{Cat}$: $-\times\cat{B}\dashv(-)^{\cat{B}}$. Aquella era una manifestació externa d'un fenomen que ara volem fer \emph{intern} a una categoria qualsevol. Una categoria cartesiana en què, per a cada objecte $B$, el functor $-\times B$ té adjunt per la dreta s'anomena \textbf{cartesiana tancada}, i és, exactament, el marc on es pot interpretar el \textbf{càlcul lambda simplement tipat}. Aquest capítol construeix aquesta noció pas a pas, la il·lustra en els exemples fonamentals i n'explica la lectura com a semàntica de tipus i termes. És el primer pont explícit del llibre cap a la lògica categòrica.

Al capítol anterior, les adjuncions van unificar construccions aparentment diverses: objectes lliures, interiors i clausures, connectius lògics i exponencials. El patró comú, una bijecció natural d'homconjunts o, equivalentment, una unitat i una counitat que satisfan les identitats triangulars, reapareixerà com a columna vertebral de tot el que segueix. En aquest capítol el trobem en l'adjunció entre el producte i l'exponencial; més endavant, en les categories amb l'estructura suficient, veurem que els \textbf{quantificadors} $\exists$ i $\forall$ es modelen, respectivament, com els adjunts per l'esquerra i per la dreta dels functors de reindexació entre subobjectes: és la formulació categòrica de la relació entre quantificació i substitució.

\section{Categories cartesianes}

Recordem que un \textbf{producte finit} és el límit d'un diagrama amb un nombre finit d'objectes i cap morfisme no trivial. El cas de zero objectes és l'objecte terminal; el de dos, el producte binari. Una categoria té \textbf{tots els productes finits} si té objecte terminal i producte binari de qualsevol parella, ja que els productes de tres o més objectes s'obtenen iterant.

\begin{definicio}[Categoria cartesiana]\index{categoria!cartesiana}
\label{def:cartesiana}
Una categoria $\cat{C}$ és \textbf{cartesiana} si té objecte terminal $1$ i producte binari $A\times B$ de qualsevol parella d'objectes. Equivalentment, si té tots els productes finits.
\end{definicio}

En una categoria cartesiana disposem, per a cada objecte $A$, del \textbf{morfisme diagonal}
\[
\delta_A=\langle\id_A,\id_A\rangle\colon A\to A\times A,
\]
l'únic morfisme les dues projeccions del qual són la identitat, i del morfisme únic $!_A\colon A\to 1$. Escrivim la diagonal amb minúscula per no confondre-la amb el functor constant $\Delta_A$ (exemple~\ref{ex:functor-constant}). Aquestes dades no són estructura addicional: queden \emph{determinades} per les propietats universals.

Els productes escollits determinen també la functorialitat del producte. Donats $f\colon A\to A'$ i $g\colon B\to B'$, posem
\[
f\times g=\langle f\comp\pi_1,\ g\comp\pi_2\rangle\colon A\times B\to A'\times B'.
\]
En particular, per a $B$ fix, l'assignació $A\mapsto A\times B$ defineix el functor $-\times B$, tal com precisem ara.

\begin{proposicio}[El producte per un objecte és functorial]\index{producte!functorialitat}
\label{prop:producte-functor}
Sigui $\cat{C}$ cartesiana i $B$ un objecte fix. L'assignació $A\mapsto A\times B$ s'estén a un functor
\[
-\times B\colon\cat{C}\to\cat{C},
\]
que envia $f\colon A\to A'$ al morfisme $f\times\id_B=\langle f\comp\pi_1,\pi_2\rangle\colon A\times B\to A'\times B$.
\end{proposicio}

\begin{prova}
Cal comprovar que $-\times B$ respecta identitats i composicions. Per a la identitat, $\id_A\times\id_B=\langle\pi_1,\pi_2\rangle=\id_{A\times B}$ per la unicitat de la propietat universal del producte. Per a la composició, siguin $f\colon A\to A'$ i $g\colon A'\to A''$. Tant $(g\comp f)\times\id_B$ com $(g\times\id_B)\comp(f\times\id_B)$ tenen com a projecció primera $g\comp f\comp\pi_1$ i com a projecció segona $\pi_2$; per unicitat, coincideixen.
\end{prova}

\section{Objectes exponencials}

En una categoria cartesiana, la pregunta clau és si el functor $-\times B$ té adjunt per la dreta. Per a $B$ i $C$ fixos, l'exponencial $C^B$, quan existeix, representa el functor
\[
A\longmapsto\Hom(A\times B,C)
\]
(capítol~\ref{cap:yoneda}); quan existeix per a tot $C$, defineix l'adjunt per la dreta de $-\times B$.

\begin{definicio}[Objecte exponencial]\index{objecte!exponencial}\index{exponencial}
\label{def:exponencial}
Siguin $B,C$ objectes d'una categoria cartesiana $\cat{C}$. Un \textbf{objecte exponencial} $C^{B}$ és un objecte dotat d'un morfisme d'\textbf{avaluació}
\[
\ev\colon C^{B}\times B\to C
\]
amb la propietat universal següent: per a tot objecte $A$ i tot morfisme $f\colon A\times B\to C$, existeix un únic morfisme $\lambda f\colon A\to C^{B}$, la \textbf{transposada} o \textbf{currificació} de $f$, tal que el triangle
\[
\begin{tikzpicture}[baseline=(m.center)]
  \matrix (m) [matrix of math nodes, column sep=3.4em, row sep=2.6em] {
    A\times B & C^{B}\times B \\
              & C \\
  };
  \draw[-{Stealth[scale=1.1]}] (m-1-1) -- node[above] {$\lambda f\times\id_B$} (m-1-2);
  \draw[-{Stealth[scale=1.1]}] (m-1-1) -- node[below left] {$f$} (m-2-2);
  \draw[-{Stealth[scale=1.1]}] (m-1-2) -- node[right] {$\ev$} (m-2-2);
\end{tikzpicture}
\]
commuta, és a dir, $\ev\comp(\lambda f\times\id_B)=f$.
\end{definicio}

L'avaluació formalitza la idea d'\enquote{aplicar una funció al seu argument}, i la currificació, la de \enquote{fixar un paràmetre}. Com tota representació, l'exponencial és únic llevat d'un únic isomorfisme compatible amb l'avaluació: si $(E,\ev)$ i $(E',\ev')$ són exponencials de $C$ per $B$, hi ha un únic isomorfisme $u\colon E\to E'$ amb $\ev'\comp(u\times\id_B)=\ev$.

\begin{observacio}[Les dues lleis de la currificació]\label{obs:lleis-currificacio}\index{currificació!lleis}
La correspondència $f\mapsto\lambda f$ és una bijecció entre $\Hom(A\times B,C)$ i $\Hom(A,C^B)$, amb inversa $g\mapsto\ev\comp(g\times\id_B)$. Les dues composicions donen identitats:
\[
\ev\comp(\lambda f\times\id_B)=f,
\qquad\qquad
\lambda\bigl(\ev\comp(g\times\id_B)\bigr)=g,
\]
per a tot $f\colon A\times B\to C$ i tot $g\colon A\to C^B$. La primera és la condició de la definició. La segona en surt per unicitat: $g$ és un morfisme $A\to C^B$ que fa commutar el triangle per a $f=\ev\comp(g\times\id_B)$, i l'únic que ho fa és $\lambda f$. Aquestes dues lleis són la versió categòrica de les lleis $\beta$ i $\eta$ del càlcul lambda, que veurem més endavant; la naturalitat de la bijecció la converteix en una adjunció.
\end{observacio}

\begin{proposicio}[L'exponencial és un adjunt per la dreta]\index{adjunció!producte-exponencial}
\label{prop:adjuncio-exp}
Si per a un objecte fix $B$ existeix l'exponencial $C^{B}$ per a tot $C$, aleshores l'assignació $C\mapsto C^{B}$ s'estén a un functor $(-)^{B}\colon\cat{C}\to\cat{C}$ i hi ha un isomorfisme natural
\[
\Hom(A\times B,\,C)\;\cong\;\Hom(A,\,C^{B}),
\]
natural en $A$ i en $C$. És a dir, $-\times B\dashv(-)^{B}$.
\end{proposicio}

\begin{prova}
La currificació $f\mapsto\lambda f$ és la bijecció buscada (observació~\ref{obs:lleis-currificacio}). La naturalitat en $A$ es comprova precomponent amb $a\colon A'\to A$ i usant la functorialitat de $-\times B$ (la proposició~\ref{prop:producte-functor}); la naturalitat en $C$ defineix, de fet, l'acció de $(-)^{B}$ sobre morfismes: donat $c\colon C\to C'$, es posa $c^{B}=\lambda\bigl(c\comp\ev\bigr)\colon C^{B}\to C'^{B}$. Que això respecta identitats i composicions se segueix, un cop més, de la unicitat de la transposada.
\end{prova}

\begin{notacio}
Escrivim indistintament $C^{B}$ o $[B,C]$ per a l'exponencial. La transposada de $f$ es denota $\lambda f$ o $\widehat f$; la seva inversa, $\ev\comp(g\times\id_B)$, es denota $\widetilde g$ i s'anomena \textbf{descurrificació} de $g$.
\end{notacio}

\section{La definició de categoria cartesiana tancada}

\begin{definicio}[Categoria cartesiana tancada]\index{categoria!cartesiana tancada}\index{CCC}
\label{def:ccc}
Una categoria $\cat{C}$ és \textbf{cartesiana tancada} (abreujat \textbf{CCC}, de l'anglès \emph{cartesian closed category}) si és cartesiana i, a més, per a cada objecte $B$, el functor $-\times B\colon\cat{C}\to\cat{C}$ té adjunt per la dreta $(-)^{B}$. Explícitament, $\cat{C}$ té:
\begin{enumerate}
\item un objecte terminal $1$;
\item productes binaris $A\times B$;
\item exponencials $C^{B}$ amb avaluació i currificació.
\end{enumerate}
\end{definicio}

\begin{observacio}[L'exponencial com a bifunctor]\label{obs:exponencial-bifunctor}\index{exponencial!bifunctor}
En una CCC, els exponencials s'organitzen en un bifunctor
\[
(-)^{(-)}\colon\cat{C}\op\times\cat{C}\longrightarrow\cat{C},
\qquad (B,C)\longmapsto C^{B},
\]
contravariant en $B$ i covariant en $C$, i la bijecció
\[
\Hom(A\times B,C)\;\cong\;\Hom(A,C^{B})
\]
és natural en les tres variables $A$, $B$ i $C$. L'acció en $C$ és la de la proposició~\ref{prop:adjuncio-exp}; la construcció de l'acció en $B$, mitjançant la transposada de $\ev\comp(\id\times f)$, es proposa als Exercicis.
\end{observacio}

La condició es pot resumir dient que en una CCC les tres operacions \enquote{no fer res} (terminal), \enquote{aparellar} (producte) i \enquote{formar espais de funcions} (exponencial) hi són totes i estan lligades per adjuncions. Aquesta és exactament l'estructura que necessita un llenguatge de tipus amb producte de tipus i tipus funció.

\begin{observacio}
La tancadura és una propietat \emph{objecte a objecte}: cal l'exponencial per a \emph{cada} base $B$. Una categoria pot tenir alguns exponencials i no ser tancada. La demanda uniforme ---un adjunt per la dreta de $-\times B$ per a tot $B$--- és el que fa que la currificació estigui sempre disponible.
\end{observacio}

\section{Exemples fonamentals}

\begin{exemple}[$\cat{Set}$]\index{Set@$\cat{Set}$!cartesiana tancada}
La categoria $\cat{Set}$ és cartesiana tancada. L'objecte terminal és qualsevol conjunt d'un element; el producte és el producte cartesià; i l'exponencial $C^{B}$ és el conjunt de \emph{totes} les aplicacions $B\to C$, amb l'avaluació $\ev(g,b)=g(b)$. La currificació $\lambda f\colon A\to C^{B}$ d'una aplicació $f\colon A\times B\to C$ és $\lambda f(a)=(b\mapsto f(a,b))$: exactament la currificació de la programació funcional. Aquest és el model prototípic i el que dona nom a tota la construcció.
\end{exemple}

\begin{exemple}[$\cat{Cat}$: transformacions naturals com a functors]\label{ex:cat-ccc}\index{Cat@$\cat{Cat}$!cartesiana tancada}\index{avaluació!a Cat@a $\cat{Cat}$}
La categoria $\cat{Cat}$ de les categories petites és cartesiana tancada: el terminal és $\cat{1}$, el producte és la categoria producte i l'exponencial de $\cat{C}$ per $\cat{B}$ és la categoria de functors $\cat{C}^{\cat{B}}$ (proposició~\ref{prop:cat-currificacio}). Fem explícits l'avaluació i la currificació, i després els llegim en un cas concret.

\emph{Avaluació.} El functor $\ev\colon\cat{C}^{\cat{B}}\times\cat{B}\to\cat{C}$ és
\[
\ev(H,b)=H(b),\qquad \ev(\eta,g)=\eta_{b'}\comp H(g)=H'(g)\comp\eta_b,
\]
per a $\eta\colon H\Rightarrow H'$ i $g\colon b\to b'$; la segona igualtat és la naturalitat de $\eta$. És la descurrificació de $\id_{\cat{C}^{\cat{B}}}$ a la demostració de la proposició~\ref{prop:cat-currificacio}, i d'allà en surt que és un functor.

\emph{Currificació i llei $\beta$.} La currificació de $F\colon\cat{A}\times\cat{B}\to\cat{C}$ és el functor $\Lambda F\colon\cat{A}\to\cat{C}^{\cat{B}}$ d'aquella demostració. El triangle de la definició~\ref{def:exponencial} commuta: sobre objectes, $\ev\bigl(\Lambda F(a),b\bigr)=F(a,b)$; sobre un morfisme $(f,g)$,
\[
\ev\bigl(\Lambda F(f),g\bigr)=\Lambda F(f)_{b'}\comp\Lambda F(a)(g)=F(f,\id_{b'})\comp F(\id_a,g)=F(f,g).
\]
Així, $\ev\comp(\Lambda F\times\id_{\cat{B}})=F$: avaluar la currificació torna la funció de dues variables.

\emph{Un cas concret: $\cat{B}=\cat{2}$.} La categoria $\cat{C}^{\cat{2}}$ és la categoria de fletxes de $\cat{C}$: els objectes són els morfismes de $\cat{C}$ i els morfismes, els quadrats commutatius (Pràctica del capítol~\ref{cap:transformacions}). L'avaluació envia $(h,0)$ al domini de $h$, $(h,1)$ al codomini i $(h,u)$ a $h$ mateix, on $u\colon0\to1$. Un functor $\Phi\colon\cat{A}\times\cat{2}\to\cat{C}$ consisteix en tres dades:
\begin{itemize}
\item dos functors $F=\Phi(-,0)$ i $G=\Phi(-,1)$ de $\cat{A}$ a $\cat{C}$;
\item una família $\alpha_a=\Phi(\id_a,u)\colon F(a)\to G(a)$;
\item la naturalitat $G(f)\comp\alpha_a=\alpha_{a'}\comp F(f)$, que no és cap condició nova: és $\Phi$ aplicat a les dues descomposicions de $(f,u)$.
\end{itemize}
És a dir, un functor $\cat{A}\times\cat{2}\to\cat{C}$ \emph{és} una transformació natural $\alpha\colon F\Rightarrow G$. La seva currificació $\Lambda\Phi\colon\cat{A}\to\cat{C}^{\cat{2}}$ envia $a$ a la fletxa $\alpha_a$ i $f$ al quadrat de naturalitat de $f$. L'avaluació recupera les tres peces: $\ev(\Lambda\Phi(a),0)=F(a)$, $\ev(\Lambda\Phi(a),1)=G(a)$ i $\ev(\Lambda\Phi(a),u)=\alpha_a$. Tenim, doncs, tres lectures d'una mateixa dada,
\[
\alpha\colon F\Rightarrow G\qquad\longleftrightarrow\qquad\cat{A}\times\cat{2}\to\cat{C}\qquad\longleftrightarrow\qquad\cat{A}\to\cat{C}^{\cat{2}},
\]
i el pas de la segona a la tercera és la currificació, amb l'avaluació com a inversa.
\end{exemple}

\begin{exemple}[Preordres cartesians tancats]\index{semireticle implicatiu}\index{àlgebra de Heyting}\index{Heyting}
Un preordre amb ínfims finits, vist com a categoria, és cartesià: el terminal és l'element màxim $\top$ i el producte és l'ínfim $x\wedge y$. Té exponencials si i només si, per a cada $b$, l'operació $-\wedge b$ té adjunt per la dreta; aquest adjunt és la \textbf{implicació} $b\Rightarrow-$, caracteritzada per
\[
a\wedge b\leq c\quad\Longleftrightarrow\quad a\leq(b\Rightarrow c).
\]
Un ordre parcial amb ínfims finits és, doncs, cartesià, i és cartesià tancat si i només si és un \textbf{semireticle implicatiu amb màxim}: un $\wedge$-semireticle amb element màxim $\top$ i una operació d'implicació que satisfà l'equivalència anterior (també anomenat \emph{semireticle de Brouwer}). Per a un preordre, la mateixa caracterització val després de passar al quocient per l'equivalència
\[
a\sim b\quad\Longleftrightarrow\quad a\leq b\ \text{i}\ b\leq a,
\]
perquè en la terminologia algebraica habitual un semireticle és un ordre parcial. Aquesta estructura captura la part $(\top,\wedge,\Rightarrow)$ de la lògica.

Convé no confondre aquesta estructura amb una àlgebra de Heyting. Una \textbf{àlgebra de Heyting}\index{àlgebra de Heyting} és un \emph{reticle distributiu acotat} ---amb ínfims finits $\top,\wedge$ i \emph{suprems finits} $\bot,\vee$--- proveït d'una implicació que satisfà l'equivalència anterior. El tancament cartesià per si sol només dona la part multiplicativa: no garanteix ni els joins $\vee$ ni el mínim $\bot$. Reservem, doncs, el nom \emph{àlgebra de Heyting} per al cas amb suprems finits; un semireticle implicatiu amb màxim que a més tingui suprems finits (i sigui distributiu, cosa automàtica en presència de la implicació) és una àlgebra de Heyting. Aquest exemple és decisiu per a la lògica: la implicació intuïcionista \emph{és} l'exponencial, i el tancament cartesià és la traducció algebraica de disposar del connectiu $\to$, mentre que la disjunció $\vee$ correspon als suprems.
\end{exemple}

\begin{observacio}[Atenció: cartesià tancat no vol dir àlgebra de Heyting]\label{obs:ccc-no-heyting}\index{àlgebra de Heyting!i preordres cartesians tancats}
En un preordre, ser cartesià tancat dona $\top$, $\wedge$ i $\Rightarrow$, i res més. Un exemple mínim mostra què pot faltar. Sigui $\mathbb Z\cup\{+\infty\}$ amb l'ordre usual. És una cadena amb màxim $+\infty$, amb ínfims $a\wedge b=\min(a,b)$, i té implicació
\[
b\Rightarrow c=\begin{cases}+\infty&\text{si } b\leq c,\\ c&\text{si } b>c.\end{cases}
\]
En efecte, si $b\leq c$ tenim $a\wedge b\leq c$ per a tot $a$; si $b>c$, tenim $\min(a,b)\leq c$ si i només si $a\leq c$. Per tant és un ordre parcial cartesià tancat. Però no té element mínim, de manera que no és una àlgebra de Heyting. Aquí no existeixen ni $\bot$, ni la negació $\neg b=(b\Rightarrow\bot)$, ni el principi \emph{ex falso}; no hi ha, doncs, cap interpretació de la falsedat. En altres exemples pot faltar el suprem binari $\vee$, i amb ell la disjunció.

La relació exacta és aquesta: una àlgebra de Heyting és un ordre parcial cartesià tancat que, a més, té coproductes finits ($\bot$ i $\vee$). La distributivitat ve aleshores gratis, perquè $-\wedge b$, com a adjunt per l'esquerra, preserva els suprems (capítol~\ref{cap:adjuncions}). L'estructura de Heyting també condiciona els morfismes: un morfisme d'àlgebres de Heyting ha de preservar $\bot$ i $\vee$, cosa que no demana un functor que només preservi l'estructura cartesiana tancada. En lògica, el tancament cartesià modela el fragment $(\top,\wedge,\to)$; la falsedat i la disjunció són estructura addicional.
\end{observacio}

\begin{exemple}[Categories de prefeixos]\index{prefeix@prefeix!cartesiana tancada}
Per a qualsevol categoria petita $\cat{C}$, la categoria de prefeixos $\widehat{\cat{C}}=\cat{Set}^{\cat{C}\op}$ (exemple~\ref{ex:categoria-prefeixos}) és cartesiana tancada. Els productes i el terminal es calculen \emph{objecte a objecte} a $\cat{Set}$ (proposició~\ref{prop:limits-punt-a-punt}), i l'exponencial es defineix, mitjançant el lema de Yoneda, per $Q^{P}(A)=\Nat(\Hom(-,A)\times P,\,Q)$. Aquestes categories seran el terreny natural dels topos de prefeixos.
\end{exemple}

\begin{exemple}[Grafs]\index{graf dirigit!cartesiana tancada}
Al capítol~\ref{cap:transformacions} vam veure que $\cat{Graph}\cong\cat{Set}^{\cat{P}}$, on $\cat{P}$ té dos objectes $E,V$ i dues fletxes $s,t\colon E\to V$. Com que $\cat{Set}^{\cat{P}}=\cat{Set}^{(\cat{P}\op)\op}$, els grafs són exactament els prefeixos sobre $\cat{P}\op$: un prefeix $G\colon(\cat{P}\op)\op\to\cat{Set}$ és la dada de dos conjunts, de vèrtexs i d'arestes, i de les aplicacions origen i destí. Per tant, com a cas particular de l'exemple anterior, $\cat{Graph}$ és cartesiana tancada. No calcularem aquí l'exponencial de dos grafs; n'hi ha prou de saber que existeix i que es pot obtenir amb la fórmula de Yoneda de l'exemple anterior.
\end{exemple}

\begin{observacio}[Un no-exemple instructiu]
La categoria $\cat{Top}$ dels espais topològics \emph{no} és cartesiana tancada: no sempre existeix una topologia sobre l'espai de funcions contínues que faci de l'avaluació un morfisme universal. Aquesta mancança va motivar històricament la introducció de categories \enquote{convenients} d'espais topològics, com la dels espais Hausdorff compactament generats, amb les construccions categòriques adequades, que sí que és cartesiana tancada. El tancament cartesià és, doncs, una propietat delicada i valuosa, no automàtica.
\end{observacio}

\section[Semàntica del càlcul lambda simplement tipat]{Semàntica del càlcul lambda\\ simplement tipat}

El motiu pel qual les CCC són centrals en lògica categòrica és la seva correspondència amb el \textbf{càlcul lambda simplement tipat}. Descrivim la traducció a nivell introductori; és una de les formes de l'anomenada \emph{correspondència de Curry--Howard--Lambek}.

\subsection{Tipus com a objectes}

Fixem un conjunt de tipus base. Els \textbf{tipus} es construeixen amb dues operacions: el \textbf{producte de tipus} $\sigma\times\tau$ i el \textbf{tipus funció} $\sigma\to\tau$, més un tipus unitat $\mathbf{1}$. En una CCC $\cat{C}$, la interpretació $\llbracket-\rrbracket$ assigna:
\[
\llbracket\mathbf{1}\rrbracket=1,\qquad
\llbracket\sigma\times\tau\rrbracket=\llbracket\sigma\rrbracket\times\llbracket\tau\rrbracket,\qquad
\llbracket\sigma\to\tau\rrbracket=\llbracket\tau\rrbracket^{\llbracket\sigma\rrbracket}.
\]
Cada tipus esdevé un objecte; el tipus funció esdevé un exponencial.

\subsection{Termes com a morfismes}

Un \textbf{context} $\Gamma=x_1:\sigma_1,\dots,x_n:\sigma_n$ s'interpreta, per recursió, com un producte:
\[
\llbracket\varnothing\rrbracket=1,
\qquad
\llbracket\Gamma,x:\sigma\rrbracket=\llbracket\Gamma\rrbracket\times\llbracket\sigma\rrbracket.
\]
Un \textbf{terme} $\Gamma\vdash t:\tau$, amb variables lliures en el context $\Gamma$, s'interpreta com un morfisme
\[
\llbracket t\rrbracket\colon\llbracket\Gamma\rrbracket\longrightarrow\llbracket\tau\rrbracket.
\]
Les regles de formació de termes es corresponen amb les operacions de la CCC:
\begin{itemize}
\item una \textbf{variable} $x_i$ s'interpreta amb la projecció $\pi_i$;
\item una \textbf{parella} $\langle t,u\rangle$, amb el morfisme aparellament $\langle\llbracket t\rrbracket,\llbracket u\rrbracket\rangle$ cap a un producte;
\item una \textbf{aplicació} $t\,u$, amb la composició de l'avaluació i l'aparellament de $\llbracket t\rrbracket$ i $\llbracket u\rrbracket$;
\item una \textbf{abstracció} $\lambda x{:}\sigma.\,t$, amb la currificació $\lambda\llbracket t\rrbracket$: si $\Gamma,x:\sigma\vdash t:\tau$, aleshores $\llbracket t\rrbracket\colon\llbracket\Gamma\rrbracket\times\llbracket\sigma\rrbracket\to\llbracket\tau\rrbracket$, i la seva transposada és un morfisme $\llbracket\Gamma\rrbracket\to\llbracket\tau\rrbracket^{\llbracket\sigma\rrbracket}$, exactament el tipus de $\lambda x{:}\sigma.\,t$.
\end{itemize}
La correspondència entre abstracció i currificació és el cor de la traducció: lligar una variable és transposar per l'adjunció, i aplicar una funció és avaluar. Les regles de conversió del càlcul es tornen igualtats de morfismes:
\[
(\beta)\quad(\lambda x.\,t)\,u = t[u/x]
\qquad\text{esdevé}\qquad
\ev\comp\langle\lambda\llbracket t\rrbracket,\llbracket u\rrbracket\rangle=\llbracket t\rrbracket\comp\langle\id,\llbracket u\rrbracket\rangle,
\]
que no és literalment l'equació de la definició~\ref{def:exponencial}, sinó que \emph{se'n dedueix}. Si $f=\llbracket t\rrbracket\colon\llbracket\Gamma\rrbracket\times\llbracket\sigma\rrbracket\to\llbracket\tau\rrbracket$, la primera llei de la currificació (observació~\ref{obs:lleis-currificacio}) dona $\ev\comp(\lambda f\times\id)=f$, i precomponent amb $\langle\id_{\llbracket\Gamma\rrbracket},\llbracket u\rrbracket\rangle\colon\llbracket\Gamma\rrbracket\to\llbracket\Gamma\rrbracket\times\llbracket\sigma\rrbracket$ s'obté
\[
\ev\comp\langle\lambda f,\llbracket u\rrbracket\rangle
=\ev\comp(\lambda f\times\id)\comp\langle\id,\llbracket u\rrbracket\rangle
=f\comp\langle\id,\llbracket u\rrbracket\rangle,
\]
que és la interpretació de la substitució $t[u/x]$. La regla $(\eta)$, $\lambda x.\,(t\,x)=t$, es correspon amb la segona llei, és a dir, amb la unicitat de la transposada.

\begin{observacio}[L'estructura cartesiana i les regles estructurals]\index{regles estructurals}
Els morfismes canònics del producte interpreten les \emph{regles estructurals} de la lògica i del càlcul. Una projecció $\llbracket\Gamma\rrbracket\times\llbracket\sigma\rrbracket\to\llbracket\Gamma\rrbracket$ permet \emph{oblidar} una variable (debilitament); la diagonal $\delta\colon\llbracket\sigma\rrbracket\to\llbracket\sigma\rrbracket\times\llbracket\sigma\rrbracket$ permet \emph{reutilitzar-la} (contracció); i l'isomorfisme de simetria $\llbracket\sigma\rrbracket\times\llbracket\tau\rrbracket\cong\llbracket\tau\rrbracket\times\llbracket\sigma\rrbracket$ permet \emph{intercanviar-ne} l'ordre (intercanvi). És per això que la diagonal, introduïda al començament del capítol, no és un detall aïllat: l'estructura \emph{cartesiana} és exactament la que fa admissibles aquestes tres regles.
\end{observacio}

\begin{observacio}[Curry--Howard--Lambek]\index{Curry--Howard--Lambek}
La correspondència es pot fer precisa: després de fixar una noció adequada de teoria del càlcul lambda simplement tipat, una estructura cartesiana tancada escollida i els morfismes que la preserven, s'obté una \emph{equivalència} de categories entre les dues presentacions. En aquest sentit, tipus, termes i equacions del càlcul lambda corresponen a objectes, morfismes i igualtats de morfismes en una CCC. Aquesta és la primera de les grans traduccions \enquote{lògica $=$ categoria} que estructuraran l'estudi posterior de la lògica categòrica; les següents afegiran, sobre aquesta base, els subobjectes per als predicats i la reindexació per a la quantificació.
\end{observacio}

\begin{resumcap}
\apartat{Idees centrals}
\begin{enumerate}
  \item L'objecte exponencial $C^{B}$ representa el functor $A\mapsto\Hom(A\times B,C)$; la currificació és la bijecció corresponent, amb les dues lleis $\ev\comp(\lambda f\times\id_B)=f$ i $\lambda(\ev\comp(g\times\id_B))=g$.
  \item Una categoria cartesiana tancada té productes finits i tots els exponencials; els exponencials s'organitzen en un bifunctor contravariant en l'exponent i covariant en la base.
  \item Un ordre parcial cartesià tancat és un semireticle implicatiu amb màxim; per a un preordre, aquesta descripció val després de passar al quocient per l'equivalència $a\sim b$. Si, a més, hi ha suprems finits, s'obté una àlgebra de Heyting; sense ells, no (per exemple, $\mathbb Z\cup\{+\infty\}$ no té $\bot$).
  \item $\cat{Set}$, $\cat{Cat}$, les categories de prefeixos i, en particular, $\cat{Graph}$ són cartesianes tancades. A $\cat{Cat}$, l'avaluació i la currificació identifiquen les transformacions naturals amb els functors $\cat{A}\times\cat{2}\to\cat{C}$ i $\cat{A}\to\cat{C}^{\cat{2}}$; $\cat{Top}$ no ho és en general.
  \item El càlcul lambda simplement tipat s'interpreta en tota CCC: els contextos com a productes, l'abstracció com a currificació, i les lleis $\beta$ i $\eta$ com a conseqüències de les dues lleis de la currificació.
\end{enumerate}
\apartat{Errors típics} Identificar tota categoria cartesiana tancada amb una àlgebra de Heyting, fins i tot quan és un ordre parcial (poden faltar $\bot$ i els joins, i amb ells la negació i la disjunció); atribuir la functorialitat del producte a la diagonal en lloc de a la seva propietat universal; i tractar la correspondència de Curry--Howard--Lambek com una igualtat en lloc d'una equivalència amb estructura triada.
\apartat{Lectura recomanada} \cite[cap.~6]{awodey} per a exponencials i CCC; \cite{lambekscott} per a la correspondència amb el càlcul lambda; \cite[cap.~1]{maclanemoerdijk} per als exemples de prefeixos; \cite[cap.~6]{barrwells}, que presenta les CCC i el càlcul lambda tipat i calcula explícitament l'exponencial de dos grafs (exemple~6.1.12); i \cite[§4.3]{riehl} per a la formulació de l'exponencial com a adjunt; \cite[§11]{streicher}, per als exponencials en categories de prefeixos (§11.1) i la semàntica del càlcul lambda tipat (§11.2); i \cite[article~V, sessions~30--31]{lawvereschanuel}, per a una introducció molt elemental als objectes de morfismes.
\end{resumcap}

\chapter[Subobjectes, imatges i factoritzacions]{\texorpdfstring{Subobjectes, imatges\\ i factoritzacions}{Subobjectes, imatges i factoritzacions}}
\label{cap:subobjectes}

\begin{objectius}
\apartat{Prerequisits} Monomorfismes i epimorfismes (cap.~\ref{cap:categories}); pullbacks, coequalitzadors, límits finits i reindexació per pullback (cap.~\ref{cap:limits}); adjuncions, especialment entre preordres (cap.~\ref{cap:adjuncions}).
\apartat{Objectius} En acabar el capítol, el lector hauria de poder:
\begin{itemize}
  \item definir els subobjectes i entendre $\Sub(A)$ com un ordre parcial de predicats sobre $A$;
  \item interpretar la reindexació $f^{*}$ com a substitució i organitzar-la, en una categoria ben potenciada, en un functor contravariant $\Sub\colon\cat{C}\op\to\cat{Poset}$;
  \item interpretar la intersecció de subobjectes com a pullback;
  \item definir epimorfisme regular, parella nucli i factorització imatge;
  \item reconèixer una categoria regular;
  \item demostrar l'adjunció $\exists_f\dashv f^{*}$.
\end{itemize}
\end{objectius}

El capítol anterior ens ha donat la semàntica dels tipus i els termes. Ara construïm la peça complementària: la interpretació de les \textbf{proposicions} i la \textbf{substitució}. La idea directriu és senzilla i profunda: un predicat sobre un objecte $A$ ---una propietat que els \enquote{elements} de $A$ poden tenir o no--- es representa categòricament per un \textbf{subobjecte} de $A$. A $\cat{Set}$, un subconjunt es representa per la seva inclusió, que és injectiva; en una categoria arbitrària, el paper de les injeccions el fan els monomorfismes. Els subobjectes d'un objecte formen un ordre parcial $\Sub(A)$, i la reindexació per pullback, que ja vam construir, hi actua com la substitució.

A la segona meitat del capítol construïm la \textbf{factorització imatge} d'un morfisme. A $\cat{Set}$, tota aplicació es factoritza com una aplicació exhaustiva sobre la seva imatge seguida de la inclusió d'aquesta imatge; en una categoria arbitrària, aquesta idea es reformula mitjançant epimorfismes regulars i monomorfismes. Després definim les \textbf{categories regulars} com aquelles on aquesta factorització és estable per canvi de base, i n'obtenim una conseqüència notable: l'\textbf{existencial} apareix com a \textbf{adjunt per l'esquerra} de la reindexació. L'universal $\forall_f$ requereix estructura addicional i es tractarà al volum de lògica categòrica. Amb això, el volum tanca la infraestructura categòrica que la lògica farà servir.

\section{Subobjectes}

Volem capturar la idea de \enquote{subconjunt} de manera purament categòrica. A $\cat{Set}$, un subconjunt $S\subseteq A$ es pot presentar per la seva injecció $S\hookrightarrow A$, que és un monomorfisme. Però dos monomorfismes diferents poden tenir la mateixa imatge: una injecció i una reparametrització bijectiva del domini defineixen \enquote{el mateix} subconjunt. Cal, doncs, identificar monomorfismes que difereixen per un isomorfisme del domini.

\begin{definicio}[Subobjecte]\index{subobjecte}
\label{def:subobjecte}
Sigui $A$ un objecte d'una categoria $\cat{C}$. Considerem els monomorfismes amb codomini $A$. Diem que dos monomorfismes $m\colon S\rightarrowtail A$ i $m'\colon S'\rightarrowtail A$ són \textbf{equivalents}, i escrivim $m\sim m'$, si existeix un isomorfisme $\varphi\colon S\to S'$ amb $m'\comp\varphi=m$:
\[
\begin{tikzpicture}[baseline=(m.center)]
  \matrix (m) [matrix of math nodes, column sep=2.8em, row sep=2.0em] {
    S & & S' \\
      & A & \\
  };
  \draw[-{Stealth[scale=1.1]}] (m-1-1) -- node[above] {$\varphi$} node[below] {$\sim$} (m-1-3);
  \draw[-{Stealth[scale=1.1]}] (m-1-1) -- node[below left] {$m$} (m-2-2);
  \draw[-{Stealth[scale=1.1]}] (m-1-3) -- node[below right] {$m'$} (m-2-2);
\end{tikzpicture}
\]
Un \textbf{subobjecte} de $A$ és una classe d'equivalència de monomorfismes per aquesta relació. Escrivim $[m]$ per al subobjecte representat per $m$.
\end{definicio}

\begin{observacio}
La relació $\sim$ és d'equivalència: la reflexivitat usa $\varphi=\id$, la simetria usa $\varphi^{-1}$ (que existeix perquè $\varphi$ és isomorfisme), i la transitivitat, la composició d'isomorfismes. Cal notar que l'isomorfisme $\varphi$, si existeix, és \emph{únic}: de $m'\comp\varphi=m=m'\comp\varphi'$ i el fet que $m'$ és mono se segueix $\varphi=\varphi'$. Els subobjectes estan, doncs, ben definits sense ambigüitat.
\end{observacio}

\begin{exemple}
A $\cat{Set}$, els subobjectes de $A$ es corresponen exactament amb els \textbf{subconjunts} de $A$. Un monomorfisme $m\colon S\rightarrowtail A$ és una aplicació injectiva, i la seva classe d'equivalència queda determinada per la imatge $m(S)\subseteq A$: dos monos són equivalents si i només si tenen la mateixa imatge. Així, $\Sub(A)\cong\mathcal{P}(A)$, el conjunt de parts de $A$.
\end{exemple}

\section{L'ordre de subobjectes}

Els subobjectes d'un objecte no formen només un conjunt: hi ha una noció natural d'inclusió entre ells.

\begin{definicio}[Ordre entre subobjectes]\index{subobjecte!ordre}\index{Sub@$\Sub(A)$}
\label{def:ordre-sub}
Donats dos monomorfismes $m\colon S\rightarrowtail A$ i $n\colon T\rightarrowtail A$, diem que $[m]\leq[n]$ si $m$ \textbf{es factoritza a través de} $n$, és a dir, si existeix un morfisme $h\colon S\to T$ amb $n\comp h=m$:
\[
\begin{tikzpicture}[baseline=(m.center)]
  \matrix (m) [matrix of math nodes, column sep=2.8em, row sep=2.0em] {
    S & & T \\
      & A & \\
  };
  \draw[-{Stealth[scale=1.1]}] (m-1-1) -- node[above] {$h$} (m-1-3);
  \draw[-{Stealth[scale=1.1]}] (m-1-1) -- node[below left] {$m$} (m-2-2);
  \draw[-{Stealth[scale=1.1]}] (m-1-3) -- node[below right] {$n$} (m-2-2);
\end{tikzpicture}
\]
\end{definicio}

\begin{proposicio}[$\Sub(A)$ és un ordre parcial]
\label{prop:sub-preordre}
La relació $\leq$ de la definició~\ref{def:ordre-sub} està ben definida sobre classes d'equivalència i és reflexiva i transitiva. A més, és antisimètrica: $[m]\leq[n]$ i $[n]\leq[m]$ impliquen $[m]=[n]$. Per tant, $\Sub(A)$ és un \textbf{ordre parcial}. (Sobre els monomorfismes mateixos, abans de passar a classes d'equivalència, la factorització només defineix un preordre.)
\end{proposicio}

\begin{prova}
El morfisme de factorització $h$, si existeix, és únic, perquè $n$ és mono: de $n\comp h=m=n\comp h'$ se segueix $h=h'$. A més, $h$ és automàticament un monomorfisme: si $h\comp u=h\comp v$, aleshores $m\comp u=n\comp h\comp u=n\comp h\comp v=m\comp v$, i com que $m$ és mono, $u=v$. La bona definició sobre classes és immediata: substituir $m$ o $n$ per equivalents compon $h$ amb isomorfismes. La reflexivitat usa $h=\id$; la transitivitat, la composició de factoritzacions.

Per a l'antisimetria, suposem $[m]\leq[n]$ i $[n]\leq[m]$, amb factoritzacions $h\colon S\to T$ i $k\colon T\to S$. Aleshores $n\comp h=m$ i $m\comp k=n$, d'on $n\comp h\comp k=m\comp k=n$; com que $n$ és mono, $h\comp k=\id_T$. Simètricament, $k\comp h=\id_S$. Per tant $h$ és un isomorfisme amb $n\comp h=m$, és a dir, $m\sim n$ i $[m]=[n]$.
\end{prova}

\begin{observacio}[Categories ben potenciades]\index{categoria!ben potenciada}\index{ben potenciada}
La petitesa local garanteix que cada $\Hom(A,B)$ és un conjunt, però \emph{no} que la col·lecció $\Sub(A)$ ho sigui: en principi podria ser massa gran. Direm que una categoria és \textbf{ben potenciada} si, per a cada objecte $A$, els subobjectes de $A$ formen un conjunt. En el marc de classes que adoptem, la formulació precisa demana un conjunt de monomorfismes representants, perquè cada classe d'equivalència pot ser una classe pròpia (apèndix~\ref{cap:mida}). Aquesta és la hipòtesi que garanteix que $\Sub(A)$ és un conjunt de debò i que, més endavant, es pugui reunir en un prefeix
\[
\Sub\colon\cat{C}\op\longrightarrow\cat{Set}.
\]
Totes les categories que trobarem seran ben potenciades; a $\cat{Set}$, per exemple, $\Sub(A)\cong\mathcal{P}(A)$ és un conjunt. De fet, l'existència d'un classificador de subobjectes en una categoria localment petita ja implica aquesta condició.
\end{observacio}

\begin{observacio}[Predicats com a subobjectes]
Aquesta és la lectura lògica fonamental. Pensem un objecte $A$ com un \textbf{tipus} i un subobjecte $[m]\in\Sub(A)$ com un \textbf{predicat} sobre aquell tipus: la propietat \enquote{pertànyer a $S$}. L'ordre $[m]\leq[n]$ és la \textbf{implicació} entre predicats: tot element que satisfà $m$ satisfà $n$. Veurem que $\Sub(A)$ té sovint estructura de reticle ---amb intersecció i unió de subobjectes--- i fins i tot d'àlgebra de Heyting, cosa que en fa l'àlgebra de proposicions sobre $A$.
\end{observacio}

\section{Reindexació de subobjectes}

Al capítol de límits vam construir, per a un morfisme $f\colon A\to B$ en una categoria amb pullbacks, el functor de reindexació $f^{*}\colon\cat{C}/B\to\cat{C}/A$ (observació~\ref{obs:reindexacio}). La clau ara és que aquesta reindexació \emph{preserva els monomorfismes}, de manera que es restringeix a una operació entre ordres parcials de subobjectes.

\begin{proposicio}[El pullback preserva monomorfismes]
\label{prop:pullback-mono}
En un quadrat cartesià
\[
\begin{tikzpicture}[baseline=(m.center)]
  \matrix (m) [matrix of math nodes, column sep=3em, row sep=2.4em] {
    P & S \\
    A & B \\
  };
  \draw[-{Stealth[scale=1.1]}] (m-1-1) -- node[above] {$f'$} (m-1-2);
  \draw[-{Stealth[scale=1.1]}] (m-1-1) -- node[left] {$m'$} (m-2-1);
  \draw[-{Stealth[scale=1.1]}] (m-1-2) -- node[right] {$m$} (m-2-2);
  \draw[-{Stealth[scale=1.1]}] (m-2-1) -- node[below] {$f$} (m-2-2);
  \node at ([xshift=0.8em,yshift=-0.7em]m-1-1) {$\scriptstyle\lrcorner$};
\end{tikzpicture}
\qquad
m\comp f'=f\comp m',
\]
si $m$ és un monomorfisme, aleshores el seu pullback $m'\colon P\to A$ al llarg de $f$ també ho és.
\end{proposicio}

\begin{prova}
Siguin $u,v\colon Z\to P$ amb $m'\comp u=m'\comp v$. Per la commutativitat del quadrat,
\[
m\comp f'\comp u=f\comp m'\comp u=f\comp m'\comp v=m\comp f'\comp v,
\]
i, com que $m$ és mono, $f'\comp u=f'\comp v$. Així, $u$ i $v$ tenen les mateixes composicions amb les dues projeccions del pullback, $m'$ i $f'$, i per la unicitat de la propietat universal del pullback, $u=v$. Per tant $m'$ és un monomorfisme.
\end{prova}

\begin{observacio}
La conseqüència que ens interessa és immediata: com que la reindexació $f^{*}$ envia monomorfismes a monomorfismes, es restringeix a una operació ben definida entre els ordres parcials de subobjectes, que és la que ara definim. A $\cat{Set}$, aquesta operació és l'antiimatge $S\mapsto f^{-1}(S)$, tal com vam veure en introduir el pullback.
\end{observacio}

\begin{definicio}[Functor de reindexació sobre subobjectes]\index{reindexació!de subobjectes}\index{substitució}
\label{def:reindexacio-sub}
Sigui $f\colon A\to B$ un morfisme en una categoria $\cat{C}$ amb pullbacks. La \textbf{reindexació} indueix una aplicació
\[
f^{*}\colon\Sub(B)\longrightarrow\Sub(A)
\]
que envia un subobjecte $[m]$, amb $m\colon S\rightarrowtail B$, al subobjecte $[m']$ del seu pullback al llarg de $f$. Per la proposició~\ref{prop:pullback-mono}, $m'$ és mono, de manera que $f^{*}[m]$ és un subobjecte de $A$ ben definit.
\end{definicio}

\begin{proposicio}[La reindexació és monòtona]
\label{prop:reindexacio-monotona}
L'aplicació $f^{*}\colon\Sub(B)\to\Sub(A)$ preserva l'ordre: si $[m]\leq[n]$ a $\Sub(B)$, aleshores $f^{*}[m]\leq f^{*}[n]$ a $\Sub(A)$. És a dir, $f^{*}$ és una aplicació monòtona.
\end{proposicio}

\begin{prova}
Suposem $[m]\leq[n]$, amb $m\colon S\rightarrowtail B$, $n\colon T\rightarrowtail B$ i $h\colon S\to T$ tal que $n\comp h=m$. Siguin $(P,m',f_m)$ un pullback de $m$ al llarg de $f$ i $(Q,n',f_n)$ un de $n$, amb $f\comp m'=m\comp f_m$ i $f\comp n'=n\comp f_n$. La parella $m'\colon P\to A$, $h\comp f_m\colon P\to T$ és un con per al pullback de $n$, perquè
\[
f\comp m'=m\comp f_m=n\comp h\comp f_m.
\]
Per tant hi ha un únic $k\colon P\to Q$ amb $n'\comp k=m'$ i $f_n\comp k=h\comp f_m$. La primera igualtat diu que $[m']\leq[n']$. Observem que l'argument val per a representants i pullbacks \emph{qualssevol}.
\end{prova}

\begin{observacio}[$f^{*}$ està ben definida]\label{obs:reindexacio-ben-definida}
La definició~\ref{def:reindexacio-sub} tria un representant $m$ de la classe i un pullback. El resultat no depèn de cap de les dues tries. Si $m\sim m_1$, aleshores $[m]\leq[m_1]$ i $[m_1]\leq[m]$, i la demostració anterior, aplicada en tots dos sentits, dona que els pullbacks corresponents són equivalents. El cas de dos pullbacks del mateix $m$ és el cas particular $m_1=m$.
\end{observacio}

\begin{proposicio}[El functor de subobjectes]\label{prop:functor-sub}\index{Sub@$\Sub$!functor}
Si $\cat{C}$ és ben potenciada i té pullbacks, les reindexacions defineixen un functor
\[
\Sub\colon\cat{C}\op\longrightarrow\cat{Poset},
\]
que envia cada objecte $A$ a l'ordre parcial $\Sub(A)$ i cada morfisme $f\colon A\to B$ a l'aplicació monòtona $f^{*}\colon\Sub(B)\to\Sub(A)$. És a dir, $\id_A^{*}=\id_{\Sub(A)}$ i $(g\comp f)^{*}=f^{*}\comp g^{*}$, com a igualtats d'aplicacions entre ordres parcials.
\end{proposicio}

\begin{prova}
Com que $\cat{C}$ és ben potenciada, cada $\Sub(A)$ és un conjunt parcialment ordenat, i cada $f^{*}$ és una aplicació monòtona (proposició~\ref{prop:reindexacio-monotona}) ben definida (observació~\ref{obs:reindexacio-ben-definida}). Per aquesta mateixa observació, per calcular $f^{*}[m]$ podem fer servir \emph{qualsevol} pullback de $m$ al llarg de $f$. Queden les dues lleis functorials.

\emph{Identitats.} Sigui $m\colon S\rightarrowtail A$. El quadrat format per $m$ als dos costats verticals i per les identitats $\id_S$ i $\id_A$ als horitzontals és un pullback de $m$ al llarg de $\id_A$. En efecte, si $x\colon Z\to A$ i $y\colon Z\to S$ compleixen $\id_A\comp x=m\comp y$, l'única fletxa $z\colon Z\to S$ amb $m\comp z=x$ i $\id_S\comp z=y$ és $z=y$. Per tant $\id_A^{*}[m]=[m]$.

\emph{Composició.} Siguin $f\colon A\to B$, $g\colon B\to C$ i $m\colon S\rightarrowtail C$. Prenem un pullback $(Q,m_1,g')$ de $m$ al llarg de $g$, amb $g\comp m_1=m\comp g'$, de manera que $g^{*}[m]=[m_1]$. Prenem després un pullback $(P,m_2,f')$ de $m_1$ al llarg de $f$, amb $f\comp m_2=m_1\comp f'$, de manera que $f^{*}(g^{*}[m])=[m_2]$. Vegem que $(P,m_2,g'\comp f')$ és un pullback de $m$ al llarg de $g\comp f$; així $(g\comp f)^{*}[m]=[m_2]$ i les dues classes coincideixen. Primer, el quadrat commuta: $(g\comp f)\comp m_2=g\comp m_1\comp f'=m\comp(g'\comp f')$. Siguin ara $x\colon Z\to A$ i $y\colon Z\to S$ amb $g\comp f\comp x=m\comp y$.
\begin{itemize}
\item \emph{Existència.} La parella $(f\comp x,\,y)$ és un con per al primer pullback, de manera que hi ha un únic $q\colon Z\to Q$ amb $m_1\comp q=f\comp x$ i $g'\comp q=y$. Aleshores $(x,q)$ és un con per al segon pullback, i hi ha un únic $p\colon Z\to P$ amb $m_2\comp p=x$ i $f'\comp p=q$. Per tant, $m_2\comp p=x$ i $(g'\comp f')\comp p=y$.
\item \emph{Unicitat.} Si $p'\colon Z\to P$ també compleix $m_2\comp p'=x$ i $g'\comp f'\comp p'=y$, aleshores $f'\comp p'$ compleix $m_1\comp(f'\comp p')=f\comp m_2\comp p'=f\comp x$ i $g'\comp(f'\comp p')=y$. Per la unicitat al primer pullback, $f'\comp p'=q$. Per la unicitat al segon, $p'=p$.
\end{itemize}
Aquest argument és el lema d'enganxament de pullbacks, en la direcció que aquí cal (exercici resolt corresponent de la Pràctica del capítol~\ref{cap:limits}). Per tant $(g\comp f)^{*}=f^{*}\comp g^{*}$.

Les dues lleis són \emph{igualtats} d'aplicacions entre ordres parcials, i no només isomorfismes. A nivell de monomorfismes, el pullback al llarg de $g\comp f$ i el pullback iterat només coincideixen llevat d'isomorfisme. En passar a classes, però, aquest isomorfisme queda absorbit.
\end{prova}

\begin{observacio}
Aquí no cal cap noció de pseudofunctor. La coherència \enquote{llevat d'isomorfisme} que vam trobar en reindexar les categories tall $\cat{C}/B$ (observació~\ref{obs:reindexacio}) desapareix en passar a classes de subobjectes. Aquest functor és precisament el que el capítol següent voldrà representar.
\end{observacio}

\begin{observacio}[La substitució]
Aquesta és la interpretació que anunciàvem. Si pensem un subobjecte $[m]\in\Sub(B)$ com un predicat $P(y)$ sobre la variable $y:B$, i $f\colon A\to B$ com un \enquote{terme} $f(x)$ amb $x:A$, aleshores $f^{*}[m]\in\Sub(A)$ és el predicat $P(f(x))$: la \textbf{substitució} de $y$ per $f(x)$ dins de $P$. A $\cat{Set}$, on $[m]$ és un subconjunt $S\subseteq B$, la reindexació és l'antiimatge:
\[
f^{*}(S)=f^{-1}(S)=\{x\in A\mid f(x)\in S\}.
\]
La monotonia de $f^{*}$ diu que la substitució respecta la implicació: si $P\Rightarrow Q$, aleshores $P(f)\Rightarrow Q(f)$.
\end{observacio}

\begin{exemple}[Subgrafs]\label{ex:subgrafs}\index{graf dirigit!subobjectes}
A $\cat{Graph}$, els monomorfismes són els morfismes de grafs injectius sobre vèrtexs i sobre arestes (vegeu els Exercicis). Un subobjecte d'un graf $G$ es pot representar, doncs, per un \textbf{subgraf}: un subconjunt de vèrtexs $V'\subseteq V_G$ i un subconjunt d'arestes $E'\subseteq E_G$ tancat per origen i destí, és a dir, tal que l'origen i el destí de tota aresta de $E'$ són a $V'$. La reindexació al llarg d'un morfisme de grafs $F\colon G'\to G$ es calcula component a component, per antiimatge de vèrtexs i d'arestes:
\[
F^{*}(V',E')=\bigl(F_V^{-1}(V'),\ F_E^{-1}(E')\bigr),
\]
que torna a ser un subgraf de $G'$, perquè $F$ preserva origen i destí. L'identificació \enquote{predicat $=$ subobjecte} no és, doncs, exclusiva dels conjunts: un predicat sobre un graf és un subgraf, i la substitució és l'antiimatge component a component.
\end{exemple}

\section{Intersecció de subobjectes}

L'ordre parcial $\Sub(A)$ té més estructura que l'ordre: quan la categoria té pullbacks, admet \textbf{ínfims binaris}, que interpreten la conjunció de predicats.

\begin{proposicio}[La intersecció és un pullback]\index{subobjecte!intersecció}
\label{prop:interseccio}
Siguin $m\colon S\rightarrowtail A$ i $n\colon T\rightarrowtail A$ dos subobjectes de $A$. El seu \textbf{pullback} sobre $A$,
\[
\begin{tikzpicture}[baseline=(m.center)]
  \matrix (m) [matrix of math nodes, column sep=3em, row sep=2.4em] {
    S\cap T & T \\
    S & A \\
  };
  \draw[-{Stealth[scale=1.1]}] (m-1-1) -- node[above] {$q$} (m-1-2);
  \draw[-{Stealth[scale=1.1]}] (m-1-1) -- node[left] {$p$} (m-2-1);
  \draw[-{Stealth[scale=1.1]}] (m-1-2) -- node[right] {$n$} (m-2-2);
  \draw[-{Stealth[scale=1.1]}] (m-2-1) -- node[below] {$m$} (m-2-2);
  \node at ([xshift=1em,yshift=-0.7em]m-1-1) {$\scriptstyle\lrcorner$};
\end{tikzpicture}
\qquad
m\comp p=n\comp q,
\]
amb el morfisme induït $m\comp p=n\comp q\colon S\cap T\to A$, que és un monomorfisme, representa l'\textbf{ínfim} $[m]\wedge[n]$ a $\Sub(A)$. Escrivim també $S\cap T=S\times_A T$.
\end{proposicio}

\begin{prova}
El morfisme $m\comp p\colon S\cap T\to A$ és mono per ser composició de monos (la projecció $p$ del pullback és mono per la proposició~\ref{prop:pullback-mono}, i $m$ és mono). Per construcció $[m\wedge n]\leq[m]$ i $[m\wedge n]\leq[n]$, ja que hi ha les projeccions cap a $S$ i $T$. I si $[k]\leq[m]$ i $[k]\leq[n]$ per a un tercer subobjecte $k\colon R\rightarrowtail A$, les factoritzacions $R\to S$ i $R\to T$ són compatibles sobre $A$, de manera que la propietat universal del pullback dona una única factorització $R\to S\cap T$; per tant $[k]\leq[m]\wedge[n]$. Això caracteritza l'ínfim.
\end{prova}

\begin{observacio}
En termes lògics, $[m]\wedge[n]$ és la \textbf{conjunció} dels predicats: l'element pertany a $S\cap T$ si i només si pertany a $S$ \emph{i} a $T$. A $\cat{Set}$ recuperem la intersecció de subconjunts. En qualsevol categoria, el subobjecte \textbf{màxim} de $\Sub(A)$ és $[\id_A]$ (el predicat sempre cert), com es comprova a la Pràctica; i quan hi ha objecte inicial estricte, el subobjecte \textbf{mínim} és la injecció de l'objecte inicial (el predicat sempre fals). Amb intersecció i aquests extrems, $\Sub(A)$ és un \emph{semireticle amb màxim}. Les \textbf{unions} ---la disjunció--- exigeixen estructura addicional: la regularitat proporciona imatges, existencials i interseccions, però no, en general, suprems finits estables a les fibres $\Sub(A)$. Aquests apareixen, per exemple, en les \textbf{categories coherents}, que es tractaran al volum de lògica categòrica.
\end{observacio}

\section{Epimorfismes regulars i imatges}

Per definir la imatge d'un morfisme, ens cal la classe adequada d'epimorfismes: els que sorgeixen com a quocients.

\begin{definicio}[Epimorfisme regular]\index{epimorfisme!regular}
\label{def:epi-regular}
Un morfisme $e\colon A\to B$ és un \textbf{epimorfisme regular} si és el coequalitzador d'alguna parella de morfismes $u,v\colon R\to A$.
\end{definicio}

\begin{definicio}[Parella nucli]\index{parella nucli}
\label{def:parella-nucli}
La \textbf{parella nucli} d'un morfisme $f\colon A\to B$ és el pullback de $f$ amb ell mateix,
\[
\begin{tikzpicture}[baseline=(q.center)]
  \matrix (q) [matrix of math nodes, column sep=3em, row sep=2.4em] {
    R & A \\
    A & B \\
  };
  \draw[-{Stealth[scale=1.1]}] (q-1-1) -- node[above] {$p_2$} (q-1-2);
  \draw[-{Stealth[scale=1.1]}] (q-1-1) -- node[left] {$p_1$} (q-2-1);
  \draw[-{Stealth[scale=1.1]}] (q-1-2) -- node[right] {$f$} (q-2-2);
  \draw[-{Stealth[scale=1.1]}] (q-2-1) -- node[below] {$f$} (q-2-2);
  \node at ([xshift=0.8em,yshift=-0.7em]q-1-1) {$\scriptstyle\lrcorner$};
\end{tikzpicture}
\qquad
f\comp p_1=f\comp p_2;
\]
els seus dos morfismes $p_1,p_2\colon R\to A$ recullen les parelles d'elements que $f$ identifica. A $\cat{Set}$, $R=\{(a,a')\mid f(a)=f(a')\}$.
\end{definicio}

\begin{lema}\label{lema:epi-regular-nucli}\index{epimorfisme!regular}\index{parella nucli}
En una categoria amb pullbacks, si $e\colon A\to B$ és un epimorfisme regular, aleshores $e$ és el coequalitzador de la seva parella nucli.
\end{lema}

\begin{prova}
Suposem que $e$ és el coequalitzador de $u,v\colon K\rightrightarrows A$, i sigui $p_1,p_2\colon R\rightrightarrows A$ la parella nucli de $e$. Com que $e\comp u=e\comp v$, la propietat universal del pullback dona $k\colon K\to R$ amb $p_1\comp k=u$ i $p_2\comp k=v$. Si $h\colon A\to Z$ satisfà $h\comp p_1=h\comp p_2$, aleshores
\[
h\comp u=h\comp p_1\comp k=h\comp p_2\comp k=h\comp v,
\]
i com que $e$ és el coequalitzador de $u$ i $v$, hi ha una única $\bar h\colon B\to Z$ amb $\bar h\comp e=h$. Com que, a més, $e\comp p_1=e\comp p_2$, això és exactament la propietat universal del coequalitzador de $p_1$ i $p_2$.
\end{prova}

\begin{definicio}[Factorització imatge]\index{imatge}\index{factorització!imatge}
\label{def:imatge}
Una \textbf{factorització imatge} de $f\colon A\to B$ és una descomposició $f=m\comp e$ amb $A\xrightarrow{\,e\,}I\xrightarrow{\,m\,}B$, on $e$ és un epimorfisme regular i $m$ un monomorfisme. El subobjecte $[m]\in\Sub(B)$ és la \textbf{imatge} de $f$, denotada $\Imatge(f)$.
\end{definicio}

\begin{proposicio}[La imatge és el mínim subobjecte]
\label{prop:imatge-unica}
Quan existeix, la factorització imatge és única llevat d'isomorfisme canònic, i $\Imatge(f)$ és el \emph{menor} subobjecte de $B$ a través del qual $f$ es factoritza: si $f$ es factoritza per un mono $n\colon T\rightarrowtail B$, aleshores $\Imatge(f)\leq[n]$.
\end{proposicio}

\begin{prova}
Sigui $f=m\comp e$ una factorització imatge. Com que $m$ és mono, qualsevol parella nucli de $e$ és també una parella nucli de $f$ (perquè $m\comp e\comp u=m\comp e\comp v$ si i només si $e\comp u=e\comp v$); en particular, les parelles nucli escollides per a $e$ i per a $f$ són canònicament isomorfes. Pel lema~\ref{lema:epi-regular-nucli}, $e$ és el coequalitzador de la seva parella nucli $p_1,p_2$. Si $f=m'\comp e'$ és una altra factorització amb $m'$ mono, aleshores $e'$ iguala les branques d'aquesta parella nucli, perquè $m'\comp e'\comp p_1=f\comp p_1=f\comp p_2=m'\comp e'\comp p_2$ i $m'$ és mono, de manera que $e'$ es factoritza per $e$ mitjançant un únic morfisme; simètricament s'obté l'invers, i les dues composicions són identitats perquè $e$ és epi. Això dona l'isomorfisme canònic. La minimalitat en segueix: si $f=n\comp g$ amb $n$ mono, aleshores $g$ iguala la parella nucli i es factoritza per $e$, d'on $\Imatge(f)=[m]\leq[n]$.
\end{prova}

\begin{observacio}
A $\cat{Set}$, $\Imatge(f)$ és el subconjunt imatge $f(A)\subseteq B$, amb $e$ l'aplicació exhaustiva $A\to f(A)$ i $m$ la injecció. Recuperem la noció usual d'imatge.
\end{observacio}

\section{Categories regulars i estabilitat}

Perquè la teoria d'imatges sigui utilitzable, cal que les imatges es comportin bé sota el canvi de base.

\begin{definicio}[Categoria regular]\index{categoria!regular}
\label{def:regular}
Una categoria $\cat{C}$ és \textbf{regular} si té tots els límits finits, tota fletxa admet una factorització imatge (epi regular seguit de mono), i els epimorfismes regulars són \textbf{estables per pullback}.
\end{definicio}

\begin{exemple}
Són regulars $\cat{Set}$, tota categoria d'àlgebres d'una teoria algebraica (grups, anells, mòduls\ldots), tota categoria abeliana i tot topos elemental. La categoria $\cat{Top}$ \emph{no} és regular, perquè les imatges no hi són estables per pullback.
\end{exemple}

\begin{proposicio}[Estabilitat de la imatge]
\label{prop:imatge-estable}
En una categoria regular, la formació d'imatges commuta amb el pullback: si es forma el pullback de $f\colon A\to B$ al llarg de $g\colon B'\to B$, aleshores $\Imatge(f')=g^{*}\bigl(\Imatge(f)\bigr)$, on $f'$ és el morfisme reindexat.
\end{proposicio}

\begin{prova}
Factoritzem $f=m\comp e$. En prendre el pullback al llarg de $g$, l'epimorfisme regular $e$ dona un epi regular $e'$ (per l'estabilitat) i el mono $m$ dona un mono $m'$ (proposició~\ref{prop:pullback-mono}). Així $f'=m'\comp e'$ és una factorització epi regular\,--\,mono, i per la unicitat de la imatge, $m'$ representa $g^{*}(\Imatge(f))$.
\end{prova}

\section{L'existencial com a adjunt de la reindexació}

De tot això en resulta que la reindexació té un adjunt per l'esquerra, que reprèn el fil del capítol d'adjuncions.

\begin{teorema}[$\exists_f\dashv f^{*}$]\index{quantificador!existencial}\index{adjunció!existencial}
\label{teo:existencial-adjunt}
Sigui $\cat{C}$ una categoria regular i $f\colon A\to B$. L'operació
\[
\exists_f[s]:=\Imatge(f\comp s)\in\Sub(B)
\]
és adjunta per l'esquerra de la reindexació $f^{*}\colon\Sub(B)\to\Sub(A)$; és a dir,
\[
\exists_f[s]\leq[t]\quad\Longleftrightarrow\quad[s]\leq f^{*}[t].
\]
\end{teorema}

\begin{observacio}
L'operació $\exists_f$ està definida sobre classes de subobjectes. Si $s\sim s'$, és a dir, $s'=s\comp\varphi$ amb $\varphi$ isomorfisme, aleshores $f\comp s'=(f\comp s)\comp\varphi$, i les factoritzacions imatge de $f\comp s$ i de $f\comp s'$ donen el mateix subobjecte de $B$. És exactament la unicitat de la imatge llevat d'isomorfisme (proposició~\ref{prop:imatge-unica}) el que fa que $\exists_f$ estigui ben definida.
\end{observacio}

\begin{prova}
Si $[s]\leq f^{*}[t]$, aleshores $S$ es factoritza a través del pullback $f^{*}T$, i component-lo amb la projecció cap a $T$ s'obté que $f\comp s$ es factoritza per $t$; com que $\Imatge(f\comp s)$ és el mínim subobjecte amb aquesta propietat, $\exists_f[s]\leq[t]$. Recíprocament, si $\exists_f[s]\leq[t]$, aleshores $f\comp s=t\comp k$ per a algun $k$, i aquesta igualtat exhibeix $S$ com un con sobre el diagrama del pullback $f^{*}T$, de manera que $[s]\leq f^{*}[t]$.
\end{prova}

\begin{observacio}[L'universal com a estructura addicional]\index{quantificador!universal}\index{adjunció!quantificadors}
En una categoria regular, la reindexació $f^{*}$ té sempre l'adjunt per l'esquerra $\exists_f$ que acabem de construir, però \emph{no} té necessàriament un adjunt per la dreta. Quan cada $f^{*}$ admet a més un adjunt per la dreta, el denotem $\forall_f$ i queda caracteritzat per
\[
f^{*}[t]\leq[s]\quad\Longleftrightarrow\quad[t]\leq\forall_f[s].
\]
Aquesta és estructura addicional, que la regularitat no garanteix: la regularitat només dona el costat existencial. Apareix, per exemple, en les categories localment cartesianes tancades amb les hipòtesis de mida adequades i, en particular, en tot topos elemental. La construcció de $\forall_f$ i el seu paper lògic es tracten al volum de lògica categòrica; no s'obté simplement dualitzant la construcció de la imatge dins de la mateixa fibra de subobjectes.
\end{observacio}

\begin{observacio}[Compatibilitat amb la substitució]\index{Beck--Chevalley}
En una categoria regular, l'estabilitat de les imatges per pullback dona una compatibilitat entre l'existencial i la substitució. Per a un quadrat cartesià
\[
\begin{tikzpicture}[baseline=(q.center)]
  \matrix (q) [matrix of math nodes, column sep=3em, row sep=2.4em] {
    A' & A \\
    B' & B \\
  };
  \draw[-{Stealth[scale=1.1]}] (q-1-1) -- node[above] {$h$} (q-1-2);
  \draw[-{Stealth[scale=1.1]}] (q-1-1) -- node[left] {$f'$} (q-2-1);
  \draw[-{Stealth[scale=1.1]}] (q-1-2) -- node[right] {$f$} (q-2-2);
  \draw[-{Stealth[scale=1.1]}] (q-2-1) -- node[below] {$g$} (q-2-2);
  \node at ([xshift=0.8em,yshift=-0.7em]q-1-1) {$\scriptstyle\lrcorner$};
\end{tikzpicture}
\]
es té la \textbf{condició de Beck--Chevalley}
\[
g^{*}\comp\exists_f=\exists_{f'}\comp h^{*}\colon\Sub(A)\to\Sub(B'):
\]
quantificar existencialment i després substituir és el mateix que substituir i després quantificar. No ho demostrem aquí; és un primer exemple del tipus de coherència que la lògica categòrica exigeix.
\end{observacio}

\begin{observacio}[Lectura lògica]
Aquesta adjunció té una lectura lògica il·lustrativa. Si $f$ és la projecció $\pi\colon X\times Y\to X$ que oblida la variable $Y$, i pensem un subobjecte de $X\times Y$ com un predicat $s(x,y)$, aleshores
\[
\exists_\pi[s]=\{x\mid\exists y,\ s(x,y)\}.
\]
L'existencial apareix, doncs, com a adjunt per l'esquerra de la substitució. Quan la categoria té també els adjunts per la dreta, l'universal es llegeix anàlogament com $\forall_\pi[s]=\{x\mid\forall y,\ s(x,y)\}$. El desenvolupament d'aquesta idea com a \emph{sistema lògic} ---amb el seu càlcul, la interpretació de l'existencial per la factorització imatge i les condicions de coherència pertinents--- és objecte de l'estudi de la lògica categòrica, i queda fora de l'abast d'aquest volum.
\end{observacio}

\begin{resumcap}
\apartat{Idees centrals}
\begin{enumerate}
  \item Un subobjecte és una classe d'equivalència de monomorfismes amb el mateix codomini; $\Sub(A)$ és un ordre parcial, l'àlgebra de predicats sobre $A$. A $\cat{Set}$ són els subconjunts, i a $\cat{Graph}$, els subgrafs.
  \item La reindexació $f^{*}$, construïda per pullback, interpreta la substitució; en una categoria ben potenciada amb pullbacks dona un functor $\Sub\colon\cat{C}\op\to\cat{Poset}$, amb igualtats estrictes.
  \item La intersecció de subobjectes és el seu pullback i interpreta la conjunció; el màxim $[\id_A]$ interpreta el predicat cert.
  \item Un epimorfisme regular és el coequalitzador de la seva parella nucli; la factorització imatge (epi regular seguit de mono) és única llevat d'isomorfisme, i en una categoria regular existeix i és estable per pullback.
  \item En una categoria regular, $f^{*}$ té un adjunt per l'esquerra $\exists_f$, que interpreta la quantificació existencial; l'universal $\forall_f$ és estructura addicional que la regularitat no garanteix.
\end{enumerate}
\apartat{Errors típics} Creure que la regularitat dona unions de subobjectes (calen categories coherents) o l'adjunt dret $\forall_f$; confondre els epimorfismes amb els epimorfismes regulars (a $\cat{Top}$, els primers són les aplicacions contínues exhaustives i els segons, les aplicacions quocient); confondre la reindexació $f^{*}$ amb la imatge $\exists_f$; i usar \enquote{imatge}, \enquote{buit} o \enquote{elements} en una categoria arbitrària sense declarar les hipòtesis.
\apartat{Lectura recomanada} \cite[cap.~5 i~\S9.5]{awodey} per a subobjectes, imatges i l'adjunció existencial; \cite[§§2.8, 2.10, 9.4 i 9.8]{barrwells} per a subobjectes, epimorfismes regulars, categories regulars i sistemes de factorització; \cite[V.7]{maclane} per al tractament clàssic dels subobjectes; \cite[A1.3]{johnstone} i \cite{borceux} per a categories regulars amb més profunditat; \cite{jacobs} per a la lectura fibracional de la reindexació.
\end{resumcap}

\chapter[Classificadors de subobjectes i topos]{\texorpdfstring{Classificadors de subobjectes\\ i introducció als topos}{Classificadors de subobjectes i introducció als topos}}
\label{cap:topos}

\begin{objectius}
\apartat{Prerequisits} Subobjectes i el functor $\Sub$ (cap.~\ref{cap:subobjectes}); representabilitat i Yoneda (cap.~\ref{cap:yoneda}); tancament cartesià (cap.~\ref{cap:cartesianes}).
\apartat{Objectius} En acabar el capítol, el lector hauria de poder:
\begin{itemize}
  \item definir el classificador de subobjectes $\Omega$ i el morfisme característic d'un subobjecte;
  \item expressar la pertinença d'un element generalitzat com una igualtat de morfismes cap a $\Omega$;
  \item demostrar que tenir classificador equival a la representabilitat del functor $\Sub$;
  \item enunciar la definició de topos elemental i reconèixer $\Omega^{E}$ com a objecte de parts;
  \item calcular morfismes característics a $\cat{Set}$ i a $\cat{Graph}$, i explicar per què la representabilitat de $\Sub$ converteix els predicats en morfismes i la substitució en composició;
  \item reconèixer com a topos $\cat{Set}$, $\cat{FinSet}$, les categories de prefeixos sobre una categoria petita, amb el classificador de sedassos, i $\cat{Graph}$, i distingir els topos elementals de les categories de prefeixos;
  \item explicar en quin sentit $\Omega$ és un objecte de valors de veritat, i per què això no fa booleà tot topos.
\end{itemize}
\end{objectius}

Al capítol anterior hem vist que els subobjectes d'un objecte $A$ formen un ordre parcial $\Sub(A)$, que podem llegir com l'ordre de predicats sobre $A$, i que la reindexació el fa functorial: en una categoria ben potenciada amb pullbacks obtenim un functor $\Sub\colon\cat{C}\op\to\cat{Poset}$ (proposició~\ref{prop:functor-sub}). Sorgeix una pregunta natural: hi ha un objecte que \emph{representi} aquest functor? És a dir, un objecte $\Omega$ tal que donar un subobjecte de $A$ equivalgui a donar un morfisme $A\to\Omega$? Quan aquest objecte existeix, s'anomena \textbf{classificador de subobjectes}, i les categories que en tenen ---juntament amb límits finits i exponencials--- són els \textbf{topos elementals}. Aquest capítol n'introdueix la definició i els exemples bàsics. No en desenvolupem la lògica interna: aquesta pertany a l'estudi de la lògica categòrica, per al qual aquest volum és la preparació.

\section{El classificador de subobjectes}

A $\cat{Set}$, un subconjunt $U\subseteq E$ es descriu per la seva \textbf{funció característica} $\chi_U\colon E\to\{0,1\}$, que val $1$ als elements de $U$ i $0$ a la resta. El conjunt $\{0,1\}$ actua com un objecte de \enquote{valors de veritat}, i triar un subconjunt equival a triar una funció cap a ell. La noció següent generalitza aquesta idea a una categoria arbitrària.

\begin{definicio}[Classificador de subobjectes]\index{classificador de subobjectes}\index{Omega@$\Omega$}
\label{def:classificador}
Sigui $\cat{E}$ una categoria amb tots els límits finits. Un \textbf{classificador de subobjectes} és un objecte $\Omega$ juntament amb un morfisme $\mathsf{true}\colon 1\to\Omega$ des de l'objecte terminal, amb la propietat universal següent: per a cada objecte $E$ i cada subobjecte $m\colon U\rightarrowtail E$, existeix un \emph{únic} morfisme $\chi_m\colon E\to\Omega$, el \textbf{morfisme característic} del subobjecte, que fa del quadrat
\[
\begin{tikzpicture}[baseline=(m.center)]
  \matrix (m) [matrix of math nodes, column sep=3em, row sep=2.4em] {
    U & 1 \\
    E & \Omega \\
  };
  \draw[-{Stealth[scale=1.1]}] (m-1-1) -- node[above] {$!$} (m-1-2);
  \draw[-{Stealth[scale=1.1]}] (m-1-1) -- node[left] {$m$} (m-2-1);
  \draw[-{Stealth[scale=1.1]}] (m-1-2) -- node[right] {$\mathsf{true}$} (m-2-2);
  \draw[-{Stealth[scale=1.1]}] (m-2-1) -- node[below] {$\chi_m$} (m-2-2);
  \node at ([xshift=0.8em,yshift=-0.7em]m-1-1) {$\scriptstyle\lrcorner$};
\end{tikzpicture}
\]
un pullback. És a dir, $U$ es recupera com el pullback de $\mathsf{true}$ al llarg de $\chi_m$. A $\cat{Set}$, el morfisme característic és la funció característica usual.
\end{definicio}

La condició diu que cada subobjecte de $E$ té un i només un morfisme característic, i que el subobjecte es reconstrueix com l'antiimatge del \enquote{punt vertader} $\mathsf{true}$. A $\cat{Set}$, $\Omega=\{0,1\}$ i $\mathsf{true}\colon 1\to\{0,1\}$ tria l'element $1$; la funció característica és la de sempre, i el pullback de $\mathsf{true}$ al llarg de $\chi_U$ recupera exactament $U$.

\begin{exemple}[Un morfisme característic calculat a $\cat{Set}$]\label{ex:caracteristic-set}\index{morfisme característic!a Set@a $\cat{Set}$}
Sigui $E=\{1,2,3,4\}$ i $U=\{2,4\}$, el predicat \enquote{$x$ és parell}, amb la injecció $m\colon U\hookrightarrow E$. El morfisme característic és
\[
\chi_U\colon E\to\{0,1\},\qquad 1\mapsto0,\quad 2\mapsto1,\quad 3\mapsto0,\quad 4\mapsto1.
\]
Comprovem les tres afirmacions de la definició sobre aquest cas.
\begin{itemize}
\item \emph{El quadrat commuta.} $\chi_U\comp m$ envia $2$ i $4$ a $1$, que és $\mathsf{true}\comp\,!_U$.
\item \emph{És un pullback.} El pullback de $\mathsf{true}$ al llarg de $\chi_U$ és $\{x\in E\mid\chi_U(x)=1\}=\{2,4\}$, amb la inclusió (el pullback a $\cat{Set}$, capítol~\ref{cap:limits}): recupera exactament $U$.
\item \emph{Unicitat.} Una altra aplicació $\chi\colon E\to\{0,1\}$ amb el mateix pullback ha de valer $1$ exactament a $\{2,4\}$, i per tant és $\chi_U$. Per exemple, l'aplicació que val $1$ a $\{2,3,4\}$ fa commutar el quadrat, però el seu pullback és $\{2,3,4\}$, no $U$: la commutativitat sola no basta.
\end{itemize}
La pertinença per elements, que formalitzarem tot seguit, hi és visible: l'element $x\colon\{\ast\}\to E$ amb $x(\ast)=4$ compleix $\chi_U\comp x=\mathsf{true}$, i el que tria $3$, no. La substitució també: si $g\colon\{a,b,c\}\to E$ és $g(a)=1$, $g(b)=2$, $g(c)=4$, aleshores $\chi_U\comp g$ envia $a\mapsto0$, $b\mapsto1$ i $c\mapsto1$, i classifica $g^{-1}(U)=\{b,c\}$. \emph{Substituir} en un predicat és \emph{precompondre} amb el seu morfisme característic.
\end{exemple}

\begin{observacio}[Pertinença per elements generalitzats]\index{element generalitzat!pertinença}
Sigui $m\colon U\rightarrowtail E$ un subobjecte amb morfisme característic $\chi_m\colon E\to\Omega$, i sigui $x\colon X\to E$ un element generalitzat. Aleshores
\[
x\ \text{factoritza a través de } m
\quad\Longleftrightarrow\quad
\chi_m\comp x=\mathsf{true}\comp\,!_X.
\]
En efecte, si $x=m\comp y$, aleshores $\chi_m\comp x=\chi_m\comp m\comp y=\mathsf{true}\comp\,!_U\comp y=\mathsf{true}\comp\,!_X$; i recíprocament, si $\chi_m\comp x=\mathsf{true}\comp\,!_X$, la propietat universal del pullback dona una factorització (única) de $x$ a través de $m$. És la lectura precisa, sense elements ordinaris, de
\[
x\in U\quad\Longleftrightarrow\quad\chi_U(x)=\mathsf{true}:
\]
el classificador converteix la pertinença a un subobjecte en una igualtat de morfismes cap a $\Omega$.
\end{observacio}

\begin{proposicio}[Classificador i representabilitat]
\label{prop:omega-representa}
Suposem $\cat{E}$ localment petita, ben potenciada i amb límits finits, de manera que el functor de subobjectes pren valors a $\cat{Set}$. Aleshores les dades d'un classificador de subobjectes $(\Omega,\mathsf{true})$ són equivalents a una representació del functor de subobjectes:
\[
\Sub_{\cat{E}}(E)\;\cong\;\Hom_{\cat{E}}(E,\Omega),
\]
natural en $E$. En conseqüència, el classificador, si existeix, és únic llevat d'un únic isomorfisme compatible amb $\mathsf{true}$.
\end{proposicio}

\begin{prova}
\emph{Del classificador a la representació.} Enviem un subobjecte $[m]$ al seu morfisme característic $\chi_m$. L'aplicació és injectiva perquè $\chi_m$ determina $[m]$ com el pullback de $\mathsf{true}$ al llarg de $\chi_m$, i és exhaustiva perquè, per a tot $\chi\colon E\to\Omega$, el pullback de $\mathsf{true}$ al llarg de $\chi$ és un subobjecte el morfisme característic del qual és $\chi$, per la unicitat de la definició~\ref{def:classificador}. La naturalitat en $E$ és la compatibilitat amb la reindexació: reindexar un subobjecte al llarg de $g\colon E'\to E$ correspon a precompondre la característica amb $g$, pel lema d'enganxament de pullbacks.

\emph{De la representació al classificador.} Suposem que $\Sub_{\cat{E}}\cong\Hom_{\cat{E}}(-,\Omega)$ naturalment, i sigui $t\colon T\rightarrowtail\Omega$ un representant de l'element universal, és a dir, del subobjecte de $\Omega$ que correspon a $\id_\Omega$. Pel lema de Yoneda, el subobjecte que correspon a un morfisme $\chi\colon E\to\Omega$ és la reindexació $\chi^{*}[t]$: tot subobjecte és el pullback de $t$ al llarg d'un únic morfisme. Vegem que $T$ és terminal. Per a cada objecte $A$, el subobjecte màxim $[\id_A]$ és el pullback de $t$ al llarg d'un únic $\chi\colon A\to\Omega$, i el pullback dona un morfisme $A\to T$. Si $g\colon A\to T$ és un morfisme qualsevol, el pullback de $t$ al llarg de $t\comp g$ és $[\id_A]$, perquè $t$ és mono; per unicitat, $t\comp g=\chi$, i com que $t$ és mono, $g$ queda determinat. Així hi ha exactament un morfisme $A\to T$, i $T\cong 1$. Identificant $T$ amb $1$, el mono $t$ és un morfisme $\mathsf{true}\colon 1\to\Omega$ que satisfà la definició~\ref{def:classificador}.

La unicitat del classificador és el principi de Yoneda aplicat a dos objectes que representen el mateix functor, amb l'element universal $\mathsf{true}$ com a dada universal.
\end{prova}

\begin{observacio}[Per què la representabilitat és un pas conceptual]\label{obs:sub-representable}
L'equivalència de la proposició no és només una reformulació. El functor $\Sub$ és una dada \emph{dispersa}: un ordre parcial per a cada objecte, i una reindexació per a cada morfisme. La representabilitat el concentra en un sol objecte $\Omega$ i en un sol subobjecte universal, $\mathsf{true}\colon 1\rightarrowtail\Omega$. Aquest canvi té tres conseqüències:
\begin{itemize}
\item \emph{Els predicats esdevenen morfismes.} Un predicat sobre $E$ ja no és una classe d'equivalència de monomorfismes, sinó una fletxa $E\to\Omega$, que es pot compondre, aparellar i currificar com qualsevol altra.
\item \emph{La substitució esdevé composició.} La naturalitat de la bijecció diu que reindexar al llarg de $g$ és precompondre amb $g$, tal com hem vist a l'exemple~\ref{ex:caracteristic-set}.
\item \emph{Tot subobjecte és un cas d'un de sol.} Cada subobjecte és el pullback de $\mathsf{true}$ al llarg del seu morfisme característic. Les operacions lògiques es poden definir, doncs, una sola vegada sobre $\Omega$, com a morfismes $\Omega\times\Omega\to\Omega$, i transportar-se a tots els objectes (exercici resolt~\ref{pr-conjuncio-omega} de la Pràctica).
\end{itemize}
És el mateix pas que va de $\Hom(A\times B,C)$ a l'exponencial $C^B$: una família de conjunts esdevé un objecte de la categoria. Per això el classificador és el que permet que la lògica \enquote{visqui dins} de la categoria. Pel lema de Yoneda, a més, el classificador es pot calcular sondejant-lo amb objectes representables, com farem a $\cat{Graph}$.
\end{observacio}

\section{Topos elementals}

\begin{definicio}[Topos elemental]\index{topos}\index{topos!elemental}
\label{def:topos}
Un \textbf{topos elemental} és una categoria $\cat{E}$ que compleix:
\begin{enumerate}
\item $\cat{E}$ té tots els límits finits;
\item $\cat{E}$ té un classificador de subobjectes $\Omega$;
\item $\cat{E}$ és cartesiana tancada (té tots els exponencials).
\end{enumerate}
\end{definicio}

Aquesta definició, compacta, resulta ser extraordinàriament rica. Es pot demostrar, per exemple, que tot topos té també tots els colímits finits, i que, si $\cat{E}$ és un topos, també ho és cada categoria $\cat{E}/A$ de $\cat{E}$ sobre un objecte $A$. No demostrem aquí aquestes propietats; les recollim per il·lustrar que la combinació \enquote{límits finits $+$ classificador $+$ exponencials} és molt més potent del que sembla a primera vista.

\begin{observacio}[Topos elemental i categoria de prefeixos]\label{obs:elemental-prefeixos}\index{topos!elemental i de prefeixos}
Convé no confondre la definició amb el seu exemple principal. Un topos \emph{elemental} es defineix per tres condicions expressables en el llenguatge de les categories, i només demana límits \emph{finits}. Les categories de prefeixos $\cat{Set}^{\cat{C}\op}$ en són una classe important, però tenen molta més estructura: tots els límits i colímits \emph{petits}, i un classificador que es pot escriure explícitament amb sedassos. No tot topos elemental és d'aquesta mena. La categoria $\cat{FinSet}$ dels conjunts finits és un topos: té límits finits, l'exponencial de dos conjunts finits és finit, i $\{0,1\}$ hi és un classificador. Però no té coproductes infinits: no hi ha cap coproducte d'una família numerable de còpies de $1$, perquè un tal objecte hauria d'admetre una aplicació cap a $\{0,1\}$ per a cada successió de zeros i uns, i un conjunt finit només n'admet un nombre finit. En canvi, tota categoria de prefeixos té tots els colímits petits, i una equivalència de categories els conserva. Per tant, $\cat{FinSet}$ no és equivalent a cap categoria de prefeixos. Entre totes dues nocions hi ha els \emph{topos de Grothendieck} (categories de feixos), que aquí no tractem.
\end{observacio}

\begin{observacio}[Objecte de parts]
En un topos, l'exponencial $\Omega^{E}$ mereix un nom propi: és l'\textbf{objecte de parts} $\mathcal{P}(E)$, l'anàleg categòric del conjunt de parts. En efecte, per la propietat de l'exponencial i la del classificador,
\[
\Hom(A,\Omega^{E})\cong\Hom(A\times E,\Omega)\cong\Sub(A\times E),
\]
de manera que els morfismes cap a $\Omega^{E}$ classifiquen relacions entre $A$ i $E$. En el cas $A=1$, com que $1\times E\cong E$, s'obté
\[
\Hom(1,\Omega^{E})\cong\Sub(E):
\]
els punts globals de $\Omega^{E}$ són exactament els subobjectes de $E$, que és el que justifica el nom d'objecte de parts. Aquesta és la porta d'entrada a la lògica d'ordre superior interna d'un topos, que és tema de l'estudi posterior.
\end{observacio}

\section{Exemples}

\begin{exemple}[Conjunts]
$\cat{Set}$ és un topos. Té límits finits, és cartesià tancat (l'exponencial és el conjunt de funcions) i té classificador $\Omega=\{0,1\}$ amb $\mathsf{true}$ l'element $1$. És el topos prototípic, i tota la teoria es pot llegir com una generalització d'aquest cas.
\end{exemple}

\begin{exemple}[Prefeixos]\index{prefeix@prefeix!topos de}\index{sedàs}
Per a qualsevol categoria petita $\cat{C}$, la categoria de prefeixos $\widehat{\cat{C}}=\cat{Set}^{\cat{C}\op}$ (exemple~\ref{ex:categoria-prefeixos}) és un topos. Ja sabíem que té límits finits, calculats objecte a objecte (proposició~\ref{prop:limits-punt-a-punt}), i que és cartesiana tancada; el que hi falta és el classificador. Aquest es construeix amb els \textbf{sedassos}. Un sedàs sobre un objecte $C$ és un conjunt $S$ de morfismes amb codomini $C$ tal que
\[
f\colon D\to C,\quad f\in S,\quad g\colon X\to D
\qquad\Longrightarrow\qquad
f\comp g\in S.
\]
El prefeix $\Omega$ envia cada objecte $C$ al conjunt $\Omega(C)$ dels sedassos sobre $C$, i cada morfisme $h\colon D\to C$ a l'aplicació
\[
\Omega(h)\colon\Omega(C)\to\Omega(D),
\qquad
S\longmapsto h^{*}S=\{g\colon X\to D\mid h\comp g\in S\},
\]
que torna a donar un sedàs. El punt distingit $\mathsf{true}\colon 1\to\Omega$ tria a cada $C$ el sedàs màxim, format per tots els morfismes cap a $C$.

Si $U\rightarrowtail F$ és un subprefeix, el seu morfisme característic envia $x\in F(C)$ al sedàs dels morfismes que porten $x$ dins de $U$:
\[
\chi_U(C)(x)=\{f\colon D\to C\mid F(f)(x)\in U(D)\}.
\]
Es comprova que aquest conjunt és un sedàs, que les aplicacions $\chi_U(C)$ són naturals en $C$, i que $U$ és el pullback de $\mathsf{true}$ al llarg de $\chi_U$. Aquest exemple cobreix una quantitat enorme de casos i és el pont cap a la teoria de feixos.
\end{exemple}

\begin{exemple}[El topos dels grafs: el classificador calculat]\label{ex:omega-graph}\index{graf dirigit!topos}\index{classificador de subobjectes!a Graph@a $\cat{Graph}$}
Al capítol~\ref{cap:transformacions} vam veure que $\cat{Graph}\cong\cat{Set}^{\cat{P}}$, i com que $\cat{Set}^{\cat{P}}=\cat{Set}^{(\cat{P}\op)\op}$, és la categoria de prefeixos sobre $\cat{P}\op$. Per l'exemple anterior, $\cat{Graph}$ és, doncs, un topos. Calculem-ne el classificador.

\emph{Els vèrtexs i les arestes de $\Omega$.} Els dos representables són el graf $\mathbf V$ d'un sol vèrtex i el graf $\mathbf E$ d'una sola aresta $e\colon s\to t$ (exemple~\ref{ex:homfunctor-grafs}): els vèrtexs d'un graf $G$ corresponen als morfismes $\mathbf V\to G$, i les arestes, als morfismes $\mathbf E\to G$. Aplicant-ho a $\Omega$ i fent servir $\Hom(-,\Omega)\cong\Sub$, obtenim
\[
V_\Omega\cong\Sub(\mathbf V),\qquad E_\Omega\cong\Sub(\mathbf E).
\]
El graf $\mathbf V$ té dos subgrafs, el buit i ell mateix; així, $\Omega$ té dos vèrtexs, que anomenem $0$ (\enquote{el vèrtex no hi és}) i $1$ (\enquote{el vèrtex hi és}). El graf $\mathbf E$ té cinc subgrafs (exemple~\ref{ex:subgrafs}): $\varnothing$, $\{s\}$, $\{t\}$, $\{s,t\}$ i $\{s,t,e\}$. Així, $\Omega$ té cinc arestes.

\emph{Origen i destí.} L'origen d'una aresta $K\subseteq\mathbf E$ de $\Omega$ s'obté reindexant $K$ al llarg del morfisme $\mathbf V\to\mathbf E$ que tria $s$; és a dir, val $1$ si $s\in K$ i $0$ altrament. El destí, igual amb $t$. Anomenant les arestes per la informació que codifiquen, queda
\[
\renewcommand{\arraystretch}{1.2}
\begin{array}{c|c|l}
\text{aresta} & \text{subgraf de }\mathbf E & \text{de}\to\text{a}\\ \hline
\bot & \varnothing & 0\to0\\
\sigma & \{s\} & 1\to0\\
\tau & \{t\} & 0\to1\\
\partial & \{s,t\} & 1\to1\\
\top & \{s,t,e\} & 1\to1
\end{array}
\]
El graf $\Omega$ té, doncs, dos vèrtexs, un bucle a $0$, una aresta en cada sentit entre $0$ i $1$, i \emph{dos} bucles a $1$ (vegeu també \cite[exemple~15.4.3]{barrwells}). El morfisme $\mathsf{true}\colon\mathbf L\to\Omega$, definit sobre el graf terminal d'un bucle (exemple~\ref{ex:producte-tres-contextos}), envia el vèrtex a $1$ i el bucle a $\top$.

\emph{El morfisme característic d'un subgraf.} Si $H\subseteq G$ és un subgraf, $\chi_H$ envia cada vèrtex a $1$ o a $0$ segons que sigui o no a $H$. Cada aresta $x\colon v\to w$ va a l'aresta de $\Omega$ que registra si $v$, $w$ i $x$ són a $H$: a $\top$ si $x\in H$; a $\partial$ si $v,w\in H$ però $x\notin H$; i a $\sigma$, $\tau$ o $\bot$ segons quins extrems hi siguin. Per exemple, sigui $G$ el cicle de vèrtexs $a,b,c$ i arestes $e\colon a\to b$, $f\colon b\to c$, $k\colon c\to a$, i sigui $H$ el subgraf amb vèrtexs $a,b$ i aresta $e$. Aleshores
\[
\chi_H\colon\quad a,b\mapsto1,\quad c\mapsto0,\qquad e\mapsto\top,\quad f\mapsto\sigma,\quad k\mapsto\tau.
\]
És un morfisme de grafs, perquè $\sigma\colon1\to0$ i $f\colon b\to c$ tenen extrems compatibles, i el mateix per a les altres arestes. A més, el pullback de $\mathsf{true}$ al llarg de $\chi_H$, que es calcula component a component, té per vèrtexs els que van a $1$, $\{a,b\}$, i per arestes les que van a $\top$, $\{e\}$: és exactament $H$. La unicitat també es veu directament. Els vèrtexs de $H$ han d'anar a $1$ i els altres a $0$. Una aresta ha d'anar a una aresta de $\Omega$ amb els extrems ja determinats, i a $\top$ exactament si és a $H$; amb aquestes condicions, en cada cas hi ha una sola opció.

\emph{Més de dos valors de veritat.} La diferència entre $\partial$ i $\top$ és el que fa interessant aquest topos. Al graf $\mathbf E$, el subgraf $\{s\}$ i el subgraf $\{t\}$ són disjunts, i $\{t\}$ és el més gran dels subgrafs disjunts de $\{s\}$, perquè tot subgraf que contingui l'aresta ha de contenir $s$. Però la seva unió és $\{s,t\}$, que no és tot $\mathbf E$: hi falta l'aresta, que és justament el que distingeix $\partial$ de $\top$. A $\Sub(\mathbf E)$, doncs, un predicat i la seva negació no cobreixen tot l'objecte, i el principi del tercer exclòs falla. El classificador d'aquest topos és un graf amb estructura pròpia, no un conjunt $\{0,1\}$. La mateixa categoria que ens ha servit per introduir objectes, functors, Yoneda i límits és, així, un exemple genuí de topos no booleà.
\end{exemple}

\begin{observacio}[$\Omega$ com a objecte de valors de veritat]\label{obs:omega-veritat}\index{Omega@$\Omega$!valors de veritat}\index{topos!booleà}
La noció de topos va aparèixer primer a l'escola de Grothendieck de geometria algebraica, com a generalització de la d'espai topològic. Un dels seus aspectes més fascinants, però, és la relació amb la lògica. Per la representabilitat de $\Sub$, un morfisme $\varphi\colon E\to\Omega$ \emph{és} un predicat sobre $E$, la seva \enquote{extensió} és el subobjecte que classifica, i $\mathsf{true}$ fa el paper del valor \enquote{cert}. En aquest sentit precís ---i només en aquest--- $\Omega$ és un objecte de valors de veritat.

Això no vol dir que $\Omega$ tingui dos valors, ni que la lògica sigui clàssica. En tot topos cada $\Sub(E)$ és una àlgebra de Heyting (nota d'ampliació següent), de manera que la lògica interna és, en general, intuïcionista. Un topos s'anomena \textbf{booleà} quan totes aquestes àlgebres de Heyting són àlgebres de Boole, és a dir, quan per a tot subobjecte $U$ es compleix $U\vee\neg U=\top$. $\cat{Set}$ i $\cat{FinSet}$ són booleans; $\cat{Graph}$ no ho és, com hem vist a l'exemple~\ref{ex:omega-graph}. Fins i tot els \enquote{punts} de $\Omega$, els valors de veritat globals $1\to\Omega$, poden ser més de dos. I, a la inversa, tenir exactament dos valors globals no fa booleà un topos: els valors globals només compten $\Hom(1,\Omega)\cong\Sub(1)$, mentre que la booleanitat és una condició sobre tots els $\Sub(E)$. Els prefeixos sobre el monoide $(\mathbb N,+)$ tenen dos valors globals i no són booleans (vegeu els Exercicis per a tots dos fenòmens). Aquesta lectura permet interpretar la lògica ---de primer ordre i d'ordre superior--- dins de qualsevol topos. El seu desenvolupament és, precisament, l'objecte de l'estudi de la lògica categòrica.
\end{observacio}

\begin{nota}[Ampliació: $\Sub(E)$ és una àlgebra de Heyting]\label{nota:sub-heyting}\index{àlgebra de Heyting!subobjectes d'un topos}
L'afirmació de l'observació anterior és un teorema de la teoria de topos que aquest volum \emph{no demostra}; la fem servir només per situar la qüestió de la booleanitat. L'enunciat és aquest: en un topos elemental, cada ordre parcial $\Sub(E)$ té màxim, mínim, ínfims i suprems binaris i una implicació $\Rightarrow$ adjunta per la dreta de $-\wedge U$, i les reindexacions $f^{*}$ preserven tota aquesta estructura.

Algunes peces ja les tenim. El màxim és $[\id_E]$, i els ínfims són pullbacks (capítol~\ref{cap:subobjectes}). Les altres requereixen resultats que no hem provat. El mínim i els suprems es construeixen amb l'objecte inicial, els coproductes i les imatges; la seva existència fa servir que tot topos té colímits finits i és una categoria regular. La implicació s'obté amb el classificador i els exponencials. A $\cat{Graph}$, les dues construccions es poden verificar a mà sobre els subgrafs, cosa que justifica independentment l'exemple~\ref{ex:omega-graph}. Per a les demostracions, vegeu \cite[cap.~IV]{maclanemoerdijk} i \cite{johnstone}; el desenvolupament correspon a l'estudi de la lògica categòrica.
\end{nota}

\begin{resumcap}
\apartat{Idees centrals}
\begin{enumerate}
  \item El classificador $\Omega$ és un objecte amb $\mathsf{true}\colon 1\to\Omega$ tal que cada subobjecte té un únic morfisme característic, del qual és el pullback; la pertinença d'un element generalitzat esdevé una igualtat de morfismes cap a $\Omega$.
  \item Tenir classificador equival a la representabilitat del functor $\Sub$, i $\Omega$ n'és l'objecte representant: els predicats esdevenen morfismes $E\to\Omega$ i la substitució, precomposició.
  \item Un topos elemental té límits finits, exponencials i classificador; en un topos, $\Omega^{E}$ és l'objecte de parts de $E$.
  \item $\cat{Set}$ és el topos prototípic, amb $\Omega=\{0,1\}$.
  \item Tota categoria de prefeixos sobre una categoria petita és un topos, amb $\Omega$ el prefeix de sedassos; en particular, $\cat{Graph}$ és un topos, amb un classificador de dos vèrtexs i cinc arestes. No tot topos elemental és de prefeixos: $\cat{FinSet}$ no ho és.
  \item $\Omega$ és un objecte de valors de veritat perquè classifica predicats, però la lògica d'un topos és en general intuïcionista: $\cat{Graph}$ no és booleà.
\end{enumerate}
\apartat{Errors típics} Oblidar la hipòtesi de petitesa local\,/\,bona potenciació en escriure $\Sub\colon\cat{C}\op\to\cat{Set}$; confondre topos elemental amb categoria de prefeixos; confondre el classificador amb un conjunt de dos elements (només a $\cat{Set}$ el model prototípic és $\{0,1\}$; en un topos general, $\Omega$ és un objecte de la categoria i pot tenir una estructura molt més rica); i suposar que la lògica interna de qualsevol topos és clàssica: en general és intuïcionista, i la booleanitat és una propietat addicional.
\apartat{Lectura recomanada} \cite[cap.~8]{awodey}, especialment §8.8, per a classificadors, topos de prefeixos i la connexió amb Yoneda; \cite[cap.~15]{barrwells} per a una presentació alternativa i l'exemple dels grafs; \cite[cap.~1--2]{maclanemoerdijk} per continuar cap a la teoria de topos i feixos; \cite[§§12 i 14]{streicher} per als topos elementals i exercicis amb topos de prefeixos; \cite[article~VI, sessions~32--33]{lawvereschanuel} per a una presentació molt elemental dels subobjectes, la veritat i els topos; \cite{johnstone} com a obra de consulta avançada.
\end{resumcap}

\chapter[Pont cap a la lògica categòrica]{\texorpdfstring{Pont cap a la lògica\\ categòrica}{Pont cap a la lògica categòrica}}
\label{cap:pont}

\begin{objectius}
\apartat{Prerequisits} Tot el volum, especialment subobjectes, reindexació i categories regulars (cap.~\ref{cap:subobjectes}) i els classificadors (cap.~\ref{cap:topos}).
\apartat{Funció del capítol} Aquest capítol és un epíleg sense demostracions ni exercicis. No desenvolupa una teoria nova, però introdueix la terminologia mínima necessària per dibuixar el mapa que porta de les construccions d'aquest volum a la lògica categòrica. Serveix d'orientació per a qui vulgui continuar cap al volum següent.
\end{objectius}

Aquest volum tanca aquí. Els capítols anteriors han construït, pas a pas, el llenguatge categòric bàsic necessari per al programa que anunciem: categories i morfismes, functors, transformacions naturals i equivalències, representabilitat i el lema de Yoneda, límits i colímits, adjuncions, tancament cartesià, subobjectes, imatges i factoritzacions, i els classificadors de subobjectes. Aquest epíleg mostra què d'aquest bagatge es prolonga en l'estudi de la \textbf{lògica categòrica}, i com: cada eina apresa aquí té una contrapartida lògica precisa a l'altra banda del pont.

\section{El que aquest volum deixa preparat}

Convé fer inventari del que ja tenim, perquè és exactament la infraestructura que la lògica categòrica pressuposa:
\begin{itemize}
\item \textbf{Categories, functors i naturalitat}: el llenguatge bàsic en què s'expressa tota la resta.
\item \textbf{La categoria prima de la deduïbilitat}: la primera aparició, al capítol inicial, de la idea que una lògica es pot veure com una categoria.
\item \textbf{Representabilitat i Yoneda}: la manera d'identificar un objecte per com es relaciona amb tots els altres; el classificador de subobjectes n'és un cas.
\item \textbf{Límits finits, i molt especialment els pullbacks}: la base per interpretar contextos, substitució i relacions.
\item \textbf{Adjuncions}: el patró que reapareixerà en operacions lògiques centrals, especialment en la implicació i en els quantificadors existencial i universal.
\item \textbf{Tancament cartesià i càlcul lambda}: la semàntica dels tipus i els termes.
\item \textbf{Subobjectes i reindexació}: la interpretació dels predicats i la substitució.
\item \textbf{Imatges i factoritzacions}: la construcció que, a l'altra banda, interpretarà la quantificació existencial.
\item \textbf{Classificadors de subobjectes i topos}: el marc on la lògica esdevé interna a la categoria.
\end{itemize}

Amb tot això, el salt cap a la lògica categòrica deixa de ser un salt: és una continuació natural.

\section{El mapa: de regular a topos}\label{sec:mapa}\index{categoria!regular}\index{categoria!coherent}\index{categoria!de Heyting}

Abans de descriure el recorregut, val la pena fixar un mapa que ordena, d'una vegada, quina estructura categòrica interpreta quin fragment lògic. És la classificació que dona sentit retrospectiu a les distincions que hem anat marcant ---per exemple, per què la regularitat basta per a l'existencial però no per a la disjunció ni per a l'universal.

La idea directriu és que cada nivell \emph{afegeix} una peça d'estructura a les fibres de subobjectes $\Sub(A)$, i que cada peça interpreta un connectiu o quantificador:
\begin{center}
\renewcommand{\arraystretch}{1.3}
\noindent\begin{tabularx}{\linewidth}{@{}>{\raggedright\arraybackslash}X >{\raggedright\arraybackslash}X >{\raggedright\arraybackslash}X@{}}
\hline
\textbf{Estructura categòrica} & \textbf{Nova estructura a $\Sub(A)$} & \textbf{Fragment lògic}\\
\hline
categoria amb límits finits & $\top$, $\wedge$ (interseccions) & conjunció, veritat\\
categoria \textbf{regular} & imatges estables: $\exists_f$ & $+\ \exists$\\
categoria \textbf{coherent} & joins finits estables: $\bot$, $\vee$ & $+\ \vee,\ \bot$\\
categoria \textbf{de Heyting} & implicació $\Rightarrow$; $\forall_f$ & $+\ \Rightarrow,\ \forall$: lògica intuïcionista de primer ordre\\
\textbf{topos} elemental & classificador $\Omega$, objectes de parts, exponencials & lògica intuïcionista d'ordre superior\\
\hline
\end{tabularx}
\renewcommand{\arraystretch}{1.0}
\end{center}

La lectura de la taula, de dalt a baix, és acumulativa i respon les preguntes que aquest volum ha deixat obertes deliberadament:
\begin{itemize}
  \item \textbf{Regular} ($\top,\wedge,\exists$). La factorització imatge estable dona l'existencial com a adjunt per l'esquerra de la reindexació. És fins on arriba el capítol de subobjectes: per això hi vam construir només $\exists_f$, i no $\forall_f$ ni la disjunció.
  \item \textbf{Coherent} ($+\vee,\bot$). S'hi afegeixen \emph{joins finits estables} a les fibres, que interpreten la disjunció i la falsedat. Aquesta és l'estructura que a la teoria de subobjectes anunciàvem com a necessària per a les \enquote{unions}, i que la regularitat sola no proporciona.
  \item \textbf{De Heyting} ($+\Rightarrow,\forall$). Les fibres esdevenen àlgebres de Heyting ---amb implicació, adjunta per la dreta de la conjunció--- i cada reindexació $f^{*}$ adquireix un adjunt per la dreta $\forall_f$. Apareixen la implicació i el quantificador universal.
  \item \textbf{Topos} elemental. El classificador $\Omega$ i els exponencials clouen el quadre amb una lògica intuïcionista d'ordre superior, on fins i tot els predicats sobre predicats tenen extensió.
\end{itemize}

\index{categoria!localment cartesiana tancada}En paral·lel, la semàntica dels \emph{tipus} segueix una altra línia. Les categories cartesianes tancades interpreten el càlcul lambda simplement tipat (capítol~\ref{cap:cartesianes}), mentre que les categories \textbf{localment cartesianes tancades} proporcionen productes dependents: per a cada $f\colon A\to B$, la reindexació $f^{*}\colon\cat{C}/B\to\cat{C}/A$ admet un adjunt per la dreta $\Pi_f$. Aquesta via és especialment important per a la teoria de tipus dependents. No és un esglaó de la cadena regular--coherent--Heyting, però totes dues vies es troben en els topos elementals, que són localment cartesians tancats i tenen lògica de Heyting als subobjectes.

Aquest mapa no és matèria d'aquest volum (no en demostrem les equivalències), però és la brúixola que orienta tot el que segueix: cada fila és, de fet, un capítol del programa de la lògica categòrica.

\section{El que continua a l'altra banda}

L'estudi de la lògica categòrica pren aquestes eines i les fa servir per construir una \textbf{semàntica de la lògica} de complexitat creixent. Sense entrar-hi ---perquè és l'objecte d'un estudi a part---, en dibuixem el recorregut, perquè el lector vegi on desemboca cada cosa que ha après:

\begin{description}
\item[Lògica d'equacions.] El punt de partida: teories algebraiques a l'estil de Lawvere, on la categoria sintàctica té productes finits i els models són functors que els preserven. És la lògica que ja s'endevina en la universalitat dels productes finits: els contextos es combinen per producte i les operacions esdevenen morfismes entre potències finites d'un objecte. El tancament cartesià pertany a una via posterior, la del càlcul lambda i la implicació.

\item[Lògica de límits finits.] El pont immediat: teories amb diverses menes, operacions parcialment definides i axiomes d'existència i unicitat. La categoria sintàctica passa a tenir \emph{límits finits}, no només productes; és el marc de les teories essencialment algebraiques.

\item[Lògica regular.] Fórmules amb igualtat, conjunció i quantificació existencial, escrites com a seqüents en context. Aquí és on les \textbf{imatges regulars} construïdes en aquest volum interpreten l'existencial, i on la factorització imatge interpreta l'eliminació de testimonis. El resultat central és que, per a una teoria regular $T$ i una categoria regular $\cat{C}$, els models de $T$ en $\cat{C}$ es descriuen mitjançant functors regulars des de la categoria sintàctica regular $\cat{C}_T$ cap a $\cat{C}$; amb la noció adequada de morfisme de models, aquesta correspondència s'organitza categòricament.

\item[Lògica coherent i geomètrica.] S'hi afegeixen disjuncions finites (coherent) i disjuncions indexades per conjunts (geomètrica). És el marc que condueix als feixos i, de nou, als topos.

\item[Lògica intuïcionista de primer ordre.] Apareixen la implicació i el quantificador universal. La idea organitzadora és que les fibres de subobjectes formen \textbf{àlgebres de Heyting}: la implicació és l'adjunt dret de la conjunció, i $\forall$ és l'adjunt dret de la substitució. Els adjunts de la reindexació que aquí hem vist com a exemple es converteixen allà en els quantificadors del sistema.

\item[Doctrines i fibracions lògiques.] La síntesi: contextos, substitució, predicats i quantificadors s'organitzen en una \emph{doctrina indexada}. Les condicions de coherència ---Beck--Chevalley i Frobenius--- expressen que substituir i quantificar interactuen correctament. És el llenguatge que unifica tots els casos anteriors.

\item[Teories classificadores i llenguatge intern.] El punt culminant: la categoria o el topos classificador d'una teoria, amb la seva propietat universal, i el llenguatge intern que permet raonar dins d'una categoria com si els seus objectes fossin conjunts, sempre que les regles emprades corresponguin a l'estructura disponible.

\item[Lògica clàssica i booleanització.] Com a epíleg, què s'afegeix en imposar el tercer exclòs, quines construccions deixen de ser constructives i com, en teoria de topos, es passa a la booleanització d'un topos. La lògica clàssica apareix així com una especialització, no com el rerefons silenciós de tota la teoria.
\end{description}

\begin{nota}[Ampliació: estructures sense conjunt subjacent]
Al capítol~\ref{cap:categories} hem partit d'estructures que tenen un conjunt subjacent. No totes les estructures matemàtiques són conjunts amb una estructura afegida, i el programa anterior ho fa visible. A cada teoria geomètrica li correspon un \emph{topos classificador}, la categoria de punts del qual és equivalent a la categoria de models de la teoria a $\cat{Set}$ \cite{caramello}. Però aquest \enquote{espai} pot ser no trivial i no tenir cap punt. Per exemple, la teoria d'una aplicació exhaustiva $\mathbb N\to\mathbb R$ no té cap model a $\cat{Set}$, pel teorema de Cantor; tanmateix, el seu \emph{locale} classificador és no trivial, i el topos de feixos sobre aquest locale, que n'és el topos classificador, no té cap punt \cite{blechschmidt}. Aquestes situacions queden fora d'aquest volum; les esmentem perquè mostren que la perspectiva de les fletxes no depèn de disposar d'elements.
\end{nota}

\section{Un exemple: la categoria sintàctica d'una teoria}\label{sec:sintactica}\index{categoria!sintàctica}

Fins ara la correspondència lògica--categòrica ha aparegut repartida en exemples solts. Val la pena reunir-la en un únic exemple pont, complet i concret, que il·lustra de cop l'eslògan \enquote{una teoria es codifica en una categoria sintàctica}. Prenem la teoria algebraica més simple no trivial: una \textbf{mena} $s$ i una operació binària $m\colon s\times s\to s$, sense cap axioma (la teoria dels \emph{magmes}). Construïm pas a pas la seva categoria sintàctica $\cat{Syn}$.

\begin{enumerate}
  \item \textbf{Contextos com a objectes.} Un \emph{context} és una llista finita de variables de la mena $s$, que escrivim $[x_1,\dots,x_n]$ i abreugem $s^n$. Els objectes de $\cat{Syn}$ són els contextos, un per a cada $n\geq 0$; el context buit $s^0$ és l'objecte terminal.
  \item \textbf{Termes com a morfismes.} Un morfisme $s^n\to s^m$ és una llista de $m$ termes $(t_1,\dots,t_m)$, cadascun construït amb les variables $x_1,\dots,x_n$ i l'operació $m$ (per exemple, $m(x_1,m(x_2,x_1))$). La composició és la \textbf{substitució}: components d'un morfisme dins d'un altre. La identitat de $s^n$ és la llista de variables $(x_1,\dots,x_n)$.
  \item \textbf{Productes com a concatenació de contextos.} El producte $s^n\times s^m$ és el context $s^{n+m}$: juxtaposar dos contextos és formar-ne el producte, i les projeccions obliden les variables sobrants. Així $\cat{Syn}$ té tots els productes finits, i l'operació $m$ del llenguatge és, ara, un morfisme distingit $s\times s\to s$ de la categoria.
  \item \textbf{Equacions com a igualtats de morfismes.} Si la teoria tingués axiomes ---posem l'associativitat de $m$---, s'imposarien identificant els dos morfismes $s^3\to s$ corresponents a $m(m(x_1,x_2),x_3)$ i $m(x_1,m(x_2,x_3))$. Una equació entre termes \emph{és}, literalment, una igualtat entre morfismes de $\cat{Syn}$.
  \item \textbf{Models com a functors.} Un model de la teoria en una categoria $\cat{C}$ amb productes finits ---per exemple, un magma ordinari a $\cat{Set}$--- és exactament un functor $\cat{Syn}\to\cat{C}$ que preserva els productes finits: tria un objecte per a $s$ i un morfisme per a $m$, respectant tota la sintaxi. Els homomorfismes de models són transformacions naturals. Aquesta és la forma precisa de \enquote{els models són functors}.
  \item \textbf{Fórmules i existencial.} Per passar de la teoria algebraica anterior a la lògica regular cal enriquir també la categoria sintàctica: la categoria $\cat{Syn}$ dels passos anteriors no conté automàticament els subobjectes que representen totes les fórmules. En la \textbf{categoria sintàctica regular} $\cat{C}_T$ d'una teoria amb predicats, una fórmula $\varphi(\bar x)$ en context determina, llevat d'equivalència demostrable, un subobjecte
  \[
  [\bar x\mid\varphi]\rightarrowtail[\bar x\mid\top].
  \]
  La substitució al llarg d'un morfisme de contextos s'interpreta per \textbf{reindexació}, és a dir, per pullback. En particular, si $\pi\colon s^{n+1}\to s^n$ és la projecció que oblida la darrera variable, $\pi^{*}$ és el debilitament, que considera un predicat sobre $s^n$ en el context ampliat. Si $P\rightarrowtail s^{n+1}$ representa $\varphi(\bar x,y)$, la \textbf{imatge} del compost
  \[
  P\rightarrowtail s^{n+1}\xrightarrow{\ \pi\ }s^n
  \]
  representa $\exists y\,\varphi(\bar x,y)$. Així, l'existencial és la imatge d'un predicat al llarg d'una projecció o, equivalentment, l'adjunt per l'esquerra de la reindexació (capítol~\ref{cap:subobjectes}).
  \item \textbf{Un càlcul complet en un model.} Fem concret tot l'anterior en un sol model, el magma $(\mathbb N,+)$, és a dir, el functor $M\colon\cat{Syn}\to\cat{Set}$ amb $M(s)=\mathbb N$ i $M(m)=+$. Prenem la fórmula $\varphi(x,y)$: $m(y,y)=x$, en el context $[x,y]$.
  \begin{itemize}
    \item \emph{Una equació és un equalitzador.} Els dos termes $m(y,y)$ i $x$ són morfismes $s^2\to s$, interpretats com les aplicacions $(x,y)\mapsto y+y$ i $(x,y)\mapsto x$ de $\mathbb N^2$ a $\mathbb N$. La fórmula s'interpreta com el seu equalitzador,
    \[
    P=\{(x,y)\in\mathbb N^2\mid y+y=x\}\rightarrowtail\mathbb N^2.
    \]
    Equivalentment, $P$ és el pullback de la diagonal $\delta_{\mathbb N}\colon\mathbb N\rightarrowtail\mathbb N\times\mathbb N$ al llarg del parell format pels dos termes: la igualtat és la diagonal.
    \item \emph{L'existencial és una imatge.} La projecció $\pi\colon\mathbb N^2\to\mathbb N$ oblida $y$. La imatge de $P\rightarrowtail\mathbb N^2\to\mathbb N$ és $\{x\mid\exists y\ y+y=x\}$, el conjunt dels nombres parells. Així, $\exists y\,(m(y,y)=x)$ s'interpreta com el predicat \enquote{$x$ és parell}, és a dir, com $\exists_\pi[P]$.
    \item \emph{El predicat és un morfisme cap a $\Omega$.} El seu morfisme característic $\mathbb N\to\{0,1\}$ val $1$ exactament als parells, com a l'exemple~\ref{ex:caracteristic-set}.
    \item \emph{La substitució és un pullback.} Substituir $x$ pel terme $m(z,m(z,z))$, en el context $[z]$, dona la fórmula $\exists y\,(m(y,y)=m(z,m(z,z)))$. Categòricament, és reindexar el predicat al llarg del morfisme $s\to s$ definit pel terme, interpretat com $z\mapsto 3z$. Per tant, s'interpreta com $\{z\mid 3z\text{ és parell}\}$, que és el conjunt dels parells. En canvi, substituir pel terme $m(z,z)$ dona $\{z\mid 2z\text{ és parell}\}=\mathbb N$, el subobjecte màxim: la fórmula $\exists y\,(m(y,y)=m(z,z))$ és demostrable, amb el testimoni $y=z$, i el model ho reflecteix.
  \end{itemize}
  Un sol exemple recorre així les peces del volum: termes com a morfismes, equacions com a equalitzadors i igualtat com a diagonal (capítols~\ref{cap:limits} i~\ref{cap:cartesianes}), existencial com a imatge (capítol~\ref{cap:subobjectes}), predicats com a morfismes cap a $\Omega$ (capítol~\ref{cap:topos}) i substitució com a pullback.
\end{enumerate}

Amb aquest exemple, cada peça del llibre troba el seu lloc: els contextos són objectes, els termes són morfismes, la substitució és composició i reindexació, els productes són concatenacions, les equacions són igualtats de fletxes, els models són functors que preserven estructura, i els quantificadors, quan existeixen, són adjunts de la reindexació. La categoria sintàctica és el punt on la sintaxi lògica i la semàntica categòrica es toquen; construir-la en general, amb l'estructura exacta que cada fragment lògic requereix, és una de les primeres tasques del volum següent.

\section{Cloenda}

El fil que recorre tot aquest programa és un de sol, i ja el coneixem: \emph{traduir entre lògica i categories}. En el marc categòric adequat a cada fragment lògic, un predicat s'interpreta com un subobjecte; la substitució, com un pullback; la conjunció, com una intersecció; l'existencial, com una imatge; la implicació i l'universal, mitjançant adjunts; una teoria es codifica en una categoria sintàctica i els seus models es descriuen per functors que preserven l'estructura pertinent. Les teories algebraiques, regulars, coherents o geomètriques tenen categories sintàctiques amb l'estructura corresponent, i per això teories sintàcticament diferents poden tenir semàntiques categòriques equivalents, o fins i tot el mateix topos classificador. Cadascuna d'aquestes traduccions descansa sobre una construcció que aquest volum ha desenvolupat amb detall.

No hem demostrat cap teorema de la lògica categòrica en aquestes pàgines, i era deliberat: la seva demostració pertany al volum següent. El que sí que hem fet és deixar el terreny preparat, de manera que quan aquells teoremes arribin, cap de les seves peces categòriques resulti desconeguda. Si el lector, en trobar més endavant que \enquote{l'existencial és un adjunt per l'esquerra de la substitució, estable per canvi de base}, reconeix en aquesta frase la factorització imatge, l'adjunció i el pullback que ha estudiat aquí, aleshores aquest volum haurà complert la seva funció.

\section{Mapa del volum següent}\label{sec:mapa-volum-seguent}

La secció~\ref{sec:mapa} ordenava la continuació per fragments lògics. Per acabar, l'ordenem per temes: per a cadascun dels cinc temes centrals del volum de lògica categòrica, què en queda establert aquí i què s'hi afegirà.
\begin{center}
\footnotesize
\renewcommand{\arraystretch}{1.3}
\noindent\begin{tabularx}{\linewidth}{@{}>{\raggedright\arraybackslash}p{0.2\linewidth}>{\raggedright\arraybackslash}X>{\raggedright\arraybackslash}X@{}}
\toprule
\textbf{Tema} & \textbf{En aquest volum} & \textbf{Al volum següent}\\
\midrule
substitució i reindexació & $f^{*}$ per pullback a $\cat{C}/B$ (capítol~\ref{cap:limits}) i a $\Sub(B)$; el functor $\Sub\colon\cat{C}\op\to\cat{Poset}$ (capítol~\ref{cap:subobjectes}) & la substitució de termes en fórmules com a reindexació en una doctrina; la condició de Beck--Chevalley\\
connectius sobre subobjectes & $\top$ i $\wedge$ com a pullbacks (capítol~\ref{cap:subobjectes}); $\Rightarrow$ als preordres cartesians tancats (capítol~\ref{cap:cartesianes}); $\wedge$ sobre $\Omega$ (capítol~\ref{cap:topos}) & $\bot$ i $\vee$ estables (categories coherents); $\Rightarrow$ a les fibres (categories de Heyting)\\
quantificadors com a adjunts & $\exists_f\dashv f^{*}\dashv\forall_f$ a $\cat{Set}$ (capítol~\ref{cap:adjuncions}); $\exists_f$ com a imatge en categories regulars (capítol~\ref{cap:subobjectes}) & $\forall_f$ en categories de Heyting; la condició de Frobenius; la quantificació al llarg de projeccions de contextos\\
igualtat & la diagonal $\delta_A$ (capítol~\ref{cap:cartesianes}); la parella nucli, que és el pullback de $\delta_B$ al llarg de $f\times f$ (capítol~\ref{cap:subobjectes}); les equacions com a equalitzadors (secció~\ref{sec:sintactica}) & la igualtat com a $\exists_{\delta_A}(\top)$, on $\exists_{\delta_A}$ és l'adjunt per l'esquerra de la contracció $\delta_A^{*}$; les regles de la igualtat com a propietats d'aquesta adjunció\\
semàntica interna & $\Omega$, la pertinença per elements generalitzats i l'objecte de parts (capítol~\ref{cap:topos}); l'exemple de la secció~\ref{sec:sintactica} & el llenguatge intern d'un topos, la semàntica de Kripke--Joyal, la lògica d'ordre superior, i la correcció i la completesa\\
\bottomrule
\end{tabularx}
\end{center}
Cap fila de la columna central no demana res que el lector no hagi vist en aquest volum. La columna de la dreta és el programa del volum següent.

\begin{resumcap}
\apartat{El fil en una frase} En el marc categòric adequat a cada fragment lògic, un predicat s'interpreta com un subobjecte; la substitució, com un pullback; la conjunció, com una intersecció; l'existencial, com una imatge; la implicació i l'universal, mitjançant adjunts; una teoria es codifica en una categoria sintàctica i els seus models es descriuen per functors que preserven l'estructura pertinent.
\apartat{Cap on continua} La lògica pròpiament dita ---sintaxi, correcció i completesa, i les condicions de coherència Beck--Chevalley i Frobenius--- es reserva per al volum de lògica categòrica, del qual aquest text és la preparació. La secció~\ref{sec:mapa-volum-seguent} en dona el mapa per temes: substitució, connectius, quantificadors, igualtat i semàntica interna.
\apartat{Lectura recomanada} \cite{lambekscott} i \cite{jacobs} per a la lògica categòrica i la teoria de tipus; \cite[§13]{streicher} per a la lògica interna dels topos; \cite{makkaireyes} i \cite{pitts} per a la lògica de primer ordre en categories regulars, coherents i de Heyting, que és el punt de partida del volum de lògica; \cite{maclanemoerdijk} per a la lògica geomètrica i els topos; \cite{johnstone} com a obra de consulta. Com a direccions complementàries: \cite{hott}, per a la teoria homotòpica de tipus, on la igualtat entre dos elements pot tenir estructura pròpia; i l'nLab \cite{nlab-classificador,nlab-locale}, com a recurs de consulta sobre topos classificadors i locales.
\end{resumcap}

\fi

% ============================================================
% PART II — PRÀCTICA
% Els capítols reprodueixen els noms i l'ordre de la part teòrica.
% ============================================================
\part{Pràctica}
\setcounter{chapter}{0}
\chapter{Categories}

En aquest capítol resolem amb detall una selecció d'exercicis corresponents al capítol \emph{Categories} de la part de Teoria. L'objectiu no és només arribar al resultat, sinó mostrar \emph{com} es raona en el llenguatge categòric: comprovar axiomes, seguir camins, traduir diagrames a equacions i reconèixer una mateixa estructura en conjunts, preordres, lògica i grafs. Convé intentar cada exercici abans de llegir-ne la solució.

\section{Definició de categoria}

\begin{exercici}\label{exr:pr-rel}
Comprova que els conjunts amb les relacions formen una categoria $\cat{Rel}$, on un morfisme $A\to B$ és una relació $R\subseteq A\times B$, la identitat de $A$ és la relació d'igualtat, i la composició de $R\subseteq A\times B$ amb $S\subseteq B\times C$ és
\[
S\comp R=\{(a,c)\in A\times C\mid \exists\,b\in B\text{ tal que }(a,b)\in R\text{ i }(b,c)\in S\}.
\]
\end{exercici}

\begin{solucio}
Hem de comprovar els dos axiomes. Comencem per les \textbf{identitats}. Sigui $\id_A=\{(a,a)\mid a\in A\}$ la relació d'igualtat sobre $A$. Per a tota relació $R\subseteq A\times B$,
\[
R\comp\id_A
=\{(a,b)\mid \exists\,a'\ (a,a')\in\id_A\text{ i }(a',b)\in R\}
=R,
\]
perquè $(a,a')\in\id_A$ obliga $a'=a$. Anàlogament $\id_B\comp R=R$.

Per a l'\textbf{associativitat}, siguin $R\subseteq A\times B$, $S\subseteq B\times C$ i $T\subseteq C\times D$. Un parell $(a,d)$ pertany a $T\comp(S\comp R)$ si i només si existeixen $b\in B$ i $c\in C$ tals que
\[
(a,b)\in R,
\qquad
(b,c)\in S,
\qquad
(c,d)\in T.
\]
La mateixa condició caracteritza $(T\comp S)\comp R$. Per tant
\[
T\comp(S\comp R)=(T\comp S)\comp R,
\]
i $\cat{Rel}$ és una categoria.
\end{solucio}

\section{Fletxes, camins i diagrames commutatius}

\begin{exercici}
Llegeix el diagrama següent i escriu l'equació que expressa que commuta:
\[
\begin{tikzpicture}[baseline=(m.center)]
  \matrix (m) [matrix of math nodes, row sep=2.7em, column sep=3.6em] {
    A & B \\
    C & D \\
  };
  \draw[-{Stealth[scale=1.1]}] (m-1-1) -- node[above] {$f$} (m-1-2);
  \draw[-{Stealth[scale=1.1]}] (m-1-1) -- node[left] {$u$} (m-2-1);
  \draw[-{Stealth[scale=1.1]}] (m-1-2) -- node[right] {$v$} (m-2-2);
  \draw[-{Stealth[scale=1.1]}] (m-2-1) -- node[below] {$g$} (m-2-2);
\end{tikzpicture}
\]
Explica també per què l'ordre de les composicions és el que escrius.
\end{exercici}

\begin{solucio}
Hi ha dos camins dirigits de $A$ a $D$. El camí superior-dret és
\[
A\xrightarrow{f}B\xrightarrow{v}D,
\]
i per tant determina $v\comp f$. El camí esquerre-inferior és
\[
A\xrightarrow{u}C\xrightarrow{g}D,
\]
i determina $g\comp u$. El quadrat commuta exactament quan
\[
\boxed{v\comp f=g\comp u}.
\]
L'ordre de cada composició es llegeix de dreta a esquerra: en $v\comp f$ seguim primer $f$ i després $v$.
\end{solucio}

\begin{exercici}
Considera a $\cat{Set}$ quatre còpies de $\mathbb Z$ i les aplicacions
\[
f(n)=n+1,
\qquad
u(n)=2n,
\qquad
v(n)=2n,
\qquad
g(n)=n+2.
\]
Comprova que el quadrat
\[
\begin{tikzpicture}[baseline=(m.center)]
  \matrix (m) [matrix of math nodes, row sep=2.8em, column sep=4.2em] {
    \mathbb Z & \mathbb Z \\
    \mathbb Z & \mathbb Z \\
  };
  \draw[-{Stealth[scale=1.05]}] (m-1-1) -- node[above] {$f$} (m-1-2);
  \draw[-{Stealth[scale=1.05]}] (m-1-1) -- node[left] {$u$} (m-2-1);
  \draw[-{Stealth[scale=1.05]}] (m-1-2) -- node[right] {$v$} (m-2-2);
  \draw[-{Stealth[scale=1.05]}] (m-2-1) -- node[below] {$g$} (m-2-2);
\end{tikzpicture}
\]
commuta.
\end{exercici}

\begin{solucio}
Cal comprovar que $v\comp f=g\comp u$. Per a qualsevol $n\in\mathbb Z$,
\[
(v\comp f)(n)=v(n+1)=2n+2,
\]
mentre que
\[
(g\comp u)(n)=g(2n)=2n+2.
\]
Les dues aplicacions coincideixen en tots els enters. Per tant
\[
v\comp f=g\comp u,
\]
i el quadrat commuta. Aquest exemple mostra com una afirmació sobre funcions es pot llegir simultàniament com una igualtat algebraica i com la commutativitat d'un diagrama.
\end{solucio}

\section{Primers exemples}

\begin{exercici}
Sigui $T$ un sistema deductiu amb una relació de deduïbilitat $\vdash_T$ reflexiva i transitiva. Comprova amb detall que les fórmules, amb una única fletxa $A\to B$ quan $A\vdash_T B$, formen una categoria prima $\cat{Prov}_T$.
\end{exercici}

\begin{solucio}
Els \textbf{objectes} són les fórmules. Entre dues fórmules $A$ i $B$ posem un únic morfisme $A\to B$ exactament quan $A\vdash_T B$.

La \textbf{identitat} existeix perquè la deduïbilitat és reflexiva:
\[
A\vdash_T A.
\]
Per tant hi ha una fletxa $\id_A\colon A\to A$.

La \textbf{composició} existeix perquè la deduïbilitat és transitiva. Si
\[
A\vdash_T B
\qquad\text{i}\qquad
B\vdash_T C,
\]
aleshores
\[
A\vdash_T C.
\]
Diagramàticament,
\[
\begin{tikzpicture}[baseline=(m.center)]
  \matrix (m) [matrix of math nodes, column sep=3.4em] { A & B & C \\ };
  \draw[-{Stealth[scale=1.05]}] (m-1-1) -- (m-1-2);
  \draw[-{Stealth[scale=1.05]}] (m-1-2) -- (m-1-3);
  \draw[-{Stealth[scale=1.0]}] (m-1-1) to[bend right=25] (m-1-3);
\end{tikzpicture}
\]

Finalment, entre dues fórmules hi ha com a màxim una fletxa. Això fa que les igualtats exigides pels axiomes d'associativitat i identitat siguin automàtiques sempre que les composicions existeixin. Per tant $\cat{Prov}_T$ és una categoria prima.

Cal remarcar què hem oblidat: si hi ha dues demostracions diferents de $B$ a partir d'$A$, aquesta categoria les identifica en una sola fletxa. Només conserva la relació de deduïbilitat.
\end{solucio}

\begin{exercici}
Siguin els grafs dirigits $G$ i $H$ següents:
\[
G:\quad 0\xrightarrow{a}1,
\qquad
H:\quad
\begin{tikzpicture}[baseline=(m.center)]
  \matrix (m) [matrix of math nodes, row sep=2.0em, column sep=2.8em] {
    & y \\
    x & z \\
  };
  \draw[-{Stealth[scale=1.0]}] (m-2-1) -- node[above left] {$r$} (m-1-2);
  \draw[-{Stealth[scale=1.0]}] (m-2-1) -- node[below] {$s$} (m-2-2);
\end{tikzpicture}
\]
Defineix $F_V(0)=x$, $F_V(1)=y$ i $F_E(a)=r$. Comprova que això defineix un morfisme de grafs $F\colon G\to H$. Explica per què substituir $F_E(a)=r$ per $F_E(a)=s$, mantenint $F_V$, no defineix un morfisme de grafs.
\end{exercici}

\begin{solucio}
En $G$ tenim
\[
s_G(a)=0,
\qquad
t_G(a)=1,
\]
i en $H$,
\[
s_H(r)=x,
\qquad
t_H(r)=y.
\]
Per tant
\[
F_V(s_G(a))=F_V(0)=x=s_H(r)=s_H(F_E(a))
\]
i
\[
F_V(t_G(a))=F_V(1)=y=t_H(r)=t_H(F_E(a)).
\]
Així es compleixen les dues equacions
\[
F_V\comp s_G=s_H\comp F_E,
\qquad
F_V\comp t_G=t_H\comp F_E,
\]
i $F$ és un morfisme de grafs.

Si poséssim $F_E(a)=s$, l'origen encara seria correcte perquè $s_H(s)=x$, però el destí fallaria:
\[
t_H(s)=z\neq y=F_V(1).
\]
Per tant el quadrat corresponent al destí no commutaria i no tindríem un morfisme de grafs.
\end{solucio}

\section{Isomorfismes}

\begin{exercici}
En una categoria qualsevol, comprova que si $f\colon A\to B$ i $g\colon B\to C$ tenen inversos, aleshores $g\comp f$ també en té, i
\[
(g\comp f)^{-1}=f^{-1}\comp g^{-1}.
\]
\end{exercici}

\begin{solucio}
La situació és
\[
A\xrightarrow{f}B\xrightarrow{g}C.
\]
Per recórrer aquest camí en sentit contrari hem d'invertir també l'ordre: primer $g^{-1}$ i després $f^{-1}$. Proposem, doncs, $f^{-1}\comp g^{-1}$ com a invers.

D'una banda,
\[
(f^{-1}\comp g^{-1})\comp(g\comp f)
=f^{-1}\comp(g^{-1}\comp g)\comp f
=f^{-1}\comp f
=\id_A,
\]
i de l'altra,
\[
(g\comp f)\comp(f^{-1}\comp g^{-1})
=g\comp(f\comp f^{-1})\comp g^{-1}
=g\comp g^{-1}
=\id_C.
\]
Per tant $(g\comp f)^{-1}=f^{-1}\comp g^{-1}$.
\end{solucio}

\begin{exercici}
Comprova que a $\cat{Rel}$ una relació $R\colon A\to B$ és un isomorfisme si i només si és el graf d'una bijecció de $A$ a $B$.
\end{exercici}

\begin{solucio}
Si $f\colon A\to B$ és una bijecció i $R$ és el seu graf, el graf de $f^{-1}$ és una relació $B\to A$ i les dues composicions són les relacions d'igualtat. Per tant $R$ és un isomorfisme.

Recíprocament, suposem que $R$ té un invers $S\colon B\to A$, és a dir,
\[
S\comp R=\id_A,
\qquad
R\comp S=\id_B.
\]
Per a cada $a\in A$, com $(a,a)\in S\comp R$, existeix algun $b\in B$ amb $(a,b)\in R$ i $(b,a)\in S$. Així cada $a$ es relaciona amb almenys un $b$.

Si $(a,b)\in R$ i $(a,b')\in R$, triem $b_0$ amb $(a,b_0)\in R$ i $(b_0,a)\in S$. Llavors
\[
(b_0,b)\in R\comp S=\id_B,
\qquad
(b_0,b')\in R\comp S=\id_B,
\]
i per tant $b=b_0=b'$. Així cada $a$ es relaciona amb un únic $b$, i $R$ és el graf d'una aplicació $f\colon A\to B$.

Aplicat a $S$, el mateix argument dona una aplicació $g\colon B\to A$. Les igualtats $S\comp R=\id_A$ i $R\comp S=\id_B$ es converteixen en
\[
g\comp f=\id_A,
\qquad
f\comp g=\id_B.
\]
Per tant $f$ és bijectiva i $R$ n'és el graf.
\end{solucio}

\begin{exercici}
Sigui $\cat{C}$ una categoria. Comprova que la relació \enquote{ésser isomorf} entre objectes de $\cat{C}$ és una relació d'equivalència.
\end{exercici}

\begin{solucio}
\textbf{Reflexivitat}: $\id_A$ és el seu propi invers, de manera que $A\cong A$.

\textbf{Simetria}: si $f\colon A\to B$ és un isomorfisme amb invers $f^{-1}\colon B\to A$, aleshores $f^{-1}$ és també un isomorfisme, amb invers $f$; per tant $B\cong A$.

\textbf{Transitivitat}: si $f\colon A\to B$ i $g\colon B\to C$ són isomorfismes, la composició $g\comp f$ també ho és pel primer exercici de la secció; per tant $A\cong C$.
\end{solucio}

\begin{exercici}
A la categoria prima $\cat{Prov}_T$, comprova que dues fórmules $A$ i $B$ són isomorfes si i només si són interdeduïbles:
\[
A\vdash_T B
\qquad\text{i}\qquad
B\vdash_T A.
\]
\end{exercici}

\begin{solucio}
Si $A\cong B$, per definició existeixen fletxes
\[
A\longrightarrow B
\qquad\text{i}\qquad
B\longrightarrow A.
\]
A $\cat{Prov}_T$ això significa exactament $A\vdash_T B$ i $B\vdash_T A$.

Recíprocament, si es donen aquestes dues deduïbilitats, existeixen les dues fletxes. Com que $\cat{Prov}_T$ és prima, qualsevol endomorfisme $A\to A$ és necessàriament l'única fletxa $\id_A$, i igualment per a $B$. Per tant les dues composicions són les identitats i les fletxes són inverses. Així $A\cong B$.
\end{solucio}

\section{Classes especials de morfismes}

\begin{exercici}
Demostra que un morfisme $f$ és un isomorfisme si i només si és un monomorfisme escindit (té retracció) i un epimorfisme.
\end{exercici}

\begin{solucio}
Si $f$ és un isomorfisme, el seu invers és alhora retracció i secció, de manera que $f$ és mono escindit; i tot isomorfisme és epi. Recíprocament, suposem que $f\colon A\to B$ té una retracció $r$ (amb $r\comp f=\id_A$) i que és epi. Partint de $r\comp f=\id_A$ i component amb $f$ per l'esquerra,
\[
(f\comp r)\comp f = f\comp(r\comp f) = f\comp\id_A = f = \id_B\comp f.
\]
Com que $f$ és epi, es pot cancel·lar el factor $f$ de la dreta, d'on
\[
f\comp r=\id_B.
\]
Juntament amb $r\comp f=\id_A$, això mostra que $r$ és també invers per la dreta de $f$: $f$ és un isomorfisme amb invers $r=f^{-1}$.
\end{solucio}

\begin{exercici}
A $\cat{Set}$, considera $f\colon\{0,1\}\to\{\ast\}$, l'única aplicació, i $g\colon\{\ast\}\to\{0,1\}$ amb $g(\ast)=0$. Comprova quina de les dues composicions $f\comp g$ i $g\comp f$ és una identitat, digues quin dels dos morfismes és una secció i quin una retracció (i de qui), i determina quin dels dos és un monomorfisme i quin un epimorfisme.
\end{exercici}

\begin{solucio}
Calculem les dues composicions. Per una banda, $f\comp g\colon\{\ast\}\to\{\ast\}$ compleix $(f\comp g)(\ast)=f(0)=\ast$, de manera que $f\comp g=\id_{\{\ast\}}$. Per l'altra, $g\comp f\colon\{0,1\}\to\{0,1\}$ envia tant $0$ com $1$ a $0$, de manera que $g\comp f\neq\id_{\{0,1\}}$ (per exemple, $(g\comp f)(1)=0\neq 1$).

Així, l'única igualtat que val és $f\comp g=\id_{\{\ast\}}$. Llegida des de cada morfisme, aquesta equació diu dues coses alhora: $g$ és una \emph{secció} de $f$ (un invers per la dreta de $f$), i $f$ és una \emph{retracció} de $g$ (un invers per l'esquerra de $g$). Convé insistir que \enquote{secció} i \enquote{retracció} són papers relatius a l'altre morfisme, no propietats absolutes: aquí $g$ no és \enquote{una secció} en abstracte, sinó una secció \emph{de $f$}.

Pel que fa a mono i epi: com que $f$ té una secció, $f$ és un epimorfisme; com a aplicació, és exhaustiva, d'acord amb la caracterització dels epimorfismes de $\cat{Set}$. Com que $g$ té una retracció, $g$ és un monomorfisme; com a aplicació, és injectiva. En canvi, $f$ no és mono, perquè l'aplicació $f$ no és injectiva ($f(0)=f(1)$), i $g$ no és epi, perquè l'aplicació $g$ no és exhaustiva ($1\notin g(\{\ast\})$). Cap dels dos és un isomorfisme: ho serien només si $g\comp f$ fos també una identitat, cosa que hem vist que no passa.
\end{solucio}

\begin{exercici}\label{exr:pr-prima-mono-epi}
Sigui $\cat{C}$ una categoria prima, és a dir, una categoria en què entre dos objectes hi ha com a màxim un morfisme. Comprova que tot morfisme de $\cat{C}$ és alhora un monomorfisme i un epimorfisme. Explica per què aquest fet mostra que, en una categoria arbitrària, \enquote{mono} i \enquote{epi} no es poden interpretar com \enquote{injectiu} i \enquote{exhaustiu}.
\end{exercici}

\begin{solucio}
Sigui $f\colon A\to B$ un morfisme de $\cat{C}$. Per veure que és un \textbf{monomorfisme}, siguin $g,h\colon X\to A$ amb $f\comp g=f\comp h$:
\[
\begin{tikzpicture}[baseline=(m.center)]
  \matrix (m) [matrix of math nodes, column sep=3.4em] { X & A & B \\ };
  \draw[-{Stealth[scale=1.0]}] ([yshift=2.6pt]m-1-1.east) -- node[above] {$g$} ([yshift=2.6pt]m-1-2.west);
  \draw[-{Stealth[scale=1.0]}] ([yshift=-2.6pt]m-1-1.east) -- node[below] {$h$} ([yshift=-2.6pt]m-1-2.west);
  \draw[-{Stealth[scale=1.1]}] (m-1-2) -- node[above] {$f$} (m-1-3);
\end{tikzpicture}
\qquad
f\comp g=f\comp h\ \Longrightarrow\ g=h.
\]
Els morfismes $g$ i $h$ tenen el mateix domini $X$ i el mateix codomini $A$: tots dos pertanyen a $\Hom(X,A)$. Com que la categoria és prima, aquest homconjunt té com a màxim un element, i per tant $g=h$. Observem que ni tan sols hem fet servir la hipòtesi $f\comp g=f\comp h$: la cancel·lació es compleix perquè no hi ha res a cancel·lar.

Per veure que és un \textbf{epimorfisme}, siguin $g,h\colon B\to Y$ amb $g\comp f=h\comp f$:
\[
\begin{tikzpicture}[baseline=(m.center)]
  \matrix (m) [matrix of math nodes, column sep=3.4em] { A & B & Y \\ };
  \draw[-{Stealth[scale=1.1]}] (m-1-1) -- node[above] {$f$} (m-1-2);
  \draw[-{Stealth[scale=1.0]}] ([yshift=2.6pt]m-1-2.east) -- node[above] {$g$} ([yshift=2.6pt]m-1-3.west);
  \draw[-{Stealth[scale=1.0]}] ([yshift=-2.6pt]m-1-2.east) -- node[below] {$h$} ([yshift=-2.6pt]m-1-3.west);
\end{tikzpicture}
\qquad
g\comp f=h\comp f\ \Longrightarrow\ g=h.
\]
Ara $g,h\in\Hom(B,Y)$, i el mateix argument dona $g=h$. També ho podem obtenir per dualitat: la categoria oposada $\cat{C}\op$ també és prima, perquè $\Hom_{\cat{C}\op}(B,A)$ està en bijecció amb $\Hom_{\cat{C}}(A,B)$; per tant $f\op$ és mono a $\cat{C}\op$, és a dir, $f$ és epi a $\cat{C}$ (teorema~\ref{teo:dualitat}).

\emph{Per què això impedeix llegir mono com a injectiu i epi com a exhaustiu.} Hi ha dues raons, una més radical que l'altra.

Primer, en una categoria prima els morfismes no són, en general, aplicacions. En un preordre, la fletxa $x\to y$ només registra que $x\le y$; no envia cap element enlloc, i preguntar si és \enquote{injectiva} no té sentit. Mono i epi, en canvi, sempre en tenen, perquè es defineixen només amb la composició.

Segon, fins i tot si volguéssim mantenir l'analogia, entraria en contradicció. A $\cat{Set}$, una aplicació injectiva i exhaustiva és bijectiva, és a dir, un isomorfisme. Si \enquote{mono $=$ injectiu} i \enquote{epi $=$ exhaustiu} valguessin en general, tot morfisme d'una categoria prima seria un isomorfisme. Però a $\cat{2}=(0\to 1)$ la fletxa $u\colon 0\to 1$ és mono i epi pel que acabem de demostrar, i no té invers, perquè no hi ha cap fletxa $1\to 0$. En un preordre qualsevol passa el mateix: la fletxa $x\to y$ és sempre mono i epi, i només és un isomorfisme quan també $y\le x$.

La conclusió és que \enquote{mono $=$ injectiu} i \enquote{epi $=$ exhaustiu} són teoremes sobre $\cat{Set}$, no definicions. Les definicions són les de cancel·lació, i en una categoria prima es compleixen de la manera més trivial possible.
\end{solucio}

\begin{exercici}\label{pr-mono-no-escindit}
Compara les aplicacions injectives $i\colon\varnothing\to\{\ast\}$ i $j\colon\{0\}\hookrightarrow\{0,1\}$ de $\cat{Set}$. Totes dues són monomorfismes? Quina d'elles és un monomorfisme escindit?
\end{exercici}

\begin{solucio}
Totes dues són injectives i, per tant, monomorfismes de $\cat{Set}$.

La inclusió $j$ té una retracció: l'aplicació $r\colon\{0,1\}\to\{0\}$ constant compleix $r\comp j=\id_{\{0\}}$. És, doncs, un monomorfisme escindit.

L'aplicació $i$, en canvi, no en té cap: una retracció seria una aplicació $\{\ast\}\to\varnothing$, i no n'hi ha cap, perquè $\ast$ no té on anar. Per tant, $i$ és un monomorfisme que no és escindit. L'exemple mostra que l'afirmació \enquote{tota aplicació injectiva té una inversa per l'esquerra} només és certa per a dominis no buits; i aquest contraexemple no depèn de l'axioma de l'elecció, a diferència del cas dual de les seccions.
\end{solucio}

\section{Construccions sobre categories}

\begin{exercici}
Descriu explícitament la categoria producte $\cat{2}\times\cat{3}$, on $\cat{2}$ és la categoria associada a l'ordre $0<1$ i $\cat{3}$ la categoria associada a $0<1<2$. Quants objectes i quants morfismes té? Com es componen?
\end{exercici}

\begin{solucio}
Els objectes són les parelles $(i,j)$ amb $i\in\{0,1\}$ i $j\in\{0,1,2\}$. N'hi ha sis i es poden disposar en una graella:
\[
\begin{tikzpicture}[baseline=(m.center)]
  \matrix (m) [matrix of math nodes, row sep=2.4em, column sep=2.8em] {
    (0,0) & (0,1) & (0,2) \\
    (1,0) & (1,1) & (1,2) \\
  };
  \draw[-{Stealth[scale=1.05]}] (m-1-1) -- (m-1-2);
  \draw[-{Stealth[scale=1.05]}] (m-1-2) -- (m-1-3);
  \draw[-{Stealth[scale=1.05]}] (m-2-1) -- (m-2-2);
  \draw[-{Stealth[scale=1.05]}] (m-2-2) -- (m-2-3);
  \draw[-{Stealth[scale=1.05]}] (m-1-1) -- (m-2-1);
  \draw[-{Stealth[scale=1.05]}] (m-1-2) -- (m-2-2);
  \draw[-{Stealth[scale=1.05]}] (m-1-3) -- (m-2-3);
\end{tikzpicture}
\]
Només hi hem dibuixat les fletxes bàsiques; les composicions produeixen també fletxes més llargues i diagonals.

Com que $\cat{2}$ i $\cat{3}$ són primes, entre $(i,j)$ i $(i',j')$ hi ha un únic morfisme exactament quan
\[
i\le i'
\qquad\text{i}\qquad
j\le j'.
\]
$\cat{2}$ té tres morfismes i $\cat{3}$ en té sis. Com que un morfisme del producte és una parella, hi ha
\[
3\cdot6=18
\]
morfismes. La composició es fa component a component. Tots els quadrats de la graella commuten perquè el producte és també una categoria prima.
\end{solucio}

\begin{exercici}
Descriu explícitament la categoria oposada $\cat{P}\op$ quan $\cat{P}$ és un preordre $(P,\leq)$ vist com a categoria. Quina relació d'ordre li correspon?
\end{exercici}

\begin{solucio}
A $\cat{P}$ hi ha un morfisme $x\to y$ exactament quan $x\leq y$. A $\cat{P}\op$ les fletxes s'inverteixen:
\[
x\longrightarrow y\text{ a }\cat{P}\op
\quad\Longleftrightarrow\quad
y\longrightarrow x\text{ a }\cat{P}
\quad\Longleftrightarrow\quad
y\leq x.
\]
Per tant $\cat{P}\op$ és el preordre invers $(P,\geq)$. En aquest exemple, invertir les fletxes és literalment invertir la desigualtat.
\end{solucio}

\begin{exercici}
Sigui $\cat{2}=(0\to 1)$ la categoria amb dos objectes i una fletxa no trivial. Prenent com a ambient la mateixa $\cat{2}$, descriu explícitament la categoria sobre l'objecte $1$, $\cat{2}/1$, i la categoria sota l'objecte $1$, $1/\cat{2}$, comptant-ne objectes i morfismes. Comprova després que $(\cat{2}/1)\op$ \emph{no} és isomorfa a $1/\cat{2}$, i identifica amb quina categoria sí que ho és.
\end{exercici}

\begin{solucio}
Escrivim la fletxa no trivial de $\cat{2}$ com $u\colon 0\to 1$.

\emph{La categoria $\cat{2}/1$} (sobre $1$). Els seus objectes són els morfismes de $\cat{2}$ amb codomini $1$: la fletxa $u\colon 0\to 1$ i la identitat $\id_1\colon 1\to 1$. Són, doncs, \emph{dos} objectes. Un morfisme de $u$ a $\id_1$ és una fletxa $h$ de $\cat{2}$ amb $\id_1\comp h=u$, és a dir $h=u$; i les identitats de cada objecte completen els morfismes. La categoria $\cat{2}/1$ té la forma d'una fletxa entre dos objectes: és isomorfa a $\cat{2}$.

\emph{La categoria $1/\cat{2}$} (sota $1$). Els seus objectes són els morfismes de $\cat{2}$ amb domini $1$: només $\id_1$, perquè de $1$ no en surt cap altra fletxa. És, doncs, \emph{un} sol objecte, amb només la seva identitat: la categoria $\cat{1}$.

\emph{El contrast.} $\cat{2}/1$ té dos objectes i $1/\cat{2}$ en té un. Per tant $(\cat{2}/1)\op$ té dos objectes i \emph{no} pot ser isomorfa a $1/\cat{2}$. El que sí que val és la fórmula general de dualitat
\[
(\cat{2}/1)\op\;\cong\;1/(\cat{2}\op).
\]
En efecte, $\cat{2}\op$ és la fletxa invertida $1\to 0$; sota l'objecte $1$ a $\cat{2}\op$ hi ha els morfismes amb domini $1$, que són $\id_1$ i la fletxa $1\to 0$: \emph{dos} objectes, com $(\cat{2}/1)\op$. La moral: en dualitzar la construcció \enquote{sobre} en \enquote{sota} cal dualitzar també la categoria ambient, no deixar-la intacta.
\end{solucio}

\section{Principi de dualitat}

\begin{exercici}\label{exr:pr-dual-mono-epi}
Demostra que la composició de dos monomorfismes és un monomorfisme. A continuació, sense repetir el càlcul, dedueix per dualitat que la composició de dos epimorfismes és un epimorfisme. Indica exactament què s'inverteix en passar a $\cat{C}\op$ i què es manté.
\end{exercici}

\begin{solucio}
\emph{Primera part: el cas dels monomorfismes.} Siguin $f\colon A\to B$ i $g\colon B\to C$ monomorfismes, i siguin $u,v\colon X\to A$ amb $(g\comp f)\comp u=(g\comp f)\comp v$:
\[
\begin{tikzpicture}[baseline=(m.center)]
  \matrix (m) [matrix of math nodes, column sep=3em] { X & A & B & C \\ };
  \draw[-{Stealth[scale=1.0]}] ([yshift=2.6pt]m-1-1.east) -- node[above] {$u$} ([yshift=2.6pt]m-1-2.west);
  \draw[-{Stealth[scale=1.0]}] ([yshift=-2.6pt]m-1-1.east) -- node[below] {$v$} ([yshift=-2.6pt]m-1-2.west);
  \draw[-{Stealth[scale=1.1]}] (m-1-2) -- node[above] {$f$} (m-1-3);
  \draw[-{Stealth[scale=1.1]}] (m-1-3) -- node[above] {$g$} (m-1-4);
\end{tikzpicture}
\qquad
(g\comp f)\comp u=(g\comp f)\comp v\ \Longrightarrow\ u=v.
\]
Per associativitat, la hipòtesi és $g\comp(f\comp u)=g\comp(f\comp v)$. Com que $g$ és mono, en cancel·lem el factor de l'esquerra i obtenim $f\comp u=f\comp v$. Com que $f$ és mono, tornem a cancel·lar i obtenim $u=v$. Per tant $g\comp f$ és mono. Observem l'ordre: primer es cancel·la $g$, el factor \emph{exterior}, i després $f$.

\emph{Segona part: dualització.} Procedim en tres passos.

\textbf{1. L'enunciat és categòric.} \enquote{Si $f\colon A\to B$ i $g\colon B\to C$ són monomorfismes, aleshores $g\comp f$ és un monomorfisme} només parla d'objectes, morfismes, composició i igualtats entre composicions (la definició de mono és d'aquesta mena). Per tant té dual, i com que l'hem demostrat en una categoria qualsevol, el principi de dualitat (teorema~\ref{teo:dualitat}) garanteix que el dual també és vàlid en tota categoria.

\textbf{2. Què s'inverteix.} Recorrem l'enunciat element per element:
\begin{center}
\begin{tabular}{@{}ll@{}}
\toprule
A $\cat{C}$ & A l'enunciat dual \\
\midrule
$f\colon A\to B$, \quad $g\colon B\to C$ & $f\colon B\to A$, \quad $g\colon C\to B$ \\
la composta $g\comp f\colon A\to C$ & la composta $f\comp g\colon C\to A$ \\
fletxes de prova $u,v\colon X\to A$ & fletxes de prova $u,v\colon A\to X$ \\
mono: $m\comp u=m\comp v\Rightarrow u=v$ & epi: $u\comp m=v\comp m\Rightarrow u=v$ \\
\bottomrule
\end{tabular}
\end{center}
Es mantenen els quantificadors (\enquote{per a tot $u,v$}), la implicació, la igualtat i l'associativitat, que és autodual. L'enunciat dual diu, doncs: si $f\colon B\to A$ i $g\colon C\to B$ són epimorfismes, aleshores $f\comp g\colon C\to A$ és un epimorfisme. Rebatejant les fletxes perquè la cadena es llegeixi d'esquerra a dreta, $p\colon A\to B$ i $q\colon B\to C$, això és: \emph{si $p$ i $q$ són epimorfismes, $q\comp p$ és un epimorfisme.}

\textbf{3. Per què és vàlid.} És l'argument del principi, fet explícit. Siguin $p\colon A\to B$ i $q\colon B\to C$ epimorfismes de $\cat{C}$. A $\cat{C}\op$, les fletxes $p\op\colon B\to A$ i $q\op\colon C\to B$ són monomorfismes, perquè ser epi a $\cat{C}$ és ser mono a $\cat{C}\op$. Per la primera part, aplicada a $\cat{C}\op$, la composta $p\op\comp q\op\colon C\to A$ és mono a $\cat{C}\op$. Però $p\op\comp q\op=(q\comp p)\op$, de manera que $(q\comp p)\op$ és mono a $\cat{C}\op$, és a dir, $q\comp p$ és epi a $\cat{C}$.

\emph{Comprovació: la prova dual.} No l'hem necessitat, però val la pena escriure-la per veure que és la primera llegida a l'inrevés. Siguin $u,v\colon C\to X$ amb $u\comp(q\comp p)=v\comp(q\comp p)$:
\[
\begin{tikzpicture}[baseline=(m.center)]
  \matrix (m) [matrix of math nodes, column sep=3em] { A & B & C & X \\ };
  \draw[-{Stealth[scale=1.1]}] (m-1-1) -- node[above] {$p$} (m-1-2);
  \draw[-{Stealth[scale=1.1]}] (m-1-2) -- node[above] {$q$} (m-1-3);
  \draw[-{Stealth[scale=1.0]}] ([yshift=2.6pt]m-1-3.east) -- node[above] {$u$} ([yshift=2.6pt]m-1-4.west);
  \draw[-{Stealth[scale=1.0]}] ([yshift=-2.6pt]m-1-3.east) -- node[below] {$v$} ([yshift=-2.6pt]m-1-4.west);
\end{tikzpicture}
\qquad
u\comp(q\comp p)=v\comp(q\comp p)\ \Longrightarrow\ u=v.
\]
Per associativitat, $(u\comp q)\comp p=(v\comp q)\comp p$; com que $p$ és epi, $u\comp q=v\comp q$; com que $q$ és epi, $u=v$. Les mateixes passes, amb cada composició en l'ordre invers: ara es cancel·la primer $p$, el factor de la \emph{dreta}, que és la fletxa per on comença el camí.

\emph{Error típic.} En dualitzar és fàcil invertir les fletxes però oblidar d'invertir l'ordre de la composició, i escriure \enquote{$p\comp q$ és epi}. Si $p\colon A\to B$ i $q\colon B\to C$, aquesta composta ni tan sols està definida, tret que $C=A$. Comprovar els dominis i els codominis de cada composta és la millor manera de detectar l'error.
\end{solucio}

\section{Grafs i categories lliures}

\begin{exercici}
Considera el graf amb un sol vèrtex i dues arestes que són bucles, $a$ i $b$. Descriu la categoria lliure que genera. Quants morfismes té?
\end{exercici}

\begin{solucio}
La categoria lliure té un sol objecte. Els morfismes són tots els camins finits:
\[
\varepsilon,
\quad a,
\quad b,
\quad aa,
\quad ab,
\quad ba,
\quad bb,
\quad aaa,
\quad\ldots
\]
on $\varepsilon$ és el camí buit i actua com a identitat. La composició és la concatenació de paraules en ordre de recorregut: $b\comp a=ab$. Com que cap relació no identifica camins diferents, hi ha una infinitat numerable de morfismes. Llevat de l'isomorfisme canònic que inverteix l'ordre de les lletres (exemple~\ref{ex:monoide-lliure-cat}), és el monoide lliure sobre l'alfabet $\{a,b\}$.
\end{solucio}

\section{Qüestions de mida}

\begin{exercici}
En una categoria petita, l'oposada $\cat{C}\op$ també és petita. Comprova-ho, i comprova a més que $\cat{C}$ i $\cat{C}\op$ tenen exactament el mateix nombre de morfismes.
\end{exercici}

\begin{solucio}
Els objectes de $\cat{C}\op$ són els mateixos que els de $\cat{C}$. Cada morfisme $f\colon A\to B$ de $\cat{C}$ correspon al mateix morfisme vist com $f\colon B\to A$ a $\cat{C}\op$. Això dona una bijecció entre les col·leccions de morfismes de totes dues categories.

Per tant, si els objectes i els morfismes de $\cat{C}$ formen conjunts, també ho fan els de $\cat{C}\op$. En particular, $\cat{C}$ és petita si i només si ho és $\cat{C}\op$.
\end{solucio}

\begin{exercici}\label{exr:pr-set-localment-petita}
Comprova que $\cat{Set}$ és localment petita: per a conjunts $A$ i $B$, mostra que la col·lecció de les aplicacions $A\to B$ és un conjunt, identificant cada aplicació amb el seu graf, que és un subconjunt d'$A\times B$.
\end{exercici}

\begin{solucio}
Una aplicació $f\colon A\to B$ queda determinada pel seu \textbf{graf}
\[
\Gamma_f=\{(a,f(a))\mid a\in A\}\subseteq A\times B.
\]
Un subconjunt $\Gamma\subseteq A\times B$ és el graf d'alguna aplicació $A\to B$ exactament quan és \emph{funcional}: per a cada $a\in A$ hi ha un únic $b\in B$ amb $(a,b)\in\Gamma$. Així, $f\mapsto\Gamma_f$ és una bijecció entre les aplicacions $A\to B$ i els subconjunts funcionals d'$A\times B$; la inversa envia $\Gamma$ a l'aplicació que assigna a cada $a$ l'únic $b$ amb $(a,b)\in\Gamma$.

Vegem ara que aquests subconjunts funcionals formen un conjunt. Com que $A$ i $B$ són conjunts, el producte $A\times B$ és un conjunt, i també ho és el seu conjunt de parts $\mathcal P(A\times B)$. La condició \enquote{ser funcional} és una propietat ben definida dels elements de $\mathcal P(A\times B)$, de manera que, per l'axioma de separació,
\[
\{\Gamma\in\mathcal P(A\times B)\mid \forall a\in A\ \exists!\,b\in B\ (a,b)\in\Gamma\}
\]
és un conjunt. Per la bijecció anterior, $\Hom_{\cat{Set}}(A,B)$ és un conjunt.

Una precisió, la mateixa de l'observació~\ref{obs:extrems-implicits} de la Teoria per a $\cat{Rel}$: el graf no determina el codomini, perquè $\Gamma_f\subseteq A\times B$ també és un subconjunt d'$A\times B'$ per a tot $B'\supseteq B$. Per això, en rigor, un morfisme de $\cat{Set}$ és la terna $(A,B,\Gamma_f)$. Això no canvia la conclusió: les ternes $(A,B,\Gamma)$ amb $\Gamma$ funcional formen el conjunt $\{A\}\times\{B\}\times\{\Gamma\in\mathcal P(A\times B)\mid \Gamma\text{ funcional}\}$, que està en bijecció amb l'anterior.

\emph{Abast.} El mateix argument explica per què $\Hom(A,B)$ és un conjunt en els exemples habituals. A $\cat{Rel}$, directament, $\Hom_{\cat{Rel}}(X,Y)$ és, llevat dels extrems, $\mathcal P(X\times Y)$. A les categories de conjunts estructurats ($\cat{Grp}$, $\cat{Top}$, $\cat{Vect}_{\mathbb K}$, $\cat{Mon}$, $\cat{Pos}$), els morfismes són aplicacions entre els conjunts subjacents que compleixen una condició addicional; formen, doncs, un subconjunt de l'homconjunt de $\cat{Set}$ corresponent, de nou per separació. En canvi, la col·lecció de \emph{tots} els objectes de $\cat{Set}$ no és un conjunt (paradoxa de Russell): $\cat{Set}$ és localment petita sense ser petita.
\end{solucio}

\chapter{Functors}\label{cap:practica-functors}

En aquest capítol resolem amb detall una selecció d'exercicis sobre \emph{functors}. Com al capítol anterior, l'objectiu no és només arribar al resultat, sinó mostrar com es comprova que una assignació és un functor, com es reconeixen els functors coneguts sota formes noves i com es distingeixen les propietats que un functor preserva de les que no. Convé intentar cada exercici abans de llegir-ne la solució.

\section{Definició de functor}

\begin{exercici}\label{prac:functor-quadrat}
Sigui $F\colon\cat{C}\to\cat{D}$ un functor i suposem que a $\cat{C}$ tenim un quadrat commutatiu, és a dir, quatre morfismes $f\colon A\to B$, $u\colon A\to C$, $v\colon B\to D$ i $g\colon C\to D$ amb
\[
v\comp f=g\comp u.
\]
Mostra que la seva imatge per $F$ torna a ser un quadrat commutatiu:
\[
F(v)\comp F(f)=F(g)\comp F(u).
\]
\end{exercici}

\begin{solucio}
El quadrat de partida i la seva imatge són:
\[
\begin{tikzpicture}[baseline=(m.center)]
  \matrix (m) [matrix of math nodes, row sep=2.6em, column sep=3em] {
    A & B \\
    C & D \\
  };
  \draw[-{Stealth[scale=1.1]}] (m-1-1) -- node[above] {$f$} (m-1-2);
  \draw[-{Stealth[scale=1.1]}] (m-1-1) -- node[left] {$u$} (m-2-1);
  \draw[-{Stealth[scale=1.1]}] (m-1-2) -- node[right] {$v$} (m-2-2);
  \draw[-{Stealth[scale=1.1]}] (m-2-1) -- node[below] {$g$} (m-2-2);
\end{tikzpicture}
\qquad\stackrel{F}{\longmapsto}\qquad
\begin{tikzpicture}[baseline=(m.center)]
  \matrix (m) [matrix of math nodes, row sep=2.6em, column sep=3.4em] {
    F(A) & F(B) \\
    F(C) & F(D) \\
  };
  \draw[-{Stealth[scale=1.1]}] (m-1-1) -- node[above] {$F(f)$} (m-1-2);
  \draw[-{Stealth[scale=1.1]}] (m-1-1) -- node[left] {$F(u)$} (m-2-1);
  \draw[-{Stealth[scale=1.1]}] (m-1-2) -- node[right] {$F(v)$} (m-2-2);
  \draw[-{Stealth[scale=1.1]}] (m-2-1) -- node[below] {$F(g)$} (m-2-2);
\end{tikzpicture}
\]
Partim de la hipòtesi $v\comp f=g\comp u$. Apliquem $F$ a aquesta igualtat; com que $F$ és una aplicació ben definida sobre els morfismes, morfismes iguals tenen imatges iguals:
\[
F(v\comp f)=F(g\comp u).
\]
Ara fem servir que $F$ preserva la composició a cada costat:
\[
F(v)\comp F(f)=F(v\comp f)=F(g\comp u)=F(g)\comp F(u).
\]
Per tant el quadrat imatge commuta. Aquest és el contingut de la idea que \emph{els functors envien diagrames commutatius a diagrames commutatius}: n'hi ha prou que ho comprovem sobre les composicions, perquè un diagrama commuta exactament quan certes composicions coincideixen.
\end{solucio}

\begin{exercici}\label{pr-parts-cov}
Comprova que l'assignació que envia cada conjunt $X$ al conjunt de parts $\mathcal P(X)$ i cada aplicació $f\colon X\to Y$ a la \textbf{imatge directa}
\[
\mathcal P_!(f)\colon\mathcal P(X)\to\mathcal P(Y),
\qquad
A\mapsto f(A)=\{f(a)\mid a\in A\},
\]
és un functor covariant
\[
\mathcal P_!\colon\cat{Set}\to\cat{Set},
\]
el \textbf{functor de parts per imatge directa}. (Reservem $\mathcal P$, sense subíndex, per al functor de parts contravariant de la Teoria, que actua per antiimatge.)
\end{exercici}

\begin{solucio}
Hem de comprovar la preservació d'identitats i de composició.

Per a la \textbf{identitat}, si $f=\id_X$ aleshores per a tot $A\subseteq X$,
\[
\mathcal P_!(\id_X)(A)=\{\id_X(a)\mid a\in A\}=A,
\]
de manera que $\mathcal P_!(\id_X)=\id_{\mathcal P(X)}$.

Per a la \textbf{composició}, siguin $f\colon X\to Y$ i $g\colon Y\to Z$. Per a tot $A\subseteq X$,
\[
\begin{aligned}
\mathcal P_!(g\comp f)(A)&=(g\comp f)(A)=g\big(f(A)\big)\\
&=\mathcal P_!(g)\big(\mathcal P_!(f)(A)\big)=\big(\mathcal P_!(g)\comp\mathcal P_!(f)\big)(A).
\end{aligned}
\]
La igualtat central, $(g\comp f)(A)=g(f(A))$, és la propietat coneguda de la imatge directa. Per tant $\mathcal P_!$ és un functor covariant. Té els mateixos objectes que $\mathcal P$, però actua sobre els morfismes en sentit contrari i per una operació diferent: són dos functors distints.
\end{solucio}

\begin{exercici}\label{pr-hom-cov}
Sigui $\cat{C}$ una categoria localment petita i $A$ un objecte fix. Comprova que l'assignació que envia cada objecte $X$ al conjunt $\Hom(A,X)$ i cada morfisme $f\colon X\to Y$ a l'aplicació
\[
f_*\colon\Hom(A,X)\to\Hom(A,Y),
\qquad
\varphi\mapsto f\comp\varphi,
\]
és un functor covariant $\Hom(A,-)\colon\cat{C}\to\cat{Set}$ (suposant $\cat{C}$ localment petita).
\end{exercici}

\begin{solucio}
Primer cal veure que $f_*$ està ben definida: si $\varphi\colon A\to X$, aleshores $f\comp\varphi\colon A\to Y$, de manera que $f_*(\varphi)\in\Hom(A,Y)$.

Per a la \textbf{identitat}, prenent $f=\id_X$,
\[
(\id_X)_*(\varphi)=\id_X\comp\varphi=\varphi,
\]
de manera que $(\id_X)_*=\id_{\Hom(A,X)}$.

Per a la \textbf{composició}, siguin $f\colon X\to Y$ i $g\colon Y\to Z$. Per a tot $\varphi\in\Hom(A,X)$, usant l'associativitat de $\cat{C}$,
\[
(g\comp f)_*(\varphi)=(g\comp f)\comp\varphi=g\comp(f\comp\varphi)=g_*\big(f_*(\varphi)\big).
\]
Per tant $(g\comp f)_*=g_*\comp f_*$ i $\Hom(A,-)$ és un functor covariant.
\end{solucio}

\section{Exemples de functors}

\begin{exercici}
Siguin $T$ i $T'$ dos sistemes deductius sobre el mateix conjunt de fórmules, amb $T'$ una \textbf{extensió} de $T$ ---és a dir, tot el que és deduïble a $T$ també ho és a $T'$: si $A\vdash_T B$, aleshores $A\vdash_{T'} B$. Comprova que la identitat sobre les fórmules indueix un functor
\[
J\colon\cat{Prov}_T\to\cat{Prov}_{T'},
\]
i explica en quins termes es podria formular que $J$ sigui ple.
\end{exercici}

\begin{solucio}
Recordem que $\cat{Prov}_T$ té les fórmules per objectes i una única fletxa $A\to B$ exactament quan $A\vdash_T B$; és una categoria prima. Definim $J$ com la identitat sobre els objectes, $J(A)=A$.

Sobre els morfismes, una fletxa $A\to B$ de $\cat{Prov}_T$ existeix exactament quan $A\vdash_T B$; la fletxa registra, doncs, el fet de la deduïbilitat, no una prova concreta. Com que $T'$ és una extensió de $T$, tenim $A\vdash_{T'} B$, de manera que hi ha una fletxa $J(A)\to J(B)$ a $\cat{Prov}_{T'}$, i és única perquè $\cat{Prov}_{T'}$ és prima. Això defineix $J$ sobre morfismes.

Les condicions de functor se satisfan automàticament perquè totes dues categories són primes: entre dos objectes hi ha com a màxim una fletxa, així que la preservació d'identitats i de composicions no imposa res més enllà de l'existència de les fletxes imatge, que ja hem comprovat.

Quant a \textbf{ser ple}: $J$ seria ple si, per a cada parella de fórmules $A,B$, tota fletxa $J(A)\to J(B)$ a $\cat{Prov}_{T'}$ provingués d'una fletxa $A\to B$ a $\cat{Prov}_T$. És a dir, si
\[
A\vdash_{T'} B\quad\Longrightarrow\quad A\vdash_T B.
\]
Dit d'una altra manera, $J$ és ple exactament quan l'extensió $T'$ \emph{no afegeix cap conseqüència nova} entre les fórmules considerades. Quan $T'$ demostra estrictament més relacions que $T$, el functor $J$ és fidel ---de fet, totes dues categories són primes, així que $J$ és sempre fidel--- però no ple.
\end{solucio}

\begin{exercici}
Comprova amb detall que un functor entre dos preordres $(P,\leq)$ i $(Q,\leq)$, vistos com a categories, és exactament una aplicació monòtona.
\end{exercici}

\begin{solucio}
Un functor $F\colon P\to Q$ dona, en primer lloc, una aplicació $x\mapsto F(x)$ dels objectes, és a dir, dels elements de $P$ als de $Q$.

La condició sobre morfismes diu que si hi ha una fletxa $x\to y$ a $P$, aleshores hi ha una fletxa $F(x)\to F(y)$ a $Q$. Com que en un preordre hi ha una fletxa $x\to y$ exactament quan $x\leq y$, això equival a
\[
x\leq y
\quad\Longrightarrow\quad
F(x)\leq F(y),
\]
que és precisament la definició d'aplicació monòtona.

Recíprocament, tota aplicació monòtona defineix un functor: envia l'única fletxa $x\to y$ (quan $x\leq y$) a l'única fletxa $F(x)\to F(y)$. Les condicions de preservació d'identitats i composició es compleixen automàticament, perquè entre dos objectes hi ha com a màxim una fletxa i totes les igualtats entre morfismes són trivials. Així, els functors $P\to Q$ i les aplicacions monòtones $P\to Q$ són el mateix.
\end{solucio}

\begin{exercici}
Sigui $M$ un monoide vist com a categoria $\cat{B}M$ amb un sol objecte $\ast$. Comprova que un functor $F\colon\cat{B}M\to\cat{Set}$ és el mateix que una acció de $M$ sobre un conjunt, és a dir, un conjunt $S$ juntament amb una operació
\[
M\times S\to S,
\qquad
(m,s)\mapsto m\cdot s,
\]
tal que $1_M\cdot s=s$ i $(mn)\cdot s=m\cdot(n\cdot s)$.
\end{exercici}
\begin{solucio}
Un functor $F\colon\cat{B}M\to\cat{Set}$ tria un conjunt $S=F(\ast)$ i envia cada element $m\in M$ ---cada morfisme $\ast\to\ast$--- a una aplicació $F(m)\colon S\to S$. Definim l'acció per
\[
m\cdot s=F(m)(s).
\]

La preservació de la identitat diu $F(1_M)=\id_S$, és a dir, $1_M\cdot s=s$ per a tot $s$. La preservació de la composició diu $F(mn)=F(m)\comp F(n)$, és a dir,
\[
(mn)\cdot s=F(mn)(s)=F(m)\big(F(n)(s)\big)=m\cdot(n\cdot s).
\]
Aquestes són exactament les dues lleis d'una acció per l'esquerra. Recíprocament, tota acció defineix, per la mateixa fórmula, un functor. Per tant els functors $\cat{B}M\to\cat{Set}$ i les accions de $M$ coincideixen.
\end{solucio}

\begin{exercici}
Comprova que la construcció de la categoria lliure és functorial sobre les fletxes: donat un morfisme de grafs $\alpha\colon G\to H$, l'assignació que envia cada camí d'arestes de $G$ a la seqüència de les seves imatges és un functor $\Free(\alpha)\colon\Free(G)\to\Free(H)$.
\end{exercici}

\begin{solucio}
Un camí a $\Free(G)$ és una seqüència finita d'arestes encadenades $a_1a_2\cdots a_n$. Definim
\[
\Free(\alpha)(a_1a_2\cdots a_n)=\alpha_E(a_1)\,\alpha_E(a_2)\cdots\alpha_E(a_n),
\]
i sobre els objectes $\Free(\alpha)=\alpha_V$. Cal comprovar que aquesta seqüència és realment un camí a $H$: com que $\alpha$ preserva origen i destí, si $a_i$ acaba on comença $a_{i+1}$, aleshores $\alpha_E(a_i)$ acaba on comença $\alpha_E(a_{i+1})$, de manera que les imatges també s'encadenen.

La \textbf{identitat} de cada vèrtex és el camí buit, que s'envia al camí buit del vèrtex imatge: $\Free(\alpha)$ preserva identitats. La \textbf{composició} és la concatenació. Siguin $p=(a_1,\dots,a_r)$ i $q=(b_1,\dots,b_s)$ dos camins componibles, amb el final de $p$ igual a l'inici de $q$; la seva composició és el camí que recorre primer $p$ i després $q$,
\[
q\comp p=(a_1,\dots,a_r,b_1,\dots,b_s).
\]
Aplicar $\alpha_E$ terme a terme i després concatenar dona el mateix que concatenar i després aplicar $\alpha_E$ terme a terme, de manera que
\[
\Free(\alpha)(q\comp p)=\Free(\alpha)(q)\comp\Free(\alpha)(p).
\]
Per tant $\Free(\alpha)$ és un functor.
\end{solucio}

\section{Composició de functors i la categoria \texorpdfstring{$\cat{Cat}$}{Cat}}

\begin{exercici}
Comprova que la composició de dos functors $F\colon\cat{C}\to\cat{D}$ i $G\colon\cat{D}\to\cat{E}$, seguida d'un tercer $H\colon\cat{E}\to\cat{F}$, és associativa:
\[
H\comp(G\comp F)=(H\comp G)\comp F.
\]
\end{exercici}

\begin{solucio}
Dues igualtats de functors es comproven sobre objectes i sobre morfismes. Sobre un objecte $A$ de $\cat{C}$,
\[
\big(H\comp(G\comp F)\big)(A)=H\big(G(F(A))\big)=\big((H\comp G)\comp F\big)(A),
\]
i sobre un morfisme $f$,
\[
\big(H\comp(G\comp F)\big)(f)=H\big(G(F(f))\big)=\big((H\comp G)\comp F\big)(f).
\]
En tots dos casos la igualtat es redueix a l'associativitat de l'aplicació successiva, primer $F$, després $G$, després $H$. Per tant la composició de functors és associativa, i amb els functors identitat com a neutres, això confirma que $\cat{Cat}$ és una categoria.
\end{solucio}

\begin{exercici}
Fixat un objecte $A$ d'una categoria $\cat{C}$, descriu el functor oblit $U_A\colon\cat{C}/A\to\cat{C}$ de la categoria de $\cat{C}$ sobre $A$ cap a $\cat{C}$, i comprova que és un functor.
\end{exercici}

\begin{solucio}
Recordem que un objecte de $\cat{C}/A$ és un morfisme $x\colon X\to A$ de $\cat{C}$, i un morfisme de $x\colon X\to A$ cap a $y\colon Y\to A$ és un morfisme $h\colon X\to Y$ de $\cat{C}$ tal que $y\comp h=x$ (el triangle commuta). Definim $U_A$ així: sobre objectes, envia $x\colon X\to A$ al seu domini $X$; sobre morfismes, envia el morfisme $h$ del triangle al mateix $h\colon X\to Y$ vist a $\cat{C}$, oblidant la condició $y\comp h=x$.

Comprovem que és un functor. \textbf{Identitats:} la identitat de l'objecte $x\colon X\to A$ a $\cat{C}/A$ és $\id_X$ (que compleix $x\comp\id_X=x$), i $U_A(\id_X)=\id_X=\id_{U_A(x)}$. \textbf{Composició:} si $h\colon x\to y$ i $k\colon y\to z$ són morfismes componibles a $\cat{C}/A$, la seva composició és $k\comp h$ com a morfisme de $\cat{C}$; per tant $U_A(k\comp h)=k\comp h=U_A(k)\comp U_A(h)$. Així $U_A$ preserva identitats i composició, i és un functor. Intuïtivament, $U_A$ \enquote{oblida} el morfisme cap a $A$ i reté només el domini i les fletxes entre dominis.
\end{solucio}

\section{Functors i isomorfismes}

\begin{exercici}
Dona un exemple de functor i d'un morfisme \emph{no} isomorf la imatge del qual \emph{sí} que és un isomorfisme. Això mostra que els functors preserven isomorfismes, però no els reflecteixen en general.
\end{exercici}

\begin{solucio}
Considerem el functor oblit $U\colon\cat{Pos}\to\cat{Set}$, que envia cada conjunt parcialment ordenat al seu conjunt subjacent i cada aplicació monòtona a ella mateixa. Siguin $P$ el conjunt $\{0,1\}$ amb l'ordre discret (cap parella d'elements diferents no és comparable) i $Q$ el mateix conjunt amb l'ordre $0<1$, i sigui
\[
f\colon P\to Q
\]
la identitat de conjunts. És monòtona, perquè a $P$ les úniques relacions són $0\le 0$ i $1\le 1$, que es conserven a $Q$. La inversa $Q\to P$, en canvi, no ho és: a $Q$ tenim $0\le 1$, però a $P$ no. Per tant $f$ \textbf{no és un isomorfisme} a $\cat{Pos}$, perquè no té cap invers dins de la categoria.

Ara bé, la seva imatge $U(f)$ és l'aplicació identitat del conjunt $\{0,1\}$, que \textbf{sí que és un isomorfisme} a $\cat{Set}$.

\emph{Un exemple anàleg, per a qui conegui la topologia.} El functor oblit $U\colon\cat{Top}\to\cat{Set}$ es comporta igual. Si $X$ és un conjunt amb dues topologies $\tau$ i $\tau'$, amb $\tau$ estrictament més fina que $\tau'$, la identitat $(X,\tau)\to(X,\tau')$ és contínua, però la seva inversa no ho és; no és, doncs, un isomorfisme a $\cat{Top}$, mentre que la seva imatge a $\cat{Set}$ és la identitat de $X$.

En tots dos casos, $U$ no reflecteix isomorfismes: la seva imatge pot ser invertible sense que l'original ho sigui. Els functors plens i fidels sí que reflecteixen isomorfismes, encara que aquesta condició no és necessària: hi ha functors que els reflecteixen sense ser plens. Per exemple, el functor oblit $\cat{Grp}\to\cat{Set}$ reflecteix isomorfismes ---un homomorfisme bijectiu de grups és un isomorfisme---, malgrat no ser ple.
\end{solucio}

\begin{exercici}\label{pr-2-set-mono-epi}
A la categoria $\cat{2}=(0\xrightarrow{u}1)$, la fletxa $u$ és mono i epi. Construeix un functor $\cat{2}\to\cat{Set}$ que no l'enviï a un monomorfisme, i un altre que no l'enviï a un epimorfisme.
\end{exercici}

\begin{solucio}
Un functor $\cat{2}\to\cat{Set}$ queda determinat per dos conjunts i una aplicació entre ells, imatges de $0$, $1$ i $u$: les úniques composicions de $\cat{2}$ involucren identitats, de manera que la functorialitat només demana enviar identitats a identitats.
\begin{itemize}
  \item $F(0)=\{0,1\}$, $F(1)=\{\ast\}$ i $F(u)$ l'única aplicació $\{0,1\}\to\{\ast\}$. No és injectiva, i per tant $F(u)$ no és mono.
  \item $G(0)=\{\ast\}$, $G(1)=\{0,1\}$ i $G(u)$ l'aplicació de valor $0$. No és exhaustiva, i per tant $G(u)$ no és epi.
\end{itemize}
Un functor, doncs, no preserva en general ni els monomorfismes ni els epimorfismes. A $\cat{2}$, que $u$ sigui mono i epi és conseqüència que la categoria és prima; a $\cat{Set}$ hi ha més fletxes amb què comparar, i aquesta trivialitat es perd.
\end{solucio}

\section{Functors covariants i contravariants}

\begin{exercici}\label{pr-parts-contra}
Comprova que el functor de parts contravariant $\mathcal P\colon\cat{Set}\op\to\cat{Set}$, que envia $f\colon X\to Y$ a l'antiimatge
\[
f^{-1}\colon\mathcal P(Y)\to\mathcal P(X),
\qquad
B\mapsto f^{-1}(B),
\]
és realment un prefeix sobre $\cat{Set}$ (definició~\ref{def:prefeix}).
\end{exercici}

\begin{solucio}
Cal comprovar la preservació d'identitats i la inversió de l'ordre en la composició.

Per a la \textbf{identitat}, $(\id_X)^{-1}(B)=B$ per a tot $B\subseteq X$, de manera que $(\id_X)^{-1}=\id_{\mathcal P(X)}$.

Per a la \textbf{composició}, siguin $f\colon X\to Y$ i $g\colon Y\to Z$. Per a tot $C\subseteq Z$,
\[
(g\comp f)^{-1}(C)=\{x\mid g(f(x))\in C\}=\{x\mid f(x)\in g^{-1}(C)\}=f^{-1}\big(g^{-1}(C)\big).
\]
Per tant
\[
(g\comp f)^{-1}=f^{-1}\comp g^{-1},
\]
amb l'ordre invertit. Aquesta és exactament la condició de prefeix, és a dir, de functor $\cat{Set}\op\to\cat{Set}$.
\end{solucio}

\begin{exercici}\label{pr-hom-contra}
Sigui $B$ un objecte fix d'una categoria localment petita $\cat{C}$. Comprova que $\Hom(-,B)$ és un prefeix sobre $\cat{C}$, és a dir, un functor $\Hom(-,B)\colon\cat{C}\op\to\cat{Set}$.
\end{exercici}

\begin{solucio}
A cada objecte $X$ li assignem el conjunt $\Hom(X,B)$. A cada morfisme $f\colon X\to Y$ li assignem l'aplicació de \textbf{precomposició}
\[
f^*\colon\Hom(Y,B)\to\Hom(X,B),
\qquad
\psi\mapsto\psi\comp f,
\]
que està ben definida perquè si $\psi\colon Y\to B$ aleshores $\psi\comp f\colon X\to B$. Observem que $f^*$ inverteix el sentit: parteix dels morfismes des de $Y$ i arriba als morfismes des de $X$.

Per a la \textbf{identitat}, $(\id_X)^*(\psi)=\psi\comp\id_X=\psi$. Per a la \textbf{composició}, siguin $f\colon X\to Y$ i $g\colon Y\to Z$; per a tot $\psi\in\Hom(Z,B)$,
\[
(g\comp f)^*(\psi)=\psi\comp(g\comp f)=(\psi\comp g)\comp f=f^*\big(g^*(\psi)\big).
\]
Per tant $(g\comp f)^*=f^*\comp g^*$, amb l'ordre invertit, i $\Hom(-,B)$ és contravariant.
\end{solucio}

\begin{exercici}\label{pr-bihom}
Sigui $\cat{C}$ localment petita. Comprova que el bifunctor
\[
\Hom(-,-)\colon\cat{C}\op\times\cat{C}\to\cat{Set},\qquad \Hom(f,g)(h)=g\comp h\comp f,
\]
respecta la composició simultània en les dues variables i les identitats. És a dir, comprova que és realment un functor, no només en cada variable per separat.
\end{exercici}

\begin{solucio}
Siguin $f\colon A'\to A$, $f'\colon A''\to A'$, $g\colon B\to B'$ i $g'\colon B'\to B''$ morfismes de $\cat{C}$. Recordem que la composició a $\cat{C}\op\times\cat{C}$ actua component a component, i que a la primera component (la de $\cat{C}\op$) l'ordre és l'invers del de $\cat{C}$. \textbf{Composició.} Per a tot $h\colon A\to B$,
\[
\begin{aligned}
\Hom(f',g')\big(\Hom(f,g)(h)\big)
&=g'\comp(g\comp h\comp f)\comp f'\\
&=(g'\comp g)\comp h\comp(f\comp f')\\
&=\Hom(f\comp f',\,g'\comp g)(h),
\end{aligned}
\]
on el pas central és només associativitat. El resultat $\Hom(f\comp f',g'\comp g)$ és precisament $\Hom$ aplicat a la composició de la parella $(f',g')$ amb $(f,g)$ a $\cat{C}\op\times\cat{C}$: a la segona component obtenim $g'\comp g$, i a la primera $f\comp f'$ (l'ordre invertit, coherent amb la contravariància). \textbf{Identitats.} Per a la parella identitat $(\id_A,\id_B)$,
\[
\Hom(\id_A,\id_B)(h)=\id_B\comp h\comp\id_A=h,
\]
és a dir $\Hom(\id_A,\id_B)=\id_{\Hom(A,B)}$. Per tant $\Hom(-,-)$ preserva composició i identitats de $\cat{C}\op\times\cat{C}$, i és un functor de ple dret.
\end{solucio}

\section{Functors fidels i plens}

\begin{exercici}
Comprova que el functor oblit $U\colon\cat{Pos}\to\cat{Set}$, que envia cada ordre parcial al seu conjunt subjacent i cada aplicació monòtona a ella mateixa, és fidel però no ple.
\end{exercici}

\begin{solucio}
\textbf{Fidel.} Siguin $f,g\colon P\to Q$ dues aplicacions monòtones amb $U(f)=U(g)$. Això vol dir que $f$ i $g$ coincideixen com a aplicacions entre els conjunts subjacents, és a dir, en tot element de $P$. Una aplicació monòtona no és res més que la seva acció sobre els elements (amb la condició afegida de preservar l'ordre), de manera que $f=g$. Així cada aplicació $U_{P,Q}\colon\Hom_{\cat{Pos}}(P,Q)\to\Hom_{\cat{Set}}(U(P),U(Q))$ és injectiva i, per tant, $U$ és fidel. (Sobre objectes, en canvi, $U$ no és injectiu: ordres diferents sobre un mateix conjunt tenen la mateixa imatge. Ser fidel no hi té res a veure.)

\textbf{No ple.} Prenem l'ordre parcial $P=\{0,1\}$ amb $0<1$ i el mateix $Q=P$. L'aplicació $\sigma\colon\{0,1\}\to\{0,1\}$ que intercanvia els dos elements ($\sigma(0)=1$, $\sigma(1)=0$) és una aplicació de conjunts perfectament vàlida, i pertany a $\Hom_{\cat{Set}}(U P,U Q)$. Però no és monòtona: de $0<1$ es passaria a $\sigma(0)=1$ i $\sigma(1)=0$, és a dir $1\le 0$, que és fals. Per tant $\sigma$ no és de la forma $U(f)$ per a cap aplicació monòtona $f$, l'aplicació $U_{P,Q}$ no és exhaustiva i $U$ no és ple.

\emph{Per a qui conegui els grups.} El functor oblit $U\colon\cat{Grp}\to\cat{Set}$ es comporta igual: és fidel perquè un homomorfisme queda determinat per la seva acció sobre els elements, i no és ple perquè hi ha aplicacions entre els conjunts subjacents que no són homomorfismes, com $x\mapsto x+1$ de $\mathbb Z$ en $\mathbb Z$, que no envia $0$ a $0$.
\end{solucio}

\begin{exercici}\label{pr-esquelet-finset}
Sigui $\cat{S}$ la subcategoria plena de $\cat{FinSet}$ que té per objectes els conjunts $\underline{0}=\varnothing$ i $\underline{n}=\{1,\dots,n\}$ per a $n\geq 1$. Comprova que la inclusió
\[
\iota\colon\cat{S}\hookrightarrow\cat{FinSet}
\]
és plena, fidel i essencialment exhaustiva, però no exhaustiva sobre objectes en sentit estricte.
\end{exercici}

\begin{solucio}
\textbf{Fidel.} La inclusió envia cada morfisme a ell mateix: per a objectes $\underline{m},\underline{n}$ de $\cat{S}$, l'aplicació $\iota_{\underline{m},\underline{n}}\colon\Hom_{\cat{S}}(\underline{m},\underline{n})\to\Hom_{\cat{FinSet}}(\underline{m},\underline{n})$ és $h\mapsto h$, que és injectiva.

\textbf{Plena.} Com que $\cat{S}$ és una subcategoria \emph{plena}, conté tots els morfismes de $\cat{FinSet}$ entre dos objectes seus: $\Hom_{\cat{S}}(\underline{m},\underline{n})=\Hom_{\cat{FinSet}}(\underline{m},\underline{n})$, el conjunt de totes les aplicacions $\underline{m}\to\underline{n}$. L'aplicació $\iota_{\underline{m},\underline{n}}$ és, doncs, la identitat d'aquest conjunt, i en particular exhaustiva.

\textbf{Essencialment exhaustiva.} Sigui $X$ un conjunt finit, i sigui $n$ el seu nombre d'elements. Per definició de cardinal finit, hi ha una bijecció $X\to\underline{n}$ (per a $n=0$, $X=\varnothing=\underline{0}$ i n'hi ha prou amb la identitat). Una bijecció és un isomorfisme de $\cat{FinSet}$, de manera que $X\cong\iota(\underline{n})$. Tot objecte de $\cat{FinSet}$ és, doncs, isomorf a una imatge de $\iota$.

\textbf{No exhaustiva sobre objectes.} El conjunt $\{2\}$ és un objecte de $\cat{FinSet}$, però no és cap dels $\underline{n}$: té un sol element, i l'únic $\underline{n}$ amb un element és $\underline{1}=\{1\}\neq\{2\}$. No és, per tant, igual a cap imatge de $\iota$, tot i ser isomorf a $\underline{1}$.

\emph{Per què és important.} Aquest functor reuneix les tres propietats ---ple, fidel i essencialment exhaustiu--- que al capítol~\ref{cap:transformacions} caracteritzaran les equivalències de categories, i alhora mostra que no cal arribar a tots els objectes en sentit estricte. Observem també que per a cada $X$ hem triat \emph{una} bijecció $X\to\underline{n}$; per construir un functor en sentit invers caldria triar-ne una per a cada conjunt finit alhora, i aquest és el punt on el capítol~\ref{cap:transformacions} recorrerà a l'axioma de l'elecció.
\end{solucio}


\chapter[Transformacions naturals]{\texorpdfstring{Transformacions naturals\\ i equivalències}{Transformacions naturals i equivalències}}

En aquest capítol resolem amb detall exercicis sobre \emph{transformacions naturals}, \emph{isomorfismes naturals} i \emph{equivalències de categories}. L'èmfasi és a comprovar la condició de naturalitat en casos concrets, a distingir un isomorfisme d'un isomorfisme natural, i a reconèixer una equivalència mitjançant les tres propietats que la caracteritzen. Convé intentar cada exercici abans de llegir-ne la solució.

\section{Definició de transformació natural}

\begin{exercici}
Siguin $F,G\colon\cat{C}\to\cat{D}$ dos functors i suposem que, per a cada objecte $A$, tenim un morfisme $\eta_A\colon F(A)\to G(A)$. Comprova que, per verificar que $\eta$ és una transformació natural, n'hi ha prou de comprovar la condició de naturalitat sobre un conjunt de morfismes que \emph{generi} $\cat{C}$ per composició, i que en el cas d'una categoria lliure $\Free(G_0)$ això vol dir comprovar-la només sobre les arestes del graf $G_0$.
\end{exercici}

\begin{solucio}
La condició de naturalitat per a un morfisme componible diu $G(f)\comp\eta_A=\eta_B\comp F(f)$. Suposem que val per a dos morfismes componibles $f\colon A\to B$ i $g\colon B\to C$. Aleshores, per a la composició $g\comp f$,
\[
G(g\comp f)\comp\eta_A=G(g)\comp G(f)\comp\eta_A=G(g)\comp\eta_B\comp F(f)=\eta_C\comp F(g)\comp F(f)=\eta_C\comp F(g\comp f),
\]
on hem fet servir la functorialitat de $G$ i $F$ i les condicions de naturalitat per a $f$ i per a $g$. Per tant, si la naturalitat val per a un conjunt de morfismes, val per a totes les seves composicions. També val trivialment per a les identitats, perquè $G(\id_A)\comp\eta_A=\eta_A=\eta_A\comp F(\id_A)$.

En una categoria lliure $\Free(G_0)$, tot morfisme és un camí, és a dir, una composició d'arestes del graf $G_0$. Per l'argument anterior, doncs, la naturalitat sobre tots els morfismes es dedueix de la naturalitat sobre les arestes. Això redueix molt la feina de comprovació.
\end{solucio}

\begin{exercici}
Considera la categoria $\cat{2}=(0\to 1)$ i dos functors $F,G\colon\cat{2}\to\cat{D}$. Descriu explícitament què és una transformació natural $\eta\colon F\Rightarrow G$.
\end{exercici}

\begin{solucio}
Un functor $F\colon\cat{2}\to\cat{D}$ és el mateix que un morfisme de $\cat{D}$: tria dos objectes $F(0),F(1)$ i un morfisme $F(u)\colon F(0)\to F(1)$, on $u\colon 0\to 1$ és l'única fletxa no trivial. Igualment $G$ dona un morfisme $G(u)\colon G(0)\to G(1)$.

Una transformació natural $\eta\colon F\Rightarrow G$ consisteix en dos components $\eta_0\colon F(0)\to G(0)$ i $\eta_1\colon F(1)\to G(1)$, subjectes a la condició de naturalitat per a $u$:
\[
G(u)\comp\eta_0=\eta_1\comp F(u).
\]
És a dir, una transformació natural entre dos morfismes de $\cat{D}$ és exactament un \textbf{quadrat commutatiu}
\[
\begin{tikzpicture}[baseline=(m.center)]
  \matrix (m) [matrix of math nodes, row sep=2.6em, column sep=3.2em] {
    F(0) & F(1) \\
    G(0) & G(1) \\
  };
  \draw[-{Stealth[scale=1.1]}] (m-1-1) -- node[above] {$F(u)$} (m-1-2);
  \draw[-{Stealth[scale=1.1]}] (m-1-1) -- node[left] {$\eta_0$} (m-2-1);
  \draw[-{Stealth[scale=1.1]}] (m-1-2) -- node[right] {$\eta_1$} (m-2-2);
  \draw[-{Stealth[scale=1.1]}] (m-2-1) -- node[below] {$G(u)$} (m-2-2);
\end{tikzpicture}
\]
Aquest exemple explica per què els quadrats commutatius apareixen tan sovint: són les transformacions naturals més simples possibles.
\end{solucio}

\section{Exemples de transformacions naturals}

\begin{exercici}\label{pr-sintaxi-semantica}
Considerem els functors sintàctic i semàntic $F,G\colon\cat{FinSet}\to\cat{Set}$ i la interpretació $\eta_X\colon F(X)\to G(X)$ de la secció~\ref{sec:cap-transformacio} de la Teoria. Comprova que $F$ i $G$ són functors i que $\eta$ és una transformació natural $F\Rightarrow G$.
\end{exercici}

\begin{solucio}
\textbf{$F$ és un functor.} Recordem que $F(f)(\varphi)=\varphi[f]$ es defineix per recursió: $F(f)(x)=f(x)$ sobre una variable, i $F(f)$ commuta amb les connectives, per exemple $F(f)(\varphi\land\psi)=F(f)(\varphi)\land F(f)(\psi)$. Comprovem les dues igualtats per recursió sobre $\varphi$. Per a la identitat, sobre una variable $F(\id_X)(x)=x$, i si $F(\id_X)$ deixa fixes $\varphi$ i $\psi$, també deixa fixa $\varphi\land\psi$ (i anàlogament per a les altres connectives); per tant $F(\id_X)=\id_{F(X)}$. Per a la composició, siguin $f\colon X\to Y$ i $g\colon Y\to Z$. Sobre una variable,
\[
F(g\comp f)(x)=g(f(x))=F(g)\bigl(F(f)(x)\bigr),
\]
i si la igualtat $F(g\comp f)=F(g)\comp F(f)$ val per a $\varphi$ i $\psi$, val per a $\varphi\land\psi$, perquè les tres aplicacions commuten amb $\land$. Així, substituir per $f$ i després per $g$ és substituir per $g\comp f$.

\textbf{$G$ és un functor.} Recordem que $G(f)(\theta)(w)=\theta(w\comp f)$ per a $\theta\in G(X)$ i $w\colon Y\to\{0,1\}$. Per a la identitat, $G(\id_X)(\theta)(v)=\theta(v\comp\id_X)=\theta(v)$. Per a la composició, si $f\colon X\to Y$, $g\colon Y\to Z$ i $u\colon Z\to\{0,1\}$,
\[
\begin{aligned}
G(g\comp f)(\theta)(u)
&=\theta\bigl(u\comp(g\comp f)\bigr)
=\theta\bigl((u\comp g)\comp f\bigr)\\
&=G(f)(\theta)(u\comp g)
=\bigl(G(g)\comp G(f)\bigr)(\theta)(u),
\end{aligned}
\]
on el pas clau és l'associativitat de la composició.

\textbf{Lema de substitució.} Per a tota $\varphi\in F(X)$ i tota $w\colon Y\to\{0,1\}$,
\[
(w\comp f)\models\varphi
\quad\Longleftrightarrow\quad
w\models\varphi[f].
\]
Per recursió sobre $\varphi$. Si $\varphi$ és una variable $x$, els dos costats diuen que $w(f(x))=1$, perquè $(w\comp f)(x)=w(f(x))$ i $x[f]=f(x)$. Si $\varphi=\psi\land\chi$, aleshores $(w\comp f)\models\varphi$ si i només si $(w\comp f)\models\psi$ i $(w\comp f)\models\chi$; per hipòtesi de recursió, si i només si $w\models\psi[f]$ i $w\models\chi[f]$, és a dir, si i només si $w\models\psi[f]\land\chi[f]=\varphi[f]$. Les altres connectives són anàlogues.

\textbf{Naturalitat.} Avaluem els dos camins del quadrat de naturalitat sobre $\varphi\in F(X)$ i $w\colon Y\to\{0,1\}$:
\[
\begin{aligned}
\bigl(\eta_Y(F(f)(\varphi))\bigr)(w)&=\eta_Y(\varphi[f])(w)=[\,w\models\varphi[f]\,],\\
\bigl(G(f)(\eta_X(\varphi))\bigr)(w)&=\eta_X(\varphi)(w\comp f)=[\,(w\comp f)\models\varphi\,],
\end{aligned}
\]
on $[\,\cdot\,]$ val $1$ o $0$ segons que la condició sigui certa o falsa. Pel lema de substitució, les dues expressions coincideixen per a tota $w$, de manera que $\eta_Y\comp F(f)=G(f)\comp\eta_X$ per a tota $f$, i $\eta$ és natural.
\end{solucio}

\begin{exercici}
Comprova en detall que l'avaluació $\eta_V\colon V\to V^{**}$, $\eta_V(v)(\varphi)=\varphi(v)$, és una transformació natural de la identitat cap al functor bidual (o doble dual), sobre la categoria dels espais vectorials sobre un cos $\mathbb K$ i les aplicacions lineals.
\end{exercici}

\begin{solucio}
Recordem primer com actua el bidual sobre morfismes. Donada una aplicació lineal $T\colon V\to W$, el seu dual $T^*\colon W^*\to V^*$ envia $\psi\mapsto\psi\comp T$, i el bidual $T^{**}\colon V^{**}\to W^{**}$ envia $\Phi\mapsto\Phi\comp T^*$.

La condició de naturalitat per a $T\colon V\to W$ és el quadrat
\[
T^{**}\comp\eta_V=\eta_W\comp T.
\]
Comprovem que tots dos costats coincideixen aplicats a un vector $v\in V$ i avaluats en una forma $\psi\in W^*$. D'una banda,
\[
\begin{aligned}
\big(T^{**}(\eta_V(v))\big)(\psi)
&=\big(\eta_V(v)\comp T^*\big)(\psi)
=\eta_V(v)\big(T^*(\psi)\big)\\
&=\eta_V(v)(\psi\comp T)
=(\psi\comp T)(v)
=\psi(T(v)).
\end{aligned}
\]
De l'altra,
\[
\big(\eta_W(T(v))\big)(\psi)=\psi\big(T(v)\big).
\]
Els dos resultats coincideixen per a tot $v$ i tot $\psi$, de manera que $T^{**}\comp\eta_V=\eta_W\comp T$. Per tant $\eta$ és natural. Cap pas no ha requerit triar una base: aquesta és la manifestació concreta de la naturalitat.
\end{solucio}

\begin{exercici}
Siguin $P$ i $Q$ preordres i $F,G,H\colon P\to Q$ aplicacions monòtones. Interpreta les transformacions naturals $F\Rightarrow G$ i comprova que la composició vertical correspon a la transitivitat de la desigualtat puntual.
\end{exercici}

\begin{solucio}
Com hem vist a la teoria, hi ha una transformació natural $F\Rightarrow G$ si i només si $F(x)\leq G(x)$ per a tot $x\in P$; i quan existeix, és única, perquè entre dos objectes d'un preordre hi ha com a màxim una fletxa.

Si tenim $\eta\colon F\Rightarrow G$ i $\theta\colon G\Rightarrow H$, això vol dir $F(x)\leq G(x)$ i $G(x)\leq H(x)$ per a tot $x$. La composició vertical $\theta\comp\eta\colon F\Rightarrow H$ existeix perquè, per transitivitat, $F(x)\leq H(x)$. Així, la composició vertical de transformacions naturals entre aplicacions monòtones és exactament la transitivitat de la relació \enquote{puntualment menor o igual}. La categoria de functors $Q^{P}$ resulta ser, en aquest cas, un preordre: el preordre de les aplicacions monòtones $P\to Q$ amb l'ordre puntual.
\end{solucio}

\section{Isomorfismes naturals}

\begin{exercici}[Ampliació, per a lectors amb àlgebra lineal]\label{ex:dual-no-natural}
Comprova que tot espai vectorial real $V$ de dimensió finita és isomorf al seu dual $V^*$, però que aquesta identificació no és \emph{canònica}: no existeix cap família d'isomorfismes $\eta_V\colon V\to V^*$ (per a $V$ espai vectorial real de dimensió finita) que sigui invariant sota tots els isomorfismes lineals, en el sentit que, per a tot isomorfisme $T\colon V\to W$, es compleixi $\eta_W\comp T=(T^{-1})^*\comp\eta_V$.
\end{exercici}

\begin{solucio}
Si $V$ té dimensió finita $n$, aleshores $V^*$ també té dimensió $n$, i dos espais reals de la mateixa dimensió finita són isomorfs. Un isomorfisme concret s'obté triant una base $\{e_1,\dots,e_n\}$ de $V$ i enviant $e_i$ a la base dual $\{e_i^*\}$. L'isomorfisme depèn de la base triada.

Cal notar primer que la comparació s'ha de plantejar amb cura, perquè la identitat és un functor covariant $\cat{Vect}_{\mathbb R}\to\cat{Vect}_{\mathbb R}$ i el dual és \emph{contravariant}, és a dir, un functor $\cat{Vect}_{\mathbb R}\op\to\cat{Vect}_{\mathbb R}$; entre functors de variància diferent no té sentit la noció de quadrat de naturalitat que hem definit. Per això formulem la invariància com a compatibilitat amb els \emph{isomorfismes} lineals, on el dual $(T^{-1})^*$ i $T$ van en el mateix sentit.

Que no existeixi una tal família es pot veure amb l'escalament. Sigui $V\neq 0$ i prenem l'homotècia $T=2\,\id_V$, que és un isomorfisme lineal (aquí fem servir que treballem sobre $\mathbb R$, on $2\neq 0$). Aleshores $T^{-1}=\tfrac12\,\id_V$ i $(T^{-1})^*=\tfrac12\,\id_{V^*}$. La condició d'invariància demanaria
\[
\eta_V\comp(2\,\id_V)=\bigl(\tfrac12\,\id_{V^*}\bigr)\comp\eta_V,
\]
és a dir, $2\,\eta_V=\tfrac12\,\eta_V$, o equivalentment $\tfrac32\,\eta_V=0$. Com que $\tfrac32\neq 0$, això força $\eta_V=0$, que no és un isomorfisme. Per tant no hi ha cap identificació $V\cong V^*$ invariant sota tots els isomorfismes lineals.

\begin{observacio}
El punt clau de l'argument és triar un escalar $\lambda$ amb $\lambda\neq\lambda^{-1}$, és a dir $\lambda^2\neq 1$: de la igualtat $\lambda\,\eta_V=\lambda^{-1}\eta_V$ només se'n dedueix $\eta_V=0$ quan $\lambda-\lambda^{-1}\neq 0$. No n'hi ha prou, doncs, amb $\lambda\neq 0,1$: sobre $\mathbb R$ l'escalar $\lambda=-1$ compliria $\lambda\neq 0,1$ però $\lambda^2=1$, i no serviria. Per això prenem $\lambda=2$, amb $\lambda^2=4\neq 1$. Sobre cossos petits pot no existir cap escalar amb $\lambda^2\neq 1$ (per exemple a $\mathbb F_2$ i $\mathbb F_3$ tots els no nuls compleixen $\lambda^2=1$), i l'argument de l'escalament s'ha d'ajustar; per això aquí ens restringim a $\mathbb R$.
\end{observacio}

Aquest contrast amb el bidual ---que, restringit als espais de dimensió finita, sí que és naturalment isomorf a la identitat--- és l'exemple canònic de la diferència entre \enquote{isomorf} i \enquote{naturalment isomorf}.
\end{solucio}

\begin{exercici}
Comprova que la relació \enquote{ser naturalment isomorf} és una relació d'equivalència sobre els functors $\cat{C}\to\cat{D}$.
\end{exercici}

\begin{solucio}
Cal veure reflexivitat, simetria i transitivitat.

\textbf{Reflexivitat}: la transformació identitat $\id_F\colon F\Rightarrow F$ té components $\id_{F(A)}$, que són isomorfismes; per tant $F\cong F$.

\textbf{Simetria}: si $\eta\colon F\Rightarrow G$ és un isomorfisme natural, hem vist a la teoria que els inversos $\eta_A^{-1}$ formen una transformació natural $\eta^{-1}\colon G\Rightarrow F$, també amb components isomorfismes; per tant $G\cong F$.

\textbf{Transitivitat}: si $\eta\colon F\Rightarrow G$ i $\theta\colon G\Rightarrow H$ són isomorfismes naturals, la composició vertical $\theta\comp\eta$ té components $\theta_A\comp\eta_A$, composició de dos isomorfismes, que és un isomorfisme. Per tant $F\cong H$.

De fet, aquesta és simplement la relació d'isomorfisme entre objectes de la categoria de functors $\cat{D}^{\cat{C}}$, que ja sabem que és una relació d'equivalència.
\end{solucio}

\section{La categoria de functors}

\begin{exercici}
Descriu la categoria de functors $\cat{D}^{\cat{1}}$, on $\cat{1}$ és la categoria amb un sol objecte i només la identitat.
\end{exercici}

\begin{solucio}
Un functor $F\colon\cat{1}\to\cat{D}$ tria un únic objecte $F(\ast)=D$ de $\cat{D}$ (i envia l'única identitat a $\id_D$). Per tant, els objectes de $\cat{D}^{\cat{1}}$ es corresponen amb els objectes de $\cat{D}$.

Una transformació natural $\eta\colon F\Rightarrow G$ entre dos tals functors té un únic component $\eta_\ast\colon F(\ast)\to G(\ast)$, és a dir, un morfisme $D\to D'$ de $\cat{D}$; la condició de naturalitat per a l'única fletxa (la identitat) es compleix trivialment. Per tant, els morfismes de $\cat{D}^{\cat{1}}$ es corresponen amb els morfismes de $\cat{D}$.

Concloem que $\cat{D}^{\cat{1}}$ és \emph{isomorfa} a $\cat{D}$: prendre functors des de la categoria terminal $\cat{1}$ recupera la categoria original.
\end{solucio}

\begin{exercici}
Comprova que un functor $F\colon\cat{2}\to\cat{D}$ és un objecte de la categoria de functors $\cat{D}^{\cat{2}}$, i que aquesta categoria és la \textbf{categoria de fletxes} de $\cat{D}$: els seus objectes són els morfismes de $\cat{D}$ i els seus morfismes són els quadrats commutatius.
\end{exercici}

\begin{solucio}
Ja hem vist que un functor $\cat{2}\to\cat{D}$ és un morfisme de $\cat{D}$, i que una transformació natural entre dos d'aquests functors és un quadrat commutatiu. Per tant, els objectes de $\cat{D}^{\cat{2}}$ són els morfismes de $\cat{D}$, i els morfismes de $\cat{D}^{\cat{2}}$ són els quadrats commutatius, amb la composició vertical corresponent a apilar quadrats. Aquesta categoria s'anomena \textbf{categoria de fletxes} de $\cat{D}$ i es denota sovint $\cat{D}^{\to}$. És un primer exemple de com les categories de functors produeixen construccions útils a partir de categories índex molt simples.
\end{solucio}

\section{Equivalència de categories}

\begin{exercici}
Comprova que la categoria dels conjunts finits és equivalent a la seva subcategoria plena formada pels conjunts $\underline{0}=\varnothing$ i $\underline{n}=\{1,\dots,n\}$ per a $n\geq 1$.
\end{exercici}

\begin{solucio}
Sigui $\cat{S}$ aquesta subcategoria plena i $\iota\colon\cat{S}\hookrightarrow\cat{FinSet}$ la inclusió. A l'exercici resolt~\ref{pr-esquelet-finset} del capítol~\ref{cap:functors} vam comprovar que $\iota$ és ple, fidel i essencialment exhaustiu. Pel teorema~\ref{teo:caract-equivalencia}, $\iota$ forma part d'una equivalència, i per tant $\cat{FinSet}\simeq\cat{S}$.

Per construir un quasiinvers $G\colon\cat{FinSet}\to\cat{S}$, la demostració del teorema tria, per a cada conjunt finit $X$ de $n$ elements, una bijecció $\varepsilon_X\colon\underline{n}\to X$, i posa $G(X)=\underline{n}$. Com que $\cat{FinSet}$ és gran, aquesta tria simultània fa servir l'axioma de l'elecció en la forma que explica l'observació que segueix el teorema. Intuïtivament: per fer teoria de categories amb conjunts finits, n'hi ha prou de considerar un conjunt de cada mida.
\end{solucio}

\begin{exercici}\label{pr-equivalencia-preserva}
Comprova que una equivalència de categories preserva monomorfismes, isomorfismes i objectes terminals. Explica per què, en canvi, \emph{no} preserva propietats estrictes, com ara la de tenir exactament un objecte.
\end{exercici}

\begin{solucio}
Sigui $F\colon\cat{C}\to\cat{D}$ part d'una equivalència, amb quasiinvers $G$ i isomorfismes naturals $\eta\colon\id_{\cat{C}}\Rightarrow G\comp F$ i $\varepsilon\colon F\comp G\Rightarrow\id_{\cat{D}}$, orientats com a la definició~\ref{def:equivalencia}. Tots dos són isomorfismes component a component, amb inversos $\eta^{-1}$ i $\varepsilon^{-1}$.

\textbf{Isomorfismes}: com que $F$ és un functor, preserva isomorfismes (proposició~\ref{prop:functor-preserva-iso}); i com que és ple i fidel ---per ser part d'una equivalència---, també els reflecteix. Així, $f$ és un isomorfisme a $\cat{C}$ si i només si $F(f)$ ho és a $\cat{D}$.

\textbf{Monomorfismes}: sigui $f\colon A\to B$ mono a $\cat{C}$; vegem que $F(f)\colon F(A)\to F(B)$ és mono a $\cat{D}$. Siguin $u,v\colon Z\to F(A)$ dos morfismes de $\cat{D}$ amb
\[
F(f)\comp u=F(f)\comp v.
\]
Apliquem el quasiinvers $G$ i fem servir la seva functorialitat:
\[
G(F(f))\comp G(u)=G\bigl(F(f)\comp u\bigr)=G\bigl(F(f)\comp v\bigr)=G(F(f))\comp G(v).
\]
Ara traslladem aquesta igualtat de $G(F(A))$ a $A$ mitjançant la naturalitat de $\eta$. Per a $f\colon A\to B$, el quadrat de naturalitat de $\eta$ dona
\[
G(F(f))\comp\eta_A=\eta_B\comp f,
\]
i, en ser $\eta_A,\eta_B$ isomorfismes, $G(F(f))=\eta_B\comp f\comp\eta_A^{-1}$. Substituint a la igualtat anterior i cancel·lant l'isomorfisme $\eta_B$ per l'esquerra,
\[
f\comp\bigl(\eta_A^{-1}\comp G(u)\bigr)=f\comp\bigl(\eta_A^{-1}\comp G(v)\bigr).
\]
Com que $f$ és mono a $\cat{C}$, en resulta $\eta_A^{-1}\comp G(u)=\eta_A^{-1}\comp G(v)$, i cancel·lant l'isomorfisme $\eta_A^{-1}$ obtenim $G(u)=G(v)$. Finalment, $G$ és fidel ---per ser part d'una equivalència---, de manera que $u=v$. Per tant $F(f)$ és mono.

\textbf{Objectes terminals}: sigui $1$ terminal a $\cat{C}$; vegem que $F(1)$ és terminal a $\cat{D}$. Sigui $D$ un objecte qualsevol de $\cat{D}$.

\emph{Existència}: com que $1$ és terminal, hi ha un únic morfisme $!\colon G(D)\to 1$ a $\cat{C}$. Aleshores el morfisme
\[
D\xrightarrow{\ \varepsilon_D^{-1}\ }F(G(D))\xrightarrow{\ F(!)\ }F(1)
\]
és un morfisme $D\to F(1)$, cosa que en garanteix l'existència.

\emph{Unicitat}: siguin $h,h'\colon D\to F(1)$ dos morfismes. Aplicant $G$, obtenim $G(h),G(h')\colon G(D)\to G(F(1))$. Component amb l'isomorfisme $\eta_1^{-1}\colon G(F(1))\to 1$, els morfismes $\eta_1^{-1}\comp G(h)$ i $\eta_1^{-1}\comp G(h')$ van de $G(D)$ a $1$ i, per tant, coincideixen per la unicitat cap a l'objecte terminal $1$. Com que $\eta_1^{-1}$ és invertible, resulta $G(h)=G(h')$, i com que $G$ és fidel, $h=h'$. Per tant $F(1)$ és terminal.

\textbf{Per què no tota propietat}: una equivalència \emph{no} preserva les propietats que depenen de la \emph{igualtat} d'objectes. Per exemple, \enquote{tenir exactament un objecte} o \enquote{ser esquelètica} no es conserven: la categoria amb un sol objecte i la seva versió amb dos objectes isomorfs són equivalents, però la primera té un objecte i la segona en té dos.

Convé distingir, a més, dues menes de nocions entre les que hem comprovat. Ser isomorfisme i ser monomorfisme són \emph{propietats de morfismes}, definides per equacions o per cancel·lació; l'equivalència les preserva, i de la mateixa manera es comprova que preserva els epimorfismes. Ser objecte terminal és, en canvi, una \emph{propietat universal} d'un objecte, que el caracteritza només llevat d'isomorfisme; el que hem vist és que la imatge d'un terminal torna a ser terminal. Més endavant trobarem altres construccions universals ---objectes inicials, límits i colímits--- i en comprovarem el comportament sota equivalència cas per cas.
\end{solucio}


\section{Composició horitzontal}

\begin{exercici}\label{pr-comp-horitzontal}
Siguin $\eta\colon F\Rightarrow G$, amb $F,G\colon\cat{C}\to\cat{D}$, i $\theta\colon K\Rightarrow L$, amb $K,L\colon\cat{D}\to\cat{E}$, com a la definició de composició horitzontal. Escriu els tipus de
\[
\theta_{G(A)}\comp K(\eta_A)
\qquad\text{i}\qquad
L(\eta_A)\comp\theta_{F(A)},
\]
explica per què coincideixen per la naturalitat de $\theta$, i verifica el quadrat de naturalitat de $\theta\ast\eta\colon K\comp F\Rightarrow L\comp G$.
\end{exercici}

\begin{solucio}
\textbf{Els tipus.} El morfisme $\eta_A\colon F(A)\to G(A)$ viu a $\cat{D}$, i $K$ el porta a $K(\eta_A)\colon K(F(A))\to K(G(A))$ a $\cat{E}$. La component de $\theta$ a l'objecte $G(A)$ de $\cat{D}$ és $\theta_{G(A)}\colon K(G(A))\to L(G(A))$. Per tant
\[
\theta_{G(A)}\comp K(\eta_A)\colon K(F(A))\longrightarrow L(G(A)).
\]
Anàlogament, $\theta_{F(A)}\colon K(F(A))\to L(F(A))$ i $L(\eta_A)\colon L(F(A))\to L(G(A))$, de manera que
\[
L(\eta_A)\comp\theta_{F(A)}\colon K(F(A))\longrightarrow L(G(A)).
\]
Tots dos morfismes tenen, doncs, el mateix origen i el mateix destí, i té sentit preguntar-se si coincideixen.

\textbf{Per què coincideixen.} Són els dos camins del quadrat de naturalitat de $\theta$ per al morfisme $\eta_A\colon F(A)\to G(A)$ de $\cat{D}$, dibuixat a l'observació~\ref{obs:comp-horitzontal} de la Teoria:
\[
\theta_{G(A)}\comp K(\eta_A)=L(\eta_A)\comp\theta_{F(A)}.
\]
La naturalitat de $\theta$ diu que aquest quadrat commuta per a \emph{tot} morfisme de $\cat{D}$; en particular, per a $\eta_A$. La component $(\theta\ast\eta)_A$ està, doncs, ben definida, sigui quina sigui l'expressió que fem servir.

\textbf{Naturalitat de $\theta\ast\eta$.} Sigui $f\colon A\to B$ un morfisme de $\cat{C}$. Cal veure que
\[
\begin{tikzpicture}[baseline=(m.center)]
  \matrix (m) [matrix of math nodes, row sep=2.8em, column sep=4.6em] {
    K(F(A)) & K(F(B)) \\
    L(G(A)) & L(G(B)) \\
  };
  \draw[-{Stealth[scale=1.1]}] (m-1-1) -- node[above] {$K(F(f))$} (m-1-2);
  \draw[-{Stealth[scale=1.1]}] (m-1-1) -- node[left] {$(\theta\ast\eta)_A$} (m-2-1);
  \draw[-{Stealth[scale=1.1]}] (m-1-2) -- node[right] {$(\theta\ast\eta)_B$} (m-2-2);
  \draw[-{Stealth[scale=1.1]}] (m-2-1) -- node[below] {$L(G(f))$} (m-2-2);
\end{tikzpicture}
\qquad
(\theta\ast\eta)_B\comp K(F(f))=L(G(f))\comp(\theta\ast\eta)_A
\]
commuta. Fent servir l'expressió $\theta_{G(-)}\comp K(\eta_-)$ per a les components,
\[
\begin{aligned}
(\theta\ast\eta)_B\comp K(F(f))
&=\theta_{G(B)}\comp K(\eta_B)\comp K(F(f))
=\theta_{G(B)}\comp K\bigl(\eta_B\comp F(f)\bigr)\\
&=\theta_{G(B)}\comp K\bigl(G(f)\comp\eta_A\bigr)
=\theta_{G(B)}\comp K(G(f))\comp K(\eta_A)\\
&=L(G(f))\comp\theta_{G(A)}\comp K(\eta_A)
=L(G(f))\comp(\theta\ast\eta)_A.
\end{aligned}
\]
Hem fet servir, per ordre: la functorialitat de $K$, la naturalitat de $\eta$ per a $f$, de nou la functorialitat de $K$, i la naturalitat de $\theta$ per al morfisme $G(f)$ de $\cat{D}$. Cada naturalitat s'utilitza una sola vegada, i cadascuna al seu nivell: la de $\eta$ dins de $\cat{D}$, i la de $\theta$ dins de $\cat{E}$.
\end{solucio}


\ifpublicacioparcial\else
\chapter[Representabilitat i Yoneda]{\texorpdfstring{Propietats universals,\\ representabilitat i Yoneda}{Propietats universals, representabilitat i Yoneda}}

Els exercicis d'aquest capítol treballen les propietats universals a través de la representabilitat i el lema de Yoneda. L'objectiu és guanyar fluïdesa amb els functors representables, amb el càlcul de transformacions naturals que en surten, i amb l'argument ---repetitiu, un cop entès--- que dedueix unicitat llevat d'isomorfisme.

\section{Objectes inicials i terminals}

\begin{exercici}\label{pr-producte-terminal}
Siguin $A$ i $B$ dos objectes d'una categoria $\cat{C}$. Un \textbf{con} sobre $A$ i $B$ és una terna $(X,f,g)$ amb $f\colon X\to A$ i $g\colon X\to B$, i un morfisme de cons $(X,f,g)\to(X',f',g')$ és un morfisme $h\colon X\to X'$ de $\cat{C}$ amb $f'\comp h=f$ i $g'\comp h=g$. Comprova que els cons sobre $A$ i $B$ formen una categoria $\mathrm{Con}(A,B)$, i que un con $(P,p_1,p_2)$ és un producte de $A$ i $B$ si i només si és un objecte terminal de $\mathrm{Con}(A,B)$. Interpreta el resultat en la categoria d'un preordre $(Q,\le)$.
\end{exercici}

\begin{solucio}
\textbf{És una categoria.} La identitat de $(X,f,g)$ és $\id_X$, que compleix $f\comp\id_X=f$ i $g\comp\id_X=g$. Si $h\colon(X,f,g)\to(X',f',g')$ i $h'\colon(X',f',g')\to(X'',f'',g'')$ són morfismes de cons, la composta $h'\comp h$ també ho és, perquè
\[
f''\comp(h'\comp h)=(f''\comp h')\comp h=f'\comp h=f,
\]
i igual amb $g$. L'associativitat i les lleis d'identitat s'hereten de $\cat{C}$, perquè la composició és la de $\cat{C}$.

\textbf{Producte i objecte terminal.} Un morfisme de cons $h\colon(X,f,g)\to(P,p_1,p_2)$ és exactament un morfisme $h\colon X\to P$ que fa commutar els dos triangles,
\[
\begin{tikzpicture}[baseline=(m.center)]
  \matrix (m) [matrix of math nodes, row sep=2.6em, column sep=2.8em] {
     & X & \\
    A & P & B \\
  };
  \draw[-{Stealth[scale=1.0]}] (m-1-2) -- node[above left] {$f$} (m-2-1);
  \draw[-{Stealth[scale=1.0]}] (m-1-2) -- node[above right] {$g$} (m-2-3);
  \draw[-{Stealth[scale=1.0]},dashed] (m-1-2) -- node[right] {$h$} (m-2-2);
  \draw[-{Stealth[scale=1.0]}] (m-2-2) -- node[below] {$p_1$} (m-2-1);
  \draw[-{Stealth[scale=1.0]}] (m-2-2) -- node[below] {$p_2$} (m-2-3);
\end{tikzpicture}
\qquad
p_1\comp h=f,\quad p_2\comp h=g.
\]
Que $(P,p_1,p_2)$ sigui terminal a $\mathrm{Con}(A,B)$ diu que, per a cada con $(X,f,g)$, hi ha exactament un morfisme de cons cap a $(P,p_1,p_2)$, és a dir, exactament un $h$ amb $p_1\comp h=f$ i $p_2\comp h=g$. Això és, paraula per paraula, la definició de producte (subsecció~\ref{subsec:producte-coproducte}). Les dues afirmacions són, doncs, la mateixa.

Com a conseqüència, la unicitat dels productes es dedueix de la dels objectes terminals (proposició~\ref{prop:unicitat-terminal}): dos productes de $A$ i $B$ són isomorfs per un únic isomorfisme \emph{de cons}, és a dir, compatible amb les projeccions. L'objecte $P$ tot sol, en canvi, pot tenir altres automorfismes que no respectin $p_1$ i $p_2$; per exemple, a $\cat{Set}$, l'intercanvi $(a,b)\mapsto(b,a)$ de $A\times A$.

\textbf{En un preordre.} A la categoria d'un preordre $(Q,\le)$, un con sobre $a$ i $b$ és un element $x$ amb $x\le a$ i $x\le b$, és a dir, una cota inferior comuna, i entre dos cons hi ha com a màxim un morfisme, que existeix quan $x\le x'$. Un objecte terminal de $\mathrm{Con}(a,b)$ és, doncs, una cota inferior $p$ tal que tota altra cota inferior $x$ compleix $x\le p$: la \emph{màxima cota inferior}, o ínfim, de $a$ i $b$. Els productes en un preordre són els ínfims, i en un ordre parcial són únics; en un preordre, únics llevat de la relació $p\le p'\le p$.
\end{solucio}

\begin{exercici}
Comprova que un objecte $Z$ és alhora inicial i terminal en una categoria si i només si, per a cada objecte $X$, hi ha exactament un morfisme $Z\to X$ i exactament un morfisme $X\to Z$. Un tal objecte s'anomena \emph{objecte zero}. Comprova que a $\cat{Pos}$ no n'hi ha cap i, per a qui conegui els grups, que a $\cat{Grp}$ el grup trivial n'és un.
\end{exercici}

\begin{solucio}
La primera afirmació és la conjunció literal de les dues definicions: $Z$ inicial dona el morfisme únic $Z\to X$, i $Z$ terminal, el morfisme únic $X\to Z$.

A $\cat{Pos}$, l'objecte inicial és el conjunt ordenat buit i el terminal és el d'un sol element. Un objecte zero hauria de ser isomorf a tots dos, i aquests dos no són isomorfs, perquè no hi ha cap aplicació $\{\ast\}\to\varnothing$. Per tant, $\cat{Pos}$ no té objecte zero.

Per al grup trivial $1=\{e\}$: donat un grup $G$, hi ha un únic homomorfisme $1\to G$ (l'element neutre ha d'anar al neutre), i un únic homomorfisme $G\to 1$ (tot va al neutre). Per tant $1$ és inicial i terminal a $\cat{Grp}$, és a dir, un objecte zero.
\end{solucio}

\section{Categories coma}

\begin{exercici}\label{pr-coma-lliure}
Sigui $G$ un graf, i sigui $U\colon\cat{Cat}\to\cat{Graph}$ el functor graf subjacent, on $\cat{Cat}$ té per objectes les categories petites. Descriu explícitament la categoria coma $(G\downarrow U)$: els seus objectes, els seus morfismes, la composició i les identitats. Demostra que la parella $(\Free(G),\eta_G)$, on $\eta_G\colon G\to U(\Free(G))$ és la inserció dels generadors, és un objecte inicial de $(G\downarrow U)$.
\end{exercici}

\begin{solucio}
\textbf{Objectes i morfismes.} Un objecte de $(G\downarrow U)$ és una parella $(\cat{D},\varphi)$ formada per una categoria petita $\cat{D}$ i un morfisme de grafs $\varphi\colon G\to U(\cat{D})$. Un morfisme $(\cat{D},\varphi)\to(\cat{D}',\varphi')$ és un functor $H\colon\cat{D}\to\cat{D}'$ que fa commutar el triangle
\[
\begin{tikzpicture}[baseline=(m.center)]
  \matrix (m) [matrix of math nodes, row sep=2.6em, column sep=2.4em] {
     & G & \\
    U(\cat{D}) & & U(\cat{D}') \\
  };
  \draw[-{Stealth[scale=1.0]}] (m-1-2) -- node[above left] {$\varphi$} (m-2-1);
  \draw[-{Stealth[scale=1.0]}] (m-1-2) -- node[above right] {$\varphi'$} (m-2-3);
  \draw[-{Stealth[scale=1.0]}] (m-2-1) -- node[below] {$U(H)$} (m-2-3);
\end{tikzpicture}
\qquad
U(H)\comp\varphi=\varphi'.
\]

\textbf{És una categoria.} La identitat de $(\cat{D},\varphi)$ és $\id_{\cat{D}}$, perquè $U(\id_{\cat{D}})\comp\varphi=\varphi$. Si $H\colon(\cat{D},\varphi)\to(\cat{D}',\varphi')$ i $H'\colon(\cat{D}',\varphi')\to(\cat{D}'',\varphi'')$, la composta $H'\comp H$ també és un morfisme, per la functorialitat de $U$:
\[
U(H'\comp H)\comp\varphi=U(H')\comp U(H)\comp\varphi=U(H')\comp\varphi'=\varphi''.
\]
L'associativitat i les lleis d'identitat s'hereten de $\cat{Cat}$.

\textbf{$(\Free(G),\eta_G)$ és inicial.} Sigui $(\cat{D},\varphi)$ un objecte qualsevol. Un morfisme $(\Free(G),\eta_G)\to(\cat{D},\varphi)$ és un functor $\bar\varphi\colon\Free(G)\to\cat{D}$ amb $U(\bar\varphi)\comp\eta_G=\varphi$, és a dir, que coincideix amb $\varphi$ sobre els vèrtexs i sobre els camins de longitud~$1$.
\begin{itemize}
  \item \emph{Existència.} Definim $\bar\varphi$ sobre els objectes com $\varphi$, sobre el camí buit en un vèrtex $x$ com $\id_{\varphi(x)}$, i sobre un camí $(a_1,\dots,a_n)$ recorregut en aquest ordre com la composta $\varphi(a_n)\comp\cdots\comp\varphi(a_1)$. Com que la composició de $\Free(G)$ és la concatenació, $\bar\varphi$ respecta composicions i identitats.
  \item \emph{Unicitat.} Tot camí és la composta de les seves arestes, i tot camí buit és una identitat. Un functor amb $U(\bar\varphi)\comp\eta_G=\varphi$ queda, per tant, determinat sobre tots els morfismes per la seva acció sobre les arestes, que està fixada per $\varphi$.
\end{itemize}
Així, des de $(\Free(G),\eta_G)$ hi ha exactament un morfisme cap a cada objecte de $(G\downarrow U)$: és un objecte inicial. Observem que no ha calgut la noció d'adjunció; és aquesta propietat la que, al capítol d'adjuncions, farà de $\Free$ un adjunt per l'esquerra de $U$.
\end{solucio}

\section{Propietats universals com a representabilitat}

\begin{exercici}
Comprova que un objecte $1$ és terminal si i només si el functor $\Hom(-,1)$ és naturalment isomorf al functor constant $\Delta_{\{\ast\}}\colon\cat{C}\op\to\cat{Set}$ (exemple~\ref{ex:functor-constant}).
\end{exercici}

\begin{solucio}
Suposem $1$ terminal. Aleshores, per a cada objecte $X$, el conjunt $\Hom(X,1)$ té exactament un element, de manera que hi ha una bijecció $\Hom(X,1)\cong\{\ast\}$. Aquestes bijeccions són automàticament naturals: l'única aplicació possible cap a un conjunt d'un punt fa commutar qualsevol quadrat. Per tant $\Hom(-,1)$ és naturalment isomorf al functor constant $\Delta_{\{\ast\}}$.

Recíprocament, si $\Hom(-,1)\cong\Delta_{\{\ast\}}$, aleshores cada $\Hom(X,1)$ té exactament un element, és a dir, hi ha exactament un morfisme $X\to 1$ per a cada $X$, i $1$ és terminal.
\end{solucio}

\begin{exercici}
Sigui $\cat{C}$ una categoria localment petita amb un objecte terminal $1$. Comprova que el \textbf{functor de punts globals}
\[
\Hom(1,-)\colon\cat{C}\longrightarrow\cat{Set}
\]
envia cada objecte $X$ a $\Hom(1,X)$. Aquests morfismes s'anomenen \emph{punts globals} de $X$. Calcula aquest conjunt a $\cat{Set}$.
\end{exercici}

\begin{solucio}
Per definició, $\Hom(1,-)$ envia $X$ a $\Hom(1,X)$, el conjunt de morfismes des de l'objecte terminal, és a dir, dels punts globals de $X$. A $\cat{Set}$, l'objecte terminal és un conjunt d'un sol element $\{\ast\}$, i una aplicació $\{\ast\}\to X$ equival a triar un element de $X$: la imatge de $\ast$. Per tant $\Hom(1,X)\cong X$ de manera natural, i a $\cat{Set}$ els punts globals recuperen exactament els elements. En general, però, $\Hom(1,-)$ pot perdre informació: hi ha categories on objectes diferents tenen els mateixos punts globals.
\end{solucio}

\section{El lema de Yoneda}

\begin{exercici}
Sigui $\cat{C}$ localment petita i $A$ un objecte. Aplicant el lema de Yoneda amb $P=\Hom(-,A)$, comprova que les úniques transformacions naturals $\Hom(-,A)\Rightarrow\Hom(-,A)$ que hi ha es corresponen amb els endomorfismes de $A$, i que la identitat correspon a $\id_A$.
\end{exercici}

\begin{solucio}
El lema de Yoneda dona una bijecció
\[
\Nat\big(\Hom(-,A),\Hom(-,A)\big)\cong\Hom(-,A)(A)=\Hom(A,A),
\]
que envia $\eta$ a $\eta_A(\id_A)$. Així, cada transformació natural de $\Hom(-,A)$ en ell mateix correspon a un endomorfisme de $A$. La transformació identitat $\id_{\Hom(-,A)}$ té components la identitat de cada $\Hom(X,A)$; el seu valor a $\id_A$ és $\id_A$, de manera que correspon a l'endomorfisme identitat. Aquest càlcul és el nucli de la demostració que la immersió de Yoneda és plena i fidel.
\end{solucio}

\begin{exercici}\label{pr-yoneda-parts}
Sigui $A$ un conjunt i sigui $\mathcal P\colon\cat{Set}\op\to\cat{Set}$ el functor de parts contravariant. Descriu explícitament, mitjançant la bijecció del lema de Yoneda, les transformacions naturals $\Hom_{\cat{Set}}(-,A)\Rightarrow\mathcal P$. Quina transformació correspon a un subconjunt $B\subseteq A$? Comprova directament que la bijecció recupera $B$, i calcula-la en el cas $A=\{0,1\}$, $B=\{1\}$, sobre l'aplicació $\varphi\colon\{a,b,c\}\to A$ amb $\varphi(a)=\varphi(c)=1$ i $\varphi(b)=0$.
\end{exercici}

\begin{solucio}
El lema de Yoneda dona una bijecció
\[
\Phi\colon\Nat\bigl(\Hom(-,A),\mathcal P\bigr)\longrightarrow\mathcal P(A),
\qquad
\Phi(\alpha)=\alpha_A(\id_A),
\]
amb inversa $\Psi$, que envia un element $x\in\mathcal P(A)$ a la transformació $\Psi(x)_X(\varphi)=\mathcal P(\varphi)(x)$. Aquí un element de $\mathcal P(A)$ és un subconjunt $B\subseteq A$, i $\mathcal P(\varphi)$ és l'antiimatge. Per tant, la transformació natural que correspon a $B$ és
\[
\Psi(B)_X\colon\Hom(X,A)\to\mathcal P(X),
\qquad
\varphi\longmapsto\varphi^{-1}(B).
\]
És a dir: tota assignació natural d'aquest tipus s'obté per antiimatge d'un únic subconjunt fix de $A$, i subconjunts diferents donen transformacions naturals diferents.

\textbf{La bijecció recupera $B$.} En un sentit, $\Phi(\Psi(B))=\Psi(B)_A(\id_A)=\id_A^{-1}(B)=B$. En l'altre, si $\alpha$ és una transformació natural qualsevol i $B=\alpha_A(\id_A)$, la naturalitat de $\alpha$ aplicada a $\varphi\colon X\to A$ i avaluada sobre $\id_A$ dona
\[
\alpha_X(\varphi)=\alpha_X(\id_A\comp\varphi)=\mathcal P(\varphi)\bigl(\alpha_A(\id_A)\bigr)=\varphi^{-1}(B),
\]
de manera que $\alpha=\Psi(B)$. La naturalitat de $\Psi(B)$ es veu també directament: per a $g\colon Y\to X$, $(\varphi\comp g)^{-1}(B)=g^{-1}\bigl(\varphi^{-1}(B)\bigr)$.

\textbf{Un càlcul.} Amb $A=\{0,1\}$ i $B=\{1\}$, la component a $X=\{a,b,c\}$ envia $\varphi$ a $\varphi^{-1}(\{1\})=\{a,c\}$. Observem que, sota l'isomorfisme natural $\mathcal P\cong\Hom(-,\{0,1\})$ de l'exemple~\ref{ex:isos-naturals-cap2}, el subconjunt $B$ correspon a la seva funció característica $\chi_B\colon A\to\{0,1\}$, i la transformació $\varphi\mapsto\varphi^{-1}(B)$ correspon a $\varphi\mapsto\chi_B\comp\varphi$: és el cas $P=\Hom(-,\{0,1\})$ del lema, en què les transformacions naturals $\Hom(-,A)\Rightarrow\Hom(-,\{0,1\})$ són les postcomposicions amb morfismes $A\to\{0,1\}$.

\textbf{Lectura lògica.} Un subconjunt $B\subseteq A$ és un \emph{predicat} sobre $A$, i $\varphi^{-1}(B)$ és el predicat sobre $X$ que s'obté substituint $\varphi$: $x\in\varphi^{-1}(B)$ si i només si $\varphi(x)\in B$. El lema de Yoneda diu, doncs, que les maneres naturals de produir predicats a partir d'aplicacions cap a $A$ són exactament les substitucions en un predicat fix sobre $A$. És la mateixa operació de substitució que, al capítol~\ref{cap:transformacions}, feia natural la interpretació de les fórmules.
\end{solucio}

\begin{exercici}\label{pr-yoneda-grafs}
Siguin $V,E\colon\cat{Graph}\to\cat{Set}$ els functors de vèrtexs i d'arestes, siguin $\mathbf V$ el graf d'un sol vèrtex sense arestes i $\mathbf E$ el graf $0\xrightarrow{e}1$. Calcula, amb la versió covariant del lema de Yoneda, $\Nat(E,V)$ i $\Nat(V,E)$, i interpreta les transformacions que hi surtin.
\end{exercici}

\begin{solucio}
Per l'exemple~\ref{ex:isos-naturals-cap2}, $V\cong\Hom_{\cat{Graph}}(\mathbf V,-)$ i $E\cong\Hom_{\cat{Graph}}(\mathbf E,-)$. La versió covariant del lema dona, per a qualsevol functor $Q\colon\cat{Graph}\to\cat{Set}$, una bijecció entre les transformacions naturals que surten de $\Hom(A,-)$ i els elements de $Q(A)$. Per tant,
\[
\Nat(E,V)\;\cong\;V(\mathbf E)=\{0,1\},
\qquad
\Nat(V,E)\;\cong\;E(\mathbf V)=\varnothing.
\]
Observem la inversió: el primer càlcul avalua el functor d'\emph{arribada}, $V$, en l'objecte que \emph{representa} el functor de sortida, $E$.

\textbf{Les dues transformacions $E\Rightarrow V$.} L'element $0\in V(\mathbf E)$ correspon a la transformació que, a cada graf $G$, envia una aresta a una aplicació $E_G\to V_G$ calculada així: l'aresta $\varepsilon$ correspon al morfisme $\mathbf E\to G$ que la tria, i aquest envia el vèrtex $0$ a l'origen de $\varepsilon$. És, doncs, l'aplicació origen $s_G$. Anàlogament, l'element $1$ dona l'aplicació destí $t_G$. La naturalitat de $s$ i de $t$ és exactament la condició que un morfisme de grafs preservi origen i destí. Encara que en un graf on totes les arestes són bucles $s_G=t_G$, les dues transformacions són diferents, perquè a $\mathbf E$ les components difereixen.

\textbf{Cap transformació $V\Rightarrow E$.} També es veu directament: una transformació natural $V\Rightarrow E$ hauria de tenir una component $V_{\mathbf V}=\{\ast\}\to E_{\mathbf V}=\varnothing$, que no existeix. No hi ha cap manera natural de triar una aresta per a cada vèrtex en tots els grafs alhora.
\end{solucio}

\section{Conseqüències}

\begin{exercici}\label{pr-yoneda-preordre}
Sigui $(Q,\le)$ un preordre vist com a categoria. Descriu el prefeix representable $\Hom(-,a)$ i calcula, amb el lema de Yoneda, les transformacions naturals $\Hom(-,a)\Rightarrow\Hom(-,b)$. Dedueix que
\[
a\le b\quad\Longleftrightarrow\quad{\downarrow}a\subseteq{\downarrow}b,
\qquad\text{on } {\downarrow}a=\{x\in Q\mid x\le a\},
\]
i interpreta el resultat a la categoria $\cat{Prov}_T$ d'un sistema deductiu.
\end{exercici}

\begin{solucio}
\textbf{El representable.} Com que entre dos objectes hi ha com a màxim una fletxa, $\Hom(x,a)$ és un conjunt d'un sol element si $x\le a$, i és buit altrament. El prefeix $\Hom(-,a)$ és, doncs, la funció característica del conjunt ${\downarrow}a$ dels elements que estan per sota de $a$: \enquote{val} en $x$ exactament quan $x\in{\downarrow}a$.

\textbf{Les transformacions naturals.} Pel lema de Yoneda (o per la immersió plena i fidel),
\[
\Nat\bigl(\Hom(-,a),\Hom(-,b)\bigr)\;\cong\;\Hom(a,b),
\]
que té un element si $a\le b$ i cap altrament. Ho podem veure també directament: una transformació natural té una component $\Hom(x,a)\to\Hom(x,b)$ per a cada $x$, i una aplicació d'un conjunt d'un sol element en un conjunt buit no existeix. Hi ha, doncs, una transformació natural exactament quan $x\le a$ implica $x\le b$ per a tot $x$, és a dir, quan ${\downarrow}a\subseteq{\downarrow}b$; i aleshores és única, i la naturalitat és automàtica, perquè tots els valors tenen com a màxim un element.

\textbf{La conclusió.} Comparant les dues descripcions, $a\le b$ si i només si ${\downarrow}a\subseteq{\downarrow}b$. La implicació d'esquerra a dreta és la transitivitat; la de dreta a esquerra s'obté prenent $x=a$, que és exactament l'avaluació a $\id_a$ del lema de Yoneda. Com a cas particular, $a\cong b$ si i només si ${\downarrow}a={\downarrow}b$: dos elements són isomorfs quan tenen els mateixos elements per sota.

\textbf{A $\cat{Prov}_T$.} Aquí ${\downarrow}p$ és el conjunt de les fórmules que dedueixen $p$, i el resultat diu que $p\vdash_T q$ si i només si tota fórmula que dedueix $p$ dedueix també $q$. Una fórmula queda determinada, llevat d'interdeduïbilitat, per les fórmules de les quals es dedueix.
\end{solucio}

\begin{exercici}\label{pr-yoneda-monoide}
Sigui $M$ un monoide i $\cat{B}M$ la categoria d'un sol objecte $\ast$ associada, amb composició $g\comp f=g\cdot f$ (escrivim $\cdot$ per a l'operació del monoide i $e$ per al neutre, que és $\id_\ast$). Comprova que un prefeix $X\colon(\cat{B}M)\op\to\cat{Set}$ és el mateix que un conjunt amb una \emph{acció per la dreta} de $M$, i descriu el prefeix representable $\Hom(-,\ast)$. Aplica el lema de Yoneda per calcular $\Nat\bigl(\Hom(-,\ast),X\bigr)$ i, en particular, $\Nat\bigl(\Hom(-,\ast),\Hom(-,\ast)\bigr)$.
\end{exercici}

\begin{solucio}
\textbf{Prefeixos i accions.} Un prefeix $X$ dona un conjunt $X(\ast)$, que escrivim $X$, i per a cada $g\in M$ una aplicació $X(g)\colon X\to X$. Escrivim $x\cdot g=X(g)(x)$. La functorialitat contravariant, $X(g\comp f)=X(f)\comp X(g)$ i $X(e)=\id$, diu exactament
\[
(x\cdot g)\cdot f=x\cdot(g\cdot f),
\qquad
x\cdot e=x:
\]
és una acció per la dreta. Una transformació natural entre dos prefeixos és una aplicació $\eta\colon X\to Y$ amb $\eta(x\cdot g)=\eta(x)\cdot g$, és a dir, una aplicació \emph{equivariant}.

\textbf{El representable.} $\Hom(\ast,\ast)=M$, i un morfisme $g$ actua per precomposició: $m\mapsto m\comp g=m\cdot g$. El prefeix representable és, doncs, $M$ amb l'acció per la dreta per multiplicació, l'\emph{acció regular}.

\textbf{Yoneda.} El lema dona una bijecció $\Nat\bigl(\Hom(-,\ast),X\bigr)\cong X$: l'aplicació equivariant $\eta\colon M\to X$ correspon a $\eta(e)$, i recíprocament $x\in X$ dona l'aplicació $m\mapsto x\cdot m$. Una aplicació equivariant des de l'acció regular queda determinada pel lloc on envia l'element neutre.

Per a $X=M$ regular, les aplicacions equivariants $M\to M$ són exactament les multiplicacions per l'esquerra $\lambda_a\colon m\mapsto a\cdot m$, una per a cada $a\in M$. Com que $\lambda_a\comp\lambda_b=\lambda_{a\cdot b}$ i $\lambda_e=\id$, la immersió de Yoneda identifica $M$ amb aquest monoide d'aplicacions de $M$ en si mateix, i el fet que sigui fidel diu que $a\mapsto\lambda_a$ és injectiva. Obtenim el \textbf{teorema de Cayley} per a monoides: tot monoide és isomorf a un monoide d'aplicacions d'un conjunt en si mateix, amb la composició.
\end{solucio}

\chapter[Límits i colímits]{Límits i colímits}

Els exercicis d'aquest capítol treballen les propietats universals dels límits i colímits concrets, i l'argument ---sempre el mateix--- que dedueix unicitat i factoritzacions. Convé dibuixar cada diagrama abans de començar.

\section{Diagrames i cons}

\begin{exercici}
Comprova que un con sobre un diagrama $D\colon\cat{I}\to\cat{C}$ amb vèrtex $A$ és el mateix que una transformació natural $\Delta_A\Rightarrow D$, on $\Delta_A$ és el functor constant amb valor $A$.
\end{exercici}

\begin{solucio}
El functor constant $\Delta_A$ envia cada objecte $i$ de $\cat{I}$ a $A$ i cada morfisme a $\id_A$. Una transformació natural $\alpha\colon\Delta_A\Rightarrow D$ té components $\alpha_i\colon\Delta_A(i)=A\to D_i$, i la condició de naturalitat per a un morfisme $u\colon i\to j$ diu
\[
D_u\comp\alpha_i=\alpha_j\comp\Delta_A(u)=\alpha_j\comp\id_A=\alpha_j.
\]
Aquesta és exactament la condició de con. Recíprocament, tota família amb aquesta propietat defineix una transformació natural. Així, els cons sobre $D$ amb vèrtex $A$ i les transformacions naturals $\Delta_A\Rightarrow D$ coincideixen.
\end{solucio}

\section{Productes i equalitzadors}

\begin{exercici}
Comprova que a $\cat{Set}$ el producte categòric coincideix amb el producte cartesià, verificant la propietat universal.
\end{exercici}

\begin{solucio}
Sigui $A\times B=\{(a,b)\mid a\in A,\ b\in B\}$ amb $\pi_1(a,b)=a$ i $\pi_2(a,b)=b$. Donats $f\colon X\to A$ i $g\colon X\to B$, definim $\langle f,g\rangle\colon X\to A\times B$ per $\langle f,g\rangle(x)=(f(x),g(x))$. Aleshores $\pi_1\comp\langle f,g\rangle=f$ i $\pi_2\comp\langle f,g\rangle=g$. I és l'únic possible: si $h\colon X\to A\times B$ compleix $\pi_1\comp h=f$ i $\pi_2\comp h=g$, aleshores $h(x)=(\pi_1(h(x)),\pi_2(h(x)))=(f(x),g(x))=\langle f,g\rangle(x)$. Per tant el producte cartesià satisfà la propietat universal del producte.
\end{solucio}

\begin{exercici}
Comprova que tot equalitzador és un monomorfisme.
\end{exercici}

\begin{solucio}
Sigui $e\colon E\to A$ l'equalitzador de $f,g\colon A\to B$, i siguin $r,s\colon Z\to E$ amb $e\comp r=e\comp s$. Volem $r=s$. Sigui $h=e\comp r=e\comp s\colon Z\to A$. Aleshores $f\comp h=f\comp e\comp r=g\comp e\comp r=g\comp h$, de manera que $h$ satisfà la condició que defineix els cons de l'equalitzador. Per la propietat universal, hi ha un \emph{únic} $k\colon Z\to E$ amb $e\comp k=h$. Però tant $r$ com $s$ compleixen aquesta equació; per unicitat, $r=s$. Així $e$ és un monomorfisme.
\end{solucio}

\section{Pullbacks}

\begin{exercici}
Comprova que a $\cat{Set}$ el pullback de $f\colon A\to C$ i $g\colon B\to C$ és $P=\{(a,b)\in A\times B\mid f(a)=g(b)\}$, i que quan $g$ és la injecció d'un subconjunt $B\subseteq C$ el pullback recupera l'antiimatge $f^{-1}(B)$.
\end{exercici}

\begin{solucio}
Amb $P=\{(a,b)\mid f(a)=g(b)\}$ i les projeccions $p_1,p_2$, es compleix $f\comp p_1=g\comp p_2$ per construcció. Donats $r_1\colon X\to A$ i $r_2\colon X\to B$ amb $f\comp r_1=g\comp r_2$, l'aplicació $h(x)=(r_1(x),r_2(x))$ té imatge dins $P$ (perquè $f(r_1(x))=g(r_2(x))$) i és l'únic morfisme amb $p_1\comp h=r_1$, $p_2\comp h=r_2$. Això és la propietat universal.

Si $B\subseteq C$ i $g$ és la injecció, la condició $f(a)=g(b)=b$ força $b=f(a)$ amb $f(a)\in B$; per tant $P\cong\{a\in A\mid f(a)\in B\}=f^{-1}(B)$. El pullback recupera l'antiimatge.
\end{solucio}

\begin{exercici}[Lema d'enganxament de pullbacks]
Considera el diagrama de dos quadrats adossats en una categoria arbitrària:
\[
\begin{tikzpicture}[baseline=(m.center)]
  \matrix (m) [matrix of math nodes, column sep=3.2em, row sep=2.6em] {
    A & B & C \\
    A' & B' & C' \\
  };
  \draw[-{Stealth[scale=1.1]}] (m-1-1) -- node[above] {$u$} (m-1-2);
  \draw[-{Stealth[scale=1.1]}] (m-1-2) -- node[above] {$v$} (m-1-3);
  \draw[-{Stealth[scale=1.1]}] (m-2-1) -- node[below] {$u'$} (m-2-2);
  \draw[-{Stealth[scale=1.1]}] (m-2-2) -- node[below] {$v'$} (m-2-3);
  \draw[-{Stealth[scale=1.1]}] (m-1-1) -- node[left] {$a$} (m-2-1);
  \draw[-{Stealth[scale=1.1]}] (m-1-2) -- node[left] {$b$} (m-2-2);
  \draw[-{Stealth[scale=1.1]}] (m-1-3) -- node[right] {$c$} (m-2-3);
\end{tikzpicture}
\]
amb tots dos quadrats commutatius. Suposem que el \textbf{quadrat dret} (sobre $B,C,B',C'$) és un pullback. Comprova que el \textbf{quadrat esquerre} és un pullback si i només si el \textbf{rectangle exterior} (sobre $A,C,A',C'$) ho és.
\end{exercici}

\begin{solucio}
Treballem en una categoria arbitrària, usant només la propietat universal.

\textbf{($\Rightarrow$) Si el quadrat esquerre és pullback, l'exterior ho és.} Sigui $T$ un objecte amb morfismes $p\colon T\to A'$ i $q\colon T\to C$ compatibles per al rectangle exterior, és a dir amb $c\comp q=v'\comp u'\comp p$. Com que el quadrat dret és un pullback i els morfismes $q\colon T\to C$ i $u'\comp p\colon T\to B'$ satisfan $v'\comp(u'\comp p)=c\comp q$, hi ha un únic $r\colon T\to B$ amb $v\comp r=q$ i $b\comp r=u'\comp p$. Ara $r$ i $p$ són compatibles per al quadrat esquerre, perquè $b\comp r=u'\comp p$; com que aquest és un pullback, hi ha un únic $t\colon T\to A$ amb $u\comp t=r$ i $a\comp t=p$. Aleshores $(v\comp u)\comp t=v\comp r=q$ i $a\comp t=p$, de manera que $t$ factoritza el con exterior, i és únic per les unicitats successives.

\textbf{($\Leftarrow$) Si el rectangle exterior és pullback, l'esquerre ho és.} Sigui $T$ amb $p\colon T\to A'$ i $r\colon T\to B$ compatibles per al quadrat esquerre, és a dir amb $b\comp r=u'\comp p$. Posem $q=v\comp r\colon T\to C$. Aleshores $p$ i $q$ són compatibles per a l'exterior, perquè
\[
c\comp q=c\comp v\comp r=v'\comp b\comp r=v'\comp u'\comp p.
\]
Com que l'exterior és un pullback, hi ha un únic $t\colon T\to A$ amb $(v\comp u)\comp t=q$ i $a\comp t=p$. Queda veure que $u\comp t=r$: component amb $v$ i amb $b$,
\[
v\comp(u\comp t)=q=v\comp r,\qquad b\comp(u\comp t)=u'\comp a\comp t=u'\comp p=b\comp r,
\]
i com que el quadrat dret és un pullback, aquestes dues igualtats forcen $u\comp t=r$. Així $t$ és la factorització única per al quadrat esquerre.

En tots dos sentits, l'ingredient és la propietat universal del quadrat dret, que transfereix les factoritzacions entre el rectangle i el quadrat esquerre. Aquest \emph{lema d'enganxament} és una de les eines més usades per raonar amb pullbacks, i el retrobarem en tractar la reindexació entre subobjectes.
\end{solucio}

\begin{exercici}\label{pr-reindexacio-set}
Sigui $f\colon A\to B$ una aplicació i $y\colon Y\to B$ un objecte de $\cat{Set}/B$. Calcula explícitament la reindexació $f^{*}Y$ (observació~\ref{obs:reindexacio}) i comprova que la fibra de $f^{*}Y\to A$ sobre cada $a\in A$ està en bijecció canònica amb la fibra $y^{-1}(f(a))$.
\end{exercici}

\begin{solucio}
Pel càlcul del pullback a $\cat{Set}$ de l'exercici anterior,
\[
f^{*}Y=A\times_B Y=\{(a,z)\in A\times Y\mid f(a)=y(z)\},
\]
amb la projecció $p_1\colon(a,z)\mapsto a$ com a objecte de $\cat{Set}/A$. La fibra de $p_1$ sobre $a$ és
\[
p_1^{-1}(a)=\{(a,z)\mid y(z)=f(a)\},
\]
i l'aplicació $(a,z)\mapsto z$ en dona una bijecció canònica amb
\[
\{z\in Y\mid y(z)=f(a)\}=y^{-1}(f(a)).
\]
Així, si interpretem $y\colon Y\to B$ com la família de fibres $(Y_b)_{b\in B}$, amb $Y_b=y^{-1}(b)$, la reindexació al llarg de $f$ s'interpreta com la família $(Y_{f(a)})_{a\in A}$: cada índex $a$ rep la fibra de l'índex $f(a)$. És la \emph{substitució} al llarg de $f$, que reapareixerà com a operació sobre predicats en el capítol de subobjectes.
\end{solucio}

\begin{exercici}\label{pr-provt-limits-finits}
A la categoria prima $\cat{Prov}_T$, suposem que el sistema deductiu $T$ té una fórmula $\top$ i una conjunció $\land$ amb les regles habituals. Comprova que $\top$ és un objecte terminal i que $p\land q$ és un producte de $p$ i $q$. Dedueix que $\cat{Prov}_T$ té tots els límits finits.
\end{exercici}

\begin{solucio}
\textbf{Terminal.} Per a tota fórmula $r$ es té $r\vdash_T\top$, i com que $\cat{Prov}_T$ és prima, aquesta fletxa és única. Per tant $\top$ és terminal.

\textbf{Producte.} Les regles d'eliminació de la conjunció donen les projeccions
\[
p\land q\vdash_T p,
\qquad
p\land q\vdash_T q.
\]
Si $r\vdash_T p$ i $r\vdash_T q$, la regla d'introducció dona $r\vdash_T p\land q$, que és la fletxa mediadora; els triangles commuten automàticament i la fletxa és única, perquè la categoria és prima. Per tant $p\land q$, amb les dues eliminacions, és un producte de $p$ i $q$.

\textbf{Límits finits.} En una categoria prima, dues fletxes paral·leles coincideixen, de manera que l'equalitzador de $x\rightrightarrows y$ és $\id_x$ i sempre existeix. Pel teorema de caracterització dels límits finits (teorema~\ref{teo:limits-finits}), objecte terminal, productes binaris i equalitzadors donen tots els límits finits. És la versió resolta de la lectura lògica de l'exemple~\ref{ex:limits-preordre}.
\end{solucio}

\section{Colímits}

\begin{exercici}
Descriu el coproducte a $\cat{Set}$ i comprova'n la propietat universal.
\end{exercici}

\begin{solucio}
El coproducte de $A$ i $B$ a $\cat{Set}$ és la unió disjunta $A+B=(\{1\}\times A)\cup(\{2\}\times B)$, amb les injeccions $\iota_1(a)=(1,a)$ i $\iota_2(b)=(2,b)$. Donats $f\colon A\to X$ i $g\colon B\to X$, l'aplicació $[f,g]\colon A+B\to X$ definida per $[f,g](1,a)=f(a)$ i $[f,g](2,b)=g(b)$ és l'únic morfisme amb $[f,g]\comp\iota_1=f$ i $[f,g]\comp\iota_2=g$. Aquesta és la definició per casos, i és exactament la propietat universal del coproducte, dual de la del producte.
\end{solucio}

\begin{exercici}
Comprova que a $\cat{Set}$ el coequalitzador de $f,g\colon A\to B$ és el quocient $B/{\sim}$ per la relació d'equivalència generada per $f(a)\sim g(a)$.
\end{exercici}

\begin{solucio}
Sigui $\sim$ la menor relació d'equivalència amb $f(a)\sim g(a)$ per a tot $a$, i $q\colon B\to B/{\sim}$ la projecció al quocient. Aleshores $q\comp f=q\comp g$, perquè $f(a)$ i $g(a)$ tenen la mateixa classe. Si $s\colon B\to S$ compleix $s\comp f=s\comp g$, aleshores $s$ és constant sobre cada classe d'equivalència, de manera que factoritza de manera única com $s=r\comp q$ amb $r([b])=s(b)$ ben definida. Aquesta és la propietat universal del coequalitzador; el quocient hi encaixa exactament.
\end{solucio}

\section{Preservació i reflexió}

\begin{exercici}
Comprova que el functor oblit $U\colon\cat{Grp}\to\cat{Set}$ preserva els productes, però no preserva tots els coproductes.
\end{exercici}

\begin{solucio}
Per als productes: el grup producte $G\times H$ té com a conjunt subjacent el producte cartesià dels conjunts subjacents, i les projeccions són les aplicacions projecció; per tant $U(G\times H)=U(G)\times U(H)$ amb les projeccions correctes, i $U$ preserva el producte.

Per als coproductes: el coproducte a $\cat{Grp}$ és el \emph{producte lliure} $G*H$. N'hi ha prou d'un contraexemple, i el més senzill és prendre dos grups trivials $1=\{e\}$. El seu producte lliure és de nou trivial, $1*1\cong 1$, de manera que $U(1*1)$ té un sol element. En canvi, el coproducte a $\cat{Set}$ dels conjunts subjacents, $U(1)+U(1)$, és la unió disjunta de dos conjunts d'un element i en té dos. El cocon imatge, amb vèrtex $U(1*1)$, no pot ser, doncs, un coproducte a $\cat{Set}$, i $U$ no preserva aquest coproducte.

Aquesta asimetria anticipa el teorema del capítol següent: $U$ és un adjunt per la dreta, i els adjunts per la dreta preserven tots els límits; no hi ha cap afirmació anàloga que els obligui a preservar tots els colímits.
\end{solucio}

\begin{exercici}
Comprova que el functor representable $\Hom(A,-)\colon\cat{Set}\to\cat{Set}$ preserva productes, és a dir, $\Hom(A,X\times Y)\cong\Hom(A,X)\times\Hom(A,Y)$.
\end{exercici}

\begin{solucio}
Per la propietat universal del producte, l'aplicació $h\mapsto(\pi_1\comp h,\ \pi_2\comp h)$ és una bijecció
\[
\Hom(A,X\times Y)\cong\Hom(A,X)\times\Hom(A,Y),
\]
natural en $X$ i $Y$ (i també en $A$, si es considera la variació contravariant en aquesta variable). Aquesta bijecció és exactament la induïda pel con imatge $\bigl(\Hom(A,X\times Y),\Hom(A,\pi_1),\Hom(A,\pi_2)\bigr)$, de manera que aquest con és un producte a $\cat{Set}$ i $\Hom(A,-)$ preserva el producte binari. El mateix argument, amb la propietat universal general del límit, mostra que $\Hom(A,-)$ preserva tots els límits (proposició~\ref{prop:representables-limits}).
\end{solucio}

\chapter[Adjuncions]{Adjuncions}

Els exercicis d'aquest capítol treballen les adjuncions des dels tres punts de vista equivalents ---bijecció d'homconjunts, unitat/counitat, identitats triangulars--- i n'apliquen les conseqüències sobre límits. Convé identificar sempre qui és l'adjunt per l'esquerra i qui per la dreta.

\section{Adjuncions entre preordres}

\begin{exercici}
Comprova que al preordre $\cat{Prov}_T$ d'un sistema proposicional amb les regles usuals de $\wedge$ i $\to$, fixada una fórmula $A$, l'operació $B\mapsto A\wedge B$ és adjunta per l'esquerra de $B\mapsto A\to B$.
\end{exercici}

\begin{solucio}
Hem de veure que $A\wedge B\vdash_T C$ si i només si $B\vdash_T A\to C$. 

Suposem $A\wedge B\vdash_T C$. Aleshores, assumint $B$, volem deduir $A\to C$: assumim també $A$; tenim $A$ i $B$, per tant $A\wedge B$, d'on $C$; descarregant la hipòtesi $A$ obtenim $A\to C$. Així $B\vdash_T A\to C$.

Recíprocament, suposem $B\vdash_T A\to C$. Assumint $A\wedge B$, en tenim $A$ i $B$; de $B$ obtenim $A\to C$, i amb $A$, per modus ponens, $C$. Així $A\wedge B\vdash_T C$.

Les dues implicacions donen l'equivalència, és a dir, $(A\wedge-)\dashv(A\to-)$. Aquesta adjunció és la lectura categòrica de la regla de la implicació.
\end{solucio}

\begin{exercici}\label{pr-quantificadors-calcul}
Sigui $f\colon X\to Y$ l'aplicació de $X=\{1,2,3,4\}$ en $Y=\{a,b,c\}$ amb $f(1)=f(2)=a$, $f(3)=b$ i $f(4)=b$, i sigui $S=\{1,2,3\}$. Calcula $\exists_f(S)$ i $\forall_f(S)$ (exemple~\ref{ex:quantificadors-set}). Després, fent servir només la cadena $\exists_f\dashv f^{*}\dashv\forall_f$, justifica que $\exists_f$ preserva unions, $\forall_f$ preserva interseccions i $f^{*}$ preserva totes dues.
\end{exercici}

\begin{solucio}
\textbf{Càlcul.} La imatge és $\exists_f(S)=\{a,b\}$. Per a $\forall_f(S)$, mirem les fibres: $f^{-1}(\{a\})=\{1,2\}\subseteq S$, de manera que $a\in\forall_f(S)$; $f^{-1}(\{b\})=\{3,4\}\not\subseteq S$, perquè $4\notin S$; i $f^{-1}(\{c\})=\varnothing\subseteq S$. Per tant
\[
\forall_f(S)=\{a,c\}.
\]
Observem dos fenòmens: $c$ hi és perquè la seva fibra és buida (la quantificació universal sobre un conjunt buit és certa), i $b$ no hi és tot i que $3\in S$.

\textbf{Preservació.} En un preordre de parts, les unions són els coproductes (suprems) i les interseccions, els productes (ínfims). Pel teorema que els adjunts per l'esquerra preserven colímits i els adjunts per la dreta preserven límits:
\begin{itemize}
  \item $\exists_f$, adjunt per l'esquerra de $f^{*}$, preserva unions: $f[S\cup S']=f[S]\cup f[S']$;
  \item $\forall_f$, adjunt per la dreta de $f^{*}$, preserva interseccions;
  \item $f^{*}$, que és adjunt per la dreta de $\exists_f$ i adjunt per l'esquerra de $\forall_f$, preserva totes dues.
\end{itemize}
En canvi, $\exists_f$ no preserva en general les interseccions: amb l'aplicació de l'enunciat, $\exists_f(\{1\}\cap\{2\})=\varnothing$, però $\exists_f(\{1\})\cap\exists_f(\{2\})=\{a\}$. Lògicament: l'existencial commuta amb la disjunció, però no amb la conjunció.
\end{solucio}

\section{Adjuncions entre categories}

\begin{exercici}\label{pr-disc-ob}
Sigui $\mathrm{Ob}\colon\cat{Cat}\to\cat{Set}$ el functor que envia cada categoria petita al seu conjunt d'objectes, i $\mathrm{Disc}\colon\cat{Set}\to\cat{Cat}$ el que envia un conjunt $X$ a la categoria discreta amb objectes $X$. Comprova que $\mathrm{Disc}\dashv\mathrm{Ob}$ i descriu-ne la unitat i la counitat.
\end{exercici}

\begin{solucio}
Un functor $H\colon\mathrm{Disc}(X)\to\cat{D}$ ha d'enviar cada identitat a una identitat, i $\mathrm{Disc}(X)$ no té cap altre morfisme; queda, doncs, determinat per una aplicació $X\to\mathrm{Ob}(\cat{D})$ entre els objectes, i qualsevol aplicació dona un functor. Això és una bijecció
\[
\Hom_{\cat{Cat}}\bigl(\mathrm{Disc}(X),\cat{D}\bigr)\;\cong\;\Hom_{\cat{Set}}\bigl(X,\mathrm{Ob}(\cat{D})\bigr),
\]
natural en $X$ i en $\cat{D}$, perquè les dues bandes es componen de la mateixa manera sobre els objectes. Per tant $\mathrm{Disc}\dashv\mathrm{Ob}$.

La \emph{unitat} $\eta_X\colon X\to\mathrm{Ob}(\mathrm{Disc}(X))=X$ és la identitat: és un isomorfisme. La \emph{counitat} $\varepsilon_{\cat{D}}\colon\mathrm{Disc}(\mathrm{Ob}(\cat{D}))\to\cat{D}$ és el functor que inclou els objectes de $\cat{D}$ amb les seves identitats, i en general no és un isomorfisme, perquè oblida tots els morfismes no identitat.

\emph{Observació.} $\mathrm{Ob}$ té també un adjunt per la dreta: el functor que envia $X$ a la categoria \emph{indiscreta}, amb exactament un morfisme entre cada parella d'objectes. Per tant $\mathrm{Ob}$ preserva límits i colímits, d'acord amb el fet que els límits i colímits de categories petites es calculen, sobre els objectes, com a $\cat{Set}$.
\end{solucio}

\section{Unitat i counitat}

\begin{exercici}
Per a l'adjunció lliure-oblit $F\dashv U$ dels monoides, descriu la unitat $\eta_X\colon X\to U(X^*)$ i la counitat $\varepsilon_M\colon (U(M))^*\to M$, i comprova que la unitat expressa la propietat universal del monoide lliure.
\end{exercici}

\begin{solucio}
La unitat $\eta_X\colon X\to U(X^*)$ envia cada element $x$ a la paraula d'una sola lletra $x$. La counitat $\varepsilon_M\colon (U(M))^*\to M$ envia una paraula d'elements de $M$ al seu producte a $M$: $\varepsilon_M(m_1\cdots m_n)=m_1\cdots m_n$ (producte calculat a $M$).

La propietat universal de la unitat diu que per a tota aplicació $\varphi\colon X\to U(M)$ hi ha un únic homomorfisme $\bar\varphi\colon X^*\to M$ amb $U(\bar\varphi)\comp\eta_X=\varphi$. Això és exactament la definició del monoide lliure: un homomorfisme des de $X^*$ queda determinat, de manera única, pels valors sobre els generadors $X$.
\end{solucio}

\begin{exercici}
Comprova que, en una adjunció $F\dashv G$, el transposat de $\varphi\colon F(A)\to B$ és $G(\varphi)\comp\eta_A$, i el transposat invers de $\psi\colon A\to G(B)$ és $\varepsilon_B\comp F(\psi)$.
\end{exercici}

\begin{solucio}
Per naturalitat de la bijecció $\theta$ en la segona variable, aplicada al morfisme $\varphi\colon F(A)\to B$ i partint de $\id_{F(A)}$,
\[
\theta(\varphi)=\theta(\varphi\comp\id_{F(A)})=G(\varphi)\comp\theta(\id_{F(A)})=G(\varphi)\comp\eta_A,
\]
on hem usat que $\theta(\id_{F(A)})=\eta_A$ per definició de la unitat. Dualment, per naturalitat en la primera variable i la definició de la counitat $\varepsilon_B=\theta^{-1}(\id_{G(B)})$,
\[
\theta^{-1}(\psi)=\theta^{-1}(\id_{G(B)}\comp\psi)=\theta^{-1}(\id_{G(B)})\comp F(\psi)=\varepsilon_B\comp F(\psi).
\]
Aquestes fórmules permeten passar de la bijecció a la unitat/counitat i viceversa.
\end{solucio}

\section{Identitats triangulars}

\begin{exercici}\label{pr-diagonal-producte}
Sigui $\cat{C}$ una categoria amb productes binaris, i siguin $\Delta\colon\cat{C}\to\cat{C}\times\cat{C}$ el functor diagonal, $\Delta(A)=(A,A)$, i $\times\colon\cat{C}\times\cat{C}\to\cat{C}$ el functor producte, $(X,Y)\mapsto X\times Y$. Comprova que $\Delta\dashv\times$, descriu-ne la unitat i la counitat, i verifica directament les dues identitats triangulars.
\end{exercici}

\begin{solucio}
\textbf{L'adjunció.} Un morfisme $\Delta(A)\to(X,Y)$ a $\cat{C}\times\cat{C}$ és una parella $(f\colon A\to X,\ g\colon A\to Y)$, i la propietat universal del producte la posa en bijecció amb els morfismes $\langle f,g\rangle\colon A\to X\times Y$:
\[
\Hom_{\cat{C}\times\cat{C}}\bigl((A,A),(X,Y)\bigr)\;\cong\;\Hom_{\cat{C}}(A,X\times Y).
\]
La bijecció és natural en $A$ i en $(X,Y)$, perquè $\langle f,g\rangle\comp a=\langle f\comp a,g\comp a\rangle$ i $(u\times v)\comp\langle f,g\rangle=\langle u\comp f,v\comp g\rangle$. Per tant $\Delta\dashv\times$.

\textbf{Unitat i counitat.} La unitat és el transposat de $\id_{(A,A)}$, és a dir, el morfisme diagonal
\[
\eta_A=\delta_A=\langle\id_A,\id_A\rangle\colon A\to A\times A,
\]
i la counitat és el transposat invers de $\id_{X\times Y}$, la parella de projeccions
\[
\varepsilon_{(X,Y)}=(\pi_1,\pi_2)\colon(X\times Y,\ X\times Y)\to(X,Y).
\]

\textbf{Primera identitat.} $\varepsilon_{\Delta(A)}\comp\Delta(\eta_A)$ és la parella
\[
(\pi_1\comp\delta_A,\ \pi_2\comp\delta_A)=(\id_A,\id_A)=\id_{\Delta(A)}:
\]
copiar $A$ i projectar a cada còpia no canvia res.

\textbf{Segona identitat.} $(\pi_1\times\pi_2)\comp\delta_{X\times Y}$ és el morfisme $X\times Y\to X\times Y$ de components $\pi_1$ i $\pi_2$, és a dir,
\[
\langle\pi_1,\pi_2\rangle=\id_{X\times Y}:
\]
duplicar una parella i quedar-se amb la primera component de la primera còpia i la segona de la segona retorna la parella original.

Observem que la unitat $\delta_A$ no és un isomorfisme (a $\cat{Set}$, només ho és si $A$ té com a màxim un element): les identitats triangulars no demanen que ho sigui.
\end{solucio}

\begin{exercici}\label{pr-free-u-triangulars}
Per a l'adjunció $\Free\dashv U$ entre $\cat{Graph}$ i $\cat{Cat}$ (exemple~\ref{ex:free-u-adjuncio}), comprova directament les dues identitats triangulars
\[
\varepsilon_{\Free(G)}\comp\Free(\eta_G)=\id_{\Free(G)},
\qquad
U(\varepsilon_{\cat{C}})\comp\eta_{U(\cat{C})}=\id_{U(\cat{C})},
\]
i explica què diu cadascuna en termes de camins.
\end{exercici}

\begin{solucio}
\textbf{Primera identitat.} El functor $\Free(\eta_G)\colon\Free(G)\to\Free(U(\Free(G)))$ envia un camí $(a_1,\dots,a_n)$ de $G$ al camí formal $\bigl((a_1),\dots,(a_n)\bigr)$, les arestes del qual són camins de longitud~$1$ de $\Free(G)$. La counitat $\varepsilon_{\Free(G)}$ avalua aquest camí formal component les seves arestes a $\Free(G)$, és a dir, concatenant:
\[
\varepsilon_{\Free(G)}\Bigl(\bigl((a_1),\dots,(a_n)\bigr)\Bigr)=(a_n)\comp\cdots\comp(a_1)=(a_1,\dots,a_n).
\]
El camí buit va al camí buit. Per tant, la composta és la identitat de $\Free(G)$: \emph{trencar un camí en arestes i tornar-lo a compondre no canvia res}.

\textbf{Segona identitat.} La unitat $\eta_{U(\cat{C})}$ envia una fletxa $f$ de $\cat{C}$, vista com a aresta del graf $U(\cat{C})$, al camí $(f)$ de longitud~$1$, i cada objecte a ell mateix. El morfisme de grafs $U(\varepsilon_{\cat{C}})$ avalua aquest camí, que té una sola fletxa, i en dona $f$. Per tant, la composta és la identitat de $U(\cat{C})$: \emph{avaluar el camí d'una sola fletxa retorna la fletxa}.

Les dues identitats expressen, doncs, que la inserció dels generadors i l'avaluació de camins són inverses en el sentit precís que demanen les identitats triangulars, sense que $\eta$ ni $\varepsilon$ siguin isomorfismes: $\Free(U(\cat{C}))$ té en general molts més morfismes que $\cat{C}$, perquè hi ha camins diferents amb la mateixa composta.
\end{solucio}

\section{Els adjunts i els límits}

\begin{exercici}\label{pr-rapl-mon}
Comprova, fent servir que els adjunts per la dreta preserven límits, que el functor oblit $U\colon\cat{Mon}\to\cat{Set}$ preserva productes. Comprova també que $U$ no preserva tots els coproductes, i explica què en dedueixes.
\end{exercici}

\begin{solucio}
El functor monoide lliure $F\colon\cat{Set}\to\cat{Mon}$, $X\mapsto X^*$, és adjunt per l'esquerra de $U$ (exemple del monoide lliure de la Teoria). Per tant $U$ és un adjunt per la dreta, i pel teorema \enquote{els adjunts per la dreta preserven límits}, $U$ preserva tots els límits petits, en particular els productes: el conjunt subjacent d'un producte de monoides és el producte dels conjunts subjacents.

El teorema, però, no diu res dels colímits, i de fet $U$ no preserva tots els coproductes. Sigui $1=\{e\}$ el monoide trivial. És un objecte inicial de $\cat{Mon}$, perquè l'únic homomorfisme $1\to M$ envia $e$ al neutre, i el coproducte d'un objecte amb un objecte inicial és l'objecte mateix; per tant, el coproducte $1+1$ a $\cat{Mon}$ és $1$, amb un sol element. En canvi, el coproducte a $\cat{Set}$ dels conjunts subjacents, $U(1)+U(1)$, en té dos, i el cocon imatge no és un coproducte. Com a conseqüència, $U$ no pot ser un adjunt per l'esquerra, perquè els adjunts per l'esquerra preserven colímits.

\emph{Ampliació, per a qui conegui els grups.} El mateix passa amb $U\colon\cat{Grp}\to\cat{Set}$. El seu adjunt per l'esquerra és el functor \emph{grup lliure}, que envia un conjunt $X$ al grup de les paraules reduïdes en les lletres de $X$ i els seus inversos formals; per tant, $U$ preserva tots els límits. El grup trivial també és inicial a $\cat{Grp}$, el coproducte (el \emph{producte lliure}) de dos grups trivials és trivial, $1*1\cong 1$, i $U$ tampoc no preserva aquest coproducte.
\end{solucio}

\begin{exercici}
Comprova que si un functor $G$ \emph{no} preserva algun límit petit, aleshores $G$ no pot tenir cap adjunt per l'esquerra.
\end{exercici}

\begin{solucio}
És el contrarecíproc del teorema. Si $G$ tingués un adjunt per l'esquerra $F$, amb $F\dashv G$, aleshores $G$ seria un adjunt per la dreta i, pel teorema, preservaria tots els límits petits que existeixen. Com que per hipòtesi $G$ no en preserva algun, no pot tenir adjunt per l'esquerra. Aquest és un criteri negatiu molt pràctic: per exemple, si un functor oblit no preserva algun coproducte, això impedeix que sigui adjunt per l'esquerra, però no impedeix que sigui adjunt per la dreta.
\end{solucio}

\section{Exponencials}

\begin{exercici}
Comprova la bijecció
\[
\Hom_{\cat{Set}}(A\times B,C)\cong\Hom_{\cat{Set}}(A,C^B)
\]
i identifica-la amb la currificació.
\end{exercici}

\begin{solucio}
A $\cat{Set}$, $C^B$ és el conjunt de les aplicacions $B\to C$. Donada $f\colon A\times B\to C$, definim $\tilde f\colon A\to C^B$ per $\tilde f(a)=(b\mapsto f(a,b))$; és a dir, $\tilde f(a)$ és la funció que fixa la primera coordenada en $a$. Recíprocament, donada $g\colon A\to C^B$, definim $\hat g\colon A\times B\to C$ per $\hat g(a,b)=g(a)(b)$. Aquestes assignacions són inverses l'una de l'altra i naturals en totes les variables. La correspondència $f\mapsto\tilde f$ és exactament la \textbf{currificació}: transformar una funció de dues variables en una funció d'una variable que retorna funcions. Aquesta bijecció expressa l'adjunció $-\times B\dashv(-)^B$.
\end{solucio}

\chapter[Cartesianes tancades]{\texorpdfstring{Categories cartesianes\\ tancades}{Categories cartesianes tancades}}

Els exercicis treballen tres aspectes del tancament cartesià: productes i diagonals, la propietat universal de l'exponencial i la interpretació de tipus i lògica. Convé tenir sempre present la bijecció
\[
\Hom(A\times B,C)\cong\Hom(A,C^{B}).
\]

\section{Productes i diagonals}

\begin{exercici}
En una categoria cartesiana, comprova que la diagonal $\delta_A\colon A\to A\times A$ satisfà
\[
\pi_1\comp\delta_A=\pi_2\comp\delta_A=\id_A
\]
i que és l'únic morfisme amb aquesta propietat.
\end{exercici}

\begin{solucio}
Per la propietat universal del producte $A\times A$, un morfisme $A\to A\times A$ està determinat per les seves dues composicions amb les projeccions. Definim $\delta_A$ com l'únic morfisme amb $\pi_1\comp\delta_A=\id_A$ i $\pi_2\comp\delta_A=\id_A$; l'existència i la unicitat són exactament el contingut de la propietat universal aplicada a la parella $(\id_A,\id_A)$. Qualsevol altre morfisme $d\colon A\to A\times A$ amb $\pi_1\comp d=\pi_2\comp d=\id_A$ compleix la mateixa condició universal i, per unicitat, $d=\delta_A$.
\end{solucio}

\section{La propietat universal de l'exponencial}

\begin{exercici}
A $\cat{Set}$, descriu explícitament l'avaluació, la currificació i la descurrificació per a l'exponencial $C^{B}$, i verifica l'equació $\ev\comp(\lambda f\times\id_B)=f$.
\end{exercici}

\begin{solucio}
A $\cat{Set}$, $C^{B}$ és el conjunt d'aplicacions $B\to C$. L'avaluació és $\ev\colon C^{B}\times B\to C$, $\ev(g,b)=g(b)$. Donada $f\colon A\times B\to C$, la currificació és $\lambda f\colon A\to C^{B}$ amb $\lambda f(a)=(b\mapsto f(a,b))$. Comprovem l'equació: per a $(a,b)\in A\times B$,
\[
\bigl(\ev\comp(\lambda f\times\id_B)\bigr)(a,b)=\ev\bigl(\lambda f(a),b\bigr)=\lambda f(a)(b)=f(a,b).
\]
Per tant $\ev\comp(\lambda f\times\id_B)=f$. La descurrificació de $g\colon A\to C^{B}$ és $\widetilde g(a,b)=g(a)(b)$, i és inversa de la currificació.
\end{solucio}

\begin{exercici}
Demostra que en qualsevol CCC hi ha un isomorfisme natural $C^{1}\cong C$ i que $1^{B}\cong 1$.
\end{exercici}

\begin{solucio}
Per a $C^{1}$: per la propietat universal, $\Hom(A,C^{1})\cong\Hom(A\times 1,C)$. Com que $A\times 1\cong A$ (l'objecte terminal és neutre per al producte), $\Hom(A\times 1,C)\cong\Hom(A,C)$, natural en $A$. Pel lema de Yoneda, $C^{1}\cong C$.

Per a $1^{B}$: $\Hom(A,1^{B})\cong\Hom(A\times B,1)$, que té exactament un element per ser $1$ terminal. Per tant $\Hom(A,1^{B})$ és un conjunt d'un punt per a tot $A$, natural en $A$, i pel lema de Yoneda $1^{B}\cong 1$. Aquestes són les versions categòriques de les identitats $c^{1}=c$ i $1^{b}=1$.
\end{solucio}

\begin{exercici}
Comprova la \enquote{llei exponencial} $C^{A\times B}\cong(C^{B})^{A}$ en una CCC.
\end{exercici}

\begin{solucio}
Fem servir la propietat universal repetidament i el lema de Yoneda. Per a qualsevol $X$,
\[
\Hom(X,C^{A\times B})\cong\Hom(X\times(A\times B),C)\cong\Hom((X\times A)\times B,C),
\]
usant l'associativitat del producte. Aplicant ara l'adjunció respecte de $B$ i després respecte de $A$,
\[
\Hom((X\times A)\times B,C)\cong\Hom(X\times A,C^{B})\cong\Hom(X,(C^{B})^{A}).
\]
Totes les bijeccions són naturals en $X$; pel lema de Yoneda, $C^{A\times B}\cong(C^{B})^{A}$. És la llei $c^{a\cdot b}=(c^{b})^{a}$, i correspon, en el càlcul lambda, al fet que una funció de dos arguments és una funció que retorna una funció (currificació iterada).
\end{solucio}

\begin{exercici}\label{pr-distributivitat}
Sigui $\cat{C}$ una CCC amb coproductes binaris i objecte inicial $0$. Demostra que el producte es distribueix sobre el coproducte:
\[
A\times(B+C)\;\cong\;(A\times B)+(A\times C),
\qquad
A\times 0\;\cong\;0.
\]
Interpreta el resultat a $\cat{Set}$ i en una àlgebra de Heyting.
\end{exercici}

\begin{solucio}
Pel tancament cartesià, $-\times A\dashv(-)^{A}$, i el producte és commutatiu llevat d'isomorfisme natural; per tant $A\times-$ també és un adjunt per l'esquerra. Els adjunts per l'esquerra preserven colímits, en particular els coproductes binaris i l'objecte inicial, que és el colímit del diagrama buit. Explícitament, el morfisme canònic
\[
[\id_A\times\iota_1,\ \id_A\times\iota_2]\colon(A\times B)+(A\times C)\longrightarrow A\times(B+C)
\]
és un isomorfisme, i $A\times 0$ és inicial, de manera que $A\times 0\cong 0$.

\emph{A $\cat{Set}$}, la primera bijecció envia una parella $(a,x)$ amb $x$ a la component $B$ o a la component $C$ de la unió disjunta al lloc corresponent; la segona diu que $A\times\varnothing=\varnothing$.

\emph{En una àlgebra de Heyting}, vista com a preordre cartesià tancat amb suprems finits, el coproducte és el suprem i l'inicial és $\bot$. Obtenim
\[
a\wedge(b\vee c)=(a\wedge b)\vee(a\wedge c),
\qquad
a\wedge\bot=\bot:
\]
tota àlgebra de Heyting és un reticle distributiu, i la distributivitat és conseqüència de l'existència de la implicació. Lògicament, la conjunció es distribueix sobre la disjunció.
\end{solucio}

\section{Heyting i lògica}

\begin{exercici}
Comprova que en una àlgebra de Heyting l'exponencial $b\Rightarrow c$ satisfà $a\wedge b\leq c\Leftrightarrow a\leq(b\Rightarrow c)$, i deriva-hi que $b\wedge(b\Rightarrow c)\leq c$ (\emph{modus ponens}).
\end{exercici}

\begin{solucio}
La condició $a\wedge b\leq c\Leftrightarrow a\leq(b\Rightarrow c)$ és, precisament, l'adjunció $-\wedge b\dashv(b\Rightarrow-)$ escrita en el llenguatge dels preordres: és la definició d'exponencial en un preordre cartesià, on $\Hom$ té com a molt un element i la bijecció d'homconjunts es redueix a una equivalència de desigualtats. Per obtenir el \emph{modus ponens}, prenem $a=(b\Rightarrow c)$. Aleshores el membre dret $a\leq(b\Rightarrow c)$ es compleix trivialment (és $b\Rightarrow c\leq b\Rightarrow c$), de manera que el membre esquerre també: $(b\Rightarrow c)\wedge b\leq c$. Aquesta és la counitat de l'adjunció, i és la forma algebraica de la regla que de $b$ i $b\to c$ se'n segueix $c$.
\end{solucio}

\begin{exercici}
En una àlgebra de Heyting, demostra que $a\leq(b\Rightarrow a)$ per a tot $a,b$, i interpreta-ho lògicament.
\end{exercici}

\begin{solucio}
Per l'adjunció, $a\leq(b\Rightarrow a)$ equival a $a\wedge b\leq a$, que és cert perquè $a\wedge b$ és un ínfim i, per tant, $a\wedge b\leq a$. Lògicament, correspon a l'axioma $a\to(b\to a)$: una proposició verdadera se segueix de qualsevol hipòtesi. És l'esquelet del combinador $\mathsf{K}$ del càlcul lambda, la funció constant.
\end{solucio}

\section{Semàntica de termes}

\begin{exercici}
Interpreta el terme $x:\sigma,\,y:\tau\vdash\langle y,x\rangle:\tau\times\sigma$ com a morfisme en una CCC i identifica'l amb el morfisme de commutació del producte.
\end{exercici}

\begin{solucio}
El context $x:\sigma,y:\tau$ s'interpreta com l'objecte $\llbracket\sigma\rrbracket\times\llbracket\tau\rrbracket$. Les variables $x$ i $y$ són les projeccions $\pi_1$ i $\pi_2$. La parella $\langle y,x\rangle$ s'interpreta amb l'aparellament de les interpretacions de $y$ i $x$, és a dir,
\[
\llbracket\langle y,x\rangle\rrbracket=\langle\pi_2,\pi_1\rangle\colon\llbracket\sigma\rrbracket\times\llbracket\tau\rrbracket\to\llbracket\tau\rrbracket\times\llbracket\sigma\rrbracket.
\]
Aquest és exactament el morfisme de \textbf{commutació} (o \emph{swap}) del producte, que intercanvia els dos factors. La seva inversa és $\langle\pi_2,\pi_1\rangle$ en sentit contrari, cosa que expressa que $\sigma\times\tau\cong\tau\times\sigma$.
\end{solucio}

\begin{exercici}
Interpreta l'abstracció $x:\sigma\vdash(\lambda y{:}\tau.\,x):\tau\to\sigma$ i comprova que és la currificació d'una projecció.
\end{exercici}

\begin{solucio}
El cos del terme, $x:\sigma,\,y:\tau\vdash x:\sigma$, s'interpreta amb la projecció $\pi_1\colon\llbracket\sigma\rrbracket\times\llbracket\tau\rrbracket\to\llbracket\sigma\rrbracket$. L'abstracció respecte de $y$ es tradueix per la currificació:
\[
\llbracket\lambda y.\,x\rrbracket=\lambda(\pi_1)\colon\llbracket\sigma\rrbracket\to\llbracket\sigma\rrbracket^{\llbracket\tau\rrbracket}.
\]
És a dir, la funció que envia cada $x$ a la funció constant de valor $x$: el combinador $\mathsf{K}$, la versió proposicional del qual és la desigualtat $a\leq(b\Rightarrow a)$ de l'exercici de Heyting anterior.
\end{solucio}

\chapter[Subobjectes, imatges i factoritzacions]{\texorpdfstring{Subobjectes, imatges\\ i factoritzacions}{Subobjectes, imatges i factoritzacions}}

Els exercicis d'aquest capítol fixen les tres idees centrals: els subobjectes com a classes de monomorfismes, l'ordre $\Sub(A)$ com a àlgebra de predicats, i la reindexació $f^{*}$ com a substitució. Molts arguments es fan \enquote{sense elements}, usant només propietats universals; convé acostumar-s'hi, perquè és el que permet traslladar-los a categories arbitràries.

\section{Subobjectes i el seu ordre}

\begin{exercici}
Comprova que a $\cat{Set}$ dos monomorfismes $m\colon S\rightarrowtail A$ i $n\colon T\rightarrowtail A$ són equivalents si i només si tenen la mateixa imatge, i dedueix que $\Sub(A)\cong\mathcal{P}(A)$.
\end{exercici}

\begin{solucio}
Si $m\sim n$, hi ha un isomorfisme $\varphi\colon S\to T$ amb $n\comp\varphi=m$; aleshores $m(S)=n(\varphi(S))=n(T)$, de manera que tenen la mateixa imatge. Recíprocament, suposem $m(S)=n(T)=:I$. Siguin $\bar m\colon S\to I$ i $\bar n\colon T\to I$ les \emph{corestriccions} de $m$ i $n$ a $I$, és a dir, les mateixes aplicacions amb codomini $I$: $\bar m(s)=m(s)$ i $\bar n(t)=n(t)$. Com que $m$ i $n$ són injectives i tenen imatge $I$, $\bar m$ i $\bar n$ són bijeccions. La composició $\varphi=\bar n^{-1}\comp\bar m\colon S\to T$ és, doncs, una bijecció, i compleix $n\comp\varphi=m$, perquè per a tot $s\in S$, $n(\varphi(s))=\bar n(\bar n^{-1}(\bar m(s)))=m(s)$. Així $m\sim n$. Per tant la imatge determina la classe, i l'assignació $[m]\mapsto m(S)$ és una bijecció entre $\Sub(A)$ i el conjunt de subconjunts $\mathcal{P}(A)$; a més respecta l'ordre, ja que $[m]\leq[n]$ equival a $m(S)\subseteq n(T)$.
\end{solucio}

\begin{exercici}
Demostra que en qualsevol categoria l'objecte $\id_A\colon A\to A$ és el subobjecte \textbf{màxim} de $\Sub(A)$.
\end{exercici}

\begin{solucio}
La identitat és un monomorfisme (és un isomorfisme). Per a qualsevol altre subobjecte $m\colon S\rightarrowtail A$, el mateix $m$ és una factorització de $m$ a través de $\id_A$, ja que $\id_A\comp m=m$. Per tant $[m]\leq[\id_A]$ per a tot $m$, cosa que fa de $[\id_A]$ el màxim. Lògicament, és el predicat \enquote{sempre cert}: tot element de $A$ el satisfà.
\end{solucio}

\begin{exercici}
Comprova que si $m\colon S\rightarrowtail A$ és un subobjecte i $\varphi\colon S'\to S$ és un isomorfisme, aleshores $m\comp\varphi$ representa el mateix subobjecte. Per què això justifica treballar amb classes i no amb monos concrets?
\end{exercici}

\begin{solucio}
Per definició de la relació d'equivalència, $m\comp\varphi\sim m$ mitjançant l'isomorfisme $\varphi$: $m\comp\varphi=m\comp\varphi$ amb $m\comp\varphi$ factoritzant. Formalment, $[m\comp\varphi]=[m]$ perquè $\varphi$ és un isomorfisme del domini que fa commutar el triangle. Això justifica treballar amb classes: la teoria dels subobjectes ha de ser invariant per reparametritzacions del domini, ja que \enquote{el subconjunt} no depèn de com l'indexem. Prendre classes elimina aquesta ambigüitat sense perdre informació geomètrica.
\end{solucio}

\section{Reindexació}

\begin{exercici}
Sigui $f\colon A\to B$ a $\cat{Set}$ i $S\subseteq B$. Comprova que $f^{*}(S)=f^{-1}(S)$ i que $f^{*}$ preserva interseccions: $f^{*}(S\cap S')=f^{*}(S)\cap f^{*}(S')$.
\end{exercici}

\begin{solucio}
El pullback de la injecció $S\hookrightarrow B$ al llarg de $f$ és $\{(a,s)\in A\times S\mid f(a)=s\}$, que s'identifica amb $\{a\in A\mid f(a)\in S\}=f^{-1}(S)$, amb la projecció cap a $A$ com a injecció. Per a les interseccions, $x\in f^{-1}(S\cap S')$ si i només si $f(x)\in S\cap S'$, és a dir, $f(x)\in S$ i $f(x)\in S'$, això és $x\in f^{-1}(S)\cap f^{-1}(S')$. La reindexació preserva, doncs, la conjunció; és una manifestació que la substitució commuta amb $\wedge$.
\end{solucio}


\begin{exercici}
Comprova que si $f\colon A\to B$ és un epimorfisme escindit (té secció $s$ amb $f\comp s=\id_B$), aleshores $f^{*}$ és injectiva sobre subobjectes.
\end{exercici}

\begin{solucio}
Si $f\comp s=\id_B$, aplicant la functorialitat de la reindexació (proposició~\ref{prop:functor-sub}), $s^{*}\comp f^{*}=(f\comp s)^{*}=(\id_B)^{*}=\id_{\Sub(B)}$. Per tant $f^{*}$ té inversa per l'esquerra $s^{*}$, cosa que la fa injectiva: si $f^{*}[m]=f^{*}[n]$, aplicant $s^{*}$ obtenim $[m]=[n]$. Com que $f$ admet una secció, els subobjectes de $B$ es poden recuperar després de reindexar-los al llarg de $f$; per això $f^{*}$ és injectiva. (No vol dir que $f$ no identifiqui elements: una epi escindida pot identificar-ne molts.)
\end{solucio}

\section{Interseccions i lògica}

\begin{exercici}
Verifica amb detall que el pullback de dos subobjectes $m\colon S\rightarrowtail A$ i $n\colon T\rightarrowtail A$ és el seu ínfim a $\Sub(A)$, escrivint les dues direccions de la propietat universal.
\end{exercici}

\begin{solucio}
Sigui $P=S\times_A T$ el pullback, amb projeccions $p\colon P\to S$ i $q\colon P\to T$ i injecció $m\comp p=n\comp q\colon P\rightarrowtail A$ (mono per ser composició de monos). \emph{Fita inferior}: $[P]\leq[m]$ via $p$, ja que $m\comp p$ és la injecció i factoritza per $m$; anàlogament $[P]\leq[n]$ via $q$. \emph{Màxima fita inferior}: sigui $k\colon R\rightarrowtail A$ amb $[k]\leq[m]$ i $[k]\leq[n]$, amb factoritzacions $a\colon R\to S$ ($m\comp a=k$) i $b\colon R\to T$ ($n\comp b=k$). Aleshores $m\comp a=k=n\comp b$, de manera que la parella $(a,b)$ és compatible amb el pullback i indueix un únic $u\colon R\to P$ amb $p\comp u=a$ i $q\comp u=b$. Aquest $u$ factoritza $k$ a través de la injecció de $P$: $(m\comp p)\comp u=m\comp a=k$. Per tant $[k]\leq[P]$, i $[P]$ és l'ínfim $[m]\wedge[n]$.
\end{solucio}

\section{Imatges i factoritzacions}

\begin{exercici}
Comprova que a $\cat{Set}$ tota aplicació $f\colon A\to B$ té factorització imatge amb $I=f(A)$, i identifica l'epimorfisme regular i el monomorfisme.
\end{exercici}

\begin{solucio}
Prenem $I=f(A)\subseteq B$. L'aplicació $e\colon A\to I$, $e(a)=f(a)$, és exhaustiva, i tota aplicació exhaustiva a $\cat{Set}$ és un epimorfisme regular: és el coequalitzador de la seva parella nucli $\{(a,a')\mid f(a)=f(a')\}$. La injecció $m\colon I\hookrightarrow B$ és injectiva, per tant un monomorfisme. Aleshores $f=m\comp e$ és la factorització imatge, i $\Imatge(f)=[m]$ és el subconjunt $f(A)$. Recuperem la noció usual d'imatge d'una aplicació.
\end{solucio}

\begin{exercici}
Demostra que si $f=m\comp e$ és una factorització imatge i $f$ és ell mateix un monomorfisme, aleshores $e$ és un isomorfisme.
\end{exercici}

\begin{solucio}
Si $f=m\comp e$ és mono i $m$ és mono, aleshores $e$ és mono: de $e\comp u=e\comp v$ se segueix $m\comp e\comp u=m\comp e\comp v$, és a dir $f\comp u=f\comp v$, d'on $u=v$. Però $e$ també és un epimorfisme regular, en particular un epimorfisme. Un morfisme que és alhora mono i epi regular és un isomorfisme: essent el coequalitzador d'una parella que ja iguala (perquè $e$ mono força la parella nucli a ser la diagonal), $e$ coequalitza la parella trivial i és, doncs, invertible. Per tant $\Imatge(f)=[f]$: la imatge d'un mono és ell mateix.
\end{solucio}

\begin{exercici}
Comprova que la imatge és el menor subobjecte pel qual es factoritza $f$: si $f=n\comp g$ amb $n\colon T\rightarrowtail B$ mono, aleshores $\Imatge(f)\leq[n]$.
\end{exercici}

\begin{solucio}
Sigui $f=m\comp e$ la factorització imatge i $f=n\comp g$ una factorització qualsevol a través del mono $n$. Aleshores $n\comp g=m\comp e$. Com que $e$ és el coequalitzador de la parella nucli de $f$ i $n$ és mono, $g$ iguala les dues branques de la parella nucli, de manera que es factoritza per $e$: hi ha $h$ amb $h\comp e=g$. Aleshores $n\comp h\comp e=n\comp g=f=m\comp e$, i com que $e$ és epi, $n\comp h=m$. Aquesta igualtat és precisament una factorització de $m$ a través de $n$, és a dir $\Imatge(f)=[m]\leq[n]$.
\end{solucio}

\section{Estabilitat i regularitat}

\begin{exercici}
A $\cat{Set}$, comprova directament que les imatges són estables per pullback: si es reindexà $f\colon A\to B$ al llarg de $g\colon B'\to B$, aleshores $\Imatge(f')=g^{-1}(\Imatge(f))$.
\end{exercici}

\begin{solucio}
El pullback de $f$ al llarg de $g$ és $A'=\{(b',a)\mid g(b')=f(a)\}$ amb $f'(b',a)=b'$. La seva imatge és $\{b'\in B'\mid\exists a,\ g(b')=f(a)\}=\{b'\mid g(b')\in f(A)\}=g^{-1}(f(A))=g^{-1}(\Imatge(f))$. Per tant $\Imatge(f')=g^{*}(\Imatge(f))$, que és l'enunciat d'estabilitat. Aquesta comprovació explícita és el model de la propietat abstracta que defineix les categories regulars.
\end{solucio}

\begin{exercici}
Demostra que en una categoria regular els epimorfismes regulars són tancats per composició: si $f\colon A\to B$ i $g\colon B\to C$ són epimorfismes regulars, aleshores $g\comp f$ també ho és.
\end{exercici}

\begin{solucio}
Siguin $f\colon A\to B$ i $g\colon B\to C$ epimorfismes regulars, i factoritzem la composició com a imatge, $g\comp f=m\comp e$, amb $e\colon A\to I$ epi regular i $m\colon I\rightarrowtail C$ mono. Sigui $p_1,p_2\colon R\rightrightarrows A$ la parella nucli de $f$. Com que $f\comp p_1=f\comp p_2$,
\[
m\comp e\comp p_1=g\comp f\comp p_1=g\comp f\comp p_2=m\comp e\comp p_2,
\]
i, com que $m$ és mono, $e\comp p_1=e\comp p_2$. Pel lema~\ref{lema:epi-regular-nucli}, $f$ és el coequalitzador de la seva parella nucli; per tant, existeix $k\colon B\to I$ amb $e=k\comp f$. Aleshores
\[
m\comp k\comp f=m\comp e=g\comp f,
\]
i com que $f$ és epi, $m\comp k=g$: el morfisme $g$ es factoritza a través de $m$. Però $g=\id_C\comp g$ és una factorització imatge de $g$, de manera que $\Imatge(g)=[\id_C]$, i per la minimalitat de la imatge (proposició~\ref{prop:imatge-unica}), $[\id_C]\leq[m]$. Com que $[\id_C]$ és el màxim de $\Sub(C)$, $[m]=[\id_C]$ i $m$ és un isomorfisme. Per tant $g\comp f=m\comp e$ és un epi regular seguit d'un isomorfisme, que torna a ser un epi regular.

A $\cat{Set}$, tot plegat es redueix al fet que la composició d'aplicacions exhaustives és exhaustiva.
\end{solucio}

\section{L'existencial}

\begin{exercici}
Verifica les dues direccions de l'adjunció $\exists_f\dashv f^{*}$ a $\cat{Set}$, on $\exists_f(S)=f(S)$ i $f^{*}(T)=f^{-1}(T)$.
\end{exercici}

\begin{solucio}
Siguin $S\subseteq A$ i $T\subseteq B$. Cal veure $f(S)\subseteq T\Leftrightarrow S\subseteq f^{-1}(T)$.

\emph{($\Rightarrow$)} Si $f(S)\subseteq T$ i $a\in S$, aleshores $f(a)\in f(S)\subseteq T$, així que $a\in f^{-1}(T)$; per tant $S\subseteq f^{-1}(T)$.

\emph{($\Leftarrow$)} Si $S\subseteq f^{-1}(T)$ i $b\in f(S)$, aleshores $b=f(a)$ amb $a\in S\subseteq f^{-1}(T)$, d'on $b=f(a)\in T$; per tant $f(S)\subseteq T$.

Aquesta és la forma conjuntista de l'adjunció imatge directa\,--\,antiimatge, i és exactament la regla d'introducció/eliminació de l'existencial quan $f$ és una projecció.
\end{solucio}

\begin{exercici}
Sigui $\pi\colon X\times Y\to X$ la projecció. Descriu $\exists_\pi[s]$ i $\forall_\pi[s]$ per a un subobjecte $[s]\subseteq X\times Y$ a $\cat{Set}$, i relaciona-ho amb els quantificadors.
\end{exercici}

\begin{solucio}
Sigui $S\subseteq X\times Y$ el subconjunt corresponent a $[s]$. Aleshores
\[
\begin{aligned}
\exists_\pi(S)&=\{x\in X\mid \exists y\in Y:\ (x,y)\in S\},\\
\forall_\pi(S)&=\{x\in X\mid \forall y\in Y:\ (x,y)\in S\}.
\end{aligned}
\]
El primer és la imatge de $S$ per $\pi$ (l'existencial); el segon, el conjunt de fibres senceres contingudes en $S$ (l'universal). Són els casos particulars, per a una projecció, de la cadena $\exists_f\dashv f^{*}\dashv\forall_f$ de l'exemple~\ref{ex:quantificadors-set} del capítol d'adjuncions: l'adjunció $\exists_\pi\dashv\pi^{*}$ es comprova com a l'exercici anterior, i $\pi^{*}\dashv\forall_\pi$ amb un argument anàleg, no per dualitat. Així, en una categoria regular tenim l'existencial; l'universal requereix, a més, adjunts per la dreta, que apareixen en contextos més rics (categories de Heyting).
\end{solucio}


\chapter[Classificadors de subobjectes i topos]{\texorpdfstring{Classificadors de subobjectes i introducció als topos}{Classificadors de subobjectes i introducció als topos}}

Els exercicis d'aquest capítol són de familiarització amb el classificador de subobjectes i la definició de topos. No entren en la lògica interna, que pertany a l'estudi posterior.

\section{El classificador}

\begin{exercici}
Comprova amb detall que a $\cat{Set}$ l'objecte $\Omega=\{0,1\}$ amb $\mathsf{true}\colon 1\to\{0,1\}$, $\mathsf{true}(\ast)=1$, és un classificador de subobjectes.
\end{exercici}

\begin{solucio}
Sigui $U\subseteq E$ un subconjunt, amb injecció $m\colon U\hookrightarrow E$. Definim $\chi_U\colon E\to\{0,1\}$ per $\chi_U(x)=1$ si $x\in U$ i $\chi_U(x)=0$ altrament. El pullback de $\mathsf{true}\colon 1\to\{0,1\}$ al llarg de $\chi_U$ és $\{x\in E\mid\chi_U(x)=1\}=U$, amb la injecció com a projecció; per tant el quadrat és cartesià. Per a la unicitat: si $\chi\colon E\to\{0,1\}$ fa el quadrat cartesià, aleshores $\chi^{-1}(1)=U$, cosa que força $\chi=\chi_U$. Així $\{0,1\}$ classifica els subobjectes de qualsevol conjunt.
\end{solucio}

\begin{exercici}\label{pr-ordre-pertinenca}
En una categoria amb classificador de subobjectes, siguin $m\colon U\rightarrowtail E$ i $n\colon V\rightarrowtail E$ dos subobjectes, amb morfismes característics $\chi_m$ i $\chi_n$. Demostra que
\[
[m]\leq[n]
\quad\Longleftrightarrow\quad
\chi_n\comp m=\mathsf{true}\comp\,!_U,
\]
és a dir, que $U$ està contingut en $V$ si i només si tots els \enquote{elements} de $U$ pertanyen a $V$. Dedueix que $\chi_m=\chi_n$ si i només si $[m]=[n]$.
\end{exercici}

\begin{solucio}
Per definició de l'ordre, $[m]\leq[n]$ vol dir que $m$ factoritza a través de $n$. Aplicant la caracterització de la pertinença per elements generalitzats de la Teoria a l'element generalitzat $m\colon U\to E$, això passa si i només si $\chi_n\comp m=\mathsf{true}\comp\,!_U$.

Per a la segona part, si $[m]=[n]$, els dos subobjectes són el mateix i tenen el mateix morfisme característic, per la unicitat de la definició. Recíprocament, si $\chi_m=\chi_n$, aleshores $\chi_n\comp m=\chi_m\comp m=\mathsf{true}\comp\,!_U$, perquè el quadrat de $m$ commuta, i per la primera part $[m]\leq[n]$; simètricament, $[n]\leq[m]$, i per l'antisimetria de $\Sub(E)$, $[m]=[n]$. Així, el morfisme característic determina el subobjecte, que és la injectivitat de la bijecció $\Sub(E)\cong\Hom(E,\Omega)$ vista ara en termes de pertinença.
\end{solucio}

\begin{exercici}\label{pr-mono-equalitzador}
En una categoria amb classificador de subobjectes, sigui $m\colon U\rightarrowtail E$ un subobjecte amb morfisme característic $\chi_m$. Demostra que $m$ és un equalitzador de $\chi_m$ i $\mathsf{true}\comp\,!_E$. Dedueix que un morfisme que és alhora mono i epi és un isomorfisme.
\end{exercici}

\begin{solucio}
\textbf{Equalitza.} Pel quadrat del classificador, $\chi_m\comp m=\mathsf{true}\comp\,!_U=\mathsf{true}\comp\,!_E\comp m$, perquè $!_E\comp m$ i $!_U$ són fletxes cap al terminal i, per tant, coincideixen.

\textbf{És universal.} Sigui $x\colon X\to E$ amb $\chi_m\comp x=\mathsf{true}\comp\,!_E\comp x$. El membre dret és $\mathsf{true}\comp\,!_X$, de manera que, per la pertinença per elements generalitzats de la Teoria, $x$ factoritza a través de $m$. La factorització és única perquè $m$ és mono. Així $m$ és un equalitzador.

\textbf{Mono i epi implica iso.} Suposem, a més, que $m$ és epi. De $\chi_m\comp m=(\mathsf{true}\comp\,!_E)\comp m$ i la cancel·lació per la dreta, $\chi_m=\mathsf{true}\comp\,!_E$. Aleshores $\id_E$ també equalitza les dues fletxes, i per la propietat de l'equalitzador hi ha $h\colon E\to U$ amb $m\comp h=\id_E$. A més, $m\comp(h\comp m)=m=m\comp\id_U$, i com que $m$ és mono, $h\comp m=\id_U$. Per tant $m$ és un isomorfisme.

Al capítol~\ref{cap:categories} vam veure que, en general, un bimorfisme no és un isomorfisme: a $\cat{Top}$, per exemple, una bijecció contínua pot no ser un homeomorfisme. Una categoria amb classificador no admet aquest fenomen. Als Exercicis es proposa treure'n la conseqüència per a $\cat{Top}$.
\end{solucio}

\section{Topos}

\begin{exercici}
Comprova que $\cat{Set}$ satisfà les tres condicions de la definició de topos.
\end{exercici}

\begin{solucio}
\emph{Límits finits}: $\cat{Set}$ té terminal (un punt), productes (producte cartesià) i equalitzadors (subconjunts definits per igualtats), i per tant tots els límits finits. \emph{Classificador}: l'exercici anterior mostra que $\{0,1\}$ n'és un. \emph{Exponencials}: l'exponencial $C^{B}$ és el conjunt d'aplicacions $B\to C$, i vam veure al capítol de tancament cartesià que $\cat{Set}$ és cartesiana tancada. Per tant $\cat{Set}$ és un topos.
\end{solucio}

\begin{exercici}
Demostra que en un topos l'objecte de parts $\mathcal{P}(E)=\Omega^{E}$ satisfà $\Hom(A,\mathcal{P}(E))\cong\Sub(A\times E)$.
\end{exercici}

\begin{solucio}
Component les dues propietats universals disponibles en un topos:
\[
\Hom(A,\Omega^{E})\cong\Hom(A\times E,\Omega)\cong\Sub(A\times E),
\]
on la primera bijecció és la de l'exponencial (currificació) i la segona és la representabilitat del classificador (\ref{prop:omega-representa}). Totes dues són naturals en $A$. Així, els morfismes cap a $\Omega^{E}$ es corresponen amb subobjectes de $A\times E$, és a dir, amb relacions entre $A$ i $E$. A $\cat{Set}$, això recupera que $\mathcal{P}(E)$ és el conjunt de parts i que una aplicació $A\to\mathcal{P}(E)$ és una relació entre $A$ i $E$.
\end{solucio}

\section{Prefeixos}

\begin{exercici}\label{pr-sedassos}
Sigui $\cat{C}$ una categoria petita, $F$ un prefeix sobre $\cat{C}$ i $U\rightarrowtail F$ un subprefeix, és a dir, $U(C)\subseteq F(C)$ per a cada $C$, amb $F(f)$ enviant $U(C)$ dins de $U(D)$ per a cada $f\colon D\to C$. Per a $x\in F(C)$, posem
\[
\chi_U(C)(x)=\{f\colon D\to C\mid F(f)(x)\in U(D)\}.
\]
Comprova que $\chi_U(C)(x)$ és un sedàs sobre $C$, que les aplicacions $\chi_U(C)$ formen una transformació natural $\chi_U\colon F\Rightarrow\Omega$, i que $U$ és el pullback de $\mathsf{true}$ al llarg de $\chi_U$.
\end{exercici}

\begin{solucio}
\textbf{És un sedàs.} Si $f\colon D\to C$ pertany a $\chi_U(C)(x)$ i $g\colon X\to D$, aleshores $F(f\comp g)(x)=F(g)\bigl(F(f)(x)\bigr)$, i com que $F(f)(x)\in U(D)$ i $U$ és un subprefeix, $F(g)$ l'envia dins de $U(X)$. Per tant $f\comp g$ pertany a $\chi_U(C)(x)$.

\textbf{Naturalitat.} Per a $h\colon D\to C$ i $x\in F(C)$, cal veure que $\chi_U(D)(F(h)(x))=h^{*}\bigl(\chi_U(C)(x)\bigr)$. Per a $g\colon X\to D$,
\[
\begin{aligned}
g\in\chi_U(D)\bigl(F(h)(x)\bigr)
&\iff F(g)\bigl(F(h)(x)\bigr)\in U(X)\\
&\iff F(h\comp g)(x)\in U(X)\\
&\iff h\comp g\in\chi_U(C)(x),
\end{aligned}
\]
i l'última condició és la definició de $g\in h^{*}\bigl(\chi_U(C)(x)\bigr)$.

\textbf{Pullback.} Com que els límits de prefeixos es calculen objecte a objecte, n'hi ha prou de veure que, per a cada $C$, $U(C)$ és el conjunt dels $x\in F(C)$ amb $\chi_U(C)(x)$ igual al sedàs màxim. Si $x\in U(C)$, tot $f\colon D\to C$ compleix $F(f)(x)\in U(D)$, de manera que el sedàs és màxim. Recíprocament, si el sedàs és màxim, conté $\id_C$, i aleshores $x=F(\id_C)(x)\in U(C)$. La unicitat de $\chi_U$ es comprova de la mateixa manera, mirant quins morfismes porten $x$ dins de $U$.

A $\cat{Set}$, vist com a prefeixos sobre la categoria d'un sol objecte i només la identitat, els sedassos sobre l'únic objecte són $\varnothing$ i $\{\id\}$, i es recupera $\Omega=\{0,1\}$.
\end{solucio}

\begin{exercici}\label{pr-conjuncio-omega}
En un topos, sigui $\wedge\colon\Omega\times\Omega\to\Omega$ el morfisme característic del subobjecte $\langle\mathsf{true},\mathsf{true}\rangle\colon 1\rightarrowtail\Omega\times\Omega$. Demostra que, per a dos subobjectes $U,V$ de $E$,
\[
\chi_{U\wedge V}=\wedge\comp\langle\chi_U,\chi_V\rangle,
\]
on $U\wedge V$ és la intersecció. Interpreta-ho a $\cat{Set}$.
\end{exercici}

\begin{solucio}
Per l'exercici~\ref{pr-ordre-pertinenca}, n'hi ha prou de veure que tots dos morfismes classifiquen el mateix subobjecte, és a dir, que un element generalitzat $x\colon X\to E$ hi pertany en un cas si i només si hi pertany en l'altre. Per la pertinença per elements generalitzats:
\begin{itemize}
  \item $x$ factoritza a través de $U\wedge V$ si i només si factoritza a través de $U$ i de $V$, perquè la intersecció és el pullback;
  \item és a dir, si i només si $\chi_U\comp x=\mathsf{true}\comp\,!_X$ i $\chi_V\comp x=\mathsf{true}\comp\,!_X$;
  \item és a dir, si i només si $\langle\chi_U,\chi_V\rangle\comp x=\langle\mathsf{true},\mathsf{true}\rangle\comp\,!_X$, que és la pertinença de $\langle\chi_U,\chi_V\rangle\comp x$ al subobjecte $\langle\mathsf{true},\mathsf{true}\rangle$;
  \item és a dir, si i només si $\wedge\comp\langle\chi_U,\chi_V\rangle\comp x=\mathsf{true}\comp\,!_X$.
\end{itemize}
Sigui $W$ el subobjecte classificat per $\wedge\comp\langle\chi_U,\chi_V\rangle$. Els dos subobjectes $U\wedge V$ i $W$ tenen, doncs, els mateixos elements generalitzats. Prenent com a $x$ un representant de $U\wedge V$, obtenim $U\wedge V\leq W$, i prenent-ne un de $W$, $W\leq U\wedge V$. Per tant són iguals, i per la unicitat del morfisme característic, $\chi_{U\wedge V}=\wedge\comp\langle\chi_U,\chi_V\rangle$.

A $\cat{Set}$, $\wedge\colon\{0,1\}\times\{0,1\}\to\{0,1\}$ val $1$ exactament en $(1,1)$: és la taula de veritat de la conjunció, i la fórmula diu que la funció característica d'una intersecció és el producte, o el mínim, de les funcions característiques.
\end{solucio}

\begin{exercici}\label{pr-reindexacio-graph}
Sigui $G$ el cicle de l'exemple~\ref{ex:omega-graph} de la Teoria, amb vèrtexs $a,b,c$, arestes $e\colon a\to b$, $f\colon b\to c$, $k\colon c\to a$, i el subgraf $H$ de vèrtexs $a,b$ i aresta $e$. Sigui $G'$ el cicle de longitud sis, amb vèrtexs $p_0,\dots,p_5$ i arestes $u_i\colon p_i\to p_{i+1}$ (índexs mòdul $6$), i sigui $F\colon G'\to G$ el morfisme que \enquote{dona dues voltes}: $p_i\mapsto a,b,c$ i $u_i\mapsto e,f,k$ segons que $i$ sigui $0,1,2$ mòdul $3$. Calcula $F^{*}H$ i comprova que $\chi_{F^{*}H}=\chi_H\comp F$.
\end{exercici}

\begin{solucio}
\textbf{La reindexació.} Es calcula per antiimatge, component a component (exemple de subgrafs del capítol~\ref{cap:subobjectes}). Els vèrtexs de $F^{*}H$ són els que van a $a$ o a $b$, és a dir, $p_0,p_1,p_3,p_4$. Les arestes són les que van a $e$, és a dir, $u_0$ i $u_3$.

\textbf{La precomposició.} Amb la taula de $\chi_H$ de la Teoria ($a,b\mapsto1$, $c\mapsto0$, $e\mapsto\top$, $f\mapsto\sigma$, $k\mapsto\tau$), el morfisme $\chi_H\comp F$ envia $p_0,p_1,p_3,p_4\mapsto1$ i $p_2,p_5\mapsto0$, i
\[
u_0,u_3\mapsto\top,\qquad u_1,u_4\mapsto\sigma,\qquad u_2,u_5\mapsto\tau.
\]

\textbf{El morfisme característic de $F^{*}H$.} Apliquem directament la regla de la Teoria a $F^{*}H\subseteq G'$. Les arestes $u_0$ i $u_3$ són al subgraf, i van a $\top$. Les arestes $u_1\colon p_1\to p_2$ i $u_4\colon p_4\to p_5$ tenen l'origen dins i el destí fora, i van a $\sigma$. Les arestes $u_2\colon p_2\to p_3$ i $u_5\colon p_5\to p_0$ tenen l'origen fora i el destí dins, i van a $\tau$. Sobre els vèrtexs coincideix també amb $\chi_H\comp F$. Així $\chi_{F^{*}H}=\chi_H\comp F$: substituir és precompondre, ara en un topos que no és $\cat{Set}$, i al llarg d'un morfisme que no és injectiu.
\end{solucio}

\section{Connectius i igualtat}

\begin{exercici}\label{pr-igualtat-omega}
En una categoria amb productes binaris i classificador, sigui $\delta_A=\langle\id_A,\id_A\rangle\colon A\to A\times A$. Comprova que $\delta_A$ és un monomorfisme i que, per a dos elements generalitzats $x,y\colon X\to A$,
\[
\chi_{\delta_A}\comp\langle x,y\rangle=\mathsf{true}\comp\,!_X
\quad\Longleftrightarrow\quad
x=y.
\]
Descriu $\chi_{\delta_A}$ a $\cat{Set}$.
\end{exercici}

\begin{solucio}
\textbf{És mono.} La projecció $\pi_1$ compleix $\pi_1\comp\delta_A=\id_A$, de manera que $\delta_A$ té una retracció i és un monomorfisme escindit.

\textbf{Pertinença a la diagonal.} Per la pertinença per elements generalitzats de la Teoria, la igualtat de l'esquerra equival a dir que $\langle x,y\rangle$ factoritza a través de $\delta_A$, és a dir, que hi ha $z\colon X\to A$ amb $\langle z,z\rangle=\langle x,y\rangle$. Component amb les projeccions, això vol dir $x=z=y$. Recíprocament, si $x=y$, n'hi ha prou de prendre $z=x$.

\textbf{A $\cat{Set}$.} $\chi_{\delta_A}(a,b)=1$ si $a=b$ i $0$ altrament: és la delta de Kronecker. En termes lògics, $\chi_{\delta_A}$ és el predicat d'igualtat \enquote{$x=y$} sobre $A$, i la diagonal n'és l'extensió. És la fila \enquote{igualtat} del mapa del capítol~\ref{cap:pont}.
\end{solucio}

\begin{exercici}\label{pr-disjuncio-negacio-set}
A $\cat{Set}$, amb $\Omega=\{0,1\}$, sigui ${\vee}\colon\Omega\times\Omega\to\Omega$ el morfisme característic del subconjunt $\{(0,1),(1,0),(1,1)\}$, i sigui $\neg\colon\Omega\to\Omega$ el del subconjunt $\{0\}$. Comprova que, per a subconjunts $U,V\subseteq E$,
\[
\chi_{U\cup V}={\vee}\comp\langle\chi_U,\chi_V\rangle,
\qquad
\chi_{E\setminus U}=\neg\comp\chi_U.
\]
Explica per què la segona igualtat no té un anàleg directe a $\cat{Graph}$.
\end{exercici}

\begin{solucio}
Per a $x\in E$, $({\vee}\comp\langle\chi_U,\chi_V\rangle)(x)=1$ si i només si $(\chi_U(x),\chi_V(x))$ és en $\{(0,1),(1,0),(1,1)\}$, és a dir, si i només si $x\in U$ o $x\in V$. Les dues aplicacions valen, doncs, $1$ exactament a $U\cup V$, i coincideixen. Anàlogament, $(\neg\comp\chi_U)(x)=1$ si i només si $\chi_U(x)=0$, és a dir, si i només si $x\notin U$. Els morfismes $\vee$ i $\neg$ són les taules de veritat de la disjunció i de la negació, i les dues fórmules diuen que la unió i el complement es calculen \enquote{punt a punt} sobre els valors de veritat, com la intersecció a l'exercici~\ref{pr-conjuncio-omega}.

A $\cat{Graph}$, el complement de vèrtexs i d'arestes d'un subgraf no és, en general, un subgraf: pot contenir una aresta sense un dels seus extrems. El millor substitut és el subgraf més gran disjunt de $U$, però, com hem vist a l'exemple~\ref{ex:omega-graph} de la Teoria, la seva unió amb $U$ pot no ser tot el graf. No hi ha, doncs, cap morfisme $\Omega\to\Omega$ que faci el paper de $\neg$ amb la propietat $U\cup\neg U=E$.
\end{solucio}

\fi

% ============================================================
% PART III — EXERCICIS PROPOSATS
% ============================================================
\part{Exercicis proposats}
\setcounter{chapter}{0}
\chapter{Categories}

Els exercicis següents completen els resolts al capítol corresponent de la part de Pràctica. Cada enunciat va acompanyat d'una indicació breu per si el lector s'encalla; convé, però, intentar-lo primer sense mirar-la. La marca \enquote{$\star$} assenyala un exercici de consolidació i \enquote{$\star\star$} un d'ampliació.

\section{Categories i axiomes}

\begin{exercici}
Comprova que un conjunt qualsevol $X$ dona lloc a una categoria \textbf{discreta}: els objectes són els elements de $X$ i els únics morfismes són les identitats. Verifica els axiomes. \hfill{\normalfont$\star$}
\begin{indicacio}
No hi ha cap morfisme entre objectes diferents; l'única composició possible en cada objecte és $\id_x\comp\id_x$.
\end{indicacio}
\end{exercici}

\begin{exercici}
Sigui $\cat{C}$ una categoria. Comprova que fixar un objecte $A$ i considerar només les fletxes $A\to A$, amb la composició de $\cat{C}$, dona un \textbf{monoide}. Quin és el seu element neutre? \hfill{\normalfont$\star$}
\begin{indicacio}
Els morfismes $A\to A$ es poden compondre entre ells i el resultat torna a anar d'$A$ a $A$. El neutre és $\id_A$.
\end{indicacio}
\end{exercici}

\begin{exercici}[Ampliació, per a lectors amb topologia]
Comprova que els espais topològics amb les aplicacions contínues formen una categoria $\cat{Top}$. Cal usar dues propietats de les aplicacions contínues: quines? \hfill{\normalfont$\star$}
\begin{indicacio}
La composició de contínues és contínua, i la identitat és contínua.
\end{indicacio}
\end{exercici}

\section{Fletxes, camins i diagrames commutatius}

\begin{exercici}
Escriu l'equació que expressa la commutativitat del triangle
\[
\begin{tikzpicture}[baseline=(m.center)]
  \matrix (m) [matrix of math nodes, column sep=3.3em, row sep=2.6em] {
    A & B \\
      & C \\
  };
  \draw[-{Stealth[scale=1.05]}] (m-1-1) -- node[above] {$f$} (m-1-2);
  \draw[-{Stealth[scale=1.05]}] (m-1-2) -- node[right] {$g$} (m-2-2);
  \draw[-{Stealth[scale=1.05]}] (m-1-1) -- node[below left] {$h$} (m-2-2);
\end{tikzpicture}
\]
i explica quins són els dos camins que es comparen. \hfill{\normalfont$\star$}
\begin{indicacio}
Tots dos camins comencen a $A$ i acaben a $C$; un passa per $B$ i l'altre és la fletxa directa.
\end{indicacio}
\end{exercici}

\begin{exercici}
Representa mitjançant un quadrat commutatiu l'equació
\[
q\comp u=v\comp p,
\]
on $u\colon A\to B$, $q\colon B\to D$, $p\colon A\to C$ i $v\colon C\to D$. \hfill{\normalfont$\star$}
\begin{indicacio}
Col·loca $A$ a l'angle superior esquerre i $D$ a l'inferior dret. Els dos membres de l'equació són dos camins d'$A$ a $D$.
\end{indicacio}
\end{exercici}

\begin{exercici}
Dona un exemple de quatre aplicacions entre conjunts que formin un quadrat que \emph{no} commuti. Indica explícitament un element sobre el qual les dues composicions donin resultats diferents. \hfill{\normalfont$\star$}
\begin{indicacio}
Pots treballar amb conjunts molt petits, per exemple $\{0,1\}$, i comparar els dos camins sobre l'element $0$.
\end{indicacio}
\end{exercici}

\section{Primers exemples}

\begin{exercici}
En un sistema deductiu $T$, suposem
\[
A\vdash_T B,
\qquad
B\vdash_T C,
\qquad
C\vdash_T A.
\]
Dibuixa el diagrama corresponent a $\cat{Prov}_T$ i comprova que $A$, $B$ i $C$ són isomorfs dos a dos. \hfill{\normalfont$\star$}
\begin{indicacio}
La transitivitat dona també $A\vdash_T C$, $B\vdash_T A$ i $C\vdash_T B$. En una categoria prima, dues fletxes en sentits oposats ja donen un isomorfisme.
\end{indicacio}
\end{exercici}

\begin{exercici}
Comprova directament que la composició de dos morfismes de grafs és un morfisme de grafs. Si $F\colon G\to H$ i $K\colon H\to L$, verifica les equacions d'origen i destí per a $K\comp F$. \hfill{\normalfont$\star$}
\begin{indicacio}
Per a l'origen, parteix de
\[
F_V\comp s_G=s_H\comp F_E,
\qquad
K_V\comp s_H=s_L\comp K_E,
\]
i compon les igualtats adequadament. Fes el mateix amb $t$.
\end{indicacio}
\end{exercici}

\begin{exercici}
Explica amb precisió quina informació té un graf dirigit que encara no té l'estructura d'una categoria. Després indica què afegeix la categoria lliure sobre aquest graf. \hfill{\normalfont$\star$}
\begin{indicacio}
Un graf té vèrtexs, arestes, origen i destí. Una categoria necessita, a més, identitats, composició i axiomes; la categoria lliure els genera mitjançant camins.
\end{indicacio}
\end{exercici}

\section{Isomorfismes}

\begin{exercici}
Comprova que a la categoria $\cat{Mon}$ dels monoides i els homomorfismes, un isomorfisme és exactament un homomorfisme bijectiu. Compara-ho amb el que pot passar a $\cat{Pos}$. \hfill{\normalfont$\star\star$}
\begin{indicacio}
Comprova que la inversa d'un homomorfisme bijectiu de monoides preserva l'operació i el neutre, de manera que torna a ser un homomorfisme. Per a $\cat{Pos}$, en canvi, la inversa d'una aplicació monòtona bijectiva pot no ser monòtona. Fes el contrast explícit amb un exemple: pren el conjunt $\{0,1\}$ amb l'ordre \emph{discret} (cap dels dos elements comparable amb l'altre) com a domini, i el mateix conjunt amb l'ordre $0<1$ com a codomini. La identitat subjacent és una aplicació monòtona (trivialment, perquè al domini no hi ha cap parella comparable) i bijectiva, però la seva inversa ---la identitat de l'ordre $0<1$ cap a l'ordre discret--- no és monòtona, ja que hauria d'enviar $0<1$ a dos elements incomparables. Per tant no és un isomorfisme a $\cat{Pos}$, tot i ser una bijecció monòtona.
\end{indicacio}
\end{exercici}

\begin{exercici}
Comprova que si $g\comp f$ és un isomorfisme, no cal que ho siguin ni $f$ ni $g$. Busca'n un exemple senzill a $\cat{Set}$ i representa'l amb fletxes. \hfill{\normalfont$\star$}
\begin{indicacio}
Pensa en una injecció seguida d'una projecció entre conjunts de mides diferents.
\end{indicacio}
\end{exercici}

\section{Classes especials de morfismes}

\begin{exercici}
Comprova que tot isomorfisme té secció. Després, prenent $A=\{0,1\}$, $B=\{\ast\}$ i $f$ l'única aplicació $A\to B$, dona una secció $s\colon B\to A$ explícita i comprova que $f$ no és isomorfisme. Observa que en aquest cas $f$ té \emph{dues} seccions distintes: quines? \hfill{\normalfont$\star$}
\begin{indicacio}
Tria $s(\ast)=0$ o $s(\ast)=1$. En tots dos casos $f\comp s=\id_B$, però $s\comp f\neq\id_A$. Les dues tries donen les dues seccions.
\end{indicacio}
\end{exercici}

\begin{exercici}
A la categoria $\cat{Set}$, exhibeix morfismes que il·lustrin els tres casos possibles pel que fa a l'existència de seccions: un morfisme sense cap secció, un amb una única secció i un amb més d'una. \hfill{\normalfont$\star\star$}
\begin{indicacio}
Per a \emph{cap} secció, cerca una aplicació no exhaustiva, com $\emptyset\to\{\ast\}$ o una injecció no exhaustiva. Per a \emph{una única}, una aplicació exhaustiva on cada element del codomini tingui una sola antiimatge (una bijecció). Per a \emph{diverses}, una aplicació exhaustiva amb alguna fibra de més d'un element: cada elecció d'antiimatges dona una secció distinta. Recorda que a $\cat{Set}$ triar una secció d'una aplicació exhaustiva és triar un element a cada fibra.
\end{indicacio}
\end{exercici}

\begin{exercici}
Demostra la caracterització dual de l'exercici resolt a la part de pràctica: un morfisme $f$ és un isomorfisme si i només si és un epimorfisme escindit (té secció) i un monomorfisme. \hfill{\normalfont$\star$}
\begin{indicacio}
És el dual de \enquote{mono escindit i epi $\Rightarrow$ iso}: parteix d'una secció $s$ amb $f\comp s=\id_B$ i fes servir que $f$ és mono per veure que $s\comp f=\id_A$.
\end{indicacio}
\end{exercici}

\begin{exercici}
Sigui $\cat{C}$ una categoria prima. Per l'exercici resolt~\ref{exr:pr-prima-mono-epi} de la Pràctica, tot morfisme de $\cat{C}$ és un bimorfisme. Demostra que un morfisme $f\colon A\to B$ de $\cat{C}$ és un isomorfisme si i només si existeix algun morfisme $B\to A$. Aplica-ho a la fletxa $u\colon 0\to 1$ de $\cat{2}$ i a dues fórmules interdeduïbles i distintes de $\cat{Prov}_T$. \hfill{\normalfont$\star$}
\begin{indicacio}
Si $g\colon B\to A$ existeix, $g\comp f$ i $\id_A$ pertanyen tots dos a $\Hom(A,A)$, que té com a màxim un element; raona igual amb $f\comp g$ i $\id_B$. A $\cat{2}$ no hi ha cap fletxa $1\to 0$, de manera que $u$ és un bimorfisme que no és iso. A $\cat{Prov}_T$, en canvi, la fletxa que expressa $p\vdash_T q$ és un isomorfisme que no és una identitat quan també $q\vdash_T p$.
\end{indicacio}
\end{exercici}

\begin{exercici}
Estudia quin factor hereta la propietat en una composició. Demostra que si $g\comp f$ és un monomorfisme, aleshores $f$ és un monomorfisme; i que, dualment, si $g\comp f$ és un epimorfisme, aleshores $g$ és un epimorfisme. Comprova amb un contraexemple que, en el primer cas, $g$ no cal que sigui mono, i en el segon, $f$ no cal que sigui epi. \hfill{\normalfont$\star\star$}
\begin{indicacio}
Per al mono: de $f\comp u=f\comp v$, compon per l'esquerra amb $g$ per obtenir $(g\comp f)\comp u=(g\comp f)\comp v$ i aplica que $g\comp f$ és mono. El cas de l'epi és el dual: compon per la dreta. Observa l'asimetria: d'un mono $g\comp f$ hereta el factor de la \emph{dreta} ($f$), i d'un epi, el de l'\emph{esquerra} ($g$). Per als contraexemples, pensa en una injecció i una projecció a $\cat{Set}$.
\end{indicacio}
\end{exercici}

\begin{exercici}
Sigui $M=\{1,e\}$ el monoide amb $e\cdot e=e$ i $1$ neutre, i sigui $\cat{B}M$ la categoria d'un sol objecte associada. Decideix si $e$ és un monomorfisme, si és un epimorfisme i si té secció o retracció. \hfill{\normalfont$\star\star$}
\begin{indicacio}
Compara $e\comp 1$ amb $e\comp e$, i $1\comp e$ amb $e\comp e$. Per a les inverses unilaterals, només hi ha dos candidats.
\end{indicacio}
\end{exercici}

\begin{exercici}
A $\cat{Set}$, demostra que tota aplicació exhaustiva $f\colon A\to B$ és un epimorfisme sense triar cap secció de $f$. Explica després per què l'enunciat \enquote{tota aplicació exhaustiva té una secció} no és una conseqüència d'aquest fet, sinó un enunciat diferent, equivalent a l'axioma de l'elecció. \hfill{\normalfont$\star\star$}
\begin{indicacio}
Si $g\comp f=h\comp f$ i $y\in B$, n'hi ha prou de tenir \emph{algun} $x$ amb $f(x)=y$ per concloure $g(y)=h(y)$; no cal triar-ne un per a cada $y$ alhora. Una secció, en canvi, és precisament una tria simultània d'una antiimatge per a cada element de $B$.
\end{indicacio}
\end{exercici}

\section{Construccions sobre categories}

\begin{exercici}
Descriu la categoria producte $\cat{3}\times\cat{3}$, on $\cat{3}$ és la categoria associada a l'ordre $0<1<2$. Quants objectes i quants morfismes té? Dibuixa les fletxes bàsiques i explica quines altres fletxes apareixen per composició. \hfill{\normalfont$\star\star$}
\begin{indicacio}
$\cat{3}$ té tres objectes i sis morfismes. En el producte, els objectes i els morfismes són parelles.
\end{indicacio}
\end{exercici}

\begin{exercici}
Descriu explícitament la categoria $\cat{3}/2$ sobre l'objecte $2$ de $\cat{3}$, on $\cat{3}$ és la categoria associada a l'ordre $0<1<2$. Enumera'n els objectes i els morfismes i identifica, si pots, una categoria finita isomorfa a aquesta categoria sobre un objecte. \hfill{\normalfont$\star$}
\begin{indicacio}
Els objectes de $\cat{3}/2$ són les fletxes de $\cat{3}$ amb codomini $2$, i un morfisme entre dos d'aquests objectes és una fletxa de $\cat{3}$ que fa commutar el triangle corresponent. Com que $\cat{3}$ és prima, tots els triangles commuten. Observa també que de cada objecte de $\cat{3}$ surt exactament una fletxa cap a $2$.
\end{indicacio}
\end{exercici}

\begin{exercici}
Compara directament $(\cat{C}\times\cat{D})\op$ i $\cat{C}\op\times\cat{D}\op$. Comprova que tenen els mateixos objectes i que els seus morfismes, identitats i composicions s'identifiquen component a component. \hfill{\normalfont$\star$}
\begin{indicacio}
Un morfisme $(A,X)\to(B,Y)$ a $(\cat{C}\times\cat{D})\op$ és, per definició, un morfisme $(B,Y)\to(A,X)$ a $\cat{C}\times\cat{D}$. Escriu què significa això en cadascuna de les dues components.
\end{indicacio}
\end{exercici}

\begin{exercici}
Per a cadascuna de les subcategories següents de $\cat{Set}$, digues si és \textbf{plena}, \textbf{ampla}, totes dues coses o cap: (a)~$\cat{FinSet}$, amb els conjunts finits per objectes i totes les aplicacions entre ells; (b)~la subcategoria amb tots els conjunts per objectes i només les aplicacions injectives; (c)~la subcategoria amb tots els conjunts per objectes i totes les aplicacions (és a dir, $\cat{Set}$ mateixa); (d)~la subcategoria amb els conjunts finits per objectes i només les aplicacions injectives entre ells. \hfill{\normalfont$\star$}
\begin{indicacio}
Recorda que \emph{plena} vol dir que es conserven \emph{tots} els morfismes entre els objectes escollits (la subcategoria queda determinada només pels objectes), i \emph{ampla} vol dir que es conserven \emph{tots} els objectes. Comprova cada cas contra aquestes dues condicions per separat: poden fallar independentment. En particular, mira si (d) és plena o ampla, i observa que una subcategoria pot no ser ni l'una ni l'altra.
\end{indicacio}
\end{exercici}

\section{Principi de dualitat}

\begin{exercici}
Escriu el dual de l'afirmació següent i justifica'l sense fer una nova demostració completa:
\begin{quote}
Si $f\colon A\to B$ té una retracció, aleshores $f$ és un monomorfisme.
\end{quote}
\hfill{\normalfont$\star$}
\begin{indicacio}
Passa a la categoria oposada i identifica quina noció és dual de \enquote{retracció} i quina és dual de \enquote{monomorfisme}. Recorda que, en dualitzar, s'inverteixen tant el sentit de les fletxes com l'ordre de les composicions.
\end{indicacio}
\end{exercici}

\section{Grafs i categories lliures}

\begin{exercici}
Classifica, \textbf{fins a isomorfisme}, els grafs dirigits que generen una categoria lliure amb exactament \textbf{quatre} morfismes. Recorda que les identitats compten com a morfismes. \hfill{\normalfont$\star\star$}
\begin{indicacio}
Cada vèrtex aporta una identitat i cada aresta un camí de longitud~$1$: quantes arestes i quants vèrtexs pot tenir, com a màxim, el graf? Què passaria si hi hagués un cicle dirigit? Després, separa els casos segons el nombre de vèrtexs.
\end{indicacio}
\end{exercici}

\begin{exercici}
Considera el graf
\[
A\xrightarrow{f}B\xrightarrow{g}C.
\]
Llista tots els morfismes de la categoria lliure que genera, dibuixa també la fletxa composta i comprova que hi ha exactament sis morfismes. \hfill{\normalfont$\star$}
\begin{indicacio}
Compta els tres camins buits, les dues arestes $f$ i $g$, i el camí compost de longitud dos.
\end{indicacio}
\end{exercici}

\begin{exercici}
Per què un graf amb un cicle, per exemple una aresta $a\colon A\to A$, genera sempre una categoria lliure amb infinits morfismes? \hfill{\normalfont$\star$}
\begin{indicacio}
El cicle es pot recórrer tantes vegades com es vulgui: $a,a^2,a^3,\ldots$, i cada repetició dona un camí diferent.
\end{indicacio}
\end{exercici}

\section{Qüestions de mida}

\begin{exercici}
Comprova que un preordre, vist com a categoria, és una categoria \textbf{petita} sempre que el conjunt subjacent $P$ sigui un conjunt. Quants morfismes hi ha, com a màxim, entre dos objectes? \hfill{\normalfont$\star$}
\begin{indicacio}
Entre dos objectes hi ha zero o un morfisme; la col·lecció total de morfismes es pot identificar amb un subconjunt de $P\times P$.
\end{indicacio}
\end{exercici}

\begin{exercici}
Explica per què el fet que tots els objectes d'una categoria siguin finits no implica que la categoria sigui petita. Aplica-ho a $\cat{FinSet}$, la categoria que té per objectes tots els conjunts finits i per morfismes totes les aplicacions entre ells. \hfill{\normalfont$\star$}
\begin{indicacio}
La mida de cada objecte i la mida de la col·lecció de tots els objectes són qüestions diferents. Per a cada conjunt $x$, el conjunt $\{x\}$ és finit: quants objectes té, doncs, $\cat{FinSet}$? Compara-ho amb la paradoxa de Russell.
\end{indicacio}
\end{exercici}

\begin{exercici}
Explica per què $\cat{Graph}$ és localment petita: donats dos grafs dirigits $G$ i $H$, per què els morfismes de grafs $G\to H$ formen un conjunt? \hfill{\normalfont$\star$}
\begin{indicacio}
Un morfisme de grafs és una parella d'aplicacions $V_G\to V_H$ i $E_G\to E_H$ que satisfà dues condicions. Aquestes parelles formen un subconjunt d'un producte de dos conjunts d'aplicacions.
\end{indicacio}
\end{exercici}

\chapter{Functors}

Els exercicis següents completen els resolts al capítol corresponent de la part de Pràctica. Cada enunciat va acompanyat d'una indicació breu per si el lector s'encalla; convé, però, intentar-lo primer sense mirar-la. La marca \enquote{$\star$} assenyala un exercici de consolidació i \enquote{$\star\star$} un d'ampliació.

\section{Definició de functor}

\begin{exercici}
Comprova que el functor identitat $\id_{\cat{C}}\colon\cat{C}\to\cat{C}$ i, per a cada objecte $D$ d'una categoria $\cat{D}$, el functor constant $\Delta_D\colon\cat{C}\to\cat{D}$ ---que envia tot objecte a $D$ i tot morfisme a $\id_D$ (exemple~\ref{ex:functor-constant})--- són functors. \hfill{\normalfont$\star$}
\begin{indicacio}
Per al functor constant, comprova que $\Delta_D(\id_A)=\id_D$ i que $\Delta_D(g\comp f)=\id_D=\id_D\comp\id_D=\Delta_D(g)\comp\Delta_D(f)$.
\end{indicacio}
\end{exercici}

\begin{exercici}\label{ex:identitat-necessaria}
Considera una assignació que respecti dominis, codominis i composició, però de la qual \emph{no} se suposi que preservi identitats. Comprova, amb un contraexemple senzill, que la preservació de la identitat no se'n dedueix, i conclou que la condició $F(\id_A)=\id_{F(A)}$ és realment necessària a la definició de functor. \hfill{\normalfont$\star$}
\begin{indicacio}
Pren el monoide de dos elements $\{1,e\}$ amb $e\comp e=e$ i $e\comp m=m\comp e=e$ per a tot $m$, vist com a categoria d'un sol objecte. L'assignació que envia cada morfisme $m$ a $e$ respecta la composició, perquè $F(m\comp n)=e=e\comp e=F(m)\comp F(n)$, però envia la identitat $1$ a $e\neq 1$. Per tant no preserva identitats malgrat respectar la composició.
\end{indicacio}
\end{exercici}

\section{Exemples de functors}

\begin{exercici}
Comprova que hi ha un functor $V\colon\cat{Graph}\to\cat{Set}$ que envia cada graf $G$ al seu conjunt de vèrtexs $V_G$ i cada morfisme de grafs $F=(F_V,F_E)\colon G\to H$ a la seva component sobre vèrtexs $F_V\colon V_G\to V_H$. Quines propietats de la composició i de les identitats de $\cat{Graph}$ en garanteixen la functorialitat? Fes el mateix per al functor d'arestes $E\colon\cat{Graph}\to\cat{Set}$. \hfill{\normalfont$\star$}
\begin{indicacio}
La identitat d'un graf té per component sobre vèrtexs la identitat de $V_G$, i la composició de morfismes de grafs es fa component a component: $(F'\comp F)_V=F'_V\comp F_V$. Quan hagis estudiat els homfunctors (subsecció~\ref{subsec:homfunctors}), compara $V$ amb $\Hom_{\cat{Graph}}(\mathbf V,-)$, on $\mathbf V$ és el graf d'un sol vèrtex.
\end{indicacio}
\end{exercici}

\begin{exercici}[Ampliació, per a lectors amb topologia]
Comprova que hi ha un functor oblit $U\colon\cat{Top}\to\cat{Set}$ i descriu explícitament què fa sobre objectes i sobre morfismes. Quines propietats de les aplicacions contínues en garanteixen la functorialitat? \hfill{\normalfont$\star$}
\begin{indicacio}
$U$ envia cada espai al seu conjunt de punts i cada aplicació contínua a la mateixa aplicació. Cal que la identitat sigui contínua i que la composició de contínues sigui contínua.
\end{indicacio}
\end{exercici}

\begin{exercici}[Ampliació, per a lectors amb àlgebra]
Sigui $G$ un grup vist com a categoria $\cat{B}G$ amb un sol objecte. Comprova que un functor $\cat{B}G\to\cat{Set}$ és el mateix que una \textbf{acció} de $G$ sobre un conjunt, i que la imatge de cada element del grup és necessàriament una \emph{bijecció}. \hfill{\normalfont$\star$}
\begin{indicacio}
Segueix el cas del monoide. La diferència és que en un grup cada element té invers; usa que un functor preserva isomorfismes i que tot morfisme de $\cat{B}G$ és un isomorfisme.
\end{indicacio}
\end{exercici}

\begin{exercici}
Descriu un functor $\cat{2}\to\cat{C}$, on $\cat{2}$ és la categoria associada a l'ordre $0<1$. Comprova que triar un tal functor equival a triar un morfisme de $\cat{C}$. \hfill{\normalfont$\star$}
\begin{indicacio}
El functor ha d'enviar l'única fletxa no trivial $0\to1$ a un morfisme de $\cat{C}$; els dos objectes es determinen com el seu domini i el seu codomini.
\end{indicacio}
\end{exercici}

\section{Composició de functors i la categoria \texorpdfstring{$\cat{Cat}$}{Cat}}

\begin{exercici}
Comprova que un isomorfisme a $\cat{Cat}$ ---és a dir, un functor $F\colon\cat{C}\to\cat{D}$ amb un functor invers $G$ tal que $G\comp F=\id_{\cat{C}}$ i $F\comp G=\id_{\cat{D}}$--- és bijectiu sobre objectes i sobre morfismes. \hfill{\normalfont$\star\star$}
\begin{indicacio}
De $G\comp F=\id_{\cat{C}}$ i $F\comp G=\id_{\cat{D}}$ dedueix que $F$ té inversa com a aplicació sobre objectes i com a aplicació sobre morfismes.
\end{indicacio}
\end{exercici}

\begin{exercici}
Siguin $F\colon\cat{E}\to\cat{C}$ i $G\colon\cat{E}\to\cat{D}$ dos functors. Defineix un functor $\langle F,G\rangle\colon\cat{E}\to\cat{C}\times\cat{D}$ per $\langle F,G\rangle(E)=(F(E),G(E))$ i $\langle F,G\rangle(h)=(F(h),G(h))$. Comprova que és un functor, que
\[
\pi_1\comp\langle F,G\rangle=F,\qquad \pi_2\comp\langle F,G\rangle=G,
\]
on $\pi_1,\pi_2$ són les projeccions de l'exemple~\ref{ex:projeccions-diagonal}, i que és l'únic functor $\cat{E}\to\cat{C}\times\cat{D}$ amb aquestes dues propietats. Quin functor s'obté amb $\cat{E}=\cat{C}=\cat{D}$ i $F=G=\id_{\cat{C}}$? \hfill{\normalfont$\star$}
\begin{indicacio}
La functorialitat es comprova component a component, com per a $\Delta$. Per a la unicitat, si $H\colon\cat{E}\to\cat{C}\times\cat{D}$ compleix les dues igualtats, escriu $H(E)$ i $H(h)$ com a parelles i mira què en diuen $\pi_1$ i $\pi_2$. És la propietat universal del producte, ara a nivell de categories.
\end{indicacio}
\end{exercici}

\section{Functors i isomorfismes}

\begin{exercici}
Comprova que si $F\colon\cat{C}\to\cat{D}$ és un functor i $A\cong B$ a $\cat{C}$, aleshores $F(A)\cong F(B)$ a $\cat{D}$. És a dir, els functors preserven la relació d'isomorfisme entre objectes. \hfill{\normalfont$\star$}
\begin{indicacio}
Aplica $F$ a un isomorfisme $A\to B$ i usa que la imatge d'un isomorfisme és un isomorfisme.
\end{indicacio}
\end{exercici}

\begin{exercici}
Dona un exemple de dos objectes no isomorfs que un functor envia a objectes isomorfs. Això mostra que els functors no reflecteixen, en general, la relació d'isomorfisme entre objectes. \hfill{\normalfont$\star$}
\begin{indicacio}
Un functor constant envia tots els objectes al mateix objecte, i per tant tots a objectes isomorfs entre ells.
\end{indicacio}
\end{exercici}

\section{Functors covariants i contravariants}

\begin{exercici}
Donat un functor $F\colon\cat{C}\to\cat{D}$, comprova que la mateixa regla sobre objectes i morfismes defineix un functor
\[
F\op\colon\cat{C}\op\to\cat{D}\op,
\]
el \textbf{functor oposat} de $F$. Comprova que $F$ és fidel (respectivament ple) si i només si $F\op$ ho és. \hfill{\normalfont$\star$}
\begin{indicacio}
Sobre objectes, $F\op=F$; sobre morfismes, $F\op(f)=F(f)$, però ara llegit entre els mateixos objectes a les categories oposades. Compondre a $\cat{C}\op$ i a $\cat{D}\op$ és compondre a $\cat{C}$ i a $\cat{D}$ en l'ordre invers, de manera que les dues condicions de functor es transporten. Els conjunts de morfismes són els mateixos, així que injectivitat i exhaustivitat es conserven.
\end{indicacio}
\end{exercici}

\begin{exercici}
Fixat un objecte $A$ d'una categoria localment petita, comprova que $\Hom(A,-)$ envia isomorfismes a bijeccions: si $f\colon X\to Y$ és un isomorfisme, aleshores
\[
f_*\colon\Hom(A,X)\to\Hom(A,Y)
\]
és una bijecció. \hfill{\normalfont$\star$}
\begin{indicacio}
$\Hom(A,-)$ és un functor i, per tant, preserva isomorfismes; a $\cat{Set}$ els isomorfismes són les bijeccions. Alternativament, comprova directament que $(f^{-1})_*$ és la inversa de $f_*$.
\end{indicacio}
\end{exercici}

\section{Functors fidels i plens}

\begin{exercici}
Comprova que la composició de dos functors fidels és fidel, i que la composició de dos functors plens és plena. Dedueix que la composició de functors plens i fidels és plena i fidel. \hfill{\normalfont$\star$}
\begin{indicacio}
Treballa amb les aplicacions $F_{A,B}$ sobre conjunts de morfismes: la composició de dues injeccions és injectiva, i la de dues exhaustives, exhaustiva.
\end{indicacio}
\end{exercici}

\begin{exercici}
Demostra que un functor fidel reflecteix monomorfismes i epimorfismes: si $F$ és fidel i $F(f)$ és mono (respectivament epi), aleshores $f$ és mono (respectivament epi). Indica on fas servir que $F$ és fidel, i comprova que no cal que $F$ sigui ple. \hfill{\normalfont$\star$}
\begin{indicacio}
Parteix d'una igualtat de composicions a la categoria d'origen, aplica-hi el functor i cancel·la a la d'arribada.
\end{indicacio}
\end{exercici}

\begin{exercici}
Comprova que el functor diagonal $\Delta\colon\cat{C}\to\cat{C}\times\cat{C}$ de l'exemple~\ref{ex:projeccions-diagonal} és sempre fidel, i que és ple si i només si $\cat{C}$ és prima. \hfill{\normalfont$\star$}
\begin{indicacio}
Sobre cada homconjunt, $\Delta$ actua com $f\mapsto(f,f)$. Quines parelles $(f,g)$ de $\cat{C}\times\cat{C}$ provenen d'una sola fletxa?
\end{indicacio}
\end{exercici}

\begin{exercici}
Comprova que el functor oblit $U\colon\cat{Cat}\to\cat{Graph}$, que envia cada categoria petita al seu graf subjacent, és fidel però \emph{no} ple. Mostra també que, per a tot graf $G$, hi ha un morfisme canònic de grafs $\eta_G\colon G\to U(\Free(G))$. \hfill{\normalfont$\star$}
\begin{indicacio}
Quant a ser fidel, un functor queda determinat per la seva acció sobre objectes i morfismes, que és justament el que reté el graf subjacent. Per veure que no és ple, observa que un morfisme de grafs entre els grafs subjacents no ha de respectar la composició, i per tant no prové necessàriament d'un functor. Per al morfisme canònic $\eta_G$, envia cada vèrtex a ell mateix i cada aresta de $G$ al camí de longitud~$1$ corresponent dins $\Free(G)$.
\end{indicacio}
\end{exercici}

\begin{exercici}[Ampliació]
Caracteritza exactament per a quins grafs $G$ el morfisme canònic $\eta_G\colon G\to U(\Free(G))$ de l'exercici anterior és un isomorfisme de grafs. \hfill{\normalfont$\star\star$}
\begin{indicacio}
Compara les arestes de $U(\Free(G))$ amb les imatges de les arestes de $G$. Quins camins de $\Free(G)$ existeixen sempre, sigui quin sigui $G$, i no tenen longitud~$1$? Tingues en compte també el cas extrem.
\end{indicacio}
\end{exercici}

\begin{exercici}
Comprova que un functor ple i fidel \textbf{reflecteix els isomorfismes entre objectes}, en el sentit següent:
\[
F(A)\cong F(B)\quad\Longrightarrow\quad A\cong B.
\]
Suposant a més que $\cat{C}$ i $\cat{D}$ són petites, dedueix que $F$ indueix una aplicació injectiva entre les classes d'isomorfisme d'objectes de $\cat{C}$ i les de $\cat{D}$. \hfill{\normalfont$\star\star$}
\begin{indicacio}
Pren un isomorfisme $F(A)\to F(B)$. Com que $F$ és ple, prové d'un morfisme $A\to B$; per la propietat que els functors plens i fidels reflecteixen isomorfismes, aquest morfisme és un isomorfisme, de manera que $A\cong B$. La implicació $F(A)\cong F(B)\Rightarrow A\cong B$ val sempre; la hipòtesi de petitesa només s'usa a la darrera part, per poder parlar del conjunt de classes d'isomorfisme: en una categoria petita els objectes formen un conjunt, i el quocient per la relació $\cong$ també, de manera que \enquote{l'aplicació entre classes d'isomorfisme} té sentit literal. En una categoria gran caldria una convenció addicional per manejar aquestes classes, que aquí evitem.
\end{indicacio}
\end{exercici}

\begin{exercici}
Considera l'assignació $\cat{C}\to\cat{C}\op$ que deixa fixos els objectes i envia cada morfisme $f$ a ell mateix, vist ara en sentit contrari. Comprova que \emph{no} és, en general, un functor covariant, però sí un functor contravariant ple i fidel. Quin és aquest functor? \hfill{\normalfont$\star\star$}
\begin{indicacio}
Sobre morfismes, aquesta assignació inverteix l'ordre de la composició; és el functor identitat vist com a functor contravariant $\cat{C}\to\cat{C}\op$. La bijecció sobre cada conjunt de morfismes és evident.
\end{indicacio}
\end{exercici}

\begin{exercici}
Comprova que si $F\colon\cat{C}\to\cat{D}$ és ple i fidel, i $\cat{C}$ és localment petita, aleshores $F$ estableix una bijecció
\[
\Hom_{\cat{C}}(A,B)\;\xrightarrow{\;\sim\;}\;\Hom_{\cat{D}}(F(A),F(B))
\]
per a cada parella d'objectes. Explica per què això no obliga $F$ a ser injectiu sobre objectes. \hfill{\normalfont$\star\star$}
\begin{indicacio}
La bijecció és la definició mateixa de ple i fidel. Per a la segona part, busca una categoria amb dos objectes isomorfs i exactament un morfisme entre cada parella ordenada, i compara-la amb $\cat{1}$.
\end{indicacio}
\end{exercici}

\begin{exercici}
Distingeix l'exhaustivitat \emph{estricta} sobre objectes (tot objecte del codomini és exactament una imatge) de l'\emph{essencial} (tot objecte del codomini és isomorf a una imatge). Sigui $\cat{I}$ la categoria amb dos objectes $a,b$ i exactament un morfisme entre cada parella ordenada d'objectes, i considera el functor $E\colon\cat{1}\to\cat{I}$ que envia l'únic objecte $\ast$ a $a$. Comprova que $E$ és essencialment exhaustiu però no exhaustiu sobre objectes. \hfill{\normalfont$\star\star$}
\begin{indicacio}
Observa que $E$ arriba estrictament només a $a$, però que $b\cong a$.
\end{indicacio}
\end{exercici}

\begin{exercici}
Sigui $T$ la lògica proposicional clàssica. Comprova que la negació $A\mapsto\neg A$ defineix un functor contravariant
\[
N\colon\cat{Prov}_T\op\to\cat{Prov}_T,
\]
i que, per tant, la contraposició ---de $A\vdash_T B$ es dedueix $\neg B\vdash_T\neg A$--- és exactament la functorialitat de $N$. Comprova que $N$ és fidel i que és ple. Explica quina llei lògica cal perquè $N$ sigui ple i per què, en lògica intuïcionista, $N$ continua sent un functor però deixa de ser ple. \hfill{\normalfont$\star\star$}
\begin{indicacio}
Les dues categories són primes, de manera que només cal comprovar que les fletxes imatge existeixen. Que $N$ sigui ple demana que de $\neg B\vdash_T\neg A$ se'n segueixi $A\vdash_T B$. Per veure que pot fallar en lògica intuïcionista, busca una fórmula que no sigui estable per doble negació.
\end{indicacio}
\end{exercici}

\chapter[Transformacions naturals]{\texorpdfstring{Transformacions naturals\\ i equivalències}{Transformacions naturals i equivalències}}

Els exercicis següents completen els resolts a la part de Pràctica. Cada enunciat va acompanyat d'una indicació breu; convé intentar-lo primer sense mirar-la. La marca \enquote{$\star$} assenyala un exercici de consolidació i \enquote{$\star\star$} un d'ampliació.

\section{Definició de transformació natural}

\begin{exercici}
Siguin $F,G,H\colon\cat{C}\to\cat{D}$ tres functors i $\eta\colon F\Rightarrow G$, $\theta\colon G\Rightarrow H$ dues transformacions naturals. Comprova directament, a partir de la definició, que la composició vertical $\theta\comp\eta$ satisfà la condició de naturalitat. \hfill{\normalfont$\star$}
\begin{indicacio}
Encadena els quadrats de naturalitat de $\eta$ i de $\theta$: $H(f)\comp\theta_A\comp\eta_A=\theta_B\comp G(f)\comp\eta_A=\theta_B\comp\eta_B\comp F(f)$.
\end{indicacio}
\end{exercici}

\begin{exercici}
Sigui $\eta\colon F\Rightarrow G$ una transformació natural i $K\colon\cat{D}\to\cat{E}$ un functor. Comprova que els morfismes $K(\eta_A)$ són les components d'una transformació natural $K\eta\colon K\comp F\Rightarrow K\comp G$ (composició d'una transformació natural amb un functor per l'esquerra). \hfill{\normalfont$\star$}
\begin{indicacio}
Aplica $K$ al quadrat de naturalitat de $\eta$ i usa la functorialitat de $K$: de $G(f)\comp\eta_A=\eta_B\comp F(f)$ dedueix $K(G(f))\comp K(\eta_A)=K(\eta_B)\comp K(F(f))$.
\end{indicacio}
\end{exercici}

\begin{exercici}
Siguin $\eta\colon F\Rightarrow G$, amb $F,G\colon\cat{C}\to\cat{D}$, i $\theta\colon K\Rightarrow L$, amb $K,L\colon\cat{D}\to\cat{E}$. Comprova que la composició horitzontal amb una transformació identitat es redueix a la composició amb un functor:
\[
\theta\ast\id_F=\theta F,
\qquad\qquad
\id_K\ast\eta=K\eta.
\]
Dedueix que la composició horitzontal generalitza les dues composicions d'una transformació natural amb un functor. \hfill{\normalfont$\star$}
\begin{indicacio}
Calcula les components amb la definició, fent servir l'expressió de $(\theta\ast\eta)_A$ que més convingui a cada cas, i recorda que un functor envia identitats a identitats.
\end{indicacio}
\end{exercici}

\begin{exercici}
Siguin $\alpha\colon F\Rightarrow G$ i $\beta\colon G\Rightarrow H$, amb $F,G,H\colon\cat{C}\to\cat{D}$, i $\sigma\colon K\Rightarrow L$ i $\tau\colon L\Rightarrow M$, amb $K,L,M\colon\cat{D}\to\cat{E}$. Escriu el tipus de cada costat i demostra la \textbf{llei d'intercanvi}
\[
(\tau\comp\sigma)\ast(\beta\comp\alpha)=(\tau\ast\beta)\comp(\sigma\ast\alpha)\colon K\comp F\Rightarrow M\comp H.
\]
\hfill{\normalfont$\star\star$}
\begin{indicacio}
Calcula les components en un objecte $A$ de les dues bandes i fes servir la naturalitat de $\sigma$ aplicada al morfisme $\beta_A$.
\end{indicacio}
\end{exercici}

\section{Exemples de transformacions naturals}

\begin{exercici}
En l'exemple sintaxi/semàntica de la secció~\ref{sec:cap-transformacio}, escriu explícitament l'equació de naturalitat de la interpretació $\eta$ per a un canvi de variables $f\colon X\to Y$, i explica en paraules què diuen els dos costats. Comprova-la per a la fórmula $\varphi=p\lor\lnot q$, amb $X=\{p,q\}$, en dos casos:
\begin{enumerate}
  \item el reanomenament $f\colon p\mapsto r,\ q\mapsto s$, amb $Y=\{r,s\}$;
  \item la substitució $f\colon p\mapsto r,\ q\mapsto r$, amb $Y=\{r\}$, que identifica les dues variables.
\end{enumerate}
En el segon cas, $\varphi$ no és una tautologia, però $\varphi[f]$ sí que ho és. Explica per què això no contradiu la naturalitat. \hfill{\normalfont$\star$}
\begin{indicacio}
Avalua tots dos camins del quadrat sobre una valoració qualsevol $w\colon Y\to\{0,1\}$: l'un passa per $w\models\varphi[f]$ i l'altre per $(w\comp f)\models\varphi$. Al segon cas, observa quines valoracions de $X$ són de la forma $w\comp f$.
\end{indicacio}
\end{exercici}

\begin{exercici}
Sigui $\cat{C}$ una categoria localment petita i siguin $X,Y$ dos objectes seus. Comprova que cada morfisme $h\colon Y\to X$ indueix una transformació natural entre els homfunctors covariants
\[
\Hom(X,-)\Rightarrow\Hom(Y,-).
\]
\hfill{\normalfont$\star$}
\begin{indicacio}
Defineix la component a $A$ per precomposició amb $h$ i usa l'associativitat per comprovar la naturalitat.
\end{indicacio}
\end{exercici}

\begin{exercici}
Considera el functor \enquote{llista} $L\colon\cat{Set}\to\cat{Set}$, que envia cada conjunt $X$ al conjunt $L(X)$ de les llistes finites d'elements de $X$. Comprova que l'aplicació que envia cada element $x$ a la llista d'un sol element $[x]$ defineix una transformació natural $\id_{\cat{Set}}\Rightarrow L$. \hfill{\normalfont$\star\star$}
\begin{indicacio}
El component a $X$ és $\eta_X\colon X\to L(X)$, $x\mapsto[x]$. La naturalitat respecte d'una aplicació $f\colon X\to Y$ diu que aplicar $f$ i després formar la llista d'un element coincideix amb formar la llista i després aplicar $f$ terme a terme.
\end{indicacio}
\end{exercici}

\section{Isomorfismes naturals}

\begin{exercici}
Considera el functor d'\textbf{intercanvi} $\sigma_{\cat{C},\cat{D}}\colon\cat{C}\times\cat{D}\to\cat{D}\times\cat{C}$ definit, sobre objectes i sobre morfismes, per
\[
\sigma_{\cat{C},\cat{D}}(A,B)=(B,A),
\qquad
\sigma_{\cat{C},\cat{D}}(f,g)=(g,f).
\]
Comprova que és un functor i que la seva composició amb l'intercanvi en sentit contrari,
\[
\sigma_{\cat{D},\cat{C}}\comp\sigma_{\cat{C},\cat{D}}=\id_{\cat{C}\times\cat{D}},
\]
és \emph{estrictament} la identitat (i igualment en l'altre sentit). Dedueix que $\sigma_{\cat{C},\cat{D}}$ és un \emph{isomorfisme de categories}, no només una equivalència, i explica per què aquí no cal passar per un isomorfisme natural com $\id\cong\sigma\comp\sigma$. \hfill{\normalfont$\star$}
\begin{indicacio}
Comprova la functorialitat component a component i calcula $\sigma_{\cat{D},\cat{C}}\comp\sigma_{\cat{C},\cat{D}}$ sobre objectes i sobre morfismes. Per a l'última pregunta, compara la definició d'isomorfisme de categories amb la d'equivalència.
\end{indicacio}
\end{exercici}

\begin{exercici}
Sigui $\cat{C}$ una categoria amb \emph{almenys un objecte}. Comprova que dos functors constants $\Delta_D,\Delta_{D'}\colon\cat{C}\to\cat{D}$ (amb valors $D$ i $D'$) són naturalment isomorfs si i només si $D\cong D'$ a $\cat{D}$. Explica què falla si $\cat{C}$ és buida. \hfill{\normalfont$\star\star$}
\begin{indicacio}
Per a ($\Rightarrow$), mira què és una component en un objecte qualsevol de $\cat{C}$. Per a ($\Leftarrow$), prova amb una família constant i escriu-ne el quadrat de naturalitat. Per al cas buit, pensa quantes transformacions naturals hi ha entre dos functors amb domini buit.
\end{indicacio}
\end{exercici}

\begin{exercici}
Sigui $\mathbb{Z}/2\mathbb{Z}=\{1,g\}$ (amb $g^2=1$) vist com a categoria d'un sol objecte $\cat{B}(\mathbb{Z}/2\mathbb{Z})$, i considera dos functors $F,G\colon\cat{B}(\mathbb{Z}/2\mathbb{Z})\to\cat{Set}$ que envien l'únic objecte al conjunt $\{0,1,2\}$. En $F$, el generador $g$ actua per la identitat; en $G$, actua per l'aplicació $\sigma$ que intercanvia $0$ i $1$ i deixa fix el $2$.
\begin{enumerate}
  \item Troba totes les transformacions naturals $F\Rightarrow G$.
  \item Dedueix que $F$ i $G$ no són naturalment isomorfs, tot i que $F(\ast)=G(\ast)$.
  \item Hi ha transformacions naturals $G\Rightarrow F$? N'hi ha alguna de bijectiva?
\end{enumerate}
Compara el resultat amb l'exemple de l'observació~\ref{obs:iso-no-natural}, on no hi havia \emph{cap} transformació natural. \hfill{\normalfont$\star\star$}
\begin{indicacio}
Escriu la condició de naturalitat per al generador $g$ i estudia si pot existir una component bijectiva que la satisfaci.
\end{indicacio}
\end{exercici}

\section{La categoria de functors}

\begin{exercici}
Comprova que si $\cat{C}$ és una categoria petita i $\cat{D}$ és localment petita, aleshores la categoria de functors $\cat{D}^{\cat{C}}$ és localment petita. \hfill{\normalfont$\star$}
\begin{indicacio}
Una transformació natural està determinada pels seus components, un per a cada objecte de $\cat{C}$; com que $\cat{C}$ és petita i cada $\Hom_{\cat{D}}(F(A),G(A))$ és un conjunt, la col·lecció de transformacions naturals $F\Rightarrow G$ és un subconjunt d'un producte de conjunts indexat per un conjunt.
\end{indicacio}
\end{exercici}

\begin{exercici}
Descriu la categoria de functors $\cat{Set}^{\cat{2}}$ i comprova que els seus isomorfismes són els quadrats commutatius amb les dues fletxes verticals bijectives. \hfill{\normalfont$\star\star$}
\begin{indicacio}
Un objecte és una aplicació $f\colon A\to B$; un morfisme $(u,v)$ és un quadrat commutatiu. Un isomorfisme a la categoria de fletxes és un quadrat amb $u$ i $v$ bijectives, ja que la inversa s'obté invertint-les.
\end{indicacio}
\end{exercici}

\section{Equivalència de categories}

\begin{exercici}
Comprova que \enquote{ser equivalents} és una relació d'equivalència entre categories. És a dir, comprova que tota categoria és equivalent a ella mateixa, que la relació és simètrica i que la composició de dues equivalències és una equivalència. \hfill{\normalfont$\star$}
\begin{indicacio}
Per a la reflexivitat, pren el functor identitat. Per a la simetria, intercanvia els papers de $F$ i $G$. Per a la transitivitat, comprova que la composició de functors plens, fidels i essencialment exhaustius torna a tenir les tres propietats, o compon directament els functors i els isomorfismes naturals.
\end{indicacio}
\end{exercici}

\begin{exercici}
Sigui $\cat{C}$ una categoria en què tot objecte és isomorf a un objecte fix $D_0$ (és a dir, hi ha un sol objecte llevat d'isomorfisme; \emph{no} se suposa que tots els morfismes siguin isomorfismes). Comprova que $\cat{C}$ és equivalent al monoide $\mathrm{End}(D_0)=\Hom(D_0,D_0)$ vist com a categoria d'un sol objecte. \hfill{\normalfont$\star\star$}
\begin{indicacio}
Considera la subcategoria plena de $\cat{C}$ que té $D_0$ com a únic objecte, i fes servir la caracterització de les equivalències. Per què no cal cap hipòtesi sobre la invertibilitat dels morfismes?
\end{indicacio}
\end{exercici}

\begin{exercici}
Dona un exemple de dues categories que siguin \emph{equivalents} però no \emph{isomorfes}, i explica per què no són isomorfes. \hfill{\normalfont$\star\star$}
\begin{indicacio}
Un isomorfisme de categories indueix una bijecció entre els objectes, mentre que una equivalència no. Busca dues categories equivalents amb un nombre diferent d'objectes; n'hi ha un exemple amb categories molt petites.
\end{indicacio}
\end{exercici}

\ifpublicacioparcial\else
\chapter[Representabilitat i Yoneda]{\texorpdfstring{Propietats universals,\\ representabilitat i Yoneda}{Propietats universals, representabilitat i Yoneda}}

Els exercicis següents completen els resolts a la part de Pràctica. Cada enunciat va acompanyat d'una indicació breu. La marca \enquote{$\star$} assenyala un exercici de consolidació i \enquote{$\star\star$} un d'ampliació.

\section{Objectes inicials i terminals}

\begin{exercici}
Comprova que si una categoria té objecte inicial, aquest és únic llevat d'isomorfisme únic, repetint l'argument fet per als objectes terminals. \hfill{\normalfont$\star$}
\begin{indicacio}
Dualitza l'argument dels objectes terminals: compara les composicions de les úniques fletxes entre dos objectes inicials amb les identitats.
\end{indicacio}
\end{exercici}

\begin{exercici}
Determina els objectes inicial i terminal a $\cat{Set}$, a $\cat{Pos}$ i en un preordre $(P,\leq)$ vist com a categoria. \hfill{\normalfont$\star$}
\begin{indicacio}
Per a cada categoria, pregunta't quins objectes reben exactament una fletxa des de cada objecte, i quins n'envien exactament una cap a cada objecte. En un preordre, tradueix-ho en termes de l'ordre.
\end{indicacio}
\end{exercici}

\begin{exercici}[Ampliació, per a lectors amb àlgebra]
Comprova que la categoria $\cat{Field}$ dels cossos i homomorfismes de cossos \emph{no} té objecte inicial ni terminal. \hfill{\normalfont$\star$}
\begin{indicacio}
Un homomorfisme de cossos preserva $0$, $1$ i les operacions, i és sempre injectiu. No hi ha cap cos que rebi un morfisme de tots els altres (pensa en característiques diferents), ni cap que n'enviï a tots.
\end{indicacio}
\end{exercici}

\section{Propietats universals com a representabilitat}

\begin{exercici}
Comprova que un objecte $0$ és inicial si i només si el functor covariant $\Hom(0,-)$ és naturalment isomorf al functor constant $\Delta_{\{\ast\}}\colon\cat{C}\to\cat{Set}$ (exemple~\ref{ex:functor-constant}). \hfill{\normalfont$\star$}
\begin{indicacio}
Compara el nombre d'elements de cada $\Hom(0,X)$ amb el de $\{\ast\}$, i pregunta't quines condicions de naturalitat queden per comprovar.
\end{indicacio}
\end{exercici}

\begin{exercici}
Sigui $\cat{C}$ localment petita amb objecte terminal $1$. Comprova que $\Hom(1,-)\colon\cat{C}\to\cat{Set}$ és un functor i demostra que envia l'objecte terminal a un conjunt terminal. Mostra amb un exemple que, tot i això, $\Hom(1,-)$ pot no distingir objectes no isomorfs: pot enviar-los a conjunts amb el mateix nombre d'elements, o fins i tot al mateix conjunt. \hfill{\normalfont$\star$}
\begin{indicacio}
Calcula $\Hom(1,1)$. Per a l'exemple, pensa en un preordre amb un element màxim i dos elements incomparables per sota.
\end{indicacio}
\end{exercici}

\begin{exercici}
Considera el prefeix de parts $\mathcal P\colon\cat{Set}\op\to\cat{Set}$ i les dues parelles candidates $(\{0,1\},\{1\})$ i $(\{0,1\},\varnothing)$. Quina de les dues és una representació de $\mathcal P$? Justifica per què l'altra no ho és. \hfill{\normalfont$\star$}
\begin{indicacio}
Per a la segona parella, intenta obtenir el subconjunt $\{\ast\}$ d'un conjunt d'un sol element com una antiimatge de $\varnothing$.
\end{indicacio}
\end{exercici}

\section{El lema de Yoneda}

\begin{exercici}
Escriu l'enunciat covariant del lema de Yoneda: per a un functor covariant $Q\colon\cat{C}\to\cat{Set}$ i un objecte $A$, hi ha una bijecció natural
\[
\Nat\bigl(\Hom(A,-),Q\bigr)\;\cong\;Q(A).
\]
Identifica la bijecció. \hfill{\normalfont$\star\star$}
\begin{indicacio}
Reprodueix la demostració del cas contravariant a $\cat{C}\op$ i avalua la transformació en $\id_A$.
\end{indicacio}
\end{exercici}

\begin{exercici}
Fes servir el lema de Yoneda per calcular
\[
\Nat\bigl(\Hom(-,A),\Hom(-,B)\bigr).
\]
Comprova que la bijecció resultant és precisament l'acció de la immersió de Yoneda sobre els morfismes $A\to B$. \hfill{\normalfont$\star\star$}
\begin{indicacio}
Aplica Yoneda al prefeix representable $P=\Hom(-,B)$ i identifica què correspon a un morfisme $A\to B$.
\end{indicacio}
\end{exercici}

\begin{exercici}
Siguin $(A,u)$ i $(B,v)$ dues representacions d'un mateix prefeix $P$. Demostra que hi ha un únic isomorfisme $f\colon A\to B$ amb $P(f)(v)=u$. Explica per què això no diu que $A$ tingui un únic automorfisme. \hfill{\normalfont$\star$}
\begin{indicacio}
Fes servir la universalitat de $v$ per obtenir $f$, i la de $u$ per obtenir $g\colon B\to A$ amb $P(g)(u)=v$. Aplica després la unicitat a $g\comp f$ i a $\id_A$.
\end{indicacio}
\end{exercici}

\begin{exercici}
Comprova que si un prefeix $P$ és representable, aleshores l'element $x\in P(A)$ que correspon a la identitat sota l'isomorfisme $\Hom(-,A)\cong P$ té la propietat universal següent: per a cada objecte $X$ i cada $y\in P(X)$, hi ha un únic morfisme $\varphi\colon X\to A$ amb $P(\varphi)(x)=y$. (Aquest $x$ s'anomena \emph{element universal}.) \hfill{\normalfont$\star$}
\begin{indicacio}
Escriu la naturalitat de l'isomorfisme $\Hom(-,A)\cong P$ per a un morfisme $\varphi\colon X\to A$, avaluada en $\id_A$.
\end{indicacio}
\end{exercici}

\section{Conseqüències}

\begin{exercici}
Sigui $\cat{C}$ una categoria petita. Comprova que la immersió de Yoneda $\yon\colon\cat{C}\to\cat{Set}^{\cat{C}\op}$ \emph{no} és, en general, essencialment exhaustiva, tot i ser plena i fidel. Quins prefeixos són de la forma $\Hom(-,A)$? \hfill{\normalfont$\star\star$}
\begin{indicacio}
Busca una propietat que tot prefeix representable $\Hom(-,A)$ compleix necessàriament en l'objecte $A$, i un prefeix que no la compleixi.
\end{indicacio}
\end{exercici}

\begin{exercici}
Dedueix del corol·lari de Yoneda que si dos objectes tenen conjunts de morfismes \enquote{naturalment iguals} des de tot objecte, aleshores són isomorfs; i explica per què això formalitza l'eslògan \enquote{un objecte es coneix per les fletxes que hi arriben}. \hfill{\normalfont$\star\star$}
\begin{indicacio}
Tradueix \enquote{naturalment iguals des de tot objecte} en termes de prefeixos representables.
\end{indicacio}
\end{exercici}

\begin{exercici}
Enuncia i comprova la versió dual del corol·lari: $A\cong B$ si i només si $\Hom(A,-)\cong\Hom(B,-)$ com a functors covariants. \hfill{\normalfont$\star\star$}
\begin{indicacio}
Aplica el corol·lari a $\cat{C}\op$, o repeteix l'argument amb la immersió covariant $\cat{C}\op\to\cat{Set}^{\cat{C}}$, que també és plena i fidel.
\end{indicacio}
\end{exercici}

\chapter[Límits i colímits]{Límits i colímits}

Els exercicis següents completen els resolts a la part de Pràctica. Cada enunciat va acompanyat d'una indicació breu. La marca \enquote{$\star$} assenyala un exercici de consolidació i \enquote{$\star\star$} un d'ampliació.

\section{Diagrames i cons}

\begin{exercici}
Descriu la categoria índex $\cat{I}$ i el diagrama corresponent per a cadascun d'aquests límits: producte binari, equalitzador, pullback i objecte terminal. \hfill{\normalfont$\star$}
\begin{indicacio}
Identifica, per a cada propietat universal, la forma mínima de diagrama que la descriu: quants objectes hi intervenen i quines fletxes cal que commutin.
\end{indicacio}
\end{exercici}

\begin{exercici}
Comprova que un morfisme de cons és un cas particular de morfisme a la categoria $\mathrm{Con}(D)$, i que la composició de morfismes de cons torna a ser-ne un. \hfill{\normalfont$\star$}
\begin{indicacio}
Si $f\colon A\to B$ i $g\colon B\to C$ satisfan $\beta_i\comp f=\alpha_i$ i $\gamma_i\comp g=\beta_i$, aleshores $\gamma_i\comp(g\comp f)=\alpha_i$.
\end{indicacio}
\end{exercici}

\section{Límits}

\begin{exercici}
Comprova que en un preordre vist com a categoria, el producte de dos elements és el seu ínfim i el coproducte és el seu suprem, quan existeixen. \hfill{\normalfont$\star$}
\begin{indicacio}
Tradueix la propietat universal del producte a desigualtats.
\end{indicacio}
\end{exercici}

\begin{exercici}
Comprova que dos límits d'un mateix diagrama són isomorfs per un únic isomorfisme compatible amb totes les projeccions. \hfill{\normalfont$\star$}
\begin{indicacio}
Treballa a la categoria $\mathrm{Con}(D)$ i fes servir la unicitat dels objectes terminals.
\end{indicacio}
\end{exercici}

\section{Productes i equalitzadors}

\begin{exercici}
Comprova, llevat d'isomorfisme, la commutativitat i l'associativitat del producte:
\[
A\times B\cong B\times A,\quad (A\times B)\times C\cong A\times(B\times C).
\] \hfill{\normalfont$\star$}
\begin{indicacio}
Fes servir la unicitat d'un objecte definit per una mateixa propietat universal.
\end{indicacio}
\end{exercici}

\begin{exercici}
Comprova que si $\cat{C}$ té objecte terminal $1$, aleshores $A\times 1\cong A$ per a tot objecte $A$. \hfill{\normalfont$\star$}
\begin{indicacio}
Compara les propietats universals de $A\times 1$ i de $A$.
\end{indicacio}
\end{exercici}

\begin{exercici}
Comprova que a $\cat{Grp}$ l'equalitzador de $f,g\colon G\to H$ és el subgrup $\{x\in G\mid f(x)=g(x)\}$. \hfill{\normalfont$\star$}
\begin{indicacio}
Comprova que aquest subconjunt és un subgrup i que la injecció satisfà la propietat universal de l'equalitzador; el functor oblit cap a $\cat{Set}$ ajuda a veure-ho.
\end{indicacio}
\end{exercici}

\section{Pullbacks}

\begin{exercici}
Comprova que en un pullback, si $g\colon B\to C$ és un monomorfisme, aleshores $p_1\colon P\to A$ també ho és. (Es diu que els monomorfismes són \emph{estables per pullback}.) \hfill{\normalfont$\star\star$}
\begin{indicacio}
Usa la propietat universal del pullback i la cancel·lació de $g$. Aquesta estabilitat serà essencial per a la teoria de subobjectes.
\end{indicacio}
\end{exercici}

\begin{exercici}
Comprova que un quadrat amb $\id$ als dos costats,
\[
\begin{array}{ccc}
A & \xrightarrow{\;\id\;} & A\\[-2pt]
\downarrow{\scriptstyle\id} & & \downarrow{\scriptstyle f}\\[-2pt]
A & \xrightarrow{\;f\;} & B
\end{array}
\]
és un pullback si i només si $f$ és un monomorfisme. \hfill{\normalfont$\star\star$}
\begin{indicacio}
El pullback d'aquest quadrat a $\cat{Set}$ és $\{(a,a')\mid f(a)=f(a')\}$; que la projecció sigui iso equival a la injectivitat de $f$. En general, tradueix-ho amb la propietat universal.
\end{indicacio}
\end{exercici}

\section{Colímits}

\begin{exercici}
Comprova que tot coequalitzador és un epimorfisme, dualitzant la demostració que tot equalitzador és un monomorfisme. \hfill{\normalfont$\star$}
\begin{indicacio}
Dualitza l'argument de la Pràctica que prova que tot equalitzador és mono.
\end{indicacio}
\end{exercici}

\begin{exercici}
Comprova que el pushout de $f\colon C\to A$ i $g\colon C\to B$ a $\cat{Set}$ és $(A+B)/{\sim}$, on $\sim$ enganxa $f(c)$ amb $g(c)$ per a cada $c\in C$. \hfill{\normalfont$\star$}
\begin{indicacio}
És el dual del pullback. Pren la unió disjunta $A+B$ i quocienta per la relació generada per $\iota_1(f(c))\sim\iota_2(g(c))$; comprova la propietat universal.
\end{indicacio}
\end{exercici}

\section{Límits finits}

\begin{exercici}
Sigui $\cat{C}$ una categoria amb productes binaris i pullbacks, i siguin $f,g\colon A\to B$. La demostració del teorema~\ref{teo:limits-finits} obté l'equalitzador de $f$ i $g$ com un pullback de dues fletxes $A\to A\times B$. Comprova que també és un pullback de la diagonal $\langle\id_B,\id_B\rangle\colon B\to B\times B$ al llarg de $\langle f,g\rangle\colon A\to B\times B$. \hfill{\normalfont$\star\star$}
\begin{indicacio}
Escriu què és un con sobre aquest pullback i tradueix la condició de con component a component, amb les dues projeccions de $B\times B$.
\end{indicacio}
\end{exercici}

\begin{exercici}
Comprova directament que una categoria amb productes binaris i equalitzadors té tots els pullbacks, amb una construcció més econòmica que la general de la demostració del teorema~\ref{teo:limits-finits}. \hfill{\normalfont$\star$}
\begin{indicacio}
Expressa el pullback com un equalitzador de dues fletxes que surten d'un producte.
\end{indicacio}
\end{exercici}

\section{Límits en categories de functors}

\begin{exercici}
Sigui $\mathbf E$ el graf amb dos vèrtexs $0,1$ i una sola aresta $e\colon 0\to 1$. Fent servir que $\cat{Graph}\cong\cat{Set}^{\cat{P}}$ i que els límits i colímits s'hi calculen punt a punt (proposició~\ref{prop:limits-punt-a-punt}), descriu el producte $\mathbf E\times\mathbf E$ i el coproducte $\mathbf E+\mathbf E$: els seus vèrtexs, les seves arestes, i l'origen i el destí de cada aresta. Dibuixa'ls. És $\mathbf E\times\mathbf E$ un quadrat? \hfill{\normalfont$\star$}
\begin{indicacio}
Calcula per separat el conjunt de vèrtexs i el d'arestes, i defineix origen i destí component a component.
\end{indicacio}
\end{exercici}

\section{Preservació i reflexió}

\begin{exercici}
Comprova que tot functor ple i fidel reflecteix límits. Dualment, reflecteix colímits. \hfill{\normalfont$\star\star$}
\begin{indicacio}
Aixeca l'existència de la factorització amb la plenitud i prova'n la unicitat amb la fidelitat.
\end{indicacio}
\end{exercici}

\begin{exercici}
Comprova que el prefeix representable
\[
\Hom(-,B)\colon\cat{C}\op\to\cat{Set}
\]
envia colímits de $\cat{C}$ a límits de $\cat{Set}$. \hfill{\normalfont$\star\star$}
\begin{indicacio}
És el dual del fet que $\Hom(A,-)$ preserva límits: un cocon amb vèrtex $B$ sobre $D$ és un con sobre $\Hom(D_-,B)$, i la universalitat es correspon.
\end{indicacio}
\end{exercici}

\chapter[Adjuncions]{Adjuncions}

Els exercicis següents completen els resolts a la part de Pràctica. Cada enunciat va acompanyat d'una indicació breu. La marca \enquote{$\star$} assenyala un exercici de consolidació i \enquote{$\star\star$} un d'ampliació.

\section{Adjuncions entre preordres}

\begin{exercici}
Comprova que en una adjunció entre preordres $f\dashv g$, l'aplicació $f$ preserva els suprems que existeixin i $g$ preserva els ínfims. \hfill{\normalfont$\star$}
\begin{indicacio}
És el cas particular, per a preordres, que els adjunts per l'esquerra preserven colímits (suprems) i els adjunts per la dreta preserven límits (ínfims). Demostra-ho directament amb la definició $f(x)\leq y\Leftrightarrow x\leq g(y)$.
\end{indicacio}
\end{exercici}

\begin{exercici}
Sigui $f\colon P\to Q$ monòtona entre \textbf{reticles complets}, és a dir, preordres en què tot subconjunt té suprem i ínfim. Comprova que $f$ té adjunt per la dreta si i només si preserva tots els suprems. \hfill{\normalfont$\star\star$}
\begin{indicacio}
Si $f$ preserva suprems, defineix $g(y)$ com el suprem d'un conjunt adequat d'elements de $P$.
\end{indicacio}
\end{exercici}

\section{Adjuncions entre categories}

\begin{exercici}\label{exr:adj-unitat-universal}
Comprova que si $F\dashv G$, aleshores per a cada objecte $A$ el morfisme $\eta_A\colon A\to G(F(A))$ és universal: tot $\psi\colon A\to G(B)$ factoritza de manera única com $G(\varphi)\comp\eta_A$. \hfill{\normalfont$\star\star$}
\begin{indicacio}
Fes servir la bijecció $\theta$ i la fórmula del transposat en termes de la unitat.
\end{indicacio}
\end{exercici}

\begin{exercici}
Sigui $\cat{C}$ una categoria amb coproductes binaris i $\Delta\colon\cat{C}\to\cat{C}\times\cat{C}$ el functor diagonal, $A\mapsto(A,A)$. Dualitzant l'exercici resolt de la Pràctica sobre $\Delta\dashv\times$, comprova que el coproducte és adjunt per l'esquerra de $\Delta$, $(+)\dashv\Delta$, i descriu-ne la unitat i la counitat. \hfill{\normalfont$\star\star$}
\begin{indicacio}
Fes servir la propietat universal del coproducte; la unitat i la counitat són les duals de la diagonal i les projeccions.
\end{indicacio}
\end{exercici}

\section{Unitat i counitat}

\begin{exercici}\label{exr:adj-fletxes-universals}
Sigui $G\colon\cat{D}\to\cat{C}$ un functor. Suposem que, per a cada objecte $A$ de $\cat{C}$, hem triat un objecte $F(A)$ de $\cat{D}$ i una fletxa universal $\eta_A\colon A\to G(F(A))$, en el sentit que tot $\psi\colon A\to G(B)$ factoritza de manera única com $G(\varphi)\comp\eta_A$. Demostra que aquestes dades determinen de manera única l'acció de $F$ sobre morfismes i defineixen un adjunt per l'esquerra $F\dashv G$. \hfill{\normalfont$\star\star$}
\begin{indicacio}
Defineix $F$ sobre objectes i sobre morfismes amb la propietat universal, i fes servir la unicitat per obtenir la functorialitat.
\end{indicacio}
\end{exercici}

\begin{exercici}
Per a l'adjunció $\Free\dashv U$ entre $\cat{Graph}$ i $\cat{Cat}$ (exemple~\ref{ex:free-u-adjuncio}), comprova que la counitat
\[
\varepsilon_{\cat{C}}\colon\Free(U(\cat{C}))\longrightarrow\cat{C}
\]
és la identitat sobre els objectes i és un functor ple. Mostra amb un exemple que, en general, no és fidel. \hfill{\normalfont$\star$}
\begin{indicacio}
Per veure que és ple, pensa en els camins de longitud~$1$. Per veure que no és fidel, busca dos camins diferents amb la mateixa composta.
\end{indicacio}
\end{exercici}

\section{Identitats triangulars}

\begin{exercici}
Comprova la segona identitat triangular,
\[
G(\varepsilon_B)\comp\eta_{G(B)}=\id_{G(B)},
\]
dualitzant la demostració de la primera. \hfill{\normalfont$\star$}
\begin{indicacio}
Dualitza la demostració de la Teoria, fent servir la fórmula del transposat en termes de la unitat.
\end{indicacio}
\end{exercici}

\begin{exercici}\label{exr:adj-triangulars-reciproc}
Comprova que dues transformacions naturals $\eta\colon\id\Rightarrow GF$ i $\varepsilon\colon FG\Rightarrow\id$ que satisfan les identitats triangulars defineixen una adjunció $F\dashv G$. \hfill{\normalfont$\star\star$}
\begin{indicacio}
Defineix els dos transposats amb $\eta$ i $\varepsilon$, i fes servir les identitats triangulars per veure que són inversos.
\end{indicacio}
\end{exercici}

\begin{exercici}
Siguin $F\colon\cat{C}\to\cat{D}$, $G\colon\cat{D}\to\cat{C}$, $F'\colon\cat{D}\to\cat{E}$ i $G'\colon\cat{E}\to\cat{D}$ functors entre categories localment petites, amb $F\dashv G$ i $F'\dashv G'$. Demostra que les adjuncions es componen:
\[
F'\comp F\;\dashv\;G\comp G'.
\]
Descriu la unitat de l'adjunció composta en termes de les unitats $\eta$ i $\eta'$. \hfill{\normalfont$\star$}
\begin{indicacio}
Encadena les dues bijeccions d'homconjunts, passant per $\cat{D}$, i comprova que la composta és natural en les dues variables. Per a la unitat, transposa la identitat de $F'(F(A))$.
\end{indicacio}
\end{exercici}

\begin{exercici}
Sigui $F\dashv G$ una adjunció amb unitat $\eta$ i counitat $\varepsilon$. Demostra que, si $\eta$ i $\varepsilon$ són isomorfismes naturals, aleshores $F$ és una equivalència de categories amb quasiinvers $G$ (definició~\ref{def:equivalencia}), i que en aquest cas $G$ també és adjunt per l'esquerra de $F$. \hfill{\normalfont$\star$}
\begin{indicacio}
Compara les dades amb la definició d'equivalència. Per a la segona part, inverteix $\eta$ i $\varepsilon$ i comprova les identitats triangulars de l'adjunció $G\dashv F$ a partir de les de $F\dashv G$.
\end{indicacio}
\end{exercici}

\section{Els adjunts i els límits}

\begin{exercici}
Comprova que el prefeix de parts $\mathcal P\colon\cat{Set}\op\to\cat{Set}$ envia coproductes a productes: $\mathcal P(A+B)\cong\mathcal P(A)\times\mathcal P(B)$. \hfill{\normalfont$\star$}
\begin{indicacio}
Descriu un subconjunt de $A+B$ a partir de les seves interseccions amb les dues components.
\end{indicacio}
\end{exercici}

\begin{exercici}[Ampliació, per a lectors amb topologia]
Sigui $U\colon\cat{Top}\to\cat{Set}$ el functor oblit. Comprova que té \emph{dos} adjunts, un per l'esquerra i un per la dreta, i digues quins són. \hfill{\normalfont$\star$}
\begin{indicacio}
Pensa en les dues topologies extremes sobre un conjunt.
\end{indicacio}
\end{exercici}

\section{Exponencials}

\begin{exercici}
Comprova que a $\cat{Set}$ es compleix $C^{A+B}\cong C^A\times C^B$ i $(C\times D)^A\cong C^A\times D^A$, i interpreta-ho en termes d'adjuncions. \hfill{\normalfont$\star$}
\begin{indicacio}
Per a la primera, fes servir que $\Hom(A+B,C)\cong\Hom(A,C)\times\Hom(B,C)$, és a dir, que el representable contravariant $\Hom(-,C)$ envia coproductes a productes. Per a la segona, fes servir que $(-)^A$ és adjunt per la dreta de $-\times A$ i, per tant, preserva productes.
\end{indicacio}
\end{exercici}

\begin{exercici}
Comprova que la currificació $\Hom(A\times B,C)\cong\Hom(A,C^B)$ és natural en les tres variables $A,B,C$. \hfill{\normalfont$\star\star$}
\begin{indicacio}
Comprova la commutativitat amb la precomposició en $A$, la postcomposició en $C$ i l'acció en $B$, aplicant les definicions de $\tilde f$ i $\hat g$.
\end{indicacio}
\end{exercici}

\chapter[Cartesianes tancades]{\texorpdfstring{Categories cartesianes\\ tancades}{Categories cartesianes tancades}}

Els exercicis següents completen els resolts a la part de Pràctica. Cada enunciat va acompanyat d'una indicació breu. La marca \enquote{$\star$} assenyala un exercici de consolidació i \enquote{$\star\star$} un d'ampliació, sensiblement més exigent.

\section{Productes finits}

\begin{exercici}
Comprova que en una categoria cartesiana el producte és associatiu i commutatiu llevat d'isomorfisme: $(A\times B)\times C\cong A\times(B\times C)$ i $A\times B\cong B\times A$. \hfill{\normalfont$\star$}
\begin{indicacio}
Construeix els morfismes en tots dos sentits amb aparellaments de projeccions i comprova que en són inversos usant la unicitat de la propietat universal. No cal calcular amb elements.
\end{indicacio}
\end{exercici}

\begin{exercici}
Demostra que si $\cat{C}$ és cartesiana, aleshores $\cat{C}\op$ té coproductes finits, i descriu-los. \hfill{\normalfont$\star$}
\begin{indicacio}
El producte a $\cat{C}$ és el coproducte a $\cat{C}\op$ per dualització de la propietat universal; el terminal esdevé inicial.
\end{indicacio}
\end{exercici}

\section{Exponencials}

\begin{exercici}
En una CCC, construeix un morfisme canònic \enquote{composició} $C^{B}\times B^{A}\to C^{A}$ i comprova que és associatiu en el sentit adient. \hfill{\normalfont$\star\star$}
\begin{indicacio}
Currifica el morfisme $C^{B}\times B^{A}\times A\to C$ obtingut avaluant primer sobre $A$ i després sobre $B$. L'associativitat es comprova amb la unicitat de la transposada.
\end{indicacio}
\end{exercici}

\begin{exercici}
Comprova que l'assignació $B\mapsto C^{B}$ és \emph{contravariant}: un morfisme $f\colon B'\to B$ indueix $C^{f}\colon C^{B}\to C^{B'}$. \hfill{\normalfont$\star$}
\begin{indicacio}
Currifica
\[
C^{B}\times B'\xrightarrow{\id\times f}C^{B}\times B\xrightarrow{\ev}C.
\]
Això mostra la contravariància en l'exponent $B$; en canvi, per a $B$ fix, el functor $(-)^{B}$ és covariant en $C$.
\end{indicacio}
\end{exercici}

\begin{exercici}
En una CCC amb coproductes, demostra la llei $C^{A+B}\cong C^{A}\times C^{B}$. \hfill{\normalfont$\star\star$}
\begin{indicacio}
Fes servir que $X\times-$ és un adjunt per l'esquerra i, per tant, preserva coproductes (exercici resolt~\ref{pr-distributivitat} de la Pràctica). Després aplica la propietat universal del coproducte i el lema de Yoneda.
\end{indicacio}
\end{exercici}

\begin{exercici}
Sabem que $\cat{Graph}$ és cartesiana tancada, perquè és una categoria de prefeixos. Fent servir la fórmula de l'exponencial de prefeixos i el lema de Yoneda, demostra que els \emph{vèrtexs} de l'exponencial $H^{G}$ de dos grafs corresponen a totes les aplicacions $V_G\to V_H$ entre els conjunts de vèrtexs, sense cap condició sobre les arestes. \hfill{\normalfont$\star\star$}
\begin{indicacio}
Els vèrtexs de $H^{G}$ es corresponen canònicament amb els morfismes de grafs $\mathbf V\to H^{G}$, on $\mathbf V$ és el graf d'un sol vèrtex. Transposa'ls i calcula el graf $\mathbf V\times G$.
\end{indicacio}
\end{exercici}

\begin{exercici}
Sigui $\cat{2}=(0\to1)$. Fent servir l'exemple de $\cat{Cat}$ de la Teoria amb $\cat{A}=\cat{2}$, comprova que un functor $\cat{2}\times\cat{2}\to\cat{C}$ és el mateix que un quadrat commutatiu de $\cat{C}$, i que la seva currificació $\cat{2}\to\cat{C}^{\cat{2}}$ és el mateix que un morfisme de la categoria de fletxes. Què fa l'avaluació sobre aquestes dades? \hfill{\normalfont$\star$}
\begin{indicacio}
Escriu els quatre objectes i les cinc fletxes no identitat de $\cat{2}\times\cat{2}$. Aplica la lectura \enquote{transformació natural $=$ functor $\cat{A}\times\cat{2}\to\cat{C}$} amb $\cat{A}=\cat{2}$.
\end{indicacio}
\end{exercici}

\section{Models i lògica}

\begin{exercici}
Comprova que tota àlgebra de Boole és una àlgebra de Heyting en què $b\Rightarrow c=\neg b\vee c$, i per tant un preordre cartesià tancat. \hfill{\normalfont$\star$}
\begin{indicacio}
Verifica l'adjunció $a\wedge b\leq c\Leftrightarrow a\leq\neg b\vee c$ usant les lleis booleanes. La lògica clàssica és el cas booleà de la intuïcionista.
\end{indicacio}
\end{exercici}

\begin{exercici}
Sigui $\cat{C}$ una CCC. Demostra que el conjunt de morfismes $1\to C^{B}$ està en bijecció amb el de morfismes $B\to C$ (els \enquote{punts} de l'exponencial són les funcions). \hfill{\normalfont$\star$}
\begin{indicacio}
És l'adjunció amb $A=1$ i la identificació $1\times B\cong B$.
\end{indicacio}
\end{exercici}

\begin{exercici}
Interpreta el combinador $\mathsf{S}$, de tipus $(\sigma\to\tau\to\rho)\to(\sigma\to\tau)\to\sigma\to\rho$, com a morfisme d'una CCC. \hfill{\normalfont$\star\star$}
\begin{indicacio}
Descurrifica successivament fins a un morfisme amb domini un producte de tres exponencials i $\llbracket\sigma\rrbracket$, i construeix-lo amb avaluacions i la diagonal. Compara'l amb la demostració intuïcionista de l'esquema d'axioma corresponent.
\end{indicacio}
\end{exercici}

\begin{exercici}
Al preordre cartesià tancat $\mathbb Z\cup\{+\infty\}$ de la Teoria, comprova el \emph{modus ponens} $b\wedge(b\Rightarrow c)\leq c$ i la llei de currificació $b\Rightarrow(c\Rightarrow d)=(b\wedge c)\Rightarrow d$, calculant els dos costats per casos. \hfill{\normalfont$\star$}
\begin{indicacio}
Fes servir la fórmula explícita de la implicació en una cadena, i distingeix els casos segons l'ordre relatiu de $b$, $c$ i $d$.
\end{indicacio}
\end{exercici}

\begin{exercici}
Afegeix a $\mathbb Z\cup\{+\infty\}$ un element mínim $-\infty$. Comprova que s'obté una àlgebra de Heyting, calcula la pseudocomplementació $\neg b=(b\Rightarrow-\infty)$ i decideix si és una àlgebra de Boole. \hfill{\normalfont$\star\star$}
\begin{indicacio}
En una cadena, els suprems finits són màxims. Per a la segona part, calcula $\neg\neg b$ i $b\vee\neg b$ per a un $b$ que no sigui un dels extrems.
\end{indicacio}
\end{exercici}

\chapter[Subobjectes, imatges i factoritzacions]{\texorpdfstring{Subobjectes, imatges\\ i factoritzacions}{Subobjectes, imatges i factoritzacions}}

Els exercicis següents completen els resolts a la part de Pràctica. Cada enunciat va acompanyat d'una indicació breu. La marca \enquote{$\star$} assenyala un exercici de consolidació i \enquote{$\star\star$} un d'ampliació.

\section{Subobjectes}

\begin{exercici}
Descriu $\Sub(A)$ quan $A$ és un objecte d'un preordre vist com a categoria. \hfill{\normalfont$\star$}
\begin{indicacio}
En un preordre hi ha com a molt un morfisme entre dos objectes, i tots són monos. Quins subobjectes té, doncs, $A$? Compara amb el conjunt d'elements $x\leq A$: quan donen dos d'aquests elements el mateix subobjecte? Què passa si el preordre és un ordre parcial?
\end{indicacio}
\end{exercici}

\begin{exercici}
Comprova que en una categoria amb objecte inicial estricte $0$, la injecció $0\rightarrowtail A$ és el subobjecte mínim de $\Sub(A)$. \hfill{\normalfont$\star$}
\begin{indicacio}
\enquote{Estricte} vol dir que tot morfisme cap a $0$ és un isomorfisme. Comprova primer que $0\to A$ és mono, i després que es factoritza a través de qualsevol subobjecte.
\end{indicacio}
\end{exercici}

\begin{exercici}[Ampliació, per a lectors amb àlgebra]
Demostra que a $\cat{Grp}$ els subobjectes d'un grup $G$ es corresponen exactament amb els seus \textbf{subgrups}. \hfill{\normalfont$\star\star$}
\begin{indicacio}
Els monomorfismes a $\cat{Grp}$ són els homomorfismes injectius. Comprova que dos monos amb codomini $G$ són equivalents si i només si tenen la mateixa imatge, i que les imatges de monos són exactament els subgrups.
\end{indicacio}
\end{exercici}

\begin{exercici}
Comprova que un morfisme de grafs $F\colon G'\to G$ és un monomorfisme de $\cat{Graph}$ si i només si les aplicacions $F_V$ i $F_E$ són injectives. Dedueix la descripció dels subobjectes d'un graf com a subgrafs de l'exemple~\ref{ex:subgrafs}. \hfill{\normalfont$\star$}
\begin{indicacio}
Per a una direcció, prova els morfismes que surten del graf d'un sol vèrtex i del graf d'una sola aresta.
\end{indicacio}
\end{exercici}

\section{Reindexació}

\begin{exercici}
Sigui $f\colon A\to B$. Comprova que $f^{*}$ envia el subobjecte màxim de $\Sub(B)$ al màxim de $\Sub(A)$. \hfill{\normalfont$\star$}
\begin{indicacio}
El pullback de $\id_B$ al llarg de $f$ és un mono equivalent a $\id_A$. Interpreta-ho: substituir dins del predicat \enquote{sempre cert} torna a donar \enquote{sempre cert}.
\end{indicacio}
\end{exercici}

\begin{exercici}
Treballa a $\cat{Set}$. Comprova que si $m\colon S\rightarrowtail B$ és un subobjecte (un subconjunt $S\subseteq B$) i $f\colon A\to B$, aleshores $f^{*}[m]$ és el subobjecte \emph{mínim} (buit) de $A$ exactament quan la imatge de $f$ no talla $S$, és a dir, quan $f(A)\cap S=\varnothing$. \hfill{\normalfont$\star$}
\begin{indicacio}
A $\cat{Set}$, $f^{-1}(S)=\varnothing$ si i només si cap element de la imatge de $f$ és a $S$. Formula-ho amb pullbacks: el pullback de $m$ al llarg de $f$ és l'objecte inicial exactament en aquest cas.
\end{indicacio}
\end{exercici}

\section{Interseccions i estructura}

\begin{exercici}
Demostra que la reindexació $f^{*}$ preserva interseccions: $f^{*}([m]\wedge[n])=f^{*}[m]\wedge f^{*}[n]$. \hfill{\normalfont$\star\star$}
\begin{indicacio}
Tant l'ínfim com la reindexació són pullbacks. Usa que els pullbacks es componen: el pullback d'un pullback sobre $B$ al llarg de $f$ es reorganitza en la intersecció dels pullbacks. És un exercici de manipulació de quadrats cartesians.
\end{indicacio}
\end{exercici}

\begin{exercici}
A $\cat{Set}$, el functor de subobjectes $\Sub\colon\cat{Set}\op\to\cat{Poset}$ (proposició~\ref{prop:functor-sub}) i el functor de parts contravariant $\mathcal P$, amb $\mathcal P(A)$ ordenat per inclusió, envien cada conjunt a ordres parcials isomorfs. Comprova que aquests isomorfismes formen un isomorfisme natural $\Sub\cong\mathcal P$. \hfill{\normalfont$\star$}
\begin{indicacio}
Fes servir la imatge d'un monomorfisme per definir les components, i compara la reindexació amb l'antiimatge.
\end{indicacio}
\end{exercici}


\section{Factorització imatge}

\begin{exercici}[Ampliació, per a lectors amb àlgebra]
Comprova que a $\cat{Grp}$ la factorització imatge d'un homomorfisme $f\colon G\to H$ és $G\twoheadrightarrow \Imatge(f)\hookrightarrow H$, on $\Imatge(f)$ és el subgrup imatge. \hfill{\normalfont$\star$}
\begin{indicacio}
Els epimorfismes regulars a $\cat{Grp}$ són les projeccions al quocient $G\to G/N$; la primera part de $f$ és $G\to G/\ker f\cong\Imatge(f)$. Fes servir el primer teorema d'isomorfia.
\end{indicacio}
\end{exercici}

\begin{exercici}
Demostra que en qualsevol categoria amb factoritzacions imatge, la imatge d'una composició satisfà $\Imatge(g\comp f)\leq\Imatge(g)$. \hfill{\normalfont$\star$}
\begin{indicacio}
$g\comp f$ es factoritza a través de $\Imatge(g)$ (via la factorització de $g$), i la imatge és el menor subobjecte pel qual això passa.
\end{indicacio}
\end{exercici}

\begin{exercici}
Comprova que un epimorfisme regular que és alhora monomorfisme és un isomorfisme, i dedueix que en una categoria regular la factorització imatge d'un mono és trivial. \hfill{\normalfont$\star$}
\begin{indicacio}
Un epi regular és el coequalitzador de la seva parella nucli; si a més és mono, la parella nucli és la diagonal, i el coequalitzador d'una parella trivial és una identitat llevat d'isomorfisme.
\end{indicacio}
\end{exercici}

\begin{exercici}
Siguin $H$ i $K$ dos subgrafs d'un graf $G$, i $F\colon G\to G''$ un morfisme de grafs. Descriu la intersecció $H\wedge K$ a $\Sub(G)$ i la imatge $\exists_F(H)$ a $\Sub(G'')$, i comprova que totes dues es calculen component a component, sobre vèrtexs i sobre arestes. \hfill{\normalfont$\star$}
\begin{indicacio}
Fes servir l'exemple~\ref{ex:subgrafs} i comprova que les operacions component a component donen subconjunts tancats per origen i destí.
\end{indicacio}
\end{exercici}

\section{Regularitat}

\begin{exercici}[Ampliació, per a lectors amb àlgebra homològica]
Demostra que tota categoria abeliana és regular. \hfill{\normalfont$\star\star$}
\begin{indicacio}
En una categoria abeliana tot morfisme es factoritza com epi (sobre la coimatge) seguit de mono (la imatge), i coimatge i imatge coincideixen. L'estabilitat per pullback se segueix de l'exactitud. Pots limitar-te a esbossar per què cada condició de la definició es compleix.
\end{indicacio}
\end{exercici}

\begin{exercici}[Ampliació, per a lectors amb topologia]
Exhibeix una aplicació quocient a $\cat{Top}$ el pullback de la qual no és una aplicació quocient, i dedueix-ne que $\cat{Top}$ no és regular. \hfill{\normalfont$\star\star$}
\begin{indicacio}
A $\cat{Top}$, els epimorfismes regulars són les aplicacions quocient (no s'han de confondre amb els epimorfismes, que són les aplicacions contínues exhaustives). Busca una aplicació quocient el pullback de la qual, al llarg d'una aplicació adequada, deixa de ser quocient.
\end{indicacio}
\end{exercici}

\section{Quantificació}

\begin{exercici}
Demostra que $\exists_f$ és monòtona. Suposa, a més, que $\Sub(A)$ té un mínim $\bot_A$ (el predicat \enquote{fals}); demostra, fent servir l'adjunció $\exists_f\dashv f^{*}$, que $\exists_f(\bot_A)$ és el mínim de $\Sub(B)$. \hfill{\normalfont$\star$}
\begin{indicacio}
Per a la monotonia, recorda que els adjunts entre preordres són monòtons. Per al mínim, escriu la condició de l'adjunció per a $\exists_f(\bot_A)\leq[t]$.
\end{indicacio}
\end{exercici}

\begin{exercici}
Usant que $\exists_f$ és adjunt per l'esquerra, demostra que preserva unions (suprems) de subobjectes quan existeixen. \hfill{\normalfont$\star\star$}
\begin{indicacio}
Els adjunts per l'esquerra preserven tots els colímits que existeixin; en un preordre, els suprems són els colímits. Escriu la cadena de bicondicionals de l'adjunció aplicada a un suprem.
\end{indicacio}
\end{exercici}


\chapter[Classificadors de subobjectes i topos]{\texorpdfstring{Classificadors de subobjectes i introducció als topos}{Classificadors de subobjectes i introducció als topos}}

Els exercicis següents completen els resolts a la part de Pràctica. La marca \enquote{$\star$} assenyala un exercici de consolidació i \enquote{$\star\star$} un d'ampliació.

\section{Classificadors}

\begin{exercici}
Comprova que a $\cat{Set}$ el morfisme característic $\chi_U$ d'un subconjunt determina i és determinat per $U$, de manera que $\Sub(E)\cong\{0,1\}^{E}$. \hfill{\normalfont$\star$}
\begin{indicacio}
Associa a cada subconjunt la seva funció característica i comprova que l'aplicació és una bijecció amb inversa $\chi\mapsto\chi^{-1}(1)$.
\end{indicacio}
\end{exercici}

\begin{exercici}
En un topos, sigui $m\colon U\rightarrowtail E$ un subobjecte i $g\colon E'\to E$ un morfisme. Demostra que el morfisme característic de la reindexació $g^{*}m$ és
\[
\chi_{g^{*}m}=\chi_m\comp g.
\]
Interpreta-ho com la compatibilitat entre la representació $\Sub\cong\Hom(-,\Omega)$ i la substitució. \hfill{\normalfont$\star$}
\begin{indicacio}
Compara els dos pullbacks de $\mathsf{true}$ i fes servir el lema d'enganxament de pullbacks.
\end{indicacio}
\end{exercici}

\begin{exercici}
Sigui $E=\{1,\dots,6\}$ i $U=\{3,6\}$, el predicat \enquote{$x$ és múltiple de $3$}. Calcula $\chi_U\colon E\to\{0,1\}$. Sigui $g\colon\{a,b,c,d\}\to E$ l'aplicació amb $g(a)=3$, $g(b)=4$, $g(c)=6$ i $g(d)=3$. Calcula $\chi_U\comp g$ i identifica el subconjunt que classifica. \hfill{\normalfont$\star$}
\begin{indicacio}
Segueix l'exemple de la Teoria. Com que $g$ no és injectiva, fixa't que dos elements diferents poden tenir la mateixa imatge i que això no afecta la reindexació.
\end{indicacio}
\end{exercici}

\begin{exercici}
Comprova que $\mathsf{true}\colon1\to\Omega$ és sempre un monomorfisme, i que $\Hom(1,\Omega)\cong\Sub(1)$. Quants valors de veritat globals té $\cat{Set}$? \hfill{\normalfont$\star$}
\begin{indicacio}
Per a la primera part, compara dues fletxes qualssevol $X\to1$. Per a la segona, particularitza la representabilitat de $\Sub$.
\end{indicacio}
\end{exercici}

\section{Topos}

\begin{exercici}
Sigui $\mathbf 2=(0\to 1)$. La categoria de fletxes $\cat{Set}^{\mathbf 2}$ té per objectes les aplicacions $X_0\to X_1$, i és una categoria de prefeixos sobre $\mathbf 2\op$, per tant un topos. Fent servir la construcció amb sedassos:
\begin{enumerate}
  \item calcula els conjunts $\Omega(0)$ i $\Omega(1)$;
  \item calcula l'aplicació de restricció entre ells, de manera que $\Omega$ és un objecte de $\cat{Set}^{\mathbf 2}$, és a dir, una fletxa;
  \item determina els elements globals $1\to\Omega$, que són els \enquote{valors de veritat globals}.
\end{enumerate}
\hfill{\normalfont$\star\star$}
\begin{indicacio}
Compta els sedassos sobre cada objecte de $\mathbf 2\op$. Observa que $\Omega$ no és un conjunt de valors, sinó una fletxa entre dos conjunts.
\end{indicacio}
\end{exercici}

\begin{exercici}
Demostra que un topos que és alhora una categoria prima és equivalent a la categoria terminal $\cat{1}$. \hfill{\normalfont$\star\star$}
\begin{indicacio}
En una categoria prima, tot morfisme és mono. Aplica-ho al morfisme $A\to 1$ i fes servir que $\Sub(1)\cong\Hom(1,\Omega)$ té com a màxim un element.
\end{indicacio}
\end{exercici}

\begin{exercici}[Ampliació]
Fent servir el classificador de $\cat{Graph}$ calculat a la Teoria, determina els valors de veritat globals, és a dir, els morfismes $\mathbf L\to\Omega$ des del graf terminal d'un bucle. Comprova que n'hi ha tres, identifica cadascun amb un subgraf de $\mathbf L$ i dedueix que $\Omega$ no és isomorf a $1+1$. \hfill{\normalfont$\star\star$}
\begin{indicacio}
Un morfisme $\mathbf L\to\Omega$ tria un bucle de $\Omega$. Contrasta el resultat amb $\Sub(\mathbf L)\cong\Hom(\mathbf L,\Omega)$ i compta els morfismes $\mathbf L\to\mathbf L+\mathbf L$.
\end{indicacio}
\end{exercici}

\begin{exercici}
Dedueix de l'exercici resolt~\ref{pr-mono-equalitzador} de la Pràctica que $\cat{Top}$ no té classificador de subobjectes, i per tant no és un topos. \hfill{\normalfont$\star$}
\begin{indicacio}
A $\cat{Top}$, els monomorfismes són les aplicacions contínues injectives i els epimorfismes, les exhaustives. Busca una bijecció contínua entre dos espais de dos punts que no sigui un homeomorfisme.
\end{indicacio}
\end{exercici}

\begin{exercici}
Sigui $\cat{B}\mathbb N$ el monoide additiu $(\mathbb N,+,0)$ vist com a categoria d'un sol objecte $\ast$. Descriu tots els sedassos sobre $\ast$ i dedueix que $\Omega(\ast)$ està en bijecció amb $\mathbb N\cup\{\infty\}$. Descriu l'acció $\Omega(h)$ de cada $h\in\mathbb N$. Comprova que només hi ha dos valors de veritat globals, però que $\Sub(\yon(\ast))$ no és una àlgebra de Boole. \hfill{\normalfont$\star\star$}
\begin{indicacio}
Un sedàs no buit és tancat per sumar qualsevol natural; mira'n l'element mínim. Un valor global és un sedàs fix per totes les accions. Per a la darrera part, fes servir que $\Sub(\yon(\ast))\cong\Omega(\ast)$ i busca un sedàs $S$ tal que el més gran dels sedassos disjunts de $S$, unit amb $S$, no doni el sedàs màxim.
\end{indicacio}
\end{exercici}

\section{Grafs i lògica}

\begin{exercici}
Sigui $G$ el graf amb vèrtexs $a,b$ i arestes $\ell\colon a\to a$, $e\colon a\to b$ i $d\colon b\to a$, i sigui $H$ el subgraf format pel vèrtex $a$ i el bucle $\ell$. Fent servir la descripció de $\Omega$ de la Teoria, calcula $\chi_H\colon G\to\Omega$ i comprova que el pullback de $\mathsf{true}$ al llarg de $\chi_H$ és $H$. \hfill{\normalfont$\star$}
\begin{indicacio}
Per a cada aresta, mira quins dels seus dos extrems són a $H$ i si ella mateixa hi és, i tria l'aresta de $\Omega$ que registra aquesta informació.
\end{indicacio}
\end{exercici}

\begin{exercici}
A $\Sub(\mathbf E)$, on $\mathbf E$ és el graf d'una sola aresta $e\colon s\to t$, sigui $\neg K$ el més gran dels subgrafs disjunts de $K$. Comprova que existeix, calcula $\neg K$ per als cinc subgrafs de $\mathbf E$ i determina per a quins es compleix $\neg\neg K=K$. Què en dedueixes sobre la doble negació? \hfill{\normalfont$\star\star$}
\begin{indicacio}
Un subgraf que conté l'aresta ha de contenir els seus dos extrems. Fes una taula amb $K$, $\neg K$ i $\neg\neg K$.
\end{indicacio}
\end{exercici}

\fi

% ============================================================
% PART IV — TESTS INTERACTIUS
% ============================================================
\part{\NomTests}
\setcounter{chapter}{0}
% En pantalla, un sol entorn Form per a tots els tests: un únic diccionari
% AcroForm. En mode imprimible les macros dels tests són estàtiques i no es
% crea cap formulari.
\ifimprimible\else\begin{Form}\fi
\chapter{Categories}

Aquest test recull els conceptes essencials del capítol. Cada pregunta té una única resposta correcta.

\iniciTest{cat}{c,b,a,d,c,d,a,b,a,d,b,c}

\begin{enumerate}

% --- Definició i composició ---
\pregunta Un lector afirma: \enquote{Com que la composició d'una categoria és associativa, l'expressió $h\comp g\comp f$ està ben definida per a qualssevol tres morfismes $f,g,h$.} On falla?
\begin{enumerate}
  \opcio{a}{No falla: l'associativitat garanteix que sempre es pot compondre.}
  \opcio{b}{Falla perquè l'associativitat només val en categories petites.}
  \opcio{c}{Falla perquè la composició és \emph{parcial}.}
  \opcio{d}{Falla perquè $h\comp g\comp f$ mai no està definit sense parèntesis.}
\end{enumerate}
\feedback
\justificacio{$h\comp g\comp f$ només té sentit si $\cod(f)=\dom(g)$ i $\cod(g)=\dom(h)$; l'associativitat només diu que, quan les dues agrupacions existeixen, coincideixen.}

% --- Diagrames commutatius ---

\pregunta Considera el quadrat format per $f\colon A\to B$, $u\colon A\to C$, $v\colon B\to D$ i $g\colon C\to D$. Quina equació expressa que el quadrat commuta?
\begin{enumerate}
  \opcio{a}{$f\comp v=u\comp g$, perquè es llegeix en l'ordre en què es recorren les fletxes.}
  \opcio{b}{$v\comp f=g\comp u$.}
  \opcio{c}{$v\comp u=f\comp g$.}
  \opcio{d}{$f=g$ i $u=v$.}
\end{enumerate}
\feedback
\justificacio{Els dos camins de $A$ a $D$ són \enquote{primer $f$, després $v$} i \enquote{primer $u$, després $g$}, i la composició s'escriu de dreta a esquerra.}

% --- Exemples: lògica ---

\pregunta A la categoria $\cat{Prov}_T$ d'un sistema deductiu $T$, siguin $p$ i $q$ dues fórmules \emph{distintes} tals que $p\vdash_T q$ i $q\vdash_T p$. Què en podem dir?
\begin{enumerate}
  \opcio{a}{Són objectes isomorfs però no iguals: $\cat{Prov}_T$ és prima, però la deduïbilitat no és antisimètrica.}
  \opcio{b}{Han de ser iguals, perquè en una categoria prima dos objectes isomorfs coincideixen.}
  \opcio{c}{Hi ha exactament dues fletxes de $p$ a $q$, una per a cada deducció.}
  \opcio{d}{No hi ha cap fletxa entre elles, perquè són fórmules diferents.}
\end{enumerate}
\feedback
\justificacio{Hi ha fletxes en tots dos sentits, i com que la categoria és prima, les composicions són identitats: $p\cong q$. Però $p$ i $q$ són fórmules diferents. Una categoria prima és un preordre, no necessàriament un ordre parcial, i aquí es veu la diferència entre igualtat i isomorfisme.}

% --- Exemples: categories primes ---

\pregunta Quina de les categories següents \emph{no} és prima?
\begin{enumerate}
  \opcio{a}{La categoria d'un preordre $(P,\le)$.}
  \opcio{b}{La categoria $\cat{Prov}_T$ d'un sistema deductiu.}
  \opcio{c}{La categoria $\cat{2}=(0\to 1)$.}
  \opcio{d}{La categoria $\cat{B}M$ del monoide $(\mathbb N,+,0)$.}
\end{enumerate}
\feedback
\justificacio{Té un sol objecte, però $\Hom(\ast,\ast)=\mathbb N$ té més d'un element.}

% --- Isomorfismes ---

\pregunta Sigui $D$ el conjunt $\{0,1\}$ amb l'ordre discret i sigui $C$ el mateix conjunt amb l'ordre $0<1$. L'aplicació identitat subjacent $f\colon D\to C$ és bijectiva i monòtona. És un isomorfisme a $\cat{Pos}$?
\begin{enumerate}
  \opcio{a}{Sí, perquè qualsevol aplicació bijectiva és un isomorfisme.}
  \opcio{b}{Sí, perquè una aplicació monòtona bijectiva té sempre inversa monòtona.}
  \opcio{c}{No, perquè la inversa subjacent $C\to D$ no és monòtona.}
  \opcio{d}{No, perquè $f$ no és monòtona.}
\end{enumerate}
\feedback
\justificacio{Per tant no és un morfisme de $\cat{Pos}$, i $f$ no té invers en la categoria.}

% --- Mono/epi: cancel·lació ---

\pregunta L'aplicació $i\colon\varnothing\to\{\ast\}$ de $\cat{Set}$ és:
\begin{enumerate}
  \opcio{a}{un isomorfisme.}
  \opcio{b}{un epimorfisme amb secció.}
  \opcio{c}{ni mono ni epi.}
  \opcio{d}{un monomorfisme sense retracció.}
\end{enumerate}
\feedback
\justificacio{És injectiva i, per tant, un monomorfisme de $\cat{Set}$; però no té retracció, perquè no hi ha cap aplicació $\{\ast\}\to\varnothing$.}

\pregunta Un morfisme $f\colon A\to B$ té una \emph{retracció} $r$ (amb $r\comp f=\id_A$). Quina conclusió és correcta i per què?
\begin{enumerate}
  \opcio{a}{$f$ és un monomorfisme: si $f\comp g=f\comp h$, component amb $r$ per l'esquerra dona $g=h$.}
  \opcio{b}{$f$ és un isomorfisme, perquè té invers per un costat.}
  \opcio{c}{$f$ és un epimorfisme, perquè $r$ el cancel·la per la dreta.}
  \opcio{d}{No se'n pot dir res sense saber la categoria.}
\end{enumerate}
\feedback
\justificacio{Si $f\comp g=f\comp h$, aleshores $g=r\comp f\comp g=r\comp f\comp h=h$. Tenir invers per un costat no fa $f$ isomorfisme: a $\cat{Set}$, la inclusió $\{0\}\hookrightarrow\{0,1\}$ té retracció però no és bijectiva.}

% --- Dualitat i oposada ---

\pregunta Considera la categoria oposada $\cat{C}\op$. Un lector diu: \enquote{Si $f$ és un monomorfisme a $\cat{C}$, aleshores $f$ és un monomorfisme a $\cat{C}\op$.} És correcte?
\begin{enumerate}
  \opcio{a}{Sí, perquè $\cat{C}$ i $\cat{C}\op$ tenen els mateixos morfismes.}
  \opcio{b}{No, perquè un mono de $\cat{C}$ és un epi de $\cat{C}\op$.}
  \opcio{c}{No, perquè a $\cat{C}\op$ no hi ha monomorfismes.}
  \opcio{d}{Sí, perquè els monomorfismes són autoduals.}
\end{enumerate}
\feedback
\justificacio{En passar a l'oposada, la cancel·lació per l'esquerra es converteix en cancel·lació per la dreta, igual que una retracció passa a ser una secció.}

% --- Construccions de categories ---

\pregunta Sigui $A$ un objecte d'una categoria $\cat{C}$. Què és un morfisme de $x\colon X\to A$ a $y\colon Y\to A$ a la categoria $\cat{C}/A$ sobre $A$?
\begin{enumerate}
  \opcio{a}{Un morfisme $h\colon X\to Y$ de $\cat{C}$ tal que $y\comp h=x$.}
  \opcio{b}{Qualsevol morfisme $h\colon X\to Y$ de $\cat{C}$, sense cap condició.}
  \opcio{c}{Un morfisme $h\colon Y\to X$ tal que $x\comp h=y$.}
  \opcio{d}{Un morfisme $A\to A$.}
\end{enumerate}
\feedback
\justificacio{És a dir, que fa commutar el triangle sobre $A$.}

% --- Construccions de categories: categoria producte ---

\pregunta Siguin $\cat{C}$ i $\cat{D}$ dues categories. Què és un morfisme de $(A,X)$ a $(B,Y)$ a la categoria producte $\cat{C}\times\cat{D}$?
\begin{enumerate}
  \opcio{a}{Un morfisme $A\times X\to B\times Y$ d'una de les dues categories.}
  \opcio{b}{Un morfisme $A\to B$ de $\cat{C}$ o bé un morfisme $X\to Y$ de $\cat{D}$.}
  \opcio{c}{L'objecte producte $A\times B$, caracteritzat per la seva propietat universal dins de $\cat{C}$.}
  \opcio{d}{Una parella $(f,g)$ amb $f\colon A\to B$ a $\cat{C}$ i $g\colon X\to Y$ a $\cat{D}$.}
\end{enumerate}
\feedback
\justificacio{La composició es fa component a component. No s'ha de confondre amb el producte d'objectes dins d'una sola categoria.}

% --- Categoria lliure ---

\pregunta Considera el graf amb un sol vèrtex $\ast$ i una sola aresta $a\colon\ast\to\ast$. Quants morfismes té la categoria lliure que genera?
\begin{enumerate}
  \opcio{a}{Dos: la identitat i $a$.}
  \opcio{b}{Infinits, un per a cada longitud de camí.}
  \opcio{c}{Un de sol, perquè el graf té un sol vèrtex.}
  \opcio{d}{Cap, perquè $a$ no es pot compondre amb si mateixa.}
\end{enumerate}
\feedback
\justificacio{Són $\id_\ast,\ a,\ a\comp a,\ a\comp a\comp a,\dots$; és la categoria $\cat{B}M$ del monoide $(\mathbb N,+,0)$.}

% --- Mida ---

\pregunta Un lector afirma: \enquote{La categoria $\cat{Set}$ és petita, perquè els seus morfismes són simplement aplicacions.} On s'equivoca?
\begin{enumerate}
  \opcio{a}{No s'equivoca: $\cat{Set}$ és petita.}
  \opcio{b}{S'equivoca perquè $\cat{Set}$ no té identitats.}
  \opcio{c}{S'equivoca perquè $\cat{Set}$ no és petita.}
  \opcio{d}{S'equivoca perquè $\cat{Set}$ ni tan sols és una categoria.}
\end{enumerate}
\feedback
\justificacio{És \emph{localment petita} (cada $\Hom(A,B)$ és un conjunt), però la col·lecció de tots els conjunts és una classe pròpia, no un conjunt.}

\end{enumerate}
\clauestatica
\finalitzaTest

\chapter{Functors}

Aquest test recull els conceptes essencials del capítol. Cada pregunta té una única resposta correcta.

\iniciTest{fun}{a,b,d,c,a,c,d,a,b,c,b}

\begin{enumerate}

\pregunta Un lector afirma: \enquote{Per definir un functor $F$ n'hi ha prou de dir on va cada objecte; l'acció sobre els morfismes queda determinada.} On falla?
\begin{enumerate}
  \opcio{a}{Falla perquè l'acció sobre morfismes no queda determinada per la dels objectes.}
  \opcio{b}{Falla només quan $\cat{C}$ és gran.}
  \opcio{c}{No falla: els objectes determinen el functor.}
  \opcio{d}{Falla perquè els functors no actuen sobre objectes.}
\end{enumerate}
\feedback
\justificacio{Un functor és \emph{dues} assignacions, sobre objectes i sobre morfismes, que han de respectar dominis, codominis, composició i identitats.}

\pregunta Un lector proposa un \enquote{functor} $\cat{B}M\to\cat{B}M$ (amb $M=\{1,e\}$, $e\comp e=e$) que fixa l'objecte i envia el bucle $1$ a $e$ i el bucle $e$ a $e$. És un functor?
\begin{enumerate}
  \opcio{a}{Sí, perquè respecta dominis i codominis.}
  \opcio{b}{No, perquè no envia la identitat a la identitat.}
  \opcio{c}{Sí, perquè $e$ és idempotent.}
  \opcio{d}{No, perquè cap assignació entre monoides és functorial.}
\end{enumerate}
\feedback
\justificacio{La identitat $1$ de $\cat{B}M$ va a parar a $e\neq 1$; és un morfisme de grafs, però no un functor.}

\pregunta Es pot afirmar que \emph{tot} functor preserva els monomorfismes?
\begin{enumerate}
  \opcio{a}{Sí, perquè preserva tota l'estructura categòrica.}
  \opcio{b}{Sí, perquè preserva la composició.}
  \opcio{c}{No, perquè cap functor preserva mai els monomorfismes.}
  \opcio{d}{No en general, tot i que alguns functors sí que els preserven.}
\end{enumerate}
\feedback
\justificacio{Un functor preserva el que està testimoniat per morfismes i equacions (isomorfismes, seccions, retraccions), però ser mono es defineix per cancel·lació sobre \emph{totes} les fletxes de la categoria d'arribada, incloses les que no provenen de la categoria d'origen.}

\pregunta Sigui $\cat{S}$ la subcategoria plena de $\cat{FinSet}$ que té per objectes els conjunts $\underline{0}=\varnothing$ i $\underline{n}=\{1,\dots,n\}$ per a $n\geq 1$, un de cada cardinal finit, i sigui $\iota\colon\cat{S}\hookrightarrow\cat{FinSet}$ la inclusió. Quina afirmació és correcta?
\begin{enumerate}
  \opcio{a}{$\iota$ és exhaustiu sobre objectes en sentit estricte.}
  \opcio{b}{$\iota$ no és ni ple ni fidel.}
  \opcio{c}{$\iota$ és essencialment exhaustiu però no exhaustiu sobre objectes.}
  \opcio{d}{$\iota$ és un isomorfisme de categories.}
\end{enumerate}
\feedback
\justificacio{Tot conjunt finit és isomorf a algun $\underline{n}$, però en general no és igual a cap d'ells.}

\pregunta Sigui $F$ un functor \emph{ple i fidel} i suposem que $F(f)$ és un isomorfisme a $\cat{D}$. Quina conclusió és correcta, i sobre quin fet reposa?
\begin{enumerate}
  \opcio{a}{Que $f$ ja era un isomorfisme a $\cat{C}$.}
  \opcio{b}{Que no es pot concloure que $f$ sigui un isomorfisme sense suposar, a més, que $F$ és exhaustiu sobre objectes.}
  \opcio{c}{Res: ser ple i fidel no diu res sobre isomorfismes.}
  \opcio{d}{Que $F$ és una equivalència.}
\end{enumerate}
\feedback
\justificacio{Com que $F$ és ple, l'invers de $F(f)$ prové d'algun $g$; com que $F$ és fidel, les identitats $g\comp f=\id$ i $f\comp g=\id$ es dedueixen de les seves imatges. Els functors plens i fidels \emph{reflecteixen} isomorfismes.}

\pregunta Siguin $f\colon X\to Y$ i $h\colon Y\to B$. Quina és l'acció de l'homfunctor contravariant $\Hom(-,B)$ sobre $f$, aplicada a $h$?
\begin{enumerate}
  \opcio{a}{$h\mapsto f\comp h$, amb valors a $\Hom(X,B)$.}
  \opcio{b}{$h\mapsto f\comp h$, amb valors a $\Hom(Y,B)$.}
  \opcio{c}{$h\mapsto h\comp f$, amb valors a $\Hom(X,B)$.}
  \opcio{d}{$h\mapsto h\comp f$, amb valors a $\Hom(Y,B)$.}
\end{enumerate}
\feedback
\justificacio{La precomposició amb $f$ va de $\Hom(Y,B)$ a $\Hom(X,B)$: inverteix el sentit de $f$.}

\pregunta Un lector diu: \enquote{El functor oblit $U\colon\cat{Grp}\to\cat{Set}$ és ple, perquè envia cada homomorfisme a la seva aplicació subjacent.} On s'equivoca?
\begin{enumerate}
  \opcio{a}{No s'equivoca: $U$ és ple.}
  \opcio{b}{S'equivoca perquè $U$ no és ni tan sols un functor.}
  \opcio{c}{S'equivoca perquè $U$ no actua sobre morfismes.}
  \opcio{d}{S'equivoca: $U$ és fidel, però no ple.}
\end{enumerate}
\feedback
\justificacio{Homomorfismes diferents donen aplicacions diferents; però entre dos grups hi ha aplicacions de conjunts que no són homomorfismes i que, per tant, no provenen de cap morfisme de $\cat{Grp}$.}

\pregunta Un lector afirma: \enquote{Com que $\cat{Cat}$ té totes les categories petites per objectes, $\cat{Cat}$ és ella mateixa una categoria petita.} És correcte?
\begin{enumerate}
  \opcio{a}{No, perquè $\cat{Cat}$ és gran, encara que localment petita.}
  \opcio{b}{Sí, perquè els functors formen conjunts.}
  \opcio{c}{Sí: una categoria de categories petites és petita.}
  \opcio{d}{No, perquè $\cat{Cat}$ no té identitats.}
\end{enumerate}
\feedback
\justificacio{La col·lecció de totes les categories petites és una classe pròpia, com la de tots els conjunts.}

\pregunta Donats un graf $G$ i una categoria petita $\cat{D}$, definir un functor $\Free(G)\to\cat{D}$ des de la categoria lliure de $G$ equival a:
\begin{enumerate}
  \opcio{a}{triar només les imatges de les arestes de $G$.}
  \opcio{b}{donar un morfisme de grafs $G\to U(\cat{D})$.}
  \opcio{c}{triar un objecte de $\cat{D}$ per a cada camí de $G$ de manera independent.}
  \opcio{d}{donar una equivalència entre $\Free(G)$ i $\cat{D}$.}
\end{enumerate}
\feedback
\justificacio{Cal triar les imatges dels vèrtexs \emph{i} de les arestes, respectant origen i destí; aquesta tria s'estén de manera única als camins.}

\pregunta Sobre el bifunctor $\Hom(-,-)\colon\cat{C}\op\times\cat{C}\to\cat{Set}$, amb $\Hom(f,g)(h)=g\comp h\comp f$, quina descripció de la variància és correcta?
\begin{enumerate}
  \opcio{a}{Covariant en totes dues variables.}
  \opcio{b}{Contravariant en totes dues variables.}
  \opcio{c}{Contravariant en la primera variable i covariant en la segona.}
  \opcio{d}{Covariant en la primera i contravariant en la segona.}
\end{enumerate}
\feedback
\justificacio{Precomposició amb $f$ en la primera variable i postcomposició amb $g$ en la segona.}

% --- Fil lògic ---

\pregunta Siguin $T\subseteq T'$ dos sistemes deductius sobre les mateixes fórmules, i sigui
\[
J\colon\cat{Prov}_T\to\cat{Prov}_{T'}
\]
el functor identitat sobre fórmules. Quan és $J$ ple?
\begin{enumerate}
  \opcio{a}{Sempre, perquè és la identitat sobre objectes.}
  \opcio{b}{Exactament quan $T'$ no introdueix cap conseqüència nova $A\vdash_{T'}B$ entre les fórmules considerades.}
  \opcio{c}{Només quan $T=T'$ com a conjunts d'axiomes.}
  \opcio{d}{Mai, perquè les categories de deduïbilitat són primes.}
\end{enumerate}
\feedback
\justificacio{Totes dues categories són primes, de manera que $J$ és ple si i només si cada fletxa $A\to B$ de $\cat{Prov}_{T'}$ prové d'una de $\cat{Prov}_T$, és a dir, si $A\vdash_{T'}B$ implica $A\vdash_T B$. Ser la identitat sobre objectes (a) no diu res sobre els morfismes.}

\end{enumerate}
\clauestatica
\finalitzaTest

\chapter[Transformacions naturals]{\texorpdfstring{Transformacions naturals\\ i equivalències}{Transformacions naturals i equivalències}}

Aquest test recull els conceptes essencials del capítol. Cada pregunta té una única resposta correcta.

\iniciTest{trn}{c,b,c,d,a,b,c,b,a,d,a,d}

\begin{enumerate}

\pregunta A $\cat{Set}$, prenem $F=G=\id_{\cat{Set}}$ i definim $\eta_A=\id_A$ per a tot conjunt $A\neq\{0,1\}$, i $\eta_{\{0,1\}}=\tau$, la transposició $\tau(0)=1$, $\tau(1)=0$. Aquesta família, \emph{és} una transformació natural $\id\Rightarrow\id$?
\begin{enumerate}
  \opcio{a}{Sí, perquè cada component $\eta_A$ és una bijecció.}
  \opcio{b}{Sí, perquè $\eta$ coincideix amb la identitat a tots els conjunts excepte un.}
  \opcio{c}{No, perquè una aplicació constant trenca el quadrat de naturalitat.}
  \opcio{d}{No, perquè la transposició $\tau$ no és un morfisme de $\cat{Set}$.}
\end{enumerate}
\feedback
\justificacio{Per a $f\colon\{0,1\}\to\{0,1\}$ constant de valor $0$, es té $f\comp\tau=f$, mentre que $\tau\comp f$ és constant de valor $1$.}

\pregunta Un lector proposa aquest argument: \enquote{Si $F(A)\cong G(A)$ per a cada objecte $A$, aleshores $F$ i $G$ són naturalment isomorfs, perquè n'hi ha prou d'agafar un isomorfisme $\eta_A\colon F(A)\to G(A)$ per a cada $A$.} On falla?
\begin{enumerate}
  \opcio{a}{No falla: l'argument és correcte.}
  \opcio{b}{Falla perquè els $\eta_A$ triats poden no complir el quadrat de naturalitat.}
  \opcio{c}{Falla perquè un isomorfisme natural exigeix que $F=G$.}
  \opcio{d}{Falla perquè $F(A)$ i $G(A)$ no poden ser mai isomorfs si $F\neq G$.}
\end{enumerate}
\feedback
\justificacio{L'existència d'isomorfismes objecte a objecte no garanteix que siguin \emph{compatibles} amb l'acció dels functors sobre els morfismes.}

\pregunta Sigui $\mathbb{Z}/2\mathbb{Z}=\{1,g\}$ vist com a categoria d'un objecte, i $F,G\colon\cat{B}(\mathbb{Z}/2\mathbb{Z})\to\cat{Set}$ que envien l'objecte a $\{0,1\}$, on $g$ actua per la identitat en $F$ i per la transposició $\tau$ en $G$. Quantes transformacions naturals $\eta\colon F\Rightarrow G$ hi ha?
\begin{enumerate}
  \opcio{a}{Exactament una, l'única component possible $\eta_\ast=\id$.}
  \opcio{b}{Dues: la identitat i la transposició.}
  \opcio{c}{Cap.}
  \opcio{d}{Infinites, una per cada aplicació $\{0,1\}\to\{0,1\}$.}
\end{enumerate}
\feedback
\justificacio{La naturalitat respecte de $g$ exigiria $\tau\comp\eta_\ast=\eta_\ast$, impossible perquè $\tau$ no té punts fixos.}

\pregunta Què és un morfisme de $f\colon A\to B$ a $g\colon A'\to B'$ a la categoria de functors $\cat{D}^{\cat{2}}$, on $\cat{2}=(0\to 1)$?
\begin{enumerate}
  \opcio{a}{Un morfisme qualsevol $A\to A'$.}
  \opcio{b}{Un functor $\cat{2}\to\cat{D}$.}
  \opcio{c}{Una parella qualsevol de morfismes $u\colon A\to A'$ i $v\colon B\to B'$.}
  \opcio{d}{Una parella $u\colon A\to A'$, $v\colon B\to B'$ amb $v\comp f=g\comp u$.}
\end{enumerate}
\feedback
\justificacio{És la condició de naturalitat per a l'única fletxa no trivial de $\cat{2}$.}

\pregunta En l'exemple sintaxi/semàntica, amb el functor sintàctic $F$, el semàntic $G$ i la interpretació $\eta$, què expressa la igualtat $G(f)\comp\eta_X=\eta_Y\comp F(f)$ per a un canvi de variables $f\colon X\to Y$?
\begin{enumerate}
  \opcio{a}{Que substituir variables i després interpretar dona la mateixa funció de veritat que interpretar i després transportar la semàntica al llarg del canvi de variables.}
  \opcio{b}{Que $\eta_X$ és una bijecció entre fórmules i funcions de veritat.}
  \opcio{c}{Que dues fórmules amb la mateixa taula de veritat són iguals.}
  \opcio{d}{Que el canvi de variables $f$ ha de ser bijectiu.}
\end{enumerate}
\feedback
\justificacio{És el quadrat de naturalitat de $\eta$ per a $f$: els dos camins de $F(X)$ a $G(Y)$ coincideixen. En termes lògics, la interpretació és compatible amb la substitució. Ni (b) ni (c) són certes, perquè fórmules diferents poden tenir la mateixa taula de veritat.}

\pregunta Un estudiant escriu que la composició vertical de $\eta\colon F\Rightarrow G$ i $\theta\colon G\Rightarrow H$ té components $(\theta\comp\eta)_A=\eta_A\comp\theta_A$. Per què és incorrecte, i quina és la forma bona?
\begin{enumerate}
  \opcio{a}{És correcte tal com està.}
  \opcio{b}{L'ordre està invertit: la composta és $\theta_A\comp\eta_A$.}
  \opcio{c}{Cap dels dos ordres té sentit; la composició vertical no existeix.}
  \opcio{d}{Tots dos ordres donen el mateix resultat, així que és igual.}
\end{enumerate}
\feedback
\justificacio{$\eta_A\colon F(A)\to G(A)$ i $\theta_A\colon G(A)\to H(A)$ només es componen en aquest ordre, i la composta va de $F(A)$ a $H(A)$.}

\pregunta Sigui $L\colon\cat{Set}\to\cat{Set}$ el functor de llistes finites, amb $L(f)$ aplicant $f$ a cada entrada, i sigui $\eta_X(x)=[x]$ la llista d'un sol element. Per a $f\colon X\to Y$, quina és l'equació de naturalitat ben tipada?
\begin{enumerate}
  \opcio{a}{$f\comp\eta_X=\eta_Y\comp L(f)$.}
  \opcio{b}{$\eta_X\comp L(f)=f\comp\eta_Y$.}
  \opcio{c}{$L(f)\comp\eta_X=\eta_Y\comp f$.}
  \opcio{d}{$\eta_Y=L(f)$.}
\end{enumerate}
\feedback
\justificacio{Tots dos costats envien $x$ a $[f(x)]$.}

\pregunta Un functor $F\colon\cat{C}\to\cat{D}$ és ple, fidel i essencialment exhaustiu. Un lector conclou: \enquote{Aleshores $F$ és bijectiu sobre els objectes.} És correcte?
\begin{enumerate}
  \opcio{a}{Sí: essencialment exhaustiu vol dir exhaustiu sobre objectes, i ple i fidel vol dir injectiu sobre objectes.}
  \opcio{b}{No: $F$ pot no ser bijectiu sobre objectes.}
  \opcio{c}{No, perquè cap functor és mai bijectiu sobre objectes.}
  \opcio{d}{Sí, sempre que $\cat{C}$ i $\cat{D}$ siguin petites.}
\end{enumerate}
\feedback
\justificacio{Essencialment exhaustiu només demana que tot objecte de $\cat{D}$ sigui \emph{isomorf} a una imatge, i ser ple i fidel afecta els morfismes; per exemple, la inclusió d'un esquelet no és bijectiva sobre objectes.}

\pregunta A la construcció del quasiinvers $G$ d'una equivalència, per a cada objecte $D$ es tria un objecte $G(D)$ i un isomorfisme $\varepsilon_D\colon F(G(D))\to D$, i sobre morfismes es defineix $G(v)$ com l'\emph{única} fletxa amb $F(G(v))=\varepsilon_{D'}^{-1}\comp v\comp\varepsilon_D$. Quina propietat de $F$ garanteix aquesta unicitat i existència?
\begin{enumerate}
  \opcio{a}{Que $F$ sigui ple i fidel.}
  \opcio{b}{Només ser fidel.}
  \opcio{c}{L'exhaustivitat essencial.}
  \opcio{d}{Que $\cat{D}$ sigui petita.}
\end{enumerate}
\feedback
\justificacio{Ser ple dona l'existència de $G(v)$, i ser fidel, la unicitat.}

\pregunta Siguin $\eta\colon F\Rightarrow G$, amb $F,G\colon\cat{C}\to\cat{D}$, i $\theta\colon K\Rightarrow L$, amb $K,L\colon\cat{D}\to\cat{E}$. Quina és la component $(\theta\ast\eta)_A$ de la composició horitzontal?
\begin{enumerate}
  \opcio{a}{$\theta_A\comp\eta_A$.}
  \opcio{b}{$K(\eta_A)\comp\theta_{G(A)}$.}
  \opcio{c}{$\eta_{K(A)}\comp\theta_A$.}
  \opcio{d}{$\theta_{G(A)}\comp K(\eta_A)$.}
\end{enumerate}
\feedback
\justificacio{Coincideix amb $L(\eta_A)\comp\theta_{F(A)}$ per la naturalitat de $\theta$ aplicada a $\eta_A$.}

\pregunta Siguin $\cat{1}$ la categoria terminal i $\cat{I}$ la categoria amb dos objectes i exactament un morfisme entre cada parella ordenada d'objectes. Quina afirmació és correcta?
\begin{enumerate}
  \opcio{a}{Són equivalents però no isomorfes.}
  \opcio{b}{Són isomorfes, perquè tots els objectes de $\cat{I}$ són isomorfs entre si.}
  \opcio{c}{No són equivalents, perquè tenen un nombre diferent d'objectes.}
  \opcio{d}{Són isomorfes, perquè totes dues són categories primes.}
\end{enumerate}
\feedback
\justificacio{La inclusió $\cat{1}\to\cat{I}$ és plena, fidel i essencialment exhaustiva, però cap functor entre elles no pot ser bijectiu sobre objectes.}

\pregunta Sigui $F\colon\cat{Mat}_{\mathbb K}\to\cat{FinVect}_{\mathbb K}$, $n\mapsto\mathbb K^n$, el functor de l'exemple de la Teoria. Quina afirmació és correcta?
\begin{enumerate}
  \opcio{a}{Les dues categories són iguals, perquè tot espai de dimensió $n$ \emph{és} $\mathbb K^n$.}
  \opcio{b}{Són isomorfes però no iguals: $F$ té un invers estricte.}
  \opcio{c}{Són equivalents i, per tant, isomorfes.}
  \opcio{d}{Són equivalents però no isomorfes; amb una tria adequada de bases, $G\comp F=\id$ literalment, però $F\comp G$ només és isomorf a la identitat.}
\end{enumerate}
\feedback
\justificacio{$F$ és ple, fidel i essencialment exhaustiu, i per tant una equivalència. No hi ha cap isomorfisme de categories, perquè $\cat{Mat}_{\mathbb K}$ és esquelètica i $\cat{FinVect}_{\mathbb K}$ no ho és. El quasiinvers tria una base de cada espai: $F(G(V))=\mathbb K^{\dim V}$, que és isomorf a $V$ però en general no hi és igual. Un espai de dimensió $n$ és isomorf a $\mathbb K^n$, no igual.}

\end{enumerate}
\clauestatica
\finalitzaTest

\ifpublicacioparcial\else
\chapter[Representabilitat i Yoneda]{\texorpdfstring{Propietats universals,\\ representabilitat i Yoneda}{Propietats universals, representabilitat i Yoneda}}

Aquest test recull els conceptes essencials del capítol. Cada pregunta té una única resposta correcta.

\iniciTest{yon}{a,c,b,d,a,d,b,c,a,c,b,d,c,d,b}

\begin{enumerate}

\pregunta Quan és \textbf{terminal} un objecte $1$ d'una categoria?
\begin{enumerate}
  \opcio{a}{Quan hi ha exactament un morfisme $X\to 1$ per a cada $X$.}
  \opcio{b}{Quan hi ha exactament un morfisme $1\to X$ per a cada $X$.}
  \opcio{c}{Quan no hi ha cap morfisme cap a ell.}
  \opcio{d}{Quan és isomorf a tots els altres objectes.}
\end{enumerate}
\feedback
\justificacio{És la definició. L'opció (b) és la d'objecte inicial, i (d) no té per què complir-se: ser terminal és una propietat universal, no un isomorfisme amb tots els objectes.}

\pregunta Siguin $S\colon\cat{A}\to\cat{C}$ i $T\colon\cat{B}\to\cat{C}$. Una parella $a\colon A\to A'$, $b\colon B\to B'$ és un morfisme $(A,B,u)\to(A',B',u')$ de la categoria coma $(S\downarrow T)$ quan:
\begin{enumerate}
  \opcio{a}{$S(a)\comp u=u'\comp T(b)$.}
  \opcio{b}{$T(b)\comp u'=u\comp S(a)$.}
  \opcio{c}{$T(b)\comp u=u'\comp S(a)$.}
  \opcio{d}{$S(a)=T(b)$.}
\end{enumerate}
\feedback
\justificacio{Comprova els tipus: $u\colon S(A)\to T(B)$ i $u'\colon S(A')\to T(B')$. L'únic quadrat ben tipat de $S(A)$ a $T(B')$ és $T(b)\comp u=u'\comp S(a)$. A (a) i (b) les composicions no existeixen en general.}

\pregunta En la formulació per representabilitat desenvolupada en aquest capítol, una \textbf{propietat universal} caracteritza una dada...
\begin{enumerate}
  \opcio{a}{per la seva constitució interna.}
  \opcio{b}{per la manera com es relaciona amb tots els altres objectes, mitjançant la representabilitat d'un functor cap a $\cat{Set}$.}
  \opcio{c}{per la seva cardinalitat.}
  \opcio{d}{només a $\cat{Set}$.}
\end{enumerate}
\feedback
\justificacio{Una propietat universal no descriu què \emph{és} la dada, sinó com es relaciona amb tota la resta: la dada és un objecte representant, amb el seu element universal.}

\pregunta Que $1$ sigui terminal equival a dir que:
\begin{enumerate}
  \opcio{a}{$\Hom(1,-)$ és constant.}
  \opcio{b}{$\Hom(-,1)$ no és functor.}
  \opcio{c}{$\Hom(-,1)$ és buit.}
  \opcio{d}{$\Hom(-,1)$ és isomorf al functor constant amb valor un punt.}
\end{enumerate}
\feedback
\justificacio{Ser terminal vol dir que cada $\Hom(X,1)$ té exactament un element. Aquestes bijeccions amb $\{\ast\}$ formen un isomorfisme natural, perquè tots els quadrats acaben en un conjunt d'un element.}

\pregunta Quan es diu que un prefeix $P\colon\cat{C}\op\to\cat{Set}$ és \textbf{representable}?
\begin{enumerate}
  \opcio{a}{Quan és isomorf a $\Hom(-,A)$ per a algun objecte $A$.}
  \opcio{b}{Quan és constant.}
  \opcio{c}{Quan és ple i fidel.}
  \opcio{d}{Quan té un sol element a cada conjunt.}
\end{enumerate}
\feedback
\justificacio{És la definició, en la forma de la proposició sobre representabilitat i isomorfisme natural. Ser constant o tenir un element a cada conjunt són casos molt particulars; \enquote{ple i fidel} és una propietat de functors entre categories, no de prefeixos.}

\pregunta Sigui $P\colon\cat{C}\op\to\cat{Set}$ un prefeix i $x\in P(A)$. Quina és la component en $X$ de la transformació natural $\Hom(-,A)\Rightarrow P$ que el lema de Yoneda associa a $x$, avaluada en $f\colon X\to A$?
\begin{enumerate}
  \opcio{a}{$P(f^{-1})(x)$.}
  \opcio{b}{$x$, sense canviar de conjunt.}
  \opcio{c}{$P(\id_A)(f)$.}
  \opcio{d}{$P(f)(x)$.}
\end{enumerate}
\feedback
\justificacio{Per contravariància, $P(f)\colon P(A)\to P(X)$; no cal que $f$ tingui invers.}

\pregunta Sigui $f\colon A\to B$ un morfisme. Les quatre descripcions següents de la component en $X$ de $\yon(f)=\Hom(-,f)$ semblen raonables. Quina és l'única en què el domini, el codomini i la composició tenen els tipus correctes?
\begin{enumerate}
  \opcio{a}{$\Hom(X,A)\to\Hom(X,B)$, $\varphi\mapsto\varphi\comp f$.}
  \opcio{b}{$\Hom(X,A)\to\Hom(X,B)$, $\varphi\mapsto f\comp\varphi$.}
  \opcio{c}{$\Hom(A,X)\to\Hom(B,X)$, $\varphi\mapsto\varphi\comp f$.}
  \opcio{d}{$\Hom(X,B)\to\Hom(X,A)$, $\varphi\mapsto\varphi\comp f$.}
\end{enumerate}
\feedback
\justificacio{Si $\varphi\colon X\to A$, l'única composició possible amb $f\colon A\to B$ és $f\comp\varphi\colon X\to B$. A (a), $\varphi\comp f$ demanaria $B=X$. A (c), amb $\varphi\colon A\to X$, $\varphi\comp f$ demanaria $B=A$: és una versió mal tipada de $\Hom(f,X)$, que va de $\Hom(B,X)$ a $\Hom(A,X)$. A (d), amb $\varphi\colon X\to B$, $\varphi\comp f$ demanaria de nou $B=X$.}

\pregunta Sigui $\mathcal P$ el prefeix de parts, $A=\{0,1\}$ i $B=\{1\}$. Quin subconjunt assigna la transformació natural $\Hom(-,A)\Rightarrow\mathcal P$ determinada per $B$ a l'aplicació $f\colon\{a,b,c\}\to A$ amb $f(a)=f(c)=1$ i $f(b)=0$?
\begin{enumerate}
  \opcio{a}{$\{b\}$.}
  \opcio{b}{$\{a,b,c\}$.}
  \opcio{c}{$\{a,c\}$.}
  \opcio{d}{$\{1\}$.}
\end{enumerate}
\feedback
\justificacio{Pel lema de Yoneda, la transformació determinada per $B\in\mathcal P(A)$ envia $f$ a $\mathcal P(f)(B)=f^{-1}(B)=\{a,c\}$. L'opció (d) és $B$ mateix, que viu a $\mathcal P(A)$ i no a $\mathcal P(\{a,b,c\})$.}

\pregunta Sigui $\Phi_{A,P}\colon\Nat(\yon(A),P)\to P(A)$, $\eta\mapsto\eta_A(\id_A)$, la bijecció del lema de Yoneda, i sigui $a\colon A'\to A$. Quina de les igualtats següents expressa la naturalitat en $A$ amb tots els tipus correctes?
\begin{enumerate}
  \opcio{a}{$\Phi_{A',P}\bigl(\eta\comp\yon(a)\bigr)=P(a)\bigl(\Phi_{A,P}(\eta)\bigr)$, com a elements de $P(A')$.}
  \opcio{b}{$\Phi_{A,P}\bigl(\eta\comp\yon(a)\bigr)=P(a)\bigl(\Phi_{A',P}(\eta)\bigr)$, com a elements de $P(A)$.}
  \opcio{c}{$\Phi_{A',P}\bigl(\yon(a)\comp\eta\bigr)=P(a)\bigl(\Phi_{A,P}(\eta)\bigr)$, com a elements de $P(A')$.}
  \opcio{d}{$\Phi_{A',P}\bigl(\eta\comp\yon(a)\bigr)=P(a)^{-1}\bigl(\Phi_{A,P}(\eta)\bigr)$, com a elements de $P(A')$.}
\end{enumerate}
\feedback
\justificacio{Tenim $\eta\colon\yon(A)\Rightarrow P$ i $\yon(a)\colon\yon(A')\Rightarrow\yon(A)$. L'única composició possible és, doncs, $\eta\comp\yon(a)$, que va de $\yon(A')$ a $P$, i s'hi aplica $\Phi_{A',P}$. Per contravariància, $P(a)\colon P(A)\to P(A')$. A (b), $\Phi_{A,P}$ no s'aplica a una transformació que surt de $\yon(A')$. A (c), $\yon(a)\comp\eta$ no és componible. A (d), $P(a)$ no té per què ser invertible.}

\pregunta A la demostració del lema, a quina dada s'envia una transformació natural $\eta\colon\Hom(-,A)\Rightarrow P$?
\begin{enumerate}
  \opcio{a}{A $\eta_A$.}
  \opcio{b}{A tot el functor $P$.}
  \opcio{c}{A l'element $\eta_A(\id_A)\in P(A)$.}
  \opcio{d}{A la identitat de $P$.}
\end{enumerate}
\feedback
\justificacio{L'aplicació del lema és $\eta\mapsto\eta_A(\id_A)$. L'opció (a) és tota la component, no un element de $P(A)$; és l'element el que determina $\eta$.}

\pregunta Una conseqüència directa del lema de Yoneda és que la immersió $\yon$ és:
\begin{enumerate}
  \opcio{a}{constant.}
  \opcio{b}{plena i fidel.}
  \opcio{c}{essencialment exhaustiva.}
  \opcio{d}{un isomorfisme de categories.}
\end{enumerate}
\feedback
\justificacio{Amb $P=\Hom(-,B)$, el lema dona $\Nat(\Hom(-,A),\Hom(-,B))\cong\Hom(A,B)$, i la bijecció és $f\mapsto\yon(f)$. No és essencialment exhaustiva: hi ha prefeixos que no són representables.}

\pregunta Segons el corol·lari de Yoneda, $A\cong B$ si i només si:
\begin{enumerate}
  \opcio{a}{$A$ i $B$ tenen la mateixa cardinalitat.}
  \opcio{b}{$A=B$.}
  \opcio{c}{hi ha un functor entre ells.}
  \opcio{d}{$\Hom(-,A)\cong\Hom(-,B)$ com a prefeixos.}
\end{enumerate}
\feedback
\justificacio{Un functor ple i fidel reflecteix isomorfismes, i tot functor els preserva; aplicat a $\yon$, dona l'equivalència. La igualtat (b) és massa forta, i la cardinalitat (a) no té sentit en una categoria general.}

\pregunta Siguin $(P,p_1,p_2)$ i $(Q,q_1,q_2)$ dos productes dels mateixos objectes. Què és únic?
\begin{enumerate}
  \opcio{a}{Res: poden ser literalment iguals o no.}
  \opcio{b}{L'isomorfisme entre els objectes $P$ i $Q$, sense cap condició.}
  \opcio{c}{L'isomorfisme $u\colon P\to Q$ que compleix $q_i\comp u=p_i$ per a $i=1,2$.}
  \opcio{d}{L'automorfisme de $P$, que només pot ser la identitat.}
\end{enumerate}
\feedback
\justificacio{La propietat universal fa que l'isomorfisme sigui únic \emph{entre els compatibles amb les projeccions}. Sense aquesta condició, pot haver-n'hi molts: a $\cat{Set}$, per exemple, qualsevol bijecció de $P$ en $Q$.}

\pregunta Siguin $S\colon\cat{A}\to\cat{C}$ i $T\colon\cat{B}\to\cat{C}$ dos functors. Un objecte de la categoria coma $(S\downarrow T)$ és:
\begin{enumerate}
  \opcio{a}{només un objecte de $\cat{C}$.}
  \opcio{b}{una transformació natural $S\Rightarrow T$.}
  \opcio{c}{una parella d'objectes isomorfs.}
  \opcio{d}{una terna $(A,B,f)$ amb $A$ de $\cat{A}$, $B$ de $\cat{B}$ i $f\colon S(A)\to T(B)$.}
\end{enumerate}
\feedback
\justificacio{És la definició: un objecte de $\cat{A}$, un de $\cat{B}$ i una fletxa de $S(A)$ a $T(B)$ a $\cat{C}$. Les categories sobre i sota un objecte en són casos particulars.}

\pregunta Si $(A,u)$ representa un prefeix $P\colon\cat{C}\op\to\cat{Set}$, què expressa la universalitat de l'element $u\in P(A)$?
\begin{enumerate}
  \opcio{a}{Que $u$ és l'únic element de $P(A)$.}
  \opcio{b}{Que per a cada objecte $X$ i cada $x\in P(X)$ hi ha un únic $f\colon X\to A$ amb $P(f)(u)=x$.}
  \opcio{c}{Que $A$ és un objecte terminal.}
  \opcio{d}{Que $P$ és un functor constant.}
\end{enumerate}
\feedback
\justificacio{És la definició de representació d'un prefeix: cada element de cada $P(X)$ s'obté de $u$ per un únic morfisme $X\to A$. Equival a l'isomorfisme natural $\Hom(-,A)\cong P$ que envia $f$ a $P(f)(u)$.}

\end{enumerate}
\clauestatica
\finalitzaTest

\chapter[Límits i colímits]{Límits i colímits}

Aquest test recull els conceptes essencials del capítol. Cada pregunta té una única resposta correcta.

\iniciTest{lim}{b,c,a,d,d,b,c,a,d,c,a,c,b,b,a}

\begin{enumerate}

\pregunta Què és un \textbf{diagrama} de forma $\cat{I}$ en una categoria $\cat{C}$?
\begin{enumerate}
  \opcio{a}{Un objecte de $\cat{C}$.}
  \opcio{b}{Un functor $D\colon\cat{I}\to\cat{C}$, on $\cat{I}$ és una categoria índex petita.}
  \opcio{c}{Una transformació natural.}
  \opcio{d}{Un morfisme de $\cat{I}$.}
\end{enumerate}
\feedback
\justificacio{És la definició: un diagrama és un functor, i la categoria índex en dona la forma. Un objecte o un morfisme sols no recullen les compatibilitats entre les fletxes.}

\pregunta Un \textbf{con} sobre $D$ amb vèrtex $A$ és el mateix que:
\begin{enumerate}
  \opcio{a}{un objecte terminal.}
  \opcio{b}{un límit.}
  \opcio{c}{una transformació natural $\Delta_A\Rightarrow D$ des del functor constant.}
  \opcio{d}{un functor $\cat{C}\to\cat{I}$.}
\end{enumerate}
\feedback
\justificacio{Les components de $\alpha\colon\Delta_A\Rightarrow D$ són fletxes $A\to D_i$, i la naturalitat per a $u\colon i\to j$ diu $D_u\comp\alpha_i=\alpha_j$, que és exactament la condició de con.}

\pregunta Un \textbf{límit} de $D$ és:
\begin{enumerate}
  \opcio{a}{un objecte terminal de la categoria de cons $\mathrm{Con}(D)$.}
  \opcio{b}{qualsevol con sobre $D$.}
  \opcio{c}{un objecte inicial de $\cat{C}$.}
  \opcio{d}{el producte de tots els objectes.}
\end{enumerate}
\feedback
\justificacio{Un con límit és un con pel qual factoritza de manera única qualsevol altre con: és a dir, un objecte terminal de $\mathrm{Con}(D)$. D'aquí ve la unicitat llevat d'un únic isomorfisme compatible amb les projeccions.}

\pregunta Quina categoria índex dona lloc al \textbf{producte} de dos objectes?
\begin{enumerate}
  \opcio{a}{La categoria buida.}
  \opcio{b}{Dues fletxes paral·leles.}
  \opcio{c}{La forma de racó.}
  \opcio{d}{La categoria discreta amb dos objectes.}
\end{enumerate}
\feedback
\justificacio{Un con sobre un parell d'objectes, sense cap fletxa entre ells, és una parella de fletxes sense cap condició: la propietat del producte. Dues fletxes paral·leles donen l'equalitzador, i el racó, el pullback.}

\pregunta A $\cat{Set}$, l'\textbf{equalitzador} de $f,g\colon A\to B$ és:
\begin{enumerate}
  \opcio{a}{el producte cartesià $A\times B$.}
  \opcio{b}{la unió disjunta.}
  \opcio{c}{el quocient de $B$.}
  \opcio{d}{el subconjunt $\{a\in A\mid f(a)=g(a)\}$.}
\end{enumerate}
\feedback
\justificacio{Una fletxa $h\colon X\to A$ amb $f\comp h=g\comp h$ pren valors dins d'aquest subconjunt, i hi factoritza de manera única. L'opció (c) és el coequalitzador, el concepte dual.}

\pregunta El \textbf{pullback} de $f\colon A\to C$ i $g\colon B\to C$ a $\cat{Set}$ és:
\begin{enumerate}
  \opcio{a}{$A\times B$ sencer.}
  \opcio{b}{$\{(a,b)\in A\times B\mid f(a)=g(b)\}$.}
  \opcio{c}{la unió $A\cup B$.}
  \opcio{d}{el quocient $A/B$.}
\end{enumerate}
\feedback
\justificacio{Una parella $x_1\colon X\to A$, $x_2\colon X\to B$ amb $f\comp x_1=g\comp x_2$ factoritza de manera única per $x\mapsto(x_1(x),x_2(x))$, que pren valors en aquest subconjunt. Tot $A\times B$ només és el pullback quan $C$ té un sol element.}

\pregunta L'\textbf{objecte terminal} és el límit de quin diagrama?
\begin{enumerate}
  \opcio{a}{El diagrama de dues fletxes paral·leles.}
  \opcio{b}{El diagrama d'un sol objecte.}
  \opcio{c}{El diagrama buit.}
  \opcio{d}{El diagrama de racó.}
\end{enumerate}
\feedback
\justificacio{Un con sobre el diagrama buit és només un vèrtex, sense projeccions. Un con límit és, doncs, un objecte al qual arriba una única fletxa des de cada objecte.}

\pregunta Un \textbf{colímit} d'un diagrama $D$ és:
\begin{enumerate}
  \opcio{a}{un objecte inicial a la categoria de cocons (un límit a la categoria oposada).}
  \opcio{b}{un altre nom per al límit.}
  \opcio{c}{sempre l'objecte terminal.}
  \opcio{d}{un con.}
\end{enumerate}
\feedback
\justificacio{És la noció dual: un cocon universal, és a dir, un objecte inicial de la categoria de cocons, o equivalentment un límit del diagrama corresponent a $\cat{C}\op$. No és un límit de la mateixa categoria.}

\pregunta Si $\cat{C}$ té límits de forma $\cat{I}$ i $\cat{A}$ és petita, com es calcula el límit d'un diagrama $F\colon\cat{I}\to\cat{C}^{\cat{A}}$?
\begin{enumerate}
  \opcio{a}{Prenent el límit només sobre els objectes terminals de $\cat{A}$.}
  \opcio{b}{Aplicant el functor de parts a cada $F(i)$.}
  \opcio{c}{No en té: les categories de functors no tenen límits.}
  \opcio{d}{Punt a punt: $\bigl(\varprojlim F\bigr)(a)\cong\varprojlim_i F(i)(a)$.}
\end{enumerate}
\feedback
\justificacio{Els límits a $\cat{C}^{\cat{A}}$ es calculen objecte per objecte a $\cat{C}$; equivalentment, cada functor d'avaluació $\mathrm{ev}_a$ els preserva.}

\pregunta Quina és la caracterització dels \textbf{límits finits}?
\begin{enumerate}
  \opcio{a}{N'hi ha prou amb tenir coproductes.}
  \opcio{b}{Cal tenir tots els límits de forma petita.}
  \opcio{c}{Equivalen a tenir objecte terminal, productes binaris i equalitzadors (o terminal i pullbacks).}
  \opcio{d}{No es poden construir a partir de casos més simples.}
\end{enumerate}
\feedback
\justificacio{És el teorema de caracterització: el límit d'un diagrama finit es construeix com un equalitzador de dues fletxes entre productes finits. Amb terminal i pullbacks s'obtenen productes i equalitzadors.}

\pregunta Quan \textbf{preserva} un functor $F$ el límit d'un diagrama $D$?
\begin{enumerate}
  \opcio{a}{Quan $(F(L),F\lambda)$ és un límit de $F\comp D$.}
  \opcio{b}{Sempre, per definició de functor.}
  \opcio{c}{Quan $F$ és constant.}
  \opcio{d}{Quan $F$ no té adjunts.}
\end{enumerate}
\feedback
\justificacio{És la definició de preservació. No tot functor preserva límits: per exemple, l'oblit $\cat{Grp}\to\cat{Set}$ no preserva els coproductes, que són límits a l'oposada.}

\pregunta El functor oblit $U\colon\cat{Grp}\to\cat{Set}$:
\begin{enumerate}
  \opcio{a}{no preserva cap límit ni cap colímit.}
  \opcio{b}{preserva tots els colímits, però no tots els límits.}
  \opcio{c}{preserva tots els límits, però no tots els colímits; per exemple, no preserva tots els coproductes.}
  \opcio{d}{preserva tots els límits i tots els colímits.}
\end{enumerate}
\feedback
\justificacio{El conjunt subjacent d'un límit de grups és el límit dels conjunts subjacents (famílies compatibles). En canvi, el coproducte de grups és el producte lliure, que no és la unió disjunta dels conjunts subjacents. Més endavant es veurà que $U$ és un adjunt per la dreta.}

\pregunta En una categoria localment petita $\cat{C}$, el functor representable
\[
\Hom(A,-)\colon\cat{C}\to\cat{Set}
\]
sempre:
\begin{enumerate}
  \opcio{a}{no és un functor.}
  \opcio{b}{preserva els límits que existeixen.}
  \opcio{c}{envia límits a colímits.}
  \opcio{d}{és constant.}
\end{enumerate}
\feedback
\justificacio{Les fletxes $A\to\varprojlim D$ es corresponen bijectivament amb els cons amb vèrtex $A$, que són les famílies compatibles de $\Hom(A,D_i)$, és a dir, els elements de $\varprojlim\Hom(A,D_i)$. L'opció (c) correspon al representable contravariant: $\Hom(-,B)$ envia colímits a límits.}

\pregunta A $\cat{Set}$, siguin $f\colon A\to B$ i $y\colon Y\to B$. Què fa la reindexació $f^{*}$?
\begin{enumerate}
  \opcio{a}{Envia $Y\to B$ al coproducte $A+Y$.}
  \opcio{b}{Envia $Y\to B$ al pullback $A\times_B Y\to A$.}
  \opcio{c}{Envia $Y\to B$ al coequalitzador de $f$ i $y$.}
  \opcio{d}{Només està definida quan $f$ és un isomorfisme.}
\end{enumerate}
\feedback
\justificacio{La fibra de $A\times_B Y\to A$ sobre $a$ està en bijecció amb la fibra $y^{-1}(f(a))$: la reindexació és la substitució al llarg de $f$.}

\pregunta En el diagrama següent, quina propietat universal fa de $P$ el \textbf{pullback} de $f$ i $g$?
\[
\begin{tikzpicture}[baseline=(m.center)]
  \matrix (m) [matrix of math nodes, row sep=2.6em, column sep=2.6em]
    { P & A \\ B & C \\ };
  \draw[-{Stealth[scale=1.1]}] (m-1-1) -- node[above] {$p_1$} (m-1-2);
  \draw[-{Stealth[scale=1.1]}] (m-1-1) -- node[left] {$p_2$} (m-2-1);
  \draw[-{Stealth[scale=1.1]}] (m-1-2) -- node[right] {$f$} (m-2-2);
  \draw[-{Stealth[scale=1.1]}] (m-2-1) -- node[below] {$g$} (m-2-2);
\end{tikzpicture}
\]
\begin{enumerate}
  \opcio{a}{Que per a tot objecte $X$ amb fletxes $x_1\colon X\to A$, $x_2\colon X\to B$ tals que $f\comp x_1=g\comp x_2$, hi ha un únic $u\colon X\to P$ amb $p_1\comp u=x_1$ i $p_2\comp u=x_2$.}
  \opcio{b}{Que $f\comp p_1=g\comp p_2$ i res més.}
  \opcio{c}{Que $P=A\times B$ sempre.}
  \opcio{d}{Que $f$ i $g$ són isomorfismes.}
\end{enumerate}
\feedback
\justificacio{És la propietat universal del pullback: el quadrat commuta i qualsevol altre quadrat commutatiu sobre $f$ i $g$ hi factoritza de manera única. L'opció (b) només diu que el quadrat commuta, i això no basta.}

\end{enumerate}
\clauestatica
\finalitzaTest

\chapter[Adjuncions]{Adjuncions}

Aquest test recull els conceptes essencials del capítol. Cada pregunta té una única resposta correcta.

\iniciTest{adj}{b,c,b,d,a,b,c,d,b,c,a,d,a,a}

\begin{enumerate}

\pregunta Què vol dir que $f\dashv g$ entre dos preordres?
\begin{enumerate}
  \opcio{a}{Que $f=g$.}
  \opcio{b}{Que $f(x)\leq y$ si i només si $x\leq g(y)$, per a tot $x,y$.}
  \opcio{c}{Que $f$ i $g$ són inverses.}
  \opcio{d}{Que $f$ i $g$ són constants.}
\end{enumerate}
\feedback
\justificacio{És la definició d'adjunció entre preordres: $f$ és adjunt per l'esquerra de $g$. No cal que siguin inverses: $f(x)\leq y\Leftrightarrow x\leq g(y)$ és més feble que $g\comp f=\id$ i $f\comp g=\id$.}

\pregunta Fixada una fórmula $A$, quina parella d'operacions forma una adjunció al preordre $\cat{Prov}_T$ d'un sistema proposicional?
\begin{enumerate}
  \opcio{a}{La negació amb ella mateixa.}
  \opcio{b}{La disjunció amb la conjunció.}
  \opcio{c}{$A\wedge-$ (esquerra) amb $A\to-$ (dreta).}
  \opcio{d}{La implicació amb la negació.}
\end{enumerate}
\feedback
\justificacio{$A\wedge B\vdash_T C$ si i només si $B\vdash_T A\to C$: és la regla de la deducció (introducció i eliminació de la implicació). Les altres parelles no satisfan cap equivalència d'aquesta forma.}

\pregunta Una \textbf{adjunció} $F\dashv G$ entre categories és:
\begin{enumerate}
  \opcio{a}{una igualtat de functors.}
  \opcio{b}{una bijecció natural $\Hom(F(A),B)\cong\Hom(A,G(B))$.}
  \opcio{c}{un isomorfisme de categories.}
  \opcio{d}{una transformació natural qualsevol.}
\end{enumerate}
\feedback
\justificacio{És la definició: una família de bijeccions $\theta_{A,B}$, natural en $A$ i en $B$. Un isomorfisme de categories és un cas molt particular, i una transformació natural qualsevol no relaciona homconjunts.}

\pregunta Sigui $\cat{I}$ una categoria petita i suposem que $\cat{C}$ té tots els límits i colímits de forma $\cat{I}$. Quina és la relació entre el functor diagonal $\Delta\colon\cat{C}\to\cat{C}^{\cat{I}}$, el límit i el colímit?
\begin{enumerate}
  \opcio{a}{$\varprojlim_{\cat{I}}\dashv\Delta\dashv\colimop_{\cat{I}}$.}
  \opcio{b}{$\Delta\dashv\colimop_{\cat{I}}$ i $\Delta\dashv\varprojlim_{\cat{I}}$.}
  \opcio{c}{El límit i el colímit són tots dos adjunts per l'esquerra de $\Delta$.}
  \opcio{d}{$\colimop_{\cat{I}}\dashv\Delta\dashv\varprojlim_{\cat{I}}$.}
\end{enumerate}
\feedback
\justificacio{Un morfisme $\Delta A\Rightarrow D$ és un con, que correspon a una fletxa $A\to\varprojlim D$; un morfisme $D\Rightarrow\Delta A$ és un cocon, que correspon a una fletxa $\colimop D\to A$.}

\pregunta La \textbf{unitat} d'una adjunció $F\dashv G$ és una transformació natural:
\begin{enumerate}
  \opcio{a}{$\id\Rightarrow G F$.}
  \opcio{b}{$F\Rightarrow G$.}
  \opcio{c}{$F G\Rightarrow\id$.}
  \opcio{d}{$G\Rightarrow F$.}
\end{enumerate}
\feedback
\justificacio{$\eta_A=\theta(\id_{F(A)})\colon A\to G(F(A))$. L'opció (c) és la counitat. Les opcions (b) i (d) no tenen sentit en general, perquè $F$ i $G$ van en sentits contraris i no comparteixen domini.}

\pregunta La \textbf{counitat} d'una adjunció $F\dashv G$ és:
\begin{enumerate}
  \opcio{a}{$\id\Rightarrow G F$.}
  \opcio{b}{$F G\Rightarrow\id$.}
  \opcio{c}{una igualtat.}
  \opcio{d}{un objecte terminal.}
\end{enumerate}
\feedback
\justificacio{$\varepsilon_B=\theta^{-1}(\id_{G(B)})\colon F(G(B))\to B$. L'opció (a) és la unitat. La counitat és una transformació natural, no una igualtat ni un objecte.}

\pregunta El transposat de $\varphi\colon F(A)\to B$ sota l'adjunció és:
\begin{enumerate}
  \opcio{a}{$\varepsilon_B\comp F(\varphi)$.}
  \opcio{b}{$\varphi$ mateix.}
  \opcio{c}{$G(\varphi)\comp\eta_A$.}
  \opcio{d}{$\eta_A\comp G(\varphi)$.}
\end{enumerate}
\feedback
\justificacio{Comprova els tipus: amb $\varphi\colon F(A)\to B$, $G(\varphi)\colon G(F(A))\to G(B)$, i compondre amb $\eta_A\colon A\to G(F(A))$ dona $A\to G(B)$. L'opció (a) és la fórmula del transposat \emph{invers}, per a $\psi\colon A\to G(B)$; l'opció (d) no és componible.}

\pregunta Les \textbf{identitats triangulars} relacionen:
\begin{enumerate}
  \opcio{a}{dos objectes terminals.}
  \opcio{b}{dos límits.}
  \opcio{c}{dues categories isomorfes.}
  \opcio{d}{la unitat i la counitat, i caracteritzen l'adjunció.}
\end{enumerate}
\feedback
\justificacio{Són $\varepsilon_{F(A)}\comp F(\eta_A)=\id_{F(A)}$ i $G(\varepsilon_B)\comp\eta_{G(B)}=\id_{G(B)}$. A més, dues transformacions naturals que les satisfan defineixen una adjunció: és la segona presentació.}

\pregunta Sigui $U\colon\cat{Mon}\to\cat{Set}$ el functor oblit. Quin fet mostra que $U$ no preserva tots els colímits?
\begin{enumerate}
  \opcio{a}{No preserva productes.}
  \opcio{b}{El monoide trivial és inicial a $\cat{Mon}$, però el seu conjunt subjacent no és $\varnothing$.}
  \opcio{c}{$\cat{Mon}$ no té objecte inicial.}
  \opcio{d}{Cap adjunt per la dreta no pot preservar colímits.}
\end{enumerate}
\feedback
\justificacio{L'objecte inicial és el colímit del diagrama buit. L'objecte inicial de $\cat{Set}$ és $\varnothing$, però $U$ envia l'inicial de $\cat{Mon}$ a un conjunt d'un element; per tant, no preserva aquest colímit.}

\pregunta Els adjunts són únics:
\begin{enumerate}
  \opcio{a}{de manera estricta (igualtat).}
  \opcio{b}{mai.}
  \opcio{c}{llevat d'isomorfisme natural.}
  \opcio{d}{només en preordres.}
\end{enumerate}
\feedback
\justificacio{Si $F\dashv G$ i $F\dashv G'$, aleshores $\Hom(-,G(B))\cong\Hom(-,G'(B))$ naturalment, i Yoneda dona $G(B)\cong G'(B)$, naturalment en $B$. No hi ha igualtat estricta: qualsevol functor isomorf a un adjunt també ho és.}

\pregunta Què diu el teorema \enquote{RAPL}?
\begin{enumerate}
  \opcio{a}{Que els adjunts per la dreta preserven límits.}
  \opcio{b}{Que tot functor és adjunt.}
  \opcio{c}{Que els límits no existeixen.}
  \opcio{d}{Que els adjunts per l'esquerra preserven límits.}
\end{enumerate}
\feedback
\justificacio{RAPL abreuja \emph{right adjoints preserve limits}. L'opció (d) confon els costats: els adjunts per l'esquerra preserven colímits.}

\pregunta Per què el functor graf subjacent $U\colon\cat{Cat}\to\cat{Graph}$ preserva límits?
\begin{enumerate}
  \opcio{a}{Perquè és ple.}
  \opcio{b}{Perquè és adjunt per l'esquerra.}
  \opcio{c}{Perquè tots els functors preserven límits.}
  \opcio{d}{Perquè és adjunt per la dreta de $\Free$.}
\end{enumerate}
\feedback
\justificacio{$\Free\dashv U$, i els adjunts per la dreta preserven límits.}

\pregunta L'adjunció \enquote{producte-exponencial} $-\times B\dashv(-)^B$ correspon a:
\begin{enumerate}
  \opcio{a}{la currificació $\Hom(A\times B,C)\cong\Hom(A,C^B)$.}
  \opcio{b}{la dualitat.}
  \opcio{c}{la negació.}
  \opcio{d}{la composició de functors.}
\end{enumerate}
\feedback
\justificacio{Fixat $B$, una fletxa $A\times B\to C$ correspon a una fletxa $A\to C^B$, naturalment en $A$ i en $C$. És exactament la bijecció de l'adjunció $-\times B\dashv(-)^B$.}

\pregunta Si $f^{*}$ admet adjunts a tots dos costats, quin paper tenen $\exists_f$ i $\forall_f$?
\begin{enumerate}
  \opcio{a}{Com l'adjunt per l'esquerra i l'adjunt per la dreta de $f^{*}$, respectivament.}
  \opcio{b}{Com objectes terminals.}
  \opcio{c}{Com functors constants.}
  \opcio{d}{Com isomorfismes de categories.}
\end{enumerate}
\feedback
\justificacio{A $\cat{Set}$, $\exists_f$ és la imatge i $\forall_f(S)$ el conjunt dels $y$ amb fibra continguda en $S$; per a una projecció, són la quantificació existencial i la universal.}

\end{enumerate}
\clauestatica
\finalitzaTest

\chapter[Cartesianes tancades]{\texorpdfstring{Categories cartesianes\\ tancades}{Categories cartesianes tancades}}

Aquest test recull els conceptes essencials del capítol. Cada pregunta té una única resposta correcta.

\iniciTest{cct}{a,c,d,b,c,a,d,a,b,c,d,b}

\begin{enumerate}

\pregunta Una categoria és \textbf{cartesiana} quan té:
\begin{enumerate}
  \opcio{a}{objecte terminal i productes binaris (és a dir, tots els productes finits).}
  \opcio{b}{tots els colímits.}
  \opcio{c}{només un objecte inicial.}
  \opcio{d}{exponencials però no productes.}
\end{enumerate}
\feedback
\justificacio{El terminal és el producte buit, i els productes de qualsevol família finita s'obtenen per inducció a partir dels binaris. Els colímits i els exponencials no formen part de la definició.}

\pregunta Quina bijecció natural expressa la propietat universal de l'exponencial $C^{B}$?
\begin{enumerate}
  \opcio{a}{$\Hom(A,C)\cong\Hom(B,C)$.}
  \opcio{b}{$\Hom(A+B,C)\cong\Hom(A,C)\times\Hom(B,C)$.}
  \opcio{c}{$\Hom(A\times B,C)\cong\Hom(A,C^{B})$.}
  \opcio{d}{$\Hom(A,C\times B)\cong\Hom(A,C^{B})$.}
\end{enumerate}
\feedback
\justificacio{És la currificació $f\mapsto\lambda f$, natural en $A$ i en $C$: expressa l'adjunció $-\times B\dashv(-)^{B}$.}

\pregunta En una CCC, quin morfisme indueix $f\colon B'\to B$ entre exponencials amb el mateix $C$?
\begin{enumerate}
  \opcio{a}{$C^{B'}\to C^{B}$, perquè l'exponencial és covariant en totes dues variables.}
  \opcio{b}{Cap: l'exponencial no depèn functorialment de $B$.}
  \opcio{c}{$B^{C}\to B'^{C}$.}
  \opcio{d}{$C^{B}\to C^{B'}$: l'exponencial és contravariant en l'exponent.}
\end{enumerate}
\feedback
\justificacio{$C^{f}$ és la transposada de $C^{B}\times B'\xrightarrow{\id\times f}C^{B}\times B\xrightarrow{\ev}C$.}

\pregunta La \textbf{currificació} $\lambda f$ d'un morfisme $f\colon A\times B\to C$ és l'únic morfisme que compleix:
\begin{enumerate}
  \opcio{a}{$\lambda f\comp\ev=f$.}
  \opcio{b}{$\ev\comp(\lambda f\times\id_B)=f$.}
  \opcio{c}{$\lambda f=f$.}
  \opcio{d}{$\lambda f\comp\pi_1=f$.}
\end{enumerate}
\feedback
\justificacio{És el triangle de la definició: currificar i després avaluar torna $f$. L'opció (a) no és ni tan sols componible, perquè $\ev$ surt de $C^B\times B$ i $\lambda f$ hi arriba només des de $A$.}

\pregunta Una categoria és \textbf{cartesiana tancada} quan és cartesiana i, a més, per a cada objecte $B$:
\begin{enumerate}
  \opcio{a}{$B$ és terminal.}
  \opcio{b}{$-\times B$ és un isomorfisme.}
  \opcio{c}{el functor $-\times B$ té adjunt per la dreta $(-)^{B}$.}
  \opcio{d}{$-\times B$ té adjunt per l'esquerra.}
\end{enumerate}
\feedback
\justificacio{Tancar cartesianament vol dir que la currificació és sempre disponible: per a cada $B$, $-\times B\dashv(-)^B$. Un adjunt per l'esquerra de $-\times B$ seria una altra cosa, i no existeix en general.}

\pregunta A $\cat{Set}$, l'exponencial $C^{B}$ és:
\begin{enumerate}
  \opcio{a}{el conjunt de \emph{totes} les aplicacions $B\to C$.}
  \opcio{b}{el producte cartesià $C\times B$.}
  \opcio{c}{la unió disjunta $C+B$.}
  \opcio{d}{el conjunt de bijeccions $B\to C$.}
\end{enumerate}
\feedback
\justificacio{L'avaluació és $\ev(g,b)=g(b)$, i la currificació de $f\colon A\times B\to C$ és $a\mapsto(b\mapsto f(a,b))$. Cal prendre totes les aplicacions, no només les bijeccions, perquè $\lambda f(a)$ pot ser qualsevol aplicació.}

\pregunta En una \textbf{àlgebra de Heyting}, l'objecte exponencial $b\Rightarrow c$ és:
\begin{enumerate}
  \opcio{a}{el suprem $b\vee c$.}
  \opcio{b}{la negació $\neg b$.}
  \opcio{c}{l'ínfim $b\wedge c$.}
  \opcio{d}{l'adjunt per la dreta de $-\wedge b$, caracteritzat per $a\wedge b\leq c\Leftrightarrow a\leq(b\Rightarrow c)$.}
\end{enumerate}
\feedback
\justificacio{En un preordre, l'exponencial és l'adjunt per la dreta del producte, que és l'ínfim. En una àlgebra de Boole, a més, $b\Rightarrow c=\neg b\vee c$, però aquesta fórmula no val en una àlgebra de Heyting qualsevol.}

\pregunta Quina categoria \emph{no} és, en general, cartesiana tancada?
\begin{enumerate}
  \opcio{a}{$\cat{Top}$, els espais topològics.}
  \opcio{b}{una àlgebra de Heyting.}
  \opcio{c}{$\cat{Set}$.}
  \opcio{d}{una categoria de prefeixos $\cat{Set}^{\cat{C}\op}$.}
\end{enumerate}
\feedback
\justificacio{A $\cat{Top}$, l'avaluació no sempre té un espai de funcions contínues universal. $\cat{Set}$, les àlgebres de Heyting com a preordres i les categories de prefeixos sí que són cartesianes tancades.}

\pregunta En la semàntica del càlcul lambda simplement tipat, el \textbf{tipus funció} $\sigma\to\tau$ s'interpreta amb:
\begin{enumerate}
  \opcio{a}{el producte $\llbracket\sigma\rrbracket\times\llbracket\tau\rrbracket$.}
  \opcio{b}{l'exponencial $\llbracket\tau\rrbracket^{\llbracket\sigma\rrbracket}$.}
  \opcio{c}{l'objecte terminal.}
  \opcio{d}{el coproducte.}
\end{enumerate}
\feedback
\justificacio{Un terme de tipus $\sigma\to\tau$ és una funció, i en una CCC les funcions de $\llbracket\sigma\rrbracket$ a $\llbracket\tau\rrbracket$ formen l'exponencial. El producte interpreta, en canvi, el tipus parella $\sigma\times\tau$.}

\pregunta En aquesta semàntica, l'\textbf{abstracció} $\lambda x.\,t$ es correspon amb:
\begin{enumerate}
  \opcio{a}{una projecció.}
  \opcio{b}{la diagonal.}
  \opcio{c}{la currificació de la interpretació de $t$.}
  \opcio{d}{l'objecte terminal.}
\end{enumerate}
\feedback
\justificacio{Si $t$ té tipus $\tau$ en el context $\Gamma,x{:}\sigma$, s'interpreta com un morfisme $\llbracket\Gamma\rrbracket\times\llbracket\sigma\rrbracket\to\llbracket\tau\rrbracket$. La seva currificació és un morfisme $\llbracket\Gamma\rrbracket\to\llbracket\tau\rrbracket^{\llbracket\sigma\rrbracket}$, que interpreta $\lambda x.\,t$.}

\pregunta La regla de conversió $(\beta)$ del càlcul lambda es correspon categòricament amb:
\begin{enumerate}
  \opcio{a}{la unicitat de l'objecte terminal.}
  \opcio{b}{l'associativitat del producte.}
  \opcio{c}{la commutativitat del coproducte.}
  \opcio{d}{l'equació de l'avaluació $\ev\comp(\lambda f\times\id_B)=f$, de la qual es dedueix.}
\end{enumerate}
\feedback
\justificacio{Aplicar una abstracció a un argument és avaluar una currificació, i l'equació $\ev\comp(\lambda f\times\id)=f$ diu que el resultat és $f$, que és el que afirma la regla $(\beta)$. La regla $(\eta)$ correspon a la segona llei, $\lambda(\ev\comp(g\times\id))=g$.}

\pregunta La correspondència de \textbf{Curry--Howard--Lambek} relaciona:
\begin{enumerate}
  \opcio{a}{grups amb anells.}
  \opcio{b}{el càlcul lambda simplement tipat amb les categories cartesianes tancades.}
  \opcio{c}{límits amb colímits.}
  \opcio{d}{monoides amb preordres.}
\end{enumerate}
\feedback
\justificacio{Els tipus s'interpreten com a objectes, els termes com a morfismes i les conversions com a igualtats. En la lectura lògica, les proposicions del fragment $(\top,\wedge,\to)$ són els tipus, i les demostracions, els termes.}

\end{enumerate}
\clauestatica
\finalitzaTest

\chapter[Subobjectes, imatges i factoritzacions]{\texorpdfstring{Subobjectes, imatges\\ i factoritzacions}{Subobjectes, imatges i factoritzacions}}

Aquest test recull els conceptes essencials del capítol. Cada pregunta té una única resposta correcta.

\iniciTest{sub}{c,b,d,b,c,d,a,c,a,d,c,b,d,a,b,a,c}

\begin{enumerate}

\pregunta Un \textbf{subobjecte} d'un objecte $A$ és:
\begin{enumerate}
  \opcio{a}{qualsevol morfisme cap a $A$.}
  \opcio{b}{un element de $A$.}
  \opcio{c}{una classe d'equivalència de monomorfismes amb codomini $A$.}
  \opcio{d}{un epimorfisme des de $A$.}
\end{enumerate}
\feedback
\justificacio{Dos monomorfismes que difereixen en un isomorfisme dels dominis representen el mateix subobjecte. Un morfisme qualsevol no basta: cal que sigui mono, el paper que a $\cat{Set}$ fan les injeccions.}

\pregunta Sigui $H$ un subgraf d'un graf $G$ i $F\colon G'\to G$ un morfisme de grafs. Què és la reindexació $F^{*}H$?
\begin{enumerate}
  \opcio{a}{La imatge de $H$ per $F$.}
  \opcio{b}{El subgraf de $G'$ format per les antiimatges dels vèrtexs i de les arestes de $H$.}
  \opcio{c}{El graf $H$ mateix, vist dins de $G'$.}
  \opcio{d}{El producte de grafs $H\times G'$.}
\end{enumerate}
\feedback
\justificacio{La reindexació es calcula component a component: $F^{*}(V',E')=(F_V^{-1}(V'),F_E^{-1}(E'))$, que és tancat per origen i destí perquè $F$ els preserva.}

\pregunta A $\cat{Set}$, l'ordre parcial $\Sub(A)$ és isomorf a:
\begin{enumerate}
  \opcio{a}{el conjunt $A$.}
  \opcio{b}{el monoide de les aplicacions $A\to A$.}
  \opcio{c}{el grup de permutacions de $A$.}
  \opcio{d}{el conjunt de parts $\mathcal{P}(A)$, ordenat per inclusió.}
\end{enumerate}
\feedback
\justificacio{Dos monomorfismes cap a $A$ són equivalents si i només si tenen la mateixa imatge, i la imatge determina la classe. L'ordre correspon a la inclusió de les imatges.}

\pregunta La \textbf{parella nucli} d'una fletxa $f\colon A\to B$, quan existeix, s'obté com:
\begin{enumerate}
  \opcio{a}{l'equalitzador de $f$ i $\id_A$.}
  \opcio{b}{el pullback de $f$ amb ella mateixa.}
  \opcio{c}{el coequalitzador de $f$ amb ella mateixa.}
  \opcio{d}{la factorització imatge de $f$.}
\end{enumerate}
\feedback
\justificacio{És el parell de projeccions $p_1,p_2\colon R\to A$ del pullback de $f$ amb $f$; a $\cat{Set}$, $R=\{(a,a')\mid f(a)=f(a')\}$. Quan existeix la parella nucli, un epimorfisme regular n'és el coequalitzador.}

\pregunta Lògicament, l'ordre $[m]\leq[n]$ a $\Sub(A)$ s'interpreta com:
\begin{enumerate}
  \opcio{a}{la negació del predicat.}
  \opcio{b}{la igualtat de predicats.}
  \opcio{c}{la implicació entre predicats sobre $A$.}
  \opcio{d}{la disjunció de predicats.}
\end{enumerate}
\feedback
\justificacio{$[m]\leq[n]$ vol dir que $m$ factoritza a través de $n$: tot el que satisfà el predicat $m$ satisfà el predicat $n$. És l'ordre de la deduïbilitat entre predicats sobre $A$.}

\pregunta El \textbf{pullback} d'un monomorfisme al llarg de qualsevol morfisme és:
\begin{enumerate}
  \opcio{a}{un epimorfisme.}
  \opcio{b}{un isomorfisme.}
  \opcio{c}{mai un monomorfisme.}
  \opcio{d}{sempre un monomorfisme.}
\end{enumerate}
\feedback
\justificacio{Si $m'\comp u=m'\comp v$ al pullback, aleshores $u$ i $v$ tenen la mateixa composició amb totes dues projeccions (la segona, per la cancel·lació de $m$), i per la unicitat del pullback $u=v$.}

\pregunta El \textbf{functor de reindexació} $f^{*}\colon\Sub(B)\to\Sub(A)$ associat a $f\colon A\to B$ envia un subobjecte al seu:
\begin{enumerate}
  \opcio{a}{pullback al llarg de $f$.}
  \opcio{b}{coproducte amb $A$.}
  \opcio{c}{producte amb $B$.}
  \opcio{d}{imatge directa.}
\end{enumerate}
\feedback
\justificacio{És la definició: el pullback d'un mono és mono, i en resulta una aplicació monòtona ben definida entre classes. La imatge directa (d) correspon a l'existencial, que va en sentit contrari.}

\pregunta A $\cat{Set}$, la reindexació $f^{*}(S)$ d'un subconjunt $S\subseteq B$ és:
\begin{enumerate}
  \opcio{a}{la imatge directa $f(S)$.}
  \opcio{b}{el complement de $S$.}
  \opcio{c}{l'antiimatge $f^{-1}(S)$.}
  \opcio{d}{el producte $S\times A$.}
\end{enumerate}
\feedback
\justificacio{El pullback de la inclusió $S\hookrightarrow B$ al llarg de $f$ és $\{a\in A\mid f(a)\in S\}$. La imatge directa (a) és l'existencial $\exists_f$, l'adjunt per l'esquerra de $f^{*}$.}

\pregunta Lògicament, la reindexació $f^{*}$ interpreta:
\begin{enumerate}
  \opcio{a}{la substitució d'una variable per un terme.}
  \opcio{b}{la quantificació existencial.}
  \opcio{c}{la negació.}
  \opcio{d}{la conjunció.}
\end{enumerate}
\feedback
\justificacio{Si $\varphi(y)$ és un predicat sobre $B$ i $f$ interpreta un terme, $f^{*}\varphi$ és $\varphi(f(x))$: substituir és fer el pullback. La quantificació existencial correspon a la imatge.}

\pregunta La \textbf{intersecció} de dos subobjectes $[m]\wedge[n]$ a $\Sub(A)$ es calcula com:
\begin{enumerate}
  \opcio{a}{el coproducte de $S$ i $T$.}
  \opcio{b}{la unió de les imatges.}
  \opcio{c}{l'exponencial $T^{S}$.}
  \opcio{d}{el pullback de $m$ i $n$ sobre $A$.}
\end{enumerate}
\feedback
\justificacio{Una fletxa factoritza a través de $m$ i de $n$ si i només si factoritza a través del seu pullback. Per tant, el pullback és el major subobjecte per sota de tots dos, i interpreta la conjunció.}

\pregunta En una categoria regular, considerem un quadrat cartesià amb $h\colon A'\to A$, $f'\colon A'\to B'$, $f\colon A\to B$ i $g\colon B'\to B$, de manera que $f\comp h=g\comp f'$. Quina igualtat és la condició de Beck--Chevalley, amb els tipus correctes?
\begin{enumerate}
  \opcio{a}{$\exists_f\comp g^{*}=h^{*}\comp\exists_{f'}$.}
  \opcio{b}{$h^{*}\comp\exists_f=\exists_{f'}\comp g^{*}\colon\Sub(B)\to\Sub(A')$.}
  \opcio{c}{$g^{*}\comp\exists_f=\exists_{f'}\comp h^{*}\colon\Sub(A)\to\Sub(B')$.}
  \opcio{d}{$g^{*}\comp\exists_f=\exists_{f'}\comp h^{*}\colon\Sub(B)\to\Sub(A')$.}
\end{enumerate}
\feedback
\justificacio{Partint de $\Sub(A)$, hi ha dos camins fins a $\Sub(B')$: quantificar al llarg de $f$ i substituir al llarg de $g$, o substituir al llarg de $h$ i quantificar al llarg de $f'$. La condició diu que coincideixen: quantificar i substituir commuten. A (a) i (b) les composicions no existeixen ($\exists_f$ surt de $\Sub(A)$, no de $\Sub(B)$), i a (d) el domini i el codomini estan intercanviats.}

\pregunta Segons el capítol, en una categoria \textbf{regular} la reindexació $f^{*}$ té garantit:
\begin{enumerate}
  \opcio{a}{cap adjunt.}
  \opcio{b}{un adjunt per l'esquerra $\exists_f$ (l'universal $\forall_f$ requereix estructura addicional).}
  \opcio{c}{un adjunt per la dreta $\forall_f$, però no per l'esquerra.}
  \opcio{d}{adjunts per tots dos costats, sempre.}
\end{enumerate}
\feedback
\justificacio{L'existencial és la imatge, i la regularitat en garanteix l'existència i l'estabilitat. L'adjunt per la dreta $\forall_f$ és estructura addicional, que tenen, per exemple, els topos.}

\pregunta Una \textbf{factorització imatge} de $f$ descompon $f$ com:
\begin{enumerate}
  \opcio{a}{un mono seguit d'un epi qualsevol.}
  \opcio{b}{dos isomorfismes.}
  \opcio{c}{dos monomorfismes.}
  \opcio{d}{un epimorfisme regular seguit d'un monomorfisme.}
\end{enumerate}
\feedback
\justificacio{A $\cat{Set}$, l'epi regular és l'aplicació exhaustiva sobre la imatge, i el mono és la inclusió de la imatge. La factorització és única llevat d'un únic isomorfisme compatible.}

\pregunta La \textbf{imatge} $\Imatge(f)$ es caracteritza com:
\begin{enumerate}
  \opcio{a}{el menor subobjecte del codomini a través del qual $f$ es factoritza.}
  \opcio{b}{el major subobjecte del domini.}
  \opcio{c}{el nucli de $f$.}
  \opcio{d}{el coproducte de domini i codomini.}
\end{enumerate}
\feedback
\justificacio{La imatge és un subobjecte pel qual factoritza $f$, i és per sota de qualsevol altre que tingui aquesta propietat. A $\cat{Set}$ és el subconjunt $f(A)\subseteq B$.}

\pregunta Una categoria és \textbf{regular} si té límits finits, factoritzacions imatge i, a més:
\begin{enumerate}
  \opcio{a}{tots els colímits.}
  \opcio{b}{els epimorfismes regulars són estables per pullback.}
  \opcio{c}{és cartesiana tancada.}
  \opcio{d}{tots els objectes són projectius.}
\end{enumerate}
\feedback
\justificacio{És la condició que fa estables les imatges per reindexació i, per tant, l'existencial compatible amb la substitució. Els colímits, el tancament cartesià i els objectes projectius no formen part de la definició.}

\pregunta En una categoria regular, la \textbf{quantificació existencial} $\exists_f[s]$ es defineix com:
\begin{enumerate}
  \opcio{a}{la imatge $\Imatge(f\comp s)$.}
  \opcio{b}{el pullback de $[s]$.}
  \opcio{c}{el complement de $[s]$.}
  \opcio{d}{l'exponencial de $[s]$.}
\end{enumerate}
\feedback
\justificacio{Es fa la composició $f\comp s$ i se'n pren la imatge a $B$; a $\cat{Set}$, $f(S)$. Aquest operador és l'adjunt per l'esquerra de $f^{*}$: $\exists_f[s]\leq[t]$ si i només si $[s]\leq f^{*}[t]$.}

\pregunta Sigui $\cat{C}$ una categoria regular i $f\colon A\to B$ un morfisme. Quina d'aquestes afirmacions \emph{no} es pot deduir només de la regularitat?
\begin{enumerate}
  \opcio{a}{Tota fletxa té una factorització imatge (única llevat d'iso).}
  \opcio{b}{La reindexació $f^{*}\colon\Sub(B)\to\Sub(A)$ té un adjunt per l'esquerra $\exists_f$.}
  \opcio{c}{Cada $\Sub(A)$ té unions binàries, i $f^{*}$ les preserva.}
  \opcio{d}{Cada $\Sub(A)$ té ínfims binaris, calculats per pullback, i $f^{*}$ els preserva.}
\end{enumerate}
\feedback
\justificacio{Les unions de subobjectes estables per reindexació no formen part de la definició de regularitat: són l'estructura de les categories coherents. El mateix passa amb l'adjunt per la dreta $\forall_f$. En canvi, (a) i (b) són part del que dona la regularitat, i (d) també val: els ínfims són pullbacks, i $f^{*}$ els preserva perquè és adjunt per la dreta de $\exists_f$.}

\end{enumerate}
\clauestatica
\finalitzaTest

\chapter[Classificadors de subobjectes i topos]{\texorpdfstring{Classificadors de subobjectes i introducció als topos}{Classificadors de subobjectes i introducció als topos}}

Aquest test recull els conceptes essencials del capítol. Cada pregunta té una única resposta correcta.

\iniciTest{top}{c,b,d,a,c,d,b,a,a,b,c,d,b}

\begin{enumerate}

\pregunta Sigui $m\colon U\rightarrowtail E$ un subobjecte classificat per $\chi_m\colon E\to\Omega$. Un element generalitzat $x\colon X\to E$ factoritza a través de $m$ exactament quan:
\begin{enumerate}
  \opcio{a}{$\chi_m\comp x$ és un isomorfisme.}
  \opcio{b}{$x$ és un monomorfisme.}
  \opcio{c}{$\chi_m\comp x=\mathsf{true}\comp\,!_X$.}
  \opcio{d}{$X$ és l'objecte terminal.}
\end{enumerate}
\feedback
\justificacio{És la propietat universal del quadrat cartesià de $m$: la pertinença d'un element generalitzat esdevé una igualtat de morfismes cap a $\Omega$.}

\pregunta A la categoria de prefeixos sobre $\cat{C}$, siguin $h\colon D\to C$ i $S$ un sedàs sobre $C$. Quin sedàs sobre $D$ és $\Omega(h)(S)$?
\begin{enumerate}
  \opcio{a}{$\{h\comp g\mid g\in S\}$.}
  \opcio{b}{$\{g\colon X\to D\mid h\comp g\in S\}$.}
  \opcio{c}{$S$ mateix.}
  \opcio{d}{El sedàs màxim sobre $D$.}
\end{enumerate}
\feedback
\justificacio{$\Omega$ actua per reindexació de sedassos: $\Omega(h)(S)=h^{*}S$, i aquest conjunt torna a ser un sedàs.}

\pregunta La propietat universal del classificador diu que cada subobjecte $U\rightarrowtail E$ determina:
\begin{enumerate}
  \opcio{a}{un isomorfisme $E\cong\Omega$.}
  \opcio{b}{cap morfisme.}
  \opcio{c}{diversos morfismes $E\to\Omega$.}
  \opcio{d}{un únic morfisme $\chi\colon E\to\Omega$ que fa un quadrat cartesià amb $\mathsf{true}$.}
\end{enumerate}
\feedback
\justificacio{És la definició de classificador: existència i unicitat de $\chi$, amb $U$ com a pullback de $\mathsf{true}$ al llarg de $\chi$. Que el quadrat només commuti no basta per a la unicitat.}

\pregunta L'existència d'un classificador de subobjectes equival a dir que el functor $\Sub(-)$ és:
\begin{enumerate}
  \opcio{a}{representable, amb objecte representant $\Omega$.}
  \opcio{b}{constant.}
  \opcio{c}{no functorial.}
  \opcio{d}{un isomorfisme.}
\end{enumerate}
\feedback
\justificacio{La bijecció $\Sub(E)\cong\Hom(E,\Omega)$, natural en $E$, és exactament una representació de $\Sub$, amb element universal $\mathsf{true}$.}

\pregunta El classificador de subobjectes, quan existeix, és únic:
\begin{enumerate}
  \opcio{a}{només a $\cat{Set}$.}
  \opcio{b}{mai.}
  \opcio{c}{llevat d'un únic isomorfisme compatible amb $\mathsf{true}$ (per Yoneda).}
  \opcio{d}{llevat d'homotopia.}
\end{enumerate}
\feedback
\justificacio{Dos objectes que representen el mateix functor són isomorfs per un únic isomorfisme compatible amb els elements universals: és el principi de Yoneda aplicat a $\Sub$.}

\pregunta Un \textbf{topos elemental} és una categoria amb límits finits, exponencials i:
\begin{enumerate}
  \opcio{a}{tots els colímits infinits.}
  \opcio{b}{un objecte inicial estricte.}
  \opcio{c}{una estructura de grup.}
  \opcio{d}{un classificador de subobjectes.}
\end{enumerate}
\feedback
\justificacio{Límits finits, exponencials i classificador de subobjectes: no es demanen colímits infinits. $\cat{FinSet}$, que no en té, és un topos.}

\pregunta Quina d'aquestes categories és un topos?
\begin{enumerate}
  \opcio{a}{la categoria dels grups.}
  \opcio{b}{$\cat{Set}$.}
  \opcio{c}{$\cat{Top}$.}
  \opcio{d}{un monoide no trivial vist com a categoria.}
\end{enumerate}
\feedback
\justificacio{$\cat{Set}$ té límits finits, exponencials i el classificador $\{0,1\}$. $\cat{Top}$ no és cartesiana tancada, i un monoide no trivial vist com a categoria no té ni tan sols tots els límits finits.}

\pregunta Per a $\cat{C}$ petita, la categoria de prefeixos $\cat{Set}^{\cat{C}\op}$ és:
\begin{enumerate}
  \opcio{a}{sempre un topos.}
  \opcio{b}{un topos només si $\cat{C}$ és un grup.}
  \opcio{c}{mai un topos.}
  \opcio{d}{cartesiana tancada però sense classificador.}
\end{enumerate}
\feedback
\justificacio{Té límits finits i exponencials (capítols de límits i de cartesianes tancades) i el classificador de sedassos. No cal cap condició sobre $\cat{C}$ més enllà de ser petita.}

\pregunta A $\cat{Graph}$, sigui $x\colon v\to w$ una aresta d'un graf $G$ i $H$ un subgraf que conté els vèrtexs $v$ i $w$ però no l'aresta $x$. A quina aresta del classificador $\Omega$ envia $\chi_H$ l'aresta $x$?
\begin{enumerate}
  \opcio{a}{Al bucle $\partial$ del vèrtex $1$, que registra \enquote{extrems sí, aresta no}.}
  \opcio{b}{Al bucle $\top$ del vèrtex $1$, perquè els dos extrems són a $H$.}
  \opcio{c}{A l'aresta $\sigma\colon1\to0$.}
  \opcio{d}{Enlloc: $\chi_H$ només actua sobre els vèrtexs.}
\end{enumerate}
\feedback
\justificacio{Com que $v,w\in H$, la imatge de $x$ ha d'anar de $1$ a $1$, i n'hi ha dues: $\partial$ i $\top$. Si fos $\top$, $x$ seria al pullback de $\mathsf{true}$, que ha de ser exactament $H$; per tant, va a $\partial$. L'opció (c) tindria el destí equivocat, i un morfisme de grafs actua també sobre les arestes.}

\pregunta En un topos, l'objecte de parts $\mathcal P(E)=\Omega^{E}$ satisfà, per a tot objecte $A$,
\[
\Hom(A,\mathcal P(E))\;\cong\;\Sub(A\times E).
\]
Quina interpretació expressa aquesta bijecció?
\begin{enumerate}
  \opcio{a}{Que tot objecte és isomorf al seu objecte de parts.}
  \opcio{b}{Que els morfismes cap a $\mathcal P(E)$ classifiquen les relacions entre $A$ i $E$.}
  \opcio{c}{Que $\mathcal P(E)$ és l'objecte terminal.}
  \opcio{d}{Que $E$ té exactament dos subobjectes.}
\end{enumerate}
\feedback
\justificacio{Combina la currificació $\Hom(A,\Omega^{E})\cong\Hom(A\times E,\Omega)$ amb la representabilitat de $\Sub$; per a $A=1$ recupera $\Hom(1,\mathcal P(E))\cong\Sub(E)$.}

\pregunta Quina afirmació sobre la categoria $\cat{FinSet}$ dels conjunts finits és correcta?
\begin{enumerate}
  \opcio{a}{No és un topos, perquè no té coproductes infinits.}
  \opcio{b}{És equivalent a $\cat{Set}^{\cat{C}\op}$ per a alguna categoria petita $\cat{C}$.}
  \opcio{c}{És un topos elemental que no és equivalent a cap categoria de prefeixos.}
  \opcio{d}{Només és un topos si s'hi afegeixen els conjunts infinits.}
\end{enumerate}
\feedback
\justificacio{Té límits finits, exponencials finits i el classificador $\{0,1\}$, de manera que és un topos. Tota categoria de prefeixos té colímits petits, i $\cat{FinSet}$ no té coproductes infinits; com que una equivalència conserva els colímits, no és equivalent a cap. Un topos elemental no demana colímits infinits.}

\pregunta A $\cat{Graph}$, $\Omega$ té dos vèrtexs i cinc arestes. Quin fet mostra que $\cat{Graph}$ no és un topos booleà?
\begin{enumerate}
  \opcio{a}{Que $\Omega$ no és un conjunt.}
  \opcio{b}{Que $\Omega$ té més de dos vèrtexs.}
  \opcio{c}{Que $\Hom(\mathbf L,\Omega)$ té exactament dos elements.}
  \opcio{d}{Que $\{t\}$ és el més gran dels subgrafs de $\mathbf E$ disjunts de $\{s\}$, però $\{s\}\cup\{t\}$ no és tot $\mathbf E$.}
\end{enumerate}
\feedback
\justificacio{Per a $U=\{s\}$, el seu \enquote{complement} $\neg U=\{t\}$ no cobreix, juntament amb $U$, tot $\mathbf E$: hi falta l'aresta. Falla, doncs, $U\vee\neg U=\top$. L'opció (b) és falsa ($\Omega$ té dos vèrtexs), i també la (c) (n'hi ha tres).}

\pregunta Un topos té exactament dos morfismes $1\to\Omega$. Què se'n pot concloure?
\begin{enumerate}
  \opcio{a}{Que és booleà, perquè $\Omega$ té exactament dos valors de veritat.}
  \opcio{b}{Res sobre la booleanitat: els morfismes $1\to\Omega$ només compten $\Sub(1)$, mentre que ser booleà és una condició sobre tots els $\Sub(E)$.}
  \opcio{c}{Que $\Omega\cong1+1$.}
  \opcio{d}{Que el topos és equivalent a $\cat{Set}$.}
\end{enumerate}
\feedback
\justificacio{$\Hom(1,\Omega)\cong\Sub(1)$. Els prefeixos sobre el monoide $(\mathbb N,+)$ tenen exactament dos valors globals, però els subobjectes del representable formen una cadena $\mathbb N\cup\{\infty\}$, que no és una àlgebra de Boole, i aquest $\Omega$ no és $1+1$. Tenir dos valors globals i ser booleà són coses diferents.}

\end{enumerate}
\clauestatica
\finalitzaTest

\fi
\ifimprimible\else\end{Form}\fi

% ============================================================
% APÈNDIXS NUMERATS
% \appendix ha d'anar dins \mainmatter (abans de \backmatter):
% a la classe book, \backmatter desactiva la numeració de capítols,
% de manera que els capítols d'apèndix es numeren A, B, ... aquí.
% ============================================================
\appendix
\part{Apèndixs}
\chapter[Axiomàtica relacional]{Una axiomàtica relacional de les categories petites}
\label{cap:axiomatica-relacional}

\begin{objectius}
\apartat{Prerequisits} La definició de categoria, la categoria oposada i la formulació operativa del principi de dualitat (capítol~\ref{cap:categories}, secció~\ref{sec:dualitat}). De lògica de primer ordre amb igualtat: la sintaxi (fórmules, sentències, variables lliures i lligades, substitució), la semàntica (estructures, satisfacció i models), la noció de deducció en un càlcul de Hilbert i les demostracions per inducció sobre fórmules. El teorema de completesa només s'esmenta. Les referències a la lògica lliure de la secció~\ref{sec:ar-relacional-funcional} són informatives i no es pressuposen.
\apartat{Objectius} En acabar l'apèndix, el lector hauria de poder:
\begin{itemize}
  \item presentar una categoria petita en un llenguatge de primer ordre \emph{amb una sola mena d'individus}, on tot individu és una fletxa i $\Dom$, $\Cod$ i $\Com$ són les úniques relacions primitives;
  \item distingir el llenguatge formal, els axiomes i el metallenguatge, i reconèixer quines nocions són primitives i quines en són derivades;
  \item deduir dels axiomes l'estabilitat dels extrems, la unicitat de la composició i la derivabilitat del tercer axioma;
  \item reconstruir els objectes com les fletxes identitat i justificar el pas a la notació funcional usual;
  \item explicar en quin sentit la presentació relacional és equivalent a la definició usual, fins a isomorfisme canònic;
  \item comparar la presentació relacional amb una presentació funcional amb operacions parcials;
  \item definir la traducció dual de fórmules, comprovar l'autodualitat dels axiomes i demostrar el principi de dualitat en les versions sintàctica i semàntica, relacionant-lo amb la versió operativa del capítol~\ref{cap:categories}.
\end{itemize}
\end{objectius}

La definició habitual d'una categoria (definició~\ref{def:categoria}) introdueix objectes, fletxes, identitats i composició. En aquest apèndix adoptem un punt de vista més auster: el llenguatge categòric primitiu només parla de fletxes. Allò que habitualment anomenem objectes es reconstrueix després com una propietat especial d'algunes fletxes. És, en un sentit ben literal, el programa de la introducció ---\emph{relaciona, no individualitza} (secció~\ref{sec:mirar-fora})--- dut fins a la seva conclusió sintàctica: si el que compta d'un objecte és el seu paper dins el sistema de fletxes, aleshores no cal introduir-lo com a dada primitiva.

L'objectiu no és substituir la presentació usual, que és molt més còmoda per al treball matemàtic ordinari. L'objectiu és separar amb precisió tres nivells: el \emph{llenguatge formal} de la teoria; els \emph{axiomes} formulats en aquest llenguatge; i el \emph{metallenguatge}, on interpretem les fórmules i introduïm notació derivada. Aquesta separació permet veure quines nocions són primitives i quines són conseqüència dels axiomes. En particular, la funcionalitat de la composició no s'incorpora a la sintaxi: es demostrarà com un teorema.

\begin{nota}
La idea de descriure una categoria \enquote{només amb fletxes}, amb els objectes identificats amb les identitats, és clàssica: apareix ja en Mac~Lane~\cite{maclane}, es formula amb una relació ternària de composició en la teoria elemental de categories de Lawvere~\cite{lawvere-ccfm}, i s'estudia sistemàticament, amb operacions parcials, en Scott~\cite{scott-identity} i en Freyd i Scedrov~\cite{freydscedrov}. Diversos d'aquests sistemes han estat comparats i verificats formalment per Benzmüller i Scott~\cite{BS2020,BSAFP}. La contribució d'aquest apèndix no és, doncs, la idea, sinó una presentació concreta: una axiomàtica de primer ordre \emph{purament relacional} ---fins i tot el domini i el codomini són relacions--- en què la funcionalitat de la composició i el comportament dels seus extrems esdevenen \emph{teoremes}.
\end{nota}

\section{El llenguatge formal}
\label{sec:ar-llenguatge}

Treballarem en lògica clàssica de primer ordre amb igualtat i amb una sola mena d'individus. La signatura no lògica és
\[
\mathcal L_{\mathrm{Cat}}
=
\{\Dom^{2},\Cod^{2},\Com^{3}\}.
\]
No hi ha constants ni símbols de funció. Per tant, els termes del llenguatge són només variables individuals $f,g,h,d,c,e,\dots$, i totes recorren un mateix univers.

Les fórmules atòmiques no lògiques són de les formes $\Dom(f,d)$, $\Cod(f,c)$ i $\Com(f,g,h)$. La seva lectura metalingüística serà:
\begin{align*}
\Dom(f,d)&:\quad \text{\enquote{el domini de $f$ és $d$}},\\
\Cod(f,c)&:\quad \text{\enquote{el codomini de $f$ és $c$}},\\
\Com(f,g,h)&:\quad \text{\enquote{primer $f$, després $g$, i la composta és $h$}}.
\end{align*}
Així, quan la notació funcional ja estigui justificada, $\Com(f,g,h)$ correspondrà metalingüísticament a $h=g\comp f$.

La paraula \emph{fletxa} pertany, en aquest punt, al metallenguatge: no hi ha cap predicat primitiu $\operatorname{Fletxa}(f)$. Tots els individus del model s'interpreten com a fletxes. Una estructura per a $\mathcal L_{\mathrm{Cat}}$ és, doncs,
\[
\mathfrak M=(M,\Dom^{\mathfrak M},\Cod^{\mathfrak M},\Com^{\mathfrak M}),
\]
on $M$ és l'univers, $\Dom^{\mathfrak M},\Cod^{\mathfrak M}\subseteq M^2$ i $\Com^{\mathfrak M}\subseteq M^3$.

\begin{observacio}[metateòrica]
\label{obs:ar-metateorica}
Amb la semàntica estàndard de primer ordre, $M$ és un conjunt \emph{no buit}; per tant, els models en sentit estricte corresponen a categories petites no buides. La categoria buida queda exclosa per la convenció habitual que l'univers d'una estructura de primer ordre no és buit; si es vol incloure-la, cal permetre explícitament estructures amb univers buit. Les categories grans, d'altra banda, exigeixen una convenció metateòrica addicional ---classes pròpies o universos de Grothendieck, en la línia de l'apèndix~\ref{cap:mida}---, però això no modifica el llenguatge categòric primitiu anterior.
\end{observacio}

\section{Els axiomes}
\label{sec:ar-axiomes}

Presentem primer el sistema en cinc blocs conceptuals. A la secció~\ref{sec:ar-redundancia} veurem que, amb una formulació prou forta de l'associativitat, el tercer bloc és derivable dels altres; la base final del sistema és, doncs, \eqref{ax:A1D}--\eqref{ax:A1C}, \eqref{ax:A2}, \eqref{ax:A4E}--\eqref{ax:A4U} i \eqref{ax:A5D}--\eqref{ax:A5C}.

\paragraph{Axioma 1: existència i unicitat de domini i codomini.}
Tota fletxa té exactament un domini i exactament un codomini:
\begin{align}
\forall f\,\exists!d\,\Dom(f,d),
\tag{A1D}\label{ax:A1D}\\
\forall f\,\exists!c\,\Cod(f,c).
\tag{A1C}\label{ax:A1C}
\end{align}
Com és habitual, $\exists!$ és només una abreviació de la lògica de primer ordre amb igualtat.

\paragraph{Axioma 2: existència de la composició.}
Dues fletxes es poden compondre exactament quan el codomini de la primera coincideix amb el domini de la segona:
\begin{equation}
\forall f\,\forall g\,
\left[
\exists e\bigl(\Cod(f,e)\land\Dom(g,e)\bigr)
\leftrightarrow
\exists h\,\Com(f,g,h)
\right].
\tag{A2}\label{ax:A2}
\end{equation}
Aquest axioma afirma existència, però no unicitat. La unicitat de la composta serà un teorema (secció~\ref{sec:ar-unicitat}).

\paragraph{Axioma 3: extrems de la composta.}
Si $h$ és una composta de $f$ amb $g$, el domini de $h$ és el de $f$ i el codomini de $h$ és el de $g$:
\begin{equation}
\begin{split}
\forall f,g,h,d,c\,\bigl[
&\Com(f,g,h)\land\Dom(f,d)\land\Cod(g,c)\\
&\longrightarrow
\Dom(h,d)\land\Cod(h,c)
\bigr].
\end{split}
\tag{A3}\label{ax:A3}
\end{equation}

\paragraph{Axioma 4: associativitat forta.}
Com que $\Com$ és una relació i la composició és parcial, la mera igualtat de dues composicions iterades ja existents no expressa tota la força de l'associativitat usual. Adoptarem una formulació en dues parts. Primer, si existeixen les dues composicions consecutives $g\comp f$ i $h\comp g$, existeixen també les dues composicions iterades:
\begin{equation}
\begin{split}
\forall f,g,h,u,v\,\bigl[
&\Com(f,g,u)\land\Com(g,h,v)\\
&\longrightarrow
\exists p\,\exists q\,
\bigl(\Com(u,h,p)\land\Com(f,v,q)\bigr)
\bigr].
\end{split}
\tag{A4E}\label{ax:A4E}
\end{equation}
Segon, qualssevol resultats obtinguts per les dues associacions coincideixen:
\begin{equation}
\begin{split}
\forall f,g,h,u,v,p,q\,\bigl[
&\Com(f,g,u)\land\Com(g,h,v)\land\\
&\Com(u,h,p)\land\Com(f,v,q)
\longrightarrow p=q
\bigr].
\end{split}
\tag{A4U}\label{ax:A4U}
\end{equation}
Els dos enunciats, \eqref{ax:A4E} i \eqref{ax:A4U}, formen el bloc d'associativitat (A4). El primer controla l'existència de les composicions iterades; el segon n'expressa la coherència. Un cop demostrada la funcionalitat de $\Com$ (secció~\ref{sec:ar-unicitat}), aquest bloc es llegeix com la llei usual d'un semigrup parcial: si $g\comp f$ i $h\comp g$ existeixen, aleshores existeixen $h\comp(g\comp f)$ i $(h\comp g)\comp f$, i són iguals.

\paragraph{Axioma 5: neutralitat.}
El domini és neutre per la dreta i el codomini és neutre per l'esquerra:
\begin{align}
\forall f\,\forall d\,
\bigl(\Dom(f,d)\rightarrow\Com(d,f,f)\bigr),
\tag{A5D}\label{ax:A5D}\\
\forall f\,\forall c\,
\bigl(\Cod(f,c)\rightarrow\Com(f,c,f)\bigr).
\tag{A5C}\label{ax:A5C}
\end{align}
Metalingüísticament, un cop recuperada la notació funcional, aquestes fórmules diran $f\comp d=f$ i $c\comp f=f$.

\section{Estabilitat dels extrems}
\label{sec:ar-estabilitat}

El primer fet rellevant és que qualsevol fletxa que apareix com a domini o codomini és estable sota les dues relacions d'extrem.

\begin{proposicio}[estabilitat dels extrems]
\label{prop:ar-estabilitat}
Si $\Dom(f,d)$, aleshores $\Dom(d,d)\land\Cod(d,d)$. Si $\Cod(f,c)$, aleshores $\Dom(c,c)\land\Cod(c,c)$.
\end{proposicio}

\begin{prova}
Suposem $\Dom(f,d)$. Per \eqref{ax:A5D}, $\Com(d,f,f)$. Per la direcció de dreta a esquerra de \eqref{ax:A2}, l'existència d'aquesta composició implica que existeix $e$ tal que $\Cod(d,e)\land\Dom(f,e)$. Com que també $\Dom(f,d)$, la unicitat del domini de $f$, donada per \eqref{ax:A1D}, implica $e=d$; per tant $\Cod(d,d)$.

Ara \eqref{ax:A5C}, aplicada a $d$, dona $\Com(d,d,d)$. Novament per \eqref{ax:A2}, existeix $e$ tal que $\Cod(d,e)\land\Dom(d,e)$. Ja sabem que $\Cod(d,d)$; per unicitat del codomini de $d$, $e=d$, i per tant $\Dom(d,d)$. Així, $\Dom(f,d)\longrightarrow \Dom(d,d)\land\Cod(d,d)$.

El cas del codomini és dual. Si $\Cod(f,c)$, llavors \eqref{ax:A5C} dona $\Com(f,c,f)$. Per \eqref{ax:A2} existeix $e$ amb $\Cod(f,e)\land\Dom(c,e)$; la unicitat del codomini de $f$ força $e=c$, de manera que $\Dom(c,c)$. Aleshores \eqref{ax:A5D} dona $\Com(c,c,c)$, i un nou ús de \eqref{ax:A2} i de la unicitat del domini de $c$ proporciona $\Cod(c,c)$.
\end{prova}

\begin{observacio}
La demostració només utilitza (A1), (A2) i (A5). No utilitza ni l'axioma dels extrems de la composta (A3) ni l'associativitat (A4).
\end{observacio}

\section{Unicitat de la composició}
\label{sec:ar-unicitat}

La relació $\Com$ no s'ha postulat com a funcional en el tercer argument. Aquesta funcionalitat es dedueix de la componibilitat, la neutralitat i l'associativitat.

\begin{proposicio}[funcionalitat de la composició]
\label{prop:ar-unicitat}
Per a totes les fletxes $f,g,h,k$,
\[
\Com(f,g,h)\land\Com(f,g,k)\longrightarrow h=k.
\]
\end{proposicio}

\begin{prova}
Suposem $\Com(f,g,h)$ i $\Com(f,g,k)$. Per \eqref{ax:A2}, existeix $e$ tal que $\Cod(f,e)\land\Dom(g,e)$. Per neutralitat, $\Com(f,e,f)$ i $\Com(e,g,g)$. Apliquem ara \eqref{ax:A4U} a la terna $f,e,g$. Prenent $u=f$, $v=g$, $p=h$, $q=k$, les quatre premisses de \eqref{ax:A4U} són precisament $\Com(f,e,f)$, $\Com(e,g,g)$, $\Com(f,g,h)$ i $\Com(f,g,k)$. Per tant, $h=k$. La demostració només fa servir (A2), (A5) i \eqref{ax:A4U}; en particular, no fa servir (A3).
\end{prova}

Així, $\Com$ és el graf d'una operació binària parcial: per a cada parella $(f,g)$ hi ha com a màxim una fletxa $h$ tal que $\Com(f,g,h)$, i (A2) determina exactament quan aquesta fletxa existeix.

\section{Identitats i objectes com a nocions derivades}
\label{sec:ar-identitats}

Definim, com a abreviació del llenguatge,
\begin{equation}
\Id(e)
\;:\Longleftrightarrow\;
\Dom(e,e)\land\Cod(e,e).
\label{def:ar-Id}
\end{equation}
El símbol $\Id$ no s'afegeix a la signatura primitiva: és només una abreviació definicional. La proposició d'estabilitat dona immediatament $\Dom(f,d)\rightarrow\Id(d)$ i $\Cod(f,c)\rightarrow\Id(c)$: tot domini i tot codomini és una fletxa identitària.

També podem introduir $\Obj(e):\Longleftrightarrow\Id(e)$. Així, un \emph{objecte} no és una segona mena d'individu: és una fletxa que té per domini i codomini ella mateixa. En particular, sota els axiomes anteriors,
\[
\Dom(e,e)
\;\Longleftrightarrow\;
\Cod(e,e)
\;\Longleftrightarrow\;
\Id(e),
\]
perquè qualsevol de les dues primeres propietats, aplicada a la proposició~\ref{prop:ar-estabilitat}, implica l'altra.

\section{El tercer axioma és derivable}
\label{sec:ar-redundancia}

La forma forta de l'associativitat permet demostrar l'axioma (A3) a partir de (A1), (A2), (A4) i (A5).

\begin{proposicio}[redundància de (A3)]
\label{prop:ar-redundancia}
Suposem (A1), (A2), (A4) i (A5). Si $\Com(f,g,h)$, $\Dom(f,d)$ i $\Cod(g,c)$, aleshores $\Dom(h,d)$ i $\Cod(h,c)$.
\end{proposicio}

\begin{prova}
Les proposicions~\ref{prop:ar-estabilitat} i~\ref{prop:ar-unicitat}, que farem servir, no depenen de (A3); l'argument no és, doncs, circular. Comencem pel domini. De $\Dom(f,d)$ i \eqref{ax:A5D} obtenim $\Com(d,f,f)$. Juntament amb $\Com(f,g,h)$, l'axioma d'existència associativa \eqref{ax:A4E}, aplicat a la terna $d,f,g$, proporciona $p,q$ tals que $\Com(f,g,p)$ i $\Com(d,h,q)$. Per \eqref{ax:A4U}, aquestes dues composicions iterades tenen el mateix resultat, és a dir, $p=q$. D'altra banda, per la unicitat de la composició (proposició~\ref{prop:ar-unicitat}), $\Com(f,g,h)$ i $\Com(f,g,p)$ impliquen $p=h$. Per tant $q=h$, i tenim $\Com(d,h,h)$. Per \eqref{ax:A2}, existeix $e$ amb $\Cod(d,e)\land\Dom(h,e)$. Com que $d$ és domini de $f$, la proposició~\ref{prop:ar-estabilitat} dona $\Cod(d,d)$. La unicitat del codomini de $d$ implica $e=d$, i obtenim $\Dom(h,d)$.

Per al codomini, de $\Cod(g,c)$ i \eqref{ax:A5C} obtenim $\Com(g,c,g)$. Aplicada a $f,g,c$, l'associativitat forta proporciona composicions iterades $\Com(h,c,p)$ i $\Com(f,g,q)$ amb $p=q$. La unicitat de la composició de $f$ amb $g$ dona $q=h$, i per tant $\Com(h,c,h)$. Per \eqref{ax:A2}, existeix $e$ amb $\Cod(h,e)\land\Dom(c,e)$. Com que $c$ és codomini de $g$, l'estabilitat dona $\Dom(c,c)$. Per unicitat del domini de $c$, $e=c$, i així $\Cod(h,c)$.
\end{prova}

\begin{corollari}
\label{cor:ar-equivalencia-sistemes}
Amb l'associativitat forta adoptada aquí, el sistema de cinc blocs $(A1),(A2),(A3),(A4),(A5)$ és equivalent al sistema reduït $(A1),(A2),(A4),(A5)$. El tercer axioma original es pot conservar en una primera presentació per la seva transparència conceptual, però no és necessari com a axioma independent.
\end{corollari}

No afirmem que els quatre blocs restants siguin independents entre si: la redundància de (A3) no estableix la minimalitat del sistema reduït. L'anàlisi d'independència de sistemes funcionals afins ha estat estudiada i formalitzada per Benzmüller i Scott~\cite{BS2020,BSAFP}, però la transferència literal d'aquells resultats al llenguatge purament relacional adoptat aquí requeriria una comparació formal addicional; una via natural seria construir models finits que separessin cadascun dels quatre blocs.

\begin{observacio}[la forma lògica de (A4) és essencial]
Aquesta reducció depèn de la forma lògica de (A4). Si l'associativitat es formula només com la igualtat de dues composicions iterades \emph{suposant que totes dues ja existeixen}, no es controla l'existència de les composicions iterades i la demostració anterior deixa de funcionar. En una teoria de composició parcial, la forma lògica de l'associativitat és, doncs, part essencial de l'axiomàtica.
\end{observacio}

\section{Recuperació de la notació funcional}
\label{sec:ar-notacio-funcional}

L'axioma (A1) permet introduir en el metallenguatge les notacions $d(f)$ i $c(f)$ per a les úniques fletxes que satisfan, respectivament, $\Dom(f,d(f))$ i $\Cod(f,c(f))$. Aquestes expressions no són termes del llenguatge primitiu $\mathcal L_{\mathrm{Cat}}$: són notació derivada, legitimada per un teorema d'existència i unicitat; en termes lògics, passem a una extensió definicional per descripcions úniques. De la mateixa manera, si $f$ i $g$ són componibles, (A2) dona l'existència d'una $h$ tal que $\Com(f,g,h)$, i la proposició~\ref{prop:ar-unicitat} en dona la unicitat; escrivim, doncs, $g\comp f$ per designar aquesta única fletxa.

Amb aquesta notació, (A2) es llegeix
\[
g\comp f\ \text{existeix}
\quad\Longleftrightarrow\quad
c(f)=d(g),
\]
i la proposició~\ref{prop:ar-redundancia} dona $d(g\comp f)=d(f)$ i $c(g\comp f)=c(g)$.

Si $A$ i $B$ són fletxes identitàries, la notació usual $f\colon A\longrightarrow B$ abrevia metalingüísticament $\Dom(f,A)\land\Cod(f,B)$; les condicions $\Id(A)$ i $\Id(B)$ són automàtiques per estabilitat. També podem escriure, metalingüísticament, $\Hom(A,B)=\{f\in M: \Dom(f,A)\land\Cod(f,B)\}$. El símbol $\Hom$ no forma part tampoc del llenguatge primitiu.

\section{Equivalència amb la definició usual}
\label{sec:ar-equivalencia}

La presentació relacional no defineix una estructura matemàtica nova: recupera la noció usual de categoria petita (definició~\ref{def:categoria}), fins a l'isomorfisme canònic que identifica cada objecte amb la seva identitat.

Abans de les demostracions, val la pena desplegar un model concret i petit, per veure com les relacions $\Dom$, $\Cod$ i $\Com$ codifiquen una categoria familiar i com els objectes hi apareixen com a fletxes identitat.

\begin{exemple}[Un model amb tres objectes i una cadena de morfismes]\label{ex:ar-model-3}
Considerem la categoria $\cat{3}$ (amb els objectes i les fletxes reanomenats), formada per una cadena de dues fletxes, $A\xrightarrow{u}B\xrightarrow{v}C$, amb la seva composta $w=v\comp u\colon A\to C$:
\[
\begin{tikzpicture}[baseline=(A.base)]
  \node (A) at (0,0) {$A$};
  \node (B) at (2.2,0) {$B$};
  \node (C) at (4.4,0) {$C$};
  \draw[-{Stealth[scale=1.1]}] (A) -- node[above] {$u$} (B);
  \draw[-{Stealth[scale=1.1]}] (B) -- node[above] {$v$} (C);
  \draw[-{Stealth[scale=1.1]}] (A) to[out=-35,in=-145] node[below] {$w=v\comp u$} (C);
\end{tikzpicture}
\]
El model relacional corresponent té per univers el conjunt de \emph{totes} les fletxes, incloses les tres identitats:
\[
M=\{\,A,\ B,\ C,\ u,\ v,\ w\,\},
\]
on $A,B,C$ són alhora els objectes i les seves fletxes identitat. Les relacions unàries derivades i els extrems de cada fletxa són:
\[
\renewcommand{\arraystretch}{1.2}
\begin{array}{c|c|c|c}
f & \Dom(f,-) & \Cod(f,-) & \Id(f) \\ \hline
A & A & A & \text{sí} \\
B & B & B & \text{sí} \\
C & C & C & \text{sí} \\
u & A & B & \text{no} \\
v & B & C & \text{no} \\
w & A & C & \text{no}
\end{array}
\]
Els objectes són exactament les fletxes $e$ amb $\Id(e)$, és a dir $A,B,C$: no calen com a dada primitiva, sinó que es reconeixen per la propietat $\Dom(e,e)\land\Cod(e,e)$. Les úniques parelles componibles $(f,g)$ ---les que compleixen $c(f)=d(g)$--- donen la taula de $\Com$, on recordem que $\Com(f,g,h)$ es llegeix \enquote{primer $f$, després $g$, amb composta $h=g\comp f$}:
\[
\renewcommand{\arraystretch}{1.2}
\begin{array}{ll}
\Com(A,A,A) & \text{(identitat de }A) \\
\Com(B,B,B) & \text{(identitat de }B) \\
\Com(C,C,C) & \text{(identitat de }C) \\
\Com(A,u,u),\ \Com(u,B,u) & (u\comp\id_A=u=\id_B\comp u) \\
\Com(B,v,v),\ \Com(v,C,v) & (v\comp\id_B=v=\id_C\comp v) \\
\Com(A,w,w),\ \Com(w,C,w) & (w\comp\id_A=w=\id_C\comp w) \\
\Com(u,v,w) & (\text{la composició no trivial }v\comp u=w)
\end{array}
\]
Val la pena llegir l'última línia amb l'ordre de $\Com$ ben present: $\Com(u,v,w)$ diu que en recórrer \emph{primer} $u$ i \emph{després} $v$ s'obté $w$, és a dir $w=v\comp u$. No hi ha cap altra parella componible; en particular, $u$ i $w$ no es componen en cap ordre, perquè $c(u)=B\neq A=d(w)$ i $c(w)=C\neq A=d(u)$. Comprovem els axiomes. (A1) és la taula d'extrems. (A2) es comprova parella per parella: les deu ternes de la taula de $\Com$ corresponen exactament a les deu parelles $(f,g)$ amb $c(f)=d(g)$. (A5) és el contingut de les sis primeres línies, que donen $\Com(d(f),f,f)$ i $\Com(f,c(f),f)$ per a cada $f$. Per a (A4), observem que una terna encadenada $(f,g,h)$, amb $\Com(f,g,u)$ i $\Com(g,h,v)$, conté sempre alguna identitat, perquè la cadena més llarga de fletxes no identitat és $u,v$, de longitud dos. Si la identitat és $f$, aleshores $u=g$ i $f$ és el domini de $v$, de manera que totes dues associacions donen $v$: $\Com(u,h,v)$ i $\Com(f,v,v)$. Els casos en què la identitat és $g$ o $h$ són anàlegs, i totes dues associacions donen la composta de les altres dues fletxes. Així es compleixen \eqref{ax:A4E} i \eqref{ax:A4U}. Reconstruint la categoria a partir d'aquest model (com farà el teorema~\ref{th:ar-model-a-categoria}) es recupera la categoria de partida, en què ara cada objecte és la seva pròpia fletxa identitat.
\end{exemple}

\begin{teorema}
\label{th:ar-model-a-categoria}
Tot model del sistema reduït $(A1),(A2),(A4),(A5)$ amb univers no buit determina una categoria petita no buida en el sentit de la definició~\ref{def:categoria}.
\end{teorema}

\begin{prova}
Sigui $\mathfrak M$ un model. Prenem com a fletxes els elements de $M$ i com a objectes $\Obj(\mathfrak M)=\{e\in M:\Id(e)\}$. Per (A1), cada fletxa $f$ té un domini únic $d(f)$ i un codomini únic $c(f)$; per estabilitat, tots dos són objectes. Si $c(f)=d(g)$, (A2) proporciona una composta i la proposició~\ref{prop:ar-unicitat} en proporciona una de sola, que denotem $g\comp f$; la proposició~\ref{prop:ar-redundancia} assegura que $d(g\comp f)=d(f)$ i $c(g\comp f)=c(g)$. Per a un objecte $A$, la seva fletxa identitat és $A$ mateix. Aquesta fletxa va de $A$ a $A$, perquè $\Dom(A,A)$ i $\Cod(A,A)$. Si $f\colon A\to B$, (A5) dona $\Com(A,f,f)$ i $\Com(f,B,f)$, és a dir, $f\comp\id_A=f=\id_B\comp f$. Per a l'associativitat, siguin $f\colon A\to B$, $g\colon B\to C$ i $h\colon C\to D$. Aplicant \eqref{ax:A4E} a $\Com(f,g,g\comp f)$ i $\Com(g,h,h\comp g)$, existeixen $p,q$ amb $\Com(g\comp f,h,p)$ i $\Com(f,h\comp g,q)$; per la unicitat de la composició, $p=h\comp(g\comp f)$ i $q=(h\comp g)\comp f$, i \eqref{ax:A4U} dona $p=q$. Així obtenim una categoria, petita perquè $M$ és un conjunt.
\end{prova}

\begin{teorema}[recíproc]
\label{th:ar-categoria-a-model}
Tota categoria petita no buida determina un model del sistema relacional.
\end{teorema}

\begin{prova}
Sigui $\cat{C}$ una categoria petita no buida. Prenem com a univers $M=\operatorname{Mor}(\cat{C})$. Si $f\colon A\to B$, definim
\[
\Dom(f,d)\Longleftrightarrow d=\id_A,
\qquad
\Cod(f,c)\Longleftrightarrow c=\id_B,
\]
i
i, per a $f\colon A\to B$ i $g\colon B'\to C$,
\[
\Com(f,g,h)\Longleftrightarrow \bigl(B=B'\ \text{i}\ h=g\comp f\bigr),
\]
on la notació de la dreta pertany ara al metallenguatge de la categoria usual. L'univers és no buit, perquè $\cat{C}$ té algun objecte i, per tant, alguna identitat. Comprovem els axiomes. (A1) és immediat. Per a (A2), $\Cod(f,e)\land\Dom(g,e)$ vol dir $e=\id_B=\id_{B'}$, i això passa si i només si $B=B'$, perquè $\id_B$ determina $B$ com a domini; i $B=B'$ és exactament la condició perquè existeixi $h$ amb $\Com(f,g,h)$. (A3) diu que $g\comp f\colon A\to C$. Per a (A4), $\Com(f,g,u)$ i $\Com(g,h,v)$ volen dir que $f,g,h$ són consecutives, amb $u=g\comp f$ i $v=h\comp g$; aleshores existeixen $h\comp u$ i $v\comp f$, que són iguals per l'associativitat de $\cat{C}$. Finalment, (A5) són les lleis d'identitat $f\comp\id_A=f$ i $\id_B\comp f=f$.
\end{prova}

\begin{observacio}
En passar d'una categoria usual $\cat{C}$ al model relacional, els individus són els morfismes de $\cat{C}$; en reconstruir després una categoria a partir del model, els seus objectes són les fletxes identitat $\id_A$, no els objectes originals $A$. La categoria reconstruïda és canònicament isomorfa a $\cat{C}$ ---la bijecció d'objectes és $A\mapsto\id_A$--- però en general no n'és literalment la mateixa estructura conjuntista. Per aquest motiu diem que les dues presentacions descriuen la mateixa estructura categòrica \emph{fins a isomorfisme canònic}, no que en siguin idèntiques. Una afirmació metateòrica més forta ---una equivalència definicional o una bi-interpretabilitat--- requeriria definir explícitament les traduccions entre les dues teories i comparar-les; no la formulem aquí.

\end{observacio}

Així, la diferència entre les dues presentacions és una diferència de \emph{llenguatge primitiu}, no de contingut categòric.

\section{Relacional i funcional}
\label{sec:ar-relacional-funcional}

Els predicats $\Dom$ i $\Cod$, sota (A1), són els grafs de dues funcions unàries totals. La relació $\Com$, sota (A2) i la proposició~\ref{prop:ar-unicitat}, és el graf d'una funció binària parcial. Aquesta observació explica dues maneres naturals de formalitzar les categories.

\paragraph{Presentació relacional.}
En lògica de primer ordre estàndard podem prendre com a primitives $\Dom(f,d)$, $\Cod(f,c)$ i $\Com(f,g,h)$. L'existència s'expressa mitjançant quantificadors i la funcionalitat s'axiomatitza o es demostra. És la via d'aquest apèndix.

\paragraph{Presentació funcional.}
Podem prendre $\dom$ i $\cod$ com a operadors unaris. El problema és la composició: com que no tota parella de fletxes és componible, una operació binària de composició no és total sobre $M\times M$. En lògica de primer ordre estàndard, en canvi, els símbols de funció s'interpreten com a funcions totals. Cal, doncs, ampliar el marc lògic o codificar explícitament la parcialitat.

Aquest és el camí seguit, en un formalisme diferent, per Benzmüller i Scott~\cite{BS2020,BSAFP}. Treballen amb una lògica lliure embeguda en Isabelle/HOL, un predicat $E$ d'existència o definició, i tres operacions parcials $\dom$, $\cod$ i composició. En el seu \emph{AxiomsSet5}, atribuït a Scott, els axiomes principals tenen la forma
\begin{align*}
E(\dom x)&\to E(x),\tag{S1}\\
E(\cod y)&\to E(y),\tag{S2}\\
E(x\cdot y)&\leftrightarrow \dom x\approx\cod y,\tag{S3}\\
x\cdot(y\cdot z)&\simeq (x\cdot y)\cdot z,\tag{S4}\\
(\cod y)\cdot y&\simeq y,\tag{S5}\\
x\cdot(\dom x)&\simeq x,\tag{S6}
\end{align*}
amb la composició escrita en ordre funcional ($x\cdot y$ és \enquote{primer $y$, després $x$}), on $\simeq$ és la \emph{igualtat de Kleene} (dos termes són iguals si tots dos són indefinits, o bé tots dos definits i iguals) i $\approx$ és la \emph{identitat existent} (tots dos termes definits i iguals). En aquesta presentació, la funcionalitat de $\dom$, $\cod$ i de la composició queda incorporada en els símbols funcionals; la parcialitat es gestiona mitjançant la lògica lliure. Els mateixos autors formalitzen també equivalències entre diversos sistemes axiomàtics i el pas, per Skolemització, de formulacions existencials cap a formulacions funcionals.

\begin{nota}
Els símbols tipogràfics per a aquestes dues igualtats varien entre fonts; aquí fixem $\simeq$ per a la igualtat de Kleene i $\approx$ per a la identitat existent, que no coincideixen necessàriament amb la notació de l'obra citada.
\end{nota}

La correspondència conceptual amb la presentació relacional és:
\[
\begin{array}{c|c}
\text{presentació relacional} & \text{presentació funcional lliure}\\
\hline
\text{(A1): domini i codomini} & \dom,\cod\ \text{i } S1,S2\\
\text{(A2): componibilitat} & S3\\
\text{(A3): extrems de la composta} & \text{derivable en sistemes reduïts}\\
\text{(A4): associativitat} & S4\\
\text{(A5): neutralitat} & S5,S6
\end{array}
\]
La correspondència no és literal fórmula per fórmula, perquè els dos llenguatges distribueixen la informació de manera diferent.

\section{El principi de dualitat}
\label{sec:ar-dualitat}\index{dualitat!principi de}

El capítol~\ref{cap:categories} formula el principi de dualitat de manera operativa (teorema~\ref{teo:dualitat}): una afirmació expressable en el llenguatge de les categories es dualitza invertint les fletxes, i, si és vàlida en totes les categories, la seva dual també ho és, perquè l'afirmació original es pot aplicar a la categoria oposada. La presentació relacional permet donar-ne una versió completament formal, en dues formes: una de \emph{sintàctica}, sobre deduccions, i una de \emph{semàntica}, sobre models. En aquesta secció escrivim $\mathrm{CT}$ per al conjunt d'axiomes \eqref{ax:A1D}--\eqref{ax:A5C}.

\paragraph{La traducció dual.}
Per a cada fórmula $\varphi$ de $\mathcal L_{\mathrm{Cat}}$ definim la seva \textbf{dual}\index{dualitat!d'un enunciat} $\varphi\op$ substituint cada fórmula atòmica no lògica segons
\begin{gather*}
\Dom(f,d)\longmapsto\Cod(f,d),\qquad
\Cod(f,c)\longmapsto\Dom(f,c),\\
\Com(f,g,h)\longmapsto\Com(g,f,h),
\end{gather*}
i deixant intactes les igualtats, els connectors i els quantificadors. Com que $\Com(f,g,h)$ es llegeix \enquote{primer $f$, després $g$, amb composta $h$}, intercanviar-ne els dos primers arguments és exactament invertir l'ordre de composició. La traducció és una involució, $(\varphi\op)\op=\varphi$, i commuta amb la substitució de variables. Un cop recuperada la notació funcional (secció~\ref{sec:ar-notacio-funcional}), aquestes substitucions són les del capítol~\ref{cap:categories}: $d(f)$ i $c(f)$ s'intercanvien i $g\comp f$ passa a $f\comp g$. Observem que la dualitat no depèn de tractar la composició com un símbol funcional total: opera directament sobre la relació ternària.

\begin{observacio}
Els connectors i quantificadors de $\varphi$ pertanyen al llenguatge $\mathcal L_{\mathrm{Cat}}$, amb què \emph{parlem} d'una categoria, i la dualitat els deixa intactes. No s'han de confondre amb la lògica que una categoria pot codificar \emph{internament}, quan els seus objectes són fórmules i els seus morfismes deduccions. En aquest cas, la dualitat de fletxes pot intercanviar la conjunció i la disjunció del sistema codificat, però el $\land$ i el $\lor$ de $\mathcal L_{\mathrm{Cat}}$ es mantenen.
\end{observacio}

\paragraph{Autodualitat dels axiomes.}

\begin{lema}\label{lema:ar-autodual}
Per a cada axioma $\alpha$ de $\mathrm{CT}$, la fórmula $\alpha\op$ és lògicament equivalent a un axioma de $\mathrm{CT}$. Concretament, \eqref{ax:A1D} i \eqref{ax:A1C} es transformen l'un en l'altre, igual que \eqref{ax:A5D} i \eqref{ax:A5C}; i cadascun dels axiomes \eqref{ax:A2}, \eqref{ax:A3}, \eqref{ax:A4E} i \eqref{ax:A4U} es transforma en una variant d'ell mateix, que només en difereix pel nom de les variables lligades, per l'ordre dels quantificadors i dels membres de les conjuncions i, en el cas de \eqref{ax:A4U}, per la simetria de la igualtat.
\end{lema}

\begin{prova}
Els casos d'\eqref{ax:A1D} i \eqref{ax:A5D} són immediats: $(\forall f\,\exists!d\,\Dom(f,d))\op$ és $\forall f\,\exists!d\,\Cod(f,d)$, que és \eqref{ax:A1C}, i $(\forall f\,\forall d\,(\Dom(f,d)\to\Com(d,f,f)))\op$ és $\forall f\,\forall d\,(\Cod(f,d)\to\Com(f,d,f))$, que és \eqref{ax:A5C}; per involució, els casos d'\eqref{ax:A1C} i \eqref{ax:A5C} són els recíprocs.

La dual d'\eqref{ax:A2} és
\[
\forall f\,\forall g\,\bigl[\exists e\bigl(\Dom(f,e)\land\Cod(g,e)\bigr)\leftrightarrow\exists h\,\Com(g,f,h)\bigr],
\]
i, intercanviant els noms de les variables lligades $f$ i $g$, s'obté \eqref{ax:A2} llevat de l'ordre dels quantificadors i dels membres de la conjunció. La dual d'\eqref{ax:A3} és
\[
\forall f,g,h,d,c\,\bigl[\Com(g,f,h)\land\Cod(f,d)\land\Dom(g,c)\longrightarrow\Cod(h,d)\land\Dom(h,c)\bigr],
\]
i l'intercanvi $f\leftrightarrow g$, $d\leftrightarrow c$ la torna en \eqref{ax:A3}. La dual d'\eqref{ax:A4E} és
\[
\forall f,g,h,u,v\,\bigl[\Com(g,f,u)\land\Com(h,g,v)\longrightarrow\exists p\,\exists q\,\bigl(\Com(h,u,p)\land\Com(v,f,q)\bigr)\bigr];
\]
amb l'intercanvi $f\leftrightarrow h$, $u\leftrightarrow v$, $p\leftrightarrow q$ esdevé
\[
\forall f,g,h,u,v\,\bigl[\Com(g,h,v)\land\Com(f,g,u)\longrightarrow\exists q\,\exists p\,\bigl(\Com(f,v,q)\land\Com(u,h,p)\bigr)\bigr],
\]
que és \eqref{ax:A4E}. El mateix reanomenament transforma la dual d'\eqref{ax:A4U} en \eqref{ax:A4U}, amb la conclusió escrita $q=p$.
\end{prova}

Intuïtivament, el lema diu que la definició de categoria no distingeix entre una fletxa i la mateixa fletxa recorreguda en sentit contrari. És l'únic ingredient específic de la teoria que necessiten els dos teoremes següents.

\paragraph{Dualitat sintàctica.}

\begin{teorema}[Dualitat sintàctica]\label{teo:ar-dualitat-sint}
Per a tota sentència $\varphi$ de $\mathcal L_{\mathrm{Cat}}$, si $\mathrm{CT}\vdash\varphi$, aleshores $\mathrm{CT}\vdash\varphi\op$.
\end{teorema}

\begin{prova}
Fixem un càlcul de Hilbert estàndard per a la lògica de primer ordre amb igualtat. La traducció dual només reanomena les relacions primitives i permuta els arguments de $\Com$; commuta amb la substitució, conserva les variables lliures i no toca ni els connectors, ni els quantificadors, ni la igualtat. Per tant, transforma cada axioma lògic (inclosos els de la igualtat) en un axioma lògic, i cada aplicació d'una regla d'inferència en una aplicació de la mateixa regla; les restriccions sobre variables de la regla de generalització es conserven, perquè les variables lliures no canvien. Aplicada fórmula a fórmula a una deducció de $\varphi$ a partir de $\mathrm{CT}$, en dona una de $\varphi\op$ a partir de $\mathrm{CT}\op=\{\alpha\op:\alpha\in\mathrm{CT}\}$. Pel lema~\ref{lema:ar-autodual}, cada $\alpha\op$ es dedueix de $\mathrm{CT}$; per tant, $\mathrm{CT}\vdash\varphi\op$.
\end{prova}

\paragraph{Dualitat semàntica.}
Per a una estructura $\mathfrak M=(M,\Dom^{\mathfrak M},\Cod^{\mathfrak M},\Com^{\mathfrak M})$ definim l'\textbf{estructura oposada}
\[
\mathfrak M\op=\bigl(M,\ \Cod^{\mathfrak M},\ \Dom^{\mathfrak M},\ \{(g,f,h):(f,g,h)\in\Com^{\mathfrak M}\}\bigr),
\]
amb el mateix univers, els papers de domini i codomini intercanviats i la composició en l'ordre invers. Si $\mathfrak M$ és el model associat a una categoria petita $\cat{C}$ pel teorema~\ref{th:ar-categoria-a-model}, aleshores l'aplicació $f\mapsto f\op$ és un isomorfisme de $\mathfrak M\op$ en el model associat a $\cat{C}\op$: per a $f\colon A\to B$, la relació $\Dom^{\mathfrak M\op}(f,\id_B)$, que és $\Cod^{\mathfrak M}(f,\id_B)$, correspon a $\Dom(f\op,(\id_B)\op)$ a $\cat{C}\op$, i $\Com^{\mathfrak M\op}(f,g,h)$, que és $\Com^{\mathfrak M}(g,f,h)$, és a dir $h=f\comp g$, correspon a $h\op=g\op\comp f\op$, és a dir, a $\Com(f\op,g\op,h\op)$ a $\cat{C}\op$. Si s'identifica cada $f\op$ amb $f$, els dos models coincideixen.

\begin{lema}\label{lema:ar-oposada}
Per a tota fórmula $\varphi$ de $\mathcal L_{\mathrm{Cat}}$ i tota assignació $s$ de les seves variables lliures,
\[
\mathfrak M\op\models\varphi[s]
\quad\Longleftrightarrow\quad
\mathfrak M\models\varphi\op[s].
\]
\end{lema}

\begin{prova}
Per inducció sobre $\varphi$. Per a les fórmules atòmiques és la definició de $\mathfrak M\op$: per exemple, $\mathfrak M\op\models\Com(f,g,h)[s]$ si i només si $(s(g),s(f),s(h))\in\Com^{\mathfrak M}$, és a dir, si i només si $\mathfrak M\models\Com(g,f,h)[s]$. Les igualtats s'interpreten igual a totes dues estructures, i els passos dels connectors i dels quantificadors són immediats perquè l'univers és el mateix.
\end{prova}

\begin{teorema}[Dualitat semàntica]\label{teo:ar-dualitat-sem}
Si $\mathfrak M$ és un model de $\mathrm{CT}$, també ho és $\mathfrak M\op$. En conseqüència, per a tota sentència $\varphi$ de $\mathcal L_{\mathrm{Cat}}$, si $\varphi$ es compleix en tots els models de $\mathrm{CT}$ ---que, pels teoremes~\ref{th:ar-model-a-categoria} i~\ref{th:ar-categoria-a-model}, corresponen a les categories petites no buides---, també s'hi compleix $\varphi\op$.
\end{teorema}

\begin{prova}
Si $\mathfrak M\models\mathrm{CT}$ i $\alpha\in\mathrm{CT}$, pel lema~\ref{lema:ar-autodual} la fórmula $\alpha\op$ és equivalent a un axioma, de manera que $\mathfrak M\models\alpha\op$, i pel lema~\ref{lema:ar-oposada}, $\mathfrak M\op\models\alpha$. Per a la segona afirmació, sigui $\mathfrak M$ un model de $\mathrm{CT}$. Com que $\mathfrak M\op$ també ho és, $\mathfrak M\op\models\varphi$, i el lema~\ref{lema:ar-oposada} dona $\mathfrak M\models\varphi\op$.
\end{prova}

Aquesta demostració és la versió formal de la del capítol~\ref{cap:categories}: allà es fa servir que $\cat{C}\op$ és una categoria; aquí, que $\mathfrak M\op$ és un model. Pel teorema de completesa, $\mathrm{CT}\vdash\varphi$ equival a $\mathrm{CT}\models\varphi$, i els dos teoremes són, doncs, dues cares d'un mateix fet. Les restriccions sobre la mida de l'observació~\ref{obs:ar-metateorica} es mantenen: la versió semàntica parla de categories petites no buides, i l'extensió a categories grans demana la mateixa convenció metateòrica addicional.

\paragraph{Abast.}
Tots dos teoremes parlen de fórmules de $\mathcal L_{\mathrm{Cat}}$. Una afirmació que mencioni elements d'un conjunt, una categoria concreta o qualsevol estructura externa no és una fórmula d'aquest llenguatge i no té dual fins que no s'hi reformula. Per això el capítol~\ref{cap:categories} enuncia el principi per a afirmacions \emph{expressables en el llenguatge de les categories}, i no per a \enquote{qualsevol afirmació sobre categories}. Els enunciats que involucren diverses categories i functors entre elles demanen un llenguatge més ric; el mecanisme és el mateix, però cal dualitzar també les categories implicades, com mostra l'exemple de les categories sobre i sota un objecte.

\paragraph{Un mecanisme general.}
El que hem fet ---una involució del llenguatge que deixa invariant el conjunt d'axiomes--- no és exclusiu de la teoria de categories. La lògica clàssica en dona una altra instància: en el llenguatge de les àlgebres booleanes, la involució que intercanvia $\wedge$ amb $\vee$ i el cert amb el fals deixa els axiomes invariants, i produeix la dualitat de De~Morgan. Un exemple encara més antic és la \textbf{dualitat projectiva}\index{dualitat!projectiva}: en el pla projectiu, intercanviar \emph{punt} per \emph{recta} ---i \enquote{passar per} per \enquote{tallar-se en}--- deixa els axiomes invariants, de manera que cada teorema té un teorema dual, com el de Desargues o la parella Pascal--Brianchon. En tots els casos actua el mateix mecanisme, aplicat a llenguatges diferents. Quan una categoria codifica lògica, la dualitat de fletxes intercanvia conjunció i disjunció si existeixen; que aquesta dualitat coincideixi exactament amb la de De~Morgan requereix, a més, un complement involutiu, i és una hipòtesi addicional sobre la categoria, no una conseqüència automàtica de la dualitat categòrica.

\section{Una observació computacional}
\label{sec:ar-computacional}

La presentació relacional és especialment transparent per a l'axiomàtica, però no és necessàriament la més còmoda per a una implementació. Un programa prefereix, habitualment, operacions que retornin valors, $\dom\colon M\to M$ i $\cod\colon M\to M$, i una composició parcial. Aquesta darrera es pot representar, per exemple, com $\operatorname{comp}\colon M\times M\longrightarrow M_{\bot}$ o mitjançant un tipus opcional que distingeixi entre una composta existent i una parella no componible.

No hi ha contradicció entre les dues perspectives. La teoria relacional demostra primer existència i unicitat; aquests resultats justifiquen després el pas a una interfície funcional. El millor llenguatge per axiomatitzar una estructura no ha de coincidir amb el millor llenguatge per implementar-la.

\section{Connexió amb una lectura estructural}
\label{sec:ar-estructural}

La presentació anterior respon a un criteri estructural molt simple: no introduir com a primitives entitats que poden reconstruir-se a partir de les relacions fonamentals de la teoria. En la presentació usual es parla de dos tipus de dades, objectes i fletxes; en la relacional elemental només hi ha una mena d'individus, i els objectes apareixen com les fletxes que satisfan $\Dom(e,e)\land\Cod(e,e)$. L'objecte no és, doncs, una entitat afegida al costat de les fletxes, sinó una posició estructural determinada per les relacions de domini i codomini ---exactament el gir que la introducció anunciava en preferir \emph{relacionar} a \emph{individualitzar} (secció~\ref{sec:mirar-fora}).

Aquesta perspectiva és compatible amb una lectura estructural de la teoria de conjunts: les realitzacions conjuntistes concretes poden quedar al metallenguatge, mentre que el llenguatge de la teoria categòrica reté només les relacions que expressen l'estructura que volem estudiar. Això no elimina la teoria de conjunts del metallenguatge ---la semàntica estàndard de primer ordre continua sent conjuntista, i el marc de mida el fixa l'apèndix corresponent--- però evita confondre una realització externa amb les nocions primitives de la teoria objecte.

La lliçó metodològica és, doncs, doble. D'una banda, la formulació relacional mostra que una categoria es pot reconstruir a partir d'un univers de fletxes i de tres relacions primitives. De l'altra, un cop feta la reconstrucció, no hi ha cap inconvenient a recuperar la notació funcional i la terminologia usual, que són molt més eficients per desenvolupar la teoria ---i que són, de fet, les que la resta d'aquest llibre empra.

\chapter{Grafs}\label{cap:grafs}

Aquest apèndix desenvolupa amb detall la noció de graf dirigit, els seus homomorfismes i els camins, i mostra com generar una categoria a partir d'un graf. El material amplia, amb molts exemples concrets, la subsecció~\ref{subsec:grafs} del capítol de categories.

\section{Grafs}

Qualsevol sistema de relacions, com les rutes de transport o les connexions d'aerolínies, es pot dibuixar com una col·lecció de punts i fletxes. En aquesta etapa no ens importa què són els punts, sinó com estan connectats. Aquesta idea es modela amb el concepte de graf. En general, un graf consta de punts (o objectes), anomenats \emph{vèrtexs}, alguns dels quals estan units per \emph{arestes}; aquí considerarem que els grafs són \emph{dirigits}, és a dir, que cada aresta apunta en una direcció, d'un punt a un altre, i per això fem servir fletxes per connectar punts. Entre dos punts hi pot haver moltes fletxes, o cap, i fins i tot pot haver-hi fletxes d'un punt a si mateix. Sovint la manera més fàcil de descriure un graf petit és dibuixar-lo, com a la figura següent.
\begin{equation}\label{graf:1}
\xymatrix{
1 \ar@(ul,ur)[]^a \ar@/^/[dr]^b
& & 2 \ar@/^/[dl]_c \ar@/_/[dr]^d \ar[ll]_e
& & 3 \ar[ll]_f \ar@(ur,dr)[]^g \\
& 4 \ar@/^/[ul]^h \ar[rr]^k
& & 5 \ar[ur]_l & 6
}
\end{equation}
El tipus de graf que tractem aquí és una versió del graf dirigit que sovint s'anomena \emph{multigraf dirigit amb bucles}, encara que farem servir la paraula \enquote{graf} per abreujar. En donem la definició formal.

\begin{definicio}\index{graf dirigit}
Un \textbf{graf} $G$ comprèn:
\begin{itemize}
\item un conjunt $V_G$ de \textbf{vèrtexs};
\item un conjunt $E_G$ d'\textbf{arestes};
\item dues aplicacions $s_G\colon E_G\to V_G$ i $t_G\colon E_G\to V_G$, anomenades \textbf{origen} i \textbf{destí}. Si $s_G(f)=a$ i $t_G(f)=b$, escrivim $a\xrightarrow{f}b$ o $f\colon a\to b$, i diem que $a$ és el vèrtex d'origen i $b$ el de destí.
\end{itemize}
La notació
\[
G=\xymatrix{
E_G \ar@<0.6ex>[r]^{s_G} \ar@<-0.6ex>[r]_{t_G} & V_G
}
\]
indica que $G$ és un graf amb vèrtexs a $V_G$ i arestes a $E_G$.
\end{definicio}

\begin{exemple}
Considerem el graf $G$ de la figura següent:
\[
\xymatrix{
	& \bullet_{a} \ar@<1ex>[dl]^f \ar@<1ex>[dr]^g &  \\
	\bullet_{b} \ar@(ul,dl)_l \ar@<1ex>[ur]^h & & \bullet_{c}  \ar@<1ex>[ll]^k \ar@(ur,dr)^m \\
}
\]
Especifica quins són els seus elements.
\end{exemple}
\begin{solucio}
Es té $V_G=\{a,b,c\}$ i $E_G=\{f,g,h,k,l,m\}$, i les dues aplicacions són
\[
    \begin{array}{c|c|c}
        E_G & s_G & t_G \\ \hline
        f & a & b  \\
        g & a & c \\
        h & b & a \\
        k & c & b \\
        l & b & b \\
        m & c & c
    \end{array}
\]
\end{solucio}

\begin{exemple}
Defineix el graf $G$ perquè la seva representació gràfica sigui \eqref{graf:1}.
\end{exemple}
\begin{solucio}
Prenem $V_G=\{1,2,3,4,5,6\}$, $E_G=\{a,b,c,d,e,f,g,h,k,l\}$ i les aplicacions
\[
    \begin{array}{c|c|c}
        E_G & s_G & t_G \\ \hline
        a & 1 & 1  \\
        b & 1 & 4 \\
        c & 2 & 4 \\
        d & 2 & 5 \\
        e & 2 & 1 \\
        f & 3 & 2 \\
        g & 3 & 3 \\
        h & 4 & 1 \\
        k & 4 & 5 \\
        l & 5 & 3
    \end{array}
\]
que reprodueixen la figura \eqref{graf:1}.
\end{solucio}

\begin{exemple}\leavevmode
\begin{enumerate}
\item Dibuixa el graf corresponent a les dades següents:
\[
    \begin{array}{c|c|c}
        E_G & s_G & t_G \\ \hline
        f & 2 & 3 \\
        g & 2 & 3 \\
        h & 2 & 3 \\
        k & 4 & 3 \\
        l & 6 & 3 \\
        m & 6 & 6
    \end{array}
    \hspace{2cm} V_G=\{1,2,3,4,5,6\}
\]
\item Escriu la taula corresponent al graf $H$ següent:
\[
\xymatrix{
	1 \ar[r]^{a} & 2 \ar[d]_{g} \ar@/^2ex/[r]^-{b} \ar[r]^-{c} & 3 \ar[r]^-{e} \ar@/^2ex/[l]^-{d} & 4 \\
	5 \ar@(ul,dl)_l & 6 \ar[l]_-{h} \ar[r]^-{k} & 7 \ar[ur]^-{f}
}
\]
\end{enumerate}
\end{exemple}
\begin{solucio}\leavevmode
\begin{enumerate}
\item
\[
\xymatrix{
	4 \ar[r]^-{k} & 3 & 6 \ar[l]_-{l} \ar@(ur,dr)\\
	1 & 2 \ar@/^2ex/[u]^-{f} \ar[u]^-{g} \ar@/_2ex/[u]_{h} & 5
}
\]
\item
\[
    \begin{array}{c|c|c}
        E_H & s_H & t_H \\ \hline
        a & 1 & 2 \\
        b & 2 & 3 \\
        c & 2 & 3 \\
        d & 3 & 2 \\
        e & 3 & 4 \\
        f & 7 & 4 \\
        g & 2 & 6 \\
        h & 6 & 5 \\
        k & 6 & 7 \\
        l & 5 & 5
    \end{array}
    \hspace{2cm} V_H=\{1,2,3,4,5,6,7\}
\]
\end{enumerate}
\end{solucio}

\begin{exemple}
En una empresa petita hi ha tres departaments: Administració, Comptabilitat i Atenció al client. Durant el funcionament diari, els documents poden circular diverses vegades entre els mateixos departaments, i alguns departaments processen documents interns sense enviar-los enlloc. En un dia es donen les situacions següents:
\begin{itemize}
\item Dos documents passen d'Administració a Comptabilitat.
\item Un document passa de Comptabilitat a Administració.
\item Tres documents passen de Comptabilitat a Atenció al client.
\item Un document passa d'Atenció al client a Comptabilitat.
\item Dos documents són processats internament a Administració.
\item Un document és processat internament a Comptabilitat.
\item Un document és processat internament a Atenció al client.
\end{itemize}
Modela aquesta situació com un graf $G$, especificant-ne els elements, interpretant-ne les arestes i representant-lo gràficament.
\end{exemple}
\begin{solucio}
El graf $G$ té com a conjunt de vèrtexs $V_G=\{a,b,c\}$, on $a,b,c$ representen Administració, Comptabilitat i Atenció al client, respectivament. El conjunt d'arestes és
\[
\begin{aligned}
E_G = \{
& a\xrightarrow{f_1}b,\ a\xrightarrow{f_2}b, \\
& b\xrightarrow{g}a, \\
& b\xrightarrow{h_1}c,\ b\xrightarrow{h_2}c,\ b\xrightarrow{h_3}c, \\
& c\xrightarrow{k}b, \\
& a\xrightarrow{j_1}a,\ a\xrightarrow{j_2}a, \\
& b\xrightarrow{i}b, \\
& c\xrightarrow{l}c
\},
\end{aligned}
\]
on els subíndexs indiquen arestes diferents entre els mateixos vèrtexs (diversos documents amb el mateix recorregut) i els bucles $a\xrightarrow{j}a$, $b\xrightarrow{i}b$ i $c\xrightarrow{l}c$ representen processament intern. La representació gràfica de $G$ és
\[
\xymatrix@R=7ex@C=8ex{
	c \ar@(ul,dl)_l \ar[r]^-{k} & b \ar@/^1.5ex/[l]^-{h_2} \ar@/_2.5ex/[l]_{h_1} \ar@/_5ex/[l]_{h_3}  \ar@(dr,ur)_i \ar@/^4ex/[d]^-{g} \\
	& a \ar@/^4ex/[u]^-{f_1} \ar@/^1ex/[u]_(.45){f_2}  \ar@(dl,dr)_{j_{1}}  \ar@(dr,ur)_{j_{2}}
}
\]
Aquest graf modela el flux real de documents dins l'empresa en un dia: les direccions indiquen el sentit del flux, el nombre d'arestes entre dos departaments indica la quantitat de documents, i els bucles indiquen activitat interna.
\end{solucio}

\begin{exercici}
Defineix el graf $G$ perquè la seva representació gràfica sigui la següent:
\[
\xymatrix{
	\bullet \ar@<1ex>[r]  & \bullet \ar@<1ex>[l] \ar@(ul,ur)  & \bullet \ar[l] \ar@(ur,dr) \ar[dl] & \\
	\bullet \ar[ur] \ar[r]  & \bullet  & \bullet
}
\]
\end{exercici}

\begin{exercici}
Dibuixa el graf corresponent a les dades següents:
\[
    \begin{array}{c|c|c}
        E_G & s_G & t_G \\ \hline
        f & v & w \\
        g & w & x \\
        h & w & x \\
        i & y & y \\
        j & y & z \\
        k & z & y
    \end{array}
    \hspace{2cm} V_G=\{v,w,x,y,z\}
\]
\end{exercici}

\begin{exercici}
En un centre logístic hi ha quatre seccions: Recepció, Magatzem, Preparació de comandes i Expedició. Els paquets es mouen entre les seccions al llarg del dia, i alguns processos es fan dins d'una mateixa secció. En una jornada s'observen els fluxos següents:
\begin{itemize}
\item Dos paquets passen de Recepció a Magatzem.
\item Un paquet passa de Magatzem a Recepció.
\item Tres paquets passen de Magatzem a Preparació de comandes.
\item Dos paquets passen de Preparació de comandes a Expedició.
\item Un paquet passa d'Expedició a Magatzem.
\item Dos paquets són reprocessats dins de Magatzem.
\item Un paquet és revisat dins de Preparació de comandes.
\item Un paquet és verificat dins d'Expedició.
\end{itemize}
Modela aquesta situació mitjançant un graf. Especifica'n els elements i representa'l gràficament.
\end{exercici}

\section{Homomorfismes de grafs}\label{sec:grafs-homomorfismes}

L'estructura d'un graf
\[
G=\xymatrix{
E_G \ar@<0.6ex>[r]^{s_G} \ar@<-0.6ex>[r]_{t_G} & V_G
}
\]
consisteix en dos conjunts $V_G$ i $E_G$ i dues aplicacions $s_G$ i $t_G$. Per relacionar dos grafs, hem de poder relacionar adequadament els seus conjunts i les seves aplicacions.

\begin{definicio}\index{morfisme de grafs}
Considerem dos grafs
\[
G=\xymatrix{
E_G \ar@<0.6ex>[r]^{s_G} \ar@<-0.6ex>[r]_{t_G} & V_G
}
\qquad\text{i}\qquad
H=\xymatrix{
E_H \ar@<0.6ex>[r]^{s_H} \ar@<-0.6ex>[r]_{t_H} & V_H
}
\]
Un \textbf{homomorfisme} $F$ de $G$ a $H$, que escrivim $F\colon G\to H$, és un parell d'aplicacions $F_V\colon V_G\to V_H$ i $F_E\colon E_G\to E_H$ tals que els diagrames següents commuten:
\begin{equation}\label{grafs:hom}
\xymatrix{
	E_G \ar[d]_-{F_E} \ar[r]^-{s_G} & V_G \ar[d]^-{F_V}\\
	E_H \ar[r]_-{s_H} & V_H
}
\hspace{2cm}
\xymatrix{
	E_G \ar[d]_-{F_E} \ar[r]^-{t_G} & V_G \ar[d]^-{F_V}\\
	E_H \ar[r]_-{t_H} & V_H
}
\end{equation}
\end{definicio}

Dit breument, un homomorfisme \emph{preserva} l'origen (diagrama de l'esquerra),
\[
F_V\comp s_G=s_H\comp F_E,
\]
i el destí (diagrama de la dreta),
\[
F_V\comp t_G=t_H\comp F_E.
\]

\begin{exemple}
Considerem els dos grafs següents:
\[
G:\ 
\xymatrix{
	1 \ar[d]_{h} \ar[r]^{f} & 2 \ar[d]^{g} \ar@(ur,ul)_{k} \\
	4 & 3
}
\qquad\text{i}\qquad
H:\ 
\xymatrix{
	b \ar@(ur,ul)_{y} \ar@(dr,ur)_{z} \\
	a \ar[u]_{x} \ar[r]_{v} & c
}
\]
Defineix un homomorfisme $F$ entre ells.
\end{exemple}
\begin{solucio}
Com que $k$ és un bucle i l'únic vèrtex de $H$ amb bucles és $b$, necessàriament $F_V(2)=b$. Com que $g$ surt de $2$, la seva imatge ha de sortir de $b$, i les úniques arestes que surten de $b$ són els bucles $y,z$; per tant $g$ també ha d'anar a un bucle i $F_V(3)=b$. Aquesta part és forçada; la resta és una elecció. Per construir un exemple concret, triem $F_V(1)=a$ i $F_V(4)=c$, i posem $F_E(k)=F_E(g)=y$, $F_E(f)=x$ i $F_E(h)=v$. Les taules següents comproven que aquesta tria preserva origen i destí.

La distinció entre una condició necessària i una elecció és important: que $F_V(2)=F_V(3)=b$ estigui forçat no vol dir que l'homomorfisme sigui únic. De fet, una enumeració finita directa de les aplicacions compatibles dona $24$ homomorfismes entre aquests dos grafs, dels quals $8$ tenen $F_V(1)=a$ i $16$ tenen $F_V(1)=b$.

Comprovem que les condicions \eqref{grafs:hom} es compleixen. Les taules dels dos grafs són
\[
    \begin{array}{c|c|c}
        E_G & s_G & t_G \\ \hline
        f & 1 & 2  \\
        g & 2 & 3 \\
        h & 1 & 4 \\
        k & 2 & 2 
    \end{array}
    \qquad
    \begin{array}{c|c|c}
        E_H & s_H & t_H \\ \hline
        v & a & c  \\
        x & a & b \\
        y & b & b \\
        z & b & b 
    \end{array}
\]
i les dues aplicacions de l'homomorfisme són
\[
F_V =
\begin{array}{c|c}
    V_G & V_H \\ \hline
     1 & a \\
     2 & b \\
     3 & b \\
     4 & c
\end{array}
\qquad
F_E =
\begin{array}{c|c}
   E_G  & E_H \\ \hline
   f & x \\
   g & y \\
   h & v \\
   k & y
\end{array}
\]
Per exemple,
\[
\begin{aligned}
(F_V\comp s_G)(f)&=F_V(s_G(f))\\
&=F_V(1)=a,\\
(s_H\comp F_E)(f)&=s_H(F_E(f))\\
&=s_H(x)=a,
\end{aligned}
\]
de manera que $(F_V\comp s_G)(f)=(s_H\comp F_E)(f)$; deixem la resta d'arestes com a exercici. Anàlogament,
\[
\begin{aligned}
(F_V\comp t_G)(h)&=F_V(t_G(h))\\
&=F_V(4)=c,\\
(t_H\comp F_E)(h)&=t_H(F_E(h))\\
&=t_H(v)=c,
\end{aligned}
\]
i per tant $(F_V\comp t_G)(h)=(t_H\comp F_E)(h)$.
\end{solucio}

\begin{exercici}
Considerem els grafs
\[
G:\ 
\xymatrix{
    1 \ar[r]^-{a} & 2 \ar[r]^-{b}& 3
}
\qquad
H:\ 
\xymatrix{
    4 \ar@<0.6ex>[r]^{d} \ar@<-0.6ex>[r]_{c} & 5 \ar@(ur,dr)^{e}
}
\]
Troba un homomorfisme $F$ de $G$ a $H$ i comprova les igualtats \eqref{grafs:hom}.
\end{exercici}

\begin{exercici}
Considerem el graf $G$ definit per
\[
    \begin{array}{c|c|c}
        E_G & s_G & t_G \\ \hline
        a & 1 & 2 \\
        b & 2 & 3 \\
        c & 1 & 4 \\
        d & 1 & 4 \\
        e & 5 & 6 
    \end{array}
    \hspace{2cm} V_G=\{1,2,3,4,5,6\}
\]
i el graf $H$ representat per
\[
\xymatrix{
	m \ar[d]_{y} \ar@<0.6ex>[r]^{w} & n \ar[l]^{x} \\
	p  & q \ar@(dr,ur)_{z}
}
\]
Es demana:
\begin{enumerate}
\item representar gràficament el graf $G$;
\item trobar un homomorfisme $F$ de $G$ a $H$ i comprovar les igualtats \eqref{grafs:hom}.
\end{enumerate}
\end{exercici}

\section{Camins en un graf}

Sigui $G$ un graf. Un \emph{camí} a $G$ és qualsevol successió d'arestes tal que el destí d'una aresta sigui l'origen de la següent. Això inclou successions de longitud~1, que són elements de $E_G$, i successions de longitud~0, que comencen i acaben al mateix vèrtex sense seguir cap aresta.

\begin{definicio}
Sigui $n\geq 1$. En un graf
\[
G=\xymatrix{
E_G \ar@<0.6ex>[r]^{s_G} \ar@<-0.6ex>[r]_{t_G} & V_G
}
\]
un \textbf{camí de longitud} $n$ del vèrtex $a$ al vèrtex $b$ és una successió d'arestes $(f_1,f_2,\dots,f_n)$, escrita en \emph{ordre de recorregut} ---primer es recorre $f_1$---, tal que
\begin{enumerate}
\item[(i)] $s_G(f_1)=a$;
\item[(ii)] $t_G(f_i)=s_G(f_{i+1})$ per a $i=1,\dots,n-1$;
\item[(iii)] $t_G(f_n)=b$.
\end{enumerate}
El conjunt de camins de longitud $n$ s'escriu $C_n$. En particular, $C_2$ és el conjunt de parells d'arestes $(f,g)$ amb $t_G(f)=s_G(g)$, anomenats parells d'\emph{arestes componibles}.
\end{definicio}

Per conveni, per a cada vèrtex $a\in V_G$ hi ha un únic camí de longitud~0 d'$a$ a si mateix, que denotem $()_a$ i anomenem \emph{camí buit} d'$a$; és el cas $n=0$, exclòs de la definició anterior perquè aleshores no hi ha cap aresta $f_1,\dots,f_n$. Si dibuixem un camí
\[
\xymatrix{
\cdot \ar[r]^{f_{1}} & \cdot \ar[r]^-{f_{2}} & \cdots \ar[r]^{f_{n-1}} & \cdot \ar[r]^{f_{n}} & \cdot
}
\]
l'aresta $f_1$ queda a l'esquerra i es recorre primer; $f_n$ queda a la dreta i es recorre última. Si $f\in E_G$, aleshores $(f)$ és un camí de longitud~1, és a dir, $(f)\in C_1$.

\begin{exemple}\label{ex:grafs-camins-G}
Considerem el graf $G$ de l'exemple d'homomorfisme de la secció anterior:
\[
G:\ 
\xymatrix{
	1 \ar[d]_{h} \ar[r]^{f} & 2 \ar[d]^{g} \ar@(ur,ul)_{k} \\
	4 & 3
}
\]
Escriu els camins de longitud $0$, $1$, $2$ i $3$ de $G$, i descriu tots els camins d'un vèrtex a un altre.
\end{exemple}
\begin{solucio}
Recordem la taula de $G$: $f\colon 1\to 2$, $g\colon 2\to 3$, $h\colon 1\to 4$ i el bucle $k\colon 2\to 2$.

\emph{Longitud $0$.} Hi ha un camí buit per a cada vèrtex: $C_0=\{()_1,\ ()_2,\ ()_3,\ ()_4\}$.

\emph{Longitud $1$.} Són les arestes: $C_1=\{(f),\ (g),\ (h),\ (k)\}$.

\emph{Longitud $2$.} Busquem els parells $(u,v)$ amb $t_G(u)=s_G(v)$. Com que $t_G(f)=t_G(k)=2$ i les arestes que surten de $2$ són $g$ i $k$, mentre que de $3$ i de $4$ no en surt cap,
\[
C_2=\{(f,g),\ (f,k),\ (k,g),\ (k,k)\}.
\]

\emph{Longitud $3$.} Cada camí de longitud $2$ que acaba a $2$ es pot allargar amb $g$ o amb $k$; els que acaben a $3$ no es poden allargar:
\[
C_3=\{(f,k,g),\ (f,k,k),\ (k,k,g),\ (k,k,k)\}.
\]

\emph{Descripció general.} Escrivim $k^n$ per a la successió $(k,\dots,k)$ de $n$ bucles, amb $k^0$ buida. Els camins d'un vèrtex a un altre són exactament:
\[
\begin{array}{c|l}
\text{d'origen a destí} & \text{camins} \\ \hline
1\to 1,\ 3\to 3,\ 4\to 4 & \text{només el camí buit} \\
1\to 2 & (f,k^n),\ n\geq 0 \\
1\to 3 & (f,k^n,g),\ n\geq 0 \\
1\to 4 & (h) \\
2\to 2 & k^n,\ n\geq 0 \ (\text{per a } n=0,\ \text{el camí buit } ()_2) \\
2\to 3 & (k^n,g),\ n\geq 0
\end{array}
\]
i no n'hi ha cap més. El bucle $k$ fa que el nombre de camins sigui infinit, però per a cada longitud $n\geq 2$ n'hi ha exactament quatre.
\end{solucio}

\begin{exemple}\label{ex:grafs-camins-H}
Fem el mateix amb el graf $H$ del mateix exemple:
\[
H:\ 
\xymatrix{
	b \ar@(ur,ul)_{y} \ar@(dr,ur)_{z} \\
	a \ar[u]_{x} \ar[r]_{v} & c
}
\]
Escriu els camins de longitud $0$, $1$ i $2$ de $H$, i compta quants camins de longitud $n$ hi ha.
\end{exemple}
\begin{solucio}
La taula de $H$ és $x\colon a\to b$, $v\colon a\to c$ i els bucles $y,z\colon b\to b$.

\emph{Longitud $0$ i $1$.} $C_0=\{()_a,\ ()_b,\ ()_c\}$ i $C_1=\{(x),\ (v),\ (y),\ (z)\}$.

\emph{Longitud $2$.} De $c$ no surt cap aresta, de manera que $v$ no es pot allargar. Després de $x$, $y$ o $z$ som a $b$, i podem continuar amb $y$ o amb $z$:
\[
C_2=\{(x,y),\ (x,z),\ (y,y),\ (y,z),\ (z,y),\ (z,z)\}.
\]

\emph{Recompte.} Un camí de longitud $n\geq 1$ de $b$ a $b$ és una successió qualsevol de $n$ bucles, triant a cada pas entre $y$ i $z$: n'hi ha $2^n$. Un camí de longitud $n\geq 2$ que surt d'$a$ ha de començar per $x$ i continuar amb $n-1$ bucles: n'hi ha $2^{n-1}$. L'únic camí que va d'$a$ a $c$ és $(v)$. Per tant, per a $n\geq 2$ hi ha
\[
2^n+2^{n-1}=3\cdot 2^{n-1}
\]
camins de longitud $n$; per exemple, $6$ de longitud $2$, d'acord amb la llista anterior.

Comparat amb l'exemple anterior, un sol bucle produeix un nombre constant de camins per a cada longitud, mentre que dos bucles al mateix vèrtex en produeixen un nombre que creix exponencialment.
\end{solucio}

\section{De grafs a categories}

Podem passar d'un graf a una categoria afegint-hi composició i identitats. La construcció següent és fonamental: mostra que tot graf genera una categoria de manera lliure.

Per a qualsevol graf $G$ hi ha una categoria $\Free(G)$ els objectes de la qual són els vèrtexs de $G$ i els morfismes de la qual són els camins de $G$. Un camí d'$a$ a $b$ és un morfisme $a\to b$. La composició s'obté concatenant camins, en \emph{ordre de recorregut}: si $p=(a_1,\dots,a_r)$ és un camí d'$A$ a $B$ i $q=(b_1,\dots,b_s)$ un camí de $B$ a $C$, la seva composició és el camí d'$A$ a $C$
\[
q\comp p=(a_1,\dots,a_r,b_1,\dots,b_s),
\]
on, d'acord amb la convenció categòrica $q\comp p$, el camí $p$ ---el que surt d'$A$--- es recorre primer. Per a cada vèrtex $A$, la identitat $\id_A$ és el camí buit $()_A$.

Comprovem que $\Free(G)$ és una categoria. La composició és associativa perquè la concatenació de successions ho és: en concatenar tres camins componibles, el resultat és la successió de totes les seves arestes en ordre, amb independència de l'ordre en què s'agrupin. El camí buit és neutre, ja que concatenar-lo no afegeix cap aresta: $(a_1,\dots,a_r)\comp()_A=(a_1,\dots,a_r)=()_B\comp(a_1,\dots,a_r)$ per a tot camí d'$A$ a $B$. Per tant, $\Free(G)$ satisfà els axiomes de categoria. S'anomena \textbf{categoria lliure generada pel graf} $G$, o també \textbf{categoria de camins} de $G$.\index{categoria!lliure}\index{categoria!de camins}

\begin{exemple}\label{ex:grafs-free-cadena}
Especifica la categoria $\Free(G)$ generada pel graf
\[
G:\ 
\xymatrix{
    1 \ar[r]^-{a} & 2 \ar[r]^-{b}& 3
}
\]
de l'exercici d'homomorfismes de la secció~\ref{sec:grafs-homomorfismes}.
\end{exemple}
\begin{solucio}
Els objectes són els vèrtexs $1$, $2$ i $3$. Els morfismes són els camins, i com que el graf no té cicles, n'hi ha un nombre finit:
\[
\begin{array}{c|l}
\Hom(1,1)=\{()_1\} & \Hom(1,2)=\{(a)\} \\
\Hom(2,2)=\{()_2\} & \Hom(2,3)=\{(b)\} \\
\Hom(3,3)=\{()_3\} & \Hom(1,3)=\{(a,b)\}
\end{array}
\]
i tots els altres homconjunts ($\Hom(2,1)$, $\Hom(3,1)$, $\Hom(3,2)$) són buits. En total, sis morfismes: les tres identitats, les dues arestes i el camí de longitud~$2$.

L'única composició que no involucra identitats és la de $(a)$ i $(b)$:
\[
\xymatrix{
1 \ar[r]^-{(a)} \ar[dr]_-{(a,b)} & 2 \ar[d]^-{(b)} \\
& 3
}
\qquad\qquad
(b)\comp(a)=(a,b).
\]
Les composicions amb identitats són les que imposen els axiomes, perquè els camins buits són neutres per a la concatenació.

Com que entre dos objectes hi ha com a màxim un morfisme, $\Free(G)$ és una categoria prima. De fet, és la categoria associada a l'ordre $1<2<3$ ---hi ha un morfisme $i\to j$ exactament quan $i\leq j$---, isomorfa, doncs, a la categoria $\cat{3}$ de l'ordre $0<1<2$ que apareix al capítol~\ref{cap:categories}.
\end{solucio}

\begin{exemple}\label{ex:grafs-free-bucle}
Especifica la categoria $\Free(H)$ generada pel graf
\[
H:\ 
\xymatrix{
    4 \ar@<0.6ex>[r]^{d} \ar@<-0.6ex>[r]_{c} & 5 \ar@(ur,dr)^{e}
}
\]
del mateix exercici.
\end{exemple}
\begin{solucio}
Els objectes són $4$ i $5$. Escrivim $e^n$ per al camí $(e,\dots,e)$ de $n$ bucles, amb $e^0=()_5$. Els camins de $H$ donen els homconjunts següents:
\[
\begin{aligned}
\Hom(4,4)&=\{()_4\},\\
\Hom(5,5)&=\{e^n : n\geq 0\},\\
\Hom(4,5)&=\{(c,e^n) : n\geq 0\}\cup\{(d,e^n) : n\geq 0\},\\
\Hom(5,4)&=\varnothing.
\end{aligned}
\]
De $5$ a $4$ no hi ha cap camí, perquè cap aresta surt de $5$ cap a $4$.

La composició és la concatenació. Per a bucles, $e^m\comp e^n=e^{m+n}$, i per als camins de $4$ a $5$,
\[
\xymatrix{
4 \ar[r]^-{(c,e^n)} \ar[dr]_-{(c,e^{n+m})} & 5 \ar[d]^-{e^m} \\
& 5
}
\qquad\qquad
e^m\comp(c,e^n)=(c,e^{n+m}),
\]
i igual amb $d$ en lloc de $c$.

Hi ha dues observacions a fer. Primera: $\Hom(5,5)$ amb la composició és el monoide $(\mathbb N,+,0)$, via $e^n\mapsto n$; la subcategoria plena sobre l'objecte $5$ és, doncs, la categoria d'un sol objecte d'aquest monoide. Segona: les arestes paral·leles $c$ i $d$ donen morfismes diferents, $(c)\neq(d)$, i també $(c,e^n)\neq(d,e^n)$ per a tot $n$. La categoria lliure no identifica mai dos camins diferents: no imposa cap relació que no vingui forçada pels axiomes.
\end{solucio}

\chapter{Convenció de mida}
\label{cap:mida}

Al llarg del llibre hem parlat de categories \emph{petites}, \emph{localment petites} i de \emph{classes pròpies} sense fixar del tot el marc conjuntista. Aquest apèndix el fa explícit, perquè el volum de lògica categòrica ---on apareixen categories de models, doctrines i fibracions--- l'exigirà.

\medskip
\noindent\textbf{Marc adoptat: classes i conjunts.} Treballem en una teoria de conjunts amb \textbf{classes pròpies}, a l'estil de von Neumann--Bernays--Gödel (NBG). Hi ha dos nivells de col·leccions: els \emph{conjunts}, que poden ser elements d'altres col·leccions, i les \emph{classes pròpies}, que no. Amb això:
\begin{itemize}
  \item una categoria és \textbf{petita} si la seva col·lecció d'objectes i la de morfismes són totes dues \emph{conjunts};
  \item és \textbf{localment petita} si tots els homconjunts $\Hom(A,B)$ són conjunts, tot i que la col·lecció d'objectes pot ser una classe pròpia;
  \item és \textbf{ben potenciada} si, per a cada objecte $A$, els subobjectes de $A$ formen un conjunt.
\end{itemize}
L'última condició demana una precisió. Un \emph{subobjecte} de $A$ no és un monomorfisme concret cap a $A$, sinó una classe d'equivalència de monomorfismes. Dos monomorfismes $m\colon S\rightarrowtail A$ i $n\colon T\rightarrowtail A$ són equivalents quan cadascun factoritza a través de l'altre, és a dir, quan difereixen en un isomorfisme entre els dominis. La distinció és essencial, i en el marc de classes cal formular-la amb cura. A $\cat{Set}$, els monomorfismes cap a un conjunt donat formen una classe pròpia, perquè cada subconjunt té moltes còpies isomorfes; i cada classe d'equivalència és, ella mateixa, una classe pròpia. Com que les classes pròpies no poden ser elements de cap col·lecció, \enquote{el conjunt de les classes d'equivalència} no té sentit literal. La definició precisa és, doncs, la següent: una categoria és \textbf{ben potenciada} si, per a cada objecte $A$, hi ha un \emph{conjunt} $R_A$ de monomorfismes cap a $A$ tal que tot monomorfisme cap a $A$ és equivalent a algun element de $R_A$. Aleshores es defineix $\Sub(A)$ com el conjunt quocient $R_A/{\sim}$, i $[m]$ designa la classe dels representants equivalents a $m$. Un altre conjunt de representants dona un quocient en bijecció canònica amb aquest, i la bijecció respecta l'ordre; per això el resultat no depèn de la tria de $R_A$. A $\cat{Set}$ es pot prendre com a $R_A$ el conjunt de les inclusions dels subconjunts de $A$, i s'obté $\Sub(A)\cong\mathcal P(A)$. La resta del llibre parla de \enquote{classes de monomorfismes} en aquest sentit.

Les categories $\cat{Set}$, $\cat{Grp}$, $\cat{Top}$, $\cat{Graph}$ i $\cat{Cat}$ són localment petites però no petites; cada preordre concret sobre un conjunt i cada categoria finita concreta són petites. Això no s'ha de confondre amb la petitesa de la categoria de \emph{tots} els preordres o de \emph{totes} les categories finites, que ---en tenir una classe pròpia de conjunts subjacents--- són grans llevat que es fixi una convenció de mida o s'esculli un conjunt de representants.

Quan alguna construcció demani triar simultàniament un objecte (o un morfisme) per a cada objecte d'una categoria \emph{gran} ---per exemple, en construir un invers d'una equivalència---, la família de tries està indexada per una classe pròpia, i l'axioma de l'elecció ordinari, que val per a famílies indexades per un conjunt, no hi arriba. En aquests casos assumim la forma \textbf{global} de l'elecció (disponible, per exemple, a NBG amb l'axioma d'elecció global). Ho indicarem allà on calgui.

\medskip
\noindent\textbf{L'alternativa dels universos.} Una segona opció, molt usada en geometria algebraica i teoria de topos, fixa un o més \textbf{universos de Grothendieck} $\mathcal{U}$ ---conjunts prou grans i tancats per les operacions habituals--- i parla de categories $\mathcal{U}$-petites i $\mathcal{U}$-grans de manera \emph{relativa} a $\mathcal{U}$. Aquesta convenció té l'avantatge que \enquote{la categoria de tots els $\mathcal{U}$-conjunts} és, ella mateixa, un objecte d'un univers més gran, cosa que simplifica certes construccions. En aquest volum \emph{no} l'adoptem, per no carregar una introducció amb la teoria d'universos; però el lector que continuï cap a la teoria de feixos la retrobarà. Les dues opcions són maneres alternatives de \emph{gestionar la mida} més que no pas fonaments idèntics: a grans trets, el paper de les \enquote{classes pròpies} de NBG el fa aquí la distinció entre $\mathcal{U}$-petit i $\mathcal{U}$-gran relativa a un univers $\mathcal{U}$ fixat. Fer precisa aquesta correspondència exigeix especificar l'univers i les hipòtesis pertinents; no cal, però, que això sigui un obstacle per a una primera lectura.

\medskip
\noindent Fixar aquesta convenció des d'ara evita l'ambigüitat que, altrament, es faria notar quan apareguin categories de models i functors classificadors: en tot moment, cada construcció d'aquest llibre declara si viu en un conjunt o en una classe pròpia.

\medskip
\noindent\textbf{Per saber-ne més.} Les qüestions de mida es tracten, al nivell d'aquest llibre, a \cite[§1.8]{awodey}, \cite[§1.1]{riehl} ---on l'observació~1.1.5 presenta la convenció d'universos---, \cite[§3.2]{leinster} i \cite[cap.~1]{perrone}. El capítol~3 de \cite{leinster} presenta, a més, la perspectiva en què el mateix llenguatge estructural fa de fonament, sense una teoria de conjunts prèvia; en són les fonts \cite{lawvere-etcs} i, com a introducció informal, \cite{leinster-rethinking}. Per a una comparació detallada de les diferents maneres de gestionar la mida ---classes, universos i altres alternatives---, adequada un cop se'n té una certa experiència, vegeu \cite{shulman-mida}.

\chapter{Solucions dels tests}
\label{cap:solucions-tests}

\ifimprimible
Aquest apèndix reuneix les respostes i les justificacions dels tests. En aquesta edició per imprimir, són la manera de corregir-los.
\else
Aquest apèndix reuneix les respostes i les justificacions dels tests interactius. Serveixen per als lectors de PDF que no admeten la correcció automàtica i per a l'ús en paper.
\fi Es recomana consultar-les només després d'haver respost tot el test: cada justificació explica per què la resposta correcta ho és i, sovint, on falla el distractor més temptador.

\imprimeixsolucionstests

\ifpublicacioparcial\else
\chapter[Quasiexemples i contraexemples]{Quasiexemples\\ i contraexemples}
\label{cap:contraexemples}

Una definició matemàtica s'entén de veritat quan se'n coneixen tant els exemples com els \emph{quasiexemples}: estructures que s'hi assemblen molt, que apareixen de manera natural, i que tanmateix fallen en un punt precís. Aquest apèndix recull quasiexemples i contraexemples per a les nocions centrals del llibre: categoria, functor, transformació natural i equivalència. Acaba amb models de la lògica que queden fora de l'explicació per topos dels capítols~\ref{cap:subobjectes}--\ref{cap:pont}.

La secció~\ref{sec:igualtat-equivalencia}, sobre igualtat, isomorfisme, isomorfisme natural i equivalència, és la més important. Confondre aquests nivells és un error freqüent, i greu, perquè porta a demostrar coses falses o a formular propietats que no tenen sentit categòric.

El nivell és el del llibre, i els exemples s'han triat perquè només demanin els coneixements que el llibre ja pressuposa. Tota afirmació que es pot justificar amb aquests coneixements es demostra. Quan un comentari remet a teories que el llibre no desenvolupa, es presenta de manera informativa i es diu explícitament.

% =====================================================================
\section{Què no és una categoria}
\label{sec:no-categoria}
% =====================================================================

\subsection{Les condicions que es poden trencar}

La definició~\ref{def:categoria} demana unes dades i dos axiomes. Per poder dir, exemple a exemple, quina condició falla, les desglossem en cinc:
\begin{enumerate}
\item[(C1)] \textbf{Extrems.} Cada morfisme $f$ té un \emph{únic} domini i un \emph{únic} codomini.
\item[(C2)] \textbf{Composició.} Per a cada parella componible $f\colon A\to B$, $g\colon B\to C$ hi ha un morfisme $g\comp f\colon A\to C$, determinat sense ambigüitat.
\item[(C3)] \textbf{Associativitat.} $h\comp(g\comp f)=(h\comp g)\comp f$.
\item[(C4)] \textbf{Existència d'identitats.} Per a cada objecte $A$ hi ha un morfisme $\id_A\colon A\to A$.
\item[(C5)] \textbf{Lleis d'identitat.} $f\comp\id_A=f=\id_B\comp f$ per a tot $f\colon A\to B$.
\end{enumerate}
Les igualtats de (C3) i (C5) són \emph{igualtats}: no isomorfismes, ni homotopies, ni cap mena d'equivalència. Gairebé tots els exemples d'aquesta secció fallen exactament aquí. Recordem també que la demostració de la unicitat de la identitat (capítol~\ref{cap:categories}) no fa servir l'associativitat: només les lleis d'identitat. Ho aprofitarem en parlar dels magmes.

\subsection{L'exemple central: els camins en un espai}

Aquest exemple fa servir la noció de continuïtat. Sigui $X$ un espai topològic. Proposem:
\begin{itemize}
\item \emph{objectes}: els punts de $X$;
\item \emph{morfismes}: els camins, és a dir, les aplicacions contínues $\gamma\colon[0,1]\to X$; escrivim $\gamma\colon x\to y$ si $\gamma(0)=x$ i $\gamma(1)=y$;
\item \emph{composició}: la concatenació. Si $\gamma\colon x\to y$ i $\delta\colon y\to z$, definim $\delta\ast\gamma\colon x\to z$ per
\[
(\delta\ast\gamma)(t)=
\begin{cases}
\gamma(2t), & 0\le t\le\tfrac12,\\
\delta(2t-1), & \tfrac12\le t\le1;
\end{cases}
\]
\item \emph{candidats a identitat}: els camins constants $c_x(t)=x$.
\end{itemize}
Escrivim $\delta\ast\gamma$ (primer $\gamma$, després $\delta$) per coherència amb $g\comp f$. Les dues branques coincideixen a $t=\frac12$, perquè $\gamma(1)=y=\delta(0)$, i per tant $\delta\ast\gamma$ és contínua. Així (C1) i (C2) es compleixen. El problema és a (C3) i (C5).

\paragraph{Fallada de l'associativitat.} Siguin $\gamma\colon x\to y$, $\delta\colon y\to z$ i $\varepsilon\colon z\to w$. Un càlcul directe dona
\[
\bigl((\varepsilon\ast\delta)\ast\gamma\bigr)(t)=
\begin{cases}
\gamma(2t), & t\in[0,\frac12],\\
\delta(4t-2), & t\in[\frac12,\frac34],\\
\varepsilon(4t-3), & t\in[\frac34,1],
\end{cases}
\qquad
\bigl(\varepsilon\ast(\delta\ast\gamma)\bigr)(t)=
\begin{cases}
\gamma(4t), & t\in[0,\frac14],\\
\delta(4t-1), & t\in[\frac14,\frac12],\\
\varepsilon(2t-1), & t\in[\frac12,1].
\end{cases}
\]
Tots dos recorren els mateixos tres trams en el mateix ordre, però a velocitats diferents.

\begin{exemple}[Tres segments a la recta]\label{ex:camins-recta}
A $X=\mathbb R$, prenem $\gamma(t)=t$, $\delta(t)=1+t$ i $\varepsilon(t)=2+t$, de manera que $\gamma\colon0\to1$, $\delta\colon1\to2$ i $\varepsilon\colon2\to3$. Avaluant a $t=\frac12$,
\[
\bigl((\varepsilon\ast\delta)\ast\gamma\bigr)(\tfrac12)=\gamma(1)=1,
\qquad
\bigl(\varepsilon\ast(\delta\ast\gamma)\bigr)(\tfrac12)=\varepsilon(0)=2.
\]
Els dos camins van de $0$ a $3$ i tenen la mateixa imatge, l'interval $[0,3]$, però són diferents: falla (C3).
\end{exemple}

\paragraph{Fallada de les identitats.} El fracàs és més radical que el de l'associativitat:

\begin{proposicio}\label{prop:camins-sense-identitat}
Sigui $\gamma\colon x\to y$ un camí injectiu. Aleshores cap camí $e\colon y\to y$ compleix $e\ast\gamma=\gamma$. Anàlogament, si $\gamma\colon y\to z$ és injectiu, cap camí $e\colon y\to y$ compleix $\gamma\ast e=\gamma$.
\end{proposicio}

\begin{prova}
Per a qualsevol $e$, $(e\ast\gamma)(\frac14)=\gamma(\frac12)$, que és diferent de $\gamma(\frac14)$ perquè $\gamma$ és injectiu. Per a l'altra afirmació, $(\gamma\ast e)(\frac34)=\gamma(\frac12)\neq\gamma(\frac34)$.
\end{prova}

A $X=\mathbb R$, entre dos punts diferents qualssevol hi ha camins injectius. Per tant, no és que s'hagin triat malament els camins constants: (C4)--(C5) no es poden satisfer de cap manera amb aquesta composició.

\paragraph{Però fallen de manera controlada.} Les igualtats que fallen valen \emph{llevat d'una reparametrització}, i les reparametritzacions són inofensives llevat d'homotopia. Dos camins $\alpha,\beta\colon x\to y$ són \emph{homòtops relativament als extrems}, i escrivim $\alpha\simeq\beta$, si hi ha una aplicació contínua $H\colon[0,1]\times[0,1]\to X$ amb $H(t,0)=\alpha(t)$, $H(t,1)=\beta(t)$, $H(0,s)=x$ i $H(1,s)=y$.

\begin{lema}[Reparametrització]\label{lema:reparametritzacio}
Sigui $\alpha\colon[0,1]\to X$ un camí i $\varphi\colon[0,1]\to[0,1]$ contínua amb $\varphi(0)=0$ i $\varphi(1)=1$. Aleshores $\alpha\comp\varphi\simeq\alpha$.
\end{lema}

\begin{prova}
L'aplicació $H(t,s)=\alpha\bigl((1-s)\varphi(t)+st\bigr)$ és contínua i ben definida, perquè $(1-s)\varphi(t)+st\in[0,1]$. Val $\alpha\comp\varphi$ per a $s=0$ i $\alpha$ per a $s=1$, i $H(0,s)=\alpha(0)$, $H(1,s)=\alpha(1)$ per a tot $s$.
\end{prova}

\begin{proposicio}\label{prop:associador-camins}
\begin{enumerate}
\item[(a)] $(\varepsilon\ast\delta)\ast\gamma=\bigl(\varepsilon\ast(\delta\ast\gamma)\bigr)\comp\varphi$, on $\varphi(t)=t/2$ a $[0,\frac12]$, $\varphi(t)=t-\frac14$ a $[\frac12,\frac34]$ i $\varphi(t)=2t-1$ a $[\frac34,1]$.
\item[(b)] $c_y\ast\gamma=\gamma\comp\psi$ amb $\psi(t)=\min(2t,1)$, i $\gamma\ast c_x=\gamma\comp\psi'$ amb $\psi'(t)=\max(2t-1,0)$.
\end{enumerate}
En conseqüència, $(\varepsilon\ast\delta)\ast\gamma\simeq\varepsilon\ast(\delta\ast\gamma)$, $c_y\ast\gamma\simeq\gamma$ i $\gamma\ast c_x\simeq\gamma$.
\end{proposicio}

\begin{prova}
Les igualtats (a) i (b) es comproven tram a tram amb les fórmules anteriors; per exemple, a $[\frac12,\frac34]$, $\varphi(t)=t-\frac14\in[\frac14,\frac12]$ i $\delta\bigl(4(t-\frac14)-1\bigr)=\delta(4t-2)$. Les funcions $\varphi$, $\psi$ i $\psi'$ són contínues i fixen $0$ i $1$, i el lema~\ref{lema:reparametritzacio} dona les homotopies.
\end{prova}

\paragraph{Primer remei: passar al quocient.} Si $\gamma\simeq\gamma'$ i $\delta\simeq\delta'$, concatenant les homotopies per a cada $s$ fix s'obté $\delta\ast\gamma\simeq\delta'\ast\gamma'$. Per tant, la concatenació indueix una composició ben definida entre classes, $[\delta]\comp[\gamma]:=[\delta\ast\gamma]$, i la proposició~\ref{prop:associador-camins} esdevé exactament (C3) i (C5) entre classes, amb $\id_x=[c_x]$. S'obté una categoria, el \textbf{grupoide fonamental} $\Pi_1(X)$\index{grupoide fonamental}: els objectes són els punts de $X$ i $\Hom(x,y)$ és el conjunt de classes d'homotopia de camins de $x$ a $y$. És un grupoide perquè el camí recorregut a l'inrevés, $\bar\gamma(t)=\gamma(1-t)$, és un invers de $[\gamma]$ (la comprovació és una homotopia explícita que no reproduïm). L'homconjunt $\Hom(x,x)$ és el que en topologia s'anomena \emph{grup fonamental} de $X$ en el punt $x$ (informatiu).

\paragraph{Segon remei: canviar els morfismes.} L'origen del problema és que tots els camins viuen en l'interval $[0,1]$, i cada concatenació ha de comprimir el temps a la meitat. Si cada camí té la seva durada, el problema desapareix. Un \textbf{camí de Moore}\index{camí de Moore} és un parell $(r,\gamma)$ amb $r\ge0$ i $\gamma\colon[0,r]\to X$ contínua, i es concatena sumant les durades:
\[
(s,\delta)\comp(r,\gamma)=(r+s,\ \delta\star\gamma),\qquad
(\delta\star\gamma)(t)=
\begin{cases}
\gamma(t), & t\in[0,r],\\
\delta(t-r), & t\in[r,r+s].
\end{cases}
\]

\begin{proposicio}
Els punts de $X$ i els camins de Moore formen una categoria, amb identitats $\id_x=(0,c_x)$.
\end{proposicio}

\begin{prova}
Les dues maneres d'agrupar tres camins $(r,\gamma)$, $(s,\delta)$ i $(u,\varepsilon)$ donen el camí de durada $r+s+u$ que val $\gamma(t)$ a $[0,r]$, $\delta(t-r)$ a $[r,r+s]$ i $\varepsilon(t-r-s)$ a $[r+s,r+s+u]$: \emph{la mateixa funció}. Les lleis d'identitat també són literals, perquè un camí de durada $0$ no ocupa temps: $(r,\gamma)\comp(0,c_x)$ té durada $r$ i val $\gamma(t-0)=\gamma(t)$.
\end{prova}

Tenim, doncs, les dues estratègies que reapareixen constantment en matemàtiques: \emph{quocientar}, és a dir, oblidar la informació que impedeix l'estrictesa (com a $\Pi_1(X)$), i \emph{estrictificar}, és a dir, substituir la composició per una altra que conservi la informació que interessa i en la qual les lleis valguin literalment (com amb els camins de Moore). Aquí no es pot parlar d'equivalència de categories, perquè l'estructura de partida no és una categoria. Els camins d'arestes d'un graf es concatenen sense problemes, amb el camí buit com a identitat ---és la categoria lliure de la definició~\ref{def:free}--- per la mateixa raó: la longitud s'hi suma i no es renormalitza mai.

\paragraph{Tercera via: fer teoria de la feblesa (informatiu).} La proposició~\ref{prop:associador-camins} no diu només que les dues associacions són homòtopes: en dona una homotopia \emph{concreta}, l'\emph{associador}. Amb quatre camins hi ha cinc maneres de posar els parèntesis, i els associadors les connecten formant un pentàgon. Si s'exigeix que els dos camins del pentàgon donin \emph{exactament} el mateix resultat, s'obté l'estructura d'una \emph{bicategoria}. Si només es demana que coincideixin llevat d'una nova homotopia, i així successivament a tots els nivells, s'arriba a les $\infty$-categories. Cap de les dues no es tracta en aquest llibre.

\subsection{Un sol objecte: els magmes}

Un \textbf{magma}\index{magma} és un conjunt $M$ amb una operació binària qualsevol, $(a,b)\mapsto ab$; un \emph{semigrup} és un magma associatiu, i un monoide, un semigrup amb element neutre. Com al capítol~\ref{cap:categories} per als monoides, a tot magma li podem associar una estructura $\cat{B}M$ amb un sol objecte $\ast$, un morfisme $\ast\to\ast$ per a cada $a\in M$ i la composició $a\comp b=ab$. Les condicions (C1) i (C2) es compleixen sempre, i $\cat{B}M$ és una categoria si i només si $M$ és un monoide. Cada magma no associatiu, o sense neutre, és un quasiexemple de categoria amb un objecte.

\begin{exemple}[La resta d'enters]
$M=(\mathbb Z,-)$. No és associatiu: $(1-1)-1=-1$, però $1-(1-1)=1$. Té neutre per la dreta, $a-0=a$, però no per l'esquerra: $e-a=a$ per a tot $a$ obligaria $e=2a$ per a tot $a$. Pel raonament de la unicitat de la identitat, un neutre bilateral hauria de coincidir amb $0$, que no ho és. Així, (C5) es compleix a mitges.
\end{exemple}

\begin{exemple}[Pedra, paper, tisora]\label{ex:ppt}
$M=\{\mathrm{Pe},\mathrm{Pa},\mathrm{Ti}\}$, on $xy$ és el guanyador de la partida entre $x$ i $y$, i $xx=x$:
\[
\begin{array}{c|ccc}
 & \mathrm{Pe} & \mathrm{Pa} & \mathrm{Ti}\\\hline
\mathrm{Pe} & \mathrm{Pe} & \mathrm{Pa} & \mathrm{Pe}\\
\mathrm{Pa} & \mathrm{Pa} & \mathrm{Pa} & \mathrm{Ti}\\
\mathrm{Ti} & \mathrm{Pe} & \mathrm{Ti} & \mathrm{Ti}
\end{array}
\]
És commutatiu i idempotent, però $(\mathrm{Pe}\,\mathrm{Pa})\,\mathrm{Ti}=\mathrm{Pa}\,\mathrm{Ti}=\mathrm{Ti}$ i $\mathrm{Pe}\,(\mathrm{Pa}\,\mathrm{Ti})=\mathrm{Pe}\,\mathrm{Ti}=\mathrm{Pe}$. Tampoc no té neutre: $\mathrm{Pe}\,e=\mathrm{Pe}$ exigeix $e\in\{\mathrm{Pe},\mathrm{Ti}\}$, $\mathrm{Pa}\,e=\mathrm{Pa}$ exigeix $e\in\{\mathrm{Pe},\mathrm{Pa}\}$ i $\mathrm{Ti}\,e=\mathrm{Ti}$ exigeix $e\in\{\mathrm{Pa},\mathrm{Ti}\}$, i la intersecció és buida. És un model finit, comprovable a mà, on fallen alhora (C3), (C4) i (C5).
\end{exemple}

\begin{exemple}[Altres magmes naturals]
\begin{itemize}
\item \emph{El punt mitjà}: $\mathbb R$ amb $ab=\frac{a+b}2$. Per exemple, $(0\cdot0)\cdot4=2$ i $0\cdot(0\cdot4)=1$; no hi ha neutre, perquè $a\cdot e=a$ força $e=a$.
\item \emph{El producte vectorial}: $\mathbb R^3$ amb $a\times b$. $(\mathbf i\times\mathbf i)\times\mathbf j=\mathbf 0$, però $\mathbf i\times(\mathbf i\times\mathbf j)=\mathbf i\times\mathbf k=-\mathbf j$. Satisfà, en lloc de l'associativitat, la identitat de Jacobi: és l'exemple bàsic d'àlgebra de Lie, una estructura perfectament bona que no és una categoria d'un objecte.
\item \emph{La coma flotant} (com a il·lustració): la suma de l'aritmètica de coma flotant d'un ordinador no és, en general, associativa. En doble precisió, una avaluació típica dona $(0.1+0.2)+0.3=0.6000000000000001$, mentre que $0.1+(0.2+0.3)=0.6$. Per convertir-ho en un magma literal caldria precisar el conjunt de valors i el tractament dels casos excepcionals (desbordaments, zeros amb signe...), cosa que no fem.
\item \emph{Associativitat sense neutre}: $(2\mathbb Z,\cdot)$ és un semigrup, però l'únic candidat a neutre seria $1\notin2\mathbb Z$. Aquesta fallada es repara de franc, afegint formalment un neutre.
\end{itemize}
\end{exemple}

\paragraph{Quocientar no sempre és un bon remei.} Davant d'un magma no associatiu podríem intentar el truc dels camins: quocientar per la relació d'equivalència més petita, compatible amb l'operació, que la faci associativa. Però el resultat pot ser catastròfic.

\begin{proposicio}
El quocient associatiu del magma pedra-paper-tisora té un sol element.
\end{proposicio}

\begin{prova}
En el quocient, $\mathrm{Ti}=(\mathrm{Pe}\,\mathrm{Pa})\,\mathrm{Ti}=\mathrm{Pe}\,(\mathrm{Pa}\,\mathrm{Ti})=\mathrm{Pe}$ i $\mathrm{Pe}=(\mathrm{Pa}\,\mathrm{Ti})\,\mathrm{Pe}=\mathrm{Pa}\,(\mathrm{Ti}\,\mathrm{Pe})=\mathrm{Pa}$.
\end{prova}

La diferència amb els camins és significativa. Allà la manca d'associativitat era \emph{coherent}, testimoniada per homotopies canòniques, i el quocient conservava la informació essencial. Aquí és arbitrària, i quocientar ho esborra tot. La jerarquia de les estructures que hem vist és aquesta:
\[
\renewcommand{\arraystretch}{1.2}
\begin{array}{lll}
\text{un sol objecte} & \text{diversos objectes} & \text{condicions}\\\hline
\text{magma} & \text{magmoide} & \text{(C1), (C2)}\\
\text{semigrup} & \text{semicategoria} & \text{(C1)--(C3)}\\
\text{monoide} & \text{categoria} & \text{(C1)--(C5)}\\
\text{grup} & \text{grupoide} & \text{(C1)--(C5) i inversos}
\end{array}
\]

\subsection{Altres maneres de fallar}

\begin{exemple}[Ordre estricte: una semicategoria]
Objectes: els nombres reals; un únic morfisme $x\to y$ si $x<y$, i cap altrament. La composició existeix per transitivitat i és associativa, perquè entre dos objectes hi ha com a molt un morfisme. Però no hi ha cap morfisme $x\to x$: falla (C4). Substituint $<$ per $\le$ s'obté una categoria.
\end{exemple}

\begin{exemple}[Grafs: no hi ha composició]
Un graf dirigit té vèrtexs i arestes, però dues arestes consecutives $x\to y\to z$ no tenen per què donar una aresta $x\to z$: falla (C2). El remei és la categoria lliure del graf (definició~\ref{def:free}).
\end{exemple}

\begin{exemple}[Aplicacions lineals no nul·les]
Objectes: els espais vectorials reals no nuls; morfismes: les aplicacions lineals \emph{no nul·les}. Les identitats hi són, i la composició és associativa quan està definida, però no sempre ho està: a $\mathbb R^2$, les projeccions $p(x,y)=(x,0)$ i $q(x,y)=(0,y)$ són no nul·les i $q\comp p=0$. Falla (C2).
\end{exemple}

\begin{exemple}[Funcions identificades amb el seu graf]
Si una funció és només el seu graf, $\{(1,1)\}$ és alhora una funció $\{1\}\to\{1\}$ i una funció $\{1\}\to\{1,2\}$: el codomini no queda determinat i falla (C1). En particular, \enquote{ser exhaustiva} deixa de ser una propietat de la funció. És el mateix fenomen que l'observació~\ref{obs:extrems-implicits} assenyala per a $\cat{Rel}$, i el remei és el mateix: el domini i el codomini formen part de les dades del morfisme.
\end{exemple}

\begin{exemple}[Un quocient que no és compatible amb la composició]
Al monoide de les funcions contínues $\mathbb R\to\mathbb R$ amb la composició, identifiquem $f\sim g$ si $f(0)=g(0)$, i intentem definir $[g]\comp[f]=[g\comp f]$. Amb $f(x)=x+1$, $g(x)=x$ i $g'(x)=2x$ tenim $g\sim g'$, però $(g\comp f)(0)=1\neq2=(g'\comp f)(0)$. La composició no està ben definida i falla (C2). Per això, en construir $\Pi_1(X)$, calia comprovar que l'homotopia és compatible amb la concatenació: no tot quocient d'una categoria és una categoria.
\end{exemple}

\subsection{Falsos contraexemples}

Per equilibrar, convé recordar estructures que \emph{semblen} fallar algun axioma però que són categories genuïnes. Ensenyen que els morfismes no han de ser funcions, ni estar definits a tot arreu.

\begin{exemple}[Funcions parcials]
Objectes: els conjunts; morfismes $f\colon A\rightharpoonup B$: les funcions definides en un subconjunt $D_f\subseteq A$. La composició $g\comp f$ és definida en $\{a\in D_f\mid f(a)\in D_g\}$. Les dues maneres d'agrupar $h$, $g$ i $f$ tenen el mateix domini de definició, $\{a\in D_f\mid f(a)\in D_g,\ g(f(a))\in D_h\}$, i hi prenen el mateix valor; les identitats són les funcions totals identitat. Obtenim la categoria $\cat{Par}$: que un morfisme no estigui definit a tot arreu no impedeix res.
\end{exemple}

Les relacions (la categoria $\cat{Rel}$ del capítol~\ref{cap:categories}) i les matrius ($\cat{Mat}_{\mathbb K}$, exemple~\ref{ex:mat-finvect}) són dos altres casos: els morfismes no són funcions, i en el segon cas ni tan sols són transformacions de res, sinó taules de nombres.

\subsection{Resum}

\[
\renewcommand{\arraystretch}{1.15}
\begin{array}{lccccl}
\text{Exemple} & \text{(C1)} & \text{(C2)} & \text{(C3)} & \text{(C4--5)} & \text{Remei}\\\hline
\text{camins sobre }[0,1] & \text{sí} & \text{sí} & \text{no} & \text{no} & \Pi_1(X)\text{ o camins de Moore}\\
(\mathbb Z,-) & \text{sí} & \text{sí} & \text{no} & \text{a mitges} & \text{---}\\
\text{pedra-paper-tisora} & \text{sí} & \text{sí} & \text{no} & \text{no} & \text{(el quocient és trivial)}\\
(2\mathbb Z,\cdot) & \text{sí} & \text{sí} & \text{sí} & \text{no} & \text{afegir un neutre}\\
(\mathbb R,<) & \text{sí} & \text{sí} & \text{sí} & \text{no} & (\mathbb R,\le)\\
\text{graf dirigit} & \text{sí} & \text{no} & \text{---} & \text{---} & \text{categoria lliure}\\
\text{lineals no nul·les} & \text{sí} & \text{no} & \text{sí} & \text{sí} & \text{admetre el }0\\
\text{funcions com a grafs} & \text{no} & \text{sí} & \text{sí} & \text{sí} & \text{morfismes amb extrems}
\end{array}
\]
Els axiomes de categoria no són convencions arbitràries: n'hi ha prou d'agafar l'estructura més natural possible ---els camins en un espai--- per veure'ls fallar. Quan fallen, hi ha tres actituds possibles: quocientar, estrictificar o acceptar la feblesa i fer-ne teoria.

% =====================================================================
\section{Què no és un functor}
\label{sec:no-functor}
% =====================================================================

Un functor ha de superar dues proves, en aquest ordre (observació~\ref{obs:comprovar-tipus}): els tipus i, després, les equacions $F(\id_A)=\id_{F(A)}$ i $F(g\comp f)=F(g)\comp F(f)$. Els gairebé-functors fallen en una de les dues.

\paragraph{Fallada de tipus.} El capítol~\ref{cap:functors} ja en dona l'exemple bàsic: l'antiimatge $f\mapsto f^{-1}$ no pot ser un functor covariant de $\cat{Set}$ en $\cat{Set}$, perquè va en sentit contrari; sí que ho és com a functor contravariant (exemple~\ref{ex:parts-dues-variancies}). Un altre exemple: l'assignació que envia cada grup $G$ al seu grup d'automorfismes $\mathrm{Aut}(G)$ no té cap manera natural d'actuar sobre els homomorfismes. Un homomorfisme $f\colon G\to H$ no indueix cap aplicació raonable $\mathrm{Aut}(G)\to\mathrm{Aut}(H)$; només ho fa quan $f$ és un isomorfisme, que és la manera de convertir-la en un functor sobre la subcategoria dels isomorfismes.

\paragraph{Una assignació sobre objectes que no admet cap functor.} El cas més instructiu és el d'un objecte que sembla \enquote{canònic} però no és functorial de cap manera.

\begin{proposicio}[El centre d'un grup no és functorial]\label{prop:centre-no-functor}
No hi ha cap functor $Z\colon\cat{Grp}\to\cat{Grp}$ que enviï cada grup $G$ al seu centre $Z(G)=\{z\in G\mid zg=gz \text{ per a tot } g\}$, sigui quina sigui la seva acció sobre els homomorfismes.
\end{proposicio}

\begin{prova}
Sigui $C_2=\{1,c\}$ el grup de dos elements i $S_3$ el grup de permutacions de tres elements. Considerem la inclusió $i\colon C_2\to S_3$, que envia $c$ a la transposició $(1\,2)$, i el signe $s\colon S_3\to C_2$, que envia les transposicions a $c$. Aleshores $s\comp i=\id_{C_2}$. Com que $C_2$ és commutatiu, $Z(C_2)=C_2$. En canvi, $Z(S_3)=\{1\}$: les tres transposicions no commuten entre si (per exemple, $(1\,2)(1\,3)=(1\,3\,2)$, mentre que $(1\,3)(1\,2)=(1\,2\,3)$), i les dues permutacions cícliques $(1\,2\,3)$ i $(1\,3\,2)$ no commuten amb $(1\,2)$, perquè $(1\,2)(1\,2\,3)(1\,2)=(1\,3\,2)$. Si $Z$ fos un functor,
\[
\id_{C_2}=Z(\id_{C_2})=Z(s\comp i)=Z(s)\comp Z(i),
\]
amb $Z(i)\colon C_2\to\{1\}$ i $Z(s)\colon\{1\}\to C_2$. Però una composició que passa per un grup d'un element és constant, i no pot ser la identitat d'un grup de dos elements.
\end{prova}

L'argument és un patró general: si $s\comp i=\id$, aleshores $F(s)\comp F(i)=\id$, de manera que tot functor converteix les retraccions en retraccions. Una assignació que envia un objecte a un de \enquote{massa petit} per factoritzar-hi una identitat no pot ser functorial. En canvi, assignacions semblants que només fan servir propietats preservades pels homomorfismes sí que són functors (vegeu els exercicis).

\paragraph{Fallada de les equacions.} Considerem el monoide de les funcions derivables $\mathbb R\to\mathbb R$ amb la composició, i l'assignació $f\mapsto f'$ cap al mateix monoide. Els tipus són correctes, però les equacions fallen. La identitat va a la funció constant $1$, que no és la identitat. I, amb $f(x)=g(x)=2x$, tenim $(g\comp f)'=4$ (constant), mentre que $g'\comp f'=2$ (constant). La regla de la cadena explica què cal canviar:

\begin{exemple}[La regla de la cadena és una llei functorial]\label{ex:regla-cadena}
Sigui $\cat{Der}$ la categoria que té per objectes els nombres reals i per morfismes $a\to b$ les funcions derivables $f\colon\mathbb R\to\mathbb R$ amb $f(a)=b$, amb la composició de funcions. Sigui $\cat{B}(\mathbb R,\cdot)$ el monoide multiplicatiu dels reals vist com a categoria d'un objecte. L'assignació
\[
D\colon\cat{Der}\to\cat{B}(\mathbb R,\cdot),\qquad (f\colon a\to b)\ \longmapsto\ f'(a)
\]
és un functor. Les identitats van a $1$, perquè $\id'(a)=1$, i la regla de la cadena, $(g\comp f)'(a)=g'(f(a))\cdot f'(a)$, diu exactament que $D(g\comp f)=D(g)\comp D(f)$, ja que $f(a)$ és el domini de $g$. El que fallava era el lloc on es mira la derivada: el functor no és $f\mapsto f'$, sinó $f\mapsto f'$ \emph{al punt de partida}.
\end{exemple}

\paragraph{Functors llevat d'isomorfisme (informatiu).} Hi ha assignacions que són functorials només llevat d'isomorfisme, exactament com els camins eren associatius llevat d'homotopia. L'exemple més important del llibre és la reindexació de l'observació~\ref{obs:reindexacio}. Per a cada $f\colon A\to B$, el pullback dona $f^{*}\colon\cat{C}/B\to\cat{C}/A$, però, un cop triats els pullbacks, $(g\comp f)^{*}$ i $f^{*}\comp g^{*}$ no són iguals, sinó canònicament isomorfs. A $\cat{Set}$, amb els pullbacks triats com a conjunts de parelles, $(g\comp f)^{*}$ envia $p\colon E\to C$ al conjunt de les parelles $(a,e)$ amb $g(f(a))=p(e)$, mentre que $f^{*}g^{*}$ l'envia al conjunt de les parelles $\bigl(a,(b,e)\bigr)$ amb $f(a)=b$ i $g(b)=p(e)$: són conjunts diferents, en bijecció canònica. Una tal assignació s'anomena \emph{pseudofunctor}. Al capítol~\ref{cap:subobjectes} el problema desapareix perquè $\Sub$ treballa amb classes de monomorfismes, i l'isomorfisme canònic queda absorbit (proposició~\ref{prop:functor-sub}): és el primer remei de la secció anterior, el quocient.

% =====================================================================
\section{Què no és una transformació natural}
\label{sec:no-natural}
% =====================================================================

El llibre ja dona els contraexemples essencials, i aquí només els ordenem. Una família de components $\eta_A\colon F(A)\to G(A)$ pot deixar de ser natural per tres motius diferents:
\begin{enumerate}
\item \emph{Les components s'han triat objecte per objecte}, sense mirar els morfismes. L'exemple~\ref{ex:no-natural} té totes les components bijectives i, tanmateix, un sol morfisme $\{0,1\}\to\{0,1\}$ trenca el quadrat.
\item \emph{La naturalitat només val per a una part dels morfismes.} L'isomorfisme entre un espai de dimensió finita i el seu dual depèn d'una base, i ni tan sols es pot plantejar com a transformació natural entre functors paral·lels, perquè la identitat és covariant i el dual contravariant. Només s'hi pot comparar restringint-se als isomorfismes, i fins i tot així no és natural (observació~\ref{obs:iso-no-natural} i exercici resolt~\ref{ex:dual-no-natural} de la Pràctica del capítol~\ref{cap:transformacions}). En canvi, $V\cong V^{**}$ és natural (exemple~\ref{ex:bidual}).
\item \emph{No hi ha cap component possible.} Amb el grup de dos elements actuant trivialment i per la transposició sobre $\{0,1\}$, no hi ha cap transformació natural entre els dos functors, ni tan sols una que no sigui un isomorfisme (observació~\ref{obs:iso-no-natural}).
\end{enumerate}
La lliçó comuna és la de la secció següent: tenir objectes isomorfs \emph{un per un} no és el mateix que tenir functors isomorfs.

% =====================================================================
\section{Igualtat, isomorfisme, isomorfisme natural i equivalència}
\label{sec:igualtat-equivalencia}
% =====================================================================

\subsection{La jerarquia}

Per a cada mena d'entitat hi ha una noció \emph{estricta} de \enquote{ser el mateix}, la igualtat, i una noció \emph{estructural}, que és la que respecta la teoria:
\[
\renewcommand{\arraystretch}{1.25}
\begin{array}{lll}
\text{Què comparem} & \text{Noció estricta} & \text{Noció estructural}\\\hline
\text{dos objectes d'una categoria }\cat{C} & A=B & A\cong B\ \text{(isomorfisme a }\cat{C})\\
\text{dos functors }F,G\colon\cat{C}\to\cat{D} & F=G & F\cong G\ \text{(isomorfisme natural)}\\
\text{dues categories} & \cat{C}=\cat{D} & \cat{C}\simeq\cat{D}\ \text{(equivalència)}
\end{array}
\]
Entre les categories hi ha, a més, un nivell intermedi: l'isomorfisme de categories $\cat{C}\cong\cat{D}$ (definició~\ref{def:iso-categories}). Les tres files tenen la mateixa forma. Un isomorfisme natural és un isomorfisme a la categoria de functors $\cat{D}^{\cat{C}}$, i per tant la segona fila és un cas de la primera. L'equivalència és el que s'obté quan, a la definició d'isomorfisme de categories, se substitueixen les igualtats $G\comp F=\id$ i $F\comp G=\id$ per isomorfismes naturals $G\comp F\cong\id$ i $F\comp G\cong\id$, perquè la manera correcta de comparar functors és l'isomorfisme natural. Cada nivell fa servir la noció estructural del nivell anterior.

Totes les implicacions possibles entre aquestes nocions van en un sol sentit:
\[
\begin{gathered}
A=B\ \Rightarrow\ A\cong B,\\
F=G\ \Rightarrow\ F\cong G\ \Rightarrow\ F(A)\cong G(A)\ \text{per a tot }A,\\
\cat{C}=\cat{D}\ \Rightarrow\ \cat{C}\cong\cat{D}\ \Rightarrow\ \cat{C}\simeq\cat{D},
\end{gathered}
\]
i cap no és reversible. Vegem-ne un contraexemple per a cada pas.

\subsection{Un contraexemple per a cada pas}

\begin{itemize}
\item \emph{Isomorfs però no iguals.} A $\cat{Set}$, $\{0,1\}\cong\{a,b\}$. A $\cat{Prov}_T$, $p\land q\cong q\land p$, que són fórmules diferents (exemple~\ref{ex:iso-equivalencia-logica}).
\item \emph{Naturalment isomorfs però no iguals.} A la categoria dels espais vectorials de dimensió finita, la identitat i el bidual són naturalment isomorfs (exemple~\ref{ex:bidual}), però no són iguals: $V^{**}$ no és $V$, sinó un espai format per funcionals sobre $V^{*}$.
\item \emph{Isomorfs objecte per objecte, però no naturalment isomorfs.} Els dos functors del grup de dos elements en $\cat{Set}$ de l'observació~\ref{obs:iso-no-natural} envien l'únic objecte al mateix conjunt $\{0,1\}$, però no són naturalment isomorfs.
\item \emph{Isomorfes però no iguals.} Una categoria i una còpia seva amb els objectes reanomenats.
\item \emph{Equivalents però no isomorfes.} $\cat{Mat}_{\mathbb K}\simeq\cat{FinVect}_{\mathbb K}$ (exemple~\ref{ex:mat-finvect}), o la categoria $\cat{1}$ i la categoria amb dos objectes isomorfs i un sol morfisme entre cada parell d'objectes.
\end{itemize}

\subsection{La invariància estructural}

El principi que governa tot el llibre és aquest: \emph{una propietat formulada només amb l'estructura ---composició, identitats, propietats universals---, sense afirmar mai que dos objectes són iguals, es transporta per la noció estructural corresponent.} Per als objectes ho diu l'observació~\ref{obs:invariancia}. Per a les categories, la peça clau és que les equivalències conserven les propietats universals.

\begin{teorema}\label{teo:equivalencia-limits}
Sigui $F\colon\cat{C}\to\cat{D}$ una equivalència. Si $(L,\lambda)$ és un límit d'un diagrama $D\colon\cat{I}\to\cat{C}$, aleshores $(F(L),F\lambda)$ és un límit de $F\comp D$. Dualment, $F$ preserva els colímits.
\end{teorema}

\begin{prova}
Pel teorema~\ref{teo:caract-equivalencia}, $F$ és ple, fidel i essencialment exhaustiu. La família $F\lambda=(F(\lambda_i))_i$ és un con sobre $F\comp D$, perquè $F$ preserva la composició. Sigui $(X,\alpha)$ un con sobre $F\comp D$. Com que $F$ és essencialment exhaustiu, hi ha un objecte $C$ de $\cat{C}$ i un isomorfisme $e\colon F(C)\to X$. Com que $F$ és ple i fidel, per a cada $i$ hi ha un únic $\beta_i\colon C\to D_i$ amb $F(\beta_i)=\alpha_i\comp e$. Aquests $\beta_i$ formen un con: per a $u\colon i\to j$,
\[
F(D_u\comp\beta_i)=F(D_u)\comp\alpha_i\comp e=\alpha_j\comp e=F(\beta_j),
\]
i, com que $F$ és fidel, $D_u\comp\beta_i=\beta_j$. Per la propietat del límit, hi ha un únic $h\colon C\to L$ amb $\lambda_i\comp h=\beta_i$. Posem $k=F(h)\comp e^{-1}\colon X\to F(L)$. Aleshores $F(\lambda_i)\comp k=F(\beta_i)\comp e^{-1}=\alpha_i$: $k$ factoritza el con.

Per a la unicitat, sigui $k'\colon X\to F(L)$ amb $F(\lambda_i)\comp k'=\alpha_i$ per a tot $i$. Com que $F$ és ple, $k'\comp e=F(h')$ per a algun $h'\colon C\to L$, i aleshores $F(\lambda_i\comp h')=\alpha_i\comp e=F(\beta_i)$. Com que $F$ és fidel, $\lambda_i\comp h'=\beta_i$, i per la unicitat del límit, $h'=h$. Per tant, $k'=F(h)\comp e^{-1}=k$. Els colímits es tracten igual, invertint totes les fletxes.
\end{prova}

D'aquest teorema, de l'exercici resolt~\ref{pr-equivalencia-preserva} de la Pràctica del capítol~\ref{cap:transformacions} (que una equivalència preserva i reflecteix isomorfismes, monomorfismes i epimorfismes) i del fet que el quasiinvers també és una equivalència, se segueix que les propietats de la columna esquerra de la taula següent les té $\cat{C}$ si i només si les té qualsevol categoria equivalent a $\cat{C}$. Per a les dues últimes (cartesiana tancada, topos), el mateix argument s'aplica a les propietats universals de l'exponencial i del classificador; no el reproduïm.
\begin{center}
\small
\renewcommand{\arraystretch}{1.2}
\begin{tabularx}{\linewidth}{@{}>{\raggedright\arraybackslash}X>{\raggedright\arraybackslash}X@{}}
\toprule
\textbf{Invariants per equivalència} & \textbf{No invariants per equivalència}\\
\midrule
tenir objecte terminal o inicial & el nombre d'objectes\\
tenir productes, equalitzadors, límits finits & tenir un sol objecte\\
ser un preordre (com a molt un morfisme entre dos objectes) & ser un ordre parcial\\
ser un grupoide (tot morfisme és invertible) & ser esquelètica\\
que tot morfisme mono i epi sigui un isomorfisme & que dos objectes concrets siguin iguals\\
ser cartesiana tancada; ser un topos & ser petita (\enquote{essencialment petita} sí que és invariant)\\
\bottomrule
\end{tabularx}
\end{center}
Les propietats de la columna de la dreta fan servir, d'una manera o altra, la igualtat d'objectes, i per això no són categòriques. Per exemple, $\cat{1}$ és un ordre parcial, i la categoria amb dos objectes isomorfs, que li és equivalent, no ho és. Quan una definició depèn de la igualtat d'objectes, convé sospitar-ne.

\subsection{Equivalència i esquelets}

El resultat següent fa veure que l'equivalència és exactament l'isomorfisme \enquote{llevat de les còpies isomorfes}.

\begin{lema}\label{lema:esqueletiques}
Una equivalència entre dues categories esquelètiques és un isomorfisme de categories.
\end{lema}

\begin{prova}
Sigui $F\colon\cat{C}\to\cat{D}$ una equivalència entre categories esquelètiques; és ple, fidel i essencialment exhaustiu. \emph{És injectiu sobre objectes}: si $F(A)=F(B)$, la identitat de $F(A)$ és $F(u)$ per a algun $u\colon A\to B$, perquè $F$ és ple, i $u$ és un isomorfisme, perquè un functor ple i fidel reflecteix els isomorfismes (proposició~\ref{prop:plenfidel-reflecteix-iso}). Com que $\cat{C}$ és esquelètica, $A=B$. \emph{És exhaustiu sobre objectes}: per a cada $D$ hi ha $C$ amb $D\cong F(C)$, i com que $\cat{D}$ és esquelètica, $D=F(C)$. A més, $F$ és bijectiu sobre cada homconjunt, perquè és ple i fidel. Per tant, $F$ és bijectiu sobre objectes i sobre morfismes, i l'assignació inversa és un functor (preserva identitats i composicions perquè $F$ les preserva). És, doncs, un isomorfisme de categories.
\end{prova}

\begin{corollari}\label{cor:equivalencia-esquelets}
Dues categories són equivalents si i només si tenen esquelets isomorfs.
\end{corollari}

\begin{prova}
En el marc d'elecció del llibre, tota categoria té esquelets (definició~\ref{def:esquelet}), i és equivalent a qualsevol d'ells: la inclusió de l'esquelet és plena i fidel per ser una subcategoria plena, i essencialment exhaustiva perquè l'esquelet conté un objecte de cada classe d'isomorfisme. Les equivalències es componen, perquè la composició de functors plens, fidels i essencialment exhaustius ho torna a ser. Si $\cat{C}\simeq\cat{D}$, els seus esquelets són equivalents, i pel lema~\ref{lema:esqueletiques} són isomorfs. Recíprocament, si els esquelets són isomorfs, $\cat{C}\simeq\mathrm{esquelet}(\cat{C})\cong\mathrm{esquelet}(\cat{D})\simeq\cat{D}$.
\end{prova}

Per als preordres, l'esquelet és l'ordre parcial que s'obté identificant els elements equivalents ($x\le y$ i $y\le x$): dos preordres són equivalents si i només si aquests ordres parcials són isomorfs.

\subsection{Com es demostra que dues categories no són equivalents}

Per demostrar que $\cat{C}\not\simeq\cat{D}$ no es poden revisar tots els functors possibles. La via és buscar una propietat invariant per equivalència que l'una tingui i l'altra no.

\begin{exemple}[$\cat{Set}\not\simeq\cat{Set}\op$]
Diem que un objecte inicial $0$ és \emph{estricte} si tot morfisme $X\to0$ és un isomorfisme. Tenir un objecte inicial estricte és invariant per equivalència. Si $F\colon\cat{C}\to\cat{D}$ és una equivalència, $F(0)$ és inicial (teorema~\ref{teo:equivalencia-limits}). Donat un morfisme $k\colon Y\to F(0)$, prenem un isomorfisme $e\colon F(C)\to Y$; com que $F$ és ple, $k\comp e=F(u)$ per a algun $u\colon C\to0$, que és un isomorfisme, i aleshores $F(u)$ i $k=F(u)\comp e^{-1}$ també ho són. A $\cat{Set}$, l'objecte inicial $\varnothing$ és estricte: només hi ha aplicacions $X\to\varnothing$ si $X=\varnothing$, i aleshores és la identitat. A $\cat{Set}\op$, l'objecte inicial és un conjunt d'un element $\{\ast\}$, i un morfisme $X\to\{\ast\}$ de $\cat{Set}\op$ és una aplicació $\{\ast\}\to X$ de $\cat{Set}$. Amb $X=\{0,1\}$ no és bijectiva, i per tant no és un isomorfisme. Així, $\cat{Set}$ no és equivalent a la seva oposada.
\end{exemple}

\begin{exemple}[$\cat{Grp}\not\simeq\cat{Ab}$]
La propietat \enquote{per a dos objectes qualssevol, el producte i el coproducte són isomorfs} és invariant per equivalència, perquè una equivalència preserva productes, coproductes i isomorfismes. 
\emph{A $\cat{Ab}$ es compleix.} Siguin $A$ i $B$ grups abelians, escrits additivament. El producte $A\times B$, amb les inclusions $a\mapsto(a,0)$ i $b\mapsto(0,b)$, és també un coproducte. Donats homomorfismes $f\colon A\to G$ i $g\colon B\to G$ cap a un grup abelià $G$, l'aplicació $(a,b)\mapsto f(a)+g(b)$ és un homomorfisme, precisament perquè $G$ és commutatiu, i restringida a les inclusions dona $f$ i $g$. És l'única, perquè $(a,b)=(a,0)+(0,b)$.

\emph{A $\cat{Grp}$ no es compleix.} Sigui $C_2=\{1,c\}$ i suposem que el coproducte $P=C_2+C_2$ és isomorf al producte $C_2\times C_2$, que és commutatiu. Els homomorfismes $C_2\to S_3$ que envien $c$ a $(1\,2)$ i a $(1\,3)$ indueixen un homomorfisme $P\to S_3$ la imatge del qual conté $(1\,2)$ i $(1\,3)$. Però la imatge d'un grup commutatiu és commutativa, i $(1\,2)$ i $(1\,3)$ no commuten. Per tant, $C_2+C_2\not\cong C_2\times C_2$.
\end{exemple}

\begin{exemple}[$\cat{FinSet}\not\simeq\cat{Set}$]
\enquote{Tota família numerable d'objectes té coproducte} és invariant per equivalència. $\cat{Set}$ la compleix i $\cat{FinSet}$ no (observació~\ref{obs:elemental-prefeixos}).
\end{exemple}

\subsection{Què té a veure tot això amb la lògica}

A la categoria $\cat{Prov}_T$ d'un sistema deductiu, els tres nivells tenen una lectura lògica precisa:
\begin{itemize}
\item la \emph{igualtat} de dues fórmules és la igualtat sintàctica: $p\land q$ i $q\land p$ són fórmules diferents;
\item l'\emph{isomorfisme} és la interdeduïbilitat, $p\vdash_T q$ i $q\vdash_T p$ (exemple~\ref{ex:iso-equivalencia-logica});
\item l'\emph{esquelet} de $\cat{Prov}_T$ és isomorf a l'ordre parcial de les classes de fórmules interdeduïbles, que en lògica s'anomena \textbf{àlgebra de Lindenbaum--Tarski}\index{àlgebra de Lindenbaum--Tarski} de $T$. Pel corol·lari~\ref{cor:equivalencia-esquelets}, $\cat{Prov}_T$ és equivalent, però en general no isomorf, a aquesta àlgebra.
\end{itemize}
En conseqüència, qualsevol propietat categòrica de $\cat{Prov}_T$ és una propietat de l'àlgebra de Lindenbaum--Tarski: tenir productes (la conjunció), ser cartesiana tancada (la implicació), ser una àlgebra de Heyting o de Boole. La lògica estudia les teories \emph{llevat d'interdeduïbilitat}, i la teoria de categories hi aporta el llenguatge que fa aquesta actitud sistemàtica: el que compta d'una teoria és la seva estructura, no la tria particular de fórmules.

\subsection{Errors típics}

\begin{itemize}
\item Escriure $A=B$ quan només es té $A\cong B$, encara que l'isomorfisme sigui canònic.
\item Deduir $F\cong G$ de $F(A)\cong G(A)$ per a cada $A$, sense comprovar la naturalitat.
\item Creure que una equivalència és un isomorfisme de categories, o que $G\comp F$ ha de ser literalment la identitat.
\item Demostrar que dues categories no són equivalents amb un argument que fa servir la igualtat d'objectes, com comptar-ne el nombre.
\item Definir una propietat d'una categoria que depèn d'objectes concrets, i després aplicar-la a una categoria equivalent.
\end{itemize}

% =====================================================================
\section{La lògica fora dels topos}
\label{sec:fora-topos}
% =====================================================================

Al llarg dels capítols~\ref{cap:cartesianes}--\ref{cap:pont}, la lògica s'ha llegit d'una manera determinada, que comença amb el tancament cartesià: una proposició sobre $E$ és un subobjecte de $E$, la conjunció és un pullback, la implicació és l'exponencial (l'adjunt per la dreta de $-\wedge U$), els quantificadors són adjunts de la reindexació, i en un topos tot això és la lògica interna de $\Omega$. Hi ha lògiques i models perfectament legítims que no encaixen en aquest esquema. Els presentem de manera sobretot informativa, però en cada cas diem amb precisió quina condició falla i, quan és fàcil, ho comprovem.

\subsection{Categories que no són cartesianes tancades}
\label{subsec:no-ccc}

Abans de mirar què hi ha entre les categories cartesianes tancades i els topos, convé veure que moltes categories molt naturals no són ni tan sols cartesianes tancades. Per demostrar-ho, n'hi ha prou de trobar una propietat que tota categoria cartesiana tancada tingui i que la categoria en qüestió no tingui. La més útil és la que dona el teorema de preservació de colímits pels adjunts per l'esquerra (capítol~\ref{cap:adjuncions}): en una categoria cartesiana tancada, $-\times B$ és adjunt per l'esquerra de $(-)^B$, i per tant preserva tots els colímits que existeixin. En particular, si $0$ és inicial, $0\times B$ també ho és.

Un \textbf{objecte zero}\index{objecte zero} és un objecte que és alhora inicial i terminal.

\begin{proposicio}\label{prop:zero-no-ccc}
Si una categoria cartesiana tancada té un objecte zero $0$, aleshores tots els seus objectes són isomorfs a $0$.
\end{proposicio}

\begin{prova}
Sigui $B$ un objecte. Com que $0$ és inicial i $-\times B$ preserva l'objecte inicial, $0\times B$ és inicial, i per tant $0\times B\cong0$. D'altra banda, com que $0$ és terminal, la projecció $\pi_2\colon0\times B\to B$ és un isomorfisme. En efecte, sigui $!_B\colon B\to0$ l'únic morfisme cap al terminal, i $s=\langle !_B,\id_B\rangle\colon B\to0\times B$. Aleshores $\pi_2\comp s=\id_B$, i
\[
s\comp\pi_2=\langle !_B\comp\pi_2,\ \pi_2\rangle=\langle\pi_1,\pi_2\rangle=\id_{0\times B},
\]
perquè $!_B\comp\pi_2$ i $\pi_1$ són dos morfismes $0\times B\to0$, que han de coincidir perquè $0$ és terminal. Per tant, $B\cong0\times B\cong0$.
\end{prova}

\begin{corollari}\label{cor:zero-no-ccc}
Les categories $\cat{Vect}_{\mathbb K}$, $\cat{Ab}$, $\cat{Grp}$, $\cat{Mon}$ i $\cat{Rel}$ no són cartesianes tancades.
\end{corollari}

\begin{prova}
Totes tenen un objecte zero: l'espai $\{0\}$; el grup, o el monoide, d'un sol element, perquè de l'element neutre surt un únic homomorfisme cap a qualsevol altre grup o monoide i cap a l'estructura trivial només n'hi ha un; i, a $\cat{Rel}$, el conjunt buit, perquè l'única relació $\varnothing\to X$ i l'única relació $X\to\varnothing$ són la relació buida. Totes tenen, però, objectes que no són isomorfs al zero: $\mathbb K$, el grup $\mathbb Z$, el monoide $(\mathbb N,+)$ o qualsevol conjunt no buit. Per la proposició~\ref{prop:zero-no-ccc}, cap no és cartesiana tancada.
\end{prova}

\paragraph{Preordres.} En un preordre amb ínfims finits, ser cartesià tancat vol dir tenir una implicació $b\Rightarrow c$ amb $a\wedge b\le c$ si i només si $a\le b\Rightarrow c$ (capítol~\ref{cap:cartesianes}). Si, a més, té suprems binaris, aleshores és distributiu, perquè $-\wedge b$ és adjunt per l'esquerra i preserva els suprems. Per tant, \emph{un reticle no distributiu no és cartesià tancat}. El reticle dels subespais de $\mathbb R^2$ que veurem tot seguit n'és un exemple.

\paragraph{Categories que ni tan sols són cartesianes.} Un monoide $M$ no trivial, vist com a categoria $\cat{B}M$, no té objecte terminal: l'únic objecte $\ast$ ho seria només si $\Hom(\ast,\ast)=M$ tingués un sol element. Un preordre amb dos elements incomparables que no tinguin cap fita inferior comuna no té el producte d'aquests dos elements. Cap d'aquestes categories no pot ser cartesiana tancada, perquè no té productes finits.

\paragraph{Els espais topològics (informatiu).} $\cat{Top}$ tampoc no és cartesiana tancada: hi ha quocients d'espais que el producte amb un espai fix no conserva. La demostració demana més topologia de la que pressuposa el llibre, i no la fem.

\subsection{Cartesianes tancades que no són topos}

La lògica de la implicació necessita només el tancament cartesià, i això és molt menys que un topos.

\begin{exemple}[$\cat{Cat}$ no és un topos]
La categoria $\cat{Cat}$ és cartesiana tancada (exemple~\ref{ex:cat-ccc}), però no té classificador de subobjectes. Sigui $\cat{2}$ la categoria $0\to1$, i $\cat{I}$ la categoria amb dos objectes i, a més de les identitats, dos morfismes $u\colon0\to1$ i $u^{-1}\colon1\to0$ inversos l'un de l'altre. La inclusió $i\colon\cat{2}\to\cat{I}$ és:
\begin{itemize}
\item \emph{mono}: és injectiva sobre objectes i morfismes, de manera que si $i\comp H=i\comp K$, aleshores $H=K$;
\item \emph{epi}: si $H,K\colon\cat{I}\to\cat{C}$ coincideixen sobre $u$, aleshores també coincideixen sobre $u^{-1}$, perquè un functor envia l'invers de $u$ a l'invers de la imatge de $u$, i l'invers és únic;
\item \emph{no és un isomorfisme}: $u^{-1}$ no és a la imatge.
\end{itemize}
En una categoria amb classificador, tot morfisme mono i epi és un isomorfisme (exercici resolt~\ref{pr-mono-equalitzador} de la Pràctica del capítol~\ref{cap:topos}). Per tant, $\cat{Cat}$ no en té cap.
\end{exemple}

L'ordre parcial $\mathbb Z\cup\{+\infty\}$ de l'observació~\ref{obs:ccc-no-heyting} és un altre cas: té conjunció i implicació, però no té falsedat.

\subsection{Lògiques on falla la distributivitat}

En una àlgebra de Heyting, $a\wedge(b\vee c)=(a\wedge b)\vee(a\wedge c)$, perquè $a\wedge-$ és adjunt per l'esquerra i preserva els suprems (capítol~\ref{cap:cartesianes}). Com que els subobjectes d'un objecte d'un topos formen una àlgebra de Heyting (nota~\ref{nota:sub-heyting}), cap reticle no distributiu pot ser-ne un.

\begin{exemple}[La lògica quàntica]
Siguin els subespais vectorials de $\mathbb R^2$, ordenats per inclusió: $\{0\}$, les rectes que passen per l'origen i $\mathbb R^2$. L'ínfim és la intersecció i el suprem és la suma. Si $a$, $b$ i $c$ són tres rectes diferents,
\[
a\wedge(b\vee c)=a\wedge\mathbb R^2=a,
\qquad
(a\wedge b)\vee(a\wedge c)=\{0\}\vee\{0\}=\{0\}.
\]
La distributivitat falla. En mecànica quàntica, les propietats observables d'un sistema es modelen amb reticles d'aquest tipus (de subespais tancats), i la \enquote{lògica quàntica} de Birkhoff i von Neumann és la lògica d'aquests reticles no distributius (informatiu). Amb aquestes operacions, aquest reticle no pot ser el reticle de subobjectes de cap objecte d'un topos.
\end{exemple}

\subsection{Lògiques amb una altra conjunció}

\begin{exemple}[La lògica de Łukasiewicz]
Els valors de veritat són els nombres de $[0,1]$, amb la conjunció forta i la implicació
\[
x\otimes y=\max(0,\,x+y-1),\qquad y\to z=\min(1,\,1-y+z).
\]
La implicació és adjunta d'aquesta conjunció, i no de l'ínfim. En efecte, si $z\ge0$, $\max(0,x+y-1)\le z$ si i només si $x\le1-y+z$, és a dir, $x\le y\to z$. En canvi, amb l'ínfim falla: per a $y=0.5$, $z=0.4$ i $x=0.6$, tenim $\min(x,y)=0.5>0.4$, però $x=0.6\le0.9=y\to z$. A més, $\otimes$ no és idempotent: $0.5\otimes0.5=0\neq0.5$. Lògicament, això vol dir que de $A$ no se'n dedueix $A\otimes A$: una hipòtesi no es pot fer servir dues vegades. Categòricament, $([0,1],\otimes,\to)$ és un preordre \emph{monoidal tancat} que no és cartesià tancat respecte de $\otimes$. Aquesta idea ---hipòtesis que són recursos i no es poden duplicar--- és la de la \emph{lògica lineal} de Girard; els seus fragments més simples es modelen en categories monoidals tancades, i la resta demana estructura addicional (informatiu).
\end{exemple}

\subsection{Lògiques on la contradicció no és el fals}

\begin{exemple}[Els tancats d'un espai]
A $\mathbb R$, els subconjunts tancats, ordenats per inclusió, formen un reticle distributiu on el \enquote{complement} natural d'un tancat $A$ és l'adherència del seu complementari, $\neg A=\overline{\mathbb R\setminus A}$. Per a $A=[0,1]$,
\[
A\wedge\neg A=[0,1]\cap\bigl((-\infty,0]\cup[1,\infty)\bigr)=\{0,1\}\neq\varnothing:
\]
una proposició i la seva negació són certes alhora a la frontera. Aquesta estructura és la duala d'una àlgebra de Heyting (una àlgebra de \emph{co-Heyting}), i serveix de model a lògiques \emph{paraconsistents}, on una contradicció no implica qualsevol cosa. Amb aquestes operacions, no pot ser el reticle de subobjectes de cap objecte d'un topos, perquè hi falla $A\wedge\neg A=\bot$.
\end{exemple}

\subsection{Lògiques amb modalitats}

Una \emph{lògica modal} afegeix operadors com \enquote{necessàriament} ($\Box$) i \enquote{possiblement} ($\Diamond$). En un model de Kripke, hi ha un conjunt de mons $W$ i una relació d'accessibilitat $R\subseteq W\times W$, i per a $A\subseteq W$ es defineix
\[
\Box A=\{x\mid \text{per a tot } y,\ xRy\Rightarrow y\in A\},
\qquad
\Diamond A=\{x\mid \text{hi ha } y \text{ amb } xRy \text{ i } y\in A\}.
\]
Aquestes modalitats són quantificadors \emph{al llarg de la relació} $R$, com els $\forall_f$ i $\exists_f$ del capítol~\ref{cap:subobjectes} ho eren al llarg d'una aplicació. Per veure-ho, sigui $\exists_R(B)=\{y\mid \text{hi ha } x\in B \text{ amb } xRy\}$, el conjunt dels mons accessibles des de $B$. Aleshores, per a $A,B\subseteq W$,
\[
\exists_R(B)\subseteq A\iff B\subseteq\Box A,
\]
és a dir, $\exists_R\dashv\Box$ entre els ordres parcials $(\mathcal P(W),\subseteq)$. En efecte, totes dues condicions diuen el mateix: per a tot $x\in B$ i tot $y$ amb $xRy$, es té $y\in A$. La modalitat $\Diamond$ és el mateix tipus de quantificador, però al llarg de la relació inversa: $\Diamond A=\{x\mid\text{hi ha } y\in A \text{ amb } yR^{-1}x\}$. Les modalitats no són, però, connectius de la lògica interna d'un topos, sinó estructura afegida (informatiu).

\subsection{Resum}

\begin{center}
\small
\renewcommand{\arraystretch}{1.2}
\begin{tabularx}{\linewidth}{@{}>{\raggedright\arraybackslash}p{0.24\linewidth}>{\raggedright\arraybackslash}X>{\raggedright\arraybackslash}X@{}}
\toprule
\textbf{Model} & \textbf{Estructura} & \textbf{Què falla respecte d'un topos}\\
\midrule
$\cat{Vect}_{\mathbb K}$, $\cat{Ab}$, $\cat{Grp}$, $\cat{Rel}$ & amb objecte zero & no són cartesianes tancades\\
$\cat{Cat}$ & cartesiana tancada & no hi ha classificador\\
$\mathbb Z\cup\{+\infty\}$ & preordre cartesià tancat & no hi ha fals\\
subespais de $\mathbb R^2$ & reticle & la distributivitat, i per tant el tancament cartesià\\
Łukasiewicz & preordre monoidal tancat & la implicació no és adjunta de $\wedge$\\
tancats de $\mathbb R$ & co-Heyting & $A\wedge\neg A\neq\bot$\\
models de Kripke & àlgebra de Boole amb modalitats & estructura afegida\\
\bottomrule
\end{tabularx}
\end{center}
Cada fila marca un camí que la lògica categòrica pot seguir més enllà dels topos: les estructures monoidals tancades per a les lògiques de recursos, els reticles no distributius per a la lògica quàntica, els operadors modals com a adjunts al llarg de relacions. El llibre s'ha centrat en l'esquema dels topos perquè és el que unifica més bé la lògica intuïcionista i la teoria de conjunts. Aquestes alternatives mostren que és una tria, no l'únic camí possible.

% =====================================================================
\section{Exercicis}
% =====================================================================

\begin{exercici}
Demostra que $c_y\ast\gamma=\gamma$ si i només si $\gamma$ és constant. \hfill{\normalfont$\star$}
\begin{indicacio}
La igualtat a $[\frac12,1]$ força $\gamma=y$ en aquest interval, i a $[0,\frac12]$ diu $\gamma(t)=\gamma(2t)$. Itera-ho i fes servir la continuïtat a $0$.
\end{indicacio}
\end{exercici}

\begin{exercici}
Comprova que, a la categoria dels camins de Moore, un camí $(r,\gamma)$ amb $r>0$ no té cap invers: no hi ha cap $(s,\delta)$ amb $(s,\delta)\comp(r,\gamma)=(0,c_x)$. \hfill{\normalfont$\star$}
\begin{indicacio}
Mira la durada de la composició.
\end{indicacio}
\end{exercici}

\begin{exercici}
Troba un magma de dos elements que no sigui associatiu, i un altre que sigui associatiu però no tingui neutre. \hfill{\normalfont$\star$}
\begin{indicacio}
Hi ha només setze operacions possibles en un conjunt de dos elements. Per al segon, prova amb una operació que no depengui del segon argument.
\end{indicacio}
\end{exercici}

\begin{exercici}
Considera els nombres reals com a objectes, i com a morfismes $x\to y$ els reals $d\ge|y-x|$, amb la composició $d'\comp d=d+d'$. És una categoria? Què canvia si s'exigeix $d=|y-x|$? \hfill{\normalfont$\star\star$}
\begin{indicacio}
Per al primer cas, fes servir la desigualtat triangular. Per al segon, mira què passa amb $x<y>z$.
\end{indicacio}
\end{exercici}

\begin{exercici}
Sigui $T(G)=\{g\in G\mid g^2=1\}$. Comprova que $G\mapsto T(G)$, amb $f\mapsto f|_{T(G)}$, és un functor $\cat{Grp}\to\cat{Set}$, i explica per què l'argument de la proposició~\ref{prop:centre-no-functor} no s'hi aplica. \hfill{\normalfont$\star$}
\begin{indicacio}
Mira primer els tipus: si $g^2=1$, què es pot dir de $f(g)^2$? Per a la segona part, calcula $T(S_3)$.
\end{indicacio}
\end{exercici}

\begin{exercici}
Demostra que un functor bijectiu sobre objectes i sobre morfismes és un isomorfisme de categories. \hfill{\normalfont$\star$}
\begin{indicacio}
Defineix l'invers sobre objectes i morfismes, i comprova que preserva la composició aplicant-hi el functor original.
\end{indicacio}
\end{exercici}

\begin{exercici}
Demostra que \enquote{ser un grupoide} és invariant per equivalència, i que \enquote{ser un ordre parcial} no ho és. \hfill{\normalfont$\star$}
\begin{indicacio}
Per a la primera part, fes servir que una equivalència és plena i fidel i que reflecteix els isomorfismes.
\end{indicacio}
\end{exercici}

\begin{exercici}
La categoria $\cat{Rel}$ és isomorfa a la seva oposada, enviant cada relació $R$ a la relació inversa. Fes-ho servir, junt amb l'exemple $\cat{Set}\not\simeq\cat{Set}\op$, per demostrar que $\cat{Set}\not\simeq\cat{Rel}$. \hfill{\normalfont$\star\star$}
\begin{indicacio}
Si $F\colon\cat{C}\to\cat{D}$ és una equivalència, $F\op\colon\cat{C}\op\to\cat{D}\op$ també ho és.
\end{indicacio}
\end{exercici}

\begin{exercici}
Sigui $\cat{Set}_\ast$ la categoria dels conjunts amb un punt distingit, amb les aplicacions que preserven el punt distingit. Comprova que té un objecte zero i dedueix que no és cartesiana tancada. \hfill{\normalfont$\star$}
\begin{indicacio}
Pensa en el conjunt d'un sol element. Per a la segona part, aplica la proposició~\ref{prop:zero-no-ccc}.
\end{indicacio}
\end{exercici}

\begin{exercici}
Comprova la fallada de la distributivitat al reticle dels subespais de $\mathbb R^2$ amb les rectes $y=0$, $x=0$ i $y=x$. Té aquest reticle un \enquote{complement} per a cada recta? \hfill{\normalfont$\star$}
\begin{indicacio}
Un complement de $a$ és un $b$ amb $a\wedge b=\{0\}$ i $a\vee b=\mathbb R^2$. És únic?
\end{indicacio}
\end{exercici}

\begin{exercici}
A $\mathbb R$, calcula $A\wedge\neg A$ per al tancat $A=[0,1]\cup\{2\}$, amb $\neg A=\overline{\mathbb R\setminus A}$. \hfill{\normalfont$\star$}
\begin{indicacio}
El resultat és la frontera de $A$.
\end{indicacio}
\end{exercici}

\fi

\backmatter
\chapter*{Bibliografia}
\addcontentsline{toc}{chapter}{Bibliografia}
\markboth{Bibliografia}{Bibliografia}

Les referències següents cobreixen tres funcions: introduccions al mateix nivell que aquest text, obres de consulta i fonts per a la transició cap a la lògica categòrica, i les fonts de la terminologia matemàtica catalana emprada.

\begin{thebibliography}{99}

\bibitem{eilenbergmaclane}
S.~Eilenberg i S.~Mac~Lane,
\emph{General theory of natural equivalences},
Transactions of the American Mathematical Society~\textbf{58} (1945), 231--294.
\newblock (Article fundacional de la teoria de categories, on es defineixen per primera vegada les nocions de categoria, functor i transformació natural. \url{https://doi.org/10.1090/S0002-9947-1945-0013131-6}.)

\bibitem{awodey}
S.~Awodey,
\emph{Category Theory},
2a edició, Oxford Logic Guides~52, Oxford University Press, 2010.
\newblock (Text introductori de referència; guia principal per a l'ordre i el nivell d'aquest llibre, i per als capítols sobre categories cartesianes tancades i lògica. Les remissions del llibre es refereixen a aquesta edició.)

\bibitem{maclane}
S.~Mac~Lane,
\emph{Categories for the Working Mathematician},
2a edició, Graduate Texts in Mathematics~5, Springer, 1998. \url{https://doi.org/10.1007/978-1-4757-4721-8}.
\newblock (Obra de consulta clàssica sobre teoria de categories. Les remissions del llibre es refereixen a aquesta edició.)

\bibitem{leinster}
T.~Leinster,
\emph{Basic Category Theory},
Cambridge Studies in Advanced Mathematics~143, Cambridge University Press, 2014. Versió en accés obert: \url{https://arxiv.org/abs/1612.09375}.
\newblock (Introducció alternativa, breu i moderna. La versió oberta es va publicar a arXiv el desembre de 2016 per acord amb Cambridge University Press; incorpora les correccions de la llista d'errates de l'edició impresa i en conserva la numeració de capítols i seccions, que és la que segueixen les remissions d'aquest llibre. Consulta: octubre de 2026.)

\bibitem{riehl}
E.~Riehl,
\emph{Category Theory in Context},
Aurora: Dover Modern Math Originals, Dover, 2016.
\newblock (Consulta complementària, amb molts exemples. Les remissions d'aquest llibre, per capítols i seccions, es refereixen a l'edició de 2016. La pàgina de l'autora ofereix una versió en accés obert, \url{https://emilyriehl.github.io/files/context.pdf}, que s'actualitza i pot no coincidir amb l'edició de 2016 en la paginació. La pàgina de l'autora distingeix actualment la primera edició de la versió en preparació de la segona edició; les remissions d'aquest llibre es refereixen a la primera edició de 2016. Consulta: octubre de 2026.)

\bibitem{barrwells}
M.~Barr i C.~Wells,
\emph{Category Theory for Computing Science},
3a edició, Les Publications CRM, Montréal, 1999. Reimprès a \emph{Reprints in Theory and Applications of Categories}~22 (2012). \url{https://www.tac.mta.ca/tac/reprints/articles/22/tr22abs.html}.
\newblock (Introducció orientada a la informàtica i la lògica, amb sistemes deductius i llenguatges funcionals com a exemples de categories. El portal de TAC n'ofereix també versions corregides posteriors a la reimpressió de 2012; les remissions d'aquest llibre es refereixen a la reimpressió de 2012; per a pàgines exactes cal indicar la versió consultada. Consulta: octubre de 2026.)

\bibitem{perrone}
P.~Perrone,
\emph{Starting Category Theory},
World Scientific, Singapur, 2024. \url{https://doi.org/10.1142/13670}.
\newblock (Introducció detallada i accessible, amb molts exemples de matemàtica bàsica.)

\bibitem{perrone-notes}
P.~Perrone,
\emph{Notes on Category Theory with examples from basic mathematics},
notes en accés obert, arXiv:1912.10642, versió~7 (abril de 2024; text no revisat, darrera actualització de febrer de 2021), \url{https://arxiv.org/abs/1912.10642}.
\newblock (Material previ al llibre \cite{perrone}, que el va ampliar i reorganitzar; les remissions a aquestes notes no es corresponen necessàriament amb els capítols del llibre. Consulta: octubre de 2026.)

\bibitem{lawvereschanuel}
F.~W.~Lawvere i S.~H.~Schanuel,
\emph{Conceptual Mathematics: A First Introduction to Categories},
2a edició, Cambridge University Press, 2009.
\newblock (Introducció molt elemental, amb conjunts, grafs i sistemes dinàmics com a exemples recurrents. Les remissions es fan per parts, articles i sessions, i s'han comprovat amb la traducció castellana de la segona edició, \emph{Matemáticas conceptuales. Una primera introducción a categorías}, de F.~Marmolejo, que en conserva la numeració.)

\bibitem{streicher}
T.~Streicher,
\emph{Introduction to Category Theory and Categorical Logic},
notes de curs, Technische Universität Darmstadt, 2003--2004.
\newblock (Introducció a la teoria de categories orientada a la lògica categòrica; referència per als capítols de categories cartesianes tancades i topos i per al pont cap al volum de lògica. Les remissions es refereixen a la versió en línia de les notes indicada a la referència.)

\bibitem{borceux}
F.~Borceux,
\emph{Handbook of Categorical Algebra},
3 volums, Encyclopedia of Mathematics and its Applications, Cambridge University Press, 1994.
\newblock (Referència exhaustiva per a límits, adjuncions i exactitud.)

\bibitem{maclanemoerdijk}
S.~Mac~Lane i I.~Moerdijk,
\emph{Sheaves in Geometry and Logic: A First Introduction to Topos Theory},
Universitext, Springer, 1992.
\newblock (Pont cap a la teoria de topos i la lògica geomètrica. Es cita l'edició estàndard impresa de 1992.)

\bibitem{lambekscott}
J.~Lambek i P.~J.~Scott,
\emph{Introduction to Higher-Order Categorical Logic},
Cambridge Studies in Advanced Mathematics~7, Cambridge University Press, 1986.
\newblock (Correspondència entre lògica d'ordre superior, càlcul lambda i categories cartesianes tancades.)

\bibitem{makkaireyes}
M.~Makkai i G.~E.~Reyes,
\emph{First Order Categorical Logic},
Lecture Notes in Mathematics~611, Springer, 1977.
\newblock (Tractament clàssic de la lògica de primer ordre en categories regulars, coherents i de Heyting, i dels topos classificadors; referència per al volum de lògica categòrica.)

\bibitem{pitts}
A.~M.~Pitts,
\emph{Categorical Logic},
dins S.~Abramsky, D.~M.~Gabbay i T.~S.~E.~Maibaum (eds.), \emph{Handbook of Logic in Computer Science}, vol.~5, Oxford University Press, 2000, pàg.~39--128.
\newblock (Introducció sistemàtica a la semàntica categòrica de la lògica, des de les teories equacionals fins a la lògica d'ordre superior; referència per al volum de lògica categòrica.)

\bibitem{jacobs}
B.~Jacobs,
\emph{Categorical Logic and Type Theory},
Studies in Logic and the Foundations of Mathematics~141, Elsevier, 1999.
\newblock (Fibracions, hiperdoctrines i teoria de tipus des del punt de vista categòric.)

\bibitem{johnstone}
P.~T.~Johnstone,
\emph{Sketches of an Elephant: A Topos Theory Compendium},
Oxford Logic Guides 43--44, Oxford University Press, 2002.
\newblock (Obra avançada de consulta sobre teoria de topos.)

\bibitem{lawvere-ccfm}
F.~W.~Lawvere,
\emph{The Category of Categories as a Foundation for Mathematics},
en \emph{Proceedings of the Conference on Categorical Algebra (La Jolla, 1965)}, Springer, 1966, pàg.~1--20.
\newblock (Teoria elemental de categories en primer ordre; precedent de la presentació relacional de l'apèndix~\ref{cap:axiomatica-relacional}.)

\bibitem{lawvere-etcs}
F.~W.~Lawvere,
\emph{An Elementary Theory of the Category of Sets},
Proceedings of the National Academy of Sciences of the USA~\textbf{52} (1964), 1506--1511. Versió llarga, amb comentaris de C.~McLarty i de l'autor, a \emph{Reprints in Theory and Applications of Categories}~11 (2005): \url{http://www.tac.mta.ca/tac/reprints/articles/11/tr11abs.html}.
\newblock (Axiomàtica de la teoria de conjunts formulada en termes de la categoria de conjunts: la teoria de categories com a fonament.)

\bibitem{leinster-rethinking}
T.~Leinster,
\emph{Rethinking Set Theory},
American Mathematical Monthly~\textbf{121} (2014), núm.~5, 403--415. Versió en accés obert: \url{https://arxiv.org/abs/1212.6543}.
\newblock (Presentació informal i breu de la teoria de conjunts estructural de Lawvere.)

\bibitem{shulman-mida}
M.~A.~Shulman,
\emph{Set Theory for Category Theory},
prepublicació, 2008. \url{https://arxiv.org/abs/0810.1279}.
\newblock (Comparació de les diferents maneres de tractar la mida en teoria de categories: classes, universos i altres alternatives.)

\bibitem{nlab-classificador}
nLab,
\emph{classifying topos},
\url{https://ncatlab.org/nlab/show/classifying+topos} (consultat el 29 de setembre de 2026).
\newblock (Wiki de referència en teoria de categories; entrada sobre el topos classificador d'una teoria geomètrica.)

\bibitem{nlab-locale}
nLab,
\emph{locale},
\url{https://ncatlab.org/nlab/show/locale} (consultat el 29 de setembre de 2026).
\newblock (Hi apareix l'exemple del locale de les aplicacions exhaustives $\mathbb N\to\mathbb R$, no trivial i sense punts.)

\bibitem{blechschmidt}
I.~Blechschmidt,
\emph{Generalized Spaces for Constructive Algebra},
prepublicació, 2020. \url{https://arxiv.org/abs/2012.13850}.
\newblock (Presentació accessible dels espais generalitzats sense punts, amb l'exemple de les aplicacions exhaustives $\mathbb N\to\mathbb R$.)

\bibitem{caramello}
O.~Caramello,
\emph{Theories, Sites, Toposes: Relating and Studying Mathematical Theories through Topos-Theoretic \enquote{Bridges}},
Oxford University Press, 2018.
\newblock (Referència sobre els toposos classificadors des del punt de vista de la lògica. Primera edició impresa de 2018; la versió electrònica es va publicar el desembre de 2017.)

\bibitem{hott}
The Univalent Foundations Program,
\emph{Homotopy Type Theory: Univalent Foundations of Mathematics},
Institute for Advanced Study, 2013. Accés obert: \url{https://homotopytypetheory.org/book/}.
\newblock (Introducció a la teoria homotòpica de tipus, llegible des de la lògica.)

\bibitem{scott-identity}
D.~S.~Scott,
\emph{Identity and Existence in Intuitionistic Logic},
en M.~Fourman, C.~Mulvey i D.~Scott (eds.), \emph{Applications of Sheaves}, Lecture Notes in Mathematics~753, Springer, 1979, pàg.~660--696.
\newblock (Operacions parcials, existència i identitat; rerefons de les axiomàtiques \enquote{només amb fletxes}.)

\bibitem{freydscedrov}
P.~J.~Freyd i A.~Scedrov,
\emph{Categories, Allegories},
North-Holland Mathematical Library~39, North-Holland, 1990.
\newblock (Presentació sistemàtica sense objectes primitius.)

\bibitem{BS2020}
C.~Benzmüller i D.~S.~Scott,
\emph{Automating Free Logic in HOL, with an Experimental Application in Category Theory},
Journal of Automated Reasoning~64, 2020, pàg.~53--72.
\newblock (Formalització en Isabelle/HOL de diverses axiomàtiques de categories amb lògica lliure.)

\bibitem{BSAFP}
C.~Benzmüller i D.~S.~Scott,
\emph{Axiom Systems for Category Theory in Free Logic},
Archive of Formal Proofs, 2018. \url{https://isa-afp.org/entries/AxiomaticCategoryTheory.html}.
\newblock (Comparació i verificació d'equivalències entre sistemes de Mac~Lane, Scott i Freyd--Scedrov.)

\bibitem{diec}
Institut d'Estudis Catalans,
\emph{Diccionari de la llengua catalana},
2a edició, Barcelona. \url{https://dlc.iec.cat}.
\newblock (Font normativa de la terminologia general catalana.)

\bibitem{termcat}
Universitat Politècnica de Catalunya; Enciclopèdia Catalana,
\emph{Diccionari de matemàtiques i estadística},
1a ed., Barcelona, 2002; 2a ed. en línia, Institut d'Estudis Catalans, 2021. \url{https://cit.iec.cat/DME/}.
\newblock (Font de la terminologia matemàtica catalana: functor, morfisme, preordre, etc.)

\end{thebibliography}

\chapter*{Glossari català--anglès}
\addcontentsline{toc}{chapter}{Glossari català--anglès}
\markboth{Glossari català--anglès}{Glossari català--anglès}

Aquest glossari recull sobretot termes conflictius, no transparents o especialment útils per llegir bibliografia internacional, i no pretén ser un vocabulari exhaustiu de teoria de categories. S'hi dona la forma anglesa i, quan escau, la notació emprada.

\medskip
\noindent\textbf{Criteri d'anglicismes.} Diversos termes s'han manllevat de l'anglès per manca d'una forma catalana establerta o per fidelitat a la bibliografia (\emph{pullback}, \emph{pushout}). En la seva primera aparició es dona, quan existeix, el terme català al costat de l'anglès i de la notació; després se n'usa una sola forma. La terminologia general segueix el \emph{Diccionari de la llengua catalana} de l'Institut d'Estudis Catalans \cite{diec}, i la matemàtica, el \emph{Diccionari de matemàtiques i estadística} del TERMCAT \cite{termcat}.

\medskip

\renewcommand{\arraystretch}{1.25}
\begin{xltabular}{\linewidth}{@{}>{\raggedright\arraybackslash}X >{\raggedright\arraybackslash}X l@{}}
\hline
\textbf{Català} & \textbf{Anglès} & \textbf{Notació}\\
\hline
\endfirsthead
\hline
\textbf{Català} & \textbf{Anglès} & \textbf{Notació}\\
\hline
\endhead
\hline
\endfoot
\endlastfoot
categoria oposada & opposite category & $\cat{C}\op$\\
categoria prima & thin category & \\
categoria sobre un objecte & slice category & $\cat{C}/A$\\
categoria sota un objecte & coslice category & $A/\cat{C}$\\
homconjunt & hom-set & $\Hom(A,B)$\\
functor representable & representable functor & p.\,ex.\ $\Hom(-,A)$\\
immersió de Yoneda & Yoneda embedding & \\
bihomfunctor (o bifunctor Hom) & hom functor, Hom bifunctor & $\Hom(-,-)$\\
component (d'una transformació natural; femení) & component & $\eta_A$\\
antiimatge & preimage, inverse image & $f^{-1}(B)$\\
exhaustiu / essencialment exhaustiu & surjective / essentially surjective & \\
\multicolumn{3}{@{}p{\linewidth}@{}}{\footnotesize En aquest llibre, \emph{exhaustiu} és el terme sistemàtic per a \emph{surjective}; \emph{essencialment exhaustiu} tradueix \emph{essentially surjective}.}\\[0.35em]
transformació natural & natural transformation & $\eta\colon F\Rightarrow G$\\
equivalència de categories & equivalence of categories & $\cat{C}\simeq\cat{D}$\\
monomorfisme / epimorfisme & monomorphism / epimorphism & mono / epi\\
objecte inicial / terminal & initial / terminal object & $0$ / $1$ (quan s'adopta aquesta notació)\\
functor ple / fidel & full / faithful functor & \\
imatge essencial & essential image & \\
categoria d'elements & category of elements & $\textstyle\int P$\\
categoria essencialment petita & essentially small category & \\
quasiinvers (d'una equivalència) & quasi-inverse / pseudo-inverse & \\
prefeix & presheaf & $\cat{Set}^{\cat{C}\op}$\\
producte fibrat (pullback) & pullback & \\
suma amalgamada (pushout) & pushout & \\
con & cone & \\
cocon & cocone & \\
límit / colímit & limit / colimit & $\varprojlim$ / $\varinjlim$\\
unitat / counitat & unit / counit & $\eta$ / $\varepsilon$\\
composició horitzontal & horizontal composition & $\theta\ast\eta$\\
composició amb un functor per l'esquerra & (left) whiskering & $K\eta$\\
composició amb un functor per la dreta & (right) whiskering & $\eta H$\\
epimorfisme regular & regular epimorphism & \\
parella nucli & kernel pair & \\
factorització imatge & image factorization & \\
reindexació & reindexing / base change & $f^{*}$\\
subobjecte & subobject & $\Sub(A)$\\
classificador de subobjectes & subobject classifier & $\Omega$\\
sedàs & sieve & \\
semireticle implicatiu & implicative meet-semilattice & \\
categoria cartesiana tancada & cartesian closed category & \\
categoria regular & regular category & \\
\hline
\end{xltabular}
\renewcommand{\arraystretch}{1.0}

\chapter*{Taula de notació}
\addcontentsline{toc}{chapter}{Taula de notació}
\markboth{Taula de notació}{Taula de notació}

Recollim la notació principal del llibre, amb el seu significat i el lloc de la primera aparició o definició. Les referències remeten al capítol o a la secció corresponent.

\medskip
\renewcommand{\arraystretch}{1.25}
\begin{xltabular}{\linewidth}{@{}l >{\raggedright\arraybackslash}X l@{}}
\hline
\textbf{Símbol} & \textbf{Significat} & \textbf{Primera aparició}\\
\hline
\endfirsthead
\hline
\textbf{Símbol} & \textbf{Significat} & \textbf{Primera aparició}\\
\hline
\endhead
\hline
\endfoot
\endlastfoot
$\cat{C},\cat{D},\dots$ & categories & cap.~\ref{cap:categories}\\
$\Hom(A,B)$, $\cat{C}(A,B)$ & homconjunt de morfismes $A\to B$ & cap.~\ref{cap:categories}\\
$\id_A$ & morfisme identitat de $A$ & cap.~\ref{cap:categories}\\
$g\comp f$ & composició ($f$ seguit de $g$) & cap.~\ref{cap:categories}\\
$\cat{C}\op$ & categoria oposada & cap.~\ref{cap:categories}\\
$f\op$ & morfisme dual de $f$ a $\cat{C}\op$ & \S\ref{sec:dualitat}\\
$\cat{Set},\cat{Grp},\cat{Top},\cat{Graph}$ & categories concretes habituals & cap.~\ref{cap:categories}\\
$\cat{Cat}$ & categoria de les categories petites & cap.~\ref{cap:functors}\\
$\cat{C}/A$, $A/\cat{C}$ & categoria sobre, sota un objecte & cap.~\ref{cap:functors}\\
$(S\downarrow T)$ & categoria coma & \S\refalt{sec:coma}{del cap.~4}\\
$F,G\colon\cat{C}\to\cat{D}$ & functors & cap.~\ref{cap:functors}\\
$\Hom(A,-)$, $\Hom(-,B)$ & homfunctors covariant i contravariant & def.~\ref{def:homfunctors}\\
$\Hom(A,f)=f_*$ & postcomposició amb $f$ (acció de $\Hom(A,-)$ sobre $f$) & not.~\ref{not:hom-morfismes}\\
$\Hom(f,B)=f^*$ & precomposició amb $f$ (acció de $\Hom(-,B)$ sobre $f$). \emph{Atenció:} el mateix símbol designa la reindexació (vegeu més avall) & not.~\ref{not:hom-morfismes}\\
$\widehat{\cat{C}}=\cat{Set}^{\cat{C}\op}$ & categoria de prefeixos sobre $\cat{C}$ & ex.~\ref{ex:categoria-prefeixos}\\
$\Hom(-,-)$, $\Hom(f,g)$ & bihomfunctor i la seva acció $h\mapsto g\comp h\comp f$ & def.~\ref{def:bihomfunctor}\\
$\alpha\colon F\Rightarrow G$ & transformació natural & cap.~\ref{cap:transformacions}\\
$\cat{D}^{\cat{C}}$, $[\cat{C},\cat{D}]$ & categoria de functors & cap.~\ref{cap:transformacions}\\
$\Hom(-,A)$ & exemple de functor representable (contravariant) & cap.~\refcap{cap:yoneda}{4}\\
$\yon$, $\Hom(-,f)$ & immersió de Yoneda i la seva acció $\yon(f)$ sobre $f$ & def.~\refalt{def:immersio-yoneda}{del cap.~4}\\
$\Delta_A$ & functor constant amb valor $A$ & ex.~\ref{ex:functor-constant}\\
$\Delta$ & functor diagonal $\cat{C}\to\cat{C}^{\cat{I}}$ & obs.~\refalt{obs:cons-coma}{del cap.~5}\\
$\delta_A$ & morfisme diagonal $\langle\id_A,\id_A\rangle\colon A\to A\times A$ & cap.~\refcap{cap:cartesianes}{7}\\
$\varprojlim D$, $\colimop D$ & límit, colímit d'un diagrama & cap.~\refcap{cap:limits}{5}\\
$A\times B$, $A+B$ & producte, coproducte & cap.~\refcap{cap:limits}{5}\\
$F\dashv G$ & $F$ és adjunt per l'esquerra de $G$ & cap.~\refcap{cap:adjuncions}{6}\\
$\eta$, $\varepsilon$ & unitat, counitat d'una adjunció & cap.~\refcap{cap:adjuncions}{6}\\
$C^{B}$, $\lambda f$ & exponencial, currificació & cap.~\refcap{cap:cartesianes}{7}\\
$\Sub(A)$ & ordre parcial de subobjectes de $A$ & cap.~\refcap{cap:subobjectes}{8}\\
$f^{*}$ & reindexació per pullback (substitució), a $\cat{C}/B$ i a $\Sub(B)$; a $\cat{Set}$, l'antiimatge $f^{-1}$. És el mateix símbol que la precomposició; el context (homconjunts o subobjectes i fibres) decideix. A $\cat{Set}$ totes dues lectures coincideixen a través de $\Sub(B)\cong\Hom(B,\{0,1\})$ & obs.~\refalt{obs:reindexacio}{del cap.~5}; cap.~\refcap{cap:subobjectes}{8}\\
$\Imatge(f)$ & imatge d'un morfisme & cap.~\refcap{cap:subobjectes}{8}\\
$\exists_f$, $\forall_f$ & quantificadors com a adjunts de $f^{*}$ & cap.~\refcap{cap:subobjectes}{8}\\
$\Omega$, $\mathsf{true}$ & classificador de subobjectes & cap.~\refcap{cap:topos}{9}\\
\hline
\end{xltabular}
\renewcommand{\arraystretch}{1.0}

\printindex

\end{document}
